\documentclass[11pt]{report}
\usepackage{amsmath,amssymb,amsthm,mathtools}
\usepackage{enumitem}
\usepackage{graphicx}
\usepackage{longtable,array,booktabs}
\usepackage[margin=2.6cm]{geometry}
\usepackage[hidelinks]{hyperref}
\usepackage[nameinlink,capitalise,noabbrev]{cleveref}

\setlist{itemsep=1pt,topsep=3pt}
\numberwithin{equation}{chapter}

\theoremstyle{plain}
\newtheorem{theorem}{Theorem}[chapter]
\newtheorem{proposition}[theorem]{Proposition}
\newtheorem{lemma}[theorem]{Lemma}
\newtheorem{corollary}[theorem]{Corollary}
\theoremstyle{definition}
\newtheorem{definition}[theorem]{Definition}
\theoremstyle{remark}
\newtheorem{remark}[theorem]{Remark}

\newcommand{\cE}{\mathcal E}
\newcommand{\cB}{\mathcal B}
\newcommand{\cP}{\mathcal P}
\newcommand{\cQ}{\mathcal Q}
\newcommand{\cT}{\mathcal T}
\newcommand{\one}{\mathbf 1}
\newcommand{\mvec}{\mathbf m}
\newcommand{\bmu}{\boldsymbol\mu}
\newcommand{\btau}{\boldsymbol\tau}
\newcommand{\ip}[2]{\langle #1,#2\rangle}
\newcommand{\oE}{\mathbb E}
\newcommand{\sG}{\mathsf G}
\newcommand{\sZ}{\mathsf Z}
\newcommand{\sB}{\mathsf B}
\newcommand{\sT}{\mathsf T}
\newcommand{\sQ}{\mathsf Q}
\newcommand{\Gst}[4]{\Gamma_{\substack{#1\\#2}}\!\begin{bmatrix}#3\\#4\end{bmatrix}}
\newcommand{\IZ}{\mathbb I_{\mathbb Z}}
\newcommand{\IR}{\mathbb I_{\mathbb R}}
\newcommand{\Ilat}[1]{\mathbb I_{#1}}
% One-scale Trailmoments at the unit exponent and the coefficient function x:
% every macro below prints the full signature of Definition def:jointtrail.
\newcommand{\Tx}[3]{(\sT_{#1;(1)}[x]\,#2)((#3,\infty))}
\newcommand{\Th}[4]{(\sT_{#1;(1)}[x]\,#2)((#3,#4])}
\newcommand{\Ttr}[4]{(\sT_{#1;I_{#2}}[y]\,#3)(\{#4<y_1<\cdots<y_{#2}\})}
\DeclareMathOperator{\re}{Re}
\DeclareMathOperator{\im}{Im}
\DeclareMathOperator{\Reg}{Reg}
\DeclareMathOperator*{\CT}{CT}
\DeclareMathOperator{\cone}{cone}
\DeclareMathOperator{\interior}{int}
\DeclareMathOperator{\Dir}{Dir}
\DeclareMathOperator{\pat}{pat}
\newcommand{\reflint}[2]{\reflectbox{$#1\int$}}
\newcommand{\pseudoint}{\mathop{\mathpalette\reflint\relax}\nolimits}
\emergencystretch=3em
\allowdisplaybreaks
% Environments and macros of the last part, whose statements are cited by their numbers.
\newtheoremstyle{handbook}{6pt}{6pt}{\itshape}{}{\bfseries}{.}{.5em}{\thmname{#1}\thmnote{ #3}}
\newtheoremstyle{handbookdef}{6pt}{6pt}{\normalfont}{}{\bfseries}{.}{.5em}{\thmname{#1}\thmnote{ #3}}
\theoremstyle{handbook}
\newtheorem*{hproposition}{Proposition}
\newtheorem*{htheorem}{Theorem}
\newtheorem*{hcorollary}{Corollary}
\newtheorem*{hlemma}{Lemma}
\theoremstyle{handbookdef}
\newtheorem*{hdefinition}{Definition}
\theoremstyle{plain}
\newtheorem{slemma}{Lemma}
\renewcommand{\theslemma}{S\arabic{slemma}}
\theoremstyle{remark}
\newtheorem*{sremark}{Remark}
\newcommand{\bsig}{\boldsymbol\sigma}
\newcommand{\bdel}{\boldsymbol\delta}
\newcommand{\hb}[1]{\ref{#1}}
\DeclareMathOperator{\wdt}{wd}
\DeclareMathOperator{\load}{load}
\title{The Foreman and the serif family\\[6pt]\large Handbook of the general definitions, constructions and theorems}
\date{29 October 2026}
\hypersetup{pdftitle={The Foreman and the serif family: handbook of the general definitions, constructions and theorems}}
\author{}
% v2: exercises and their solutions (Part XI).
\theoremstyle{definition}
\newtheorem{exercise}{Exercise}[chapter]
\newenvironment{solution}{\begin{proof}[Solution]}{\end{proof}}
\theoremstyle{plain}
\begin{document}
\maketitle

\chapter*{How this handbook is written}
\addcontentsline{toc}{chapter}{How this handbook is written}

This handbook states the general theory of the Foreman and the serif family: the members of the serif family and their relations; the tables and families they evaluate; the Foreman's classification of a family and his tree; the generalized integrals the Foreman sets; the constructions and compositions of families; the motive of a family and the arithmetic of its values; rays of data and the laws their values obey; the relations among values; and the family read as a model of data. Every definition, construction, theorem and procedure of the program is stated here, once, at the place where it is first needed.

\emph{The order of the book.} The serif family comes first (Part~I), because every value the theory produces is closed in it. Then the path of a computation: a list of marks and the tables and families built on it (Part~II); the classification of a family by its realizations, its reading, its bunches and the Foreman tree (Part~III); the integrals the Foreman procedure sets, and their evaluation along the tree (Part~IV); the constructions that build families from cells, holds and bunches, and price them (Part~V); the motive of a family and the coaction along its tree, positivity and the gap (Part~VI); rays of data, their modules, their two ends and the reading of their values (Part~VII); relations among values and the content they live in (Part~VIII); and the family as a model of data (Part~IX). No statement and no proof uses a statement of a later chapter, and a later chapter is mentioned only as a pointer; inside a chapter a statement may point ahead to a result proved below it, and then says so.

\emph{Proofs.} A statement that is more general than its counterpart in the reference \emph{The Foreman and the serif family} (16 October 2026), or that has no counterpart there, is proved in Part~\ref{part:proofs}, one section per chapter in the order of the chapters; the other statements are proved in the reference. In the source file every statement is preceded by a comment that says which: either the label and line of the statement in the reference, or that it is proved in the last part.

\emph{Exercises.} Computations on particular lists, tables and rays are the exercises of Part~\ref{part:exercises}. Each is a guided example with its solution worked in full: every step names the statement or procedure it uses, and every quantity the next step needs is displayed in exact arithmetic.

\emph{Names.} The text uses the program's own names. The classical names of its objects, and the classical names of results that the program reaches by its own route, are collected in the Aliases at the end.

\emph{One governing definition per object.} Each object of the theory is defined once, with all of its slots. Every later appearance of the object is that definition with particular data in its slots, written through the full signature of its governing symbol: a specialization displays the data it chooses, and it never receives a symbol of its own. Thus a Breta on one scale is a Breta with one bunch, a Trailmoment on a line is a Trailmoment whose cut region is a half-line, and a value of interest at particular exponents is the interest with those exponents in its slots.

\emph{Abbreviations are local.} Inside one statement an abbreviation may be declared, once and in words, for data that the statement holds fixed; the declaration displays the full symbol it abbreviates, and it lapses at the end of the statement.

\emph{Relations are stated at their height.} A relation holds on the values of the slots on which its proof establishes it (Definition~\ref{def:height}). A hypothesis that restricts a slot is written into the statement that needs it, and nothing outside that statement restricts the object.

\emph{Final answers} are written with $\exp$, $\log$, the stacked Gamma and the serif family (Definition~\ref{def:conv}(vi)).

\tableofcontents

\part{The serif family}

\chapter{Conventions, the stacked Gamma and the finite part}\label{ch:conv}

% Reference 20: def:conv; source line 71
\begin{definition}[Conventions]\label{def:conv}
\begin{enumerate}[label=(\roman*)]
\item $\cE$ and $B$ are finite nonempty sets and $n=|\cE|$. For $x\in\mathbb R^{\cE}_{>0}$ and $v\in\mathbb C^{\cE}$, $x^{v}=\prod_ex_e^{v_e}$ with $x_e^{v_e}=\exp(v_e\log x_e)$. For a vector $v$, $|v|=\sum_ev_e$. $\one$ is the all-ones vector, $\one_S$ the indicator of $S$, $x^{S}=x^{\one_S}$, $v_S=\sum_{e\in S}v_e$, $v|_S$ the restriction of $v$ to $S$, and $x^{-1}=(x_e^{-1})_e$. For a set $X$ and elements $a,b$, $X-a=X\setminus\{a\}$ and $X+b=X\cup\{b\}$. $\mathbb N=\{0,1,2,\dots\}$.
\item For $\alpha\in\mathbb R^{\cE}_{>0}$ the \emph{logarithmic coordinates} are $y=\log\alpha$ and the \emph{signed coordinates} are $z=\tfrac12y$, coordinatewise, so $\alpha=e^{y}=e^{2z}$.
\item $\Delta_B=\{u\in\mathbb R^B_{>0}:|u|=1\}$ is the \emph{shape simplex} of $B$, and $d\omega_B(u)$ is Lebesgue measure in the coordinates $(u_e)_{e\neq e_0}$ for a fixed $e_0\in B$; a change of $e_0$ is affine of determinant $\pm1$. For $|B|=1$, $\omega_B$ is the Dirac measure at the point of $\Delta_B$.
\item An integral \emph{converges} when it converges absolutely.
\item $(a)_j=\Gamma(a+j)/\Gamma(a)$ for $j\in\mathbb N$; $d_n=\operatorname{lcm}(1,\dots,n)$ with $d_0=1$; $e_k$ is the elementary symmetric polynomial of degree $k$; for a coordinate $x$ of the unit interval, $\bar x=1-x$; $[X]\,Y$ is the coefficient of $X$ in $Y$.
\item \emph{Final answers} are written with $\exp$, $\log$, the stacked Gamma and the serif family: the Gramma $\sG$ and the Zamma $\sZ$ (Chapter~\ref{ch:gramma}), the Breta $\sB$ (Chapter~\ref{ch:breta}) and the Trailmoment $\sT$ (Chapter~\ref{ch:trail}). A value of interest is left in another form only with a proof that it has no such expression. A \emph{word} is an expression built from rationals, the data and the variables by sums, products, powers with rational exponents or exponents among the data, these functions, integrals over simplices, half-lines and fibres, and the iterate $\mathcal Q^M$ of an integral operator whose kernel is a word. A finite sequence $(w_1,\dots,w_k)$ is a \emph{string} of length $k$. An \emph{insertion} is a function multiplied into an integrand.
\item For masses $\mvec$, $d\alpha_{\mvec}:=\alpha^{\mvec-\one}d\alpha$ is a measure on the orthant, and a form $f\,d\alpha_{\mvec}$ unchanged by $\alpha\mapsto\lambda\alpha$, $f$ homogeneous of degree $-m_{\cE}$, is integrated over the rays of the orthant, that is over any section (Chapter~\ref{ch:expectations}), as the form $W\,d\alpha_{\mvec}$ of a family word $W$ (Chapter~\ref{ch:families}) is on its balance wall; a mark or a set of marks as a subscript, as in $d\alpha_B$, denotes coordinates. The law of a random variable, or of a measure under a map, is its image measure; an identity between members of the serif family is a relation (Chapter~\ref{ch:trail}). An expectation names its random variable and that variable's law, $\oE_{U\sim P}[g(U)]=\int g(u)\,P(du)$; inside a statement that has named them once, the subscript may be dropped for the rest of that statement. The \emph{Gamma law} of index $\nu>0$ is $\operatorname{Gamma}_\nu(dx)=x^{\nu-1}e^{-x}dx/\Gamma(\nu)$ on $(0,\infty)$, and for $\mvec\in\mathbb R^B_{>0}$, $\operatorname{Gamma}_{\mvec}=\bigotimes_{e\in B}\operatorname{Gamma}_{m_e}$ is the law on $\mathbb R^B_{>0}$ of independent Gamma coordinates of indices $m_e$. The generalized integrals of Chapter~\ref{ch:expectations} are unnormalized, and their expectations divide them by the period of the centre: on a connected family every expectation of a form is taken under the law of its centre, $\mathbb P_{\mathrm c}$, of the departure that its evaluation sets.
\item The \emph{root rate} of terms $a_\eta$ indexed by sizes $\eta\to\infty$, as the terms of a lattice by the real parts of their exponents or the terms of a closing by its steps, is $\limsup|a_\eta|^{1/\eta}$. When the number of terms of size at most $N$ grows at most polynomially in $N$, terms of root rate below one are absolutely summable, and terms of root rate above one do not tend to zero.
\end{enumerate}
\end{definition}

\section{The stacked Gamma}

% Reference 20: def:stacked; source line 83
\begin{definition}[The stacked Gamma]\label{def:stacked}
For $\alpha\in\mathbb C^p$ and $\beta\in\mathbb C^q$,
\[
\Gst{p}{q}{\alpha_1,\ \dots,\ \alpha_p}{\beta_1,\ \dots,\ \beta_q}:=\frac{\Gamma(\alpha_1)\cdots\Gamma(\alpha_p)}{\Gamma(\beta_1)\cdots\Gamma(\beta_q)} ,
\]
wherever no $\alpha_i$ is a pole of $\Gamma$; it is a meromorphic function of its entries. The stacked Gamma is \emph{balanced} when $\sum_i\alpha_i=\sum_j\beta_j$.
\end{definition}

The subscripts of the bracket record the numbers of its upper and lower entries, so a quotient of Gammas is read off the symbol with no translation. The stacked Gamma $\Gst{|B|}{1}{\mvec}{m_B}=\prod_{e\in B}\Gamma(m_e)/\Gamma(m_B)$ of masses $\mvec\in\mathbb C^B$ over their total $m_B$ is, for $\re m_e>0$, the integral of $u^{\mvec-\one}$ over the shape simplex $\Delta_B$ (Theorem~\ref{thm:dirichlet}).

% Reference 20: prop:stacked; source line 91
\begin{proposition}[Relations of the stacked Gamma]\label{prop:stacked}
\begin{enumerate}[label=(\roman*)]
\item \emph{Concatenation.} $\Gst pq{\boldsymbol\alpha}{\boldsymbol\beta}\Gst{p'}{q'}{\boldsymbol\alpha'}{\boldsymbol\beta'}=\Gst{p+p'}{q+q'}{\boldsymbol\alpha,\ \boldsymbol\alpha'}{\boldsymbol\beta,\ \boldsymbol\beta'}$.
\item \emph{Cancellation.} $\Gst pq{\boldsymbol\alpha,\ c}{\boldsymbol\beta,\ c}=\Gst{p-1}{q-1}{\boldsymbol\alpha}{\boldsymbol\beta}$.
\item \emph{Shift.} $\Gst pq{\alpha_1+1,\ \boldsymbol\alpha}{\boldsymbol\beta}=\alpha_1\Gst pq{\alpha_1,\ \boldsymbol\alpha}{\boldsymbol\beta}$ and $\Gst pq{\boldsymbol\alpha}{\beta_1+1,\ \boldsymbol\beta}=\beta_1^{-1}\Gst pq{\boldsymbol\alpha}{\beta_1,\ \boldsymbol\beta}$.
\item \emph{Residue.} For $j\in\mathbb N$, $\displaystyle\operatorname*{Res}_{\alpha_1=-j}\Gst pq{\alpha_1,\ \boldsymbol\alpha}{\boldsymbol\beta}=\frac{(-1)^j}{j!}\Gst{p-1}{q}{\boldsymbol\alpha}{\boldsymbol\beta}\Big|_{\alpha_1=-j}$.
\item \emph{Reflection.} $\Gst20{a,\ 1-a}{}=\pi/\sin\pi a$.
\item \emph{Laurent data at a pole.} For $j\in\mathbb N$,
\[
\Gamma(-j+\varepsilon)=\frac{(-1)^j}{j!}\Big(\frac1\varepsilon+(\log\Gamma)'(j+1)+O(\varepsilon)\Big),\qquad (\log\Gamma)'(\nu+1)=(\log\Gamma)'(\nu)+\frac1\nu .
\]
\end{enumerate}
\end{proposition}

\begin{remark}[Balance is a property of the written stacked Gamma]\label{rem:sl-balance}
Balance is not a property of the value of a stacked Gamma: the shift (iii) and the reflection (v) change the sums of the entries, multiplying the value by an algebraic number when the entries are rational. For rational $a\notin\mathbb Z$, (v) at $a$ and at $\frac12$, with $\Gamma(1)=1$, gives
\[
\Gst20{a,\ 1-a}{}=\frac1{\sin\pi a}\,\Gst21{\frac12,\ \frac12}{1},
\]
the left side with the sums $1$ and $0$, the right side balanced, and $\sin\pi a$ algebraic.
\end{remark}

\section{The finite part}

% Reference 20: def:Reg; source line 113 (real exponents).
% Proved in the last part: def:Reg.
\begin{definition}[The finite part]\label{def:Reg}
A function $h$ on $(0,\infty)$ has a \emph{log-power expansion} at $0$ if for every $N$ there are finitely many coefficients $c_{k,j}$, with exponents $k\in\mathbb C$ of real part below $N$ and $j\in\mathbb N$, such that $h(x)-\sum c_{k,j}x^k\log^jx=O(x^N)$ as $x\to0$; at $\infty$ the same holds with $x^{-k}$ and $O(x^{-N})$. The coefficients are determined by $h$. Then $\Reg_{x\to0}h$ and $\Reg_{x\to\infty}h$ are the coefficients of $x^0\log^0x$ in these expansions, zero if absent. If $f$ has a primitive $h$ with log-power expansions at both ends,
\[
\Reg\!\int_0^\infty f\,dx:=\Reg_{x\to\infty}h-\Reg_{x\to0}h ,
\]
which does not depend on the primitive; on $(0,b]$, $\Reg\!\int_0^bf:=h(b)-\Reg_{x\to0}h$. A sum is regularized the same way, $\Reg_{N\to\infty}$ of its partial sums. The coordinate is part of the data of a finite part, and for iterated finite parts so is the order.
\end{definition}

% Reference 20: prop:Regconv; source line 121 (real exponents).
% Proved in the last part: prop:Regconv.
\begin{proposition}[A convergent integral is its own finite part]\label{prop:Regconv}
If $\int_0^\infty f\,dx$ converges and $f$ has a primitive with log-power expansions at both ends, then $\Reg\!\int_0^\infty f\,dx=\int_0^\infty f\,dx$.
\end{proposition}

% Reference 20: thm:RegCT; source line 128 (real exponents).
% Proved in the last part: thm:RegCT.
\begin{theorem}[A finite part is a constant term]\label{thm:RegCT}
If $f$ has log-power expansions at $0$ and at $\infty$, then
\[
\Reg\!\int_0^\infty f\,dx=\CT_{\varepsilon=0}\int_0^\infty x^{\varepsilon}f(x)\,dx ,
\]
the right side continued meromorphically in $\varepsilon$ from $\int_0^1$ and from $\int_1^\infty$ separately. Termwise, $x^b\log^jx$ in the expansion at $0$ contributes $\int_0^1x^{b+\varepsilon}\log^jx\,dx=(-1)^jj!/(b+1+\varepsilon)^{j+1}$, whose constant term is $(-1)^jj!/(b+1)^{j+1}$ for $b\neq-1$ and $0$ for $b=-1$; the end at $\infty$ is the same after $x\mapsto1/x$.
\end{theorem}

% Reference 20: cor:RegMellin; source line 143 ($f$ smooth on $[0,1]$).
% Proved in the last part: cor:RegMellin.
\begin{corollary}[Finite parts at a pole of a Mellin integral]\label{cor:RegMellin}
Let $f$ have a log-power expansion at $0$ and be integrable on $[\delta,1]$ for every $\delta>0$, and let $I(a)=\int_0^1x^{a-1}f(x)\,dx$, continued meromorphically. Then $\Reg\!\int_0^1x^{a-1}f(x)\,dx=\CT_{a'=a}I(a')$ for every $a\in\mathbb C$: it is $I(a)$ at every $a$ that is not a pole, and the poles lie at the points $a=-k$ for the exponents $k$ of the expansion, a term $x^k\log^jx$ giving a pole of order at most $j+1$.
\end{corollary}

% Reference 20: lem:Regscale; source line 150 (the instance $\int_0^\infty dx/(1+cx)$).
% Proved in the last part: lem:Regscale.
\begin{lemma}[A finite part depends on the coordinate]\label{lem:Regscale}
Let $h$ have a log-power expansion at $\infty$, and let $a_j$ be its coefficient of $x^0\log^jx$, $j\in\mathbb N$. For $c,p>0$, $h(cx^p)$ has a log-power expansion at $\infty$, and
\[
\Reg_{x\to\infty}h(cx^p)=\sum_{j\ge0}a_j\,(\log c)^j ;
\]
the same holds at $0$. Consequently, if $f$ has log-power expansions at $0$ and at $\infty$, with coefficients $f^0_j$ and $f^\infty_j$ of $x^{-1}\log^jx$, then for $c,p>0$
\[
\Reg\!\int_0^\infty f(cx^p)\,c\,p\,x^{p-1}\,dx=\Reg\!\int_0^\infty f(x)\,dx+\sum_{j\ge0}\frac{(\log c)^{j+1}}{j+1}\big(f^\infty_j-f^0_j\big):
\]
a power of the coordinate leaves a finite part unchanged, and a dilation moves it by a polynomial in $\log c$ read from the terms $x^{-1}\log^jx$ of the integrand at its two ends.
\end{lemma}

% Reference 20: rem:Regprimitives; source line 157
A finite part is the difference of the regularized limits of one primitive at the two ends, not the integral of a measure. It is linear, it moves with the coordinate (Lemma~\ref{lem:Regscale}), and iterated in several variables it moves with the order. The generalized integrals of Chapter~\ref{ch:expectations} are built from it: a quasi-integral is an iterated finite part along a chart and an order, and the pseudointegral is the value of the native completion (Definition~\ref{def:pseudoint}).


\chapter{The Gramma and the Zamma}\label{ch:gramma}

\section{The Gramma}

% Reference 20: def:gramma; source line 164
\begin{definition}[Gramma]\label{def:gramma}
Let $\cE$ be a finite set of \emph{scales}. A \emph{configuration} on $\cE$ is a finite family $t=(t_1,\dots,t_N)$ of \emph{exponents} $t_k\in\mathbb R^{\cE}$; $t$ also denotes the $\cE\times N$ matrix with these columns. For \emph{coefficients} $\boldsymbol\xi\in(0,\infty)^N$ and an \emph{index} $\nu\in\mathbb C^{\cE}$, the \emph{Gramma} is
\[
\sG_{\nu;t}(\boldsymbol\xi):=\int_{\mathbb R^{\cE}_{>0}}x^{\nu-\one}\exp\Big(-\sum_{k=1}^N\xi_k\,x^{t_k}\Big)\,dx
\]
wherever the integral converges. Put $\cone(t)=\{\sum_kc_kt_k:\ c\in[0,\infty)^N\}$ and $\cT_t=\{\nu:\re\nu\in\interior\cone(t)\}$. In $y=\log x$,
\[
\sG_{\nu;t}(\boldsymbol\xi)=\int_{\mathbb R^{\cE}}e^{\ip\nu y}\exp\Big(-\sum_k\xi_ke^{\ip{t_k}y}\Big)\,dy ,
\]
the moment generating function, in the index, of the density $\exp\big(-\sum_k\xi_ke^{\ip{t_k}y}\big)$. Exponents and coefficients combine by concatenation, $t\oplus t'$ and $\boldsymbol\xi\oplus\boldsymbol\xi'$ (configurations on two sets of scales extended by zero to their union), and a shape $u$ dresses coefficients by $u^{t}\boldsymbol\xi:=(u^{t_k}\xi_k)_k$. The \emph{unit configuration} on $\cE$ is $I_{\cE}=(\mathbf e_e)_{e\in \cE}$, one unit column per scale, and on one scale the configuration of the single exponent $1$ is written $(1)$.
\end{definition}

The three slots of the Gramma (Definition~\ref{def:height}) are its index, its configuration and its coefficients. Every structural relation of the Gramma is a statement about how these slots move: the index moves along the exponents, the coefficients are rescaled by the torus of the scales, and the configuration passes to its facets at its walls.

% Reference 20: thm:grammadomain; source line 178
\begin{theorem}[Domain]\label{thm:grammadomain}
$\sG_{\nu;t}(\boldsymbol\xi)$ converges if and only if $\nu\in\cT_t$, and it is holomorphic on $\cT_t$. If $\cone(t)=\mathbb R^{\cE}$ it is entire; if $\cone(t)$ lies in a hyperplane it converges nowhere.
\end{theorem}

% Reference 20: prop:grammareflect; source line 201 (the reflection; the linear change and torus homogeneity stated apart).
% Proved in the last part: prop:grammareflect.
\begin{proposition}[Affine changes of the logarithmic coordinates]\label{prop:grammareflect}
For $A\in GL_{\cE}(\mathbb R)$ and $\mathbf h\in\mathbb R^{\cE}$: $\cT_{A^{\mathsf T}t}=A^{\mathsf T}\cT_t$, and for $\nu\in\cT_t$
\[
\sG_{\nu;t}(\boldsymbol\xi)=|\det A|\,e^{\ip\nu{\mathbf h}}\,\sG_{A^{\mathsf T}\nu;A^{\mathsf T}t}\big((e^{\ip{t_k}{\mathbf h}}\xi_k)_k\big).
\]
At $A=-I_{\cE}$ and $\mathbf h=0$ this is the reflection, $\cT_{-t}=-\cT_t$ and $\sG_{\nu;t}(\boldsymbol\xi)=\sG_{-\nu;-t}(\boldsymbol\xi)$.
\end{proposition}

% Reference 20: thm:grammalaws; source line 185
\begin{theorem}[Relations of the Gramma]\label{thm:grammalaws}
Let $\nu\in\cT_t$; then $\nu+t_k\in\cT_t$ for every $k$.
\begin{enumerate}[label=(\roman*)]
\item \emph{Ladder.} $\partial_{\xi_k}\sG_{\nu;t}(\boldsymbol\xi)=-\sG_{\nu+t_k;t}(\boldsymbol\xi)$; for $\mathbf w\in\mathbb N^N$, $\partial^{\mathbf w}\sG_{\nu;t}(\boldsymbol\xi)=(-1)^{|\mathbf w|}\sG_{\nu+t\mathbf w;t}(\boldsymbol\xi)$.
\item \emph{Exchange, one relation per scale.} For each $e\in \cE$: $\ \nu_e\,\sG_{\nu;t}(\boldsymbol\xi)=\sum_kt_{k,e}\,\xi_k\,\sG_{\nu+t_k;t}(\boldsymbol\xi)$.
\item \emph{Torus homogeneity.} For $\lambda\in\mathbb R^{\cE}_{>0}$: $\ \sG_{\nu;t}\big((\lambda^{t_k}\xi_k)_k\big)=\lambda^{-\nu}\sG_{\nu;t}(\boldsymbol\xi)$ (Proposition~\ref{prop:grammareflect} at $A=I_{\cE}$, $\mathbf h=\log\lambda$).
\item \emph{Toric relation.} For every $\mathbf v\in\mathbb Z^N$ with $t\mathbf v=0$, written $\mathbf v=\mathbf v^+-\mathbf v^-$ with $\mathbf v^\pm\in\mathbb N^N$ of disjoint supports,
\[
\partial^{\mathbf v^+}\sG_{\nu;t}(\boldsymbol\xi)=(-1)^{|\mathbf v^+|-|\mathbf v^-|}\,\partial^{\mathbf v^-}\sG_{\nu;t}(\boldsymbol\xi).
\]
\end{enumerate}
\end{theorem}

The ladder moves the index along one exponent; the exchange relation, one for each scale, is linear in the Grammas at the index and at its $N$ neighbours $\nu+t_k$; the toric relation records the linear dependencies among the exponents themselves. On one scale with integer exponents the exchange relation is a recurrence on the indices $\nu+\mathbb Z$ linking the indices from $\nu+\min(0,\min_kt_k)$ to $\nu+\max(0,\max_kt_k)$.

% Reference 20: thm:grammaproduct; source line 220 (two blocks).
% Proved in the last part: thm:grammaproduct.
\begin{theorem}[Product]\label{thm:grammaproduct}
If $\cE=\cE_1\sqcup\cdots\sqcup\cE_q$ and each $t_k$ is supported in one $\cE_i$ (zero vectors assigned to $\cE_1$), then $\cT_t=\cT_{t^{(1)}}\times\cdots\times\cT_{t^{(q)}}$ and
\[
\sG_{\nu;t}(\boldsymbol\xi)=\prod_{i=1}^q\sG_{\nu|_{\cE_i};t^{(i)}}(\boldsymbol\xi^{(i)}),
\]
with $t^{(i)}$ the exponents supported in $\cE_i$ and $\boldsymbol\xi^{(i)}$ their coefficients.
\end{theorem}

% Reference 20: cor:kernelconf; source line 245; cor:oneexp; source line 252 (stated for every square configuration)
\begin{corollary}[Square configurations]\label{cor:square}
Let $N=|\cE|$ and let the matrix $t$ be invertible, and put $\mathbf b=t^{-1}\nu$. Then $\cT_t=\{\nu:\re\mathbf b\in(0,\infty)^N\}$ and
\[
\sG_{\nu;t}(\boldsymbol\xi)=|\det t|^{-1}\prod_{k=1}^N\Gamma(b_k)\,\xi_k^{-b_k}.
\]
In particular $\sG_{\nu;I_{\cE}}(\boldsymbol\xi)=\prod_e\Gamma(\nu_e)\,\xi_e^{-\nu_e}$; on one scale $\sG_{\nu;(t_1)}(\xi)=t_1^{-1}\Gamma(\nu/t_1)\,\xi^{-\nu/t_1}$ for $t_1>0$; and $\Gamma(\nu)=\sG_{\nu;(1)}(1)$, whose exchange relation is $\nu\Gamma(\nu)=\Gamma(\nu+1)$.
\end{corollary}

The Gamma function is thus the Gramma of one scale and one unit exponent, and at every invertible square configuration the Gramma is a stacked Gamma times powers, divided by $|\det t|$. A configuration with more exponents than scales is where the Gramma becomes more than a stacked Gamma times powers: the difference $N-|\cE|$ is the dimension of the share side of Proposition~\ref{prop:shareform}.

% Reference 20: lem:scaleshape; source line 227
\begin{lemma}[Scale and shape]\label{lem:scaleshape}
The map $(0,\infty)\times\Delta_B\to\mathbb R^B_{>0}$, $(\rho,u)\mapsto\rho u$, is a bijection with inverse $x\mapsto(|x|,x/|x|)$, and $dx=\rho^{|B|-1}\,d\rho\,d\omega_B(u)$. Consequently $(\rho u)^{\pi}=\rho^{\ip\pi{\one_B}}u^{\pi}$ for $\pi\in\mathbb R^B$, and $x^{\mvec-\one}dx=\rho^{m_B-1}u^{\mvec-\one}\,d\rho\,d\omega_B(u)$ for $\mvec\in\mathbb C^B$.
\end{lemma}

% Reference 20: thm:scaleshapeset; source line 234 (one set $B$).
% Proved in the last part: thm:scaleshapeset.
\begin{theorem}[Scale and shape along sets]\label{thm:scaleshapeset}
Let $B_1,\dots,B_q$ be pairwise disjoint nonempty subsets of $\cE$, $R=\cE\setminus(B_1\cup\cdots\cup B_q)$ and $\nu\in\cT_t$. With the \emph{levels} $\ell_{k,i}=\ip{t_k}{\one_{B_i}}$, the configuration $\ell=(\ell_1,\dots,\ell_N)$ on $q$ scales and the coefficients $\mathbf b(u,x_R)=\big(\xi_k\prod_iu_i^{t_k|_{B_i}}x_R^{t_k|_R}\big)_k$ for $u=(u_i)_i\in\prod_i\Delta_{B_i}$,
\[
\sG_{\nu;t}(\boldsymbol\xi)=\int_{\mathbb R^R_{>0}}x_R^{\nu|_R-\one}\int_{\prod_i\Delta_{B_i}}\prod_iu_i^{\nu|_{B_i}-\one}\ \sG_{(\nu_{B_i})_i;\ell}\big(\mathbf b(u,x_R)\big)\,d\omega(u)\,dx_R ,
\]
with $d\omega=\prod_id\omega_{B_i}$, the inner Gramma, on $q$ scales at the index $(\nu_{B_i})_i$, $\nu_{B_i}=\sum_{e\in B_i}\nu_e$, converging for almost every $(u,x_R)$; when $R$ is empty there is no outer integral.
\end{theorem}

Passing to scale and shape on disjoint sets of scales replaces each set by one scale, its total, carrying the levels of the exponents on the set; what is left of the exponents dresses the coefficients by the shapes. The Breta of Chapter~\ref{ch:breta} is this passage made into a symbol.

% Reference 20: thm:grammawalls; source line 287
\begin{theorem}[The walls of a Gramma]\label{thm:grammawalls}
Let $\cE=\cE'\sqcup\{e\}$, write $t_k=(t'_k,t_{k,e})$ and $\nu=(\nu',\nu_e)$, and suppose $t_{k,e}\ge0$ for every $k$, with some $t_{k,e}>0$. Let $H=\{k:t_{k,e}=0\}$ be the \emph{facet terms}, suppose $\cone(t'_H)$ has nonempty interior in $\mathbb R^{\cE'}$ and contains $t'_k$ for every $k\notin H$, and let $\re\nu'\in\interior\cone(t'_H)$. Then $\sG_{\nu;t}(\boldsymbol\xi)$ converges for $0<\re\nu_e<c$, for some $c\in(0,\infty]$, and extends to a meromorphic function of $\nu_e\in\mathbb C$ with at most simple poles at $\nu_e=-p$, $p\in P=\{\sum_{k\notin H}n_kt_{k,e}:\ \mathbf n\in\mathbb N^{H^c}\}$, where
\[
\operatorname*{Res}_{\nu_e=-p}\sG_{\nu;t}(\boldsymbol\xi)=\sum_{\mathbf n\in\mathbb N^{H^c}:\ \sum_kn_kt_{k,e}=p}\ \prod_{k\notin H}\frac{(-\xi_k)^{n_k}}{n_k!}\ \sG_{\nu'+\sum_{k\notin H}n_kt'_k;\ t'_H}(\boldsymbol\xi_H).
\]
At $p=0$ the residue is the Gramma of the facet terms: at its first wall a Gramma factors onto its facet. A linear change of the logarithmic coordinates, $y=Ay'$ with $A\in GL_{\cE}(\mathbb R)$, gives $\sG_{\nu;t}(\boldsymbol\xi)=|\det A|\,\sG_{A^{\mathsf T}\nu;A^{\mathsf T}t}(\boldsymbol\xi)$ (Proposition~\ref{prop:grammareflect} at $\mathbf h=0$), and with the last column of $A$ the inward normal of a facet of $\cone(t)$ it puts that facet in the position of $H$.
\end{theorem}

The poles of a Gramma sit on the walls of its cone, and each residue is a Gramma of the facet, of one scale fewer, at an index moved by the exponents that leave the facet. On one scale with positive exponents this is the continuation of the Gramma to the whole index line, with at most simple poles at the points $-p$, $p=\ip{\mathbf n}t$ for $\mathbf n\in\mathbb N^N$, the residue at $-p$ being $\sum_{\ip{\mathbf n}t=p}\prod_k(-\xi_k)^{n_k}/n_k!$.

% Reference 20: prop:expansion; source line 298
\begin{proposition}[The expansion lattice]\label{prop:expansion}
Let $S$ be a set of terms such that the integral with the signs of those terms flipped converges,
\[
\int_{\mathbb R^{\cE}}\exp\Big(\ip{\re\nu}y+\sum_{k\in S}\xi_ke^{\ip{t_k}y}-\sum_{k\notin S}\xi_ke^{\ip{t_k}y}\Big)dy<\infty .
\]
Then, absolutely,
\[
\sG_{\nu;t}(\boldsymbol\xi)=\sum_{\mathbf n\in\mathbb N^S}\ \prod_{k\in S}\frac{(-\xi_k)^{n_k}}{n_k!}\ \sG_{\nu+\sum_{k\in S}n_kt_k;\ t_{S^c}}(\boldsymbol\xi_{S^c}).
\]
\end{proposition}

% Reference 20: thm:decay; source line 313
\begin{theorem}[The decay polytope]\label{thm:decay}
Let $\nu_0\in\interior\cone(t)$ be real, $\eta\in\mathbb R^{\cE}$, and $D_t=\{y\in\mathbb R^{\cE}:\ |\ip{t_k}y|<\frac\pi2\ \forall k\}$. For every $y\in D_t$,
\[
\big|\sG_{\nu_0+i\eta;t}(\boldsymbol\xi)\big|\le e^{-\ip\eta y}\ \sG_{\nu_0;t}\big((\xi_k\cos\ip{t_k}y)_k\big).
\]
Consequently, with $h_t(\eta)=\sup_{y\in D_t}\ip\eta y$, for every $\varepsilon\in(0,1)$
\[
\big|\sG_{\nu_0+i\eta;t}(\boldsymbol\xi)\big|\le\sG_{\nu_0;t}\big(\sin(\tfrac{\pi\varepsilon}2)\boldsymbol\xi\big)\,e^{-(1-\varepsilon)h_t(\eta)},
\]
and $h_t$ is a norm when the $t_k$ span $\mathbb R^{\cE}$.
\end{theorem}

\section{The Zamma}

% Reference 20: def:zamma; source line 330
\begin{definition}[Zamma]\label{def:zamma}
Wherever $\sG_{\nu;t}(\boldsymbol\xi)$ converges and no $\nu_e$ lies in $-\mathbb N$, the \emph{Zamma} is the Gramma divided by the Gammas of its own index,
\[
\sZ_{\nu;t}(\boldsymbol\xi):=\frac{\sG_{\nu;t}(\boldsymbol\xi)}{\prod_{e\in \cE}\Gamma(\nu_e)} .
\]
At the unit configuration it is $\sZ_{\nu;I_{\cE}}(\boldsymbol\xi):=\prod_e\xi_e^{-\nu_e}$ for every $\nu\in\mathbb C^{\cE}$; where $\re\nu_e>0$ for every $e$ this is the quotient $\sG_{\nu;I_{\cE}}(\boldsymbol\xi)/\prod_e\Gamma(\nu_e)$. The Zamma carries the continuations of the Gramma wherever the Gramma has them.
\end{definition}

The Zamma divides out of the Gramma what is first order in its index, the Gammas that the unit configuration produces on its own; what remains is the content of the configuration. It has the three slots of the Gramma.

% Reference 20: prop:zammafacts; source line 338 (one scale in (i) and (iv)).
% Proved in the last part: prop:zammafacts.
\begin{proposition}[First facts about the Zamma]\label{prop:zammafacts}
\begin{enumerate}[label=(\roman*)]
\item \emph{First order.} At the unit configuration the Zamma is the monomial of Definition~\ref{def:zamma}; at an invertible square configuration $t$, $N=|\cE|$, with $\mathbf b=t^{-1}\nu$,
\[
\sZ_{\nu;t}(\boldsymbol\xi)=|\det t|^{-1}\,\Gst{N}{N}{\mathbf b}{\nu}\prod_k\xi_k^{-b_k}.
\]
\item \emph{A power of a sum is a Zamma.} With $N$ unit exponents on one scale, $\sZ_{\nu;(1,\dots,1)}(\boldsymbol\xi)=\big(\sum_k\xi_k\big)^{-\nu}$.
\item \emph{Relations.} Where the terms converge, with quotients of Gammas in place of the factors the index moves,
\[
\partial_{\xi_k}\sZ_{\nu;t}(\boldsymbol\xi)=-\prod_e\Gst11{\nu_e+t_{k,e}}{\nu_e}\sZ_{\nu+t_k;t}(\boldsymbol\xi),
\]
\[
\sZ_{\nu;t}(\boldsymbol\xi)=\sum_kt_{k,e}\,\xi_k\,\Gst11{\nu_e+t_{k,e}}{\nu_e+1}\prod_{f\neq e}\Gst11{\nu_f+t_{k,f}}{\nu_f}\sZ_{\nu+t_k;t}(\boldsymbol\xi)\qquad(e\in \cE).
\]
The other relations of Theorem~\ref{thm:grammalaws} are read with the same normalization, and identities persist on the meromorphic continuations.
\item \emph{As an expectation.} If a configuration on $\cE$ contains the unit configuration $I_{\cE}$ with unit coefficients, then for real $\nu\in\mathbb R^{\cE}_{>0}$
\[
\sZ_{\nu;I_{\cE}\oplus t'}(\one\oplus\boldsymbol\xi')=\oE_{U\sim\operatorname{Gamma}_\nu}\Big[\exp\Big(-\sum_k\xi'_kU^{t'_k}\Big)\Big];
\]
it lies in $(0,1]$, decreases in every coefficient, and, where every $\nu+t'_k$ lies in $\mathbb R^{\cE}_{>0}$,
\[
\sZ_{\nu;I_{\cE}\oplus t'}(\one\oplus\boldsymbol\xi')\ \ge\ \exp\Big(-\sum_k\xi'_k\prod_{e\in\cE}\Gst11{\nu_e+t'_{k,e}}{\nu_e}\Big).
\]
\end{enumerate}
\end{proposition}

% Reference 20: prop:levelpoly; source line 364 (the power $L(\rho)^{-\lambda}$ in (i)--(iii)).
% Proved in the last part: prop:levelpoly.
\begin{proposition}[Powers of a polynomial in the scales]\label{prop:levelpoly}
Let $t$ be a configuration and $\nu\in\cT_t$ with no $\nu_e$ in $-\mathbb N$, or $t=I_{\cE}$ and $\nu\in\mathbb C^{\cE}$; let $\xi_k(\rho)=\sum_{l\in\mathbb N^d}\rho^l\xi^{(l)}_k$ be polynomials in $\rho\in\mathbb R^d$ with real coefficients and $\boldsymbol\xi^{(0)}\in(0,\infty)^N$, and let $c^\nu_q=[\rho^q]\,\sZ_{\nu;t}(\boldsymbol\xi(\rho))$ be the Taylor coefficients at $\rho=0$. At the unit configuration, where $\sZ_{\nu;I_{\cE}}(\boldsymbol\xi)=\boldsymbol\xi^{-\nu}$ is read with the principal branch near $\boldsymbol\xi^{(0)}$, the coefficients $\xi^{(l)}_k$ with $l\neq0$ may be complex and $\rho\in\mathbb C^d$.
\begin{enumerate}[label=(\roman*)]
\item $c^\nu_q=\sum_{\mathbf n}(-1)^{|\mathbf n|}\prod_{e\in\cE}\Gst11{\nu_e+(t\mathbf w)_e}{\nu_e}\,\sZ_{\nu+t\mathbf w;t}(\boldsymbol\xi^{(0)})\prod_{l\neq0,\,k}\frac{(\xi^{(l)}_k)^{n_{l,k}}}{n_{l,k}!}$, the sum over the arrays $\mathbf n=(n_{l,k})_{l\neq0,\,k}$ in $\mathbb N$ with $\sum_{l,k}l\,n_{l,k}=q$, where $\mathbf w=(\sum_ln_{l,k})_k$; at the unit configuration the factor $\prod_e\Gst11{\nu_e+(t\mathbf w)_e}{\nu_e}$ is the polynomial $\prod_e(\nu_e)_{w_e}$, and away from it each product $\prod_e\Gst11{\nu_e+(t\mathbf w)_e}{\nu_e}\sZ_{\nu+t\mathbf w;t}(\boldsymbol\xi^{(0)})$ is $\sG_{\nu+t\mathbf w;t}(\boldsymbol\xi^{(0)})/\prod_e\Gamma(\nu_e)$.
\item At one unit exponent, $t=(1)$, with $\boldsymbol\xi(\rho)=L(\rho)$ and the index $\lambda$: for every $i$, $\sum_lL^{(l)}\,(q_i-l_i+1+\lambda l_i)\,c^\lambda_{q-l+\mathbf e_i}=0$, coefficients with a negative index being $0$.
\item $\partial_{\xi^{(l)}_k}c^\nu_q=-\prod_{e}\Gst11{\nu_e+t_{k,e}}{\nu_e}\,c^{\nu+t_k}_{q-l}$, read in the same way.
\item For a configuration $t$, an index $\nu\in\mathbb C^{\cE}$ and coefficients $\xi_k(\rho)$ depending smoothly on $\rho$ (Proposition~\ref{prop:zammafacts}(iii) and the chain rule),
\[
\rho_i\partial_{\rho_i}\sZ_{\nu;t}(\boldsymbol\xi(\rho))=-\sum_k\prod_e\Gst11{\nu_e+t_{k,e}}{\nu_e}\big(\rho_i\partial_{\rho_i}\xi_k\big)\sZ_{\nu+t_k;t}(\boldsymbol\xi(\rho)).
\]
\end{enumerate}
\end{proposition}


\chapter{The Breta}\label{ch:breta}

\section{Definition and the tower}

% Reference 20: def:breta; source line 381
\begin{definition}[Breta]\label{def:breta}
The data of a Breta are:
\begin{itemize}[leftmargin=1.5em,itemsep=1pt]
\item a finite set $\cE$ of scales and a partition $\cP$ of $\cE$ into \emph{bunches} $Q$; each bunch is read in scale and shape, $x_Q=\rho_Qu_Q$ with $\rho_Q>0$ and $u_Q\in\Delta_Q$;
\item a \emph{shape index} $\mvec\in\mathbb C^{\cE}$ and a \emph{scale index} $\boldsymbol\nu=(\nu_Q)_{Q\in\cP}$;
\item terms $k=1,\dots,N$, each with a \emph{shape exponent} $\pi_k\in\mathbb R^{\cE}$, a \emph{level} $\ell_k=(\ell_{k,Q})_{Q}\in\mathbb R^{\cP}$ and a coefficient $\xi_k>0$.
\end{itemize}
With $u^{\pi_k}=\prod_Qu_Q^{\pi_k|_Q}$ and $d\omega=\prod_Qd\omega_Q$, the \emph{Breta} is
\[
\sB^{\cP}_{\mvec,\boldsymbol\nu;(\pi,\ell)}(\boldsymbol\xi):=\int_{\prod_Q\Delta_Q}u^{\mvec-\one}\ \sZ_{\boldsymbol\nu;\ell}\big((u^{\pi_k}\xi_k)_k\big)\,d\omega(u),
\]
the Zamma on the bunch scales with configuration $\ell$, wherever the integral converges. The data are \emph{homogeneous} when every level is the degree of its shape exponent on each bunch, $\ell_{k,Q}=\ip{\pi_k}{\one_Q}$; then the shape exponents form a configuration $t=(\pi_k)_k$ on $\cE$ and the Breta is written $\sB^{\cP}_{\mvec,\boldsymbol\nu;t}(\boldsymbol\xi)$. When $\cP=\{B\}$ has one bunch, the scale index has one entry and is written as a number, $\sB^{\{B\}}_{\mvec,\nu;t}(\boldsymbol\xi)$.
\end{definition}

The Breta has five slots: the partition, the shape index, the scale index, the terms and the coefficients. It is the Gramma read in scale and shape on every bunch at once (Theorem~\ref{thm:scaleshapeset}), with the scale integrals closed into a Zamma and the shapes left as an integral; the scale index is free, and for homogeneous data fixing it at the bunch totals of the shape index, $\nu_Q=m_Q$, makes the Breta the Gramma itself up to Gammas.

% Reference 20: thm:tower; source line 398
\begin{theorem}[The tower]\label{thm:tower}
Let the data be homogeneous and $\mvec\in\cT_t$. For every partition $\cP$ of the scales,
\[
\sG_{\mvec;t}(\boldsymbol\xi)=\prod_{Q\in\cP}\Gamma(m_Q)\ \sB^{\cP}_{\mvec,(m_Q)_Q;t}(\boldsymbol\xi).
\]
With the finest partition the Breta is $\sZ_{\mvec;t}(\boldsymbol\xi)$; with the coarsest, $\Gamma(m_{\cE})$ times the Breta is the integral of Theorem~\ref{thm:scaleshapeset} with $R=\varnothing$. At the unit configuration $I_B$ of a bunch $B$ and the coarsest partition $\{B\}$, the Zamma inside is $\big(\sum_e\xi_eu_e\big)^{-m_B}$, and
\[
\prod_{e\in B}\Gamma(m_e)\,\xi_e^{-m_e}=\Gamma(m_B)\,\sB^{\{B\}}_{\mvec,m_B;I_B}(\boldsymbol\xi)\qquad(\re m_e>0).
\]
\end{theorem}

% Reference 20: cor:coarsen; source line 427
\begin{corollary}[Coarsening]\label{cor:coarsen}
Let the data be homogeneous, and let $\cP$ refine $\cQ$. Then
\[
\sB^{\cQ}_{\mvec,(m_Q)_Q;t}(\boldsymbol\xi)=\prod_{Q\in\cQ}\Gst{k_Q}{1}{(m_P)_{P\subseteq Q}}{m_Q}\ \sB^{\cP}_{\mvec,(m_P)_P;t}(\boldsymbol\xi),
\]
with $k_Q$ the number of bunches of $\cP$ inside $Q$.
\end{corollary}

% Reference 20: prop:bretaorthant; source line 438 (homogeneous data).
% Proved in the last part: prop:bretaorthant.
\begin{proposition}[The Breta on the orthant]\label{prop:bretaorthant}
For any data,
\[
\sB^{\cP}_{\mvec,\boldsymbol\nu;(\pi,\ell)}(\boldsymbol\xi)=\frac1{\prod_Q\Gamma(\nu_Q)}\int_{\mathbb R^{\cE}_{>0}}x^{\mvec-\one}\prod_Q|x_Q|^{\nu_Q-m_Q}\exp\Big(-\sum_k\xi_k\,x^{\pi_k}\prod_Q|x_Q|^{\ell_{k,Q}-\ip{\pi_k}{\one_Q}}\Big)\,dx ,
\]
either side converging when the other does: a free scale index is one more power of each bunch's total, a level that is not the degree of its shape exponent on a bunch is one more power of that bunch's total inside its term, and for homogeneous data at $\nu_Q=m_Q$ this is Theorem~\ref{thm:tower}.
\end{proposition}

% Reference 20: thm:bretalaws; source line 449 ((v) for homogeneous data).
% Proved in the last part: thm:bretalaws.
\begin{theorem}[Relations of the Breta]\label{thm:bretalaws}
Each identity holds wherever every Breta in it converges. The terms are $(\pi,\ell)$; homogeneous data have $t_k=\pi_k$ and $\ell_{k,Q}=\ip{t_k}{\one_Q}$. In this statement write $g_k=\prod_{Q}\Gst11{\nu_Q+\ell_{k,Q}}{\nu_Q}$.
\begin{enumerate}[label=(\roman*)]
\item \emph{Shape split}, per bunch $Q$:
\[
\sB^{\cP}_{\mvec,\boldsymbol\nu;(\pi,\ell)}(\boldsymbol\xi)
=\sum_{e\in Q}\sB^{\cP}_{\mvec+\mathbf e_e,\boldsymbol\nu;(\pi,\ell)}(\boldsymbol\xi).
\]
\item \emph{Ladder}, per term $k$:
\[
\partial_{\xi_k}\sB^{\cP}_{\mvec,\boldsymbol\nu;(\pi,\ell)}(\boldsymbol\xi)
=-g_k\,\sB^{\cP}_{\mvec+\pi_k,\boldsymbol\nu+\ell_k;(\pi,\ell)}(\boldsymbol\xi).
\]
\item \emph{Scale exchange}, per bunch $Q$:
\[
\sB^{\cP}_{\mvec,\boldsymbol\nu;(\pi,\ell)}(\boldsymbol\xi)
=\sum_k\ell_{k,Q}\,\xi_k\,\Gst11{\nu_Q+\ell_{k,Q}}{\nu_Q+1}\prod_{Q'\neq Q}\Gst11{\nu_{Q'}+\ell_{k,Q'}}{\nu_{Q'}}\ \sB^{\cP}_{\mvec+\pi_k,\boldsymbol\nu+\ell_k;(\pi,\ell)}(\boldsymbol\xi).
\]
\item \emph{Shape exchange}, for scales $e\neq f$ of one bunch, with $\re m_e,\re m_f>0$:
\[
\begin{split}
m_e\,\sB^{\cP}_{\mvec+\mathbf e_f,\boldsymbol\nu;(\pi,\ell)}(\boldsymbol\xi)
-m_f\,\sB^{\cP}_{\mvec+\mathbf e_e,\boldsymbol\nu;(\pi,\ell)}(\boldsymbol\xi)
={}&\sum_k\xi_k\,g_k\Big[
\pi_{k,e}\,\sB^{\cP}_{\mvec+\pi_k+\mathbf e_f,\boldsymbol\nu+\ell_k;(\pi,\ell)}(\boldsymbol\xi)\\
&\hspace{3.2em}-\pi_{k,f}\,\sB^{\cP}_{\mvec+\pi_k+\mathbf e_e,\boldsymbol\nu+\ell_k;(\pi,\ell)}(\boldsymbol\xi)\Big].
\end{split}
\]
\item \emph{Shape exchange on the orthant}, for scales $e\neq f$ of one bunch $Q$:
\[
\begin{split}
&(m_e-m_f)\,\sB^{\cP}_{\mvec,\boldsymbol\nu;(\pi,\ell)}(\boldsymbol\xi)
+(\nu_Q-m_Q)\Big(\sB^{\cP}_{\mvec+\mathbf e_e,\boldsymbol\nu;(\pi,\ell)}(\boldsymbol\xi)
-\sB^{\cP}_{\mvec+\mathbf e_f,\boldsymbol\nu;(\pi,\ell)}(\boldsymbol\xi)\Big)\\
&\quad=\sum_k\xi_k\,g_k\Big[(\pi_{k,e}-\pi_{k,f})\,\sB^{\cP}_{\mvec+\pi_k,\boldsymbol\nu+\ell_k;(\pi,\ell)}(\boldsymbol\xi)\\
&\qquad\qquad+\big(\ell_{k,Q}-\ip{\pi_k}{\one_Q}\big)\Big(\sB^{\cP}_{\mvec+\pi_k+\mathbf e_e,\boldsymbol\nu+\ell_k;(\pi,\ell)}(\boldsymbol\xi)-\sB^{\cP}_{\mvec+\pi_k+\mathbf e_f,\boldsymbol\nu+\ell_k;(\pi,\ell)}(\boldsymbol\xi)\Big)\Big];
\end{split}
\]
for homogeneous data the second line of the right side is absent.
\end{enumerate}
For homogeneous data a term raises the shape total of each bunch and its scale index by the same level, so the ladder, the scale exchange and the shape exchange (v) keep the condition $\nu_Q=m_Q$; in (v) the shifts of the shape index alone carry the factor $\nu_Q-m_Q$, which vanishes there. The shape split and the shape exchange (iv) shift the shape index alone.
\end{theorem}

\section{Letters}\label{sec:letters}

A \emph{letter} is a function that enters an integrand through a power $P^{-s}$, $s$ its \emph{letter exponent}. On bunches read in scale and shape, a polynomial letter with positive coefficients is a sum of patterns, each a monomial in the scales and the shapes; the letter presentation writes the integral of a kernel against a product of letters as one Breta, by giving every letter a scale of its own.

% Reference 20: prop:letters; source line 475 (one bunch $B$).
% Proved in the last part: prop:letters.
\begin{proposition}[Letter presentation]\label{prop:letters}
Let the scales $B$ carry a partition $\cQ$ into bunches, each read in scale and shape, $x_Q=\rho_Qu_Q$, with $\rho=(\rho_Q)_{Q\in\cQ}$ and $u=(u_Q)_{Q\in\cQ}$; let \emph{kernel terms} $g_j\rho^{\lambda_j}u^{\pi_j}$ ($g_j>0$, $\lambda_j\in\mathbb R^{\cQ}$, $\pi_j\in\mathbb R^B$) and letters
\[
P_i(\rho,u)=\sum_{p}C_{i,p}\,\rho^{\ell_{i,p}}u^{\pi_{i,p}},\qquad C_{i,p}>0,\qquad i=1,\dots,N,
\]
be given, the sum over the \emph{patterns} $p$ of the letter $P_i$, with $\ell_{i,p}\in\mathbb R^{\cQ}$, $\pi_{i,p}\in\mathbb R^B$ and letter exponents $s_i$, $\re s_i>0$. The array $(C_{i,p})$ lists the coefficients of the patterns. Introduce one scale $v_i$ per letter, put $V=\{v_1,\dots,v_N\}$, and take the partition of the scales $B\sqcup V$
\[
\cP=\cQ\cup\{\{v_1\},\dots,\{v_N\}\}.
\]
On $B\sqcup V$ the shape index is $(\mvec,\mathbf s)$, and on the bunches the scale index is $(\boldsymbol\nu,\mathbf s)$, $\boldsymbol\nu\in\mathbb C^{\cQ}$. The term arrays $(\widetilde\pi,\widetilde\ell)$ have one column per kernel term and one per pattern:
\[
\begin{array}{c|c|c|c}
\text{column}&\text{shape exponent on }B\sqcup V&\text{level on }\cP&\text{coefficient}\\ \hline
j&(\pi_j,0_V)&(\lambda_j,0_V)&g_j\\
(i,p)&(\pi_{i,p},\mathbf e_i)&(\ell_{i,p},\mathbf e_i)&C_{i,p}
\end{array}
\]
with $\mathbf e_i$ the unit vector at $v_i$ in the shape column and at the singleton bunch $\{v_i\}$ in the level column, and $\mathbf g\oplus\mathbf C$ lists the coefficients in this column order. Then
\begin{multline*}
\int_{\prod_Q\Delta_Q}u^{\mvec-\one}\int_{\mathbb R^{\cQ}_{>0}}\rho^{\boldsymbol\nu-\one}
\exp\Big(-\sum_jg_j\rho^{\lambda_j}u^{\pi_j}\Big)
\prod_iP_i(\rho,u)^{-s_i}\,d\rho\,d\omega(u)\\
=\prod_{Q\in\cQ}\Gamma(\nu_Q)\,
\sB^{\cP}_{(\mvec,\mathbf s),(\boldsymbol\nu,\mathbf s);(\widetilde\pi,\widetilde\ell)}
(\mathbf g\oplus\mathbf C).
\end{multline*}
The shapes of the singleton bunches are points, so the entries of the shape index on $V$ do not enter the integral; they move with the scale entries on $V$, which keeps the relations of Theorem~\ref{thm:bretalaws} literal. In the rest of this statement the partition, the term arrays and the coefficients are fixed, and
\[
b(\mvec,\boldsymbol\nu;\mathbf s)\ \text{ stands for }\ \sB^{\cP}_{(\mvec,\mathbf s),(\boldsymbol\nu,\mathbf s);(\widetilde\pi,\widetilde\ell)}(\mathbf g\oplus\mathbf C).
\]
The relations of Theorem~\ref{thm:bretalaws} read: the scale exchange at $\{v_i\}$ is the \emph{letter split}
\[
b(\mvec,\boldsymbol\nu;\mathbf s)=\sum_{p}C_{i,p}\prod_{Q\in\cQ}\Gst11{\nu_Q+\ell_{i,p,Q}}{\nu_Q}\,b(\mvec+\pi_{i,p},\boldsymbol\nu+\ell_{i,p};\mathbf s+\mathbf e_i);
\]
the ladders are
\begin{gather*}
\partial_{g_j}b(\mvec,\boldsymbol\nu;\mathbf s)=-\prod_{Q\in\cQ}\Gst11{\nu_Q+\lambda_{j,Q}}{\nu_Q}\,b(\mvec+\pi_j,\boldsymbol\nu+\lambda_j;\mathbf s),\\
\partial_{C_{i,p}}b(\mvec,\boldsymbol\nu;\mathbf s)=-s_i\prod_{Q\in\cQ}\Gst11{\nu_Q+\ell_{i,p,Q}}{\nu_Q}\,b(\mvec+\pi_{i,p},\boldsymbol\nu+\ell_{i,p};\mathbf s+\mathbf e_i);
\end{gather*}
and the scale exchange on each $Q\in\cQ$ is
\begin{multline*}
\nu_Q\,b(\mvec,\boldsymbol\nu;\mathbf s)=\sum_j\lambda_{j,Q}\,g_j\prod_{Q'\in\cQ}\Gst11{\nu_{Q'}+\lambda_{j,Q'}}{\nu_{Q'}}\,b(\mvec+\pi_j,\boldsymbol\nu+\lambda_j;\mathbf s)\\
+\sum_{i,p}s_i\,\ell_{i,p,Q}\,C_{i,p}\prod_{Q'\in\cQ}\Gst11{\nu_{Q'}+\ell_{i,p,Q'}}{\nu_{Q'}}\,b(\mvec+\pi_{i,p},\boldsymbol\nu+\ell_{i,p};\mathbf s+\mathbf e_i).
\end{multline*}
\end{proposition}

% Reference 20: prop:letterorthant; source line 497 (homogeneous letter data on one bunch).
% Proved in the last part: prop:letterorthant.
\begin{proposition}[Letter data on the orthant]\label{prop:letterorthant}
Use the partition and the term arrays of Proposition~\ref{prop:letters}, and for $\alpha\in\mathbb R^B_{>0}$ write $P_i(\alpha)$ for the letter at $\rho_Q=|\alpha_Q|$, $u_Q=\alpha_Q/|\alpha_Q|$:
\[
P_i(\alpha)=\sum_pC_{i,p}\,\alpha^{\pi_{i,p}}\prod_{Q\in\cQ}|\alpha_Q|^{\ell_{i,p,Q}-\ip{\pi_{i,p}}{\one_Q}} .
\]
Then, either side converging when the other does,
\begin{multline*}
\prod_{Q\in\cQ}\Gamma(\nu_Q)\,
\sB^{\cP}_{(\mvec,\mathbf s),(\boldsymbol\nu,\mathbf s);(\widetilde\pi,\widetilde\ell)}
(\mathbf g\oplus\mathbf C)\\
=\int_{\mathbb R^B_{>0}}\alpha^{\mvec-\one}\prod_{Q\in\cQ}|\alpha_Q|^{\nu_Q-m_Q}
\exp\Big(-\sum_jg_j\,\alpha^{\pi_j}\prod_{Q\in\cQ}|\alpha_Q|^{\lambda_{j,Q}-\ip{\pi_j}{\one_Q}}\Big)\prod_iP_i(\alpha)^{-s_i}\,d\alpha :
\end{multline*}
each bunch total $|\alpha_Q|$ is one more letter, at the exponent $m_Q-\nu_Q$, and it enters every kernel term and pattern whose level on $Q$ is not its degree there. For homogeneous letter data, $\lambda_{j,Q}=\ip{\pi_j}{\one_Q}$ and $\ell_{i,p,Q}=\ip{\pi_{i,p}}{\one_Q}$, the letters are $P_i(\alpha)=\sum_pC_{i,p}\alpha^{\pi_{i,p}}$, and at $\nu_Q=m_Q$ the totals are absent.
\end{proposition}

% Reference 20: prop:collapse; source line 508
\begin{proposition}[Equal levels collapse]\label{prop:collapse}
In a Breta the dressed coefficients of all terms with one level add: the Zamma inside has one term per distinct level, whose coefficient is that level's polynomial in the shape. For homogeneous letters of a bunch the level of a pattern is its degree.
\end{proposition}

\section{The shape integral}

% Reference 20: thm:dirichlet; source line 517
\begin{theorem}[The shape integral]\label{thm:dirichlet}
For $\re m_e>0$ ($e\in B$),
\[
\int_{\Delta_B}u^{\mvec-\one}\,d\omega_B(u)=\sB^{\{B\}}_{\mvec,m_B;I_B}(\one)=\Gst{|B|}{1}{\mvec}{m_B}:
\]
the tower (Theorem~\ref{thm:tower}) at the unit configuration $I_B$ with unit coefficients, where the Zamma inside the Breta is $\big(\sum_eu_e\big)^{-m_B}=1$.
\end{theorem}

The stacked Gamma of the masses over their total is therefore the integral of the shape density $u^{\mvec-\one}$ on the simplex, and the Dirichlet law is that density normalised.

% Reference 20: def:dirlaw; source line 535
\begin{definition}[Dirichlet law]\label{def:dirlaw}
For $\mvec\in\mathbb R^B_{>0}$, $\Dir_{\mvec}(du):=u^{\mvec-\one}\,d\omega_B(u)\big/\Gst{|B|}{1}{\mvec}{m_B}$, a probability measure on $\Delta_B$.
\end{definition}

% Reference 20: thm:dirmoments; source line 539
\begin{theorem}[Moments and steps]\label{thm:dirmoments}
For $\mvec\in\mathbb R^B_{>0}$ and $\mathbf k\in\mathbb C^B$ with $\re(m_e+k_e)>0$ for every $e\in B$,
\[
\oE_{U\sim\Dir_{\mvec}}\big[U^{\mathbf k}\big]=\frac{\Gst{|B|}{1}{\mvec+\mathbf k}{m_B+|\mathbf k|}}{\Gst{|B|}{1}{\mvec}{m_B}},
\]
which for $\mathbf k\in\mathbb N^B$ is $\prod_e(m_e)_{k_e}/(m_B)_{|\mathbf k|}$. For every $\mvec\in\mathbb C^B$ off the poles the stacked Gamma $\Gst{|B|}{1}{\mvec}{m_B}$ satisfies the shape split and the shape exchange of its Breta $\sB^{\{B\}}_{\mvec,m_B;I_B}(\one)$:
\begin{gather*}
\Gst{|B|}{1}{\mvec}{m_B}
=\sum_{e\in B}\Gst{|B|}{1}{\mvec+\mathbf e_e}{m_B+1},\qquad
m_e\Gst{|B|}{1}{\mvec+\mathbf e_f}{m_B+1}
=m_f\Gst{|B|}{1}{\mvec+\mathbf e_e}{m_B+1},\\
\Gst{|B|}{1}{\mvec+\mathbf e_e}{m_B+1}
=\frac{m_e}{m_B}\Gst{|B|}{1}{\mvec}{m_B}.
\end{gather*}
\end{theorem}

% Reference 20: thm:aggregation; source line 561
\begin{theorem}[Aggregation]\label{thm:aggregation}
Let $B=B_1\sqcup\dots\sqcup B_q$ with nonempty blocks and $J=\{1,\dots,q\}$. The map $(w,v^{(1)},\dots,v^{(q)})\mapsto u$, $u|_{B_j}=w_jv^{(j)}$, is a bijection $\Delta_J\times\prod_j\Delta_{B_j}\to\Delta_B$, and
\[
u^{\mvec-\one}d\omega_B=\prod_jw_j^{m_{B_j}-1}\,d\omega_J(w)\prod_j(v^{(j)})^{\mvec|_{B_j}-\one}d\omega_{B_j}(v^{(j)}).
\]
Consequently $\Gst{|B|}{1}{\mvec}{m_B}=\Gst{q}{1}{(m_{B_j})_j}{m_B}\prod_j\Gst{|B_j|}{1}{\mvec|_{B_j}}{m_{B_j}}$, and under $\Dir_{\mvec}$ the block totals and the block shapes are independent, with laws $\Dir_{(m_{B_j})_j}$ and $\Dir_{\mvec|_{B_j}}$.
\end{theorem}

\section{Products of shape integrals}

% Reference 20: thm:shruffle; source line 591 (two bunches).
% Proved in the last part: thm:shruffle.
\begin{theorem}[The Shruffle]\label{thm:shruffle}
For a bunch $Q$ and a total order $\varsigma$ on $Q$ let $\Delta^\varsigma_Q=\{u\in\Delta_Q:u_{\varsigma(1)}>u_{\varsigma(2)}>\cdots\}$ and $I^\varsigma_Q(f;\mvec)=\int_{\Delta^\varsigma_Q}u^{\mvec-\one}f(u)\,d\omega_Q$. For pairwise disjoint bunches $A_1,\dots,A_q$ with total orders $\varsigma_i$ on $A_i$, functions $f_i$ on $\Delta_{A_i}$, $A=A_1\sqcup\cdots\sqcup A_q$ and $\bigotimes_if_i(w)=\prod_if_i(w_{A_i}/|w_{A_i}|)$ on $\Delta_A$, with $\re m_{A_i}>0$ and the integrals absolutely convergent,
\[
\prod_{i=1}^qI^{\varsigma_i}_{A_i}(f_i;\mvec|_{A_i})=\Gst1q{m_A}{m_{A_1},\ \dots,\ m_{A_q}}\sum_{\varsigma}I^{\varsigma}_{A}\Big(\bigotimes_if_i;\mvec\Big),
\]
the sum over the orders $\varsigma$ of $A$ that restrict to every $\varsigma_i$.
\end{theorem}

The product of shape integrals on disjoint bunches $A_1,\dots,A_q$ is $\Gst1q{m_A}{m_{A_1},\ \dots,\ m_{A_q}}$ times the sum of the shape integrals on the joined bunch over the interleavings of their orders: this is the Shruffle, the product rule of the shape side, whose counterpart for Trailmoments at a common cut is the Struffle (Theorem~\ref{thm:strufflefree}).

\section{The Beta port}

% Reference 20: def:betaport; source line 604
\begin{definition}[Beta port]\label{def:betaport}
For $q\in\mathbb N$ and $a,b\in\mathbb C$,
\begin{align*}
\Pi^{a,b}_q(x)&=\frac{(-1)^qx^{1-a}(1-x)^{1-b}}{q!}\,\partial_x^q\big[x^{a+q-1}(1-x)^{b+q-1}\big]\\
&=\sum_{k=0}^q(-1)^{q+k}\binom{a+q-1}{q-k}\binom{b+q-1}{k}x^k(1-x)^{q-k},
\end{align*}
a polynomial of degree at most $q$. For positive integers $a,b$ its coefficients are integers; for fixed $q$ and $a,b$ affine in an integer parameter, its coefficients are integers of polynomial growth in that parameter.
\end{definition}

% Reference 20: thm:betaport; source line 613 ((iii) for $d_i\ge0$; (v) at $a=b=n+1$).
% Proved in the last part: thm:betaport.
\begin{theorem}[Relations of the Beta port]\label{thm:betaport}
\begin{enumerate}[label=(\roman*)]
\item \emph{Action on a monomial.} For $\re c>0$ and $\re b>0$: $\int_0^1x^{c-1}(1-x)^{b-1}\,\Pi^{a,b}_q(x)\,dx=\binom{c-a}{q}\Gst21{c,\ b+q}{c+b+q}$. With $c=a+j$, $j\in\mathbb N$, the binomial vanishes for $j<q$: the Beta port is orthogonal, for the density $x^{a-1}(1-x)^{b-1}$, to every polynomial of degree below $q$.
\item \emph{Integration by parts.} For $\re a,\re b>0$ and $g$ with $q$ continuous derivatives on $[0,1]$,
\[
\int_0^1x^{a-1}(1-x)^{b-1}\,\Pi^{a,b}_q(x)\,g(x)\,dx=\frac1{q!}\int_0^1x^{a+q-1}(1-x)^{b+q-1}\,g^{(q)}(x)\,dx ;
\]
with one Beta port per scale, of orders $\mathbf n=(n_e)_{e\in \cE}$,
\begin{multline*}
\int_{[0,1]^{\cE}}\prod_ex_e^{a_e-1}(1-x_e)^{b_e-1}\Pi^{a_e,b_e}_{n_e}(x_e)\,g(\mathbf x)\,d\mathbf x\\
=\frac1{\prod_en_e!}\int_{[0,1]^{\cE}}\prod_ex_e^{a_e+n_e-1}(1-x_e)^{b_e+n_e-1}\,\partial^{\mathbf n}g(\mathbf x)\,d\mathbf x .
\end{multline*}
\item \emph{Letter ladder.} For letters $L_i=c_i+d_ix$ positive on $[0,1]$ ($c_i>0$ and $c_i+d_i>0$) and exponents $w_i\in\mathbb C$,
\[
\frac{(-1)^q}{q!}\,\partial_x^q\prod_iL_i^{-w_i}=\prod_iL_i^{-w_i}\sum_{|v|=q}\prod_i\frac{(w_i)_{v_i}}{v_i!}\Big(\frac{d_i}{L_i}\Big)^{v_i},
\]
so that, for $\re a,\re b>0$,
\begin{multline*}
\int_0^1x^{a-1}(1-x)^{b-1}\,(-1)^q\Pi^{a,b}_q\prod_iL_i^{-w_i}\,dx\\
=\sum_{|v|=q}\prod_i\frac{(w_i)_{v_i}}{v_i!}\int_0^1x^{a-1}(1-x)^{b-1}\prod_iL_i^{-w_i}\prod_i\Big(\frac{x(1-x)d_i}{L_i}\Big)^{v_i}dx .
\end{multline*}
Let $a,b>0$ and every $w_i\ge0$. If every $d_i\ge0$ the value is nonnegative, and if every $d_i\le0$ it is $(-1)^q$ times a nonnegative number; whatever the signs of the $d_i$, if $x(1-x)|d_i|/L_i\le\varepsilon$ on $[0,1]$ for every $i$, its modulus is at most $\frac{(|\mathbf w|)_q}{q!}\varepsilon^q\int_0^1x^{a-1}(1-x)^{b-1}\prod_iL_i^{-w_i}dx$.
\item \emph{Reflection.} $\Pi^{a,b}_q(1-x)=(-1)^q\Pi^{b,a}_q(x)$.
\item \emph{Three-term and derivative relations.} For $q\ge1$ and $a,b\in\mathbb C$ with $a+b+2q-i\neq0$ for $i=0,1,2$,
\begin{gather*}
x\,\Pi^{a,b}_q=\frac{(q+1)(a+b+q-1)}{(a+b+2q)(a+b+2q-1)}\,\Pi^{a,b}_{q+1}+\Big(\frac{(q+1)(a+q)}{a+b+2q}-\frac{q(a+q-1)}{a+b+2q-2}\Big)\Pi^{a,b}_q\\
\hspace{14em}+\frac{(a+q-1)(b+q-1)}{(a+b+2q-1)(a+b+2q-2)}\,\Pi^{a,b}_{q-1},\\
x(1-x)\,\partial_x\Pi^{a,b}_q=\Big(\frac{q(b+q-1)}{a+b+2q-2}-q\,x\Big)\Pi^{a,b}_q+\frac{(a+q-1)(b+q-1)}{a+b+2q-2}\,\Pi^{a,b}_{q-1};
\end{gather*}
for every $q\in\mathbb N$ the coefficient of $x^j$ in $\Pi^{a,b}_q$ is $(-1)^{q+j}\binom{a+q-1}{q-j}\binom{a+b+q+j-2}{j}$, and that of $x^q$ is $\binom{a+b+2q-2}{q}$.
\end{enumerate}
\end{theorem}

The Beta port moves derivatives of an integrand onto its shape density: integrating against $\Pi^{a,b}_q$ is integrating the $q$-th derivative against the density with both exponents raised by $q$. Against a product of linear letters whose slopes have one sign it keeps the sign of the integrand, up to $(-1)^q$ when the slopes are negative, which is what the arithmetic of Part~\ref{part:arith} uses. The Beta port of order $q$ is orthogonal, for the density $x^{a-1}(1-x)^{b-1}$, to every polynomial of lower degree (Theorem~\ref{thm:betaport}(i)), and at every pair of exponents multiplication by $x$ and the operator $x(1-x)\partial_x$ move each Beta port by at most one step in its order (Theorem~\ref{thm:betaport}(v)): finite coefficient rules, which is what turns a mark relation (Theorem~\ref{thm:markrelation}) into a recurrence of bounded width.


\chapter{Fibre integrals and closing}\label{ch:fibre}

% Reference 20: def:fibre; source line 645
\begin{definition}[Fibre integral]\label{def:fibre}
Let $\mathbf b\in\mathbb C^N$ be \emph{shares}, pinned by $P\mathbf b=\boldsymbol\pi$ with $P$ a real $q\times N$ matrix with linearly independent rows and $\boldsymbol\pi\in\mathbb C^q$. Let $f$ be holomorphic on a tube $\{\re\mathbf b\in U\}$, and let $\mathbf c\in U$ be an \emph{admissible real point}, $P\mathbf c=\re\boldsymbol\pi$. The \emph{fibre integral} is
\[
\int_{P\mathbf b=\boldsymbol\pi}f:=(2\pi)^{q-N}\int_{\mathbb R^N}f(\mathbf c+i\mathbf y)\,\delta^{q}\big(P\mathbf y-\im\boldsymbol\pi\big)\,d\mathbf y=\frac{(2\pi)^{q-N}}{\sqrt{\det PP^{\mathsf T}}}\int_{P\mathbf y=\im\boldsymbol\pi}f(\mathbf c+i\mathbf y)\,d\mathrm{vol}_{N-q}(\mathbf y)
\]
whenever it converges, $\mathrm{vol}_{N-q}$ being the $(N-q)$-dimensional volume on the affine subspace; by Cauchy's theorem it does not depend on $\mathbf c$ within one connected set of admissible real points. Its \emph{dimension} is $N-q$. With one pin, $P=(1,\dots,1)$ and $\boldsymbol\pi=s$, it is the \emph{share integral}
\[
\int_{a_1+\cdots+a_k=s}f(\mathbf a)=\frac{1}{(2\pi i)^{k-1}}\int_{c_1+i\mathbb R}\cdots\int_{c_{k-1}+i\mathbb R}f\Big(a_1,\dots,a_{k-1},\,s-\sum_{i<k}a_i\Big)\,da_1\cdots da_{k-1}.
\]
At dimension zero, $q=N$, the value is $|\det P|^{-1}f(P^{-1}\boldsymbol\pi)$. When $\boldsymbol\pi$ is real, with $V=\{\mathbf y:\ P\mathbf y=0\}$ and the measure $d\boldsymbol\vartheta=(\det PP^{\mathsf T})^{-1/2}\,d\mathrm{vol}_{N-q}$ on it, the fibre integral is written $\int_{\mathbf c+iV}f=(2\pi)^{q-N}\int_Vf(\mathbf c+i\boldsymbol\vartheta)\,d\boldsymbol\vartheta$, its admissible real point shown.
\end{definition}

% Reference 20: prop:shareform; source line 668
\begin{proposition}[The share side of the Gramma]\label{prop:shareform}
For $\nu\in\cT_t$,
\[
\sG_{\nu;t}(\boldsymbol\xi)=\int_{t\mathbf b=\nu}\ \prod_{k=1}^N\Gamma(b_k)\,\xi_k^{-b_k},
\]
the fibre integral with one share per term, pinned by the configuration ($|\cE|$ pins), of dimension $N-|\cE|$; its admissible real points are the $\mathbf c\in(0,\infty)^N$ with $t\mathbf c=\re\nu$, and they exist exactly on the domain. At dimension zero it is the value of Corollary~\ref{cor:square}.
\end{proposition}

On the share side the configuration is the matrix of the pins and the index their values, one pin per scale; each term carries one share, with its Gamma and its coefficient's power as factors.

% Reference 20: thm:configgauge; source line 683 (the index $\lambda\nu$ in (iii)).
% Proved in the last part: thm:configgauge.
\begin{theorem}[The native gauge of a configuration]\label{thm:configgauge}
Let the $t_k$ span $\mathbb R^{\cE}$, $\nu\in\interior\cone(t)$ real and $\boldsymbol\xi>0$.
\begin{enumerate}[label=(\roman*)]
\item $\Phi_{\boldsymbol\xi}(y)=\sum_k\xi_ke^{\ip{t_k}y}-\ip\nu y$ is strictly convex and tends to $+\infty$ along every ray; its unique critical point $y^\star$ solves $\sum_k\xi_ke^{\ip{t_k}{y^\star}}t_k=\nu$, and $e^{y^\star}$ is the \emph{native gauge} of the configuration.
\item The shares $w^\star_k=\xi_ke^{\ip{t_k}{y^\star}}$ are an admissible real point of the share side (Proposition~\ref{prop:shareform}). On the real fibre $\{\mathbf w>0:\ t\mathbf w=\nu\}$ the function $S(\mathbf w)=\sum_k\big(w_k-w_k\log(w_k/\xi_k)\big)$ is strictly concave with unique critical point $\mathbf w^\star$, and $S(\mathbf w^\star)=\Phi_{\boldsymbol\xi}(y^\star)$.
\item For $\nu'\in\mathbb C^{\cE}$, as $\lambda\to\infty$,
\[
\sG_{\lambda\nu+\nu';t}(\lambda\boldsymbol\xi)=e^{\ip{\nu'}{y^\star}}\,e^{-\lambda\Phi_{\boldsymbol\xi}(y^\star)}\Big(\frac{2\pi}\lambda\Big)^{|\cE|/2}\det\Big(\sum_kw^\star_kt_kt_k^{\mathsf T}\Big)^{-1/2}\big(1+O(\lambda^{-1})\big).
\]
\end{enumerate}
\end{theorem}

The native gauge of a configuration is read on both of its sides at once: on the orthant side its logarithmic coordinate $y^\star$ minimizes the convex function $\Phi_{\boldsymbol\xi}$, on the share side the shares $\mathbf w^\star$ it determines maximize the entropy-like function $S$ on the real fibre, and the minimum equals the maximum.

% Reference 20: thm:E1; source line 705
\begin{corollary}[Splitting a coefficient]\label{thm:E1}
Let $k\ge2$, $X\in(0,\infty)^k$, $\re s>0$ and $\mathbf c\in(0,\infty)^{k-1}$ with $\sum_ic_i<\re s$. On this contour, every share having positive real part,
\[
\sZ_{s;(1)}(X_1+\cdots+X_k)=\Big(\sum_{i=1}^kX_i\Big)^{-s}=\int_{a_1+\cdots+a_k=s}\Gst k1{\mathbf a}{s}\prod_{i=1}^kX_i^{-a_i},
\]
equivalently $\Gamma(s)(X_1+\cdots+X_k)^{-s}=\int_{|\mathbf a|=s}\prod_i\Gamma(a_i)X_i^{-a_i}$: the share side of the one-scale Gramma $\sG_{s;(1,\dots,1)}(X)$ (Proposition~\ref{prop:shareform}). The stacked Gamma $\Gst k1{\mathbf a}{s}$ is the \emph{share kernel} of the letter $X_1+\cdots+X_k$ at the exponent $s$.
\end{corollary}

% Reference 20: thm:closing; source line 716
\begin{theorem}[Closing]\label{thm:closing}
Let a fibre integral have the form $\int_{P\mathbf b=\boldsymbol\pi}\Gamma\big(\ell(\mathbf b)\big)G(\mathbf b)$ with $\ell$ affine and not constant on the fibre. Suppose that, moving the contour along a direction of the fibre in which $\ell$ decreases, $\Gamma(\ell)$ is the only factor with poles crossed, and the shifted contour's contribution tends to zero. Then
\[
\int_{P\mathbf b=\boldsymbol\pi}\Gamma\big(\ell(\mathbf b)\big)\,G(\mathbf b)=\sum_{j\ge0}\frac{(-1)^j}{j!}\int_{P\mathbf b=\boldsymbol\pi,\ \ell(\mathbf b)=-j}G(\mathbf b),
\]
each term a fibre integral with one more pin. When several factors have poles crossed, each contributes its own lattice; where several share a pole hyperplane, the residue is a Laurent coefficient of the Gamma factors (Chapter~\ref{ch:trail}).
\end{theorem}

Closing trades one dimension of the fibre for a sum over the poles crossed: the Gamma whose argument falls is replaced by its residues. Where closing can be repeated until the fibre has dimension zero, every term is the value $|\det P|^{-1}G(P^{-1}\boldsymbol\pi)$ of the remaining factors at one point.

% Reference 20: thm:sharemerge; source line 737
\begin{theorem}[The share kernels form a semigroup]\label{thm:sharemerge}
For a letter $\sum_{i=1}^kX_i$, shares $\mathbf a,\mathbf b,\mathbf g\in\mathbb C^k$ and $|\mathbf g|=s+s'$, on admissible contours,
\[
\int_{\mathbf a+\mathbf b=\mathbf g,\ |\mathbf a|=s}\Gst k1{\mathbf a}{s}\Gst k1{\mathbf b}{s'}=\Gst k1{\mathbf g}{s+s'}.
\]
\end{theorem}

Splitting the letter at the exponent $s$ and at the exponent $s'$ and merging the shares is splitting it at $s+s'$: the share kernels of Corollary~\ref{thm:E1} compose like the powers they open.

% Reference 20: cor:fourgammas; source line 747 (the display alone).
% Proved in the last part: cor:fourgammas.
\begin{corollary}[Four Gammas on a line]\label{cor:fourgammas}
On a contour separating the poles of $\Gamma(a+z)\Gamma(b+z)$ from those of $\Gamma(c-z)\Gamma(d-z)$,
\[
\frac1{2\pi i}\int\Gamma(a+z)\,\Gamma(b+z)\,\Gamma(c-z)\,\Gamma(d-z)\,dz=\Gst41{a+c,\ b+d,\ a+d,\ b+c}{a+b+c+d}:
\]
Theorem~\ref{thm:sharemerge} at $k=2$, multiplied by $\Gamma(a+c)\Gamma(b+d)$, with the shares $(a+z,\,c-z)$ at the exponent $a+c$ and $(d-z,\,b+z)$ at the exponent $b+d$, merged to $(a+d,\,b+c)$.
\end{corollary}

The four Gammas are the Gamma factors of the four shares of a Gramma on three scales, pinned by the three scales, and the one free share is the contour variable; integrating the orthant side scale by scale leaves $\Gamma(a+c)\Gamma(b+d)$ times the Beta corner $\Gst21{a+d,\ b+c}{a+b+c+d}$, the shape integral of two scales (Theorem~\ref{thm:dirichlet}).


\chapter{The Trailmoment}\label{ch:trail}

\section{Definition}

% Reference 20: def:jointtrail; source line 945 (with def:trailmoment, source line 759)
\begin{definition}[Trailmoment]\label{def:jointtrail}
The data of a Trailmoment are an index $\nu$ and a configuration $t$, those of a Zamma; a measure $\phi$, the \emph{input}, on a space $Y$, the \emph{source}; a measurable \emph{coefficient function} $\mathbf c:Y\to(0,\infty)^N$; and a measurable \emph{cut region} $Y'\subseteq Y$. The \emph{Trailmoment} is
\[
(\sT_{\nu;t}[\mathbf c]\,\phi)(Y')=\int_{Y'}\sZ_{\nu;t}\big(\mathbf c(y)\big)\,\phi(dy)
\]
wherever it converges: a normalized Zamma kernel read against the input over the cut region.
\end{definition}

The Trailmoment has five slots: its index and configuration, which are the kernel's, its coefficient function, its input and its cut region. On a line, $Y=\mathbb R$, the cut region $(N,\infty)$ is the \emph{cut} at $N$, and the region $(N,N']$ between two cuts is a \emph{head}; on $\mathbb R^d$ the region $X_N=\{N<y_1<\cdots<y_d\}$ is a \emph{trail region}, the nested tails that give the Trailmoment its name; the trails of Section~\ref{sec:trails} are read over it. The input is written as a density with its domain marked by an indicator, so that the measure is part of the data.

% Reference 20: def:trailmoment; source line 759 (the corners of the line)
\begin{definition}[Lattice inputs; the index-one value]\label{def:lattice}
The \emph{lattice input} on the coset $b+q\mathbb Z$, $b\in\mathbb R$, of step $q>0$ and frequency $a\in\mathbb C$ is $e^{2\pi iax}\,\Ilat{b+q\mathbb Z}=\sum_{m\in b+q\mathbb Z}e^{2\pi iam}\,\delta_m$; at frequency $0$ it is counting measure on $b+q\mathbb Z$, written $\Ilat{b+q\mathbb Z}$, and for $\im a>0$ the values $e^{2\pi iam}$ decay geometrically toward $+\infty$. The input $f\,\Ilat{b+q\mathbb Z}$ has the density $f$ on the lattice, and $\IZ$ is the \emph{counting input}. At one unit exponent and the coefficient function $x$ the kernel is $\sZ_{\nu;(1)}(x)=x^{-\nu}$, so at a cut beyond which every lattice point is positive
\[
(\sT_{\nu;(1)}[x]\,f\,\Ilat{b+\mathbb Z})((N,\infty))=\sum_{m\in b+\mathbb Z,\ m>N}f(m)\,m^{-\nu}.
\]
At index one the lattice input is regularized at its upper end:
\[
\Tx{1}{\Ilat{b+\mathbb Z}}{N}:=\lim_{M\to\infty}\Big(\Th{1}{\Ilat{b+\mathbb Z}}{N}{N+M}-\log M\Big),
\]
which exists, the heads being $\log M+\Tx{1}{\Ilat{b+\mathbb Z}}{N}+O(M^{-1})$: the regularized value is the coefficient of $M^0\log^0M$ in the heads.
\end{definition}

% Reference 20: prop:traillaws(i); source line 781
\begin{proposition}[The cut relation]\label{prop:cutlaw}
The Trailmoment is additive in its cut region: for disjoint measurable $Y'_1,Y'_2\subseteq Y$, $(\sT_{\nu;t}[\mathbf c]\,\phi)(Y'_1\sqcup Y'_2)=(\sT_{\nu;t}[\mathbf c]\,\phi)(Y'_1)+(\sT_{\nu;t}[\mathbf c]\,\phi)(Y'_2)$, either side converging when the other does; on a line, for $N<N'$,
\[
(\sT_{\nu;t}[\mathbf c]\,\phi)((N,\infty))=(\sT_{\nu;t}[\mathbf c]\,\phi)((N,N'])+(\sT_{\nu;t}[\mathbf c]\,\phi)((N',\infty)),
\]
and for a lattice input at index one it holds with the regularized values.
\end{proposition}

% Reference 20: def:inputs; source line 11728
\begin{definition}[Inputs]\label{def:inputs}
An input is a nonnegative $\sigma$-finite measure $\mu$, or such a measure with a complex density, $\phi=f\mu$ with $|\phi|=|f|\mu$; a Trailmoment converges when $\int|\sZ_{\nu;t}(\mathbf c)|\,d|\phi|$ over the cut region is finite. The image $\tau_*\phi$ of an input with a complex density under a measurable map $\tau$ is used only in integrals $\int g\,d(\tau_*\phi):=\int g\circ\tau\,d\phi$ with $\int|g\circ\tau|\,d|\phi|<\infty$; $|\tau_*\phi|$ denotes $\tau_*|\phi|$, and tameness and convergence of $\tau_*\phi$ are meant with it. The \emph{Lebesgue input} $\IR$ is Lebesgue measure on $\mathbb R$, $f\,\IR$ is the input with density $f$, $\Ilat{D}$ is Lebesgue measure on an interval $D$, and $\Ilat{\mathbb R^{\cE}}$ is Lebesgue measure on $\mathbb R^{\cE}$; the lattice inputs are those of Definition~\ref{def:lattice}, and at integer frequency the lattice input on $\mathbb Z$ is the counting input.

An input is \emph{tame beyond the cut $N$} when its part beyond $N$ is carried by $[x_0,\infty)$ for some $x_0>0$ and $\int_{(N,\infty)}x^{-\beta}\,|\phi|(dx)<\infty$ for some real $\beta$. Its Laplace transform beyond the cut is
\[
\mathcal F_N(u)=\int_{(N,\infty)}e^{-ux}\,\phi(dx),
\]
which converges for $\re u>0$, with $|\mathcal F_N(u)|\le e^{-(u-1)x_0}\int_{(N,\infty)}e^{-x}\,|\phi|(dx)$ for $u\ge1$. When $\mathcal F_N$ is moreover meromorphic at $u=0$, $\Tx{\nu}{\phi}{N}$ denotes, for every $\nu$ off its poles, the continuation in which $\Gamma(\nu)\,\Tx{\nu}{\phi}{N}$ is the meromorphic continuation of $\int_0^\infty u^{\nu-1}\mathcal F_N(u)\,du$, the Laplace side of Theorem~\ref{thm:laplaceside} below at one unit exponent. At index one on a lattice input it is the regularized value of Definition~\ref{def:lattice}, which is the constant term of the continuation at $\nu=1$.
\end{definition}

% Reference 20: prop:traillaws(vi); source line 798, and the first step of the proof of thm:twosidedtrail, source line 930
\begin{proposition}[The continued Trailmoment on a line]\label{prop:continuedline}
Let $\phi$ be tame beyond the cut $N$ with $\mathcal F_N(u)=\sum_{k\ge-k_0}c_ku^k$ near $u=0$. Then $\Gamma(\nu)\,\Tx{\nu}{\phi}{N}$ is meromorphic in $\nu$ with at most simple poles, at $\nu=-k$, $k\ge-k_0$, with residue $c_k$, and at the nonpositive integers
\[
\Tx{-j}{\phi}{N}=(-1)^jj!\,c_j\qquad(j\in\mathbb N).
\]
\end{proposition}

At a nonpositive integer index the Trailmoment of a tame input is a Taylor coefficient of its Laplace transform, so a cut that releases a head changes these values by the head read at the same index.

\section{Slots, corners and heights}

% Reference 20: def:height; source line 11722
\begin{definition}[Slots, corners, height]\label{def:height}
A \emph{slot} is a place of a symbol of the serif family where one of its data stands. A Gramma or a Zamma has three slots: its index, its configuration and its coefficients. A Breta has five: its partition into bunches, its shape index, its scale index, its terms (shape exponents and levels) and its coefficients. A Trailmoment has five: its index $\nu$, its configuration $t$, its coefficient function $\mathbf c$, its input $\phi$ and its cut region. A \emph{corner} of a member of the serif family is that member with some of its slots given particular values. A \emph{relation} is an identity between members of the serif family. Its \emph{height} is the set of values of the slots on which one proof establishes it; a slot is \emph{free} in a relation when the height does not restrict it and \emph{pinned} otherwise. For a pinned slot, the \emph{mechanism} is the step of the proof that uses the pinned value.
\end{definition}

A relation recited at a corner is understood once its height is known: which slots the proof leaves free, which it uses, and at which step. A relation whose height is every value of every slot is a property of the operator. A pin records what a proof uses; it does not say that the relation fails at the other values of the slot, and where the converse is proved, so that the relation is a property of the pinned value, it is stated with the relation.

\section{Relations at free slots}

% Reference 20: thm:gaugecoef; source line 11736 (one unit exponent in the last sentence).
% Proved in the last part: thm:gaugecoef.
\begin{theorem}[The gauge of the coefficient function]\label{thm:gaugecoef}
Let $(\nu,t,\mathbf c,\phi,Y')$ be the data of a Trailmoment on a source $Y$ and let $\ell:Y\to(0,\infty)^{\cE}$ be measurable. Then
\[
\big(\sT_{\nu;t}[(\ell^{t_k}c_k)_k]\,\phi\big)(Y')=\big(\sT_{\nu;t}[\mathbf c]\,\ell^{-\nu}\phi\big)(Y'),
\]
either side converging when the other does. If the columns $t_k$, $k\in S$, form a basis of $\mathbb R^{\cE}$, the choice $\log\ell=-(t_S^{\mathsf T})^{-1}\log\mathbf c_S$ makes $\ell^{t_k}c_k=1$ for every $k\in S$: the gauge moves $|\cE|$ of the $N$ coefficient functions into the input, and the kernel then depends on the source point only through the $N-|\cE|$ functions $\ell^{t_k}c_k$, $k\notin S$. At the unit configuration, $(\sT_{\nu;I_{\cE}}[\mathbf c]\,\phi)(Y')=\int_{Y'}\mathbf c^{-\nu}\,d\phi$.
\end{theorem}

At the unit configuration the kernel is a density on the input: index, coefficient function and input enter only through the measure $\mathbf c^{-\nu}\phi$. With $N>|\cE|$ terms the kernel depends on the source point through $N-|\cE|$ gauge-invariant functions as well as through a density, and this count is the dimension of the share side (Theorem~\ref{thm:sharetrail}).

% Reference 20: thm:exchangetrail; source line 11752 (one unit exponent in (iii)).
% Proved in the last part: thm:exchangetrail.
\begin{theorem}[Exchange and ladder of the Trailmoment]\label{thm:exchangetrail}
Each identity holds wherever every Trailmoment in it converges.
\begin{enumerate}[label=(\roman*)]
\item \emph{Exchange, one relation per scale.} For each $e\in \cE$,
\[
(\sT_{\nu;t}[\mathbf c]\,\phi)(Y')=\sum_{k=1}^Nt_{k,e}\ \Gst11{\nu_e+t_{k,e}}{\nu_e+1}\prod_{f\neq e}\Gst11{\nu_f+t_{k,f}}{\nu_f}\ \big(\sT_{\nu+t_k;t}[\mathbf c]\,c_k\phi\big)(Y').
\]
\item \emph{Ladder.} If $\mathbf c=\mathbf c_\zeta$ depends differentiably on a real parameter $\zeta$ and the differentiated integrand is dominated by an integrable function locally uniformly in $\zeta$,
\[
\partial_\zeta(\sT_{\nu;t}[\mathbf c_\zeta]\,\phi)(Y')=-\sum_{k=1}^N\prod_{e}\Gst11{\nu_e+t_{k,e}}{\nu_e}\ \big(\sT_{\nu+t_k;t}[\mathbf c_\zeta]\,(\partial_\zeta c_k)\,\phi\big)(Y').
\]
\item \emph{The unit configuration.} At $t=I_{\cE}$ the exchange relation at each scale has one term, and on the interior of the set of indices where the Trailmoment converges the index is differentiated through the input: for every $e\in\cE$,
\[
(\sT_{\nu;I_{\cE}}[\mathbf c]\,\phi)(Y')=(\sT_{\nu+\mathbf e_e;I_{\cE}}[\mathbf c]\,c_e\phi)(Y'),\qquad\partial_{\nu_e}(\sT_{\nu;I_{\cE}}[\mathbf c]\,\phi)(Y')=-(\sT_{\nu;I_{\cE}}[\mathbf c]\,(\log c_e)\,\phi)(Y').
\]
\end{enumerate}
\end{theorem}

At the unit configuration a step of the index at a scale is a factor of that coordinate of the coefficient function on the input, in both directions. With several terms on a scale the exchange relation is linear in $N+1$ Trailmoments whose inputs differ by the factors $c_k$, one relation per scale, and a step of the index is no longer a density.

% Reference 20: thm:transportfree; source line 11775
\begin{theorem}[Transport]\label{thm:transportfree}
For a measurable map $\tau:Y\to Y_2$, an input $\phi$ on $Y$, a coefficient function $\mathbf c$ on $Y_2$ and a cut region $Y_2'\subseteq Y_2$,
\[
(\sT_{\nu;t}[\mathbf c]\,\tau_*\phi)(Y_2')=(\sT_{\nu;t}[\mathbf c\circ\tau]\,\phi)(\tau^{-1}Y_2'),
\]
whenever the right side converges; for a nonnegative input either side converges when the other does. In particular, transporting by the coefficient function itself, a convergent Trailmoment is
\[
(\sT_{\nu;t}[\mathbf c]\,\phi)(Y')=\big(\sT_{\nu;t}[\mathrm{id}]\,\mathbf c_*(\phi|_{Y'})\big)\big((0,\infty)^N\big):
\]
coefficient function, input and cut region enter only through the image $\mathbf c_*(\phi|_{Y'})$ of the input on the orthant of coefficients.
\end{theorem}

A Trailmoment of a nonnegative input is therefore a function of three things: an index, a configuration, and a measure on the orthant of coefficients. Two descriptions of that measure, by different coefficient functions, inputs and cut regions, give the same member of the serif family. An increasing bijection of the line carries half-lines to half-lines: with $\tau(x)=x+h$ it is the \emph{shift} of the cut, and with $\tau(x)=qx$ at one unit exponent it gives $(\sT_{\nu;(1)}[x]\,\tau_*\phi)((qN,\infty))=q^{-\nu}\Tx{\nu}{\phi}{N}$.

% Reference 20: thm:sharetrail; source line 11794
\begin{theorem}[The share side of the Trailmoment]\label{thm:sharetrail}
Let $\re\nu\in\interior\cone(t)$, and let $\mathbf s\in(0,\infty)^N$ be an admissible real point, $t\mathbf s=\re\nu$, with $\int_{Y'}\prod_kc_k(y)^{-s_k}\,|\phi|(dy)<\infty$. On the contour through $\mathbf s$,
\begin{gather*}
\prod_{e\in \cE}\Gamma(\nu_e)\ (\sT_{\nu;t}[\mathbf c]\,\phi)(Y')=\int_{t\mathbf b=\nu}\ \prod_{k=1}^N\Gamma(b_k)\ (\sT_{\mathbf b;I_N}[\mathbf c]\,\phi)(Y'),\\
(\sT_{\mathbf b;I_N}[\mathbf c]\,\phi)(Y')=\int_{Y'}\prod_{k=1}^Nc_k(y)^{-b_k}\,\phi(dy),
\end{gather*}
a fibre integral of dimension $N-|\cE|$ of Trailmoments at the unit configuration $I_N$, with one scale per term of $t$. Where some $\nu_e$ lies in $-\mathbb N$ the left side is read as $\int_{Y'}\sG_{\nu;t}(\mathbf c)\,d\phi$.
\end{theorem}

Every configuration is written through the unit configuration: a Trailmoment at any configuration is a superposition of Trailmoments at the unit configuration over the fibre of its configuration, each multiplied by the Gammas of its shares. Closing the fibre integral (Theorem~\ref{thm:closing}) replaces the superposition by residues, and the residues are Trailmoments at the unit configuration at shifted indices; that is why their values appear in the closings of every configuration.

% Reference 20: thm:lettertrail; source line 11808
\begin{theorem}[The Laplace side of the Trailmoment]\label{thm:laplaceside}
The joint Laplace transform of the data $(\mathbf c,\phi,Y')$ is $\mathcal F(v)=\int_{Y'}e^{-\ip v{\mathbf c(y)}}\,\phi(dy)$ for $v\in(0,\infty)^N$. For $\nu\in\cT_t$, wherever the double integral converges,
\[
\prod_{e\in \cE}\Gamma(\nu_e)\ (\sT_{\nu;t}[\mathbf c]\,\phi)(Y')=\int_{\mathbb R^{\cE}_{>0}}x^{\nu-\one}\,\mathcal F\big(x^{t_1},\dots,x^{t_N}\big)\,dx=\big(\sT_{\one-\nu;I_{\cE}}[x]\ \mathcal F(x^{t})\,\Ilat{\mathbb R^{\cE}}\big)\big(\mathbb R^{\cE}_{>0}\big),
\]
where $\mathcal F(x^t)=\mathcal F(x^{t_1},\dots,x^{t_N})$ and the kernel of the last Trailmoment is $\prod_ex_e^{\nu_e-1}$, the Zamma of the unit configuration at the index $\one-\nu$; where some $\nu_e$ lies in $-\mathbb N$ the left side is read as $\int_{Y'}\sG_{\nu;t}(\mathbf c)\,d\phi$. The exponential is the Zamma of a term at exponent zero, $\sZ_{1;(1,0)}(1,\xi)=e^{-\xi}$, so the joint Laplace transform is a Trailmoment of the same input: $\mathcal F(v)=(\sT_{1;(1,0)}[(1,\ip v{\mathbf c})]\,\phi)(Y')$.
\end{theorem}

Index $\nu$ on the input side is index $\one-\nu$ on the Laplace side. Besides the input itself, a Trailmoment has two sides: the share side, which resolves the configuration into unit configurations, and the Laplace side, which replaces the input by its Laplace transform; the reflections of the index (Theorems~\ref{thm:reflectpoles} and~\ref{thm:reflectinvert}), the Shruffle (Theorem~\ref{thm:shrufflefree}) and the Cremona (Theorem~\ref{thm:cremonafree}) are read on the Laplace side, and the lattice relations and the Struffle on the input itself.

\section{Lattices}

% Reference 20: thm:latticefree; source line 11823
\begin{theorem}[The lattice relations]\label{thm:latticefree}
Let $b\in\mathbb R$, let $q\ge1$ be an integer, let $(\nu,t,\mathbf c)$ be the data of a kernel on $\mathbb R$, $f$ a density and $N$ a cut.
\begin{enumerate}[label=(\roman*)]
\item \emph{Cosets.} Wherever the right side converges,
\[
\sum_{r=0}^{q-1}\big(\sT_{\nu;t}[\mathbf c]\,f\,\Ilat{\frac{b+r}q+\mathbb Z}\big)((N,\infty))=\big(\sT_{\nu;t}[\mathbf c(\,\cdot\,/q)]\,f(\,\cdot\,/q)\,\Ilat{b+\mathbb Z}\big)((qN,\infty)).
\]
\item \emph{Frequencies.} For $a\in\mathbb R$, wherever the Trailmoments on the left converge,
\[
\sum_{r=0}^{q-1}\big(\sT_{\nu;t}[\mathbf c]\,e^{2\pi i(a+\frac rq)(x-b)}f\,\Ilat{b+\mathbb Z}\big)((N,\infty))=q\,\big(\sT_{\nu;t}[\mathbf c]\,e^{2\pi ia(x-b)}f\,\Ilat{b+q\mathbb Z}\big)((N,\infty)),
\]
the right side summing over the lattice of step $q$.
\end{enumerate}
\end{theorem}

A lattice is assembled from a sublattice in two ways: by cosets, the $q$ cosets of $\mathbb Z$ in $q^{-1}\mathbb Z$; by frequencies, the $q$ exponentials $e^{2\pi irx/q}$, $0\le r<q$. Both statements are free in the index, the configuration, the coefficient function, the density, the coset and the cut. At one unit exponent with the coefficient function $x$ the rescaled coefficient function of (i) multiplies the kernel by the scalar $q^{\nu}$, $\sZ_{\nu;(1)}(x/q)=q^{\nu}x^{-\nu}$; at index one, with the regularized values, the same assembly carries the further term $q\log q$, the finite part read in two coordinates (Lemma~\ref{lem:Regscale}).

\section{The Gamma read by Trailmoments}

% Reference 20: prop:headstails; source line 848 (with thm:trailcorners(c), source line 805)
\begin{proposition}[Heads and higher orders at lattice data]\label{prop:headstails}
For real $b>-1$ and $|z|<1+b$,
\[
\log\Gamma(1+b-z)=\log\Gamma(1+b)+\sum_{i\ge1}\frac{z^i}{i}\,\Tx{i}{\Ilat{b+\mathbb Z}}{b},
\]
the index-one Trailmoment regularized; equivalently, for $j\in\mathbb N$,
\[
(\log\Gamma)^{(j+1)}(1+b)=(-1)^{j+1}j!\,\Tx{j+1}{\Ilat{b+\mathbb Z}}{b}.
\]
Further $\Gamma(-k+z)=\dfrac{(-1)^k\,\Gamma(1+z)\,\Gamma(1-z)}{z\,\Gamma(k+1-z)}$ for $k\in\mathbb N$. In this statement write $h=\prod_a\Gamma(1+b_a-z_a)\big/\prod_c\Gamma(1+b'_c-z'_c)$, a stacked Gamma whose constants $b_a,b'_c$ lie in one translate $b_0+\mathbb N$ with $b_0>-1$, with $\sum_az_a=\sum_cz'_c$, $|z_a|<1+b_a$ and $|z'_c|<1+b'_c$. Then
\begin{multline*}
\log\frac{h}{h|_{z=0}}=\sum_cz'_c\,\Th{1}{\Ilat{b_0+\mathbb Z}}{b_0}{b'_c}\\
-\sum_az_a\,\Th{1}{\Ilat{b_0+\mathbb Z}}{b_0}{b_a}\\
+\sum_{i\ge2}\frac1i\Big(\sum_az_a^{\,i}\,\Tx{i}{\Ilat{b_0+\mathbb Z}}{b_a}\\
-\sum_cz_c'^{\,i}\,\Tx{i}{\Ilat{b_0+\mathbb Z}}{b'_c}\Big):
\end{multline*}
the regularized index-one values cancel, the first order is a finite sum over a head of the lattice, rational when $b_0$ is, and the orders $i\ge2$ are Trailmoments of the lattice input cut at the data. At integer data, $b_0=0$, the first order is rational and the higher orders are Trailmoments of the counting input.
\end{proposition}

% Reference 20: prop:resonance; source line 824 (the residue of $\Gamma(b)^2h(b)$).
% Proved in the last part: prop:resonance.
\begin{proposition}[Laurent data of the Gamma at its poles]\label{prop:resonance}
Let $j\in\mathbb N$. For $|\varepsilon|<1$, with $(-1)^j\varepsilon\,\Gamma(-j+\varepsilon)=\Gamma(1+\varepsilon)\,\Gamma(1-\varepsilon)/\Gamma(1+j-\varepsilon)$,
\begin{multline*}
\log\big((-1)^j\varepsilon\,\Gamma(-j+\varepsilon)\big)=-\log j!-\Tx{1}{\IZ}{j}\,\varepsilon\\
-\sum_{k\ge2}\frac{\Tx{k}{\IZ}{j}}k\,\varepsilon^k+\sum_{k\ge1}\frac{\Tx{2k}{\IZ}{0}}k\,\varepsilon^{2k},
\end{multline*}
and $(\log\Gamma)'(j+1)=-\Tx{1}{\IZ}{j}$. Let $h$ be holomorphic near $b=-j$, and let $k_1,\dots,k_p$ be integers with $j_i=j-k_i\in\mathbb N$, so that every $\Gamma(b+k_i)$ has a pole at $-j$. Then
\[
\operatorname*{Res}_{b=-j}\Big(h(b)\prod_{i=1}^p\Gamma(b+k_i)\Big)=(-1)^{j_1+\cdots+j_p}\sum_{l=0}^{p-1}\frac{h^{(l)}(-j)}{l!}\,[\varepsilon^{p-1-l}]\prod_{i=1}^p\big((-1)^{j_i}\varepsilon\,\Gamma(-j_i+\varepsilon)\big),
\]
each factor of the product being the exponential of the display at $j_i$; at $p=2$ and $k_1=k_2=0$ it is $\frac1{(j!)^2}\big(h'(-j)+2\,(\log\Gamma)'(j+1)\,h(-j)\big)$. Every Laurent coefficient of $\prod_i\Gamma(b+k_i)$ at $-j$ is $\prod_i(j_i!)^{-1}$ times a polynomial in Trailmoments of the counting input cut at the $j_i$ and at $0$.
\end{proposition}

% Reference 20: prop:grade; source line 864
\begin{proposition}[Selecting an index inside a word]\label{prop:grade}
Let $w\ge2$ be an integer and $h(z)=\prod_i\Gamma(1+b_i-a_iz)^{n_i}$, a finite product with real $b_i>-1$, $a_i\in\mathbb C$ and $n_i\in\mathbb Z$. Its $w$-th multiplicative difference satisfies
\[
\prod_{j=0}^wh(jz)^{(-1)^{w-j}\binom wj}=1+(w-1)!\,\Big[\sum_in_ia_i^w\,\Tx{w}{\Ilat{b_i+\mathbb Z}}{b_i}\Big]z^w+O(z^{w+1}).
\]
At integer values $b_i=k_i\in\mathbb N$ with rational $a_i$, an integer $\ell\ge1$ with every $\ell a_i\in\mathbb Z$ and $k_\ast=\max_ik_i$, the coefficient of $z^w$ is $A\,\Tx{w}{\IZ}{0}-B$ with $A=(w-1)!\sum_in_ia_i^w$ and $B=(w-1)!\sum_in_ia_i^w\,\Th{w}{\IZ}{0}{k_i}$, both cleared by $\ell^wd_{k_\ast}^{\,w}$. For the one factor $\Gamma(k+1-z)$, the coefficient of $z^w$ in $\prod_{j=0}^w\Gamma(k+1-jz)^{(-1)^{w-j}\binom wj}$ is $(w-1)!\,\Tx{w}{\IZ}{k}$, and passing from $k$ to $k+1$ multiplies this product by $\prod_{j=0}^w(k+1-jz)^{(-1)^{w-j}\binom wj}$. The difference acts multiplicatively on one word, not on the sums of a closing; along a ray from the origin whose law of $[1]$ is reflexive of order two, the coefficient of the coordinate is read without any sum, from the closings of the ray (Chapter~\ref{ch:rayarith}).
\end{proposition}

The $w$-th difference annihilates every power of $z$ below the $w$-th in the logarithm of a Gamma word, so it removes the regularized index-one values together with all Trailmoments of lower index and keeps exactly those of index $w$: from a word in the Gammas of the data it selects one Trailmoment index, with a rational head.

% Reference 20: lem:resonantpair; source line 875
\begin{lemma}[Two Gammas and their reflections]\label{lem:resonantpair}
The stacked Gamma $\Gst40{a,\ a,\ 1-a,\ 1-a}{}=\pi^2/\sin^2\pi a$ has at every integer $m$ the principal part $(a-m)^{-2}$ and no simple part, it equals $\sum_{m\in\mathbb Z}(a+m)^{-2}$, and on the line $a=\frac12+iy$ it equals $4\pi^2/(e^{\pi y}+e^{-\pi y})^2$.
\begin{enumerate}[label=(\roman*)]
\item \emph{Right closing.} For $0<c<1$ and $g$ rational, holomorphic on $\{\re a\ge c\}$, with $g(a)=O(|a|^{-2})$: $\displaystyle\frac1{2\pi i}\int_{(c)}\Gst40{a,\ a,\ 1-a,\ 1-a}{}\,g(a)\,da=-\sum_{m\ge1}g'(m)$.
\item \emph{Left closing.} For $0<X<1$: $\displaystyle\frac1{2\pi i}\int_{(1/2)}\Gst40{a,\ a,\ 1-a,\ 1-a}{}\,X^{-a}\,da=\frac{\log X}{X-1}$.
\item \emph{Fourier read.} For real $t$,
\[
\frac1{2\pi}\int_{\mathbb R}\Gst40{\frac12+iy,\ \frac12+iy,\ \frac12-iy,\ \frac12-iy}{}\,e^{ity}\,dy=\Gst20{1+\frac{it}{2\pi},\ 1-\frac{it}{2\pi}}{}=\frac{t}{e^{t/2}-e^{-t/2}}.
\]
\end{enumerate}
\end{lemma}

A closing whose poles coincide (Theorem~\ref{thm:closing}) takes Laurent coefficients of its Gammas, and a balanced word near lattice data has heads and higher orders (Proposition~\ref{prop:headstails}): in both cases the Gammas contribute words in Trailmoments of lattice inputs, read at the cuts that the data set. The stacked Gamma of Lemma~\ref{lem:resonantpair} is the simplest coincidence, two Gammas and their two reflections meeting at every integer: its double poles have no simple parts, so its right closing reads the derivative of the rest of the integrand at the positive integers, and on the line $\re a=\frac12$ it is $2\pi$ times the density of the critical-line law (Chapter~\ref{ch:share}).

% Reference 20: thm:hparity; source line 12025
\begin{theorem}[Reflected pairs and matched parity]\label{thm:hparity}
\begin{enumerate}[label=(\roman*)]
\item \emph{Cosets.} For $0<b<1$ and every integer $n\ge1$,
\[
\Tx{n}{\Ilat{b+\mathbb Z}}{b-1}+(-1)^n\,\Tx{n}{\Ilat{-b+\mathbb Z}}{-b}=\frac{(-1)^{n-1}}{(n-1)!}\,\partial_b^{\,n-1}\big(\pi\cot\pi b\big),
\]
the index-one values regularized. The left side is $\sum_{m\in b+\mathbb Z}m^{-n}$ over the whole lattice, in symmetric order at $n=1$, and the right side is $\pi^n$ times a polynomial of degree $n$ in $\cot\pi b$ with rational coefficients.
\item \emph{Frequencies.} For $0<b\le1$ and every integer $n\ge2$, the Trailmoment $\Tx{n}{\sin\!\big(2\pi bx-\tfrac{\pi(n-1)}2\big)\,\IZ}{0}$ is $\frac{(-1)^{n-1}(2\pi)^n}{2}$ times the coefficient of $u^{n-1}$ in $e^{-bu}/(1-e^{-u})$. The density is $(-1)^{n/2}\cos(2\pi bx)$ at even $n$ and $(-1)^{(n-1)/2}\sin(2\pi bx)$ at odd $n$, and the value is $\pi^n$ times a polynomial of degree $n$ in $b$ with rational coefficients.
\item \emph{Periodic densities.} For integers $0<r<q$ and $n\ge1$,
\[
\Tx{n}{\big(\Ilat{r+q\mathbb Z}+(-1)^n\,\Ilat{-r+q\mathbb Z}\big)}{0}=\frac{(-1)^{n-1}}{q^{n}(n-1)!}\,\partial_b^{\,n-1}\big(\pi\cot\pi b\big)\Big|_{b=r/q},
\]
at $n=1$ with the points of the two lattices taken in increasing order.
\item \emph{Folding.} At $b=\frac12$ the two lattices of (i) are one, and (i) reads
\[
\big(1+(-1)^n\big)\Tx{n}{\Ilat{\frac12+\mathbb Z}}{0}=\frac{(-1)^{n-1}}{(n-1)!}\,\partial_b^{\,n-1}(\pi\cot\pi b)\Big|_{b=1/2}:
\]
at odd $n$ both sides vanish, and at even $n=2j$, with the lattice relation at the step $2$ (Theorem~\ref{thm:latticefree}(i)),
\[
\Tx{2j}{\IZ}{0}=\frac{-1}{2\,(2^{2j}-1)\,(2j-1)!}\ \partial_b^{\,2j-1}\big(\pi\cot\pi b\big)\Big|_{b=1/2}.
\]
At $b=1$, (ii) gives $\Tx{2j}{\IZ}{0}=(-1)^{j+1}\frac{(2\pi)^{2j}}2$ times the coefficient of $u^{2j-1}$ in $1/(e^u-1)$, and the identity $0=0$ at odd $n$.
\end{enumerate}
\end{theorem}

At an integer index $n$ the kernel $\sZ_{n;(1)}(x)=x^{-n}$ has the parity of $n$ under $x\mapsto-x$. What closes, for every coset $b+\mathbb Z$ and every integer index, is the value of the whole coset, points of both signs: the reflection of the Gamma evaluates it, and at $n=2$ and $0<b<1$ it is $\Gst40{b,\ b,\ 1-b,\ 1-b}{}$ (Lemma~\ref{lem:resonantpair}). A Trailmoment reads the points beyond its cut only, so this relation gives it when the input, continued to the whole lattice, has the parity of the index: an even periodic density closes at even indices and an odd one at odd indices, by (iii), and on the coset $q\mathbb Z$ by (iv). On a lattice equal to its own reflection the pair of (i) folds onto one Trailmoment with the factor $1+(-1)^n$; at even $n$ that Trailmoment closes, and at odd $n$ the relation is the identity $0=0$ and carries no information. Statement (ii) is the same rule read on frequencies. The mismatched parities, the sine at an even index and the cosine at an odd index, are reached by neither statement on any coset or at any frequency; the counting input at an odd index is the cosine at frequency zero, and its Trailmoments are the ones that the positive words of Part~\ref{part:arith} carry.

\section{Reflections of the index}

% Reference 20: thm:reflectpoles; source line 11875
\begin{theorem}[Reflection through the poles of the Laplace transform]\label{thm:reflectpoles}
Let $\mathcal F$ be a function on $(0,\infty)$ that extends to a meromorphic function on $\mathbb C$ with no pole on $(0,\infty)$, such that
\begin{enumerate}[label=(\alph*)]
\item its poles other than $u=0$ are simple, with residue $r_p$ at the pole $p$, and its pole at $0$ has order $k_0\ge0$, with $k_0=0$ when there is none;
\item there are radii $R_m\to\infty$ and a real $A$ with $\sup_{|u|=R_m}|\mathcal F(u)|=O(R_m^{A})$.
\end{enumerate}
Write $\Sigma_m(\nu)=\sum_{0<|p|<R_m}r_p\,(-p)^{\nu-1}$, with $(-p)^{\nu-1}=\exp\big((\nu-1)\log(-p)\big)$ and $|\arg(-p)|<\pi$.
\begin{enumerate}[label=(\roman*)]
\item \emph{On the strip.} If moreover $\mathcal F(u)=O(u^{A})$ as $u\to+\infty$, then for $k_0<\re\nu<-A$ with $\nu\notin\mathbb Z$,
\[
\int_0^\infty u^{\nu-1}\mathcal F(u)\,du=\frac{\pi}{\sin\pi\nu}\,\lim_{m\to\infty}\Sigma_m(\nu).
\]
\item \emph{Beyond the strip.} If $\mathcal F=\mathcal F_N$ is the Laplace transform of an input $\phi$ tame beyond the cut $N$, its poles lie in $\re u\le0$, and for $\re\nu<-A$ with $\nu\notin\{1,2,3,\dots\}$,
\[
\Tx{\nu}{\phi}{N}=\Gamma(1-\nu)\,\lim_{m\to\infty}\Sigma_m(\nu).
\]
\end{enumerate}
\end{theorem}

Reflection of the index replaces an input by the poles of its Laplace transform: each simple pole $p$ becomes the point measure $r_p\,\delta_{-p}$, read by the kernel $\sZ_{1-\nu;(1)}$. The relation is free in the input within (a) and (b), and pinned at one unit exponent with the coefficient function $x$; the mechanism is the Laplace side of the kernel $\sZ_{\nu;(1)}$, where the index is the exponent of the scale and a full turn of the scale multiplies the integrand by $e^{2\pi i\nu}$. For the input $e^{-x}\,\IR$ cut at $0$, whose Laplace transform $1/(1+u)$ has the one pole $-1$, statement (i) is the reflection of the Gamma (Proposition~\ref{prop:stacked}(v)); for a lattice input the strip of (i) is empty, and the reflection is a statement about the continued Trailmoment, which is (ii).

% Reference 20: thm:reflectinvert; source line 11949
\begin{theorem}[Reflection by inversion of the Laplace transform]\label{thm:reflectinvert}
Let $\phi$ and $\psi$ be tame beyond the cut $0$, with Laplace transforms $\mathcal F$ and $\mathcal G$, and suppose that for some $k\in\mathbb R$, $\lambda>0$, $C\neq0$ and $f_0,g_0\in\mathbb C$,
\[
\mathcal F(u)+f_0=C\,u^{-k}\Big(\mathcal G\big(\tfrac\lambda u\big)+g_0\Big)\qquad(u>0).
\]
Then $\Gamma(\nu)\,\Tx{\nu}{\phi}{0}$, defined for large $\re\nu$ by the Laplace side, continues meromorphically to $\mathbb C$,
\begin{multline*}
\Gamma(\nu)\,\Tx{\nu}{\phi}{0}=\int_{\sqrt\lambda}^\infty u^{\nu-1}\mathcal F(u)\,du+C\lambda^{\nu-k}\int_{\sqrt\lambda}^\infty u^{k-\nu-1}\mathcal G(u)\,du\\
+\frac{C\,g_0\,\lambda^{(\nu-k)/2}}{\nu-k}-\frac{f_0\,\lambda^{\nu/2}}{\nu},
\end{multline*}
whose only possible poles are simple, at $\nu=k$ and $\nu=0$. The same holds with $\psi$ in the place of $\phi$, and
\[
\Gamma(\nu)\,\Tx{\nu}{\phi}{0}=C\,\lambda^{\nu-k}\,\Gamma(k-\nu)\,\Tx{k-\nu}{\psi}{0}.
\]
\end{theorem}

The inversion $u\mapsto\lambda/u$ of the scale is the reflection of the Gramma (Proposition~\ref{prop:grammareflect}) on the Laplace side. The relation is free in the two inputs, which need not be lattices and whose Laplace transforms need not be meromorphic; it uses one property of the pair, that the inversion of the scale carries the Laplace transform of one to that of the other. The constants $f_0$ and $g_0$ are the multiples of $\delta_0$ that the two inputs would carry at the origin, where the kernel $\sZ_{\nu;(1)}$ is not defined, and for $k\neq0$ they give the residues of the continuation, $-f_0$ at $\nu=0$ and $Cg_0$ at $\nu=k$.

% Reference 20: thm:massreflect; source line 12271
\begin{theorem}[Reflection at the coefficient function $x+\mu$]\label{thm:massreflect}
Let $\phi$, $\psi$, $\mathcal F$, $\mathcal G$, $k$, $\lambda$, $C$, $f_0$ and $g_0$ be as in Theorem~\ref{thm:reflectinvert}, and let $\mu>0$. For $\re s>\max(k,0)$ with $\int_{(0,\infty)}(x+\mu)^{-\re s}\,|\phi|(dx)<\infty$,
\[
\Gamma(s)\Big[(\sT_{s;(1)}[x+\mu]\,\phi)((0,\infty))+f_0\,\mu^{-s}\Big]=C\,\lambda^{s-k}\int_{(0,\infty)}\sG_{k-s;(1,-1)}(y,\lambda\mu)\,\psi(dy)+C\,g_0\,\Gamma(s-k)\,\mu^{k-s}.
\]
The integral on the right converges for every $s\in\mathbb C$ and is entire in $s$.
\end{theorem}

\section{Joint sources}

% Reference 20: thm:jointclosure; source line 953
\begin{theorem}[Closure with the source kept joint]\label{thm:jointclosure}
\begin{enumerate}[label=(\roman*)]
\item $\sZ_{\nu;t}(\mathbf c)\,\sZ_{\eta;s}(\mathbf d)=\sZ_{(\nu,\eta);\,t\oplus s}(\mathbf c\oplus\mathbf d)$, the configuration $t\oplus s$ acting on the disjoint union of the two sets of scales.
\item A finite expression built from Zammas with positive indices by products and by nested reads against nonnegative measures is one Trailmoment on a source: its configuration is the concatenation $\oplus$ of the factors' configurations, and its input is the joint measure $\phi(dy)\,\psi_y(dw)$ of the source variables, every shared variable, conditional dependence and cut retained. The same holds for complex data where every integral converges absolutely.
\end{enumerate}
\end{theorem}

% Reference 20: prop:jointmoments; source line 967 (polynomials).
% Proved in the last part: prop:jointmoments.
\begin{proposition}[The joint Laplace transform of Dirichlet inputs]\label{prop:jointmoments}
Let $U=(U_Q)_{Q\in\cP}$ have independent blocks $U_Q\sim\Dir_{\mvec|_Q}$ on a partition $\cP$ of the scales, let $P_1,\dots,P_q$ be finite sums $P_i(u)=\sum_{\mathbf a}[u^{\mathbf a}]P_i\,u^{\mathbf a}$ over exponents $\mathbf a\in[0,\infty)^{\cE}$, polynomials when the exponents are integers, and let $t$ be a configuration of $q$ columns on $d$ scales and $\nu\in\mathbb C^d$. Then $\mathcal F(v)=\oE_{U\sim\bigotimes_Q\Dir_{\mvec|_Q}}\big[\exp\big(-\sum_iv_iP_i(U)\big)\big]$ is entire and
\begin{gather*}
\mathcal F(v)=\sum_{\mathbf k\in\mathbb N^q}\frac{(-v)^{\mathbf k}}{\mathbf k!}M(\mathbf k),\\
M(\mathbf k)=\sum_{\mathbf a}\big[u^{\mathbf a}\big]\big(P_1(u)^{k_1}\cdots P_q(u)^{k_q}\big)\prod_e\Gst11{m_e+a_e}{m_e}\prod_Q\Gst11{m_Q}{m_Q+a_Q},\qquad a_Q=\sum_{e\in Q}a_e ;
\end{gather*}
for integer exponents the stacked Gammas are $\prod_e(m_e)_{a_e}/\prod_Q(m_Q)_{a_Q}$. If every $P_i$ is positive on the interior of the source and $\re\nu_j>0$, then, wherever the read converges,
\[
\oE_{U\sim\bigotimes_Q\Dir_{\mvec|_Q}}\big[\sZ_{\nu;t}\big(P_1(U),\dots,P_q(U)\big)\big]=\frac1{\prod_j\Gamma(\nu_j)}\int_{(0,\infty)^d}x^{\nu-\one}\,\mathcal F\big(x^{t_1},\dots,x^{t_q}\big)\,dx ,
\]
the Laplace side of the Trailmoment $(\sT_{\nu;t}[(P_i)_i]\,\otimes_Q\Dir_{\mvec|_Q})\big(\prod_Q\Delta_Q\big)$. For real coefficients, two choices of the data with equal joint moment arrays $M$ give the vectors $(P_i(U))_i$ equal laws, hence equal reads of every kernel.
\end{proposition}

\section{Trails}\label{sec:trails}

% Reference 20: def:trail; source line 889
\begin{definition}[Trail data]\label{def:trail}
Take $d$ scales, an index $\boldsymbol\nu\in\mathbb C^d$, the \emph{trail columns} $I_d=(\mathbf e_1,\dots,\mathbf e_d)$, and \emph{constant columns} $(t_k)_{k\in C}$ with constant coefficients $\lambda_k>0$ (a constant column of coefficient zero is absent). Put
\[
t=I_d\oplus(t_k)_{k\in C},\qquad
\mathbf c(y)=(y_1,\dots,y_d,(\lambda_k)_{k\in C}),\qquad
X_N=\{N<y_1<\cdots<y_d\}.
\]
The \emph{trail} at the cut $N$ of an input $\phi$ on $\mathbb R^d$, carried where every trail coefficient is positive, is the Trailmoment of Definition~\ref{def:jointtrail} with this coefficient function and cut region,
\[
(\sT_{\boldsymbol\nu;\,I_d\oplus(t_k)_{k\in C}}[(y_1,\dots,y_d,(\lambda_k)_{k\in C})]\,\phi)(X_N);
\]
with no constant columns it is $(\sT_{\boldsymbol\nu;I_d}[y]\,\phi)(X_N)$. For lattices $\mathbb L_1,\dots,\mathbb L_d\subset\mathbb R$ the input $\Ilat{\mathbb L_1\times\cdots\times\mathbb L_d}$ is counting measure on their product, of which the cut region selects the points in $X_N$. The input is \emph{given by increments} $\psi_1,\dots,\psi_d$, measures on $(0,\infty)$, when it is the image $\phi_N$ of $\bigotimes_j\psi_j$ under $s\mapsto y$, $y_i=N+s_1+\cdots+s_i$; their Laplace transforms are $L_j(U)=\int e^{-sU}\psi_j(ds)$. By transport (Theorem~\ref{thm:transportfree}) the trail of such an input is the Trailmoment on the increment source $(0,\infty)^d$ with the input $\bigotimes_j\psi_j$ and the coefficient function $s\mapsto(N+s_1,\dots,N+s_1+\cdots+s_d,(\lambda_k)_{k\in C})$, over the whole source: there the cut $N$ is carried by the coefficient function.
\end{definition}

% Reference 20: thm:trailreading; source line 897
\begin{theorem}[The trail by its increments]\label{thm:trailreading}
Let the input be given by increments, put $t=I_d\oplus(t_k)_{k\in C}$, and let $U_i=x_i+\cdots+x_d$ be the sums of the scales from the $i$-th on, $U_{d+1}=0$. Wherever the right side converges absolutely, with $x_i=U_i-U_{i+1}$,
\begin{multline*}
(\sT_{\boldsymbol\nu;t}[(y,(\lambda_k)_{k\in C})]\,\phi_N)(X_N)\\
=\int_{U_1>\cdots>U_d>0}\prod_{i=1}^d\frac{(U_i-U_{i+1})^{\nu_i-1}}{\Gamma(\nu_i)}\,
\exp\Big(-\sum_{k\in C}\lambda_kx^{t_k}\Big)e^{-NU_1}\prod_{j=1}^dL_j(U_j)\,dU .
\end{multline*}
Its relations hold wherever the reads in them converge, the derivatives under the integral being dominated:
\begin{enumerate}[label=(\roman*)]
\item \emph{Moving the cut}, with the increments fixed:
\[
\partial_N(\sT_{\boldsymbol\nu;t}[(y,(\lambda_k)_{k\in C})]\,\phi_N)(X_N)
=-\sum_i\nu_i\,(\sT_{\boldsymbol\nu+\mathbf e_i;t}[(y,(\lambda_k)_{k\in C})]\,\phi_N)(X_N).
\]
\item \emph{Ladder.}
\[
\partial_{\lambda_k}(\sT_{\boldsymbol\nu;t}[(y,(\lambda_h)_{h\in C})]\,\phi_N)(X_N)
=-\Gst dd{\boldsymbol\nu+t_k}{\boldsymbol\nu}\,
(\sT_{\boldsymbol\nu+t_k;t}[(y,(\lambda_h)_{h\in C})]\,\phi_N)(X_N).
\]
\item \emph{Exchange, per scale.}
\begin{multline*}
(\sT_{\boldsymbol\nu;t}[(y,(\lambda_k)_{k\in C})]\,\phi_N)(X_N)=(\sT_{\boldsymbol\nu+\mathbf e_i;t}[(y,(\lambda_k)_{k\in C})]\,y_i\phi_N)(X_N)\\
+\frac1{\nu_i}\sum_{k\in C}t_{k,i}\lambda_k\Gst dd{\boldsymbol\nu+t_k}{\boldsymbol\nu}\,
(\sT_{\boldsymbol\nu+t_k;t}[(y,(\lambda_h)_{h\in C})]\,\phi_N)(X_N).
\end{multline*}
\item \emph{Composition.} If every constant column is supported on one scale, partition $C$ into sets $C_i$ with $\operatorname{supp}(t_k)\subseteq\{i\}$ for $k\in C_i$, assigning every zero column to one scale, and put $t^{(i)}=(1)\oplus(t_{k,i})_{k\in C_i}$. Then
\[
\sZ_{\boldsymbol\nu;t}(y,(\lambda_k)_{k\in C})
=\prod_i\sZ_{\nu_i;t^{(i)}}(y_i,(\lambda_k)_{k\in C_i}),
\]
and the trail is the iterated read over $s_1,\dots,s_d>0$, in that order, of $\prod_i\sZ_{\nu_i;t^{(i)}}(N+s_1+\cdots+s_i,(\lambda_k)_{k\in C_i})$ against $\bigotimes_i\psi_i$: each one-scale Trailmoment has the next read in its input. With counting increments, every $\psi_j$ the counting input on the positive integers, no constant columns and $N\in b+\mathbb Z$, this is the nest of one-scale Trailmoments at successive running cuts, with indices $\nu_1,\dots,\nu_d$ and lattice input $\Ilat{b+\mathbb Z}$, and the trail is $\Ttr{\boldsymbol\nu}{d}{\Ilat{(b+\mathbb Z)^d}}{N}$.
\end{enumerate}
\end{theorem}

\section{Products of Trailmoments}

% Reference 20: thm:strufflefree; source line 12084
\begin{theorem}[The Struffle]\label{thm:strufflefree}
\begin{enumerate}[label=(\roman*)]
\item Let $\phi,\psi$ be inputs on $\mathbb R$ with any indices, configurations and coefficient functions. If the product below converges with $|\phi|$, $|\psi|$ and the moduli of the kernels, then
\begin{align*}
&(\sT_{\nu;t}[\mathbf c]\,\phi)((N,\infty))\,(\sT_{\nu';t'}[\mathbf c']\,\psi)((N,\infty))\\
&\quad=\Big(\sT_{\nu;t}[\mathbf c]\,\big[x\mapsto(\sT_{\nu';t'}[\mathbf c']\,\psi)((x,\infty))\big]\phi\Big)((N,\infty))\\
&\qquad+\Big(\sT_{\nu';t'}[\mathbf c']\,\big[x\mapsto(\sT_{\nu;t}[\mathbf c]\,\phi)((x,\infty))\big]\psi\Big)((N,\infty))\\
&\qquad+\big(\sT_{(\nu,\nu');\,t\oplus t'}[\mathbf c\oplus\mathbf c']\,\phi\wedge\psi\big)((N,\infty)),
\end{align*}
where the \emph{diagonal input} is $\phi\wedge\psi=\sum_x\phi(\{x\})\psi(\{x\})\delta_x$, the sum over the points that both inputs charge, and the diagonal kernel is the Zamma of the concatenation $t\oplus t'$ (Theorem~\ref{thm:jointclosure}(i)).
\item For inputs $\phi_1,\dots,\phi_r$ on $\mathbb R$ with kernel data $(\nu_i,t_i,\mathbf c_i)$, if $\prod_i(\sT_{\nu_i;t_i}[\mathbf c_i]\,\phi_i)((N,\infty))$ converges with the total variations and the moduli of the kernels, it is the sum over the ordered partitions $(B_1,\dots,B_s)$ of $\{1,\dots,r\}$ into nonempty blocks of
\[
\int_{N<x_1<\cdots<x_s}\prod_{j=1}^s\Big(\prod_{i\in B_j}\sZ_{\nu_i;t_i}(\mathbf c_i(x_j))\Big)
\Big(\bigwedge_{i\in B_j}\phi_i\Big)(dx_j),
\]
with $\bigwedge_{i\in B}\phi_i=\sum_x\prod_{i\in B}\phi_i(\{x\})\delta_x$ for a block of at least two elements and $\phi_i$ for the block $\{i\}$.
\end{enumerate}
\end{theorem}

The Struffle orders the points of the inputs along the line; the diagonal is there exactly when the inputs charge common points beyond the cut. It is free in the index, the configuration and the coefficient function of both factors and in their inputs on the line; the cut region is pinned at one common cut. At one unit exponent with one coefficient function on both factors the diagonal kernel is $\sZ_{\nu+\nu';(1)}$.

% Reference 20: thm:shrufflefree; source line 12104 (two strings in (i) and (iii)).
% Proved in the last part: thm:shrufflefree.
\begin{theorem}[The Shruffle on a line]\label{thm:shrufflefree}
For a string $w=(w_1,\dots,w_n)$ of measurable functions on $(0,\infty)$ put $\operatorname{chain}(w)=\int_{x_1>\cdots>x_n>0}\prod_iw_i(x_i)\,dx$, defined when $\operatorname{chain}(|w|)<\infty$.
\begin{enumerate}[label=(\roman*)]
\item If $\operatorname{chain}(|w^{(1)}|),\dots,\operatorname{chain}(|w^{(r)}|)$ are finite, $\prod_{i=1}^r\operatorname{chain}(w^{(i)})=\sum_\varsigma\operatorname{chain}(\varsigma)$, the sum over the $(n_1+\cdots+n_r)!/(n_1!\cdots n_r!)$ interleavings of the $r$ strings, $n_i$ the length of $w^{(i)}$.
\item Let a trail have its input given by nonnegative increments, with Laplace transforms $L_1,\dots,L_d$, and a cut $N\in\mathbb R$ such that the image of the first increment $\psi_1$ under $s\mapsto s+N$ is carried by $(0,\infty)$, as it is when $N\ge0$; let it have no constant columns and positive integer indices $a_1,\dots,a_d$. Then, both sides finite or infinite together, with $1^{a}$ for $a$ copies of the constant function $1$, the trail is the chain integral of its string:
\[
(\sT_{\mathbf a;I_d}[y]\,\phi_N)(X_N)=\operatorname{chain}\big(e^{-N\,\cdot}L_1,\ 1^{a_1-1},\ L_2,\ 1^{a_2-1},\ \dots,\ L_d,\ 1^{a_d-1}\big).
\]
\item The product of finitely many such trails is the sum, over the interleavings of their strings, of trails of the same kind at the cut $0$: every interleaving begins with a Laplace transform and is the string of the trail whose increments have the Laplace transforms of the interleaving in their order, each index being one more than the number of constant functions $1$ that follow its Laplace transform, $e^{-N\,\cdot}L$ being the Laplace transform of the increment moved by $N$.
\end{enumerate}
\end{theorem}

The Shruffle orders the scales on the Laplace side. It is free in the nonnegative increments, in the number of scales, and in the cut as long as the first point beyond it is positive; it is pinned at positive integer indices without constant columns, where $(U_i-U_{i+1})^{a-1}/\Gamma(a)$ is the volume of a chain of points. Where a nest of Trailmoments on the input side is a trail given by increments all equal to one measure at the cut $0$, as for the counting input cut at $0$, the Struffle and the Shruffle expand one product into the same members, and comparing the two expansions gives linear relations among the trails.

% New: proved in the last part.
\begin{theorem}[Trails at frequencies are iterated integrals]\label{thm:ct-trailchain}
Let $d\ge1$, $\boldsymbol\nu\in\{1,2,\dots\}^d$ and $\mathbf a\in\mathbb R^d$, put $\zeta_i=e^{2\pi ia_i}$, and let $\zeta_d\neq1$ when $\nu_d=1$. Put $c_i=\zeta_i\zeta_{i+1}\cdots\zeta_d$, $w=|\boldsymbol\nu|$, and let $(\ell_1,\dots,\ell_w)$ be the word $c_1^{-1},0^{\nu_1-1},c_2^{-1},0^{\nu_2-1},\dots,c_d^{-1},0^{\nu_d-1}$, $0^r$ standing for $r$ letters $0$. Then the heads of the trail over $\{0<y_1<\cdots<y_d\le N\}$ of the lattice input $e^{2\pi i\ip{\mathbf a}{y}}\,\Ilat{\mathbb Z^d}$ at the indices $\boldsymbol\nu$ converge as $N\to\infty$, to the trail itself when $\nu_d\ge2$, and with the trail read as this limit
\[
\big(\sT_{\boldsymbol\nu;I_d}[y]\ e^{2\pi i\ip{\mathbf a}{y}}\,\Ilat{\mathbb Z^d}\big)\big(\{0<y_1<\cdots<y_d\}\big)=(-1)^d\int_{0<t_1<\cdots<t_w<1}\ \prod_{j=1}^w\frac{dt_j}{t_j-\ell_j},
\]
an absolutely convergent integral. For $|z|<1$ the trail with the factor $z^{y_d}$ in its density is $(-1)^d$ times the same integral over $0<t_1<\cdots<t_w<z$ along the segment from $0$ to $z$.
\end{theorem}

On the Laplace side the input of the trail, given by its increments, has the letters $c_iz/(1-c_iz)$ in the variable $z=e^{-U}$, singular at the $c_i^{-1}$, and the index of each scale adds the letters $0$ of the factor $1^{\nu_i-1}$ of the Shruffle (Theorem~\ref{thm:shrufflefree}(ii)): the theorem is the chain integral of that string, read in $z$, for inputs with complex densities.

% Reference 20: thm:splitsum; source line 12123
\begin{theorem}[The sum over the splittings of an index]\label{thm:splitsum}
Let $b\in\mathbb R$, let $N$ be a cut, and let $m_0>0$ be the first point of $b+\mathbb Z$ beyond $N$. For an integer $w\ge3$,
\begin{gather*}
\sum_{a=1}^{w-2}\Ttr{(a,w-a)}{2}{\Ilat{(b+\mathbb Z)^2}}{N}=\Tx{w-1}{\eta\,\Ilat{b+\mathbb Z}}{N},\\
\eta(m)=\Tx{1}{\IZ}{m-m_0}-\Tx{1}{\Ilat{b+\mathbb Z}}{m}=\sum_{k\ge1}\Big(\frac1{m-m_0+k}-\frac1{m+k}\Big),
\end{gather*}
the index-one values regularized. At every point $m$ of the lattice beyond the cut,
\[
\eta(m)-\frac1m=(\log\Gamma)'(m)-(\log\Gamma)'(m-m_0+1),
\]
which is positive for $m_0>1$, zero for $m_0=1$ and negative for $m_0<1$. So the sum over the splittings is $\Tx{w}{\Ilat{b+\mathbb Z}}{N}$ if and only if $m_0=1$, that is, for the counting input at a cut in $[0,1)$; it is larger for $m_0>1$ and smaller for $m_0<1$. For the counting input cut at an integer $c\ge0$, $m_0=c+1$ and $\eta(m)=\Th{1}{\IZ}{m-c-1}{m}$, a head of $c+1$ terms. For a lattice $b+q\mathbb Z$ of step $q>0$ with a positive first point $m_0$ beyond the cut, the sum over $1\le a\le w-2$ of the trails with the indices $(a,w-a)$ on its points is $\Tx{w}{\Ilat{b+q\mathbb Z}}{N}$ if and only if $m_0=q$.
\end{theorem}

Summing the trails on two scales over the splittings of an index telescopes in the inner variable, and what is left is one Trailmoment of index one less, with the density $\eta$. The index rises back by one exactly when $\eta(m)=\frac1m$ at every point, which is exactly when the first point beyond the cut is $1$.

% Reference 20: thm:cremonafree; source line 12141
\begin{theorem}[The Cremona of an increment]\label{thm:cremonafree}
Let $L>0$ be continuous on $(0,\infty)$ with $\int_x^\infty L<\infty$ for $x>0$ and $\int_0^\infty L=\infty$, and let $\operatorname{tail}_L(x)=\int_x^\infty L(u)\,du$, the \emph{tail function}, a decreasing bijection of $(0,\infty)$. For a string $w$ whose entries are $L(x)\,dx$ or $dx$, let $\operatorname{chain}_L(w)$ be its chain integral over $x_1>\cdots>x_n>0$, largest variable first, and $w^\ast$ the reversed string with $L(x)\,dx$ and $dx$ swapped. Then $\operatorname{tail}_L\circ\operatorname{tail}_L=\mathrm{id}$ if and only if $L(\operatorname{tail}_L(x))\,L(x)=1$ for every $x$, and in that case $\operatorname{chain}_L(w)=\operatorname{chain}_L(w^\ast)$ for every string, finite or infinite together.
\end{theorem}

The Cremona uses one property of the increment: its tail function $\operatorname{tail}_L$ is an involution of the scale, so that the substitution $y=\operatorname{tail}_L(x)$ carries $dx$ to $L(y)\,dy$ as well as $L(x)\,dx$ to $dy$. For the counting increments, $L(u)=1/(e^u-1)$ and the involution is the curve $e^{-x}+e^{-y}=1$.

\section{Convolutions of inputs}

% Reference 20: def:convolutions; source line 12164
\begin{definition}[The two convolutions]\label{def:convolutions}
Let $\phi$ and $\psi$ be inputs carried by $[0,\infty)$, and $\phi\otimes\psi$ their product on $[0,\infty)^2$. The \emph{additive convolution} $\phi\boxplus\psi$ is the image of $\phi\otimes\psi$ under $(x,y)\mapsto x+y$. When both inputs are carried by $(0,\infty)$, the \emph{multiplicative convolution} $\phi\boxtimes\psi$ is the image of $\phi\otimes\psi$ under $(x,y)\mapsto xy$. Both are images in the sense of Definition~\ref{def:inputs}: for inputs with complex densities they are used only in integrals that converge absolutely on the product $\phi\otimes\psi$. The multiplicative convolution of two $\sigma$-finite measures need not be $\sigma$-finite; it is used only where an integral of Theorem~\ref{thm:twoconv}(ii) below converges absolutely, and there it is $\sigma$-finite, its kernel being positive and finite on $(0,\infty)$. For lattice inputs with densities, $(f\,\IZ)\boxplus(g\,\IZ)$ and $(f\,\IZ)\boxtimes(g\,\IZ)$ denote the convolutions of the parts beyond the cut $0$: for densities $f$ and $g$ on the positive integers, the first has the density $\sum_{j+k=n}f(j)g(k)$ at $n$ and the second the density $\sum_{jk=n}f(j)g(k)$; likewise for lattice inputs on other cosets. For an input carried by $[0,\infty)$ write $\mathcal F_\phi(u)=\int_{[0,\infty)}e^{-ux}\,\phi(dx)$, the kernel being $1$ at the point $x=0$: it is the Laplace transform beyond every cut $N<0$, and for an input carried by $(0,\infty)$ the Laplace transform beyond the cut $0$.
\end{definition}

% Reference 20: thm:twoconv; source line 12168 (one scale in (ii) and (iii)).
% Proved in the last part: thm:twoconv.
\begin{theorem}[Two convolutions]\label{thm:twoconv}
Let $\phi$ on $Y_1$ and $\psi$ on $Y_2$ be inputs with coefficient functions $c$ and $c'$ and cut regions $Y_1'$ and $Y_2'$, and write $c+c'$ and $c\,c'$ for the functions $(y_1,y_2)\mapsto c(y_1)+c'(y_2)$ and $(y_1,y_2)\mapsto c(y_1)\,c'(y_2)$ on $Y_1\times Y_2$.
\begin{enumerate}[label=(\roman*)]
\item \emph{The kernel $\sZ_{1;(1,0)}(1,\cdot)$ on a sum.} If the two Trailmoments on the right converge,
\[
\big(\sT_{1;(1,0)}[(1,c+c')]\ \phi\otimes\psi\big)(Y_1'\times Y_2')=(\sT_{1;(1,0)}[(1,c)]\,\phi)(Y_1')\ (\sT_{1;(1,0)}[(1,c')]\,\psi)(Y_2').
\]
For inputs carried by $[0,\infty)$ and $u\in\mathbb C$ such that $e^{-x\re u}$ is integrable for $|\phi|$ and for $|\psi|$, $\mathcal F_{\phi\boxplus\psi}(u)=\mathcal F_\phi(u)\,\mathcal F_\psi(u)$.
\item \emph{The kernel $\sZ_{\nu;I_{\cE}}$ on a product.} For coefficient functions $\mathbf c$ and $\mathbf c'$ with values in $(0,\infty)^{\cE}$, $\mathbf c\,\mathbf c'$ their coordinatewise product and $\nu\in\mathbb C^{\cE}$, if the two Trailmoments on the right converge,
\[
\big(\sT_{\nu;I_{\cE}}[\mathbf c\,\mathbf c']\ \phi\otimes\psi\big)(Y_1'\times Y_2')=(\sT_{\nu;I_{\cE}}[\mathbf c]\,\phi)(Y_1')\ (\sT_{\nu;I_{\cE}}[\mathbf c']\,\psi)(Y_2').
\]
For inputs carried by $(0,\infty)$,
\begin{gather*}
\Tx{\nu}{\phi\boxtimes\psi}{0}=\Tx{\nu}{\phi}{0}\ \Tx{\nu}{\psi}{0},\\
\mathcal F_{\phi\boxtimes\psi}(u)=\int_{(0,\infty)}\mathcal F_\psi(ux)\,\phi(dx),
\end{gather*}
the first when its right side converges and the second when $\int\!\!\int e^{-xy\re u}\,|\phi|(dx)\,|\psi|(dy)<\infty$.
\item \emph{Converses.} Let $g$ be a continuous complex-valued function on $(0,\infty)^{\cE}$, not identically zero, with $g(s+s')=g(s)\,g(s')$ for all $s,s'\in(0,\infty)^{\cE}$. Then $g(s)=e^{-\ip us}$ for one $u\in\mathbb C^{\cE}$. Let $h$ be a continuous complex-valued function on $(0,\infty)^{\cE}$, not identically zero, with $h(x\,x')=h(x)\,h(x')$ for the coordinatewise product. Then $h=\sZ_{\nu;I_{\cE}}$ for one $\nu\in\mathbb C^{\cE}$. The two hypotheses are the product relations of (i) and of (ii), for a kernel $g$ or $h$ in the place of the kernel there, on every pair of Dirac measures.
\end{enumerate}
\end{theorem}

The two convolutions add and multiply the points of two inputs, and each has one kernel that turns it into a product: on sums the kernel $\sZ_{1;(1,0)}(1,\cdot)$ of a term at exponent zero beside a unit term, which is $e^{-c}$; on products the kernel $\sZ_{\nu;I_{\cE}}$ of the unit configuration, the power $\mathbf c\mapsto\mathbf c^{-\nu}$. Both statements are free in the two inputs, the two coefficient functions and the two cut regions, and each is pinned in its kernel; (iii) is the converse of that pin.


\chapter{Rational words and live inputs}\label{ch:transport}

A rational word acting on a Trailmoment through its input is moved, by partial fractions and the shift of the cut, onto Trailmoments of shifted inputs at the unit configuration; the pole columns of the word then separate what is a Trailmoment from what is a finite difference. An input is \emph{live} when the values of its density are themselves read from the serif family, a Breta or a Gramma at a shifting slot, and the relations of that member give the input a recurrence of finite order. Every identity for rational words and live inputs is a finite identity first, at finite cuts; its infinite form uses the displayed convergence, or the common finite upper cut of the closing it comes from.

% Reference 20: def:ratword; source line 1051
\begin{definition}[Rational words and their transport]\label{def:ratword}
A \emph{rational word} is a function
\[
R(x)=\sum_{k=0}^{k_{\max}}\sum_{j=1}^{J_k}c_{k,j}(x+k)^{-j},
\]
kept with its partial-fraction presentation. For the data $(\nu,t,\mathbf c,\phi)$ of a Trailmoment on a line and the cut region $(N,\infty)$, the rational word acts through the input,
\[
(\sT_{\nu;t}[\mathbf c]\,(R\phi))((N,\infty))=\int_{(N,\infty)}R(x)\,\sZ_{\nu;t}(\mathbf c(x))\,\phi(dx).
\]
Write $(\operatorname{Sh}^hf)(m)=f(m+h)$, $\Delta H(m)=H(m+1)-H(m)$, $\tau_k(x)=x+k$ and $\phi_k=(\tau_k)_*\phi$. On the counting input at index zero, where $\sZ_{0;(1)}(x)=1$, this is $\Tx{0}{Rf\,\IZ}{N}=\sum_{m>N}R(m)f(m)$ for an integer cut $N\ge0$.
\end{definition}

% Reference 20: thm:shiftedinputs; source line 1063
\begin{theorem}[Shifted inputs]\label{thm:shiftedinputs}
With $g_k(y)=\sZ_{\nu;t}(\mathbf c(y-k))f(y-k)$,
\[
(\sT_{\nu;t}[\mathbf c]\,(Rf\phi))((N,\infty))=\sum_{k,j}c_{k,j}\,(\sT_{j;(1)}[x]\,(g_k\phi_k))((N+k,\infty)).
\]
On the counting input at integer cuts $M>N\ge0$,
\[
\Th{0}{Rf\,\IZ}{N}{M}=\sum_{k,j}c_{k,j}\,\Th{j}{(\operatorname{Sh}^{-k}f)\,\IZ}{N+k}{M+k},
\]
and when the Trailmoments converge individually
\[
\Tx{0}{Rf\,\IZ}{N}=\sum_{k,j}c_{k,j}\,\Tx{j}{(\operatorname{Sh}^{-k}f)\,\IZ}{N+k}.
\]
When only the combination converges, the common finite upper cut is kept and the limit is taken in the combination. For $f=\one$, $\sum_kc_{k,1}=0$ and an integer $m_0\ge1$, with the index-one Trailmoments regularized,
\begin{align*}
&\Tx{0}{R\,\IZ}{m_0-1}=\sum_{k,s}c_{k,s}\,\Tx{s}{\IZ}{m_0-1+k}\\
&\qquad=\sum_{s\ge2}\sum_kc_{k,s}\,\Tx{s}{\IZ}{m_0-1+k}-\sum_kc_{k,1}\,\Th{1}{\IZ}{0}{m_0-1+k}\\
&\qquad=\int_0^\infty\Big[\sum_{k,s}\frac{c_{k,s}}{\Gamma(s)}u^{s-1}e^{-ku}\Big]\frac{e^{-(m_0-1)u}}{e^u-1}\,du .
\end{align*}
\end{theorem}

% Reference 20: thm:commoncut; source line 1087
\begin{theorem}[Common cut and heads]\label{thm:commoncut}
For a common cut $N'\ge N+k_{\max}$,
\begin{multline*}
\Tx{j}{(\operatorname{Sh}^{-k}f)\IZ}{N+k}=\Tx{j}{(\operatorname{Sh}^{-k}f)\IZ}{N'}\\
+\Th{j}{(\operatorname{Sh}^{-k}f)\IZ}{N+k}{N'},
\end{multline*}
the head being $\sum_{N+k<\ell\le N'}\ell^{-j}f(\ell-k)$ (Proposition~\ref{prop:cutlaw}). For $R=-g'$ with $g(a)=\sum_{h,s}c_{h,s}(a+h)^{-s}=O(a^{-2})$ and $C_s=\sum_hc_{h,s}$, under individual convergence,
\begin{multline*}
\Tx{0}{(-g')f\,\IZ}{N}=\sum_ss\,C_s\,\Tx{s+1}{f\,\IZ}{N'}\\
+\sum_{h,s}s\,c_{h,s}\,\Tx{s+1}{(\operatorname{Sh}^{-h}f-f)\,\IZ}{N'}\\
+\sum_{h,s}s\,c_{h,s}\,\Th{s+1}{(\operatorname{Sh}^{-h}f)\,\IZ}{N+h}{N'}.
\end{multline*}
\end{theorem}

% Reference 20: thm:columns; source line 1102 (rational coefficients and $R=O(t^{-2})$).
% Proved in the last part: thm:columns.
\begin{theorem}[Pole columns are a rational difference quotient]\label{thm:columns}
Let $R(t)=\sum_{k=0}^{k_{\max}}\sum_{s=1}^wc_{k,s}(t+k)^{-s}$ with coefficients in a subfield $\mathbb K$ of $\mathbb C$, $C_s=\sum_kc_{k,s}$, and $H(t)=\sum_{k,s}c_{k,s}\sum_{j=0}^{k-1}(t+j)^{-s}$, a rational function with coefficients in $\mathbb K$. Then
\[
R(t)=\sum_{s=1}^wC_st^{-s}+\Delta H(t),
\]
and $R(t)=O(t^{-2})$ exactly when $C_1=0$. Conversely a difference $\Delta H$ of a proper rational function with poles on the nonpositive integer lattice has every column sum zero; modulo such differences, the rational words of maximum order $w$ have the basis $t^{-1},\dots,t^{-w}$, and those that are $O(t^{-2})$ the basis $t^{-2},\dots,t^{-w}$. If the poles fall in several cosets of $\mathbb Z$, each coset has its own columns.
\end{theorem}

% Reference 20: thm:telescoper; source line 1113 ($C_1=0$ in (iv)).
% Proved in the last part: thm:telescoper.
\begin{theorem}[Summation by parts with a live input]\label{thm:telescoper}
Let $H$ be a function on the integers above $N$ and $f$ a density on them.
\begin{enumerate}[label=(\roman*)]
\item $\Delta(Hf)(m)=\Delta H(m)f(m)+H(m+1)\Delta f(m)$, so at finite cuts $M>N$
\[
\sum_{N<m\le M}\Delta H(m)f(m)=H(M+1)f(M+1)-H(N+1)f(N+1)-\sum_{N<m\le M}H(m+1)\Delta f(m).
\]
\item If the upper boundary vanishes and the series converges, $\sum_{m>N}\Delta H(m)f(m)=-H(N+1)f(N+1)-\sum_{m>N}H(m+1)\Delta f(m)$.
\item For a Trailmoment read as a function of its cut, $\Delta\,\Tx{p}{g\,\IZ}{m}=-(m+1)^{-p}g(m+1)$ (Proposition~\ref{prop:cutlaw}), and hence
\begin{multline*}
\sum_{m>N}\Delta H(m)\,\Tx{p}{g\,\IZ}{m}\\
=-H(N+1)\,\Tx{p}{g\,\IZ}{N+1}+\Tx{p}{Hg\,\IZ}{N+1},
\end{multline*}
with one fewer Trailmoment in the input.
\item With $R$, $C_s$, $H$ as in Theorem~\ref{thm:columns} and the index-one Trailmoment regularized (Definition~\ref{def:lattice}),
\begin{equation*}
\lim_{M\to\infty}\Big(\Th{0}{R\,\IZ}{N}{N+M}-C_1\log M\Big)=\sum_{s=1}^wC_s\,\Tx{s}{\IZ}{N}-H(N+1);
\end{equation*}
when $C_1=0$, that is when $R=O(t^{-2})$, the left side is $\Tx{0}{R\,\IZ}{N}$ and the sum starts at $s=2$. For a general density $f$, when the terms converge and the upper boundary vanishes,
\begin{multline*}
\Tx{0}{Rf\,\IZ}{N}=\sum_{s=1}^wC_s\,\Tx{s}{f\,\IZ}{N}\\
-H(N+1)f(N+1)-\sum_{m>N}H(m+1)\Delta f(m).
\end{multline*}
\end{enumerate}
\end{theorem}

% Reference 20: thm:descent; source line 1132
\begin{theorem}[Translation and the terminating descent]\label{thm:descent}
For the counting input, each inner Trailmoment below is read at its displayed running cut.
\begin{enumerate}[label=(\roman*)]
\item $\Tx{p}{g\,\IZ}{y-k}=\Tx{p}{g\,\IZ}{y}+\Th{p}{g\,\IZ}{y-k}{y}$, the head being $\sum_{a=0}^{k-1}(y-a)^{-p}g(y-a)$ (Proposition~\ref{prop:cutlaw}).
\item
\begin{multline*}
\Tx{j}{\big[x\mapsto\Tx{p}{g\,\IZ}{x-k}\big]\,\IZ}{N+k}\\
=\Tx{j}{\big[x\mapsto\Tx{p}{g\,\IZ}{x}\big]\,\IZ}{N+k}\\
+\sum_{a=0}^{k-1}\Tx{p}{(x+a)^{-j}g\,\IZ}{N+k-a}.
\end{multline*}
\item For positive integers $p,j$, each product $z^{-p}(z+a)^{-j}$ has finite rational partial fractions with poles at $0$ and $-a$ only. The recursion therefore terminates on every finite nest of Trailmoments of the counting input, each taking the next as a function of its cut: it returns finitely many such nests, finite shifted copies of $g$ and their heads, keeps the selected cuts, lowers at each correction the number of Trailmoments nested in the remaining input, and uses only rational operations on the pole data.
\item If $R(-L-t)=-R(t)$, then $\sum_kc_{k,j}=0$ for every even $j$, and for a density $f$ that is a nest of $d$ Trailmoments of the counting input, at a common cut $C$,
\[
\Tx{0}{Rf\,\IZ}{N}=\sum_j\Big(\sum_kc_{k,j}\Big)\Tx{j}{f\,\IZ}{C}+E_d ,
\]
with $E_d$ a sum of terms with at most $d$ nested Trailmoments, so the terms with $d+1$ nested Trailmoments carry only odd outer indices, starting at three when $\sum_kc_{k,1}=0$. For $f=\one$, an integer $m_0\ge1$ and $R(-k_{\max}-t)=-R(t)$, $c_{k_{\max}-k,s}=-(-1)^sc_{k,s}$, and at every even index $s$ the Trailmoments combine into the finite sum
\begin{multline*}
\sum_kc_{k,s}\Big[\Tx{s}{\IZ}{m_0-1+k}-\Tx{s}{\IZ}{m_0-1}\Big]\\
=-\sum_kc_{k,s}\,\Th{s}{\IZ}{m_0-1}{m_0-1+k}.
\end{multline*}
\end{enumerate}
\end{theorem}

% Reference 20: prop:transportmeasure; source line 1149
\begin{proposition}[Zamma kernels and general measures]\label{prop:transportmeasure}
For counting measure the product rule applies to the kernel times the density, $g(m)=\sZ_{\nu;t}(\mathbf c(m))f(m)$, with $\Delta g(m)=\sZ_{\nu;t}(\mathbf c(m+1))\Delta f(m)+f(m)\Delta[\sZ_{\nu;t}(\mathbf c(m))]$. For a general measure $\mu$ restricted to $(N,M]$, $\int\Delta H\,g\,d\mu_{N,M}=\int H\,d\big((\tau_1)_*(g\mu_{N,M})-g\mu_{N,M}\big)$.
\end{proposition}

\section{Live inputs}

% Reference 20: thm:bretainput; source line 1156 ($a,b>0$, $0<q_i<1$, real $s_i$).
% Proved in the last part: thm:bretainput.
\begin{theorem}[A Breta input with a recurrence of finite order]\label{thm:bretainput}
For $\re a,\re b>0$, real $q_i<1$ with $q_i\neq0$, and $s_i\in\mathbb C$, let
\[
f(m)=\int_0^1x^{m+a-1}(1-x)^{b-1}\prod_{i=1}^r(1-q_ix)^{-s_i}\,dx\qquad(m\ge0).
\]
For $\re s_i>0$ this is a Breta in which $m$ sits in the shape index: on the bunch $\{1,2\}$ with the shape $u=(x,1-x)$, the kernel terms $\rho u_1$ and $\rho u_2$ with unit coefficients and the letters $P_i(\rho,u)=\rho\big((1-q_i)u_1+u_2\big)$, Proposition~\ref{prop:letters} at the scale index $1+|\mathbf s|$ gives
\[
f(m)=\Gamma(1+|\mathbf s|)\,\sB^{\cP}_{((m+a,\,b),\,\mathbf s),\,(1+|\mathbf s|,\,\mathbf s);\,(\widetilde\pi,\widetilde\ell)}(\mathbf g\oplus\mathbf C).
\]
Then $f(m)=O(m^{-\re b})$, and with $Q=\prod_i(1-q_ix)$, $Q_i=Q/(1-q_ix)$ and
\[
B_m(x)=\big[(m+a)(1-x)-bx\big]Q(x)+x(1-x)\sum_i(s_i-1)q_iQ_i(x)=\sum_{\ell=0}^{r+1}B_\ell(m)x^\ell ,
\]
the density obeys $\sum_{\ell=0}^{r+1}B_\ell(m)f(m+\ell)=0$, with $B_0(m)=m+a$ and $B_{r+1}(m)=(-1)^{r+1}(\prod_iq_i)(m+a+b+r-\sum_is_i)$. If $a+b+r-\sum_is_i\notin-\mathbb N$, both extreme coefficients are nonzero for every $m\ge0$, and all shifts reduce to $(f(m),\dots,f(m+r))$ with coefficients rational in $m$.
\end{theorem}

% Reference 20: thm:fixedblock; source line 1167 (the data of thm:bretainput with $a,b>0$, $0<q_i<1$, real $s_i$).
% Proved in the last part: thm:fixedblock.
\begin{theorem}[A fixed set of input-dependent evaluations]\label{thm:fixedblock}
With $f$, $a$, $b$, $r$, $q_i$ and $s_i$ as in Theorem~\ref{thm:bretainput}, write in this statement $Q(x)=\prod_i(1-q_ix)$, $C(x)=(1-x)Q(x)$, $\eta(x)=x^{a-1}(1-x)^{b-1}\prod_i(1-q_ix)^{-s_i}$ and
\[
B(x)=\big[a(1-x)-bx\big]Q(x)+x(1-x)\sum_i(s_i-1)\,q_i\,\frac{Q(x)}{1-q_ix},
\]
so that $(xC\eta)'=B\eta$; let $B_\ell$, $C_\ell$ and $Q_\ell$ be the coefficients of $x^\ell$.
\begin{enumerate}[label=(\roman*)]
\item For each maximum index $A$, every shift reduces to $f(N+1),\dots,f(N+r+1)$ and the Trailmoments $(\sT_{p;(1)}[x+k]\,f\,\IZ)((N,\infty))=\sum_{m>N}f(m)(m+k)^{-p}$, $1\le p\le A$, $0\le k\le r$, together with $(\sT_{p;(1)}[x+a]\,f\,\IZ)((N,\infty))$ when $a$ is an integer at least $r+1$.
\item For $k\ge r+1$ and $p\ge2$,
\begin{multline*}
\sum_{\ell=0}^{r+1}(B_\ell-kC_\ell)\,(\sT_{p;(1)}[x+k-\ell]\,f\,\IZ)((N,\infty))+\sum_{\ell=0}^{r+1}C_\ell\,(\sT_{p-1;(1)}[x+k-\ell]\,f\,\IZ)((N,\infty))\\
=\sum_{\ell=0}^{r+1}\sum_{v=1}^{\ell}\Big[\frac{B_\ell-kC_\ell}{(N+v+k-\ell)^p}+\frac{C_\ell}{(N+v+k-\ell)^{p-1}}\Big]f(N+v),
\end{multline*}
and at $p=1$,
\begin{multline*}
\sum_{\ell=0}^{r+1}(B_\ell-kC_\ell)\,(\sT_{1;(1)}[x+k-\ell]\,f\,\IZ)((N,\infty))\\
=\sum_{\ell=0}^{r+1}\sum_{v=1}^{\ell}\frac{B_\ell-kC_\ell}{N+v+k-\ell}f(N+v)-\sum_{\ell=0}^{r}Q_\ell f(N+1+\ell).
\end{multline*}
\item Cuts move by $(\sT_{p;(1)}[x+k]\,f\,\IZ)((N+1,\infty))=(\sT_{p;(1)}[x+k]\,f\,\IZ)((N,\infty))-f(N+1)/(N+1+k)^p$ (Proposition~\ref{prop:cutlaw}).
\item For rational data, if $c_\ast$ clears the coefficients of $B$, $C$ and $Q$, then $d_{N+J}^{A}\prod_{r+1\le k\le J,\ k\neq a}|c_\ast(a-k)|^{A}$ clears all shifts through $J$ and indices through $A$.
\end{enumerate}
\end{theorem}

% New: proved in the last part.
\begin{theorem}[The row of a rational function against an input with a recurrence]\label{thm:ra-liverow}
Let $f$ be a density on the integers $m>N$, $N\in\mathbb Z$, with $\sum_{m>N}|f(m)|(1+|m|)^d<\infty$, that satisfies $\sum_{i=0}^\tau p_i(m)\,f(m+i)=0$ for every integer $m>N$, with $\tau\ge1$ and $p_i\in\mathbb Z[m]$ of degree at most $d$. Let $u\ge1$ be an integer and $h_\ast\in\frac1u\mathbb Z$ with $N+1+h_\ast>0$. In this statement write, for integers $s\ge1$, $r\ge0$ and $h\in h_\ast+\mathbb N$,
\begin{gather*}
T_s(h)=(\sT_{s;(1)}[x+h]\,f\,\IZ)((N,\infty))=\sum_{m>N}\frac{f(m)}{(m+h)^s},\\
M_r=(\sT_{-r;(1)}[x]\,f\,\IZ)((N,\infty))=\sum_{m>N}m^rf(m).
\end{gather*}
\begin{enumerate}[label=(\roman*)]
\item For a rational function $R(m)=P(m)+\sum_{h,s}c_{h,s}(m+h)^{-s}$, $P$ a polynomial of degree at most $d$ and finitely many $h\in h_\ast+\mathbb N$, $\sum_{m>N}R(m)f(m)=\sum_r([m^r]P)\,M_r+\sum_{h,s}c_{h,s}\,T_s(h)$: one coordinate for each index and shift.
\item Let $J\ge\tau$ with $p_0(-h_\ast-j)\neq0$ for $\tau\le j\le J$. Then for $s\ge1$ and $\tau\le j\le J$, $T_s(h_\ast+j)$ is a rational combination of the $T_{s'}(h_\ast+i)$, $1\le s'\le s$, $0\le i<\tau$, of the moments $M_r$, $0\le r\le d-1$, and of the values $f(N+1),\dots,f(N+\tau)$.
\item For every $w\ge1$ the coefficients of (ii) of all the $T_s(h_\ast+j)$, $1\le s\le w$, $\tau\le j\le J$, are cleared by
\[
d_{u(N+h_\ast+J)}^{\,w}\prod_{j=\tau}^{J}\big|u^dp_0(-h_\ast-j)\big|^w .
\]
\end{enumerate}
A rational function with poles in several cosets of $\mathbb Z$ takes one such reduction for each coset.
\end{theorem}

Theorem~\ref{thm:fixedblock} is the case of the Breta input of Theorem~\ref{thm:bretainput}, read through its differential equation; here the input is any density with a recurrence, the live inputs of the closings along a ray among them (Theorem~\ref{thm:ra-oneshare}). The clearing has two parts: the heads released by the moves of the cut, cleared by a power of $d_n$ on the scale $u$ of the coset, and the values of the first coefficient of the recurrence at the shifts used, raised to the power of the index, which along a ray is the frame of the recurrence (Proposition~\ref{prop:frame}).

% Reference 20: thm:grammainput; source line 1188
\begin{theorem}[A Gramma input with a recurrence of finite order]\label{thm:grammainput}
Let a one-scale configuration $t$ have positive integer exponents with greatest exponent $L$, the term at exponent $L$ present. For $\lambda,c>0$ put $g(m)=\sG_{\lambda+m;t}(\boldsymbol\xi)/\Gamma(c+m)$. Then
\[
(m+\lambda)g(m)=\sum_kt_k\xi_k(m+c)_{t_k}\,g(m+t_k),
\]
a shift relation of order $L$, the exchange relation of the Gramma (Theorem~\ref{thm:grammalaws}(ii)) divided by $\Gamma(c+m)$; for $L\ge2$, $\log g(m)\le-(1-1/L)m\log m+O(m)$, so $g$ decays faster than every inverse power.
\end{theorem}


\part{Tables and families}

\chapter{Tables}\label{ch:tables}

\section{Characters, tables, lists}

% Reference 20: def:table; source line 1208
\begin{definition}[Characters, tables, lists]\label{def:table}
The \emph{characters} are a finite set $\cE$, whose elements are the \emph{marks}; each mark $e$ carries a positive variable $\alpha_e$, and its name carries nothing. A \emph{table} on $\cE$ is a nonempty family $\cB$ of $r$-element subsets of $\cE$, the \emph{bases}, with the \emph{exchange property}: for bases $T,T'$ and $a\in T\setminus T'$ there is $b\in T'\setminus T$ with $(T\setminus\{a\})\cup\{b\}$ a basis. Each basis carries a weight $\theta_T>0$; $r$ is the \emph{rank} and $r^\ast=n-r$ the \emph{corank}. A \emph{list} is a table whose weights are all $1$; it is specified by its characters and its bases, $(\cE,\cB)$, and by nothing else; for a set $S$ and $0\le k\le|S|$, $\binom Sk$ is the list on $S$ whose bases are the $k$-subsets of $S$. A mark in no basis is a \emph{loop}, a mark in every basis a \emph{coloop}. The \emph{dual table} has the bases $\cE\setminus T$ with the weights $\theta_T$. The \emph{Foreman polynomials} of a table are its basis and cobasis polynomials,
\[
K(\alpha)=\sum_{T\in\cB}\theta_T\,\alpha^{T},\qquad \Psi(\alpha)=\sum_{T\in\cB}\theta_T\,\alpha^{\cE\setminus T}\qquad(\alpha\in\mathbb R^{\cE}_{>0});
\]
$K(\one)=\Psi(\one)=\sum_T\theta_T$ is the total weight. $K$ and $\Psi$ are the two letters of the table (Section~\ref{sec:letters}): they enter a family word through powers (Chapter~\ref{ch:families}). The \emph{table's gauge} is $\alpha=\one$.
\end{definition}

A table carries two descriptors of one object, what is there (the bases, read by $K$) and what is not (their complements, read by $\Psi$), and the marks on which both sit; any two of these determine the third, and the theory keeps both polynomials at every step.

% Reference 20: prop:identity; source line 1216
\begin{proposition}[One function read twice]\label{prop:identity}
For every table, $\Psi(\alpha)=\alpha^{\cE}K(\alpha^{-1})$, and the Foreman polynomials of the dual table are $K^\ast=\Psi$, $\Psi^\ast=K$. Inverting the gauge swaps $K$ and $\Psi$ up to the factor $\alpha^{\cE}$, and up to scale the gauges fixed by $\alpha\mapsto\alpha^{-1}$ are $\alpha\propto\one$.
\end{proposition}

% Reference 20: def:delcon; source line 1223
\begin{definition}[Deletion and contraction]\label{def:delcon}
For a mark $e$ and $R=\cE\setminus\{e\}$, the \emph{deletion} has the bases $T\not\ni e$ and the \emph{contraction} the bases $T\setminus\{e\}$, $T\ni e$, both on $R$ with the inherited weights; the deletion is a table when $e$ is not a coloop, and the contraction when $e$ is not a loop. With $K_{\setminus e},\Psi_{\setminus e}$ and $K_{/e},\Psi_{/e}$ the sums over these bases, a sum without terms being $0$,
\[
K=K_{\setminus e}+\alpha_eK_{/e},\qquad \Psi=\alpha_e\Psi_{\setminus e}+\Psi_{/e}.
\]
\end{definition}

\section{Circuits and molecules}

% Reference 20: def:molecule; source line 1232
\begin{definition}[Independence, circuits, linking, molecules]\label{def:molecule}
A set $S\subseteq\cE$ is \emph{independent} in a table if it lies in a basis; its \emph{rank} $r(S)$ is the largest size of an independent subset. A \emph{circuit} is a minimal set that is not independent, and a \emph{cocircuit} is a minimal set that meets every basis. Two marks are \emph{linked} if they lie in a common circuit of some table of a family on $\cE$, a finite sequence of tables on $\cE$ (Chapter~\ref{ch:families}). The \emph{molecules} are the smallest nonempty sets of marks that contain, with each mark, every mark linked to it; a molecule with one mark is a \emph{unit}, and the set of molecules is the \emph{polycule}. A family whose marks form more than one molecule is a \emph{compound}.
\end{definition}

% Reference 20: def:patterns (minors, restriction, contraction); source line 1996
\begin{definition}[Minors, restriction, contraction]\label{def:minors}
For disjoint $A,D\subseteq\cE$ such that some basis contains $A$ and avoids $D$, the \emph{minor} contracting $A$ and deleting $D$ is the table on $\cE\setminus(A\cup D)$ with the bases $T\setminus A$, $T\supseteq A$, $T\cap D=\varnothing$, and the weights $\theta_T$; deletion and contraction of one mark are the minors with $D=\{e\}$ and with $A=\{e\}$. For $X\subseteq\cE$ the \emph{restriction} $\cB|X$ is the table on $X$ with the bases $T\cap X$, $|T\cap X|=r(X)$, and the \emph{contraction} $\cB/X$ is the table on $\cE\setminus X$ with the bases $T\setminus X$ of the same $T$; when the weights of these $T$ factor as $\theta_T=\theta'_{T\cap X}\theta''_{T\setminus X}$, $\cB|X$ and $\cB/X$ carry $\theta'$ and $\theta''$. When $e$ is not a loop, $\cB/\{e\}$ is the contraction of $e$ (Definition~\ref{def:delcon}); at a loop $\cB/\{e\}$ is the deletion of $e$, while $K_{/e}=\Psi_{/e}=0$.
\end{definition}

% New: proved in the last part.
\begin{proposition}[Duality of minors]\label{prop:mot-dualminors}
For a table $(\cB,\theta)$ on $\cE$ with the dual table $\cB^\ast$ and $X\subseteq\cE$, $X'=\cE\setminus X$,
\[
(\cB|X)^\ast=\cB^\ast/X',\qquad(\cB/X)^\ast=\cB^\ast|X',
\]
the duals taken on $X$ and on $X'$: a basis is at the top of $X$ exactly when its complement is at the top of $X'$ in the dual table. When the weights of the bases at the top of $X$ factor, $\theta_T=\theta'_{T\cap X}\theta''_{T\setminus X}$, the minors of $\cB^\ast$ carry the same weights, $\cB^\ast/X'$ those of $\cB|X$ and $\cB^\ast|X'$ those of $\cB/X$ on the complemented bases. So $K_{\cB^\ast/X'}=\Psi_{\cB|X}$, $\Psi_{\cB^\ast/X'}=K_{\cB|X}$, $K_{\cB^\ast|X'}=\Psi_{\cB/X}$ and $\Psi_{\cB^\ast|X'}=K_{\cB/X}$.
\end{proposition}

% Reference 20: lem:linkexchange; source line 1236
\begin{lemma}[Linking read on the bases]\label{lem:linkexchange}
Marks $e\neq f$ lie in a common circuit of a table exactly when some basis $T$ has $e\in T$, $f\notin T$ and $(T\setminus\{e\})\cup\{f\}$ a basis.
\end{lemma}

% Reference 20: lem:circuit; source line 1243
\begin{lemma}[The circuit lemma]\label{lem:circuit}
If $w\in\mathbb R^{\cE}$ has $\ip w{\one_T}$ constant over the bases of a table, then $w_a=w_b$ for any two marks $a,b$ linked by a circuit of that table. If $\ip w{\one_T}$ is constant over the bases of every table of a family, $w$ is constant on every molecule.
\end{lemma}

% Reference 20: prop:directsum; source line 1250
\begin{proposition}[Molecules split every table]\label{prop:directsum}
Each molecule is a union of molecules of every table, those of a table being the molecules of that table alone, and every table is the \emph{direct sum} of its restrictions to the molecules: its bases are exactly the unions of one basis of each restriction. When the weights factor across the molecules, $\theta_T=\prod_M\theta_M(T\cap\cE_M)$, as they do for a list, the Foreman polynomials factor, $K(\alpha)=\prod_MK_M(\alpha_{\cE_M})$ and $\Psi(\alpha)=\prod_M\Psi_M(\alpha_{\cE_M})$.
\end{proposition}

% Reference 20: lem:separator; source line 6842
\begin{lemma}[Unions of molecules]\label{lem:separator}
For a table on $\cE$ of rank $r$ and $S\subseteq\cE$ the following are equivalent: (a) $r(S)+r(\cE\setminus S)=r$; (b) every basis meets $S$ in the same number of marks; (c) $S$ is a union of molecules of the table (Definition~\ref{def:molecule}).
\end{lemma}

% Reference 20: thm:irreducible; source line 1257
\begin{theorem}[The Foreman polynomials of a linked table do not factor]\label{thm:irreducible}
If the circuits of a table link all of its $n\ge2$ marks, then $K$ and $\Psi$ are irreducible polynomials.
\end{theorem}

\section{The Skeleton}

% Reference 20: def:skeleton; source line 1266
\begin{definition}[Skeleton counts]\label{def:skeleton}
The \emph{cocircuit layer} of a table is its set of cocircuits and the \emph{circuit layer} its set of circuits; the two layers together are the table's \emph{Skeleton}. For $S\subseteq\cE$, with $r^\ast(S)=|S|+r(\cE\setminus S)-r$ the dual rank,
\[
d(S)=|S|-r^\ast(S),\qquad d^\ast(S)=|S|-r(S),\qquad c(S)=|S|-d(S)-d^\ast(S),\qquad \varepsilon(S)=d(S)-d^\ast(S):
\]
$c(S)$ is the \emph{width} and $\varepsilon(S)$ the \emph{slant} of $S$.
\end{definition}

% Reference 20: thm:window; source line 1274
\begin{theorem}[The window]\label{thm:window}
For every $S\subseteq\cE$: $d(S)=\min_T|T\cap S|$, $d^\ast(S)=\min_T|S\setminus T|$, $r(S)=\max_T|T\cap S|$, and $\{|T\cap S|:T\in\cB\}=\{d(S),d(S)+1,\dots,r(S)\}$. Hence $c(S)=r(S)-d(S)$ is the width of this window. Moreover $c(\cE\setminus S)=c(S)$ and $\varepsilon(S)+\varepsilon(\cE\setminus S)=2r-n$.
\end{theorem}

The number of marks of $S$ that a basis occupies ranges over a window of consecutive integers; its floor is $d(S)$, its width is the width $c(S)$, and the same window read on the complements has the floor $d^\ast(S)$. Every count that the integrals of a family need is read from these two layers.

% Reference 20: def:Lambda; source line 1281
\begin{definition}[Signed indicators and the centred polynomial]\label{def:Lambda}
For $T\subseteq\cE$ the \emph{signed indicator} is $\chi_T=\one_T-\one_{\cE\setminus T}$. For $\alpha=e^{2z}$ the \emph{centred polynomial} is
\[
\Lambda(z)=\sum_{T\in\cB}\theta_T\,e^{\ip z{\chi_T}} .
\]
\end{definition}

% Reference 20: thm:tworeadings; source line 1288 (the direction $\one$ in (iii)).
% Proved in the last part: thm:tworeadings.
\begin{theorem}[One function, two Foreman polynomials]\label{thm:tworeadings}
For $\alpha=e^{2z}$:
\begin{enumerate}[label=(\roman*)]
\item $K(\alpha)=e^{\ip\one z}\Lambda(z)$ and $\Psi(\alpha)=e^{\ip\one z}\Lambda(-z)$;
\item the dual table has the centred polynomial $\Lambda^\ast(z)=\Lambda(-z)$: duality and inversion of the gauge are $z\mapsto-z$;
\item for every $v\in\mathbb R^{\cE}$ constant on each molecule of the table alone, with value $v_M$ on the molecule $M$, $\ip v{\chi_T}=\sum_Mv_M(2r(M)-|M|)$ for every basis and
\[
\Lambda(z+v)=\exp\Big(\sum_Mv_M\big(2r(M)-|M|\big)\Big)\Lambda(z);
\]
these are the only $v$ for which $\Lambda(z+\lambda v)/\Lambda(z)$ does not depend on $z$ for every $\lambda$, and at $v=\lambda\one$ the factor is $e^{\lambda(2r-n)}$.
\end{enumerate}
\end{theorem}

% Reference 20: thm:innerMGF; source line 1300 (the even statistic at the table's gauge).
% Proved in the last part: thm:innerMGF.
\begin{theorem}[The inner moment generating function]\label{thm:innerMGF}
For a gauge $\alpha=e^{2z}$ the \emph{basis law} and the \emph{cobasis law} at $\alpha$ are the laws on $\cB$
\[
\mathbb P_\alpha(T)=\frac{\theta_T\alpha^T}{K(\alpha)}=\frac{\theta_Te^{\ip z{\chi_T}}}{\Lambda(z)},\qquad \mathbb P^\ast_\alpha(T)=\frac{\theta_T\alpha^{\cE\setminus T}}{\Psi(\alpha)}=\mathbb P_{\alpha^{-1}}(T).
\]
Then
\[
\frac{\Lambda(z+w)}{\Lambda(z)}=\oE_{T\sim\mathbb P_\alpha}\big[e^{\ip w{\chi_T}}\big],\qquad \nabla\log\Lambda(z)=\oE_{T\sim\mathbb P_\alpha}[\chi_T]=2R(\alpha)-\one ,
\]
where $R_e(\alpha)=\alpha_e\partial_e\log K(\alpha)=\mathbb P_\alpha(e\in T)$ is the \emph{occupation}, and the derivatives of $\log\Lambda$ of order $k$ at $z$ are the joint cumulants of order $k$ of $\chi_T$ under $\mathbb P_\alpha$. The \emph{even} and \emph{odd statistics} $E(z)=\log\Lambda(z)+\log\Lambda(-z)$ and $O(z)=\log\Lambda(z)-\log\Lambda(-z)$ satisfy $e^{E(z)}=K\Psi/\alpha^{\cE}$, $e^{O(z)}=K/\Psi$, and, for every $w\in\mathbb R^{\cE}$,
\[
e^{E(z+w)}=e^{E(z)}\,\oE_{(T,T')\sim\mathbb P_\alpha\otimes\mathbb P^\ast_\alpha}\big[e^{\ip{2w}{\one_T-\one_{T'}}}\big]:
\]
$E(z+\cdot\,)-E(z)$ is the cumulant generating function of $\chi_T-\chi_{T'}$ for independent $T\sim\mathbb P_\alpha$ and $T'\sim\mathbb P^\ast_\alpha$; at the table's gauge $e^{E(0)}=K(\one)^2$.
For table $j$ of a family the two laws at $\alpha$ are $\mathbb P_{j,\alpha}$ and $\mathbb P^\ast_{j,\alpha}=\mathbb P_{j,\alpha^{-1}}$, the occupation is $R_j(\alpha)$, and the \emph{dual occupation} $R^\ast_j(\alpha)$, with $R^\ast_{j,e}(\alpha)=\mathbb P^\ast_{j,\alpha}(e\notin T)=1-R_{j,e}(\alpha^{-1})$, is the occupation of the dual table. At the table's gauge the two laws coincide, $\mathbb P^\ast_{j,\one}=\mathbb P_{j,\one}$.
\end{theorem}

The centred polynomial is a moment generating function: its terms at $z$, divided by $\Lambda(z)$, are the law of a random basis, and its logarithm generates the cumulants of the signed indicator. The basis law and the cobasis law at a gauge are one law read at $\alpha$ and at $\alpha^{-1}$, and at every gauge the product of the two evaluations generates the difference of a basis drawn from the basis law and an independent basis drawn from the cobasis law.

% Reference 20: def:weightlaw; source line 1625
\begin{definition}[Weight law, indicators, levels, cumulants]\label{def:weightlaw}
The \emph{weight law} of table $j$ is its basis law at the table's gauge, $\mathbb P_{j,\one}(T)=\theta_{j,T}/K_j(\one)$ on $\cB_j$; for a list it is uniform. For a random basis $T$ the \emph{indicators} are $X_e=\one[e\in T]$, and for a set $C$ of marks the \emph{level} is $N_C=|T\cap C|=\sum_{e\in C}X_e$. For functions $Y_1,\dots,Y_k$ of $T$, $\kappa_{j,\alpha}(Y_1,\dots,Y_k)$ is their joint cumulant under $\mathbb P_{j,\alpha}$, and $Y^{\times a}$ in its argument lists $Y$ $a$ times; along a string $w=(e_1,\dots,e_k)$ of marks, $\kappa_{j,\alpha}(w)=\kappa_{j,\alpha}(X_{e_1},\dots,X_{e_k})=\partial_{y_{e_1}}\cdots\partial_{y_{e_k}}\log K_j(e^y)$ at $y=\log\alpha$. With one table the index $j$ is dropped: $\mathbb P_\alpha$, $R(\alpha)$, $\kappa_\alpha$.
\end{definition}

% Reference 20: thm:newton; source line 1311
\begin{theorem}[The Skeleton is the Newton data of $\Lambda$]\label{thm:newton}
For every $S\subseteq\cE$ the values $\ip{\one_S}{\chi_T}$, $T\in\cB$, fill $\{2j-|S|:\ d(S)\le j\le r(S)\}$. The Newton polytope of $\Lambda$ is $2P-\one$ with $P=\operatorname{conv}\{\one_T\}$ the base polytope; its width in the direction $\one_S$ is $2c(S)$ and the midpoint of its extent in that direction is $\varepsilon(S)$.
\end{theorem}

% Reference 20: prop:content; source line 1318
\begin{proposition}[Loops and coloops]\label{prop:content}
If $e$ is a loop, $\Lambda(z)=e^{-z_e}\Lambda'(z|_{\cE\setminus e})$ with $\Lambda'$ the centred polynomial of the same bases on $\cE\setminus e$; if $e$ is a coloop, $\Lambda(z)=e^{z_e}\Lambda'(z|_{\cE\setminus e})$ with $\Lambda'$ that of the bases $T\setminus e$.
\end{proposition}

\section{Multivectors of any tensor rank}

% Reference 20: def:form; source line 1327
\begin{definition}[Multivectors and tensor rank]\label{def:multivector}
Let $\mathbb F$ be $\mathbb R$ or $\mathbb C$. A \emph{multivector} of a table of rank $r$ is an $\boldsymbol\omega\in\textstyle\bigwedge^r\mathbb F^{\cE}$ with $|\boldsymbol\omega_T|^2=\theta_T$ for $T\in\cB$ and $\boldsymbol\omega_T=0$ otherwise; $\boldsymbol\omega_T=\sqrt{\theta_T}$ is one. Its \emph{tensor rank} $\operatorname{rk}(\boldsymbol\omega)$ is the least number of decomposable elements of $\textstyle\bigwedge^r\mathbb F^{\cE}$ with sum $\boldsymbol\omega$, written $\boldsymbol\omega=\sum_{k=1}^{\operatorname{rk}(\boldsymbol\omega)}\boldsymbol\omega^{(k)}$, $\boldsymbol\omega^{(k)}_T=\det A^{(k)}_T$ with \emph{slices} $A^{(k)}\in\mathbb F^{r\times n}$. The table is \emph{tensorminental of tensor rank $q$} over $\mathbb F$ when $q$ is the least tensor rank of its multivectors, and \emph{determinantal} when $q=1$. The \emph{dual multivector} is the Hodge complement $\boldsymbol\omega^\ast\in\textstyle\bigwedge^{n-r}\mathbb F^{\cE}$, $|\boldsymbol\omega^\ast_{\cE\setminus T}|=|\boldsymbol\omega_T|$, and $\operatorname{rk}(\boldsymbol\omega^\ast)=\operatorname{rk}(\boldsymbol\omega)$.
\end{definition}

Every table has multivectors, and every multivector has a tensor rank; tensor rank one, the determinantal case, is where a single matrix carries the table, and it is one case of the statements on multivectors, which hold at every tensor rank.

% Reference 20: prop:tensorletters; source line 1331
\begin{proposition}[The letters of a tensorminental table]\label{prop:tensorletters}
For a decomposition $\boldsymbol\omega=\sum_k\boldsymbol\omega^{(k)}$ into decomposable elements with slices $A^{(k)}$ and $D=\operatorname{diag}(\alpha)$,
\[
K(\alpha)=\sum_{k,l}\det\big(A^{(k)}DA^{(l)\mathsf H}\big)=\one^{\mathsf H}\mathbf G(\alpha)\one,\qquad \mathbf G_{kl}(\alpha)=\det\big(A^{(k)}DA^{(l)\mathsf H}\big),
\]
and $\Psi$ is the same expression for the slices of the dual multivector. $\mathbf G(\alpha)$ is the Gram matrix of the $\boldsymbol\omega^{(k)}$ for $\ip\eta\xi_\alpha=\sum_T\alpha^T\eta_T\overline{\xi_T}$, Hermitian and positive semidefinite.
\end{proposition}

% Reference 20: thm:rankhered; source line 1342
\begin{theorem}[Tensor rank is hereditary]\label{thm:rankhered}
For disjoint $A,D\subseteq\cE$ with $A$ extending to a basis avoiding $D$, $\boldsymbol\omega_{A|D}:=(\mathbf e_A\lrcorner\boldsymbol\omega)|_{\cE\setminus(A\cup D)}$, the restriction of the interior product of $\boldsymbol\omega$ with the wedge $\mathbf e_A$ of the unit vectors of the marks of $A$, $(\mathbf e_A\lrcorner\boldsymbol\omega)_X=\pm\boldsymbol\omega_{X\cup A}$, is a multivector of the minor that contracts $A$ and deletes $D$, with the weight $\theta_{X\cup A}$ on each basis $X$ of the minor, and $\operatorname{rk}(\boldsymbol\omega_{A|D})\le\operatorname{rk}(\boldsymbol\omega)$.
\end{theorem}

% Reference 20: thm:interference; source line 2071
\begin{theorem}[The pattern law of a tensorminental table is a finite interference]\label{thm:interference}
Let the table have a multivector of tensor rank $q$ with slices $A^{(1)},\dots,A^{(q)}$, let $B\subseteq\cE$ and $R=\cE\setminus B$, and let $\alpha$ be a gauge at which every $M_{kl}=A^{(k)}D_\alpha A^{(l)\mathsf H}$, $D_\alpha=\operatorname{diag}(\alpha)$, with $\mathbf G_{kl}\not\equiv0$ is invertible; pairs with $\mathbf G_{kl}\equiv0$ as a polynomial in the gauge are omitted (Proposition~\ref{prop:tensorletters}). For $x\in\mathbb C^B$, with $x\circ\alpha_B=(x_e\alpha_e)_{e\in B}$,
\[
\oE_{T\sim\mathbb P_\alpha}\Big[\prod_{e\in B}x_e^{\one[e\in T]}\Big]=\frac{K(x\circ\alpha_B,\alpha_R)}{K(\alpha)}=\sum_{k,l}w_{kl}\det\big(I+(D_x-I)\Pi^{(kl)}_{BB}\big),
\]
where $\mathbb P_\alpha$ is the basis law at $\alpha$, $D_x=\operatorname{diag}(x)$, $w_{kl}=\mathbf G_{kl}(\alpha)/K(\alpha)$ and $\Pi^{(kl)}=A^{(l)\mathsf H}M_{kl}^{-1}A^{(k)}D_\alpha$. Here $\sum_{k,l}w_{kl}=1$, $w_{lk}=\overline{w_{kl}}$, every $\Pi^{(kl)}$ is idempotent, and $\Pi^{(kk)}$ is the $D_\alpha$-orthogonal projection onto the row space of $A^{(k)}$. The hypothesis holds off a proper algebraic set of gauges.
\end{theorem}

% Reference 20: cor:interference; source line 2082 (orders one to three in (ii)).
% Proved in the last part: cor:interference.
\begin{corollary}[Moments and cumulants]\label{cor:interference}
In the setting of Theorem~\ref{thm:interference}:
\begin{enumerate}[label=(\roman*)]
\item $\mathbb P_\alpha(S\subseteq T)=\sum_{k,l}w_{kl}\det\Pi^{(kl)}_{SS}$;
\item the law of $T\cap B$ is the $w$-mixture of $q^2$ determinantal components, whose cumulants are $\kappa^{(kl)}(e)=\Pi^{(kl)}_{ee}$ and, for $p\ge2$ distinct marks,
\[
\kappa^{(kl)}(e_1,\dots,e_p)=(-1)^{p-1}\sum_\varsigma\prod_{i=1}^p\Pi^{(kl)}_{e_{\varsigma(i)}e_{\varsigma(i+1)}},
\]
the sum over the cyclic orders $\varsigma$ of $\{1,\dots,p\}$, with $\varsigma(p+1)=\varsigma(1)$; the joint cumulants $\kappa_\alpha$ of the indicators $X_e=\one[e\in T]$ under $\mathbb P_\alpha$ (Definition~\ref{def:weightlaw}) are their total cumulants over the mixing index: for distinct marks $e_1,\dots,e_p$ of $B$,
\[
\kappa_\alpha(e_1,\dots,e_p)=\sum_{\mathfrak p}\kappa_w\Big(\big(\kappa^{(kl)}(e_C)\big)_{C\in\mathfrak p}\Big),
\]
the sum over the partitions $\mathfrak p$ of $\{1,\dots,p\}$ into nonempty blocks, $e_C=(e_c)_{c\in C}$, where for finitely many functions $(Y_i)_{i\in I}$ of the index $(k,l)$
\[
\kappa_w\big((Y_i)_{i\in I}\big)=\sum_{\mathfrak q}(-1)^{|\mathfrak q|-1}(|\mathfrak q|-1)!\prod_{D\in\mathfrak q}\ \sum_{k,l}w_{kl}\prod_{i\in D}Y_i(k,l),
\]
the sum over the partitions $\mathfrak q$ of $I$, is their joint cumulant over the index with the mixing coefficients $w_{kl}$; at one mark this is the occupation $R_e(\alpha)=\sum_{k,l}w_{kl}\Pi^{(kl)}_{ee}$. This is an identity of generating functions, valid for complex mixing coefficients summing to one.
\end{enumerate}
\end{corollary}

A table with a multivector of tensor rank $q$ reads the law of its patterns as an interference of $q^2$ determinantal components, mixed with complex coefficients that sum to one. At tensor rank one there is a single component and the mixture is a determinantal law; at every higher tensor rank the joint cumulants over the mixing index of the components' cumulants are the part that no single component carries.


\chapter{Families}\label{ch:families}

\section{The interest and its walls}

% Reference 20: def:family; source line 1375
\begin{definition}[Family, interest, walls]\label{def:family}
A \emph{family} on $\cE$ is a finite sequence $F=\big((\cB_j,\theta_j);\,s_j,t_j\big)_{j=1}^N$ of tables on $\cE$ with exponents $s_j,t_j\in\mathbb C$; $K_j,\Psi_j,\Lambda_j,E_j,O_j$, $r_j$ and $d_j,d_j^\ast,c_j,\varepsilon_j$ belong to table $j$. A table is \emph{active} when its exponents are not both zero, and a letter when its exponent is not zero. For \emph{masses} $\mvec\in\mathbb C^{\cE}$ the \emph{interest} is
\[
J_F(\mvec;\mathbf s,\mathbf t)=\int_{\mathbb R^{\cE}_{>0}}\alpha^{\mvec-\one}e^{-|\alpha|}\prod_jK_j(\alpha)^{-s_j}\Psi_j(\alpha)^{-t_j}\,d\alpha ,
\]
with its slots: the family's tables, the masses $\mvec$ and the exponents $\mathbf s=(s_j)_j$, $\mathbf t=(t_j)_j$. The \emph{family word} $W=\prod_jK_j^{-s_j}\Psi_j^{-t_j}$ is homogeneous of degree $-\sum_j(s_jr_j+t_jr^\ast_j)$, and with the measure $d\alpha_{\mvec}=\alpha^{\mvec-\one}d\alpha$ (Definition~\ref{def:conv}(vii)), $J_F(\mvec;\mathbf s,\mathbf t)=\int W\,e^{-|\alpha|}d\alpha_{\mvec}$. The \emph{pair coordinates} are the \emph{strengths} $\sigma_j=\frac{s_j+t_j}2$ and the \emph{leans} $\delta_j=\frac{s_j-t_j}2$, with the total strength $\sigma=\sum_j\sigma_j$, so that $s_j=\sigma_j+\delta_j$ and $t_j=\sigma_j-\delta_j$; the \emph{offsets} are $\bmu=\mvec-\sigma\one$. The \emph{walls} and \emph{top walls} are
\begin{gather*}
\kappa_S=m_S-\sum_j\big(s_jd_j(S)+t_jd^\ast_j(S)\big),\\
\kappa^{\mathrm{top}}_S=\sum_j\big(s_jr_j(S)+t_j(|S|-d_j(S))\big)-m_S\qquad(\varnothing\neq S\subseteq\cE).
\end{gather*}
The \emph{imbalance} is $\kappa_{\cE}=m_{\cE}-\sum_j(s_jr_j+t_jr^\ast_j)$, and the \emph{balance wall} is the hyperplane $\kappa_{\cE}=0$. On it the data are fixed by their pair coordinates $(\bmu_\perp;\boldsymbol\sigma,\boldsymbol\delta)$, with $\bmu_\perp=\bmu-\frac{\mu_{\cE}}n\one$ the projection of the offsets to $\one^\perp$: the exponents are $(\boldsymbol\sigma+\boldsymbol\delta,\boldsymbol\sigma-\boldsymbol\delta)$ and the masses are the \emph{balanced masses}
\[
\mvec_F(\bmu_\perp;\boldsymbol\sigma,\boldsymbol\delta)=\bmu_\perp+\Big(\sigma+\frac1n\sum_j\delta_j(2r_j-n)\Big)\one ,
\]
linear in $(\bmu_\perp;\boldsymbol\sigma,\boldsymbol\delta)$. $F$ is \emph{connected} if every nonempty proper $S$ has $c_j(S)>0$ for some active table $j$. A family of one table is named by its table: for a table $\mathcal A$ with its weights $\theta^{\mathcal A}$, every object of the family $((\mathcal A,\theta^{\mathcal A});s,t)$ carries $\mathcal A$ in place of $F$, as $J_{\mathcal A}(\mvec;s,t)$. A table may carry two systems of positive weights, $\theta^K$ for its letter $K$ and $\theta^\Psi$ for its letter $\Psi$; at $(s,t)$ it names the family of the two tables $((\cB,\theta^K);s,0)$ and $((\cB,\theta^\Psi);0,t)$, whose word is $K_{\theta^K}^{-s}\Psi_{\theta^\Psi}^{-t}$.
\end{definition}

The exponents of a table enter in two coordinates: $\sigma_j$, how much the family weighs the table, and $\delta_j$, which side it leans to, toward the bases ($K$) or toward their complements ($\Psi$). The pair $(K_j\Psi_j,K_j/\Psi_j)$ carries $\sigma_j$ and $\delta_j$ separately, and the walls are read in either coordinates.

% Reference 20: thm:pairform; source line 1387
\begin{theorem}[The interest is a partition function in the pair coordinates]\label{thm:pairform}
\begin{enumerate}[label=(\roman*)]
\item In the offsets,
\[
J_F(\mvec;\mathbf s,\mathbf t)=\int_{\mathbb R^{\cE}_{>0}}\alpha^{\bmu-\one}e^{-|\alpha|}\prod_j\Lambda_j(\tfrac12\log\alpha)^{-s_j}\Lambda_j(-\tfrac12\log\alpha)^{-t_j}\,d\alpha .
\]
\item In signed coordinates,
\[
J_F(\mvec;\mathbf s,\mathbf t)=2^{n}\int_{\mathbb R^{\cE}}\exp\Big(2\ip{\bmu}{z}-\sum_j\big(\sigma_jE_j(z)+\delta_jO_j(z)\big)-\sum_ee^{2z_e}\Big)\,dz ,
\]
the partition function of the exponential family with natural parameters $(\bmu;\sigma_j,\delta_j)$ and sufficient statistics $(2z;-E_j,-O_j)$.
\item $\kappa_S=\mu_S+\sum_j\big(\sigma_j\,c_j(S)-\delta_j\,\varepsilon_j(S)\big)$ and $\kappa_{\cE}=\mu_{\cE}-\sum_j\delta_j(2r_j-n)$.
\end{enumerate}
\end{theorem}

% Reference 20: lem:greedy; source line 1429
\begin{lemma}[Greedy bases on a sector]\label{lem:greedy}
Let $\pi$ be a total order of $\cE$ with tails $U_k=\{\pi(k),\dots,\pi(n)\}$. There are bases $T^\pm_\pi$ with $|T^+_\pi\cap U_k|=d(U_k)$ and $|U_k\setminus T^-_\pi|=d^\ast(U_k)$ for all $k$, and on the sector $H_\pi=\{\alpha_{\pi(1)}\ge\dots\ge\alpha_{\pi(n)}\}$ every basis has $\alpha^T\le\alpha^{T^+_\pi}$ and $\alpha^{\cE\setminus T}\le\alpha^{\cE\setminus T^-_\pi}$. For a nonempty set of $r$-element subsets of $\cE$ this conclusion holds exactly when the set has the exchange property.
\end{lemma}

The sectors are where the exchange property is used: on each sector one basis dominates every monomial of $K$ and one cobasis every monomial of $\Psi$, and that is what gives the walls their form. What does not use the exchange property is the opening of every monomial by its shares (Corollary~\ref{thm:E1}).

% Reference 20: lem:cube; source line 1440
\begin{lemma}[Cube integrals]\label{lem:cube}
Let $h(q,\lambda)$ be smooth in $q$ near $[0,1]^m$, with every $q$-derivative continuous in $(q,\lambda)$ and holomorphic in $\lambda$ on an open set. Then $\operatorname{cube}(x,\lambda;h)=\int_{[0,1]^m}\prod_kq_k^{x_k-1}h(q,\lambda)\,dq$, defined for $\re x_k>0$, extends meromorphically to $\mathbb C^m$, jointly in $\lambda$, with poles on the hyperplanes $x_k=-j$, $j\in\mathbb N$: with $h=\sum_{j<M}h_j(q',\lambda)q_1^j+q_1^Mr_M$,
\[
\operatorname{cube}(x,\lambda;h)=\sum_{j<M}\frac{\operatorname{cube}(x',\lambda;h_j)}{x_1+j}+\operatorname{cube}\big((x_1+M,x'),\lambda;r_M\big),
\]
repeated in each coordinate.
\end{lemma}

% Reference 20: thm:exist; source line 1451
\begin{theorem}[Existence and continuation]\label{thm:exist}
\begin{enumerate}[label=(\roman*)]
\item $J_F(\mvec;\mathbf s,\mathbf t)$ converges if and only if $\re\kappa_S>0$ for every nonempty $S\subseteq\cE$, whatever $\mathbf s,\mathbf t\in\mathbb C^N$; it is holomorphic there.
\item On that region
\[
J_F(\mvec;\mathbf s,\mathbf t)=\Gamma(\kappa_{\cE})\sum_\pi\int_{[0,1]^{n-1}}\prod_{k=2}^nq_k^{\kappa_{U_k}-1}h_\pi(q)\,dq,\qquad h_\pi=L_\pi^{-\kappa_{\cE}}\prod_jk_{j,\pi}^{-s_j}\psi_{j,\pi}^{-t_j},
\]
where in the sector coordinates $\alpha_{\pi(1)}=\rho$, $\alpha_{\pi(k)}=\rho q_2\cdots q_k$ one has
\[
K_j=\rho^{r_j}\prod_kq_k^{d_j(U_k)}k_{j,\pi}(q),\qquad \Psi_j=\rho^{n-r_j}\prod_kq_k^{d^\ast_j(U_k)}\psi_{j,\pi}(q),\qquad |\alpha|=\rho L_\pi(q);
\]
the \emph{sector letters} $k_{j,\pi},\psi_{j,\pi},L_\pi$ have nonnegative coefficients and positive constant terms $\theta_{j,T^+_{j,\pi}}$, $\theta_{j,T^-_{j,\pi}}$, $T^\pm_{j,\pi}$ the greedy bases of table $j$ (Lemma~\ref{lem:greedy}), and $1$, and $L_\pi\in[1,n]$.
\item $J_F(\mvec;\mathbf s,\mathbf t)$ extends to a meromorphic function of $(\mvec,\mathbf s,\mathbf t)$ with poles on the hyperplanes $\kappa_S=-j$, $S\neq\varnothing$, $j\in\mathbb N$; identities valid on an open subset of the region persist.
\end{enumerate}
\end{theorem}

% Reference 20: cor:pastwalls; source line 1470
\begin{corollary}[Past several walls at once]\label{cor:pastwalls}
For an order $\pi$ write $\kappa_k=\kappa_{U_k}$ and let
\[
I_\pi=\int_{[0,1]^{n-1}}\prod_kq_k^{\kappa_k-1}h_\pi(q)\,dq=\int_{[0,\infty)^{n-1}}e^{-\sum_k\kappa_ku_k}h_\pi(e^{-u})\,du
\]
be the term of Theorem~\ref{thm:exist}(ii), written also in the gaps $u_k=-\log q_k$; in the log-ratios $\eta_k=\log(\alpha_{\pi(1)}/\alpha_{\pi(k)})=u_2+\cdots+u_k$ the sector is the trail region $0<\eta_2<\cdots<\eta_n$. Let no $\kappa_k$ be a nonpositive integer and $L_k\in\mathbb N$ with $\re\kappa_k+L_k>0$. Then the continuation of $I_\pi$ in the masses is
\[
\sum_{A\subseteq\{2,\dots,n\}}\ \sum_{\mathbf l:\ l_k<L_k\ (k\in A)}\ \prod_{k\in A}\frac1{\kappa_k+l_k}\int_{[0,1]^{A^c}}\prod_{k\notin A}q_k^{\kappa_k-1}\Big[\prod_{k\notin A}(1-\operatorname{Tay}_k)\,h^{(A,\mathbf l)}_\pi\Big]dq_{A^c},
\]
where $h^{(A,\mathbf l)}_\pi=\partial^{\mathbf l}_{q_A}h_\pi/\mathbf l!$ at $q_A=0$ and $\operatorname{Tay}_k$ takes the Taylor polynomial in $q_k$ of degree below $L_k$: the term $A$ collapses the tails $U_k$, $k\in A$, together, at the levels $\mathbf l$.
\end{corollary}

% Reference 20: thm:projective; source line 1504
\begin{theorem}[Projective form; the period]\label{thm:projective}
Put
\begin{align*}
P_F(\mvec;\mathbf s,\mathbf t)&=\int_{\Delta_{\cE}}u^{\bmu-\one}\prod_j\Lambda_j(\tfrac12\log u)^{-s_j}\Lambda_j(-\tfrac12\log u)^{-t_j}\,d\omega_{\cE}(u)\\
&=\int_{\Delta_{\cE}}u^{\mvec-\one}\prod_jK_j(u)^{-s_j}\Psi_j(u)^{-t_j}\,d\omega_{\cE}(u).
\end{align*}
It converges if and only if $\re\kappa_S>0$ for every nonempty proper $S$, and $J_F(\mvec;\mathbf s,\mathbf t)=\Gamma(\kappa_{\cE})P_F(\mvec;\mathbf s,\mathbf t)$ as meromorphic functions; $P_F(\mvec;\mathbf s,\mathbf t)$ is the \emph{period}. On the balance wall the period is the residue of $J_F(\mvec;\mathbf s,\mathbf t)$ there; for every positive $h$ homogeneous of degree one, the integral of $\alpha^{\mvec-\one}\prod_jK_j^{-s_j}\Psi_j^{-t_j}d\alpha$ over $\{h=1\}$ (Definition~\ref{def:conv}(vii)) is the period.
\end{theorem}

% Reference 20: def:continued; source line 1485
\begin{definition}[The continued interest]\label{def:continued}
The continuation of Theorem~\ref{thm:exist}(iii) keeps the symbol $J_F(\mvec;\mathbf s,\mathbf t)$ away from its poles. Among them is the whole balance wall, where $J_F(\mvec;\mathbf s,\mathbf t)$ has the residue $P_F(\mvec;\mathbf s,\mathbf t)$ (Theorem~\ref{thm:projective}). The factor $e^{-|\alpha|}$ damps the scaling mode, so no completion enters the continued interest.
\end{definition}

% Reference 20: thm:confJ; source line 1402
\begin{theorem}[The interest is a Gramma; its letter and mark relations]\label{thm:confJ}
For $\re s_j,\re t_j>0$, on the scales $\cE\sqcup\{v_j\}\sqcup\{w_j\}$ let $t_F$ be the configuration of the columns $(\mathbf e_e;0;0)$ with coefficient $1$, and $(\chi_T/2;\mathbf e_j;0)$, $(-\chi_T/2;0;\mathbf e_j)$ with coefficient $\theta_{j,T}$ for $T\in\cB_j$. Then
\[
\prod_j\Gamma(s_j)\Gamma(t_j)\,J_F(\mvec;\mathbf s,\mathbf t)=\sG_{(\bmu;\mathbf s;\mathbf t);t_F}(\one\oplus\boldsymbol\theta\oplus\boldsymbol\theta),
\]
either side converging with the other. As identities of meromorphic functions:
\begin{enumerate}[label=(\roman*)]
\item (letter relations) for each $j$,
\[
J_F(\mvec;\mathbf s,\mathbf t)=\sum_{T\in\cB_j}\theta_{j,T}J_F(\mvec+\one_T;\mathbf s+\mathbf e_j,\mathbf t)=\sum_{T\in\cB_j}\theta_{j,T}J_F(\mvec+\one_{\cE\setminus T};\mathbf s,\mathbf t+\mathbf e_j);
\]
\item (mark relations) for each $e$,
\begin{multline*}
m_eJ_F(\mvec;\mathbf s,\mathbf t)=J_F(\mvec+\mathbf e_e;\mathbf s,\mathbf t)+\sum_js_j\sum_{T\ni e}\theta_{j,T}J_F(\mvec+\one_T;\mathbf s+\mathbf e_j,\mathbf t)\\
+\sum_jt_j\sum_{T\not\ni e}\theta_{j,T}J_F(\mvec+\one_{\cE\setminus T};\mathbf s,\mathbf t+\mathbf e_j);
\end{multline*}
\item the same for the period $P_F(\mvec;\mathbf s,\mathbf t)$ of Theorem~\ref{thm:projective}, with $J_F(\mvec;\mathbf s,\mathbf t)=\Gamma(\kappa_{\cE})P_F(\mvec;\mathbf s,\mathbf t)$: the letter relations hold verbatim, and in the mark relation the pure mass shift carries the factor $\kappa_{\cE}$, $m_eP_F(\mvec;\mathbf s,\mathbf t)=\kappa_{\cE}P_F(\mvec+\mathbf e_e;\mathbf s,\mathbf t)+\cdots$.
\end{enumerate}
\end{theorem}

The interest is not a new kind of object: it is one Gramma, on the marks and one scale for each letter, and its relations are that Gramma's exchange relations. The letter relation is the exchange relation on a letter's scale, and the mark relation the exchange relation on a mark's scale, rewritten with the letter relations.

% Reference 20: def:shapelaw; source line 3407
\begin{definition}[Breta shape law; Breta law]\label{def:shapelaw}
For real data with every proper wall positive, the \emph{Breta shape law} $\mathrm{Breta}^\circ(F;\mvec;\mathbf s,\mathbf t)$ is the probability measure
\[
u^{\mvec-\one}\prod_jK_j(u)^{-s_j}\Psi_j(u)^{-t_j}\,d\omega_{\cE}(u)\big/P_F(\mvec;\mathbf s,\mathbf t)\quad\text{on }\Delta_{\cE}.
\]
With every wall positive, the \emph{Breta law} $\mathrm{Breta}(F;\mvec;\mathbf s,\mathbf t)$ is the law on $\mathbb R^{\cE}_{>0}$ with density $\alpha^{\mvec-\one}e^{-|\alpha|}\prod_jK_j(\alpha)^{-s_j}\Psi_j(\alpha)^{-t_j}/J_F(\mvec;\mathbf s,\mathbf t)$.
\end{definition}

% Reference 20: prop:latent; source line 1418 (over $\mathbb R$ with $s_j=t_j$).
% Proved in the last part: prop:latent.
\begin{proposition}[The interest over latent vectors]\label{prop:latent}
Let every table $j$ of a family be determinantal over $\mathbb F_j\in\{\mathbb R,\mathbb C\}$ (Definition~\ref{def:multivector}), put $\beta_j=\dim_{\mathbb R}\mathbb F_j$, and let $2s_j/\beta_j$ and $2t_j/\beta_j$ be nonnegative integers. Let $V_j\subseteq\mathbb F_j^{\cE}$ be the row space of a slice of table $j$ and $V_j^\perp$ its orthogonal complement, each with its Lebesgue measure, and for the \emph{latent vectors} $x_{j,i}\in V_j$, $1\le i\le2s_j/\beta_j$, and $x^\ast_{j,i}\in V_j^\perp$, $1\le i\le2t_j/\beta_j$, put $\mathfrak q_e=\sum_j\big(\sum_i|(x_{j,i})_e|^2+\sum_i|(x^\ast_{j,i})_e|^2\big)$. Then
\[
J_F(\mvec;\mathbf s,\mathbf t)=\frac{\prod_e\Gamma(m_e)}{\prod_jK_j(\one)^{s_j+t_j}\pi^{r_js_j+r^\ast_jt_j}}\int\prod_{e\in\cE}\big(1+\mathfrak q_e\big)^{-m_e}\prod_{j,i}dx_{j,i}\,dx^\ast_{j,i},
\]
both sides finite or infinite together. Where they are finite, the Breta law $\mathrm{Breta}(F;\mvec;\mathbf s,\mathbf t)$ (Definition~\ref{def:shapelaw}) is the law of $\alpha$ under the joint law of $\alpha$ and the latent vectors with density proportional to $\alpha^{\mvec-\one}\exp\big(-\sum_e\alpha_e(1+\mathfrak q_e)\big)$, under which, given the latent vectors, the scales $\alpha_e$ are independent and $\alpha_e(1+\mathfrak q_e)$ has the law $\operatorname{Gamma}_{m_e}$. Over $\mathbb R$ with $s_j=t_j$ the sums $u_{j,i}=x_{j,i}+x^\ast_{j,i}$ run over $\mathbb R^{\cE}$, with $|(x_{j,i})_e|^2+|(x^\ast_{j,i})_e|^2=(\operatorname{pr}_ju_{j,i})_e^2+\big((I-\operatorname{pr}_j)u_{j,i}\big)_e^2$, $\operatorname{pr}_j$ the orthogonal projection onto $V_j$.
\end{proposition}

The letters of a determinantal table open as Gaussian integrals: its basis polynomial is a Gram determinant, whose inverse powers are Gaussian integrals over the row space of its slice, real or complex, and its cobasis polynomial is the same over the orthogonal complement, each letter with as many latent vectors as its exponent asks; over $\mathbb R$ at equal exponents the two pair into Gaussian vectors on $\mathbb R^{\cE}$ read through the two projections. The integrals over the scales are then Gamma integrals, and the interest becomes a rational integral in the latent vectors.

% Reference 20: prop:gammadraws; source line 1489
\begin{proposition}[The interest under Gamma draws]\label{prop:gammadraws}
Let the data be real with $\mvec>0$.
\begin{enumerate}[label=(\roman*)]
\item Where every proper wall is positive, $P_F(\mvec;\mathbf s,\mathbf t)=\Gst{n}{1}{\mvec}{m_{\cE}}\,\oE_{U\sim\Dir_{\mvec}}[W(U)]$. With $\phi_{\mvec}$ the joint law of $\big(K_1(U),\dots,K_N(U),\Psi_1(U),\dots,\Psi_N(U)\big)$ for $U\sim\Dir_{\mvec}$, an input on the source $(0,\infty)^{2N}$, $P_F(\mvec;\mathbf s,\mathbf t)$ is a Trailmoment at the unit configuration $I_{2N}$:
\[
P_F(\mvec;\mathbf s,\mathbf t)=\Gst{n}{1}{\mvec}{m_{\cE}}\,\big(\sT_{(\mathbf s,\mathbf t);I_{2N}}[\mathrm{id}]\,\phi_{\mvec}\big)\big((0,\infty)^{2N}\big).
\]
\item Where moreover $\kappa_{\cE}>0$, so that every wall is positive, let $\alpha\sim\operatorname{Gamma}_{\mvec}$, with total $|\alpha|$ and shape $U=\alpha/|\alpha|$. The scale and the shape separate: $|\alpha|\sim\operatorname{Gamma}_{m_{\cE}}$ and $U\sim\Dir_{\mvec}$ are independent, and
\begin{gather*}
J_F(\mvec;\mathbf s,\mathbf t)=\prod_e\Gamma(m_e)\,\oE_{\alpha\sim\operatorname{Gamma}_{\mvec}}[W(\alpha)],\\
\oE_{\alpha\sim\operatorname{Gamma}_{\mvec}}[W(\alpha)]=\frac{\Gamma(\kappa_{\cE})}{\Gamma(m_{\cE})}\,\oE_{U\sim\Dir_{\mvec}}[W(U)].
\end{gather*}
\item In the exponents, $J_F(\mvec;\mathbf s,\mathbf t)/\prod_e\Gamma(m_e)$ is the moment generating function of the letters' logarithms $-\log K_j(\alpha)$, $-\log\Psi_j(\alpha)$ for $\alpha\sim\operatorname{Gamma}_{\mvec}$, and $P_F(\mvec;\mathbf s,\mathbf t)/\Gst{n}{1}{\mvec}{m_{\cE}}$ that of $-\log K_j(U)$, $-\log\Psi_j(U)$ for $U\sim\Dir_{\mvec}$. Their ratio $\Gamma(\kappa_{\cE})/\Gamma(m_{\cE})$ has its pole on the balance wall, where the moment generating function under $\operatorname{Gamma}_{\mvec}$ diverges and that under $\Dir_{\mvec}$ is the period over $\Gst{n}{1}{\mvec}{m_{\cE}}$.
\end{enumerate}
\end{proposition}

% Reference 20: thm:wallscremona; source line 1515
\begin{theorem}[Walls under Cremona]\label{thm:wallscremona}
As functions of $(\bmu;\boldsymbol\sigma,\boldsymbol\delta)$: $\kappa^{\mathrm{top}}_S(\bmu;\boldsymbol\sigma,\boldsymbol\delta)=\kappa_S(-\bmu;\boldsymbol\sigma,-\boldsymbol\delta)$, $\kappa_S+\kappa^{\mathrm{top}}_S=2\sum_j\sigma_jc_j(S)$, and $\kappa_{\cE\setminus S}=\kappa_{\cE}+\kappa^{\mathrm{top}}_S$. In pair coordinates $\kappa^{\mathrm{top}}_S=\sum_j\big(\sigma_jc_j(S)+\delta_j\varepsilon_j(S)\big)-\mu_S$.
\end{theorem}

% Reference 20: thm:cremona; source line 1522
\begin{theorem}[Cremona on the balance wall]\label{thm:cremona}
For a connected family, on the hyperplane $\kappa_{\cE}=0$,
\[
P_F(\bmu+\sigma\one;\boldsymbol\sigma+\boldsymbol\delta,\boldsymbol\sigma-\boldsymbol\delta)
=P_F(-\bmu+\sigma\one;\boldsymbol\sigma-\boldsymbol\delta,\boldsymbol\sigma+\boldsymbol\delta),
\]
the walls of the right side at $S$ being those of the left at $\cE\setminus S$; the identity continues meromorphically along the hyperplane.
\end{theorem}

On the balance wall the period is unchanged when the offsets and the leans change sign together: Cremona (Definition~\ref{def:realization}) swaps the two exponents of each table, $\delta_j\mapsto-\delta_j$, negates the offsets, and swaps every wall with the wall of the complement.

% Reference 20: prop:dualfamily; source line 1529
\begin{proposition}[Dual tables]\label{prop:dualfamily}
Replacing every table by its dual, with the same exponents, gives the interest at the exponents $(t_j,s_j)$: in pair coordinates $\delta_j\mapsto-\delta_j$.
\end{proposition}

% Reference 20: prop:compoundwall; source line 1536 ($c_j(S)=0$ in every table).
% Proved in the last part: prop:compoundwall.
\begin{proposition}[The wall of a compound]\label{prop:compoundwall}
For every nonempty proper $S$, $\kappa_S+\kappa_{\cE\setminus S}=\kappa_{\cE}+2\sum_j\sigma_jc_j(S)$. So if $c_j(S)=0$ in every table with $\sigma_j\neq0$, then $\kappa_S+\kappa_{\cE\setminus S}=\kappa_{\cE}$, and on the balance wall a compound split along $S$ has a wall with real part at most zero, and its period diverges.
\end{proposition}

\section{The native gauge}

% Reference 20: def:native; source line 1545 (real data with $s_j,t_j\ge0$, not all zero, and $\mvec>0$; a positive gauge).
% Proved in the last part: def:native.
\begin{definition}[Target; native gauge]\label{def:native}
For data with $m_{\cE}\neq0$ the \emph{target} is $\btau=\frac{\sum_j(s_jr_j+t_jr^\ast_j)}{m_{\cE}}\mvec$; on the balance wall it is $\mvec$. A point $\alpha^\star\in(\mathbb C^\times)^{\cE}$ with $K_j(\alpha^\star)\Psi_j(\alpha^\star)\neq0$ for every $j$ is \emph{native} when its \emph{usage} equals the target,
\[
\sum_j\big(s_j\,R_j(\alpha^\star)+t_j\,R^\ast_j(\alpha^\star)\big)=\btau ,
\]
where $R_{j,e}=\alpha_e\partial_eK_j/K_j$ is the occupation of table $j$ and $R^\ast_{j,e}=\alpha_e\partial_e\Psi_j/\Psi_j$ that of its dual (Theorem~\ref{thm:innerMGF}), rational functions homogeneous of degree zero; so nativity is a property of the point up to scale. A native point in $\mathbb R^{\cE}_{>0}$, written $\alpha^\star=e^{2z^\star}$, is a \emph{native gauge}; the other native points are the \emph{native gauges at complex points}.
\end{definition}

% Reference 20: def:potential; source line 1638
\begin{definition}[Potential; its even and odd parts]\label{def:potential}
In logarithmic coordinates the \emph{potential} of a family with its exponent pairs is
\[
\Phi(y)=\sum_j\Big(s_j\log K_j(e^{y})+t_j\log\Psi_j(e^{y})\Big)=\Phi_{\mathbf E}(y)+\Phi_{\mathbf O}(y),
\]
\[
\Phi_{\mathbf E}(y)=\sum_j\sigma_j\log\mathbf E_j(e^{y}),\qquad \Phi_{\mathbf O}(y)=\sum_j\delta_j\log\mathbf O_j(e^{y}),\qquad \mathbf E_j=K_j\Psi_j,\quad\mathbf O_j=K_j/\Psi_j ,
\]
the \emph{even} and \emph{odd letters} of table $j$, so that $\mathbf E_j(e^{2z})=e^{2\ip\one z}e^{E_j(z)}$ and $\mathbf O_j(e^{2z})=e^{O_j(z)}$. The pair $(\Phi_{\mathbf E},\Phi_{\mathbf O})$, either part written $\Phi_{\mathbf X}$ with $\mathbf X\in\{\mathbf E,\mathbf O\}$, is read whole: neither part is switched off at any lean, and the two are never added before they are read.
\end{definition}

% Reference 20: thm:saddle; source line 1553
% Proved in the last part: thm:saddle (with the native gauges of def:native).
\begin{theorem}[The native gauge is the saddle]\label{thm:saddle}
The potential satisfies $\partial_{y_e}\Phi=\sum_j\big(s_jR_{j,e}+t_jR^\ast_{j,e}\big)$, the usage at $\alpha=e^y$, and the integrand of the interest is $e^{\ip\mvec y-\Phi(y)-|\alpha|}$ in $y$. In the signed coordinates $z=\frac12y$, $\Phi(y)=\sigma\ip\one y+\Phi^\circ(z)$ with the \emph{centred potential} $\Phi^\circ(z)=\sum_j\big(s_j\log\Lambda_j(z)+t_j\log\Lambda_j(-z)\big)=\sum_j\big(\sigma_jE_j(z)+\delta_jO_j(z)\big)$.
\begin{enumerate}[label=(\roman*)]
\item A gauge $\alpha^\star$ is native iff $\nabla\Phi(\log\alpha^\star)=\btau$; on the balance wall, iff $\log\alpha^\star$ is a critical point of $\ip\mvec y-\Phi(y)$.
\item For real data with $s_j,t_j\ge0$, not all zero, and $\mvec>0$, $\Phi$ is convex, and a native gauge exists iff $\btau$ lies in the relative interior of the Newton sum $\mathcal N_F=\sum_j\big(s_jP_j+t_jP^\ast_j\big)$, with $P_j$ the base polytope of table $j$ (Theorem~\ref{thm:newton}) and $P^\ast_j$ that of its dual; it is then unique modulo the space of directions along which $\Phi$ is affine.
\item For such data $\mathcal N_F=\{x:\ x_{\cE}=g(\cE),\ x_S\ge g(S)\ \forall S\}$ with the set function $g=\sum_j(s_jd_j+t_jd^\ast_j)$. For a connected family that space is $\mathbb R\one$, and the condition is $\tau_S-g(S)>0$ for every nonempty proper $S$: the walls at the target are positive. The native gauge is then the unique point, up to scale, at which $\alpha^{\btau}W(\alpha)$, constant along each ray, is largest; on the balance wall $\btau=\mvec$.
\end{enumerate}
\end{theorem}

A native gauge is a gauge at which the family uses each mark exactly as much as the target asks, and the same equation at a complex point gives the native gauges at complex points that the rays of Chapter~\ref{ch:rays} use. With real data, nonnegative exponents and positive masses, for a connected family it is unique up to scale and exists exactly when every proper wall at the target is positive (Theorem~\ref{thm:saddle}(iii)); on a compound it exists exactly when every molecule of the active family is balanced at the target with its proper walls positive there, with one free scale for each such molecule (Proposition~\ref{prop:compoundgauge}). It is a point at which $\alpha^{\btau}W(\alpha)$ is largest, the saddle of the potential, and it is the measuring stick of every computation on the family.

% Reference 20: prop:compoundgauge; source line 1569 (weights factoring across the molecules).
% Proved in the last part: prop:compoundgauge.
\begin{proposition}[The native gauges of a compound]\label{prop:compoundgauge}
Let $F$ have real data, $s_j,t_j\ge0$ not all zero, $\mvec>0$, and target $\btau$, and let $\cE_1,\dots,\cE_q$ be the molecules of its active family. A native gauge exists exactly when every $\cE_i$ is balanced at the target, $\tau_{\cE_i}=\sum_j\big(s_jr_j(\cE_i)+t_jr^\ast_j(\cE_i)\big)$, and every proper wall of every $\cE_i$ is positive there, $\tau_S>\sum_j\big(s_jd_j(S)+t_jd^\ast_j(S)\big)$ for $\varnothing\neq S\subsetneq\cE_i$. The native gauges are then the gauges $\beta$ with $\beta/\alpha^\star$ constant on every $\cE_i$, $\alpha^\star$ any one of them: one free scale for each molecule. When the weights factor across the $\cE_i$, the native gauges are the products of native gauges of the molecules. A mark that is a loop or a coloop in every table is a unit, whose letters are monomials.
\end{proposition}

% Reference 20: prop:centrefamily; source line 1587
\begin{proposition}[The centre of an arbitrary family]\label{prop:centrefamily}
Take tables on $\cE$, $s_j=t_j=\sigma_j\ge0$ with $\sigma=\sum_j\sigma_j>0$, and $\mvec=\sigma\one$. Then $\delta_j=0$, $\bmu=0$, the centred potential $\Phi^\circ(z)=\sum_j\sigma_jE_j(z)$ is even, and:
\begin{enumerate}[label=(\roman*)]
\item the gauge $\one$ is native, unique modulo one scale for each molecule of the active family (Proposition~\ref{prop:compoundgauge});
\item the walls are $\kappa_S=\sum_j\sigma_jc_j(S)$, positive on every proper nonempty $S$ exactly when the active family is connected;
\item every odd derivative of $\Phi^\circ$ vanishes at $z=0$, so every odd derivative of order at least three of the potential $\Phi$ vanishes at $y=0$.
\end{enumerate}
\end{proposition}

The family at the centre, every $\delta_j=0$ (Definition~\ref{def:realization}), with the masses $\sigma\one$ has its offsets zero. It exists for every family, its native gauge is the table's gauge, and its walls are the combinations $\kappa_S=\sum_j\sigma_jc_j(S)$ of the widths; it is where the family is read before any lean is asked of it.

% Reference 20: prop:prescribe; source line 1599
\begin{proposition}[Prescribing a gauge]\label{prop:prescribe}
For any positive gauge $a$ and positive exponent pairs, the usage $\tau_e=\sum_j\big(s_jR_{j,e}(a)+t_jR^\ast_{j,e}(a)\big)$ has total $\sum_j(s_jr_j+t_jr^\ast_j)$, and every mass vector on the ray $c\btau$, $c>0$, makes $a$ native. Fixing values $K^0_j>0$ and $\Psi^0_j>0$ of the letters at $a$, the constraints
\[
\sum_T\theta_{j,T}a^T=K^0_j,\qquad\sum_T\theta_{j,T}a^{\cE\setminus T}=\Psi^0_j,\qquad\sum_{j,T}\theta_{j,T}\Big[\frac{s_ja^T}{K^0_j}\one_T+\frac{t_ja^{\cE\setminus T}}{\Psi^0_j}\one_{\cE\setminus T}\Big]=\btau
\]
are linear in the weights, and every positive deformation in their homogeneous kernel preserves them exactly.
\end{proposition}


\part{Classification}

\chapter{Realizations and the still}\label{ch:still}

\section{Realizations}

% Reference 20: def:realization; source line 1617
\begin{definition}[Realization, lean, carrier]\label{def:realization}
Let a family carry tables $j$ on its marks $\cE$, $n=|\cE|$, its potential read as the pair $(\Phi_{\mathbf E},\Phi_{\mathbf O})$ of Definition~\ref{def:potential}. A \emph{realization} of the family is $\Upsilon=\big((s_j,t_j)_{j};\,\alpha\big)$: exponents $s_j,t_j\ge0$ with $s_j+t_j>0$, and a gauge $\alpha\in\mathbb R^{\cE}_{>0}$. Its pair coordinates are the strengths $\sigma_j=\frac12(s_j+t_j)>0$ and the leans $\delta_j=\frac12(s_j-t_j)$ of Definition~\ref{def:family}, $|\delta_j|\le\sigma_j$, and in them the realization is written
\[
\Upsilon_{\boldsymbol\sigma,\boldsymbol\delta;\alpha}=\big((\sigma_j+\delta_j,\ \sigma_j-\delta_j)_j;\ \alpha\big).
\]
$\boldsymbol\delta=(\delta_j)_j$ is the \emph{lean}, and the \emph{centre} is $\boldsymbol\delta=0$; at the table's gauge the realization at the centre has the masses $\sigma\one$ (Proposition~\ref{prop:centrefamily}), and a family at the centre is taken with these masses. Table $j$ has a \emph{basis lean} when $\delta_j>0$, a \emph{cobasis lean} when $\delta_j<0$, and \emph{rides the line} when $\delta_j=0$; its lean is \emph{full} when $|\delta_j|=\sigma_j$ (then the factor of table $j$ in the family word is $K_j^{-s_j}$ alone, or $\Psi_j^{-t_j}$ alone) and \emph{partial} when $0<|\delta_j|<\sigma_j$. The family has a basis lean when every $\delta_j\ge0$ and some $\delta_j>0$, a cobasis lean when every $\delta_j\le0$ and some $\delta_j<0$, and otherwise, off the centre, a mixed lean. The realization is \emph{at one pair} $(\sigma,\delta)$, $\sigma=\sum_j\sigma_j$ and $\delta=\sum_j\delta_j$, when its pairs lie on one ray, $(\sigma_j,\delta_j)=\frac{\sigma_j}{\sigma}(\sigma,\delta)$ for every $j$; its \emph{lean sign} is the sign of $\delta$. A family with one table is at one pair. The \emph{masses} of $\Upsilon$ are its usage,
\[
\tau_e(\Upsilon)=\sum_j\Big(s_j\,\frac{\alpha_e\partial_eK_j}{K_j}(\alpha)+t_j\,\frac{\alpha_e\partial_e\Psi_j}{\Psi_j}(\alpha)\Big),
\]
for which $\alpha$ is the native gauge (Proposition~\ref{prop:prescribe}). The \emph{carrier} of $\Upsilon$ is $(\alpha;\,K_j,\Psi_j;\,\boldsymbol\tau;\,s_j,t_j)$. A gauge is written as a deviation from the table's gauge, $\alpha=e^{y_0}$. Inverting the gauge and swapping the two exponents of every table, $\big((s_j,t_j)_j;\alpha\big)\mapsto\big((t_j,s_j)_j;\alpha^{-1}\big)$, sends the offsets and the lean to $-\bmu$ and $-\boldsymbol\delta$, the change of data of Theorem~\ref{thm:cremona}; so \emph{Cremona} swaps basis and cobasis leans, and the table's gauge is its fixed point up to scale (Proposition~\ref{prop:identity}).
\end{definition}

A realization is read off its exponents: the masses are the usage at its gauge, so the gauge is native by definition, and a realization sits on a wall only along a union of molecules, so that a realization of a molecule never does (Proposition~\ref{prop:wallsrealization}). A molecule is read at all of its realizations, and a particular integral is one of them, fixed by its exponents and its native gauge.

% Reference 20: prop:wallsrealization; source line 1576 (a connected family).
% Proved in the last part: prop:wallsrealization.
\begin{proposition}[A realization sits on a wall only along its molecules]\label{prop:wallsrealization}
Let a family carry real exponents $s_j,t_j\ge0$, $s_j+t_j>0$, and a gauge $\alpha$, with masses its usage $\mvec=\sum_j\big(s_jR_j(\alpha)+t_jR^\ast_j(\alpha)\big)$, that is, a realization (Definition~\ref{def:realization}). With $N_S=|T\cap S|$,
\[
\kappa_S=\sum_j\Big(s_j\big(\oE_{T\sim\mathbb P_{j,\alpha}}[N_S]-d_j(S)\big)+t_j\big(r_j(S)-\oE_{T\sim\mathbb P^\ast_{j,\alpha}}[N_S]\big)\Big)\ \ge\ 0\qquad(\varnothing\neq S\subseteq\cE),
\]
with equality exactly when $S$ is a union of molecules (Definition~\ref{def:molecule}). In particular $\kappa_{\cE}=0$, and every proper wall of a connected family is positive.
\end{proposition}

% Reference 20: prop:massesgauge; source line 1629 (at the table's gauge).
% Proved in the last part: prop:massesgauge.
\begin{proposition}[Masses at a gauge]\label{prop:massesgauge}
At $\alpha=e^{y_0}$, in this statement write $R^\pm_{j,e}=\frac12\big(R_{j,e}(e^{y_0})\pm R_{j,e}(e^{-y_0})\big)$; then
\[
\tau_e=\sum_j\big(s_jR_{j,e}(e^{y_0})+t_j(1-R_{j,e}(e^{-y_0}))\big)=\sum_j\big(\sigma_j(1+2R^-_{j,e})+\delta_j(2R^+_{j,e}-1)\big),
\]
and at $\alpha=\one$, where $R^-_{j,e}=0$ and $R^+_{j,e}=R_{j,e}(\one)$, $\tau_e=\sum_j\big(\sigma_j+\delta_j(2R_{j,e}(\one)-1)\big)$.
\end{proposition}

\section{The potential and the still}

% Reference 20: def:still; source line 1656
\begin{definition}[Still, rungs, tallies, line functions]\label{def:still}
The \emph{still} of $\Upsilon$ is the sequence of symmetric tensors
\[
\kappa_\Upsilon(e_1,\dots,e_k)=2^k\,\partial_{y_{e_1}}\cdots\partial_{y_{e_k}}\Phi(\log\alpha),\qquad k\ge2 ,
\]
the Taylor tensor of $\Phi$ in the signed coordinates $z=\frac12y$ from order two on; its entries of order $k$ form \emph{rung} $k$. \emph{Rung one} is the gradient $\partial_{y_e}\Phi(\log\alpha)=\tau_e(\Upsilon)$, the masses (Theorem~\ref{thm:saddle}), taken without the factor $2$. Each part has its own still, $\kappa^{\mathbf E}_\Upsilon$ and $\kappa^{\mathbf O}_\Upsilon$, with $\kappa_\Upsilon=\kappa^{\mathbf E}_\Upsilon+\kappa^{\mathbf O}_\Upsilon$. The \emph{tally} of a mark $e$ at rung $k$ is the multiset $\operatorname{tally}^{\mathbf X}_k(e)=\{\!\{\kappa^{\mathbf X}_\Upsilon(e,w):\ w\in\cE^{k-1}\}\!\}$. The \emph{self-line} and \emph{partner-line functions} of $e$ are
\[
\varphi^{\mathbf X}_e(v)=\partial_{y_e}\Phi_{\mathbf X}\big(\log\alpha+v\one_e\big),\qquad \varphi^{\mathbf X}_{e,g}(v)=\partial_{y_e}\Phi_{\mathbf X}\big(\log\alpha+v\one_g\big)\quad(g\neq e).
\]
\end{definition}

The still is the Taylor tensor of the potential at the realization from rung two on, one object, recording how the marks covary under the two laws; beside it, rung one is the masses. The letters are affine in every variable, so they have finitely many nontrivial derivatives, and every rung is read from those through the logarithm; the readings of Chapter~\ref{ch:reading} take the potential as one function, of which the still is the expansion.

% Reference 20: prop:compound; source line 1911
% Proved in the last part: prop:compound (the last sentence).
\begin{proposition}[The potential of a compound splits by molecules]\label{prop:compound}
Let the weights factor across the molecules. Then $\Phi_{\mathbf X}(y)=\sum_M\Phi_{\mathbf X,M}(y_{\cE_M})$ with $\Phi_{\mathbf X,M}$ the potential of the restriction to $M$, and every entry of the still on a string that meets two molecules is $0$. A mark that is a loop or a coloop in every table is a unit whose letters are monomials (Proposition~\ref{prop:compoundgauge}): every entry of the still on a string through it is $0$, and its rung one is its mass.
\end{proposition}

% Reference 20: prop:deviation; source line 1699
\begin{proposition}[A deviation reweights the bases]\label{prop:deviation}
At the realization $\Upsilon=\big((s_j,t_j)_j;e^{y_0}\big)$, for $k\ge2$ and every string $w$ of length $k$,
\[
\kappa_\Upsilon(w)=2^k\sum_j\Big(s_j\,\kappa_{j,e^{y_0}}(w)+t_j(-1)^k\kappa_{j,e^{-y_0}}(w)\Big).
\]
In this statement write $\kappa^\pm_j(w)=\frac12\big(\kappa_{j,e^{y_0}}(w)\pm\kappa_{j,e^{-y_0}}(w)\big)$; then
\[
\kappa_\Upsilon(w)=\begin{cases}2^{k+1}\sum_j\big(\sigma_j\kappa^+_j(w)+\delta_j\kappa^-_j(w)\big),&k\text{ even},\\[2pt] 2^{k+1}\sum_j\big(\sigma_j\kappa^-_j(w)+\delta_j\kappa^+_j(w)\big),&k\text{ odd},\end{cases}
\]
and the cumulants at the deviation are read from those at the table's gauge, $\kappa_{j,e^{y_0}}(w)=\sum_{i\ge0}\frac1{i!}\sum_{u\in\cE^i}\kappa_{j,\one}(w,u)\prod_{l=1}^iy_{0,u_l}$, for $y_0$ in the domain of convergence of the Taylor series of $y\mapsto\log K_j(e^y)$ at $0$.
\end{proposition}

% Reference 20: prop:stillgauge; source line 1681
\begin{proposition}[The still at the table's gauge]\label{prop:stillgauge}
At every realization $\Upsilon=\big((s_j,t_j)_j;\one\big)$, for $k\ge2$ and every string $w$ of length $k$,
\[
\kappa_\Upsilon(w)=2^k\sum_j\big(s_j+(-1)^kt_j\big)\kappa_{j,\one}(w)=\begin{cases}2^{k+1}\sum_j\sigma_j\,\kappa_{j,\one}(w),&k\text{ even},\\[2pt] 2^{k+1}\sum_j\delta_j\,\kappa_{j,\one}(w),&k\text{ odd}.\end{cases}
\]
This is Proposition~\ref{prop:deviation} at $y_0=0$. At the table's gauge $\Phi_{\mathbf E}$ is even in $y$ and $\Phi_{\mathbf O}$ is odd, each up to a linear term; so the even rungs are those of $\kappa^{\mathbf E}_\Upsilon$, independent of the lean, and the odd rungs are those of $\kappa^{\mathbf O}_\Upsilon$, linear in the lean and zero at the centre.
\end{proposition}

% Reference 20: prop:balanced; source line 1692
\begin{proposition}[Balanced tallies]\label{prop:balanced}
At every realization $\Upsilon$, for each part, every rung $k\ge2$ and every mark $e$, the tally of $e$ sums to $0$: $\sum_{w\in\cE^{k-1}}\kappa^{\mathbf X}_\Upsilon(e,w)=0$.
\end{proposition}

% Reference 20: prop:channels; source line 1710
\begin{proposition}[The two occupations of a mark]\label{prop:markoccupations}
For a mark $e$ of table $j$, neither a loop nor a coloop, with room $R=\cE-e$, let $K_{j,/e}(\alpha_R)$ and $K_{j,\setminus e}(\alpha_R)$ be the basis polynomials of its contraction and deletion at $e$ (Definition~\ref{def:delcon}), and
\[
\mathrm D_{j,e}(y_R)=\log K_{j,/e}(e^{y_R})-\log K_{j,\setminus e}(e^{y_R}),\qquad \mathrm D^\pm_{j,e}(y_R)=\tfrac12\big(\mathrm D_{j,e}(y_R)\pm\mathrm D_{j,e}(-y_R)\big).
\]
At the gauge $e^y$ the occupation and the dual occupation of $e$ (Theorem~\ref{thm:innerMGF}) are
\begin{gather*}
\frac{\alpha_e\partial_eK_j}{K_j}=\frac1{1+e^{-\lambda_{j,e}-\eta_{j,e}}},\qquad \frac{\alpha_e\partial_e\Psi_j}{\Psi_j}=\frac1{1+e^{-\lambda_{j,e}+\eta_{j,e}}},\\
\lambda_{j,e}=y_e+\mathrm D^-_{j,e}(y_R),\qquad \eta_{j,e}=\mathrm D^+_{j,e}(y_R),
\end{gather*}
and at the table's gauge $\lambda_{j,e}=0$: the two occupations are $(1+e^{-\eta_{j,e}})^{-1}$ and $(1+e^{\eta_{j,e}})^{-1}$.
\end{proposition}

% Reference 20: prop:occupations; source line 1725 (at the table's gauge).
% Proved in the last part: prop:occupations.
\begin{proposition}[Occupations at a gauge]\label{prop:occupations}
At $\alpha=e^{y_0}$, with $R^\pm_{j,e}$ as in Proposition~\ref{prop:massesgauge}, $\varphi^{\mathbf E}_e(0)=\sum_j\sigma_j(1+2R^-_{j,e})$ and $\varphi^{\mathbf O}_e(0)=\sum_j\delta_j(2R^+_{j,e}-1)$. At $\alpha=\one$, $\varphi^{\mathbf E}_e(0)=\sum_j\sigma_j$ at every mark and $\varphi^{\mathbf O}_e(0)=\sum_j\delta_j(2R_{j,e}(\one)-1)$.
\end{proposition}


\chapter{The reading, classes, atoms}\label{ch:reading}

\section{Views and the reading}

% Reference 20: def:view; source line 1737
\begin{definition}[Views; the reading; classes; atoms]\label{def:view}
Let $M$ be a molecule with $n$ marks carrying tables $j$, and $\Upsilon$ a realization of $M$ (Definition~\ref{def:realization}). The room of a mark $e$ is $\cE\setminus\{e\}$ (Chapter~\ref{ch:bunch} calls it the room of the bunch $\{e\}$), and an \emph{order} of a set is a listing of its elements. For a mark $e$ and distinct marks $\mathbf g=(g_1,\dots,g_k)$ of its room, $0\le k\le n-1$, the \emph{view} of $e$ with the partners $\mathbf g$ is the pair $(\Phi^{\mathbf E}_{e,\mathbf g},\Phi^{\mathbf O}_{e,\mathbf g})$,
\[
\Phi^{\mathbf X}_{e,\mathbf g}(v_0,v_1,\dots,v_k)=\Phi_{\mathbf X}\Big(\log\alpha+v_0\one_e+\sum_{i=1}^kv_i\one_{g_i}\Big),
\]
the potential on the scales of $e$ and its partners with every other mark at the gauge; the variable $v_i$ sits at \emph{position} $i$. The \emph{reading of $e$ with $k$ partners} is the multiset of its views over all $k$-tuples of distinct partners,
\[
\operatorname{read}^{(k)}_\Upsilon(e)=\{\!\{(\Phi^{\mathbf E}_{e,\mathbf g},\Phi^{\mathbf O}_{e,\mathbf g}):\ \mathbf g\in(\cE\setminus e)^{(k)}\}\!\}:
\]
the functions are kept and the names of the partners are not. With $k=n-1$, $\mathbf g$ is an order of the room and the pair $(\Phi^{\mathbf E}_{e,\mathbf g},\Phi^{\mathbf O}_{e,\mathbf g})$ is the potential \emph{written from $e$} in that order; $\operatorname{read}^{(n-1)}_\Upsilon(e)$ is the \emph{reading} of $e$. The \emph{classes} of $\Upsilon$ are the level sets of $e\mapsto\operatorname{read}^{(n-1)}_\Upsilon(e)$; they form the partition $\operatorname{Cl}(\Upsilon)$. The realization is \emph{atomic}, an \emph{atom}, when $\operatorname{Cl}(\Upsilon)=\{\cE\}$. For a family whose weights factor across its molecules, the classes of a realization are, by definition, the classes of the realizations of its molecules, each read on its own marks. With the gauge $\alpha$ and $\boldsymbol\sigma$ fixed, the partition at the lean $\boldsymbol\delta$ is $\operatorname{Cl}(\Upsilon_{\boldsymbol\sigma,\boldsymbol\delta;\alpha})$.
\end{definition}

The reading of a mark is the whole potential seen from that mark, its partners unnamed, and the reading is a function on the marks whose values are functions. Classes are its level sets, and an atom is a realization at which it is constant, so that the integral cannot tell its marks apart at any order.

% Reference 20: prop:viewform; source line 1749 (the inclusion laws at the table's gauge).
% Proved in the last part: prop:viewform.
\begin{proposition}[A view through the patterns]\label{prop:viewform}
For table $j$ and positions $0,\dots,k$ held by $e,g_1,\dots,g_k$, let the pattern of a basis be the set $p\subseteq\{0,\dots,k\}$ of positions it holds (its trace $T\cap\{e,g_1,\dots,g_k\}$ read in positions, the pattern of the basis on the bunch $\{e,g_1,\dots,g_k\}$ in Chapter~\ref{ch:bunch}), $\bar p$ its complement, and
\[
K_{j;e,\mathbf g}(x)=\sum_p\vartheta^K_{j,p}\,x^p,\qquad \Psi_{j;e,\mathbf g}(x)=\sum_p\vartheta^\Psi_{j,p}\,x^{\bar p},\qquad \vartheta^K_{j,p}=\sum_{T:\,p}\theta_{j,T}\alpha^T,\quad \vartheta^\Psi_{j,p}=\sum_{T:\,p}\theta_{j,T}\alpha^{\cE\setminus T},
\]
sums over the bases of pattern $p$. With $x_i=e^{v_i}$,
\[
\Phi^{\mathbf E}_{e,\mathbf g}=\sum_j\sigma_j\log\big(K_{j;e,\mathbf g}\Psi_{j;e,\mathbf g}\big),\qquad \Phi^{\mathbf O}_{e,\mathbf g}=\sum_j\delta_j\log\big(K_{j;e,\mathbf g}/\Psi_{j;e,\mathbf g}\big).
\]
At every gauge
\begin{gather*}
K_{j;e,\mathbf g}(\one+x)=K_j(\alpha)\sum_{A\subseteq\{0,\dots,k\}}\mathbb P_{j,\alpha}\big(p(T)\supseteq A\big)\,x^A,\\
\Psi_{j;e,\mathbf g}(\one+x)=\Psi_j(\alpha)\sum_{A\subseteq\{0,\dots,k\}}\mathbb P^\ast_{j,\alpha}\big(p(T)\cap A=\varnothing\big)\,x^A,
\end{gather*}
with $p(T)$ the pattern of $T$, so the view with $k$ partners is read from the joint inclusion laws of $e,g_1,\dots,g_k$ under the basis and cobasis laws. At the table's gauge $\vartheta^K_{j,p}=\vartheta^\Psi_{j,p}=\vartheta_{j,p}$, the weight of the bases of pattern $p$, both laws are the weight law $\mathbb P_{j,\one}$, and $\Psi_{j;e,\mathbf g}(x)=x^{\{0,\dots,k\}}K_{j;e,\mathbf g}(x^{-1})$. For one table, two views of one realization agree at the centre exactly when the polynomials $K_{e,\mathbf g}\Psi_{e,\mathbf g}$ agree, and at a lean exactly when $K_{e,\mathbf g}$ and $\Psi_{e,\mathbf g}$ both agree; at the table's gauge $K_{e,\mathbf g}$ alone decides at a lean.
\end{proposition}

% Reference 20: prop:refine; source line 1768
\begin{proposition}[More partners refine; the tallies are inside]\label{prop:refine}
\begin{enumerate}[label=(\roman*)]
\item $\operatorname{read}^{(k)}_\Upsilon(e)$ determines $\operatorname{read}^{(k')}_\Upsilon(e)$ for every $k'\le k$; so the level sets of the reading with $k$ partners refine those with $k'$.
\item For $k\ge1$, $\operatorname{read}^{(k)}_\Upsilon(e)$ determines the self-line functions $\varphi^{\mathbf X}_e$, the multisets $\{\!\{\varphi^{\mathbf X}_{e,g}\}\!\}_{g\neq e}$ of partner-line functions, and for every rung the multiset of the entries $\kappa^{\mathbf X}_\Upsilon(e,w)$ over the strings $w$ of marks with at most $k$ distinct marks other than $e$; in particular the tallies of $e$ at the rungs $2,\dots,k+1$.
\item The reading determines the tallies of $e$ at every rung $k\ge2$, its self-line and partner-line functions for both parts, and its mass $\tau_e(\Upsilon)$.
\end{enumerate}
\end{proposition}

% Reference 20: prop:onewriting; source line 1779
\begin{proposition}[One writing determines the reading]\label{prop:onewriting}
The reading of $e$ is determined by any one writing of the potential from $e$. So the reading is constant on a set $C$ of marks exactly when one pair of functions is the potential written from every mark of $C$: there are orders $\mathbf g_e$ of the rooms, $e\in C$, with $\Phi^{\mathbf X}_{e,\mathbf g_e}=\Phi^{\mathbf X}_{f,\mathbf g_f}$ for all $e,f\in C$ and both parts. In particular $\Upsilon$ is an atom exactly when one pair of functions is the potential written from every mark.
\end{proposition}

\section{Identifying writings}

% Reference 20: thm:onetable; source line 1788
\begin{theorem}[One table]\label{thm:onetable}
Let $M$ carry one table, $n\ge2$, at the table's gauge, and let $C$ be a set of marks. The reading is constant on $C$
\begin{enumerate}[label=(\roman*)]
\item at a lean, exactly when one polynomial is $K$ written from every mark of $C$;
\item at the centre, exactly when one polynomial is, from every mark of $C$, either $K$ or $\Psi$ written from it.
\end{enumerate}
Here $K$ written from $e$ in the order $\mathbf g$ is $K_{e,\mathbf g}$ with $k=n-1$, a polynomial in the position variables.
\end{theorem}

% Reference 20: thm:poles; source line 1800
\begin{theorem}[Several tables: the potential is read factor by factor]\label{thm:poles}
At a realization $\Upsilon$ write every $K_j(\alpha\circ x)$ and $\Psi_j(\alpha\circ x)$ as a positive constant times a product of irreducible polynomials in $x$, each positive on the orthant and equal to $1$ at $x=\one$. Give each distinct such polynomial $L$ the exponents
\[
w_{\mathbf E}(L)=\sum\sigma_j>0,\qquad w_{\mathbf O}(L)=\sum\pm\delta_j,
\]
summed over its occurrences with multiplicity, with $+$ in a $K_j$ and $-$ in a $\Psi_j$. Then $\Phi_{\mathbf X}(\log\alpha+\log x)=\Phi_{\mathbf X}(\log\alpha)+\sum_Lw_{\mathbf X}(L)\log L(x)$, and two pairs $(\Phi_{\mathbf E},\Phi_{\mathbf O})$ of this form, in the same variables, are equal exactly when they have the same polynomials with the same pairs of exponents $(w_{\mathbf E}(L),w_{\mathbf O}(L))$. So the reading is constant on $C$ exactly when one set of polynomials with exponents $\{(L,w_{\mathbf E}(L),w_{\mathbf O}(L))\}$, with the variables placed in positions, is the family written from every mark of $C$. For tables whose weights factor across their molecules the polynomials $L$ are the normalized Foreman polynomials of their restrictions to those molecules.
\end{theorem}

% Reference 20: cor:blocktilts; source line 1811
\begin{corollary}[Moving blocks of partners reads alike]\label{cor:blocktilts}
Let $\mathbf g,\mathbf g'$ be orders of the rooms of $e$ and $f$ with $(\Phi^{\mathbf E}_{e,\mathbf g},\Phi^{\mathbf O}_{e,\mathbf g})=(\Phi^{\mathbf E}_{f,\mathbf g'},\Phi^{\mathbf O}_{f,\mathbf g'})$, and let $\beta$ be the bijection of $\cE$ with $\beta(e)=f$ and $\beta(g_i)=g'_i$. Then for all disjoint sets $S_1,\dots,S_m$ of the room of $e$,
\[
\Phi_{\mathbf X}\Big(\log\alpha+v_0\one_e+\sum_{i=1}^mv_i\one_{S_i}\Big)=\Phi_{\mathbf X}\Big(\log\alpha+v_0\one_f+\sum_{i=1}^mv_i\one_{\beta(S_i)}\Big).
\]
So on a class every move of the gauge that shifts a mark and several sets of partners, each by one amount, reads alike: for every partition $S_1,\dots,S_m$ of the room, the potential moved by $v_0$ on the mark and by $v_i$ on $S_i$ reads alike from every mark of the class.
\end{corollary}

% Reference 20: prop:computeclasses; source line 1822 (at the table's gauge).
% Proved in the last part: prop:computeclasses.
\begin{proposition}[Computing the classes]\label{prop:computeclasses}
Read every mark with $k=0,1,2,\dots$ partners and group the marks by value. The level sets only refine as $k$ grows (Proposition~\ref{prop:refine}). A level set is a class as soon as one pair of functions is written from each of its marks (Proposition~\ref{prop:onewriting}); the orders are built one partner at a time, keeping every view equal to the restriction of that one function, which loses nothing since a view with fewer partners is a restriction of a writing. Where no one function serves a whole level set, the marks from which a given function can be written form one class and the rest are read again. The computation ends by $k=n-1$, and at a gauge $\alpha$ every step is an identity of polynomials with coefficients in the field generated by the weights and the coordinates of $\alpha$; at the table's gauge that is the field generated by the weights, rational for rational weights and integers for lists (Proposition~\ref{prop:viewform}, Theorem~\ref{thm:poles}).
\end{proposition}

\section{What a class holds}

% Reference 20: prop:massclass; source line 1831 (the total mass).
% Proved in the last part: prop:massclass.
\begin{proposition}[Masses and occupations on a class]\label{prop:massclass}
At every realization $\Upsilon$ the masses $\tau_e(\Upsilon)$ and each of their two parts $\varphi^{\mathbf E}_e(0)$ and $\varphi^{\mathbf O}_e(0)$ (Proposition~\ref{prop:occupations}) are constant on every class. At the table's gauge:
\begin{enumerate}[label=(\roman*)]
\item $\sum_j\delta_j(2R_{j,e}(\one)-1)$ is constant on every class, and $\kappa^{\mathbf E}_\Upsilon(e,e)=8\sum_j\sigma_jR_{j,e}(\one)(1-R_{j,e}(\one))$;
\item with one table, the occupation $R_e(\one)$ is constant on every class at a lean, and at the centre two marks $e,f$ of one class have $R_f(\one)=R_e(\one)$ or $R_f(\one)=1-R_e(\one)$.
\end{enumerate}
\end{proposition}


\chapter{Lean, gauge and label}\label{ch:lean}

\section{The lean}

% Reference 20: thm:leansplits; source line 1846
\begin{theorem}[The lean only splits]\label{thm:leansplits}
At every gauge and every $\boldsymbol\sigma$, the partition at a lean $\boldsymbol\delta$ refines the partition at the centre. An atom at a lean is an atom at the centre, and a realization that is not an atom at the centre is an atom at no lean.
\end{theorem}

% Reference 20: prop:leanray; source line 1853
\begin{proposition}[What of the lean is read]\label{prop:leanray}
With the gauge $\alpha$ fixed, $\operatorname{Cl}(\Upsilon_{\boldsymbol\sigma,\boldsymbol\delta;\alpha})$ depends on $(\boldsymbol\sigma,\boldsymbol\delta)$ only through the rays $\{c\boldsymbol\sigma\}$ and $\{c'\boldsymbol\delta\}$, $c,c'>0$, and $\operatorname{Cl}(\Upsilon_{\boldsymbol\sigma,-\boldsymbol\delta;\alpha})=\operatorname{Cl}(\Upsilon_{\boldsymbol\sigma,\boldsymbol\delta;\alpha})$. So at one pair, $\boldsymbol\delta=\frac\delta\sigma\boldsymbol\sigma$, there are two partitions: $\operatorname{Cl}(\Upsilon_{\boldsymbol\sigma,0;\alpha})$ at the centre, and $\operatorname{Cl}(\Upsilon_{\boldsymbol\sigma,\boldsymbol\delta;\alpha})$, the same for every $\delta\neq0$ and every $\boldsymbol\sigma$ on its ray, which refines it.
\end{proposition}

% Reference 20: prop:oneside; source line 1860
\begin{proposition}[A lean reads one side]\label{prop:oneside}
Let $\Upsilon$ be at one pair $(\sigma,\delta)$, at the table's gauge, and put
\[
\Phi_K(y)=\frac1\sigma\sum_j\sigma_j\log K_j(e^y),\qquad \Phi_\Psi(y)=\frac1\sigma\sum_j\sigma_j\log\Psi_j(e^y)=\Phi_K(-y)+\ip\one y .
\]
Then $\Phi_{\mathbf E}=\sigma(\Phi_K+\Phi_\Psi)$ and $\Phi_{\mathbf O}=\delta(\Phi_K-\Phi_\Psi)$, and:
\begin{enumerate}[label=(\roman*)]
\item at every lean the reading of a mark determines, and is determined by, the multiset of writings of $\Phi_K$ from it; so for $\delta\neq0$, $\operatorname{Cl}(\Upsilon)$ is the partition of the marks into the level sets of the multiset of writings of $\Phi_K$, and equally of $\Phi_\Psi$;
\item with the tensors $\mathcal K_k(w)=\frac1\sigma\sum_j\sigma_j\kappa_{j,\one}(w)$ on the strings of length $k$, rung $k$ is $2^{k+1}\sigma\mathcal K_k$ for even $k$ and $2^{k+1}\delta\mathcal K_k$ for odd $k$: at a basis lean every rung is a positive multiple of the cumulant sequence $(\mathcal K_k)_k$ of the bases, at a cobasis lean of the cumulant sequence $((-1)^k\mathcal K_k)_k$ of the cobases, and at the centre only the even rungs, common to the two sequences, are read.
\end{enumerate}
\end{proposition}

A lean reads the molecule with a preference: at a basis lean it reads the cumulant sequence of the bases, at a cobasis lean that of the cobases, and at the centre only what the two sequences share. The lean can only split the classes of the centre, never join them.

\section{The gauge}

% Reference 20: thm:gaugestep; source line 1877
\begin{theorem}[The gauge step]\label{thm:gaugestep}
Let $\Upsilon=((s_j,t_j)_j;\alpha)$ be a realization of a molecule. If the self-line functions $\varphi^{\mathbf E}_e$ of all marks agree up to additive constants, in particular if $\Upsilon$ is an atom, then $\alpha\propto\one$.
\end{theorem}

% Reference 20: cor:gaugescale; source line 1902
\begin{corollary}[Scale changes nothing]\label{cor:gaugescale}
At $\alpha=c\one$ every writing is the writing at $\one$ plus a constant that depends on neither the mark nor the order. So $\operatorname{Cl}(\Upsilon_{\boldsymbol\sigma,\boldsymbol\delta;c\one})=\operatorname{Cl}(\Upsilon_{\boldsymbol\sigma,\boldsymbol\delta;\one})$, and every atom is an atom at the table's gauge.
\end{corollary}

\section{Answers and the label}

% Reference 20: prop:assembly; source line 3674
\begin{proposition}[Assembly from arbitrary seed tables]\label{prop:assembly}
Let $H$ be a set of bijections of $\cE$ closed under composition and inverses that takes every mark to every other. Start from seed tables on $\cE$, and for each seed include its image under every $h\in H$, weights transported with the bases and repetitions kept. Give every image of seed $a$ the exponents $(s_a,t_a)$; let $q_a$ be the number of its images with multiplicity, $r_a$ its rank, $D=\sum_aq_a(s_ar_a+t_a(n-r_a))$, $\mvec=\frac Dn\one$, and suppose the active assembled family is connected. Then for every positive exponent pairs the gauge $\one$ is native, every proper wall is positive, and the realization $((s_a,t_a)_a;\one)$ of the assembled family is an atom.
\begin{enumerate}[label=(\roman*)]
\item At $\one$ the occupations summed over the images of seed $a$ are the same at every mark, $q_ar_a/n$, and the dual occupations summed likewise are $q_a(n-r_a)/n$; their sum with the coefficients $s_a$ and $t_a$ is $D/n$.
\item The images under all permutations of a seed with at least two bases form a connected family.
\item The seed weights may vary over any positive values before assembly, giving atoms at the one gauge $\one$.
\end{enumerate}
\end{proposition}

% Reference 20: def:label; source line 1918
\begin{definition}[Answers and the label]\label{def:label}
A molecule \emph{answers yes} when some realization of it is atomic, and \emph{no} otherwise. Let $k$ be the number of molecules of a family with at least two marks and $i$ the number of them that answer yes; units are atomic at every realization and are not counted. The family is \emph{Foreman $k$-symmetric} when $i=k$ and \emph{Foreman asymmetric} when $i<k$; an asymmetric family with $i>0$ is \emph{sub-symmetric ($i$ of $k$)}. The \emph{label} of the family is the pair $(i,k)$.
\end{definition}

% Reference 20: thm:centre; source line 1922
\begin{theorem}[The answer is read at the centre]\label{thm:centre}
A molecule answers yes if and only if some centre $\Upsilon_{\boldsymbol\sigma,0;\one}$ is atomic. With one table this holds if and only if one polynomial is, from every mark, either $K$ or $\Psi$ written from it; the condition does not involve $\sigma$.
\end{theorem}

% Reference 20: cor:label; source line 1929
\begin{corollary}[The label is read at the centres]\label{cor:label}
By Theorem~\ref{thm:centre} the label of a family is read at the centres $\Upsilon_{\boldsymbol\sigma,0;\one}$ of the table's gauge of its molecules, and a molecule carrying one table is decided by reading it once, at any one such centre.
\end{corollary}

% Reference 20: cor:twoquestions; source line 1936
\begin{corollary}[The leaning realizations at one pair]\label{cor:twoquestions}
Let a molecule be read at one pair at the table's gauge, with the strengths $\boldsymbol\sigma$. Its leaning realizations $\Upsilon_{\boldsymbol\sigma,\boldsymbol\delta;\one}$, $\boldsymbol\delta=\frac\delta\sigma\boldsymbol\sigma$, $\delta\neq0$, are atomic all together or not at all; if they are, so is the centre, and an atomic centre need not make them atomic.
\end{corollary}

Whether a molecule can be atomic and whether a particular realization of it is atomic are read differently. The first is read at the centres of the table's gauge, at any one centre when the molecule carries one table; the second at the realization itself, and a lean may split what the centre held together.

% Reference 20: lem:nodegenerate; source line 1943 (a molecule with $n\ge3$ marks).
% Proved in the last part: lem:nodegenerate.
\begin{lemma}[No degenerate pair]\label{lem:nodegenerate}
At every realization $\Upsilon$ of a family and for marks $e\neq f$, $\kappa^{\mathbf E}_\Upsilon(e,f)^2\le\kappa^{\mathbf E}_\Upsilon(e,e)\,\kappa^{\mathbf E}_\Upsilon(f,f)$, with equality exactly when $e$ or $f$ is a unit or $\{e,f\}$ is a molecule. In particular the inequality is strict for any two marks of a molecule with at least three marks.
\end{lemma}

% Reference 20: lem:equaldiag; source line 1950 (at an atom, rung two of the even part).
% Proved in the last part: lem:equaldiag.
\begin{lemma}[Equal diagonals on a class]\label{lem:equaldiag}
At every realization $\Upsilon$, for each part and every $k\ge2$, the diagonal entry $\kappa^{\mathbf X}_\Upsilon(e,\dots,e)=2^k(\varphi^{\mathbf X}_e)^{(k-1)}(0)$ of rung $k$ is the same at every mark of a class, the reading determining the self-line functions (Proposition~\ref{prop:refine}(iii)); in particular $\kappa^{\mathbf E}_\Upsilon(e,e)=4\,(\varphi^{\mathbf E}_e)'(0)$ is constant on every class, and at an atom it is the same at every mark.
\end{lemma}

\section{Maps as certificates}

% Reference 20: def:carry; source line 1959
\begin{definition}[Carrying and exchanging]\label{def:carry}
Let $F$ and $F'$ be families on $\cE$ and $\cE'$ with tables indexed by $\{1,\dots,N\}$ and $\{1,\dots,N'\}$. A bijection $g:\cE\to\cE'$ \emph{carries} $F$ to $F'$ when some bijection $h:\{1,\dots,N\}\to\{1,\dots,N'\}$ has $g(\cB_j)=\cB'_{h(j)}$, $\theta'_{h(j),g(T)}=\theta_{j,T}$ and $(\sigma'_{h(j)},\delta'_{h(j)})=(\sigma_j,\delta_j)$; it \emph{exchanges} $F$ with $F'$ when instead $T\mapsto\cE'\setminus g(T)$ maps $\cB_j$ onto $\cB'_{h(j)}$ with $\theta'_{h(j),\cE'\setminus g(T)}=\theta_{j,T}$ and the same condition on the pairs. With $F'=F$, $g$ carries or exchanges the family. Write $(g^\ast y)_e=y_{g(e)}$ for $y\in\mathbb R^{\cE'}$.
\end{definition}

% Reference 20: prop:certify; source line 1963
\begin{proposition}[What a map certifies]\label{prop:certify}
At gauges $\alpha$ of $F$ and $\alpha'$ of $F'$ with $g^\ast\log\alpha'=\log\alpha$, the tables' gauges in particular, and in this statement writing $\Phi^{\boldsymbol\delta}_{\mathbf O}$ for the odd potential at the lean $\boldsymbol\delta$: if $g$ carries $F$ to $F'$ then $\Phi_{\mathbf X}\circ g^\ast=\Phi'_{\mathbf X}$, and the reading of $g(e)$ in $F'$ equals that of $e$ in $F$; if $g$ exchanges $F$ with $F'$ then $\Phi_{\mathbf E}\circ g^\ast=\Phi'_{\mathbf E}$ and $\Phi^{-\boldsymbol\delta}_{\mathbf O}\circ g^\ast=\Phi'^{\,\boldsymbol\delta}_{\mathbf O}$, and the reading of $g(e)$ in $F'$ at $-\boldsymbol\delta$ equals that of $e$ in $F$ at $\boldsymbol\delta$. So a map carrying the family takes classes onto classes, and an exchange of the family takes the classes at $\boldsymbol\delta$ onto those at $-\boldsymbol\delta$ and the classes of the centre onto the classes of the centre.
\end{proposition}

A map is a certificate, not a definition: classes are the level sets of the reading, and a map carrying or exchanging the family is one way to exhibit that two marks read alike. With one table the writings and the maps say the same thing (Corollary~\ref{cor:certifyone}); with several tables the reading is the definition.

% Reference 20: cor:exchangecumulants; source line 1970 (a one-table molecule).
% Proved in the last part: cor:exchangecumulants.
\begin{corollary}[An exchange negates odd cumulants]\label{cor:exchangecumulants}
If $g$ exchanges a family, with $h$ the bijection of its tables (Definition~\ref{def:carry}), then under the weight laws $\kappa_{h(j),\one}(g(w))=(-1)^k\kappa_{j,\one}(w)$ for every table $j$ and every string $w$ of length $k\ge2$; so the stills at the table's gauge satisfy $\kappa_{\Upsilon_{\boldsymbol\sigma,\boldsymbol\delta;\one}}(g(w))=\kappa_{\Upsilon_{\boldsymbol\sigma,-\boldsymbol\delta;\one}}(w)$, and $g$ maps every class of the centre onto itself.
\end{corollary}

% Reference 20: cor:certifyone; source line 1977 (a one-table molecule at the table's gauge).
\begin{corollary}[With one table the writings are certificates]\label{cor:certifyone}
In a one-table molecule at the table's gauge, $e$ and $f$ lie in one class at a lean exactly when some bijection taking $e$ to $f$ carries the table, and in one class at the centre exactly when some bijection taking $e$ to $f$ carries or exchanges it.
\end{corollary}

% Reference 20: thm:acrosscentre; source line 1984
\begin{theorem}[A class across the centre]\label{thm:acrosscentre}
In a one-table molecule at the table's gauge, a class $C$ of the centre is, at every lean, one class or two classes with equally many marks. In the second case the two are $C_K$ and $C_\Psi$, the marks of $C$ from which $K$, respectively $\Psi$, is written as one and the same polynomial, and an exchange $\beta$ built from two such writings has $\beta(C)=C$ and $\beta(C_K)=C_\Psi$.
\end{theorem}

Leaning away from the centre, a class of a one-table molecule either stays whole or splits into the half from which the basis polynomial is written and the half from which the cobasis polynomial is: the lean sign decides which half the Foreman meets first (Corollary~\ref{cor:acrossbunch}).


\chapter{Restriction to a bunch}\label{ch:bunch}

\section{Patterns and the Taylor cube}

% Reference 20: def:patterns; source line 1996
\begin{definition}[Patterns and children]\label{def:patterns}
Let $(\cB,\theta)$ be a table of rank $r$ on $\cE$. A bunch (Definition~\ref{def:breta}) is here any nonempty set $B\subseteq\cE$, and $R=\cE\setminus B$ is its \emph{room}; for a family the objects of a bunch in table $j$ carry the index $j$. The \emph{patterns} of $B$ are the traces of the bases on it, $\pat(B)=\{T\cap B:\ T\in\cB\}$, and $\pat_j(B)$ for table $j$ of a family. For $p\in\pat(B)$ the \emph{signed pattern} is $\chi_p=\one_p-\one_{B\setminus p}\in\mathbb R^B$ and the \emph{signed level} is $\ell_p=|p|-\frac{|B|}2$. The \emph{child} at $p$ is the table on $R$ with the bases $\cB_{(p)}=\{T\setminus p:\ T\cap B=p\}$ and the weights $\theta_{(p),X}=\theta_{X\cup p}$; its Foreman polynomials and its centred polynomial are $K_{(p)}$, $\Psi_{(p)}$ and $\Lambda_{(p)}$.
\end{definition}

% Reference 20: lem:minor; source line 2000
\begin{lemma}[Children are tables]\label{lem:minor}
$\cB_{(p)}$ is the set of bases of the minor contracting $p$ and deleting $B\setminus p$ (Definition~\ref{def:minors}); it has the exchange property and rank $r-|p|$.
\end{lemma}

% New: proved in the last part.
\begin{lemma}[Children do not widen]\label{lem:mot-childwidth}
Let $B$ be a bunch of a table with the room $R$, $p$ a pattern of $B$, $A\subseteq R$ and $W\subseteq B$. The width of $A$ in the child at $p$ is at most the width of $A\cup W$ in the table: $c_{(p)}(A)\le c(A\cup W)$.
\end{lemma}

% Reference 20: def:communicatroid; source line 2007
\begin{definition}[Communicatroid; interchange]\label{def:communicatroid}
The \emph{communicatroid} of $B$ is the indexed family of its children $(\cB_{(p)},\theta_{(p)})_{p\in\pat(B)}$, tables on the room indexed by their patterns; its \emph{gluing} is $\{U\sqcup p:\ p\in\pat(B),\ U\in\cB_{(p)}\}$ with the weights $\theta_{(p),U}$. An indexed family $(\mathcal A_p)_p$ of tables on $R$, indexed by sets $p\subseteq B$ with $\mathcal A_p$ of rank $r-|p|$, has \emph{interchange} when, for $U\sqcup p$ and $U'\sqcup p'$ in its gluing and $x\in(U\sqcup p)\setminus(U'\sqcup p')$, there is $y\in(U'\sqcup p')\setminus(U\sqcup p)$ with
\begin{gather*}
U-x+y\in\mathcal A_p\ \ (x,y\in R),\qquad U+y\in\mathcal A_{p-x}\ \ (x\in B,\ y\in R),\\
U-x\in\mathcal A_{p+y}\ \ (x\in R,\ y\in B),\qquad U\in\mathcal A_{p-x+y}\ \ (x,y\in B):
\end{gather*}
a step of the exchange property inside a child, one level down, one level up, or along a level.
\end{definition}

% Reference 20: prop:interchange; source line 2016
\begin{proposition}[Interchange is the parent's exchange property]\label{prop:interchange}
The gluing of the communicatroid of $B$ is the table, weights included, and the communicatroid has interchange. Conversely, an indexed family of tables on $R$ with interchange is the communicatroid of $B$ in its gluing, which is a table: the children remember the parent because the parent is the glued children.
\end{proposition}

% Reference 20: prop:adjacent; source line 2023
\begin{proposition}[Adjacent siblings]\label{prop:adjacent}
If $p$ and $p+b$ are patterns of $B$, the children at $p$ and at $p+b$ are the deletion and the contraction of $b$ in one table on $R+b$, the minor contracting $p$ and deleting $B\setminus(p+b)$. So the child at $p+b$ is a quotient of the child at $p$ of rank one less: every set of room marks that is maximal for its rank in the child at $p+b$ is maximal for its rank in the child at $p$, every basis of the child at $p+b$ extends by one room mark to a basis of the child at $p$, and every basis of the child at $p$ less one suitable mark is a basis of the child at $p+b$.
\end{proposition}

% Reference 20: thm:taylorcube; source line 2030
\begin{theorem}[The Taylor cube]\label{thm:taylorcube}
For $\alpha=(\alpha_B,\alpha_R)$ and $z=(z_B,z_R)$,
\begin{gather*}
K(\alpha)=\sum_{p\in\pat(B)}\alpha_B^{\,p}K_{(p)}(\alpha_R),\qquad \Psi(\alpha)=\sum_{p\in\pat(B)}\alpha_B^{\,B\setminus p}\Psi_{(p)}(\alpha_R),\\
\Lambda(z)=\sum_{p\in\pat(B)}e^{\ip{z_B}{\chi_p}}\Lambda_{(p)}(z_R),
\end{gather*}
with $K_{(p)}=(\partial^pK)|_{\alpha_B=0}$. Along the scale of $B$, $\alpha_B=e^{2z_B}=\rho u$ with $u\in\Delta_B$, the factor of $p$ in $\Lambda$ is $e^{\ip{z_B}{\chi_p}}=\rho^{\ell_p}u^{\chi_p/2}$; the values $|p|$ fill $\{d(B),\dots,r(B)\}$, and the signed levels fill a window of width $c(B)$ centred at $\varepsilon(B)/2$. For $B=\{g\}$ this is the expansion $K=K_{\setminus g}+\alpha_gK_{/g}$, $\Psi=\alpha_g\Psi_{\setminus g}+\Psi_{/g}$ of Definition~\ref{def:delcon}.
\end{theorem}

The Taylor cube is the bunch read on its own coordinates: each pattern contributes a monomial on the bunch, at the level $|p|$, times the polynomial of its child on the room. The levels that occur fill the window of the bunch (Theorem~\ref{thm:window}), and nothing about the room is lost, since the children glued back are the table.

% Reference 20: thm:tensor; source line 2041
\begin{theorem}[The pattern tensor]\label{thm:tensor}
Under the basis law at $\alpha$, $\mathbb P_\alpha(T\cap B=p)=\alpha_B^pK_{(p)}(\alpha_R)/K(\alpha)$, and under the cobasis law $\mathbb P^\ast_\alpha(T\cap B=p)=\alpha_B^{B\setminus p}\Psi_{(p)}(\alpha_R)/\Psi(\alpha)$. The moments of the bunch are the upper sums $\mathbb P_\alpha(T\cap B\supseteq S)=\sum_{p\supseteq S}\mathbb P_\alpha(T\cap B=p)$, and the joint cumulants $\kappa_\alpha(w)$ on the strings $w$ of marks of $B$ are their image under the moment--cumulant relation. Scaling $\alpha_B$ multiplies the entry of $p$ by its own monomial, and inverting the gauge swaps the two laws pattern by pattern (Proposition~\ref{prop:splitforced}(ii) below).
\end{theorem}

% Reference 20: prop:tinkers; source line 2062
\begin{proposition}[Tinkers along a bunch; selecting patterns]\label{prop:tinkers}
For a \emph{tinker} $v=(v_p)_{p\in\pat(B)}$ of positive numbers, the weights $v_{T\cap B}\theta_T$ define the \emph{tinkered table}, with the same bases and the Foreman polynomials
\[
K_v(\alpha)=\sum_pv_p\,\alpha_B^pK_{(p)}(\alpha_R),\qquad \Psi_v(\alpha)=\sum_pv_p\,\alpha_B^{B\setminus p}\Psi_{(p)}(\alpha_R).
\]
Selecting the pattern $p$ on $B$, that is keeping the bases with $T\cap B=p$, and then $p'$ on $B'$ selects $p\cup p'$ on $B\cup B'$ when $p\cap B'=p'\cap B$, and nothing otherwise; so the selections on a table with $n$ marks are at most the $3^n$ conditions that require each mark present, absent, or nothing. At the centre every tinkered table keeps $\one$ native (Proposition~\ref{prop:centrefamily}).
\end{proposition}

\section{The Skeleton polynomials of a bunch}

% Reference 20: def:skelpoly; source line 2115
\begin{definition}[Skeleton polynomials, values and coins]\label{def:skelpoly}
At a gauge $\alpha$, with the level $N_B=|T\cap B|$ and the \emph{colevel} $N^\ast_B=|B\setminus T|$, the \emph{Skeleton polynomials} of $B$ are
\begin{gather*}
\sQ^K_B(x;\alpha)=\oE_{T\sim\mathbb P_\alpha}\big[x^{N_B-d(B)}\big]=x^{-d(B)}\frac{K(x\alpha_B,\alpha_R)}{K(\alpha)},\\
\sQ^\Psi_B(x;\alpha)=\oE_{T\sim\mathbb P^\ast_\alpha}\big[x^{N^\ast_B-d^\ast(B)}\big]=x^{-d^\ast(B)}\frac{\Psi(x\alpha_B,\alpha_R)}{\Psi(\alpha)},
\end{gather*}
and $\sQ^K_{j,B}$, $\sQ^\Psi_{j,B}$ for table $j$ of a family. Each has degree $c(B)$, positive coefficients and value $1$ at $x=1$, so its roots avoid $[0,\infty)$; writing them $-\xi_1,\dots,-\xi_{c(B)}$,
\[
\sQ^K_B(x;\alpha)=\prod_{i=1}^{c(B)}\frac{x+\xi_i}{1+\xi_i}=\prod_{i=1}^{c(B)}(1-p_i+p_ix),\qquad p_i=\frac1{1+\xi_i}.
\]
The $\xi_i$ are the \emph{Skeleton values}, their multiset $\operatorname{spec}^K_B(\alpha)\subset\mathbb C\setminus(-\infty,0]$ is the \emph{Skeleton spectrum}, and the $p_i$ are the \emph{coins}; likewise $\operatorname{spec}^\Psi_B(\alpha)$ and the coins of $\sQ^\Psi_B$. The \emph{cumulants of a coin} are $\varkappa_1(p)=p$ and $\varkappa_{k+1}(p)=p(1-p)\,\varkappa_k'(p)$.
\end{definition}

The Skeleton polynomial is the law of the level of the bunch, shifted to start at its floor: its degree is the width of the bunch, and its roots are the negatives of the Skeleton values. When the spectrum is positive the level less its floor $d(B)$ is a sum of $c(B)$ independent $\{0,1\}$-valued variables whose means are the coins.

% Reference 20: prop:sklaws; source line 2128
\begin{proposition}[Relations of the Skeleton polynomials]\label{prop:sklaws}
\begin{enumerate}[label=(\roman*)]
\item \emph{Scaling.} Scaling $\alpha_B$ by $\lambda>0$ divides the spectrum by $\lambda$:
\[
\sQ^K_B\big(x;(\lambda\alpha_B,\alpha_R)\big)=\sQ^K_B(\lambda x;\alpha)/\sQ^K_B(\lambda;\alpha);
\]
likewise for $\Psi$; scaling all of $\alpha$ changes nothing.
\item \emph{Complement.} $\sQ^K_R(x;\alpha)=x^{c(B)}\sQ^K_B(x^{-1};\alpha)$: the spectrum of the room is the reciprocal of that of the bunch, and each coin $p_i$ becomes $1-p_i$; likewise for $\Psi$.
\item \emph{Cremona.} $\sQ^\Psi_B(x;\alpha)=x^{c(B)}\sQ^K_B(x^{-1};\alpha^{-1})$; at the table's gauge, $\sQ^\Psi_B(x;\one)=x^{c(B)}\sQ^K_B(x^{-1};\one)$.
\item \emph{Coins.} $\log\sQ^K_B(e^v;\alpha)=\sum_i\log(1-p_i+p_ie^v)$; under $\mathbb P_\alpha$ the level $N_B$ has mean $d(B)+\sum_ip_i=\sum_{e\in B}R_e(\alpha)$ and $k$-th cumulant $\sum_i\varkappa_k(p_i)$ for $k\ge2$. When the spectrum is positive the coins lie in $(0,1)$ and $N_B-d(B)$ is distributed as a sum of independent $\{0,1\}$-valued variables whose means are the coins.
\end{enumerate}
\end{proposition}

% Reference 20: prop:stillcoins; source line 2145
\begin{proposition}[The still read along a bunch]\label{prop:stillcoins}
Let a family be at a gauge $\alpha$, let $p_{j,i}$ be the coins of $\sQ^K_{j,B}(\cdot;\alpha)$ and $p^\ast_{j,i}$ those of $\sQ^\Psi_{j,B}(\cdot;\alpha)$, and put $\partial_B=\sum_{e\in B}\partial_{y_e}$. Then
\begin{gather*}
\sum_j\Big(s_j\sum_ip_{j,i}+t_j\sum_ip^\ast_{j,i}\Big)=\partial_B\Phi(\log\alpha)-\sum_j\big(s_jd_j(B)+t_jd^\ast_j(B)\big),\\
\partial_B^k\Phi(\log\alpha)=\sum_j\Big(s_j\sum_i\varkappa_k(p_{j,i})+t_j\sum_i\varkappa_k(p^\ast_{j,i})\Big)\qquad(k\ge2).
\end{gather*}
At a native gauge the first total is the wall $\kappa_B$ at the masses $\btau$, and on the balance wall it is $\kappa_B$; the entries of rung $k$ of the still of $\big((s_j,t_j)_j;\alpha\big)$, summed over $B^k$, are $2^k$ times the second total.
\end{proposition}

% Reference 20: prop:coinsrank; source line 2157 (rank one over $\mathbb R$).
% Proved in the last part: prop:coinsrank.
\begin{proposition}[Coins of a tensorminental table]\label{prop:coinsrank}
In the setting of Theorem~\ref{thm:interference}, let $p^{(kl)}_i$, $1\le i\le|B|$, be the eigenvalues of $\Pi^{(kl)}_{BB}$. Then
\[
x^{d(B)}\sQ^K_B(x;\alpha)=\sum_{k,l}w_{kl}\prod_{i=1}^{|B|}\big(1-p^{(kl)}_i+p^{(kl)}_ix\big):
\]
the Skeleton polynomial is the $w$-mixture of the coin products of the components. At tensor rank one over $\mathbb F\in\{\mathbb R,\mathbb C\}$ there is one component, and $\Pi^{(11)}_{BB}$ is similar to the compression to $\mathbb F^B$ of the orthogonal projection onto the row space of $AD_\alpha^{1/2}$, $A$ the slice; its eigenvalues are $\cos^2\vartheta_i$ for the principal angles $\vartheta_i$ between $\mathbb F^B$ and that row space. Of these, $d(B)$ are $1$ and $d^\ast(B)$ are $0$, and the $c(B)$ coins are the values $\cos^2\vartheta_i\in(0,1)$, with the positive Skeleton values $\tan^2\vartheta_i$.
\end{proposition}

% Reference 20: def:pairsk; source line 2179
\begin{definition}[The pair of Skeleton functions]\label{def:pairsk}
The \emph{even} and \emph{odd Skeleton functions} of $B$ are $\sQ^{\mathbf E}_B(x;\alpha)=\sQ^K_B(x;\alpha)\,\sQ^\Psi_B(x;\alpha)$ and $\sQ^{\mathbf O}_B(x;\alpha)=\sQ^K_B(x;\alpha)/\sQ^\Psi_B(x;\alpha)$. So
\[
\sQ^{\mathbf E}_B(x;\alpha)=\oE_{(T,T')\sim\mathbb P_\alpha\otimes\mathbb P^\ast_\alpha}\Big[x^{|T\cap B|-d(B)+|B\setminus T'|-d^\ast(B)}\Big],
\]
and its roots are the negatives of the \emph{joint spectrum} $\operatorname{spec}^K_B(\alpha)\cup\operatorname{spec}^\Psi_B(\alpha)$.
\end{definition}

% Reference 20: prop:setscale; source line 2183
\begin{proposition}[A set on its own scale is read by its Skeleton functions]\label{prop:setscale}
For a family at a gauge $\alpha$, a set $C$ of marks and $x=e^{v}$,
\begin{gather*}
\Phi_{\mathbf E}(\log\alpha+v\one_C)-\Phi_{\mathbf E}(\log\alpha)=\sum_j\sigma_j\Big(\log\sQ^{\mathbf E}_{j,C}(x;\alpha)+\big(d_j(C)+d^\ast_j(C)\big)v\Big),\\
\Phi_{\mathbf O}(\log\alpha+v\one_C)-\Phi_{\mathbf O}(\log\alpha)=\sum_j\delta_j\Big(\log\sQ^{\mathbf O}_{j,C}(x;\alpha)+\big(d_j(C)-d^\ast_j(C)\big)v\Big).
\end{gather*}
Along the scale of a set the $K$-half of each table's joint spectrum carries the exponent $s_j=\sigma_j+\delta_j$ and the $\Psi$-half carries $t_j=\sigma_j-\delta_j$.
\end{proposition}

\section{Levels}

% Reference 20: def:levels; source line 2199
\begin{definition}[Levels]\label{def:levels}
For $d(B)\le k\le r(B)$ the \emph{level-$k$ patterns} are $\pat^{(k)}(B)=\{p\in\pat(B):\ |p|=k\}$, the \emph{level-$k$ bases} are $\cB^{(k)}_B=\{T\in\cB:\ |T\cap B|=k\}$, and the \emph{level letter} is
\[
K^{(k)}_B(\alpha)=\sum_{T\in\cB^{(k)}_B}\theta_T\alpha^T=\sum_{p\in\pat^{(k)}(B)}\alpha_B^pK_{(p)}(\alpha_R).
\]
For $d^\ast(B)\le k\le|B|-d(B)$ the \emph{colevel letter} is $\Psi^{(k)}_B(\alpha)=\sum_{T:\,|B\setminus T|=k}\theta_T\alpha^{\cE\setminus T}$. The \emph{excess} of the level $k$, of its patterns and of its bases is $k-d(B)$, and that of the colevel $k$ is $k-d^\ast(B)$. Then $K=\sum_kK^{(k)}_B$, $\Psi=\sum_k\Psi^{(k)}_B$,
\[
x^{d(B)}\sQ^K_B(x;\alpha)=\sum_kx^k\,\frac{K^{(k)}_B(\alpha)}{K(\alpha)},\qquad x^{d^\ast(B)}\sQ^\Psi_B(x;\alpha)=\sum_kx^k\,\frac{\Psi^{(k)}_B(\alpha)}{\Psi(\alpha)}.
\]
\end{definition}

% Reference 20: prop:higgs; source line 2203
\begin{proposition}[Pattern levels are lifts]\label{prop:higgs}
Let $\cB|B$ be the restriction to $B$ and $\cB/R$ the contraction of the room, a table on $B$ with rank function $r_{\cB/R}(X)=r(X\cup R)-r(R)$ and rank $d(B)$.
\begin{enumerate}[label=(\roman*)]
\item $p\subseteq B$ is a pattern if and only if $p$ is independent and $r_{\cB/R}(p)=d(B)$.
\item For $i=k-d(B)$, the independent sets $X\subseteq B$ with $|X|-r_{\cB/R}(X)\le i$ are the independent sets of a table on $B$ of rank $k$, the \emph{$i$-th lift} of $\cB/R$ toward $\cB|B$, whose bases are the level-$k$ patterns. So every pattern level $\pat^{(k)}(B)$ is a table on $B$.
\item The room sides $T\cap R$ of the level-$k$ bases are the bases of the $(r(B)-k)$-th lift of $\cB/B$ toward $\cB|R$.
\end{enumerate}
\end{proposition}

% Reference 20: def:interaction; source line 11356 (one table).
% Proved in the last part: def:interaction (the family pair).
\begin{definition}[Interaction]\label{def:interaction}
Let a table have a bunch $B$ with room $R$ and $\pat(B)=2^B$: every $p\subseteq B$ is a pattern, equivalently $c(B)=|B|$, that is $B$ is independent and some basis avoids it (Proposition~\ref{prop:higgs}(i)). The \emph{interaction} and the \emph{dual interaction} of $B$ are
\[
\iota_B(\alpha_R)=\sum_{p\subseteq B}(-1)^{|B\setminus p|}\log K_{(p)}(\alpha_R),\qquad \iota^\ast_B(\alpha_R)=\sum_{p\subseteq B}(-1)^{|p|}\log\Psi_{(p)}(\alpha_R),
\]
the second the interaction of the dual table, and for a split of $B$ into blocks $B_1,\dots,B_k$,
\[
\iota_{(B_1,\dots,B_k)}(\alpha_R)=\sum_{A\subseteq\{1,\dots,k\}}(-1)^{k-|A|}\log K_{(\cup_{a\in A}B_a)}(\alpha_R).
\]
For two marks, $\iota_{\{e,f\}}=\log\frac{K_{(\{e,f\})}K_{(\varnothing)}}{K_{(\{e\})}K_{(\{f\})}}$, the logarithm of the odds ratio of the presence of the two marks. For a family with $\pat_j(B)=2^B$ in every active table $j$, the \emph{interaction} of $B$ is the pair
\[
\iota^{\mathbf E}_B=\sum_j\sigma_j\big(\iota_{j,B}+\iota^\ast_{j,B}\big),\qquad \iota^{\mathbf O}_B=\sum_j\delta_j\big(\iota_{j,B}-\iota^\ast_{j,B}\big),\qquad \iota^{\mathbf E}_B+\iota^{\mathbf O}_B=\sum_j\big(s_j\iota_{j,B}+t_j\iota^\ast_{j,B}\big),
\]
the interactions of the even and the odd part of the potential (Definition~\ref{def:potential}).
\end{definition}

% Reference 20: thm:levels; source line 2215
\begin{theorem}[Levels are sums of children]\label{thm:levels}
$\cB^{(k)}_B=\bigsqcup_{p\in\pat^{(k)}(B)}\cB_{[p]}$, where $\cB_{[p]}=\{T\in\cB:\ T\cap B=p\}$ is the set of bases of the table on $\cE$ that has $p$ as coloops, $B\setminus p$ as loops and the child at $p$ on $R$, with the parent's weights; its basis polynomial is $\alpha_B^pK_{(p)}(\alpha_R)$. A bijection $g$ of $\cE$ that carries the table and keeps $B$ takes $\cB_{[p]}$ to $\cB_{[g(p)]}$. The extreme levels have constant children:
\[
\cB^{(d(B))}_B=\{p\sqcup X:\ p\in\cB/R,\ X\in\cB|R\},\qquad \cB^{(r(B))}_B=\{p\sqcup X:\ p\in\cB|B,\ X\in\cB/B\};
\]
for a table determinantal over $\mathbb R$ or $\mathbb C$ the weights of the extreme levels factor along these products, and with them the extreme letters, $K^{(d(B))}_B(\alpha)=K_{\cB/R}(\alpha_B)\,K_{\cB|R}(\alpha_R)$ and $K^{(r(B))}_B(\alpha)=K_{\cB|B}(\alpha_B)\,K_{\cB/B}(\alpha_R)$.
\end{theorem}

The bases of one level need not form a table: a step of the exchange property may leave the level, and it then moves to an adjacent one (Proposition~\ref{prop:interchange}). What a level keeps is its patterns, which form a table (Proposition~\ref{prop:higgs}); its bases are the glued children of those patterns, and children on adjacent levels are linked by Proposition~\ref{prop:adjacent}. The level is the coordinate the bunch keeps for itself, and the children hand the rest down to the room.

% New: proved in the last part.
\begin{proposition}[The bottom of a chain factors over its steps]\label{prop:mot-chainbottom}
Let $(\cB,\theta)$ be a table on $\cE$ and $\varnothing=S_0\subsetneq S_1\subsetneq\cdots\subsetneq S_\ell\subsetneq S_{\ell+1}=\cE$, $\ell\ge1$, a chain with the steps $Z_i=S_i\setminus S_{i-1}$, $1\le i\le\ell+1$. Some bases are at the bottom of every $S_i$, $|T\cap S_i|=d(S_i)$ for $1\le i\le\ell$. For any one of them, $T_0$, and every $i$, the traces on $Z_i$ of these bases are the bases of the child $\cB_{(T_0\setminus Z_i)}$ at the pattern $T_0\setminus Z_i$ of $\cE\setminus Z_i$ (Definition~\ref{def:patterns}), the \emph{step table at the bottom} at $Z_i$, which depends on $S_{i-1}$ and $S_i$ alone; and the bases at the bottom of every $S_i$ are exactly the unions $\bigsqcup_iX_i$ of one basis $X_i$ of each step table. Likewise the bases at the top of every $S_i$, $|T\cap S_i|=r(S_i)$, with one of them $T_1$, give the \emph{step tables at the top} $\cB_{(T_1\setminus Z_i)}$.

For a list, the terms of $K$ of the least degree in $\alpha_{S_i}$ for every $i$ at once are the product $\prod_iK_{\cB_{(T_0\setminus Z_i)}}(\alpha_{Z_i})$, and those of $\Psi$ the product $\prod_i\Psi_{\cB_{(T_1\setminus Z_i)}}(\alpha_{Z_i})$; the factors are the \emph{step letters} of $K$ and of $\Psi$, the same products holding for a table whose weights on these bases are products of one factor for each step. For the linear letter $\alpha_J$, the basis polynomial of $\binom J1$ read on $\cE$ with the marks outside $J$ as loops, the step letter is $\alpha_{J\cap Z_i}$ at the first step $Z_i$ with $J\subseteq S_i$ and $1$ at the other steps. Along a maximal chain, every step one mark, exactly one basis is at the bottom of every $S_i$, and for a list the terms of $K$ of the least degree in every $\alpha_{S_i}$ at once are one monomial with the coefficient $1$.
\end{proposition}

% Reference 20: thm:sharedwalls; source line 2248
\begin{theorem}[Shared walls]\label{thm:sharedwalls}
Let $L_i=\sum_hK_{i,h}$, $i=1,\dots,q$, each a sum of the basis polynomials $K_{i,h}$ of tables $\cB_{i,h}$ on $\cE$, with exponents $w_i>0$, and put $d_{i,h}(S)=\min_{T\in\cB_{i,h}}|T\cap S|$. Then
\[
\int_{\mathbb R^{\cE}_{>0}}\alpha^{\mvec-\one}e^{-|\alpha|}\prod_iL_i(\alpha)^{-w_i}\,d\alpha
\]
converges if and only if there are $w_{i,h}>0$ with $\sum_hw_{i,h}=w_i$ and $m_S>\sum_{i,h}w_{i,h}\,d_{i,h}(S)$ for every nonempty $S\subseteq\cE$; equivalently $\mvec\in\interior\big(\sum_iw_i\operatorname{conv}\bigcup_hP_{i,h}+\mathbb R^{\cE}_{\ge0}\big)$, with $P_{i,h}$ the base polytope of $\cB_{i,h}$.
\end{theorem}

% Reference 20: prop:onelevel; source line 2255
\begin{proposition}[The walls of one level are exact]\label{prop:onelevel}
With $P$ the base polytope, $P^{(k)}_B=\operatorname{conv}\{\one_T:\ T\in\cB^{(k)}_B\}$ and $d^{(k)}_B(S)=\min_{T\in\cB^{(k)}_B}|T\cap S|$:
\[
P^{(k)}_B=P\cap\{x_B=k\},\qquad P^{(k)}_B+\mathbb R^{\cE}_{\ge0}=\{x:\ x_S\ge d^{(k)}_B(S)\ \text{for all }S\}.
\]
So for $w>0$ the integral $\int_{\mathbb R^{\cE}_{>0}}\alpha^{\mvec-\one}e^{-|\alpha|}K^{(k)}_B(\alpha)^{-w}d\alpha$ converges exactly when $m_S>w\,d^{(k)}_B(S)$ for every nonempty $S$. The one-dimensional faces of $P^{(k)}_B$ that do not join two bases differing by a single step of the exchange property are diagonals of square two-dimensional faces of $P$, one step into $B$ and one out of it.
\end{proposition}

For several levels together the $0/1$ walls need not be exact: the facets of $\sum_kw_kP^{(k)}_B+\mathbb R^{\cE}_{\ge0}$ are normal to rays of the common refinement of the normal fans of the $P^{(k)}_B$, which can lie outside $\{0,1\}^{\cE}$, and the exact domain of several levels together is then the one of Theorem~\ref{thm:sharedwalls}.

% Reference 20: cor:levelsplit; source line 2266 (some table with a positive exponent connected).
% Proved in the last part: cor:levelsplit.
\begin{corollary}[The native level split is interior]\label{cor:levelsplit}
Let a family with real data, $s_j,t_j\ge0$ not all zero and $\mvec>0$, have a native gauge $\alpha^\star$, and give the level letters $K^{(k)}_{j,B}$ of table $j$ the exponents $s_j\,\mathbb P_{j,\alpha^\star}(N_B=k)$ and its colevel letters $\Psi^{(k)}_{j,B}$ the exponents $t_j\,\mathbb P^\ast_{j,\alpha^\star}(N^\ast_B=k)$. Then $\btau$ lies in the relative interior of the Minkowski sum of the level and colevel polytopes, each multiplied by the exponent of its letter, and every wall of these letters at the masses $\btau$,
\[
\tau_S-\sum_{j,k}\Big(s_j\,\mathbb P_{j,\alpha^\star}(N_B=k)\,d^{(k)}_{j,B}(S)+t_j\,\mathbb P^\ast_{j,\alpha^\star}(N^\ast_B=k)\,d^{\ast(k)}_{j,B}(S)\Big),
\]
$d^{(k)}_{j,B}$ the level counts of Proposition~\ref{prop:onelevel} and $d^{\ast(k)}_{j,B}(S)$ the least of $|S\setminus T|$ over the colevel-$k$ bases, is positive except on the sets $S$ on which $x_S$ is constant over every level and colevel polytope, among them $B$ and $R$. If the marks of $B$ are joined to one another, and the marks of $R$ to one another, by chains in which consecutive marks lie in a common circuit of a table with a positive exponent, as they are when such a table is connected, then the directions $v$ with $\ip vx$ constant on that Minkowski sum are spanned by $\one_B$ and $\one_R$, there are no exceptions besides these and $\varnothing,\cE$, and
\[
\int_{\Delta_B\times\Delta_R}u^{\btau|_B-\one}v^{\btau|_R-\one}\prod_{j,k}K^{(k)}_{j,B}(u,v)^{-s_j\mathbb P_{j,\alpha^\star}(N_B=k)}\Psi^{(k)}_{j,B}(u,v)^{-t_j\mathbb P^\ast_{j,\alpha^\star}(N^\ast_B=k)}\,d\omega_B(u)\,d\omega_R(v)
\]
converges.
\end{corollary}

\section{Shares and the native split}

% Reference 20: def:shares; source line 2275
\begin{definition}[Shares, loads, children of a family]\label{def:shares}
For a family and a bunch $B$, \emph{shares} are $\mathbf a_j=(a_{j,p})_{p\in\pat_j(B)}$ with $|\mathbf a_j|=s_j$ and $\mathbf b_j=(b_{j,p})_{p\in\pat_j(B)}$ with $|\mathbf b_j|=t_j$; the \emph{pair shares} are $\sigma_{j,p}=\frac{a_{j,p}+b_{j,p}}2$ and $\delta_{j,p}=\frac{a_{j,p}-b_{j,p}}2$. A $K$-share of pattern $p$ loads the marks of $p$, and a $\Psi$-share loads those of $B\setminus p$: the \emph{load} and the \emph{odd load} on $e\in B$ are
\[
\operatorname{load}_e=\sum_j\Big(\sum_{p\ni e}a_{j,p}+\sum_{p\not\ni e}b_{j,p}\Big)=\sigma+\langle\boldsymbol\delta\rangle_e,\qquad \langle\boldsymbol\delta\rangle_e=\sum_{j,p}\delta_{j,p}\chi_p(e).
\]
The \emph{children} are the family
\[
F_B(\mathbf a,\mathbf b)=\big((\cB_{j,(p)},\theta_{j,(p)});\,a_{j,p},b_{j,p}\big)_{j,p}\quad\text{on }R .
\]
\end{definition}

% Reference 20: def:split; source line 2279
\begin{definition}[Pattern probabilities; the native split]\label{def:split}
At a gauge $\beta$ the \emph{pattern probabilities} of table $j$ are $\pi_{j,p}(\beta)=\mathbb P_{j,\beta}(T\cap B=p)$ and $\pi^\ast_{j,p}(\beta)=\mathbb P^\ast_{j,\beta}(T\cap B=p)$. At a native gauge $\alpha^\star$ the \emph{native split} is the choice of shares $\mathbf a^\star$, $\mathbf b^\star$ with $a^\star_{j,p}=s_j\pi_{j,p}(\alpha^\star)$ and $b^\star_{j,p}=t_j\pi^\ast_{j,p}(\alpha^\star)$, and $F_B(\mathbf a^\star,\mathbf b^\star)$ are the children at the native split.
\end{definition}

% Reference 20: thm:hered; source line 2283
\begin{theorem}[Nativity is hereditary]\label{thm:hered}
$\alpha^\star|_R$ is a native gauge for the children $F_B(\mathbf a^\star,\mathbf b^\star)$ with target $\btau|_R$, and they lie on their own balance wall at the masses $\btau|_R$; every set $S\subseteq R$ with $c_{j,(p)}(S)=0$ for every child table is balanced, $\tau_S=\sum_{j,p}\big(a^\star_{j,p}r_{j,(p)}(S)+b^\star_{j,p}r^\ast_{j,(p)}(S)\big)$.
\end{theorem}

% Reference 20: prop:herecentre; source line 2290
\begin{proposition}[The centre is hereditary]\label{prop:herecentre}
Let $\Upsilon=((s_j,t_j)_j;\one)$ be at the table's gauge, with its masses. The native split gives the child $(j,p)$ the pair
\[
\sigma_{j,p}=\pi_{j,p}(\one)\,\sigma_j,\qquad \delta_{j,p}=\pi_{j,p}(\one)\,\delta_j,\qquad \pi_{j,p}(\one)=\mathbb P_{j,\one}(T\cap B=p)>0,
\]
and the gauge $\one$ on $R$, the children's own table's gauge, native for them. At a centre every child sits at a centre; at one pair the children are at the same pair; a full lean is handed down full.
\end{proposition}

% Reference 20: prop:splitforced; source line 2301
\begin{proposition}[The split is forced]\label{prop:splitforced}
\begin{enumerate}[label=(\roman*)]
\item For every gauge $\beta$ and every probability vector $w$ on $\pat_j(B)$ with every $w_p>0$,
\[
K_j(\beta)\ \ge\ \prod_p\Big(\frac{\beta_B^pK_{j,(p)}(\beta_R)}{w_p}\Big)^{w_p},
\]
with equality if and only if $\pi_{j,p}(\beta)=w_p$ for every $p$; likewise for $\Psi_j$ with the pattern terms $\beta_B^{B\setminus p}\Psi_{j,(p)}(\beta_R)$ and $\pi^\ast_{j,p}(\beta)$. So the shares $a_{j,p}=s_jw_p$ replace $K_j^{-s_j}$ by the product of its pattern terms exactly at the gauges where the basis law puts the proportions $w$ on the patterns; at the gauge of a realization the only such shares are its native split.
\item $\beta_B^{-p}K_{j,(p)}(\beta_R^{-1})=\beta^{-\cE}\beta_B^{B\setminus p}\Psi_{j,(p)}(\beta_R)$ and $\pi_{j,p}(\beta^{-1})=\pi^\ast_{j,p}(\beta)$: Cremona turns the $K$-term of each pattern into its $\Psi$-term, acting on each child as the child's own Cremona. At the table's gauge $\pi^\ast_{j,p}(\one)=\pi_{j,p}(\one)$: the forced shares split $s_j$ and $t_j$ in the same proportions.
\end{enumerate}
\end{proposition}

A bunch hands its room a family of children, one for each pattern of each table, and the native split is the one choice of shares under which the room inherits everything: the gauge stays native, the masses restrict, and a centre hands down centres. The inequality of (i) is why: at a gauge, the pattern terms of a letter can stand in for the letter with no loss only in the proportions in which the law at that gauge already uses them.


\chapter{The order of classes and the Foreman bunch}\label{ch:order}

\section{The rises of a class}

% Reference 20: def:classviews; source line 2322
\begin{definition}[The rises of a class]\label{def:classviews}
Let a molecule $M$ with $n$ marks be read at realizations $\Upsilon_{\boldsymbol\sigma,\boldsymbol\delta;\one}$ at the table's gauge with $\boldsymbol\sigma$ fixed, and let $C$ be a class of the realization at hand, with $m$ marks, and $R$ its room. For a function $f$ of $v=(v_0,\dots,v_k)$ analytic at $0$, $\mathring f(v)=f(v)-f(0)-\sum_iv_i\,\partial_{v_i}f(0)$ is $f$ less its value and its tangent.
\begin{enumerate}[label=(\roman*)]
\item For $0\le k\le n-m$ and distinct partners $\mathbf g=(g_1,\dots,g_k)$ in $R$, the \emph{rise of $C$ with the partners $\mathbf g$} is the increase of the potential along the class, each partner on its own scale:
\[
V^{\mathbf X}_{C,\mathbf g}(v_0,\dots,v_k)=\Phi_{\mathbf X}\Big(v_0\one_C+\sum_{i=1}^kv_i\one_{g_i}\Big)-\Phi_{\mathbf X}\Big(\sum_{i=1}^kv_i\one_{g_i}\Big).
\]
\emph{Stage} $k$ is the multiset $\mathcal V_k(C)$ of the pairs $(\mathring V^{\mathbf E}_{C,\mathbf g},\mathring V^{\mathbf O}_{C,\mathbf g})$ over the $k$-tuples of distinct marks of $R$; stage $0$ is the class on its own scale.
\item For an order $w=(c_1,\dots,c_m;r_1,\dots,r_{n-m})$ of the marks that lists those of $C$ first and then those of $R$, the \emph{rise from the class} is
\[
\overline V^{\mathbf X}_{C,w}(u_1,\dots,u_m,v_1,\dots,v_{n-m})=\Phi_{\mathbf X}\Big(\sum_iu_i\one_{c_i}+\sum_jv_j\one_{r_j}\Big)-\Phi_{\mathbf X}(0),
\]
with positions $0,\dots,n-1$ in the order of $w$, and the \emph{last stage} is the multiset $\mathcal V_\ast(C)$ of the pairs $(\overline V^{\mathbf E}_{C,w},\overline V^{\mathbf O}_{C,w})$ over the $m!\,(n-m)!$ such orders.
\end{enumerate}
$\mathcal V_k(C)$ has $(n-m)!/(n-m-k)!$ elements; $v_k=0$ turns the rise with $(g_1,\dots,g_k)$ into the rise with $(g_1,\dots,g_{k-1})$; for a single mark an element of $\mathcal V_\ast(\{e\})$ is a writing of the potential from $e$, less $\Phi_{\mathbf X}(0)$.
\end{definition}

A class is first read on its own scale, then with one partner of the room at a time, then with more, and last from the class with every mark on its own scale. Every stage before the last is taken less its value and tangent; the last stage holds the whole potential.

% Reference 20: prop:skeletonviews; source line 2338
\begin{proposition}[Every rise is read by the Skeleton functions of the class]\label{prop:skeletonviews}
With $x=e^{v_0}$ and $\alpha=e^{\sum_iv_i\one_{g_i}}$, the gauge the partners move to,
\begin{gather*}
V^{\mathbf E}_{C,\mathbf g}(v_0,\dots,v_k)=\sum_j\sigma_j\Big(\log\sQ^{\mathbf E}_{j,C}(x;\alpha)+\big(d_j(C)+d^\ast_j(C)\big)v_0\Big),\\
V^{\mathbf O}_{C,\mathbf g}(v_0,\dots,v_k)=\sum_j\delta_j\Big(\log\sQ^{\mathbf O}_{j,C}(x;\alpha)+\big(d_j(C)-d^\ast_j(C)\big)v_0\Big).
\end{gather*}
Stage $0$ reads the Skeleton functions of $C$ at the table's gauge, and stage $k$ reads them at every gauge that $k$ partners move to.
\end{proposition}

% Reference 20: prop:viewcumulants; source line 2350
\begin{proposition}[The rises hold the cumulants of the class with its partners]\label{prop:viewcumulants}
For $\mathbf a=(a_0,\dots,a_k)\in\mathbb N^{k+1}$ write, in this statement, $\kappa_j(\mathbf a)=\kappa_{j,\one}\big(N_C^{\times a_0},X_{g_1}^{\times a_1},\dots,X_{g_k}^{\times a_k}\big)$, the joint cumulant, under the weight law $\mathbb P_{j,\one}$, of the sequence of variables in which the level $N_C$ appears $a_0$ times and the indicator $X_{g_i}$ appears $a_i$ times. For $a_0\ge1$ and $|\mathbf a|\ge2$,
\[
[v^{\mathbf a}]\,V^{\mathbf E}_{C,\mathbf g}=\frac{1+(-1)^{|\mathbf a|}}{\mathbf a!}\sum_j\sigma_j\kappa_j(\mathbf a),\qquad [v^{\mathbf a}]\,V^{\mathbf O}_{C,\mathbf g}=\frac{1-(-1)^{|\mathbf a|}}{\mathbf a!}\sum_j\delta_j\kappa_j(\mathbf a),
\]
and the coefficients with $a_0=0$ vanish. By multilinearity $\kappa_j(\mathbf a)$ is the sum, over the strings $u\in C^{a_0}$, of the cumulants $\kappa_{j,\one}$ along the strings that continue $u$ with each $g_i$ taken $a_i$ times: the entries of the still on those strings, summed over the class and read table by table.
\end{proposition}

% Reference 20: cor:firststages; source line 2361
\begin{corollary}[The first two stages in coins]\label{cor:firststages}
\begin{enumerate}[label=(\roman*)]
\item For $a\ge2$, $[v_0^a]\,\mathring V^{\mathbf E}_{C,()}=\frac2{a!}\sum_j\sigma_j\kappa_{j,\one}(N_C^{\times a})$ for even $a$ and $[v_0^a]\,\mathring V^{\mathbf O}_{C,()}=\frac2{a!}\sum_j\delta_j\kappa_{j,\one}(N_C^{\times a})$ for odd $a$, the other coefficients being zero; and $\kappa_{j,\one}(N_C^{\times a})=\sum_i\varkappa_a(p_{j,i})$ over the coins $p_{j,i}$ of $\sQ^K_{j,C}(\cdot;\one)$.
\item For a partner $g$ and $q\ge1$, with $R_{j,g}(e^{\pm v_0\one_C})$ the occupation of $g$ in table $j$ at the gauge $e^{\pm v_0\one_C}$,
\begin{gather*}
[v_1^q]\,V^{\mathbf E}_{C,(g)}=\frac1{q!}\sum_j\sigma_j\Big(\varkappa_q\big(R_{j,g}(e^{v_0\one_C})\big)+(-1)^q\varkappa_q\big(R_{j,g}(e^{-v_0\one_C})\big)-\big(1+(-1)^q\big)\varkappa_q\big(R_{j,g}(\one)\big)\Big),\\
[v_1^q]\,V^{\mathbf O}_{C,(g)}=\frac1{q!}\sum_j\delta_j\Big(\varkappa_q\big(R_{j,g}(e^{v_0\one_C})\big)-(-1)^q\varkappa_q\big(R_{j,g}(e^{-v_0\one_C})\big)-\big(1-(-1)^q\big)\varkappa_q\big(R_{j,g}(\one)\big)\Big).
\end{gather*}
\end{enumerate}
\end{corollary}

\section{The order}

% Reference 20: def:order; source line 2379
\begin{definition}[Comparing]\label{def:order}
\begin{enumerate}[label=(\roman*)]
\item \emph{Exponents.} For $\mathbf a,\mathbf a'\in\mathbb N^{k+1}$, $\mathbf a\sqsubset\mathbf a'$ when $a_i<a'_i$ at the largest position $i$ with $a_i\neq a'_i$.
\item \emph{Functions.} For functions $f\neq f'$ of the positions $0,\dots,k$, analytic at $0$: $f\triangleleft f'$ when $[v^{\mathbf a}]f<[v^{\mathbf a}]f'$ at the $\sqsubset$-least exponent at which their coefficients differ.
\item \emph{Pairs.} $(f^{\mathbf E},f^{\mathbf O})\triangleleft(f'^{\mathbf E},f'^{\mathbf O})$ when $f^{\mathbf E}\triangleleft f'^{\mathbf E}$, or $f^{\mathbf E}=f'^{\mathbf E}$ and $f^{\mathbf O}\triangleleft f'^{\mathbf O}$.
\item \emph{Multisets.} Two different multisets of pairs with equally many elements are read from the least element upward in $\triangleleft$, and the one that first holds the smaller element comes first.
\item \emph{Classes.} For classes $C\neq D$ of one realization with $m$ marks each, $C\prec D$ when $\mathcal V_s(C)\triangleleft\mathcal V_s(D)$ at the first $s$ in $0,1,\dots,n-m,\ast$ with $\mathcal V_s(C)\neq\mathcal V_s(D)$.
\end{enumerate}
\end{definition}

% Reference 20: def:bunch; source line 2389
\begin{definition}[The Foreman bunch]\label{def:bunch}
At the centre, the \emph{Foreman bunch} $B(M;\boldsymbol\sigma,0)$ is the class of $\Upsilon_{\boldsymbol\sigma,0;\one}$ with the fewest marks and, among those, the $\prec$-least. At a lean $\boldsymbol\delta$ each class $C$ of $\Upsilon_{\boldsymbol\sigma,\boldsymbol\delta;\one}$ lies in one class $C^\circ$ of the centre (Theorem~\ref{thm:leansplits}), and $B(M;\boldsymbol\sigma,\boldsymbol\delta)$ is the class with the least
\[
\operatorname{key}_{\boldsymbol\sigma,\boldsymbol\delta}(C)=\big(|C|,\ |C^\circ|,\ C^\circ,\ C\big),
\]
compared entry by entry: sizes as integers, $C^\circ$ by $\prec$ at the centre, $C$ by $\prec$ at the lean. At every realization the Foreman bunch is a class with the fewest marks among the classes of that realization; the classes of the centre and the order only choose among those.
\end{definition}

% Reference 20: thm:order; source line 2397
\begin{theorem}[What the order is]\label{thm:order}
\begin{enumerate}[label=(\roman*)]
\item \emph{Total.} For every $m$, $\prec$ is a strict total order on the classes with $m$ marks; $B(M;\boldsymbol\sigma,0)$ and $B(M;\boldsymbol\sigma,\boldsymbol\delta)$ are each one class.
\item \emph{Directions.} With the classes given, $\prec$ depends on the realization only through the rays of $\boldsymbol\sigma$ and $\boldsymbol\delta$; at one pair only through the lean sign.
\item \emph{Mirror.} If $g$ exchanges the family (Definition~\ref{def:carry}), every stage of $g(C)$ at $-\boldsymbol\delta$ equals that stage of $C$ at $\boldsymbol\delta$; so $C\prec D$ at $\boldsymbol\delta$ exactly when $g(C)\prec g(D)$ at $-\boldsymbol\delta$, and $B(M;\boldsymbol\sigma,-\boldsymbol\delta)=g\big(B(M;\boldsymbol\sigma,\boldsymbol\delta)\big)$. If $g$ carries the family, the same holds at one lean.
\item \emph{Nesting.} When some class at $\boldsymbol\delta$ inside $B(M;\boldsymbol\sigma,0)$ has the fewest marks among the classes at $\boldsymbol\delta$, $B(M;\boldsymbol\sigma,\boldsymbol\delta)\subseteq B(M;\boldsymbol\sigma,0)$, and it is the class at $\boldsymbol\delta$ inside $B(M;\boldsymbol\sigma,0)$ with the fewest marks and, among those, the $\prec$-least. In a one-table molecule this is so whenever $B(M;\boldsymbol\sigma,0)$ splits at $\boldsymbol\delta$. Otherwise $B(M;\boldsymbol\sigma,\boldsymbol\delta)$ need not lie in $B(M;\boldsymbol\sigma,0)$.
\end{enumerate}
\end{theorem}

Every rise is a linear combination of logarithms of polynomials in $e^{v_0},\dots,e^{v_k}$, less the same combination at $v_0=0$ (Proposition~\ref{prop:skeletonviews}), so equality of two rises is an identity of polynomials with exponents (Theorem~\ref{thm:poles}). When two rises differ, the first differing exponent is found column by column: $v_k=0$ is the rise with one partner fewer, the coefficient of $v_k^b$ is a rational function of the scales at the other positions (Corollary~\ref{cor:firststages}(ii) for $k=1$), each column is certified equal by a polynomial identity, and the deciding coefficient is a joint cumulant of Proposition~\ref{prop:viewcumulants}, multiplied by the $\sigma_j$ or the $\delta_j$, and rational when the weights of the tables are.

% Reference 20: prop:rungsorder; source line 2416
\begin{proposition}[Rung one enters only at the last stage]\label{prop:rungsorder}
Every stage before the last reads the still from rung two on; the masses, rung one, enter only at the last stage, and at a centre they are the same at every mark.
\end{proposition}

\section{What the first stages read}

% Reference 20: prop:firstnumber; source line 2425
\begin{proposition}[At the centre the first number is how far the class is from a scale of its own]\label{prop:firstnumber}
$\operatorname{Var}_{T\sim\mathbb P_{j,\one}}|T\cap U|=0$ in every table exactly when $U$ is a union of molecules. So for a class $C$ of a molecule other than the whole,
\[
[v_0^2]\,\mathring V^{\mathbf E}_{C,()}=\sum_j\sigma_j\operatorname{Var}_{T\sim\mathbb P_{j,\one}}|T\cap C|>0 .
\]
\end{proposition}

% Reference 20: prop:mirror; source line 2435 (levels mirrored table by table; the exchange case for a one-table molecule).
% Proved in the last part: prop:mirror.
\begin{proposition}[Mirror classes]\label{prop:mirror}
Let $C,D$ be classes with equally many marks whose levels are mirrored across the tables: there are a bijection $h$ of the tables with $(\sigma_{h(j)},\delta_{h(j)})=(\sigma_j,\delta_j)$ and integers $k_j$ with $\mathbb P_{h(j),\one}(|T\cap D|=a)=\mathbb P_{j,\one}(|T\cap C|=k_j-a)$ for every $j$ and every $a$. This holds with $h$ the identity for two classes that together make up the molecule ($k_j=r_j$), and for $C$ and $g(C)$ under an exchange $g$ of the family, $h$ the bijection of its tables ($k_j=|C|$). Then:
\begin{enumerate}[label=(\roman*)]
\item $\kappa_{h(j),\one}(N_D^{\times a})=(-1)^a\kappa_{j,\one}(N_C^{\times a})$ for $a\ge2$, so $\sum_j\sigma_j\kappa_{j,\one}(N_D^{\times a})=(-1)^a\sum_j\sigma_j\kappa_{j,\one}(N_C^{\times a})$, and likewise with the $\delta_j$: the $\mathbf E$ parts of $\mathcal V_0(C)$ and $\mathcal V_0(D)$ are equal and their $\mathbf O$ parts are negatives;
\item at a lean, if some odd $a\ge3$ has $\sum_j\delta_j\kappa_{j,\one}(N_C^{\times a})\neq0$ and $a$ is the least such, then $C\prec D$ exactly when $\sum_j\delta_j\kappa_{j,\one}(N_C^{\times a})<0$, and at the opposite lean the order is the other one;
\item at the centre stage $0$ does not separate $C$ and $D$, and a later stage does.
\end{enumerate}
\end{proposition}

% Reference 20: cor:acrossbunch; source line 2447
\begin{corollary}[A bunch across the centre]\label{cor:acrossbunch}
In a one-table molecule, if $B(M;\sigma,0)$ splits at a lean $\delta$ into $C_K$ and $C_\Psi$ (Theorem~\ref{thm:acrosscentre}), then $B(M;\sigma,\delta)\in\{C_K,C_\Psi\}$ and $B(M;\sigma,-\delta)=\beta\big(B(M;\sigma,\delta)\big)$ is the other one, $\beta$ the exchange of Theorem~\ref{thm:acrosscentre}; the two are mirror classes, ordered at a lean by Proposition~\ref{prop:mirror}(ii).
\end{corollary}

The Foreman bunch is the least of the smallest classes, and the order that chooses it reads the class first on its own scale. At the centre the first number read is the spread of the level of the class under the weight laws, which vanishes only for a union of molecules; a lean adds the odd cumulants of that level, and two mirror classes are told apart by the first of them that does not vanish, the opposite lean reversing the choice.


\chapter{Peel, children and the Foreman tree}\label{ch:tree}

\section{The peel and the tree}

% Reference 20: def:peel; source line 2459
\begin{definition}[Peel; childcules]\label{def:peel}
Families are read at realizations at the table's gauge, and for a molecule $M$ of a family $F$, $F|M$ is the family of the restrictions of its tables to $M$. The \emph{peel} of a family $F$ at a bunch $B$ is the family of its children at the native split of the table's gauge (Proposition~\ref{prop:herecentre}),
\[
F_B(\mathbf a^\star,\mathbf b^\star),\qquad a^\star_{j,p}=\pi_{j,p}(\one)\,s_j,\quad b^\star_{j,p}=\pi_{j,p}(\one)\,t_j,
\]
on the room $R$; children with the same bases and weights are one table, with their exponents added. A \emph{childcule} is a molecule of the peel.
\end{definition}

% Reference 20: def:record; source line 2463
\begin{definition}[The record of a peel]\label{def:record}
The \emph{record} of the peel of $F$ at $B$ is the bunch $B$ with, for every table $j$, the communicatroid of $B$ (Definition~\ref{def:communicatroid}) and the pattern probabilities $\pi_{j,p}(\one)$: the indexed children $\big(j,p;\ \cB_{j,(p)},\theta_{j,(p)};\ \pi_{j,p}(\one)s_j,\ \pi_{j,p}(\one)t_j\big)$. The peel is the record with the indices $(j,p)$ forgotten and equal children merged.
\end{definition}

% Reference 20: prop:record; source line 2467
\begin{proposition}[The record keeps the parent]\label{prop:record}
The record determines $F$: table $j$ is the gluing of its communicatroid (Proposition~\ref{prop:interchange}), and its exponents are the sums of its children's, $\sum_p\pi_{j,p}(\one)(s_j,t_j)=(s_j,t_j)$.
\end{proposition}

Merging equal children forgets which pattern each share came from, and the Gammas $\Gamma\big(\mu_e-\langle\boldsymbol\delta\rangle_e\big)$ of the port (Theorem~\ref{thm:port}) depend on it. The record keeps it. A bunch of width one whose communicatroid is that of one mark, two children on two levels, may instead be pooled into one mark of the room, its shape integrated and its scale kept (Theorem~\ref{thm:pooling}).

% Reference 20: def:tree; source line 2478
\begin{definition}[The Foreman tree]\label{def:tree}
For a family $F$ read at $\Upsilon_{\boldsymbol\sigma,\boldsymbol\delta;\one}$, the \emph{Foreman tree} $\operatorname{Tree}(F;\boldsymbol\sigma,\boldsymbol\delta)$ splits $F$ into molecules. A molecule with one mark is a unit (Definition~\ref{def:molecule}), and a molecule $M$ whose realization is atomic is an atom (Definition~\ref{def:view}); otherwise the tree records the Foreman bunch $B(M;\boldsymbol\sigma_M,\boldsymbol\delta_M)$, $(\boldsymbol\sigma_M,\boldsymbol\delta_M)$ the pair coordinates of the tables of $M$, and continues with the peel of $F|M$ at that bunch, on its room.
\end{definition}

The tree asks each molecule whether it is an atom; where it is not, the tree takes the least of its smallest classes as the bunch, hands the room the children of the bunch at the native split, and asks again. Nothing is lost on the way down: the record keeps the parent, and the native split keeps the gauge native and the centre a centre.

% Reference 20: prop:twosteps; source line 2482
\begin{proposition}[Peeling in two steps is peeling at once]\label{prop:twosteps}
For $B=B_1\sqcup B_2$, the peel at $B_2$ of the peel of $F$ at $B_1$ is the peel of $F$ at $B$: the child of $F$ at $(j,p)$ and the child at $\big((j,p\cap B_1),p\cap B_2\big)$ of the peel at $B_1$ are one table on $\cE\setminus B$ with the same exponents. So if a molecule is peeled at $B_1$, its room is one childcule and the bunch there is $B_2$, the family reached is the peel at $B_1\sqcup B_2$.
\end{proposition}

% Reference 20: thm:threetrees; source line 2489
\begin{theorem}[Three trees]\label{thm:threetrees}
$\operatorname{Tree}(F;\boldsymbol\sigma,\boldsymbol\delta)$ depends on $(\boldsymbol\sigma,\boldsymbol\delta)$ only through their rays; at one pair only through the lean sign. So a molecule at one pair has three Foreman trees, $\operatorname{Tree}(F;\boldsymbol\sigma,0)$ and $\operatorname{Tree}(F;\boldsymbol\sigma,\pm\boldsymbol\delta)$ for a basis lean $\boldsymbol\delta$, fixed by its tables.
\end{theorem}

% Reference 20: lem:descent; source line 2496
\begin{lemma}[Maps descend]\label{lem:descent}
Let $g:\cE\to\cE'$ be a bijection, $B\subseteq\cE$ with room $R$, and $B'=g(B)$.
\begin{enumerate}[label=(\roman*)]
\item If $g$ carries $F$ to $F'$, then $g|_R$ carries the peel of $F$ at $B$ to the peel of $F'$ at $B'$, the child $(j,p)$ going to $(h(j),g(p))$, $h$ the bijection of the tables of Definition~\ref{def:carry}.
\item If $g$ exchanges $F$ with $F'$, then $g|_R$ exchanges the peel of $F$ at $B$ with the peel of $F'$ at $B'$, the child $(j,p)$ going to $(h(j),B'\setminus g(p))$.
\item In both cases $g$ maps the molecules of $F$ onto those of $F'$, and on a molecule $M$ its restriction carries, respectively exchanges, $F|M$ with $F'|g(M)$.
\end{enumerate}
\end{lemma}

% Reference 20: thm:mirrortree; source line 2508
\begin{theorem}[The mirror]\label{thm:mirrortree}
If $g$ exchanges $F$ with $F'$, then $g\big(\operatorname{Tree}(F;\boldsymbol\sigma,\boldsymbol\delta)\big)=\operatorname{Tree}(F';\boldsymbol\sigma,-\boldsymbol\delta)$: molecules go onto molecules, units onto units, atoms onto atoms, each recorded bunch onto the bunch recorded at the same place, rooms onto rooms. If $g$ carries $F$ to $F'$ the same holds at the same lean.
\end{theorem}

% Reference 20: cor:onetableexchange; source line 2515 (a one-table molecule).
% Proved in the last part: cor:onetableexchange.
\begin{corollary}[A molecule with an exchange]\label{cor:onetableexchange}
If $g$ exchanges a molecule $M$ (Definition~\ref{def:carry}), then $g$ maps $\operatorname{Tree}(M;\boldsymbol\sigma,0)$ onto itself, every class of the centre onto itself (Corollary~\ref{cor:exchangecumulants}), and $\operatorname{Tree}(M;\boldsymbol\sigma,\boldsymbol\delta)$ onto $\operatorname{Tree}(M;\boldsymbol\sigma,-\boldsymbol\delta)$: $B(M;\boldsymbol\sigma,-\boldsymbol\delta)=g\big(B(M;\boldsymbol\sigma,\boldsymbol\delta)\big)$ and $g\big(B(M;\boldsymbol\sigma,0)\big)=B(M;\boldsymbol\sigma,0)$. If $g$ carries $M$, it maps every class of $\Upsilon_{\boldsymbol\sigma,\boldsymbol\delta;\one}$ onto itself and each of the trees at $0$, $\boldsymbol\delta$ and $-\boldsymbol\delta$ onto itself.
\end{corollary}

% Reference 20: lem:windowturn; source line 2522 (an exchange of one table).
% Proved in the last part: lem:windowturn.
\begin{lemma}[An exchange turns a window over]\label{lem:windowturn}
If $g$ exchanges a family $F$ with $F'$, with the bijection $h$ of the tables (Definition~\ref{def:carry}), and $S'=g(S)$, then for every table $j$ the window of $S'$ in table $h(j)$ of $F'$ is $\{|S|-r_j(S),\dots,|S|-d_j(S)\}$: $S$ in table $j$ and $S'$ in table $h(j)$ have the same width, and the bases of table $j$ at level $k$ on $S$ correspond to those of table $h(j)$ at level $|S|-k$ on $S'$.
\end{lemma}

% Reference 20: cor:belowno; source line 2529
\begin{corollary}[Below a no]\label{cor:belowno}
A molecule that answers no is peeled at its Foreman bunch read at a centre of the table's gauge; each childcule is a molecule at a centre of the table's gauge, and whether it answers yes is read there in turn.
\end{corollary}

\section{Patterns of one kind}

% Reference 20: prop:levelform; source line 2563 (one table).
% Proved in the last part: prop:levelform.
\begin{proposition}[Level form]\label{prop:levelform}
A function of the patterns of a bunch $B$ has \emph{level form} when it depends on $|p|$ alone, and $B$ is in level form in a table when its children are, and in a family when it is in level form in every table. $B$ is in level form in table $j$, $K_{j,(p)}=K_{j,(p')}$ whenever $|p|=|p'|$, if and only if every transposition of two marks of $B$, fixing every other mark, maps the bases of table $j$ onto themselves with their weights; so $B$ is in level form in a family exactly when every such transposition carries the family with the identity on its tables.
\end{proposition}

% Reference 20: def:kinds; source line 2570
\begin{definition}[Kinds of pattern]\label{def:kinds}
Let a one-table molecule be read at the table's gauge and $B$ a set of marks. A \emph{writing from a pattern} $p$ is $K$ written in an order whose first $|p|$ positions hold $p$, the next $|B\setminus p|$ hold $B\setminus p$, and the rest the room; a writing of $\Psi$ from $B\setminus p$ likewise. Patterns $p,p'$ are of one \emph{kind at a lean} when one polynomial is written from $p$ and from $p'$, and of one \emph{kind at the centre} when moreover a writing from $p$ may equal a writing of $\Psi$ from $B\setminus p'$.
\end{definition}

% Reference 20: thm:kinds; source line 2574
\begin{theorem}[One child per kind]\label{thm:kinds}
\begin{enumerate}[label=(\roman*)]
\item $p,p'$ are of one kind at a lean exactly when a bijection keeping $B$ and carrying the table maps $p$ to $p'$, and of one kind at the centre exactly when, in addition, an exchanging bijection keeping $B$ may map $p$ to $B\setminus p'$. Being of one kind is reflexive, symmetric and transitive.
\item Patterns of one kind have one child between them: the child at $p'$ is the image of the child at $p$ under the bijection restricted to the room, or the dual of that image when the bijection exchanges; their pattern probabilities are equal.
\item So the level letters $K^{(k)}_B(\alpha)=\sum_{p\in\pat^{(k)}(B)}\alpha_B^pK_{(p)}(\alpha_R)$ are determined by one child per kind, the patterns of each kind, and the bijections.
\end{enumerate}
\end{theorem}

% Reference 20: cor:kindsfamily; source line 2585
\begin{corollary}[A family with maps in hand]\label{cor:kindsfamily}
Let $\mathcal G$ be any set of bijections keeping $B$ that carry or exchange a family. Patterns reached from one another by the moves $p\mapsto g(p)$, $g\in\mathcal G$ carrying, and $p\mapsto B\setminus g(p)$, $g\in\mathcal G$ exchanging, have children that are one table up to the bijections on the room and duality, with equal pattern probabilities.
\end{corollary}

Kinds are read on the writings of the two polynomials of one table; with several tables there is no single pair of polynomials to write, and the corollary takes the bijections as given.

% Reference 20: thm:patterntensor; source line 2592
\begin{theorem}[The pattern tensor of a class]\label{thm:patterntensor}
Let $B$ be a class of a one-table molecule at a realization at the table's gauge, $e,f\in B$, and $\beta$ the bijection taking an order in which the class's one function is written from $e$ to one in which it is written from $f$ (Theorem~\ref{thm:onetable}). Then:
\begin{enumerate}[label=(\roman*)]
\item $\beta$ carries the table, or at the centre possibly exchanges it, and maps every class onto itself; in particular $\beta(B)=B$ and $\beta(R)=R$;
\item the child at $\beta(p)$ is the child at $p$ with the room mapped by $\beta$, or, when $\beta$ exchanges, the child at $B\setminus\beta(p)$ is the dual of that image; so $p$ and $\beta(p)$, respectively $p$ and $B\setminus\beta(p)$, are of one kind;
\item when $\beta$ carries, the patterns through $e$ correspond to those through $f$; when it exchanges, the patterns through $e$ correspond to those avoiding $f$; in both cases with equal pattern probabilities and with children equal up to the bijection on the room and duality.
\end{enumerate}
So the pattern tensor of a class is the same from each of its marks, at the centre up to complements.
\end{theorem}

\section{What the room keeps}\label{sec:roomkeeps}

% Reference 20: def:keeppot; source line 2608 (at a realization at the table's gauge).
% Proved in the last part: def:keeppot.
\begin{definition}[Keeping the potential; twins; descents; copies]\label{def:keeppot}
Let a family be read at a realization $\Upsilon=((s_j,t_j)_j;\alpha)$, with the potential $(\Phi_{\mathbf E},\Phi_{\mathbf O})$ of Definition~\ref{def:potential} and $(g^\ast y)_e=y_{g(e)}$ as in Definition~\ref{def:carry}. A bijection $g$ of the marks of a molecule $M$ \emph{keeps the potential} when $\Phi_{\mathbf X}(\log\alpha+g^\ast y)=\Phi_{\mathbf X}(\log\alpha+y)$ for every $y$ and both parts; at the table's gauge this is $\Phi_{\mathbf X}\circ g^\ast=\Phi_{\mathbf X}$. A bijection $g:M\to M'$ onto the marks of a molecule $M'$, read at a realization with the gauge $\alpha'$, \emph{takes the potential} of $M$ to that of $M'$ when $\Phi_{\mathbf X}(\log\alpha+g^\ast y)=\Phi'_{\mathbf X}(\log\alpha'+y)+c_{\mathbf X}$ for every $y$ and both parts, with constants $c_{\mathbf X}$. Twins, descents and copies are read at the table's gauge, $\alpha=\one$: two childcules of the peel of $M$ are \emph{twins} when a bijection keeping the potential of $M$ maps one onto the other. A \emph{descent} of a Foreman tree is a sequence $M_0\supset M_1\supset\cdots\supset M_d$ of its nodes from the molecule $M_0$ to a leaf $M_d$, each $M_{k+1}$ a childcule of the peel of $M_k$ at its bunch. On a descent, with $B_i$ the bunch of $M_i$, the \emph{copies inside $M_k$} of $Y=M_i$ or $Y=B_i$, $k\le i$, are the sets $g_k(g_{k+1}(\cdots g_{i-1}(Y)\cdots))$, each $g_l$ a bijection keeping the potential of $M_l$; for $k=i$ they are $Y$ alone.
\end{definition}

% Reference 20: lem:certifypot; source line 2612 (at a realization at the table's gauge).
% Proved in the last part: lem:certifypot.
\begin{lemma}[The writings certify the potential]\label{lem:certifypot}
Let a molecule be read at a realization $\Upsilon=((s_j,t_j)_j;\alpha)$.
\begin{enumerate}[label=(\roman*)]
\item If $e$ and $f$ read alike, the bijection $\beta$ that takes, position by position, an order with $e$ first to an order with $f$ first in which one pair of functions is written from both (Proposition~\ref{prop:onewriting}) keeps the potential, and $\beta(e)=f$.
\item A bijection keeping the potential maps every class onto itself, and a bijection taking the potential of $M$ to that of $M'$ maps the classes of $M$ onto those of $M'$, at the centre and at the realization's lean.
\end{enumerate}
So two marks read alike exactly when a bijection keeping the potential takes one to the other.
\end{lemma}

% Reference 20: lem:peelfactors; source line 2624
\begin{lemma}[The peel reads the factors of the letters]\label{lem:peelfactors}
Write the potential of a family by its irreducible factors with their exponents $(L,w_{\mathbf E}(L),w_{\mathbf O}(L))$, each normalized by $L(\one)=1$ (Theorem~\ref{thm:poles}), and for a bunch $B$ let $L_{[q]}$ be the coefficient of $\alpha_B^{\,q}$ in $L$, a polynomial on the room $R$. In this statement write $\Phi^{\mathrm{peel}}_{\mathbf X}$ for the potential of the peel $F_B(\mathbf a^\star,\mathbf b^\star)$. Up to a constant,
\[
\Phi^{\mathrm{peel}}_{\mathbf X}(y_R)=\sum_Lw_{\mathbf X}(L)\sum_qL_{[q]}(\one)\log L_{[q]}\big(e^{y_R}\big),
\]
the inner sum over the $q$ with $L_{[q]}\neq0$, whose values $L_{[q]}(\one)$ add to $1$.
\end{lemma}

% Reference 20: thm:roomkeeps; source line 2639
\begin{theorem}[The room keeps its classes]\label{thm:roomkeeps}
Let a molecule $M$ be read at a realization at the table's gauge, let $B$ be a union of its classes, and let the children's weights factor across the childcules of the peel. Every bijection $g$ keeping the potential of $M$ keeps $B$; its restriction to the room keeps the potential of the peel, maps each childcule $Y$ onto a childcule $g(Y)$, and takes the potential of $Y$ to that of $g(Y)$, the constant being $0$ when $g(Y)=Y$. Consequently, for a class $C$ of $M$ inside the room:
\begin{enumerate}[label=(\roman*)]
\item the childcules that $C$ meets are twins of one another, and $C$ meets each of them in equally many marks;
\item inside each childcule, $C$ lies in one class.
\end{enumerate}
\end{theorem}

% Reference 20: prop:treepot; source line 2651
\begin{proposition}[The tree is read off the potential]\label{prop:treepot}
Let a molecule $M$ of $F$ and a molecule $M'$ of $F'$ be read at realizations at the table's gauge, with pair coordinates $(\boldsymbol\sigma,\boldsymbol\delta)$ and $(\boldsymbol\sigma',\boldsymbol\delta')$, the children's weights factoring across the childcules at every node of their trees, and let $g:M\to M'$ take the potential of $M$ to that of $M'$. Then $g$ maps $\operatorname{Tree}(F|M;\boldsymbol\sigma,\boldsymbol\delta)$ onto $\operatorname{Tree}(F'|M';\boldsymbol\sigma',\boldsymbol\delta')$: classes onto classes, units onto units, atoms onto atoms, each bunch onto the bunch recorded at the same place, and each childcule onto a childcule, its potential to that of its image.
\end{proposition}

% Reference 20: thm:leafglue; source line 2658
\begin{theorem}[The leaf bounds the bunches]\label{thm:leafbounds}
Let a molecule $F$ be read at a realization at the table's gauge, the centre or a lean, the children's weights factoring across the childcules at every node, and let $M_0\supset M_1\supset\cdots\supset M_d$ be a descent of its Foreman tree at that realization, $B_k$ the bunch of $M_k$ and $R_k=M_k\setminus B_k$. Let $h_k$ be the number of twins of $M_{k+1}$ among the childcules of $M_k$, $M_{k+1}$ itself included, and $h_{<k}=h_0h_1\cdots h_{k-1}$.
\begin{enumerate}[label=(\roman*)]
\item A class of $M_k$ inside $R_k$ that meets $M_{k+1}$ meets the twins of $M_{k+1}$ and no other childcule, each in equally many marks, and lies in one class of each. The copies inside $M_k$ of $M_i$, $k\le i\le d$, respectively of $B_i$, $k\le i<d$, are $h_{<i}/h_{<k}$ disjoint nodes, respectively bunches, of the subtree at $M_k$, leaves when $i=d$; and $M_d$ together with its copies inside $M_k$ is a union of whole classes of $M_k$.
\item A class $C$ of $M_k$ that meets $M_{k+1}$ has at least $|B_k|$ marks and meets $M_{k+1}$ in $|C|/h_k\ge|B_k|/h_k$ of them, all in one class of $M_{k+1}$; the bunch $B_{k+1}$, and the leaf when $k+1=d$, contains such a part. So
\[
|B_0|\le h_{<1}|B_1|\le h_{<2}|B_2|\le\cdots\le h_{<d-1}|B_{d-1}|\le h_{<d}|M_d|,
\]
the $k$-th term being the number of marks of the copies of $B_k$ inside $M_0$, and the last that of the copies of $M_d$: down the descent the bunches with their copies never shrink, and the leaf with its copies has at least as many marks as every bunch with its copies.
\end{enumerate}
\end{theorem}

Theorem~\ref{thm:leafbounds} rests on two things: the peel splits a class only across twin childcules, each part inside one class (Theorem~\ref{thm:roomkeeps}), and the bunch is a class with the fewest marks of the realization at which its node is read. Peeled at another bunch, a molecule can end at a leaf smaller than the bunch above it; at a lean the count is read on the classes of the lean, which can split a class of the centre into parts smaller than the least class of the centre; and the copies are needed when a room falls into twins.


\part{Evaluation}

\chapter{The integral at a bunch}\label{ch:closing}

\section{The scale of a bunch}

% Reference 20: def:skbreta; source line 2705
\begin{proposition}[A scale integral is a Breta of one bunch]\label{prop:scalebreta}
Let $J=\{0,\dots,L\}$ carry coefficients $\boldsymbol\xi=(\xi_\ell)_{\ell\in J}$ and letter exponents $\mathbf w=(w_\ell)_{\ell\in J}$ with total $|\mathbf w|=\sum_\ell w_\ell$, and let $I_J$ be the unit configuration on $J$. For $\boldsymbol\xi>0$, $\re w_\ell>0$ and $0<\re\kappa<\re|\mathbf w|$,
\[
\int_0^\infty x^{\kappa-1}\prod_{\ell\in J}(x+\xi_\ell)^{-w_\ell}\,dx=\Gst{2}{|J|}{\kappa,\ |\mathbf w|-\kappa}{\mathbf w}\,\sB^{\{J\}}_{\mathbf w,\,|\mathbf w|-\kappa;\,I_J}(\boldsymbol\xi):
\]
the Breta of Definition~\ref{def:breta} with the one bunch $J$, the shape index $\mathbf w$, the scale index $|\mathbf w|-\kappa$, the unit configuration as its terms and $\boldsymbol\xi$ as its coefficients. For $\xi_\ell\in\mathbb C\setminus(-\infty,0]$ the left side, with principal branches, converges exactly when $0<\re\kappa<\re|\mathbf w|$ and is holomorphic in its data there, and it continues the right side from positive coefficients and from letter exponents with positive real part; wherever the left side converges, the product on the right denotes its value. Coincident coefficients merge into one with the summed letter exponent.
\end{proposition}

The scale of a bunch is integrated against a product of linear letters, one for each point of the joint spectrum, and that integral is a Breta: the points of the spectrum are the coefficients, the letter exponents the shape index, and the index of the scale integral sits in the scale index as $|\mathbf w|-\kappa$. Nothing beyond the general Breta is needed to write it. With one coefficient the Breta is the power $\xi_0^{\kappa-w_0}$, and the identity is the integral of a linear letter: for $c_0,c_1>0$ and $0<\re\nu<\re s$, $\int_0^\infty\rho^{\nu-1}(c_0+c_1\rho)^{-s}\,d\rho=\Gst21{\nu,\ s-\nu}{s}c_0^{\nu-s}c_1^{-\nu}$. A stacked Gamma $\Gst21{a,\ b}{a+b}$ is a \emph{Beta corner}.

% New: proved in the last part.
\begin{proposition}[Where the Breta of one bunch converges]\label{prop:mot-bretadomain}
In Proposition~\ref{prop:scalebreta} let $\re w_\ell>0$ for every $\ell$ and $0<\re\kappa<\re|\mathbf w|$, and let $\Xi_J$ be the set of the $\boldsymbol\xi\in(\mathbb C\setminus(-\infty,0])^J$ whose convex hull avoids $(-\infty,0]$. It is open, star-shaped about $\one$, and contains every $\boldsymbol\xi$ whose coordinates have positive real parts. On $\Xi_J$ the integral that defines the Breta at positive coefficients (Definition~\ref{def:breta}),
\[
\int_{\Delta_J}u^{\mathbf w-\one}\Big(\sum_{\ell\in J}\xi_\ell u_\ell\Big)^{-(|\mathbf w|-\kappa)}d\omega_J(u),
\]
with the principal branch, converges, is holomorphic in $\boldsymbol\xi$, and continues the Breta $\sB^{\{J\}}_{\mathbf w,\,|\mathbf w|-\kappa;\,I_J}(\boldsymbol\xi)$ from positive coefficients; so the identity of Proposition~\ref{prop:scalebreta} holds on $\Xi_J$ with both sides convergent integrals.
\end{proposition}

% Reference 20: thm:scale; source line 2713
\begin{theorem}[The scale of a bunch is one Breta]\label{thm:scale}
Let a family $F$ on $\cE$ be at a gauge $\alpha$ with masses $\mvec$, let $B$ be a proper bunch with room $R$, its Skeleton polynomials, values and coins those of Definition~\ref{def:skelpoly}, and let $\xi_0,\dots,\xi_L$ be the distinct points of the joint spectrum of the letters with nonzero exponent, the union of the spectra $\operatorname{spec}^K_{j,B}(\alpha)$ with $s_j\neq0$ and $\operatorname{spec}^\Psi_{j,B}(\alpha)$ with $t_j\neq0$, with the letter exponents
\[
w_\ell=\sum_j\big(s_j\operatorname{mult}^K_{j,\ell}+t_j\operatorname{mult}^\Psi_{j,\ell}\big)=\sum_j\big(\sigma_j(\operatorname{mult}^K_{j,\ell}+\operatorname{mult}^\Psi_{j,\ell})+\delta_j(\operatorname{mult}^K_{j,\ell}-\operatorname{mult}^\Psi_{j,\ell})\big),
\]
$\operatorname{mult}^K_{j,\ell}$ and $\operatorname{mult}^\Psi_{j,\ell}$ the multiplicities of $\xi_\ell$ in $\operatorname{spec}^K_{j,B}(\alpha)$ and $\operatorname{spec}^\Psi_{j,B}(\alpha)$, and $J=\{0,\dots,L\}$. The \emph{scale integral}
\[
\operatorname{Sc}_B(\alpha;\mvec;\mathbf s,\mathbf t)=\int_0^\infty x^{m_B-1}\prod_jK_j(x\alpha_B,\alpha_R)^{-s_j}\Psi_j(x\alpha_B,\alpha_R)^{-t_j}\,dx
\]
converges exactly when $\re\kappa_B>0$ and $\re\kappa^{\mathrm{top}}_B>0$, and then
\[
\operatorname{Sc}_B(\alpha;\mvec;\mathbf s,\mathbf t)=\prod_jK_j(\alpha)^{-s_j}\Psi_j(\alpha)^{-t_j}\prod_\ell(1+\xi_\ell)^{w_\ell}\ \Gst{2}{|J|}{\kappa_B,\ \kappa^{\mathrm{top}}_B}{\mathbf w}\,\sB^{\{J\}}_{\mathbf w,\,\kappa^{\mathrm{top}}_B;\,I_J}(\boldsymbol\xi),
\]
with $|\mathbf w|=\sum_j(s_j+t_j)c_j(B)=\kappa_B+\kappa^{\mathrm{top}}_B$, so that the scale index $|\mathbf w|-\kappa_B$ of Proposition~\ref{prop:scalebreta} is the top wall; the Breta is read by that proposition at complex spectra.
\end{theorem}

% Reference 20: cor:stage; source line 2728 (one mark, neither a loop nor a coloop).
% Proved in the last part: cor:stage.
\begin{corollary}[The scale of a bunch of width one]\label{cor:stage}
Let $B$ be a bunch with room $R$. Along $B$ each letter of a table $j$ with $c_j(B)=1$ has one Skeleton value, the level letter at the floor over the level letter at the largest level, colevel letters for $\Psi_j$ (Definition~\ref{def:levels}), and its coin is the probability of the largest level, respectively colevel; the letters of a table with $c_j(B)=0$ have none:
\begin{gather*}
\operatorname{spec}^K_{j,B}(\alpha)=\Big\{\frac{K^{(d_j(B))}_{j,B}(\alpha)}{K^{(d_j(B)+1)}_{j,B}(\alpha)}\Big\},\ \text{coin }\mathbb P_{j,\alpha}\big(N_B=d_j(B)+1\big);\\
\operatorname{spec}^\Psi_{j,B}(\alpha)=\Big\{\frac{\Psi^{(d^\ast_j(B))}_{j,B}(\alpha)}{\Psi^{(d^\ast_j(B)+1)}_{j,B}(\alpha)}\Big\},\ \text{coin }\mathbb P^\ast_{j,\alpha}\big(N^\ast_B=d^\ast_j(B)+1\big).
\end{gather*}
At $B=\{e\}$, $e$ neither a loop nor a coloop of table $j$, the two levels are $K_{j,\setminus e}(\alpha_R)$ and $\alpha_eK_{j,/e}(\alpha_R)$ and the two colevels $\Psi_{j,/e}(\alpha_R)$ and $\alpha_e\Psi_{j,\setminus e}(\alpha_R)$: each Skeleton value is a deletion--contraction ratio, and the coins are the occupation $R_{j,e}(\alpha)$ and the dual occupation $R^\ast_{j,e}(\alpha)$.
\end{corollary}

% Reference 20: cor:slots; source line 2738
\begin{corollary}[One share per coefficient beyond the first]\label{cor:sharepercoefficient}
Let $\re w_\ell>0$ for all $\ell$, $0<\re\kappa<\re|\mathbf w|$, and choose $c_\ell\in(0,\re w_\ell)$, $\ell=1,\dots,L$, with $\re\kappa-\re w_0<\sum_\ell c_\ell<\re\kappa$, possible exactly under these conditions. With the share integrals on $\re a_\ell=c_\ell$ and $|\mathbf a|=a_1+\dots+a_L$,
\begin{multline*}
\Gst{2}{|J|}{\kappa,\ |\mathbf w|-\kappa}{\mathbf w}\,\sB^{\{J\}}_{\mathbf w,\,|\mathbf w|-\kappa;\,I_J}(\boldsymbol\xi)\\
=\int_{a_1+b_1=w_1}\!\!\cdots\int_{a_L+b_L=w_L}\ \prod_{\ell=1}^L\Gst21{a_\ell,\ b_\ell}{w_\ell}\xi_\ell^{-b_\ell}\cdot\xi_0^{\kappa-|\mathbf a|-w_0}\,\Gst21{\kappa-|\mathbf a|,\ w_0-\kappa+|\mathbf a|}{w_0}.
\end{multline*}
Any coefficient may serve as $\xi_0$; $L+1$ distinct values need $L$ shares.
\end{corollary}

% Reference 20: thm:gaps; source line 2749 ($0<\xi_\ell\le\xi_0$, positive letter exponents, about the largest coefficient).
% Proved in the last part: thm:gaps.
\begin{theorem}[The scale in its spectral gaps]\label{thm:gaps}
Let $\xi_\ell\in\mathbb C\setminus(-\infty,0]$ and $w_\ell\in\mathbb C$, $\ell\in J$, with $0<\re\kappa<\re|\mathbf w|$, and let $r>0$ have $|\xi_\ell-r|<r$ for every $\ell$; such an $r$ exists exactly when every $\xi_\ell$ has positive real part, since $|\xi_\ell-r|<r$ reads $\re\xi_\ell>|\xi_\ell|^2/(2r)$. Put $g_\ell=1-\xi_\ell/r$, the \emph{spectral gaps} at $r$, so that $|g_\ell|<1$. Then
\[
\Gst{2}{|J|}{\kappa,\ |\mathbf w|-\kappa}{\mathbf w}\,\sB^{\{J\}}_{\mathbf w,\,|\mathbf w|-\kappa;\,I_J}(\boldsymbol\xi)=r^{\kappa-|\mathbf w|}\sum_{\mathbf n\in\mathbb N^J}\prod_{\ell\in J}\frac{(w_\ell)_{n_\ell}}{n_\ell!}g_\ell^{\,n_\ell}\ \Gst21{\kappa,\ |\mathbf w|+|\mathbf n|-\kappa}{|\mathbf w|+|\mathbf n|},
\]
a convergent series whose terms, grouped by the total order $|\mathbf n|$, have root rate at most $\max_\ell|g_\ell|$. A coefficient equal to $r$ has spectral gap zero and drops out. For $0<\xi_\ell\le r$, $w_\ell>0$ and $0<\kappa<|\mathbf w|$ the gaps lie in $[0,1)$ and every term is positive. For positive coefficients the bound $\max_\ell|g_\ell|$ is least at $r=\frac12(\max_\ell\xi_\ell+\min_\ell\xi_\ell)$, where it is $\frac{\max_\ell\xi_\ell-\min_\ell\xi_\ell}{\max_\ell\xi_\ell+\min_\ell\xi_\ell}$, against $1-\frac{\min_\ell\xi_\ell}{\max_\ell\xi_\ell}$ at $r=\max_\ell\xi_\ell$.
\end{theorem}

% Reference 20: cor:inversion; source line 2760
\begin{corollary}[Inversion is Cremona on the scale]\label{cor:inversion}
In Theorem~\ref{thm:scale}:
\begin{gather*}
\sB^{\{J\}}_{\mathbf w,\,|\mathbf w|-\kappa;\,I_J}(\boldsymbol\xi)=\prod_\ell\xi_\ell^{-w_\ell}\,\sB^{\{J\}}_{\mathbf w,\,\kappa;\,I_J}(\boldsymbol\xi^{-1}),\\
\operatorname{Sc}_B(\alpha;\mvec;\mathbf s,\mathbf t)=(\alpha^{\cE})^{-2\sigma}\,\operatorname{Sc}_B(\alpha^{-1};2\sigma\one-\mvec;\mathbf t,\mathbf s):
\end{gather*}
the joint spectrum at $(\alpha^{-1};\mathbf t,\mathbf s)$ is the reciprocal of that at $(\alpha;\mathbf s,\mathbf t)$ with the same letter exponents, and the index of the scale integral moves from $\kappa_B$ to $\kappa^{\mathrm{top}}_B=|\mathbf w|-\kappa_B$.
\end{corollary}

\section{Closing}

% Reference 20: def:tail; source line 2789
\begin{definition}[Exceedance polynomials]\label{def:tail}
For a bunch $B$ of a table the \emph{exceedance polynomials} are $\sQ^K_{B,>}(x;\alpha)=\sum_{k\ge0}\mathbb P_\alpha(N_B-d(B)>k)\,x^k$ and $\sQ^\Psi_{B,>}(x;\alpha)=\sum_{k\ge0}\mathbb P^\ast_\alpha(N^\ast_B-d^\ast(B)>k)\,x^k$. They have nonnegative coefficients and degree $c(B)-1$, and $1-\sQ^K_B=(1-x)\sQ^K_{B,>}$, $1-\sQ^\Psi_B=(1-x)\sQ^\Psi_{B,>}$.
\end{definition}

% Reference 20: thm:sectors; source line 2793
\begin{theorem}[The scale closes on its two sectors]\label{thm:sectors}
Let $B$ be a proper bunch of a family with room $R$, and let the data be real with $s_j,t_j\ge0$, $\kappa_B>0$ and $\kappa^{\mathrm{top}}_B>0$. For $S\in\{B,R\}$ write the polynomials $\sQ^K_{j,S}(\cdot;\alpha)$, with the exponents $s_j$, and $\sQ^\Psi_{j,S}(\cdot;\alpha)$, with the exponents $t_j$, as $\sQ_{S,i}(\cdot;\alpha)$ with exponents $w_i$, $1\le i\le2N$ for a family of $N$ tables, $\mathbf w=(w_i)_i$, and their exceedance polynomials as $\sQ_{S,i,>}$.
\begin{enumerate}[label=(\roman*)]
\item The scale integral is the sum of the Trailmoments of its two sectors, with the coefficient functions $x\mapsto(\sQ_{S,i}(x;\alpha))_i$ and the inputs $x^{\kappa-1}\Ilat{(0,1)}$:
\begin{multline*}
\operatorname{Sc}_B(\alpha;\mvec;\mathbf s,\mathbf t)=\prod_jK_j(\alpha)^{-s_j}\Psi_j(\alpha)^{-t_j}\Big(\big(\sT_{\mathbf w;I_{2N}}[(\sQ_{B,i}(x;\alpha))_i]\,x^{\kappa_B-1}\Ilat{(0,1)}\big)\big((0,1)\big)\\
+\big(\sT_{\mathbf w;I_{2N}}[(\sQ_{R,i}(x;\alpha))_i]\,x^{\kappa^{\mathrm{top}}_B-1}\Ilat{(0,1)}\big)\big((0,1)\big)\Big).
\end{multline*}
\item With $\lambda_S=\max_i\big(1-\sQ_{S,i}(0;\alpha)\big)<1$ and $\kappa>0$,
\begin{multline*}
\big(\sT_{\mathbf w;I_{2N}}[(\sQ_{S,i}(x;\alpha))_i]\,x^{\kappa-1}\Ilat{(0,1)}\big)\big((0,1)\big)=\int_0^1x^{\kappa-1}\prod_i\sQ_{S,i}(x;\alpha)^{-w_i}\,dx\\
=\sum_{\mathbf n\in\mathbb N^{2N}}\prod_i\frac{(w_i)_{n_i}}{n_i!}\sum_{k\ge0}\Big([x^k]\prod_i\sQ_{S,i,>}(x;\alpha)^{n_i}\Big)\Gst21{\kappa+k,\ |\mathbf n|+1}{\kappa+k+|\mathbf n|+1},
\end{multline*}
a series of nonnegative terms, those with $|\mathbf n|=M$ summing to at most $\frac{(|\mathbf w|)_M}{M!}\lambda_S^M\kappa^{-1}$.
\end{enumerate}
For complex exponents with $\re\kappa>0$ the series converges absolutely to the same value. No Skeleton value is used, and a complex spectrum closes as a real one does.
\end{theorem}

% Reference 20: prop:axispoint; source line 2809
\begin{proposition}[Closing around a point of the axis]\label{prop:axispoint}
The expansion $x+\xi=(x+r)\big(1-\frac{r-\xi}{x+r}\big)$, $r>0$, converges uniformly on $x\ge0$ exactly when $|\xi-r|<r$; one $r$ serves every value of a finite spectrum exactly when every value has positive real part, that is, when every irreducible real factor of every Skeleton polynomial has positive coefficients.
\end{proposition}

% Reference 20: cor:cubeclose; source line 2816 (real data with $s_j,t_j\ge0$).
% Proved in the last part: cor:cubeclose.
\begin{corollary}[The period closes on the ordering sectors]\label{cor:cubeclose}
Let the data be complex with $\kappa_{\cE}=0$ and $\re\kappa_S>0$ for every nonempty proper $S$. For an order $\pi$ with tails $U_k$ and sector letters $k_{j,\pi},\psi_{j,\pi}$ (Theorem~\ref{thm:exist}), with $N_S=|T\cap S|$, on $[0,1]^{n-1}$
\[
1-\frac{k_{j,\pi}(q)}{K_j(\one)}=\sum_{k=2}^n(1-q_k)\,\oE_{T\sim\mathbb P_{j,\one}}\Big[\prod_{k'<k}q_{k'}^{N_{U_{k'}}-d_j(U_{k'})}\sum_{0\le i<N_{U_k}-d_j(U_k)}q_k^{\,i}\Big],
\]
the same for $\psi_{j,\pi}$ with the colevels and $d^\ast_j$; both sides lie in $[0,\lambda]$, $\lambda<1$ the largest of the numbers $1-\theta_{j,T^\pm_{j,\pi}}/K_j(\one)$. Consequently
\begin{multline*}
P_F(\mvec;\mathbf s,\mathbf t)=\prod_jK_j(\one)^{-s_j-t_j}\sum_\pi\sum_{\mathbf n,\mathbf n'}\prod_j\frac{(s_j)_{n_j}(t_j)_{n'_j}}{n_j!\,n'_j!}\\
\times\int_{[0,1]^{n-1}}\prod_kq_k^{\kappa_{U_k}-1}\prod_j\Big(1-\frac{k_{j,\pi}(q)}{K_j(\one)}\Big)^{n_j}\Big(1-\frac{\psi_{j,\pi}(q)}{K_j(\one)}\Big)^{n'_j}\,dq,
\end{multline*}
a convergent series, each integral a finite positive combination of products
\[
\prod_k\Gst21{\kappa_{U_k}+\beta_k,\ \gamma_k+1}{\kappa_{U_k}+\beta_k+\gamma_k+1},\qquad \beta_k,\gamma_k\in\mathbb N ;
\]
for real data with $s_j,t_j\ge0$ every term is nonnegative. With $w=\sum_j(|s_j|+|t_j|)$, the sum of the moduli of the exponents, the terms with $|\mathbf n|+|\mathbf n'|=M$ have moduli summing to at most $\frac{(w)_M}{M!}\lambda^M\prod_jK_j(\one)^{-\re(s_j+t_j)}\sum_\pi\prod_k(\re\kappa_{U_k})^{-1}$.
\end{corollary}

% Reference 20: prop:sectorcompiler; source line 2832 (rational data with $s_j,t_j\ge0$).
% Proved in the last part: prop:sectorcompiler.
\begin{proposition}[An enclosure of the period in the field of its data]\label{prop:sectorcompiler}
In Corollary~\ref{cor:cubeclose} let the data be real, and let the weights, masses and exponents lie in a subfield $k$ of $\mathbb R$. For an order $\pi$ put $a_{j,\pi}=1-k_{j,\pi}/K_j(\one)$ and $b_{j,\pi}=1-\psi_{j,\pi}/K_j(\one)$, let $\lambda_\pi<1$ be the largest of the numbers $1-\theta_{j,T^\pm_{j,\pi}}/K_j(\one)$, and let $f_{\pi,M}(q)$ be the coefficient of $\zeta^M$ in the power series $\prod_j(1-\zeta a_{j,\pi})^{-s_j}(1-\zeta b_{j,\pi})^{-t_j}$.
\begin{enumerate}[label=(\roman*)]
\item $\prod_jK_j(\one)^{s_j+t_j}\,P_F(\mvec;\mathbf s,\mathbf t)=\sum_\pi\sum_{M\ge0}I_{\pi,M}$ with $I_{\pi,M}=\int_{[0,1]^{n-1}}\prod_kq_k^{\kappa_{U_k}-1}f_{\pi,M}(q)\,dq$, an element of $k$, nonnegative when $s_j,t_j\ge0$, by $\int\prod_kq_k^{\kappa_{U_k}-1+\beta_k}dq=\prod_k(\kappa_{U_k}+\beta_k)^{-1}$.
\item With $w=\sum_j(|s_j|+|t_j|)$, if $r_{\pi,Q}=\lambda_\pi\max\big(1,\frac{w+Q+1}{Q+2}\big)<1$, then $\big|\sum_{M>Q}I_{\pi,M}\big|\le\operatorname{err}_{\pi,Q}=\prod_k\kappa_{U_k}^{-1}\,\frac{(w)_{Q+1}}{(Q+1)!}\frac{\lambda_\pi^{Q+1}}{1-r_{\pi,Q}}$, the sum being nonnegative when $s_j,t_j\ge0$, and $\operatorname{err}_{\pi,Q}\to0$ as $Q\to\infty$: the partial sums minus and plus $\sum_\pi\operatorname{err}_{\pi,Q}$, and for $s_j,t_j\ge0$ the partial sums and the partial sums plus $\sum_\pi\operatorname{err}_{\pi,Q}$, are elements of $k$ enclosing the left side of (i) within any prescribed width.
\item The coefficients $f_{\pi,M}$ satisfy a linear recurrence in $M$ with coefficients polynomial in $M$ and in the $a_{j,\pi},b_{j,\pi}$: the power series $f(\zeta)=\prod_i(1-\zeta A_i)^{-w_i}$, with $(A_i,w_i)$ running over the pairs $(a_{j,\pi},s_j)$ and $(b_{j,\pi},t_j)$, satisfies $\prod_i(1-\zeta A_i)\,f'(\zeta)=\big(\sum_iw_iA_i\prod_{i'\neq i}(1-\zeta A_{i'})\big)f(\zeta)$.
\end{enumerate}
\end{proposition}

% Reference 20: thm:peak; source line 2847
\begin{theorem}[The peak closing of a period]\label{thm:peak}
Let the data be real with $s_j,t_j\ge0$, $\kappa_{\cE}=0$ and every proper wall positive. Write the letters $K_j$ (exponent $s_j$, degree $r_j$, counts $d_j$) and $\Psi_j$ (exponent $t_j$, degree $n-r_j$, counts $d^\ast_j$) as $L_i$ with exponents $w_i$, degrees $g_i$ and counts $d_i(S)$, choose references $M_i\ge\max_{\Delta_{\cE}}L_i$, and put $b_{\mathbf n}=\int_{\Delta_{\cE}}u^{\mvec-\one}\prod_i(M_i-L_i(u))^{n_i}d\omega_{\cE}(u)\ge0$.
\begin{enumerate}[label=(\roman*)]
\item $P_F(\mvec;\mathbf s,\mathbf t)=\sum_{\mathbf n}\prod_i\frac{(w_i)_{n_i}}{n_i!}M_i^{-w_i-n_i}b_{\mathbf n}$, a series of nonnegative terms. With the references $M_i=\max_T\theta_{j,T}/g_i!$, $j$ the table of the letter $L_i$, every $b_{\mathbf n}$ is a finite positive combination of the stacked Gammas $\Gst{n}{1}{\mvec+\boldsymbol\gamma}{m_{\cE}+|\boldsymbol\gamma|}$, $\boldsymbol\gamma\in\mathbb N^{\cE}$; the same series, with every $b_{\mathbf n}$ such a finite positive combination, holds for $\int_{\Delta_{\cE}}u^{\mvec-\one}\prod_iL_i^{-w_i}\,d\omega_{\cE}$ with any letters $L_i$ homogeneous of degree $g_i$ with nonnegative coefficients and squarefree monomials, $M_i$ being the largest coefficient of $L_i$ over $g_i!$.
\item For each letter let $\varpi_i$ be the least of $\kappa_S/d_i(S)$ over nonempty proper $S$ with $d_i(S)>0$, and $k_i$ the largest length of a chain $S_1\subsetneq\cdots\subsetneq S_{k_i}$ with $\kappa_{S_l}=\varpi_id_i(S_l)$; put $\varpi=\min_i\varpi_i$ and $k=\max\{k_i:\varpi_i=\varpi\}$. The partial sums $P^{\le Q}$ of (i) over $\max_in_i\le Q$ satisfy $P_F(\mvec;\mathbf s,\mathbf t)-P^{\le Q}=\Theta\big(Q^{-\varpi}(\log Q)^{k-1}\big)$, and if a single letter attains $\varpi$,
\[
P_F(\mvec;\mathbf s,\mathbf t)-P^{\le Q}\sim\frac{c_i\,\Gamma(w_i+\varpi)}{\varpi\,\Gamma(k)\,\Gamma(w_i)}\,Q^{-\varpi}(\log Q)^{k-1},
\]
with $c_i>0$ the coefficient of $(z+\varpi_i)^{-k_i}$ in $M_i^{-z}\int_{\Delta_{\cE}}u^{\mvec-\one}L_i^z\prod_{i'}L_{i'}^{-w_{i'}}d\omega_{\cE}$, whose pole at $-\varpi_i$ has order exactly $k_i$. A letter that vanishes on no proper face closes geometrically.
\item The fastest start is the sharp reference $M_i=\max_{\Delta_{\cE}}L_i$; where it is attained inside the simplex, $u^\star_e\partial_e\log L_i(u^\star)=g_iu^\star_e$ for every mark. It can be attained on a proper face of the simplex.
\item At the target of a molecule with a native gauge every proper wall is positive (Theorem~\ref{thm:saddle}(iii)), so the series applies there and its rate $\varpi$ is positive.
\end{enumerate}
\end{theorem}

% Reference 20: thm:contcorners; source line 2870 (one real index between two consecutive negative integers).
% Proved in the last part: thm:contcorners.
\begin{theorem}[Continued Beta corners]\label{thm:contcorners}
In Theorem~\ref{thm:sectors} and Corollary~\ref{cor:cubeclose} let the data be complex and let no index $\kappa$, respectively no $\kappa_{U_k}$, be a nonpositive integer; for an index with $\re\kappa\le0$ let $h\in\mathbb N$ have $-h-1<\re\kappa\le-h$. Then the series converge, with each Beta corner $\Gst21{\kappa+k,\ |\mathbf n|+1}{\kappa+k+|\mathbf n|+1}$ in Theorem~\ref{thm:sectors} and $\Gst21{\kappa_{U_k}+\beta_k,\ \gamma_k+1}{\kappa_{U_k}+\beta_k+\gamma_k+1}$ in Corollary~\ref{cor:cubeclose} read as the meromorphic stacked Gamma, to the $\Reg$ of the integral, its continuation in the indices; the terms of total order $M$ are at most $\frac{(w)_M}{M!}\lambda^M$ times a polynomial in $M$ whose degree is the sum of the integers $h+1$ over the indices with real part at most $0$ (over the proper walls with $\re\kappa_S\le0$ in Corollary~\ref{cor:cubeclose}), with $w=\sum_j(|s_j|+|t_j|)$ the sum of the moduli of the exponents. The sector series needs no positivity of the walls and computes the continued value at every point off the pole hyperplanes.
\end{theorem}

The closings expand each letter around a reference and integrate term by term: on the two sectors of a bunch's scale the reference is the letter's value at the end of the sector, on the ordering sectors of the simplex it is the total weight, and at the peak it is a bound on the simplex. For real data with nonnegative exponents every term is a positive combination of stacked Gammas, nonnegative where the walls are positive, and the rate at which they fall is read from the walls.

\section{The port}

% Reference 20: thm:port; source line 2879
\begin{theorem}[The port]\label{thm:port}
Let $B$ be a bunch of a family $F$ on $\cE$, with room $R$, and suppose there are reals $c_{j,p},c'_{j,p}>0$, one for each pattern $p$ of $B$ in table $j$, with $\sum_pc_{j,p}=\re s_j$ and $\sum_pc'_{j,p}=\re t_j$, such that (a) $\re\mu_e-\sum_{j,p}\frac{c_{j,p}-c'_{j,p}}2\chi_p(e)>0$ for $e\in B$, and (b) the children $F_B(\mathbf c,\mathbf c')$ converge at the masses $\re\mvec|_R$. Then $J_F(\mvec;\mathbf s,\mathbf t)$ converges and, with the share integrals on $\re a_{j,p}=c_{j,p}$, $\re b_{j,p}=c'_{j,p}$,
\begin{multline*}
J_F(\mvec;\mathbf s,\mathbf t)=\int_{|\mathbf a_j|=s_j}\int_{|\mathbf b_j|=t_j}\prod_j\Gst{|\pat_j(B)|}{1}{\mathbf a_j}{s_j}\Gst{|\pat_j(B)|}{1}{\mathbf b_j}{t_j}\\
\times\prod_{e\in B}\Gamma\big(\mu_e-\langle\boldsymbol\delta\rangle_e\big)\,J_{F_B(\mathbf a,\mathbf b)}(\mvec|_R;\mathbf a,\mathbf b),
\end{multline*}
with the odd loads $\langle\boldsymbol\delta\rangle_e$ of the shares (Definition~\ref{def:shares}).
For $B=\cE$ this is the \emph{atom formula}: every child is $\{\varnothing\}$ with the weight $\theta_{j,T}$, and $J_{F_{\cE}(\mathbf a,\mathbf b)}(\varnothing;\mathbf a,\mathbf b)=\prod_{j,T}\theta_{j,T}^{-a_{j,T}-b_{j,T}}$.
\end{theorem}

The port opens every letter over the patterns of the bunch, one share for each pattern and letter. On the bunch the shares leave a product of Gammas, one per mark, whose arguments are the offsets less the odd loads; on the room they leave the interest of the children at the shares as their exponents.

% Reference 20: cor:content; source line 2890 (one mark).
% Proved in the last part: cor:content.
\begin{corollary}[A loop or coloop in every table]\label{cor:content}
Let every mark of a set $A$ be in every table a loop or a coloop, with $\eta_{j,e}=-1$ for a loop and $\eta_{j,e}=+1$ for a coloop. Then $J_F(\mvec;\mathbf s,\mathbf t)=\prod_{e\in A}\Gamma\big(\mu_e-\sum_j\eta_{j,e}\delta_j\big)\,J_{F'}(\mvec|_{\cE\setminus A};\mathbf s,\mathbf t)$, with $F'$ the family on $\cE\setminus A$ of the deletions of its tables at the loops and the contractions at the coloops (Definition~\ref{def:delcon}).
\end{corollary}

% Reference 20: cor:portbreta; source line 2897
\begin{corollary}[The damped bunch integral is a Breta]\label{cor:portbreta}
In Theorem~\ref{thm:port} fix $\alpha_R$, let $\re s_j,\re t_j>0$ (a letter with exponent zero is $1$ and gets no scale), put $z=\frac12\log\alpha$, introduce one singleton scale per letter, $V=\{v_{j,+},v_{j,-}\}_j$, put $\mathbf q=(\mathbf s,\mathbf t)$ on $V$, and take the partition $\cP_{\mathrm{port}}=\{B\}\cup\{\{v\}:v\in V\}$. The term arrays $(\pi^{\mathrm{port}},\ell^{\mathrm{port}})$ of the letter presentation (Proposition~\ref{prop:letters}) have the columns:
\begin{itemize}
\item for each $e\in B$, the kernel column with shape exponent $(\mathbf e_e,0_V)$, level $(1,0_V)$ and coefficient $1$;
\item for every $(j,p)$, a column with shape exponent $(\chi_p/2,\mathbf e_{j,+})$, level $(\ell_p,\mathbf e_{j,+})$ and coefficient $\Lambda_{j,(p)}(z_R)$;
\item for every $(j,p)$, a column with shape exponent $(-\chi_p/2,\mathbf e_{j,-})$, level $(-\ell_p,\mathbf e_{j,-})$ and coefficient $\Lambda_{j,(p)}(-z_R)$;
\end{itemize}
and $\boldsymbol\xi_{\mathrm{port}}(\alpha_R)$ lists the coefficients in this order. Then, either side converging with the other,
\[
\int_{\mathbb R^B_{>0}}\alpha_B^{\bmu|_B-\one}e^{-|\alpha_B|}\prod_j\Lambda_j(z)^{-s_j}\Lambda_j(-z)^{-t_j}\,d\alpha_B=\Gamma(\mu_B)\,\sB^{\cP_{\mathrm{port}}}_{(\bmu|_B,\mathbf q),\,(\mu_B,\mathbf q);\,(\pi^{\mathrm{port}},\ell^{\mathrm{port}})}\big(\boldsymbol\xi_{\mathrm{port}}(\alpha_R)\big).
\]
\end{corollary}

The scale of Theorem~\ref{thm:scale} keeps the letters closed and integrates the scale alone, one share per distinct Skeleton value beyond the first (Corollary~\ref{cor:sharepercoefficient}).

% Reference 20: thm:admissible; source line 2911
\begin{theorem}[The native split is an admissible contour]\label{thm:admissible}
For real data with $s_j,t_j\ge0$, not all zero, $\mvec>0$ and $\kappa_{\cE}>0$, a native gauge $\alpha^\star$ and a bunch $B$, the real parts $c_{j,p}=s_j\pi_{j,p}(\alpha^\star)$, $c'_{j,p}=t_j\pi^\ast_{j,p}(\alpha^\star)$ of the native split satisfy the hypotheses of Theorem~\ref{thm:port}; a letter with exponent zero is $1$ and is not opened.
\end{theorem}

% Reference 20: cor:portcomplete; source line 2918
\begin{corollary}[The port decides convergence]\label{cor:portcomplete}
For real data with $s_j,t_j\ge0$ and any bunch $B$, $J_F(\mvec;\mathbf s,\mathbf t)$ converges if and only if some real parts satisfy the hypotheses of Theorem~\ref{thm:port} at $B$.
\end{corollary}

\section{Signed letters}

% Reference 20: def:signed; source line 2927
\begin{definition}[Signed letters]\label{def:signed}
A \emph{signed letter} is $L=L_+-L_-$ with $L_\pm$ positive combinations of monomials of one degree $g$; adjoining $\zeta|u|^g$ gives $L+\zeta$ on the simplex, positive for $\zeta>-\min_{\Delta_{\cE}}L$.
\end{definition}

% Reference 20: prop:signed; source line 2931
\begin{proposition}[A signed letter is a boundary value]\label{prop:signed}
Let a signed letter $L$ enter a form with exponent $s$ beside positive letters $L_i$ with exponents $w_i$, and for $\zeta>-\min_{\Delta_{\cE}}L$ put $G(\zeta)=\int_{\Delta_{\cE}}u^{\mvec-\one}(L+\zeta)^{-s}\prod_iL_i^{-w_i}\,d\omega_{\cE}$.
\begin{enumerate}[label=(\roman*)]
\item With $\nu$ the image of $u^{\mvec-\one}\prod_iL_i^{-w_i}d\omega_{\cE}$ under $u\mapsto L(u)$, $G(\zeta)=\int(y+\zeta)^{-s}\,\nu(dy)$ is holomorphic off $(-\infty,-\min_{\Delta_{\cE}}L]$, continues across that half-line except at the negatives of the critical values of $L$ on the simplex and its faces, and satisfies a linear differential equation in $\zeta$ with polynomial coefficients.
\item $G(0\pm i0)=\int_{L>0}u^{\mvec-\one}L^{-s}\prod_iL_i^{-w_i}d\omega_{\cE}+e^{\mp i\pi s}\int_{L<0}u^{\mvec-\one}|L|^{-s}\prod_iL_i^{-w_i}d\omega_{\cE}$, the region integrals convergent for $\re s<1$ near regular zeros and taken by Lemma~\ref{lem:cube} in the coordinate $y=L(u)$ otherwise.
\item $G(0+i0)-G(0-i0)=-2i\sin(\pi s)\int_{L<0}u^{\mvec-\one}|L|^{-s}\prod_iL_i^{-w_i}d\omega_{\cE}$.
\end{enumerate}
\end{proposition}

On the scale of a bunch a zero of a signed letter is a split point: on each of the intervals it bounds, rescaled to $[0,1]$, the zero is at an endpoint, where it is a wall of the integral over that interval and Theorem~\ref{thm:contcorners} applies, and the phase $e^{\mp i\pi s}$ attaches to the interval on which the letter is negative.


\chapter{Integrals, expectations and the Foreman procedure}\label{ch:expectations}

\section{Forms, sections, completions}

% Reference 20: def:diffform; source line 2950
\begin{definition}[Differential forms, charts, schedules]\label{def:diffform}
The \emph{differential form} of a family at masses $\mvec$ is
\[
\omega=W\,d\alpha_{\mvec}=\alpha^{\mvec-\one}\prod_jK_j(\alpha)^{-s_j}\Psi_j(\alpha)^{-t_j}d\alpha ,
\]
the family word against the measure $d\alpha_{\mvec}$ (Definition~\ref{def:family}). It is \emph{projective} when $\kappa_{\cE}=0$, and then it is integrated over the rays of the orthant, over any section (Lemma~\ref{lem:chart}, next). The \emph{chart} at a mark $g$ is the section $\alpha_g=1$ with the coordinates $\alpha_e$, $e\neq g$, and the chart integral is $\int_{\alpha_g=1}\omega=\int_{\mathbb R^{\cE-g}_{>0}}(\cdots)|_{\alpha_g=1}\prod_{e\neq g}d\alpha_e$. Passing to the chart at $g$ is the \emph{snip} at $g$. A \emph{schedule} $\mathfrak s=(g;\varsigma)$ is the chart at a mark $g$ and an order $\varsigma(1),\dots,\varsigma(n-1)$ of the other marks.
\end{definition}

% Reference 20: lem:chart; source line 2958
\begin{lemma}[Every section gives the period]\label{lem:chart}
If $\omega$ is projective with convergent period, $f$ its integrand, $h>0$ homogeneous of degree one and $\varphi$ a function on $(0,\infty)$ with $\int_0^\infty|\varphi(\rho)|\,d\rho/\rho<\infty$, then $\int_{\mathbb R^{\cE}_{>0}}f\,\varphi(h)\,d\alpha=\big(\int_0^\infty\varphi(\rho)\frac{d\rho}\rho\big)P_F(\mvec;\mathbf s,\mathbf t)$. In particular every chart integral converges and equals $P_F(\mvec;\mathbf s,\mathbf t)$.
\end{lemma}

% Reference 20: prop:sectiontable; source line 11347
\begin{proposition}[A section is a table]\label{prop:sectiontable}
Let $W=\prod_jK_j^{-s_j}\Psi_j^{-t_j}$ be a family word, in this statement of the degree $-d_W$, $d_W=\sum_j(s_jr_j+t_jr^\ast_j)$, and $\mvec$ masses with every proper wall positive. Then $\int_{\Delta_{\cE}}u^{\mvec-\one}W\,d\omega_{\cE}$ is the value on its balance wall of the family with the list $\binom{\cE}1$, whose bases are the singletons and whose letter is $K_0=|\alpha|$, added at the exponents $(m_{\cE}-d_W,0)$; its proper walls are those of the family. In particular, for the family $F_0$ of the one table $\binom{\cE}1$, $\Dir_{\mvec}=\mathrm{Breta}^\circ(F_0;\mvec;m_{\cE},0)$, the Breta shape law of $\binom{\cE}1$ at a full basis lean; and for two families $F_1,F_2$ on $\cE$ with words $W_1,W_2$, the logarithms of the letters of $F_2$ under the law on $\Delta_{\cE}$ with density proportional to $u^{\mvec-\one}W_1$ have the Laplace transform
\[
\int_{\Delta_{\cE}}u^{\mvec-\one}W_1W_2\,d\omega_{\cE}\Big/\int_{\Delta_{\cE}}u^{\mvec-\one}W_1\,d\omega_{\cE}
\]
in the exponents of $W_2$, a ratio of two values of the family $F_1\cup F_2\cup F_0$.
\end{proposition}

% Reference 20: def:completion; source line 2965
\begin{definition}[Completions]\label{def:completion}
For $w\in\mathbb R^{\cE}$ with $w_{\cE}=1$ the \emph{completion} of $\omega=W\,d\alpha_{\mvec}$ along $w$ is $\omega_{\mvec'}=\alpha^{-\kappa_{\cE}w}\omega=W\,d\alpha_{\mvec'}$, $\mvec'=\mvec-\kappa_{\cE}w$, which is projective: a completion acts on the measure and leaves the word. The \emph{chart completion} at $g$ has $w=\mathbf e_g$ and the masses $\mvec^{(g)}=\mvec-\kappa_{\cE}\mathbf e_g$. For real data the \emph{native completion} has $w=\mvec/m_{\cE}$; its measure is the differential
\[
d\alpha_{\btau}=\alpha^{-(\kappa_{\cE}/m_{\cE})\mvec}\,d\alpha_{\mvec},\qquad \btau=\frac{n\sigma+\sum_j\delta_j(2r_j-n)}{n\sigma+\mu_{\cE}}\,\mvec ,
\]
at the target, with $\kappa_{\cE}=\mu_{\cE}-\sum_j\delta_j(2r_j-n)$ and $m_{\cE}=n\sigma+\mu_{\cE}$ in pair coordinates (Theorem~\ref{thm:pairform}(iii)), and the completed form is $\omega_{\btau}=W\,d\alpha_{\btau}$.
\end{definition}

% Reference 20: prop:snip; source line 2973 ($-\kappa_{\cE}\in\mathbb N$ in (iii)).
% Proved in the last part: prop:snip.
\begin{proposition}[Snips are completions]\label{prop:snip}
\begin{enumerate}[label=(\roman*)]
\item $\int_{\alpha_g=1}\omega=\int_{\alpha_g=1}\omega_{\mvec^{(g)}}$, the period $P_F(\mvec^{(g)};\mathbf s,\mathbf t)$ wherever it converges.
\item $K_j|_{\alpha_g=1}=K_{j,\setminus g}+K_{j,/g}$ and $\Psi_j|_{\alpha_g=1}=\Psi_{j,\setminus g}+\Psi_{j,/g}$, and for $\re s_j,\re t_j>0$, on the contours of Corollary~\ref{thm:E1},
\begin{multline*}
\prod_jK_j^{-s_j}\Psi_j^{-t_j}\Big|_{\alpha_g=1}\\
=\int\!\!\cdots\!\!\int\prod_j\Gst21{a_j,\ a'_j}{s_j}\Gst21{b_j,\ b'_j}{t_j}K_{j,\setminus g}^{-a_j}K_{j,/g}^{-a'_j}\Psi_{j,\setminus g}^{-b_j}\Psi_{j,/g}^{-b'_j}
\end{multline*}
over $a_j+a'_j=s_j$, $b_j+b'_j=t_j$: a snip opens every letter into its deletion and its contraction, with stacked-Gamma shares.
\item For a bunch $B$, the integral of $\omega$ over the bunch section $|\alpha_B|=1$ is the integral over the rays of the projective form $|\alpha_B|^{-\kappa_{\cE}}\omega$. In this statement write $F'$ for $F$ with one table appended at the exponents $(\kappa_{\cE},0)$: the list $\binom B1$ read on $\cE$, whose bases are the singletons of $B$, with the marks outside $B$ as loops, and whose letter is $|\alpha_B|$. Then $|\alpha_B|^{-\kappa_{\cE}}\omega$ is the form of $F'$ at the masses $\mvec$, $F'$ is balanced there, its proper walls are $\kappa_S-\kappa_{\cE}\,\one[B\subseteq S]$, and
\[
\int_{|\alpha_B|=1}\omega=P_{F'}\big(\mvec;(\mathbf s,\kappa_{\cE}),(\mathbf t,0)\big)
\]
wherever $\re\kappa_S-\re\kappa_{\cE}\,\one[B\subseteq S]>0$ for every nonempty proper $S$. At $B=\{g\}$ the appended letter is the monomial $\alpha_g$, and this is (i); at $B=\cE$ the appended table is the list $\binom{\cE}1$ (Proposition~\ref{prop:sectiontable}). For $k=-\kappa_{\cE}\in\mathbb N$, $k\ge1$, $|\alpha_B|^k=\sum_{|\mathbf n|=k}\binom k{\mathbf n}\alpha_B^{\mathbf n}$, so $\int_{|\alpha_B|=1}\omega=\sum_{\mathbf n\in\mathbb N^B,\,|\mathbf n|=k}\binom k{\mathbf n}\int_{|\alpha_B|=1}\omega_{\mvec+\mathbf n}$, each term the period $P_F(\mvec+\mathbf n;\mathbf s,\mathbf t)$ wherever it converges.
\end{enumerate}
\end{proposition}

% Reference 20: cor:fiber; source line 2772 (on the balance wall, at a chart of the room).
% Proved in the last part: cor:fiber.
\begin{corollary}[The period fibres over the scale of a bunch]\label{cor:fiber}
Let $B$ be a proper bunch of a family $F$ on $\cE$, with room $R$, and let $h>0$ be a function of $\alpha_R$ homogeneous of degree one. Wherever the integral of $W\,d\alpha_{\mvec}$ over the section $\{h(\alpha_R)=1\}$ of the orthant converges, both sections taken with the measures of Lemma~\ref{lem:chart},
\[
\int_{h(\alpha_R)=1}W\,d\alpha_{\mvec}=\int_{\{\alpha_R\in\mathbb R^R_{>0}:\ h(\alpha_R)=1\}}\int_{\Delta_B}u^{\mvec|_B-\one}\alpha_R^{\mvec|_R-\one}\,\operatorname{Sc}_B(u,\alpha_R;\mvec;\mathbf s,\mathbf t)\,d\omega_B(u),
\]
and at every point $\operatorname{Sc}_B$ is the product of Theorem~\ref{thm:scale}, a stacked Gamma times a Breta of one bunch. On the balance wall the left side is $P_F(\mvec;\mathbf s,\mathbf t)$. At the chart $h=\alpha_g$, $g\in R$, the left side is $P_F(\mvec^{(g)};\mathbf s,\mathbf t)$ (Proposition~\ref{prop:snip}(i)) and the outer integral is over $\alpha_{R-g}$ at $\alpha_g=1$; at $h=|\alpha_R|$ it is over $\Delta_R$ against $d\omega_R$.
\end{corollary}

% Reference 20: prop:scaling; source line 2988
\begin{proposition}[The scaling mode]\label{prop:scaling}
Along $z\mapsto z+c\one$, $E_j(z+c\one)=E_j(z)$, $O_j(z+c\one)=O_j(z)+2c(2r_j-n)$, and the exponent $2\ip{\bmu}z-\sum_j(\sigma_jE_j+\delta_jO_j)$ of the pair form moves by $2c\kappa_{\cE}$. Off the balance wall the form is not integrable along this direction; it is \emph{damped} by $e^{-|\alpha|}$ (the interest, $J_F(\mvec;\mathbf s,\mathbf t)=\Gamma(\kappa_{\cE})P_F(\mvec;\mathbf s,\mathbf t)$), \emph{conditioned} by a section (at a chart a completion, Proposition~\ref{prop:snip}; in general another projective form, Proposition~\ref{prop:sectionsoff} below), or \emph{slid} by moving the masses along their ray to the balance wall (the native completion).
\end{proposition}

% Reference 20: thm:transport; source line 2995
\begin{theorem}[The slide]\label{thm:transport}
Let $\omega$ have real data on a connected family, with $s_j,t_j\ge0$, not all zero, and $\mvec>0$, native gauge $\alpha^\star$ and target $\btau$.
\begin{enumerate}[label=(\roman*)]
\item If $\kappa_{\cE}\neq0$, exactly one completion of $\omega$ has $\alpha^\star$ as a native gauge: $\omega_{\btau}$. If $\kappa_{\cE}=0$, every completion is $\omega$ and $\btau=\mvec$.
\item As the masses slide along their ray, $\mvec\mapsto\lambda\mvec$, the data $(\lambda\mvec;\mathbf s,\mathbf t)$ lie on the balance wall exactly when $\lambda\mvec=\btau$.
\item Every proper wall of $\omega_{\btau}$ is positive, so its period converges.
\item $\btau$, $\omega_{\btau}$ and $\alpha^\star$ are the same at every point of the ray of $\mvec$.
\end{enumerate}
\end{theorem}

Each of the three treatments of the free scaling mode (Proposition~\ref{prop:scaling}) is an evaluation, and sliding keeps the native gauge.

\section{Three integrals and their expectations}

% New: proved in the last part.
\begin{lemma}[The period is meromorphic in the data]\label{lem:periodmero}
For every family, $P_F$ is meromorphic in the masses and the exponents, with poles only where a proper wall is a nonpositive integer, and so is its restriction to the balance wall.
\end{lemma}

% Reference 20: def:threeE; source line 3010
\begin{definition}[Three integrals and their expectations]\label{def:threeE}
Let $\omega=W\,d\alpha_{\mvec}$ have real data, $s_j,t_j\ge0$ not all zero and $\mvec>0$, so that $W$ is homogeneous of degree $-\tau_{\cE}$, $\tau_{\cE}=n\sigma+\sum_j\delta_j(2r_j-n)$ with $\sigma=\sum_j\sigma_j$. On a connected family the centre of $\omega$ is the form of the family at the centre (Definition~\ref{def:realization}), at the exponents $(\boldsymbol\sigma,\boldsymbol\sigma)$ and the masses $\sigma\one$,
\[
\omega_{\mathrm c}=\prod_j\big(K_j\Psi_j\big)^{-\sigma_j}d\alpha_{\sigma\one}.
\]
It is projective, its proper walls $\sum_j\sigma_jc_j(S)$ are positive and the table's gauge is native for it (Proposition~\ref{prop:centrefamily}), so its period $P_{\mathrm c}=P_F(\sigma\one;\boldsymbol\sigma,\boldsymbol\sigma)$ is a positive convergent integral, and $\mathbb P_{\mathrm c}=\mathrm{Breta}^\circ(F;\sigma\one;\boldsymbol\sigma,\boldsymbol\sigma)$, the measure $\omega_{\mathrm c}/P_{\mathrm c}$ read on the rays, is the \emph{law of the centre}. For masses $\mvec'$ of total $\tau_{\cE}$ the \emph{departure} of $W\,d\alpha_{\mvec'}$ from the centre is the ratio of the two forms,
\[
\frac{W\,d\alpha_{\mvec'}}{\omega_{\mathrm c}}(\alpha)=\alpha^{\mvec'-\sigma\one}\prod_j\Big(\frac{K_j(\alpha)}{\Psi_j(\alpha)}\Big)^{-\delta_j},
\]
homogeneous of degree $0$: an \emph{offset factor} $\alpha^{\mvec'-\sigma\one}$ and a \emph{lean factor} $\prod_j(K_j/\Psi_j)^{-\delta_j}$. Each integral below is unnormalized. On a connected family its expectation divides it by $P_{\mathrm c}$: it is the expectation, under the law of the centre, of the departure that its evaluation sets.
\begin{enumerate}[label=(\roman*)]
\item The \emph{integral} $\int W\,d\alpha_{\mvec}$ is defined when $\omega$ is projective with convergent period: it is the integral of $\omega$ over any section, $P_F(\mvec;\mathbf s,\mathbf t)$, and $\int W\,d\alpha_{\mvec}=P_{\mathrm c}\,\oE_{U\sim\mathbb P_{\mathrm c}}\big[(\omega/\omega_{\mathrm c})(U)\big]$.
\item For a schedule $\mathfrak s=(g;\varsigma)$ the \emph{quasi-integral} is the iterated finite part
\[
\Reg^{\mathfrak s}\!\!\int W\,d\alpha_{\mvec}=\Reg\!\int_0^\infty dx_{n-1}\cdots\Reg\!\int_0^\infty dx_1\Big(\alpha^{\mvec-\one}\prod_jK_j^{-s_j}\Psi_j^{-t_j}\Big)\Big|_{\alpha_g=1},\qquad x_k=\alpha_{\varsigma(k)},
\]
defined whenever the integrand of every one of these iterated integrals has log-power expansions at $0$ and at $\infty$. It sets the departure of the chart completion, and the \emph{quasi-expectation} is $\mathfrak E^{\mathfrak s}_{U\sim\mathbb P_{\mathrm c}}\big[(\omega_{\mvec^{(g)}}/\omega_{\mathrm c})(U)\big]=\Reg^{\mathfrak s}\!\!\int W\,d\alpha_{\mvec}\big/P_{\mathrm c}$.
\item For a connected family of at least two marks, the \emph{pseudointegral} is
\[
\pseudoint W\,d\alpha_{\mvec}=\CT_{\eta=0}P_F(\btau+\eta\mvec;\mathbf s,\mathbf t),
\]
the constant term at the target along the ray of the masses, $P_F$ being meromorphic in the masses with poles only where a proper wall is a nonpositive integer, on the \emph{pole hyperplanes} $\kappa_S=-q$, $S$ proper, $q\in\mathbb N$ (Lemma~\ref{lem:periodmero}). Off those hyperplanes it is $P_F(\btau;\mathbf s,\mathbf t)$, the value at the target of the meromorphic function $P_F$ on the balance wall. With a native gauge every proper wall is positive at $\btau$ and $\pseudoint W\,d\alpha_{\mvec}=\int W\,d\alpha_{\btau}$: the Foreman sets the measure to the native completion. It sets the departure of the native completion, and the \emph{pseudoexpectation} is
\[
\mathsf E_{U\sim\mathbb P_{\mathrm c}}\Big[\frac{\omega_{\btau}}{\omega_{\mathrm c}}(U)\Big]=\pseudoint W\,d\alpha_{\mvec}\Big/P_{\mathrm c};
\]
with a native gauge it is the expectation $\oE_{U\sim\mathbb P_{\mathrm c}}[(\omega_{\btau}/\omega_{\mathrm c})(U)]$: the Foreman sets the departure and the centre sets the law. On a family whose weights factor across its molecules, the pseudointegral and the pseudoexpectation are products over the molecules, each at its own target and with its own centre; a one-mark factor contributes the value of its integrand at $\alpha=\one$ to the pseudointegral and $1$ to the pseudoexpectation, its integrand and its centre's agreeing there.
\end{enumerate}
\end{definition}

A point $(\bmu;\boldsymbol\sigma,\boldsymbol\delta)$ splits into an even part, the strengths $\sigma_j$, and an odd part, the offsets and the leans; the fold (Theorem~\ref{thm:fold}) negates the odd part and keeps the even part. The centre is the point with the same even part and no odd part. It lies on the balance wall, it is inside its walls at every strength on a connected family, and its native gauge is the table's gauge, so its period is a convergent integral of the family itself, the same for every evaluation. Dividing by it prices the odd part of a point against its even part: with a native gauge the pseudoexpectation is the exponential of a divergence from the law of the centre, at least $1$, equal to $1$ exactly when the native completion is the centre's form, and unchanged by the fold (Proposition~\ref{prop:centrelaw}).

% Reference 20: thm:quasi; source line 3042
\begin{theorem}[Quasi-integrals extend the integral]\label{thm:quasi}
Each $\Reg^{\mathfrak s}\!\!\int$ is linear in the integrand, and $\Reg^{\mathfrak s}\!\!\int W\,d\alpha_{\mvec}=\int W\,d\alpha_{\mvec}$ for every schedule wherever $\int W\,d\alpha_{\mvec}$ is defined; there every chart completion is $\omega$, $\mvec^{(g)}=\mvec$, and on a connected family $\mathfrak E^{\mathfrak s}_{U\sim\mathbb P_{\mathrm c}}[(\omega_{\mvec^{(g)}}/\omega_{\mathrm c})(U)]=\oE_{U\sim\mathbb P_{\mathrm c}}[(\omega/\omega_{\mathrm c})(U)]$. Elsewhere the value can depend on the chart, on the order and on the coordinates (Lemma~\ref{lem:Regscale}).
\end{theorem}

% Reference 20: prop:compoundlaw; source line 3049 (two parts).
% Proved in the last part: prop:compoundlaw.
\begin{proposition}[The form of a compound at one point]\label{prop:compoundlaw}
Let $F$ be the direct sum of families $F_1,\dots,F_q$ on $\cE_1,\dots,\cE_q$, $q\ge2$, $\cE=\cE_1\sqcup\cdots\sqcup\cE_q$, every table a direct sum across the split with factoring weights, so that $K_j=\prod_iK_{j,i}$ and $\Psi_j=\prod_i\Psi_{j,i}$; let the data be real, the masses $\mvec=(\mvec_1,\dots,\mvec_q)$ on the balance wall of $F$, and $\kappa_i$ the imbalance of $F_i$ at $\mvec_i$, so $\sum_i\kappa_i=0$. On the section $|\alpha_q|=1$ with $\alpha_i=\rho_iu_i$, $u_i\in\Delta_{\cE_i}$, for $i<q$ and $u_q=\alpha_q$, the form of $F$ is
\[
\prod_{i<q}\rho_i^{\kappa_i-1}d\rho_i\ \cdot\ \prod_{i=1}^qu_i^{\mvec_i-\one}\prod_jK_{j,i}^{-s_j}\Psi_{j,i}^{-t_j}\,d\omega_{\cE_i}.
\]
\begin{enumerate}[label=(\roman*)]
\item Its integral diverges for all $\kappa_1,\dots,\kappa_q$.
\item The finite part of each $\rho_i$-integral is $0$ for every $\kappa_i$.
\item On the masses $\mvec(\boldsymbol\eta)=(\mvec_i^0+i\eta_iw_i)_{i=1}^q$, $\boldsymbol\eta\in\mathbb R^q$ with $\sum_i\eta_i=0$, with $\mvec^0_i$ balanced for $F_i$ inside its walls and $w_i\in\mathbb R^{\cE_i}$ of total $1$, so that $\kappa_i=i\eta_i$, the integral of the form is, as a distribution in $(\eta_1,\dots,\eta_{q-1})$, $(2\pi)^{q-1}\prod_{i=1}^qP_{F_i}(\mvec^0_i;\mathbf s,\mathbf t)$ times the Dirac measure at $0$.
\end{enumerate}
So the compound carries nothing beyond the points of its parts, and only where each molecule is balanced on its own.
\end{proposition}

% Reference 20: thm:invariance; source line 3060
\begin{theorem}[Presentations]\label{thm:invariance}
Let $\omega=W\,d\alpha_{\mvec}$ have real data on a connected family, $s_j,t_j\ge0$ not all zero and $\mvec>0$. A \emph{presentation} of $\omega$ is a bijection $\varphi$ of the marks together with coordinates $x=\lambda^{-1}\alpha$, $\lambda\in\mathbb R^{\cE}_{>0}$, in which the letters are $K_j(\lambda x)$, $\Psi_j(\lambda x)$ and the form carries the constant $\lambda^{\mvec}$. The split, the targets, the native gauges as points of the torus, the stills, the classes, the native splits and the children computed from a presentation are the images under $\varphi$ of those computed from $\omega$, and, when the target lies on no pole hyperplane, the pseudointegral and the pseudoexpectation of the presentation, its centre presented with it, are $(\lambda^{\mvec})^{\kappa_{\cE}/m_{\cE}}$ times those of $\omega$, the factor being $1$ on the balance wall; on a pole hyperplane the constant term moves with the coordinates as well (Proposition~\ref{prop:polecoords} below). In the coordinates $x$ the table's gauge is $x=\lambda^{-1}$, and the Foreman tree is the same. On a family whose weights factor across its molecules the factor is the product of the molecules' factors.
\end{theorem}

% Reference 20: def:procedure; source line 3067
\begin{definition}[The Foreman procedure]\label{def:procedure}
\emph{Classification}, from the tables alone:
\begin{enumerate}[label=(\arabic*)]
\item \emph{Split} the marks into molecules and units (Proposition~\ref{prop:compound}).
\item \emph{Read} each molecule $M$ at a centre $\Upsilon_{\boldsymbol\sigma,0;\one}$ of the table's gauge (Chapter~\ref{ch:reading}).
\item \emph{Answer}: a constant reading is an atom; otherwise take the Foreman bunch $B(M;\boldsymbol\sigma,0)$ (Definition~\ref{def:bunch}).
\item \emph{Peel} at $B(M;\boldsymbol\sigma,0)$; the childcules sit at the centres they inherit (Proposition~\ref{prop:herecentre}).
\item \emph{Recurse} to atoms and units: the output is the Foreman tree and, at the root, the label.
\end{enumerate}
\emph{Evaluation} at a real point $(\mvec;\mathbf s,\mathbf t)$ with $s_j,t_j\ge0$, not all zero, and $\mvec>0$:
\begin{enumerate}[label=(\arabic*),resume]
\item \emph{Native gauge} of each molecule, and its target $\btau$ (Theorem~\ref{thm:saddle}).
\item \emph{Slide} the masses to the target (Theorem~\ref{thm:transport}).
\item \emph{Ports} along the tree, every share contour through the native split (Theorems~\ref{thm:port}, \ref{thm:admissible} and~\ref{thm:hered}); on the balance wall the port of Chapter~\ref{ch:law}.
\item \emph{Atoms}: the atom formula (Theorem~\ref{thm:port} with $B=\cE$; on the balance wall, Chapter~\ref{ch:law}).
\end{enumerate}
On the orthant side the ports integrate the Foreman bunch of each node inside the integral over its room, all its marks at once in scale and shape (Corollary~\ref{cor:portbreta}): the root's bunch first and the atoms last. This is the \emph{Foreman's order}.
\end{definition}

% Reference 20: thm:foremanE; source line 3086
\begin{theorem}[The Foreman computes the pseudointegral]\label{thm:foremanE}
Let $\omega=W\,d\alpha_{\mvec}$ have real data on a connected family, with $s_j,t_j\ge0$, not all zero, and $\mvec>0$, and a native gauge. For every $\eta>0$ the procedure at the masses $\btau+\eta\mvec$ runs on admissible contours and computes $J_F(\btau+\eta\mvec;\mathbf s,\mathbf t)$, and
\[
\pseudoint W\,d\alpha_{\mvec}=\lim_{\eta\to0^+}\frac{J_F(\btau+\eta\mvec;\mathbf s,\mathbf t)}{\Gamma(\eta m_{\cE})}.
\]
No finite part, chart, order or closing rule enters: the regulator $\eta$ leaves by a limit.
\end{theorem}

% Reference 20: cor:onemember; source line 3097
\begin{corollary}[One member among the generalized integrals]\label{cor:onemember}
Wherever $\int W\,d\alpha_{\mvec}$ is defined for real data on a connected family, $\pseudoint W\,d\alpha_{\mvec}=\int W\,d\alpha_{\mvec}$, and with Theorem~\ref{thm:quasi} every evaluation sets the one departure $\omega/\omega_{\mathrm c}$ on the balance wall. Off the balance wall, with $s_j,t_j\ge0$, not all zero, and $\mvec>0$, the period at the target converges exactly when the family has a native gauge (Theorem~\ref{thm:saddle}(iii)), and otherwise $\pseudoint W\,d\alpha_{\mvec}$ is its continuation in the masses, on a pole hyperplane the constant term of the continuation along the ray.
\end{corollary}

% Reference 20: thm:offdomain; source line 3104 (a connected family with real data, $s_j,t_j\ge0$).
% Proved in the last part: thm:offdomain.
\begin{theorem}[The pseudointegral off the domain]\label{thm:offdomain}
For every family, if some nonempty proper $S$ has $c_j(S)=0$ in every active table, then, first where every proper wall has positive real part and then as meromorphic functions (Lemma~\ref{lem:periodmero}),
\[
P_F(\mvec;\mathbf s,\mathbf t)=\Gst21{\kappa_S,\ \kappa_{\cE\setminus S}}{\kappa_{\cE}}\int_{\Delta_S\times\Delta_{\cE\setminus S}}\alpha^{\mvec-\one}\prod_jK_j(\alpha)^{-s_j}\Psi_j(\alpha)^{-t_j}\,d\omega_S\,d\omega_{\cE\setminus S},
\]
and the restriction of $P_F$ to the balance wall is $0$. Let now $F$ be connected with real data, $s_j,t_j\ge0$ not all zero; on the balance wall the value of $P_F$ at a target depends on no path.
\begin{enumerate}[label=(\roman*)]
\item The continuation in the masses of $f\mapsto\int_{\Delta_{\cE}}f\,u^{\mvec-\one}\prod_jK_j^{-s_j}\Psi_j^{-t_j}d\omega_{\cE}$, taken by Lemma~\ref{lem:cube} on each ordering sector, is a distribution on $\Delta_{\cE}$ with normaliser $P_F(\mvec;\mathbf s,\mathbf t)$ and monomial moments $P_F(\mvec+\boldsymbol\gamma;\mathbf s,\mathbf t)/P_F(\mvec;\mathbf s,\mathbf t)$; where every proper wall is positive it is $P_F(\mvec;\mathbf s,\mathbf t)$ times the Breta shape law. Just past the first wall crossed it is not a positive measure: $P_F$ passes from $+\infty$ to $-\infty$ there, while $P_F(\mvec+\boldsymbol\gamma;\mathbf s,\mathbf t)$ is a positive convergent integral for $\gamma_S$ large.
\item Theorem~\ref{thm:cremona} and the recurrences of Theorem~\ref{thm:confJ} hold for these values as identities of meromorphic functions.
\end{enumerate}
\end{theorem}

% Reference 20: thm:fold; source line 3116
\begin{theorem}[The fold: Cremona off the balance wall]\label{thm:fold}
Let $\omega=W\,d\alpha_{\mvec}$, with the exponents $(\mathbf s,\mathbf t)$, have real data on a connected family, with $s_j,t_j\ge0$, not all zero, and $\mvec>0$, native gauge $\alpha^\star$ and target $\btau$. Then $2\sigma\one-\btau$ is positive and balanced for the exponents $(\mathbf t,\mathbf s)$, with native gauge $(\alpha^\star)^{-1}$, and for every $\lambda>0$
\[
\pseudoint\prod_jK_j^{-s_j}\Psi_j^{-t_j}\,d\alpha_{\mvec}=\pseudoint\prod_jK_j^{-t_j}\Psi_j^{-s_j}\,d\alpha_{\lambda(2\sigma\one-\btau)}.
\]
The map from the ray of $\mvec$ to the ray of $2\sigma\one-\btau$ is an involution on rays.
\end{theorem}

The fold swaps the two letters of every table and reflects the target through $\sigma\one$, negating the offsets and the leans; it keeps the strengths, so the two forms have one centre, and it carries uniform masses to uniform masses.

% Reference 20: prop:readingpoint; source line 3134
\begin{proposition}[Points and realizations]\label{prop:readingpoint}
A real point $(\mvec;\mathbf s,\mathbf t)$ of a connected family, with $s_j,t_j\ge0$, not all zero, and $\mvec>0$ and native gauge $\alpha^\star$, determines the realization $((s_j,t_j)_j;\alpha^\star)$, whose masses are the target; every realization arises so (Proposition~\ref{prop:prescribe}). The reading of a point is the reading of its realization and may differ from point to point; the label is read at centres of the table's gauge, the Foreman tree at the realizations of the table's gauge at the centre and at the leans (Theorem~\ref{thm:threetrees}), and both are functions of the tables. An atomic point is at the table's gauge up to scale (Theorem~\ref{thm:gaugestep}), and its centre is atomic (Theorem~\ref{thm:leansplits}).
\end{proposition}

% Reference 20: thm:parity; source line 3141
\begin{theorem}[Parity]\label{thm:parity}
Let $g$ be an exchange of a connected family, with the bijection $h$ of its tables (Definition~\ref{def:carry}); then $r_{h(j)}=n-r_j$ for every $j$. With $(g^\ast z)_e=z_{g(e)}$ (Definition~\ref{def:carry}), let $\Phi^\circ$ be the centred potential (Theorem~\ref{thm:saddle}).
\begin{enumerate}[label=(\roman*)]
\item $\Lambda_j(g^\ast z)=\Lambda_{h(j)}(-z)$; so $E_j(g^\ast z)=E_{h(j)}(z)$, $O_j(g^\ast z)=-O_{h(j)}(z)$ and $\Phi^\circ(-g^\ast z)=\Phi^\circ(z)$.
\item If the data are real with $s_j,t_j\ge0$, not all zero, and $\mvec>0$, have a native gauge, and $\tau_{g(e)}=2\sigma-\tau_e$ for every $e$, the native gauge $e^{2z^\star}$ with $\ip\one{z^\star}=0$ has $-g^\ast z^\star=z^\star$, and the still of the realization $\Upsilon=((s_j,t_j)_j;e^{2z^\star})$ has $\kappa_\Upsilon(g(e_1),\dots,g(e_k))=(-1)^k\kappa_\Upsilon(e_1,\dots,e_k)$ for every $k$.
\item Under (ii) even rungs never separate $e$ from $g(e)$, and at odd rungs their tallies (Definition~\ref{def:still}) are negatives of each other.
\item If moreover $\boldsymbol\delta=0$ and $\btau=\sigma\one$, then $z^\star=0$, every odd rung vanishes, and no rung separates $e$ from $g(e)$.
\item At $(\mathbf t,\mathbf s)$ and the target $2\sigma\one-\btau$ the native gauge is $e^{-2z^\star}$ and rung $k$ is $(-1)^k$ times rung $k$ at $(\mathbf s,\mathbf t)$; under (ii) this is the still at $(\mathbf s,\mathbf t)$ carried by $g$, and the Foreman bunch at $(\mathbf t,\mathbf s)$ is the image under $g^{-1}$ of the Foreman bunch at $(\mathbf s,\mathbf t)$.
\end{enumerate}
\end{theorem}

\section{Quasi-integrals as Laurent data}

% Reference 20: def:regular; source line 3157
\begin{definition}[Regular schedules]\label{def:regular}
For a schedule $\mathfrak s=(g;\varsigma)$ put $x_k=\alpha_{\varsigma(k)}$ and $v_k=\mathbf e_{\varsigma(k)}-\mathbf e_g$. Let $\mathcal J_0$ be the chart integrand and $\mathcal J_k(\varepsilon_1,\dots,\varepsilon_k;x_{>k})=\int_0^\infty x_k^{\varepsilon_k}\mathcal J_{k-1}\,dx_k$, with $\int_0^1$ and $\int_1^\infty$ each continued meromorphically in the \emph{regulator} $\varepsilon_k$. The schedule is \emph{regular} when for every $k$, as $x_k\to0$ and as $x_k\to\infty$, $\mathcal J_{k-1}$ is a finite sum of terms $c(\varepsilon,x_{>k})\,x_k^{a(\varepsilon)}$ plus $O(x_k^{\pm M})$ for every $M$, the exponents affine in $\varepsilon_{<k}$ and independent of $x_{>k}$, uniformly for the regulators in a small polydisc away from their poles.
\end{definition}

% Reference 20: thm:bridge; source line 3161
\begin{theorem}[Bridge]\label{thm:bridge}
For real data on a connected family and a regular schedule, $\Reg^{\mathfrak s}\!\!\int W\,d\alpha_{\mvec}$ is defined and
\[
\Reg^{\mathfrak s}\!\!\int W\,d\alpha_{\mvec}=\CT_{\varepsilon_{n-1}}\cdots\CT_{\varepsilon_1}P_F\Big(\mvec^{(g)}+\sum_{k=1}^{n-1}\varepsilon_kv_k;\mathbf s,\mathbf t\Big),
\]
the constant terms taken innermost first, $|\varepsilon_1|\ll\cdots\ll|\varepsilon_{n-1}|$.
\end{theorem}

% Reference 20: thm:last; source line 3172
\begin{theorem}[A wall is closed by its last variable]\label{thm:last}
For a schedule $\mathfrak s=(g;\varsigma)$ put $\mathbf p=\mvec^{(g)}$ and $I(\varepsilon)=P_F\big(\mathbf p+\sum_k\varepsilon_kv_k;\mathbf s,\mathbf t\big)$. A proper $U$ is \emph{polar} when $\kappa_U(\mathbf p)\in-\mathbb N$; its \emph{members} are the $k$ with $\varsigma(k)\in U$ if $g\notin U$, and with $\varsigma(k)\notin U$ if $g\in U$; $\operatorname{last}(U)$ is its last member.
\begin{enumerate}[label=(\roman*)]
\item Near $0$, $I(\varepsilon)=\sum_C f_C(\varepsilon)\big/\prod_{U\in C}\ell_U(\varepsilon)+H(\varepsilon)$, $\ell_U=\pm\sum_{k\text{ member}}\varepsilon_k$, over the nonempty sets $C$ of polar walls that are tails of one order, with $f_C,H$ holomorphic; the member sets in one $C$ are nested or disjoint.
\item Under $\CT_{\varepsilon_k}$, $1/\ell_U$ is holomorphic at $\varepsilon_k=0$ unless $k=\operatorname{last}(U)$; for one polar wall with numerator $f_C$ the iterated constant term of $f_C/\ell_U$ is $\pm\partial_{\varepsilon_a}f_C(0)$, $a=\operatorname{last}(U)$.
\item For a fixed chart, the quasi-integral of a regular schedule depends on the order only through $U\mapsto\operatorname{last}(U)$ on the polar walls.
\end{enumerate}
\end{theorem}

With one regulator $\varepsilon_D$ for each class $D$ of a realization, taken in some order of the classes, $\ell_U=\pm\sum_D|U\cap D|\,\varepsilon_D$; a wall is closed by the last class it meets, which contributes $|U\cap D|^{-1}\partial_{\varepsilon_D}$, and the value depends only on the order of the classes, which can be taken from the order of classes (Theorem~\ref{thm:order}).

% Reference 20: prop:coordscale; source line 3188
\begin{proposition}[The coordinates set the log-scale]\label{prop:coordscale}
In the coordinates $u=\alpha/\beta$, $\beta_g=1$, with the regulators $\varepsilon_k$ entering as $u_{\varsigma(k)}^{\varepsilon_k}$, the regulated function is $e^{-\sum_k\varepsilon_k\lambda_k}I(\varepsilon)$, $\lambda_k=\log\beta_{\varsigma(k)}$, and for one polar wall with numerator $f_C$ and last member $a$ the iterated constant term becomes $\pm\big(\partial_{\varepsilon_a}f_C(0)-\lambda_af_C(0)\big)$: changing coordinates shifts the finite part by $-\log\beta_{\varsigma(a)}$ times the signed residue of each wall closed by a mark whose gauge differs from the chart's.
\end{proposition}

% Reference 20: thm:wallres; source line 3195
\begin{theorem}[Residue at a wall]\label{thm:wallres}
Let $\omega$ be projective, $S$ a proper bunch with room $R$, $g\in R$ and $\alpha_S=\rho u$. In this statement write $K^{[l]}_j=K^{(d_j(S)+l)}_{j,S}(u,\alpha_R)$ and $\Psi^{[l]}_j=\Psi^{(d^\ast_j(S)+l)}_{j,S}(u,\alpha_R)$ for the level and colevel letters of Definition~\ref{def:levels} counted from their floors, so that $K_j(\rho u,\alpha_R)=\rho^{d_j(S)}\sum_{l\ge0}\rho^lK^{[l]}_j$ and $\Psi_j(\rho u,\alpha_R)=\rho^{d^\ast_j(S)}\sum_{l\ge0}\rho^l\Psi^{[l]}_j$. At points of $\kappa_S=-q$, $q\in\mathbb N$, where every other proper wall is positive,
\[
\operatorname*{Res}_{\kappa_S=-q}P_F(\mvec;\mathbf s,\mathbf t)=\int_{\alpha_g=1}\int_{\Delta_S}u^{\mvec|_S-\one}\alpha_R^{\mvec|_R-\one}\prod_j\big(K^{[0]}_j\big)^{-s_j}\big(\Psi^{[0]}_j\big)^{-t_j}\,[\rho^q]h\ d\omega_S\,d\alpha_{R-g},
\]
where
\begin{gather*}
h=\prod_j\Big(\sum_l\rho^lK^{[l]}_j/K^{[0]}_j\Big)^{-s_j}\Big(\sum_l\rho^l\Psi^{[l]}_j/\Psi^{[0]}_j\Big)^{-t_j},\\
[\rho^q]\Big(1+\sum_{l\ge1}\rho^lX_l\Big)^{-s}=\sum_{\sum_lln_l=q}(-1)^{|\mathbf n|}\frac{\Gamma(s+|\mathbf n|)}{\Gamma(s)}\prod_l\frac{X_l^{n_l}}{n_l!}.
\end{gather*}
If the weights factor on the bottom levels, then with the \emph{faces} of the wall, the families
\[
F^S=\big((\cB_j/R;\,s_j,0),\ (\cB_j|S;\,0,t_j)\big)_j\ \text{on }S,\qquad F^R=\big((\cB_j|R;\,s_j,0),\ (\cB_j/S;\,0,t_j)\big)_j\ \text{on }R
\]
(Definition~\ref{def:minors}), the residue at $\kappa_S=0$ is the product of their periods at the masses $\mvec|_S$ and $\mvec|_R$; both faces are balanced, the walls of $F^S$ are the $\kappa_U$, $U\subseteq S$, and those of $F^R$ are the $\kappa_{S\cup U}$, $U\subseteq R$, on the wall. Here $P_F$ is the meromorphic function of Lemma~\ref{lem:periodmero}; on a connected family with real data and $s_j,t_j\ge0$, not all zero, it is the normaliser of the distribution of Theorem~\ref{thm:offdomain}(i).
\end{theorem}

% Reference 20: prop:atwall; source line 3215
\begin{proposition}[A molecule at a nonpositive wall]\label{prop:atwall}
Let a molecule carry real data with $s_j,t_j\ge0$ and target $\btau$, and let $S$ be a proper bunch with room $R$, $\alpha_S=\rho u$ on a chart $\alpha_g=1$, $g\in R$, with $H(\rho)=H(\rho;u,\alpha_R)$ the native completion $W\,d\alpha_{\btau}$ divided by $\rho^{\kappa_S(\btau)-1}d\rho\,d\omega_S\,d\alpha_{R-g}$ and $H_l=[\rho^l]H$ as in Theorem~\ref{thm:wallres}.
\begin{enumerate}[label=(\roman*)]
\item Along the collapse of $S$, $\Phi(y_0-\lambda\one_S)-\ip{\btau}{y_0-\lambda\one_S}=\lambda\,\kappa_S(\btau)+O(1)$ as $\lambda\to\infty$, the left side being the negative logarithm of the density in $y=\log\alpha$ (Theorem~\ref{thm:saddle}); so the density does not decay along the collapse exactly when $\kappa_S(\btau)\le0$, and then the form is not integrable at $\rho=0$.
\item With $\langle\cdot\rangle$ the integral over $u$ and $\alpha_R$, if every other proper wall at $\btau$ is positive, $\kappa_S(\btau)$ is not a nonpositive integer, and $L\in\mathbb N$ has $\kappa_S(\btau)+L>0$,
\[
P_F(\btau;\mathbf s,\mathbf t)=\sum_{l<L}\frac{\langle H_l\rangle}{\kappa_S(\btau)+l}+\int_0^1\rho^{\kappa_S(\btau)-1}\Big(\langle H\rangle-\sum_{l<L}\rho^l\langle H_l\rangle\Big)d\rho+\int_1^\infty\rho^{\kappa_S(\btau)-1}\langle H\rangle\,d\rho ,
\]
the continuation of the period past the wall of $S$; its poles at $\kappa_S=-q$ have the residues $\langle H_q\rangle$ of Theorem~\ref{thm:wallres}. When the weights factor on the bottom levels, $\langle H_0\rangle$ is the product of the integrals of the faces $F^S$ at $\btau|_S$ and $F^R$ at $\btau|_R$, whose imbalances are $\kappa_S(\btau)$ and $-\kappa_S(\btau)$: the leading term is the compound of Proposition~\ref{prop:compoundlaw}.
\end{enumerate}
\end{proposition}

% Reference 20: cor:corner; source line 3230
\begin{corollary}[The vertex]\label{cor:vertex}
At $S=\cE-e$, with $e$ neither a loop nor a coloop in any table, the face $F^R$ is the single mark $e$ and contributes $1$, and
\[
\langle H_0\rangle=\int_{\Delta_{\cE-e}}u^{\btau|_{\cE-e}-\one}\prod_jK_{j,/e}(u)^{-s_j}\Psi_{j,\setminus e}(u)^{-t_j}\,d\omega_{\cE-e}(u),
\]
the integral of the \emph{vertex child}: the contractions at $e$ on the basis side and the deletions on the cobasis side. Proposition~\ref{prop:atwall}(ii) at $S=\cE-e$ removes the leading terms of the vertex from the period exactly.
\end{corollary}

\section{The pseudointegral}

% Reference 20: def:perioddatum; source line 4498 (lists).
% Proved in the last part: def:perioddatum.
\begin{definition}[Period datum; moves]\label{def:perioddatum}
Let a family of tables with positive rational weights have rational exponents and masses and lie on the balance wall. Its \emph{period datum} is the positive orthant of $\mathbb R^{\cE}$, read projectively, with the form $W\,d\alpha_{\mvec}=\alpha^{\mvec-\one}\prod_jK_j^{-s_j}\Psi_j^{-t_j}d\alpha$, algebraic over $\mathbb Q$ and homogeneous of degree zero, the hypersurface $\{\prod_jK_j\Psi_j=0\}$ and the coordinate hyperplanes; the face $\{\alpha_S=0\}$ carries the wall $\kappa_S$ (Theorem~\ref{thm:exist}). A \emph{move} between period data is a change of variables by a birational map defined over $\mathbb Q$ carrying the open domain bijectively onto the open domain, faces to faces and form to form; a product of data; a linear combination; or the boundary term of an exact form.
\end{definition}

% Reference 20: def:pseudoint; source line 3239 (lists in (i)).
% Proved in the last part: def:pseudoint.
\begin{definition}[The pseudointegral over a region]\label{def:pseudoint}
Let $F$ be a family with real data, $s_j,t_j\ge0$ not all zero, masses $\mvec>0$ and weights factoring across its molecules, with form $\omega=W\,d\alpha_{\mvec}$.
\begin{enumerate}[label=(\roman*)]
\item A \emph{region} of $\omega$ is a domain $\mathcal R$ in coordinates $x$ with a change of variables mapping it bijectively onto a section $\{h=1\}$ of the orthant, $h>0$ homogeneous of degree one: a chart $\alpha_g=1$, the simplex, a bunch section $|\alpha_B|=1$. The \emph{density} of $\omega$ on $\mathcal R$ is the function $f$ for which $f(x)\,dx$ is the form that $\omega$ induces on the section (Lemma~\ref{lem:chart}). A region of the damped form $e^{-|\alpha|}\omega$ is the same with the orthant in place of the section. For a family of tables with positive rational weights and rational data on the balance wall the regions reached from these by moves (Definition~\ref{def:perioddatum}) are domains cut out by polynomial inequalities over $\mathbb Q$.
\item The \emph{pseudointegral} of $\omega$ over a region is
\[
\pseudoint\!\!\cdots\!\!\pseudoint_{\mathcal R}f(x)\,dx:=\pseudoint W\,d\alpha_{\mvec}
\]
(Definition~\ref{def:threeE}(iii)), and over a region of the damped form, away from the poles of $J_F$, it is the continued interest, $J_F(\mvec;\mathbf s,\mathbf t)=\pseudoint W\,e^{-|\alpha|}d\alpha_{\mvec}$ (Definition~\ref{def:continued}), written $\pseudoint\alpha^{\mvec-\one}e^{-|\alpha|}\prod_jK_j^{-s_j}\Psi_j^{-t_j}d\alpha$ when the integrand is displayed. The value belongs to $\omega$: the native completion and the ray of Definition~\ref{def:threeE}(iii) are taken in the table's coordinates $\alpha$, and a presentation (Theorem~\ref{thm:invariance}) declares them anew and is another form; the region and its coordinates say how the integration is presented, and a density determines $\omega$ only together with the family, a chart $\alpha_g=1$ not seeing the mass $m_g$ and the simplex not seeing a power of $|\alpha|$, the letter of the list $\binom{\cE}1$ (Proposition~\ref{prop:sectiontable}). Divided by the period of the centre, $P_{\mathrm c}$ on a molecule, it is the pseudoexpectation $\mathsf E_{U\sim\mathbb P_{\mathrm c}}[(\omega_{\btau}/\omega_{\mathrm c})(U)]$.
\item The pseudointegral of $\omega$ is read in the Foreman's order (Definition~\ref{def:procedure}). The tree is drawn from the tables at the centre of the table's gauge; on each molecule the measure $d\alpha_{\mvec}$ is set to $d\alpha_{\btau}$; the bunches are integrated in the Foreman's order, each at once in scale and shape; where the target lies outside a wall the value is the continuation in the masses, and on a pole hyperplane its constant term along the ray of the masses. The damped form takes no completion, and the Foreman computes it on the ray beyond the target (Theorem~\ref{thm:foremanE}).
\end{enumerate}
\end{definition}

% Reference 20: thm:pseudoint; source line 3252
\begin{theorem}[The pseudointegral and the native completion]\label{thm:pseudoint}
Let $\omega=W\,d\alpha_{\mvec}$ have real data on a connected family, $s_j,t_j\ge0$ not all zero and $\mvec>0$, with target $\btau$.
\begin{enumerate}[label=(\roman*)]
\item With a native gauge $\alpha^\star$ the pseudointegral is read at the realization $\Upsilon=((s_j,t_j)_j;\alpha^\star)$, whose masses are $\btau$ and whose proper walls are positive. The form $W\,d\alpha_{\btau}$ is positive with convergent period, and every iterated integral of it over every section, in any coordinates, in the Foreman's order and in every other, converges absolutely to
\[
\pseudoint W\,d\alpha_{\mvec}=\int W\,d\alpha_{\btau}=P_F(\btau;\mathbf s,\mathbf t):
\]
no finite part is taken, and the pseudoexpectation is the expectation of a positive function, $\mathsf E_{U\sim\mathbb P_{\mathrm c}}[(\omega_{\btau}/\omega_{\mathrm c})(U)]=\oE_{U\sim\mathbb P_{\mathrm c}}[(\omega_{\btau}/\omega_{\mathrm c})(U)]$.
\item If $\btau$ lies on no pole hyperplane, every regular schedule of the native completion returns the pseudointegral,
\[
\Reg^{\mathfrak s}\!\!\int W\,d\alpha_{\btau}=\pseudoint W\,d\alpha_{\mvec}=P_F(\btau;\mathbf s,\mathbf t),
\]
each iterated integral a finite part (Definition~\ref{def:Reg}), and the value does not depend on the order; its quasi-expectation is the pseudoexpectation, $\mathfrak E^{\mathfrak s}_{U\sim\mathbb P_{\mathrm c}}[(\omega_{\btau}/\omega_{\mathrm c})(U)]=\mathsf E_{U\sim\mathbb P_{\mathrm c}}[(\omega_{\btau}/\omega_{\mathrm c})(U)]$.
\item On a pole hyperplane $\pseudoint W\,d\alpha_{\mvec}=\CT_{\eta=0}P_F(\btau+\eta\mvec;\mathbf s,\mathbf t)$ exists and depends on $\mvec$ only through its ray. It takes no chart and no order: its regulator is $\eta$, along the ray of the masses, and its section is the simplex that the damping singles out, $P_F(\mvec;\mathbf s,\mathbf t)=J_F(\mvec;\mathbf s,\mathbf t)/\Gamma(\kappa_{\cE})$ (Theorem~\ref{thm:projective}).
\end{enumerate}
\end{theorem}

% Reference 20: prop:polecoords; source line 3274
\begin{proposition}[A pole target in other coordinates]\label{prop:polecoords}
In Theorem~\ref{thm:pseudoint}(iii) let $x=\lambda^{-1}\alpha$ be the coordinates of a presentation (Theorem~\ref{thm:invariance}), $W^\lambda(x)=W(\lambda x)$, and
\[
G_k(\eta)=\int_{\Delta_{\cE}}\big(\log|\lambda u|\big)^k\,u^{\btau+\eta\mvec-\one}W^\lambda(u)\,d\omega_{\cE}(u)\qquad(k\in\mathbb N),
\]
continued in $\eta$, so that $G_0(\eta)$ is the integral over the simplex of the presented word $W^\lambda$ at the masses $\btau+\eta\mvec$. Each $G_k$ has at $\eta=0$ a pole of order at most $n-1$, and the presentation has pseudointegral
\[
(\lambda^{\mvec})^{\kappa_{\cE}/m_{\cE}}\Big(\pseudoint W\,d\alpha_{\mvec}-\lambda^{\btau}\sum_{1\le i+k\le n-1}\frac{\ip{\mvec}{\log\lambda}^i}{i!}\,\frac{(-m_{\cE})^k}{k!}\,[\eta^{-i-k}]\,G_k\Big),
\]
$[\eta^{-i-k}]$ the Laurent coefficient at $\eta=0$. Off the pole hyperplanes every $G_k$ is holomorphic at $0$ and this is Theorem~\ref{thm:invariance}; on them the finite part is taken in the coordinates of the table, as every finite part is taken in its own coordinate (Lemma~\ref{lem:Regscale}).
\end{proposition}

% Reference 20: prop:sectionsoff; source line 3297 (a bunch section with $-\kappa_{\cE}\in\mathbb N$).
% Proved in the last part: prop:sectionsoff.
\begin{proposition}[Sections off the balance wall]\label{prop:sectionsoff}
Let $\omega=W\,d\alpha_{\mvec}$ have real data, $s_j,t_j\ge0$ not all zero, $\mvec>0$ and imbalance $\kappa_{\cE}\neq0$.
\begin{enumerate}[label=(\roman*)]
\item The integral of $\omega$ over a section $\{h=1\}$ is the integral over every section of the projective form $h^{-\kappa_{\cE}}\omega$, wherever it converges. Over a chart it is the period at the chart completion (Proposition~\ref{prop:snip}(i)), over a bunch section $|\alpha_B|=1$ the period of the family with the list $\binom B1$, the marks outside $B$ as loops, appended at the exponents $(\kappa_{\cE},0)$, a multinomial sum of completions when $-\kappa_{\cE}\in\mathbb N$ (Proposition~\ref{prop:snip}(iii)), and over the simplex, where the proper walls at $\mvec$ are positive, $P_F(\mvec;\mathbf s,\mathbf t)=\Gst{n}{1}{\mvec}{m_{\cE}}\,\oE_{U\sim\Dir_{\mvec}}[W(U)]$ (Proposition~\ref{prop:gammadraws}(i)). Different sections integrate different projective forms; the pseudointegral integrates the native completion over all of them.
\item The walls of the chart completion at $g$ and of the native completion are
\[
\kappa_S(\mvec^{(g)})=\kappa_S(\mvec)-\kappa_{\cE}\,\one_S(g),\qquad\kappa_S(\btau)=\kappa_S(\mvec)-\frac{\kappa_{\cE}}{m_{\cE}}\,m_S :
\]
the snip moves the walls through $g$ by the whole imbalance and leaves the others, and the native completion moves each wall by its share. For $\kappa_{\cE}>0$ a proper wall through $g$ is lower at the chart completion than at the target by $\kappa_{\cE}(1-m_S/m_{\cE})$, and one avoiding $g$ is higher by $\kappa_{\cE}m_S/m_{\cE}$. So a snip can put the walls through $g$ on pole hyperplanes, or past them, while the target is inside every wall.
\end{enumerate}
\end{proposition}

% Reference 20: thm:flatsection; source line 3312
\begin{theorem}[Integrals over sections]\label{thm:flatsection}
Let $\omega=W\,d\alpha_{\mvec}$ have real data, $s_j,t_j\ge0$ not all zero and $\mvec>0$, and for $h>0$ homogeneous of degree one let $\int_{h=1}\omega$ denote the integral over the section $\{h=1\}$ with the measure of Lemma~\ref{lem:chart}.
\begin{enumerate}[label=(\roman*)]
\item By Proposition~\ref{prop:sectionsoff}(i) the integral of $\omega$ over a section $\{h=1\}$ is the integral over the rays of $h^{-\kappa_{\cE}}\omega$, and for $\kappa_{\cE}\neq0$, $h\mapsto h^{-\kappa_{\cE}}$ carries the sections bijectively onto the positive multipliers of degree $-\kappa_{\cE}$. The monomial sections $\{\alpha^w=1\}$, $w_{\cE}=1$, are the completions, $(\alpha^w)^{-\kappa_{\cE}}\omega=W\,d\alpha_{\mvec-\kappa_{\cE}w}$: they act on the measure and keep the word.
\item The native completion is the section $\{\alpha^{\mvec}=1\}$, the hyperplane $\ip{\mvec}{y}=0$ in $y=\log\alpha$, where the density $e^{\ip{\mvec}y}$ of $d\alpha_{\mvec}$ against $dy$ is $1$; it is the one monomial section with this property. With a native gauge on a connected family of at least two marks, for every mark $g$,
\[
\pseudoint W\,d\alpha_{\mvec}=\frac{m_{\cE}}{m_g}\int_{\mathbb R^{\cE-g}}W(e^{y})\,dy_{\cE-g},\qquad y_g=-\frac1{m_g}\sum_{e\neq g}m_ey_e ,
\]
and $\mathsf E_{U\sim\mathbb P_{\mathrm c}}[(\omega_{\btau}/\omega_{\mathrm c})(U)]$ is this integral over the integral of $\omega_{\mathrm c}$ over the same section, which is $P_{\mathrm c}$.
\item The simplex $\{|\alpha|=1\}$ is not a monomial section for $n\ge2$: its multiplier is $|\alpha|^{-\kappa_{\cE}}$, which keeps the measure and multiplies the word by a power of the letter $K_0=|\alpha|$ of the list $\binom{\cE}1$. The integral over it is the integral of the projective form $K_0^{-\kappa_{\cE}}W\,d\alpha_{\mvec}$, the family with that table added at the exponents $(\kappa_{\cE},0)$ (Proposition~\ref{prop:sectiontable}, Proposition~\ref{prop:snip}(iii) at $B=\cE$), and wherever every proper wall at $\mvec$ is positive it is the period $P_F(\mvec;\mathbf s,\mathbf t)$ of Proposition~\ref{prop:gammadraws}(i), the law $\Dir_{\mvec}$ there being the law of $\binom{\cE}1$ at a full basis lean.
\item On the balance wall every section reads $\int W\,d\alpha_{\mvec}$ (Theorem~\ref{thm:projective}).
\end{enumerate}
\end{theorem}

% Reference 20: prop:centrelaw; source line 3335
\begin{proposition}[The centre sets the law]\label{prop:centrelaw}
Let $\omega=W\,d\alpha_{\mvec}$ have real data on a connected family of $n\ge2$ marks, $s_j,t_j\ge0$ not all zero and $\mvec>0$, with target $\btau$, centre $\omega_{\mathrm c}$ and law of the centre $\mathbb P_{\mathrm c}$ (Definition~\ref{def:threeE}). For laws $\mathbb P,\mathbb P'$ with $\mathbb P$ absolutely continuous with respect to $\mathbb P'$ write $\mathrm{KL}(\mathbb P\,\|\,\mathbb P')=\oE_{U\sim\mathbb P}\big[\log\frac{d\mathbb P}{d\mathbb P'}(U)\big]$, the Kullback--Leibler divergence.
\begin{enumerate}[label=(\roman*)]
\item The native completion is the centre times the departure, $\omega_{\btau}=\frac{\omega_{\btau}}{\omega_{\mathrm c}}\,\omega_{\mathrm c}$, and with a native gauge the Foreman's law, the shape law $\mathrm{Breta}^\circ(F;\btau;\mathbf s,\mathbf t)$ of the native completion, is the law of the centre reweighted by the departure:
\[
\mathrm{Breta}^\circ(F;\btau;\mathbf s,\mathbf t)(du)=\frac{(\omega_{\btau}/\omega_{\mathrm c})(u)}{\mathsf E_{U\sim\mathbb P_{\mathrm c}}[(\omega_{\btau}/\omega_{\mathrm c})(U)]}\,\mathbb P_{\mathrm c}(du).
\]
\item With a native gauge,
\[
\log\mathsf E_{U\sim\mathbb P_{\mathrm c}}\Big[\frac{\omega_{\btau}}{\omega_{\mathrm c}}(U)\Big]=\mathrm{KL}\big(\mathbb P_{\mathrm c}\,\big\|\,\mathrm{Breta}^\circ(F;\btau;\mathbf s,\mathbf t)\big)\ \ge\ 0,
\]
with equality exactly when $\mathrm{Breta}^\circ(F;\btau;\mathbf s,\mathbf t)=\mathbb P_{\mathrm c}$. So every pseudointegral with a native gauge is at least the period of its centre, $\pseudoint W\,d\alpha_{\mvec}\ge P_{\mathrm c}$, and among the forms of given strengths $\sigma_j$ with a native gauge the centre has the least value.
\item The pseudoexpectation is kept by the fold: the two forms of Theorem~\ref{thm:fold}, at the masses $\mvec$ with the exponents $(\mathbf s,\mathbf t)$ and at the masses $\lambda(2\sigma\one-\btau)$ with $(\mathbf t,\mathbf s)$, have one centre and one pseudoexpectation, and the same holds for continued values off the pole hyperplanes where $2\sigma\one-\btau>0$.
\item Let $\omega$ be the form of a realization at the table's gauge, $\alpha^\star=\one$ (Definition~\ref{def:realization}). In signed coordinates on $\one^\perp$ the departure is
\[
\frac{\omega_{\btau}}{\omega_{\mathrm c}}(e^{2z})=\exp\Big(-\sum_j\delta_j\tilde O_j(z)\Big),\qquad\tilde O_j(z)=O_j(z)-\ip{\nabla O_j(0)}z ,
\]
each $\tilde O_j$ odd with value, gradient and Hessian $0$ at $0$. The Foreman's law and the law of the centre have their mode at the table's gauge and the same Hessian of the logarithm of the density there, and the logarithm of their ratio is, up to a constant, $-\sum_j\delta_j\tilde O_j$: minus the odd part of the centred potential less its linear term, whose Taylor coefficients at the table's gauge are minus the odd rungs of order at least three of the potential, the rungs that vanish at the centre (Proposition~\ref{prop:centrefamily}(iii)).
\end{enumerate}
\end{proposition}

% Reference 20: prop:completionlaws; source line 3366
\begin{proposition}[The list $\binom{\cE}1$ as a normaliser]\label{prop:completionlaws}
Let $\omega=W\,d\alpha_{\mvec}$ have real data on a connected family of $n\ge2$ marks, $s_j,t_j\ge0$ not all zero and $\mvec>0$, so that $W$ is homogeneous of degree $-\tau_{\cE}$; let $K_0=|\alpha|$ be the letter of the list $\binom{\cE}1$, whose bases are the singletons.
\begin{enumerate}[label=(\roman*)]
\item Every completion $W\,d\alpha_{\mvec'}$, $\mvec'=\mvec-\kappa_{\cE}w$, has total mass $m'_{\cE}=\tau_{\cE}$. The native completion is the one whose masses keep the proportions of $\mvec$, $\btau/\tau_{\cE}=\mvec/m_{\cE}$, and they are positive, while the chart completion at $g$ has the mass $m_g-\kappa_{\cE}$ at $g$.
\item For $\mvec'>0$ with $m'_{\cE}=\tau_{\cE}$ the form $K_0^{-\tau_{\cE}}d\alpha_{\mvec'}$ is projective, with integral $\Gst{n}{1}{\mvec'}{\tau_{\cE}}$ over every section: the period of $\binom{\cE}1$ at the full basis lean $(\tau_{\cE},0)$, whose law is $\Dir_{\mvec'}$ (Proposition~\ref{prop:sectiontable}). So $\Gst{n}{1}{\btau}{\tau_{\cE}}=\pseudoint K_0^{-\tau_{\cE}}d\alpha_{\mvec}$; for a schedule with $\mvec^{(g)}>0$, $\Gst{n}{1}{\mvec^{(g)}}{\tau_{\cE}}=\Reg^{\mathfrak s}\!\!\int K_0^{-\tau_{\cE}}d\alpha_{\mvec}$, a convergent integral; and with a native gauge
\[
\pseudoint W\,d\alpha_{\mvec}\Big/\pseudoint K_0^{-\tau_{\cE}}d\alpha_{\mvec}=\oE_{U\sim\Dir_{\btau}}[W(U)] .
\]
\item With a native gauge, and for continued values off the pole hyperplanes where $2\sigma\one-\btau>0$, the fold keeps the pseudointegral and carries the target $\btau$ to $2\sigma\one-\btau$ (Theorem~\ref{thm:fold}), so the quotients of (ii) for $\omega$ and for its image are $P_F(\btau;\mathbf s,\mathbf t)$ divided by $\Gst{n}{1}{\btau}{\tau_{\cE}}$ and by $\Gst{n}{1}{2\sigma\one-\btau}{\tau^\ast_{\cE}}$, $\tau^\ast_{\cE}=2n\sigma-\tau_{\cE}$. The second quotient over the first is $\Gst{n}{1}{\btau}{\tau_{\cE}}\big/\Gst{n}{1}{2\sigma\one-\btau}{\tau^\ast_{\cE}}$, which is not $1$ in general. The period of $\binom{\cE}1$ is that of a full basis lean of a table outside the family, and it prices a lean and its mirror differently; the centre, which the fold fixes, prices them alike (Proposition~\ref{prop:centrelaw}(iii)).
\item Along the ray of the masses, wherever the proper walls at $\lambda\mvec$ are positive, the integral of the projective form $K_0^{-(\lambda m_{\cE}-\tau_{\cE})}W\,d\alpha_{\lambda\mvec}$ is $\Gst{n}{1}{\lambda\mvec}{\lambda m_{\cE}}\,\oE_{U\sim\Dir_{\lambda\mvec}}[W(U)]$ (Theorem~\ref{thm:flatsection}(iii), Proposition~\ref{prop:gammadraws}(i)). The power of $K_0$ vanishes exactly at $\lambda=\tau_{\cE}/m_{\cE}$, where the form is the native completion $W\,d\alpha_{\btau}$.
\end{enumerate}
\end{proposition}

% Reference 20: cor:pseudoquasi; source line 3389
\begin{corollary}[The pseudointegral is the quasi-integral that always exists]\label{cor:pseudoquasi}
For real data with $s_j,t_j\ge0$ not all zero, $\mvec>0$ and weights factoring across the molecules, the pseudointegral $\pseudoint W\,d\alpha_{\mvec}$ and the pseudoexpectation $\mathsf E_{U\sim\mathbb P_{\mathrm c}}[(\omega_{\btau}/\omega_{\mathrm c})(U)]$ are defined. On a molecule whose target lies on no pole hyperplane the pseudointegral is the quasi-integral of the native completion and the pseudoexpectation its quasi-expectation, the same along every regular schedule; on a pole hyperplane, where the regular schedules of the native completion can depend on the order (Theorem~\ref{thm:last}), it is the one constant term along the ray of the masses. On each molecule every generalized expectation divides by the one period of its centre, a positive convergent integral (Proposition~\ref{prop:centrefamily}); when the molecule has a native gauge, the native completion is the one completion that keeps it (Theorem~\ref{thm:transport}(i)). The quasi-integrals of $\omega$ itself are Laurent data at its chart completions (Theorem~\ref{thm:bridge}), which off the balance wall can depend on the chart and on the order.
\end{corollary}

The pseudointegral names a procedure as well as a value. Written for a family whose realization is specified later, $\pseudoint\!\cdots\!\pseudoint_{\mathcal R}$ commits to the Foreman's order at every realization: the tree is drawn once from the tables, and the realization $((s_j,t_j)_j;\alpha^\star)$ of the point fixes the measure $d\alpha_{\btau}$ and where the contours of the ports run. An iterated integral in another order, or after a snip, computes a quasi-integral and not the pseudointegral, even where the numbers agree; when they agree, that is Theorem~\ref{thm:pseudoint}(i)--(ii), not the order that was followed.

On a pole hyperplane the regulator $\eta$ of the pseudointegral enters as $u^{\eta\mvec}$, as in Theorem~\ref{thm:RegCT}. Past a wall the pseudointegral normalizes the distribution of Theorem~\ref{thm:offdomain}(i).


\chapter{The family as one Laplace transform}\label{ch:questions}

On its balance wall a family is one Laplace transform. In the pair coordinates of its data and the signed coordinates of the shape, its period is the Laplace transform, over the tilts, of the even and odd statistics against the offsets (Definition~\ref{def:question}; for the interest, Theorem~\ref{thm:pairform}); the Breta shape laws at the points of the balance wall form an exponential family in these coordinates, and every ratio of periods is a moment generating function of the statistics. On the sector of every order the exponent of a question is minus twice its wall at the tilt, the walls of the tails weighted by the differences of the tilt along the order, less the logarithms of the sector letters, which are bounded. The Laplace transform is read through that one reading of the walls: where the exponent is largest and how it falls, how the period behaves at low and at high strength, how far a question can lean, and how fast the tails of the statistics fall.

\section{Questions and the law of the tilt}

% Reference 20: def:question; source line 11287
\begin{definition}[Tilts; questions]\label{def:question}
The space of \emph{tilts} is $\mathcal T=\mathbb R^{\cE}/\mathbb R\one$, with the Lebesgue measure $dz$ of a chart $z_g=0$, which does not depend on $g$. The \emph{tilt} of a shape $u\in\Delta_{\cE}$ is the image in $\mathcal T$ of its signed coordinates $\frac12\log u$, and a tilt is read at its representative in $\one^\perp$ when an odd statistic is applied to it, since $O_j(z+c\one)=O_j(z)+2c(2r_j-n)$ (Theorem~\ref{thm:tworeadings}(iii)). A \emph{question} to a family $F$ is a point of its balance wall with real data, written in its pair coordinates $(\bmu_\perp;\boldsymbol\sigma,\boldsymbol\delta)$: the data $(\mvec;\mathbf s,\mathbf t)=\big(\mvec_F(\bmu_\perp;\boldsymbol\sigma,\boldsymbol\delta);\,\boldsymbol\sigma+\boldsymbol\delta,\,\boldsymbol\sigma-\boldsymbol\delta\big)$ (Definition~\ref{def:family}). In the coordinates $z=\frac12\log u$ of Theorem~\ref{thm:projective},
\[
P_F(\mvec;\mathbf s,\mathbf t)=2^{n-1}\int_{\mathcal T}\exp\Big(2\ip{\bmu_\perp}z-\sum_j\big(\sigma_jE_j(z)+\delta_jO_j(z)\big)\Big)\,dz .
\]
When every proper wall is positive, the tilt $Z$ of $U\sim\mathrm{Breta}^\circ(F;\mvec;\mathbf s,\mathbf t)$ has the density
\[
2^{n-1}\exp\Big(2\ip{\bmu_\perp}z-\sum_j\big(\sigma_jE_j(z)+\delta_jO_j(z)\big)\Big)\Big/P_F(\mvec;\mathbf s,\mathbf t)
\]
with respect to $dz$: this is the \emph{law of the question}. A direction $v=(\dot\mvec;\dot{\mathbf s},\dot{\mathbf t})$ of the data is \emph{balanced} when $\dot m_{\cE}=\sum_j(\dot s_jr_j+\dot t_jr^\ast_j)$; its pair coordinates $(c;\mathbf a,\mathbf b)$ have $\dot\mvec=\mvec_F(c;\mathbf a,\mathbf b)$, $\dot{\mathbf s}=\mathbf a+\mathbf b$ and $\dot{\mathbf t}=\mathbf a-\mathbf b$, and $D_v$ is the derivative along $v$. The derivatives $\partial_{\sigma_j}$, $\partial_{\delta_j}$ and $\nabla_{\bmu_\perp}$ of a function of the question are those along the balanced directions with the pair coordinates $(0;\mathbf e_j,\mathbf 0)$, $(0;\mathbf 0,\mathbf e_j)$ and $(\mathbf e_e-\frac1n\one;\mathbf 0,\mathbf 0)$, that is $(\one;\mathbf e_j,\mathbf e_j)$, $\big(\frac{2r_j-n}n\one;\mathbf e_j,-\mathbf e_j\big)$ and $\big(\mathbf e_e-\frac1n\one;\mathbf 0,\mathbf 0\big)$. The centre of a question is the question $(\sigma\one;\boldsymbol\sigma,\boldsymbol\sigma)$ with the same strengths, the data of the family at the centre (Definition~\ref{def:realization}), of pair coordinates $(0;\boldsymbol\sigma,\mathbf 0)$; on a connected family its law is the law of the centre $\mathbb P_{\mathrm c}$ of Definition~\ref{def:threeE} read on the tilts.
\end{definition}

% Reference 20: cor:letterstransform; source line 11295
\begin{corollary}[The Laplace transform of the tilt and the statistics]\label{cor:letterstransform}
Let a question $(\mvec;\mathbf s,\mathbf t)$ have every proper wall positive. In this statement $\oE$ is the expectation of $U\sim\mathrm{Breta}^\circ(F;\mvec;\mathbf s,\mathbf t)$ and $Z$ is the tilt of $U$. For $c\in\one^\perp\otimes\mathbb C$ and $\mathbf a,\mathbf b\in\mathbb C^N$ such that the real parts of the shifted data $\big(\mvec+\mvec_F(c;\mathbf a,\mathbf b);\,\mathbf s+\mathbf a+\mathbf b,\,\mathbf t+\mathbf a-\mathbf b\big)$ have every proper wall positive,
\[
\oE\Big[e^{2\ip cZ-\sum_j(a_jE_j(Z)+b_jO_j(Z))}\Big]=\frac{P_F\big(\mvec+\mvec_F(c;\mathbf a,\mathbf b);\,\mathbf s+\mathbf a+\mathbf b,\,\mathbf t+\mathbf a-\mathbf b\big)}{P_F(\mvec;\mathbf s,\mathbf t)},
\]
analytic in $(c,\mathbf a,\mathbf b)$ there: the shift moves the pair coordinates of the question by $(c;\mathbf a,\mathbf b)$. So the joint cumulants of order $k$ of $(E_j(Z),O_j(Z))_j$ are $(-1)^k$ times the derivatives of order $k$ of $\log P_F(\mvec;\mathbf s,\mathbf t)$ in $(\sigma_j,\delta_j)_j$, and their characteristic function is the same ratio at $c=0$ and purely imaginary $(\mathbf a,\mathbf b)$: the family on the fibre of imaginary exponents over the question.
\end{corollary}

% Proved in the last part: def:tiltwall (new).
\begin{definition}[The wall at a tilt]\label{def:tiltwall}
For an order $\pi$ of $\cE$ with tails $U_k=\{\pi(k),\dots,\pi(n)\}$, the sector $H_\pi$ of Lemma~\ref{lem:greedy} is read in the tilts as $\{z\in\mathcal T:\ z_{\pi(1)}\ge\cdots\ge z_{\pi(n)}\}$. On it the measure $dz$ of Definition~\ref{def:question} is the Lebesgue measure in the differences $z_{\pi(k-1)}-z_{\pi(k)}\ge0$, $2\le k\le n$, which at the gauge $\alpha=e^{2z}$ are half the gaps $\log(\alpha_{\pi(k-1)}/\alpha_{\pi(k)})$ of Corollary~\ref{cor:pastwalls}. For data on the balance wall the \emph{wall at a tilt} $z\in H_\pi$ is
\[
\kappa(z)=\sum_{k=2}^n\kappa_{U_k}\big(z_{\pi(k-1)}-z_{\pi(k)}\big),
\]
the walls of the tails of $\pi$ weighted by the differences of $z$ along $\pi$, half the walls weighted by the gaps. It does not depend on the order whose sector holds $z$, the tails with a positive difference being the sets $\{e:\ z_e<z_f\}$ for the $f\in\cE$ with $z_f>\min_ez_e$; it is positively homogeneous of degree one and linear on every sector, and $\kappa(-\one_S)=\kappa_S$ for every nonempty proper $S$, with $-\one_S$ read in $\mathcal T$, the direction along which $S$ collapses (Proposition~\ref{prop:atwall}(i)). In pair coordinates $\kappa(z)$ is linear in the question.
\end{definition}

% Reference 20: prop:evenfloor; source line 11306 (one table at a centre).
% Proved in the last part: prop:evenfloor.
\begin{proposition}[The exponent of a question reads its walls; the even statistic is least at the table's gauge]\label{prop:evenfloor}
Let a family carry a question $(\bmu_\perp;\boldsymbol\sigma,\boldsymbol\delta)$, its exponents of any signs, and in this statement write $X(z)=2\ip{\bmu_\perp}z-\sum_j\big(\sigma_jE_j(z)+\delta_jO_j(z)\big)$ for the exponent of the law of the question (Definition~\ref{def:question}) and $\kappa(z)$ for its wall at the tilt $z$ (Definition~\ref{def:tiltwall}).
\begin{enumerate}[label=(\roman*)]
\item On the sector $H_\pi$ of every order $\pi$, with the sector letters $k_{j,\pi},\psi_{j,\pi}$ of Theorem~\ref{thm:exist}(ii) read at the sector coordinates $q_k=\alpha_{\pi(k)}/\alpha_{\pi(k-1)}=e^{-2(z_{\pi(k-1)}-z_{\pi(k)})}$ of $\alpha=e^{2z}$,
\[
X(z)=-2\kappa(z)-\sum_j\big(s_j\log k_{j,\pi}(q)+t_j\log\psi_{j,\pi}(q)\big).
\]
With $h_j(z)=\max_{T\in\cB_j}\ip{\chi_T}z$, the support function of the Newton polytope $2P_j-\one$ of $\Lambda_j$ (Theorem~\ref{thm:newton}), $2\kappa(z)=\sum_j\big(s_jh_j(z)+t_jh_j(-z)\big)-2\ip{\bmu_\perp}z$ on $\one^\perp$, and for $s_j,t_j\ge0$ the wall at a tilt is the support function of $\mathcal N_F-\mvec$, $\mathcal N_F$ the Newton sum of Theorem~\ref{thm:saddle}(ii). The sector letters are polynomials with nonnegative coefficients, with the constant terms $\theta_{j,T^+_{j,\pi}}$ and $\theta_{j,T^-_{j,\pi}}$ and the value $K_j(\one)$ at $q=\one$ (Theorem~\ref{thm:exist}(ii)), so on the cube they lie in $[\theta_{j,T^+_{j,\pi}},K_j(\one)]$ and $[\theta_{j,T^-_{j,\pi}},K_j(\one)]$; with $\theta_{j,\min}=\min_T\theta_{j,T}$, $|X+2\kappa|\le\sum_j(|s_j|+|t_j|)\max\big(|\log\theta_{j,\min}|,|\log K_j(\one)|\big)$, and for $s_j,t_j\ge0$
\[
2\kappa(z)+\sum_j(s_j+t_j)\log\theta_{j,\min}\le-X(z)\le2\kappa(z)+\sum_j(s_j+t_j)\log K_j(\one).
\]
\item $\kappa>0$ on $\mathcal T\setminus0$ exactly when every proper wall is positive, and then $\kappa(z)\ge\kappa_{\min}\big(\max_ez_e-\min_ez_e\big)$, $\kappa_{\min}$ the least proper wall.
\item At the centre of strengths $\boldsymbol\sigma$, $\sigma_j\ge0$, where $\bmu_\perp=0$ and $\boldsymbol\delta=0$, the exponent is minus the centred potential $\sum_j\sigma_jE_j$ (Proposition~\ref{prop:centrefamily}) and $2\kappa(z)=\sum_j\sigma_j\operatorname{wd}_j(z)$, with $\operatorname{wd}_j(z)=h_j(z)+h_j(-z)$ the width of $2P_j-\one$ in the direction $z$, twice the width of $P_j$, so that $\operatorname{wd}_j(\one_S)=2c_j(S)$. Every $E_j$ is convex and even, $E_j(z)\ge E_j(0)=2\log K_j(\one)$, and
\[
\sum_j\sigma_j\operatorname{wd}_j(z)+2\sum_j\sigma_j\log\theta_{j,\min}\le\sum_j\sigma_jE_j(z)\le\sum_j\sigma_j\operatorname{wd}_j(z)+2\sum_j\sigma_j\log K_j(\one).
\]
If the family is connected, $\sum_j\sigma_jE_j$ is strictly convex on $\mathcal T$ and $\sum_j\sigma_j\operatorname{wd}_j>0$ on $\mathcal T\setminus0$.
\end{enumerate}
\end{proposition}

% Reference 20: thm:severity; source line 11317 (one table, along the centre direction and at offsets zero).
% Proved in the last part: thm:severity.
\begin{theorem}[The two ends of strength]\label{thm:severity}
Let a question $q=(\bmu_\perp;\boldsymbol\sigma,\boldsymbol\delta)$ of a family have every proper wall positive, and for $\lambda\in\mathbb C$ with $\re\lambda>0$ read the data $\lambda q$ of pair coordinates $(\lambda\bmu_\perp;\lambda\boldsymbol\sigma,\lambda\boldsymbol\delta)$, a question when $\lambda$ is real, whose walls $\lambda\kappa_S$ have positive real parts; in this statement $P_F(\lambda q)$ is its period, $2^{n-1}\int_{\mathcal T}e^{\lambda X(z)}dz$ by the change of variables of Definition~\ref{def:question} with the powers of the positive letters principal (Theorem~\ref{thm:projective}), and $X$ and $\kappa$ are those of $q$ (Proposition~\ref{prop:evenfloor}).
\begin{enumerate}[label=(\roman*)]
\item \emph{Low strength.} As $\lambda\downarrow0$,
\[
\lambda^{n-1}P_F(\lambda q)\to\sum_\pi\prod_{k=2}^n\kappa_{U_k}^{-1}=(n-1)!\,\operatorname{vol}\{z\in\mathcal T:\ \kappa(z)\le1\},
\]
and
\[
\lambda^{n-1}P_F(\lambda q)=\sum_\pi\prod_{k=2}^n\kappa_{U_k}^{-1}\Big(1-\lambda\sum_j\big(s_j\log\theta_{j,T^+_{j,\pi}}+t_j\log\theta_{j,T^-_{j,\pi}}\big)\Big)+o(\lambda),
\]
with the greedy bases $T^\pm_{j,\pi}$ of Lemma~\ref{lem:greedy}: the walls set the period at low strength, and the weights enter at the next order, through the greedy bases.
\item \emph{High strength, in the half-plane.} Let $s_j,t_j\ge0$. Then $-X$ is strictly convex on $\mathcal T$, its unique critical point $z^\star$ is the native gauge of every question $\lambda q$, $\lambda>0$, and as $|\lambda|\to\infty$ with $\re\lambda\ge\lambda_0$ for a fixed $\lambda_0>0$, in particular as $\lambda\to+\infty$ and along $\lambda=\lambda_0+i\eta$, $|\eta|\to\infty$,
\[
P_F(\lambda q)=2^{n-1}e^{\lambda X(z^\star)}\Big(\frac{2\pi}\lambda\Big)^{\frac{n-1}2}\det\big(-\nabla^2X(z^\star)\big)^{-\frac12}\big(1+O(|\lambda|^{-1})\big),
\]
the power principal and the Hessian taken on a chart $z_g=0$, where it is rung two of the still of the realization $\big((s_j,t_j)_j;e^{2z^\star}\big)$ of the active tables (Definition~\ref{def:still}) with the row and the column of $g$ deleted. At the centre of strengths $\boldsymbol\sigma$, $z^\star=0$, $X(0)=-2\sum_j\sigma_j\log K_j(\one)$, and on the chart $z_g=0$, $-\nabla^2X(0)=8\sum_j\sigma_j\big(\kappa_{j,\one}(X_e,X_f)\big)_{e,f\neq g}$.
\end{enumerate}
\end{theorem}

\begin{remark}[The volume of the walls is the volume of a polar body]\label{rem:sl-polar}
At the centre of strengths $\boldsymbol\sigma$, $\sigma_j\ge0$, Proposition~\ref{prop:evenfloor}(iii) gives $2\kappa(z)=\sum_j\sigma_j\operatorname{wd}_j(z)$, and $\operatorname{wd}_j(z)=h_j(z)+h_j(-z)$, with $h_j$ the support function of $2P_j-\one$ and $P_j$ the base polytope of table $j$ (Theorem~\ref{thm:newton}), is twice the support function of $P_j-P_j$ at $z$. So $\kappa$ is the support function of the centrally symmetric polytope $\mathcal Q=\sum_j\sigma_j(P_j-P_j)$, a Minkowski sum in $\one^\perp$, read on $\mathcal T=\mathbb R^{\cE}/\mathbb R\one$, and $\{z\in\mathcal T:\ \kappa(z)\le1\}$ is its polar body $\mathcal Q^\circ=\{z:\ \ip xz\le1\text{ for every }x\in\mathcal Q\}$, bounded when every proper wall is positive (Proposition~\ref{prop:evenfloor}(ii)). By Theorem~\ref{thm:severity}(i), $\sum_\pi\prod_{k=2}^n\kappa_{U_k}^{-1}=(n-1)!\operatorname{vol}(\mathcal Q^\circ)$: the volume $V_1$ of the direction of the centre ray at the strengths $\boldsymbol\sigma$ (Proposition~\ref{prop:principal}), whose walls are those of this question, is $(n-1)!$ times the volume of the polar body of $\sum_j\sigma_j(P_j-P_j)$.
\end{remark}

At low strength the walls alone set the period, through the volume of the region where the wall at a tilt is at most one, and the weights enter at the next order through the greedy bases. At high strength the period concentrates at the native gauge, where it is a Gaussian integral whose Hessian is rung two of the still.

% Reference 20: prop:reach; source line 11332 (one table at offsets zero).
% Proved in the last part: prop:reach.
\begin{proposition}[How far a question can lean; the values of the statistics]\label{prop:reach}
For every family and every nonempty proper $S$ put $\varepsilon^\perp_j(S)=\varepsilon_j(S)-|S|\frac{2r_j-n}n$, so that $\varepsilon^\perp_j(\cE\setminus S)=-\varepsilon^\perp_j(S)$, and for a question $(\bmu_\perp;\boldsymbol\sigma,\boldsymbol\delta)$ let
\[
\kappa^{\mathbf E}_S(\boldsymbol\sigma)=\sum_j\sigma_jc_j(S),\qquad \kappa^{\mathbf O}_S(\bmu_\perp,\boldsymbol\delta)=\mu_{\perp,S}-\sum_j\delta_j\varepsilon^\perp_j(S),
\]
the even and the odd part of the wall at $S$: on the balance wall $\kappa_S=\kappa^{\mathbf E}_S+\kappa^{\mathbf O}_S$ and $\kappa_{\cE\setminus S}=\kappa^{\mathbf E}_S-\kappa^{\mathbf O}_S$, and the fold, which negates the odd part (Theorem~\ref{thm:fold}), swaps the two. So a question is inside its walls exactly when $|\kappa^{\mathbf O}_S(\bmu_\perp,\boldsymbol\delta)|<\kappa^{\mathbf E}_S(\boldsymbol\sigma)$ for every $S$. When the centre of strengths $\boldsymbol\sigma$ is inside its walls, the \emph{slope} of an odd direction $(c,\mathbf b)\in\one^\perp\times\mathbb R^N$ at $\boldsymbol\sigma$ is
\[
\ell^\ast_{\boldsymbol\sigma}(c,\mathbf b)=\min\Big\{\frac{\kappa^{\mathbf E}_S(\boldsymbol\sigma)}{|\kappa^{\mathbf O}_S(c,\mathbf b)|}:\ \kappa^{\mathbf O}_S(c,\mathbf b)\neq0\Big\},
\]
$\ell^\ast_{\boldsymbol\sigma}=\infty$ if every $\kappa^{\mathbf O}_S(c,\mathbf b)$ vanishes; it is homogeneous of degree one in $\boldsymbol\sigma$, and even and positively homogeneous of degree $-1$ in $(c,\mathbf b)$. The questions of strengths $\boldsymbol\sigma$ inside their walls are those whose odd part has $\ell^\ast_{\boldsymbol\sigma}(\bmu_\perp,\boldsymbol\delta)>1$, an open convex polyhedron symmetric under the fold, bounded exactly when no nonzero odd direction has $\kappa^{\mathbf O}_S(c,\mathbf b)=0$ for every $S$. For one connected table at offsets zero they form the wedge $\sigma>0$, $|\delta|<\ell^\ast_1(0,1)\,\sigma$, with $\ell^\ast_1(0,1)=\min\{c(S)/|\varepsilon^\perp(S)|:\ \varepsilon^\perp(S)\neq0\}$.
\begin{enumerate}[label=(\roman*)]
\item For the tilt $Z$ of $U\sim\mathrm{Breta}^\circ(F;\sigma\one;\boldsymbol\sigma,\boldsymbol\sigma)$, the centre of strengths $\boldsymbol\sigma$, and real $c\in\one^\perp$, $\mathbf a,\mathbf b\in\mathbb R^N$, the expectation of $e^{2\ip cZ-\sum_j(a_jE_j(Z)+b_jO_j(Z))}$ is finite exactly when every proper wall of the question $(c;\boldsymbol\sigma+\mathbf a,\mathbf b)$ is positive, $|\kappa^{\mathbf O}_S(c,\mathbf b)|<\kappa^{\mathbf E}_S(\boldsymbol\sigma+\mathbf a)$ for every $S$, and it is then the ratio of the periods of that question and of the centre (Corollary~\ref{cor:letterstransform}); in particular $\oE_{U\sim\mathrm{Breta}^\circ(F;\sigma\one;\boldsymbol\sigma,\boldsymbol\sigma)}\big[e^{2\ip cZ-\sum_jb_jO_j(Z)}\big]$ is finite exactly when $\ell^\ast_{\boldsymbol\sigma}(c,\mathbf b)>1$.
\item For a question inside its walls with $|\delta_j|\le\sigma_j$ the native gauge $z^\star$ minimizes $-2\ip{\bmu_\perp}z+\sum_j(\sigma_jE_j+\delta_jO_j)$ on $\one^\perp$, so $\big(z^\star;E_j(z^\star),O_j(z^\star)\big)_j$ is the point of the set of values $\{(z;E_j(z),O_j(z))_j:\ z\in\one^\perp\}$ of the statistics at which $-2\ip{\bmu_\perp}x+\sum_j(\sigma_jy_j+\delta_jw_j)$ is least; at the centre it is $\big(0;2\log K_j(\one),0\big)_j$.
\end{enumerate}
\end{proposition}

% Proved in the last part: prop:tails (new).
\begin{proposition}[The tails of the statistics are read at the walls]\label{prop:tails}
Let a question $q$ of a family have every proper wall positive, let $Z$ be the tilt of $U\sim\mathrm{Breta}^\circ(F;\mvec;\mathbf s,\mathbf t)$ at its data, and for a real balanced direction $v$ of pair coordinates $(c;\mathbf a,\mathbf b)$ (Definition~\ref{def:question}) let $Y=2\ip cZ-\sum_j\big(a_jE_j(Z)+b_jO_j(Z)\big)$. On the balance wall the walls are linear in the pair coordinates, $\kappa_S(q+\lambda v)=\kappa_S(q)+\lambda\kappa_S(v)$ with $\kappa_S(v)=\kappa^{\mathbf E}_S(\mathbf a)+\kappa^{\mathbf O}_S(c,\mathbf b)$ (Proposition~\ref{prop:reach}). Let
\[
\lambda^\ast=\min\Big\{\frac{\kappa_S(q)}{-\kappa_S(v)}:\ \kappa_S(v)<0\Big\},
\]
the value of $\lambda$ at which the ray $q+\lambda v$, $\lambda>0$, first meets a wall, with $\lambda^\ast=\infty$ if no $\kappa_S(v)$ is negative. For real $\lambda\ge0$, $\oE[e^{\lambda Y}]$ is finite exactly for $\lambda<\lambda^\ast$, where it is the ratio of the periods of the questions $q+\lambda v$ and $q$ (Corollary~\ref{cor:letterstransform}); if $\lambda^\ast<\infty$,
\[
\mathbb P(Y>x)=e^{-\lambda^\ast x+o(x)}\qquad(x\to\infty),
\]
and if $\lambda^\ast=\infty$, $Y$ is bounded above. Under the law of the centre of strengths $\boldsymbol\sigma$ an odd direction, of pair coordinates $(c;\mathbf 0,\mathbf b)$, has $\lambda^\ast=\ell^\ast_{\boldsymbol\sigma}(c,\mathbf b)$ for $Y$ and for $-Y$ alike.
\end{proposition}

The period of a question is a Laplace transform in the tilts, so the shape laws form an exponential family whose natural parameters are the offsets and the pair coordinates of the exponents and whose statistics are the tilt and the even and odd statistics $E_j$, $O_j$. On the sector of each order the exponent of a question is minus twice its wall at the tilt, the walls of the tails weighted by the differences of the tilt along the order, up to the bounded sector letters. So the walls alone fix how the period blows up as the strength goes to zero, the weights entering only at the next order and only through the greedy bases; at high strength, real or complex with its real part bounded below, the period is read at the native gauge of the ray. At a centre the exponent is minus the centred potential $\sum_j\sigma_jE_j$, which is convex, smallest at the table's gauge and within a constant of the widths of the Newton polytopes. A question can lean only as far as its strengths allow: the odd part of every wall is bounded by its even part, the slope of its odd part exceeds one, and the same slope bounds the exponential moments of the odd statistics at the centre; along every direction the tail of a statistic falls at the rate of the first wall the direction meets. Integrating over the simplex is integrating against the list $\binom{\cE}1$ at a full basis lean, so every Dirichlet expectation of letters is a ratio of periods (Proposition~\ref{prop:sectiontable}).

\chapter{The Breta law and the port on the balance wall}\label{ch:law}

\section{The laws}

% Reference 20: thm:normalizer; source line 3411
\begin{theorem}[The pseudointegral is a normalizer]\label{thm:normalizer}
Let $F$ be connected on $n\ge2$ marks with real data, $s_j,t_j\ge0$, not all zero, and $\mvec>0$, a native gauge $\alpha^\star$, target $\btau$ and form $\omega=W\,d\alpha_{\mvec}$.
\begin{enumerate}[label=(\roman*)]
\item $\pseudoint W\,d\alpha_{\mvec}=P_F(\btau;\mathbf s,\mathbf t)$ is the normalizer of $\mathrm{Breta}^\circ(F;\btau;\mathbf s,\mathbf t)$, and the pseudoexpectation $\mathsf E_{U\sim\mathbb P_{\mathrm c}}[(\omega_{\btau}/\omega_{\mathrm c})(U)]=\pseudoint W\,d\alpha_{\mvec}/P_{\mathrm c}$ is the normalizer of its density against the law of the centre (Proposition~\ref{prop:centrelaw}(i)). With $\phi_{\mvec'}$ the joint law of $\big(K_1(U),\dots,K_N(U),\allowbreak\Psi_1(U),\dots,\Psi_N(U)\big)$ for $U\sim\Dir_{\mvec'}$, both periods are stacked Gammas of masses over their total times Trailmoments (Proposition~\ref{prop:gammadraws}(i)):
\[
\mathsf E_{U\sim\mathbb P_{\mathrm c}}\big[(\omega_{\btau}/\omega_{\mathrm c})(U)\big]=\frac{\Gst{n}{1}{\btau}{\tau_{\cE}}\,\big(\sT_{(\mathbf s,\mathbf t);I_{2N}}[\mathrm{id}]\,\phi_{\btau}\big)\big((0,\infty)^{2N}\big)}{\Gst{n}{1}{\sigma\one}{n\sigma}\,\big(\sT_{(\boldsymbol\sigma,\boldsymbol\sigma);I_{2N}}[\mathrm{id}]\,\phi_{\sigma\one}\big)\big((0,\infty)^{2N}\big)} .
\]
\item In the coordinates $y=\log u$ modulo $\mathbb R\one$ the law $\mathrm{Breta}^\circ(F;\btau;\mathbf s,\mathbf t)$ has density proportional to $e^{\ip{\btau}y-\Phi(y)}$, whose unique mode is $\log\alpha^\star$.
\item (Cremona) If $U\sim\mathrm{Breta}^\circ(F;\mvec;\mathbf s,\mathbf t)$ on the balance wall, then $U^{-1}/|U^{-1}|\sim\mathrm{Breta}^\circ(F;2\sigma\one-\mvec;\mathbf t,\mathbf s)$.
\item (The even--odd pair) On the balance wall the inversion of (iii) carries $\mathrm{Breta}^\circ(F;\mvec;\mathbf s,\mathbf t)$ at the data $(\bmu_\perp;\boldsymbol\sigma,\boldsymbol\delta)$ to the law at $(-\bmu_\perp;\boldsymbol\sigma,-\boldsymbol\delta)$; at $(0;\boldsymbol\sigma,\mathbf 0)$ it fixes the law, and with $Z$ the signed coordinate of $U$ the odd statistics $(O_j(Z))_j$ have a law symmetric under $O\mapsto-O$. The joint Laplace transform of the even and odd statistics is Corollary~\ref{cor:letterstransform}.
\end{enumerate}
\end{theorem}

% Reference 20: prop:marklaw; source line 3431
\begin{proposition}[The mark relation of the Breta shape law]\label{prop:marklaw}
Let $F$ have real data, masses $\mvec$ on the balance wall and every proper wall positive, and in this statement write $\mathbf u(\alpha)=\sum_j\big(s_jR_j(\alpha)+t_jR^\ast_j(\alpha)\big)$ for the usage at $\alpha$ and $\oE$ for the expectation of $U\sim\mathrm{Breta}^\circ(F;\mvec;\mathbf s,\mathbf t)$. For $Y=\log U$ modulo $\mathbb R\one$, every mark $e$ and every observable $g$, a smooth function of degree $0$ in $\alpha$ such that $g$ and $\partial_{y_e}g$ grow at most exponentially at a rate below the walls,
\[
\oE\big[(m_e-\mathbf u_e(U))\,g(U)+\partial_{y_e}g(U)\big]=0 .
\]
At $g=1$ the masses are the mean usage, $\oE[\mathbf u(U)]=\mvec$; with a native gauge $\alpha^\star$ the usage equals the target $\btau=\mvec$ at the mode $\log\alpha^\star$ (Theorem~\ref{thm:normalizer}(ii)).
\end{proposition}

% Reference 20: prop:stick; source line 3442 (on the balance wall).
% Proved in the last part: prop:stick.
\begin{proposition}[Scale and shapes along a bunch]\label{prop:stick}
For a proper bunch $B$ with room $R$, masses $\mvec>0$ and every proper wall of positive real part,
\begin{multline*}
P_F(\mvec;\mathbf s,\mathbf t)=\Gst{|B|}{1}{\mvec|_B}{m_B}\Gst{|R|}{1}{\mvec|_R}{m_R}\\
\times\oE_{(U'_B,U'_R)\sim\Dir_{\mvec|_B}\otimes\Dir_{\mvec|_R}}\Big[\int_0^\infty x^{m_B-1}(1+x)^{-\kappa_{\cE}}\prod_jK_j(xU'_B,U'_R)^{-s_j}\Psi_j(xU'_B,U'_R)^{-t_j}\,dx\Big],
\end{multline*}
the inner integral being the scale integral of the family with the list $\binom{\cE}1$ appended at the exponents $(\kappa_{\cE},0)$, whose letter at $(xU'_B,U'_R)$ is $1+x$ (Proposition~\ref{prop:sectiontable}); on the balance wall it is $\operatorname{Sc}_B(U'_B,U'_R;\mvec;\mathbf s,\mathbf t)$ (Theorem~\ref{thm:scale}), and the display is Corollary~\ref{cor:fiber} at $h=|\alpha_R|$. For real data with every proper wall positive and $U\sim\mathrm{Breta}^\circ(F;\mvec;\mathbf s,\mathbf t)$, with the internal shapes $V_B=U_B/|U_B|$ and $V_R=U_R/|U_R|$ and the totals $|U_B|+|U_R|=1$ (Theorem~\ref{thm:aggregation}), the scale $\rho=|U_B|/|U_R|$ of $B$ given the internal shapes has density proportional to $\rho^{m_B-1}(1+\rho)^{-\kappa_{\cE}}\prod_jK_j(\rho V_B,V_R)^{-s_j}\Psi_j(\rho V_B,V_R)^{-t_j}$.
\end{proposition}

% Reference 20: thm:bretalawthm; source line 3453
\begin{theorem}[Structure of the Breta law]\label{thm:bretalawthm}
Let $X\sim\mathrm{Breta}(F;\mvec;\mathbf s,\mathbf t)$; in this statement $\oE$ is the expectation of $X$.
\begin{enumerate}[label=(\roman*)]
\item $|X|\sim\operatorname{Gamma}_{\kappa_{\cE}}$ and $X/|X|\sim\mathrm{Breta}^\circ(F;\mvec;\mathbf s,\mathbf t)$, independently.
\item The laws form an exponential family with respect to $e^{-|\alpha|}d\alpha/\alpha^{\one}$, with statistics $(\log X_e)_e$ and $(-\log K_j(X),-\log\Psi_j(X))_j$, natural parameters $(\mvec,\mathbf s,\mathbf t)$ and log-partition $\log J_F(\mvec;\mathbf s,\mathbf t)$: $\oE\log X_e=\partial_{m_e}\log J_F(\mvec;\mathbf s,\mathbf t)$, $\oE\log(K_j\Psi_j)(X)=-D_{(\mathbf 0;\mathbf e_j,\mathbf e_j)}\log J_F(\mvec;\mathbf s,\mathbf t)$, $\oE\log(K_j/\Psi_j)(X)=-D_{(\mathbf 0;\mathbf e_j,-\mathbf e_j)}\log J_F(\mvec;\mathbf s,\mathbf t)$, $D_v$ the derivative along the direction $v$ of the data, the masses fixed, and the covariances are second derivatives.
\item $\oE[X^{\mathbf k}]=J_F(\mvec+\mathbf k;\mathbf s,\mathbf t)/J_F(\mvec;\mathbf s,\mathbf t)$ and $\oE\big[\prod_jK_j(X)^{-a_j}\Psi_j(X)^{-b_j}\big]=J_F(\mvec;\mathbf s+\mathbf a,\mathbf t+\mathbf b)/J_F(\mvec;\mathbf s,\mathbf t)$.
\end{enumerate}
\end{theorem}

% Reference 20: thm:race; source line 3465
\begin{theorem}[The race representation]\label{thm:race}
Let $X\sim\mathrm{Breta}(F;\mvec;\mathbf s,\mathbf t)$. Given $X=\alpha$, run for each table $j$ a \emph{basis race}, the basis $T$ ringing with intensity $\theta_{j,T}\alpha^T$, and a \emph{cobasis race}, with intensity $\theta_{j,T}\alpha^{\cE\setminus T}$, all clocks independent; let $Y_j$ and $Y^\ast_j$ be the first ringing times and $T_j$, $T'_j$ the winners. In this statement $\oE$ and $\mathbb P$ are taken over the clocks of the races at $\alpha$.
\begin{enumerate}[label=(\roman*)]
\item In each race the time and the winner are independent: $K_j(\alpha)Y_j\sim\operatorname{Gamma}_1$ and $T_j\sim\mathbb P_{j,\alpha}$; $\Psi_j(\alpha)Y^\ast_j\sim\operatorname{Gamma}_1$ and $T'_j\sim\mathbb P^\ast_{j,\alpha}$.
\item $K_j(\alpha)^{-s_j}\Psi_j(\alpha)^{-t_j}=\oE\big[Y_j^{s_j}(Y^\ast_j)^{t_j}\big]\big/\big(\Gamma(1+s_j)\Gamma(1+t_j)\big)$.
\item With $\mathrm g_j=-\log(K_jY_j)$ and $\mathrm g'_j=-\log(\Psi_jY^\ast_j)$, independent, $-\log(Y_jY^\ast_j)=\log(K_j\Psi_j)+\mathrm g_j+\mathrm g'_j$ and $-\log(Y_j/Y^\ast_j)=\log(K_j/\Psi_j)+\mathrm g_j-\mathrm g'_j$, and, for real $a,b$ with $1+a\pm b>0$, $\oE\big[e^{-a(\mathrm g_j+\mathrm g'_j)-b(\mathrm g_j-\mathrm g'_j)}\big]=\Gamma(1+a+b)\,\Gamma(1+a-b)$.
\item $\mathbb P(Y_j<Y^\ast_j)=K_j/(K_j+\Psi_j)=1/(1+e^{-O_j(z)})$, $\alpha=e^{2z}$.
\item $D=\chi_{T_j}-\chi_{T'_j}$ has $\oE[e^{\ip wD}]=e^{E_j(z+w)-E_j(z)}$: the even statistic $E_j$ is the cumulant function of the difference of the two winners (by (i), Theorem~\ref{thm:innerMGF} for table $j$).
\end{enumerate}
\end{theorem}

The two exponents of a table are the two clocks of its races: its even letter reads the sum $\mathrm g_j+\mathrm g'_j$ of the two logarithmic clocks of (iii) and its odd letter their difference, so that the pair coordinates $(\sigma_j,\delta_j)$ are the exponents of $\mathrm g_j+\mathrm g'_j$ and $\mathrm g_j-\mathrm g'_j$.

% Reference 20: thm:conjugacy; source line 3479
\begin{theorem}[Conjugacy]\label{thm:conjugacy}
Let $X\sim\mathrm{Breta}(F;\mvec;\mathbf s,\mathbf t)$.
\begin{enumerate}[label=(\roman*)]
\item Given the winner $T$ of the $j$-th basis race, its time unobserved, $X\sim\mathrm{Breta}(F;\mvec+\one_T;\mathbf s+\mathbf e_j,\mathbf t)$; given the winner $T'$ of the $j$-th cobasis race, $X\sim\mathrm{Breta}(F;\mvec+\one_{\cE\setminus T'};\mathbf s,\mathbf t+\mathbf e_j)$. Both keep $\kappa_{\cE}$.
\item Under the joint law of $X$ and the races of Theorem~\ref{thm:race}, $\mathbb P(T_j=T)=\theta_{j,T}P_F(\mvec+\one_T;\mathbf s+\mathbf e_j,\mathbf t)/P_F(\mvec;\mathbf s,\mathbf t)$ and $\mathbb P(T'_j=T')=\theta_{j,T'}P_F(\mvec+\one_{\cE\setminus T'};\mathbf s,\mathbf t+\mathbf e_j)/P_F(\mvec;\mathbf s,\mathbf t)$, each summing to one.
\item Given the time $Y_j=\lambda$ as well, $X$ has the density of $\mathrm{Breta}(F;\mvec+\one_T;\mathbf s,\mathbf t)$ multiplied by $e^{-\lambda K_j(\alpha)}$ and renormalized: the kernel terms $\lambda\theta_{j,T}\alpha^T$ join those of $e^{-|\alpha|}$, and the law exists only on the orthant.
\end{enumerate}
\end{theorem}

% New: proved in the last part.
\begin{corollary}[Conjugacy for the continued period]\label{cor:conjugacyoff}
For every family and every table $j$ of it, with its bases $\cB_j$ and weights $\theta_j$,
\[
\sum_{T\in\cB_j}\theta_{j,T}\,P_F(\mvec+\one_T;\mathbf s+\mathbf e_j,\mathbf t)=P_F(\mvec;\mathbf s,\mathbf t)=\sum_{T'\in\cB_j}\theta_{j,T'}\,P_F(\mvec+\one_{\cE\setminus T'};\mathbf s,\mathbf t+\mathbf e_j)
\]
as identities of meromorphic functions of the masses and the exponents (Lemma~\ref{lem:periodmero}), and so on the balance wall: the identities of Theorem~\ref{thm:conjugacy}(ii) hold for the continued values of Theorem~\ref{thm:offdomain}.
\end{corollary}

% Reference 20: thm:winnersrank; source line 3491
\begin{theorem}[The winners of a tensorminental table]\label{thm:winnersrank}
If table $j$ is tensorminental of tensor rank $q$ with slices $A^{(k)}$, then for a bunch $B$ with room $R$ and at a gauge as in Theorem~\ref{thm:interference}, the pattern $T_j\cap B$ of the basis winner has the generating function $\sum_{k,l}w_{kl}\det\big(I+(D_x-I)\Pi^{(kl)}_{BB}\big)$ of that theorem, and the cobasis winner's is the same with the slices of the dual multivector. Conditioning on the winner's pattern leaves tables with multivectors of tensor rank at most $q$ on the room.
\end{theorem}

% Reference 20: prop:laplace; source line 3498 (data scaled from the origin).
% Proved in the last part: prop:laplace.
\begin{proposition}[The period along a ray of data]\label{prop:laplace}
Let $F$ be connected, let $(\btau;\mathbf s,\mathbf t)$ be real data on its balance wall with $\btau>0$, $s_j,t_j\ge0$ not all zero, and a native gauge $\alpha^\star$, let $(\mvec_0;\mathbf s_0,\mathbf t_0)$ be complex data on the balance wall, and in this statement write $W_0=\prod_jK_j^{-s_{0,j}}\Psi_j^{-t_{0,j}}$. As $\lambda\to+\infty$ the data $(\mvec_0+\lambda\btau;\mathbf s_0+\lambda\mathbf s,\mathbf t_0+\lambda\mathbf t)$ have every proper wall of positive real part, and, with $\Phi$ the potential at $(\mathbf s,\mathbf t)$ and a chart $g$,
\begin{multline*}
P_F(\mvec_0+\lambda\btau;\mathbf s_0+\lambda\mathbf s,\mathbf t_0+\lambda\mathbf t)\\
\sim(\alpha^\star)^{\mvec_0}W_0(\alpha^\star)\Big(\frac{2\pi}\lambda\Big)^{\frac{n-1}2}\det\big(\nabla^2\Phi(\log\alpha^\star)\big)_{(g)}^{-1/2}\,e^{-\lambda(\Phi(\log\alpha^\star)-\ip{\btau}{\log\alpha^\star})},
\end{multline*}
$(g)$ deleting the row and the column of $g$, where $\nabla^2\Phi$ is $\frac14$ times rung two of the still and $(\alpha^\star)^{\mvec_0}W_0(\alpha^\star)$ does not depend on the scale of $\alpha^\star$. At $(\mvec_0;\mathbf s_0,\mathbf t_0)=0$ the display is the leading term of Theorem~\ref{thm:severity}(ii) at real $\lambda$ for the question of the data $(\btau;\mathbf s,\mathbf t)$. For real $(\mvec_0;\mathbf s_0,\mathbf t_0)$, under $\mathrm{Breta}^\circ(F;\mvec_0+\lambda\btau;\mathbf s_0+\lambda\mathbf s,\mathbf t_0+\lambda\mathbf t)$, $\sqrt\lambda(Y-\log\alpha^\star)$ tends to the centred normal law on $\mathbb R^{\cE}/\mathbb R\one$ with covariance $(\nabla^2\Phi)^{-1}$, $Y=\log U$.
\end{proposition}

A base point added to the data moves the ray off the origin: along $(\mvec_0+\lambda\btau;\mathbf s_0+\lambda\mathbf s,\mathbf t_0+\lambda\mathbf t)$, with complex $(\mvec_0;\mathbf s_0,\mathbf t_0)$ on the balance wall, the leading term of Theorem~\ref{thm:severity}(ii) gains the factor $(\alpha^\star)^{\mvec_0}\prod_jK_j(\alpha^\star)^{-s_{0,j}}\Psi_j(\alpha^\star)^{-t_{0,j}}$ at the native gauge $\alpha^\star=e^{2z^\star}$.

% Reference 20: def:skellaw; source line 3509 (positive letter exponents).
% Proved in the last part: def:skellaw.
\begin{definition}[Skeleton law]\label{def:skellaw}
For coefficients $\xi_\ell\in\mathbb C\setminus(-\infty,0]$, $\ell\in J=\{0,\dots,L\}$, closed under conjugation, real nonzero letter exponents $w_\ell$, equal on conjugate coefficients, with total $|\mathbf w|$, and $0<\kappa<|\mathbf w|$, the \emph{Skeleton law} $\mathrm{Skel}_{\mathbf w,\kappa}(\boldsymbol\xi)$ is the law on $(0,\infty)$ with density
\[
y^{\kappa-1}\prod_\ell(y+\xi_\ell)^{-w_\ell}\Big/\Gst{2}{|J|}{\kappa,\ |\mathbf w|-\kappa}{\mathbf w}\,\sB^{\{J\}}_{\mathbf w,\,|\mathbf w|-\kappa;\,I_J}(\boldsymbol\xi),
\]
its normalizer being Proposition~\ref{prop:scalebreta}; the density is positive, and integrable exactly when $0<\kappa<|\mathbf w|$.
\end{definition}

% Reference 20: thm:skellaw; source line 3513 (on the balance wall; the mixture about the largest coefficient).
% Proved in the last part: thm:skellaw.
\begin{theorem}[The Skeleton law]\label{thm:skellaw}
\begin{enumerate}[label=(\roman*)]
\item Under $\mathrm{Breta}^\circ(F;\mvec;\mathbf s,\mathbf t)$, the scale of a proper bunch $B$ given the internal shapes of $B$ and of its room has the law $\mathrm{Skel}_{\mathbf w,\kappa_B}(\boldsymbol\xi)$: $\boldsymbol\xi$ the joint spectrum of $B$ at that point and $\mathbf w$ its letter exponents of Theorem~\ref{thm:scale}, and off the balance wall the coefficient $1$ besides, with the letter exponent $\kappa_{\cE}$, added to the letter exponent of $1$ when $1$ is a point of the joint spectrum; coefficients of letter exponent zero are discarded.
\item For $Y\sim\mathrm{Skel}_{\mathbf w,\kappa}(\boldsymbol\xi)$ and $-\kappa<\re k<|\mathbf w|-\kappa$,
\[
\oE_{Y\sim\mathrm{Skel}_{\mathbf w,\kappa}(\boldsymbol\xi)}[Y^k]=\Gst{2}{2}{\kappa+k,\ |\mathbf w|-\kappa-k}{\kappa,\ |\mathbf w|-\kappa}\,\frac{\sB^{\{J\}}_{\mathbf w,\,|\mathbf w|-\kappa-k;\,I_J}(\boldsymbol\xi)}{\sB^{\{J\}}_{\mathbf w,\,|\mathbf w|-\kappa;\,I_J}(\boldsymbol\xi)},
\]
the quotient of the values of Proposition~\ref{prop:scalebreta} at the indices $\kappa+k$ and $\kappa$, read as that quotient when a letter exponent is not positive; and $1/Y\sim\mathrm{Skel}_{\mathbf w,|\mathbf w|-\kappa}(\boldsymbol\xi^{-1})$.
\item With real coefficients, positive letter exponents and a point $r\ge\max_\ell\xi_\ell$ of the axis, $Y$ is the mixture over $\mathbf n\in\mathbb N^J$ of $rY_{\mathbf n}$, $Y_{\mathbf n}\sim\mathrm{Skel}_{|\mathbf w|+|\mathbf n|,\kappa}(1)$ the law with the one coefficient $1$, mixed in proportion to $\prod_{\ell\in J}\frac{(w_\ell)_{n_\ell}}{n_\ell!}g_\ell^{\,n_\ell}\Gst21{\kappa,\ |\mathbf w|+|\mathbf n|-\kappa}{|\mathbf w|+|\mathbf n|}$, $g_\ell=1-\xi_\ell/r$ the spectral gaps at $r$ of Theorem~\ref{thm:gaps}; at $r=\max_\ell\xi_\ell$ the largest coefficient drops out.
\item In the situation of (i) on the balance wall, with $s_j,t_j\ge0$: given $Y\le1$, $Y$ is the mixture of the first coordinates of the laws $\Dir_{(\kappa_B+k,|\mathbf n|+1)}$, and given $Y\ge1$, $1/Y$ is the mixture of the first coordinates of $\Dir_{(\kappa^{\mathrm{top}}_B+k,|\mathbf n|+1)}$, mixed in proportion to the terms of Theorem~\ref{thm:sectors}(ii); the two sectors carry probabilities proportional to the sector Trailmoments of that theorem at the indices $\kappa_B$ and $\kappa^{\mathrm{top}}_B$.
\end{enumerate}
\end{theorem}

\section{The port on the balance wall}

% Reference 20: def:wallcontour; source line 3528 (real data).
% Proved in the last part: def:wallcontour.
\begin{definition}[Wall contour; joint period]\label{def:wallcontour}
On the balance wall, with $\re s_j>0$ or $s_j=0$, and $\re t_j>0$ or $t_j=0$, for every $j$, a \emph{wall contour} at a bunch $B$ puts the shares of Theorem~\ref{thm:port} of the letters with nonzero exponent on real parts $c_{j,p},c'_{j,p}>0$ with $\sum_pc_{j,p}=\re s_j$, $\sum_pc'_{j,p}=\re t_j$ and real odd load equal to the real part of the offset on the bunch, $\sum_{j,p}\frac{c_{j,p}-c'_{j,p}}2\chi_p(e)=\re\mu_e$; equivalently the real part of every mass of the bunch is carried exactly, $\re\operatorname{load}_e=\re m_e$. The \emph{fibre integral at $B$} is the fibre integral (Definition~\ref{def:fibre}) with the pins $|\mathbf a_j|=s_j$, $|\mathbf b_j|=t_j$ and $\langle\boldsymbol\delta\rangle_e=\mu_e$, $e\in B$. For a family $G=((\cB_i,\theta_i);a_i,b_i)_i$ whose molecules are $R_1,\dots,R_q$, the \emph{joint period} $P^\otimes_G(\mvec;\mathbf a,\mathbf b)$ is the integral of $\alpha^{\mvec-\one}\prod_iK_i(\alpha)^{-a_i}\Psi_i(\alpha)^{-b_i}$ over $\alpha\in\prod_{k=1}^q\Delta_{R_k}$ against $\prod_kd\omega_{R_k}$, wherever it converges; when the weights factor across the molecules it is the product of the molecules' periods.
\end{definition}

% Reference 20: thm:wallsplit; source line 3532 (real data with $s_j,t_j\ge0$).
% Proved in the last part: thm:wallsplit.
\begin{theorem}[The port on the balance wall]\label{thm:wallsplit}
Let $F$ have data on the balance wall with $\re s_j>0$ or $s_j=0$, and $\re t_j>0$ or $t_j=0$, for every $j$, let $P_F(\mvec;\mathbf s,\mathbf t)$ converge, and let $B$ be a proper bunch with room $R$; put $z=\frac12\log\alpha$.
\begin{enumerate}[label=(\roman*)]
\item The \emph{bunch factor} $\Xi_B(\alpha_R;\bmu|_B;\mathbf s,\mathbf t)=\int_{\mathbb R^B_{>0}}\alpha_B^{\bmu|_B-\one}\prod_j\Lambda_j(z)^{-s_j}\Lambda_j(-z)^{-t_j}\,d\alpha_B$ converges for every $\alpha_R$, $P_F(\mvec;\mathbf s,\mathbf t)=\int_{\Delta_R}v^{\bmu|_R-\one}\,\Xi_B(v;\bmu|_B;\mathbf s,\mathbf t)\,d\omega_R(v)$, and on every wall contour
\begin{multline*}
\Xi_B(\alpha_R;\bmu|_B;\mathbf s,\mathbf t)=\int_{|\mathbf a_j|=s_j,\ |\mathbf b_j|=t_j,\ \langle\boldsymbol\delta\rangle=\bmu|_B}\ \prod_j\Gst{|\pat_j(B)|}1{\mathbf a_j}{s_j}\Gst{|\pat_j(B)|}1{\mathbf b_j}{t_j}\\
\times\prod_{j,p}\Lambda_{j,(p)}(z_R)^{-a_{j,p}}\Lambda_{j,(p)}(-z_R)^{-b_{j,p}}.
\end{multline*}
\item Let $R_1,\dots,R_q$ be the molecules of the children and $\kappa_{R_i}(\mathbf a,\mathbf b)$ the wall at $R_i$ of the children $F_B(\mathbf a,\mathbf b)$ with the masses $\mvec|_R$ (Definition~\ref{def:family}); as $R_i$ is a molecule, $\kappa_{R_i}(\mathbf a,\mathbf b)=m_{R_i}-\sum_{j,p}\big(a_{j,p}r_{j,(p)}(R_i)+b_{j,p}(|R_i|-r_{j,(p)}(R_i))\big)$, with $r_{j,(p)}$ the rank function of the child $(j,p)$. Let $\mathcal X$ be the fibre of the shares cut by the pins $|\mathbf a_j|=s_j$, $|\mathbf b_j|=t_j$, $\langle\boldsymbol\delta\rangle=\bmu|_B$ and $\kappa_{R_1}=\cdots=\kappa_{R_{q-1}}=0$, with $|B|+q-1$ pins beyond the share sums; $\kappa_{R_q}=0$ holds on it. On a wall contour with $\re\kappa_{R_i}=0$ for $i<q$ at which $P^\otimes_{F_B(\mathbf c,\mathbf c')}(\re\mvec|_R;\mathbf c,\mathbf c')$ converges,
\[
P_F(\mvec;\mathbf s,\mathbf t)=\int_{\mathcal X}\ \prod_j\Gst{|\pat_j(B)|}1{\mathbf a_j}{s_j}\Gst{|\pat_j(B)|}1{\mathbf b_j}{t_j}\,P^\otimes_{F_B(\mathbf a,\mathbf b)}(\mvec|_R;\mathbf a,\mathbf b).
\]
\item The native split of the real parts $(\re\mvec;\re\mathbf s,\re\mathbf t)$, which have a native gauge, is a wall contour whose real parts are those of the points of $\mathcal X$, and $P^\otimes$ converges there; for real data it is the native split, and it lies on the fibre.
\end{enumerate}
\end{theorem}

The bunch is integrated on its orthant, scale and shape at once: writing $\alpha_B=\rho u$ before opening the letters would split the $|B|$ pins into a pin on the scale and a stacked Gamma over the shape whose indices sum to zero, carrying $1/\Gamma(0)=0$, and neither has a value alone. For the same reason $B=\cE$ is excluded: the imaginary part of the odd load then lies in the hyperplane $\sum_{e\in\cE}\im\langle\boldsymbol\delta\rangle_e=\im\mu_{\cE}$, the scaling mode, and one mark stays in the room.

% Reference 20: lem:conchild; source line 3555
\begin{lemma}[Connected tables have connected children]\label{lem:conchild}
If some table with $s_j+t_j>0$ is connected, then at every bunch the children at a split with positive shares form a connected family.
\end{lemma}

At the masses $\btau+\eta\mvec$ the port on the orthant carries the stacked Gamma with the upper entries $\nu_e=\mu_e-\langle\boldsymbol\delta\rangle_e$, $e\in B$, and $\kappa_{R_i}$, and the lower entry $\eta m_{\cE}$; its $|B|+q$ upper entries sum to $\eta m_{\cE}$ and have real parts positive multiples of $\eta$ on the wall contour of the native split (Theorem~\ref{thm:wallsplit}(iii)); as $\eta\to0$ it tends to $(2\pi)^{|B|+q-1}$ times the delta of the pins. So the port on the balance wall is the limit of the Foreman procedure on the ray (Theorem~\ref{thm:foremanE}): the Gammas are the pins.

% Reference 20: prop:peelmol; source line 3566
\begin{proposition}[Peeling a molecule]\label{prop:peelmol}
In Theorem~\ref{thm:wallsplit} with a native gauge, let $R_1$ be a molecule of the children and $g\in R_1$. Each child letter is $\sum_X\alpha_{R_1}^XK_{[X]}(\alpha_{R\setminus R_1})$ over the bases $X$ of its restriction to $R_1$, with $K_{[X]}$ the letter on $R\setminus R_1$ with the bases $Y$ and the weights $\theta_{X\cup Y}$ (for $\Psi$, $\alpha_{R_1}^{R_1\setminus X}\Psi_{[X]}$). Opening these sums with one share per part, $P^\otimes$ is the fibre integral with the $|R_1|-1$ pins $\operatorname{load}_e=m_e$, $e\in R_1-g$, of the stacked Gammas of these shares against the joint period of the new children on the remaining molecules; the pin at $g$ holds by the balance of $R_1$, the native split lies on the fibre, and the joint period converges there. With weights multiplicative across $R_1$ every $K_{[X]}$ is a multiple of one letter and this is the factorization of Theorem~\ref{thm:wallsplit}(ii).
\end{proposition}

% Reference 20: cor:atomwall; source line 3573
\begin{corollary}[Atoms on the balance wall]\label{cor:atomwall}
In Theorem~\ref{thm:wallsplit}, fix $g$. On a wall contour for $\cE-g$, with one share per basis and letter,
\[
P_F(\mvec;\mathbf s,\mathbf t)=\int_{|\mathbf a_j|=s_j,\ |\mathbf b_j|=t_j,\ \langle\boldsymbol\delta\rangle=\bmu|_{\cE-g}}\ \prod_j\Gst{|\cB_j|}1{\mathbf a_j}{s_j}\Gst{|\cB_j|}1{\mathbf b_j}{t_j}\prod_{j,T}\theta_{j,T}^{-a_{j,T}-b_{j,T}},
\]
independent of $g$; the pin at $g$ holds by itself.
\end{corollary}

Running the ports on the balance wall down to rooms of one mark writes $P_F(\mvec;\mathbf s,\mathbf t)$ as an iterated fibre integral with $n-1$ pins in all; with one share per pattern and letter its dimension is $\sum(|\cB_j|-1)-(n-1)$, one term for each active letter, $j$ its table, independent of the tree when no letter is split across molecules.

% New: proved in the last part.
\begin{proposition}[The pins reach every integer point]\label{prop:ct-pins}
In Corollary~\ref{cor:atomwall} take one share for each basis and each active letter, the $K_j$ with $s_j\neq0$ and the $\Psi_j$ with $t_j\neq0$, and write the pins in the shares as the share sums $|\mathbf a_j|=s_j$, $|\mathbf b_j|=t_j$ of the active letters and the loads $\operatorname{load}_e=m_e$, $e\in\cE-g$ (Definition~\ref{def:shares}). Let every two marks be joined by a sequence of marks in which any two consecutive marks lie in a common circuit of a table with an active letter. Then the pins map the integer points of the share space onto the integer points of their target. So at integer exponents and masses the fibre has integer points, its integer points are $x_0+\Lambda$ for one of them $x_0$ and the lattice $\Lambda$ of the integer points of its directions, which has the rank of the dimension of the fibre, and a basis of $\Lambda$ gives integral coordinates on the fibre in which the integer points are exactly the points with integer coordinates.
\end{proposition}

% Reference 20: cor:closure; source line 3599
\begin{corollary}[Closure]\label{cor:closure}
If the weights of a form $W\,d\alpha_{\mvec}$ with real data, $s_j,t_j\ge0$ not all zero, factor across its molecules and every molecule has a native gauge, $\pseudoint W\,d\alpha_{\mvec}$ is a product over the molecules of iterated fibre integrals over affine fibres of share space, with integrands products of stacked Gammas of the shares and monomials in the weights, $n-1$ pins on a molecule of $n$ marks. The period of each molecule at its target is also the nonnegative series of Corollary~\ref{cor:cubeclose} and of Theorem~\ref{thm:peak}. Every quasi-integral is an iterated constant term of the continuation of the share integral of Theorem~\ref{thm:port} at its chart completion.
\end{corollary}

% Reference 20: cor:marginal; source line 3606
\begin{corollary}[The marginal law of the shape]\label{cor:marginal}
In Theorem~\ref{thm:wallsplit} let $U\sim\mathrm{Breta}^\circ(F;\mvec;\mathbf s,\mathbf t)$, with the internal shapes $V_B=U_B/|U_B|$ and $V_R=U_R/|U_R|$ of $B$ and $R$ (Proposition~\ref{prop:stick}).
\begin{enumerate}[label=(\roman*)]
\item $\oE_{U'\sim\Dir_{\mvec|_B}}\big[\operatorname{Sc}_B(U',v;\mvec;\mathbf s,\mathbf t)\big]=v^{-\sigma\one}\,\Xi_B(v;\bmu|_B;\mathbf s,\mathbf t)\big/\Gst{|B|}1{\mvec|_B}{m_B}$.
\item Given $V_R=v$, for $\mathbf k\in\mathbb C^B$ with $|\mathbf k|=0$, $\oE_{U\sim\mathrm{Breta}^\circ(F;\mvec;\mathbf s,\mathbf t)}[V_B^{\mathbf k}\mid V_R=v]=\Xi_B(v;\bmu|_B+\mathbf k;\mathbf s,\mathbf t)/\Xi_B(v;\bmu|_B;\mathbf s,\mathbf t)$, the bunch factor at shifted offsets, whose fibre is $\langle\boldsymbol\delta\rangle=\bmu|_B+\mathbf k$.
\item $V_R$ has density proportional to $v^{\bmu|_R-\one}\,\Xi_B(v;\bmu|_B;\mathbf s,\mathbf t)$, the fibre mixture of the children's shape laws in proportion to the stacked Gammas times $P_{F_B(\mathbf a,\mathbf b)}(\mvec|_R;\mathbf a,\mathbf b)$ when the children are connected (Lemma~\ref{lem:conchild}).
\item Given both shapes, the scale of $B$ has the Skeleton law of Theorem~\ref{thm:skellaw}(i).
\end{enumerate}
\end{corollary}

\section{Cycles}

% Reference 20: def:cycle; source line 3625
\begin{definition}[Cycle of a family]\label{def:cycle}
A \emph{cycle} of a family is a bunch $C$ that is a circuit of one of its tables. That table's cobasis polynomial has floor one on $C$, and on its bottom colevel the letter is linear on $C$; a letter linear on a bunch gives the bunch the \emph{kernel} $e^{-h}$, $h$ the letter divided by the sum of its coefficients (Theorem~\ref{thm:cyclekernel}(ii) below).
\end{definition}

% Reference 20: thm:cyclekernel; source line 3629
\begin{theorem}[A cycle carries its own kernel]\label{thm:cyclekernel}
Let $F$ have real data with $s_j,t_j\ge0$, $\kappa_{\cE}=0$, $P_F(\mvec;\mathbf s,\mathbf t)$ convergent, and $C\neq\cE$ a circuit of table $j_0$ with $t_{j_0}>0$ and room $R$.
\begin{enumerate}[label=(\roman*)]
\item The patterns of table $j_0$ on $C$ of colevel one are exactly $p_e=C-e$, $e\in C$, and the child at $p_e$ is the contraction $\cB_{j_0}/C$ on $R$ with the weights $\theta_{j_0,p_e\cup X}$. So $\Psi_{j_0}=\Psi^\flat+\Psi^\sharp$ with $\Psi^\flat=\sum_{e\in C}\alpha_e\Psi_{j_0,(p_e)}(\alpha_R)$ linear in $\alpha_C$. The circuits are the only bunches whose colevel-one patterns are one per mark.
\item With $G(\alpha_R)=\sum_{e\in C}\Psi_{j_0,(p_e)}(\alpha_R)$ and $h=\Psi^\flat/G$, if the family with $\Psi_{j_0}^{-t_{j_0}}$ replaced by $(\Psi^\flat)^{-c}(\Psi^\sharp)^{-(t_{j_0}-c)}$ converges, then on $\re\beta=c$
\begin{multline*}
P_F(\mvec;\mathbf s,\mathbf t)=\frac1{\Gamma(t_{j_0})}\int_{\beta+\beta'=t_{j_0}}\Gamma(\beta')\\
\times\int_{\mathbb R^{\cE}_{>0}}\alpha^{\mvec-\one}e^{-h(\alpha)}G(\alpha_R)^{-\beta}\Psi^\sharp(\alpha)^{-\beta'}K_{j_0}^{-s_{j_0}}\prod_{j\neq j_0}K_j^{-s_j}\Psi_j^{-t_j}\,d\alpha .
\end{multline*}
When the colevel-one children are multiples $\vartheta_e$ of one letter, as for unit weights and for determinantal tables, the kernel is $e^{-\ip\vartheta{\alpha_C}/|\vartheta|}$.
\item On a wall contour at $C$ the pins at the marks of $C$ hold exactly when $b_{j_0,p_e}=m_e-\lambda_e$, $\lambda_e$ the load on $e$ of all the other shares, and $\sum_{j,p}\big(a_{j,p}(|p|-d_j(C))+b_{j,p}(|C\setminus p|-d^\ast_j(C))\big)=\kappa_C$, the coefficients being the excesses of the level and the colevel of $p$ in table $j$ (Definition~\ref{def:levels}); the fibre integral at $C$ is over the other shares with the single pin that their total excess is the wall of the cycle, the Gammas $\Gamma(m_e-\lambda_e)$ of the cycle standing in the places of the $\Gamma(b_{j_0,p_e})$.
\end{enumerate}
Dually, a cocircuit has floor one in the basis polynomial, and its bottom level gives its kernel.
\end{theorem}

% Reference 20: cor:cycle; source line 3648
\begin{corollary}[A table that is one circuit]\label{cor:cycle}
Let $F=((\cB,\theta);s,t)$ with $s,t>0$, where the bases of $\cB$ are the sets $\cE-e$, $e\in\cE$, with the weights $\theta_e$; let the masses lie on the balance wall, $m_{\cE}=(n-1)s+t$, with $m_S>s(|S|-1)$ for every nonempty proper $S$, and put $\nu_e=m_e-s$. Then
\begin{align*}
P_F(\mvec;s,t)&=\frac{\prod_e\theta_e^{-\nu_e}}{\Gamma(s)\Gamma(t)}\int_0^\infty y^{s-1}\prod_{e\in\cE}\sG_{\nu_e;(1,-1)}\big(1,\theta_e^2y\big)\,dy\\
&=\frac{\prod_e\theta_e^{-\nu_e}}{\Gamma(s)\Gamma(t)}\int_{\sum_ea_e=s}\ \prod_{e\in\cE}\Gamma(a_e)\,\Gamma(a_e+\nu_e)\,\theta_e^{-2a_e}:
\end{align*}
the cobasis polynomial $\Psi(\alpha)=\ip\theta\alpha$ is linear on the whole set, which so carries its own kernel (Definition~\ref{def:cycle}), and the fibre of the atom formula (Corollary~\ref{cor:atomwall}), of dimension $n-1$, closes to one scale.
\end{corollary}

\chapter{Bunches of width one: pooling and cleaning}\label{ch:blocks}

A bunch whose bases meet it on two levels only is seen by its room as one mark: the room reads the bunch through two ratios of polynomials on its marks, one for each Foreman polynomial, and the bunch is a table on its own marks and one more. A bunch of width one is read as a mark of its room, its shape integrated while its scale is kept, its walls glued, and its scale closed at a wall, at the gauge and on the totals of the bunch and the room. Cleaning, the repeated pooling of such bunches, is a move between period data and has one result.

\section{Bunches of width one}

% Reference 20: def:width; source line 3734
\begin{definition}[Parallel and series bunches]\label{def:width}
A \emph{parallel bunch} of a table (a parallel class) is a set of at least two marks, none a loop, no two of which lie in a common basis; a \emph{series bunch} (a series class) is a set of at least two marks, none a coloop, every basis containing all but at most one of them. A bunch $B$ of width $c(B)=r(B)-d(B)$ (Definition~\ref{def:skeleton}, Theorem~\ref{thm:window}) has its patterns on the $c(B)+1$ levels $d(B),\dots,r(B)$ (Theorem~\ref{thm:taylorcube}), and its room has the same width.
\end{definition}

% Reference 20: def:readings; source line 3772 (the factorization fixed tacitly).
\begin{definition}[Hinge values; their ratio]\label{def:readings}
Let a bunch $B$ of a table $(\cB,\theta)$ have width one, the weights factoring across $B$ and its room $R$, $\theta_T=\theta'_{T\cap B}\,\theta''_{T\cap R}$, so that $\cB|B$ and $\cB/R$ carry $\theta'$ and $\cB|R$ and $\cB/B$ carry $\theta''$ (Definition~\ref{def:minors}). The \emph{hinge value of $B$ by $K$} and its \emph{hinge value by $\Psi$} are
\[
\gamma^K_B=\frac{K_{\cB|B}}{K_{\cB/R}},\qquad \gamma^\Psi_B=\frac{\Psi_{\cB/R}}{\Psi_{\cB|B}},
\]
functions of $\alpha_B$ homogeneous of degree one, the values at which $K$ and $\Psi$ see the bunch replaced by one mark, and their \emph{ratio}
\[
\varrho_B=\frac{\gamma^\Psi_B}{\gamma^K_B}=\frac{\mathbf E_{\cB/R}}{\mathbf E_{\cB|B}},
\]
the quotient of the even letters $\mathbf E=K\Psi$ (Definition~\ref{def:potential}) of the two tables on $B$, is homogeneous of degree zero, a function of the shape $u=\alpha_B/|\alpha_B|$ alone. The room has width one as well, and its hinge values are those of $R$ as a bunch: $\gamma^K_R=K_{\cB|R}/K_{\cB/B}$, $\gamma^\Psi_R=\Psi_{\cB/B}/\Psi_{\cB|R}$ and $\varrho_R=\mathbf E_{\cB/B}/\mathbf E_{\cB|R}$. Statements about hinge values and their ratios are read at one fixed factorization. For table $j$ of a family they are $\gamma^K_{j,B}$, $\gamma^\Psi_{j,B}$ and $\varrho_{j,B}$, each table with its own factorization.
\end{definition}

% Reference 20: prop:classreadings; source line 3780 (lists).
% Proved in the last part: prop:classreadings.
\begin{proposition}[Hinge values of parallel and series bunches]\label{prop:classreadings}
A parallel or series bunch has width at most one, and width zero exactly when it is a union of molecules. Let a parallel bunch $B$ of width one have its weights factoring across $B$ and $R$ (Definition~\ref{def:readings}); then $\cB/R$ has the single basis $\varnothing$, with the weight $\theta'_\varnothing$, and $\cB|B$ the singletons, with the weights $\theta'_{\{e\}}$, and in this statement write $w_e=\theta'_{\{e\}}/\theta'_\varnothing$. So $K_{\cB/R}=\theta'_\varnothing$, $K_{\cB|B}=\theta'_\varnothing\sum_{e\in B}w_e\alpha_e$, $\Psi_{\cB|B}=\theta'_\varnothing\,e_{|B|}(\alpha_B)\sum_{e\in B}w_e\alpha_e^{-1}$ and $\Psi_{\cB/R}=\theta'_\varnothing\,e_{|B|}(\alpha_B)$, and
\[
\gamma^K_B=\sum_{e\in B}w_e\alpha_e,\qquad \gamma^\Psi_B=\Big(\sum_{e\in B}w_e\alpha_e^{-1}\Big)^{-1}:
\]
$K$ adds the marks of the class with the factors $w_e$ and $\Psi$ combines them as reciprocals with the same factors. A series bunch $B$ of width one, its weights factoring, swaps the two roles: $\cB|B$ has the single basis $B$, with the weight $\theta'_B$, and $\cB/R$ the sets $B-e$, with the weights $\theta'_{B-e}$; with $w_e=\theta'_{B-e}/\theta'_B$, $K_{\cB|B}=\theta'_B\,e_{|B|}(\alpha_B)$, $K_{\cB/R}=\theta'_B\,e_{|B|}(\alpha_B)\sum_ew_e\alpha_e^{-1}$, $\Psi_{\cB|B}=\theta'_B$ and $\Psi_{\cB/R}=\theta'_B\sum_ew_e\alpha_e$, so $\gamma^\Psi_B=\sum_ew_e\alpha_e$ and $\gamma^K_B=\big(\sum_ew_e\alpha_e^{-1}\big)^{-1}$. In a list, the case used for cleaning (Definition~\ref{def:cleaning} below), with the factorization $\theta'=\theta''=1$, every $w_e$ is $1$, and a pair with the shape $(x,1-x)$ has $\varrho_B=x(1-x)$ when it is parallel and $\varrho_B=1/(x(1-x))$ when it is in series.
\end{proposition}

% Reference 20: thm:widthone; source line 3738 (the converse at two homogeneous parts).
% Proved in the last part: thm:widthone.
\begin{theorem}[A bunch of width one is one mark to its room]\label{thm:widthone}
Let a bunch $B$ of a table $(\cB,\theta)$ have width one, the weights factoring across $B$ and $R$, $\theta_T=\theta'_{T\cap B}\,\theta''_{T\cap R}$, and let $g$ be a new mark. The \emph{cell of the bunch} $\cB^+_B$ is the table on $B+g$ whose bases avoiding $g$ are those of $\cB|B$ and whose bases through $g$ are the $X+g$, $X\in\cB/R$, with the weights $\theta'$; the \emph{cell of the room} $\cB^+_R$ is the table on $R+g$ whose bases avoiding $g$ are those of $\cB|R$ and whose bases through $g$ are the $Y+g$, $Y\in\cB/B$, with the weights $\theta''$. Both are tables, lists when $(\cB,\theta)$ is one, in which $g$, the \emph{hinge}, is neither a loop nor a coloop, and the bases of $\cB$ are theirs glued along $g$ with $g$ on exactly one side:
\[
\cB=\{X\sqcup Y:\ X\in\cB|B,\ Y\in\cB/B\}\ \sqcup\ \{X\sqcup Y:\ X\in\cB/R,\ Y\in\cB|R\}.
\]
With the first argument of the polynomials of a cell the variable of its hinge and the hinge values of Definition~\ref{def:readings},
\[
K(\alpha)=K_{\cB/R}(\alpha_B)\,K_{\cB^+_R}\big(\gamma^K_B(\alpha_B),\alpha_R\big),\qquad \Psi(\alpha)=\Psi_{\cB|B}(\alpha_B)\,\Psi_{\cB^+_R}\big(\gamma^\Psi_B(\alpha_B),\alpha_R\big).
\]
Conversely, if $K(\alpha)=\sum_{l=0}^{k}A_l(\alpha_B)\,C_l(\alpha_R)$ with every $C_l$ homogeneous, then $c(B)\le k$; in particular $K(\alpha)=A(\alpha_B)\big(C_0(\alpha_R)+G(\alpha_B)\,C_1(\alpha_R)\big)$ with $C_0$ and $C_1$ homogeneous gives $c(B)\le1$.
\end{theorem}

The deletions and the contractions of the hinge are the four minors $\cB^+_B\setminus g=\cB|B$ and $\cB^+_B/g=\cB/R$, $\cB^+_R\setminus g=\cB|R$ and $\cB^+_R/g=\cB/B$ of the table. So the hinge values of $B$ are those of $B$ in its own cell, and across a bunch of width one the bunch and the room see each other as one mark: the room reads the hinge at $\alpha_g=\gamma^K_B(\alpha_B)$ in $K$ and at $\alpha_g=\gamma^\Psi_B(\alpha_B)$ in $\Psi$.

% Reference 20: cor:widthonestage; source line 3762
\begin{corollary}[The scale of a bunch of width one is that of one mark]\label{cor:widthonestage}
Along the scale of a bunch $B$ of width one of a table whose weights factor across $B$ and $R$, each Foreman polynomial has one Skeleton value, the room's hinge value over the bunch's:
\begin{gather*}
\operatorname{spec}^K_B(\alpha)=\Big\{\frac{\gamma^K_R(\alpha_R)}{\gamma^K_B(\alpha_B)}\Big\}=\Big\{\frac{K_{\cB/R}(\alpha_B)\,K_{\cB|R}(\alpha_R)}{K_{\cB|B}(\alpha_B)\,K_{\cB/B}(\alpha_R)}\Big\},\\
\operatorname{spec}^\Psi_B(\alpha)=\Big\{\frac{\gamma^\Psi_R(\alpha_R)}{\gamma^\Psi_B(\alpha_B)}\Big\}=\Big\{\frac{\Psi_{\cB|B}(\alpha_B)\,\Psi_{\cB/B}(\alpha_R)}{\Psi_{\cB/R}(\alpha_B)\,\Psi_{\cB|R}(\alpha_R)}\Big\},
\end{gather*}
the Skeleton value of the hinge as a one-mark bunch of the cell of the room (Corollary~\ref{cor:stage}), read at $\alpha_g=\gamma^K_B(\alpha_B)$, respectively at $\alpha_g=\gamma^\Psi_B(\alpha_B)$. The value of $K$ over that of $\Psi$ is $\varrho_B(\alpha_B)/\varrho_R(\alpha_R)$, the even letters of the minors of the level $d(B)$ over those of the level $r(B)$, and their product is the same quotient of the odd letters.
\end{corollary}

% New: proved in the last part.
\begin{lemma}[One hinge value of any degree is width at most one]\label{lem:mot-onehinge}
Let $B$ be a bunch of a table $(\cB,\theta)$ with the room $R$, and suppose that on the open orthant
\[
K(\alpha)=A(\alpha_B)\,P\big(\alpha_R,\gamma(\alpha_B)\big),
\]
with $A$ and $\gamma$ positive functions of $\alpha_B$, homogeneous, $\gamma$ of degree one, and $P$ a polynomial in $\alpha_R$ and one more variable with real coefficients, of any degree in that variable. Then $c(B)\le1$. Conversely, when the weights factor across $B$ and $R$ and $c(B)\le1$, $K$ has this form: at width zero with $P$ free of its last variable, and at width one with $A=K_{\cB/R}$, $\gamma=\gamma^K_B$ and $P(\alpha_R,y)=K_{\cB^+_R}(y,\alpha_R)$ (Theorem~\ref{thm:widthone}). The same holds for $\Psi$, the basis polynomial of the dual table, in which $B$ has the same width.
\end{lemma}

% New: proved in the last part.
\begin{proposition}[The freedom of the factorization]\label{prop:hingefactor}
Let a bunch $B$ of a table $(\cB,\theta)$ have width one, the weights factoring across $B$ and $R$, $\theta_T=\theta'_{T\cap B}\,\theta''_{T\cap R}$. The factorization is fixed up to one positive constant on each of the two levels: multiplying $\theta'$ by $\lambda_d$ on the patterns of the level $d(B)$ and by $\lambda_r$ on those of the level $r(B)$, and dividing $\theta''$ by $\lambda_d$ and by $\lambda_r$ on the room sides of the bases of these levels, multiplies the hinge values $\gamma^K_B$ and $\gamma^K_R$ by $c=\lambda_r/\lambda_d$, $\gamma^\Psi_B$ and $\gamma^\Psi_R$ by $c^{-1}$, and the ratios $\varrho_B$ and $\varrho_R$ by $c^{-2}$ (Definition~\ref{def:readings}), and leaves $\gamma^K_R/\gamma^K_B$, $\gamma^\Psi_R/\gamma^\Psi_B$ and $\varrho_B/\varrho_R$ unchanged. In a family each table has its own two constants.
\end{proposition}

\section{Pooling}

% Reference 20: thm:cleaning; source line 3793
\begin{theorem}[Pooling integrates the shape and keeps the scale]\label{thm:pooling}
Let $F=((\cB_j,\theta_j);s_j,t_j)_{j=1}^N$ be a family and $B$ a bunch of width one in every table, the weights of every table factoring across $B$ and $R$, and read $\alpha_B=\rho u$, $\rho>0$, $u\in\Delta_B$ (Lemma~\ref{lem:scaleshape}). For $u\in\Delta_B$ the \emph{pooled family} $F_{B\to g}(u)$ is the family on $R+g$ with, for every $j$, two tables with the bases of the cell of the room $\cB^+_{j,R}$ (Theorem~\ref{thm:widthone}): the first with its weights multiplied by $\gamma^K_{j,B}(u)$ on its bases through $g$ and the exponents $(s_j,0)$, the second with its weights multiplied by $\gamma^\Psi_{j,B}(u)$ on its bases avoiding $g$ and the exponents $(0,t_j)$. Its letters are $K_{\cB^+_{j,R}}\big(\gamma^K_{j,B}(u)\alpha_g,\alpha_R\big)^{-s_j}$ and $\Psi_{\cB^+_{j,R}}\big(\gamma^\Psi_{j,B}(u)\alpha_g,\alpha_R\big)^{-t_j}$; listing the $N$ first tables and then the $N$ second ones, its exponents are $(\mathbf s,\mathbf 0)$ and $(\mathbf 0,\mathbf t)$. To \emph{pool} $B$ is to pass from $F$ to its pooled families $F_{B\to g}(u)$, keeping the letters $K_{\cB_j/R}$ and $\Psi_{\cB_j|B}$ of the shape of $B$, its cell kernel (Chapter~\ref{ch:construct}), and the hinge values of every table. Give the hinge the mass $\kappa_B$, the wall of $B$. Then $F_{B\to g}(u)$ has the imbalance of $F$, and
\begin{multline*}
J_F(\mvec;\mathbf s,\mathbf t)=\int_{\Delta_B}u^{\mvec|_B-\one}\prod_jK_{\cB_j/R}(u)^{-s_j}\Psi_{\cB_j|B}(u)^{-t_j}\\
\times J_{F_{B\to g}(u)}\big((\kappa_B,\mvec|_R);(\mathbf s,\mathbf 0),(\mathbf 0,\mathbf t)\big)\,d\omega_B(u),
\end{multline*}
both sides converging together; the same holds with $P$ for $J$, on the balance wall as well.
\end{theorem}

The port (Theorem~\ref{thm:port}) integrates the scale and the shape of a bunch together; pooling integrates the shape of a bunch of width one and keeps its scale, which becomes the hinge of the room. For a list the pooled families again have weights that factor across every bunch of $R$ and across every bunch containing $g$, so poolings can be repeated. Proposition~\ref{prop:stick} is the same split with the scale integral left whole. Read backwards, pooling is gluing: the table is its cell of the bunch substituted into its cell of the room at the hinge (Proposition~\ref{prop:substitution}), and every cell of a substitution is pooled at once by Theorem~\ref{thm:corereads}.

\section{Walls glue}

% Reference 20: thm:wallsglue; source line 3825 (one table; in (iii) at the exponents (0,t)).
% Proved in the last part: thm:wallsglue.
\begin{theorem}[Walls glue]\label{thm:wallsglue}
Let a bunch $B$ of a table $(\cB,\theta)$ have width one, the weights factoring across $B$ and $R$, with the cells $\cB^+_B$ and $\cB^+_R$ of Theorem~\ref{thm:widthone}. For $S_B\subseteq B$ and $S_R\subseteq R$ the counts of the table glue from those of the cells, the hinge counted on exactly one side, as in the gluing of the bases:
\begin{align*}
d(S_B\sqcup S_R)&=\min\big\{d_{\cB^+_B}(S_B+g)+d_{\cB^+_R}(S_R),\ d_{\cB^+_B}(S_B)+d_{\cB^+_R}(S_R+g)\big\},\\
d^\ast(S_B\sqcup S_R)&=\min\big\{d^\ast_{\cB^+_B}(S_B+g)+d^\ast_{\cB^+_R}(S_R),\ d^\ast_{\cB^+_B}(S_B)+d^\ast_{\cB^+_R}(S_R+g)\big\}.
\end{align*}
For a family $F=((\cB_j,\theta_j);s_j,t_j)_{j=1}^N$ with masses $\mvec$, $B$ having width one in every table and the weights of every table factoring across $B$ and $R$:
\begin{enumerate}[label=(\roman*)]
\item the walls of the pooled family $F_{B\to g}(u)$ at the masses $(\kappa_B,\mvec|_R)$ are walls of $F$: its wall at $S_R$ is $\kappa_{S_R}$ and its wall at $S_R+g$ is $\kappa_{B\sqcup S_R}$;
\item the density $u^{\mvec|_B-\one}\prod_jK_{\cB_j/R}(u)^{-s_j}\Psi_{\cB_j|B}(u)^{-t_j}$ of the cell kernel of $B$ has at the face $u_{S_B}=0$ of $\Delta_B$ the wall of $F$, $m_{S_B}-\sum_j\big(s_j\,d_{\cB_j/R}(S_B)+t_j\,d^\ast_{\cB_j|B}(S_B)\big)=\kappa_{S_B}$;
\item on the balance wall, with the cell families $G_B=((\cB^+_{j,B},\theta'_j);s_j,t_j)_j$ at the masses $(\mvec|_B,2\sigma-\kappa_B)$ and $G_R=((\cB^+_{j,R},\theta''_j);s_j,t_j)_j$ at the masses $(\kappa_B,\mvec|_R)$, $\sigma$ the total strength, whose walls are written $\kappa^{G_B}_S$ and $\kappa^{G_R}_S$, both cell families lie on their balance walls and every wall of a cell family is a wall of $F$,
\[
\kappa^{G_B}_{S_B}=\kappa_{S_B},\qquad \kappa^{G_B}_{S_B+g}=\kappa_{S_B\sqcup R},\qquad \kappa^{G_R}_{S_R}=\kappa_{S_R},\qquad \kappa^{G_R}_{S_R+g}=\kappa_{B\sqcup S_R};
\]
if moreover $s_j,t_j\ge0$ and every active letter reads $B$ alike, its hinge value a constant multiple of one function $\gamma$ on $\Delta_B$, as when one letter is active, every wall of $F$ is glued from walls of the cell families,
\[
\kappa_{S_B\sqcup S_R}=\min\big\{\kappa^{G_B}_{S_B}+\kappa^{G_R}_{S_R},\ \kappa^{G_B}_{S_B+g}+\kappa^{G_R}_{S_R+g}\big\},
\]
and $F$ has every proper wall positive exactly when both cell families do.
\end{enumerate}
\end{theorem}

\section{Full leans}

% Reference 20: thm:cornerfactor; source line 3871 (one table at a full lean).
% Proved in the last part: thm:cornerfactor.
\begin{theorem}[Pooling factors at one hinge value]\label{thm:cornerfactor}
Let $F=((\cB_j,\theta_j);s_j,t_j)_{j=1}^N$ have real exponents and be on the balance wall with every proper wall positive, and let a bunch $B$ have width one in every table, the weights of every table factoring across $B$ and $R$. Suppose that every active letter reads $B$ alike, its hinge value a constant multiple of one function $\gamma$ on $\Delta_B$, the factorizations being chosen so that every constant is $1$: every active letter has the one hinge value $\gamma$. (A change of the factorization of table $j$ multiplies $\gamma^K_{j,B}$ by a positive constant and $\gamma^\Psi_{j,B}$ by its inverse (Proposition~\ref{prop:hingefactor}). So a table with one active letter can always be so normalized, and when every table is at a full lean the hypothesis is a condition across the tables alone; one table with both letters active can be so normalized exactly when $\varrho_B$ is constant.) Then $0<\kappa_B<2\sigma$, $\sigma$ the total strength (Definition~\ref{def:family}), and, with the cell families $G_B=((\cB^+_{j,B},\theta'_j);s_j,t_j)_j$ and $G_R=((\cB^+_{j,R},\theta''_j);s_j,t_j)_j$,
\[
\Gst21{\kappa_B,\ 2\sigma-\kappa_B}{2\sigma}\,P_F(\mvec;\mathbf s,\mathbf t)=P_{G_B}\big((\mvec|_B,2\sigma-\kappa_B);\mathbf s,\mathbf t\big)\,P_{G_R}\big((\kappa_B,\mvec|_R);\mathbf s,\mathbf t\big),
\]
both cell families on their balance walls with every proper wall positive. One table at the exponents $(0,t)$ or $(s,0)$ has one active letter, and then $2\sigma$ is $t$, respectively $s$.
\end{theorem}

When every active letter has one hinge value, as for one table at a full lean, where one letter is active, the hinge carries the mass $\kappa_B$ in the room and $2\sigma-\kappa_B$ in the bunch, and the period of the family is the product of the periods of its two cell families divided by the Beta corner of the hinge.

% Reference 20: prop:cleanrealization; source line 3898 (one list at the exponents (0,t)).
% Proved in the last part: prop:cleanrealization.
\begin{proposition}[Pooling a realization]\label{prop:cleanrealization}
Let $F=((\cB_j,\theta_j);s_j,t_j)_{j=1}^N$ be a family, $B$ a bunch of width one in every table, the weights of every table factoring across $B$ and $R$, and let $\Upsilon=((s_j,t_j)_j;\alpha)$ be a realization of $F$ with masses $\btau$ (Definition~\ref{def:realization}); write $\alpha_B=\rho u$, $\rho>0$, $u\in\Delta_B$ (Lemma~\ref{lem:scaleshape}), and read the two tables of index $j$ of the pooled family $F_{B\to g}(u)$ of Theorem~\ref{thm:pooling} as one table with two systems of weights at the exponents $(s_j,t_j)$ (Definition~\ref{def:family}). Then $((s_j,t_j)_j;(\rho,\alpha_R))$ is a realization of $F_{B\to g}(u)$, with the masses $\big(\tau_B-\sum_j(s_jd_j(B)+t_jd^\ast_j(B)),\ \btau|_R\big)$, the first at the hinge: pooling sends a realization to a realization, and the mass of the hinge is the wall $\kappa_B$ at the masses $\btau$. With one table at the exponents $(0,t)$ the letter $\Psi$ of $F_{B\to g}(u)$ at $(\rho,\alpha_R)$ is that of the cell of the room $(\cB^+_R,\theta'')$ at $(\gamma^\Psi_B(\alpha_B),\alpha_R)$, and $((0,t);(\gamma^\Psi_B(\alpha_B),\alpha_R))$ is a realization of $(\cB^+_R,\theta'')$ with the masses $(\tau_B-t\,d^\ast(B),\ \btau|_R)$.
\end{proposition}

% Reference 20: prop:readingsagree; source line 3905 (one table with $\varrho_B=1$).
% Proved in the last part: prop:readingsagree.
\begin{proposition}[Hinge values that agree]\label{prop:readingsagree}
Let $F=((\cB_j,\theta_j);s_j,t_j)_{j=1}^N$ have real exponents and be on the balance wall with every proper wall positive, and let a bunch $B$ have width one in every table, the weights of every table factoring across $B$ and $R$. Suppose that every active letter reads $B$ alike, its hinge value a constant multiple of one function $\gamma$ on $\Delta_B$: $\gamma^K_{j,B}=c_j\gamma$ whenever $s_j\neq0$ and $\gamma^\Psi_{j,B}=c^\ast_j\gamma$ whenever $t_j\neq0$, with constants $c_j,c^\ast_j>0$; put $c_j=1$ when $s_j=0$ and $c^\ast_j=1$ when $t_j=0$. Let $G^{\mathbf c}_R$ be the pooled family of Theorem~\ref{thm:pooling} with these constants in place of the hinge values: for every $j$ the table $\cB^+_{j,R}$ with its weights multiplied by $c_j$ on its bases through $g$, at the exponents $(s_j,0)$, and with its weights multiplied by $c^\ast_j$ on its bases avoiding $g$, at the exponents $(0,t_j)$. Then
\begin{multline*}
P_F(\mvec;\mathbf s,\mathbf t)=P_{G^{\mathbf c}_R}\big((\kappa_B,\mvec|_R);(\mathbf s,\mathbf 0),(\mathbf 0,\mathbf t)\big)\\
\times\int_{\Delta_B}u^{\mvec|_B-\one}\prod_jK_{\cB_j/R}(u)^{-s_j}\Psi_{\cB_j|B}(u)^{-t_j}\,\gamma(u)^{-\kappa_B}\,d\omega_B(u):
\end{multline*}
the shape and the room separate at every pair of exponents. For one table with $\varrho_B$ constant, a factorization with $\varrho_B=1$ (Proposition~\ref{prop:hingefactor}) gives $c=c^\ast=1$ and $\gamma=\gamma^K_B$, and $G^{\mathbf c}_R$ is then $((\cB^+_R,\theta'');s,t)$.
\end{proposition}

\section{Partial leans}

% Reference 20: thm:twosides; source line 3947 ($s_j,t_j>0$).
% Proved in the last part: thm:twosides.
\begin{theorem}[Gluing across a bunch of width one]\label{thm:twosides}
Let $F=((\cB_j,\theta_j);s_j,t_j)_{j=1}^N$ have real data with $s_j,t_j\ge0$, $s_j+t_j>0$, and masses $\mvec$ on the balance wall with every proper wall positive, and let a bunch $B$ have width one in every table, the weights factoring across $B$ and $R$. In table $j$ call $d_j(B)$ the \emph{level $d$} and $r_j(B)=d_j(B)+1$ the \emph{level $r$}. \emph{Level shares} $\mathbf a=(a^d_j,a^r_j)_j$ and $\mathbf b=(b^d_j,b^r_j)_j$ split the exponents over the two levels, $a^d_j+a^r_j=s_j$ and $b^d_j+b^r_j=t_j$, a letter of exponent zero being opened by nothing, its two level shares $0$; they are the shares of Definition~\ref{def:shares} summed over the patterns of each level, and their pair shares (Definition~\ref{def:shares}) are the \emph{even share} $\sigma^\ell_j=\frac12(a^\ell_j+b^\ell_j)$ and the \emph{odd share} $\delta^\ell_j=\frac12(a^\ell_j-b^\ell_j)$, $\ell\in\{d,r\}$. The \emph{bunch family} and the \emph{room family} are
\begin{gather*}
F^{\mathrm{bunch}}(\mathbf a,\mathbf b)=\big((\cB_j/R;\,a^d_j,b^d_j),\ (\cB_j|B;\,a^r_j,b^r_j)\big)_j\quad\text{on }B,\\
F^{\mathrm{room}}(\mathbf a,\mathbf b)=\big((\cB_j|R;\,a^d_j,b^d_j),\ (\cB_j/B;\,a^r_j,b^r_j)\big)_j\quad\text{on }R:
\end{gather*}
each level carries one table of the bunch and one of the room, at the same shares. The scale of $B$ pins the shares of the minors it multiplies,
\[
\sum_j\big(a^r_j+b^d_j\big)=\kappa_B,\qquad\text{equivalently}\qquad \sum_j\delta^d_j=\tfrac12\Big(\sum_js_j-\kappa_B\Big),
\]
and, for pattern shares with these level sums, equivalently $\sum_{e\in B}\langle\boldsymbol\delta\rangle_e=\mu_B$, the sum over the bunch of the pins $\langle\boldsymbol\delta\rangle_e=\mu_e$ of a wall contour (Definition~\ref{def:wallcontour}). On this pin both families lie on their balance walls at $\mvec|_B$ and $\mvec|_R$, and on every contour on which the stacked Gammas below have arguments of positive real part and both periods converge absolutely,
\begin{multline*}
P_F(\mvec;\mathbf s,\mathbf t)=2^{N-1}\int_{\sum_j\delta^d_j=\frac12(\sum_js_j-\kappa_B)}\ \prod_j\Gst42{a^d_j,\ b^d_j,\ a^r_j,\ b^r_j}{s_j,\ t_j}\\
\times P_{F^{\mathrm{bunch}}(\mathbf a,\mathbf b)}(\mvec|_B;\mathbf a,\mathbf b)\,P_{F^{\mathrm{room}}(\mathbf a,\mathbf b)}(\mvec|_R;\mathbf a,\mathbf b),
\end{multline*}
the fibre integral (Definition~\ref{def:fibre}) in the pair coordinates $(\sigma^d_j,\delta^d_j)_j$, with $\sigma^r_j=\sigma_j-\sigma^d_j$ and $\delta^r_j=\delta_j-\delta^d_j$: its dimension is $2N-1$, and its $N$ even shares $\sigma^d_j$ are free. This display needs every $s_j$ and $t_j$ positive; in general, with $p$ the number of active letters,
\begin{multline*}
P_F(\mvec;\mathbf s,\mathbf t)=\int_{\sum_j(a^r_j+b^d_j)=\kappa_B}\ \prod_{s_j\neq0}\Gst21{a^d_j,\ a^r_j}{s_j}\prod_{t_j\neq0}\Gst21{b^d_j,\ b^r_j}{t_j}\\
\times P_{F^{\mathrm{bunch}}(\mathbf a,\mathbf b)}(\mvec|_B;\mathbf a,\mathbf b)\,P_{F^{\mathrm{room}}(\mathbf a,\mathbf b)}(\mvec|_R;\mathbf a,\mathbf b),
\end{multline*}
the share integral (Definition~\ref{def:fibre}) over the shares $a^r_j$, $s_j\neq0$, and $b^d_j$, $t_j\neq0$, of the minors that the scale of $B$ multiplies, pinned by their sum $\kappa_B$, of dimension $p-1$; with every exponent positive it is the display above, and at $p=1$ it is the value at the one admissible point. For one table with $s,t>0$ the odd share is fixed, $2\delta^d=s-\kappa_B=\delta\,(1+\varepsilon(B))-\mu_B$, and the even share alone is free: on a line $\re\sigma^d=\sigma^d_\ast$ of such a contour,
\[
P_F(\mvec;s,t)=\frac1{2\pi i}\int_{\re\sigma^d=\sigma^d_\ast}\Gst42{a^d,\ b^d,\ a^r,\ b^r}{s,\ t}\,P_{F^{\mathrm{bunch}}(\mathbf a,\mathbf b)}(\mvec|_B;\mathbf a,\mathbf b)\,P_{F^{\mathrm{room}}(\mathbf a,\mathbf b)}(\mvec|_R;\mathbf a,\mathbf b)\,d\sigma^d .
\]
Swapping $B$ and $R$ swaps the two levels and leaves the formula unchanged.
\end{theorem}

The scale of the bunch is opened letter by letter, one share on each level of each letter; the scale integrates to the pin, and what is left is the form of the bunch family on its simplex times the form of the room family on its own. The level shares are the pattern shares of the port (Theorem~\ref{thm:port}) summed over each level: at width one the children of the level $d$ are all $\cB_j|R$ and those of the level $r$ all $\cB_j/B$, up to constant factors of their weights (Theorem~\ref{thm:levels}), so the room family is the children of the port with equal children merged. At the level shares $(a^d_j,b^d_j,a^r_j,b^r_j)=(s_j,0,0,t_j)$, on the wall $\kappa_B=0$, the two families are the faces of Theorem~\ref{thm:wallres}. Left unopened on the room, as in Theorem~\ref{thm:pooling}, the shape enters the pooled family only through the hinge values of $B$ by its letters, and a rescaling of the hinge variable absorbs one of them: the shares that separate the shape from the room are one fewer than the hinge values of $B$ by the active letters counted up to constant factors, none when every active letter reads $B$ alike (Proposition~\ref{prop:readingsagree}), as for one table at a full lean (Theorem~\ref{thm:cornerfactor}).

% Reference 20: rem:pinnedproduct; source line 4946 (one table at a full lean).
% Proved in the last part: cor:pinnedproduct.
\begin{corollary}[One Skeleton value pins the share]\label{cor:pinnedproduct}
Let $F=((\cB_j,\theta_j);s_j,t_j)_{j=1}^N$ have real exponents and be on the balance wall with every proper wall positive, and let a bunch $B$ have width one in every table, the weights factoring across $B$ and $R$. Suppose that every active letter reads $B$ alike, its hinge value a constant multiple of one function $\gamma$ on $\Delta_B$, and reads $R$ alike with the same constant, its hinge value of $R$ that constant times one function $\gamma_R$ on $\Delta_R$: the active letters have the one Skeleton value $\gamma_R(v)/\gamma(u)$ on the scale of $B$ at every gauge $(u,v)\in\Delta_B\times\Delta_R$ (Corollary~\ref{cor:widthonestage}). Then $0<\kappa_B<2\sigma$, $\sigma$ the total strength, and for all level shares $(\mathbf a,\mathbf b)$ of Theorem~\ref{thm:twosides}, those of the inactive letters $0$, with $a^d_j+a^r_j=s_j$, $b^d_j+b^r_j=t_j$ and $\sum_j(a^r_j+b^d_j)=\kappa_B$,
\[
P_F(\mvec;\mathbf s,\mathbf t)=\Gst21{\kappa_B,\ 2\sigma-\kappa_B}{2\sigma}\,P_{F^{\mathrm{bunch}}(\mathbf a,\mathbf b)}(\mvec|_B;\mathbf a,\mathbf b)\,P_{F^{\mathrm{room}}(\mathbf a,\mathbf b)}(\mvec|_R;\mathbf a,\mathbf b),
\]
the product of the two periods not depending on the level shares: the Beta corner of the hinge times the periods of the bunch family and the room family. One active letter has one Skeleton value: for one table at $(0,t)$, at the level shares $(a^d,b^d,a^r,b^r)=(0,\kappa_B,0,t-\kappa_B)$, the two periods are those of the forms $u^{\mvec|_B-\one}\Psi_{\cB/R}^{-\kappa_B}\Psi_{\cB|B}^{-(t-\kappa_B)}$ on $\Delta_B$ and $v^{\mvec|_R-\one}\Psi_{\cB|R}^{-\kappa_B}\Psi_{\cB/B}^{-(t-\kappa_B)}$ on $\Delta_R$, and at $(s,0)$, at the level shares $(s-\kappa_B,0,\kappa_B,0)$, the same holds with the letters of $K$.
\end{corollary}

% Reference 20: lem:slopesrange; source line 4192 (with the slopes of cor:closetwosides(i), source line 3982; one table).
% Proved in the last part: lem:slopesrange.
\begin{lemma}[Slopes and the range of the ratios]\label{lem:slopesrange}
In Theorem~\ref{thm:twosides} with every $s_j$ and $t_j$ positive, the even share $\sigma^d_j$ of table $j$ enters the two families only through $\big(\varrho_{j,B}(u)/\varrho_{j,R}(v)\big)^{-\sigma^d_j}$, $(u,v)\in\Delta_B\times\Delta_R$, the ratio of the hinge values of the bunch over that of the room in table $j$ (Definition~\ref{def:readings}). The wall at $S$ of either family is affine in the shares, with the \emph{slope} in $\sigma^d_j$ the width of $S$ in the table of the level $d$ of table $j$ less its width in the table of the level $r$ of table $j$: with $r^+_j$ the rank function of the cell of table $j$ on the side of $S$,
\begin{gather*}
\big[r^+_j(S+g)-r^+_j(S)\big]+\big[r^+_j(B\setminus S+g)-r^+_j(B\setminus S)\big]-1\qquad(S\subseteq B),\\
1-\big[r^+_j(S+g)-r^+_j(S)\big]-\big[r^+_j(R\setminus S+g)-r^+_j(R\setminus S)\big]\qquad(S\subseteq R).
\end{gather*}
$\sup_{\Delta_B}\varrho_{j,B}<\infty$ exactly when no slope of table $j$ on $B$ is positive and $\inf_{\Delta_B}\varrho_{j,B}>0$ exactly when none on $B$ is negative; $\inf_{\Delta_R}\varrho_{j,R}>0$ exactly when no slope of table $j$ on $R$ is positive and $\sup_{\Delta_R}\varrho_{j,R}<\infty$ exactly when none on $R$ is negative. So a wall of either family falls as $\sigma^d_j$ moves toward the level $d$ only if $\sup\varrho_{j,B}/\inf\varrho_{j,R}=\infty$, and as it moves toward the level $r$ only if $\inf\varrho_{j,B}/\sup\varrho_{j,R}=0$. These conditions, like the quotient $\varrho_{j,B}/\varrho_{j,R}$, do not depend on the factorization. For one table the index $j$ is omitted.
\end{lemma}

A wall that falls as the even share moves toward the level being closed on is therefore the same thing as the two levels coming to overlap at that wall (Proposition~\ref{prop:overlap}): the quotient $\varrho_B/\varrho_R$ of the even letters of the two levels, below one where the closing starts, is unbounded near the face where the wall sits.

\section{Closing at a wall and closing at the gauge}\label{sec:wallgauge}

A share between the levels of a letter can be closed in two places. \emph{At the gauge} every letter is expanded about its value at a point of the scale; the terms are nonnegative and nothing is asked of the Skeleton values. That closing is Theorem~\ref{thm:sectors}, and it closes every letter on the scale at once. \emph{At a wall} a letter is expanded about one of its two extreme coefficients, the other coefficients entering as integer powers; the terms then separate over the bunch and the room, and the expansion holds as far as the letters are separated along the scale. The closing at a wall holds for any number of letters and any width; across a bunch of width one it is a lattice of products of periods of the bunch family and the room family, and it stops where the values no longer stand in one order, which for one table is where the two levels overlap (Proposition~\ref{prop:overlap}).

% Reference 20: section preamble of sec:wallgauge; source line 4012
\begin{definition}[Letters on the scale]\label{def:scaleletters}
Let $F=((\cB_j,\theta_j);s_j,t_j)_{j=1}^N$ have real masses and real exponents of either sign, let $B$ be a proper bunch with room $R$, and read $\alpha_B=\rho u$ with $u\in\Delta_B$. At the gauge $(u,\alpha_R)$ the \emph{letters on the scale} of $B$ are the $2N$ polynomials in $\rho$
\[
\rho^{-d_j(B)}K_j(\rho u,\alpha_R)=\sum_{l=0}^{c_j(B)}\rho^l\,K^{(d_j(B)+l)}_{j,B}(u,\alpha_R),\qquad \rho^{-d^\ast_j(B)}\Psi_j(\rho u,\alpha_R)=\sum_{l=0}^{c_j(B)}\rho^l\,\Psi^{(d^\ast_j(B)+l)}_{j,B}(u,\alpha_R),
\]
with the letter exponents $s_j$ and $t_j$: their coefficients are the level and colevel letters of Definition~\ref{def:levels} counted from their floors, and divided by their values at $\rho=1$ they are the Skeleton polynomials $\sQ^K_{j,B}(\rho;(u,\alpha_R))$ and $\sQ^\Psi_{j,B}(\rho;(u,\alpha_R))$ (Definition~\ref{def:skelpoly}). Listed as $L_i(\rho)=\sum_{l=0}^{c_i}\rho^lL^{(l)}_i$, $i=1,\dots,2N$, with the letter exponent $w_i$ and the width $c_i=c_j(B)$ of their table, every coefficient is positive (Theorem~\ref{thm:taylorcube}), a letter of width zero is constant, and the scale integral of Theorem~\ref{thm:scale} at the gauge $(u,\alpha_R)$ is
\[
\operatorname{Sc}_B(u,\alpha_R;\mvec;\mathbf s,\mathbf t)=\int_0^\infty\rho^{\kappa_B-1}\prod_iL_i(\rho)^{-w_i}\,d\rho ,
\]
convergent exactly when $0<\kappa_B<\sum_iw_ic_i$, where $\sum_iw_ic_i=\kappa_B+\kappa^{\mathrm{top}}_B$.
\end{definition}

% Reference 20: def:readends; source line 4022
\begin{definition}[Letters read from an end of the scale; separated groups]\label{def:readends}
Fix $\lambda\in(0,1)$. A letter $L(\rho)=\sum_{l=0}^{c}\rho^lL^{(l)}$ with positive coefficients is \emph{read from the bottom up to $a$} when $\sum_{l\ge1}a^lL^{(l)}\le\lambda\,L^{(0)}$, and \emph{read from the top down to $a$} when $\sum_{l<c}a^lL^{(l)}\le\lambda\,a^cL^{(c)}$. A partition of the letters on the scale of $B$ into \emph{groups} $\mathcal G_1,\dots,\mathcal G_q$ is \emph{separated} by the scales $0<a_1<\dots<a_{q-1}$ when every letter of $\mathcal G_k$ is read from the bottom up to $a_{k-1}$ (for $k\ge2$) and from the top down to $a_k$ (for $k\le q-1$); a letter whose exponent is a nonpositive integer is a polynomial factor of the integrand, may be put in any group, and nothing is asked of it. A group $\mathcal G$ has the product $f_{\mathcal G}=\prod_{i\in\mathcal G}L_i^{-w_i}$, the \emph{scale degree} $w_{\mathcal G}=\sum_{i\in\mathcal G}w_ic_i$, and the \emph{bottom} and \emph{top coefficients} $\beta^\downarrow_{\mathcal G,n}$ and $\beta^\uparrow_{\mathcal G,n}$,
\[
f_{\mathcal G}(\rho)=\sum_{n\ge0}\beta^\downarrow_{\mathcal G,n}\,\rho^{n}\quad(\rho\text{ small}),\qquad f_{\mathcal G}(\rho)=\rho^{-w_{\mathcal G}}\sum_{n\ge0}\beta^\uparrow_{\mathcal G,n}\,\rho^{-n}\quad(\rho\text{ large}).
\]
\end{definition}

Below the scale $a$ a letter read from the bottom is its lowest coefficient up to the factor $1+\lambda$, and above $a$ a letter read from the top is its top term $\rho^cL^{(c)}$ up to the same factor. At width one, $L(\rho)=L^{(0)}+\rho L^{(1)}$ has the single Skeleton value $\xi=L^{(0)}/L^{(1)}$ (Corollary~\ref{cor:widthonestage}); it is read from the bottom up to $a$ exactly when $a\le\lambda\xi$ and from the top down to $a$ exactly when $a\ge\xi/\lambda$. So single letters of width one with the values $\xi_1,\dots,\xi_q$ are separated, by $a_k=\sqrt{\xi_k\xi_{k+1}}$, exactly when $\xi_k\le\lambda^2\xi_{k+1}$ for every $k$: their values stand in one order with room between them.

% Reference 20: thm:groupsclose; source line 4031
\begin{theorem}[The scale closes between separated groups]\label{thm:groupsclose}
At the gauge $(u,\alpha_R)$ let the letters on the scale of $B$ be separated into groups $\mathcal G_1,\dots,\mathcal G_q$, let $0<\kappa<\sum_iw_ic_i$, and suppose $\kappa-\sum_{k'\le k}w_{\mathcal G_{k'}}\notin\mathbb Z$ for $1\le k<q$. For $1\le k\le q$ and $\mathbf n=(n_{k'})_{k'\neq k}\in\mathbb N^{q-1}$ put
\[
\nu_k(\mathbf n)=\kappa-\sum_{k'<k}\big(w_{\mathcal G_{k'}}+n_{k'}\big)+\sum_{k'>k}n_{k'} .
\]
\begin{enumerate}[label=(\roman*)]
\item
\[
\int_0^\infty\rho^{\kappa-1}\prod_iL_i(\rho)^{-w_i}\,d\rho=\sum_{k=1}^{q}\ \sum_{\mathbf n}\ \prod_{k'<k}\beta^\uparrow_{\mathcal G_{k'},n_{k'}}\prod_{k'>k}\beta^\downarrow_{\mathcal G_{k'},n_{k'}}\ \Reg\!\int_0^\infty\rho^{\nu_k(\mathbf n)-1}f_{\mathcal G_k}(\rho)\,d\rho ;
\]
at $\kappa=\kappa_B$ the left side is $\operatorname{Sc}_B(u,\alpha_R;\mvec;\mathbf s,\mathbf t)$.
\item The series converges absolutely, and the sum of the absolute values of its terms is at most $C$ times the left side, where $C$ depends only on $\lambda$, the letter exponents, the widths, $\kappa$ and the distances of the numbers $\kappa-\sum_{k'\le k}w_{\mathcal G_{k'}}$ from $\mathbb Z$.
\item $\beta^\downarrow_{\mathcal G,n}=\prod_{i\in\mathcal G}\big(L^{(0)}_i\big)^{-w_i}\,[\rho^n]\prod_{i\in\mathcal G}\big(1+\sum_{l\ge1}\rho^lL^{(l)}_i/L^{(0)}_i\big)^{-w_i}$: the bottom coefficients of a group are the integrand of Theorem~\ref{thm:wallres} for the residue at $\kappa_B=-n$, taken for the letters of the group alone; the top coefficients are the same with the coefficients $L^{(c_i)}_i$ in place of $L^{(0)}_i$, for the residue at $\kappa^{\mathrm{top}}_B=-n$.
\end{enumerate}
\end{theorem}

Each group in turn keeps the scale. The groups that lie above it on the scale are read from the bottom, the groups below it from the top, and what is left to integrate is the keeping group alone, at an index shifted by the integers of the others, and as a finite part when that index has left the interval $(0,w_{\mathcal G})$ in which the group's own integral converges. The coefficients are the wall residues the family already has: for each letter read, a power of its extreme coefficient times a polynomial in its other coefficients. Every coefficient of $K_j$ is a sum of the pattern terms $u^pK_{j,(p)}(\alpha_R)$ and every coefficient of $\Psi_j$ of the terms $u^{B\setminus p}\Psi_{j,(p)}(\alpha_R)$, so the children enter with their pattern index, as in the record of a peel (Definition~\ref{def:record}).

% Reference 20: cor:twogroups; source line 4089
\begin{corollary}[Two groups: the walls of one against the continuation of the other]\label{cor:twogroups}
Let the letters be separated into a lower group $\mathcal L$ and an upper group $\mathcal U$, with $0<\kappa<\sum_iw_ic_i$ and $\kappa-w_{\mathcal L}\notin\mathbb Z$. Then
\[
\int_0^\infty\rho^{\kappa-1}f_{\mathcal L}f_{\mathcal U}\,d\rho=\sum_{n\ge0}\beta^\downarrow_{\mathcal U,n}\,\Reg\!\int_0^\infty\rho^{\kappa+n-1}f_{\mathcal L}\,d\rho+\sum_{n\ge0}\beta^\uparrow_{\mathcal L,n}\,\Reg\!\int_0^\infty\rho^{\kappa-w_{\mathcal L}-n-1}f_{\mathcal U}\,d\rho .
\]
For any group, $\Reg\!\int_0^\infty\rho^{\nu-1}f_{\mathcal G}\,d\rho$ is the value at $\nu$ of a meromorphic function, and at a pole of that function its constant term there (Theorem~\ref{thm:RegCT}); the poles are at most simple and lie among the points $\nu=-n$, of residue $\beta^\downarrow_{\mathcal G,n}$, and $\nu=w_{\mathcal G}+n$, of residue $-\beta^\uparrow_{\mathcal G,n}$, $n\in\mathbb N$; the two families of points meet only when $w_{\mathcal G}\in-\mathbb N$, and at a common point the two residues add.
\end{corollary}

The scale integral of a product closes on the walls of each factor. The residues of the upper group at the wall of $B$ are paired with the scale integral of the lower group continued past its top, and the residues of the lower group at the top wall with the scale integral of the upper group continued past its bottom. No bound on a moved contour is used: both expansions hold on the stretch of the scale between the groups, and the scale is divided there.

% Reference 20: thm:jointclosing; source line 4106
\begin{theorem}[Joint closing across a bunch of width one]\label{thm:jointclosing}
Let $F=((\cB_j,\theta_j);s_j,t_j)_{j=1}^N$ have real data, its masses on the balance wall with $P_F(\mvec;\mathbf s,\mathbf t)$ convergent, and let a bunch $B$ have width at most one in every table, the weights factoring across $B$ and $R$. A table in which $B$ has width zero is the direct sum of its restrictions to $B$ and to $R$ (Proposition~\ref{prop:directsum}); its two letters on the scale are constants, and it enters the bunch family and the room family of Theorem~\ref{thm:twosides} as those restrictions with its own exponents. For the $N'$ tables in which $B$ has width one the letters on the scale (Definition~\ref{def:scaleletters}) are, on $\Delta_B\times\Delta_R$,
\[
K_{\cB_j/R}(u)K_{\cB_j|R}(v)+\rho\,K_{\cB_j|B}(u)K_{\cB_j/B}(v),\qquad \Psi_{\cB_j|B}(u)\Psi_{\cB_j/B}(v)+\rho\,\Psi_{\cB_j/R}(u)\Psi_{\cB_j|R}(v),
\]
of the letter exponents $s_j$ and $t_j$, with the Skeleton values of Corollary~\ref{cor:widthonestage}; in each the coefficient multiplying $\rho$ is its \emph{scale coefficient}, the level letter of the level $r$ for $K_j$ and the colevel letter of the level $d$ for $\Psi_j$, and the other its \emph{constant coefficient}. Suppose the $2N'$ letters can be listed as $\ell=1,\dots,2N'$, with letter exponents $w_\ell$ and Skeleton values $\xi_\ell$, so that for some $\vartheta<1$
\[
\xi_\ell\le\vartheta\,\xi_{\ell+1}\quad\text{on }\Delta_B\times\Delta_R\qquad(1\le\ell<2N'),
\]
and that $\kappa_B-\sum_{\ell'\le\ell}w_{\ell'}\notin\mathbb Z$ for $1\le\ell<2N'$. For a letter $i$ and $\mathbf n=(n_\ell)_{\ell\neq i}\in\mathbb N^{2N'-1}$ put $\nu_i(\mathbf n)=\kappa_B-\sum_{\ell<i}(w_\ell+n_\ell)+\sum_{\ell>i}n_\ell$, and let $(\mathbf a,\mathbf b)_{i,\mathbf n}$ be the level shares in which
\begin{itemize}
\item a letter $\ell<i$ has its constant coefficient at $-n_\ell$ and its scale coefficient at $w_\ell+n_\ell$,
\item a letter $\ell>i$ has its scale coefficient at $-n_\ell$ and its constant coefficient at $w_\ell+n_\ell$,
\item the letter $i$ has its scale coefficient at $\nu_i(\mathbf n)$ and its constant coefficient at $w_i-\nu_i(\mathbf n)$.
\end{itemize}
Then:
\begin{enumerate}[label=(\roman*)]
\item at every $(\mathbf a,\mathbf b)_{i,\mathbf n}$ the bunch family and the room family lie on their balance walls and their periods converge;
\item with both periods at $(\mathbf a,\mathbf b)=(\mathbf a,\mathbf b)_{i,\mathbf n}$,
\begin{multline*}
P_F(\mvec;\mathbf s,\mathbf t)=\sum_{i=1}^{2N'}\ \sum_{\mathbf n\in\mathbb N^{2N'-1}}\ \prod_{\ell\neq i}\frac{(-1)^{n_\ell}(w_\ell)_{n_\ell}}{n_\ell!}\ \Gst21{\nu_i(\mathbf n),\ w_i-\nu_i(\mathbf n)}{w_i}\\
\times P_{F^{\mathrm{bunch}}(\mathbf a,\mathbf b)}(\mvec|_B;\mathbf a,\mathbf b)\,P_{F^{\mathrm{room}}(\mathbf a,\mathbf b)}(\mvec|_R;\mathbf a,\mathbf b),
\end{multline*}
where the stacked Gamma is $0$ when $w_i\in-\mathbb N$; when no letter exponent lies in $-\mathbb N$ the coefficient is $\prod_{\ell\neq i}\frac{(-1)^{n_\ell}}{n_\ell!}\,\Gst{2N'+1}{2N'}{(w_\ell+n_\ell)_{\ell\neq i},\ \nu_i(\mathbf n),\ w_i-\nu_i(\mathbf n)}{(w_\ell)_\ell}$;
\item the series converges absolutely, the absolute values of all its terms add up to at most $C\,P_F(\mvec;\mathbf s,\mathbf t)$, and those with $|\mathbf n|=M$ to at most $C'(1+M)^D\vartheta^M\,P_F(\mvec;\mathbf s,\mathbf t)$, where $C$, $C'$ and $D$ depend only on $\vartheta$, the exponents and $\kappa_B$.
\end{enumerate}
\end{theorem}

The $2N'$ Skeleton values stand in one order on the whole shape domain, and each letter in turn keeps the scale. A letter whose value lies below the keeper's is read from the top of the scale, so its constant coefficient is opened; a letter whose value lies above is read from the bottom, so its scale coefficient is opened: each letter opens the one of its two terms that is the smaller at the keeper's place on the scale. An opened coefficient sits at the share $-n$, which puts the $n$-th power of that coefficient into the forms of the two families, a moment; the keeper carries the Beta corner $\Gst21{\nu,\ w-\nu}{w}$ of a linear letter (Proposition~\ref{prop:scalebreta} with one coefficient), continued when $\nu$ has left $(0,w)$. So closing across the bunch pairs moments of the bunch family with moments of the room family on $2N'$ lattices of dimension $2N'-1$, the free shares of Theorem~\ref{thm:twosides} closed at once. Nothing is asked of the walls of the two families, (i) being a conclusion, and nothing uniform in $\mathbf n$ is needed beyond (iii), which bounds the whole lattice by the period itself.

The values are quotients of hinge values, the room's over the bunch's by each letter (Corollary~\ref{cor:widthonestage}). The value of $K_j$ over that of $\Psi_j$ is $\varrho_{j,B}(u)/\varrho_{j,R}(v)$, the even letters of its two levels, and their product is the corresponding quotient of the odd letters: the two values of a table sit on either side of the geometric mean set by its odd letters, at the distance set by its even letters. The $N'$ quotients of even letters and the $N'-1$ quotients of the products of one table by those of another are what the $N'$ even shares and the $N'-1$ free odd shares of Theorem~\ref{thm:twosides} carry; for one table only $\varrho_B/\varrho_R$ is left.

% Reference 20: cor:closetwosides; source line 3979
\begin{corollary}[Closing the even share]\label{cor:closetwosides}
For one table in Theorem~\ref{thm:twosides} let $f(\sigma^d)$ be the integrand of its one-table display.
\begin{enumerate}[label=(\roman*)]
\item If $\sup_{\Delta_B}\varrho_B<\inf_{\Delta_R}\varrho_R$, then $P_F(\mvec;s,t)=\sum\operatorname{Res}f$ over the poles of the Gamma pair $\Gamma(a^d)\Gamma(b^d)$ of the level $d$, at $\sigma^d=\delta^d-n$ and $\sigma^d=-\delta^d-n$, $n\in\mathbb N$, the $n$-th terms of the size $n^{s+t-2}\big(\sup\varrho_B/\inf\varrho_R\big)^n$; every period at these points converges, and when $2\delta^d\notin\mathbb Z$ the $n$-th terms are at most $C(1+n)^D\big(\sup\varrho_B/\inf\varrho_R\big)^nP_F(\mvec;s,t)$, with $C$ and $D$ depending only on that ratio, the exponents and $\kappa_B$ (Theorem~\ref{thm:jointclosing}(i),(iii) with one table, its letters listed as $K$, then $\Psi$). If $\inf_{\Delta_B}\varrho_B>\sup_{\Delta_R}\varrho_R$, then $P_F(\mvec;s,t)=-\sum\operatorname{Res}f$ over the poles of the Gamma pair of the level $r$, at $\sigma^r=\pm\delta^r-n$. The share closes toward the level whose even letter is the smaller.
\item At $\sigma^d=\delta^d-n$ the level shares are $(a^d,b^d,a^r,b^r)=(2\delta^d-n,\,-n,\,s-2\delta^d+n,\,t+n)$, and the residue is
\[
\frac{(-1)^n}{n!}\Gst32{2\delta^d-n,\ s-2\delta^d+n,\ t+n}{s,\ t}\,P_{F^{\mathrm{bunch}}(\mathbf a,\mathbf b)}(\mvec|_B;\mathbf a,\mathbf b)\,P_{F^{\mathrm{room}}(\mathbf a,\mathbf b)}(\mvec|_R;\mathbf a,\mathbf b),
\]
where $P_{F^{\mathrm{bunch}}(\mathbf a,\mathbf b)}=\int_{\Delta_B}u^{\mvec|_B-\one}K_{\cB/R}^{-2\delta^d}K_{\cB|B}^{-(s-2\delta^d)}\Psi_{\cB|B}^{-t}\,\varrho_B^{\,n}\,d\omega_B$ is the $n$-th moment of $\varrho_B$ against a form that does not depend on $n$, and $P_{F^{\mathrm{room}}(\mathbf a,\mathbf b)}$ likewise the $n$-th moment of $\varrho_R^{-1}$. At $\sigma^d=-\delta^d-n$ the level shares are $(-n,\,-2\delta^d-n,\,s+n,\,t+2\delta^d+n)$, and the same holds with the form $\Psi_{\cB/R}^{2\delta^d}K_{\cB|B}^{-s}\Psi_{\cB|B}^{-t-2\delta^d}$ on $B$ and its counterpart on $R$. Closing pairs the moments of the ratios $\varrho_B$ and $\varrho_R$ of the bunch and of the room. The two lattices of a level meet exactly when $2\delta^\ell\in\mathbb Z$, at the centre exactly when $\mu_B\in\mathbb Z$. At a common point, where $\Gamma(a^d)$ sits at $-k$ and $\Gamma(b^d)$ at $-n$, the pole is double: with $h=\Gst22{s+k,\ t+n}{s,\ t}P_{F^{\mathrm{bunch}}(\mathbf a,\mathbf b)}P_{F^{\mathrm{room}}(\mathbf a,\mathbf b)}$,
\begin{multline*}
\operatorname*{Res}f=\frac{(-1)^{n+k}}{n!\,k!}\,h\,\Big((\log\Gamma)'(n+1)-(\log\Gamma)'(t+n)+(\log\Gamma)'(k+1)-(\log\Gamma)'(s+k)\\
+\partial_{\sigma^d}\log\big(P_{F^{\mathrm{bunch}}(\mathbf a,\mathbf b)}P_{F^{\mathrm{room}}(\mathbf a,\mathbf b)}\big)\Big).
\end{multline*}
At integer exponents the two differences of $(\log\Gamma)'$ are heads, finite sums of reciprocals (Proposition~\ref{prop:headstails}), and $(\log\Gamma)'(j+1)=-\Tx{1}{\IZ}{j}$ for $j\in\mathbb N$ (Proposition~\ref{prop:resonance}).
\end{enumerate}
\end{corollary}

Closed on a lattice, the even share writes the period as a row over finitely many data of the room whenever the room's periods along the lattice satisfy a recurrence of finite order: the room's initial periods, and at a meeting of the two lattices their derivatives, against sums along the lattice of the bunch's data and the room's fundamental solutions (Proposition~\ref{prop:ra-hingerow}).

Off resonance every pole is simple and the closed value is a sum over two lattices of products of moments of the two ratios $\varrho_B$ and $\varrho_R$, with stacked Gamma coefficients: no heads. Each lattice's series has a pole where the two lattices of the level meet, from $\Gamma(2\delta^d-n)$ and $\Gamma(-2\delta^d-n)$, and their sum tends to the double-pole residues of (ii); the heads at the centre are this meeting. When the ranges of the even letters of the two levels overlap, the quotient $\varrho_B(u)/\varrho_R(v)$ taking values on both sides of $1$, neither closing of (i) applies; the residues of a family at one of its walls are those of Theorem~\ref{thm:wallres}, periods of the family's faces.

% Reference 20: cor:nested; source line 4205
\begin{corollary}[Joint closing, nested]\label{cor:nested}
In Theorem~\ref{thm:jointclosing} let $B'$ be a bunch of the bunch family, of width at most one in each of its tables, with factoring weights, and let the letters of the bunch family on the scale of $B'$ have their Skeleton values in one order with a ratio $\vartheta'<1$. Then at every $(\mathbf a,\mathbf b)_{i,\mathbf n}$ at which the non-integrality condition of Theorem~\ref{thm:jointclosing} holds for the bunch family at $B'$, $P_{F^{\mathrm{bunch}}(\mathbf a,\mathbf b)}(\mvec|_B;\mathbf a,\mathbf b)$ is given by Theorem~\ref{thm:jointclosing} for the family $F^{\mathrm{bunch}}(\mathbf a,\mathbf b)$ at $B'$; likewise for the room family at a bunch of $R$. Continuing down a tree of splits of width one, each node splitting its family at a bunch of width one into the bunch family and the room family, with the values ordered at every node, $P_F(\mvec;\mathbf s,\mathbf t)$ is a nested lattice sum of stacked Gammas and of the periods of the families at the leaves, each leaf at exponents shifted by the integers of the lattices above it. The sum is taken from the leaves up, and at each node it converges absolutely at that node's own root rate.
\end{corollary}

The tree of the nesting is a tree of splits of width one and not the Foreman tree of Definition~\ref{def:tree}: a node may be any bunch of width one at which the values are ordered. At a lattice point of the node above, a letter of the family below sits at an exponent $-n$ or $w+n$; the letters below are the same letters at shifted powers, and the next split applies to them because whether the values are in one order is read off the tables. Nothing uniform in the outer integers is used: at each lattice point the inner sum is an identity. The constant $C$ of Theorem~\ref{thm:jointclosing}(iii) at an inner node depends on the exponents there and grows at most geometrically in the outer integers; that matters only for summing in another order. The multiple series converges absolutely when the outer ratio $\vartheta$ beats that growth, and otherwise the sum is the iterated one.

% Reference 20: prop:weakorder; source line 4214
\begin{proposition}[Values in one order without room]\label{prop:weakorder}
In Theorem~\ref{thm:jointclosing} let $\xi_\ell<\xi_{\ell+1}$ on $\Delta_B\times\Delta_R$ for every $\ell$, with no $\vartheta<1$. For $0<\zeta<1$ let $P^{(\zeta)}$ be the integral obtained from $P_F(\mvec;\mathbf s,\mathbf t)$ by multiplying the constant coefficient of the $\ell$-th letter by $\zeta^{-\ell}$; it is the period of the family in which $K_j$ and $\Psi_j$ are read from two copies of table $j$ whose level weights are scaled accordingly. Then every period of Theorem~\ref{thm:jointclosing}(i) converges, and with $A_{i,\mathbf n}$ the term of index $(i,\mathbf n)$ of the series of Theorem~\ref{thm:jointclosing}(ii) and $x_i=\sum_{\ell>i}\ell w_\ell+i\sum_{\ell\le i}w_\ell-i\kappa_B$,
\[
P^{(\zeta)}=\sum_{i,\mathbf n}\zeta^{\,x_i+\sum_{\ell}|\ell-i|\,n_\ell}A_{i,\mathbf n}\quad\text{absolutely},\qquad P^{(\zeta)}\to P_F(\mvec;\mathbf s,\mathbf t)\ \ (\zeta\uparrow1).
\]
So the lattice sum of Theorem~\ref{thm:jointclosing}(ii) is summable to the period by scaling the level weights, and equals the period whenever it converges absolutely.
\end{proposition}

% Reference 20: prop:gaugeclose; source line 4229 (nonnegative exponents).
% Proved in the last part: prop:gaugeclose.
\begin{proposition}[The same shares closed at the gauge]\label{prop:gaugeclose}
Let $F=((\cB_j,\theta_j);s_j,t_j)_{j=1}^N$ have complex exponents and masses $\mvec$ and let a bunch $B$ have width one in every table, the weights factoring across $B$ and $R$, with $\re\kappa_B>0$ and $\re\kappa^{\mathrm{top}}_B>0$. Write the Skeleton polynomials at the gauge $(u,v)\in\Delta_B\times\Delta_R$ as
\[
\sQ^K_{j,B}\big(x;(u,v)\big)=\pi_{j,d}+\pi_{j,r}\,x,\qquad \sQ^\Psi_{j,B}\big(x;(u,v)\big)=\pi^\ast_{j,r}+\pi^\ast_{j,d}\,x,
\]
the chances of the two levels under the basis law and the cobasis law at $(u,v)$, each the sum over its level of the pattern probabilities of Definition~\ref{def:split}. Then, with $M=|\mathbf n|+|\mathbf n'|$,
\begin{multline*}
\operatorname{Sc}_B(u,v;\mvec;\mathbf s,\mathbf t)=\prod_jK_j(u,v)^{-s_j}\Psi_j(u,v)^{-t_j}\sum_{\mathbf n,\mathbf n'\in\mathbb N^N}\prod_j\frac{(s_j)_{n_j}(t_j)_{n'_j}}{n_j!\,n'_j!}\\
\times\Big(\prod_j\pi_{j,r}^{\,n_j}\pi_{j,d}^{\ast\,n'_j}\ \Gst21{\kappa_B,\ M+1}{\kappa_B+M+1}+\prod_j\pi_{j,d}^{\,n_j}\pi_{j,r}^{\ast\,n'_j}\ \Gst21{\kappa^{\mathrm{top}}_B,\ M+1}{\kappa^{\mathrm{top}}_B+M+1}\Big),
\end{multline*}
a series converging absolutely, of nonnegative terms when the data are real with $s_j,t_j\ge0$, whatever the order of the Skeleton values: it is Theorem~\ref{thm:sectors} at width one, whose exceedance polynomials are the constants $\pi_{j,r}$, $\pi^\ast_{j,d}$ on the bunch and $\pi_{j,d}$, $\pi^\ast_{j,r}$ on the room. Without the prefactor the terms with $M$ fixed add up in absolute value to at most $\frac{(\sum_j(|s_j|+|t_j|))_M}{M!}\big(\lambda_B^M/\re\kappa_B+\lambda_R^M/\re\kappa^{\mathrm{top}}_B\big)$, with $\lambda_B=\max_j(\pi_{j,r},\pi^\ast_{j,d})<1$ and $\lambda_R=\max_j(\pi_{j,d},\pi^\ast_{j,r})<1$.
\end{proposition}

At the gauge every letter is closed about its value at one point of the scale, both levels together. The lattice has one index for each letter, and what is raised to the powers are the chances of the levels at that point. Nothing is asked of the order of the values, and the root rate is the largest chance of a level. The price is that the closed letters $K_j(u,v)$ and $\Psi_j(u,v)$ stay in front, so the terms do not separate over the bunch and the room; integrated over the shapes on their ordering sectors the same closing is Corollary~\ref{cor:cubeclose}, with Theorem~\ref{thm:contcorners} past the walls.

% Reference 20: prop:crossing; source line 4243
\begin{proposition}[Overlapping levels]\label{prop:overlap}
In Theorem~\ref{thm:twosides} with one table let $s,t>0$ and $2\delta^d=s-\kappa_B\notin\mathbb Z$, and let $\mathcal M(x)$ be the integral over $\Delta_B\times\Delta_R$ of
\[
u^{\mvec|_B-\one}v^{\mvec|_R-\one}\big(K_{\cB/R}(u)K_{\cB|R}(v)\big)^{\kappa_B-s}\big(K_{\cB|B}(u)K_{\cB/B}(v)\big)^{-\kappa_B}\big(\Psi_{\cB|B}(u)\Psi_{\cB/B}(v)\big)^{-t}\Big(\frac{\varrho_B(u)}{\varrho_R(v)}\Big)^{x}.
\]
The terms of the two lattices of the level $d$ in Corollary~\ref{cor:closetwosides}(ii) are $A_n=a_n\mathcal M(n)$ and $A'_n=a'_n\mathcal M(n+2\delta^d)$, with
\begin{gather*}
a_n=\Gst22{2\delta^d,\ 1-2\delta^d}{s,\ t}\Gst22{t+n,\ \kappa_B+n}{1+n,\ 1-2\delta^d+n},\\
a'_n=-\Gst22{2\delta^d,\ 1-2\delta^d}{s,\ t}\Gst22{s+n,\ t+2\delta^d+n}{1+n,\ 1+2\delta^d+n}.
\end{gather*}
Suppose these integrals converge for every $n$. If $\varrho_B(u)>\varrho_R(v)$ on a set of positive measure, then $|A_n+A'_n|\to\infty$; likewise for the level $r$ if $\varrho_B(u)<\varrho_R(v)$ on a set of positive measure. So when the two levels \emph{overlap}, both sets having positive measure, the lattices of neither level sum to the period by convergent periods.
\end{proposition}

The two lattices of a level have opposite signs and coefficients whose ratio tends to $-1$. When $\sup\varrho_B<\inf\varrho_R$ each of them converges by itself. When $\varrho_B(u)<\varrho_R(v)$ throughout with $\sup\varrho_B=\inf\varrho_R$ neither need converge; the moments at $n$ and at $n+2\delta^d$ then agree to leading order, so the two lattices cancel to that order, and Proposition~\ref{prop:weakorder} sums them to the period by scaling the level weights. Past that, the two lattices part and the series diverges. By Lemma~\ref{lem:slopesrange} the overlap reaches a wall when $\varrho_B/\varrho_R$ is unbounded, and then the periods themselves stop converging along the lattice. An overlap is closed on the totals of the bunch and the room (Section~\ref{sec:pencil}), and, read in one scale, on the lattice of continued periods of a level whenever the reach of the bunch on its side lies below the decay of the other (Theorem~\ref{thm:nodeintervals}).

% Reference 20: prop:wallthengauge; source line 4264
\begin{proposition}[Between groups at the wall, inside a group at the gauge]\label{prop:wallthengauge}
In Theorem~\ref{thm:groupsclose} fix a group $\mathcal G$, an index $\nu$ with $\nu\notin-\mathbb N$ and $\nu-w_{\mathcal G}\notin\mathbb N$, and a point $\rho_0$ of the scale. For the letter $i$ of $K_j$ in $\mathcal G$ put $\sQ_i=\sQ^K_{j,B}(\cdot;(\rho_0u,\alpha_R))$ and $\sQ^\top_i=\sQ^K_{j,R}(\cdot;(\rho_0u,\alpha_R))$, and for that of $\Psi_j$ the same with $\Psi$; then $\sQ_i(x)=L_i(x\rho_0)/L_i(\rho_0)$ and $\sQ^\top_i(x)=x^{c_i}\sQ_i(1/x)$ (Proposition~\ref{prop:sklaws}(i),(ii)), and
\begin{multline*}
\Reg\!\int_0^\infty\rho^{\nu-1}f_{\mathcal G}(\rho)\,d\rho=\rho_0^{\,\nu}f_{\mathcal G}(\rho_0)\Big(\Reg\!\int_0^1x^{\nu-1}\prod_{i\in\mathcal G}\sQ_i(x)^{-w_i}dx\\
+\Reg\!\int_0^1x^{w_{\mathcal G}-\nu-1}\prod_{i\in\mathcal G}\sQ^\top_i(x)^{-w_i}dx\Big).
\end{multline*}
Each of the two finite parts is the series of Theorem~\ref{thm:sectors}(ii) for these polynomials, absolutely convergent, with its Beta corners continued as in Theorem~\ref{thm:contcorners}. So for every separated partition every share of the scale is closed: between the groups on the lattices of Theorem~\ref{thm:groupsclose}, inside a group at the gauge. With one group this is Theorem~\ref{thm:sectors}; with single letters of width one in one order it is Theorem~\ref{thm:jointclosing}.
\end{proposition}

A group that cannot be split, because the values of its letters do not stand in one order, is kept as one Breta of one bunch (Proposition~\ref{prop:scalebreta}); it is closed about a point of the scale, where no order is needed, and the letters outside it are closed about its ends. The two closings do not compete: the wall closing does as much as the scale separates, and the gauge closing does the rest. When its index lies in $(0,w_{\mathcal G})$ and its letter exponents are nonnegative, a group read on the totals of the bunch and the room closes about constant references, its terms pairing the bunch with the room at width one (Section~\ref{sec:pencil}).

\section{The scale on the totals}\label{sec:pencil}

The scale of a bunch is the quotient of two totals, that of the bunch and that of the room, and the two totals lie on the simplex $\Delta_{\{B,R\}}$ of the two blocks (Theorem~\ref{thm:aggregation}). Read there, a letter on the scale is homogeneous in the totals, with its level or colevel letters as coefficients. At width one it is linear, with its constant coefficient at one vertex of the simplex and its scale coefficient at the other, and its Skeleton value is the point of the scale where its two terms are equal. The free shares of Theorem~\ref{thm:twosides} are the share side of this integral over the totals, the pin being the shape integral over the totals (Theorem~\ref{thm:dirichlet}), and the lattices of Section~\ref{sec:wallgauge} close them. The integral over the totals can also be closed as it stands. On the simplex a linear letter is at most the larger of its two coefficients, whatever their ratio, so it closes about a constant reference with nonnegative terms, which pair moments of the bunch with moments of the room. No order of the values is asked, and the overlap that stops the lattices of a level (Proposition~\ref{prop:overlap}) is not seen.

% Reference 20: def:totalsreading; source line 4284
\begin{definition}[The totals and their letters]\label{def:totalsreading}
The \emph{totals} of a bunch $B$ and its room $R$ are the point $x=(x_B,x_R)$ of the shape simplex $\Delta_{\{B,R\}}$ of the two blocks. Let $\mvec^{\mathrm{tot}}=(m^{\mathrm{tot}}_B,m^{\mathrm{tot}}_R)$ have positive entries, let $\phi$ be a nonnegative measure on $\Delta_B\times\Delta_R$, and let finitely many \emph{letters on the totals}
\[
\widehat L_i(x)=\sum_{l=0}^{c_i}x_B^{\,l}\,x_R^{\,c_i-l}L^{(l)}_i,\qquad i=1,\dots,q,
\]
with letter exponents $w_i\in\mathbb R$, have coefficients $L^{(l)}_i\ge0$ that are measurable functions of the shapes $(u,v)\in\Delta_B\times\Delta_R$, each letter positive on the interior of $\Delta_{\{B,R\}}$ at $\phi$-almost every shape. The \emph{integral over the totals} at the shapes $(u,v)$ is the inner integral below, over $\Delta_{\{B,R\}}$, and its integral against $\phi$ is the Trailmoment on the source $\Delta_{\{B,R\}}\times\Delta_B\times\Delta_R$ (Definition~\ref{def:jointtrail}) with the coefficient function $(x,u,v)\mapsto(\widehat L_i(x))_i$ and the input $x^{\mvec^{\mathrm{tot}}-\one}d\omega(x)\,\phi(du,dv)$, read over the whole source,
\begin{multline*}
\big(\sT_{\mathbf w;I_q}[(\widehat L_i)_i]\;x^{\mvec^{\mathrm{tot}}-\one}d\omega(x)\,\phi(du,dv)\big)\big(\Delta_{\{B,R\}}\times\Delta_B\times\Delta_R\big)\\
=\int_{\Delta_B\times\Delta_R}\int_{\Delta_{\{B,R\}}}x^{\mvec^{\mathrm{tot}}-\one}\prod_i\widehat L_i(x)^{-w_i}\,d\omega(x)\,\phi(du,dv),
\end{multline*}
wherever it converges; its kernel $\sZ_{\mathbf w;I_q}((\widehat L_i)_i)=\prod_i\widehat L_i^{-w_i}$ is the product of the one-scale kernels $\sZ_{w_i;(1)}(\widehat L_i)$ (Definition~\ref{def:zamma}, Theorem~\ref{thm:jointclosure}(i)). For one linear letter $x_RL^{(0)}+x_BL^{(1)}$ the integral over the totals is the Breta with the one bunch $\{B,R\}$, the shape index $\mvec^{\mathrm{tot}}$, the scale index $w$ and the unit configuration (Definition~\ref{def:breta}):
\[
\int_{\Delta_{\{B,R\}}}x^{\mvec^{\mathrm{tot}}-\one}\big(x_BL^{(1)}+x_RL^{(0)}\big)^{-w}\,d\omega(x)=\sB^{\{\{B,R\}\}}_{\mvec^{\mathrm{tot}},\,w;\,I_{\{B,R\}}}\big(L^{(1)},L^{(0)}\big).
\]
\end{definition}

% Reference 20: thm:scaletotals; source line 4300
\begin{theorem}[The scale on the totals]\label{thm:scaletotals}
At the gauge $(u,\alpha_R)$ let $\kappa_B>0$ and $\kappa^{\mathrm{top}}_B>0$, and read every letter on the scale of $B$ (Definition~\ref{def:scaleletters}) on the totals,
\[
\widehat L_i(x)=x_R^{\,c_i}L_i(x_B/x_R)=\sum_{l=0}^{c_i}x_B^{\,l}\,x_R^{\,c_i-l}L^{(l)}_i ,
\]
a letter of width zero being a constant on the totals.
\begin{enumerate}[label=(\roman*)]
\item At every point $\rho_0>0$ of the scale, which may depend on the shapes,
\[
\operatorname{Sc}_B(u,\alpha_R;\mvec;\mathbf s,\mathbf t)=\rho_0^{\kappa_B}\int_{\Delta_{\{B,R\}}}x_B^{\kappa_B-1}x_R^{\kappa^{\mathrm{top}}_B-1}\prod_i\widehat L_i(\rho_0x_B,\ x_R)^{-w_i}\,d\omega(x),
\]
both sides converging together. At $\rho_0=1$ and $\alpha_R=v\in\Delta_R$ the point $(\rho u,v)$ is projectively $(x_Bu,x_Rv)\in\Delta_{\cE}$, whose totals on the blocks $B$ and $R$ are $x_B$ and $x_R$ (Theorem~\ref{thm:aggregation}), with $\rho=x_B/x_R$. On the balance wall $P_F(\mvec;\mathbf s,\mathbf t)$ is the Trailmoment of Definition~\ref{def:totalsreading} for the letters $\widehat L_i$ at $\mvec^{\mathrm{tot}}=(\kappa_B,\kappa^{\mathrm{top}}_B)$ with $\phi=u^{\mvec|_B-\one}v^{\mvec|_R-\one}d\omega_B(u)\,d\omega_R(v)$; both sides converge together.
\item Let $B$ have width at most one in every table, the weights factoring across $B$ and $R$; write the letters of width one $L_i(\rho)=L^{(0)}_i+\rho L^{(1)}_i$, with their constant and scale coefficients and their Skeleton values $\xi_i=L^{(0)}_i/L^{(1)}_i$ (Theorem~\ref{thm:jointclosing}), and fix one of them, $i_0$. At the point $\rho_0=\xi_{i_0}$ the two terms of $i_0$ are equal and its letter on the totals is the constant $L^{(0)}_{i_0}$. The products in the displays run over the letters of width one, and both right sides carry the further factor $\prod L_i^{-w_i}$ over the constant letters of the tables of width zero:
\begin{align*}
\operatorname{Sc}_B(u,v;\mvec;\mathbf s,\mathbf t)&=\big(L^{(0)}_{i_0}\big)^{\kappa_B-w_{i_0}}\big(L^{(1)}_{i_0}\big)^{\kappa^{\mathrm{top}}_B-w_{i_0}}\\
&\qquad\times\int_{\Delta_{\{B,R\}}}x_B^{\kappa_B-1}x_R^{\kappa^{\mathrm{top}}_B-1}\prod_{i\neq i_0}\big(x_RL^{(0)}_iL^{(1)}_{i_0}+x_BL^{(0)}_{i_0}L^{(1)}_i\big)^{-w_i}\,d\omega(x)\\
&=\xi_{i_0}^{\,\kappa_B}\prod_i\big(L^{(0)}_i\big)^{-w_i}\int_{\Delta_{\{B,R\}}}x_B^{\kappa_B-1}x_R^{\kappa^{\mathrm{top}}_B-1}\prod_{i\neq i_0}\Big(x_R+x_B\,\frac{\xi_{i_0}}{\xi_i}\Big)^{-w_i}\,d\omega(x).
\end{align*}
On the balance wall $P_F(\mvec;\mathbf s,\mathbf t)$ is the Trailmoment of Definition~\ref{def:totalsreading} for these linear letters at $\mvec^{\mathrm{tot}}=(\kappa_B,\kappa^{\mathrm{top}}_B)$, with $\phi$ of (i) multiplied by $(L^{(0)}_{i_0})^{\kappa_B-w_{i_0}}(L^{(1)}_{i_0})^{\kappa^{\mathrm{top}}_B-w_{i_0}}$ and by the further factor of the constant letters: the coefficients of each linear letter are the constant coefficient of one letter times the scale coefficient of the other, and their ratio is the ratio of the two Skeleton values.
\item In (ii) let every $w_i>0$, and open each linear letter by Corollary~\ref{thm:E1}, with the share $\nu_i$ on $x_BL^{(0)}_{i_0}L^{(1)}_i$ and $w_i-\nu_i$ on $x_RL^{(0)}_iL^{(1)}_{i_0}$, on a contour with $0<\re\nu_i<w_i$ and $\kappa_B-w_{i_0}<\sum_{i\neq i_0}\re\nu_i<\kappa_B$. The integral over the totals is then the Beta corner $\Gst21{\nu_{i_0},\ w_{i_0}-\nu_{i_0}}{w_{i_0}}$ of the letter $i_0$ at $\nu_{i_0}=\kappa_B-\sum_{i\neq i_0}\nu_i$, so that the shares on the scale coefficients add to $\kappa_B$ identically, and, with the further factor of (ii) on the right,
\[
\operatorname{Sc}_B(u,v;\mvec;\mathbf s,\mathbf t)=\int_{\sum_i\nu_i=\kappa_B}\ \prod_i\Gst21{\nu_i,\ w_i-\nu_i}{w_i}\big(L^{(0)}_i\big)^{-(w_i-\nu_i)}\big(L^{(1)}_i\big)^{-\nu_i},
\]
the share side of the scale integral with one share per letter (Corollary~\ref{cor:sharepercoefficient}). Integrated over the shapes it is Theorem~\ref{thm:twosides}: the pin of the shares is the shape integral over the totals.
\end{enumerate}
\end{theorem}

The scale integral is the integral over the totals, and the shares of Theorem~\ref{thm:twosides} are its share side: opening every letter but one leaves over the totals the Beta corner of the last, and the pin is what that integral asks of the shares. The point $\rho_0$ is free. At $\rho_0=1$ the totals are those of Theorem~\ref{thm:aggregation}, the coordinates of a node in Theorem~\ref{thm:drawtree}; at the value of a letter that letter leaves the integral, and for one table a single linear letter is left.

% Reference 20: cor:pencilone; source line 4335
\begin{corollary}[One table: the even share on the totals]\label{cor:pencilone}
For one table in Theorem~\ref{thm:twosides} write, on $\Delta_B\times\Delta_R$,
\[
\mathbf E^d=\mathbf E_{\cB/R}(u)\,\mathbf E_{\cB|R}(v),\qquad \mathbf E^r=\mathbf E_{\cB|B}(u)\,\mathbf E_{\cB/B}(v)
\]
for the even letters of the two levels, so that $\mathbf E^d/\mathbf E^r=\varrho_B(u)/\varrho_R(v)$, and let $2\delta^d=s-\kappa_B$ and $2\delta^r=\kappa_B-t$; here $\kappa^{\mathrm{top}}_B=s+t-\kappa_B$. Then
\begin{align*}
\operatorname{Sc}_B(u,v;\mvec;s,t)&=\big(K_{\cB/R}(u)K_{\cB|R}(v)\big)^{-2\delta^d}\big(K_{\cB|B}(u)K_{\cB/B}(v)\big)^{-2\delta^r}\,\sB^{\{\{B,R\}\}}_{(\kappa_B,\kappa^{\mathrm{top}}_B),\,t;\,I_{\{B,R\}}}\big(\mathbf E^d,\mathbf E^r\big)\\
&=\big(\Psi_{\cB/R}(u)\Psi_{\cB|R}(v)\big)^{2\delta^d}\big(\Psi_{\cB|B}(u)\Psi_{\cB/B}(v)\big)^{2\delta^r}\,\sB^{\{\{B,R\}\}}_{(\kappa_B,\kappa^{\mathrm{top}}_B),\,s;\,I_{\{B,R\}}}\big(\mathbf E^r,\mathbf E^d\big),
\end{align*}
and $\sB^{\{\{B,R\}\}}_{(\kappa_B,\kappa^{\mathrm{top}}_B),t;I_{\{B,R\}}}(\mathbf E^d,\mathbf E^r)=(\mathbf E^r)^{-t}\,\sB^{\{\{B,R\}\}}_{(\kappa_B,\kappa^{\mathrm{top}}_B),t;I_{\{B,R\}}}\big(\varrho_B(u)/\varrho_R(v),1\big)$. So $P_F(\mvec;s,t)$ is the Trailmoment of Definition~\ref{def:totalsreading} for the single linear letter $x_R\mathbf E^r+x_B\mathbf E^d$ at the exponent $t$ and at $\mvec^{\mathrm{tot}}=(\kappa_B,\kappa^{\mathrm{top}}_B)$, with $\phi=\big(K_{\cB/R}K_{\cB|R}\big)^{-2\delta^d}\big(K_{\cB|B}K_{\cB/B}\big)^{-2\delta^r}u^{\mvec|_B-\one}v^{\mvec|_R-\one}d\omega_B\,d\omega_R$, and the shapes enter the Breta only through the ratio $\varrho_B(u)/\varrho_R(v)$ of its coefficients. Opened by Corollary~\ref{thm:E1}, the linear letter has the share $b^d=\sigma^d-\delta^d$ on $x_B\mathbf E^d$, and the integrand is that of Theorem~\ref{thm:twosides}: the free even share is the share of this letter, and the integral over the totals is the even share integrated.
\end{corollary}

% Reference 20: thm:pencilpeaks; source line 4378
\begin{theorem}[Closing at the peaks, one reference per coefficient]\label{thm:pencilpeaks}
In the Trailmoment of Definition~\ref{def:totalsreading} let every $w_i>0$, and choose constants $\bar L_{i,l}\ge L^{(l)}_i$, one for each coefficient, at $\phi$-almost every shape; the \emph{reference letters} are $\bar L_i(x)=\sum_{l=0}^{c_i}x_B^{\,l}x_R^{\,c_i-l}\bar L_{i,l}$. For an array $\mathbf k=(k_{i,l})$ with the rows $\mathbf k_i=(k_{i,l})_{0\le l\le c_i}$ put
\[
\boldsymbol\beta(\mathbf k)=\big(\beta_B(\mathbf k),\beta_R(\mathbf k)\big)=\Big(m^{\mathrm{tot}}_B+\sum_{i,l}l\,k_{i,l},\ m^{\mathrm{tot}}_R+\sum_{i,l}(c_i-l)\,k_{i,l}\Big),\qquad \mathbf w(\mathbf k)=\big(w_i+|\mathbf k_i|\big)_i .
\]
Introduce one singleton scale $\mathfrak v_i$ per letter, $\mathfrak V=\{\mathfrak v_1,\dots,\mathfrak v_q\}$, take the partition $\cP_{\mathrm{tot}}=\{\{B,R\}\}\cup\{\{\mathfrak v_i\}:1\le i\le q\}$, and let $(\pi^{\mathrm{tot}},\ell^{\mathrm{tot}})$ be the term arrays of the letter presentation (Proposition~\ref{prop:letters}) with the kernel column of shape exponent $0$, level $(1,0_{\mathfrak V})$ and coefficient $1$, and for every $(i,l)$ with $\bar L_{i,l}>0$ a column of shape exponent $\big((l,c_i-l),\mathbf e_i\big)$, level $(0,\mathbf e_i)$ and coefficient $\bar L_{i,l}$; $1\oplus\bar{\mathbf L}$ lists the coefficients in this order. Then, with the arrays $\mathbf k$ supported on these columns,
\begin{multline*}
\int_{\Delta_B\times\Delta_R}\int_{\Delta_{\{B,R\}}}x^{\mvec^{\mathrm{tot}}-\one}\prod_i\widehat L_i(x)^{-w_i}\,d\omega(x)\,\phi(du,dv)\\
=\sum_{\mathbf k}\ \prod_i\frac{(w_i)_{|\mathbf k_i|}}{\mathbf k_i!}\ \sB^{\cP_{\mathrm{tot}}}_{(\boldsymbol\beta(\mathbf k),\mathbf w(\mathbf k)),\,(1,\mathbf w(\mathbf k));\,(\pi^{\mathrm{tot}},\ell^{\mathrm{tot}})}\big(1\oplus\bar{\mathbf L}\big)\int\prod_{i,l}\big(\bar L_{i,l}-L^{(l)}_i\big)^{k_{i,l}}\,d\phi,
\end{multline*}
a series of nonnegative terms, both sides finite or infinite together. The Breta coefficient is a number: the kernel integrates its scale to $1$, and
\[
\sB^{\cP_{\mathrm{tot}}}_{(\boldsymbol\beta(\mathbf k),\mathbf w(\mathbf k)),\,(1,\mathbf w(\mathbf k));\,(\pi^{\mathrm{tot}},\ell^{\mathrm{tot}})}\big(1\oplus\bar{\mathbf L}\big)=\int_{\Delta_{\{B,R\}}}x^{\boldsymbol\beta(\mathbf k)-\one}\prod_i\bar L_i(x)^{-w_i-|\mathbf k_i|}\,d\omega(x).
\]
Each moment is a read of $\phi$ against the one-scale kernels $\sZ_{-k_{i,l};(1)}\big(\bar L_{i,l}-L^{(l)}_i\big)=\big(\bar L_{i,l}-L^{(l)}_i\big)^{k_{i,l}}$.
\end{theorem}

% Reference 20: thm:pencilpeak; source line 4350
\begin{corollary}[Closing at the peak on the totals]\label{cor:pencilpeak}
In Theorem~\ref{thm:pencilpeaks} take the references $\bar L_{i,l}=\binom{c_i}{l}\bar L_i$ with constants $\bar L_i>0$ such that the deviations $h_{i,l}=\binom{c_i}{l}\bar L_i-L^{(l)}_i$ are nonnegative. The reference letter is then the constant $\bar L_i(x_B+x_R)^{c_i}=\bar L_i$ on the totals, the Breta coefficient is $\prod_i\bar L_i^{\,-w_i-|\mathbf k_i|}$ times the Beta corner $\Gst21{\beta_B(\mathbf k),\ \beta_R(\mathbf k)}{\beta_B(\mathbf k)+\beta_R(\mathbf k)}$, and that Trailmoment is
\[
\sum_{\mathbf k}\ \prod_i\frac{(w_i)_{|\mathbf k_i|}}{\mathbf k_i!}\,\bar L_i^{\,-w_i-|\mathbf k_i|}\ \Gst21{\beta_B(\mathbf k),\ \beta_R(\mathbf k)}{\beta_B(\mathbf k)+\beta_R(\mathbf k)}\int\prod_{i,l}h_{i,l}^{\,k_{i,l}}\,d\phi .
\]
For linear letters $\widehat L_i(x)=x_RL^{(0)}_i+x_BL^{(1)}_i$ with $\bar L_i\ge L^{(0)}_i,L^{(1)}_i$, with $n_i$ counting the constant coefficient and $n'_i$ the scale coefficient of the letter $i$, it is
\[
\sum_{\mathbf n,\mathbf n'}\ \prod_i\frac{(w_i)_{n_i+n'_i}}{n_i!\,n'_i!}\,\bar L_i^{\,-w_i-n_i-n'_i}\ \Gst21{m^{\mathrm{tot}}_B+|\mathbf n'|,\ m^{\mathrm{tot}}_R+|\mathbf n|}{|\mvec^{\mathrm{tot}}|+|\mathbf n|+|\mathbf n'|}\int\prod_i\big(\bar L_i-L^{(0)}_i\big)^{n_i}\big(\bar L_i-L^{(1)}_i\big)^{n'_i}\,d\phi .
\]
If a coefficient is a product $f'(u)f''(v)$ and its reference a product $\bar f'\bar f''$ of constants $\bar f'\ge f'$ and $\bar f''\ge f''$, then $\bar f'\bar f''-f'f''=\bar f'(\bar f''-f'')+f''(\bar f'-f')$, both terms nonnegative and each a function on the bunch times a function on the room; so when every coefficient, every reference and $\phi$ are products, every moment is a finite nonnegative combination of products of moments over $\Delta_B$ and over $\Delta_R$.
\end{corollary}

No order of the values is used, only bounds on the coefficients, and the terms are nonnegative. Applied to every letter on the whole simplex with the references $M_i=\max_T\theta_{j,T}/g_i!$, Theorem~\ref{thm:peak} is the draw series (Proposition~\ref{prop:jointdraw}); applied on the totals it closes the scale of the bunch alone, and the shapes of the bunch and the room enter only through the moments of the deviations of the coefficients from their references. Across a bunch of width one every coefficient is a letter of a minor on the bunch times a letter of a minor on the room (Theorem~\ref{thm:jointclosing}), so with references that are products the moments pair the bunch with the room, as the lattices of Theorem~\ref{thm:jointclosing} do, with nonnegative terms. If every $\widehat L_i$ is at least $(1-\lambda)\bar L_i$ on the totals at every shape, with $\lambda<1$, the terms of total order $\sum_i|\mathbf k_i|=k$ add up to at most $\frac{(|\mathbf w|)_k}{k!}\lambda^k$ times the term of order $0$; when a coefficient vanishes on a face of a shape simplex the terms need not decay geometrically.

For one table read at $\rho_0=1$ the two letters on the totals are linear, and so are their reference letters, one constant for each coefficient. Rescaling the scale by the ratio of the references of the letter $K$ carries its reference letter to a multiple of $x_R+x_B$, and that of $\Psi$ to a multiple of the linear letter whose scale coefficient over its constant coefficient is the reference of the even letter of the level $d$ over that of the level $r$; the rescaling multiplies the Breta coefficient by a power of the ratio of the references of $K$. Two letters on the line of the totals have two zeros besides its two endpoints, and four points of a line have one modulus, the ratio of the references of the even letters of the two levels: the Breta coefficients depend on the references only through their scales and this modulus, defined up to its inverse by the naming of the two levels. What one reference for each coefficient adds is that every deviation can vanish at its own peak, at the price of the one modulus in the coefficients. When the modulus is $1$ both reference letters are multiples of $x_R+x_B$, which is $1$ on the simplex, and Corollary~\ref{cor:pencilpeak} applies. With $N$ tables there are $2N$ letters on the totals and $2N-1$ such moduli.

% Reference 20: cor:pencilladder; source line 4420
\begin{corollary}[The ladder on the totals]\label{cor:pencilladder}
In Theorem~\ref{thm:scaletotals}(ii) let every $w_i>0$, put $w'=\sum_{i\neq i_0}w_i$, and let $\zeta>0$.
\begin{enumerate}[label=(\roman*)]
\item At shapes where every ratio $\xi_{i_0}/\xi_i$ lies in $(0,2\zeta)$,
\begin{multline*}
\int_{\Delta_{\{B,R\}}}x_B^{\kappa_B-1}x_R^{\kappa^{\mathrm{top}}_B-1}\prod_{i\neq i_0}\Big(x_R+x_B\,\frac{\xi_{i_0}}{\xi_i}\Big)^{-w_i}d\omega(x)\\
=\sum_{\mathbf n}\ \prod_{i\neq i_0}\frac{(w_i)_{n_i}}{n_i!}\Big(\zeta-\frac{\xi_{i_0}}{\xi_i}\Big)^{n_i}\ \sB^{\{\{B,R\}\}}_{(\kappa_B+|\mathbf n|,\kappa^{\mathrm{top}}_B),\,w'+|\mathbf n|;\,I_{\{B,R\}}}(\zeta,1),
\end{multline*}
absolutely: the linear letters are moved to the common ratio $\zeta$, where they merge into one, and the coefficients do not depend on the shapes. At $\zeta=1$ they are the Beta corners $\Gst21{\kappa_B+|\mathbf n|,\ \kappa^{\mathrm{top}}_B}{\kappa_B+\kappa^{\mathrm{top}}_B+|\mathbf n|}$. For one linear letter this is the ladder of the Breta in its first coefficient (Theorem~\ref{thm:bretalaws}(ii)), $\partial_\zeta\sB^{\{\{B,R\}\}}_{(\kappa_B,\kappa^{\mathrm{top}}_B),w;I_{\{B,R\}}}(\zeta,1)=-w\,\sB^{\{\{B,R\}\}}_{(\kappa_B+1,\kappa^{\mathrm{top}}_B),w+1;I_{\{B,R\}}}(\zeta,1)$, summed as a Taylor series.
\item Let $P_F(\mvec;\mathbf s,\mathbf t)$ converge and let every ratio $\xi_{i_0}/\xi_i$ lie in $[\zeta_-,\zeta_+]\subset(0,2\zeta)$ on $\Delta_B\times\Delta_R$. Then (i) integrates term by term against $\phi$ of Theorem~\ref{thm:scaletotals}(i) multiplied by $\xi_{i_0}^{\,\kappa_B}\prod_i(L^{(0)}_i)^{-w_i}$ and by the further factor of Theorem~\ref{thm:scaletotals}(ii), and the terms of $P_F(\mvec;\mathbf s,\mathbf t)$ with $|\mathbf n|=n$ add up in absolute value to at most $C\frac{(w')_n}{n!}\vartheta^nP_F(\mvec;\mathbf s,\mathbf t)$, where $\vartheta=\max(1-\zeta_-/\zeta,\ \zeta_+/\zeta-1)<1$ and $C=\sB^{\{\{B,R\}\}}_{(\kappa_B,\kappa^{\mathrm{top}}_B),w';I_{\{B,R\}}}(\zeta,1)\big/\sB^{\{\{B,R\}\}}_{(\kappa_B,\kappa^{\mathrm{top}}_B),w';I_{\{B,R\}}}(\zeta_+,1)$. If $\zeta\ge\sup\xi_{i_0}/\xi_i$ for every $i$ the terms are nonnegative and no lower bound on the ratios is needed; both sides are then finite or infinite together.
\item For one table with $i_0$ the letter of $K$, under the hypotheses of (ii) the only other ratio is $\varrho_B(u)/\varrho_R(v)$, and
\[
P_F(\mvec;s,t)=\sum_{n\ge0}\frac{(t)_n}{n!}\,\sB^{\{\{B,R\}\}}_{(\kappa_B+n,\kappa^{\mathrm{top}}_B),\,t+n;\,I_{\{B,R\}}}(\zeta,1)\int_{\Delta_B\times\Delta_R}\big(\zeta\mathbf E^r-\mathbf E^d\big)^n\,(\mathbf E^r)^{-t-n}\,\phi(du,dv),
\]
with $\mathbf E^d$, $\mathbf E^r$ and $\phi$ of Corollary~\ref{cor:pencilone}. With $i_0$ the letter of $\Psi$ the ratio is $\varrho_R(v)/\varrho_B(u)$, and the same holds with $s$ for $t$, the moments $\int(\zeta\mathbf E^d-\mathbf E^r)^n(\mathbf E^d)^{-s-n}$ and the measure $\phi=\big(\Psi_{\cB/R}\Psi_{\cB|R}\big)^{2\delta^d}\big(\Psi_{\cB|B}\Psi_{\cB/B}\big)^{2\delta^r}u^{\mvec|_B-\one}v^{\mvec|_R-\one}d\omega_B\,d\omega_R$. So one table closes on the totals with nonnegative terms about any $\zeta\ge\sup\varrho_B/\inf\varrho_R$ when no slope is positive, and with $i_0$ the letter of $\Psi$ about any point at least $\sup\varrho_R/\inf\varrho_B$ when no slope is negative (Lemma~\ref{lem:slopesrange}). With $\zeta=\zeta'\zeta''$, $\zeta'\ge\sup_{\Delta_B}\varrho_B$ and $\zeta''\ge1/\inf_{\Delta_R}\varrho_R$,
\[
\zeta\mathbf E^r-\mathbf E^d=\zeta'\mathbf E_{\cB|B}(u)\big(\zeta''\mathbf E_{\cB/B}(v)-\mathbf E_{\cB|R}(v)\big)+\mathbf E_{\cB|R}(v)\big(\zeta'\mathbf E_{\cB|B}(u)-\mathbf E_{\cB/R}(u)\big),
\]
both terms nonnegative and each a function on the bunch times a function on the room.
\item Since $x_R+x_B\xi_{i_0}/\xi_i=(\xi_{i_0}/\xi_i)\big(x_B+x_R\xi_i/\xi_{i_0}\big)$, the integral of (i) is $\prod_{i\neq i_0}(\xi_i/\xi_{i_0})^{w_i}$ times the same integral with $\kappa_B$ and $\kappa^{\mathrm{top}}_B$ swapped and every ratio replaced by its inverse. So (i) and (ii) hold in the inverse ratios, the factor $\prod_{i\neq i_0}(\xi_i/\xi_{i_0})^{w_i}$ joining $\phi$, with nonnegative terms about any common point at least $\sup\xi_i/\xi_{i_0}$: the ladder closes whenever the ratios are bounded on one side.
\end{enumerate}
\end{corollary}

The three closings of the even share read the same linear letters. At a wall the even share closes toward the level whose even letter is the smaller, which needs $\varrho_B/\varrho_R$ on one side of $1$ with room (Corollary~\ref{cor:closetwosides}(i)); the ladder moves the linear letters to a common ratio, which needs the quotient bounded on one side, and its coefficients are constants of the scale; the peak closing moves each linear letter to a constant reference, which needs the coefficients bounded and asks nothing of their ratio. In all three the terms pair moments of the bunch with moments of the room. With $\varrho_B(u)<\varrho_R(v)$ throughout and $\sup\varrho_B=\inf\varrho_R$, the ladder about $\zeta=1$ is a series of nonnegative terms whose coefficients are Beta corners, where Proposition~\ref{prop:weakorder} sums the lattices by scaling the level weights. Where the quotient takes every positive value, of the three only the peak closing applies; the gauge closing applies as well, its terms not separating over the bunch and the room (Proposition~\ref{prop:gaugeclose}).

\section{Width zero}

% Reference 20: prop:loopsimplex; source line 4485 (one table and a loop).
% Proved in the last part: prop:loopsimplex.
\begin{proposition}[A loop or a coloop on the simplex]\label{prop:loopsimplex}
Let a mark $e$ be in every table a loop ($\eta_j=-1$) or a coloop ($\eta_j=+1$), let $F'$ be the family of Corollary~\ref{cor:content} on $\cE-e$, and put $\nu_e=\mu_e-\sum_j\eta_j\delta_j$, so that $\kappa'=\kappa_{\cE}-\nu_e$ is the imbalance of $F'$ at the masses $\mvec|_{\cE-e}$. Where $\nu_e>0$ and $\kappa'>0$,
\[
P_F(\mvec;\mathbf s,\mathbf t)=\Gst21{\nu_e,\ \kappa'}{\kappa_{\cE}}\,P_{F'}(\mvec|_{\cE-e};\mathbf s,\mathbf t),
\]
the mark rescaling every other mark by the common factor $1-u_e$; it is the compound formula of Theorem~\ref{thm:offdomain} at $S=\{e\}$, with the walls $\kappa_{\{e\}}=\nu_e$ and $\kappa_{\cE\setminus e}=\kappa'$. Marks $e_1,\dots,e_k$ of this kind combine first among themselves and then with the rest, $\Gst{k+1}{1}{(\nu_{e_i})_i,\ \kappa'}{\kappa_{\cE}}=\Gst{k}{1}{(\nu_{e_i})_i}{a}\,\Gst21{a,\ \kappa'}{\kappa_{\cE}}$ with $a=\sum_i\nu_{e_i}$ and $\kappa'$ the imbalance once all of them are removed (Theorem~\ref{thm:aggregation}). On the balance wall $\kappa'=-\nu_e$ and the period of the compound diverges (Proposition~\ref{prop:compoundwall}): the Foreman tree splits molecules first (Definition~\ref{def:tree}).
\end{proposition}

The bunch $\{e\}$ has width zero: the room does not see it at all, and the period integrates it out as the Beta corner $\Gst21{\nu_e,\ \kappa'}{\kappa_{\cE}}$, whose Gamma $\Gamma(\nu_e)$ is the Gamma $\Gamma(\mu_e-\langle\boldsymbol\delta\rangle_e)$ of the port at $e$ (Theorem~\ref{thm:port}).

\section{Pooling is a move}

% Reference 20: thm:cleaningmove; source line 4502 (lists; (ii) for one list at a full lean).
% Proved in the last part: thm:poolingmove.
\begin{theorem}[Pooling is a move]\label{thm:poolingmove}
Let a bunch $B$ have width one in every table of a family of tables with positive rational weights, the weights of every table factoring across $B$ and $R$ with rational factors, at rational exponents and masses on the balance wall.
\begin{enumerate}[label=(\roman*)]
\item Theorem~\ref{thm:pooling} is one change of variables, $(\rho,u,\alpha_R)\mapsto(\rho u,\alpha_R)$, the scale and the shape of $B$: it is birational over $\mathbb Q$ with the inverse $\alpha\mapsto(|\alpha_B|,\alpha_B/|\alpha_B|,\alpha_R)$, bijective on the open domains, and it carries the face of $S\subseteq\cE$ to $\{\rho=0\}\times\{\alpha_{S\cap R}=0\}$ if $S\supseteq B$ and to $\{u_{S\cap B}=0\}\times\{\alpha_{S\cap R}=0\}$ otherwise, with the walls matched by Theorem~\ref{thm:wallsglue}: by (i) there when $S\supseteq B$ or $S\subseteq R$, by (ii) there when $S\subseteq B$, and by the gluing of the counts otherwise.
\item When the factorizations, with rational factors, give every active letter the one hinge value $\gamma$ of Theorem~\ref{thm:cornerfactor}, that theorem is a composite of moves: the change of variables (i); the change $\rho\mapsto\rho\,\gamma(u)$, birational over $\mathbb Q$ and positive on the open domain; the product of the datum of the shape with that of the cell family of the room; and the integral of the hinge, the Beta corner of a linear letter (Proposition~\ref{prop:scalebreta} with one coefficient), itself a change of variables and a product with the datum of the list on two marks whose bases are the two singletons, at the exponents $(0,2\sigma)$. Every integral in the chain converges exactly when $P_F(\mvec;\mathbf s,\mathbf t)$ does (Theorem~\ref{thm:cornerfactor}).
\end{enumerate}
So the identities of Theorems~\ref{thm:pooling} and~\ref{thm:cornerfactor} hold between period data, not only between their values: when every active letter has one hinge value the period datum of the family times that of the list on two marks is the product of the data of its two cell families.
\end{theorem}

In the interior Theorem~\ref{thm:pooling} writes the period as an integral over the shape simplex, whose faces carry the walls of Theorem~\ref{thm:wallsglue}(ii), of the periods of the pooled families $F_{B\to g}(u)$, whose walls are those of Theorem~\ref{thm:wallsglue}(i).

\section{Cleaning a molecule}

% Reference 20: def:cleaning; source line 4520
\begin{definition}[Cleaning]\label{def:cleaning}
A \emph{cleaning} of a list is a sequence of poolings of parallel and series bunches of width one, each in the list left by the previous ones, until at most two marks or no such class remain.
\end{definition}

% Reference 20: prop:poolscommute; source line 4524
\begin{proposition}[Poolings of disjoint classes commute]\label{prop:poolscommute}
Let $X$ and $Y$ be disjoint parallel or series bunches of width one of a list. Pooling $X$ keeps $Y$ a class of the same kind and width, with the same restriction to $Y$ and the same contraction of its room, and the two orders of pooling give the same list on $(\cE\setminus(X\cup Y))\sqcup\{g_X,g_Y\}$, the same cell kernels and the same hinge values. A cleaning ends after at most $|\cE|-2$ poolings.
\end{proposition}

% Reference 20: prop:cleanmaps; source line 4550 (lists).
% Proved in the last part: prop:cleanmaps.
\begin{proposition}[Cleaning commutes with maps]\label{prop:cleanmaps}
A bijection that carries one table to another (Definition~\ref{def:carry}) maps the parallel and series bunches of the first onto those of the second, with the same widths and, the weights factoring across a bunch and its room and the factorization carried along, the same hinge values. An exchange maps parallel bunches onto series bunches and swaps the hinge values by $K$ and by $\Psi$, the weight $\theta'_p$ of a pattern $p$ of a bunch $B$ becoming that of the pattern $g(B\setminus p)$ of its image. Between lists it maps the cleaning of the first list at the exponents $(s,t)$ onto the cleaning of the second at $(t,s)$.
\end{proposition}

% Reference 20: thm:cleanunique; source line 4531
\begin{theorem}[Cleaning has one result]\label{thm:cleanunique}
Let a list be one molecule. Every cleaning of it ends at the same list up to the names of its marks. If that list has more than two marks, every cleaning ends with the same marks, each the same set of marks of the original list; a list that ends with two marks may divide its marks between them in more than one way. Pooling the parallel bunches first, then the series bunches, and repeating, is one cleaning, and every other order ends at the same place.
\end{theorem}

Pooling keeps a molecule a molecule, and two poolings at one list can be continued to one list; so the end of a cleaning does not depend on the order of the poolings. The cleaned list has no parallel or series bunch of width one left, each of its marks gathers a set of marks of the molecule, and the period of the molecule is the integral, over the shape simplices of the pooled bunches against their cell kernels, of the periods of the pooled families of the cleaned list.


\chapter{Words of a tree}\label{ch:words}

A period is read back from what the Foreman procedure leaves: the bunches peeled at the nodes of a tree of bunches and the leaves at its bottom. The period of a family is the word of such a tree, an iterated fibre integral of stacked Gammas and of the periods of the leaves; opened at the atoms along a chain, or wherever no table is opened on two childcules, the word is one Gramma on the marks and the letters. The word closes in three ways: by draws, with nonnegative terms at nonnegative exponents; at the centre, through the half-sum and the difference of the two letters of a table; and node by node in one scale, where each side of a node of width one is an interval from the decay of its transform to its reach. Read on one tree, the period shows what each node adds, what the leaves contain, and how the integrations of the simplex are shared out.

\section{The two letters of a table at once}

% Reference 20: prop:pairladder; source line 4560
\begin{proposition}[The pair ladder is the expansion about the midpoint]\label{prop:pairladder}
For $a,b>0$, $s,t>0$ and $x=\frac{a-b}{a+b}\in(-1,1)$,
\[
a^{-s}b^{-t}=\Big(\frac{a+b}2\Big)^{-s-t}(1+x)^{-s}(1-x)^{-t},
\]
and the coefficients of the power series $(1+x)^{-s}(1-x)^{-t}$ are
\[
[x^k](1+x)^{-s}(1-x)^{-t}=\sum_{i+i'=k}(-1)^i\frac{(s)_i\,(t)_{i'}}{i!\,i'!}=\frac{(s+t)_k}{k!}\Gst12{s+t}{s,\ t}\int_0^1w^{s-1}(1-w)^{t-1}(1-2w)^k\,dw .
\]
So $\Gst21{s,\ t}{s+t+k}[x^k](1+x)^{-s}(1-x)^{-t}=\frac1{k!}\int_0^1w^{s-1}(1-w)^{t-1}(1-2w)^k\,dw$; swapping $s$ and $t$ multiplies the $k$-th coefficient by $(-1)^k$; and at $s=t$ only even $k$ occur, $[x^{2l}](1-x^2)^{-s}=(s)_l/l!$.
\end{proposition}

The two letters of a table enter a period as $\Lambda(z)^{-s}\Lambda(-z)^{-t}$. The pair ladder expands such a pair about its midpoint; under one power, the shape of a two-element bunch in the tower decides how much of each letter is taken (Proposition~\ref{prop:sideexchange}).

% Reference 20: prop:sideexchange and thm:junctionpattern; source lines 4572 and 4585 (one table).
% Proved in the last part: def:sideshape.
\begin{definition}[Side shape; half-sum and difference; pattern coefficients]\label{def:sideshape}
For table $j$ of a family and $w\in[0,1]$ the letter of the \emph{side shape} $w$ is $L_{j,w}(z)=w\Lambda_j(z)+(1-w)\Lambda_j(-z)$; $L_{j,1/2}(z)=\frac12\big(\Lambda_j(z)+\Lambda_j(-z)\big)$ is the \emph{half-sum} and $\partial_wL_{j,w}(z)=\Lambda_j(z)-\Lambda_j(-z)$ the \emph{difference} of the two letters of table $j$. For a bunch $B$ with room $R$ and the Taylor cube $\Lambda_j(z)=\sum_{p\subseteq B}e^{\ip{z_B}{\chi_p}}\Lambda_{j,(p)}(z_R)$ (Theorem~\ref{thm:taylorcube}), $\Lambda_{j,(p)}=0$ when $p$ is not a pattern of table $j$, the \emph{pattern coefficients} of the side shape are
\[
c_{j,p}(w;z_R)=w\,\Lambda_{j,(p)}(z_R)+(1-w)\,\Lambda^\ast_{j,(p)}(z_R),\qquad \Lambda^\ast_{j,(p)}(z_R)=\Lambda_{j,(B\setminus p)}(-z_R),
\]
so that $L_{j,w}(z)=\sum_pe^{\ip{z_B}{\chi_p}}c_{j,p}(w;z_R)$. For one table the index $j$ is omitted.
\end{definition}

% Reference 20: prop:sideexchange; source line 4572 (one table).
% Proved in the last part: prop:sideexchange.
\begin{proposition}[The side shape, and what an exchange does to it]\label{prop:sideexchange}
Let a bijection $g$ carry or exchange a family $F$, with the bijection $h$ of its tables (Definition~\ref{def:carry}). If $g$ carries $F$, then $L_{h(j),w}(z)=L_{j,w}(g^\ast z)$; if it exchanges $F$, then $L_{h(j),w}(z)=L_{j,1-w}(g^\ast z)$ (Definition~\ref{def:sideshape}). For $s_j,t_j>0$, by the tower at the unit configuration on a two-element set $J$ (Theorem~\ref{thm:tower}),
\begin{align*}
\Lambda_j(z)^{-s_j}\Lambda_j(-z)^{-t_j}&=\Gst12{s_j+t_j}{s_j,\ t_j}\,\sB^{\{J\}}_{(s_j,t_j),\,s_j+t_j;\,I_J}\big(\Lambda_j(z),\Lambda_j(-z)\big)\\
&=\Gst12{s_j+t_j}{s_j,\ t_j}\int_0^1w^{s_j-1}(1-w)^{t_j-1}\,\sZ_{s_j+t_j;(1)}\big(L_{j,w}(z)\big)\,dw ,
\end{align*}
and for one table an exchange together with $w\mapsto1-w$ keeps $\sZ_{s+t;(1)}(L_w(z))$, keeping the factor $w^{s-1}(1-w)^{t-1}$ exactly when $s=t$.
\end{proposition}

The side shape weighs the basis letter against the cobasis letter, and an exchange of the table turns it over. The half-sum is even in $z$ and the difference odd: $\Lambda(z)\Lambda(-z)=L_{1/2}(z)^2-\frac14\big(\partial_wL_w(z)\big)^2$.

\section{The junction in pattern form}

% Reference 20: thm:junctionpattern; source line 4585 (one table).
% Proved in the last part: thm:junctionpattern.
\begin{theorem}[The junction in pattern form]\label{thm:junctionpattern}
Let $F=((\cB_j,\theta_j);s_j,t_j)_{j=1}^N$ be a molecule on $n$ marks with real exponents $s_j,t_j\ge0$, its masses on the balance wall with every proper wall positive, and let $B$ be a proper bunch with room $R$ and $g\in R$. For a table at a full lean put $w_j=1$ when $t_j=0$ and $w_j=0$ when $s_j=0$. With the pattern coefficients $c_{j,p}(w;z_R)$ of the side shapes (Definition~\ref{def:sideshape}),
\begin{multline*}
P_F(\mvec;\mathbf s,\mathbf t)=\frac{2^{n-1}}{\prod_{j:\,s_j>0}\Gamma(s_j)\prod_{j:\,t_j>0}\Gamma(t_j)}\int\prod_{j:\,s_j,t_j>0}w_j^{s_j-1}(1-w_j)^{t_j-1}\int_{\mathbb R^{R-g}}e^{2\ip{\bmu|_R}{z_R}}\\
\times\int\prod_{j\text{ active}}y_j^{s_j+t_j-1}\,\sG_{2\bmu|_B;(\chi_p)_p}\Big(\Big(\sum_{j\text{ active}}y_j\,c_{j,p}(w_j;z_R)\Big)_p\Big)\,d\mathbf y\,dz_{R-g}\Big|_{z_g=0}\,d\mathbf w ,
\end{multline*}
with $d\mathbf y=\prod_jdy_j$ over the active tables, $y_j\in(0,\infty)$, and $d\mathbf w=\prod_jdw_j$ over the tables with $s_j,t_j>0$, $w_j\in[0,1]$; the Gramma on the scales $e^{z_B}$ with one term for each $p\subseteq B$ whose coefficient $\sum_jy_jc_{j,p}(w_j;z_R)$ is not $0$. A table at a full lean has no integral over its side shape: at $(s_j,0)$ its coefficients are $y_j\Lambda_{j,(p)}(z_R)$, at $(0,t_j)$ they are $y_j\Lambda^\ast_{j,(p)}(z_R)$.
\end{theorem}

The two letters of every table become one letter, the side shapes are integrated last, and the bunch is integrated on its own orthant against the patterns of all the tables at once: each pattern is one term of one Gramma on the bunch, whose coefficient adds the pattern coefficients of the tables, each at its own scale $y_j$, and the room sits in the coefficients. A map that keeps the bunch moves the coefficients with the side shapes and the tables.

% Reference 20: prop:mapcoeff; source line 4597 (one table).
% Proved in the last part: prop:mapcoeff.
\begin{proposition}[What a map keeping the bunch gives the coefficients]\label{prop:mapcoeff}
In Theorem~\ref{thm:junctionpattern} let a bijection $\tilde g$ of $\cE$ keep $B$, with $\tilde g_R$ its restriction to the room, and carry or exchange $F$ with the bijection $h$ of its tables. If $\tilde g$ carries $F$, then $c_{h(j),\tilde g(p)}(w;z_R)=c_{j,p}(w;\tilde g_R^\ast z_R)$ for every $j$, $p$ and $w$; if it exchanges $F$, then $c_{h(j),\tilde g(p)}(1-w;z_R)=c_{j,p}(w;\tilde g_R^\ast z_R)$.
\end{proposition}

% Reference 20: prop:levelformQ; source line 4604 (one table).
% Proved in the last part: prop:levelformQ.
\begin{proposition}[Level form of the half-sum]\label{prop:levelformQ}
In Theorem~\ref{thm:junctionpattern} the half-sum and the difference of table $j$ have the pattern coefficients $c_{j,p}(\frac12;z_R)=\frac12\big(\Lambda_{j,(p)}+\Lambda^\ast_{j,(p)}\big)(z_R)$ and $\partial_wc_{j,p}(w;z_R)=\big(\Lambda_{j,(p)}-\Lambda^\ast_{j,(p)}\big)(z_R)$.
\begin{enumerate}[label=(\roman*)]
\item For every table, $c_{j,B\setminus p}(\frac12;-z_R)=c_{j,p}(\frac12;z_R)$ and $\partial_wc_{j,B\setminus p}(w;-z_R)=-\partial_wc_{j,p}(w;z_R)$.
\item The half-sum of table $j$ has level form on $B$ (Proposition~\ref{prop:levelform}), $c_{j,p}(\frac12;\cdot)$ depending on $|p|$ alone, exactly when every transposition of two marks of $B$ fixing the rest carries the multiset $\cB_j\uplus\{\cE\setminus T:\ T\in\cB_j\}$, with its weights, onto itself; if $2r_j\neq n$ this is the condition that every such transposition carries table $j$ (Proposition~\ref{prop:levelform}).
\item If every such transposition carries or exchanges table $j$, its half-sum has level form and the square of its difference is unchanged by every permutation of $B$.
\end{enumerate}
\end{proposition}

\section{The word of a tree}

% Reference 20: the trees of cor:treeword and thm:drawtree; source lines 4625 and 4703
\begin{definition}[Tree of bunches]\label{def:bunchtree}
A \emph{tree of bunches} of a family $F$ on $\cE$ has \emph{vertices} $v$, each with marks $M_v\subseteq\cE$ and tables on $M_v$ with their exponents, the root having $M_v=\cE$ and the active tables of $F$. A vertex whose tables form one molecule is a \emph{node}. A node $v$ that is not a leaf carries a proper bunch $B_v\subsetneq M_v$; the tables at $v$ have their children at $B_v$, indexed by the table and the pattern as in the record of the peel (Definition~\ref{def:record}) and not merged, the \emph{tables made at} $v$, each carrying the letters its parent carries, its exponent in each letter a positive multiple of the parent's, as at every native split (Definition~\ref{def:split}). The molecules $R_{v,1},\dots,R_{v,q_v}$ of the tables made at $v$ on $M_v\setminus B_v$, in a fixed order, are the childcules of $v$ (Definition~\ref{def:peel}), each a node whose tables are the restrictions to it of the tables made at $v$. The \emph{room} of $v$ is the one childcule when $q_v=1$; when $q_v\ge2$ it is a \emph{split vertex}, with the marks $M_v\setminus B_v$ and the tables made at $v$, whose children are the childcules. The root of a family whose active tables form two or more molecules is likewise a split vertex, whose children are the molecules with the restrictions of its tables. A leaf $v$ has $B_v=M_v$. The \emph{width} of a node is the largest width $c(B_v)$ of its bunch in the tables at $v$, zero at a leaf. Merging equal children with their exponents added, as the peel does (Definition~\ref{def:peel}), keeps every letter when the exponents are nonnegative, and can cancel a letter when exponents of both signs meet; the tables made at a node are therefore indexed. The tables at a node are thus the children of the tables of $F$ at the patterns shown on the bunches above it, restricted to its marks (Proposition~\ref{prop:twosteps}); with a native gauge they carry as exponents the shares of the native splits at the nodes above (Theorem~\ref{thm:hered}). Its descents are the sequences of nodes from the root to a leaf, each a childcule of the one before, as on the Foreman tree (Definition~\ref{def:keeppot}), which is a tree of bunches (Definition~\ref{def:tree}).
\end{definition}

% Reference 20: cor:treechain; source line 4618
\begin{corollary}[The tree of a linked molecule is a chain]\label{cor:treechain}
Let some active table of a molecule $F$ be connected, as for a molecule of one table. Every tree of bunches of $F$ is then a chain $M_0\supset M_1\supset\cdots\supset M_d$ with $M_{k+1}=M_k\setminus B_k$, above one leaf: the tables at $M_k$ are the children of the tables of $F$ at $S_k=B_0\sqcup\cdots\sqcup B_{k-1}$ (Proposition~\ref{prop:twosteps}), which form one molecule (Lemma~\ref{lem:conchild}). With a native gauge their exponents are the shares of the native split, and on the Foreman tree the family at $M_k$ is the peel $F_{S_k}(\mathbf a^\star,\mathbf b^\star)$ (Definition~\ref{def:peel}).
\end{corollary}

% Reference 20: cor:leafchain; source line 2674
\begin{corollary}[A linked molecule]\label{cor:leafchain}
Let some active table of $F$ be connected, as for a molecule of one table. Every peel then leaves one childcule, so the tree is a chain (Corollary~\ref{cor:treechain}), there are no twins, and the room of each node is the next node. For every chain that peels a class at each node, every class of a node other than its bunch lies inside one class of the next node, and the leaf is a union of whole classes of every node above it (Theorems~\ref{thm:roomkeeps}(ii) and~\ref{thm:leafbounds}(i)). For the Foreman tree, which is then a chain, at every realization, the centre and every lean, Theorem~\ref{thm:leafbounds}(ii) with no twins gives
\[
|B_0|\le|B_1|\le\cdots\le|B_{d-1}|\le|M_d|:
\]
the leaf alone has at least as many marks as every bunch.
\end{corollary}

% Reference 20: cor:treeword; source line 4625
\begin{corollary}[The word of a tree]\label{cor:treeword}
Let $F$ have real data with $s_j,t_j\ge0$, a native gauge and its masses on the balance wall, and let the children's weights factor across the childcules at every node of a tree of bunches. Run the port of Theorem~\ref{thm:wallsplit}(ii) at every node, on the native splits. A table $\cB'$ made at a node, with the exponents $(a_{\cB'},b_{\cB'})$, is the direct sum of its restrictions to the childcules there (Proposition~\ref{prop:directsum}); let $q_{\cB'}$ be the number of those childcules, other than leaves, on which its letter is not a monomial. Then
\[
P_F(\mvec;\mathbf s,\mathbf t)=\prod_j\frac1{\Gamma(s_j)\Gamma(t_j)}\int_{\mathcal X}\ \prod_{\cB'}\big(\Gamma(a_{\cB'})\Gamma(b_{\cB'})\big)^{1-q_{\cB'}}\ \prod_{\text{leaves }v}P_{F_v}(\mvec|_{M_v};\mathbf a_v,\mathbf b_v),
\]
the \emph{word of the tree}: an iterated fibre integral through the native splits, over the fibre $\mathcal X$ cut out by the share sums and by the $|B_v|+q_v-1$ pins of every node $v$ that is not a leaf (Theorem~\ref{thm:wallsplit}(ii)), a monomial restriction entering the pins of its node as a load; $F_v$ is the family of the tables at the leaf $v$ with their exponents $(\mathbf a_v,\mathbf b_v)$, and a letter with exponent zero carries no shares and no Gamma. A table opened on one childcule carries no factor, and when an active table is connected every $q_{\cB'}$ is $0$ or $1$ (Corollary~\ref{cor:treechain}).
\end{corollary}

% Reference 20: thm:buildup; source line 4636 (with one table at the centre in the last sentence of (iv)).
% Proved in the last part: thm:buildup.
\begin{theorem}[Building back up]\label{thm:buildup}
In Corollary~\ref{cor:treeword} open every leaf by the atom formula on the balance wall (Corollary~\ref{cor:atomwall}), a unit being a monomial, and write the shares $(a,b)$ of one pattern of one table in pair coordinates, $\sigma=\frac{a+b}2$ and $\delta=\frac{a-b}2$.
\begin{enumerate}[label=(\roman*)]
\item \emph{Polycules multiply.} At every node and at every point of its fibre on the wall contour of the native split (Theorem~\ref{thm:wallsplit}(iii)), the integral over the shares below the node is the product over its childcules of their words: the polycule of the childcules is evaluated as the product of its molecules, a unit contributing a monomial in the weights.
\item \emph{The pins are odd.} If every table has $s_j,t_j>0$, every pin other than the share sums involves the odd shares alone: the pins $\langle\boldsymbol\delta\rangle_e=\mu_e$ of the bunch, the pins of the childcules,
\[
\kappa_{R_i}=\mu_{R_i}-\sum_{j,p}\delta_{j,p}\big(2r_{j,(p)}(R_i)-|R_i|\big)=0,
\]
the sum over the tables made at the node, indexed as in Theorem~\ref{thm:wallsplit} by the table $j$ at the node and the pattern $p$, and the pins of the atoms. The share sums are $\sum_p\sigma_{j,p}=\sigma_j$ and $\sum_p\delta_{j,p}=\delta_j$, so the fibre is the product of an even fibre, cut by the even sums alone, and an odd fibre.
\item \emph{One Gramma.} If at every node every table made there is a monomial on all its childcules but at most one, in particular along a chain (Corollary~\ref{cor:treechain}), the leaf shares are one for each basis $T\in\cB_j$ and letter, the pins of the tree are, for a fixed mark $g$, the pins $\langle\boldsymbol\delta\rangle_e=\mu_e$, $e\in\cE-g$, of Corollary~\ref{cor:atomwall} after a change of determinant $\pm1$, and the word is the atom formula, one Gramma on the marks and the letters:
\[
\prod_{j:\,s_j>0}\Gamma(s_j)\prod_{j:\,t_j>0}\Gamma(t_j)\,P_F(\mvec;\mathbf s,\mathbf t)=\sG_{(\bmu|_{\cE-g};\mathbf s;\mathbf t);\,t^{(g)}_F}(\boldsymbol\theta\oplus\boldsymbol\theta),
\]
with $t^{(g)}_F$ the configuration $t_F$ of Theorem~\ref{thm:confJ} without its unit columns and without the row of $g$: on the scales $(\cE-g)\sqcup\{v_j\}\sqcup\{w_j\}$, the columns $(\frac12\chi_T|_{\cE-g};\mathbf e_j;0)$ and $(-\frac12\chi_T|_{\cE-g};0;\mathbf e_j)$, $T\in\cB_j$, each with the coefficient $\theta_{j,T}$, a letter with exponent zero omitted: one scale for each mark but $g$, which carries an odd pin, and one for each letter.
\item \emph{The even shares close jointly.} In (iii) let every table have $s_j,t_j>0$, and write $c_{j,T}=b_{j,T}-a_{j,T}=-2\delta_{j,T}$. At fixed differences the $K$-shares of table $j$ close into one scale,
\[
\int_{\sum_Ta_{j,T}=s_j}\ \prod_{T\in\cB_j}\Gamma(a_{j,T})\,\Gamma(a_{j,T}+c_{j,T})\,\theta_{j,T}^{-2a_{j,T}}=\int_0^\infty y^{s_j-1}\prod_{T\in\cB_j}\sG_{c_{j,T};(1,-1)}\big(1,\theta_{j,T}^2y\big)\,dy ,
\]
and $\prod_j\Gamma(s_j)\Gamma(t_j)P_F(\mvec;\mathbf s,\mathbf t)$ is the fibre integral over the differences, with the pins $\sum_Tc_{j,T}=t_j-s_j$ and $-\frac12\sum_{j,T}c_{j,T}\chi_T(e)=\mu_e$, $e\in\cE-g$, of $\prod_j\prod_T\theta_{j,T}^{-c_{j,T}}$ times these scale integrals; on the wall contour of the native split (Theorem~\ref{thm:wallsplit}(iii)) every integral converges absolutely. The fibre of differences has dimension $\sum_j(|\cB_j|-1)-(n-1)$. At the centre, every table riding the line, $s_j=t_j$, with the masses $\sigma\one$, the pins are $\sum_{T\in\cB_j}c_{j,T}=0$ for every $j$ and $\sum_jM_j^{\mathsf T}\mathbf c_j=0$, with $\mathbf c_j=(c_{j,T})_{T\in\cB_j}$ and the incidence matrices $(M_j)_{Te}=\one_T(e)$, and the contour is $\re\mathbf c=0$.
\end{enumerate}
\end{theorem}

Built back up, the word is the period read from its atoms. At every node the childcules multiply and the pins see only the odd shares; along a chain nothing is opened twice, and the word is the atom formula itself, one Gramma on the marks and the letters whose configuration is that of the interest with the masses' own columns removed. The even shares are free but for their sums, and they close table by table into one scale of Grammas $\sG_{\nu;(1,-1)}$, as at the centre of a molecule (Theorem~\ref{thm:buildup}(iv)) and for the table that is one circuit (Corollary~\ref{cor:cycle}).

\section{Closing the tree by draws}

% Reference 20: prop:jointdraw; source line 4681 (with cor:drawseries, source line 4670; the worked reading for one table).
% Proved in the last part: prop:jointdraw.
\begin{proposition}[The draw series]\label{prop:jointdraw}
Let $U=(U_Q)_{Q\in\cP}$ have independent blocks $U_Q\sim\Dir_{\mvec|_Q}$ on a partition $\cP$ of the marks, and let $P_1,\dots,P_q$ be nonzero polynomials with nonnegative coefficients, $P_i$ homogeneous of degree $g_{i,Q}$ in the variables of each block $Q$. For an exponent vector $\mathbf a$ with block sums $g_{i,Q}$ put $b_{\mathbf a}=\prod_Qg_{i,Q}!/\prod_ea_e!$ and $C_i=\max_{\mathbf a}\big([u^{\mathbf a}]P_i\big)/b_{\mathbf a}$. Then $\prod_Q\big(\sum_{e\in Q}u_e\big)^{g_{i,Q}}-P_i/C_i$ has nonnegative coefficients and equals $1-P_i/C_i\in[0,1)$ on the interior of the source, and for $\nu_i\ge0$
\begin{multline*}
\oE_{U\sim\bigotimes_Q\Dir_{\mvec|_Q}}\Big[\prod_iP_i(U)^{-\nu_i}\Big]=\prod_iC_i^{-\nu_i}\sum_{\mathbf k\in\mathbb N^q}\prod_i\frac{(\nu_i)_{k_i}}{k_i!}\\
\times\oE_{U\sim\bigotimes_Q\Dir_{\mvec|_Q}}\Big[\prod_i\Big(1-\frac{P_i(U)}{C_i}\Big)^{k_i}\Big],
\end{multline*}
every term nonnegative, both sides finite or infinite together. The left side is the Trailmoment of Proposition~\ref{prop:jointmoments},
\[
\Big(\sT_{\boldsymbol\nu;I_q}[(P_i)_i]\,\bigotimes_Q\Dir_{\mvec|_Q}\Big)\Big(\prod_Q\Delta_Q\Big),
\]
and each expectation on the right is a finite nonnegative combination of the moments of Theorem~\ref{thm:dirmoments}. For a family $F=((\cB_j,\theta_j);s_j,t_j)_{j=1}^N$ with $s_j,t_j\ge0$, one block, $P_{2j-1}=K_j$ and $P_{2j}=\Psi_j$ with the indices $(s_j,t_j)$, the constants are $C_{2j-1}=\theta_{j,\max}/r_j!$ and $C_{2j}=\theta_{j,\max}/(n-r_j)!$, $\theta_{j,\max}=\max_T\theta_{j,T}$; the left side times $\Gst{n}{1}{\mvec}{m_{\cE}}$ is $\int_{\Delta_{\cE}}u^{\mvec-\one}\prod_jK_j(u)^{-s_j}\Psi_j(u)^{-t_j}\,d\omega_{\cE}(u)$, which on the balance wall is $P_F(\mvec;\mathbf s,\mathbf t)$; and the expectation of the term $\mathbf k=(k_j,l_j)_j$ is the chance that, in the urn that draws a mark with probability proportional to its current mass and then adds one to that mass, for $j=1,\dots,N$ in turn $k_j$ rounds of $r_j$ draws and then $l_j$ rounds of $n-r_j$ draws all fail for table $j$, a round failing when its marks are not a basis, respectively not a cobasis, of table $j$, and, when they are the basis $T$, respectively the cobasis $\cE\setminus T$, with probability $1-\theta_{j,T}/\theta_{j,\max}$: a finite nonnegative combination of Dirichlet moments, a finite sum of them for lists. This series is Theorem~\ref{thm:peak}(i) with the references $M_{2j-1}=C_{2j-1}$ and $M_{2j}=C_{2j}$.
\end{proposition}

% Reference 20: prop:cremonatrail; source line 4815 (one table).
% Proved in the last part: cor:cremonatrail.
\begin{corollary}[At $\delta=0$ and at a full lean the period reads the law of one letter for each table]\label{cor:cremonatrail}
Let a family $F=((\cB_j,\theta_j);s_j,t_j)_{j=1}^N$, every table active, carry its masses $\mvec$ on the balance wall with every proper wall positive, every table riding the line, $(s_j,t_j)=(\sigma_j,\sigma_j)$, or at a full lean, $(s_j,0)$ or $(0,t_j)$. Put $L_j=\mathbf E_j=K_j\Psi_j$ and $\nu_j=\sigma_j$ when table $j$ rides the line, $L_j=K_j$ and $\nu_j=s_j$ at $(s_j,0)$, and $L_j=\Psi_j$ and $\nu_j=t_j$ at $(0,t_j)$, and let $\phi_{\mvec}$ be the joint law of $(L_1(U),\dots,L_N(U))$ for $U\sim\Dir_{\mvec}$ (Proposition~\ref{prop:gammadraws}). The kernel $\prod_ja_j^{-s_j}b_j^{-t_j}$ of the letters $a_j=K_j$ and $b_j=\Psi_j$ is the kernel $\sZ_{\boldsymbol\nu;I_N}\big((L_j)_j\big)$, one scale for each table, so
\[
P_F(\mvec;\mathbf s,\mathbf t)=\Gst{n}{1}{\mvec}{m_{\cE}}\,\big(\sT_{\boldsymbol\nu;I_N}[\mathrm{id}]\,\phi_{\mvec}\big)\big((0,\infty)^N\big);
\]
with one table $\phi_{\mvec}$ is a law on the half-line and the right side is $\Gst{n}{1}{\mvec}{m_{\cE}}\,\Tx{\nu}{\phi_{\mvec}}{0}$, at $\delta=0$ the law $\mathbf E_\ast\Dir_{\mvec}$ of the even letter $\mathbf E(U)=K(U)\Psi(U)$ (Theorem~\ref{thm:transportfree}). The law $\phi_{\mvec}$ is carried by $\prod_j(0,\max_{\Delta_{\cE}}L_j]$, it is determined by its moments, and its moments are finite sums of Dirichlet moments (Theorem~\ref{thm:dirmoments}).
\end{corollary}

% Reference 20: thm:drawchain; source line 4692 (one list).
% Proved in the last part: thm:drawchain.
\begin{theorem}[The draw series along a chain]\label{thm:drawchain}
Let a family $F=((\cB_j,\theta_j);s_j,t_j)_{j=1}^N$ have the exponents $s_j,t_j\ge0$ and masses $\mvec\in\mathbb R^{\cE}_{>0}$, and let $M_0=\cE\supset M_1\supset\cdots\supset M_d$ be a chain of sets with $B_i=M_i\setminus M_{i+1}$, $i<d$, and $S_i=B_0\sqcup\cdots\sqcup B_{i-1}$, $S_0=\varnothing$, $S_{d+1}=\cE$. A round of draws \emph{extends on} $S$ for table $j$ when its draws inside $S$ are distinct and are $T\cap S$ for some basis $T$ of table $j$. In this statement write $\theta_j(S;P)$ for the largest weight $\theta_{j,T}$ over the bases $T$ of table $j$ with $T\cap S=P$, and say that a round that extends on $S_i$ and on $S_{i+1}$, with the traces $P_i$ and $P_{i+1}$, is \emph{accepted at} $S_{i+1}$ with probability $\theta_j(S_{i+1};P_{i+1})/\theta_j(S_i;P_i)$, independently for each $i$. Let $\mathfrak f_{j,i}(u)$, $0\le i\le d$, be the chance that a round of $r_j$ draws from $u$ extends on $S_i$ and is accepted at $S_1,\dots,S_i$ but does not both extend on $S_{i+1}$ and get accepted there, and $\mathfrak f^\ast_{j,i}$ the same for rounds of $n-r_j$ draws and the cobases $\cE\setminus T$ with the weights $\theta_{j,T}$. Then, with $\theta_{j,\max}$ the largest weight of table $j$, $1-r_j!\,K_j/\theta_{j,\max}=\sum_i\mathfrak f_{j,i}$, $1-(n-r_j)!\,\Psi_j/\theta_{j,\max}=\sum_i\mathfrak f^\ast_{j,i}$, and
\begin{multline*}
\int_{\Delta_{\cE}}u^{\mvec-\one}\prod_jK_j^{-s_j}\Psi_j^{-t_j}\,d\omega_{\cE}=\Gst{n}{1}{\mvec}{m_{\cE}}\prod_j\Big(\frac{r_j!}{\theta_{j,\max}}\Big)^{s_j}\Big(\frac{(n-r_j)!}{\theta_{j,\max}}\Big)^{t_j}\\
\times\sum_{\mathbf k,\mathbf l}\ \prod_j\frac{(s_j)_{|\mathbf k_j|}\,(t_j)_{|\mathbf l_j|}}{\mathbf k_j!\,\mathbf l_j!}\ \oE_{U\sim\Dir_{\mvec}}\Big[\prod_{j,i}\mathfrak f_{j,i}(U)^{k_{j,i}}\,\mathfrak f^\ast_{j,i}(U)^{l_{j,i}}\Big],
\end{multline*}
over $\mathbf k_j,\mathbf l_j\in\mathbb N^{d+1}$, every term a finite nonnegative combination of Dirichlet moments, a finite sum of them for lists, whose rounds are always accepted. The numbers $\sum_{i'\ge i}k_{j,i'}$ of failing basis rounds of table $j$ that pass $B_0,\dots,B_{i-1}$ decrease with $i$: they are nested, as the coordinates of the region $X_N$ of a trail are (Definition~\ref{def:trail}). Let a bijection $g$ keep each $B_i$ and the masses, let $h$ be a bijection of the tables, and put $(g_\ast u)_{g(e)}=u_e$. If $g$ maps the bases of every table $j$ onto those of table $h(j)$ with their weights, then $\mathfrak f_{h(j),i}(g_\ast u)=\mathfrak f_{j,i}(u)$ and $\mathfrak f^\ast_{h(j),i}(g_\ast u)=\mathfrak f^\ast_{j,i}(u)$; if $T\mapsto\cE\setminus g(T)$ does, then $\mathfrak f^\ast_{h(j),i}(g_\ast u)=\mathfrak f_{j,i}(u)$ and $\mathfrak f_{h(j),i}(g_\ast u)=\mathfrak f^\ast_{j,i}(u)$, and the array of expectations is unchanged when every $(\mathbf k_j,\mathbf l_j)$ is replaced by $(\mathbf l_{h^{-1}(j)},\mathbf k_{h^{-1}(j)})$; for one table, where an exchange forces $2r=n$, this is the symmetry under $\mathbf k\leftrightarrow\mathbf l$.
\end{theorem}

% Reference 20: thm:drawtree; source line 4703 (a molecule of lists).
% Proved in the last part: thm:drawtree.
\begin{theorem}[The draw series along a tree]\label{thm:drawtree}
Let a molecule $F=((\cB_j,\theta_j);s_j,t_j)_j$ on $\cE$ carry $s_j,t_j\ge0$ and masses $\mvec\in\mathbb R^{\cE}_{>0}$, and take a tree of bunches (Definition~\ref{def:bunchtree}) at every node of which the children's weights factor across the childcules, as they do for lists: for a table $\cB'$ at a node $v$ and a pattern $p$ of $\cB'$ at $B_v$, the restrictions $\cB'_{(p)}|R_{v,i}$ of the child to the childcules (Proposition~\ref{prop:directsum}) carry weights whose products are the weights of the child, $\theta_{(p),X}=\prod_i\theta_{(p),i}(X\cap R_{v,i})$. Write $\theta_{\max}(\cB')$ for the largest weight of a table $\cB'$, and say that a round of draws that is a basis $T$ of $\cB'$ is \emph{accepted} with probability $\theta_T/\theta_{\max}(\cB')$. A cobasis round of table $j$ is a basis round of its dual table, whose children are the duals of the children, with the same weights (Lemma~\ref{lem:descent}(ii) with the identity). The \emph{coordinates of the tree} are, at each node $v$, the totals $x^{(v)}\in\Delta_{\{B_v,R_{v,1},\dots,R_{v,q_v}\}}$ of the shape of $M_v$ on its blocks and the shape $u^{(v)}\in\Delta_{B_v}$ of the bunch; under $\Dir_{\mvec}$ they are independent, with the laws $\Dir_{(m_{B_v},m_{R_{v,1}},\dots,m_{R_{v,q_v}})}$ and $\Dir_{\mvec|_{B_v}}$ (Theorem~\ref{thm:aggregation}). For a table $\cB'$ at a node write its chances as functions of the coordinates of the subtree of that node.
\begin{enumerate}[label=(\roman*)]
\item \emph{A round passes a node, then each childcule.} For a table $\cB'$ of rank $r'$ at a node $v$, the chance $\mathfrak s_{\cB'}=r'!\,K_{\cB'}/\theta_{\max}(\cB')$ that $r'$ draws from the shape of $M_v$ are a basis of $\cB'$ and are accepted is
\begin{gather*}
\mathfrak s_{\cB'}=\sum_p\mathfrak p_{\cB',p}\prod_i\mathfrak s_{\cB'_{(p)}|R_{v,i}},\\
\mathfrak p_{\cB',p}=\frac{\prod_i\theta_{\max}\big(\cB'_{(p)}|R_{v,i}\big)}{\theta_{\max}(\cB')}\,\frac{r'!}{\prod_ir'_{p,i}!}\,\big(x^{(v)}_{B_v}\big)^{|p|}\big(u^{(v)}\big)^{p}\prod_i\big(x^{(v)}_{R_{v,i}}\big)^{r'_{p,i}},
\end{gather*}
over the patterns $p$ of $\cB'$ at $B_v$, with $r'_{p,i}$ the rank on $R_{v,i}$ of the child $\cB'_{(p)}$: $\mathfrak p_{\cB',p}$ is the chance that the round \emph{passes} $v$ showing $p$, its draws in $B_v$ being the marks of $p$ and its numbers of draws in the childcules the $r'_{p,i}$, and is accepted there, with the probability given by its first factor, which is at most $1$; and $\mathfrak s_{\cB'_{(p)}|R_{v,i}}$ is the chance that its draws in $R_{v,i}$ are a basis of the restriction of the child there and are accepted. At a leaf $\mathfrak s_{\cB'}=r'!\,K_{\cB'}/\theta_{\max}(\cB')$.
\item \emph{A failure is met at one node.} The chance $1-\mathfrak s_{\cB'}$ that the round fails, by not being a basis or by not being accepted, is $\sum_w\mathfrak f_{\cB'\to w}$ over the nodes and leaves $w$ of the subtree of $v$, where $\mathfrak f_{\cB'\to v}=1-\sum_p\mathfrak p_{\cB',p}$ and, for $w$ in the subtree of $R_{v,i}$,
\[
\mathfrak f_{\cB'\to w}=\sum_p\mathfrak p_{\cB',p}\prod_{i'<i}\mathfrak s_{\cB'_{(p)}|R_{v,i'}}\ \mathfrak f_{\cB'_{(p)}|R_{v,i}\to w},
\]
the chance that the round first fails at $w$. Each $\mathfrak f_{\cB'\to w}$ is a finite sum of products, over the blocks of the tree, of polynomials with nonnegative coefficients in one block each.
\item \emph{The series.} With $\mathfrak f_{j,w}$ and $\mathfrak f^\ast_{j,w}$ these chances for the basis and the cobasis rounds of table $j$ of $F$ from the root, and $\theta_{j,\max}$ the largest weight of table $j$,
\begin{multline*}
\int_{\Delta_{\cE}}u^{\mvec-\one}\prod_jK_j^{-s_j}\Psi_j^{-t_j}\,d\omega_{\cE}=\Gst{n}{1}{\mvec}{m_{\cE}}\prod_j\Big(\frac{r_j!}{\theta_{j,\max}}\Big)^{s_j}\Big(\frac{(n-r_j)!}{\theta_{j,\max}}\Big)^{t_j}\\
\times\sum_{\mathbf k,\mathbf l}\ \prod_j\frac{(s_j)_{|\mathbf k_j|}\,(t_j)_{|\mathbf l_j|}}{\mathbf k_j!\,\mathbf l_j!}\ \oE_{U\sim\Dir_{\mvec}}\Big[\prod_{j,w}\mathfrak f_{j,w}(U)^{k_{j,w}}\,\mathfrak f^\ast_{j,w}(U)^{l_{j,w}}\Big],
\end{multline*}
every term nonnegative, both sides finite or infinite together; on the balance wall the left side is $P_F(\mvec;\mathbf s,\mathbf t)$. Each expectation is a finite nonnegative combination of products over the blocks of Dirichlet moments, a finite sum of them for lists: at every node, the node's own moments times the product over its childcules of theirs.
\item \emph{Nesting along the tree.} The numbers $k_{j,\ge w}=\sum_{w'}k_{j,w'}$, over the $w'$ in the subtree of $w$, of the basis rounds of table $j$ that first fail in the subtree of $w$, satisfy $k_{j,\ge w}=k_{j,w}+\sum_{w'}k_{j,\ge w'}$ over the childcules $w'$ of $w$, and likewise for the cobasis rounds: along every descent they are nested, as the coordinates of the region $X_N$ of a trail are (Definition~\ref{def:trail}). On a chain this is Theorem~\ref{thm:drawchain}.
\end{enumerate}
\end{theorem}

At nonnegative exponents every share of the word closes at once. The draw series closes each letter about its largest coefficient on the simplex, and the tree only sorts its failures: a round that fails is charged to the first node at which it stops passing, the chances of passing are monomials in the coordinates of that node, and the coordinates of different nodes are independent. So every term is a nonnegative combination of products, over the blocks of the tree, of Dirichlet moments, and the counts of failures nest along the descents as the coordinates of the region $X_N$ of a trail do.

% Reference 20: cor:chaincoords; source line 2685
\begin{corollary}[What each node of a chain integrates]\label{cor:chaincoords}
Let the tree of a molecule on $n$ marks be a chain $M_0\supset M_1\supset\cdots\supset M_d$ with bunches $B_k$ and rooms $R_k=M_{k+1}$. The simplex of the molecule is the product
\[
\Delta_{M_0}\cong\prod_{k=0}^{d-1}\Big(\Delta_{B_k}\times\Delta_{\{B_k,R_k\}}\Big)\times\Delta_{M_d}
\]
of the shape of each bunch, the totals of the bunch and its room at each node, and the shape of the leaf, and under $\Dir_{\mvec}$ these coordinates are independent: they are the coordinates of the tree of Theorem~\ref{thm:drawtree}, which use no weights. The node of $B_k$ has $|B_k|$ coordinates, $|B_k|-1$ of shape and one total, the leaf has $|M_d|-1$, and they add up to $n-1$. When some active table is connected, the Foreman tree, which is then a chain, has, at the centre and at every lean, no node with more coordinates than the leaf has marks (Corollary~\ref{cor:leafchain}).
\end{corollary}

% Reference 20: prop:centrerates; source line 4735 (one connected table).
% Proved in the last part: prop:centrerates.
\begin{proposition}[Rates at the centre]\label{prop:centrerates}
Let a family on $n\ge3$ marks be at the centre, every table with the exponents $(\sigma_j,\sigma_j)$, with the masses $\sigma\one$, $\sigma=\sum_j\sigma_j$, and let every active table be connected.
\begin{enumerate}[label=(\roman*)]
\item The walls are $\kappa_S=\sum_j\sigma_jc_j(S)$ (Proposition~\ref{prop:centrefamily}(ii)), so $\kappa_S=\kappa_{\cE\setminus S}$.
\item The rate of the peak closing (Theorem~\ref{thm:peak}(ii)) is $\varpi=\sigma/\max_j\max(r_j-1,n-r_j-1)$, the maximum over the active tables, attained at the faces $S=\cE-e$, where every active table has $c_j=1$, $d_j=r_j-1$ and $d^\ast_j=n-r_j-1$.
\item Near the vertex of $e$, with $u_e=1-x$ and $u_{\cE-e}=xv$, $v\in\Delta_{\cE-e}$, every active table has
\begin{gather*}
K_j(u)=(1-x)\,x^{r_j-1}K_{j,/e}(v)+x^{r_j}K_{j,\setminus e}(v),\\
\Psi_j(u)=(1-x)\,x^{n-r_j-1}\Psi_{j,\setminus e}(v)+x^{n-r_j}\Psi_{j,/e}(v);
\end{gather*}
the \emph{vertex product} of $e$ in table $j$ is $K_{j,/e}\Psi_{j,\setminus e}$, the product of the letters of the two minors that dominate as $x\to0$.
\item For two marks $e,f$ and a bijection $\beta$ with $\beta(e)=f$ that carries or exchanges the family, with the bijection $h$ of its tables, the restriction $\beta:\cE-e\to\cE-f$ takes the vertex product of $e$ in table $j$ to that of $f$ in table $h(j)$. For one table such a $\beta$ exists exactly when $e$ and $f$ lie in one class at the centre (Corollary~\ref{cor:certifyone}).
\end{enumerate}
\end{proposition}

\section{Devices at the centre}

% Reference 20: prop:halfletters; source line 4750
\begin{proposition}[The half-sum and the difference of a table with $2r=n$]\label{prop:halfletters}
Let $F=((\cB,\theta);\sigma,\sigma)$, $\sigma>0$, have $2r=n$ and masses $\mvec$ on its balance wall. The half-sum $\frac12(K+\Psi)=\frac12\sum_T\theta_T\big(\alpha^T+\alpha^{\cE\setminus T}\big)$ is homogeneous of degree $r$ and unchanged by $T\mapsto\cE\setminus T$; the difference $K-\Psi$ vanishes at the table's gauge; in signed coordinates they are $e^{\ip\one z}L_{1/2}(z)$ and $e^{\ip\one z}\partial_wL_w(z)$ (Definition~\ref{def:sideshape}). Then $K\Psi=\big(\frac12(K+\Psi)\big)^2-\frac14(K-\Psi)^2$ with $|K-\Psi|<K+\Psi$ on the orthant, and by Proposition~\ref{prop:pairladder}
\[
(K\Psi)^{-\sigma}=\sum_{l\ge0}\frac{(\sigma)_l}{l!\,4^l}\,(K-\Psi)^{2l}\Big(\frac{K+\Psi}2\Big)^{-2\sigma-2l},
\]
a series of nonnegative terms. For a mark $e$ with room $R=\cE-e$ write $\frac12(K+\Psi)=h_0(\alpha_R)+\alpha_eh_1(\alpha_R)$, so that, by Definition~\ref{def:delcon}, $h_0=\frac12(K_{\setminus e}+\Psi_{/e})$, $h_1=\frac12(K_{/e}+\Psi_{\setminus e})$ and $h_1(u)=u^{R}h_0(u^{-1})$. Expanding $(K-\Psi)^{2l}=\sum_{\mathbf a}\big([\alpha^{\mathbf a}](K-\Psi)^{2l}\big)\,\alpha^{\mathbf a}$ and integrating the mark $e$ of each term as the Beta corner of a linear letter (Proposition~\ref{prop:scalebreta} with one coefficient), with $w_l=2\sigma+2l$,
\begin{multline*}
P_F(\mvec;\sigma,\sigma)=\sum_{l\ge0}\frac{(\sigma)_l}{l!\,4^l}\sum_{\mathbf a}\big([\alpha^{\mathbf a}](K-\Psi)^{2l}\big)\,\Gst21{m_e+a_e,\ w_l-m_e-a_e}{w_l}\\
\times\int_{\Delta_R}u^{(\mvec+\mathbf a)|_R-\one}h_0(u)^{-(w_l-m_e-a_e)}h_1(u)^{-(m_e+a_e)}\,d\omega_R(u),
\end{multline*}
whenever the terms converge absolutely.
\end{proposition}

The half-sum is a positive polynomial whose support is closed under complements, but it is not in general the basis polynomial of a table: the bases together with the cobases need not have the exchange property, even for connected tables, so the integrals left over the room are integrals of a positive polynomial against its complement and not periods of a family. When complementation maps the bases onto themselves the difference vanishes and the room is a period.

% Reference 20: cor:mergedslots; source line 4765 (at the centre).
% Proved in the last part: cor:mergedvalues.
\begin{corollary}[Merged Skeleton values]\label{cor:mergedvalues}
If $T\mapsto\cE\setminus T$ maps the bases of a table $(\cB,\theta)$ onto themselves with $\theta_{\cE\setminus T}=\theta_T$, then $2r=n$, no mark is a loop or a coloop, $\Psi=K$, $\Psi_{\setminus e}=K_{/e}$ and $\Psi_{/e}=K_{\setminus e}$ for every mark $e$, and the two Skeleton values of the scale of $e$ coincide at every gauge (Corollary~\ref{cor:stage}). For $F=((\cB,\theta);s,t)$ with $s,t\ge0$, $s+t>0$, its masses on the balance wall with every proper wall positive, and a mark $e$, one has $0<m_e<s+t$ and
\[
P_F(\mvec;s,t)=\Gst21{m_e,\ s+t-m_e}{s+t}\,P_{\cB\setminus e}\big(\mvec|_{\cE-e};\,s+t-m_e,\,m_e\big),
\]
the period of the deletion $(\cB\setminus e,\theta)$ on $\cE-e$ at the exponents $(s+t-m_e,m_e)$, on its balance wall; so $P_F(\mvec;s,t)$ depends on the exponents only through $s+t$, and at the centre, $s=t=\sigma$ and $\mvec=\sigma\one$, the identity reads $P_F(\sigma\one;\sigma,\sigma)=\Gst21{\sigma,\ \sigma}{2\sigma}P_{\cB\setminus e}(\sigma\one|_{\cE-e};\sigma,\sigma)$.
\end{corollary}

% Reference 20: thm:atomforward; source line 4772 (one table).
% Proved in the last part: thm:atomforward.
\begin{theorem}[The scale integral of one mark]\label{thm:atomforward}
Let a family $F=((\cB_j,\theta_j);s_j,t_j)_j$ carry real exponents $s_j,t_j\ge0$, its masses on the balance wall with every proper wall positive, and let $e$ be a mark that is neither a loop nor a coloop in any active table, with room $R=\cE-e$; write $\Lambda_{j,\setminus e}$ and $\Lambda_{j,/e}$ for the centred polynomials of the deletion and the contraction of table $j$, so that $\Lambda_j(z)=e^{z_e}\Lambda_{j,/e}(z_R)+e^{-z_e}\Lambda_{j,\setminus e}(z_R)$ (Theorem~\ref{thm:taylorcube}).
\begin{enumerate}[label=(\roman*)]
\item At the gauge $(1,e^{2z_R})$ the scale integral of the bunch $\{e\}$ (Theorem~\ref{thm:scale}) is
\begin{multline*}
\operatorname{Sc}_{\{e\}}\big((1,e^{2z_R});\mvec;\mathbf s,\mathbf t\big)\\
=e^{-2\sigma\ip\one{z_R}}\int_0^\infty x^{m_e-1}\prod_j\big(\Lambda_{j,\setminus e}(z_R)+x\Lambda_{j,/e}(z_R)\big)^{-s_j}\big(\Lambda_{j,/e}(-z_R)+x\Lambda_{j,\setminus e}(-z_R)\big)^{-t_j}dx ,
\end{multline*}
$\sigma=\sum_j\sigma_j$ the total strength: up to the factor $e^{-2\sigma\ip\one{z_R}}$ it depends only on $m_e$, the exponents and the functions $\Lambda_{j,\setminus e}$ and $\Lambda_{j,/e}$ on the room, and Corollary~\ref{cor:fiber} at the bunch $\{e\}$ reads, for every $g\in R$,
\[
P_F(\mvec;\mathbf s,\mathbf t)=2^{n-2}\int_{\mathbb R^{R-g}}e^{2\ip{\mvec|_R}{z_R}}\,\operatorname{Sc}_{\{e\}}\big((1,e^{2z_R});\mvec;\mathbf s,\mathbf t\big)\,dz_{R-g}\Big|_{z_g=0}.
\]
\item At an atom $\Upsilon=((s_j,t_j)_j;\alpha^\star)$ with its target $\btau$, under $\mathrm{Breta}^\circ(F;\btau;\mathbf s,\mathbf t)$ the law of the log-ratios $(y_{g_k}-y_e)_{k=1}^{n-1}$, read in an order $\mathbf g_e$ of the room that writes the potential from $e$ as the one pair of functions of Proposition~\ref{prop:onewriting}, does not depend on $e$.
\item For one table at the centre, $s=t=\sigma$ and $\bmu=0$, with the even statistics $E_{\setminus e}$ and $E_{/e}$ of the deletion and the contraction (Theorem~\ref{thm:innerMGF}) and a two-element set $J$, the functions of $z_R$ understood,
\begin{multline*}
\int_0^\infty x^{\sigma-1}\big(\Lambda_{\setminus e}+x\Lambda_{/e}\big)^{-\sigma}\big(\Lambda_{/e}(-\cdot)+x\Lambda_{\setminus e}(-\cdot)\big)^{-\sigma}dx\\
=e^{-\frac\sigma2(E_{\setminus e}+E_{/e})}\,\sB^{\{J\}}_{(\sigma,\sigma),\,\sigma;\,I_J}\big(e^{\frac12(E_{\setminus e}-E_{/e})},\,e^{\frac12(E_{/e}-E_{\setminus e})}\big),
\end{multline*}
and the Breta is $\Gst21{\sigma,\ \sigma}{2\sigma}$ where the two even statistics agree.
\end{enumerate}
\end{theorem}

For one table at the centre the scale integral of a mark is the factor $e^{-\frac\sigma2(E_{\setminus e}+E_{/e})}$ times a Breta that sees the room through one number, the difference of the even statistics of the deletion and the contraction; at an atom the law of the log-ratios read from a mark does not depend on the mark. Proposition~\ref{prop:markbreta} reads the scale integral of the mark with the exponents in the shape index.

% Reference 20: prop:markbreta; source line 4792
\begin{proposition}[A mark is the Breta of its children's even letters]\label{prop:markbreta}
For one table, a mark $e$ that is neither a loop nor a coloop, exponents $s,t>0$, $0<m_e<s+t$ and the top wall $\kappa^{\mathrm{top}}_{\{e\}}=s+t-m_e$, integrating $\alpha_e=x$ with the room fixed,
\[
\int_0^\infty x^{m_e-1}K^{-s}\Psi^{-t}\,dx=K_{/e}^{\,t-m_e}\Psi_{\setminus e}^{\,s-m_e}\,\Gst22{m_e,\ \kappa^{\mathrm{top}}_{\{e\}}}{s,\ t}\ \sB^{\{J\}}_{(s,t),\,\kappa^{\mathrm{top}}_{\{e\}};\,I_J}\big(\mathbf E_{\setminus e},\mathbf E_{/e}\big),
\]
with $\mathbf E_{\setminus e}=K_{\setminus e}\Psi_{\setminus e}$ and $\mathbf E_{/e}=K_{/e}\Psi_{/e}$ the even letters of the deletion and the contraction, $J$ a two-element set, and
\[
\sB^{\{J\}}_{(s,t),\,\kappa;\,I_J}(a,b)=\int_0^1w^{s-1}(1-w)^{t-1}\,\sZ_{\kappa;(1)}\big(wa+(1-w)b\big)\,dw ,
\]
the shape $w$ of $J$ weighing the deletion against the contraction. At the centre, $s=t=m_e=\sigma$, the prefactor and the stacked Gamma are $1$.
\end{proposition}

The same scale integral is the Breta of Corollary~\ref{cor:pencilone} for the bunch $\{e\}$, whose levels $d$ and $r$ have the children $\cB\setminus e$ and $\cB/e$: there the walls $(m_e,\kappa^{\mathrm{top}}_{\{e\}})$ fill the shape index and the exponent $t$ the scale index, here the exponents $(s,t)$ fill the shape index and the top wall the scale index. Both read the scale integral of the mark through the even letters of the two children, the deletion and the contraction.

\section{The junction under level form}

% Reference 20: lem:hn; source line 4828
\begin{lemma}[An even power of $e^v-e^{-v}$ on a symmetric two-sided scale]\label{lem:hn}
For $\xi>0$ and $l\in\mathbb N$,
\[
\int_{\mathbb R}\big(e^v-e^{-v}\big)^{2l}e^{-\xi(e^v+e^{-v})}\,dv=4^l\big(\tfrac12\big)_l\,\xi^{-l}\,\sG_{l;(1,-1)}(\xi,\xi).
\]
\end{lemma}

% Reference 20: thm:uniformjunction; source line 4838 (one table).
% Proved in the last part: thm:uniformjunction.
\begin{theorem}[The junction under level form]\label{thm:uniformjunction}
Let $F=((\cB_j,\theta_j);\sigma_j,\sigma_j)_{j=1}^N$, every $\sigma_j>0$, be a molecule on $n$ marks with the masses $\sigma\one$, $\sigma=\sum_j\sigma_j$, let $B$, $b=|B|$, be a proper bunch with room $R$ such that every transposition of two marks of $B$ fixing the rest carries or exchanges every table, and let $g\in R$. By Proposition~\ref{prop:levelformQ} every half-sum has level form; in this statement write $q_{j,k}(z_R)$ for the common value of $c_{j,p}(\frac12;z_R)$ over $|p|=k$ and $\mathbf q_j=(q_{j,0},\dots,q_{j,b})$, so that $q_{j,b-k}(-z_R)=q_{j,k}(z_R)$ and
\[
L_{j,1/2}(z)=\sum_{k=0}^be^{-\ip\one{z_B}}e_k\big(e^{2z_B}\big)\,q_{j,k}(z_R),\qquad e^{-\ip\one{z_B}}e_k\big(e^{2z_B}\big)=\sum_{|p|=k}e^{\ip{z_B}{\chi_p}},
\]
the centred polynomial of the list on $B$ whose bases are the $k$-subsets of $B$; and for each $\mathbf l=(l_j)_j\in\mathbb N^N$
\[
\prod_j\big(\partial_wL_{j,w}(z)\big)^{2l_j}=e^{-2|\mathbf l|\ip\one{z_B}}\sum_{|\mathbf h|=2|\mathbf l|}\ \prod_{k=0}^be_k\big(e^{2z_B}\big)^{h_k}\,c^{(\mathbf l)}_{\mathbf h}(z_R),
\]
with $c^{(\mathbf l)}_{\mathbf h}$ Laurent polynomials in $e^{z_R}$.
\begin{enumerate}[label=(\roman*)]
\item With the Gramma $\sG_{0;(\chi_p)_{p\subseteq B}}\big((\upsilon_{|p|})_{p\subseteq B}\big)$, $\boldsymbol\upsilon=(\upsilon_0,\dots,\upsilon_b)$, whose coefficients depend on $|p|$ alone and which depends on $B$ through $b$ alone,
\begin{multline*}
P_F(\sigma\one;\boldsymbol\sigma,\boldsymbol\sigma)=2^{n-1}\sum_{\mathbf l\in\mathbb N^N}\prod_j\frac{(\sigma_j)_{l_j}}{l_j!\,4^{l_j}\,\Gamma(2\sigma_j+2l_j)}\int_{\mathbb R^{R-g}}\int_{(0,\infty)^N}\prod_jy_j^{2\sigma_j+2l_j-1}\\
\times\sum_{|\mathbf h|=2|\mathbf l|}c^{(\mathbf l)}_{\mathbf h}(z_R)\Big((-\partial_{\boldsymbol\upsilon})^{\mathbf h}\,\sG_{0;(\chi_p)_p}\big((\upsilon_{|p|})_p\big)\Big)\Big|_{\boldsymbol\upsilon=\sum_jy_j\mathbf q_j(z_R)}\,d\mathbf y\,dz_{R-g}\Big|_{z_g=0},
\end{multline*}
the inner sum being $\int_{\mathbb R^B}\prod_j\big(\partial_wL_{j,w}(z)\big)^{2l_j}e^{-\sum_jy_jL_{j,1/2}(z)}\,dz_B\ge0$, a level with $\sum_jy_jq_{j,k}=0$ absent.
\item \emph{One mark at a time.} For $e\in B$, $B'=B-e$ and $\boldsymbol\upsilon>0$,
\[
\sG_{0;(\chi_p)_{p\subseteq B}}\big((\upsilon_{|p|})_p\big)=\int_0^\infty\sG_{0;(\chi_p)_{p\subseteq B'}}\big((x\,\upsilon_{|p|+1}+x^{-1}\upsilon_{|p|})_{p}\big)\,\frac{dx}x,
\]
and for $B=\{e\}$ this Gramma is $\sG_{0;(1,-1)}(\upsilon_1,\upsilon_0)$.
\item \emph{The share side.} Let $G$ be the family on $B$ of the lists whose bases are the $k$-subsets of $B$, $1\le k\le b-1$, with the exponents $(a'_k,0)$, at the masses $\frac1b\sum_{k=1}^{b-1}ka'_k\,\one_B$ on its balance wall. Where every $\re a'_k>0$ its period converges, is unchanged when the exponents of the lists of ranks $k$ and $b-k$ are swapped, and is read on the hyperplane $\ip\one w=0$:
\[
P_G\Big(\tfrac1b\sum_{k=1}^{b-1}ka'_k\,\one_B;(a'_k)_{k=1}^{b-1},\mathbf 0\Big)=2^{b-1}\sqrt b\int_{\ip\one w=0}\ \prod_{k=1}^{b-1}e_k\big(e^{2w}\big)^{-a'_k}\,d\mathrm{vol}_{b-1}(w).
\]
For $\boldsymbol\upsilon>0$ and $\mathbf h\in\mathbb N^{b+1}$, on a contour with $\re a_k>h_k$ for every $k$ and with $a'_k=a_k-h_k$,
\begin{multline*}
\Big((-\partial_{\boldsymbol\upsilon})^{\mathbf h}\,\sG_{0;(\chi_p)_p}\big((\upsilon_{|p|})_p\big)\Big)(\boldsymbol\upsilon)\\
=2^{1-b}\int_{\sum_{k=0}^b(2k-b)(a_k-h_k)=0}\ \prod_{k=0}^b\Gamma(a_k)\,\upsilon_k^{-a_k}\ P_G\Big(\tfrac1b\sum_{k=1}^{b-1}ka'_k\,\one_B;(a'_k)_{k=1}^{b-1},\mathbf 0\Big).
\end{multline*}
\item \emph{Two marks.} If $B=\{e,f\}$ and the swap of $e$ and $f$ exchanges every table, then with $z_\pm=z_e\pm z_f$, for every $j$,
\begin{gather*}
L_{j,1/2}(z)=e^{-z_+}c_{j,\varnothing}(\tfrac12;z_R)+e^{z_+}c_{j,\varnothing}(\tfrac12;-z_R)+\big(e^{z_-}+e^{-z_-}\big)c_{j,\{e\}}(\tfrac12;z_R),\\
\partial_wL_{j,w}(z)=\big(e^{z_-}-e^{-z_-}\big)\,\partial_wc_{j,\{e\}}(w;z_R),
\end{gather*}
the functions $c_{j,\{e\}}(\frac12;\cdot)$ and $\partial_wc_{j,\{e\}}$ are even, and, writing in this statement $C_\varnothing(z_R)=\sum_jy_jc_{j,\varnothing}(\frac12;z_R)$ and $C_e(z_R)=\sum_jy_jc_{j,\{e\}}(\frac12;z_R)$,
\begin{multline*}
P_F(\sigma\one;\boldsymbol\sigma,\boldsymbol\sigma)=2^{n-2}\sum_{\mathbf l\in\mathbb N^N}\big(\tfrac12\big)_{|\mathbf l|}\prod_j\frac{(\sigma_j)_{l_j}}{l_j!\,\Gamma(2\sigma_j+2l_j)}\int_{\mathbb R^{R-g}}\prod_j\big(\partial_wc_{j,\{e\}}\big)^{2l_j}\\
\times\int_{(0,\infty)^N}\prod_jy_j^{2\sigma_j+2l_j-1}\,C_e^{-|\mathbf l|}\,\sG_{0;(1,-1)}\big(C_\varnothing(-z_R),\,C_\varnothing(z_R)\big)\\
\times\sG_{|\mathbf l|;(1,-1)}\big(C_e(z_R),\,C_e(z_R)\big)\,d\mathbf y\,dz_{R-g}\Big|_{z_g=0},
\end{multline*}
a series of nonnegative terms.
\end{enumerate}
\end{theorem}

Under level form the bunch enters only through its levels: one Gramma with one coefficient for each level, which depends on the number of its marks alone and is computed one mark at a time, is differentiated at the level coefficients of the half-sums, added over the tables with their scales, and its share side is a period of the lists of $k$-subsets of the bunch's own marks. For a pair whose swap exchanges every table the bunch splits along $z_e+z_f$ and $z_e-z_f$ into two Grammas $\sG_{\nu;(1,-1)}$.

\section{The node in one scale}\label{sec:onescale}

At a bunch of width one the integrand of Theorem~\ref{thm:twosides} is, in the even share, the product of two Mellin transforms, one carrying the bunch and one the room. Read back on the half-line, the period is one integral over one scale of two functions, the bunch transform and the room transform, each an integral of a Gramma $\sG_{\nu;(1,-1)}$ over its own shape. A closing toward a level is the expansion of one of them in powers of the scale, integrated term by term against the other. The expansion of a bunch is measured by one number, its reach: every bunch has a reach, no reach lies below the least value of the bunch's ratio, and a node closes on its lattice toward every level at which the growths of its two families multiply to less than one. Each side of a node is an interval, from the decay of its transform to its reach, and, against a room and a bunch bounded below, a node closes on a geometric lattice exactly when the two intervals are disjoint.

% Reference 20: def:nodetransforms; source line 5120 (with the setting of its section, source line 5114)
\begin{definition}[Node of width one; its two transforms]\label{def:nodetransforms}
A \emph{node of width one}, a node of a tree of bunches (Definition~\ref{def:bunchtree}) whose bunch has width one, is a family $F=((\cB,\theta);s,t)$ with real data, $s,t>0$, its masses on the balance wall with every proper wall positive, and a bunch $B$ of width one with room $R$, the weights factoring across them: the setting of Theorem~\ref{thm:twosides} for one table. Along the level $d$ its level shares are $(a^d,b^d,a^r,b^r)=(s-\kappa_B+b^d,\ b^d,\ \kappa_B-b^d,\ t-b^d)$, the odd share being fixed at $2\delta^d=s-\kappa_B$; its ratios $\varrho_B$ on $\Delta_B$ and $\varrho_R$ on $\Delta_R$ are those of Definition~\ref{def:readings}; and $\sigma^d_\ast$ is a real value of the even share at which both families converge and the four arguments of its stacked Gamma are positive, so that the line $\re\sigma^d=\sigma^d_\ast$ is a contour of Theorem~\ref{thm:twosides}, with $b_\ast=\sigma^d_\ast-\delta^d>0$ the real part of $b^d$ there. The node carries the nonnegative measures
\begin{gather*}
\phi_B(du)=u^{\mvec|_B-\one}K_{\cB/R}^{-(s-\kappa_B)}K_{\cB|B}^{-\kappa_B}\Psi_{\cB|B}^{-t}\,d\omega_B(u)\quad\text{on }\Delta_B,\\
\phi_R(dv)=v^{\mvec|_R-\one}K_{\cB|R}^{-(s-\kappa_B)}K_{\cB/B}^{-\kappa_B}\Psi_{\cB/B}^{-t}\,d\omega_R(v)\quad\text{on }\Delta_R,
\end{gather*}
the forms of the bunch family and of the room family at the level share $b^d=0$; Chapter~\ref{ch:hinged} calls $\phi_B$ the base form of the cell of the bunch $\cB^+_B$ at the exponents $(s,t)$. At the share $b^d$ the periods of the two families are
\[
P_{F^{\mathrm{bunch}}(\mathbf a,\mathbf b)}(\mvec|_B;\mathbf a,\mathbf b)=\int_{\Delta_B}\varrho_B^{-b^d}\,d\phi_B,\qquad P_{F^{\mathrm{room}}(\mathbf a,\mathbf b)}(\mvec|_R;\mathbf a,\mathbf b)=\int_{\Delta_R}\varrho_R^{\,b^d}\,d\phi_R,
\]
continued meromorphically in $b^d$ by their sector forms (Theorem~\ref{thm:exist}, Lemma~\ref{lem:cube}): the \emph{continued periods along the level $d$}. For $y>0$ the \emph{bunch transform} and the \emph{room transform} are
\begin{gather*}
f_B(y)=\int_{\Delta_B}\sG_{s-\kappa_B;(1,-1)}\big(1,\varrho_B(u)\,y\big)\,\phi_B(du),\\
f_R(y)=\int_{\Delta_R}\big(\varrho_R(v)\,y\big)^{\kappa_B}\,\sG_{t-\kappa_B;(1,-1)}\big(1,\varrho_R(v)\,y\big)\,\phi_R(dv),
\end{gather*}
with values in $[0,\infty]$. Swapping $B$ and $R$ gives the \emph{swapped node}, with the bunch $R$, the room $B$ and the two levels swapped (Theorem~\ref{thm:twosides}).
\end{definition}

The bunch family at the share $b^d$ has the form $\phi_B$ times $(\mathbf E_{\cB/R}/\mathbf E_{\cB|B})^{-b^d}=\varrho_B^{-b^d}$, and the room family the form $\phi_R$ times $(\mathbf E_{\cB|R}/\mathbf E_{\cB/B})^{-b^d}=\varrho_R^{\,b^d}$: the even share enters through the ratio of the hinge values of each side alone (Lemma~\ref{lem:slopesrange}).

% Reference 20: lem:twoletters; source line 5139
\begin{lemma}[Two letters in one scale]\label{lem:twoletters}
For linear letters $L_i(\rho)=L^{(0)}_i+\rho L^{(1)}_i$ with positive coefficients and the Skeleton values $\xi_i=L^{(0)}_i/L^{(1)}_i$, $i=1,2$, exponents $s,t>0$ and $0<\kappa<s+t$,
\begin{multline*}
\int_0^\infty\rho^{\kappa-1}L_1(\rho)^{-s}L_2(\rho)^{-t}\,d\rho\\
=\frac{(L^{(0)}_1)^{\kappa-s}(L^{(1)}_1)^{-\kappa}(L^{(0)}_2)^{-t}}{\Gamma(s)\Gamma(t)}\int_0^\infty x^{\kappa-1}\,\sG_{s-\kappa;(1,-1)}\Big(1,\frac{\xi_1}{\xi_2}\,x\Big)\,\sG_{t-\kappa;(1,-1)}(1,x)\,dx .
\end{multline*}
\end{lemma}

% Reference 20: thm:onescale; source line 5201
\begin{theorem}[The node in one scale]\label{thm:onescale}
Let $F$ and $B$ be a node of width one (Definition~\ref{def:nodetransforms}).
\begin{enumerate}[label=(\roman*)]
\item The period is one integral over one scale of the two transforms,
\[
P_F(\mvec;s,t)=\frac1{\Gamma(s)\Gamma(t)}\int_0^\infty f_B(y)\,f_R(y)\,\frac{dy}y ,
\]
the integrand nonnegative.
\item With the hinge values of Definition~\ref{def:readings},
\[
f_B(y)=\int_{\Delta_B}u^{\mvec|_B-\one}K_{\cB|B}^{-s}\Psi_{\cB|B}^{-t}\,\sG_{s-\kappa_B;(1,-1)}\big(\gamma^K_B(u)^{-1},\,y\,\gamma^\Psi_B(u)\big)\,d\omega_B(u):
\]
the bunch meets the room through its two hinge values, as in pooling (Theorem~\ref{thm:pooling}).
\item For real $b^d>0$ with $a^d>0$, and for real $b^d$ with $a^r>0$ and $b^r>0$, respectively, in $[0,\infty]$,
\begin{gather*}
\int_0^\infty y^{b^d-1}f_B(y)\,dy=\Gamma(b^d)\,\Gamma(a^d)\int_{\Delta_B}\varrho_B^{-b^d}\,d\phi_B,\\
\int_0^\infty y^{-b^d-1}f_R(y)\,dy=\Gamma(a^r)\,\Gamma(b^r)\int_{\Delta_R}\varrho_R^{\,b^d}\,d\phi_R .
\end{gather*}
Where the integrals over the shapes are finite they are the continued periods of the bunch family and of the room family at the share $b^d$, with their Gamma pairs, and (i), written through these two integrals in the share $b^d$, is Theorem~\ref{thm:twosides} for one table.
\item If no slope on the room is positive, so that $\inf\varrho_R>0$ (Lemma~\ref{lem:slopesrange}), then $\phi_R(\Delta_R)<\infty$, and with $a=\frac{t+\kappa_B}2-\frac14$
\[
f_R(y)\le\sqrt\pi\,e^{(t-\kappa_B)^2/4}\,\phi_R(\Delta_R)\,\big(y\inf\varrho_R\big)^a\,e^{-2\sqrt{y\inf\varrho_R}}\qquad\Big(y\ge\frac{\max(1,a^2)}{\inf\varrho_R}\Big):
\]
the room transform decays at the least ratio of the room.
\end{enumerate}
\end{theorem}

% Reference 20: lem:orthant; source line 5244
\begin{lemma}[The bunch transform on the orthant]\label{lem:orthant}
For a node of width one and $y>0$,
\[
f_B(y)=\int_{\mathbb R^B_{>0}}x^{\mvec|_B-\one}K_{\cB|B}(x)^{-s}\Psi_{\cB|B}(x)^{-t}\exp\Big(-\frac1{\gamma^K_B(x)}-y\,\gamma^\Psi_B(x)\Big)\,dx .
\]
\end{lemma}

The bunch transform is a Laplace transform in the hinge value $\gamma^\Psi_B$, damped by the hinge value $\gamma^K_B$: the room meets the bunch through its two hinge values and nothing else. Its expansion in powers of $y$ is measured by its reach.

% Reference 20: def:latticeexp; source line 5149
\begin{definition}[Lattice expansion; reach]\label{def:latticeexp}
A function $f$ on $(0,\infty)$ has a \emph{lattice expansion of reach} $x\ge0$ when there are finitely many reals $\eta_1,\dots,\eta_q$, an integer $J\ge0$ and coefficients $\beta_{\eta,j}$, for $\eta$ in $A=\bigcup_i(\eta_i+\mathbb N)$ and $0\le j\le J$, such that
\[
f(y)=\sum_{\eta\in A}\sum_{j=0}^J\beta_{\eta,j}\,y^\eta(-\log y)^j\quad(y>0),\qquad |\beta_{\eta,j}|\le C(1+\eta)^D\,\frac{x^\eta}{\Gamma(1+\eta)^2}\quad(\eta\in A,\ \eta\ge0),
\]
with constants $C$ and $D$; the series then converges absolutely at every $y>0$, and $\sum_{\eta,j}|\beta_{\eta,j}|<\infty$. Its \emph{least exponent} is the least $\eta$ with some $\beta_{\eta,j}\neq0$. A function has at most one lattice expansion. The \emph{reach of the bunch} of a node of width one, $\operatorname{reach}(B)$, is the infimum of the $x$ for which the bunch transform has a lattice expansion of reach $x$, and the \emph{reach of the room}, $\operatorname{reach}(R)$, is the reach of $R$ as the bunch of the swapped node.
\end{definition}

% Reference 20: lem:onemark; source line 5157
\begin{lemma}[The lattice of one mark]\label{lem:onemark}
Let $\nu$ be real.
\begin{enumerate}[label=(\roman*)]
\item For $x>0$, $\sG_{\nu;(1,-1)}(1,x)$ is the sum of the residues of $\Gamma(z)\Gamma(z+\nu)x^{-z}$ at the poles of $\Gamma(z)\Gamma(z+\nu)$. Off the integers
\[
\sG_{\nu;(1,-1)}(1,x)=\sum_{k\ge0}\frac{(-1)^k}{k!}\Big(\Gamma(\nu-k)\,x^k+\Gamma(-\nu-k)\,x^{\nu+k}\Big);
\]
at $\nu=l\in\mathbb N$
\begin{multline*}
\sG_{l;(1,-1)}(1,x)=\sum_{k<l}\frac{(-1)^k}{k!}\Gamma(l-k)\,x^k\\
+(-1)^l\sum_{k\ge l}\frac{x^k}{k!\,(k-l)!}\Big((\log\Gamma)'(1+k)+(\log\Gamma)'(1+k-l)-\log x\Big),
\end{multline*}
and $\sG_{-l;(1,-1)}(1,x)=x^{-l}\sG_{l;(1,-1)}(1,x)$.
\item For $x\ge1$, $0<\sG_{\nu;(1,-1)}(1,x)\le\sqrt\pi\,e^{\nu^2/4}\,x^{\nu/2-1/4}e^{-2\sqrt x}$.
\item For $w>0$, $y\mapsto\sG_{\nu;(1,-1)}(1,w^2y)$ has a lattice expansion of reach $w^2$, with exponents in $\mathbb N\cup(\nu+\mathbb N)$, in $-l+\mathbb N$ at $\nu=-l$, and at most one power of $-\log y$.
\end{enumerate}
\end{lemma}

% Reference 20: lem:latticeproduct; source line 5190
\begin{lemma}[Products of lattice expansions]\label{lem:latticeproduct}
If $f_1$ and $f_2$ have lattice expansions of reaches $x_1$ and $x_2$, then $f_1f_2$ has one of reach $(\sqrt{x_1}+\sqrt{x_2})^2$, whose terms are the products of theirs collected by exponent, and for real $a$, $y^af_1$ has one of reach $x_1$.
\end{lemma}

% Reference 20: lem:latticepairing; source line 5254
\begin{lemma}[Closing a lattice expansion against a decaying function]\label{lem:latticepairing}
Let $f$ have a lattice expansion $\sum_{\eta,j}\beta_{\eta,j}y^\eta(-\log y)^j$ of reach $x$ with least exponent $\eta_0$, and let $g$ be measurable on $(0,\infty)$ and integrable on its compact subsets, with $\int_0^1y^{\eta_0-1-\epsilon}|g(y)|\,dy<\infty$ for some $\epsilon>0$ and $|g(y)|\le C'y^ae^{-2\sqrt{x'y}}$ for $y\ge y_0$, where $x<x'$. Then
\[
\int_0^\infty f(y)\,g(y)\,\frac{dy}y=\sum_{\eta,j}\beta_{\eta,j}\int_0^\infty y^{\eta-1}(-\log y)^j\,g(y)\,dy,
\]
every integral on the right and the series converging absolutely, and the terms with $\eta\ge N$ add up in absolute value to at most $C''(1+N)^{D''}(x/x')^N$, with $C''$ and $D''$ independent of $N$.
\end{lemma}

% Reference 20: lem:cubegrowth; source line 5269
\begin{lemma}[Continued cube integrals grow at most exponentially]\label{lem:cubegrowth}
Let $G$ and $h$ be analytic on a neighbourhood of $[0,1]^m$ in $\mathbb C^m$, with $h>0$ on $[0,1]^m$, and let $a_k,b_k$ be real with $a_k>0$ wherever $b_k=0$. The integral
\[
I(z)=\int_{[0,1]^m}\prod_kq_k^{a_k+b_kz-1}\,G(q)\,h(q)^{-z}\,dq ,
\]
convergent where every $a_k+b_k\re z>0$, continues meromorphically to $z\in\mathbb C$, with poles only in the set $\mathcal W$ of the real points at which some $a_k+b_kz$ lies in $-\mathbb N$ (Lemma~\ref{lem:cube}). For every $\vartheta>0$ there is $M$, and then for every $\epsilon>0$ there are $C$ and $D$, with
\[
|I(z)|\le C(1+|z|)^D\,M^{|\re z|}\,e^{\vartheta|\im z|}
\]
at every $z$ at distance at least $\epsilon$ from $\mathcal W$.
\end{lemma}

% Reference 20: the growth rate (unlabelled paragraph); source line 5292; the growth at the integers, thm:simpletop, source line 5612
\begin{definition}[Growth along a level; growth at the integers]\label{def:growth}
For a node of width one, the \emph{growth} $\operatorname{growth}(B)$ of the bunch family along the level $d$ is the infimum of the $M\ge0$ for which, for every $\epsilon>0$, there are $C$, $D$ and $\vartheta<\pi$ with
\[
\big|P_{F^{\mathrm{bunch}}(\mathbf a,\mathbf b)}(\mvec|_B;\mathbf a,\mathbf b)\big|\le C(1+|b^d|)^D\,M^{-\re b^d}\,e^{\vartheta|\im b^d|}
\]
for the continued period at the share $b^d$, at every $b^d$ with $\re b^d\le0$ at distance at least $\epsilon$ from the set $\mathcal W_B$ of the real points at which a proper wall of the bunch family lies in $-\mathbb N$, which holds its poles; the growth $\operatorname{growth}(R)$ of the room family is the same with $P_{F^{\mathrm{room}}(\mathbf a,\mathbf b)}(\mvec|_R;\mathbf a,\mathbf b)$ and the set $\mathcal W_R$ of the room family. Along the level $r$ the growths are those of the swapped node. The \emph{growth at the integers} of the bunch is
\[
\operatorname{growth}_{\mathbb Z}(B)=\limsup_{k\to\infty}\big|P_{F^{\mathrm{bunch}}(\mathbf a,\mathbf b)}(\mvec|_B;\mathbf a,\mathbf b)\big|^{1/k},
\]
the continued period of the bunch family taken at $b^d=-k$, $k\in\mathbb N$.
\end{definition}

% Reference 20: thm:nodeclosing; source line 5298
\begin{theorem}[A node closes at the growth of its two families]\label{thm:nodeclosing}
Let $F$ and $B$ be a node of width one.
\begin{enumerate}[label=(\roman*)]
\item \emph{Every bunch has a reach.} The growths are finite; $\operatorname{growth}(B)\le\sup\varrho_B$ and $\operatorname{growth}(R)\le1/\inf\varrho_R$ when these are finite; and $\operatorname{reach}(B)\le\operatorname{growth}(B)$. The exponents of the lattice expansion of $f_B$ are the $\eta$ for which $b^d=-\eta$ is a pole of $\Gamma(b^d)\Gamma(a^d)P_{F^{\mathrm{bunch}}(\mathbf a,\mathbf b)}(\mvec|_B;\mathbf a,\mathbf b)$ to the left of the line $\re b^d=b_\ast$, and $\sum_jj!\,\beta_{\eta,j}(b^d+\eta)^{-j-1}$ is its principal part there: the lattice of the bunch is the lattice of its continued period with its Gamma pair, the wall residues joining.
\item \emph{Against a room bounded below.} Let no slope on the room be positive, and let $\operatorname{reach}(B)<\inf\varrho_R$. Then
\[
P_F(\mvec;s,t)=\frac1{\Gamma(s)\Gamma(t)}\sum_{\eta,j}\beta_{\eta,j}\,\partial_{b^d}^{\,j}\Big[\Gamma(a^r)\Gamma(b^r)\,P_{F^{\mathrm{room}}(\mathbf a,\mathbf b)}(\mvec|_R;\mathbf a,\mathbf b)\Big]_{b^d=-\eta},
\]
absolutely, every room period in it convergent, and for every $x>\operatorname{reach}(B)$ the terms with $\eta\ge N$ add up in absolute value to at most $C(1+N)^D(x/\inf\varrho_R)^N$. The terms are the residues of the integrand of Theorem~\ref{thm:twosides} at its poles toward the level $d$: the even share closes on the lattice of continued periods of the bunch family at the root rate $\operatorname{reach}(B)/\inf\varrho_R$.
\item \emph{Both families continued.} Whatever the slopes, if $\operatorname{growth}(B)\operatorname{growth}(R)<1$, then $P_F(\mvec;s,t)$ is the sum of the residues of the integrand
\[
\Gst42{a^d,\ b^d,\ a^r,\ b^r}{s,\ t}\,P_{F^{\mathrm{bunch}}(\mathbf a,\mathbf b)}(\mvec|_B;\mathbf a,\mathbf b)\,P_{F^{\mathrm{room}}(\mathbf a,\mathbf b)}(\mvec|_R;\mathbf a,\mathbf b)
\]
of Theorem~\ref{thm:twosides}, continued in $b^d$, at its poles with $\re b^d<b_\ast$; for every $M_B>\operatorname{growth}(B)$ and $M_R>\operatorname{growth}(R)$ the residues at the poles with real part in $(-N-1,-N]$ add up in absolute value to at most $C(1+N)^D(M_BM_R)^N$. The even share closes on the lattices of continued periods of both families, the wall residues of both joining.
\item \emph{Toward the level $r$.} Swapping $B$ and $R$ leaves the period and Theorem~\ref{thm:twosides} unchanged; the swapped node has the bunch $R$ with the ratio $\varrho_R$ and the room $B$ with the ratio $\varrho_B$. Applied to it, (i) to (iii) close the even share toward the level $r$: against a bunch with $\inf\varrho_B>0$ at the root rate $\operatorname{reach}(R)/\inf\varrho_B$, and in general at a root rate at most the product of the growths of the two families along the level $r$.
\end{enumerate}
\end{theorem}

A node is therefore closed on its lattice toward the level at which the growths of its two families multiply to less than one, each growth read on its family alone; where neither level has a product below one, the peak closing (Theorem~\ref{thm:pencilpeaks}) closes the node with nonnegative terms and needs no root rate. Against a room bounded below the bunch enters through its reach alone, and the domination is a reach against a decay.

% Reference 20: prop:reachfloor; source line 5332
\begin{proposition}[No reach is below the least ratio]\label{prop:reachfloor}
Let $F$ and $B$ be a node of width one.
\begin{enumerate}[label=(\roman*)]
\item If $\inf\varrho_B>0$, there are $C$ and $a$ with $f_B(y)\le Cy^a\,e^{-2\sqrt{y\inf\varrho_B}}$ for $y\ge1$, and for every $x'>\inf\varrho_B$ there is $C_->0$ with $f_B(y)\ge C_-e^{-2\sqrt{x'y}}$ for $y\ge1$: the bunch transform decays at the least ratio of the bunch.
\item $\operatorname{reach}(B)\ge\inf\varrho_B$: no lattice expansion of the bunch transform has a reach below the least ratio.
\item When $\varrho_B$ is bounded, $\operatorname{reach}(B)=\operatorname{growth}(B)=\sup\varrho_B$: the reach of a bounded bunch is its largest ratio.
\end{enumerate}
\end{proposition}

Against a room bounded below, the root rate $\operatorname{reach}(B)/\inf\varrho_R$ of Theorem~\ref{thm:nodeclosing}(ii) therefore lies between $\inf\varrho_B/\inf\varrho_R$ and $\sup\varrho_B/\inf\varrho_R$. The series bunches take the lower end, and the bounded bunches the upper, which is $\sup(\varrho_B/\varrho_R)$, the root rate of Corollary~\ref{cor:closetwosides}(i).

% Reference 20: prop:seriesproduct; source line 5364
\begin{proposition}[A series bunch is its marks]\label{prop:seriesproduct}
Let the bunch $B$ of a node of width one be a series bunch, $\theta'_B$ the weight of its pattern $B$ and $\theta'_{B-e}$ that of its pattern $B-e$ (Definition~\ref{def:readings}, Proposition~\ref{prop:classreadings}), and in this statement write $w_e=\theta'_{B-e}/\theta'_B$.
\begin{enumerate}[label=(\roman*)]
\item The bunch transform is the product over the marks of $B$:
\begin{align*}
f_B(y)&=(\theta'_B)^{-s-t}\prod_{e\in B}\sG_{m_e-s;(1,-1)}(w_ey,\,w_e)\\
&=(\theta'_B)^{-s-t}\,y^{s-\kappa_B}\prod_{e\in B}w_e^{s-m_e}\,\sG_{m_e-s;(1,-1)}\big(1,\,w_e^2y\big).
\end{align*}
\item $\varrho_B(u)=\big(\sum_ew_e/u_e\big)\big(\sum_ew_eu_e\big)\ge\big(\sum_ew_e\big)^2$, with equality exactly at the point of equal coordinates of $\Delta_B$, and $\varrho_B$ is unbounded: whatever the weights, the least ratio of a series bunch is taken at the point of equal coordinates of its shape.
\item $f_B$ has the lattice expansion of reach $\big(\sum_ew_e\big)^2=\inf\varrho_B$ that is the product of the lattices of its marks, so $\operatorname{reach}(B)=\inf\varrho_B$. With $\boldsymbol\nu=\mvec|_B-s\one_B$ and no $\nu_e$ an integer, its exponents lie in the lattices $-\nu_S+\mathbb N$, $S\subseteq B$: $S=\varnothing$ and $S=B$ are the lattices of the Gamma pair, and for every other $S$ the lattice is where the wall of the bunch family at $B\setminus S$, continued, meets $-\mathbb N$. Finally $f_B(y)\le C\,y^a\,e^{-2(\sum_ew_e)\sqrt y}$ for $y\ge\max_ew_e^{-2}$, with constants $C$ and $a$.
\end{enumerate}
\end{proposition}

% Reference 20: prop:seriesconnection; source line 5387
\begin{proposition}[A bunch in series is the product of its parts]\label{prop:seriesconnection}
Let the bunch of a node of width one be $B=B_1\sqcup\cdots\sqcup B_q$, and on each $B_i$ let $\mathcal X_i$ and $\mathcal Y_i$ be tables, $\mathcal Y_i$ of rank one less than $\mathcal X_i$, with the hinge values $\gamma^K_i=K_{\mathcal X_i}/K_{\mathcal Y_i}$ and $\gamma^\Psi_i=\Psi_{\mathcal Y_i}/\Psi_{\mathcal X_i}$ and the ratio $\varrho_i=\mathbf E_{\mathcal Y_i}/\mathbf E_{\mathcal X_i}$, such that for constants $C,C'>0$
\[
K_{\cB|B}=C\prod_iK_{\mathcal X_i},\qquad \Psi_{\cB|B}=C'\prod_i\Psi_{\mathcal X_i},\qquad \frac1{\gamma^K_B}=\sum_i\frac1{\gamma^K_i},\qquad \gamma^\Psi_B=\sum_i\gamma^\Psi_i :
\]
$B$ is the \emph{series connection} of the parts $B_i$. A series bunch is the series connection of its marks, with $K_{\mathcal X_e}=\alpha_e$, $\Psi_{\mathcal X_e}=1$, $K_{\mathcal Y_e}=w_e$, $\Psi_{\mathcal Y_e}=w_e\alpha_e$ and $\varrho_e=w_e^2$ (Proposition~\ref{prop:seriesproduct}).
\begin{enumerate}[label=(\roman*)]
\item $f_B=C^{-s}C'^{-t}\prod_if^\circ_i$, with the \emph{orthant transforms} of the parts
\[
f^\circ_i(y)=\int_{\mathbb R^{B_i}_{>0}}x^{\mvec|_{B_i}-\one}K_{\mathcal X_i}(x)^{-s}\Psi_{\mathcal X_i}(x)^{-t}\exp\Big(-\frac1{\gamma^K_i(x)}-y\,\gamma^\Psi_i(x)\Big)\,dx ;
\]
and $\sqrt{\inf\varrho_B}=\sum_i\sqrt{\inf\varrho_i}$.
\item If every $f^\circ_i$ has a lattice expansion of reach $x_i$, then $f_B$ has one of reach $\big(\sum_i\sqrt{x_i}\big)^2$.
\item Let every part be a single mark, or have a bounded ratio and a finite measure $u^{\mvec|_{B_i}-\one}K_{\mathcal X_i}^{-s}\Psi_{\mathcal X_i}^{-t}(\gamma^\Psi_i)^{-a_i}\,d\omega_{B_i}$ on $\Delta_{B_i}$, where $a_i=m_{B_i}-s\,r_i-t\,(|B_i|-r_i)$ with $r_i$ the rank of $\mathcal X_i$; and let no sum of the $a_i$ over a nonempty set of parts be an integer. Then $f^\circ_i$ has the reach $\sup\varrho_i$, and $\operatorname{reach}(B)=\big(\sum_i\sqrt{\sup\varrho_i}\big)^2$: in a series connection the square roots of the reaches add.
\end{enumerate}
\end{proposition}

% Reference 20: lem:centrestationary; source line 5435
\begin{lemma}[The point of equal coordinates is stationary]\label{lem:centrestationary}
Let $\mathcal X$ be a table on a set $S$, with any weights. At every point $u$ of $(0,\infty)^S$ with equal coordinates, $u_e\,\partial_{u_e}\log\mathbf E_{\mathcal X}(u)=1$ for every $e\in S$. So a ratio of even letters of tables on $S$ is stationary at the point of equal coordinates of $\Delta_S$: $\varrho_B$ on $\Delta_B$, $\varrho_R$ on $\Delta_R$, and $\varrho_B/\varrho_R$ at the pair of these points.
\end{lemma}

% Reference 20: thm:nodeintervals; source line 5442
\begin{theorem}[The two intervals of a node]\label{thm:nodeintervals}
For a node of width one, the bunch transform of the swapped node is $y^{-t}f_R(y)$, so that $\inf\varrho_B$ and $\inf\varrho_R$ are the decays of the two transforms, and $\operatorname{reach}(B)$ and $\operatorname{reach}(R)$ their reaches (Proposition~\ref{prop:reachfloor}(i), Lemma~\ref{lem:latticeproduct}).
\begin{enumerate}[label=(\roman*)]
\item \emph{No transform reaches less far than it decays}: $\inf\varrho_B\le\operatorname{reach}(B)$ and $\inf\varrho_R\le\operatorname{reach}(R)$.
\item \emph{Toward the level $d$.} Let no slope on the room be positive, so that $\inf\varrho_R>0$, and index the terms of the lattice of Theorem~\ref{thm:nodeclosing}(ii) by $(\eta,j)$. Their root rate is $\operatorname{reach}(B)/\inf\varrho_R$: for every $x>\operatorname{reach}(B)$ they are at most $C(1+\eta)^D(x/\inf\varrho_R)^\eta$ in modulus, and for every $\epsilon>0$ infinitely many of them exceed $\big(\operatorname{reach}(B)/\inf\varrho_R-\epsilon\big)^\eta$. So the lattice converges absolutely when $\operatorname{reach}(B)<\inf\varrho_R$, and its terms do not tend to zero when $\operatorname{reach}(B)>\inf\varrho_R$.
\item \emph{Toward the level $r$.} Let no slope on the bunch be negative, so that $\inf\varrho_B>0$. The terms of the lattice of continued periods toward the level $r$ have the root rate $\operatorname{reach}(R)/\inf\varrho_B$ in the same sense: the lattice converges absolutely when $\operatorname{reach}(R)<\inf\varrho_B$, and its terms do not tend to zero when $\operatorname{reach}(R)>\inf\varrho_B$.
\item \emph{One level at most.} $\operatorname{reach}(B)<\inf\varrho_R$ and $\operatorname{reach}(R)<\inf\varrho_B$ never hold together. Under both hypotheses of (ii) and (iii) the node therefore closes on a geometrically convergent lattice exactly when the intervals $[\inf\varrho_B,\operatorname{reach}(B)]$ of the bunch and $[\inf\varrho_R,\operatorname{reach}(R)]$ of the room are disjoint, toward the level of the lower one, at the ratio of the upper end of the lower interval to the lower end of the upper. When each interval extends beyond the lower end of the other, neither lattice converges, and the share closes at its peaks (Theorem~\ref{thm:pencilpeaks}); when they touch, the root rate is one and domination decides nothing.
\end{enumerate}
\end{theorem}

Each side of a node is read by two numbers of its own transform, its decay and the reach of its expansion, and no expansion reaches less far than its transform decays. A series bunch has both at its least ratio, and its interval is the point $(\sum_ew_e)^2$ (Proposition~\ref{prop:seriesproduct}); a bounded bunch spreads over its whole range $[\inf\varrho_B,\sup\varrho_B]$ (Proposition~\ref{prop:reachfloor}(iii)); in a series connection of marks and bounded parts the square roots of the reaches add (Proposition~\ref{prop:seriesconnection}(iii)). When the bunch and the room are both series bunches the two intervals are points: the node closes toward the level of the smaller point at the ratio of the two points, its terms toward the other level grow at the inverse ratio, and where the points meet domination decides nothing. Where a far side is unbounded, the level toward it can still close with both families continued, at a root rate at most the product of their growths (Theorem~\ref{thm:nodeclosing}(iii)). By Lemma~\ref{lem:centrestationary} the value of $\varrho_B/\varrho_R$ at the points of equal coordinates of $\Delta_B$ and $\Delta_R$ is a stationary value of the node's ratio, but the root rate is read at the extremes of the two transforms. The reach of a bunch that is neither bounded nor built from parallel and series bunches is read on the law of its ratio (Section~\ref{sec:reachlaw}).

% Reference 20: thm:classcloses; source line 5415 (series and parallel bunches).
% Proved in the last part: thm:classcloses.
\begin{theorem}[A series connection closes at its reach]\label{thm:classcloses}
\begin{enumerate}[label=(\roman*)]
\item Let the bunch $B$ of a node of width one be the series connection of parts $B_1,\dots,B_q$ (Proposition~\ref{prop:seriesconnection}), every orthant transform $f^\circ_i$ having a lattice expansion of reach $x_i$, and let no slope on the room be positive. If $\big(\sum_i\sqrt{x_i}\big)^2<\inf\varrho_R$, the even share closes toward the level $d$ on the lattice of continued periods of the bunch family (Theorem~\ref{thm:nodeclosing}(ii)): for every $x>\big(\sum_i\sqrt{x_i}\big)^2$ its terms with $\eta\ge N$ add up in absolute value to at most a constant times $(1+N)^D(x/\inf\varrho_R)^N$, and its coefficients $\beta_{\eta,j}$ are those of $C^{-s}C'^{-t}\prod_if^\circ_i$, finite sums of products, over the parts, of the coefficients of their lattice expansions. Under the hypotheses of Proposition~\ref{prop:seriesconnection}(iii) the root rate of the terms is $\big(\sum_i\sqrt{\sup\varrho_i}\big)^2/\inf\varrho_R$ (Theorem~\ref{thm:nodeintervals}(ii)). For a series bunch, the series connection of its marks with $w_e$ as in Proposition~\ref{prop:seriesproduct} and $x_e=w_e^2$, the root rate is $\inf\varrho_B/\inf\varrho_R=\big(\sum_ew_e\big)^2/\inf\varrho_R$ and the coefficients are those of $(\theta'_B)^{-s-t}y^{s-\kappa_B}\prod_ew_e^{s-m_e}\sG_{m_e-s;(1,-1)}(1,w_e^2y)$: finite sums of products, over the marks of $B$, of the coefficients of Lemma~\ref{lem:onemark}, Laurent coefficients of the Gamma pairs $\Gamma(z)\Gamma(z+m_e-s)$.
\item Let $B$ be the series connection of parts $B_1,\dots,B_q$ in the dual table at the exponents $(t,s)$ (Proposition~\ref{prop:dualfamily}), its orthant transforms there having lattice expansions of reaches $x^\ast_i$, and let no slope on the room be negative. If $\big(\sum_i\sqrt{x^\ast_i}\big)^2\sup\varrho_R<1$, the even share closes toward the level $r$ on the lattice of continued periods of the bunch family, at a root rate at most $\big(\sum_i\sqrt{x^\ast_i}\big)^2\sup\varrho_R$. A parallel bunch $B$ is a series bunch of the dual table (Proposition~\ref{prop:cleanmaps}), with $w_e=\theta'_{\{e\}}/\theta'_\varnothing$, $\theta'_\varnothing$ and $\theta'_{\{e\}}$ the weights of its patterns $\varnothing$ and $\{e\}$; then $x^\ast_e=w_e^2$ and $\sup\varrho_B=\big(\sum_ew_e\big)^{-2}$, taken at the point of equal coordinates of $\Delta_B$, so the share closes when $\sup\varrho_R<\sup\varrho_B$, at the root rate $\sup\varrho_R/\sup\varrho_B$.
\end{enumerate}
\end{theorem}

No order of the values is asked: the room's ratio need only be bounded on the side away from the closing, and the bunch's ratio may be unbounded, its walls falling toward the level closed on. When the bunch and the room are both series bunches, the root rate $\inf\varrho_B/\inf\varrho_R$ is read at the points of equal coordinates of both shapes. When the extreme ratio of the room sits at a vertex of its shape simplex the root rate is not read at the points of equal coordinates, which are only stationary.

% Reference 20: thm:grammanested; source line 5465
\begin{theorem}[Products of Grammas $\sG_{\nu;(1,-1)}$ close on nested lattices]\label{thm:grammanested}
For a finite set $S$, numbers $w_e>0$ and real indices $\nu_e$, put $w_S=\sum_{e\in S}w_e$ and $\nu^-_S=\sum_{e\in S}\min(0,\nu_e)$, and in this statement write, for $\re z>-\nu^-_S$,
\[
\mathcal M_S(z)=\int_0^\infty y^{z-1}\prod_{e\in S}\sG_{\nu_e;(1,-1)}\big(1,w_e^2y\big)\,dy .
\]
\begin{enumerate}[label=(\roman*)]
\item $\mathcal M_{\{e\}}(z)=w_e^{-2z}\,\Gamma(z)\Gamma(z+\nu_e)$, and for two marks with $w_e=w_f=w$
\[
\mathcal M_{\{e,f\}}(z)=w^{-2z}\,\Gst41{z,\ z+\nu_e,\ z+\nu_f,\ z+\nu_e+\nu_f}{2z+\nu_e+\nu_f}.
\]
\item For $S=S_1\sqcup S_2$ with $w_{S_1}<w_{S_2}$, let $\prod_{e\in S_1}\sG_{\nu_e;(1,-1)}(1,w_e^2y)=\sum\beta_{\eta,j}y^\eta(-\log y)^j$ be the product of the lattices of the marks of $S_1$ (Lemmas~\ref{lem:onemark} and~\ref{lem:latticeproduct}). For $\re z>-\nu^-_S$ and $i\in\mathbb N$,
\[
(-\partial_z)^i\mathcal M_S(z)=\sum_{\eta,j}\beta_{\eta,j}\,(-\partial_z)^{i+j}\mathcal M_{S_2}(z+\eta),
\]
every $z+\eta$ in the domain of $\mathcal M_{S_2}$, absolutely, the terms with $\eta\ge N$ adding up in absolute value to at most $C(1+N)^D(w_{S_1}/w_{S_2})^{2N}$.
\item Removing at each step a mark $e$ of least $w_e$, until one mark is left or two with equal $w_e$, writes $\mathcal M_S(z)$ and its derivatives as nested lattice sums of products of Laurent coefficients of Gamma pairs at the points of the lattices, ending at (i): at most $|S|-1$ nested lattices, summed from the inside out, the one that removes $e$ from $S'$ converging absolutely at every point of the lattices above it, at the root rate $\big(w_e/(w_{S'}-w_e)\big)^2$, its constant depending on that point.
\end{enumerate}
\end{theorem}

% Reference 20: cor:cyclenested; source line 5487
\begin{corollary}[A table that is one circuit is a nested lattice of Gamma data]\label{cor:cyclenested}
In Corollary~\ref{cor:cycle} the integral over $y$ is the integral of Theorem~\ref{thm:grammanested} with $S=\cE$, $w_e=\theta_e$ and the indices $\nu_e=m_e-s$, taken at $z=s$, and $s>-\sum_e\min(0,\nu_e)$. So the period of every table that is one circuit, at all exponents $s,t>0$, all masses of Corollary~\ref{cor:cycle} and all weights, is a nested lattice sum of products of Laurent coefficients of Gamma pairs, with at most $n-1$ nested lattices, each converging absolutely at every point of the lattices above it and summed from the inside out. With unit weights, removing one mark at a time, the root rates are $(n-1)^{-2},(n-2)^{-2},\dots,\frac14$, and the last two marks give four Gammas over one; splitting the marks into a lighter part $S_1$ and a heavier part $S_2$ converges at the root rate $(\theta_{S_1}/\theta_{S_2})^2$. At integer data with unit weights every coefficient is a rational combination of products of values $(\log\Gamma)^{(i)}(k)$, $i,k\ge1$, which are Trailmoments of the counting input (Propositions~\ref{prop:headstails} and~\ref{prop:resonance}).
\end{corollary}

\section{The reach read on the law}\label{sec:reachlaw}

The reach of a bunch is defined by the expansion of its transform (Definition~\ref{def:latticeexp}) and is bounded below by the least ratio and above by the growth of the continued family (Proposition~\ref{prop:reachfloor}, Theorem~\ref{thm:nodeclosing}). Read on one measure on the half-line, the law of the ratio under the base form, the reach is the point beyond which the law has a density equal to its own wall series (Definition~\ref{def:lawbunch}), the series whose terms are the poles of the continued family; every point at which the law fails to be analytic therefore bounds the reach from below, and the least ratio and the largest ratio of a bounded bunch are such points. From above the reach is bounded by the ceiling of a continuation (Definition~\ref{def:keyhole}): in each sector the falling variables are carried around zero on circles, as far as the sector letters allow, and the largest modulus of the ratio met there bounds the reach. From below it is bounded by the growth of the continued family at the integers; a continuation whose largest values sit at simple tops that do not cancel has that growth as its ceiling, so that the reach is at least the modulus of a stationary value of the ratio on a flat face, and equal to it when the ceilings come down to it.

% Reference 20: def:lawbunch; source line 5500
\begin{definition}[The law of a bunch; wall series]\label{def:lawbunch}
The \emph{law} of the bunch of a node of width one is the image $(\varrho_B)_\ast\phi_B$ of its base form under its ratio, a measure on $(0,\infty)$ (Definition~\ref{def:inputs}). The continued period of the bunch family at the share $b^d$ is $\int_0^\infty x^{-b^d}\,d\big((\varrho_B)_\ast\phi_B\big)(x)$ wherever the integral converges, which is for $\re b^d$ in an open interval containing $b_\ast$, the \emph{domain interval} of the law, the walls being affine in $b^d$ and the conditions of Theorem~\ref{thm:projective} strict. A function $g$ on $(x_0,\infty)$, $x_0\ge0$, is a \emph{wall series of radius} $x_0$ when there are finitely many reals $\eta_i$, an integer $J\ge0$ and coefficients $a_{\eta,j}$, for $\eta\in A=\bigcup_i(\eta_i+\mathbb N)$ and $0\le j\le J$, with $\sum_{\eta,j}|a_{\eta,j}|x'^{-\eta}<\infty$ for every $x'>x_0$, such that
\[
g(x)=\sum_{\eta,j}a_{\eta,j}\,x^{-\eta-1}(\log x)^j\qquad(x>x_0).
\]
The series then converges absolutely and locally uniformly on $\{x\in\mathbb C:\ |x|>x_0,\ |\arg x|<\pi\}$, so a wall series is real-analytic.
\end{definition}

% Reference 20: thm:lawreach; source line 5508
\begin{theorem}[The reach is read on the law]\label{thm:lawreach}
Let $F$ and $B$ be a node of width one.
\begin{enumerate}[label=(\roman*)]
\item On $(\operatorname{reach}(B),\infty)$ the law $(\varrho_B)_\ast\phi_B$ has a density, and the density is a wall series of radius $\operatorname{reach}(B)$. Its terms are the poles of the continued period of the bunch family to the left of the domain interval: a pole $b^d=-\eta$ with the principal part $\sum_jj!\,a_{\eta,j}(b^d+\eta)^{-j-1}$ gives the terms $a_{\eta,j}x^{-\eta-1}(\log x)^j$, and there are no others.
\item If on some $(x_0,\infty)$ the law has a density that is a wall series of radius $x_0$, then $\operatorname{reach}(B)\le x_0$.
\item So the reach is the least $x_0\ge0$ for which the law has on $(x_0,\infty)$ a density that is a wall series of radius $x_0$. In particular $\operatorname{reach}(B)\ge x_1$ at every $x_1>0$ near which $(\varrho_B)_\ast\phi_B$ has no real-analytic density: at the least ratio when it is positive, and at the largest ratio of a bounded bunch.
\end{enumerate}
\end{theorem}

So no reach lies below a point where the law of its ratio is singular, and Proposition~\ref{prop:reachfloor}(ii), with the lower half of (iii), is the case of the two ends of the law. The law of a family at a node is likewise its own wall series beyond its growth, so that $\operatorname{reach}(B)\le\operatorname{growth}(B)$ of Theorem~\ref{thm:nodeclosing}(i) is an equality (Corollary~\ref{cor:growthreach}); along a direction of the shares the law reads the reach of a piece of any width (Theorem~\ref{thm:dirreach}). From above the reach is bounded by the ceilings of continuations.

% Reference 20: lem:keyhole; source line 5541
\begin{lemma}[Finite parts on a keyhole]\label{lem:keyhole}
Let $0<r\le1$, $r'>r$, $a$ real, and let $G$ be holomorphic on $U\times\mathbb C$ for an open $U$ containing $\{|q|\le r'\}\cup[r,1]$. Then $I(z)=\int_0^1q^{a+z-1}G(q,z)\,dq$, defined for $\re z>-a$, continues meromorphically to $\mathbb C$ with poles only where $a+z\in-\mathbb N$, and
\[
I(z)=\sum_{l\ge0}g_l(z)\,\frac{r^{a+z+l}}{a+z+l}+\int_r^1q^{a+z-1}G(q,z)\,dq,\qquad g_l(z)=\frac1{2\pi i}\oint_{|\zeta|=r'}\zeta^{-l-1}G(\zeta,z)\,d\zeta .
\]
Where $|a+z+l|\ge\epsilon$ for every $l\in\mathbb N$,
\[
|I(z)|\le\frac{r'}{\epsilon(r'-r)}\,r^{a+\re z}\max_{|\zeta|=r'}|G(\zeta,z)|+\int_r^1q^{a+\re z-1}|G(q,z)|\,dq .
\]
If $0<q_0\le1$ and $U$ contains the closed disc $D$ with diameter $[-q_0,1]$, then off the points with $a+z\in\mathbb Z$
\[
I(z)=\frac1{e^{2\pi i(a+z)}-1}\oint_{\partial D}q^{a+z-1}G(q,z)\,dq,
\]
the circle run once counterclockwise from $1$, with $q^{a+z-1}$ continued along it from its positive value at $1$.
\end{lemma}

% Reference 20: the sector reading before def:keyhole; source line 5560
\begin{definition}[Falling and base variables; keyholes]\label{def:falling}
For a node of width one and an order $\pi$ of its bunch $B$ let $q=(q_k)_k$ be its sector coordinates and $U_k$ its tails, and let $k_{\cB|B,\pi}$, $\psi_{\cB|B,\pi}$, $k_{\cB/R,\pi}$, $\psi_{\cB/R,\pi}$ be the sector letters of the two tables of the bunch family (Theorem~\ref{thm:exist}(ii)). Along the level $d$ the wall of the bunch family at $U_k$ is $\kappa_k+\operatorname{sl}_kb^d$ at the share $b^d$, with its \emph{slope} $\operatorname{sl}_k\in\{-1,0,1\}$ (Lemma~\ref{lem:slopesrange}). On the sector of $\pi$ the base form is
\[
\phi_\pi=\prod_kq_k^{\kappa_k-1}\,k_{\cB/R,\pi}^{-(s-\kappa_B)}\,k_{\cB|B,\pi}^{-\kappa_B}\,\psi_{\cB|B,\pi}^{-t}\,dq,
\]
and the ratio is
\[
\varrho_\pi=\frac{k_{\cB/R,\pi}\,\psi_{\cB/R,\pi}}{k_{\cB|B,\pi}\,\psi_{\cB|B,\pi}}\prod_kq_k^{-\operatorname{sl}_k},
\]
continued to the complex points of the coordinates at which the sector letters do not vanish; so the continued period of the bunch family at the share $b^d$ is $\sum_\pi\int_{[0,1]^{|B|-1}}\varrho_\pi^{-b^d}\,\phi_\pi$. The $q_k$ with $\operatorname{sl}_k=1$ are the \emph{falling variables} of $\pi$, whose walls fall as $\re b^d$ decreases, and the others its \emph{base variables}, \emph{flat} when $\operatorname{sl}_k=0$ and \emph{rising} when $\operatorname{sl}_k=-1$. Where some flat variables vanish the sector reaches a \emph{flat face} of the shape, on which every wall crossed has slope $0$ and $\varrho_\pi$ has the finite limit given by the leading parts of the letters; where some falling variable vanishes it reaches a \emph{falling face}. For $0<r\le b\le1$ the \emph{keyhole} and its \emph{outline} are $Y_r[0,b]=\{q\in\mathbb C:|q|\le r\}\cup[r,b]$ and $\mathcal O_r[0,b]=\{|q|=r\}\cup[r,b]$, with $Y_r=Y_r[0,1]$ and $\mathcal O_r=\mathcal O_r[0,1]$.
\end{definition}

% Reference 20: def:keyhole; source line 5562
\begin{definition}[Keyhole continuations; ceiling]\label{def:keyhole}
A \emph{keyhole continuation} of the bunch family of a node of width one assigns to every order $\pi$ of $B$ finitely many closed boxes $\Omega$ with disjoint interiors covering the cube of its base variables, and to every such box and falling variable $q_k$ a radius $r_{\pi,\Omega,k}\in(0,1]$, such that none of the four sector letters vanishes on the set $Y=\Omega\times\prod_kY_{r_{\pi,\Omega,k}}$, the product over the falling variables. The boxes may also divide the sides of the falling variables: they then cover the whole cube of the order, a box being the product of its part in the base variables and of its sides in the falling variables, with radii of its own; a falling variable with the side $[0,b_k]$ carries a radius $r\le b_k$ and the sets $Y_r[0,b_k]$ and $\mathcal O_r[0,b_k]$, and one with the side $[a_k,b_k]$, $a_k>0$, is kept real, its sets $Y$ and $\mathcal O$ being $[a_k,b_k]$. A box, divided or not, is written $\Omega$ with its radii and its sets $Y$ and $\mathcal O=\Omega\times\prod_k\mathcal O_{r_{\pi,\Omega,k}}$, where $\Omega$ stands for its part in the variables without keyholes and the products run over the falling variables with keyholes. The \emph{ceiling} $H_{\mathcal O}$ of the continuation is the largest value of $|\varrho_\pi|$ on the sets $\mathcal O$ of its boxes, an order without falling variables contributing the largest value of $\varrho_\pi$ on its cube.
\end{definition}

The sector letters are positive on the closed cube, with positive constant terms (Theorem~\ref{thm:exist}(ii)), so small radii give keyhole continuations, and every bunch has one of finite ceiling. Whether a letter vanishes on a keyhole is decided on its outline and at one point.

% Reference 20: lem:keyholezeros; source line 5568
\begin{lemma}[Zeros on a keyhole are decided on its outline and at one point]\label{lem:keyholezeros}
Let $\Omega$ be a connected set, $J$ a finite set, $0<r_j\le b_j\le1$ for $j\in J$, and $L$ a continuous function on $\Omega\times\mathbb C^J$, a polynomial in $q=(q_j)_{j\in J}$. If $L$ has no zero on $\Omega\times\prod_j\mathcal O_{r_j}[0,b_j]$, and for some $v_0\in\Omega$ and each $j$ the polynomial in $q_j$ obtained from $L(v_0,\cdot)$ by putting $q_k=r_k$ for $k\neq j$ has no zero in the disc $|q_j|<r_j$, then $L$ has no zero on $\Omega\times\prod_jY_{r_j}[0,b_j]$.
\end{lemma}

So a box and its radii keep a letter from vanishing on the set $Y$ of the box when the letter has no zero on its set $\mathcal O$, where the ceiling is read, and, at one point of $\Omega$ with the other falling coordinates at their radii, none in the disc of each keyhole; the falling variables kept real count among the coordinates of $\Omega$.

% Reference 20: thm:heightbound; source line 5577
\begin{theorem}[No reach exceeds the ceiling]\label{thm:heightbound}
For a node of width one, the reach of its bunch is at most the ceiling of every keyhole continuation: $\inf\varrho_B\le\operatorname{reach}(B)\le H_{\mathcal O}$ for the ceiling $H_{\mathcal O}$ of each. When $\varrho_B$ is bounded, every keyhole continuation has the ceiling $\sup\varrho_B$, and $\operatorname{reach}(B)=\sup\varrho_B$.
\end{theorem}

The letters of the contraction enter the keyhole continuations because the family is continued at points off the integers, where $\varrho_\pi^{-b^d}$ branches at their zeros; at the integers $-k$, $k\in\mathbb N$, when $s-\kappa_B$ is zero or a negative integer, only the letters of the restriction $\cB|B$ matter. When a zero of a letter of the contraction is met by the keyholes that would bring $|\varrho_\pi|$ down, the ceiling can exceed the reach. From below, the reach is read at the integers, where $\psi_{\cB/R,\pi}$ does not enter, nor $k_{\cB/R,\pi}$ when $s-\kappa_B$ is an integer, and a continuation whose largest values sit at simple tops that do not cancel cannot be lowered.

% Reference 20: def:intkeyhole; source line 5598
\begin{definition}[Integer keyhole continuations; simple tops]\label{def:intkeyhole}
For a node of width one let $\nu=s-\kappa_B$. An \emph{integer keyhole continuation} is defined as in Definition~\ref{def:keyhole} with only $k_{\cB|B,\pi}$ and $\psi_{\cB|B,\pi}$, and $k_{\cB/R,\pi}$ when $\nu$ is not an integer, required not to vanish on the sets $Y$ of its boxes; its ceiling is defined in the same way. For an order $\pi$ and a box $\Omega$ write $J$ for its falling variables with keyholes and $r_j$ for their radii, and let $C_j$ be the circle $|q_j|=r_j$ run once counterclockwise from $r_j$. The \emph{keyhole contour} of the continuation is the formal sum over the orders and their boxes of
\[
\Omega\times\prod_{j\in J}\Big(\frac{C_j}{e^{2\pi i\kappa_j}-1}+[r_j,1]\Big),
\]
the segment being $[r_j,b_j]$ on a divided side $[0,b_j]$, a falling variable kept real on $[a_j,b_j]$ entering with that segment, and an order without falling variables contributing its cube; on the keyhole contour of $\pi$ it carries the form $\phi_\pi$, the powers continued along the circles from their positive values at $r_j$. A \emph{top} is a point of the keyhole contour at which $|\varrho_\pi|$ equals the ceiling. A top $p$ is \emph{simple} when it lies on no face of its box other than faces $q_j=0$ of flat variables, its falling coordinates lie on the circles away from $r_j$ or inside the segments, and, with $\mathbf x$ the coordinates that vanish at $p$ and $\mathbf y$ real coordinates of the keyhole contour for the others, vanishing at $p$ (the angle on a circle),
\[
\log\frac{\varrho_\pi}{\varrho_\pi(p)}=-\ip{\boldsymbol\lambda}{\mathbf x}-\tfrac12\,\mathbf y^{\mathsf T}\mathbf Q\,\mathbf y+O\big(|\mathbf x|^2+|\mathbf x|\,|\mathbf y|+|\mathbf y|^3\big)
\]
on the keyhole contour near $p$, with $\re\lambda_j>0$ for every $j$ and $\re\mathbf Q$ positive definite. The \emph{exponent} of $p$ is $\operatorname{ex}(p)=\sum_j\kappa_j+\frac12\dim\mathbf y$, the sum over the coordinates $\mathbf x$, and its \emph{coefficient} is
\[
A(p)=\phi_p\prod_j\Gamma(\kappa_j)\,\lambda_j^{-\kappa_j}\,(2\pi)^{\dim\mathbf y/2}(\det\mathbf Q)^{-1/2},
\]
$\phi_p$ the value at $p$ of the density of $\phi_\pi$ with the coefficient $1/(e^{2\pi i\kappa}-1)$ of each circle of the keyhole contour through $p$, $\kappa$ the wall of its falling variable, against $\prod_jx_j^{\kappa_j-1}\,d\mathbf x\,d\mathbf y$, the powers and the root continued from real $\boldsymbol\lambda$ and real $\mathbf Q$ along segments.
\end{definition}

The coordinates $\mathbf x$ belong to flat variables, whose walls $\kappa_j$ are positive, so the Gammas are finite.

% Reference 20: thm:simpletop; source line 5612
\begin{theorem}[A simple top is the growth at the integers]\label{thm:simpletop}
For a node of width one let the walls $\kappa_j$ of the falling variables of every order be off the integers, and let $\operatorname{growth}_{\mathbb Z}(B)$ be the growth at the integers (Definition~\ref{def:growth}).
\begin{enumerate}[label=(\roman*)]
\item For every integer keyhole continuation and every integer $k\ge\max(0,s-\kappa_B)$, the continued period at $b^d=-k$ is the integral of $\varrho_\pi^k\phi_\pi$ over its keyhole contour. So its modulus is at most $CH_{\mathcal O}^k$, $H_{\mathcal O}$ the ceiling, and $\operatorname{growth}_{\mathbb Z}(B)\le H_{\mathcal O}$.
\item $\operatorname{reach}(B)\ge\operatorname{growth}_{\mathbb Z}(B)$.
\item Let an integer keyhole continuation of ceiling $H_{\mathcal O}$ have finitely many tops, all simple, let $\operatorname{ex}_0$ be the least of their exponents, and for $|\zeta|=1$ let $A_\zeta$ be the sum of the coefficients $A(p)$ of the tops with $\operatorname{ex}(p)=\operatorname{ex}_0$ and $\varrho_\pi(p)=H_{\mathcal O}\zeta$. Each $A(p)$ is nonzero, and each top is a stationary point of the restriction of $\varrho_\pi$ to the flat face $\mathbf x=0$, the sector itself when there are no coordinates $\mathbf x$. If some $A_\zeta\neq0$, then $\operatorname{growth}_{\mathbb Z}(B)=H_{\mathcal O}\le\operatorname{reach}(B)$: the reach is at least the modulus of a stationary value of the ratio on a flat face.
\item If moreover there are keyhole continuations with ceilings arbitrarily close to $H_{\mathcal O}$, then $\operatorname{reach}(B)=H_{\mathcal O}$. In particular, if a keyhole continuation, with all four letters kept from vanishing, has finitely many tops, all simple, and some $A_\zeta\neq0$, the reach is its ceiling.
\end{enumerate}
\end{theorem}

So $\operatorname{growth}_{\mathbb Z}(B)\le\operatorname{reach}(B)$ always, with equality to the ceiling whenever the tops of an integer keyhole continuation are finitely many, all simple, and do not cancel, and the keyhole continuations come down to it: the reach is then the modulus of a stationary value of the ratio on a flat face, the one at which the best continuation is pinned. The stationary point need not be a positive point of the face.

The upper end needs the letters of the contraction only where they meet the falling faces. Elsewhere a zero of $k_{\cB/R,\pi}$ or $\psi_{\cB/R,\pi}$ is a zero of $\varrho_\pi$, and the law is read near infinity on the fibres $\varrho_\pi=x$ with $|x|$ large, which never pass through it.

% Reference 20: lem:degreeone; source line 5647
\begin{lemma}[A falling variable takes each large value once]\label{lem:degreeone}
Let $g$ be meromorphic on a neighbourhood of the closed disc $|q|\le r$, with a simple pole at $0$ as its only pole there, and $|g|\le M$ on the circle $|q|=r$. Then $g$ takes every value $x$ with $|x|>M$ exactly once in the disc, with $g'\neq0$ there, and maps the open set $\{|q|<r:\ |g(q)|>M\}\cup\{0\}$ biholomorphically onto $\{|x|>M\}\cup\{\infty\}$.
\end{lemma}

% Reference 20: lem:wallcontinue; source line 5654
\begin{lemma}[A wall series continued is a wall series]\label{lem:wallcontinue}
If the sum of a wall series of radius $x_0$ continues analytically along every path in $\{|x|>x_1\}$, with $x_1<x_0$, then the series is a wall series of radius $x_1$.
\end{lemma}

% Reference 20: def:fibrekeyhole; source line 5661
\begin{definition}[Fibre continuations; stationary bound]\label{def:fibrekeyhole}
A \emph{fibre continuation} is defined as in Definition~\ref{def:keyhole}, divided sides allowed, with only $k_{\cB|B,\pi}$ and $\psi_{\cB|B,\pi}$ required not to vanish on the sets $Y$ of its boxes, and $k_{\cB/R,\pi}$ and $\psi_{\cB/R,\pi}$ only at their points with a falling coordinate $0$; its ceiling $H_{\mathcal O}$ is defined in the same way. Its \emph{stationary bound} $S$ is the supremum of $|\varrho_\pi|$ over the points of these sets with no falling coordinate $0$ at which at least two falling coordinates lie inside their discs, $|q_j|<r_j$, and the derivatives of $\varrho_\pi$ in all the falling coordinates inside their discs vanish; $S=0$ when there is no such point.
\end{definition}

With one falling coordinate inside its disc and the others on their outlines, $\varrho_\pi$ is in that coordinate as in Lemma~\ref{lem:degreeone}, with the bound $H_{\mathcal O}$ on the circle, or identically $0$ when a rising base coordinate is $0$; so it has no stationary point there with $|\varrho_\pi|>H_{\mathcal O}$, and the stationary bound concerns two or more falling coordinates at once.

% Reference 20: thm:fibrebound; source line 5667
\begin{theorem}[The contraction enters only on the falling faces]\label{thm:fibrebound}
For a node of width one and every fibre continuation, of ceiling $H_{\mathcal O}$ and stationary bound $S$, $\operatorname{reach}(B)\le\max(H_{\mathcal O},S)$. So, under the hypotheses of Theorem~\ref{thm:simpletop}, if a fibre continuation that is also an integer keyhole continuation has $S\le H_{\mathcal O}$ and finitely many tops, all simple, with some $A_\zeta\neq0$, then $\operatorname{reach}(B)=\operatorname{growth}_{\mathbb Z}(B)=H_{\mathcal O}$: the reach is read at the integers.
\end{theorem}

The letters of the contraction, which Theorem~\ref{thm:heightbound} keeps from vanishing on whole keyholes, matter for the upper bound only on the falling faces, where the coefficient of the pole of the ratio vanishes, and through the stationary points of the ratio in two or more falling variables. The two bounds from above work box by box, so they can share one continuation.

% Reference 20: thm:mixedbound; source line 5696
\begin{theorem}[Keyhole boxes and fibre boxes mix]\label{thm:mixedbound}
Let a continuation assign to every order $\pi$ of the bunch $B$ of a node of width one boxes and radii as in Definition~\ref{def:keyhole}, the sides of the falling variables divided or not, each box either a \emph{keyhole box}, on whose sets $Y$ none of the four sector letters vanishes, or a \emph{fibre box}, on whose sets $Y$ the conditions of Definition~\ref{def:fibrekeyhole} hold. Let $H_{\mathcal O}$ be its ceiling and $S$ the stationary bound of Definition~\ref{def:fibrekeyhole} taken over its fibre boxes alone. Then $\operatorname{reach}(B)\le\max(H_{\mathcal O},S)$.
\end{theorem}

% Reference 20: cor:onefalling; source line 5705
\begin{corollary}[One falling variable needs only the restriction]\label{cor:onefalling}
In an order with at most one falling variable, a box on whose sets $Y$ the letters of the restriction $\cB|B$ do not vanish is a fibre box, and it adds nothing to the stationary bound. So in a continuation of Theorem~\ref{thm:mixedbound} such boxes need only the letters of the restriction, and the reach is at most $\max(H_{\mathcal O},S)$ with $S$ taken over the fibre boxes of the orders with two or more falling variables.
\end{corollary}

So the stationary bound enters only where a letter of the contraction vanishes inside the keyholes of a sector with two or more falling variables. The zeros of the contraction on the falling faces can still keep every such bound above the reach.

These theorems bound the right end of each interval $[\inf\varrho_B,\operatorname{reach}(B)]$ of Theorem~\ref{thm:nodeintervals} from both sides. The reach need not be read on the positive shape: the ratio is real at the real points of a flat face whatever their signs, and a keyhole of a falling variable crosses the negative axis nearer to $0$ than every pole of the ratio on it, so a simple top can sit at a point with a negative coordinate.

\section{One tree}\label{sec:treecoaction}

Read on one tree, the period of a point, built back up, is the periods of the leaves multiplied by the share kernels of the nodes and integrated over the free shares; its $n-1$ integrations are shared out among the bunches and the leaves; at integer data the share kernels bring the Trailmoments of the counting input order by order; and where the free shares close on integer points the period is the sum of the residues along the lattice of the tree.

% Reference 20: def:pieces; source line 4900
\begin{definition}[Share kernels, contents, integration count]\label{def:contents}
Let a molecule of lists with a native gauge be read on a tree of bunches, its Foreman tree or another, with the port of the balance wall run at every node on the native splits (Corollary~\ref{cor:treeword}). The share kernels of the tree (Corollary~\ref{thm:E1}) are what the ports place on the shares: one stacked Gamma $\Gst{k}{1}{\mathbf a}{a}$ for each letter opened at a node, $k$ the number of its patterns, $\mathbf a$ its shares and $a$ its exponent, which in Corollary~\ref{cor:treeword} leave the factors $\big(\Gamma(a_{\cB'})\Gamma(b_{\cB'})\big)^{1-q_{\cB'}}$ and $\prod_j\Gamma(s_j)^{-1}\Gamma(t_j)^{-1}$; each is balanced (Definition~\ref{def:stacked}), its shares adding to its exponent on the fibre. The \emph{free shares} are coordinates of the fibre of the word. Let $\mathcal Z$ be the $\mathbb Q$-algebra generated by the Trailmoments $\Tx{k}{\IZ}{0}$, $k\ge2$, of the counting input; it maps into the content $\mathcal C^{(1)}$ of Chapter~\ref{ch:content}, the algebra of trails of lattice inputs, which also holds the sums along the lattices of paired words (Chapters~\ref{ch:rayarith} and~\ref{ch:content}). The \emph{content} $\mathcal C_M$ of a childcule $M$ of a node, a leaf in particular, is the $\mathcal Z$-module spanned by the values, at the integer exponents and masses at which it converges, of the period of a family on $M$ whose tables are those of $M$, and of its derivatives in the exponents and masses along the balance wall; the content of the molecule is the same with its own tables. The content $\mathcal C_B$ of the bunch $B$ of a node with room $R$ is the same for the families on $B$ whose tables, each with both letters, are the minors $\cB'|B$ and $\cB'/R$ of the tables $\cB'$ at the node: those of the face of $B$ (Theorem~\ref{thm:wallres}), at width one those of the bunch family (Theorem~\ref{thm:twosides}). The content $\mathcal C_R$ of the room is the same for the families on $R$ whose tables are the $\cB'|R$ and $\cB'/B$.

The \emph{integration count} of an integral as written is the number of its integrations, assigned to the written expression by these rules. The period of a family on $k$ marks has the integration count $k-1$, its simplex. A trail $\Ttr{\mathbf a}{d}{\Ilat{(b+\mathbb Z)^d}}{b}$ of a lattice input with integer indices, $b>-1$ (Definition~\ref{def:trail}), has the integration count $a_1+\cdots+a_d$, its total index; the counting trails are the case $b=0$. A Gamma has the integration count one, $\Gamma(a)=\int_0^\infty x^{a-1}e^{-x}\,dx$ being one integration, so a stacked Gamma with $k$ Gammas above and $l$ below has the integration count $k-l$. A derivative in an exponent or a mass adds one, the logarithm of a letter or of a mark. A product has the sum of the integration counts of its factors, and a fibre integral of dimension $D$ (Definition~\ref{def:fibre}) the integration count of its integrand less $D$, each share integrated inverting one Mellin transform. A real number has \emph{integration count at most $w$} when it is a rational combination of values, at rational data with unit weights, of integrals of integration count at most $w$.
\end{definition}

The openings keep the integration count. The power $(X_1+\cdots+X_k)^{-s}$ is the share integral of $\Gst{k}{1}{\mathbf a}{s}\prod_iX_i^{-a_i}$, of integration count $(k-1)-(k-1)=0$ (Corollary~\ref{thm:E1}), and $e^{-y}$ is the integral of $\Gamma(b)y^{-b}$ over one share, of integration count $1-1=0$. A Beta corner $\Gst21{a,\ b}{a+b}$ has the integration count one, as the period of the list on two marks whose bases are the two singletons; the stacked Gamma $\Gst{|B|}{1}{\mvec}{m_B}$ of the shape integral has the integration count $|B|-1$, its simplex (Theorem~\ref{thm:dirichlet}); the interest $J_F=\Gamma(\kappa_{\cE})P_F$ has the integration count $n$, its orthant (Theorem~\ref{thm:projective}). The atom formula of a family on $n$ marks (Corollary~\ref{cor:atomwall}) has the integration count $\sum(|\cB_j|-1)$ of its stacked Gammas, over the active letters, less its dimension $\sum(|\cB_j|-1)-(n-1)$: $n-1$, the integration count of the period. Taking the residue at a pole of one Gamma removes that Gamma and adds one pin (Theorem~\ref{thm:closing}), so every term of a closing has the integration count of the integral closed.

% Reference 20: thm:onetree; source line 4910 (a molecule of one list in the chain clause of (iii)).
% Proved in the last part: thm:onetree.
\begin{theorem}[One tree suffices]\label{thm:onetree}
Let $F$ be a molecule of lists on $n$ marks with real exponents $s_j,t_j\ge0$ and masses $\mvec$ on the balance wall, with $P_F(\mvec;\mathbf s,\mathbf t)$ convergent, and let its Foreman tree be read at its exponent pairs, at the centre or at a lean (Proposition~\ref{prop:readingpoint}, Theorem~\ref{thm:threetrees}). Run the port of the balance wall at every node on the native splits of the point, as in Corollary~\ref{cor:treeword}. Parts (i), (ii), (iv) and (v) hold as well for every tree of bunches in place of the Foreman tree.
\begin{enumerate}[label=(\roman*)]
\item \emph{The word starts at the leaves.} $P_F(\mvec;\mathbf s,\mathbf t)$ is the word of the tree (Corollary~\ref{cor:treeword}): the periods of the leaves, each read at its shares and each entering once on its descent, multiplied by the share kernels and integrated over the free shares. At a node the words of the childcules multiply, and along a chain, its leaf opened by the atom formula, the word is one Gramma on the marks and the letters (Theorem~\ref{thm:buildup}).
\item \emph{The integrations are counted.} The word has the integration count $n-1$, that of the period, and at rational data $P_F(\mvec;\mathbf s,\mathbf t)$ has integration count at most $n-1$. At a node with bunch $B$ and $q$ childcules the integrations of the simplex are the $|B|+q-1$ pins of the node, $|B|-1$ for the shape of the bunch and $q$ for the totals of the bunch and of the childcules, $q+1$ numbers that add to one; in a leaf $L$ they are the $|L|-1$ integrations of its period.
\item \emph{The leaf bounds the bunches.} Along every descent of the Foreman tree, at the centre and at every lean, the copies of every bunch have at most as many marks as the copies of the leaf (Theorem~\ref{thm:leafbounds}). When some active list is connected, as for a molecule of one list, the tree is a chain and $|B_0|\le\cdots\le|B_{d-1}|\le|L|$ (Corollary~\ref{cor:leafchain}): the shape of every bunch has at most as many integrations as the leaf, and its node adds one, the total of the bunch against its room.
\item \emph{At integer data the share kernels bring $\mathcal Z$.} Let the exponents and masses be integers and let $x$ be a point of the fibre at which every share is an integer. The pins and the share sums have integer coefficients, so the directions of the fibre have a rational basis; write the fibre near $x$ as $x+V\mathbf v$ with $V$ rational. Near $x$ the share kernels are
\[
c\,\prod_i\ell_i(\mathbf v)^{-\epsilon_i}\exp h(\mathbf v),
\]
with $c$ rational, $\ell_i(\mathbf v)$ the deviation from its value at $x$ of each argument of the share kernels that is a nonpositive integer at $x$, $\epsilon_i$ equal to $1$ for a Gamma above and to $-1$ for one below, and $h$ a power series without constant term whose part of degree $k$ has its coefficients in $\mathbb Q+\mathbb Q\,\Tx{k}{\IZ}{0}$ for $k\ge2$ and in $\mathbb Q$ for $k=1$: the terms of order $k$ of the share kernels carry total index at most $k$. Where the periods of the leaves converge at $x$, every residue of the word at $x$, taken in turn along rational directions spanning the fibre, is therefore a finite $\mathcal Z$-combination of products, over the leaves, of elements of their contents.
\item \emph{Summed along the tree.} A free share closed on the poles of the share kernels, where they alone have poles on the way and the closing converges, writes the word as a sum over a lattice of terms of its own integration count, each a fibre integral with one more pin (Theorem~\ref{thm:closing}); where every free share closes so, on points at which every share is an integer and the periods of the leaves converge, the period at integer data is the sum along that lattice of the residues of (iv), finite $\mathcal Z$-combinations of products of the contents of the leaves. Across a bunch of width one, grouping its patterns by level puts the period of the bunch family in the word in place of the bunch (Theorem~\ref{thm:twosides}); for one table the even share then closes on its lattice whenever $\sup\varrho_B<\inf\varrho_R$ or $\inf\varrho_B>\sup\varrho_R$, and at integer data its residues are finite $\mathcal Z$-combinations of products of an element of $\mathcal C_B$ and one of $\mathcal C_R$ (Corollary~\ref{cor:closetwosides}).
\end{enumerate}
\end{theorem}

Whatever the order of the even letters, with $s,t>0$, the even share of a node of width one closes on the lattices of continued periods of the bunch family and the room family toward a level at which their growths multiply to less than one (Theorem~\ref{thm:nodeclosing}); a series bunch closes toward the level $d$ against every room with $\inf\varrho_R>\inf\varrho_B$, and a series connection against every room with $\sqrt{\inf\varrho_R}$ above the sum of the square roots of the reaches of its parts, the room's ratio bounded only on the side away from the closing (Theorem~\ref{thm:classcloses}); and against a room bounded below and a bunch bounded below the node closes on a geometric lattice exactly when the intervals of the bunch and of the room are disjoint, at its peaks otherwise (Theorem~\ref{thm:nodeintervals}). Off the integers the level shares of several tables close in the same way across bunches of width one whose values stand in one order, node after node down a tree of splits of width one (Theorem~\ref{thm:jointclosing}, Corollary~\ref{cor:nested}). At nonnegative exponents the period is in every case the sum, along the draws nested on the descents of the tree, of products over its blocks of Dirichlet moments (Theorem~\ref{thm:drawtree}).

The pins need not be unimodular, so a closing can reach points at which some share is not an integer, even at integer data. When the active tables link all the marks the pins reach every integer point of their target (Proposition~\ref{prop:ct-pins}), and the fibre has integral coordinates through an integer point; the shares that a closing leaves after pinning the others at the poles of their Gammas are solved from a square block of the pins, whose determinant can exceed $1$ in modulus: for the list $\binom{[4]}2$ at the data $(\one;1,1)$ (Exercise~\ref{ex:cls-pins}), the fibre has points at which every share off such a block is an integer and some share is $\frac12$. There the share kernels have rational arguments, and their expansion runs in heads and Trailmoments of the lattice inputs of those arguments (Proposition~\ref{prop:headstails}); that is why (iv) and (v) speak of integer points. Where a node has no free share its sum has one term. For a list whose bunch has width one, at a full lean, $(s,0)$ or $(0,t)$, the level shares are pinned and the period is a Beta corner times the periods of the bunch family and the room family (Corollary~\ref{cor:pinnedproduct}); at rational data this is an identity of period data (Theorem~\ref{thm:poolingmove}), at integer data the Beta corner is rational, and with the room a leaf the period is a rational multiple of a product of an element of $\mathcal C_B$ and one of the content of the leaf. When the two hinge values of a bunch of width one agree the shape and the room separate at every pair of exponents (Proposition~\ref{prop:readingsagree}), and at the wall of a set, with the weights factoring on its bottom levels, the residue of the period is the product of the periods of its two faces (Theorem~\ref{thm:wallres}).

The integration count of a period bounds that of its value, and the theorem is stated for lists because the weights of a table can carry integrations of their own, which the integration count of a period leaves out. With unit weights and rational data nothing is carried in. On an atom the tree is the atom and the theorem says no more than its own content; its atom formula is a word of the same kind whose leaves are single marks, and the argument of (iv) puts each of its residues at integer points in $\mathcal Z$. The lattice can make what no bunch or leaf shows: a tree whose bunches and leaves have their contents in $\mathcal Z$ can sum, along the lattice of one free share, to a Trailmoment of a lattice input with a frequency (Definition~\ref{def:inputs}) that none of them carries. The finite form of the theorem, the period in the $\mathcal Z$-algebra generated by the contents of the bunches and the leaves, holds where the shares are pinned at an integer point; where they are free the tree gives the sum along its lattice. For a word of one free share made of Gammas alone, as the atom formula leaves it at integer data, that sum lies in the content of Definition~\ref{def:ct-content} when the word is paired (Propositions~\ref{prop:ra-paired} and~\ref{prop:ct-level}), and otherwise it is a combination of Trailmoments of live inputs (Theorem~\ref{thm:ra-oneshare}).

Theorem~\ref{thm:onetree} is the Foreman procedure read, step by step, as a statement about numbers. \emph{Split first}: a compound is split into its polycule before any molecule is asked whether it answers yes, and at every node the polycule of the childcules is evaluated as the product of its molecules (Theorem~\ref{thm:buildup}(i)). \emph{Read at the table's gauge, then ask}: each molecule is read at its table's gauge, native for the masses it uses and inherited by its children (Proposition~\ref{prop:herecentre}), whether it answers yes is read at the centre (Theorem~\ref{thm:centre}), and a realization that is not atomic is peeled. \emph{Peel small}: the peel is taken at a class with the fewest marks, it splits a class only across twin childcules, each part inside one class (Theorem~\ref{thm:roomkeeps}), so the leaf with its copies holds whole classes of every node above it and is at least as large as every bunch with its copies (Theorem~\ref{thm:leafbounds}). \emph{The share kernels reapply}: what a node adds is its port, stacked Gammas of the shares of its bunch and the pins of its shape and its totals; near integer data a balanced stacked Gamma brings $\mathcal Z$ order by order in total index, and in the draw series each node contributes its own Dirichlet moments to every product over the blocks. \emph{One tree}: every tree of bunches computes the same period and shares out the same $n-1$ integrations, and the Foreman tree is one in which the copies of no bunch have more marks than the copies of the leaf, and when some active list is connected no bunch has more marks than the leaf. \emph{Pinned shares}: where the shares of a node are pinned, as for one list at a full lean across a bunch of width one, the node contributes an identity of period data (Theorem~\ref{thm:poolingmove}).

Chapter~\ref{ch:motive} lifts the statements of this section from periods to motivic periods: at integer data the left factors of the coaction of the motivic period of a family are read along every tree of bunches, the Foreman tree among them, through the steps attached to its vertices (Theorem~\ref{thm:mot-smallsteps}); at its nodes, at any width, through the joint spectrum of the bunch (Theorem~\ref{thm:mot-node}); and where the shares of a node are pinned the identities of period data above are identities of motivic periods (Theorem~\ref{thm:mot-pinned}, Corollary~\ref{cor:mot-dualfull}).



\part{Construction and composition}\label{part:construction}

\chapter{Gluing and what crosses it}\label{ch:glue}

Every table read at a bunch is the gluing of what its bunch and its room carry (Proposition~\ref{prop:interchange}), and at a bunch of width one the gluing is along one mark, the hinge (Theorem~\ref{thm:widthone}). Gluing is an operation with its own calculus. A table is built from cells, tables with one marked hinge, substituted into the marks of a core; the bases, the weights, the multivectors, the letters and the walls of the result are read off the cells and the core; the blocks of Theorem~\ref{thm:block} and the series connections of Section~\ref{sec:onescale} are the cells and cores of the simplest kinds; two cell families glued along their hinges have hinge walls that add to $2\sigma$ on the balance wall; and every connected table whose weights factor at its bunches of width one, every list and every determinantal table among them, is a tree of such substitutions whose nodes are prime tables and series and parallel nodes. Cells with several hinges splice in pairs along any graph of pairs, closed paths included, and stay tables, their channels composing multilinearly. What a room receives from its bunch, at any width, for any family and any multivector, is the joint law of finitely many bunch letters: as many as the ranks of the level weights, which for a multivector are bounded by the squares of its interface ranks and for a determinantal table by $\binom{c(B)}{i}^2$ at the level $d(B)+i$, whatever the numbers of marks; for a determinantal table the bunch crosses its gluing as one positive definite matrix of size $c(B)$, whose determinant at width one is the hinge value $\gamma^K_B$, and one scalar.

\section{Gluing along a mark}\label{sec:gluemark}

% Reference 20: def:substitution; source line 5731
\begin{definition}[Cells, cores, substitution]\label{def:substitution}
A \emph{cell} is a table $\mathcal A$ on $B\sqcup\{g\}$ in which the mark $g$, its \emph{hinge}, is neither a loop nor a coloop; it may carry two systems of positive weights on its bases, $\theta^K$ for its letter $K_{\mathcal A}$ and $\theta^\Psi$ for its letter $\Psi_{\mathcal A}$ (Definition~\ref{def:family}), and it is \emph{plain} when they agree, as for one table. Its \emph{hinge minors} are the deletion $\mathcal A\setminus g$ and the contraction $\mathcal A/g$, tables on $B$ of ranks $r_{\mathcal A}$ and $r_{\mathcal A}-1$ with the inherited weights (Definition~\ref{def:delcon}), so that
\[
K_{\mathcal A}=K_{\mathcal A\setminus g}+\alpha_gK_{\mathcal A/g},\qquad \Psi_{\mathcal A}=\alpha_g\Psi_{\mathcal A\setminus g}+\Psi_{\mathcal A/g};
\]
its hinge values and its ratio, as for a bunch of width one (Definition~\ref{def:readings}), are
\[
\gamma^K_{\mathcal A}=\frac{K_{\mathcal A\setminus g}}{K_{\mathcal A/g}},\qquad \gamma^\Psi_{\mathcal A}=\frac{\Psi_{\mathcal A/g}}{\Psi_{\mathcal A\setminus g}},\qquad \varrho_{\mathcal A}=\frac{\gamma^\Psi_{\mathcal A}}{\gamma^K_{\mathcal A}}=\frac{\mathbf E_{\mathcal A/g}}{\mathbf E_{\mathcal A\setminus g}},
\]
functions of $\alpha_B$ homogeneous of the degrees $1$, $1$ and $0$. The cell of a bunch $\cB^+_B$ of Theorem~\ref{thm:widthone} is a cell, with the hinge minors $\cB|B$ and $\cB/R$, the hinge values $\gamma^K_B$, $\gamma^\Psi_B$ and the ratio $\varrho_B$ of Definition~\ref{def:readings}. A \emph{core} is a table $\mathcal H$ on marks $h_1,\dots,h_q$, none of them a loop or a coloop, plain or with two systems of weights as a cell. For cells $\mathcal A_i$ on $B_i\sqcup\{g_i\}$, the $B_i$ disjoint, the \emph{substitution} $\mathcal H\langle\mathcal A_1,\dots,\mathcal A_q\rangle$ is the family of sets $\bigsqcup_iX_i\subseteq\bigsqcup_iB_i$ for which some basis $T$ of $\mathcal H$ has $X_i\in\mathcal A_i\setminus g_i$ when $h_i\in T$ and $X_i\in\mathcal A_i/g_i$ when $h_i\notin T$, with the weight, in each letter,
\[
\theta_{\sqcup_iX_i}=\theta^{\mathcal H}_T\prod_{h_i\in T}\theta^{\mathcal A_i}_{X_i}\prod_{h_i\notin T}\theta^{\mathcal A_i}_{X_i+g_i},
\]
every $\theta$ read as the weight $\theta^K$ of the letter $K$ throughout, or as $\theta^\Psi$ throughout; $T$ is determined by the sizes, $h_i\in T$ exactly when $|X_i|=r_{\mathcal A_i}$. A mark $h$ is \emph{left in place} by the cell $\binom{\{h,g\}}1$, the list on $\{h,g\}$ with the bases $\{h\}$ and $\{g\}$, whose hinge minors have the single bases $\{h\}$ and $\varnothing$ and whose hinge values are both $\alpha_h$. A plain cell is a \emph{dual pair} when its contraction $\mathcal A/g$ is the dual table of its deletion $\mathcal A\setminus g$, weights included: the bases of $\mathcal A/g$ are the $B\setminus X$, $X$ a basis of $\mathcal A\setminus g$, with $\theta^{\mathcal A}_{(B\setminus X)+g}=\theta^{\mathcal A}_X$; a mark left in place is one. One cell $\mathcal A$ at one mark $h$ of $\mathcal H$, the other marks left in place, is the \emph{gluing of $\mathcal A$ and $\mathcal H$ along $h$}: the gluing (Definition~\ref{def:communicatroid}) of the communicatroid of $B$ whose children are $\theta^{\mathcal A}_X$ times $\mathcal H/h$ at the bases $X$ of $\mathcal A\setminus g$ and $\theta^{\mathcal A}_{X+g}$ times $\mathcal H\setminus h$ at the bases $X$ of $\mathcal A/g$. A family is \emph{substituted} table by table, $(\cB_j,\theta_j)=\mathcal H_j\langle\mathcal A_{1j},\dots,\mathcal A_{qj}\rangle$, the cells of all tables at the mark $h_i$ sharing the marks $B_i\sqcup\{g_i\}$; its masses are on $\cE=\bigsqcup_iB_i$.
\end{definition}

% Reference 20: lem:gluemark; source line 5739
\begin{lemma}[Gluing along a mark gives a table]\label{lem:gluemark}
Let $\mathcal A$ be a cell on $B\sqcup\{g\}$ and $\mathcal A'$ a cell on $R\sqcup\{g'\}$. The sets $X\sqcup Y$ with $X\in\mathcal A\setminus g$ and $Y\in\mathcal A'/g'$, or $X\in\mathcal A/g$ and $Y\in\mathcal A'\setminus g'$, are the bases of a table on $B\sqcup R$ of rank $r_{\mathcal A}+r_{\mathcal A'}-1$, in which $B$ is a bunch of width one whose cell of the bunch and cell of the room (Theorem~\ref{thm:widthone}) are $\mathcal A$ and $\mathcal A'$, the hinges read as the one mark $g$.
\end{lemma}

% Reference 20: prop:substitution; source line 5746
\begin{proposition}[Substitution]\label{prop:substitution}
\begin{enumerate}[label=(\roman*)]
\item $\mathcal H\langle\mathcal A_1,\dots,\mathcal A_q\rangle$ is a table of rank $r_{\mathcal H}+\sum_i(r_{\mathcal A_i}-1)$; it is the gluing of the cells to $\mathcal H$ along its marks one at a time, in any order; each $B_i$ is a bunch of it of width one whose cell of the bunch is $\mathcal A_i$, its weights factoring across $B_i$ and its room.
\item Substituting into a mark of a cell is substituting the substituted cell: for a mark $x$ of $B_1$ and a cell $\mathcal A'$ at $x$,
\[
\big(\mathcal H\langle\mathcal A_1,\dots,\mathcal A_q\rangle\big)\langle\mathcal A'\text{ at }x\rangle=\mathcal H\big\langle\mathcal A_1\langle\mathcal A'\text{ at }x\rangle,\mathcal A_2,\dots,\mathcal A_q\big\rangle .
\]
\item Conversely, a table with a bunch $B$ of width one, the weights factoring across $B$ and $R$, is the gluing of its cell of the bunch $\cB^+_B$ and its cell of the room $\cB^+_R$ along $g$ (Theorem~\ref{thm:widthone}): every gluing along a mark is that decomposition read backwards, and every such decomposition is a gluing.
\item If $\mathcal H\langle\mathcal A_1,\dots,\mathcal A_q\rangle$ is connected, so are $\mathcal H$ and every cell. A connected table with a bunch $B$ of width one, $|B|,|R|\ge2$, the weights factoring across $B$ and $R$, is the gluing along a mark of two connected tables on fewer marks.
\end{enumerate}
\end{proposition}

% Reference 20: thm:substletters; source line 5761
\begin{theorem}[The letters of a substitution]\label{thm:substletters}
Let $K$ and $\Psi$ be the letters of $\mathcal H\langle\mathcal A_1,\dots,\mathcal A_q\rangle$, each with its own weights, and read the hinge minors and the hinge values of each cell on its own marks. Then
\begin{gather*}
K(\alpha)=\prod_iK_{\mathcal A_i/g_i}(\alpha_{B_i})\ K_{\mathcal H}\big(\gamma^K_{\mathcal A_1}(\alpha_{B_1}),\dots,\gamma^K_{\mathcal A_q}(\alpha_{B_q})\big),\\
\Psi(\alpha)=\prod_i\Psi_{\mathcal A_i\setminus g_i}(\alpha_{B_i})\ \Psi_{\mathcal H}\big(\gamma^\Psi_{\mathcal A_1}(\alpha_{B_1}),\dots,\gamma^\Psi_{\mathcal A_q}(\alpha_{B_q})\big).
\end{gather*}
For one cell, the other marks left in place, this is Theorem~\ref{thm:widthone}.
\end{theorem}

The core reads each of its cells as one mark, at the hinge value of the cell by the letter: in $K$ the mark $h_i$ stands at $\gamma^K_{\mathcal A_i}$ and in $\Psi$ at $\gamma^\Psi_{\mathcal A_i}$, and what is left over is a product of one letter of a hinge minor of every cell. Nothing in the core sees more of a cell than its two hinge values.

A family can be generated rather than found: seed tables assembled under a set of bijections give atoms (Proposition~\ref{prop:assembly}), and a mark that is neither a loop nor a coloop, replaced by a block of new marks, gives a bunch of width one, which the room reads as the mark it replaced.

% Reference 20: thm:block; source line 3686 (one mark).
% Proved in the last part: thm:block.
\begin{theorem}[Replacement of marks by blocks]\label{thm:block}
Let $(\cB,\theta)$ be a table of rank $r$ on $\cE=\cE_1\sqcup\cE_2$, and replace every mark $h$ by a \emph{block} $X_h$ of $q_h\ge1$ new marks, by the \emph{first rule} at the marks of $\cE_1$ and by the \emph{complementary rule}, the first rule read on the dual table, at those of $\cE_2$: a basis $T$ gives the sets that contain one mark of $X_h$ for each $h\in T\cap\cE_1$, none for $h\in\cE_1\setminus T$, all of $X_h$ for $h\in T\cap\cE_2$ and all of $X_h$ but one for $h\in\cE_2\setminus T$, each with the weight of $T$. They are the bases of a table on $\bigsqcup_hX_h$ of rank $r+\sum_{h\in\cE_2}(q_h-1)$; a block of one mark leaves its mark in place, and when no mark is a loop or a coloop the table is the substitution into $(\cB,\theta)$ of the \emph{parallel blocks} $\binom{X_h\sqcup\{g_h\}}1$ at $\cE_1$ and the \emph{series blocks} $\binom{X_h\sqcup\{g_h\}}{q_h}$ at $\cE_2$ (Definition~\ref{def:substitution}). Its Foreman polynomials are, each $e_k$ read on the marks of its block,
\begin{gather*}
K'=\prod_{h\in\cE_2}e_{q_h-1}\cdot K\Big(\big(e_1\big)_{h\in\cE_1},\big(\tfrac{e_{q_h}}{e_{q_h-1}}\big)_{h\in\cE_2}\Big),\qquad \Psi'=\prod_{h\in\cE_1}e_{q_h-1}\cdot\Psi\Big(\big(\tfrac{e_{q_h}}{e_{q_h-1}}\big)_{h\in\cE_1},\big(e_1\big)_{h\in\cE_2}\Big);
\end{gather*}
for one block $X=\{x_1,\dots,x_q\}$ at the mark $h$, the other marks left in place, the first rule gives $K'=K_{\setminus h}+e_1(\alpha_X)\,K_{/h}$ and $\Psi'=e_q(\alpha_X)\,\Psi_{\setminus h}+e_{q-1}(\alpha_X)\,\Psi_{/h}$, and the complementary rule gives $K'=e_q(\alpha_X)\,K_{/h}+e_{q-1}(\alpha_X)\,K_{\setminus h}$ and $\Psi'=\Psi_{/h}+e_1(\alpha_X)\,\Psi_{\setminus h}$ (Definition~\ref{def:delcon}). In scale and shape, $\alpha_X=\rho u$ with $u\in\Delta_X$: $e_1(\alpha_X)=\rho$, $e_q(\alpha_X)=\rho^qe_q(u)$, $e_{q-1}(\alpha_X)=\rho^{q-1}e_{q-1}(u)$, and $\alpha_X^{\mvec|_X-\one}d\alpha_X=\rho^{m_X-1}d\rho\;u^{\mvec|_X-\one}d\omega_X(u)$ (Lemma~\ref{lem:scaleshape}).
\end{theorem}

% Reference 20: thm:blockcounts; source line 3693 (uniform replacement; the atom at the centre).
% Proved in the last part: thm:blockcounts.
\begin{theorem}[Counts, walls and atoms under replacement]\label{thm:blockcounts}
Let $(\cB,\theta)$ be a table of rank $r$ on $\cE$, $n=|\cE|$, and replace every mark $h$ by a block $X_h$ of $q_h\ge1$ new marks as in Theorem~\ref{thm:block}, by the first rule at the marks of $\cE_1$ and by the complementary rule at those of $\cE_2$, $\cE=\cE_1\sqcup\cE_2$.
\begin{enumerate}[label=(\roman*)]
\item The marks become $\sum_hq_h$; the rank becomes $r+\sum_{h\in\cE_2}(q_h-1)$; the total weight becomes
\[
\sum_{T\in\cB}\theta_T\prod_{h\in T\cap\cE_1}q_h\prod_{h\in\cE_2\setminus T}q_h=\prod_{h\in\cE_2}q_h\cdot K\Big((q_h)_{h\in\cE_1},(q_h^{-1})_{h\in\cE_2}\Big),
\]
which is $K((q_h)_h)$ under the first rule alone and $\Psi((q_h)_h)$ under the complementary rule alone.
\item For a set $Y$ of new marks let $\lceil Y\rceil=\{h:\ X_h\cap Y\neq\varnothing\}$ and $\lfloor Y\rfloor=\{h:\ X_h\subseteq Y\}$, and let $r'$, $d'$, $d^{\ast\prime}$, $c'$ be the counts of the new table (Definition~\ref{def:skeleton}). Then
\begin{gather*}
d'(Y)=\sum_{h\in\cE_2}\max\big(|X_h\cap Y|-1,0\big)+d\big((\lfloor Y\rfloor\cap\cE_1)\sqcup(\lceil Y\rceil\cap\cE_2)\big),\\
d^{\ast\prime}(Y)=\sum_{h\in\cE_1}\max\big(|X_h\cap Y|-1,0\big)+d^\ast\big((\lceil Y\rceil\cap\cE_1)\sqcup(\lfloor Y\rfloor\cap\cE_2)\big),
\end{gather*}
and $c'(Y)=|Y|-d'(Y)-d^{\ast\prime}(Y)$. Under the first rule alone $r'(Y)=r(\lceil Y\rceil)$, $d'(Y)=r-r(\cE\setminus\lfloor Y\rfloor)$ and $c'(Y)=r(\lceil Y\rceil)+r(\cE\setminus\lfloor Y\rfloor)-r$; under the complementary rule alone $c'(Y)=|\lceil Y\rceil|-|\lfloor Y\rfloor|+r(\cE\setminus\lceil Y\rceil)+r(\lfloor Y\rfloor)-r$. If $Y=\bigcup_{h\in S}X_h$, then $c'(Y)=c(S)$; when the seed table is connected on at least two marks, a proper nonempty part of one block has width one.
\item Let a family of seed tables with positive exponent pairs be connected on at least two marks and carried to itself by a set $H$ of bijections as in Proposition~\ref{prop:assembly}, and replace every mark by a block of the same $q$ marks, all by the first rule or all by the complementary rule. Then the bijections of $H$ lift to the blocks, the permutations inside a block carry every replaced table to itself, the replaced family is connected, and for every $q$ its realization at the table's gauge is an atom.
\end{enumerate}
\end{theorem}

The shape of a block is integrated by Theorem~\ref{thm:pooling}, and its scale is the mark of the seed. A block $X_h$ of Theorem~\ref{thm:block} with $q_h\ge2$ whose mark is neither a loop nor a coloop of the seed is a parallel bunch of width one under the first rule and a series bunch of width one under the complementary rule (Definition~\ref{def:width}).

% Reference 20: cor:blocksclosures; source line 5773
\begin{corollary}[Blocks and closures]\label{cor:blocksclosures}
\begin{enumerate}[label=(\roman*)]
\item The series block on a set $B$ of $\ell$ marks (Theorem~\ref{thm:block}) is the cell $\binom{B\sqcup\{g\}}{\ell}$, whose hinge minors have $K_{\setminus g}=e_\ell(\alpha_B)$, $K_{/g}=e_{\ell-1}(\alpha_B)$, $\Psi_{\setminus g}=1$ and $\Psi_{/g}=e_1(\alpha_B)$; the parallel block is the cell $\binom{B\sqcup\{g\}}1$, with $K_{\setminus g}=e_1(\alpha_B)$, $K_{/g}=1$, $\Psi_{\setminus g}=e_{\ell-1}(\alpha_B)$ and $\Psi_{/g}=e_\ell(\alpha_B)$. Series blocks at the marks $h_i$, $i\in S$, of a core $\mathcal H$ and parallel blocks at the others replace the marks of $\mathcal H$ by blocks as in Theorem~\ref{thm:block}, the series blocks by its complementary rule and the parallel blocks by its first rule, and give the letters of that theorem, each $e_k$ read on the marks of its block: a series block is seen by $\Psi$ through its total and by $K$ through the hinge value $e_\ell/e_{\ell-1}=\big(\sum_x\alpha_x^{-1}\big)^{-1}$, a parallel block the other way round.
\item The \emph{series closure} of cells $\mathcal A_1,\dots,\mathcal A_q$ is the substitution $\binom{\{h_1,\dots,h_q\}}{q-1}\langle\mathcal A_1,\dots,\mathcal A_q\rangle$ into the core of unit weights; its letters are
\[
K=\prod_iK_{\mathcal A_i\setminus g_i}\cdot\sum_i\frac{K_{\mathcal A_i/g_i}}{K_{\mathcal A_i\setminus g_i}},\qquad \Psi=\prod_i\Psi_{\mathcal A_i\setminus g_i}\cdot\sum_i\frac{\Psi_{\mathcal A_i/g_i}}{\Psi_{\mathcal A_i\setminus g_i}},
\]
so that the bunch $B_1\sqcup\dots\sqcup B_{q-1}$ is the series connection of the parts $B_i$ with the tables $\mathcal A_i\setminus g_i$ and $\mathcal A_i/g_i$ (Proposition~\ref{prop:seriesconnection}, with $C=C'=1$). The \emph{parallel closure} $\binom{\{h_1,\dots,h_q\}}1\langle\mathcal A_1,\dots,\mathcal A_q\rangle$ has the letters $K=\prod_iK_{\mathcal A_i/g_i}\cdot\sum_iK_{\mathcal A_i\setminus g_i}/K_{\mathcal A_i/g_i}$ and $\Psi=\prod_i\Psi_{\mathcal A_i/g_i}\cdot\sum_i\Psi_{\mathcal A_i\setminus g_i}/\Psi_{\mathcal A_i/g_i}$; it is the dual of the series closure of the dual cells.
\item The gluing of two cells along their hinges is $\binom{\{h_1,h_2\}}1\langle\mathcal A_1,\mathcal A_2\rangle$, with the letters $K=K_{\mathcal A_1\setminus g_1}K_{\mathcal A_2/g_2}+K_{\mathcal A_1/g_1}K_{\mathcal A_2\setminus g_2}$ and $\Psi=\Psi_{\mathcal A_1/g_1}\Psi_{\mathcal A_2\setminus g_2}+\Psi_{\mathcal A_1\setminus g_1}\Psi_{\mathcal A_2/g_2}$.
\end{enumerate}
\end{corollary}

% New: proved in the last part.
\begin{proposition}[Dual pairs]\label{prop:mot-dualletters}
Let $\mathcal A$ be a plain cell on $B\sqcup\{g\}$ of rank $r_{\mathcal A}$ that is a dual pair (Definition~\ref{def:substitution}), and put $k=K_{\mathcal A\setminus g}$ and $\psi=\Psi_{\mathcal A\setminus g}$.
\begin{enumerate}[label=(\roman*)]
\item A plain cell is a dual pair exactly when its bases are closed under complement with equal weights; then $|B|=2r_{\mathcal A}-1$, so that $k$ has the degree $r_{\mathcal A}$ and $\psi$ the degree $r_{\mathcal A}-1$.
\item $K_{\mathcal A/g}=\psi$ and $\Psi_{\mathcal A/g}=k$; the two hinge values agree, $\gamma^K_{\mathcal A}=\gamma^\Psi_{\mathcal A}=k/\psi$, and $\varrho_{\mathcal A}=1$.
\item Let $\mathcal H\langle\mathcal A_1,\dots,\mathcal A_q\rangle$ be a substitution of dual pairs into a plain core $\mathcal H$, of rank $r_{\mathcal H}$, with $k_i$ and $\psi_i$ as above and $\pi=\psi_1\cdots\psi_q$; a mark left in place is a dual pair with $k_i=\alpha_h$ and $\psi_i=1$. Then, as identities of polynomials, with $\pi k_i/\psi_i=k_i\prod_{l\neq i}\psi_l$ and $K$, $\Psi$ the letters of the substitution,
\[
K_{\mathcal H}\big((\pi k_i/\psi_i)_i\big)=\pi^{\,r_{\mathcal H}-1}K,\qquad\Psi_{\mathcal H}\big((\pi k_i/\psi_i)_i\big)=\pi^{\,q-r_{\mathcal H}-1}\Psi .
\]
\end{enumerate}
\end{proposition}

% Reference 20: prop:substform; source line 5794
\begin{proposition}[The multivector of a substitution]\label{prop:substform}
Let $\mathcal H$ and the cells carry multivectors $\boldsymbol\omega_{\mathcal H}$ and $\boldsymbol\omega_i$ (Definition~\ref{def:multivector}), the hinge first in the order of the marks of each cell, and order the marks of the substitution by putting $B_i$, in its order in $\mathcal A_i$, in the place of $h_i$ in the order of $\mathcal H$. Then
\[
\boldsymbol\omega\Big(\bigsqcup_iX_i\Big)=\boldsymbol\omega_{\mathcal H}(T)\prod_{h_i\in T}\boldsymbol\omega_i(X_i)\prod_{h_i\notin T}\boldsymbol\omega_i(g_i+X_i)
\]
is a multivector of $\mathcal H\langle\mathcal A_1,\dots,\mathcal A_q\rangle$; it is multilinear in $(\boldsymbol\omega_{\mathcal H},\boldsymbol\omega_1,\dots,\boldsymbol\omega_q)$, a substitution of decomposable multivectors is decomposable, and
\[
\operatorname{rk}(\boldsymbol\omega)\le\operatorname{rk}(\boldsymbol\omega_{\mathcal H})\prod_i\operatorname{rk}(\boldsymbol\omega_i).
\]
A substitution of determinantal tables is determinantal.
\end{proposition}

% Reference 20: thm:substwalls; source line 5805
\begin{theorem}[The walls of a substitution]\label{thm:substwalls}
For $S=\bigsqcup_iS_i$, $S_i\subseteq B_i$, let $D(S)$ be the set of the marks $h_i$ of the core with $d_{\mathcal A_i}(S_i+g_i)>d_{\mathcal A_i}(S_i)$ and $D^\ast(S)$ that of the $h_i$ with $d^\ast_{\mathcal A_i}(S_i+g_i)>d^\ast_{\mathcal A_i}(S_i)$, the two differences being $0$ or $1$. The counts of the substitution are
\[
d(S)=\sum_id_{\mathcal A_i}(S_i)+d_{\mathcal H}\big(D(S)\big),\qquad d^\ast(S)=\sum_id^\ast_{\mathcal A_i}(S_i)+d^\ast_{\mathcal H}\big(D^\ast(S)\big),
\]
and for a substituted family, with $D_j(S)$ and $D^\ast_j(S)$ read in table $j$, its walls are
\[
\kappa_S=\sum_i\Big(m_{S_i}-\sum_j\big(s_jd_{\mathcal A_{ij}}(S_i)+t_jd^\ast_{\mathcal A_{ij}}(S_i)\big)\Big)-\sum_j\Big(s_jd_{\mathcal H_j}\big(D_j(S)\big)+t_jd^\ast_{\mathcal H_j}\big(D^\ast_j(S)\big)\Big).
\]
At $S=B_i$, $D(S)=D^\ast(S)=\{h_i\}$ and the wall is $\kappa_{B_i}=m_{B_i}-\sum_j\big(s_j(r_{\mathcal A_{ij}}-1)+t_j(|B_i|-r_{\mathcal A_{ij}})\big)$, the mass that pooling gives the hinge of $B_i$ (Theorem~\ref{thm:pooling}).
\end{theorem}

% New: proved in the last part.
\begin{corollary}[The width of the marks left in place]\label{cor:mot-widthbound}
Let a substitution have the core $\mathcal H$ on $L\sqcup H'$, the marks of $L$ left in place and cells with at least two marks at the marks of $H'$. Every $S\subseteq L$ has the width it has in the core, and
\[
c(S)=c_{\mathcal H}(S)\le|L\setminus S|+|H'| .
\]
So a bunch of width $c$ that is all the marks left in place needs at least $c$ other cells, and every width occurs: in $\binom{B\sqcup\{h_1,\dots,h_c\}}c$ with $|B|=c$, the marks of $B$ left in place, $B$ has the width $c$, whatever the cells at $h_1,\dots,h_c$.
\end{corollary}

% Reference 20: lem:twosumpieces; source line 5820 (one table).
% Proved in the last part: lem:twosumpieces.
\begin{lemma}[Two-sums of cells; hinge walls]\label{lem:twosumpieces}
Let $\mathcal A$ be a cell on $B\sqcup\{g\}$ and $\mathcal A'$ a cell on $R\sqcup\{g\}$ with the one hinge $g$, $B\cap R=\varnothing$, and let $\mathcal A\oplus_g\mathcal A'$ be their gluing along the hinge, $\binom{\{h_1,h_2\}}1\langle\mathcal A,\mathcal A'\rangle$ (Corollary~\ref{cor:blocksclosures}(iii)), each basis carrying in each letter the product of the weights of the two bases of $\mathcal A$ and $\mathcal A'$ it is made of: a table on $B\sqcup R$ of rank $r_{\mathcal A}+r_{\mathcal A'}-1$, plain when $\mathcal A$ and $\mathcal A'$ are, in which $B$ is a bunch of width one with the cell of the bunch $\mathcal A$ and the cell of the room $\mathcal A'$ (Lemma~\ref{lem:gluemark}); every table with a bunch of width one, the weights factoring across it, is the two-sum of its two cells (Proposition~\ref{prop:substitution}(iii)). Let the cell families $F_B=((\mathcal A_j,\theta^{\mathcal A_j});s_j,t_j)_{j=1}^N$ on $B\sqcup\{g\}$ and $F_R=((\mathcal A'_j,\theta^{\mathcal A'_j});s_j,t_j)_{j=1}^N$ on $R\sqcup\{g\}$ be glued table by table into the family $F=((\mathcal A_j\oplus_g\mathcal A'_j);s_j,t_j)_j$ on $B\sqcup R$, with masses $\mvec$. The walls of $F$ at $B$ and at $R$ are the hinge walls of the two cell families, for one table those of the two cells, $\kappa_{\mathcal A}$ and $\kappa_{\mathcal A'}$,
\[
\kappa_{F_B}=m_B-\sum_j\big(s_j(r_{\mathcal A_j}-1)+t_j(|B|-r_{\mathcal A_j})\big),\qquad \kappa_{F_R}=m_R-\sum_j\big(s_j(r_{\mathcal A'_j}-1)+t_j(|R|-r_{\mathcal A'_j})\big),
\]
and $\kappa_{F_B}+\kappa_{F_R}=2\sigma+\kappa_{\cE}$, $\sigma$ the total strength and $\kappa_{\cE}$ the imbalance of $F$ (Definition~\ref{def:family}): on the balance wall the two hinge walls add to $2\sigma=\sum_j(s_j+t_j)$, and the mass $2\sigma-\kappa_{F_B}$ on $g$ puts the cell family $F_B$ itself on its balance wall.
\end{lemma}

Each letter of a two-sum is a letter of a hinge minor of one cell times the letter of the other cell read at one hinge value of the first (Theorem~\ref{thm:substletters}): $K$ sees $\mathcal A$ through $\gamma^K_{\mathcal A}$ and $\Psi$ through $\gamma^\Psi_{\mathcal A}$. The hinge wall of a cell is a number of the cell alone, its masses and its exponents, and balance makes the two hinge walls of a two-sum complementary; the period of a two-sum is read through them in Theorem~\ref{thm:glue}.

\section{Splicing hinges in pairs}\label{sec:hingepairs}

A cell has one hinge. Gluing along several marks at once, or around a closed path of cells, needs cells with several hinges and a splicing of hinges in pairs. The splicing stays inside tables because the levels of a bunch are tables (Proposition~\ref{prop:higgs}); one pair joining two cells is the gluing of Lemma~\ref{lem:gluemark}.

% Reference 20: def:channels; source line 5849
\begin{definition}[Cells with several hinges; channels]\label{def:channels}
A \emph{cell with the hinges} $G$ is a table $\mathcal A$ on $B\sqcup G$ with two systems of positive weights $\theta^K$, $\theta^\Psi$ on its bases; $n_{\mathcal A}=|B|$ is its number of marks other than its hinges. Its \emph{channels} are the letters of the children of the bunch $G$ (Definition~\ref{def:patterns}): for a pattern $p\in\pat(G)$, the basis polynomial $K_{\mathcal A,(p)}$ of the child at $p$ with the weights $\theta^K$ and its cobasis polynomial $\Psi_{\mathcal A,(p)}$ with the weights $\theta^\Psi$, polynomials on $B$, so that by the Taylor cube (Theorem~\ref{thm:taylorcube})
\[
K_{\mathcal A}=\sum_{p\in\pat(G)}\alpha_G^{\,p}\,K_{\mathcal A,(p)}(\alpha_B),\qquad \Psi_{\mathcal A}=\sum_{p\in\pat(G)}\alpha_G^{\,G\setminus p}\,\Psi_{\mathcal A,(p)}(\alpha_B):
\]
$p$ is the set of hinges a basis occupies, in both letters. A channel has positive coefficients on its support and the degree
\[
\deg K_{\mathcal A,(p)}=r_{\mathcal A}-|p|,\qquad \deg\Psi_{\mathcal A,(p)}=n_{\mathcal A}-r_{\mathcal A}+|p| .
\]
The \emph{channel vector} $\operatorname{ch}(\mathcal A)$ lists the channels by their \emph{tags}, the pairs of a letter and a set $p\subseteq G$, $2^{|G|+1}$ of them; a tag whose set is not a pattern is an absent channel, counted as zero. A cell of Definition~\ref{def:substitution} is the case of one hinge $g$ that is neither a loop nor a coloop, with $K_{\mathcal A,(\varnothing)}=K_{\mathcal A\setminus g}$, $K_{\mathcal A,(\{g\})}=K_{\mathcal A/g}$, $\Psi_{\mathcal A,(\varnothing)}=\Psi_{\mathcal A\setminus g}$ and $\Psi_{\mathcal A,(\{g\})}=\Psi_{\mathcal A/g}$. \emph{Splicing} a pair $\{a,b\}$ of hinges keeps the bases that occupy exactly one of them and forgets the pair; on either letter it is
\[
L\longmapsto(\partial_{\alpha_a}+\partial_{\alpha_b})L\big|_{\alpha_a=\alpha_b=0}.
\]
Cells are spliced along pairs of their hinges by splicing the pairs in their direct sum (Proposition~\ref{prop:directsum}).
\end{definition}

% Reference 20: thm:hingepairs; source line 5865 (two cells in (iii)).
% Proved in the last part: thm:hingepairs.
\begin{theorem}[Splicing hinges in pairs]\label{thm:hingepairs}
\begin{enumerate}[label=(\roman*)]
\item Splicing a pair $\{a,b\}$ of hinges of a cell $\mathcal A$ at which some basis occupies exactly one of $a$, $b$ gives a cell on $B\sqcup(G\setminus\{a,b\})$ of rank $r_{\mathcal A}-1$: its bases are the sets $T\setminus\{a,b\}$ with $|T\cap\{a,b\}|=1$, each carrying in each letter the sum of the weights of the bases $T$ that give it.
\item Let a direct sum of cells carry a set of disjoint pairs of hinges, forming any graph on the cells: closed paths, a pair inside one cell and several pairs between two cells are allowed. If some choice of one basis in each cell occupies exactly one hinge of every pair, splicing the pairs one at a time, in any order, gives one cell; its rank is the sum of the ranks less the number of pairs, its bases are the unions of the unpaired parts of such choices, and its weight in each letter at a basis is the sum, over the choices that give it, of the products of the weights chosen. Plain cells give a plain cell.
\item Under the hypothesis of (ii), let the cells be $\mathcal A_1,\dots,\mathcal A_m$ with the hinges $G_1,\dots,G_m$, and in this statement let $G'$ be the set of the hinges that lie in no pair, $k$ the number of pairs and, for $p\subseteq G'$, let $J$ run over the sets of paired hinges that contain exactly one hinge of every pair. The spliced cell $\mathcal A$ has the channels
\[
K_{\mathcal A,(p)}=\sum_J\prod_{i=1}^mK_{\mathcal A_i,((p\sqcup J)\cap G_i)},\qquad \Psi_{\mathcal A,(p)}=\sum_J\prod_{i=1}^m\Psi_{\mathcal A_i,((p\sqcup J)\cap G_i)},
\]
absent channels counting as zero; it has $\sum_in_{\mathcal A_i}$ marks besides its hinges $G'$ and the rank $\sum_ir_{\mathcal A_i}-k$, and every product on the right has the degree of the channel on the left. The map $\big(\operatorname{ch}(\mathcal A_1),\dots,\operatorname{ch}(\mathcal A_m)\big)\mapsto\operatorname{ch}(\mathcal A)$ is multilinear. For two cells $\mathcal A_1$ and $\mathcal A_2$ spliced along $k$ pairs, each pair named once, so that the paired hinges of either cell form one set $I$ of $k$ names and the other hinges form $G'_1$ and $G'_2$, it is the bilinear map
\begin{gather*}
K_{\mathcal A_{12},(p_1\sqcup p_2)}=\sum_{I'\subseteq I}K_{\mathcal A_1,(p_1\sqcup I')}\,K_{\mathcal A_2,(p_2\sqcup(I\setminus I'))},\\
\Psi_{\mathcal A_{12},(p_1\sqcup p_2)}=\sum_{I'\subseteq I}\Psi_{\mathcal A_1,(p_1\sqcup I')}\,\Psi_{\mathcal A_2,(p_2\sqcup(I\setminus I'))},
\end{gather*}
$p_1\subseteq G'_1$, $p_2\subseteq G'_2$; with $k=1$ and no other hinges it is Corollary~\ref{cor:blocksclosures}(iii).
\end{enumerate}
\end{theorem}

Chapter~\ref{ch:interfaces} turns the bilinearity of Theorem~\ref{thm:hingepairs}(iii) into finite moment vectors and transfer matrices.

\section{What crosses a gluing}\label{sec:crossing}

% Reference 20: def:levelweights; source line 5886
\begin{definition}[Level weights, level factorizations, interface]\label{def:levelweights}
Let $(\cB,\theta)$ be a table of rank $r$ with a bunch $B$, of any width, and room $R$. For a level $k$, $d(B)\le k\le r(B)$, the \emph{level weights} are the matrix $\Theta^{(k)}_B=(\theta_{p\sqcup X})_{p,X}$ indexed by the $k$-sets $p\subseteq B$ and the $(r-k)$-sets $X\subseteq R$, with $\theta_{p\sqcup X}=0$ when $p\sqcup X$ is not a basis. A \emph{level factorization} writes every level letter (Definition~\ref{def:levels}) as a sum of products over a finite set $C_k$,
\[
K^{(k)}_B(\alpha)=\sum_{c\in C_k}V_{k,c}(\alpha_B)\,L_{k,c}(\alpha_R),
\]
each $V_{k,c}$ a real linear combination of the monomials $\alpha_B^p$, $|p|=k$, and each $L_{k,c}$ a real linear combination of the monomials $\alpha_R^X$, $|X|=r-k$: equivalently, a writing of $\Theta^{(k)}_B$ as a sum of $|C_k|$ products of a column and a row, so that $|C_k|\ge\operatorname{rank}\Theta^{(k)}_B$, with equality for some factorization. The $V_{k,c}$ are the \emph{bunch letters} and the $L_{k,c}$ the \emph{room letters}. With $V^\ast_{k,c}(\alpha_B)=\alpha_B^{\,B}V_{k,c}(\alpha_B^{-1})$ and $L^\ast_{k,c}(\alpha_R)=\alpha_R^{\,R}L_{k,c}(\alpha_R^{-1})$ the same factorization writes the colevel letter of the colevel $|B|-k$,
\[
\Psi^{(|B|-k)}_B(\alpha)=\sum_{c\in C_k}V^\ast_{k,c}(\alpha_B)\,L^\ast_{k,c}(\alpha_R):
\]
each level is one function read twice (Proposition~\ref{prop:identity}). The \emph{stations} of $B$ are the maximal sets of patterns whose children are positive multiples of one another, each with a fixed representative, the \emph{child of the station}; every station lies in one level. The \emph{children factorization} lets $c$ run over the stations of level $k$, with $V_{k,c}=\sum_{p\in c}w_p\,\alpha_B^p$ for the child at $p$ equal to $w_p$ times the child of the station, and $L_{k,c}$ the basis polynomial of the child of the station: all its terms are nonnegative. The \emph{interface dimension} of $B$ is $\sum_k\operatorname{rank}\Theta^{(k)}_B$; for a family it is summed over the tables.
\end{definition}

Along the scale of $B$, $\alpha_B=\rho u$ with $u\in\Delta_B$ (Lemma~\ref{lem:scaleshape}), $V_{k,c}(\rho u)=\rho^kV_{k,c}(u)$ and $V^\ast_{k,c}(\rho u)=\rho^{|B|-k}V^\ast_{k,c}(u)$; since $|B|-r(B)=d^\ast(B)$,
\begin{gather*}
K(\rho u,\alpha_R)=\rho^{d(B)}\sum_k\rho^{k-d(B)}\sum_cV_{k,c}(u)\,L_{k,c}(\alpha_R),\\
\Psi(\rho u,\alpha_R)=\rho^{d^\ast(B)}\sum_k\rho^{r(B)-k}\sum_cV^\ast_{k,c}(u)\,L^\ast_{k,c}(\alpha_R),
\end{gather*}
each letter counted from its own bottom level, both powers running over $\rho^0,\dots,\rho^{c(B)}$.

% Reference 20: thm:crossing; source line 5897
\begin{theorem}[What crosses a gluing]\label{thm:crossing}
Let $F=((\cB_j,\theta_j);s_j,t_j)_{j=1}^N$ be a family on $\cE$ with real data and masses $\mvec$ on its balance wall, $B$ a bunch with room $R$, and fix a level factorization of every table, with the bunch letters $V_{j,k,c}$, the room letters $L_{j,k,c}$ and the starred letters $V^\ast_{j,k,c}$, $L^\ast_{j,k,c}$. Let $\mathbf V(u)=\big(V_{j,k,c}(u),V^\ast_{j,k,c}(u)\big)_{j,k,c}$ be the \emph{letter vector} of $u\in\Delta_B$. Then, in $[0,\infty]$,
\begin{multline*}
P_F(\mvec;\mathbf s,\mathbf t)=\int_{\Delta_R}v^{\mvec|_R-\one}\int_{\Delta_B}u^{\mvec|_B-\one}\int_0^\infty\rho^{\kappa_B-1}\prod_j\Big(\sum_{k,c}\rho^{k-d_j(B)}V_{j,k,c}(u)\,L_{j,k,c}(v)\Big)^{-s_j}\\
\times\Big(\sum_{k,c}\rho^{r_j(B)-k}V^\ast_{j,k,c}(u)\,L^\ast_{j,k,c}(v)\Big)^{-t_j}d\rho\,d\omega_B(u)\,d\omega_R(v),
\end{multline*}
$\kappa_B$ the wall of $B$, and the integrand, apart from the measure $u^{\mvec|_B-\one}d\omega_B(u)$, depends on $u$ only through $\mathbf V(u)$. So the period sees the bunch only through its wall $\kappa_B$, its widths $c_j(B)$ and the image $\mathbf V_\ast\big(u^{\mvec|_B-\one}d\omega_B(u)\big)$, the joint law of its bunch letters, unnormalized: two families that agree on the room, with the same room letters, and whose bunches have equal walls, widths and laws of their letter vectors, have equal periods, whatever the numbers of marks of the bunches; the bunches are interchangeable in every room. The same holds for the interest $J_F(\mvec;\mathbf s,\mathbf t)$ at any real masses, with $\alpha_R\in\mathbb R^R_{>0}$ in place of $v\in\Delta_R$ and the factor $e^{-\rho-|\alpha_R|}$ in the integrand.
\end{theorem}

At width one with factoring weights the children factorization has one set at each level, and the letter vector is $(K_{\cB/R},K_{\cB|B};\Psi_{\cB|B},\Psi_{\cB/R})$, each letter from its bottom level: Theorem~\ref{thm:pooling}, which keeps the scale as the hinge $g$, keeps of it the cell kernel (Definition~\ref{def:twistedcore}) and the two hinge values $\gamma^K_B$, $\gamma^\Psi_B$. The theorem holds whatever the weights; what the weights decide is how many bunch letters there are.

% Reference 20: thm:cutrank; source line 5911 (determinantal tables in (iii)).
% Proved in the last part: thm:interfacerank.
\begin{theorem}[Multivectors bound what crosses]\label{thm:interfacerank}
Let the table carry a multivector $\boldsymbol\omega$ of any tensor rank (Definition~\ref{def:multivector}), the marks of $B$ ordered before those of $R$, so that $\boldsymbol\omega_{p\sqcup X}$ is the coefficient of $\mathbf e_p\wedge\mathbf e_X$, $\mathbf e_p$ the wedge of the unit vectors of the marks of $p$ in their order. At the level $k$ let the \emph{interface rank} $\operatorname{ir}_k(\boldsymbol\omega)$ of $\boldsymbol\omega$ at $B$ be the rank of the matrix $(\boldsymbol\omega_{p\sqcup X})_{p,X}$, and write $\boldsymbol\omega_{p\sqcup X}=\sum_{i=1}^{\operatorname{ir}_k(\boldsymbol\omega)}f_i(p)\,f'_i(X)$.
\begin{enumerate}[label=(\roman*)]
\item $\operatorname{rank}\Theta^{(k)}_B\le\operatorname{ir}_k(\boldsymbol\omega)^2$, and $\operatorname{rank}\Theta^{(k)}_B\le\frac12\operatorname{ir}_k(\boldsymbol\omega)\big(\operatorname{ir}_k(\boldsymbol\omega)+1\big)$ when $\boldsymbol\omega$ is real. With the positive semidefinite Hermitian matrices
\[
\mathbf C_k(\alpha_B)=\sum_p\alpha_B^p\,f(p)f(p)^\ast,\qquad \mathbf C'_k(\alpha_R)=\sum_X\alpha_R^X\,f'(X)f'(X)^\ast,
\]
with $f(p)=\big(f_i(p)\big)_i$ and $f'(X)=\big(f'_i(X)\big)_i$,
the level letter is $K^{(k)}_B(\alpha)=\operatorname{tr}\big(\mathbf C_k(\alpha_B)\,\overline{\mathbf C'_k(\alpha_R)}\big)$, and the real and imaginary parts of the entries of $\mathbf C_k(\alpha_B)$ are the bunch letters of a real level factorization with at most $\operatorname{ir}_k(\boldsymbol\omega)^2$ terms.
\item If $\boldsymbol\omega=\sum_{l=1}^{\operatorname{rk}(\boldsymbol\omega)}\boldsymbol\omega^{(l)}$ with slices $A^{(l)}$, let the \emph{width of the slice} $c_l$ be the dimension of the intersection of the column spaces of the columns $A^{(l)}_B$ of $B$ and $A^{(l)}_R$ of $R$, and its \emph{floor} $d_l=r-\dim\big(\text{column space of }A^{(l)}_R\big)$, which at tensor rank one are $c(B)$ and $d(B)$. Then $\operatorname{ir}_k(\boldsymbol\omega)\le\sum_l\binom{c_l}{k-d_l}$.
\item So $\operatorname{rank}\Theta^{(k)}_B\le\big(\sum_l\binom{c_l}{k-d_l}\big)^2$, and the interface dimension of $B$ is at most $\sum_k\big(\sum_l\binom{c_l}{k-d_l}\big)^2$, a bound in the widths and floors of the slices alone, whatever the numbers of marks of $B$ and $R$; for a real multivector each square may be replaced by $\frac12i_k(i_k+1)$, $i_k=\sum_l\binom{c_l}{k-d_l}$. For a determinantal table, of tensor rank one with $c_1=c(B)$ and $d_1=d(B)$, $\operatorname{ir}_k(\boldsymbol\omega)\le\binom{c(B)}{k-d(B)}$, $\operatorname{rank}\Theta^{(k)}_B\le\binom{c(B)}{k-d(B)}^2$ and the interface dimension is at most $\binom{2c(B)}{c(B)}$; at width one each level has one bunch letter, as Theorem~\ref{thm:levels} says.
\end{enumerate}
\end{theorem}

% Reference 20: thm:schurinterface; source line 5923
\begin{theorem}[The interface of a determinantal table]\label{thm:schurinterface}
Let a table be determinantal over $\mathbb F$ with slice $A$, $\theta_T=|\det A_T|^2$, and let $B$ have width $c(B)=w$, the minors $\cB|B$ and $\cB/R$ on $B$ and $\cB|R$ and $\cB/B$ on $R$ carrying the factoring weights of Theorem~\ref{thm:levels}. There are Hermitian $w\times w$ matrices $\mathbf S_B(\alpha_B)$ and $\mathbf S_R(\alpha_R)$, positive definite on the open orthant, each entry a rational function homogeneous of degree one, with
\[
K(\alpha)=K_{\cB/R}(\alpha_B)\,K_{\cB/B}(\alpha_R)\,\det\big(\mathbf S_B(\alpha_B)+\mathbf S_R(\alpha_R)\big),\qquad \det\mathbf S_B=\frac{K_{\cB|B}}{K_{\cB/R}},\qquad \det\mathbf S_R=\frac{K_{\cB|R}}{K_{\cB/B}},
\]
and $\Psi(\alpha)=\alpha^{\cE}K(\alpha^{-1})$. Along the scale of $B$, $K(\rho\alpha_B,\alpha_R)=\rho^{d(B)}K_{\cB/R}(\alpha_B)K_{\cB/B}(\alpha_R)\det\big(\rho\,\mathbf S_B+\mathbf S_R\big)$, so the Skeleton values of $B$ (Definition~\ref{def:skelpoly}) are the eigenvalues of $\mathbf S_B^{-1/2}\mathbf S_R\mathbf S_B^{-1/2}$, real and positive. At width one $\mathbf S_B=\gamma^K_B$, the hinge value by $K$ (Definition~\ref{def:readings}). So a bunch of width $w$ of a determinantal table crosses its gluing as one positive definite $w\times w$ matrix and one scalar, and the same read at $\alpha^{-1}$ for $\Psi$.
\end{theorem}

For a determinantal table a bunch $B$ is thus read by its room through one positive definite matrix, $\binom{c(B)+1}2$ real entries when it is symmetric and $c(B)^2$ real parameters when it is Hermitian, and one scalar for each letter, and the bunch letters of Theorem~\ref{thm:interfacerank} are the compound determinants of that matrix. For a multivector of tensor rank $\operatorname{rk}(\boldsymbol\omega)$ the level letters carry the entries of the matrices $\mathbf C_k$ of Theorem~\ref{thm:interfacerank}(i), whose size is fixed by the interface ranks. In every case Theorem~\ref{thm:crossing} says what crosses, and Theorems~\ref{thm:interfacerank} and~\ref{thm:schurinterface} say how much.


\chapter{Periods of constructions}\label{ch:construct}

A construction (Definition~\ref{def:construction}) hands its period a structure: the period of a substitution is read from the core and the cells, for any family. The core reads each cell through the joint law of its hinge values, which is pooling done at every cell at once; opened on the core's own shares, the cells ride on the core's fibre, each entering through the period of its own family at the exponents the shares assign it; and where every active letter reads a cell alike, the cell's share is pinned; where the active letters read every cell by one hinge value, as at a full lean of one table, the period factors through the core. Chains of cells are the powers of one operator, and the letters of any chain of bunches, and their integer joint moments, are products of finite matrices. A bunch of any width opens on one share per term of its level factorization; in level form, one term per level, one table at a full lean keeps $c(B)-1$ free shares. The same periods are read through the laws of the cell families, one measure on $(0,\infty)^{2N-1}$ for each (Chapter~\ref{ch:hinged}), and through the periods at the full leans of the nodes (Chapters~\ref{ch:holding} and~\ref{ch:leantree}).

\section{The core reads its cells}\label{sec:corereads}

% Reference 20: def:twist; source line 6157 (with def:hold, source line 6603; the reference rescales a twist at one mark and, for one table, a twist of the cobasis letter).
% Proved in the last part: def:twist.
\begin{definition}[Twists]\label{def:twist}
A \emph{twist} of a family $F$ on $\cE$ assigns to every table $j$ two vectors $\mathbf x^K_j,\mathbf x^\Psi_j\in(0,\infty)^{\cE}$. The \emph{twisted family} $F[\mathbf x]$, $\mathbf x=(\mathbf x^K_j,\mathbf x^\Psi_j)_j$, has the bases and the exponents of $F$, the basis weights $\theta^K_{j,T}(\mathbf x^K_j)^{T}$ and the cobasis weights $\theta^\Psi_{j,T}(\mathbf x^\Psi_j)^{\cE\setminus T}$: its letters are $K_j(\mathbf x^K_j\circ\alpha)$ and $\Psi_j(\mathbf x^\Psi_j\circ\alpha)$, with $\mathbf x\circ\alpha=(x_e\alpha_e)_e$. A twist \emph{at a mark $g$} has all its coordinates $1$ off $g$ and is written by its values $\mathbf x=(x^K_j,x^\Psi_j)_j$ at $g$; then $F[\mathbf x@g]$ multiplies, in table $j$, the weights in $K_j$ of the bases through $g$ by $x^K_j$ and the weights in $\Psi_j$ of the bases avoiding $g$ by $x^\Psi_j$. A twist of one table $\mathcal H$ in its cobasis letter alone, $\mathbf x^K=\one$, is written $\mathcal H[\mathbf x]$ with $\mathbf x=\mathbf x^\Psi$: $K_{\mathcal H[\mathbf x]}=K_{\mathcal H}$ and $\Psi_{\mathcal H[\mathbf x]}(\alpha)=\Psi_{\mathcal H}(\mathbf x\circ\alpha)$. On the balance wall the period obeys the \emph{torus relation}: for every twist $\mathbf x$, every $\boldsymbol\lambda\in(0,\infty)^{\cE}$ and all constants $c_j,c^\ast_j>0$,
\[
P_{F[(c_j\boldsymbol\lambda\circ\mathbf x^K_j,\ c^\ast_j\boldsymbol\lambda\circ\mathbf x^\Psi_j)_j]}(\mvec;\mathbf s,\mathbf t)=\boldsymbol\lambda^{-\mvec}\prod_jc_j^{-r_js_j}\,c_j^{\ast\,-r^\ast_jt_j}\,P_{F[\mathbf x]}(\mvec;\mathbf s,\mathbf t),
\]
both sides converging together, in $[0,\infty]$ for real data; at the scalar $\boldsymbol\lambda=\mu\one$ with $c_j=c^\ast_j=\mu^{-1}$ the twist is unchanged and the factor is $1$, by balance. So a twist at a mark $g$, $\boldsymbol\lambda$ equal to $1$ off $g$, has $P_{F[\lambda\mathbf x@g]}(\mvec;\mathbf s,\mathbf t)=\lambda^{-m_g}P_{F[\mathbf x@g]}(\mvec;\mathbf s,\mathbf t)$ and is read with $x^K_1=1$, a point of $(0,\infty)^{2N-1}$; for one table it is one number, $x=x^\Psi_1$, and only the cobasis letter moves. A twist of every mark of one letter by one factor multiplies the period by that factor to the power $-r_js_j$ or $-r^\ast_jt_j$; for one table in its cobasis letter, $P_{\mathcal H[\mu\mathbf x]}(\mvec;s,t)=\mu^{-r^\ast t}P_{\mathcal H[\mathbf x]}(\mvec;s,t)$.
\end{definition}

% Reference 20: def:twistedcore; source line 5947
\begin{definition}[Cell kernels; the core family read at a shape]\label{def:twistedcore}
Let $F=((\cB_j,\theta_j);s_j,t_j)_{j=1}^N$ be a substituted family (Definition~\ref{def:substitution}), table $j$ being $\mathcal H_j\langle\mathcal A_{1j},\dots,\mathcal A_{qj}\rangle$ with the cells $\mathcal A_{ij}$ on $B_i\sqcup\{g_i\}$. The \emph{cell kernel} of $B_i$ is the measure
\[
u^{\mvec|_{B_i}-\one}\prod_jK_{\mathcal A_{ij}/g_i}(u)^{-s_j}\Psi_{\mathcal A_{ij}\setminus g_i}(u)^{-t_j}\,d\omega_{B_i}(u)\qquad\text{on }\Delta_{B_i}.
\]
The \emph{core family} is $F_{\mathcal H}=((\mathcal H_j,\theta^{\mathcal H_j});s_j,t_j)_j$ on the marks $h_1,\dots,h_q$. For $u=(u_i)_i\in\prod_i\Delta_{B_i}$ the \emph{core family read at $u$} is its twist $F_{\mathcal H}[\boldsymbol\gamma(u)]$ (Definition~\ref{def:twist}) by the hinge values, $\boldsymbol\gamma(u)$ putting $\gamma^K_{\mathcal A_{ij}}(u_i)$ in $K_j$ and $\gamma^\Psi_{\mathcal A_{ij}}(u_i)$ in $\Psi_j$ at the mark $h_i$: table $j$ has the basis weights $\theta^{\mathcal H_j}_T\prod_{h_i\in T}\gamma^K_{\mathcal A_{ij}}(u_i)$ and the cobasis weights $\theta^{\mathcal H_j}_T\prod_{h_i\notin T}\gamma^\Psi_{\mathcal A_{ij}}(u_i)$, and the letters $K_{\mathcal H_j}\big((\gamma^K_{\mathcal A_{ij}}(u_i)\,\alpha_{h_i})_i\big)$ and $\Psi_{\mathcal H_j}\big((\gamma^\Psi_{\mathcal A_{ij}}(u_i)\,\alpha_{h_i})_i\big)$.
\end{definition}

% Reference 20: thm:corereads; source line 5951
\begin{theorem}[The core reads its cells]\label{thm:corereads}
For a substituted family $F$ the core family read at $u$ has the imbalance of $F$, and
\begin{multline*}
P_F(\mvec;\mathbf s,\mathbf t)=\int_{\prod_i\Delta_{B_i}}\prod_i\Big(u_i^{\mvec|_{B_i}-\one}\prod_jK_{\mathcal A_{ij}/g_i}(u_i)^{-s_j}\Psi_{\mathcal A_{ij}\setminus g_i}(u_i)^{-t_j}\,d\omega_{B_i}(u_i)\Big)\\
\times P_{F_{\mathcal H}[\boldsymbol\gamma(u)]}\big((\kappa_{B_1},\dots,\kappa_{B_q});\mathbf s,\mathbf t\big),
\end{multline*}
both sides converging together; the same holds with $J$ for $P$, and on the balance wall. So the period of a substitution depends on each cell only through its wall $\kappa_{B_i}$ and the joint law, under its cell kernel, of its hinge values $\big(\gamma^K_{\mathcal A_{ij}},\gamma^\Psi_{\mathcal A_{ij}}\big)_j$: two cells with equal walls and equal laws are interchangeable in every core. For one cell this is Theorem~\ref{thm:pooling}.
\end{theorem}

Pooling, done at every cell at once, leaves the core: each cell is integrated over its shape with its kernel and collapses to one mark of the core, which keeps the mass $\kappa_{B_i}$ and is read through the cell's hinge values as weights. The cells meet only in the core.

% Reference 20: thm:coreshares; source line 5965 (real data with $s_j,t_j\ge0$).
% Proved in the last part: thm:coreshares.
\begin{theorem}[The cells ride on the core's fibre]\label{thm:coreshares}
Let $F$ be a substituted family with complex data, every exponent $0$ or of positive real part. Open $K_{\mathcal H_j}$ with one share $a_{j,T}$ and $\Psi_{\mathcal H_j}$ with one share $b_{j,T}$ per basis $T$ of $\mathcal H_j$ (in this statement $|\mathcal H_j|$ is the number of bases of $\mathcal H_j$), $\sum_Ta_{j,T}=s_j$ and $\sum_Tb_{j,T}=t_j$, a letter with exponent zero opened by nothing, and let the \emph{loads} of the mark $h_i$ be $a_j(h_i)=\sum_{T\ni h_i}a_{j,T}$ and $b_j(h_i)=\sum_{T\not\ni h_i}b_{j,T}$. The \emph{load family} $F_i(\mathbf a,\mathbf b)$ is the family on $B_i$ with, for every $j$, the table $\mathcal A_{ij}\setminus g_i$ at the exponents $\big(a_j(h_i),\,t_j-b_j(h_i)\big)$ and the table $\mathcal A_{ij}/g_i$ at $\big(s_j-a_j(h_i),\,b_j(h_i)\big)$; its imbalance is $\kappa_{B_i}-\sum_j\big(a_j(h_i)+b_j(h_i)\big)$. In this statement write $\mathbf s^{(i)}=\big((a_j(h_i))_j,(s_j-a_j(h_i))_j\big)$ and $\mathbf t^{(i)}=\big((t_j-b_j(h_i))_j,(b_j(h_i))_j\big)$ for its exponents, the tables $\mathcal A_{ij}\setminus g_i$ listed first. On every contour of positive real parts of the shares at which all load families converge absolutely,
\begin{multline*}
J_F(\mvec;\mathbf s,\mathbf t)=\int_{|\mathbf a_j|=s_j,\ |\mathbf b_j|=t_j}\ \prod_j\Gst{|\mathcal H_j|}1{\mathbf a_j}{s_j}\Gst{|\mathcal H_j|}1{\mathbf b_j}{t_j}\prod_{j,T}\big(\theta^{\mathcal H_j}_T\big)^{-a_{j,T}-b_{j,T}}\\
\times\prod_iJ_{F_i(\mathbf a,\mathbf b)}\big(\mvec|_{B_i};\mathbf s^{(i)},\mathbf t^{(i)}\big),
\end{multline*}
with $J_{F_i(\mathbf a,\mathbf b)}\big(\mvec|_{B_i};\mathbf s^{(i)},\mathbf t^{(i)}\big)=\Gamma\big(\kappa_{B_i}-\sum_j(a_j(h_i)+b_j(h_i))\big)P_{F_i(\mathbf a,\mathbf b)}\big(\mvec|_{B_i};\mathbf s^{(i)},\mathbf t^{(i)}\big)$. The integrand without the periods of the load families is that of the interest $J_{F_{\mathcal H}}\big((\kappa_{B_i})_i;\mathbf s,\mathbf t\big)$ of the core family (Definition~\ref{def:twistedcore}) opened with one share per basis (Corollary~\ref{thm:E1}).
\end{theorem}

The cells ride on the core's share space: the core's letters are opened over its bases, every $K$-share of a basis loads the marks it occupies and every $\Psi$-share the marks it leaves out, and each cell enters through the period of its own two hinge minors at the exponents that the loads of its mark assign them.

\section{Pinned cells}\label{sec:pinned}

% Reference 20: thm:pinned; source line 5978 (the hypothesis on the letters with positive exponents).
% Proved in the last part: thm:pinned.
\begin{theorem}[A cell read alike is pinned]\label{thm:pinned}
Let $F$ be a substituted family at every cell of which every active letter reads alike: for every $i$ there are a function $\gamma_i$ on $\Delta_{B_i}$ and constants $c_{ij},c^\ast_{ij}>0$ with $\gamma^K_{\mathcal A_{ij}}=c_{ij}\gamma_i$ whenever $s_j\neq0$ and $\gamma^\Psi_{\mathcal A_{ij}}=c^\ast_{ij}\gamma_i$ whenever $t_j\neq0$, and $c_{ij}=1$ when $s_j=0$, $c^\ast_{ij}=1$ when $t_j=0$. Let $F_{\mathcal H}[\mathbf c]$ be the core family twisted by these constants at the marks $h_i$, $c_{ij}$ in $K_j$ and $c^\ast_{ij}$ in $\Psi_j$ (Definitions~\ref{def:twist} and~\ref{def:twistedcore}). Then on the balance wall
\[
P_F(\mvec;\mathbf s,\mathbf t)=P_{F_{\mathcal H}[\mathbf c]}\big((\kappa_{B_i})_i;\mathbf s,\mathbf t\big)\prod_i\int_{\Delta_{B_i}}\gamma_i^{-\kappa_{B_i}}\,u^{\mvec|_{B_i}-\one}\prod_jK_{\mathcal A_{ij}/g_i}^{-s_j}\Psi_{\mathcal A_{ij}\setminus g_i}^{-t_j}\,d\omega_{B_i}(u),
\]
the core at the walls of the cells times the \emph{pinned shape integrals} of the cells. In general the free shares at the cell $i$ are one fewer than the distinct hinge values of $B_i$ by the active letters, as at a bunch of width one (Theorem~\ref{thm:twosides}).
\end{theorem}

% Reference 20: cor:fullleancore; source line 5989 (one table at $(0,t)$).
% Proved in the last part: cor:fullleancore.
\begin{corollary}[Read by one hinge value, the period factors through the core]\label{cor:fullleancore}
Let $F$ be a substituted family with real data on its balance wall (Definition~\ref{def:substitution}), table $j$ being $\mathcal H_j\langle\mathcal A_{1j},\dots,\mathcal A_{qj}\rangle$, at every cell of which the active letters read the cell by one hinge value: for every $i$ there is a function $\gamma_i$ on $\Delta_{B_i}$ with $\gamma^K_{\mathcal A_{ij}}=\gamma_i$ whenever $s_j\neq0$ and $\gamma^\Psi_{\mathcal A_{ij}}=\gamma_i$ whenever $t_j\neq0$, as at a mark left in place, whose two hinge values are both $\alpha_h$. Let $\sigma$ be the total strength (Definition~\ref{def:family}) and $F_i=((\mathcal A_{ij},\theta^{\mathcal A_{ij}});s_j,t_j)_j$ the cell family of the cells at $B_i$, on $B_i\sqcup\{g_i\}$. Then, in $[0,\infty]$,
\[
P_F(\mvec;\mathbf s,\mathbf t)=P_{F_{\mathcal H}}\big((\kappa_{B_i})_i;\mathbf s,\mathbf t\big)\prod_i\int_{\Delta_{B_i}}\gamma_i^{-\kappa_{B_i}}\,u^{\mvec|_{B_i}-\one}\prod_jK_{\mathcal A_{ij}/g_i}^{-s_j}\Psi_{\mathcal A_{ij}\setminus g_i}^{-t_j}\,d\omega_{B_i}(u),
\]
and at a cell with $0<\kappa_{B_i}<2\sigma$ the family $F_i$ at the masses $(\mvec|_{B_i},2\sigma-\kappa_{B_i})$ is on its balance wall and its pinned shape integral is $P_{F_i}\big((\mvec|_{B_i},2\sigma-\kappa_{B_i});\mathbf s,\mathbf t\big)\big/\Gst21{\kappa_{B_i},\ 2\sigma-\kappa_{B_i}}{2\sigma}$. If $F$ has every proper wall positive, then $0<\kappa_{B_i}<2\sigma$ for every cell, and
\[
P_F(\mvec;\mathbf s,\mathbf t)\prod_i\Gst21{\kappa_{B_i},\ 2\sigma-\kappa_{B_i}}{2\sigma}=P_{F_{\mathcal H}}\big((\kappa_{B_i})_i;\mathbf s,\mathbf t\big)\prod_iP_{F_i}\big((\mvec|_{B_i},2\sigma-\kappa_{B_i});\mathbf s,\mathbf t\big),
\]
every factor on its balance wall with every proper wall positive, the hinge $g_i$ at the mass $2\sigma-\kappa_{B_i}$ and the mark $h_i$ of the core family at $\kappa_{B_i}$. For one cell this is Theorem~\ref{thm:cornerfactor}, and the first display is the shape-integral form of Proposition~\ref{prop:readingsagree}.
\end{corollary}

% New: proved in the last part.
\begin{corollary}[The pinned shape integral is the period of a load family]\label{cor:mot-loadpinned}
In Corollary~\ref{cor:fullleancore}, let the loads $a_{ij},b_{ij}$ of the cell $i$ be real, with $a_{ij}=0$ when $s_j=0$, $b_{ij}=0$ when $t_j=0$ and $\sum_j(a_{ij}+b_{ij})=\kappa_{B_i}$. The load family $F_i(\mathbf a,\mathbf b)$ of Theorem~\ref{thm:coreshares} at the loads $a_j(h_i)=a_{ij}$ and $b_j(h_i)=b_{ij}$, with the tables $\mathcal A_{ij}\setminus g_i$ at $(a_{ij},t_j-b_{ij})$ and $\mathcal A_{ij}/g_i$ at $(s_j-a_{ij},b_{ij})$ and the masses $\mvec|_{B_i}$, has the imbalance zero, and its form does not depend on the loads: in $[0,\infty]$ its period is the pinned shape integral of the cell,
\[
P_{F_i(\mathbf a,\mathbf b)}\big(\mvec|_{B_i};\mathbf s^{(i)},\mathbf t^{(i)}\big)=\int_{\Delta_{B_i}}\gamma_i^{-\kappa_{B_i}}\,u^{\mvec|_{B_i}-\one}\prod_jK_{\mathcal A_{ij}/g_i}^{-s_j}\Psi_{\mathcal A_{ij}\setminus g_i}^{-t_j}\,d\omega_{B_i}(u),
\]
and $P_F(\mvec;\mathbf s,\mathbf t)=P_{F_{\mathcal H}}\big((\kappa_{B_i})_i;\mathbf s,\mathbf t\big)\prod_iP_{F_i(\mathbf a,\mathbf b)}\big(\mvec|_{B_i};\mathbf s^{(i)},\mathbf t^{(i)}\big)$. If every proper wall of $F$ is positive and $s_j,t_j\ge0$, there are such loads with $0\le a_{ij}\le s_j$ and $0\le b_{ij}\le t_j$, integers when the data are; the load family then has nonnegative exponents, and every proper wall of it is positive when $P_F$ is finite.
\end{corollary}

\section{Chains and transfer}\label{sec:transfer}

% Reference 20: thm:chainoperator; source line 6055 (exponents $s_j,t_j\ge0$).
% Proved in the last part: thm:chainoperator.
\begin{theorem}[A chain is the power of one operator]\label{thm:chainoperator}
Let $F=((\mathcal A_j,\theta_j);s_j,t_j)_{j=1}^N$ be a family on $\cE$ with two marks $g_{\mathrm{in}}\neq g_{\mathrm{out}}$, each neither a loop nor a coloop in every table, real exponents, masses on the balance wall with $m_{g_{\mathrm{in}}}+m_{g_{\mathrm{out}}}=\sum_j(s_j+t_j)$, and let $F^{\rhd k}$ be $k$ copies of $F$ glued table by table along marks, $g_{\mathrm{out}}^{(i)}$ with $g_{\mathrm{in}}^{(i+1)}$ (Lemma~\ref{lem:twosumpieces}). For $\mathbf x=(x^K_j,x^\Psi_j)_j\in(0,\infty)^{2N}$ let $F^{\rhd k}[\mathbf x@g_{\mathrm{in}}^{(1)}]$ be $F^{\rhd k}$ twisted at $g_{\mathrm{in}}^{(1)}$ (Definition~\ref{def:twist}), and in this statement write $P_k(\mathbf x)$ for its period at the masses of the copies and $\mathcal A^{\mathbf x}_j$ for table $j$ of $F[\mathbf x@g_{\mathrm{in}}]$, read as a cell with the hinge $g_{\mathrm{out}}$. With $B=\cE\setminus\{g_{\mathrm{out}}\}$ define on nonnegative functions $H$ of the twist
\begin{gather*}
(\mathcal QH)(\mathbf x)=\int_{\Delta_B}H\Big(\big(\gamma^K_{\mathcal A^{\mathbf x}_j}(u),\gamma^\Psi_{\mathcal A^{\mathbf x}_j}(u)\big)_j\Big)\,u^{\mvec|_B-\one}\prod_jK_{\mathcal A^{\mathbf x}_j/g_{\mathrm{out}}}(u)^{-s_j}\Psi_{\mathcal A^{\mathbf x}_j\setminus g_{\mathrm{out}}}(u)^{-t_j}\,d\omega_B(u),\\
P_0(\mathbf x)=\int_0^1v^{m_{g_{\mathrm{in}}}-1}(1-v)^{m_{g_{\mathrm{out}}}-1}\prod_j\big(1-v+x^K_jv\big)^{-s_j}\big(1-v+x^\Psi_jv\big)^{-t_j}\,dv .
\end{gather*}
Then $P_k=\mathcal QP_{k-1}$, and in $[0,\infty]$
\[
P_{F^{\rhd k}}(\mvec;\mathbf s,\mathbf t)=(\mathcal Q^kP_0)(\one).
\]
The operator acts on functions of the $2N$ twists, of which one is removed by the scale of the hinge. Where the twist is pinned, the active letters reading the cell by one hinge value, $\gamma^K_{\mathcal A_j}$ for $s_j\neq0$ and $\gamma^\Psi_{\mathcal A_j}$ for $t_j\neq0$ all equal, as at a full lean of one table, every $P_k$ is a multiple of one power of the twist along the diagonal $\mathbf x=\lambda\one$, $\mathcal Q$ multiplies it by $P_F(\mvec;\mathbf s,\mathbf t)/\Gst21{m_{g_{\mathrm{in}}},\ m_{g_{\mathrm{out}}}}{m_{g_{\mathrm{in}}}+m_{g_{\mathrm{out}}}}$, and
\[
P_{F^{\rhd k}}(\mvec;\mathbf s,\mathbf t)=\frac{P_F(\mvec;\mathbf s,\mathbf t)^k}{\Gst21{m_{g_{\mathrm{in}}},\ m_{g_{\mathrm{out}}}}{m_{g_{\mathrm{in}}}+m_{g_{\mathrm{out}}}}^{k-1}}.
\]
Hinge values that are only proportional, with constants other than one, still pin each cell (Theorem~\ref{thm:pinned}) but carry the twist off the diagonal.
\end{theorem}

A chain is pooled one copy at a time: the first copy, read as the cell of its bunch, collapses to its hinge and twists the next copy by its hinge values, so the period of the chain is one operator on functions of the twist applied once per copy, and the two-mark table left at the end supplies the starting function. Where the hinge values pin the twist the operator is a multiplication, and the chain costs a power.

% Reference 20: thm:lettertransfer; source line 6071 (the count of states bounded for a determinantal table).
% Proved in the last part: thm:lettertransfer.
\begin{theorem}[Letters by transfer]\label{thm:lettertransfer}
Let $(\cB,\theta)$ be a table with any weights and $\varnothing=S_0\subset S_1\subset\dots\subset S_L=\cE$ a chain of bunches, with the steps $B_{l-1}=S_l\setminus S_{l-1}$. Group the patterns of each $S_l$ into its stations (Definition~\ref{def:levelweights}), the children being tables on $\cE\setminus S_l$ and the child at $p$ being $w_p$ times the child of its station $c$. For a station $c$ of $S_{l-1}$ and a station $c'$ of $S_l$ let
\[
\mathbf M^K_l(c,c')=\sum_{p'}w(c,p')\,\alpha_{B_{l-1}}^{\,p'},\qquad \mathbf M^\Psi_l(c,c')=\sum_{p'}w(c,p')\,\alpha_{B_{l-1}}^{\,B_{l-1}\setminus p'},
\]
the sums over the patterns $p'$ of $B_{l-1}$ in the child of the station $c$ whose child is $w(c,p')$ times the child of the station $c'$. Then, with $\mathbf e$ the unit vector of the station of the empty pattern of $S_0$ and $\boldsymbol\theta_\varnothing$ the vector of the weights of the one basis $\varnothing$ of the children of the stations of $S_L$,
\[
K=\mathbf e^{\mathsf T}\mathbf M^K_1\mathbf M^K_2\cdots\mathbf M^K_L\,\boldsymbol\theta_\varnothing,\qquad \Psi=\mathbf e^{\mathsf T}\mathbf M^\Psi_1\cdots\mathbf M^\Psi_L\,\boldsymbol\theta_\varnothing .
\]
The number of states at $S_l$ is the number of its stations, at least the interface dimension of $S_l$ (Definition~\ref{def:levelweights}), and the states can be reduced to that many linearly independent ones; at a bunch $S_l$ of width one with factoring weights it is two. For a table carrying a multivector $\boldsymbol\omega$ of any tensor rank (Definition~\ref{def:multivector}), with the widths $c_i$ and the floors $d_i$ of its slices at $S_l$ (Theorem~\ref{thm:interfacerank}(ii)), the interface dimension of $S_l$ is at most $\sum_k\big(\sum_i\binom{c_i}{k-d_i}\big)^2$, each square $q^2$ replaced by $\frac12q(q+1)$ when $\boldsymbol\omega$ is real; for a determinantal table it is at most $\binom{2c(S_l)}{c(S_l)}$ (Theorem~\ref{thm:interfacerank}).
\end{theorem}

% Reference 20: thm:momenttransfer; source line 6086 (one table and identical steps).
% Proved in the last part: thm:momenttransfer.
\begin{theorem}[Integer joint moments by fixed matrices]\label{thm:momenttransfer}
Let the letters of a family $F=((\cB_j,\theta_j);s_j,t_j)_{j=1}^N$ be products of matrix functions of the variables of the steps $B_0,\dots,B_{L-1}$ of one chain of bunches, $K_j=\mathbf e_j^{\mathsf T}\mathbf M^K_{j,0}(\alpha_{B_0})\cdots\mathbf M^K_{j,L-1}(\alpha_{B_{L-1}})\boldsymbol\theta_{j,\varnothing}$ and $\Psi_j=\mathbf e_j^{\mathsf T}\mathbf M^\Psi_{j,0}(\alpha_{B_0})\cdots\mathbf M^\Psi_{j,L-1}(\alpha_{B_{L-1}})\boldsymbol\theta_{j,\varnothing}$ with constant vectors $\mathbf e_j$ and $\boldsymbol\theta_{j,\varnothing}$, as Theorem~\ref{thm:lettertransfer} gives for every table along any chain of bunches; and let the steps carry independent positive measures $\phi_0,\dots,\phi_{L-1}$ with the required moments finite. For integers $a_j,b_j\ge0$ put
\[
\mathbf A_l=\int\bigotimes_j\Big(\mathbf M^K_{j,l}(\alpha)^{\otimes a_j}\otimes\mathbf M^\Psi_{j,l}(\alpha)^{\otimes b_j}\Big)\,\phi_l(d\alpha),
\]
one $\alpha$ in all the factors. Then
\[
\int\prod_jK_j^{a_j}\Psi_j^{b_j}\prod_{l<L}\phi_l(d\alpha_{B_l})=\Big(\bigotimes_j\mathbf e_j^{\otimes(a_j+b_j)}\Big)^{\mathsf T}\mathbf A_0\mathbf A_1\cdots\mathbf A_{L-1}\Big(\bigotimes_j\boldsymbol\theta_{j,\varnothing}^{\otimes(a_j+b_j)}\Big).
\]
For a chain of identical steps, with one pair of matrix functions $\mathbf M^K_j,\mathbf M^\Psi_j$ for each table, the states of all the $S_l$ indexed by one set, $\boldsymbol\theta_{j,\varnothing}$ vanishing at the states that do not occur at $S_L$, and one measure $\phi$, every $\mathbf A_l$ is one matrix $\mathbf A$: the moment sequence along $L$ is the power of one matrix, of size $\prod_j(\text{number of states of table }j)^{a_j+b_j}$, and satisfies a linear recurrence with constant coefficients of that order. With independent laws $\operatorname{Gamma}_{m_x}$ on the marks (Definition~\ref{def:conv}(vii)) the moment $\oE_{U\sim\Dir_{\mvec}}\big[\prod_jK_j(U)^{a_j}\Psi_j(U)^{b_j}\big]$ is that integral divided by the rising factorial $(m_{\cE})_{\sum_j(a_jr_j+b_jr^\ast_j)}$; with the law $\Dir_{\mvec|_{B_l}}$ on each step the entries are the Dirichlet moments $\prod_{x\in B_l}(m_x)_{k_x}/(m_{B_l})_{|\mathbf k|}$ of the monomials $\alpha_{B_l}^{\mathbf k}$ (Theorem~\ref{thm:dirmoments}).
\end{theorem}

\section{A bunch of any width}\label{sec:anywidth}

% Reference 20: thm:anywidth; source line 6099 (the display at width two, and the free share at width one for one table).
% Proved in the last part: thm:anywidth.
\begin{theorem}[Gluing across a bunch of any width]\label{thm:anywidth}
Let $F$ be a family with real data on its balance wall with every proper wall positive, $B$ a bunch of any width with room $R$, and fix level factorizations with nonnegative terms (Definition~\ref{def:levelweights}), the children factorization for instance, with the bunch letters $V_{j,k,c}$ and the room letters $L_{j,k,c}$; drop the terms whose bunch letter or room letter vanishes identically, and let $C_j$ be the set of the terms $(k,c)$ left for table $j$. Open every active letter with one share per term, $a_{j,k,c}$ for the term $V_{j,k,c}L_{j,k,c}$ of $K_j$ with $\sum_{k,c}a_{j,k,c}=s_j$ and $b_{j,k,c}$ for the term $V^\ast_{j,k,c}L^\ast_{j,k,c}$ of $\Psi_j$ with $\sum_{k,c}b_{j,k,c}=t_j$, pinned by the scale of $B$,
\[
\sum_{j,k,c}\Big(\big(k-d_j(B)\big)a_{j,k,c}+\big(r_j(B)-k\big)b_{j,k,c}\Big)=\kappa_B .
\]
On every contour of positive real parts at which the two integrals below converge,
\begin{multline*}
P_F(\mvec;\mathbf s,\mathbf t)=\int_{|\mathbf a_j|=s_j,\ |\mathbf b_j|=t_j,\ \text{pin}}\ \prod_j\Gst{|C_j|}1{\mathbf a_j}{s_j}\Gst{|C_j|}1{\mathbf b_j}{t_j}\\
\times\int_{\Delta_B}u^{\mvec|_B-\one}\prod_{j,k,c}V_{j,k,c}^{-a_{j,k,c}}V^{\ast}_{j,k,c}{}^{-b_{j,k,c}}\,d\omega_B(u)\int_{\Delta_R}v^{\mvec|_R-\one}\prod_{j,k,c}L_{j,k,c}^{-a_{j,k,c}}L^{\ast}_{j,k,c}{}^{-b_{j,k,c}}\,d\omega_R(v),
\end{multline*}
the two inner integrals projective integrals of degree zero, the periods of the \emph{bunch word} and the \emph{room word} at the shares. Their letters are the bunch and room letters of the level factorization, which need not be the letters of tables, so their convergence is that of the integrals themselves. The free shares are the terms less the active letters less one. For one table at a full lean $(0,t)$ in level form, one station per level, with the share $b_k$ and the bunch and room letters $V^\ast_k$, $L^\ast_k$ of the station of the level $k$, $d(B)\le k\le r(B)$, the pins are $\sum_kb_k=t$ and $\sum_k(r(B)-k)\,b_k=\kappa_B$, the free shares are the $c(B)-1$ shares of the levels strictly between $d(B)$ and $r(B)$, and on a contour of such shares
\begin{multline*}
P_F(\mvec;0,t)=\frac1{(2\pi i)^{c(B)-1}}\int\frac{\prod_{k=d(B)}^{r(B)}\Gamma(b_k)}{c(B)\,\Gamma(t)}\int_{\Delta_B}u^{\mvec|_B-\one}\prod_kV^\ast_k{}^{-b_k}\,d\omega_B(u)\\
\times\int_{\Delta_R}v^{\mvec|_R-\one}\prod_kL^\ast_k{}^{-b_k}\,d\omega_R(v)\prod_{d(B)<k<r(B)}db_k,
\end{multline*}
with $b_{d(B)}=\frac1{c(B)}\big(\kappa_B-\sum_{d(B)<k<r(B)}(r(B)-k)\,b_k\big)$ and $b_{r(B)}=t-b_{d(B)}-\sum_{d(B)<k<r(B)}b_k$. At $c(B)=2$, with $b_{d(B)+1}=z$, the shares are $\frac{\kappa_B-z}2$, $z$ and $\frac{2t-\kappa_B-z}2$; at $c(B)=1$ no share is free, and the period is $\Gst21{\kappa_B,\ t-\kappa_B}{t}$ times the two inner periods (Corollary~\ref{cor:pinnedproduct}). At width one with factoring weights, every letter active, the free shares are $2N-1$, and this is Theorem~\ref{thm:twosides}.
\end{theorem}

For one table in level form at a full lean $(0,t)$ the scale can also be integrated before the shapes. Writing $\Psi(\rho u,v)=\rho^{d^\ast(B)}\sum_{i=0}^{c(B)}\rho^iV^\ast_{r(B)-i}(u)L^\ast_{r(B)-i}(v)$ and rescaling $\rho$ so that the two extreme terms become $1$ and $x^{c(B)}$, the period is the integral over $\Delta_B\times\Delta_R$ of a form that factors into a bunch part and a room part, times the scale integral $\int_0^\infty x^{\kappa_B-1}\big(1+\sum_{i=1}^{c(B)-1}\zeta_ix^i+x^{c(B)}\big)^{-t}dx$, whose coefficients
\[
\zeta_i=\frac{V^\ast_{r(B)-i}(u)}{V^\ast_{r(B)}(u)^{1-i/c(B)}\,V^\ast_{d(B)}(u)^{i/c(B)}}\cdot\frac{L^\ast_{r(B)-i}(v)}{L^\ast_{r(B)}(v)^{1-i/c(B)}\,L^\ast_{d(B)}(v)^{i/c(B)}}
\]
are each a bunch part times a room part: the room receives $c(B)-1$ ratios. At $c(B)=1$ the scale integral is $\Gst21{\kappa_B,\ t-\kappa_B}{t}$, which is Theorem~\ref{thm:cornerfactor}; its Mellin transform in $(\zeta_1,\dots,\zeta_{c(B)-1})$ at $(z_1,\dots,z_{c(B)-1})$ is $\prod_{k=d(B)}^{r(B)}\Gamma(b_k)/\big(c(B)\Gamma(t)\big)$ at the shares $b_{r(B)-i}=z_i$ of Theorem~\ref{thm:anywidth}, its share kernel with the factor $1/c(B)$ of the pin, at $c(B)=2$ the kernel $\Gamma(z)\Gamma\big(\frac{\kappa_B-z}2\big)\Gamma\big(\frac{2t-\kappa_B-z}2\big)/\big(2\Gamma(t)\big)$. With the factorization taken to be the canonical one of the seam of the bunch, the bunch word is written through one law, which with the wall is exactly what the bunch contributes (Chapter~\ref{ch:width}).


\chapter{The law of a cell}\label{ch:hinged}

A table with a bunch of width one across which its weights factor is two cells glued along a hinge that neither of them shows to the other (Lemma~\ref{lem:twosumpieces}). What a cell family must carry so that its contribution to every family built on it is fixed is its hinge wall and one measure on $(0,\infty)^{2N-1}$, the law of its ratio vector under its base form, for one cell a measure on the half-line, and nothing less (Theorem~\ref{thm:retained}). A cell family enters a period only through that law (Theorem~\ref{thm:glue}); the law is read in two ways, by the transform of the cell family and by its dual, and the two transforms expand one function on the two sides of one strip (Proposition~\ref{prop:faces}). The step that makes all of this compose is Lemma~\ref{lem:absorb}: a cell sitting at a mark of any family can be integrated away, leaving the family with that mark twisted by the ratio vector of the cell family; pooling (Theorem~\ref{thm:pooling}) and Theorem~\ref{thm:corereads}, in which the core sees each cell as one mark, are its cases. The series closure of any cells, at any positive exponents, is one integral over the shared scales of the product of the transforms of the cells, so identical cells enter as one transform to a power (Theorem~\ref{thm:seriesclosure}). A core with cells at its markers is read the same way: the law of a substitution is the mixture, over the laws of its cell families, of the laws of the twisted core family (Theorem~\ref{thm:node}); at a series node the transforms multiply and at a parallel node the dual transforms multiply (Theorem~\ref{thm:seriesparallel}); and the period of a family built along a common hinge tree is read from the leaves up, with one law per edge, computed once for each distinct subtree (Theorem~\ref{thm:treeperiod}, Proposition~\ref{prop:work}). The calculus extends to bunches of every width (Chapter~\ref{ch:width}).

\section{Cells and their laws}

% Reference 20: def:hinged; source line 6128 (with the family objects of lem:absorb, source line 6161; the transforms and the hinge measures of one plain cell).
% Proved in the last part: def:hinged.
\begin{definition}[The law of a cell family]\label{def:hinged}
A \emph{cell family} is a family $F=((\mathcal A_j,\theta_j);s_j,t_j)_{j=1}^N$ on $B\sqcup\{g\}$ whose tables are cells with the one hinge $g$ (Definition~\ref{def:substitution}), with real exponents, carrying masses $\mvec$ on $B$; one cell $\mathcal A$ at the exponents $(s,t)$ is the cell family $((\mathcal A,\theta);s,t)$, and its objects carry the subscript $\mathcal A$. Its \emph{hinge wall} is
\[
\kappa_F=m_B-\sum_j\big(s_j(r_{\mathcal A_j}-1)+t_j(|B|-r_{\mathcal A_j})\big),
\]
for one cell the hinge wall $\kappa_{\mathcal A}$ of Lemma~\ref{lem:twosumpieces}. Its \emph{base form} and its \emph{ratio vector} are
\begin{gather*}
\phi_F(du)=u^{\mvec-\one}\prod_jK_{\mathcal A_j/g}(u)^{-s_j}\Psi_{\mathcal A_j\setminus g}(u)^{-t_j}\,\gamma^K_{\mathcal A_1}(u)^{-\kappa_F}\,d\omega_B(u)\quad\text{on }\Delta_B,\\
\boldsymbol\varrho_F=\Big(\frac{\gamma^K_{\mathcal A_j}}{\gamma^K_{\mathcal A_1}},\ \frac{\gamma^\Psi_{\mathcal A_j}}{\gamma^K_{\mathcal A_1}}\Big)_j:\ \Delta_B\to(0,\infty)^{2N-1},
\end{gather*}
the coordinate $\gamma^K_{\mathcal A_1}/\gamma^K_{\mathcal A_1}=1$ omitted, and its \emph{law} is the image $(\boldsymbol\varrho_F)_\ast\phi_F$, a positive measure on $(0,\infty)^{2N-1}$. For $\mathbf z=(z_{j,K},z_{j,\Psi})_j$, also written as the pair $(\mathbf z_K,\mathbf z_\Psi)$ of $\mathbf z_K=(z_{j,K})_j$ and $\mathbf z_\Psi=(z_{j,\Psi})_j$, write
\[
\boldsymbol\varrho_F^{-\mathbf z}=\prod_{j\ge2}\Big(\frac{\gamma^K_{\mathcal A_j}}{\gamma^K_{\mathcal A_1}}\Big)^{-z_{j,K}}\prod_j\Big(\frac{\gamma^\Psi_{\mathcal A_j}}{\gamma^K_{\mathcal A_1}}\Big)^{-z_{j,\Psi}},
\]
the coordinates of the ratio vector to the powers $-z_{j,K}$ ($j\ge2$) and $-z_{j,\Psi}$; the \emph{Mellin transform} of the law, $\int_{\Delta_B}\boldsymbol\varrho_F^{-\mathbf z}\,d\phi_F$, lies in $[0,\infty]$ for real $\mathbf z$ and is defined for complex $\mathbf z$ at whose real part it is finite, with modulus at most that value. For one cell the ratio vector is the ratio $\varrho_{\mathcal A}$ (Definition~\ref{def:substitution}), the base form is
\[
\phi_{\mathcal A}(du)=u^{\mvec-\one}K_{\mathcal A/g}(u)^{\kappa_{\mathcal A}-s}K_{\mathcal A\setminus g}(u)^{-\kappa_{\mathcal A}}\Psi_{\mathcal A\setminus g}(u)^{-t}\,d\omega_B(u),
\]
the law $(\varrho_{\mathcal A})_\ast\phi_{\mathcal A}$ is a measure on $(0,\infty)$, and, with $z=z_{1,\Psi}$, the Mellin transform is $\int_{\Delta_B}\varrho_{\mathcal A}^{-z}\,d\phi_{\mathcal A}=\int_0^\infty x^{-z}\,d\big((\varrho_{\mathcal A})_\ast\phi_{\mathcal A}\big)(x)$. For a cell family with $s_j,t_j>0$ the \emph{transform}, the \emph{dual transform} and the two \emph{hinge measures} are, for $\mathbf x,\mathbf y\in(0,\infty)^N$,
\begin{gather*}
f_F(\mathbf x,\mathbf y)=\int_{\mathbb R^B_{>0}}\alpha^{\mvec-\one}\prod_jK_{\mathcal A_j\setminus g}^{-s_j}\Psi_{\mathcal A_j\setminus g}^{-t_j}\exp\Big(-\sum_j\frac{x_j}{\gamma^K_{\mathcal A_j}}-\sum_jy_j\,\gamma^\Psi_{\mathcal A_j}\Big)\,d\alpha,\\
\check f_F(\mathbf x,\mathbf y)=\int_{\mathbb R^B_{>0}}\alpha^{\mvec-\one}\prod_jK_{\mathcal A_j/g}^{-s_j}\Psi_{\mathcal A_j/g}^{-t_j}\exp\Big(-\sum_j\frac{x_j}{\gamma^\Psi_{\mathcal A_j}}-\sum_jy_j\,\gamma^K_{\mathcal A_j}\Big)\,d\alpha,
\end{gather*}
with values in $[0,\infty]$, the measure $\mathfrak r_F$, the image of $\alpha^{\mvec-\one}\prod_jK_{\mathcal A_j\setminus g}^{-s_j}\Psi_{\mathcal A_j\setminus g}^{-t_j}\,d\alpha$ under $\alpha\mapsto\big((1/\gamma^K_{\mathcal A_j}(\alpha))_j,(\gamma^\Psi_{\mathcal A_j}(\alpha))_j\big)$, and the measure $\check{\mathfrak r}_F$, the image of $\alpha^{\mvec-\one}\prod_jK_{\mathcal A_j/g}^{-s_j}\Psi_{\mathcal A_j/g}^{-t_j}\,d\alpha$ under $\alpha\mapsto\big((1/\gamma^\Psi_{\mathcal A_j}(\alpha))_j,(\gamma^K_{\mathcal A_j}(\alpha))_j\big)$, positive measures on $(0,\infty)^{2N}$. For $c>0$, $f_F(\mathbf x,\mathbf y)=c^{\kappa_F-|\mathbf s|}f_F(\mathbf x/c,c\mathbf y)$ and $\check f_F(\mathbf x,\mathbf y)=c^{\kappa_F-|\mathbf t|}\check f_F(\mathbf x/c,c\mathbf y)$, so each is fixed by its values at $x_1=1$. For one cell they are read there: the transform and the dual transform are $f_{\mathcal A}(y)=f_F(1,y)$ and $\check f_{\mathcal A}(y)=\check f_F(1,y)$, $y>0$, and the hinge measures $\mathfrak r^\Psi_{\mathcal A}$ and $\mathfrak r^K_{\mathcal A}$ are the images of $e^{-\lambda}\mathfrak r_F(d\lambda,d\mu)$ and of $e^{-\lambda}\check{\mathfrak r}_F(d\lambda,d\mu)$ under $(\lambda,\mu)\mapsto\mu$, positive measures on $(0,\infty)$.
\end{definition}

The base form of a cell is the restriction to $\Delta_B$ of the orthant form $x^{\mvec-\one}K_{\mathcal A/g}^{\kappa_{\mathcal A}-s}K_{\mathcal A\setminus g}^{-\kappa_{\mathcal A}}\Psi_{\mathcal A\setminus g}^{-t}dx$, homogeneous of degree $m_B+(\kappa_{\mathcal A}-s)(r_{\mathcal A}-1)-\kappa_{\mathcal A}r_{\mathcal A}-t(|B|-r_{\mathcal A})=0$, so it is read on any section of a positive function of degree one. When $\mathcal A$ is the cell of a bunch $B$ of width one in a table, its hinge minors are the restriction $\cB|B$ and the contraction $\cB/R$, $\kappa_{\mathcal A}$ is the wall $\kappa_B$, $\phi_{\mathcal A}$ is the base form $\phi_B$, $(\varrho_{\mathcal A})_\ast\phi_{\mathcal A}$ is the law of the bunch, $\int\varrho_{\mathcal A}^{-b^d}d\phi_{\mathcal A}$ is the continued period of the bunch family at the share $b^d$, and $f_{\mathcal A}$ is the bunch transform $f_B$ (Definitions~\ref{def:nodetransforms} and~\ref{def:lawbunch}, Lemma~\ref{lem:orthant}); the cell kernel of Definition~\ref{def:twistedcore} is $(\gamma^K_{\mathcal A})^{\kappa_{\mathcal A}}\phi_{\mathcal A}$, and for a cell family $(\gamma^K_{\mathcal A_1})^{\kappa_F}\phi_F$. None of these objects mentions a room: a cell carries them by itself, and they are all that a room ever sees of it (Theorem~\ref{thm:retained}). The hinge wall is read from the other side of the hinge: with the mass $s+t-\kappa_{\mathcal A}$ on $g$ the table $\mathcal A$ itself is on its balance wall, and in a two-sum on its balance wall that mass is the hinge wall of the other cell (Lemma~\ref{lem:twosumpieces}).

\section{Absorbing a cell}



% Reference 20: lem:absorb; source line 6161
\begin{lemma}[Absorbing a cell]\label{lem:absorb}
Let $F=((\mathcal A_j,\theta_j);s_j,t_j)_{j=1}^N$ be a cell family on $B\sqcup\{g\}$ with masses $\mvec$ on $B$, its hinge wall $\kappa_F$, base form $\phi_F$ and ratio vector $\boldsymbol\varrho_F$, whose coordinates are written $x_{j,K}$ and $x_{j,\Psi}$, $x_{1,K}=1$. Let $Z$ be a finite set disjoint from $B\sqcup\{g\}$, $\mvec_Z$ masses on $Z$, and $f_1,\dots,f_p$ positive continuous functions on $(0,\infty)^Z\times(0,\infty)$, each assigned to one letter of $F$, with real exponents $a_i$ such that the exponents assigned to each letter add to its exponent, $s_j$ for $K_{\mathcal A_j}$ and $t_j$ for $\Psi_{\mathcal A_j}$. In this statement write $L_i$, $\gamma_i$ and $x_i$ for the letter of the hinge minor, the hinge value and the coordinate of the letter of $f_i$: $L_i=K_{\mathcal A_j/g}$, $\gamma_i=\gamma^K_{\mathcal A_j}$ and $x_i=x_{j,K}$ when it is $K_{\mathcal A_j}=K_{\mathcal A_j/g}\,(\gamma^K_{\mathcal A_j}+\alpha_g)$, and $L_i=\Psi_{\mathcal A_j\setminus g}$, $\gamma_i=\gamma^\Psi_{\mathcal A_j}$ and $x_i=x_{j,\Psi}$ when it is $\Psi_{\mathcal A_j}=\Psi_{\mathcal A_j\setminus g}\,(\gamma^\Psi_{\mathcal A_j}+\alpha_g)$. Let $H\ge0$ be a Borel function on $(0,\infty)^p$. If the form
\[
\omega=\alpha_B^{\mvec-\one}\alpha_Z^{\mvec_Z-\one}\,H\Big(\big(f_i(\alpha_Z,\gamma_i(\alpha_B))\big)_i\Big)\prod_{i=1}^p\Big[L_i(\alpha_B)\,f_i\big(\alpha_Z,\gamma_i(\alpha_B)\big)\Big]^{-a_i}d\alpha_B\,d\alpha_Z
\]
is homogeneous of degree zero on $(0,\infty)^{B\sqcup Z}$, then, in $[0,\infty]$,
\begin{gather*}
\int_{\{h=1\}}\omega=\int_{(0,\infty)^{2N-1}}\Big(\int_{\{h'=1\}}\omega_{\mathbf x}\Big)\,d\big((\boldsymbol\varrho_F)_\ast\phi_F\big)(\mathbf x),\\
\omega_{\mathbf x}=\alpha_Z^{\mvec_Z-\one}\alpha_g^{\kappa_F-1}\,H\Big(\big(f_i(\alpha_Z,x_i\alpha_g)\big)_i\Big)\prod_if_i(\alpha_Z,x_i\alpha_g)^{-a_i}\,d\alpha_Z\,d\alpha_g,
\end{gather*}
the integrals taken over any sections of positive functions $h$, $h'$ of degree one on the two orthants.
\end{lemma}

For one table the ratio vector is the ratio of the cell. The lemma says that whatever the rest of a family does with a cell, provided it sees the cell through the cell's letters at the hinge, it sees it only through the law: the cell collapses to its hinge, which keeps the mass $\kappa_F$, and the letters that see the hinge see it stretched by the coordinates of the ratio vector. Pooling (Theorem~\ref{thm:pooling}) is the case of one cell family read on the shape before the law is formed; Theorem~\ref{thm:corereads}, in which the core sees each cell as one mark, is the lemma at every cell of a core, with the hinge values of the cells kept as weights of the core instead of their ratios.

\section{Closures}\label{sec:closures}

% Reference 20: thm:seriesclosure; source line 6036 (the responses written through one Gramma at the centre of one table).
% Proved in the last part: thm:seriesclosure.
\begin{theorem}[The series closure at positive exponents]\label{thm:seriesclosure}
Let every table of $F$ be a series closure, $(\cB_j,\theta_j)=\binom{\{h_1,\dots,h_q\}}{q-1}\langle\mathcal A_{1j},\dots,\mathcal A_{qj}\rangle$ (Corollary~\ref{cor:blocksclosures}(ii)), with exponents $s_j,t_j>0$, and let $f_{F_i}$ be the transform of the cell family $F_i=((\mathcal A_{ij},\theta^{\mathcal A_{ij}});s_j,t_j)_j$ of the cells at $B_i$, on $B_i\sqcup\{g_i\}$ (Definition~\ref{def:hinged}). On the balance wall, in $[0,\infty]$,
\[
P_F(\mvec;\mathbf s,\mathbf t)=\prod_j\frac1{\Gamma(s_j)\Gamma(t_j)}\int\prod_jx_j^{s_j-1}\prod_{j\ge2}y_j^{t_j-1}\ \prod_{i=1}^qf_{F_i}\big(\mathbf x,(1,y_2,\dots,y_N)\big)\ d\mathbf x\,dy_2\cdots dy_N,
\]
the integral over $(0,\infty)^N\times(0,\infty)^{N-1}$: the cells communicate only through $2N-1$ shared scales. A cell entering $k$ times enters as its transform to the power $k$: one transform per kind of cell, its multiplicity a power, the order of the cells immaterial; for one table the closure is $\frac1{\Gamma(s)\Gamma(t)}\int_0^\infty y^{t-1}\prod_if_{\mathcal A_i}(y)\,dy$. With $|\mathbf s|=\sum_js_j$ and the wall $\kappa_{B_i}$ of the cell $B_i$, its hinge wall $\kappa_{F_i}$ (Theorem~\ref{thm:substwalls}), each transform is one Gramma integrated over the shape of the cell,
\begin{multline*}
f_{F_i}(\mathbf x,\mathbf y)=\int_{\Delta_{B_i}}\sG_{|\mathbf s|-\kappa_{B_i};(1,-1)}\Big(\sum_j\frac{x_j}{\gamma^K_{\mathcal A_{ij}}(u)},\ \sum_jy_j\,\gamma^\Psi_{\mathcal A_{ij}}(u)\Big)\\
\times u^{\mvec|_{B_i}-\one}\prod_jK_{\mathcal A_{ij}\setminus g_i}(u)^{-s_j}\Psi_{\mathcal A_{ij}\setminus g_i}(u)^{-t_j}\,d\omega_{B_i}(u),
\end{multline*}
against the form of the family of the hinge minors $F_i\setminus g_i=((\mathcal A_{ij}\setminus g_i,\theta^{\mathcal A_{ij}});s_j,t_j)_j$ at the masses $\mvec|_{B_i}$. When that family is on its balance wall, $\kappa_{B_i}=|\mathbf s|$, as at every cell of one table at the centre, $s=t=\sigma$ with every mass $\sigma$, the index is $0$, the Gramma depends on its two coefficients through their product, $\sG_{0;(1,-1)}(a,b)=\sG_{0;(1,-1)}(1,ab)$ (Theorem~\ref{thm:grammalaws}(iii)), and, when the family of the hinge minors has every proper wall positive,
\[
f_{F_i}(\mathbf x,\mathbf y)=P_{F_i\setminus g_i}(\mvec|_{B_i};\mathbf s,\mathbf t)\ \oE_{U\sim\mathrm{Breta}^\circ(F_i\setminus g_i;\mvec|_{B_i};\mathbf s,\mathbf t)}\Big[\sG_{0;(1,-1)}\Big(1,\ \sum_j\frac{x_j}{\gamma^K_{\mathcal A_{ij}}(U)}\sum_jy_j\,\gamma^\Psi_{\mathcal A_{ij}}(U)\Big)\Big],
\]
for one table $P_{\mathcal A_i\setminus g_i}(\mvec|_{B_i};s,t)\,\oE\big[\sG_{0;(1,-1)}\big(1,xy\,\varrho_{\mathcal A_i}(U)\big)\big]$: the image under the ratio of the form of the hinge minor, read through $\sG_{0;(1,-1)}$. So $k\ge2$ copies of one cell family closed in series have the period $\prod_j(\Gamma(s_j)\Gamma(t_j))^{-1}\int\prod_jx_j^{s_j-1}\prod_{j\ge2}y_j^{t_j-1}f_{F_1}\big(\mathbf x,(1,y_2,\dots,y_N)\big)^k\,d\mathbf x\,dy_2\cdots dy_N$.
\end{theorem}

For one table the closure of $q$ cells is a node of width one (Definition~\ref{def:nodetransforms}) with the bunch $B_1\sqcup\dots\sqcup B_{q-1}$, a series connection (Proposition~\ref{prop:seriesconnection}), against the room $B_q$. Theorem~\ref{thm:onescale} reads it as $\frac1{\Gamma(s)\Gamma(t)}\int_0^\infty f_B(y)f_R(y)\frac{dy}y$, and Theorem~\ref{thm:seriesclosure} is the same integral written symmetrically in all the cells, $\frac1{\Gamma(s)\Gamma(t)}\int_0^\infty y^{t-1}\prod_if_{\mathcal A_i}(y)\,dy$ with the transforms of the cells (Definition~\ref{def:hinged}), which are the orthant transforms $f^\circ_i$ of Proposition~\ref{prop:seriesconnection}; for $k$ equal cells at the centre the formula for $k$ copies is Theorem~\ref{thm:seriesrepeat}(i), with $k-1$ copies at a series node and one copy as the room. The table that is one circuit, $\binom{\cE}{n-1}$, is the series closure of $n$ single marks, each with the transform $\sG_{s-m_e;(1,-1)}(1,y)$ up to its weight (Corollary~\ref{cor:cycle}). The parallel closure is the dual statement (Proposition~\ref{prop:dualfamily}).

\section{Cores and their markers}\label{sec:cores}

A table is built from cells by two-sums at the markers of cores. 

% Reference 20: def:node; source line 6316
\begin{definition}[Markers; the twisted core]\label{def:node}
Let the exponents be $s,t>0$. A \emph{core with markers} is a plain table $\mathcal H$ on $\cE_{\mathcal H}\sqcup G_{\mathcal H}$ none of whose marks is a loop or a coloop, a core of Definition~\ref{def:substitution}, whose \emph{markers} are the marks of $G_{\mathcal H}$, the marks at which cells may be substituted (the marks $h_i$ of Definition~\ref{def:substitution}). \emph{Cells at the markers} are plain cells $\mathcal A_1,\dots,\mathcal A_k$ on $B_i\sqcup\{g_i\}$ whose hinges are identified with distinct markers $g_1,\dots,g_k$, the sets $\cE_{\mathcal H}$ and $B_i$ disjoint, and $\mathcal H\langle\mathcal A_1,\dots,\mathcal A_k\rangle$ is their substitution (Definition~\ref{def:substitution}), the marks of $\cE_{\mathcal H}$ and the markers without a cell left in place. For $\mathbf x\in(0,\infty)^k$ the \emph{twisted core} $\mathcal H[\mathbf x]$ is $\mathcal H$ twisted in its cobasis letter by $x_i$ at every $g_i$ (Definition~\ref{def:twist}), carrying the masses $\kappa_{\mathcal A_i}$ on the markers $g_i$ and the given masses $\mvec|_{\cE_{\mathcal H}}$ on $\cE_{\mathcal H}$; it has the basis weights $\theta_T$ and the cobasis weights $\theta_T\prod_{g_i\notin T}x_i$, and its period is that of a table with two systems of weights (Definition~\ref{def:family}). With the cells at the shapes $u=(u_i)_i$, the twisted core $\mathcal H[(\varrho_{\mathcal A_i}(u_i))_i]$ is the core read at $u$ of Definition~\ref{def:twistedcore} with the scale of every marker $g_i$ divided by $\gamma^K_{\mathcal A_i}(u_i)$. A core with one marker $g_0$ left free is a \emph{hinged core}; then $\mathcal H[\mathbf x]$ is a cell with the hinge $g_0$, not plain, and its hinge wall, base form, ratio and law are those of Definition~\ref{def:hinged}.
\end{definition}

% Reference 20: lem:nodesum; source line 6320
\begin{lemma}[Substitution at the markers]\label{lem:nodesum}
The substitution $\mathcal H\langle\mathcal A_1,\dots,\mathcal A_k\rangle$ does not depend on the order of the two-sums (Proposition~\ref{prop:substitution}(i)). Its bases are the sets $(T\cap\cE_{\mathcal H})\sqcup\bigsqcup_iC_i$ with $T$ a basis of $\mathcal H$ and, for each $i$, $C_i$ a basis of $\mathcal A_i$ avoiding $g_i$ when $g_i\in T$ and $C_i+g_i$ a basis of $\mathcal A_i$ when $g_i\notin T$, with the product weights. Its rank is $r_{\mathcal H}+\sum_i(r_{\mathcal A_i}-1)$, and with $G'$ the markers left free (Theorem~\ref{thm:substletters})
\begin{gather*}
K_{\mathcal H\langle\mathcal A_1,\dots,\mathcal A_k\rangle}=\prod_iK_{\mathcal A_i/g_i}\cdot K_{\mathcal H}\big(\alpha_{\cE_{\mathcal H}},(\gamma^K_{\mathcal A_i})_i,\alpha_{G'}\big),\\
\Psi_{\mathcal H\langle\mathcal A_1,\dots,\mathcal A_k\rangle}=\prod_i\Psi_{\mathcal A_i\setminus g_i}\cdot\Psi_{\mathcal H}\big(\alpha_{\cE_{\mathcal H}},(\gamma^\Psi_{\mathcal A_i})_i,\alpha_{G'}\big).
\end{gather*}
For a hinged core the letters of the two hinge minors $\mathcal H\langle\mathcal A_1,\dots,\mathcal A_k\rangle\setminus g_0$ and $\mathcal H\langle\mathcal A_1,\dots,\mathcal A_k\rangle/g_0$ are the same products with the letters of $\mathcal H\setminus g_0$ and $\mathcal H/g_0$ in place of $K_{\mathcal H}$ and $\Psi_{\mathcal H}$.
\end{lemma}

% Reference 20: thm:node; source line 6332 (one table; the gluing of a family in rem:familiesnodes).
% Proved in the last part: thm:node.
\begin{theorem}[The law of a substitution]\label{thm:node}
Let $F$ be a family with real data whose table $j$ is $\mathcal H_j\langle\mathcal A_{1j},\dots,\mathcal A_{kj}\rangle$, the cores $\mathcal H_j$ plain on $\cE_{\mathcal H}\sqcup G_{\mathcal H}$ with $\{g_1,\dots,g_k\}\subseteq G_{\mathcal H}$, the cells of all tables at $g_i$ forming the cell family $F_i=((\mathcal A_{ij},\theta^{\mathcal A_{ij}});s_j,t_j)_j$ on $B_i\sqcup\{g_i\}$ (Definition~\ref{def:hinged}), with masses on $\cE_{\mathcal H}$ and on the marks of the cells. For $\mathbf x=(\mathbf x_i)_i$, $\mathbf x_i=(x_{i,j,K},x_{i,j,\Psi})_j\in(0,\infty)^{2N-1}$ with $x_{i,1,K}=1$, let $F_{\mathcal H}[\mathbf x]$ be the family $((\mathcal H_j,\theta^{\mathcal H_j});s_j,t_j)_j$ twisted at every $g_i$ by $\mathbf x_i$, $x_{i,j,K}$ in $K_j$ and $x_{i,j,\Psi}$ in $\Psi_j$ (Definition~\ref{def:twist}), with the masses $\kappa_{F_i}$ on the markers $g_i$ and the given masses on $\cE_{\mathcal H}$.
\begin{enumerate}[label=(\roman*)]
\item If $G_{\mathcal H}=\{g_1,\dots,g_k\}$ and $F$ is on its balance wall, then, in $[0,\infty]$,
\[
P_F(\mvec;\mathbf s,\mathbf t)=\int P_{F_{\mathcal H}[\mathbf x]}\big((\mvec|_{\cE_{\mathcal H}},\kappa_{F_1},\dots,\kappa_{F_k});\mathbf s,\mathbf t\big)\prod_id\big((\boldsymbol\varrho_{F_i})_\ast\phi_{F_i}\big)(\mathbf x_i).
\]
With the core $\binom{\{g_1,g_2\}}1$ in every table this is the gluing of two cell families $F_1$ and $F_2$, $\kappa_{F_1}+\kappa_{F_2}=2\sigma$, $\sigma$ the total strength (Definition~\ref{def:family}): the period of the twisted core family is $\int_0^\infty v^{\kappa_{F_1}-1}\prod_j(x_{1,j,K}v+x_{2,j,K})^{-s_j}(x_{1,j,\Psi}v+x_{2,j,\Psi})^{-t_j}\,dv$, a Breta of one bunch with at most $2N$ coefficients (Proposition~\ref{prop:scalebreta}), unchanged when $(\kappa_{F_1},\mathbf x_1)$ and $(\kappa_{F_2},\mathbf x_2)$ are swapped.
\item If $G_{\mathcal H}=\{g_0,g_1,\dots,g_k\}$, $F$ is a cell family with the hinge $g_0$; its hinge wall is that of $F_{\mathcal H}[\mathbf x]$, which does not depend on $\mathbf x$,
\[
\kappa_F=m_{\cE_{\mathcal H}}+\sum_i\kappa_{F_i}-\sum_j\big(s_j(r_{\mathcal H_j}-1)+t_j(|\cE_{\mathcal H}|+k-r_{\mathcal H_j})\big),
\]
and its law is the mixture of the twisted laws: for every Borel $H\ge0$ on $(0,\infty)^{2N-1}$,
\[
\int H\,d\big((\boldsymbol\varrho_F)_\ast\phi_F\big)=\int\Big(\int H\,d\big((\boldsymbol\varrho_{F_{\mathcal H}[\mathbf x]})_\ast\phi_{F_{\mathcal H}[\mathbf x]}\big)\Big)\prod_id\big((\boldsymbol\varrho_{F_i})_\ast\phi_{F_i}\big)(\mathbf x_i).
\]
In particular $\int\boldsymbol\varrho_F^{-\mathbf z}d\phi_F=\int\big(\int\boldsymbol\varrho_{F_{\mathcal H}[\mathbf x]}^{-\mathbf z}d\phi_{F_{\mathcal H}[\mathbf x]}\big)\prod_id\big((\boldsymbol\varrho_{F_i})_\ast\phi_{F_i}\big)(\mathbf x_i)$ for every real $\mathbf z$.
\end{enumerate}
For one table with $s,t>0$ the laws are measures on $(0,\infty)$ and the twisted core family is the twisted core $\mathcal H[\mathbf x]$ of Definition~\ref{def:node}.
\end{theorem}

The core is the only place where cells interact. Below a core each cell family is one measure on $(0,\infty)^{2N-1}$, for one table on a half-line; at the core the cell families are fed into the core family as twists of both letters of every table at the markers, of the cobasis letter alone for one table (Lemma~\ref{lem:absorb}), and the core family, with the masses $\kappa_{F_i}$ on its markers, is integrated once for each value of the twists. So the work at a core is governed by the core itself and by the number of its markers, not by the sizes of the cells: a cell with a thousand marks and a cell with one mark cost the same once their laws are known. Theorem~\ref{thm:twosides} is the gluing of two cell families read through Mellin transforms (Theorem~\ref{thm:glue}(iv)).

\section{Two cells glued}

% Reference 20: lem:twomarks; source line 6192
\begin{lemma}[Two marks]\label{lem:twomarks}
For $s,t>0$, $\kappa,\kappa'>0$ with $\kappa+\kappa'=s+t$, and $x,x'>0$, let $\binom{\{g,g'\}}1[(x,x')]$ be the list $\binom{\{g,g'\}}1$ twisted in its cobasis letter by $x$ at $g$ and by $x'$ at $g'$ (Definition~\ref{def:twist}). With the masses $\kappa$ on $g$ and $\kappa'$ on $g'$ its period is
\begin{align*}
P_{\binom{\{g,g'\}}1[(x,x')]}\big((\kappa,\kappa');s,t\big)&=\int_0^\infty v^{\kappa-1}(1+v)^{-s}(xv+x')^{-t}\,dv\\
&=x^{-t}\,\Gst22{\kappa,\ \kappa'}{s,\ t}\,\sB^{\{\{0,1\}\}}_{(s,t),\,\kappa';\,I_{\{0,1\}}}\big((1,x'/x)\big)\\
&=\frac1{\Gamma(s)\Gamma(t)}\int_0^\infty y^{t-1}\,\sG_{s-\kappa;(1,-1)}(1,xy)\,\sG_{s-\kappa';(1,-1)}(1,x'y)\,dy
\end{align*}
(Proposition~\ref{prop:scalebreta}, the two coefficients merging into one of letter exponent $s+t$ when $x=x'$), unchanged when $(\kappa,x)$ and $(\kappa',x')$ are swapped. For a family of tables the kernel is the period of the twisted core family of Theorem~\ref{thm:node}(i) with the core $\binom{\{g_1,g_2\}}1$ in every table.
\end{lemma}

% Reference 20: prop:faces; source line 6236 (one plain cell).
% Proved in the last part: prop:faces.
\begin{proposition}[Faces of one law]\label{prop:faces}
Let $F=((\mathcal A_j,\theta_j);s_j,t_j)_{j=1}^N$ be a cell family on $B\sqcup\{g\}$ with plain tables, $s_j,t_j>0$, masses $\mvec$ on $B$ and the hinge wall $\kappa=\kappa_F$; in this statement the points of $(0,\infty)^{2N-1}$ are written $\mathbf z=(z_{j,K},z_{j,\Psi})_j$, $z_{1,K}=1$, and $\mathbf x'=(x_2,\dots,x_N)$.
\begin{enumerate}[label=(\roman*)]
\item \emph{Law forms.} For $\mathbf x,\mathbf y\in(0,\infty)^N$,
\begin{gather*}
f_F(\mathbf x,\mathbf y)=\int\sG_{|\mathbf s|-\kappa;(1,-1)}\Big(\sum_j\frac{x_j}{z_{j,K}},\ \sum_jy_jz_{j,\Psi}\Big)\prod_jz_{j,K}^{-s_j}\,d\big((\boldsymbol\varrho_F)_\ast\phi_F\big)(\mathbf z),\\
\check f_F(\mathbf x,\mathbf y)=\int\sG_{|\mathbf t|-\kappa;(1,-1)}\Big(\sum_jx_j\frac{z_{1,\Psi}}{z_{j,\Psi}},\ \sum_jy_j\frac{z_{j,K}}{z_{1,\Psi}}\Big)\prod_j\Big(\frac{z_{j,\Psi}}{z_{1,\Psi}}\Big)^{-t_j}z_{1,\Psi}^{-\kappa}\,d\big((\boldsymbol\varrho_F)_\ast\phi_F\big)(\mathbf z):
\end{gather*}
the two transforms depend on $F$ only through its hinge wall and its law. For one cell, for $y>0$,
\[
f_{\mathcal A}(y)=\int_{\Delta_B}\sG_{s-\kappa;(1,-1)}\big(1,\varrho_{\mathcal A}y\big)\,d\phi_{\mathcal A},\qquad \check f_{\mathcal A}(y)=\int_{\Delta_B}\sG_{t-\kappa;(1,-1)}\big(1,y/\varrho_{\mathcal A}\big)\,\varrho_{\mathcal A}^{-\kappa}\,d\phi_{\mathcal A}.
\]
\item \emph{Hinge measures.}
\[
f_F(\mathbf x,\mathbf y)=\int e^{-\ip{\mathbf x}{\boldsymbol\lambda}-\ip{\mathbf y}{\boldsymbol\mu}}\,d\mathfrak r_F(\boldsymbol\lambda,\boldsymbol\mu),\qquad \check f_F(\mathbf x,\mathbf y)=\int e^{-\ip{\mathbf x}{\boldsymbol\lambda}-\ip{\mathbf y}{\boldsymbol\mu}}\,d\check{\mathfrak r}_F(\boldsymbol\lambda,\boldsymbol\mu),
\]
and $\mathfrak r_F$ is the image of the product of $\prod_jz_{j,K}^{-s_j}\,(\boldsymbol\varrho_F)_\ast\phi_F(d\mathbf z)$ and $p^{\kappa-|\mathbf s|-1}dp$ under $(\mathbf z,p)\mapsto\big((1/(pz_{j,K}))_j,(pz_{j,\Psi})_j\big)$. So $f_F$ and $\check f_F$ are completely monotone in their $2N$ variables, $(-1)^{|\mathbf n|}\partial^{\mathbf n}f_F\ge0$ for every multi-index $\mathbf n$ on the interior of the set where $f_F$ is finite and the same for $\check f_F$, and log-convex along every segment on which they are finite. For one cell, $f_{\mathcal A}(y)=\int e^{-y\lambda}d\mathfrak r^\Psi_{\mathcal A}(\lambda)$ and $\check f_{\mathcal A}(y)=\int e^{-y\lambda}d\mathfrak r^K_{\mathcal A}(\lambda)$; $\mathfrak r^\Psi_{\mathcal A}$ is the image of $\alpha^{\mvec-\one}K_{\mathcal A\setminus g}^{-s}\Psi_{\mathcal A\setminus g}^{-t}e^{-1/\gamma^K_{\mathcal A}}d\alpha$ under $\gamma^\Psi_{\mathcal A}$ and of $(\varrho_{\mathcal A})_\ast\phi_{\mathcal A}(dx)\otimes p^{s-\kappa-1}e^{-p}dp$ under $(x,p)\mapsto x/p$, and $\mathfrak r^K_{\mathcal A}$ the image of $\alpha^{\mvec-\one}K_{\mathcal A/g}^{-s}\Psi_{\mathcal A/g}^{-t}e^{-1/\gamma^\Psi_{\mathcal A}}d\alpha$ under $\gamma^K_{\mathcal A}$ and of $x^{-\kappa}(\varrho_{\mathcal A})_\ast\phi_{\mathcal A}(dx)\otimes p^{t-\kappa-1}e^{-p}dp$ under $(x,p)\mapsto1/(xp)$; so $f_{\mathcal A}$ and $\check f_{\mathcal A}$ are decreasing and log-convex wherever finite.
\item \emph{Mellin transforms.} For real $\mathbf a,\mathbf b\in(0,\infty)^N$ with $\sum_jb_j-\sum_ja_j=\kappa-|\mathbf s|$, in $[0,\infty]$,
\begin{multline*}
\int_{(0,\infty)^{2N-1}}\prod_{j\ge2}x_j^{a_j-1}\prod_jy_j^{b_j-1}\,f_F\big((1,\mathbf x'),\mathbf y\big)\,d\mathbf x'\,d\mathbf y\\
=\prod_j\Gamma(a_j)\Gamma(b_j)\int_{\Delta_B}\boldsymbol\varrho_F^{-(\mathbf s-\mathbf a,\mathbf b)}\,d\phi_F,
\end{multline*}
and for real $\mathbf a,\mathbf b\in(0,\infty)^N$ with $\sum_jb_j-\sum_ja_j=\kappa-|\mathbf t|$ the same holds for $\check f_F$ with $\int_{\Delta_B}\boldsymbol\varrho_F^{-(\mathbf b,\mathbf t-\mathbf a)}\,d\phi_F$ on the right; both Mellin transforms of the law are taken on the hyperplane $\sum_j(z_{j,K}+z_{j,\Psi})=\kappa$ of their arguments. For one cell, with $z=b$, these are
\begin{gather*}
\int_0^\infty y^{z-1}f_{\mathcal A}(y)\,dy=\Gamma(z)\Gamma(z+s-\kappa)\int_{\Delta_B}\varrho_{\mathcal A}^{-z}\,d\phi_{\mathcal A}\qquad\big(z>\max(0,\kappa-s)\big),\\
\int_0^\infty y^{z-1}\check f_{\mathcal A}(y)\,dy=\Gamma(z)\Gamma(z+t-\kappa)\int_{\Delta_B}\varrho_{\mathcal A}^{z-\kappa}\,d\phi_{\mathcal A}\qquad\big(z>\max(0,\kappa-t)\big).
\end{gather*}
\item \emph{The dual cell family.} The dual tables $\mathcal A^\ast_j$, with the hinge $g$, at the exponents $(t_j,s_j)$ form a cell family $F^\ast$ with the hinge wall $\kappa$, the hinge values $\gamma^K_{\mathcal A^\ast_j}=\gamma^\Psi_{\mathcal A_j}$ and $\gamma^\Psi_{\mathcal A^\ast_j}=\gamma^K_{\mathcal A_j}$, the base form $(\gamma^\Psi_{\mathcal A_1}/\gamma^K_{\mathcal A_1})^{-\kappa}\phi_F$, the law the image of $z_{1,\Psi}^{-\kappa}\,(\boldsymbol\varrho_F)_\ast\phi_F(d\mathbf z)$ under $\mathbf z\mapsto(z_{j,\Psi}/z_{1,\Psi},\,z_{j,K}/z_{1,\Psi})_j$, the transform $\check f_F$ and the dual transform $f_F$. For one cell the ratio is $1/\varrho_{\mathcal A}$, the base form $\varrho_{\mathcal A}^{-\kappa}\phi_{\mathcal A}$, the law the image of $x^{-\kappa}(\varrho_{\mathcal A})_\ast\phi_{\mathcal A}(dx)$ under $x\mapsto1/x$, and $\int\varrho_{\mathcal A^\ast}^{-z}d\phi_{\mathcal A^\ast}=\int\varrho_{\mathcal A}^{z-\kappa}d\phi_{\mathcal A}$.
\item \emph{Change of transform.} For real $\mathbf a,\mathbf b$ with $0<a_j<s_j$, $0<b_j<t_j$ and $\sum_jb_j-\sum_ja_j=\kappa-|\mathbf s|$, the left side of (iii) for $\check f_F$ at $(\mathbf t-\mathbf b,\mathbf s-\mathbf a)$ times $\prod_j\Gamma(a_j)\Gamma(b_j)$ equals that for $f_F$ at $(\mathbf a,\mathbf b)$ times $\prod_j\Gamma(t_j-b_j)\Gamma(s_j-a_j)$, in $(0,\infty]$. For one cell, with $z=s-a$: for real $z$ with $\max(0,\kappa-t)<z<\min(s,\kappa)$, in $(0,\infty]$,
\[
\int_0^\infty y^{z-1}\check f_{\mathcal A}(y)\,dy=\frac{\Gamma(z)\Gamma(z+t-\kappa)}{\Gamma(\kappa-z)\Gamma(s-z)}\int_0^\infty y^{\kappa-z-1}f_{\mathcal A}(y)\,dy .
\]
When $\mathcal A$ is the cell of the bunch of a node of width one (Definition~\ref{def:nodetransforms}) and $\int\varrho_{\mathcal A}^{-b}d\phi_{\mathcal A}<\infty$ at some $b$ with $\max(0,\kappa_{\mathcal A}-s)<b<\min(t,\kappa_{\mathcal A})$, the expansion of $f_{\mathcal A}$ at $y\to0$ is the sum of the residues of $y^{-z}\Gamma(z)\Gamma(z+s-\kappa_{\mathcal A})\int\varrho_{\mathcal A}^{-z}d\phi_{\mathcal A}$ to the left of the line $\re z=b$ (Theorem~\ref{thm:nodeclosing}(i)), and that of $\check f_{\mathcal A}$, the transform of $\mathcal A^\ast$ read as the bunch of the dual node of width one at the exponents $(t,s)$ (Proposition~\ref{prop:dualfamily}), is the sum of the residues of $y^{-z}\Gamma(z)\Gamma(z+t-\kappa_{\mathcal A})\int\varrho_{\mathcal A}^{z-\kappa_{\mathcal A}}d\phi_{\mathcal A}$ to the left of the line $\re z=\kappa_{\mathcal A}-b$, at the poles of the Gamma pair and at the points $\kappa_{\mathcal A}-w$, $w$ a pole of the continued $\int\varrho_{\mathcal A}^{-w}d\phi_{\mathcal A}$ to the right of the line $\re w=b$: the two transforms expand one function on the two sides of its strip.
\end{enumerate}
\end{proposition}

So one law has four faces. The base form and its image, the law, are where the cell family lives; the transform reads the law through Grammas of index $|\mathbf s|-\kappa_F$ and is the Laplace transform of the hinge measure $\mathfrak r_F$, carried by the reciprocals of the hinge values by $K$ and the hinge values by $\Psi$; the dual transform reads the law turned over through Grammas of index $|\mathbf t|-\kappa_F$ and is the Laplace transform of $\check{\mathfrak r}_F$, carried by the reciprocals of the hinge values by $\Psi$ and the hinge values by $K$; and on the Mellin side both are the one function $\int\boldsymbol\varrho_F^{-\mathbf z}d\phi_F$ with two Gamma products. Passing from one transform to the other costs nothing on the Mellin side, a reflection and a ratio of Gammas, and costs integral transforms on the side of the transforms; for one cell, on the half-line, the two expansions at $y\to0$ are read on opposite sides of the strip: the transform expands where the law is large, the dual transform where it is small.

% Reference 20: thm:glue; source line 6210 (two plain cells of one table).
% Proved in the last part: thm:glue.
\begin{theorem}[Two cell families glued]\label{thm:glue}
Let $F_1$ and $F_2$ be cell families with plain tables on $B\sqcup\{g\}$ and $R\sqcup\{g\}$, at the exponents $(s_j,t_j)_j$ with $s_j,t_j>0$, and let $F_1\oplus_gF_2$, whose table $j$ is the two-sum of the tables $j$ of $F_1$ and $F_2$ (Lemma~\ref{lem:twosumpieces}), carry masses $\mvec$ on its balance wall, so that $\kappa_{F_1}+\kappa_{F_2}=2\sigma$, $\sigma$ the total strength (Definition~\ref{def:family}); let $0<\kappa_{F_1}<2\sigma$. In this statement the points of the laws of $F_1$ and $F_2$ are written $\mathbf z$ and $\mathbf z'$, $\mathcal G(\mathbf z,\mathbf z')=\int_0^\infty v^{\kappa_{F_1}-1}\prod_j(z_{j,K}v+z'_{j,K})^{-s_j}(z_{j,\Psi}v+z'_{j,\Psi})^{-t_j}\,dv$ is the period of the twisted core family of Theorem~\ref{thm:node}(i) with the core $\binom{\{g_1,g_2\}}1$ in every table, the transforms are read at $x_1=1$, with $\mathbf x'=(x_2,\dots,x_N)$, and the hinge measures of $F_1$ act in the variables $(\boldsymbol\lambda,\boldsymbol\mu)$, those of $F_2$ in $(\boldsymbol\lambda',\boldsymbol\mu')$. Then, in $[0,\infty]$,
\begin{align*}
P_{F_1\oplus_gF_2}(\mvec;\mathbf s,\mathbf t)&=\iint\mathcal G(\mathbf z,\mathbf z')\,d\big((\boldsymbol\varrho_{F_1})_\ast\phi_{F_1}\big)(\mathbf z)\,d\big((\boldsymbol\varrho_{F_2})_\ast\phi_{F_2}\big)(\mathbf z') \tag{i}\\
&=\prod_j\frac1{\Gamma(s_j)\Gamma(t_j)}\int\prod_{j\ge2}x_j^{s_j-1}\prod_jy_j^{t_j-1}\,f_{F_1}f_{F_2}\,d\mathbf x'\,d\mathbf y\\
&=\prod_j\frac1{\Gamma(s_j)\Gamma(t_j)}\int\prod_{j\ge2}x_j^{t_j-1}\prod_jy_j^{s_j-1}\,\check f_{F_1}\check f_{F_2}\,d\mathbf x'\,d\mathbf y \tag{ii}\\
&=\frac1{\Gamma(s_1)}\iint e^{-\lambda_1-\lambda'_1}\prod_{j\ge2}(\lambda_j+\lambda'_j)^{-s_j}\prod_j(\mu_j+\mu'_j)^{-t_j}\,d\mathfrak r_{F_1}\,d\mathfrak r_{F_2}\\
&=\frac1{\Gamma(t_1)}\iint e^{-\lambda_1-\lambda'_1}\prod_{j\ge2}(\lambda_j+\lambda'_j)^{-t_j}\prod_j(\mu_j+\mu'_j)^{-s_j}\,d\check{\mathfrak r}_{F_1}\,d\check{\mathfrak r}_{F_2}, \tag{iii}
\end{align*}
and these are finite exactly when every proper wall of $F_1\oplus_gF_2$ is positive. On every contour of the level shares of Theorem~\ref{thm:twosides} on which its stacked Gammas have arguments of positive real part and the two Mellin transforms below converge absolutely (Definition~\ref{def:hinged}),
\begin{multline*}
P_{F_1\oplus_gF_2}(\mvec;\mathbf s,\mathbf t)=2^{N-1}\int_{\sum_j\delta^d_j=\frac12(|\mathbf s|-\kappa_{F_1})}\ \prod_j\Gst42{a^d_j,\ b^d_j,\ a^r_j,\ b^r_j}{s_j,\ t_j}\\
\times\int_{\Delta_B}\boldsymbol\varrho_{F_1}^{-(\mathbf a^r,\mathbf b^d)}\,d\phi_{F_1}\int_{\Delta_R}\boldsymbol\varrho_{F_2}^{-(\mathbf a^d,\mathbf b^r)}\,d\phi_{F_2}, \tag{iv}
\end{multline*}
the formula of Theorem~\ref{thm:twosides}, the periods of its bunch family and room family at the level shares being these two Mellin transforms. For one table the laws are $(\varrho_{\mathcal A})_\ast\phi_{\mathcal A}$ and $(\varrho_{\mathcal A'})_\ast\phi_{\mathcal A'}$ and the kernel of (i) is the period of Lemma~\ref{lem:twomarks}; (ii) reads $\frac1{\Gamma(s)\Gamma(t)}\int_0^\infty y^{t-1}f_{\mathcal A}f_{\mathcal A'}\,dy=\frac1{\Gamma(s)\Gamma(t)}\int_0^\infty y^{s-1}\check f_{\mathcal A}\check f_{\mathcal A'}\,dy$, its first form Theorem~\ref{thm:onescale}(i) with the room transform $f_R(y)=y^tf_{\mathcal A'}(y)$; (iii) reads $\frac1{\Gamma(s)}\iint(\lambda+\lambda')^{-t}\,d\mathfrak r^\Psi_{\mathcal A}\,d\mathfrak r^\Psi_{\mathcal A'}=\frac1{\Gamma(t)}\iint(\lambda+\lambda')^{-s}\,d\mathfrak r^K_{\mathcal A}\,d\mathfrak r^K_{\mathcal A'}$; and (iv) holds on a line $\re z=z_0$ with $\max(0,\kappa_{\mathcal A}-s)<z_0<\min(t,\kappa_{\mathcal A})$ on which $\int\varrho_{\mathcal A}^{-z_0}d\phi_{\mathcal A}$ and $\int\varrho_{\mathcal A'}^{z_0-t}d\phi_{\mathcal A'}$ are finite, as
\begin{multline*}
P_{\mathcal A\oplus_g\mathcal A'}(\mvec;s,t)=\frac1{\Gamma(s)\Gamma(t)}\,\frac1{2\pi i}\int_{\re z=z_0}\Gamma(z)\Gamma(z+s-\kappa_{\mathcal A})\Gamma(t-z)\Gamma(\kappa_{\mathcal A}-z)\\
\times\int_{\Delta_B}\varrho_{\mathcal A}^{-z}\,d\phi_{\mathcal A}\int_{\Delta_R}\varrho_{\mathcal A'}^{z-t}\,d\phi_{\mathcal A'}\ dz .
\end{multline*}
\end{theorem}

The period of a two-sum is a pairing of two laws by one kernel, the period of the twisted core family on two marks, for one table the period of two marks. Read through the Grammas the pairing is one integral over the scales of the two transforms, symmetric in the two sides; for one table, read through the hinge measures, it is the moment of order $-t$ of a sum: under $\mathfrak r^\Psi_{\mathcal A}$ the variable is the hinge value $\gamma^\Psi_{\mathcal A}$ of the cell (Proposition~\ref{prop:faces}(ii)), and gluing the two cells into one table adds their hinge values by $\Psi$, as $\Psi=\Psi_{\mathcal A\setminus g}\Psi_{\mathcal A'\setminus g}(\gamma^\Psi_{\mathcal A}+\gamma^\Psi_{\mathcal A'})$ says.

\section{What a cell retains}

% Reference 20: def:onemark; source line 6269 (one table).
% Proved in the last part: def:onemark.
\begin{definition}[One weighted mark]\label{def:onemark}
For $\mathbf x_0=(x_{0,j,K},x_{0,j,\Psi})_{j=1}^N$ with $x_{0,1,K}=1$, a point of $(0,\infty)^{2N-1}$, and a mark $e$ with the mass $m_e$, the \emph{one-mark cell family of ratio vector $\mathbf x_0$} has in table $j$ the bases $\{e\}$ and $\{g\}$, with the weights $x_{0,j,K}$ and $1$ in its letter $K_j$ and the weights $1$ and $x_{0,j,\Psi}$ in its letter $\Psi_j$ (Definition~\ref{def:substitution}): $K_j=x_{0,j,K}\alpha_e+\alpha_g$ and $\Psi_j=\alpha_g+x_{0,j,\Psi}\alpha_e$. Its hinge values are $x_{0,j,K}\alpha_e$ and $x_{0,j,\Psi}\alpha_e$, its ratio vector is the constant $\mathbf x_0$, its hinge wall is $m_e$, and its law is $\delta_{\mathbf x_0}$. For one table and $x_0>0$, the \emph{one-mark cell of ratio $x_0$} is the plain table $\big(\binom{\{e,g\}}1,(1,x_0^{1/2})\big)$ on $\{e,g\}$, the bases $\{e\}$ and $\{g\}$ with the weights $1$ and $x_0^{1/2}$. Its hinge minors have $K_{\setminus g}=\alpha_e$, $K_{/g}=x_0^{1/2}$, $\Psi_{\setminus g}=1$ and $\Psi_{/g}=x_0^{1/2}\alpha_e$, its ratio is the constant $x_0$, its hinge wall is $m_e$, and its law is $x_0^{(m_e-s)/2}\,\delta_{x_0}$.
\end{definition}

% Reference 20: thm:retained; source line 6273 (plain cells of one table).
% Proved in the last part: thm:retained.
\begin{theorem}[The law is what a cell family retains]\label{thm:retained}
Let $F$ and $F'$ be cell families with plain tables, at the exponents $(s_j,t_j)_j$ with $s_j,t_j>0$, with masses on their marks other than the hinge and one hinge wall $\kappa\in(0,2\sigma)$, $\sigma$ the total strength (Definition~\ref{def:family}), and let there be a real $\mathbf z=(z_{j,K},z_{j,\Psi})_j$ with $0<z_{j,K}<s_j$, $0<z_{j,\Psi}<t_j$ and $\sum_j(z_{j,K}+z_{j,\Psi})=\kappa$ at which $\int\boldsymbol\varrho_F^{-\mathbf z}d\phi_F$ and $\int\boldsymbol\varrho_{F'}^{-\mathbf z}d\phi_{F'}$ are finite (Definition~\ref{def:hinged}). The following are equivalent.
\begin{enumerate}[label=(\alph*)]
\item $(\boldsymbol\varrho_F)_\ast\phi_F=(\boldsymbol\varrho_{F'})_\ast\phi_{F'}$.
\item $P_{F\oplus_gF''}=P_{F'\oplus_gF''}$, in $[0,\infty]$, for every cell family $F''$ with the hinge wall $2\sigma-\kappa$, the masses of $F$, $F'$ and $F''$ kept.
\item (b) holds for the one-mark cell families $F''$ of ratio vector $\mathbf x_0$ and mass $2\sigma-\kappa$ (Definition~\ref{def:onemark}), for every $\mathbf x_0$ in a nonempty open subset of $(0,\infty)^{2N-1}$.
\item $\int\boldsymbol\varrho_F^{-\mathbf w}d\phi_F=\int\boldsymbol\varrho_{F'}^{-\mathbf w}d\phi_{F'}$ at every complex $\mathbf w$ with real part $\mathbf z$.
\end{enumerate}
So the hinge wall and the law are the whole of what a cell family contributes to the families built on it, and the one-mark cell families, whose laws are single points, already separate any two laws. For one table the plain one-mark cells of ratio $x_0$ suffice in (c), and (d) is the agreement of $\int\varrho_{\mathcal A}^{-w}d\phi_{\mathcal A}$ and $\int\varrho_{\mathcal A'}^{-w}d\phi_{\mathcal A'}$ on the line $\re w=z_{1,\Psi}$, $\max(0,\kappa-s)<z_{1,\Psi}<\min(t,\kappa)$.
\end{theorem}

What determines the contribution of a cell family to a larger construction is its hinge wall, which balance hands to the other side of the hinge, and its law, one positive measure on $(0,\infty)^{2N-1}$, on the half-line for one cell. Anything finer about the cells, their marks, their masses one by one, their letters, is invisible from outside, and anything coarser fails already against one weighted mark. A cell family is therefore replaceable, in every family it can sit in, by any other cell family of the same wall and the same law, and a cell family whose law is one point is a mark. Theorem~\ref{thm:crossing} says the same at any width for the joint law of the bunch letters; at width one the joint law collapses to the law of the ratio vector.

% Reference 20: cor:piecesmarks; source line 6295 (one plain cell).
% Proved in the last part: cor:piecesmarks.
\begin{corollary}[Cells that are marks]\label{cor:piecesmarks}
Let $F$ be a cell family with plain tables and $s_j,t_j>0$, and $\mathcal A$ a plain cell with exponents $s,t>0$.
\begin{enumerate}[label=(\roman*)]
\item The law of $F$ is carried by a single point, $(\boldsymbol\varrho_F)_\ast\phi_F=\phi_F(\Delta_B)\,\delta_{\mathbf x_0}$, exactly when the ratio vector is constant, $\boldsymbol\varrho_F\equiv\mathbf x_0$; then $\phi_F(\Delta_B)\in(0,\infty]$, finite whenever $F$ sits in a gluing of finite period, and in every gluing $F$ acts as $\phi_F(\Delta_B)$ times the one-mark cell family of ratio vector $\mathbf x_0$ and mass $\kappa_F$. For one cell with $\varrho_{\mathcal A}\equiv x_0$ this is $\phi_{\mathcal A}(\Delta_B)\,x_0^{(s-\kappa_{\mathcal A})/2}$ times the one-mark cell of ratio $x_0$ and mass $\kappa_{\mathcal A}$.
\item If $\mathcal A/g$ is the dual table of $\mathcal A\setminus g$ with the identity on $B$ and the same weights, then $\varrho_{\mathcal A}\equiv1$ and $\phi_{\mathcal A}(\Delta_B)=P_{\mathcal A\setminus g}(\mvec;\kappa_{\mathcal A},s+t-\kappa_{\mathcal A})$.
\item If a table $(\cB,\theta)$ is \emph{identically self-dual}, its bases closed under complement with $\theta_{\cE\setminus T}=\theta_T$, then for every mark $g$ and masses on the balance wall with every proper wall positive,
\[
P_F(\mvec;s,t)=\Gst21{m_g,\ s+t-m_g}{s+t}\,P_{F\setminus g}\big(\mvec|_{\cE\setminus g};s+t-m_g,m_g\big),
\]
$F=((\cB,\theta);s,t)$ and $F\setminus g=((\cB\setminus g,\theta);s+t-m_g,m_g)$ the family of the deletion.
\end{enumerate}
\end{corollary}

So (iii) splits off any mark of an identically self-dual table, by one Beta corner, at every pair of exponents, and Proposition~\ref{prop:readingsagree}, every letter active, is (i) read across a bunch of width one, the ratio vector being constant. The relation is a consequence of structure alone: the mark $g$ and the rest of the table see each other through a ratio that is identically one, so the hinge integral separates. For a table that is its own dual only up to a relabelling of its marks the hypothesis of (iii) fails, and (iii) does not apply.

\section{Series and parallel nodes}

% Reference 20: def:spnodes; source line 6362
\begin{definition}[Series and parallel cores]\label{def:spnodes}
For $q\ge2$ and the markers $G=\{g_0,g_1,\dots,g_q\}$ a \emph{series core} is a table with the bases of the list $\binom Gq$ and a \emph{parallel core} a table with the bases of the list $\binom G1$, with any positive weights. With cells $\mathcal A_1,\dots,\mathcal A_q$ at $g_1,\dots,g_q$ and the hinge $g_0$, $\binom Gq\langle\mathcal A_1,\dots,\mathcal A_q\rangle$ and $\binom G1\langle\mathcal A_1,\dots,\mathcal A_q\rangle$, the cores with unit weights, are the \emph{series connection} of the cells, whose bunch $\bigsqcup_iB_i$ is the series connection of the parts $B_i$ of Proposition~\ref{prop:seriesconnection}, and their \emph{parallel connection}; at the vertices of a hinge tree series and parallel cores are \emph{series nodes} and \emph{parallel nodes}.
\end{definition}

A series bunch of marks is the series connection of its marks, each a one-mark cell, and a parallel bunch the parallel connection (Proposition~\ref{prop:classreadings}); and the series closure $\binom{\{h_1,\dots,h_q\}}{q-1}\langle\mathcal A_1,\dots,\mathcal A_q\rangle$ of Corollary~\ref{cor:blocksclosures}(ii) is the series connection of $\mathcal A_1,\dots,\mathcal A_{q-1}$ glued to $\mathcal A_q$ along $g_0$.

% Reference 20: thm:seriesparallel; source line 6368 (plain cells of one table).
% Proved in the last part: thm:seriesparallel.
\begin{theorem}[Series and parallel nodes]\label{thm:seriesparallel}
Let $F_1,\dots,F_q$, $q\ge2$, be cell families with plain tables at the exponents $(s_j,t_j)_j$, $s_j,t_j>0$, the table $j$ of $F_i$ being $\mathcal A_{ij}$, and $G=\{g_0,\dots,g_q\}$.
\begin{enumerate}[label=(\roman*)]
\item \emph{Series.} The cell family $F$ whose table $j$ is $\mathcal A_j=\binom Gq\langle\mathcal A_{1j},\dots,\mathcal A_{qj}\rangle$ has $\kappa_F=\sum_i\kappa_{F_i}-(q-1)|\mathbf s|$, its hinge values combine in every table as $1/\gamma^K_{\mathcal A_j}=\sum_i1/\gamma^K_{\mathcal A_{ij}}$ and $\gamma^\Psi_{\mathcal A_j}=\sum_i\gamma^\Psi_{\mathcal A_{ij}}$, and
\[
f_F=\prod_{i=1}^qf_{F_i},\qquad \mathfrak r_F=\mathfrak r_{F_1}\ast\cdots\ast\mathfrak r_{F_q},
\]
the convolution on $(0,\infty)^{2N}$. For one table, $\mathcal A=\binom Gq\langle\mathcal A_1,\dots,\mathcal A_q\rangle$ has $f_{\mathcal A}=\prod_if_{\mathcal A_i}$ and $\mathfrak r^\Psi_{\mathcal A}=\mathfrak r^\Psi_{\mathcal A_1}\ast\cdots\ast\mathfrak r^\Psi_{\mathcal A_q}$, the convolution on $(0,\infty)$, and the series core twisted by $\mathbf x$ at $g_1,\dots,g_q$ has, on the simplex of its markers $g_1,\dots,g_q$ with the coordinates $v$, the ratio and the base form
\[
\Big(\sum_i\frac1{v_i}\Big)\Big(\sum_ix_iv_i\Big),\qquad \prod_iv_i^{\kappa_{\mathcal A_i}-1}\,e_{q-1}(v)^{\kappa_{\mathcal A}-s}\Big(\prod_iv_i\Big)^{-\kappa_{\mathcal A}}d\omega_{\{g_1,\dots,g_q\}}(v).
\]
\item \emph{Parallel.} The cell family $F$ whose table $j$ is $\mathcal A_j=\binom G1\langle\mathcal A_{1j},\dots,\mathcal A_{qj}\rangle$ has $\kappa_F=\sum_i\kappa_{F_i}-(q-1)|\mathbf t|$, its hinge values combine as $\gamma^K_{\mathcal A_j}=\sum_i\gamma^K_{\mathcal A_{ij}}$ and $1/\gamma^\Psi_{\mathcal A_j}=\sum_i1/\gamma^\Psi_{\mathcal A_{ij}}$, and
\[
\check f_F=\prod_{i=1}^q\check f_{F_i},\qquad \check{\mathfrak r}_F=\check{\mathfrak r}_{F_1}\ast\cdots\ast\check{\mathfrak r}_{F_q};
\]
for one table $\check f_{\mathcal A}=\prod_i\check f_{\mathcal A_i}$ and $\mathfrak r^K_{\mathcal A}=\mathfrak r^K_{\mathcal A_1}\ast\cdots\ast\mathfrak r^K_{\mathcal A_q}$.
\item \emph{Extremes of the ratio.} For one table, with plain cells $\mathcal A_1,\dots,\mathcal A_q$ at the exponents $(s,t)$: in a series connection $\sqrt{\inf\varrho_{\mathcal A}}=\sum_i\sqrt{\inf\varrho_{\mathcal A_i}}$ (Proposition~\ref{prop:seriesconnection}(i)) and $\sup\varrho_{\mathcal A}=\infty$; in a parallel connection $1/\sqrt{\sup\varrho_{\mathcal A}}=\sum_i1/\sqrt{\sup\varrho_{\mathcal A_i}}$ and $\inf\varrho_{\mathcal A}=0$, with $1/\sqrt\infty=0$.
\item \emph{Reaches.} For one table: if every $f_{\mathcal A_i}$ has a lattice expansion of reach $\zeta_i$ (Definition~\ref{def:latticeexp}), the series connection has $f_{\mathcal A}$ with a lattice expansion of reach $\big(\sum_i\sqrt{\zeta_i}\big)^2$ (Proposition~\ref{prop:seriesconnection}(ii)); if every $\check f_{\mathcal A_i}$ has one of reach $\zeta_i$, the parallel connection has $\check f_{\mathcal A}$ with one of reach $\big(\sum_i\sqrt{\zeta_i}\big)^2$.
\end{enumerate}
\end{theorem}

A series node adds the hinge values by $\Psi$ of its cells and a parallel node their hinge values by $K$; in the matching transform each is a product, and in the matching hinge measure a convolution. In the other transform neither node is a product: the dual orthant form of a series connection does not factor, since $K_{\mathcal A/g_0}=\prod_iK_{\mathcal A_i/g_i}\,e_{q-1}\big((\gamma^K_{\mathcal A_i})_i\big)$ and $1/\gamma^K_{\mathcal A}=\sum_i1/\gamma^K_{\mathcal A_i}$ mix the cells. What changes at the edges where the kind of node changes is the transform, not the cells (Proposition~\ref{prop:faces}(v)). The ceilings of Section~\ref{sec:reachlaw} are read on a bunch as a whole, where the letters $K_{\mathcal A/g_0}$ and $\Psi_{\mathcal A/g_0}$ of the hinge minors of a series connection mix the cells and their zeros can raise a ceiling above the reach; at a node the cells are separate factors and the reach is read cell by cell (Lemma~\ref{lem:latticeproduct}). So ceilings are to be applied to the prime nodes of a hinge tree, and reaches combined across series and parallel nodes by (iv).

\section{Hinge trees}

% Reference 20: def:hingetree; source line 6402 (one table).
% Proved in the last part: def:hingetree.
\begin{definition}[Hinge tree]\label{def:hingetree}
A \emph{hinge tree} of a table on $\cE$ is a finite tree $\mathfrak T$ with a core $\mathcal H_v$ on $\cE_v\sqcup G_v$ at each vertex $v$, its \emph{node}, where the $\cE_v$ partition $\cE$ and $G_v=\{g_{v,\mathfrak e}\}$ has one marker for each edge $\mathfrak e$ at $v$, such that the table is obtained from the $\mathcal H_v$ by the two-sums identifying $g_{v,\mathfrak e}$ with $g_{w,\mathfrak e}$ for every edge $\mathfrak e=\{v,w\}$. Rooted at a vertex $o$, every vertex $v\neq o$ has a parent edge $\mathfrak e_v$, and the \emph{subtree cell} $\mathcal A_v$ is the substitution of the $\mathcal H_w$ over the vertices $w$ below $v$, $v$ included, with the hinge $g_{v,\mathfrak e_v}$; its cells are the subtree cells of the vertices directly below $v$. The tables of a family $F=((\cB_j,\theta_j);s_j,t_j)_j$ have a \emph{common hinge tree} when one tree $\mathfrak T$, with the same sets $\cE_v$ and markers $G_v$, is a hinge tree of every table, table $j$ with the nodes $\mathcal H_{v,j}$; the nodes then form the \emph{node families} $F_{\mathcal H_v}=((\mathcal H_{v,j},\theta^{\mathcal H_{v,j}});s_j,t_j)_j$ and, rooted at $o$, the subtree cells $\mathcal A_{v,j}$ the cell families $F_v=((\mathcal A_{v,j},\theta^{\mathcal A_{v,j}});s_j,t_j)_j$ of the vertices $v\neq o$ (Definition~\ref{def:hinged}).
\end{definition}

% Reference 20: thm:primetree; source line 5838
\begin{theorem}[A connected table is a tree of prime tables]\label{thm:primetree}
Call a connected table on at least four marks \emph{prime} when it has no bunch of width one with at least two marks on each side. Every connected table on at least three marks whose weights factor across each bunch of width one at which it is decomposed, a list or a determinantal table for instance (Theorem~\ref{thm:levels}), is obtained by substitutions along a tree from prime tables and the series and parallel nodes $\binom S{|S|-1}$ and $\binom S1$, $|S|\ge3$ (Definition~\ref{def:spnodes}), the weights carried along. With no two series nodes and no two parallel nodes adjacent in the tree, the tree, the \emph{canonical hinge tree} (Definition~\ref{def:hingetree}), is unique up to the names of the markers. The prime nodes of a list are therefore invariants of the list.
\end{theorem}

By Lemma~\ref{lem:nodesum}, applied leaf by leaf, the substitution along a tree does not depend on the order of the two-sums: its bases are the sets $\bigsqcup_v(T_v\cap\cE_v)$ with $T_v$ a basis of $\mathcal H_v$ and every edge's marker in exactly one of the two bases at its ends. Every connected table has a hinge tree with one vertex, and decomposing at bunches of width one with at least two marks on each side, the weights factoring across each such bunch, markers counted as marks, refines it until no node has such a bunch. A node on at least four marks without one is prime (Theorem~\ref{thm:primetree}); the nodes on three marks, with the bases of $\binom G2$ or of $\binom G1$, $|G|=3$, are series and parallel nodes, and merging adjacent series nodes, and adjacent parallel nodes, gives the series and parallel nodes of any size; with no two series and no two parallel nodes adjacent this is the canonical hinge tree of Theorem~\ref{thm:primetree}.

% Reference 20: thm:contraction; source line 6408 (one table with $s,t>0$).
% Proved in the last part: thm:treeperiod.
\begin{theorem}[The period along a hinge tree]\label{thm:treeperiod}
Let $F=((\cB_j,\theta_j);s_j,t_j)_{j=1}^N$ have real data on its balance wall and its tables a common hinge tree rooted at $o$, with the node families $F_{\mathcal H_v}$ and the cell families $F_v$ of the vertices $v\neq o$ (Definition~\ref{def:hingetree}), and let $F_{\mathcal H_v}[\mathbf x]$ be $F_{\mathcal H_v}$ twisted at the markers of the edges to the vertices $w$ directly below $v$ by $\mathbf x_w\in(0,\infty)^{2N-1}$, with the masses $\kappa_{F_w}$ there and $\mvec|_{\cE_v}$ on $\cE_v$ (Theorem~\ref{thm:node}).
\begin{enumerate}[label=(\roman*)]
\item For $v\neq o$, $\kappa_{F_v}$ is the hinge wall of $F_{\mathcal H_v}[\mathbf x]$, and
\[
\int H\,d\big((\boldsymbol\varrho_{F_v})_\ast\phi_{F_v}\big)=\int\Big(\int H\,d\big((\boldsymbol\varrho_{F_{\mathcal H_v}[\mathbf x]})_\ast\phi_{F_{\mathcal H_v}[\mathbf x]}\big)\Big)\prod_wd\big((\boldsymbol\varrho_{F_w})_\ast\phi_{F_w}\big)(\mathbf x_w)
\]
for every Borel $H\ge0$ on $(0,\infty)^{2N-1}$.
\item The period is
\[
P_F(\mvec;\mathbf s,\mathbf t)=\int P_{F_{\mathcal H_o}[\mathbf x]}\big((\mvec|_{\cE_o},(\kappa_{F_w})_w);\mathbf s,\mathbf t\big)\prod_wd\big((\boldsymbol\varrho_{F_w})_\ast\phi_{F_w}\big)(\mathbf x_w),
\]
the product over the vertices $w$ directly below the root.
\item The law of $F_v$ is determined by the tables $\mathcal H_{w,j}$ and the masses on $\cE_w$ for the vertices $w$ below $v$, and by the exponents; nothing above $v$ enters it.
\end{enumerate}
For one table the cell families are the subtree cells $\mathcal A_v$, with laws on $(0,\infty)$.
\end{theorem}

Where the active letters read every subtree cell by one hinge value, as at a full lean of one table, every cell of the tree is pinned (Theorem~\ref{thm:pinned}), and Chapter~\ref{ch:leantree} writes the period along the tree as a product and carries it to every load.

% Reference 20: prop:work; source line 6422 (one table).
% Proved in the last part: prop:work.
\begin{proposition}[What the work depends on]\label{prop:work}
In Theorem~\ref{thm:treeperiod}:
\begin{enumerate}[label=(\roman*)]
\item Two vertices whose subtree cells are carried onto each other, in every table, by one bijection keeping the tables, the weights, the masses and the hinge have equal hinge walls and equal laws. So the laws to be computed are one for each orbit of subtree cells under such bijections.
\item Let $s_j,t_j>0$. At a series node $v$ whose cell families directly below fall into orbits $c_1,\dots,c_p$ with multiplicities $k_1,\dots,k_p$, $f_{F_v}=\prod_lf_{c_l}^{\,k_l}$ and $\kappa_{F_v}=\sum_lk_l\kappa_{c_l}-\big(\sum_lk_l-1\big)|\mathbf s|$, $f_{c_l}$ and $\kappa_{c_l}$ the transform and the hinge wall of a member of $c_l$; at a parallel node $\check f_{F_v}=\prod_l\check f_{c_l}^{\,k_l}$ and $\kappa_{F_v}=\sum_lk_l\kappa_{c_l}-\big(\sum_lk_l-1\big)|\mathbf t|$. The multiplicities enter as exponents.
\item A change at a vertex $v_0$ of the weights, of the masses keeping their total, or of the tables $\mathcal H_{v_0,j}$ keeping their ranks and their numbers of marks, changes no hinge wall and changes the laws of the cell families $F_v$ only for the vertices $v$ on the path from $v_0$ to the root. A change of masses at two vertices $v_0,v_1$ that keeps the total mass changes the hinge wall $\kappa_{F_v}$ exactly for the vertices $v$ whose subtree contains one of $v_0,v_1$ and not the other, the edges on the path between $v_0$ and $v_1$, by the change of the total mass below $\mathfrak e_v$.
\end{enumerate}
\end{proposition}

A hinge tree is read from the leaves up by a computation whose reusable data are the laws on the edges. Its cost is governed by the prime nodes, where cells interact through the node's own integral, and by the number of distinct subtrees, since equal subtrees have equal laws; series and parallel nodes cost no integration in the matching transform, and repeated subtrees enter as powers, as the hinge tree procedure makes explicit.

% Reference 20: def:contractproc; source line 6436 (one table).
% Proved in the last part: def:treeproc.
\begin{definition}[The hinge tree procedure]\label{def:treeproc}
\emph{Given}: a common hinge tree of the tables of a family (Definition~\ref{def:hingetree}), the masses on the marks of its nodes, exponents $s_j,t_j>0$, and a root. \emph{Reusable data}: for every orbit of subtree cells (Proposition~\ref{prop:work}(i)) its hinge wall $\kappa$ and its law, kept in one of three representations: (a) the transform $f$ and the dual transform $\check f$ on a grid of $(\mathbf x,\mathbf y)$ with $x_1=1$, of $y$ for one table; (b) the Mellin transform $\int\boldsymbol\varrho^{-\mathbf z}d\phi$ of the law on a vertical subspace of its domain of convergence, for one table $\int\varrho^{-z}d\phi$ on a vertical line of its strip; (c) for one table, the lattice expansion of $f$ at $y\to0$, the residues of $y^{-z}\Gamma(z)\Gamma(z+s-\kappa)\int\varrho^{-z}d\phi$ to the left of the line (Theorem~\ref{thm:nodeclosing}(i)). \emph{Operations}, from the leaves up:
\begin{enumerate}[label=(\arabic*)]
\item a one-mark leaf, the one-mark cell family of ratio vector $\mathbf x_0$ (Definition~\ref{def:onemark}), has $\kappa=m_e$, the law $\delta_{\mathbf x_0}$ and the transforms
\begin{gather*}
f(\mathbf x,\mathbf y)=\prod_jx_{0,j,K}^{-s_j}\,\sG_{|\mathbf s|-m_e;(1,-1)}\Big(\sum_j\frac{x_j}{x_{0,j,K}},\ \sum_jy_jx_{0,j,\Psi}\Big),\\
\check f(\mathbf x,\mathbf y)=\prod_jx_{0,j,\Psi}^{-t_j}\,\sG_{|\mathbf t|-m_e;(1,-1)}\Big(\sum_j\frac{x_j}{x_{0,j,\Psi}},\ \sum_jy_jx_{0,j,K}\Big);
\end{gather*}
for one table the one-mark cell of ratio $x_0$ has the law $x_0^{(m_e-s)/2}\,\delta_{x_0}$ and
\[
f(y)=x_0^{(m_e-s)/2}\sG_{s-m_e;(1,-1)}(1,x_0y),\qquad \check f(y)=x_0^{(m_e-s)/2-m_e}\sG_{t-m_e;(1,-1)}(1,y/x_0);
\]
\item at a series node multiply the transforms of the vertices directly below, in (a) pointwise, in (b) by Mellin convolution of the Mellin transforms of Proposition~\ref{prop:faces}(iii), for one table $\Gamma(z)\Gamma(z+s-\kappa)\int\varrho^{-z}d\phi$, in (c) by the product of lattice expansions (Lemma~\ref{lem:latticeproduct}); at a parallel node the same with the dual transforms;
\item along an edge where the kind of node changes, change the transform by Proposition~\ref{prop:faces}(v): in (b) a reflection and a ratio of Gammas, in (a) a Mellin transform followed by an inverse Mellin transform, each one integral on $(0,\infty)^{2N-1}$, on the half-line for one table;
\item at a prime node evaluate the Mellin transform of the law of the twisted node family, or its transform, as a period of the twisted node family, and integrate it against the laws of the vertices directly below (Theorem~\ref{thm:node}(ii));
\item at the root evaluate Theorem~\ref{thm:node}(i); for a series root, with the cell families $F_i$ directly below and their transforms read at $x_1=1$, $\mathbf x'=(x_2,\dots,x_N)$, the period is
\begin{gather*}
\prod_j\frac1{\Gamma(s_j)\Gamma(t_j)}\int\prod_{j\ge2}x_j^{s_j-1}\prod_jy_j^{t_j-1}\prod_if_{F_i}\,d\mathbf x'\,d\mathbf y,\\
\text{for one table}\quad\frac1{\Gamma(s)\Gamma(t)}\int_0^\infty y^{t-1}\prod_if_{\mathcal A_i}(y)\,dy,
\end{gather*}
and for a parallel root the same with $s_j$ and $t_j$ swapped and the dual transforms (Theorem~\ref{thm:glue}(ii)).
\end{enumerate}
\emph{Output}: the period, and with it every law on every edge, ready to be reused when the tree is changed (Chapter~\ref{ch:repetition}).
\end{definition}

% Reference 20: cor:sptables; source line 6448 (one table).
% Proved in the last part: cor:sptables.
\begin{corollary}[Series-parallel families]\label{cor:sptables}
Let the tables of a family with exponents $s_j,t_j>0$ have a common hinge tree whose nodes are series and parallel nodes and whose leaves are one-mark cell families, and at every edge joining a series node to a parallel node let the law of the subtree below have a Mellin transform finite at a real $\mathbf z$ with $0<z_{j,K}<s_j$ and $0<z_{j,\Psi}<t_j$ for every $j$ (Proposition~\ref{prop:faces}(v)). Then its period is obtained from the transforms of the leaves, Grammas of the marks (Definition~\ref{def:treeproc}(1)), by pointwise products, by two integral transforms on $(0,\infty)^{2N-1}$ for each edge joining a series node to a parallel node, and by one final integral on $(0,\infty)^{2N-1}$; the number of integrations depends on the alternations of the tree and not on the number of marks, and subtrees repeated at a node enter as powers. For one table the integral transforms are on the half-line.
\end{corollary}


\chapter{Repetition and the change of query}\label{ch:repetition}

A construction that repeats a cell should not repeat the work, and a query that changes one weight or one mass should not restart the computation. A cell family repeated $k$ times at a series node enters its period as the $k$-th power of its transform, and repetition without bound forces its hinge wall to be $\sum_js_j$; for one cell the growth in $k$ is then factorial, with the base $A_1/\min(s,t)$, $A_1$ the period of the cell with its hinge deleted, and a prefactor read from one further constant of the cell and one leading coefficient of the room (Theorems~\ref{thm:repeatwall} and~\ref{thm:seriesrepeat}); a chain of repeated cell families is a power of one positive operator (Proposition~\ref{prop:chain}); the weight of a hinge enters through a Mellin convolution whose two expansions are the two closings of the node of Section~\ref{sec:onescale} (Theorem~\ref{thm:hingeweight}); and a change of data at one node moves only the laws on one path of the tree (Proposition~\ref{prop:local}).

\section{Repetition at a node}

For a cell family $F$ with plain tables $\mathcal A_j$ and $k\ge1$, $F^{\mathrm s}_k$ and $F^{\mathrm p}_k$ are the cell families whose table $j$ is the series and the parallel connection of $k$ copies of $\mathcal A_j$ (Definition~\ref{def:spnodes}), with $F^{\mathrm s}_1=F^{\mathrm p}_1=F$; for one plain cell $\mathcal A$ they are written $\mathcal A^{\mathrm s}_k$ and $\mathcal A^{\mathrm p}_k$. The cell family is \emph{series-balanced} when $\kappa_F=|\mathbf s|$ and \emph{parallel-balanced} when $\kappa_F=|\mathbf t|$; by Definition~\ref{def:hinged} the first says that the deletions $\mathcal A_j\setminus g$, at the exponents $(s_j,t_j)$, form a family on its balance wall at the masses $\mvec$, $m_B=\sum_j\big(s_jr_{\mathcal A_j}+t_j(|B|-r_{\mathcal A_j})\big)$, and the second that the contractions $\mathcal A_j/g$ do; for one cell they say $\kappa_{\mathcal A}=s$ and $\kappa_{\mathcal A}=t$.

% Reference 20: thm:repeatwall; source line 6462 (one plain cell).
% Proved in the last part: thm:repeatwall.
\begin{theorem}[How often a cell family repeats]\label{thm:repeatwall}
Let $F$ be a cell family with plain tables and real exponents.
\begin{enumerate}[label=(\roman*)]
\item $\kappa_{F^{\mathrm s}_k}=|\mathbf s|+k(\kappa_F-|\mathbf s|)$ and $\kappa_{F^{\mathrm p}_k}=|\mathbf t|+k(\kappa_F-|\mathbf t|)$.
\item If $F'$ is a cell family with $F^{\mathrm s}_k\oplus_gF'$ on its balance wall and of convergent period, then $0<|\mathbf s|+k(\kappa_F-|\mathbf s|)<2\sigma$, $\sigma$ the total strength (Definition~\ref{def:family}). So series connections of any number of copies of $F$ occur in convergent periods only if $F$ is series-balanced; otherwise $k<|\mathbf s|/(|\mathbf s|-\kappa_F)$ when $\kappa_F<|\mathbf s|$ and $k<|\mathbf t|/(\kappa_F-|\mathbf s|)$ when $\kappa_F>|\mathbf s|$. For parallel connections the same holds with $|\mathbf t|$ in place of $|\mathbf s|$ and $|\mathbf s|$ in place of $|\mathbf t|$.
\end{enumerate}
\end{theorem}

The hinge wall of a series connection drifts by $\kappa_F-|\mathbf s|$ for every copy, and the room must absorb the drift; only a series-balanced cell family repeats without limit. It is the cell family whose deletions at the hinge form a balanced family, and for one cell the period of that deletion is the number that governs the repetition.

% Reference 20: thm:seriesrepeat; source line 6475 (one plain cell).
% Proved in the last part: thm:seriesrepeat.
\begin{theorem}[Series repetition]\label{thm:seriesrepeat}
Let $F$ be a series-balanced cell family on $B\sqcup\{g\}$ and $F'=((\mathcal A'_j,\theta'_j);s_j,t_j)_j$ a cell family on $R\sqcup\{g\}$ with $\kappa_{F'}=|\mathbf t|$, plain tables, exponents $s_j,t_j>0$, so that $F^{\mathrm s}_k\oplus_gF'$ is on its balance wall for every $k$; in this statement write $P(k)$ for its period, at the masses of the copies and of $F'$, and read the transforms at $x_1=1$, $\mathbf x'=(x_2,\dots,x_N)$.
\begin{enumerate}[label=(\roman*)]
\item In $[0,\infty]$,
\[
P(k)=\prod_j\frac1{\Gamma(s_j)\Gamma(t_j)}\int\prod_{j\ge2}x_j^{s_j-1}\prod_jy_j^{t_j-1}\,f_F^{\,k}f_{F'}\,d\mathbf x'\,d\mathbf y,
\]
the $k$-th moment of the image under $f_F$ of the measure
\[
\prod_j\frac1{\Gamma(s_j)\Gamma(t_j)}\prod_{j\ge2}x_j^{s_j-1}\prod_jy_j^{t_j-1}\,f_{F'}(\mathbf x,\mathbf y)\,d\mathbf x'\,d\mathbf y .
\]
Where finite, $P(k)^2\le P(k-1)P(k+1)$, and for $z\ge0$, with $P(0)$ the period at $\mvec|_R$ of the family of the contractions $((\mathcal A'_j/g,\theta'_j);s_j,t_j)_j$, the hinge of $F'$ contracted,
\[
\sum_{k\ge0}P(k)\frac{z^k}{k!}=\prod_j\frac1{\Gamma(s_j)\Gamma(t_j)}\int\prod_{j\ge2}x_j^{s_j-1}\prod_jy_j^{t_j-1}\,e^{zf_F}f_{F'}\,d\mathbf x'\,d\mathbf y .
\]
For one cell, $P(k)=\frac1{\Gamma(s)\Gamma(t)}\int_0^\infty y^{t-1}f_{\mathcal A}(y)^kf_{\mathcal A'}(y)\,dy$ and $P(0)=\int\varrho_{\mathcal A'}^{-t}d\phi_{\mathcal A'}=P_{\mathcal A'/g}(\mvec|_R;s,t)$.
\item For one cell $\mathcal A$ and one cell $\mathcal A'$, at the exponents $(s,t)$, let $\mathcal A\setminus g$ have every proper wall positive at the masses of $\mathcal A$ and the exponents $(s,t)$, and put
\[
A_1=\int_{\Delta_B}d\phi_{\mathcal A}=P_{\mathcal A\setminus g}(\mvec|_B;s,t),\qquad A_0=2A_1(\log\Gamma)'(1)-\int_{\Delta_B}\log\varrho_{\mathcal A}\,d\phi_{\mathcal A}.
\]
Then $\int\varrho_{\mathcal A}^{-z}d\phi_{\mathcal A}$ is finite on a neighbourhood of $z=0$, and for every $\varepsilon\in(0,1)$ at which $\int\varrho_{\mathcal A}^{\varepsilon}d\phi_{\mathcal A}<\infty$ there is $C$ with $\big|f_{\mathcal A}(y)-A_1(-\log y)-A_0\big|\le Cy^\varepsilon$ for all $y>0$.
\item Let moreover $\int\varrho_{\mathcal A'}^{-z}d\phi_{\mathcal A'}$ be finite on a neighbourhood of $z=\max(0,t-s)$, and let the \emph{polynomial of the room} be
\[
\mathcal R(v)=\begin{cases}\Gamma(s-t)\int d\phi_{\mathcal A'}&(s>t),\\ \Gamma(t-s)\int\varrho_{\mathcal A'}^{s-t}d\phi_{\mathcal A'}&(s<t),\\ v\int d\phi_{\mathcal A'}+2(\log\Gamma)'(1)\int d\phi_{\mathcal A'}-\int\log\varrho_{\mathcal A'}\,d\phi_{\mathcal A'}&(s=t),\end{cases}
\]
with the coefficients $\mathcal R(v-A_0/A_1)=\sum_i\mathcal R_iv^i$; let $P(1)<\infty$. Then there is $\varepsilon>0$ such that for every $\vartheta\in\big(\min(s,t)/(\min(s,t)+\varepsilon),1\big)$
\[
P(k)=\frac{A_1^k\,e^{\min(s,t)A_0/A_1}}{\Gamma(s)\Gamma(t)}\sum_i\mathcal R_i\,\frac{\Gamma(k+i+1)}{\min(s,t)^{k+i+1}}\,\big(1+O(\vartheta^k)\big).
\]
In particular $P(k)\sim\frac{c\,e^{\min(s,t)A_0/A_1}}{\Gamma(s)\Gamma(t)}\Gamma(k+1+j)A_1^k\min(s,t)^{-k-1-j}$, with $j=1$ and $c=\int d\phi_{\mathcal A'}$ when $s=t$ and $j=0$ and $c=\mathcal R$ otherwise: the growth is factorial, with the base $A_1/\min(s,t)$.
\item Under the hypotheses of (iii) the series of (i) converges for $0\le z<\min(s,t)/A_1$ and diverges at $z=\min(s,t)/A_1$.
\end{enumerate}
\end{theorem}

So a repeated cell family enters its construction as a power, and the whole sequence $P(k)$ is the sequence of moments of one positive measure: log-convex, with an exponential generating function that is one integral. For one cell the factorial growth is not a property of the room: the transform of a series-balanced cell grows like $A_1\log(1/y)$ at small scales, its $k$-th power concentrates where $\log(1/y)$ is of the order $k/\min(s,t)$, and there only two numbers of the cell survive, the period $A_1$ of the cell with its hinge deleted and the constant $A_0$, and only the first lattice term of the room, through $\min(s,t)$ and $\mathcal R$. The error is exponentially small, not a power of $1/k$: the next terms of both lattice expansions sit at distance $\varepsilon$ to the right and contribute in the ratio $\big(\min(s,t)/(\min(s,t)+\varepsilon)\big)^k$. Theorem~\ref{thm:seriesclosure} reads the series closure of any cells, for any family, through their transforms; for $k$ equal cells at the centre it is (i) with $k-1$ copies at the series node and one copy as the room.

% Reference 20: cor:parallelrepeat; source line 6524 (one plain cell).
% Proved in the last part: cor:parallelrepeat.
\begin{corollary}[Parallel repetition]\label{cor:parallelrepeat}
Let $F$ be a parallel-balanced cell family on $B\sqcup\{g\}$ and $F'=((\mathcal A'_j,\theta'_j);s_j,t_j)_j$ a cell family on $R\sqcup\{g\}$ with $\kappa_{F'}=|\mathbf s|$, plain tables, exponents $s_j,t_j>0$, and in this statement write $P(k)$ for the period of $F^{\mathrm p}_k\oplus_gF'$, the transforms read at $x_1=1$, $\mathbf x'=(x_2,\dots,x_N)$. Then
\[
P(k)=\prod_j\frac1{\Gamma(s_j)\Gamma(t_j)}\int\prod_{j\ge2}x_j^{t_j-1}\prod_jy_j^{s_j-1}\,\check f_F^{\,k}\check f_{F'}\,d\mathbf x'\,d\mathbf y,
\]
with $P(0)$ the period at $\mvec|_R$ of the family of the deletions $((\mathcal A'_j\setminus g,\theta'_j);s_j,t_j)_j$; for one cell, $P(k)=\frac1{\Gamma(s)\Gamma(t)}\int_0^\infty y^{s-1}\check f_{\mathcal A}(y)^k\check f_{\mathcal A'}(y)\,dy$. For one cell $\mathcal A$ and one cell $\mathcal A'$ let $\mathcal A/g$ have every proper wall positive at the masses of $\mathcal A$ and the exponents $(s,t)$, and put
\[
A_1=\int\varrho_{\mathcal A}^{-t}d\phi_{\mathcal A}=P_{\mathcal A/g}(\mvec|_B;s,t),\qquad A_0=2A_1(\log\Gamma)'(1)+\int\varrho_{\mathcal A}^{-t}\log\varrho_{\mathcal A}\,d\phi_{\mathcal A}.
\]
Then Theorem~\ref{thm:seriesrepeat}(ii) to (iv) hold with $s$ and $t$ swapped and with these $A_1$ and $A_0$, the dual transform $\check f_{\mathcal A}$ in place of the transform $f_{\mathcal A}$, the finiteness of $\int\varrho_{\mathcal A}^{-z}d\phi_{\mathcal A}$ near $z=t$ in place of near $z=0$, the condition $\int\varrho_{\mathcal A}^{-t-\varepsilon}d\phi_{\mathcal A}<\infty$ in place of $\int\varrho_{\mathcal A}^{\varepsilon}d\phi_{\mathcal A}<\infty$, the finiteness of $\int\varrho_{\mathcal A'}^{-z}d\phi_{\mathcal A'}$ near $z=\min(s,t)$ as the hypothesis on $\mathcal A'$, and the polynomial of the room
\[
\mathcal R(v)=\begin{cases}\Gamma(t-s)\int\varrho_{\mathcal A'}^{-s}d\phi_{\mathcal A'}&(s<t),\\ \Gamma(s-t)\int\varrho_{\mathcal A'}^{-t}d\phi_{\mathcal A'}&(s>t),\\ v\int\varrho_{\mathcal A'}^{-s}d\phi_{\mathcal A'}+2(\log\Gamma)'(1)\int\varrho_{\mathcal A'}^{-s}d\phi_{\mathcal A'}+\int\varrho_{\mathcal A'}^{-s}\log\varrho_{\mathcal A'}\,d\phi_{\mathcal A'}&(s=t),\end{cases}
\]
whose leading coefficient is the constant $c$ of the equivalent in (iii).
\end{corollary}

\section{Chains}

% Reference 20: prop:chain; source line 6537 (one table).
% Proved in the last part: prop:chain.
\begin{proposition}[A chain is a power of one operator]\label{prop:chain}
Let a \emph{chain cell family} $U$ have in every table a hinged core $\mathcal U_j$ (Definition~\ref{def:node}) with one further marker $g_{\mathrm{in}}$, its hinge $g_{\mathrm{out}}$, own marks $\cE_{\mathcal U}$ with masses, and real exponents $(s_j,t_j)$; further fixed cell families may be substituted into $U$ beforehand. For a cell family $F_0$ let $F_{l+1}=U\langle F_l\rangle$, table by table, the cell at $g_{\mathrm{in}}$.
\begin{enumerate}[label=(\roman*)]
\item $\kappa_{F_{l+1}}=\kappa_{F_l}+\kappa_U$, with $\kappa_U=m_{\cE_{\mathcal U}}-\sum_j\big(s_j(r_{\mathcal U_j}-1)+t_j(|\cE_{\mathcal U}|+1-r_{\mathcal U_j})\big)$ the hinge wall of $U$ with the mass $0$ on $g_{\mathrm{in}}$. For a fixed cell family $F'$, the families $F_l\oplus_gF'$ are all on their balance walls exactly when $F_0\oplus_gF'$ is and $\kappa_U=0$; then $\kappa_{F_l}=\kappa_{F_0}$.
\item Let $\kappa_U=0$, and in this statement write $\lambda_l$ for the law of $F_l$ and $\lambda^U_{\mathbf x}$ for the law of $U$ twisted at $g_{\mathrm{in}}$ by $\mathbf x\in(0,\infty)^{2N-1}$ (Definition~\ref{def:twist}) with the mass $\kappa_{F_0}$ there. On positive measures $\mu$ and nonnegative functions $H$ on $(0,\infty)^{2N-1}$ put
\begin{gather*}
\mathcal Q\mu=\int\lambda^U_{\mathbf x}\,d\mu(\mathbf x),\qquad(\mathcal Q^\ast H)(\mathbf x)=\int H\,d\lambda^U_{\mathbf x},\\
h_{F'}(\mathbf x)=\int\Big(\int_0^\infty v^{\kappa_{F_0}-1}\prod_j(x_{j,K}v+x'_{j,K})^{-s_j}(x_{j,\Psi}v+x'_{j,\Psi})^{-t_j}\,dv\Big)\,d\big((\boldsymbol\varrho_{F'})_\ast\phi_{F'}\big)(\mathbf x'),
\end{gather*}
the inner integral the period of the twisted core family of Theorem~\ref{thm:node}(i) with the core $\binom{\{g_1,g_2\}}1$ in every table. Then $\lambda_l=\mathcal Q^l\lambda_0$ and $P_{F_l\oplus_gF'}=\int(\mathcal Q^{\ast l}h_{F'})\,d\lambda_0$; for $z\ge0$, $\sum_lz^lP_{F_l\oplus_gF'}=\int\sum_lz^l(\mathcal Q^{\ast l}h_{F'})\,d\lambda_0$, in $[0,\infty]$.
\item If $H>0$ and $M>0$ have $\mathcal Q^\ast H\le MH$ and $h_{F'}\le CH$, then $P_{F_l\oplus_gF'}\le CM^l\int H\,d\lambda_0$; if $H'\ge0$ and $M'>0$ have $\mathcal Q^\ast H'\ge M'H'$ and $h_{F'}\ge c\,H'$, then $P_{F_l\oplus_gF'}\ge c\,M'^l\int H'\,d\lambda_0$.
\item Let the tables be plain and $s_j,t_j>0$. If $U$ is the series core on $\{g_{\mathrm{out}},g_{\mathrm{in}},g'\}$ with a fixed cell family $F''$ at $g'$, then $f_{F_{l+1}}=f_{F''}f_{F_l}$ and $\kappa_U=\kappa_{F''}-|\mathbf s|$; if it is the parallel core, then $\check f_{F_{l+1}}=\check f_{F''}\check f_{F_l}$ and $\kappa_U=\kappa_{F''}-|\mathbf t|$.
\end{enumerate}
For one table the laws are measures on $(0,\infty)$, the twist is one number, and the inner integral of $h_{F'}$ is the period of Lemma~\ref{lem:twomarks}.
\end{proposition}

A chain is one kernel applied $l$ times: the law of the chain cell family twisted at $g_{\mathrm{in}}$, computed once, carries the law of the chain from one length to the next, and the period at length $l$ pairs the $l$-th iterate with the room. The work for all lengths up to $l$ is $l$ applications of one operator, and for one length $\lceil\log_2l\rceil$ compositions of kernels. A series chain cell family is diagonal in the transform and a parallel one in the dual transform, so for them the operator is a multiplication and the length is an exponent; a chain cell family that alternates the two kinds mixes the two hinge values, and its operator is not diagonal in either. Theorem~\ref{thm:chainoperator} is the same chain read on functions of the twists, a family glued to copies of itself at two marks, with the pinned case $P_F(\mvec;\mathbf s,\mathbf t)^k/\Gst21{m_{g_{\mathrm{in}}},\ m_{g_{\mathrm{out}}}}{m_{g_{\mathrm{in}}}+m_{g_{\mathrm{out}}}}^{k-1}$; Proposition~\ref{prop:chain} reads it on laws, and positive test functions bound it from both sides. Every polynomial Gamma moment of a chain is the power of one finite matrix (Theorem~\ref{thm:transferpower}).

\section{The weight of a hinge}

% Reference 20: thm:hingeweight; source line 6558 (one plain cell in (i)).
% Proved in the last part: thm:hingeweight.
\begin{theorem}[The weight of a hinge]\label{thm:hingeweight}
For plain cells $\mathcal A$ on $B\sqcup\{g\}$ and $\mathcal A'$ on $R\sqcup\{g\}$ with exponents $s,t>0$, and $w>0$, let $w\mathcal A$ be $\mathcal A$ with the weights of its bases through $g$ multiplied by $w$, so that in $(w\mathcal A)\oplus_g\mathcal A'$ the bases made of a basis of $\mathcal A/g$ and a basis of $\mathcal A'\setminus g$, the second kind of Lemma~\ref{lem:gluemark}, carry the factor $w$; let it be on its balance wall, put $q=w^2$, and in this statement write $\Pi(q)=q^{(s-\kappa_{\mathcal A})/2}P_{(w\mathcal A)\oplus_g\mathcal A'}(\mvec;s,t)$.
\begin{enumerate}[label=(\roman*)]
\item For a cell family $F$ with plain tables $\mathcal A_j$ and $\mathbf w\in(0,\infty)^N$, let $\mathbf wF$ be $F$ with the weights of the bases through $g$ multiplied by $w_j$ in table $j$. Its hinge wall is $\kappa_F$, its ratio vector has the coordinates $(w_1/w_j)\,\gamma^K_{\mathcal A_j}/\gamma^K_{\mathcal A_1}$ and $w_1w_j\,\gamma^\Psi_{\mathcal A_j}/\gamma^K_{\mathcal A_1}$, the hinge values of $F$ being read, $\phi_{\mathbf wF}=w_1^{\kappa_F}\prod_jw_j^{-s_j}\,\phi_F$, and, for $s_j,t_j>0$, $f_{\mathbf wF}(\mathbf x,\mathbf y)=f_F(\mathbf w\circ\mathbf x,\mathbf w\circ\mathbf y)$ and $\check f_{\mathbf wF}(\mathbf x,\mathbf y)=\prod_jw_j^{-s_j-t_j}\,\check f_F(\mathbf x/\mathbf w,\mathbf y/\mathbf w)$, the products and quotients taken coordinatewise. For one cell, $\varrho_{w\mathcal A}=w^2\varrho_{\mathcal A}$, $\phi_{w\mathcal A}=w^{\kappa_{\mathcal A}-s}\phi_{\mathcal A}$, $f_{w\mathcal A}(y)=w^{\kappa_{\mathcal A}-s}f_{\mathcal A}(w^2y)$ and $\check f_{w\mathcal A}(y)=w^{-\kappa_{\mathcal A}-s}\check f_{\mathcal A}(y/w^2)$.
\item $\Pi(q)=\frac1{\Gamma(s)\Gamma(t)}\int_0^\infty y^{t-1}f_{\mathcal A}(qy)f_{\mathcal A'}(y)\,dy$, a decreasing function of $q$, and for real $z$ with $\max(0,\kappa_{\mathcal A}-s)<z<\min(t,\kappa_{\mathcal A})$, in $(0,\infty]$,
\[
\int_0^\infty q^{z-1}\Pi(q)\,dq=\frac{\Gamma(z)\Gamma(z+s-\kappa_{\mathcal A})\int\varrho_{\mathcal A}^{-z}d\phi_{\mathcal A}\cdot\Gamma(t-z)\Gamma(\kappa_{\mathcal A}-z)\int\varrho_{\mathcal A'}^{z-t}d\phi_{\mathcal A'}}{\Gamma(s)\Gamma(t)} .
\]
\item If this is finite at a real $c$ in that interval, then $\Pi$ is finite and continuous on $(0,\infty)$ and is the value there of the function $\frac1{2\pi i}\int_{(c)}q^{-z}(\cdots)\,dz$, with $(\cdots)$ the right side of the display in (ii), holomorphic on the sector $|\arg q|<2\pi$ of the logarithmic cover of $\mathbb C\setminus\{0\}$.
\item In the setting of Section~\ref{sec:onescale}, with $\mathcal A$ the cell of the bunch and $\mathcal A'$ that of the room, let $f_{\mathcal A}$ have a lattice expansion $\sum_{\eta,j}\beta_{\eta,j}y^\eta(-\log y)^j$ of every reach above $\operatorname{reach}(B)$ (Definition~\ref{def:latticeexp}), and let $\inf\varrho_{\mathcal A'}>0$. Then for $0<q<\inf\varrho_{\mathcal A'}/\operatorname{reach}(B)$, absolutely,
\begin{multline*}
\Pi(q)=\frac1{\Gamma(s)\Gamma(t)}\sum_{\eta,j}\beta_{\eta,j}\,q^\eta\sum_{i=0}^j\binom ji(-\log q)^{j-i}\\
\times(-\partial_z)^i\Big[\Gamma(z)\Gamma(z+s-\kappa_{\mathcal A'})\int\varrho_{\mathcal A'}^{-z}d\phi_{\mathcal A'}\Big]_{z=t+\eta};
\end{multline*}
at $q=1$ this is the closing of Theorem~\ref{thm:nodeclosing}(ii) toward the level $d$. Swapping $\mathcal A$ and $\mathcal A'$, $q^{-t}$ times the same expansion in $1/q$ holds for $q>\operatorname{reach}(R)/\inf\varrho_{\mathcal A}$, and is the closing toward the level $r$. The two ranges never overlap.
\end{enumerate}
\end{theorem}

The closings of a node of Section~\ref{sec:onescale} are therefore not two separate devices but the two expansions, at $q\to0$ and at $q\to\infty$, of one function of the weight of its hinge, holomorphic on two sheets around the origin and given everywhere by the contour integral of (iii), whose integrand is the product of the Mellin transforms of the two sides over $\Gamma(s)\Gamma(t)$ (Proposition~\ref{prop:faces}(iii)). The unweighted period is the value at $q=1$; a closing reaches it when $1$ lies inside its range, its terms do not tend to zero when $1$ lies beyond it, and at the common end of the two ranges the root rate is one and domination decides nothing (Theorem~\ref{thm:nodeintervals}(ii) to (iv), the ranges being read at the reaches of the two sides); and the coefficients of each expansion are the reusable lattice data of one side paired with the Mellin data of the other, the same for every weight.

\section{Changing masses and weights}

% Reference 20: prop:local; source line 6582 (one table).
% Proved in the last part: prop:local.
\begin{proposition}[Carrying a calculation through a change]\label{prop:local}
Let the tables of a family have a common hinge tree rooted at $o$ (Definition~\ref{def:hingetree}), with every law computed by Definition~\ref{def:treeproc}.
\begin{enumerate}[label=(\roman*)]
\item A change at a vertex $v'$ of the weights, of the masses keeping their total, or of the tables $\mathcal H_{v',j}$ keeping their ranks and their numbers of marks, is carried through by recomputing the laws at the vertices on the path from $v'$ to $o$, each from the new law of the vertex directly below it on the path and the old laws of the other vertices directly below it, and then the root.
\item A change of the weight of the hinge of the edge $\mathfrak e_v$ by the factor $w_j$ in table $j$ replaces the law of $F_v$ by $w_1^{\kappa_{F_v}}\prod_jw_j^{-s_j}$ times its image under $(x_{j,K},x_{j,\Psi})_j\mapsto\big((w_1/w_j)\,x_{j,K},\,w_1w_j\,x_{j,\Psi}\big)_j$ (Theorem~\ref{thm:hingeweight}(i)), for one table by $w^{\kappa_{\mathcal A_v}-s}$ times its image under $x\mapsto w^2x$; nothing below $v$ changes, and above $v$ (i) applies along the path from $v$ to $o$.
\item A change of masses at two vertices $v',v''$ keeping the total mass changes the hinge walls named in Proposition~\ref{prop:work}(iii), those on the edges of the path between $v'$ and $v''$, and is carried through by recomputing the laws on the paths from $v'$ and from $v''$ to $o$.
\end{enumerate}
\end{proposition}

A query about one part of a construction is answered by the part of the computation that sees it. The laws on the edges off the path are kept, and with them the work that made them; a construction with many equal subtrees, or one asked many queries about one of its nodes, pays for each subtree once and for each query along one path.


\chapter{Holds and the lean convolution}\label{ch:holding}

A table built along a hinge tree is read from the laws of its subtrees, and at a prime node the twisted core is integrated over its own simplex once for every value of the twists (Definition~\ref{def:treeproc}(4)). Holding the marks removes that integral, and with it the dependence of the reusable data on the query. Every mark of every letter of a family can be twisted, and a twist can be \emph{held} at a load, by a Mellin transform. Twists separate any word in positive letters into its factors, its period becoming one fibre integral of the periods of the factors (Theorem~\ref{thm:factors}), with the convergence read off the Newton sums (Theorem~\ref{thm:poswordconv}); the chapter opens with these words. Holding every mark separates the letters of the family completely: the hold is the product of the periods of its letters alone, the periods of its tables at their full leans, each read at its own load (Lemma~\ref{lem:holdall}). Undoing the hold writes the period at all exponents, every mass and every twist as one fibre integral of dimension $(2N-1)(n-1)$ over the splits of the masses among the letters (Theorem~\ref{thm:leanconv}). The periods at the full leans of a table built along a hinge tree are the products of those of its nodes over Beta corners at its edges, with no integral at all (Theorem~\ref{thm:leantree}), so the period of a family of such tables is one fibre integral of node data (Theorem~\ref{thm:network}).

\section{Newton words}

Twists separate any product of powers of positive letters into its factors, and nothing in this uses the exchange property. The convergence criterion of Theorem~\ref{thm:exist}, whose sectors rest on the greedy bases that need the exchange property (Lemma~\ref{lem:greedy}), extends to such products through the Gramma and its domain (Theorem~\ref{thm:poswordconv}), and on its domain the separation holds (Theorem~\ref{thm:factors}). The holds of the next section are the case of the letters of a family (Lemma~\ref{lem:holdall}).

% Reference 20: def:posword; source line 6772 (Newton words, positive exponents).
% Proved in the last part: def:posword.
\begin{definition}[Positive letters; Newton words and sums]\label{def:posword}
A \emph{positive letter} on $\cE$ is a nonzero homogeneous polynomial $L=\sum_kc_kx^{t_k}$ of degree $\deg L\ge1$ with positive coefficients $c_k$ on distinct exponents $t_k\in\mathbb N^{\cE}$, its \emph{support}; its Newton polytope is $\mathcal N_L=\operatorname{conv}\{t_k\}$, in the hyperplane $|x|=\deg L$. A \emph{product of powers} of positive letters $W=\prod_iL_i^{-w_i}$, with exponents $w_i\in\mathbb C$, has the degree $D=\sum_iw_i\deg L_i$, and for masses $\mvec\in\mathbb C^{\cE}$ with $m_{\cE}=D$ its period $P_W(\mvec)$ is the integral of the form $\alpha^{\mvec-\one}W\,d\alpha$ over a section (Lemma~\ref{lem:chart}), wherever it converges absolutely; it does not depend on the section. It is a \emph{Newton word} when every $w_i>0$, with the \emph{Newton sum} $\mathcal N(W)=\sum_iw_i\mathcal N_{L_i}$, in the hyperplane $|x|=D$. A Newton word is \emph{split into factors} when it is written $W=W_1\cdots W_p$, each factor a product of some of its letters with their exponents. The family word of a family is a product of powers of its letters, with $P_W(\mvec)=P_F(\mvec;\mathbf s,\mathbf t)$, and a Newton word in its active letters when its exponents are nonnegative; for a table $\mathcal N_K=P(\mathcal H)$ and $\mathcal N_\Psi=P(\mathcal H^\ast)$, and for a family $\mathcal N(W)$ is the Newton sum $\mathcal N_F$ of Theorem~\ref{thm:saddle}.
\end{definition}

% Reference 20: thm:poswordconv; source line 6776
\begin{theorem}[Convergence of a Newton word]\label{thm:poswordconv}
$P_W(\mvec)$ converges exactly when $\mathcal N(W)$ has dimension $n-1$ and $\re\mvec\in\mathrm{ri}\,\mathcal N(W)$; it is then holomorphic in $\mvec$ on the balance hyperplane. For a family of tables this is the condition that every proper wall be positive (Theorems~\ref{thm:projective} and~\ref{thm:saddle}(iii)).
\end{theorem}

% Proved in the last part: prop:oneletter (new).
\begin{proposition}[One positive letter]\label{prop:oneletter}
Let $L=\sum_kc_kx^{t_k}$ be a positive letter with $\dim\mathcal N_L=n-1$, $\mathbf a\in\mathbb C^{\cE}$ with $\re|\mathbf a|>0$, and $w=|\mathbf a|/\deg L$.
\begin{enumerate}[label=(\roman*)]
\item $P_{L^{-w}}(\mathbf a)=\deg L\,\sG_{\mathbf a;(t_k)_k}\big((c_k)_k\big)/\Gamma(w)$, both sides converging together.
\item $P_{L^{-w}}(\mathbf a)$ converges exactly when $\re\mathbf a\in(\re w)\,\mathrm{ri}\,\mathcal N_L$.
\item For real $\mathbf a$ at which it converges, every $\boldsymbol\eta\in\mathbb R^{\cE}$ with $|\boldsymbol\eta|=0$ and every $\varepsilon\in(0,1)$,
\[
\big|P_{L^{-w}}(\mathbf a+i\boldsymbol\eta)\big|\le\min\Big(1,\ \big(\sin\tfrac{\pi\varepsilon}2\big)^{-w}e^{-(1-\varepsilon)h_L(\boldsymbol\eta)}\Big)\,P_{L^{-w}}(\mathbf a),
\]
with $h_L$ the support function of $\{y:\ |\ip{t_k}y|<\frac\pi2\ \text{for every }k\}$, a norm on $\mathbb R^{\cE}$.
\end{enumerate}
\end{proposition}

% Reference 20: thm:blocks; source line 6787
\begin{theorem}[Separation into factors]\label{thm:factors}
Let $W=W_1\cdots W_p$ be split into factors with $\dim\mathcal N(W_i)=n-1$ for every $i$, of degrees $D_i$, and let $\mvec\in\mathbb C^{\cE}$ have $m_{\cE}=D$. Points $(\mathbf c_i)_i$ with $\mathbf c_i\in\mathrm{ri}\,\mathcal N(W_i)$ and $\sum_i\mathbf c_i=\re\mvec$ exist exactly when $\re\mvec\in\mathrm{ri}\,\mathcal N(W)$, and on each of them, for a fixed mark $e_0$ and all twists $\mathbf x_i\in(0,\infty)^{\cE}$ of the factors,
\[
\int_{\alpha_{e_0}=1}\alpha^{\mvec-\one}\prod_iW_i(\mathbf x_i\circ\alpha)\,d\alpha'=\int_{\mathbf a_1+\cdots+\mathbf a_p=\mvec,\ |\mathbf a_i|=D_i\ (i<p)}\ \prod_iP_{W_i}(\mathbf a_i)\,\mathbf x_i^{-\mathbf a_i},
\]
a fibre integral of dimension $(p-1)(n-1)$, absolutely convergent, $P_{W_i}$ the period of the factor $W_i$ at its own masses.
\end{theorem}

% Reference 20: cor:exponentsadd; source line 6798
\begin{corollary}[Exponents add on one letter]\label{cor:exponentsadd}
For a positive letter $L$ with $\dim\mathcal N_L=n-1$, exponents $w,w'>0$ and $\re\mvec\in(w+w')\,\mathrm{ri}\,\mathcal N_L$,
\[
P_{L^{-w-w'}}(\mvec)=\int_{\mathbf a+\mathbf a'=\mvec,\ |\mathbf a|=w\deg L}P_{L^{-w}}(\mathbf a)\,P_{L^{-w'}}(\mathbf a').
\]
At $L=e_1(x)=\sum_{e\in\cE}x_e$, $P_{L^{-w}}(\mathbf a)=\Gst n1{\mathbf a}{w}$ and this is Theorem~\ref{thm:sharemerge}; for an identically self-dual table $\Psi=K$, and the lean convolution is this identity at $L=K$, $w=s$, $w'=t$, $P_{\mathcal H}(\mvec;s,t)=P_{\mathcal H}(\mvec;s+t,0)$.
\end{corollary}

% Reference 20: prop:simplicial; source line 6809
\begin{proposition}[Simplicial letters]\label{prop:simplicial}
If the support of a positive letter $L$ consists of $n$ linearly independent exponents $t_1,\dots,t_n$, so that $L$ is \emph{simplicial}, with the configuration $t=(t_1,\dots,t_n)$, then for $\mathbf a$ with $|\mathbf a|=w\deg L$ and the shares $\mathbf b=t^{-1}\mathbf a$,
\[
P_{L^{-w}}(\mathbf a)=\frac{\deg L}{|\det t|}\,\Gst n1{\mathbf b}{w}\,\prod_kc_k^{-b_k},\qquad |\mathbf b|=w,
\]
converging exactly when $\re\mathbf b>0$: the Gramma of Proposition~\ref{prop:oneletter}(i) at a square configuration (Corollary~\ref{cor:square}). A table that is one molecule has a simplicial basis letter exactly when it has $n$ bases.
\end{proposition}

% Reference 20: lem:concat; source line 6820 (two parts).
% Proved in the last part: lem:concat.
\begin{lemma}[Concatenation of Grammas]\label{lem:concat}
Let $t^{(1)},\dots,t^{(p)}$, $p\ge2$, be configurations on $\cE$ with $t^{(2)},\dots,t^{(p)}$ spanning $\mathbb R^{\cE}$, $\boldsymbol\xi^{(1)},\dots,\boldsymbol\xi^{(p)}$ coefficients, and $\mathbf c_i\in\interior\cone(t^{(i)})$ for $i\ge2$ with $\re\nu-\sum_{i\ge2}\mathbf c_i\in\interior\cone(t^{(1)})$. Then
\begin{multline*}
\sG_{\nu;t^{(1)}\oplus\cdots\oplus t^{(p)}}\big(\boldsymbol\xi^{(1)}\oplus\cdots\oplus\boldsymbol\xi^{(p)}\big)\\
=\frac1{(2\pi i)^{(p-1)n}}\int_{\mathbf c_2+i\mathbb R^{\cE}}\!\!\cdots\int_{\mathbf c_p+i\mathbb R^{\cE}}\sG_{\nu-\sum_{i\ge2}\mathbf w_i;t^{(1)}}(\boldsymbol\xi^{(1)})\prod_{i\ge2}\sG_{\mathbf w_i;t^{(i)}}(\boldsymbol\xi^{(i)})\,d\mathbf w_2\cdots d\mathbf w_p ,
\end{multline*}
absolutely convergent: a Gramma of a concatenated configuration is the Mellin convolution of the Grammas of its parts.
\end{lemma}

For one table the concatenation $(\one_T)_T\oplus(\one_{\cE\setminus T})_T$ with the coefficients $\boldsymbol\theta\oplus\lambda\boldsymbol\theta$ is the Gramma $\sG_{\mvec;(\one_T)_T\oplus(\one_{\cE\setminus T})_T}(\boldsymbol\theta\oplus\lambda\boldsymbol\theta)=\int x^{\mvec-\one}e^{-K(x)-\lambda\Psi(x)}dx$, and Lemma~\ref{lem:concat} writes it as $(2\pi i)^{-n}\int\sG_{\mvec-\mathbf b;(\one_T)_T}(\boldsymbol\theta)\,\sG_{\mathbf b;(\one_{\cE\setminus T})_T}(\boldsymbol\theta)\,\lambda^{-|\mathbf b|/r^\ast}d\mathbf b$, by the torus homogeneity of $\sG_{\mathbf b;(\one_{\cE\setminus T})_T}$. The Mellin transform in $\lambda$ at $t$ pins $|\mathbf b|=r^\ast t$, the pin contributing the factor $r^\ast$, and through Proposition~\ref{prop:fullleans}(i) returns $\Gamma(s)\Gamma(t)/r$ times the lean convolution of Theorem~\ref{thm:leanconv} at $\mathbf x=\one$; Proposition~\ref{prop:leanmoment} states the resulting identity $\int_0^\infty\lambda^{t-1}\sG_{\mvec;(\one_T)_T\oplus(\one_{\cE\setminus T})_T}(\boldsymbol\theta\oplus\lambda\boldsymbol\theta)\,d\lambda=\Gamma(s)\Gamma(t)P_{\mathcal H}(\mvec;s,t)/r$. So at fixed masses the lean is the Mellin variable of the relative scale $\lambda$ of the two letters, and the two-step integration separates what the masses fix, the Gramma of the concatenated configuration, from what the exponents choose, one Mellin transform in one variable.


\section{Holds and the periods at the full leans}

% Reference 20: def:hold; source line 6603 (one table, twisted in its cobasis letter).
% Proved in the last part: def:hold.
\begin{definition}[Holding]\label{def:hold}
Let $F=((\cB_j,\theta_j);s_j,t_j)_{j=1}^N$ be a family on $\cE$, $n=|\cE|\ge2$, with exponents $s_j,t_j>0$, whose tables are each one molecule, table $j$ of rank $r_j$ and corank $r^\ast_j=n-r_j$, so that $0<r_j<n$ and the bases of every table span $\mathbb R^{\cE}$: a vector orthogonal to every $\one_T$, $T\in\cB_j$, is constant by Lemma~\ref{lem:circuit}, and then zero since $r_j>0$. $P_F(\mvec;\mathbf s,\mathbf t)$ is the period (Theorem~\ref{thm:projective}) on the balance wall $m_{\cE}=\sum_j(r_js_j+r^\ast_jt_j)$, and $F[\mathbf x]$ is $F$ twisted by $\mathbf x=(\mathbf x^K_j,\mathbf x^\Psi_j)_j$ (Definition~\ref{def:twist}). For $G\subseteq\cE$, masses $\mvec$ and real loads $\boldsymbol\beta=(\boldsymbol\beta^K_j,\boldsymbol\beta^\Psi_j)_j$, $\boldsymbol\beta^K_j,\boldsymbol\beta^\Psi_j\in\mathbb R^G$, with $\boldsymbol\beta^K_1$ omitted, \emph{holding the marks of $G$ at $\boldsymbol\beta$} is the Mellin transform in the twists of every letter but $K_1$ at the marks of $G$, the twist of $K_1$ and the twists at the other marks kept at $\one$,
\[
\operatorname{Hold}_G(\boldsymbol\beta)=\int_{((0,\infty)^G)^{2N-1}}\prod_{j\ge2}(\mathbf x^K_j)^{\boldsymbol\beta^K_j-\one}\prod_{j=1}^N(\mathbf x^\Psi_j)^{\boldsymbol\beta^\Psi_j-\one}\,P_{F[\mathbf x]}(\mvec;\mathbf s,\mathbf t)\,d\mathbf x ,
\]
in $[0,\infty]$ for real masses and loads; for complex masses or loads it is taken where it converges absolutely, that is where it is finite at the real parts.
\end{definition}

The twists of every letter but $K_1$ are held. By the torus relation of Definition~\ref{def:twist}, a twist of every mark of one letter by one factor multiplies the period by that factor to the power $-r_js_j$, or $-r^\ast_jt_j$ for a cobasis letter, so that a hold of every mark is taken with one twist of each held letter fixed; and a twist of every letter by one vector $\boldsymbol\lambda$ multiplies the period by $\boldsymbol\lambda^{-\mvec}$, so that the twist of $K_1$ is kept at $\one$. A twist is a change of weights: Proposition~\ref{prop:weightstwists} carries product weights on the bases into twists.

% Reference 20: def:fullleanfns; source line 6617 (one table).
% Proved in the last part: def:fullleanfns.
\begin{definition}[Loads; admissible points]\label{def:fullleanfns}
For a split of the masses among the letters of $F$, $\mvec=\sum_j(\mathbf a_j+\mathbf b_j)$, $\mathbf a_j$ is the \emph{$K_j$-load} and $\mathbf b_j$ the \emph{$\Psi_j$-load}, as a $K$-share of a pattern loads its marks and a $\Psi$-share the marks it leaves out (Definition~\ref{def:shares}); for one table, $\mvec=\mathbf a+\mathbf b$ with the $K$-load $\mathbf a$ and the $\Psi$-load $\mathbf b$. The periods of a table $\mathcal H$ of $F$, with the weights $\boldsymbol\theta$, of rank $r$ and corank $r^\ast$, at its two full leans are read as functions of their loads, $P_{\mathcal H}(\mathbf a;|\mathbf a|/r,0)$ and $P_{\mathcal H}(\mathbf b;0,|\mathbf b|/r^\ast)$, the exponent of the one active letter being the total load divided by its degree. Its letters give the Grammas (Definition~\ref{def:gramma})
\[
\sG_{\mathbf a;(\one_T)_{T\in\mathcal H}}(\boldsymbol\theta)=\int_{\mathbb R^{\cE}_{>0}}x^{\mathbf a-\one}e^{-K_{\mathcal H}(x)}dx,\qquad \sG_{\mathbf b;(\one_{\cE\setminus T})_{T\in\mathcal H}}(\boldsymbol\theta)=\int_{\mathbb R^{\cE}_{>0}}x^{\mathbf b-\one}e^{-\Psi_{\mathcal H}(x)}dx .
\]
A real load $\mathbf a$ with $|\mathbf a|=rs$, $s>0$, is \emph{admissible for the full basis lean} when $\mathbf a\in s\cdot\mathrm{ri}\,P(\mathcal H)$, $P(\mathcal H)$ the base polytope of $\mathcal H$ (Theorem~\ref{thm:newton}) and $\mathrm{ri}$ the relative interior, and a real $\mathbf b$ with $|\mathbf b|=r^\ast t$, $t>0$, is \emph{admissible for the full cobasis lean} when $\mathbf b\in t\cdot\mathrm{ri}\,P(\mathcal H^\ast)$, $P(\mathcal H^\ast)$ the base polytope of the dual table. A real point $\mathbf c=(\mathbf c^K_j,\mathbf c^\Psi_j)_j$ is \emph{admissible} for $(\mvec;\mathbf s,\mathbf t)$ when $\sum_j(\mathbf c^K_j+\mathbf c^\Psi_j)=\re\mvec$, every $\mathbf c^K_j$ is admissible for the full basis lean of table $j$ at $s_j$ and every $\mathbf c^\Psi_j$ for its full cobasis lean at $t_j$; for one table it is named by its $\Psi$-load $\mathbf c=\mathbf c^\Psi_1$, with $\re\mvec-\mathbf c$ admissible for the full basis lean and $\mathbf c$ for the full cobasis lean.
\end{definition}

% Reference 20: prop:fullleans; source line 6629 (one table).
% Proved in the last part: prop:fullleans.
\begin{proposition}[The periods at the two full leans]\label{prop:fullleans}
In (i), (ii) and (iv) $\mathcal H$ is a table of $F$, with the weights $\boldsymbol\theta$, of rank $r$ and corank $r^\ast$.
\begin{enumerate}[label=(\roman*)]
\item For $\mathbf a\in\mathbb C^{\cE}$ with $\re|\mathbf a|>0$ and $s=|\mathbf a|/r$, $P_{\mathcal H}(\mathbf a;s,0)=r\,\sG_{\mathbf a;(\one_T)_T}(\boldsymbol\theta)/\Gamma(s)$, both sides converging together; likewise $P_{\mathcal H}(\mathbf b;0,t)=r^\ast\sG_{\mathbf b;(\one_{\cE\setminus T})_T}(\boldsymbol\theta)/\Gamma(t)$ with $t=|\mathbf b|/r^\ast$. By Cremona, $P_{\mathcal H}(\mathbf b;0,t)=P_{\mathcal H}(t\one-\mathbf b;t,0)$.
\item $P_{\mathcal H}(\mathbf a;s,0)$ converges exactly when $\re a_S>(\re s)\,d(S)$ for every nonempty proper $S$, which for real $s>0$ is $\re\mathbf a\in s\cdot\mathrm{ri}\,P(\mathcal H)$; $P_{\mathcal H}(\mathbf b;0,t)$ converges exactly when $\re b_S>(\re t)\,d^\ast(S)$ for every nonempty proper $S$, which for real $t>0$ is $\re\mathbf b\in t\cdot\mathrm{ri}\,P(\mathcal H^\ast)$.
\item At every gauge $\alpha$ the occupation $R_j(\alpha)$ of table $j$ lies in $\mathrm{ri}\,P(\cB_j)$ and its dual occupation $R^\ast_j(\alpha)$ in $\mathrm{ri}\,P(\cB_j^\ast)$ (Theorem~\ref{thm:innerMGF}). At a native gauge $\alpha^\star$ of real masses $\mvec$ on the balance wall (Definition~\ref{def:native}), which exists when every proper wall is positive, the \emph{native loads} $\mathbf a^\star_j=s_jR_j(\alpha^\star)$ and $\mathbf b^\star_j=t_jR^\ast_j(\alpha^\star)$ are a split of $\mvec$ among the letters of $F$, and $\mathbf c=(\mathbf a^\star_j,\mathbf b^\star_j)_j$ is admissible.
\item On a vertical line the periods at the full leans decay exponentially: for admissible real $\mathbf a$, every $\mathbf y\in\mathbb R^{\cE}$ with $|\mathbf y|=0$ and every $\varepsilon\in(0,1)$,
\[
\big|P_{\mathcal H}(\mathbf a+i\mathbf y;s,0)\big|\le\min\Big(1,\ \big(\sin\tfrac{\pi\varepsilon}2\big)^{-s}e^{-(1-\varepsilon)h(\mathbf y)}\Big)\,P_{\mathcal H}(\mathbf a;s,0),
\]
with $h$ the support function of the set $\{y:\ |\ip{\one_T}y|<\frac\pi2\ \text{for every basis }T\}$, a norm on $\mathbb R^{\cE}$; likewise at the full cobasis lean with the vectors $\one_{\cE\setminus T}$.
\end{enumerate}
Items (i), (ii) and (iv) hold for every positive letter with a Newton polytope of dimension $n-1$, in the form of Proposition~\ref{prop:oneletter}.
\end{proposition}

% Reference 20: lem:holdall; source line 6652 (one table).
% Proved in the last part: lem:holdall.
\begin{lemma}[Holding every mark separates the letters]\label{lem:holdall}
Fix $e_0\in\cE$ and let $\mvec\in\mathbb C^{\cE}$ have $m_{\cE}=\sum_j(r_js_j+r^\ast_jt_j)$. For real loads $\boldsymbol\beta=(\boldsymbol\beta^K_j,\boldsymbol\beta^\Psi_j)_j$ on $\cE\setminus e_0$, $\boldsymbol\beta^K_1$ omitted, let $\mathbf a_j$ ($j\ge2$) extend $\boldsymbol\beta^K_j$ by $a_{j,e_0}=r_js_j-|\boldsymbol\beta^K_j|$, let $\mathbf b_j$ extend $\boldsymbol\beta^\Psi_j$ by $b_{j,e_0}=r^\ast_jt_j-|\boldsymbol\beta^\Psi_j|$, and put $\mathbf a_1=\mvec-\sum_{j\ge2}\mathbf a_j-\sum_j\mathbf b_j$. Then, in $[0,\infty]$ for real $\mvec$,
\[
\operatorname{Hold}_{\cE\setminus e_0}(\boldsymbol\beta)=\prod_{j=1}^NP_{\cB_j}(\mathbf a_j;s_j,0)\,P_{\cB_j}(\mathbf b_j;0,t_j),
\]
the twists at $e_0$ being $1$; for complex $\mvec$ and complex $\boldsymbol\beta$ the same holds whenever the integral converges at the real parts.
\end{lemma}

The lemma is the whole mechanism. Every letter but $K_1$ sees its own twists and $K_1$ sees none; a Mellin transform in all of them at once absorbs every mark of each twisted letter into its own load, and what is left of $K_1$ is its own period at the remaining load. Nothing is integrated against anything: the hold is a product.

\section{The lean convolution}

% Reference 20: thm:leanconv; source line 6671 (one table).
% Proved in the last part: thm:leanconv.
\begin{theorem}[The lean convolution]\label{thm:leanconv}
Let $\mvec\in\mathbb C^{\cE}$ have $m_{\cE}=\sum_j(r_js_j+r^\ast_jt_j)$.
\begin{enumerate}[label=(\roman*)]
\item Admissible points for $(\mvec;\mathbf s,\mathbf t)$ exist exactly when every proper wall of $\re\mvec$ is positive; they form a convex set, and for real $\mvec$ the native loads (Proposition~\ref{prop:fullleans}(iii)) are one of them.
\item On every admissible point and for every twist $\mathbf x=(\mathbf x^K_j,\mathbf x^\Psi_j)_j$ (Definition~\ref{def:twist}),
\[
P_{F[\mathbf x]}(\mvec;\mathbf s,\mathbf t)=\int_{\substack{\sum_j(\mathbf a_j+\mathbf b_j)=\mvec\\ |\mathbf a_j|=r_js_j\ (j\ge2),\ |\mathbf b_j|=r^\ast_jt_j}}\ \prod_{j=1}^NP_{\cB_j}(\mathbf a_j;s_j,0)\,P_{\cB_j}(\mathbf b_j;0,t_j)\,(\mathbf x^K_j)^{-\mathbf a_j}(\mathbf x^\Psi_j)^{-\mathbf b_j},
\]
the fibre integral of Definition~\ref{def:fibre}, of dimension $(2N-1)(n-1)$, absolutely convergent; in the coordinates of the loads of every letter but $K_1$ at the marks other than a fixed mark $e_0$ it is $(2\pi i)^{-(2N-1)(n-1)}\int\cdots$ over the imaginary translate through the admissible point, the loads at $e_0$ and $\mathbf a_1$ being fixed by the pins. At $\mathbf x=\one$ it is the period.
\item The modulus of the integrand is at most its value at the real point,
\[
\prod_jP_{\cB_j}(\mathbf c^K_j;s_j,0)\,P_{\cB_j}(\mathbf c^\Psi_j;0,t_j)\,(\mathbf x^K_j)^{-\mathbf c^K_j}(\mathbf x^\Psi_j)^{-\mathbf c^\Psi_j},
\]
and it decays exponentially in the imaginary part of the load of every letter but $K_1$, at the rate of Proposition~\ref{prop:fullleans}(iv) for its table.
\end{enumerate}
With exponents $s_j,t_j\ge0$, not all zero, the same holds over the active letters, $p$ of them, any one of them in the place of $K_1$, in dimension $(p-1)(n-1)$.
\end{theorem}

For one table, at $\mathcal H=\binom{\{1,2\}}1$ the lean convolution is Corollary~\ref{cor:fourgammas}, and at the circuit $\binom{\cE}{n-1}$ it is Corollary~\ref{cor:cycle}, the periods at the full leans being those of Proposition~\ref{prop:spleans}. Opening each of them on its share side (Proposition~\ref{prop:shareform}) gives the atom formula of Corollary~\ref{cor:atomwall}, one share per basis and letter: the atom formula opens every letter over its bases and every mass over the letters at once, while the lean convolution keeps each letter closed, read at its full lean, and opens only the split of each mass among the letters.

% Reference 20: cor:retaintable; source line 6697 (one table).
% Proved in the last part: cor:retaintable.
\begin{corollary}[What a family retains]\label{cor:retaintable}
The Grammas $\sG_{\mathbf a;(\one_T)_{T\in\cB_j}}(\boldsymbol\theta_j)$, $j=1,\dots,N$, each one function of the $n$ variables $\mathbf a$, determine every period of $F$ and of all its twists: at all exponents $s_j,t_j\ge0$, not all zero, at all masses on the balance wall with every proper wall positive, and at every twist. At the exponents $(\mathbf s,\mathbf t)$ only the values of the Gramma of table $j$ on the hyperplanes $|\mathbf a|=r_js_j$ and $|t_j\one-\mathbf b|=r_jt_j$ enter, through $P_{\cB_j}(\mathbf a;s_j,0)=r_j\,\sG_{\mathbf a;(\one_T)_T}(\boldsymbol\theta_j)/\Gamma(s_j)$ and $P_{\cB_j}(\mathbf b;0,t_j)=r_j\,\sG_{t_j\one-\mathbf b;(\one_T)_T}(\boldsymbol\theta_j)/\Gamma(t_j)$. Conversely the periods of table $j$ at its full basis leans determine its Gramma, through Proposition~\ref{prop:fullleans}(i).
\end{corollary}

A cell of a hinge tree retains, for one pair of exponents, one measure on the half-line (Theorem~\ref{thm:retained}); a family retains, for all exponents and all masses at once, one function of the loads for each of its tables. The two retentions serve different queries; Lemma~\ref{lem:hingesplit} relates them, and Section~\ref{sec:holdnodes} compares them.

\section{Holding some marks}

% Reference 20: thm:pinnedhold; source line 6708 (one table).
% Proved in the last part: thm:pinnedhold.
\begin{theorem}[Pinned holds]\label{thm:pinnedhold}
Let every proper wall of $\re\mvec$ be positive, $e_0$ a fixed mark, $\varnothing\neq G\subseteq\cE\setminus e_0$, and $\boldsymbol\beta=(\boldsymbol\beta^K_j,\boldsymbol\beta^\Psi_j)_j$, $\boldsymbol\beta^K_1$ omitted, the restriction to $G$ of the loads of an admissible point. Then for every real $\boldsymbol\eta=(\boldsymbol\eta^K_j,\boldsymbol\eta^\Psi_j)_j$ of the same shape the hold $\operatorname{Hold}_G(\boldsymbol\beta+i\boldsymbol\eta)$ converges absolutely and equals the fibre integral of Theorem~\ref{thm:leanconv}(ii) at $\mathbf x=\one$ with the $(2N-1)|G|$ further pins $\mathbf a_j|_G=\boldsymbol\beta^K_j+i\boldsymbol\eta^K_j$ ($j\ge2$) and $\mathbf b_j|_G=\boldsymbol\beta^\Psi_j+i\boldsymbol\eta^\Psi_j$,
\[
\operatorname{Hold}_G(\boldsymbol\beta+i\boldsymbol\eta)=\int_{\sum_j(\mathbf a_j+\mathbf b_j)=\mvec,\ |\mathbf a_j|=r_js_j\ (j\ge2),\ |\mathbf b_j|=r^\ast_jt_j,\ \text{the pins on }G}\ \prod_{j=1}^NP_{\cB_j}(\mathbf a_j;s_j,0)\,P_{\cB_j}(\mathbf b_j;0,t_j),
\]
of dimension $(2N-1)(n-1-|G|)$, in the coordinates of the loads of every letter but $K_1$ at the marks off $G\cup e_0$ the integral $(2\pi i)^{-(2N-1)(n-1-|G|)}\int\cdots$. At $G=\cE\setminus e_0$ it is Lemma~\ref{lem:holdall}.
\end{theorem}

For one mark the hold has a closed form in the minors at that mark. For a mark $g$ that is neither a loop nor a coloop of any table, the patterns of the bunch $\{g\}$ in table $j$ are $\varnothing$ and $\{g\}$, with the children the deletion $\cB_j\setminus g$, of rank $r_j$ and corank $r^\ast_j-1$, and the contraction $\cB_j/g$, of rank $r_j-1$ and corank $r^\ast_j$ (Definition~\ref{def:patterns}). For $K$-loads $a_j$ and $\Psi$-loads $b_j$ at $g$ the children family of $\{g\}$ (Definition~\ref{def:shares}) at the $K$-shares $(s_j-a_j,a_j)$ and the $\Psi$-shares $(b_j,t_j-b_j)$ of the patterns $\varnothing$ and $\{g\}$ is
\[
F_{\{g\}}\big((s_j-a_j,a_j),(b_j,t_j-b_j)\big)_j=\Big(\big((\cB_j\setminus g,\theta_j);\ s_j-a_j,\ b_j\big),\ \big((\cB_j/g,\theta_j);\ a_j,\ t_j-b_j\big)\Big)_j,
\]
a family of $2N$ tables on $\cE\setminus g$ whose word is $\prod_jK_{j\setminus g}^{a_j-s_j}\Psi_{j\setminus g}^{-b_j}K_{j/g}^{-a_j}\Psi_{j/g}^{b_j-t_j}$. At masses $\mvec|_{\cE\setminus g}$ with $\sum_j(a_j+b_j)=m_g$ it is on its balance wall: its degree is $\sum_j\big((s_j-a_j)r_j+b_j(r^\ast_j-1)+a_j(r_j-1)+(t_j-b_j)r^\ast_j\big)=\sum_j(r_js_j+r^\ast_jt_j)-m_g$.

% Reference 20: prop:holdone; source line 6725 (one table).
% Proved in the last part: prop:holdone.
\begin{proposition}[Holding one mark]\label{prop:holdone}
Let $g$ be a mark that is neither a loop nor a coloop of any table, $\mvec$ real masses on the balance wall, $a_j$ ($j\ge2$) and $b_j$ ($j\le N$) real, and $a_1=m_g-\sum_{j\ge2}a_j-\sum_jb_j$. The hold
\[
\operatorname{Hold}_{\{g\}}\big((a_j)_{j\ge2},(b_j)_j\big)=\int_{(0,\infty)^{2N-1}}\prod_{j\ge2}(x^K_j)^{a_j-1}\prod_{j=1}^N(x^\Psi_j)^{b_j-1}\,P_{F[\mathbf x@g]}(\mvec;\mathbf s,\mathbf t)\,d\mathbf x,\qquad x^K_1=1,
\]
is finite exactly when $0<a_j<s_j$ and $0<b_j<t_j$ for every $j$ and every proper wall of $F_{\{g\}}\big((s_j-a_j,a_j),(b_j,t_j-b_j)\big)_j$ at $\mvec|_{\cE\setminus g}$ is positive, and then
\[
\operatorname{Hold}_{\{g\}}\big((a_j)_{j\ge2},(b_j)_j\big)=\prod_{j=1}^N\Gst21{a_j,\ s_j-a_j}{s_j}\Gst21{b_j,\ t_j-b_j}{t_j}\,P_{F_{\{g\}}((s_j-a_j,a_j),(b_j,t_j-b_j))_j}\big(\mvec|_{\cE\setminus g}\big);
\]
the identity holds as well for complex loads and complex masses whose real parts are such data. Where the loads are those at $g$ of an admissible point, Theorem~\ref{thm:pinnedhold} writes the same value as the lean convolution with the pins at $g$. If $m_g>0$, then for every letter of $F$ whose exponent $u$ exceeds $m_g$, as the twists at $g$ of all the other letters decrease to $0$, the twist of that letter staying $1$, $P_{F[\mathbf x@g]}(\mvec;\mathbf s,\mathbf t)$ increases, in $(0,\infty]$, to $\Gst21{m_g,\ u-m_g}{u}$ times the period at $\mvec|_{\cE\setminus g}$ of the children family of $\{g\}$ at the shares that put the whole mass $m_g$ on that letter: $(s_j-m_g,m_g)$ if it is $K_j$, $(m_g,t_j-m_g)$ if it is $\Psi_j$, $(s_{j'},0)$ for every other letter $K_{j'}$ and $(0,t_{j'})$ for every other letter $\Psi_{j'}$.
\end{proposition}

The ends of the twist are the faces of the children family at which every letter but one has the mark $g$ switched off: the chosen letter keeps both its minors and every other letter keeps one. For one table the two limits are those of $P_{\mathcal H[x@g]}$ as $x\downarrow0$ and of $x^{m_g}P_{\mathcal H[x@g]}$ as $x\uparrow\infty$ (Definition~\ref{def:twist}), and a weight on a mark driven to one end of its range leaves the deletion and the contraction of the mark, two tables carrying three active letters between them. With Proposition~\ref{prop:weightstwists} and Definition~\ref{def:twist}, a weight $w$ on the bases through $g$ is the twist $w^{-2}$ up to the factor $w^{-m_g-t}$, so the period behaves like $w^{-m_g-t}$ times the first limit as $w\to\infty$ and like $w^{m_g-t}$ times the second as $w\to0$, where these are finite.

% Reference 20: lem:hingesplit; source line 6742 (one cell).
% Proved in the last part: lem:hingesplit.
\begin{lemma}[The Mellin transform of a law is the period of a children family]\label{lem:hingesplit}
Let $F=((\mathcal A_j,\theta_j);s_j,t_j)_{j=1}^N$ be a cell family on $B\sqcup\{g\}$ with masses $\mvec$ on $B$, hinge wall $\kappa_F$, base form $\phi_F$ and ratio vector $\boldsymbol\varrho_F$ (Definition~\ref{def:hinged}). In this statement let $\mathbf z=(z_{j,K},z_{j,\Psi})_j$ have $\sum_j(z_{j,K}+z_{j,\Psi})=\kappa_F$, and write $\boldsymbol\varrho_F^{-\mathbf z}=\prod_{j\ge2}\big(\gamma^K_{\mathcal A_j}/\gamma^K_{\mathcal A_1}\big)^{-z_{j,K}}\prod_j\big(\gamma^\Psi_{\mathcal A_j}/\gamma^K_{\mathcal A_1}\big)^{-z_{j,\Psi}}$, the coordinates of the ratio vector to the powers $-z_{j,K}$ ($j\ge2$) and $-z_{j,\Psi}$. For real $\mathbf z$, in $[0,\infty]$,
\[
\int_{\Delta_B}\boldsymbol\varrho_F^{-\mathbf z}\,d\phi_F=P_{F_{\{g\}}((z_{j,K},\,s_j-z_{j,K}),(t_j-z_{j,\Psi},\,z_{j,\Psi}))_j}(\mvec),
\]
the period of the children family of the bunch $\{g\}$ of $F$ (Definition~\ref{def:shares}), its tables $\mathcal A_j\setminus g$ at the exponents $(z_{j,K},t_j-z_{j,\Psi})$ and $\mathcal A_j/g$ at $(s_j-z_{j,K},z_{j,\Psi})$, on its balance wall at $\mvec$. Let moreover $s_j,t_j>0$, and put the mass $\sum_j(s_j+t_j)-\kappa_F$ on $g$, which puts the family of the tables $\mathcal A_j$ on its balance wall. For complex $\mathbf z$ whose real part has $0<\re z_{j,K}<s_j$ and $0<\re z_{j,\Psi}<t_j$ for every $j$ and is a point at which $\int\boldsymbol\varrho_F^{-\re\mathbf z}d\phi_F$ is finite,
\begin{multline*}
\int_{(0,\infty)^{2N-1}}\prod_{j\ge2}(x^K_j)^{s_j-z_{j,K}-1}\prod_{j=1}^N(x^\Psi_j)^{t_j-z_{j,\Psi}-1}\,P_{F[\mathbf x@g]}\Big(\big(\mvec,\textstyle\sum_j(s_j+t_j)-\kappa_F\big);\mathbf s,\mathbf t\Big)\,d\mathbf x\\
=\prod_{j=1}^N\Gst21{s_j-z_{j,K},\ z_{j,K}}{s_j}\Gst21{t_j-z_{j,\Psi},\ z_{j,\Psi}}{t_j}\int_{\Delta_B}\boldsymbol\varrho_F^{-\mathbf z}\,d\phi_F ,
\end{multline*}
$F[\mathbf x@g]$ the family of the tables $\mathcal A_j$ twisted at $g$, $x^K_1=1$ (Definition~\ref{def:twist}). So the Mellin transform of the law of a cell family is the Mellin transform of its own period in the twists of its hinge, divided by $2N$ Beta corners.
\end{lemma}

% Reference 20: cor:selfdualtwist; source line 6757
\begin{corollary}[Identically self-dual tables twisted at one mark]\label{cor:selfdualtwist}
Let the family be one table $\mathcal H$ at the exponents $(s,t)$, identically self-dual (Corollary~\ref{cor:piecesmarks}(iii)), $g$ a mark, and $\mvec$ on the balance wall with every proper wall positive. For every $x>0$,
\[
P_{\mathcal H[x@g]}(\mvec;s,t)=P_{\mathcal H\setminus g}\big(\mvec|_{\cE\setminus g};\ s+t-m_g,\ m_g\big)\,P_{\binom{\{g,g'\}}1[(x,1)]}\big((m_g,s+t-m_g);s,t\big),
\]
the second factor the period $\int_0^\infty v^{m_g-1}(1+v)^{-s}(1+xv)^{-t}dv$ of the twisted list on two marks (Lemma~\ref{lem:twomarks}); and the hold at $g$ is $\Gst21{m_g-b,\ s-m_g+b}{s}\Gst21{b,\ t-b}{t}$ times the same period of $\mathcal H\setminus g$ for every $b$ in its strip: the children family does not depend on the split. At $x=1$ this is Corollary~\ref{cor:piecesmarks}(iii), and Corollary~\ref{cor:fourgammas} recovers it from the hold,
\[
\frac1{2\pi i}\int\Gst21{m_g-b,\ s-m_g+b}{s}\Gst21{b,\ t-b}{t}\,db=\Gst21{m_g,\ s+t-m_g}{s+t}.
\]
\end{corollary}

\chapter{Periods at the full leans and networks along a hinge tree}\label{ch:leantree}

At a full lean the period of a two-sum is the product of the periods of its two cells over one Beta corner (Theorem~\ref{thm:cornerfactor}), and along any hinge tree this is Corollary~\ref{cor:primelean}, which for one table at a full lean is Theorem~\ref{thm:leantree} at the loads $\mvec$. Along a hinge tree this needs only bookkeeping, and the bookkeeping is affine: every marker receives a load that is a fixed affine function of the loads of the marks, the same function at every pair of exponents. So the periods of a substitution at its full leans are products of those of its nodes, read at affine loads, over Beta corners at its edges, with no integral. Where the active letters read every cell by one hinge value, the period along a hinge tree is the product of the periods of the nodes over Beta corners at the edges: at a full lean of a table whose weights factor at its bunches of width one, a product of stacked Gammas, monomials in the weights and the periods of the prime nodes of its canonical hinge tree, and $k$ copies glued along any tree cost one period and a power. Blocks aggregate exactly where every active letter reads them by one hinge value, their total or the inverse of the sum of the reciprocals of their marks (Theorem~\ref{thm:blockagg}). The nodes whose periods at the full leans cost nothing are characterized by counting bases. The period of a family of such tables, at all exponents, every mass and every twist, is then one fibre integral of node data, of dimension $(2N-1)(n-1)$: the lean convolution of Chapter~\ref{ch:holding} with the periods of every letter at its full lean read along the hinge tree of its table (Theorem~\ref{thm:network}).

\section{Molecules and marker loads}

% Reference 20: lem:twosummol; source line 6849 (two cells).
% Proved in the last part: lem:twosummol.
\begin{lemma}[Substitutions of molecules]\label{lem:twosummol}
For a core $\mathcal H$ and cells $\mathcal A_1,\dots,\mathcal A_q$ at its marks (Definition~\ref{def:substitution}), the substitution $\mathcal H\langle\mathcal A_1,\dots,\mathcal A_q\rangle$ is one molecule exactly when $\mathcal H$ and the $\mathcal A_i$, as tables on their marks and hinges, are each one molecule. In particular the two-sum $\mathcal A\oplus_g\mathcal A'$ of two cells is one molecule exactly when $\mathcal A$ and $\mathcal A'$ are.
\end{lemma}

So a substitution of nodes that are molecules is one molecule, and a table with every proper wall positive, which is one molecule by Proposition~\ref{prop:compoundwall}, has every cell of its hinge trees one molecule.

% Reference 20: hinge-tree conventions; source line 6835
\begin{definition}[Branches of a hinge tree]\label{def:branches}
Let $\cB$ be a table on $\cE$, $|\cE|=n$, of rank $r$ and corank $r^\ast$, with a hinge tree $\mathfrak T$ (Definition~\ref{def:hingetree}) whose nodes $\mathcal H_v$ on $\cE_v\sqcup G_v$ are each one molecule, of rank $r_v$ and corank $r^\ast_v=|\cE_v|+|G_v|-r_v$; the weights of $\cB$ are the products of the weights of the nodes (Lemma~\ref{lem:twosumpieces}). For an edge $\mathfrak e=\{v,w\}$ of $\mathfrak T$ and its end $v$, the \emph{branch} $\mathfrak T^{\mathfrak e}_v$ is the component of $\mathfrak T-\mathfrak e$ containing $v$, $\cE^{\mathfrak e}_v$ the union of the $\cE_u$ over its vertices, and $\mathcal A^{\mathfrak e}_v$ the substitution of its nodes, a cell on $\cE^{\mathfrak e}_v\sqcup\{g_{v,\mathfrak e}\}$ of rank $r^{\mathfrak e}_v$ and corank $r^{\ast\mathfrak e}_v=|\cE^{\mathfrak e}_v|+1-r^{\mathfrak e}_v$; $\cB$ is the two-sum of $\mathcal A^{\mathfrak e}_v$ and $\mathcal A^{\mathfrak e}_w$.
\end{definition}

% Reference 20: def:markerloads; source line 6858
\begin{definition}[Marker loads; node loads]\label{def:markerloads}
For a table with a hinge tree as in Definition~\ref{def:branches} and loads $\mathbf a,\mathbf b\in\mathbb C^{\cE}$ with $|\mathbf a|=rs$ and $|\mathbf b|=r^\ast t$, the \emph{marker loads} at the end $v$ of the edge $\mathfrak e$ are
\[
k_{v,\mathfrak e}=s\,r^{\mathfrak e}_v-a_{\cE^{\mathfrak e}_v},\qquad l_{v,\mathfrak e}=t\,r^{\ast\mathfrak e}_v-b_{\cE^{\mathfrak e}_v},
\]
and the \emph{node loads} of $v$ are $\mathbf a_v=\big(\mathbf a|_{\cE_v},(k_{v,\mathfrak e})_{\mathfrak e\ni v}\big)$ and $\mathbf b_v=\big(\mathbf b|_{\cE_v},(l_{v,\mathfrak e})_{\mathfrak e\ni v}\big)$, the load of the marker $g_{v,\mathfrak e}$ being $k_{v,\mathfrak e}$, respectively $l_{v,\mathfrak e}$.
\end{definition}

% Reference 20: lem:nodeloads; source line 6866
\begin{lemma}[Node loads]\label{lem:nodeloads}
In Definition~\ref{def:markerloads}:
\begin{enumerate}[label=(\roman*)]
\item $k_{v,\mathfrak e}+k_{w,\mathfrak e}=s$ and $l_{v,\mathfrak e}+l_{w,\mathfrak e}=t$ on every edge.
\item $|\mathbf a_v|=r_vs$ and $|\mathbf b_v|=r^\ast_vt$ at every node.
\item If $\mathbf a+\mathbf b=\mvec$ with $\mvec$ on the balance wall of $\cB$, then $k_{v,\mathfrak e}+l_{v,\mathfrak e}=s+t-\kappa_{\mathcal A^{\mathfrak e}_v}=\kappa_{\mathcal A^{\mathfrak e}_w}$: the node masses $\mathbf a_v+\mathbf b_v$ are $\mvec|_{\cE_v}$ on the marks and the hinge walls of the branches on the markers, the masses of the twisted core in Theorem~\ref{thm:node}.
\end{enumerate}
\end{lemma}

\section{The periods of a substitution at its full leans}

% Reference 20: thm:leantree; source line 6879
\begin{theorem}[Periods at the full leans along a hinge tree]\label{thm:leantree}
In Definition~\ref{def:markerloads}, for real loads $\mathbf a$ with $|\mathbf a|=rs$, if every marker load $k_{v,\mathfrak e}$ lies in $(0,s)$, then, in $(0,\infty]$,
\[
P_{\cB}(\mathbf a;s,0)=\frac{\prod_vP_{\mathcal H_v}(\mathbf a_v;s,0)}{\prod_{\mathfrak e=\{v,w\}}\Gst21{k_{v,\mathfrak e},\ k_{w,\mathfrak e}}{s}},
\]
and if some marker load lies outside $(0,s)$, then $P_{\cB}(\mathbf a;s,0)=+\infty$. So $P_{\cB}(\mathbf a;s,0)$ is finite exactly when every marker load lies in $(0,s)$ and the period of every node at its full basis lean is finite at its node loads. The same holds at the full cobasis leans, with $\mathbf b$, the loads $l_{v,\mathfrak e}$ and $t$. For complex loads both identities hold on the tube over the real loads at which they are finite, both sides being holomorphic there.
\end{theorem}

In the language of bunches this is Theorem~\ref{thm:cornerfactor} at every bunch of width one of the tree at once: the bunch $\cE^{\mathfrak e}_v$ has, at $(s,0)$, the wall $k_{w,\mathfrak e}$, and its two cells are $\mathcal A^{\mathfrak e}_v$ and $\mathcal A^{\mathfrak e}_w$ (Lemma~\ref{lem:twosumpieces}). What the tree adds is that the loads of all markers are affine in the loads of the marks, fixed once for every pair of exponents, so that the periods of the substitution at its full leans are evaluated at any loads by evaluating each node once.

% Reference 20: cor:nodeconv; source line 6898 (one table).
% Proved in the last part: cor:nodeconv.
\begin{corollary}[Convergence is decided at the nodes]\label{cor:nodeconv}
Let $F=((\cB_j,\theta_j);s_j,t_j)_{j=1}^N$ have exponents $s_j,t_j>0$, and let table $j$ carry a hinge tree $\mathfrak T_j$ as in Definition~\ref{def:branches}, its nodes $\mathcal H_{j,v}$ on $\cE_{j,v}\sqcup G_{j,v}$ each one molecule, with the node loads of Definition~\ref{def:markerloads} read in table $j$. For masses $\mvec$ on the balance wall of $F$ every proper wall of $F$ is positive exactly when there is a real split $\mvec=\sum_j(\mathbf a_j+\mathbf b_j)$ with $|\mathbf a_j|=r_js_j$ and $|\mathbf b_j|=r^\ast_jt_j$ at which every node of every table has
\[
(\mathbf a_{j,v})_S>s_j\,d_{j,v}(S)\quad\text{and}\quad(\mathbf b_{j,v})_S>t_j\,d^\ast_{j,v}(S)\qquad(\varnothing\neq S\subsetneq\cE_{j,v}\sqcup G_{j,v}),
\]
$d_{j,v}$ and $d^\ast_{j,v}$ the counts of $\mathcal H_{j,v}$. This is a system of $\sum_j\sum_v2(2^{|\cE_{j,v}|+|G_{j,v}|}-2)$ strict linear inequalities in $(2N-1)(n-1)$ unknowns, a number linear in the number of nodes when the nodes have bounded size, against the $2^n-2$ walls of $F$; the native loads of $F$ solve it.
\end{corollary}

\section{What a node costs}

% Reference 20: lem:basecount; source line 6911
\begin{lemma}[Bases outnumber marks]\label{lem:basecount}
A table on $n\ge2$ marks that is one molecule has at least $n$ bases, and exactly $n$ only when its rank is $1$ or $n-1$, that is, as a list, when it is $\binom{\cE}1$ or $\binom{\cE}{n-1}$.
\end{lemma}

% Reference 20: cor:nodecost; source line 6920
\begin{corollary}[The cost of a node]\label{cor:nodecost}
For a node $\mathcal H_v$ that is one molecule on $n_v=|\cE_v|+|G_v|$ marks with $|\mathcal H_v|$ bases, the share side of its period at each of its two full leans (Proposition~\ref{prop:shareform}) is a fibre integral of dimension $|\mathcal H_v|-n_v\ge0$. The dimension is zero exactly when $\mathcal H_v$ has the bases of $\binom{\cE_v\sqcup G_v}{n_v-1}$ or of $\binom{\cE_v\sqcup G_v}1$, with any weights: the series and the parallel nodes up to their weights, whose periods at the full leans are stacked Gammas times monomials in the weights (Proposition~\ref{prop:simplicial}; Proposition~\ref{prop:spleans}, next); a prime node on at least four marks has $|\mathcal H_v|>n_v$.
\end{corollary}

% Reference 20: prop:spleans; source line 6927 (the closed forms alone).
% Proved in the last part: prop:spleans.
\begin{proposition}[Series and parallel nodes at their full leans]\label{prop:spleans}
On a set $\mathcal X$ of $q+1$ marks let $\binom{\mathcal X}q$ carry the weight $\theta_e$ on its basis $\mathcal X\setminus e$ and $\binom{\mathcal X}1$ the weight $\theta_e$ on its basis $\{e\}$, the series and the parallel node with weights (Definition~\ref{def:spnodes}). Then
\begin{gather*}
P_{\binom{\mathcal X}q}(\mathbf a;s,0)=\Gst{q+1}1{(s-a_e)_e}{s}\prod_e\theta_e^{a_e-s},\qquad P_{\binom{\mathcal X}q}(\mathbf b;0,t)=\Gst{q+1}1{(b_e)_e}{t}\prod_e\theta_e^{-b_e},\\
P_{\binom{\mathcal X}1}(\mathbf a;s,0)=\Gst{q+1}1{(a_e)_e}{s}\prod_e\theta_e^{-a_e},\qquad P_{\binom{\mathcal X}1}(\mathbf b;0,t)=\Gst{q+1}1{(t-b_e)_e}{t}\prod_e\theta_e^{b_e-t},
\end{gather*}
with $|\mathbf a|=qs$, $|\mathbf b|=t$ at the series node and $|\mathbf a|=s$, $|\mathbf b|=qt$ at the parallel node; each converges exactly when every Gamma above has an argument of positive real part. These are Proposition~\ref{prop:simplicial} at the configurations $(\mathbf e_x)_{x\in\mathcal X}$, of the letters $K$ of $\binom{\mathcal X}1$ and $\Psi$ of $\binom{\mathcal X}q$, and $(\one-\mathbf e_x)_{x\in\mathcal X}$, of the letters $K$ of $\binom{\mathcal X}q$ and $\Psi$ of $\binom{\mathcal X}1$.
\end{proposition}

% Reference 20: cor:spfulllean; source line 6939 (unit weights).
% Proved in the last part: cor:spfulllean.
\begin{corollary}[Series-parallel tables at their full leans]\label{cor:spfulllean}
In Definition~\ref{def:branches} let every node of $\mathfrak T$ be a series or a parallel node (Definition~\ref{def:spnodes}), the node $v$ on $\cE_v\sqcup G_v$ carrying the weight $\theta_{v,x}$ on its basis that omits $x$ if it is a series node and on its basis $\{x\}$ if it is a parallel node, and write $(\mathbf a_v)_x$ and $(\mathbf b_v)_x$ for the node loads at $x\in\cE_v\sqcup G_v$ (Definition~\ref{def:markerloads}). Then
\begin{align*}
P_{\cB}(\mathbf a;s,0)&=\frac1{\Gamma(s)}\prod_{e\text{ at series}}\Gamma(s-a_e)\prod_{e\text{ at parallel}}\Gamma(a_e)\prod_{\mathfrak e\text{ series--parallel}}\frac{\Gamma(s-k_{\mathfrak e})}{\Gamma(k_{\mathfrak e})}\\
&\qquad\times\prod_{v\text{ series}}\prod_x\theta_{v,x}^{(\mathbf a_v)_x-s}\prod_{v\text{ parallel}}\prod_x\theta_{v,x}^{-(\mathbf a_v)_x},\\
P_{\cB}(\mathbf b;0,t)&=\frac1{\Gamma(t)}\prod_{e\text{ at series}}\Gamma(b_e)\prod_{e\text{ at parallel}}\Gamma(t-b_e)\prod_{\mathfrak e\text{ series--parallel}}\frac{\Gamma(l_{\mathfrak e})}{\Gamma(t-l_{\mathfrak e})}\\
&\qquad\times\prod_{v\text{ series}}\prod_x\theta_{v,x}^{-(\mathbf b_v)_x}\prod_{v\text{ parallel}}\prod_x\theta_{v,x}^{(\mathbf b_v)_x-t},
\end{align*}
the products over the marks by the kind of their node, over the edges joining a series node to a parallel node, with $k_{\mathfrak e}$, $l_{\mathfrak e}$ the marker loads at the series end, and over the nodes and their marks and markers $x$; an edge joining two nodes of one kind contributes no Gamma. Each of these periods is finite exactly when every Gamma above its line has an argument with positive real part.
\end{corollary}

The period of a series-parallel table at a full lean is written down from its tree in time linear in the size of the tree, one Gamma for each mark, two for each change of kind along an edge and a monomial in the weights of each node, while the table may have a number of bases exponential in $n$. The atom formula of the same period (Corollary~\ref{cor:atomwall}) has one share for each basis.

% Reference 20: cor:primelean; source line 6000 (one table at $(0,t)$ along its canonical hinge tree).
% Proved in the last part: cor:primelean.
\begin{corollary}[Periods along a hinge tree read by one hinge value]\label{cor:primelean}
Let $F=((\cB_j,\theta_j);s_j,t_j)_{j=1}^N$ have real data on its balance wall with every proper wall positive, let its tables have a common hinge tree $\mathfrak T$ rooted at $o$, with the node families $F_{\mathcal H_v}$ (Definition~\ref{def:hingetree}), and let the active letters read every subtree cell by one hinge value: for every $v\neq o$ there is a function $\gamma_v$ with $\gamma^K_{\mathcal A_{v,j}}=\gamma_v$ whenever $s_j\neq0$ and $\gamma^\Psi_{\mathcal A_{v,j}}=\gamma_v$ whenever $t_j\neq0$. With $\sigma$ the total strength (Definition~\ref{def:family}),
\[
P_F(\mvec;\mathbf s,\mathbf t)=\frac{\prod_vP_{F_{\mathcal H_v}}(\mvec_v;\mathbf s,\mathbf t)}{\prod_{\text{edges }\mathfrak e}\Gst21{\kappa_{\mathfrak e},\ 2\sigma-\kappa_{\mathfrak e}}{2\sigma}},
\]
each node family on its balance wall with every proper wall positive, its masses $\mvec_v$ those of its marks and, on the marker at an end of the edge $\mathfrak e$, the wall of the marks across $\mathfrak e$: $\kappa_{\mathfrak e}$ on one side and $2\sigma-\kappa_{\mathfrak e}$ on the other. For one table at a full lean, $(0,t)$ or $(s,0)$, one letter is active, the hypothesis holds on every hinge tree, and this is Theorem~\ref{thm:leantree} at the loads $\mvec$. With the canonical hinge tree of a connected table whose weights factor across each bunch of width one at which it is decomposed, a list or a determinantal table for instance (Theorem~\ref{thm:primetree}), a series or a parallel node, with any weights on its bases, contributes at a full lean a stacked Gamma times a monomial in its weights (Proposition~\ref{prop:spleans}), so the period at a full lean of such a table is a product of stacked Gammas, monomials in the weights and the periods of its prime nodes. In particular connected tables $\mathcal A_1,\dots,\mathcal A_k$ at $(0,t)$, each on its balance wall with every proper wall positive, glued along any tree, every gluing joining a mark of mass $\kappa_{\mathfrak e}$ of one to a mark of mass $t-\kappa_{\mathfrak e}$ of another, have the period
\[
\frac{\prod_vP_{\mathcal A_v}(\mvec_v;0,t)}{\prod_{\mathfrak e}\Gst21{\kappa_{\mathfrak e},\ t-\kappa_{\mathfrak e}}{t}},
\]
whatever the shape of the tree; for $k$ copies of one table glued at the masses $(\kappa,t-\kappa)$ this is $\Gst21{\kappa,\ t-\kappa}{t}\big(P_{\mathcal A}(\mvec;0,t)/\Gst21{\kappa,\ t-\kappa}{t}\big)^k$. On the chain $F^{\rhd k}$ of Theorem~\ref{thm:chainoperator}, with every proper wall positive and the active letters reading each copy by one hinge value, the product is $P_F^k/\Gst21{m_{g_{\mathrm{in}}},\ m_{g_{\mathrm{out}}}}{2\sigma}^{k-1}$, the display of that theorem.
\end{corollary}

Where the active letters read every cell through one hinge value, as at a full lean of one table, where only one letter sees a hinge, $\Psi$ at $(0,t)$ and $K$ at $(s,0)$, every cell is pinned, every hinge integral is the Beta corner of a linear letter (Proposition~\ref{prop:scalebreta} with one coefficient), and the period of the whole construction is the product of the periods of its nodes over the Beta corners of its hinges. Nothing is integrated against anything.

\section{Blocks aggregate}\label{sec:blockaggregate}

% Reference 20: thm:blockagg; source line 6017 (blocks read through their totals).
% Proved in the last part: thm:blockagg.
\begin{theorem}[Blocks read by one hinge value aggregate exactly]\label{thm:blockagg}
Let $F$ be a family substituted with blocks (Corollary~\ref{cor:blocksclosures}(i)) into the tables $\mathcal H_j$ of a core family $F_{\mathcal H}=((\mathcal H_j,\theta^{\mathcal H_j});s_j,t_j)_j$, with unit weights in the blocks, the block $B_i$ of $\ell_i$ marks. The active letters read $B_i$ \emph{through its total} when every table with $s_j\neq0$ has a parallel block and every table with $t_j\neq0$ a series block at $B_i$, and through $e_{\ell_i}/e_{\ell_i-1}=\big(\sum_{x\in B_i}\alpha_x^{-1}\big)^{-1}$ when every table with $s_j\neq0$ has a series block and every table with $t_j\neq0$ a parallel block (Corollary~\ref{cor:blocksclosures}(i)). Let $U\sim\Dir_{\mvec}$ (Definition~\ref{def:dirlaw}), with the block totals $|U_{B_i}|$.
\begin{enumerate}[label=(\roman*)]
\item If the active letters read every block through its total, the joint law of all the active letters, $\big(K_j(U)\big)_{s_j\neq0}$ and $\big(\Psi_j(U)\big)_{t_j\neq0}$, with the totals $(|U_{B_i}|)_i$ themselves, is that of $\big(K_{\mathcal H_j}(U')\big)_{s_j\neq0}$, $\big(\Psi_{\mathcal H_j}(U')\big)_{t_j\neq0}$ and $U'$ for $U'\sim\Dir_{(m_{B_i})_i}$: every nonnegative function of these variables, and every absolutely integrable one, has the same expectation on both sides.
\item Let the active letters read $B_i$ through its total for $i\in I$ and through $e_{\ell_i}/e_{\ell_i-1}$ for $i\notin I$, and let $\sigma$ be the total strength (Definition~\ref{def:family}). The wall of a block is $\kappa_{B_i}=m_{B_i}$ for $i\in I$ and $\kappa_{B_i}=m_{B_i}-(\ell_i-1)\,2\sigma$ for $i\notin I$, and on the balance wall, in $[0,\infty]$,
\[
P_F(\mvec;\mathbf s,\mathbf t)=\prod_{i\in I}\Gst{\ell_i}1{\mvec|_{B_i}}{m_{B_i}}\prod_{i\notin I}\Gst{\ell_i}1{(2\sigma-m_x)_{x\in B_i}}{2\sigma-\kappa_{B_i}}\ P_{F_{\mathcal H}}\big((\kappa_{B_i})_i;\mathbf s,\mathbf t\big),
\]
a stacked Gamma read as $+\infty$ when one of its upper arguments is not positive; both sides are finite together when these are positive.
\item If the active letters read every block through its total, then at fixed pooled masses $m_{B_i}$, dividing each among any number of marks leaves the expectation $\oE_{U\sim\Dir_{\mvec}}[W(U)]$ of the family word unchanged and multiplies the period by the stacked Gammas $\Gst{|B_i|}1{\mvec|_{B_i}}{m_{B_i}}$ alone: lengthening the blocks creates no new integral.
\end{enumerate}
A block read through its total by some active letters and through $e_\ell/e_{\ell-1}$ by others keeps the block as one mark of Theorem~\ref{thm:corereads}, with its kernel and its hinge values.
\end{theorem}

\section{The network theorem}\label{sec:network}

The lean convolution of a family of substitutions, the periods of its letters at their full leans written along the hinge trees of its tables, is one fibre integral whose integrand is a product of node data at affine loads. Read as one integral over all the loads, it has dimension $(2N-1)(n-1)$ for every family whose tables carry hinge trees, and for series-parallel tables its integrand is a written-down stacked Gamma times an exponential of the loads. Integrated from the leaves up, when the tables share one hinge tree, the same integral is computed node by node with one rule at every node, prime or not (Theorem~\ref{thm:holdnode}): a node enters through the periods of its tables at their full leans and its cell families through the Mellin transforms of their laws; for one table this is the hinge tree procedure of Chapter~\ref{ch:hinged}.

% Reference 20: thm:network; source line 6957 (one table).
% Proved in the last part: thm:network.
\begin{theorem}[The network]\label{thm:network}
Let $F=((\cB_j,\theta_j);s_j,t_j)_{j=1}^N$ be a family on $\cE$ with exponents $s_j,t_j>0$ whose table $j$ carries a hinge tree $\mathfrak T_j$ as in Definition~\ref{def:branches}, its nodes $\mathcal H_{j,v}$ each one molecule, and let $\mvec$ be on the balance wall of $F$ with every proper wall positive. Then every table is one molecule, and on every admissible point (Definition~\ref{def:fullleanfns}) and for every twist $\mathbf x$ (Definition~\ref{def:twist}),
\[
P_{F[\mathbf x]}(\mvec;\mathbf s,\mathbf t)=\int\ \prod_{j=1}^N\frac{\prod_{v\in\mathfrak T_j}P_{\mathcal H_{j,v}}(\mathbf a_{j,v};s_j,0)\,P_{\mathcal H_{j,v}}(\mathbf b_{j,v};0,t_j)}{\prod_{\mathfrak e=\{v,w\}\in\mathfrak T_j}\Gst21{k_{j;v,\mathfrak e},\ k_{j;w,\mathfrak e}}{s_j}\Gst21{l_{j;v,\mathfrak e},\ l_{j;w,\mathfrak e}}{t_j}}\ (\mathbf x^K_j)^{-\mathbf a_j}(\mathbf x^\Psi_j)^{-\mathbf b_j},
\]
the fibre integral of Theorem~\ref{thm:leanconv} over the loads $(\mathbf a_j,\mathbf b_j)_j$ of the marks, of dimension $(2N-1)(n-1)$, with the node loads $\mathbf a_{j,v},\mathbf b_{j,v}$ and the marker loads $k_{j;v,\mathfrak e},l_{j;v,\mathfrak e}$ of Definition~\ref{def:markerloads} read in table $j$ at its loads $(\mathbf a_j,\mathbf b_j)$ and its exponents $(s_j,t_j)$. At every admissible point the periods of every node at both full leans are finite and every marker load of table $j$ lies in $(0,s_j)$, respectively $(0,t_j)$; the native loads of $F$ are an admissible point. The nodes enter only through their periods at the two full leans, as functions of their loads, the same functions for all exponents, all masses and all twists.
\end{theorem}

% Reference 20: cor:spnetwork; source line 6968 (one table, unit weights).
% Proved in the last part: cor:spnetwork.
\begin{corollary}[Series-parallel periods]\label{cor:spnetwork}
In Theorem~\ref{thm:network} let every node of every table be a series or a parallel node with the weights of Corollary~\ref{cor:spfulllean}, $\theta_{j,v,x}$ at the node $v$ of $\mathfrak T_j$. Then
\begin{multline*}
\prod_j\Gamma(s_j)\Gamma(t_j)\,P_{F[\mathbf x]}(\mvec;\mathbf s,\mathbf t)=\int\ \prod_j\Bigg[\prod_{e\text{ at series}}\Gamma(s_j-a_{j,e})\Gamma(b_{j,e})\prod_{e\text{ at parallel}}\Gamma(a_{j,e})\Gamma(t_j-b_{j,e})\\
\times\prod_{\mathfrak e\text{ series--parallel}}\frac{\Gamma(s_j-k_{j,\mathfrak e})\,\Gamma(l_{j,\mathfrak e})}{\Gamma(k_{j,\mathfrak e})\,\Gamma(t_j-l_{j,\mathfrak e})}\prod_{v\text{ series}}\prod_x\theta_{j,v,x}^{(\mathbf a_{j,v})_x-(\mathbf b_{j,v})_x-s_j}\\
\times\prod_{v\text{ parallel}}\prod_x\theta_{j,v,x}^{(\mathbf b_{j,v})_x-(\mathbf a_{j,v})_x-t_j}\,(\mathbf x^K_j)^{-\mathbf a_j}(\mathbf x^\Psi_j)^{-\mathbf b_j}\Bigg],
\end{multline*}
the kinds of the nodes and the marker loads $k_{j,\mathfrak e}$, $l_{j,\mathfrak e}$ at the series end read in $\mathfrak T_j$, the fibre integral of Theorem~\ref{thm:network}, of dimension $(2N-1)(n-1)$, on any point at which every Gamma above the line has an argument of positive real part. Every period of a family of series-parallel tables with any weights on their nodes, at all exponents, at all masses with positive walls and at every twist, is a fibre integral of dimension $(2N-1)(n-1)$ of an explicit stacked Gamma times an exponential of the loads, whatever the numbers of bases, and at a full lean of one table no integral is left (Corollary~\ref{cor:spfulllean}).
\end{corollary}

% Reference 20: prop:poolpair; source line 6980 (unit weights).
% Proved in the last part: prop:poolpair.
\begin{proposition}[Pooling a pair]\label{prop:poolpair}
In Corollary~\ref{cor:spnetwork}, for one table twisted in its cobasis letter by $\mathbf x$ (Definition~\ref{def:twist}), let $e,f$ be two marks of one series node $v$ whose twists times the squares of their weights agree, $x_e\theta_{v,e}^2=x_f\theta_{v,f}^2$, and $c_e=s-m_e$, $c_f=s-m_f$. Their loads enter the integrand through $\Gamma(c_e+b_e)\Gamma(b_e)\Gamma(c_f+b_f)\Gamma(b_f)$ and otherwise through $b_{\{e,f\}}=b_e+b_f$ alone, and at fixed $b_{\{e,f\}}$
\begin{multline*}
\frac1{2\pi i}\int\Gamma(c_e+b_e)\Gamma(b_e)\Gamma\big(c_f+b_{\{e,f\}}-b_e\big)\Gamma\big(b_{\{e,f\}}-b_e\big)\,db_e\\
=\Gst41{c_e+c_f+b_{\{e,f\}},\ c_e+b_{\{e,f\}},\ c_f+b_{\{e,f\}},\ b_{\{e,f\}}}{c_e+c_f+2b_{\{e,f\}}}.
\end{multline*}
For two marks of one parallel node $v$ with $x_e\theta_{v,e}^{-2}=x_f\theta_{v,f}^{-2}$, at fixed $b_{\{e,f\}}$, the pair gives
\[
\Gst41{m_e+m_f-b_{\{e,f\}},\ m_e+t-b_{\{e,f\}},\ m_f+t-b_{\{e,f\}},\ 2t-b_{\{e,f\}}}{m_e+m_f+2t-2b_{\{e,f\}}}.
\]
Each pooled pair lowers the dimension of the fibre by one, and disjoint pairs can be pooled at once.
\end{proposition}

The pooled pair is a series pair read as one cell: its factor is $\Gst21{c_e+b_{\{e,f\}},\ c_f+b_{\{e,f\}}}{c_e+c_f+2b_{\{e,f\}}}\Gamma(c_e+c_f+b_{\{e,f\}})\Gamma(b_{\{e,f\}})$, the Beta corner being the Mellin transform of the law of the pair at $b_{\{e,f\}}$ and the two Gammas those of its hinge in the periods of the node at its full leans (Theorem~\ref{thm:holdnode}). A series bunch of three or more marks at one node can be pooled pair by pair only once: the pooled factor of a pair is no longer of the form $\Gamma(c+b)\Gamma(b)$, and the class is a circuit of its marks (Corollary~\ref{cor:cycle}).

\section{Nodes in hold form}\label{sec:holdnodes}

% Reference 20: thm:holdnode; source line 6999 (one table).
% Proved in the last part: thm:holdnode.
\begin{theorem}[Nodes in hold form]\label{thm:holdnode}
Let $F$ be a substituted family (Definition~\ref{def:substitution}) whose tables are each one molecule, with exponents $s_j,t_j>0$, table $j$ being $\mathcal H_j\langle\mathcal A_{1j},\dots,\mathcal A_{kj}\rangle$ with the cores $\mathcal H_j$ on $\cE_{\mathcal H}\sqcup G_{\mathcal H}$ and the cells $\mathcal A_{ij}$ on $B_i\sqcup\{g_i\}$ at the markers $g_1,\dots,g_k$ (Definition~\ref{def:node}), and let $F_i=((\mathcal A_{ij},\theta^{\mathcal A_{ij}});s_j,t_j)_j$ be the cell family at $g_i$ (Definition~\ref{def:hinged}). For loads $(\mathbf a_j,\mathbf b_j)_j$ on the marks and markers of the cores write $\mathbf z_i=(a_{j,g_i},b_{j,g_i})_j$ for their values at $g_i$, and in this statement write $\int\boldsymbol\varrho_{F_i}^{-\mathbf z_i}d\phi_{F_i}$ for the Mellin transform of the law of $F_i$ of Lemma~\ref{lem:hingesplit} at $\mathbf z_i$.
\begin{enumerate}[label=(\roman*)]
\item \emph{Root.} If $G_{\mathcal H}=\{g_1,\dots,g_k\}$ and $F$ is on its balance wall with every proper wall positive, let $\mathcal F$ be the fibre of the loads of the cores with the pins $\sum_j(\mathbf a_j+\mathbf b_j)=(\mvec|_{\cE_{\mathcal H}},\kappa_{F_1},\dots,\kappa_{F_k})$, $|\mathbf a_j|=r_{\mathcal H_j}s_j$ ($j\ge2$) and $|\mathbf b_j|=r^\ast_{\mathcal H_j}t_j$. On the node loads at the cores (Definition~\ref{def:markerloads}, in every table) of any admissible point of $F$ and for twists $\mathbf x$ of $F$ with every coordinate $1$ off $\cE_{\mathcal H}$,
\[
P_{F[\mathbf x]}(\mvec;\mathbf s,\mathbf t)=\int_{\mathcal F}\prod_{j=1}^NP_{\mathcal H_j}(\mathbf a_j;s_j,0)\,P_{\mathcal H_j}(\mathbf b_j;0,t_j)\,(\mathbf x^K_j)^{-\mathbf a_j|_{\cE_{\mathcal H}}}(\mathbf x^\Psi_j)^{-\mathbf b_j|_{\cE_{\mathcal H}}}\prod_{i=1}^k\int\boldsymbol\varrho_{F_i}^{-\mathbf z_i}\,d\phi_{F_i},
\]
of dimension $(2N-1)(|\cE_{\mathcal H}|+k-1)$, every $\int\boldsymbol\varrho_{F_i}^{-\mathbf z_i}d\phi_{F_i}$ finite at the real point.
\item \emph{Hinged core.} If $G_{\mathcal H}=\{g_0,g_1,\dots,g_k\}$, let $F'$ be the cell family with the hinge $g_0$ whose table $j$ is $\mathcal H_j\langle\mathcal A_{1j},\dots,\mathcal A_{kj}\rangle$, with the hinge wall $\kappa_{F'}$, and for $\mathbf z=(z_{j,K},z_{j,\Psi})_j$ with $\sum_j(z_{j,K}+z_{j,\Psi})=\kappa_{F'}$ let $\mathcal F_{\mathbf z}$ be the fibre of the loads of the cores with the pins $\sum_j(\mathbf a_j+\mathbf b_j)=\big(\mvec|_{\cE_{\mathcal H}},\kappa_{F_1},\dots,\kappa_{F_k},\sum_j(s_j+t_j)-\kappa_{F'}\big)$, $|\mathbf a_j|=r_{\mathcal H_j}s_j$ ($j\ge2$), $|\mathbf b_j|=r^\ast_{\mathcal H_j}t_j$, $a_{j,g_0}=s_j-z_{j,K}$ ($j\ge2$) and $b_{j,g_0}=t_j-z_{j,\Psi}$. For $\mathbf z$ whose real part is $(s_j-c^K_{j,g_0},t_j-c^\Psi_{j,g_0})_j$, $\mathbf c$ the node loads at the cores of an admissible point of the family of the tables of $F'$ with the mass $\sum_j(s_j+t_j)-\kappa_{F'}$ on $g_0$,
\begin{multline*}
\prod_{j=1}^N\Gst21{s_j-z_{j,K},\ z_{j,K}}{s_j}\Gst21{t_j-z_{j,\Psi},\ z_{j,\Psi}}{t_j}\int\boldsymbol\varrho_{F'}^{-\mathbf z}\,d\phi_{F'}\\
=\int_{\mathcal F_{\mathbf z}}\prod_{j=1}^NP_{\mathcal H_j}(\mathbf a_j;s_j,0)\,P_{\mathcal H_j}(\mathbf b_j;0,t_j)\prod_{i=1}^k\int\boldsymbol\varrho_{F_i}^{-\mathbf z_i}\,d\phi_{F_i},
\end{multline*}
of dimension $(2N-1)(|\cE_{\mathcal H}|+k-1)$.
\end{enumerate}
The cores and the cells are molecules by Lemma~\ref{lem:twosummol}.
\end{theorem}

% Reference 20: cor:holdrules; source line 7024 (one table).
% Proved in the last part: cor:holdrules.
\begin{corollary}[The glue and the series and parallel rules in hold form]\label{cor:holdrules}
On the lines of Theorem~\ref{thm:holdnode}:
\begin{enumerate}[label=(\roman*)]
\item For one table with $\mathcal H=\binom{\{g_1,g_2\}}1$ and the cells $\mathcal A$ and $\mathcal A'$, Theorem~\ref{thm:holdnode}(i) is the integral of Theorem~\ref{thm:glue}(iv), with the same integrand; for a family with the core $\binom{\{g_1,g_2\}}1$ in every table it is the integral of Theorem~\ref{thm:twosides} over the scale shares.
\item For one table and a series node $\binom Gq$, Theorem~\ref{thm:holdnode}(ii) reads
\[
\Gamma(z)\Gamma(z+s-\kappa_{\mathcal A})\int\varrho_{\mathcal A}^{-z}d\phi_{\mathcal A}=\int_{z_1+\cdots+z_q=z}\ \prod_{i=1}^q\Gamma(z_i)\Gamma(z_i+s-\kappa_{\mathcal A_i})\int\varrho_{\mathcal A_i}^{-z_i}d\phi_{\mathcal A_i},
\]
the Mellin transform of $f_{\mathcal A}$ as the Mellin convolution of those of the $f_{\mathcal A_i}$ (Proposition~\ref{prop:faces}(iii)): the Mellin transform, on these lines, of the product rule $f_{\mathcal A}=\prod_if_{\mathcal A_i}$ of Theorem~\ref{thm:seriesparallel}(i).
\item For one table and a parallel node $\binom G1$ it reads, with $z'=\kappa_{\mathcal A}-z$ and $z'_i=\kappa_{\mathcal A_i}-z_i$,
\[
\Gamma(z')\Gamma(z'+t-\kappa_{\mathcal A})\int\varrho_{\mathcal A}^{z'-\kappa_{\mathcal A}}d\phi_{\mathcal A}=\int_{z'_1+\cdots+z'_q=z'}\ \prod_{i=1}^q\Gamma(z'_i)\Gamma(z'_i+t-\kappa_{\mathcal A_i})\int\varrho_{\mathcal A_i}^{z'_i-\kappa_{\mathcal A_i}}d\phi_{\mathcal A_i},
\]
the Mellin convolution of the transforms of the $\check f_{\mathcal A_i}$: the Mellin transform, on these lines, of the product rule $\check f_{\mathcal A}=\prod_i\check f_{\mathcal A_i}$ of Theorem~\ref{thm:seriesparallel}(ii).
\end{enumerate}
\end{corollary}

So the three rules of the hinge tree procedure, the glue and the series and parallel products (Definition~\ref{def:treeproc}), are one rule read at three nodes, on the lines that the network certifies, the rules themselves holding on their own wider ranges; and the same rule applies at a prime node, where the procedure integrates the twisted core over its simplex for every value of the twists (Definition~\ref{def:treeproc}(4)). In Theorem~\ref{thm:holdnode} the prime node enters through the periods of its tables at their full leans, functions of its loads that do not depend on the twists, on the exponents or on the masses, and the cell families enter through the Mellin transforms of their laws on imaginary translates, for one table the representation (b) of Definition~\ref{def:treeproc}; Theorem~\ref{thm:holdnode}(i) is Theorem~\ref{thm:node}(i) with the twisted core replaced by its lean convolution and the integration against the laws carried out through Mellin transforms.

\section{The procedure and its cost}

% Reference 20: def:networkproc; source line 7050 (one table).
% Proved in the last part: def:networkproc.
\begin{definition}[The network procedure]\label{def:networkproc}
\emph{Given}: a hinge tree of every table of a family (Definition~\ref{def:hingetree}) with the weights and the marks of its nodes; exponents $s_j,t_j>0$ and masses $\mvec$ on the balance wall, or several of each; optionally twists $\mathbf x$, or product weights (Proposition~\ref{prop:weightstwists} below). \emph{Reusable data}: for every orbit of nodes, over all the tables, under bijections keeping the table and its weights, the Gramma $\sG_{\mathbf a;(\one_T)_T}(\boldsymbol\theta)$ of its table, a function of the loads alone (Corollary~\ref{cor:retaintable}), kept as (a) the closed form at series and parallel nodes (Proposition~\ref{prop:spleans}); (b) the share side, a fibre integral of dimension $|\mathcal H_v|-n_v$ (Corollary~\ref{cor:nodecost}); (c) for an identically self-dual node, the reduction to the deletion of one mark (Corollaries~\ref{cor:piecesmarks}(iii) and~\ref{cor:selfdualtwist}); (d) values on the imaginary translates in use. \emph{Operations}:
\begin{enumerate}[label=(\arabic*)]
\item \emph{Contour.} Find a split satisfying the node inequalities of Corollary~\ref{cor:nodeconv}, a linear feasibility problem with $\sum_j\sum_v2(2^{|\cE_{j,v}|+|G_{j,v}|}-2)$ strict inequalities in $(2N-1)(n-1)$ unknowns; the native loads solve it, and infeasibility proves divergence.
\item \emph{Assemble.} The node loads of every table by Definition~\ref{def:markerloads}, the integrand of Theorem~\ref{thm:network}, at series-parallel nodes that of Corollary~\ref{cor:spnetwork}, with $\prod_j(\mathbf x^K_j)^{-\mathbf a_j}(\mathbf x^\Psi_j)^{-\mathbf b_j}$ for twists.
\item \emph{Pool}, for one table, disjoint pairs of marks of one series or parallel node whose twists and weights satisfy the condition of Proposition~\ref{prop:poolpair}, by that proposition.
\item \emph{Integrate} the remaining loads on the contour; in the imaginary parts the integrand is analytic on a strip about the contour, the admissible points forming an open set, and decays exponentially there (Theorem~\ref{thm:leanconv}(iii)), so the trapezoidal rule with step $\varsigma$ has an error of the order $e^{-2\pi d/\varsigma}$, $d$ the half-width of the strip; or close contours on lattices (Theorem~\ref{thm:closing}).
\item \emph{Recursive Mellin evaluation}, an alternative to (2) to (4) when the hinge trees of the tables are one rooted tree with the same marks $\cE_v$ at every vertex, so that at every vertex $v$ other than the root the subtree cells of the tables (Definition~\ref{def:hingetree}) form a cell family, the \emph{subtree cell family} of $v$ (Definition~\ref{def:hinged}): from the leaves up, compute the Mellin transform of the law of each orbit of subtree cell families on an imaginary translate by Theorem~\ref{thm:holdnode}(ii), one fibre integral of dimension $(2N-1)(|\cE_v|+|G_v|-2)$ per point, and finish at the root by Theorem~\ref{thm:holdnode}(i); for one table, at series and parallel nodes with repeated subtrees, equal and with equal masses, pass to the transforms, where repetition is a power (Proposition~\ref{prop:work}(ii)).
\end{enumerate}
\emph{Output}: the period, and with the same reusable data the period at other exponents, masses and twists (Proposition~\ref{prop:carry} below).
\end{definition}

% Reference 20: prop:networkcost; source line 7062 (one table).
% Proved in the last part: prop:networkcost.
\begin{proposition}[What the network costs]\label{prop:networkcost}
In Definition~\ref{def:networkproc}:
\begin{enumerate}[label=(\roman*)]
\item The fibre of the loads in (1)--(4) has dimension $(2N-1)(n-1)$, for one table less one for each pooled pair; opening every prime node on the share sides of the periods of its table at the two full leans adds $2(|\mathcal H_{j,v}|-n_{j,v})$ for each, and series and parallel nodes add nothing. The dimension does not depend on the numbers $|\cB_j|$ of bases of the tables, against $\sum_j2|\cB_j|-2N-n+1$ for the atom formula (Corollary~\ref{cor:atomwall}) and $\sum_j|\cB_j|-n+1$ after its even shares close (Theorem~\ref{thm:buildup}(iv)).
\item The reusable data are one function of $n_v$ variables per orbit of nodes, the same for every node of the orbit, in every table, and for every query; a value of the period of a node at a full lean at given loads, once computed, serves every node of the orbit and every query at which those loads occur.
\item In the recursive Mellin evaluation the work is one Mellin transform on an imaginary translate for each orbit of subtree cell families (as in Proposition~\ref{prop:work}(i)), each point costing one fibre integral of dimension $(2N-1)(|\cE_v|+|G_v|-2)$ at its top node, and at the root one of dimension $(2N-1)(|\cE_o|+|G_o|-1)$, whose integrand needs the periods of the tables of the node at their full leans, at the loads of the fibre: closed forms at series and parallel nodes, and at a prime node an integral of its own, of dimension $|\mathcal H_{j,v}|-n_v$ on its share side, unless another representation of Definition~\ref{def:networkproc} is cheaper.
\end{enumerate}
\end{proposition}

In the recursive Mellin evaluation equal subtrees at a series node enter as a power of one transform (Theorems~\ref{thm:seriesparallel}(i) and~\ref{thm:glue}(ii)), whose growth in the number of copies Theorem~\ref{thm:seriesrepeat} reads when the subtree is series-balanced.

\section{Changing the query in hold form}\label{sec:holdquestion}

A query about a period is a choice of exponents, of masses, of weights and of the kind of answer asked for. In the hold form each of them enters at one place: exponents and masses move the pins and the affine marker loads, twists and product weights multiply the integrand by an exponential of the loads, and a replaced node, or a general change of the weights of one node, replaces that node's two functions and nothing else. Product weights are read as twists, and the exponents at fixed masses as the Mellin variables of the relative scales of the letters.

% Reference 20: prop:weightstwists; source line 7080 (one table).
% Proved in the last part: prop:weightstwists.
\begin{proposition}[Product weights are twists]\label{prop:weightstwists}
For a family $F$ on $\cE$ and $\boldsymbol\zeta=(\boldsymbol\zeta_j)_j$, $\boldsymbol\zeta_j\in(0,\infty)^{\cE}$, let $F^{\boldsymbol\zeta}$ be $F$ with the weights $\theta_{j,T}\boldsymbol\zeta_j^T$ in table $j$. Then $K_{F^{\boldsymbol\zeta},j}(\alpha)=K_j(\boldsymbol\zeta_j\circ\alpha)$ and $\Psi_{F^{\boldsymbol\zeta},j}(\alpha)=\boldsymbol\zeta_j^{\cE}\Psi_j(\alpha/\boldsymbol\zeta_j)$: $F^{\boldsymbol\zeta}$ has the letters of the twist $F[\mathbf x]$ with $\mathbf x^K_j=\boldsymbol\zeta_j$ and $\mathbf x^\Psi_j=\boldsymbol\zeta_j^{-1}$ (Definition~\ref{def:twist}), its cobasis letters multiplied by the constants $\boldsymbol\zeta_j^{\cE}$, and
\begin{gather*}
P_{F^{\boldsymbol\zeta}}(\mvec;\mathbf s,\mathbf t)=\prod_j\boldsymbol\zeta_j^{-t_j\one}\,P_{F[\mathbf x]}(\mvec;\mathbf s,\mathbf t),\\
P_{\cB_j^{\boldsymbol\zeta_j}}(\mathbf a;s_j,0)=\boldsymbol\zeta_j^{-\mathbf a}P_{\cB_j}(\mathbf a;s_j,0),\qquad P_{\cB_j^{\boldsymbol\zeta_j}}(\mathbf b;0,t_j)=\boldsymbol\zeta_j^{\mathbf b-t_j\one}P_{\cB_j}(\mathbf b;0,t_j),
\end{gather*}
so that the lean convolution of $F^{\boldsymbol\zeta}$ is that of $F$ with the exponential $\prod_j\boldsymbol\zeta_j^{-t_j\one}\boldsymbol\zeta_j^{\mathbf b_j-\mathbf a_j}$ of the loads. With one $\boldsymbol\zeta$ for every table, $P_{F^{\boldsymbol\zeta}}(\mvec;\mathbf s,\mathbf t)=\boldsymbol\zeta^{-\mvec-|\mathbf t|\one}P_{F[\mathbf x']}(\mvec;\mathbf s,\mathbf t)$, $\mathbf x'$ the twist of every cobasis letter by $\boldsymbol\zeta^{-2}$ and $|\mathbf t|=\sum_jt_j$. Weights $\zeta_{j,g}$ on the bases through one mark $g$, in each table $j$, are the twist at $g$ with $x^K_j=\zeta_{j,g}$ and $x^\Psi_j=\zeta_{j,g}^{-1}$, up to the factor $\prod_j\zeta_{j,g}^{-t_j}$: the weight of a hinge of Theorem~\ref{thm:hingeweight}, at any mark.
\end{proposition}

% Reference 20: prop:carry; source line 7092 (one table).
% Proved in the last part: prop:carry.
\begin{proposition}[Carrying a calculation through a change, in hold form]\label{prop:carry}
In Theorem~\ref{thm:network}, with its node data computed:
\begin{enumerate}[label=(\roman*)]
\item a change of the exponents $(s_j,t_j)$ changes the pins $|\mathbf a_j|=r_js_j$, $|\mathbf b_j|=r^\ast_jt_j$ and the marker loads of table $j$, which are affine in $(s_j,t_j)$, and nothing else; the native loads of the new data are a new contour;
\item a change of the masses changes the load $\mathbf a_1=\mvec-\sum_{j\ge2}\mathbf a_j-\sum_j\mathbf b_j$ in the coordinates of Theorem~\ref{thm:leanconv}(ii), and through it the marker loads $k_{1;v,\mathfrak e}$ of table $1$, and in the integrand nothing else; in (i) and (ii) the contour must move when the old one is no longer admissible;
\item a change of twists, or of product weights on the bases of the tables, multiplies the integrand by an exponential of the loads (Proposition~\ref{prop:weightstwists}); product weights on the bases of one node multiply its periods at the two full leans by exponentials of its node loads;
\item replacing a node of a table by another table on the same marks and markers, of the same rank, one molecule with its markers neither loops nor coloops, or changing the weights of a node in any way, replaces the periods of that node at its two full leans and nothing else;
\item in the recursive Mellin evaluation (Definition~\ref{def:networkproc}(5)), a change at one node changes the Mellin transforms of the laws of the subtree cell families on the path from it to the root only, each recomputed by Theorem~\ref{thm:holdnode}(ii) from the new one of the vertex directly below it on the path and the kept ones of the other vertices directly below it, as Proposition~\ref{prop:local} says of laws.
\end{enumerate}
\end{proposition}

The reusable data of Chapter~\ref{ch:hinged}, the laws of the subtrees, depend on the exponents and on the masses below each edge (Theorem~\ref{thm:treeperiod}(iii)); the reusable data of the network do not, and a sweep over exponents, over masses, over twists or over product weights recomputes no node. What does change is the contour, and Corollary~\ref{cor:nodeconv} decides at the nodes whether the new query has an answer.

% Reference 20: prop:leanmoment; source line 7108 (one table, along its line of leans).
% Proved in the last part: prop:leanmoment.
\begin{proposition}[The exponents are moment indices]\label{prop:leanmoment}
Let $F=((\cB_j,\theta_j);s_j,t_j)_{j=1}^N$ be a family on $\cE$, fix real masses $\mvec$, and let $I$ be the set of exponents $(\mathbf s,\mathbf t)\in(0,\infty)^{2N}$ with $\sum_j(r_js_j+r^\ast_jt_j)=m_{\cE}$ at which every proper wall of $(\mvec;\mathbf s,\mathbf t)$ is positive: a relatively open convex subset of that hyperplane, possibly empty. In this statement write $L_1,\dots,L_{2N}$ for the letters $K_j$ and $\Psi_j$ of $F$, $w_i$ for their exponents and $g_i$ for their degrees, and for each $i$ and $\boldsymbol\lambda\in(0,\infty)^{2N-1}$ put
\[
\sG_i(\mvec;\boldsymbol\lambda)=\int_{\mathbb R^{\cE}_{>0}}x^{\mvec-\one}\exp\Big(-L_i(x)-\sum_{i'\neq i}\lambda_{i'}L_{i'}(x)\Big)\,dx,
\]
the Gramma (Definition~\ref{def:gramma}) of the concatenated configuration of all the letters, with the coefficients of $L_i$ and $\lambda_{i'}$ times those of $L_{i'}$. For $(\mathbf s,\mathbf t)\in I$ and every $i$,
\begin{align*}
\Gamma(w_i)\,P_F(\mvec;\mathbf s,\mathbf t)&=g_i\int_{\mathbb R^{\cE}_{>0}}x^{\mvec-\one}e^{-L_i(x)}\prod_{i'\neq i}L_{i'}(x)^{-w_{i'}}\,dx\\
&=\frac{g_i}{\prod_{i'\neq i}\Gamma(w_{i'})}\int_{(0,\infty)^{2N-1}}\prod_{i'\neq i}\lambda_{i'}^{w_{i'}-1}\,\sG_i(\mvec;\boldsymbol\lambda)\,d\boldsymbol\lambda .
\end{align*}
At fixed masses the numbers $\Gamma(w_i)P_F(\mvec;\mathbf s,\mathbf t)/g_i$ on $I$, functions of the exponents of the other letters, are the integrals of $\prod_{i'\neq i}L_{i'}^{-w_{i'}}$ against the one positive measure $x^{\mvec-\one}e^{-L_i(x)}dx$, and $\prod_{i'}\Gamma(w_{i'})\,P_F(\mvec;\mathbf s,\mathbf t)/g_i$ is the Mellin transform at $(w_{i'})_{i'\neq i}$ of the one function $\sG_i(\mvec;\cdot)$ of the relative scales of the other letters to $L_i$, completely monotone in each of them. On $I$ the period $P_F(\mvec;\mathbf s,\mathbf t)$ is log-convex, by Proposition~\ref{prop:gammadraws}(iii) and H\"older's inequality.
\end{proposition}

So the period of a family at fixed masses is one integral over the orthant with any one of its letters closed: the closed letter leaves one positive measure on the orthant, the integrand of its Gramma at the index $\mvec$, and the period is a joint moment of the other letters under it, their exponents with the sign changed being its orders.


\chapter{Finite interfaces and certified enclosures}\label{ch:interfaces}

A cell enters every construction through its hinge wall and its law (Theorem~\ref{thm:retained}), and a cell with several hinges through the joint law of its channels (Theorem~\ref{thm:crossing}, Definition~\ref{def:channels}). A law is not a finite list of numbers. Finitely many numbers are retained for every prescribed polynomial moment, the joint moments of the channels under independent Gamma draws on the marks, up to a degree $D$: they compose by the bilinear map of Theorem~\ref{thm:hingepairs}, a repeated chain cell acts on them by one finite matrix, so that every polynomial Gamma moment of a repetition satisfies a linear recurrence with constant coefficients, and the walls turn a moment of finite degree into a certified enclosure of the period. The period is retained as well over a neighbourhood of the query: the walls give the neighbourhood explicitly, the Taylor coefficients in the weights are periods with walls at least as high, and a finite parameter bank with two explicit errors answers every query inside it. Exact representation size, arithmetic work and accuracy are separate quantities.

\section{The joint law of the channels and its moments}\label{sec:channellaw}

Let $\mathcal A$ be a cell with hinges $G$ on its marks $B$, or cells with the hinges $G$ on the same marks $B$ kept together (Definition~\ref{def:channellaw}), with channel vector $\operatorname{ch}(\mathcal A)$ (Definition~\ref{def:channels}), with positive masses $\mvec$ on $B$; the hinges carry no mass. Its total mass $m_B$ and the product $\prod_{e\in B}\Gamma(m_e)$ are kept as data: at a splicing the total masses add and these products multiply. For a multi-index $\mathbf k$ on the channels, $\operatorname{ch}(\mathcal A)^{\mathbf k}$ is the product of the channels with these multiplicities, a polynomial of the degree $\deg\mathbf k=\sum_ck_c\deg c$, the sum over the channels $c$.

% Reference 20: def:channellaw; source line 7129 (one cell).
% Proved in the last part: def:channellaw.
\begin{definition}[The channel law; moment vectors]\label{def:channellaw}
Cells $\mathcal A_1,\dots,\mathcal A_N$ with the hinges $G$ on the same marks $B$ are \emph{kept together} when their channels are read jointly: $\mathcal A=(\mathcal A_1,\dots,\mathcal A_N)$ has the channel vector $\operatorname{ch}(\mathcal A)=(\operatorname{ch}(\mathcal A_j))_j$, listing the channels of all of them by the tags $(j,\text{letter},p)$, $N2^{|G|+1}$ of them, absent channels counting as zero; one cell is the case $N=1$, and the tables of a family, with $G=\varnothing$, are kept together with their letters as channels. The \emph{channel law} of $\mathcal A$ is the image $\operatorname{ch}(\mathcal A)_\ast\Dir_{\mvec}$ of the Dirichlet law of the shape (Definition~\ref{def:dirlaw}) under the channel vector. The \emph{Gamma draws} on the marks are $\alpha\sim\operatorname{Gamma}_{\mvec}$, independent $\operatorname{Gamma}_{m_e}$ variables (Proposition~\ref{prop:gammadraws}), and the \emph{moments} and the \emph{moment vector of degree $D$} of $\mathcal A$ are
\[
\operatorname{Mom}_{\mathcal A}(\mathbf k)=\oE_{\alpha\sim\operatorname{Gamma}_{\mvec}}\big[\operatorname{ch}(\mathcal A)(\alpha)^{\mathbf k}\big],\qquad \operatorname{Mom}^{\le D}_{\mathcal A}=\big(\operatorname{Mom}_{\mathcal A}(\mathbf k)\big)_{|\mathbf k|\le D},
\]
at most $\binom{q+D}{D}$ numbers, $q$ the number of channels. A \emph{polynomial Gamma moment} of degree at most $D$ of $\mathcal A$ is $\oE_{\alpha\sim\operatorname{Gamma}_{\mvec}}\big[p(\operatorname{ch}(\mathcal A)(\alpha))\big]$ for a polynomial $p$ of degree at most $D$ in the channels, and a \emph{shape moment} is the same with $U\sim\Dir_{\mvec}$ in place of $\alpha$. A cell without marks has constant channels, a point as its law and the products of its channels as its moments, with total mass $0$ and the empty product of Gammas.
\end{definition}

The channels must be kept jointly: the moments of the channels one at a time do not determine a product of different channels exposed to the same room.

% Reference 20: prop:channellaw; source line 7139 (one cell).
% Proved in the last part: prop:channellaw.
\begin{proposition}[The channel law composes]\label{prop:channellaw}
Let $\mathcal A_1$ and $\mathcal A_2$ be cells, or cells kept together, with disjoint marks $B_1$, $B_2$ and masses $\mvec_1$, $\mvec_2$, spliced table by table along pairs of their hinges to $\mathcal A_{12}$ (Theorem~\ref{thm:hingepairs}), both with marks. Let $U_1\sim\Dir_{\mvec_1}$, $U_2\sim\Dir_{\mvec_2}$ and $(V_1,V_2)\sim\Dir_{(m_{B_1},m_{B_2})}$ be independent; replace every channel $c$ of $\mathcal A_1$ by $V_1^{\deg c}c(U_1)$ and every channel $c$ of $\mathcal A_2$ by $V_2^{\deg c}c(U_2)$. The law of the image of the two channel vectors so replaced under the bilinear map of Theorem~\ref{thm:hingepairs}(iii), applied table by table, is the channel law of $\mathcal A_{12}$. For the tables of a family, kept together without hinges, the channels are their letters, and the period of the family is Proposition~\ref{prop:gammadraws}(i) read on the channel law: one Trailmoment of the channel law on its source space (Definition~\ref{def:jointtrail}, Theorem~\ref{thm:jointclosure}). For a cell family $F=((\mathcal A_j,\theta_j);s_j,t_j)_{j=1}^N$ with the one hinge $g$ (Definition~\ref{def:hinged}), its cells kept together, with the coordinates $(x_{j,0},x_{j,1},x^\ast_{j,0},x^\ast_{j,1})$ of the channels $K_{\mathcal A_j\setminus g}$, $K_{\mathcal A_j/g}$, $\Psi_{\mathcal A_j\setminus g}$, $\Psi_{\mathcal A_j/g}$, the law of $F$ is the weighted image
\[
(\boldsymbol\varrho_F)_\ast\phi_F=\Gst{|B|}{1}{\mvec}{m_B}\ (R_F)_\ast\Big[\prod_jx_{j,1}^{-s_j}x_{j,0}^{\ast\,-t_j}\Big(\frac{x_{1,0}}{x_{1,1}}\Big)^{-\kappa_F}\ \operatorname{ch}(\mathcal A)_\ast\Dir_{\mvec}\Big],
\]
with $R_F=\big((x_{j,0}x_{1,1}/(x_{j,1}x_{1,0}))_{j\ge2},(x^\ast_{j,1}x_{1,1}/(x^\ast_{j,0}x_{1,0}))_j\big)$, an equality of positive measures, possibly infinite: the channel law, which knows nothing of the exponents, gives the law of the cell family at all exponents by a weight and an image.
\end{proposition}

% Reference 20: thm:gammalift; source line 7155 (one cell).
% Proved in the last part: thm:gammalift.
\begin{theorem}[Gamma draws close the moments]\label{thm:gammalift}
\begin{enumerate}[label=(\roman*)]
\item \emph{Leaves.} If $\operatorname{ch}(\mathcal A)(x)^{\mathbf k}=\sum_{\mathbf h}c_{\mathbf h}x^{\mathbf h}$, then $\operatorname{Mom}_{\mathcal A}(\mathbf k)=\sum_{\mathbf h}c_{\mathbf h}\prod_e(m_e)_{h_e}$, rational at rational masses and weights.
\item \emph{Assembly.} Let $\operatorname{ch}(\mathcal A_{12})=\operatorname{sp}\big(\operatorname{ch}(\mathcal A_1),\operatorname{ch}(\mathcal A_2)\big)$ with the bilinear map $\operatorname{sp}$ of Theorem~\ref{thm:hingepairs}(iii), applied table by table to cells kept together, and write $\operatorname{sp}(c,c')^{\mathbf k}=\sum_{\mathbf k_1,\mathbf k_2}b_{\mathbf k;\mathbf k_1,\mathbf k_2}c^{\mathbf k_1}c'^{\mathbf k_2}$, a sum over $|\mathbf k_1|=|\mathbf k_2|=|\mathbf k|$. Then
\[
\operatorname{Mom}_{\mathcal A_{12}}(\mathbf k)=\sum_{\mathbf k_1,\mathbf k_2}b_{\mathbf k;\mathbf k_1,\mathbf k_2}\operatorname{Mom}_{\mathcal A_1}(\mathbf k_1)\operatorname{Mom}_{\mathcal A_2}(\mathbf k_2),
\]
and no moment of degree above $|\mathbf k|$ of either cell is needed; a splicing inside one cell is linear and keeps the degree.
\item \emph{Return to the shape.} For a cell with marks, $\oE_{U\sim\Dir_{\mvec}}\big[\operatorname{ch}(\mathcal A)(U)^{\mathbf k}\big]=\operatorname{Mom}_{\mathcal A}(\mathbf k)/(m_B)_{\deg\mathbf k}$: the joint moments of the channel law through degree $D$ are read off $\operatorname{Mom}^{\le D}_{\mathcal A}$, the rank, the sizes and the tags fixing the denominators.
\item \emph{Sufficiency.} Two cells with the same hinges, channel tags and degrees and the same $\operatorname{Mom}^{\le D}$ give every final polynomial Gamma moment of degree at most $D$ in the channels, in any assembly with cells on disjoint marks, the same value, even where an intermediate cell exposes more hinges than the replaced one; with the same total mass as well, they give every final shape moment of degree at most $D$ of the channels of the root, under its Dirichlet law, the same value.
\item \emph{The whole law.} For a cell with marks, the $\operatorname{Mom}^{\le D}_{\mathcal A}$ for all $D$, with $m_B$ and the degrees, determine its channel law, and with it the expectation of every integrable function of the channels (Proposition~\ref{prop:channellaw}).
\end{enumerate}
\end{theorem}

The draws are used for nonnegative moments of the channels. A negative moment of a period integrand under the draws diverges on the balance wall: the scale contributes $\int_0^\infty S^{-1}e^{-S}dS$, the pole of $\Gamma(\kappa_{\cE})$ in Proposition~\ref{prop:gammadraws}. The return to the shape, followed by Proposition~\ref{prop:channellaw} or by the peak closing (Theorem~\ref{thm:peak}(i)) with its stopping rule (Theorem~\ref{thm:certstop}), is the change of representation that the period needs.

% Reference 20: cor:interfacework; source line 7174
\begin{corollary}[The work follows the distinct interfaces]\label{cor:interfacework}
Fix a degree $D$, an assembly that exposes at most $h$ hinges at every cell it keeps, and identifications of its repeated subassemblies keeping their data and hinges. Every requested moment of the root is computed exactly by keeping at most $\binom{2^{h+1}+D}{D}$ numbers for each distinct cell, $\binom{N2^{h+1}+D}{D}$ for $N$ tables kept together. The arithmetic is $O\big(\sum_vE_v(D)\big)$ operations once the local expansions and the leaf moments are built, $E_v(D)$ the number of nonzero coefficients $b_{\mathbf k;\mathbf k_1,\mathbf k_2}$ at a distinct splicing $v$ of two cells, at most the product of the sizes of its three moment vectors, or of two for a splicing inside one cell. Equal copies share one moment vector and remain independent draws: their moments are multiplied as in Theorem~\ref{thm:gammalift}(ii), never identified.
\end{corollary}

This counts operations. Building the local expansions, finding an assembly of small width, the bit lengths of exact coefficients and the degree a moment needs are separate parts of the cost; finding the cheapest assembly is a search, as for constructions (Section~\ref{sec:constructionprice}).

\section{Repetition is a finite matrix}\label{sec:transfer2}

% Reference 20: thm:transferpower; source line 7185
\begin{theorem}[A finite matrix for every polynomial Gamma moment]\label{thm:transferpower}
Let a chain cell $\mathcal U$ (Proposition~\ref{prop:chain}, here with several hinges) be a cell with two disjoint sets of hinges $G_{\mathrm{in}}$ and $G_{\mathrm{out}}$, each in a fixed bijection with the hinges of a cell $\mathcal A_0$, and let $\mathcal A_{l+1}$ be $\mathcal U$, with fresh marks and the same masses and weights at every use, spliced onto $\mathcal A_l$ by pairing $G_{\mathrm{in}}$ with the hinges of $\mathcal A_l$; so $\mathcal A_{l+1}$ exposes $G_{\mathrm{out}}$, named through the bijection as the hinges of $\mathcal A_l$, with the same channel tags, and its channel vector is the image of $\operatorname{ch}(\mathcal A_l)$ and that of the copy of $\mathcal U$ under the map $\operatorname{sp}$. Keep one list of channel tags for every $l$, absent channels being zero. For every $D$ there is a matrix $\mathbf T_D$, independent of $l$ and computed from $\mathcal U$ by finite expansion and Theorem~\ref{thm:gammalift}(i), with
\[
\operatorname{Mom}^{\le D}_{\mathcal A_l}=\mathbf T_D^{\,l}\,\operatorname{Mom}^{\le D}_{\mathcal A_0}.
\]
Every fixed polynomial Gamma moment of degree at most $D$ of $\mathcal A_l$, a fixed terminal cell spliced on it included, is $\mathbf e^{\mathsf T}\mathbf T_D^{\,l}\operatorname{Mom}^{\le D}_{\mathcal A_0}$ for a fixed row $\mathbf e^{\mathsf T}$; so its generating series in $l$ is $\mathbf e^{\mathsf T}(I-\zeta\mathbf T_D)^{-1}\operatorname{Mom}^{\le D}_{\mathcal A_0}$, rational in $\zeta$, the sequence satisfies the linear recurrence with constant coefficients given by the characteristic polynomial of $\mathbf T_D$, of order at most $\binom{q+D}{D}$, and one length $l$ costs $O\big(\binom{q+D}{D}^3\log(l+1)\big)$ operations by repeated squaring. $\mathbf T_D$ is block diagonal by the degree, and a moment of one degree uses its block alone; a repeated sequence of different chain cells is the power of the product of their matrices.
\end{theorem}

On a cell with one hinge $g$, a chain cell $\mathcal U$ with the hinges $i$ and $o$, $i$ spliced to $g$, and the channels $K_{\mathcal U,(P)}$, $\Psi_{\mathcal U,(P)}$, $P\subseteq\{i,o\}$, acts by
\[
\begin{pmatrix}K_{\mathcal A_{l+1},(\varnothing)}\\K_{\mathcal A_{l+1},(\{o\})}\end{pmatrix}=\begin{pmatrix}K_{\mathcal U,(\{i\})}&K_{\mathcal U,(\varnothing)}\\K_{\mathcal U,(\{i,o\})}&K_{\mathcal U,(\{o\})}\end{pmatrix}\begin{pmatrix}K_{\mathcal A_l,(\varnothing)}\\K_{\mathcal A_l,(\{g\})}\end{pmatrix},
\]
and by the same matrix of the $\Psi_{\mathcal U,(P)}$ on the $\Psi$ channels, which is Theorem~\ref{thm:hingepairs}(iii) with one pair. A moment of degree $a$ in the $K$ channels and $b$ in the $\Psi$ channels lives on $(a+1)(b+1)$ monomials, and the block of $\mathbf T_D$ on them is the joint expectation over the draws of $\mathcal U$ of the expanded actions: the two matrices share the marks of $\mathcal U$ and are not independent. For a family at integer exponents, the joint moments of its letters along a chain of bunches with matrix letters are the case of Theorem~\ref{thm:momenttransfer}. An iterate of an integral operator gives no finite closing relation by itself (Theorem~\ref{thm:chainoperator}, Proposition~\ref{prop:chain}); every polynomial Gamma moment of a repetition does. The period is recovered from a sequence of such moments with certified tails, in Section~\ref{sec:certstop}.

% Reference 20: cor:shaperec; source line 7206
\begin{corollary}[Shape moments have recurrences]\label{cor:shaperec}
For the repetition of Theorem~\ref{thm:transferpower} let the total masses be $m_{\mathcal A_l}=m_{\mathcal A_0}+l\,m_{\mathcal U}$, $m_{\mathcal U}$ that of $\mathcal U$, let $p$ be a polynomial in the channels, the same for every $l$, homogeneous of degree $d_l=d_0+lv$ in the marks, $v$ an integer, and let $M_l$ be the polynomial Gamma moment of $\mathcal A_l$ for $p$. The shape moment for $p$ is $M_l/(m_{\mathcal A_l})_{d_l}$, and every recurrence $\sum_ic_iM_{l+i}=0$ gives $\sum_ic_i\frac{(m_{\mathcal A_{l+i}})_{d_{l+i}}}{(m_{\mathcal A_l})_{d_l}}\frac{M_{l+i}}{(m_{\mathcal A_{l+i}})_{d_{l+i}}}=0$. If the masses are rational and $h\ge1$ is an integer with $hm_{\mathcal U}\in\mathbb Z$, then on each residue of $l$ modulo $h$, with $\mathbf T_D^h$ for $\mathbf T_D$, the coefficients are rational functions of $l$ over the field of the weights, polynomials after clearing denominators, in $\mathbb Q[l]$ at rational weights.
\end{corollary}

\section{A certified enclosure of the period}\label{sec:certstop}

% Reference 20: thm:certstop; source line 7217
\begin{theorem}[Explicit stopping from the walls]\label{thm:certstop}
Let a family $F$ on $n\ge2$ marks have real masses on its balance wall and exponents $s_j,t_j\ge0$, with every proper wall positive. Write its active letters, the $K_j$ with $s_j>0$ and the $\Psi_j$ with $t_j>0$, as $L_i$ with exponents $w_i>0$, degrees $g_i$ and counts $d_i(S)$, as in Theorem~\ref{thm:peak}, so that $m_{\cE}=\sum_ig_iw_i$ and $\kappa_S=m_S-\sum_iw_id_i(S)$. Choose references $M_i\ge\max_{\Delta_{\cE}}L_i$, and in this statement put $Y_i=L_i/M_i\in(0,1]$ on the open simplex. For an order $\pi$ of the marks let $U_{\pi,k}$ be its tails, $L_{i,\pi}$ the sector letter of $L_i$, the $k_{j,\pi}$ or $\psi_{j,\pi}$ of Theorem~\ref{thm:exist}(ii), and $\underline L_{i,\pi}>0$ its constant term. For integers $Q_i\ge0$ put
\[
P^{\le\mathbf Q}=\Gst{n}{1}{\mvec}{m_{\cE}}\prod_iM_i^{-w_i}\sum_{0\le k_i\le Q_i}\prod_i\frac{(w_i)_{k_i}}{k_i!}\ \oE_{U\sim\Dir_{\mvec}}\Big[\prod_i\big(1-Y_i(U)\big)^{k_i}\Big],
\]
the partial sums of the peak closing (Theorem~\ref{thm:peak}(i)), computed exactly from the moments of the letters through the degree $\sum_iQ_i$, each monomial with its own rising factorial (Theorem~\ref{thm:gammalift}(iii)). For every choice of $\epsilon_i>0$ with
\[
\kappa_S-\epsilon_id_i(S)>0\qquad(\varnothing\neq S\subsetneq\cE),
\]
which is exactly $\epsilon_i<\varpi_i$, $\varpi_i$ the least of the ratios $\kappa_S/d_i(S)$ in Theorem~\ref{thm:peak}(ii) (any $\epsilon_i>0$ when no proper $S$ has $d_i(S)>0$), put
\[
C_i=M_i^{\epsilon_i}\,n^{g_i\epsilon_i}\sum_\pi\underline L_{i,\pi}^{-\epsilon_i}\prod_{i'}\underline L_{i',\pi}^{-w_{i'}}\prod_{k=2}^n\frac1{\kappa_{U_{\pi,k}}-\epsilon_id_i(U_{\pi,k})}.
\]
Then
\[
0\le P_F(\mvec;\mathbf s,\mathbf t)-P^{\le\mathbf Q}\le\sum_iC_i\,\frac{\Gamma(w_i+\epsilon_i)\,\Gamma(w_i+Q_i+1)}{\Gamma(w_i)\,\Gamma(w_i+Q_i+1+\epsilon_i)},
\]
and for $\epsilon_i\le1$ the last factor is at most $(w_i/(w_i+Q_i+1))^{\epsilon_i}$. So finitely many moments and finitely many walls enclose the period to any prescribed accuracy; the same moments answer every other choice of exponents $w_i$ with $\sum_ig_iw_i=m_{\cE}$ at the same masses, only the coefficients of $P^{\le\mathbf Q}$, the walls and the certificate changing.
\end{theorem}

At rational data the constants of the certificate are positive rationals and their rational powers, and no period other than the one sought enters it; the stacked Gamma $\Gst{n}{1}{\mvec}{m_{\cE}}$ and the Gamma ratios stay exact products and quotients of values of the Gamma function, so the enclosure is exact without being rational. The certificate can itself be expensive, $C_i$ being a sum over all $n!$ orders, and it is crude: it replaces every sector letter by its least value. The true tail is of the order $Q^{-\varpi}(\log Q)^{k-1}$ (Theorem~\ref{thm:peak}(ii)); the certified decay $Q_i^{-\epsilon_i}$ tends to it as $\epsilon_i\to\varpi_i$, at the cost of the factor $\prod_k(\kappa_{U_{\pi,k}}-\epsilon_id_i(U_{\pi,k}))^{-1}$, which grows without bound there. Sharper references and sharper sector bounds lower the degree needed without changing the moments. The degree needed can grow with the size of the assembly, so this is a terminating certified procedure, not a uniform fast approximation.

\section{Changing the query with finite data}\label{sec:symbolicmoments}

% Reference 20: prop:symbolicmoments; source line 7248
\begin{proposition}[Parameters and insertions]\label{prop:symbolicmoments}
Fix the local supports and the assembly.
\begin{enumerate}[label=(\roman*)]
\item For each $D$ the leaf moments and the splicing and transfer coefficients are polynomials in the local weights and in the masses; they can be kept symbolic and specialized at any positive values.
\item Carrying the leaf moments, Theorem~\ref{thm:gammalift}(ii) and Theorem~\ref{thm:transferpower} in the truncated ring $\mathbb Q(\boldsymbol\theta,\mvec)[\boldsymbol\zeta]/(\boldsymbol\zeta)^{h+1}$ of $p$ local parameters $\boldsymbol\zeta$, over the field generated by the weights $\boldsymbol\theta$ and the masses, gives every derivative of the moments through the order $h$, at most $\binom{p+h}{h}$ coefficients per entry.
\item A further polynomial observable, a mark total for instance, can be made a channel when every local update is separately affine in each independent cell vector, a constant channel adjoined.
\item Raising the masses by nonnegative integers $\boldsymbol\gamma$ is a polynomial insertion: for every observable $g$,
\[
\oE_{\alpha\sim\operatorname{Gamma}_{\mvec+\boldsymbol\gamma}}[g(\alpha)]=\oE_{\alpha\sim\operatorname{Gamma}_{\mvec}}\big[\alpha^{\boldsymbol\gamma}g(\alpha)\big]\Big/\prod_e(m_e)_{\gamma_e}.
\]
\end{enumerate}
\end{proposition}

The moments do not depend on the exponents asked for: changing the exponents reuses every moment, and changing a local weight or mass changes the local data and their ancestors in the assembly only, as Proposition~\ref{prop:local} says of laws. Shared parameters may appear in many copies, whose draws stay independent.

\section{Certified parameter banks}\label{sec:banks}

A family $F$ on its balance wall, its letters $K_j$ and $\Psi_j$ written $L_i$, $i=1,\dots,\ell$, with the exponents $w_i$ (the $s_j$ and $t_j$), the degrees $g_i$ and the counts $d_i(S)=\min\{\gamma_S:\ [u^{\boldsymbol\gamma}]L_i>0\}$, so that $P_F(\mvec;\mathbf s,\mathbf t)=\int_{\Delta_{\cE}}u^{\mvec-\one}\prod_iL_i(u)^{-w_i}d\omega_{\cE}(u)$, with $m_{\cE}=\sum_ig_iw_i$ and the walls $\kappa_S=m_S-\sum_iw_id_i(S)>0$ at the starting data, is asked a neighbourhood of queries: $p$ complex parameters $\boldsymbol\zeta$, and a bank of Taylor coefficients through the total order $h$, at most $\binom{h+p}{p}$ numbers. The order of a bank is the number of differentiations, not the degree of a polynomial moment of Section~\ref{sec:channellaw}.

% Reference 20: thm:respwalls; source line 7267 (nonnegative exponents).
% Proved in the last part: thm:respwalls.
\begin{theorem}[Taylor coefficients in the weights are periods with walls at least as high]\label{thm:respwalls}
Let the exponents $w_i$ be real and change the coefficients of the letters, $L_i(\boldsymbol\zeta)=L_i+\sum_{j=1}^p\zeta_jD_{ij}$, each $D_{ij}$ homogeneous of degree $g_i$ with support in that of $L_i$ and coefficients of either sign, and in this statement write $P(\boldsymbol\zeta)=\int_{\Delta_{\cE}}u^{\mvec-\one}\prod_iL_i(\boldsymbol\zeta)(u)^{-w_i}\,d\omega_{\cE}(u)$, so that $P(\mathbf 0)=P_F(\mvec;\mathbf s,\mathbf t)$. Then, with $l_i=\sum_jl_{ij}$,
\[
[\boldsymbol\zeta^{\mathbf k}]P=\sum_{\substack{l_{ij}\ge0\\ \sum_il_{ij}=k_j}}\ \prod_i\frac{(-1)^{l_i}(w_i)_{l_i}}{\prod_jl_{ij}!}\int_{\Delta_{\cE}}u^{\mvec-\one}\prod_i\Big(L_i^{-w_i-l_i}\prod_jD_{ij}^{l_{ij}}\Big)\,d\omega_{\cE},
\]
and after the numerators are expanded every term is a convergent period of the family at masses $\mvec+\boldsymbol\gamma$ and exponents $w_i+l_i$, balanced, with every proper wall $\kappa'_S=\kappa_S+\gamma_S-\sum_il_id_i(S)\ge\kappa_S$.
\end{theorem}

% Reference 20: thm:weightbank; source line 7278 (nonnegative exponents).
% Proved in the last part: thm:weightbank.
\begin{theorem}[A uniform certificate for weights]\label{thm:weightbank}
In the setting of Theorem~\ref{thm:respwalls}, with $P(\boldsymbol\zeta)$ as there, write $L_i=\sum_{\boldsymbol\gamma}c_{i,\boldsymbol\gamma}u^{\boldsymbol\gamma}$ with $c_{i,\boldsymbol\gamma}>0$ and $D_{ij}=\sum_{\boldsymbol\gamma}c_{i,\boldsymbol\gamma}v_{i,\boldsymbol\gamma,j}u^{\boldsymbol\gamma}$, $V_{ij}=\max_{\boldsymbol\gamma}|v_{i,\boldsymbol\gamma,j}|$, and choose radii $R_j>0$ with $r_i=\sum_jR_jV_{ij}<1$. On $|\zeta_j|\le R_j$, $|P(\boldsymbol\zeta)|\le P(\mathbf 0)\prod_i\max\big((1-r_i)^{-w_i},(1+r_i)^{-w_i}\big)=:M$, made explicit by $P(\mathbf 0)\le\sum_\pi\prod_i\max\big(L_{i,\pi}(\mathbf 0)^{-w_i},L_{i,\pi}(\one)^{-w_i}\big)\prod_{k=2}^n\kappa_{U_{\pi,k}}^{-1}$, with the sector letters $L_{i,\pi}$ of the letters on the sector of the order $\pi$, the $k_{j,\pi}$ or $\psi_{j,\pi}$ of Theorem~\ref{thm:exist}(ii), or by the upper end of the enclosure of Proposition~\ref{prop:sectorcompiler}; $P$ is holomorphic near the closed polydisc, $|[\boldsymbol\zeta^{\mathbf k}]P|\le M\mathbf R^{-\mathbf k}$, and its Taylor polynomial $P^{[h]}$ through the order $h$ has, on $|\zeta_j|\le\vartheta R_j$ with $0\le\vartheta<1$,
\[
|P(\boldsymbol\zeta)-P^{[h]}(\boldsymbol\zeta)|\le M\sum_{k>h}\binom{k+p-1}{p-1}\vartheta^k\le M\,\frac{\binom{h+p}{p-1}\vartheta^{h+1}}{1-\vartheta(h+p+1)/(h+2)},
\]
the second bound when $\vartheta(h+p+1)<h+2$.
\end{theorem}

% Reference 20: thm:wallbank; source line 7289
\begin{theorem}[Wall-certified parameter banks]\label{thm:wallbank}
Keep the letters and change the masses and exponents affinely, $m_e(\boldsymbol\zeta)=m^0_e+\sum_ja_{ej}\zeta_j$ and $w_i(\boldsymbol\zeta)=w^0_i+\sum_jb_{ij}\zeta_j$ with real slopes and balance in every direction, so that $\kappa_S(\boldsymbol\zeta)=\kappa^0_S+\sum_jc_{Sj}\zeta_j$ with $c_{Sj}=\sum_{e\in S}a_{ej}-\sum_ib_{ij}d_i(S)$, and in this statement write $P(\boldsymbol\zeta)=P_F(\mvec(\boldsymbol\zeta);\mathbf s(\boldsymbol\zeta),\mathbf t(\boldsymbol\zeta))$, the exponents $w_i(\boldsymbol\zeta)$ on the letters. Choose $R_j>0$ with
\[
\underline\kappa_S:=\kappa^0_S-\sum_j|c_{Sj}|R_j>0\qquad(\varnothing\neq S\subsetneq\cE).
\]
On the sector of an order $\pi$ every letter is its monomial times its sector letter $L_{i,\pi}$, the $k_{j,\pi}$ or $\psi_{j,\pi}$ of Theorem~\ref{thm:exist}(ii), with nonnegative coefficients in the sector coordinates; put $\underline L_{i,\pi}=L_{i,\pi}(\mathbf 0)$, $\overline L_{i,\pi}=L_{i,\pi}(\one)$, $\eta_{i,\pi}=\max(|\log\underline L_{i,\pi}|,|\log\overline L_{i,\pi}|)$, $\nu_i=\sum_j|b_{ij}|R_j$,
\[
B_\pi=\prod_i\max\big(\underline L_{i,\pi}^{-w^0_i},\overline L_{i,\pi}^{-w^0_i}\big)e^{\eta_{i,\pi}\nu_i},\qquad M_{\mathbf R}=\sum_\pi B_\pi\prod_{k=2}^n\underline\kappa_{U_{\pi,k}}^{-1}.
\]
Then $P$ is holomorphic near $|\zeta_j|\le R_j$ with $|P|\le M_{\mathbf R}$, every $|[\boldsymbol\zeta^{\mathbf k}]P|\le M_{\mathbf R}\mathbf R^{-\mathbf k}$, the bound of Theorem~\ref{thm:weightbank} holds with $M_{\mathbf R}$ for $M$, and
\[
[\boldsymbol\zeta^{\mathbf k}]P=\frac1{\mathbf k!}\int_{\Delta_{\cE}}u^{\mvec^0-\one}\prod_iL_i^{-w^0_i}\prod_j\Big(\sum_ea_{ej}\log u_e-\sum_ib_{ij}\log L_i\Big)^{k_j}\,d\omega_{\cE},
\]
absolutely integrable and of degree zero. The starting exponents may be any reals at which the proper walls are positive.
\end{theorem}

% Reference 20: thm:twoerrors; source line 7308
\begin{theorem}[A finite compiler with two errors]\label{thm:twoerrors}
In the setting of Theorem~\ref{thm:wallbank}, with $P(\boldsymbol\zeta)$ as there, put $y_{i,\pi}=1-L_{i,\pi}/\overline L_{i,\pi}\in[0,\lambda_\pi]$, $\lambda_\pi=\max_i(1-\underline L_{i,\pi}/\overline L_{i,\pi})<1$, $\beta_i=|w^0_i|+\nu_i$ and $A_\pi=\prod_i\overline L_{i,\pi}^{-w^0_i}e^{|\log\overline L_{i,\pi}|\nu_i}\prod_{k=2}^n\underline\kappa_{U_{\pi,k}}^{-1}$. Expand $\prod_iL_{i,\pi}^{-w_i(\boldsymbol\zeta)}=\prod_i\overline L_{i,\pi}^{-w_i(\boldsymbol\zeta)}\sum_{\mathbf l}\prod_i\frac{(w_i(\boldsymbol\zeta))_{l_i}}{l_i!}y_{i,\pi}^{l_i}$, truncate at $|\mathbf l|\le Q$, and integrate the finitely many terms exactly, each $\prod_iy_{i,\pi}^{l_i}=\sum_{\boldsymbol\gamma}c_{\boldsymbol\gamma}q^{\boldsymbol\gamma}$ in the sector coordinates $q$ contributing
\[
\prod_i\overline L_{i,\pi}^{-w_i(\boldsymbol\zeta)}\prod_i\frac{(w_i(\boldsymbol\zeta))_{l_i}}{l_i!}\sum_{\boldsymbol\gamma}c_{\boldsymbol\gamma}\prod_{k=2}^n\frac1{\kappa_{U_{\pi,k}}(\boldsymbol\zeta)+\gamma_k},
\]
to obtain $\widetilde P_Q(\boldsymbol\zeta)$. Then on the whole polydisc
\[
|P(\boldsymbol\zeta)-\widetilde P_Q(\boldsymbol\zeta)|\le\operatorname{err}_Q:=\sum_\pi A_\pi\sum_{k>Q}\frac{(|\boldsymbol\beta|)_k}{k!}\lambda_\pi^k\le\sum_\pi A_\pi\frac{(|\boldsymbol\beta|)_{Q+1}\lambda_\pi^{Q+1}}{(Q+1)!\,(1-r_{\pi,Q})},
\]
the second bound when $r_{\pi,Q}=\lambda_\pi\max\big(1,\frac{|\boldsymbol\beta|+Q+1}{Q+2}\big)<1$ for every $\pi$; and the Taylor polynomial $\widetilde P_Q^{[h]}$ of $\widetilde P_Q$ through the order $h$, computed by finite arithmetic in a truncated ring, has on $|\zeta_j|\le\vartheta R_j$
\[
|P(\boldsymbol\zeta)-\widetilde P_Q^{[h]}(\boldsymbol\zeta)|\le\frac{\operatorname{err}_Q}{(1-\vartheta)^p}+M_{\mathbf R}\sum_{k>h}\binom{k+p-1}{p-1}\vartheta^k .
\]
With no parameters, $p=0$, so that $\nu_i=0$ and $\underline\kappa_S=\kappa_S$, and with every $w^0_i\ge0$, the first bound is that of Proposition~\ref{prop:sectorcompiler}(ii): there $\overline L_{i,\pi}=L_i(\one)$ for every $\pi$, $\widetilde P_Q=\prod_iL_i(\one)^{-w^0_i}\sum_\pi\sum_{M\le Q}I_{\pi,M}$, and the explicit tail is $\prod_iL_i(\one)^{-w^0_i}\sum_\pi\operatorname{err}_{\pi,Q}$, the letters being the $K_j$ and $\Psi_j$ with the exponents $s_j$ and $t_j$.
\end{theorem}

For fixed data and $\vartheta<1$ the orders $Q$ and $h$ grow like the logarithm of the inverse error, up to the polynomial factors of the scalar tails; expanding the sector polynomials still costs work that depends on the dimension and the supports. The walls also cover paths: along a segment on which every proper wall is positive, with $\kappa_-$ the least proper wall at the two ends and $a$ the largest change of a proper wall along it, every point of the segment, parametrized by $\lambda\in[0,1]$, has the disc of radius $\kappa_-/(2a)$ in the complex parameter on which every proper wall keeps a real part at least $\kappa_-/2$, the walls being affine, and a compact path has a finite cover by such discs. A bank at one point does not give the coefficients at another; what is reused everywhere is the compiled sector data, the coefficient procedure and the wall test.

\section{Procedures}\label{sec:interfaceproc}

% Reference 20: def:interfaceproc; source line 7332
\begin{definition}[The interface procedure]\label{def:interfaceproc}
\emph{Given}: cells with their weights and masses; the pairs of hinges to splice and an assembly order with its identifications of repeated parts; a polynomial moment, a peak truncation or an accuracy, with, for an accuracy, references and a wall certificate. \emph{Steps}: (1) read the channels of the leaves by coefficient extraction, with their degrees, the sizes, the total masses and the products $\prod_e\Gamma(m_e)$; (2) choose $D$, the degree of the moment or $\sum_iQ_i$ from Theorem~\ref{thm:certstop}; (3) compute the leaf moments by Theorem~\ref{thm:gammalift}(i), only at the multi-indices the moment pulls back to; (4) compile every distinct splicing by Theorem~\ref{thm:gammalift}(ii), and every repeated chain cell once into $\mathbf T_D$ and its power (Theorem~\ref{thm:transferpower}), reusing a vector only when its tags and data match; (5) splice the pairs at the root, return to the shape by Theorem~\ref{thm:gammalift}(iii), and evaluate $P^{\le\mathbf Q}$ with its certificate. \emph{Output}: an exact polynomial moment, a recurrence or rational generating series for a repetition, or a certified enclosure of the period, with the vectors, the splicing coefficients, the parameter expressions and the certificate kept for the next query.
\end{definition}

% Reference 20: def:bankproc; source line 7336
\begin{definition}[The bank procedure]\label{def:bankproc}
\emph{Given}: a family on its balance wall with positive walls, and affine directions of change of its masses and exponents, or of the weights of its letters. \emph{Steps} along masses and exponents: compute the walls and a box of radii satisfying the strict inequalities of Theorem~\ref{thm:wallbank}; form the sector constants and $M_{\mathbf R}$; choose $Q$ and $h$ until the bound of Theorem~\ref{thm:twoerrors} is below the tolerance; compile the retained sector terms in the truncated ring. Along weights, where the sector letters move with $\boldsymbol\zeta$ and Theorem~\ref{thm:twoerrors} does not apply: choose radii by Theorem~\ref{thm:weightbank}, with its bound made explicit there; each coefficient through the order $h$ is the finite signed sum of periods with walls at least as high of Theorem~\ref{thm:respwalls}, each enclosed by Proposition~\ref{prop:sectorcompiler}, and the tail beyond $h$ is bounded by Theorem~\ref{thm:weightbank}. \emph{Output}: the parameter bank, its domain and its certified error; every query inside the smaller box is a polynomial evaluation, and the exact coefficients remain available as the periods of Theorem~\ref{thm:respwalls} or the integrals of Theorem~\ref{thm:wallbank}.
\end{definition}


\chapter{Bunches of any width}\label{ch:width}

A bunch of width one is one mark to its room (Theorem~\ref{thm:widthone}): its patterns fill two levels, each level is one child of the bunch against one table of the room, and what the room sees of the bunch is its wall and the law of its ratio vector, one ratio for one table (Lemma~\ref{lem:absorb}, Theorem~\ref{thm:retained}). At every width a room sees its bunch through the wall, the widths and the joint law of the bunch letters of a level factorization (Theorem~\ref{thm:crossing}); that statement has a canonical form, and the calculus of width one carries to every width. At the width $c$ the patterns fill $c+1$ levels, and a level between the extremes pairs several stations of patterns with several stations of room sides: which pairs meet, and with which constants, is the seam of the bunch (Definition~\ref{def:seam}), graded by the excess, which no description by marks or by level totals contains. The room sees the bunch through the letters of the tables it presents to it, one for each station of room sides, and the law of these letters is what the bunch retains: it suffices against every partner, and the tinkers of one partner already separate any two laws (Theorem~\ref{thm:retainwidth}). A bunch of width $c$ is absorbed into its partner as one stand-in mark of degree $c$ (Lemma~\ref{lem:absorbwidth}), a mark of a table exactly at width one (Theorem~\ref{thm:storeone}), and a nonempty family of sets of one size is the family of bases of a table exactly when, at one bunch, its children and room children are tables, adjacent ones are quotients, the extreme ones are one, and its levels are lifts (Theorem~\ref{thm:gluing}). The circuits of the seam are the directions of the shares of the compatible pairs that neither side sees, and the period integrates over them through the Gramma of the compatible pairs (Proposition~\ref{prop:seamtransfer}). Width one is the case of the hinge throughout: the seam of a bunch of width one is its hinge (Proposition~\ref{prop:ladder}(ii)), and its law is the law of its ratio vector read through one fixed change of scale (Proposition~\ref{prop:widthonelaw}). Bunches compose along chains and trees (Chapter~\ref{ch:composebunches}), the glue closes along the directions of its shares (Chapter~\ref{ch:closeglue}), and a law has regions of equality and a reach at its ends (Chapter~\ref{ch:lawreach}).

\section{The seam of a bunch}\label{sec:seam}

% Reference 20: def:seam; source line 7347
% Proved in the last part: def:seam (the pair constant).
\begin{definition}[Seam]\label{def:seam}
Let a table $(\cB,\theta)$ have the bunch $B$ and the room $R$. A level $k$ of $B$ is named by its excess $i=k-d(B)$ (Definition~\ref{def:levels}), $0\le i\le c(B)$: its patterns have $d(B)+i$ marks and its room sides $d(R)+c(B)-i$, the room having the width of the bunch (Definition~\ref{def:width}). Fix an excess $i$.
\begin{enumerate}[label=(\roman*)]
\item Two patterns of excess $i$ lie in one station when their children (Definition~\ref{def:patterns}) are positive multiples of one another, weights included, as at the bunches $S_l$ of Theorem~\ref{thm:lettertransfer}. The \emph{room child} at a room side $X$ is the child of the bunch $R$ at its pattern $X$: the table on $B$ with the bases $\cB_{(X)}=\{p:\ p\sqcup X\in\cB\}$ and the weights $\theta_{(X),p}=\theta_{p\sqcup X}$ (Lemma~\ref{lem:minor}). Two room sides of excess $i$ lie in one station when their room children are positive multiples of one another. The stations of patterns of excess $i$ form the set $\mathcal K_i$ and those of room sides the set $\mathcal D_i$: the stations of the bunch and of the room are $\mathcal K=\bigsqcup_i\mathcal K_i$ and $\mathcal D=\bigsqcup_i\mathcal D_i$, each station carrying its excess $i(C)$ or $i(D)$.
\item Within a station $C$ the child at $p\in C$ is $w_p$ times one weighted table on $R$, the \emph{child of the station} $C$, with the bases $\cB_{(C)}$ and the weights $\theta_{(C)}$, the \emph{multipliers} $w_p>0$ being normalized by $\max_{p\in C}w_p=1$; likewise the room child at $X\in D$ is $w'_X$ times the \emph{room child of the station} $D$, a table on $B$ with the bases $\cB_{(D)}$ and the weights $\theta_{(D)}$, with $\max_{X\in D}w'_X=1$. The \emph{pair constant} of $C\in\mathcal K_i$ and $D\in\mathcal D_i$ is $h_{CD}=\theta_{p\sqcup X}/(w_pw'_X)$ for any $p\in C$ and $X\in D$ with $p\sqcup X\in\cB$, which does not depend on the choice, and $h_{CD}=0$ when there is none; $(C,D)$ is \emph{compatible} when $h_{CD}>0$. The \emph{seam matrix} of the excess $i$ is $H_i=(h_{CD})_{C\in\mathcal K_i,\,D\in\mathcal D_i}$, and the \emph{seam} of $B$ is $(H_0,\dots,H_{c(B)})$ with its stations; a family has one seam in each table.
\item The \emph{station letters} are $\mathsf x_C=\sum_{p\in C}w_p\,\alpha_B^{\,p}$ and $\mathsf x^\ast_C=\sum_{p\in C}w_p\,\alpha_B^{\,B\setminus p}$ on $B$, and $\mathsf y_D=\sum_{X\in D}w'_X\,\alpha_R^{\,X}$ and $\mathsf y^\ast_D=\sum_{X\in D}w'_X\,\alpha_R^{\,R\setminus X}$ on $R$. The \emph{seam letters} of $B$ are the letters $K_{(D)}$ and $\Psi_{(D)}$ of the room children of the stations $D\in\mathcal D$, polynomials on $B$; those of $R$ are the letters $K_{(C)}$ and $\Psi_{(C)}$ of the children of the stations $C\in\mathcal K$, polynomials on $R$.
\end{enumerate}
For a list every multiplier and every pair constant is $1$. The normalization by the largest multiplier chooses nothing and commutes with every bijection of the marks; any other normalization multiplies $\mathsf x_C$ and $\mathsf y_D$ by positive constants and $h_{CD}$ by their inverses, and leaves every product $h_{CD}\,\mathsf x_C\,\mathsf y_D$ unchanged. In table $j$ of a family the stations, the letters and the constants carry the index $j$: $\mathcal K_j$, $\mathcal D_j=\bigsqcup_i\mathcal D_{j,i}$, $\mathsf x_{jC}$, $\mathsf y_{jD}$, $h_{j,CD}$, $K_{j,(D)}$ and so on.
\end{definition}

% Reference 20: prop:seamidentity; source line 7357 (single entries when the weights of the extreme levels factor).
% Proved in the last part: prop:seamidentity.
\begin{proposition}[The seam identity]\label{prop:seamidentity}
In Definition~\ref{def:seam}:
\begin{enumerate}[label=(\roman*)]
\item The seam letters are the station letters through the seam:
\begin{gather*}
K_{(D)}=\sum_Ch_{CD}\,\mathsf x_C,\qquad\Psi_{(D)}=\sum_Ch_{CD}\,\mathsf x^\ast_C,\\
K_{(C)}=\sum_Dh_{CD}\,\mathsf y_D,\qquad\Psi_{(C)}=\sum_Dh_{CD}\,\mathsf y^\ast_D:
\end{gather*}
the seam letters of a side are the letters of the tables that the other side presents to it.
\item If $(C,D)$ is compatible, every $p\in C$ and $X\in D$ make a basis, with $\theta_{p\sqcup X}=h_{CD}\,w_p\,w'_X$, and the bases of excess $i$ are the disjoint union of these \emph{rectangles} $C\times D$ over the compatible pairs.
\item With $\mathsf x_i$, $\mathsf x^\ast_i$, $\mathsf y_i$, $\mathsf y^\ast_i$ the vectors of the station letters of excess $i$,
\begin{gather*}
K=\sum_{i=0}^{c(B)}\mathsf x_i^{\mathsf T}H_i\,\mathsf y_i=\sum_{D\in\mathcal D}K_{(D)}\,\mathsf y_D=\sum_{C\in\mathcal K}\mathsf x_C\,K_{(C)},\\
\Psi=\sum_{i=0}^{c(B)}\mathsf x_i^{\ast\mathsf T}H_i\,\mathsf y^\ast_i=\sum_{D\in\mathcal D}\Psi_{(D)}\,\mathsf y^\ast_D=\sum_{C\in\mathcal K}\mathsf x^\ast_C\,\Psi_{(C)},
\end{gather*}
and along the scale of $B$, $\alpha_B=\rho u$ with $u\in\Delta_B$ (Lemma~\ref{lem:scaleshape}),
\begin{gather*}
K(\rho u,\alpha_R)=\rho^{d(B)}\sum_{i=0}^{c(B)}\rho^i\sum_{D\in\mathcal D_i}K_{(D)}(u)\,\mathsf y_D(\alpha_R),\\
\Psi(\rho u,\alpha_R)=\rho^{d^\ast(B)}\sum_{i=0}^{c(B)}\rho^{c(B)-i}\sum_{D\in\mathcal D_i}\Psi_{(D)}(u)\,\mathsf y^\ast_D(\alpha_R).
\end{gather*}
The writing by the stations of the room, with the bunch letters $K_{(D)}$ and the room letters $\mathsf y_D$, is a level factorization with nonnegative terms (Definition~\ref{def:levelweights}), the \emph{room-station factorization}; the writing by the stations of the bunch is the children factorization.
\item Within one excess no two stations of $\mathcal K_i$ have proportional rows of $H_i$, and no two stations of $\mathcal D_i$ proportional columns.
\item At the excesses $0$ and $c(B)$ every pattern meets every room side in a basis (Theorem~\ref{thm:levels}). $H_0$, respectively $H_{c(B)}$, is a single positive entry exactly when the weights of that level factor along these products, as for unit weights, for weights factoring across $B$ and $R$ and for tables determinantal over $\mathbb R$ or $\mathbb C$; when both do, the room child of the one station of excess $0$ is $\cB/R$ and that of excess $c(B)$ is $\cB|B$ (Proposition~\ref{prop:higgs}), these tables carrying the weights of the room children of the two stations, and the seam letters at the two extremes are $K_{\cB/R}$, $\Psi_{\cB/R}$ and $K_{\cB|B}$, $\Psi_{\cB|B}$.
\item Swapping $B$ and $R$ transposes every $H_i$ and reverses the excesses, $i\mapsto c(B)-i$; the seam letters of $R$ are then those of the bunch $R$.
\end{enumerate}
\end{proposition}

A compatible pair fills its rectangle: inside it every tuple of the bunch meets every tuple of the room, with weights of rank one. What the marks do not show is the pattern of the rectangles, which stations of the bunch meet which stations of the room. At an extreme excess every station of the bunch meets every station of the room, in one rectangle when the weights of that level factor; between the extremes a station need not meet every station of the other side, and the seam matrix $H_i$, not the marks, says how they meet.

% Reference 20: prop:ladder; source line 7382 (one rectangle when the weights factor; level form for $c(B)\ge2$).
% Proved in the last part: prop:ladder.
\begin{proposition}[The width ladder]\label{prop:ladder}
In Definition~\ref{def:seam}:
\begin{enumerate}[label=(\roman*)]
\item $c(B)=0$ exactly when $B$ is a union of molecules of the table; there is one excess, every pattern meets every room side, and $K=\mathsf x_0^{\mathsf T}H_0\,\mathsf y_0$, $\Psi=\mathsf x_0^{\ast\mathsf T}H_0\,\mathsf y^\ast_0$. The seam is one rectangle $C\times D$ exactly when the weights factor across $B$ and $R$, as for lists; then $K=h_{CD}\,\mathsf x_C\,\mathsf y_D$ and $\Psi=h_{CD}\,\mathsf x^\ast_C\,\mathsf y^\ast_D$: the letters factor over the bunch and the room.
\item $c(B)=1$ exactly when there are two excesses, both extreme. Each is one rectangle exactly when the weights factor across $B$ and $R$; then for the stations $D_0$ and $D_1$ of the excesses $0$ and $1$ the seam letters are $K_{(D_0)}=K_{\cB/R}$, $K_{(D_1)}=K_{\cB|B}$, $\Psi_{(D_0)}=\Psi_{\cB/R}$ and $\Psi_{(D_1)}=\Psi_{\cB|B}$, the room station letters are $\mathsf y_{D_0}=K_{\cB|R}$ and $\mathsf y_{D_1}=K_{\cB/B}$, these two tables carrying the multipliers, the largest $1$, and the room-station factorization is the writing
\[
K=K_{\cB/R}\,K_{\cB|R}+K_{\cB|B}\,K_{\cB/B}
\]
of Theorem~\ref{thm:widthone}: the letter adds over two levels, and the seam is the hinge.
\item For $c(B)\ge2$ the excesses $0<i<c(B)$ can carry several rectangles.
\end{enumerate}
At every width an excess with one station on one side has one station on the other, and the children at the patterns of each excess are equal, $B$ being in level form, exactly when every transposition of two marks of $B$, fixing every other mark, carries the table (Proposition~\ref{prop:levelform}); then every excess is one rectangle.
\end{proposition}

\section{Pieces, their laws, and the glue}\label{sec:seampieces}

The seam identity writes a family as a pairing of what the bunch presents to the room and what the room presents to the bunch, station by station. The pairing uses nothing of the exchange property and holds for pieces of any kind.

% Reference 20: def:seampiece; source line 7411
\begin{definition}[Pieces with stations; their glue]\label{def:seampiece}
Let $N\ge1$ and $c_1,\dots,c_N\ge0$ be integers. A set of \emph{stations} for $N$ tables of the widths $c_j$ is a family $\mathcal D=(\mathcal D_j)_{j=1}^N$ of finite sets graded by excesses, $\mathcal D_j=\bigsqcup_{i=0}^{c_j}\mathcal D_{j,i}$, with $\mathcal D_{j,0}$ and $\mathcal D_{j,c_j}$ nonempty; a station $D\in\mathcal D_{j,i}$ has the excess $i(D)=i$. A \emph{piece} with the stations $\mathcal D$ on a finite set $B$, at exponents $(s_j,t_j)_j$ with $s_j,t_j\ge0$, carries for every $j$ and $D\in\mathcal D_j$ two nonzero polynomials $V_{jD}$ and $V^\ast_{jD}$ on $B$ with nonnegative coefficients, its \emph{seam letters} (positive letters, Definition~\ref{def:posword}, where their degree is positive), homogeneous of the degrees $d_j+i(D)$ and $d^\ast_j+c_j-i(D)$ for integers $d_j,d^\ast_j\ge0$, and masses $\mvec_B$; its \emph{wall} and its \emph{top wall} are
\[
\kappa_B=m_B-\sum_j\big(s_jd_j+t_jd^\ast_j\big),\qquad\kappa^{\mathrm{top}}_B=\sum_j(s_j+t_j)c_j-\kappa_B .
\]
A \emph{partner} is a piece with the same stations on a set $R$ disjoint from $B$, read with the reversed excesses $c_j-i(D)$: its letters, the \emph{room letters} $L_{jD}$ and $L^\ast_{jD}$ (as in Definition~\ref{def:levelweights}), have the degrees $e_j+c_j-i(D)$ and $e^\ast_j+i(D)$, and its wall is $\kappa_R=m_R-\sum_j(s_je_j+t_je^\ast_j)$. A \emph{family of Newton words} is a sequence $(K_j,\Psi_j;s_j,t_j)_{j=1}^N$ of pairs of positive letters on one set with exponents $s_j,t_j\ge0$, carrying masses; its word is $\prod_jK_j^{-s_j}\Psi_j^{-t_j}$ and its period is the period of that word (Definition~\ref{def:posword}). The \emph{glue} of the piece and the partner is the family of Newton words on $B\sqcup R$ with the letters
\[
K_j=\sum_{D\in\mathcal D_j}V_{jD}\,L_{jD},\qquad\Psi_j=\sum_{D\in\mathcal D_j}V^\ast_{jD}\,L^\ast_{jD},
\]
the exponents $(s_j,t_j)$ and the masses $(\mvec_B,\mvec_R)$; its period is written $P_{B\sqcup R}(\mvec)$, and its balance wall is $\kappa_R=\kappa^{\mathrm{top}}_B$. A piece is named by its set, or by a calligraphic letter $\mathcal X$ when several pieces are compared, and its objects carry that name.
\end{definition}

By Proposition~\ref{prop:seamidentity}(iii) a bunch of a family of tables, with its seam letters $V_{jD}=K_{j,(D)}$ and $V^\ast_{jD}=\Psi_{j,(D)}$, is a piece with the stations of its room, $d_j=d_j(B)$ and $d^\ast_j=d^\ast_j(B)$; its room, with the room station letters $L_{jD}=\mathsf y_{jD}$ and $L^\ast_{jD}=\mathsf y^\ast_{jD}$, is a partner, $e_j=d_j(R)$ and $e^\ast_j=d^\ast_j(R)$; and their glue is the family. By (vi) the same holds with $B$ and $R$ swapped, on the stations $\mathcal K$. The walls are the walls $\kappa_B$ and $\kappa_R$ of the bunch and of the room, and the top wall is the top wall $\kappa^{\mathrm{top}}_B$ of the bunch (Definition~\ref{def:family}). A glue need not be a family of tables; its period converges exactly when its Newton sum has the dimension $|B\sqcup R|-1$ and the masses lie in its relative interior (Theorem~\ref{thm:poswordconv}), and every statement about pieces is a statement about the period.

% Reference 20: def:seamlaw; source line 7421
\begin{definition}[The law of a piece]\label{def:seamlaw}
Let a piece on $B$ with the stations $\mathcal D$ have exponents with $\sum_j(s_j+t_j)c_j>0$. Its active letters (Definition~\ref{def:family}) are $K_j$ when $s_j>0$ and $\Psi_j$ when $t_j>0$; every coordinate, share, pin, product and count indexed by a letter runs over the active letters only, an inactive letter entering no period. The \emph{excess sums} are $V^{(i)}_j=\sum_{D\in\mathcal D_{j,i}}V_{jD}$ and $V^{\ast(i)}_j=\sum_{D\in\mathcal D_{j,i}}V^\ast_{jD}$, the superscript counting excesses. On $\Delta_B$ put
\begin{gather*}
\gamma^\rho_B=\prod_j\Big(\frac{V^{(c_j)}_j}{V^{(0)}_j}\Big)^{s_j/(\kappa_B+\kappa^{\mathrm{top}}_B)}\Big(\frac{V^{\ast(0)}_j}{V^{\ast(c_j)}_j}\Big)^{t_j/(\kappa_B+\kappa^{\mathrm{top}}_B)},\\
\phi^\rho_B(du)=u^{\mvec_B-\one}\prod_j\big(V^{(0)}_j\big)^{-s_j}\big(V^{\ast(c_j)}_j\big)^{-t_j}\,(\gamma^\rho_B)^{-\kappa_B}\,d\omega_B(u),\\
\varrho^\rho_{jD}=\frac{V_{jD}}{V^{(0)}_j\,(\gamma^\rho_B)^{i(D)}},\qquad\varrho^{\rho\ast}_{jD}=\frac{V^\ast_{jD}}{V^{\ast(c_j)}_j\,(\gamma^\rho_B)^{c_j-i(D)}},
\end{gather*}
the \emph{scale of the piece}, homogeneous of degree one, the \emph{scale form}, and the coordinates of the \emph{seam vector} $\boldsymbol\varrho^\rho_B=(\varrho^\rho_{jD},\varrho^{\rho\ast}_{jD})_{j,D}$, homogeneous of degree zero, with a coordinate $\varrho^\rho_{jD}$ for each active $K_j$ and $D\in\mathcal D_j$ and a coordinate $\varrho^{\rho\ast}_{jD}$ for each active $\Psi_j$ and $D\in\mathcal D_j$. Its values, the \emph{seam points}, lie on the \emph{seam torus}: the set of the points $\mathbf x=(x_{jD},x^\ast_{jD})_{j,D}$ with positive coordinates on which $\sum_{D\in\mathcal D_{j,0}}x_{jD}=1$ for every active $K_j$, $\sum_{D\in\mathcal D_{j,c_j}}x^\ast_{jD}=1$ for every active $\Psi_j$, and
\[
\prod_j\Big(\sum_{D\in\mathcal D_{j,c_j}}x_{jD}\Big)^{s_j}\Big(\sum_{D\in\mathcal D_{j,0}}x^\ast_{jD}\Big)^{t_j}=1,
\]
of dimension $n_{\mathcal D}$, the number of the coordinates less the number of the active letters less one, which is $2|\mathcal D|-2N-1$, $|\mathcal D|=\sum_j|\mathcal D_j|$, when every letter is active. The \emph{law} of the piece is the image $(\boldsymbol\varrho^\rho_B)_\ast\phi^\rho_B$, a positive measure carried by the seam torus. Its Mellin transform, at shares $(\mathbf a,\mathbf b)=(a_{jD},b_{jD})_{j,D}$ of the active letters,
\[
\int_{\Delta_B}(\boldsymbol\varrho^\rho_B)^{-(\mathbf a,\mathbf b)}\,d\phi^\rho_B=\int\mathbf x^{-(\mathbf a,\mathbf b)}\,d\big((\boldsymbol\varrho^\rho_B)_\ast\phi^\rho_B\big)(\mathbf x),\qquad\mathbf x^{-(\mathbf a,\mathbf b)}=\prod_{j,D}x_{jD}^{-a_{jD}}\,x^\ast_{jD}{}^{-b_{jD}},
\]
is the \emph{seam function}, in $[0,\infty]$ at real shares and defined at complex shares wherever it converges absolutely. A partner has its scale, scale form, seam vector, law and seam function by the same definition, with its room letters and the reversed excesses.
\end{definition}

The scale form is projective: its orthant form has the degree $m_B-\sum_j(s_jd_j+t_jd^\ast_j)-\kappa_B=0$, so it is read on any section (Lemma~\ref{lem:chart}). The seam points lie on the seam torus: the sums at the excesses $0$ and $c_j$ are $1$ by definition, and the product is $(\gamma^\rho_B)^{\kappa_B+\kappa^{\mathrm{top}}_B}(\gamma^\rho_B)^{-\sum_j(s_j+t_j)c_j}=1$. The scale of the piece is the geometric mean of the oriented extreme ratios of the active letters, each taken to the root of its width: $K_j$ measures the scale by $(V^{(c_j)}_j/V^{(0)}_j)^{1/c_j}$ and $\Psi_j$ by $(V^{\ast(0)}_j/V^{\ast(c_j)}_j)^{1/c_j}$, in the proportions $s_jc_j$ and $t_jc_j$; a letter of width zero does not see the scale. Any other positive scale $q\,\gamma^\rho_B$ of degree one, $q$ a positive function of the seam vector, multiplies $\phi^\rho_B$ by $q^{-\kappa_B}$ and the coordinates of the seam vector by powers of $q$: a fixed bijection of the tori with a density, so the law is determined up to it. The scale named is canonical: it uses the seam letters of the active letters alone, treats the two letters alike, and is kept at the centre by the bijections exchanging the family (Proposition~\ref{prop:seamsym}(iii)). With weights factoring at the extreme levels, $V^{(0)}_j=K_{\cB_j/R}$, $V^{(c_j)}_j=K_{\cB_j|B}$, $V^{\ast(0)}_j=\Psi_{\cB_j/R}$ and $V^{\ast(c_j)}_j=\Psi_{\cB_j|B}$ (Proposition~\ref{prop:seamidentity}(v)); so
\begin{gather*}
\gamma^\rho_B=\prod_j\Big(\frac{K_{\cB_j|B}}{K_{\cB_j/R}}\Big)^{s_j/(\kappa_B+\kappa^{\mathrm{top}}_B)}\Big(\frac{\Psi_{\cB_j/R}}{\Psi_{\cB_j|B}}\Big)^{t_j/(\kappa_B+\kappa^{\mathrm{top}}_B)},\\
\phi^\rho_B(du)=u^{\mvec_B-\one}\prod_jK_{\cB_j/R}^{-s_j}\Psi_{\cB_j|B}^{-t_j}\,(\gamma^\rho_B)^{-\kappa_B}\,d\omega_B(u).
\end{gather*}

The law is read off the measure of Theorem~\ref{thm:crossing}. For the room-station factorization the letter vector is $\mathbf V(u)=\big(K_{j,(D)}(u),\Psi_{j,(D)}(u)\big)_{j,D}$; the density of $\phi^\rho_B$ against $u^{\mvec_B-\one}d\omega_B$ and the seam vector are functions of $\mathbf V(u)$, so the law is an image of the measure $\mathbf V_\ast\big(u^{\mvec_B-\one}d\omega_B(u)\big)$ with a density. The seam vector forgets the seam letters of the inactive letters and, of the others, one constant for each active letter and the scale, $V_{jD}\mapsto\lambda^{i(D)}V_{jD}$ and $V^\ast_{jD}\mapsto\lambda^{c_j-i(D)}V^\ast_{jD}$: the directions that the pins of Lemma~\ref{lem:seamword} fix, on which the scale form changes by constant factors.

% Reference 20: prop:seamsym; source line 7393 (an exchange at the centre).
% Proved in the last part: prop:seamsym.
\begin{proposition}[Symmetries of a seam]\label{prop:seamsym}
Let a bijection $g$ of the marks $\cE$ of a family map $B$ onto $B'$ and $R$ onto $R'$.
\begin{enumerate}[label=(\roman*)]
\item If $g$ carries table $j$ to table $j'$ (Definition~\ref{def:carry}), it maps the patterns and the room sides of $B$ of each excess in table $j$ onto those of $B'$ of the same excess in table $j'$, stations onto stations and compatible pairs onto compatible pairs, and the station letters and seam letters of the image stations are those of the stations read at $g^\ast\alpha$, the normalization of the multipliers commuting with $g$.
\item If $g$ exchanges table $j$ with table $j'$, the windows of $B$ and $B'$ correspond as in Lemma~\ref{lem:windowturn}, and $p\mapsto B'\setminus g(p)$, $X\mapsto R'\setminus g(X)$ map the stations of $B$ of excess $i$ in table $j$ onto those of $B'$ of excess $c(B)-i$ in table $j'$ and compatible pairs onto compatible pairs, the station letter $\mathsf x_C$ and the seam letter $K_{(D)}$ going to the station letter $\mathsf x^\ast_{C'}$ and the seam letter $\Psi_{(D')}$ of the image stations $C'$, $D'$, and back.
\item So the bijections that keep $B$ and carry or exchange the family permute the stations of its seams, with their station letters and seam letters, up to the bijection. Let such a bijection also keep the masses. A carrying one maps the seam torus of $B$ onto itself, permuting the coordinates, and keeps the law $(\boldsymbol\varrho^\rho_B)_\ast\phi^\rho_B$. An exchanging one, sending table $j$ to table $j'$, maps the seam torus of $B$ at the lean $\boldsymbol\delta$ onto the seam torus of $B$ at the lean $-\boldsymbol\delta$, the coordinates $\varrho^\rho_{jD}$ of table $j$ going to the coordinates $\varrho^{\rho\ast}_{j'D'}$ of table $j'$ and back, and maps the law of $B$ at $\boldsymbol\delta$ onto the law of $B$ at $-\boldsymbol\delta$; at the centre it keeps the law. At the centre the potential is kept by the carrying and the exchanging maps alike, at a lean by the carrying maps (Proposition~\ref{prop:certify}): the centre does not tell apart stations joined by a bijection carrying or exchanging the family, nor a lean stations joined by a carrying one.
\end{enumerate}
\end{proposition}

At the centre two stations can lie in one orbit through an exchange and in different orbits at every lean: the centre sees them as one, a lean separates them. These are the confusable stations of the seam, the analogue for stations of the halves $C_K$, $C_\Psi$ of a class of marks (Theorem~\ref{thm:acrosscentre}); the seam carries them down the tree with the classes of marks (Theorem~\ref{thm:roomkeeps}).

% Reference 20: lem:seamword; source line 7445
\begin{lemma}[The seam function is the period of the seam word]\label{lem:seamword}
On the \emph{pins of the seam}, the letter totals $\sum_Da_{jD}=s_j$ for every active $K_j$ and $\sum_Db_{jD}=t_j$ for every active $\Psi_j$ and the \emph{scale pin}
\[
\sum_{j,D}\Big(i(D)\,a_{jD}+\big(c_j-i(D)\big)\,b_{jD}\Big)=\kappa_B,
\]
the seam function is the period at the masses $\mvec_B$ of the \emph{seam word} $W_B(\mathbf a,\mathbf b)=\prod_{j,D}V_{jD}^{-a_{jD}}V^\ast_{jD}{}^{-b_{jD}}$ (Definition~\ref{def:posword}):
\[
\int_{\Delta_B}(\boldsymbol\varrho^\rho_B)^{-(\mathbf a,\mathbf b)}\,d\phi^\rho_B=P_{W_B(\mathbf a,\mathbf b)}(\mvec_B)=\int_{\Delta_B}u^{\mvec_B-\one}\prod_{j,D}V_{jD}(u)^{-a_{jD}}V^\ast_{jD}(u)^{-b_{jD}}\,d\omega_B(u).
\]
\end{lemma}

For the piece of a bunch of a family the seam word is the bunch word of the room-station factorization (Theorem~\ref{thm:anywidth}), and its period on the pins is the Mellin transform of one law.

% Reference 20: lem:absorbwidth; source line 7456 (functions linear in the seam letters).
% Proved in the last part: lem:absorbwidth.
\begin{lemma}[Absorbing a piece of any width]\label{lem:absorbwidth}
Let a piece on $B$ with the stations $\mathcal D$ carry masses $\mvec_B$ and exponents $(s_j,t_j)_j$ with $\sum_j(s_j+t_j)c_j>0$, its wall $\kappa_B$ and its seam vector $\boldsymbol\varrho^\rho_B$, whose values are written $\mathbf x=(x_{jD},x^\ast_{jD})_{j,D}$. Let $g$ be a mark outside $B$, the \emph{stand-in} of the piece, $Z$ a finite set disjoint from $B\sqcup\{g\}$, $\mvec_Z$ masses on $Z$, and $f_1,\dots,f_p$ positive continuous functions, each assigned to one active letter of the piece, $f_l$ defined on $(0,\infty)^Z\times(0,\infty)^{\mathcal D_j}$ when its letter is $K_j$ or $\Psi_j$ and homogeneous of degree one in its second argument, $f_l(\alpha_Z,\lambda\mathbf y)=\lambda\,f_l(\alpha_Z,\mathbf y)$ for $\lambda>0$, with real exponents $a_l$ such that the exponents assigned to each letter add to its exponent. In this statement write $\mathbf V_l$ for the vector of the seam letters of the letter of $f_l$, $(V_{jD})_{D\in\mathcal D_j}$ for $K_j$ and $(V^\ast_{jD})_{D\in\mathcal D_j}$ for $\Psi_j$; $V_l$ for its extreme excess sum, $V^{(0)}_j$ for $K_j$ and $V^{\ast(c_j)}_j$ for $\Psi_j$; and $\mathbf x_l(\alpha_g)$ for the vector $(x_{jD}\,\alpha_g^{\,i(D)})_D$ for $K_j$ and $(x^\ast_{jD}\,\alpha_g^{\,c_j-i(D)})_D$ for $\Psi_j$. Let $H\ge0$ be a Borel function on $(0,\infty)^p$. If the form
\[
\omega=\alpha_B^{\mvec_B-\one}\alpha_Z^{\mvec_Z-\one}\,H\Big(\Big(\frac{f_l\big(\alpha_Z,\mathbf V_l(\alpha_B)\big)}{V_l(\alpha_B)}\Big)_l\Big)\prod_{l=1}^pf_l\big(\alpha_Z,\mathbf V_l(\alpha_B)\big)^{-a_l}\,d\alpha_B\,d\alpha_Z
\]
is homogeneous of degree zero on $(0,\infty)^{B\sqcup Z}$, then, in $[0,\infty]$,
\begin{gather*}
\int_{\{h=1\}}\omega=\int\Big(\int_{\{h'=1\}}\omega_{\mathbf x}\Big)\,d\big((\boldsymbol\varrho^\rho_B)_\ast\phi^\rho_B\big)(\mathbf x),\\
\omega_{\mathbf x}=\alpha_Z^{\mvec_Z-\one}\alpha_g^{\kappa_B-1}\,H\Big(\big(f_l(\alpha_Z,\mathbf x_l(\alpha_g))\big)_l\Big)\prod_lf_l\big(\alpha_Z,\mathbf x_l(\alpha_g)\big)^{-a_l}\,d\alpha_Z\,d\alpha_g,
\end{gather*}
the integrals taken over any sections of positive functions $h$, $h'$ of degree one on the two orthants.
\end{lemma}

The piece collapses to its stand-in $g$, which carries the mass $\kappa_B$; a letter of table $j$ that sees the piece sees the stand-in with the degree $c_j$, at its stations twisted by the seam point. At width one this is Lemma~\ref{lem:absorb} read through Proposition~\ref{prop:widthonelaw}.

% Reference 20: lem:nativestations; source line 10292 (the split at the stations of one bunch).
\begin{lemma}[The native split at the stations]\label{lem:nativestations}
Let a family of tables have real data on its balance wall, with $s_j,t_j\ge0$, not all zero, and $\mvec>0$, and a native gauge $\alpha^\star$ (Definition~\ref{def:native}), $B$ a bunch with $\sum_j(s_j+t_j)c_j(B)>0$, and put, for the active letters,
\[
a_{jD}=s_j\,\mathbb P_{j,\alpha^\star}(T\cap R\in D),\qquad b_{jD}=t_j\,\mathbb P^\ast_{j,\alpha^\star}(T\cap R\in D),\qquad D\in\mathcal D_j,
\]
the native split (Definition~\ref{def:split}) aggregated to the stations of room sides. Then $(\mathbf a,\mathbf b)$ is a point of the pins of the seam with positive coordinates, at which $\mvec_B$ and $\mvec_R$ lie in the relative interiors of the Newton sums of the bunch word and of the room word (Theorem~\ref{thm:anywidth}); the two words converge there exactly when these Newton sums have the full dimensions $|B|-1$ and $|R|-1$ (Theorem~\ref{thm:poswordconv}).
\end{lemma}

% Reference 20: thm:gluewidth; source line 7472
\begin{theorem}[The glue across a seam]\label{thm:gluewidth}
Let a piece on $B$ and a partner on $R$ with the stations $\mathcal D$ be glued, with masses on the balance wall of the glue and $\sum_j(s_j+t_j)c_j>0$.
\begin{enumerate}[label=(\roman*)]
\item \emph{The piece absorbed.} In $[0,\infty]$,
\[
P_{B\sqcup R}(\mvec)=\int P_{R^{(\mathbf x)}}(\mvec_R,\kappa_B)\,d\big((\boldsymbol\varrho^\rho_B)_\ast\phi^\rho_B\big)(\mathbf x),
\]
with $R^{(\mathbf x)}$ the \emph{partner with its stand-in at} $\mathbf x$: the family of Newton words on $R\sqcup\{g\}$ with the letters $\sum_Dx_{jD}\,\alpha_g^{\,i(D)}L_{jD}$ and $\sum_Dx^\ast_{jD}\,\alpha_g^{\,c_j-i(D)}L^\ast_{jD}$, the exponents $(s_j,t_j)$ and the mass $\kappa_B$ on $g$.
\item \emph{Both absorbed.} On the scales $\{g\}$ and one scale for each active letter let $t_{\mathcal D}$ be the configuration whose column at the coordinate $(j,D)$ of $K_j$ is $i(D)\,\mathbf e_g+\mathbf e_{K_j}$ and at that of $\Psi_j$ is $(c_j-i(D))\,\mathbf e_g+\mathbf e_{\Psi_j}$. In $[0,\infty]$,
\[
P_{B\sqcup R}(\mvec)=\int\!\!\int\Gamma(\kappa_B)\,\sZ_{(\kappa_B,\mathbf s,\mathbf t);t_{\mathcal D}}(\mathbf x\circ\mathbf x')\,d\big((\boldsymbol\varrho^\rho_B)_\ast\phi^\rho_B\big)(\mathbf x)\,d\big((\boldsymbol\varrho^\rho_R)_\ast\phi^\rho_R\big)(\mathbf x'),
\]
$\mathbf x\circ\mathbf x'$ the coordinatewise product and the index $(\kappa_B,\mathbf s,\mathbf t)$ read on the active letters (Definition~\ref{def:zamma}). The kernel is the period of the \emph{stand-in pair}, the marks $g$ and $g'$ with the masses $\kappa_B$ and $\kappa_R$ and the letters $\sum_Dx_{jD}x'_{jD}\,\alpha_g^{i(D)}\alpha_{g'}^{c_j-i(D)}$ and $\sum_Dx^\ast_{jD}x'^\ast_{jD}\,\alpha_g^{c_j-i(D)}\alpha_{g'}^{i(D)}$, read on one scale:
\[
\Gamma(\kappa)\,\sZ_{(\kappa,\mathbf s,\mathbf t);t_{\mathcal D}}(\mathbf x\circ\mathbf x')=\int_0^\infty v^{\kappa-1}\prod_j\Big(\sum_{D\in\mathcal D_j}x_{jD}x'_{jD}\,v^{i(D)}\Big)^{-s_j}\Big(\sum_{D\in\mathcal D_j}x^\ast_{jD}x'^\ast_{jD}\,v^{c_j-i(D)}\Big)^{-t_j}dv .
\]
\item \emph{One share per station.} On every contour of shares $a_{jD}$, $b_{jD}$ of the active letters with positive real parts, on the pins of the seam (Lemma~\ref{lem:seamword}), which on the balance wall are the pins of the partner as well, at which the seam functions of the piece and of the partner converge,
\[
P_{B\sqcup R}(\mvec)=\int_{\text{pins}}\prod_{j:\,s_j>0}\Gst{|\mathcal D_j|}1{\mathbf a_j}{s_j}\prod_{j:\,t_j>0}\Gst{|\mathcal D_j|}1{\mathbf b_j}{t_j}\int_{\Delta_B}(\boldsymbol\varrho^\rho_B)^{-(\mathbf a,\mathbf b)}\,d\phi^\rho_B\int_{\Delta_R}(\boldsymbol\varrho^\rho_R)^{-(\mathbf a,\mathbf b)}\,d\phi^\rho_R,
\]
a fibre integral of the dimension $n_{\mathcal D}$ of the seam torus: the two seam functions are evaluated at the same shares, one per station and letter.
\end{enumerate}
Parts (i) and (ii) hold in every case; (iii) needs a point of the pins at which both seam functions are finite. Full dimension of the Newton sums of the seam word and of the room word is necessary for it (Theorem~\ref{thm:poswordconv}), and when the glue is a family of tables with a native gauge, every proper wall positive, it is sufficient, the native split at the stations being such a point (Lemma~\ref{lem:nativestations}); without it the share side carries one pin more for each direction that a side does not see (Theorem~\ref{thm:transversepins} below).
\end{theorem}

The glue is one pairing at every width: the bunch is a law on its seam torus, the room another, and the kernel is the period of two stand-in marks whose letters pair the stations. At width one the seam torus has one dimension less than the number of active letters, a half-line for one table with $s,t>0$ and $(0,\infty)^{2N-1}$ when every letter is active (Proposition~\ref{prop:widthonelaw}), and the kernel is the period of two marks in every table (Lemma~\ref{lem:twomarks}, Theorem~\ref{thm:node}(i)); at the width $c$ the seam torus has the dimension $n_{\mathcal D}$ and the kernel is one scale with at most $c_j+1$ levels in the letters of table $j$, exactly $c_j+1$ for a table. Opened on its shares, the glue is the product of the two seam functions against the share kernels of Corollary~\ref{thm:E1}, one for each active letter.

% Reference 20: ex:widthone; source line 14535 (stated as a proposition, at the scope of a family; the reference reads (iv) for one table, its cells carrying the weights of the room children).
% Proved in the last part: prop:widthonelaw.
\begin{proposition}[Width one is the hinge]\label{prop:widthonelaw}
Let $B$ have width one in every table of a family with exponents $s_j,t_j>0$, the weights factoring across $B$ and $R$, $\theta_{j,T}=\theta'_{j,T\cap B}\,\theta''_{j,T\cap R}$, with $\theta''_j$ of largest value $1$ on the room sides of each excess, so that $\theta''_j$ gives the multipliers of the room sides (Definition~\ref{def:seam}), with the stations $D_0$ and $D_1$ of the excesses $0$ and $1$ in every table (Proposition~\ref{prop:ladder}(ii)), the hinge values $\gamma^K_{j,B}$, $\gamma^\Psi_{j,B}$ of Definition~\ref{def:readings}, and the cell family $F^+_B=((\cB^+_{j,B},\theta'_j);s_j,t_j)_j$ of the bunch (Definition~\ref{def:hinged}), with its hinge wall $\kappa_B$, base form $\phi_{F^+_B}$ and ratio vector $\boldsymbol\varrho_{F^+_B}$, whose values are written $\mathbf x=(x_{j,K},x_{j,\Psi})_j$, $x_{1,K}=1$. In this statement write $\psi(\mathbf x)=\prod_jx_{j,K}^{\,s_j/(2\sigma)}x_{j,\Psi}^{\,t_j/(2\sigma)}$, $\sigma=\sum_j\sigma_j$ the total strength (Definition~\ref{def:family}).
\begin{enumerate}[label=(\roman*)]
\item $\gamma^\rho_B=\prod_j(\gamma^K_{j,B})^{s_j/(2\sigma)}(\gamma^\Psi_{j,B})^{t_j/(2\sigma)}=\gamma^K_{1,B}\,\psi(\boldsymbol\varrho_{F^+_B})$ and $\phi^\rho_B=\psi(\boldsymbol\varrho_{F^+_B})^{-\kappa_B}\,\phi_{F^+_B}$.
\item The seam vector of the piece has $\varrho^\rho_{jD_0}=\varrho^{\rho\ast}_{jD_1}=1$, $\varrho^\rho_{jD_1}=x_{j,K}/\psi(\mathbf x)$ and $\varrho^{\rho\ast}_{jD_0}=x_{j,\Psi}/\psi(\mathbf x)$ at $\mathbf x=\boldsymbol\varrho_{F^+_B}$; the map $\mathbf x\mapsto\big(x_{j,K}/\psi(\mathbf x),\,x_{j,\Psi}/\psi(\mathbf x)\big)_j$ is a bijection of $(0,\infty)^{2N-1}$ onto the seam torus, and under it the law $(\boldsymbol\varrho^\rho_B)_\ast\phi^\rho_B$ is the image of $\psi^{-\kappa_B}$ times the law $(\boldsymbol\varrho_{F^+_B})_\ast\phi_{F^+_B}$ of the cell family.
\item The partner with its stand-in at the seam point of $\mathbf x$ is, after $\alpha_g\mapsto\psi(\mathbf x)\,\alpha_g$, the family $((\cB^+_{j,R},\theta''_j);s_j,t_j)_j$ of the cells of the room twisted at $g$ by $\mathbf x$, $x_{j,K}$ in $K_j$ and $x_{j,\Psi}$ in $\Psi_j$ (Definition~\ref{def:twist}); the change of variable multiplies $\alpha_g^{\kappa_B-1}d\alpha_g$ by $\psi(\mathbf x)^{\kappa_B}$, which cancels the density of (ii), and Lemma~\ref{lem:absorbwidth} is Lemma~\ref{lem:absorb} for the cell family $F^+_B$.
\item The room, read as the partner, has the same description through the cell family $F^+_R=((\cB^+_{j,R},\theta''_j);s_j,t_j)_j$ of the room, with its hinge wall $\kappa_R$ and ratio vector $\boldsymbol\varrho_{F^+_R}$, whose values are written $\mathbf x'$, $x'_{1,K}=1$: its scale is $\gamma^K_{1,R}\,\psi(\boldsymbol\varrho_{F^+_R})$, its scale form $\psi(\boldsymbol\varrho_{F^+_R})^{-\kappa_R}\phi_{F^+_R}$, its seam vector has $\varrho^\rho_{jD_1}=\varrho^{\rho\ast}_{jD_0}=1$, $\varrho^\rho_{jD_0}=x'_{j,K}/\psi(\mathbf x')$ and $\varrho^{\rho\ast}_{jD_1}=x'_{j,\Psi}/\psi(\mathbf x')$ at $\mathbf x'=\boldsymbol\varrho_{F^+_R}$, and its law is the image of $\psi^{-\kappa_R}$ times the law $(\boldsymbol\varrho_{F^+_R})_\ast\phi_{F^+_R}$ under the bijection of (ii). At the seam points the coordinate ratios $\varrho^{\rho\ast}_{jD_0}/\varrho^\rho_{jD_1}$ of the piece and $\varrho^{\rho\ast}_{jD_1}/\varrho^\rho_{jD_0}$ of the partner are the ratios $\varrho_{j,B}$ and $\varrho_{j,R}$ of Definition~\ref{def:readings}. At the seam points $\mathbf p$ of the piece and $\mathbf p'$ of the partner that are the images of $\mathbf x$ and $\mathbf x'$,
\begin{align*}
\Gamma(\kappa_B)\,\sZ_{(\kappa_B,\mathbf s,\mathbf t);t_{\mathcal D}}(\mathbf p\circ\mathbf p')&=\int_0^\infty v^{\kappa_B-1}\prod_j\Big(\frac{x_{j,K}\,v}{\psi(\mathbf x)}+\frac{x'_{j,K}}{\psi(\mathbf x')}\Big)^{-s_j}\Big(\frac{x_{j,\Psi}\,v}{\psi(\mathbf x)}+\frac{x'_{j,\Psi}}{\psi(\mathbf x')}\Big)^{-t_j}dv\\
&=\psi(\mathbf x)^{\kappa_B}\psi(\mathbf x')^{\kappa_R}\,P_{F_{\mathcal H}[(\mathbf x,\mathbf x')]}\big((\kappa_B,\kappa_R);\mathbf s,\mathbf t\big),
\end{align*}
$F_{\mathcal H}[(\mathbf x,\mathbf x')]$ the core $\binom{\{g,g'\}}1$ in every table twisted at $g$ by $\mathbf x$ and at $g'$ by $\mathbf x'$, with the masses $\kappa_B$ and $\kappa_R$ (Theorem~\ref{thm:node}(i)). So Theorem~\ref{thm:gluewidth}(i) is Lemma~\ref{lem:absorb} at the hinge, by part (iii); Theorem~\ref{thm:gluewidth}(ii) is Theorem~\ref{thm:node}(i) with that core and the cell families $F^+_B$ at $g$ and $F^+_R$ at $g'$, the factor $\psi(\mathbf x)^{\kappa_B}\psi(\mathbf x')^{\kappa_R}$ cancelling the densities of the two laws; and Theorem~\ref{thm:gluewidth}(iii) is Theorem~\ref{thm:twosides}.
\end{enumerate}
\end{proposition}

At width one the scale of the piece is the geometric mean of the hinge values in the proportions of the exponents, and the general law is the law of the cell family of the bunch with a fixed power of its ratio vector as density. Every statement about pieces therefore specializes to its statement about cells through this one fixed change, and the calculus of cells (Chapter~\ref{ch:hinged}) is the case of width one of the calculus of pieces.

% Reference 20: thm:transversepins; source line 7495
\begin{theorem}[The share side with transverse pins]\label{thm:transversepins}
In Theorem~\ref{thm:gluewidth}(iii), the period of the glue converging, so that its Newton sum has full dimension (Theorem~\ref{thm:poswordconv}), let $\vec{\mathcal N}_B\subseteq\mathbb R^B$ and $\vec{\mathcal N}_R\subseteq\mathbb R^R$ be the spaces of the directions of the Newton sums of the seam word of the piece and of the room word of the partner at shares with positive real parts, which do not depend on the shares, and $M=\vec{\mathcal N}_B^\perp\oplus\vec{\mathcal N}_R^\perp\subseteq\mathbb R^{B\sqcup R}$, which contains $\one$; let $\mathbf v_1,\dots,\mathbf v_q$ complete $\one$ to a basis of the lattice $M\cap\mathbb Z^{B\sqcup R}$. For every active letter and station fix the exponents $\mathbf n_{jD}$, $\mathbf n^\ast_{jD}$, $\mathbf n'_{jD}$ and $\mathbf n'^\ast_{jD}$ of one monomial of each of $V_{jD}$, $V^\ast_{jD}$, $L_{jD}$ and $L^\ast_{jD}$, and in this statement write
\[
\mathbf n(\mathbf a,\mathbf b)=\sum_{j,D}\Big(a_{jD}\big(\mathbf n_{jD}+\mathbf n'_{jD}\big)+b_{jD}\big(\mathbf n^\ast_{jD}+\mathbf n'^\ast_{jD}\big)\Big),
\]
which modulo $\vec{\mathcal N}_B\oplus\vec{\mathcal N}_R$ does not depend on the monomials fixed. For a word $W$ on $B$ whose Newton sum has the space of directions $\vec{\mathcal N}_B$, its \emph{reduced period} is
\[
P^\flat_W(\mvec_B)=\int_{\mathbb R^B/\vec{\mathcal N}_B^\perp}e^{\ip{\mvec_B}y}\,W(e^y)\,dy,
\]
$dy$ the quotient of Lebesgue measure by the measure on $\vec{\mathcal N}_B^\perp$ that gives its integer lattice the covolume one; likewise on $R$. On every contour of shares of the active letters with positive real parts, pinned by the letter totals and by $\ip{\mathbf n(\mathbf a,\mathbf b)}{\mathbf v_l}=\ip{\mvec}{\mathbf v_l}$, $l\le q$, pins that are independent by the full dimension, at which the reduced periods of the two words converge,
\[
P_{B\sqcup R}(\mvec)=\int_{\text{pins}}\prod_{j:\,s_j>0}\Gst{|\mathcal D_j|}1{\mathbf a_j}{s_j}\prod_{j:\,t_j>0}\Gst{|\mathcal D_j|}1{\mathbf b_j}{t_j}\ P^\flat_{W_B(\mathbf a,\mathbf b)}(\mvec_B)\,P^\flat_{W_R(\mathbf a,\mathbf b)}(\mvec_R),
\]
$W_R(\mathbf a,\mathbf b)=\prod_{j,D}L_{jD}^{-a_{jD}}L^\ast_{jD}{}^{-b_{jD}}$ the room word of the partner; on the pins the integrands of the reduced periods are invariant under $\vec{\mathcal N}_B^\perp$ and $\vec{\mathcal N}_R^\perp$. When $\vec{\mathcal N}_B$ and $\vec{\mathcal N}_R$ are the hyperplanes of sum zero, $q=1$, the pin is the scale pin, the reduced periods are the periods, and this is Theorem~\ref{thm:gluewidth}(iii). When every room letter of the partner is a monomial, as for the room of a table whose stations of room sides are single room sides, with the coefficients $1$, $\vec{\mathcal N}_R=0$, the pins on $R$ say that the exponent of the partner's monomials is $\mvec_R$, and the reduced period of the partner is the product of the coefficients of these monomials raised to minus the shares.
\end{theorem}

So a degenerate seam word or room word costs pins, not convergence: the directions that a side does not see are fixed by the masses, one pin each, and the integral over them is gone.

\section{The Gramma of a set of pairs}\label{sec:relgramma}

The station letters are what the bunch is; the seam letters are what the room sees. The seam joins the two, and its compatible pairs enter the period through one Gramma.

% Reference 20: def:relgramma; source line 7520
\begin{definition}[The Gramma of a set of pairs]\label{def:relgramma}
Let $\mathfrak R\subseteq\mathcal K\times\mathcal D$ be a set of pairs of stations of two finite sets, every station in some pair, with coefficients $h_{CD}>0$ on its pairs. Its \emph{station map} sends $\mathbf e=(e_{CD})\in\mathbb C^{\mathfrak R}$ to $\partial\mathbf e$, with $(\partial\mathbf e)_C=\sum_De_{CD}$ and $(\partial\mathbf e)_D=\sum_Ce_{CD}$. Two pairs lie in one \emph{part} when a chain of pairs, each sharing a station with the next, joins them; a \emph{circuit} is a minimal set of pairs whose columns $\mathbf e_C+\mathbf e_D$ are linearly dependent, which is the set of pairs of a closed chain $C_1D_1C_2D_2\cdots C_mD_mC_1$, $m\ge2$, of distinct stations with $(C_l,D_l)$ and $(C_{l+1},D_l)$ in $\mathfrak R$. Choose one station of $\mathcal D$ in each part, and call the other stations \emph{kept}. For \emph{loads} $(\mathbf a,\mathbf a')\in\mathbb C^{\mathcal K}\oplus\mathbb C^{\mathcal D}$ whose totals on the two sides of every part agree, the \emph{Gramma of $\mathfrak R$} is the Gramma $\sG_{(\mathbf a,\mathbf a');\partial\mathfrak R}(\mathbf h)$ (Definition~\ref{def:gramma}) on the scales of the kept stations whose configuration $\partial\mathfrak R$ has at the pair $(C,D)$ the column $\mathbf e_C+\mathbf e_D$, a station that is not kept read as $0$, whose index is the loads of the kept stations, and whose coefficients are the $h_{CD}$. By Proposition~\ref{prop:shareform} it is the fibre integral
\[
\sG_{(\mathbf a,\mathbf a');\partial\mathfrak R}(\mathbf h)=\int_{\partial\mathbf e=(\mathbf a,\mathbf a')}\ \prod_{(C,D)\in\mathfrak R}\Gamma(e_{CD})\,h_{CD}^{-e_{CD}},
\]
pinned by the sums at the kept stations, which imply the others. Another choice of the stations not kept changes the pins by an integer matrix of determinant $\pm1$, which leaves the fibre integral unchanged; its dimension is $|\mathfrak R|-|\mathcal K|-|\mathcal D|$ plus the number of parts. The \emph{seam Gramma} of the excess $i$ is the Gramma of the set $\mathfrak R_i$ of the compatible pairs of $H_i$, with the pair constants as its coefficients.
\end{definition}

% Reference 20: prop:seamtransfer; source line 7528 (letter by letter, at the scope of a family)
\begin{proposition}[The seam transfers the station loads of the bunch to those of the room]\label{prop:seamtransfer}
For a bunch of a family let the \emph{station word} of the bunch be $\prod_{j,C}\mathsf x_{jC}^{-a_{jC}}\mathsf x^\ast_{jC}{}^{-b_{jC}}$, with the station loads $a_{jC}$ of the active $K_j$ and $b_{jC}$ of the active $\Psi_j$ on the pins $\sum_Ca_{jC}=s_j$, $\sum_Cb_{jC}=t_j$ and $\sum_{j,C}\big(i(C)a_{jC}+(c_j-i(C))b_{jC}\big)=\kappa_B$, and let the station loads $a_{jD}$ and $b_{jD}$ of the room be shares on the pins of the seam (Lemma~\ref{lem:seamword}). On contours of positive real parts at which the station word of the bunch converges,
\begin{multline*}
\prod_{j:\,s_j>0}\prod_{D\in\mathcal D_j}\Gamma(a_{jD})\prod_{j:\,t_j>0}\prod_{D\in\mathcal D_j}\Gamma(b_{jD})\ P_{W_B(\mathbf a,\mathbf b)}(\mvec_B)\\
=\int\prod_{j:\,s_j>0}\prod_{i=0}^{c_j}\sG_{(\mathbf a_{j,\mathcal K_i},\mathbf a_{j,\mathcal D_i});\partial\mathfrak R_{j,i}}(\mathbf h_j)\prod_{j:\,t_j>0}\prod_{i=0}^{c_j}\sG_{(\mathbf b_{j,\mathcal K_i},\mathbf b_{j,\mathcal D_i});\partial\mathfrak R_{j,i}}(\mathbf h_j)\ P_{\prod_{j,C}\mathsf x_{jC}^{-a_{jC}}\mathsf x^\ast_{jC}{}^{-b_{jC}}}(\mvec_B),
\end{multline*}
the integral over the station loads of the bunch with equal totals on the two sides of every part of every set $\mathfrak R_{j,i}$ of the compatible pairs of table $j$ at the excess $i$, $\mathbf h_j$ its pair constants.
\end{proposition}

Along a circuit the pair terms $h_{CD}\,\mathsf x_C\,\mathsf y_D$ of a level have the alternating product $\prod_lh_{C_lD_l}/h_{C_{l+1}D_l}$, a constant: a direction $\mathbf e$ of the pair shares with $\partial\mathbf e=0$ moves neither the station loads of the bunch, which the bunch sees, nor those of the room, which the room sees, and the period integrates over it through the seam Gramma alone. A part without circuits has a fibre of dimension zero, and its Gramma is the product $\prod\Gamma(e_{CD})h_{CD}^{-e_{CD}}$ at the one solution of $\partial\mathbf e=(\mathbf a,\mathbf a')$ (Corollary~\ref{cor:square}).

% Reference 20: prop:relgrammaforeman; source line 7539 (one part).
% Proved in the last part: prop:relgrammaforeman.
\begin{proposition}[The Gramma of a set of pairs in the Foreman's order]\label{prop:relgrammaforeman}
Let $\mathfrak R$ have the parts $\mathfrak R_1,\dots,\mathfrak R_\nu$, the part $\mathfrak R_l$ on the stations $\mathcal K_l\sqcup\mathcal D_l$, and loads with $|\mathbf a_{\mathcal K_l}|=|\mathbf a'_{\mathcal D_l}|$ for every $l$; in this statement write $L_C(v)=\sum_{D:\,(C,D)\in\mathfrak R}h_{CD}\,v_D$ for $C\in\mathcal K_l$ and $v\in\Delta_{\mathcal D_l}$. Then
\begin{multline*}
\sG_{(\mathbf a,\mathbf a');\partial\mathfrak R}(\mathbf h)=\prod_{C\in\mathcal K}\Gamma(a_C)\prod_{l=1}^{\nu}\int_{\Delta_{\mathcal D_l}}v^{\mathbf a'_{\mathcal D_l}-\one}\prod_{C\in\mathcal K_l}L_C(v)^{-a_C}\,d\omega_{\mathcal D_l}(v)\\
=\prod_{C\in\mathcal K}\Gamma(a_C)\prod_{l=1}^{\nu}P_{\prod_{C\in\mathcal K_l}L_C^{-a_C}}(\mathbf a'_{\mathcal D_l}),
\end{multline*}
both sides converging together, the period of the part $l$ being that of the word $\prod_{C\in\mathcal K_l}L_C^{-a_C}$ on $\mathcal D_l$ at the masses $\mathbf a'_{\mathcal D_l}$; with $\re a_C>0$ it converges exactly when, for every $l$, $\re\mathbf a'_{\mathcal D_l}$ lies in the relative interior of $\sum_{C\in\mathcal K_l}\re a_C\,\operatorname{conv}\{\mathbf e_D:\ (C,D)\in\mathfrak R\}$ (Theorem~\ref{thm:poswordconv}).
\end{proposition}

The stations of the bunch are integrated first, each on its half-line, and those of the room last, on the simplex of each part: the Foreman's order of Definition~\ref{def:procedure}, the bunch inside the integral over its room. For a seam the inner step returns, at each station $C$ of the bunch, its child letter $K_{(C)}=\sum_Dh_{CD}\,\mathsf y_D$ with every room station letter $\mathsf y_D$ replaced by a scale $v_D$. No station is singled out and no order is chosen.

% Reference 20: prop:relgrammaclose; source line 7552 (one part, real loads).
% Proved in the last part: prop:relgrammaclose.
\begin{proposition}[Closing the Gramma of a set of pairs]\label{prop:relgrammaclose}
In Proposition~\ref{prop:relgrammaforeman} let the loads be real and positive in the domain for (i) and (iii); in this statement write, for a part $\mathfrak R_l$, $L_C(v)=\sum_{D:\,(C,D)\in\mathfrak R}h_{CD}\,v_D$ for $C\in\mathcal K_l$ and $v\in\Delta_{\mathcal D_l}$, $M_C=\max_Dh_{CD}$ for its largest coefficient, its largest value on the simplex, and $\sG_{(\mathbf a,\mathbf a')}$ for $\sG_{(\mathbf a,\mathbf a');\partial\mathfrak R}(\mathbf h)$.
\begin{enumerate}[label=(\roman*)]
\item \emph{Peak closing.} For every part $\mathfrak R_l$, with $\bar L_C=1-L_C/M_C\ge0$ on $\Delta_{\mathcal D_l}$,
\begin{multline*}
\int_{\Delta_{\mathcal D_l}}v^{\mathbf a'_{\mathcal D_l}-\one}\prod_{C\in\mathcal K_l}L_C(v)^{-a_C}\,d\omega_{\mathcal D_l}(v)\\
=\prod_{C\in\mathcal K_l}M_C^{-a_C}\,\Gst{|\mathcal D_l|}1{\mathbf a'_{\mathcal D_l}}{|\mathbf a'_{\mathcal D_l}|}\sum_{\mathbf n\in\mathbb N^{\mathcal K_l}}\prod_{C\in\mathcal K_l}\frac{(a_C)_{n_C}}{n_C!}\ \oE_{V\sim\Dir_{\mathbf a'_{\mathcal D_l}}}\Big[\prod_{C\in\mathcal K_l}\bar L_C(V)^{n_C}\Big],
\end{multline*}
a series of nonnegative terms, each a finite nonnegative combination of Dirichlet moments (Theorem~\ref{thm:dirmoments}), rational at rational loads and coefficients; it is Theorem~\ref{thm:peak}(i) with the sharp references $M_C$, and $\sG_{(\mathbf a,\mathbf a')}$ is $\prod_C\Gamma(a_C)$ times the product of these right sides over the parts.
\item \emph{Relations at shifted loads.} At all loads in the domain, real or complex, for every $C$ and $D$,
\begin{gather*}
a_C\,\sG_{(\mathbf a,\mathbf a')}=\sum_{D':\,(C,D')\in\mathfrak R}h_{CD'}\,\sG_{(\mathbf a+\mathbf e_C,\,\mathbf a'+\mathbf e_{D'})},\\
a'_D\,\sG_{(\mathbf a,\mathbf a')}=\sum_{C':\,(C',D)\in\mathfrak R}h_{C'D}\,\sG_{(\mathbf a+\mathbf e_{C'},\,\mathbf a'+\mathbf e_D)}.
\end{gather*}
\item \emph{Residues.} A part that is one circuit $C_1D_1\cdots C_mD_m$ has one share $z$ moving round the circuit with alternating signs: for a real point $\mathbf q$ of the fibre the pairs $(C_l,D_l)$ carry $q_l+z$ and the pairs $(C_{l+1},D_l)$ carry $q'_l-z$. With $x=\prod_lh_{C_lD_l}/h_{C_{l+1}D_l}$ and $\Sigma=\sum_l(q_l+q'_l)$, the total load of the part, its Gramma is a constant times
\[
\frac1{2\pi i}\int_{(0)}x^{-z}\prod_{l=1}^m\Gamma(q_l+z)\,\Gamma(q'_l-z)\,dz .
\]
If $x<1$, or $x=1$ and $\Sigma<m$, the contour moves to the left across the poles of the $\Gamma(q_l+z)$, and if $x>1$, or $x=1$ and $\Sigma<m$, to the right across those of the $\Gamma(q'_l-z)$ (Theorem~\ref{thm:closing}): the Gramma is the limit of the sums of the residues at the poles crossed, a series over $m$ lattices, absolutely convergent when $x\ne1$ or $\Sigma<m-1$. The residue at a simple pole is a product of the other Gammas at the shifted point, and where poles of several factors meet, the $q_l$ being congruent modulo $\mathbb Z$, it is a Laurent coefficient. At $x=1$ with $\Sigma\ge m$ the shifted contours need not tend to zero and the residues need not sum.
\end{enumerate}
\end{proposition}

% Reference 20: thm:blockshares; source line 7576
\begin{theorem}[Shares inside the rectangles close]\label{thm:blockshares}
Fix a table and an excess $i$, and give every basis $T$ of excess $i$ a share $a_T$. For pattern loads $\mathbf a^B=(a^B_p)_p$ and room-side loads $\mathbf a^R=(a^R_X)_X$ whose station loads $a^B_C=\sum_{p\in C}a^B_p$ and $a^R_D=\sum_{X\in D}a^R_X$ have equal totals on every part of $\mathfrak R_i$, on contours of positive real parts,
\begin{multline*}
\int_{\sum_Xa_{p\sqcup X}=a^B_p,\ \sum_pa_{p\sqcup X}=a^R_X}\ \prod_T\Gamma(a_T)\,\theta_T^{-a_T}\\
=\prod_{C\in\mathcal K_i}\Gst{|C|}1{\mathbf a^B|_C}{a^B_C}\prod_pw_p^{-a^B_p}\prod_{D\in\mathcal D_i}\Gst{|D|}1{\mathbf a^R|_D}{a^R_D}\prod_Xw'_X{}^{-a^R_X}\ \sG_{(\mathbf a^B_{\mathcal K_i},\mathbf a^R_{\mathcal D_i});\partial\mathfrak R_i}(\mathbf h),
\end{multline*}
$\mathbf a^B_{\mathcal K_i}$ and $\mathbf a^R_{\mathcal D_i}$ the vectors of the station loads, the fibre on the left pinned by these sums at all patterns and at all room sides but one in each part of $\mathfrak R_i$, which imply the others.
\end{theorem}

Inside the rectangle of a compatible pair the shares spread freely and close to stacked Gammas, one for each station; what does not close is how the stations share one another, the seam Gramma.

\section{What a piece retains}\label{sec:retainwidth}

% Reference 20: def:tinkers; source line 10050
\begin{definition}[Tinkers]\label{def:tinkers}
A \emph{tinker} of a partner $\mathcal Y$ with the stations $\mathcal D$ is a vector $\mathbf v=(v_{jD},v^\ast_{jD})_{j,D}$ of positive numbers, one for each coordinate of the seam vector; the \emph{tinkered partner} $\mathcal Y_{\mathbf v}$ has the room letters $v_{jD}L_{jD}$ and $v^\ast_{jD}L^\ast_{jD}$. A tinker is \emph{plain} when $v_{jD}=v^\ast_{jD}$ for every coordinate. On the room of a table a plain tinker multiplies the weights of the room sides of each station by one constant, which keeps the bases and gives a table (Proposition~\ref{prop:tinkers}, read for the bunch $R$); a tinker that is not plain gives the room two systems of weights, one for each letter, as a cell may carry (Definition~\ref{def:substitution}).
\end{definition}

% Reference 20: thm:retainwidth; source line 10054 (the plain tinkers when every $t_j=0$ or every $s_j=0$, and for one table at width one).
% Proved in the last part: thm:retainwidth.
\begin{theorem}[The law is what a piece retains, at every width]\label{thm:retainwidth}
Let $\mathcal X$ and $\mathcal X'$ be pieces with the stations $\mathcal D$, on sets $B$ and $B'$ of any sizes, with the exponents $(s_j,t_j)_j$, $\sum_j(s_j+t_j)c_j>0$, and one wall $\kappa$. Let $\mathcal Y_0$ be a partner of both with the wall $\sum_j(s_j+t_j)c_j-\kappa$, and let some real point $\mathbf c$ of the pins of the seam, with positive coordinates, lie in the interior, relative to the pins, of the domains of the seam functions of $\mathcal X$, $\mathcal X'$ and $\mathcal Y_0$. The following are equivalent.
\begin{enumerate}[label=(\alph*)]
\item $(\boldsymbol\varrho^\rho_{\mathcal X})_\ast\phi^\rho_{\mathcal X}=(\boldsymbol\varrho^\rho_{\mathcal X'})_\ast\phi^\rho_{\mathcal X'}$.
\item $P_{\mathcal X\sqcup\mathcal Y}=P_{\mathcal X'\sqcup\mathcal Y}$, in $[0,\infty]$, for every partner $\mathcal Y$ with the wall $\sum_j(s_j+t_j)c_j-\kappa$.
\item $P_{\mathcal X\sqcup\mathcal Y_{0,\mathbf v}}=P_{\mathcal X'\sqcup\mathcal Y_{0,\mathbf v}}$ for every tinker $\mathbf v$ of $\mathcal Y_0$.
\item The seam functions of $\mathcal X$ and $\mathcal X'$ agree on the contour of the pins through $\mathbf c$.
\end{enumerate}
When every table $j$ with both letters active has $|\mathcal D_j|=1$, except at most one, which has $c_j\ge1$ and $|\mathcal D_j|=2$, the plain tinkers suffice in (c). For one table at width one with factoring weights, $|\mathcal D|=2$, and the theorem is Theorem~\ref{thm:retained} read through Proposition~\ref{prop:widthonelaw}: the one-mark cell of ratio $1$ (Definition~\ref{def:onemark}), glued at the hinge, is a partner whose plain tinkers $(1,q^{1/2})$ are the one-mark cells of ratio $q$.
\end{theorem}

So the wall and the law are the whole of what a piece contributes to the families glued on it, and the tinkers of one partner already separate any two laws. A plain tinker multiplies the two room letters of a station by one constant, so with every letter active, and some table of positive width, the plain tinkers of one partner move only $|\mathcal D|-N$ of the $2|\mathcal D|-2N-1$ directions of the seam torus. The other directions are reached by partners with one system of weights whose weights and masses vary.

% Reference 20: def:plainpartner; source line 10075
\begin{definition}[Plain partners]\label{def:plainpartner}
A partner on $R$ with the stations $\mathcal D$ is \emph{plain} when for every table $j$ with an active letter and every $D\in\mathcal D_j$ there are a nonempty family $\mathcal S_{jD}$ of subsets of $R$ of $e_j+c_j-i(D)$ marks, its \emph{room sides}, and positive weights $w'_{jD,X}$, $X\in\mathcal S_{jD}$, with
\[
L_{jD}=\sum_{X\in\mathcal S_{jD}}w'_{jD,X}\,\alpha_R^{X},\qquad L^\ast_{jD}=\sum_{X\in\mathcal S_{jD}}w'_{jD,X}\,\alpha_R^{R\setminus X},
\]
one system of weights for the two letters of a station; then $e^\ast_j=|R|-e_j-c_j$. The room of a family of tables is plain (Proposition~\ref{prop:seamidentity}(iii)), its stations of room sides being disjoint, with $w'_{jD,X}$ the multiplier of $X$ in table $j$ (Definition~\ref{def:seam}).
\end{definition}

% Reference 20: def:seamtransform; source line 10083
\begin{definition}[The station transform of a piece]\label{def:seamtransform}
For a piece $\mathcal X$ on $B$ with the stations $\mathcal D$, $\sum_j(s_j+t_j)c_j>0$, the wall $\kappa_{\mathcal X}$ and the masses $\mvec_B$, and for $\mathbf p=(p_{jD},p^\ast_{jD})_{j,D}$ with positive coordinates, one for each coordinate of the seam vector, the \emph{station transform} is
\begin{align*}
\mathcal I_{\mathcal X}(\mathbf p)&=\int\Gamma(\kappa_{\mathcal X})\,\sZ_{(\kappa_{\mathcal X},\mathbf s,\mathbf t);t_{\mathcal D}}(\mathbf x\circ\mathbf p)\,d\big((\boldsymbol\varrho^\rho_{\mathcal X})_\ast\phi^\rho_{\mathcal X}\big)(\mathbf x)\\
&=\int_{(0,\infty)^B}\alpha^{\mvec_B-\one}\prod_j\Big(\sum_Dp_{jD}V_{jD}\Big)^{-s_j}\Big(\sum_Dp^\ast_{jD}V^\ast_{jD}\Big)^{-t_j}d\alpha,
\end{align*}
in $[0,\infty]$: the kernel of Theorem~\ref{thm:gluewidth}(ii) read against the law, the period of the piece against the constant room letters $\mathbf p$, the second expression by the computation of Lemma~\ref{lem:absorbwidth} with $\alpha=\rho u$ and $v=\rho\,\gamma^\rho_{\mathcal X}(u)$.
\end{definition}

% Reference 20: thm:plainretain; source line 10092
\begin{theorem}[Plain partners separate laws]\label{thm:plainretain}
Let $\mathcal X$ and $\mathcal X'$ be pieces as in Theorem~\ref{thm:retainwidth}, with a real point $\mathbf c$ of the pins, with positive coordinates, in the interior of the domains of their seam functions on the pins. Fix room sides $\mathcal S_{jD}$ on a set $R$, every family with at least two members, such that the room word of the plain partners with these room sides has at $\mathbf c$ a Newton sum of the dimension $|R|-1$, as when one family of an active letter has a Newton polytope of that dimension, and masses $\mvec^0_R$ in its relative interior, with $\kappa_R=\sum_j(s_j+t_j)c_j-\kappa$. Then $(\boldsymbol\varrho^\rho_{\mathcal X})_\ast\phi^\rho_{\mathcal X}=(\boldsymbol\varrho^\rho_{\mathcal X'})_\ast\phi^\rho_{\mathcal X'}$ exactly when $P_{\mathcal X\sqcup\mathcal Y}=P_{\mathcal X'\sqcup\mathcal Y}$ for every plain partner $\mathcal Y$ with these room sides, at every positive weights and every masses in a neighbourhood of $\mvec^0_R$ on the hyperplane $\kappa_R=\sum_j(s_j+t_j)c_j-\kappa$. Such room sides exist for all stations and degrees: $e_j=1$, $|R|\ge\max_jc_j+2$, and every $\mathcal S_{jD}$ all the subsets of $1+c_j-i(D)$ marks.
\end{theorem}

The room of a family of tables is a plain partner, and multiplying the weights of the bases of each table by positive numbers that depend on the room side alone keeps the bases and the stations, and multiplies the members of each station of room sides by these numbers. So when a bunch of a family has at least two room sides in every station of every table with an active letter and the room word of its room has a Newton sum of full dimension, these families, with the masses of the room near the family's and those of the bunch kept, separate its law from every other law whose seam function converges, as its own does, at the native split (Lemma~\ref{lem:nativestations}, which puts the masses of the room in the relative interior there).

% New: proved in the last part.
\begin{lemma}[Partners on one mark see a point mass]\label{lem:mot-probe}
Let $\mathcal X$ be a piece on $B$ with the stations $\mathcal D$ (Definition~\ref{def:seampiece}), the exponents $(s_j,t_j)_j$ with $w=\sum_j(s_j+t_j)c_j>0$, the masses $\mvec_B$ and the wall $\kappa=\kappa_{\mathcal X}$ with $0<\kappa<w$. For a point $\mathbf x$ of its seam torus (Definition~\ref{def:seamlaw}), the \emph{one-mark piece at $\mathbf x$} is the piece on one mark $g$ with the seam letters $x_{jD}\alpha_g^{\,i(D)}$ and $x^\ast_{jD}\alpha_g^{\,c_j-i(D)}$ and the mass $\kappa$; its scale is $\alpha_g$, its seam vector is $\mathbf x$, its law is $\delta_{\mathbf x}$, and its station transform (Definition~\ref{def:seamtransform}) is $\Gamma(\kappa)\sZ_{(\kappa,\mathbf s,\mathbf t);t_{\mathcal D}}(\mathbf x\circ\mathbf p)$. If for some $c>0$ and $\mathbf x_0$ on the seam torus
\[
\mathcal I_{\mathcal X}(\mathbf p)=c\,\Gamma(\kappa)\,\sZ_{(\kappa,\mathbf s,\mathbf t);t_{\mathcal D}}(\mathbf x_0\circ\mathbf p)\qquad\text{for every }\mathbf p,
\]
then the law of $\mathcal X$ is $c\,\delta_{\mathbf x_0}$. The station transform at $\mathbf p$ is the period of the glue of $\mathcal X$ with the partner on one mark $b$ with the room letters $p_{jD}\alpha_b^{\,c_j-i(D)}$ and $p^\ast_{jD}\alpha_b^{\,i(D)}$ and the mass $w-\kappa$, the tinker $\mathbf p$ of one partner on one mark (Definition~\ref{def:tinkers}). So a piece that acts on these partners as $c$ times one mark has the law of $c$ times one mark, with no condition on the domains of the seam functions as in Theorem~\ref{thm:retainwidth}.
\end{lemma}

\section{Storage, and what replaces it}\label{sec:storage}

% Reference 20: thm:storeone; source line 10113 (its cells carrying the weights of the room children).
% Proved in the last part: thm:storeone.
\begin{theorem}[A bunch is stored on one mark exactly at width one]\label{thm:storeone}
In Theorem~\ref{thm:gluewidth}(i) let the partner be the room of a family of tables. The partner with its stand-in at $\mathbf x$, $R^{(\mathbf x)}$, is a family of tables on $R\sqcup\{g\}$, with a system of weights for each letter as a cell may carry (Definition~\ref{def:substitution}), at every point of the seam torus exactly when every table with an active letter has $c_j\le1$, and $c_j=1$ when both of its letters are active. At width one with factoring weights, normalized as in Proposition~\ref{prop:widthonelaw}, it is the family of the cells $\cB^+_{j,R}$ of the room with $g$ twisted in both letters, which the change of variable of Proposition~\ref{prop:widthonelaw}(iii) turns into the twist of Definition~\ref{def:twist}, and Lemma~\ref{lem:absorbwidth} is Lemma~\ref{lem:absorb} with the scale changed from $\gamma^K_{1,B}$ to $\gamma^\rho_B$, a fixed bijection of the tori with a density (Proposition~\ref{prop:widthonelaw}). A table with $c_j=0$ does not see the stand-in, in either letter, and at the width $c_j\ge2$ the stand-in enters the letters of table $j$ with the degree $c_j$, the coefficient of $\alpha_g^{c_j}$ in $K_j$ being $\sum_{D\in\mathcal D_{j,c_j}}x_{jD}L_{jD}\neq0$: the partner with its stand-in is a family of Newton words, whose period converges exactly when its Newton sum has the dimension $|R|$ and the masses lie in its relative interior (Theorem~\ref{thm:poswordconv}).
\end{theorem}

So the single mark of width one is the case $c=1$ of a stand-in of degree $c$. What replaces storage at every width is the stand-in itself, the station-wise pairing of Theorem~\ref{thm:gluewidth}: every statement about pieces is made on pieces with stations, whose glue is a family of Newton words, and table structure is never used to compose.

% New: proved in the last part.
\begin{theorem}[A point mass is a pinned set]\label{thm:mot-pointmass}
Let $F=((\cB_j,\theta_j);s_j,t_j)_{j=1}^N$ be a family on $\cE$ with real exponents $s_j,t_j\ge0$ and masses on its balance wall with every proper wall positive, and let $\varnothing\neq S\subsetneq\cE$, with the room $S'=\cE\setminus S$, be read as a piece with the stations of its room and $S'$ as a partner (Definition~\ref{def:seampiece}). Call $S$ \emph{pinned} when every active table has width at most one at $S$, with its weights factoring across $S$ and $S'$, and the active letters of the tables of width one read $S$ alike (Theorem~\ref{thm:pinned}): there are a function $\gamma$ on $\Delta_S$ and constants $c_j,c^\ast_j>0$ with $\gamma^K_{j,S}=c_j\gamma$ whenever $s_j>0$ and $\gamma^\Psi_{j,S}=c^\ast_j\gamma$ whenever $t_j>0$ (Definition~\ref{def:readings}). Then $w=\sum_j(s_j+t_j)c_j(S)>0$ and $0<\kappa_S<w$, and the following are equivalent.
\begin{enumerate}[label=(\roman*)]
\item The law of $S$ is carried by one point: $(\boldsymbol\varrho^\rho_S)_\ast\phi^\rho_S=c\,\delta_{\mathbf x_0}$ (Definition~\ref{def:seamlaw}) with $c>0$ and $\mathbf x_0$ on the seam torus.
\item $S$ is pinned.
\item $S$ \emph{splits off}: there are $c>0$ and $\mathbf x_0$ on the seam torus such that every partner $\mathcal Y$ of $S$ with the wall $w-\kappa_S$ has $P_{S\sqcup\mathcal Y}=c\,P_{\mathcal Y^{(\mathbf x_0)}}$ in $[0,\infty]$, $\mathcal Y^{(\mathbf x_0)}$ the partner with its stand-in at $\mathbf x_0$ carrying the mass $\kappa_S$ (Theorem~\ref{thm:gluewidth}(i)).
\item $S$ splits off against the partners on one mark: its station transform is $c$ times that of the one-mark piece at $\mathbf x_0$ with the mass $\kappa_S$ (Lemma~\ref{lem:mot-probe}).
\end{enumerate}
When they hold, the seam vector is the constant $\mathbf x_0$, $c=\int_{\Delta_S}\phi^\rho_S$, and $P_F(\mvec;\mathbf s,\mathbf t)=c\,P_{S'^{(\mathbf x_0)}}$. For a family of lists, read at the factorization with the weights $1$, the active letters of the tables of width one read $S$ alike exactly when their hinge values are equal, $c_j=c^\ast_j=1$, and then $\mathbf x_0=\one$. For one list $\cB$ with $s,t>0$ and $|S|\ge2$, $S$ is pinned exactly when it has width one and $\varrho_S=1$, that is $\mathbf E_{\cB|S}=\mathbf E_{\cB/S'}$.
\end{theorem}
When $\cB|S$ and $\cB/S'$ are connected, Chapter~\ref{ch:fusion} reads the last condition as a dual pair (Corollary~\ref{cor:mot-pinneddual}).

% New: proved in the last part.
\begin{corollary}[Which sets split off]\label{cor:mot-whatsplits}
Let $F$ and $S$ be as in Theorem~\ref{thm:mot-pointmass}.
\begin{enumerate}[label=(\roman*)]
\item If some active table has width two or more at $S$, then $S$ does not split off, and its law is carried by no point.
\item If $F=((\cB,\theta);s,t)$ is one prime table (Theorem~\ref{thm:primetree}) with $s+t>0$, no set with at least two marks on each side splits off.
\item If $F=(\cB;s,0)$ is one list with $s>0$ and $e$ is a mark that is neither a loop nor a coloop, then $\cE\setminus\{e\}$ splits off, and
\[
P_F(\mvec;s,0)=\Gst21{s-m_e,\ m_e}{s}\,P_{((\cB/e;\,m_e,0),\,(\cB\setminus e;\,s-m_e,0))}\big(\mvec|_{\cE\setminus\{e\}}\big),
\]
the period of the bunch family of Corollary~\ref{cor:pinnedproduct}; at $(0,t)$, $t>0$, the same holds with $\Gst21{t-m_e,\ m_e}{t}$ and the family $((\cB/e;\,0,t-m_e),(\cB\setminus e;\,0,m_e))$.
\item If $F=(\cB;s,t)$ is one list on at least three marks with $s,t>0$ and $e$ is neither a loop nor a coloop, then $\cE\setminus\{e\}$ splits off exactly when $\mathbf E_{\cB\setminus e}=\mathbf E_{\cB/e}$.
\end{enumerate}
\end{corollary}

% Reference 20: lem:higgsmajor; source line 10122
\begin{lemma}[The level cell of a bunch]\label{lem:higgsmajor}
Let $B$ have the width $c=c(B)\ge1$ in a table, and let $G$ be a set of $c$ new marks. The sets $p\sqcup F$, $F\subseteq G$ and $p$ a pattern of excess $c-|F|$, are the bases of a table $\cB^\natural_B$ on $B\sqcup G$ of rank $r(B)$, the \emph{level cell} of $B$, with any positive weights; $\cB^\natural_B\setminus G=\cB|B$ and $\cB^\natural_B/G=\cB/R$. At width one $\cB^\natural_B$ is the cell of the bunch $\cB^+_B$ of Theorem~\ref{thm:widthone}.
\end{lemma}

% Reference 20: thm:envelope; source line 10129
\begin{theorem}[The level envelope]\label{thm:envelope}
Let the weights of a table factor across $B$ and $R$, $\theta_T=\theta'_{T\cap B}\,\theta''_{T\cap R}$, with $c=c(B)\ge1$, and let $\cB^\natural_B$ on $B\sqcup G$ and $\cB^\natural_R$ on $R\sqcup G'$ be the level cells of the bunch and of the room, $G'$ a copy of $G$, with the weights $\theta'_p$ on $p\sqcup F$ and $\theta''_X\binom c{|F'|}^{-1}$ on $X\sqcup F'$. Splicing the $c$ pairs of a mark of $G$ and its copy (Theorem~\ref{thm:hingepairs}) gives the table $\cB^{\mathrm{env}}$ on $\cE$, the \emph{level envelope} of the bunch, whose bases are the sets $p\sqcup X$ of a pattern $p$ and a room side $X$ of one excess, with the weights $\theta'_p\,\theta''_X$. It has the patterns and the room sides of the table at every excess, its seam has one rectangle at every excess, and $\cB\subseteq\cB^{\mathrm{env}}$, with equality exactly when every excess of the seam of $B$ is one rectangle.
\end{theorem}

A bunch in level form is therefore stored on $c(B)$ pairs of hinges, the table being the splicing of its two level cells, and every bunch is read off its envelope by its seam: the table keeps, at each excess, the rectangles of the compatible pairs.

% Reference 20: prop:gluenec; source line 10136
\begin{proposition}[When a glue of seamed pieces is a table: what is necessary]\label{prop:gluenec}
Let $\cB$ be a nonempty family of $r$-sets of $B\sqcup R$, with children and room children as in Definition~\ref{def:seam}. If $\cB$ is the family of bases of a table, then:
\begin{enumerate}[label=(\alph*)]
\item every child and every room child is a table;
\item the children at adjacent patterns $p$ and $p+b$ are quotients, the child at $p+b$ of the child at $p$, and likewise the room children at adjacent room sides;
\item at the least and at the largest number of marks of a pattern all the children are one table, and likewise all the room children;
\item with $d$ and $d+c$ the least and the largest numbers of marks of a pattern, for every $k$ with $d\le k\le d+c$ the patterns of $k$ marks are the $k$-sets independent in the table of the patterns of $d+c$ marks that span the table of the patterns of $d$ marks, the bases of the $(k-d)$-th lift of the latter toward the former (Proposition~\ref{prop:higgs}); and likewise for the room sides.
\end{enumerate}
\end{proposition}

% Reference 20: thm:gluing; source line 10143
\begin{theorem}[The gluing criterion]\label{thm:gluing}
Conditions (a)--(d) of Proposition~\ref{prop:gluenec} suffice: a nonempty family of $r$-sets of $B\sqcup R$ satisfying them is the family of bases of a table. So a nonempty family of sets of one size is the family of bases of a table exactly when, at one bunch, its children and room children are tables, adjacent ones are quotients, the extreme ones are one, and its levels are lifts (Proposition~\ref{prop:higgs}).
\end{theorem}

Each of (c) and (d) is needed, and (d) on each side. The criterion is what storage of a bunch on new marks of a table asks, and nothing in the periods depends on it: the glue of Theorem~\ref{thm:gluewidth} composes pieces whether or not their glue is a table.


\chapter{Composing bunches}\label{ch:composebunches}

The pieces of Chapter~\ref{ch:width} compose. A family with several bunches is read through a core in which every bunch is replaced by its stand-in at its seam point, and its period is the mixture of the periods of these cores over the laws of the bunches (Theorem~\ref{thm:coreswidth}); a bunch inside a bunch makes the law of the larger one a mixture of the laws of its nodes (Theorem~\ref{thm:nestedlaw}); and along a chain of bunches the letters are products of transfer matrices indexed by the stations of room sides, the laws are carried from bunch to bunch by one transfer each, and the period opens as a flow of shares through the stations (Theorem~\ref{thm:chainwidth}, Corollary~\ref{cor:chainword}). Every operation is a rule on the seam functions on contours or on the station transforms (Theorem~\ref{thm:lawreadings}), a closed chain of cells is read the same way (Theorem~\ref{thm:cyclicchain}), and the Foreman tree with nodes of any width is evaluated by one procedure whose reusable data are the laws on its bunches (Definition~\ref{def:seamproc}). Read backward, the period as a function of the masses determines the potential, so the classes of marks and the tree are read off the periods, and with retention at every node this is the correspondence between a realization and its word (Theorems~\ref{thm:periodpotential} and~\ref{thm:correspondence}).

\section{Cores, nested pieces and chains of bunches}\label{sec:treeswidth}

% Reference 20: thm:coreswidth; source line 10172
\begin{theorem}[The core reads its pieces, at any width]\label{thm:coreswidth}
Let $B_1,\dots,B_q$ be disjoint bunches of a family $F$ on its balance wall, each with $\sum_j(s_j+t_j)c_j(B_l)>0$, with their scales $\gamma^\rho_{B_l}$, scale forms $\phi^\rho_{B_l}$ and seam vectors $\boldsymbol\varrho^\rho_{B_l}$ (Definition~\ref{def:seamlaw}), and let $Z$ be the marks outside them. For $u=(u_l)_l\in\prod_l\Delta_{B_l}$ read the letters of $F$ at $\alpha_{B_l}=\alpha_{g_l}u_l/\gamma^\rho_{B_l}(u_l)$, for new marks $g_l$, and divide each letter of table $j$ by the product over $l$ of the extreme excess sums of the bunches $B_l$ in that letter (Definition~\ref{def:seamlaw}), $V^{(0)}_j$ for $K_j$ and $V^{\ast(c_j(B_l))}_j$ for $\Psi_j$, read at the same point. The letters obtained depend on $u$ only through the seam points $\mathbf x_l=\boldsymbol\varrho^\rho_{B_l}(u_l)$; with them, the exponents of $F$ and the masses $\kappa_{B_l}$ on $g_l$ and $\mvec_Z$ on $Z$ they form the \emph{stand-in core} $F^{(\mathbf x_1,\dots,\mathbf x_q)}$, a family of Newton words on $Z\sqcup\{g_1,\dots,g_q\}$ (Definition~\ref{def:seampiece}) in which the stand-in $g_l$ enters the letters $K_j$ and $\Psi_j$ with the degree $c_j(B_l)$. In $[0,\infty]$,
\begin{align*}
P_F(\mvec;\mathbf s,\mathbf t)&=\int_{\prod_l\Delta_{B_l}}\prod_l\phi^\rho_{B_l}(du_l)\ P_{F^{(\boldsymbol\varrho^\rho_{B_1}(u_1),\dots,\boldsymbol\varrho^\rho_{B_q}(u_q))}}\big(\mvec_Z,\kappa_{B_1},\dots,\kappa_{B_q}\big)\\
&=\int P_{F^{(\mathbf x_1,\dots,\mathbf x_q)}}\big(\mvec_Z,\kappa_{B_1},\dots,\kappa_{B_q}\big)\prod_ld\big((\boldsymbol\varrho^\rho_{B_l})_\ast\phi^\rho_{B_l}\big)(\mathbf x_l).
\end{align*}
At width one with factoring weights the stand-in core is the twisted family $F_{\mathcal H}[\mathbf x]$ of Theorem~\ref{thm:node}(i), up to the change of scale of Proposition~\ref{prop:widthonelaw}, and this is Theorem~\ref{thm:node}(i); before the laws are formed it is Theorem~\ref{thm:corereads} with the scale of every marker divided by the scale $\gamma^\rho_{B_l}$ of its piece.
\end{theorem}

The core meets its pieces only through their seam points: every bunch collapses to one stand-in mark carrying the mass of its wall, and enters each letter with the degree of its width, at its stations twisted by its seam point. The cores of Chapter~\ref{ch:construct} are the case in which every piece is a cell.

% Reference 20: thm:nestedlaw; source line 10183
\begin{theorem}[The law of a piece from its sub-pieces]\label{thm:nestedlaw}
Let $S$ be a bunch of a family with $\sum_j(s_j+t_j)c_j(S)>0$, and $S_1,\dots,S_k\subseteq S$ disjoint bunches with $\sum_j(s_j+t_j)c_j(S_l)>0$.
\begin{enumerate}[label=(\roman*)]
\item The seam letters of $S$ are linear in those of each $S_l$, with nonnegative coefficients: in each table, for $S_1$, every station $D$ of room sides of $S$ and any member $X\in D$,
\[
K_{(D)}=\sum_q\frac{w'_{q\sqcup X}}{w'_X}\,\alpha^{\,q}\,K_{(D_1(q))},\qquad\Psi_{(D)}=\sum_q\frac{w'_{q\sqcup X}}{w'_X}\,\alpha^{(S\setminus S_1)\setminus q}\,\Psi_{(D_1(q))},
\]
the sums over the sets $q\subseteq S\setminus S_1$ for which $q\sqcup X$ is a room side of $S_1$, $D_1(q)$ its station, $w'_{q\sqcup X}$ its multiplier there and $w'_X$ that of $X$ at $S$; the coefficients and the stations $D_1(q)$ do not depend on the member $X$.
\item Let the \emph{node} $S^{(\mathbf x_1,\dots,\mathbf x_k)}$, the node at $S$ of the tree of the bunches $S\supseteq S_1,\dots,S_k$, be the piece on $(S\setminus\bigcup_lS_l)\sqcup\{g_1,\dots,g_k\}$ whose seam letters are those of $S$ with each $S_l$ replaced by its stand-in at $\mathbf x_l$, as in Theorem~\ref{thm:coreswidth}, with the masses $\kappa_{S_l}$ on $g_l$. The node has the wall of $S$, and, in this statement writing $\lambda_S$, $\lambda_{S_l}$ and $\lambda_{\mathbf x}$ for the laws of $S$, of $S_l$ and of the node at $\mathbf x=(\mathbf x_l)_l$, the law of $S$ is the mixture of the laws of its nodes:
\[
\int H\,d\lambda_S=\int\Big(\int H\,d\lambda_{\mathbf x}\Big)\prod_ld\lambda_{S_l}(\mathbf x_l)\qquad\text{for every Borel }H\ge0,
\]
and so are the seam functions, $\int\mathbf z^{-(\mathbf a,\mathbf b)}\,d\lambda_S(\mathbf z)=\int\big(\int\mathbf z^{-(\mathbf a,\mathbf b)}\,d\lambda_{\mathbf x}(\mathbf z)\big)\prod_ld\lambda_{S_l}(\mathbf x_l)$ at all shares $(\mathbf a,\mathbf b)$.
\end{enumerate}
At width one, with factoring weights, this is Theorem~\ref{thm:node}(ii), up to the change of scale of Proposition~\ref{prop:widthonelaw}.
\end{theorem}

% Reference 20: thm:chainwidth; source line 10198
\begin{theorem}[A chain of bunches]\label{thm:chainwidth}
Let $\varnothing=S_0\subset S_1\subset\cdots\subset S_L=\cE$, $L\ge2$, be a chain of bunches of a family with the steps $B_{l-1}=S_l\setminus S_{l-1}$, and for $0<l<L$ let $\mathcal D(S_l)$ be the stations of room sides of the bunch $S_l$, in each table, the index $j$ dropped.
\begin{enumerate}[label=(\roman*)]
\item \emph{Letters.} For $2\le l\le L-1$, $D'\in\mathcal D(S_{l-1})$ and $D\in\mathcal D(S_l)$ put, for any $X\in D$,
\[
\mathbf N^K_l(D',D)=\sum_q\frac{w'_{q\sqcup X}}{w'_X}\,\alpha_{B_{l-1}}^{\,q},\qquad\mathbf N^\Psi_l(D',D)=\sum_q\frac{w'_{q\sqcup X}}{w'_X}\,\alpha_{B_{l-1}}^{\,B_{l-1}\setminus q},
\]
the sums over the $q\subseteq B_{l-1}$ for which $q\sqcup X$ is a room side of $S_{l-1}$ in the station $D'$. With $\mathbf k=(K_{(D)})_{D\in\mathcal D(S_1)}$ and $\mathbf k^\ast=(\Psi_{(D)})_{D\in\mathcal D(S_1)}$ the seam letters of $S_1$ and $\mathsf y=(\mathsf y_D)_{D\in\mathcal D(S_{L-1})}$ and $\mathsf y^\ast=(\mathsf y^\ast_D)_{D\in\mathcal D(S_{L-1})}$ the room station letters of $S_{L-1}$,
\[
K=\mathbf k^{\mathsf T}\mathbf N^K_2\cdots\mathbf N^K_{L-1}\,\mathsf y,\qquad\Psi=\mathbf k^{\ast\mathsf T}\mathbf N^\Psi_2\cdots\mathbf N^\Psi_{L-1}\,\mathsf y^\ast,
\]
products of matrices of polynomials with nonnegative coefficients, indexed by the stations of room sides of the bunches $S_l$: Theorem~\ref{thm:lettertransfer} with these stations as the states.
\item \emph{Laws.} Let $\sum_j(s_j+t_j)c_j(S_l)>0$ for every $0<l<L$, and in this statement write $\lambda_l$ for the law of $S_l$. For $2\le l\le L-1$ let $S_l^{(\mathbf x)}$ be the node of $S_l$ with its sub-bunch $S_{l-1}$ at $\mathbf x$ (Theorem~\ref{thm:nestedlaw}(ii)), the piece on $B_{l-1}\sqcup\{g\}$ with the seam letters $\sum_{D'}x_{D'}\,\alpha_g^{\,i(D')}\mathbf N^K_l(D',D)$ and $\sum_{D'}x^\ast_{D'}\,\alpha_g^{\,c(S_{l-1})-i(D')}\mathbf N^\Psi_l(D',D)$, $D\in\mathcal D(S_l)$, and the mass $\kappa_{S_{l-1}}$ on $g$, and let $\lambda^{(l)}_{\mathbf x}$ be its law. Then $\lambda_l$ is the image of $\lambda_{l-1}$ under the \emph{transfer} $\mathcal Q_l$,
\[
\lambda_l=\mathcal Q_l\lambda_{l-1}=\int\lambda^{(l)}_{\mathbf x}\,d\lambda_{l-1}(\mathbf x),
\]
and $P_F(\mvec;\mathbf s,\mathbf t)=\int P_{B_{L-1}^{(\mathbf x)}}(\mvec_{B_{L-1}},\kappa_{S_{L-1}})\,d\lambda_{L-1}(\mathbf x)$, the room $B_{L-1}$ of the last bunch $S_{L-1}$ with its stand-in at $\mathbf x$ (Theorem~\ref{thm:gluewidth}(i)).
\item \emph{Repetition.} If the nodes $S_l^{(\cdot)}$, $2\le l\le L-1$, are one piece up to bijections of marks keeping seam letters and masses, the transfers are one operator $\mathcal Q$ and $\lambda_l=\mathcal Q^{l-1}\lambda_1$: the chain is the power of one operator, whatever the widths of its bunches. At width one with factoring weights this is Theorem~\ref{thm:chainoperator} and Proposition~\ref{prop:chain}, the seam point being the twist up to the change of scale of Proposition~\ref{prop:widthonelaw}.
\end{enumerate}
\end{theorem}

% Reference 20: lem:nativestations; source line 10292 (the chain clause).
\begin{lemma}[The native split along a chain]\label{lem:nativechain}
Let a family of tables have real data on its balance wall, with $s_j,t_j\ge0$, not all zero, and $\mvec>0$, and a native gauge $\alpha^\star$ (Definition~\ref{def:native}), and let $\varnothing=S_0\subset S_1\subset\cdots\subset S_L=\cE$ be a chain of bunches as in Theorem~\ref{thm:chainwidth}, with the steps $B_{l-1}=S_l\setminus S_{l-1}$. For the active letters put
\[
e^{(l)}_{jD'D}=s_j\,\mathbb P_{j,\alpha^\star}\big(T\setminus S_{l-1}\in D',\ T\setminus S_l\in D\big),\qquad f^{(l)}_{jD'D}=t_j\,\mathbb P^\ast_{j,\alpha^\star}\big(T\setminus S_{l-1}\in D',\ T\setminus S_l\in D\big),
\]
and $a^{(l)}_{jD}=s_j\,\mathbb P_{j,\alpha^\star}(T\setminus S_l\in D)$, $b^{(l)}_{jD}=t_j\,\mathbb P^\ast_{j,\alpha^\star}(T\setminus S_l\in D)$, the native split at the stations of the bunch $S_l$ (Lemma~\ref{lem:nativestations}). Then the entry shares conserve the loads, $\sum_{D'}e^{(l)}_{jD'D}=a^{(l)}_{jD}=\sum_{D''}e^{(l+1)}_{jDD''}$, and likewise with $f$ and $b$; and $\mvec_{B_{l-1}}$ lies in the relative interior of the Newton sum of the word of the node $l$ at its shares, which converges there exactly when that Newton sum has the dimension $|B_{l-1}|-1$.
\end{lemma}

% Reference 20: cor:chainword; source line 10218
\begin{corollary}[The word of a chain of bunches]\label{cor:chainword}
In Theorem~\ref{thm:chainwidth} read the room $B_{L-1}$ of the last bunch $S_{L-1}$ as a last node whose transfer matrix from $\mathcal D(S_{L-1})$ to one terminal station is the column $\mathsf y$, and $\mathsf y^\ast$ for $\Psi$, and open every active letter with one share per nonzero entry of every transfer matrix: $e^{(l)}_{jD'D}$ for the entry $(D',D)$ of $\mathbf N^K_l$ and $f^{(l)}_{jD'D}$ for that of $\mathbf N^\Psi_l$, $2\le l\le L$. A station $D$ of a bunch $S_l$, $2\le l\le L-1$, carries the load $a^{(l)}_{jD}=\sum_{D'}e^{(l)}_{jD'D}$ in and $\sum_{D''}e^{(l+1)}_{jDD''}$ out, and a station $D$ of the first bunch $S_1$ the load $a^{(1)}_{jD}=\sum_{D''}e^{(2)}_{jDD''}$; likewise $b$ with $f$. In this statement write $W_l(\mathbf e,\mathbf f)=\prod_{D',D}\mathbf N^K_l(D',D)^{-e_{D'D}}\,\mathbf N^\Psi_l(D',D)^{-f_{D'D}}$ for the word of the node $l$ on $B_{l-1}$, the products over the active letters of every table understood. On contours of positive real parts at which the seam function of $S_1$ and the words of the nodes converge, which needs their Newton sums to have full dimension and, for a family with a native gauge, holds as soon as they do (Lemmas~\ref{lem:nativestations} and~\ref{lem:nativechain}),
\begin{multline*}
P_F(\mvec;\mathbf s,\mathbf t)=\int\prod_{j:\,s_j>0}\frac{\prod_{l=2}^L\prod_{D',D}\Gamma(e^{(l)}_{jD'D})}{\Gamma(s_j)\prod_{l=2}^{L-1}\prod_D\Gamma(a^{(l)}_{jD})}\prod_{j:\,t_j>0}\frac{\prod_{l=2}^L\prod_{D',D}\Gamma(f^{(l)}_{jD'D})}{\Gamma(t_j)\prod_{l=2}^{L-1}\prod_D\Gamma(b^{(l)}_{jD})}\\
\times P_{W_{S_1}(\mathbf a^{(1)},\mathbf b^{(1)})}(\mvec_{B_0})\prod_{l=2}^LP_{W_l(\mathbf e^{(l)},\mathbf f^{(l)})}(\mvec_{B_{l-1}}),
\end{multline*}
$W_{S_1}$ the seam word of the first bunch $S_1$ (Lemma~\ref{lem:seamword}), the fibre pinned by the letter totals, the equality of the loads in and out at every station of the bunches $S_2,\dots,S_{L-1}$, and the homogeneity of the first bunch and of every node but one. At a bunch or a node whose Newton sum is degenerate the pins and the reduced period of Theorem~\ref{thm:transversepins} take the place of the period. At $L=2$ this is Theorem~\ref{thm:gluewidth}(iii).
\end{corollary}

The word of a chain of bunches is a flow: one share per transfer entry, conserved at every station, each node contributing the period of its own word at the loads of its entries. Every station through which the flow passes is paid once by a Gamma in the denominator, except at the first bunch, where the seam function of the piece takes the loads. At width one with factoring weights every bunch $S_l$ has two stations, and this is the word of the tree (Corollary~\ref{cor:treeword}) grouped by the two stations of each bunch.

% Reference 20: thm:lawreadings; source line 10236 (nesting at one sub-piece).
% Proved in the last part: thm:lawreadings.
\begin{theorem}[Two transforms of a law, and their rules]\label{thm:lawreadings}
Let a piece $\mathcal X$ on $B$ with the stations $\mathcal D$ have $\sum_j(s_j+t_j)c_j>0$, and a real point $\mathbf c$ of the pins of the seam, with positive coordinates, in the interior of the domain of its seam function on the pins.
\begin{enumerate}[label=(\roman*)]
\item \emph{Two transforms.} The law $(\boldsymbol\varrho^\rho_{\mathcal X})_\ast\phi^\rho_{\mathcal X}$ is determined by its seam function on the contour through $\mathbf c$, and by its station transform $\mathcal I_{\mathcal X}$ (Definition~\ref{def:seamtransform}) on any nonempty open set of points $\mathbf p$; the two are related by
\[
\mathcal I_{\mathcal X}(\mathbf p)=\int_{\text{pins}}\prod_{j:\,s_j>0}\Gst{|\mathcal D_j|}1{\mathbf a_j}{s_j}\prod_{j:\,t_j>0}\Gst{|\mathcal D_j|}1{\mathbf b_j}{t_j}\ \mathbf p^{-(\mathbf a,\mathbf b)}\int_{\Delta_B}(\boldsymbol\varrho^\rho_{\mathcal X})^{-(\mathbf a,\mathbf b)}\,d\phi^\rho_{\mathcal X},
\]
$\mathbf p^{-(\mathbf a,\mathbf b)}=\prod_{j,D}p_{jD}^{-a_{jD}}p^{\ast}_{jD}{}^{-b_{jD}}$.
\item \emph{The glue.} For a partner $\mathcal Y$ with the wall $\sum_j(s_j+t_j)c_j-\kappa_{\mathcal X}$, $P_{\mathcal X\sqcup\mathcal Y}=\int\mathcal I_{\mathcal X}(\mathbf x')\,d\big((\boldsymbol\varrho^\rho_{\mathcal Y})_\ast\phi^\rho_{\mathcal Y}\big)(\mathbf x')$, the station transform read against the law of the partner; on the pins, at a point where the seam function of the partner converges too, it is the share side of Theorem~\ref{thm:gluewidth}(iii). It extends the node of width one in one scale (Theorem~\ref{thm:onescale}) to any width.
\item \emph{Nesting.} In Theorem~\ref{thm:nestedlaw}, $\mathcal I_S(\mathbf p)=\int\mathcal I_{S^{(\mathbf x_1,\dots,\mathbf x_k)}}(\mathbf p)\prod_ld\big((\boldsymbol\varrho^\rho_{S_l})_\ast\phi^\rho_{S_l}\big)(\mathbf x_l)$. With one sub-piece $S_1$, in this statement write $\mathbf N^K_{2,j}(\alpha)$ and $\mathbf N^\Psi_{2,j}(\alpha)$, $\alpha=\alpha_{S\setminus S_1}$, for the matrices $\mathbf N^K_2$ and $\mathbf N^\Psi_2$ in table $j$ of the chain $\varnothing\subset S_1\subset S\subset\cE$ (Theorem~\ref{thm:chainwidth}(i)), the transfer matrices from the bunch $S_1$ to the bunch $S$. Then
\[
\mathcal I_S(\mathbf p)=\int_{(0,\infty)^{S\setminus S_1}}\alpha^{\mvec_{S\setminus S_1}-\one}\,\mathcal I_{S_1}\big(\big(\mathbf N^K_{2,j}(\alpha)\,\mathbf p_j\big)_j,\big(\mathbf N^\Psi_{2,j}(\alpha)\,\mathbf p^\ast_j\big)_j\big)\,d\alpha,
\]
$\mathbf p_j=(p_{jD})_D$ and $\mathbf p^\ast_j=(p^\ast_{jD})_D$; and at shares $(\mathbf a,\mathbf b)$ on the pins of the seam of $S$, on contours of positive real parts at which the seam function of $S_1$ and the word of the node converge,
\begin{multline*}
\int_{\Delta_S}(\boldsymbol\varrho^\rho_S)^{-(\mathbf a,\mathbf b)}\,d\phi^\rho_S=\int\prod_{j:\,s_j>0}\frac{\prod_{D_1,D}\Gamma(e_{jD_1D})}{\prod_D\Gamma(a_{jD})}\prod_{j:\,t_j>0}\frac{\prod_{D_1,D}\Gamma(f_{jD_1D})}{\prod_D\Gamma(b_{jD})}\\
\times\int_{\Delta_{S_1}}(\boldsymbol\varrho^\rho_{S_1})^{-(\mathbf a^{(1)},\mathbf b^{(1)})}\,d\phi^\rho_{S_1}\ P_{W(\mathbf e,\mathbf f)}(\mvec_{S\setminus S_1}),
\end{multline*}
the shares $e_{jD_1D}$ and $f_{jD_1D}$ one per nonzero entry of the transfer matrices, pinned by $\sum_{D_1}e_{jD_1D}=a_{jD}$, $\sum_{D_1}f_{jD_1D}=b_{jD}$ and the homogeneity of the node, with the loads $a^{(1)}_{jD_1}=\sum_De_{jD_1D}$ and $b^{(1)}_{jD_1}=\sum_Df_{jD_1D}$, and $W(\mathbf e,\mathbf f)$ the word of the entries on $S\setminus S_1$ as in Corollary~\ref{cor:chainword}.
\item \emph{Repetition.} The rule of (iii) for a node repeated along a chain is one operator on seam functions, and on station transforms one operator on functions of $\mathbf p$; a chain of repeated nodes is its power applied to the first bunch.
\end{enumerate}
\end{theorem}

So every operation along a chain is a rule on the seam functions on contours, or on the station transforms: the representations (b) and (a) of Definition~\ref{def:treeproc}, at every width.

% Reference 20: def:cyclicchain; source line 10262
\begin{definition}[Cyclic chains]\label{def:cyclicchain}
A family on $\cE=B_1\sqcup\cdots\sqcup B_k$, $k\ge2$, the \emph{segments} $B_i$ indexed by $\mathbb Z_k$, is a \emph{cyclic chain} with the finite sets of \emph{states} $\mathcal J_0,\dots,\mathcal J_{k-1}$, $\mathcal J_k=\mathcal J_0$, and the \emph{top layer} $h\in\mathbb N$, when every active letter $L$ is
\[
L=[\varepsilon^h]\operatorname{tr}\prod_{i=1}^k\mathbf M_{L,i}(\alpha_{B_i};\varepsilon),
\]
the coefficient of $\varepsilon^h$ in the trace of a product of matrices $\mathbf M_{L,i}$ indexed by $\mathcal J_{i-1}\times\mathcal J_i$, the \emph{segment matrices}, whose entries are polynomials in $\alpha_{B_i}$ and $\varepsilon$ with nonnegative coefficients.
\end{definition}

% Reference 20: thm:cyclicchain; source line 10270 (one table in (i)).
% Proved in the last part: thm:cyclicchain.
\begin{theorem}[Cyclic chains]\label{thm:cyclicchain}
\begin{enumerate}[label=(\roman*)]
\item Let every table $j$ of a family be a closed chain of cells $\mathcal A_{j,1},\dots,\mathcal A_{j,k}$ with several hinges spliced along hinge pairs (Theorem~\ref{thm:hingepairs}(ii)), the cells $\mathcal A_{j,i}$ of all tables having the marks $B_i$ besides their hinges, and the pairs $I_i$ joining the cell $i$ to the cell $i+1$ named alike in every table, each pair named once. The family is a cyclic chain with the top layer $0$: its states at the boundary $i$ are the sets $J\subseteq I_i$ of the pairs whose hinge on the side of the cell $i$ is occupied, and the segment matrices of the letters of table $j$ are the channels of its cells (Definition~\ref{def:channels}),
\[
\mathbf M_{K_j,i}[J,J']=K_{\mathcal A_{j,i},((I_{i-1}\setminus J)\sqcup J')},\qquad\mathbf M_{\Psi_j,i}[J,J']=\Psi_{\mathcal A_{j,i},((I_{i-1}\setminus J)\sqcup J')}.
\]
\item Widen the states by layers, $\mathcal J_i\times\{0,\dots,h\}$, a term in $\varepsilon^\mu$ of an entry raising the layer by $\mu$, and write $\widetilde{\mathbf M}_{L,i}$ for the widened segment matrices. Then $L=\sum_{J\in\mathcal J_0}\mathbf e_{(J,0)}^{\mathsf T}\prod_i\widetilde{\mathbf M}_{L,i}\,\mathbf e_{(J,h)}$: for every starting state a chain in the sense of Theorem~\ref{thm:chainwidth}(i), the segment matrices as transfer matrices. In this statement write $w_L$ for the exponent of the active letter $L$. Open every active letter first over the starting states $J$ at which its term $\mathbf e_{(J,0)}^{\mathsf T}\prod_i\widetilde{\mathbf M}_{L,i}\,\mathbf e_{(J,h)}$ is not zero, with shares $c_{L,J}$, $\sum_Jc_{L,J}=w_L$, and then every chain as in Corollary~\ref{cor:chainword}: one share $e^{(J)}_{L,i}$ per starting state $J$ and per nonzero entry of a widened segment matrix that lies on a path from $(J,0)$ to $(J,h)$, the shares of each starting state a flow from $(J,0)$, with the total $c_{L,J}$, to $(J,h)$, conserved at every boundary, $a^{(J)}_{L,i}$ the loads it leaves at the boundary after the segment $i$, and $e^{(J)}_{L,1}=a^{(J)}_{L,1}$. On contours of positive real parts at which the words of the segments converge,
\[
P_F(\mvec;\mathbf s,\mathbf t)=\int\prod_L\frac1{\Gamma(w_L)}\prod_{L,J}\frac{\prod_{i=2}^k\prod\Gamma(e^{(J)}_{L,i})}{\prod_{i=2}^{k-1}\prod\Gamma(a^{(J)}_{L,i})}\ \prod_{i=1}^kP_{W_i}(\mvec_{B_i}),
\]
the inner products over the entries and over the states, $W_i$ the word on $B_i$ of the entries of the segment $i$ at their total shares $\sum_Je^{(J)}_{L,i}$, and the fibre pinned by the conservation and by the homogeneity of every segment but one.
\item For $k$ copies of one segment the integrand of (ii), read at each boundary as a function of the loads of all the starting states, is the kernel of one operator $\mathcal C$, restricted at each boundary to the loads of the states that lie on a path there, which integrates the loads out over the fibre pinned by the homogeneity of the segment. With $\mathcal C_\circ$ the same kernel without that pin, which for one segment is implied by those of the others, the letter totals and the balance, the period is
\[
P_F(\mvec;\mathbf s,\mathbf t)=\int\prod_L\frac{\prod_J\Gamma(c_{L,J})}{\Gamma(w_L)}\,\big(\mathcal C^{k-1}\mathcal C_\circ\big)(\mathbf c,\mathbf c)\,d\mathbf c,
\]
a power of one operator taken from the loads $\mathbf c$ at the starting states back to them. Its price is that of one segment and the binary digits of $k$ (Chapter~\ref{ch:price}); evaluated, the power takes at most $\lfloor\log_2k\rfloor$ squarings of $\mathcal C$ and as many products.
\end{enumerate}
\end{theorem}

A closed chain is a chain whose last boundary is glued back to its first: the trace sums over the states at one boundary, and each starting state carries its own flow around the closed chain. Cells spliced along any graph of hinge pairs stay tables (Theorem~\ref{thm:hingepairs}); around a closed chain their letters are traces of products of channel matrices, and repeated cells give the power of one operator.

\section{The procedure}\label{sec:widthproc}

The Newton sums of Lemmas~\ref{lem:nativestations} and~\ref{lem:nativechain} can be degenerate. When every station of room sides is a single room side, every room station letter is a monomial, the room word is a monomial, and its period diverges at every share; the period is then read by Theorem~\ref{thm:gluewidth}(i) or (ii), or on the share side of Theorem~\ref{thm:transversepins} with $\vec{\mathcal N}_R=0$.

% Reference 20: def:seamproc; source line 10308
\begin{definition}[The seam procedure]\label{def:seamproc}
\emph{Given}: a family $F$ with real data on its balance wall, every proper wall positive. \emph{Classification}, from the tables alone: the molecules, the readings at the centre and the Foreman tree (Definition~\ref{def:procedure}, steps (1)--(5)); for a molecule with a connected active table the tree is a chain whose bunches $B_0,\dots,B_{d-1}$ (Corollary~\ref{cor:treechain}) are the steps of the chain of bunches $S_k=B_0\sqcup\cdots\sqcup B_{k-1}$, the leaf being the last step. \emph{Seams}, from the tables alone, at every bunch $S_k$:
\begin{enumerate}[label=(\arabic*)]
\item group the patterns of $S_k$ of each excess by their weighted children up to positive multiples, and the room sides by their weighted room children, with the multipliers normalized by their largest value (Definition~\ref{def:seam}); the stations, the station letters and the pair constants are determined;
\item record the seam matrices $H_i$ and the transfer matrices $\mathbf N^K_k$, $\mathbf N^\Psi_k$ of Theorem~\ref{thm:chainwidth}(i);
\item record the orbits of the stations under the bijections that keep $S_k$ and keep the potential (Proposition~\ref{prop:seamsym}, Lemma~\ref{lem:certifypot}).
\end{enumerate}
\emph{Evaluation} at a point $(\mvec;\mathbf s,\mathbf t)$:
\begin{enumerate}[label=(\arabic*),resume]
\item the law of $S_1=B_0$ (Definition~\ref{def:seamlaw}), with its canonical scale;
\item the transfers of the laws from bunch to bunch (Theorem~\ref{thm:chainwidth}(ii)), one operator for repeated nodes (Theorem~\ref{thm:chainwidth}(iii)), the transfer entries of pairs of stations in one orbit being images of one another (Proposition~\ref{prop:seamsym}(i));
\item the period at the leaf with its stand-in at the seam point of the last bunch $S_d$, integrated against the law of that bunch (Theorem~\ref{thm:gluewidth}(i)); equivalently the word of Corollary~\ref{cor:chainword} on the contour through the native shares of Lemma~\ref{lem:nativechain}, with the transverse pins of Theorem~\ref{thm:transversepins} at every bunch or node whose Newton sum is degenerate, every node a Newton word with multilinear letters closed by its peak series with the references of Theorem~\ref{thm:peak}(i) and every seam Gramma by its peak closing with the sharp references (Proposition~\ref{prop:relgrammaclose}(i)).
\end{enumerate}
\emph{Output}: the period, with the seams, the transfer matrices and the law on every bunch, ready to be reused when a node is changed. At a general Foreman tree the bunches of a node with several childcules are composed by Theorem~\ref{thm:coreswidth}, the childcules of a peel being the pieces of the core.
\end{definition}

Nothing in the procedure is chosen. The tree is the Foreman tree (Definition~\ref{def:tree}); the stations are fixed by the children, and their multipliers by the normalization of the largest one; the scale of a piece is the geometric mean of the oriented extreme ratios of the active letters, each taken to the root of its width, kept by the symmetries of the family (Proposition~\ref{prop:seamsym}(iii)); inside a seam Gramma the stations of the bunch are integrated before those of the room, the Foreman's order (Proposition~\ref{prop:relgrammaforeman}), and another choice of the stations that are not kept changes the pins by a unimodular matrix; the real point of the share contour, where the word is used, is the native split at the stations and the entries (Lemmas~\ref{lem:nativestations} and~\ref{lem:nativechain}); the references of the closings are those of Theorem~\ref{thm:peak}(i) and the sharp ones, read off the coefficients. Repetition is paid once: stations in one orbit carry data that are images of one another, and repeated nodes one operator.

% Reference 20: prop:seamwork; source line 10326
\begin{proposition}[What the work depends on]\label{prop:seamwork}
\begin{enumerate}[label=(\roman*)]
\item The law of a piece depends on the piece only through the joint law of the seam letters of its active letters under the measure $u^{\mvec_B-\one}d\omega_B$, with its exponents and wall: two pieces, on sets of any sizes, whose seam letters have the same joint law have equal laws and are interchangeable in every glue (Theorem~\ref{thm:retainwidth}); in particular so are two pieces whose seam letters and masses correspond under a bijection of their marks.
\item Along a chain the letters cost one product of matrices of the sizes $|\mathcal D(S_l)|$, and a seam Gramma of a part without circuits costs no integration (Definition~\ref{def:relgramma}, Corollary~\ref{cor:square}).
\item The law of a bunch depends on the room only through the room children of the bunch, the tables on the bunch that the room presents to it: changing the room while keeping its room children and their stations changes no law below the bunch.
\end{enumerate}
\end{proposition}

\section{The correspondence}\label{sec:correspondence}

The procedure reads a realization forward, from its tables to its word. The other direction asks what the periods distinguish, and it has two answers: as a function of the masses the period determines the potential, which carries the classes of marks and the tree; at one realization each bunch is seen by its room through its wall and its law, and only through them.

% Reference 20: thm:periodpotential; source line 10341
\begin{theorem}[The period function determines the potential]\label{thm:periodpotential}
Let a family have real exponents $s_j,t_j\ge0$ with $s_j+t_j>0$ for every $j$ and a period that converges at some masses on its balance wall, and let $U$ be the set of the complex masses on the balance wall at which it converges.
\begin{enumerate}[label=(\roman*)]
\item With the measure $dy$ on the hyperplane $\one^\perp=\{y\in\mathbb R^{\cE}:\ \ip\one y=0\}$ of the section $\prod_e\alpha_e=1$ (Lemma~\ref{lem:chart}),
\[
P_F(\mvec;\mathbf s,\mathbf t)=\int_{\one^\perp}e^{\ip{\mvec}y-\Phi(y)}\,dy,\qquad\mvec\in U,
\]
$\Phi$ the potential at the table's gauge (Definition~\ref{def:potential}). The function $P_F(\cdot;\mathbf s,\mathbf t)$ on $U$ determines $\Phi$, and is determined by it.
\item Two families on the same marks with the same exponents have the same period function exactly when they have the same irreducible factors with their exponents $(L,w_{\mathbf E}(L),w_{\mathbf O}(L))$ (Theorem~\ref{thm:poles}) and the same constant $\Phi(0)=\sum_j(s_j+t_j)\log K_j(\one)$; at the centre, every $\delta_j=0$, exactly when they have the same $(L,w_{\mathbf E}(L))$ and the same constant. In particular a family and the family of its dual tables (Proposition~\ref{prop:dualfamily}), the bases replaced by their complements, have the same period function at the centre.
\item A bijection $g$ of the marks keeps the period function, $P_F(g^\ast\mvec;\mathbf s,\mathbf t)=P_F(\mvec;\mathbf s,\mathbf t)$ on $U$, exactly when it keeps the potential, both parts. For a molecule, two marks lie in one class exactly when a bijection keeping the period function takes one to the other (Lemma~\ref{lem:certifypot}), and when its children's weights factor across the childcules at every node, every such bijection carries the Foreman tree onto itself, with its classes, units, atoms, bunches and childcules (Proposition~\ref{prop:treepot}).
\item At the centre the bijections that carry the family and those that exchange it keep the period function; at a lean the carrying ones keep it, and an exchanging one sends it to the period function of the opposite lean (Proposition~\ref{prop:certify}). For a molecule of one table two marks lie in one class at the centre exactly when a bijection carrying or exchanging the table takes one to the other, and at a lean exactly when a carrying one does (Corollary~\ref{cor:certifyone}), and a class of the centre that is not a class at a lean is there two classes of equal size, which an exchange swaps (Theorem~\ref{thm:acrosscentre}).
\end{enumerate}
\end{theorem}

So two marks of a molecule are told apart by the period function of the centre, or of a lean, exactly when no bijection keeping that potential joins them; for one table these are, at the centre, the bijections carrying or exchanging the table, and at a lean the carrying ones. The confusable classes, the classes of a lean that share a class of the centre, are exactly what the period function of the centre does not separate and that of a lean does. The tree is read off the period function but does not determine it; the irreducible factors with their exponents, and the constant $\Phi(0)$, do.

% Reference 20: thm:correspondence; source line 10360
\begin{theorem}[The correspondence]\label{thm:correspondence}
Let a molecule be read at a realization with masses on its balance wall and every proper wall positive, its children's weights factoring across the childcules at every node.
\begin{enumerate}[label=(\roman*)]
\item \emph{Forward.} Its period is the reassembly of the laws of the bunches of its Foreman tree with the leaves, by the seam procedure: the law of the first bunch on its seam torus, transferred along the chain by the later bunches as nodes, the childcules of a node composed at its core, and glued to the leaves by the stand-ins; nothing in it is chosen (Definition~\ref{def:seamproc}, Theorems~\ref{thm:coreswidth} and~\ref{thm:chainwidth}).
\item \emph{Backward at the realization.} At the node reached after $k$ peels the room sees the bunch $S_k$, the union of the bunches peeled down to it, through its wall and its law, and only through them: bunches with the same stations, exponents, walls and laws are interchangeable in every family glued on that room, and under the hypotheses of Theorem~\ref{thm:retainwidth} two different laws are told apart by the tinkers of one room and by plain partners (Theorems~\ref{thm:retainwidth} and~\ref{thm:plainretain}); the laws of nested bunches compose by Theorems~\ref{thm:nestedlaw} and~\ref{thm:chainwidth}.
\end{enumerate}
\end{theorem}

Whether the molecule has an atomic realization is read at a centre: it has one exactly when some centre is atomic, and an atom at a lean is an atom at the centre (Theorems~\ref{thm:centre} and~\ref{thm:leansplits}). The word of the tree is the forward procedure written out, and retention is the statement that it loses nothing that a room could detect.


\chapter{Foreman fusion}\label{ch:fusion}

A table of rank $r$ and corank $r^\ast$ is the contraction of $c=r^\ast-r$ new marks in an identically self-dual table, in which the new marks form a bunch in level form and their deletion leaves the dual table, exactly when every cocircuit of the table is a union of circuits and, at $c=0$, the table has the weights of its dual. The bases of the larger table are then unique, and at $c\le1$ the larger table itself: they are those of the level cell of the pair formed by the dual table and the table, the new marks are its fusion bunch, and their children on the marks of the table are the lifts of the table toward its dual, one for each level, whose weights between the two extremes, the middle weights, are a choice paired by complement, all $1$ unless others are named. This is Foreman fusion, and the table is Foreman-embedded in its fusion. The fused table is identically self-dual, so each of its marks comes off by one Beta corner at every pair of exponents (Corollary~\ref{cor:mergedvalues}); its automorphisms contain those of the table that keep the middle weights, all of them when the middle weights are $1$, times the permutations of the fusion bunch, and are often more; and every identically self-dual table is the fusion of each of its one-mark contractions. At the fusion width one the fusion bunch is one mark $g$, the letter of the fused table is $\Psi+\alpha_gK$, and the table is read off it by one Beta corner: the law of the fusion bunch is a point. At the width $c$ the seam of the fusion bunch is one rectangle at every excess, its seam letters are the elementary symmetric polynomials of the new marks up to constants, all $1$ for a list, and the fused table is the glue (Theorem~\ref{thm:gluewidth}) of a universal piece of width $c$, which does not depend on the table, with the lift partner, whose room letters are the letters of the lifts. The law of the universal piece is the image of a measure on the torus of the fusion bunch, of dimension $c-1$, and the law of the lift partner is a point only when its lift ratios are constant, which they are not at $c\ge2$: these two laws are the price of width. Down the partial fusions the first step is a Beta corner and every later step is a cell whose law is not a point. Every one-mark contraction of the fused table has the fusion width one and fuses back into it: these cousins share one period up to Beta corners, and an automorphism of the fused table carries the cousin at a mark onto the cousin at the image of that mark. Across a two-sum of lists the widths add, the condition is read at the hinge by four conditions on each cell, and at the total width one a two-sum fuses exactly when both cells fuse, the cell of the width zero being then identically self-dual. A table of the fusion width $c\ge0$ that does not fuse has uncovered marks, or has the width zero and the bases but not the weights of its dual; doubling the uncovered marks in parallel, or one mark in the second case, gives a table that fuses. The chapter ends with the branch these statements set at every node of the Foreman tree.

\section{Quotients, lifts and the level cell of a pair}\label{sec:fusionpair}

% New: proved in the last part.
\begin{definition}[Quotients; lifts; the level cell of a pair]\label{def:quotientlift}
Let $\mathcal M$ and $\mathcal N$ be tables on one set of marks $\cE$, of the ranks $r_{\mathcal M}$ and $r_{\mathcal N}$, with the rank functions $r_{\mathcal M}(Y)$ and $r_{\mathcal N}(Y)$ (Definition~\ref{def:molecule}). $\mathcal N$ is a \emph{quotient} of $\mathcal M$ when every set that is maximal for its rank in $\mathcal N$ is maximal for its rank in $\mathcal M$, as the child at $p+b$ is of the child at $p$ (Proposition~\ref{prop:adjacent}); its \emph{rank difference} is $c=r_{\mathcal M}-r_{\mathcal N}$. For $0\le j\le c$ the \emph{$j$-th lift of $\mathcal N$ toward $\mathcal M$} is the table $\mathcal L_j(\mathcal M,\mathcal N)$ on $\cE$ with the rank function
\[
r_{\mathcal L_j(\mathcal M,\mathcal N)}(Y)=\min\big(r_{\mathcal M}(Y),\ r_{\mathcal N}(Y)+j\big)\qquad(Y\subseteq\cE),
\]
a table by Proposition~\ref{prop:paircell}(iii) below; for a bunch $B$ of a table and $(\mathcal M,\mathcal N)=(\cB|B,\cB/R)$ these are the lifts of Proposition~\ref{prop:higgs}. For a set $G$ of $c$ new marks the \emph{level cell} of the pair $(\mathcal M,\mathcal N)$ is the table on $\cE\sqcup G$ whose bases are the sets $X\sqcup Q$ with $Q\subseteq G$ and $X$ a basis of $\mathcal L_{c-|Q|}(\mathcal M,\mathcal N)$; for $(\cB|B,\cB/R)$ it is the level cell of the bunch $B$ (Lemma~\ref{lem:higgsmajor}). It carries any positive weights, and $G$ is in level form in it (Proposition~\ref{prop:levelform}) when the weight of $X\sqcup Q$ depends on $X$ and $|Q|$ alone.
\end{definition}

% New: proved in the last part.
\begin{proposition}[The level cell of a pair]\label{prop:paircell}
Let $\mathcal N$ be a quotient of $\mathcal M$ on $\cE$ with the rank difference $c$, let $G$ be a set of $c$ new marks, and in this statement write $\mathcal L_j$ for $\mathcal L_j(\mathcal M,\mathcal N)$.
\begin{enumerate}[label=(\roman*)]
\item For tables $\mathcal M$ and $\mathcal N$ on $\cE$, $\mathcal N$ is a quotient of $\mathcal M$ exactly when $r_{\mathcal M}(Y)-r_{\mathcal M}(X)\ge r_{\mathcal N}(Y)-r_{\mathcal N}(X)$ for all $X\subseteq Y\subseteq\cE$. Then $c\ge0$, every set independent in $\mathcal N$ is independent in $\mathcal M$, every set spanning $\mathcal M$ spans $\mathcal N$, and a quotient of $\mathcal N$ is a quotient of $\mathcal M$.
\item The level cell of $(\mathcal M,\mathcal N)$ is a table of rank $r_{\mathcal M}$ with the rank function
\[
r^\natural(Y\sqcup P)=\min\big(r_{\mathcal M}(Y)+|P|,\ r_{\mathcal N}(Y)+c\big)\qquad(Y\subseteq\cE,\ P\subseteq G).
\]
In it every subset of $G$ is a pattern of the bunch $G$, $d(G)=0$ and $c(G)=c$; the child at $Q\subseteq G$ is $\mathcal L_{c-|Q|}$; the deletion of $G$ is $\mathcal M$ and the contraction of $G$ is $\mathcal N$; and every permutation of $G$ maps its bases onto themselves.
\item Each $\mathcal L_j$ is a table of rank $r_{\mathcal N}+j$ whose bases are the sets of $r_{\mathcal N}+j$ marks that are independent in $\mathcal M$ and span $\mathcal N$; $\mathcal L_0=\mathcal N$, $\mathcal L_c=\mathcal M$, and $\mathcal L_j$ is a quotient of $\mathcal L_{j'}$ for $j\le j'$. Lifts compose: for $0\le a\le b\le c$ and $0\le i\le b-a$,
\[
\mathcal L_i(\mathcal L_b,\mathcal L_a)=\mathcal L_{a+i}.
\]
\item For disjoint $C,D\subseteq G$ the minor of the level cell contracting $C$ and deleting $D$ is the level cell of the pair $(\mathcal L_{c-|C|},\mathcal L_{|D|})$ on $\cE\sqcup\big(G\setminus(C\cup D)\big)$.
\end{enumerate}
\end{proposition}

The level cell stores a whole chain of lifts on $c$ new marks: one level of the bunch $G$ for each lift, the deletion of $G$ at the top and its contraction at the bottom, and every interval of the chain as a minor at the marks of $G$. Nothing in its bases is chosen except the names of the new marks, which every permutation of $G$ moves without moving the bases; its weights are a choice.

\section{The fused table}\label{sec:fused}

% New: proved in the last part.
\begin{lemma}[Complements of the lifts toward the dual]\label{lem:liftcomplement}
Let $\cB$ be a table of rank $r$ on $\cE$, $n=|\cE|$, which is a quotient of its dual table $\cB^\ast$, of rank $r^\ast=n-r$ (Definition~\ref{def:quotientlift}), let $c=r^\ast-r$, and let $\mathcal L_j=\mathcal L_j(\cB^\ast,\cB)$, $0\le j\le c$. Then $X$ is a basis of $\mathcal L_j$ exactly when $\cE\setminus X$ is a basis of $\mathcal L_{c-j}$: the dual table of $\mathcal L_j$ has the bases of $\mathcal L_{c-j}$.
\end{lemma}

% New: proved in the last part.
\begin{definition}[Fusion width; Foreman fusion; majors]\label{def:fusion}
Let $(\cB,\theta)$ be a table of rank $r$ and corank $r^\ast$ on $\cE$, $n=|\cE|$, with its dual table $(\cB^\ast,\theta^\ast)$, $\theta^\ast_{\cE\setminus T}=\theta_T$ (Definition~\ref{def:table}), and with the rank function $r(Y)$ and the dual rank function $r^\ast(Y)$ (Definition~\ref{def:skeleton}). Its \emph{fusion width} is $c=r^\ast-r=n-2r$, and $(\cB,\theta)$ is \emph{fusable} when $\cB$ is a quotient of $\cB^\ast$ (Definition~\ref{def:quotientlift}) and, at the fusion width zero, $\theta^\ast=\theta$. A quotient forces $c\ge0$ (Proposition~\ref{prop:paircell}(i)), and a quotient of the same rank has the same bases, so a fusable table of the fusion width zero is identically self-dual (Corollary~\ref{cor:piecesmarks}(iii)), and a list of the fusion width zero is fusable exactly when it is identically self-dual. Let $c\ge0$, and let $G$ be a set of $c$ new marks, the \emph{fusion bunch}.
\begin{enumerate}[label=(\roman*)]
\item A \emph{major} of $(\cB,\theta)$ is a table on $\cE\sqcup G$ whose contraction of $G$ is $(\cB,\theta)$ and whose deletion of $G$ is $(\cB^\ast,\theta^\ast)$.
\item The \emph{lifts} of a fusable $(\cB,\theta)$ are $\mathcal L_j=\mathcal L_j(\cB^\ast,\cB)$, $0\le j\le c$, with the weights $\theta^{(j)}$: $\theta^{(0)}=\theta$ on $\mathcal L_0=\cB$, $\theta^{(c)}=\theta^\ast$ on $\mathcal L_c=\cB^\ast$, and for $0<j<c$ the \emph{middle weights}, positive weights with $\theta^{(c-j)}_{\cE\setminus X}=\theta^{(j)}_X$, complement mapping the bases of $\mathcal L_j$ onto those of $\mathcal L_{c-j}$ (Lemma~\ref{lem:liftcomplement}), a choice, all $1$ unless others are named; so the lifts of a list are lists unless other middle weights are named. Its \emph{lift letters} are $L_l=K_{\mathcal L_{c-l}}$, $0\le l\le c$, read with these weights, so that $L_0=\Psi$ and $L_c=K$, and its \emph{lift ratios} are $z_j=K_{\mathcal L_j}^2/(K_{\mathcal L_{j-1}}K_{\mathcal L_{j+1}})$, $0<j<c$.
\item The \emph{Foreman fusion} of a fusable $(\cB,\theta)$, with its middle weights, is the level cell of $(\cB^\ast,\cB)$ on $\cE\sqcup G$ with the weight $\theta^{(c-|Q|)}_X$ on $X\sqcup Q$, written $(\cB^\sharp,\theta^\sharp)$; $(\cB,\theta)$, the contraction of the fusion bunch in it, is its \emph{Foreman embedding}. For $G'\subseteq G$ the \emph{partial fusion} at $G'$ is $\cB^\sharp/(G\setminus G')$, written $\cB^\sharp_i$ for $|G'|=i$ when the marks of $G'$ are not named; $\cB^\sharp_0=\cB$ and $\cB^\sharp_c=\cB^\sharp$.
\end{enumerate}
At the fusion width zero $G$ is empty and $\cB^\sharp=\cB$. A table of negative fusion width $c$ has a dual of the fusion width $-c$, and the dual is the table that is fused.
\end{definition}

% New: proved in the last part.
\begin{theorem}[Foreman fusion]\label{thm:fusion}
Let $(\cB,\theta)$ be fusable, of the fusion width $c\ge0$, rank $r$ and corank $r^\ast$, with its Foreman fusion $(\cB^\sharp,\theta^\sharp)$ on $\cE\sqcup G$ and its lifts $\mathcal L_j$.
\begin{enumerate}[label=(\roman*)]
\item $\cB^\sharp$ is a table of rank $r^\ast$ on $n+c=2r^\ast$ marks with the rank function
\[
r^\sharp(Y\sqcup P)=\min\big(r^\ast(Y)+|P|,\ r(Y)+c\big)\qquad(Y\subseteq\cE,\ P\subseteq G),
\]
its contraction of $G$ is $(\cB,\theta)$ and its deletion of $G$ is $(\cB^\ast,\theta^\ast)$: it is a major of $(\cB,\theta)$.
\item The dual of each lift is a lift, weights included: $\mathcal L_j^\ast=\mathcal L_{c-j}$. So $(\cB^\sharp,\theta^\sharp)$ is identically self-dual (Corollary~\ref{cor:piecesmarks}(iii)), and
\[
K^\sharp=\Psi^\sharp=\sum_{l=0}^ce_l(\alpha_G)\,L_l(\alpha_{\cE}),\qquad L_l=K_{\mathcal L_{c-l}}=\Psi_{\mathcal L_l}.
\]
\item Every basis of every major of $(\cB,\theta)$ is a basis of $\cB^\sharp$.
\item Every major of $(\cB,\theta)$ in which $G$ is in level form, equivalently that every permutation of $G$ carries, has the bases of $\cB^\sharp$, the weights $\theta$ on the bases $X\sqcup G$ and $\theta^\ast$ on the bases $X\sqcup\varnothing$, and on the others positive weights depending on $X$ and $|Q|$ alone; it is identically self-dual exactly when these weights, read as middle weights $\theta^{(c-|Q|)}_X$ on the bases $X\sqcup Q$, are paired as in Definition~\ref{def:fusion}(ii), and then it is the Foreman fusion with these middle weights. So the identically self-dual majors in level form are the Foreman fusions with all the choices of middle weights, and for a list $\cB^\sharp$ is the only one that is a list. For every bijection $\beta$ of $\cE$ carrying $(\cB,\theta)$ with $\theta^{(j)}\circ\beta=\theta^{(j)}$ for $0<j<c$, as for the middle weights $1$, and every permutation $\pi$ of $G$, $\beta\sqcup\pi$ carries $(\cB^\sharp,\theta^\sharp)$.
\item At the fusion width one, $G=\{g\}$, the bases of $\cB^\sharp$ are the sets $T+g$ and $\cE\setminus T$, $T\in\cB$, both with the weight $\theta_T$, $K^\sharp=\Psi+\alpha_gK$, and $(\cB^\sharp,\theta^\sharp)$ is the only major of $(\cB,\theta)$.
\end{enumerate}
\end{theorem}

The fusion stores the table and its dual as the two extreme children of one bunch, and fills the levels between them with the lifts, which is the only way to fill them in level form that the exchange property allows (Proposition~\ref{prop:gluenec}(d)): the containment (iii) says that no major has a basis the fusion lacks, and level form at $G$ selects the bases of the fusion among the majors, its weights between the extremes being a choice paired by complement. Self-duality comes from the chain of lifts turning over under duality, the $j$-th lift onto the $(c-j)$-th. At the fusion width one there is no level between the extremes, and the fusion is the only major; at the width $c$ the middle levels admit majors that are not in level form, among them identically self-dual ones, and the bases of the fusion are the ones that every permutation of the new marks keeps.

% New: proved in the last part.
\begin{theorem}[The tell]\label{thm:fusiontell}
For a table $(\cB,\theta)$ of the fusion width $c$, conditions (a), (b) and (c) are equivalent and imply $c\ge0$. The table is fusable exactly when they hold and, at $c=0$, $\theta^\ast=\theta$; so a list, or a table of the fusion width $c\ge1$, is fusable exactly when (a) holds. For $c\ge0$ the table is fusable exactly when (d) holds.
\begin{enumerate}[label=(\alph*)]
\item $\cB$ is a quotient of $\cB^\ast$.
\item Every cocircuit of $\cB$ is a union of circuits of $\cB$ (Definition~\ref{def:molecule}).
\item For every mark $e$ and every split $\cE\setminus\{e\}=A\sqcup A'$, the rank of $A$ or the rank of $A'$ does not grow when $e$ is added to it.
\item $(\cB,\theta)$ has a major (Definition~\ref{def:fusion}(i)); for $c\ge1$, equivalently, some table on $\cE\sqcup G$, $|G|=c$, has $\cB^\ast$ and $\cB$ as the children of the empty pattern and of the full pattern of its bunch $G$.
\end{enumerate}
A table that satisfies (b) has no coloop. For $c<0$ the dual table has the fusion width $-c$, and the theorem read on it has (b) in the form: every circuit of $\cB$ is a union of cocircuits.
\end{theorem}

The tell is read on the Skeleton of the table (Definition~\ref{def:skeleton}): the cocircuit layer must be covered by the circuit layer, cocircuit by cocircuit, before anything is integrated. In the form (c) it is a statement about one mark at a time: however the other marks are split in two, one half spans the mark. A coloop fails it at once, being a cocircuit of one mark; a loop lies in no cocircuit and never fails it. The weights enter only at the width zero, where the bases of a table that passes are those of its dual and the table fuses when its weights are those of its dual as well.

% New: proved in the last part.
\begin{corollary}[The fused tables are the identically self-dual tables]\label{cor:fusionisd}
\begin{enumerate}[label=(\roman*)]
\item Let $(\mathcal H,\theta)$ be identically self-dual on $\cE'$ and $C\subseteq\cE'$ independent, every transposition of two marks of $C$ carrying $(\mathcal H,\theta)$. Then $\mathcal H/C$ is fusable of the fusion width $|C|$, and $(\mathcal H,\theta)$ is its Foreman fusion with the fusion bunch $C$ and the middle weights read off $\theta$. In particular every identically self-dual table on at least one mark is the Foreman fusion of each of its one-mark contractions, and the fused tables are exactly the identically self-dual tables.
\item If $\cB$ is connected and fusable, $\cB^\sharp$ is connected.
\end{enumerate}
\end{corollary}

% New: proved in the last part.
\begin{proposition}[Partial fusions and inner lifts]\label{prop:fusionminors}
Let $(\cB,\theta)$ be fusable of the fusion width $c\ge1$, with its lifts $\mathcal L_j$.
\begin{enumerate}[label=(\roman*)]
\item For disjoint $C,D\subseteq G$, $\cB^\sharp/C\setminus D$ is the level cell of $(\mathcal L_{c-|C|},\mathcal L_{|D|})$ on $\cE\sqcup\big(G\setminus(C\cup D)\big)$, with the weights of $\cB^\sharp$.
\item The partial fusion at $G'$, $|G'|=i$, is the level cell of $(\mathcal L_i,\cB)$ on $\cE\sqcup G'$. It has the fusion width $c-i$, it is fusable, and its Foreman fusion with the middle weights read off $\theta^\sharp$ is $(\cB^\sharp,\theta^\sharp)$ with the fusion bunch $G\setminus G'$. So $\cB=\cB^\sharp_0,\cB^\sharp_1,\dots,\cB^\sharp_c=\cB^\sharp$, each the contraction of the next at one mark, all fuse into one table.
\item $\mathcal L_j$ has the fusion width $c-2j$. For $2j\le c$ it is fusable, with the Foreman fusion $\cB^\sharp/C\setminus D$, $|C|=|D|=j$, its middle weights read off $\theta^\sharp$. For even $c$ the middle lift $\mathcal L_{c/2}$ is identically self-dual; for odd $c$ the lift $\mathcal L_{(c-1)/2}$ has the fusion width one and fuses into $\cB^\sharp/C\setminus D$ with $|C|=|D|=(c-1)/2$.
\end{enumerate}
\end{proposition}

So the fused table holds, at its minors on the fusion bunch, the fusion of every table in the chain between $\cB$ and $\cB^\ast$: each partial fusion fuses into it again, and each inner lift fuses into the minor that deletes as many marks of $G$ as it contracts.

\begin{remark}[Examples]\label{rem:fusionexamples}
(i) The uniform list $\binom{\cE}r$ with $n\ge2r$ is fusable, its cocircuits having $r^\ast+1\ge r+1$ marks; its lifts are the uniform lists $\binom{\cE}{r+j}$, and its fusion is $\binom{\cE\sqcup G}{r^\ast}$: the fusion of $\binom{[n]}r$ is $\binom{[2r^\ast]}{r^\ast}$, and the fusion of $\binom{[5]}2$ is $\binom{[6]}3$; for $\binom{[2r+1]}r$, $r\ge2$, the fusion $\binom{[2r+2]}{r+1}$ is the table one mark larger of Proposition~\ref{prop:onemarkhidden}(ii).

(ii) A list of rank two with no loop has as its bases the pairs of marks in different classes of a partition of its marks, the classes of two or more marks being its parallel classes (Definition~\ref{def:width}). With $k\ge3$ classes it is fusable exactly when $k\ge4$, or $k=3$ and every class has at least two marks: its cocircuits are the complements of the classes, and a mark of a class $P'$ lies in a circuit inside the complement of another class $P$ exactly when $P'$ has a second mark or two classes other than $P$ and $P'$ exist. So the list $\binom{[3]}2$ and every such list with three classes one of which is a single mark fail, and those with four or more classes pass at every fusion width (receipt \texttt{t03\_rank2\_census.py}, every connected list of rank two on $5$ to $9$ marks).

(iii) The list of rank two on $n\ge5$ marks with one parallel pair $P$ and every other mark in a class of its own has the fusion width $n-4$, and its fusion has exactly two sets of $n-2$ marks that are not bases, the complementary halves $P\sqcup G$ and $\cE\setminus P$; its automorphisms are the $2((n-2)!)^2$ bijections keeping this pair of halves, transitive on the marks. At $n=5$ the fusion is the two-sum of two copies of $\binom{[4]}2$.

(iv) Of the $108$ lists of rank three on seven marks, up to the names of the marks, $22$ are fusable, of the fusion width one, and they fuse into the $17$ identically self-dual lists on eight marks; of the $18$ that are prime (Theorem~\ref{thm:primetree}), $10$ are fusable, those whose lines all have three marks and meet pairwise (receipts \texttt{census7.py}, \texttt{atoms7.py}).

(v) The majors of $\binom{[4]}1$ that are lists, of the fusion width two, are $79$ lists on six marks; $7$ of them are identically self-dual and one is in level form at $G$, and that one is the fusion $\binom{[6]}3$ (receipt \texttt{t02\_majors\_U14.py}). At rank three the following are fusable, of the fusion widths two and three: $\binom{[8]}3$ and $\binom{[9]}3$; the list $\mathrm F_7$ on the marks $\mathbb Z/7$ whose dependent sets of three marks are the seven sets $\{i,i+1,i+3\}$, with a free mark added; the list $\mathrm A_9$ on the marks $\mathbb F_3^2$ whose dependent sets of three marks are the twelve sets $\{a,a+d,a+2d\}$, $d\neq0$, and $\mathrm A_9$ less a mark; the list $\mathrm Z_9$ on the marks $a_i,b_i,c_i$, $i=1,2,3$, whose dependent sets of three marks are $\{a_1,a_2,a_3\}$, $\{b_1,b_2,b_3\}$, $\{c_1,c_2,c_3\}$ and the six sets $\{a_i,b_j,c_{i+j-2}\}$ with $i\neq j$; and five further lists of rank three whose dependent sets of three marks form configurations of lines; and their fusions have the seam of Proposition~\ref{prop:fusionseam} (receipt \texttt{w01\_seam\_rank3.py}).
\end{remark}

\section{The seam of the fusion bunch and the glue}\label{sec:fusionseam}

% New: proved in the last part.
\begin{proposition}[The seam of the fusion bunch]\label{prop:fusionseam}
Let $(\cB,\theta)$ be fusable of the fusion width $c\ge1$, with its fusion $(\cB^\sharp,\theta^\sharp)$, its lifts $\mathcal L_j$ and its lift letters $L_l$.
\begin{enumerate}[label=(\roman*)]
\item In $\cB^\sharp$ the bunch $G$ has $\pat(G)=2^G$, $d(G)=0$ and $c(G)=c$, and it is in level form. At the excess $l$ (Definition~\ref{def:seam}) there is one station of patterns, all $Q\subseteq G$ with $|Q|=l$, with the multipliers $1$ and the station letters $e_l(\alpha_G)$ and $e_{c-l}(\alpha_G)$, and one station $D_l$ of room sides, the bases of $\mathcal L_{c-l}$, whose room children are positive multiples of $\binom Gl$. Every excess is one rectangle, $\cB^\sharp$ is its own level envelope (Theorem~\ref{thm:envelope}), and the room-station factorization (Proposition~\ref{prop:seamidentity}(iii)) is
\[
K^\sharp=\sum_{l=0}^ce_l(\alpha_G)\,L_l(\alpha_{\cE}),\qquad \Psi^\sharp=\sum_{l=0}^ce_{c-l}(\alpha_G)\,L_{c-l}(\alpha_{\cE}).
\]
The seam letters of $G$ are $K_{(D_l)}=\mu_le_l(\alpha_G)$ and $\Psi_{(D_l)}=\mu_le_{c-l}(\alpha_G)$, with $\mu_l$ the largest weight of $\mathcal L_{c-l}$, the room-station letters are $L_l/\mu_l$ and $L_{c-l}/\mu_l$, and the pair constants are $\mu_l$; for a list every $\mu_l$ is $1$.
\item The interaction of $G$ (Definition~\ref{def:interaction}) and its dual interaction are
\[
\iota_G=\iota^\ast_G=\sum_{j=0}^c(-1)^j\binom cj\log K_{\mathcal L_j}(\alpha_{\cE}),
\]
and $\iota_G=-\log z_1$ at $c=2$.
\item At the table's gauge the pattern probabilities of the excess $l$ add up to $\Pi_l=\binom cl\,L_l(\one)/K^\sharp(\one)$, with $\Pi_l=\Pi_{c-l}$, and the peel of $((\cB^\sharp,\theta^\sharp);s,t)$ at $G$ (Definition~\ref{def:peel}) is the family of the lifts $((\mathcal L_j,\theta^{(j)});\Pi_js,\Pi_jt)_{0\le j\le c}$, whose word is $\prod_jK_{\mathcal L_j}^{-2\sigma\Pi_j}$.
\end{enumerate}
\end{proposition}

Between the extremes the seam of a bunch can carry several rectangles (Proposition~\ref{prop:ladder}(iii)); the fusion bunch carries one at every excess, with the multipliers $1$ on its patterns and the pair constants $\mu_l$, and what it presents to its room at the excess $l$ is the elementary symmetric polynomial of degree $l$ of its own marks. These letters do not depend on the table. Everything the table contributes sits on the other side of the seam, in the letters of its lifts.

% New: proved in the last part.
\begin{definition}[The universal piece; the lift partner]\label{def:universalpiece}
For $c\ge1$ the \emph{universal piece of width $c$} is the piece on a set $G$ of $c$ marks for one table (Definition~\ref{def:seampiece}) with the stations $D_0,\dots,D_c$, $D_l$ of the excess $l$, the seam letters $V_{D_l}=e_l(\alpha_G)$ and $V^\ast_{D_l}=e_{c-l}(\alpha_G)$, and $d=d^\ast=0$. The \emph{lift partner} of a fusable $(\cB,\theta)$ of the fusion width $c$ is the partner on $\cE$ with the same stations, the room letters $L_{D_l}=L_l$ and $L^\ast_{D_l}=L_{c-l}$, and $e=e^\ast=r$.
\end{definition}

By Proposition~\ref{prop:fusionseam}(i) and Proposition~\ref{prop:seamidentity}(iii) the glue of the universal piece of width $c$ and the lift partner of $(\cB,\theta)$ is the fusion $(\cB^\sharp,\theta^\sharp)$. Since $\cB^\sharp$ is identically self-dual its word at the exponents $(s,t)$ is $K^{\sharp\,-2\sigma}$, $\sigma=\frac{s+t}2$ the strength (Definition~\ref{def:family}), and the fused table is read below as the family $((\cB^\sharp,\theta^\sharp);2\sigma,0)$, with the one active letter $K^\sharp$; its walls are those at $(s,t)$, $d$ and $d^\ast$ agreeing on an identically self-dual table. Its balance wall is $m_{\cE}+m_G=2\sigma r^\ast$. In this section $H_G=\{w\in\mathbb R^G:\ \sum_gw_g=0\}$ is the torus of $G$ in logarithmic coordinates, with the measure $d\lambda_G(w)=dw_{g_1}\cdots dw_{g_{c-1}}$ for an order $g_1,\dots,g_c$ of $G$, and $\mathbf e(y)=(e_0(y),\dots,e_c(y))$ for $y\in(0,\infty)^G$.

% New: proved in the last part.
\begin{theorem}[The glue of a fused table]\label{thm:fusionglue}
Let $(\cB,\theta)$ be fusable of the fusion width $c\ge1$, with its fusion $(\cB^\sharp,\theta^\sharp)$ at a strength $\sigma>0$ and masses $\mvec$ on $\cE\sqcup G$ on its balance wall. For $\mathbf x=(x_0,\dots,x_c)\in(0,\infty)^{c+1}$ let $W[\mathbf x]$ be the Newton word on $\cE\sqcup\{g\}$, $g$ the stand-in of the universal piece, with the one letter $\sum_{l=0}^cx_l\,\alpha_g^{\,l}\,L_l(\alpha_{\cE})$ at the exponent $2\sigma$ (Definition~\ref{def:posword}).
\begin{enumerate}[label=(\roman*)]
\item \emph{The universal law.} The scale of the universal piece is $\gamma^\rho_G=e_c(\alpha_G)^{1/c}$, its seam vector at the excess $l$ is $e_l(\alpha_G)/\gamma^\rho_G(\alpha_G)^l$, and its law is the image of the measure $c\,e^{\ip{\mvec|_G}w}\,d\lambda_G(w)$ on $H_G$ under $w\mapsto\mathbf e(e^w)$.
\item \emph{The torus form.} In $[0,\infty]$,
\[
P_{\cB^\sharp}(\mvec;s,t)=c\int_{H_G}e^{\ip{\mvec|_G}w}\,P_{W[\mathbf e(e^w)]}\big(\mvec|_{\cE},m_G\big)\,d\lambda_G(w),
\]
the stand-in carrying the mass $m_G$, the wall of the universal piece.
\item \emph{The scale of the stand-in.} At a fixed $\alpha_{\cE}$ write $\sum_lx_lL_l\,y^l=x_cL_c\prod_{i=1}^c(y+\xi_i)$; the $\xi_i$, in $\mathbb C\setminus(-\infty,0]$, are the Skeleton values of the stand-in, read on the letter of $W[\mathbf x]$ as in Definition~\ref{def:skelpoly}. For $0<m_G<2\sigma c$,
\[
\int_0^\infty\rho^{m_G-1}\Big(\sum_lx_l\rho^lL_l\Big)^{-2\sigma}d\rho=(x_cL_c)^{-2\sigma}\,\Gst{2}{c}{m_G,\ 2\sigma c-m_G}{2\sigma,\ \dots,\ 2\sigma}\,\sB^{\{J\}}_{(2\sigma)_{i\in J},\,2\sigma c-m_G;\,I_J}(\boldsymbol\xi),
\]
the Breta of one bunch $J=\{1,\dots,c\}$ of Proposition~\ref{prop:scalebreta}, coincident Skeleton values merging.
\item \emph{One share per level.} On every contour of shares $\mathbf a=(a_0,\dots,a_c)$ with positive real parts on the pins $\sum_la_l=2\sigma$ and $\sum_ll\,a_l=m_G$ at which the two periods below converge, as at the native split at the stations when every proper wall of $\cB^\sharp$ is positive (Theorem~\ref{thm:gluewidth}(iii)),
\[
P_{\cB^\sharp}(\mvec;s,t)=\int_{\sum_la_l=2\sigma,\ \sum_ll\,a_l=m_G}\Gst{c+1}1{\mathbf a}{2\sigma}\,P_{\prod_{l=1}^ce_l^{-a_l}}\big(\mvec|_G\big)\,P_{\prod_{l=0}^cL_l^{-a_l}}\big(\mvec|_{\cE}\big),
\]
a fibre integral of the dimension $c-1$.
\item \emph{The universal factor.} For $c\ge2$, $P_{\prod_{l=1}^ce_l^{-a_l}}(\mvec|_G)=P_{\prod_{l=1}^{c-1}e_l^{-a_l}}(\mvec|_G-a_c\one)$, the period on $G$ of the uniform lists $\binom Gl$, $0<l<c$, at the exponents $(a_l,0)$ and the masses $\mvec|_G-a_c\one$; at $c=1$ it is $1$.
\end{enumerate}
\end{theorem}

The universal piece is absorbed into its stand-in at its own scale, the geometric mean of the scales of $G$, and what is left of it is a point of the torus of $G$, entering the letter of the stand-in through the elementary symmetric polynomials of the exponentials of that point. The stand-in enters with the degree $c$, so $W[\mathbf x]$ is not the word of a table for $c\ge2$ (Theorem~\ref{thm:storeone}), and integrating its scale is one Breta over the $c$ Skeleton values of the stand-in. Opened on its shares, the fused table carries one share per level and two pins, the strength and the scale, so $c-1$ shares are free; the universal factor is the period of the uniform lists on $G$, and everything that depends on the table is the period of the family of its lifts, in which the table itself is the member with $a_1=\dots=a_{c-1}=0$, $P_{L_0^{-a_0}L_c^{-a_c}}(\mvec|_{\cE})=P_{\cB}(\mvec|_{\cE};a_c,a_0)$.

% New: proved in the last part.
\begin{corollary}[Width one]\label{cor:fusionone}
At the fusion width one, on the balance wall of $\cB^\sharp$ with every proper wall positive,
\[
P_{\cB^\sharp}(\mvec|_{\cE},m_g;s,t)=\Gst21{m_g,\ 2\sigma-m_g}{2\sigma}\,P_{\cB}\big(\mvec|_{\cE};m_g,2\sigma-m_g\big):
\]
the exponents of the table are read off the mass of the fusion bunch. The universal piece is the mark $g$ itself, the laws of the universal piece and of the lift partner are points, and the share form of Theorem~\ref{thm:fusionglue}(iv) has no free share.
\end{corollary}

% New: proved in the last part.
\begin{corollary}[Width two]\label{cor:fusiontwo}
Let $c=2$, $G=\{g_1,g_2\}$, with the masses $m_1,m_2$ on $g_1,g_2$.
\begin{enumerate}[label=(\roman*)]
\item On a line of one free share $a_2$, with $a_1=m_1+m_2-2a_2$ and $a_0=2\sigma-m_1-m_2+a_2$,
\[
P_{\cB^\sharp}(\mvec;s,t)=\frac1{2\pi i}\int\frac{\Gamma(a_0)\,\Gamma(a_2)\,\Gamma(m_1-a_2)\,\Gamma(m_2-a_2)}{\Gamma(2\sigma)}\,P_{L_0^{-a_0}L_1^{-a_1}L_2^{-a_2}}\big(\mvec|_{\cE}\big)\,da_2:
\]
the universal factor is the Beta corner $\Gst21{m_1-a_2,\ m_2-a_2}{a_1}$, and the Gamma of the middle level cancels.
\item The universal law is the image of $2v^{m_1-m_2-1}dv$ on $(0,\infty)$ under $v\mapsto(1,v+v^{-1},1)$, and
\[
P_{\cB^\sharp}(\mvec;s,t)=2\int_0^\infty v^{m_1-m_2-1}\,P_{(L_0+(v+v^{-1})\alpha_gL_1+\alpha_g^2L_2)^{-2\sigma}}\big(\mvec|_{\cE},m_1+m_2\big)\,dv .
\]
\item At a fixed $\alpha_{\cE}$, for $0<m_1<2\sigma$ and $0<m_2<2\sigma$,
\begin{multline*}
\int_{(0,\infty)^2}\alpha_{g_1}^{m_1-1}\alpha_{g_2}^{m_2-1}K^{\sharp\,-2\sigma}\,d\alpha_G\\
=\Gst21{m_2,\ 2\sigma-m_2}{2\sigma}L_0^{m_1+m_2-2\sigma}L_1^{-m_1-m_2}\int_0^\infty x^{m_1-1}(1+x)^{m_2-2\sigma}\big(1+z_1^{-1}x\big)^{-m_2}dx,
\end{multline*}
a Beta corner times a Breta of one bunch with the two coefficients $1$ and $z_1$ (Proposition~\ref{prop:scalebreta}). At $m_1=m_2=1$, $2\sigma=2$ it is $\log z_1/(L_1^2-L_0L_2)$ where $z_1\neq1$ and $L_1^{-2}$ where $z_1=1$, and at integer $m_1,m_2,2\sigma$, where $z_1\neq1$, it is $L_0^{m_1+m_2-2\sigma}L_1^{-m_1-m_2}\big(R_1(z_1)+R_2(z_1)\log z_1\big)$ with rational functions $R_1,R_2$.
\end{enumerate}
\end{corollary}

\section{The price of width}\label{sec:fusionprice}

% New: proved in the last part.
\begin{lemma}[Products of lifts]\label{lem:liftproducts}
Let $\mathcal N$ be a quotient of $\mathcal M$ with the rank difference $c$, its lifts $\mathcal L_j=\mathcal L_j(\mathcal M,\mathcal N)$ carrying any positive weights. For $0\le i'<i\le j<j'\le c$ with $i+j=i'+j'$, $K_{\mathcal L_i}K_{\mathcal L_j}$ has a monomial in which exactly $r_{\mathcal N}+i$ exponents equal $2$, while in every monomial of $K_{\mathcal L_{i'}}K_{\mathcal L_{j'}}$ at most $r_{\mathcal N}+i'$ do; so the two products are not proportional. In particular, for $c\ge2$, no ratio $K_{\mathcal L_j}^2/(K_{\mathcal L_{j-1}}K_{\mathcal L_{j+1}})$, $0<j<c$, is constant on the orthant.
\end{lemma}

% New: proved in the last part.
\begin{theorem}[The price of width]\label{thm:fusionprice}
Let $(\cB,\theta)$ be fusable of the fusion width $c\ge1$, with its fusion read at a strength $\sigma>0$ as in Theorem~\ref{thm:fusionglue}, its lift letters $L_l$ and its lift ratios $z_j$.
\begin{enumerate}[label=(\roman*)]
\item The law of the universal piece is carried by one point exactly when $c=1$. The law of the lift partner is carried by one point exactly when every lift ratio is constant, that is exactly when $c=1$.
\item \emph{The Skeleton braid.} Write $\sum_l\binom clL_l\,x^l=L_0\prod_{i=1}^c(1+\eta_ix)$; the $\eta_i$ are the reciprocals of the Skeleton values of $G$ at the gauge $(\one_G,\alpha_{\cE})$ (Definition~\ref{def:skelpoly}). Then
\[
K^\sharp=\frac{L_0}{c!}\sum_{\pi\in S_c}\prod_{i=1}^c\big(1+\alpha_{g_i}\eta_{\pi(i)}\big)=\frac{L_0}{c!}\operatorname{per}\big(\one\one^{\mathsf T}+\alpha_G\boldsymbol\eta^{\mathsf T}\big),
\]
$\one\one^{\mathsf T}$ the $c\times c$ matrix of ones and $\operatorname{per}$ the permanent: $K^\sharp$ is $L_0$ times the average, over the bijections from the marks of $G$ to the Skeleton values, of products of $c$ letters of one mark. At a gauge $\alpha_{\cE}$ the $\eta_i$ coincide exactly when every lift ratio is $1$ there, and for $c\ge2$ they coincide on no open set of the orthant: $K^\sharp$ is not $L_0\prod_{g\in G}(1+\alpha_g\eta)$ for any one function $\eta$.
\item \emph{The partial fusions.} For $1\le i\le c$, $G'\subseteq G$ with $|G'|=i$ and $g\in G'$, read the partial fusion at $G'$ as a cell with the hinge $g$ (Definition~\ref{def:substitution}). For $i=c$ its ratio is $1$, and on the balance wall of $\cB^\sharp$ with every proper wall positive
\[
P_{\cB^\sharp}(\mvec;s,t)=\Gst21{m_g,\ 2\sigma-m_g}{2\sigma}\,P_{\cB^\sharp/g}\big(\mvec|_{\cE\sqcup G-g};\,m_g,\,2\sigma-m_g\big);
\]
for $i<c$ its ratio is not constant, and at any exponents $s,t>0$ and masses its law is not carried by one point (Corollary~\ref{cor:piecesmarks}(i)).
\end{enumerate}
\end{theorem}

At the fusion width one both laws are points and the glue is a product: one Beta corner. At the width $c\ge2$ neither is: the universal law spreads over the $(c-1)$-dimensional torus of $G$, and the lift partner's seam vector moves with the room because the lifts are not a geometric chain, the lift ratios being not constant. Read on the Skeleton values of the fusion bunch, the letter of the fused table is an average over the bijections from the marks of $G$ to the Skeleton values, and it is one product of $c$ letters of one mark on no open set, and at no gauge for a list (Remark~\ref{rem:fusionlorentz}(i)); one product is what $c$ independent fusions of width one would give. Read down the partial fusions, only the step at the top is free.

Among the plain extensions of a table by one mark with connected hinge minors, the fusion is the only one that is free in this way, up to one constant on the weights of the bases avoiding the new mark.

% New: proved in the last part.
\begin{proposition}[The free extensions by one mark]\label{prop:freeextension}
Let $\mathcal A$ be a plain cell on $\cE\sqcup\{g\}$, $|\cE|\ge2$, whose hinge minors $\mathcal A\setminus g$ and $\mathcal A/g$ are connected. Its ratio $\varrho_{\mathcal A}$ (Definition~\ref{def:substitution}) is constant exactly when $\mathcal A/g$ is the dual of $\mathcal A\setminus g$ with its weights multiplied by one constant $\mu>0$; then $\varrho_{\mathcal A}\equiv\mu^2$, $\mathcal A/g$ has the fusion width one, and $\mu K_{\mathcal A}(\alpha_{\cE},\alpha_g)=K^\sharp(\alpha_{\cE},\mu\alpha_g)$, $K^\sharp$ the letter of the Foreman fusion of $\mathcal A/g$: $\mathcal A$ is that fusion with the weights of its bases avoiding $g$ divided by $\mu$. For a list, $\varrho_{\mathcal A}$ is constant exactly when $\mathcal A$ is identically self-dual, that is when $\mathcal A$ is the Foreman fusion of $\mathcal A/g$. So, at any exponents $s,t>0$ and masses, the law of $\mathcal A$ is a point (Corollary~\ref{cor:piecesmarks}(i)) exactly when $\mathcal A$ is the Foreman fusion of $\mathcal A/g$ up to one constant on the weights of the bases avoiding $g$, and for a list exactly when it is that fusion.
\end{proposition}

% New: proved in the last part.
\begin{corollary}[A pinned set of one list is a dual pair]\label{cor:mot-pinneddual}
Let $\cB$ be a list on $\cE$, let $s,t>0$, and let $S\subsetneq\cE$, $|S|\ge2$, with $S'=\cE\setminus S$, be such that $\cB|S$ and $\cB/S'$ are connected. Then $S$ is pinned (Theorem~\ref{thm:mot-pointmass}) exactly when it has width one and $\cB/S'=(\cB|S)^\ast$: the cell of the bunch $\cB^+_S$ (Theorem~\ref{thm:widthone}) is a dual pair.
\end{corollary}

\begin{remark}[The lift ratios at a gauge; width three]\label{rem:fusionlorentz}
(i) Let $\cB$ be a list. The room sides of $G$ in $\cB^\sharp$ are the sets independent in $\cB^\ast$ that span $\cB$: the bases of the identity map from $\cB^\ast$ to its quotient $\cB$ read as a morphism of matroids, the arena property that Theorem~\ref{thm:fusiontell}(a) supplies. For such a morphism Eur and Huh (\emph{Logarithmic concavity for morphisms of matroids}, Corollary~4.6) show that the basis generating polynomial, each basis counted once and homogenized by one more variable to the power $n$ less the size of the basis, is Lorentzian, and a linear change of the variables with nonnegative coefficients keeps that property (Br\"and\'en and Huh), so at every gauge $\alpha_{\cE}$ the values $K_{\mathcal L_j}(\alpha_{\cE})\big/\binom n{r+j}$ form a log-concave sequence:
\[
z_j(\alpha_{\cE})\ \ge\ \frac{(k+1)(n-k+1)}{k(n-k)}>1,\qquad k=r+j,\ 0<j<c .
\]
So for a list at $c\ge2$ the lift ratios exceed one at every gauge and the $\eta_i$ of Theorem~\ref{thm:fusionprice}(ii) never coincide. The bound holds on the $6720$ exact ratios of the receipt \texttt{t12\_ratios\_and\_tell.py}, with equality in $35$ of them, and it is attained at the table's gauge whenever $\mathcal L_{j-1}$, $\mathcal L_j$ and $\mathcal L_{j+1}$ are uniform lists, as for every uniform list: there $K_{\mathcal L_i}(\one)=\binom n{r+i}$. The receipt \texttt{t08\_logconcave.py} finds every lift ratio above one on $6720$ exact ratios of its own. With weights the lift ratios can fall below one: $\binom{[4]}1$ with the weights $2$ and the middle weights $1$ has $K_{\mathcal L_0}(\one)=K_{\mathcal L_2}(\one)=8$ and $K_{\mathcal L_1}(\one)=6$, so $z_1=\frac9{16}$ at the table's gauge.

(ii) At $c=2$ the discriminant of $L_0+2L_1x+L_2x^2$ is $4L_0L_2(z_1-1)$, so the Skeleton values of $G$ at $(\one_G,\alpha_{\cE})$ are real and distinct where $z_1>1$, at every gauge for a list, and complex conjugate where $z_1<1$.

(iii) At $c=3$, with the scale $a$ of one mark of $G$ held fixed, $K^\sharp$ is the letter of width two with the lift letters $L'_l=L_l+aL_{l+1}$, $0\le l\le2$; so at the masses $1$ on $G$ and $2\sigma=2$ the integral over the scales of $G$ at a fixed room point is, by Corollary~\ref{cor:fusiontwo}(iii),
\[
\int_0^\infty\frac{\log\big(L'^2_1/(L'_0L'_2)\big)}{L'^2_1-L'_0L'_2}\,da,\qquad L'^2_1-L'_0L'_2=-\big((L_1L_3-L_2^2)a^2+(L_0L_3-L_1L_2)a+(L_0L_2-L_1^2)\big),
\]
the coefficients of the quadratic being the $2\times2$ minors of the matrix with the rows $(L_0,L_1,L_2)$ and $(L_1,L_2,L_3)$, and its discriminant $-\frac1{27}$ times that of the diagonal level cubic $L_0+3L_1x+3L_2x^2+L_3x^3$ (receipts \texttt{t10\_weight.py}, a rendering, and \texttt{t13\_faces\_and\_cubic.py}). The cubic is real-rooted at all $560$ gauges sampled on $8$ lists of the fusion width three (receipts \texttt{w02\_skeleton\_braid.py}, \texttt{w03\_real\_c3.py}).
\end{remark}

\section{Cousins and hidden symmetry}\label{sec:fusionbuys}

% New: proved in the last part.
\begin{proposition}[Cousins]\label{prop:cousins}
The \emph{cousins} of an identically self-dual table are its one-mark contractions. Let $(\cB,\theta)$ be fusable of the fusion width $c\ge1$, and $x$ a mark of $\cE\sqcup G$.
\begin{enumerate}[label=(\roman*)]
\item The cousin $\cB^\sharp/x$ has the fusion width one and is fusable, $\cB^\sharp\setminus x$ is its dual, and its Foreman fusion is $(\cB^\sharp,\theta^\sharp)$ with the fusion bunch $\{x\}$. For $x\in G$ the cousin is the partial fusion $\cB^\sharp_{c-1}$, and at $c=1$ the table $\cB$ is the cousin at $g$.
\item On the balance wall of $\cB^\sharp$ with every proper wall positive, for every two marks $x,y$,
\begin{align*}
P_{\cB^\sharp}(\mvec;s,t)&=\Gst21{m_x,\ 2\sigma-m_x}{2\sigma}\,P_{\cB^\sharp/x}\big(\mvec_{-x};m_x,2\sigma-m_x\big)\\
&=\Gst21{m_y,\ 2\sigma-m_y}{2\sigma}\,P_{\cB^\sharp/y}\big(\mvec_{-y};m_y,2\sigma-m_y\big),
\end{align*}
$\mvec_{-x}$ the masses of the marks other than $x$: the cousins share one period up to Beta corners, each reading its exponents off the mass of the mark it lacks.
\item For every bijection $\beta$ of $\cE\sqcup G$ carrying $(\cB^\sharp,\theta^\sharp)$, $P_{\cB^\sharp}(\beta^\ast\mvec;s,t)=P_{\cB^\sharp}(\mvec;s,t)$, with $(\beta^\ast\mvec)_e=m_{\beta(e)}$. At $c=1$ the function
\[
(\mvec|_{\cE},s')\ \mapsto\ \Gst21{s',\ 2\sigma-s'}{2\sigma}\,P_{\cB}\big(\mvec|_{\cE};s',2\sigma-s'\big)
\]
on the part of the hyperplane $m_{\cE}+s'=2\sigma r^\ast$ where every proper wall of $\cB^\sharp$ is positive, which the automorphisms keep, is invariant under the automorphisms of $\cB^\sharp$ acting on $\cE\sqcup\{g\}$, $s'$ read as the mass of $g$: an automorphism $\beta$ with $\beta(g)=e\neq g$ makes the mass $m_e$ the basis exponent and gives $s'$ to the mark $\beta^{-1}(g)$ as its mass, and the transposition of $g$ and $e$, when it is an automorphism, exchanges $s'$ and $m_e$.
\end{enumerate}
\end{proposition}

What fusion buys at every width is one step below the fused table. The fused table is identically self-dual and its automorphisms permute its marks; each cousin is the fused table with one mark integrated by a Beta corner, so the cousins are tables of fusion width one, often not isomorphic to one another, whose periods are one function. At the fusion width one, the table itself is a cousin and the hidden symmetry acts on its exponent and its masses together; at the width $c\ge2$ the partial fusion $\cB^\sharp_{c-1}$ is the cousin, and the table lies $c-1$ cells below it.

\begin{remark}[Hidden symmetry]\label{rem:fusionsymmetry}
The index of the automorphisms of $\cB$ times the permutations of $G$ in those of $\cB^\sharp$ measures the symmetry the fusion adds. For $\binom{\cE}r$ it is $\binom{n+c}c$, all bijections of $\cE\sqcup G$ carrying $\binom{\cE\sqcup G}{r^\ast}$; for the list of rank two with one parallel pair of Remark~\ref{rem:fusionexamples}(iii) it is $(n-2)(n-3)$; for the $10$ fusable prime lists of rank three on seven marks the fusion is transitive and its group is $8$ times that of the list (receipts \texttt{census\_extra.py}, \texttt{atoms7.py}). At the fusion width two the index is $45$ for $\binom{[8]}3$, $20$ for one line of three marks and five free marks, $15$ for one line of four marks and four free marks, $6$ for two lines of three marks meeting in a mark and three free marks, $2$ for two disjoint lines of three marks and two free marks, and $1$ for $\mathrm F_7$ with a free mark, for $\mathrm A_9$ less a mark and for two disjoint lines of four marks; at the width three it is $220$ for $\binom{[9]}3$ and $1$ for $\mathrm A_9$ and for $\mathrm Z_9$ (receipt \texttt{w01\_seam\_rank3.py}). At $\binom{[2r+1]}r$, $r\ge2$, the symmetric group of $\binom{[2r+2]}{r+1}$ accounts for a nonzero submodule that the closings miss at its centre, with every symmetry of the four kinds of its centre ray in the group (Proposition~\ref{prop:onemarkhidden}); at $\binom{[5]}2$ its part in the span of a box is a line (Exercise~\ref{ex:closingssee}). Among $40$ identically self-dual lists of rank five on ten marks drawn at random, in each of which the sets of five marks that are not bases form a family closed under complement whose members meet pairwise in at most three marks, $23$ have ten cousins no two of which are isomorphic (receipt \texttt{random10.py}), and all the cousins of these $23$ are prime (receipt \texttt{random10b.py}).
\end{remark}

\section{Two-sums and repair}\label{sec:fusiontwosum}

% New: proved in the last part.
\begin{definition}[Covered marks; the conditions at a hinge]\label{def:hingeconditions}
Let $\cB$ be a table on $\cE$. A mark $e$ of a cocircuit $D$ is \emph{covered in} $D$ when some circuit $C$ has $e\in C\subseteq D$, and $e$ is \emph{uncovered} when some cocircuit contains it without covering it. For a cell $\mathcal A$ on $B\sqcup\{g\}$ with the hinge $g$ (Definition~\ref{def:substitution}) the \emph{conditions at the hinge} are:
\begin{enumerate}[label=(\alph*)]
\item[(a)] every cocircuit avoiding $g$ is a union of circuits;
\item[(u)] every cocircuit $D$ through $g$ covers every mark of $D\setminus\{g\}$;
\item[(s)] every cocircuit $D$ through $g$ covers every mark of $D\setminus\{g\}$ by circuits avoiding $g$;
\item[(h)] every cocircuit through $g$ covers $g$.
\end{enumerate}
For two cells $\mathcal A_1$, $\mathcal A_2$ the conditions are written $a_1,u_1,s_1,h_1$ and $a_2,u_2,s_2,h_2$.
\end{definition}

% New: proved in the last part.
\begin{theorem}[Fusion across a two-sum]\label{thm:fusiontwosum}
Let $\mathcal A_1$ and $\mathcal A_2$ be cells that are lists, on $\cE_1\sqcup\{g\}$ and $\cE_2\sqcup\{g\}$ with $n_1$ and $n_2$ marks, of the ranks $r_1,r_2$ and the fusion widths $c_1,c_2$, and $\mathcal S=\mathcal A_1\oplus_g\mathcal A_2$ their two-sum on $\cE_1\sqcup\cE_2$ (Lemma~\ref{lem:twosumpieces}).
\begin{enumerate}[label=(\roman*)]
\item $\mathcal S$ has the fusion width $c_1+c_2$. Its circuits are the circuits of $\mathcal A_1$ and of $\mathcal A_2$ that avoid $g$ and the sets $(C_1\setminus g)\cup(C_2\setminus g)$ with $C_i$ a circuit of $\mathcal A_i$ through $g$; its cocircuits are the same with cocircuits in place of circuits; and $\mathcal S^\ast=\mathcal A_1^\ast\oplus_g\mathcal A_2^\ast$.
\item $\mathcal S$ is fusable exactly when
\[
a_1\ \text{and}\ u_1\ \text{and}\ a_2\ \text{and}\ u_2\ \text{and}\ (s_1\ \text{or}\ h_2)\ \text{and}\ (s_2\ \text{or}\ h_1),
\]
and a cell, read as a table, is fusable exactly when $a$, $u$ and $h$ hold. So two fusable cells give a fusable two-sum.
\item If $\mathcal S$ is fusable, then $c_1,c_2\ge0$, and in $\mathcal S^\sharp$ the sets $\cE_1$ and $\cE_2$ are bunches of the widths $c_1+1$ and $c_2+1$.
\item If both cells are fusable, the two-sum $\mathcal A_1^\sharp\oplus_g\mathcal A_2^\sharp$, with the fusion bunch $G_1\sqcup G_2$, is an identically self-dual major of $\mathcal S$. It is $\mathcal S^\sharp$ when $c_1=0$ or $c_2=0$. When $c_1,c_2\ge1$ its children at one mark of $G_1$ and at one mark of $G_2$ have the ranks $r^\ast_1-1$ and $r^\ast_1$ on $\cE_1$, $r^\ast_1=n_1-r_1$ the corank of $\mathcal A_1$, so $G_1\sqcup G_2$ is not in level form in it and it is a major other than $\mathcal S^\sharp$.
\item If $c_1+c_2=1$, $\mathcal S$ is fusable exactly when both cells are fusable; then, for $c_2=0$, $\mathcal A_2$ is identically self-dual, $\mathcal S^\sharp=\mathcal A_1^\sharp\oplus_g\mathcal A_2$, and the hinge of $\mathcal S$ is a hinge of $\mathcal S^\sharp$.
\end{enumerate}
\end{theorem}

The condition is local at the hinge: away from it each cell answers for itself, and through it a cocircuit of the two-sum joins a cocircuit of each cell, so a mark is covered either by a circuit of its own cell avoiding the hinge or by a circuit through the hinge continued by a circuit through the hinge on the other side; that continuation exists exactly when the other cell covers its hinge. In the fused table the hinge of a cell of positive width thickens into a bunch of that width plus one. Fusing cell by cell gives a major, identically self-dual, that is not the fusion as soon as two cells have positive width: the canonical fusion of a two-sum is the global level cell. At the total width one the hinge survives the fusion, with the identically self-dual cell on one side and the fusion of the other cell on the other.

% New: proved in the last part.
\begin{corollary}[A two-sum with an identically self-dual cell]\label{cor:isdcell}
Let $\mathcal A_1$ be identically self-dual and $s,t>0$, with the masses $\mvec$ of $\mathcal S=\mathcal A_1\oplus_g\mathcal A_2$ on its balance wall with every proper wall positive, and let $\kappa=m_{\cE_1}-s(r_1-1)-t(|\cE_1|-r_1)$ be the hinge wall of $\mathcal A_1$ (Definition~\ref{def:hinged}). Then
\[
P_{\mathcal S}(\mvec;s,t)=P_{\mathcal A_1}\big(\mvec|_{\cE_1},2\sigma-\kappa;s,t\big)\,P_{\mathcal A_2}\big(\mvec|_{\cE_2},\kappa;s,t\big)\Big/\Gst21{2\sigma-\kappa,\ \kappa}{2\sigma},
\]
the hinge carrying $2\sigma-\kappa$ in $\mathcal A_1$ and $\kappa$ in $\mathcal A_2$. When the cells are lists and $\mathcal A_2$ is fusable of the fusion width one, this is the period of $\mathcal S$ read across the hinge that $\mathcal S^\sharp=\mathcal A_1\oplus_g\mathcal A_2^\sharp$ keeps.
\end{corollary}

By Corollary~\ref{cor:mergedvalues} at the hinge, the first factor divided by the Beta corner is the period $P_{\mathcal A_1\setminus g}(\mvec|_{\cE_1};\kappa,2\sigma-\kappa)$ of the deletion at the hinge. So identically self-dual cells attached at their hinges to one cell that is not identically self-dual cost the period of that cell and, for each of them, the period of its deletion at the hinge; two cells that are not identically self-dual on the two sides of a hinge pair two laws, and the hinge is a convolution (Theorem~\ref{thm:glue}).

% New: proved in the last part.
\begin{proposition}[Uncovered marks and repair]\label{prop:repair}
Let $(\cB,\theta)$ be a table on $\cE$ of the fusion width $c$.
\begin{enumerate}[label=(\roman*)]
\item $\cB$ has no uncovered mark (Definition~\ref{def:hingeconditions}) exactly when every cocircuit of $\cB$ is a union of circuits. So a list, or a table of the fusion width $c\neq0$, is fusable exactly when it has no uncovered mark, and a table of the fusion width zero exactly when it is identically self-dual.
\item For a set $S$ of marks, none a loop, let $(\cB_S,\theta)$ be the table in which each $h\in S$ is replaced by a block $\{h,h'\}$ of two marks by the first rule of Theorem~\ref{thm:block}, a parallel pair, every new basis keeping the weight of its basis, and $S'=\{h':\ h\in S\}$. Then $\cB_S$ has the fusion width $c+|S|$, $(\cB_S,\theta)\setminus S'=(\cB,\theta)$, and $\cB_S$ has no uncovered mark exactly when $S$ contains every uncovered mark of $\cB$; when $\cB$ is a list or $c+|S|\neq0$, $\cB_S$ is fusable exactly then.
\item If $\cB_S$ is fusable, $(\cB,\theta)=(\cB_S)^\sharp/G\setminus S'$ is a minor of an identically self-dual table. If $c\ge0$ and $(\cB,\theta)$ is not fusable, $\cB_S$ is fusable for $S$ the set of the uncovered marks when there is one, and for $S$ any one mark when there is none, no mark being then a loop. Doubling every mark of a list that is not a loop gives a fusable list.
\item At the full basis lean, where every other proper wall is positive, the residue of $P_{\cB_S}(\mvec;s,0)$ at the wall $m_{h'}=0$ of one copy $h'$ is $P_{\cB_{S-h}}(\mvec_{-h'};s,0)$ (Theorem~\ref{thm:wallres}); so $P_{\cB}(\cdot;s,0)$ is the residue of $P_{\cB_S}(\cdot;s,0)$ at the walls of the copies, taken one at a time.
\end{enumerate}
\end{proposition}

A mark is uncovered when a cocircuit through it has no circuit through it inside: the cocircuit is a set the bases must meet, and the mark is not tied by any circuit to the other marks there. A parallel copy is a circuit of two marks through it inside every cocircuit through it, so doubling an uncovered mark covers it, and each doubled mark adds one to the fusion width. The list $\binom{[3]}2$, of the fusion width $-1$, has every mark uncovered: one or two doubled marks do not fuse, all three do, at the width two (receipts \texttt{t04\_twosums.py}, \texttt{w04\_repair.py}); its dual $\binom{[3]}1$, which the procedure below reads first, fuses as it is, into $\binom{[4]}2$.

\section{The procedure}\label{sec:fusionproc}

The statements of this chapter set a branch at every node of the Foreman tree (Definition~\ref{def:tree}), read before anything is integrated.
\begin{enumerate}[label=(\arabic*)]
\item \emph{Peel.} If the realization of the node is not an atom, peel it at its Foreman bunch and continue on the room (Definition~\ref{def:tree}).
\item \emph{Go down.} If the table of the node is identically self-dual, every mark comes off by one Beta corner at every pair of exponents (Corollary~\ref{cor:mergedvalues}); its cousins share its period (Proposition~\ref{prop:cousins}).
\item \emph{Go up.} Otherwise read the fusion width $c=r^\ast-r$, passing to the dual table when it is negative, and the tell on the Skeleton (Theorem~\ref{thm:fusiontell}). If the table is fusable, fuse it (Theorem~\ref{thm:fusion}); at $c=0$ it is not, a fusable table of the width zero being identically self-dual. At $c=1$ the fusion is free, one Beta corner (Corollary~\ref{cor:fusionone}): the automorphisms of the identically self-dual table $\cB^\sharp$ act on the exponent and the masses of the table, its cousins share its period, and its module carries over (Chapter~\ref{ch:twoends}). At $c\ge2$ the table is the child of the full pattern of the fusion bunch of $\cB^\sharp$; the fused table is the universal piece of width $c$ glued to the lift partner (Theorem~\ref{thm:fusionglue}), both laws spread (Theorem~\ref{thm:fusionprice}), the step from $\cB^\sharp$ to $\cB^\sharp_{c-1}$ is a Beta corner and the $c-1$ steps below it are cells whose laws are not points.
\item \emph{Two-sums.} On a two-sum of lists the tell is read at the hinge (Theorem~\ref{thm:fusiontwosum}(ii)); at the total width one the hinge survives the fusion and the period is a product over one Beta corner (Corollary~\ref{cor:isdcell}); when both cells are fusable and one has the width zero the fusion is the two-sum of the fusions of the cells, and when both cells have positive width it is the global level cell, not that two-sum (Theorem~\ref{thm:fusiontwosum}(iv)).
\item \emph{Repair.} If the table is not fusable and has uncovered marks, they are the obstruction: doubling exactly those in parallel gives a fusable table with the table as its deletion at the copies (Proposition~\ref{prop:repair}), at the fusion width $c$ plus the number of uncovered marks; a table of the width zero with no uncovered mark that is not identically self-dual, which needs weights, becomes fusable when any one mark is doubled.
\end{enumerate}
The fusion width is the price, and it is read off the rank and the number of marks before any integral is taken: at the width one the fusion bunch is a single mark with a point for its law; at the width $c\ge2$ it is the universal piece, whose law spreads over a torus of the dimension $c-1$ against a lift partner whose law is not a point.


\chapter{Closing a glue along a direction}\label{ch:closeglue}

A glue opened on its shares is a fibre integral of Gammas and two seam functions (Theorem~\ref{thm:gluewidth}(iii)). It is closed by moving the contour along a direction of the shares, collecting a fibre integral of residues at every polar hyperplane crossed and showing that what is left tends to zero. The glue closes on the lattices of its Gammas along every direction whose rate, read off the two laws through the largest values of the monomials of their letter vectors, is negative (Theorem~\ref{thm:closecone}), and past the walls of the piece, or of the piece and the partner both, at the growths of the continued seam functions (Theorem~\ref{thm:pastwalls}). Moving the contour is exact at every step, and the limit needs only a bound that travels with it (Lemma~\ref{lem:pathclose}): the series depends on the direction only through its chamber (Theorem~\ref{thm:chamber}), the laws can be closed one source point at a time and averaged, at the rate read on the support of the source measure rather than on its convex hull (Theorem~\ref{thm:averaging}), and past the walls of a piece the bound is read on the outline of a keyhole continuation, the modulus and the phase of the sector letters taken at one point (Theorem~\ref{thm:keyholebound}). At simple nonresonant crossings of a rational ray, a crossing term kept with its falling shares integrated on their strips is a finite sum of joint Trailmoments of fixed functions, a falling Gamma becoming a positive share and one Beta variable and a continued source factor a segment and circle integrals (Theorem~\ref{thm:retainedchain}); for a positive kernel its root rate is the exponential of the least Gamma rate of its chamber in every fibre dimension and, averaged, of the essential supremum over the source of that least rate, whose sign, when it is not zero, decides whether the terms tend to zero (Theorem~\ref{thm:retainedchamberrate}, Corollary~\ref{cor:retainedaveragedrate}); and an averaged closing can converge where no direction of its chamber certifies it, at a family of tables whose seam has single stations (Theorem~\ref{thm:canonicaltablegap}). Coincident crossings are packets of factors polar on one crossed hyperplane, assembled before a sign or a top is read (Theorem~\ref{thm:crossingjet}), and at the root rate one a positive moment sequence is decided by its law near its top (Theorem~\ref{thm:unitrate}).

\section{Closing along a direction}\label{sec:closecone}

% Reference 20: conventions of the closing; source line 7588
\begin{definition}[Letter vectors]\label{def:readingvectors}
Let a piece on $B$ and a partner on $R$ with the stations $\mathcal D$ be glued with masses on the balance wall of the glue (Definition~\ref{def:seampiece}), the glue read on its shares as in Theorem~\ref{thm:gluewidth}(iii) or Theorem~\ref{thm:transversepins}. The \emph{letter vectors} of the piece and of the partner are $\mathbf V_B=(V_{jD},V^\ast_{jD})_{j,D}$ and $\mathbf V_R=(L_{jD},L^\ast_{jD})_{j,D}$, and for a vector $\mathbf v$ of the coordinates of the shares $\mathbf V_B^{\mathbf v}=\prod_{j,D}V_{jD}^{v_{jD}}V^\ast_{jD}{}^{v^\ast_{jD}}$, and likewise $\mathbf V_R^{\mathbf v}$.
\end{definition}

% Reference 20: thm:closecone; source line 7588 (the largest values read on the law of the piece).
% Proved in the last part: thm:closecone.
\begin{theorem}[Closing the glue along a direction]\label{thm:closecone}
In Theorem~\ref{thm:transversepins} let $V$ be the space of the free directions of the shares, so that for $\mathbf v\in V$ the monomials $\mathbf V_B^{\mathbf v}$ and $\mathbf V_R^{\mathbf v}$ (Definition~\ref{def:readingvectors}) are homogeneous of degree zero, and their suprema over the orthants do not change under $\vec{\mathcal N}_B^\perp$ and $\vec{\mathcal N}_R^\perp$. When the seam word of the piece has a Newton sum of full dimension, $\sup\mathbf V_B^{\mathbf v}=\sup_{u\in\Delta_B}\boldsymbol\varrho^\rho_B(u)^{\mathbf v}$, the largest value of the monomial $\mathbf x^{\mathbf v}$ on the support of the law of the piece, and when the room word of the partner has a Newton sum of full dimension, $\sup\mathbf V_R^{\mathbf v}=\sup_{v\in\Delta_R}\boldsymbol\varrho^\rho_R(v)^{\mathbf v}$, the largest value on the support of the law of the partner; where the room letters of the partner are monomials, $\sup\mathbf V_R^{\mathbf v}=\boldsymbol\xi^{\mathbf v}$, $\boldsymbol\xi$ their coefficients. For $\boldsymbol\sigma\in V$ put
\[
h(\boldsymbol\sigma)=\sum_{j,D}\big(\sigma_{jD}\log|\sigma_{jD}|+\sigma^\ast_{jD}\log|\sigma^\ast_{jD}|\big)+\log\sup\mathbf V_B^{-\boldsymbol\sigma}+\log\sup\mathbf V_R^{-\boldsymbol\sigma},
\]
with $0\log0=0$, the first sum over the shares of the active letters, the \emph{rate} of the direction; it is homogeneous of degree one. Let $h(\boldsymbol\sigma)<0$, and let the real point $\mathbf c$ of the pins lie off the null set of the points whose ray $\mathbf c+\tau\boldsymbol\sigma$, $\tau\ge0$, meets two polar hyperplanes of the Gammas that are not parallel on the fibre at one point, or, for arbitrarily large $\tau$, at points less than $\tau^{-3}$ apart. Then the ray meets only the hyperplanes $a_{jD}=-n$ with $\sigma_{jD}<0$ and $b_{jD}=-n$ with $\sigma^\ast_{jD}<0$, $n\in\mathbb N$, and
\[
P_{B\sqcup R}(\mvec)=\sum_{\text{crossings}}\mathcal R,
\]
the sum, in the order of crossing, of the fibre integrals $\mathcal R$ of the residues of the integrand at the hyperplanes crossed, each pinned also by its hyperplane, of the dimension $\dim V-1$ and of the same kind, at the point where the ray meets it; a hyperplane of the fibre on which several Gammas have poles is crossed once, for every $\mathbf c$, and its residue is a Laurent coefficient (Theorem~\ref{thm:closing}). The sum converges absolutely: the residue at a crossing at $\tau$, and the integral left after it, are $O\big(e^{\tau(h(\boldsymbol\sigma)+\epsilon)}\big)$ for every $\epsilon>0$.
\end{theorem}

So the glue closes on the lattices of its Gammas along every direction of the cone $\{h<0\}$ of $V$, the terms falling at least by the factor $e^{h(\boldsymbol\sigma)}$ per unit of $\boldsymbol\sigma$. The cone and the rate are read off the laws through the largest values of the monomials of their letter vectors, and the term $\sum\sigma\log|\sigma|$ of the rate is the cost of moving shares between stations: it vanishes at width one with the stations of one excess merged, every letter then moving its two shares oppositely, and along the one direction of a node of width one the condition is that the ranges of the two ratios be separated, as in Section~\ref{sec:onescale}. A direction along which a law is unbounded, its supremum infinite, crosses the walls of its seam function as well. The seam function of a piece continues past its walls with a growth that its sector integrals bound, and the partner can be read in its monomials, which have no walls, or continued past its own walls in the same way.

\section{Closing along a path}\label{sec:pathclose}

The closings of Theorems~\ref{thm:closecone} and~\ref{thm:pastwalls} move the contour of the share fibre along a ray, collect a fibre integral of residues at every polar hyperplane crossed, and show that what is left tends to zero. The two steps come apart. Moving the contour is exact at every step: the period is the sum of the crossing terms met so far and of the fibre integral at the current point, by Cauchy's theorem on bounded contours. Only the last step is a limit, and it asks for a bound on the remaining fibre integral at the points that the path passes through, not for a domination of the double family of crossings and fibres, which is what interchanging the sum over the crossings with the fibre integrals would ask, and what fails past the walls. The step is Lemma~\ref{lem:pathclose}, with three consequences: a closing depends on its direction only through its chamber (Theorem~\ref{thm:chamber}); the laws can be closed one source point at a time and averaged (Theorem~\ref{thm:averaging}); and past the walls of a piece the bound that travels is read on the outline of a keyhole continuation (Theorem~\ref{thm:keyholebound}). The bounds multiply over products, survive averaging and pass to the crossing terms (Proposition~\ref{prop:pathcompose}).

% Reference 20: lem:pathclose; source line 7644
\begin{lemma}[Closing along a path]\label{lem:pathclose}
Let $A$ be the affine space of the real points of the real pins of a fibre integral and $V$ its space of directions, of dimension $n\ge1$, with the measure $d\boldsymbol\vartheta$, the fibre integral through a real point $p$ being $\int_{p+iV}I$ (Definition~\ref{def:fibre}) wherever it converges absolutely; and let $\Omega\subseteq A$ be open. Let $I$ be meromorphic on the tube $\Omega+iV$, with its poles on the complex hyperplanes $\{\ell_H=0\}$ of a family $\mathfrak H$ of real hyperplanes $H=\{\ell_H=0\}$ of $A$, locally finite in $\Omega$, each $\ell_H$ real, affine and not constant. Let $\gamma:[0,\infty)\to\Omega$ be continuous and piecewise linear, with finitely many breaks on every bounded interval and $\gamma(0)=\mathbf c$ off $\bigcup\mathfrak H$, and assume:
\begin{enumerate}[label=(H\arabic*),start=0]
\item \emph{Clean crossings.} $\gamma$ meets $\bigcup\mathfrak H$ at a discrete set of times, its \emph{crossings}, none of them a break of $\gamma$; at a crossing $\tau^\ast$ the hyperplanes of $\mathfrak H$ through $\gamma(\tau^\ast)$ are one hyperplane $H$, and $\ell_H\circ\gamma$ changes sign at $\tau^\ast$ with a nonzero slope.
\item \emph{Local decay.} For every compact $Q\subset\Omega$ and $\delta>0$ there are $C,c>0$ with $|I(x+i\boldsymbol\vartheta)|\le Ce^{-c|\boldsymbol\vartheta|}$ whenever $x\in Q$ and $x+i\boldsymbol\vartheta$ lies at distance at least $\delta$ from every hyperplane $\{\ell_H=0\}$ of the tube.
\item \emph{A travelling bound.} There are times $\tau_1<\tau_2<\cdots$ tending to infinity, none a crossing, at which
\[
\beta(\tau_k)=\int_V\big|I(\gamma(\tau_k)+i\boldsymbol\vartheta)\big|\,d\boldsymbol\vartheta
\]
is finite and tends to zero.
\end{enumerate}
Then for every $k$
\[
\int_{\mathbf c+iV}I=\sum_{\tau^\ast<\tau_k}\mathcal R_{\tau^\ast}+\int_{\gamma(\tau_k)+iV}I,\qquad\Big|\int_{\gamma(\tau_k)+iV}I\Big|\le(2\pi)^{-n}\beta(\tau_k),
\]
so that $\int_{\mathbf c+iV}I=\lim_k\sum_{\tau^\ast<\tau_k}\mathcal R_{\tau^\ast}$. The \emph{crossing term} at a crossing $\tau^\ast$ of the hyperplane $H$ is
\[
\mathcal R_{\tau^\ast}=\pm(2\pi)^{1-n}\int_{V_H}\operatorname{Res}_{w=0}I\big(\gamma(\tau^\ast)+w\,\mathbf e+i\boldsymbol\vartheta'\big)\,d\mu_H(\boldsymbol\vartheta'),
\]
with $V_H\subset V$ the kernel of the linear part $\ell^{\mathrm{lin}}_H$ of $\ell_H$, $\mu_H$ the measure on $V_H$ with $d\boldsymbol\vartheta=d\mu_H\,d\ell^{\mathrm{lin}}_H(\boldsymbol\vartheta)$, $\mathbf e\in V$ any vector with $\ell^{\mathrm{lin}}_H(\mathbf e)=1$, on which the term does not depend, and the sign $+$ when $\ell_H$ decreases along $\gamma$ and $-$ when it increases. The residue is that of the whole integrand, a Laurent coefficient where several of its factors are polar on $H$, and the crossing term is its fibre integral, pinned also by $\ell_H=0$. If some $\tau_k$ lies between any two consecutive crossings and $\sum_k\beta(\tau_k)<\infty$, the crossing terms are absolutely summable, a crossing between $\tau_k$ and $\tau_{k+1}$, the only one there, having $|\mathcal R_{\tau^\ast}|\le(2\pi)^{-n}\big(\beta(\tau_k)+\beta(\tau_{k+1})\big)$.
\end{lemma}

% Reference 20: thm:pastwalls; source line 7609 (a piece of tables, the partner in its monomials).
% Proved in the last part: thm:pastwalls.
\begin{theorem}[Closing past the walls]\label{thm:pastwalls}
Let the seam function of the piece converge at a real point of the pins of the seam with positive shares, so that its seam word has a Newton sum of full dimension (Theorem~\ref{thm:poswordconv}). Read the partner in its monomials, $L_{jD}=\sum_k\xi_{jDk}\,\alpha_R^{\mathbf n_{jDk}}$ and $L^\ast_{jD}=\sum_k\xi^\ast_{jDk}\,\alpha_R^{\mathbf n^\ast_{jDk}}$, with one share $a_{jDk}$ or $b_{jDk}$ for each monomial, and write $\bar a_{jD}=\sum_ka_{jDk}$ and $\bar b_{jD}=\sum_kb_{jDk}$; bars mark sums over the monomials.
\begin{enumerate}[label=(\roman*)]
\item On every contour of shares with positive real parts, pinned by the letter totals and by $\sum_{j,D,k}(a_{jDk}\mathbf n_{jDk}+b_{jDk}\mathbf n^\ast_{jDk})=\mvec_R$, at which the seam function of the piece converges, the period of the glue converging,
\[
P_{B\sqcup R}(\mvec)=\int_{\text{pins}}\prod_{j:\,s_j>0}\frac{\prod_{D,k}\Gamma(a_{jDk})\,\xi_{jDk}^{-a_{jDk}}}{\Gamma(s_j)}\prod_{j:\,t_j>0}\frac{\prod_{D,k}\Gamma(b_{jDk})\,\xi^\ast_{jDk}{}^{-b_{jDk}}}{\Gamma(t_j)}\ P_{W_B(\bar{\mathbf a},\bar{\mathbf b})}(\mvec_B).
\]
\item Fix $e_0\in B$ and a \emph{fan of the seam word}: finitely many rational simplicial cones of dimension $|B|-1$ with disjoint interiors covering $\mathbb R^{B\setminus e_0}$, each contained, for every seam letter $V$ of an active letter, in a set $\{\mathbf v:\ \ip{\mathbf t}{\mathbf v}\le\ip{\mathbf t'}{\mathbf v}\ \text{for every exponent}\ \mathbf t'\ \text{of the support of}\ V\}$, $\mathbf t$ one exponent of that support, the pairings taken on $B\setminus e_0$. On a cone $\pi$ of it, a \emph{sector}, with the primitive integral generators $\mathbf v_{\pi,1},\dots,\mathbf v_{\pi,|B|-1}$, write the logarithmic coordinates $x=\log\alpha_{B\setminus e_0}$ of the section $\alpha_{e_0}=1$ as $x=-\sum_ku_k\mathbf v_{\pi,k}$, $u_k\ge0$; the $q_k=e^{-u_k}\in(0,1]$ are the \emph{sector coordinates}. On the section a seam letter $V$ of an active letter is $V=\prod_kq_k^{\ip{\mathbf t}{\mathbf v_{\pi,k}}}\,k_{V,\pi}(q)$, $\mathbf t$ the exponent of its support at which every generator of $\pi$ pairs least, and its \emph{sector letter} $k_{V,\pi}$ is a polynomial in $q$ with nonnegative coefficients and a positive constant term. Write $k_{jD,\pi}$ and $\psi_{jD,\pi}$ for the sector letters of $V_{jD}$ and $V^\ast_{jD}$, $\mathbf t_{jD,\pi}$ and $\mathbf t^\ast_{jD,\pi}$ for their exponents, and for a generator $\mathbf v$ of $\pi$ call
\[
\kappa_{\mathbf v}(\mathbf z)=\Big\langle\mvec_B-\sum_{j,D}\big(\bar a_{jD}\,\mathbf t_{jD,\pi}+\bar b_{jD}\,\mathbf t^\ast_{jD,\pi}\big),\ \mathbf v\Big\rangle,\qquad\mathbf z=(\bar{\mathbf a},\bar{\mathbf b}),
\]
the pairing on $B\setminus e_0$, the \emph{wall} of $\mathbf v$: affine in the shares with a rational linear part, and the same for every sector of which $\mathbf v$ is a generator. On the pins of the seam, at every real point at which every wall is positive,
\begin{multline*}
P_{W_B(\mathbf z)}(\mvec_B)\\
=\sum_\pi|\det(\mathbf v_{\pi,1},\dots,\mathbf v_{\pi,|B|-1})|\int_{[0,1]^{|B|-1}}\prod_kq_k^{\kappa_{\mathbf v_{\pi,k}}(\mathbf z)-1}\prod_{j,D}k_{jD,\pi}(q)^{-\bar a_{jD}}\,\psi_{jD,\pi}(q)^{-\bar b_{jD}}\,dq,
\end{multline*}
and the seam function converges at a real point exactly when every wall is positive there. It continues meromorphically to all complex shares, with poles, of order at most $|B|-1$, only on the hyperplanes on which a wall lies in $-\mathbb N$, and for every $\theta>0$ there are $M$, $C_0$ and $n_1$ with
\[
\big|P_{W_B(\mathbf z)}(\mvec_B)\big|\le C_0\,\delta^{-(|B|-1)}(1+|\mathbf z|)^{n_1}M^{|\re\mathbf z|_1}e^{\theta|\im\mathbf z|_1},
\]
$\delta$ the smaller of $1$ and the distance of $\mathbf z$ from these hyperplanes. The piece is \emph{of tables} when at every station the seam letters of the active letters are the basis and the cobasis polynomials of one table on $B$ with the bases $\cB_{jD}$ and positive weights, as for a bunch of a family of tables (Proposition~\ref{prop:seamidentity}(i)). Its seam word then has a Newton sum of full dimension exactly when the family of these tables is connected; the sectors of the orders $\pi$ of $B$ (Theorem~\ref{thm:exist}(ii)) form a fan of it, every determinant being $1$, the sector letters being those of the tables $\cB_{jD}$ and the generators $\one_{U_k}-(\one_{U_k})_{e_0}\one$ on $B\setminus e_0$, $U_k$ the tails; the wall of the generator of $U$ is the proper wall $\kappa_U=m_U-\sum_{j,D}\big(\bar a_{jD}\,d_{jD}(U)+\bar b_{jD}\,d^\ast_{jD}(U)\big)$, $d_{jD}$ and $d^\ast_{jD}$ the counts of $\cB_{jD}$.
\item For a direction $\boldsymbol\sigma$ of the shares of (i) and $\theta\ge0$ let $g(\boldsymbol\sigma;\theta)$, at the real point $\bar{\mathbf c}$ of the pins of the seam, be the infimum of the $M$ for which there are $C$, $n_1$ and $n_2$ with
\[
\big|P_{W_B(\bar{\mathbf c}+\tau\bar{\boldsymbol\sigma}+i\bar{\boldsymbol\vartheta})}(\mvec_B)\big|\le C\,\delta^{-n_2}(1+\tau+|\boldsymbol\vartheta|)^{n_1}M^\tau e^{\theta|\boldsymbol\vartheta|_1}
\]
for $\tau\ge0$ and $\boldsymbol\vartheta$ in the directions of the shares, $\delta$ the smaller of $1$ and the distance from the polar hyperplanes: the \emph{growth} of the continued seam function along $\boldsymbol\sigma$ at the angle $\theta$, the growth of Definition~\ref{def:growth} read along any direction with the angle fixed. It is finite for $\theta>0$, and $g(\boldsymbol\sigma;0)=\sup\mathbf V_B^{-\bar{\boldsymbol\sigma}}$ when no wall of the seam word falls along $\bar{\boldsymbol\sigma}$, which is exactly when that supremum is finite. Put
\[
h_\Gamma(\boldsymbol\sigma)=\sum_l\sigma_l\log|\sigma_l|-\sum_l\sigma_l\log\xi_l,\qquad h_\theta(\boldsymbol\sigma)=h_\Gamma(\boldsymbol\sigma)+\log g(\boldsymbol\sigma;\theta)+\frac{2\theta^2}\pi\sum_{l:\,\sigma_l>0}\sigma_l,
\]
$l$ running over the shares and $\xi_l$ being the coefficient of the monomial of the share $l$; $h_\Gamma$ is the \emph{Gamma rate} of the direction. Let the real point $\mathbf c$ of the pins lie off the null set of Theorem~\ref{thm:closecone}, taken with the polar hyperplanes of the seam function as well, and let $g$ and $h_\theta$ be taken at its sums $\bar{\mathbf c}$. If $h_\theta(\boldsymbol\sigma)<0$ for some $\theta\in[0,\pi/8)$, then
\[
P_{B\sqcup R}(\mvec)=\sum_{\text{crossings}}\mathcal R,
\]
the sum, in the order of crossing, of the fibre integrals of the residues of the integrand at the polar hyperplanes met by the ray $\mathbf c+\tau\boldsymbol\sigma$: the hyperplanes $a_l=-n$ of the shares with $\sigma_l<0$, $n\in\mathbb N$, and the hyperplanes on which a wall of (ii) that falls along $\bar{\boldsymbol\sigma}$ lies in $-\mathbb N$, a hyperplane polar for several factors being crossed once, with its full residue. The sum converges absolutely, the terms at $\tau$ and the integrals left after them being $O\big(e^{\tau(h_\theta(\boldsymbol\sigma)+\epsilon)}\big)$ for every $\epsilon>0$.
\item \emph{Both sides continued.} Let the glue be read on the share side of Theorem~\ref{thm:gluewidth}(iii), one share $a_{jD}$ or $b_{jD}$ per station, in place of the reading of (i), and let the seam functions of the piece and of the partner both converge at a real point $\mathbf c$ of the pins with positive shares. By (ii), applied to the piece and to the partner, a piece on $R$ read with the reversed excesses (Definition~\ref{def:seampiece}), both continue meromorphically to all complex shares of the pins, with their poles on the hyperplanes on which a wall of the seam word of the piece or of the room word of the partner lies in $-\mathbb N$, each with the bound of (ii). For a direction $\boldsymbol\sigma$ of the shares and $\theta\ge0$ let $g_B(\boldsymbol\sigma;\theta)$ and $g_R(\boldsymbol\sigma;\theta)$ be the growths of the two continued seam functions along $\boldsymbol\sigma$ at $\mathbf c$, defined as in (iii) with the shares in place of their sums, and put
\[
h_{\theta_B,\theta_R}(\boldsymbol\sigma)=\sum_l\sigma_l\log|\sigma_l|+\log g_B(\boldsymbol\sigma;\theta_B)+\log g_R(\boldsymbol\sigma;\theta_R)+\frac{2(\theta_B+\theta_R)^2}\pi\sum_{l:\,\sigma_l>0}\sigma_l .
\]
Let $\mathbf c$ lie off the null set of Theorem~\ref{thm:closecone}, taken with the polar hyperplanes of both seam functions as well. If $h_{\theta_B,\theta_R}(\boldsymbol\sigma)<0$ for some $\theta_B,\theta_R\ge0$ with $\theta_B+\theta_R<\pi/8$, then $P_{B\sqcup R}(\mvec)=\sum_{\text{crossings}}\mathcal R$, the sum, in the order of crossing, of the fibre integrals of the residues of the integrand at the polar hyperplanes met by the ray $\mathbf c+\tau\boldsymbol\sigma$: the hyperplanes $a_l=-n$ of the shares with $\sigma_l<0$, $n\in\mathbb N$, and those on which a wall of either word that falls along $\boldsymbol\sigma$ lies in $-\mathbb N$, a hyperplane polar for several factors being crossed once, with its full residue. The sum converges absolutely, the terms at $\tau$ and the integrals left after them being $O\big(e^{\tau(h_{\theta_B,\theta_R}(\boldsymbol\sigma)+\epsilon)}\big)$ for every $\epsilon>0$. A side along whose direction no wall of its word falls has the growth $\sup\mathbf V^{-\boldsymbol\sigma}$ of its letter vector at every angle, the largest value of the monomial on the support of its law (Theorem~\ref{thm:closecone}).
\end{enumerate}
\end{theorem}

Along a direction in which the seam function converges, (iii) at $\theta=0$ is the bound of Theorem~\ref{thm:closecone} for the glue read in the monomials of the partner, and past a wall the residues of the seam function join those of the Gammas at the cost $\log g$ of its continuation, as the walls of the two families join the lattice at width one (Theorem~\ref{thm:nodeclosing}(iii)). Read in its monomials as well, any piece contributes only its coefficients: the integrand is then Gammas and coefficients, every supremum of a monomial of the letter vectors is a monomial in the coefficients, and every direction of the shares with $h_\Gamma(\boldsymbol\sigma)<0$ closes, by Theorem~\ref{thm:closecone}. In which representation to close is a choice of rate: a law bounded along the direction closes at the largest value of its monomial, a side past its walls at the growth of its continuation, a side read in monomials at its coefficients, at the price of one pin for each of its marks. So (iii) and (iv) read each side at its own rate, the coefficients $-\sum_l\sigma_l\log\xi_l$ for a side read in its monomials and the logarithm of the growth of its continued seam function for a side read in its seam letters; at width one, where the share side has one free share, (iv) continues both families of the node, as Theorem~\ref{thm:nodeclosing}(iii) does.

Expanding each active letter $L$, with its exponent $e_L$, in some of its terms, $L=Y_L+X_L$ with $X_L=\sum_kX_{L,k}$ the chosen terms, gives the multiple lattice whose terms are $(-1)^{|\mathbf n_L|}(e_L)_{|\mathbf n_L|}/\mathbf n_L!$ times the periods of $\alpha^{\mvec-\one}\prod_LY_L^{-e_L-|\mathbf n_L|}\prod_kX_{L,k}^{n_{L,k}}$, positive words up to sign. Summed term by term in absolute value, it converges absolutely exactly when $X_L\le Y_L$ on the orthant for every active letter and the integral of $\alpha^{\mvec-\one}\prod_L(Y_L-X_L)^{-e_L}$ over a section is finite (Proposition~\ref{prop:expansion}). This needs the Newton polytope of $X_L$ inside that of $Y_L$, so it never holds for a letter of a table, all of whose monomials are vertices of its Newton polytope: for families of tables the lattices close only in the order of a direction.

Theorems~\ref{thm:closecone} and~\ref{thm:pastwalls} are the lemma along a ray: (H0) holds off a null set of starting points, (H1) from the exponential decay of the Gamma factors in the imaginary directions (Theorem~\ref{thm:decay} at one scale, carried to negative real parts by the shift of Proposition~\ref{prop:stacked}(iii)), and (H2) from the largest values of the monomials on the supports of the laws or from the growth of the continued seam function. Two closings of one fibre integral along two paths agree at all times up to their remainders,
\[
\sum_{\tau^\ast<\tau_k}\mathcal R_{\tau^\ast}-\sum_{\tau'^\ast<\tau'_k}\mathcal R'_{\tau'^\ast}=\int_{\gamma'(\tau'_k)+iV}I-\int_{\gamma(\tau_k)+iV}I,
\]
and one of them may converge while the other does not. Each crossing term is a fibre integral of one dimension less, of an integrand of the same kind, which the lemma can close again along a path of its own when a bound travels with it there; this is how multiple lattices arise.

% Reference 20: def:chamber; source line 7685 (the closings of thm:closecone and thm:pastwalls(iii)).
% Proved in the last part: def:chamber.
\begin{definition}[Chambers of a closing]\label{def:chamber}
Let $I$ be meromorphic on the tube $\Omega+iV$ of Lemma~\ref{lem:pathclose} over an open convex $\Omega\subseteq A$, with its poles on the complex hyperplanes on which one of finitely many real affine functions $\ell$ on $A$, not constant, lies in $-\mathbb N$: its \emph{polar functions}. A direction $\boldsymbol\sigma\in V$ \emph{falls} a polar function $\ell$ when $\partial_{\boldsymbol\sigma}\ell<0$, at the \emph{falling rate} $-\partial_{\boldsymbol\sigma}\ell$. For a real point $\mathbf c\in\Omega$ at which every polar function is positive, the \emph{chamber} $C(\boldsymbol\sigma)$ of a direction $\boldsymbol\sigma$ whose ray $\mathbf c+\tau\boldsymbol\sigma$, $\tau\ge0$, lies in $\Omega$ is the set of the $\boldsymbol\sigma'\in V$ whose ray lies in $\Omega$ and which fall exactly the polar functions that $\boldsymbol\sigma$ falls, each at the same falling rate; it is convex. In Theorem~\ref{thm:closecone} the polar functions are the shares of the active letters, and $\Omega$ is the interior of the set of the real points of the pins at which the seam functions of the piece and of the partner converge, their reduced periods in Theorem~\ref{thm:transversepins}, which is convex by H\"older's inequality and contains the ray $\mathbf c+\tau\boldsymbol\sigma$, $\tau\ge0$, from a point $\mathbf c\in\Omega$ whenever $\sup\mathbf V_B^{-\boldsymbol\sigma}$ and $\sup\mathbf V_R^{-\boldsymbol\sigma}$ are finite. In Theorem~\ref{thm:pastwalls}(iii) they are the shares and the walls of the seam word, read as affine functions of the shares through their sums (Theorem~\ref{thm:pastwalls}(ii)), and $\Omega=A$. In a fibre integral of Gammas and coefficients, as the share side of a Gramma (Proposition~\ref{prop:shareform}) and the kernel of Theorem~\ref{thm:averaging} below, they are the shares and $\Omega=A$, so that
\[
C(\boldsymbol\sigma)=\{\boldsymbol\sigma'\in V:\ \sigma'_l=\sigma_l\ \text{where}\ \sigma_l<0,\ \ \sigma'_l\ge0\ \text{elsewhere}\},
\]
a compact convex polytope when the pins fix the total of the shares of every letter; the walls of source factors that are cube integrals (Definition~\ref{def:retainedcohort} below) are polar functions as well.
\end{definition}

% Reference 20: thm:chamber; source line 7689 (the closings of thm:closecone and thm:pastwalls(iii)).
% Proved in the last part: thm:chamber.
\begin{theorem}[A closing depends only on its chamber]\label{thm:chamber}
Let $I$, $\mathbf c$ and $\boldsymbol\sigma$ be as in Definition~\ref{def:chamber}, let $I$ have the local decay (H1) of Lemma~\ref{lem:pathclose} on $\Omega$, and let $\boldsymbol\sigma'\in C(\boldsymbol\sigma)$. The rays $\mathbf c+\tau\boldsymbol\sigma$ and $\mathbf c+\tau\boldsymbol\sigma'$ meet the same polar hyperplanes at the same times, and (H0) holds along one of them exactly when it holds along the other; at every other time the fibre integrals through their two points are equal, and at a crossing their crossing terms are equal. So the closings along the two directions are one series, with the same remainders: the closing along $\boldsymbol\sigma$ converges, with the bounds of its terms and of its remainders, as soon as the closing along some direction of its chamber does. In Theorem~\ref{thm:closecone}, or in Theorem~\ref{thm:pastwalls}(iii), if $\mathbf c$ lies off the null set of Theorem~\ref{thm:closecone}, taken for Theorem~\ref{thm:pastwalls}(iii) with the polar hyperplanes of the seam function as well, the closing along $\boldsymbol\sigma$ converges absolutely as soon as
\[
\inf_{\boldsymbol\sigma'\in C(\boldsymbol\sigma)}h(\boldsymbol\sigma')<0,
\]
$h$ the rate of Theorem~\ref{thm:closecone}, the infimum being taken for Theorem~\ref{thm:pastwalls}(iii) over the rates $h_\theta$, $\theta\in[0,\pi/8)$, as well; its terms and the integrals left after them are then $O(e^{\tau(h'+\epsilon)})$ for every $\epsilon>0$, $h'$ that infimum. On the chamber the rate $h$ of Theorem~\ref{thm:closecone} is convex.
\end{theorem}

So the closing along a direction converges as soon as some direction of its chamber is in the cone of Theorem~\ref{thm:closecone}, and its best certified rate is a convex problem in the free coordinates of the chamber. A chamber has positive dimension when the shares that do not fall have free directions on the pins, which stations of one excess supply. Theorem~\ref{thm:closecone} reads both laws along one direction, through the largest values of the monomials of their letter vectors. The glue is also the average over the laws of the stand-in kernel (Theorem~\ref{thm:gluewidth}(ii)), whose integrand at one source pair has no walls, and, with the partner read in its monomials, the average over the law of the piece alone of an integrand without walls; each source point can therefore close along its own best direction of the chamber, and the averages of the crossing terms close the glue.

% Reference 20: thm:pool; source line 7704 (the share sides of the kinds (A) and (B)).
% Proved in the last part: thm:averaging.
\begin{theorem}[Averaging the laws]\label{thm:averaging}
Take a share side, on pins that fix the total of the shares of every active letter, with $V$ the space of the directions of the shares on the pins, whose integrand is
\[
\prod_{j:\,s_j>0}\frac{\prod\Gamma(a_l)}{\Gamma(s_j)}\prod_{j:\,t_j>0}\frac{\prod\Gamma(a_l)}{\Gamma(t_j)}\int e^{-\ip{\mathbf a}{y(\zeta)}}\,d\lambda(\zeta),
\]
the products over the shares $l$ of $K_j$ and of $\Psi_j$, $\mathbf a$ the vector of all the shares, $\lambda$ a positive measure on a source, the \emph{source measure}, whose points are the \emph{source points} $\zeta$, and $y=(y_l)_l$ a continuous map of the source to the space of the shares. In Theorem~\ref{thm:gluewidth} two share sides are of this kind, their shares $l$ those of the active letters:
\begin{enumerate}[label=(\Alph*)]
\item that of Theorem~\ref{thm:gluewidth}(iii), with one share per station; $\lambda=(\boldsymbol\varrho^\rho_B)_\ast\phi^\rho_B\otimes(\boldsymbol\varrho^\rho_R)_\ast\phi^\rho_R$, a source point $\zeta=(\mathbf x,\mathbf x')$ is a seam point of the piece and one of the partner, and $y_l(\zeta)=\log(x_{jD}x'_{jD})$ or $\log(x^\ast_{jD}x'^\ast_{jD})$ for the share of the station $D$ of the letter;
\item that of Theorem~\ref{thm:transversepins} with the partner read in its monomials, one share per monomial, written as in Theorem~\ref{thm:pastwalls}(i), the seam word of the piece having a Newton sum of full dimension; $\lambda=(\boldsymbol\varrho^\rho_B)_\ast\phi^\rho_B$ is the law of the piece, a source point is a seam point $\zeta=\mathbf x$ of the piece, and $y_l(\zeta)=\log x_{jD}+\log\xi_l$ or $\log x^\ast_{jD}+\log\xi_l$ for the share of the monomial of the coefficient $\xi_l$ at the station $D$.
\end{enumerate}
Let $S_\lambda$ be the support of the image of $\lambda$ under $y$, the \emph{log-support} of the source measure, $\mathbf c$ a real point of the pins with positive coordinates at which $\int e^{-\ip{\mathbf c}{y}}\,d\lambda$ converges (in (A) the seam functions of the piece and of the partner converge there, in (B) that of the piece), $\boldsymbol\sigma\in V$, and $C=C(\boldsymbol\sigma)$ the chamber of Definition~\ref{def:chamber} of the kernel below, whose polar functions are the shares:
\[
C=\{\boldsymbol\sigma'\in V:\ \sigma'_l=\sigma_l\ \text{where}\ \sigma_l<0,\ \ \sigma'_l\ge0\ \text{elsewhere}\},
\]
a compact convex polytope. Two shares are in one \emph{bundle} when they belong to one letter and have one column in the pins: in (A) when their stations have one excess, in (B) when moreover they carry one monomial of the partner. For $y$ in the space of the shares put
\[
h(\boldsymbol\sigma';y)=\sum_l\sigma'_l\log|\sigma'_l|-\ip{\boldsymbol\sigma'}y,\qquad h_C(y)=\min_{\boldsymbol\sigma'\in C}h(\boldsymbol\sigma';y),\qquad\hat h_C=\sup_{y\in S_\lambda}h_C(y)\in(-\infty,\infty],
\]
the \emph{rate} of $\boldsymbol\sigma'$ at $y$, the \emph{chamber rate} of the kernel at $y$, concave in $y$, and the \emph{averaged rate}.

Let $\mathbf c$ be interior, relative to the pins, to the set of the real points at which the Laplace transform $\int e^{-\ip{\mathbf a}y}\,d(y_\ast\lambda)(y)$ of $y_\ast\lambda$ converges; equivalently, for a basis $\mathbf v_1,\dots,\mathbf v_d$ of $V$ there is $\eta>0$ with
\begin{equation}\label{eq:avgcert}
\sum_{\boldsymbol\epsilon\in\{-1,1\}^d}\int\exp\Big(-\Big\langle\mathbf c-\eta\sum_i\epsilon_i\mathbf v_i,\,y\Big\rangle\Big)\,d(y_\ast\lambda)(y)<\infty,
\end{equation}
a finite certificate that $\mathbf c$ is interior to the set where that Laplace transform converges, that set being convex. Put $|y|_V=\sum_i|\ip{\mathbf v_i}y|$, let $p$ be the largest total order of the Gamma factors polar on one hyperplane that the ray of $\boldsymbol\sigma$ crosses, at most their number, and in this statement write $w(y)=e^{-\ip{\mathbf c}y}(1+|y|_V)^{p-1}$. Every power of $1+|y|_V$ is at most a constant times $e^{\eta|y|_V}\le\sum_{\boldsymbol\epsilon}\exp(\eta\sum_i\epsilon_i\ip{\mathbf v_i}y)$, so \eqref{eq:avgcert} gives $\int w\,d(y_\ast\lambda)<\infty$.
\begin{enumerate}[label=(\roman*)]
\item \emph{The kernel.} The \emph{kernel} at a source point $\zeta$ is the fibre integral on the pins
\[
\int_{\text{pins}}\prod_{j:\,s_j>0}\frac{\prod\Gamma(a_l)}{\Gamma(s_j)}\prod_{j:\,t_j>0}\frac{\prod\Gamma(a_l)}{\Gamma(t_j)}\ e^{-\ip{\mathbf a}{y(\zeta)}},
\]
of an integrand without walls, and the fibre integral of the share side is its integral against $\lambda$, which in (A) and (B) is $P_{B\sqcup R}(\mvec)$; in (A) the kernel is $\Gamma(\kappa_B)\,\sZ_{(\kappa_B,\mathbf s,\mathbf t);t_{\mathcal D}}(\mathbf x\circ\mathbf x')$ of Theorem~\ref{thm:gluewidth}(ii). Let $\mathbf c$ lie off the null set of Theorem~\ref{thm:closecone} for the direction $\boldsymbol\sigma$. For every $\epsilon>0$ there is $C_\epsilon$, not depending on the source point, such that the crossing terms $\mathcal R_k(\zeta)$ of the closing of the kernel along $\boldsymbol\sigma$, at their crossing times $\tau_k$, satisfy
\[
|\mathcal R_k(\zeta)|\le C_\epsilon\,w(y)\,e^{\tau_k(h_C(y)+\epsilon)},\qquad y=y(\zeta),
\]
and the integrals left at times separating the crossings satisfy the same bound at the time $\tau$ with $e^{-\ip{\mathbf c}y}$ in place of $w(y)$. When every crossed hyperplane carries a simple pole, $p=1$ and $w(y)=e^{-\ip{\mathbf c}y}$.
\item \emph{The averaged closing.} If $\hat h_C<0$, the fibre integral of the share side is
\[
\sum_k\int\mathcal R_k\,d\lambda,
\]
the sum in the order of crossing, absolutely convergent, its terms $O(e^{\tau_k(\hat h_C+\epsilon)})$ for every $\epsilon>0$. At a crossing whose point along some $\boldsymbol\sigma'\in C$ lies in the interior of the set where $\int e^{-\ip{\mathbf a}y}\,d\lambda$ converges, the averaged term is the crossing term of the share side there; so when it converges along the ray of $\boldsymbol\sigma$, the averaged closing is the closing of the share side along $\boldsymbol\sigma$, in (A) and (B) that of Theorem~\ref{thm:closecone}.
\item \emph{One direction reads the convex hull.} $\min_{\boldsymbol\sigma'\in C}\sup_{y\in S_\lambda}h(\boldsymbol\sigma';y)=\sup_{y\in\operatorname{conv}S_\lambda}h_C(y)\ge\hat h_C$, and $\sup_{y\in S_\lambda}h(\boldsymbol\sigma';y)$ is the rate of the closing of the share side along $\boldsymbol\sigma'$ read through the largest value of $e^{-\ip{\boldsymbol\sigma'}y}$ on the log-support, in (A) and (B) the rate $h(\boldsymbol\sigma')$ of Theorem~\ref{thm:closecone}: the best single direction of the chamber reads the convex hull of the log-support, the averaged closing the log-support itself. In a bundle $G$ of shares that $\boldsymbol\sigma$ does not fall, the least value of $\sum_{l\in G}(\sigma'_l\log\sigma'_l-\sigma'_ly_l)$ over the $\sigma'_l\ge0$ with $\sum_G\sigma'_l=g$ is $g\log g-g\log\sum_{l\in G}e^{y_l}$, attained at $\sigma'_l=g\,e^{y_l}/\sum_{l'\in G}e^{y_{l'}}$: at each source point the best direction merges the bundle into the sum $e^{y_G}=\sum_{l\in G}e^{y_l}$, which in (A) is the sum $\sum_Dx_{jD}x'_{jD}$ that the kernel reads at the excess, splitting $g$ in the proportions of its terms.
\end{enumerate}
\end{theorem}

At width one the share side of a node merged by its bundles has at most one free share and the seam torus is a half-line, so the log-support is an interval, its own hull, and one direction loses nothing. A gap between the averaged rate and the best single direction needs a chamber of positive dimension and a log-support that is not convex in its directions: the best split at each source point moves with it, and one direction must take one split for all of them. Averaging can reach a rate that no single direction certifies, and an averaged closing is in general no closing of one share side along one ray: its terms are averages of the crossing terms of the kernel, and it can converge past walls that every ray of the chamber meets. Where the averaged rate is infinite, a closing past the walls can still converge, at the rate of the best source point: past its walls the law is read through its expansion at its falling end, not through its moments, as at width one (Theorem~\ref{thm:lawreach}). Past the walls of a piece the bound that travels is read on the sector letters (Theorem~\ref{thm:keyholebound}).

% Reference 20: thm:canonicaltablegap; source line 7810
\begin{theorem}[An averaging gap at a family of tables]\label{thm:canonicaltablegap}
There are three tables with the bases $\binom{\{1,2,3,4\}}2$, at the exponents $s_j=1$, $t_j=0$ and the masses $\frac32$ on every mark, whose bunch $\{1,2\}$, with the room $\{3,4\}$, read through its seam with the product of the two laws as the source measure, the kind (A) of Theorem~\ref{thm:averaging}, has one station at each excess of each letter, so that every bundle of Theorem~\ref{thm:averaging} is a single share, and a chamber $C$ with
\[
\min_{\boldsymbol\sigma\in C}\sup_{y\in S_\lambda}h(\boldsymbol\sigma;y)=0,\qquad\hat h_C=\sup_{y\in S_\lambda}\min_{\boldsymbol\sigma\in C}h(\boldsymbol\sigma;y)<0:
\]
no direction of the chamber certifies the closing (Theorem~\ref{thm:closecone}), and the averaged closing converges geometrically (Theorem~\ref{thm:averaging}(ii)). Two directions of the chamber, the better one taken at each source point, already give the bound.
\end{theorem}

% Reference 20: rem:gapscope; source line 7867
\begin{remark}[What the gap separates]\label{rem:gapscope}
The gap lies between the averaged closing and the directions of one share side, the seam of the bunch with its bundles merged, and no merging of stations removes it, every bundle being a single share. It is not a gap against every closing of the period: closing each source point along its own direction, as the bound does with two directions, is itself a closing of the period, and so is completing another inner read before closing. A gap in every closing is a statement about a specified set of closings; for the closings of one share side by the directions of its chamber, the family of Theorem~\ref{thm:canonicaltablegap} has one, while at a seam whose two stations of one excess multiply the same monomial of the partner, merging the two stations can remove such a gap.
\end{remark}

% Reference 20: def:dirkeyhole; source line 7871 (the orders of a piece of tables).
% Proved in the last part: def:dirkeyhole.
\begin{definition}[Keyhole continuations along a direction; ceiling and phase]\label{def:dirkeyhole}
In Theorem~\ref{thm:pastwalls} let $\boldsymbol\sigma$ be a direction of the shares and $\bar{\boldsymbol\sigma}$ its sums over the monomials of the partner, and fix a fan of the seam word (Theorem~\ref{thm:pastwalls}(ii)). For a sector $\pi$ of it let $q=(q_k)_k$ be its sector coordinates, $\kappa_k=\kappa_{\mathbf v_{\pi,k}}$ the walls of its generators and $\mathbf k_\pi=(k_{jD,\pi},\psi_{jD,\pi})_{j,D}$ the vector of its sector letters, so that the seam function on the pins is the sum over the sectors of
\[
|\det(\mathbf v_{\pi,1},\dots,\mathbf v_{\pi,|B|-1})|\int_{[0,1]^{|B|-1}}\prod_kq_k^{\kappa_k(\bar{\mathbf a},\bar{\mathbf b})-1}\,\mathbf k_\pi(q)^{-(\bar{\mathbf a},\bar{\mathbf b})}\,dq,
\]
the walls affine in the shares; for a piece of tables and the sectors of the orders $\pi$ of $B$ these are the walls $\kappa_{U_k}$ of the tails and the sector letters of the tables $\cB_{jD}$ (Theorem~\ref{thm:exist}(ii)). Put $\dot\kappa_k=\partial_{\bar{\boldsymbol\sigma}}\kappa_k$, the derivative of the wall along $\bar{\boldsymbol\sigma}$, and
\[
\varrho_{\pi,\bar{\boldsymbol\sigma}}(q)=\prod_kq_k^{\dot\kappa_k}\ \mathbf k_\pi(q)^{-\bar{\boldsymbol\sigma}},
\]
the \emph{directional ratio in the sector} of $\pi$; its \emph{falling variables} along $\bar{\boldsymbol\sigma}$ are the $q_k$ with $\dot\kappa_k<0$. A \emph{keyhole continuation along} $\bar{\boldsymbol\sigma}$ is defined as in Definition~\ref{def:keyhole}, with the sectors in place of the orders, with these falling variables, and with the sector letters of $\mathbf k_\pi$ required not to vanish on the sets $Y$ of its boxes; its \emph{outline} $\mathcal O$ is the union of the sets $\mathcal O$ of its boxes, and its \emph{ceiling along} $\bar{\boldsymbol\sigma}$ is the largest value of $|\varrho_{\pi,\bar{\boldsymbol\sigma}}|$ on the outline, $\pi$ the sector of the box. Each set $Y$ is a product of a box and of keyholes, simply connected, and on a neighbourhood of it the sector letters have logarithms real at the real points; the \emph{phase} of the continuation at a point $q$ of the outline of a box of $\pi$ is the linear function $\boldsymbol\vartheta\mapsto\ip{\bar{\boldsymbol\vartheta}}{\arg\mathbf k_\pi(q)}$ on the directions of the shares, $\arg$ the imaginary part of the logarithm and the bars marking sums over the monomials of the partner.
\end{definition}

The modulus $|\varrho_{\pi,\bar{\boldsymbol\sigma}}(q)|$ is the factor by which the part of the integrand that grows with $\tau$ along $\bar{\boldsymbol\sigma}$ is multiplied per unit of $\tau$ at the point; the phase is the linear function by which the imaginary parts of the shares enter it. With no falling variables the outline is the real cube, the letters are positive there, the phase vanishes, and the ceiling is $\sup\mathbf V_B^{-\bar{\boldsymbol\sigma}}$, the largest value of the monomial of Theorem~\ref{thm:closecone}; at width one, along the fall of the share $b^d$ of Section~\ref{sec:onescale}, the directional ratio in the sector is the ratio $\varrho_\pi$ of Definition~\ref{def:falling}.

% Reference 20: thm:keyholebound; source line 7886 (a piece of tables).
% Proved in the last part: thm:keyholebound.
\begin{theorem}[The travelling bound past the walls]\label{thm:keyholebound}
In Theorem~\ref{thm:pastwalls} let $\boldsymbol\sigma$ be a direction of the shares and take a keyhole continuation along $\bar{\boldsymbol\sigma}$ (Definition~\ref{def:dirkeyhole}) with the outline $\mathcal O$.
\begin{enumerate}[label=(\roman*)]
\item For every $\epsilon>0$ there is $C$ such that for $\tau\ge0$ and $\boldsymbol\vartheta$ in the directions of the shares
\[
\big|P_{W_B(\bar{\mathbf c}+\tau\bar{\boldsymbol\sigma}+i\bar{\boldsymbol\vartheta})}(\mvec_B)\big|\le C\,\delta^{-(|B|-1)}\sup_{q\in\mathcal O}\Big(|\varrho_{\pi,\bar{\boldsymbol\sigma}}(q)|^\tau\,e^{\tau\epsilon+\ip{\bar{\boldsymbol\vartheta}}{\arg\mathbf k_\pi(q)}}\Big),
\]
$\delta$ the smaller of $1$ and the distance of $\bar{\mathbf c}+\tau\bar{\boldsymbol\sigma}$ from the hyperplanes on which a wall of the fan of Definition~\ref{def:dirkeyhole} lies in $-\mathbb N$, which hold the poles of the seam function.
\item Let $\boldsymbol\upsilon$ be a bounded function on $\mathcal O$ with values in the vectors indexed by the shares, with $\ip{\bar{\boldsymbol\vartheta}}{\arg\mathbf k_\pi(q)}=\sum_l\upsilon_l(q)\vartheta_l$ for $\boldsymbol\vartheta$ in the directions of the shares and $|\upsilon_l(q)|\le\frac\pi2-\epsilon_0$ for some $\epsilon_0>0$; representatives differing by a vector orthogonal to the directions give the same phase on the fibre. Put
\[
h_{\boldsymbol\upsilon}(\boldsymbol\sigma)=h_\Gamma(\boldsymbol\sigma)+\sup_{q\in\mathcal O}\Big(\log|\varrho_{\pi,\bar{\boldsymbol\sigma}}(q)|-\sum_{l:\,\sigma_l>0}\sigma_l\log\cos\upsilon_l(q)\Big),
\]
$h_\Gamma$ the Gamma rate of Theorem~\ref{thm:pastwalls}(iii). If $h_{\boldsymbol\upsilon}(\boldsymbol\sigma)<0$ and the real point $\mathbf c$ of the pins lies off the null set of Theorem~\ref{thm:closecone}, taken with these hyperplanes as well, the conclusion of Theorem~\ref{thm:pastwalls}(iii) holds along $\boldsymbol\sigma$, its terms and the integrals left after them being $O(e^{\tau(h_{\boldsymbol\upsilon}(\boldsymbol\sigma)+\epsilon)})$ for every $\epsilon>0$.
\end{enumerate}
\end{theorem}

% Reference 20: cor:phasefree; source line 7924 (a piece of tables).
% Proved in the last part: cor:phasefree.
\begin{corollary}[When the phase costs nothing]\label{cor:phasefree}
In Theorem~\ref{thm:keyholebound}(ii) let $H_{\mathcal O}$ be the ceiling of the continuation along $\bar{\boldsymbol\sigma}$. Then $h_{\boldsymbol\upsilon}(\boldsymbol\sigma)=h_\Gamma(\boldsymbol\sigma)+\log H_{\mathcal O}$ in each of two cases:
\begin{enumerate}[label=(\roman*)]
\item $\upsilon_l(q)=0$ for every $q\in\mathcal O$ and every share with $\sigma_l>0$: the phase sits on shares that fall or stay, whose Gammas absorb it;
\item $-\sum_{l:\,\sigma_l>0}\sigma_l\log\cos\upsilon_l(q)\le\log H_{\mathcal O}-\log|\varrho_{\pi,\bar{\boldsymbol\sigma}}(q)|$ for every $q\in\mathcal O$: the modulus lost away from the highest points pays for the phase.
\end{enumerate}
With no falling variables (i) holds with $\boldsymbol\upsilon=0$, and the rate is that of Theorem~\ref{thm:closecone} for the glue read in the monomials of the partner.
\end{corollary}

Theorem~\ref{thm:pastwalls}(iii) pays for an allowed growth $e^{\theta|\im\mathbf z|_1}$ of the seam function with the term $\frac{2\theta^2}\pi\sum_{\sigma_l>0}\sigma_l$; Theorem~\ref{thm:keyholebound} pays only for the phase that the letters have at the points where the ceiling is read, and at the exact price $-\log\cos\upsilon_l$ of the rotation. Enlarging the radii of the keyholes, as far as the sector letters do not vanish, can lower the ceiling; it comes down to the reach at simple tops that do not cancel (Theorem~\ref{thm:dirheights}(iii)). The joint words of Section~\ref{sec:retainedcrossings} complete the retained fibre on its strips before a modulus is taken, and their bound $r_Gc_{\max}$ has no phase term; the ceiling $c_{\max}$ of the joint word need not be the ceiling of this outline.

% Reference 20: prop:pathcompose; source line 7938 (the orders of a piece of tables).
% Proved in the last part: prop:pathcompose.
\begin{proposition}[How travelling bounds compose]\label{prop:pathcompose}
\begin{enumerate}[label=(\roman*)]
\item \emph{Products.} When the share integrand carries several continued factors, each bounded as in Theorem~\ref{thm:keyholebound}(i) on its own outline, their product is bounded on the product of the outlines by the products of their moduli and of the exponentials of their phases, and Theorem~\ref{thm:keyholebound}(ii) holds with these on the product outline: the phases of different factors at one point of it may cancel.
\item \emph{Averaging.} Let integrands $I_\zeta$, for $\zeta$ in a source with a measure $\lambda$, have fibre integrals that close along one path with their crossings at the same times, with the crossing terms $\mathcal R_k(\zeta)$, and with crossing terms and remainders, at times $\tau_k$ separating the crossings, at most $M(\zeta)\beta(\tau_k)$, where $M$ is integrable against $\lambda$ and $\sum_k\beta(\tau_k)<\infty$. Then $\int\big(\int_{\mathbf c+iV}I_\zeta\big)\,d\lambda(\zeta)=\sum_k\int\mathcal R_k\,d\lambda$, absolutely convergent, the $k$-th term at most $\|M\|_1\beta(\tau_k)$. When for some $\hat h<0$ the bound holds with $\beta(\tau)=e^{\tau(\hat h+\epsilon)}$ for every small $\epsilon>0$, $M$ depending on $\epsilon$, the terms are $O(e^{\tau(\hat h+\epsilon)})$: the rate is the bound $\hat h$ shared by the source points, read on the support of $\lambda$, the uniformity over the source being the hypothesis on $M$. Where the fibre integral and the source integral of a crossing term converge absolutely together, $\int\mathcal R_k\,d\lambda$ is the crossing term of the averaged integrand $\int I_\zeta\,d\lambda$.
\item \emph{Crossing.} The crossing terms of the closings of Theorems~\ref{thm:pastwalls}(iii) and~\ref{thm:keyholebound} are fibre integrals of one dimension less of integrands of the same kind. At a share crossing $a_l=-n$ the Gamma becomes $(-1)^n/n!$ and the other factors are read on the hyperplane. At a wall crossing where the pole of the seam function is simple, its residue is the \emph{face coefficient} of the walls through the hyperplane: the sum, over the generators $\mathbf v$ of a fan of the seam word (Theorem~\ref{thm:pastwalls}(ii)) whose wall $\kappa_{\mathbf v}$ takes a value $-n_{\mathbf v}\in-\mathbb N$ on the hyperplane and over the sectors $\pi$ of which they are generators, of $|\det(\mathbf v_{\pi,1},\dots,\mathbf v_{\pi,|B|-1})|$ times the cube integral over the other coordinates of the sector of the $n_{\mathbf v}$-th Taylor coefficient at $q_{\mathbf v}=0$ of the rest of the integrand (Lemma~\ref{lem:cube}), $q_{\mathbf v}$ the sector coordinate of $\mathbf v$, each divided by the derivative of $\kappa_{\mathbf v}$ along the vector $\mathbf e$ of Lemma~\ref{lem:pathclose}; its letters are the sector letters at $q_{\mathbf v}=0$, positive on the face, and its numerators are polynomials in the shares and in the derivatives of the letters. For a piece of tables and the sectors of the orders the generators are those of the tails $U$, and the sectors the orders of which they are tails. Where the pole is multiple, because several factors, or two walls of one sector, are polar on the hyperplane, the residue is a Laurent coefficient, a finite sum of derivatives in the shares of products of such terms, which brings powers of the logarithms of the letters and of the $q_k$. A keyhole continuation of the seam function restricts to one of each face coefficient, $q_{\mathbf v}=0$ lying in the disc of the keyhole of $q_{\mathbf v}$, and Cauchy's estimate on its circle bounds the Taylor coefficients by the values on the outline; so each crossing term can be closed again by Lemma~\ref{lem:pathclose} whenever a travelling bound of its own tends to zero.
\end{enumerate}
\end{proposition}

Completing an inner read is not a crossing. Merging the stations of one excess by Corollary~\ref{thm:E1}, integrating a source before closing, or finishing a transverse fibre replaces the integrand by a completed one, with a smaller fibre; the identity between the two integrands holds where the completed read converges and persists under continuation, but the travelling bound does not pass from one to the other, and the completed object can grow along the outer path where the uncompleted integrand does not. Completing an inner read is free when the completed object obeys, along the outer path, a bound already in hand. Merging the stations of one excess reads the law through the sum of their seam letters, whose expansion at the falling end of the law can stop at an edge of the support nearer than the far edge that the stations read, and the merged closing can then diverge where the station closing converges. The two closings agree at every step up to their remainders (Lemma~\ref{lem:pathclose}), and only the remainder of the station closing need tend to zero. On the shares, merging the stations is conditioning on their merged seam letter: integration commutes with completing the fibre, and continuation commutes with it only while the level sets of the merged seam letter stay inside the region where the expansion of the law at its falling end represents it. Completing the inner fibre before closing and keeping its inherited cut are two representations, whose finite heads agree together with their boundary words; the regions on which a law equals its expansion, with their cuts and boundary words, are built in Section~\ref{sec:lawregions}, and a convergent germ alone does not equal the law.

\section{Retained crossings as joint words}\label{sec:retainedcrossings}

\subsection{Completing the retained fibre before taking its size}

Theorem~\ref{thm:keyholebound} bounds the share integrand before its fibre is completed, and pays for the phase of the letters where it reads the ceiling. The procedure for joint words completes the retained shares first. It keeps the sectors and coordinate cuts of every continued source and its finite heads, and its auxiliary variables are part of one joint input of a Trailmoment (Definition~\ref{def:jointtrail}), their number not growing along a retained cohort (Definition~\ref{def:retainedcohort}): a falling Gamma kept inside a crossing term becomes a positive share and one Beta variable, and a continued source strip becomes a segment and a few circle integrals whose powers on the circles have integral exponents.

% Reference 20: lem:gammastriplift; source line 7983
\begin{lemma}[A Gamma strip as a Beta input]\label{lem:gammastriplift}
Let $\delta\ge1$ be an integer, $b=a+\delta$, and $0<\re b<1$. For a nonzero coefficient $\xi$, with one fixed logarithm,
\[
\Gamma(a)\,\xi^{-a}=\frac{(-1)^\delta\xi^\delta}{\Gamma(\delta)}\int_0^1(1-u)^{\delta-1}\,\Gamma(b)\,(\xi u)^{-b}\,du .
\]
The identity is pointwise in every remaining share. So on a fibre $P\mathbf a=\boldsymbol\pi$ replacing $a_l$ by $b_l=a_l+\delta_l$ changes the pins to $P\mathbf b=\boldsymbol\pi+P(\delta_l)_l$ and does not move the contour. A share with positive real part is left unchanged, with $\delta_l=0$ and no auxiliary variable. In a family $\re a_l(n)=a^0_l-d_ln$ with $d_l\in\mathbb N$ and $a^0_l\notin\mathbb Z$ one may take $\delta_l(n)=d_ln+\delta_l(0)$ with $0<a^0_l+\delta_l(0)<1$, the lift being used at the $n$ with $\delta_l(n)\ge1$, all but finitely many when $d_l\ge1$: one Beta input serves at every $n$.
\end{lemma}

% Reference 20: lem:sourcestriplift; source line 8009
\begin{lemma}[A continued source strip in six integrals]\label{lem:sourcestriplift}
Let $g$ be holomorphic near $\{|q|\le r_2\}\cup[r,1]$, and choose $0<r<r_1<rT<r_2$ and $T>1$. If $\ell\ge1$ is an integer and $0<\re(\kappa+\ell)<1$, the continued integral $I_\kappa=\Reg\int_0^1q^{\kappa-1}g(q)\,dq$ (Definition~\ref{def:Reg}) is
\[
I_\kappa=\int_r^1q^{\kappa-1}g(q)\,dq+R_\ell+R'_\ell,
\]
all circles run counterclockwise, with
\begin{align*}
R_\ell&=\frac1{2\pi i}\int_0^rq^{\kappa+\ell-1}\oint_{|\zeta|=r_1}\frac{g(\zeta)\,\zeta^{-\ell}}{\zeta-q}\,d\zeta\,dq,\\
R'_\ell&=-\frac{r^\kappa}{2\pi i}\Big(\int_1^Tu^{\kappa-1}\oint_{|\zeta|=r_2}+\int_T^\infty u^{\kappa-1}\oint_{|\zeta|=r_1}\Big)\,g(\zeta)\,\frac{1-(ru/\zeta)^\ell}{\zeta-ru}\,d\zeta\,du .
\end{align*}
Splitting the numerator of $R'_\ell$ gives four absolutely convergent integrals, six in all. This remains true with holomorphic parameters and other source variables on compact sets, uniformly when the strip and the radius inequalities have positive margins. For $\kappa(\mathbf a,n)=\kappa_0(\mathbf a)-dn$ and $\ell(n)=\ell_0+dn$, $d\in\mathbb N$, each of the six has a fixed domain and is a fixed function times an $n$-th power; the exponent of $\zeta$ that contains $\ell(n)$ is integral and independent of the shares that move. At a simple wall $\kappa=-N$, $N\in\mathbb N$, the residue is the single coefficient input
\[
[q^N]g(q)=\frac1{2\pi i}\oint_{|\zeta|=r_1}g(\zeta)\,\zeta^{-N-1}\,d\zeta .
\]
A cut upper end $q_0>0$ is kept by applying the same formula after $q=q_0v$, the factor $q_0^\kappa$ included; a positive lower cut is an ordinary segment throughout.
\end{lemma}

Two of the six integrals simplify by Cauchy's formula: the integral of $R'_\ell$ over $[1,T]$ with the numerator $1$ is $-\int_r^{rT}q^{\kappa-1}g(q)\,dq$, and the integral over $(T,\infty)$ with the numerator $1$ vanishes, its pole $ru$ lying outside the smaller circle. So $I_\kappa$ is the segment $\int_{rT}^1q^{\kappa-1}g$, the remainder $R_\ell$ on the circle of radius $r_1$, and the two integrals of $R'_\ell$ with the numerator $(ru/\zeta)^\ell$; along a family $\kappa_0-dn$ the segment carries the $n$-th power of $q^{-d}$, at most $(rT)^{-d}$, and the circle integrals that of $\zeta^{-d}$, of modulus $r_1^{-d}$ or $r_2^{-d}$. The radius $r$ splits the contour, not the source: the sum does not depend on it.

% Reference 20: def:retainedcohort; source line 8078 (the share form of thm:pastwalls).
% Proved in the last part: def:retainedcohort.
\begin{definition}[A rational retained cohort]\label{def:retainedcohort}
A \emph{retained cohort} is one sequence of simple crossing terms, or one sequence of travelling fibre integrals, of a fibre integral of Gammas, coefficients and source factors whose pins carry a \emph{grading}, a rational vector $\mathbf v$ with $\ip{t_l}{\mathbf v}=1$ for the column $t_l$ of every share in the pins, such as the share side of Theorem~\ref{thm:pastwalls} and the kernel of Theorem~\ref{thm:averaging}, graded by their letter totals, and the share side of a Gramma whose configuration has a grading (Proposition~\ref{prop:shareform}), with the following data retained: the pins; the inherited contour; the crossed hyperplane and its orientation, if any; the source sectors and their coordinate cuts. The contours have the real points $\mathbf a^0+n\mathbf d$ with $\mathbf d$ rational. Nonresonance is tested against the full displayed arrangement of the polar hyperplanes of the Gammas and of the source walls, including candidate walls which might cancel after the sectors are added: the real points lie on none of these hyperplanes other than the one being crossed, and at a source crossing each contributing sector has just one wall on that hyperplane. The constant terms of the shares and walls and the indices are real. Finitely many initial terms may be removed and kept as a finite head. Passing to finitely many arithmetic subsequences makes the changes of all falling shares and walls integral; this is the \emph{integer step} of the cohort, and its index $n$ always refers to that step. The same definition permits a finite product of source factors that are cube integrals in their charts (Lemma~\ref{lem:cube})
\[
\Reg\int_{(0,1)^k}q^{\boldsymbol\kappa(\mathbf a)-\one}\prod_iP_i(q)^{-w_i(\mathbf a)}\,dq,
\]
the $P_i$ with nonnegative coefficients and positive constant terms and the affine functions $\kappa_i$, $w_i$ with rational linear parts. A growing wall is eventually positive; a stationary wall may be negative and off the integers, and is then continued with a constant integer $\ell$ in Lemma~\ref{lem:sourcestriplift}. Finite coordinate box cuts are part of the retained data. The sector integrals of a fan of a seam word (Theorem~\ref{thm:pastwalls}(ii)), those of tables among them (Corollary~\ref{cor:pastwalls}), have these properties.
\end{definition}

% Reference 20: prop:cohortselection; source line 8112
\begin{proposition}[Computing the cohorts of a rational ray]\label{prop:cohortselection}
Let the pins and the finitely many candidate polar functions $\ell_i$ be rational, and let a rational open box in the affine space of the pins be given on which the required initial walls and shares are positive. For a prescribed rational direction a rational starting point in that box can be chosen such that distinct polar hyperplanes never meet the ray at the same time; factors on the same hyperplane are kept together. After a finite head, the crossing order and the separating gaps are a finite repeated pattern, and its crossing cohorts, endpoint cohorts, the integers $\ell$ of Lemma~\ref{lem:sourcestriplift} and a uniform positive separation are computed by finite rational arithmetic. More precisely, with the direction multiplied by a common denominator so that every falling polar function decreases at an integer falling rate $d_i>0$, there are at most $\sum_id_i$ crossings in each interval of length one of this parameter; neither that number nor the stored source sectors depend on how many such intervals are read.
\end{proposition}

% Reference 20: thm:retainedchain; source line 8164 (the share form of thm:pastwalls, graded by its letter totals).
% Proved in the last part: thm:retainedchain.
\begin{theorem}[A finite procedure for a retained cohort]\label{thm:retainedchain}
For a rational retained cohort there is a finite procedure, using the pins and the monomials of the source factors, with the following output, the \emph{joint word} of the cohort, after a finite head and a possible change to the integer step:
\begin{equation}\label{eq:retainedmomentword}
\mathcal R_n=G_n\sum_i\int_{Y_i}c_i(x)^n\,\phi_i(dx)=G_n\sum_i\big(\sT_{-n;(1)}[c_i]\,\phi_i\big)(Y_i),\qquad n\in\mathbb N,
\end{equation}
a finite sum over $i$. Here $Y_i$ is a compact cube times finitely many circles, the coefficient function $c_i$ is continuous on it, and the input $\phi_i$ is a complex measure of finite total variation; the coefficients, their branches, the orientations and every cut are retained. On every $Y_i$ each coefficient of a letter is a real Laurent monomial times a \emph{coefficient factor}, holomorphic there, with its real part bounded below by a positive constant and its modulus bounded above, uniformly on the compact variables; these two bounds are the \emph{margins} of the coefficient factors, and the radii $r<r_1<rT<r_2$ of Lemma~\ref{lem:sourcestriplift} are admissible while the margins hold. At the integer index $-n$ the Zamma of the one-term configuration is the polynomial $\sZ_{-n;(1)}(c)=c^n$, at $c=0$ as well, so the second expression is the Trailmoment of Definition~\ref{def:jointtrail} extended to a continuous complex coefficient function and a complex input of finite variation, as Theorem~\ref{thm:jointclosure}(ii) extends its reads to complex data where every integral converges absolutely.

The scalar is explicit. Let $\mathbf v$ be the grading of the pins (Definition~\ref{def:retainedcohort}) and $w=\ip{\mathbf v}{\boldsymbol\nu}$, $\boldsymbol\nu$ the values of the pins, the total of the shares on the pins, which on the share sides of Theorems~\ref{thm:pastwalls} and~\ref{thm:averaging} is $|\mathbf s|+|\mathbf t|$, the sum of the exponents of the active letters. Let $\mathcal S_-$ be the shares treated by Lemma~\ref{lem:gammastriplift}, with $\delta_l(n)=\delta_l(0)+d_ln$. At a Gamma crossing let $-N(n)$, with $N(n)=N_0+d_0n$, be the pole of its Gamma that is crossed; at a source crossing or for a travelling fibre omit $N$ and the factor $\Gamma(N+1)$ of $G_n$, and put $d_0=0$. Up to fixed nonzero factors of the share side, as $1/\Gamma(s_j)$ and $1/\Gamma(t_j)$, and the inherited crossing sign,
\[
G_n=\frac{(-1)^{N(n)+\sum_{l\in\mathcal S_-}\delta_l(n)}\,\Gamma\big(w+N(n)+\sum_{l\in\mathcal S_-}\delta_l(n)\big)}{\Gamma(N(n)+1)\prod_{l\in\mathcal S_-}\Gamma(\delta_l(n))},
\]
fixed nonzero constants may equally be put in the inputs. With $\bar d=d_0+\sum_{l\in\mathcal S_-}d_l$,
\[
\lim_{n\to\infty}|G_n|^{1/n}=r_G=\frac{\bar d^{\,\bar d}}{d_0^{\,d_0}\prod_{l\in\mathcal S_-}d_l^{\,d_l}},\qquad0^0=1;
\]
when $\bar d=0$ the scalar is constant and $r_G=1$, and in all cases $|G_n|\le C(n+1)^pr_G^{\,n}$ for computable $C$ and $p$. The procedure introduces at most six integrals per continued source wall and one Beta variable per negative retained Gamma, whatever $n$. Its remaining work is a finite common Newton fan and simplicial subdivision; this fan, the inputs and the coefficient functions $c_i$ serve for every $n$, and the procedure does not expand the Taylor head of degree $\ell(n)-1$ of Lemma~\ref{lem:sourcestriplift} or a residue coefficient of degree $N(n)$ into its monomials.
\end{theorem}

% Reference 20: cor:retainedheight; source line 8412
\begin{corollary}[Finite bounds and the exact positive root rate]\label{cor:retainedheight}
For the joint word of Theorem~\ref{thm:retainedchain} put $c_{\max}=\max_i\max_{Y_i}|c_i|$ and $M=\sum_i\|\phi_i\|$. Then
\[
|\mathcal R_n|\le M\,|G_n|\,c_{\max}^{\,n}\le CM(n+1)^p(r_Gc_{\max})^n,
\]
and $c_{\max}$ is the \emph{ceiling} of the joint word. Every quantity in this bound has a finite certificate: the margins of the coefficient factors and the sector walls bound $M$, and compact subdivision bounds $c_{\max}$; with computable input coefficients a strict inequality $r_Gc_{\max}<r<1$ is certified by finite subdivision and gives a computable geometric tail. If the original source has a convergent positive representation on every contour of the cohort, no source continuation or coefficient input is needed, and the remaining constant terms are real, then the output of the procedure has, after its common scalar sign is removed, one positive input $\phi$ and a coefficient function $c\ge0$, and
\[
\lim_{n\to\infty}|\mathcal R_n|^{1/n}=r_G\operatorname*{ess\,sup}_\phi c .
\]
For a positive read on an unbounded domain, before it is compactified, the same holds with an infinite essential supremum whenever the moments exist. So averaging over a positive source measure takes the essential supremum on the joint input of the source, Beta and shape variables; it does not reverse the order of an infimum over directions and a supremum over source points.
\end{corollary}

% Reference 20: def:retainedtops; source line 8460
\begin{definition}[Simple tops of the joint word]\label{def:retainedtops}
For \eqref{eq:retainedmomentword} with $c_{\max}>0$, tops are read as in Definition~\ref{def:intkeyhole}, with the coefficient function $c_i$ in place of $\varrho_\pi$, the oriented $Y_i$ in place of the keyhole contour and the input $\phi_i$ in place of $\phi_\pi$: a simple top is a point $p$ with $|c_i(p)|=c_{\max}$ for which local real coordinates on the oriented $Y_i$ are $(\mathbf x,\mathbf y)\in[0,\epsilon)^{b_1}\times(-\epsilon,\epsilon)^{b_2}$, and
\[
\log\frac{c_i(\mathbf x,\mathbf y)}{c_i(p)}=-\ip{\boldsymbol\lambda}{\mathbf x}-\tfrac12\,\mathbf y^{\mathsf T}\mathbf Q\,\mathbf y+O\big(|\mathbf x|^2+|\mathbf x|\,|\mathbf y|+|\mathbf y|^3\big),\qquad\phi_i=\prod_{j=1}^{b_1}x_j^{\kappa_j-1}\,g(\mathbf x,\mathbf y)\,d\mathbf x\,d\mathbf y,
\]
with $\kappa_j>0$, $g(0,0)\neq0$, $\re\lambda_j>0$ and $\re\mathbf Q$ positive definite, $\mathbf Q$ complex symmetric, and $g$, continuous and bounded near $(0,0)$, including the orientation and every factor of the circle integrals. Its exponent and its coefficient are
\[
\operatorname{ex}(p)=\sum_j\kappa_j+\tfrac{b_2}2,\qquad A(p)=g(0,0)\,(2\pi)^{b_2/2}(\det\mathbf Q)^{-1/2}\prod_j\Gamma(\kappa_j)\,\lambda_j^{-\kappa_j},
\]
the square root continued from positive definite real matrices. The subdivisions of the procedure, the splits of the Beta intervals and of the scales and the cones of the fan, divide one domain into parts on which the integrand is the restriction of one function; tops are read with these parts reassembled, a point on a common boundary of the parts being a simple top when the expansion holds on its reassembled neighbourhood, in the coordinates before the subdivision. Parts that are different integrals, as the segment and the circle integrals of Lemma~\ref{lem:sourcestriplift}, keep their own coefficients, and the contributions at one point are added before noncancellation is tested.
\end{definition}

% Reference 20: thm:retainedtops; source line 8487
\begin{theorem}[The joint tops give the root rate of a complex word]\label{thm:retainedtops}
Let the largest values of $|c_i|$ on the joint word of Theorem~\ref{thm:retainedchain} sit at finitely many simple tops. Let $\operatorname{ex}_0$ be the least of their exponents, and for $|\zeta|=1$ let $A_\zeta$ be the sum of the coefficients $A(p)$ of the tops with $\operatorname{ex}(p)=\operatorname{ex}_0$ and $c_i(p)=c_{\max}\,\zeta$. If some $A_\zeta$ is not zero, then
\[
\limsup_{n\to\infty}|\mathcal R_n|^{1/n}=r_G\,c_{\max} .
\]
No separate price of the phase is added to $c_{\max}$: the completed fibre and its source phases belong to the one coefficient function and the oriented input. A largest value of $|c_i|$ alone is not the hypothesis; the expansion includes the stationarity of the complex phase in every tangential variable.
\end{theorem}

This is Theorem~\ref{thm:simpletop} for joint words: the coordinates $\mathbf x$ that vanish at a top carry the Gammas of the boundary, the coordinates $\mathbf y$ the factor $(2\pi)^{b_2/2}(\det\mathbf Q)^{-1/2}$, and the coefficients of the tops of least exponent and one phase add before the root rate is read.

% Reference 20: rem:continuedtops; source line 8530
\begin{remark}[Where the tops of a continued source factor sit]\label{rem:continuedtops}
In the integrals of Lemma~\ref{lem:sourcestriplift} a falling wall carries the $n$-th power of $q^{-d}$ on the segment and of $\zeta^{-d}$ on the circles. On a circle the modulus of $\zeta^{-d}$ is constant and its phase turns, so a largest value of $|c_i|$ there is a simple top only where the other factors make the phase stationary, a saddle in the radius. The radii $r$, $r_1$, $r_2$ and $T$ are free within the margins of the coefficient factors, and the sum does not depend on them; Theorem~\ref{thm:retainedtops} is applied after the radii are chosen. When the largest values of $|c_i|$ sit at simple tops of the segment, of the shape or of the Beta inputs, or at saddles of the circles, and do not cancel, the root rate is $r_Gc_{\max}$; a largest value spread over a circle, or one at the end $q=rT$ of the segment, which moves with the radii, certifies the upper bound only, and the radii may be moved, within the margins, to lower it. Where the law has log expansions the continuation is made on the law, and its bound has no phase and no margin (Proposition~\ref{prop:logexpansions}).
\end{remark}

% Reference 20: cor:pathclosevalues; source line 8534
\begin{corollary}[Closing along a path under an endpoint condition]\label{cor:pathclosevalues}
In Lemma~\ref{lem:pathclose} keep (H0) and (H1) and replace (H2) by the weaker condition $\int_{\gamma(\tau_k)+iV}I\to0$ along times $\tau_k\to\infty$ off the poles. Then the sums of the signed crossing terms before $\tau_k$ converge to $\int_{\mathbf c+iV}I$. If the crossing terms are absolutely summable, this is their absolutely convergent ordered sum; (H2) remains a sufficient condition for the endpoint condition.
\end{corollary}

% Reference 20: cor:retainedclosing; source line 8556
\begin{corollary}[Closing with the joint words]\label{cor:retainedclosing}
For a rational ray with simple nonresonant crossings, compute the finitely many arithmetic cohorts of crossings and one cohort of endpoints off the poles in each repeating gap. If each crossing cohort and each required endpoint cohort has a certified upper root rate $r_Gc_{\max}<1$, the crossing terms are absolutely summable and their ordered sum is the initial fibre integral. The finite head is retained, and the bounds of Corollary~\ref{cor:retainedheight} give an effective stopping rule. A cohort with the positive hypothesis of Corollary~\ref{cor:retainedheight}, or the noncancelling simple tops of Theorem~\ref{thm:retainedtops}, has the stated exact root rate and cannot tend to zero when it exceeds one.
\end{corollary}

The procedure keeps the falling rates, which fix the crossed heads, and the cut source; within the chamber the growing shares can still be changed by Theorem~\ref{thm:chamber}, the completed word being the same. The bound $r_Gc_{\max}$ has no phase term: $c_{\max}$ is the ceiling of this joint word, which need not equal the ceiling of a keyhole continuation read before the fibre is completed, nor the least ceiling over all joint words.

% Reference 20: lem:positivecapacity; source line 8587
\begin{lemma}[Positive capacity and the growing shares]\label{lem:positivecapacity}
Let $t_g\in\mathbb R^r$, $g\in G$, be a finite nonempty configuration with a grading $\ip{t_g}{\mathbf v}=1$, let $\xi_g>0$, $w>0$ and $\ip{\mathbf v}\nu=w$, and in this statement write $L(X)=\sum_{g\in G}\xi_gX^{t_g}$ and $\Pi(\nu,w)=\{\mathbf b\in[0,\infty)^G:\ \sum_gb_g=w,\ \sum_gb_gt_g=\nu\}$. With $0^0=1$ and the infimum of the empty set equal to $+\infty$,
\[
\sup_{X\in(0,\infty)^r}\frac{X^\nu}{L(X)^w}=\inf_{\mathbf b\in\Pi(\nu,w)}\prod_{g\in G}\Big(\frac{b_g}{w\,\xi_g}\Big)^{b_g},
\]
the left side being the \emph{capacity} of $L$ at $(\nu,w)$.
The supremum may equally be taken on any logarithmic section transverse to $\mathbf v$. When $\nu/w$ lies in the relative interior of the convex hull of the $t_g$, the infimum has a unique positive minimizer and is attained by a gauge $X$, unique modulo the directions constant on all the $t_g$; on a proper face the same computation on that face gives the value, approached by suppressing the other monomials.
\end{lemma}

The supremum on the left is the largest value of a monomial of the letter vector, and the logarithm of the infimum on the right is the least Gamma rate over the polytope $\Pi(\nu,w)$ of shares, less $w\log w$: the lemma is the duality that turns the chamber rate of a positive kernel into a capacity of its letter.

% Reference 20: thm:retainedchamberrate; source line 8650 (fixed rational pins).
% Proved in the last part: thm:retainedchamberrate.
\begin{theorem}[The exact positive rate of a retained chamber]\label{thm:retainedchamberrate}
Take a Gramma $\sG_{\nu;t}(\boldsymbol\xi)$ with rational configuration and real index read on its share side (Proposition~\ref{prop:shareform}), the pins $t\mathbf a=\nu$ with the grading of the letter totals and positive coefficients $\boldsymbol\xi$. Let a nonzero rational ray start at a positive real point of the pins, and consider any simple nonresonant Gamma crossing cohort of that ray, or any cohort of travelling fibres off the poles, with its integer step as in Definition~\ref{def:retainedcohort}. Write its direction per step as $\boldsymbol\sigma$, and put
\[
\mathcal S_-=\{l:\ \sigma_l<0\},\qquad d_l=-\sigma_l\ \ (l\in\mathcal S_-),\qquad\mathcal S_+=\mathcal S_-^c,
\]
$t_l$ the columns of $t$, with the chamber of $\boldsymbol\sigma$ (Definition~\ref{def:chamber}), the polar functions of the share side being the shares,
\[
C=\{\boldsymbol\sigma':\ t\boldsymbol\sigma'=0,\ \sigma'_l=-d_l\ (l\in\mathcal S_-),\ \sigma'_l\ge0\ (l\in\mathcal S_+)\}.
\]
Then the retained terms have a common sign at each index, independent of $\boldsymbol\xi$, after their crossing orientation is fixed, and
\[
\lim_{n\to\infty}|\mathcal R_n(\boldsymbol\xi)|^{1/n}=\exp\min_{\boldsymbol\sigma'\in C}h_\Gamma(\boldsymbol\sigma'),
\]
$h_\Gamma$ the Gamma rate of Theorem~\ref{thm:pastwalls}(iii), the coefficients being those of the Gramma: the exponential of the least Gamma rate of the chamber. This minimizes over the growing shares, the falling rates $d_l$ being fixed; if the integer step has the length $\Delta$ in the parameter of the ray, the logarithm is $\Delta$ times the rate per unit of that parameter.
\end{theorem}

% Reference 20: cor:retainedpooledrate; source line 8751
\begin{corollary}[Exact positive averaging at every retained dimension]\label{cor:retainedaveragedrate}
In Theorem~\ref{thm:retainedchamberrate} let the coefficients be measurable positive functions $\xi_l(\zeta)$ of a source with a nonzero positive measure $\lambda$, the configuration, the index, the chamber and the integer step fixed across the source, and put
\[
h_C(\zeta)=\min_{\boldsymbol\sigma'\in C}\Big(\sum_l\sigma'_l\log|\sigma'_l|-\sum_l\sigma'_l\log\xi_l(\zeta)\Big),\qquad\hat h_C=\operatorname*{ess\,sup}_\lambda h_C .
\]
Let one averaged retained term be absolutely integrable, so that after multiplication by $\prod_l\xi_l(\zeta)^{-a^0_l}$, $\mathbf a^0$ the real point at which the cohort starts (Definition~\ref{def:retainedcohort}), the joint positive input is finite. If $\hat h_C<\infty$, every later averaged term is integrable and
\[
\lim_{n\to\infty}\Big|\int\mathcal R_n(\zeta)\,d\lambda(\zeta)\Big|^{1/n}=e^{\hat h_C};
\]
if all these terms are finite and $\hat h_C=+\infty$, their root rate is $+\infty$. No estimate uniform over the source and no retained fibre of dimension one is needed. This applies to the kernel of Theorem~\ref{thm:averaging} on rational pins along a rational ray whose crossings are simple and nonresonant, under a finite averaged rate and the certificate \eqref{eq:avgcert}, whose bound supplies the finite first term. For a continuous coefficient map the essential supremum is the supremum on the support of the image of $\lambda$; it is $\operatorname{ess\,sup}_\zeta\min_C$, with no interchange to $\min_C\operatorname{ess\,sup}_\zeta$.
\end{corollary}

% Reference 20: thm:pool; source line 7704 (sharpness at fibre dimension one).
% Proved in the last part: thm:averagingsharp.
\begin{theorem}[Averaging is sharp]\label{thm:averagingsharp}
In the setting of Theorem~\ref{thm:averaging}, with its chamber $C$, its averaged rate $\hat h_C$, its source measure $\lambda$ and the crossing terms $\mathcal R_k$ of the closing of its kernel at the crossing times $\tau_k$, let the pins be rational, as in the kinds (A) and (B), let $\boldsymbol\sigma$ have rational coordinates, let the crossings of its ray be simple and nonresonant (Definition~\ref{def:retainedcohort}), and let $\hat h_C<\infty$. Along each cohort of crossings, with its integer step of length $\Delta$ in $\tau$, the averaged terms have one sign at each index and $\lim_{n\to\infty}\big|\int\mathcal R_n\,d\lambda\big|^{1/n}=e^{\Delta\hat h_C}$ (Theorem~\ref{thm:retainedchamberrate}, Corollary~\ref{cor:retainedaveragedrate}); so for every $\epsilon>0$, $\big|\int\mathcal R_k\,d\lambda\big|\ge e^{\tau_k(\hat h_C-\epsilon)}$ for all large $k$, and the terms of the averaged closing, which are those of the closing of the share side along $\boldsymbol\sigma$ when $\int e^{-\ip{\mathbf a}y}\,d\lambda$ converges along it, tend to zero when $\hat h_C<0$ and do not when $\hat h_C>0$. When every share that $\boldsymbol\sigma$ falls is alone in its bundle and the share side read with one share per bundle (Corollary~\ref{thm:E1} merging each bundle) has a fibre of dimension one, a positive multiple of $\boldsymbol\sigma$ has a direction with rational coordinates in its chamber, along which the closing is the same series after a change of scale of $\tau$ (Theorem~\ref{thm:chamber}); the crossings are simple and nonresonant as soon as no two lattices of the falling shares meet; and every crossing term is
\[
\mathcal R_k(\zeta)=r_k\prod_Ge^{-\bar a_{k,G}\,y_G(\zeta)},\qquad y_G=\log\sum_{l\in G}e^{y_l},
\]
with $r_k$ real and not depending on the source point and $\bar{\mathbf a}_k$ the crossing point of the merged shares.
\end{theorem}

So Theorem~\ref{thm:averagingsharp} holds in every retained dimension, with no condition on the bundles: the falling shares need not be alone in them, and the share side merged by its bundles may have any number of free shares. The chamber minimum is a finite convex problem, shared by all the cohorts of the ray once their common integer step is chosen, and the completed source average is one joint Trailmoment. The positivity belongs to the completed source average; a continued source factor, with its oriented circle integrals, is read by Theorem~\ref{thm:retainedtops}, its phases not being replaced by a supremum over a positive support. A completed crossing and the tail it inherits are different series: closing a retained share only through finitely many of its poles leaves an exact Taylor remainder whose rate depends on how many poles are taken with each crossing, and it can grow while the completed crossings fall.

\subsection{Coincident crossings and the boundary at root rate one}

A residue is taken after the factors on its hyperplane have been assembled. A deformation may separate those factors for computation, but its separate residues are then one packet. The finite operation of Theorem~\ref{thm:crossingjet} retains that packet, including its vanishing factors, and extends Theorem~\ref{thm:retainedchain} to coincident crossings.

% Reference 20: lem:sourcepolebridge; source line 8929
\begin{lemma}[A source pole with a uniform parameter margin]\label{lem:sourcepolebridge}
Use the radii of Lemma~\ref{lem:sourcestriplift}. Near $\kappa=-N$, $N\in\mathbb N$, put $z=\kappa+N$ and fix $0<\eta<1$. On $|\re z|\le\eta$ the same continued source factor is exactly
\[
I_\kappa=\int_r^1q^{\kappa-1}g(q)\,dq+R_{N+1}+R'_N+g_N\,\frac{r^z}z,\qquad g_N=\frac1{2\pi i}\oint g(\zeta)\,\zeta^{-N-1}\,d\zeta,
\]
$R_{N+1}$ and $R'_N$ given by the formulas of that lemma with these indices, $R'_0=0$. The first three terms are uniformly absolutely convergent on the whole strip; splitting $R'_N$ gives at most seven integrals in all, the circle integral $g_N$ carrying the explicit pole $r^z/z$. If $N=N_0+dn$, the domains stay fixed and their dependence on $n$ is again in fixed powers; in particular a small complex circle of the parameter about $z=0$ has a common integrable source representation after its displayed pole is removed.
\end{lemma}

% Reference 20: thm:crossingjet; source line 8959
\begin{theorem}[A finite jet for a coincident crossing]\label{thm:crossingjet}
Consider a rational affine family of crossing fibres, indexed by $n$, with $n$ in one fixed arithmetic progression. Fix the inherited contour and the whole packet of factors polar on its crossed hyperplane, let the other polar hyperplanes be separated from this packet on the real contour by a positive constant, and let the total order of the packet be at most $p$, independent of $n$. Perturb the constant terms of all shares and walls by one parameter $\varepsilon$, with a fixed rational vector of perturbations, keeping the affine pins and their prescribed continuation; the vector can be chosen outside finitely many rational hyperplanes so that distinct members of the packet separate, factors forced to remain coincident being kept as one member. Let $\mathcal R_{n,j}(\varepsilon)$ be the inherited crossing terms of the separated members. Then
\[
\mathcal R_n=[\varepsilon^0]\sum_j\mathcal R_{n,j}(\varepsilon).
\]
Each member has a pole of order at most $p-1$ in $\varepsilon$ if its separated pole is simple, and the same bound suffices after keeping forced multiple members intact; so the constant coefficient uses only the Taylor coefficients through the degree $p-1$ of the holomorphic parts after their poles are removed. Further meromorphic normalizing factors raise this order by their known pole orders. Every coefficient is a jet of the Gammas, of the source or of the parameter; no limit in $n$ is part of the computation. More generally, whenever a finite exact procedure returns terms $\mathcal R_{n,j}=\varepsilon^{-o_j}\widetilde{\mathcal R}_{n,j}$ with $\widetilde{\mathcal R}_{n,j}$ holomorphic, the required coefficients are exactly $[\varepsilon^{o_j}]\widetilde{\mathcal R}_{n,j}$; positive orders of a factor that vanishes at $\varepsilon=0$ are retained just as negative orders of a polar factor are. The local decay hypothesis of Lemma~\ref{lem:pathclose} is used on the common enclosing contour, uniformly on a small disc of the parameter. For the Gamma words and the continued sector integrals a uniform procedure is available: factors whose Gamma poles lie on the crossed hyperplane are constant on its residual fibre and are put in the scalar Gamma word, each source wall coincident on that hyperplane is treated by Lemma~\ref{lem:sourcepolebridge}, and the other walls and Gamma strips have strict margins.
\end{theorem}

The auxiliary parameter separates the poles of one assembled crossing; it does not identify the independent regulators of a product of words, which keep their separate constant terms (Proposition~\ref{prop:regulators}). If the value asked for is a finite part, its prescribed Mellin regulator is retained and its constant coefficient taken as in Theorem~\ref{thm:RegCT} and Corollary~\ref{cor:RegMellin}; finite parts in different coordinates are not interchangeable (Lemma~\ref{lem:Regscale}: a dilation of the coordinate moves a finite part by a polynomial in the logarithm of the dilation). Several such regulators give finitely many multivariate Laurent coefficients, with the coefficient order fixed by the value asked for, the retention of finitely many numbers at every corner (Theorem~\ref{thm:corner}) applied to a whole crossing fibre. The separation hypothesis concerns the actual residual contour: a transverse meeting of different hyperplanes needs the finite ordered packet of the chosen crossing flag, with its iterated transverse contours, and does not allow an arbitrary order of point residues. Nor does a Laurent coefficient keep a sharp root rate by itself: the coefficient can vanish.

% Reference 20: prop:topjets; source line 9050
\begin{proposition}[Bounds and top coefficients pass to jets]\label{prop:topjets}
In Theorem~\ref{thm:crossingjet} let the procedure of Theorem~\ref{thm:retainedchain} for the separated terms give a uniform bound on a circle of the parameter, $|\varepsilon^{o_j}\mathcal R_{n,j}(\varepsilon)|\le C_\delta(1+n)^{p'}(q+\delta)^n$ for every $\delta>0$, the circle allowed to shrink with $\delta$. Then $\limsup_n|\mathcal R_n|^{1/n}\le q$. For perturbations of the constant terms only, the slopes in $n$, the linear parts of the pins, the letters and the coefficients fixed, Lemma~\ref{lem:sourcepolebridge} and the extraction of the scalar give more: one common fan serves on a small disc of the parameter, each coefficient function $c_i$ is independent of the parameter, only the fixed amplitudes, the explicit polar factors and the constant terms of the Gammas varying, and their endpoint powers have positive margins once the explicit poles of Lemma~\ref{lem:sourcepolebridge} are removed; so the finitely many Laurent coefficients of each order have the same upper root rate as their common collection of coefficient functions and Gamma slopes, with only polynomial and logarithmic losses in $n$. If the joint word has finitely many simple tops (Definition~\ref{def:retainedtops}), uniformly separated from the rest of its domain, its finitely many Laurent coefficients supply every prescribed finite order of the expansion of $\mathcal R_n$ at the tops: differentiate the exact integrand to the required order in the parameter, expand its phase and amplitude to the required Taylor order at each top, use $\mathbf x=\mathbf X/n$ in its decaying boundary coordinates and $\mathbf y=\mathbf Y/\sqrt n$ in its coordinates $\mathbf y$, integrate the resulting monomials against the Gamma densities in $\mathbf X$ and the density $e^{-\frac12\mathbf Y^{\mathsf T}\mathbf Q\mathbf Y}$ in $\mathbf Y$, and expand the Gamma prefactors to the same order. Equal top phases are added before any test of nonvanishing. A surviving coefficient at the root rate $q$ proves $\limsup_n|\mathcal R_n|^{1/n}=q$; if the tested coefficients cancel, the upper bound remains and no sharp lower bound follows.
\end{proposition}

At the root rate one the boundary is decided exactly, without simple or isolated tops, by the law of the coefficient function that the procedure retains in the positive case; its geometry is computed once and answers several queries at the boundary.

% Reference 20: thm:unitrate; source line 9128
\begin{theorem}[The test at root rate one for a positive moment sequence]\label{thm:unitrate}
Let $\phi$ be a finite positive measure, $0\le c\le1$ on its support, $m_n=\int c^n\,d\phi$, and
\[
b_n=G_n\,m_n,\qquad G_n>0,\qquad G_n\sim C\,n^\beta,\quad C>0,\ \beta\in\mathbb R:
\]
the form of Corollary~\ref{cor:retainedheight} at the root rate one, the coefficient function divided by its positive largest value $c_{\max}$ and the scalar multiplied by $c_{\max}^{\,n}$, the exponent $\beta$ read from the explicit Gamma quotient and the scalar sign a fixed number of modulus one to the power $n$. Finitely many initial terms have no effect.
\begin{enumerate}[label=(\roman*)]
\item The series $\sum_nb_n$ is finite exactly when: $\beta<-1$, always; $\beta=-1$ and $\int\log\frac1{1-c}\,d\phi<\infty$; $\beta>-1$ and $\int(1-c)^{-\beta-1}\,d\phi<\infty$; the integrands of the last two conditions having the value $+\infty$ at $c=1$. These are also the exact conditions for absolute convergence after multiplication by any $\zeta^n$, $|\zeta|=1$.
\item Let moreover $G_n$ be a balanced Gamma quotient with its exponential factor removed, or more generally let its forward differences satisfy $\Delta^kG_n=O(n^{\beta-k})$ for every fixed $k$. For $|\zeta|=1$, $\zeta\neq1$, the ordered series $\sum_n\zeta^nb_n$ converges exactly when $b_n\to0$, and with $M(t)=\phi\{c>1-t\}$ this is: $\beta<0$, always; $\beta=0$ and $\phi\{c=1\}=0$; $\beta>0$ and $M(t)=o(t^\beta)$ as $t\downarrow0$.
\end{enumerate}
So (i) and (ii) tell absolute convergence, conditional convergence and the failure of the term test apart with one retained law.
\end{theorem}

A certified law near its largest value of the form $M(t)\asymp t^a(\log(1/t))^j$ turns these criteria into comparisons of powers and logarithms; the criteria stay exact for nonisolated or degenerate maxima, Dirac measures and mixtures. They do not change the crossing order: a periodic family must be reassembled in its original order before a conditional verdict is used, since splitting an alternating series into its even and odd subseries loses exactly that information.

% Reference 20: prop:boundaryjets; source line 9224
\begin{proposition}[Finite certificates at the boundary from the jets of the tops]\label{prop:boundaryjets}
Let the assembled terms at the root rate one have a certified finite expansion
\[
\mathcal R_n=\sum_{\zeta\in Z}\zeta^n\sum_{\beta\ge-1}\sum_{j=0}^{J_\beta}c_{\zeta,\beta,j}\,n^\beta(\log n)^j+O(n^{-1-\delta}),\qquad\delta>0,
\]
$Z$ a finite subset of the unit circle, the real exponents $\beta$ a finite set and equal triples combined; the jets of the tops of Proposition~\ref{prop:topjets} produce such a certificate to every order needed when its hypotheses on simple tops and uniformity hold. Then:
\begin{enumerate}[label=(\roman*)]
\item the series is absolutely convergent if and only if every displayed coefficient vanishes;
\item it converges in its given order if and only if every coefficient with $\zeta=1$ vanishes and every coefficient with $\beta\ge0$ vanishes;
\item when (ii) holds and (i) does not, the convergence is conditional.
\end{enumerate}
The verdict uses exact certificates of zero or nonzero for the finitely many assembled coefficients; where those are unavailable, the finitely many coefficients are the output, a numerical near cancellation not being a certificate of zero.
\end{proposition}


\chapter{Laws at their ends: regions of equality and reach}\label{ch:lawreach}

A closing past the walls reads a law where its moments diverge, at its falling ends, through its expansion there. The regions on which a law equals its expansion at an end are built from the sector letters, with the map from the motion of the shares to the exponents, the coefficients, the truncation bounds and exact boundary words (Theorems~\ref{thm:lawendchart} and~\ref{thm:lawregionatlas}); the domain where a series converges and the region where it equals the law are kept apart (Proposition~\ref{prop:lawreadingboundary}). The reach is read on the law along every rational direction of the shares, as at width one: along a rational direction the seam function is the Mellin transform of one measure on the half-line, the image of the law, multiplied by $\mathbf x^{-\mathbf w_0}$ for the starting point $\mathbf w_0$, under the directional ratio, and the reach is the far boundary of the region on which that measure is its own wall series (Theorem~\ref{thm:dirreach}). Run on the laws of the two families at a node of width one, the same argument reads the growth of each family on its law and makes that of the bunch its reach (Corollary~\ref{cor:growthreach}). For a piece the reach is bounded by the ceilings of the keyhole continuations along the direction, with no phase, and attained at their simple tops that do not cancel (Theorem~\ref{thm:dirheights}). A closing past the walls converges when the exponential of the Gamma rate times the reach of its tube, the imaginary translates of the line through its starting point, is less than one, its terms then having at most that root rate, and a law with log expansions bounds that reach in every direction with nothing paid for the phase; with one free share the tube is the line, and the terms reach that root rate under conditions that keep the poles apart (Theorems~\ref{thm:dirclosing} and~\ref{thm:onefreeshare}, Corollary~\ref{cor:anyfibre}).

\section{Regions on which a law is represented}\label{sec:lawregions}

A region of equality belongs to an expansion of a law: to the coordinates advanced, the coordinates retained, and the grouping of its terms. A disc in a sector variable is not by itself a region of equality of the law: the map to the seam torus can identify points, have critical values, or have an image of smaller dimension than its torus. The procedure of Theorem~\ref{thm:lawregionatlas} keeps that map, and gives a law, with its tangential source retained, before any assertion about a density on the whole torus.

% Reference 20: def:lawendchart; source line 9278
\begin{definition}[A chart at an end of a law]\label{def:lawendchart}
Let a part of the source of a law have the coordinates $(q,z)\in(0,r_1)\times\cdots\times(0,r_k)\times Z$ and the form
\[
q^{\boldsymbol\kappa-\one}A(q,z)\,dq\,\phi(dz),\qquad\kappa_i>0,
\]
$\phi$ a finite positive measure retained as an input and $A$ holomorphic in $q$ on a neighbourhood of a closed polydisc, uniformly bounded there for $z\in Z$; complex densities are allowed, and a fixed monomial in the seam coordinates can be included in the integrand and removed after integration. The selected seam coordinates of the law at the end are $\check x_i^{-1}$, with
\[
\check x_i=\prod_{j=1}^kq_j^{N_{ij}}\,\mathfrak u_i(q,z),\qquad1\le i\le k,
\]
$\mathbf N=(N_{ij})$ an invertible integer matrix and the \emph{chart factors} $\mathfrak u_i$ positive on the real source and holomorphic and nonzero on the complex neighbourhood; the other coordinates of the seam point are a retained map $\mathbf C(q,z)$, and branches are continued from the positive source. The \emph{chart coordinates} are
\[
\mathfrak t_j=\prod_i\check x_i^{(\mathbf N^{-1})_{ji}}=q_j\prod_i\mathfrak u_i(q,z)^{(\mathbf N^{-1})_{ji}}.
\]
A certificate of the chart supplies positive radii $R_i$ and a uniform analytic inverse $q_j=\mathfrak t_j\,g_j(\mathfrak t,z)$ for $|\mathfrak t_i|<R_i$, with bounds for the inverse, its Jacobian and the holomorphic functions used in Theorem~\ref{thm:lawendchart} and Corollary~\ref{cor:lawfacejets} below; the cuts on $Z$ are retained. Write $\one_{\mathrm{cut}}(\mathfrak t,z)$ for the indicator of any further retained cut of the source, $\one_{\mathrm{cut}}=1$ when there is none; if a cut of the source moves in these coordinates, its indicator is retained in the integrand and is neither differentiated nor included in a holomorphic coefficient.
\end{definition}

% Reference 20: thm:lawendchart; source line 9312
\begin{theorem}[Computing and using an end chart]\label{thm:lawendchart}
For a chart of Definition~\ref{def:lawendchart} the following computations are effective.
\begin{enumerate}[label=(\roman*)]
\item A smaller certified polydisc exists whenever the holomorphy and nonvanishing conditions hold uniformly on a compact tangential source. In its coordinates the source form is
\[
\mathfrak t^{\boldsymbol\kappa-\one}\widetilde A(\mathfrak t,z)\,d\mathfrak t\,\phi(dz),\qquad\widetilde A=A(\mathfrak tg,z)\prod_jg_j^{\kappa_j-1}\det\frac{\partial(\mathfrak tg)}{\partial\mathfrak t},
\]
the determinant positive on a small enough positive chart. The seam map is $(\mathfrak t,z)\mapsto\big(\mathfrak t^{\mathbf N},\mathbf C(\mathfrak tg(\mathfrak t,z),z)\big)$; with $\mathbf C_{\mathfrak t}(z)=\mathbf C(\mathfrak tg(\mathfrak t,z),z)$, its density in the normal coordinates $\check x$, a measure in the tangential ones, is
\[
\frac{\check x^{-\one}\mathfrak t^{\boldsymbol\kappa}}{|\det\mathbf N|}\,(\mathbf C_{\mathfrak t})_\ast\big(\one_{\mathrm{cut}}(\mathfrak t,\cdot)\widetilde A(\mathfrak t,\cdot)\,\phi\big)\,d\check x,
\]
a density only in the selected normal coordinates, whose tangential value may be a singular measure.
\item Write $\widetilde A=\sum_{\mathbf n\ge0}\widetilde A_{\mathbf n}(z)\,\mathfrak t^{\mathbf n}$. The coefficients are tangential words computed by finite jet arithmetic on the original letters; the inverse certifies that they represent the source and is not evaluated. If $|\widetilde A|\le M$ on $|\mathfrak t_i|=R_i$, then $|\widetilde A_{\mathbf n}(z)|\le M\prod_iR_i^{-n_i}$. For every bounded measurable test function $H$, the law on the chart is, for real $0<\mathfrak t_i<R_i$,
\begin{equation}\label{eq:retainedwallseries}
\mathfrak t^{\boldsymbol\kappa-\one}\sum_{\mathbf n\ge0}\mathfrak t^{\mathbf n}\int_Z\widetilde A_{\mathbf n}(z)\,\one_{\mathrm{cut}}(\mathfrak t,z)\,H\big(\mathfrak t^{\mathbf N},\mathbf C(\mathfrak tg(\mathfrak t,z),z)\big)\,\phi(dz)\,d\mathfrak t .
\end{equation}
The same seam map and the same retained cut occur in every term. With rational data the inverse branch is algebraic, specified by the positive branch and the certificate; alternatively every exact pairing with a test function is kept on the original source $(q,z)$ with that cut pulled back, and an inverse function need not be added to the operations. The series converges in total variation on compact subcharts, the normal power included. If the tangential map and the cut do not depend on $\mathfrak t$, its coefficients are the fixed face laws $\mathbf C_\ast(\widetilde A_{\mathbf n}\phi)$; otherwise \eqref{eq:retainedwallseries} is a series of pairings on the one tangential source, not an assertion that those images are fixed.
\item With $\vartheta_i=|\mathfrak t_i|/R_i<1$, truncating to $0\le n_i<k_i$ has, before the normal power and for $\|H\|_\infty\le1$, the error at most
\begin{equation}\label{eq:lawchartremainder}
M\,\phi(Z)\prod_i(1-\vartheta_i)^{-1}\sum_i\vartheta_i^{k_i}.
\end{equation}
It also has an exact boundary word: with $T_i$ the Taylor operator through the degree $k_i-1$ in $\mathfrak t_i$,
\begin{gather}
\widetilde A=\prod_iT_i\widetilde A+\sum_{\varnothing\neq I\subseteq\{1,\dots,k\}}\prod_{i\in I}(1-T_i)\prod_{j\notin I}T_j\widetilde A,\label{eq:lawboundaryword}\\
(1-T_i)\widetilde A=\frac{\mathfrak t_i^{k_i}}{2\pi i}\oint_{|\zeta_i|=R_i}\frac{\widetilde A(\mathfrak t_1,\dots,\zeta_i,\dots,\mathfrak t_k,z)}{\zeta_i^{k_i}(\zeta_i-\mathfrak t_i)}\,d\zeta_i:\notag
\end{gather}
there are $2^k-1$ boundary words, whatever the truncation orders $k_i$.
\item The chart records the map from the motion of the shares to the exponents. The following statements about walls and continuation concern normal Mellin powers with an analytic tangential composite, or a further word with its own sector certificate; the bounded measurable test functions of (ii) have only the convergent identity and the error bound, and no continuation is asserted for them. For the Mellin power in the seam coordinates $\check x_i^{-1}$ at the shares $\mathbf z$, the normal exponent of a term is $\boldsymbol\kappa+\mathbf N^{\mathsf T}\mathbf z+\mathbf n-\one$. For a fixed tangential cut and a normal Mellin power, integrating these powers gives the walls $\kappa_j+(\mathbf N^{\mathsf T}\mathbf z)_j+n_j=0$; a moving cut keeps its indicator in the integrand, with no such set of poles. A motion $\boldsymbol\sigma$ of the shares in the coordinates of these Mellin powers has the wall slopes $\mathbf N^{\mathsf T}\boldsymbol\sigma$; for the original coordinates of the shares one first applies to $\boldsymbol\sigma$ the recorded linear change of the exponents of the Mellin powers. Other requested coordinates remain in the tangential word, and are not replaced by their values on the face unless that replacement is part of the stated test function.
\end{enumerate}
\end{theorem}

% Reference 20: cor:lawfacejets; source line 9456
\begin{corollary}[Face jets for an analytic test function]\label{cor:lawfacejets}
Let the tangential cut be fixed, the map $\mathbf C(\mathfrak tg(\mathfrak t,z),z)$ holomorphic on a specified polydisc about $\mathfrak t=0$, that point included, and $H$ a bounded holomorphic tangential test function on a neighbourhood of its image. Expand $\widetilde A(\mathfrak t,z)\,H\big(\mathbf C(\mathfrak tg(\mathfrak t,z),z)\big)$ instead of $\widetilde A$ alone: its coefficient of degree $\mathbf n$ is a finite sum of tangential source integrals of derivatives of $H$ at $\mathbf C(0,z)$, multiplied by jets of $\widetilde A$, $\mathbf C$ and $g$ of degrees at most $\mathbf n$. These are fixed face functionals, and their series and boundary words obey the same coefficient bounds as in Theorem~\ref{thm:lawendchart}(ii), with the bound of this composite. A change of analytic tangential test function needs its jets, not a new continuation of the source. The normal monomial $\mathfrak t^{\mathbf N}$ remains unexpanded in the integrand; a normal test function may be expanded as well when its composite is holomorphic and bounded on the polydisc and its jets are supplied by finite operations, a further condition, since negative entries of $\mathbf N$ can give poles at $\mathfrak t=0$.
\end{corollary}

% Reference 20: thm:lawregionatlas; source line 9487
\begin{theorem}[Certified regions of equality and composition]\label{thm:lawregionatlas}
For rational positive polynomial letters and rational monomial seam coordinates the following is a procedure.
\begin{enumerate}[label=(\roman*)]
\item The finite sector data determine the exponent matrices and the chart factors of Definition~\ref{def:lawendchart}. Every invertible normal minor gives a proposed end chart, with the other variables and test functions retained; enumerating rational tangential boxes and radii and verifying the inequalities of the certificate and of nonvanishing produces certified charts. Every compact tangential set on the end face $q=0$ on which the proposed chart has nonvanishing chart factors and an invertible normal matrix has a small enough normal neighbourhood covered by finitely many charts produced in this way. Extending a chart away from that face needs its certificate, or a new certificate of inverse functions for the full Jacobian; invertibility of $\mathbf N$ alone does not supply it.
\item The region of equality supplied by one chart is the image of its certified polydisc in $\mathfrak t$ under $\check x=\mathfrak t^{\mathbf N}$, with the specified branches and cuts; its positive part satisfies $\sum_i(\mathbf N^{-1})_{ji}\log\check x_i<\log R_j$, $1\le j\le k$, and need not be a rectangle in the original seam coordinates. A finite atlas is certified for a requested region of equality only after its source patches cover the inverse image of that region; a patch not yet covered is retained as its original source word. So at every step the output is an exact finite head together with the boundary words \eqref{eq:lawboundaryword} and the uncovered source word, with bounds for the parts certified small; an uncovered word is never omitted.
\item Atlases may be restricted to common source patches, added, transported by their seam maps, and composed with a further test function. On compact certified subregions every order of the absolutely convergent reads has the same value. At finite truncations the equality uses the boundary words \eqref{eq:lawboundaryword}, and moving a cut retains the source between the old and the new cut. For words whose sector certificate supplies a meromorphic continuation, the same finite identities hold there, the full Laurent coefficient taken before pairing with a test function at coincident poles; an arbitrary bounded test function is not thereby continued. An infinite read in another grouping is asserted only with its own bound of convergence or as the continuation of the retained word.
\end{enumerate}
\end{theorem}

The retained data are the sector syntax, the matrix $\mathbf N$, the cuts, the recursion for the inverse jets and its certificate, and the tangential coefficient words, computed once per distinct source chart. For a prescribed error on $|\mathfrak t_i|\le\vartheta_iR_i$ the orders $k_i$ are read off \eqref{eq:lawchartremainder}; the finite head has $\prod_ik_i$ coefficient words and its exact remainder $2^k-1$ boundary words. A repeated chart reuses its coefficient words, a new analytic test function supplies its finite face jets, and the certificate and the retained boundary remain valid when only the truncation changes.

% Reference 20: prop:lawdensityregion; source line 9582
\begin{proposition}[When a retained chart carries a density]\label{prop:lawdensityregion}
A retained chart is in general not a chart of an ambient density. It becomes one in the selected normal coordinates and further tangential coordinates $\mathbf y$ when their image has a certified local product parametrization over $(\mathfrak t,\mathbf y)$, with compact fibres, fixed cuts, a holomorphic form and an integrable uniform bound; integration over those fibres then gives a holomorphic coefficient function, and the series of Theorem~\ref{thm:lawendchart}, or of Corollary~\ref{cor:lawfacejets}, represents that density with the same bounds. For a source form made of integrable powers of the boundary coordinates of a cube times a holomorphic density, such parametrizations are built near every compact regular fibre of a rational map $\mathbf C$ on the cube, when the map is defined on a neighbourhood of the closed cube, its restriction to every face meeting the fibre has a surjective differential, and its letters are nonzero. To assert a density for the whole law, the target patch must have no other or uncovered preimages, and in particular no omitted pole face may supply further source patches; otherwise the procedure describes the covered part of the law only, and the further preimages remain uncovered source words under Theorem~\ref{thm:lawregionatlas}. With the cube coordinates $v_i$ put
\[
\mathbf J_{ji}=v_i(1-v_i)\,\partial_{v_i}C_j,\qquad X_j=\sum_iv_i(1-v_i)\big[\mathbf J^{\mathsf T}(\mathbf J\mathbf J^{\mathsf T})^{-1}\big]_{ij}\,\partial_{v_i};
\]
then $D\mathbf C(X_j)=\mathbf e_j$. On a small enough neighbourhood the flows of these vector fields carry the reference fibre to the neighbouring fibres, keep all faces and give the product parametrization. Uniform nonvanishing of $\det(\mathbf J\mathbf J^{\mathsf T})$, of the letters and of the relevant Jacobians, together with a bound for the vector fields on a larger neighbourhood, is a finite certificate; the iteration of the integral equations of the flows, certified to converge geometrically, computes the local flows and their holomorphic extensions.
\end{proposition}

The possible failures of this procedure are explicit: zeros of letters, loss of properness at an omitted face of the source, or loss of surjectivity of the differential of the seam map on a face. Their images are candidate boundaries, not automatically singularities of the law, since different contributions can cancel. If the source has smaller dimension than the number of requested seam coordinates, the retained law and its atlas still exist but no ambient density is asserted. A genuine edge of the support is an obstruction: if the law vanishes on an open set on one side of an edge and has positive measure arbitrarily close on the other, it has no analytic density across that edge, by the identity theorem, as in Theorem~\ref{thm:lawreach}(iii).

% Reference 20: prop:lawreadingboundary; source line 9655
\begin{proposition}[Convergence domains and equality regions]\label{prop:lawreadingboundary}
Let the fixed face laws of a certified chart, or the face functionals for a fixed bounded set of analytic test functions, be $\widetilde A_{\mathbf n}$ in a Banach space. The interior $\Omega_{\mathrm{abs}}$ of
\[
\Big\{\mathfrak t:\ \sum_{\mathbf n\ge0}\|\widetilde A_{\mathbf n}\|\,|\mathfrak t^{\mathbf n}|<\infty\Big\}
\]
is the logarithmically convex \emph{domain of absolute convergence} of that series, which contains with each of its points $\mathfrak t$ every $\mathfrak t'$ with $|\mathfrak t'_i|\le|\mathfrak t_i|$ for all $i$, and every certified polydisc lies in it. Its \emph{region of equality} is retained separately: the chart patches covered by the source on which the series is proved to read the law, together with the patches reached by a chain of transports of charts that keep the equality. The series represents the law only on the intersection of the domain $\Omega_{\mathrm{abs}}$ and the region of equality; analytic continuation of the sum alone is not a transport of the equality across an edge of the support or a singular value of the law. For a positive integer vector $\mathbf d$ and a fixed point $\mathfrak t$ of the chart retain the grouping
\[
\widetilde A_{(N)}(\mathfrak t)=\sum_{\ip{\mathbf d}{\mathbf n}=N}\widetilde A_{\mathbf n}\,\mathfrak t^{\mathbf n};
\]
the grouped analytic germ in one variable $v$ on the curve $(\mathfrak t_iv^{d_i})_i$ has the exact radius $R$ with $R^{-1}=\limsup_N\|\widetilde A_{(N)}(\mathfrak t)\|^{1/N}$: inside it the grouped series converges, and outside it its grouped terms do not tend to zero. These statements concern the germ, and they become statements about the law only where its region of equality is certified as well; on the boundary of convergence the coefficients decide. Completing another read replaces these grouped terms by their images, whose radius can be larger, the equality with the new expansion of the law still using the transported source and cuts. Neither radius is a reach of the law by definition, nor is it replaced by the supremum of an unintegrated modulus of the source.
\end{proposition}

The comparison of radii for a bounded read coefficient by coefficient concerns the same normal coordinates and the same bank of coefficients. Merging a law and then choosing a new coordinate is another operation: the source is carried by that map and its expansion is regraded in the new coordinate, whose coefficients and grouped terms need not be bounded images of the old ones, and whose boundary of equality can come earlier. The old and the new expansion agree through their common source, with the cuts, the finite heads and the boundary words transported; their domains of convergence and atlases of equality are recorded in their own coordinates. Maximal analytic continuation, and deciding every cancellation at a singular boundary, are neither used by the procedure nor concluded from a ceiling. In one selected coordinate with a fixed tangential test function the procedure certifies local patches of equality with the wall series, and the global inversion of Theorem~\ref{thm:dirreach} (Theorem~\ref{thm:lawreach} at width one) identifies the reach of the law along that coordinate; in several coordinates the certified atlas of equality, the domain of convergence of the series and its retained grouping are stored separately.

\section{The reach along a direction}\label{sec:reachdirection}

At width one the reach of a bunch is read on the law of its ratio, as the point beyond which the law is its own wall series (Theorem~\ref{thm:lawreach}); it is bounded above by the ceilings of the keyhole continuations and attained at their simple tops that do not cancel (Theorems~\ref{thm:heightbound} and~\ref{thm:simpletop}), and against a room bounded below the terms of the node have the root rate of the reach over the least ratio of the room, which infinitely many of them reach (Theorem~\ref{thm:nodeintervals}). A closing past the walls at any width advances along one direction of the shares, and along a rational direction the seam function is the Mellin transform of one measure on the half-line, the directional law (Definition~\ref{def:dirreach}). The theory of width one applies to that measure line by line, and its counterparts are stated along a direction.

% Reference 20: the properties (W1) and (W2); source line 9738
\begin{definition}[Walls and growth bounds]\label{def:wallsgrowth}
Let a piece have its seam function $P_{W_B(\mathbf w)}(\mvec_B)$ on the pins of the seam, an affine space $A_B$ of real shares $\mathbf w=(\mathbf a,\mathbf b)$ of the active letters with the space of directions $\vec A_B$, of dimension $n_{\mathcal D}$; for $\mathbf w\in A_B+i\vec A_B$ it is $\int\mathbf x^{-\mathbf w}\,d\big((\boldsymbol\varrho^\rho_B)_\ast\phi^\rho_B\big)(\mathbf x)$ where this converges, the law being carried by the closure of the image of the open simplex under the seam vector. The seam function has \emph{walls} when
\begin{enumerate}[label=(W\arabic*)]
\item there are finitely many affine functions $\kappa_U$ on $A_B$ with rational linear parts, its \emph{walls}, such that it converges at a real point exactly when every wall is positive there, which happens on an open convex set $\Omega_B\subseteq A_B$, and it continues meromorphically to $A_B+i\vec A_B$ with its poles on the hyperplanes on which a wall lies in $-\mathbb N$;
\end{enumerate}
and it has \emph{growth bounds} when
\begin{enumerate}[label=(W\arabic*),resume]
\item for every $\theta>0$ there are $M$, $C_0$, $n_1$ and $n_2$ with $|P_{W_B(\mathbf z)}(\mvec_B)|\le C_0\,\delta^{-n_2}(1+|\mathbf z|)^{n_1}M^{|\re\mathbf z|_1}e^{\theta|\im\mathbf z|_1}$, $\delta$ the smaller of $1$ and the distance of $\mathbf z$ from these hyperplanes.
\end{enumerate}
The same definitions apply to the period of any word whose exponents are affine functions of the shares.
\end{definition}

Every piece whose seam function converges at a real point has walls and growth bounds, and so does the seam function of its seam letters merged by bundles: the seam word of a piece is a word in positive letters, and so is the word of the merged seam letters of Theorem~\ref{thm:averaging}, sums of seam letters of several stations of one excess, which are no longer letters of tables; such a word continues past its walls as the seam word of every piece does (Theorem~\ref{thm:pastwalls}(ii)).

% Reference 20: lem:newtonwalls; source line 9745
\begin{lemma}[A word in positive letters continues past its walls]\label{lem:newtonwalls}
\begin{enumerate}[label=(\roman*)]
\item Let $W=\prod_lL_l^{-w_l}$ be a product of powers of positive letters on a set $B$ (Definition~\ref{def:posword}), whose exponents $w_l$ and masses $\mvec_B$ are affine functions of $\mathbf z\in\mathbb C^p$ with rational linear parts, the masses on the balance hyperplane of $W$ for every $\mathbf z$, and let its period converge at a real point $\mathbf z_0$, the exponents there being of any sign. Then there are finitely many affine functions $\kappa_U$ on $\mathbb C^p$ with rational linear parts, its walls, such that the period converges at a real point exactly when every $\kappa_U$ is positive there, and $P_W(\mathbf z)$ continues meromorphically to $\mathbb C^p$, with its poles, of order at most $|B|-1$, on the hyperplanes $\kappa_U(\mathbf z)\in-\mathbb N$; and for every $\theta>0$ there are $M$, $C_0$ and $n_1$ with
\[
|P_W(\mathbf z)|\le C_0\,\delta^{-(|B|-1)}(1+|\mathbf z|)^{n_1}M^{|\re\mathbf z|_1}e^{\theta|\im\mathbf z|_1},
\]
$\delta$ the smaller of $1$ and the distance of $\mathbf z$ from these hyperplanes.
\item For a piece and a partner read in its monomials, as in Theorem~\ref{thm:averaging}(B), merge the shares of each bundle $G$ of Theorem~\ref{thm:averaging} into $\bar a_G$, and put $L_G=\sum_{l\in G}\xi_l\,V_{jD(l)}$ for a bundle of shares of $K_j$, and $L_G$ alike with the seam letters $V^\ast_{jD}$ for a bundle of shares of $\Psi_j$. On the pins of the merged shares the \emph{merged seam word} $W^{\mathrm m}_B(\bar{\mathbf a})=\prod_GL_G^{-\bar a_G}$ is such a product at the masses $\mvec_B$, and its period, the \emph{merged seam function},
\[
P_{W^{\mathrm m}_B(\bar{\mathbf a})}(\mvec_B)=\int\prod_Ge^{-\bar a_G\,y_G}\,d\lambda,\qquad y_G=\log\sum_{l\in G}e^{y_l},
\]
reads the merged values $y_G$ at the source points, against the source measure $\lambda$ (Theorem~\ref{thm:averagingsharp}); it converges at the sums over the bundles of every real point with positive coordinates at whose sums over the monomials the seam function of the piece converges.
\end{enumerate}
\end{lemma}

For the seam word of a piece, at the masses $\mvec_B$ and with the shares as its exponents, Lemma~\ref{lem:newtonwalls}(i) gives walls and growth bounds whenever the seam function converges at a real point, the functions $\kappa_U$ of the lemma serving as the walls, with $n_2=|B|-1$; when the shares are positive at that point the seam word has a Newton sum of full dimension (Theorem~\ref{thm:poswordconv}), and the walls of the generators of a fan of the seam word serve, with the sector integrals of Theorem~\ref{thm:pastwalls}(ii), the proper walls for a piece of tables. A piece whose seam function on the pins is a product of Beta corners $\Gst21{x,\ y}{x+y}$ of affine functions of the shares with rational linear parts and of real exponentials of the shares, converging exactly where those affine functions are positive, has walls, those functions, and growth bounds, by the reflection formula (Proposition~\ref{prop:stacked}(v)) and the asymptotic expansion of $\log\Gamma$ away from the negative real axis, and growth bounds even at $\theta=0$ when for each Beta corner the linear parts of its two arguments on $\vec A_B$ are proportional, as when one argument does not depend on the shares.

% Reference 20: def:dirreach; source line 9770
\begin{definition}[Directional ratios and laws; the reach along a direction]\label{def:dirreach}
Let a piece have walls and growth bounds (Definition~\ref{def:wallsgrowth}), and let $\bar{\boldsymbol\sigma}\in\vec A_B\setminus\{0\}$ have rational coordinates. Its \emph{directional ratio} is the positive function $\mathbf x^{-\bar{\boldsymbol\sigma}}$ on the seam torus. For $\mathbf w_0\in\Omega_B$ and $\bar{\boldsymbol\vartheta}\in\vec A_B$ the \emph{directional law} at $\mathbf w_0+i\bar{\boldsymbol\vartheta}$ is the image
\[
\lambda_{\bar{\boldsymbol\sigma},\mathbf w_0+i\bar{\boldsymbol\vartheta}}=\big(\mathbf x^{-\bar{\boldsymbol\sigma}}\big)_\ast\Big(\mathbf x^{-\mathbf w_0-i\bar{\boldsymbol\vartheta}}\,(\boldsymbol\varrho^\rho_B)_\ast\phi^\rho_B\Big),
\]
a complex measure on $(0,\infty)$ whose variation is at most the positive directional law $\lambda_{\bar{\boldsymbol\sigma},\mathbf w_0}$, which integrates to the seam function at $\mathbf w_0$; it is \emph{oscillating} when $\bar{\boldsymbol\vartheta}\neq0$, and the line $\{\mathbf w_0+i\bar{\boldsymbol\vartheta}+z\bar{\boldsymbol\sigma}:z\in\mathbb R\}$ is the \emph{imaginary translate} by $\bar{\boldsymbol\vartheta}$ of the line through $\mathbf w_0$. Its \emph{line function} is
\[
\ell(z)=P_{W_B(\mathbf w_0+i\bar{\boldsymbol\vartheta}+z\bar{\boldsymbol\sigma})}(\mvec_B)=\int_0^\infty X^z\,d\lambda_{\bar{\boldsymbol\sigma},\mathbf w_0+i\bar{\boldsymbol\vartheta}}(X),
\]
absolutely convergent for $\re z$ in the open interval of the $x$ with $\mathbf w_0+x\bar{\boldsymbol\sigma}\in\Omega_B$, which contains $0$, and continued meromorphically by (W1). A wall falls along $\bar{\boldsymbol\sigma}$ when its falling rate $-\dot\kappa_U$, $\dot\kappa_U=\partial_{\bar{\boldsymbol\sigma}}\kappa_U$, is positive (Definition~\ref{def:chamber}). The poles of $\ell$ with $\re z\ge0$ lie among the \emph{candidate poles}
\[
z=\big(\kappa_U(\mathbf w_0)+n+i\,\partial_{\bar{\boldsymbol\vartheta}}\kappa_U\big)/(-\dot\kappa_U),\qquad n\in\mathbb N,
\]
of the falling walls: on finitely many horizontal lines, with real parts in finitely many cosets of one lattice $N_0^{-1}\mathbb Z$, the falling rates being rational. Let $\Sigma$ be the half-strip of the $z$ with $\re z\ge0$ and $\im z$ within $1$ of the interval spanned by $0$ and these lines. The \emph{reach along} $\bar{\boldsymbol\sigma}$ at $\mathbf w_0+i\bar{\boldsymbol\vartheta}$ is
\[
r(\bar{\boldsymbol\sigma};\mathbf w_0+i\bar{\boldsymbol\vartheta})=\inf\big\{M\ge0:\ \text{for every }\delta>0\text{ there are }C,n'\text{ with }|\ell(z)|\le C(1+|z|)^{n'}M^{\re z}\big\},
\]
the bound asked at every $z\in\Sigma$ at distance at least $\delta$ from the candidate poles; it is finite by (W2), and by the maximum modulus principle on circles about the candidates it does not change when the actual poles are used. A \emph{wall series} is as in Definition~\ref{def:lawbunch}, the $\eta_i$ allowed to be complex and its condition read as $\sum_{\eta,j}|a_{\eta,j}|x'^{-\re\eta}<\infty$ for every $x'>x_0$, $x_0$ its radius. At $\bar{\boldsymbol\vartheta}=0$ the directional law is the \emph{law along} $\bar{\boldsymbol\sigma}$. The same definitions apply to the merged seam function of Lemma~\ref{lem:newtonwalls}(ii) on the pins of the merged shares, the merged values $e^{y_G}$ in place of the seam coordinates and the source measure in place of the law.
\end{definition}

The falling end of the directional law is $X\to\infty$, where the walls that fall along $\bar{\boldsymbol\sigma}$ make its moments diverge. In an end chart of Theorem~\ref{thm:lawendchart} the Mellin power at the shares $\mathbf z$ is $\mathfrak t^{\mathbf N^{\mathsf T}\mathbf z}$, so the directional ratio is $\mathfrak t^{-\mathbf d}$, $\mathbf d=-\mathbf N^{\mathsf T}\bar{\boldsymbol\sigma}$ the falling rates of the walls of the chart, times a power of the tangential coordinates: the wall series of the directional law groups the terms of the chart by their exponents $(\kappa_j+n_j)/d_j$ along the direction, and the reach is the far boundary, in $X$, of the region on which the law equals that expansion. At width one with one table and factoring weights, normalized as in Proposition~\ref{prop:widthonelaw}, the seam torus is a curve, $\vec A_B$ is the line of $\bar{\boldsymbol\sigma}=(-1,1,-1,1)$ in the shares $(a_{D_0},a_{D_1},b_{D_0},b_{D_1})$, and the directional ratio is $x^\ast_{D_0}/x_{D_1}$, which is $\varrho_B$ at the seam points (Proposition~\ref{prop:widthonelaw}); every coordinate of the torus is a power of $\varrho_B$, so the directional law is a power of $\varrho_B$ times the law of Definition~\ref{def:lawbunch}, and by Theorems~\ref{thm:dirreach} and~\ref{thm:lawreach} the reach along $\bar{\boldsymbol\sigma}$ is the reach $\operatorname{reach}(B)$ of the bunch.

% Reference 20: thm:dirreach; source line 9792 (items (i) to (iv)).
% Proved in the last part: thm:dirreach.
\begin{theorem}[The reach along a direction is read on the law]\label{thm:dirreach}
Let $\bar{\boldsymbol\sigma}$, $\mathbf w_0$, $\bar{\boldsymbol\vartheta}$, the directional law $\lambda=\lambda_{\bar{\boldsymbol\sigma},\mathbf w_0+i\bar{\boldsymbol\vartheta}}$, its line function $\ell$ and its reach $r$ be as in Definition~\ref{def:dirreach}.
\begin{enumerate}[label=(\roman*)]
\item On $(r,\infty)$ the measure $\lambda$ has a density which is a wall series of radius $r$. Its terms are the poles of $\ell$ with $\re z>0$: a pole $\eta$ with the principal part $\sum_j(-1)^{j+1}j!\,a_{\eta,j}(z-\eta)^{-j-1}$ gives the terms $a_{\eta,j}X^{-\eta-1}(\log X)^j$, and there are no others.
\item If on some $(x_0,\infty)$ the measure $\lambda$ has a density which is a wall series of radius $x_0$, then $r\le x_0$; and for every $x'>x_0$ there is $C$ with $|\ell(z)|\le C\delta^{-J-1}x'^{\,\re z}$ at every $z$ with $\re z\ge0$ at distance at least $\delta\le1$ from the poles, whatever $\im z$, $J$ the largest power of the logarithm in the series.
\item So $r$ is the least $x_0\ge0$ for which $\lambda$ has on $(x_0,\infty)$ a density which is a wall series of radius $x_0$. It does not change when $\mathbf w_0$ moves along $\bar{\boldsymbol\sigma}$, $r(c\bar{\boldsymbol\sigma};\cdot)=r(\bar{\boldsymbol\sigma};\cdot)^c$ for rational $c>0$, and $r$ is at least every point near which $\lambda$ has no real-analytic density. At $\bar{\boldsymbol\vartheta}=0$ the law $\lambda$ is positive and carried by the closure of the values of $\mathbf x^{-\bar{\boldsymbol\sigma}}$ on the support of the law of the piece: $r\ge X_{\min}$, the least of these values, when it is positive, the least value of the directional ratio taking the place of the least ratio (Proposition~\ref{prop:reachfloor}); and no wall falls along $\bar{\boldsymbol\sigma}$ exactly when $\mathbf x^{-\bar{\boldsymbol\sigma}}$ is bounded on that support, and then $r=X_{\max}$, the largest. When the seam word has a Newton sum of full dimension, as for a piece of tables whose tables form a connected family (Theorem~\ref{thm:pastwalls}(ii)), $X_{\min}=(\sup\mathbf V_B^{\bar{\boldsymbol\sigma}})^{-1}$ and $X_{\max}=\sup\mathbf V_B^{-\bar{\boldsymbol\sigma}}$, the suprema of the monomials of the letter vector of Theorem~\ref{thm:closecone} (Definition~\ref{def:readingvectors}).
\item At $\bar{\boldsymbol\vartheta}=0$ let $\mathbb L=\{z_0+kz_1:\ k\in\mathbb N\}$, $z_1>0$ rational, be a progression at positive distance from the poles of $\ell$, and put $r_{\mathrm{res}}=\limsup_\eta\max_j|a_{\eta,j}|^{1/\eta}$ over the poles and $r_{\mathbb L}=\limsup_k|\ell(z_0+kz_1)|^{1/(z_0+kz_1)}$. Then $r=\max(r_{\mathrm{res}},r_{\mathbb L})$: the reach is carried by the residues of the line or by its values on every progression.
\item For every $\theta\ge0$, $r$ is the infimum of the $M\ge0$ for which, for every $\delta>0$, there are $C$ and $n'$ with $|\ell(z)|\le C(1+|z|)^{n'}M^{\re z}e^{\theta|\im z|}$ at every $z$ with $\re z\ge0$ at distance at least $\delta$ from the candidate poles: the growth of the line function on its whole half-plane, at any allowance in the imaginary direction, is the reach.
\end{enumerate}
\end{theorem}

At width one the same argument, run on the law of a family at a node, reads the growth of the family on its law, and that of the bunch is its reach.

% Reference 20: cor:growthreach; source line 9831
\begin{corollary}[At width one the growth of a family is read on its law]\label{cor:growthreach}
In Section~\ref{sec:onescale} let $\lambda_B$ be the image of the base form $\phi_B$ under $\varrho_B$ and $\lambda_R$ the image of $\phi_R$ under $1/\varrho_R$. For $X\in\{B,R\}$ the growth $\operatorname{growth}(X)$ along the level $d$ (Definition~\ref{def:growth}) is the least $x_0\ge0$ for which $\lambda_X$ has on $(x_0,\infty)$ a density that is a wall series of radius $x_0$. In particular $\operatorname{growth}(B)=\operatorname{reach}(B)$: the inequality $\operatorname{reach}(B)\le\operatorname{growth}(B)$ of Theorem~\ref{thm:nodeclosing}(i) is an equality, and Theorem~\ref{thm:nodeclosing}(iii) closes a node on the lattices of both families when the two points beyond which their laws are their own wall series multiply to less than one.
\end{corollary}

A closing of the share side moves the imaginary parts of the shares over a whole fibre, so its bounds read the seam function on every imaginary translate of a line as well as on the line itself.

% Reference 20: cor:dirtube; source line 9840
\begin{corollary}[The reach of a tube]\label{cor:dirtube}
Let $\bar{\boldsymbol\sigma}$ and $\mathbf w_0$ be as in Definition~\ref{def:dirreach} and $V'\subseteq\vec A_B$ a subspace containing $\bar{\boldsymbol\sigma}$. For $\theta\ge0$ let $g_{V'}(\bar{\boldsymbol\sigma};\theta)$ be the infimum of the $M$ for which there are $C$, $n_1$ and $n_2$ with
\[
\big|P_{W_B(\mathbf w_0+\tau\bar{\boldsymbol\sigma}+i\bar{\boldsymbol\vartheta})}(\mvec_B)\big|\le C\,\delta^{-n_2}(1+\tau+|\bar{\boldsymbol\vartheta}|)^{n_1}M^\tau e^{\theta|\bar{\boldsymbol\vartheta}|_1}\qquad(\tau\ge0,\ \bar{\boldsymbol\vartheta}\in V'),
\]
$\delta$ the smaller of $1$ and the distance from the hyperplanes of (W1), and let $r^\sharp_{V'}(\bar{\boldsymbol\sigma};\mathbf w_0)=\lim_{\theta\downarrow0}g_{V'}(\bar{\boldsymbol\sigma};\theta)$ be the \emph{reach of the tube} over $V'$.
\begin{enumerate}[label=(\roman*)]
\item $r(\bar{\boldsymbol\sigma};\mathbf w_0+i\bar{\boldsymbol\vartheta})\le r^\sharp_{V'}$ for every $\bar{\boldsymbol\vartheta}\in V'$; in particular $r(\bar{\boldsymbol\sigma};\mathbf w_0)\le r^\sharp_{V'}$, with equality when $V'=\mathbb R\bar{\boldsymbol\sigma}$.
\item \emph{A certificate on the real fibres.} Let $\mathbf v_1,\dots,\mathbf v_{n_{\mathcal D}-1}$ complete $\bar{\boldsymbol\sigma}$ to a basis of $\vec A_B$, and $y_i=\ip{\mathbf v_i}{\log\mathbf x}$ on the torus. If the image of the measure $\mathbf x^{-\mathbf w_0}(\boldsymbol\varrho^\rho_B)_\ast\phi^\rho_B$ under $(\mathbf x^{-\bar{\boldsymbol\sigma}},\mathbf y)$, on the part where $\mathbf x^{-\bar{\boldsymbol\sigma}}>x_0$, is $\sum_{\eta,j}X^{-\eta-1}(\log X)^jdX\otimes\phi_{\eta,j}(d\mathbf y)$, $j\le J$, with finite signed measures $\phi_{\eta,j}$, the $\eta$ in finitely many sets $\eta_i+\mathbb N$, and $\sum_{\eta,j}\|\phi_{\eta,j}\|x'^{-\eta}<\infty$ for every $x'>x_0$, then $g_{\vec A_B}(\bar{\boldsymbol\sigma};0)\le x_0$, so that $r^\sharp_{\vec A_B}\le x_0$: an expansion at the face over the real fibres certifies the tube, with no allowance for the phase.
\end{enumerate}
\end{corollary}

The certificate asks the expansion to hold over whole fibres of the directional ratio, with fixed measures on them, along one direction. Read in the coordinates in which the shares pair with the law linearly, expansions with fixed tangential measures at the falling ends certify the tube in every direction at once: in those coordinates a complex share multiplies the law by a factor of modulus one, and the continuation is made on the law itself.

% Reference 20: def:logcharts; source line 9861
\begin{definition}[Log expansions of a law]\label{def:logexpansions}
For a piece, with the space $\vec A_B$ of the directions of its shares of Definition~\ref{def:wallsgrowth} and its dual $\vec A_B^\ast$, the \emph{log-ratio map} sends a point $\mathbf x$ of the seam torus to the linear function $\mathbf z\mapsto\ip{\mathbf z}{\log\mathbf x}$ on $\vec A_B$, written $\log\mathbf x|_{\vec A_B}$; for a real point $\bar{\mathbf c}$ of the pins at which the seam function converges, $\lambda_{\bar{\mathbf c}}=(\log\mathbf x|_{\vec A_B})_\ast\big(\mathbf x^{-\bar{\mathbf c}}(\boldsymbol\varrho^\rho_B)_\ast\phi^\rho_B\big)$ is a finite positive measure on $\vec A_B^\ast$, and the seam function at $\bar{\mathbf c}+\mathbf z$ is $\int e^{-\ip{\mathbf z}f}\,d\lambda_{\bar{\mathbf c}}(f)$ for $\mathbf z\in\vec A_B+i\vec A_B$ wherever the integral converges absolutely. A \emph{log expansion} is the image of the complex measure
\[
\sum_{\boldsymbol\eta,\mathbf j}\mathbf p^{\mathbf j}e^{-\ip{\boldsymbol\eta}{\mathbf p}}\,d\mathbf p\otimes\phi_{\boldsymbol\eta,\mathbf j}(df)\qquad\text{on }[0,\infty)^k\times Y_Q
\]
under $(\mathbf p,f)\mapsto f-\sum_{h\le k}p_hu_h$, where $u_1,\dots,u_k\in\vec A_B^\ast$ are linearly independent, $\boldsymbol\eta$ runs over the union of finitely many translates of $\mathbb N^k$ by positive vectors, $\mathbf j\le\mathbf J$, and the $\phi_{\boldsymbol\eta,\mathbf j}$ are complex measures on a compact set $Y_Q\subset\vec A_B^\ast$, its \emph{inner face}, with $\sum\|\phi_{\boldsymbol\eta,\mathbf j}\|<\infty$. The law \emph{has log expansions at} $\bar{\mathbf c}$ when $\lambda_{\bar{\mathbf c}}$ is the sum of finitely many log expansions $Q$ and of a complex measure $\phi_0$ carried by a compact set $Y_0$. Its \emph{inner bound} along a direction $\bar{\boldsymbol\sigma}\in\vec A_B$ is the largest value of the directional ratio $e^{-\ip{\bar{\boldsymbol\sigma}}f}$ on $Y_0$ and on the inner faces,
\[
X_{\mathrm{in}}(\bar{\boldsymbol\sigma})=\max\Big(\max_{f\in Y_0}e^{-\ip{\bar{\boldsymbol\sigma}}f},\ \max_Q\max_{f\in Y_Q}e^{-\ip{\bar{\boldsymbol\sigma}}f}\Big),
\]
whose logarithm is convex in $\bar{\boldsymbol\sigma}$.
\end{definition}

An end chart of Theorem~\ref{thm:lawendchart} on a compact tangential source, whose tangential map and cut do not depend on $\mathfrak t$ and whose normal coordinates enter the log-ratio map through $k$ linearly independent elements of $\vec A_B^\ast$, gives finitely many log expansions and a compact part: the log-ratio map is linear in $\log\mathfrak t$ and in the logarithms of the tangential coordinates; with $p_h=\log(R'_h/\mathfrak t_h)$ on the polydisc of the radii $R'_h<R_h$ the exponents are $\boldsymbol\kappa+\mathbf n$ and the coefficients $\widetilde A_{\mathbf n}R'^{\,\boldsymbol\kappa+\mathbf n}$, summable by Theorem~\ref{thm:lawendchart}(ii); where some $\mathfrak t_h$ lie in $[R'_h,R_h)$ these join the tangential coordinates, the others keeping their expansions, summable uniformly by Cauchy's estimate on the closed polydisc; and where all of them do, the part is compact. A law with independent coordinates has such log expansions from the expansions of its factors at their ends.

% Reference 20: prop:logcharts; source line 9875
\begin{proposition}[A law with log expansions continues without phase]\label{prop:logexpansions}
Let the law of a piece have log expansions at $\bar{\mathbf c}$ (Definition~\ref{def:logexpansions}). Then for $\mathbf z\in\vec A_B+i\vec A_B$ off the hyperplanes $\ip{\mathbf z}{u_h}=\eta_h$ of the log expansions,
\[
P_{W_B(\bar{\mathbf c}+\mathbf z)}(\mvec_B)=\int e^{-\ip{\mathbf z}f}\,d\phi_0(f)+\sum_Q\sum_{\boldsymbol\eta,\mathbf j}\ \prod_{h\le k}\frac{j_h!}{(\eta_h-\ip{\mathbf z}{u_h})^{j_h+1}}\int_{Y_Q}e^{-\ip{\mathbf z}f}\,d\phi_{\boldsymbol\eta,\mathbf j}(f),
\]
continued: the seam function at $\bar{\mathbf c}+\cdot$ is meromorphic on $\vec A_B+i\vec A_B$ with its poles on these hyperplanes, and there are $C$ and $n$ with
\[
\big|P_{W_B(\bar{\mathbf c}+\mathbf z)}(\mvec_B)\big|\le C\,\delta^{-n}\max\Big(\max_{f\in Y_0}e^{-\ip{\re\mathbf z}f},\ \max_Q\max_{f\in Y_Q}e^{-\ip{\re\mathbf z}f}\Big),
\]
$\delta$ the smaller of $1$ and the distance of $\mathbf z$ from these hyperplanes; the same holds, with another constant, for $\delta$ the smaller of $1$ and the distance from the polar hyperplanes of the seam function alone. On the contour through a point $\bar{\mathbf c}+\tau\bar{\boldsymbol\sigma}$, $\tau\ge0$, the bound is $C\delta^{-n}X_{\mathrm{in}}(\bar{\boldsymbol\sigma})^\tau$, with no factor in the imaginary parts and no power of $\tau$. So, for a piece whose seam word has a Newton sum of full dimension and a rational $\bar{\boldsymbol\sigma}$, $r(\bar{\boldsymbol\sigma};\bar{\mathbf c})\le r^\sharp_{\vec A_B}(\bar{\boldsymbol\sigma};\bar{\mathbf c})\le g_{\vec A_B}(\bar{\boldsymbol\sigma};0)\le X_{\mathrm{in}}(\bar{\boldsymbol\sigma})$, and the growth of Theorem~\ref{thm:pastwalls}(iii) is $g(\boldsymbol\sigma;0)\le X_{\mathrm{in}}(\bar{\boldsymbol\sigma})$ for every direction $\boldsymbol\sigma$ of the shares of Theorem~\ref{thm:pastwalls}(i) whose sums over the monomials are $\bar{\boldsymbol\sigma}$.
\end{proposition}

The certificate of Corollary~\ref{cor:dirtube}(ii) is one expansion along the direction itself, with fibre measures of any extent and no condition on the rest of the law: on the contours through the points of a ray of $\bar{\boldsymbol\sigma}$ the factor $\mathbf x^{-i\bar{\boldsymbol\vartheta}}$ has modulus one and only the values of the directional ratio enter, not the compactness of the faces. Neither certificate applies as stated where the tangential coordinates of a face move with the normal ones, as when the support of the law curves in its log-ratio coordinates, or where a cut of the fibres moves, as in Theorem~\ref{thm:lawendchart}: the measures on the fibres are then not fixed, and the imaginary translates are read one at a time by Theorem~\ref{thm:dirreach}, or bounded by the ceilings of Theorem~\ref{thm:dirheights}, which carry the phase of the letters on these lines through their constants.

Without such a certificate a tube can reach farther than every one of its lines. There are sector integrals in two shares, with a single wall and growth bounds, for which every line of the tube along the fall of the wall has a finite reach while the reach of the tube over all the directions of the shares is infinite: the factor $\mathbf x^{-i\bar{\boldsymbol\vartheta}}$ enters the Taylor coefficients of a letter at the falling face through the powers of $\bar{\boldsymbol\vartheta}$. The seam function of a piece can do the same: when its seam function on the pins is a Beta corner whose two arguments move independently on the pins, it has growth bounds only at positive $\theta$, and on an imaginary translate along which one argument alone moves in the imaginary direction the Beta corner grows like a power of the imaginary part whose degree grows with the real parts, so that its tube reaches infinitely far while each of its lines has a finite reach. The sector letters of a piece of tables are never monomials in the falling variables, through which alone the factor $\mathbf x^{-i\bar{\boldsymbol\vartheta}}$ would enter without reaching the Taylor coefficients at the falling face. Theorem~\ref{thm:dirclosing}(i) is the limit $\theta\to0$ of Theorem~\ref{thm:pastwalls}(iii); when the reach of the tube is infinite, the bound of that theorem at a positive $\theta$, which pays for the phase, is the one that travels.

% Reference 20: thm:dirheights; source line 9908 (a piece of tables).
% Proved in the last part: thm:dirheights.
\begin{theorem}[Ceilings, simple tops and fibres along a direction]\label{thm:dirheights}
Let the seam function of the piece converge at a real point of the pins with positive shares, let $\bar{\boldsymbol\sigma}$ and $\mathbf w_0$ be as in Definition~\ref{def:dirreach}, the walls being those of a fan of the seam word (Theorem~\ref{thm:pastwalls}(ii)), and for a sector $\pi$ of it let $\varrho_{\pi,\bar{\boldsymbol\sigma}}=\prod_kq_k^{\dot\kappa_k}\,\mathbf k_\pi^{-\bar{\boldsymbol\sigma}}$ be the directional ratio in the sector (Definition~\ref{def:dirkeyhole}), the sector letters $L$ of $\mathbf k_\pi$ entering with the coordinates $\bar\sigma_L$ of $\bar{\boldsymbol\sigma}$ at their stations and $w_{0,L}$ of $\mathbf w_0$. The keyhole continuations along a direction $\boldsymbol\sigma$ depend on $\boldsymbol\sigma$ only through $\bar{\boldsymbol\sigma}$, and the falling variables of $\pi$ are those with $\dot\kappa_k<0$; for a piece of tables the sectors of the orders $\pi$ of $B$ may be taken.
\begin{enumerate}[label=(\roman*)]
\item \emph{No reach exceeds the ceiling.} For every keyhole continuation along $\bar{\boldsymbol\sigma}$, of ceiling $H_{\mathcal O}$, and every $\bar{\boldsymbol\vartheta}\in\vec A_B$, $r(\bar{\boldsymbol\sigma};\mathbf w_0+i\bar{\boldsymbol\vartheta})\le H_{\mathcal O}$: the phase of the letters does not enter the reach.
\item \emph{The values at the integers.} Let $\bar{\boldsymbol\sigma}$ be scaled so that every $\dot\kappa_k$ is an integer, and let the walls of the falling variables of every sector be off the integers at $\mathbf w_0$. An \emph{integer keyhole continuation along} $\bar{\boldsymbol\sigma}$ is defined as in Definition~\ref{def:dirkeyhole} with every sector letter kept from vanishing except those $L$ whose exponent $-w_{0,L}-m\bar\sigma_L$ at $\mathbf w_0+m\bar{\boldsymbol\sigma}$ is a nonnegative integer for every large integer $m$, its ceiling defined in the same way; its keyhole contour, its tops and its simple tops, with their exponents $\operatorname{ex}(p)$ and coefficients $A(p)$, are those of Definition~\ref{def:intkeyhole} with $\varrho_{\pi,\bar{\boldsymbol\sigma}}$ in place of $\varrho_\pi$, the sector form $\phi_{\pi,\mathbf w_0}=|\det(\mathbf v_{\pi,1},\dots,\mathbf v_{\pi,|B|-1})|\prod_kq_k^{\kappa_k(\mathbf w_0)-1}\mathbf k_\pi^{-\mathbf w_0}dq$ in place of $\phi_\pi$, the walls $\kappa_j(\mathbf w_0)$ of the sector in place of $\kappa_j$, and the faces $q_j=0$ with $\dot\kappa_j=0$ in place of the faces of the level variables. Then for every large integer $m$ the seam function at $\mathbf w_0+m\bar{\boldsymbol\sigma}$ is the integral of $\varrho_{\pi,\bar{\boldsymbol\sigma}}^{\,m}\,\phi_{\pi,\mathbf w_0}$ over the keyhole contour, and the growth at the integers $\operatorname{growth}_{\mathbb Z}(\bar{\boldsymbol\sigma};\mathbf w_0)=\limsup_m|P_{W_B(\mathbf w_0+m\bar{\boldsymbol\sigma})}(\mvec_B)|^{1/m}$ is at most the ceiling of every integer keyhole continuation along $\bar{\boldsymbol\sigma}$, and at most $r(\bar{\boldsymbol\sigma};\mathbf w_0)$.
\item \emph{Simple tops.} If an integer keyhole continuation along $\bar{\boldsymbol\sigma}$ of ceiling $H_{\mathcal O}$ has finitely many tops, all simple, and some $A_\zeta\neq0$ (Theorem~\ref{thm:simpletop}(iii)), then $\operatorname{growth}_{\mathbb Z}(\bar{\boldsymbol\sigma};\mathbf w_0)=H_{\mathcal O}\le r(\bar{\boldsymbol\sigma};\mathbf w_0)$; if moreover some keyhole continuation along $\bar{\boldsymbol\sigma}$ has the ceiling $H_{\mathcal O}$, then $r(\bar{\boldsymbol\sigma};\mathbf w_0)=H_{\mathcal O}$.
\item \emph{Fibres.} Let $\bar{\boldsymbol\sigma}$ be scaled so that its coordinates and every $\dot\kappa_k$ are integers. Call a box of a continuation along $\bar{\boldsymbol\sigma}$ a \emph{fibre box} when on its sets $Y$ the letters with $\bar\sigma_L\ge0$ do not vanish, and those with $\bar\sigma_L<0$ do not vanish at the points with a falling coordinate $0$; its stationary bound, as in Definition~\ref{def:fibrekeyhole} with $\varrho_{\pi,\bar{\boldsymbol\sigma}}$ in place of $\varrho_\pi$, is the supremum of $|\varrho_{\pi,\bar{\boldsymbol\sigma}}|$ over the points of these sets with no falling coordinate $0$ at which some falling coordinates lie inside their discs and the derivatives of $\varrho_{\pi,\bar{\boldsymbol\sigma}}$ in them vanish, the points with only one such coordinate, of falling rate one, left out. For every continuation along $\bar{\boldsymbol\sigma}$ whose boxes are keyhole boxes or fibre boxes, of ceiling $H_{\mathcal O}$, with the stationary bound $S$ of its fibre boxes, $r(\bar{\boldsymbol\sigma};\mathbf w_0)\le\max(H_{\mathcal O},S)$.
\end{enumerate}
\end{theorem}

So the reach along a direction sits between the growth at the integers and the ceiling of every keyhole continuation along it, and the two meet at simple tops that do not cancel when a keyhole continuation along the direction has the same ceiling, as at width one; and, by Theorem~\ref{thm:dirreach}, without any top it is the far boundary of the region on which the directional law is its own wall series. The phase of the letters, which Theorem~\ref{thm:keyholebound} pays for at the points where the ceiling is read, does not enter the reach: on a line the imaginary parts stay bounded. It enters only through the transverse fibre of a closing, which Theorem~\ref{thm:dirclosing} reads.

The integer scaling in (iv) is needed: a sector with one falling variable of falling rate one and a letter entering with a coordinate that is not an integer can have a fibre box of stationary bound $0$ and a reach above its ceiling, the directional ratio not being single valued in the disc, so that Lemma~\ref{lem:degreeone} does not apply; scaled to integers, the stationary point that the bound left out is counted.

% Reference 20: thm:dirclosing; source line 9939 (the partner in its monomials).
% Proved in the last part: thm:dirclosing.
\begin{theorem}[The closing past the walls is read along its direction]\label{thm:dirclosing}
Read the glue as in Theorem~\ref{thm:pastwalls}(i), the partner in its monomials, for a piece with walls and growth bounds whose seam word has a Newton sum of full dimension. Let $V$ be the space of the directions of the shares, of dimension $n$, the number of free shares, $\boldsymbol\sigma\in V$ a direction with rational coordinates, $\bar{\boldsymbol\sigma}$ its sums, $V'\subseteq\vec A_B$ the space of the sums of the directions in $V$, and $\mathbf c$ the real starting point, with $\bar{\mathbf c}\in\Omega_B$, off the null set of Theorem~\ref{thm:closecone} taken with the polar hyperplanes of the seam function. Let $h_\Gamma$ be the Gamma rate of Theorem~\ref{thm:pastwalls}(iii), so that the rate of that theorem, which holds for such a piece (Theorem~\ref{thm:pastwalls}, its seam function converging at the sums $\bar{\mathbf c}\in\Omega_B$ of the positive shares of the starting point), is $h_\theta(\boldsymbol\sigma)=h_\Gamma(\boldsymbol\sigma)+\log g(\boldsymbol\sigma;\theta)+\frac{2\theta^2}\pi\sum_{\sigma_l>0}\sigma_l$, and let $I$ be the integrand of the share side.
\begin{enumerate}[label=(\roman*)]
\item \emph{The tube bounds the closing.} If $h_\Gamma(\boldsymbol\sigma)+\log r^\sharp_{V'}(\bar{\boldsymbol\sigma};\bar{\mathbf c})<0$, the closing along $\boldsymbol\sigma$ converges absolutely, its terms and the integrals left after them being $O\big(e^{\tau(h_\Gamma(\boldsymbol\sigma)+\log r^\sharp_{V'}+\epsilon)}\big)$ for every $\epsilon>0$; the closing along $\boldsymbol\sigma$ is the closing along every direction of its chamber (Theorem~\ref{thm:chamber}), whose rates bound it as well. For the share side of Theorem~\ref{thm:pastwalls}(iv), with $\boldsymbol\sigma$ a direction of its shares with rational coordinates and $\mathbf c$ off the null set of that part, the closing along $\boldsymbol\sigma$ converges absolutely when $\sum_l\sigma_l\log|\sigma_l|+\log r^\sharp_{B}+\log r^\sharp_{R}<0$, in this sentence $r^\sharp_{B}$ and $r^\sharp_{R}$ being the reaches of the tubes over the space of the directions of its shares (Corollary~\ref{cor:dirtube}) of the seam functions of the piece and of the partner along $\boldsymbol\sigma$ at $\mathbf c$; its terms and the integrals left after them are then $O\big(e^{\tau(\sum_l\sigma_l\log|\sigma_l|+\log r^\sharp_{B}+\log r^\sharp_{R}+\epsilon)}\big)$ for every $\epsilon>0$.
\item \emph{One free share.} Let $n=1$, and put $r=r(\bar{\boldsymbol\sigma};\bar{\mathbf c})$ and $\ell(z)=P_{W_B(\bar{\mathbf c}+z\bar{\boldsymbol\sigma})}(\mvec_B)$. Then $r^\sharp_{V'}=r$, and for every $r'>r$ the terms are $O\big((1+\tau)^{n'}e^{\tau(h_\Gamma(\boldsymbol\sigma)+\log r')}\big)$; so the closing converges absolutely when $h_\Gamma(\boldsymbol\sigma)+\log r<0$. Let $r>0$, let $\mathbb L$ be the progression of the crossings of a falling share $l_0$, and let the data keep the poles apart: from some point on, $\mathbb L$ meets no pole of $\ell$ and no pole of another Gamma, and no pole of $\ell$ is a pole of a Gamma. If the poles of $\ell$ crossed by the ray are simple from some point on, or the values of $\ell$ on $\mathbb L$ carry the reach (Theorem~\ref{thm:dirreach}(iv)), then for every $\epsilon>0$ infinitely many terms have modulus at least $e^{\tau(h_\Gamma(\boldsymbol\sigma)+\log r-\epsilon)}$: the rate $h_\Gamma(\boldsymbol\sigma)+\log r$ is exact, and the terms do not tend to zero when it is positive. Without either proviso, if those poles have order at most $p$, the same holds for all but at most $p-1$ of the partners obtained by multiplying the coefficients by $e^{\tau_0v_l}$, $\tau_0\in\mathbb R$, for a fixed $\mathbf v$ with $\ip{\mathbf v}{\boldsymbol\sigma}=1$, the rate of the partner of $\tau_0$ being $h_\Gamma(\boldsymbol\sigma)-\tau_0+\log r$. The same holds for the share side of Theorem~\ref{thm:gluewidth}(iii) with one free share and a partner whose seam function converges along the ray, with $\sum_l\sigma_l\log|\sigma_l|+\log\sup\mathbf V_R^{-\bar{\boldsymbol\sigma}}$ in place of $h_\Gamma(\boldsymbol\sigma)$, the partners being the tinkered partners that multiply the room letters of each station $D$ by $e^{\tau_0v_D}$ (Definition~\ref{def:tinkers}); at a node of width one, where the share side has one free share, this rate is $\log\big(\operatorname{reach}(B)/\inf\varrho_R\big)$, the logarithm of the root rate of Theorem~\ref{thm:nodeintervals}(ii).
\item \emph{Several free shares.} Let $n\ge2$. Fix a complement $V_\perp$ of $\mathbb R\boldsymbol\sigma$ in $V$, a Lebesgue measure $d\boldsymbol\vartheta_\perp$ on it and $k_\sigma>0$ with $d\boldsymbol\vartheta=k_\sigma\,dy\,d\boldsymbol\vartheta_\perp$ for $\boldsymbol\vartheta=y\boldsymbol\sigma+\boldsymbol\vartheta_\perp$, $d\boldsymbol\vartheta$ the measure of Definition~\ref{def:fibre}. Let $T_\sigma$ be the space of the \emph{transverse tinkers}, the tinkers of the partner read in its monomials (Definition~\ref{def:tinkers}, with one multiplier $e^{v_l}$ for each monomial in place of each coordinate of the seam vector), multiplying the coefficients read on $V$ as $\xi_l\mapsto\xi_le^{v_l}$, with $\ip{\mathbf v}{\boldsymbol\sigma}=0$, identified with the dual of $V_\perp$ and carrying the dual Lebesgue measure $d\mathbf v$. The crossings of the closing along $\boldsymbol\sigma$ are the same for every tinker, and at a crossing at the time $\tau^\ast$
\[
\int_{T_\sigma}e^{\ip{\mathbf v}{\mathbf c}}\,\mathcal R_{\tau^\ast}\big(\boldsymbol\xi e^{\mathbf v}\big)\,d\mathbf v=-k_\sigma\operatorname*{Res}_{z=\tau^\ast}I(\mathbf c+z\boldsymbol\sigma),
\]
$\mathcal R_{\tau^\ast}(\boldsymbol\xi e^{\mathbf v})$ the crossing term of the partner with the coefficients $\boldsymbol\xi e^{\mathbf v}$. Integrated over the transverse tinkers against the factor $e^{\ip{\mathbf v}{\mathbf c}}$, the terms of the closing are $-k_\sigma$ times the residues of the integrand on the real line through $\mathbf c$ in the direction $\boldsymbol\sigma$; these are $O\big((1+\tau)^{n'}e^{\tau(h_\Gamma(\boldsymbol\sigma)+\log r')}\big)$ for every $r'>r=r(\bar{\boldsymbol\sigma};\bar{\mathbf c})$, and under the conditions of (ii) on that line infinitely many of them have modulus at least $e^{\tau(h_\Gamma(\boldsymbol\sigma)+\log r-\epsilon)}$.
\end{enumerate}
\end{theorem}

% Reference 20: thm:onefreeshare; source line 9987
\begin{theorem}[A closing with one free share is read on the law along it]\label{thm:onefreeshare}
Take, in Theorem~\ref{thm:averaging}, the share side of the kind (B), its bundles merged by Corollary~\ref{thm:E1}, the bars marking sums over the bundles, and in this statement write
\[
I(\bar{\mathbf a})=\frac{\prod_G\Gamma(\bar a_G)}{\prod_j\Gamma(s_j)\prod_j\Gamma(t_j)}\ P_{W^{\mathrm m}_B(\bar{\mathbf a})}(\mvec_B),
\]
the products over the bundles and over the active letters, with the merged seam function of Lemma~\ref{lem:newtonwalls}. Let the pins of the merged shares leave a fibre of dimension one, spanned by a rational direction $\bar{\boldsymbol\sigma}$, and let $\bar{\mathbf c}$ be a real point of the merged pins with positive coordinates at which the merged seam function converges. Let the merged law along $\bar{\boldsymbol\sigma}$, the law along $\bar{\boldsymbol\sigma}$ of the merged seam function (Definition~\ref{def:dirreach}), be the image of $\prod_Ge^{-\bar c_Gy_G}\lambda$ under $X=\prod_Ge^{-\bar\sigma_Gy_G}$, $\lambda$ the source measure, its line function $\ell(\zeta)=P_{W^{\mathrm m}_B(\bar{\mathbf c}+\zeta\bar{\boldsymbol\sigma})}(\mvec_B)$ and its reach $r$, the least $x_0\ge0$ for which it has on $(x_0,\infty)$ a density that is a wall series of radius $x_0$ (Definition~\ref{def:dirreach}, Theorem~\ref{thm:dirreach}(iii)); and in this statement write $h^{\mathrm m}_\Gamma(\bar{\boldsymbol\sigma})=\sum_G\bar\sigma_G\log|\bar\sigma_G|$, the Gamma rate of the merged share side, whose coefficients are inside the merged values.
\begin{enumerate}[label=(\roman*)]
\item There is $k_P>0$, depending only on the pins, with $P_{B\sqcup R}(\mvec)=\frac{k_P}{2\pi i}\int_{\re\zeta=0}I(\bar{\mathbf c}+\zeta\bar{\boldsymbol\sigma})\,d\zeta$, and the crossing term (Lemma~\ref{lem:pathclose}) at a crossing $\zeta^\ast$ of the ray is $-k_P\operatorname{Res}_{\zeta=\zeta^\ast}I(\bar{\mathbf c}+\zeta\bar{\boldsymbol\sigma})$: at the crossing $\bar a_f=-n$ of the falling share $f$, at which no other factor of $I$ is polar, it is
\[
k_P\,\frac{(-1)^n}{n!\,|\bar\sigma_f|}\,\frac{\prod_{G\neq f}\Gamma(\bar a_G)}{\prod_j\Gamma(s_j)\prod_j\Gamma(t_j)}\ P_{W^{\mathrm m}_B(\bar{\mathbf a})}(\mvec_B)
\]
at $\bar{\mathbf a}=\bar{\mathbf c}+\zeta^\ast\bar{\boldsymbol\sigma}$, and at a simple pole of the merged seam function at which no Gamma is polar it is $-k_P\prod_G\Gamma(\bar a_G)/\big(\prod_j\Gamma(s_j)\prod_j\Gamma(t_j)\big)$ there times $\operatorname{Res}_{\zeta^\ast}\ell$.
\item For every $r'>r$ there are $C$ and $n'$ such that the crossing terms at their times $\tau$, and the integrals left at the times at a fixed distance from the crossings, are at most $C(1+\tau)^{n'}e^{\tau h^{\mathrm m}_\Gamma(\bar{\boldsymbol\sigma})}r'^{\,\tau}$.
\item Let the ray $\bar{\mathbf c}+\tau\bar{\boldsymbol\sigma}$, $\tau\ge0$, cross its polar points simply: no crossing point of a falling merged share, one with $\bar\sigma_G<0$, is a pole of the merged seam function or a crossing point of another falling share, and the poles of the merged seam function on the ray are simple. Then for each falling share $f$ the crossing terms of $f$ and those at the poles of the merged seam function have together the root rate $e^{h^{\mathrm m}_\Gamma(\bar{\boldsymbol\sigma})}r$: the upper limit of $|\mathcal R|^{1/\tau}$ over them is $e^{h^{\mathrm m}_\Gamma(\bar{\boldsymbol\sigma})}r$. The same holds when the data keep the poles apart only from some point on, as in Theorem~\ref{thm:dirclosing}(ii), read on $\ell$ and on the progression of the crossings of $f$, and the poles of $\ell$ crossed by the ray are simple from some point on or the values of $\ell$ on that progression carry the reach; without either proviso, if those poles have order at most $p$, it holds for all but at most $p-1$ of the partners obtained by multiplying the coefficients of the shares of each bundle $G$ by $e^{\tau_0v_G}$, $\tau_0\in\mathbb R$, for a fixed $\mathbf v$ with $\sum_Gv_G\bar\sigma_G=1$, whose merged laws are constant multiples of the images of the merged law under $X\mapsto e^{-\tau_0}X$, with the root rates $e^{h^{\mathrm m}_\Gamma(\bar{\boldsymbol\sigma})-\tau_0}r$.
\item So the closing of the merged share side along $\bar{\boldsymbol\sigma}$, $P_{B\sqcup R}(\mvec)=\sum\mathcal R$ over the crossings of the ray in their order, converges absolutely when $e^{h^{\mathrm m}_\Gamma(\bar{\boldsymbol\sigma})}r<1$, and under the conditions of (iii) its terms do not tend to zero when $e^{h^{\mathrm m}_\Gamma(\bar{\boldsymbol\sigma})}r>1$. No hypothesis on tops or on the phase of the letters enters, only the separations of the poles in (iii): the rate is read on the merged law along the direction. Where $X$ is bounded on the support of the law of the piece, $r$ is its largest value there, and when moreover the falling shares are alone in their bundles, $h^{\mathrm m}_\Gamma(\bar{\boldsymbol\sigma})+\log r$ is the averaged rate $\hat h_C$ of Theorem~\ref{thm:averagingsharp}; where $X$ is unbounded, $\hat h_C$ is infinite and the averaged terms diverge past the walls, and (iii) carries the sharpness of that theorem past them, with the reach of the merged law in place of the largest value of $X$. When every bundle is a single share, $\boldsymbol\sigma=\bar{\boldsymbol\sigma}$, the merged seam function is the product of $\boldsymbol\xi^{-\mathbf a}$ and of the seam function at the sums over the monomials (Theorem~\ref{thm:pastwalls}(i)), $r$ is $e^{-\sum_l\sigma_l\log\xi_l}$ times the reach of the piece along the sums of $\boldsymbol\sigma$ (Definition~\ref{def:dirreach}), and $e^{h^{\mathrm m}_\Gamma(\bar{\boldsymbol\sigma})}r=e^{h_0(\boldsymbol\sigma)}$, $h_0$ the rate of Theorem~\ref{thm:pastwalls}(iii) at $\theta=0$, whose growth $g(\boldsymbol\sigma;0)$ is on a fibre of dimension one that reach: on such a fibre the rate of Theorem~\ref{thm:pastwalls}(iii) at $\theta=0$ is exact under the conditions of (iii), and (iii) is Theorem~\ref{thm:dirclosing}(ii).
\end{enumerate}
\end{theorem}

The merged share side and the share side at the stations, before the stations of one excess are merged (Corollary~\ref{thm:E1}), are two representations of one period, with two laws along their directions; completing the splits of a bundle can move the edge of the support at which the law stops being its own wall series. The merged law can stop being its own wall series at the image of a nearer edge of the support, where a level set of the merged seam letter ends the expansion of the law at its falling end, while the station closings, with more free shares, can read a farther edge: the merged closing can diverge where the station closing converges.

% Reference 20: cor:anyfibre; source line 10018 (the law of the piece).
% Proved in the last part: cor:anyfibre.
\begin{corollary}[A law with log expansions bounds a closing of any fibre dimension]\label{cor:anyfibre}
In Theorem~\ref{thm:pastwalls}(iii), with the starting point $\mathbf c$ off its null set, let the law of the piece have log expansions at the sums $\bar{\mathbf c}$ of $\mathbf c$ (Definition~\ref{def:logexpansions}), with the inner bound $X_{\mathrm{in}}$. The function $h_\Gamma(\boldsymbol\sigma')+\log X_{\mathrm{in}}(\bar{\boldsymbol\sigma}')$ is convex on the chamber $C(\boldsymbol\sigma)$ of Definition~\ref{def:chamber}. If its infimum $h'$ on the chamber is negative, the closing along $\boldsymbol\sigma$ converges absolutely to $P_{B\sqcup R}(\mvec)$, its terms and the integrals left after them being $O(e^{\tau(h'+\epsilon)})$ for every $\epsilon>0$. The fibre may have any dimension, and no margin of coefficient factors and no phase enter: the rate is read on the compact part of the law and on the inner faces of its log expansions. In Theorem~\ref{thm:pastwalls}(iv), with $\mathbf c$ off the null set of that part, if the laws of the piece and of the partner have log expansions at $\mathbf c$, with the inner bounds written $X^B_{\mathrm{in}}$ and $X^R_{\mathrm{in}}$ in this sentence, the same holds with $\sum_l\sigma'_l\log|\sigma'_l|+\log X^B_{\mathrm{in}}(\boldsymbol\sigma')+\log X^R_{\mathrm{in}}(\boldsymbol\sigma')$, convex on the chamber $C(\boldsymbol\sigma)$, in place of $h_\Gamma(\boldsymbol\sigma')+\log X_{\mathrm{in}}(\bar{\boldsymbol\sigma}')$.
\end{corollary}

So the counterparts of the reach theorems of width one hold along every rational direction. The tinkers of Definition~\ref{def:tinkers} play two parts: all the tinkers of one partner separate the laws (Theorem~\ref{thm:retainwidth}), and the integral over the transverse tinkers of a closing, the partner read in its monomials, reads the reach of the real line (Theorem~\ref{thm:dirclosing}(iii)).

For a positive kernel the root rate of a crossing term is the exponential of the least Gamma rate of its chamber in every fibre dimension (Theorem~\ref{thm:retainedchamberrate}), and averaging reads the log-support of the source measure (Theorem~\ref{thm:averaging}); Corollary~\ref{cor:anyfibre} reads, past the walls, the compact part and the inner faces of the log expansions. It is an upper bound: with two or more free shares a crossing term is a fibre integral of positive dimension, whose size the law along one direction does not fix, and the best direction of a chamber can lose against the averaged closing of the same source (Theorem~\ref{thm:canonicaltablegap}); the exact rate is then read on a joint word at a noncancelling simple top (Theorem~\ref{thm:retainedtops}). One partner does not read the real line: it averages over its transverse fibre, and its terms can fall below the rate of every real line of its chamber, the integral of the terms over the transverse tinkers then being carried by the other tinkered partners. Whatever the closing, the reach along the direction of the closing is at least the least value of the directional ratio on the support of the law, and at least every point near which the directional law, or after merging the merged law, has no real-analytic density (Theorem~\ref{thm:dirreach}(iii)).

\section{What the calculus of bunches settles}\label{sec:widthopen}

% Reference 20: rem:widthopen; source line 10375
\begin{remark}[What the calculus of bunches settles, and what it leaves open]\label{rem:widthopen}
A closing of the glue along a direction of its share fibre is exact at every step and depends only on its chamber (Lemma~\ref{lem:pathclose}, Theorem~\ref{thm:chamber}); Theorems~\ref{thm:closecone}, \ref{thm:pastwalls} and~\ref{thm:keyholebound} give sufficient conditions for it, and the closings can be averaged over the source (Theorem~\ref{thm:averaging}). The cases that these conditions leave are settled as follows.
\begin{enumerate}[label=(\roman*)]
\item \emph{Crossing terms of positive dimension.} Along a rational ray, at simple nonresonant crossings and after a finite head, a crossing term kept with its falling shares integrated on their strips is, along its cohort, a scalar $G_n$ times a finite sum $\sum_i\big(\sT_{-n;(1)}[c_i]\,\phi_i\big)(Y_i)$ of joint Trailmoments of fixed functions on fixed compact domains, the joint word (Theorem~\ref{thm:retainedchain}): a falling Gamma becomes a positive share and one Beta variable (Lemma~\ref{lem:gammastriplift}), a continued source strip a segment and circle integrals (Lemma~\ref{lem:sourcestriplift}), and nothing is extracted from the term. For a positive kernel its rate is the least Gamma rate of the chamber, in every fibre dimension and with no hypothesis on the tops (Theorem~\ref{thm:retainedchamberrate}); averaged over a positive source it is the averaged rate $\hat h_C$ (Corollary~\ref{cor:retainedaveragedrate}), so a nonzero averaged rate decides by its sign whether the averaged terms tend to zero, with no condition on the bundles or the dimension (Theorem~\ref{thm:averagingsharp}), and at $\hat h_C=0$ Theorem~\ref{thm:unitrate} decides. The joint word has the root rate $r_Gc_{\max}$ at noncancelling simple tops, $r_G$ the root rate of the scalar and $c_{\max}$ the ceiling of the word (Theorem~\ref{thm:retainedtops}). Coincident crossings are packets of factors polar on one crossed hyperplane, assembled before a sign or a top is read (Theorem~\ref{thm:crossingjet}, Proposition~\ref{prop:topjets}); at the root rate one a positive moment sequence is decided by its law near the top (Theorem~\ref{thm:unitrate}), and an expansion by its coefficients (Proposition~\ref{prop:boundaryjets}).
\item \emph{The phase.} A joint word is bounded by $|G_n|\sum_i\|\phi_i\|\,c_{\max}^{\,n}$ (Corollary~\ref{cor:retainedheight}): the fibre is completed before a modulus is taken, and the phase of the letters enters the coefficient functions, where it can cancel, and costs no term of the bound. It confines the radii of the procedure to the margins of the coefficient factors (Theorem~\ref{thm:retainedchain}, Remark~\ref{rem:continuedtops}). Read on the law, the phase costs nothing where the law has log expansions: in its log-ratio coordinates the shares pair with the law linearly, and the compact part and the inner faces of its log expansions bound the continued seam function on whole contours (Proposition~\ref{prop:logexpansions}). Where the tangential coordinates move with the normal coordinates at a falling end, the phase enters the coefficients of the expansion there, and a tube can reach infinitely far while each of its lines has a finite reach.
\item \emph{Regions of equality and bounds at any width.} The regions on which a law equals its expansion at a falling end are certified from the sector letters, with the map from the motion of the shares to the exponents, the coefficients, the truncation bounds and exact boundary words (Theorems~\ref{thm:lawendchart} and~\ref{thm:lawregionatlas}); the domain where a series converges and the region where it equals the law are kept apart (Proposition~\ref{prop:lawreadingboundary}), and density charts are built at regular fibres (Proposition~\ref{prop:lawdensityregion}). The terms of a closing have at most the root rate $r_Gc_{\max}$ of every joint word whose radii keep within the holomorphy of its source and the margins of its coefficient factors, $c_{\max}$ the ceiling of the word, and have that root rate when the word has a noncancelling simple top.
\item \emph{The reach at any width.} Along every rational direction of the shares the reach is read on the law: it is the far boundary of the region on which the directional law, the law multiplied by $\mathbf x^{-\mathbf w_0}$ for the starting point $\mathbf w_0$ and read through the directional ratio, is its own wall series, whose exponents lie among the values $(\kappa_U(\mathbf w_0)+n)/(-\dot\kappa_U)$ of the walls divided by their falling rates, and it is at least every point at which that law is not analytic, the least value of the directional ratio among them (Theorem~\ref{thm:dirreach}). The walls and the growth bounds that this uses hold for every seam function of a piece that converges at a real point and for its seam letters merged by bundles (Lemma~\ref{lem:newtonwalls}), and at width one the same argument reads the growth of a family on its law, that of the bunch being its reach (Corollary~\ref{cor:growthreach}). For a piece the reach is at most the ceiling of every keyhole continuation along the direction, with no phase, equal to it at simple tops that do not cancel when a keyhole continuation has the same ceiling, and bounded by fibres where letters vanish inside the keyholes, the direction scaled to integers (Theorem~\ref{thm:dirheights}). A closing past the walls converges when the exponential of the Gamma rate times the reach of its tube is less than one, its terms then having at most that root rate (Theorem~\ref{thm:dirclosing}(i)); an expansion at the face over the real fibres bounds the reach of the tube along one direction and log expansions in every direction, with no allowance for the phase (Corollary~\ref{cor:dirtube}, Proposition~\ref{prop:logexpansions}), and otherwise the tube can reach farther than every one of its lines; a law with log expansions bounds a closing of any fibre dimension by the rate read on the compact part and the inner faces of its log expansions (Corollary~\ref{cor:anyfibre}). With one free share, as at a node of width one, the tube is the line, and the terms attain the rate that the reach along the direction gives when the poles of the line and of the Gammas keep apart and the poles of the line are eventually simple or its values on a progression of crossings carry the reach (Theorem~\ref{thm:dirclosing}); the closing of a share side merged by its bundles to one free share, a series other than the station closings, has under the same conditions the root rate of Theorem~\ref{thm:onefreeshare}(iii), read on the merged law along its direction, its terms not tending to zero beyond it, which carries the sharpness of Theorem~\ref{thm:averagingsharp} past the walls (Theorem~\ref{thm:onefreeshare}). In every number of free shares the terms integrated over the tinkers of the partner transverse to the direction are the residues of the real line, which attain its rate under the same conditions (Theorem~\ref{thm:dirclosing}), while the terms of one partner average over its transverse fibre and can fall below every real line of its chamber.
\item \emph{A gap at a family of tables.} A family of three tables whose seam has single stations has a chamber in which no direction certifies the closing while the averaged closing converges geometrically (Theorem~\ref{thm:canonicaltablegap}); the gap is relative to the closings of one share side by the directions of its chamber (Remark~\ref{rem:gapscope}).
\end{enumerate}
At width one the reach is read on the law, as the point beyond which the law is its own wall series, and it is the ceiling of a keyhole continuation whose tops are simple and do not cancel (Theorems~\ref{thm:lawreach} and~\ref{thm:simpletop}). At any width the first holds along every rational direction (Theorem~\ref{thm:dirreach}), and the second along every rational direction for every piece, under the conditions of Theorem~\ref{thm:dirheights}(ii) and~(iii), and for joint words (Theorem~\ref{thm:retainedtops}). What a node of width one does not have is a transverse fibre. With two or more free shares the crossing terms of one partner integrate over the transverse fibre the imaginary translates of the line through its starting point. Theorem~\ref{thm:dirclosing}(i) bounds them by the reach of the tube when the exponential of the Gamma rate times that reach is less than one; a law with log expansions bounds the tube by the directional ratio on the compact part and the inner faces of its log expansions, and a closing of any fibre dimension with nothing paid for the phase (Corollary~\ref{cor:anyfibre}), while without log expansions the reach of the tube can be infinite though every line of the tube has a finite reach, and the bound that travels pays for the phase at a positive $\theta$ (Theorem~\ref{thm:pastwalls}(iii)). Theorem~\ref{thm:dirclosing}(iii) gives the integral of the terms over the transverse tinkers the rate of the real line, under the conditions of its part (ii) on that line; Theorem~\ref{thm:retainedtops} gives the rate of one partner at the noncancelling simple tops of its joint words. The real line does not give the rate of one partner, whose terms can fall below every real line of its chamber, and the best direction of a chamber can lose against averaging (Theorem~\ref{thm:canonicaltablegap}). Before the walls, under the hypotheses of Theorem~\ref{thm:averagingsharp}, the rate of one partner is averaged over the source, and past them the closing of a share side merged by its bundles to one free share has, under the conditions of Theorem~\ref{thm:onefreeshare}(iii), the rate read on the merged law along its direction. Two problems are open. The first is the rate past the walls of one partner whose joint words have no noncancelling simple top: at two or more free shares, where it asks for an averaged closing of the continued law, and at one free share when the data do not keep the poles apart or no proviso of Theorems~\ref{thm:dirclosing}(ii) and~\ref{thm:onefreeshare}(iii) holds for that partner. The second is to read the growth $g(\boldsymbol\sigma;\theta)$ on the law, as Theorem~\ref{thm:dirreach} reads the reach of a line. At $\theta=0$ the reach of the line bounds it below and log expansions bound it above by their inner bound (Corollary~\ref{cor:dirtube}, Proposition~\ref{prop:logexpansions}), and the two meet when, for every $\epsilon>0$, the law has log expansions whose compact part and inner faces lie where the directional ratio is at most the reach plus $\epsilon$; in general, and without log expansions, where it can be infinite at $\theta=0$, it is known only through these bounds and those of Theorems~\ref{thm:pastwalls}(iii) and~\ref{thm:keyholebound} and of Corollary~\ref{cor:retainedheight}, within the margins of the coefficient factors (Theorem~\ref{thm:retainedchain}). Degenerate and cancelling tops are open at width one as well.
\end{remark}


\chapter{Price and complexity}\label{ch:price}

A construction is a description of a family, and the exact theorems of Part~\ref{part:construction} turn it into a description of the period; both have a price. The price of a construction is the length of its description, a repeated subtree described once and its multiplicity written in binary; the work of a word is the dimension of the integrals it asks for, an integral that recurs counted once. Repetition then costs what it should: $M$ copies of a cell glued along a tree at a full lean, or closed in series at any exponents, or chained, have periods whose words cost the word of the cell and the binary length of $M$, and a departure from the repetition costs the cells that depart. The laws of cells (Chapter~\ref{ch:hinged}), the networks of holds (Chapters~\ref{ch:holding} and~\ref{ch:leantree}) and the finite interfaces (Chapter~\ref{ch:interfaces}) are priced the same way: a law is computed once for each distinct subtree, a node of a network costs its periods at the two full leans whatever the number of bases of the whole table, and a polynomial Gamma moment of a repetition is the power of one finite matrix. What a value can contain is bounded the same way along a ray (Theorem~\ref{thm:rayspace}).

\section{Constructions and their price}\label{sec:constructionprice}

% Reference 20: def:construction; source line 10392
\begin{definition}[Construction]\label{def:construction}
A \emph{construction} $\mathfrak C$ is a finite rooted tree with a family at each leaf, its marks, bases, weights, masses and exponents given, and at each inner node one of the operations: the direct sum of its subtrees; the substitution of its subtrees, each with a named hinge, into the marks of a core family named at the node, a mark left without a subtree being left in place (Definition~\ref{def:substitution}); and the \emph{powers} of a subtree $\mathfrak C_1$ with an integer $M$: $\mathfrak C_1^{\oplus M}$, $M$ copies in direct sum; $\circ(\mathfrak C_1^M)$, the series closure of $M$ copies of the cell $\mathfrak C_1$ (Corollary~\ref{cor:blocksclosures}(ii)); and $\mathfrak C_1^{\rhd M}$, the chain of $M$ copies glued along two named marks (Theorem~\ref{thm:chainoperator}). A subtree may be given by a reference to one already described. A construction builds one family, by Proposition~\ref{prop:substitution} at each substitution.
\end{definition}

% Reference 20: def:price; source line 10396
\begin{definition}[Price; work]\label{def:price}
The \emph{price} $\operatorname{price}(\mathfrak C)$ of a construction is the number of bits of its description as written: for every leaf its marks, bases, weights, masses and exponents, the numbers written as rationals in binary; for every node its operation, its core and its attachments; for every power the $\lceil\log_2(M+1)\rceil$ bits of $M$; a subtree may be written once and referred to by an index wherever it recurs. The price of a word (Definition~\ref{def:conv}(vi)) is defined in the same way, and its \emph{work} $\operatorname{work}(w)$ is the sum of the dimensions of its distinct integrals, an integral that recurs, under a power or an iterate in particular, counted once. The integration count of Definition~\ref{def:contents} counts every occurrence: the work of $w^M$ is that of $w$, and its integration count is $M$ times that of $w$.
\end{definition}

The work counts the distinct integrals a word is written with, and the price the length of what has to be written down; neither counts the digits of an evaluated value, nor the cost of evaluating: a scalar power costs one exponentiation, a power of a $k\times k$ matrix costs $O(k^3\log M)$ operations by repeated squaring, and an iterate $\mathcal Q^M$ of an operator costs $M$ applications of $\mathcal Q$, nested integrals of $M$ times its dimension, or $O(\log M)$ matrix products when $\mathcal Q$ leaves a finite-dimensional space of functions invariant. Every polynomial Gamma moment of degree $D$ of a repetition (Chapter~\ref{ch:interfaces}) is such a matrix power, a fixed row times a power of the transfer matrix $\mathbf T_D$ applied to the moment vector of the start (Theorem~\ref{thm:transferpower}), and the period itself is reached through such moments with the certified stopping rule of Theorem~\ref{thm:certstop}, at a degree that depends on the accuracy asked.

% Reference 20: thm:pricerepeat; source line 10402
\begin{theorem}[Pricing by repetition]\label{thm:pricerepeat}
Let $\mathcal A$ be a family, a cell where a hinge is needed, described by a construction of price $\operatorname{price}(\mathcal A)$, with the period $P_{\mathcal A}$ at its data, and let $M\ge1$.
\begin{enumerate}[label=(\roman*)]
\item \emph{Direct sum.} The molecules of $\mathcal A^{\oplus M}$, $\mathcal A$ connected, are its $M$ copies, and the product of their periods is $P_{\mathcal A}^M$ (Proposition~\ref{prop:directsum}; the period of the direct sum itself diverges on the balance wall, Proposition~\ref{prop:compoundwall}): price at most $\operatorname{price}(P_{\mathcal A})+\lceil\log_2(M+1)\rceil+O(1)$, work that of $P_{\mathcal A}$.
\item \emph{A tree of copies at a full lean.} For a table $\mathcal A$ at the exponents $(0,t)$, $M$ copies glued along any tree at the masses $(\kappa,t-\kappa)$ on the two marks of every gluing have the word $\Gst21{\kappa,\ t-\kappa}{t}\big(P_{\mathcal A}(\mvec;0,t)/\Gst21{\kappa,\ t-\kappa}{t}\big)^M$ (Corollary~\ref{cor:primelean}): price at most $\operatorname{price}(P_{\mathcal A})+\lceil\log_2(M+1)\rceil+O(1)$, work that of $P_{\mathcal A}$, whatever the shape of the tree.
\item \emph{A series closure.} The series closure of $n_k$ copies of cells $\mathcal A_k$, $k\le d$, at positive exponents of $N$ tables on the balance wall has the word of Theorem~\ref{thm:seriesclosure}, $\prod_j(\Gamma(s_j)\Gamma(t_j))^{-1}\int\prod_jx_j^{s_j-1}\prod_{j\ge2}y_j^{t_j-1}\prod_k\mathfrak L_k(x,y)^{n_k}\,dx\,dy$ with the responses $\mathfrak L_k$ of the cells: price $O\big(\sum_k\operatorname{price}(\mathcal A_k)\big)+\sum_k\lceil\log_2(n_k+1)\rceil+O(N)$, the letters of a constructed cell being written through Theorem~\ref{thm:substletters}, and work at most $\sum_k|B_k|+2N-1$, $B_k$ the marks of $\mathcal A_k$ besides its hinge, independent of the multiplicities and of the order of the cells.
\item \emph{A chain.} $\mathcal A^{\rhd M}$ has the word $(\mathcal Q^MP_0)(\one)$ of Theorem~\ref{thm:chainoperator}, with the operator $\mathcal Q$ and the function $P_0$ of that theorem: price at most $\operatorname{price}(\mathcal A)+\lceil\log_2(M+1)\rceil+O(1)$, work at most $|\cE_{\mathcal A}|-1$; evaluated, the iterate unrolls to nested integrals of total dimension $M(|\cE_{\mathcal A}|-2)+1$, that of the expanded period, unless $\mathcal Q$ leaves a finite-dimensional space of functions invariant; where the twist is pinned the word is $P_{\mathcal A}^M/\Gst21{m_{g_{\mathrm{in}}},\ m_{g_{\mathrm{out}}}}{m_{g_{\mathrm{in}}}+m_{g_{\mathrm{out}}}}^{M-1}$, $m_{g_{\mathrm{in}}},m_{g_{\mathrm{out}}}$ the masses of the two named marks $g_{\mathrm{in}},g_{\mathrm{out}}$, of work $|\cE_{\mathcal A}|-1$ evaluated as well.
\item \emph{Integer joint moments} of a repeated transfer matrix, with Gamma laws on the marks or Dirichlet laws on the steps, have as word a fixed row times the power $\mathbf A^M$ of the one matrix of Theorem~\ref{thm:momenttransfer} times a fixed column, with rational entries at rational masses (Theorem~\ref{thm:momenttransfer}): price at most that of the matrices of the cell plus $\lceil\log_2(M+1)\rceil+O(1)$, work zero.
\item \emph{Blocks} read through their totals cost one stacked Gamma each, $\Gst{|B_i|}1{\mvec|_{B_i}}{m_{B_i}}$ at the block $B_i$ (Theorem~\ref{thm:blockagg}): the work of the period of the family with blocks is that of the period of its core family, whatever the lengths of the blocks.
\item \emph{A series repetition.} $M$ copies of a series-balanced cell $\mathcal A$ at a series node against a cell $\mathcal A'$, at the exponents $(s,t)$, have the word $(\Gamma(s)\Gamma(t))^{-1}\int_0^\infty y^{t-1}f_{\mathcal A}(y)^Mf_{\mathcal A'}(y)\,dy$, with the transforms $f_{\mathcal A},f_{\mathcal A'}$ of Theorem~\ref{thm:seriesrepeat}: price at most $\operatorname{price}(\mathcal A)+\operatorname{price}(\mathcal A')+\lceil\log_2(M+1)\rceil+O(1)$, work $|B|+|R|-1$ for the marks $B$ of $\mathcal A$ and $R$ of $\mathcal A'$ besides the hinge, whatever $M$; at a series or parallel node of a hinge tree the multiplicities of equal subtrees enter as exponents (Proposition~\ref{prop:work}).
\item \emph{Series-parallel tables.} A table with a hinge tree of series and parallel nodes over marks has periods at the full leans that are stacked Gammas times monomials in the weights of its nodes, with one Gamma for each mark and two for each change of kind along an edge, and every period, at every lean and every twist, a fibre integral of dimension $n-1$ of an explicit stacked Gamma times an exponential of the loads, and a family of $N$ such tables one of dimension $(2N-1)(n-1)$ (Corollaries~\ref{cor:spfulllean} and~\ref{cor:spnetwork}): price linear in the size of the tree, work at most $n-1$, less one for each pooled pair (Proposition~\ref{prop:poolpair}), and at most $(2N-1)(n-1)$ for the family, whatever the numbers of bases, which can be exponential in $n$.
\item \emph{Polynomial Gamma moments of a repetition} have as word a fixed row times a power of the transfer matrix $\mathbf T_D$ applied to a fixed vector, with rational entries at rational data (Theorem~\ref{thm:transferpower}): price that of the matrix of the chain cell plus $\lceil\log_2(M+1)\rceil+O(1)$, work zero.
\end{enumerate}
In the cases with a multiplicity $M$ the expanded family has of the order of $M$ times as many marks, and its period written as one integral over its simplex has the work of one less than its number of marks.
\end{theorem}

% Reference 20: thm:departures; source line 10421
\begin{theorem}[Repeated subtrees and departures]\label{thm:departures}
\begin{enumerate}[label=(\roman*)]
\item In a construction, replacing identical subtrees, with matching hinges and data, by one subtree and references leaves the family built, the period and every exact word of Part~\ref{part:construction} unchanged, and lowers the price by the price of the copies removed, less that of the references.
\item In (ii), (iii) and (iv) of Theorem~\ref{thm:pricerepeat}, replacing $q$ of the $M$ copies by cells $\mathcal A'_1,\dots,\mathcal A'_q$ changes the word to $\Gst21{\kappa,\ t-\kappa}{t}\big(P_{\mathcal A}/\Gst21{\kappa,\ t-\kappa}{t}\big)^{M-q}\prod_i\big(P_{\mathcal A'_i}/\Gst21{\kappa_i,\ t-\kappa_i}{t}\big)$ at a full lean, to the product of $\mathfrak L_{\mathcal A}^{M-q}$ and of the responses of the new cells in a series closure, and to the composite of $\mathcal Q^{M-q}$ with the $q$ operators of the new cells in a chain: the price grows by $\sum_i\operatorname{price}(\mathcal A'_i)+O(q\log M)$ and the work by at most the dimensions of the integrals of the new cells. A change of $q$ leaves in a tree with at most $h$ nodes on every path from its root to a leaf touches at most $qh$ of its nodes.
\end{enumerate}
\end{theorem}

Repetition is thus priced at the binary length of its count, as $2^M$ is written with $M$ rather than with its digits, and the work follows the distinct subtrees. The construction is the certificate of the compression: a family given expanded, without its construction, can still be read for its bunches of width one, its prime nodes (Theorem~\ref{thm:primetree}) and its twins (Theorem~\ref{thm:leafbounds}), but finding its cheapest construction is a search. The prime nodes of the canonical hinge tree of a list are invariants of the list (Theorem~\ref{thm:primetree}); at a full lean its period is a product of stacked Gammas and the periods at the full leans of those nodes (Corollary~\ref{cor:primelean}, Theorem~\ref{thm:leantree}), and at every lean a fibre integral of their periods at the full leans, of dimension $n-1$ (Theorem~\ref{thm:network}); a prime node $\mathcal H_v$ on $n_v$ marks with $|\mathcal H_v|$ bases has periods at the full leans of dimension $|\mathcal H_v|-n_v>0$ on their share sides (Corollary~\ref{cor:nodecost}). Symmetry is a second source of cheapness, distinct from repetition: a table assembled from a seed by a set of bijections (Proposition~\ref{prop:assembly}) is described by the seed and the bijections, and the table that is one circuit on $n$ marks, prime for no $n$, is a series closure of single marks whose period is one integral over one scale of a product of $n$ Grammas (Corollary~\ref{cor:cycle}), of work one whatever $n$; the list $\binom{\cE}r$ with $2\le r\le|\cE|-2$ is prime and cheap to describe, and its period at each full lean is one fibre integral of dimension $\binom{|\cE|}r-|\cE|$ on its share side (Corollary~\ref{cor:nodecost}).

\part{The arithmetic of the output}\label{part:arith}

\chapter{The motive of a family and the coaction along a tree}\label{ch:motive}

A family of lists at integer data whose proper walls are all at least one has as its period the period of a mixed motive: the relative cohomology of the toric variety of the braid fan, less the closures of the zero loci of its letters and of the linear letters, relative to its toric boundary. The walls are the orders of the form along the boundary divisors, and every boundary stratum is the product of the models of the steps of its flag. Brown's coaction acts on the motivic period, and birational changes of variables, Cremona among them, are identities of motivic periods. Along any tree of bunches the left factors of the coaction of weight at most $k$ are sums of products of open periods of steps of at most $k+1$ marks; each step is attached to one vertex of the tree, its letters are letters of children of the tables there, and its open periods are rational combinations of constant terms of continued periods of families of lists on the step, hence of convergent periods of such families and of their derivatives in the masses. At a node of width $c$ the bunch and its room meet only through the Skeleton values of the letters, at most $c$ for each letter: every open period of a node step is a shape integral of rational functions against finite parts of scale integrals, constant terms of stacked Gammas times Bretas of one bunch at the joint spectrum, and at the root so is every left factor; at width one each Skeleton value is a quotient of hinge values and the share side separates the bunch from the room. Cells read through one hinge value split off as identities of motivic periods, whatever the widths of the bunches of the core; so do the cells of dual pairs at every pair of exponents, and at a full lean the motivic period of a list is a rational multiple of the product of those of its prime nodes. A pinned set of at least two marks at which every active table has width one splits off motivically.

\section{The model}\label{sec:mot-model}

The motive of a family lives on the toric variety of the braid fan. Its charts are the sector cubes of Theorem~\ref{thm:exist}(ii), and its boundary strata are indexed by chains of sets of marks.

% New: proved in the last part.
\begin{definition}[The model of a family]\label{def:mot-model}
Let $F=((\cB_j;s_j,t_j))_{j=1}^N$ be a family of lists on $\cE$, $n=|\cE|\ge2$ (Definition~\ref{def:family}), with integer exponents and integer masses $\mvec$ on its balance wall; it is \emph{convergent} when $\kappa_S\ge1$ for every $\varnothing\neq S\subsetneq\cE$. Its \emph{letters} are $\operatorname{Let}(F)=\{K_j:\ s_j\neq0\}\cup\{\Psi_j:\ t_j\neq0\}$; the \emph{linear letters} are the $\alpha_J=\sum_{e\in J}\alpha_e$, $J\subseteq\cE$, the basis polynomials of the lists $\binom J1$; and the \emph{augmented letters} are $\operatorname{Let}^+(F)=\operatorname{Let}(F)\cup\{\alpha_J:\ J\subseteq\cE,\ |J|\ge2\}$. For an order $e_1,\dots,e_n$ of $\cE$ put
\[
\Omega_\cE=\sum_{i=1}^n(-1)^{i-1}\alpha_{e_i}\,d\alpha_{e_1}\wedge\cdots\wedge\widehat{d\alpha_{e_i}}\wedge\cdots\wedge d\alpha_{e_n},
\]
which another order changes by its sign; the \emph{form} of $F$ is
\[
\omega_F=\alpha^{\mvec}\prod_jK_j^{-s_j}\Psi_j^{-t_j}\,\frac{\Omega_\cE}{\alpha^{\cE}},
\]
a rational form of degree $n-1$ on $\mathbb P(\mathbb Q^{\cE})$. Let $X_\cE$ be the toric variety over $\mathbb Q$ of the \emph{braid fan} in $\mathbb R^{\cE}/\mathbb R\one$, whose rays are the images $\rho_S$ of the $\one_S$, $\varnothing\neq S\subsetneq\cE$, and whose cones are spanned by the rays of the chains $S_1\subsetneq\cdots\subsetneq S_\ell$; $D_S$ is the divisor of $\rho_S$, $D=\bigcup_SD_S$ the \emph{boundary}, and the \emph{positive cell} $\widetilde\sigma_\cE$ is the closure in $X_\cE(\mathbb R)$ of the set $\widetilde\sigma^\circ_\cE$ of the points of the torus with every $\alpha_e>0$. An order $e_1,\dots,e_n$ of $\cE$ gives the \emph{chart} $U\cong\mathbb A^{n-1}$ of the cone of the chain $S_i=\{e_1,\dots,e_i\}$, $i<n$, with the coordinates $\varepsilon_i=\alpha_{e_i}/\alpha_{e_{i+1}}$, so that $\alpha_{e_j}/\alpha_{e_n}=\prod_{i\ge j}\varepsilon_i$. For a polynomial $f$ on $\cE$, $\widetilde Y_f$ is the closure in $X_\cE$ of its zero locus in the torus. The \emph{model} of $F$ is the pair
\[
V_F=X_\cE\setminus\bigcup_{f\in\operatorname{Let}^+(F)}\widetilde Y_f,\qquad D_F=D\cap V_F,
\]
with the relative cohomology $H_F=H^{n-1}(V_F,D_F)$.
\end{definition}

% New: proved in the last part.
\begin{definition}[Flags and steps]\label{def:mot-flag}
A \emph{flag} of $\cE$ is a chain $\mathfrak f:\ \varnothing=S_0\subsetneq S_1\subsetneq\cdots\subsetneq S_\ell\subsetneq S_{\ell+1}=\cE$, $\ell\ge0$. Its \emph{steps} are the sets $Z_i=S_i\setminus S_{i-1}$, $1\le i\le\ell+1$, as in Theorem~\ref{thm:lettertransfer}, and its \emph{dimension} is $\sum_i(|Z_i|-1)=n-1-\ell$. In the model of a family $F$ (Definition~\ref{def:mot-model}) its \emph{stratum} is $D_{\mathfrak f}=D_{S_1}\cap\cdots\cap D_{S_\ell}\cap V_F$, the model $V_F$ itself at $\ell=0$. The \emph{letters of the step} $(Z_i;\mathfrak f)$ are the step letters at $Z_i$ (Proposition~\ref{prop:mot-chainbottom}) of the augmented letters of $F$ that are not monomials: $K$ of the step list of $\cB_j$ at $Z_i$ at the bottom of the chain for every $K_j\in\operatorname{Let}(F)$, $\Psi$ of its step list at the top for every $\Psi_j\in\operatorname{Let}(F)$, and the $\alpha_J$ with $J\subseteq Z_i$, $|J|\ge2$; for $\ell=0$ the step lists at the one step $\cE$ are the lists of $F$ themselves. They form a set $\operatorname{Let}(Z_i;\mathfrak f)$ of polynomials on $Z_i$, which depends on $S_{i-1}$ and $S_i$ alone. The \emph{model of the step} is
\[
V_{Z_i;\mathfrak f}=X_{Z_i}\setminus\bigcup_{g\in\operatorname{Let}(Z_i;\mathfrak f)}\widetilde Y_g,
\]
a point when $|Z_i|=1$, with its toric boundary $\partial V_{Z_i;\mathfrak f}$.
\end{definition}

% New: proved in the last part.
\begin{lemma}[Charts, strata and the letters in a chart]\label{lem:mot-charts}
In Definition~\ref{def:mot-model}:
\begin{enumerate}[label=(\roman*)]
\item $X_\cE$ is smooth and projective of dimension $n-1$, and the charts of the orders cover it. In the chart of the order $e_1,\dots,e_n$, $D_{S_i}\cap U=\{\varepsilon_i=0\}$ for $i<n$ and $D_S\cap U=\varnothing$ for every other $S$; $D$ is a divisor with simple normal crossings and smooth components, and $D_{S_1}\cap\cdots\cap D_{S_\ell}$ is not empty exactly when the $S_i$ form a chain. In the chart $\widetilde\sigma_\cE\cap U=[0,\infty)^{n-1}$, and the closure of the set of the points of $\widetilde\sigma^\circ_\cE$ with $\alpha_{e_1}<\cdots<\alpha_{e_n}$, the \emph{sector} of the order, is the cube $[0,1]^{n-1}$; the sectors of the orders cover $\widetilde\sigma_\cE$, two of them meeting in a set of measure zero.
\item Let $f$ be a letter $K$ or $\Psi$ of a list $\cB$ on $\cE$, or a linear letter $\alpha_J$, of the degree $\delta$, with the \emph{floor} $d_f(S)$: $d(S)$ for $K$, $d^\ast(S)$ for $\Psi$ (Theorem~\ref{thm:window}), and $1$ or $0$ for $\alpha_J$ as $J\subseteq S$ or not. In the chart of an order,
\[
f(\alpha)\big/\alpha_{e_n}^{\delta}=\prod_{i<n}\varepsilon_i^{\,d_f(S_i)}\,\widetilde f(\varepsilon),
\]
with $\widetilde f$ a polynomial with nonnegative integer coefficients and constant term $1$, the sector letter of Theorem~\ref{thm:exist}(ii); so $\widetilde f\ge1$ where $\varepsilon\ge0$. For $I\subseteq\{1,\dots,n-1\}$, the restriction of $\widetilde f$ to $\{\varepsilon_i=0:\ i\in I\}$ is the product, over the steps $Z$ of the chain $(S_i)_{i\in I}$, of the polynomials that this rule attaches to the step letters of $f$ at $Z$ (Proposition~\ref{prop:mot-chainbottom}) in the chart of $X_Z$ of the order induced on $Z$, whose coordinates are the $\varepsilon_k$ with $e_k,e_{k+1}\in Z$. Moreover $\widetilde Y_f\cap U=\{\widetilde f=0\}$.
\item In the chart, with the walls $\kappa_S$ of $F$,
\[
\omega_F=\pm\prod_{i<n}\varepsilon_i^{\,\kappa_{S_i}-1}\prod_j\widetilde K_j^{\,-s_j}\widetilde\Psi_j^{\,-t_j}\,d\varepsilon_1\wedge\cdots\wedge d\varepsilon_{n-1}.
\]
\end{enumerate}
\end{lemma}

% New: proved in the last part.
\begin{proposition}[The model]\label{prop:mot-model}
Let $F$ be a family as in Definition~\ref{def:mot-model}.
\begin{enumerate}[label=(\roman*)]
\item $\widetilde\sigma_\cE$ is compact, $\widetilde\sigma_\cE\setminus\widetilde\sigma^\circ_\cE\subseteq D$, and $\widetilde\sigma_\cE$ misses $\widetilde Y_f$ for every letter $f$ of a list on $\cE$ and every linear letter; so $\widetilde\sigma_\cE\subseteq V_F(\mathbb R)$.
\item $\omega_F$ has the order $\kappa_S-1$ along $D_S$ for every $\varnothing\neq S\subsetneq\cE$ and neither zeros nor poles on the points of the torus in $V_F$. So $\omega_F$ is a regular form on $V_F$ exactly when $F$ is convergent, and then its integral over $\widetilde\sigma_\cE$, oriented so that it is positive, converges and is the period $P_F(\mvec;\mathbf s,\mathbf t)$ (Theorem~\ref{thm:projective}).
\item $V_F$ is a smooth affine variety, $D_F$ is a divisor in it with simple normal crossings and smooth components, and every stratum is smooth and affine. More generally $X_\cE\setminus\bigcup_{g\in\mathcal G}\widetilde Y_g$ is affine for every finite set $\mathcal G$ of polynomials on $\cE$ that contains the linear letters $\alpha_J$, $|J|\ge2$.
\item The stratum of a flag $\mathfrak f$ with the steps $Z_1,\dots,Z_{\ell+1}$ is isomorphic to $\prod_iV_{Z_i;\mathfrak f}$ (Definition~\ref{def:mot-flag}), compatibly with the boundaries, and $\widetilde\sigma_\cE$ meets it in $\prod_i\widetilde\sigma_{Z_i}$.
\item For the tables of $F$ and complex masses and exponents on the balance wall, wherever every proper wall has positive real part,
\[
P_F(\mvec;\mathbf s,\mathbf t)=\sum_\pi\operatorname{cube}\big(\boldsymbol\kappa_\pi(\mvec);h_\pi\big),\qquad\boldsymbol\kappa_\pi(\mvec)=\big(\kappa_{S^\pi_i}(\mvec)\big)_{i<n},\quad h_\pi=\prod_j\widetilde K_j^{\,-s_j}\widetilde\Psi_j^{\,-t_j},
\]
the sum over the orders $\pi$ of $\cE$, with the chains $S^\pi_i$ and the sector letters of their charts, and $\operatorname{cube}$ the cube integral of Lemma~\ref{lem:cube}. The right side continues meromorphically along the balance wall, jointly in the exponents, with poles only on the hyperplanes $\kappa_S=-j$, $\varnothing\neq S\subsetneq\cE$, $j\in\mathbb N$: it is the \emph{continued period} of Theorems~\ref{thm:exist}(iii) and~\ref{thm:projective}.
\end{enumerate}
\end{proposition}

The walls are thus the orders of the form along the boundary, and convergence is regularity. The linear letters are removed so that the model is affine; they are the letters of the lists $\binom J1$, and the motivic period does not see them (Lemma~\ref{lem:mot-birational}).

\section{Motivic periods}\label{sec:mot-periods}

% New: proved in the last part.
\begin{definition}[Motivic periods; Galois conjugates; open spaces]\label{def:mot-mperiod}
An object $M$ of Brown's category of realizations over $\mathbb Q$ (F.~Brown, \emph{Notes on motivic periods}, Commun.\ Number Theory Phys.\ 11 (2017), 557--655, \S3.1) has a Betti realization $M_{\mathrm B}$ and a de Rham realization $M_{\mathrm{dR}}$; its \emph{motivic periods} $[M,\gamma,\omega]^{\mathrm{mot}}$, $\gamma\in M_{\mathrm B}^\vee$ and $\omega\in M_{\mathrm{dR}}$, span the ring $\mathcal P^{\mathrm{mot}}$, bilinear in $(\gamma,\omega)$, with $[M,\phi^{\mathsf t}\gamma',\omega]^{\mathrm{mot}}=[M',\gamma',\phi(\omega)]^{\mathrm{mot}}$ for every morphism $\phi:M\to M'$; its de Rham periods $[M,f,\omega]^{\mathrm{dR}}$, $f\in M_{\mathrm{dR}}^\vee$, span the ring $\mathcal P^{\mathrm{dR}}$; and the \emph{coaction} $\Delta:\mathcal P^{\mathrm{mot}}\to\mathcal P^{\mathrm{mot}}\otimes\mathcal P^{\mathrm{dR}}$ is
\[
\Delta[M,\gamma,\omega]^{\mathrm{mot}}=\sum_i[M,\gamma,e_i]^{\mathrm{mot}}\otimes[M,e_i^\vee,\omega]^{\mathrm{dR}}
\]
for a basis $(e_i)$ of $M_{\mathrm{dR}}$ with the dual basis $(e_i^\vee)$; the $[M,\gamma,e_i]^{\mathrm{mot}}$ are its \emph{left factors}. For $\xi\in\mathcal P^{\mathrm{mot}}$ the space of its \emph{Galois conjugates} is
\[
\mathrm C(\xi)=\big\{(\mathrm{id}\otimes\varphi)\Delta\xi:\ \varphi\in\operatorname{Hom}_{\mathbb Q}(\mathcal P^{\mathrm{dR}},\mathbb Q)\big\};
\]
for an object $M$ and $\gamma\in M_{\mathrm B}^\vee$, $L(M,\gamma)=\{[M,\gamma,\omega]^{\mathrm{mot}}:\ \omega\in M_{\mathrm{dR}}\}$; and $\mathcal A\cdot\mathcal A'$ is the span of the products of the elements of two subspaces. For a convergent family $F$ (Definition~\ref{def:mot-model}) the \emph{motivic period} is
\[
P^{\mathrm{mot}}_F=[H_F,\widetilde\sigma_\cE,\omega_F]^{\mathrm{mot}},
\]
the positive cell, oriented so that $\Omega_\cE/\alpha^\cE$ is positive on $\widetilde\sigma^\circ_\cE$, read as a class of relative homology, and the form as a class of relative de Rham cohomology; its period is $P_F(\mvec;\mathbf s,\mathbf t)$ (Proposition~\ref{prop:mot-model}(ii)). For a step $(Z;\mathfrak f)$ of a flag of $\cE$ (Definition~\ref{def:mot-flag}) with $|Z|\ge2$, its \emph{open space} is the set of motivic periods
\[
\operatorname{Open}(Z;\mathfrak f)=\Big\{\big[H^{|Z|-1}(V_{Z;\mathfrak f},\partial V_{Z;\mathfrak f}),\widetilde\sigma_Z,\eta\big]^{\mathrm{mot}}:\ \eta\ \text{a regular form of degree }|Z|-1\text{ on }V_{Z;\mathfrak f}\Big\},
\]
$\widetilde\sigma_Z$ oriented as $\widetilde\sigma_\cE$ is, and $\operatorname{Open}(Z;\mathfrak f)=\mathbb Q$ when $|Z|=1$.
\end{definition}

% New: proved in the last part.
\begin{lemma}[Galois conjugates]\label{lem:mot-conj}
For $\xi,\xi'\in\mathcal P^{\mathrm{mot}}$, an object $M$ and $\gamma\in M_{\mathrm B}^\vee$:
\begin{enumerate}[label=(\roman*)]
\item $\xi\in\mathrm C(\xi)$;
\item $\mathrm C([M,\gamma,\omega]^{\mathrm{mot}})\subseteq L(M,\gamma)$, and $\Delta L(M,\gamma)\subseteq L(M,\gamma)\otimes\mathcal P^{\mathrm{dR}}$;
\item $\mathrm C(\xi+\xi')\subseteq\mathrm C(\xi)+\mathrm C(\xi')$, $\mathrm C(\xi\xi')\subseteq\mathrm C(\xi)\cdot\mathrm C(\xi')$, and $\mathrm C(q\,\xi)=\mathrm C(\xi)$ for $q\in\mathbb Q$, $q\neq0$.
\end{enumerate}
\end{lemma}

% New: proved in the last part.
\begin{lemma}[Birational invariance]\label{lem:mot-birational}
For $i=1,2$ let $X_i$ be a smooth projective variety over $\mathbb Q$ of dimension $m$, $D_i\subset X_i$ a divisor with simple normal crossings and smooth components, $Y_i\subset X_i$ a closed set, $X_i^\circ=X_i\setminus Y_i$ and $D_i^\circ=D_i\setminus Y_i$. Let $\sigma_i\subset X_i(\mathbb R)$ be compact, the closure of an oriented open semialgebraic set $\sigma_i^\circ\subseteq(X_i\setminus(D_i\cup Y_i))(\mathbb R)$ of dimension $m$, with $\sigma_i\cap Y_i=\varnothing$ and $\sigma_i\setminus\sigma_i^\circ\subseteq D_i$, and let $\omega_i$ be a regular form of degree $m$ on $X_i^\circ$. If an isomorphism $\Phi:U_1\to U_2$ of dense open sets $U_i\subseteq X_i\setminus(D_i\cup Y_i)$ with $\sigma_i^\circ\subseteq U_i(\mathbb R)$ maps $\sigma_1^\circ$ onto $\sigma_2^\circ$ preserving the orientations and has $\Phi^\ast\omega_2=\omega_1$, then
\[
\big[H^m(X_1^\circ,D_1^\circ),\sigma_1,\omega_1\big]^{\mathrm{mot}}=\big[H^m(X_2^\circ,D_2^\circ),\sigma_2,\omega_2\big]^{\mathrm{mot}}.
\]
In particular, with $\Phi$ the identity, removing a larger closed set that misses $\sigma_1$ leaves the motivic period unchanged; the motivic period of a convergent family is the same in the model without the linear letters, $X_\cE\setminus\bigcup_{f\in\operatorname{Let}(F)}\widetilde Y_f$; and two convergent families on $\cE$ with one form have one motivic period.
\end{lemma}

% New: proved in the last part.
\begin{corollary}[Cremona between motivic periods]\label{cor:mot-cremona}
Let $F=((\cB_j;s_j,t_j))_j$ be a convergent family with the masses $\mvec$ and the total strength $\sigma$, and $F^{\mathrm c}$ the family of the same tables with the exponents $(t_j,s_j)$ and the masses $2\sigma\one-\mvec$. Then $F^{\mathrm c}$ is convergent, its wall at $S$ is the wall of $F$ at $\cE\setminus S$, and $P^{\mathrm{mot}}_{F^{\mathrm c}}=P^{\mathrm{mot}}_F$: Theorem~\ref{thm:cremona} is an identity of motivic periods.
\end{corollary}

\section{Steps along a tree of bunches}\label{sec:mot-steps}

The boundary strata are indexed by flags, and a stratum is a product over the steps of its flag (Proposition~\ref{prop:mot-model}(iv)). A tree of bunches (Definition~\ref{def:bunchtree}) assigns every step to one of its vertices.

% New: proved in the last part.
\begin{definition}[Attachment]\label{def:mot-attach}
Let $\mathcal T$ be a tree of bunches of a family $F$ (Definition~\ref{def:bunchtree}) and $(Z;\mathfrak f)$ a step of a flag of $\cE$. Starting at the root, $Z$ descends through vertices whose marks contain it:
\begin{enumerate}[label=(\roman*)]
\item at a split vertex, if $Z$ lies in one of its children it descends to that child; otherwise $Z$ is \emph{attached} there as a \emph{node step};
\item at a node $v$ that is not a leaf, with its bunch $B_v$: if $Z\subseteq B_v$, $Z$ is attached to $v$ as a \emph{bunch step}; if $Z$ meets $B_v$ and $M_v\setminus B_v$, as a \emph{node step}; and if $Z\subseteq M_v\setminus B_v$, it descends to the room of $v$;
\item at a leaf, $Z$ is attached as a \emph{leaf step}.
\end{enumerate}
Every step is attached to exactly one vertex.
\end{definition}

% New: proved in the last part.
\begin{lemma}[The letters of an attached step]\label{lem:mot-attach}
Let $\mathcal T$ be a tree of bunches of a family $F$ as in Definition~\ref{def:mot-model}, and let the step $(Z;\mathfrak f)$ be attached to the vertex $v$, with the marks $M_v$. Every letter of the step other than the linear letters is $K$ or $\Psi$ of the child, at a pattern of $M_v\setminus Z$ (Definition~\ref{def:patterns}), of a table at $v$ that carries that letter.
\end{lemma}

Every family has trees of bunches: at every node that is not a leaf choose a bunch and pass to its room, which has fewer marks. For nonnegative exponents the Foreman tree is one of them (Definition~\ref{def:tree}); with nonnegative exponents merging equal children keeps every letter, so the record and the peel give the same tree, and it is fixed by the tables and the rays of the exponent pairs (Theorem~\ref{thm:threetrees}). The statements below hold along every tree of bunches, and along the Foreman tree the vertex at which each step is attached, and the node form there, are fixed by the family.

\section{The weights of the coaction}\label{sec:mot-small}

% New: proved in the last part.
\begin{lemma}[One form carries every period]\label{lem:mot-openall}
Let $V$ be a smooth affine variety over $\mathbb Q$ of dimension $m$ and $D\subset V$ a divisor with simple normal crossings and smooth components. Every class of the relative de Rham cohomology $H^m_{\mathrm{dR}}(V,D)$ is the class of a regular form of degree $m$ on $V$. So, for a convergent family $F$ (Definition~\ref{def:mot-model}), $L(H_F,\widetilde\sigma_\cE)$ is the open space of the one step $\cE$ of the flag $\varnothing\subsetneq\cE$, and for every step, $L\big(H^{|Z|-1}(V_{Z;\mathfrak f},\partial V_{Z;\mathfrak f}),\widetilde\sigma_Z\big)=\operatorname{Open}(Z;\mathfrak f)$.
\end{lemma}

So the open space of the whole set of marks already holds every period of $H_F$ on its cell, and a decomposition of the left factors into products over steps carries information only when it is graded. The steps enter through the weight filtration $W$ of the de Rham realization, by Brown's weight-graded face maps.

% New: proved in the last part.
\begin{proposition}[Weight-graded faces]\label{prop:mot-weights}
Let $F$ be convergent (Definition~\ref{def:mot-model}), $k<n-1$ and $e\in W_k(H_F)_{\mathrm{dR}}$. Then
\[
[H_F,\widetilde\sigma_\cE,e]^{\mathrm{mot}}\in\sum_{\mathfrak f\ \text{of dimension }k}\ \prod_{\text{steps }Z\text{ of }\mathfrak f}\operatorname{Open}(Z;\mathfrak f).
\]
\end{proposition}

% New: proved in the last part.
\begin{theorem}[Small steps along the tree]\label{thm:mot-smallsteps}
Let $F$ be convergent, $\mathcal T$ a tree of bunches of $F$, and $(e_i)$ a basis of $(H_F)_{\mathrm{dR}}$ adapted to the weight filtration $W$. For $k<n-1$ every left factor $[H_F,\widetilde\sigma_\cE,e_i]^{\mathrm{mot}}$ of the coaction of $P^{\mathrm{mot}}_F$ with $e_i\in W_k$ satisfies
\[
[H_F,\widetilde\sigma_\cE,e_i]^{\mathrm{mot}}\in\sum_{\mathfrak f\ \text{of dimension }k}\ \prod_{\text{steps }Z\text{ of }\mathfrak f}\operatorname{Open}(Z;\mathfrak f).
\]
Only steps of at most $k+1$ marks occur. Each is attached to exactly one vertex $v$ of $\mathcal T$, and every letter of it other than the linear letters is $K$ or $\Psi$ of the child, at a pattern of $M_v\setminus Z$, of a table at $v$ that carries that letter.
\end{theorem}

The theorem is read on Brown's affine model, which removes the linear letters $\alpha_J$ to make the complement affine. They are the letters of the lists $\binom J1$, and they enter the families of the steps below as tables. Each open period is a constant term of continued periods, and those constant terms are read on convergent data.

% New: proved in the last part.
\begin{proposition}[Laurent coefficients are read on convergent data]\label{prop:mot-convergentdata}
Let $F=((\cB_j;s_j,t_j))_{j=1}^N$ be a family of lists on $\cE$ with integer exponents and integer masses $\mvec$ on its balance wall, and $\mathbf v\in\mathbb Q^{\cE}$ with $v_{\cE}=0$ and $v_S\neq0$ for every $\varnothing\neq S\subsetneq\cE$. Then $\varepsilon\mapsto P_F(\mvec+\varepsilon\mathbf v;\mathbf s,\mathbf t)$, the continued period (Proposition~\ref{prop:mot-model}(v)), is meromorphic near $\varepsilon=0$, and there are finitely many convergent families $F_k=((\cB_j;s^{(k)}_j,t^{(k)}_j))_j$ on the same tables, with $\mathbf s^{(k)}-\mathbf s,\ \mathbf t^{(k)}-\mathbf t\in\mathbb N^N$ and integer masses $\mvec^{(k)}$ on their balance walls, and rational functions $R_k\in\mathbb Q(\varepsilon)$, with
\[
P_F(\mvec+\varepsilon\mathbf v;\mathbf s,\mathbf t)=\sum_kR_k(\varepsilon)\,P_{F_k}\big(\mvec^{(k)}+\varepsilon\mathbf v;\mathbf s^{(k)},\mathbf t^{(k)}\big)
\]
near $\varepsilon=0$. So every coefficient of the Laurent expansion at $\varepsilon=0$ is a rational combination of the absolutely convergent integrals
\[
\partial_\varepsilon^{\,l}P_{F_k}\big(\mvec^{(k)}+\varepsilon\mathbf v;\mathbf s^{(k)},\mathbf t^{(k)}\big)\Big|_{\varepsilon=0}=\int_{\Delta_{\cE}}u^{\mvec^{(k)}-\one}\ip{\mathbf v}{\log u}^l\prod_jK_j(u)^{-s^{(k)}_j}\Psi_j(u)^{-t^{(k)}_j}\,d\omega_{\cE}(u),
\]
$l\in\mathbb N$. If $\mathbf s,\mathbf t\ge0$, every $F_k$ has nonnegative exponents and positive masses.
\end{proposition}

% New: proved in the last part.
\begin{corollary}[Open periods from families of lists]\label{cor:mot-openfamilies}
Let $(Z;\mathfrak f)$ be a step of a flag of $\cE$ with $|Z|\ge2$ and $\eta$ a regular form of degree $|Z|-1$ on $V_{Z;\mathfrak f}$. On the torus,
\[
\eta=\sum_tc_t\,\omega_{F_t},\qquad c_t\in\mathbb Q,
\]
a finite sum over families $F_t$ of lists on $Z$ at integer data on their balance walls: the tables of $F_t$ are the lists of the letters of the step, the step list for its $K$ or its $\Psi$ and $\binom J1$ for $\alpha_J$, each carrying its letter with an exponent $k_{t,g}\in\mathbb N$, and the masses are some $\mathbf a_t\in\mathbb Z^Z$. For every $\mathbf v\in\mathbb Q^Z$ with $v_Z=0$ and $v_S\neq0$ for every $\varnothing\neq S\subsetneq Z$, each continued period $P_{F_t}(\mathbf a_t+\varepsilon\mathbf v)$ (Proposition~\ref{prop:mot-model}(v)) is meromorphic in $\varepsilon$ near $0$, and
\[
\int_{\widetilde\sigma_Z}\eta=\sum_tc_t\,\CT_{\varepsilon=0}P_{F_t}(\mathbf a_t+\varepsilon\mathbf v).
\]
By Proposition~\ref{prop:mot-convergentdata} each constant term is a rational combination of convergent periods of families of lists on $Z$ with nonnegative integer exponents and positive integer masses, values of the serif family (Proposition~\ref{prop:gammadraws}), and of their derivatives along $\mathbf v$ in the masses. So every element of $\operatorname{Open}(Z;\mathfrak f)$ has such a number as its period, and by Theorem~\ref{thm:mot-smallsteps} every left factor of weight at most $k$ has as its period a sum of products of such numbers over the steps of the flags of dimension $k$, whose letters are letters of children of the tables at the vertices where the steps are attached.
\end{corollary}

\section{The node}\label{sec:mot-node}

% New: proved in the last part.
\begin{theorem}[The node form]\label{thm:mot-node}
Let $F$ be convergent (Definition~\ref{def:mot-model}), $\mathcal T$ a tree of bunches of $F$, $v$ a node of $\mathcal T$ that is not a leaf, with its bunch $B_v$ of width $c$ (Definition~\ref{def:bunchtree}), and $(Z;\mathfrak f)$ a node step of $v$ (Definition~\ref{def:mot-attach}), with $Z_B=Z\cap B_v$ and $Z_R=Z\cap(M_v\setminus B_v)$. Read the bunch side in scale and shape and the room side on its simplex, $\alpha_{Z_B}=\rho u$ with $\rho>0$ and $u\in\Delta_{Z_B}$, and $\alpha_{Z_R}=u'\in\Delta_{Z_R}$ (Lemma~\ref{lem:scaleshape}).
\begin{enumerate}[label=(\roman*)]
\item Every letter $g$ of the step is the basis polynomial of a list $\mathcal N_g$ on $Z$: the step list at the bottom for a letter $K$, the dual of the step list at the top for a letter $\Psi$ (Proposition~\ref{prop:identity}), and $\binom J1$, its marks outside $J$ loops, for $\alpha_J$. The width $c_g$ of $Z_B$ in $\mathcal N_g$ is at most $c$, and $c\ge1$.
\item With $d_g$ the floor of $Z_B$ in $\mathcal N_g$ and the Skeleton values $\xi_{g,1},\dots,\xi_{g,c_g}$ of $K_{\mathcal N_g}$ along $Z_B$ at $(u,u')$ (Definition~\ref{def:skelpoly}),
\[
g(\rho u,u')=\rho^{d_g}\,K_{\mathcal N_g|Z_B}(u)\,K_{\mathcal N_g/Z_B}(u')\prod_{i=1}^{c_g}\big(\rho+\xi_{g,i}(u,u')\big),
\]
the leading coefficient being the top level of $Z_B$ (Theorem~\ref{thm:levels}); these Skeleton values, each with its letter, form the joint spectrum of $Z_B$ at $(u,u')$ (Definition~\ref{def:pairsk}).
\item For every regular form $\eta$ of degree $|Z|-1$ on $V_{Z;\mathfrak f}$ there are finitely many terms $t$, with rational functions $A_t$ on $\Delta_{Z_B}$ and $B_t$ on $\Delta_{Z_R}$, integers $\nu_t$ and letter exponents $k_{t,g}\in\mathbb N$, such that
\[
\int_{\widetilde\sigma_Z}\eta=\int_{\Delta_{Z_B}\times\Delta_{Z_R}}\sum_tA_t(u)\,B_t(u')\,\mathrm{Sc}_t(u,u')\,d\omega_{Z_B}(u)\,d\omega_{Z_R}(u'),
\]
the integrand absolutely integrable. Here $\mathrm{Sc}_t=0$ when no letter with $k_{t,g}>0$ has a Skeleton value, and otherwise
\begin{align*}
\mathrm{Sc}_t&=\Reg\!\int_0^\infty\rho^{\nu_t-1}\prod_{g,i}\big(\rho+\xi_{g,i}\big)^{-k_{t,g}}\,d\rho\\
&=\CT_{\kappa=\nu_t}\ \Gst{2}{|J_t|}{\kappa,\ |\mathbf w_t|-\kappa}{\mathbf w_t}\,\sB^{\{J_t\}}_{\mathbf w_t,\,|\mathbf w_t|-\kappa;\,I_{J_t}}\big(\boldsymbol\xi(u,u')\big),
\end{align*}
$J_t$ indexing the Skeleton values $\xi_{g,i}$ of the letters with $k_{t,g}>0$, each with the letter exponent $w_{t,(g,i)}=k_{t,g}$, so that $|\mathbf w_t|=\sum_gk_{t,g}c_g$: the finite part in $\rho$ (Definition~\ref{def:Reg}) is the constant term at $\nu_t$ of the stacked Gamma times the Breta of one bunch of Proposition~\ref{prop:scalebreta}, continued from the strip $0<\re\kappa<|\mathbf w_t|$. This is the \emph{node form} of $\eta$ at $v$: the bunch and the room meet only through the joint spectrum, as the coefficients of Bretas of one bunch at the unit configuration.
\item (\emph{Width one.}) If $c=1$, every letter has at most one Skeleton value, the room's hinge value over the bunch's, $\xi_g=\gamma^g_R(u')/\gamma^g_B(u)$ with $\gamma^g_B=K_{\mathcal N_g|Z_B}/K_{\mathcal N_g/Z_R}$ and $\gamma^g_R=K_{\mathcal N_g|Z_R}/K_{\mathcal N_g/Z_B}$ (Corollary~\ref{cor:widthonestage}), positive. For $\kappa$ in the strip, the share side (Corollary~\ref{cor:sharepercoefficient}) writes the stacked Gamma times the Breta of the term $t$ as an integral over $|J_t|-1$ shares of stacked Gammas times powers $\xi_g^{-b_g}=\gamma^g_R(u')^{-b_g}\gamma^g_B(u)^{\,b_g}$, the $b_g$ affine in the shares and in $\kappa$, with real parts constant along the contour. On every such contour on which these stacked Gammas have arguments with positive real parts and
\[
\int_{\Delta_{Z_B}}|A_t|\prod_g(\gamma^g_B)^{\re b_g}\,d\omega_{Z_B}<\infty,\qquad\int_{\Delta_{Z_R}}|B_t|\prod_g(\gamma^g_R)^{-\re b_g}\,d\omega_{Z_R}<\infty,
\]
the integral over $\Delta_{Z_B}\times\Delta_{Z_R}$ of $A_tB_t$ times the stacked Gamma times the Breta is the share integral of the stacked Gammas times $\int_{\Delta_{Z_B}}A_t\prod_g(\gamma^g_B)^{b_g}d\omega_{Z_B}$ times $\int_{\Delta_{Z_R}}B_t\prod_g(\gamma^g_R)^{-b_g}d\omega_{Z_R}$: the bunch and the room separate. When $0<\nu_t<|\mathbf w_t|$, $\mathrm{Sc}_t$ is the convergent scale integral at $\kappa=\nu_t$, and the term itself separates.
\item If the root of $\mathcal T$ is a node that is not a leaf, the period of every element of $L(H_F,\widetilde\sigma_\cE)$, every left factor of the coaction of $P^{\mathrm{mot}}_F$ among them, is given by the node form at the root.
\item At a split vertex and a node step $Z$ of it that meets the child $M_1$, the same holds with $Z_B=Z\cap M_1$, $Z_R=Z\setminus M_1$ and $c=1$: the letters other than the linear letters have width zero at $Z_B$, and only the $\alpha_J$ with $J$ meeting $Z_B$ and $Z_R$ have a Skeleton value.
\end{enumerate}
\end{theorem}

At width two or more the Skeleton values are the negatives of the roots of Skeleton polynomials of degree $c_g$, whose extreme coefficients are products of a bunch letter and a room letter and whose middle coefficients are sums over the patterns of their levels (Theorem~\ref{thm:levels}); the node form reads them as the coefficients of a Breta of one bunch at any width, while hinge values are defined at width one only (Definition~\ref{def:readings}). For the period itself, $\eta=\omega_F$ at the root, the glue reads the bunch and the room through their two laws (Theorem~\ref{thm:gluewidth}(ii)), opened with one share per station (Theorem~\ref{thm:gluewidth}(iii)) or with one more pin for every direction a side does not see (Theorem~\ref{thm:transversepins}); at width one in one table this is Theorem~\ref{thm:twosides}, and with one Skeleton value Corollary~\ref{cor:pinnedproduct}. Theorem~\ref{thm:mot-node} extends the node form from the period to every form on every node step, and so to every left factor; width one enters it only in (iv).

\section{Pinned cells}\label{sec:mot-pinned}

% New: proved in the last part.
\begin{theorem}[Pinned cells split off motivically]\label{thm:mot-pinned}
Let $F=((\cB_j;s_j,t_j))_{j=1}^N$ be a convergent family (Definition~\ref{def:mot-model}) on $\cE=L\sqcup R_1\sqcup\cdots\sqcup R_q$, $q\ge1$, $|R_i|\ge2$, whose active tables are substituted table by table (Definition~\ref{def:substitution}): $\cB_j=\mathcal H_j\langle\mathcal A_{1j},\dots,\mathcal A_{qj}\rangle$, the marks of $L$ left in place, with the cores $\mathcal H_j$, lists on $L\sqcup\{h_1,\dots,h_q\}$, and the cells $\mathcal A_{ij}$, lists on $R_i\sqcup\{g_i\}$; here and below $j$ runs over the active tables, the others carrying no letter. Let the active letters read every cell through one hinge value: for every $i$ there is a function $\gamma_i$ on $\Delta_{R_i}$ with $\gamma^K_{\mathcal A_{ij}}=\gamma_i$ whenever $s_j\neq0$ and $\gamma^\Psi_{\mathcal A_{ij}}=\gamma_i$ whenever $t_j\neq0$. Let $\kappa_i=\kappa_{R_i}$, the wall of the cell (Theorem~\ref{thm:substwalls}), put $2\sigma=\sum_j(s_j+t_j)$, and fix integer loads $a_{ij},b_{ij}$ with $a_{ij}=0$ when $s_j=0$, $b_{ij}=0$ when $t_j=0$ and $\sum_j(a_{ij}+b_{ij})=\kappa_i$.
\begin{enumerate}[label=(\roman*)]
\item $1\le\kappa_i\le2\sigma-1$. The \emph{core family} $F_{\mathcal H}=((\mathcal H_j;s_j,t_j))_j$ on $L\sqcup\{h_1,\dots,h_q\}$, with the masses $\mvec|_L$ and $\kappa_i$ at $h_i$, is convergent, and so is, for every $i$, the \emph{pinned family} $G_i$ on $R_i$, the load family at zero imbalance of Corollary~\ref{cor:mot-loadpinned}: the tables $\mathcal A_{ij}\setminus g_i$ at $(a_{ij},t_j-b_{ij})$ and $\mathcal A_{ij}/g_i$ at $(s_j-a_{ij},b_{ij})$, with the masses $\mvec|_{R_i}$.
\item The motivic period factors through the core:
\[
P^{\mathrm{mot}}_F=P^{\mathrm{mot}}_{F_{\mathcal H}}\prod_{i=1}^qP^{\mathrm{mot}}_{G_i},\qquad\mathrm C(P^{\mathrm{mot}}_F)\subseteq\mathrm C(P^{\mathrm{mot}}_{F_{\mathcal H}})\cdot\prod_{i=1}^q\mathrm C(P^{\mathrm{mot}}_{G_i}).
\]
\item The \emph{cell family} $F_i=((\mathcal A_{ij};s_j,t_j))_j$ on $R_i\sqcup\{g_i\}$, with the masses $\mvec|_{R_i}$ and $2\sigma-\kappa_i$ at $g_i$, is convergent, and
\[
P^{\mathrm{mot}}_{F_i}=\Gst21{\kappa_i,\ 2\sigma-\kappa_i}{2\sigma}\,P^{\mathrm{mot}}_{G_i},
\]
the Beta corner a positive rational number.
\item The form, the period and the motivic period of $G_i$ do not depend on the loads: its form is $\gamma_i^{-\kappa_i}\alpha^{\mvec|_{R_i}}\prod_jK_{\mathcal A_{ij}/g_i}^{-s_j}\Psi_{\mathcal A_{ij}\setminus g_i}^{-t_j}\,\Omega_{R_i}/\alpha^{R_i}$, and its period is the pinned shape integral of the cell (Theorem~\ref{thm:pinned}).
\item Every $S\subseteq L$ has in $\cB_j$ the width it has in $\mathcal H_j$ (Corollary~\ref{cor:mot-widthbound}): a bunch among the marks left in place keeps the width it has in the core, and no condition is placed on it.
\end{enumerate}
\end{theorem}

One hinge value confines the set it reads to width at most one (Lemma~\ref{lem:mot-onehinge}), and the marks of a cell always form a set of width one (Proposition~\ref{prop:substitution}(i)); so it is the cells that have width one, and a bunch among the marks left in place has any width up to the number of the other marks and cells (Corollary~\ref{cor:mot-widthbound}). With $q$ cells the letters read the room through $q$ hinge values, one for each cell.

% New: proved in the last part.
\begin{corollary}[Dual pairs and full leans]\label{cor:mot-dualfull}
Let $F=(\cB;s,t)$ be a convergent family of one list (Definition~\ref{def:mot-model}).
\begin{enumerate}[label=(\alph*)]
\item (\emph{Dual pairs.}) Let $\cB=\mathcal H\langle\mathcal A_1,\dots,\mathcal A_q\rangle$ as in Theorem~\ref{thm:mot-pinned}, with every cell a dual pair (Definition~\ref{def:substitution}). Then every cell is read through one hinge value at all integers $s,t$; the pinned family of the cell $i$ has the form of the family $(\cB|R_i;\ \kappa_i,\ s+t-\kappa_i)$ at the masses $\mvec|_{R_i}$, with $\kappa_i=m_{R_i}-(s+t)(r(R_i)-1)$; and
\[
P^{\mathrm{mot}}_F=P^{\mathrm{mot}}_{(\mathcal H;s,t)}\prod_{i=1}^qP^{\mathrm{mot}}_{(\cB|R_i;\,\kappa_i,\,s+t-\kappa_i)},
\]
the core at the masses $\mvec|_L$ and $\kappa_i$ at $h_i$.
\item (\emph{Full leans.}) Let $t=0$ or $s=0$. Split $\cB$ at a set of width one with at least two marks on each side into the substitution of its two cells (Proposition~\ref{prop:substitution}(iii)), and repeat at the nodes so obtained until no node has such a set. The nodes on at least four marks are then prime (Theorem~\ref{thm:primetree}), and those on two or three marks are the lists $\binom S1$ and $\binom S{|S|-1}$. Give each prime node $\mathcal H_v$ the family $F_v=(\mathcal H_v;s,t)$ with the masses of its marks and, at each marker, the wall of $F$ at the marks on the other side of the split that made it. Every $F_v$ is convergent, and
\[
P^{\mathrm{mot}}_F=c\prod_vP^{\mathrm{mot}}_{F_v},\qquad c\in\mathbb Q_{>0};
\]
so $\mathrm C(P^{\mathrm{mot}}_F)\subseteq\prod_v\mathrm C(P^{\mathrm{mot}}_{F_v})$, and a list without prime nodes has a rational motivic period at its full leans.
\end{enumerate}
\end{corollary}

Part (b) is the motivic form of Corollary~\ref{cor:primelean} and Theorem~\ref{thm:leantree}: the series and parallel nodes contribute stacked Gammas at integer arguments (Proposition~\ref{prop:spleans}), which are rational, and their motivic periods are these rationals. At a full lean the width of a bunch plays no part, and the two full leans of one list are exchanged by Cremona, as motivic periods (Corollary~\ref{cor:mot-cremona}). In a prime list no set with two marks on each side splits off (Corollary~\ref{cor:mot-whatsplits}(ii)), so the splitting of (b) ends at the prime nodes; their only cells are the complements of single marks, which at a full lean split off by one Beta corner and, with both letters active, exactly where the deletion and the contraction have one even letter (Corollary~\ref{cor:mot-whatsplits}(iii),(iv)).

% New: proved in the last part.
\begin{theorem}[A pinned set splits off motivically]\label{thm:mot-splitoff}
Let $F=((\cB_j;s_j,t_j))_j$ be a convergent family (Definition~\ref{def:mot-model}) with $s_j,t_j\ge0$, and let $S$, $|S|\ge2$, be pinned (Theorem~\ref{thm:mot-pointmass}) with width one in every active table, $S'=\cE\setminus S$. Every active table is then the substitution of its cell of the bunch $\cB^+_{j,S}$ into its cell of the room $\cB^+_{j,S'}$ (Theorem~\ref{thm:widthone}, Proposition~\ref{prop:substitution}(iii)), the marks of $S'$ left in place, and the active letters read the cell through the one hinge value $\gamma=\gamma^K_{j,S}=\gamma^\Psi_{j,S}$. The core family of Theorem~\ref{thm:mot-pinned} is the room with its stand-in at $\one$, $S'^{(\one)}$ (Theorem~\ref{thm:gluewidth}(i)), the family $((\cB^+_{j,S'};s_j,t_j))_j$ of the active tables with the masses $(\mvec|_{S'},\kappa_S)$; the period of its pinned family $G$ is $c=\int_{\Delta_S}\phi^\rho_S$, the total mass of the law of $S$; and
\[
P^{\mathrm{mot}}_F=P^{\mathrm{mot}}_{S'^{(\one)}}\,P^{\mathrm{mot}}_G,\qquad\mathrm C(P^{\mathrm{mot}}_F)\subseteq\mathrm C(P^{\mathrm{mot}}_{S'^{(\one)}})\cdot\mathrm C(P^{\mathrm{mot}}_G):
\]
the identity $P_F=c\,P_{S'^{(\one)}}$ of Theorem~\ref{thm:mot-pointmass} is an identity of motivic periods.
\end{theorem}

Every period of a family glued on a set is the integral, against the law of the set, of the period of the partner with a stand-in for it (Theorem~\ref{thm:gluewidth}(i)). Where the law is a point the integral is one evaluation (Theorem~\ref{thm:mot-pointmass}); where the set is moreover a cell, the change of variables behind the evaluation is birational and the evaluation is an identity of motivic periods. No other set splits off, and a set of width two or more enters every period glued on it through a law that is not a point (Corollary~\ref{cor:mot-whatsplits}(i)); at every node, of every width, the open periods of node steps have the node form of Theorem~\ref{thm:mot-node}.

\chapter{Positive words}\label{ch:positive}

The output of the theory is a value: a period, a Breta, a closing of fibre integrals, and such values are read arithmetically. A closing at rational data writes a value as a rational combination of Trailmoments of lattice inputs and rational heads (Proposition~\ref{prop:headstails}, Lemma~\ref{lem:polecount}); a sequence of such values that do not vanish, cleared by integers that grow more slowly than the values decay, shows that some Trailmoment it carries is irrational (Theorem~\ref{thm:gap}, Chapter~\ref{ch:relations}). What keeps a value from vanishing is positivity, and positive words are made in four ways: second jets of a period in its weights, even powers of the contrast of two bases, Beta ports on the pairs of a manufactured family, and Beta ports read against one mark with the room integrated. The jets and the contrasts are combinations of periods of the same family at shifted data on its balance wall, bounded by the period times a factor that their statements control; the port words of a manufactured family are combinations of its interests at shifted masses over one normalization, and the reads of the Beta ports are bounded by the read of the constant insertion.

\section{Coefficient jets and basis contrasts}

% Reference 20: thm:jets; source line 10463 (rational positive weights).
% Proved in the last part: thm:jets.
\begin{theorem}[Positive coefficient jets]\label{thm:jets}
Let $F$ have real exponents $s_j,t_j\ge0$ and real masses on its balance wall, every proper wall positive. For an integer $n\ge1$, the \emph{depth}, the data $(n\mvec;n\mathbf s,n\mathbf t)$ are again on the balance wall, with walls $n\kappa_S$; in this statement $\oE$ is the expectation of $U\sim\mathrm{Breta}^\circ(F;n\mvec;n\mathbf s,n\mathbf t)$. Choose \emph{weight directions} $\mathbf b_{j,T}\in\mathbb Q^k$, one for each table $j$ and basis $T\in\cB_j$, and put
\[
\partial_{\mathbf b}K_j=\sum_{T\in\cB_j}\mathbf b_{j,T}\,\alpha^T,\qquad \partial_{\mathbf b}\Psi_j=\sum_{T\in\cB_j}\mathbf b_{j,T}\,\alpha^{\cE\setminus T},
\]
the derivatives of the letters of table $j$ along the weights in these directions.
\begin{enumerate}[label=(\roman*)]
\item For $v\in\mathbb R^k$ and small real $\varepsilon$, replace every table $j$ by the two tables $(\cB_j,\theta_j\pm\varepsilon\,v^{\mathsf T}\mathbf b_j)$, each with the exponents $(ns_j/2,nt_j/2)$. The period of this \emph{doubled family} at the masses $n\mvec$ is
\[
P_F(n\mvec;n\mathbf s,n\mathbf t)\ \oE\Big[\Big(\prod_j\Big(1-\varepsilon^2\Big(\frac{v^{\mathsf T}\partial_{\mathbf b}K_j}{K_j}\Big)^2\Big)^{-\frac{ns_j}2}\Big(1-\varepsilon^2\Big(\frac{v^{\mathsf T}\partial_{\mathbf b}\Psi_j}{\Psi_j}\Big)^2\Big)^{-\frac{nt_j}2}\Big)(U)\Big],
\]
whose coefficient of $\varepsilon^2$ is $v^{\mathsf T}\mathbf H_nv$, with the \emph{jet matrix}
\[
\mathbf H_n=\frac n2\,P_F(n\mvec;n\mathbf s,n\mathbf t)\ \oE\Big[\sum_j\Big(s_j\,\frac{\partial_{\mathbf b}K_j\,\partial_{\mathbf b}K_j^{\mathsf T}}{K_j^2}+t_j\,\frac{\partial_{\mathbf b}\Psi_j\,\partial_{\mathbf b}\Psi_j^{\mathsf T}}{\Psi_j^2}\Big)(U)\Big].
\]
$\mathbf H_n$ is positive semidefinite, and positive definite when the directions $\mathbf b_{j,T}$ of the active tables span $\mathbb R^k$.
\item For a rational positive semidefinite $k\times k$ matrix $Q$, $\operatorname{tr}(Q\mathbf H_n)>0$ exactly when $\mathbf b_{j,T}^{\mathsf T}Q\,\mathbf b_{j,T}>0$ for some active table $j$ and basis $T$. Completing squares writes $Q=\sum_lc_lv_lv_l^{\mathsf T}$ with rational $c_l>0$ and rational $v_l$, and so writes $\operatorname{tr}(Q\mathbf H_n)=\sum_lc_l\,v_l^{\mathsf T}\mathbf H_nv_l$ as a sum of directional second jets with positive rational coefficients.
\item In periods of $F$ (Theorem~\ref{thm:respwalls} at order two, the two copies of each table merged),
\begin{multline*}
\operatorname{tr}(Q\mathbf H_n)=\frac n2\sum_j\sum_{T,T'\in\cB_j}\mathbf b_{j,T}^{\mathsf T}Q\,\mathbf b_{j,T'}\Big[s_j\,P_F\big(n\mvec+\one_T+\one_{T'};\,n\mathbf s+2\mathbf e_j,\,n\mathbf t\big)\\
+t_j\,P_F\big(n\mvec+\one_{\cE\setminus T}+\one_{\cE\setminus T'};\,n\mathbf s,\,n\mathbf t+2\mathbf e_j\big)\Big]:
\end{multline*}
every term stays on its balance wall, a basis insertion changes the wall of $S$ by $|T\cap S|+|T'\cap S|-2d_j(S)\ge0$, and a cobasis insertion by $|S\setminus T|+|S\setminus T'|-2d^\ast_j(S)\ge0$.
\item Whenever (ii) gives positivity, $0<\operatorname{tr}(Q\mathbf H_n)\le n\,C(Q)\,P_F(n\mvec;n\mathbf s,n\mathbf t)$, with $C(Q)=\frac12\sum_j(s_j+t_j)\max_T\mathbf b_{j,T}^{\mathsf T}Q\,\mathbf b_{j,T}/\theta_{j,T}^2$.
\item With the basis weights and the cobasis weights of every table taken as independent variables, the part of $\mathbf H_n$ carried by $K_j$ is the Hessian of $P_F(n\mvec;n\mathbf s,n\mathbf t)$ in the basis weights of table $j$, paired with the directions $\mathbf b_{j,T}$, divided by $2(ns_j+1)$, and the part carried by $\Psi_j$ is the same in the cobasis weights divided by $2(nt_j+1)$. On the share side of the period (Corollary~\ref{cor:atomwall}) a weight derivative multiplies the integrand by minus the share divided by its weight, so every pin and every candidate pole hyperplane stays fixed.
\item At the centre, $s_j=t_j$ for every $j$ and $\mvec=\sigma\one$, the doubled family is again at its centre, and the table's gauge $\one$ stays native (Proposition~\ref{prop:centrefamily}).
\end{enumerate}
\end{theorem}

A weight direction moves the two copies of a table apart symmetrically, so the period moves at second order only, by the expectation of a square under the Breta shape law: a jet of the period in its weights is a positive word. Under the basis law $\partial_{\mathbf b}K_j/K_j$ is the mean of the directions $\mathbf b_{j,T}/\theta_{j,T}$, $\oE_{T\sim\mathbb P_{j,\alpha}}[\mathbf b_{j,T}/\theta_{j,T}]$, so the jet is the second moment, under the shape law, of the mean weight direction that the bases drawn at the shape carry. Table by table it is the period's own Hessian in the weights divided by $2(ns_j+1)$ or $2(nt_j+1)$, it is a finite combination of periods of the same family on the balance wall at data whose walls are no smaller, and it costs at most a factor linear in the depth against the period.

% Reference 20: lem:purityround; source line 10848
\begin{lemma}[Rounding inside a rational subspace]\label{lem:purityround}
Let $V$ be the subspace of the real symmetric $k\times k$ matrices on which finitely many linear functionals $Q\mapsto\operatorname{tr}(MQ)$, $M$ rational, vanish, spanned over $\mathbb R$ by integer symmetric matrices $G_1,\dots,G_r$, and let $Q_\ast=\sum_ju_jG_j\in V$ with $Q_\ast\succeq\varepsilon I$, $\varepsilon>0$. For an integer $\ell\ge1$ put $Q=\sum_j\operatorname{round}(\ell u_j)\,G_j$ and $\beta=\frac12\sum_j\big(\operatorname{tr}G_j^2\big)^{1/2}$. Then $Q$ is an integer matrix in $V$ with $Q\succeq(\ell\varepsilon-\beta)I$ and $\operatorname{tr}Q\le\ell\operatorname{tr}Q_\ast+k\beta$; for $\ell>\beta/\varepsilon$ it is positive definite.
\end{lemma}

The jet word $\operatorname{tr}(Q\mathbf H_n)$ is linear in $Q$, so conditions on its coefficients of Trailmoments, such as the vanishing of chosen ones, define a rational subspace of matrices $Q$. A positive definite matrix in that subspace rounds to an integer one in the same subspace, and the jet word keeps its positivity and the vanishing of the chosen coefficients together at integer data.

% Reference 20: thm:twolists; source line 10482 (a connected family with unit weights at its centre).
% Proved in the last part: thm:twolists.
\begin{theorem}[Two bases generate positive contrast words of any depth]\label{thm:twolists}
Let $F$ have real data $(\mvec;\mathbf s,\mathbf t)$ on its balance wall with every proper wall positive, with the pair coordinates $(\bmu_\perp;\boldsymbol\sigma,\boldsymbol\delta)$ (Definition~\ref{def:family}) and the centred potential $\Phi^\circ$ (Theorem~\ref{thm:saddle}). Let $T\neq T'$ be two bases of a table $j$, and for integers $n,h\ge1$ let the \emph{contrast word} be
\[
\frac12\,P_F(n\mvec;n\mathbf s,n\mathbf t)\ \oE_{U\sim\mathrm{Breta}^\circ(F;n\mvec;n\mathbf s,n\mathbf t)}\Big[\big(\mathbb P_{j,U}(T)-\mathbb P_{j,U}(T')\big)^{2h}+\big(\mathbb P^\ast_{j,U}(T)-\mathbb P^\ast_{j,U}(T')\big)^{2h}\Big]>0 ;
\]
the two contrasts have modulus below one on the open orthant.
\begin{enumerate}[label=(\roman*)]
\item The first half of the contrast word is
\[
\frac12\sum_{l=0}^{2h}(-1)^l\binom{2h}{l}\theta_{j,T}^{2h-l}\theta_{j,T'}^{\,l}\,P_F\big(n\mvec+(2h-l)\one_T+l\one_{T'};\,n\mathbf s+2h\mathbf e_j,\,n\mathbf t\big),
\]
and the second is the same with $\one_{\cE\setminus T}$, $\one_{\cE\setminus T'}$ in the masses and $2h\mathbf e_j$ added to the cobasis exponents; every term is balanced with its walls weakly raised. The inversion $\alpha\mapsto\alpha^{-1}$ carries the second half at $(\mvec;\mathbf s,\mathbf t)$ to the first half at $(2\sigma\one-\mvec;\mathbf t,\mathbf s)$, the data of Theorem~\ref{thm:cremona}; at the centre the two halves are equal.
\item For rational $\eta>0$ and $h=\eta n$, along the integers $n$ with $\eta n\in\mathbb N$, $\frac1n\log$ of the contrast word tends to
\[
-\inf_{z\in\one^\perp}\Big(\Phi^\circ(z)-2\ip{\bmu_\perp}z-2\eta\log\max\Big(\big|\mathbb P_{j,e^{2z}}(T)-\mathbb P_{j,e^{2z}}(T')\big|,\ \big|\mathbb P^\ast_{j,e^{2z}}(T)-\mathbb P^\ast_{j,e^{2z}}(T')\big|\Big)\Big);
\]
the infimum exceeds the minimum of $\Phi^\circ-2\ip{\bmu_\perp}{\cdot}$ on $\one^\perp$ and tends to infinity with $\eta$; that minimum is $-\lim_n\frac1n\log P_F(n\mvec;n\mathbf s,n\mathbf t)$.
\item Let the exponents be nonnegative, and write $q=2h$ and $H_q(x,y)=\sum_{l=0}^q(-1)^l\binom ql(x)_{q-l}(y)_l$. When $s_j>0$, the first half is the share side of $P_F(n\mvec;n\mathbf s,n\mathbf t)$ (Corollary~\ref{cor:atomwall}) with its integrand multiplied by $H_q(a_{j,T},a_{j,T'})/\big(2(ns_j)_q\big)$, $a$ the basis shares; when $t_j>0$, the second half is the same with $H_q(b_{j,T},b_{j,T'})/\big(2(nt_j)_q\big)$, $b$ the cobasis shares: polynomial insertions that keep every pin and every pole hyperplane.
\end{enumerate}
\end{theorem}

Two bases of one table give a contrast of modulus below one, which stays away from one on the compact sets where the Breta shape law concentrates; so its even powers are positive words that become exponentially small against the period, at a rate that grows without bound with the ratio of the power to the depth. At the centre with unit weights the basis law gives every basis the same chance at the table's gauge, where the law concentrates, and the contrast vanishes there. On the share side the contrast is a polynomial in the shares, so the contrast words of every depth keep the fibre, the pins and the poles of the period.

\section{Manufactured tables and their port words}

% Reference 20: def:manufactured; source line 10499
\begin{definition}[Manufactured tables]\label{def:manufactured}
Let \emph{seed letters} on $\cE$ be polynomials $L_i(x)=\sum_{P\in\operatorname{supp}L_i}\theta_{i,P}\,x^P$, $i\in I$, with supports $\operatorname{supp}L_i\subseteq2^{\cE}$, weights $\theta_{i,P}>0$ and exponents $w_i\ge0$; a table $(\cB_j,\theta_j)$ of a family with exponents $(s_j,t_j)$ gives the two seed letters $K_j$, with exponent $s_j$, and $\Psi_j$, with exponent $t_j$. Choose $\beta_i\in(0,1)$ and replace every mark $e$ by a pair $\{e^+,e^-\}$ of new marks, on $\cE^\pm=\{e^+,e^-:e\in\cE\}$. The polynomials
\begin{gather*}
\widetilde K_i(\alpha)=L_i(\one)\prod_{e\in\cE}(\alpha_{e^+}+\alpha_{e^-})-\beta_i\sum_P\theta_{i,P}\prod_{e\in P}\alpha_{e^+}\prod_{e\notin P}(\alpha_{e^+}+\alpha_{e^-}),\\
\widetilde\Psi_i(\alpha)=L_i(\one)\prod_{e\in\cE}(\alpha_{e^+}+\alpha_{e^-})-\beta_i\sum_P\theta_{i,P}\prod_{e\in P}\alpha_{e^-}\prod_{e\notin P}(\alpha_{e^+}+\alpha_{e^-})
\end{gather*}
are the basis and cobasis polynomials of a table on $\cE^\pm$ of rank $|\cE|$, whose bases $U^+\cup(\cE\setminus U)^-$, $U\subseteq\cE$, take one mark of each pair, with the weights $\widetilde\theta_{i,U}=L_i(\one)-\beta_i\sum_{P\subseteq U}\theta_{i,P}\ge L_i(\one)(1-\beta_i)>0$; the seed weights are recovered from $\big(L_i(\one)-\widetilde\theta_{i,U}\big)/\beta_i=\sum_{P\subseteq U}\theta_{i,P}$ by inversion over the subsets. The \emph{manufactured family} $\widetilde F$ has these tables with the exponents $(w_i,0)$; for masses $\mvec$ on $\cE^\pm$ its \emph{pair walls} are $\kappa_{\{e^+,e^-\}}=m_{e^+}+m_{e^-}-|\mathbf w|$, $|\mathbf w|=\sum_iw_i$, and its \emph{pair shapes} are $x_e=\alpha_{e^+}/(\alpha_{e^+}+\alpha_{e^-})$. For orders $\mathbf n\in\mathbb N^{\cE}$, positive masses and positive pair walls, its \emph{port word} is
\[
M_{\widetilde F,\mathbf n}(\mvec)=\frac{J_{\widetilde F}(\mvec;\mathbf w,\mathbf 0)}{\prod_e\Gamma\big(\kappa_{\{e^+,e^-\}}\big)}\ \oE_{\alpha\sim\mathrm{Breta}(\widetilde F;\mvec;\mathbf w,\mathbf 0)}\Big[\prod_{e\in\cE}\Pi^{m_{e^+},m_{e^-}}_{n_e}(x_e)\Big],
\]
and $M^\ast_{\widetilde F,\mathbf n}(\mvec)$ is the same at the exponents $(0,w_i)$, with $J_{\widetilde F}(\mvec;\mathbf 0,\mathbf w)$ and $\mathrm{Breta}(\widetilde F;\mvec;\mathbf 0,\mathbf w)$. An \emph{occupancy} is a family $\mathbf o=(o_{i,P})$ of integers $o_{i,P}\ge0$, zero when $w_i=0$, with the totals $|o_i|=\sum_Po_{i,P}$ and $|\mathbf o|=\sum_i|o_i|$, the \emph{loads} $\operatorname{load}_e(\mathbf o)=\sum_{i}\sum_{P\ni e}o_{i,P}$, and the \emph{coefficient}
\[
\operatorname{cf}(\mathbf o)=\prod_i\Big[\frac{(w_i)_{|o_i|}}{L_i(\one)^{w_i}\prod_Po_{i,P}!}\prod_P\Big(\frac{\beta_i\,\theta_{i,P}}{L_i(\one)}\Big)^{o_{i,P}}\Big].
\]
\end{definition}

% Reference 20: thm:manufinterest; source line 10507
\begin{theorem}[Interest of a manufactured family; the occupancy sum]\label{thm:manufinterest}
In the pair scales and shapes, $(\alpha_{e^+},\alpha_{e^-})=\rho_e(x_e,1-x_e)$,
\[
\widetilde K_i=\prod_e\rho_e\,\big(L_i(\one)-\beta_iL_i(x)\big),\qquad \widetilde\Psi_i=\prod_e\rho_e\,\big(L_i(\one)-\beta_iL_i(\one-x)\big).
\]
Let the masses be positive and the pair walls positive.
\begin{enumerate}[label=(\roman*)]
\item Every wall of $\widetilde F$ at the exponents $(w_i,0)$ is positive, and
\[
J_{\widetilde F}(\mvec;\mathbf w,\mathbf 0)=\prod_e\Gamma\big(\kappa_{\{e^+,e^-\}}\big)\int_{[0,1]^{\cE}}\prod_ex_e^{m_{e^+}-1}(1-x_e)^{m_{e^-}-1}\prod_i\big(L_i(\one)-\beta_iL_i(x)\big)^{-w_i}\,dx ,
\]
the integral convergent at all positive masses; the simplex integral $P_{\widetilde F}(\mvec;\mathbf w,\mathbf 0)$ is $\Gst{|\cE|}{1}{(\kappa_{\{e^+,e^-\}})_e}{\sum_e\kappa_{\{e^+,e^-\}}}$ times the same integral. Under $\mathrm{Breta}(\widetilde F;\mvec;\mathbf w,\mathbf 0)$ the pair scales are independent, $\rho_e\sim\operatorname{Gamma}_{\kappa_{\{e^+,e^-\}}}$, and independent of the pair shapes, whose density on the cube is proportional to the integrand. So $M_{\widetilde F,\mathbf 0}(\mvec)$ is the integral displayed, and $M_{\widetilde F,\mathbf n}(\mvec)$ is that integral with $\prod_e\Pi^{m_{e^+},m_{e^-}}_{n_e}(x_e)$ inserted; with $L_i(\one-x)$ in place of $L_i(x)$ the same holds for $M^\ast_{\widetilde F,\mathbf n}(\mvec)$.
\item The integral is the series of nonnegative terms
\[
M_{\widetilde F,\mathbf 0}(\mvec)=\sum_{\mathbf o}\operatorname{cf}(\mathbf o)\prod_e\Gst21{m_{e^+}+\operatorname{load}_e(\mathbf o),\ m_{e^-}}{m_{e^+}+m_{e^-}+\operatorname{load}_e(\mathbf o)}
\]
over the occupancies.
\end{enumerate}
\end{theorem}

% Reference 20: thm:manufports; source line 10526
\begin{theorem}[The port word of a manufactured family]\label{thm:manufports}
Let the masses be positive, the pair walls positive and $\mathbf n\in\mathbb N^{\cE}$. Then
\begin{align*}
M_{\widetilde F,\mathbf n}(\mvec)&=\sum_{\mathbf o}\operatorname{cf}(\mathbf o)\prod_e\binom{\operatorname{load}_e(\mathbf o)}{n_e}\Gst21{m_{e^+}+\operatorname{load}_e(\mathbf o),\ m_{e^-}+n_e}{m_{e^+}+m_{e^-}+\operatorname{load}_e(\mathbf o)+n_e}\\
&=\frac1{\prod_en_e!}\int_{[0,1]^{\cE}}\prod_ex_e^{m_{e^+}+n_e-1}(1-x_e)^{m_{e^-}+n_e-1}\ \partial_x^{\mathbf n}\prod_i\big(L_i(\one)-\beta_iL_i(x)\big)^{-w_i}\,dx .
\end{align*}
Every occupancy with $\operatorname{load}_e(\mathbf o)<n_e$ for some $e$ is annihilated, every surviving term is positive, every mixed derivative of $\prod_i(L_i(\one)-\beta_iL_i(x))^{-w_i}$ is nonnegative on the cube, and $M_{\widetilde F,\mathbf n}(\mvec)>0$ exactly when every mark $e$ with $n_e>0$ lies in a member of the support of a seed letter with $w_i>0$. At integer masses the Beta ports have integer coefficients, and $M_{\widetilde F,\mathbf n}(\mvec)$ is an integer combination of the $M_{\widetilde F,\mathbf 0}$ at the masses shifted by $(k_e,n_e-k_e)$ on the pair of $e$, $0\le k_e\le n_e$: of the interests of $\widetilde F$ at these masses, all of pair walls $\kappa_{\{e^+,e^-\}}+n_e$, over the one normalization $\prod_e\Gamma(\kappa_{\{e^+,e^-\}}+n_e)$.
\end{theorem}

% Reference 20: thm:suppression; source line 10538
\begin{theorem}[Suppression selected by incidence]\label{thm:suppression}
In Theorem~\ref{thm:manufports} write $\operatorname{occ}(\mathbf n)=\min\{|\mathbf o|:\operatorname{load}_e(\mathbf o)\ge n_e\text{ for every }e\}$ over the occupancies, and $\beta=\max_i\beta_i<1$ over the seed letters with $w_i>0$. Then
\[
0\le M_{\widetilde F,\mathbf n}(\mvec)\le\prod_e\Big[\Gst21{m_{e^+},\ m_{e^-}}{m_{e^+}+m_{e^-}}\frac{(m_{e^-})_{n_e}}{n_e!}\Big]\prod_iL_i(\one)^{-w_i}\sum_{l\ge\operatorname{occ}(\mathbf n)}\frac{(|\mathbf w|)_l}{l!}\,\beta^l .
\]
For an integer demand $\mathbf d\in\mathbb N^{\cE}$ let $\operatorname{fc}(\mathbf d)=\min\big\{\sum_{i,P}\lambda_{i,P}:\ \lambda\ge0,\ \sum_{i,P}\lambda_{i,P}\one_P\ge\mathbf d\big\}$, the \emph{fractional cover} of the demand by the members of the supports of the seed letters with $w_i>0$. Then $\operatorname{occ}(q\mathbf d)=q\operatorname{fc}(\mathbf d)+O(1)$ as $q\to\infty$, and $\limsup_q\frac1q\log M_{\widetilde F,q\mathbf d}(\mvec)\le\operatorname{fc}(\mathbf d)\log\beta<0$ for a nonzero demand that the supports cover. When the seed letters are the letters $K_j,\Psi_j$ of the tables of a family, every surviving occupancy satisfies $\sum_j\big(r_j(S)\,|o_{K_j}|+(|S|-d_j(S))\,|o_{\Psi_j}|\big)\ge\sum_{e\in S}n_e$ for every $S\subseteq\cE$, $|o_{K_j}|$ and $|o_{\Psi_j}|$ the total occupancies of the two letters of table $j$.
\end{theorem}

% Reference 20: thm:manufreflection; source line 10549
\begin{theorem}[Reflection of the port words of a manufactured family]\label{thm:manufreflection}
The swap $\alpha_{e^+}\leftrightarrow\alpha_{e^-}$ on every pair, $x_e\leftrightarrow1-x_e$, swaps $\widetilde K_i$ with $\widetilde\Psi_i$ and $m_{e^+}$ with $m_{e^-}$, and keeps every seed weight and every $\beta_i$: with $\mvec'$ the masses swapped on every pair, $M_{\widetilde F,\mathbf n}(\mvec)=(-1)^{|\mathbf n|}M^\ast_{\widetilde F,\mathbf n}(\mvec')$.
\end{theorem}

A manufactured table turns any positive polynomials, the two letters of every table of a family among them, into a table whose bases are the transversals of the pairs. On the pair shapes its letter is the constant $L_i(\one)$ less a copy of the seed scaled by $\beta_i$, so the port word is a series of nonnegative terms indexed by occupancies of the members of the seed, and a Beta port of order $n_e$ on a pair keeps only the occupancies that load the mark at least $n_e$ times. The port words are therefore positive exactly when the incidence of the supports can meet the demand, and they are suppressed at the exponential rate of the cheapest fractional cover of the demand: the incidence structure of the seed decides which port words survive and how small they are.

\section{Mark relations with the room integrated}

% Reference 20: def:markwords; source line 10582
\begin{definition}[Mark reads]\label{def:markwords}
Select a mark with coordinate $x$, and let the coordinates $\alpha_R$ of its room range over a measure space with a density $\Omega(\alpha_R)\ge0$. The letters (Section~\ref{sec:letters}) in the selected mark are polynomials $L_i(x,\alpha_R)$ in $x$, $i=1,\dots,h$, positive where $x$ ranges, with exponents $w_i$; they are \emph{linear} when $L_i=c_i(\alpha_R)+d_i(\alpha_R)\,x$ with $c_i>0$ and $d_i\ge0$, and a linear letter \emph{sees the mark} when $w_id_i>0$ on a set of positive measure. For $a,b>0$ the \emph{mark read} of an insertion $g(x,\alpha_R)$ on the box is
\[
\langle g\rangle_{a,b;\mathbf w}=\int\!\!\int_0^1g\,x^{a-1}(1-x)^{b-1}\prod_iL_i^{-w_i}\,dx\ \Omega(\alpha_R)\,d\alpha_R ,
\]
and for $a,\lambda>0$ its read on the half-line is
\[
\langle g\rangle^{\lambda}_{a;\mathbf w}=\int\!\!\int_0^\infty g\,x^{a-1}e^{-\lambda x}\prod_iL_i^{-w_i}\,dx\ \Omega(\alpha_R)\,d\alpha_R ,
\]
wherever these converge; at a fixed point of the room, $\Omega(\alpha_R)\,d\alpha_R$ is the Dirac measure there. A sequence of reads at the data $a_n+j$, $b_n$, $\mathbf w(n)$ and $\Omega_n$ keeps the depth $n$ (Theorem~\ref{thm:jets}) apart from the \emph{moment index} $j$ of the selected mark.
\end{definition}

% Reference 20: thm:markrelation; source line 10558
\begin{theorem}[The mark relation]\label{thm:markrelation}
Let the letters be polynomial in the selected mark, with exponents $w_i$, and let $A(x,\alpha_R)$ be a polynomial in $x$ with coefficients depending on the room.
\begin{enumerate}[label=(\alph*)]
\item On the box, for $a,b>0$:
\[
\Big\langle x(1-x)\,\partial_xA+\big(a-(a+b)x\big)A-x(1-x)\,A\sum_iw_i\,\partial_x\log L_i\Big\rangle_{a,b;\mathbf w}=0 .
\]
\item On the half-line, for $a,\lambda>0$, whenever the reads converge:
\[
\Big\langle x\,\partial_xA+(a-\lambda x)A-x\,A\sum_iw_i\,\partial_x\log L_i\Big\rangle^{\lambda}_{a;\mathbf w}=0 .
\]
\item On a restricted interval the difference of the flux at its two ends is retained, and a kernel coupled to $x$ contributes its product-rule term.
\end{enumerate}
Multiplying by $\prod_iL_i$ gives a polynomial numerator; on a box of several marks the relation holds in each coordinate with the others as room.
\end{theorem}

% Reference 20: thm:momentblock; source line 10571
\begin{theorem}[The moments of a selected mark close in a finite span]\label{thm:momentblock}
On a family select a mark $e$, $x=\alpha_e$, with room $R=\cE-e$; splitting the bases and the cobases by $e$ gives the two coefficients of every letter in $x$. Pull the letters free of $x$ into the room, absorb pure powers of $x$ into its mass, and write the remaining letters $L_i=c_i(\alpha_R)+d_i(\alpha_R)\,x$, $c_i,d_i>0$, $i=1,\dots,h$, with exponents $w_i$; read them on the half-line at a fixed point of the room (Definition~\ref{def:markwords}), with $m,\lambda>0$. With $\prod_iL_i=\sum_{k=0}^hp_kx^k$, $p_{-1}=0$, $q_0=0$ and $q_k=\sum_iw_id_i\,[x^{k-1}]\prod_{i'\neq i}L_{i'}$,
\[
\lambda p_h\,\langle x^{n+h+1}\rangle^{\lambda}_{m;\mathbf w}=\sum_{k=0}^h\big[(m+n+k)\,p_k-q_k-\lambda p_{k-1}\big]\,\langle x^{n+k}\rangle^{\lambda}_{m;\mathbf w},
\]
a recurrence whose coefficients are polynomials in the room coordinates, with leading coefficient $\lambda p_h$, $p_h=\prod_id_i>0$. For a mark $f$ of the room,
\[
\partial_{\alpha_f}\langle x^n\rangle^{\lambda}_{m;\mathbf w}=-\sum_iw_i\Big[(\partial_{\alpha_f}d_i)\,\langle x^{n+1}\rangle^{\lambda}_{m;\mathbf w+\mathbf e_i}+(\partial_{\alpha_f}c_i)\,\langle x^n\rangle^{\lambda}_{m;\mathbf w+\mathbf e_i}\Big],
\]
and the moments with at most $\ell$ such derivatives lie in the span, over rational functions of the room, of $\{\langle x^k\rangle^{\lambda}_{m;\mathbf w+v}:0\le k\le h,\ v\in\mathbb N^h,\ |v|\le\ell\}$: at most $(h+1)\binom{h+\ell}{\ell}$ states.
\end{theorem}

% Reference 20: thm:iterdefect; source line 10586
\begin{theorem}[Positive reads of the Beta ports]\label{thm:portreads}
Let the letters be linear, $L_i=c_i+d_ix$, with exponents $w_i\ge0$, let $a,b>0$ and $q\in\mathbb N$, and let $J$ be a set of letters, with $w_J=\sum_{i\in J}w_i$.
\begin{enumerate}[label=(\roman*)]
\item The Beta port of order $q$ reads (Theorem~\ref{thm:betaport}(iii), integrated over the room)
\[
\big\langle(-1)^q\Pi^{a,b}_q\big\rangle_{a,b;\mathbf w}=\sum_{|v|=q}\prod_i\frac{(w_i)_{v_i}}{v_i!}\,\Big\langle\big(x(1-x)\big)^q\prod_i\Big(\frac{d_i}{L_i}\Big)^{v_i}\Big\rangle_{a,b;\mathbf w}\in\Big[0,\frac{(|\mathbf w|)_q}{q!}\,\langle1\rangle_{a,b;\mathbf w}\Big].
\]
\item The part of this sum over the $v$ supported in $J$ lies in $\big[0,\frac{(w_J)_q}{q!}\langle1\rangle_{a,b;\mathbf w}\big]$, and for $q>0$ it is positive when some letter of $J$ sees the mark.
\item At $q=1$,
\begin{multline*}
a\,\langle1\rangle_{a,b;\mathbf w}-(a+b)\,\langle x\rangle_{a,b;\mathbf w}-\sum_{i\notin J}w_i\Big\langle\frac{x(1-x)\,d_i}{L_i}\Big\rangle_{a,b;\mathbf w}\\
=\sum_{i\in J}w_i\Big\langle\frac{x(1-x)\,d_i}{L_i}\Big\rangle_{a,b;\mathbf w}\in\big[0,\,w_J\langle1\rangle_{a,b;\mathbf w}\big].
\end{multline*}
\end{enumerate}
Along a sequence at the data $a_n+j$, $b_n$, $\mathbf w(n)$ the bounds hold with $\langle1\rangle_{a_n,b_n;\mathbf w(n)}\ge\langle1\rangle_{a_n+j,b_n;\mathbf w(n)}$ in place of the read at the moment index $j$, uniformly in $j$.
\end{theorem}

% Reference 20: prop:finiteband; source line 10601
\begin{proposition}[The finite-band polynomial]\label{prop:finiteband}
In Theorem~\ref{thm:portreads} write $L_{J^c}=\prod_{i\notin J}L_i$. The polynomial
\[
A=\big[a-(a+b)x\big]L_{J^c}+x(1-x)\Big(\partial_xL_{J^c}-L_{J^c}\sum_{i\notin J}\frac{w_id_i}{L_i}\Big),
\]
of degree at most $h-|J|+1$ in $x$, has the read
\[
\langle A\rangle_{a,b;\mathbf w}=\sum_{i\in J}w_i\Big\langle\frac{x(1-x)\,d_i\,L_{J^c}}{L_i}\Big\rangle_{a,b;\mathbf w},
\]
positive when a letter of $J$ sees the mark and the reads are finite. On the half-line, $A=a\,L_{J^c}+x\,\partial_xL_{J^c}-\lambda xL_{J^c}-xL_{J^c}\sum_{i\notin J}w_id_i/L_i$ has $\langle A\rangle^{\lambda}_{a;\mathbf w}=\sum_{i\in J}w_i\big\langle x\,d_i\,L_{J^c}/L_i\big\rangle^{\lambda}_{a;\mathbf w}$.
\end{proposition}

% Reference 20: thm:portsequence; source line 10608 (a selected mark at mass one).
% Proved in the last part: thm:portsequence.
\begin{theorem}[A positive sequence of Beta ports on an interval of a selected mark]\label{thm:portsequence}
In a projective chart take a selected mark at a mass $m_e\in(0,1]$, restricted to an interval $0<\lambda_0<\alpha_e<\lambda_1$, and write $\alpha_e=\lambda_0+(\lambda_1-\lambda_0)x$, $x\in[0,1]$; its letters are linear in $\alpha_e$ with nonnegative coefficients, $L_i=c_i+d_ix$ with $d_i\ge0$ and $c_i\ge\lambda_0d_i/(\lambda_1-\lambda_0)$, and exponents $w_i\ge0$, the letters constant in $x$ going into the room density, and the factor $\alpha_e^{m_e-1}$ of the mass is counted among them as the letter $\lambda_0+(\lambda_1-\lambda_0)x$ at the exponent $1-m_e$. So $x(1-x)\,d_i/L_i\le(\sqrt{\lambda_1}-\sqrt{\lambda_0})/(\sqrt{\lambda_1}+\sqrt{\lambda_0})$ on $[0,1]$ for every letter. Let $\langle1\rangle_{1,1;\mathbf w}$ be finite. By Theorem~\ref{thm:portreads}(i) at $a=b=1$ and the bound of Theorem~\ref{thm:betaport}(iii), integrated over the room, for $n\in\mathbb N$,
\[
0<\big\langle(-1)^n\Pi^{1,1}_n\big\rangle_{1,1;\mathbf w}\le\frac{(|\mathbf w|)_n}{n!}\Big(\frac{\sqrt{\lambda_1}-\sqrt{\lambda_0}}{\sqrt{\lambda_1}+\sqrt{\lambda_0}}\Big)^n\langle1\rangle_{1,1;\mathbf w},
\]
the first inequality when a letter sees the mark, as the letter of the mass does when $m_e<1$.
\end{theorem}

% Reference 20: thm:portrecurrence; source line 10619 (on the diagonal $a=b=n+1$).
% Proved in the last part: thm:portrecurrence.
\begin{theorem}[The finite-band recurrence in the order of the Beta ports]\label{thm:portrecurrence}
At a fixed point of the room let the letters be linear, $L_i=c_i+d_ix$ with $c_i>0$ and $d_i\ge0$, $i=1,\dots,h$, with exponents $w_i(n)$ depending on $n\in\mathbb N$, let $a_n,b_n>0$, and in this statement write $L=\prod_iL_i=\sum_kp_kx^k$ and
\[
A_n=\big(a_n-(a_n+b_n)x\big)\,L+x(1-x)\,\partial_xL-x(1-x)\,L\sum_i\frac{w_i(n)\,d_i}{L_i}.
\]
For every polynomial $g$, $\big\langle L\,x(1-x)\,\partial_xg+A_ng\big\rangle_{a_n,b_n;\mathbf w(n)}=0$ (Theorem~\ref{thm:markrelation}(a) with the polynomial $Lg$). Expanding $L\,x(1-x)\,\partial_xg+A_ng$ at $g=(-1)^r\Pi^{a_n,b_n}_r$ in the Beta ports $(-1)^{r+s}\Pi^{a_n,b_n}_{r+s}$ by the finite rules of Theorem~\ref{thm:betaport}(v) gives the relation
\[
\sum_{s=-h-1}^{h+1}\mathfrak b_{n,r,s}\,\big\langle(-1)^{r+s}\Pi^{a_n,b_n}_{r+s}\big\rangle_{a_n,b_n;\mathbf w(n)}=0
\]
of width at most $2h+3$, whose coefficients, written $\mathfrak b_{n,r,s}$ in this statement, are rational in $n$ and $r$ when $a_n$, $b_n$ and the $w_i(n)$ are, and polynomial in the room data. With $|\mathbf w(n)|=\sum_iw_i(n)$ and $p_h=\prod_id_i$, the furthest forward coefficient is
\[
\mathfrak b_{n,r,h+1}=\big[|\mathbf w(n)|-a_n-b_n-h-r\big]\,p_h\,(-1)^{h+1}\binom{a_n+b_n+2r-2}{r}\Big/\binom{a_n+b_n+2r+2h}{r+h+1}.
\]
\end{theorem}

% Reference 20: lem:recurrencesign; source line 10626
\begin{lemma}[A rational eventual sign certificate]\label{lem:recurrencesign}
Let $Q_n(X)=n^dq_d(X)+\sum_{k<d}n^kq_k(X)$ with polynomials $q_k$, and let rationals $X_1\ge X_{\max}>0$ and $\delta>0$ satisfy $q_d\ge\delta$ on $[0,X_1]$. With $B_k=\sum_i\big|[X^i]q_k\big|\,X_1^i$, $Q_n\ge\delta n^d/2$ on $[0,X_{\max}]$ for every $n\ge\max\big(1,2\sum_{k<d}B_k/\delta\big)$. Composition of recurrences with variable coefficients keeps its shifts: with $\operatorname{Sh}$ the shift $n\mapsto n+1$, $\big(\sum_jb_j(n)\operatorname{Sh}^j\big)\big(\sum_ia_i(n)\operatorname{Sh}^i\big)=\sum_{i,j}b_j(n)\,a_i(n+j)\operatorname{Sh}^{i+j}$.
\end{lemma}

With the room held fixed, a mark is integrated alone against its letters, and the mark relation is the integration by parts that this allows: the flux of the density $x^{a-1}(1-x)^{b-1}$ times a polynomial vanishes at both ends of the box, and that of the Gamma density at both ends of the half-line. Two uses follow. Against the Gamma density a mark of a family closes its moments into a recurrence of bounded order whose coefficients are polynomials in the room, and the derivatives in the room stay in a finite span: one mark closes all its moments. Against a Beta port the same relation gives positive reads, sums of the ratios $x(1-x)d_i/L_i$ that a letter seeing the mark keeps positive and the total exponent bounds, and at every pair $a_n,b_n>0$ the Beta ports $\Pi^{a_n,b_n}_r$ obey a recurrence of bounded width in their order $r$ whose furthest coefficient is explicit: when every $d_i$ is positive it vanishes at the one order $r=|\mathbf w(n)|-a_n-b_n-h$ at most.

\chapter{The critical-line law}\label{ch:share}

Let $\cE$ have two marks, and let $F$ be the family of the one list $\binom{\cE}1$, whose bases are its two singletons, at the exponents $(1,1)$ and the masses $\one$ (Theorem~\ref{thm:sharelaw}). Its period is one, $P_F(\one;1,1)=1$. On the fibre of the atom formula on the balance wall (Corollary~\ref{cor:atomwall}), one share per basis and letter, the pins leave one free share: the basis share and the cobasis share of one basis are equal, $a=b=\frac12+iy$, and the fibre is the line $\re a=\frac12$, the critical line. The law of $y$, the imaginary part of that share, is the critical-line law $\mathbb P_{\mathrm{crit}}$ (Definition~\ref{def:pairlaw}). Its characteristic function is a stacked Gamma, its moments are rational, and its negative moments, read as finite parts, are rational multiples of the odd Trailmoments of the lattice input on the coset $\frac12+\mathbb Z$ (Theorem~\ref{thm:sharelaw}). A function $R$ of the share, read against an odd derivative of $\mathbb P_{\mathrm{crit}}$ along the critical line, gives the critical-line word $\operatorname{crit}_k[R]$ (Definition~\ref{def:critical}): it carries one Trailmoment of odd index with a rational part, and it is positive whenever $R$ is the Fourier transform of a positive measure (Theorem~\ref{thm:positiveslot}).

% Reference 20: thm:sharelaw (the density); source line 10641
\begin{definition}[The critical-line law]\label{def:pairlaw}
The \emph{critical-line law} is the probability law $\mathbb P_{\mathrm{crit}}(dy)=p_{\mathrm{crit}}(y)\,dy$ on $\mathbb R$ with
\[
p_{\mathrm{crit}}(y)=\frac1{2\pi}\,\Gst40{\frac12+iy,\ \frac12+iy,\ \frac12-iy,\ \frac12-iy}{}=\frac{2\pi}{(e^{\pi y}+e^{-\pi y})^2},
\]
the stacked Gamma of Lemma~\ref{lem:resonantpair} on the line $\re a=\frac12$ divided by $2\pi$.
\end{definition}

% Reference 20: thm:sharelaw; source line 10641
\begin{theorem}[The critical-line law and its moments]\label{thm:sharelaw}
Let $\cE$ have two marks and let $F$ be the family of the one list $\binom{\cE}1$, whose bases are its two singletons, at the exponents $(1,1)$ and the masses $\one$. On the fibre of Corollary~\ref{cor:atomwall} the basis and cobasis shares of one basis are equal, $a=b=\frac12+iy$, the fibre integral is $P_F(\one;1,1)=1$, and the imaginary part $Y$ of the share has the law $\mathbb P_{\mathrm{crit}}$. In this statement $\oE$ is the expectation of $Y\sim\mathbb P_{\mathrm{crit}}$.
\begin{enumerate}[label=(\roman*)]
\item $\oE[e^{i\xi Y}]$ is the Fourier read of Lemma~\ref{lem:resonantpair}(iii) at $t=\xi$; every odd moment of $Y$ vanishes and every even moment is rational.
\item For $j\ge1$ the finite part at $y=0$ (Definition~\ref{def:Reg}) of $\oE[Y^{-2j}]$ is
\begin{align*}
\Reg\oE\big[Y^{-2j}\big]&=(-1)^j\,2j\,\Tx{2j+1}{\Ilat{\frac12+\mathbb Z}}{0}\\
&=(-1)^j\,2j\,(2^{2j+1}-1)\,\Tx{2j+1}{\IZ}{0},
\end{align*}
the second line by the lattice relation (Theorem~\ref{thm:latticefree}(i)) at the step $q=2$, the lattice $\frac12\mathbb Z$ splitting into $\mathbb Z$ and $\frac12+\mathbb Z$.
\item With $X=\frac12+iY$, $\oE[e^{tX}]=\dfrac{te^t}{e^t-1}=t+\dfrac t{e^t-1}$, and $d_{m+1}\,\oE[X^m]\in\mathbb Z$ for every $m\in\mathbb N$.
\end{enumerate}
\end{theorem}

The list $\binom{\cE}1$ on two marks is the smallest table with a share that moves: at the centre the pins hold its basis and cobasis shares equal, on the critical line, and the period, one, is the integral of the law $\mathbb P_{\mathrm{crit}}$ of their imaginary part. Everything about that law is read off the stacked Gamma $\Gst40{a,\ a,\ 1-a,\ 1-a}{}$ of Lemma~\ref{lem:resonantpair}. Its characteristic function is the stacked Gamma $\Gst20{1+\frac{i\xi}{2\pi},\ 1-\frac{i\xi}{2\pi}}{}$, its moments are rational with denominators bounded by the clearing integers, and its negative moments, which diverge at the origin, have finite parts that are rational multiples of the odd Trailmoments of the lattice input on the coset $\frac12+\mathbb Z$, and so of the counting input.

\section{Critical-line words}

% Reference 20: prop:critical (the word); source line 10660
\begin{definition}[Critical-line words]\label{def:critical}
For an integer $k\ge1$ and a function $R$ on the line $\re u=\frac12$, the \emph{critical-line word} of order $k$ is
\[
\operatorname{crit}_k[R]=-\frac2{(2k)!}\,\frac1{2\pi i}\int_{(1/2)}\frac{R(u)}{1-2u}\,\partial_u^{2k-1}\Gst40{u,\ u,\ 1-u,\ 1-u}{}\,du ,
\]
wherever it converges; the integrand is regular at $u=\frac12$, where the odd derivative of the stacked Gamma, which is even about $\frac12$, vanishes. On the line, $u=\frac12+iy$,
\[
\operatorname{crit}_k[R]=\frac{(-1)^k}{(2k)!}\int_{\mathbb R}R(\tfrac12+iy)\,\frac{p_{\mathrm{crit}}^{(2k-1)}(y)}{y}\,dy .
\]
\end{definition}

% Reference 20: prop:critical; source line 10660
\begin{proposition}[Critical-line words of polynomials]\label{prop:critical}
For an integer $k\ge1$ and a polynomial $R(u)=(1-2u)\,Q(u)+R(\frac12)$, with $\oE$ the expectation of $Y\sim\mathbb P_{\mathrm{crit}}$,
\begin{align*}
\operatorname{crit}_k[R]&=R(\tfrac12)\,\Tx{2k+1}{\Ilat{\frac12+\mathbb Z}}{0}+\frac2{(2k)!}\,\oE\big[Q^{(2k-1)}(\tfrac12+iY)\big]\\
&=(2^{2k+1}-1)\,R(\tfrac12)\,\Tx{2k+1}{\IZ}{0}+\frac2{(2k)!}\,\oE\big[Q^{(2k-1)}(\tfrac12+iY)\big],
\end{align*}
the second line by Theorem~\ref{thm:latticefree}(i), as in Theorem~\ref{thm:sharelaw}(ii).
\end{proposition}

% Reference 20: rows, before lem:lcm; source line 10720
\begin{definition}[Rows; clearing]\label{def:row}
A value written as $x=\sum_{i=1}^mA_iv_i-B$, a rational combination of fixed real numbers $v_i$, its \emph{coordinates}, such as Trailmoments of the counting input, less a rational number $B$, its \emph{rational part}, has the \emph{row} $(A_1,\dots,A_m;B)$; a positive integer \emph{clears} the row when it makes all these coefficients integers.
\end{definition}

% Reference 20: cor:criticalden; source line 10671
\begin{corollary}[The rational part of an integer polynomial argument]\label{cor:criticalden}
For $P\in\mathbb Z[x]$ of degree $D$ and $R(u)=P(2u)$, the rational part of $\operatorname{crit}_k[R]$ (Definition~\ref{def:row}) is cleared by $\frac{(2k)!}2\,d_{\max(1,D-2k+1)}$.
\end{corollary}

% New: proved in the last part.
\begin{proposition}[Binomial moments of the critical-line law]\label{prop:ra-binomial}
Let $X=\frac12+iY$ with $Y\sim\mathbb P_{\mathrm{crit}}$, $\oE$ its expectation, and $H_m=\Th{1}{\IZ}{0}{m}$.
\begin{enumerate}[label=(\roman*)]
\item For complex $z$ with $|z|<1$, $\oE\big[(1+z)^X\big]=(1+z)\log(1+z)/z$.
\item $\oE\big[\binom X0\big]=1$ and $\oE\big[\binom Xi\big]=\dfrac{(-1)^{i+1}}{i(i+1)}$ for $i\ge1$.
\item For $i\ge1$, $\oE\Big[\dfrac{d}{du}\dbinom ui\Big|_{u=X}\Big]=2(-1)^i\,\dfrac{H_{i-1}-1}{i(i+1)}$. So for every polynomial $Q=\sum_iq_i\binom ui$,
\[
\operatorname{crit}_1\big[(1-2u)\,Q\big]=2\sum_{i\ge1}q_i\,(-1)^i\,\frac{H_{i-1}-1}{i(i+1)} .
\]
\end{enumerate}
\end{proposition}

In the binomial basis the critical-line word of a polynomial has explicit coefficients, harmonic numbers over $i(i+1)$, whose denominators a prime $q$ divides only to bounded powers; this is the basis in which the clearings of the critical-line words of polynomials are counted, prime by prime.

% Reference 20: thm:positiveslot; source line 10678
\begin{theorem}[Positive arguments for the critical-line word]\label{thm:positiveslot}
For an integer $k\ge1$ and real $t$, with $\oE$ the expectation of $Y\sim\mathbb P_{\mathrm{crit}}$,
\begin{gather*}
\operatorname{crit}_k\big[e^{t(u-\frac12)}\big]=\frac1{(2k)!}\int_{|t|}^\infty x^{2k-1}\,\oE\big[e^{ixY}\big]\,dx>0,\\
\operatorname{crit}_k[1]=\Tx{2k+1}{\Ilat{\frac12+\mathbb Z}}{0}=(2^{2k+1}-1)\,\Tx{2k+1}{\IZ}{0}\quad\text{(Theorem~\ref{thm:latticefree}(i))}.
\end{gather*}
If $R(\frac12+iy)=\int e^{ity}\,\phi(dt)$ for a finite nonzero positive measure $\phi$ on $\mathbb R$, then $\operatorname{crit}_k[R]=\int\operatorname{crit}_k\big[e^{t(u-\frac12)}\big]\,\phi(dt)>0$.
\end{theorem}

On the critical line the word reads its argument against an odd derivative of the density of the critical-line law divided by the imaginary part of the share, a kernel without a fixed sign. Its Fourier transform, the word of an exponential argument, is the integral beyond $|t|$ of the characteristic function times a power, a positive function, so an argument that is the Fourier transform of a positive measure has a positive word. This is how positive words carrying one odd Trailmoment are made.

% Reference 20: thm:pureslot; source line 10690
\begin{theorem}[A positive rational argument with one odd index]\label{thm:pureslot}
Let $M\in\mathbb N$ and $p(v)=\sum_{j=0}^Db_jv^j\in\mathbb Q[v]$, nonnegative on $[0,1]$ and not identically zero, and let
\[
R_{M,p}(u)=\sum_{j=0}^Db_j\Big(\frac1{u+M+j}+\frac1{1-u+M+j}\Big),\qquad R_{M,p}(\tfrac12)=4\sum_j\frac{b_j}{2M+2j+1}>0 .
\]
Then for every $k\ge1$
\begin{align*}
\operatorname{crit}_k[R_{M,p}]&=\frac2{(2k)!}\int_0^1\!\!\int_0^1\frac{v^Mx^{-1/2}\,p(v)\,[-\log(vx)]^{2k}}{1-vx}\,dx\,dv\\
&=R_{M,p}(\tfrac12)\,\Tx{2k+1}{\Ilat{\frac12+\mathbb Z}}{0}\\
&\qquad-4\sum_j\frac{b_j}{2M+2j+1}\,\Tx{2k+1}{\IZ}{M+j}\\
&=(2^{2k+1}-2)\,R_{M,p}(\tfrac12)\,\Tx{2k+1}{\IZ}{0}\\
&\qquad+4\sum_j\frac{b_j}{2M+2j+1}\,\Th{2k+1}{\IZ}{0}{M+j}>0 .
\end{align*}
If $Cb_j\in\mathbb Z$ for every $j$, then $C\,d_{2M+2D+1}\,d_{M+D}^{\,2k+1}$ clears both rational coefficients; the Trailmoments of $\Ilat{\frac12+\mathbb Z}$ are read in those of $\IZ$ by Theorem~\ref{thm:latticefree}(i).
\end{theorem}

% Reference 20: cor:twobretaslot; source line 10703
\begin{corollary}[Arguments of two Beta corners]\label{cor:twobretaslot}
For $p(v)=(1-v)^q$, $q\in\mathbb N$, the argument of Theorem~\ref{thm:pureslot} is
\[
R_{M,(1-v)^q}(u)=\Gst21{u+M,\ q+1}{u+M+q+1}+\Gst21{1-u+M,\ q+1}{2-u+M+q} .
\]
Its words are positive and carry one Trailmoment of the counting input, of the odd index $2k+1$, at every $k$. For fixed $M$, $\operatorname{crit}_k[R_{M,(1-v)^q}]\sim\frac{4\,\Gamma(M+1)}{(2k)!}\,q^{-M-1}(\log q)^{2k}$ as $q\to\infty$.
\end{corollary}

% New: proved in the last part.
\begin{proposition}[Poles at the half-integers]\label{prop:ra-halfpoles}
For integers $k\ge1$ and $l\ge0$,
\[
\frac1{2\pi i}\int_{(1/2)}\frac1{u-\frac12-l}\sum_{m\in\mathbb Z}\frac{du}{(u+m)^{2k+1}}=-\Tx{2k+1}{\Ilat{\frac12+\mathbb Z}}{l},
\]
the integrand being regular on the line at $l=0$; and for $l\ge1$, with the critical-line word of Definition~\ref{def:critical},
\[
\operatorname{crit}_k\Big[\frac1{u-\frac12-l}\Big]=-\operatorname{crit}_k\Big[\frac1{u-\frac12+l}\Big]=-\frac1l\,\Th{2k+1}{\Ilat{\frac12+\mathbb Z}}{0}{l}.
\]
\end{proposition}

So the critical-line word of a simple pole at a half-integer other than $\frac12$ is a head of the lattice input on the coset $\frac12+\mathbb Z$, a rational number, and with Proposition~\ref{prop:critical} the word is read on every polynomial plus simple poles at these points. The rational arguments of Theorem~\ref{thm:pureslot} have their simple poles at the integers; the half-integers give heads and tails on the coset $\frac12$.

% Reference 20: prop:slotrate; source line 10710
\begin{proposition}[The decay of the word of an exponential argument]\label{prop:slotrate}
For an integer $k\ge1$ and $t>0$, with the polynomial $c_k(t)=\sum_{l=0}^{2k}\frac{(2k)!}{l!}\,2^{2k-l+1}\,t^l$,
\[
\frac{e^{-t/2}}{(2k)!}\,c_k(t)\le\operatorname{crit}_k\big[e^{t(u-\frac12)}\big]\le\frac{e^{-t/2}}{(2k)!\,(1-e^{-t})}\,c_k(t).
\]
\end{proposition}

A rational argument whose positive measure is the exponential profile of a polynomial nonnegative on $[0,1]$ gives a positive word with one odd Trailmoment of the counting input and a rational part cleared by $d_{2M+2D+1}d_{M+D}^{2k+1}$. The words of exponential arguments decay like $e^{-t/2}$ times a polynomial of degree $2k$, so the words of the arguments of two Beta corners decay like a power of their exponent $q$: they are positive, carry one Trailmoment and are controlled, the raw material of the relation blocks of Chapter~\ref{ch:relations}.

\chapter{The gap and the shape words}\label{ch:arith}

A value read at rational data has a row, its coefficients on Trailmoments of the counting input and a rational part (Definition~\ref{def:row}). The gap theorem turns a sequence of nonzero values whose rows are cleared by integers that grow more slowly than the values decay into irrationality, and the shape words supply such sequences. They are integrals over a cube against one letter; opened by one share they close on the stacked Gamma $\Gst40{a,\ a,\ 1-a,\ 1-a}{}$ of Lemma~\ref{lem:resonantpair}, so their rows are explicit, cleared by powers of the clearing integers $d_n$ and carried by the Trailmoments $\Tx{s}{\IZ}{0}$ of index at most one more than the number of shapes. Beta ports inserted on one shape keep the words positive while they move the row, filters built from recurrences remove unwanted coordinates, and Beta ports whose order grows with the depth trade a slower decay for a selected coordinate.

\section{Clearing and the gap}

% Reference 20: lem:lcm; source line 10723
\begin{lemma}[The growth of the clearing integers]\label{lem:lcm}
For integers $n\ge0$ and $s\ge1$, $d_n^{\,s}$ is the least positive integer $D$ with $D\,\Th{s}{f\,\IZ}{0}{n}\in\mathbb Z$ for every integer density $f$, and
\[
\limsup_{n\to\infty}\frac{\log d_n}n\le c_{\mathrm{lcm}}=\frac{14}{25}\log2+\frac9{25}\log3+\frac15\log5=1.10555\ldots<1.106 .
\]
\end{lemma}

% Reference 20: thm:gap; source line 10733
\begin{theorem}[The gap]\label{thm:gap}
Let $v_1,\dots,v_m$ be real, and let real numbers $x_n=\sum_iA_{n,i}v_i-B_n$, with $A_{n,i},B_n\in\mathbb Q$, have positive integers $D_n$ clearing all these coefficients, with $0<|D_nx_n|\to0$.
\begin{enumerate}[label=(\roman*)]
\item Some $v_i$ is irrational.
\item For $m=1$, $A_{n,1}\neq0$ for all large $n$.
\end{enumerate}
The hypothesis $x_n\neq0$ holds in particular when $x_n$ is the value of a word with a strictly positive read.
\end{theorem}

An integer of modulus below one is zero; nothing else is used. What the theory contributes is the two other hypotheses: positivity, which keeps the cleared values from vanishing, and rows whose clearing integers grow more slowly than the values decay.

\section{Shape words with polynomial insertions}

% Reference 20: thm:insertclosing (the word); source line 10747
\begin{definition}[Shape words]\label{def:shapeword}
For integers $n\ge0$ and $k\ge2$ let $L=\bar x_0+x_0x_1\cdots x_k$ on the cube $[0,1]^{k+1}$ with coordinates $x_0,\dots,x_k$. The \emph{shape word} of depth $n$ with $k$ shapes of a polynomial insertion $P\in\mathbb Q[x_0,\dots,x_k]$ is
\[
I_n(k;P)=\int_{[0,1]^{k+1}}P(x)\,\frac{\prod_{i=0}^k(x_i\bar x_i)^n}{L^{n+1}}\,dx=\int_{[0,1]^{k+1}}P\,f_k^n\,\frac{dx}L,\qquad f_k=\frac{\prod_{i=0}^kx_i\bar x_i}{L},
\]
with the \emph{base} $f_k$. For a monomial $x^p$, $p=(p_0,\dots,p_k)$, the \emph{share function} is
\[
g_{n,p}(a)=\frac{(n!)^{k-2}\,(1-a)_n^2\,(a)_{p_0}}{(n+1)_{p_0}\prod_{i=1}^k(a+p_i)_{n+1}},
\]
and for $P=\sum_p\big([x^p]P\big)x^p$ it is $g_{n,P}=\sum_p\big([x^p]P\big)\,g_{n,p}$, with the coefficients $c_{n,h,s}(P)$ of its principal parts $c_{n,h,s}(P)(a+h)^{-s}$ at the poles $a=-h$ and the \emph{column sums} $C_{n,s}(P)=\sum_hc_{n,h,s}(P)$. The \emph{port factor} is $\mathfrak a_k=x_0\,x_1\bar x_1\,x_2\cdots x_k/L$, and $M_k(\eta)=\sup_{(0,1)^{k+1}}f_k\mathfrak a_k^\eta$ for $\eta\ge0$.
\end{definition}

% Reference 20: thm:insertclosing; source line 10747
\begin{theorem}[Shape words: the one-share closing]\label{thm:insertclosing}
Let $P$ have degree at most $D$, with $p_0\le(k-2)(n+1)$ on every monomial $x^p$ of $P$.
\begin{enumerate}[label=(\roman*)]
\item \emph{Share side.} For a monomial, on $0<\re a=c<n+1$,
\begin{align*}
I_n(k;x^p)&=\frac1{2\pi i}\int_{(c)}\Gst21{a,\ n+1-a}{n+1}\Gst21{a+p_0,\ n+1-a}{n+1+p_0}\prod_{i=1}^k\Gst21{a+p_i,\ n+1}{a+p_i+n+1}\,da\\
&=\frac1{2\pi i}\int_{(c)}\Gst40{a,\ a,\ 1-a,\ 1-a}{}\,g_{n,p}(a)\,da .
\end{align*}
\item \emph{Closing.} $I_n(k;P)=-\sum_{m\ge1}g_{n,P}'(m)$ (Lemma~\ref{lem:resonantpair}(i)). Every pole of $g_{n,P}$ lies among $0,-1,\dots,-n-D$ with order at most $k$; the integers $1,\dots,n$ are zeros of order at least two of each $g_{n,p}$; $g_{n,p}(a)=O(|a|^{-2})$.
\item \emph{The row.} $C_{n,1}(P)=0$ and
\begin{align*}
I_n(k;P)&=\sum_{h=0}^{n+D}\sum_{s=1}^ks\,c_{n,h,s}(P)\,\Tx{s+1}{\IZ}{h}\\
&=\sum_{s=2}^ks\,C_{n,s}(P)\,\Tx{s+1}{\IZ}{0}-\sum_{h,s}s\,c_{n,h,s}(P)\,\Th{s+1}{\IZ}{0}{h}.
\end{align*}
\end{enumerate}
\end{theorem}

% Reference 20: prop:insertden; source line 10771
\begin{proposition}[An integer certificate for inserted words]\label{prop:insertden}
In Theorem~\ref{thm:insertclosing} let $P$ have integer coefficients. Then $(n+1)_D\,d_{n+D}^{\,k-s}\,c_{n,h,s}(P)\in\mathbb Z$, and $(n+1)_D\,d_{n+D}^{\,k+1}$ clears every coefficient of the row, its rational part included.
\end{proposition}

% Reference 20: thm:intclearing; source line 10778
\begin{theorem}[Integer-valued insertion clearing]\label{thm:intclearing}
In Theorem~\ref{thm:insertclosing} let $P$ have integer coefficients, let $p_{\max}$ be its largest exponent of a shape $x_i$, $i\ge1$, and $p_{0,\max}$ its largest exponent of $x_0$, and let $b_P=\operatorname{lcm}\big\{\binom{n+j}j:\ j\text{ an exponent of }x_0\text{ in }P\big\}$. Then $b_P\,d_{\max(n+p_{\max},\,p_{0,\max})}^{\,k-s}\,c_{n,h,s}(P)\in\mathbb Z$ and $b_P\,d_{\max(n+p_{\max},\,p_{0,\max})}^{\,k+1}$ clears the whole row, and $b_P$ divides $d_{n+p_{0,\max}}$; insertions free of $x_0$ are cleared by $d_{n+p_{\max}}^{\,k+1}$ whatever their degrees. For $k\ge3$ the insertion $x_0^n$ retains the top coordinate:
\[
C_{n,k}(x_0^n)=(-1)^{n(k+1)}\binom{2n}n,\qquad \big[\Tx{k+1}{\IZ}{0}\big]\,I_n(k;x_0^n)=k\,(-1)^{n(k+1)}\binom{2n}n.
\]
\end{theorem}

% Reference 20: thm:portgauge; source line 10921
\begin{theorem}[The exact maximum under a growing Beta port]\label{thm:portgauge}
Fix $k\ge2$, with the base $f_k$ and the port factor $\mathfrak a_k$ of Definition~\ref{def:shapeword}, so that $I_n(k;P)=\int Pf_k^n\,dx/L$. For $q\in\mathbb N$,
\[
I_n\big(k;(-1)^q\Pi^{n+1,n+1}_q(x_1)\big)=\binom{n+q}q\int_{[0,1]^{k+1}}f_k^n\,\mathfrak a_k^q\,\frac{dx}L>0 .
\]
For every real $\eta\ge0$, $f_k\mathfrak a_k^\eta$ has a unique maximum on the open cube: with $a$ the unique root in $0<a<1/(2+\eta)$ of
\[
(1-2a)\big(1-(2+\eta)a\big)(1+\eta a)^{k-1}=(1+\eta)^ka^{k+2},
\]
it is at $x_1=a$, $x_2=\cdots=x_k=\frac{(1+\eta)a}{1+\eta a}$, $x_0=\frac{(1+\eta)a}{1-a}$, with value
\[
M_k(\eta)=\frac{(1-a)^{k-2}}{(1+\eta a)^{k-1}}\,(1-2a)^{1+\eta}\,\big(1-(2+\eta)a\big),
\]
and $0<I_n\big(k;(-1)^q\Pi^{n+1,n+1}_q(x_1)\big)\le\binom{n+q}q\,M_k(q/n)^n\,I_0(k;1)$.
\end{theorem}

% Reference 20: prop:kdepth; source line 10788 (the maximizer at $k=2$ and at $k=4$).
% Proved in the last part: prop:kdepth.
\begin{proposition}[The root rate of the shape words]\label{prop:kdepth}
For $k\ge2$:
\begin{enumerate}[label=(\roman*)]
\item $0<I_n(k;1)\le M_k(0)^n\,I_0(k;1)$ (Theorem~\ref{thm:portgauge} at $q=0$);
\item $\lim_nI_n(k;1)^{1/n}=M_k(0)$, where
\[
M_k(0)=\max_{(0,1)^{k+1}}f_k=\max_{0<a<1}\frac{a^k(1-a)^k}{(1+a^{k/2})^2}\in\Big[\frac{4^{-k}}{(1+2^{-k/2})^2},\ 4^{-k}\Big].
\]
\end{enumerate}
The maximum over $a$ is at the root $a\in(0,\frac12)$ of $a^{k/2+1}+2a=1$, the root of Theorem~\ref{thm:portgauge} at $\eta=0$. For every $k\ge2$, $-\log M_k(0)>(k+1)\,c_{\mathrm{lcm}}$, so that $d_n^{\,k+1}I_n(k;1)\to0$.
\end{proposition}

% Reference 20: thm:T3; source line 10807
\begin{theorem}[An irrational Trailmoment of the counting input]\label{thm:T3}
The shape word $I_n(2;1)$ is positive, and with $a_l=(-1)^l\binom nl\binom{n+l}l$, the coefficient of $x^l$ in $(-1)^n\Pi^{1,1}_n(x)$ (Theorem~\ref{thm:betaport}(v)), and $b_l=2a_l\sum_{j\neq l}\frac{a_j}{j-l}$,
\begin{gather*}
I_n(2;1)=\sum_{l=0}^n\Big(2a_l^2\,\Tx{3}{\IZ}{l}+b_l\,\Tx{2}{\IZ}{l}\Big)\\
=2\big(A_n\,\Tx{3}{\IZ}{0}-B_n\big),\\
A_n=\sum_la_l^2\in\mathbb Z,\qquad d_n^{\,3}B_n\in\mathbb Z,\\
B_n=\sum_la_l^2\,\Th{3}{\IZ}{0}{l}+\sum_la_l\,\Th{2}{\IZ}{0}{l}\sum_{j\neq l}\frac{a_j}{j-l},
\end{gather*}
$I_n(2;1)\le2\,\Tx{3}{\IZ}{0}\,(\sqrt2-1)^{4n}$, and $\Tx{3}{\IZ}{N}\notin\mathbb Q$ for every integer cut $N\ge0$.
\end{theorem}

A shape word is the cube against one letter, and one share opens the letter: the Beta corners of the shapes and of $x_0$ leave a single line integral of stacked Gammas, the stacked Gamma $\Gst40{a,\ a,\ 1-a,\ 1-a}{}$ times a rational share function, which closes on the positive integers. The share function has poles of order at most $k$ at the nonpositive integers, so the row carries the Trailmoments of the counting input of index at most $k+1$ with rational parts cleared by powers of $d_{n+D}$, and the coefficients of the simple poles sum to zero because the share function decays quadratically. At two shapes the row has one coordinate besides the rationals, the decay $(\sqrt2-1)^{4n}$ beats the clearing $d_n^3$, and the gap gives the irrationality of $\Tx{3}{\IZ}{N}$ at every integer cut $N$. With more shapes the decay still beats the clearing, but the row carries several coordinates, and the filters and the growing Beta ports select one (Proposition~\ref{prop:positivefilter}, Theorem~\ref{thm:portbudget}).

\section{Rows}

% Reference 20: prop:terminalrow; source line 10822
\begin{proposition}[The row of a closing in rational terms]\label{prop:terminalrow}
Let a closing give $x_n=\sum_{t\ge t_0}r_n(t)$, the sum over the integers $t\ge t_0$, with $r_n(t)=\sum_{k,s}\tilde c_{n,k,s}(t+k)^{-s}=O(t^{-2})$, $\tilde c_{n,k,s}\in\mathbb Q$, integer pole positions and $h_k=t_0+k-1\ge0$. With $\tilde C_{n,s}=\sum_k\tilde c_{n,k,s}$,
\begin{gather*}
x_n=\sum_{k,s}\tilde c_{n,k,s}\,\Tx{s}{\IZ}{h_k}=\sum_{s\ge2}\tilde C_{n,s}\,\Tx{s}{\IZ}{0}-B_n,\\
B_n=\sum_{k,s}\tilde c_{n,k,s}\,\Th{s}{\IZ}{0}{h_k},
\end{gather*}
the index-one Trailmoments regularized, and $D_n=\operatorname{lcm}_{k,s}\big(\operatorname{den}(\tilde c_{n,k,s})\,d_{h_k}^{\,s}\big)$ clears the row.
\end{proposition}

% Reference 20: prop:primitive; source line 10834
\begin{proposition}[The primitive row]\label{prop:primitive}
For a nonzero rational pair $(A,B)$, a common clearing integer $D$ and $g=\gcd(DA,DB)$, the pair $g^{-1}(DA,DB)\in\mathbb Z^2$ is, up to sign, independent of $D$ and of every nonzero rational rescaling of $(A,B)$, and with $(a,b)=g^{-1}(DA,DB)$, $av-b=(D/g)(Av-B)$ for every real $v$.
\end{proposition}

% Reference 20: cor:insertgap; source line 10841
\begin{corollary}[The top-index criterion]\label{cor:insertgap}
Let $k\ge3$, and suppose there are integer polynomials $P_n$ of degree at most a fixed $D$ with $C_{n,s}(P_n)=0$ for $2\le s<k$ and $I_n(k;P_n)>0$ for all large $n$, $P_n$ possibly changing sign on the cube. If the sums $\|P_n\|_1$ of the absolute values of the coefficients of $P_n$ satisfy
\[
\limsup_n\frac1n\log\|P_n\|_1<-\log M_k(0)-(k+1)\limsup_n\frac1n\log d_n ,
\]
then $\Tx{k+1}{\IZ}{N}\notin\mathbb Q$ for every integer cut $N\ge0$.
\end{corollary}

A closing in rational terms has a row whose clearing is read off its pole positions and orders, and a row is a primitive integer pair up to the clearing. With $k$ shapes the row of $I_n(k;P)$ carries $\Tx{3}{\IZ}{0},\dots,\Tx{k+1}{\IZ}{0}$; an insertion that kills the lower column sums and keeps the word positive at a cost below the margin between the decay and the clearing separates $\Tx{k+1}{\IZ}{0}$. The insertion may change sign: positivity is asked of the word, not of the insertion.

\section{The Beta port ladder and the filters}

% Reference 20: thm:fourshapeladder; source line 10857
\begin{theorem}[The Beta port ladder of the shape words]\label{thm:fourshapeladder}
For the shape word with $k\ge2$ shapes at depth $n\ge1$ write its letter as $L=\bar x_0+x_0x_2\cdots x_k\,x_1$, and in this statement write $P_r=(-1)^r\Pi^{n+1,n+1}_r(x_1)$, $g_{n,r}=g_{n,P_r}$ and $g_n=g_{n,0}$, the share function of the insertion $1$ (Definition~\ref{def:shapeword}).
\begin{enumerate}[label=(\roman*)]
\item $0<I_n(k;P_r)\le\binom{n+r}r\,I_n(k;1)$ (Theorem~\ref{thm:portgauge}, the port factor being at most one).
\item $g_{n,r}(a)=\binom{n+r}r\dfrac{(n+1-a)_r}{(n+1+a)_r}\,g_n(a)$ with $g_n(a)=\dfrac{(n!)^{k-2}(1-a)_n^2}{(a)_{n+1}^k}$, and with $b=n+r+1$, $(r+1)(a+b)\,g_{n,r+1}(a)=b\,(b-a)\,g_{n,r}(a)$.
\item For a function $\varphi$ on the integers and a finite set $\mathcal M$ of positive integers write $S_{n,r}[\varphi]=\sum_{m\in\mathcal M}g_{n,r}(m)\varphi(m)$ and $S'_{n,r}[\varphi]=-\sum_{m\in\mathcal M}g_{n,r}'(m)\varphi(m)$, the sums over the common finite set $\mathcal M$, with functions of the summation variable $m$ inside the brackets. Then, with $b=n+r+1$,
\[
(r+1)\,S'_{n,r+1}\big[(m+b)\varphi\big]=b\,S'_{n,r}\big[(b-m)\varphi\big]+(r+1)\,S_{n,r+1}[\varphi]+b\,S_{n,r}[\varphi].
\]
\item $\dfrac{g_{n+1,r}(a)}{g_{n,r}(a)}=\dfrac{(n+1)^{k-3}(n+r+1)(n+1-a)(n+r+1-a)}{(a+n+1)^{k-1}(a+n+r+1)}$, and with $U(m)=(m+n+1)^{k-1}(m+n+r+1)$ and $V(m)=(n+1)^{k-3}(n+r+1)(n+1-m)(n+r+1-m)$,
\[
S_{n+1,r}[U\varphi]=S_{n,r}[V\varphi],\qquad S'_{n+1,r}[U\varphi]-S_{n+1,r}[U'\varphi]=S'_{n,r}[V\varphi]-S_{n,r}[V'\varphi].
\]
\item For $0\le r\le h$ the insertions $P_r$ have degree at most $h$, their rows are cleared by $d_{n+h}^{\,k+1}$ (Theorem~\ref{thm:intclearing}), and their positive reads cost at most $\binom{n+h}h$ times the base word $I_n(k;1)$.
\item In the selected mark $x_1$ the furthest forward coefficient of Theorem~\ref{thm:portrecurrence}, for the one letter $L$ of exponent $n+1$, is nonzero on the open cube.
\end{enumerate}
\end{theorem}

% Reference 20: prop:portindexstates; source line 10872 (four shapes).
% Proved in the last part: prop:portindexstates.
\begin{proposition}[The insertion index supplies auxiliary states]\label{prop:portindexstates}
Fix $k\ge3$ and $n\ge1$ and, in this statement, write $P_r=(-1)^r\Pi^{n+1,n+1}_r(x_1)$, an insertion of the shape word with $k$ shapes. As functions of $r$, $C_{n,k}(P_r)\in\mathbb Q[r]$ and $C_{n,k-1}(P_r)=-2\,\Th{1}{\IZ}{0}{n+r}\,C_{n,k}(P_r)+q_n(r)$ with $q_n\in\mathbb Q(r)$. Consequently a scalar finite shift recurrence over $\mathbb Q(r)$ annihilating $C_{n,k-1}(P_r)$ for all large $r$ also annihilates $C_{n,k}(P_r)$.
\end{proposition}

% Reference 20: prop:filtershapestep; source line 10879
\begin{proposition}[The shape step is a two-state operator relation]\label{prop:filtershapestep}
For $k\ge2$ and $n\ge0$ let $g_n$ be the share function of the insertion $1$, $g_n(a)=(n!)^{k-2}(1-a)_n^2/(a)_{n+1}^k$, and for integer cuts $0\le N_0<N_1$ write $S_n[\varphi]=\sum_{N_0<m\le N_1}g_n(m)\varphi(m)$ and $S'_n[\varphi]=-\sum_{N_0<m\le N_1}g_n'(m)\varphi(m)$. Then
\begin{gather*}
S_{n+1}\big[(m+n+1)^k\varphi\big]=(n+1)^{k-2}\,S_n\big[(m-n-1)^2\varphi\big],\\
\begin{split}
S'_{n+1}\big[(m+n+1)^k\varphi\big]={}&k\,S_{n+1}\big[(m+n+1)^{k-1}\varphi\big]+(n+1)^{k-2}\,S'_n\big[(m-n-1)^2\varphi\big]\\
&-2(n+1)^{k-2}\,S_n\big[(m-n-1)\varphi\big],
\end{split}
\end{gather*}
and the same holds over all $m>N_0$ whenever the sums converge one by one: for bounded $\varphi$, when $(k-2)(n+1)>1$ in the first identity and $(k-2)(n+1)>0$ in the second.
\end{proposition}

% Reference 20: prop:localfilter; source line 10894
\begin{proposition}[A polynomial filter from a bounded local recurrence]\label{prop:localfilter}
Let unwanted closing coefficients satisfy $c_{n,j}=\ell(n)\,y_{n,j}$ and $y_{n+1,j}=\mathbf A(n)\,y_{n,j}$, $1\le j\le r$, with state vectors $y_{n,j}$ of a fixed length $h$, a $1\times h$ matrix $\ell(n)$ and an $h\times h$ matrix $\mathbf A(n)$ over $\mathbb Q(n)$. Write $\mathbf A=\widehat{\mathbf A}/q$ and $\ell=\hat\ell/s$ with integer polynomial matrices $\widehat{\mathbf A}$, $\hat\ell$ and integer polynomials $q$, $s$, put $\mathbf A^{[0]}=I$, $\mathbf A^{[i]}(n)=\mathbf A(n+i-1)\cdots\mathbf A(n)$, $Q(n)=\prod_{t=0}^{h-1}q(n+t)\prod_{t=0}^hs(n+t)$ and the $1\times h$ matrices $\mathbf r_i(n)=Q(n)\,\ell(n+i)\,\mathbf A^{[i]}(n)$, $0\le i\le h$. If the $\mathbf r_i$ span $\mathbb Q(n)^{1\times h}$ and $a_i(n)$ is $(-1)^i$ times the determinant of the $h\times h$ matrix stacking the $\mathbf r_{i'}$, $i'\neq i$, in order, then the $a_i$ are integer polynomials, not all zero, and
\[
\sum_{i=0}^ha_i(n)\,c_{n+i,j}=0\qquad(1\le j\le r)
\]
away from the integer zeros of the denominators; a target sequence $x_n$ is retained at $n$ exactly when $\sum_ia_i(n)\,x_{n+i}\neq0$. With $D_{\mathbf A}=\max(\deg\widehat{\mathbf A},\deg q)$, $D_s=\deg s$ and $D_\ell=\deg\hat\ell$: $\deg a_i\le h(hD_{\mathbf A}+hD_s+D_\ell)$. For two states, with $\ell=(1,0)$ and $\mathbf A=(\beta_{il})$, the filter is $\beta_{12}c_{n+2}-(\beta^+_{11}\beta_{12}+\beta^+_{12}\beta_{22})c_{n+1}+\beta^+_{12}\det\mathbf A(n)\,c_n=0$, $\beta^+_{il}=\beta_{il}(n+1)$; for one state, $q(n)c_{n+1}=p(n)c_n$ gives the filter $q(n)\operatorname{Sh}-p(n)$, $\operatorname{Sh}$ the shift $n\mapsto n+1$.
\end{proposition}

% Reference 20: prop:jetrecurrence; source line 10905
\begin{proposition}[The finite jet recurrence]\label{prop:jetrecurrence}
If a state has the recurrence $U_{n+1}(\varepsilon)=\mathbf A(n,\varepsilon)\,U_n(\varepsilon)$, regular at $\varepsilon=0$, then with $U_n(\varepsilon)=\sum_jU_{n,j}\varepsilon^j$ and $\mathbf A(n,\varepsilon)=\sum_j\mathbf A_j(n)\varepsilon^j$, $U_{n+1,j}=\sum_{i=0}^j\mathbf A_i(n)\,U_{n,j-i}$ for $0\le j\le J$: a lower block-triangular recurrence for the jets of order at most $J$, with diagonal blocks $\mathbf A_0(n)$.
\end{proposition}

% Reference 20: prop:positivefilter; source line 10912
\begin{proposition}[A recurrence filters the positive word]\label{prop:positivefilter}
Let $x_n=\int X^n\,d\mu$ with $0<X\le X_{\max}<1$ on the support of a finite positive measure $\mu$, and $x_n=\sum_{j=1}^r\tilde C_{n,j}v_j-B_n$ with $\tilde C_{n,j},B_n\in\mathbb Q$. Let $a_0(n),\dots,a_h(n)\in\mathbb Z[n]$, $Q_n(X)=\sum_ia_i(n)X^i$ and $\tilde x_n=\sum_ia_i(n)\,x_{n+i}$, and suppose that for all large $n$, $\sum_ia_i(n)\tilde C_{n+i,j}=0$ for $j\neq j_\ast$ and $Q_n$ has one sign on the support of $\mu$, strictly on a set of positive measure. Then $\tilde x_n=\int X^nQ_n\,d\mu=A_nv_{j_\ast}-\sum_ia_i(n)B_{n+i}$ with $A_n=\sum_ia_i(n)\tilde C_{n+i,j_\ast}$, and $0<|\tilde x_n|\le O(n^{\max_i\deg a_i})\,X_{\max}^n\mu(\mathbb R)$. If positive integers $D_{n,h}$ clear the rows of $x_n,\dots,x_{n+h}$, and $\limsup_n\frac1n\log D_{n,h}<-\log X_{\max}$, then $v_{j_\ast}\notin\mathbb Q$ and $A_n\neq0$ for all large $n$. The sign of $Q_n$ on $(0,X_{\max}]$ is certified when $Q_n(X)=\sum_{l=0}^h\beta_l(n)X^l(X_1-X)^{h-l}$ for a rational $X_1>X_{\max}$ with all $\beta_l(n)\ge0$ and one of them positive.
\end{proposition}

The Beta ports of one shape move the row of a shape word along a ladder: the share function of the Beta port of order $r+1$ is that of order $r$ times a rational factor, so the sums over a common finite set of integers satisfy finite operator relations in the order of the Beta ports and in the depth. Unwanted coordinates whose coefficients follow a bounded local recurrence in the depth are annihilated by a filter with integer polynomial coefficients, read off the determinants of that recurrence, and jets in a regulator follow a block-triangular recurrence with the same diagonal. A filter applied to a positive word stays positive as long as its polynomial keeps one sign on the support of the word's measure, which a nonnegative expansion in the products $X^l(X_1-X)^{h-l}$ certifies, and then the filtered word separates the coordinate it keeps, at the cost of clearing the rows of $h+1$ consecutive values.

\section{Growing Beta ports}

% Reference 20: thm:portbudget; source line 10936 (the coefficient of the top coordinate for even $k$).
% Proved in the last part: thm:portbudget.
\begin{theorem}[An arithmetic budget for growing Beta ports]\label{thm:portbudget}
For $k\ge2$, $\eta>0$ and $q_n=\lfloor\eta n\rfloor$, in this statement write $\ell_k(\eta)=\log M_k(\eta)+(1+\eta)\log(1+\eta)-\eta\log\eta$. Then
\begin{gather*}
\lim_n\frac1n\log I_n\big(k;(-1)^{q_n}\Pi^{n+1,n+1}_{q_n}(x_1)\big)=\ell_k(\eta),\\
\limsup_n\frac1n\log\Big(d_{n+q_n}^{\,k+1}\,I_n\big(k;(-1)^{q_n}\Pi^{n+1,n+1}_{q_n}(x_1)\big)\Big)\le\ell_k(\eta)+(k+1)(1+\eta)\,c_{\mathrm{lcm}} .
\end{gather*}
The integer $d_{n+q_n}^{\,k+1}$ clears the whole row, which stays on the coordinates $\Tx{3}{\IZ}{0}$, \dots, $\Tx{k+1}{\IZ}{0}$ and a rational part, and the coefficient of $\Tx{k+1}{\IZ}{0}$ is
\[
k\binom{n+q_n}{q_n}\sum_{j=0}^n(-1)^{jk}\binom nj^k\binom{n+j}n^2\frac{(n+1+j)_{q_n}}{(n+1-j)_{q_n}},
\]
positive for even $k$.
\end{theorem}

% Reference 20: cor:growingport; source line 10947 (four shapes at the ratio $\eta=3/50$).
% Proved in the last part: cor:growingport.
\begin{corollary}[A growing Beta port with a certified gap]\label{cor:growingport}
For every $k\ge2$ there is a rational $\eta_0>0$ such that $\ell_k(\eta)+(k+1)(1+\eta)\,c_{\mathrm{lcm}}<0$ for every rational $\eta\in(0,\eta_0]$ (Theorem~\ref{thm:portbudget}), the root of Theorem~\ref{thm:portgauge} at each such $\eta$ being certified as a root of an integer polynomial in a rational interval. So with $q_n=\lfloor\eta n\rfloor$ the positive words $I_n\big(k;(-1)^{q_n}\Pi^{n+1,n+1}_{q_n}(x_1)\big)$ satisfy
\[
\limsup_n\frac1n\log\Big(d_{n+q_n}^{\,k+1}\,I_n\big(k;(-1)^{q_n}\Pi^{n+1,n+1}_{q_n}(x_1)\big)\Big)<0,
\]
on the fixed coordinates $\Tx{3}{\IZ}{0}$, \dots, $\Tx{k+1}{\IZ}{0}$ and a rational part, with a positive coefficient of $\Tx{k+1}{\IZ}{0}$ when $k$ is even.
\end{corollary}

A Beta port whose order grows in proportion to the depth pulls the word toward the port factor: the word decays at the maximum value $M_k(\eta)$ of $f_k\mathfrak a_k^\eta$, the binomial of the Beta port adds its entropy $(1+\eta)\log(1+\eta)-\eta\log\eta$, and the clearing grows to $d_{n+q_n}^{k+1}$. Where the sum is negative the cleared words tend to zero on a fixed set of coordinates, and for even $k$ the top coordinate keeps a positive coefficient. For every number of shapes the margin is positive at small rational ratios, by continuity from the base word (Proposition~\ref{prop:kdepth}), each certified by an explicit algebraic root, and the cleared positive words decay exponentially on the coordinates $\Tx{3}{\IZ}{0},\dots,\Tx{k+1}{\IZ}{0}$ with a rational part.

\part{Rays}\label{part:rays}

\chapter{Rays, native gauges and tiers}\label{ch:rays}

A ray moves the data of a family along an integral direction. Along it the values of every insertion obey two relations, integration by parts in every mark and the shift by the stride, wherever the integrals converge and, when the walls of the direction are positive, at every strength as identities of meromorphic functions; these relations make the values a holonomic sequence: they satisfy a recurrence of order at most the dimension of the ray. The complex points at which the direction is native, the native gauges of the ray, carry formal solutions of the relations and bound the dimension from below, and the frame of any recurrence is a power times a stacked Gamma. Every ray with real origin is one function, carried by a pencil, with a branch at the native gauge of its direction, and on the centre ray of a family every sign is a candidate native gauge; a ray along a wall carries the product of the rays of its two faces. What a value can contain is bounded: finitely many values span every value along a ray, and a closing in rational terms carries Trailmoments of index at most the largest order of its poles (Theorem~\ref{thm:rayspace}, Lemma~\ref{lem:polecount}).

\section{The relations along a ray}

% Reference 20: def:ray; source line 10957
\begin{definition}[Ray, stride, insertions, values]\label{def:ray}
Let $F$ be a family on $\cE$. A \emph{ray} of $F$ is the sequence of data $(\mvec_0+k\mvec_1;\ \mathbf s_0+k\mathbf s_1;\ \mathbf t_0+k\mathbf t_1)$, $k\in\mathbb N$, the points at the integer times of the half-line from its \emph{origin} $(\mvec_0;\mathbf s_0,\mathbf t_0)$, which is complex, along its \emph{direction} $(\mvec_1;\mathbf s_1,\mathbf t_1)$, which is integral, $\mvec_1\in\mathbb Z^{\cE}$ and $s_{1j},t_{1j}\in\mathbb N$, both on the balance wall; the walls along it are $\kappa_{0,S}+k\kappa_{1,S}$, with $\kappa_{0,S}$ and $\kappa_{1,S}$ the walls of the origin and of the direction. On the chart $y_{e_0}=0$ of a fixed mark $e_0$, with $dy=\prod_{e\neq e_0}d\alpha_e/\alpha_e$, the \emph{stride} and the \emph{origin form} are
\[
f=\alpha^{\mvec_1}\prod_jK_j^{-s_{1j}}\Psi_j^{-t_{1j}},\qquad \omega_0=\alpha^{\mvec_0}\prod_jK_j^{-s_{0j}}\Psi_j^{-t_{0j}}\,dy ,
\]
$f$ a rational function of degree $0$. The \emph{letter complement} of $F$ is $X_F=\{\alpha\in(\mathbb C^\times)^{\cE}/\mathbb C^\times:\ K_j(\alpha)\Psi_j(\alpha)\neq0\text{ for every }j\}$, and the insertions along the ray are its regular functions, the algebra $\mathbb C[X_F]$, which is generated by the monomials $\alpha^v$, $v\in\mathbb Z^{\cE}$ with $v_{\cE}=0$, and the chances of the bases under the basis and cobasis laws, $\alpha\mapsto\mathbb P_{j,\alpha}(T)=\theta_{j,T}\alpha^T/K_j(\alpha)$ and $\alpha\mapsto\mathbb P^\ast_{j,\alpha}(T)=\theta_{j,T}\alpha^{\cE\setminus T}/\Psi_j(\alpha)$, $T\in\cB_j$. The \emph{values} along the ray are $\langle g\rangle_k=\int_{y_{e_0}=0}g\,f^k\omega_0$ for $g\in\mathbb C[X_F]$, and over a compact real $(n-1)$-dimensional chain $D$ in $X_F$, $\langle g\rangle_{k;D}=\int_Dg\,f^k\omega_0$. The \emph{defects} of the origin and of the direction are the masses less the usages,
\[
c^{(i)}_e=m_{i,e}-\sum_j\big(s_{ij}R_{j,e}+t_{ij}R^\ast_{j,e}\big),\qquad i=0,1,
\]
with the occupation $R_{j,e}$ and the dual occupation $R^\ast_{j,e}$ of table $j$ (Theorem~\ref{thm:innerMGF}).
\end{definition}

% Reference 20: lem:insertions; source line 10965
\begin{lemma}[Insertions]\label{lem:insertions}
\begin{enumerate}[label=(\roman*)]
\item $\mathbb C[X_F]$ is closed under every $\partial_{y_e}$: $\partial_{y_e}\alpha^v=v_e\alpha^v$, $\partial_{y_e}\mathbb P_{j,\alpha}(T)=(\one_T(e)-R_{j,e})\,\mathbb P_{j,\alpha}(T)$ and $\partial_{y_e}\mathbb P^\ast_{j,\alpha}(T)=(\one_{\cE\setminus T}(e)-R^\ast_{j,e})\,\mathbb P^\ast_{j,\alpha}(T)$; the occupations and the defects lie in $\mathbb C[X_F]$.
\item The stride lies in $\mathbb C[X_F]$, and $c^{(1)}_e=\partial_{y_e}\log f$, $c^{(0)}_e+k\,c^{(1)}_e=\partial_{y_e}\log(f^k\omega_0/dy)$.
\item Every value $\langle g\rangle_k$ is a finite combination, with products of weights as coefficients, of periods $P_F$ at data $(\mvec_0+k\mvec_1+v;\ \mathbf s_0+k\mathbf s_1+\mathbf a;\ \mathbf t_0+k\mathbf t_1+\mathbf b)$ on the balance wall, $v\in\mathbb Z^{\cE}$ and $\mathbf a,\mathbf b\in\mathbb N^N$.
\end{enumerate}
\end{lemma}

% Reference 20: thm:raylaws; source line 10976
\begin{theorem}[The relations along a ray]\label{thm:raylaws}
If $\kappa_{1,S}>0$ for every nonempty proper $S$, then for every insertion $g$ the value $\langle g\rangle_k$ converges for all large $k$, and for every mark $e$
\[
\big\langle\big(c^{(0)}_e+k\,c^{(1)}_e+\partial_{y_e}\big)g\big\rangle_k=0,\qquad \langle fg\rangle_k=\langle g\rangle_{k+1} .
\]
\end{theorem}

In Theorem~\ref{thm:raylaws} the first relation is integration by parts in the mark $e$ against the form $f^k\omega_0$, whose logarithmic derivative is the defect; the second is the definition of the stride. Proposition~\ref{prop:rays-strength} carries both to every strength: wherever the integrals converge, whatever the walls of the direction, and, when these walls are positive, as identities of meromorphic functions of $k$. Both act on the insertions, which are closed under the derivatives, so the values of all insertions along a ray form one module over the operators in $k$, and every relation among insertions modulo these two is a recurrence among values.

\section{Native gauges carry solutions}

% Reference 20: def:raygauge; source line 10988
\begin{definition}[Native gauges of a ray; the dimension]\label{def:raygauge}
A \emph{native gauge} of a ray is a point $\alpha_c$ of $X_F$ at which the direction is native, $c^{(1)}(\alpha_c)=0$: a native gauge of the direction at a complex point (Definition~\ref{def:native}); its \emph{native value} is $f(\alpha_c)$, and it is \emph{nondegenerate} when the Hessian of $\log f$ in $y$ modulo $\mathbb C\one$ is nonsingular there. Near a nondegenerate native gauge there are holomorphic coordinates $v\in\mathbb C^{n-1}$ with $\log f=\log f(\alpha_c)-\frac12\sum_iv_i^2$; the image $D_c$ of a real ball $\{v\in\mathbb R^{n-1}:|v|\le\varepsilon\}$ is a \emph{steepest-descent disc}. The \emph{relation space} $\mathcal L\subseteq\mathbb C[X_F]\otimes\mathbb C(k)$ is the $\mathbb C(k)$-span of the elements $(c^{(0)}_e+k\,c^{(1)}_e+\partial_{y_e})g$, and the \emph{dimension} of the ray is $\dim_{\mathbb C(k)}\big(\mathbb C[X_F]\otimes\mathbb C(k)\big)/\mathcal L$.
\end{definition}

% New: proved in the last part.
\begin{proposition}[The relations at every strength]\label{prop:rays-strength}
Let a ray of a family $F$ be given. For an insertion $g=\alpha^u\prod_jK_j^{-a_j}\Psi_j^{-b_j}$ with $u\in\mathbb Z^{\cE}$, $\mathbf a,\mathbf b\in\mathbb N^N$ and $u_{\cE}=\sum_j(a_jr_j+b_jr^\ast_j)$, such insertions spanning $\mathbb C[X_F]$, write $\kappa_U(g)=u_U-\sum_j\big(a_jd_j(U)+b_jd^\ast_j(U)\big)$ for the walls of its data $(u;\mathbf a,\mathbf b)$ (Definition~\ref{def:family}).
\begin{enumerate}[label=(\roman*)]
\item Whatever the walls of the direction, at every complex $k$ with $\re\kappa_{0,U}+\kappa_U(g)+\re(k)\,\kappa_{1,U}>0$ for every nonempty proper $U$ the values $\langle g\rangle_k$ and $\big\langle(c^{(0)}_e+k\,c^{(1)}_e+\partial_{y_e})g\big\rangle_k$ converge absolutely, and the second is $0$ for every mark $e$. So if polynomials $p_0,\dots,p_h\in\mathbb C[k]$, an insertion $g'$ and finite combinations $G_e$ of insertions of the above kind, with coefficients in $\mathbb C(k)$, satisfy
\[
\sum_{j=0}^hp_j(k)\,f^jg'=\sum_{e\in\cE}\big(c^{(0)}_e+k\,c^{(1)}_e+\partial_{y_e}\big)G_e ,
\]
an explicit element of the relation space $\mathcal L$ (Definition~\ref{def:raygauge}), a \emph{certificate} of the recurrence, then $\sum_jp_j(k)\langle g'\rangle_{k+j}=0$ at every $k$ off the poles of those coefficients at which the values $\langle g'\rangle_{k+j}$ converge absolutely and the condition above holds for every insertion in the $G_e$.
\item If $\kappa_{1,S}>0$ for every nonempty proper $S$, every value $\langle g\rangle_k$ continues to a meromorphic function of $k\in\mathbb C$, whose poles have order at most $n-1$ and lie among the $k$ with $\kappa_{0,U}+\kappa_U(g)+k\,\kappa_{1,U}\in\{0,-1,-2,\dots\}$ for some nonempty proper $U$. For every element $\sum_lr_l(k)\,g_l$ of $\mathcal L$ the continuation of $\sum_lr_l(k)\langle g_l\rangle_k$ vanishes identically, and $\langle fg\rangle_k=\langle g\rangle_{k+1}$ identically. So a recurrence $\sum_jp_j(k)\langle g\rangle_{k+j}=0$ that follows from the relations, its element $\sum_jp_j(k)f^jg$ lying in $\mathcal L$, holds as an identity of meromorphic functions of $k$, and at every $k$ at which its terms converge.
\end{enumerate}
\end{proposition}

% Reference 20: thm:gaugesolution; source line 10992
\begin{theorem}[A native gauge carries a solution]\label{thm:gaugesolution}
Let $\alpha_c$ be a nondegenerate native gauge of a ray and $D_c$ a steepest-descent disc through it.
\begin{enumerate}[label=(\roman*)]
\item As $k\to\infty$,
\[
\langle g\rangle_{k;D_c}\sim f(\alpha_c)^k\,k^{-(n-1)/2}\sum_{i\ge0}a_i(g)\,k^{-i},
\]
where $a_0(g)=(2\pi)^{(n-1)/2}(\omega_0/dv)(\alpha_c)\,g(\alpha_c)$ with $(\omega_0/dv)(\alpha_c)\neq0$, and each $a_i(g)$ is a finite combination of Taylor coefficients of $g\,\omega_0/dv$ and of $\log f$ at $\alpha_c$. Further $\langle fg\rangle_{k;D_c}=\langle g\rangle_{k+1;D_c}$, and $\big\langle(c^{(0)}_e+k\,c^{(1)}_e+\partial_{y_e})g\big\rangle_{k;D_c}$ is at most a constant times $|f(\alpha_c)|^ke^{-k\varepsilon^2/2}$; so the expansions satisfy both relations to all orders.
\item For a finite set $C$ of distinct nondegenerate native gauges $\alpha_c$, the maps $A^{(c)}:\mathbb C[X_F]\otimes\mathbb C(k)\to\mathbb C((1/k))$, $g\mapsto\sum_la^{(c)}_l(g)\,k^{-l}$, extended $\mathbb C(k)$-linearly, vanish on $\mathcal L$ and are linearly independent over $\mathbb C((1/k))$.
\item The dimension of the ray is at least the number of its nondegenerate native gauges.
\end{enumerate}
\end{theorem}

% Reference 20: cor:symbol; source line 11004
\begin{corollary}[The symbol]\label{cor:symbol}
Let $g$ be an insertion and $p_0,\dots,p_h\in\mathbb C[k]$, not all zero, with $\sum_jp_j(k)\,f^jg\in\mathcal L$. Under the hypothesis of Theorem~\ref{thm:raylaws}, $\sum_jp_j(k)\,\langle g\rangle_{k+j}=0$ for all large $k$. With $d$ the largest degree of the $p_j$ and $c_j$ the coefficient of $k^d$ in $p_j$, the \emph{symbol} $\sum_jc_j\lambda^j$ vanishes at the native value of every nondegenerate native gauge $\alpha_c$ with $g(\alpha_c)\neq0$. So $h$ is at least the number of distinct native values of such native gauges.
\end{corollary}

% New: proved in the last part.
\begin{proposition}[The top and the bottom of a certificate]\label{prop:rays-certtop}
Let $g$ be an insertion and $\sum_{j=0}^hp_j(k)\,f^jg=\sum_e\big(c^{(0)}_e+k\,c^{(1)}_e+\partial_{y_e}\big)G_e$ a certificate (Proposition~\ref{prop:rays-strength}(i)) in which the $p_j\in\mathbb C[k]$ have the largest degree $d\ge1$ and every $G_e=\sum_{i<d}k^i\,G^{[i]}_e$ is a polynomial in $k$ of degree less than $d$ with coefficients $G^{[i]}_e\in\mathbb C[X_F]$. With $c_j$ the coefficient of $k^d$ in $p_j$ and the symbol $\chi(\lambda)=\sum_jc_j\lambda^j$ (Corollary~\ref{cor:symbol}),
\[
\chi(f)\,g=\sum_ec^{(1)}_e\,G^{[d-1]}_e,\qquad \sum_{j=0}^hp_j(0)\,f^jg=\sum_e\big(c^{(0)}_e+\partial_{y_e}\big)G^{[0]}_e .
\]
So $\chi$ vanishes at the native value of every native gauge $\alpha_c$ of the ray with $g(\alpha_c)\neq0$, nondegenerate or not, and for the origin $0$ the bottom piece $\sum_jp_j(0)f^jg$ is the divergence $\sum_e\partial_{y_e}G^{[0]}_e$.
\end{proposition}

% Reference 20: prop:nativeorder; source line 11011 (a connected family with real data and nonnegative exponents).
% Proved in the last part: prop:nativeorder.
\begin{proposition}[The native value bounds the order]\label{prop:nativeorder}
Let a ray of a family have $\kappa_{1,S}>0$ for every nonempty proper $S$. Then $\langle1\rangle_k\neq0$ for all large $k$, positive when the origin is real, and $\langle1\rangle_{k+1}/\langle1\rangle_k\to f(\alpha^\star)$, the native value of the native gauge $\alpha^\star$ of the direction on $\mathbb R^{\cE}_{>0}$, which is the largest value of $f$ there. If $\sum_{j=0}^hp_j(k)\,\langle1\rangle_{k+j}=0$ for all large $k$, with $p_j\in\mathbb Q[k]$ not all zero, then the symbol vanishes at $f(\alpha^\star)$: $f(\alpha^\star)$ is algebraic, $h\ge[\mathbb Q(f(\alpha^\star)):\mathbb Q]$, and every conjugate of $f(\alpha^\star)$ is a root of the symbol. In particular, when $f(\alpha^\star)$ is irrational the values $\langle1\rangle_k$ satisfy no recurrence of order one over $\mathbb Q(k)$, and when it is transcendental none at all.
\end{proposition}

% New: proved in the last part.
\begin{proposition}[The conjugates of the native value of the direction]\label{prop:rays-conjugates}
Let a ray of a family whose weights lie in a subfield $\Bbbk$ of $\mathbb C$ have $\kappa_{1,S}>0$ for every nonempty proper $S$, and let $\alpha^\star$ be the native gauge of its direction on the positive orthant, $\alpha^\star_{e_0}=1$. Then the coordinates of $\alpha^\star$ are algebraic over $\Bbbk$, and for every embedding $\tau$ of $\Bbbk(\alpha^\star)$ in $\mathbb C$ over $\Bbbk$ the point $\tau(\alpha^\star)$ is a nondegenerate native gauge of the ray, of native value $\tau\big(f(\alpha^\star)\big)$. So every conjugate of $f(\alpha^\star)$ over $\Bbbk$ is the native value of a nondegenerate native gauge, the dimension of the ray is at least $[\Bbbk(\alpha^\star):\Bbbk]$, and the symbol of every recurrence for $\langle1\rangle_k$ that follows from the relations vanishes at every conjugate of $f(\alpha^\star)$.
\end{proposition}

% Reference 20: prop:frame; source line 11018
\begin{proposition}[The frame of a recurrence is a power times a stacked Gamma]\label{prop:frame}
Let $\sum_{j=0}^hp_j(k)\,x_{k+j}=0$ for $k\ge k_0$, with $p_j\in\mathbb C[k]$ and $p_0,p_h$ without zeros at the integers $k\ge k_0$, and let $x^{(1)},\dots,x^{(h)}$ be solutions. Their \emph{frame} $C_k=\det\big[x^{(i)}_{k+l}\big]_{1\le i\le h,\ 0\le l<h}$ satisfies $C_{k+1}=(-1)^h\frac{p_0(k)}{p_h(k)}\,C_k$, and writing $\frac{p_0(k)}{p_h(k)}=c\,\frac{\prod_{i=1}^{\deg p_0}(k-a_i)}{\prod_{i=1}^{\deg p_h}(k-b_i)}$,
\[
C_k=C_{k_0}\big((-1)^hc\big)^{k-k_0}\,\Gst{\deg p_0}{\deg p_h}{k-a_1,\ \dots}{k-b_1,\ \dots}\Big/\Gst{\deg p_0}{\deg p_h}{k_0-a_1,\ \dots}{k_0-b_1,\ \dots}.
\]
The \emph{tier} of a sequence over a field $\Bbbk\subseteq\mathbb C$ is the least order of a recurrence with coefficients in $\Bbbk[k]$ and $p_0p_h\neq0$ that it satisfies for large $k$; a sequence of tier one is a power times a stacked Gamma, and so is every frame.
\end{proposition}

% New: proved in the last part.
\begin{proposition}[The growth of the frame]\label{prop:ra-framegrowth}
In Proposition~\ref{prop:frame} let $C_{k_0}\neq0$, let $d_0$ and $d_h$ be the degrees of $p_0$ and $p_h$, and let $c$ be the quotient of their leading coefficients. As $k\to\infty$,
\[
\log|C_k|=k\log|c|+(d_0-d_h)(k\log k-k)+O(\log k).
\]
When $d_0=d_h$ is the largest degree of the $p_j$, the symbol $\sum_jc_j\lambda^j$ (Corollary~\ref{cor:symbol}) has the degree $h$ and $c_0\neq0$, the product of its roots is $(-1)^hc$, and $|C_k|^{1/k}\to|c|$.
\end{proposition}

A native gauge of a ray is a complex point at which the direction uses every mark exactly as its masses ask, the complex counterpart of the native gauge of a family. Through it a steepest-descent disc carries a value that satisfies both relations up to exponentially small terms, and its asymptotic expansion is a formal solution; distinct native gauges give independent solutions, so their number bounds the dimension, and their native values are roots of the symbol of every recurrence that follows from the relations. On the real orthant the native gauge of the direction governs the growth of the values, and an irrational native value forbids a recurrence of order one. Whatever the order, the frame of a recurrence is a power times a stacked Gamma: the sequences of tier one are the bottom of the hierarchy of sequences that the relations produce.

\section{The centre ray of a family}

% Reference 20: centre ray, before thm:pencil; source line 11030 (one connected table).
% Proved in the last part: def:centreray.
\begin{definition}[The centre ray]\label{def:centreray}
Let $F$ be a family of $N$ tables on $n\ge2$ marks, and let strengths $\sigma_j\in\mathbb N$, not all zero, make it connected: every nonempty proper $S$ has $c_j(S)>0$ for some $j$ with $\sigma_j>0$ (Definition~\ref{def:family}); put $\boldsymbol\sigma=(\sigma_j)_j$ and $\sigma=\sum_j\sigma_j$. The \emph{centre ray} of $F$ at the strengths $\boldsymbol\sigma$ is its ray with origin $0$ and direction $(\sigma\one;\boldsymbol\sigma,\boldsymbol\sigma)$. Its stride is
\[
f=\alpha^{\sigma\one}\prod_j\big(K_j\Psi_j\big)^{-\sigma_j}=\prod_j\big(K_j(e^y)K_j(e^{-y})\big)^{-\sigma_j},
\]
its values are $\langle1\rangle_k=\int_{y_{e_0}=0}f^k\,dy=P_F(k\sigma\one;k\boldsymbol\sigma,k\boldsymbol\sigma)$, $k\ge1$, and its direction walls are $\kappa_{1,S}=\sum_j\sigma_jc_j(S)>0$ (Proposition~\ref{prop:centrefamily}).
\end{definition}

% Reference 20: thm:pencil; source line 11033 (the centre ray of one table).
% Proved in the last part: thm:pencil.
\begin{theorem}[A ray is one function]\label{thm:pencil}
Let a ray of a family on $n\ge2$ marks have real origin and $\kappa_{1,S}>0$ for every nonempty proper $S$, let $\alpha^\star$ be the native gauge of its direction on $\mathbb R^{\cE}_{>0}$, with $\alpha^\star_{e_0}=1$ (Theorem~\ref{thm:saddle}), and let $k_0\ge0$ be an integer with $\kappa_{0,S}+k_0\kappa_{1,S}>0$ for every nonempty proper $S$.
\begin{enumerate}[label=(\roman*)]
\item $0<f\le f(\alpha^\star)$ on $\mathbb R^{\cE}_{>0}$, with equality only at $\alpha\propto\alpha^\star$, and $f\le C'e^{-c'\|y\|}$ on the chart for some $c',C'>0$; every value $\langle1\rangle_k$, $k\ge k_0$, converges and is positive.
\item For $|q|<f(\alpha^\star)^{-1}$ the generating function $G(q)=\sum_{k\ge k_0}\langle1\rangle_kq^{k-k_0}$ is
\begin{gather*}
G(q)=\int_{y_{e_0}=0}\frac{f^{k_0}\omega_0}{1-qf}=\int_{y_{e_0}=0}\frac{\alpha^{\mvec_0+k_0\mvec_1}\prod_jK_j^{-s_{0j}-(k_0-1)s_{1j}}\Psi_j^{-t_{0j}-(k_0-1)t_{1j}}}{\prod_jK_j^{s_{1j}}\Psi_j^{t_{1j}}-q\,\alpha^{\mvec_1}}\,dy,\\
G^{(l)}(q)=l!\int_{y_{e_0}=0}\frac{f^{k_0+l}\omega_0}{(1-qf)^{l+1}} ,
\end{gather*}
and the positive integral converges for every real $q<f(\alpha^\star)^{-1}$: the whole ray is carried by the pencil $\prod_jK_j^{s_{1j}}\Psi_j^{t_{1j}}-q\,\alpha^{\mvec_1}$.
\item For $k\ge k_0$ and $h\ge1$,
\[
\det\big[\langle1\rangle_{k+i+i'}\big]_{i,i'=0}^{h-1}=\frac1{h!}\int\prod_{a=1}^hf(y_a)^k\,\frac{\omega_0}{dy}(y_a)\prod_{a<b}\big(f(y_a)-f(y_b)\big)^2\prod_ady_a>0,
\]
and $G$ is not a rational function of $q$.
\end{enumerate}
On the centre ray of a family (Definition~\ref{def:centreray}) one may take $k_0=1$, and $\alpha^\star=\one$.
\end{theorem}

% Reference 20: thm:branch; source line 11053 (the centre ray of one table).
% Proved in the last part: thm:branch.
\begin{theorem}[The branch at the native gauge of the direction]\label{thm:branch}
For a ray as in Theorem~\ref{thm:pencil}, with its generating function $G$ and the native gauge $\alpha^\star$ of its direction, put $A=f(\alpha^\star)^{-1}$, $\varepsilon=1-q/A$ and $a=\frac{n-1}2$, and let $H$ be the Hessian of $-\log f$ at $\log\alpha^\star$ on the chart,
\[
H=\Big(\sum_j\big(s_{1j}\,\kappa_{j,\alpha^\star}(X_e,X_{e'})+t_{1j}\,\kappa_{j,(\alpha^\star)^{-1}}(X_e,X_{e'})\big)\Big)_{e,e'\neq e_0},
\]
the covariances of the indicators under the basis laws and the cobasis laws at $\alpha^\star$ (Definition~\ref{def:weightlaw}, Theorem~\ref{thm:innerMGF}); it is positive definite. As $q\uparrow A$ there are functions $U,V$ holomorphic at $\varepsilon=0$ with
\begin{gather*}
G(q)=U(\varepsilon)+\varepsilon^{a-1}V(\varepsilon)\quad(n\text{ even}),\qquad G(q)=U(\varepsilon)+\varepsilon^{a-1}\log\varepsilon\;V(\varepsilon)\quad(n\text{ odd}),\\
V(0)=\frac{\Gamma(1-a)\,(2\pi)^a\,(f^{k_0-1}\omega_0/dy)(\alpha^\star)}{A\sqrt{\det H}}\quad(n\text{ even}),\\
V(0)=\frac{(-1)^a(2\pi)^a\,(f^{k_0-1}\omega_0/dy)(\alpha^\star)}{\Gamma(a)\,A\sqrt{\det H}}\quad(n\text{ odd}).
\end{gather*}
The derivatives $G^{(l)}(A)$ are finite exactly for $l<a-1$, and
\[
G(q)=\sum_{l<a-1}\frac{G^{(l)}(A)}{l!}\big(q-A\big)^l+V(0)\,\varepsilon^{a-1}\big(1\text{ or }\log\varepsilon\big)+O\big(\varepsilon^{\lceil a-1\rceil}\big).
\]
On the centre ray of a family, with $k_0=1$, $A=\prod_jK_j(\one)^{2\sigma_j}$, the density $f^{k_0-1}\omega_0/dy$ is $1$, and $H=2\sum_j\sigma_j\big(\kappa_{j,\one}(X_e,X_{e'})\big)_{e,e'\neq e_0}$, the covariances under the weight laws.
\end{theorem}

% New: proved in the last part.
\begin{proposition}[The native gauges of a centre ray]\label{prop:rays-centrenative}
On the centre ray of a family at the strengths $\boldsymbol\sigma$ (Definition~\ref{def:centreray}) the defects of the direction are
\[
c^{(1)}_e(\alpha)=\sum_j\sigma_j\big(R_{j,e}(\alpha^{-1})-R_{j,e}(\alpha)\big),
\]
so a point $\alpha$ of $X_F$ is a native gauge of the ray exactly when the occupations weighted by the strengths agree at $\alpha$ and at $\alpha^{-1}$, $\sum_j\sigma_jR_j(\alpha)=\sum_j\sigma_jR_j(\alpha^{-1})$. Its native value is $\prod_j\big(K_j(\alpha)K_j(\alpha^{-1})\big)^{-\sigma_j}$.
\end{proposition}

% Reference 20: thm:signgauges; source line 11067 (the centre ray of one table).
% Proved in the last part: thm:signgauges.
\begin{theorem}[Native gauges at the signs]\label{thm:signgauges}
For the centre ray of a family at the strengths $\boldsymbol\sigma$ (Definition~\ref{def:centreray}), on the chart $\alpha_{e_0}=1$, with the family word $W=\prod_jK_j^{-\sigma_j}\Psi_j^{-\sigma_j}$ at the exponents $(\boldsymbol\sigma,\boldsymbol\sigma)$, so that $f=\alpha^{\sigma\one}W$:
\begin{enumerate}[label=(\roman*)]
\item $f(\alpha^{-1})=f(\alpha)$; the points fixed by $\alpha\mapsto\alpha^{-1}$ are the $2^{n-1}$ \emph{signs} $\eta\in\{\pm1\}^{\cE}$, $\eta_{e_0}=1$; every native gauge of the centre ray is a sign or is paired with its image, a native gauge of the same native value.
\item A sign with $K_j(\eta)=\sum_T\theta_{j,T}\eta^T\neq0$ for every $j$ is a native gauge, of native value $\prod_jK_j(\eta)^{-2\sigma_j}$. In the chart $\alpha_e=\eta_ee^{y_e}$, $y_{e_0}=0$, the Hessian of $-\log f$ at $y=0$ is $2\sum_j\sigma_j\big(\kappa_{j,\eta}(X_e,X_{e'})\big)_{e,e'\neq e_0}$, with $\kappa_{j,\eta}$ the covariances of the indicators under the signed measure $\mathbb P_{j,\eta}(T)=\theta_{j,T}\eta^T/K_j(\eta)$ of total mass one at the sign.
\item The critical polynomials $\sigma W^{-1}-\alpha_e\partial_eW^{-1}$, $e\neq e_0$, vanish identically modulo $\alpha_e^2=1$, and $W^{-1}\equiv\alpha^{\sigma\one}\prod_jK_j^{2\sigma_j}$ there. In $\mathbb Q(\boldsymbol\theta)[\alpha_e:e\neq e_0]/(\alpha_e^2-1)$, $\mathbb Q(\boldsymbol\theta)$ the field of the weights, localized at $\prod_jK_j$, the idempotents $2^{1-n}\prod_{e\neq e_0}(1+\eta_e\alpha_e)$ of the signs with $K_j(\eta)\neq0$ for every $j$ form a basis, and multiplication by $W^{-1}/\alpha^{\sigma\one}$ acts on the idempotent of $\eta$ by $\prod_jK_j(\eta)^{2\sigma_j}$.
\end{enumerate}
\end{theorem}

% Reference 20: thm:separatesigns; source line 11079 (one table).
% Proved in the last part: thm:separatesigns.
\begin{theorem}[Weights that separate the signs]\label{thm:separatesigns}
Let a family of $N$ tables on $n\ge2$ marks be connected at the strengths $\sigma_j\in\mathbb N$, of total $\sigma$, with its centre ray (Definition~\ref{def:centreray}). There are weights $\theta_{j,T}\in\big\{1,\dots,2^{n-1}\big((2n-1)N+(2^{n-1}-1)\sigma\big)+1\big\}$ for which every sign has $K_j(\eta)\neq0$ for every $j$ and is a nondegenerate native gauge of the centre ray, and the $2^{n-1}$ native values $\prod_jK_j(\eta)^{-2\sigma_j}$ are distinct. For such weights the dimension of the centre ray is at least $2^{n-1}$, and every recurrence for $\langle1\rangle_k$ that follows from the relations (Corollary~\ref{cor:symbol}) has order at least $2^{n-1}$.
\end{theorem}

Along a ray with real origin the whole sequence of values is one function of the pencil parameter, holomorphic up to the reciprocal of the native value of the direction. The determinants $\det\big[\langle1\rangle_{k+i+i'}\big]$ of the values are positive integrals of squared differences, so the generating function is not rational; at the native gauge of the direction it branches, with the exponent set by the number of marks and the coefficient by the value of $f^{k_0-1}\omega_0/dy$ there and the determinant of the covariances of the basis and cobasis laws. The centre ray of a family is its period at the centre read at every multiple of its strengths; its native gauge is the table's gauge, where it branches at $\prod_jK_j(\one)^{2\sigma_j}$ with the covariances of the weight laws. Every sign is a candidate native gauge, the signed measures $\mathbb P_{j,\eta}$ playing the part of the weight laws, and with weights chosen to separate the signs every recurrence that the relations produce for the centre ray of a family on $n$ marks has order at least $2^{n-1}$.

\section{What a wall carries}

% Reference 20: lem:wallgauges; source line 11088 (one table).
% Proved in the last part: lem:wallgauges.
\begin{lemma}[Where native gauges can reach a wall]\label{lem:wallgauges}
Let $B$ be a proper bunch with room $R\ni e_0$, $\alpha_B=\rho u$, and in this statement write $K^{[l]}_j=K^{(d_j(B)+l)}_{j,B}(u,\alpha_R)$ and $\Psi^{[l]}_j=\Psi^{(d^\ast_j(B)+l)}_{j,B}(u,\alpha_R)$, so that $K_j=\rho^{d_j(B)}\sum_l\rho^lK^{[l]}_j$ and $\Psi_j=\rho^{d^\ast_j(B)}\sum_l\rho^l\Psi^{[l]}_j$ (Theorem~\ref{thm:wallres}), with the bottom letters $K^{[0]}_j$ and $\Psi^{[0]}_j$.
\begin{enumerate}[label=(\roman*)]
\item At every native gauge $\alpha$ of a direction $(\mvec_1;\mathbf s_1,\mathbf t_1)$ with $s_{1j},t_{1j}\ge0$, $m_{1,B}=\sum_j\big(s_{1j}\sum_T|T\cap B|\,\mathbb P_{j,\alpha}(T)+t_{1j}\sum_T|B\setminus T|\,\mathbb P^\ast_{j,\alpha}(T)\big)$: the direction's mass on $B$ is its usage of $B$.
\item Let native gauges of such directions $\mathbf v_i\to\mathbf v$ have bunch scales $\rho_i\to0$ and $(u_i,\alpha_{R,i})\to(u^\ast,\alpha^\ast_R)$ in the torus, with $\kappa_{1,B}(\mathbf v)>0$. Then some bottom letter vanishes at $(u^\ast,\alpha^\ast_R)$, and if only one does, its gradient in $(u,\alpha_R)$ vanishes there as well: the product $\prod_jK^{[0]}_j\Psi^{[0]}_j$ of the bottom letters vanishes at $(u^\ast,\alpha^\ast_R)$ together with its gradient in $(u,\alpha_R)$.
\end{enumerate}
So the native gauges of a family leave for the wall of a bunch only at singular points of the zero set of the product of its bottom letters.
\end{lemma}

% Reference 20: thm:raywall; source line 11100
\begin{theorem}[A ray along a wall carries its faces]\label{thm:raywall}
Let $B$ be a proper bunch with room $R\ni e_0$ on whose bottom levels the weights factor, and let a ray have its direction on the wall, $\kappa_{1,B}=0$, and $\re\kappa_{0,B}>0$. With $\alpha_B=\rho u$, and writing in this statement $K^{[l]}_j=K^{(d_j(B)+l)}_{j,B}(u,\alpha_R)$ and $\Psi^{[l]}_j=\Psi^{(d^\ast_j(B)+l)}_{j,B}(u,\alpha_R)$ as in Theorem~\ref{thm:wallres},
\[
f=f_B(u)\,f_R(\alpha_R)\,\big(1+\rho\,h_1+O(\rho^2)\big),\qquad h_1=-\sum_j\Big(s_{1j}\frac{K^{[1]}_j}{K^{[0]}_j}+t_{1j}\frac{\Psi^{[1]}_j}{\Psi^{[0]}_j}\Big),
\]
where $f_B$ and $f_R$ are the strides along the ray of the faces $F^B$ and $F^R$ of Theorem~\ref{thm:wallres}, and $\omega_0=\rho^{\kappa_{0,B}-1}d\rho\wedge\omega^B_0\wedge\omega^R_0\,(1+O(\rho))$ with the origin forms of the faces. Let $\alpha^B_c$ and $\alpha^R_c$ be nondegenerate native gauges of the rays of the faces with $h_1(\alpha^B_c,\alpha^R_c)\neq0$, let $\re\big(e^{i\vartheta}h_1\big)<0$ near $(\alpha^B_c,\alpha^R_c)$, and let $D=e^{i\vartheta}[0,\rho_1]\times D_B\times D_R$ with steepest-descent discs $D_B,D_R$ of the faces. Then the values $\langle g\rangle_{k;D}$ satisfy both relations up to terms exponentially smaller than $|f_B(\alpha^B_c)f_R(\alpha^R_c)|^k$, and
\[
\langle1\rangle_{k;D}=\Gamma(\kappa_{0,B})\,\big(-k\,h_1(\alpha^B_c,\alpha^R_c)\big)^{-\kappa_{0,B}}\,\langle1\rangle^{F^B}_{k;D_B}\,\langle1\rangle^{F^R}_{k;D_R}\,\big(1+O(1/k)\big),
\]
with $\langle1\rangle^{F^B}_{k;D_B}$ and $\langle1\rangle^{F^R}_{k;D_R}$ the values of the rays of the faces over their discs (Theorem~\ref{thm:gaugesolution}), and the power taken continuously along $e^{i\vartheta}(0,\infty)$.
\end{theorem}

% Reference 20: def:clusterport; source line 11115
\begin{definition}[The whole cluster port]\label{def:clusterport}
On a bunch scale $\rho$ take letters with \emph{lower coefficients} $\boldsymbol\lambda_i\in\mathbb C^{r_i}$,
\[
L_i(\rho,\boldsymbol\lambda_i)=\sum_{l<r_i}\lambda_{i,l}\,\rho^l+\rho^{r_i}C_i(\rho),\qquad C_i(0)\neq0 ,
\]
with exponents $w_i$, and a word $\rho^{\kappa-1}G(\rho)\prod_iL_i^{-w_i}$, with $G$ and the $C_i$ holomorphic on a neighbourhood of $[0,\infty)$, each $|C_i|$ bounded below by a positive constant there, and the word integrable at $\infty$; put $\kappa^\flat=\kappa-\sum_ir_iw_i$. Fix $b>0$ with every $C_i$ zero-free on $|\rho|\le b$; for $\boldsymbol\lambda$ in a small polydisc each $L_i$ has exactly $r_i$ roots in $|\rho|<b$, the \emph{cluster}, none on $|\rho|=b$ and none on $[b,\infty)$. The \emph{cluster port} of the word is its integral over the \emph{cluster contour}, which comes in along the half-line $[b,\infty)$ of the scale, goes once around $|\rho|=b$ and goes out on the new sheet:
\[
\operatorname{Port}_\kappa\Big[G\prod_iL_i^{-w_i}\Big](\boldsymbol\lambda)=\frac1{2\pi i}\Big[(e^{2\pi i\kappa^\flat}-1)\int_b^\infty\rho^{\kappa-1}G\prod_iL_i^{-w_i}\,d\rho+\oint_{|\rho|=b}\rho^{\kappa-1}G\prod_iL_i^{-w_i}\,d\rho\Big].
\]
\end{definition}

% Reference 20: thm:clusterport; source line 11126
\begin{theorem}[Relations of the whole cluster port]\label{thm:clusterport}
\begin{enumerate}[label=(\roman*)]
\item One turn around $|\rho|=b$ multiplies the word by $e^{2\pi i\kappa^\flat}$; the cluster port does not depend on $b$, is holomorphic in all the lower coefficients jointly on the polydisc, and the contour has no end at which the word is nonzero, so every integral of a $\rho$-derivative over it vanishes.
\item On the stratum $\boldsymbol\lambda=0$ the word is $\rho^{\kappa^\flat-1}h$, with the \emph{stratum word} $h=G\prod_iC_i^{-w_i}$, and
\begin{gather*}
\operatorname{Port}_\kappa\Big[G\prod_iL_i^{-w_i}\Big](0)=\frac{e^{2\pi i\kappa^\flat}-1}{2\pi i}\,\Reg\!\int_0^\infty\rho^{\kappa^\flat-1}h\,d\rho\quad(\kappa^\flat\notin\mathbb Z),\\
\operatorname{Port}_\kappa\Big[G\prod_iL_i^{-w_i}\Big](0)=[\rho^{-\kappa^\flat}]\,h\quad(\kappa^\flat\in\mathbb Z_{\le0}),\qquad\operatorname{Port}_\kappa\Big[G\prod_iL_i^{-w_i}\Big](0)=0\quad(\kappa^\flat\in\mathbb Z_{>0}).
\end{gather*}
\item For multi-indices $\mathbf a_i$, with $A_i=|\mathbf a_i|$ and $D=\sum_i\sum_l(r_i-l)\,a_{i,l}$, the derivative $\partial^{\mathbf a}_{\boldsymbol\lambda}$ of the cluster port at $\boldsymbol\lambda=0$ is $\prod_i(-1)^{A_i}(w_i)_{A_i}$ times the value of (ii) for the stratum word $G\prod_iC_i^{-w_i-A_i}$ at the index $\kappa^\flat-D$; at $\kappa^\flat=0$ it is $\prod_i(-1)^{A_i}(w_i)_{A_i}\cdot\frac1{D!}\,\partial_\rho^D\big[G\prod_iC_i^{-w_i-A_i}\big]_{\rho=0}$.
\end{enumerate}
\end{theorem}

% Reference 20: prop:oneslot; source line 11142 (one letter).
% Proved in the last part: prop:oneslot.
\begin{proposition}[Small lower coefficients on one parameter]\label{prop:oneslot}
Let $C_1,\dots,C_p$ be bounded below by a positive constant on $[0,\infty)$, $G$ and the $C_i$ holomorphic near $0$ and continuous on $[0,\infty)$, $\xi_1,\dots,\xi_p\in\mathbb C\setminus(-\infty,0]$, $\re\kappa>0$ and $\kappa^\flat=\kappa-|\mathbf w|\notin\mathbb Z$, with $\rho^{\kappa-1}G\prod_i(\rho C_i)^{-w_i}$ integrable at $\infty$ together with every $\rho^{\kappa^\flat-|\mathbf a|-1}G\prod_iC_i^{-w_i-a_i}$, $\mathbf a\in\mathbb N^p$.
\begin{enumerate}[label=(\roman*)]
\item For small $\lambda>0$,
\[
\int_0^\infty\rho^{\kappa-1}G(\rho)\prod_i\big(\lambda\xi_i+\rho\,C_i(\rho)\big)^{-w_i}d\rho=A(\lambda)+\lambda^{\kappa^\flat}B(\lambda)
\]
with $A,B$ holomorphic at $0$,
\begin{gather*}
A(\lambda)=\sum_{\mathbf a\in\mathbb N^p}\prod_i\binom{-w_i}{a_i}(\lambda\xi_i)^{a_i}\,\Reg\!\int_0^\infty\rho^{\kappa^\flat-|\mathbf a|-1}G\prod_iC_i^{-w_i-a_i}\,d\rho,\\
B(0)=G(0)\int_0^\infty u^{\kappa-1}\prod_i\big(\xi_i+u\,C_i(0)\big)^{-w_i}\,du ,
\end{gather*}
the last integral continued in $\kappa$ and $\mathbf w$; where every $\re w_i>0$ it is
\[
\prod_iC_i(0)^{-w_i}\,\Gst{2}{p}{\kappa,\ |\mathbf w|-\kappa}{\mathbf w}\,\sB^{\{J\}}_{\mathbf w,\,|\mathbf w|-\kappa;\,I_J}\big((\xi_i/C_i(0))_i\big),\qquad J=\{1,\dots,p\}
\]
(Proposition~\ref{prop:scalebreta}). The analytic part opens with the finite part of the stratum word along the whole scale, and the branch part with $G(0)$ times the scale integral of the letters linearized at the wall.
\item When $G$ and the $C_i$ are holomorphic on a neighbourhood of $[0,\infty)$, the cluster port of Definition~\ref{def:clusterport} for these letters, $r_i=1$ with the lower coefficients $\lambda\xi_i$, is $\frac{e^{2\pi i\kappa^\flat}-1}{2\pi i}A(\lambda)$.
\end{enumerate}
\end{proposition}

At a wall the ray splits: the stride is the product of the strides of the two faces up to a first-order correction in the scale of the bunch, and integrating that scale against the correction contributes one Gamma of the origin wall, so a ray along a wall carries the product of the native values of its faces. The native gauges of a family reach a wall only where the product of its bottom letters is singular. When the bottom letters of a bunch are perturbed by lower coefficients, the scale integral is read on the cluster contour, which surrounds the whole cluster of roots at once: the cluster port is holomorphic in the lower coefficients, its value on the stratum where they vanish is the finite part of the stratum word, and its derivatives there are the finite parts of the stratum words with raised exponents. With small lower coefficients on one parameter the scale integral splits into an analytic part, the cluster port up to a constant, and a branch part led by the scale integral of the letters linearized at the wall, a stacked Gamma times a Breta of one bunch.

\section{What a value can contain}\label{sec:valuespace}

% Reference 20: thm:rayspace; source line 10435 (rational weights, origin and direction).
% Proved in the last part: thm:rayspace.
\begin{theorem}[The values along a ray]\label{thm:rayspace}
Let a ray of a family have weights and origin in a subfield $\Bbbk$ of $\mathbb C$, with $\kappa_{1,S}>0$ for every nonempty proper $S$.
\begin{enumerate}[label=(\roman*)]
\item The dimension of the ray is finite: it is the dimension over $\mathbb C(k)$ of the top cohomology of the twisted algebraic de Rham complex of the letter complement $X_F$, which is finite. It is at least the number of nondegenerate native gauges of the ray (Theorem~\ref{thm:gaugesolution}). The native gauges are the critical points of the stride $f$, a product of powers of the letters with the exponents of the direction; for a direction outside a countable union of proper subvarieties they are $(-1)^{n-1}\operatorname{eu}(X_F)$ in number, $\operatorname{eu}$ the Euler characteristic, and all nondegenerate, so that the dimension is then at least $(-1)^{n-1}\operatorname{eu}(X_F)$. An integral direction can lie in that exceptional set.
\item For every insertion $g$ with coefficients in $\Bbbk$ there are $p_0,\dots,p_h\in\Bbbk[k]$, $p_h\neq0$, $h$ at most the dimension, with $\sum_jp_j(k)\,\langle g\rangle_{k+j}=0$ for all large $k$.
\item So there is $k_1$ such that every value $\langle g\rangle_k$, $k\ge k_1$, is a $\Bbbk$-linear combination of the $h$ values $\langle g\rangle_{k_1},\dots,\langle g\rangle_{k_1+h-1}$: the values along a ray span a $\Bbbk$-space of dimension at most that of the ray, and whatever Trailmoments those $h$ values contain, no later value contains others.
\end{enumerate}
\end{theorem}

% Reference 20: lem:polecount; source line 10447 ($r\in\mathbb Q(k)$).
% Proved in the last part: lem:polecount.
\begin{lemma}[Index by pole count]\label{lem:polecount}
Let $\Bbbk$ be a subfield of $\mathbb C$, and let $r\in\Bbbk(k)$ have all its poles at negative integers, with largest order $p$, and $r(k)=O(k^{-2})$. Then $\sum_{k\ge0}r(k)\in\Bbbk+\sum_{j=2}^p\Bbbk\,\Tx{j}{\IZ}{0}$. If moreover $r(-M-k)=-r(k)$ for an integer $M$, only odd $j$ occur. If the poles lie in one coset $-b+\mathbb Z$, $b\in\mathbb Q$, the same holds with the Trailmoments of the lattice input $\Ilat{b+\mathbb Z}$ at its cuts; poles in several cosets leave, besides, differences of the index-one values of the cosets.
\end{lemma}

A closing whose terms are rational functions of the lattice index brings Trailmoments of the counting input of index at most the largest order of its poles, and a reflection of the terms keeps only the odd indices. With Theorem~\ref{thm:rayspace} this is what bounds the complexity of a family along a ray: the space of values is finite-dimensional, spanned by finitely many of them, and where a closing writes the values as sums of rational terms the index is read off the orders of the poles, which along the ray can be bounded once for all its points.

% ---- New chapter: the two ends of a ray (inserted after ch:rays) ----
\chapter{The two ends of a ray}\label{ch:twoends}

The values of the insertions along a ray form one module over the rational functions of the strength $k$, and masters write the step from one strength to the next as one rational matrix, the one-step matrix. Finitely many relations of the period write the step on the Laurent part of the module, which is all of it when the direction uses every letter: on a box of insertions bounded by the powers of the letters in their denominators and by their growth against the reciprocal stride, one elimination over the rational functions of the strength gives a spanning set of that part, the matrix of the step on it and a certificate, and a search through the boxes stops. A family of functionals that respects the two relations of the ray is a solution of the module, and the one-step matrix is fixed by its solutions: by a full set of them outright, and by fewer exactly when they annihilate no submodule. The solutions come from the two ends of strength. At the high end the sources of the stride, components of its critical locus, carry formal solutions: one for a nondegenerate native gauge, one for each class of paths at a native gauge of corank one, one for every cycle on a clean family of native gauges. At the low end the values continue to the nonpositive integers, where their top principal parts are the constant terms of the powers of the reciprocal stride: the closings, one exact solution read on the compact torus. The symmetries of the ray, built from the automorphisms and the exchanges of the family and from the Cremona, act on the module; the values live on its invariant part, every orbit of nondegenerate native gauges contributes to it, and a degenerate native gauge only through the classes on which its stabilizer acts trivially. At a lean of a family with an exchange the exchange survives only composed with the Cremona: the integral keeps a symmetry that carries the basis half of a class of the centre onto its cobasis half, while the reading, which is oriented, separates the halves. Over the formal series at the high end the module splits by native values into pieces, each mapped to itself by the symmetries, and the closings glue the pieces back into the rational matrix. The closings, the values and the expansion at the native gauge of the direction vanish on the same elements, the positive orthant and the torus through that gauge meeting there and nowhere else: the closings determine the one-step matrix exactly where the values do, which is where that gauge sees every submodule, and a submodule that it does not see is invisible to the integral itself. Such a submodule can come from a table one mark larger: when the list is the contraction of an identically self-dual table, the integral over the contracted scale makes the automorphisms of the larger table act on the module, and the closings see only what they fix; each orbit of nondegenerate native gauges of the larger table that the symmetries of the ray split into $s$ orbits adds $s-1$ to the dimension of what the closings miss, and at the centre of every $\binom{[2r+1]}r$, $r\ge2$, the signs of $\binom{[2r+2]}{r+1}$ with two scales $-1$ split in two, with every symmetry of the ray of the four kinds in the group. On a wheel the letters are traces of one holonomy in $SL_2$, and the sources of its centre ray are read off the critical locus of the pair of traces and the stratum where both holonomies are $-I$.

\section{The module of a ray}\label{sec:module}

% New: proved in the last part.
\begin{definition}[The module of a ray; masters; the one-step matrix]\label{def:raymodule}
Let a ray of a family $F$ have the stride $f$, the origin form $\omega_0$ and the relation space $\mathcal L$ (Definitions~\ref{def:ray} and~\ref{def:raygauge}). The \emph{module} of the ray is the $\mathbb C(k)$-vector space
\[
\mathfrak M=\big(\mathbb C[X_F]\otimes\mathbb C(k)\big)/\mathcal L ,
\]
whose dimension is the dimension of the ray, finite when $\kappa_{1,S}>0$ for every nonempty proper $S$ (Theorem~\ref{thm:rayspace}(i)); $[g]$ is the class of $g\otimes1$. The \emph{shift} is $\operatorname{Sh}[g\otimes r(k)]=[fg\otimes r(k+1)]$; it is well defined and bijective, and $\operatorname{Sh}(r(k)\,m)=r(k+1)\operatorname{Sh}(m)$. A \emph{submodule} is a $\mathbb C(k)$-subspace $\mathfrak N$ with $\operatorname{Sh}(\mathfrak N)=\mathfrak N$, and $\mathfrak N\neq0$ is \emph{irreducible} when its only submodules are $0$ and $\mathfrak N$. \emph{Masters} of a finite-dimensional submodule $\mathfrak N$ are insertions $b_1,\dots,b_{\mathfrak m}$ whose classes form a basis of $\mathfrak N$, $\mathfrak m=\dim_{\mathbb C(k)}\mathfrak N$, and their \emph{one-step matrix} $\mathbf B_{\mathbf b}\in GL_{\mathfrak m}(\mathbb C(k))$ is the matrix of $\operatorname{Sh}^{-1}$ on them,
\[
[f^{-1}b_i]=\sum_{l=1}^{\mathfrak m}\mathbf B_{\mathbf b,il}(k)\,[b_l],\qquad 1\le i\le\mathfrak m .
\]
The \emph{cyclic submodule} of $[1]$ is the span of the $\operatorname{Sh}^i[1]$, $i\in\mathbb Z$. When it is finite-dimensional, of dimension $h$, $[1],\operatorname{Sh}[1],\dots,\operatorname{Sh}^{h-1}[1]$ is a basis of it, and the \emph{law} of $[1]$ is the relation
\[
\sum_{i=0}^hp_i(k)\operatorname{Sh}^i[1]=0,\qquad p_i\in\mathbb C[k]\text{ without a common factor},
\]
unique up to a constant factor; $p_0p_h\neq0$.
\end{definition}

The law of $[1]$ is the shortest relation among the strides of $1$ that the relations of the ray produce, so by Corollary~\ref{cor:symbol} the values obey it, $\sum_ip_i(k)\langle1\rangle_{k+i}=0$ for all large $k$; its order $h$ is at most the dimension of the ray.

% New: proved in the last part.
\begin{definition}[Solutions]\label{def:raysolution}
The \emph{spaces of germs} are: the sequences on the integers $k\ge k_0$, for some $k_0$, two of them identified when they agree for all large $k$, the \emph{germs at $+\infty$}; the same on the integers $k\le k_0$, the \emph{germs at $-\infty$}; and, for $\lambda\in\mathbb C^\times$ and $\rho\in\mathbb C$, the \emph{formal series} $\lambda^kk^\rho\,\mathbb C((1/k))$. Each carries the multiplication by $\mathbb C(k)$ and the shift $\varsigma$: $(\varsigma x)(k)=x(k+1)$ on germs, and $\varsigma\big(\lambda^kk^\rho\varphi(k)\big)=\lambda^kk^\rho\cdot\lambda\big(1+\tfrac1k\big)^\rho\varphi(k+1)$ on formal series, re-expanded in powers of $1/k$; in each case $\varsigma(rx)=r(k+1)\,\varsigma(x)$ for $r\in\mathbb C(k)$. A \emph{solution} of the ray with values in a space of germs $\mathcal G$, or in a finite direct sum of spaces of formal series, is a $\mathbb C(k)$-linear map $V:\mathfrak M\to\mathcal G$ with
\[
V\circ\operatorname{Sh}=\varsigma\circ V .
\]
A solution with values in $\lambda^kk^\rho\,\mathbb C((1/k))$ has the \emph{type} $(\lambda,\rho+\mathbb Z)$ and the \emph{normalization} $\hat V=\lambda^{-k}k^{-\rho}\,V:\mathfrak M\to\mathbb C((1/k))$. For masters $\mathbf b$ the \emph{solution vector} of $V$ is $v_V=(V[b_1],\dots,V[b_{\mathfrak m}])$. The \emph{annihilator} of a set $\mathcal S$ of solutions is $\operatorname{Ann}(\mathcal S)=\{m\in\mathfrak M:\ V(m)=0\text{ for every }V\in\mathcal S\}$.
\end{definition}

When $\kappa_{1,S}>0$ for every nonempty proper $S$, the values along the ray are a solution with values in germs at $+\infty$, $V[g]=\langle g\rangle_k$ (Theorem~\ref{thm:raylaws}). At a nondegenerate native gauge $\alpha_c$ the expansion $f(\alpha_c)^kk^{-(n-1)/2}\sum_ia_i(g)k^{-i}$ of Theorem~\ref{thm:gaugesolution} is a solution of type $\big(f(\alpha_c),-\frac{n-1}2+\mathbb Z\big)$, written $A_c$; its normalization is the map $A^{(c)}$ there.

% New: proved in the last part.
\begin{lemma}[Solutions of distinct types]\label{lem:types}
Let $A_1,\dots,A_r$ be solutions with values in formal series such that, for every type, those of that type among them are linearly independent over $\mathbb C$. Then their normalizations are linearly independent over $\mathbb C((1/k))$, and $\dim_{\mathbb C(k)}\mathfrak N\ge r$ for every submodule $\mathfrak N$ on which their restrictions are, type by type, linearly independent over $\mathbb C$.
\end{lemma}

% New: proved in the last part.
\begin{theorem}[What solutions determine]\label{thm:determine}
Let $\mathbf b$ be masters of a submodule $\mathfrak N$, $\mathbf B=\mathbf B_{\mathbf b}$, and let $\mathcal S$ be a set of solutions.
\begin{enumerate}[label=(\roman*)]
\item Every solution $V$ has $\varsigma^{-1}v_V=\mathbf B\,v_V$, that is $v_V(k-1)=\mathbf B(k)\,v_V(k)$. When the values along the ray are a solution, $\langle b_i\rangle_{k-1}=\sum_l\mathbf B_{il}(k)\,\langle b_l\rangle_k$ for all large $k$.
\item $\operatorname{Ann}(\mathcal S)$ is a submodule.
\item A matrix $\mathbf B'\in\operatorname{Mat}_{\mathfrak m}(\mathbb C(k))$ has $\varsigma^{-1}v_V=\mathbf B'v_V$ for every $V\in\mathcal S$ if and only if every row of $\mathbf B'-\mathbf B$, read as the element $\sum_l(\mathbf B'-\mathbf B)_{il}[b_l]$ of $\mathfrak N$, lies in $\operatorname{Ann}(\mathcal S)$. So $\mathbf B$ is the only such matrix exactly when $\operatorname{Ann}(\mathcal S)\cap\mathfrak N=0$.
\item Let $\operatorname{Ann}(\mathcal S)\cap\mathfrak N=0$, write $\mathbf B=P_0/Q_0$ with $Q_0$ the monic least common denominator of its entries, and put $D_0=\max(\deg P_0,\deg Q_0)$. For $D\ge0$ the equations $Q(k)\,\varsigma^{-1}v_V=P(k)\,v_V$, $V\in\mathcal S$, in pairs $(P,Q)$ of a matrix and a scalar of polynomials of degree at most $D$, read at every integer beyond a bound past which (i) holds for a germ and at every power of $k$ for a formal series, are linear and homogeneous in the coefficients of $(P,Q)$. Their solutions are the pairs $(cP_0,cQ_0)$ with $c\in\mathbb C[k]$ of degree at most $D-D_0$: only $0$ for $D<D_0$, one line at $D=D_0$. For each $D$ finitely many of the equations have the same solutions as all of them.
\end{enumerate}
\end{theorem}

% New: proved in the last part.
\begin{corollary}[One solution of an irreducible module]\label{cor:onesolution}
If $\mathfrak N$ is irreducible, every solution that does not vanish on $\mathfrak N$ determines $\mathbf B_{\mathbf b}$ in the sense of Theorem~\ref{thm:determine}(iii). If the cyclic submodule of $[1]$ is finite-dimensional and irreducible, then for every solution $V$ with $V[1]\neq0$, $V[1]$ satisfies the law of $[1]$, $\sum_ip_i(k)\,\varsigma^iV[1]=0$, and satisfies no relation $\sum_{i<h}q_i(k)\,\varsigma^iV[1]=0$ with $q_i\in\mathbb C(k)$ not all zero.
\end{corollary}

Masters are a basis of a submodule, the module itself or its invariant part (Theorem~\ref{thm:invariant}), while a box of insertions reduces to a spanning set of a submodule that contains $[1]$, which is a set of masters exactly when it is independent (Section~\ref{sec:box}); the one-step matrix is the matrix of the inverse shift, so that the values of the masters step down in strength by $\mathbf B_{\mathbf b}$. A solution is anything that respects the two relations: the values over the orthant, the expansion at a native gauge, the closings on the torus (Theorem~\ref{thm:closings}). Each solution gives $\mathfrak m$ linear equations on the numerator and the denominator of $\mathbf B_{\mathbf b}$ for every integer or every power of $k$, and the equations of a set of solutions determine $\mathbf B_{\mathbf b}$ exactly when the solutions annihilate no submodule: a submodule annihilated by all of them is a direction in which a row of the matrix can move without any of them noticing. An irreducible module is determined by any one of its nonzero solutions, and the law of $[1]$ by any solution that does not vanish at $[1]$ when the cyclic submodule of $[1]$ is irreducible.

% New: proved in the last part.
\begin{proposition}[The law of {$[1]$} read from a one-step matrix]\label{prop:rays-lawrows}
Let insertions $b_1=1,b_2,\dots,b_{\mathfrak m}$ have classes that span a submodule $\mathfrak N$, and let $\mathbf B\in\operatorname{Mat}_{\mathfrak m}(\mathbb C(k))$ have $[f^{-1}b_i]=\sum_l\mathbf B_{il}(k)[b_l]$ for every $i$; for masters, $\mathbf B=\mathbf B_{\mathbf b}$. Put $w_0=(1,0,\dots,0)$ and $w_i(k)=w_{i-1}(k-1)\,\mathbf B(k)$, so that $w_i(k)=w_0\,\mathbf B(k-i+1)\cdots\mathbf B(k)$.
\begin{enumerate}[label=(\roman*)]
\item $\operatorname{Sh}^{-i}[1]=\sum_lw_{i,l}(k)[b_l]$. So every solution $V$ has $\varsigma^{-i}V[1]=w_i\,v_V$, with $v_V=(V[b_1],\dots,V[b_{\mathfrak m}])$ as in Definition~\ref{def:raysolution}, and when $\kappa_{1,S}>0$ for every nonempty proper $S$, $\langle1\rangle_{k-i}=\sum_lw_{i,l}(k)\langle b_l\rangle_k$ as meromorphic functions of $k$.
\item Let $h$ be the least index at which $w_0,\dots,w_h$ are linearly dependent over $\mathbb C(k)$, so that $h\le\mathfrak m$, let $\sum_{i=0}^hq_i(k)\,w_i(k)=0$ with $q_i\in\mathbb C[k]$ without a common factor, and put $p_j(k)=q_{h-j}(k+h)$. Then $\sum_{j=0}^hp_j(k)\operatorname{Sh}^j[1]=0$, that is $\sum_jp_j(k)f^j\in\mathcal L$; $p_0\neq0$, and $p_h\neq0$ when $\mathbf B$ is invertible. The order $h$ is at least the dimension of the cyclic submodule of $[1]$, with equality when $\mathbf b$ are masters, and then the relation is the law of $[1]$.
\item Let $\mathbf W=[w_0;\dots;w_{h-1}]$ be the matrix of the first $h$ rows and $J$ a set of $h$ columns with $\det\mathbf W_J\neq0$. Then $(q_i)_i$ is proportional over $\mathbb C(k)$ to the signed maximal minors $\big((-1)^i\det[w_0;\dots;\widehat{w_i};\dots;w_h]_J\big)_i$. With $Q$ the monic least common denominator of the entries of $\mathbf B$, $q_h$ divides the polynomial $\det\mathbf W(k)_J\prod_{j=0}^{h-1}Q(k-j)^{h-j}$. So at a zero $k_0$ of $q_h$ with $Q(k_0-j)\neq0$ for $0\le j<h$ the rows $w_0(k_0),\dots,w_{h-1}(k_0)$ are linearly dependent: there $\langle1\rangle_{k_0},\dots,\langle1\rangle_{k_0-h+1}$, written by (i) as linear forms in the values of $\mathbf b$ at $k_0$, are dependent forms.
\item If $h=\mathfrak m$, then $q_0(k)\det\mathbf W(k)=(-1)^hq_h(k)\det\mathbf W(k-1)\det\mathbf B(k)$.
\item The same holds for every element $m=\sum_lc_l(k)[b_l]$ of $\mathfrak N$, with the row $c$ in place of $w_0$, $m$ in place of $[1]$, and the cyclic submodule of $m$, the span of the $\operatorname{Sh}^im$, in place of that of $[1]$.
\end{enumerate}
\end{proposition}

The first dependence among the rows $w_i$, which carry $\langle1\rangle$ to the lower strengths in the coordinates of a spanning set, is a recurrence of the values of $1$, and on masters it is the law of $[1]$. Its end coefficient $p_0(k)=q_h(k+h)$ vanishes only at shifted zeros of the denominator of $\mathbf B$ and at strengths where the rows line up, and when the rows reach the whole span, $h=\mathfrak m$, the end coefficients are tied to the determinant of $\mathbf B$. A spanning set that is not independent gives a relation whose order may exceed that of the law of $[1]$.

\section{Symmetries of a ray}\label{sec:raysym}

% New: proved in the last part.
\begin{proposition}[The symmetries of a ray]\label{prop:raysym}
Let a ray of $F$ have the origin $(\mvec_0;\mathbf s_0,\mathbf t_0)$ and the direction $(\mvec_1;\mathbf s_1,\mathbf t_1)$, and put $\sigma_i=\frac12\sum_j(s_{ij}+t_{ij})$, $i=0,1$. On the torus $(\mathbb C^\times)^{\cE}/\mathbb C^\times$ put, for a bijection $\beta$ of the marks, $\phi_\beta(\alpha)=\alpha\circ\beta$ with $(\alpha\circ\beta)_e=\alpha_{\beta(e)}$; the \emph{Cremona} $\operatorname{cr}(\alpha)=\alpha^{-1}$; and the \emph{reversed} map $\bar\phi_\beta=\operatorname{cr}\circ\phi_\beta$, $\bar\phi_\beta(\alpha)=(\alpha\circ\beta)^{-1}$. A \emph{symmetry} of the ray is a map $\phi$ among the $\phi_\beta$ and the $\bar\phi_\beta$ that keeps the stride and the density of the origin form, $f\circ\phi=f$ and $(\omega_0/dy)\circ\phi=\omega_0/dy$, the latter on the positive orthant with the principal powers when the origin is not integral; the symmetries form a group. Let $h$ be a bijection of the tables, write $u\circ\beta^{-1}$ for $e\mapsto u_{\beta^{-1}(e)}$ and $\mathbf a\circ h^{-1}$ for $j\mapsto a_{h^{-1}(j)}$, and say that $\beta$ \emph{carries}, respectively \emph{exchanges}, the tables with $h$ when it satisfies the conditions of Definition~\ref{def:carry} on the bases and the weights, without its condition on the pairs.
\begin{enumerate}[label=(\roman*)]
\item If $\beta$ carries the tables, $K_j\circ\phi_\beta=K_{h(j)}$ and $\Psi_j\circ\phi_\beta=\Psi_{h(j)}$. If $\beta$ exchanges them, $K_j\circ\phi_\beta=\Psi_{h(j)}$, $\Psi_j\circ\phi_\beta=K_{h(j)}$, $K_j\circ\bar\phi_\beta=\alpha^{-\cE}K_{h(j)}$ and $\Psi_j\circ\bar\phi_\beta=\alpha^{-\cE}\Psi_{h(j)}$. Always $K_j\circ\operatorname{cr}=\alpha^{-\cE}\Psi_j$ and $\Psi_j\circ\operatorname{cr}=\alpha^{-\cE}K_j$.
\item The following maps are symmetries when, for $i=0,1$, every mark $e$ and every table $j$:
\begin{center}
\begin{tabular}{lll}
\toprule
map & masses & exponents\\
\midrule
$\phi_\beta$, $\beta$ carrying & $m_{i,\beta(e)}=m_{i,e}$ & $(s_{i,h(j)},t_{i,h(j)})=(s_{ij},t_{ij})$\\
$\phi_\beta$, $\beta$ exchanging & $m_{i,\beta(e)}=m_{i,e}$ & $(s_{i,h(j)},t_{i,h(j)})=(t_{ij},s_{ij})$\\
$\operatorname{cr}$ & $\mvec_i=\sigma_i\one$ & $s_{ij}=t_{ij}$\\
$\bar\phi_\beta$, $\beta$ exchanging & $m_{i,\beta(e)}+m_{i,e}=2\sigma_i$ & $(s_{i,h(j)},t_{i,h(j)})=(s_{ij},t_{ij})$\\
\bottomrule
\end{tabular}
\end{center}
A symmetry maps $X_F$ and the positive orthant onto themselves, has $\phi^\ast\omega_0=\pm\omega_0$ when the origin is integral, and keeps the values, $\langle g'\circ\phi\rangle_k=\langle g'\rangle_k$ for every insertion $g'$.
\item On the insertions $\alpha^u\prod_jK_j^{-a_j}\Psi_j^{-b_j}$, written $(u;\mathbf a,\mathbf b)$, the four kinds act by $(u;\mathbf a,\mathbf b)\mapsto$
\begin{gather*}
(u\circ\beta^{-1};\mathbf a\circ h^{-1},\mathbf b\circ h^{-1})\ \ (\phi_\beta,\ \beta\text{ carrying}),\qquad(-u+|\mathbf a+\mathbf b|\one;\mathbf b,\mathbf a)\ \ (\operatorname{cr}),\\
(u\circ\beta^{-1};\mathbf b\circ h^{-1},\mathbf a\circ h^{-1})\ \ (\phi_\beta,\ \beta\text{ exchanging}),\\
(-u\circ\beta^{-1}+|\mathbf a+\mathbf b|\one;\mathbf a\circ h^{-1},\mathbf b\circ h^{-1})\ \ (\bar\phi_\beta,\ \beta\text{ exchanging}),
\end{gather*}
with $|\mathbf a+\mathbf b|=\sum_j(a_j+b_j)$.
\end{enumerate}
\end{proposition}

The four kinds are the monomial maps of the torus that a list can certify: a permutation of the marks that carries the bases, one that carries them onto the cobases, the inversion of every scale, and the composite of the last two. The first keeps each letter, the second and the third swap the two letters of every table, and the fourth keeps each letter up to the monomial $\alpha^{-\cE}$, which the masses absorb when the masses of a mark and of its image add up to $2\sigma_i$, twice the total strength. Composites of symmetries are symmetries, so a group generated by maps of the four kinds can be used whole.

% New: proved in the last part.
\begin{theorem}[The invariant module]\label{thm:invariant}
Let $G$ be a finite group of symmetries of a ray.
\begin{enumerate}[label=(\roman*)]
\item The maps $\phi^\ast[g']=[g'\circ\phi]$, $\phi\in G$, are well defined on $\mathfrak M$ and commute with the shift. With $e_G=|G|^{-1}\sum_{\phi\in G}\phi^\ast$, $\mathfrak M=\mathfrak M^G\oplus(1-e_G)\mathfrak M$ as submodules, and $\mathfrak M^G=e_G\mathfrak M$ consists of the classes of the invariant insertions.
\item A solution $V$ with $V\circ\phi^\ast=V$ for every $\phi\in G$, the values along the ray among them, vanishes on $(1-e_G)\mathfrak M$: $\langle g'\rangle_k=\langle e_Gg'\rangle_k$. The cyclic submodule of $[1]$ lies in $\mathfrak M^G$, so the law of $[1]$ is read in $\mathfrak M^G$, on masters among the invariant insertions.
\item If $\phi\in G$ fixes a nondegenerate native gauge $\alpha_c$, then $A_c\circ\phi^\ast=A_c$. For the $G$-orbit $O$ of $\alpha_c$, with stabilizer $G_c$, the sum $A_O=\sum_{\phi G_c\in G/G_c}A_c\circ\phi^\ast$ is invariant and does not vanish on $\mathfrak M^G$, and $\dim_{\mathbb C(k)}\mathfrak M^G$ is at least the number of $G$-orbits of nondegenerate native gauges.
\end{enumerate}
\end{theorem}

% New: proved in the last part.
\begin{theorem}[A reversed exchange at a lean]\label{thm:reversed}
Let $\beta$ exchange a family with itself (Definition~\ref{def:carry}), with the bijection $h$ of its tables, and let $\Upsilon=\Upsilon_{\boldsymbol\sigma,\boldsymbol\delta;\one}$ be a realization at the table's gauge with masses $\btau$ (Definition~\ref{def:realization}).
\begin{enumerate}[label=(\roman*)]
\item $R_{h(j),\beta(e)}(\one)=1-R_{j,e}(\one)$ for every table and mark, and $\tau_{\beta(e)}+\tau_e=2\sigma$. So $\bar\phi_\beta$ keeps the integrand $\alpha^{\btau}\prod_jK_j^{-s_j}\Psi_j^{-t_j}$ of the period at the data of $\Upsilon$, and it is a symmetry of the ray with origin $0$ and direction $(c\btau;c\mathbf s,c\mathbf t)$ whenever that direction is integral. The exchange $\phi_\beta$ and the Cremona each take that integrand to the integrand at the data of $\Upsilon_{\boldsymbol\sigma,-\boldsymbol\delta;\one}$.
\item For every mark $e$ and partners $\mathbf g=(g_1,\dots,g_q)$, with $\beta(\mathbf g)=(\beta(g_1),\dots,\beta(g_q))$, the views (Definition~\ref{def:view}) satisfy
\[
\Phi^{\mathbf E}_{\beta(e),\beta(\mathbf g)}=\Phi^{\mathbf E}_{e,\mathbf g},\qquad \Phi^{\mathbf O}_{\beta(e),\beta(\mathbf g)}=-\Phi^{\mathbf O}_{e,\mathbf g}:
\]
the reading of $\beta(e)$ at $\boldsymbol\delta$ is the reading of $e$ at $-\boldsymbol\delta$. At the centre $\beta$ maps every class onto itself; at a lean it permutes the classes of $\Upsilon$.
\item In a one-table molecule, $\beta$ maps every class $C$ of the centre onto itself, and when $C$ splits at the lean into $C_K$ and $C_\Psi$ (Theorem~\ref{thm:acrosscentre}), $\beta(C_K)=C_\Psi$ and $\beta(C_\Psi)=C_K$.
\end{enumerate}
\end{theorem}

An exchange carries the bases of a family onto the cobases, and the integrand at a lean weighs the two unequally: the exchange alone takes the lean to its opposite, and so does the Cremona, which inverts every scale. Their composite keeps the integrand at every lean at the table's gauge, because there the masses of a mark and of its image add up to $2\sigma$. So the integral at a lean keeps a symmetry carrying each mark to its image under the exchange, and insertions related by it have equal values; but the symmetry reverses every scale, $y\mapsto-\beta^\ast y$, and the reading, whose odd part changes sign under a reversal, sees $\beta(e)$ as $e$ read at the opposite lean. A bijection that keeps the potential fixes every class (Lemma~\ref{lem:certifypot}); the reversed exchange permutes the classes, and on a class of the centre that the lean splits it swaps the basis half and the cobasis half. At the centre the odd part is absent and the exchange, the Cremona and their composite all keep the integrand: the reversal is invisible to the reading there, and the classes it joins at the centre are those the lean separates.

\begin{remark}\label{rem:reversedgauge}
The condition $m_{\beta(e)}+m_e=2\sigma$ of Proposition~\ref{prop:raysym} is met off the table's gauge as well, for instance by the uniform masses $\sigma\one$ at every lean of a family with $\sum_j\delta_j(2r_j-n)=0$, as for one table with $2r=n$. The native gauge of such a realization on the positive orthant is fixed by $\bar\phi_\beta$ up to scale, and the views of $\beta(e)$ are those of $e$ at the reflected arguments $v\mapsto-v$ up to affine terms in $v$, by the identities $\Phi_{\mathbf O}(-\beta^\ast y)=\Phi_{\mathbf O}(y)$ and $\Phi_{\mathbf E}(-\beta^\ast y)=\Phi_{\mathbf E}(y)-2\sigma\ip\one y$ of Section~\ref{sec:chtwoends}. When that gauge is not proportional to $\one$ the realization is not an atom (Theorem~\ref{thm:gaugestep}), though the integral keeps the reversed exchange.
\end{remark}

\section{The box of the walls}\label{sec:box}

The module of a ray is the quotient by the relations of the period, and finitely many of them give a spanning set of its Laurent part and the step on it. An insertion is measured by its level, the powers of the letters in its denominator, and by its wall degree, the power of the reciprocal stride that dominates it, read on the walls of the direction. A box is the finite set of symbols of insertions within given levels and wall degree; its relations, the letter identities among them and the mark relations of its stems, are relations of the module, so they hold as identities of meromorphic functions of the strength. One elimination of a box leaves a spanning set of a submodule that contains $[1]$ and the matrix of the inverse shift on it, with a certificate supported on the stems, and a search through the boxes stops by the finite dimension of the ray. When the origin carries no letter exponents, the relations $(c^{(0)}_e+k\,c^{(1)}_e+\partial_{y_e})(\alpha^wf^{-1})$ involve Laurent insertions alone, and they do not suffice: the evaluation at a singular point of the zero set of $f^{-1}$ on the torus, when there is one, vanishes on all of them and not on the law of $[1]$. Insertions are the regular functions on the letter complement $X_F$ of Definition~\ref{def:ray}, on which every letter is inverted whatever its exponent along the ray; a box need carry no unit of a letter whose exponent along the ray vanishes identically, and its relations are relations of the module all the same.

% New: proved in the last part.
\begin{definition}[Levels, wall degree, boxes]\label{def:raybox}
Let a ray of a family $F$ of $N$ tables on $n\ge2$ marks have $\kappa_{1,S}>0$ for every nonempty proper $S$, and let $\Bbbk\subseteq\mathbb C$ be a field that contains the weights and the origin.
\begin{enumerate}[label=(\roman*)]
\item The \emph{symbols} are the data $(u;\mathbf a,\mathbf b)$ with $u\in\mathbb Z^{\cE}$, $\mathbf a,\mathbf b\in\mathbb N^N$ and $u_{\cE}=\sum_j(a_jr_j+b_jr^\ast_j)$. The \emph{evaluation} of $(u;\mathbf a,\mathbf b)$ is the insertion $\alpha^u\prod_jK_j^{-a_j}\Psi_j^{-b_j}$ (Proposition~\ref{prop:raysym}(iii)); its \emph{level} is $(\mathbf a,\mathbf b)$, and the symbols at the level $(\mathbf 0,\mathbf 0)$ are the \emph{Laurent} symbols $\alpha^u$, $u_{\cE}=0$. A \emph{formal combination} is a finite combination of symbols with coefficients in $\Bbbk(k)$, and its evaluation is the same combination of the insertions, an element of $\mathbb C[X_F]\otimes\mathbb C(k)$. The evaluations of the symbols span $\mathbb C[X_F]$.
\item The \emph{walls} of a symbol $g=(u;\mathbf a,\mathbf b)$ are those of its data, $\kappa_U(g)=u_U-\sum_j\big(a_jd_j(U)+b_jd^\ast_j(U)\big)$, and its \emph{wall degree} is
\[
\deg_\kappa g=\max_{\varnothing\neq U\subsetneq\cE}\frac{-\kappa_U(g)}{\kappa_{1,U}},
\]
the least real $x$ with $\kappa_U(g)+x\,\kappa_{1,U}\ge0$ for every nonempty proper $U$.
\item The \emph{mark relation} of a symbol $g=(u;\mathbf a,\mathbf b)$ at a mark $e$ is the formal combination
\begin{multline*}
\mathrm{MR}_e(g)=\big(m_{0,e}+k\,m_{1,e}+u_e\big)\,g-\sum_j\big(s_{0j}+k\,s_{1j}+a_j\big)\sum_{T\ni e}\theta_{j,T}\,(u+\one_T;\mathbf a+\mathbf e_j,\mathbf b)\\
-\sum_j\big(t_{0j}+k\,t_{1j}+b_j\big)\sum_{T\not\ni e}\theta_{j,T}\,(u+\one_{\cE\setminus T};\mathbf a,\mathbf b+\mathbf e_j),
\end{multline*}
the sums over the bases $T$ of table $j$, the terms whose coefficient vanishes identically in $k$ omitted. Its evaluation is $(c^{(0)}_e+k\,c^{(1)}_e+\partial_{y_e})$ applied to that of $g$. The \emph{letter identities} of $g$ are $g-\sum_{T\in\cB_j}\theta_{j,T}(u+\one_T;\mathbf a+\mathbf e_j,\mathbf b)$, for $K_j$, and $g-\sum_{T\in\cB_j}\theta_{j,T}(u+\one_{\cE\setminus T};\mathbf a,\mathbf b+\mathbf e_j)$, for $\Psi_j$; they evaluate to $0$.
\item A letter $K_j$, or $\Psi_j$, is \emph{active along the ray} when its exponent $s_{0j}+k\,s_{1j}$, or $t_{0j}+k\,t_{1j}$, does not vanish identically. \emph{Level bounds} are $\bar{\mathbf a},\bar{\mathbf b}\in\mathbb N^N$ with $\bar a_j\ge1$ when $K_j$ is active and $\bar b_j\ge1$ when $\Psi_j$ is active; put $\bar a^-_j=\max(\bar a_j-1,0)$ and $\bar b^-_j=\max(\bar b_j-1,0)$. For $\mathfrak t\in\mathbb N$ the \emph{box} $\mathfrak B=\mathfrak B_{\mathfrak t}(\bar{\mathbf a},\bar{\mathbf b})$ is the set of the symbols of wall degree at most $\mathfrak t$ at the levels $(\mathbf a,\mathbf b)$ with $\mathbf a\le\bar{\mathbf a}$ and $\mathbf b\le\bar{\mathbf b}$, and its \emph{stems} are its symbols at the level $(\bar{\mathbf a}^-,\bar{\mathbf b}^-)$. Its formal combinations form the space $\mathcal F_{\mathfrak B}$, and its \emph{relations} $\mathcal R_{\mathfrak B}\subseteq\mathcal F_{\mathfrak B}$ are the $\Bbbk(k)$-span of the letter identities of its symbols $(u;\mathbf a,\mathbf b)$, for $K_j$ when $a_j<\bar a_j$ and for $\Psi_j$ when $b_j<\bar b_j$, and of the mark relations of its stems at every mark. The box with \emph{one unit of each active letter} has $\bar a_j=1$ exactly when $K_j$ is active and $\bar b_j=1$ exactly when $\Psi_j$ is; its stems are the Laurent symbols of wall degree at most $\mathfrak t$.
\item The \emph{Laurent part} $\mathfrak M_{\mathrm L}$ of the module $\mathfrak M$ (Definition~\ref{def:raymodule}) is the span of the classes of the Laurent insertions.
\item A group of symmetries generated by maps of the four kinds of Proposition~\ref{prop:raysym}(ii) carries every letter to a letter times a monomial (Proposition~\ref{prop:raysym}(i)): each of its elements $\phi$ permutes the letters, and the symbols, $g\mapsto\phi^\ast g$ (Proposition~\ref{prop:raysym}(iii)), the evaluation of $\phi^\ast g$ being that of $g$ composed with $\phi$. Such a finite group $G$ \emph{keeps} the box when its permutations of the letters keep the level bounds, $\bar a_j$ being the bound of $K_j$ and $\bar b_j$ that of $\Psi_j$.
\end{enumerate}
\end{definition}

% New: proved in the last part.
\begin{proposition}[The wall degree]\label{prop:box-walldegree}
For a ray as in Definition~\ref{def:raybox} and its symbols $g$:
\begin{enumerate}[label=(\roman*)]
\item For real $x$, $|g|\,f^x$ is bounded on the positive orthant exactly when $x\ge\deg_\kappa g$: the symbols of wall degree at most $x$ are those whose evaluation is dominated by $f^{-x}$. For the origin $0$ the value $\langle g\rangle_k$ converges absolutely exactly when $\re k>\deg_\kappa g$, and for any origin it converges absolutely at every $k$ with $\re\kappa_{0,U}+(\re k-\deg_\kappa g)\,\kappa_{1,U}>0$ for every nonempty proper $U$.
\item A Laurent symbol $\alpha^u$ has the wall degree $\max_U(-u_U)/\kappa_{1,U}\ge0$, zero exactly for $u=0$. Every monomial of the Laurent polynomial $f^{-1}$ has wall degree at most $1$, and every monomial of $f^{-1}\alpha^u$ wall degree at most $\deg_\kappa\alpha^u+1$.
\item The terms of the mark relations and of the letter identities of $g$ have wall degree at most $\deg_\kappa g$.
\item At every level and for every real $x$ there are finitely many symbols of wall degree at most $x$.
\item For every element $\phi$ of a group as in Definition~\ref{def:raybox}(vi), $\deg_\kappa\phi^\ast g=\deg_\kappa g$, and $\phi^\ast$ maps Laurent symbols to Laurent symbols.
\end{enumerate}
\end{proposition}

% New: proved in the last part.
\begin{lemma}[The Laurent part]\label{lem:box-laurent}
For a ray as in Definition~\ref{def:raybox}, the Laurent part $\mathfrak M_{\mathrm L}$ is a submodule. It contains $[1]$, and the class of every insertion $g$ with $f^{-l}g$ a Laurent polynomial for some $l\in\mathbb N$; when every letter exponent of the direction is positive, $s_{1j},t_{1j}\ge1$, it is $\mathfrak M$. For a finite group $G$ of symmetries, $e_G\mathfrak M_{\mathrm L}=\mathfrak M_{\mathrm L}\cap\mathfrak M^G$ is a submodule, spanned by the classes of the invariant Laurent insertions.
\end{lemma}

% New: proved in the last part.
\begin{proposition}[The relations of a box]\label{prop:box-relations}
Let $\mathfrak B=\mathfrak B_{\mathfrak t}(\bar{\mathbf a},\bar{\mathbf b})$ be a box of a ray as in Definition~\ref{def:raybox}.
\begin{enumerate}[label=(\roman*)]
\item Every term of every relation of the box lies in the box. The letter identities evaluate to $0$ and the mark relation $\mathrm{MR}_e(g)$ of a stem to the element $(c^{(0)}_e+k\,c^{(1)}_e+\partial_{y_e})g$ of $\mathcal L$: every relation evaluates into $\mathcal L$, and its value vanishes identically in $k$ (Proposition~\ref{prop:rays-strength}(ii)).
\item A formal combination of symbols of the box whose evaluation is $0$ is a combination of the letter identities of the box.
\item The mark relation at every mark of every symbol of wall degree at most $\mathfrak t$ at a level at most $(\bar{\mathbf a}^-,\bar{\mathbf b}^-)$ lies in $\mathcal R_{\mathfrak B}$.
\item If a finite group $G$ of symmetries keeps the box, every $\phi\in G$ maps the box onto itself, its stems onto its stems and $\mathcal R_{\mathfrak B}$ onto itself; so, with $e_G=|G|^{-1}\sum_{\phi\in G}\phi^\ast$ and $\mathcal F^G_{\mathfrak B}=e_G\mathcal F_{\mathfrak B}$, $e_G\mathcal R_{\mathfrak B}=\mathcal R_{\mathfrak B}\cap\mathcal F^G_{\mathfrak B}$.
\end{enumerate}
\end{proposition}

% New: proved in the last part.
\begin{proposition}[The closing relations have more solutions than the ray]\label{prop:box-closing}
Let a ray as in Definition~\ref{def:raybox} have an origin with $\mathbf s_0=\mathbf t_0=\mathbf 0$, and write $f^{-1}=\sum_v\lambda_v\alpha^v$. For a Laurent insertion $\alpha^w$ and a mark $e$ the \emph{closing relation}
\[
\big(c^{(0)}_e+k\,c^{(1)}_e+\partial_{y_e}\big)\big(\alpha^wf^{-1}\big)=\sum_v\lambda_v\big(m_{0,e}+w_e+(1-k)\,v_e\big)\,\alpha^{w+v}
\]
is an element of $\mathcal L$ whose terms are Laurent insertions. Let $z\in(\mathbb C^\times)^{\cE}$ be a point at which $f^{-1}$ and every $\partial_{y_e}f^{-1}$ vanish, a singular point of the zero set of $f^{-1}$ on the torus.
\begin{enumerate}[label=(\roman*)]
\item The evaluation at $z$, $\alpha^w\mapsto z^w$ on the Laurent insertions, extended $\mathbb C(k)$-linearly, vanishes on every closing relation.
\item It takes the value $q_0(k)$ on every element $\sum_{i=0}^hq_i(k)f^{-i}$, $q_i\in\mathbb C(k)$. So no element of $\mathcal L$ of this form with $q_0\neq0$ is a combination of closing relations, while $\mathcal L$ contains such elements: the law $\sum_jp_j(k)\operatorname{Sh}^j[1]=0$ of $[1]$, with $p_0p_h\neq0$ (Definition~\ref{def:raymodule}), gives $\sum_{i=0}^hp_{h-i}(k-h)f^{-i}\in\mathcal L$. No $\mathbb C(k)$-linear form that vanishes on $\mathcal L$ agrees with the evaluation at $z$ on the Laurent insertions.
\end{enumerate}
\end{proposition}

The point $z$ lies off $X_F$, a letter that the direction uses vanishing there. The relations of $\mathcal L$ that the closing relations lack involve the insertions with inverse letters, through the letter identities and the mark relations of the occupations, and a box supplies them through its units of the letters.

% New: proved in the last part.
\begin{definition}[The reduction of a box]\label{def:box-reduction}
Let a finite group $G$ of symmetries keep a box $\mathfrak B=\mathfrak B_{\mathfrak t}(\bar{\mathbf a},\bar{\mathbf b})$ (Definition~\ref{def:raybox}); $G$ may be trivial.
\begin{enumerate}[label=(\roman*)]
\item The \emph{orbit average} of a symbol $g$ of the box is $e_Gg$, with $e_G$ as in Proposition~\ref{prop:box-relations}(iv); the orbit averages of the orbits of symbols form a basis of $\mathcal F^G_{\mathfrak B}$, and $\mathcal R^G_{\mathfrak B}=e_G\mathcal R_{\mathfrak B}$.
\item A \emph{ranking} is a total order of these orbits in which every orbit at a positive level ranks above every Laurent orbit and, of two Laurent orbits, the one of larger wall degree ranks higher; the level of the symbols of an orbit is $(\mathbf 0,\mathbf 0)$ for all of them or for none, and they all have the same wall degree (Proposition~\ref{prop:box-walldegree}(v)). The \emph{leading} orbit of a nonzero element of $\mathcal F^G_{\mathfrak B}$ is the highest-ranked orbit in it; an orbit is a \emph{pivot} when it leads an element of $\mathcal R^G_{\mathfrak B}$, and \emph{free} otherwise. The \emph{reduction} of $x\in\mathcal F^G_{\mathfrak B}$ is the unique combination of the averages of the free orbits that is congruent to $x$ modulo $\mathcal R^G_{\mathfrak B}$; if $x$ involves no orbit ranked above a given orbit, neither does its reduction.
\item For $\mathfrak s\in\mathbb N$ with $\mathfrak s<\mathfrak t$, the \emph{spanning set} of the box at $\mathfrak s$ is $\mathbf b=(b_1,\dots,b_{\mathfrak m})$, the averages of the free Laurent orbits of wall degree at most $\mathfrak s$, listed with $b_1=1$, the orbit of $1$ being free. Its \emph{targets} are the invariant Laurent polynomials $f^{-1}b_i$, elements of $\mathcal F^G_{\mathfrak B}$. The box is \emph{accepted} at $\mathfrak s$ when the reduction of every target is a combination of the spanning set, $f^{-1}b_i\equiv\sum_l\mathbf B_{il}\,b_l$ modulo $\mathcal R^G_{\mathfrak B}$ with $\mathbf B_{il}\in\Bbbk(k)$; $\mathbf B$ is then the \emph{matrix} of the box.
\end{enumerate}
\end{definition}

% New: proved in the last part.
\begin{theorem}[The box gives a spanning set, the one-step matrix and a certificate]\label{thm:raybox}
Let a finite group $G$ of symmetries keep a box $\mathfrak B=\mathfrak B_{\mathfrak t}(\bar{\mathbf a},\bar{\mathbf b})$ of a ray as in Definition~\ref{def:raybox}, and let the box be accepted at $\mathfrak s$ (Definition~\ref{def:box-reduction}), with the spanning set $\mathbf b$ and the matrix $\mathbf B$.
\begin{enumerate}[label=(\roman*)]
\item In the module, $[f^{-1}b_i]=\sum_l\mathbf B_{il}(k)[b_l]$ for every $i$. The classes of $\mathbf b$ span a submodule $\mathfrak N$ of $\mathfrak M^G$ that contains $[1]$: the span of the classes of the invariant combinations of Laurent symbols of wall degree at most $\mathfrak s$, which lies in $e_G\mathfrak M_{\mathrm L}$ (Lemma~\ref{lem:box-laurent}). The spanning set is a set of masters of $\mathfrak N$ exactly when its classes are linearly independent over $\mathbb C(k)$, and then $\mathbf B=\mathbf B_{\mathbf b}$ (Definition~\ref{def:raymodule}).
\item \emph{The certificate.} There are formal combinations $G_{i,e}$ of the stems, with coefficients in $\Bbbk(k)$, such that
\[
f^{-1}b_i-\sum_l\mathbf B_{il}(k)\,b_l=\sum_{e\in\cE}\big(c^{(0)}_e+k\,c^{(1)}_e+\partial_{y_e}\big)G_{i,e}
\]
as functions on $X_F$, the $G_{i,e}$ evaluated: a certificate of the one-step relation (Proposition~\ref{prop:rays-strength}(i)).
\item Every solution $V$ has $\varsigma^{-1}v_V=\mathbf B\,v_V$, with $v_V=(V[b_1],\dots,V[b_{\mathfrak m}])$, and $\langle b_i\rangle_{k-1}=\sum_l\mathbf B_{il}(k)\langle b_l\rangle_k$ as an identity of meromorphic functions of $k$. At every real $k$ off the poles of the coefficients with $\re\kappa_{0,U}+(k-\mathfrak t)\,\kappa_{1,U}>0$ for every nonempty proper $U$, all the values in (ii) converge absolutely and every term $(c^{(0)}_e+k\,c^{(1)}_e+\partial_{y_e})s$, $s$ a stem, has the value $0$, so the certificate integrates over the chart term by term.
\item Proposition~\ref{prop:rays-lawrows} applies to $\mathbf b$ and $\mathbf B$: the first dependence of the rows gives a relation $\sum_jp_j(k)f^j\in\mathcal L$, with a certificate assembled from those of (ii), whose order is at least that of the law of $[1]$, with equality when the spanning set is a set of masters, and which then is the law of $[1]$.
\end{enumerate}
\end{theorem}

% New: proved in the last part.
\begin{theorem}[The search stops]\label{thm:rayboxstop}
Let $G$ be a finite group of symmetries of a ray as in Definition~\ref{def:raybox}, generated by maps of the four kinds of Proposition~\ref{prop:raysym}(ii) and possibly trivial, and fix a ranking for every box that $G$ keeps. There are $\mathfrak t\ge1$, level bounds kept by $G$ and $\mathfrak s<\mathfrak t$ at which the box $\mathfrak B_{\mathfrak t}(\bar{\mathbf a},\bar{\mathbf b})$ modulo $G$ is accepted, with a spanning set whose classes span $e_G\mathfrak M_{\mathrm L}$, which is $\mathfrak M^G$ when every letter exponent of the direction is positive. So a search that runs through all the $(\mathfrak t,\bar{\mathbf a},\bar{\mathbf b},\mathfrak s)$ with level bounds kept by $G$, in any order that reaches each, and tests the acceptance by one elimination over $\Bbbk(k)$ at each, stops.
\end{theorem}

So the construction is finite and every step is an identity: a box is finite (Proposition~\ref{prop:box-walldegree}(iv)), its relations are relations of the period read at the data of the ray (Proposition~\ref{prop:box-relations}), so that whatever the box, the one-step relation and the recurrence it produces hold at every strength, and a search through the boxes stops (Theorem~\ref{thm:rayboxstop}). The spanning set spans a submodule that contains $[1]$, which is all a recurrence of the values of $1$ needs; it is a set of masters exactly when it is independent, and then the recurrence that its rows give is the law of $[1]$ and its matrix is $\mathbf B_{\mathbf b}$.

\section{Sources at the high end}\label{sec:sources}

% New: proved in the last part.
\begin{definition}[Sources; regular points; the normal density]\label{def:source}
The \emph{critical locus} of a ray is the set of its native gauges, $\{\alpha\in X_F:\ c^{(1)}(\alpha)=0\}$, the critical points of $\log f$. A \emph{source} is an irreducible component $Z$ of it; $f$ is constant on $Z$, with value $f_Z$, its \emph{native value}. A point $z\in Z$ is \emph{regular} when it lies on no other source, $Z$ is smooth at $z$ of dimension $d=\dim Z$, and the Hessian of $\log f$ in $y$ at $z$ has rank $n-1-d$ on the tangent space of the torus; its kernel is then the tangent space of $Z$, and it induces a nondegenerate quadratic form on the normal space. $Z$ is \emph{clean} when it has a regular point; the regular points then form a dense open subset $Z^\circ$. In local holomorphic coordinates $(u,w)$ at a regular point, with $Z=\{w=0\}$ and $-\log f=-\log f_Z+\frac12w^{\mathsf T}A(u)w+O(|w|^3)$, the \emph{normal density} is the $d$-form
\[
\nu=\det A(u)^{-1/2}\,\frac{\omega_0}{du\wedge dw}\Big|_{w=0}\,du ,
\]
which does not depend on the coordinates up to its sign. For an integral origin $\omega_0$ is single-valued on $X_F$, and the \emph{normal cover} $\tilde Z^\circ\to Z^\circ$ is the double cover on which $\nu$ is single-valued, its deck involution changing the sign of $\nu$. A source of dimension $0$ is an isolated native gauge, clean exactly when it is nondegenerate.
\end{definition}

% New: proved in the last part.
\begin{theorem}[A closed cycle on a source carries a solution]\label{thm:closedcycle}
Let a ray have an integral origin, let $Z$ be a clean source of dimension $d$, $m=n-1-d$, and let $\gamma$ be a compact oriented cycle of real dimension $d$ in $\tilde Z^\circ$.
\begin{enumerate}[label=(\roman*)]
\item Over $\gamma$ there is a chain $N(\gamma)$ of real dimension $n-1$ in $X_F$, a union of discs, one through each point $z$ of $\gamma$, transverse to $Z$, on which $|f|\le|f_Z|e^{-c|v|^2}$ at the displacement $v$ from $z$ for some $c>0$, oriented by the sheet of $\nu$; its boundary lies where $|f|\le|f_Z|e^{-\varepsilon}$ for some $\varepsilon>0$. Two such chains over homologous cycles differ by a boundary and by pieces on which $|f|\le|f_Z|e^{-\varepsilon}$.
\item As $k\to\infty$,
\[
\int_{N(\gamma)}g\,f^k\omega_0\sim f_Z^{\,k}\,k^{-m/2}\sum_{i\ge0}c_i(\gamma;g)\,k^{-i},\qquad c_0(\gamma;g)=(2\pi)^{m/2}\int_\gamma g\,\nu .
\]
Where one family of holomorphic normal slices $\{u=\mathrm{const}\}$, coordinates $(u,w)$ as in Definition~\ref{def:source}, covers $Z^\circ$, $c_i(\gamma;g)=\int_\gamma a_i(\cdot\,;g)$, with $f_Z^{\,k}k^{-m/2}\sum_ik^{-i}a_i(u;g)$ the Laplace expansion of the integral over the disc in the slice at $u$: $a_0(\cdot\,;g)=(2\pi)^{m/2}g\,\nu$, and each $a_i(\cdot\,;g)$ is a holomorphic $d$-form given by finitely many $w$-derivatives of $g$, of $\omega_0/(du\wedge dw)$ and of $\log f$ at $w=0$.
\item The map $A_\gamma:g\mapsto f_Z^{\,k}k^{-m/2}\sum_ic_i(\gamma;g)\,k^{-i}$ is a solution of type $\big(f_Z,-\frac m2+\mathbb Z\big)$. It depends only on the homology class of $\gamma$ in $\tilde Z^\circ$, and the image of $\gamma$ under the deck involution gives $-A_\gamma$.
\item If $\gamma_1,\dots,\gamma_r$ are such cycles and $g_1,\dots,g_r$ insertions with $\det\big[\int_{\gamma_a}g_b\,\nu\big]_{a,b}\neq0$, the $A_{\gamma_a}$ are linearly independent over $\mathbb C$, and their normalizations over $\mathbb C((1/k))$.
\end{enumerate}
For $d=0$, $\gamma$ is the point with a sheet and this is Theorem~\ref{thm:gaugesolution}. For an origin that is not integral, $\nu$ is multivalued along $Z^\circ$, and the same holds for cycles with coefficients in the local system of its branches.
\end{theorem}

% New: proved in the last part.
\begin{corollary}[A curve of native gauges: residues and periods]\label{cor:curvesources}
Let a ray have an integral origin, let $Z$ be a clean source of dimension one, and let $\bar Z$ be the smooth compact curve that contains $\tilde Z^\circ$ as the complement of finitely many \emph{punctures}: the points over the points of $Z$ that are not regular, among them its crossings with other sources and the points where the normal form degenerates, and the points where $Z$ leaves $X_F$. The solutions of $Z$ are spanned by those of a small loop around each puncture and those of a basis of the cycles of $\bar Z$. When $Z$ is the closure of $\{\alpha_0b^{\mathbf v}:b\in\mathbb C^\times\}$, $\mathbf v\in\mathbb Z^{\cE}$, with $b\mapsto\alpha_0b^{\mathbf v}$ injective, and the normal coordinates are those of one fixed complement $N$ of $\mathbb C\mathbf v$ in $\mathbb C^{\cE}/\mathbb C\one$, $\alpha=\alpha_0b^{\mathbf v}e^w$ with $w\in N$, the determinant $\det A(b)$ is rational in $b$, every $a_i(\cdot\,;g)$ is $r_i(b;g)\det A(b)^{-1/2}\,db/b$ with $r_i$ rational in $b$, $\tilde Z^\circ$ is an open part of the curve $Y^2=\det A(b)$, and the loop around a puncture $q$ gives $2\pi i$ times the residue at $q$ of $f_Z^{\,k}k^{-m/2}\sum_ik^{-i}a_i(\cdot\,;g)$.
\end{corollary}

% New: proved in the last part.
\begin{theorem}[A native gauge of corank at most one]\label{thm:amu}
Let $\alpha_c$ be an isolated native gauge of a ray at which the Hessian of $\log f$ in $y$ has corank at most one. There are holomorphic coordinates $(X,w)\in\mathbb C\times\mathbb C^{n-2}$ at $\alpha_c$ with
\[
\log f=\log f(\alpha_c)-\tfrac12\sum_iw_i^2+c\,X^{\mu+1},\qquad c\neq0,
\]
for an integer $\mu\ge1$, the \emph{multiplicity}; $\mu=1$ exactly when $\alpha_c$ is nondegenerate, and otherwise $\partial_X$ spans the kernel of the Hessian at $\alpha_c$. Then $\alpha_c$ is an \emph{$A_\mu$ point}, and:
\begin{enumerate}[label=(\roman*)]
\item There are $\mu$ solutions $A_{c,i}$, $0\le i<\mu$, $A_{c,i}$ of type $\big(f(\alpha_c),-\frac{n-2}2-\frac{i+1}{\mu+1}+\mathbb Z\big)$: the expansions of the integrals of $g\,f^k\omega_0$ over the chains (real normal disc)$\times\Gamma_i$, where $\Gamma_i$ is the combination of paths between the $\mu+1$ sectors on which $\re(cX^{\mu+1})<0$ with $\int_{\Gamma_i}X^je^{cX^{\mu+1}}dX=\delta_{ij}$ for $0\le j<\mu$. The leading coefficient of $A_{c,i}(g)$ is $(2\pi)^{(n-2)/2}$ times the coefficient of $X^i$ in the Taylor expansion of $g\,\omega_0/(dX\wedge dw)$ at $\alpha_c$. They are linearly independent, their exponents lying in distinct classes modulo $\mathbb Z$, and $A_{c,0}=(2\pi)^{-1/2}A_c$ for $\mu=1$.
\item If a symmetry $\phi$ of the ray fixes $\alpha_c$ and $\mu\ge2$, its differential multiplies the kernel by a root of unity $\zeta_\phi$ with $\zeta_\phi^{\mu+1}=1$, and $A_{c,i}\circ\phi^\ast=\zeta_\phi^{\,i}A_{c,i}$.
\end{enumerate}
\end{theorem}

% New: proved in the last part.
\begin{corollary}[Signs at a centre]\label{cor:signparity}
On the centre ray of a family (Definition~\ref{def:centreray}) let a sign $\eta$ lie in $X_F$ (Theorem~\ref{thm:signgauges}). The Cremona is a symmetry of the ray that fixes $\eta$ and acts on the tangent space there by $-1$.
\begin{enumerate}[label=(\roman*)]
\item If $\eta$ is nondegenerate, $A_\eta\circ\operatorname{cr}^\ast=A_\eta$.
\item If $\eta$ is isolated and the Hessian at $\eta$ has corank one, $\eta$ is an $A_\mu$ point with $\mu\ge3$ odd, and $A_{\eta,i}\circ\operatorname{cr}^\ast=(-1)^iA_{\eta,i}$: the $\frac{\mu+1}2$ solutions of even $i$ do not vanish on the Cremona-invariant insertions, and the $\frac{\mu-1}2$ of odd $i$ vanish on them. For $\mu=3$ the three exponents are $-\frac{n-2}2-\frac14$, $-\frac{n-1}2$ and $-\frac{n-2}2-\frac34$, the middle one odd.
\end{enumerate}
\end{corollary}

% New: proved in the last part.
\begin{proposition}[Symmetries on the cycles of a source]\label{prop:sourcesym}
Let a ray have an integral origin, and let a symmetry $\phi$, with $\phi^\ast\omega_0=\varepsilon_\phi\,\omega_0$, map a clean source $Z$ onto a source. It maps $Z^\circ$ onto $\phi(Z)^\circ$ and lifts to $\tilde\phi:\tilde Z^\circ\to\widetilde{\phi(Z)}^\circ$, carrying each oriented normal disc to its image, and
\[
A_\gamma\circ\phi^\ast=\varepsilon_\phi\,A_{\tilde\phi_\ast\gamma}.
\]
For a finite group $G$ of symmetries, on $\mathfrak M^G$ the solution of $\gamma$ is that of the averaged cycle $|G|^{-1}\sum_{\phi\in G}\varepsilon_\phi\,\tilde\phi_\ast\gamma$; a cycle class with $\sum_{\phi\in G_Z}\varepsilon_\phi\,\tilde\phi_\ast[\gamma]=0$, $G_Z$ the stabilizer of $Z$, gives a solution that vanishes on $\mathfrak M^G$.
\end{proposition}

% New: proved in the last part.
\begin{theorem}[What the invariant module sees at the high end]\label{thm:invariantcount}
Let $G$ be a finite group of symmetries of a ray, and let $\alpha_c$ be an $A_\mu$ point with stabilizer $G_c$. The solutions $A_{c,i}$ for which $\phi\mapsto\zeta_\phi^{\,i}$ is trivial on $G_c$ (all of them when $\mu=1$, with $\zeta_\phi^0=1$) average over $G$ to invariant solutions that do not vanish on $\mathfrak M^G$; the others vanish on $\mathfrak M^G$. So
\[
\dim_{\mathbb C(k)}\mathfrak M^G\ \ge\ \sum_O\#\big\{0\le i<\mu_O:\ \zeta_\phi^{\,i}=1\text{ for every }\phi\in G_{c_O}\big\},
\]
the sum over the $G$-orbits $O$ of the isolated native gauges of corank at most one, $\alpha_{c_O}$ a point of $O$ and $\mu_O$ its multiplicity.
\end{theorem}

Where the critical locus of the stride is not a set of nondegenerate points, its components still carry solutions. Over a clean family of native gauges the integral over a normal disc is a Gaussian, whose expansion on holomorphic normal slices is a holomorphic form of top degree on the family, so a closed one, with values in the orientation of the normal discs; every closed cycle carries a solution, because integration by parts is exact under the normal Gaussian and integrates to zero over a cycle along the family. On a curve the cycles are the loops around its punctures, among them its crossings with other sources and its degenerate points, and the cycles of its compactification: residues and periods, on the double cover where the determinant of the normal form has a square root. A native gauge of corank one carries one solution for each class of paths between the sectors where its reduced function decays, with exponents spaced by $\frac1{\mu+1}$. A symmetry that fixes a native gauge acts on its solutions: trivially at a nondegenerate one, and through the powers of its eigenvalue on the kernel at a degenerate one, so that at a sign of a centre, fixed by the Cremona, a point of multiplicity three carries an odd solution that every Cremona-invariant insertion misses. The invariant module counts each orbit of sources by the solutions its stabilizer leaves alone.

\section{Closings at the low end}\label{sec:torus}

% New: proved in the last part.
\begin{definition}[Closings]\label{def:closings}
Let a ray have an integral origin, $\mvec_0\in\mathbb Z^{\cE}$ and $s_{0j},t_{0j}\in\mathbb Z$, and a direction in which every letter has a positive exponent, $s_{1j},t_{1j}\ge1$. For an insertion $g$ let $j(g)$ be the least integer such that $g\,f^{-j}\omega_0/dy$ is a Laurent polynomial for every $j\ge j(g)$. The \emph{closings} of $g$ are the constant terms
\[
\mathsf C_j(g)=\CT\big(g\,f^{-j}\,\omega_0/dy\big)=\frac1{(2\pi i)^{n-1}}\oint_{|\alpha_e|=1,\ e\neq e_0}g\,f^{-j}\,\omega_0,\qquad j\ge j(g),
\]
taken on the chart $\alpha_{e_0}=1$, where they do not depend on $e_0$.
\end{definition}

% New: proved in the last part.
\begin{proposition}[One form on the orthant and on the torus]\label{prop:rays-pencil}
Let a ray have origin $0$ and $\kappa_{1,S}>0$ for every nonempty proper $S$, let $\alpha^\star$ be the native gauge of its direction on the positive orthant, $\alpha^\star_{e_0}=1$, and put $A=f(\alpha^\star)^{-1}$. On the chart $\alpha_{e_0}=1$, with $dy=\prod_{e\neq e_0}d\alpha_e/\alpha_e$, the rational form $\varpi_q=dy/(f^{-1}-q)$ has
\begin{gather*}
\int_{y_{e_0}=0}\varpi_q=\sum_{k\ge1}\langle1\rangle_k\,q^{k-1}\qquad(|q|<A),\\
\frac1{(2\pi i)^{n-1}}\oint_{|\alpha_e|=\alpha^\star_e,\ e\neq e_0}\varpi_q=-\sum_{j\ge0}\CT\big(f^{-j}\big)\,q^{-j-1}\qquad(|q|>A),
\end{gather*}
where $\CT(f^{-j})$ is the constant term of the Laurent polynomial $f^{-j}$, the closing $\mathsf C_j(1)$ when every letter exponent of the direction is positive (Definition~\ref{def:closings}). So the generating function of the values near $q=0$ and that of the constant terms near $q=\infty$ are the integrals of one form, whose denominator is the pencil $\alpha^{-\mvec_1}\big(\prod_jK_j^{s_{1j}}\Psi_j^{t_{1j}}-q\,\alpha^{\mvec_1}\big)$ of Theorem~\ref{thm:pencil}(ii), over the orthant and over the torus through $\alpha^\star$.
\end{proposition}

% New: proved in the last part.
\begin{proposition}[The constant terms of a centre ray and its signs]\label{prop:rays-signclosings}
On the centre ray of a family at the strengths $\boldsymbol\sigma$ (Definition~\ref{def:centreray}), with the tables indexed by $l$, for every $j\in\mathbb N$
\[
\frac1{2^{n-1}}\sum_\eta\prod_lK_l(\eta)^{2j\sigma_l}=\sum_{v\in2\mathbb Z^{\cE}}[\alpha^v]\,f^{-j}\ \ge\ \CT\big(f^{-j}\big),
\]
the sum over the $2^{n-1}$ signs $\eta$ of Theorem~\ref{thm:signgauges}: the mean of the $j$-th powers of the values $f^{-1}(\eta)=\prod_lK_l(\eta)^{2\sigma_l}$ at the signs is the sum of the coefficients of $f^{-j}$ at the monomials with even exponents.
\end{proposition}

% New: proved in the last part.
\begin{theorem}[The closings are a solution]\label{thm:closings}
For a ray as in Definition~\ref{def:closings}, $\mathsf C_j\big((c^{(0)}_e-j\,c^{(1)}_e+\partial_{y_e})g\big)=0$ and $\mathsf C_j(fg)=\mathsf C_{j-1}(g)$ for $j>j(g)$, and
\[
V_{\mathsf C}\Big(\sum_ig_i\otimes r_i\Big)(k)=\sum_ir_i(k)\,\mathsf C_{-k}(g_i),
\]
at the integers $k\le-\max_ij(g_i)$ off the poles of the $r_i$, is a well defined solution with values in germs at $-\infty$. So for masters $\mathbf b$, $\mathsf C_{j+1}(\mathbf b)=\mathbf B_{\mathbf b}(-j)\,\mathsf C_j(\mathbf b)$ for all large $j$, with $\mathsf C_j(\mathbf b)=(\mathsf C_j(b_1),\dots,\mathsf C_j(b_{\mathfrak m}))$, and the law of $[1]$ gives $\sum_ip_i(-j)\,\mathsf C_{j-i}(1)=0$ for all large $j$.
\end{theorem}

% New: proved in the last part.
\begin{proposition}[The closings are the top principal parts]\label{prop:principal}
Let a ray as in Definition~\ref{def:closings} have origin $0$, so that $\omega_0=dy$ and $\langle g\rangle_k=\int_{y_{e_0}=0}g\,f^k\,dy$, and $\kappa_{1,S}>0$ for every nonempty proper $S$, and put
\[
V_1=\sum_\pi\prod_{l=2}^n\kappa_{1,U_l}^{-1}>0,
\]
the sum over the orders $\pi$ of the marks with tails $U_l$ (Theorem~\ref{thm:severity}(i)). For every insertion $g$ the value $\langle g\rangle_k$ continues to a meromorphic function of $k$ whose poles have order at most $n-1$, and for every integer $j\ge j(g)$
\[
\lim_{k\to-j}(k+j)^{n-1}\langle g\rangle_k=V_1\,\mathsf C_j(g).
\]
\end{proposition}

% New: proved in the last part.
\begin{proposition}[The low end, for every direction]\label{prop:rays-lowpoles}
Let a ray have $\kappa_{1,S}>0$ for every nonempty proper $S$; its letter exponents $s_{1j},t_{1j}\in\mathbb N$ may vanish.
\begin{enumerate}[label=(\roman*)]
\item For an integral origin, every pole of the value of an insertion $g=\alpha^u\prod_jK_j^{-a_j}\Psi_j^{-b_j}$ is a $k$ with $\kappa_{0,U}+\kappa_U(g)+k\,\kappa_{1,U}\in\{0,-1,-2,\dots\}$ for some nonempty proper $U$ (Proposition~\ref{prop:rays-strength}(ii)), and so lies on one of the lines $\kappa_{1,U}\,k\in\mathbb Z$ set by the walls of the direction; its order is at most $n-1$.
\item For the origin $0$, for every insertion $g$ and every $j\in\mathbb N$ with $f^{-j}g$ a Laurent polynomial, with $\CT$ its constant term as in Definition~\ref{def:closings},
\[
\lim_{k\to-j}(k+j)^{n-1}\langle g\rangle_k=V_1\,\CT\big(f^{-j}g\big),\qquad V_1=\sum_\pi\prod_{l=2}^n\kappa_{1,U_l}^{-1}
\]
as in Proposition~\ref{prop:principal}; in particular $\lim_{k\to-j}(k+j)^{n-1}\langle1\rangle_k=V_1\CT(f^{-j})$ for every $j\in\mathbb N$.
\end{enumerate}
\end{proposition}

% New: proved in the last part.
\begin{corollary}[The law read at the low end]\label{cor:reflected}
For a ray as in Proposition~\ref{prop:principal}, with the law $\sum_{i=0}^hp_i(k)\operatorname{Sh}^i[1]=0$ of $[1]$, the closings of $1$, $\mathsf C_j(1)=\CT(f^{-j})$, satisfy
\[
\sum_{i=0}^{\min(h,j)}p_i(-j)\,\mathsf C_{j-i}(1)=0\qquad\text{for every }j\ge0 .
\]
On the centre ray of a family with every $\sigma_j\ge1$, $f^{-1}=\prod_j\big(K_j(\alpha)K_j(\alpha^{-1})\big)^{\sigma_j}$, and $\mathsf C_j(1)$ is the constant term of its $j$-th power.
\end{corollary}

% New: proved in the last part.
\begin{definition}[The low-end coefficients]\label{def:rays-lowend}
For a ray with origin $0$ and $\kappa_{1,S}>0$ for every nonempty proper $S$, the \emph{low-end coefficients} $\mathsf L_i$, $i\in\mathbb N$, are the Laurent coefficients of the values of $1$ at $k=0$,
\[
\langle1\rangle_k=\sum_{i\ge0}\mathsf L_i\,k^{i-n+1}\qquad\text{near }k=0;
\]
$\mathsf L_0=V_1>0$.
\end{definition}

% New: proved in the last part.
\begin{corollary}[The low-end relation]\label{cor:rays-lowend}\label{cor:lowendrelation}
Let a ray have origin $0$ and $\kappa_{1,S}>0$ for every nonempty proper $S$, whatever its letter exponents, and let $p_0,\dots,p_h\in\mathbb C[k]$ have $\sum_ip_i(k)f^i\in\mathcal L$: the law of $[1]$, or a relation given by Proposition~\ref{prop:rays-lawrows}, by a box among them (Theorem~\ref{thm:raybox}(iv)).
\begin{enumerate}[label=(\roman*)]
\item $\sum_ip_i(k)\langle1\rangle_{k+i}=0$ as meromorphic functions of $k$, and at every $k$ with $\re k>0$; and $\sum_{i=0}^{\min(h,j)}p_i(-j)\,\CT\big(f^{i-j}\big)=0$ for every $j\in\mathbb N$.
\item $p_0$ vanishes to order at least $n-1$ at $k=0$, and the \emph{low-end relation}
\[
[k^{n-1}]p_0\cdot\mathsf L_0+\sum_{i=1}^hp_i(0)\,\langle1\rangle_i=0
\]
holds.
\item If $\sum_ip_i(k)f^i=\sum_e\big(k\,c^{(1)}_e+\partial_{y_e}\big)G_e$ is a certificate whose $G_e$ have no pole at $k=0$, then $\sum_{i\ge1}p_i(0)f^i=\sum_e\partial_{y_e}G_e(0)$, and on the chart the flux of the field $\big(G_e(0)-G_{e_0}(0)\big)_{e\neq e_0}$ out of the cube $\{|y_e|\le R,\ e\neq e_0\}$ tends to $-[k^{n-1}]p_0\cdot\mathsf L_0$ as $R\to\infty$: the low-end relation is the flux at infinity of the bottom piece of the certificate (Proposition~\ref{prop:rays-certtop}).
\end{enumerate}
\end{corollary}

% New: proved in the last part.
\begin{proposition}[The low-end coefficients are integrals over the faces of the sectors]\label{prop:rays-faces}
Let a ray have origin $0$ and $\kappa_{1,S}>0$ for every nonempty proper $S$. For an order $\pi$ of $\cE$ with tails $U_l$ write $\kappa_l=\kappa_{1,U_l}$ and, in the sector coordinates $q\in[0,1]^{n-1}$ (Theorem~\ref{thm:exist}(ii)), $f=\prod_{l=2}^nq_l^{\kappa_l}\,\varphi_\pi(q)$ with $\varphi_\pi=\prod_jk_{j,\pi}^{-s_{1j}}\psi_{j,\pi}^{-t_{1j}}$, positive and smooth on a neighbourhood of the cube. For $A\subseteq\{2,\dots,n\}$ and a function $h$ on the cube put
\[
h_{[A]}=\sum_{B\subseteq A}(-1)^{|A\setminus B|}\,h\big|_{q_l=0,\ l\notin B},
\]
a function of $q_A$ that vanishes where some $q_l$, $l\in A$, vanishes. Then
\[
\mathsf L_i=\sum_\pi\sum_{A}\prod_{l\notin A}\kappa_l^{-1}\sum_{m+|\boldsymbol\gamma|=i-|A|}\frac1{m!}\int_{[0,1]^A}\prod_{l\in A}\frac{(\kappa_l\log q_l)^{\gamma_l}}{\gamma_l!\,q_l}\,\big(\log^m\varphi_\pi\big)_{[A]}\,dq_A ,
\]
the second sum over $A\subseteq\{2,\dots,n\}$ and the third over $m\in\mathbb N$ and $\boldsymbol\gamma\in\mathbb N^A$, every integral converging. So $\mathsf L_0=V_1$ and
\[
\mathsf L_1=\sum_\pi\prod_l\kappa_l^{-1}\log\varphi_\pi(0)=-\sum_\pi\prod_l\kappa_l^{-1}\sum_j\big(s_{1j}\log\theta_{j,T^+_{j,\pi}}+t_{1j}\log\theta_{j,T^-_{j,\pi}}\big),
\]
with the greedy bases of Lemma~\ref{lem:greedy}; up to $i=n-1$ only the faces $A\neq\{2,\dots,n\}$ of the cube enter, $1_{[A]}$ being $0$ for $A\neq\varnothing$, and the whole cube enters from $i=n$ on.
\end{proposition}

At the low end of strength the values of a ray from the origin have poles at the nonpositive integers, and the leading coefficient of each pole is one constant, the volume read by the walls of the direction, times the constant term of a Laurent polynomial: the closings. Integration by parts lowers the pole of every monomial but the constant one, so only the constant term survives at the top order. The closings are the values over the compact torus, a cycle on which every negative power of the stride is single-valued, and they obey every relation of the ray; read at $k=-j$, the law of $[1]$ is a recurrence for the constant terms of the powers of $f^{-1}$, from $j=0$ on.

\begin{remark}\label{rem:threads}
The constant term of a Laurent polynomial is its mean over every torus $|\alpha_e|=\rho_e$. When the walls of the direction are positive, on the torus through the native gauge $\alpha^\star$ of the direction $|f^{-1}|$ is largest at $\alpha^\star$ alone, and as $j\to\infty$ the closings localize there, with the expansion of $A_\star$ read at $k=-j$ (Theorem~\ref{thm:onekernel}). Where another torus can be deformed into descent cycles of $|f^{-1}|$, the closings decompose over the sources met on the way, with integer multiplicities, a family of native gauges among them. The decomposition does not enlarge what they see: they annihilate exactly what $A_\star$ annihilates (Section~\ref{sec:closingssee}), and what they supply is a cheap stream of equations with the strength of that one solution.
\end{remark}

\section{Separating the sources and gluing them}\label{sec:glue}

% New: proved in the last part.
\begin{proposition}[The high end reads the formal type]\label{prop:formaltype}
Let $\mathbf B=\mathbf B_{\mathbf b}$ be regular at $k=\infty$, $\mathbf B(k)=\mathbf B_0+\mathbf B_1k^{-1}+O(k^{-2})$, and let $V$ be a solution of type $(\lambda,\rho+\mathbb Z)$ whose solution vector is $v_V=\lambda^kk^{\rho}(a_0+a_1k^{-1}+\cdots)$ with $a_0\neq0$. Then $\mathbf B_0a_0=\lambda^{-1}a_0$, and for every row vector $w$ with $w\mathbf B_0=\lambda^{-1}w$ and $wa_0\neq0$, $\rho=-\lambda\,w\mathbf B_1a_0/(wa_0)$. So the native value of every source with a solution that does not vanish on $\mathfrak N$ is the reciprocal of an eigenvalue of $\mathbf B(\infty)$; its exponent, congruent modulo $\mathbb Z$ to $-\frac{n-1-d}2$ for a clean source of dimension $d$ and to $-\frac{n-2}2-\frac{i+1}{\mu+1}$ for the solution $A_{c,i}$ of an $A_\mu$ point, is read from $\mathbf B_1$ when a left eigenvector of $\mathbf B_0$ for $\lambda^{-1}$ does not annihilate $a_0$, as when $\lambda^{-1}$ is a simple eigenvalue.
\end{proposition}

% New: proved in the last part.
\begin{proposition}[The spectrum at infinite strength; changes of masters]\label{prop:rays-spectrum}
Let insertions $b_1=1,b_2,\dots,b_{\mathfrak m}$ have classes that span a submodule $\mathfrak N$, and let $\mathbf B\in\operatorname{Mat}_{\mathfrak m}(\mathbb C(k))$ have $[f^{-1}b_i]=\sum_l\mathbf B_{il}[b_l]$: masters with $\mathbf B_{\mathbf b}$, or the spanning set and the matrix of an accepted box (Theorem~\ref{thm:raybox}). For a nondegenerate native gauge $\alpha_c$ put $\beta_c=\big(b_l(\alpha_c)\big)_l$, nonzero as $b_1=1$, and let $\hat A_c$ be the normalization of the expansion $A_c$ (Definition~\ref{def:raysolution}).
\begin{enumerate}[label=(\roman*)]
\item If the vectors $\beta_c$ of a set $C$ of nondegenerate native gauges span $\mathbb C^{\mathfrak m}$, the classes of the $b_l$ are linearly independent: they are masters of $\mathfrak N$, $\dim_{\mathbb C(k)}\mathfrak N=\mathfrak m$, and $\mathbf B=\mathbf B_{\mathbf b}$.
\item If $\mathbf B$ is regular at $k=\infty$, then $\mathbf B(\infty)\,\beta_c=f(\alpha_c)^{-1}\beta_c$ for every nondegenerate native gauge $\alpha_c$: the reciprocal of every nondegenerate native value is an eigenvalue of $\mathbf B(\infty)$, with the values of the $b_l$ at the gauge as an eigenvector. So gauges of distinct native values have independent vectors $\beta_c$, there are at most $\mathfrak m$ distinct nondegenerate native values, and under (i) $\mathbf B(\infty)$ is diagonalizable with the eigenvalues $f(\alpha_c)^{-1}$, $\alpha_c\in C$. In particular, when two nondegenerate native gauges of distinct native values have proportional vectors $\beta_c$, $\mathbf B$ is not regular at $k=\infty$.
\item Let $T\in GL_{\mathfrak m}(\mathbb C(k))$ make $\tilde{\mathbf B}(k)=T(k-1)\,\mathbf B(k)\,T(k)^{-1}$, the matrix of $\operatorname{Sh}^{-1}$ on the elements $\tilde b_i=\sum_lT_{il}(k)[b_l]$, regular at $k=\infty$. Then $f(\alpha_c)^{-1}$ is an eigenvalue of $\tilde{\mathbf B}(\infty)$ for every nondegenerate native gauge $\alpha_c$, with the leading coefficient of the series vector $T(k)\big(\hat A_c(b_l)\big)_l$ as an eigenvector; and if $T(k-1)T(k)^{-1}\to I$ as $k\to\infty$, the eigenvalues of $\mathbf B(k)$ tend to those of $\tilde{\mathbf B}(\infty)$, with their multiplicities.
\item The condition $T(k-1)T(k)^{-1}\to I$ cannot be dropped. For a constant $\mathbf B_\circ=\begin{bmatrix}x_{11}&x_{12}\\x_{21}&x_{22}\end{bmatrix}\in GL_2(\mathbb C)$ with $x_{12}\neq0$ and $S(k)=\begin{bmatrix}1&0\\k&1\end{bmatrix}$, the matrix $\mathbf B(k)=S(k-1)\,\mathbf B_\circ\,S(k)^{-1}$ is not regular at $k=\infty$, the change $T=S^{-1}$ makes it $\mathbf B_\circ$, with $T(k-1)T(k)^{-1}=\begin{bmatrix}1&0\\1&1\end{bmatrix}$, and at every $k$ the eigenvalues of $\mathbf B(k)$ are the roots of $\lambda^2-(x_{11}+x_{22}-x_{12})\lambda+x_{11}x_{22}-x_{12}x_{21}$, which are not those of $\mathbf B_\circ$.
\end{enumerate}
\end{proposition}

% New: proved in the last part.
\begin{theorem}[Separating the sources]\label{thm:separate}
Let $\mathbf B=\mathbf B_{\mathbf b}$ be regular at $k=\infty$, and let the spectrum of $\mathbf B(\infty)$ be the disjoint union of nonempty sets $\Lambda_1,\dots,\Lambda_s$. Over $\mathbb C((1/k))$, with the shift $\varsigma$,
\[
\mathfrak N\otimes\mathbb C((1/k))=\bigoplus_{a=1}^s\hat{\mathfrak N}_a,
\]
uniquely, with every $\hat{\mathfrak N}_a$ stable under the shift and the matrix of $\operatorname{Sh}^{-1}$ on it regular at $\infty$ in a suitable basis, with value at $\infty$ of spectrum $\Lambda_a$; equivalently there is $T\in GL_{\mathfrak m}(\mathbb C[[1/k]])$ with $T(k-1)\,\mathbf B(k)\,T(k)^{-1}$ block diagonal. A solution of type $(\lambda,\rho+\mathbb Z)$ with $\lambda^{-1}\in\Lambda_a$ vanishes on every $\hat{\mathfrak N}_b$ with $b\neq a$, and every group of symmetries that maps $\mathfrak N$ onto itself maps each $\hat{\mathfrak N}_a$ onto itself.
\end{theorem}

% New: proved in the last part.
\begin{corollary}[Gluing at the two ends]\label{cor:glue}
Let $\mathcal S$ consist of formal solutions and the closings $V_{\mathsf C}$, and let $\mathfrak N_{\mathsf C}=\operatorname{Ann}(V_{\mathsf C})\cap\mathfrak N$ be the submodule the closings miss. Then $\mathcal S$ determines $\mathbf B_{\mathbf b}$ exactly when no nonzero submodule of $\mathfrak N_{\mathsf C}$ is annihilated by all the formal solutions in $\mathcal S$. The closings alone determine it exactly when $\mathfrak N_{\mathsf C}=0$, as they do when $\mathfrak N$ is irreducible and $V_{\mathsf C}$ does not vanish on it; formal solutions alone determine it when $\mathfrak m$ of them are linearly independent over $\mathbb C$ on $\mathfrak N$, type by type.
\end{corollary}

At the high end the module comes apart: over the formal series in $1/k$ it is the direct sum of one piece for each cluster of native values, the pieces are canonical, and every symmetry of the ray maps each piece to itself. This is where distinguishable sources are separated before the symmetric procedure runs: a piece is read with its own solutions, each class at a degenerate point and each cycle on a family counted by its character under the stabilizer, and the formal type of the one-step matrix at infinity, its eigenvalues and its exponents, contains the native values and the exponents of the sources that the module sees. The pieces need not be rational. The rational one-step matrix glues them, and the closings, an exact solution at the other end of strength, supply the glue: they determine the matrix on every submodule that meets their annihilator only in $0$, every irreducible submodule on which they do not vanish among them, and formal solutions are needed only on the submodule that the closings miss.

\section{What the closings see}\label{sec:closingssee}

% New: proved in the last part.
\begin{theorem}[The closings, the values and the native gauge]\label{thm:onekernel}
Let a ray have an integral origin and an integral direction in which every letter has a positive exponent, $s_{1j},t_{1j}\ge1$, with $\kappa_{1,S}>0$ for every nonempty proper $S$, and let $\alpha^\star$ be the native gauge of the direction on the positive orthant, where $f$ is largest (Theorem~\ref{thm:pencil}(i)). Then $\alpha^\star$ is a nondegenerate native gauge of the ray; write $A_\star$ for the expansion over the orthant near it. For every $m\in\mathfrak M$ the following are equivalent:
\begin{enumerate}[label=(\roman*)]
\item the closings vanish on $m$: $V_{\mathsf C}(m)=0$;
\item the values vanish on $m$: $\langle m\rangle_k=0$ for all large integers $k$;
\item the expansion at $\alpha^\star$ vanishes on $m$: $A_\star(m)=0$;
\end{enumerate}
and when they hold, the continuation of $\langle m\rangle_k$ in $k$ vanishes identically. So $\operatorname{Ann}(V_{\mathsf C})=\operatorname{Ann}(\langle\cdot\rangle)=\operatorname{Ann}(A_\star)$, and the closings, the values and $A_\star$ each determine $\mathbf B_{\mathbf b}$ on exactly the same submodules (Theorem~\ref{thm:determine}(iii)). Moreover, for every insertion $g$,
\[
j^{(n-1)/2}f(\alpha^\star)^j\,\mathsf C_j(g)\ \sim\ (2\pi)^{-(n-1)}\,\hat A_\star(g)\big|_{k=-j}\qquad(j\to\infty).
\]
\end{theorem}

The positive orthant and the torus through $\alpha^\star$ meet only at $\alpha^\star$. The values are read on the orthant, where $f$ is largest at $\alpha^\star$. The closings are means of Laurent polynomials over any torus, and on the torus through $\alpha^\star$ the reciprocal stride has its largest modulus at $\alpha^\star$ alone, because the letters have positive coefficients and the walls of the direction are positive. So the two ends of strength expand at one point, and their expansions are one formal solution, $A_\star$, read as $k\to\infty$ and at $k=-j$. An element on which $A_\star$ vanishes has a density of values that vanishes near the top of its range, and analyticity carries the vanishing to the whole range; and the values reach the closings through their principal parts (Proposition~\ref{prop:principal}). The closings see exactly what the values see, and both see exactly what the native gauge of the direction sees.

\begin{definition}[The value module; what a solution sees]\label{def:valuemodule}
For a ray with $\kappa_{1,S}>0$ for every nonempty proper $S$, the \emph{kernel of the values} is $\operatorname{Ann}(\langle\cdot\rangle)$, the submodule of the elements whose values vanish for all large $k$ (Theorem~\ref{thm:determine}(ii)), and the \emph{value module} is $\mathfrak M/\operatorname{Ann}(\langle\cdot\rangle)$. A solution \emph{sees} a submodule when it does not vanish on it, and a nondegenerate native gauge $\alpha_c$ sees it when $A_c$ does.
\end{definition}

% New: proved in the last part.
\begin{corollary}[When the closings miss nothing]\label{cor:seeall}
Under the hypotheses of Theorem~\ref{thm:onekernel}, let $\mathfrak N$ be a submodule. The closings annihilate no nonzero submodule of $\mathfrak N$ exactly when the values vanish on no nonzero element of $\mathfrak N$, and exactly when $\alpha^\star$ sees every nonzero submodule of $\mathfrak N$. Every symmetry of the ray fixes $\alpha^\star$, so for every finite group $G$ of symmetries the closings, the values and $A_\star$ vanish on $(1-e_G)\mathfrak M$: the condition is to be read on $\mathfrak M^G$. Let $\mathbf B=\mathbf B_{\mathbf b}$, for masters of $\mathfrak N$, be regular at $k=\infty$ with $1/f(\alpha^\star)$ a simple eigenvalue of $\mathbf B(\infty)$, let $\hat{\mathfrak N}_\star$ be the line of Theorem~\ref{thm:separate} for that eigenvalue, and let $\alpha^\star$ see $\mathfrak N$. Then $\alpha^\star$ sees a submodule $\mathfrak N'$ exactly when $\hat{\mathfrak N}_\star\subseteq\mathfrak N'\otimes\mathbb C((1/k))$; when it sees every nonzero submodule, $\mathfrak N$ has exactly one minimal nonzero submodule; and when $\mathfrak N=\mathfrak N_1\oplus\mathfrak N_2$ with $\mathfrak N_1,\mathfrak N_2$ nonzero submodules, the closings annihilate one of the two.
\end{corollary}

% New: proved in the last part.
\begin{corollary}[What the closings determine]\label{cor:valuemodule}
Under the hypotheses of Theorem~\ref{thm:onekernel}, the closings determine the one-step matrix of the value module: for insertions whose classes modulo $\operatorname{Ann}(\langle\cdot\rangle)$ form a basis of the image of a submodule, Theorem~\ref{thm:determine}(iii) and~(iv) hold on the value module with the closings alone. For polynomials $p_0,\dots,p_h$, $\sum_ip_i(k)\langle1\rangle_{k+i}=0$ for all large $k$ exactly when $\sum_ip_i(-j)\,\mathsf C_{j-i}(1)=0$ for all large $j$: the values of $1$ and the closings of $1$ satisfy the same recurrences, read at $k$ and at $-j$, and have the same tier (Proposition~\ref{prop:frame}).
\end{corollary}

% New: proved in the last part.
\begin{proposition}[The closings measure what they miss]\label{prop:blinddim}
Let $\mathbf b$ be masters of a submodule $\mathfrak N$, $\mathfrak m$ of them, under the hypotheses of Theorem~\ref{thm:closings}; let $d=\dim_{\mathbb C(k)}\big(\operatorname{Ann}(V_{\mathsf C})\cap\mathfrak N\big)$, the dimension of the largest submodule of $\mathfrak N$ that the closings annihilate, and for $D\ge0$ let $\mathcal S_D$ be the space of pairs $(P,Q)$ of a matrix and a scalar of polynomials of degree at most $D$ with $Q(-j)\,\mathsf C_{j+1}(\mathbf b)=P(-j)\,\mathsf C_j(\mathbf b)$ for all large $j$. Then $\dim\mathcal S_D=(\mathfrak md+1)D+O(1)$ as $D\to\infty$. When $d=0$, $\dim\mathcal S_D=\max(0,D-D_0+1)$ with $D_0$ as in Theorem~\ref{thm:determine}(iv); so a single step $\dim\mathcal S_{D+1}-\dim\mathcal S_D>1$ proves $d>0$.
\end{proposition}

The closings therefore certify a one-step matrix by themselves exactly when the native gauge of the direction sees every nonzero submodule. When it does not, the closings show it, by a step of more than one in $\dim\mathcal S_D$, and the slope of $\dim\mathcal S_D$ is $\mathfrak md+1$; a slope of one over the degrees computed is evidence that nothing is missed, not a certificate, the degrees of the missed part not being bounded here. What the closings miss is missed by the values as well. It is no defect of the torus: a submodule that the native gauge of the direction does not see is invisible to the integral itself, and an element of it is a relation among values that the relations of the ray do not produce. The symmetries of the ray account for some of these, on $(1-e_G)\mathfrak M$, and so can symmetries that are not maps of its torus at all: when the list is the contraction of an identically self-dual table, the integral over the contracted scale carries the module of the larger table onto that of the ray, and the automorphisms of the larger table act through it.

% New: proved in the last part.
\begin{proposition}[A table one mark larger]\label{prop:onemark}
Let $(\cB',\theta')$ be a connected identically self-dual table on $\cE'=\cE\sqcup\{g\}$, $|\cE|\ge2$ (Corollary~\ref{cor:piecesmarks}(iii)), and $(\cB,\theta)$ its contraction at $g$ (Definition~\ref{def:delcon}), with letters $K,\Psi$; then $(\cB,\theta)$ is connected, and the letter of $\cB'$ is $K'=\Psi+\alpha_gK$. Let $F'$ and $F$ be the two tables at the strength $1$, read on their centre rays (Definition~\ref{def:centreray}), $F'$ on a chart $y_{e_0}=0$ with $e_0\in\cE$; their strides are $f'=\alpha^{\one}\alpha_g/K'^2$ and $f=\alpha^{\one}/(K\Psi)$, their modules $\mathfrak M'$ and $\mathfrak M$. For a Laurent polynomial $h$ on $\cE$, $a\in\mathbb Z$ and $j\in\mathbb N$ with $\alpha_g^aK'^{-j}h$ of degree $0$, put
\[
\pi_\ast\big(\alpha_g^aK'^{-j}h\big)=r_{a,j}(k)\Big(\frac\Psi K\Big)^{a}\Psi^{-j}h,\qquad r_{a,j}(k)=\frac{\Gamma(k+a)\,\Gamma(k+j-a)\,\Gamma(2k)}{\Gamma(k)^2\,\Gamma(2k+j)}\in\mathbb Q(k).
\]
\begin{enumerate}[label=(\roman*)]
\item For every insertion $g'$ of $F'$, at the points of the positive orthant and for all large $k$, $\int_0^\infty g'f'^k\,d\alpha_g/\alpha_g=\Gst21{k,\ k}{2k}f^k\,\pi_\ast g'$. So $\langle g'\rangle'_k=\Gst21{k,\ k}{2k}\langle\pi_\ast g'\rangle_k$ for all large $k$, and $\mathsf C'_j(g')=\binom{2j}j\,\mathsf C_j\big(\pi_\ast g'|_{k=-j}\big)$ for all large $j$. The native gauges of $F'$ are the points $\alpha_{c'}=\big(\alpha_c,\Psi(\alpha_c)/K(\alpha_c)\big)$ over the native gauges $\alpha_c$ of $F$, with $f'(\alpha_{c'})=f(\alpha_c)/4$, and $\alpha_{c'}$ is nondegenerate exactly when $\alpha_c$ is; then $A'_{c'}(g')=\Gst21{k,\ k}{2k}A_c(\pi_\ast g')$, the stacked Gamma read as its expansion at $k=\infty$ and the disc at $\alpha_{c'}$ taken over the disc at $\alpha_c$.
\item $\pi_\ast$ induces a bijection $\mathfrak M'\to\mathfrak M$, with $\pi_\ast\operatorname{Sh}'=\frac k{2(2k+1)}\operatorname{Sh}\pi_\ast$.
\item Every automorphism $\beta$ of $(\cB',\theta')$ acts on $\mathfrak M$ by $T_\beta=\pi_\ast\phi_\beta^\ast\pi_\ast^{-1}$, which commutes with $\operatorname{Sh}$; $T_\beta=\phi_{\beta}^\ast$ when $\beta(g)=g$, and when the transposition $\beta$ of a mark $e\in\cE$ and $g$ is an automorphism,
\[
T_\beta[\alpha^u]=\frac{\Gamma(k+u_e)\,\Gamma(k-u_e)}{\Gamma(k)^2}\,\Big[\alpha^{u-u_e\one_e}\Big(\frac\Psi K\Big)^{u_e}\Big].
\]
With the Cremona, which $\pi_\ast$ carries to the Cremona, the $T_\beta$ generate a finite group $G'$.
\item The values, the closings and $A_\star$, at the native gauge $\alpha^\star=\one$, are invariant under $G'$ and vanish on $(1-e_{G'})\mathfrak M$. For a group $G\subseteq G'$ of symmetries of $F$, $\mathfrak M^G=\mathfrak M^{G'}\oplus(1-e_{G'})\mathfrak M^G$ as submodules, and $\pi_\ast$ maps the invariant module of $F'$ under its automorphisms and the Cremona onto $\mathfrak M^{G'}$.
\end{enumerate}
\end{proposition}

% New: proved in the last part.
\begin{corollary}[The cousins of a fused table]\label{cor:cousinsonemark}
Let $(\cB,\theta)$ be a connected fusable table on $n$ marks, of the fusion width $c\ge1$, with $n+c\ge3$, and $(\cB^\sharp,\theta^\sharp)$ its Foreman fusion (Definition~\ref{def:fusion}). Then $\cB^\sharp$ is connected, and Proposition~\ref{prop:onemark} applies to $\cB^\sharp$ and each of its cousins (Proposition~\ref{prop:cousins}): the automorphisms of $\cB^\sharp$ act on the module of the centre ray of every cousin.
\end{corollary}

The scale of $g$ enters the letter of the larger table linearly, so the integral over it is a Beta corner, and it carries the centre of $F'$ onto the centre of $F$ insertion by insertion: a power of $\alpha_g$ becomes the same power of $\Psi/K$, the ratio at which the scale of $g$ is native, times a rational function of $k$. The relations at the scale of $g$ are exactly what that integral forgets, so the two modules are one, with the shift twisted by the Beta corner. An automorphism of $\cB'$ that moves $g$ is no symmetry of $F$, nor a map of its torus: it moves insertions, not points, and on native gauges it is the lift $\alpha_c\mapsto\alpha_{c'}$ followed by the permutation and the projection. The values, the closings and the expansion at $\alpha^\star$ are read on the orthant, the torus and the point $\one$ of $F'$, which every automorphism keeps; so they see only the part of $\mathfrak M$ that $G'$ fixes, and when $\mathfrak M^G$ is larger than $\mathfrak M^{G'}$ the closings miss the difference, with every symmetry of $F$ in $G$. A sign $\eta$ of $F$ lifts to the sign $(\eta,\eta^{\cE})$ of $F'$, since $\Psi(\eta)=\eta^{\cE}K(\eta)$, with $f(\eta)=4/K'(\eta,\eta^{\cE})^2$; signs that the automorphisms of $\cB'$ carry to one another have one native value. Proposition~\ref{prop:onemarkhidden} measures what is missed in this way: an orbit of nondegenerate native gauges of $F'$ that the symmetries of $F$ split into several orbits gives, orbit against orbit, solutions that vanish on $\mathfrak M^{G'}$ and not on $\mathfrak M^G$.



% New: proved in the last part.
\begin{proposition}[What a table one mark larger hides from the closings]\label{prop:onemarkhidden}
Let $(\cB',\theta')$, $g$, $F'$, $F$, $\pi_\ast$ and $G'$ be as in Proposition~\ref{prop:onemark}. Let $\Gamma'$ be the group of symmetries of the centre ray of $F'$ generated by the maps $\phi_\beta$, $\beta$ an automorphism of $(\cB',\theta')$, and the Cremona, and $G$ the group of symmetries of $F$ generated by the maps $\phi_\beta$, $\beta$ an automorphism of $(\cB',\theta')$ with $\beta(g)=g$ read on $\cE$, and the Cremona, so that $G\subseteq G'$ (Proposition~\ref{prop:onemark}(iii)). A native gauge $\alpha_c$ of $F$ lies \emph{below} the native gauge $\alpha_{c'}=\big(\alpha_c,\Psi(\alpha_c)/K(\alpha_c)\big)$ of $F'$ (Proposition~\ref{prop:onemark}(i)).
\begin{enumerate}[label=(\roman*)]
\item Let $O'_1,\dots,O'_q$ be distinct orbits of $\Gamma'$ of nondegenerate native gauges of $F'$, and let the native gauges of $F$ below $O'_a$ form $s_a$ orbits of $G$. Then
\[
\dim_{\mathbb C(k)}(1-e_{G'})\mathfrak M^G\ \ge\ \sum_{a=1}^q(s_a-1),
\]
and the closings, the values and $A_\star$ vanish on the submodule $(1-e_{G'})\mathfrak M^G$. For an orbit $O'$ fix an expansion $A'_{c'_0}$ at one of its points, put $A'_{\gamma(c'_0)}=A'_{c'_0}\circ\gamma^\ast$ for $\gamma\in\Gamma'$, and for $\alpha_c$ below $O'$ put $A_c=\Gst21{k,\ k}{2k}^{-1}A'_{c'}\circ\pi_\ast^{-1}$, the stacked Gamma read as its expansion at $k=\infty$. If $O_1$ and $O_2$ are two orbits of $G$ below $O'$, the solution $|O_1|^{-1}\sum_{c\in O_1}A_c-|O_2|^{-1}\sum_{c\in O_2}A_c$ is invariant under $G$, vanishes on $\mathfrak M^{G'}$ and does not vanish on $(1-e_{G'})\mathfrak M^G$.
\item For $r\ge2$ let $\cB'=\binom{[2r+2]}{r+1}$ with unit weights and $g=2r+2$, so that $F$ is $\binom{[2r+1]}r$ read on its centre ray at the strength $1$; $G$ is then the group generated by the symmetries of the four kinds of that ray (Proposition~\ref{prop:raysym}): the maps $\phi_\beta$, $\beta$ a bijection of $[2r+1]$, and the Cremona. With $C_r=\frac1{r+1}\binom{2r}r$, the signs of $F$ with one and with two scales $-1$, $2r+1$ and $\binom{2r+1}2$ of them, are nondegenerate native gauges of the native value $C_r^{-2}$ and form two orbits of $G$, and the native gauges above them are the $\binom{2r+2}2$ signs of $F'$ with two scales $-1$, one orbit of $\Gamma'$. At the sign of $F'$ with the scales $-1$ on $D=\{a_1,a_2\}$, with $C=[2r+2]\setminus D$, the Hessian $H'$ of $\log f'$ is such that $-H'$ has, modulo $\one$, the eigenvalue $r+1$ on $\one_{a_1}-\one_{a_2}$, the eigenvalue $\frac{r+1}{2r-1}$ on every $\one_c-\one_{c'}$ with $c,c'\in C$, and the eigenvalue $-(r+1)$ on $r\one_D-\one_C$; at the signs of $F$ with one and with two scales $-1$, on a chart, $\det(-H)=-(r+1)^{2r}/(2r-1)^{2r-1}$. So $(1-e_{G'})\mathfrak M^G\neq0$: at the centre of every $\binom{[2r+1]}r$, $r\ge2$, the closings miss a nonzero submodule of $\mathfrak M^G$, with every symmetry of the four kinds in $G$.
\end{enumerate}
\end{proposition}

The two kinds of signs of $\binom{[2r+1]}r$ share their native value because the signs above them are one orbit of the symmetric group of $\binom{[2r+2]}{r+1}$, while the symmetries of the ray keep the two kinds apart; the mean of the solutions over the one kind less the mean over the other is a solution that sees what the larger table hides. Exercise~\ref{ex:closingssee} reads the small centres with these statements; at $r=2$ the part of the span of a box that the larger table hides from the closings is a line.

\section{Wheels: the letters as a holonomy}\label{sec:wheels}

% New: proved in the last part.
\begin{proposition}[The holonomy of a wheel]\label{prop:wheelholonomy}
For $p\ge3$ let $W_p$ be the list of the wheel graph: a hub and rim vertices $1,\dots,p$ in cyclic order, the marks the spokes $s_i$, joining the hub to $i$, and the rims $r_i$, joining $i$ to $i+1$, indices modulo $p$; its bases are the spanning trees, of rank $p$ on $n=2p$ marks. With
\[
\mathbf U(a)=\begin{bmatrix}1&a\\0&1\end{bmatrix},\qquad \mathbf L(b)=\begin{bmatrix}1&0\\b&1\end{bmatrix},\qquad \mathbf h(\alpha)=\prod_{i=1}^p\mathbf U(\alpha_{s_i})\,\mathbf L(\alpha_{r_i}^{-1})\in SL_2 ,
\]
the product taken in the order $i=1,\dots,p$,
\[
K(\alpha)=\prod_i\alpha_{r_i}\,\big(\operatorname{tr}\mathbf h(\alpha)-2\big),\qquad \Psi(\alpha)=\prod_i\alpha_{s_i}\,\big(\operatorname{tr}\mathbf h(\alpha^{-1})-2\big).
\]
The bijection $x$ with $x(s_i)=r_{p+1-i}$ and $x(r_i)=s_{p+1-i}$ has $\operatorname{tr}\mathbf h(\alpha\circ x)=\operatorname{tr}\mathbf h(\alpha^{-1})$, so $K\circ\phi_x=\Psi$ and $x$ exchanges $W_p$. At the table's gauge $\mathbf h(\one)=\begin{bmatrix}2&1\\1&1\end{bmatrix}^p$, whose eigenvalues are $\lambda_\pm^p$ with $\lambda_\pm=\frac{3\pm\sqrt5}2$, the squares of the golden ratio and of its inverse; so $K(\one)=\lambda_+^p+\lambda_-^p-2$, the Lucas number of index $2p$ less $2$.
\end{proposition}

% New: proved in the last part.
\begin{proposition}[The centre ray of a wheel]\label{prop:wheelcentre}
Put $\mathfrak d(\alpha)=\operatorname{tr}\mathbf h(\alpha)-2$ and $\mathfrak d^\ast(\alpha)=\operatorname{tr}\mathbf h(\alpha^{-1})-2=\mathfrak d(\alpha\circ x)$. The centre ray of $W_p$ at the strength $\sigma$ has the stride $f=(\mathfrak d\,\mathfrak d^\ast)^{-\sigma}$, and $X_F=\{\mathfrak d\,\mathfrak d^\ast\neq0\}$.
\begin{enumerate}[label=(\roman*)]
\item Its native gauges are the points of $X_F$ with $\mathfrak d^\ast\,d\mathfrak d+\mathfrak d\,d\mathfrak d^\ast=0$. At every one $d\mathfrak d$ and $d\mathfrak d^\ast$ are linearly dependent: the native gauges lie on the critical locus of the pair of traces $(\mathfrak d,\mathfrak d^\ast)$, where its differential has rank at most one. The native value is $(\mathfrak d\,\mathfrak d^\ast)^{-\sigma}$.
\item At a sign $\eta$, $\mathfrak d^\ast(\eta)=\mathfrak d(\eta)$ and the native value is $\mathfrak d(\eta)^{-2\sigma}$; at the table's gauge it is $(\lambda_+^p+\lambda_-^p-2)^{-2\sigma}$.
\item Every point with $\mathbf h(\alpha)=-I$ has $d\mathfrak d=0$. The \emph{torsion stratum} $\{\mathbf h(\alpha)=\mathbf h(\alpha^{-1})=-I\}$ lies in $X_F$ and consists of native gauges of native value $16^{-\sigma}$ at which $d\mathfrak d=d\mathfrak d^\ast=0$; it is cut out of the torus of dimension $2p-1$ by six equations, so that each of its components has dimension at least $2p-7$.
\end{enumerate}
\end{proposition}

% New: proved in the last part.
\begin{proposition}[The classes of a wheel]\label{prop:wheelclasses}
At the table's gauge the occupation of a spoke is
\[
R_{\mathrm s}=\frac{\lambda_+^p-\lambda_-^p}{\sqrt5\,\big(\lambda_+^p+\lambda_-^p-2\big)},
\]
equal to $\frac12$ exactly for $p=3$; the occupation of a rim is $R_{\mathrm r}=1-R_{\mathrm s}$, and the masses of $\Upsilon_{\sigma,\delta;\one}$ are $\sigma+\delta(2R_{\mathrm s}-1)$ on the spokes and $\sigma-\delta(2R_{\mathrm s}-1)$ on the rims. $W_p$ is an atom at the centre. For $p\ge4$ its classes at every lean are the spokes and the rims, which the reversed exchange $\bar\phi_x$, a symmetry of the integrand at every lean (Theorem~\ref{thm:reversed}), interchanges; $W_3=M(K_4)$ is an atom at every lean.
\end{proposition}

\begin{remark}\label{rem:wheeluniform}
The uniform masses $\sigma\one$ lie on the balance wall of $W_p$ at every lean, since $2r=n$, and they meet the condition of Proposition~\ref{prop:raysym} for $\bar\phi_x$, which therefore keeps the integrand of every such realization; its native gauge is constant on the spokes and on the rims. For $p=3$ the uniform masses are the masses of the table's gauge at every lean, $R_{\mathrm s}$ being $\frac12$, and the native gauge stays at $\one$; for $p\ge4$ they are not, and the native gauge leaves $\one$ as soon as the realization leans. For $W_4$, at the gauge $c$ on the spokes and $1$ on the rims, $K=c(c+2)^2(c+4)$ and $\Psi=(2c+1)^2(4c+1)$. At the full basis lean nativity asks every occupation to be $\frac12$, which on a spoke reads $c^2+2c-4=0$: the native gauge on the orthant has spokes $\sqrt5-1$, twice the inverse of the golden ratio, against rims $1$, and at the full cobasis lean, where the condition reads $4c^2-2c-1=0$, spokes $\frac{\sqrt5+1}4$, half the golden ratio, its reciprocal. The other roots, $-\sqrt5-1$ and $\frac{1-\sqrt5}4$, are native gauges as well, with negative spokes and in $X_F$: sources off the orthant. At the basis lean the integrand is $c^{4\sigma}K^{-2\sigma}=\big(c/((c+2)^2(c+4))\big)^{2\sigma}$, whose base is $\frac{5\sqrt5-11}8$ at the positive root and $-\frac{5\sqrt5+11}8$ at the negative one, a quarter of the inverse fifth power of the golden ratio and minus a quarter of its fifth power; the cobasis lean gives the same two numbers. Off the centre the gauge on the orthant is not proportional to $\one$, so the realization is not an atom (Theorem~\ref{thm:gaugestep}); its classes are the spokes and the rims, and $\bar\phi_x$ interchanges them. The torsion stratum of $W_4$ is a union of curves.
\end{remark}

The letters of a wheel are traces of one product of unipotent matrices, the holonomy of a transfer around the rim, and the planar self-duality of the wheel is the transpose of that product. On the centre ray the stride is the inverse product of the two traces, at $\alpha$ and at $\alpha^{-1}$, so a native gauge lies on the critical locus of the pair of traces, and wherever both holonomies are $-I$ both traces are critical at once: a stratum of native gauges of native value $16^{-\sigma}$, of dimension at least $2p-7$ wherever it is not empty. At a sign $\eta$ the two traces agree, in an integer $\mathfrak d(\eta)$, and the native value is $\mathfrak d(\eta)^{-2\sigma}$; at the table's gauge the holonomy is the $p$-th power of the step $\begin{bmatrix}2&1\\1&1\end{bmatrix}$, whose eigenvalues are the squares of the golden ratio and of its inverse. The wheel is an atom at the centre and splits into spokes and rims at every lean once $p\ge4$, the reversed exchange interchanging them; with uniform masses its native gauge then leaves the table's gauge as soon as the realization leans, and for $W_4$ at a full lean the ratio of the scales of spokes and rims is twice the inverse of the golden ratio or half the golden ratio.

\section{The procedure}\label{sec:machineproc}

The statements of this chapter assemble into a procedure that reads the one-step matrix and the law of $[1]$ of a ray from its two ends, with no box (Definition~\ref{def:raybox}) reduced.
\begin{enumerate}[label=(\arabic*)]
\item \emph{Symmetries.} From the list read the automorphisms and the exchanges, and keep the maps of Proposition~\ref{prop:raysym} that are symmetries of the ray; at a lean of a self-dual family at the table's gauge these include the reversed exchanges (Theorem~\ref{thm:reversed}). Work in $\mathfrak M^G$ with invariant insertions (Theorem~\ref{thm:invariant}). When the list is the contraction of an identically self-dual table at a mark, the automorphisms of that table act as well, through the integral over the scale of the mark, and the values live on the part they fix, which the centre of the larger table reads directly (Proposition~\ref{prop:onemark}).
\item \emph{Sources.} Read the critical locus of the stride modulo $G$: orbits of nondegenerate native gauges, one solution each; $A_\mu$ points, with the classes of trivial character (Theorems~\ref{thm:amu} and~\ref{thm:invariantcount}); clean families, with their averaged cycles (Theorem~\ref{thm:closedcycle}, Corollary~\ref{cor:curvesources}, Proposition~\ref{prop:sourcesym}). Their native values and exponents enter the formal type of the one-step matrix at infinity (Proposition~\ref{prop:formaltype}).
\item \emph{Closings.} Compute the closings of the candidate masters, constant terms of Laurent polynomials (Definition~\ref{def:closings}).
\item \emph{Masters.} Choose invariant insertions whose solution vectors over the supplied sources and the closings are linearly independent; there are at most $\dim\mathfrak M^G$ of them, and the count of Theorem~\ref{thm:invariantcount} bounds $\dim\mathfrak M^G$ below.
\item \emph{The one-step matrix.} Solve the homogeneous linear equations of Theorem~\ref{thm:determine}(iv) at increasing degree $D$. When the supplied solutions annihilate no submodule, the first degree with a nonzero solution is $D_0$, the solution is unique up to a constant factor, and it is $\mathbf B_{\mathbf b}$.
\item \emph{The law.} The law of $[1]$ is the least relation among $[1],\operatorname{Sh}[1],\dots$, computed with $\mathbf B_{\mathbf b}$; on a ray from the origin it is read again on the closings of $1$ at $-j$ (Corollary~\ref{cor:reflected}).
\end{enumerate}
A case is certified when the supplied solutions annihilate no submodule: by Corollary~\ref{cor:glue} this holds when $\mathfrak m$ supplied formal solutions are linearly independent over $\mathbb C$ on $\mathfrak M^G$ type by type, a full set from the high end; and when the walls of the direction are positive the closings alone certify exactly when the native gauge of the direction sees every nonzero submodule of $\mathfrak M^G$ (Corollary~\ref{cor:seeall}), as when $\mathfrak M^G$ is irreducible and they do not vanish on it. The closings show it when they miss something: the dimension of the solutions of their equations then has slope $\mathfrak md+1$, with $d$ the dimension of the largest missed submodule $\operatorname{Ann}(V_{\mathsf C})\cap\mathfrak M^G$, and a single step of more than one proves $d>0$ (Proposition~\ref{prop:blinddim}); a slope of one over the degrees computed is evidence that nothing is missed, not a certificate. When they miss something they still give the one-step matrix of the value module and the law of the values (Corollary~\ref{cor:valuemodule}), and the missing part is read from sources whose solutions see it: when it comes from a table one mark larger, the centre of that table reads $\mathfrak M^{G'}$ (Proposition~\ref{prop:onemark}), and the differences of the solutions over the orbits of native gauges that the symmetries of the ray split do not vanish on the rest (Proposition~\ref{prop:onemarkhidden}). With fewer formal solutions and no such knowledge of $\mathfrak M^G$, a unique solution at the least degree that also meets equations held out is evidence, not a certificate: a submodule annihilated by every supplied solution would leave a different matrix that meets all of them. Masters that span $\mathfrak M^G$ come from a box when every letter exponent of the direction is positive: the spanning set of a box accepted at $\mathfrak s$ spans the classes of the invariant Laurent combinations of wall degree at most $\mathfrak s$ (Theorem~\ref{thm:raybox}(i)), which increase with $\mathfrak s$, and some box is accepted with a spanning set whose classes span $\mathfrak M^G$ (Theorem~\ref{thm:rayboxstop}). That spanning set is a set of masters of $\mathfrak M^G$ when it is independent: when as many solutions as it has elements are linearly independent on its span, type by type (Lemma~\ref{lem:types}), or when its values at nondegenerate native gauges span $\mathbb C^{\mathfrak m}$ (Proposition~\ref{prop:rays-spectrum}(i)).


% ---- New chapter: Foreman fusion (inserted after ch:twoends) ----
% ---- New chapter: the stride's law and the gauge end (after ch:twoends) ----
% No new macros. Local symbols: \mathfrak D_{\mathbf p} (the density operator of a recurrence),
% \mu_f, \varrho_f (the stride's law and its density), \nu^{\mathbb T}_f, \varrho^{\mathbb T}_f and
% \langle1\rangle^{\mathbb T}_k (the torus law, its density and the torus value).
\chapter{The stride's law and the gauge end}\label{ch:strideslaw}

The values of a ray with real origin and positive walls are the moments of one measure, the stride's law: the image of the origin form on the positive part of the chart under the stride. It is carried by $(0,f(\alpha^\star)]$, and the values are read from it three ways, as a Trailmoment on a line, through one Gramma on the Laplace side, and as the log-partition of the exponential family that the Breta shape laws form along the ray. Its two ends are the two ends of strength: near $0$ its mass grows like a power of the logarithm, with the volume of the direction as coefficient, and at the native value $f(\alpha^\star)$ it ends in a power $(f(\alpha^\star)-t)^{\frac{n-3}2}$ times a function holomorphic there, whose value is read off the Hessian at the native gauge. Its density is real-analytic in between. Every recurrence of the values is a differential equation for the density, the recurrence read backwards, and holds at every strength once it holds at the large integers. At the native value of the direction, when it is a simple root of the symbol, the germs of solutions that the equation lets end there form a line; the density spans it, its exponent is forced to $\frac{n-3}2$, and its continuation along the open support is the density: so the recurrence, the native gauge of the direction and the Hessian there give the values and their derivatives in the strength exactly, with no datum from the low end and no inequality. At a centre the same native gauge lies on the unit torus, where the stride is real and at least $f(\alpha^\star)$; the torus value continues the closings in the strength, and the densities of the chart and of the torus at the native value are the two cut solutions of one function, on its two sides.

\section{The stride's law}\label{sec:sl-law}

% New: proved in the last part.
\begin{definition}[The stride's law]\label{def:strideslaw}
Let a ray of a family on $n\ge2$ marks have a real origin and $\kappa_{1,S}>0$ for every nonempty proper $S$ (Definition~\ref{def:ray}). On the positive part of the chart, $\alpha_{e_0}=1$ and $y=\log\alpha\in\mathbb R^{\cE\setminus e_0}$, the origin form $\omega_0=(\omega_0/dy)\,dy$ is a positive measure, its density $\alpha^{\mvec_0}\prod_jK_j^{-s_{0j}}\Psi_j^{-t_{0j}}$ taken with the principal powers. The \emph{stride's law} of the ray is its image under the stride,
\[
\mu_f(A)=\int_{\{y:\ f(e^y)\in A\}}\omega_0\qquad(A\subseteq(0,\infty)\ \text{a Borel set}),
\]
a measure on $(0,\infty)$, and $\varrho_f$ is its density where it has one. For the origin $0$, $\omega_0=dy$, and $\mu_f$ is the image under $f$ of the Lebesgue input of the chart (Definition~\ref{def:inputs}).
\end{definition}

% New: proved in the last part.
\begin{proposition}[A ray is one Trailmoment on a line and one Gramma]\label{prop:sl-moments}
Let a ray be as in Definition~\ref{def:strideslaw}, let $\alpha^\star$ be the native gauge of its direction on the positive orthant, $\alpha^\star_{e_0}=1$ (Theorem~\ref{thm:pencil}), and write in this statement $f^\star=f(\alpha^\star)$.
\begin{enumerate}[label=(\roman*)]
\item $\mu_f$ is carried by $(0,f^\star]$, it is finite on $[t,f^\star]$ for every $t>0$, and at every complex $k$ at which $\langle1\rangle_k$ converges
\[
\langle1\rangle_k=\int_{(0,f^\star]}t^k\,\mu_f(dt)=\big(\sT_{k;(1)}[t^{-1}]\,\mu_f\big)\big((0,\infty)\big)=\Tx{k}{\iota_\ast\mu_f}{0},\qquad\iota(t)=t^{-1}:
\]
the values are the Trailmoments on a line of index $k$, one unit exponent and the coefficient function $x$ of one input, the image of $\mu_f$ under the reciprocal.
\item Let the origin carry no letter exponents, $\mathbf s_0=\mathbf t_0=\mathbf 0$, so that $\omega_0=\alpha^{\mvec_0}dy$, and write $f^{-1}=\sum_iw_i\,\alpha^{\mathbf v_i}$ on the chart, a Laurent polynomial in the scales $\alpha_e$, $e\neq e_0$, with positive coefficients $w_i$ and exponents $\mathbf v_i\in\mathbb Z^{\cE\setminus e_0}$; put $\mathbf v=(\mathbf v_i)_i$ and $\mathbf w=(w_i)_i$. The cone of $\mathbf v$ is $\mathbb R^{\cE\setminus e_0}$, so the Gramma $\sG_{\nu;\mathbf v}(u\mathbf w)$ is entire in $\nu$ (Theorem~\ref{thm:grammadomain}), and at every $k$ with $\re k>0$ at which $\langle1\rangle_k$ converges
\[
\Gamma(k)\,\langle1\rangle_k=\int_0^\infty u^{k-1}\,\sG_{\mvec_0|_{\cE\setminus e_0};\,\mathbf v}(u\,\mathbf w)\,du ,
\]
the Laplace side (Theorem~\ref{thm:laplaceside}) of (i), whose Laplace transform beyond the cut $0$ is the Gramma of the reciprocal stride.
\item For real $k$ at which $\langle1\rangle_k$ converges, let $U\sim\mathrm{Breta}^\circ(F;\mvec_0+k\mvec_1;\mathbf s_0+k\mathbf s_1,\mathbf t_0+k\mathbf t_1)$ (Definition~\ref{def:shapelaw}). The law of $f(U)$ is $t^k\mu_f(dt)/\langle1\rangle_k$: along the ray the shape laws are the exponential family with the sufficient statistic $\log f$, the base measure $\mu_f$ and the log-partition $\log\langle1\rangle_k$, so that
\[
\frac{d}{dk}\log\langle1\rangle_k=\oE\big[\log f(U)\big],\qquad\frac{d^2}{dk^2}\log\langle1\rangle_k=\operatorname{Var}\big[\log f(U)\big]>0 .
\]
\end{enumerate}
\end{proposition}

% New: proved in the last part.
\begin{proposition}[The two ends of the stride's law]\label{prop:sl-ends}
Let a ray be as in Definition~\ref{def:strideslaw}, with $\alpha^\star$ as in Proposition~\ref{prop:sl-moments} and $H$ the Hessian of $-\log f$ at $\log\alpha^\star$ on the chart (Theorem~\ref{thm:branch}); write in this statement $f^\star=f(\alpha^\star)$ and $w^\star=(\omega_0/dy)(\alpha^\star)$.
\begin{enumerate}[label=(\roman*)]
\item \emph{The open support.} $\mu_f$ has a density $\varrho_f$, real-analytic and positive on $(0,f^\star)$.
\item \emph{The native value.} Near $f^\star$, $\varrho_f(t)=(f^\star-t)^{\frac{n-3}2}g(f^\star-t)$ for $t<f^\star$, with $g$ holomorphic at $0$ and
\[
g(0)=\frac{w^\star}{\Gamma(\frac{n-1}2)\sqrt{\det H}}\Big(\frac{2\pi}{f^\star}\Big)^{\frac{n-1}2}.
\]
\item \emph{The low end.} Let the origin be $0$, let $\kappa$ be the wall at a tilt of the question of the direction (Definitions~\ref{def:question} and~\ref{def:tiltwall}), read at $y$, and put $C_-=\prod_jK_j(\one)^{-s_{1j}-t_{1j}}$, $C_+=\prod_j\theta_{j,\min}^{-s_{1j}-t_{1j}}$, $\theta_{j,\min}$ the least weight of table $j$. Then $C_-e^{-\kappa(y)}\le f(e^y)\le C_+e^{-\kappa(y)}$ on the chart, $\int e^{-\kappa(y)}dy=V_1$ with the volume $V_1=\sum_\pi\prod_{l=2}^n\kappa_{1,U_l}^{-1}$ of the direction (Proposition~\ref{prop:principal}, Theorem~\ref{thm:severity}(i)), and for $0<t\le C_-$
\[
\frac{V_1}{(n-1)!}\Big(\log\frac{C_-}t\Big)^{n-1}\le\mu_f\big([t,f^\star]\big)\le\frac{V_1}{(n-1)!}\Big(\log\frac{C_+}t\Big)^{n-1}.
\]
So $\mu_f([t,f^\star])\sim\frac{V_1}{(n-1)!}(\log\frac1t)^{n-1}$ as $t\downarrow0$, and $C_-^kV_1\le k^{n-1}\langle1\rangle_k\le C_+^kV_1$ for every real $k>0$.
\end{enumerate}
\end{proposition}

The two ends of the stride's law are the two ends of strength. At the low end the mass of $\mu_f$ near $0$ is set by the walls of the direction alone, through the volume of the direction, and the values blow up at $k=0$ with that volume as coefficient (Proposition~\ref{prop:rays-lowpoles}). At the high end $\mu_f$ ends at the native value of the direction in a power fixed by the number of marks, with a coefficient read off the Hessian of the stride at its native gauge; the branch of the generating function in Theorem~\ref{thm:branch} is this end read through the pencil.

\section{The gauge end}\label{sec:sl-gauge}

% New: proved in the last part.
\begin{definition}[The density operator of a recurrence]\label{def:sl-densityop}
Let $\mathbf p=(p_0,\dots,p_h)$ be polynomials in $\mathbb C[k]$, not all zero, of largest degree $D$, with the symbol $\chi(\lambda)=\sum_ic_i\lambda^i$, $c_i=[k^D]p_i$ (Corollary~\ref{cor:symbol}), and the \emph{second symbol} $\chi_1(\lambda)=\sum_i[k^{D-1}]p_i\,\lambda^i$. Its \emph{density operator} is
\[
\mathfrak D_{\mathbf p}=\sum_{i=0}^ht^i\,p_i(-1-i-\theta),\qquad\theta=t\frac d{dt},
\]
a differential operator with polynomial coefficients acting on functions and distributions on $(0,\infty)$; written $\sum_{j=0}^Da_j(t)\,\partial_t^j$, its leading coefficient is $a_D=(-1)^Dt^D\chi(t)$. At a root $\lambda\neq0$ of $\chi$, in the variable $s=\lambda-t$, write $\mathfrak D_{\mathbf p}=\sum_jb_j(s)\,\partial_s^j$ and $b_{j,l}=[s^l]\,b_j$. A function $u$ on an interval $(\lambda-\delta,\lambda)$, or $(\lambda,\lambda+\delta)$, \emph{extends by zero} at $\lambda$ when it is integrable there, $\mathfrak D_{\mathbf p}u=0$ there, and the distribution $u\,\one_{(\lambda-\delta,\lambda)}$, or $u\,\one_{(\lambda,\lambda+\delta)}$, is annihilated by $\mathfrak D_{\mathbf p}$ on $(\lambda-\delta,\lambda+\delta)$; this is a property of the germ of $u$ at $\lambda$.
\end{definition}

The density operator is the recurrence read on the other side of the moments: $\mathfrak D_{\mathbf p}$ is the operator whose formal adjoint sends $t^k$ to $\sum_ip_i(k)\,t^{k+i}$, so a measure on $(0,\infty)$ annihilated by it has moments that satisfy the recurrence, and its singular points in $\mathbb C^\times$ are the roots of the symbol.

% New: proved in the last part.
\begin{proposition}[A recurrence of the values is an equation for the density]\label{prop:sl-density}
Let a ray be as in Definition~\ref{def:strideslaw}, let $f^\star$ be as in Proposition~\ref{prop:sl-moments} and $\mathbf p=(p_0,\dots,p_h)$ as in Definition~\ref{def:sl-densityop}. The following are equivalent:
\begin{enumerate}[label=(\alph*)]
\item $\sum_ip_i(k)\langle1\rangle_{k+i}=0$ for all large integers $k$;
\item $\sum_ip_i(k)\langle1\rangle_{k+i}=0$ at every complex $k$ at which every $\langle1\rangle_{k+i}$ converges, and as an identity of meromorphic functions of $k$ (Proposition~\ref{prop:rays-strength}(ii));
\item the distribution $\mathfrak D_{\mathbf p}\,\mu_f$ on $(0,\infty)$ vanishes.
\end{enumerate}
When they hold, $\mathfrak D_{\mathbf p}\varrho_f=0$ on $(0,f^\star)$ and the germ of $\varrho_f$ at $f^\star$ extends by zero; and for the origin $0$, the conclusions of Corollary~\ref{cor:rays-lowend}(i) and~(ii) hold for $p_0,\dots,p_h$: $\sum_{i\le\min(h,j)}p_i(-j)\CT(f^{i-j})=0$ for every $j\in\mathbb N$, and $p_0$ vanishes to order at least $n-1$ at $0$ with $[k^{n-1}]p_0\cdot\mathsf L_0+\sum_{i\ge1}p_i(0)\langle1\rangle_i=0$.
\end{proposition}

So a recurrence that the values satisfy at the large integers, whether or not it follows from the relations of the ray, holds at every strength: no bound on the pieces of a certificate is needed down to $k\to0$, and the low-end relation follows for it as for the law of $[1]$.

% New: proved in the last part.
\begin{theorem}[The gauge end]\label{thm:gaugeend}
Let a ray be as in Definition~\ref{def:strideslaw}, with $\alpha^\star$, $f^\star=f(\alpha^\star)$, the density $\varrho_f$ and the function $g$ of Proposition~\ref{prop:sl-ends}, and let $\mathbf p=(p_0,\dots,p_h)$ satisfy the equivalent conditions of Proposition~\ref{prop:sl-density}; let $f^\star$ be a simple root of the symbol $\chi$, and use the notation of Definition~\ref{def:sl-densityop} at $\lambda=f^\star$. Put $[x]_j=x(x-1)\cdots(x-j+1)$ and $R_m(x)=\sum_jb_{j,m-D+1+j}\,[x]_j$, with $b_{j,l}=0$ for $l<0$.
\begin{enumerate}[label=(\roman*)]
\item $D\ge1$, $b_{D,0}=0$, $b_{D,1}=-(f^\star)^D\chi'(f^\star)\neq0$, and $s^{D-1}\mathfrak D_{\mathbf p}=\sum_{m\ge0}s^mR_m(s\partial_s)$ with
\[
R_0(x)=b_{D,1}\,[x]_{D-1}\,(x-\gamma),\qquad\gamma=D-1-\frac{b_{D-1,0}}{b_{D,1}}=-1+\frac{\chi_1(f^\star)}{f^\star\chi'(f^\star)} :
\]
$s=0$ is a regular singular point of $\mathfrak D_{\mathbf p}$ with the exponents $0,1,\dots,D-2$ and $\gamma$.
\item \emph{The line.} The germs at $f^\star$ of functions on the chart's side, $t<f^\star$, that extend by zero form a space of dimension at most one.
\item The germ of $\varrho_f$ extends by zero, so it spans that space, and $\gamma=\frac{n-3}2$.
\item \emph{The germ.} If $\gamma\notin\{0,1,\dots,D-2\}$, then $g(s)=g(0)\,\Phi(s)$, where $\Phi=\sum_N\Phi_Ns^N$, $\Phi_0=1$, is fixed by $R_0(\gamma+N)\Phi_N=-\sum_{m=1}^NR_m(\gamma+N-m)\Phi_{N-m}$, $N\ge1$. If $\gamma\in\{0,1,\dots,D-2\}$, then $\varrho_f=\sum_ru_rs^r$ near $f^\star$ with $u_r=0$ for $r<\gamma$, $u_\gamma=g(0)$, the $u_r$ with $\gamma<r\le D-2$ fixed one at a time by $\mathsf p_r(u)=0$, where
\[
\mathsf p_r(u)=\sum_{i\le r}\ \sum_{i<j\le D}(-1)^l\binom{j-1-i}{l}\,i!\,l!\,b_{j,l}\,u_i,\qquad l=j+r+1-D-i\ge0,
\]
whose coefficient of $u_r$ is $r!\,b_{D,1}(r-\gamma)$, and the $u_N$ with $N\ge D-1$ by $R_0(N)u_N=-\sum_{m=1}^NR_m(N-m)u_{N-m}$.
\item \emph{The values.} $\varrho_f$ extends holomorphically to a neighbourhood of $(0,f^\star)$ in $\mathbb C$ on which $\mathfrak D_{\mathbf p}\varrho_f=0$: it is the continuation of the germ of (iv) along $(0,f^\star)$, around the roots of $\chi$ on the segment by small arcs on either side. For every $j\in\mathbb N$ and every $k$ at which $\langle1\rangle_k$ converges,
\[
\frac{\partial^j}{\partial k^j}\langle1\rangle_k=\int_0^{f^\star}t^k\,(\log t)^j\,\varrho_f(t)\,dt .
\]
\end{enumerate}
\end{theorem}

The theorem reads a ray from its gauge end. The recurrence fixes the operator, and the native gauge of the direction fixes the rest: $f^\star$, the exponent, which the recurrence must reproduce, and the scale $g(0)$, through the Hessian. For odd $n$ the exponent is an integer and the line lies among the holomorphic solutions; when the exponent is at most $D-2$, the equations $\mathsf p_r(u)=0$, the absence of point masses at $f^\star$, single it out. For even $n$ it is the half-integer class. Nothing from the low end enters, nor any inequality, nor the native values on the segment $(0,f^\star)$, which the density passes analytically.

% New: proved in the last part.
\begin{proposition}[The law knows the dimension of its chart]\label{prop:sl-exponents}
Let a ray have $\kappa_{1,S}>0$ for every nonempty proper $S$, let $\mathbf p=(p_0,\dots,p_h)$ have $\sum_ip_i(k)f^i\in\mathcal L$ (Definition~\ref{def:raygauge}), as the law of $[1]$ has, and let $\lambda$ be a simple root of its symbol that is the native value of a nondegenerate native gauge. Then the formal solutions of $\sum_ip_i(k)x(k+i)=0$ of type $(\lambda,\rho+\mathbb Z)$ (Definition~\ref{def:raysolution}) form a line, spanned by $\lambda^kk^{\rho_\lambda}(1+\sum_{m\ge1}e_mk^{-m})$ with $\rho_\lambda=-\chi_1(\lambda)/(\lambda\chi'(\lambda))=-\frac{n-1}2$, and $\lambda$ is a regular singular point of $\mathfrak D_{\mathbf p}$ with the exponents $0,1,\dots,D-2$ and $-1-\rho_\lambda=\frac{n-3}2$.
\end{proposition}

So a recurrence read off closings, integers indexed by one strength, has to carry the dimension $n-1$ of the chart at every simple root of its symbol that is a native value: a test that a candidate law must pass. A root of the symbol need not be a native value, and at such a root the exponent is whatever the recurrence makes it.

% New: proved in the last part.
\begin{proposition}[The germ at the gauge end is local data at the native gauge]\label{prop:sl-jet}
Let a ray have the origin $0$ and $\kappa_{1,S}>0$ for every nonempty proper $S$, with $\alpha^\star$, $f^\star=f(\alpha^\star)$, $H$ and $g$ as in Proposition~\ref{prop:sl-ends}; in the coordinates $y$ of the chart, centred at $\log\alpha^\star$, write $\log f=\log f^\star-\frac12H_{ab}y^ay^b+\frac16h_{abc}y^ay^by^c+\frac1{24}h_{abcd}y^ay^by^cy^d+\cdots$, with the summation over repeated indices in $\cE\setminus e_0$ and $G=H^{-1}$.
\begin{enumerate}[label=(\roman*)]
\item For $m\ge2$ the Taylor tensor of order $m$ of $\log f$ at $\log\alpha^\star$ is $-2^{-m}$ times rung $m$ of the still of the realization $\big((s_{1j},t_{1j})_j;\alpha^\star\big)$ (Definition~\ref{def:still}), that is $-\sum_j\big(s_{1j}\kappa_{j,\alpha^\star}+(-1)^mt_{1j}\kappa_{j,(\alpha^\star)^{-1}}\big)$ in the cumulants of the basis laws (Proposition~\ref{prop:deviation}).
\item As $k\to\infty$,
\[
\langle1\rangle_k=(f^\star)^k\Big(\frac{2\pi}k\Big)^{\frac{n-1}2}\frac1{\sqrt{\det H}}\Big(1+\frac{c_1}k+O(k^{-2})\Big),
\]
with
\[
c_1=\tfrac18h_{abcd}G^{ab}G^{cd}+\tfrac18h_{abc}h_{def}G^{ab}G^{cd}G^{ef}+\tfrac1{12}h_{abc}h_{def}G^{ad}G^{be}G^{cf}.
\]
\item With $g(s)=\sum_mg_ms^m$, the asymptotic expansion of $\langle1\rangle_k$ is $\sum_mg_m(f^\star)^{k+\frac{n-1}2+m}\mathrm B\big(k+1,\frac{n-1}2+m\big)$, read term by term, and the term of $g_m$ begins at $(f^\star)^kk^{-\frac{n-1}2-m}$ with the coefficient $(f^\star)^{\frac{n-1}2+m}\Gamma(\frac{n-1}2+m)g_m$. So the coefficients of the expansion of $(f^\star)^{-k}k^{\frac{n-1}2}\langle1\rangle_k$ in powers of $k^{-1}$ and the Taylor coefficients of $g$ determine each other, the $m$-th by the first $m+1$.
\end{enumerate}
\end{proposition}

% New: proved in the last part.
\begin{proposition}[Slopes from the gauge end]\label{prop:sl-slopes}
Let a ray and $\mathbf p$ be as in Theorem~\ref{thm:gaugeend}, with $\varrho_f(t)=(f^\star-t)^{\frac{n-3}2}\sum_mg_m(f^\star-t)^m$ near $f^\star$.
\begin{enumerate}[label=(\roman*)]
\item For every $j\in\mathbb N$, $\sum_{l=0}^j\binom jl\sum_ip_i^{(j-l)}(k)\,\partial_k^l\langle1\rangle_{k+i}=0$ identically in $k$: the $j$-th derivative of the values satisfies the recurrence, with an inhomogeneity made of the lower derivatives.
\item If the series $\sum_mg_ms^m$ converges on $[0,f^\star)$ and $\sum_m|g_m|(f^\star)^m\mathrm B\big(K'+1,\frac{n-1}2+m\big)<\infty$ for a real $K'$, then for every real $K>K'$ and every $j\in\mathbb N$
\[
\partial_k^j\langle1\rangle_k\big|_{k=K}=\sum_{m\ge0}g_m\,\frac{\partial^j}{\partial a^j}\Big[(f^\star)^{a+\frac{n-1}2+m}\,\mathrm B\big(a+1,\tfrac{n-1}2+m\big)\Big]_{a=K}.
\]
\end{enumerate}
\end{proposition}

The derivatives of the values in the strength, the slopes, need no datum from the low end either: they are the logarithmic moments of the same density, and the recurrence differentiated carries them down from large strengths once they are read there.

% New: proved in the last part.
\begin{definition}[The procedure that reads a ray from its gauge end]\label{def:sl-gaugeread}
Input: a ray as in Definition~\ref{def:strideslaw} with its data in a field $\Bbbk\subseteq\mathbb R$, and polynomials $\mathbf p=(p_0,\dots,p_h)$ over $\Bbbk$ with $\sum_ip_i(k)\langle1\rangle_{k+i}=0$ for all large integers $k$, $f^\star=f(\alpha^\star)$ a simple root of their symbol. Output: the germ of the density at $f^\star$, and the values and slopes at the positive integers.
\begin{enumerate}[label=(\arabic*)]
\item \emph{The native gauge.} $\alpha^\star$, the unique critical point of $\log f$ on the positive part of the chart, where $-\log f$ is strictly convex (Theorem~\ref{thm:severity}(ii)); $f^\star$, $\det H$ and $(\omega_0/dy)(\alpha^\star)$; the constant $g(0)$ (Proposition~\ref{prop:sl-ends}(ii)).
\item \emph{The operator.} $\mathfrak D_{\mathbf p}$ and its coefficients $b_{j,l}$ at $f^\star$ (Definition~\ref{def:sl-densityop}), and $\gamma$ (Theorem~\ref{thm:gaugeend}(i)); a recurrence with $\gamma\neq\frac{n-3}2$ is not a recurrence of the values (Theorem~\ref{thm:gaugeend}(iii)).
\item \emph{The germ.} The Taylor coefficients of $g$, by Theorem~\ref{thm:gaugeend}(iv), over $\Bbbk(f^\star)$ up to the factor $g(0)$.
\item \emph{Large strengths.} For $K$ large and $0<\delta$ less than the radius of convergence of $g$, $\langle1\rangle_{K+i}=\int_{f^\star-\delta}^{f^\star}t^{K+i}\varrho_f(t)\,dt+\epsilon_{K+i}$, $0\le i<h$, the integral read from the germ, and $0\le\epsilon_k\le(f^\star-\delta)^{k-k'}\langle1\rangle_{k'}$ for every $k'\le k$ at which $\langle1\rangle_{k'}$ converges; for the origin $0$, $\langle1\rangle_{k'}\le C_+^{k'}V_1k'^{1-n}$ (Proposition~\ref{prop:sl-ends}(iii)).
\item \emph{The run back.} The recurrence, which holds at every positive integer (Proposition~\ref{prop:sl-density}), solved for $\langle1\rangle_k$ at the $k<K$ with $p_0(k)\neq0$, from $k=K-1$ down; at a positive integer zero of $p_0$, and below it, the values are read as moments of $\varrho_f$ continued along $(0,f^\star)$ (Theorem~\ref{thm:gaugeend}(v)).
\item \emph{Slopes.} The same with $(\log t)^j$ in the integrals and the differentiated recurrence (Proposition~\ref{prop:sl-slopes}).
\item \emph{Checks.} For the origin $0$, the low-end relation (Corollary~\ref{cor:rays-lowend}(ii)) read on the output returns $V_1$; two depths $K$ give the same values.
\end{enumerate}
\end{definition}

Every step but the truncations in step~(4) is an identity: the germ is exact over $\Bbbk(f^\star)$ up to $g(0)$, the error $\epsilon_k$ of the cut at $f^\star-\delta$ is bounded by step~(4), and the run back is exact arithmetic on exact coefficients. The truncation of the series of $g$ on $[f^\star-\delta,f^\star]$ is estimated, by the size of the last terms kept and by the agreement of two depths. The run back multiplies an error at $K$ by the growth of the other solutions of the recurrence from $K$ down to $1$, which is why the depth $K$ grows with the precision asked. The exercises read the centres of three and five marks this way (Exercises~\ref{ex:smallc-gauge} and~\ref{ex:cfive-gauge}), and a ray of a list on four marks whose symbol has a root on the segment $(0,f^\star)$ (Exercise~\ref{ex:lista-hslit}).

\section{The torus value}\label{sec:sl-torus}

% New: proved in the last part.
\begin{definition}[The torus value]\label{def:torusvalue}
Let the centre ray of a family at the strengths $\boldsymbol\sigma$ (Definition~\ref{def:centreray}) be given, with its stride $f=\prod_j\big(K_j(\alpha)K_j(\alpha^{-1})\big)^{-\sigma_j}$ and $f^\star=f(\one)$. On the torus $\mathbb T=\{\alpha:\ |\alpha_e|=1\ (e\in\cE),\ \alpha_{e_0}=1\}$, $\alpha_e=e^{i\vartheta_e}$, with the normalized measure $d\vartheta/(2\pi)^{n-1}$, the stride is $f=\prod_j|K_j(\alpha)|^{-2\sigma_j}$, a function with values in $[f^\star,\infty]$, equal to $\infty$ exactly where an active letter vanishes. The \emph{torus law} $\nu^{\mathbb T}_f$ is the image of $d\vartheta/(2\pi)^{n-1}$ under $f$, a probability measure on $[f^\star,\infty)$, $\varrho^{\mathbb T}_f$ its density where it has one, and the \emph{torus value} is
\[
\langle1\rangle^{\mathbb T}_k=\int t^k\,\nu^{\mathbb T}_f(dt)=\frac1{(2\pi)^{n-1}}\int_{\mathbb T}f^k\,d\vartheta\qquad(\re k\le0).
\]
When no letter $K_j$ of the family vanishes on $\mathbb T$, so that $\mathbb T$ lies in $X_F$, it is $(2\pi i)^{1-n}\langle1\rangle_{k;\mathbb T}$ in the notation of Definition~\ref{def:ray}, $\mathbb T$ oriented by $d\vartheta$.
\end{definition}

% New: proved in the last part.
\begin{proposition}[The torus is the second chain]\label{prop:sl-torus}
For the centre ray of a family at the strengths $\boldsymbol\sigma$, with the notation of Definition~\ref{def:torusvalue}, $\varrho_f$ and $g$ as in Proposition~\ref{prop:sl-ends}:
\begin{enumerate}[label=(\roman*)]
\item On $\mathbb T$, $f>f^\star$ off $\one$; the zero set of every active letter on $\mathbb T$ has measure zero. $\langle1\rangle^{\mathbb T}_k$ is continuous on $\re k\le0$ and holomorphic on $\re k<0$, and $\langle1\rangle^{\mathbb T}_{-j}=\CT(f^{-j})$ for $j\in\mathbb N$, the closing $\mathsf C_j(1)$ when every $\sigma_j\ge1$ (Definition~\ref{def:closings}).
\item Let every $\sigma_j\ge1$, and let $\mathbf p=(p_0,\dots,p_h)$ satisfy $\sum_ip_i(k)\langle1\rangle_{k+i}=0$ for all large integers $k$, as the law of $[1]$ does. Then $\sum_ip_i(k)\langle1\rangle^{\mathbb T}_{k+i}=0$ for every $k$ with $\re k\le-h$, and the distribution $\mathfrak D_{\mathbf p}\,\nu^{\mathbb T}_f$ on $(0,\infty)$ vanishes.
\item Near $f^\star$, $\nu^{\mathbb T}_f$ has a density, and with the one function $g$
\[
\varrho_f(t)=(f^\star-t)^{\frac{n-3}2}g(f^\star-t)\quad(t<f^\star),\qquad\varrho^{\mathbb T}_f(t)=(2\pi)^{1-n}\,(t-f^\star)^{\frac{n-3}2}g(f^\star-t)\quad(t>f^\star).
\]
\item For $\mathbf p$ as in (ii) with $f^\star$ a simple root of its symbol, the germ of $\varrho_f$ on the chart's side spans the line of Theorem~\ref{thm:gaugeend}(ii), the germ of $\varrho^{\mathbb T}_f$ extends by zero on the torus's side, and the two are the cut solutions $|f^\star-t|^{\frac{n-3}2}g(f^\star-t)$ of one germ on its two sides, the second with the factor $(2\pi)^{1-n}$.
\end{enumerate}
\end{proposition}

At a centre the native gauge of the direction lies on the positive orthant and on the unit torus, and the two meet only there (Theorem~\ref{thm:onekernel}). The torus value continues the closings in the strength, as the values continue them at the poles of the low end (Proposition~\ref{prop:principal}), and it satisfies every recurrence of the values. Its density is the density of the chart seen from the other side of the native value: the Morse function of the stride at $\one$ changes sign between the orthant and the torus, so the one function $g$ enters both, and the asymptotic statement at the end of Theorem~\ref{thm:onekernel} is the shadow of this exact one. Exercise~\ref{ex:smallc-torus} works the three-mark centre, where the torus value, the values and the solution at the sign are tied by an explicit connection.

% ---- New chapter: reading a ray (after ch:strideslaw) ----
% No new macros.
\chapter{Reading a ray}\label{ch:readray}

A ray is read from its list and its direction. A box of its relations gives a spanning set of a submodule that contains $[1]$, the one-step matrix on it with a certificate, and a recurrence of the values that holds at every strength (Section~\ref{sec:box}); the gauge end then gives the values themselves (Chapter~\ref{ch:strideslaw}). This chapter assembles the reading and says what each step certifies. At the low end a recurrence of the values meets the values at the positive integers through the poles at the nonpositive integers: the low-end relation at $0$ and, where the coefficients of the recurrence vanish to the order of the poles, second low-end relations at the other poles, with coefficients in the field of the recurrence, the weights and the volume of the direction. At the high end the generating function of the values is holomorphic off the reciprocals of $(0,f(\alpha^\star)]$, and its jump across them is the density of the stride's law. When every other root of the symbol lies outside the closed disc of radius $f(\alpha^\star)$, the values are, up to scale, the only solution of the recurrence whose exponential rate of growth is at most $f(\alpha^\star)$; the low-end relation fixes the scale, and solutions over the field of the recurrence approximate the values at the rate $f(\alpha^\star)/\lambda'$, $\lambda'$ the least modulus of the other roots. Positivity is the moment problem of the stride's law on $[0,f(\alpha^\star)]$. Against a solution carried by a segment that ends at a real point $\lambda$ off $[0,f(\alpha^\star)]$ its brackets shrink at the rate of the Green function of the plane cut along $[0,f(\alpha^\star)]$, read at $\lambda$; it cannot see a solution carried by $[0,f(\alpha^\star)]$ and dominated by the stride's law, which the analyticity of the density excludes; and for a recurrence of order two it gives the value at $k=1$ as the volume of the direction times a continued fraction in the coefficients. The procedure that reads a ray from its list runs the box by evaluation at strengths drawn at random, modulo primes for a list, and bounds the chance of a wrong output by degrees that the box fixes, or removes it with enough strengths. Its cost is set by the stencils of the relations on a lattice, by the degrees of the one-step matrix, which the number of pivots bounds, and by the order and the degree of the recurrence. The nondegenerate native gauges bound the sizes of a reading from below; over all exponents at once they form the graph of one map, through which they are found by continuation and certified one at a time.

\section{The low end and the slit}\label{sec:rd-lowend}

The low-end relation of a recurrence of the values is Corollary~\ref{cor:rays-lowend}(ii), with the low-end coefficients of Definition~\ref{def:rays-lowend}; by Proposition~\ref{prop:sl-density} it holds for every recurrence that the values satisfy at the large integers, whether or not it follows from the relations of the ray. The poles at the negative integers (Proposition~\ref{prop:rays-lowpoles}) give more relations.

% New: proved in the last part.
\begin{proposition}[The low-end relations at every pole]\label{prop:rd-second}
Let a ray have the origin $0$ and $\kappa_{1,S}>0$ for every nonempty proper $S$, with the volume $V_1$ of its direction (Proposition~\ref{prop:principal}), and let $p_0,\dots,p_h\in\mathbb C[k]$ satisfy $\sum_ip_i(k)\langle1\rangle_{k+i}=0$ for all large integers $k$. Let $j\in\mathbb N$, and let $m_0\in\{0,\dots,n-1\}$ be such that every $p_i$ with $i\le\min(h,j)$ vanishes to order at least $m_0$ at $k=-j$.
\begin{enumerate}[label=(\roman*)]
\item The function $\sum_{i\le\min(h,j)}p_i(k)\langle1\rangle_{k+i}$ is holomorphic at $k=-j$, with the value $-\sum_{i=j+1}^hp_i(-j)\langle1\rangle_{i-j}$ there.
\item If $m_0<n-1$, then $\sum_{i\le\min(h,j)}[(k+j)^{m_0}]p_i\cdot\CT\big(f^{i-j}\big)=0$.
\item If $m_0=n-1$, then
\[
V_1\sum_{i\le\min(h,j)}[(k+j)^{n-1}]p_i\cdot\CT\big(f^{i-j}\big)+\sum_{i=j+1}^hp_i(-j)\,\langle1\rangle_{i-j}=0 .
\]
At $j=0$ this is the low-end relation; at $j\ge1$ it is the \emph{second low-end relation} at $-j$, a linear relation among the values at the positive integers whose coefficients lie in the field generated by those of the $p_i$, the weights and $V_1$.
\end{enumerate}
\end{proposition}

At $m_0=0$, (ii) is the recurrence of the constant terms, $\sum_ip_i(-j)\CT(f^{i-j})=0$ (Corollary~\ref{cor:rays-lowend}(i)). At larger $m_0$ it ties the derivatives of order $m_0$ of $p_0,\dots,p_j$ at $-j$, and (i) ties the values at the positive integers to the Laurent coefficients of the values at the poles $-j,\dots,0$, those at $0$ being integrals over the faces of the sector cubes (Proposition~\ref{prop:rays-faces}).

The generating function of the values is read on the plane cut along the reciprocals of $(0,f(\alpha^\star)]$. The segment $(0,f(\alpha^\star))$ of the plane of native values is the \emph{slit}.

% New: proved in the last part.
\begin{proposition}[The values live on the slit plane]\label{prop:rd-slit}
Let a ray have a real origin and $\kappa_{1,S}>0$ for every nonempty proper $S$, with $\alpha^\star$, $f^\star=f(\alpha^\star)$, $\mu_f$ and $\varrho_f$ as in Proposition~\ref{prop:sl-ends}, and let $k_0$ be an integer at which $\langle1\rangle_{k_0}$ converges.
\begin{enumerate}[label=(\roman*)]
\item For $|q|<1/f^\star$, $G(q)=\sum_{k\ge k_0}\langle1\rangle_kq^{k-k_0}=\int t^{k_0}\mu_f(dt)/(1-qt)$, and the right side is holomorphic on $\mathbb C\setminus[1/f^\star,\infty)$: $G$ continues to the plane cut along the reciprocals of $(0,f^\star]$.
\item For $x>1/f^\star$, $G$ continues holomorphically across a neighbourhood of $x$ in $(1/f^\star,\infty)$ from either side, with the limits $G(x\pm i0)$, and $G(x+i0)-G(x-i0)=2\pi i\,x^{-1-k_0}\,\varrho_f(1/x)$: the jump of $G$ across the cut is the density of the stride's law.
\item On the centre ray of a family, every native gauge on the torus $\mathbb T$ of Definition~\ref{def:torusvalue} has its native value in $[f^\star,\infty)$, off the slit.
\item Every conjugate of $f^\star$ over the field of the weights is the native value of a nondegenerate native gauge (Proposition~\ref{prop:rays-conjugates}), and a root of the symbol of every recurrence that the values satisfy at the large integers with coefficients in that field; the real conjugates in $(0,f^\star)$ lie on the slit.
\end{enumerate}
\end{proposition}

A root $\lambda$ of a symbol on the slit sits on the cut, at $1/\lambda\in(1/f^\star,\infty)$, where the generating function of the values continues from either side and jumps by the density of the stride's law; a root off $[0,f^\star]$ sits where it is holomorphic. Not every root of a symbol is a native value: whenever the values satisfy $\sum_ip_i(k)\langle1\rangle_{k+i}=0$, they satisfy the recurrence with the coefficients $p_{i-1}(k+1)-\lambda_0\,p_i(k)$, $0\le i\le h+1$, $p_{-1}=p_{h+1}=0$, which follows from the relations when the first does and has the symbol $(\lambda-\lambda_0)\chi(\lambda)$, for every $\lambda_0$. In the other direction, the native value of every nondegenerate native gauge is a root of the symbol of every recurrence that follows from the relations (Corollary~\ref{cor:symbol}), and that of every native gauge when the recurrence has a certificate as in Proposition~\ref{prop:rays-certtop}. Exercises~\ref{ex:lista-hcentre} and~\ref{ex:lista-hslit} work second low-end relations on a list of four marks.

\section{Positive selection}\label{sec:rd-selection}

% New: proved in the last part.
\begin{theorem}[Positive selection]\label{thm:positiveselection}
Let a ray have a real origin and $\kappa_{1,S}>0$ for every nonempty proper $S$, with $f^\star=f(\alpha^\star)$, and let $p_0,\dots,p_h\in\mathbb C[k]$, $h\ge1$, satisfy $\sum_ip_i(k)\langle1\rangle_{k+i}=0$ for all large integers $k$, so that the recurrence holds at every integer at which the values converge (Proposition~\ref{prop:sl-density}). Let $p_0$ and $p_h$ have the largest degree $D$, so that the symbol $\chi$ has degree $h$ and $\chi(0)\neq0$; then $f^\star$ is a root of $\chi$, and let the roots of $\chi$ be distinct, $\lambda_1=f^\star,\lambda_2,\dots,\lambda_h$. Let $k_1$ be an integer at which the values converge, with $p_0(k)p_h(k)\neq0$ for every integer $k\ge k_1$. A \emph{solution} is a sequence $(x(k))_{k\ge k_1}$ with $\sum_ip_i(k)x(k+i)=0$ for every $k\ge k_1$; the solutions form a space $\mathcal V$ of dimension $h$, and the values are one of them.
\begin{enumerate}[label=(\roman*)]
\item For every solution, $X(q)=\sum_{k\ge k_1}x(k)q^{k-k_1}$ satisfies $L_qX=R_x$ as a formal power series, with
\[
L_q=\sum_{i=0}^hq^{h-i}\,p_i\Big(q\frac d{dq}+k_1-i\Big)
\]
and $R_x$ a polynomial of degree less than $h$ that depends linearly on $x(k_1),\dots,x(k_1+h-1)$. The coefficient of $(d/dq)^D$ in $L_q$ is $q^{D+h}\chi(1/q)$, so the singular points of $L_q$ in $\mathbb C^\times$ are the points $1/\lambda_c$. For the values, $X$ continues holomorphically to $\mathbb C\setminus[1/f^\star,\infty)$ (Proposition~\ref{prop:rd-slit}).
\item The solutions with $x(k)=O(r^k)$ for every $r>f^\star$ form a subspace of $\mathcal V$ whose dimension is the number of roots in the closed disc $|\lambda|\le f^\star$; the values lie in it, $\langle1\rangle_k=C\,(f^\star)^kk^{-\frac{n-1}2}\big(1+o(1)\big)$ with $C>0$, and the exponent $-\chi_1(f^\star)/(f^\star\chi'(f^\star))$ of the root $f^\star$ is $-\frac{n-1}2$.
\item Let every root other than $f^\star$ have modulus greater than $f^\star$, and put $\lambda'=\min_{c\ge2}|\lambda_c|$. Then the solutions with $x(k)=O(r^k)$ for some $r<\lambda'$ are the multiples of the values, and so are those whose generating function converges on $|q|<1/f^\star$: the values are, up to a constant factor, the \emph{solution of least growth}. Let moreover the origin be $0$, $k_1=1$ and $a=[k^{n-1}]p_0\neq0$, and put $\ell(x)=-a^{-1}\sum_{i=1}^hp_i(0)\,x(i)$. Then $\ell(\langle1\rangle)=V_1$, and $\langle1\rangle_k=V_1\,x(k)/\ell(x)$ for every nonzero solution $x$ of least growth: the values divided by the volume of the direction are a function of the recurrence alone.
\item Under the hypotheses of (iii), let $x^{(0)}\in\mathcal V$ have $\ell(x^{(0)})=V_1$ and let $x^{(1)},\dots,x^{(h-1)}$ be a basis of $\{\ell=0\}$, so that $\langle1\rangle=x^{(0)}+\sum_{i\ge1}M_ix^{(i)}$ for one $M\in\mathbb C^{h-1}$. For all large $K$ the equations $x^{(0)}(K+l)+\sum_{i\ge1}M^{(K)}_ix^{(i)}(K+l)=0$, $0\le l\le h-2$, have exactly one solution $M^{(K)}$, and $|M^{(K)}-M|=O\big((f^\star/\lambda')^KK^{C'}\big)$ for some $C'$. When the $p_i$ have their coefficients in a field $\Bbbk$ that contains $V_1$, the $x^{(i)}$ can be taken with values in $\Bbbk$, and then every $M^{(K)}$ lies in $\Bbbk^{h-1}$.
\end{enumerate}
\end{theorem}

The selection runs on growth alone: among the solutions, the values are the one that grows like the native value of the direction, and the other roots of the symbol, outside the closed disc of radius $f^\star$, give solutions that grow faster and are cut away. The hypothesis of (iii) is what makes the values the solution of least growth: a root in the closed disc, in particular one on the slit, gives a solution whose exponential rate of growth is at most that of the values, and growth cannot remove it (Section~\ref{sec:rd-positivity}). The low-end relation then fixes the scale through the volume of the direction, so the values over that volume are a function of the recurrence alone.

\section{Positivity as a moment problem}\label{sec:rd-positivity}

% New: proved in the last part.
\begin{proposition}[Positivity is the moment problem of the stride's law]\label{prop:rd-moment}
Let a ray have a real origin and $\kappa_{1,S}>0$ for every nonempty proper $S$, with $f^\star=f(\alpha^\star)$ and $\mu_f$ as in Proposition~\ref{prop:sl-ends}. For a real polynomial $p=\sum_mp_mt^m$, a real $K$ and a sequence $x$ defined at $K+\mathbb N$, put $L_{K,p}(x)=\sum_mp_m\,x(K+m)$.
\begin{enumerate}[label=(\roman*)]
\item If $\langle1\rangle_K$ converges, $p\ge0$ on $[0,f^\star]$ and $p\neq0$, then $L_{K,p}(\langle1\rangle)=\int t^Kp(t)\,\mu_f(dt)>0$.
\item Every real polynomial of degree at most $N$ that is nonnegative on $[0,f^\star]$ is a sum of polynomials $q^2$, $t\,q^2$, $(f^\star-t)\,q^2$ and $t(f^\star-t)\,q^2$ of degree at most $N$, with real $q$. So for a real sequence $x$, $L_{K,p}(x)\ge0$ for every such $p$ exactly when the Hankel matrices
\begin{gather*}
\big[x(K+i+i')\big],\qquad\big[x(K+1+i+i')\big],\\
\big[f^\star x(K+i+i')-x(K+1+i+i')\big],\qquad\big[f^\star x(K+1+i+i')-x(K+2+i+i')\big],
\end{gather*}
with $0\le i,i'<\lfloor\frac N2\rfloor+1$, $\lfloor\frac{N+1}2\rfloor$, $\lfloor\frac{N+1}2\rfloor$ and $\lfloor\frac N2\rfloor$ respectively, are positive semidefinite.
\item A real sequence $(x(K+m))_{m\in\mathbb N}$ has $L_{K,p}(x)\ge0$ for every real $p\ge0$ on $[0,f^\star]$ exactly when $x(K+m)=\int t^m\nu(dt)$ for every $m$, for a positive measure $\nu$ on $[0,f^\star]$, which is then unique; for the values, $\nu=t^K\mu_f$.
\end{enumerate}
\end{proposition}

So the positivity of the values is the statement that they are the moments of a positive measure on $[0,f^\star]$, and Theorem~\ref{thm:pencil}(iii) is its first family, the Hankel determinants of the values. For a family $x^{(0)}+\sum_iM_ix^{(i)}$ of solutions of a recurrence, as in Theorem~\ref{thm:positiveselection}(iv), each condition $L_{K,p}\ge0$ is one linear inequality in the coordinates $M$, and the conditions of degree at most $N$ at the strength $K$ cut out a closed convex set, the \emph{bracket} of degree $N$, which contains the coordinates of the values.

% New: proved in the last part.
\begin{theorem}[The rate of positivity off the segment]\label{thm:rd-rate}
Let a ray be as in Proposition~\ref{prop:rd-moment}, $K$ an integer at which $\langle1\rangle_K$ converges, $\lambda\in\mathbb R\setminus[0,f^\star]$ and $I_\lambda$ the closed segment between $0$ and $\lambda$. Let $u(k)=\int_{I_\lambda}t^k\,\nu(t)\,dt$ for the integers $k\ge K$, where $\nu$ is a real function on $I_\lambda$ with $\int_{I_\lambda}|t|^K|\nu(t)|\,dt<\infty$ and $\nu(t)=c\,|\lambda-t|^a\big(1+o(1)\big)$ as $t\to\lambda$ in $I_\lambda$, with $c\neq0$ and $a>-1$. Let $E_N$ be the set of the real $\varepsilon$ with $L_{K,p}(\langle1\rangle+\varepsilon u)\ge0$ for every real $p\ge0$ on $[0,f^\star]$ of degree at most $N$. Then $E_N$ is a closed interval that contains $0$, and there is a constant $C$ with its length at most $C\,N^{a+1}e^{-Ng(\lambda)}$ for all large $N$, where
\[
g(\lambda)=\log\big|w+\sqrt{w^2-1}\big|>0,\qquad w=\frac{2\lambda}{f^\star}-1,
\]
the branch of the square root with $|w+\sqrt{w^2-1}|>1$. For $\lambda>f^\star$, with $\rho=f^\star/\lambda$,
\[
e^{-g(\lambda)}=\frac{1-\sqrt{1-\rho}}{1+\sqrt{1-\rho}}=\frac\rho4+O(\rho^2).
\]
\end{theorem}

The rate is read off the Green function $g$ of the plane cut along $[0,f^\star]$, with its pole at infinity, at the end $\lambda$ of the segment that carries $u$: the farther $\lambda$ lies from $[0,f^\star]$, the faster positivity separates the values from the values plus a multiple of $u$. For a recurrence of order two, with the coordinate $M$ of Theorem~\ref{thm:positiveselection}(iv) and $u=x^{(1)}$ of this form, $E_N$ is the bracket of degree $N$ of $M$ less the value of $M$.

% New: proved in the last part.
\begin{proposition}[What positivity cannot see]\label{prop:rd-blind}
Let a ray be as in Proposition~\ref{prop:rd-moment}, $K$ a real at which $\langle1\rangle_K$ converges, and $\nu$ a real signed measure on $[0,f^\star]$ with $|\nu|\le C\mu_f$ for a constant $C$; put $u(k)=\int t^k\nu(dt)$ for $k\ge K$.
\begin{enumerate}[label=(\roman*)]
\item For $|\varepsilon|\le1/C$, $\langle1\rangle+\varepsilon u$ is from $K$ on the moment sequence of the positive measure $\mu_f+\varepsilon\nu$, and $L_{K',p}(\langle1\rangle+\varepsilon u)\ge0$ for every real $K'\ge K$ and every real $p\ge0$ on $[0,f^\star]$.
\item If on an open interval $J\subseteq(0,f^\star)$ the measure $\nu$ has a density that agrees almost everywhere with no real-analytic function on $J$, then for $\varepsilon\neq0$ the measure $t^K(\mu_f+\varepsilon\nu)$, the only finite measure on $[0,f^\star]$ whose moments are the $\langle1\rangle_{K+m}+\varepsilon u(K+m)$, $m\in\mathbb N$, has no density on $(0,f^\star)$ that is real-analytic, while $t^K\mu_f$ has one (Proposition~\ref{prop:sl-ends}(i)).
\item Let the origin be $0$, let $\mathbf p=(p_0,\dots,p_h)$ satisfy the conditions of Proposition~\ref{prop:sl-density}, and let $\nu$ have a density that is bounded near $0$, with $\mathfrak D_{\mathbf p}\nu=0$ as distributions on $(0,\infty)$. Then $u(k)=\int t^k\nu(dt)$ converges at every $k$ with $\re k>-1$, $\sum_ip_i(k)u(k+i)=0$ there, and $\sum_{i=1}^hp_i(0)\,u(i)=0$: $u$ satisfies the recurrence and contributes nothing to the low-end relation.
\end{enumerate}
\end{proposition}

So a sequence $u$ as in (i) and (iii), a solution of the recurrence carried by $[0,f^\star]$ and dominated by the stride's law, is invisible to the recurrence, to the low-end relation and to positivity together: read as conditions on candidates, they leave the direction of $u$ free. For a root $\lambda$ of the symbol on the slit, a density $\nu$ on $(0,\lambda)$ with $\mathfrak D_{\mathbf p}\big(\nu\,\one_{(0,\lambda)}\big)=0$ as distributions on $(0,\infty)$, a solution of the density equation on $(0,\lambda)$ that extends by zero at $\lambda$ (Definition~\ref{def:sl-densityop}), bounded near $0$ and with $|\nu|\le C\varrho_f$ on $(0,\lambda)$, gives such a $u$. The gauge end does not choose among candidates: it constructs the density, and by (ii) the analyticity of the density on the open support is as much a property of the stride's law as its positivity.

% New: proved in the last part.
\begin{proposition}[A recurrence of order two is the volume of the direction times a continued fraction]\label{prop:rd-onemaster}
Let a ray have the origin $0$ and $\kappa_{1,S}>0$ for every nonempty proper $S$, with $f^\star=f(\alpha^\star)$, and let $p_0,p_1,p_2\in\mathbb R[k]$ satisfy $p_0(k)\langle1\rangle_k+p_1(k)\langle1\rangle_{k+1}+p_2(k)\langle1\rangle_{k+2}=0$ for all large integers $k$; put $a=[k^{n-1}]p_0$ and $\rho_k=\langle1\rangle_{k+1}/\langle1\rangle_k$ for the integers $k\ge1$.
\begin{enumerate}[label=(\roman*)]
\item $0<\rho_k<f^\star$, and $\rho_k\big(p_1(k)+p_2(k)\rho_{k+1}\big)=-p_0(k)$ for every integer $k\ge1$.
\item $\langle1\rangle_1\big(p_1(0)+p_2(0)\rho_1\big)=-a\,V_1$. When $-a/p_1(0)$ and $-a/(p_1(0)+f^\star p_2(0))$ are finite and positive, $\langle1\rangle_1/V_1$ lies between them, strictly when they differ.
\item Let the hypotheses of Theorem~\ref{thm:positiveselection}(iii) hold with $h=2$ and $k_1=1$, so that $\lambda'=|\lambda_2|$, and let $a\neq0$; put $\mathsf M_k(r)=-p_0(k)/\big(p_1(k)+p_2(k)\,r\big)$. There are $K_0$, $C$ and $C'$ such that for every $K\ge K_0$ and every $r\in[0,f^\star]$ the composite $\mathsf M_1\circ\cdots\circ\mathsf M_K(r)$ is defined and lies within $C(f^\star/\lambda')^KK^{C'}$ of $\rho_1$. So $\rho_1$ is the value of a continued fraction in the coefficients of the recurrence, and $\langle1\rangle_1=-aV_1/(p_1(0)+p_2(0)\rho_1)$. The sets of the $\mathsf M_1\circ\cdots\circ\mathsf M_K(r)$ over the $r\in(0,f^\star)$ whose intermediate values $\mathsf M_k\circ\cdots\circ\mathsf M_K(r)$, $2\le k\le K$, lie in $(0,f^\star)$ are nested in $K$ and contain $\rho_1$.
\end{enumerate}
\end{proposition}

The low-end relation fixes the value at $k=1$ through the volume of the direction once the ratio $\rho_1$ is known, and positivity alone places every $\rho_k$ in $(0,f^\star)$. The bracket of (ii) is the image of $\rho_1\in(0,f^\star)$ under $\rho\mapsto-a/(p_1(0)+p_2(0)\rho)$; the nested sets of (iii) sharpen that condition with the conditions $\langle1\rangle_k>0$ and $f^\star\langle1\rangle_k>\langle1\rangle_{k+1}$ at every $k\le K+1$, and they shrink to $\rho_1$ at the rate $(f^\star/\lambda')^K$, up to a power of $K$. The exercises work the brackets and the rate of Theorem~\ref{thm:rd-rate} on the centre of three marks (Exercises~\ref{ex:smallc-bracket} and~\ref{ex:cthree-positivity}), and the continued fraction on $W_3$ at $(\one;0,2)$ (Exercise~\ref{ex:wthree-bracket}).

\section{The procedure}\label{sec:rd-procedure}

% New: proved in the last part.
\begin{definition}[The procedure that reads a ray from its list]\label{def:readprocedure}
Input: a family whose weights lie in a field $\Bbbk\subseteq\mathbb R$, $\mathbb Q$ for a list; a ray of it with $\kappa_{1,S}>0$ for every nonempty proper $S$ and its origin in $\Bbbk$, the origin $0$ unless stated; the precision wanted. Output: a spanning set and its one-step matrix with a certificate, a recurrence of the values that holds at every strength, the values at the positive integers, and brackets for them.
\begin{enumerate}[label=(\arabic*)]
\item \emph{Walls and symmetries.} The walls $\kappa_{1,S}$ and the volume $V_1=\sum_\pi\prod_{l=2}^n\kappa_{1,U_l}^{-1}$ of the direction (Proposition~\ref{prop:principal}); a finite group $G$ generated by maps of the four kinds of Proposition~\ref{prop:raysym}(ii) that are symmetries of the ray.
\item \emph{The box.} Through the boxes $\mathfrak B_{\mathfrak t}(\bar{\mathbf a},\bar{\mathbf b})$ with level bounds kept by $G$, each with a ranking, and the $\mathfrak s<\mathfrak t$, in an order that reaches each, the first box accepted at some $\mathfrak s$, with its spanning set and its stems (Definition~\ref{def:box-reduction}); the search stops (Theorem~\ref{thm:rayboxstop}). Acceptance is decided by eliminations of the relations at strengths drawn at random, for a list modulo a prime (Proposition~\ref{prop:rd-certain}(iii) and~(v), below).
\item \emph{The one-step matrix.} At strengths drawn at random, for a list modulo primes, the reductions of the targets, from the relations they need (Proposition~\ref{prop:rd-certain}(vi)); the entries of $\mathbf B$ over $\Bbbk(k)$ by rational reconstruction from these values (Proposition~\ref{prop:rd-certain}(ii)), for a list with Chinese remaindering and rational reconstruction of the coefficients over the primes; a check against eliminations of the whole box at strengths and primes drawn afresh (Proposition~\ref{prop:rd-certain}(iv) and~(v)). A certificate supported on the stems exists (Theorem~\ref{thm:raybox}(ii)), and every solution, the values among them, satisfies $\varsigma^{-1}v_V=\mathbf Bv_V$ (Theorem~\ref{thm:raybox}(iii)).
\item \emph{The recurrence.} The first dependence of the rows $w_i(k)=w_{i-1}(k-1)\,\mathbf B(k)$ (Proposition~\ref{prop:rays-lawrows}(ii)), read at points and reconstructed as in step~(3), its coefficients being proportional to signed maximal minors (Proposition~\ref{prop:rays-lawrows}(iii)): $\sum_jp_j(k)f^j\in\mathcal L$, the law of $[1]$ when the spanning set is independent (Theorem~\ref{thm:raybox}(iv)), as it is when the values of the spanning set at the nondegenerate native gauges span (Proposition~\ref{prop:rays-spectrum}(i)). Its end coefficient $p_0$ vanishes only at the shifted zeros of the denominator of $\mathbf B$ and at the strengths where the rows line up (Proposition~\ref{prop:rays-lawrows}(iii)).
\item \emph{The values.} Either the gauge end, from the recurrence, $\alpha^\star$ and $\det H$, when $f(\alpha^\star)$ is a simple root of its symbol (Definition~\ref{def:sl-gaugeread}); or, when the hypotheses of Theorem~\ref{thm:positiveselection}(iii) hold, positive selection with the low-end relation (Theorem~\ref{thm:positiveselection}(iv)), with brackets from positivity (Proposition~\ref{prop:rd-moment}).
\item \emph{Checks, outside the construction.} When the origin is integral and every letter exponent of the direction is at least $1$, the closings of the spanning set satisfy $\mathsf C_{j+1}(\mathbf b)=\mathbf B(-j)\,\mathsf C_j(\mathbf b)$ for all large $j$ (Theorem~\ref{thm:closings}, Theorem~\ref{thm:raybox}(iii)); for the origin $0$, the constant terms of the $f^{-j}$ satisfy the reflected recurrence (Corollary~\ref{cor:rays-lowend}(i)); when $\mathbf B$ is regular at $\infty$, the native values of nondegenerate native gauges found by other means are reciprocal eigenvalues of $\mathbf B(\infty)$ (Proposition~\ref{prop:rays-spectrum}(ii)); the exponent at $f(\alpha^\star)$ is $\frac{n-3}2$ (Theorem~\ref{thm:gaugeend}(iii)); the low-end relation returns $V_1$ from the values read; two readings from different depths agree.
\end{enumerate}
\end{definition}

Every relation of a box is a relation of the module, so every identity that the elimination gives over $\Bbbk(k)$ holds at every strength (Proposition~\ref{prop:box-relations}, Theorem~\ref{thm:raybox}), and the recurrence of step~(4) holds for the values at every strength at which they converge (Proposition~\ref{prop:sl-density}). Chance enters only through the strengths and primes at which steps (2) to (4) evaluate: Proposition~\ref{prop:rd-certain} bounds it by degrees that the box fixes, and enough strengths remove it. The procedure needs no bound on the box: the search stops (Theorem~\ref{thm:rayboxstop}).

The procedure of Section~\ref{sec:machineproc} reads the same module from its solutions, with no box reduced. Its output is certified when the solutions supplied annihilate no submodule (Corollary~\ref{cor:glue}), and the closings alone certify it exactly when the native gauge of the direction sees every nonzero submodule (Corollary~\ref{cor:seeall}); otherwise a unique solution of its equations is evidence. The procedure above certifies the matrix of its spanning set by construction: the certificate of the box is an identity of the relations, whatever any solution sees, and the matrix is the one-step matrix of the module exactly when the spanning set is independent. The two check each other: every solution that Section~\ref{sec:machineproc} uses satisfies $\varsigma^{-1}v_V=\mathbf Bv_V$ on the spanning set of the box (Theorem~\ref{thm:raybox}(iii)), which is step~(6) for the closings, and where the supplied solutions annihilate a submodule, the matrix of the box is still certified by its certificate. Section~\ref{sec:machineproc} ends with the law of $[1]$; the procedure above goes on to the values, at the gauge end or by positive selection. A recurrence read off the closings is a recurrence of the values when it holds for all large $j$ and every letter exponent of the direction is positive (Corollary~\ref{cor:valuemodule}); finitely many closings do not show that, while the recurrence of step~(4) carries a certificate.

% New: proved in the last part.
\begin{proposition}[The relations of a box are difference equations]\label{prop:rd-lattice}
Let a box $\mathfrak B=\mathfrak B_{\mathfrak t}(\bar{\mathbf a},\bar{\mathbf b})$ of a ray be given (Definition~\ref{def:raybox}).
\begin{enumerate}[label=(\roman*)]
\item The symbols at a level $(\mathbf a,\mathbf b)$ are the points of the affine lattice $\{u\in\mathbb Z^{\cE}:\ u_{\cE}=\sum_j(a_jr_j+b_jr^\ast_j)\}$, of rank $n-1$. There is a constant $C$, depending on the ray and the level, such that the coordinates $u_e$, $e\neq e_0$, of the symbols of wall degree at most $\mathfrak t$ at that level lie in intervals of length at most $C(1+\mathfrak t)$; so they number at most $\big(C(1+\mathfrak t)+1\big)^{n-1}$.
\item A letter identity of table $j$ ties a symbol to the $|\cB_j|$ symbols one unit of that letter up, with weights as coefficients; the mark relation of a symbol at a mark ties it to at most $\sum_j|\cB_j|$ symbols one unit of a letter up, with coefficients affine in $k$. The shifts of $(u;\mathbf a,\mathbf b)$ that occur lie in one finite set that does not depend on the symbol.
\item In the basis of the symbols the relations of the box form a matrix $A(k)=A_0+kA_1$ with entries in the field of the weights and the origin; $A_1$ vanishes outside the rows of the mark relations, where its entries are $m_{1,e}$, $-s_{1j}\theta_{j,T}$ and $-t_{1j}\theta_{j,T}$. The same holds for the orbit averages of the relations modulo a finite group of symmetries that keeps the box, in the basis of the orbit averages of the symbols.
\end{enumerate}
\end{proposition}

% New: proved in the last part.
\begin{proposition}[What a reading by evaluation certifies]\label{prop:rd-certain}
Let a finite group $G$ of symmetries keep a box $\mathfrak B$ of a ray, with the field $\Bbbk$ of Definition~\ref{def:raybox} and a ranking (Definition~\ref{def:box-reduction}). Let $A(k)$ be the matrix whose rows are the orbit averages of the letter identities and of the mark relations that span $\mathcal R_{\mathfrak B}$, in the basis of the orbit averages ordered by the ranking, so that its rows span $\mathcal R^G_{\mathfrak B}$ and its entries are affine in $k$ (Proposition~\ref{prop:rd-lattice}(iii)). Let $P$ be the pivot orbits and $F$ the free orbits, $r=|P|$, and let $A_P(k)$ and $A_F(k)$ be $r$ rows of $A$ that are independent over $\Bbbk(k)$, read on the columns $P$ and $F$.
\begin{enumerate}[label=(\roman*)]
\item $r$ is the rank of $A$, $A_P(k)$ is invertible over $\Bbbk(k)$, and the reduction of $x\in\mathcal F^G_{\mathfrak B}$ with coefficients in $\Bbbk$ is $x_F-x_PA_P(k)^{-1}A_F(k)$; $\det A_P$ is a nonzero polynomial of degree at most $r$. At $k_0\in\Bbbk$ the elimination of all the rows of $A(k_0)$ in the order of the ranking has the pivots $P$ exactly when some minor of order $r$ of $A(k_0)$ on the columns $P$ is nonzero, as when $\det A_P(k_0)\neq0$; then the reduction it gives of every such $x$ is the value at $k_0$ of the reduction over $\Bbbk(k)$.
\item When the box is accepted, every entry of its matrix $\mathbf B$ is a polynomial of degree at most $r$ divided by $\det A_P$. A rational function whose numerator and denominator have degrees at most $d_1$ and $d_2$ is fixed, among such functions, by its values at any $d_1+d_2+1$ points at which the denominators do not vanish; with a common denominator known, a numerator of degree at most $d_1$ is fixed by $d_1+1$ values.
\item The verdict of the elimination at $k_0$ on acceptance at $\mathfrak s$, whether the reductions of the targets are combinations of the free Laurent orbits of wall degree at most $\mathfrak s$, is the verdict over $\Bbbk(k)$ except at no more than $2r$ strengths $k_0$.
\item Let the box be accepted, and let $\mathbf B'$ be a matrix over $\Bbbk(k)$ whose entries have a common denominator $Q'$, with $Q'$ and the numerators of degree at most $d$. If $\mathbf B'\neq\mathbf B$, then $\mathbf B'(k_0)$ agrees with the matrix that the elimination at $k_0$ gives at no more than $r+d$ strengths $k_0$ with $Q'(k_0)\neq0$ at which the elimination has the pivots $P$. So agreement at $r+d+1$ such strengths proves $\mathbf B'=\mathbf B$, and a wrong $\mathbf B'$ passes the check at a strength drawn uniformly from a finite set $S\subseteq\Bbbk$ with probability at most $(r+d)/|S|$.
\item Let $\Bbbk=\mathbb Q$, and let $q$ be a prime that divides no denominator of an entry of $A$, with $\det A_P$ not identically $0$ modulo $q$. Then $A$ modulo $q$ has the rank $r$ and the pivots $P$ over $\mathbb F_q(k)$, its reductions are those over $\mathbb Q(k)$ taken modulo $q$, and (i) to (iv) hold over $\mathbb F_q$ with strengths in $\mathbb F_q$. A $\mathbf B'$ over $\mathbb Q(k)$ with $\mathbf B'\neq\mathbf B$ passes the check modulo $q$ at a strength drawn uniformly from $\mathbb F_q$ with probability at most $(r+d)/q$, for every prime $q$ outside a finite set fixed by $A$ and $\mathbf B'$.
\item If a subspace $\mathcal R'\subseteq\mathcal R^G_{\mathfrak B}$ contains $x-\bar x$ for every target $x$, $\bar x$ its reduction, then the reduction of every target modulo $\mathcal R'$, defined as in Definition~\ref{def:box-reduction}(ii) with $\mathcal R'$ in place of $\mathcal R^G_{\mathfrak B}$, is $\bar x$.
\end{enumerate}
\end{proposition}

An elimination over $\Bbbk(k)$ carries polynomials whose degrees are bounded only by the number $r$ of pivots, which grows with the box. By (ii) a reconstruction from evaluations needs only as many strengths as the degrees of the entries of $\mathbf B$ ask, and by (iv) a check at one fresh strength drawn from a set $S$ bounds the chance of a wrong result by $(r+d)/|S|$, while $r+d+1$ checks remove it.

A ray with a real origin is read in the same way: the box carries the origin in its mark relations (Definition~\ref{def:raybox}(iii)), and the stride's law, its density and the gauge end are those of Chapter~\ref{ch:strideslaw}, which allow a real origin, the constant $g(0)$ carrying the factor $(\omega_0/dy)(\alpha^\star)$ (Proposition~\ref{prop:sl-ends}(ii)); the low-end relation and Theorem~\ref{thm:positiveselection}(iii) and~(iv) ask for the origin $0$. Data with rational masses are read through an origin.

% New: proved in the last part.
\begin{proposition}[Data with rational masses]\label{prop:rd-rational}
Let data $(\mvec;\mathbf s,\mathbf t)$ of a family on the balance wall have $\mvec\in\mathbb Q^{\cE}$, $\mathbf s,\mathbf t\in\mathbb N^N$ and every wall $\kappa_U$, $U$ nonempty and proper, positive, and put $f_{\mvec}=\alpha^{\mvec}\prod_jK_j^{-s_j}\Psi_j^{-t_j}$ on the positive orthant.
\begin{enumerate}[label=(\roman*)]
\item Let $\mvec'\in\mathbb Z^{\cE}$ put $(\mvec';\mathbf s,\mathbf t)$ on the balance wall with every wall positive. Then the ray with the origin $(\mvec-\mvec';\mathbf 0,\mathbf 0)$ and the direction $(\mvec';\mathbf s,\mathbf t)$ has positive walls at every real $k\ge1$, its values converge there, and $\langle1\rangle_1=P_F(\mvec;\mathbf s,\mathbf t)$.
\item Let $N\ge1$ be an integer with $N\mvec\in\mathbb Z^{\cE}$. Then the ray with the origin $0$ and the direction $N(\mvec;\mathbf s,\mathbf t)$ has positive walls, its stride is $f_{\mvec}^N$, $\langle1\rangle_{1/N}=P_F(\mvec;\mathbf s,\mathbf t)$, and the wall degree of a symbol $g$ for it is $\frac1N\max_U\big(-\kappa_U(g)/\kappa_U\big)$, the maximum over the nonempty proper $U$.
\end{enumerate}
\end{proposition}

So data with rational masses are read at $k=1$ on a ray with an origin, by the gauge end, or at $k=\frac1N$ on a ray from the origin $0$, whose box of wall degree $\mathfrak t$ holds the symbols of wall degree $N\mathfrak t$ read on the walls of the data.

\section{What a reading costs}\label{sec:rd-cost}

A box has at most $\big(C(1+\mathfrak t)+1\big)^{n-1}$ symbols at each level (Proposition~\ref{prop:rd-lattice}(i)), and modulo a group $G$ of symmetries its orbits number at least its symbols divided by $|G|$; one evaluation in step~(3) of Definition~\ref{def:readprocedure} is one elimination of its relations, and the number of evaluations is set by the degrees of $\mathbf B$ (Proposition~\ref{prop:rd-certain}(ii)). The relations are affine in the strength, so one factorization at one strength gives every Taylor coefficient there.

% New: proved in the last part.
\begin{proposition}[One factorization serves every coefficient]\label{prop:rd-lift}
In the setting of Proposition~\ref{prop:rd-certain}, let $k_0\in\Bbbk$ have $\det A_P(k_0)\neq0$, and put $M=A_P(k_0)$ and $E=A_P(k_0+1)-A_P(k_0)$, the part of $A_1$ on the rows and columns of $A_P$, so that $A_P(k_0+\epsilon)=M+\epsilon E$. For every row $x_P$ over $\Bbbk$, as power series in $\epsilon$ convergent for small $|\epsilon|$,
\[
x_P\,A_P(k_0+\epsilon)^{-1}=\sum_{j\ge0}\lambda_j\epsilon^j,\qquad\lambda_0=x_PM^{-1},\quad\lambda_j=-\lambda_{j-1}EM^{-1}.
\]
So the Taylor coefficients at $k_0$ of the reductions, and of the entries of $\mathbf B$, up to any order $J$ come from one factorization of $M$ and $J$ products by $EM^{-1}$; and a rational function whose numerator and denominator have degrees at most $d_1$ and $d_2$, the denominator not vanishing at $k_0$, is fixed, among such functions, by its Taylor coefficients of order at most $d_1+d_2$ at $k_0$.
\end{proposition}

So one factorization at one strength can replace many evaluations; each step of the series is one solve with the factorization of $M$, and what it saves depends on the fill of that factorization.

% New: proved in the last part.
\begin{proposition}[The densities of a spanning set]\label{prop:rd-masterdensities}
Let a ray have a real origin and $\kappa_{1,S}>0$ for every nonempty proper $S$, with $\alpha^\star$, $f^\star=f(\alpha^\star)$ and the function $g$ of Proposition~\ref{prop:sl-ends}, and let $b_1,\dots,b_{\mathfrak m}$ be Laurent insertions with real coefficients and $Q$, $\tilde P_{il}$ polynomials with $Q(k)\langle b_i\rangle_{k-1}=\sum_l\tilde P_{il}(k)\langle b_l\rangle_k$ as meromorphic functions, as a box gives with $\mathbf B=\tilde P/Q$ (Theorem~\ref{thm:raybox}(iii)). For a Laurent insertion $b$ with real coefficients let $\mu_{f,b}$ be the image under $f$ of the signed measure $b\,\omega_0$ on the positive part of the chart, a signed measure on $(0,f^\star]$, finite on $[t,f^\star]$ for $t>0$.
\begin{enumerate}[label=(\roman*)]
\item $\langle b\rangle_k=\int t^k\mu_{f,b}(dt)$ wherever $\langle b\rangle_k$ converges absolutely; $\mu_{f,b}$ has a density $\varrho_{f,b}$, real-analytic on $(0,f^\star)$, and near $f^\star$, $\varrho_{f,b}(t)=(f^\star-t)^{\frac{n-3}2}g_b(f^\star-t)$ with $g_b$ holomorphic at $0$ and $g_b(0)=b(\alpha^\star)\,g(0)$.
\item As distributions on $(0,\infty)$, $Q(-1-\theta)\big(t^{-1}\mu_{f,b_i}\big)=\sum_l\tilde P_{il}(-1-\theta)\,\mu_{f,b_l}$ for every $i$, $\theta=t\frac d{dt}$: a system of order at most $\max(\deg Q,\deg\tilde P_{il})$ for the $\mathfrak m$ densities.
\item With $v(k)=(\langle b_l\rangle_k)_l$ and every $K\ge2$ at which $\mathbf B(2),\dots,\mathbf B(K)$ are finite and $v(K)$ converges, the continuation of $v$ is finite at $1$ and $v(1)=\mathbf B(2)\mathbf B(3)\cdots\mathbf B(K)\,v(K)$: the values of the spanning set run back from $K$ with no scalar recurrence.
\end{enumerate}
\end{proposition}

The density operator of a scalar recurrence has the order of the degree of its coefficients, which carry the strengths where the rows line up (Proposition~\ref{prop:rays-lawrows}(iii)); the system of (ii) has the order of the degrees of the entries of $\mathbf B$, for $\mathfrak m$ densities whose germs at $f^\star$ start from the values $b_l(\alpha^\star)$ times the one constant $g(0)$.

% New: proved in the last part.
\begin{proposition}[The native gauges of all exponents]\label{prop:rd-sizes}
Let a family of $N$ tables on $n\ge2$ marks have the letter complement $X_F$ (Definition~\ref{def:ray}), and for $\alpha\in X_F$ and complex exponents $(\mathbf s,\mathbf t)\in\mathbb C^{2N}$ write $\mathbf u(\alpha;\mathbf s,\mathbf t)=\sum_j\big(s_jR_j(\alpha)+t_jR^\ast_j(\alpha)\big)\in\mathbb C^{\cE}$ for the usage (Definition~\ref{def:native}).
\begin{enumerate}[label=(\roman*)]
\item The data $\big(\mathbf u(\alpha;\mathbf s,\mathbf t);\mathbf s,\mathbf t\big)$ lie on the balance wall. Call $\alpha\in X_F$ \emph{native for} data $(\mvec;\mathbf s,\mathbf t)$ on the balance wall when $\mvec=\mathbf u(\alpha;\mathbf s,\mathbf t)$: the native gauges of a ray are the points native for its direction (Definition~\ref{def:raygauge}), and for $m_{\cE}\neq0$ the points native for data are its native points (Definition~\ref{def:native}). So the set of the $(\alpha,\mvec,\mathbf s,\mathbf t)$ with $\alpha$ native for $(\mvec;\mathbf s,\mathbf t)$ is the graph of the map $(\alpha,\mathbf s,\mathbf t)\mapsto\mathbf u(\alpha;\mathbf s,\mathbf t)$ on $X_F\times\mathbb C^{2N}$, a smooth irreducible variety of dimension $n-1+2N$, and every $\alpha\in X_F$ is native for the data $\big(\mathbf u(\alpha;\mathbf s,\mathbf t);\mathbf s,\mathbf t\big)$, whatever $(\mathbf s,\mathbf t)$.
\item A ray of $F$ with $\kappa_{1,S}>0$ for every nonempty proper $S$ has at most $(-1)^{n-1}\operatorname{eu}(X_F)$ nondegenerate native gauges, the number of native gauges of the directions outside the exceptional set of Theorem~\ref{thm:rayspace}(i).
\end{enumerate}
\end{proposition}

The nondegenerate native gauges bound the sizes of a reading from below: the dimension of the invariant module is at least the number of their orbits (Theorem~\ref{thm:invariant}(iii)), and the order of the law of $[1]$ at least the number of their distinct native values (Corollary~\ref{cor:symbol} at $g=1$); Proposition~\ref{prop:rd-sizes}(ii) bounds their number from above. These are inequalities, with $X_F$ inverting every letter whatever its exponent along the ray (Definition~\ref{def:ray}).

% New: proved in the last part.
\begin{definition}[The procedure that finds native gauges by continuation]\label{def:rd-continuation}
Input: a family and a ray of it with $\kappa_{1,S}>0$ for every nonempty proper $S$. Output: nondegenerate native gauges of the ray, each with a certificate, and their native values.
\begin{enumerate}[label=(\arabic*)]
\item \emph{A start.} Draw $(\alpha_0,\mathbf s_0,\mathbf t_0)\in X_F\times\mathbb C^{2N}$ at random; $\alpha_0$ is native for the data $w_0=\big(\mathbf u(\alpha_0;\mathbf s_0,\mathbf t_0);\mathbf s_0,\mathbf t_0\big)$ (Proposition~\ref{prop:rd-sizes}(i)).
\item \emph{More points native for $w_0$.} Follow the points native for $w_0$ found so far along closed loops $w_0\to w_1\to w_2\to w_0$ of random data on the balance wall, by Newton's method along the segments, and add every new endpoint; stop when several loops in a row add nothing.
\item \emph{To the direction.} Follow each point native for $w_0$ by Newton's method along a path of data from $w_0$ to the direction of the ray, and keep the endpoints that lie in $X_F$.
\item \emph{Certificates.} Around each endpoint, a box in the coordinates $\alpha_e$, $e\neq e_0$, of the chart on which interval enclosures of the letters exclude $0$; Krawczyk's interval Newton test on the equations $c^{(1)}_e=0$, $e\neq e_0$, in the box, and an interval enclosure of the determinant of the Hessian of $\log f$ on the box that excludes $0$. When both succeed, the box holds a native gauge of the ray, and every native gauge in it is nondegenerate. Disjoint boxes hold distinct native gauges, and disjoint enclosures of $f$ on them prove distinct native values.
\end{enumerate}
\end{definition}

The certified native gauges of step~(4) are distinct nondegenerate native gauges of the ray, so their number bounds the dimension of the ray from below (Theorem~\ref{thm:gaugesolution}(iii)), the number of their orbits under a finite group of symmetries bounds $\dim\mathfrak M^G$ from below (Theorem~\ref{thm:invariant}(iii)), and the number of their distinct native values bounds the order of every recurrence that follows from the relations (Corollary~\ref{cor:symbol}); by Proposition~\ref{prop:rd-sizes}(ii) their number is at most $(-1)^{n-1}\operatorname{eu}(X_F)$. The stopping rule of step~(2) is not a certificate that every point native for $w_0$ was found, so the procedure gives lower bounds. Nothing in it writes a relation of a box. Exercise~\ref{ex:listsb-table} reads the dimensions and the orders of the laws by the box, ray by ray.

% New chapter: inputs along a ray (after ch:readray).
% No new macros.
\chapter{Inputs along a ray}\label{ch:rayinputs}

A function that satisfies a recurrence with polynomial coefficients in its argument, above all the values of $1$ along a ray continued in the strength, gives two live inputs on a lattice: its values and its slopes. The slopes satisfy the same recurrence, inhomogeneous by the values read through the derivatives of its coefficients. A Trailmoment at index zero of a polynomial word against either input is closed by summation by parts in the adjoint form of the recurrence: what crosses the cut is read off the Laurent data of the function at the low end, for a ray the low-end relation and its derivative, and when the symbol does not vanish at $1$ everything else lands on a fixed block of $2D$ Trailmoments, $D$ the degree of the recurrence. A mark of a list opened on its share side and closed at an integer strength is such a pairing: a polynomial word against the slope input of the room, less its derivative against the values input. The recurrence read at the poles of an input ties its initial slopes to residues and finite parts, and minimality, the root rate of the input, removes the components along the faster solutions. A recurrence of the values that is self-adjoint under the reflection of the strength pairs the values with the closings in one constant, the low-end relation, whose limit at high strength ties the volume of the direction to the Hessian at its native gauge through the native values. The last section reads all of this on the centre ray of the list $\binom{\cE}1$, for every number $n$ of marks: its native gauges are the signs, and at a lean the points whose occupations take two values; a parallel class enters a list through the closings and the signs of $\binom{P}1$; a mark is added by closing the critical-line kernel against the slopes of $n$ marks; the law of the values is the $n$-th symmetric power of the second-order relation that the Gramma of one mark satisfies in its coefficient, read in the index, of order $\lceil n/2\rceil$, with the native values $(n-2j)^{-2}$ as the roots of its symbol, and one identity carries it from $n$ marks to $n+1$; the lowered settings are its solutions without low-end coefficient; and the exact pairing holds at every strength. The $n$-th power of the Gramma of one mark is a mixture of its dilations by the stride's law, and the balanced directions with one total and one multiset of offsets have one value up to Beta corners.

\section{Inputs read from a ray}\label{sec:ri-inputs}

% New: proved in the last part.
\begin{definition}[Recurrences; the values input and the slope input]\label{def:ri-inputs}
A \emph{recurrence} of order $h$ is an operator $(L\varphi)(x)=\sum_{i=0}^hc_i(x)\,\varphi(x+i)$ with polynomials $c_i\in\mathbb C[x]$ and $c_0c_h\neq0$. Its \emph{adjoint} and its \emph{derivative} are
\[
(L^\ast\varphi)(x)=\sum_{i=0}^hc_i(x-i)\,\varphi(x-i),\qquad (L'\varphi)(x)=\sum_{i=0}^hc_i'(x)\,\varphi(x+i),
\]
$L'^\ast$ is the adjoint of $L'$, its \emph{degree} is $D=\max_i\deg c_i$, and its \emph{symbol} is $\chi(\lambda)=\sum_i[x^D]c_i\,\lambda^i$, as in Corollary~\ref{cor:symbol}. Let $p$ be meromorphic on $\mathbb C$ with $Lp=0$ identically, holomorphic at the points of a coset $b+\mathbb Z$ beyond a cut $N$, and put $\tilde p=-p'$. The \emph{values input} of $p$ is $p\,\Ilat{b+\mathbb Z}$ and its \emph{slope input} is $\tilde p\,\Ilat{b+\mathbb Z}$ (Definition~\ref{def:inputs}), both read beyond the cut $N$. When $p$ continues the values $\langle1\rangle_k$ of a ray and $L$ is a recurrence of these values, they are the \emph{inputs read from the ray}.
\end{definition}

% New: proved in the last part.
\begin{proposition}[The values of a ray as an input]\label{prop:ri-rayinput}
Let a ray of a family on $n\ge2$ marks have origin $0$, every letter exponent of its direction positive and $\kappa_{1,S}>0$ for every nonempty proper $S$, and let $L$ be a recurrence of order $h$ with $\sum_{i=0}^hc_i(k)\langle1\rangle_{k+i}=0$ for all large integers $k$.
\begin{enumerate}[label=(\roman*)]
\item $\langle1\rangle_k$ is holomorphic in $k$ on $\re k>0$ and continues to a meromorphic function $p$ on $\mathbb C$ with $Lp=0$ identically. At every $j\in\mathbb N$, $\lim_{k\to-j}(k+j)^{n-1}p(k)=V_1\,\mathsf C_j(1)$, with $V_1$ of Proposition~\ref{prop:principal}, and $\sum_{i=0}^{\min(h,j)}c_i(-j)\,\mathsf C_{j-i}(1)=0$.
\item The slope input satisfies $L\tilde p=L'p$ identically: a recurrence of the same order, inhomogeneous by the values read through the derivative of $L$.
\item Along the positive integers $p$ has the root rate $f(\alpha^\star)$, the native value of the native gauge $\alpha^\star$ of the direction on the positive orthant, and $\tilde p$ a root rate at most $f(\alpha^\star)$.
\end{enumerate}
\end{proposition}

The values of a ray are therefore an input of the kind of Chapter~\ref{ch:transport}: their density is read from the serif family, and the recurrence gives the values input and the slope input recurrences of one finite order. The slope input carries what the derivative in the strength adds, and the values input enters its recurrence only through $L'$.

\section{Closing an input}\label{sec:ri-close}

% New: proved in the last part.
\begin{proposition}[Summation by parts at the cut]\label{prop:ri-sbp}
Let $p$ and $L$ be as in Definition~\ref{def:ri-inputs}, $p$ holomorphic at the integers $x\ge1$, with $p$ and $\tilde p$ of root rate below one along them, and for $0\le j<h$ put $\beta_j(x)=\sum_{i=0}^jc_i(x)\,p(x+i)$, which is holomorphic at $x=-j$ because $\beta_j=-\sum_{i>j}c_i(x)\,p(x+i)$. For every polynomial $R$,
\begin{gather*}
\Tx0{(L^\ast R)\,p\,\IZ}0=-\sum_{j=0}^{h-1}R(-j)\,\beta_j(-j),\\
\Tx0{(L^\ast R)\,\tilde p\,\IZ}0=\Tx0{(L'^\ast R)\,p\,\IZ}0+\sum_{j=0}^{h-1}R(-j)\,\beta_j'(-j).
\end{gather*}
If every $c_i$ vanishes to order two at the integers $-j$ with $1\le j<i$, only $j=0$ remains, with $\beta_0=c_0\,p$: the boundary is the value and the derivative of $c_0p$ at the cut.
\end{proposition}

What a summation by parts puts on the inputs beyond the cut is the adjoint word; what it leaves at the cut is read off $p$ at $0,-1,\dots,-(h-1)$, where for the values of a ray the poles sit. The same move closes the slope input, once through $L$ and once through $L'$.

% New: proved in the last part.
\begin{proposition}[The fixed block]\label{prop:ri-block}
Let $L$ be a recurrence of degree $D$ whose symbol has $\chi(1)\neq0$, and let $\Bbbk$ be a field containing the coefficients of the $c_i$.
\begin{enumerate}[label=(\roman*)]
\item $L^\ast$ raises the degree of every nonzero polynomial by exactly $D$, and $\Bbbk[x]=L^\ast(\Bbbk[x])\oplus\{r\in\Bbbk[x]:\ \deg r<D\}$.
\item In the setting of Proposition~\ref{prop:ri-sbp}, for every polynomial $W\in\Bbbk[x]$ the two Trailmoments
\[
\Tx0{W\,p\,\IZ}0,\qquad \Tx0{W\,\tilde p\,\IZ}0
\]
lie in the $\Bbbk$-span of the $2D$ \emph{block values}
\[
\Tx0{(x-1)^a\,p\,\IZ}0,\qquad \Tx0{(x-1)^a\,\tilde p\,\IZ}0\qquad(0\le a<D)
\]
and of the boundary numbers $\beta_j(-j)$ and $\beta_j'(-j)$, $0\le j<h$.
\end{enumerate}
\end{proposition}

This is the counterpart at the index zero of Theorem~\ref{thm:fixedblock}: the words of a closing are polynomials, the shifts of the input are moved onto the word, and the input is read on one block of evaluations fixed once for the recurrence. When the boundary numbers are known, as they are at the cut of a ray from the low-end relation, a closing lands on the block plus a number of $\Bbbk$.

\section{A mark opened at an integer strength}\label{sec:ri-mark}

A mark that is neither a loop nor a coloop is a bunch of width one, and Theorem~\ref{thm:twosides} opens the centre on it: at the strength $s$ the period is the integral over one even share $d$ of $\Gst42{d,\ d,\ s-d,\ s-d}{s,\ s}$ against the period of the room family, the deletion at the share $d$ and the contraction at $s-d$. Closing the share (Corollary~\ref{cor:closetwosides}) at an integer strength turns its heads into one polynomial word.

% New: proved in the last part.
\begin{corollary}[A mark opened at an integer strength]\label{cor:ri-openmark}
Let $\cB$ be a connected list on $\cE$, $|\cE|\ge3$, and $e$ a mark that is neither a loop nor a coloop; on the chart $y_{e_0}=0$ of $R=\cE-e$ let $X_\setminus=K_{\cB\setminus e}(\alpha)K_{\cB\setminus e}(\alpha^{-1})$ and $X_/=K_{\cB/e}(\alpha)K_{\cB/e}(\alpha^{-1})$, the even letters of the deletion and of the contraction over $\alpha^R$. Suppose that one lies below the other: $\{X_\flat,X_\sharp\}=\{X_\setminus,X_/\}$ with $\rho=\sup X_\flat/X_\sharp<1$ on the chart. For $s\in\{1,2,\dots\}$ and $x\ge0$ put
\[
q_s(x)=\int_{y_{e_0}=0}\Big(\frac{X_\flat}{X_\sharp}\Big)^xX_\sharp^{-s}\,dy,\qquad \tilde q_s=-\partial_xq_s,\qquad V_s(x)=\big((x+1)(x+2)\cdots(x+s-1)\big)^2 .
\]
Then $q_s$ converges and is holomorphic on a neighbourhood of $[0,\infty)$, and
\[
\Gamma(s)^2\,P_{\cB}(s\one;s,s)=\sum_{x\in\mathbb N}\big(V_s(x)\,\tilde q_s(x)-V_s'(x)\,q_s(x)\big),
\]
the terms of root rate at most $\rho$: the centre at the strength $s$ is a polynomial word against the slope input of $q_s$ less its derivative against its values input, read on $\mathbb N$.
\end{corollary}

The function $q_s$ is read along the direction of the ratio $X_\flat/X_\sharp$ of the two even letters, with the origin form $X_\sharp^{-s}dy$, which decays where the ratio does not; where it satisfies a recurrence in $x$, Sections~\ref{sec:ri-close} and~\ref{sec:ri-boundary} close it, and its values half $\sum_xV_s(x)\,q_s(x)$ with it.

\section{Boundaries and minimality}\label{sec:ri-boundary}

% New: proved in the last part.
\begin{proposition}[Boundary relations]\label{prop:ri-boundary}
Let $p$ and $L$ be as in Definition~\ref{def:ri-inputs}, let $x_0\in\mathbb C$ and $0\le j<h$, and let $p$ be holomorphic at $x_0+i$ for $j<i\le h$. Then $\beta_{x_0,j}(x)=\sum_{i=0}^jc_i(x)\,p(x+i)$ is holomorphic at $x_0$, its value and its derivative there are read off the Laurent expansions of $p$ at $x_0,\dots,x_0+j$ and the Taylor expansions of $c_0,\dots,c_j$ at $x_0$, and
\begin{gather*}
\sum_{i=j+1}^hc_i(x_0)\,p(x_0+i)=-\beta_{x_0,j}(x_0),\\
\sum_{i=j+1}^h\big(c_i(x_0)\,\tilde p(x_0+i)-c_i'(x_0)\,p(x_0+i)\big)=\beta_{x_0,j}'(x_0).
\end{gather*}
Where $c_h$ has no zero at the integers $x\ge1$, the slope input on these integers is fixed by its initial slopes $\tilde p(1),\dots,\tilde p(h)$ through $L\tilde p=L'p$, and every second relation at an integer $x_0$ with $x_0+j\ge0$ is an affine relation among them; at $x_0=-j$ its right side is read off the Laurent data of $p$ at the nonpositive integers $-j,\dots,0$ alone. At the cut of the values of a ray (Proposition~\ref{prop:ri-rayinput}), with $c_0(x)=\sum_m\gamma_mx^m$ and $p(x)=\sum_{i\ge0}\mathsf L_ix^{i-n+1}$ near $x=0$, the low-end coefficients of Definition~\ref{def:rays-lowend}, $c_0$ vanishes to order at least $n-1$ at $0$, $\mathsf L_0=V_1$, and at $x_0=0$, $j=0$ the two relations are
\[
\sum_{i=1}^hc_i(0)\,p(i)=-\gamma_{n-1}V_1,\qquad \sum_{i=1}^h\big(c_i(0)\,\tilde p(i)-c_i'(0)\,p(i)\big)=\gamma_{n-1}\mathsf L_1+\gamma_nV_1 ,
\]
the low-end relation (Corollary~\ref{cor:rays-lowend}(ii), which holds for every recurrence of the values at the large integers by Proposition~\ref{prop:sl-density}) and its derivative; when every table of the family has unit weights, $\mathsf L_1=0$.
\end{proposition}

At a pole of the input that a wall or a zero share puts on the lattice, the right sides are its principal part and its finite part: at a simple pole $x_0$ where $c_0$ has a simple zero, $\beta_{x_0,0}(x_0)=c_0'(x_0)\operatorname{Res}_{x_0}p$ and $\beta_{x_0,0}'(x_0)=c_0'(x_0)\,a+\frac12c_0''(x_0)\operatorname{Res}_{x_0}p$, $a$ the constant term of $p$ at $x_0$. Each independent relation removes one direction from the initial slopes; what the relations leave free, the root rate of the input fixes in the directions that grow faster than the input.

% New: proved in the last part.
\begin{proposition}[Minimality]\label{prop:ri-minimal}
Let $p$ and $L$ be as in Definition~\ref{def:ri-inputs}, $p$ holomorphic at the integers $x\ge1$, with $c_0(x)c_h(x)\neq0$ at these integers, every $c_i$ of degree at most $\deg c_0=\deg c_h$, and the roots $\lambda_1,\dots,\lambda_h$ of the symbol of distinct moduli. Let $\lambda>0$ be at least the root rates of $p$ and of $\tilde p$ along the positive integers (Definition~\ref{def:conv}(vii)).
\begin{enumerate}[label=(\roman*)]
\item The solutions $u$ of $Lu=0$ on the integers $x\ge1$ of root rate at most $\lambda$ form a space $U_\lambda$ of dimension $\#\{c:\ |\lambda_c|\le\lambda\}$, and every nonzero $u\in U_\lambda$ has $u(x+1)/u(x)\to\lambda_c$ for one of these roots, so the root rate $|\lambda_c|$.
\item The functions $\varphi$ on the integers $x\ge1$ with $L\varphi=L'p$ and root rate at most $\lambda$ are the $\tilde p+u$, $u\in U_\lambda$, and those with $L\varphi=0$ are the $p+u$. So the slope input is the only one of them that satisfies a set of second relations of Proposition~\ref{prop:ri-boundary} whose left sides, read on $u$, vanish together on no nonzero $u\in U_\lambda$; the values input is singled out by the first relations in the same way.
\end{enumerate}
\end{proposition}

So the initial slopes are fixed by as many boundary relations as there are roots of the symbol of modulus at most the root rate of the input; the roots beyond it are removed by the root rate alone.

\section{The two ends paired}\label{sec:ri-pairing}

% New: proved in the last part.
\begin{proposition}[The two ends paired]\label{prop:ri-pairing}
Let a ray be as in Proposition~\ref{prop:ri-rayinput}, with a recurrence $L$ of its values whose coefficients satisfy $c_i(-x-i)=\epsilon\,c_i(x)$ for every $i$, $\epsilon\in\{1,-1\}$, and put $\mathsf C_j=\mathsf C_j(1)$ for $j\ge0$ and $\mathsf C_j=0$ for $j<0$.
\begin{enumerate}[label=(\roman*)]
\item $\sum_{i=0}^hc_i(\sigma-i)\,\mathsf C_{\sigma-i}=0$ for every integer $\sigma$: the closings, read at the strength $\sigma$, solve the adjoint recurrence.
\item The \emph{concomitant}
\[
B(\sigma)=\sum_{i=1}^h\sum_{j=1}^ic_i(\sigma-j)\,\mathsf C_{\sigma-j}\,\langle1\rangle_{\sigma-j+i}
\]
takes one value at every integer $\sigma\ge1$, $B(\sigma)=\sum_{i=1}^hc_i(0)\langle1\rangle_i=-\gamma_{n-1}V_1$ with $\gamma_{n-1}=[x^{n-1}]c_0$: the low-end relation, carried to every strength.
\item On the centre ray of a family (Definition~\ref{def:centreray}), if every $c_i$ has degree at most $n-1$, then $c_0=\chi(0)\,x^{n-1}$, $\lambda_\star=f(\alpha^\star)$ is a simple root of the symbol, and
\[
V_1\,\det H=-\frac{\lambda_\star\,\chi'(\lambda_\star)}{\chi(0)}=\prod_{c\neq\star}\Big(1-\frac{\lambda_\star}{\lambda_c}\Big),
\]
with $H$ the Hessian of $-\log f$ at $\one$ on the chart (Theorem~\ref{thm:branch}) and $\lambda_c$ the roots of the symbol, counted with multiplicity.
\end{enumerate}
\end{proposition}

The values at the high end and the closings at the low end are two solutions of one recurrence and of its adjoint, and their concomitant is the constant that the low end fixes. At high strength the two decay and grow at the reciprocal rates $f(\alpha^\star)^{\pm k}$ with the Gaussian constants of the native gauge of the direction, and the constant becomes an identity between the volume $V_1$ of the direction, read by its walls, and the Hessian at its native gauge, through the other native values.

\section{The centre ray of \texorpdfstring{$\binom{\cE}1$}{(E 1)}}\label{sec:ri-centre}

For the list $\binom{\cE}1$ on $n$ marks, whose bases are the singletons, $K(\alpha)=\sum_e\alpha_e$ and $\Psi(\alpha)=\alpha^{\cE}K(\alpha^{-1})$, and the stride of the centre ray at the strength $\sigma$ is $\big(K(\alpha)K(\alpha^{-1})\big)^{-\sigma}=\big(\sum_{e,e'}\alpha_e/\alpha_{e'}\big)^{-\sigma}$. Its widths are all one, so $V_1=n!$ at the strength one. Its share side is one Gramma read $n$ times: by Corollary~\ref{cor:cycle} for the dual table, whose bases are the sets $\cE-e$, and Proposition~\ref{prop:dualfamily}, for masses $\mvec$ on the balance wall, $m_{\cE}=s+(n-1)t$, with $m_S>t(|S|-1)$ for every nonempty proper $S$,
\[
\Gamma(s)\,\Gamma(t)\,P_{\binom{\cE}1}(\mvec;s,t)=\int_0^\infty y^{t-1}\prod_{e\in\cE}\sG_{m_e-t;(1,-1)}(1,y)\,dy ,
\]
and at the centre $\Gamma(s)^2P_{\binom{\cE}1}(s\one;s,s)=\int_0^\infty y^{s-1}\sG_{0;(1,-1)}(1,y)^n\,dy$. With $\vartheta=y\,d/dy$, the ladder and the exchange relation of the Gramma (Theorem~\ref{thm:grammalaws}) give $\vartheta^2\mathsf g=y\,\mathsf g$ for $\mathsf g(y)=\sG_{0;(1,-1)}(1,y)$, and the partner $\mathsf h(y)=\sum_{m\ge0}y^m/m!^2$ of its lattice expansion (Lemma~\ref{lem:onemark}(i), $\mathsf g=-\mathsf h\log y+\sum_m2(\log\Gamma)'(1+m)\,y^m/m!^2$) satisfies the same relation. The statements below read the centre ray from these two functions.

% New: proved in the last part.
\begin{proposition}[The native gauges of $\binom{\cE}1$]\label{prop:ri-signs}
Let $n\ge2$, let $s,t\ge1$ and $m$ be integers with $nm=s+(n-1)t$, and read a ray of $\binom{\cE}1$ in the direction $(m\one;s,t)$; in this statement $u_e(\alpha)=\alpha_e/K(\alpha)$ is the occupation of the mark $e$ and $\lambda(\alpha)=\big(K(\alpha)K(\alpha^{-1})\big)^{-1}$.
\begin{enumerate}[label=(\roman*)]
\item A point $\alpha$ of $X_F$ is a native gauge exactly when its occupations take at most two values, $u_+$ on $n-j$ marks and $u_-$ on $j$ marks, $0\le j\le n$, $2j\neq n$, with
\[
(n-2j)\,u_+=1-j\,\frac{s-t}{ns},\qquad u_++u_-=\frac{s-t}{ns},\qquad u_+u_-\neq0 ;
\]
then $\lambda(\alpha)=-s\,u_+u_-/t$ and the native value is $f(\alpha)=\big(\prod_eu_e\big)^{m-t}\lambda(\alpha)^t$. Up to scale these points are $\alpha=(u_e)_e$, one for each set of $j$ marks, a set and its complement giving the same point.
\item At the centre, $s=t=\sigma$, $m=\sigma$, the native gauges are exactly the signs $\eta$ with $K(\eta)=\sum_e\eta_e\neq0$, of the native values $K(\eta)^{-2\sigma}=(n-2j)^{-2\sigma}$, $j$ the number of marks at $-1$: there are $\lceil n/2\rceil$ native values. Each of these signs is nondegenerate: on the chart $\alpha_e=\eta_ee^{y_e}$, $y_{e_0}=0$, the Hessian of $-\log f$ at $y=0$ has the determinant $(2\sigma)^{n-1}\prod_e\eta_e/K(\eta)^n$.
\end{enumerate}
\end{proposition}

A lean moves every sign to a real point with the same pattern: the occupations keep two values, on the same marks, and only the two values change. Every list has the signs with $K(\eta)\neq0$ among the native gauges of its centre ray (Theorem~\ref{thm:signgauges}); the list $\binom{\cE}1$ has nothing else.

% New: proved in the last part.
\begin{proposition}[A parallel class enters through its list]\label{prop:ri-parallel}
Let $\mathcal H$ be a connected list on the marks $1,\dots,q$, $q\ge2$, and replace every mark $i$ by a parallel class $P_i$ of $s_i\ge1$ marks by the first rule of Theorem~\ref{thm:block}: the list $\cB$ on $\cE=\bigsqcup_iP_i$ whose bases are the sets with one mark in $P_i$ for every $i\in T$ and none in the other classes, $T$ a basis of $\mathcal H$. Write $A_i=\sum_{e\in P_i}\alpha_e$.
\begin{enumerate}[label=(\roman*)]
\item $K_{\cB}(\alpha)=K_{\mathcal H}(A_1,\dots,A_q)$, and for $v\in\mathbb N^{\cE}$, with $w_i=v_{P_i}$, $[\alpha^v]K_{\cB}^k=\big([A^w]K_{\mathcal H}^k\big)\prod_i\binom{w_i}{v|_{P_i}}$, with the multinomial coefficients $\binom{w_i}{v|_{P_i}}$.
\item The closings of the centre ray of $\cB$ at the strength one (Definition~\ref{def:closings}) are
\[
\mathsf C^{\cB}_k(1)=\sum_{w\in\mathbb N^q}\big([A^w]K_{\mathcal H}^k\big)^2\prod_{i=1}^q\mathsf C^{\binom{P_i}1}_{w_i}(1),\qquad \mathsf C^{\binom{P}1}_w(1)=\sum_{v\in\mathbb N^P,\ |v|=w}\binom wv^2 ,
\]
the second the closings of the centre ray of $\binom{P}1$ at the strength one, read as $1$ when $|P|=1$.
\item At a sign $\eta$ of $\cB$, $K_{\cB}(\eta)=K_{\mathcal H}(a_1,\dots,a_q)$ with $a_i=\sum_{e\in P_i}\eta_e\in\{s_i,s_i-2,\dots,-s_i\}$, the sign values of $\binom{P_i}1$: on the centre ray of $\cB$ at the strength $\sigma$ the native values of the signs with $K_{\cB}(\eta)\neq0$ (Theorem~\ref{thm:signgauges}(ii)) are the $K_{\mathcal H}(\mathbf a)^{-2\sigma}$ at these points $\mathbf a$.
\end{enumerate}
\end{proposition}

A parallel class enters the closings of a list only through the closings of its own list $\binom{P}1$, and the signs only through the sign values of that list, as a parallel block is read by $K$ through its total and by $\Psi$ through its hinge value (Corollary~\ref{cor:blocksclosures}(i)).

% New: proved in the last part.
\begin{proposition}[A mark added is a closing of the critical-line kernel]\label{prop:ri-addmark}
Let $n\ge2$; in this statement $p_n(s)=P_{\binom{\cE}1}(s\one;s,s)=\int_{y_{e_0}=0}\big(K(\alpha)K(\alpha^{-1})\big)^{-s}dy$ for $|\cE|=n$ and $\re s>0$, $\tilde p_n=-p_n'$, and $W_s(x)=\prod_{j=1}^{s-1}(x-j)^2$.
\begin{enumerate}[label=(\roman*)]
\item For $\re s>0$ and $0<c<\re s$, $\Gamma(s)^2p_{n+1}(s)=\frac1{2\pi i}\int_{(c)}\Gamma(d)^2\,\Gamma(s-d)^2\,p_n(d)\,dd$. At $s=1$ the kernel on the line $\re d=\frac12$ is $2\pi$ times the critical-line density (Definition~\ref{def:pairlaw}): $p_{n+1}(1)=\oE\big[p_n(\tfrac12+iY)\big]$, $Y\sim\mathbb P_{\mathrm{crit}}$.
\item For $s\in\{1,2,\dots\}$,
\begin{gather*}
\Gamma(s)^2p_{n+1}(s)=\Tx0{(W_s\tilde p_n-W_s'p_n)\,\IZ}0,\\
\int_0^\infty y^{s-1}\sG_{0;(1,-1)}(1,y)^n\sum_{m\ge0}\frac{y^m}{m!^2}\,dy=\Tx0{W_s\,p_n\,\IZ}0 ;
\end{gather*}
at $s=1$, $p_{n+1}(1)=\Tx0{\tilde p_n\,\IZ}0$: the centre of $n+1$ marks at the strength one is the slope input of $n$ marks summed over the positive integers.
\item On the share side every mark carries on average the share $\frac1{n+1}$ of the strength:
\[
(n+1)\Big(\Tx0{p_n\,\IZ}0-\Tx0{(x-1)\,\tilde p_n\,\IZ}0\Big)=\Tx0{\tilde p_n\,\IZ}0 .
\]
\end{enumerate}
\end{proposition}

% New: proved in the last part.
\begin{theorem}[The law of the values of $\binom{\cE}1$]\label{thm:ri-law}
Let $n\ge1$ and $h=\lceil n/2\rceil$; in this statement $\vartheta=y\,d/dy$, and the operators $\mathbf R_i$ are $\mathbf R_0=1$, $\mathbf R_1=\vartheta$ and $\mathbf R_{i+1}=\vartheta\mathbf R_i-i(n+1-i)\,y\,\mathbf R_{i-1}$.
\begin{enumerate}[label=(\roman*)]
\item $\mathbf Q_n=\mathbf R_{n+1}$ annihilates every product of $n$ solutions of $\vartheta^2u=yu$. Written $\mathbf Q_n=\sum_{k\ge0}y^kq_{n,k}(\vartheta)$ with polynomials $q_{n,k}$, $q_{n,0}=\vartheta^{n+1}$ and $\deg q_{n,k}\le n+1-2k$, so that $q_{n,k}=0$ for $k>h$; with $a_{n,k}$ the coefficient of $\vartheta^{n+1-2k}$ in $q_{n,k}$,
\[
\sum_ka_{n,k}\,\eta^k\xi^{n+1-2k}=\prod_{j=0}^n\big(\xi-(n-2j)\sqrt\eta\big)=\xi^{n+1-2h}\prod_{0\le j<n/2}\big(\xi^2-(n-2j)^2\eta\big),
\]
so $q_{n,h}\neq0$. The formal adjoint for the measure $dy/y$ is $\mathbf Q_n^\ast=(-1)^{n+1}\mathbf Q_n$: $q_{n,k}(-\vartheta-k)=(-1)^{n+1}q_{n,k}(\vartheta)$.
\item Put $c_{n,0}(s)=(-1)^{n+1}s^{n-1}$ and $c_{n,k}(s)=q_{n,k}(-s-k)\big((s+1)(s+2)\cdots(s+k-1)\big)^2$ for $1\le k\le h$. Every $c_{n,k}$ has degree $n-1$, $c_{n,k}(-s-k)=(-1)^{n+1}c_{n,k}(s)$, $c_{n,k}$ vanishes to order two at $s=-1,\dots,-(k-1)$, $c_{n,h}$ has no zero at the integers $s\ge0$, and the symbol is
\[
\chi_n(\lambda)=(-1)^{n+1}\prod_{0\le j<n/2}\big(1-(n-2j)^2\lambda\big),
\]
whose roots are the native values $(n-2j)^{-2}$ of the signs (Proposition~\ref{prop:ri-signs}). For $n\ge2$ the values satisfy
\[
\sum_{k=0}^hc_{n,k}(s)\,P_{\binom{\cE}1}\big((s+k)\one;s+k,s+k\big)=0
\]
identically in $s$.
\item For $n\ge2$ the values $P_{\binom{\cE}1}(k\one;k,k)$ satisfy no recurrence of order less than $h$: their tier (Proposition~\ref{prop:frame}) is $h$, and the recurrence of (ii) is their least recurrence. The closings $\mathsf C_j(1)=\sum_{|v|=j}\binom jv^2$ satisfy $\sum_{k=0}^{\min(h,j)}c_{n,k}(-j)\,\mathsf C_{j-k}(1)=0$ for every $j\in\mathbb N$ and no recurrence of order less than $h$, and the cyclic submodule of $[1]$ in the value module (Definition~\ref{def:valuemodule}) has dimension $h$.
\end{enumerate}
\end{theorem}

The law of the values of $\binom{\cE}1$ is the relation of one mark raised to the $n$-th symmetric power and read in the index, where $\vartheta$ becomes $-s$ and $y$ the shift. Its leading coefficients are those of the product of the $n+1$ factors $\xi-(n-2j)\sqrt\eta$, which pair into the native values of the signs, the tie $n=2j$ at even $n$ leaving the factor $\xi$; every coefficient $c_{n,k}$ carries the square of a rising factorial, which is why the boundaries below the cut vanish. One identity carries it from $n$ marks to $n+1$.

% New: proved in the last part.
\begin{proposition}[A mark added carries the law of the values]\label{prop:ri-markadded}
In the notation of Theorem~\ref{thm:ri-law}, for every $n\ge1$:
\begin{enumerate}[label=(\roman*)]
\item For every solution $u$ of $\vartheta^2u=yu$ and every smooth function $F$ on $(0,\infty)$,
\[
\mathbf Q_{n+1}(uF)=u\,\vartheta\,\mathbf Q_n(F)+(n+2)\,(\vartheta u)\,\mathbf Q_n(F).
\]
\item With $\binom ab=\Gamma(a+1)/\big(\Gamma(b+1)\Gamma(a-b+1)\big)$ and $A_s(x)=\big((n+1)x-(n+2)s\big)\binom xs^2$,
\[
\sum_{k=0}^{\lceil(n+1)/2\rceil}c_{n+1,k}(s)\binom{x-1}{s+k-1}^2=\sum_{j=0}^{\lceil n/2\rceil}c_{n,j}(x-j)\,A_s(x-j),
\]
an identity of rational functions of $(x,s)$ after division by $\binom xs^2$. At integers $s,x\ge1$ the words $\binom{x-1}{s+k-1}^2$ are the words $W_{s+k}(x)/\Gamma(s+k)^2$ of the closing of Proposition~\ref{prop:ri-addmark}(ii): weighted by the coefficients $c_{n+1,k}(s)$, they are the adjoint of the recurrence of $n$ marks applied to the one word $A_s$, which vanishes at $x=0,1,\dots,s-1$.
\end{enumerate}
\end{proposition}

% New: proved in the last part.
\begin{proposition}[The lowered settings]\label{prop:ri-lowered}
Let $n\ge2$ and $h=\lceil n/2\rceil$, with $c_{n,k}$ of Theorem~\ref{thm:ri-law}; in this statement $\mathsf g(y)=\sG_{0;(1,-1)}(1,y)$ and $\mathsf h(y)=\sum_{m\ge0}y^m/m!^2$, and for $0\le j<n/2$ the \emph{lowered setting} with $j$ marks lowered is
\[
\xi_{n,j}(s)=\Gamma(s)^{-2}\int_0^\infty y^{s-1}\,\mathsf g(y)^{n-j}\,\mathsf h(y)^j\,dy\qquad(\re s>0),
\]
so that $\xi_{n,0}(s)=P_{\binom{\cE}1}(s\one;s,s)$, $|\cE|=n$.
\begin{enumerate}[label=(\roman*)]
\item Every $\xi_{n,j}$ satisfies the recurrence of Theorem~\ref{thm:ri-law}(ii) on $\re s>0$.
\item For $j\ge1$, $s^{n-1}\xi_{n,j}(s)\to0$ as $s\downarrow0$, and $\sum_{k=1}^hc_{n,k}(0)\,\xi_{n,j}(k)=0$, against $\sum_{k=1}^hc_{n,k}(0)\,\xi_{n,0}(k)=(-1)^nn!$: the lowered settings have no low-end coefficient.
\item $\xi_{n,j}(k+1)/\xi_{n,j}(k)\to(n-2j)^{-2}$ along the integers, so $\xi_{n,j}$ has the root rate $(n-2j)^{-2}$: the setting with $j$ marks lowered carries the native value of the signs with $j$ marks at $-1$.
\item The $h$ sequences $(\xi_{n,j}(k))_{k\ge1}$ are a basis of the solutions of the recurrence on the integers $k\ge1$, and its solutions $x$ with $\sum_{k=1}^hc_{n,k}(0)\,x(k)=(-1)^nn!$ are the $\xi_{n,0}+\sum_{j=1}^{h-1}b_j\,\xi_{n,j}$, $b_j\in\mathbb C$.
\end{enumerate}
\end{proposition}

So the recurrence and the low-end relation fix the values only up to the $h-1$ lowered settings with $j\ge1$, which satisfy both with the right side zero; $h-1$ further numbers fix them, for instance the first $h-1$ values (Corollary~\ref{cor:ri-span}). The integrand $\mathsf g^{n-j}\mathsf h^j$ of the setting with $j$ marks lowered decays like $e^{-2(n-2j)\sqrt y}$ up to powers of $y$, and its integral converges on all of $\re s>0$ exactly when $j<n/2$: at even $n$ the tie $j=n/2$, the sign with $K(\eta)=0$, is no native gauge and gives no setting.

% New: proved in the last part.
\begin{corollary}[The value span of $\binom{\cE}1$]\label{cor:ri-span}
Let $n\ge2$, $h=\lceil n/2\rceil$ and $p_n(s)=P_{\binom{\cE}1}(s\one;s,s)$, $|\cE|=n$, with $c_{n,k}$ of Theorem~\ref{thm:ri-law}.
\begin{enumerate}[label=(\roman*)]
\item $\sum_{k=1}^hc_{n,k}(0)\,p_n(k)=(-1)^nn!$, and every value $p_n(k)$, $k\ge1$, lies in $\mathbb Q+\sum_{i=1}^{h-1}\mathbb Q\,p_n(i)$; so do the values $p_n(k\sigma)$ of the centre ray at every strength $\sigma$.
\item For even $n$ the symbol has $\chi_n(1)\neq0$, and for every polynomial $W$ with rational coefficients, $\Tx0{W\,p_n\,\IZ}0$ and $\Tx0{W\,\tilde p_n\,\IZ}0$, $\tilde p_n=-p_n'$, lie in $\mathbb Q$ plus the $\mathbb Q$-span of the $2(n-1)$ block values $\Tx0{(x-1)^a\,p_n\,\IZ}0$ and $\Tx0{(x-1)^a\,\tilde p_n\,\IZ}0$, $a<n-1$. In particular so do $\Gamma(s)^2P_{\binom{[n+1]}1}(s\one;s,s)$ and $\Gamma(s)^2\xi_{n+1,1}(s)$ at every integer $s\ge1$.
\end{enumerate}
\end{corollary}

% New: proved in the last part.
\begin{corollary}[The exact pairing of $\binom{\cE}1$]\label{cor:ri-pairing}
Let $n\ge2$, $h=\lceil n/2\rceil$, and $p_n$ and $c_{n,k}$ as in Corollary~\ref{cor:ri-span}; put $\mathsf C_l=\sum_{v\in\mathbb N^{\cE},\ |v|=l}\binom lv^2$ for $l\ge0$ and $\mathsf C_l=0$ for $l<0$. For every integer $\sigma\ge1$,
\[
\sum_{i=1}^h\sum_{j=1}^ic_{n,i}(\sigma-j)\,\mathsf C_{\sigma-j}\,p_n(\sigma-j+i)=(-1)^nn!,
\]
and its limit at high strength is the identity
\[
n!\cdot\frac{2^{n-1}}{n^n}=\prod_{j=1}^{h-1}\Big(1-\frac{(n-2j)^2}{n^2}\Big)=\prod_{j=1}^{h-1}\frac{4j(n-j)}{n^2}
\]
of Proposition~\ref{prop:ri-pairing}(iii), with $V_1=n!$, $\det H=2^{n-1}/n^n$, $\lambda_\star=n^{-2}$ and the other roots $(n-2j)^{-2}$.
\end{corollary}

% New: proved in the last part.
\begin{proposition}[The power of one mark as a mixture]\label{prop:ri-mixture}
Let $n\ge2$; in this statement $\mathsf g(y)=\sG_{0;(1,-1)}(1,y)$, and $f_n$, $f_{n+1}$ are the strides $\big(K(\alpha)K(\alpha^{-1})\big)^{-1}$ of the centre rays of $\binom{\cE}1$ at the strength one for $|\cE|=n$ and for $|\cE|=n+1$, with the stride's laws $\mu_{f_n}$ and $\mu_{f_{n+1}}$ (Definition~\ref{def:strideslaw}): the images of the Lebesgue measure $dy$ of the chart $y_{e_0}=0$ under the strides, carried by $(0,n^{-2}]$ and by $(0,(n+1)^{-2}]$.
\begin{enumerate}[label=(\roman*)]
\item $\mathsf g(y)^n=\int\mathsf g(y/t)\,\mu_{f_n}(dt)$ for every $y>0$: the $n$-th power of the Gramma of one mark is a mixture of its dilations by the stride's law.
\item For $t>0$ let $\phi_t(x)=x/\big((x+t^{-1/2})(1+t^{-1/2}x)\big)$ on $(0,\infty)$, the stride of the centre ray at the strength one of the table $\binom{\{1,2\}}1$ with the weights $t^{-1/2}$ on its basis $\{1\}$ and $1$ on $\{2\}$, on the chart $\alpha_1=1$, $\alpha_2=x$; its stride's law $\mu_{\phi_t}$ is the image of $dx/x$ under $\phi_t$. Then $\mu_{f_{n+1}}=\int\mu_{\phi_t}\,\mu_{f_n}(dt)$: adding a mark mixes the stride's law of two marks over the stride's law of $n$.
\end{enumerate}
\end{proposition}

% New: proved in the last part.
\begin{proposition}[Balanced directions of $\binom{\cE}1$]\label{prop:ri-directions}
Let $|\cE|=n\ge2$, let $s,t\ge1$ be integers and $\mvec\in\mathbb Z^{\cE}$ with $m_{\cE}=s+(n-1)t$ and $m_S>t(|S|-1)$ for every nonempty proper $S$, and put $c_e=m_e-t$, so that $\sum_ec_e=s-t$. For real $k>0$, with $P(k)=P_{\binom{\cE}1}(k\mvec;ks,kt)$,
\[
\Gamma(sk)\,\Gamma(tk)\,P(k)=\int_0^\infty y^{\frac{(s+t)k}2-1}\prod_{e\in\cE}\sG_{c_ek;(1,-1)}\big(\sqrt y,\sqrt y\big)\,dy ,
\]
which depends on the direction only through $s+t$ and the multiset $\{|c_e|\}_e$. Two such directions $(\mvec;s,t)$ and $(\mvec';s',t')$ with $s+t=s'+t'$ and the same multiset have, with $P'(k)=P_{\binom{\cE}1}(k\mvec';ks',kt')$,
\[
\frac{P'(k)}{P(k)}=\Gst21{sk,\ tk}{(s+t)k}\Big/\Gst21{s'k,\ t'k}{(s+t)k},
\]
a quotient of Beta corners at one total, rational at the integers $k$, so that their values span one $\mathbb Q$-space; the inversion $\alpha\mapsto\alpha^{-1}$ carries the integrand at $(\mvec;s,t)$ to the integrand at $((s+t)\one-\mvec;t,s)$, which is $c\mapsto-c$.
\end{proposition}

The exercises work these statements on the lists of a few marks: the closings and the law of the values of $\binom{[n]}1$ for $n=3,4,5$ (Chapter~\ref{ch:exrays}), the value at the centre of five marks as a Trailmoment of a ray input (Exercise~\ref{ex:cfive-slope}), and the lean $(2\one;4,1)$ of three marks (Exercise~\ref{ex:smallc-lean}).

\chapter{Arithmetic along a ray}\label{ch:rayarith}

A free share closed along a ray brings at every lattice point a rational number of tier one in the point times Trailmoments of the counting input at cuts that move with the point; summed, these are Trailmoments of live inputs, and along the ray the same finitely many live inputs carry every value, the step from one point of the ray to the next entering as rational words. The terms are rational in the lattice point, up to a sign, exactly when the lower Gammas of the word whose slope is at least two in modulus pair off; the word is then a product of sines times a rational function, and its closings are combinations of Trailmoments of lattice inputs with coefficients polynomial in $\Tx{2}{\IZ}{0}$. The frame of a law of the values bounds every clearing of their rows from below, by the product of the roots of the symbol, and the merit of the rows, the decay at the native gauge of the direction against the growth of the clearing, decides the gap along the ray. On a ray from the origin the law of $[1]$ vanishes to the order $n-1$ at its low end, which gives one linear relation among the first values; at order two every value is then a row in the one coordinate $\langle1\rangle_1$, and when the law is reflexive the coefficient of that coordinate is a closing of the ray, the constant term of a power of the reciprocal stride, read with no lattice sum. A hinge closing is a row over finitely many periods of the room as soon as the room's periods along the lattice of the closing satisfy a recurrence of finite order. The critical primes of a native gauge, the primes at which its native value is not integral, are primes at which the letters degenerate, and when every root of the symbol of a law is such a native value they are the primes of its leading coefficient.

\section{Closing a free share along a ray}\label{sec:ra-closings}

% New: proved in the last part.
\begin{theorem}[Closing one free share along a ray]\label{thm:ra-oneshare}
Let $A$ and $B$ be finite sets with integers $b_\gamma$ and $u_\gamma$ for $\gamma\in A\sqcup B$, \emph{slopes} $\epsilon_\alpha\in\{1,-1\}$ for $\alpha\in A$ and $\epsilon_\beta\in\mathbb Z\setminus\{0\}$ for $\beta\in B$, balanced in the slopes, $\sum_{\alpha\in A}\epsilon_\alpha=\sum_{\beta\in B}\epsilon_\beta$. In this statement write, for $k\in\mathbb N$ and a share $a$,
\[
W_k(a)=\frac{\prod_{\alpha\in A}\Gamma(\epsilon_\alpha a+b_\alpha+ku_\alpha)}{\prod_{\beta\in B}\Gamma(\epsilon_\beta a+b_\beta+ku_\beta)},
\]
the stacked Gamma (Definition~\ref{def:stacked}) of one free share at the $k$-th point of the ray along the integral direction $(u_\gamma)_\gamma$; write $A_\pm$ for the $\alpha\in A$ with $\epsilon_\alpha=\pm1$, $B_\pm$ for the $\beta\in B$ with $\pm\epsilon_\beta>0$, $w=|A_+|-|B_+|$, and
\[
r_W=\prod_{\beta\in B_+}\epsilon_\beta^{\,\epsilon_\beta}\Big/\prod_{\beta\in B_-}|\epsilon_\beta|^{|\epsilon_\beta|}.
\]
\begin{enumerate}[label=(\roman*)]
\item \emph{The step factor.} $W_k=W_0\,g_k$ with
\[
g_k(a)=\frac{\prod_{\alpha\in A}(\epsilon_\alpha a+b_\alpha)_{ku_\alpha}}{\prod_{\beta\in B}(\epsilon_\beta a+b_\beta)_{ku_\beta}},\qquad (x)_{-r}=\frac1{(x-r)_r}\quad(r\in\mathbb N),
\]
a rational function of the share with rational coefficients.
\item \emph{The terms.} The poles of $W_k$ lie at integers. Let $n_1$ be an integer beyond which, at $a=-n$, the Gammas of $A_+$ and of $B_+$ have nonpositive integer arguments and those of $A_-$ and of $B_-$ positive ones, and for $\gamma\in A\sqcup B$ let $N_\gamma(n)$ be minus the argument of its Gamma at $a=-n$ in the first case and that argument less one in the second; $N_\gamma(n)$ is affine in $n$ with the slope $|\epsilon_\gamma|$. For every integer $n>n_1$ and small $\varepsilon$,
\[
W_k(-n+\varepsilon)=C_n\,\varepsilon^{-w}\exp\Big(\sum_{i\ge1}\Lambda_{n,i}\,\varepsilon^i\Big),
\]
where $C_n$ is a nonzero rational number, a signed product of the factorials $N_\gamma(n)!$, to the power $1$ for $\gamma\in A_-\cup B_+$ and $-1$ for $\gamma\in A_+\cup B_-$, and of the slopes of $B_+$; $C_{n+1}=\rho(n)\,C_n$ for a rational function $\rho$ with rational coefficients, so that $(C_n)_n$ has tier at most one (Proposition~\ref{prop:frame}), and $|C_n|^{1/n}\to r_W$, so that $r_W$ is the root rate (Definition~\ref{def:conv}(viii)) of the $C_n$; $\Lambda_{n,1}$ is a rational combination of the heads $\Th{1}{\IZ}{0}{N_\gamma(n)}$; and for $i\ge2$, $\Lambda_{n,i}$ is a rational combination of $\Tx{i}{\IZ}{0}$ and the heads $\Th{i}{\IZ}{0}{N_\gamma(n)}$. The coefficients of these combinations do not depend on $n$ or $k$. The residue of $W_k$ at $-n$ is $C_n$ times the coefficient of $\varepsilon^{w-1}$ in the exponential, a polynomial with rational coefficients, not depending on $n$, in the Trailmoments $\Tx{i}{\IZ}{0}$, $2\le i\le w-1$, and the heads $\Th{i}{\IZ}{0}{N_\gamma(n)}$, $1\le i\le w-1$, each of its monomials of total index $w-1$; it is zero when $w\le0$.
\item \emph{The closing.} Let $W_k$ close to the left on a line $\re a=c$ that separates the poles of the Gammas of $A_+$ from those of $A_-$ (Theorem~\ref{thm:closing}). Then $\frac1{2\pi i}\int_{(c)}W_k(a)\,da$ is the sum of the residues of $W_k$ at the integers $-n<c$. When $r_W<1$ it is a finite rational combination of products of Trailmoments $\Tx{i}{\IZ}{0}$, $i\ge2$, with $1$ and with the absolutely convergent sums
\[
\Tx{0}{\psi\,\IZ}{n_1}=\sum_{n>n_1}\psi(n),\qquad \psi(n)=C_n\prod_{l=1}^{q}\Th{i_l}{\IZ}{0}{N_{\gamma_l}(n)},
\]
over the finitely many products with $\gamma_l\in A\sqcup B$, $i_l\ge1$ and $i_1+\cdots+i_q\le w-1$; each such $\psi$ is a live input, its density read from the serif family at the moving cut, and it satisfies a recurrence of order at most $2^q$ with polynomial coefficients and nonzero extreme coefficients.
\item \emph{Along the ray.} For two steps $k$ and $k'$ and every $n$ beyond both bounds $n_1$, $C_n$ at $k$ is $C_n$ at $k'$ times a rational function of $n$ with rational coefficients, and every $N_\gamma(n)$ at $k$ is $N_\gamma(n)$ at $k'$ plus an integer that does not depend on $n$. So every input $\psi$ of (iii) at the step $k$ is a combination, with coefficients rational functions of $n$ with rational coefficients, of the inputs $\psi$ of (iii) at the step $k'$: when $r_W<1$ every value along the ray is a finite rational combination of products of Trailmoments $\Tx{i}{\IZ}{0}$ with $1$ and with Trailmoments at index zero of the finitely many live inputs of one step against rational functions of the point with rational coefficients, whose rows Theorem~\ref{thm:ra-liverow} reads.
\end{enumerate}
Closing to the right is the same statement after the share $a$ is replaced by $-a$.
\end{theorem}

A word of one free share arises from the atom formula (Corollary~\ref{cor:atomwall}) when its other shares are closed or merged, for instance by the four Gammas on a line (Corollary~\ref{cor:fourgammas}): integrating $\Gamma(a+z)\Gamma(b+z)\Gamma(c-z)\Gamma(d-z)$ over $z$ leaves four upper Gammas and the lower Gamma $\Gamma(a+b+c+d)$, whose slope in a remaining share is the sum of the slopes of $a$, $b$, $c$ and $d$. A lower Gamma of slope $v$ contributes $v^v$ to the root rate of the terms and one of slope $-v$ contributes $v^{-v}$, so a lone lower Gamma of slope $-2$ makes the terms of a left closing fall like $4^{-n}$, while lower Gammas that pair off leave terms that do not fall geometrically and are rational in $n$ up to a sign.

% New: proved in the last part.
\begin{proposition}[Paired words]\label{prop:ra-paired}
In Theorem~\ref{thm:ra-oneshare} the following are equivalent.
\begin{enumerate}[label=(\alph*)]
\item The lower Gammas of slope at least two in modulus pair off: for every integer $v\ge2$ as many $\beta\in B$ have $\epsilon_\beta=v$ as have $\epsilon_\beta=-v$.
\item For some $k$, and then for every $k$, there are $\varsigma\in\{0,1\}$ and a rational function $R$ with $C_n=(-1)^{\varsigma n}R(n)$ for every $n>n_1$.
\end{enumerate}
When they hold, put $p=|A_+|-\#\{\beta\in B:\epsilon_\beta=1\}$ and let $B^{\ge2}$ be the set of the $\beta$ with $\epsilon_\beta\ge2$, so that $w=p-|B^{\ge2}|$. Then
\[
W_k(a)=\pm\,\pi^{w}\,\frac{\prod_{\beta\in B^{\ge2}}\sin(\pi\epsilon_\beta a)}{\sin(\pi a)^{p}}\,R_k(a),
\]
with $R_k$ a rational function with rational coefficients, and for $n>n_1$ the residue of $W_k$ at $-n$ is
\[
\pm(-1)^{\varsigma n}\sum_{0\le l\le(w-1)/2}\tau_l\,\pi^{2l}\,\frac{R_k^{(w-1-2l)}(-n)}{(w-1-2l)!},
\]
with $\varsigma\equiv p+\sum_{\beta\in B^{\ge2}}\epsilon_\beta\pmod 2$, rational $\tau_l$ that depend only on the slopes, and $\tau_0=\prod_{\beta\in B^{\ge2}}\epsilon_\beta$. So a closing of Theorem~\ref{thm:ra-oneshare}(iii) is, up to finitely many terms in $\mathbb Q[\pi^2]$, a sum over the lattice of $(-1)^{\varsigma n}r(n)$ with $r$ a rational function of $n$ whose coefficients lie in $\mathbb Q[\pi^2]$, and Lemma~\ref{lem:polecount} over the field $\mathbb Q(\pi^2)$, applied to $m\mapsto r(n_1+1+m)$ when $\varsigma=0$ and to $m\mapsto r(n_1+1+2m)-r(n_1+2+2m)$ when $\varsigma=1$, writes it over $1$, the Trailmoments of the lattice inputs on the cosets of the poles at their cuts and differences of the index-one values of these inputs, with coefficients in $\mathbb Q[\pi^2]$; here $\pi^2=6\,\Tx{2}{\IZ}{0}$ (Theorem~\ref{thm:hparity}(iv)).
\end{proposition}

So the pairing of the lower Gammas, not their number, keeps the terms rational. The two-mark word $\Gst40{a,\ a,\ 1-a,\ 1-a}{}$ of Lemma~\ref{lem:resonantpair} has no lower Gamma; the word of two merged pairs on four marks, with one lower Gamma of slope $2$ and one of slope $-2$, pairs off, and the poles of its rational factor lie on the coset $\frac12+\mathbb Z$; the word of one merged pair on three marks, with one lower Gamma of slope $-2$, does not pair off, and its terms carry a live input of root rate $\frac14$. Exercise~\ref{ex:cfour-halflattice} works the word on four marks.

\section{The clearing along a ray}\label{sec:ra-clearing}

% New: proved in the last part.
\begin{theorem}[The frame bounds the clearing from below]\label{thm:ra-frameclearing}
Let $p_0,\dots,p_h\in\mathbb Q[k]$, $h\ge1$, with $p_0$ and $p_h$ without zeros at the integers $k\ge k_0$; let $d_0$, $d_h$ and $c$ be the degrees of $p_0$ and $p_h$ and the quotient of their leading coefficients.
\begin{enumerate}[label=(\roman*)]
\item Let rational sequences $x^{(1)},\dots,x^{(h)}$ on the integers $k\ge k_0$ satisfy $\sum_{j=0}^hp_j(k)\,x^{(i)}_{k+j}=0$ and be linearly independent, with frame $C_k$ (Proposition~\ref{prop:frame}). If a positive integer $D_k$ makes every $D_k\,x^{(i)}_{k+l}$, $1\le i\le h$, $0\le l<h$, an integer, then $D_k^{\,h}\,|C_k|\ge1$, and as $k\to\infty$
\[
\frac{\log D_k}k\ \ge\ \frac{d_h-d_0}{h}\,(\log k-1)+\frac1h\log\frac1{|c|}+O\Big(\frac{\log k}k\Big).
\]
When $d_0=d_h$ is the largest degree of the $p_j$, $\liminf_k\frac1k\log D_k\ge\frac1h\log\frac1{|c|}$, and $|c|$ is the modulus of the product of the roots of the symbol (Proposition~\ref{prop:ra-framegrowth}).
\item Let $x_k=\sum_{l=1}^mA_{k,l}v_l-B_k$, $k\ge k_0$, with rational $A_{k,l}$ and $B_k$ and real $v_l$ such that $1,v_1,\dots,v_m$ are linearly independent over $\mathbb Q$, satisfy $\sum_jp_j(k)\,x_{k+j}=0$. Then each of the $m+1$ sequences $(A_{k,l})_k$ and $(B_k)_k$ satisfies the same recurrence, and if $h$ of them are linearly independent, every $D_k$ clearing the rows of $x_k,\dots,x_{k+h-1}$ (Definition~\ref{def:row}) satisfies (i).
\end{enumerate}
\end{theorem}

The bound reads the clearing on $h$ consecutive rows together, through one determinant: the frame grows like the $k$-th power of the product of the roots of the symbol, and the clearing raised to the power $h$ must undo it. For the values $\langle1\rangle_k$ along a ray with $\kappa_{1,S}>0$ for every nonempty proper $S$, the roots of the symbol of every law of the values with rational coefficients include the native value $f(\alpha^\star)$ of the native gauge of the direction and its conjugates (Proposition~\ref{prop:nativeorder}). The frame does not say how the clearing is shared among the coordinates; on the rays of Section~\ref{sec:ra-ordertwo} the denominator of the ratio of a reflexive law falls on the coefficient of the one coordinate, at its full power up to what the closings absorb.

\section{The merit of a ray}\label{sec:ra-merit}

% New: proved in the last part.
\begin{definition}[The merit of rows along a ray]\label{def:ra-merit}
Let a ray of a family have real origin and $\kappa_{1,S}>0$ for every nonempty proper $S$, with the native gauge $\alpha^\star$ of its direction on $\mathbb R^{\cE}_{>0}$ (Theorem~\ref{thm:pencil}), and let its values have rows $\langle1\rangle_k=\sum_{l=1}^mA_{k,l}v_l-B_k$, $k\ge k_0$, over fixed real numbers $v_1,\dots,v_m$ (Definition~\ref{def:row}), cleared by positive integers $D_k$. The \emph{merit} of these rows is
\[
-\log f(\alpha^\star)-\limsup_{k\to\infty}\frac1k\log D_k\ \in[-\infty,\infty).
\]
\end{definition}

% New: proved in the last part.
\begin{theorem}[The merit decides the gap along a ray]\label{thm:ra-merit}
In Definition~\ref{def:ra-merit} let $k_1\ge k_0$ be an integer with $\kappa_{0,S}+k_1\kappa_{1,S}>0$ for every nonempty proper $S$.
\begin{enumerate}[label=(\roman*)]
\item For every $k\ge k_1$, $0<D_k\langle1\rangle_k\le D_k\,f(\alpha^\star)^{k-k_1}\,\langle1\rangle_{k_1}$.
\item If the merit is positive, then $D_k\langle1\rangle_k\to0$ and some $v_l$ is irrational; for $m=1$, $v_1\notin\mathbb Q$ and $A_{k,1}\neq0$ for all large $k$.
\item If $D_k\le C\gamma^k\,d_{\lfloor ak+b\rfloor}^{\,w}$ for all large $k$, with constants $C,\gamma>0$, $a,w\ge0$ and $b$, and $\lfloor ak+b\rfloor\ge0$, the merit is at least $-\log f(\alpha^\star)-\log\gamma-a\,w\,c_{\mathrm{lcm}}$ (Lemma~\ref{lem:lcm}), and $\gamma\,e^{a w c_{\mathrm{lcm}}}f(\alpha^\star)<1$ gives (ii).
\item If each $D_k$ also clears the rows of $\langle1\rangle_{k+1},\dots,\langle1\rangle_{k+h-1}$, $1,v_1,\dots,v_m$ are linearly independent over $\mathbb Q$, the values satisfy a law $\sum_{j=0}^hp_j(k)\langle1\rangle_{k+j}=0$, $k\ge k_0$, as in Theorem~\ref{thm:ra-frameclearing} with $p_0$ and $p_h$ of the largest degree, and $h$ of the sequences $(A_{k,l})_k$, $(B_k)_k$ are linearly independent, then the merit is at most $-\log f(\alpha^\star)+\frac1h\log|c|$, with $|c|$ the modulus of the product of the roots of the symbol.
\end{enumerate}
\end{theorem}

On the centre ray of a family, $f(\alpha^\star)=\prod_jK_j(\one)^{-2\sigma_j}$ (Theorem~\ref{thm:branch}), and a clearing $D_k=\gamma^{k-1}d_{k-1}^{\,w}$ passes the gap when $\gamma\,e^{wc_{\mathrm{lcm}}}<\prod_jK_j(\one)^{2\sigma_j}$, the minimum of the reciprocal stride on the orthant. Only the upper bound of the values is used: positivity and the largest value of the stride on the orthant, both read before any closing. By (iv), with the native value among the roots of the symbol, a law of order $h$ leaves positive merit only to a ray whose symbol has roots of modulus larger than $f(\alpha^\star)$.

\section{Laws at the low end and their closings}\label{sec:ra-ordertwo}

% New: proved in the last part.
\begin{proposition}[The low end of the law]\label{prop:ra-lowend}
Let a ray be as in Proposition~\ref{prop:principal}, on $n$ marks, with origin $0$, every letter of positive exponent in the direction, $\kappa_{1,S}>0$ for every nonempty proper $S$, and the volume $V_1$ of the direction; let the law of $[1]$ (Definition~\ref{def:raymodule}) be $\sum_{i=0}^hp_i(k)\operatorname{Sh}^i[1]=0$ with $p_0,\dots,p_h\in\mathbb Q[k]$, and write $x_k=\langle1\rangle_k$.
\begin{enumerate}[label=(\roman*)]
\item The values satisfy the law at every integer $k\ge1$. Further $k^{n-1}$ divides $p_0$, and with $\varpi$ the coefficient of $k^{n-1}$ in $p_0$ the values satisfy the \emph{low-end relation}
\[
\varpi\,V_1+\sum_{i=1}^hp_i(0)\,x_i=0 ;
\]
$V_1$ is rational. For a solution $y$ of the law on the integers $k\ge1$, $\sum_ip_i(k)y_{k+i}=0$ for every $k\ge1$, call $-\sum_{i=1}^hp_i(0)\,y_i$ its \emph{low-end coefficient}; that of the values is $\varpi V_1$, the value at $k=0$ of the continuation of $p_0(k)\langle1\rangle_k$.
\item Let $h=2$ and $p_2(k)\neq0$ at every integer $k\ge0$, and let $c$ and $r$ be the solutions of the law on $k\ge1$ with $(c_1,c_2)=\big(1,-p_1(0)/p_2(0)\big)$ and $(r_1,r_2)=\big(0,-\varpi V_1/p_2(0)\big)$. They are rational, the low-end coefficient of $c$ vanishes, that of $r$ is $\varpi V_1$, and
\[
x_k=x_1\,c_k+r_k\qquad(k\ge1):
\]
every value has a row in the one coordinate $x_1$, its coefficient the solution whose low-end coefficient vanishes and whose first term is $1$.
\end{enumerate}
\end{proposition}

% New: proved in the last part.
\begin{definition}[Reflexive laws]\label{def:ra-reflexive}
A recurrence $\sum_{i=0}^hp_i(k)\,x_{k+i}=0$ with $p_i\in\mathbb Q[k]$ is \emph{reflexive} with the \emph{ratio} $\lambda\in\mathbb Q^\times$ when there is a constant $\gamma\neq0$ with
\[
p_{h-i}(-k-h+1)=\gamma\,\lambda^i\,p_i(k)\qquad(0\le i\le h)
\]
as polynomials in $k$. When $p_0$ and $p_h$ have the largest degree $d$ of the $p_i$, the roots of the symbol of a reflexive recurrence are exchanged by $\varrho\mapsto\lambda/\varrho$, and at $h=2$ their product is $\pm\lambda$.
\end{definition}

% New: proved in the last part.
\begin{theorem}[The closings read the coordinate of a reflexive law]\label{thm:ra-closingcoordinate}
Let a ray be as in Proposition~\ref{prop:ra-lowend}, with the law of $[1]$ reflexive with the ratio $\lambda$, and let $\mathsf C_j(1)=\CT(f^{-j})$ be the closings of $1$ (Definition~\ref{def:closings}).
\begin{enumerate}[label=(\roman*)]
\item $y_k=\lambda^{k-1}\,\mathsf C_{k-1}(1)$, $k\ge1$, is a solution of the law on $k\ge1$, and its low-end coefficient vanishes.
\item If $h=2$ and $p_2(k)\neq0$ at every integer $k\ge0$, then $y$ is the solution $c$ of Proposition~\ref{prop:ra-lowend}(ii):
\[
\langle1\rangle_k=\lambda^{k-1}\,\mathsf C_{k-1}(1)\,\langle1\rangle_1+r_k\qquad(k\ge1),
\]
the coefficient of the coordinate read from the closings of the ray, with no lattice sum. If moreover $\langle1\rangle_1=A_1v-B_1$ is a row in one coordinate $v$, the coefficient of $v$ in $\langle1\rangle_k$ is $A_1\lambda^{k-1}\mathsf C_{k-1}(1)$.
\item If moreover the weights are integers, $\mathsf C_j(1)\in\mathbb Z$, and with $\lambda=a/b$ in lowest terms, $b>0$, the coefficient of $\langle1\rangle_1$ in $\langle1\rangle_k$ has the denominator $b^{k-1}/\gcd\big(b^{k-1},\mathsf C_{k-1}(1)\big)$.
\end{enumerate}
\end{theorem}

% New: proved in the last part.
\begin{corollary}[The merit of a reflexive ray of order two]\label{cor:ra-meritbound}
In Theorem~\ref{thm:ra-closingcoordinate}(iii), every clearing $D_k$ of the rows $\langle1\rangle_k=\lambda^{k-1}\mathsf C_{k-1}(1)\,\langle1\rangle_1+r_k$ in the coordinate $\langle1\rangle_1$ has merit (Definition~\ref{def:ra-merit}) at most
\[
-\log f(\alpha^\star)-\log b+\liminf_{k\to\infty}\frac1k\log\gcd\big(b^{k-1},\mathsf C_{k-1}(1)\big).
\]
\end{corollary}

The closings are the constant terms of the powers of the reciprocal stride, finite sums of products of weights; on a reflexive ray of order two they are the coefficient of the coordinate at every step, up to the power $\lambda^{k-1}$, and the clearing pays the full power $b^{k-1}$ of the denominator of the ratio except for what the closings absorb. The ray is then judged by the divisibility of its closings. The centre rays of the lists $\binom{[3]}1$ and $\binom{[4]}1$ have reflexive laws of order two with the ratios $\frac19$ and $\frac1{64}$; the exercises derive both laws (Exercises~\ref{ex:smallc-box} and~\ref{ex:cfour-law}), and on four marks the rows, the coefficient of the coordinate and the clearing (Exercises~\ref{ex:cfour-row}, \ref{ex:cfour-closing} and~\ref{ex:cfour-clearing}).

\section{The hinge closing as a row over the room}\label{sec:ra-hinge}

% New: proved in the last part.
\begin{proposition}[The hinge closing as a row over the room]\label{prop:ra-hingerow}
For one table in Theorem~\ref{thm:twosides} let $\sup_{\Delta_B}\varrho_B<\inf_{\Delta_R}\varrho_R$, so that the even share closes on the lattices of the level $d$ (Corollary~\ref{cor:closetwosides}(i)). Group the residues by the index $n$ of their points and write the closed value as $P_F(\mvec;s,t)=\sum_{n\ge0}\ip{\Phi_n}{\mathcal X_n}$, where $\mathcal X_n$ lists the data of the room family that the residues at the $n$-th points contain, its periods at the level shares of those points and, where the two lattices meet, their derivatives in the even share (Corollary~\ref{cor:closetwosides}(ii)), and $\Phi_n$ lists the coefficients that Corollary~\ref{cor:closetwosides}(ii) puts in front of them.
\begin{enumerate}[label=(\roman*)]
\item Let there be numbers $\chi_1,\dots,\chi_h$ and lists $\mathbf c_1(n),\dots,\mathbf c_h(n)$ with $\mathcal X_n=\sum_{i=1}^h\chi_i\,\mathbf c_i(n)$ for every $n$ and $\sum_n|\ip{\Phi_n}{\mathbf c_i(n)}|<\infty$ for every $i$. Then
\[
P_F(\mvec;s,t)=\sum_{i=1}^h\chi_i\,\Xi_i,\qquad \Xi_i=\sum_{n\ge0}\ip{\Phi_n}{\mathbf c_i(n)}:
\]
a row over $h$ data of the room, with coefficients that are sums along the lattice of the bunch's data against the sequences $\mathbf c_i$.
\item When $2\delta^d\notin\mathbb Z$ every pole is simple and $\mathcal X_n$ lists the two room periods $x_n^\pm$ at the points $\sigma^d=\pm\delta^d-n$. If each sequence $x^\pm$ satisfies a recurrence $\sum_{i=0}^{h_\pm}p^\pm_i(n)\,x^\pm_{n+i}=0$, $n\ge0$, with $p^\pm_{h_\pm}(n)\neq0$ for $n\ge0$, let $c^\pm_i$, $0\le i<h_\pm$, be its solutions with the initial values $c^\pm_i(l)=\delta_{il}$, $0\le l<h_\pm$, rational when the $p^\pm_i$ have rational coefficients. Then $x^\pm_n=\sum_ix^\pm_i\,c^\pm_i(n)$ for every $n$, and whenever the sums $\sum_n\Phi^\pm_n\,c^\pm_i(n)$ converge absolutely, (i) holds with the $h_++h_-$ initial periods $x^\pm_0,\dots,x^\pm_{h_\pm-1}$.
\item When $2\delta^d\in\mathbb Z$ the poles lie on the one lattice $\sigma^d=|\delta^d|-n$, $n\in\mathbb N$. Let the room period at $\sigma^d=\sigma-n$, a function $x_n(\sigma)$ of $\sigma$ near $\sigma_0=|\delta^d|$, satisfy a recurrence $\sum_{i=0}^hp_i(n,\sigma)\,x_{n+i}(\sigma)=0$ whose coefficients are polynomials in $n$ and $\sigma$, with $p_h(n,\sigma_0)\neq0$ for $n\ge0$. Then the values $x_n(\sigma_0)$ and the derivatives $x'_n(\sigma_0)$ satisfy the block-triangular recurrence of Proposition~\ref{prop:jetrecurrence}, they are combinations of the $2h$ initial data $x_i(\sigma_0)$, $x'_i(\sigma_0)$, $i<h$, with coefficients that do not depend on the room's periods, and whenever the resulting sums converge absolutely (i) holds with these $2h$ data.
\end{enumerate}
The same holds with the bunch and the room exchanged when $\inf_{\Delta_B}\varrho_B>\sup_{\Delta_R}\varrho_R$, the share closing on the level $r$.
\end{proposition}

Read down a tree of bunches, a hinge closing writes a period as a row over finitely many periods of its room, the coefficients being sums along the lattice of the hinge of the bunch's data against the fundamental solutions of the room's recurrence. The room's periods along the lattice are the moments of one ratio against one form (Corollary~\ref{cor:closetwosides}(ii)); when they have tier one, each lattice contributes one room period times one sum along the lattice. Proposition~\ref{prop:ct-tiertwo} gives a condition on the roots of the symbols under which the sums converge.

\section{Critical primes}\label{sec:ra-critical}

% New: proved in the last part.
\begin{definition}[Critical primes]\label{def:ra-criticalprime}
Let a ray of a family have an integral direction with the stride $f=\alpha^{\mvec_1}\prod_LL^{-w_L}$, the product over the distinct letters $L$ of positive exponent $w_L$, and let its weights lie in a number field $\Bbbk$ that contains the coordinates of a native gauge $\alpha_c$ (Definition~\ref{def:raygauge}) on the chart $\alpha_{e_0}=1$. A prime $\mathfrak p$ of $\Bbbk$ at which the weights are integral and the coordinates of $\alpha_c$ are units is a \emph{critical prime} of $\alpha_c$ when the native value is not integral there, $v_{\mathfrak p}\big(f(\alpha_c)\big)<0$.
\end{definition}

% New: proved in the last part.
\begin{theorem}[The letters degenerate at a critical prime]\label{thm:ra-degenerate}
Let $\mathfrak p$ be a critical prime of a native gauge $\alpha_c$, with the residue field $\kappa$ and the reduction $\bar\alpha_c\in(\kappa^\times)^{\cE}$, and let $S$ be the set of the letters of positive exponent whose valuation $v_{\mathfrak p}(L(\alpha_c))$ is the largest, a positive number.
\begin{enumerate}[label=(\roman*)]
\item Every letter of $S$ vanishes at $\bar\alpha_c$ after reduction, and there are units $u_L$ of the local ring at $\mathfrak p$ with
\[
\sum_{L\in S}\overline{w_Lu_L}\;\nabla\bar L(\bar\alpha_c)=0,
\]
$\nabla$ the gradient in the coordinates $\alpha_e$, $e\in\cE$, and the bar the reduction.
\item If the characteristic of $\kappa$ divides none of the exponents $w_L$, $L\in S$, the zero sets of the reduced letters of $S$ are not smooth hypersurfaces meeting transversally at $\bar\alpha_c$: when $S$ has one letter, that letter is singular at $\bar\alpha_c$, and when it has several, their gradients there are linearly dependent.
\end{enumerate}
\end{theorem}

% New: proved in the last part.
\begin{proposition}[Critical primes and the symbol]\label{prop:ra-symbolprimes}
Let $\sigma(\lambda)=\sum_{j=0}^hc_j\lambda^j\in\mathbb Z[\lambda]$, $c_h\neq0$, $h\ge1$, be a symbol (Corollary~\ref{cor:symbol}) with coprime coefficients, and $K$ a number field containing its roots.
\begin{enumerate}[label=(\roman*)]
\item For every root $\lambda$ of $\sigma$, $c_h\lambda$ is an algebraic integer.
\item A prime $p$ divides $c_h$ exactly when some root of $\sigma$ has negative valuation at some prime of $K$ above $p$; and then at every prime of $K$ above $p$ some root has negative valuation.
\item Let every root of $\sigma$ be the native value of a native gauge whose coordinates lie in $K$, at which the weights are integral and the coordinates are units at the primes above $p$. Then $p$ divides $c_h$ exactly when some prime of $K$ above $p$ is a critical prime of one of these native gauges. When $\sigma$ is the symbol of a recurrence whose $p_0$ and $p_h$ have the largest degree, $c_0/c_h$ is the quotient of their leading coefficients, whose $k$-th power enters the frame (Proposition~\ref{prop:frame}).
\end{enumerate}
\end{proposition}

At a critical prime the native gauge reduces to a point at which the letters of largest valuation meet badly, so the critical primes of a ray are read from the stride and its critical points before any value is computed, and under (iii) they are the primes of the leading coefficient of the symbol, which the frame carries. The value span of a ray can see fewer primes: on the centre ray of $\binom{[4]}1$ the critical prime $2$ sits in the frame and in the denominators of the coefficients, and the coordinate is the Trailmoment $\Tx{3}{\IZ}{0}$ of the counting input. On the centre ray of $\binom{[3]}1$ the native value $\frac19$ reduces modulo $3$ to a tangency of the letters, and on that of $\binom{[4]}1$ the native value $\frac1{16}$ reduces modulo $2$ to the collapse of three intersection points of the letters onto the native gauge.

\part{Relations and the content}\label{part:relations}

\chapter{Relations and separation}\label{ch:relations}

A value is the evaluation of a word, and two words of equal value either are related by the proved relations of the theory or are not, a problem that relation blocks make finite. The admitted words form an algebra with an evaluation, the proved relations generate an ideal inside its kernel, and completeness asks for equality. Relation blocks reduce a finite space of words to representatives, each word with a certificate; when the representatives have values independent over the rationals, the block is complete. Words of polynomials on a box are complete outright, a finite census of insertions separates numerators, and the critical-line words of order $k$ of polynomials form a block whose values are the rational combinations of $1$ and $\Tx{2k+1}{\IZ}{0}$; its two representatives have values independent over the rationals exactly when $\Tx{2k+1}{\IZ}{0}\notin\mathbb Q$, and the block is then complete. Independence of values is proved by separators: families of rational matrices whose rows are small on the vector of values.

\section{Words and relations}

% Reference 20: def:words; source line 11163
\begin{definition}[Words, evaluation, relations, completeness]\label{def:words}
An \emph{admitted word} is a word (Definition~\ref{def:conv}(vi)) given by finite rational data (tables and weights, coefficients, masses, exponents, inputs with a finite presentation, cuts) together with the rule that assigns its value: a convergent integral, a finite part or a continuation, as proved for it. The admitted words span a rational algebra $\mathfrak W$, whose product tensors the data, with independent integration variables and separate regulators, and \emph{evaluation} is the map $\operatorname{ev}:\mathfrak W\to\mathbb C$. The \emph{relation ideal} $\mathfrak I\subseteq\mathfrak W$ is generated by the proved relations, each on its stated domain: linearity; changes of variables and of order; scale and shape; split, port and gluing; the Gramma relations; the letter relations and the mark relations (Theorems~\ref{thm:confJ} and~\ref{thm:markrelation}); the subtraction of a common boundary term from the two sides of a relation; the cut relation of the Trailmoment (Proposition~\ref{prop:cutlaw}); the closings. A set of words is \emph{complete} when each of its rational combinations of value zero lies in $\mathfrak I$.
\end{definition}

% Reference 20: prop:regulators; source line 11167
\begin{proposition}[Products keep separate regulators]\label{prop:regulators}
If $A(\varepsilon)$ and $A'(\varepsilon')$ have Laurent expansions in disjoint sets of regulators, then $\CT_\varepsilon\CT_{\varepsilon'}\big(A(\varepsilon)A'(\varepsilon')\big)=\big(\CT_\varepsilon A\big)\big(\CT_{\varepsilon'}A'\big)$, and evaluation is multiplicative on $\mathfrak W$. Identifying regulators is a further operation: $\CT(\varepsilon^{-1})=\CT(\varepsilon)=0$, while $\CT(\varepsilon^{-1}\varepsilon)=1$.
\end{proposition}

% Reference 20: prop:sound; source line 11174
\begin{proposition}[Soundness]\label{prop:sound}
$\mathfrak I\subseteq\ker(\operatorname{ev})$.
\end{proposition}

% Reference 20: prop:relationmatrix; source line 11181 (one letter).
% Proved in the last part: prop:relationmatrix.
\begin{proposition}[A finite matrix of mark relations]\label{prop:relationmatrix}
On $[0,1]^d$ let $a_l,b_l>0$, let $D_1,\dots,D_h$ be rational polynomials positive on the closed box with exponents $p_1,\dots,p_h\in\mathbb Q$, and in this statement write $D=\prod_kD_k=\sum_\gamma D_\gamma x^\gamma$ and $x_i\sum_kp_k\,\partial_iD_k\prod_{k'\neq k}D_{k'}=\sum_\gamma D^{(i)}_\gamma x^\gamma$. For every polynomial $A$ and every coordinate $i$,
\[
\operatorname{mr}_i(A)=D\big[x_i(1-x_i)\,\partial_iA+\big(a_i-(a_i+b_i)x_i\big)A\big]-x_i(1-x_i)\,A\sum_kp_k\,\partial_iD_k\prod_{k'\neq k}D_{k'}
\]
has $\int_{[0,1]^d}\operatorname{mr}_i(A)\prod_lx_l^{a_l-1}(1-x_l)^{b_l-1}\prod_kD_k^{-p_k-1}\,dx=0$ (Theorem~\ref{thm:markrelation}(a) with the letters $D_k$ at the exponents $p_k+1$ and the polynomial $AD$), and
\[
\operatorname{mr}_i(x^\beta)=\sum_\gamma\Big[\big((\beta_i+a_i)D_\gamma-D^{(i)}_\gamma\big)x^{\beta+\gamma}-\big((\beta_i+a_i+b_i)D_\gamma-D^{(i)}_\gamma\big)x^{\beta+\gamma+\mathbf e_i}\Big].
\]
For a degree bound these vectors form a finite rational relation matrix; a rational solution of its column equation for a numerator writes the numerator as a sum of polynomials $\operatorname{mr}_i(A_i)$, so that its word vanishes.
\end{proposition}

% Reference 20: thm:betastage; source line 11196
\begin{theorem}[Words of polynomials on a box are complete]\label{thm:betastage}
For $P\in\mathbb Q[x]$ and rational $a_i,b_i>0$ the words $\int_{[0,1]^d}P(x)\prod_ix_i^{a_i-1}(1-x_i)^{b_i-1}\,dx$ have the mark relations $x^{\beta+\mathbf e_i}\equiv\frac{\beta_i+a_i}{\beta_i+a_i+b_i}\,x^\beta$, and every $P$ reduces to
\[
\sum_\beta\big([x^\beta]P\big)\prod_i\frac{(a_i)_{\beta_i}}{(a_i+b_i)_{\beta_i}},
\]
the expectation of $P$ under the Dirichlet laws $\Dir_{(a_i,b_i)}$ of the pairs $(x_i,1-x_i)$, times the base word, of value $\prod_i\Gst21{a_i,\ b_i}{a_i+b_i}>0$. A word vanishes exactly when its reduction is zero. The same holds for every finite tensor product of such words.
\end{theorem}

% Reference 20: thm:census; source line 11203
\begin{theorem}[A finite insertion census separates numerators]\label{thm:census}
Let $\phi$ be a measure on $[0,1]^d$ with an integrable density positive on the interior, and $D$ a rational polynomial positive on the closed box. For $|\beta|,|\gamma|\le h$ the matrix $\mathbf G_{\beta\gamma}=\int x^{\beta+\gamma}D^{-1}\,d\phi$ is positive definite. For $A=\sum_\beta a_\beta x^\beta$ of degree at most $h$ the insertion values $v_\beta=\int x^\beta AD^{-1}\,d\phi$ satisfy $\mathbf v=\mathbf G\mathbf a$ and $\operatorname{adj}(\mathbf G)\,\mathbf v=\det(\mathbf G)\,\mathbf a$: these $\binom{d+h}d$ values determine $A$. Equivalently, the Laplace transform $t\mapsto\int AD^{-1}e^{-\ip tx}\,d\phi$ has all its Taylor coefficients of degree at most $h$ zero at $t=0$ only for $A=0$.
\end{theorem}

Evaluation is multiplicative because a product of words keeps the regulators of its factors apart; identifying two regulators is a new operation, not a relation. Every relation the theory proves lies in the kernel of evaluation, and completeness asks whether the kernel holds anything more. On a box against positive letters the mark relations form a finite rational matrix, so that vanishing is certified by linear algebra; for polynomials on a box without a letter the relations reduce every insertion to one rational multiple of the base word, and with a letter the values of finitely many monomial insertions, through the positive definite matrix $\mathbf G$ of Theorem~\ref{thm:census}, recover the insertion itself.

\section{Relation blocks}

% Reference 20: def:relblock; source line 11212
\begin{definition}[Relation block]\label{def:relblock}
A \emph{relation block} is a finite-dimensional rational space $V$ of admitted words with a matrix $C$ whose columns are proved relations in $V$, each with its derivation. Rational row reduction gives representatives $\mathfrak w_0,\dots,\mathfrak w_s$ of $V/\operatorname{im}C$ and for every $\mathfrak w\in V$ a certificate $\mathfrak w=C\mathbf v+\sum_jc_j\mathfrak w_j$.
\end{definition}

% Reference 20: prop:blockseparation; source line 11216
\begin{proposition}[Separation completes a block]\label{prop:blockseparation}
If $\operatorname{ev}(\mathfrak w_0),\dots,\operatorname{ev}(\mathfrak w_s)$ are linearly independent over $\mathbb Q$, then $\ker(\operatorname{ev}|_V)=\operatorname{im}C$.
\end{proposition}

% Reference 20: thm:exhaustion; source line 11223
\begin{theorem}[Completeness from separated blocks]\label{thm:exhaustion}
If every finite set of words of a subalgebra $\mathfrak W'\subseteq\mathfrak W$ lies in a relation block whose representatives have values linearly independent over $\mathbb Q$, then $\mathfrak W'$ is complete, and each of its words of value zero has a finite derivation, its certificate.
\end{theorem}

% Reference 20: prop:tensorrelations; source line 11230
\begin{proposition}[Tensoring relation blocks]\label{prop:tensorrelations}
For blocks $(V_1,C_1)$ and $(V_2,C_2)$ the induced relations on $V_1\otimes V_2$ have the matrix $\big[\,C_1\otimes I\ \big|\ I\otimes C_2\,\big]$, with quotient $(V_1/\operatorname{im}C_1)\otimes(V_2/\operatorname{im}C_2)$, and the values of the representatives are the products of the values (Proposition~\ref{prop:regulators}).
\end{proposition}

% Reference 20: prop:criticalstage; source line 11237 (order two).
% Proved in the last part: prop:criticalstage.
\begin{proposition}[The critical polynomial block]\label{prop:criticalstage}
For an integer $k\ge1$ and $P\in\mathbb Q[x]$ put $Q_P(u)=\frac{P(2u)-P(1)}{1-2u}$ and $h(P)=\frac2{(2k)!}\,\oE_{Y\sim\mathbb P_{\mathrm{crit}}}\big[Q_P^{(2k-1)}(\frac12+iY)\big]$. Then $h(P)\in\mathbb Q$ and
\[
\operatorname{crit}_k[P(2u)]=h(P)+(2^{2k+1}-1)\,P(1)\,\Tx{2k+1}{\IZ}{0}.
\]
So the values of the rational constants and of the critical-line words of order $k$ of polynomials factor through the surjective coordinate map $(c,P)\mapsto\big(c+h(P),(2^{2k+1}-1)P(1)\big)\in\mathbb Q^2$ followed by $(a,b)\mapsto a+b\,\Tx{2k+1}{\IZ}{0}$, which is injective exactly when $\Tx{2k+1}{\IZ}{0}\notin\mathbb Q$; moving the cut of the Trailmoment subtracts a rational head.
\end{proposition}

A relation block is the finite form of completeness: inside a finite space of words, rational row reduction by the proved relations leaves representatives and certificates, and independence of the values of the representatives over the rationals makes the block complete (Proposition~\ref{prop:blockseparation}). Blocks tensor, their quotients tensoring and the values of their representatives multiplying, so a product of blocks is complete as soon as the products of those values are independent over the rationals. The critical-line words of order $k$ of polynomials have a block whose quotient has dimension at most two: every such word is a rational number plus a rational multiple of $\Tx{2k+1}{\IZ}{0}$, read off two coordinates of the polynomial, and the values of the two representatives are independent over the rationals exactly when that Trailmoment is irrational, the block being then complete.

\section{Separators}

% Reference 20: thm:exterior; source line 11250
\begin{theorem}[A primitive exterior separator]\label{thm:exterior}
Let $v=(v_1,\dots,v_s,1)\in\mathbb R^{s+1}$. For a rational $s\times(s+1)$ matrix $A$ of rank $s$ with rows $r_i$, let $w\in\mathbb Q^{s+1}$ be its cofactor vector, $\det[x;A]=\ip xw$ for all $x$, and $z=Mw$ its primitive integral multiple, $M>0$. If $D_i$ clears $r_i$, then $M=\prod_iD_i/c$ with $c$ the greatest common divisor of the maximal minors of $\operatorname{diag}(D_i)A$. Put $y_i=\ip{r_i}v$, let $\hat A_i$ be the other rows restricted to the first $s$ columns, and
\[
\beta(A)=M\sum_i|y_i|\sum_{J\subseteq\{1,\dots,s\},\ |J|=s-1}\big|\det\hat A_i[:,J]\big| .
\]
If for all large $n$ finitely many such matrices $A^{(n)}_1,\dots,A^{(n)}_{k_n}$ have primitive cofactor vectors spanning $\mathbb Q^{s+1}$ and $\max_l\beta(A^{(n)}_l)\to0$, then $1,v_1,\dots,v_s$ are linearly independent over $\mathbb Q$.
\end{theorem}

% Reference 20: cor:amplified; source line 11261
\begin{corollary}[Separation with bounded integer entries]\label{cor:amplified}
If integer matrices $A^{(n)}_l$ with spanning cofactor vectors have entries at most $H_n\ge1$ and $|\ip{r_i}v|\le\epsilon_n$ for every row, with $H_n^{s-1}\epsilon_n\to0$, then $1,v_1,\dots,v_s$ are linearly independent over $\mathbb Q$. In rates: if rational matrices with spanning cofactor vectors have entries at most $e^{an+o(n)}$, a common clearing integer at most $e^{bn+o(n)}$ and row values at most $e^{-cn+o(n)}$, with $c>sb+(s-1)a$, the same conclusion holds.
\end{corollary}

% Reference 20: cor:consecutive; source line 11268
\begin{corollary}[Consecutive rows]\label{cor:consecutive}
Let rows $r_k\in\mathbb Q^{s+1}$ be such that the square matrix of the rows $r_k,\dots,r_{k+s}$ is invertible for large $k$. Then for each such $k$ the $s+1$ matrices omitting one of these rows have spanning primitive cofactor vectors, and Theorem~\ref{thm:exterior} applies with these matrices. If the coordinate sequences of the rows satisfy one recurrence of order $s+1$ whose extreme coefficients have no integer zeros from $k_0$ on, invertibility of that matrix at $k_0$ gives it for every $k\ge k_0$.
\end{corollary}

% Reference 20: thm:pin; source line 11275
\begin{theorem}[Function level: pole pinning]\label{thm:pin}
Let $f$ and $g$ be meromorphic on a connected open subset of $\mathbb C$. If $f$ has infinitely many poles at which $g$ is holomorphic, and the values of $g$ at those poles include infinitely many distinct numbers, then no nonzero polynomial $Q\in\mathbb C[X,Y]$ has $Q(f,g)\equiv0$.
\end{theorem}

A separator is a family of rational matrices whose rows are nearly orthogonal to the vector of values: a primitive integral multiple of the cofactor vector of each matrix pairs with any integer relation to an integer, and if the rows are small enough against the clearing and the bounds on the entries, no rational relation among the values survives. Consecutive rows of a recurrence give such matrices as soon as one square matrix of $s+1$ consecutive rows is invertible. At the level of functions, a function with infinitely many poles at which a second function takes infinitely many values satisfies no polynomial relation with it.

\chapter{The content and its coaction}\label{ch:content}

The content is the algebra of the trails of lattice inputs at rational frequencies, on rational cosets and at rational cuts, read modulo the proved relations. At the level $M$ it is spanned by the roots of unity of order $M$ times the trails of the counting lattice at the cut $0$ whose frequencies have the denominator $M$; it is closed under products, it holds every trail with rational data at the level that its denominators fix, and it holds the differences of the index-one values of lattice inputs on rational cosets, so the sums along the lattices of paired words lie in it. Its values are the chains of level $M$: the iterated integrals from $0$ to $1$ whose letters are $0$ and the roots of unity of order $M$, each a trail whose frequencies are the ratios of its consecutive nonzero letters; read between $0$ and a root of unity in either order a chain is again a value of the content. The cutting rule keeps a subsequence of the letters of a chain and reads every run of the letters it drops between the two letters that bound it; on the symbols of chains it is a coassociative coaction, graded by the weight and compatible with the rotations and the reversal of chains. A live input has a letter too, its generating function, which satisfies a linear differential equation singular at the reciprocals of the roots of the symbol of its recurrence; the letter of a product of two live inputs continues across every point that is not a product of singular points of the letters of the factors, and the row of a hinge closing is the room's periods against the letters of the live inputs built on the room's fundamental solutions.

\section{The content}\label{sec:ct-content}

% New: proved in the last part.
\begin{definition}[The content]\label{def:ct-content}
Let $M\ge1$ be an integer and $\mu_M$ the group of the roots of unity of order dividing $M$. A \emph{trail of level $M$} is the trail (Definition~\ref{def:trail}) at the cut $0$ of a lattice input on $\mathbb Z^d$ at frequencies in $\frac1M\mathbb Z$, with positive integer indices and no constant columns,
\[
T=\big(\sT_{\boldsymbol\nu;I_d}[y]\ e^{2\pi i\ip{\mathbf a}{y}}\,\Ilat{\mathbb Z^d}\big)\big(\{0<y_1<\cdots<y_d\}\big),\qquad d\ge1,\ \boldsymbol\nu\in\{1,2,\dots\}^d,\ \mathbf a\in\big(\tfrac1M\mathbb Z\big)^d,
\]
with $e^{2\pi ia_d}\neq1$ when $\nu_d=1$, read as the limit of its heads over $\{0<y_1<\cdots<y_d\le N\}$ as $N\to\infty$, which exists and is the Trailmoment itself when $\nu_d\ge2$ (Theorem~\ref{thm:ct-trailchain}); $|\boldsymbol\nu|$ is its \emph{weight} and $d$ its \emph{depth}. The \emph{content of level $M$}, $\mathcal C^{(M)}$, is the $\mathbb Q$-span in $\mathfrak W/\mathfrak I$ (Definition~\ref{def:words}) of the classes of the words $\zeta$ and $\zeta T$, for $\zeta\in\mu_M$, read as $e^{2\pi ir/M}$, and $T$ a trail of level $M$; $\mathcal C^{(M)}_w$ is the span of those with $T$ of weight at most $w$, and the \emph{content} is $\mathcal C=\bigcup_M\mathcal C^{(M)}$.
\end{definition}

The content is not the content $\mathcal C_M$ of a childcule $M$ in Definition~\ref{def:contents}, a $\mathcal Z$-module of values of periods; the level is written as an upper index to keep the two apart. The algebra $\mathcal Z$ of Definition~\ref{def:contents} maps onto the subalgebra of $\mathcal C^{(1)}$ generated by its trails of depth one, the Trailmoments $\Tx{k}{\IZ}{0}$, $k\ge2$ (Proposition~\ref{prop:ct-algebra}(i)).

% New: proved in the last part.
\begin{proposition}[The content is an algebra]\label{prop:ct-algebra}
\begin{enumerate}[label=(\roman*)]
\item $\mathcal C^{(M)}$ contains the classes of $1$ and of every product of Trailmoments $\Tx{k}{\IZ}{0}$, $k\ge2$, so it contains the image of $\mathcal Z$ (Definition~\ref{def:contents}); and $\mathcal C^{(M)}\subseteq\mathcal C^{(M')}$ when $M\mid M'$.
\item Let $T$ and $T'$ be trails of level $M$ with the data $(\nu_i,a_i)_{i\le d}$ and $(\nu'_j,a'_j)_{j\le d'}$. For a surjection $\phi$ from the disjoint union $\{1,\dots,d\}\sqcup\{1',\dots,d'\}$ onto $\{1,\dots,e\}$, strictly increasing on each of the two parts, let $T_\phi$ be the trail of level $M$ of depth $e$ whose $r$-th point has the index $\sum_{\phi(x)=r}\nu_x$ and the frequency $\sum_{\phi(x)=r}a_x$. Then
\[
TT'=\sum_\phi T_\phi,
\]
the sum over these surjections, of depths between $\max(d,d')$ and $d+d'$, every $T_\phi$ of the weight $|\boldsymbol\nu|+|\boldsymbol\nu'|$.
\item So $\mathcal C^{(M)}$ is a $\mathbb Q$-algebra with $\mathcal C^{(M)}_w\,\mathcal C^{(M)}_{w'}\subseteq\mathcal C^{(M)}_{w+w'}$, and $\mathcal C$ is a $\mathbb Q$-algebra.
\end{enumerate}
\end{proposition}

% New: proved in the last part.
\begin{proposition}[Cosets, cuts and index one]\label{prop:ct-level}
Let $M,M'\ge1$ be integers.
\begin{enumerate}[label=(\roman*)]
\item Let a trail have a lattice input on $\prod_{i=1}^d(b_i+\mathbb Z)$ with $b_i\in\frac1M\mathbb Z$, at frequencies $a_i\in\frac1{M'}\mathbb Z$, with positive integer indices, no constant columns and a cut $N\in\frac1M\mathbb Z$, $N\ge0$, and with $e^{2\pi ia_d}\neq1$ when $\nu_d=1$, read as the limit of its heads. Then its class lies in $\mathcal C^{(MM')}$, and the class of every head of it lies in the span of the classes of $\mu_{MM'}$; its value is a combination, with coefficients in $\mathbb Q(\mu_{MM'})$, of $1$ and the trails of level $MM'$ of depths at most $d$.
\item For $b,b'\in\frac1M\mathbb Z$ and cuts $N,N'\in\frac1M\mathbb Z$, $N,N'\ge0$, the difference $\Tx{1}{\Ilat{b+\mathbb Z}}{N}-\Tx{1}{\Ilat{b'+\mathbb Z}}{N'}$ of index-one values (Definition~\ref{def:lattice}) has its class in $\mathcal C^{(M)}$.
\end{enumerate}
\end{proposition}

So the sum along a lattice of a paired word at integer data (Proposition~\ref{prop:ra-paired}), a combination of Trailmoments of lattice inputs on rational cosets at their cuts and differences of their index-one values with coefficients in $\mathbb Q[\Tx{2}{\IZ}{0}]$, lies in the content. The values met on the centre rays of the lists of three and four marks are trails of depth one: at the frequency $\frac13$ and the index $2$, and on the coset $\frac12$ at the index $3$, which the coset relation returns to the counting input; Exercise~\ref{ex:ct-values} reads them.

\section{Chains}\label{sec:ct-chains}

% New: proved in the last part.
\begin{definition}[Chains]\label{def:ct-chain}
For distinct points $a_0,a_{w+1}\in\mathbb C$ and points $a_1,\dots,a_w$, $w\ge0$, none on the open segment from $a_0$ to $a_{w+1}$, the \emph{chain} is the iterated integral along that segment, $\gamma(t)=a_0+t(a_{w+1}-a_0)$,
\[
\mathrm I(a_0;a_1,\dots,a_w;a_{w+1})=\int_{0<t_1<\cdots<t_w<1}\ \prod_{i=1}^w\frac{\gamma'(t_i)\,dt_i}{\gamma(t_i)-a_i},
\]
with $\mathrm I(a_0;\varnothing;a_{w+1})=1$, wherever it converges. The $a_1,\dots,a_w$ are its \emph{letters}, $a_0$ and $a_{w+1}$ its \emph{ends} and $w$ its \emph{weight}. It converges when its first letter differs from its first end and its last letter from its last end (Proposition~\ref{prop:ct-between}(i), below). A chain is \emph{of level $M$} when its ends and letters lie in $\{0\}\cup\mu_M$.
\end{definition}

% New: proved in the last part.
\begin{proposition}[The content as chains]\label{prop:ct-chains}
\begin{enumerate}[label=(\roman*)]
\item A trail $T$ of level $M$ with the data $(\nu_i,a_i)_{i\le d}$ is
\[
T=(-1)^d\,\mathrm I\big(0;\ c_1^{-1},0^{\nu_1-1},\ \dots,\ c_d^{-1},0^{\nu_d-1};\ 1\big),\qquad c_i=e^{2\pi i(a_i+\cdots+a_d)}\in\mu_M,
\]
a chain of level $M$ from $0$ to $1$ whose first letter is not $0$ and whose last letter is not $1$ (Theorem~\ref{thm:ct-trailchain}).
\item Conversely, let $\mathrm I(0;a_1,\dots,a_w;1)$ be a chain of level $M$ with $a_1\neq0$ and $a_w\neq1$, let $j_1<\cdots<j_d$ be the positions of its nonzero letters, $j_{d+1}=w+1$, and put $c_i=a_{j_i}^{-1}$, $\nu_i=j_{i+1}-j_i$, $c_{d+1}=1$. With frequencies $a'_i\in\frac1M\mathbb Z$ such that $e^{2\pi ia'_i}=c_i/c_{i+1}$, the chain is $(-1)^d$ times the trail of level $M$ with the data $(\nu_i,a'_i)_{i\le d}$.
\item So the values of $\mathcal C^{(M)}$ form the $\mathbb Q(\mu_M)$-span of $1$ and of the chains of level $M$ from $0$ to $1$ with first letter not $0$ and last letter not $1$, and those of $\mathcal C^{(M)}_w$ the span of $1$ and of such chains of weight at most $w$.
\end{enumerate}
\end{proposition}

% New: proved in the last part.
\begin{proposition}[Chains with an end at zero]\label{prop:ct-between}
\begin{enumerate}[label=(\roman*)]
\item A chain whose first letter differs from its first end and whose last letter differs from its last end converges absolutely.
\item For $\lambda\in\mathbb C^\times$ and $\eta\in\mathbb C$, $\mathrm I(\lambda a_0+\eta;\lambda a_1+\eta,\dots,\lambda a_w+\eta;\lambda a_{w+1}+\eta)=\mathrm I(a_0;a_1,\dots,a_w;a_{w+1})$.
\item $\mathrm I(a_{w+1};a_w,\dots,a_1;a_0)=(-1)^w\,\mathrm I(a_0;a_1,\dots,a_w;a_{w+1})$.
\item Every chain of level $M$ with one end $0$ and the other in $\mu_M$, whose first letter differs from its first end and whose last letter differs from its last end, is $\pm1$ times a chain of level $M$ from $0$ to $1$ with first letter not $0$ and last letter not $1$, so its value lies in the values of $\mathcal C^{(M)}$.
\end{enumerate}
\end{proposition}

A chain between two singular points, two points of $\{0\}\cup\mu_M$, is the verb that the cutting rule uses on its right factors; with one end at $0$ it is a value of the content, and on the symbols below it is read as it stands.

\section{The cutting rule}\label{sec:ct-cutting}

% New: proved in the last part.
\begin{definition}[Symbols of chains]\label{def:ct-symbols}
For an integer $M\ge1$ the \emph{symbols of level $M$} form the commutative $\mathbb Q(\mu_M)$-algebra $\mathfrak S_M$ of polynomials in the \emph{formal chains} $\mathrm I^{\mathfrak s}(a_0;a_1,\dots,a_w;a_{w+1})$, $w\ge1$, with $a_0,\dots,a_{w+1}\in\{0\}\cup\mu_M$ and $a_0\neq a_{w+1}$, as independent variables. In $\mathfrak S_M$ one writes $\mathrm I^{\mathfrak s}(a;\varnothing;b)=1$ for all $a,b$, and $\mathrm I^{\mathfrak s}(a;a_1,\dots,a_w;a)=0$ for $w\ge1$. The \emph{weight} of $\mathrm I^{\mathfrak s}(a_0;a_1,\dots,a_w;a_{w+1})$ is $w$, and $\mathfrak S_M$ is graded by the weight, the weight of a product being the sum. The \emph{symbol} of a trail of level $M$ is $(-1)^d\,\mathrm I^{\mathfrak s}\big(0;c_1^{-1},0^{\nu_1-1},\dots,c_d^{-1},0^{\nu_d-1};1\big)$ (Proposition~\ref{prop:ct-chains}(i)), and the \emph{realization} of a formal chain with one end $0$, the other in $\mu_M$, first letter not its first end and last letter not its last end, is the chain with the same entries (Proposition~\ref{prop:ct-between}(i)).
\end{definition}

% New: proved in the last part.
\begin{definition}[The cutting rule]\label{def:ct-cutting}
The \emph{cutting rule} is the homomorphism of $\mathbb Q(\mu_M)$-algebras $\Delta:\mathfrak S_M\to\mathfrak S_M\otimes_{\mathbb Q(\mu_M)}\mathfrak S_M$ given on the formal chains by
\[
\Delta\,\mathrm I^{\mathfrak s}(a_0;a_1,\dots,a_w;a_{w+1})=\sum_{S\subseteq\{1,\dots,w\}}\mathrm I^{\mathfrak s}(a_0;a_S;a_{w+1})\otimes\prod_{p=0}^{|S|}\mathrm I^{\mathfrak s}\big(a_{s_p};a_{s_p+1},\dots,a_{s_{p+1}-1};a_{s_{p+1}}\big),
\]
where $s_1<\cdots<s_{|S|}$ lists $S$, $s_0=0$, $s_{|S|+1}=w+1$ and $a_S=(a_{s_1},\dots,a_{s_{|S|}})$: the left factor keeps the letters of $S$ between the same ends, and each right factor is a run of dropped letters read between the two kept letters, or ends, that bound it. The \emph{counit} $\varepsilon:\mathfrak S_M\to\mathbb Q(\mu_M)$ is the homomorphism that vanishes on the formal chains. For $\lambda\in\mu_M$ the \emph{rotation} $\rho_\lambda$ multiplies every entry of every formal chain by $\lambda$, and the \emph{reversal} $\varrho$ maps $\mathrm I^{\mathfrak s}(a_0;a_1,\dots,a_w;a_{w+1})$ to $(-1)^w\,\mathrm I^{\mathfrak s}(a_{w+1};a_w,\dots,a_1;a_0)$; both extend to automorphisms of the algebra.
\end{definition}

% New: proved in the last part.
\begin{theorem}[The cutting rule is a coaction]\label{thm:ct-coaction}
\begin{enumerate}[label=(\roman*)]
\item $\Delta$ is coassociative and counital: $(\Delta\otimes\mathrm{id})\circ\Delta=(\mathrm{id}\otimes\Delta)\circ\Delta$ and $(\varepsilon\otimes\mathrm{id})\circ\Delta=\mathrm{id}=(\mathrm{id}\otimes\varepsilon)\circ\Delta$.
\item $\Delta$ preserves the weight: it maps a symbol of weight $w$ to a sum of tensors of weights $w'$ and $w''$ with $w'+w''=w$; and for every formal chain $x$, $\Delta x-x\otimes1-1\otimes x$ is a sum of tensors whose two factors have positive weights.
\item $\Delta\circ\rho_\lambda=(\rho_\lambda\otimes\rho_\lambda)\circ\Delta$ for $\lambda\in\mu_M$, and $\Delta\circ\varrho=(\varrho\otimes\varrho)\circ\Delta$.
\end{enumerate}
\end{theorem}

On the symbol of a trail of level $M$ the left factors are the symbols of the chains from $0$ to $1$ of the letters kept, up to sign the symbol of a trail of level $M$ whenever the kept letters begin with a letter other than $0$ and end with one other than $1$ (Proposition~\ref{prop:ct-chains}(ii)), and the right factors are runs of the trail's letters read between singular points, which at tier one are $0$ and roots of unity of order $M$. Exercise~\ref{ex:ct-cutting} computes the rule on the symbols of the trails $\Tx{3}{\IZ}{0}$ and of index $2$ at the frequency $\frac13$.

\section{The letter of a live input}\label{sec:ct-live}

% New: proved in the last part.
\begin{definition}[The letter of an input]\label{def:ct-letter}
The \emph{letter} of the input with density $\psi$ on the integers $n\ge n_0$ of the counting input, $n_0\in\mathbb N$, is the power series $L_\psi(z)=\sum_{n\ge n_0}\psi(n)\,z^n$. At $z=e^{-u}$ it is the Laplace transform of the input beyond the cut $n_0-1$ (Definition~\ref{def:inputs}), and where the series converges at $z=1$ its value there is the Trailmoment at index zero $\Tx{0}{\psi\,\IZ}{n_0-1}=\sum_{n\ge n_0}\psi(n)$. The letter of the lattice input at the frequency $a$ on $\mathbb Z$ beyond the cut $0$ is $e^{2\pi ia}z/(1-e^{2\pi ia}z)$, singular at the letter $e^{-2\pi ia}$ of its chains (Proposition~\ref{prop:ct-chains}).
\end{definition}

% New: proved in the last part.
\begin{proposition}[The letter of a live input]\label{prop:ct-letter}
Let $\psi$ satisfy $\sum_{i=0}^hp_i(n)\,\psi(n+i)=0$ for the integers $n\ge n_0$, with $p_i\in\mathbb C[n]$, let $d$ be the largest degree of the $p_i$ and $\sigma(\lambda)=\sum_ic_i\lambda^i$ the symbol, $c_i$ the coefficient of $n^d$ in $p_i$; write $\theta=z\frac{d}{dz}$.
\begin{enumerate}[label=(\roman*)]
\item As formal power series, $\sum_{i=0}^hz^{h-i}p_i(\theta-i)\,L_\psi=Q$ for a polynomial $Q$ of degree less than $n_0+h$: a linear differential equation of order $d$ whose coefficient of $\frac{d^d}{dz^d}$ is $z^{d}\sum_ic_iz^{h-i}=z^{d+h}\sigma(1/z)$. Its singular points in $\mathbb C^\times$ are the reciprocals of the nonzero roots of $\sigma$.
\item Let $c_h\neq0$, that is, $p_h$ of the largest degree, and let $\varrho$ be the largest modulus of a root of $\sigma$. Then for every $\eta>0$, $|\psi(n)|\le C_\eta(\varrho+\eta)^n$ for $n\ge n_0$; $L_\psi$ converges on $|z|<1/\varrho$ and continues analytically along every path in $\mathbb C^\times$ that starts in that disc and avoids the reciprocals of the nonzero roots of $\sigma$, and to the plane cut along the rays $\{t/\lambda:t\ge1\}$, $\lambda$ a nonzero root of $\sigma$.
\item The termwise product of two sequences that satisfy, for large $n$, recurrences of orders $h$ and $h'$ with polynomial coefficients and nonzero extreme coefficients satisfies, for large $n$, one of order at most $hh'$ with nonzero extreme coefficients.
\item Let $\psi$ and $\psi'$ be as in (ii), with symbols $\sigma$ and $\sigma'$. Then the letter of $\psi\psi'$ converges on $|z|<1/(\varrho\varrho')$ and continues analytically to the plane cut along the rays $\{t/(\lambda\lambda'):t\ge1\}$, $\lambda$ and $\lambda'$ nonzero roots of $\sigma$ and $\sigma'$.
\end{enumerate}
\end{proposition}

% New: proved in the last part.
\begin{proposition}[The row of a hinge closing read by letters]\label{prop:ct-tiertwo}
In Proposition~\ref{prop:ra-hingerow}(ii) let the recurrences of the room have nonzero extreme coefficients and $p^\pm_{h_\pm}$ of the largest degree, with $\varrho^\pm$ the largest modulus of a root of their symbols, and let the coefficients $\Phi^\pm_n$ of the room periods $x^\pm_n$ satisfy, for large $n$, recurrences of the same kind, with $\varrho^\pm_\Phi$ the largest modulus of a root of their symbols. If $\varrho^\pm_\Phi\varrho^\pm<1$, then:
\begin{enumerate}[label=(\roman*)]
\item every sum $\Xi^\pm_i=\sum_n\Phi^\pm_n\,c^\pm_i(n)$ converges absolutely, it is the value $L_{\psi^\pm_i}(1)$ of the letter of the live input $\psi^\pm_i(n)=\Phi^\pm_n\,c^\pm_i(n)$, and
\[
P_F(\mvec;s,t)=\sum_\pm\sum_{i<h_\pm}x^\pm_i\,L_{\psi^\pm_i}(1):
\]
the room's $h_++h_-$ initial periods against the letters of the live inputs built on its fundamental solutions;
\item each $L_{\psi^\pm_i}$ satisfies a linear differential equation with polynomial coefficients and continues analytically to the plane cut along the rays $\{t/(\lambda\lambda'):t\ge1\}$, $\lambda$ a nonzero root of the symbol of the room's recurrence and $\lambda'$ one of the symbol of the coefficients' recurrence.
\end{enumerate}
\end{proposition}

At tier one the room's recurrence has order one, its fundamental solution is a power times a stacked Gamma (Proposition~\ref{prop:frame}), and each lattice carries one letter; at higher tier the room's periods open into its $h$ fundamental solutions, and the letters of the row are singular only at products of the reciprocals of the roots of the two symbols.

\chapter{Finite algebra for scale reads}\label{ch:finitealgebra}

A bunch read on its scale against linear letters is a Breta of one bunch at the unit configuration (Proposition~\ref{prop:scalebreta}), and a Breta of unit columns on one bunch is, up to a stacked Gamma, such a scale integral. Its values at integer shifts of its index and its letter exponents satisfy two relations, the letter split and the scale exchange, and these relations are complete: every linear relation between the shifted values, with coefficients rational in the index and the letter exponents, is a combination of shifted letter splits and scale exchanges, and the values modulo the relations form a space whose dimension is the number of letters. The same holds at every corner, with the coefficients $\xi_\ell$ of the letters $x+\xi_\ell$ as the variables. For polynomial letters on one scale the same two relations reduce every insertion to one read per distinct root of the product of the letters, counted without finding a root; at a numerical corner three Laurent coefficients per read suffice, every integer ray has a recurrence with a certificate, and the jets in the coefficients of the letters are reads of rational insertions.

\section{Complete relations for linear letters and unit columns}

% Reference 20: def:shifts; source line 12312
\begin{definition}[Shifts, relations between shifts, the relation space]\label{def:shifts}
Let $J=\{0,\dots,L\}$ carry distinct coefficients $\xi_\ell\in\mathbb C\setminus(-\infty,0]$, with the unit configuration $I_J$, and let $\mathbb K=\mathbb C(\kappa,w_0,\dots,w_L)$ be the field of rational functions of an index $\kappa$ and letter exponents $\mathbf w$. Let $\mathcal M$ be the vector space over $\mathbb K$ with one basis element $[a,\mathbf b]$ for each \emph{shift} $(a,\mathbf b)\in\mathbb Z\times\mathbb Z^J$. The \emph{value} of $[a,\mathbf b]$ is the scale integral at the shifted index and letter exponents,
\begin{multline*}
\int_0^\infty x^{\kappa+a-1}\prod_{\ell\in J}(x+\xi_\ell)^{-w_\ell-b_\ell}\,dx\\
=\Gst{2}{|J|}{\kappa+a,\ |\mathbf w|+|\mathbf b|-\kappa-a}{\mathbf w+\mathbf b}\,\sB^{\{J\}}_{\mathbf w+\mathbf b,\,|\mathbf w|+|\mathbf b|-\kappa-a;\,I_J}(\boldsymbol\xi)
\end{multline*}
which converges on $0<\re\kappa+a<\re|\mathbf w|+|\mathbf b|$, the second form holding where moreover $\re(w_\ell+b_\ell)>0$ for every $\ell$ (Proposition~\ref{prop:scalebreta}); finitely many shifts converge together on a nonempty convex open set. An element $\sum_ic_i[a_i,\mathbf b_i]$ of $\mathcal M$ is a \emph{relation between shifts} when the combination of the values of the $[a_i,\mathbf b_i]$ with the coefficients $c_i(\kappa,\mathbf w)$ vanishes at every point where all its terms converge and every $c_i$ is defined. At the shift $(a,\mathbf b)$ the letter split of the letter $x+\xi_\ell$ and the scale exchange are
\begin{align*}
\operatorname{split}_\ell(a,\mathbf b)&=[a,\mathbf b]-[a+1,\mathbf b+\mathbf e_\ell]-\xi_\ell\,[a,\mathbf b+\mathbf e_\ell],\\
\operatorname{exch}(a,\mathbf b)&=(\kappa+a)\,[a,\mathbf b]-\sum_{\ell\in J}(w_\ell+b_\ell)\,[a+1,\mathbf b+\mathbf e_\ell],
\end{align*}
the letter split and the scale exchange of Proposition~\ref{prop:letters} for the letters $x+\xi_\ell$. The \emph{relation space} $\mathcal L_J\subseteq\mathcal M$ is the span over $\mathbb K$ of all letter splits and scale exchanges, and the \emph{dimension} of the shifts is $\dim_{\mathbb K}\mathcal M/\mathcal L_J$; on one scale these play the part of the relation space and the dimension of a ray (Definition~\ref{def:raygauge}).
\end{definition}

% Reference 20: thm:skcomplete; source line 12327
\begin{theorem}[The relations of a scale integral are complete]\label{thm:skcomplete}
Let the coefficients range over distinct points of $\mathbb C\setminus(-\infty,0]$.
\begin{enumerate}[label=(\roman*)]
\item Every element of $\mathcal L_J$ is a relation between shifts.
\item $\mathcal M=\mathcal L_J+\sum_{\ell\in J}\mathbb K\,[0,\mathbf e_\ell]$: the dimension is at most $|J|$.
\item The values of the $|J|$ shifts $[0,\mathbf e_\ell]$,
\[
\Gst{2}{|J|}{\kappa,\ |\mathbf w|+1-\kappa}{\mathbf w+\mathbf e_\ell}\,\sB^{\{J\}}_{\mathbf w+\mathbf e_\ell,\,|\mathbf w|+1-\kappa;\,I_J}(\boldsymbol\xi),\qquad \ell\in J,
\]
are linearly independent over $\mathbb K$: a combination of them with coefficients in $\mathbb K$ that vanishes on a nonempty open subset of $\{0<\re\kappa<\re|\mathbf w|+1\}$ has every coefficient $0$.
\item The dimension is $|J|$, and every relation between shifts lies in $\mathcal L_J$: it is a combination over $\mathbb K$ of shifted letter splits and scale exchanges.
\item \emph{Coefficients as variables.} Let the coefficients be rational functions of parameters $\mathbf u=(u_1,\dots,u_d)$, distinct and not $0$ as rational functions, defined with distinct values in $\mathbb C\setminus(-\infty,0]$ on a nonempty open subset of $\mathbb R^d$ or of $\mathbb C^d$. Then (i), (ii) and (iv) hold with the field $\mathbb C(\mathbf u,\kappa,\mathbf w)$ in the place of $\mathbb K$, a relation between shifts being an element whose combination of values vanishes at every $\mathbf u$ of that set and every $(\kappa,\mathbf w)$ where all its terms converge and every coefficient is defined.
\end{enumerate}
\end{theorem}

The letter splits are the partial fractions of the integrand, and the scale exchange is one integration by parts along the scale. Part (ii) reduces every shift to the $|J|$ values with one letter exponent raised by one, by these two relations alone, and part (iii) says that those values are linearly independent over $\mathbb K$. So a relation between shifts that is not a combination of shifted letter splits and scale exchanges does not exist, and the number of independent values is the number of letters; for one letter the dimension is one: every shift of the Beta corner of a linear letter is a rational multiple of one value. By Theorem~\ref{thm:scale} the scale integral of a family at a bunch is such an integral, with the points of the joint spectrum as its coefficients, so its relations in its index and its letter exponents are complete; relations between scale integrals of a family at shifted masses and exponents, where the letter exponents are tied to the exponents of the family, are a further problem.

What decides the relations at a corner, where the index and the letter exponents are numbers and the coefficients are the variables, is the behaviour of the values as one coefficient tends to $0$: its letter becomes a power of the scale and joins the index.

% Reference 20: lem:slotzero; source line 12375
\begin{lemma}[A coefficient tends to $0$]\label{lem:coefzero}
Let $L\ge1$, let the coefficients be positive, and let the index and the letter exponents be numbers with $\re w_\ell>0$ for every $\ell$ and $0\le\re\kappa'<1$, where $\kappa'=\kappa-w_0$. In this statement $J'=J\setminus\{0\}$, $\mathbf w'=\mathbf w|_{J'}$, $\boldsymbol\xi'=\boldsymbol\xi|_{J'}$, $f_\ell$ is the value of $[0,\mathbf e_\ell]$ (Definition~\ref{def:shifts}), and for letter exponents $\mathbf v$ on $J'$ with $\re|\mathbf v|>0$ and $h_{\mathbf v}(x)=\prod_{j\in J'}(x+\xi_j)^{-v_j}$,
\[
S'(z;\mathbf v)=\frac{h_{\mathbf v}(0)}z+\int_0^1x^{z-1}\big(h_{\mathbf v}(x)-h_{\mathbf v}(0)\big)\,dx+\int_1^\infty x^{z-1}h_{\mathbf v}(x)\,dx
\]
for $-1<\re z<\re|\mathbf v|$, $z\neq0$. For $\re z>0$ and $\re v_j>0$ for every $j$, $S'(z;\mathbf v)=\Gst{2}{|J'|}{z,\ |\mathbf v|-z}{\mathbf v}\,\sB^{\{J'\}}_{\mathbf v,\,|\mathbf v|-z;\,I_{J'}}(\boldsymbol\xi')$, the scale integral with the coefficients $\boldsymbol\xi'$, so $S'$ continues it in its index across the wall $\re z=0$, with a simple pole at $z=0$ of residue $h_{\mathbf v}(0)$; its constant term $\CT_{z=0}S'(z;\mathbf v)$ is the sum of the two integrals at $z=0$, and its values off the pole, at complex $z$ and complex letter exponents, are the finite parts $\Reg\int_0^\infty x^{z-1}h_{\mathbf v}\,dx$, its constant term at $z=0$ being $\Reg\int_0^\infty x^{-1}h_{\mathbf v}\,dx$ (Theorem~\ref{thm:RegCT}). Fix $\boldsymbol\xi'$ and let $\xi_0\to0^+$. Then, with $j\in J'$ in (ii)--(v):
\begin{enumerate}[label=(\roman*)]
\item $\xi_0^{\,1-\kappa'}f_0\to\Gst21{\kappa,\ 1-\kappa'}{w_0+1}\,h_{\mathbf w'}(0)$, which is not $0$.
\item If $\re\kappa'>0$, then $f_j\to S'(\kappa';\mathbf w'+\mathbf e_j)$.
\item If $\re\kappa'=0$ and $\kappa'\neq0$, then
\[
f_j=\Gst21{\kappa,\ -\kappa'}{w_0}\,\frac{h_{\mathbf w'}(0)}{\xi_j}\,\xi_0^{\,\kappa'}+S'(\kappa';\mathbf w'+\mathbf e_j)+o(1).
\]
\item If $\kappa'=0$, then
\[
f_j=\frac{h_{\mathbf w'}(0)}{\xi_j}\Big(\log\frac1{\xi_0}+(\log\Gamma)'(1)-(\log\Gamma)'(\kappa)\Big)+\CT_{z=0}S'(z;\mathbf w'+\mathbf e_j)+o(1).
\]
\item \emph{The relations across the wall.} For $-1<\re z\le0$, $z\neq0$,
\[
S'(z;\mathbf v)=S'(z+1;\mathbf v+\mathbf e_j)+\xi_j\,S'(z;\mathbf v+\mathbf e_j),\qquad z\,S'(z;\mathbf v)=\sum_{i\in J'}v_i\,S'(z+1;\mathbf v+\mathbf e_i),
\]
and $\CT_{z=0}S'(z;\mathbf v)=S'(1;\mathbf v+\mathbf e_j)+\xi_j\,\CT_{z=0}S'(z;\mathbf v+\mathbf e_j)$.
\end{enumerate}
\end{lemma}

As the coefficient $\xi_0$ tends to $0$ its letter $x+\xi_0$ becomes the power $x$. The value with the exponent of that letter raised grows like $\xi_0^{\,\kappa'-1}$, and when $\re\kappa'>0$ the others tend to the scale integrals of the remaining coefficients at the index $\kappa'=\kappa-w_0$, the index left when the letter exponent $w_0$ has joined it. When $\re\kappa'=0$ that index lies on the wall of the remaining integral, and a second term of the same size stands next to the integral continued across its wall: the oscillating power $\xi_0^{\,\kappa'}$, and at $\kappa'=0$ a logarithm, which grows. The letter split and the scale exchange hold across the wall, and at the pole $z=0$ the scale exchange keeps only its residue, the first identity of Theorem~\ref{thm:skcorner}(ii) for the remaining coefficients. These are the first terms of the closing between two separated groups (Corollary~\ref{cor:twogroups}), the letter of $\xi_0$ being the lower group: one lattice carries the powers $\xi_0^{\,\kappa'+n}$, with the wall residues of the other letters as coefficients, and the other carries the integer powers $\xi_0^{\,n}$, with the scale integral of the other coefficients continued in its index; at $\kappa'=0$ the two lattices meet, and what their poles leave is the logarithm of (iv).

% Reference 20: thm:skcorner; source line 12441
\begin{theorem}[The relations at a corner]\label{thm:skcorner}
Let the index $\kappa$ and the letter exponents $\mathbf w$ be numbers with $0<\re\kappa<\re|\mathbf w|+1$, so that the values of the $|J|$ shifts $[0,\mathbf e_\ell]$ converge, and let the coefficients be distinct points of $\mathbb C\setminus(-\infty,0]$. A shift $[a,\mathbf b]$ \emph{converges at the corner} when $0<\re\kappa+a<\re|\mathbf w|+|\mathbf b|$, and a letter split or a scale exchange at the shift $(a,\mathbf b)$, with the index and the letter exponents at their numbers, converges at the corner when $[a,\mathbf b]$ does; all its terms then converge. A \emph{linear relation} between shifts that converge at the corner is a combination $\sum_ir_i[a_i,\mathbf b_i]$ of such shifts, with coefficients $r_i$ rational in $\boldsymbol\xi$, whose value vanishes for $\boldsymbol\xi$ in a nonempty open set; it is a combination of letter splits and scale exchanges when it is one as a formal combination of shifts. The values are holomorphic in $\boldsymbol\xi$ on $(\mathbb C\setminus(-\infty,0])^J$, coincident coefficients included, and functions of $\boldsymbol\xi$ are \emph{independent over $\boldsymbol\xi$} when no combination of them with coefficients polynomial in $\boldsymbol\xi$, not all $0$, vanishes at all $\boldsymbol\xi>0$, equivalently for $\boldsymbol\xi$ in a nonempty open set.
\begin{enumerate}[label=(\roman*)]
\item \emph{Reduction.} If no letter exponent lies in $\{-1,-2,\dots\}$, the value of every shift that converges at the corner is a combination of the values of the $[0,\mathbf e_\ell]$ with coefficients rational in $\boldsymbol\xi$; the combination follows from letter splits and scale exchanges that converge at the corner, taken at shifts $(a',\mathbf b')$ with $b'_\ell\ge\min(b_\ell,0)$ for every $\ell$.
\item \emph{The two walls.} If $\re|\mathbf w|>0$, the values of the $[0,\mathbf e_\ell]$ at the index $1$ and at the index $|\mathbf w|$ satisfy
\begin{gather*}
\sum_{\ell\in J}w_\ell\int_0^\infty\prod_{\ell'\in J}(x+\xi_{\ell'})^{-w_{\ell'}-\one[\ell'=\ell]}\,dx=\prod_{\ell\in J}\xi_\ell^{-w_\ell},\\
\sum_{\ell\in J}w_\ell\,\xi_\ell\int_0^\infty x^{|\mathbf w|-1}\prod_{\ell'\in J}(x+\xi_{\ell'})^{-w_{\ell'}-\one[\ell'=\ell]}\,dx=1,
\end{gather*}
the integrals being, where every $\re w_\ell>0$, the stacked Gammas times the Bretas of Definition~\ref{def:shifts}.
\item \emph{Independence.} The values of the $|J|$ shifts $[0,\mathbf e_\ell]$ are independent over $\boldsymbol\xi$ if and only if no letter exponent lies in $\{-1,-2,\dots\}$.
\item \emph{Completeness.} Whatever the letter exponents, every linear relation between shifts that converge at the corner is a combination, with coefficients rational in $\boldsymbol\xi$, of letter splits and scale exchanges that converge at the corner, and over the rational functions of $\boldsymbol\xi$ the values of these shifts span a space of dimension $|J|$: at every corner the dimension is $|J|$.
\item \emph{A coefficient equal to $1$.} Parts (i), (iii) and (iv) hold as well when one coefficient is the constant $1$ and the others are the variables.
\end{enumerate}
\end{theorem}

At a corner the relations are the same, and when no letter exponent lies in $\{-1,-2,\dots\}$ they reduce every shift that converges to the $|J|$ values with one letter exponent raised without leaving the corner: every relation of the reduction has all its terms convergent, and the leading coefficients $\kappa-r$, $r\ge1$, and $\kappa+j-|\mathbf w|$, $j\ge0$, by which the descents at the origin and at infinity divide vanish only at shifts that have left the domain. The two identities of (ii) are the relations at the two walls of the domain: the first is the scale exchange at the shift $(-1,\mathbf 0)$ at the index $1$, whose member on the wall is replaced by its limit $\prod_\ell\xi_\ell^{-w_\ell}$, and the second is the scale exchange less the letter splits multiplied by their letter exponents at the shift $(0,\mathbf 0)$, at the index $|\mathbf w|$, where $|\mathbf w|-\kappa$ times the value of $[0,\mathbf 0]$ tends to $1$; at a wall one combination of the $|J|$ values is a monomial in $\boldsymbol\xi$. A letter exponent in $\{-1,-2,\dots\}$ is not on a wall: its letter has become a polynomial factor of the integrand, the $|J|$ values with one letter exponent raised are dependent and no longer span the shifts that converge; the corner with that letter exponent moved up to $0$ has the same shifts, and its $|J|$ values with one letter exponent raised span them. Independence is read from the coefficients one at a time: as $\xi_0$ tends to $0$ the value with the letter exponent $w_0$ raised grows faster than the others (Lemma~\ref{lem:coefzero}), so in a relation the polynomial multiplying it vanishes at $\xi_0=0$, and what is left is a relation between scale integrals of the remaining coefficients, at the index $\kappa-w_0$, or on their wall, where the relations across the wall carry it to the next index.

% Reference 20: cor:unitcomplete; source line 12535 (with lem:unitskeleton, source line 12523)
\begin{corollary}[Unit columns on one bunch]\label{cor:unitcomplete}
Let $B$ be a bunch with the unit configuration $I_B$ and distinct coefficients $\xi_e>0$. By Proposition~\ref{prop:scalebreta} with $J=B$, $\mathbf w=\mvec$ and $\kappa=m_B-\nu$, the value $\sB^{\{B\}}_{\mvec+\mathbf a,\,\nu+n;\,I_B}(\boldsymbol\xi)$, $\mathbf a\in\mathbb Z^B$, $n\in\mathbb Z$, is the value of the shift $[|\mathbf a|-n,\mathbf a]$ of Definition~\ref{def:shifts} divided by $\Gst{2}{|B|}{m_B-\nu+|\mathbf a|-n,\ \nu+n}{\mvec+\mathbf a}$, wherever that proposition applies; this stacked Gamma is the one at $(\mvec,\nu)$ times a rational function of $(\mvec,\nu)$. Each relation below is read where all its Bretas converge.
\begin{enumerate}[label=(\roman*)]
\item \emph{The mark relation.} For each $e\in B$,
\[
m_e\,\sB^{\{B\}}_{\mvec,\,\nu;\,I_B}(\boldsymbol\xi)=(m_B-\nu)\,\sB^{\{B\}}_{\mvec+\mathbf e_e,\,\nu;\,I_B}(\boldsymbol\xi)+\nu\,\xi_e\,\sB^{\{B\}}_{\mvec+\mathbf e_e,\,\nu+1;\,I_B}(\boldsymbol\xi):
\]
the sum over $f\in B$ of the shape exchanges of Theorem~\ref{thm:bretalaws}(iv) for the pairs $(e,f)$, reduced by the shape split and the scale exchange; it is the letter split of the letter $x+\xi_e$.
\item The shape split, the scale exchange and the shape exchanges of Theorem~\ref{thm:bretalaws}, at all integer shifts of $(\mvec,\nu)$, span over $\mathbb C(\mvec,\nu)$ every linear relation with coefficients in $\mathbb C(\mvec,\nu)$ between the values $\sB^{\{B\}}_{\mvec+\mathbf a,\,\nu+n;\,I_B}(\boldsymbol\xi)$, and the quotient of the space of shifts by these relations has dimension $|B|$. The same holds over $\mathbb C(\boldsymbol\xi,\mvec,\nu)$ for relations whose coefficients are rational in $\boldsymbol\xi$ as well and which hold at all distinct $\boldsymbol\xi>0$ at which their coefficients are defined.
\item \emph{At a corner.} Let $\mvec$ and $\nu$ be numbers with $\re m_e>0$ for every $e$ and $\re\nu>0$, and call a value $\sB^{\{B\}}_{\mvec+\mathbf a,\,\nu+n;\,I_B}(\boldsymbol\xi)$ \emph{interior} when $\re(m_e+a_e)>0$ for every $e$ and $\re\nu+n>0$; the Breta converges there. Over the rational functions of $\boldsymbol\xi$ the interior values span a space of dimension $|B|$, and every linear relation between interior values with coefficients rational in $\boldsymbol\xi$ is a combination, with such coefficients, of shape splits, scale exchanges and shape exchanges all of whose terms are interior. If $\re\nu<\re m_B$, the $|B|$ values $\sB^{\{B\}}_{\mvec+\mathbf e_e,\,\nu+1;\,I_B}(\boldsymbol\xi)$, $e\in B$, are a basis of that space.
\end{enumerate}
\end{corollary}

Completeness holds for the scale integral of one bunch against linear letters and for unit columns on one bunch, which are the same integrals: the coefficients are the same, the shape index is the letter exponents, and the index of the scale integral is the amount $m_B-\nu$ by which the scale index falls short of the mass $m_B$. It holds at a free index and free letter exponents, and at every corner with the coefficients as the variables, one of them possibly the constant $1$. For a Breta with a kernel, with letters of more than two patterns or of levels other than $0$ and $1$, or with several bunches, the relations of Theorem~\ref{thm:bretalaws} hold, and their completeness is a further problem.

\section{Polynomial letters on one scale}

A word on one scale with polynomial letters, such as the scale integral of a bunch at fixed shapes (Theorems~\ref{thm:scale} and~\ref{thm:scaletotals}), has the same two relations, the letter splits becoming partial fractions over the distinct factors of the letters, and the same finiteness: the retained data are one read per distinct root of the product of the letters, counted without finding a root, and at a numerical corner three Laurent coefficients per read. Every rational insertion, every integer shift of the index and the letter exponents, every derivative in the coefficients of the letters and every integer ray is then finite algebra, with a certificate. 

% Reference 20: sec:scalesources, opening; source line 12568 (letters with nonnegative coefficients and positive constant terms).
% Proved in the last part: def:insertions.
\begin{definition}[Insertions and reads on one scale]\label{def:insertions}
Let $\mathbb F\subseteq\mathbb C$ be a field, the field of the letters, and let $P_1,\dots,P_N\in\mathbb F[x]$ be nonconstant letters without roots on $[0,\infty)$, a constant letter being moved into a scalar; put $\mathbf P=(P_1,\dots,P_N)$. For an index $\kappa$ and letter exponents $w_1,\dots,w_N$ in a field $\mathbb F'\supseteq\mathbb F(\kappa,\mathbf w)$ embedded in $\mathbb C$, the insertions, the regular functions of the letter complement as in Definition~\ref{def:ray}, are the elements of the ring $\mathfrak A=\mathbb F'[x,x^{-1},P_1^{-1},\dots,P_N^{-1}]$, and the \emph{read} of an insertion $R$ is
\[
\langle R\rangle_{\kappa,\mathbf w;\mathbf P}=\int_0^\infty x^{\kappa-1}\prod_jP_j(x)^{-w_j}\,R(x)\,dx
\]
where it converges absolutely, each power $P_j^{-w_j}$ being taken on $[0,\infty)$ by continuity from its principal value at $0$, and continued in the index and the letter exponents elsewhere. The shift of the index by $a$ and of the letter exponents by $b_j$, integers, is the insertion $x^a\prod_jP_j^{-b_j}$, and a derivative in the coefficients of a letter is an insertion. A squarefree decomposition with gcd splitting gives pairwise coprime squarefree $q_1,\dots,q_M\in\mathbb F[x]$, none vanishing at $0$, with $P_j=c_j\prod_kq_k^{m_{jk}}$, without roots and without a field extension; $Q=q_1\cdots q_M$ is the monic radical of $P_1\cdots P_N$, of degree $D>0$. The \emph{scale exchange} on insertions is
\[
\eth g=xg'+\Big(\kappa-\sum_jw_j\frac{xP_j'}{P_j}\Big)g,\qquad\text{so that}\qquad\Big(x^{\kappa}\prod_jP_j^{-w_j}\,g\Big)'=x^{\kappa-1}\prod_jP_j^{-w_j}\,\eth g .
\]
For the linear letters $x+\xi_\ell$ with distinct $\xi_\ell\in\mathbb C\setminus(-\infty,0]$, the reads of the shifts are the values of Definition~\ref{def:shifts}, $D=|J|$, and $\eth$ applied to $x^a\prod_\ell(x+\xi_\ell)^{-b_\ell}$ is the scale exchange $\operatorname{exch}(a,\mathbf b)$ read under the integral.
\end{definition}

% Reference 20: thm:scalereduce; source line 12577 (letters with nonnegative coefficients and positive constant terms).
% Proved in the last part: thm:scalereduce.
\begin{theorem}[Reduction by the scale exchange]\label{thm:scalereduce}
Let the numbers
\[
\kappa-r\quad(r\ge1),\qquad \kappa+a-\sum_jw_j\deg P_j\quad(a\ge0),\qquad \sum_jm_{jk}w_j+h\quad(h\ge1,\ 1\le k\le M)
\]
be nonzero in $\mathbb F'$. Every $R\in\mathfrak A$ has a unique decomposition
\[
R=\eth g+\sum_{i=0}^{D-1}c_i\,\frac{x^i}{Q},\qquad g\in\mathfrak A,\quad c_i\in\mathbb F',
\]
obtained by polynomial arithmetic, partial fractions and the extended Euclidean algorithm; so $\dim_{\mathbb F'}\mathfrak A/\eth\mathfrak A=D$, and on the domain where the reads converge $\langle R\rangle_{\kappa,\mathbf w;\mathbf P}=\sum_ic_i\,\langle x^i/Q\rangle_{\kappa,\mathbf w;\mathbf P}$. For the linear letters $x+\xi_\ell$ with distinct $\xi_\ell\in\mathbb C\setminus(-\infty,0]$, $D=|J|$ and this is Theorem~\ref{thm:skcomplete}(ii) over $\mathbb K$, the $x^i/Q$ spanning the same space as the $(x+\xi_\ell)^{-1}$.
\end{theorem}

The dimension is that of the quotient by the scale exchange, the dimension of the ray on one scale (Theorem~\ref{thm:rayspace}); it does not assert that the reads of the $D$ insertions $x^i/Q$ are independent once parameters are tied, and symmetries of the letters can relate them further. For linear letters Theorem~\ref{thm:skcomplete}(iii) gives the independence. The three conditions are the leading coefficients of the reduction at the origin, at infinity and at a repeated factor. Exact computation needs exact arithmetic with a zero test in $\mathbb F'$, numerical exponents being taken in an effective extension; the analytic statements hold at all complex exponents.

% Reference 20: sec:scalesources, Corners; source line 12594 (the one direction $(\frac12;\one)$).
% Proved in the last part: def:masters.
\begin{definition}[Deformation; masters]\label{def:masters}
At numerical exponents $(\kappa,\mathbf w)$ fix a direction $(\lambda_0;\boldsymbol\lambda)\in\mathbb R\times\mathbb R^N$ of the index and the letter exponents that is \emph{admissible}: $\lambda_0>0$, $\sum_j\lambda_j\deg P_j>\lambda_0$ and $\sum_jm_{jk}\lambda_j\neq0$ for every $k$. The \emph{deformation} moves the data along it, to $(\kappa+\lambda_0\varepsilon,\ \mathbf w+\varepsilon\boldsymbol\lambda)$; the three conditions of Theorem~\ref{thm:scalereduce} hold over $\mathbb F'(\varepsilon)$, and for an insertion $R$ the coefficients $c_i(\varepsilon)$ are those of its decomposition there. The \emph{masters} are the deformed reads
\[
H_i(\varepsilon)=\big\langle x^i/Q\big\rangle_{\kappa+\lambda_0\varepsilon,\ \mathbf w+\varepsilon\boldsymbol\lambda;\ \mathbf P},\qquad 0\le i<D .
\]
For large real $\varepsilon$ every finite set of deformed reads converges, the exponent at the origin growing with the slope $\lambda_0$ and the decay at infinity with the slope $\sum_j\lambda_j\deg P_j-\lambda_0>0$, so the reductions hold there and then as identities of meromorphic functions. The masters are finite combinations of reads of shifts of the original letters: $1/Q=\big(P_1\cdots P_N/Q\big)\prod_jP_j^{-1}$ with $P_1\cdots P_N/Q$ a polynomial, so each term of $x^iP_1\cdots P_N/Q$ shifts the index by an integer and raises every letter exponent by one. The admissible directions form a nonempty open cone; the direction $(\frac12;\one)$ is admissible.
\end{definition}

% Reference 20: lem:simplepoles; source line 12600 (the direction $(\frac12;\one)$, letters with nonnegative coefficients).
% Proved in the last part: lem:simplepoles.
\begin{lemma}[Simple poles at a corner]\label{lem:simplepoles}
For every insertion $R$ independent of $\varepsilon$, or holomorphic in $\varepsilon$, the coefficients $c_i(\varepsilon)$ have at most a simple pole at $\varepsilon=0$, and every master $H_i(\varepsilon)$ is meromorphic near $0$ with at most a simple pole.
\end{lemma}

% Reference 20: thm:corner; source line 12607 (the direction $(\frac12;\one)$, letters with nonnegative coefficients).
% Proved in the last part: thm:corner.
\begin{theorem}[Finitely many numbers at every corner]\label{thm:corner}
Along an admissible direction (Definition~\ref{def:masters}), for an insertion $R$ write $H_i(\varepsilon)=h_{i,-1}\varepsilon^{-1}+h_{i,0}+h_{i,1}\varepsilon+O(\varepsilon^2)$ and $c_i(\varepsilon)=c_{i,-1}\varepsilon^{-1}+c_{i,0}+c_{i,1}\varepsilon+O(\varepsilon^2)$. The $3D$ \emph{corner numbers} $(h_{i,-1},h_{i,0},h_{i,1})_{i<D}$ determine every absolutely convergent read:
\[
\langle R\rangle_{\kappa,\mathbf w;\mathbf P}=\sum_{i=0}^{D-1}\big(c_{i,-1}h_{i,1}+c_{i,0}h_{i,0}+c_{i,1}h_{i,-1}\big).
\]
They depend on the letters, the exponents and the direction only, not on the insertion or the size of the shift. The reads of $R\,\big(\lambda_0\log x-\sum_j\lambda_j\log P_j\big)^b$, the derivatives $\partial_\varepsilon^b$ at $\varepsilon=0$ of the deformed reads $\langle R\rangle_{\kappa+\lambda_0\varepsilon,\,\mathbf w+\varepsilon\boldsymbol\lambda;\,\mathbf P}$, need the masters through the order $b+1$: $(b+3)D$ numbers. When $\langle R\rangle_{\kappa,\mathbf w;\mathbf P}$ converges absolutely, a read of $R\,(\log x)^{k_0}\prod_j(\log P_j)^{k_j}$, $b=k_0+\sum_jk_j$, is a rational combination of the reads of $R\,\big(\lambda_0\log x-\sum_j\lambda_j\log P_j\big)^b$ along finitely many admissible directions, each with its own corner numbers.
\end{theorem}

The residues are finite extractions of endpoint coefficients, the other numbers finite parts and derivatives in the deformation. The $D$ values of the masters suffice when every pivot of the reduction is regular and the masters are holomorphic at the corner, as for linear letters at a corner with no letter exponent in $\{-1,-2,\dots\}$ and convergent masters (Theorem~\ref{thm:skcorner}(i)); regular pivots alone do not suffice, and at an arbitrary corner the corner numbers are what is retained.

% Reference 20: thm:rayrec; source line 12620 (letters with nonnegative coefficients and positive constant terms).
% Proved in the last part: thm:rayrec.
\begin{theorem}[Integer rays have certified recurrences]\label{thm:rayrec}
Fix integers $a,b_1,\dots,b_N$, and let $\eth_n$ be the scale exchange at the deformed data $\big(\kappa+an+\frac\varepsilon2,\ \mathbf w+n\mathbf b+\varepsilon\one\big)$. A finite algorithm gives polynomials $p_0(n,\varepsilon),\dots,p_D(n,\varepsilon)$ and a certificate $g$, rational in $x$, $n$ and $\varepsilon$, with
\[
\eth_ng=\sum_{k=0}^Dp_k(n,\varepsilon)\,x^{ak}\prod_jP_j^{-b_jk},
\]
the polynomials $p_k(n)=p_k(n,0)$ not all zero; then
\[
\sum_{k=0}^Dp_k(n)\,\langle 1\rangle_{\kappa+a(n+k),\ \mathbf w+(n+k)\mathbf b;\ \mathbf P}=0
\]
at every integer $n$ at which the displayed reads converge.
\end{theorem}

The recurrence need not have the least order, and its leading coefficient may vanish at finitely many integers, read off that polynomial; there further values are given or another dependence is formed. On one scale this supplies, for every ray, the certificate that a recurrence asks for: an insertion $g$ whose scale exchange is the recurrence.

% Reference 20: thm:scalejets; source line 12633 (letters with nonnegative coefficients and positive constant terms).
% Proved in the last part: thm:scalejets.
\begin{theorem}[Coefficient jets on one scale]\label{thm:scalejets}
Let $P_j(x;u)=P_j(x)+\sum_{a=1}^pu_aP_{ja}(x)$ with $\deg P_{ja}\le\deg P_j$, choose bounds $M_{ja}\ge\sup_{x>0}|P_{ja}/P_j|$, finite since the quotient is continuous on $[0,\infty]$, and let $\langle 1\rangle_{\kappa,\mathbf w;\mathbf P}$ converge absolutely. On the polydisc where $\sum_a|u_a|M_{ja}<1$ for every $j$ the read $\langle 1\rangle_{\kappa,\mathbf w;\mathbf P(\,\cdot\,;u)}$, the powers of the changed letters being continued in $u$ from $u=0$, is holomorphic and
\[
\langle 1\rangle_{\kappa,\mathbf w;\mathbf P(\,\cdot\,;u)}=\sum_{(n_{ja})\ge0}\ \prod_j\frac{(-1)^{|n_j|}(w_j)_{|n_j|}}{\prod_an_{ja}!}\ \prod_au_a^{\sum_jn_{ja}}\ \Big\langle \prod_{j,a}\Big(\frac{P_{ja}}{P_j}\Big)^{n_{ja}}\Big\rangle_{\kappa,\mathbf w;\mathbf P},
\]
$|n_j|=\sum_an_{ja}$, absolutely; every coefficient is the read of a rational insertion, given by the corner numbers of Theorem~\ref{thm:corner}.
\end{theorem}

Splitting of repeated roots under the change costs nothing, the expansion being made against the original letters; a change that raises a degree is outside the statement. For coefficients that are symbolic parameters with fixed gcds and multiplicities, differentiating the masters gives rational insertions, so Theorem~\ref{thm:scalereduce} gives matrices $\partial_{u_a}H=\mathbf A_aH$, with certificates, for the vector $H$ of the masters, whose mixed derivatives agree because they agree before the reduction; at a corner where a leading coefficient of the reduction vanishes the corner numbers replace them. The structural data reused are the squarefree base, $Q$ and the inverses of $xq'_k$ modulo $q_k$; the analytic data are the $D$ masters or the $3D$ corner numbers. A rational insertion with polynomial part of degree $A_+$, a pole of order $A_-$ at the origin and of order $h_k$ at $q_k$ takes at most $(A_++1)+A_-+\sum_k\max(h_k-1,0)$ elimination rounds, counted in operations, the growth of exact coefficients charged separately; a repeated letter adds $m_{jk}$ times its letter exponent to the exponent of each squarefree factor $q_k$ and leaves $Q$ and $D$ unchanged. A recurrence in a count of repetitions needs the integer shifts of Theorem~\ref{thm:rayrec}, and a multiplicity written in binary, which is what the price of Chapter~\ref{ch:price} counts, does not by itself make the evaluation logarithmic in it.

\part{The family as inference}\label{part:inference}

\chapter{The family as inference}\label{ch:inference}

A family is a model of data in two ways. A sample of bases, drawn from the basis and cobasis laws of the tables of a family at an unknown gauge, is itself a question: its counts are the masses, the sizes of its parts are the exponents, and its chance as a function of the tilt is the integrand of the period. The law of a question is then the conjugate prior of the tilt, and every posterior, every predictive chance and every evidence is a ratio of periods; read at a bunch, a sample is a sample of the children at the shares it realizes, and each of its walls is a sum of independent excesses of the levels and colevels of its draws over their floors, whose law the Skeleton polynomials give. A sample of tilts drawn from the law of a question makes the period the cumulant function of an exponential family. Between the two sits the family word read at Dirichlet draws: the negative logarithms of its letters have the Grey law, whose Laplace transform in the exponents is the period divided by the stacked Gamma of the masses over their total; the pseudo-MGF and the pseudoexpectation are that Laplace transform fixed at the Foreman's question and at the centre. A draw of the shape law, read as a vector of numbers, has an order that is a sector of the period, best and worst bases that are functions of the order, and entropies and moments that are steps of the data. What a bunch all of whose subsets are patterns reads of its room is its interaction (Definition~\ref{def:interaction}), the total over the scales of the bunch of the joint cumulants of the use of its marks under the basis and cobasis laws, weighted by the exponents; for a regression design the interaction of a split of the prior marks into two blocks is minus the logarithm of the generalized variance inflation factor.

\section{Samples of bases}

% Reference 20: prop:samplequestion; source line 11409
\begin{proposition}[A sample of bases is a question]\label{prop:samplequestion}
For a family at a gauge $\alpha=e^{2z}$ draw, for every table $j$, $s_j$ bases $T_{j,1},\dots,T_{j,s_j}$ from its basis law $\mathbb P_{j,\alpha}$ and $t_j$ bases $T'_{j,1},\dots,T'_{j,t_j}$ from its cobasis law $\mathbb P^\ast_{j,\alpha}$, all independent, let $m_e$ count the $T_{j,i}$ that hold $e$ and the $T'_{j,i}$ that miss $e$, and let $\theta_{\mathrm{draws}}=\prod_{j,i}\theta_{j,T_{j,i}}\prod_{j,i}\theta_{j,T'_{j,i}}$ be the product of the weights of the drawn bases.
\begin{enumerate}[label=(\roman*)]
\item As a function of the tilt $z$, the chance of the sample is
\[
\theta_{\mathrm{draws}}\cdot\exp\Big(2\ip{\bmu_\perp}z-\sum_j\big(\sigma_jE_j(z)+\delta_jO_j(z)\big)\Big),
\]
$2^{1-n}\theta_{\mathrm{draws}}$ times the integrand of the period of the question $(\mvec;\mathbf s,\mathbf t)$, of pair coordinates $(\bmu_\perp;\boldsymbol\sigma,\boldsymbol\delta)$ (Definition~\ref{def:question}): the counts are the masses, the sizes of the parts are the exponents, and $m_{\cE}=\sum_j(s_jr_j+t_jr^\ast_j)$ puts the question on the balance wall.
\item Its logarithmic gradient in $z$ is $2\big(\mvec-\sum_j(s_jR_j(\alpha)+t_jR^\ast_j(\alpha))\big)$, twice the excess of the masses over the usage at $\alpha$. For a connected family the maximum likelihood tilt is the native gauge of the question, and it exists, unique on $\mathcal T$, exactly when every proper wall of the question is positive.
\item $\kappa_S=\sum_j\Big(\sum_{i=1}^{s_j}\big(|T_{j,i}\cap S|-d_j(S)\big)+\sum_{i=1}^{t_j}\big(|S\setminus T'_{j,i}|-d^\ast_j(S)\big)\Big)\ge0$: a wall is the excess of the sample's use of $S$ over its least use.
\end{enumerate}
\end{proposition}

% Reference 20: prop:posterior; source line 11421
\begin{proposition}[The conjugate posterior]\label{prop:posterior}
Let $F$ be a connected family. Take as prior on the tilt of the gauge the law of a question $(\mvec_0;\mathbf s_0,\mathbf t_0)$ with every proper wall positive, and observe a sample as in Proposition~\ref{prop:samplequestion}, with counts $\mvec$, sizes $\mathbf s,\mathbf t$ and the product $\theta_{\mathrm{draws}}$ of the weights of the drawn bases.
\begin{enumerate}[label=(\roman*)]
\item The posterior is the law of the question $(\mvec_0+\mvec;\mathbf s_0+\mathbf s,\mathbf t_0+\mathbf t)$: the masses add and the exponents add.
\item The chance of the sample with the tilt integrated out is
\[
\theta_{\mathrm{draws}}\cdot\frac{P_F(\mvec_0+\mvec;\mathbf s_0+\mathbf s,\mathbf t_0+\mathbf t)}{P_F(\mvec_0;\mathbf s_0,\mathbf t_0)} .
\]
\item At the posterior data $(\mvec';\mathbf s',\mathbf t')$ the chance that one more draw from the basis law of table $j$ is $T$ is $\theta_{j,T}\,P_F(\mvec'+\one_T;\mathbf s'+\mathbf e_j,\mathbf t')/P_F(\mvec';\mathbf s',\mathbf t')$, and from its cobasis law $\theta_{j,T}\,P_F(\mvec'+\one_{\cE\setminus T};\mathbf s',\mathbf t'+\mathbf e_j)/P_F(\mvec';\mathbf s',\mathbf t')$.
\item The posterior mean of the usage $\sum_j(s'_jR_j+t'_jR^\ast_j)$ at the posterior exponents is the posterior masses $\mvec_0+\mvec$.
\item Every wall of the posterior is the wall of the prior plus the excess of Proposition~\ref{prop:samplequestion}(iii): a prior inside its walls gives a posterior inside its walls for every sample, the empty one included.
\end{enumerate}
\end{proposition}

% Reference 20: prop:windowsdata; source line 11435 (draws of one table).
% Proved in the last part: prop:windowsdata.
\begin{proposition}[What the data see of the bases]\label{prop:windowsdata}
\begin{enumerate}[label=(\roman*)]
\item For a circuit $C$ of $k\ge2$ marks of a table and each $e\in C$ some basis has $T\cap C=C\setminus e$: $r(C)=k-1$, every $(k-1)$-subset of $C$ is a pattern of $C$ and $C$ is not; for $k\ge3$ every two marks of $C$ lie in a common basis, and $C$ in none.
\item For a sample of a family at a gauge $\alpha$ as in Proposition~\ref{prop:samplequestion} and $S\subseteq\cE$, the level $|T_{j,i}\cap S|$ of every basis draw of table $j$ lies in $[d_j(S),r_j(S)]$ and the colevel $|S\setminus T'_{j,i}|$ of every cobasis draw in $[d^\ast_j(S),|S|-d_j(S)]$; the chance that no basis draw of table $j$ reaches the largest level $r_j(S)$ of its window is $(1-p_j)^{s_j}$ with $p_j=\mathbb P_{j,\alpha}(|T\cap S|=r_j(S))>0$, and likewise at the floor $d_j(S)$ and for the cobasis draws at the ends of their window, with the cobasis law and $t_j$.
\end{enumerate}
\end{proposition}

% Reference 20: prop:samplelaws; source line 11445 (draws of one table).
% Proved in the last part: prop:samplelaws (the family form and (iii)).
\begin{proposition}[The exact laws of a sample of bases]\label{prop:samplelaws}
Let a sample of a family be drawn at a gauge $\alpha$ as in Proposition~\ref{prop:samplequestion}, $s_j$ basis and $t_j$ cobasis draws of table $j$, with counts $\mvec$, and let $B$ be a bunch with room $R$.
\begin{enumerate}[label=(\roman*)]
\item For a basis draw of table $j$, $\mathbb P_{j,\alpha}(T)=\pi_{j,p}(\alpha)\,\mathbb P_{j,(p),\alpha_R}(T\setminus p)$ with $p=T\cap B$, and for a cobasis draw $\mathbb P^\ast_{j,\alpha}(T)=\pi^\ast_{j,p}(\alpha)\,\mathbb P^\ast_{j,(p),\alpha_R}(T\setminus p)$, with the pattern probabilities $\pi_{j,p}(\alpha)=\alpha_B^pK_{j,(p)}(\alpha_R)/K_j(\alpha)$ and $\pi^\ast_{j,p}(\alpha)=\alpha_B^{B\setminus p}\Psi_{j,(p)}(\alpha_R)/\Psi_j(\alpha)$ (Theorem~\ref{thm:tensor}) and the basis and cobasis laws of the child $(j,p)$ at $\alpha_R$. So the numbers $a_{j,p}$ of the basis draws and $b_{j,p}$ of the cobasis draws of table $j$ that meet $B$ in $p$ are independent multinomial vectors with the chances $\pi_{j,p}(\alpha)$ and $\pi^\ast_{j,p}(\alpha)$; they are shares of $B$, $|\mathbf a_j|=s_j$ and $|\mathbf b_j|=t_j$ (Definition~\ref{def:shares}); the counts on the bunch are their loads, $m_e=\operatorname{load}_e$ for $e\in B$; and given them the parts in the room are a sample of the children $F_B(\mathbf a,\mathbf b)$ at $\alpha_R$, $a_{j,p}$ basis and $b_{j,p}$ cobasis draws of the child $(j,p)$, with the counts $\mvec|_R$. The mean shares are $s_j\pi_{j,p}(\alpha)$ and $t_j\pi^\ast_{j,p}(\alpha)$, at a native gauge the native split (Definition~\ref{def:split}).
\item Given $\mvec$, each sample with these counts has the probability $\theta_{\mathrm{draws}}$ divided by the coefficient of $\alpha^{\mvec}$ in $\prod_jK_j(\alpha)^{s_j}\Psi_j(\alpha)^{t_j}$, free of the gauge.
\item For every nonempty proper $S$ the wall $\kappa_S$ of the sample (Proposition~\ref{prop:samplequestion}(iii)) is the sum of the independent excesses of the levels of its basis draws and of the colevels of its cobasis draws over their floors, with the generating function
\[
\oE\big[x^{\kappa_S}\big]=\prod_j\sQ^K_{j,S}(x;\alpha)^{s_j}\,\sQ^\Psi_{j,S}(x;\alpha)^{t_j}
\]
of the Skeleton polynomials (Definition~\ref{def:skelpoly}). Its values fill $\{0,1,\dots,\sum_j(s_j+t_j)c_j(S)\}$; its mean is the excess $\sum_j\big(s_jR_j(\alpha)+t_jR^\ast_j(\alpha)\big)_S-\sum_j\big(s_jd_j(S)+t_jd^\ast_j(S)\big)$ of the usage at $\alpha$ over the least use, and its cumulants of order $k\ge2$ are $\partial_S^k\Phi(\log\alpha)$, $\partial_S=\sum_{e\in S}\partial_{y_e}$ (Proposition~\ref{prop:stillcoins}); it vanishes, every draw sitting at the floor of its window on $S$, with the chance $\prod_j\sQ^K_{j,S}(0;\alpha)^{s_j}\sQ^\Psi_{j,S}(0;\alpha)^{t_j}>0$; and when the Skeleton spectra of $S$ are positive it is distributed as a sum of independent $\{0,1\}$-valued variables whose means are the coins, $s_j$ copies of the coins of $\sQ^K_{j,S}$ and $t_j$ of those of $\sQ^\Psi_{j,S}$ (Proposition~\ref{prop:sklaws}(iv)).
\end{enumerate}
\end{proposition}

% Reference 20: prop:fibretests; source line 11457 (one list).
% Proved in the last part: prop:fibretests.
\begin{proposition}[Tests that do not see the gauge]\label{prop:fibretests}
For a family $F$ draw a sample at a gauge $\alpha$ as in Proposition~\ref{prop:samplequestion}, with counts $\mvec$, and let $\Theta_F(\mvec)$ be the coefficient of $\alpha^{\mvec}$ in $\prod_jK_j(\alpha)^{s_j}\Psi_j(\alpha)^{t_j}$, the total of $\theta_{\mathrm{draws}}$ over the set $\mathfrak F_F(\mvec)$ of the samples with these counts.
\begin{enumerate}[label=(\roman*)]
\item Given $\mvec$, the sample has the law $\theta_{\mathrm{draws}}/\Theta_F(\mvec)$ on $\mathfrak F_F(\mvec)$, free of the gauge; for a family of lists it is uniform, and $|\mathfrak F_F(\mvec)|=\Theta_F(\mvec)$.
\item A test of $F$ that rejects on an event $A$ whose samples with counts $\mvec$ carry at most $\varepsilon\,\Theta_F(\mvec)$ of the weight $\theta_{\mathrm{draws}}$, for every $\mvec$, has size at most $\varepsilon$ at every gauge.
\item If every basis of every table of a family $F'$ on $\cE$ is a basis of the same table of $F$, with the same weight, then given $\mvec$ the chance under $F$ that every draw is a basis of its table in $F'$ is $p_{F'|F}(\mvec)=\Theta_{F'}(\mvec)/\Theta_F(\mvec)$, and rejecting $F$ when this happens and $p_{F'|F}(\mvec)\le\varepsilon$ has size at most $\varepsilon$ at every gauge; one draw that is a basis of its table in $F$ and not in $F'$ rejects $F'$ without error.
\item If each family of a collection has a test as in (ii), the families not rejected contain the true family with probability at least $1-\varepsilon$, at every gauge.
\end{enumerate}
\end{proposition}

% Reference 20: cor:evidence; source line 11470 (lists).
% Proved in the last part: cor:evidence.
\begin{corollary}[The evidence of a family]\label{cor:evidence}
Let every family $F$ of a collection on $\cE$, each with $N$ tables, carry as prior the law of its centre of strengths $\boldsymbol\sigma_0=(\sigma_{0,j})_j$, the question $(\sigma_0\one;\boldsymbol\sigma_0,\boldsymbol\sigma_0)$ with $\sigma_0=\sum_j\sigma_{0,j}$, inside its walls. For a sample as in Proposition~\ref{prop:fibretests}, with the sizes $(s_j,t_j)_j$ and counts $\mvec$, the probability of the sample under $F$ with the gauge integrated out is
\[
\operatorname{evidence}(F)=\theta^F_{\mathrm{draws}}\;\one[\text{every draw is a basis of its table of }F]\;\frac{P_F(\sigma_0\one+\mvec;\boldsymbol\sigma_0+\mathbf s,\boldsymbol\sigma_0+\mathbf t)}{P_F(\sigma_0\one;\boldsymbol\sigma_0,\boldsymbol\sigma_0)},
\]
$\theta^F_{\mathrm{draws}}$ the product of the weights in $F$ of the drawn bases, $1$ for lists; and if the family is drawn from the collection with the chances $p(F)$, the posterior chance of $F$ is proportional to $p(F)\operatorname{evidence}(F)$.
\end{corollary}

A sample of bases asks the family a question: the counts are its masses and the sample sizes its exponents, and a wall of the question is how far the sample's use of a set exceeds the least use any basis makes of it, a sum of the independent excesses of the levels and colevels of the draws over their floors, whose generating function is the product of the Skeleton polynomials of the set raised to the sizes. Read at a bunch, a sample is a sample of the children at the shares it realizes: the loads of these shares are the counts of the bunch, and their means at a native gauge are the native split. The law of a question is the conjugate prior of the tilt, so the posterior, the predictive chances and the evidence of a family are ratios of periods, and a prior inside its walls stays inside them. Given the counts, a sample has a law that does not see the gauge, uniform for a family of lists, so a test built on the samples with those counts has a size that no gauge can change.

\section{Samples of tilts}

% Reference 20: prop:sampletilts; source line 11483
\begin{proposition}[A sample of tilts: the period is the cumulant function]\label{prop:sampletilts}
Let a question $(\mvec;\mathbf s,\mathbf t)$ of pair coordinates $(\bmu_\perp;\boldsymbol\sigma,\boldsymbol\delta)$ have every proper wall positive, let $Z_1,\dots,Z_q$ be the tilts of independent draws from $\mathrm{Breta}^\circ(F;\mvec;\mathbf s,\mathbf t)$, and let $\bar Z$, $\bar E_j$, $\bar O_j$ be the sample means of $Z_i$, $E_j(Z_i)$, $O_j(Z_i)$. In this statement $\oE$ is the expectation of $U\sim\mathrm{Breta}^\circ(F;\mvec;\mathbf s,\mathbf t)$ and $Z$ is its tilt.
\begin{enumerate}[label=(\roman*)]
\item The sample has the density
\[
2^{q(n-1)}\exp\Big(q\Big[2\ip{\bmu_\perp}{\bar Z}-\sum_j\big(\sigma_j\bar E_j+\delta_j\bar O_j\big)\Big]\Big)\Big/P_F(\mvec;\mathbf s,\mathbf t)^q
\]
on $\mathcal T^q$, so $(\bar Z;\bar E_j,\bar O_j)_j$ is sufficient for the question.
\item The means are $\oE[Z]=\frac12\nabla_{\bmu_\perp}\log P_F(\mvec;\mathbf s,\mathbf t)$, $\oE[E_j(Z)]=-\partial_{\sigma_j}\log P_F(\mvec;\mathbf s,\mathbf t)$ and $\oE[O_j(Z)]=-\partial_{\delta_j}\log P_F(\mvec;\mathbf s,\mathbf t)$, the derivatives of Definition~\ref{def:question}, and the likelihood equations equate them to the sample means.
\item Against a second question $(\mvec';\mathbf s',\mathbf t')$ of pair coordinates $(\bmu'_\perp;\boldsymbol\sigma',\boldsymbol\delta')$ the logarithm of the likelihood ratio is
\[
q\Big[2\ip{\bmu'_\perp-\bmu_\perp}{\bar Z}-\sum_j\big((\sigma'_j-\sigma_j)\bar E_j+(\delta'_j-\delta_j)\bar O_j\big)\Big]-q\log\frac{P_F(\mvec';\mathbf s',\mathbf t')}{P_F(\mvec;\mathbf s,\mathbf t)} .
\]
\item The sums $q\bar Z$, $q\bar E_j$, $q\bar O_j$ have as Laplace transform that of Corollary~\ref{cor:letterstransform} raised to the power $q$.
\end{enumerate}
\end{proposition}

% Reference 20: def:pseudoMGF; source line 11496
\begin{definition}[The Foreman's question; the pseudo-MGF]\label{def:pseudoMGF}
Let $\omega=W\,d\alpha_{\mvec}$ have real data on a connected family, $s_j,t_j\ge0$ not all zero and $\mvec>0$, with target $\btau$. The \emph{Foreman's question} of $\omega$ is the question $(\btau;\mathbf s,\mathbf t)$ at its target, of pair coordinates $\big(\btau-\frac{\tau_{\cE}}n\one;\boldsymbol\sigma,\boldsymbol\delta\big)$; off the pole hyperplanes its period $P_F(\btau;\mathbf s,\mathbf t)$ is the pseudointegral $\pseudoint W\,d\alpha_{\mvec}$ (Definition~\ref{def:threeE}(iii)). Where the target lies on no pole hyperplane and $P_F(\btau;\mathbf s,\mathbf t)\neq0$, the \emph{pseudo-MGF} of $\omega$ is
\[
\mathsf M_\omega(c;\mathbf a,\mathbf b)=\frac{P_F\big(\btau+\mvec_F(c;\mathbf a,\mathbf b);\,\mathbf s+\mathbf a+\mathbf b,\,\mathbf t+\mathbf a-\mathbf b\big)}{P_F(\btau;\mathbf s,\mathbf t)}\qquad(c\in\one^\perp\otimes\mathbb C,\ \mathbf a,\mathbf b\in\mathbb C^N),
\]
continued meromorphically (Lemma~\ref{lem:periodmero}), and the \emph{pseudo-cumulants} of $\omega$ are the derivatives of $\log\mathsf M_\omega$ at $0$: the derivatives of $\log P_F$ at the Foreman's question in its pair coordinates.
\end{definition}

% Reference 20: prop:pseudoMGF; source line 11504
\begin{proposition}[The pseudo-MGF inside the walls]\label{prop:pseudoMGF}
In Definition~\ref{def:pseudoMGF} let $\omega$ have a native gauge, and let $Z$ be the tilt of $U\sim\mathrm{Breta}^\circ(F;\btau;\mathbf s,\mathbf t)$, the Foreman's law of Proposition~\ref{prop:centrelaw}(i); the law of $Z$ is the law of the Foreman's question.
\begin{enumerate}[label=(\roman*)]
\item Wherever the real parts of the shifted data $\big(\btau+\mvec_F(c;\mathbf a,\mathbf b);\,\mathbf s+\mathbf a+\mathbf b,\,\mathbf t+\mathbf a-\mathbf b\big)$ have every proper wall positive, a region containing a neighbourhood of $0$, $\mathsf M_\omega(c;\mathbf a,\mathbf b)$ is the moment generating function of the statistics $(2Z;-E_j(Z),-O_j(Z))_j$ under the Foreman's law (Corollary~\ref{cor:letterstransform} at the Foreman's question); the pseudo-cumulants are their joint cumulants.
\item For a balanced shift $(\dot\mvec;\dot{\mathbf s},\dot{\mathbf t})$ the function $U^{\dot\mvec}\prod_jK_j(U)^{-\dot s_j}\Psi_j(U)^{-\dot t_j}$ has degree $0$, a function of the ray of $U$, and where the shifted data have every proper wall positive its expectation under the Foreman's law is $P_F(\btau+\dot\mvec;\mathbf s+\dot{\mathbf s},\mathbf t+\dot{\mathbf t})/P_F(\btau;\mathbf s,\mathbf t)$ (the case of Proposition~\ref{prop:momentsteps}(i) below at the Foreman's question); $\mathsf M_\omega(c;\mathbf a,\mathbf b)$ is this ratio at the shift of pair coordinates $(c;\mathbf a,\mathbf b)$, so $\mathsf M_\omega$ is the moment generating function of $(\log U;-\log K_j(U),-\log\Psi_j(U))_j$ along the balanced shifts.
\item The normaliser is $P_F(\btau;\mathbf s,\mathbf t)=P_{\mathrm c}\,\mathsf E_{U\sim\mathbb P_{\mathrm c}}[(\omega_{\btau}/\omega_{\mathrm c})(U)]$ (Theorem~\ref{thm:normalizer}(i)). The centre $\omega_{\mathrm c}$ has the target $\sigma\one$ and the Foreman's question $(\sigma\one;\boldsymbol\sigma,\boldsymbol\sigma)$, and the pseudoexpectation is the pseudo-MGF of the centre at the departure:
\[
\mathsf E_{U\sim\mathbb P_{\mathrm c}}\Big[\frac{\omega_{\btau}}{\omega_{\mathrm c}}(U)\Big]=\mathsf M_{\omega_{\mathrm c}}\Big(\btau-\frac{\tau_{\cE}}n\one;\ \mathbf 0,\ \boldsymbol\delta\Big)=\oE_{U\sim\mathbb P_{\mathrm c}}\Big[e^{2\ip{\btau-\frac{\tau_{\cE}}n\one}{Z}-\sum_j\delta_jO_j(Z)}\Big],
\]
$Z$ the tilt of $U$: at fixed strengths the pseudoexpectations of all the questions are values of one moment generating function, that of $(2Z;-O_j(Z))_j$ under the law of the centre.
\end{enumerate}
Past the walls $\mathsf M_\omega$ is the continuation, as the moments of the distribution of Theorem~\ref{thm:offdomain}(i) are, and $\mathsf M_\omega(0)=1$ throughout.
\end{proposition}

% Reference 20: prop:twolayers; source line 11529
\begin{proposition}[Two layers: the word evaluated at the draws]\label{prop:twolayers}
Let a family have real data with $\mvec>0$, and let $U_1,\dots,U_q$ be independent draws from $\Dir_{\mvec}$, the shapes $\alpha/|\alpha|$ of $\alpha\sim\operatorname{Gamma}_{\mvec}$ (Proposition~\ref{prop:gammadraws}(ii)).
\begin{enumerate}[label=(\roman*)]
\item By Proposition~\ref{prop:gammadraws}(i), $\oE_{U\sim\Dir_{\mvec}}[W(U)]=P_F(\mvec;\mathbf s,\mathbf t)\big/\Gst{n}{1}{\mvec}{m_{\cE}}$, and it is finite exactly when every proper wall at $(\mvec;\mathbf s,\mathbf t)$ is positive.
\item $\oE_{U\sim\Dir_{\mvec}}[W(U)^k]=P_F(\mvec;k\mathbf s,k\mathbf t)\big/\Gst{n}{1}{\mvec}{m_{\cE}}$, and the sample mean $\overline W_q=\frac1q\sum_iW(U_i)$ has the mean $P_F(\mvec;\mathbf s,\mathbf t)/\Gst{n}{1}{\mvec}{m_{\cE}}$ and the variance
\[
\frac1q\bigg[\frac{P_F(\mvec;2\mathbf s,2\mathbf t)}{\Gst{n}{1}{\mvec}{m_{\cE}}}-\bigg(\frac{P_F(\mvec;\mathbf s,\mathbf t)}{\Gst{n}{1}{\mvec}{m_{\cE}}}\bigg)^2\bigg],
\]
finite exactly when every proper wall at $(\mvec;2\mathbf s,2\mathbf t)$ is positive.
\item Where every proper wall at $(\mvec;\mathbf s,\mathbf t)$ is positive, $\mathrm{Breta}^\circ(F;\mvec;\mathbf s,\mathbf t)$ has the density $W\,\Gst{n}{1}{\mvec}{m_{\cE}}/P_F(\mvec;\mathbf s,\mathbf t)$ with respect to $\Dir_{\mvec}$, and under it a second word $W'=\prod_jK_j^{-s'_j}\Psi_j^{-t'_j}$ of the same tables has $\oE_{U\sim\mathrm{Breta}^\circ(F;\mvec;\mathbf s,\mathbf t)}[W'(U)^k]=P_F(\mvec;\mathbf s+k\mathbf s',\mathbf t+k\mathbf t')/P_F(\mvec;\mathbf s,\mathbf t)$.
\end{enumerate}
\end{proposition}

% Reference 20: def:greylaw; source line 11541
\begin{definition}[The Grey law]\label{def:greylaw}
For a family $F$ and masses $\mvec>0$ the \emph{Grey law} $\mathrm{Grey}(F;\mvec)$ is the joint law of
\[
(G_j,G^\ast_j)_j=\big(-\log K_j(U),\,-\log\Psi_j(U)\big)_j\qquad\text{for }U\sim\Dir_{\mvec};
\]
it depends on the tables of $F$ and on $\mvec$, and not on the exponents.
\end{definition}

% Reference 20: prop:greylaw; source line 11545 ((iii) as a limit superior, for $a_j,b_j\ge0$).
% Proved in the last part: prop:greylaw(iii).
\begin{proposition}[The Grey law and its exponential family]\label{prop:greylaw}
Let $\mvec>0$ and let $(G_j,G^\ast_j)_j\sim\mathrm{Grey}(F;\mvec)$.
\begin{enumerate}[label=(\roman*)]
\item For complex exponents $(\mathbf s,\mathbf t)$ whose real parts keep every proper wall at $\mvec$ positive,
\[
\oE_{(G_j,G^\ast_j)_j\sim\mathrm{Grey}(F;\mvec)}\Big[e^{\sum_j(s_jG_j+t_jG^\ast_j)}\Big]=\frac{P_F(\mvec;\mathbf s,\mathbf t)}{\Gst{n}{1}{\mvec}{m_{\cE}}}\quad\text{(Proposition~\ref{prop:gammadraws}(i),(iii))}:
\]
on real exponents the moment generating function, on imaginary exponents the characteristic function; the joint cumulants are the derivatives of $\log P_F(\mvec;\mathbf s,\mathbf t)$ in the exponents at $(\mathbf 0,\mathbf 0)$.
\item For real exponents with every proper wall positive, under $\mathrm{Breta}^\circ(F;\mvec;\mathbf s,\mathbf t)$ the vector $\big(-\log K_j(U),-\log\Psi_j(U)\big)_j$ has the law
\[
e^{\sum_j(s_jg_j+t_jg^\ast_j)}\,\frac{\Gst{n}{1}{\mvec}{m_{\cE}}}{P_F(\mvec;\mathbf s,\mathbf t)}\,\mathrm{Grey}(F;\mvec)(dg\,dg^\ast):
\]
over the exponents these laws form the exponential family with base $\mathrm{Grey}(F;\mvec)$ and cumulant function $\log\big(P_F(\mvec;\mathbf s,\mathbf t)/\Gst{n}{1}{\mvec}{m_{\cE}}\big)$.
\item For a real direction $(\mathbf a,\mathbf b)\neq0$ of the exponents let $g(S)=\sum_j\big(a_jd_j(S)+b_jd^\ast_j(S)\big)$ and $k^\ast=\min\{m_S/g(S):\ S\text{ nonempty proper},\ g(S)>0\}$, the value of $\lambda$ at which the ray $\lambda(\mathbf a,\mathbf b)$, $\lambda>0$, first meets a wall at $\mvec$, with $k^\ast=\infty$ if no $g(S)$ is positive. If $k^\ast<\infty$,
\[
\mathbb P\Big(\sum_j\big(a_jG_j+b_jG^\ast_j\big)>x\Big)=e^{-k^\ast x+o(x)}\qquad(x\to\infty),
\]
and if $k^\ast=\infty$ the sum is bounded above.
\item For $q$ independent copies the sum has the Laplace transform $\big(P_F(\mvec;\mathbf s,\mathbf t)/\Gst{n}{1}{\mvec}{m_{\cE}}\big)^q$, and the sample mean $\overline W_q$ of Proposition~\ref{prop:twolayers}(ii) is the empirical counterpart of the transform of (i) at the exponents $(\mathbf s,\mathbf t)$.
\end{enumerate}
\end{proposition}

% Reference 20: cor:greytarget; source line 11564
\begin{corollary}[The Grey law at the target]\label{cor:greytarget}
Let a connected family carry real exponents $s_j,t_j\ge0$, not all zero, and masses $\mvec>0$ with a native gauge and the target $\btau$. Then
\begin{align*}
\oE_{U\sim\Dir_{\btau}}[W(U)]&=\oE_{(G_j,G^\ast_j)_j\sim\mathrm{Grey}(F;\btau)}\Big[e^{\sum_j(s_jG_j+t_jG^\ast_j)}\Big]=\frac{P_F(\btau;\mathbf s,\mathbf t)}{\Gst{n}{1}{\btau}{\tau_{\cE}}}\\
&=\frac{P_{\mathrm c}}{\Gst{n}{1}{\btau}{\tau_{\cE}}}\,\mathsf E_{U\sim\mathbb P_{\mathrm c}}\Big[\frac{\omega_{\btau}}{\omega_{\mathrm c}}(U)\Big]:
\end{align*}
the moment generating function of the Grey law at the target, read at the exponents, is the pseudoexpectation times the period of the centre divided by $\Gst{n}{1}{\btau}{\tau_{\cE}}$. With $\btau$ held at the target, the logarithmic derivatives of $(\mathbf s',\mathbf t')\mapsto P_F(\btau;\mathbf s',\mathbf t')$ at $(\mathbf s,\mathbf t)$ are the joint cumulants of $\big(-\log K_j(U),-\log\Psi_j(U)\big)_j$ under the Foreman's law $\mathrm{Breta}^\circ(F;\btau;\mathbf s,\mathbf t)$.
\end{corollary}

% Reference 20: prop:movingtarget; source line 11575
\begin{proposition}[A Laplace transform with a moving target]\label{prop:movingtarget}
Let a connected family carry real exponents $s_j,t_j\ge0$, not all zero, and masses $\mvec>0$ whose target $\btau$ has every proper wall positive, and let $g(S)=\sum_j\big(s_jd_j(S)+t_jd^\ast_j(S)\big)$ be the set function of Theorem~\ref{thm:saddle}. With the tails $U_k$ of the orders $\pi$ of $\cE$ (Lemma~\ref{lem:greedy}),
\[
\lim_{\lambda\to0^+}\oE_{V\sim\Dir_{\lambda\btau}}\Big[\prod_jK_j(V)^{-\lambda s_j}\Psi_j(V)^{-\lambda t_j}\Big]=\frac{\sum_\pi\prod_{k=2}^n\big(\tau_{U_k}-g(U_k)\big)^{-1}}{\sum_\pi\prod_{k=2}^n\tau_{U_k}^{-1}}\ \ge\ 1,
\]
with equality exactly when $g(S)=0$ for every nonempty proper $S$. The limit is the ratio of the limits of Theorem~\ref{thm:severity}(i) for the question of data $(\btau;\mathbf s,\mathbf t)$ and for the list $\binom{\cE}1$ at a full basis lean at the masses $\btau$, whose walls are the $\tau_S$ (Proposition~\ref{prop:sectiontable}).
\end{proposition}

The word read at Dirichlet draws is a random variable whose expectation, as a function of the exponents, is the period divided by the stacked Gamma of the masses over their total: the Laplace transform of the Grey law of the family. At fixed masses the Grey law is the base of an exponential family over the exponents, and its tail along a direction of exponents is set by the first wall met along that direction. Read at the target of each exponent, as the normaliser $\binom{\cE}1$ of Proposition~\ref{prop:completionlaws}(ii) reads it, the expectation of the word is not a moment generating function in the exponents: its limit at the origin exceeds $1$ as soon as the word has a floor on a proper set, since at each exponent it reads the Grey law at another target. The pseudo-MGF fixes the target once, at the Foreman's question, and the pseudoexpectation fixes the law once for given strengths, at the centre.

% Reference 20: prop:dial; source line 11591 ($s_j,t_j\ge0$).
% Proved in the last part: prop:dial.
\begin{proposition}[The exponent is a dial between the mean and the maximum]\label{prop:dial}
In Proposition~\ref{prop:twolayers} let the exponents be real, not all zero, and put $Y_i=\log W(U_i)$ and $Y=\log W(U)$ for $U\sim\Dir_{\mvec}$; in this statement $\oE$ is the expectation of independent $U,U_1,\dots,U_q\sim\Dir_{\mvec}$. Let $g(S)=\sum_j\big(s_jd_j(S)+t_jd^\ast_j(S)\big)$, $k^\ast=\min\{m_S/g(S):\ S\text{ nonempty proper},\ g(S)>0\}$, with $k^\ast=\infty$ if no $g(S)$ is positive, and for $k>0$ let $A_k=\frac1k\log\big(\frac1q\sum_ie^{kY_i}\big)$.
\begin{enumerate}[label=(\roman*)]
\item $A_k$ is nondecreasing in $k$, tends to the sample mean of the $Y_i$ as $k\to0$ and to $\max_iY_i$ as $k\to\infty$, and $\max_iY_i-\frac{\log q}k\le A_k\le\max_iY_i$.
\item For $k<k^\ast$, $\oE[e^{kA_k}]=P_F(\mvec;k\mathbf s,k\mathbf t)\big/\Gst{n}{1}{\mvec}{m_{\cE}}$, and $e^{kA_k}$ has a finite variance exactly when $2k<k^\ast$.
\item For $k\ge k^\ast$, $\oE[e^{kY}]=\infty$; and $\mathbb P(\max_iY_i\le x)=\mathbb P(Y\le x)^q$.
\item If $k^\ast<\infty$, $\mathbb P(Y>x)=e^{-k^\ast x+o(x)}$ as $x\to\infty$, and if $k^\ast=\infty$, $Y$ is bounded above (Proposition~\ref{prop:greylaw}(iii)). For $s_j,t_j\ge0$ and a set $S$ attaining $k^\ast$ the scale $\rho=|U_S|$ has $(\rho,1-\rho)\sim\Dir_{(m_S,\,m_{\cE\setminus S})}$, so $\rho$ has density proportional to $\rho^{m_S-1}(1-\rho)^{m_{\cE\setminus S}-1}$, and $W(U)\ge\prod_jK_j(\one)^{-s_j-t_j}\rho^{-g(S)}$.
\end{enumerate}
\end{proposition}

% Reference 20: prop:greycentre; source line 11604 (one table).
% Proved in the last part: prop:greycentre.
\begin{proposition}[The even and odd statistics from a centre]\label{prop:greycentre}
Let a family carry the centre of strengths $\boldsymbol\sigma=(\sigma_j)_j$, $\sigma_j\ge0$, inside its walls, let $Z$ be the tilt of $U\sim\mathrm{Breta}^\circ(F;\sigma\one;\boldsymbol\sigma,\boldsymbol\sigma)$, with density proportional to $e^{-\sum_j\sigma_jE_j(z)}$ on $\mathcal T$, and let $\ell^\ast_{\boldsymbol\sigma}$ be the slope of Proposition~\ref{prop:reach}. In this statement $\oE$ is the expectation of $U\sim\mathrm{Breta}^\circ(F;\sigma\one;\boldsymbol\sigma,\boldsymbol\sigma)$.
\begin{enumerate}[label=(\roman*)]
\item $\sum_j\sigma_jE_j(Z)\ge2\sum_j\sigma_j\log K_j(\one)$, with equality only at $Z=0$ (Proposition~\ref{prop:evenfloor}), where $\sum_j\sigma_jE_j$ has the Hessian $8\sum_j\sigma_j\big(\kappa_{j,\one}(X_e,X_f)\big)_{e,f}$, rung two of the still of the realization $\big((\sigma_j,\sigma_j)_j;\one\big)$ of the active tables, eight times the covariance of the indicators under the weight laws; and each $E_j(Z)\ge2\log K_j(\one)$, with equality exactly where $Z$ is constant on the molecules of table $j$. For real $\mathbf k=(k_j)_j$,
\[
\oE\Big[e^{-\sum_jk_jE_j(Z)}\Big]=\frac{P_F\big((\sigma+|\mathbf k|)\one;\boldsymbol\sigma+\mathbf k,\boldsymbol\sigma+\mathbf k\big)}{P_F(\sigma\one;\boldsymbol\sigma,\boldsymbol\sigma)},
\]
the ratio of the periods of the centres of strengths $\boldsymbol\sigma+\mathbf k$ and $\boldsymbol\sigma$, finite exactly when $\sum_j(\sigma_j+k_j)c_j(S)>0$ for every nonempty proper $S$: along the even direction $\mathbf k=\lambda\boldsymbol\sigma$ every wall closes at once, at $\lambda=-1$, and along any other even direction $\lambda\mathbf k$, $\lambda>0$, the first wall closes at the least of the ratios $\kappa^{\mathbf E}_S(\boldsymbol\sigma)/(-\kappa^{\mathbf E}_S(\mathbf k))$ with $\kappa^{\mathbf E}_S(\mathbf k)<0$.
\item $\big(Z,(E_j(Z))_j,(O_j(Z))_j\big)$ and $\big(-Z,(E_j(Z))_j,(-O_j(Z))_j\big)$ have one law, so every joint moment of odd total degree in $Z$ and the $O_j(Z)$ vanishes. For real $c\in\one^\perp$ and $\mathbf b\in\mathbb R^N$, $\oE\big[e^{2\ip cZ-\sum_jb_jO_j(Z)}\big]$ is finite exactly when $\ell^\ast_{\boldsymbol\sigma}(c,\mathbf b)>1$ (Proposition~\ref{prop:reach}(i)), where it is
\[
P_F\big(\mvec_F(c;\boldsymbol\sigma,\mathbf b);\,\boldsymbol\sigma+\mathbf b,\,\boldsymbol\sigma-\mathbf b\big)\big/P_F(\sigma\one;\boldsymbol\sigma,\boldsymbol\sigma),
\]
for $|b_j|\le\sigma_j$ the pseudoexpectation of the question of pair coordinates $(c;\boldsymbol\sigma,\mathbf b)$. If $\ell^\ast_{\boldsymbol\sigma}(c,\mathbf b)<\infty$, both tails of $Y=2\ip cZ-\sum_jb_jO_j(Z)$ have the rate $\ell^\ast_{\boldsymbol\sigma}(c,\mathbf b)$: $\mathbb P(\pm Y>x)=e^{-\ell^\ast_{\boldsymbol\sigma}(c,\mathbf b)x+o(x)}$ as $x\to\infty$; if it is infinite, $Y$ is bounded (Proposition~\ref{prop:tails}).
\item At $\alpha=e^{2z}$, $K_j(\alpha)^{-s_j}=e^{-s_j\ip\one z}e^{-\frac{s_j}2(E_j(z)+O_j(z))}$ and $\Psi_j(\alpha)^{-t_j}=e^{-t_j\ip\one z}e^{-\frac{t_j}2(E_j(z)-O_j(z))}$: every letter is its table at a full lean, $K_j^{-s_j}$ at the pair coordinates $(\sigma,\delta)=(\frac{s_j}2,\frac{s_j}2)$ and $\Psi_j^{-t_j}$ at $(\frac{t_j}2,-\frac{t_j}2)$, with the masses shifted by $-\frac{s_j}2\one$, resp.\ $-\frac{t_j}2\one$; a family is the family of its letters, and its pair coordinates are the sums of theirs.
\end{enumerate}
\end{proposition}

% Reference 20: prop:twostatistics; source line 11616 (one table).
% Proved in the last part: prop:twostatistics.
\begin{proposition}[What the two statistics are]\label{prop:twostatistics}
For table $j$ of a family and a tilt $z$, $\log\Lambda_j(z)=\log\sum_T\theta_{j,T}e^{\ip{\chi_T}z}$ is a soft maximum and $-\log\Lambda_j(-z)$ a soft minimum of the scores $\ip{\chi_T}z$ of its bases, weighted by $\theta_{j,T}$.
\begin{enumerate}[label=(\roman*)]
\item $E_j(z)$ is the soft maximum less the soft minimum, and $O_j(z)$ is their sum; the centred potential $\Phi^\circ=\sum_j(\sigma_jE_j+\delta_jO_j)$ is the sum over the tables of $s_j$ times the soft maximum less $t_j$ times the soft minimum.
\item With $Y_j=\ip{\chi_T}z$ for $T$ under the weight law of table $j$, near $z=0$
\begin{gather*}
E_j(z)=2\log K_j(\one)+\sum_{k\text{ even}}\frac2{k!}\,\kappa_{j,\one}(Y_j^{\times k}),\\
O_j(z)=\sum_{k\text{ odd}}\frac2{k!}\,\kappa_{j,\one}(Y_j^{\times k}),
\end{gather*}
so $E_j(z)=2\log K_j(\one)+\kappa_{j,\one}(Y_j,Y_j)+\frac1{12}\kappa_{j,\one}(Y_j^{\times4})+\cdots$ and $O_j(z)=2\kappa_{j,\one}(Y_j)+\frac13\kappa_{j,\one}(Y_j^{\times3})+\cdots$, with $\kappa_{j,\one}(Y_j)=\sum_e(2R_{j,e}(\one)-1)z_e$ and, for $k\ge2$, $\kappa_{j,\one}(Y_j^{\times k})=2^k\sum_{w\in\cE^k}\kappa_{j,\one}(w)\,z^w$, $z^w=z_{w_1}\cdots z_{w_k}$. Summed with the strengths and the leans, the terms of order $k\ge2$ of $\Phi^\circ$ are $\frac1{k!}\sum_{w\in\cE^k}2^k\sum_j\big(s_j+(-1)^kt_j\big)\kappa_{j,\one}(w)\,z^w$, which for the exponents of a realization $\Upsilon=((s_j,t_j)_j;\one)$ is $\frac1{k!}\sum_w\kappa_\Upsilon(w)\,z^w$, rung $k$ of its still (Proposition~\ref{prop:stillgauge}): the even statistics carry the even rungs and the odd statistics the odd rungs.
\item As $\lambda\to\infty$, $E_j(\lambda z)/\lambda\to\operatorname{wd}_j(z)$, the width of Proposition~\ref{prop:evenfloor}, and $O_j(\lambda z)/\lambda\to h_j(z)-h_j(-z)$, the maximum of the scores being at the greedy basis $T^+_{j,\pi}$ and the minimum at $T^-_{j,\pi}$ for an order $\pi$ along which the coordinates of $z$ decrease (Lemma~\ref{lem:greedy}). So $\Phi^\circ(\lambda z)/\lambda\to\ip{2v_\pi-2\sigma\one}z$ with $v_\pi=\sum_j\big(s_j\one_{T^+_{j,\pi}}+t_j\one_{\cE\setminus T^-_{j,\pi}}\big)$, for $s_j,t_j\ge0$ a vertex of the Newton sum at which $\ip{\cdot}z$ is largest, the \emph{greedy vertex} of $\pi$; and for a question the exponent of Proposition~\ref{prop:evenfloor} has $X(\lambda z)/\lambda\to-2\kappa(z)=2\ip{\mvec-v_\pi}z$ on $\one^\perp$: the walls of the tails of $\pi$ are the tail sums of the masses less the greedy vertex, $\kappa_{U_k}=m_{U_k}-(v_\pi)_{U_k}$.
\end{enumerate}
\end{proposition}

% Reference 20: prop:expfamily; source line 11628
\begin{proposition}[The laws of the questions are an exponential family]\label{prop:expfamily}
Let a connected family be such that $1$, the coordinates $z_e$, $e\neq g$, of a chart $z_g=0$ of $\mathcal T$ and the statistics $E_j$, $O_j$ are linearly independent functions on $\mathcal T$. Then the laws of its questions with every proper wall positive form a regular exponential family of dimension $n-1+2N$, with natural parameter $(2\bmu_\perp;-\boldsymbol\sigma,-\boldsymbol\delta)$, statistic $(z;E_j(z),O_j(z))_j$ and cumulant function $\log P_F\big(\mvec_F(\bmu_\perp;\boldsymbol\sigma,\boldsymbol\delta);\boldsymbol\sigma+\boldsymbol\delta,\boldsymbol\sigma-\boldsymbol\delta\big)$. For the sample means $(\bar Z;\bar E_j,\bar O_j)_j$ of $q$ independent tilts, as in Proposition~\ref{prop:sampletilts}:
\begin{enumerate}[label=(\roman*)]
\item $(\bar Z;\bar E_j,\bar O_j)_j$ is sufficient, and a function of it with expectation zero under every law of the family vanishes almost surely;
\item it is the unbiased estimator of least variance of the mean parameters
\[
\Big(\frac12\nabla_{\bmu_\perp}\log P_F;\ -\partial_{\sigma_j}\log P_F,\ -\partial_{\delta_j}\log P_F\Big)_j ,
\]
with covariance $\frac1q$ times the Hessian of $\log P_F$ in the natural parameter;
\item the likelihood equations have a solution, unique, exactly when $(\bar Z;\bar E_j,\bar O_j)_j$ is interior to the convex hull of the image of $z\mapsto(z;E_j(z),O_j(z))_j$;
\item for $\delta_j\le\delta_0$ against $\delta_j>\delta_0$, with the other parameters unknown, rejecting for small $\bar O_j$ given the other coordinates of the sample means is uniformly most powerful among unbiased tests, with a conditional null law free of the other parameters; likewise for $\sigma_j$ and $\bar E_j$;
\item the law of the sums $q\bar Z$, $q\bar E_j$, $q\bar O_j$ has the Laplace transform of Proposition~\ref{prop:sampletilts}(iv), so the power of every test is a computation in the family.
\end{enumerate}
When the complement of every basis of table $j$ is a basis with the same weight, $O_j\equiv0$ and the statements hold with $\delta_j$ removed.
\end{proposition}

A word read at Dirichlet draws is a dial: its power means run from the mean of the logarithms to their maximum as the exponent grows, and they keep a finite expectation up to the first wall met along the ray of exponents, where the tail of the logarithm takes over with the rate of that wall. At a centre the centred potential $\sum_j\sigma_jE_j$ is least at the table's gauge and its exponential moments along the strengths close every wall at once, while the odd statistics are jointly symmetric and every odd direction has both tails at the rate of its slope. The two statistics of a table are the soft maximum and the soft minimum of the scores of its bases, read as their difference and their sum: at the table's gauge their Taylor terms are the rungs of the still, the even ones carried by the strengths and the odd ones by the leans, and at infinity they read the greedy vertex of the Newton sum, whose tail sums against the masses are the walls. The laws of the questions are a regular exponential family in the tilt and these statistics, so the classical optimal estimators and tests are computations with periods.

\section{Draws, orders and races}

% Reference 20: prop:sortedrow; source line 11649 (a connected family with $s_j,t_j\ge0$).
% Proved in the last part: prop:sortedrow.
\begin{proposition}[A sorted draw lives on the sector of its order]\label{prop:sortedrow}
Let a family have real data on its balance wall with every proper wall positive, and let $U\sim\mathrm{Breta}^\circ(F;\mvec;\mathbf s,\mathbf t)$. For an order $\pi$ of $\cE$ let $\kappa_k$ be the walls of its tails and $I_\pi$ its sector integral (Corollary~\ref{cor:pastwalls}), so that $P_F(\mvec;\mathbf s,\mathbf t)=\sum_\pi I_\pi$ (Theorem~\ref{thm:exist}(ii)).
\begin{enumerate}[label=(\roman*)]
\item $\mathbb P\big(U_{\pi(1)}>\cdots>U_{\pi(n)}\big)=I_\pi/P_F(\mvec;\mathbf s,\mathbf t)$.
\item Given that order, the gaps $u_k=\log\big(U_{\pi(k-1)}/U_{\pi(k)}\big)$, $2\le k\le n$, have the density
\[
I_\pi^{-1}\exp\Big(-\sum_{k=2}^n\kappa_ku_k\Big)\prod_jk_{j,\pi}(e^{-u})^{-s_j}\psi_{j,\pi}(e^{-u})^{-t_j}
\]
on $[0,\infty)^{n-1}$, with the sector letters $k_{j,\pi},\psi_{j,\pi}$ of Theorem~\ref{thm:exist}(ii).
\item $\theta_{j,T^+_{j,\pi}}\le k_{j,\pi}\le K_j(\one)$ and $\theta_{j,T^-_{j,\pi}}\le\psi_{j,\pi}\le K_j(\one)$ on the cube, so the density of (ii) lies between constant multiples of the density of independent exponential gaps of the rates $\kappa_k$; for $s_j,t_j\ge0$, $I_\pi$ lies between $\prod_jK_j(\one)^{-s_j-t_j}\prod_k\kappa_k^{-1}$ and $\prod_j\theta_{j,T^+_{j,\pi}}^{-s_j}\theta_{j,T^-_{j,\pi}}^{-t_j}\prod_k\kappa_k^{-1}$.
\item On the sector of $\pi$ the largest monomial $U^T$ over the bases of table $j$ is at its greedy basis $T^+_{j,\pi}$ and the smallest at $T^-_{j,\pi}$ (Lemma~\ref{lem:greedy}): the best and the worst basis of a draw are functions of its order, and the best basis of table $j$ at $U$ is $T$ with probability $\sum_{\pi:\,T^+_{j,\pi}=T}I_\pi/P_F(\mvec;\mathbf s,\mathbf t)$.
\end{enumerate}
\end{proposition}

% Reference 20: prop:orderedsample; source line 11662
\begin{proposition}[An ordered sample and a cut]\label{prop:orderedsample}
Let $Y_1,\dots,Y_q$ be independent with a probability law $\phi$ on $\mathbb R$, an input (Definition~\ref{def:inputs}) that charges no point.
\begin{enumerate}[label=(\roman*)]
\item The ordered sample has the law $q!\,\phi^{\otimes q}$ on the region $\{y_1<\cdots<y_q\}$.
\item The number of observations above a cut $N$ is binomial with $q$ trials and the chance $\phi((N,\infty))$, which for $\phi$ living on $(0,\infty)$ and $N\ge0$ is the Trailmoment $\Tx{0}{\phi}{N}$ at index $0$; given that it is $d$, the ordered observations above the cut have the law $d!\,\phi^{\otimes d}/\phi((N,\infty))^d$ on the trail region $X_N=\{N<y_1<\cdots<y_d\}$.
\item If $\phi$ lives on $(0,\infty)$ and $N\ge0$, the empirical input $\hat\phi=\frac1q\sum_i\delta_{Y_i}$ has $\Tx{\nu}{\hat\phi}{N}=\frac1q\sum_{Y_i>N}Y_i^{-\nu}$, with the mean $\Tx{\nu}{\phi}{N}$ and the variance
\[
\frac1q\Big[\Tx{2\nu}{\phi}{N}-\Tx{\nu}{\phi}{N}^2\Big].
\]
\end{enumerate}
\end{proposition}

% Reference 20: prop:dilation; source line 11674
\begin{proposition}[The races and the dilation relation]\label{prop:dilation}
For table $j$ of a family and a tilt $z$, with $\alpha=e^{2z}$, the laws of the winners of the basis race and the cobasis race (Theorem~\ref{thm:race}) are $\mathbb P_{j,\alpha}$ and $\mathbb P^\ast_{j,\alpha}$; let
\[
\operatorname{Ent}_j(z)=-\sum_T\mathbb P_{j,\alpha}(T)\log\frac{\mathbb P_{j,\alpha}(T)}{\theta_{j,T}},\qquad \operatorname{Ent}^\ast_j(z)=-\sum_T\mathbb P^\ast_{j,\alpha}(T)\log\frac{\mathbb P^\ast_{j,\alpha}(T)}{\theta_{j,T}}
\]
be their entropies relative to the weights.
\begin{enumerate}[label=(\roman*)]
\item $E_j(z)=\ip z{\nabla E_j(z)}+\operatorname{Ent}_j(z)+\operatorname{Ent}^\ast_j(z)$ and $O_j(z)=\ip z{\nabla O_j(z)}+\operatorname{Ent}_j(z)-\operatorname{Ent}^\ast_j(z)$, with
\begin{gather*}
\ip z{\nabla E_j(z)}=\oE_{T\sim\mathbb P_{j,\alpha}}\big[\ip{\chi_T}z\big]-\oE_{T\sim\mathbb P^\ast_{j,\alpha}}\big[\ip{\chi_T}z\big],\\
\ip z{\nabla O_j(z)}=\oE_{T\sim\mathbb P_{j,\alpha}}\big[\ip{\chi_T}z\big]+\oE_{T\sim\mathbb P^\ast_{j,\alpha}}\big[\ip{\chi_T}z\big].
\end{gather*}
\end{enumerate}
Let the family be connected and let a question $(\mvec;\mathbf s,\mathbf t)$ of pair coordinates $(\bmu_\perp;\boldsymbol\sigma,\boldsymbol\delta)$ have every proper wall positive. In the rest of this statement $\oE$ is the expectation of $U\sim\mathrm{Breta}^\circ(F;\mvec;\mathbf s,\mathbf t)$, $Z$ is its tilt and $\mathbf u(\alpha)=\sum_j\big(s_jR_j(\alpha)+t_jR^\ast_j(\alpha)\big)$ the usage.
\begin{enumerate}[label=(\roman*),resume]
\item \emph{Dilation relation.}
\[
2\ip{\bmu_\perp}{\oE[Z]}-\sum_j\big(\sigma_j\oE[E_j(Z)]+\delta_j\oE[O_j(Z)]\big)=-(n-1)-\sum_j\big(s_j\oE[\operatorname{Ent}_j(Z)]+t_j\oE[\operatorname{Ent}^\ast_j(Z)]\big);
\]
at a centre, $\sum_j\sigma_j\oE[E_j(Z)]=n-1+\sum_j\sigma_j\oE[\operatorname{Ent}_j(Z)+\operatorname{Ent}^\ast_j(Z)]$.
\item $\operatorname{Cov}\big(\mathbf u_e(U),Z_f\big)=\frac12\big(\one[e=f]-\frac1n\big)$, $Z$ read on $\one^\perp$.
\item $\log P_F$ is convex in the pair coordinates of the question, so for a balanced direction $v=(\dot\mvec;\dot{\mathbf s},\dot{\mathbf t})$ with the questions at the data $(\mvec\pm\dot\mvec;\mathbf s\pm\dot{\mathbf s},\mathbf t\pm\dot{\mathbf t})$ inside their walls,
\begin{gather*}
\log\frac{P_F(\mvec;\mathbf s,\mathbf t)}{P_F(\mvec-\dot\mvec;\mathbf s-\dot{\mathbf s},\mathbf t-\dot{\mathbf t})}\le D_v\log P_F(\mvec;\mathbf s,\mathbf t)\le\log\frac{P_F(\mvec+\dot\mvec;\mathbf s+\dot{\mathbf s},\mathbf t+\dot{\mathbf t})}{P_F(\mvec;\mathbf s,\mathbf t)},
\end{gather*}
bracketing each mean of Proposition~\ref{prop:sampletilts}(ii) by neighbouring values of the period.
\end{enumerate}
\end{proposition}

% Reference 20: prop:momentsteps; source line 11687
\begin{proposition}[Moments are steps; entropies are derivatives one column up]\label{prop:momentsteps}
Let a connected family carry a question $(\mvec;\mathbf s,\mathbf t)$ with every proper wall positive, and in this statement let $\oE$ be the expectation of $U\sim\mathrm{Breta}^\circ(F;\mvec;\mathbf s,\mathbf t)$.
\begin{enumerate}[label=(\roman*)]
\item For a real balanced shift $(\dot\mvec;\dot{\mathbf s},\dot{\mathbf t})$ with every proper wall of the shifted question positive,
\[
\oE\Big[U^{\dot\mvec}\prod_jK_j(U)^{-\dot s_j}\Psi_j(U)^{-\dot t_j}\Big]=\frac{P_F(\mvec+\dot\mvec;\mathbf s+\dot{\mathbf s},\mathbf t+\dot{\mathbf t})}{P_F(\mvec;\mathbf s,\mathbf t)} .
\]
\item For a basis $T$ of table $j$ the \emph{columns} $v_{j,T}=(\one_T;\mathbf e_j,\mathbf 0)$ and $v^\ast_{j,T}=(\one_{\cE\setminus T};\mathbf 0,\mathbf e_j)$ are balanced, and for $a\ge0$
\begin{gather*}
\oE\Big[\Big(\frac{\mathbb P_{j,U}(T)}{\theta_{j,T}}\Big)^a\Big]=\frac{P_F(\mvec+a\one_T;\,\mathbf s+a\mathbf e_j,\,\mathbf t)}{P_F(\mvec;\mathbf s,\mathbf t)},\\
\oE\Big[\Big(\frac{\mathbb P^\ast_{j,U}(T)}{\theta_{j,T}}\Big)^a\Big]=\frac{P_F(\mvec+a\one_{\cE\setminus T};\,\mathbf s,\,\mathbf t+a\mathbf e_j)}{P_F(\mvec;\mathbf s,\mathbf t)},
\end{gather*}
the walls only growing along a column.
\item With the entropies of Proposition~\ref{prop:dilation},
\begin{gather*}
\oE[\operatorname{Ent}_j(Z)]=-\sum_T\theta_{j,T}\,\frac{(D_{v_{j,T}}P_F)(\mvec+\one_T;\,\mathbf s+\mathbf e_j,\,\mathbf t)}{P_F(\mvec;\mathbf s,\mathbf t)},\\
\oE[\operatorname{Ent}^\ast_j(Z)]=-\sum_T\theta_{j,T}\,\frac{(D_{v^\ast_{j,T}}P_F)(\mvec+\one_{\cE\setminus T};\,\mathbf s,\,\mathbf t+\mathbf e_j)}{P_F(\mvec;\mathbf s,\mathbf t)},
\end{gather*}
$Z$ the tilt of $U$: an entropy is a derivative of the period one column up.
\end{enumerate}
\end{proposition}

% Reference 20: cor:derivsum; source line 11703
\begin{corollary}[A derivative adds one finite sum]\label{cor:derivsum}
The series of Corollary~\ref{cor:cubeclose} may be differentiated term by term along a balanced direction $v=(\dot\mvec;\dot{\mathbf s},\dot{\mathbf t})$ at a point inside the walls, with $\dot\kappa_S$ the change of the walls; each factor of a term has an elementary logarithmic derivative,
\begin{gather*}
D_v\log\frac{(s_j)_{n_j}}{n_j!}=\dot s_j\sum_{i<n_j}\frac1{s_j+i},\qquad D_v\log\frac{(t_j)_{n'_j}}{n'_j!}=\dot t_j\sum_{i<n'_j}\frac1{t_j+i},\\
D_v\log K_j(\one)^{-s_j-t_j}=-(\dot s_j+\dot t_j)\log K_j(\one),\\
D_v\log\Gst21{\kappa_{U_k}+\beta_k,\ \gamma_k+1}{\kappa_{U_k}+\beta_k+\gamma_k+1}=-\dot\kappa_{U_k}\sum_{i=0}^{\gamma_k}\frac1{\kappa_{U_k}+\beta_k+i},
\end{gather*}
so every quantity of Proposition~\ref{prop:momentsteps} is the same positive series, at the question or one column up, with each term multiplied by a finite sum of reciprocals and logarithms of total weights.
\end{corollary}

Sorting a draw of the shape law lands it in one sector of the period: the chance of an order is the fraction of the period carried by its sector, the gaps of the sorted draw are exponential at the rates of the walls of the tails up to bounded factors, and the best and the worst basis of the draw are the greedy bases of its order. The races read the statistics as entropies, and the dilation relation is the balance between the strength of a question and the entropies of its winners. Every moment of a letter, of a winner's chance or of the shape is a step of the data, and every entropy is a derivative one column up, so on the closing series of the period each is a positive series whose terms carry one extra finite sum.

\section{The interaction of a bunch}

% Reference 20: prop:interactionroom; source line 11364 (one table).
% Proved in the last part: prop:interactionroom.
\begin{proposition}[The interaction belongs to the room]\label{prop:interactionroom}
Let a family have a bunch $B$, $|B|\ge2$, with $\pat_j(B)=2^B$ in every active table.
\begin{enumerate}[label=(\roman*)]
\item $\iota_{j,B}$ and $\iota^\ast_{j,B}$, and so $\iota^{\mathbf E}_B$ and $\iota^{\mathbf O}_B$, are homogeneous of degree $0$ in $\alpha_R$.
\item At a gauge $\alpha$ put, for $P\subseteq B$, $w(P)=\prod_j\pi_{j,P}(\alpha)^{s_j}\pi^\ast_{j,B\setminus P}(\alpha)^{t_j}$, with the pattern probabilities $\pi_{j,p}(\alpha)=\alpha_B^pK_{j,(p)}(\alpha_R)/K_j(\alpha)$ and $\pi^\ast_{j,p}(\alpha)=\alpha_B^{B\setminus p}\Psi_{j,(p)}(\alpha_R)/\Psi_j(\alpha)$ of Definition~\ref{def:split}; for a sample as in Proposition~\ref{prop:samplequestion} it is the chance that every draw uses exactly $P$ on $B$, its basis draws meeting $B$ in $P$ and its cobasis draws missing exactly $P$. In the expansion $\log w(P)=\sum_{A\subseteq P}\lambda_A$ over the subsets of $B$, the top term is $\lambda_B=\iota^{\mathbf E}_B+\iota^{\mathbf O}_B$, and for $A\subsetneq B$ with $|A|\ge2$, $\lambda_A$ is the sum $\iota^{\mathbf E}_A+\iota^{\mathbf O}_A$ of the interaction of $A$ in the face $F^{\cE\setminus(B\setminus A)}$ of the wall at $B\setminus A$ (Theorem~\ref{thm:wallres}), whose tables are the restrictions $\cB_j|\big(\cE\setminus(B\setminus A)\big)$, the deletions of $B\setminus A$, at the exponents $(s_j,0)$ and the contractions $\cB_j/(B\setminus A)$ at $(0,t_j)$; none of them depends on $\alpha_B$. For one table at the exponents $(1,0)$, $w(P)=\pi_P(\alpha)$ and $\lambda_B=\iota_B$, the top term of the log-linear expansion of the pattern probabilities.
\end{enumerate}
\end{proposition}

% Reference 20: thm:interactiontotal; source line 11375 (one table).
% Proved in the last part: thm:interactiontotal.
\begin{theorem}[The interaction totals the joint cumulants over the bunch's scales]\label{thm:interactiontotal}
Let a family with real exponents have a bunch $B=\{e_1,\dots,e_k\}$, $k\ge2$, with $\pat_j(B)=2^B$ in every active table. Then for both parts $\mathbf X\in\{\mathbf E,\mathbf O\}$
\[
\iota^{\mathbf X}_B(\alpha_R)=\int_{\mathbb R^B}\partial_{y_{e_1}}\cdots\partial_{y_{e_k}}\Phi_{\mathbf X}(y_B,\log\alpha_R)\,dy_B ,
\]
converging absolutely; at a realization the integrand of the sum of the parts is $2^{-k}\kappa_\Upsilon(e_1,\dots,e_k)=\sum_j\big(s_j\kappa_{j,\alpha}(e_1,\dots,e_k)+(-1)^kt_j\kappa_{j,\alpha^{-1}}(e_1,\dots,e_k)\big)$, with $\kappa_\Upsilon$ the still at the gauge $\alpha=(e^{y_B},\alpha_R)$ (Proposition~\ref{prop:deviation}). For one table at the exponents $(1,0)$ the integrand is the joint cumulant of the indicators of the marks of $B$ under the basis law at $(e^{y_B},\alpha_R)$ (Definition~\ref{def:weightlaw}), and the total is $\iota_B(\alpha_R)$.
\end{theorem}

% Reference 20: prop:interactionparity; source line 11386 (one table).
% Proved in the last part: prop:interactionparity.
\begin{proposition}[Both sides, and parity]\label{prop:interactionparity}
For a table with $\pat(B)=2^B$, with $\Psi=\sum_p\alpha_B^{B\setminus p}\Psi_{(p)}(\alpha_R)$, as functions of the tilt $z_R$ of the room $\iota^\ast_B(z_R)=(-1)^{|B|}\iota_B(-z_R)$. So for a family
\[
\iota^{\mathbf E}_B(z_R)=\sum_j\sigma_j\big(\iota_{j,B}(z_R)+(-1)^{|B|}\iota_{j,B}(-z_R)\big),\qquad \iota^{\mathbf O}_B(z_R)=\sum_j\delta_j\big(\iota_{j,B}(z_R)-(-1)^{|B|}\iota_{j,B}(-z_R)\big):
\]
for $|B|$ even the even part of the interaction is even in $z_R$ and the odd part is odd, for $|B|$ odd the reverse, and at the table's gauge of the room the interaction is $\big(2\sum_j\sigma_j\iota_{j,B},0\big)$ for $|B|$ even and $\big(0,2\sum_j\delta_j\iota_{j,B}\big)$ for $|B|$ odd: the strengths read the even bunches and the leans the odd ones, as they read the even and the odd rungs of the still at the table's gauge (Proposition~\ref{prop:stillgauge}). The total of Theorem~\ref{thm:interactiontotal} for the logarithm $-\Phi$ of the family word is $-\sum_j\big(s_j\iota_{j,B}(z_R)+(-1)^{|B|}t_j\iota_{j,B}(-z_R)\big)$.
\end{proposition}

% Reference 20: prop:design; source line 11393
\begin{proposition}[A design]\label{prop:design}
Let $\mathbf X\in\mathbb R^{\cE\times C}$ have full column rank, with rows $x_e$, and for $J\subseteq C$ let $K_J(\alpha)=\det\big(\mathbf X_J^{\mathsf T}\operatorname{diag}(\alpha)\mathbf X_J\big)$, $\mathbf X_J$ the columns in $J$, with $K_\varnothing=1$.
\begin{enumerate}[label=(\roman*)]
\item The marks are the observations of the determinantal table with the slice $\mathbf X^{\mathsf T}$ (Definition~\ref{def:multivector}) and the weights $\theta_T=\det(\mathbf X_{T,C})^2$, $|T|=|C|$: $K=K_C$, and the occupation of a mark is its leverage, $R_e(\alpha)=\alpha_e\,x_e^{\mathsf T}\big(\mathbf X^{\mathsf T}\operatorname{diag}(\alpha)\mathbf X\big)^{-1}x_e$.
\item Add one \emph{prior mark} $\gamma^{\mathrm p}$ for each column $\gamma\in C$, with the rows of the identity and the scales $\beta$: $K(\alpha,\beta)=\det\big(\mathbf X^{\mathsf T}\operatorname{diag}(\alpha)\mathbf X+\operatorname{diag}(\beta)\big)=\sum_{J\subseteq C}\beta^{C\setminus J}K_J(\alpha)$. The prior marks form a bunch $C^{\mathrm p}$ with $\pat(C^{\mathrm p})=2^{C^{\mathrm p}}$, and the child of the pattern $q^{\mathrm p}$, $q\subseteq C$, is $K_{C\setminus q}$.
\item For a split $C=J_1\sqcup J_2$ the generalized variance inflation factor is $\frac{K_{J_1}K_{J_2}}{K_CK_\varnothing}=\exp\big(-\iota_{(J_1^{\mathrm p},J_2^{\mathrm p})}\big)$, and for two columns $\gamma,\gamma'$, one minus the squared partial correlation given the rest is $\exp(\iota_{\{\gamma^{\mathrm p},\gamma'^{\mathrm p}\}})$ in the table with the other prior marks deleted.
\item The dual multivector of the table of (ii) has the slice $[\,I_{\cE}\mid-\mathbf X\,]$, and the basis polynomial of its dual table, $\det\big(\operatorname{diag}(\alpha')+\mathbf X\operatorname{diag}(\beta')\mathbf X^{\mathsf T}\big)$, is the polynomial of (ii) for the matrix $\mathbf X^{\mathsf T}$: the columns of $\mathbf X$ are its observations, at the scales $\beta'$, and the marks of $\cE$ its prior marks, at the scales $\alpha'$. Its child with every prior mark deleted, $\det\big(\mathbf X\operatorname{diag}(\beta')\mathbf X^{\mathsf T}\big)$, is nonzero exactly when $\mathbf X$ has rank $|\cE|$.
\end{enumerate}
\end{proposition}

A bunch all of whose subsets are patterns sees the room through its children, one for each subset of the bunch, and the alternating sum of their logarithms is the interaction: it is homogeneous of degree zero in the room, it is the top term of the log-linear expansion of the chance that a sample uses one pattern of the bunch throughout, and it is the total, over the bunch's scales, of the joint cumulants of the use of the bunch's marks under the basis and cobasis laws, weighted by the exponents. The dual table reads the same interaction at the opposite tilt with the sign of the parity of the bunch, so a family reads the interaction of an even bunch through its strengths and that of an odd bunch through its leans, as it reads the rungs of the same parities. In a regression design the prior marks of the columns are such a bunch: the interactions are minus the logarithms of the generalized variance inflation factors and the logarithms of one minus the squared partial correlations, and the occupations are the leverages.


\part{Proofs of the general forms}\label{part:proofs}

\chapter{Proofs of the general forms}\label{ch:proofs}

This part proves the statements that are more general than their counterparts in the reference \emph{The Foreman and the serif family} (16 October 2026), and those with no counterpart there; in the source file the comment before each such statement says so. Statements are cited by their numbers, and every symbol and term is that of the chapter the statement belongs to. A statement of the reference that is a special case of a statement proved here is not proved again; where reading it off needs a word, the word follows the proof. Lemmas that belong to this part alone are numbered S1, S2, \dots. The chapters are taken in order, one section each, and the last section records the checks run on small families.

\section{Chapter~\hb{ch:conv}: conventions, the stacked Gamma and the finite part}\label{sec:ch01}

\begin{slemma}\label{lem:conv-average}
Let $\Theta\subset\mathbb R$ be finite and $\phi(y)=\sum_{\theta\in\Theta}a_\theta e^{i\theta y}$ with $a_\theta\in\mathbb C$. If $\phi(y)\to0$ as $y\to-\infty$, then every $a_\theta$ is $0$.
\end{slemma}

\begin{proof}
Fix $\theta_0\in\Theta$. For $Y>0$,
\[
\frac1Y\int_{-2Y}^{-Y}\phi(y)\,e^{-i\theta_0y}\,dy=a_{\theta_0}+\sum_{\theta\neq\theta_0}a_\theta\,\frac1Y\int_{-2Y}^{-Y}e^{i(\theta-\theta_0)y}\,dy ,
\]
and each integral on the right has modulus at most $2/|\theta-\theta_0|$, so the sum is $O(Y^{-1})$. The left side has modulus at most $\sup_{y\le-Y}|\phi(y)|$, which tends to $0$. Letting $Y\to\infty$ gives $a_{\theta_0}=0$.
\end{proof}

\begin{slemma}\label{lem:conv-leading}
Let $g(x)=\sum_{(k,j)\in F}c_{k,j}\,x^k\log^jx$ for $x\in(0,1)$, where $F\subset\mathbb C\times\mathbb N$ is finite and nonempty and every $c_{k,j}$ is nonzero. Let $\kappa$ be the least real part of the exponents $k$, $(k,j)\in F$, and $J$ the largest $j$ such that $(k,j)\in F$ for some $k$ with $\re k=\kappa$. Then $g(x)/(x^\kappa\log^Jx)$ does not tend to $0$ as $x\to0$.
\end{slemma}

\begin{proof}
Put $y=\log x$, so that $y\to-\infty$ as $x\to0$ and $x^k=e^{ky}$. Then
\[
\frac{g(x)}{x^\kappa\log^Jx}=\sum_{\theta}c_{\kappa+i\theta,J}\,e^{i\theta y}+\sum_{(k,j)}c_{k,j}\,e^{(k-\kappa)y}\,y^{j-J},
\]
the first sum over the real $\theta$ with $(\kappa+i\theta,J)\in F$ and the second over the other elements of $F$. Every term of the second sum tends to $0$: if $\re k=\kappa$, then $j<J$, $|e^{(k-\kappa)y}|=1$ and $|y|^{j-J}\to0$; if $\re k>\kappa$, then $|e^{(k-\kappa)y}y^{j-J}|=e^{(\re k-\kappa)y}|y|^{j-J}\to0$. The first sum is a function $\phi(y)$ as in Lemma~\ref{lem:conv-average} with a nonzero coefficient, so it does not tend to $0$ as $y\to-\infty$, and neither does the left side.
\end{proof}

\begin{hdefinition}[\hb{def:Reg} (The finite part)]
As stated: the exponents of a log-power expansion are complex numbers, and its coefficients are determined by the function.
\end{hdefinition}

\begin{proof}
What is to be shown is that the coefficients of a log-power expansion are determined by $h$, so that the finite parts are well defined; the rest of the definition is that of the reference, with complex exponents.

\emph{(1) One remainder.} Let two finite sums $S_1$ and $S_2$ of terms $c\,x^k\log^jx$, every exponent of real part below $N$, both differ from $h$ by $O(x^N)$ as $x\to0$. Their difference $g=S_1-S_2$ is a finite sum of the same kind with $g(x)=O(x^N)$. If $g$ had a nonzero coefficient, then, with $\kappa<N$ and $J$ as in Lemma~\ref{lem:conv-leading}, $|g(x)/(x^\kappa\log^Jx)|\le Cx^{N-\kappa}|\log x|^{-J}\to0$, which that lemma excludes. So $S_1$ and $S_2$ have the same coefficients.

\emph{(2) Two remainders.} Let $N<N'$, and let $S$ and $S'$ be finite sums for $N$ and for $N'$. A term $x^k\log^jx$ of $S'$ with $\re k\ge N$ is $O(x^{N-\delta})$ for every $\delta>0$, so the part $S'_{<N}$ of $S'$ with exponents of real part below $N$ has $h-S'_{<N}=O(x^{N-\delta})$. Choose $\delta>0$ smaller than $N-\re k$ for every exponent $k$ of $S$ and of $S'_{<N}$. Then $S$ and $S'_{<N}$ both differ from $h$ by $O(x^{N-\delta})$, with all exponents of real part below $N-\delta$, so they coincide by (1). Hence the coefficient of a term $x^k\log^jx$ is the same for every $N>\re k$, and the coefficient of $x^0\log^0x$, read with any $N>0$, is determined by $h$. At $\infty$ the same holds for $h(1/x)$ at $0$: the term $x^{-k}\log^jx$ is $(-1)^jx'^k\log^jx'$ in $x'=1/x$, and $x^0\log^0x$ is $x'^0\log^0x'$.

\emph{(3) The primitive.} Two primitives of $f$ differ by a constant $C$, whose expansion at either end is the single term $C\,x^0\log^0x$; by (2) the finite parts of the second primitive at $0$ and at $\infty$ are those of the first plus $C$. So $\Reg_{x\to\infty}h-\Reg_{x\to0}h$ and $h(b)-\Reg_{x\to0}h$ do not depend on the primitive.
\end{proof}

\begin{hproposition}[\hb{prop:Regconv} (A convergent integral is its own finite part)]
As stated, for log-power expansions with complex exponents.
\end{hproposition}

\begin{proof}
Let $h$ be a primitive of $f$ with log-power expansions at both ends. Since $\int_0^\infty f\,dx$ converges, $h(x)=h(1)-\int_x^1f$ has a limit $L_0$ as $x\to0$, $h$ has a limit $L_\infty$ as $x\to\infty$, and $\int_0^\infty f\,dx=L_\infty-L_0$. Take the expansion at $0$ for some $N>0$, and let $G$ be its part with exponents of real part at most $0$. The other terms and the remainder tend to $0$, so $G(x)\to L_0$. Write $G-L_0$ as a finite sum of terms $c\,x^k\log^jx$, the constant $-L_0$ joining the coefficient of $x^0\log^0x$, and suppose it has a nonzero coefficient; let $\kappa\le0$ and $J$ be as in Lemma~\ref{lem:conv-leading}. If $(\kappa,J)\neq(0,0)$, then $|x^\kappa\log^Jx|\to\infty$, so $(G-L_0)/(x^\kappa\log^Jx)\to0$. If $(\kappa,J)=(0,0)$, every exponent of $G-L_0$ has real part $0$ and every power of the logarithm is $0$, and $(G-L_0)/(x^0\log^0x)=G-L_0\to0$. Both are excluded by Lemma~\ref{lem:conv-leading}. So $G$ is the constant $L_0$: the expansion at $0$ has no term with an exponent of real part at most $0$ other than $L_0\,x^0\log^0x$, and $\Reg_{x\to0}h=L_0$. In the same way, for $h(1/x)$ at $0$, $\Reg_{x\to\infty}h=L_\infty$. Hence $\Reg\!\int_0^\infty f\,dx=L_\infty-L_0=\int_0^\infty f\,dx$.
\end{proof}

\begin{htheorem}[\hb{thm:RegCT} (A finite part is a constant term)]
As stated, for log-power expansions with complex exponents.
\end{htheorem}

\begin{proof}
The steps are those of the reference, proof of \hb{thm:RegCT}; they hold at complex exponents as follows.

\emph{The end at $0$.} On $(0,1]$ write $f=\sum c_{b,j}x^b\log^jx+r_N$, the sum over the finitely many terms of the expansion at $0$ with $\re b<N$, where $N>0$ and $r_N=O(x^N)$ as $x\to0$; $r_N$ is integrable on $(0,1]$. For $b\neq-1$ put
\[
Q_{b,j}(y)=\sum_{i=0}^j(-1)^{j-i}\frac{j!}{i!\,(b+1)^{j-i+1}}\,y^i ;
\]
then $(b+1)Q_{b,j}+Q_{b,j}'=y^j$, so $\frac d{dx}\big(x^{b+1}Q_{b,j}(\log x)\big)=x^b\log^jx$, and $Q_{b,j}(0)=(-1)^jj!/(b+1)^{j+1}$. Hence
\[
\int_x^1t^b\log^jt\,dt=Q_{b,j}(0)-x^{b+1}Q_{b,j}(\log x)\quad(b\neq-1),\qquad\int_x^1t^{-1}\log^jt\,dt=-\frac{\log^{j+1}x}{j+1},
\]
and $\int_x^1r_N=\int_0^1r_N-\int_0^xr_N$ with $\int_0^xr_N=O(x^{N+1})$. So $\int_x^1f$ has a log-power expansion at $0$, and its coefficient of $x^0\log^0x$ is $\sum_{b\neq-1}c_{b,j}Q_{b,j}(0)+\int_0^1r_N$: a term $x^{b+1}\log^ix$ has the exponent $b+1\neq0$, and a term $\log^{j+1}x$ has a positive power of the logarithm. For a primitive $h$ of $f$, $h(1)-\Reg_{x\to0}h=\Reg_{x\to0}\int_x^1f$, so by \hb{def:Reg}
\[
\Reg\!\int_0^1f\,dx=\sum_{b\neq-1}c_{b,j}\frac{(-1)^jj!}{(b+1)^{j+1}}+\int_0^1r_N\,dx .
\]
On the other side, by $\int_0^1x^{s-1}\log^jx\,dx=(-1)^jj!/s^{j+1}$ for $\re s>0$,
\[
\int_0^1x^\varepsilon f\,dx=\sum c_{b,j}\frac{(-1)^jj!}{(b+1+\varepsilon)^{j+1}}+\int_0^1x^\varepsilon r_N\,dx
\]
for $\re\varepsilon$ large, the last term holomorphic on $\re\varepsilon>-N-1$; this is the meromorphic continuation there. Its constant term at $\varepsilon=0$ is the sum displayed above: a term with $b\neq-1$ gives its value, and a term with $b=-1$ is the pure pole $(-1)^jj!\,\varepsilon^{-j-1}$, of constant term $0$.

\emph{The end at $\infty$.} By $x\mapsto1/x$, $\int_1^\infty x^\varepsilon f\,dx=\int_0^1x^{-\varepsilon}\,x^{-2}f(1/x)\,dx$ and $\int_1^xf=\int_{1/x}^1t^{-2}f(1/t)\,dt$; the function $x^{-2}f(1/x)$ has a log-power expansion at $0$, the coefficient of $x^0\log^0x$ at $\infty$ is that of $x'^0\log^0x'$ at $0$ in $x'=1/x$, and a constant term at $\varepsilon=0$ does not change under $\varepsilon\mapsto-\varepsilon$. So $\Reg_{x\to\infty}\int_1^xf=\CT_{\varepsilon=0}\int_1^\infty x^\varepsilon f\,dx$ by the end at $0$.

Adding, $\Reg\!\int_0^\infty f\,dx=\Reg_{x\to\infty}h-\Reg_{x\to0}h=\Reg_{x\to\infty}\int_1^xf+\Reg_{x\to0}\int_x^1f$ is the constant term of the sum of the two continuations.
\end{proof}

\begin{hcorollary}[\hb{cor:RegMellin} (Finite parts at a pole of a Mellin integral)]
As stated: for every $f$ with a log-power expansion at $0$, at every $a\in\mathbb C$.
\end{hcorollary}

\begin{proof}
Fix $a$ and $N>1-\re a$, and write $f=\sum c_{k,j}x^k\log^jx+r_N$ on $(0,1]$, the finitely many terms of the expansion with $\re k<N$, $r_N=O(x^N)$ at $0$ and integrable on $[\delta,1]$ for every $\delta>0$. For $\re a'$ large,
\[
I(a')=\sum c_{k,j}\frac{(-1)^jj!}{(a'+k)^{j+1}}+\int_0^1x^{a'-1}r_N\,dx ,
\]
the last term holomorphic on $\re a'>-N$. This is the meromorphic continuation of $I$ on $\re a'>-N$; its poles lie at the points $-k$, a term $x^k\log^jx$ giving a pole of order at most $j+1$, and since $N$ is arbitrary this describes the continuation on $\mathbb C$. Its constant term at $a'=a$ is
\[
\CT_{a'=a}I(a')=\sum_{a+k\neq0}c_{k,j}\frac{(-1)^jj!}{(a+k)^{j+1}}+\int_0^1x^{a-1}r_N\,dx ,
\]
a term with $a+k=0$ being the pure pole $(-1)^jj!\,(a'-a)^{-j-1}$; at an $a$ that is not a pole it is $I(a)$. On the other side, $x^{a-1}f$ has the log-power expansion at $0$ with the terms $c_{k,j}x^{a+k-1}\log^jx$ and the remainder $x^{a-1}r_N=O(x^{N+\re a-1})$, $N+\re a-1>0$. The end at $0$ in the proof of \hb{thm:RegCT} above, with $b=a+k-1$ (so $b\neq-1$ exactly when $a+k\neq0$, and $b+1=a+k$), gives for $\Reg\!\int_0^1x^{a-1}f\,dx$ the same sum.
\end{proof}

\begin{hlemma}[\hb{lem:Regscale} (A finite part depends on the coordinate)]
As stated: the dilation rule for every function with log-power expansions.
\end{hlemma}

\begin{proof}
\emph{(a) The finite part of $h(cx^p)$.} Let $h(y)=\sum c_{k,j}\,y^{-k}\log^jy+O(y^{-N})$ as $y\to\infty$, finitely many terms with $\re k<N$, so that $a_j=c_{0,j}$. For $y=cx^p$, $y^{-k}=c^{-k}x^{-pk}$ and $\log y=\log c+p\log x$, so
\[
h(cx^p)=\sum c_{k,j}\,c^{-k}x^{-pk}\sum_{i=0}^j\binom ji(\log c)^{j-i}p^i\log^ix+O(x^{-pN}),
\]
a log-power expansion at $\infty$ with the exponents $pk$, of real parts below $pN$; as $N$ is arbitrary, so is $pN$. The term $x^0\log^0x$ comes from $pk=0$, that is $k=0$, and $i=0$; by \hb{def:Reg} its coefficient $\sum_jc_{0,j}(\log c)^j=\sum_ja_j(\log c)^j$ is $\Reg_{x\to\infty}h(cx^p)$. At $0$ the computation is the same.

\emph{(b) The finite part of the integral.} Let $F(x)=\int_1^xf$. At $\infty$, $f(t)=\sum f_{k,j}t^{-k}\log^jt+O(t^{-N})$ with $N>1$; integrating from $1$ to $x$, a term with $k\neq1$ gives $x^{1-k}Q_{-k,j}(\log x)$ up to a constant (with $Q_{b,j}$ as in the proof of \hb{thm:RegCT} above), a term with $k=1$ gives $\log^{j+1}x/(j+1)$, and the remainder gives a constant plus $O(x^{1-N})$. So $F$ has a log-power expansion at $\infty$, and for $i\ge1$ its coefficient of $x^0\log^ix$ comes from the term $k=1$, $j=i-1$ alone: it is $f^\infty_{i-1}/i$. In the same way, at $0$ the coefficient of $x^0\log^ix$ in $F$ is $f^0_{i-1}/i$ for $i\ge1$. Write $A^\infty_i$ and $A^0_i$ for the coefficients of $x^0\log^ix$ of $F$ at $\infty$ and at $0$; $A^\infty_0-A^0_0=\Reg\!\int_0^\infty f\,dx$. The function $G(x)=F(cx^p)$ is a primitive of $f(cx^p)\,c\,p\,x^{p-1}$, and by (a) at both ends
\[
\Reg\!\int_0^\infty f(cx^p)\,c\,p\,x^{p-1}dx=\sum_{i\ge0}\big(A^\infty_i-A^0_i\big)(\log c)^i=\Reg\!\int_0^\infty f\,dx+\sum_{i\ge1}\frac{f^\infty_{i-1}-f^0_{i-1}}i(\log c)^i ,
\]
which is the display with $i=j+1$. At $c=1$ the correction vanishes: a power of the coordinate leaves the finite part unchanged.
\end{proof}

\begin{sremark}
The reference states the instance $f(x)=1/(1+cx)$, $c>0$. This $f$ has the coefficient $f^\infty_0=1/c$ of $x^{-1}$ at $\infty$ and no term $x^{-1}\log^jx$ at $0$, and with the dilation $c^{-1}$ and $p=1$ the lemma gives $\frac1c\Reg\!\int_0^\infty\frac{dx'}{1+x'}=\Reg\!\int_0^\infty\frac{dx}{1+cx}-\frac{\log c}c$. The left side is $0$, the primitive $\log(1+x')=\log x'+\log(1+1/x')$ having no constant term at either end; so $\Reg\!\int_0^\infty\frac{dx}{1+cx}=\frac{\log c}c$, the value of the reference. A second check, with a logarithm: $f(x)=\log(1+x)/(1+x)$ has the primitive $\frac12\log^2(1+x)$, with no constant term at either end, so $\Reg\!\int_0^\infty f\,dx=0$; at $\infty$, $f=x^{-1}\log x+O(x^{-2}\log x)$, so $f^\infty_1=1$ and $f^\infty_0=0$, and at $0$ it has no term $x^{-1}\log^jx$. The lemma gives $\Reg\!\int_0^\infty f(cx)\,c\,dx=\frac12\log^2c$; directly, the primitive $\frac12\log^2(1+cx)=\frac12\big(\log c+\log x+\log(1+\tfrac1{cx})\big)^2$ has the constant term $\frac12\log^2c$ at $\infty$ and none at $0$.
\end{sremark}

\section{Chapter~\hb{ch:gramma}: the Gramma and the Zamma}\label{sec:ch02}

\begin{hproposition}[\hb{prop:grammareflect} (Affine changes of the logarithmic coordinates)]
As stated: one statement for the reflection, the linear changes of the coordinates and the torus homogeneity.
\end{hproposition}

\begin{proof}
Let $A\in GL_{\cE}(\mathbb R)$ and $\mathbf h\in\mathbb R^{\cE}$. The map $y'\mapsto y=Ay'+\mathbf h$ is a bijection of $\mathbb R^{\cE}$ with $dy=|\det A|\,dy'$, and $\ip\nu y=\ip{A^{\mathsf T}\nu}{y'}+\ip\nu{\mathbf h}$, $\ip{t_k}y=\ip{A^{\mathsf T}t_k}{y'}+\ip{t_k}{\mathbf h}$. In the logarithmic form of \hb{def:gramma},
\[
e^{\ip\nu y}\exp\Big(-\sum_k\xi_ke^{\ip{t_k}y}\Big)\,dy=|\det A|\,e^{\ip\nu{\mathbf h}}\,e^{\ip{A^{\mathsf T}\nu}{y'}}\exp\Big(-\sum_k\big(e^{\ip{t_k}{\mathbf h}}\xi_k\big)\,e^{\ip{A^{\mathsf T}t_k}{y'}}\Big)\,dy' .
\]
The moduli of the two sides are the same expressions with $\nu$ replaced by $\re\nu$, so the integral of the left side converges exactly when that of the right side does. The left side is the integrand of $\sG_{\nu;t}(\boldsymbol\xi)$ and the right side, up to the constant factor, that of $\sG_{A^{\mathsf T}\nu;A^{\mathsf T}t}\big((e^{\ip{t_k}{\mathbf h}}\xi_k)_k\big)$, with positive coefficients; by \hb{thm:grammadomain}, $\nu\in\cT_t$ exactly when $A^{\mathsf T}\nu\in\cT_{A^{\mathsf T}t}$, that is $\cT_{A^{\mathsf T}t}=A^{\mathsf T}\cT_t$, and integrating gives the display. At $A=-I_{\cE}$ and $\mathbf h=0$, $|\det A|=1$ and $A^{\mathsf T}t=-t$: the reflection. At $A=I_{\cE}$ and $\mathbf h=\log\lambda$, $\lambda\in\mathbb R^{\cE}_{>0}$, $e^{\ip{t_k}{\mathbf h}}=\lambda^{t_k}$ and $e^{\ip\nu{\mathbf h}}=\lambda^\nu$: the torus homogeneity.
\end{proof}

\begin{htheorem}[\hb{thm:grammaproduct} (Product)]
As stated, for any number of blocks.
\end{htheorem}

\begin{proof}
Write $x=(x^{(1)},\dots,x^{(q)})$ with $x^{(i)}\in\mathbb R^{\cE_i}_{>0}$. Each $t_k$ is supported in one $\cE_i$, so $x^{t_k}=(x^{(i)})^{t_k|_{\cE_i}}$ for that $i$, and
\[
x^{\nu-\one}\exp\Big(-\sum_k\xi_kx^{t_k}\Big)=\prod_{i=1}^q(x^{(i)})^{\nu|_{\cE_i}-\one}\exp\Big(-\sum_{k\in K_i}\xi_k\,(x^{(i)})^{t_k|_{\cE_i}}\Big),
\]
$K_i$ the set of the terms whose exponents are supported in $\cE_i$, the zero vectors counted in $K_1$. The moduli of the factors are positive functions of separate variables, so by Tonelli the modulus of the product is integrable exactly when every factor is integrable over $\mathbb R^{\cE_i}_{>0}$; by \hb{thm:grammadomain} this says $\cT_t=\cT_{t^{(1)}}\times\cdots\times\cT_{t^{(q)}}$, and Fubini gives the product.
\end{proof}

\begin{htheorem}[\hb{thm:scaleshapeset} (Scale and shape along sets)]
As stated, for several disjoint sets at once.
\end{htheorem}

\begin{proof}
Apply \hb{lem:scaleshape} on each $B_i$: $x_{B_i}=\rho_iu_i$ with $\rho_i>0$ and $u_i\in\Delta_{B_i}$, and $x_{B_i}^{\nu|_{B_i}-\one}dx_{B_i}=\rho_i^{\nu_{B_i}-1}u_i^{\nu|_{B_i}-\one}\,d\rho_i\,d\omega_{B_i}(u_i)$. With $x_R$ the coordinates on $R$, the map $(x_R,(\rho_i)_i,(u_i)_i)\mapsto x$ is a bijection onto $\mathbb R^{\cE}_{>0}$, and since $(\rho_iu_i)^{t_k|_{B_i}}=\rho_i^{\ip{t_k}{\one_{B_i}}}u_i^{t_k|_{B_i}}$,
\[
x^{\nu-\one}dx=x_R^{\nu|_R-\one}\prod_i\rho_i^{\nu_{B_i}-1}u_i^{\nu|_{B_i}-\one}\,d\rho\,d\omega(u)\,dx_R,\qquad \xi_kx^{t_k}=b_k(u,x_R)\prod_i\rho_i^{\ell_{k,i}} .
\]
At $\nu\in\cT_t$ the Gramma converges (\hb{thm:grammadomain}), that is, its integrand is absolutely integrable. By Tonelli and Fubini the integral may be taken with the integral over $\rho=(\rho_i)_i\in\mathbb R^q_{>0}$ innermost, that integral converging absolutely for almost every $(u,x_R)$, and it is
\[
\int_{\mathbb R^q_{>0}}\rho^{(\nu_{B_i})_i-\one}\exp\Big(-\sum_kb_k(u,x_R)\,\rho^{\ell_k}\Big)\,d\rho=\sG_{(\nu_{B_i})_i;\ell}\big(\mathbf b(u,x_R)\big),
\]
the Gramma of \hb{def:gramma} on $q$ scales with the configuration $\ell$. When $R$ is empty there are no coordinates $x_R$.
\end{proof}

\begin{hproposition}[\hb{prop:zammafacts} (First facts about the Zamma)]
As stated: (i) at every invertible square configuration and (iv) on any number of scales; (ii) and (iii) are those of the reference.
\end{hproposition}

\begin{proof}
(i) By \hb{cor:square}, $\sG_{\nu;t}(\boldsymbol\xi)=|\det t|^{-1}\prod_k\Gamma(b_k)\,\xi_k^{-b_k}$ on $\cT_t=\{\nu:\re\mathbf b\in(0,\infty)^N\}$; dividing by $\prod_e\Gamma(\nu_e)$ (\hb{def:zamma}) gives the display, the quotient of Gammas being the stacked Gamma of \hb{def:stacked}, and the display carries the continuation. At the unit configuration $\mathbf b=\nu$ and the display is the monomial; on one scale with $t_1>0$, $|\det t|=t_1$ and $b_1=\nu/t_1$, which is the formula of the reference.

(iv) Write $t''=I_{\cE}\oplus t'$ and $\boldsymbol\xi''=\one\oplus\boldsymbol\xi'$. Since $\cone(t'')\supseteq\cone(I_{\cE})=[0,\infty)^{\cE}$, every $\nu\in\mathbb R^{\cE}_{>0}$ lies in $\interior\cone(t'')$, hence in $\cT_{t''}$. By \hb{def:gramma},
\[
\sG_{\nu;t''}(\boldsymbol\xi'')=\int_{\mathbb R^{\cE}_{>0}}x^{\nu-\one}e^{-\sum_ex_e}\exp\Big(-\sum_k\xi'_kx^{t'_k}\Big)\,dx=\prod_e\Gamma(\nu_e)\,\oE_{U\sim\operatorname{Gamma}_\nu}\Big[\exp\Big(-\sum_k\xi'_kU^{t'_k}\Big)\Big],
\]
$\operatorname{Gamma}_\nu$ having the density $\prod_ex_e^{\nu_e-1}e^{-x_e}/\Gamma(\nu_e)$ (\hb{def:conv}(vii)); dividing by $\prod_e\Gamma(\nu_e)$ gives the display. The variable $\exp(-\sum_k\xi'_kU^{t'_k})$ lies in $(0,1]$ and decreases in each $\xi'_k$, and so does its expectation. Put $X=\sum_k\xi'_kU^{t'_k}\ge0$. Where every $\nu+t'_k$ lies in $\mathbb R^{\cE}_{>0}$, the coordinates of $U$ being independent, $\oE[U^{t'_k}]=\prod_e\oE[U_e^{t'_{k,e}}]=\prod_e\Gamma(\nu_e+t'_{k,e})/\Gamma(\nu_e)$, so $\oE X=\sum_k\xi'_k\prod_e\Gst11{\nu_e+t'_{k,e}}{\nu_e}$ is finite. The function $s\mapsto e^{-s}$ is convex, so $e^{-s}\ge e^{-s_0}-e^{-s_0}(s-s_0)$ for all real $s,s_0$; at $s=X$ and $s_0=\oE X$, taking expectations, $\oE[e^{-X}]\ge e^{-\oE X}$, which is the bound.
\end{proof}

\begin{hproposition}[\hb{prop:levelpoly} (Powers of a polynomial in the scales)]
As stated: (i) and (iii) for every Zamma; (ii) and (iv) are those of the reference.
\end{hproposition}

\begin{proof}
Write $\boldsymbol\xi(\rho)=\boldsymbol\xi^{(0)}+\boldsymbol\Delta(\rho)$ with $\Delta_k(\rho)=\sum_{l\neq0}\rho^l\xi^{(l)}_k$, a polynomial vanishing at $\rho=0$.

\emph{(1) Derivatives.} Away from the unit configuration, $\sZ_{\nu;t}=\sG_{\nu;t}/\prod_e\Gamma(\nu_e)$ on $(0,\infty)^N$. By \hb{thm:grammalaws}(i), $\partial^{\mathbf w}\sG_{\nu;t}=(-1)^{|\mathbf w|}\sG_{\nu+t\mathbf w;t}$, every $\nu+t\mathbf w$ lying in $\cT_t$ because $\interior\cone(t)+\cone(t)\subseteq\interior\cone(t)$; so $\sZ_{\nu;t}$ is infinitely differentiable on $(0,\infty)^N$ and
\[
\partial^{\mathbf w}\sZ_{\nu;t}=(-1)^{|\mathbf w|}\frac{\sG_{\nu+t\mathbf w;t}}{\prod_e\Gamma(\nu_e)}=(-1)^{|\mathbf w|}\prod_e\Gst11{\nu_e+(t\mathbf w)_e}{\nu_e}\,\sZ_{\nu+t\mathbf w;t},
\]
the second form wherever no entry of $\nu+t\mathbf w$ lies in $-\mathbb N$, the product being read as the first form otherwise, as in the statement. At the unit configuration $\sZ_{\nu;I_{\cE}}(\boldsymbol\xi)=\boldsymbol\xi^{-\nu}$ (\hb{def:zamma}), holomorphic near $\boldsymbol\xi^{(0)}$ with the principal branch, and $\partial^{\mathbf w}\boldsymbol\xi^{-\nu}=(-1)^{|\mathbf w|}\prod_e(\nu_e)_{w_e}\,\boldsymbol\xi^{-\nu-\mathbf w}$ with $(\nu_e)_{w_e}=\nu_e(\nu_e+1)\cdots(\nu_e+w_e-1)$.

\emph{(2) Taylor coefficients.} Each monomial of $\boldsymbol\Delta(\rho)^{\mathbf w}$ has total degree at least $|\mathbf w|$ in $\rho$, every $l\neq0$ having $|l|\ge1$. So Taylor's formula for $\sZ_{\nu;t}$ at $\boldsymbol\xi^{(0)}$ to the order $|q|$, whose remainder is $O(|\boldsymbol\Delta|^{|q|+1})=O(|\rho|^{|q|+1})$, gives
\[
c^\nu_q=\sum_{|\mathbf w|\le|q|}\partial^{\mathbf w}\sZ_{\nu;t}(\boldsymbol\xi^{(0)})\,[\rho^q]\,\frac{\boldsymbol\Delta(\rho)^{\mathbf w}}{\mathbf w!} .
\]
By the multinomial theorem, $\Delta_k^{w_k}/w_k!=\sum\prod_{l\neq0}(\rho^l\xi^{(l)}_k)^{n_{l,k}}/n_{l,k}!$, the sum over the $(n_{l,k})_l$ in $\mathbb N$ with $\sum_ln_{l,k}=w_k$; multiplying over $k$ and taking the coefficient of $\rho^q$ leaves the arrays $\mathbf n$ with $\sum_{l,k}l\,n_{l,k}=q$ and $\mathbf w=(\sum_ln_{l,k})_k$, $|\mathbf w|=|\mathbf n|$. With (1) this is (i). At one unit exponent, $t=(1)$, $\boldsymbol\xi=L$ and the index $\lambda$, the product in (i) is $\Gst11{\lambda+|\mathbf n|}{\lambda}(L^{(0)})^{-\lambda-|\mathbf n|}\prod_l(L^{(l)})^{n_l}/n_l!$, which is the formula of the reference.

\emph{(3) The coefficients as variables.} The function $(\rho,(\xi^{(l)}_k))\mapsto\sZ_{\nu;t}(\boldsymbol\xi(\rho))$ is infinitely differentiable, so derivatives in the coefficients commute with taking Taylor coefficients in $\rho$. Since $\partial\xi_k(\rho)/\partial\xi^{(l)}_k=\rho^l$, (1) with $\mathbf w=\mathbf e_k$ gives
\[
\partial_{\xi^{(l)}_k}\sZ_{\nu;t}(\boldsymbol\xi(\rho))=-\rho^l\prod_e\Gst11{\nu_e+t_{k,e}}{\nu_e}\,\sZ_{\nu+t_k;t}(\boldsymbol\xi(\rho)),
\]
and the coefficient of $\rho^q$ is $-\prod_e\Gst11{\nu_e+t_{k,e}}{\nu_e}\,c^{\nu+t_k}_{q-l}$, which is (iii). At $t=(1)$ the factor is $\Gst11{\lambda+1}{\lambda}=\lambda$, the formula of the reference.
\end{proof}

\section{Chapter~\hb{ch:breta}: the Breta}\label{sec:ch03}

\begin{hproposition}[\hb{prop:bretaorthant} (The Breta on the orthant)]
As stated, for any data.
\end{hproposition}

\begin{proof}
By \hb{def:breta} and \hb{def:zamma}, with the dressed coefficients $\eta_k=u^{\pi_k}\xi_k$ and the bunch scales $\rho=(\rho_Q)_{Q\in\cP}$,
\[
\sZ_{\boldsymbol\nu;\ell}(\boldsymbol\eta)=\frac1{\prod_Q\Gamma(\nu_Q)}\int_{\mathbb R^{\cP}_{>0}}\rho^{\boldsymbol\nu-\one}\exp\Big(-\sum_k\xi_k\,u^{\pi_k}\rho^{\ell_k}\Big)\,d\rho .
\]
On each bunch, \hb{lem:scaleshape} with $x_Q=\rho_Qu_Q$ gives $x_Q^{\mvec|_Q-\one}dx_Q=\rho_Q^{m_Q-1}u_Q^{\mvec|_Q-\one}d\rho_Q\,d\omega_Q(u_Q)$, hence
\[
u_Q^{\mvec|_Q-\one}\rho_Q^{\nu_Q-1}\,d\rho_Q\,d\omega_Q(u_Q)=x_Q^{\mvec|_Q-\one}|x_Q|^{\nu_Q-m_Q}\,dx_Q ,
\]
and each dressed term is
\[
\xi_k\,u^{\pi_k}\rho^{\ell_k}=\xi_k\prod_Q\Big(\frac{x_Q}{|x_Q|}\Big)^{\pi_k|_Q}|x_Q|^{\ell_{k,Q}}=\xi_k\,x^{\pi_k}\prod_Q|x_Q|^{\ell_{k,Q}-\ip{\pi_k}{\one_Q}} .
\]
So, under the bijection $x\leftrightarrow(\rho_Q,u_Q)_Q$ of \hb{lem:scaleshape}, the Breta with its Zamma written as the scale integral is the integral of the right side over $\mathbb R^{\cE}_{>0}$. The moduli of the two integrands are the same functions with $\mvec$ and $\boldsymbol\nu$ replaced by their real parts, so, as in the reference, proof of \hb{prop:bretaorthant}, Tonelli gives that either side converges when the other does, and Fubini gives the equality. For homogeneous data the inner powers of the bunch totals are absent, and at $\nu_Q=m_Q$ so are the outer ones: the right side is then $\sG_{\mvec;t}(\boldsymbol\xi)/\prod_Q\Gamma(m_Q)$, which is \hb{thm:tower}.
\end{proof}

\begin{htheorem}[\hb{thm:bretalaws} (Relations of the Breta)]
As stated: (v) for any terms; (i) to (iv) are those of the reference.
\end{htheorem}

\begin{proof}
The proofs of (i) to (iv) in the reference, proof of \hb{thm:bretalaws}, use no homogeneity. For (v), let $e\neq f$ be scales of the bunch $Q$, and in this proof write, with the products over the bunches $Q'$ of $\cP$,
\[
F(x)=x^{\mvec-\one}\prod_{Q'}|x_{Q'}|^{\nu_{Q'}-m_{Q'}}\exp\Big(-\sum_k\xi_k\varphi_k(x)\Big),\qquad \varphi_k(x)=x^{\pi_k}\prod_{Q'}|x_{Q'}|^{\ell_{k,Q'}-\ip{\pi_k}{\one_{Q'}}} ,
\]
the integrand of \hb{prop:bretaorthant} at $(\mvec,\boldsymbol\nu)$ without the factor $1/\prod_{Q'}\Gamma(\nu_{Q'})$.

\emph{(1) Derivatives.} For $e\in Q$, $x_e\partial_{x_e}|x_Q|=x_e$ and $x_e\partial_{x_e}|x_{Q'}|=0$ for $Q'\neq Q$, so
\[
x_e\partial_{x_e}F=F\Big[(m_e-1)+(\nu_Q-m_Q)\frac{x_e}{|x_Q|}-\sum_k\xi_k\varphi_k\Big(\pi_{k,e}+\big(\ell_{k,Q}-\ip{\pi_k}{\one_Q}\big)\frac{x_e}{|x_Q|}\Big)\Big],
\]
and the same holds with $f$ in place of $e$.

\emph{(2) The integrands.} By \hb{prop:bretaorthant}: $F$ is the integrand of $\prod_{Q'}\Gamma(\nu_{Q'})\,\sB^{\cP}_{\mvec,\boldsymbol\nu;(\pi,\ell)}(\boldsymbol\xi)$; $F\,x_e/|x_Q|$ is that of $\prod_{Q'}\Gamma(\nu_{Q'})\,\sB^{\cP}_{\mvec+\mathbf e_e,\boldsymbol\nu;(\pi,\ell)}(\boldsymbol\xi)$, the power of $|x_Q|$ being $\nu_Q-(m_Q+1)$; $F\varphi_k$ is that of $\prod_{Q'}\Gamma(\nu_{Q'}+\ell_{k,Q'})\,\sB^{\cP}_{\mvec+\pi_k,\boldsymbol\nu+\ell_k;(\pi,\ell)}(\boldsymbol\xi)$, the power of each $|x_{Q'}|$ being $(\nu_{Q'}+\ell_{k,Q'})-(m_{Q'}+\ip{\pi_k}{\one_{Q'}})$; and $F\varphi_k\,x_e/|x_Q|$ is that of $\prod_{Q'}\Gamma(\nu_{Q'}+\ell_{k,Q'})\,\sB^{\cP}_{\mvec+\pi_k+\mathbf e_e,\boldsymbol\nu+\ell_k;(\pi,\ell)}(\boldsymbol\xi)$. Here $\prod_{Q'}\Gamma(\nu_{Q'}+\ell_{k,Q'})=g_k\prod_{Q'}\Gamma(\nu_{Q'})$. These Bretas, and those with $f$ in place of $e$, are the Bretas of (v), which converge by hypothesis; so all these functions are integrable on $\mathbb R^{\cE}_{>0}$.

\emph{(3) No boundary terms.} By (1) and (2), $\partial_{x_e}(x_eF)=F+x_e\partial_{x_e}F$ is integrable. By Fubini, for almost every value $x'$ of the coordinates other than $x_e$ the function $G(s)=s\,F(x',s)$ of $s=x_e>0$ has $G'$ and $G(s)/s$ integrable on $(0,\infty)$. Then $G(s)=G(1)+\int_1^sG'$ has limits at $0$ and at $\infty$, and both are $0$: a nonzero limit $L$ at one end would make $|G(s)/s|\ge|L|/(2s)$ near that end, which is not integrable there. So $\int_0^\infty G'(s)\,ds=0$, and by Fubini $\int\partial_{x_e}(x_eF)\,dx=0$; likewise for $f$. Hence
\[
\int_{\mathbb R^{\cE}_{>0}}\big(x_e\partial_{x_e}-x_f\partial_{x_f}\big)F\,dx=\int\big[\partial_{x_e}(x_eF)-F\big]dx-\int\big[\partial_{x_f}(x_fF)-F\big]dx=0 .
\]

\emph{(4) Collecting.} By (1), $(x_e\partial_{x_e}-x_f\partial_{x_f})F$ is $F$ times
\[
(m_e-m_f)+(\nu_Q-m_Q)\frac{x_e-x_f}{|x_Q|}-\sum_k\xi_k\varphi_k\Big((\pi_{k,e}-\pi_{k,f})+\big(\ell_{k,Q}-\ip{\pi_k}{\one_Q}\big)\frac{x_e-x_f}{|x_Q|}\Big).
\]
Integrating, using (2) and dividing by $\prod_{Q'}\Gamma(\nu_{Q'})$ gives (v). For homogeneous data $\ell_{k,Q}=\ip{\pi_k}{\one_Q}$ and the second line of the right side of (v) is absent.

\emph{(5) The closing sentences.} A term $k$ raises the shape total of a bunch $Q'$ by $\ip{\pi_k}{\one_{Q'}}$ and its scale index by $\ell_{k,Q'}$, which are equal for homogeneous data; so the ladder, the scale exchange and the shifts $(\mvec+\pi_k,\boldsymbol\nu+\ell_k)$ of (v) keep the condition $\nu_Q=m_Q$, and in (v) the shifts $\mvec+\mathbf e_e$, $\mvec+\mathbf e_f$ of the shape index alone carry the factor $\nu_Q-m_Q$. The shape split and (iv) shift the shape index alone.
\end{proof}

\begin{sremark}
Relation (v) was checked in exact arithmetic on one bunch of two scales with one term whose level is not the degree of its shape exponent, where the Breta has a closed form, for shape exponents $(1,0)$, $(\tfrac32,-\tfrac13)$ and symbolic (\texttt{sweep/s0107/breta\_v.py}).
\end{sremark}

\begin{hproposition}[\hb{prop:letters} (Letter presentation)]
As stated, for a partition of the scales $B$ into any number of bunches.
\end{hproposition}

\begin{proof}
\emph{(1) The presentation.} The bunches of $\cP$ are those of $\cQ$, with the scales $\rho_Q$, and the singletons $\{v_i\}$, whose scales are written $v_i$ and whose shapes are the point $u_{v_i}=1$ of $\Delta_{\{v_i\}}$, with $\omega_{\{v_i\}}$ the Dirac measure (\hb{def:conv}(iii)). The dressed coefficients of the columns are $g_ju^{\pi_j}$ and $C_{i,p}u^{\pi_{i,p}}$, and $u^{(\mvec,\mathbf s)-\one}=u^{\mvec-\one}$. By \hb{def:gramma}, the Gramma on the bunch scales with the configuration $\widetilde\ell$ and these coefficients is
\[
\int\rho^{\boldsymbol\nu-\one}\prod_iv_i^{s_i-1}\exp\Big(-\sum_jg_ju^{\pi_j}\rho^{\lambda_j}-\sum_{i,p}C_{i,p}u^{\pi_{i,p}}\rho^{\ell_{i,p}}v_i\Big)\,d\rho\,dv ,
\]
the level of the pattern $(i,p)$ on $\{v_i\}$ being $1$. For $\re s_i>0$, integrating $v_i$ gives $\Gamma(s_i)\,P_i(\rho,u)^{-s_i}$ (\hb{cor:square} on one scale, $P_i>0$). Dividing by $\prod_Q\Gamma(\nu_Q)\prod_i\Gamma(s_i)$ gives the Zamma (\hb{def:zamma}), and integrating against $u^{\mvec-\one}d\omega(u)$ over $\prod_Q\Delta_Q$ gives $1/\prod_Q\Gamma(\nu_Q)$ times the left side of the display; Tonelli on the moduli and Fubini justify the order of the integrals, either side converging when the other does.

\emph{(2) The relations.} They are \hb{thm:bretalaws}(ii),(iii) for the terms $(\widetilde\pi,\widetilde\ell)$. For the scale exchange at $\{v_i\}$ the columns of nonzero level there are the patterns $(i,p)$, of level $1$, with the factor $\Gst11{s_i+1}{s_i+1}=1$ at $\{v_i\}$; the other singletons give $\Gst11{s_{i'}}{s_{i'}}=1$ and each $Q\in\cQ$ gives $\Gst11{\nu_Q+\ell_{i,p,Q}}{\nu_Q}$: this is the letter split. In the ladders the factor $g_k$ of \hb{thm:bretalaws} is $\prod_Q\Gst11{\nu_Q+\lambda_{j,Q}}{\nu_Q}$ for a kernel column and $s_i\prod_Q\Gst11{\nu_Q+\ell_{i,p,Q}}{\nu_Q}$ for a pattern, by $\Gst11{s_i+1}{s_i}=s_i$. The scale exchange on $Q\in\cQ$ is \hb{thm:bretalaws}(iii) multiplied by $\nu_Q$, with $\nu_Q\Gst11{\nu_Q+\lambda}{\nu_Q+1}=\Gst11{\nu_Q+\lambda}{\nu_Q}$ (\hb{prop:stacked}(iii)), the singletons contributing the factors $1$ for a kernel column and $s_i$ for a pattern.
\end{proof}

\begin{hproposition}[\hb{prop:letterorthant} (Letter data on the orthant)]
As stated, for any levels and any number of bunches.
\end{hproposition}

\begin{proof}
In the left side of the display of \hb{prop:letters} apply \hb{lem:scaleshape} on each bunch $Q\in\cQ$, $\alpha_Q=\rho_Qu_Q$: $u_Q^{\mvec|_Q-\one}\rho_Q^{\nu_Q-1}d\rho_Q\,d\omega_Q=\alpha_Q^{\mvec|_Q-\one}|\alpha_Q|^{\nu_Q-m_Q}d\alpha_Q$, and for $\lambda\in\mathbb R^{\cQ}$, $\pi\in\mathbb R^B$, $\rho^\lambda u^\pi=\alpha^\pi\prod_Q|\alpha_Q|^{\lambda_Q-\ip\pi{\one_Q}}$. So each kernel term becomes $g_j\alpha^{\pi_j}\prod_Q|\alpha_Q|^{\lambda_{j,Q}-\ip{\pi_j}{\one_Q}}$ and each letter $P_i(\alpha)$. Tonelli on the moduli gives the equivalence of convergence and Fubini the equality, and \hb{prop:letters} identifies the left side with $\prod_Q\Gamma(\nu_Q)$ times the Breta. For homogeneous letter data every exponent $\lambda_Q-\ip\pi{\one_Q}$ is $0$, and at $\nu_Q=m_Q$ so is $\nu_Q-m_Q$.
\end{proof}

\begin{htheorem}[\hb{thm:shruffle} (The Shruffle)]
As stated, for any number of bunches.
\end{htheorem}

\begin{proof}
Apply \hb{thm:aggregation} to $A=A_1\sqcup\cdots\sqcup A_q$ with $J=\{1,\dots,q\}$: $u|_{A_i}=w_iv^{(i)}$ with $w\in\Delta_J$, $v^{(i)}\in\Delta_{A_i}$, and
\[
u^{\mvec-\one}d\omega_A=\prod_iw_i^{m_{A_i}-1}\,d\omega_J(w)\prod_i(v^{(i)})^{\mvec|_{A_i}-\one}d\omega_{A_i}(v^{(i)}),\qquad \bigotimes_if_i(u)=\prod_if_i(v^{(i)}) .
\]
Scaling by $w_i>0$ keeps the order of the coordinates of $u|_{A_i}$, which is that of $v^{(i)}$. Up to the null set where two coordinates of $u$ are equal, the points $u$ whose restrictions follow every $\varsigma_i$ form the disjoint union of the sets $\Delta^\varsigma_A$ over the orders $\varsigma$ of $A$ that restrict to every $\varsigma_i$, and they are the images of $\Delta_J\times\prod_i\Delta^{\varsigma_i}_{A_i}$. The integrals converge absolutely, so by Fubini
\begin{multline*}
\sum_\varsigma I^\varsigma_A\Big(\bigotimes_if_i;\mvec\Big)=\int_{\Delta_J}\prod_iw_i^{m_{A_i}-1}\,d\omega_J(w)\ \prod_iI^{\varsigma_i}_{A_i}(f_i;\mvec|_{A_i})\\=\Gst q1{m_{A_1},\ \dots,\ m_{A_q}}{m_A}\prod_iI^{\varsigma_i}_{A_i}(f_i;\mvec|_{A_i}),
\end{multline*}
by \hb{thm:dirichlet} on $J$, where $\re m_{A_i}>0$. Dividing by the stacked Gamma, whose reciprocal is $\Gst1q{m_A}{m_{A_1},\ \dots,\ m_{A_q}}$, gives the statement; at $q=2$ it is the factor $\Gst12{m_A+m_B}{m_A,\ m_B}$ of the reference.
\end{proof}

\begin{htheorem}[\hb{thm:betaport} (Relations of the Beta port)]
As stated: (iii) for linear letters of either slope and complex exponents, and (v) at every pair of exponents; (i), (ii) and (iv) are those of the reference.
\end{htheorem}

\begin{proof}
\emph{(iii)} On $[0,1]$ every letter $L_i=c_i+d_ix$ is positive, since it is affine with the positive values $c_i$ and $c_i+d_i$ at the ends, so $L_i^{-w_i}$ is infinitely differentiable there for every $w_i\in\mathbb C$, with $\partial_x^vL_i^{-w_i}=(-1)^v(w_i)_vd_i^vL_i^{-w_i-v}$, $d_i$ real. Leibniz's rule gives
\[
\partial_x^q\prod_iL_i^{-w_i}=\sum_{|v|=q}\frac{q!}{\prod_iv_i!}\prod_i(-1)^{v_i}(w_i)_{v_i}d_i^{v_i}L_i^{-w_i-v_i},
\]
which is the first display. Then (ii), with $g=\prod_iL_i^{-w_i}$ and $\re a,\re b>0$, gives the second display, by $x^{a+q-1}(1-x)^{b+q-1}=x^{a-1}(1-x)^{b-1}\prod_i(x(1-x))^{v_i}$ for $|v|=q$. Let $a,b>0$ and $w_i\ge0$. Each integrand on the right is $x^{a-1}(1-x)^{b-1}\prod_iL_i^{-w_i}>0$ times $\prod_i(x(1-x)d_i/L_i)^{v_i}$, which is nonnegative if every $d_i\ge0$, and is $(-1)^{|v|}=(-1)^q$ times a nonnegative function if every $d_i\le0$; the coefficients $\prod_i(w_i)_{v_i}/v_i!$ are nonnegative. This gives the two signs. If $x(1-x)|d_i|/L_i\le\varepsilon$ on $[0,1]$ for every $i$, each integrand has modulus at most $\varepsilon^q x^{a-1}(1-x)^{b-1}\prod_iL_i^{-w_i}$, and
\[
\sum_{|v|=q}\prod_i\frac{(w_i)_{v_i}}{v_i!}=\frac{(|\mathbf w|)_q}{q!},
\]
the coefficient of $y^q$ on both sides of $\prod_i(1-y)^{-w_i}=(1-y)^{-|\mathbf w|}$, $|y|<1$, by $(1-y)^{-w}=\sum_v\frac{(w)_v}{v!}y^v$. This gives the bound.

\emph{(v)} Fix $q\ge1$, write $s=a+b$, and first let $\re a,\re b>0$.

(1) Coefficients. Expanding $x^k(1-x)^{q-k}$ in the second form of \hb{def:betaport},
\[
[x^j]\,\Pi^{a,b}_q=(-1)^{q+j}\sum_k\binom{a+q-1}{q-k}\binom{q-k}{q-j}\binom{b+q-1}{k}.
\]
By $\binom NM\binom ML=\binom NL\binom{N-L}{M-L}$, a polynomial identity in $N$, at $N=a+q-1$, $M=q-k$, $L=q-j$, this is $(-1)^{q+j}\binom{a+q-1}{q-j}\sum_k\binom{a+j-1}{j-k}\binom{b+q-1}{k}$, and $\sum_k\binom{A}{j-k}\binom{B}{k}=\binom{A+B}{j}$, a polynomial identity in $A$ and $B$ that holds at nonnegative integers by counting the $j$-subsets of a disjoint union, gives $\binom{s+q+j-2}{j}$: the coefficient formula. At $j=q$, $\lambda_q:=[x^q]\Pi^{a,b}_q=\binom{s+2q-2}q=(s+q-1)_q/q!$, which is nonzero since $\re s>0$. Also
\[
\frac{[x^{q-1}]\Pi^{a,b}_q}{\lambda_q}=-\frac{q(a+q-1)}{s+2q-2},\qquad\frac{[x^q]\Pi^{a,b}_{q+1}}{\lambda_{q+1}}=-\frac{(q+1)(a+q)}{s+2q}.
\]

(2) Pairing. Put $\ip gh_\circ=\int_0^1x^{a-1}(1-x)^{b-1}g\,h\,dx$ for polynomials $g,h$. By (i) at $c=a+j$, $\ip{\Pi^{a,b}_k}{x^j}_\circ=\binom jk\Gst21{a+j,\ b+k}{s+j+k}$, which is $0$ for $j<k$ and $\beta_k:=\Gst21{a+k,\ b+k}{s+2k}\neq0$ for $j=k$.

(3) Span. Let $g$ be a polynomial of degree at most $q+1$ with $\ip g{x^j}_\circ=0$ for $j\le q-2$. Since $\deg\Pi^{a,b}_k=k$ by (1), $g=\sum_{k\le q+1}\gamma_k\Pi^{a,b}_k$; pairing with $x^0,x^1,\dots,x^{q-2}$ in turn, by (2) the pairing with $x^j$ involves only $\gamma_0,\dots,\gamma_j$, with $\gamma_j$ multiplied by $\beta_j\neq0$, so $\gamma_0=\dots=\gamma_{q-2}=0$ and $g\in\operatorname{span}(\Pi^{a,b}_{q-1},\Pi^{a,b}_q,\Pi^{a,b}_{q+1})$.

(4) Multiplication by $x$. $\ip{x\Pi^{a,b}_q}{x^j}_\circ=\ip{\Pi^{a,b}_q}{x^{j+1}}_\circ=0$ for $j\le q-2$, so by (3) $x\Pi^{a,b}_q=A\Pi^{a,b}_{q+1}+B\Pi^{a,b}_q+C\Pi^{a,b}_{q-1}$. The coefficients of $x^{q+1}$ give $A=\lambda_q/\lambda_{q+1}=\frac{(q+1)(s+q-1)}{(s+2q)(s+2q-1)}$. Pairing with $x^{q-1}$ gives $\beta_q=C\beta_{q-1}$, so $C=\frac{(a+q-1)(b+q-1)}{(s+2q-1)(s+2q-2)}$. The coefficients of $x^q$ give $[x^{q-1}]\Pi^{a,b}_q=A\,[x^q]\Pi^{a,b}_{q+1}+B\lambda_q$, so by (1) $B=\frac{(q+1)(a+q)}{s+2q}-\frac{q(a+q-1)}{s+2q-2}$.

(5) The operator $D=x(1-x)\partial_x$. For a polynomial $g$, integrating by parts, the boundary terms $x^{a+j}(1-x)^bg$ vanishing at both ends,
\[
\ip{Dg}{x^j}_\circ=\int_0^1x^{a+j}(1-x)^bg'\,dx=-(a+j)\ip g{x^j}_\circ+(s+j)\ip g{x^{j+1}}_\circ ,
\]
which is $0$ at $g=\Pi^{a,b}_q$ for $j\le q-2$. By (3), $D\Pi^{a,b}_q\in\operatorname{span}(\Pi^{a,b}_{q-1},\Pi^{a,b}_q,\Pi^{a,b}_{q+1})$, and eliminating $\Pi^{a,b}_{q+1}$ by (4), $A\neq0$, $D\Pi^{a,b}_q=(e\,x+e_0)\Pi^{a,b}_q+e_1\Pi^{a,b}_{q-1}$. Since $Dx^k=kx^k-kx^{k+1}$, the coefficients of $x^{q+1}$ give $e=-q$. Pairing with $x^{q-1}$: the left side is $(s+q-1)\beta_q$ and the right side $-q\beta_q+e_1\beta_{q-1}$, so $e_1=(s+2q-1)\beta_q/\beta_{q-1}=\frac{(a+q-1)(b+q-1)}{s+2q-2}$. The coefficients of $x^q$ give $q\lambda_q-(q-1)[x^{q-1}]\Pi^{a,b}_q=-q[x^{q-1}]\Pi^{a,b}_q+e_0\lambda_q$, so $e_0=q+[x^{q-1}]\Pi^{a,b}_q/\lambda_q=\frac{q(b+q-1)}{s+2q-2}$.

(6) Continuation. The three displays compare polynomials in $x$ whose coefficients are rational functions of $(a,b)$, the coefficients of the Beta ports being polynomials in $(a,b)$. They agree on the open set $\re a,\re b>0$, hence wherever the coefficients are finite, which is the condition $a+b+2q-i\neq0$, $i=0,1,2$; the coefficient formula is a polynomial identity and holds for all $a,b$ and every $q\in\mathbb N$.
\end{proof}

\begin{sremark}
At $a=b=n+1$ the middle coefficient of the three-term relation is $\tfrac12$, and the relations become those of the reference: $(1-2x)\Pi^{n+1,n+1}_r=-2A\,\Pi^{n+1,n+1}_{r+1}-2C\,\Pi^{n+1,n+1}_{r-1}$ and $x(1-x)\partial_x\Pi^{n+1,n+1}_r=\frac r2(1-2x)\Pi^{n+1,n+1}_r+\frac{r+n}2\Pi^{n+1,n+1}_{r-1}$, with the leading coefficient $\binom{2n+2r}r$, and at $a=b=1$ the coefficient formula gives $(-1)^r\Pi^{1,1}_r=\sum_k(-1)^k\binom rk\binom{r+k}kx^k$. The three displays and the coefficient formula were checked in exact arithmetic with $a$ and $b$ symbolic for $q\le9$, together with these diagonal forms (\texttt{sweep/s0107/betaport.py}).
\end{sremark}

\section{Chapter~\hb{ch:fibre}: fibre integrals and closing}\label{sec:ch04}

\begin{htheorem}[\hb{thm:configgauge} (The native gauge of a configuration)]
As stated: (iii) with a fixed offset $\nu'$ of the index; (i) and (ii) are those of the reference.
\end{htheorem}

\begin{proof}
By (i), $\Phi_{\boldsymbol\xi}$ is smooth and strictly convex with its minimum at $y^\star$, where its Hessian $H^\star=\sum_kw^\star_kt_kt_k^{\mathsf T}$ is positive definite.

\emph{(1) The integral.} In the logarithmic form of \hb{def:gramma},
\[
\sG_{\lambda\nu+\nu';t}(\lambda\boldsymbol\xi)=\int_{\mathbb R^{\cE}}e^{\ip{\nu'}y}\,e^{-\lambda\Phi_{\boldsymbol\xi}(y)}\,dy ,
\]
which converges when $\lambda\nu+\nu'\in\cT_t$ (\hb{thm:grammadomain}). Since $\interior\cone(t)$ is an open cone containing $\nu$, it contains $\nu+\lambda^{-1}\re\nu'$ for $\lambda$ large, so the integral converges for $\lambda\ge\lambda_0$, some $\lambda_0>0$.

\emph{(2) Away from $y^\star$.} Fix $\delta>0$ and put $\mu_\delta=\min_{|y-y^\star|=\delta}\big(\Phi_{\boldsymbol\xi}(y)-\Phi_{\boldsymbol\xi}(y^\star)\big)>0$. For a unit vector $d$ the function $\sigma\mapsto\Phi_{\boldsymbol\xi}(y^\star+\sigma d)-\Phi_{\boldsymbol\xi}(y^\star)$ is convex and vanishes at $\sigma=0$, so for $\sigma\ge\delta$ it is at least $\sigma\mu_\delta/\delta$. With $|e^{\ip{\nu'}y}|\le e^{|\re\nu'|\,|y^\star|}e^{|\re\nu'|\,|y-y^\star|}$,
\[
\Big|\int_{|y-y^\star|\ge\delta}e^{\ip{\nu'}y}e^{-\lambda(\Phi_{\boldsymbol\xi}(y)-\Phi_{\boldsymbol\xi}(y^\star))}\,dy\Big|\le e^{|\re\nu'|\,|y^\star|}\int_{|u|\ge\delta}e^{-(\lambda\mu_\delta/\delta-|\re\nu'|)|u|}\,du ,
\]
which, once $\lambda\mu_\delta/\delta\ge2|\re\nu'|$ and $\lambda\mu_\delta\ge2\delta$, is at most a constant times $e^{-\lambda\mu_\delta/4}$: on $|u|\ge\delta$, $e^{-\Lambda|u|}\le e^{-\Lambda\delta/2}e^{-\Lambda|u|/2}$ for $\Lambda=\lambda\mu_\delta/\delta-|\re\nu'|\ge\lambda\mu_\delta/(2\delta)\ge1$. So it is $O(\lambda^{-|\cE|/2-1})$.

\emph{(3) Near $y^\star$.} Laplace's method at the nondegenerate minimum $y^\star$, with the smooth amplitude $e^{\ip{\nu'}y}$, gives
\[
\int_{|y-y^\star|<\delta}e^{\ip{\nu'}y}e^{-\lambda(\Phi_{\boldsymbol\xi}(y)-\Phi_{\boldsymbol\xi}(y^\star))}\,dy=\Big(\frac{2\pi}\lambda\Big)^{|\cE|/2}\det(H^\star)^{-1/2}\Big(e^{\ip{\nu'}{y^\star}}+O(\lambda^{-1})\Big),
\]
as in the reference, proof of \hb{thm:configgauge}, where the amplitude is $1$: the amplitude enters the leading term through its value at $y^\star$, and its derivatives only the term $O(\lambda^{-1})$.

Adding (2) and (3), multiplying by $e^{-\lambda\Phi_{\boldsymbol\xi}(y^\star)}$ and factoring out $e^{\ip{\nu'}{y^\star}}\neq0$ gives (iii).
\end{proof}

\begin{hcorollary}[\hb{cor:fourgammas} (Four Gammas on a line)]
As stated: the display is \hb{thm:sharemerge} on two coefficients.
\end{hcorollary}

\begin{proof}
\emph{(1) A straight contour.} Let first $\max(-\re a,-\re b)<\min(\re c,\re d)$, and take the line $\re z=\gamma$ with $\max(-\re a,-\re b)<\gamma<\min(\re c,\re d)$; it separates the two sets of poles. Apply \hb{thm:sharemerge} with $k=2$, the exponents $a+c$ and $b+d$, and $(a+d,\,b+c)$, of total $a+b+c+d$, in the place of the merged shares. In the shares $(\alpha_1,\alpha_2,\beta_1,\beta_2)$ the fibre has the pins $\alpha_1+\alpha_2=a+c$, $\alpha_1+\beta_1=a+d$, $\alpha_2+\beta_2=b+c$, and it is parametrized by $\alpha=(a+z,\,c-z)$, $\beta=(d-z,\,b+z)$. The matrix $P$ of the pins has the rows $(1,1,0,0)$, $(1,0,1,0)$, $(0,1,0,1)$, so $\det PP^{\mathsf T}=4$, and its kernel is spanned by $(1,-1,-1,1)$, of length $2$. By \hb{def:fibre}, with the admissible real point on the line $\re z=\gamma$, where all four shares have positive real parts, the fibre integral of a function $G$ is $(2\pi)^{-1}\cdot4^{-1/2}\cdot2\int G\,d\tau=\frac1{2\pi i}\int_{\gamma+i\mathbb R}G\,dz$, $z=\gamma+i\tau$. The integrand of \hb{thm:sharemerge} is
\[
\Gst21{a+z,\ c-z}{a+c}\Gst21{d-z,\ b+z}{b+d}=\frac{\Gamma(a+z)\Gamma(b+z)\Gamma(c-z)\Gamma(d-z)}{\Gamma(a+c)\Gamma(b+d)},
\]
and its right side is $\Gst21{a+d,\ b+c}{a+b+c+d}$. Multiplying by $\Gamma(a+c)\Gamma(b+d)$ (\hb{prop:stacked}(i)) gives the display.

\emph{(2) Continuation.} A contour separating the poles exists exactly when no pole $-a-n$ or $-b-n$ of the first two factors coincides with a pole $c+n'$ or $d+n'$ of the last two, $n,n'\in\mathbb N$, that is on the set $U$ of $(a,b,c,d)\in\mathbb C^4$ with $a+c,a+d,b+c,b+d\notin-\mathbb N$. Near a point of $U$ one fixed contour separates the poles, and by \hb{prop:stacked}(iii) and \hb{thm:decay} on one scale the four Gammas decay exponentially along it, locally uniformly in $(a,b,c,d)$; so the left side is holomorphic on $U$, and by Cauchy's theorem it does not depend on the separating contour. The right side is holomorphic on $U$. The set $U$ is connected, being the complement in $\mathbb C^4$ of countably many complex hyperplanes, and contains the open set of (1), so the identity holds on $U$.
\end{proof}

\section{Chapter~\hb{ch:trail}: the Trailmoment}\label{sec:ch05}

\begin{htheorem}[\hb{thm:gaugecoef} (The gauge of the coefficient function)]
As stated: the last sentence at the unit configuration on any number of scales.
\end{htheorem}

\begin{proof}
The first two sentences are those of the reference (the reference, proof of \hb{thm:gaugecoef}). At the unit configuration $\sZ_{\nu;I_{\cE}}(\mathbf c)=\prod_ec_e^{-\nu_e}=\mathbf c^{-\nu}$ for every $\nu\in\mathbb C^{\cE}$ (\hb{def:zamma}), so by \hb{def:jointtrail} $(\sT_{\nu;I_{\cE}}[\mathbf c]\,\phi)(Y')=\int_{Y'}\mathbf c^{-\nu}\,d\phi$, either side converging when the other does.
\end{proof}

\begin{htheorem}[\hb{thm:exchangetrail} (Exchange and ladder of the Trailmoment)]
As stated: (iii) at the unit configuration on any number of scales; (i) and (ii) are those of the reference.
\end{htheorem}

\begin{proof}
At $t=I_{\cE}$ the column $t_k$ is $\mathbf e_k$, $k\in\cE$, so in (i) at the scale $e$ only $k=e$ has $t_{k,e}\neq0$, with $t_{e,e}=1$ and the factor $\Gst11{\nu_e+1}{\nu_e+1}\prod_{f\neq e}\Gst11{\nu_f}{\nu_f}=1$. Directly, for every $\nu$, $\sZ_{\nu;I_{\cE}}(\mathbf c)=\mathbf c^{-\nu}=c_e\,\mathbf c^{-\nu-\mathbf e_e}=c_e\,\sZ_{\nu+\mathbf e_e;I_{\cE}}(\mathbf c)$ (\hb{def:zamma}); integrating against $\phi$ over $Y'$ gives the first identity.

For the derivative, let $D\subseteq\mathbb C^{\cE}$ be the set of indices where the Trailmoment converges, $\int_{Y'}\mathbf c^{-\re\nu}\,d|\phi|<\infty$, and $\nu_0$ an interior point of $D$. Take $\delta>0$ such that $\nu_0\pm\delta\mathbf e_e\in D$. For $\nu=\nu_0+\zeta\mathbf e_e$ with $|\zeta|<\delta/2$, $|\mathbf c^{-\nu}\log c_e|=\mathbf c^{-\re\nu_0}c_e^{-\re\zeta}|\log c_e|$. For $\eta>0$ and $y>0$, $|\log y|\le\eta^{-1}(y^\eta+y^{-\eta})$; with $\eta=\delta/2$, and $y^{\theta}\le y^{\delta}+y^{-\delta}$ for $|\theta|\le\delta$,
\[
\big|\mathbf c^{-\nu}\log c_e\big|\le\frac4\delta\Big(\mathbf c^{-(\re\nu_0-\delta\mathbf e_e)}+\mathbf c^{-(\re\nu_0+\delta\mathbf e_e)}\Big),
\]
an integrable function independent of $\zeta$. So $\zeta\mapsto\int_{Y'}\mathbf c^{-\nu_0-\zeta\mathbf e_e}d\phi$ may be differentiated under the integral sign, and $\partial_\zeta\mathbf c^{-\nu}=-(\log c_e)\,\mathbf c^{-\nu}$ gives the second identity.
\end{proof}

\begin{hproposition}[\hb{prop:resonance} (Laurent data of the Gamma at its poles)]
As stated: the residue and the Laurent coefficients of a product of any number of Gammas with a common pole.
\end{hproposition}

\begin{proof}
The quotient form of $(-1)^j\varepsilon\Gamma(-j+\varepsilon)$, the expansion of its logarithm and $(\log\Gamma)'(j+1)=-\Tx{1}{\IZ}{j}$ are those of the reference (the reference, proof of \hb{prop:resonance}). On $|\varepsilon|<1$ each $\phi_i(\varepsilon):=(-1)^{j_i}\varepsilon\,\Gamma(-j_i+\varepsilon)=\Gamma(1+\varepsilon)\Gamma(1-\varepsilon)/\Gamma(1+j_i-\varepsilon)$ is holomorphic and nonzero, and it is the exponential of the expansion at $j_i$. Put $b=-j+\varepsilon$; then $b+k_i=-j_i+\varepsilon$ and $\Gamma(b+k_i)=(-1)^{j_i}\varepsilon^{-1}\phi_i(\varepsilon)$, so
\[
h(b)\prod_{i=1}^p\Gamma(b+k_i)=(-1)^{j_1+\cdots+j_p}\,\varepsilon^{-p}\,h(-j+\varepsilon)\prod_{i=1}^p\phi_i(\varepsilon).
\]
The residue at $\varepsilon=0$ is $(-1)^{\sum_ij_i}$ times the coefficient of $\varepsilon^{p-1}$ in $h(-j+\varepsilon)\prod_i\phi_i(\varepsilon)$, and with $h(-j+\varepsilon)=\sum_lh^{(l)}(-j)\varepsilon^l/l!$ the Cauchy product gives the displayed sum. At $p=2$ and $k_1=k_2=0$, $\phi_1\phi_2=\exp(2\log\phi_1)=(j!)^{-2}\big(1-2\Tx{1}{\IZ}{j}\varepsilon+O(\varepsilon^2)\big)$, so the residue is $(j!)^{-2}\big(h'(-j)-2\Tx{1}{\IZ}{j}h(-j)\big)$, which is the formula of the reference. The Laurent coefficients of $\prod_i\Gamma(b+k_i)$ at $-j$ are $(-1)^{\sum_ij_i}$ times the Taylor coefficients of $\prod_i\phi_i=\prod_i(j_i!)^{-1}\exp\big(\sum_i(\log\phi_i+\log j_i!)\big)$; the series $\sum_i(\log\phi_i+\log j_i!)$ has no constant term and its coefficients are rational multiples of the Trailmoments $\Tx{k}{\IZ}{j_i}$ and $\Tx{2k}{\IZ}{0}$, so the Taylor coefficients of its exponential are polynomials in these Trailmoments with rational coefficients.
\end{proof}

\begin{sremark}
The residue formula was checked in exact arithmetic for $p=2,3,4$ at several $j$ and $(k_i)$ (\texttt{sweep/s0107/resonance.py}).
\end{sremark}

\begin{hproposition}[\hb{prop:jointmoments} (The joint Laplace transform of Dirichlet inputs)]
As stated, for finite sums of monomials with nonnegative real exponents.
\end{hproposition}

\begin{proof}
\emph{(1) The series.} Every exponent $\mathbf a$ lies in $[0,\infty)^{\cE}$, so $0<u^{\mathbf a}\le1$ on the source $\prod_Q\Delta_Q$ and $|P_i(u)|\le\|P_i\|:=\sum_{\mathbf a}\big|[u^{\mathbf a}]P_i\big|$. Hence $\sum_{\mathbf k}\frac{|v|^{\mathbf k}}{\mathbf k!}\prod_i|P_i(U)|^{k_i}\le\exp\big(\sum_i|v_i|\,\|P_i\|\big)$, the exponential series of $\exp(-\sum_iv_iP_i(U))$ converges uniformly on the source for each $v\in\mathbb C^q$ and integrates term by term, and $\mathcal F(v)=\sum_{\mathbf k}\frac{(-v)^{\mathbf k}}{\mathbf k!}\oE[\prod_iP_i(U)^{k_i}]$, a power series converging on $\mathbb C^q$, hence entire.

\emph{(2) The moments.} Multiplying out, $\prod_iP_i^{k_i}=\sum_{\mathbf a}[u^{\mathbf a}]\big(\prod_iP_i^{k_i}\big)u^{\mathbf a}$, a finite sum with exponents in $[0,\infty)^{\cE}$. The blocks of $U$ are independent, so $\oE[U^{\mathbf a}]=\prod_Q\oE_{U_Q\sim\Dir_{\mvec|_Q}}[U_Q^{\mathbf a|_Q}]$, and by \hb{thm:dirmoments} at the exponents $\mathbf a|_Q$, whose real parts satisfy $m_e+a_e>0$, together with \hb{prop:stacked}(i),(ii),
\[
\oE_{U_Q\sim\Dir_{\mvec|_Q}}\big[U_Q^{\mathbf a|_Q}\big]=\frac{\Gst{|Q|}1{\mvec|_Q+\mathbf a|_Q}{m_Q+a_Q}}{\Gst{|Q|}1{\mvec|_Q}{m_Q}}=\prod_{e\in Q}\Gst11{m_e+a_e}{m_e}\ \Gst11{m_Q}{m_Q+a_Q} .
\]
Summing gives $M(\mathbf k)$; for integer exponents \hb{thm:dirmoments} gives the rising factorials.

\emph{(3) The read.} Where every $P_i$ is positive on the interior of the source, which carries the law, \hb{def:zamma} and \hb{def:gramma} give $\sZ_{\nu;t}(\mathbf P(U))=\frac1{\prod_j\Gamma(\nu_j)}\int_{(0,\infty)^d}x^{\nu-\one}\exp\big(-\sum_iP_i(U)x^{t_i}\big)dx$ where this converges. By Tonelli on the moduli and Fubini, wherever the double integral converges, the expectation of the left side is the integral of $\oE[\exp(-\sum_iP_i(U)x^{t_i})]=\mathcal F(x^{t_1},\dots,x^{t_q})$ against $x^{\nu-\one}dx/\prod_j\Gamma(\nu_j)$. The left side is the Trailmoment $(\sT_{\nu;t}[(P_i)_i]\,\otimes_Q\Dir_{\mvec|_Q})(\prod_Q\Delta_Q)$ by \hb{def:jointtrail}.

\emph{(4) Uniqueness.} With real coefficients the vector $(P_i(U))_i$ takes its values in the compact box $\prod_i[-\|P_i\|,\|P_i\|]$, and its moments $\oE[\prod_iP_i(U)^{k_i}]=M(\mathbf k)$ are the integrals of the monomials against its law. Polynomials are dense in the continuous functions on the box, so a law on the box is determined by these integrals; and the read $\oE[\sZ_{\nu;t}(\mathbf P(U))]$ depends on the law of $\mathbf P(U)$ alone.
\end{proof}

\begin{htheorem}[\hb{thm:shrufflefree} (The Shruffle on a line)]
As stated: (i) and (iii) for any number of strings and of trails; (ii) is that of the reference.
\end{htheorem}

\begin{proof}
(i) Let $w^{(i)}$ have length $n_i$, $n=\sum_in_i$, and $S_i=\{x^{(i)}_1>\cdots>x^{(i)}_{n_i}>0\}$. By Tonelli the function $\prod_i\prod_jw^{(i)}_j(x^{(i)}_j)$ is integrable on $S_1\times\cdots\times S_r$, with integral $\prod_i\operatorname{chain}(w^{(i)})$ by Fubini. The points with two equal coordinates form a null set. At every other point the $n$ coordinates are distinct, and their decreasing order is a total order of the $n$ positions $(i,j)$ that restricts to the order of each string: an interleaving $\varsigma$, the string of the $n$ functions in that order. The points of one interleaving are carried by a permutation of the coordinates, which keeps Lebesgue measure, onto $\{y_1>\cdots>y_n>0\}$, the integrand becoming $\prod_m\varsigma_m(y_m)$; so their integral is $\operatorname{chain}(\varsigma)$. The interleavings are the ways of distributing the $n$ places among the strings, $n!/(n_1!\cdots n_r!)$ of them.

(iii) Let $r$ trails be as in (ii), with strings $\sigma^{(1)},\dots,\sigma^{(r)}$ and cuts $N_1,\dots,N_r$. By (ii) the trail $i$ is $\operatorname{chain}(\sigma^{(i)})$, and by (i) the product is $\sum_\varsigma\operatorname{chain}(\varsigma)$ over the interleavings. An interleaving begins with the first function of one of the strings, a Laplace transform $e^{-N_i\,\cdot}L^{(i)}_1$, which is $\int e^{-(s+N_i)\,\cdot}\,\psi(ds)$ for the first increment $\psi$ of trail $i$: the Laplace transform of its image under $s\mapsto s+N_i$, carried by $(0,\infty)$ by the hypothesis of (ii). The other Laplace transforms of $\varsigma$ are those of nonnegative increments or of such images. Reading $\varsigma$ as in (ii), each Laplace transform followed by its run of constant functions $1$ gives an increment with the index one more than the length of the run; the trail at the cut $0$ with these increments has no constant columns and positive integer indices, and by (ii) it is $\operatorname{chain}(\varsigma)$.
\end{proof}

\begin{htheorem}[\hb{thm:ct-trailchain} (Trails at frequencies are iterated integrals)]
As stated.
\end{htheorem}

\begin{proof}
For a word $\ell=(\ell_1,\dots,\ell_w)$ with $\ell_1\neq0$ and letters in $\{0\}\cup\{|z|=1\}$ put $F_\varnothing=1$ and $F_\ell(z)=\int_0^zF_{\ell'}(t)\,dt/(t-\ell_w)$ on the disc $|z|<1$, $\ell'$ the word without its last letter, along the segment. By induction $F_\ell$ is holomorphic on the disc with $F_\ell(0)=0$ for $w\ge1$: for $\ell_w\neq0$ the factor $1/(t-\ell_w)$ is holomorphic there, and for $\ell_w=0$, $F_{\ell'}(t)/t$ is, as $F_{\ell'}(0)=0$, the word $\ell'$ being nonempty since $\ell_1\neq0$. Unwinding the recursion, $F_\ell(z)$ is the iterated integral over $0<t_1<\cdots<t_w<z$ along the segment.

For $|z|<1$ let $S(z)=\sum_{0<n_1<\cdots<n_d}\prod_i\zeta_i^{n_i}n_i^{-\nu_i}z^{n_d}$, absolutely convergent, the trail with $z^{y_d}$ in its density, $S(0)=0$. If $\nu_d\ge2$, $zS'(z)$ is the same series with $\nu_d-1$ in place of $\nu_d$, whose word is $\ell$ without its last letter $0$, and $S(z)=\int_0^zS_{\nu_d-1}(t)\,dt/t$. If $\nu_d=1$, summing the geometric series over $n_d>n_{d-1}$,
\[
S'(z)=\sum_{0<n_1<\cdots<n_{d-1}}\prod_{i<d}\zeta_i^{n_i}n_i^{-\nu_i}\ \frac{\zeta_d\,(\zeta_dz)^{n_{d-1}}}{1-\zeta_dz}=\frac{\zeta_d}{1-\zeta_dz}\,S^\flat(z),
\]
where $S^\flat$ is the series of depth $d-1$ with the last frequency $a_{d-1}+a_d$, so $\zeta_{d-1}\zeta_d$, and $S^\flat=1$ when $d=1$. Its $c_i$ are those of $S$ for $i<d$, its word is $\ell$ without the last letter $c_d^{-1}=\zeta_d^{-1}$, and $\zeta_d/(1-\zeta_dz)=-1/(z-\zeta_d^{-1})$, so $S(z)=-\int_0^zS^\flat(t)\,dt/(t-\zeta_d^{-1})$. By induction on $w$, $S=(-1)^dF_\ell$ on the disc, which is the last sentence.

Write $S(z)=\sum_nb_nz^n$ with $b_n=A(n)\zeta_d^nn^{-\nu_d}$, $A(n)=\sum_{0<n_1<\cdots<n_{d-1}<n}\prod_{i<d}\zeta_i^{n_i}n_i^{-\nu_i}$, $A=1$ when $d=1$; then $|A(n)|\le(1+\log n)^{d-1}$, and the head over $y_d\le N$ is $\sum_{n\le N}b_n$. For $\nu_d\ge2$ the series $\sum b_n$ converges absolutely, so the heads converge to the trail. For $\nu_d=1$ the partial sums of $\zeta_d^n$ are bounded by $2/|1-\zeta_d|$, and $A(n)/n\to0$ with $\sum_n|A(n+1)/(n+1)-A(n)/n|\le\sum_nC(1+\log n)^{d-1}n^{-2}<\infty$, since $A(n+1)-A(n)$ is the sum of the terms with $n_{d-1}=n$; summation by parts makes $\sum b_n$ converge. By Abel's theorem on power series, $S(z)\to\sum_nb_n$ as $z\to1$ along $[0,1)$.

On the other side, the integrand $\prod_j1/|t_j-\ell_j|$ on the simplex of $[0,1]$ is at most $C\prod_{\ell_j=0}t_j^{-1}\prod_{\ell_j=1}(1-t_j)^{-1}$, the other letters lying off $[0,1]$. Split the simplex by the number of the $t_j$ below $\frac12$; on the part below, with $\ell_1\neq0$, the iterated integrals $G_j(t)$ of the first $j$ factors over $0<t_1<\cdots<t_j<t$ satisfy $G_1(t)\le Ct$ and $G_j(t)\le\int_0^tG_{j-1}(s)\,C(1+1/s)\,ds\le C_jt$; on the part above the same holds after $t\mapsto1-t$, as $\ell_w\neq1$ ($\ell_w=0$ when $\nu_d\ge2$, $\ell_w=\zeta_d^{-1}\neq1$ when $\nu_d=1$). So the integral over $0<t_1<\cdots<t_w<1$ converges absolutely, and by dominated convergence $F_\ell(z)$ tends to it as $z\to1$ along $[0,1)$. Hence the limit of the heads is $(-1)^d$ times that integral.
\end{proof}

\begin{htheorem}[\hb{thm:twoconv} (Two convolutions)]
As stated: (ii) and (iii) at the unit configuration on any number of scales; (i) and the displays of (ii) for inputs carried by $(0,\infty)$ are those of the reference.
\end{htheorem}

\begin{proof}
(ii) By \hb{def:zamma}, $\sZ_{\nu;I_{\cE}}(\mathbf c\,\mathbf c')=\prod_e(c_ec'_e)^{-\nu_e}=\sZ_{\nu;I_{\cE}}(\mathbf c)\,\sZ_{\nu;I_{\cE}}(\mathbf c')$ at every point of $Y_1\times Y_2$. The two Trailmoments on the right converge, so by Tonelli the kernel is integrable for $|\phi|\otimes|\psi|$ on $Y_1'\times Y_2'$, and Fubini gives the product (\hb{def:jointtrail}).

(iii) \emph{The kernel on sums.} First, $g$ has no zero: if $g(s_0)=0$, then $g(s_0/2)^2=g(s_0)=0$, so $g(s_0/2^n)=0$ for every $n$; every $s\in(0,\infty)^{\cE}$ has $s-s_0/2^n\in(0,\infty)^{\cE}$ for $n$ large, and $g(s)=g(s_0/2^n)\,g(s-s_0/2^n)=0$, against the hypothesis. Fix $e\in\cE$. For $s,s''\in(0,\infty)^{\cE}$ and $\tau>0$, $g(s+\tau\mathbf e_e)\,g(s'')=g(s+s''+\tau\mathbf e_e)=g(s''+\tau\mathbf e_e)\,g(s)$, so $\varphi_e(\tau)=g(s+\tau\mathbf e_e)/g(s)$ does not depend on $s$. It is continuous on $(0,\infty)$, has no zero, and $\varphi_e(\tau+\tau')=\varphi_e(\tau)\,\varphi_e(\tau')$, by the independence of the base point applied at $s$ and at $s+\tau\mathbf e_e$. By the converse on one scale, the case $|\cE|=1$ (the reference, proof of \hb{thm:twoconv}), $\varphi_e(\tau)=e^{-u_e\tau}$ for one $u_e\in\mathbb C$, and then $g(s+\tau\mathbf e_e)=e^{-u_e\tau}g(s)$ also for $\tau<0$ with $s+\tau\mathbf e_e\in(0,\infty)^{\cE}$, applying it at the base point $s+\tau\mathbf e_e$. Moving one coordinate at a time from a fixed point $s_1$, $g(s)=C\,e^{-\ip us}$ on $(0,\infty)^{\cE}$ with $C=g(s_1)e^{\ip u{s_1}}\neq0$, and $g(s+s')=g(s)\,g(s')$ gives $C=C^2$, so $C=1$.

\emph{The kernel on products.} If $h(x)\neq0$, then $h(x)=h(x)\,h(\one)$ gives $h(\one)=1$, and $h(x)\,h(x^{-1})=h(\one)=1$ shows that $h$ has no zero. The function $s\mapsto h(e^s)$ on $(0,\infty)^{\cE}$ is continuous, not identically zero and satisfies the hypothesis on sums, so $h(e^s)=e^{-\ip us}$, that is $h(x)=x^{-u}$ for every $x$ with all coordinates above $1$. Any $x\in(0,\infty)^{\cE}$ is $x'/x''$ with $x''_e>\max(1,1/x_e)$ and $x'=x\,x''$, both with all coordinates above $1$, and $h(x)=h(x')/h(x'')=x^{-u}$. So $h=\sZ_{u;I_{\cE}}$.
\end{proof}

\section{Chapter~\hb{ch:transport}: rational words and live inputs}\label{sec:ch06}

\begin{htheorem}[\hb{thm:columns} (Pole columns are a rational difference quotient)]
As stated: coefficients in any subfield of $\mathbb C$, and no decay condition.
\end{htheorem}

\begin{proof}
\emph{(1) The decomposition.} For $k\ge0$ and $s\ge1$, telescoping gives $\Delta\sum_{j=0}^{k-1}(t+j)^{-s}=\sum_{j=0}^{k-1}\big((t+j+1)^{-s}-(t+j)^{-s}\big)=(t+k)^{-s}-t^{-s}$. So $\Delta H(t)=\sum_{k,s}c_{k,s}\big((t+k)^{-s}-t^{-s}\big)=R(t)-\sum_sC_st^{-s}$, and $H$ has its coefficients in $\mathbb K$.

\emph{(2) The decay.} As $t\to\infty$, $(t+k)^{-s}=t^{-s}(1+k/t)^{-s}=t^{-s}+O(t^{-s-1})$, so $R(t)=C_1t^{-1}+O(t^{-2})$, which is $O(t^{-2})$ exactly when $C_1=0$.

\emph{(3) The converse.} A proper rational function $H$ with poles on the nonpositive integer lattice has one partial-fraction presentation $H(t)=\sum_{k,s}d_{k,s}(t+k)^{-s}$, finitely many terms with $k\in\mathbb N$, $s\ge1$. Then $\Delta H(t)=\sum_{k,s}d_{k,s}\big((t+k+1)^{-s}-(t+k)^{-s}\big)$, whose coefficient of $(t+k)^{-s}$ is $d_{k-1,s}-d_{k,s}$ (with $d_{-1,s}=0$), and the column sum over $k$ telescopes to $0$. The column sums $R\mapsto(C_s)_s$ are linear on rational words, presentations being unique, and vanish on such differences. A combination $\sum_sa_st^{-s}$ has the column sums $(a_s)_s$, so it is a difference only if it is $0$: the functions $t^{-1},\dots,t^{-w}$ are independent modulo differences, and by (1) they span the rational words of maximum order $w$ modulo differences. A difference is a rational word with column sums zero, hence $O(t^{-2})$ by (2); and a word that is $O(t^{-2})$ has $C_1=0$, so it is congruent to $\sum_{s\ge2}C_st^{-s}$: modulo differences these words have the basis $t^{-2},\dots,t^{-w}$.

\emph{(4) Cosets.} Translation by one keeps each coset $b+\mathbb Z$, so the steps above apply coset by coset, with $(t+b)^{-s}$ in the place of $t^{-s}$.
\end{proof}

\begin{htheorem}[\hb{thm:telescoper} (Summation by parts with a live input)]
As stated: (iv) without the condition $C_1=0$; (i) to (iii) are those of the reference.
\end{htheorem}

\begin{proof}
(iv) For an integer $M\ge1$ the head is $\Th{0}{R\,\IZ}{N}{N+M}=\sum_{N<m\le N+M}R(m)$, the kernel at index $0$ being $1$. By \hb{thm:columns}, $R(m)=\sum_sC_sm^{-s}+H(m+1)-H(m)$, so
\[
\Th{0}{R\,\IZ}{N}{N+M}=\sum_{s=1}^wC_s\,\Th{s}{\IZ}{N}{N+M}+H(N+M+1)-H(N+1).
\]
As $M\to\infty$: $H(N+M+1)\to0$, $H$ being proper; for $s\ge2$, $\Th{s}{\IZ}{N}{N+M}\to\Tx{s}{\IZ}{N}$; and $\Th{1}{\IZ}{N}{N+M}-\log M\to\Tx{1}{\IZ}{N}$, the regularized value of \hb{def:lattice}. This gives the first display. When $C_1=0$, $R(m)=O(m^{-2})$ by \hb{thm:columns}, the series $\sum_{m>N}R(m)$ converges absolutely, and the left side is $\Tx{0}{R\,\IZ}{N}$, the term $s=1$ being absent. For a density $f$, when the terms converge, $\sum_{m>N}R(m)f(m)=\sum_{s=1}^wC_s\sum_{m>N}m^{-s}f(m)+\sum_{m>N}\Delta H(m)f(m)$, and (ii) gives the last two terms of the second display when the upper boundary vanishes.
\end{proof}

\begin{htheorem}[\hb{thm:bretainput} (A Breta input with a recurrence of finite order)]
As stated: complex $a,b$ with positive real parts, real $q_i<1$ of either sign, and complex $s_i$.
\end{htheorem}

\begin{proof}
\emph{(1) The letters.} For real $q_i<1$ and $x\in[0,1]$, $1-q_ix$ lies between $1$ and $1-q_i$, so in $[\min(1,1-q_i),\max(1,1-q_i)]\subset(0,\infty)$. So $\varrho(x)=\prod_i(1-q_ix)^{-s_i}$ is infinitely differentiable on $[0,1]$, and $|\varrho|=\prod_i(1-q_ix)^{-\re s_i}$ lies between two positive constants.

\emph{(2) The size.} For $m\ge1$, by $x\le e^{-(1-x)}$ on $[0,1]$ and $y=1-x$,
\begin{align*}
|f(m)|&\le\max|\varrho|\int_0^1x^{m+\re a-1}(1-x)^{\re b-1}dx\\&\le\max|\varrho|\int_0^\infty e^{-(m+\re a-1)y}y^{\re b-1}dy=\max|\varrho|\,\Gamma(\re b)\,(m+\re a-1)^{-\re b},
\end{align*}
so $f(m)=O(m^{-\re b})$.

\emph{(3) The recurrence.} Differentiating the logarithm, $\frac d{dx}\big[x^{m+a}(1-x)^b\prod_i(1-q_ix)^{1-s_i}\big]$ is $x^{m+a-1}(1-x)^{b-1}\varrho(x)$ times
\[
(m+a)(1-x)Q(x)-bxQ(x)+x(1-x)\sum_i(s_i-1)q_iQ_i(x)=B_m(x).
\]
The bracket tends to $0$ at both ends, since $\re(m+a)>0$ and $\re b>0$, and $x^{m+a-1}(1-x)^{b-1}\varrho B_m$ is integrable; so integrating over $[0,1]$ gives $\sum_\ell B_\ell(m)f(m+\ell)=0$.

\emph{(4) The extreme coefficients.} $B_0(m)=B_m(0)=(m+a)Q(0)=m+a$. Every $q_i\neq0$, so $Q$ has degree $r$ and leading coefficient $(-1)^r\prod_iq_i$, and $Q_i$ has degree $r-1$ and leading coefficient $(-1)^{r-1}\prod_{i'\neq i}q_{i'}$. The coefficient of $x^{r+1}$ in $B_m$ is therefore
\[
(-1)^{r+1}\prod_iq_i\Big((m+a)+b+\sum_i(1-s_i)\Big)=(-1)^{r+1}\Big(\prod_iq_i\Big)\Big(m+a+b+r-\sum_is_i\Big).
\]
$B_0(m)\neq0$ since $\re a>0$, and $B_{r+1}(m)\neq0$ for every integer $m\ge0$ exactly when $a+b+r-\sum_is_i\notin-\mathbb N$. Then the recurrence solved for $f(m+r+1)$, and for $f(m)$, expresses every shift through $(f(m),\dots,f(m+r))$ with coefficients rational in $m$.

\emph{(5) The Breta.} Let $\re s_i>0$. On the bunch $\{1,2\}$ with the scale $\rho$ and the shape $u=(x,1-x)$, $\omega_{\{1,2\}}$ being $dx$ in the coordinate $x=u_1$, the kernel terms $\rho u_1$ and $\rho u_2$ with unit coefficients sum to $\rho$, and the letter $P_i=\rho\big((1-q_i)u_1+u_2\big)=\rho(1-q_ix)$ has the positive pattern coefficients $1-q_i$ and $1$. At the shape index $(m+a,b)$ and the scale index $1+|\mathbf s|$ the left side of \hb{prop:letters} is
\[
\int_0^1x^{m+a-1}(1-x)^{b-1}\int_0^\infty\rho^{|\mathbf s|}e^{-\rho}\rho^{-|\mathbf s|}\prod_i(1-q_ix)^{-s_i}\,d\rho\,dx=f(m),
\]
the inner integral being $\Gamma(1)=1$; \hb{prop:letters} identifies it with $\Gamma(1+|\mathbf s|)$ times the Breta.
\end{proof}

\begin{htheorem}[\hb{thm:fixedblock} (A fixed set of input-dependent evaluations)]
As stated, for the data of \hb{thm:bretainput}.
\end{htheorem}

\begin{proof}
The steps of the reference, proof of \hb{thm:fixedblock}, hold for these data; they are as follows. Write $F_{p,k}(N)=(\sT_{p;(1)}[x+k]\,f\,\IZ)((N,\infty))=\sum_{m>N}f(m)(m+k)^{-p}$, and, for $0\le x<1$, $S_p(x)=\sum_{m>N}x^m(m+k)^{-p}$, with an integer cut $N\ge0$ and $k\ge0$; $S_0(x)=x^{N+1}/(1-x)$.

\emph{(1) The identity behind (ii).} As in the proof of \hb{thm:bretainput}, $(xC\eta)'=B\eta$, with $xC\eta=x^a(1-x)^b\prod_i(1-q_ix)^{1-s_i}$. Since $xS_p'(x)=\sum_{m>N}m\,x^m(m+k)^{-p}=S_{p-1}(x)-kS_p(x)$,
\[
\frac d{dx}\big[xC\eta\,S_p\big]=\eta\big[(B-kC)S_p+C\,S_{p-1}\big]\qquad(p\ge1).
\]
At $x=0$ the bracket is $O(|x^{a+N+1}|)\to0$, since $\re a>0$. At $x=1$ it is bounded by a constant times $|(1-x)^b|$ for $p\ge2$, where $S_p(1)$ is finite, and times $|(1-x)^b|\,|\log(1-x)|$ for $p=1$, where $S_1(x)=O(|\log(1-x)|)$; it tends to $0$ since $\re b>0$. The functions $|\eta|\,S_p$ and $|\eta|\,|\log(1-x)|$ are integrable on $[0,1)$, so integrating over $[0,1)$ and summing term by term (Tonelli on the moduli) gives, for $p\ge2$,
\[
\sum_{m>N}(m+k)^{-p}\int_0^1x^m(B-kC)\eta\,dx+\sum_{m>N}(m+k)^{1-p}\int_0^1x^mC\eta\,dx=0 .
\]
With $\int_0^1x^{m+\ell}\eta\,dx=f(m+\ell)$ and, for $k\ge r+1\ge\ell$, $\sum_{m>N}f(m+\ell)(m+k)^{-p}=F_{p,k-\ell}(N)-\sum_{v=1}^{\ell}f(N+v)(N+v+k-\ell)^{-p}$, this is (ii) for $p\ge2$. At $p=1$ the term with $S_0$ is $\int_0^1\eta\,C\,S_0\,dx=\int_0^1\eta\,x^{N+1}Q\,dx=\sum_{\ell=0}^rQ_\ell f(N+1+\ell)$, by $C\,S_0=x^{N+1}Q$, which gives (ii) at $p=1$.

\emph{(2) The reduction (i).} In (ii) the coefficient of $F_{p,k}(N)$, the term $\ell=0$, is $B_0-kC_0=a-k$, since $B(0)=aQ(0)=a$ and $C(0)=1$. For $k\ge r+1$ with $k\neq a$, (ii) expresses $F_{p,k}(N)$ through the $F_{p,k-\ell}(N)$, $1\le\ell\le r+1$, the $F_{p-1,k-\ell}(N)$ (none at $p=1$), and the values $f(N+v)$, $1\le v\le r+1$. Taking the indices $p$ in increasing order and, for each, the shifts $k$ in increasing order, every $F_{p,k}(N)$ with $p\le A$ reduces to the $F_{p',k'}(N)$ with $p'\le p$ and $k'\le r$, the $F_{p',a}(N)$ when $a$ is an integer at least $r+1$, and $f(N+1),\dots,f(N+r+1)$.

\emph{(3) The cuts (iii)} are \hb{prop:cutlaw} for the counting input.

\emph{(4) Denominators (iv).} For rational data, let $c_\ast$ clear the coefficients of $B$, $C$ and $Q$, so that $c_\ast(a-k)$, $c_\ast(B_\ell-kC_\ell)$, $c_\ast C_\ell$ and $c_\ast Q_\ell$ are integers for integer $k$. By induction along the order of (2), the coefficients of $F_{p,k}(N)$, $k\ge r+1$, in the reduced form are cleared by $D_{p,k}=d_{N+k}^{\,p}\prod_{r+1\le k'\le k,\ k'\neq a}|c_\ast(a-k')|^{p}$. Multiply (ii) by $c_\ast$ and divide by $c_\ast(a-k)$: the terms $F_{p,k-\ell}(N)$ and $F_{p-1,k-\ell}(N)$ carry integer factors, and their reduced forms are cleared by $D_{p,k-\ell}$ and $D_{p-1,k-\ell}$, which divide $D_{p,k}/|c_\ast(a-k)|$; the terms in $f(N+v)$ carry the denominators $(N+v+k-\ell)^p$ and $(N+v+k-\ell)^{p-1}$, which divide $d_{N+k}^{\,p}$, since $1\le N+v+k-\ell\le N+k$. So $D_{p,k}$ clears $F_{p,k}(N)$, and $d_{N+J}^{A}\prod_{r+1\le k\le J,\ k\neq a}|c_\ast(a-k)|^{A}$, a multiple of every $D_{p,k}$ with $p\le A$ and $k\le J$, clears all shifts through $J$ and indices through $A$.
\end{proof}

\begin{htheorem}[\hb{thm:ra-liverow} (The row of a rational function against an input with a recurrence)]
As stated.
\end{htheorem}

\begin{proof}
All the sums converge absolutely: $m+h\ge N+1+h_\ast>0$ for $m>N$ and $h\ge h_\ast$, and $\sum_{m>N}|f(m)|(1+|m|)^d<\infty$. (i) is linearity, $P(m)=\sum_r([m^r]P)m^r$.

(ii) Expand $p_i$ about $-h$: $p_i(m)=\sum_e\pi_{i,e}(h)(m+h)^e$ with $\pi_{i,e}(h)=p_i^{(e)}(-h)/e!$. Multiply the recurrence at $m$ by $(m+h)^{-s}$, sum over $m>N$ and put $m'=m+i$:
\[
\sum_{i=0}^\tau\sum_e\pi_{i,e}(h)\sum_{m'>N+i}f(m')\,(m'+h-i)^{e-s}=0 .
\]
Moving the cut of each inner sum to $N$ (\hb{prop:cutlaw}) and writing, for $e\ge s$, $(m'+h-i)^{e-s}=\sum_{r\le e-s}\binom{e-s}r(h-i)^{e-s-r}m'^r$,
\begin{multline*}
p_0(-h)\,T_s(h)=\sum_{i=1}^\tau\sum_{v=1}^{i}\sum_e\pi_{i,e}(h)\,f(N+v)\,(N+v+h-i)^{e-s}\\
-\sum_{(i,e)\neq(0,0),\ e<s}\pi_{i,e}(h)\,T_{s-e}(h-i)-\sum_{e\ge s}\sum_i\pi_{i,e}(h)\sum_{r\le e-s}\binom{e-s}r(h-i)^{e-s-r}M_r ,
\end{multline*}
with $\pi_{0,0}(h)=p_0(-h)$. At $h=h_\ast+j$, $\tau\le j\le J$, divide by $p_0(-h)\neq0$: the right side holds $T_{s'}$ with $s'<s$ at the same shift, or with $s'\le s$ at the smaller shifts $h-i\ge h_\ast$, moments $M_r$ with $r\le d-s\le d-1$, and the values $f(N+v)$, $v\le\tau$. Induction on $j$, and for fixed $j$ on $s$, gives (ii); all coefficients are rational.

(iii) Put $Q(s,j)=d_{u(N+h_\ast+j)}^{\,s}\prod_{j'=\tau}^{j}|u^dp_0(-h_\ast-j')|^s$; we show by the same induction that $Q(s,j)$ clears the coefficients of $T_s(h_\ast+j)$, and $Q(s,j)$ divides the integer of (iii). Since $h\in\frac1u\mathbb Z$ and $p_i\in\mathbb Z[m]$, $\pi_{i,e}(h)$ has a denominator dividing $u^{\deg p_i-e}$, and $u^dp_0(-h)$ is an integer. Write the relation as $T_s(h)=\big(u^dp_0(-h)\big)^{-1}\,u^d\,[\cdots]$ and take its terms in turn.
\begin{itemize}
\item $T_{s-e}(h-i)$ with $i\ge1$ or $e\ge1$, $s-e\ge1$: the factor $u^d\pi_{i,e}(h)$ is an integer, and the coefficients of $T_{s-e}(h-i)$ are $1$ if it is a coordinate and are cleared by $Q(s-e,j-i)$ otherwise, by induction; $|u^dp_0(-h)|\,Q(s-e,j-i)$ divides $Q(s,j)$, as $s-e\le s$, $j-i\le j$, and when $i=0$, $e\ge1$ leaves the factor $|u^dp_0(-h)|^e$ free in $Q(s,j)/Q(s-e,j)$.
\item The moments: $u^d\pi_{i,e}(h)(h-i)^{e-s-r}$ has a denominator dividing $u^{-d+\deg p_i-e+e-s-r}$, an integer power of $u$ with exponent at most $0$, so it is an integer.
\item The values $f(N+v)$: for $e<s$, $(N+v+h-i)^{e-s}=u^{s-e}\big(u(N+v+h-i)\big)^{e-s}$, where $u(N+v+h-i)$ is an integer with $1\le u(N+v+h-i)\le u(N+h)$, positive since $h-i\ge h_\ast$ and $N+1+h_\ast>0$; so its power $s-e\le s$ divides $d_{u(N+h_\ast+j)}^{\,s}$, and $u^d\pi_{i,e}(h)u^{s-e}$ is an integer. For $e\ge s$ the term $u^d\pi_{i,e}(h)(N+v+h-i)^{e-s}$ has a denominator dividing $u^{\deg p_i-e+(e-s)-d}$, a power of $u$ with exponent at most $0$, so it is an integer.
\end{itemize}
Every coefficient of $T_s(h_\ast+j)$ is a sum of such products, so $Q(s,j)$ clears it; and for $s\le w$, $j\le J$, $Q(s,j)$ divides $d_{u(N+h_\ast+J)}^{\,w}\prod_{j'=\tau}^J|u^dp_0(-h_\ast-j')|^w$, as $d_n\mid d_{n'}$ for $n\le n'$. Poles in several cosets are reduced coset by coset, each with its own $u$ and $h_\ast$.
\end{proof}

\section{Chapter~\hb{ch:tables}: tables}\label{sec:ch07}

\begin{hproposition}[\hb{prop:mot-dualminors} (Duality of minors)]
As stated.
\end{hproposition}

\begin{proof}
Every basis has $r$ marks, so $|T\cap X|$ is largest exactly when $|T\cap X'|$ is least, that is exactly when the cobasis $\cE\setminus T$ has $|(\cE\setminus T)\cap X'|=|X'|-|T\cap X'|$ largest: $T$ is at the top of $X$ exactly when $\cE\setminus T$ is at the top of $X'$ in $\cB^\ast$ (\hb{def:minors}). By \hb{def:minors} the bases of $\cB^\ast/X'$ are the sets $(\cE\setminus T)\setminus X'=X\setminus(T\cap X)$ for these $T$, the complements in $X$ of the bases $T\cap X$ of $\cB|X$, and the bases of $\cB^\ast|X'$ are the sets $(\cE\setminus T)\cap X'=X'\setminus(T\setminus X)$, the complements in $X'$ of the bases $T\setminus X$ of $\cB/X$. The weight of the cobasis $\cE\setminus T$ is $\theta_T=\theta'_{T\cap X}\theta''_{T\setminus X}$, the factor $\theta''_{T\setminus X}$ belonging to its part $X'\setminus(T\setminus X)$ and the factor $\theta'_{T\cap X}$ to its part $X\setminus(T\cap X)$; these are the weights of the dual tables of \hb{def:table}. The letters follow from \hb{prop:identity}.
\end{proof}

\begin{htheorem}[\hb{thm:tworeadings} (One function, two Foreman polynomials)]
As stated: (iii) along every direction constant on the molecules of the table, with the converse; (i) and (ii) are those of the reference.
\end{htheorem}

\begin{proof}
\emph{The directions.} By \hb{prop:directsum} for the family whose one table is this table, the table is the direct sum of its restrictions to its molecules $M$: every basis $T$ is the union of one basis of each restriction $\cB|M$, and a basis of $\cB|M$ has $r(M)$ marks (\hb{def:minors}). So $|T\cap M|=r(M)$, and for $v$ constant on each $M$, with value $v_M$,
\[
\ip v{\chi_T}=\sum_Mv_M\big(|T\cap M|-|M\setminus T|\big)=\sum_Mv_M\big(2r(M)-|M|\big)
\]
for every basis. By \hb{def:Lambda}, $\Lambda(z+v)=\sum_T\theta_Te^{\ip z{\chi_T}}e^{\ip v{\chi_T}}=\exp\big(\sum_Mv_M(2r(M)-|M|)\big)\Lambda(z)$. At $v=\lambda\one$ the exponent is $\lambda(2r-n)$, the $r(M)$ adding up to $r$, since a basis meets each molecule in $r(M)$ marks, and the $|M|$ to $n$.

\emph{The converse.} Let $v\in\mathbb R^{\cE}$ and real $\lambda\neq0$ be such that $\Lambda(z+\lambda v)=c\,\Lambda(z)$ for every $z\in\mathbb R^{\cE}$, with $c$ independent of $z$. Then $\sum_T\theta_T\big(e^{\lambda\ip v{\chi_T}}-c\big)e^{\ip z{\chi_T}}=0$ for every $z$. Distinct bases have distinct signed indicators, and exponentials $e^{\ip z\chi}$ with finitely many distinct vectors $\chi$ are linearly independent functions of $z$: on a line $z=\sigma y$ with the $\ip y\chi$ distinct, as $\sigma\to\infty$ the term with the largest $\ip y\chi$ dominates, so its coefficient vanishes, and the others follow in turn. So $\theta_T(e^{\lambda\ip v{\chi_T}}-c)=0$, and since $\theta_T>0$, $e^{\lambda\ip v{\chi_T}}=c$ for every basis. Hence $\ip v{\chi_T}=2\ip v{\one_T}-\ip v\one$ is constant over the bases, and so is $\ip v{\one_T}$. By \hb{lem:circuit}, $v_a=v_b$ for any two marks linked by a circuit of the table; the molecules of the table alone are the smallest nonempty sets of marks containing, with each mark, every mark so linked to it (\hb{def:molecule}), so $v$ is constant on each of them.
\end{proof}

\begin{htheorem}[\hb{thm:innerMGF} (The inner moment generating function)]
As stated: the even statistic at every gauge, with the basis law and the cobasis law there.
\end{htheorem}

\begin{proof}
For real $z$ and $w$, $e^{E(z+w)}=\Lambda(z+w)\,\Lambda(-z-w)$. By the first display of the theorem, $\Lambda(z+w)/\Lambda(z)=\oE_{T\sim\mathbb P_\alpha}[e^{\ip w{\chi_T}}]$, and the same display at the gauge $\alpha^{-1}=e^{-2z}$, whose basis law is $\mathbb P_{\alpha^{-1}}=\mathbb P^\ast_\alpha$, gives $\Lambda(-z-w)/\Lambda(-z)=\oE_{T'\sim\mathbb P^\ast_\alpha}[e^{-\ip w{\chi_{T'}}}]$. Multiplying, and using the independence of $T$ and $T'$ under $\mathbb P_\alpha\otimes\mathbb P^\ast_\alpha$,
\[
\frac{e^{E(z+w)}}{e^{E(z)}}=\oE_{(T,T')\sim\mathbb P_\alpha\otimes\mathbb P^\ast_\alpha}\big[e^{\ip w{\chi_T-\chi_{T'}}}\big],\qquad \chi_T-\chi_{T'}=2(\one_T-\one_{T'}) ,
\]
which is the display; its logarithm says that $w\mapsto E(z+w)-E(z)$ is the cumulant generating function of $\chi_T-\chi_{T'}$. At $z=0$, $e^{E(0)}=\Lambda(0)^2=K(\one)^2$ by (i) of \hb{thm:tworeadings}, and $\mathbb P^\ast_{\one}=\mathbb P_{\one}$; with $w$ written $z$ this is the display of the reference.
\end{proof}

\begin{sremark}
The display at a general gauge, and (iii) of \hb{thm:tworeadings}, were checked in exact arithmetic on the table whose bases are the unions of a $2$-subset of $\{0,1,2\}$ and a $1$-subset of $\{3,4\}$, with arbitrary positive weights (\texttt{sweep/s0107/tables.py}).
\end{sremark}

\begin{hcorollary}[\hb{cor:interference} (Moments and cumulants)]
As stated: the joint cumulants of the table at every order; (i) and the cumulants of the components are those of the reference.
\end{hcorollary}

\begin{proof}
For a function $Y$ of the index $(k,l)$ write $\oE_w[Y]=\sum_{k,l}w_{kl}Y(k,l)$; $\oE_w[1]=1$ by \hb{thm:interference}.

\emph{(1) Generating functions.} By \hb{thm:interference} at $x=e^v$, $v\in\mathbb C^B$,
\[
\oE_{T\sim\mathbb P_\alpha}\Big[\exp\Big(\sum_{e\in B}v_eX_e\Big)\Big]=\sum_{k,l}w_{kl}\,M_{kl}(v),\qquad M_{kl}(v)=\det\big(I+(D_{e^v}-I)\Pi^{(kl)}_{BB}\big),
\]
with $M_{kl}(0)=1$. So $C_{kl}(v)=\log M_{kl}(v)$ is a power series in $v$ without constant term, and for distinct marks the cumulant $\kappa^{(kl)}(e_1,\dots,e_p)$ of the component is the coefficient of $v_{e_1}\cdots v_{e_p}$ in $C_{kl}$; the reference computes it by the cyclic formula. The joint cumulant $\kappa_\alpha(e_1,\dots,e_p)$ is the coefficient of $v_{e_1}\cdots v_{e_p}$ in $\log\sum_{k,l}w_{kl}e^{C_{kl}(v)}$, a power series since $\sum_{k,l}w_{kl}=1$.

\emph{(2) Cumulants over the index.} Let $(Y_i)_{i\in I}$ be finitely many functions of the index, $\lambda=(\lambda_i)_{i\in I}$ formal variables, and $\mu(\lambda)=\oE_w[\exp(\sum_i\lambda_iY_i)]$, a power series with constant term $1$. For $D\subseteq I$ write $\lambda^D=\prod_{i\in D}\lambda_i$; the coefficient of $\lambda^D$ in $\mu$ is $\oE_w[\prod_{i\in D}Y_i]$. In $\log\mu=\sum_{j\ge1}\frac{(-1)^{j-1}}j(\mu-1)^j$ the coefficient of $\lambda^I$ comes from the products of $j$ monomials of $\mu-1$ without repeated variables and with disjoint supports covering $I$; each partition $\mathfrak q$ of $I$ into $j$ blocks arises from $j!$ such ordered products. So
\[
[\lambda^I]\log\mu=\sum_{j\ge1}\frac{(-1)^{j-1}}j\,j!\sum_{|\mathfrak q|=j}\ \prod_{D\in\mathfrak q}\oE_w\Big[\prod_{i\in D}Y_i\Big]=\kappa_w\big((Y_i)_{i\in I}\big),
\]
and, setting the variables outside a subset $D$ to $0$, $[\lambda^D]\log\mu=\kappa_w((Y_i)_{i\in D})$.

\emph{(3) Substitution.} Fix distinct marks $e_1,\dots,e_p$ of $B$; for nonempty $C\subseteq\{1,\dots,p\}$ put $v^C=\prod_{c\in C}v_{e_c}$ and $Y_C(k,l)=\kappa^{(kl)}(e_C)$, the coefficient of $v^C$ in $C_{kl}$. Setting the variables $v_e$ with $e\notin\{e_1,\dots,e_p\}$ to $0$ and dropping every monomial divisible by some $v_{e_c}^2$ is a ring homomorphism onto the ring $\mathcal R$ of polynomials in $v_{e_1},\dots,v_{e_p}$ with $v_{e_c}^2=0$; it keeps the coefficient of $v_{e_1}\cdots v_{e_p}$ and commutes with $\exp$ and $\log$ of series without constant term. In $\mathcal R$, $C_{kl}$ becomes $\sum_{C\neq\varnothing}Y_C(k,l)\,v^C$, the image of $\sum_C\lambda_CY_C(k,l)$ under $\lambda_C\mapsto v^C$. So the coefficient of $v_{e_1}\cdots v_{e_p}$ in $\log\sum_{k,l}w_{kl}e^{C_{kl}}$ is that of the image of $\log\mu(\lambda)$, with $I$ the set of nonempty $C$. A monomial $\prod_C\lambda_C^{n_C}$ maps to $\prod_C(v^C)^{n_C}$, which is $v_{e_1}\cdots v_{e_p}$ exactly when the $C$ with $n_C>0$ have $n_C=1$ and form a partition $\mathfrak p$ of $\{1,\dots,p\}$, and is $0$ or another monomial otherwise. By (2),
\[
\kappa_\alpha(e_1,\dots,e_p)=\sum_{\mathfrak p}\,[\lambda^{\mathfrak p}]\log\mu=\sum_{\mathfrak p}\kappa_w\Big(\big(\kappa^{(kl)}(e_C)\big)_{C\in\mathfrak p}\Big).
\]
At $p=1$ the single partition gives $\kappa_\alpha(e)=\oE_w[\Pi^{(kl)}_{ee}]$, and $\kappa_\alpha(e)=\mathbb P_\alpha(e\in T)=R_e(\alpha)$. Every step is an identity of power series in which the $w_{kl}$ enter linearly inside $\mu$, with $\sum_{k,l}w_{kl}=1$, so it holds for complex mixing coefficients.
\end{proof}

\begin{sremark}
At $p=2$ and $p=3$ the formula gives $\kappa_\alpha(e,f)=\oE_w[\kappa^{(kl)}(e,f)]+\kappa_w(\Pi^{(kl)}_{ee},\Pi^{(kl)}_{ff})$ and $\kappa_\alpha(e,f,g)=\oE_w[\kappa^{(kl)}(e,f,g)]+\sum_3\kappa_w(\Pi^{(kl)}_{ee},\kappa^{(kl)}(f,g))+\kappa_w(\Pi^{(kl)}_{ee},\Pi^{(kl)}_{ff},\Pi^{(kl)}_{gg})$, the sum over the three choices of the mark that stands alone: the covariances and the third cumulant over the index of the reference. The formula was checked in exact arithmetic against the generating function of a mixture of three components with complex mixing coefficients for $p=2,3,4,5$ (\texttt{sweep/s0107/cumulants2.py}).
\end{sremark}

\section{Chapter~\hb{ch:families}: families}\label{sec:ch08}

\begin{hproposition}[\hb{prop:latent} (The interest over latent vectors)]
As stated: every table determinantal over $\mathbb R$ or over $\mathbb C$, each letter opened with as many latent vectors as its own exponent asks.
\end{hproposition}

\begin{proof}
The masses are real and positive. For a mass $m_e\le0$ the interest diverges by \hb{thm:exist}(i), the wall $\kappa_{\{e\}}=m_e-\sum_j\big(s_jd_j(\{e\})+t_jd^\ast_j(\{e\})\big)$ being at most $m_e$, and so does the integral over $\alpha_e$ in step~(4).

Fix a table $j$, write $\mathbb F=\mathbb F_j$, $\beta=\beta_j$, $r=r_j$, and let $A\in\mathbb F^{r\times n}$ be a slice of it: $\det A_T=\boldsymbol\omega_T$, so that $|\det A_T|^2=\theta_{j,T}$ for $T\in\cB_j$ and $\det A_T=0$ for every other set $T$ of $r$ marks (\hb{def:multivector}). For a matrix $X$ with $n$ columns write $X^{\mathsf H}$ for its conjugate transpose, and $D_\alpha=\operatorname{diag}(\alpha)$. Expanding by the maximal minors of $X$,
\[
\det\big(XD_\alpha X^{\mathsf H}\big)=\sum_W|\det X_W|^2\,\alpha^W ,
\]
the sum over the sets $W$ of as many marks as $X$ has rows; this is the expansion behind \hb{prop:tensorletters}, and with $X=A$ it gives $K_j(\alpha)=\det(AD_\alpha A^{\mathsf H})$. So $K_j(\one)=\det(AA^{\mathsf H})>0$, $A$ has rank $r$, and its row space $V_j=\{A^{\mathsf H}y:\ y\in\mathbb F^r\}$, the space onto which $\Pi^{(11)}$ of \hb{thm:interference} projects, has dimension $r$ over $\mathbb F$.

\emph{(1) The basis letter.} The matrix $\hat A=(AA^{\mathsf H})^{-1/2}A$ has orthonormal rows and $\{\hat A^{\mathsf H}y\}=V_j$, and $\det\hat A_T=\det(AA^{\mathsf H})^{-1/2}\det A_T$. So $\theta_{j,T}=K_j(\one)\,|\det\hat A_T|^2$ and $K_j(\alpha)=K_j(\one)\det(\hat AD_\alpha\hat A^{\mathsf H})$.

\emph{(2) The cobasis letter.} Choose $\hat C\in\mathbb F^{(n-r)\times n}$ with orthonormal rows and $\{\hat C^{\mathsf H}y\}=V_j^\perp$, the orthogonal complement for the form $\sum_eu_e\overline{v_e}$. Then $\hat A\hat C^{\mathsf H}=0$, and the square matrix $U$ whose rows are those of $\hat A$ followed by those of $\hat C$ is unitary. For an invertible matrix, $\det U[I,J]=\pm\det U\cdot\det U^{-1}[J^c,I^c]$ for every set $I$ of rows and set $J$ of columns of equal size; with $U^{-1}=U^{\mathsf H}$ and $|\det U|=1$ this gives $|\det U[I,J]|=|\det U[I^c,J^c]|$. With $I$ the rows of $\hat A$ and $J=T$, $|\det\hat C_{\cE\setminus T}|=|\det\hat A_T|$ for every set $T$ of $r$ marks. The expansion above, for $X=\hat C$ and the sets $W=\cE\setminus T$, gives
\[
\det\big(\hat CD_\alpha\hat C^{\mathsf H}\big)=\sum_T|\det\hat A_T|^2\alpha^{\cE\setminus T}=\frac1{K_j(\one)}\sum_{T\in\cB_j}\theta_{j,T}\,\alpha^{\cE\setminus T}=\frac{\Psi_j(\alpha)}{K_j(\one)} .
\]

\emph{(3) The Gaussian opening.} For a Hermitian positive definite matrix $G$ of size $d$ over $\mathbb F$, a unitary change of variables, which keeps the Lebesgue measure of $\mathbb F^d\cong\mathbb R^{\beta d}$, makes $G$ diagonal, and $\int_{\mathbb F}e^{-\lambda|z|^2}dz=(\pi/\lambda)^{\beta/2}$ for $\lambda>0$; so $\int_{\mathbb F^d}e^{-z^{\mathsf H}Gz}dz=\pi^{\beta d/2}\det(G)^{-\beta/2}$. For an integer $k\ge0$ and $s=\beta k/2$ it follows that
\[
\det(G)^{-s}=\pi^{-ds}\int_{(\mathbb F^d)^k}\exp\Big(-\sum_{i=1}^kz_i^{\mathsf H}Gz_i\Big)\,dz_1\cdots dz_k .
\]
Take $G=\hat AD_\alpha\hat A^{\mathsf H}$, $d=r$ and $k=2s_j/\beta$. The map $z\mapsto x=\hat A^{\mathsf H}z$ is an $\mathbb F$-linear isometry of $\mathbb F^r$ onto $V_j$, so it carries the Lebesgue measure of $\mathbb F^r$ to that of $V_j$, and $z^{\mathsf H}\hat AD_\alpha\hat A^{\mathsf H}z=\sum_e\alpha_e|x_e|^2$. With step~(1),
\[
K_j(\alpha)^{-s_j}=K_j(\one)^{-s_j}\pi^{-r_js_j}\int_{V_j^{2s_j/\beta}}\exp\Big(-\sum_{i,e}\alpha_e|(x_{j,i})_e|^2\Big)\prod_idx_{j,i},
\]
and in the same way, with step~(2), $G=\hat CD_\alpha\hat C^{\mathsf H}$, $d=n-r_j=r^\ast_j$ and $k=2t_j/\beta$,
\[
\Psi_j(\alpha)^{-t_j}=K_j(\one)^{-t_j}\pi^{-r^\ast_jt_j}\int_{(V_j^\perp)^{2t_j/\beta}}\exp\Big(-\sum_{i,e}\alpha_e|(x^\ast_{j,i})_e|^2\Big)\prod_idx^\ast_{j,i}.
\]
An exponent $0$ asks for no vectors, and the formula then reads $1=1$.

\emph{(4) The scales.} Multiplying over $j$, with $c_F=\prod_jK_j(\one)^{s_j+t_j}\pi^{r_js_j+r^\ast_jt_j}$,
\[
W(\alpha)=\frac1{c_F}\int\exp\Big(-\sum_e\alpha_e\mathfrak q_e\Big)\prod_{j,i}dx_{j,i}\,dx^\ast_{j,i}.
\]
Insert this into $J_F(\mvec;\mathbf s,\mathbf t)=\int\alpha^{\mvec-\one}e^{-|\alpha|}W(\alpha)\,d\alpha$. Every integrand is nonnegative, so the integrals may be taken in either order, and
\[
\int_0^\infty\alpha_e^{m_e-1}e^{-\alpha_e(1+\mathfrak q_e)}\,d\alpha_e=\Gamma(m_e)\,(1+\mathfrak q_e)^{-m_e}.
\]
Integrating the scales first gives the right side of the statement, and integrating the latent vectors first gives $J_F(\mvec;\mathbf s,\mathbf t)$: both are the integral of one nonnegative function, finite or infinite together.

\emph{(5) The Breta law.} Where the two sides are finite, the function
\[
f(\alpha,x,x^\ast)=\alpha^{\mvec-\one}\exp\Big(-\sum_e\alpha_e(1+\mathfrak q_e)\Big)
\]
is integrable, with integral $c_FJ_F(\mvec;\mathbf s,\mathbf t)$ by step~(4). Integrating out the latent vectors, its marginal in $\alpha$ is $c_F\,\alpha^{\mvec-\one}e^{-|\alpha|}W(\alpha)$, proportional to the density $\alpha^{\mvec-\one}e^{-|\alpha|}W(\alpha)/J_F(\mvec;\mathbf s,\mathbf t)$ of $\mathrm{Breta}(F;\mvec;\mathbf s,\mathbf t)$ (\hb{def:shapelaw}; the interest is finite exactly when every wall is positive, \hb{thm:exist}(i)). Given the latent vectors, $f$ is the product over the marks of $\alpha_e^{m_e-1}e^{-\alpha_e(1+\mathfrak q_e)}$: the scales are independent, and $\alpha_e(1+\mathfrak q_e)$ has the density $u^{m_e-1}e^{-u}/\Gamma(m_e)$ on $(0,\infty)$, the law $\operatorname{Gamma}_{m_e}$.

\emph{(6) Equal exponents over $\mathbb R$.} If $\mathbb F_j=\mathbb R$ and $s_j=t_j$, the vectors of the two letters come in pairs $(x_{j,i},x^\ast_{j,i})$, $i\le2s_j$. The map $V_j\times V_j^\perp\to\mathbb R^{\cE}$, $(x,x^\ast)\mapsto x+x^\ast$, is a linear isometry, so it carries the product of the Lebesgue measures to that of $\mathbb R^{\cE}$, and for $u=x+x^\ast$, $x=\operatorname{pr}_ju$ and $x^\ast=(I-\operatorname{pr}_j)u$. The constant of table $j$ is then $K_j(\one)^{2s_j}\pi^{ns_j}$.
\end{proof}

\begin{sremark}
The reference, proof of \hb{prop:latent}, treats $s_j=t_j=k_j/2$ over $\mathbb R$ and opens the two letters of table $j$ with one set of $k_j$ vectors $u=\hat A^{\mathsf T}y+\hat C^{\mathsf T}y^\ast$; that is step~(6). In the proof above each letter is opened on its own space, with as many vectors as its exponent asks, and over $\mathbb C$ by the complex Gaussian integral, $\beta=2$.
\end{sremark}

\begin{slemma}\label{lem:families-separators}
For a family on $\cE$ and a nonempty set $S\subseteq\cE$, $c_j(S)=0$ for every table $j$ of the family exactly when $S$ is a union of molecules.
\end{slemma}

\begin{proof}
The molecules are the smallest nonempty sets of marks that contain, with each mark, every mark linked to it (\hb{def:molecule}); so a set of marks with that property is a union of molecules, and a union of molecules has it.

Let $S$ be a union of molecules. By \hb{prop:directsum} every table is the direct sum of its restrictions to the molecules: its bases are the unions of one basis of each restriction, and a basis of the restriction to a molecule $M$ has $r_j(M)$ marks (\hb{def:minors}). So $|T\cap S|=\sum_{M\subseteq S}r_j(M)$ on every basis $T$ of table $j$, and by \hb{thm:window}, $d_j(S)=r_j(S)$, that is $c_j(S)=0$.

Let $S$ not be a union of molecules. Then some $e\in S$ is linked to some $f\notin S$: $e$ and $f$ lie in a common circuit of some table $j$. By \hb{lem:linkexchange} some basis $T$ of table $j$ has $e\in T$, $f\notin T$ and $T'=(T\setminus\{e\})\cup\{f\}$ a basis. Then $|T'\cap S|=|T\cap S|-1$, so $|T\cap S|$ takes two values on the bases of table $j$, and $c_j(S)\ge1$ by \hb{thm:window}.
\end{proof}

\begin{sremark}
Since the weight law charges every basis, Lemma~\ref{lem:families-separators} is the first sentence of \hb{prop:firstnumber} read on the bases: $\operatorname{Var}_{T\sim\mathbb P_{j,\one}}|T\cap U|=0$ exactly when $|T\cap U|$ is constant on $\cB_j$, that is when $c_j(U)=0$.
\end{sremark}

\begin{hproposition}[\hb{prop:compoundwall} (The wall of a compound)]
As stated: the identity holds for every nonempty proper $S$, and only the tables with $\sigma_j\neq0$ enter it.
\end{hproposition}

\begin{proof}
\hb{thm:wallscremona} gives, as identities of functions of the data, $\kappa_S+\kappa^{\mathrm{top}}_S=2\sum_j\sigma_jc_j(S)$ and $\kappa_{\cE\setminus S}=\kappa_{\cE}+\kappa^{\mathrm{top}}_S$. Adding the two,
\[
\kappa_S+\kappa_{\cE\setminus S}=\kappa_{\cE}+2\sum_j\sigma_jc_j(S).
\]
If $c_j(S)=0$ in every table with $\sigma_j\neq0$, the sum vanishes and $\kappa_S+\kappa_{\cE\setminus S}=\kappa_{\cE}$; this holds in particular when $S$ is a union of molecules, where $c_j(S)=0$ in every table (Lemma~\ref{lem:families-separators}). On the balance wall $\kappa_{\cE}=0$, so $\re\kappa_S+\re\kappa_{\cE\setminus S}=0$ and one of the two walls has real part at most zero. By \hb{thm:projective} the period converges only where every nonempty proper wall has positive real part, so it diverges.
\end{proof}

\begin{sremark}
The reference, \hb{prop:compoundwall}, takes $c_j(S)=0$ in every table. With complex exponents a table with $s_j=-t_j\neq0$ is active and has $\sigma_j=0$; it does not enter the identity.
\end{sremark}

\begin{hdefinition}[\hb{def:native} (Target; native gauge)]
As stated: the target for all data with $m_{\cE}\neq0$ and on the balance wall, and native points of $(\mathbb C^\times)^{\cE}$ off the zeros of the letters, the native gauges being the native points in $\mathbb R^{\cE}_{>0}$.
\end{hdefinition}

\begin{proof}
The definition is consistent.

\emph{(1) Scale.} $K_j$ and $\Psi_j$ are homogeneous polynomials of degrees $r_j$ and $r^\ast_j=n-r_j$, every basis having $r_j$ marks and every complement $n-r_j$ (\hb{def:table}). So for $c\in\mathbb C^\times$, $K_j(c\alpha)=c^{r_j}K_j(\alpha)$ and $(\partial_eK_j)(c\alpha)=c^{r_j-1}(\partial_eK_j)(\alpha)$, whence $R_{j,e}(c\alpha)=R_{j,e}(\alpha)$ wherever $K_j(\alpha)\neq0$, and likewise $R^\ast_{j,e}(c\alpha)=R^\ast_{j,e}(\alpha)$ wherever $\Psi_j(\alpha)\neq0$. The condition $K_j(\alpha)\Psi_j(\alpha)\neq0$ is unchanged by scaling and the target does not involve the point, so the usage, and with it nativity, is a property of the point up to scale.

\emph{(2) The two clauses of the target.} On the balance wall $m_{\cE}=\sum_j(s_jr_j+t_jr^\ast_j)$ (\hb{def:family}); where this is not zero the formula gives $\btau=\mvec$, the value assigned on the balance wall. So the two clauses agree where both apply, and the second also covers the balance wall with $m_{\cE}=0$.

\emph{(3) Positive points.} At $\alpha\in\mathbb R^{\cE}_{>0}$ the letters are positive, the weights being positive, so every positive point lies off their zeros, and $R_{j,e}(\alpha)=\mathbb P_{j,\alpha}(e\in T)$ and $R^\ast_{j,e}(\alpha)=\mathbb P^\ast_{j,\alpha}(e\notin T)$ (\hb{thm:innerMGF}).
\end{proof}

\begin{sremark}
The direction of a ray lies on the balance wall (\hb{def:ray}), so its target is its masses $\mvec_1$, and the native gauges at complex points of the direction are the points $\alpha_c$ of the letter complement with $c^{(1)}(\alpha_c)=0$: the native gauges of the ray (\hb{def:raygauge}).
\end{sremark}

\begin{htheorem}[\hb{thm:saddle} (The native gauge is the saddle)]
As stated, with the native gauges of Definition~\hb{def:native}.
\end{htheorem}

\begin{proof}
At a gauge $\alpha=e^y$ the usage of Definition~\hb{def:native} is $\nabla\Phi(y)$: by the chain rule $\partial_{y_e}\log K_j(e^y)=\alpha_e\partial_eK_j(\alpha)/K_j(\alpha)=R_{j,e}(\alpha)$ and $\partial_{y_e}\log\Psi_j(e^y)=R^\ast_{j,e}(\alpha)$, and $\Phi$ is $\sum_j\big(s_j\log K_j+t_j\log\Psi_j\big)$ read at $e^y$ (\hb{def:potential}). For real data with $s_j,t_j\ge0$, not all zero, and $\mvec>0$, $m_{\cE}>0$, so the target of that definition is the one given by its formula, and its native gauges are the gauges of (i). With this reading the statement is the reference, proof of \hb{thm:saddle}.
\end{proof}

\begin{hproposition}[\hb{prop:compoundgauge} (The native gauges of a compound)]
As stated: the weights need not factor except in the sentence on products, and the molecules are those of the active family.
\end{hproposition}

\begin{proof}
Write $g=\sum_j(s_jd_j+t_jd^\ast_j)$. An inactive table enters neither the potential, nor the target, nor $g$, nor the Newton sum $\mathcal N_F$, its terms carrying the factor zero; every statement below concerns the active family, whose molecules are $\cE_1,\dots,\cE_q$.

\emph{(1) The criterion of the saddle.} By \hb{thm:saddle}(ii) a native gauge exists exactly when $\btau$ lies in the relative interior of $\mathcal N_F=\sum_j(s_jP_j+t_jP^\ast_j)$, and it is then unique modulo the space $V$ of directions along which $\Phi$ is affine. By \hb{thm:saddle}(iii), whose first sentence uses no connectedness, $\mathcal N_F=\{x:\ x_{\cE}=g(\cE),\ x_S\ge g(S)\text{ for all }S\}$.

\emph{(2) The extent of $x_S$ on $\mathcal N_F$.} The least value of a linear function on a Minkowski sum with nonnegative coefficients is the same sum of its least values on the summands. The least value of $x_S$ on $P_j=\operatorname{conv}\{\one_T\}$ is $\min_T|T\cap S|=d_j(S)$, and on $P^\ast_j=\operatorname{conv}\{\one_{\cE\setminus T}\}$ it is $\min_T|S\setminus T|=d^\ast_j(S)$ (\hb{thm:window}). So the least value of $x_S$ on $\mathcal N_F$ is $g(S)$, and that of $x_{\cE\setminus S}$ is $g(\cE\setminus S)$; since $x_{\cE}=g(\cE)$ on $\mathcal N_F$ and $\mathcal N_F$ is convex, $x_S$ takes on $\mathcal N_F$ exactly the values of $[g(S),\,g(\cE)-g(\cE\setminus S)]$. At any masses $\kappa_S+\kappa_{\cE\setminus S}-\kappa_{\cE}=g(\cE)-g(S)-g(\cE\setminus S)$, the masses cancelling, so by \hb{prop:compoundwall} the length of this interval is $2\sum_j\sigma_jc_j(S)=\sum_j(s_j+t_j)c_j(S)$.

\emph{(3) The relative interior.} An active table has $s_j+t_j>0$. So the interval of step~(2) is a point exactly when $c_j(S)=0$ for every active table, that is, by Lemma~\ref{lem:families-separators} applied to the active family, when $S$ is a union of the $\cE_i$. The relative interior of a polyhedron given by linear equalities and inequalities consists of its points at which every inequality that is not an equality on the whole polyhedron holds strictly. So $\btau$ lies in the relative interior of $\mathcal N_F$ exactly when $\tau_S=g(S)$ for every union $S$ of the $\cE_i$ and $\tau_S>g(S)$ for every other nonempty $S$.

\emph{(4) Molecule by molecule.} By \hb{prop:directsum} for the active family, every active table is the direct sum of its restrictions to the $\cE_i$, its bases being the unions of one basis of each restriction. Taking the least of $|T\cap S|$, and of $|S\setminus T|$, restriction by restriction, $d_j(S)=\sum_id_j(S\cap\cE_i)$ and $d^\ast_j(S)=\sum_id^\ast_j(S\cap\cE_i)$ for every active $j$; so $g(S)=\sum_ig(S\cap\cE_i)$ and $\tau_S-g(S)=\sum_i\big(\tau_{S\cap\cE_i}-g(S\cap\cE_i)\big)$. Hence the conditions of step~(3) hold exactly when $\tau_{\cE_i}=g(\cE_i)$ for every $i$ and $\tau_S>g(S)$ for every $S$ with $\varnothing\neq S\subsetneq\cE_i$ for some $i$: given these, a union of the $\cE_i$ has $\tau_S-g(S)=0$, and any other nonempty $S$ meets some $\cE_i$ in a nonempty proper subset, whose term is positive while the others are nonnegative; conversely each $\cE_i$ is a union of the $\cE_i$, and a nonempty proper subset of one $\cE_i$ is not. Since $c_j(\cE_i)=0$ for every active $j$, $d_j(\cE_i)=r_j(\cE_i)$, so $r_j(\cE\setminus\cE_i)=r_j-r_j(\cE_i)$ (\hb{def:skeleton}, $d_j(S)=r_j-r_j(\cE\setminus S)$) and $d^\ast_j(\cE_i)=|\cE_i|-r_j(\cE_i)=r^\ast_j(\cE_i)$. So $g(\cE_i)=\sum_j\big(s_jr_j(\cE_i)+t_jr^\ast_j(\cE_i)\big)$, and the conditions are the balance of every $\cE_i$ at the target and the positivity of its proper walls there. This proves the criterion.

\emph{(5) The free scales.} By \hb{def:weightlaw} the second derivative of $\log K_j(e^y)$ along a direction $v$ is the variance of $\ip v{\one_T}$ under $\mathbb P_{j,e^y}$, and that of $\log\Psi_j(e^y)=\ip\one y+\log K_j(e^{-y})$ (\hb{prop:identity}) is the variance of $\ip v{\one_T}$ under $\mathbb P_{j,e^{-y}}$. Both laws charge every basis, so each variance vanishes exactly when $\ip v{\one_T}$ is constant over $\cB_j$. The second derivative of $\Phi$ along $v$ is the sum of these variances with the coefficients $s_j,t_j\ge0$, and every active table has a positive coefficient; so $\Phi$ is affine along $v$ exactly when $\ip v{\one_T}$ is constant over the bases of every active table. By \hb{lem:circuit} for the active family such a $v$ is constant on every $\cE_i$. Conversely, if $v$ is constant on every $\cE_i$, with value $v_i$, then $\ip v{\one_T}=\sum_iv_i|T\cap\cE_i|=\sum_iv_ir_j(\cE_i)$ on every basis of an active table $j$, by step~(4). So $V=\operatorname{span}\{\one_{\cE_1},\dots,\one_{\cE_q}\}$, and the native gauges are the gauges $e^{v}\alpha^\star$, $v\in V$: the gauges $\beta$ with $\beta/\alpha^\star$ constant on every $\cE_i$.

\emph{(6) Factoring weights.} When the weights factor across the $\cE_i$, the letters of every active table factor,
\[
K_j(\alpha)=\prod_iK_{j,i}(\alpha_{\cE_i}),\qquad \Psi_j(\alpha)=\prod_i\Psi_{j,i}(\alpha_{\cE_i}),
\]
with the Foreman polynomials of the restrictions (\hb{prop:directsum}). So $\Phi(y)=\sum_i\Phi_i(y_{\cE_i})$, with $\Phi_i$ the potential of the restriction of the family to $\cE_i$, and the condition $\nabla\Phi=\btau$ splits into the conditions $\nabla\Phi_i=\btau|_{\cE_i}$, each the native condition of the restriction to $\cE_i$ at the masses $\btau|_{\cE_i}$, which lie on its balance wall by step~(4), so that its target is $\btau|_{\cE_i}$ (\hb{def:native}). The native gauges are the gauges whose restriction to every $\cE_i$ is native there: the products of native gauges of the molecules.

\emph{(7) Units.} A loop is a circuit by itself, so it lies in no other circuit; and a coloop $e$ lies in no circuit, since a circuit $C$ through $e$ would lie in every basis containing the independent set $C\setminus\{e\}$, such a basis containing $e$, while a circuit lies in no basis. So a mark $e$ that is a loop or a coloop in every table is linked to no other mark, and $\{e\}$ is a unit. Every basis of table $j$ contains $e$, or none does, so $K_j(\alpha)=\alpha_e^{\epsilon_j}\tilde K_j(\alpha_{\cE\setminus e})$ and $\Psi_j(\alpha)=\alpha_e^{1-\epsilon_j}\tilde\Psi_j(\alpha_{\cE\setminus e})$ with $\epsilon_j\in\{0,1\}$: the letters of $\{e\}$ are monomials.
\end{proof}

\begin{sremark}
The criterion is stated for the molecules of the active family because linking (\hb{def:molecule}) runs through the circuits of every table, while an inactive table enters neither the potential, nor the target, nor the walls. So a molecule of the family may join several molecules of the active family, and then a proper wall of that molecule vanishes at the target although a native gauge exists. For instance, on $\cE=\{1,2,3,4\}$ let the list with the bases $\{1,3\},\{1,4\},\{2,3\},\{2,4\}$ carry $(s,t)=(1,0)$ and the list $\binom{\cE}{1}$ carry $(0,0)$, with the masses $\frac12\one$ on the balance wall. Through the second table $\cE$ is one molecule of the family, while the active family has the molecules $\{1,2\}$ and $\{3,4\}$. Every mark lies in half of the bases of the first table, so the gauge $\one$ is native, and the wall at $\{1,2\}$ is $\tau_{\{1,2\}}-d_1(\{1,2\})=1-1=0$.
\end{sremark}

\begin{sremark}
At the centre of \hb{prop:centrefamily} the gauge $\one$ is native, and step~(5) shows that the native gauges are the gauges constant on every molecule of the active family: this is the uniqueness clause of \hb{prop:centrefamily}(i).
\end{sremark}

\section{Chapter~\hb{ch:still}: realizations and the still}\label{sec:ch09}

\begin{hproposition}[\hb{prop:wallsrealization} (A realization sits on a wall only along its molecules)]
As stated, for every family.
\end{hproposition}

\begin{proof}
\emph{(1) The walls as expectations.} By \hb{thm:innerMGF}, $R_{j,e}(\alpha)=\mathbb P_{j,\alpha}(e\in T)$ and $R^\ast_{j,e}(\alpha)=\mathbb P^\ast_{j,\alpha}(e\notin T)$. Summing the usage over $e\in S$,
\[
m_S=\sum_j\Big(s_j\oE_{T\sim\mathbb P_{j,\alpha}}[N_S]+t_j\oE_{T\sim\mathbb P^\ast_{j,\alpha}}\big[|S|-N_S\big]\Big),
\]
and with $d^\ast_j(S)=|S|-r_j(S)$ (\hb{def:skeleton}) the wall $\kappa_S=m_S-\sum_j\big(s_jd_j(S)+t_jd^\ast_j(S)\big)$ is the displayed sum.

\emph{(2) The sign.} By \hb{thm:window}, $d_j(S)\le N_S\le r_j(S)$ on every basis of table $j$. So both brackets are nonnegative, and so is $\kappa_S$, the exponents being nonnegative.

\emph{(3) Equality.} The laws $\mathbb P_{j,\alpha}(T)=\theta_{j,T}\alpha^T/K_j(\alpha)$ and $\mathbb P^\ast_{j,\alpha}(T)=\theta_{j,T}\alpha^{\cE\setminus T}/\Psi_j(\alpha)$ charge every basis. So the first bracket of table $j$ vanishes exactly when $N_S=d_j(S)$ on every basis, and the second exactly when $N_S=r_j(S)$ on every basis; since the values of $N_S$ fill $\{d_j(S),\dots,r_j(S)\}$ (\hb{thm:window}), each happens exactly when $c_j(S)=0$. Every table has $s_j>0$ or $t_j>0$, so $\kappa_S=0$ exactly when $c_j(S)=0$ for every table, that is, by Lemma~\ref{lem:families-separators}, when $S$ is a union of molecules.

\emph{(4) The last sentence.} $\cE$ is a union of molecules, so $\kappa_{\cE}=0$. At a realization every table is active, so the family is connected (\hb{def:family}) exactly when no nonempty proper $S$ has $c_j(S)=0$ for every table, that is, by Lemma~\ref{lem:families-separators}, when no nonempty proper set is a union of molecules; then every proper wall is positive by step~(3).
\end{proof}

\begin{sremark}
The reference, proof of \hb{prop:wallsrealization}, has steps~(1) and~(2) and the positivity of the bracket of an active table with $c_j(S)\ge1$, which connectedness supplies for every nonempty proper $S$. Step~(3) reads off the equality case, which holds for every family.
\end{sremark}

\begin{hproposition}[\hb{prop:massesgauge} (Masses at a gauge)]
As stated, at every gauge $e^{y_0}$.
\end{hproposition}

\begin{proof}
By \hb{def:realization} the masses are $\tau_e=\sum_j\big(s_jR_{j,e}(\alpha)+t_jR^\ast_{j,e}(\alpha)\big)$, with $R_{j,e}=\alpha_e\partial_eK_j/K_j$ the occupation and $R^\ast_{j,e}=\alpha_e\partial_e\Psi_j/\Psi_j$ the dual occupation, and $R^\ast_{j,e}(\alpha)=1-R_{j,e}(\alpha^{-1})$ (\hb{thm:innerMGF}). At $\alpha=e^{y_0}$ this is the first expression. Put $a_j=R_{j,e}(e^{y_0})$ and $b_j=R_{j,e}(e^{-y_0})$, so that $R^\pm_{j,e}=\frac12(a_j\pm b_j)$. With $s_j=\sigma_j+\delta_j$ and $t_j=\sigma_j-\delta_j$,
\[
s_ja_j+t_j(1-b_j)=\sigma_j(1+a_j-b_j)+\delta_j(a_j+b_j-1)=\sigma_j(1+2R^-_{j,e})+\delta_j(2R^+_{j,e}-1).
\]
At $y_0=0$, $a_j=b_j=R_{j,e}(\one)$: $R^-_{j,e}=0$ and $R^+_{j,e}=R_{j,e}(\one)$.
\end{proof}

\begin{sremark}
By \hb{def:weightlaw}, $R_{j,e}(e^{y})=\partial_{y_e}\log K_j(e^y)$, whose Taylor series at $0$ is
\[
\sum_{i\ge0}\frac1{i!}\sum_{u\in\cE^i}\kappa_{j,\one}(e,u)\prod_{l=1}^iy_{u_l}
\]
in its domain of convergence, as in \hb{prop:deviation}. At $y=\pm y_0$ the term of order $i$ carries $(\pm1)^i$, so $R^+_{j,e}$ collects the terms of even order and $R^-_{j,e}$ those of odd order. Rung one at a deviation has the parity of the odd rungs of \hb{prop:deviation}, $\sigma_j$ with the odd part of the occupations and $\delta_j$ with the even part, besides the constant $\sum_j(\sigma_j-\delta_j)=\sum_jt_j$.
\end{sremark}

\begin{hproposition}[\hb{prop:compound} (The potential of a compound splits by molecules)]
As stated; the last sentence is new.
\end{hproposition}

\begin{proof}
The first two sentences are the reference, proof of \hb{prop:compound}. Let $e$ be a loop or a coloop in every table. By \hb{prop:compoundgauge} $\{e\}$ is a unit, and every basis of table $j$ contains $e$, or none does, so $K_j(\alpha)=\alpha_e^{\epsilon_j}\tilde K_j(\alpha_{\cE\setminus e})$ and $\Psi_j(\alpha)=\alpha_e^{1-\epsilon_j}\tilde\Psi_j(\alpha_{\cE\setminus e})$ with $\epsilon_j\in\{0,1\}$. So $\Phi(y)=\sum_j\big(s_j\epsilon_j+t_j(1-\epsilon_j)\big)y_e$ plus a function of $y_{\cE\setminus e}$ (\hb{def:potential}), every derivative of $\Phi$ of order at least two that involves $y_e$ vanishes, every entry of the still on a string through $e$ is zero, and rung one at $e$ is the constant $\sum_j\big(s_j\epsilon_j+t_j(1-\epsilon_j)\big)$, the mass of $e$ (\hb{def:still}). The same holds for the two parts $\Phi_{\mathbf E}$ and $\Phi_{\mathbf O}$, which are $\sum_j\sigma_j\log(K_j\Psi_j)$ and $\sum_j\delta_j\log(K_j/\Psi_j)$ read at $e^y$.
\end{proof}

\begin{hproposition}[\hb{prop:occupations} (Occupations at a gauge)]
As stated, at every gauge $e^{y_0}$.
\end{hproposition}

\begin{proof}
By \hb{def:still}, $\varphi^{\mathbf X}_e(0)=\partial_{y_e}\Phi_{\mathbf X}(\log\alpha)$, and by \hb{def:potential}, $\Phi_{\mathbf E}(y)=\sum_j\sigma_j\big(\log K_j(e^y)+\log\Psi_j(e^y)\big)$ and $\Phi_{\mathbf O}(y)=\sum_j\delta_j\big(\log K_j(e^y)-\log\Psi_j(e^y)\big)$. By \hb{thm:saddle}, $\partial_{y_e}\log K_j(e^y)=R_{j,e}$ and $\partial_{y_e}\log\Psi_j(e^y)=R^\ast_{j,e}$, and $R^\ast_{j,e}(\alpha)=1-R_{j,e}(\alpha^{-1})$ (\hb{thm:innerMGF}). At $\alpha=e^{y_0}$, with $a_j=R_{j,e}(e^{y_0})$ and $b_j=R_{j,e}(e^{-y_0})$,
\begin{gather*}
\varphi^{\mathbf E}_e(0)=\sum_j\sigma_j(a_j+1-b_j)=\sum_j\sigma_j(1+2R^-_{j,e}),\\
\varphi^{\mathbf O}_e(0)=\sum_j\delta_j(a_j-1+b_j)=\sum_j\delta_j(2R^+_{j,e}-1).
\end{gather*}
At $\alpha=\one$, $R^-_{j,e}=0$ and $R^+_{j,e}=R_{j,e}(\one)$. The two parts add to the mass of \hb{prop:massesgauge}.
\end{proof}

\section{Chapter~\hb{ch:reading}: the reading, classes, atoms}\label{sec:ch10}

\begin{hproposition}[\hb{prop:viewform} (A view through the patterns)]
As stated; the inclusion laws are read at every gauge.
\end{hproposition}

\begin{proof}
The expression of the views through $K_{j;e,\mathbf g}$ and $\Psi_{j;e,\mathbf g}$ and the closing sentence on one table are proved in the reference, proof of \hb{prop:viewform}. What follows is the expansion at a gauge $\alpha$ and its reading at the table's gauge. Since $\prod_{i\in p}(1+x_i)=\sum_{A\subseteq p}x^A$,
\begin{gather*}
K_{j;e,\mathbf g}(\one+x)=\sum_p\vartheta^K_{j,p}\sum_{A\subseteq p}x^A=\sum_Ax^A\sum_{p\supseteq A}\vartheta^K_{j,p},\\
\sum_{p\supseteq A}\vartheta^K_{j,p}=\sum_{T:\ p(T)\supseteq A}\theta_{j,T}\alpha^T=K_j(\alpha)\,\mathbb P_{j,\alpha}\big(p(T)\supseteq A\big),
\end{gather*}
the basis law at $\alpha$ being $\mathbb P_{j,\alpha}(T)=\theta_{j,T}\alpha^T/K_j(\alpha)$ (\hb{thm:innerMGF}). Likewise $(\one+x)^{\bar p}=\prod_{i\notin p}(1+x_i)=\sum_{A\cap p=\varnothing}x^A$, so
\begin{gather*}
\Psi_{j;e,\mathbf g}(\one+x)=\sum_Ax^A\sum_{p\cap A=\varnothing}\vartheta^\Psi_{j,p},\\
\sum_{p\cap A=\varnothing}\vartheta^\Psi_{j,p}=\sum_{T:\ p(T)\cap A=\varnothing}\theta_{j,T}\alpha^{\cE\setminus T}=\Psi_j(\alpha)\,\mathbb P^\ast_{j,\alpha}\big(p(T)\cap A=\varnothing\big),
\end{gather*}
with $\mathbb P^\ast_{j,\alpha}(T)=\theta_{j,T}\alpha^{\cE\setminus T}/\Psi_j(\alpha)$. So the polynomials $K_{j;e,\mathbf g}$ and $\Psi_{j;e,\mathbf g}$, and with them the view, are determined by $K_j(\alpha)$, $\Psi_j(\alpha)$ and the joint inclusion laws of $e,g_1,\dots,g_k$ under $\mathbb P_{j,\alpha}$ and $\mathbb P^\ast_{j,\alpha}$. At $\alpha=\one$, $\alpha^T=\alpha^{\cE\setminus T}=1$: $\vartheta^K_{j,p}=\vartheta^\Psi_{j,p}=\vartheta_{j,p}$, both laws are the weight law (\hb{thm:innerMGF}), and $x^{\bar p}=x^{\{0,\dots,k\}}(x^{-1})^p$ gives $\Psi_{j;e,\mathbf g}(x)=x^{\{0,\dots,k\}}K_{j;e,\mathbf g}(x^{-1})$.
\end{proof}

\begin{hproposition}[\hb{prop:computeclasses} (Computing the classes)]
As stated; at a gauge $\alpha$ the identities have their coefficients in the field generated by the weights and the coordinates of $\alpha$.
\end{hproposition}

\begin{proof}
The steps of the computation compare views and writings, as in the reference, proof of \hb{prop:computeclasses}; only the field is new. Let $k$ be the field generated by the weights and the coordinates of $\alpha$. By \hb{prop:viewform}, a view is the pair $\sum_j\sigma_j\log(K_{j;e,\mathbf g}\Psi_{j;e,\mathbf g})$, $\sum_j\delta_j\log(K_{j;e,\mathbf g}/\Psi_{j;e,\mathbf g})$ in $x_i=e^{v_i}$, and the coefficients $\vartheta^K_{j,p}=\sum_{T:\,p}\theta_{j,T}\alpha^T$ and $\vartheta^\Psi_{j,p}=\sum_{T:\,p}\theta_{j,T}\alpha^{\cE\setminus T}$ of these polynomials are polynomials with integer coefficients in the weights and the coordinates of $\alpha$, so they lie in $k$. Factor every $K_{j;e,\mathbf g}$ and $\Psi_{j;e,\mathbf g}$ into irreducible polynomials over $k$. Each factor of a polynomial positive on the orthant vanishes nowhere on it and so has constant sign there, and it is normalized to the value $1$ at $x=\one$ by a constant of $k$. Two different normalized polynomials irreducible over $k$ have no common factor over $\mathbb C$, the greatest common divisor of two polynomials over $k$ being their greatest common divisor over $\mathbb C$. So the factors over $\mathbb C$ of \hb{thm:poles} are the factors over $\mathbb C$ of the factors over $k$, and a factor over $\mathbb C$ divides exactly one normalized factor over $k$. If two views carry over $\mathbb C$ the same factors with the same pairs of exponents, a factor over $k$ of one and a factor over $k$ of the other that share a factor over $\mathbb C$ are equal, and comparing the exponents of that shared factor gives equal pairs of exponents. So by \hb{thm:poles} two views, or two writings, are equal exactly when they carry the same normalized polynomials irreducible over $k$ with the same pairs of exponents: every step is an identity in $k[x]$. At $\alpha=\one$, $k$ is the field generated by the weights, $\mathbb Q$ for rational weights, and the coefficients are integers for lists.
\end{proof}

\begin{hproposition}[\hb{prop:massclass} (Masses and occupations on a class)]
As stated; the first sentence holds for each of the two parts of the mass.
\end{hproposition}

\begin{proof}
By \hb{def:still}, $\varphi^{\mathbf X}_e(0)=\partial_{y_e}\Phi_{\mathbf X}(\log\alpha)$ is the part $\mathbf X$ of rung one, and $\tau_e(\Upsilon)=\partial_{y_e}\Phi(\log\alpha)=\varphi^{\mathbf E}_e(0)+\varphi^{\mathbf O}_e(0)$. By \hb{prop:refine}(iii) the reading of $e$ determines its self-line functions $\varphi^{\mathbf E}_e$ and $\varphi^{\mathbf O}_e$, hence their values at $0$, and the classes are the level sets of the reading (\hb{def:view}). So the two parts, and their sum, are constant on every class. For a family whose weights factor across its molecules, $\Phi_{\mathbf X}$ is the sum of the potentials of the restrictions to the molecules (\hb{prop:compound}), so $\varphi^{\mathbf X}_e(0)$ is the same in the family and in the realization of the molecule of $e$, on which its classes are read. Items~(i) and~(ii) are proved in the reference, proof of \hb{prop:massclass}.
\end{proof}

\section{Chapter~\hb{ch:lean}: lean, gauge and label}\label{sec:ch11}

\begin{hlemma}[\hb{lem:nodegenerate} (No degenerate pair)]
As stated, at every realization of every family, with the equality case.
\end{hlemma}

\begin{proof}
\emph{(1) A sum of covariances.} By \hb{def:still} and \hb{def:potential}, $\kappa^{\mathbf E}_\Upsilon(e,f)=4\sum_j\sigma_j\big(\partial_{y_e}\partial_{y_f}\log K_j(e^y)+\partial_{y_e}\partial_{y_f}\log\Psi_j(e^y)\big)$ at $y=\log\alpha$. By \hb{def:weightlaw} the first derivative is $\kappa_{j,\alpha}(e,f)$, and since $\log\Psi_j(e^y)=\ip\one y+\log K_j(e^{-y})$ (\hb{prop:identity}) the second is $\kappa_{j,\alpha^{-1}}(e,f)$. So $\kappa^{\mathbf E}_\Upsilon=4\sum_j\sigma_j\big(C_{j,\alpha}+C_{j,\alpha^{-1}}\big)$, $C_{j,\beta}$ the covariance matrix of the indicators $(X_e)_e$ under $\mathbb P_{j,\beta}$, and for $w\in\mathbb R^{\cE}$
\[
w^{\mathsf T}\kappa^{\mathbf E}_\Upsilon w=4\sum_j\sigma_j\Big(\operatorname{Var}_{T\sim\mathbb P_{j,\alpha}}\ip w{\one_T}+\operatorname{Var}_{T\sim\mathbb P_{j,\alpha^{-1}}}\ip w{\one_T}\Big)\ \ge\ 0,
\]
since $\sigma_j>0$ at a realization.

\emph{(2) The inequality.} The principal $2\times2$ submatrix of $\kappa^{\mathbf E}_\Upsilon$ on $\{e,f\}$ is positive semidefinite, so its determinant $\kappa^{\mathbf E}_\Upsilon(e,e)\kappa^{\mathbf E}_\Upsilon(f,f)-\kappa^{\mathbf E}_\Upsilon(e,f)^2$ is nonnegative. It vanishes exactly when that submatrix is singular, that is when some $w=\lambda\one_e+\mu\one_f\neq0$ has $w^{\mathsf T}\kappa^{\mathbf E}_\Upsilon w=0$, a positive semidefinite matrix annihilating exactly the vectors at which its form vanishes.

\emph{(3) The vectors of form zero.} Both laws in step~(1) charge every basis (\hb{thm:innerMGF}), so $w^{\mathsf T}\kappa^{\mathbf E}_\Upsilon w=0$ exactly when $\ip w{\one_T}$ is constant over the bases of every table. By \hb{lem:circuit} such a $w$ is constant on every molecule. Conversely, if $w$ is constant on every molecule, with value $w_M$ on the molecule $M$, then $\ip w{\one_T}=\sum_Mw_Mr_j(M)$ on every basis of table $j$, each table being the direct sum of its restrictions to the molecules (\hb{prop:directsum}).

\emph{(4) The cases.} If $e$ is a unit, $w=\one_e$ is constant on every molecule, and equality holds; likewise if $f$ is a unit. If neither is a unit and they lie in different molecules, the molecule of $e$ has a mark other than $e$, not $f$, at which $w=\lambda\one_e+\mu\one_f$ vanishes, so $\lambda=0$; likewise $\mu=0$, and the inequality is strict. If $e$ and $f$ lie in one molecule $M$: when $M$ has a third mark, $w$ vanishes there, so $\lambda=\mu=0$ and the inequality is strict; when $M=\{e,f\}$, $w=\one_e+\one_f$ is constant on every molecule, and equality holds. A molecule with at least three marks contains neither a unit nor a molecule of two marks, so the inequality is strict for any two of its marks.
\end{proof}

\begin{sremark}
The reference, proof of \hb{lem:nodegenerate}, has steps~(1) to~(3) and concludes for a molecule with at least three marks; the case analysis of step~(4) gives the equality case for every family.
\end{sremark}

\begin{hlemma}[\hb{lem:equaldiag} (Equal diagonals on a class)]
As stated: on every class, for both parts and every rung.
\end{hlemma}

\begin{proof}
By \hb{def:still}, $\varphi^{\mathbf X}_e(v)=\partial_{y_e}\Phi_{\mathbf X}(\log\alpha+v\one_e)$, so $(\varphi^{\mathbf X}_e)^{(k-1)}(0)=\partial^k_{y_e}\Phi_{\mathbf X}(\log\alpha)$, and the still of each part gives
\[
\kappa^{\mathbf X}_\Upsilon(e,\dots,e)=2^k\,\partial^k_{y_e}\Phi_{\mathbf X}(\log\alpha)=2^k(\varphi^{\mathbf X}_e)^{(k-1)}(0).
\]
By \hb{prop:refine}(iii) the reading of $e$ determines $\varphi^{\mathbf E}_e$ and $\varphi^{\mathbf O}_e$; the marks of a class have one reading (\hb{def:view}), so they have the same self-line functions, and the same derivatives of every order at $0$. At an atom there is one class.
\end{proof}

\begin{sremark}
The reference states \hb{lem:equaldiag} at an atom of a molecule with at least three marks, for $\kappa^{\mathbf E}(e,e)$, and proves it through the rung-two tallies and \hb{lem:nodegenerate}. The self-line functions carry every order and both parts, and make the statement one about classes.
\end{sremark}

\begin{hcorollary}[\hb{cor:exchangecumulants} (An exchange negates odd cumulants)]
As stated, for an exchange of any family.
\end{hcorollary}

\begin{proof}
\emph{(1) The weight laws.} By \hb{def:carry}, $T\mapsto T'=\cE\setminus g(T)$ maps $\cB_j$ onto $\cB_{h(j)}$ with $\theta_{h(j),T'}=\theta_{j,T}$, and it is injective because $g$ is. So $K_{h(j)}(\one)=K_j(\one)$, and for $T\sim\mathbb P_{j,\one}$ the basis $T'$ has the law $\mathbb P_{h(j),\one}$.

\emph{(2) The cumulants.} $X_{g(e)}(T')=\one[g(e)\notin g(T)]=1-X_e(T)$. Joint cumulants of order at least two are multilinear and unchanged when constants are added to the variables, so for a string $w=(e_1,\dots,e_k)$, $k\ge2$,
\begin{align*}
\kappa_{h(j),\one}(g(w))&=\kappa\big(X_{g(e_1)}(T'),\dots,X_{g(e_k)}(T')\big)\\
&=\kappa\big(1-X_{e_1}(T),\dots,1-X_{e_k}(T)\big)=(-1)^k\kappa_{j,\one}(w).
\end{align*}

\emph{(3) The stills.} By \hb{prop:stillgauge}, $\kappa_{\Upsilon_{\boldsymbol\sigma,\boldsymbol\delta;\one}}(w)=2^k\sum_j\big(s_j+(-1)^kt_j\big)\kappa_{j,\one}(w)$. The condition on the pairs in \hb{def:carry} gives $(s_{h(j)},t_{h(j)})=(s_j,t_j)$, and $h$ is a bijection of the tables; reindexing the sum by $h$ and using step~(2),
\begin{align*}
\kappa_{\Upsilon_{\boldsymbol\sigma,\boldsymbol\delta;\one}}(g(w))&=2^k\sum_j\big(s_{h(j)}+(-1)^kt_{h(j)}\big)\kappa_{h(j),\one}(g(w))\\
&=2^k\sum_j\big(t_j+(-1)^ks_j\big)\kappa_{j,\one}(w),
\end{align*}
the still at $\Upsilon_{\boldsymbol\sigma,-\boldsymbol\delta;\one}$, whose exponents are $(t_j,s_j)$.

\emph{(4) The classes of the centre.} By \hb{prop:certify} with $F'=F$ at the table's gauge, the reading of $g(e)$ at $-\boldsymbol\delta$ equals that of $e$ at $\boldsymbol\delta$; at the centre $-\boldsymbol\delta=\boldsymbol\delta=0$, so $g(e)$ lies in the class of $e$. Hence $g$ maps every class of the centre into itself, and onto itself, $g$ being injective and the class finite.
\end{proof}

\begin{sremark}
The reference, proof of \hb{cor:exchangecumulants}, is the case of one table, where $h$ is the identity.
\end{sremark}

\section{Chapter~\hb{ch:bunch}: restriction to a bunch}\label{sec:ch12}

\begin{hlemma}[\hb{lem:mot-childwidth} (Children do not widen)]
As stated.
\end{hlemma}

\begin{proof}
The child at $p$ is a table on $R$ (\hb{lem:minor}), and a set $X\subseteq R$ is one of its bases exactly when $p\sqcup X$ is a basis of the table. Since $p,W\subseteq B$ and $X,A\subseteq R$,
\[
|X\cap A|=|(p\sqcup X)\cap(A\cup W)|-|p\cap W| .
\]
So the traces of $A$ on the bases of the child are traces of $A\cup W$ on bases of the table, all shifted by the one constant $|p\cap W|$. The width is the largest trace less the least (\hb{thm:window}), so $c_{(p)}(A)\le c(A\cup W)$.
\end{proof}

\begin{hproposition}[\hb{prop:coinsrank} (Coins of a tensorminental table)]
As stated: tensor rank one over $\mathbb R$ or over $\mathbb C$.
\end{hproposition}

\begin{proof}
The first display is proved in the reference, proof of \hb{prop:coinsrank}. At tensor rank one there is one slice $A\in\mathbb F^{r\times n}$, of rank $r$ since $K(\one)=\det(AA^{\mathsf H})>0$ (\hb{prop:tensorletters}); the only pair is $(1,1)$, with $w_{11}=1$, and $\Pi^{(11)}=A^{\mathsf H}M^{-1}AD_\alpha$, $M=AD_\alpha A^{\mathsf H}$ (\hb{thm:interference}). So $x^{d(B)}\sQ^K_B(x;\alpha)=\prod_{i=1}^{|B|}(1-p_i+p_ix)$ with $p_i$ the eigenvalues of $\Pi^{(11)}_{BB}$.

\emph{(1) The projection.} Put $P=D_\alpha^{1/2}\Pi^{(11)}D_\alpha^{-1/2}=D_\alpha^{1/2}A^{\mathsf H}M^{-1}AD_\alpha^{1/2}$. Then $P^{\mathsf H}=P$ and $P^2=D_\alpha^{1/2}A^{\mathsf H}M^{-1}(AD_\alpha A^{\mathsf H})M^{-1}AD_\alpha^{1/2}=P$, and $P\,D_\alpha^{1/2}A^{\mathsf H}=D_\alpha^{1/2}A^{\mathsf H}$; so $P$ is the orthogonal projection onto $\{D_\alpha^{1/2}A^{\mathsf H}y\}$, the row space of $AD_\alpha^{1/2}$ in the sense of \hb{thm:interference}.

\emph{(2) Similarity.} $D_\alpha$ is diagonal, so $P_{BB}=D_B^{1/2}\Pi^{(11)}_{BB}D_B^{-1/2}$, with $D_B$ the diagonal block on $B$, and $P_{BB}$ has the eigenvalues $p_i$.

\emph{(3) Principal angles.} $P_{BB}$ is the compression of $P$ to $\mathbb F^B$, the vectors supported on $B$: $v^{\mathsf H}P_{BB}v=\|Pv\|^2$ for $v\in\mathbb F^B$. With $E_B$ the inclusion of $\mathbb F^B$ and $Q$ a matrix whose columns are an orthonormal basis of the range of $P$, $P_{BB}=(E_B^{\mathsf H}Q)(E_B^{\mathsf H}Q)^{\mathsf H}$. Its eigenvalues are the squares of the singular values of $E_B^{\mathsf H}Q$, completed by zeros, which are the cosines of the principal angles between $\mathbb F^B$ and the range of $P$; they lie in $[0,1]$.

\emph{(4) The eigenvalue $1$.} For $v\in\mathbb F^B$, $P_{BB}v=v$ exactly when $\|Pv\|=\|v\|$, that is when $v$ lies in the range of $P$. So the eigenspace is $\{D_\alpha^{1/2}A^{\mathsf H}y:\ (D_\alpha^{1/2}A^{\mathsf H}y)_R=0\}=\{D_\alpha^{1/2}A^{\mathsf H}y:\ A_R^{\mathsf H}y=0\}$, of dimension $r-\operatorname{rank}A_R$, the map $y\mapsto D_\alpha^{1/2}A^{\mathsf H}y$ being injective. The bases of the table are the sets $T$ with $\det A_T\neq0$ (\hb{def:multivector}), so $r(R)=\operatorname{rank}A_R$, and $r-r(R)=d(B)$ (\hb{def:skeleton}, $d(B)=|B|-r^\ast(B)$ with $r^\ast(B)=|B|+r(R)-r$).

\emph{(5) The eigenvalue $0$.} For $v\in\mathbb F^B$, $P_{BB}v=0$ exactly when $Pv=0$, that is $AD_\alpha^{1/2}v=A_BD_B^{1/2}v_B=0$. So the eigenspace has dimension $|B|-\operatorname{rank}A_B=|B|-r(B)=d^\ast(B)$.

\emph{(6) The coins.} $P_{BB}$ is Hermitian, so of its $|B|$ eigenvalues $d(B)$ are $1$, $d^\ast(B)$ are $0$, and the other $|B|-d(B)-d^\ast(B)=c(B)$ lie in $(0,1)$. In $x^{d(B)}\sQ^K_B(x;\alpha)=\prod_i(1-p_i+p_ix)$ the eigenvalues $1$ give the factor $x^{d(B)}$ and the eigenvalues $0$ the factor $1$; so the coins of $\sQ^K_B$ are the eigenvalues $p_i=\cos^2\vartheta_i\in(0,1)$, and the Skeleton values are $\xi_i=(1-p_i)/p_i=\tan^2\vartheta_i>0$ (\hb{def:skelpoly}).
\end{proof}

\begin{sremark}
Over $\mathbb R$ this is the reference, proof of \hb{prop:coinsrank}. Over $\mathbb C$ the projection is Hermitian, and each step reads the same with conjugate transposes.
\end{sremark}

\begin{hdefinition}[\hb{def:interaction}: the claims]
$\pat(B)=2^B$ holds exactly when $B$ is independent ($B$ is a pattern) and some basis avoids $B$ ($\varnothing$ is a pattern): if both hold, $d(B)=0$, so the contraction $\cB/R$ has rank $0$, and every subset of $B$ is independent with $r_{\cB/R}=0=d(B)$, a pattern by Proposition~\hb{prop:higgs}(i). Equivalently $r(B)=|B|$ and $d(B)=0$, that is $c(B)=|B|$. The dual table has the bases $\cE\setminus T$ and the patterns $B\setminus p$, with the children of cobasis letters $\Psi_{(p)}$; so every subset of $B$ is a pattern of the dual table when every subset is a pattern of the table, and its interaction is $\sum_{p}(-1)^{|B\setminus(B\setminus p)|}\log\Psi_{(p)}=\iota^\ast_B$. The parts: $\Phi_{\mathbf E}=\sum_j\sigma_j(\log K_j+\log\Psi_j)$ and $\Phi_{\mathbf O}=\sum_j\delta_j(\log K_j-\log\Psi_j)$ (Definition~\hb{def:potential}), and $\sigma_j\pm\delta_j$ are $s_j$ and $t_j$.
\end{hdefinition}

\begin{hproposition}[\hb{prop:mot-chainbottom} (The bottom of a chain factors over its steps)]
As stated.
\end{hproposition}

\begin{proof}
\emph{Existence.} Order the marks of each step arbitrarily and the steps as they come; the sets $S_i$ are initial segments of this order of $\cE$, which are the tails of the reversed order, and \hb{lem:greedy} gives a basis $T^+$ with $|T^+\cap U|=d(U)$ for every tail $U$, in particular at the bottom of every $S_i$.

\emph{One step is exchanged at a time.} Let $T$ and $T'$ be at the bottom of every $S_k$, and $1\le i\le\ell+1$. For $i\ge2$, by \hb{thm:levels} at the bunch $S_{i-1}$ the bases at its bottom are all the unions of a pattern of that level and a room side of that level, so $T''=(T'\cap S_{i-1})\cup(T\setminus S_{i-1})$ is a basis; for $i=1$ put $T''=T$. Its traces are $|T''\cap S_k|=|T'\cap S_k|=d(S_k)$ for $k<i$ and $d(S_{i-1})+|T\cap S_k|-|T\cap S_{i-1}|=d(S_k)$ for $k\ge i$. For $i\le\ell$, \hb{thm:levels} at the bunch $S_i$, applied to $T''$ and $T'$, both at its bottom, gives the basis $(T''\cap S_i)\cup(T'\setminus S_i)=(T'\setminus Z_i)\cup(T\cap Z_i)$; for $i=\ell+1$ this set is $T''$ itself. Its trace on $S_k$ is that of $T'$ for $k<i$, and for $k\ge i$ it is
\[
d(S_{i-1})+\big(d(S_i)-d(S_{i-1})\big)+\big(d(S_k)-d(S_i)\big)=d(S_k),
\]
since $T$ and $T'$ meet $Z_i$ in $d(S_i)-d(S_{i-1})$ marks and $S_k\setminus S_i$ in $d(S_k)-d(S_i)$. So the part on one step of a basis at the bottom of every $S_k$ can be replaced by that of any other, and these bases are exactly the unions of their traces on the steps, chosen independently.

\emph{The step tables.} If $(T_0\setminus Z_i)\sqcup X$ is a basis with $X\subseteq Z_i$, then $|X|=|T_0\cap Z_i|$, its traces on the $S_k$ are those of $T_0$, and it is at the bottom of every $S_k$. Conversely, for $T$ at the bottom of every $S_k$ the exchange gives the basis $(T_0\setminus Z_i)\cup(T\cap Z_i)$. So the traces on $Z_i$ are the bases of the child at the pattern $T_0\setminus Z_i$, a table (\hb{lem:minor}), the same set for every $T_0$. A basis $T$ at the bottom of $S_{i-1}$ and of $S_i$ alone has the same trace on $Z_i$ as one at the bottom of every $S_k$: the two exchanges above with $T'=T_0$ use only that $T$ is at the bottom of $S_{i-1}$ and of $S_i$, and give the basis $(T_0\setminus Z_i)\cup(T\cap Z_i)$, whose traces were computed above. So the step table depends on $S_{i-1}$ and $S_i$ alone.

\emph{The top.} $|T\cap S_i|$ is largest exactly when the cobasis $\cE\setminus T$ meets $S_i$ least, so the bases at the top of every $S_i$ are the complements of the bases at the bottom of every $S_i$ in the dual table (\hb{def:table}), and the above applies there. For $q$ a pattern of $\cE\setminus Z_i$ in the dual table, $X\subseteq Z_i$ is a basis of its child exactly when $\big((\cE\setminus Z_i)\setminus q\big)\sqcup(Z_i\setminus X)$ is a basis of $\cB$: the child of the dual table at $q$ is the dual of the child of $\cB$ at $(\cE\setminus Z_i)\setminus q$. At $q=(\cE\setminus T_1)\setminus Z_i$ this is the dual of $\cB_{(T_1\setminus Z_i)}$.

\emph{The letters.} The degree of $\alpha^T$ in $\alpha_{S_i}$ is $|T\cap S_i|$, whose least value is $d(S_i)$, attained for every $i$ at once by the bases at the bottom of every $S_i$; for a list their sum is $\sum_{X_1,\dots,X_{\ell+1}}\prod_i\alpha^{X_i}=\prod_iK_{\cB_{(T_0\setminus Z_i)}}(\alpha_{Z_i})$, and with weights $\theta_T=\prod_i\theta_i(T\cap Z_i)$ the same with the weights $\theta_i$. The degree of $\alpha^{\cE\setminus T}$ in $\alpha_{S_i}$ is $|S_i\setminus T|$, least exactly when $T$ is at the top of $S_i$, and by the dual statement the terms of $\Psi$ of the least degree for every $i$ at once are $\prod_iK_{(\cB_{(T_1\setminus Z_i)})^\ast}(\alpha_{Z_i})=\prod_i\Psi_{\cB_{(T_1\setminus Z_i)}}(\alpha_{Z_i})$ (\hb{prop:identity}). The bases of $\binom J1$ on $\cE$ are the $\{j\}$, $j\in J$, and $d(S)$ is $1$ or $0$ as $J\subseteq S$ or not; so $\{j\}$ is at the bottom of every $S_k$ exactly when $j\notin S_k$ for every $k$ with $J\not\subseteq S_k$, that is when $j\in J\cap Z_{i_0}$, $Z_{i_0}$ the first step with $J\subseteq S_{i_0}$. Its step table at $Z_{i_0}$ has the bases $\{j\}$, $j\in J\cap Z_{i_0}$, and its step tables at the other steps the one basis $\varnothing$.

\emph{A maximal chain.} When every step is one mark, a basis at the bottom of every $S_i$ meets the step $\{e_i\}$ in $d(S_i)-d(S_{i-1})$ marks, so it is determined by the chain, and it exists; its term in $K$ is one monomial, with the coefficient $1$ for a list.
\end{proof}

\begin{hcorollary}[\hb{cor:levelsplit} (The native level split is interior)]
As stated: the directions come from chains of circuits joining the marks of each side, which a connected table with a positive exponent supplies.
\end{hcorollary}

\begin{proof}
The first two sentences, the relative interior and the walls with their exceptions, are proved in the reference, proof of \hb{cor:levelsplit}, without any connectedness. Write $\mathcal L$ for the set of level and colevel letters with positive exponent, $c_L>0$ for the exponent of $L\in\mathcal L$, $N_L$ for the convex hull of its exponent vectors, and $Q=\sum_{L\in\mathcal L}c_LN_L$ for the Minkowski sum of the statement.

\emph{(1) A circuit with two marks on one side.} Let $a\neq a'$ lie both in $B$ or both in $R$, and in a common circuit of a table $j$ with a positive exponent. By \hb{lem:linkexchange} some basis $T$ of table $j$ has $a\in T$, $a'\notin T$ and $T'=(T\setminus\{a\})\cup\{a'\}$ a basis, and $|T\cap B|=|T'\cap B|=:k$, the two marks lying on one side. At the native gauge every basis has positive probability under both laws (\hb{thm:innerMGF}). If $s_j>0$, the level letter $K^{(k)}_{j,B}$ has the positive exponent $s_j\mathbb P_{j,\alpha^\star}(N_B=k)$, and $N_{K^{(k)}_{j,B}}$ contains $\one_T$ and $\one_{T'}$, whose difference is $\one_a-\one_{a'}$. If $t_j>0$, the colevel letter $\Psi^{(|B|-k)}_{j,B}$ has the positive exponent $t_j\mathbb P^\ast_{j,\alpha^\star}(N^\ast_B=|B|-k)$, and its polytope contains $\one_{\cE\setminus T}$ and $\one_{\cE\setminus T'}$, whose difference is $\one_{a'}-\one_a$. Fixing points of the other summands, $c_L(\one_a-\one_{a'})$ is the difference of two points of $Q$ for one $L\in\mathcal L$, so every $v$ with $\ip vx$ constant on $Q$ has $v_a=v_{a'}$.

\emph{(2) The directions.} Under the hypothesis the chains join all the marks of $B$ and all the marks of $R$, so by step~(1) such a $v$ is constant on $B$ and on $R$. Conversely, a monomial $\alpha^T$ of a level letter of level $k$ has $\ip{\one_B}{\one_T}=k$, and a monomial $\alpha^{\cE\setminus T}$ of a colevel letter of colevel $k$ has $\ip{\one_B}{\one_{\cE\setminus T}}=|B\setminus T|=k$; so $\ip{\one_B}x$ is constant on every $N_L$, hence on $Q$, and so is $\ip{\one_R}x=\ip\one x-\ip{\one_B}x$, each letter being homogeneous. So the directions $v$ with $\ip vx$ constant on $Q$ are $\operatorname{span}\{\one_B,\one_R\}$.

\emph{(3) The exceptions.} A linear function is constant on a Minkowski sum with positive coefficients exactly when it is constant on every summand. So $S$ is an exception exactly when $\one_S\in\operatorname{span}\{\one_B,\one_R\}$, that is $S\in\{\varnothing,B,R,\cE\}$.

\emph{(4) Convergence.} Write $\alpha_B=\rho_Bu$ with $\rho_B=\ip{\one_B}{\alpha_B}$ and $\alpha_B=e^{y_B}$; then $d\alpha_B=\rho_B^{|B|-1}d\rho_B\,d\omega_B(u)=\prod_{e\in B}\alpha_e\,dy_B$, and splitting $y_B$ into its component $t$ along the unit vector $|B|^{-1/2}\one_B$, which changes $\rho_B$ by the factor $e^{t|B|^{-1/2}}$ and leaves $u$ unchanged, and its component orthogonal to $\one_B$, on which $u$ depends, gives $d\rho_B/\rho_B=|B|^{-1/2}dt$ and so $d\omega_B(u)=|B|^{1/2}\prod_{e\in B}u_e\,dy_B$ on the hyperplane $\ip{y_B}{\one_B}=0$; likewise on $R$. So with $Y=\{y:\ \ip{y_B}{\one_B}=\ip{y_R}{\one_R}=0\}$ the integral of the statement is $(|B||R|)^{1/2}\int_Yu^{\btau|_B}v^{\btau|_R}\prod_{L\in\mathcal L}L(u,v)^{-c_L}dy$, with $u=e^{y_B}/\rho_B$ and $v=e^{y_R}/\rho_R$. Each $L$ is homogeneous in $\alpha_B$ and in $\alpha_R$, and the degrees weighted by the exponents add to $\tau_B$ and to $\tau_R$: for the level letters of table $j$ they add to $s_j\oE_{T\sim\mathbb P_{j,\alpha^\star}}[N_B]=s_j\sum_{e\in B}R_{j,e}(\alpha^\star)$ in $\alpha_B$, for its colevel letters to $t_j\oE_{T\sim\mathbb P^\ast_{j,\alpha^\star}}[N^\ast_B]=t_j\sum_{e\in B}R^\ast_{j,e}(\alpha^\star)$, and the usage at the native gauge is $\btau$ (\hb{def:native}); likewise on $R$. So the powers of $\rho_B$ and $\rho_R$ cancel, and the integrand is $e^{\ip{\btau}y}\prod_LL(e^y)^{-c_L}$. Each $L(e^y)$ is at least its least coefficient times $e^{\max_{x\in N_L}\ip xy}$, the maximum being attained at an exponent vector; so the integrand is at most $C\exp\big(-(h(y)-\ip{\btau}y)\big)$, with $h(y)=\max_{x\in Q}\ip xy=\sum_Lc_L\max_{x\in N_L}\ip xy$. By step~(2) the affine hull of $Q$ has the direction space $\operatorname{span}\{\one_B,\one_R\}^\perp=Y$, and $\btau$ lies in the relative interior of $Q$; so for some $\varepsilon>0$, $\btau+w\in Q$ for all $w\in Y$ with $\|w\|\le\varepsilon$, and $h(y)\ge\ip{\btau}y+\varepsilon\|y\|$ for $y\in Y$. The integrand is at most $Ce^{-\varepsilon\|y\|}$, integrable over $Y$.

\emph{(5) A connected table.} In a connected table any two marks lie in a common circuit, the fact used in the reference, proof of \hb{cor:levelsplit}; so a connected table with a positive exponent joins any two marks of $B$, and any two marks of $R$, by a chain of length one.
\end{proof}

\begin{sremark}
The hypothesis cannot be replaced by connectedness of the family. On $\cE=\{a,a',b,b'\}$ with $B=\{a,a'\}$, let the three tables of rank one with the bases $\{a\},\{b\}$, then $\{a'\},\{b'\}$, then $\{a\},\{b'\}$ carry positive exponents, read at a native gauge. The family is connected, but every level of every table holds one basis, so every level and colevel polytope is a point, every set is an exception, and the wall at $\{a\}$ vanishes. Each circuit $\{a,b\}$, $\{a',b'\}$, $\{a,b'\}$ meets both sides, and no chain on one side joins $a$ to $a'$.
\end{sremark}

\section{Chapter~\hb{ch:order}: the order of classes and the Foreman bunch}\label{sec:ch13}

\begin{hproposition}[\hb{prop:mirror} (Mirror classes)]
As stated: levels mirrored up to a bijection of the tables, and the exchange of any family.
\end{hproposition}

\begin{proof}
\emph{(i)} Fix $j$. Under $\mathbb P_{h(j),\one}$ the level $N_D$ takes the value $a$ with probability $\mathbb P_{j,\one}(N_C=k_j-a)$, so it has the law of $k_j-N_C$ under $\mathbb P_{j,\one}$. A cumulant of order $a\ge2$ is unchanged by adding a constant and is homogeneous of degree $a$, so $\kappa_{h(j),\one}(N_D^{\times a})=(-1)^a\kappa_{j,\one}(N_C^{\times a})$. Since $h$ is a bijection of the tables with $\sigma_{h(j)}=\sigma_j$,
\[
\sum_j\sigma_j\kappa_{j,\one}(N_D^{\times a})=\sum_j\sigma_{h(j)}\kappa_{h(j),\one}(N_D^{\times a})=(-1)^a\sum_j\sigma_j\kappa_{j,\one}(N_C^{\times a}),
\]
and likewise with the $\delta_j$. By \hb{cor:firststages}(i) the coefficients of $v_0^a$, $a\ge2$, in the $\mathbf E$ part of stage $0$ are $\frac2{a!}\sum_j\sigma_j\kappa_{j,\one}(N_C^{\times a})$ at even $a$ and zero at odd $a$, and in its $\mathbf O$ part $\frac2{a!}\sum_j\delta_j\kappa_{j,\one}(N_C^{\times a})$ at odd $a$ and zero at even $a$; so the $\mathbf E$ parts of $\mathcal V_0(C)$ and $\mathcal V_0(D)$ are equal and their $\mathbf O$ parts are negatives.

\emph{(ii)} Stage $0$ holds one pair. By (i) the $\mathbf E$ parts agree, so the pairs are compared by their $\mathbf O$ parts (\hb{def:order}(iii)). Their coefficients vanish at every even $a$ and at every odd $a$ below the least odd $a$ with $c=\frac2{a!}\sum_j\delta_j\kappa_{j,\one}(N_C^{\times a})\neq0$, where $C$ has $c$ and $D$ has $-c$. So stage $0$ separates $C$ and $D$, and $C\prec D$ exactly when $c<-c$, that is $c<0$. At the opposite lean $c$ changes sign, being linear in $\boldsymbol\delta$, and the order is the other one.

\emph{(iii)} At the centre the $\mathbf O$ parts vanish and the $\mathbf E$ parts agree by (i), so stage $0$ does not separate $C$ and $D$. They are different classes with equally many marks and $\prec$ is a strict total order on them (\hb{thm:order}(i)), so a later stage does.

\emph{The two instances.} If $C$ and $D$ together make up the molecule, $|T\cap D|=r_j-|T\cap C|$ on every basis of table $j$: $h$ is the identity and $k_j=r_j$. If $g$ exchanges the family, with $h$ the bijection of its tables, then for $T\sim\mathbb P_{j,\one}$ the basis $T'=\cE\setminus g(T)$ has the law $\mathbb P_{h(j),\one}$ (proof of \hb{cor:exchangecumulants}, step~(1)) and $|T'\cap g(C)|=|g(C)\setminus g(T)|=|C\setminus T|=|C|-|T\cap C|$, so $\mathbb P_{h(j),\one}(|T\cap g(C)|=a)=\mathbb P_{j,\one}(|T\cap C|=|C|-a)$ and $k_j=|C|$. Moreover $g(C)$ is a class at the lean with $|C|$ marks: $g$ maps the classes at $\boldsymbol\delta$ onto those at $-\boldsymbol\delta$ (\hb{prop:certify}), which are the classes at $\boldsymbol\delta$ (\hb{prop:leanray}).
\end{proof}

\begin{sremark}
The reference, proof of \hb{prop:mirror}, takes the mirror table by table, $h$ the identity, and the exchange of a one-table molecule. Its conclusions use only the sums over the tables weighted by the $\sigma_j$ and the $\delta_j$, which a bijection of the tables preserving the pairs leaves unchanged.
\end{sremark}

\section{Chapter~\hb{ch:tree}: peel, children and the Foreman tree}\label{sec:S1417-ch14}

\begin{hcorollary}[\hb{cor:onetableexchange} (A molecule with an exchange)]
As stated: a molecule with any number of tables; the classes of the centre are those of \hb{cor:exchangecumulants}.
\end{hcorollary}

\begin{proof}
Read $M$ as the family $F|M$ of the restrictions of the tables to $M$ (\hb{def:peel}), with the pair coordinates $(\sigma_j,\delta_j)$ of its tables. By \hb{def:carry} an exchange $g$ of $M$ comes with a bijection $h$ of the tables such that $T\mapsto M\setminus g(T)$ maps the bases of table $j$ onto those of table $h(j)$, with their weights, and $(\sigma_{h(j)},\delta_{h(j)})=(\sigma_j,\delta_j)$. The first two conditions concern the tables alone, and the third holds with $\boldsymbol\delta$ replaced by $-\boldsymbol\delta$ or by $0$ when it holds at $\boldsymbol\delta$. So $g$ exchanges $M$ at each of the leans $\boldsymbol\delta$, $-\boldsymbol\delta$ and $0$; likewise a bijection carrying $M$ at $\boldsymbol\delta$ carries it at these three leans.

\emph{Trees and bunches.} \hb{thm:mirrortree} with $F=F'=M$ gives $g\big(\operatorname{Tree}(M;\boldsymbol\sigma,\boldsymbol\delta)\big)=\operatorname{Tree}(M;\boldsymbol\sigma,-\boldsymbol\delta)$, each recorded bunch going onto the bunch recorded at the same place. At the root the recorded bunches are $B(M;\boldsymbol\sigma,\boldsymbol\delta)$ and $B(M;\boldsymbol\sigma,-\boldsymbol\delta)$ (\hb{def:tree}), so $g\big(B(M;\boldsymbol\sigma,\boldsymbol\delta)\big)=B(M;\boldsymbol\sigma,-\boldsymbol\delta)$. At the lean $0$ the two trees are one tree, which gives $g\big(\operatorname{Tree}(M;\boldsymbol\sigma,0)\big)=\operatorname{Tree}(M;\boldsymbol\sigma,0)$ and $g\big(B(M;\boldsymbol\sigma,0)\big)=B(M;\boldsymbol\sigma,0)$.

\emph{Classes of the centre.} \hb{cor:exchangecumulants}, applied to the family $F|M$.

\emph{A carrying map.} Let $g$ carry $M$. \hb{prop:certify} at the tables' gauges, where $g^\ast\log\one=\log\one$, says that the reading of $g(e)$ equals that of $e$ at $\Upsilon_{\boldsymbol\sigma,\boldsymbol\delta;\one}$, for every mark $e$. So $g(e)$ lies in the class of $e$, that is $g(C)\subseteq C$ for every class $C$, and $g(C)=C$ since $g$ is injective and $C$ is finite. The last sentence of \hb{thm:mirrortree}, at each of the three leans at which $g$ carries $M$, maps each of the trees at $0$, $\boldsymbol\delta$ and $-\boldsymbol\delta$ onto itself.
\end{proof}

\begin{hlemma}[\hb{lem:windowturn} (An exchange turns a window over)]
As stated: an exchange of one family with another, table by table through the bijection $h$ of the tables.
\end{hlemma}

\begin{proof}
Let $F$ be on $\cE$ and $F'$ on $\cE'$. By \hb{def:carry} the map $T\mapsto\cE'\setminus g(T)$ is a bijection from the bases of table $j$ of $F$ onto those of table $h(j)$ of $F'$. For a basis $T$, $g$ being a bijection,
\[
\big|(\cE'\setminus g(T))\cap g(S)\big|=|g(S)|-|g(T)\cap g(S)|=|S|-|T\cap S| .
\]
By \hb{thm:window} the numbers $|T\cap S|$, $T$ a basis of table $j$, fill $\{d_j(S),\dots,r_j(S)\}$. So the numbers $|T'\cap S'|$, $T'$ a basis of table $h(j)$, fill $\{|S|-r_j(S),\dots,|S|-d_j(S)\}$, a window of the same width $c_j(S)=r_j(S)-d_j(S)$, and the bijection takes the bases of table $j$ at level $k$ on $S$ to those of table $h(j)$ at level $|S|-k$ on $S'$. This is the reference, proof of \hb{lem:windowturn}, with $h$.
\end{proof}

\begin{hproposition}[\hb{prop:levelform} (Level form)]
As stated: table by table in a family, and for the family through the bijections that carry it with the identity on its tables.
\end{hproposition}

\begin{proof}
Fix a table $j$ and a transposition $\tau$ of two marks of $B$ fixing every other mark. A basis $T$ with $T\cap B=p$ is $p\sqcup X$ with $X$ a basis of the child at $p$, of weight $\theta_{j,(p),X}=\theta_{j,T}$, and $K_{j,(p)}(\alpha_R)=\sum_X\theta_{j,(p),X}\alpha_R^X$ (\hb{def:patterns}, \hb{thm:taylorcube}). Distinct bases of one child give distinct monomials, so $K_{j,(p)}$ determines the child at $p$ with its weights, and $K_{j,(p)}=0$ when $p$ is not a pattern.

Let $K_{j,(p)}=K_{j,(p')}$ whenever $|p|=|p'|$. For a basis $T=p\sqcup X$, $\tau(T)=\tau(p)\sqcup X$, since $\tau$ fixes $R$, and $|\tau(p)|=|p|$. So $X$ is a basis of the child at $\tau(p)$ with the weight $\theta_{j,(p),X}$: $\tau(T)$ is a basis of table $j$ and $\theta_{j,\tau(T)}=\theta_{j,T}$. Hence $\tau$ maps the bases of table $j$ into themselves with their weights, and onto them, being injective on a finite set.

Conversely, let every such transposition map the bases of table $j$ onto themselves with their weights, and let $p,p'\subseteq B$ have $|p|=|p'|$. The transpositions of marks of $B$ generate the permutations of $B$, so a product $\varsigma$ of them has $\varsigma(p)=p'$; it maps the bases onto themselves with their weights and fixes $R$. So $p\sqcup X\mapsto p'\sqcup X$ is a bijection from the bases with trace $p$ onto those with trace $p'$ keeping the weights, and $K_{j,(p')}=K_{j,(p)}$. This is the reference, proof of \hb{prop:levelform}, for table $j$, where \hb{lem:descent}(i) is used with the identity on $R$.

$B$ is in level form in the family when it is in level form in every table, that is, by the above, when every such transposition maps the bases of every table onto themselves with their weights. That is the condition of \hb{def:carry} with $F'=F$ and $h$ the identity, whose condition on the pairs, $(\sigma_{h(j)},\delta_{h(j)})=(\sigma_j,\delta_j)$, then holds by itself.
\end{proof}

\begin{hdefinition}[\hb{def:keeppot} (Keeping the potential; twins; descents; copies)]
As stated: keeping and taking the potential at every realization; twins, descents and copies at the table's gauge.
\end{hdefinition}

\begin{hlemma}[\hb{lem:certifypot} (The writings certify the potential)]
As stated, at every realization $\Upsilon=((s_j,t_j)_j;\alpha)$.
\end{hlemma}

\begin{proof}
Let $M$ have $n$ marks, and let $(\Phi_{\mathbf E},\Phi_{\mathbf O})$ be its potential. For a mark $x$ and an order $w=(w_1,\dots,w_{n-1})$ of its room, the potential written from $x$ in the order $w$ is the pair of functions
\[
\Phi^{\mathbf X}_{x,w}(v_0,\dots,v_{n-1})=\Phi_{\mathbf X}\Big(\log\alpha+v_0\one_x+\sum_{i=1}^{n-1}v_i\one_{w_i}\Big)
\]
(\hb{def:view}). For a bijection $g$ and a mark $a$, $g^\ast\one_{g(a)}=\one_a$, since $(g^\ast\one_{g(a)})_e=(\one_{g(a)})_{g(e)}=\one[e=a]$.

(i) By \hb{prop:onewriting} there are orders $(a_1,\dots,a_{n-1})$ of the room of $e$ and $(b_1,\dots,b_{n-1})$ of the room of $f$ with $\Phi^{\mathbf X}_{e,(a_i)_i}=\Phi^{\mathbf X}_{f,(b_i)_i}$ for both parts. Put $a_0=e$ and $b_0=f$, and let $\beta$ be the bijection with $\beta(a_i)=b_i$ for $0\le i\le n-1$; so $\beta(e)=f$. For $y\in\mathbb R^M$ put $v_i=y_{b_i}$. Then $\sum_iv_i\one_{b_i}=y$ and $\sum_iv_i\one_{a_i}=\sum_iy_{\beta(a_i)}\one_{a_i}=\beta^\ast y$, and the equality of the two writings at $(v_0,\dots,v_{n-1})$ reads
\[
\Phi_{\mathbf X}(\log\alpha+\beta^\ast y)=\Phi_{\mathbf X}(\log\alpha+y)
\]
for every $y$ and both parts: $\beta$ keeps the potential (\hb{def:keeppot}).

(ii) Let $g$ keep the potential. For a mark $x$, an order $w$ of its room and $v\in\mathbb R^n$, put $y=v_0\one_{g(x)}+\sum_iv_i\one_{g(w_i)}$; then $g^\ast y=v_0\one_x+\sum_iv_i\one_{w_i}$, and
\[
\Phi^{\mathbf X}_{x,w}(v)=\Phi_{\mathbf X}(\log\alpha+g^\ast y)=\Phi_{\mathbf X}(\log\alpha+y)=\Phi^{\mathbf X}_{g(x),g(w)}(v),\qquad g(w)=(g(w_1),\dots,g(w_{n-1})).
\]
As $w$ runs over the orders of the room of $x$, $g(w)$ runs over those of the room of $g(x)$, once each. So the readings of $x$ and of $g(x)$, the multisets of these pairs of functions, are equal, and $g(x)$ lies in the class of $x$. Thus $g(C)\subseteq C$ for every class $C$, and $g(C)=C$, $g$ being injective and $C$ finite.

Let $g:M\to M'$ take the potential of $M$, read at $\Upsilon$, to that of $M'$, read at a realization with the gauge $\alpha'$. The same computation gives $\Phi^{\mathbf X}_{x,w}(v)=\Phi'^{\mathbf X}_{g(x),g(w)}(v)+c_{\mathbf X}$: the reading of $x$ in $M$ is the reading of $g(x)$ in $M'$ with every function shifted by the constants $c_{\mathbf X}$. So two marks of $M$ have equal readings exactly when their images have equal readings in $M'$, and $g$ maps the classes of $M$ onto those of $M'$.

The centre of $\Upsilon=\Upsilon_{\boldsymbol\sigma,\boldsymbol\delta;\alpha}$ is $\Upsilon_{\boldsymbol\sigma,0;\alpha}$, read at the same gauge (\hb{def:view}). Its potential is $(\Phi_{\mathbf E},0)$, $\Phi_{\mathbf E}=\sum_j\sigma_j\log\mathbf E_j(e^y)$ being the same function at every lean (\hb{def:potential}). A bijection keeping the potential at $\Upsilon$ keeps it at the centre, and one taking the potential of $M$ to that of $M'$ at their realizations takes it at their centres, with the constants $(c_{\mathbf E},0)$; so the two preceding paragraphs hold at the centre as well.

For the last sentence: if $e$ and $f$ read alike, (i) gives a bijection keeping the potential that takes $e$ to $f$; conversely, a bijection keeping the potential maps the class of $e$ onto itself by (ii), so if it takes $e$ to $f$, then $f$ lies in that class.
\end{proof}

\begin{sremark}
At the table's gauge $\log\alpha=0$, the two conditions of \hb{def:keeppot} read $\Phi_{\mathbf X}\circ g^\ast=\Phi_{\mathbf X}$ and $\Phi_{\mathbf X}\circ g^\ast=\Phi'_{\mathbf X}+c_{\mathbf X}$, and the proof above is the reference, proof of \hb{lem:certifypot}. The statements that keep or take the potential at the table's gauge, \hb{thm:roomkeeps}, \hb{prop:treepot}, \hb{thm:leafbounds}, \hb{def:seamproc} and \hb{thm:periodpotential}, and the twins, descents and copies of \hb{def:keeppot}, read these conditions as in the reference.
\end{sremark}

\section{Chapter~\hb{ch:closing}: the integral at a bunch}\label{sec:S1417-ch15}

\begin{hproposition}[\hb{prop:mot-bretadomain} (Where the Breta of one bunch converges)]
As stated.
\end{hproposition}

\begin{proof}
\emph{The set.} If a convex combination of the coordinates of $t\boldsymbol\xi+(1-t)\one$, $0<t\le1$, were a point $p\in(-\infty,0]$, the same combination of the $\xi_\ell$ would be $(p-1+t)/t\in(-\infty,0]$; at $t=0$ the combination is $1$. So $\Xi_J$ is star-shaped about $\one$, hence connected. The convex hull of the $\xi_\ell$ is compact and avoids the closed set $(-\infty,0]$, so it lies at a distance $\delta>0$ from it, and moving every $\xi_\ell$ by less than $\delta$ moves every convex combination by less than $\delta$: $\Xi_J$ is open. Convex combinations of points with positive real parts have positive real parts.

\emph{The integral.} Put $\nu=|\mathbf w|-\kappa$, so $\re\nu>0$, and $z(u)=\sum_\ell\xi_\ell u_\ell$, which for $u$ in the closed simplex lies in the convex hull of the $\xi_\ell$. On a compact neighbourhood of a point of $\Xi_J$ inside $\Xi_J$ these hulls stay at a distance at least some $\delta'>0$ from $(-\infty,0]$, so the principal branch $z^{-\nu}=\exp(-\nu\log z)$ is holomorphic in $\boldsymbol\xi$ and $|z^{-\nu}|=|z|^{-\re\nu}e^{\im\nu\arg z}\le\delta'^{-\re\nu}e^{\pi|\im\nu|}$ there. Since $u^{\mathbf w-\one}$ is integrable on $\Delta_J$ for $\re w_\ell>0$ (\hb{thm:dirichlet}), the integral converges absolutely, locally uniformly, and is holomorphic in $\boldsymbol\xi$ on $\Xi_J$. At positive coefficients it is the Breta: with the unit configuration on one bunch the Zamma is $\Gamma(\nu)^{-1}\int_0^\infty x^{\nu-1}e^{-x\sum_\ell\xi_\ell u_\ell}dx=z(u)^{-\nu}$ (\hb{def:gramma}, \hb{def:zamma}).

\emph{The identity.} The left side of \hb{prop:scalebreta} is holomorphic on $(\mathbb C\setminus(-\infty,0])^J\supseteq\Xi_J$, and the stacked Gamma does not depend on $\boldsymbol\xi$; the two sides agree on $(0,\infty)^J$. A holomorphic function on the connected open set $\Xi_J$ that vanishes on the nonempty open set $(0,\infty)^J$ of its real points vanishes identically, its Taylor coefficients at a real point being its real derivatives there; so they agree on $\Xi_J$.
\end{proof}

\begin{hcorollary}[\hb{cor:stage} (The scale of a bunch of width one)]
As stated: every bunch, table by table, with one Skeleton value in each table of width one.
\end{hcorollary}

\begin{proof}
Fix a table $j$ and write $d=d_j(B)$, $d^\ast=d^\ast_j(B)$ and $c=c_j(B)$. By \hb{thm:window} the levels $|T\cap B|$ fill $\{d,\dots,d+c\}$, and the colevels $|B\setminus T|$ fill $\{|B|-d-c,\dots,|B|-d\}=\{d^\ast,\dots,d^\ast+c\}$, since $d^\ast=|B|-r_j(B)=|B|-d-c$ (\hb{def:skeleton}). Each level letter $K^{(k)}_{j,B}$ and each colevel letter $\Psi^{(k)}_{j,B}$ in these windows is a sum of positive terms over a nonempty set of bases, positive on the orthant. By \hb{def:levels},
\[
x^{d}\sQ^K_{j,B}(x;\alpha)=\sum_{k=d}^{d+c}x^k\frac{K^{(k)}_{j,B}(\alpha)}{K_j(\alpha)},\qquad x^{d^\ast}\sQ^\Psi_{j,B}(x;\alpha)=\sum_{k=d^\ast}^{d^\ast+c}x^k\frac{\Psi^{(k)}_{j,B}(\alpha)}{\Psi_j(\alpha)} .
\]
If $c=0$ both Skeleton polynomials are the constant $1$, of degree $0$, and the spectra are empty (\hb{def:skelpoly}). If $c=1$, then $K_j=K^{(d)}_{j,B}+K^{(d+1)}_{j,B}$, and with $\xi=K^{(d)}_{j,B}(\alpha)/K^{(d+1)}_{j,B}(\alpha)>0$,
\[
\sQ^K_{j,B}(x;\alpha)=\frac{K^{(d)}_{j,B}(\alpha)+xK^{(d+1)}_{j,B}(\alpha)}{K^{(d)}_{j,B}(\alpha)+K^{(d+1)}_{j,B}(\alpha)}=\frac{x+\xi}{1+\xi} .
\]
So the Skeleton value is $\xi$, the level letter at the floor over the level letter at the largest level, and the coin is $1/(1+\xi)=K^{(d+1)}_{j,B}(\alpha)/K_j(\alpha)$, the probability under $\mathbb P_{j,\alpha}(T)=\theta_{j,T}\alpha^T/K_j(\alpha)$ of the bases at level $d+1$ (\hb{thm:innerMGF}), that is $\mathbb P_{j,\alpha}(N_B=d+1)$. The same computation with the colevel letters gives the Skeleton value $\Psi^{(d^\ast)}_{j,B}(\alpha)/\Psi^{(d^\ast+1)}_{j,B}(\alpha)$ and the coin $\Psi^{(d^\ast+1)}_{j,B}(\alpha)/\Psi_j(\alpha)=\mathbb P^\ast_{j,\alpha}(N^\ast_B=d^\ast+1)$, since $\mathbb P^\ast_{j,\alpha}(T)=\theta_{j,T}\alpha^{\cE\setminus T}/\Psi_j(\alpha)$.

Let $B=\{e\}$ with $e$ neither a loop nor a coloop of table $j$. Then $d=0$ and $r_j(\{e\})=1$, so $c=1$ and $d^\ast=0$. The bases at level $0$ avoid $e$ and those at level $1$ contain it, so $K^{(0)}_{j,\{e\}}=K_{j,\setminus e}(\alpha_R)$ and $K^{(1)}_{j,\{e\}}=\alpha_eK_{j,/e}(\alpha_R)$; the bases at colevel $0$ contain $e$ and those at colevel $1$ avoid it, so $\Psi^{(0)}_{j,\{e\}}=\Psi_{j,/e}(\alpha_R)$ and $\Psi^{(1)}_{j,\{e\}}=\alpha_e\Psi_{j,\setminus e}(\alpha_R)$ (\hb{def:delcon}). The coins are $\alpha_eK_{j,/e}/K_j=\mathbb P_{j,\alpha}(e\in T)=R_{j,e}(\alpha)$ and $\alpha_e\Psi_{j,\setminus e}/\Psi_j=\mathbb P^\ast_{j,\alpha}(e\notin T)=R^\ast_{j,e}(\alpha)$ (\hb{thm:innerMGF}). This is the reference, proof of \hb{cor:stage}.
\end{proof}

\begin{htheorem}[\hb{thm:gaps} (The scale in its spectral gaps)]
As stated: complex coefficients and letter exponents, expanded about any point $r$ of the axis with $|\xi_\ell-r|<r$ for every $\ell$.
\end{htheorem}

\begin{proof}
The left side is the value of the scale integral $\int_0^\infty x^{\kappa-1}\prod_{\ell\in J}(x+\xi_\ell)^{-w_\ell}dx$, which converges since $0<\re\kappa<\re|\mathbf w|$ (\hb{prop:scalebreta}). For $x\ge0$ put $q_\ell(x)=(r-\xi_\ell)/(x+r)$, so that $x+\xi_\ell=(x+r)(1-q_\ell(x))$ and $|q_\ell(x)|\le|r-\xi_\ell|/r=|g_\ell|<1$.

\emph{Expansion.} Since $x+r>0$ and $\re(1-q_\ell(x))>0$, the principal argument of $x+\xi_\ell$ is that of $1-q_\ell(x)$, which lies in $(-\pi/2,\pi/2)$; so $(x+\xi_\ell)^{-w_\ell}=(x+r)^{-w_\ell}(1-q_\ell(x))^{-w_\ell}$ with principal powers, and the binomial series $(1-q)^{-w}=\sum_{n\ge0}\frac{(w)_n}{n!}q^n$ converges absolutely for $|q|<1$. Hence
\[
x^{\kappa-1}\prod_\ell(x+\xi_\ell)^{-w_\ell}=\sum_{\mathbf n\in\mathbb N^J}\prod_\ell\frac{(w_\ell)_{n_\ell}}{n_\ell!}(r-\xi_\ell)^{n_\ell}\;x^{\kappa-1}(x+r)^{-|\mathbf w|-|\mathbf n|} .
\]

\emph{Term by term.} $|(w)_n|=\prod_{i<n}|w+i|\le\prod_{i<n}(|w|+i)=(|w|)_n$, and $(x+r)^{-|\mathbf n|}\le r^{-|\mathbf n|}$. So the moduli of the terms sum to at most
\[
x^{\re\kappa-1}(x+r)^{-\re|\mathbf w|}\prod_\ell\sum_{n\ge0}\frac{(|w_\ell|)_n}{n!}|g_\ell|^n=x^{\re\kappa-1}(x+r)^{-\re|\mathbf w|}\prod_\ell(1-|g_\ell|)^{-|w_\ell|},
\]
which is integrable on $(0,\infty)$ since $0<\re\kappa<\re|\mathbf w|$; so the series may be integrated term by term. For $0<\re\kappa<\re a$ the integral of a linear letter, in the paragraph after \hb{prop:scalebreta} with $c_0=r$ and $c_1=1$, is
\[
\int_0^\infty x^{\kappa-1}(x+r)^{-a}dx=\Gst21{\kappa,\ a-\kappa}{a}r^{\kappa-a},
\]
and at $a=|\mathbf w|+|\mathbf n|$, with $(r-\xi_\ell)^{n_\ell}r^{-n_\ell}=g_\ell^{n_\ell}$, the term of $\mathbf n$ is the term of the statement.

\emph{Root rate.} Let $A_M$ be the sum of the terms with $|\mathbf n|=M$, $v=\sum_\ell|w_\ell|$ and $\bar g=\max_\ell|g_\ell|$. Comparing the coefficients of $z^M$ in $\prod_\ell(1-z)^{-|w_\ell|}=(1-z)^{-v}$ gives $\sum_{|\mathbf n|=M}\prod_\ell\frac{(|w_\ell|)_{n_\ell}}{n_\ell!}=\frac{(v)_M}{M!}$. The Beta corner of order $M$ is the integral $\int_0^1x^{\kappa-1}(1-x)^{|\mathbf w|+M-\kappa-1}dx$ (\hb{thm:dirichlet}), of modulus at most $b_0=\int_0^1x^{\re\kappa-1}(1-x)^{\re|\mathbf w|-\re\kappa-1}dx<\infty$. So $|A_M|\le|r^{\kappa-|\mathbf w|}|\,b_0\frac{(v)_M}{M!}\bar g^M$, and $\frac{(v)_M}{M!}=\prod_{i=1}^M\big(1+\frac{v-1}i\big)\le e^{|v-1|(1+\log M)}$ grows at most polynomially in $M$. Hence $\limsup_M|A_M|^{1/M}\le\bar g$ (\hb{def:conv}(viii)).

\emph{The remaining sentences.} A coefficient $\xi_\ell=r$ has $g_\ell=0$, and every term with $n_\ell>0$ vanishes. For $0<\xi_\ell\le r$, $w_\ell>0$ and $0<\kappa<|\mathbf w|$ every factor of every term is positive, a factor $g_\ell^{0}$ being $1$. For positive coefficients write $\xi_{\min}$ and $\xi_{\max}$ for the least and the largest; $|1-\xi/r|$ is largest over $\xi\in[\xi_{\min},\xi_{\max}]$ at an end, so $\max_\ell|g_\ell|=\max\big(|1-\xi_{\min}/r|,|1-\xi_{\max}/r|\big)$. For $r\le\xi_{\min}$ this is $\xi_{\max}/r-1$, decreasing in $r$; for $r\ge\xi_{\max}$ it is $1-\xi_{\min}/r$, increasing in $r$; for $\xi_{\min}\le r\le\xi_{\max}$ it is the larger of $1-\xi_{\min}/r$, increasing, and $\xi_{\max}/r-1$, decreasing, which are equal at $r=\frac12(\xi_{\min}+\xi_{\max})$, with the value $\frac{\xi_{\max}-\xi_{\min}}{\xi_{\max}+\xi_{\min}}$. At that $r$, $|\xi_\ell-r|\le\frac12(\xi_{\max}-\xi_{\min})<r$ for every $\ell$. At $r=\xi_{\max}$ the bound is $1-\xi_{\min}/\xi_{\max}$, the largest coefficient drops out, and the statement is the reference, proof of \hb{thm:gaps}.
\end{proof}

\begin{hcorollary}[\hb{cor:cubeclose} (The period closes on the ordering sectors)]
As stated: complex data with $\kappa_{\cE}=0$ and every proper wall of positive real part.
\end{hcorollary}

\begin{proof}
\emph{Sector form.} Let $\Omega$ be the set of data at which every proper wall has positive real part. It is open and convex, the walls being affine in the data with real coefficients, and it contains points at which $\re\kappa_{\cE}>0$ as well (raise every mass). By \hb{thm:exist}(ii), valid for all complex exponents, and \hb{thm:projective}, wherever every wall has positive real part
\[
P_F(\mvec;\mathbf s,\mathbf t)=\sum_\pi I_\pi,\qquad I_\pi=\int_{[0,1]^{n-1}}\prod_{k=2}^nq_k^{\kappa_{U_k}-1}h_\pi(q)\,dq,\quad h_\pi=L_\pi^{-\kappa_{\cE}}\prod_jk_{j,\pi}^{-s_j}\psi_{j,\pi}^{-t_j}.
\]
The sector letters are polynomials with nonnegative coefficients and positive constant terms, and $L_\pi\in[1,n]$ on the cube, so $h_\pi$ is smooth on a neighbourhood of the cube, with every $q$-derivative continuous in $q$ and the data and entire in the data, and $|h_\pi|$ is bounded on the cube locally uniformly in the data. The tails $U_k$, $k\ge2$, of the orders are the nonempty proper sets of marks. So every $I_\pi$ converges and is holomorphic on $\Omega$, and so is $P_F$ (\hb{thm:projective}); $\Omega$ being connected, $P_F=\sum_\pi I_\pi$ on $\Omega$. At $\kappa_{\cE}=0$, $h_\pi=\prod_jk_{j,\pi}^{-s_j}\psi_{j,\pi}^{-t_j}$.

\emph{The identity and the range.} In the sector coordinates $\alpha^T=\rho^{|T|}\prod_kq_k^{N_{U_k}(T)}$, so $k_{j,\pi}(q)=\sum_T\theta_{j,T}\prod_kq_k^{N_{U_k}(T)-d_j(U_k)}$ (\hb{thm:exist}(ii)), with exponents in $\mathbb N$, all zero only at $T=T^+_{j,\pi}$ (\hb{lem:greedy}). So $k_{j,\pi}/K_j(\one)$ increases in each $q_k$ on the cube, from $\theta_{j,T^+_{j,\pi}}/K_j(\one)$ at $q=0$ to $1$ at $q=\one$, and $0\le1-k_{j,\pi}/K_j(\one)\le\lambda$. The identity is the expectation under $\mathbb P_{j,\one}$ of $1-\prod_kq_k^{e_k}=\sum_k\prod_{k'<k}q_{k'}^{e_{k'}}(1-q_k^{e_k})$, $e_k=N_{U_k}-d_j(U_k)$, with $1-q^e=(1-q)\sum_{0\le i<e}q^i$. The same holds for $\psi_{j,\pi}$ with $N^\ast$, $d^\ast_j$ and $T^-_{j,\pi}$, and $\Psi_j(\one)=K_j(\one)$. These are statements about the tables alone, as in the reference, proof of \hb{cor:cubeclose}.

\emph{Expansion.} With $a_j=1-k_{j,\pi}/K_j(\one)$ and $b_j=1-\psi_{j,\pi}/K_j(\one)$ in $[0,\lambda]$, $k_{j,\pi}^{-s_j}=K_j(\one)^{-s_j}\sum_{n\ge0}\frac{(s_j)_n}{n!}a_j^n$ and $\psi_{j,\pi}^{-t_j}=K_j(\one)^{-t_j}\sum_{n\ge0}\frac{(t_j)_n}{n!}b_j^n$, with principal powers of positive numbers. By $|(s)_n|\le(|s|)_n$ the moduli of the terms of the product of these series sum to at most $\prod_jK_j(\one)^{-\re(s_j+t_j)}(1-\lambda)^{-|s_j|-|t_j|}$ on the cube, and $\int\prod_kq_k^{\re\kappa_{U_k}-1}dq=\prod_k(\re\kappa_{U_k})^{-1}<\infty$. So the series may be integrated term by term over the cube, which gives the displayed series, convergent.

\emph{Terms.} By the identity, $a_j$ and $b_j$ are sums, with nonnegative coefficients, of products $(1-q_k)q^{\boldsymbol\beta}$, $\boldsymbol\beta\in\mathbb N^{n-1}$; so $\prod_ja_j^{n_j}b_j^{n'_j}$ is a finite sum, with nonnegative coefficients, of products $\prod_k(1-q_k)^{\gamma_k}q_k^{\beta_k}$, and for $\re\kappa>0$, $\int_0^1q^{\kappa+\beta-1}(1-q)^{\gamma}dq=\Gst21{\kappa+\beta,\ \gamma+1}{\kappa+\beta+\gamma+1}$ (\hb{thm:dirichlet}). For real data with $s_j,t_j\ge0$ the coefficients $(s_j)_n/n!$ and $(t_j)_n/n!$ and the integrands are nonnegative.

\emph{Bound.} The terms with $|\mathbf n|+|\mathbf n'|=M$ have moduli summing to at most
\[
\prod_jK_j(\one)^{-\re(s_j+t_j)}\sum_{|\mathbf n|+|\mathbf n'|=M}\prod_j\frac{(|s_j|)_{n_j}(|t_j|)_{n'_j}}{n_j!\,n'_j!}\,\lambda^M\sum_\pi\prod_k(\re\kappa_{U_k})^{-1},
\]
and comparing coefficients of $z^M$ in $\prod_j(1-z)^{-|s_j|}(1-z)^{-|t_j|}=(1-z)^{-w}$ gives $\frac{(w)_M}{M!}$ for the inner sum, $w=\sum_j(|s_j|+|t_j|)$.
\end{proof}

\begin{hproposition}[\hb{prop:sectorcompiler} (An enclosure of the period in the field of its data)]
As stated: weights, masses and exponents in a subfield $k$ of $\mathbb R$, the exponents of any sign.
\end{hproposition}

\begin{proof}
The data are real, so by the hypotheses of \hb{cor:cubeclose} every proper wall is positive, and every $\kappa_{U_k}>0$.

(i) The weights lie in $k$, so the coefficients of the sector letters, sums of weights (\hb{thm:exist}(ii)), and $K_j(\one)$ lie in $k$, and $a_{j,\pi},b_{j,\pi}\in k[q]$. For $s\in k$ and $n\in\mathbb N$, $(s)_n/n!\in k$. So $f_{\pi,M}=\sum_{|\mathbf n|+|\mathbf n'|=M}\prod_j\frac{(s_j)_{n_j}(t_j)_{n'_j}}{n_j!\,n'_j!}a_{j,\pi}^{n_j}b_{j,\pi}^{n'_j}$, the coefficient of $\zeta^M$ in the product of the binomial series, lies in $k[q]$. Each $\kappa_{U_k}$ is an integer combination of masses and exponents and lies in $k$, and $\int_{[0,1]^{n-1}}\prod_kq_k^{\kappa_{U_k}-1+\beta_k}dq=\prod_k(\kappa_{U_k}+\beta_k)^{-1}\in k$ for $\boldsymbol\beta\in\mathbb N^{n-1}$; so $I_{\pi,M}\in k$. Grouping the absolutely convergent series of \hb{cor:cubeclose} by $M=|\mathbf n|+|\mathbf n'|$ gives the identity. When $s_j,t_j\ge0$ the coefficients are nonnegative, and $a_{j,\pi},b_{j,\pi}\ge0$ on the cube, so $I_{\pi,M}\ge0$.

(ii) On the cube $0\le a_{j,\pi},b_{j,\pi}\le\lambda_\pi$ (\hb{cor:cubeclose}), so by $|(s)_n|\le(|s|)_n$ and the coefficient identity of the proof of \hb{cor:cubeclose}, $|f_{\pi,M}|\le\frac{(w)_M}{M!}\lambda_\pi^M$ and $|I_{\pi,M}|\le\prod_k\kappa_{U_k}^{-1}\frac{(w)_M}{M!}\lambda_\pi^M$. The ratio of consecutive bounds is $\frac{w+M}{M+1}\lambda_\pi$, and $M\mapsto\frac{w+M}{M+1}$ is nonincreasing for $w\ge1$ and below $1$ for $w<1$; so for $M\ge Q+1$ the ratio is at most $r_{\pi,Q}<1$, and summing the geometric majorant gives $\big|\sum_{M>Q}I_{\pi,M}\big|\le\operatorname{err}_{\pi,Q}$. The numbers $w$, $\lambda_\pi$, $r_{\pi,Q}$ and $\operatorname{err}_{\pi,Q}$ lie in $k$, an ordered subfield of $\mathbb R$ being closed under $|\cdot|$ and under maxima of finitely many elements. As $Q\to\infty$, $r_{\pi,Q}\to\lambda_\pi<1$, and $\frac{(w)_{Q+1}}{(Q+1)!}\lambda_\pi^{Q+1}\to0$, the factor $\frac{(w)_{Q+1}}{(Q+1)!}$ growing at most polynomially; so $\operatorname{err}_{\pi,Q}\to0$. With $P^{\le Q}=\sum_\pi\sum_{M\le Q}I_{\pi,M}\in k$, the left side of (i) lies between $P^{\le Q}-\sum_\pi\operatorname{err}_{\pi,Q}$ and $P^{\le Q}+\sum_\pi\operatorname{err}_{\pi,Q}$, and for $s_j,t_j\ge0$, the remainders being nonnegative, between $P^{\le Q}$ and $P^{\le Q}+\sum_\pi\operatorname{err}_{\pi,Q}$.

(iii) $\partial_\zeta\log f=\sum_iw_iA_i/(1-\zeta A_i)$; multiply by $f\prod_i(1-\zeta A_i)$. Comparing the coefficients of $\zeta^M$ on the two sides gives a linear recurrence for the coefficients of $f$, with coefficients polynomial in $M$ and in the $A_i$. With rational data this is the reference, proof of \hb{prop:sectorcompiler}.
\end{proof}

\begin{htheorem}[\hb{thm:contcorners} (Continued Beta corners)]
As stated: complex data, every index off the nonpositive integers, several indices of nonpositive real part at once.
\end{htheorem}

\begin{proof}
\emph{A continuation as a linear functional.} For $\kappa\in\mathbb C$ and $h\in\mathbb N$ with $\re\kappa>-h-1$ and $\kappa\notin\{0,-1,\dots,-h\}$, and $f$ smooth on $[0,1]$, \hb{lem:cube} with $m=1$ and $M=h+1$ gives the continuation in $\kappa$ of $\int_0^1x^{\kappa-1}f\,dx$:
\[
\mathcal C(\kappa;f)=\int_0^1x^{\kappa-1}\big(f-\operatorname{Tay}_hf\big)dx+\sum_{l\le h}\frac{f^{(l)}(0)}{l!\,(\kappa+l)},
\]
$\operatorname{Tay}_h$ taking the Taylor polynomial of degree $h$ at $0$; it is the finite part $\Reg\int_0^1x^{\kappa-1}f\,dx$ (\hb{cor:RegMellin}). By Taylor's formula $|f-\operatorname{Tay}_hf|\le\frac{x^{h+1}}{(h+1)!}\sup|f^{(h+1)}|$, so
\[
|\mathcal C(\kappa;f)|\le\frac{\sup|f^{(h+1)}|}{(h+1)!\,(\re\kappa+h+1)}+\sum_{l\le h}\frac{|f^{(l)}(0)|}{l!\,|\kappa+l|}:
\]
$\mathcal C(\kappa;\cdot)$ is linear and bounded by a constant times the largest of the sup-norms of $f,\dots,f^{(h+1)}$. On a cube $[0,1]^m$, with a number $L_k\in\mathbb N$ for each coordinate, $\re x_k+L_k>0$ and $x_k\notin\{0,\dots,1-L_k\}$, \hb{lem:cube} repeated in each coordinate gives the continuation of $\int\prod_kq_k^{x_k-1}f\,dq$ as a finite sum, over the sets $A$ of coordinates and the $\mathbf l$ with $l_k<L_k$ for $k\in A$, of $\prod_{k\in A}(x_k+l_k)^{-1}$ times the integral over the other coordinates of $\prod_{k\notin A}q_k^{x_k-1}$ against $\prod_{k\notin A}(1-\operatorname{Tay}_{L_k-1})$ applied to $\partial^{\mathbf l}_{q_A}f/\mathbf l!$ at $q_A=0$ (as in \hb{cor:pastwalls}), with $\operatorname{Tay}_{-1}=0$. Taylor's formula in each coordinate outside $A$ bounds each term by a constant, depending only on the $x_k$, times the largest sup-norm of the derivatives of $f$ of total order at most $L=\sum_kL_k$; so this continuation is again linear and bounded on $C^L([0,1]^m)$, and it is the iterated finite part, by \hb{cor:RegMellin} in one coordinate at a time.

\emph{Leibniz bounds.} Let finitely many functions $f_i$ be smooth on a cube, with $0\le f_i\le\lambda$, $0<\lambda<1$, and let $D\ge1$ bound their derivatives of order at most $L$. For a multi-index $\mathbf n$ over the $i$ and $f_{\mathbf n}=\prod_if_i^{n_i}$, a derivative of total order $l\le L$ distributes over the $|\mathbf n|$ factors in at most $|\mathbf n|^l$ ways, each product having at most $l$ differentiated factors, each bounded by $D$, and the others bounded by $\lambda$; so $|\partial^{\boldsymbol\beta}f_{\mathbf n}|\le D^l|\mathbf n|^l\lambda^{|\mathbf n|-l}$ for $|\boldsymbol\beta|=l$. Hence the sum, over the $\mathbf n$ with $|\mathbf n|=M$, of $\prod_i\frac{|(w_i)_{n_i}|}{n_i!}$ times the $C^L$-norm of $f_{\mathbf n}$ is at most $\frac{(w)_M}{M!}D^L(1+M)^L\lambda^{M-L}$, $w=\sum_i|w_i|$, by the coefficient identity of the proof of \hb{thm:gaps}. These bounds are summable over $M$, so $\sum_{\mathbf n}\prod_i\frac{(w_i)_{n_i}}{n_i!}f_{\mathbf n}$ converges in $C^L$, to $\prod_i(1-f_i)^{-w_i}$; a bounded linear functional on $C^L$ commutes with the sum, and applied to $f_{\mathbf n}$ it is at most a constant times $D^L(1+|\mathbf n|)^L\lambda^{|\mathbf n|-L}$. If $\lambda=0$ every $f_i$ vanishes and only $\mathbf n=0$ contributes.

\emph{The two sectors of a scale.} In \hb{thm:sectors} take $f_i=1-\sQ_{S,i}(\cdot;\alpha)$ on $[0,1]$: each $\sQ_{S,i}$ has nonnegative coefficients, positive constant term and value $1$ at $x=1$ (\hb{def:skelpoly}), so $0\le f_i\le1-\sQ_{S,i}(0;\alpha)\le\lambda_S$ there, and $\prod_i\sQ_{S,i}^{-w_i}=\sum_{\mathbf n}\prod_i\frac{(w_i)_{n_i}}{n_i!}f_{\mathbf n}$. Let $\kappa$ be the index of the sector, not a nonpositive integer. If $\re\kappa>0$, take $\mathcal C(\kappa;f)=\int_0^1x^{\kappa-1}f\,dx$, bounded by $\sup|f|/\re\kappa$, with $L=0$; otherwise take $h$ as in the statement and $L=h+1$. By the preceding paragraph,
\[
\mathcal C\Big(\kappa;\prod_i\sQ_{S,i}^{-w_i}\Big)=\sum_{\mathbf n}\prod_i\frac{(w_i)_{n_i}}{n_i!}\,\mathcal C(\kappa;f_{\mathbf n}),
\]
the terms of total order $M$ having moduli summing to at most $\frac{(w)_M}{M!}\lambda_S^M$ times a polynomial in $M$ of degree $L$, $w=\sum_j(|s_j|+|t_j|)$. The left side is the finite part of the sector integral and its continuation in $\kappa$. By \hb{def:tail}, $f_{\mathbf n}=(1-x)^{|\mathbf n|}\prod_i\sQ_{S,i,>}^{n_i}$, a polynomial, and $\mathcal C\big(\kappa;x^k(1-x)^{|\mathbf n|}\big)$ is the continuation of $\int_0^1x^{\kappa+k-1}(1-x)^{|\mathbf n|}dx=\Gst21{\kappa+k,\ |\mathbf n|+1}{\kappa+k+|\mathbf n|+1}$ (\hb{thm:dirichlet}), the meromorphic stacked Gamma, finite since $\kappa+k\notin-\mathbb N$. So the series of \hb{thm:sectors}(ii), with these Beta corners, converges to the finite part, with the stated bound. By \hb{thm:sectors}(i) the scale integral is the product of the letters at $\alpha$ and the sum of the two sectors at the indices $\kappa_B$ and $\kappa^{\mathrm{top}}_B$, and its continuation in the data is the same expression with the sectors continued.

\emph{The ordering sectors.} By the first step of the proof of \hb{cor:cubeclose} and \hb{lem:cube}, $P_F=\sum_\pi I_\pi$ as meromorphic functions of the data, each $I_\pi$ continued by \hb{lem:cube}, whose formula expresses the continuation at a point through the values of $h_\pi$ at that point; on the balance wall $h_\pi=\prod_jk_{j,\pi}^{-s_j}\psi_{j,\pi}^{-t_j}$. At a point of the balance wall at which no proper wall is a nonpositive integer, take for an order $\pi$ the numbers $L_k=h_k+1$ when $\re\kappa_{U_k}\le0$, with $-h_k-1<\re\kappa_{U_k}\le-h_k$, and $L_k=0$ otherwise. With $f$ running over the $a_j,b_j$ of the proof of \hb{cor:cubeclose}, in $[0,\lambda]$ on the cube, the preceding paragraphs give the continuation of $I_\pi$ as the series of \hb{cor:cubeclose} for $\pi$, each integral replaced by its continuation; and the continuation of $\int\prod_kq_k^{\kappa_{U_k}+\beta_k-1}(1-q_k)^{\gamma_k}dq$ is the product over $k$ of the continued Beta corners, each finite since $\kappa_{U_k}+\beta_k\notin-\mathbb N$. The terms of total order $M$ have moduli summing to at most $\frac{(w)_M}{M!}\lambda^M$ times a polynomial in $M$ of degree $\sum_kL_k$, and the tails of one order being distinct proper sets, $\sum_kL_k$ is at most the sum of the integers $h+1$ over the proper walls of real part at most $0$. Every proper set of marks is a tail of some order, so the sector series computes $P_F$ at every point of the balance wall off the pole hyperplanes, with no positivity of the walls. For one real index in $(-h-1,-h)$ this is the reference, proof of \hb{thm:contcorners}.
\end{proof}

\begin{hcorollary}[\hb{cor:content} (A loop or coloop in every table)]
As stated: a set of marks, each a loop or a coloop in every table.
\end{hcorollary}

\begin{proof}
For a table $j$ let $p_j$ be the set of the marks of $A$ that are coloops of table $j$. Every basis contains every coloop and avoids every loop, so $T\cap A=p_j$ for every basis $T$ of table $j$. The sets $T\setminus A$ are then the bases of the table on $\cE\setminus A$ obtained from table $j$ by contracting the marks of $p_j$ and deleting those of $A\setminus p_j$, one mark at a time (\hb{def:delcon}; a loop is never a coloop, so each step is a table), with the weights $\theta_{j,T}$; with $K'_j,\Psi'_j$ its Foreman polynomials and $\alpha'=\alpha_{\cE\setminus A}$,
\[
K_j(\alpha)=\alpha_A^{p_j}K'_j(\alpha'),\qquad\Psi_j(\alpha)=\alpha_A^{A\setminus p_j}\Psi'_j(\alpha'),
\]
since $\cE\setminus T=(A\setminus p_j)\sqcup\big((\cE\setminus A)\setminus(T\setminus A)\big)$. So the integrand of $J_F(\mvec;\mathbf s,\mathbf t)$ is the product of
\[
\prod_{e\in A}\alpha_e^{m_e-1-\sum_j(s_j\one[e\in p_j]+t_j\one[e\notin p_j])}e^{-\alpha_e}
\]
and of the integrand of $J_{F'}(\mvec|_{\cE\setminus A};\mathbf s,\mathbf t)$. With $s_j=\sigma_j+\delta_j$, $t_j=\sigma_j-\delta_j$ and $m_e=\mu_e+\sigma$, the exponent of $\alpha_e$ is $\mu_e-\sum_j\eta_{j,e}\delta_j-1$. The integrand is a product of functions of disjoint sets of variables, so by Tonelli on the modulus its integral converges exactly when every factor's does, and then, by Fubini, it is the product of the factors' integrals, with $\int_0^\infty\alpha_e^{\nu-1}e^{-\alpha_e}d\alpha_e=\Gamma(\nu)$ for $\re\nu>0$. As meromorphic functions the identity persists (\hb{thm:exist}(iii)). For one mark this is the reference, proof of \hb{cor:content}, there read through \hb{thm:port} at $B=\{e\}$.
\end{proof}

\section{Chapter~\hb{ch:expectations}: integrals, expectations and the Foreman procedure}\label{sec:S1417-ch16}

\begin{hproposition}[\hb{prop:snip} (Snips are completions)]
As stated: (iii) at every imbalance, through the appended list $\binom B1$ of the singletons of the bunch.
\end{hproposition}

\begin{proof}
Items (i) and (ii) are the reference, proof of \hb{prop:snip}. For (iii) let $R=\cE\setminus B$.

\emph{Projectivity.} The family word $W$ is homogeneous of degree $-\sum_j(s_jr_j+t_jr^\ast_j)=\kappa_{\cE}-m_{\cE}$ (\hb{def:family}), so $\omega=\alpha^{\mvec-\one}W\,d\alpha$ is multiplied by $\lambda^{\kappa_{\cE}}$ under $\alpha\mapsto\lambda\alpha$, and $|\alpha_B|^{-\kappa_{\cE}}\omega$ is unchanged: it is projective (\hb{def:conv}(vii)). On the section $|\alpha_B|=1$ the multiplier is $1$, so the two forms induce the same form there.

\emph{The appended table.} The sets $\{e\}$, $e\in B$, have one element each and the exchange property, two of them differing in one mark; so they are the bases of a table on $\cE$ of rank $1$ (\hb{def:table}), in which the marks of $R$ lie in no basis, and with unit weights its basis polynomial is $\sum_{e\in B}\alpha_e=|\alpha_B|$. So $|\alpha_B|^{-\kappa_{\cE}}W$ is the family word of $F'$ at the exponents $\big((\mathbf s,\kappa_{\cE}),(\mathbf t,0)\big)$, the appended table carrying $(\kappa_{\cE},0)$, and $|\alpha_B|^{-\kappa_{\cE}}\omega$ is the form of $F'$ at the masses $\mvec$.

\emph{Walls.} For the appended table $d_0(S)=\min_{e\in B}|\{e\}\cap S|=\one[B\subseteq S]$ (\hb{thm:window}), and its exponent on $\Psi$ is $0$; so by \hb{def:family} the walls of $F'$ at $\mvec$ are $\kappa_S-\kappa_{\cE}\one[B\subseteq S]$, and its imbalance is $\kappa_{\cE}-\kappa_{\cE}\cdot1=0$, the appended table having rank $1$.

\emph{Value.} $|\alpha_B|$ is positive and homogeneous of degree one on the orthant. Where every proper wall of $F'$ has positive real part, the period of $F'$ converges, and the integral of its form over the section $\{|\alpha_B|=1\}$ is that period (\hb{thm:projective}); this is the display.

\emph{The cases.} For $B=\{g\}$ the appended table has the one basis $\{g\}$ and the letter $\alpha_g$, and $\alpha_g^{-\kappa_{\cE}}W\,d\alpha_{\mvec}=W\,d\alpha_{\mvec-\kappa_{\cE}\mathbf e_g}=\omega_{\mvec^{(g)}}$ (\hb{def:completion}): this is (i). For $B=\cE$ the appended table is the list $\binom{\cE}1$, and the statement is \hb{prop:sectiontable}. For $k=-\kappa_{\cE}\in\mathbb N$, $k\ge1$: on the section $\omega=|\alpha_B|^k\omega$, $|\alpha_B|^k=\sum_{|\mathbf n|=k}\binom k{\mathbf n}\alpha_B^{\mathbf n}$, and $\alpha_B^{\mathbf n}\omega=W\,d\alpha_{\mvec+\mathbf n}$ is the completion along $\mathbf n/k$ (\hb{def:completion}: $\mvec-\kappa_{\cE}\mathbf n/k=\mvec+\mathbf n$), projective, whose integral over the section is $P_F(\mvec+\mathbf n;\mathbf s,\mathbf t)$ wherever that period converges (\hb{thm:projective}). This last case is the reference, proof of \hb{prop:snip}(iii).
\end{proof}

\begin{hcorollary}[\hb{cor:fiber} (The period fibres over the scale of a bunch)]
As stated: every section $\{h(\alpha_R)=1\}$ of the room, and off the balance wall.
\end{hcorollary}

\begin{proof}
For a positive function $h$ on an orthant $\mathbb R^X_{>0}$, homogeneous of degree one, write $\lambda_h$ for the measure on $\Sigma_h=\{\beta:\ h(\beta)=1\}$ with $d\alpha=\rho^{|X|-1}d\rho\,d\lambda_h(\beta)$ under the bijection $(\rho,\beta)\mapsto\rho\beta$ of $(0,\infty)\times\Sigma_h$ onto $\mathbb R^X_{>0}$: the measure of \hb{lem:chart} (the reference, proof of \hb{lem:chart}). For $h=\alpha_g$ it is the Lebesgue measure in the coordinates $\alpha_e$, $e\neq g$, at $\alpha_g=1$, and for $h=|\alpha|$ it is $d\omega$ (\hb{lem:scaleshape}).

Let $h$ be a function of $\alpha_R$ as in the statement, and put $\tilde h(\alpha)=h(\alpha_R)$ on $\mathbb R^{\cE}_{>0}$, positive and homogeneous of degree one. The map $(\rho,\beta_B,\beta_R)\mapsto(\rho\beta_B,\rho\beta_R)$, $\beta_B\in\mathbb R^B_{>0}$, $\beta_R\in\Sigma_h$, is a bijection onto $\mathbb R^{\cE}_{>0}$, with $d\alpha_R=\rho^{|R|-1}d\rho\,d\lambda_h(\beta_R)$ and $d\alpha_B=\rho^{|B|}d\beta_B$ at fixed $\rho$. So $d\alpha=\rho^{n-1}d\rho\,d\beta_B\,d\lambda_h(\beta_R)$, and the measure of the section $\{\tilde h=1\}$ is $d\beta_B\,d\lambda_h(\beta_R)$. The integral of $W\,d\alpha_{\mvec}$ over it is
\[
\int_{\Sigma_h}\beta_R^{\mvec|_R-\one}\int_{\mathbb R^B_{>0}}\beta_B^{\mvec|_B-\one}\prod_jK_j(\beta_B,\beta_R)^{-s_j}\Psi_j(\beta_B,\beta_R)^{-t_j}\,d\beta_B\,d\lambda_h(\beta_R),
\]
by Tonelli on the modulus, which is the same integrand at the real parts of the data, and then Fubini. With $\beta_B=xu$, $x>0$, $u\in\Delta_B$, $\beta_B^{\mvec|_B-\one}d\beta_B=x^{m_B-1}u^{\mvec|_B-\one}dx\,d\omega_B(u)$ (\hb{lem:scaleshape}), so the inner integral is $\int_{\Delta_B}u^{\mvec|_B-\one}\operatorname{Sc}_B(u,\beta_R;\mvec;\mathbf s,\mathbf t)\,d\omega_B(u)$, with the scale integral of \hb{thm:scale} at the gauge $(u,\beta_R)$. This is the display.

The scale integral converges exactly when $\re\kappa_B>0$ and $\re\kappa^{\mathrm{top}}_B>0$ (\hb{thm:scale}), conditions on the data alone; by Tonelli it converges at almost every $(u,\beta_R)$, hence at every point, and there it is the product of \hb{thm:scale}. On the balance wall the form is projective and its integral over every section is the period (\hb{thm:projective}). At $h=\alpha_g$ the integral over $\{\alpha_g=1\}$ is $P_F(\mvec^{(g)};\mathbf s,\mathbf t)$ wherever it converges (\hb{prop:snip}(i)), and $\lambda_h$ is the Lebesgue measure in $\alpha_{R-g}$ at $\alpha_g=1$; at $h=|\alpha_R|$, $\Sigma_h=\Delta_R$ and $\lambda_h=d\omega_R$. At $h=\alpha_g$ on the balance wall this is the reference, proof of \hb{cor:fiber}.
\end{proof}

\begin{hlemma}[\hb{lem:periodmero} (The period is meromorphic in the data)]
As stated, for every family and all complex data.
\end{hlemma}

\begin{proof}
\emph{Meromorphy.} By the first step of the proof of \hb{cor:cubeclose} above, which uses neither connectedness nor real data, $P_F=\sum_\pi I_\pi$ on the open convex set $\Omega$ of the data at which every proper wall has positive real part, with $I_\pi=\int_{[0,1]^{n-1}}\prod_{k=2}^nq_k^{\kappa_{U_k}-1}h_\pi\,dq$ and $h_\pi=L_\pi^{-\kappa_{\cE}}\prod_jk_{j,\pi}^{-s_j}\psi_{j,\pi}^{-t_j}$ smooth near the closed cube, its derivatives continuous in the variables and the data and entire in the data. By \hb{lem:cube} each $I_\pi$ extends to a meromorphic function of the data, a quotient of a holomorphic function by products of affine functions $\kappa_{U_k}+i$, $i\in\mathbb N$, with poles only on the hyperplanes $\kappa_{U_k}=-i$. Every nonempty proper set of marks is a tail of some order, and no tail is $\cE$. So $P_F$ is meromorphic in the masses and the exponents, with poles only where a proper wall is a nonpositive integer.

\emph{The restriction to the balance wall.} For a nonempty proper $U$ the affine function $\kappa_U+i$ of the data has the part $m_U$ in the masses, which is not proportional to the part $m_{\cE}$ of $\kappa_{\cE}$; so the hyperplane $\kappa_U=-i$ is not the balance wall $\{\kappa_{\cE}=0\}$, and it meets the balance wall in a hyperplane of it. Near every point of the balance wall $P_F$ is a quotient of a holomorphic function by a product of such affine functions, none vanishing identically on the balance wall; so the restriction of $P_F$ to the balance wall is meromorphic, with poles only where a proper wall is a nonpositive integer.
\end{proof}

\begin{hproposition}[\hb{prop:compoundlaw} (The form of a compound at one point)]
As stated, for any number $q\ge2$ of parts.
\end{hproposition}

\begin{proof}
\emph{The form.} Write $r_{j,i}$ for the rank of the restriction of table $j$ to $\cE_i$, so that $K_{j,i}$ is homogeneous of degree $r_{j,i}$, $\Psi_{j,i}$ of degree $|\cE_i|-r_{j,i}$, and $\kappa_i=m_{\cE_i}-\sum_j\big(s_jr_{j,i}+t_j(|\cE_i|-r_{j,i})\big)$; summing over $i$, $\sum_i\kappa_i=\kappa_{\cE}=0$. The section $\{|\alpha_q|=1\}$ carries the measure $\prod_{i<q}d\alpha_i\,d\omega_{\cE_q}(u_q)$, by the computation in the proof of \hb{cor:fiber} above with the room there replaced by $\cE_q$. For $i<q$, $\alpha_i=\rho_iu_i$ gives $\alpha_i^{\mvec_i-\one}d\alpha_i=\rho_i^{m_{\cE_i}-1}u_i^{\mvec_i-\one}d\rho_i\,d\omega_{\cE_i}(u_i)$ (\hb{lem:scaleshape}), $K_{j,i}(\rho_iu_i)=\rho_i^{r_{j,i}}K_{j,i}(u_i)$ and $\Psi_{j,i}(\rho_iu_i)=\rho_i^{|\cE_i|-r_{j,i}}\Psi_{j,i}(u_i)$. The powers of $\rho_i$ collect to $\rho_i^{\kappa_i-1}$, which gives the displayed form.

(i) For real $\kappa$, $\int_0^\infty\rho^{\kappa-1}d\rho$ diverges, at $\infty$ if $\kappa\ge0$ and at $0$ if $\kappa\le0$. The form is positive, and by Tonelli its integral is the product of the $q-1$ integrals $\int_0^\infty\rho_i^{\kappa_i-1}d\rho_i$, each infinite, and of the $q$ shape integrals, each positive; so it is infinite.

(ii) A primitive of $\rho^{\kappa-1}$ is $\rho^\kappa/\kappa$, or $\log\rho$ at $\kappa=0$. It is its own log-power expansion at $0$ and at $\infty$, with no term $\rho^0\log^0\rho$, so both regularized limits are $0$ and the finite part is $0$ (\hb{def:Reg}).

(iii) At $\mvec(\boldsymbol\eta)$, $\kappa_i=\kappa_{F_i}(\mvec^0_i)+i\eta_i|w_i|=i\eta_i$, and $\sum_i\kappa_i=0$ since $\sum_i\eta_i=0$. The walls of $F_i$ at $\mvec^0_i+i\eta_iw_i$ have the real parts of those at $\mvec^0_i$, positive, so $f(\boldsymbol\eta')=\prod_{i=1}^qP_{F_i}(\mvec^0_i+i\eta_iw_i;\mathbf s,\mathbf t)$, with $\boldsymbol\eta'=(\eta_1,\dots,\eta_{q-1})$ and $\eta_q=-\sum_{i<q}\eta_i$, is a smooth function of $\boldsymbol\eta'$, each period being holomorphic in the masses there (\hb{thm:projective}). In $y_i=\log\rho_i$, $\rho_i^{\kappa_i-1}d\rho_i=e^{i\eta_iy_i}dy_i$. For a test function $H$ of $\boldsymbol\eta'$ of compact support, the integral of the form against $H$, taken over the shapes, then over $\boldsymbol\eta'$, then over $\mathbf y=(y_1,\dots,y_{q-1})$, is
\[
\int_{\mathbb R^{q-1}}\int_{\mathbb R^{q-1}}H(\boldsymbol\eta')f(\boldsymbol\eta')\,e^{i\ip{\boldsymbol\eta'}{\mathbf y}}\,d\boldsymbol\eta'\,d\mathbf y=(2\pi)^{q-1}H(0)f(0),
\]
by Fourier inversion for the smooth function $Hf$ of compact support, whose Fourier transform is integrable; and $f(0)=\prod_iP_{F_i}(\mvec^0_i;\mathbf s,\mathbf t)$. For $q=2$ this is the reference, proof of \hb{prop:compoundlaw}.
\end{proof}

\begin{slemma}\label{lem:ch16-compound}
Let $S$ be a nonempty proper set of marks of a family $F$ on $\cE$ with $c_j(S)=0$ in every active table, and put $S'=\cE\setminus S$.
\begin{enumerate}[label=(\roman*)]
\item In every active table $j$, every basis meets $S$ in $r_j(S)$ marks and $S'$ in $r_j(S')$ marks; the bases of $\cB_j|S$ and of $\cB_j|S'$ are the sets $T\cap S$ and $T\cap S'$, $T\in\cB_j$; and $X\sqcup Y\in\cB_j$ for every basis $X$ of $\cB_j|S$ and every basis $Y$ of $\cB_j|S'$.
\item In this lemma write
\[
I_S(\mvec;\mathbf s,\mathbf t)=\int_{\Delta_S\times\Delta_{S'}}\alpha^{\mvec-\one}\prod_jK_j(\alpha)^{-s_j}\Psi_j(\alpha)^{-t_j}\,d\omega_S\,d\omega_{S'} .
\]
It converges, and is holomorphic in the data, wherever $\re\kappa_U>0$ for every nonempty proper subset $U$ of $S$ and every nonempty proper subset $U$ of $S'$; and it extends to a meromorphic function of the data whose poles lie on the hyperplanes $\kappa_U=-i$, $i\in\mathbb N$, $U$ a nonempty proper subset of $S$ or of $S'$.
\item Wherever every proper wall has positive real part, $\kappa_{\cE}=\kappa_S+\kappa_{S'}$ has positive real part and $J_F(\mvec;\mathbf s,\mathbf t)=\Gamma(\kappa_S)\Gamma(\kappa_{S'})\,I_S(\mvec;\mathbf s,\mathbf t)$.
\end{enumerate}
\end{slemma}

\begin{proof}
(i) Fix an active table $j$. By \hb{thm:window} the numbers $|T\cap S|$ fill $\{d_j(S),\dots,r_j(S)\}$, a single value since $c_j(S)=0$; so every basis meets $S$ in $r_j(S)$ marks and $S'$ in $r_j-r_j(S)$ marks, a constant, which is then $r_j(S')$. By \hb{def:minors} the bases of $\cB_j|S$ are the sets $T\cap S$ with $|T\cap S|=r_j(S)$, that is all of them, and likewise for $S'$. Let $X=T_1\cap S$ and $Y=T_2\cap S'$ for bases $T_1,T_2$; induct on $k=|(T_2\cap S)\setminus X|$. If $k=0$, then $T_2\cap S\subseteq X$, both of size $r_j(S)$, so $T_2=X\sqcup Y$. If $k>0$, take $a\in(T_2\cap S)\setminus X$; then $a\notin T_1$. The exchange property (\hb{def:table}) gives $b\in T_1\setminus T_2$ with $T_3=(T_2\setminus\{a\})\cup\{b\}$ a basis. As $|T_3\cap S|=r_j(S)=|T_2\cap S|$ and $a\in S$, $b\in S$, so $b\in X\setminus T_2$. Then $T_3\cap S'=Y$ and $|(T_3\cap S)\setminus X|=k-1$, and by induction $X\sqcup Y$ is a basis.

\emph{Homogeneity in each block.} By (i), for an active table $j$, $K_j(\rho\alpha_S,\rho'\alpha_{S'})=\rho^{r_j(S)}\rho'^{r_j(S')}K_j(\alpha)$ and $\Psi_j(\rho\alpha_S,\rho'\alpha_{S'})=\rho^{|S|-r_j(S)}\rho'^{|S'|-r_j(S')}\Psi_j(\alpha)$. For $U\in\{S,S'\}$, $c_j(U)=0$ (\hb{thm:window}: $c_j(S')=c_j(S)$), so $d_j(U)=r_j(U)$ and $d^\ast_j(U)=|U|-r_j(U)$ (\hb{def:skeleton}). An inactive table enters with the exponents $0$. So with $g(U)=\sum_j\big(s_jd_j(U)+t_jd^\ast_j(U)\big)=m_U-\kappa_U$, the function $\phi(\alpha)=\alpha^{\mvec-\one}\prod_jK_j(\alpha)^{-s_j}\Psi_j(\alpha)^{-t_j}$ has
\[
\phi(\rho\alpha_S,\rho'\alpha_{S'})=\rho^{\kappa_S-|S|}\rho'^{\kappa_{S'}-|S'|}\phi(\alpha),
\]
and $|\phi|$ is the same function at the real parts of the data.

(iii) With $\alpha_S=\rho u$ and $\alpha_{S'}=\rho'u'$ (\hb{lem:scaleshape} in each block), $d\alpha=\rho^{|S|-1}\rho'^{|S'|-1}d\rho\,d\rho'\,d\omega_S(u)\,d\omega_{S'}(u')$, and the integrand of $J_F$ becomes $\rho^{\kappa_S-1}e^{-\rho}\rho'^{\kappa_{S'}-1}e^{-\rho'}\phi(u,u')$. Wherever every proper wall has positive real part, $\re\kappa_{\cE}=\re\kappa_S+\re\kappa_{S'}>0$, by \hb{prop:compoundwall}, $c_j(S)=0$ holding in every table with $\sigma_j\neq0$, these being active. Then every wall has positive real part, and $J_F$ at the real parts of the data is finite (\hb{thm:exist}(i)); by Tonelli it is $\Gamma(\re\kappa_S)\Gamma(\re\kappa_{S'})$ times the integral of $|\phi|$ over $\Delta_S\times\Delta_{S'}$, which is therefore finite, and Fubini gives $J_F=\Gamma(\kappa_S)\Gamma(\kappa_{S'})I_S$.

(ii) \emph{Charts.} For a finite set $X$, a mark $x\in X$ and a function $f$ on $\mathbb R^X_{>0}$, homogeneous of degree $-|X|$ and nonnegative or integrable on $\Delta_X$, $\int_{\Delta_X}f\,d\omega_X=\int f|_{\beta_x=1}\,d\beta_{X-x}$, and the same holds on the cones over subsets of $\Delta_X$ and their traces on the chart: for a nonnegative function $\upsilon$ on $(0,\infty)$ with $\int_0^\infty\upsilon(\rho)\,d\rho/\rho=1$, both sides equal $\int_{\mathbb R^X_{>0}}f(\alpha)\upsilon(|\alpha|)\,d\alpha$, by \hb{lem:scaleshape} on the left, and on the right by the change of variables $\alpha=\rho(1,\beta)$, $d\alpha=\rho^{|X|-1}d\rho\,d\beta$, with $\int_0^\infty\upsilon(\rho|(1,\beta)|)\,d\rho/\rho=1$. The function $\psi(\alpha)=|\alpha_S|^{-\kappa_S}|\alpha_{S'}|^{-\kappa_{S'}}\phi(\alpha)$ is homogeneous of degree $-|S|$ in $\alpha_S$ and $-|S'|$ in $\alpha_{S'}$, and equal to $\phi$ on $\Delta_S\times\Delta_{S'}$.

For an order $\pi$ of $S$, with tails $U_k=\{\pi(k),\dots,\pi(|S|)\}$, $2\le k\le|S|$, use on the chart $\beta_{\pi(1)}=1$ the coordinates $\beta_{\pi(k)}=q_2\cdots q_k$, $q\in(0,1]^{|S|-1}$, which cover the trace of the sector $\{\beta_{\pi(1)}\ge\cdots\ge\beta_{\pi(|S|)}\}$; the exponent of $q_k$ in $\beta^X$ is $|X\cap U_k|$, $d\beta_{S-\pi(1)}=\prod_kq_k^{|S|-k}dq$, $\beta^{\mvec|_S-\one}d\beta_{S-\pi(1)}=\prod_kq_k^{m_{U_k}-1}dq$, and $|\beta|=L_\pi(q)=\sum_{i=1}^{|S|}\prod_{k=2}^iq_k\in[1,|S|]$. Likewise $\pi'$, $U'_l$, $q'$ and $L_{\pi'}$ on $S'$. Decomposing each simplex into its sectors, and applying the identity above in each block, by Tonelli on the modulus and then Fubini,
\[
I_S=\sum_{\pi,\pi'}\int\prod_kq_k^{m_{U_k}-1}\prod_lq'^{\,m_{U'_l}-1}_l\,L_\pi(q)^{-\kappa_S}L_{\pi'}(q')^{-\kappa_{S'}}\prod_jK_j(\beta,\beta')^{-s_j}\Psi_j(\beta,\beta')^{-t_j}\,dq\,dq',
\]
the integral over $(0,1]^{|S|-1}\times(0,1]^{|S'|-1}$.

\emph{The letters on a product of sectors.} Fix an active table $j$. By (i) and \hb{def:minors}, $\cB_j|S$ is a table on $S$, and by \hb{lem:greedy} for $\pi$ it has bases $X^\pm$ with $|X^+\cap U_k|=\min_X|X\cap U_k|$ and $|U_k\setminus X^-|=\min_X|U_k\setminus X|$, the minima over its bases; by (i) these minima are $\min_T|T\cap U_k|=d_j(U_k)$ and $\min_T|U_k\setminus T|=d^\ast_j(U_k)$ (\hb{thm:window}). Likewise $\cB_j|S'$ has bases $Y^\pm$ for $\pi'$. For a basis $T$, $\beta^{T\cap S}/\beta^{X^+}=\prod_kq_k^{|T\cap U_k|-d_j(U_k)}$ with exponents in $\mathbb N$, and likewise on $S'$. So
\begin{gather*}
K_j(\beta,\beta')=\prod_kq_k^{d_j(U_k)}\prod_lq'^{\,d_j(U'_l)}_l\,k_{j,\pi,\pi'}(q,q'),\\ k_{j,\pi,\pi'}=\sum_T\theta_{j,T}\prod_kq_k^{|T\cap U_k|-d_j(U_k)}\prod_lq'^{\,|T\cap U'_l|-d_j(U'_l)}_l,
\end{gather*}
a polynomial with nonnegative coefficients. Its constant term is the sum of $\theta_{j,T}$ over the bases with $|T\cap U_k|=d_j(U_k)$ for $k\ge2$ and $|T\cap U'_l|=d_j(U'_l)$ for $l\ge2$. As $|T\cap S|=|X^+|$, these counts determine $T\cap S=X^+$, the mark $\pi(k)$ lying in $T$ exactly when $|T\cap U_k|-|T\cap U_{k+1}|=1$, with $U_1=S$ and $U_{|S|+1}=\varnothing$; likewise $T\cap S'=Y^+$. By (i) $X^+\sqcup Y^+$ is a basis, so the constant term is $\theta_{j,X^+\sqcup Y^+}>0$. In the same way $\Psi_j(\beta,\beta')=\prod_kq_k^{d^\ast_j(U_k)}\prod_lq'^{\,d^\ast_j(U'_l)}_l\,\psi_{j,\pi,\pi'}(q,q')$, with nonnegative coefficients and the constant term $\theta_{j,X^-\sqcup Y^-}>0$. As $m_U-g(U)=\kappa_U$,
\[
I_S=\sum_{\pi,\pi'}\int\prod_{k=2}^{|S|}q_k^{\kappa_{U_k}-1}\prod_{l=2}^{|S'|}q'^{\,\kappa_{U'_l}-1}_l\,h_{\pi,\pi'}(q,q')\,dq\,dq',\qquad h_{\pi,\pi'}=L_\pi^{-\kappa_S}L_{\pi'}^{-\kappa_{S'}}\prod_jk_{j,\pi,\pi'}^{-s_j}\psi_{j,\pi,\pi'}^{-t_j},
\]
the product over the active tables.

\emph{Continuation.} The polynomials $k_{j,\pi,\pi'}$, $\psi_{j,\pi,\pi'}$, $L_\pi$ and $L_{\pi'}$ are at least their positive constant terms on the closed cube, hence positive on a neighbourhood of it; so $h_{\pi,\pi'}$ is smooth there, with every derivative continuous in the variables and the data and entire in the data, and $|h_{\pi,\pi'}|$, which is $h_{\pi,\pi'}$ at the real parts of the data, lies between positive constants on the cube, locally uniformly in the data. So the term of $(\pi,\pi')$ converges, and is holomorphic, exactly where $\re\kappa_{U_k}>0$ and $\re\kappa_{U'_l}>0$ for all $k,l$; the tails $U_k$, $k\ge2$, of the orders of $S$ are the nonempty proper subsets of $S$, and likewise for $S'$. This gives the convergence and the holomorphy of $I_S$. By \hb{lem:cube}, each term is a quotient of a holomorphic function of the data by products of the affine functions $\kappa_{U_k}+i$ and $\kappa_{U'_l}+i$, $i\in\mathbb N$, each a nonconstant function of the data (the coefficient of $m_e$, $e\in U$, being $1$); so it extends meromorphically with poles only on these hyperplanes, and so does $I_S$.
\end{proof}

\begin{htheorem}[\hb{thm:offdomain} (The pseudointegral off the domain)]
As stated: the compound formula, with the value $0$ on the balance wall, for every family and all complex data; the clause on targets, (i) and (ii) under the hypotheses of the reference.
\end{htheorem}

\begin{proof}
\emph{The compound formula.} Let $S$ be nonempty and proper with $c_j(S)=0$ in every active table, $S'=\cE\setminus S$, and $I_S$ the integral of Lemma~\ref{lem:ch16-compound}, the integral of the statement. On the open convex set $\Omega$ of the data at which every proper wall has positive real part, $\re\kappa_{\cE}>0$ and $J_F=\Gamma(\kappa_S)\Gamma(\kappa_{S'})I_S$ (Lemma~\ref{lem:ch16-compound}(iii)), and $J_F=\Gamma(\kappa_{\cE})P_F$ (\hb{thm:projective}); so $P_F=\Gst21{\kappa_S,\ \kappa_{S'}}{\kappa_{\cE}}I_S$ on $\Omega$, with $\kappa_S+\kappa_{S'}=\kappa_{\cE}$ (\hb{prop:compoundwall}). Both sides are meromorphic in the data, by \hb{lem:periodmero}, by Lemma~\ref{lem:ch16-compound}(ii) and by \hb{def:stacked}, and they agree on the nonempty open set $\Omega$; so they agree as meromorphic functions.

\emph{The value $0$ on the balance wall.} Put $G=\Gamma(\kappa_S)\Gamma(\kappa_{S'})I_S$, meromorphic in the data. The poles of $\Gamma(\kappa_S)\Gamma(\kappa_{S'})$ lie on the hyperplanes $\kappa_S=-i$ and $\kappa_{S'}=-i$, $i\in\mathbb N$, and those of $I_S$ on the hyperplanes $\kappa_U=-i$ with $U$ a nonempty proper subset of $S$ or of $S'$ (Lemma~\ref{lem:ch16-compound}(ii)). In each case $U$ is a nonempty proper set of marks, so, as in the proof of \hb{lem:periodmero}, none of these hyperplanes contains the balance wall, and $G$ restricts to a meromorphic function on the balance wall, holomorphic off a proper analytic subset of it. The function $1/\Gamma$ is entire and vanishes at $0$, so $1/\Gamma(\kappa_{\cE})$ is entire in the data and vanishes on the balance wall. Hence $P_F=G\cdot\Gamma(\kappa_{\cE})^{-1}$ vanishes at every point of the balance wall at which $G$ is holomorphic, and the meromorphic restriction of $P_F$ to the balance wall is $0$.

\emph{The clause on targets, (i) and (ii)} are under the hypotheses of the reference, connected with real data and $s_j,t_j\ge0$ not all zero: the reference, proof of \hb{thm:offdomain}.
\end{proof}

\begin{hdefinition}[\hb{def:perioddatum} (Period datum; moves)]
As stated, for tables with positive rational weights.
\end{hdefinition}

\begin{proof}
The definition asks of the family only that its form be algebraic over $\mathbb Q$ and homogeneous of degree zero. The letters $K_j$ and $\Psi_j$ are polynomials whose coefficients are the weights, here rational; with rational exponents and masses $W\,d\alpha_{\mvec}$ is a product of rational powers of polynomials over $\mathbb Q$ and of the coordinates, times $d\alpha$, hence algebraic over $\mathbb Q$, and on the balance wall its degree is $\kappa_{\cE}=0$ (\hb{def:family}). The walls of the faces are those of \hb{thm:exist}, for any weights.
\end{proof}

\begin{hproposition}[\hb{prop:sectionsoff} (Sections off the balance wall)]
As stated: a bunch section at every imbalance.
\end{hproposition}

\begin{proof}
The clause on a bunch section is \hb{prop:snip}(iii), proved above; the rest is the reference, proof of \hb{prop:sectionsoff}. In (ii) the walls of the chart completion at $g$ are those of \hb{prop:snip}(iii) at $B=\{g\}$, $\kappa_S-\kappa_{\cE}\one_S(g)$.
\end{proof}


\begin{hdefinition}[\hb{def:pseudoint} (The pseudointegral over a region)]
As stated; the last sentence of (i) is stated for families of tables with positive rational weights, where the reference takes a list.
\end{hdefinition}

\begin{proof}
Let the tables have positive rational weights and the data be rational, on the balance wall. The letters $K_j$ and $\Psi_j$ are polynomials whose coefficients are the weights, so the form $W\,d\alpha_{\mvec}$ is algebraic over $\mathbb Q$ and the family has a period datum (\hb{def:perioddatum}), whose moves are defined over $\mathbb Q$ by that definition. The three regions named in (i) are cut out by polynomial inequalities over $\mathbb Q$: the chart $\alpha_g=1$ by $\alpha_e>0$ for $e\neq g$; the simplex, in the coordinates $(u_e)_{e\neq e_0}$, by $u_e>0$ and $1-\sum_{e\neq e_0}u_e>0$; a bunch section $|\alpha_B|=1$ by these inequalities on $B$ together with $\alpha_e>0$ on the room.

A change of variables of \hb{def:perioddatum} keeps this property. Let $\mathcal R$ lie in the open domain and be cut out there by $p_1>0,\dots,p_k>0$ with $p_i\in\mathbb Q[x]$, and let $\phi$ be birational over $\mathbb Q$ and carry the open domain bijectively onto the open domain. Each coordinate of $\phi^{-1}$ is a quotient of polynomials over $\mathbb Q$ whose denominator does not vanish on the open domain, where $\phi^{-1}$ is defined; so $p_i\circ\phi^{-1}=A_i/D_i$ with $A_i,D_i\in\mathbb Q[y]$, $D_i$ a product of these denominators, and
\[
\phi(\mathcal R)=\{y\ \text{in the open domain}:\ A_i(y)\,D_i(y)>0\ \text{for}\ 1\le i\le k\},
\]
the open domain being an orthant, cut out by the positivity of its coordinates. The region of a product of data is the product of the regions, cut out by the inequalities of the two factors in disjoint coordinates; linear combinations and boundary terms are formed on the regions already reached, as for a list. Of the family the argument of the reference uses only that its period datum and its moves are defined over $\mathbb Q$, which holds for positive rational weights.
\end{proof}

\section{Chapter~\hb{ch:questions}: the family as one Laplace transform}\label{sec:ch38}

\subsection{What the one-table proofs read, and what a family reads}\label{sec:famread}

The one-table statements read the exponent of the law of a question, $\sigma E+\delta O$ at offsets zero, through the width $\wdt$ of the Newton polytope of $\Lambda$ and through the convexity of $E$; the walls of the wedge, its slope and the volume at low strength are all values of the width. For a family the exponent is
\[
X(z)=2\ip{\bmu_\perp}z-\sum_j\big(\sigma_jE_j(z)+\delta_jO_j(z)\big),
\]
and what replaces the width is the wall at a tilt (Definition~\hb{def:tiltwall}): on the sector of every order the exponent is minus twice the wall at the tilt less the logarithms of the sector letters (Proposition~\hb{prop:evenfloor}(i)), which is Theorem~\hb{thm:exist}(ii) read in the signed coordinates. Every statement below is read off that identity, the convexity of the centred potential, and, for samples of bases, the factorization of the chance of a sample over its draws. Throughout, $n=|\cE|\ge2$, the data are real and on the balance wall unless said otherwise, a question is written $q=(\bmu_\perp;\bsig,\bdel)$, its data are $\big(\mvec_F(\bmu_\perp;\bsig,\bdel);\bsig+\bdel,\bsig-\bdel\big)$, and $P_F(q)$ denotes its period. On the sector $H_\pi$ of an order $\pi$, read in the tilts (Definition~\hb{def:tiltwall}), the differences of a tilt along $\pi$ are written $v_k=z_{\pi(k-1)}-z_{\pi(k)}\ge0$, $2\le k\le n$, so that $\kappa(z)=\sum_k\kappa_{U_k}v_k$; at $\alpha=e^{2z}$ the gaps of Corollary~\hb{cor:pastwalls} are $2v_k$.

\begin{slemma}[The walls on the balance wall]\label{lem:walls}
On the balance wall, for every nonempty proper $S$,
\[
\kappa_S=\kappa^{\mathbf E}_S(\bsig)+\kappa^{\mathbf O}_S(\bmu_\perp,\bdel),\qquad \kappa^{\mathbf E}_S(\bsig)=\sum_j\sigma_jc_j(S),\qquad \kappa^{\mathbf O}_S(\bmu_\perp,\bdel)=\mu_{\perp,S}-\sum_j\delta_j\varepsilon^\perp_j(S),
\]
with $\varepsilon^\perp_j(S)=\varepsilon_j(S)-|S|\frac{2r_j-n}n$; and $\kappa^{\mathbf E}_{\cE\setminus S}=\kappa^{\mathbf E}_S$, $\kappa^{\mathbf O}_{\cE\setminus S}=-\kappa^{\mathbf O}_S$, so $\kappa_{\cE\setminus S}=\kappa^{\mathbf E}_S-\kappa^{\mathbf O}_S$. In particular every wall is a linear function of the pair coordinates $(\bmu_\perp;\bsig,\bdel)$ of the question.
\end{slemma}

\begin{proof}
By Theorem~\hb{thm:pairform}(iii), $\kappa_S=\mu_S+\sum_j\big(\sigma_jc_j(S)-\delta_j\varepsilon_j(S)\big)$ and $\kappa_{\cE}=\mu_{\cE}-\sum_j\delta_j(2r_j-n)$. On the balance wall $\mu_{\cE}=\sum_j\delta_j(2r_j-n)$, and $\bmu=\bmu_\perp+\frac{\mu_{\cE}}n\one$ gives $\mu_S=\mu_{\perp,S}+\frac{|S|}n\sum_j\delta_j(2r_j-n)$; substituting gives the decomposition. By Theorem~\hb{thm:window}, $c_j(\cE\setminus S)=c_j(S)$ and $\varepsilon_j(\cE\setminus S)=2r_j-n-\varepsilon_j(S)$, so $\varepsilon^\perp_j(\cE\setminus S)=2r_j-n-\varepsilon_j(S)-\frac{n-|S|}n(2r_j-n)=-\varepsilon^\perp_j(S)$, and $\mu_{\perp,\cE\setminus S}=-\mu_{\perp,S}$ as $\bmu_\perp\in\one^\perp$.
\end{proof}

\subsection{The exponent of a question}

\begin{hdefinition}[\hb{def:tiltwall}: the claims]
Let $z\in H_\pi$. A difference $v_k$ is positive exactly when $z_{\pi(k-1)}>z_{\pi(k)}$, and then $U_k=\{e:\ z_e<z_{\pi(k-1)}\}$, a set $\{e:\ z_e<z_f\}$ with $z_f>\min_ez_e$; so the tails with a positive difference, and their differences, are read off $z$ alone, and $\kappa(z)$ does not depend on the order $\pi$ whose sector holds $z$. On $H_\pi$ the differences are linear in $z$, so $\kappa$ is linear there, and $\kappa(cz)=c\kappa(z)$ for $c\ge0$. The tilt $-\one_S$ lies in the sectors of the orders that list $\cE\setminus S$ first; there its only positive difference is $1$, at the tail $U_k=S$, so $\kappa(-\one_S)=\kappa_S$. In the chart $z_{\pi(1)}=0$ the coordinates are $z_{\pi(i)}=-(v_2+\cdots+v_i)$, a triangular change of coordinates with determinant $\pm1$, so $dz=dv_2\cdots dv_n$ on $H_\pi$; the chart does not matter (Definition~\hb{def:question}). Finally $\kappa(z)$ is linear in the walls, which are linear in the question by Lemma~\ref{lem:walls}.
\end{hdefinition}

\begin{hproposition}[\hb{prop:evenfloor} (The exponent of a question reads its walls; the even statistic is least at the table's gauge)]
As stated, for a question with real exponents of any signs.
\end{hproposition}

\begin{proof}
(i) Write $\bmu=\mvec-\sigma\one$ for the offsets, so that $\bmu-\bmu_\perp\in\mathbb R\one$ and $\mu_{\cE}=\sum_j\delta_j(2r_j-n)$ on the balance wall. The function $\tilde X(z)=2\ip{\bmu}z-\sum_j\big(\sigma_jE_j(z)+\delta_jO_j(z)\big)$ on $\mathbb R^{\cE}$ is unchanged by $z\mapsto z+c\one$, since $E_j$ is unchanged and $O_j$ moves by $2c(2r_j-n)$ (Theorem~\hb{thm:tworeadings}(iii)), the move of the whole being $2c\kappa_{\cE}=0$ (Proposition~\hb{prop:scaling}); on $\one^\perp$ it is $X$. So $X$ may be read at any representative of a tilt. At $\alpha=e^{2z}$, Theorem~\hb{thm:tworeadings}(i) gives $\sigma_jE_j+\delta_jO_j=s_j\log\Lambda_j(z)+t_j\log\Lambda_j(-z)=s_j\log K_j(\alpha)+t_j\log\Psi_j(\alpha)-(s_j+t_j)\ip\one z$, and $2\ip{\bmu}z+2\sigma\ip\one z=\ip{\mvec}{\log\alpha}$, so
\[
X(z)=\log\Big(\alpha^{\mvec}\prod_jK_j(\alpha)^{-s_j}\Psi_j(\alpha)^{-t_j}\Big),
\]
the logarithm of the integrand of the period against the measure $d\alpha/\alpha^{\one}$. Take the sector coordinates of Theorem~\hb{thm:exist}(ii) for $\pi$: $\rho=\alpha_{\pi(1)}=e^{2z_{\pi(1)}}$ and $q_k=\alpha_{\pi(k)}/\alpha_{\pi(k-1)}=e^{-2v_k}$. Since $\alpha_{\pi(i)}=\rho q_2\cdots q_i$ and $\pi(i)\in U_k$ exactly when $k\le i$, $\alpha^{\mvec}=\rho^{m_{\cE}}\prod_kq_k^{m_{U_k}}$; with $K_j=\rho^{r_j}\prod_kq_k^{d_j(U_k)}k_{j,\pi}(q)$ and $\Psi_j=\rho^{n-r_j}\prod_kq_k^{d^\ast_j(U_k)}\psi_{j,\pi}(q)$,
\[
\alpha^{\mvec}\prod_jK_j^{-s_j}\Psi_j^{-t_j}=\rho^{\kappa_{\cE}}\prod_{k=2}^nq_k^{\kappa_{U_k}}\prod_jk_{j,\pi}(q)^{-s_j}\psi_{j,\pi}(q)^{-t_j},
\]
and $\rho^{\kappa_{\cE}}=1$ on the balance wall. The logarithm, with $\log q_k=-2v_k$, is the identity of (i).

On $H_\pi$ the greedy basis $T^+_{j,\pi}$ has the largest monomial $\alpha^T$ among the bases of table $j$, and $T^-_{j,\pi}$ the largest $\alpha^{\cE\setminus T}$ (Lemma~\hb{lem:greedy}); as $-\chi_T=\chi_{\cE\setminus T}$, $h_j(z)=\ip{\chi_{T^+_{j,\pi}}}z$ and $h_j(-z)=\ip{\chi_{\cE\setminus T^-_{j,\pi}}}z$, whatever the signs of the exponents. Write $z=z_{\pi(1)}\one-\sum_kv_k\one_{U_k}$; for $A\subseteq\cE$, $\ip{\chi_A}z=z_{\pi(1)}(2|A|-n)-\sum_kv_k(2|A\cap U_k|-|U_k|)$, and $|T^+_{j,\pi}\cap U_k|=d_j(U_k)$, $|U_k\setminus T^-_{j,\pi}|=d^\ast_j(U_k)$. So
\[
\sum_j\big(s_jh_j(z)+t_jh_j(-z)\big)=2z_{\pi(1)}\sum_j\delta_j(2r_j-n)+2\sum_kv_k\Big(\sigma|U_k|-\sum_j\big(s_jd_j(U_k)+t_jd^\ast_j(U_k)\big)\Big),
\]
while $2\ip{\bmu}z=2z_{\pi(1)}\mu_{\cE}-2\sum_kv_k\mu_{U_k}$. The difference is $2\sum_kv_k\big(m_{U_k}-\sum_j(s_jd_j(U_k)+t_jd^\ast_j(U_k))\big)=2\kappa(z)$, the term in $z_{\pi(1)}$ vanishing on the balance wall; on $\one^\perp$, $\ip{\bmu}z=\ip{\bmu_\perp}z$. For $s_j,t_j\ge0$: the support functions of $P_j$ and $P^\ast_j$ are $\frac12(h_j(z)+\ip\one z)$ and $\frac12(h_j(-z)+\ip\one z)$, support functions add over a Minkowski sum with nonnegative coefficients, and so $2\kappa(z)=2h_{\mathcal N_F}(z)-2\sigma\ip\one z-2\ip{\bmu}z=2h_{\mathcal N_F}(z)-2\ip{\mvec}z$, twice the support function of $\mathcal N_F-\mvec$, which lies in $\one^\perp$ and so reads tilts.

The sector letters are polynomials in $q$ with nonnegative coefficients and the constant terms $\theta_{j,T^+_{j,\pi}}$, $\theta_{j,T^-_{j,\pi}}$, and their value at $q=\one$ is $K_j(\one)=\Psi_j(\one)$; on $[0,1]^{n-1}$ they increase in every $q_k$ and so lie between these (Theorem~\hb{thm:exist}(ii)), whatever the signs of the exponents. Taking logarithms and weighting by $s_j$ and $t_j$ gives both bounds.

(ii) If every proper wall is positive, on $H_\pi$, $\kappa(z)=\sum_k\kappa_{U_k}v_k\ge\kappa_{\min}\sum_kv_k=\kappa_{\min}(z_{\pi(1)}-z_{\pi(n)})$, and $\max_ez_e-\min_ez_e$ vanishes on $\cT$ only at $0$. Conversely $\kappa(-\one_S)=\kappa_S$.

(iii) At the centre $\bmu_\perp=0$ and $\bdel=0$, so $X=-\sum_j\sigma_jE_j$, and the walls are $\kappa_S=\sum_j\sigma_jc_j(S)$ (Proposition~\hb{prop:centrefamily}(ii)). By (i), $2\kappa=\sum_j\sigma_j(h_j(z)+h_j(-z))=\sum_j\sigma_j\wdt_j(z)$, and $\wdt_j(z)=\max_T\ip{\chi_T}z-\min_T\ip{\chi_T}z=2\big(\max_T\ip{\one_T}z-\min_T\ip{\one_T}z\big)$ is twice the width of $P_j$ in the direction $z$, equal to $2(r_j(S)-d_j(S))=2c_j(S)$ at $z=\one_S$ (Theorem~\hb{thm:window}). The bounds of (i) at $s_j=t_j=\sigma_j$ are the displayed ones. The Hessian of $\log\Lambda_j$ at $z$ is the covariance of $\chi_T$ under $\mathbb P_{j,e^{2z}}$ (Theorem~\hb{thm:innerMGF}), so $\log\Lambda_j$ is convex, and so is $E_j(z)=\log\Lambda_j(z)+\log\Lambda_j(-z)$, which is even; hence $E_j(z)=\frac12(E_j(z)+E_j(-z))\ge E_j(0)=2\log\Lambda_j(0)=2\log K_j(\one)$.

For $v\in\one^\perp$,
\[
v^{\mathsf T}\nabla^2\Big(\sum_j\sigma_jE_j\Big)(z)\,v=\sum_j\sigma_j\Big(\operatorname{Var}_{T\sim\mathbb P_{j,e^{2z}}}\ip{\chi_T}v+\operatorname{Var}_{T\sim\mathbb P_{j,e^{-2z}}}\ip{\chi_T}v\Big).
\]
Every basis has positive chance under both laws, so the form vanishes only if $\ip{\chi_T}v=2\ip{\one_T}v-\ip\one v$ is constant over the bases of every table with $\sigma_j>0$. Then $v$ is constant on every molecule of the family of these tables (Lemma~\hb{lem:circuit}), and every such molecule $M$ has $c_j(M)=0$ in each of them, the table being the direct sum of its restrictions to the molecules (Proposition~\hb{prop:directsum}), so that $|T\cap M|=r_j(M)$ for every basis. If $v\notin\mathbb R\one$, the union $S$ of the molecules on which $v$ is largest is nonempty and proper with $c_j(S)=0$ for every table with $\sigma_j>0$: the family is not connected. So on a connected family the form is positive on $\cT\setminus0$ and $\sum_j\sigma_jE_j$ is strictly convex on $\cT$. Finally $\sum_j\sigma_j\wdt_j=2\kappa$ is positive on $\cT\setminus0$ by (ii), the walls $\sum_j\sigma_jc_j(S)$ of the centre being positive exactly when the family is connected (Proposition~\hb{prop:centrefamily}(ii)).
\end{proof}

\begin{sremark}
For one table at the centre of strength $\sigma$, (iii) is the statement of the reference, with $2\kappa=\sigma\wdt$; the family statement adds (i), in which the width of one table is replaced by the wall at the tilt of any question, and (ii). The identity of (i) is the sector form of the period read pointwise: integrated against $dz=dv$ on each sector it is $P_F(q)=2^{n-1}\sum_\pi\int_{[0,\infty)^{n-1}}e^{-2\sum_k\kappa_{U_k}v_k}h_\pi(e^{-2v})\,dv$, with $h_\pi=\prod_jk_{j,\pi}^{-s_j}\psi_{j,\pi}^{-t_j}$, which is $\sum_\pi I_\pi$ of Corollary~\hb{cor:pastwalls} in the gaps $2v_k$.
\end{sremark}

\subsection{The two ends of strength}

The high end needs one fact about integrals of $e^{-\lambda f}$ with a convex exponent, for complex $\lambda$ in a half-plane; it is proved here once, in the form in which the period uses it.

\begin{slemma}[A convex exponent at high strength]\label{lem:highstrength}
Let $f$ be a smooth function on a real vector space $V$ of dimension $d\ge1$ with Lebesgue measure $dz$ and a norm $|\cdot|$, such that
\begin{enumerate}[label=(\alph*)]
\item every derivative of $f$ of order at least one is bounded;
\item $f$ is convex, with a critical point $z^\star$ at which $A_0=\nabla^2f(z^\star)$ is positive definite;
\item $f(z)-f(z^\star)\ge c|z|-C$ for some $c>0$ and $C$.
\end{enumerate}
Then $z^\star$ is the only critical point, and as $|\lambda|\to\infty$ with $\re\lambda\ge\lambda_0$, for a fixed $\lambda_0>0$,
\[
\int_Ve^{-\lambda f(z)}\,dz=e^{-\lambda f(z^\star)}\Big(\frac{2\pi}\lambda\Big)^{\frac d2}\det(A_0)^{-\frac12}\big(1+O(|\lambda|^{-1})\big),
\]
the determinant in the coordinates of $dz$ and the power principal.
\end{slemma}

\begin{proof}
Use a Euclidean norm, equivalent to the given one. \emph{The gradient away from $z^\star$.} By convexity $\ip{\nabla f(z)}{z-z^\star}\ge f(z)-f(z^\star)\ge c|z|-C$, so $|\nabla f(z)|\ge(c|z|-C)/|z-z^\star|$, which tends to $c$ as $|z|\to\infty$. A convex function with a positive definite Hessian at a critical point has its strict minimum there, and its critical points form the convex set of its minima, which is $\{z^\star\}$. So for every neighbourhood $V_0$ of $z^\star$, $|\nabla f|\ge c_0>0$ on $V\setminus V_0$: on a large ball by compactness, outside it by the bound.

\emph{A chart at $z^\star$.} Put $A(z)=2\int_0^1(1-t)\,\nabla^2f\big(z^\star+t(z-z^\star)\big)\,dt$, smooth and symmetric with $A(z^\star)=A_0$. Since $\nabla f(z^\star)=0$, $f(z)=f(z^\star)+\frac12\ip{A(z)(z-z^\star)}{z-z^\star}$. On a ball $V_1$ about $z^\star$ on which $A(z)$ is positive definite put $w(z)=A(z)^{1/2}(z-z^\star)$, the positive square root being smooth there; then $f=f(z^\star)+\frac12|w|^2$, and $Dw(z^\star)=A_0^{1/2}$ is invertible, so $w$ is a diffeomorphism of a smaller ball $V_2$ onto a neighbourhood of $0$.

\emph{Near $z^\star$.} Take a smooth $\chi$ with compact support in $V_2$ and $\chi=1$ near $z^\star$. In the coordinates $w$, $\int\chi e^{-\lambda f}dz=e^{-\lambda f(z^\star)}\int_{\mathbb R^d}e^{-\lambda|w|^2/2}g(w)\,dw$, with $g$ smooth of compact support and $g(0)=|\det Dw(z^\star)|^{-1}=\det(A_0)^{-1/2}$. For $\re\lambda>0$ and $\xi\in\mathbb R^d$, $\int e^{-\lambda|w|^2/2+i\ip\xi w}dw=(2\pi/\lambda)^{d/2}e^{-|\xi|^2/(2\lambda)}$: both sides are holomorphic in $\lambda$ on the half-plane and agree for $\lambda>0$, completing the square in each coordinate. Writing $g(w)=(2\pi)^{-d}\int\hat g(\xi)e^{i\ip\xi w}d\xi$ with $\hat g$ rapidly decreasing,
\[
\int e^{-\lambda|w|^2/2}g(w)\,dw=\Big(\frac{2\pi}\lambda\Big)^{\frac d2}(2\pi)^{-d}\int\hat g(\xi)\,e^{-|\xi|^2/(2\lambda)}\,d\xi=\Big(\frac{2\pi}\lambda\Big)^{\frac d2}\big(g(0)+\varrho(\lambda)\big),
\]
and since $\re(1/\lambda)\ge0$ and $|e^{-x}-1|\le|x|$ for $\re x\ge0$, $|\varrho(\lambda)|\le(2|\lambda|)^{-1}(2\pi)^{-d}\int|\xi|^2|\hat g(\xi)|\,d\xi$.

\emph{Away from $z^\star$.} On the support of $1-\chi$, $|\nabla f|\ge c_0$. The operator $L=-\lambda^{-1}|\nabla f|^{-2}\ip{\nabla f}{\nabla}$ has $Le^{-\lambda f}=e^{-\lambda f}$, and its transpose is $L^{\mathsf t}\varphi=\lambda^{-1}\operatorname{div}\big(\varphi\,|\nabla f|^{-2}\nabla f\big)$. Integrating by parts $N$ times, $\int(1-\chi)e^{-\lambda f}dz=\int e^{-\lambda f}(L^{\mathsf t})^N(1-\chi)\,dz$, with no boundary terms: $1-\chi$ vanishes near $z^\star$, and on large spheres $|e^{-\lambda f}|=e^{-\re\lambda\,f}$ decays exponentially by (c) while the coefficients stay bounded. Now $(L^{\mathsf t})^N(1-\chi)=\lambda^{-N}G_N$ with $G_N$ a polynomial in the derivatives of $\chi$, in those of $f$ of orders $1$ to $N+1$ and in $|\nabla f|^{-2}$, bounded by (a) and $|\nabla f|\ge c_0$. Since $f\ge f(z^\star)$ and $\re\lambda\ge\lambda_0$,
\[
\Big|\int(1-\chi)e^{-\lambda f}dz\Big|\le|\lambda|^{-N}\sup|G_N|\;e^{-\re\lambda\,f(z^\star)}\int_Ve^{-\lambda_0(f-f(z^\star))}dz ,
\]
the last integral finite by (c). With $N>\frac d2+1$ this is $O(|\lambda|^{-1})$ times the modulus $(2\pi/|\lambda|)^{d/2}e^{-\re\lambda f(z^\star)}\det(A_0)^{-1/2}$ of the leading term, and the two parts add to the statement. (On the real half-line this is Laplace's method, along the imaginary direction stationary phase, and $w$ is the Morse chart.)
\end{proof}

\begin{htheorem}[\hb{thm:severity} (The two ends of strength)]
As stated.
\end{htheorem}

\begin{proof}
The walls and the exponents of $\lambda q$ are $\lambda$ times those of $q$ (Lemma~\ref{lem:walls}), so its walls $\lambda\kappa_S$ have positive real parts for $\re\lambda>0$. At such data the period is the convergent integral of Theorem~\hb{thm:projective}, the powers of the positive letters principal, and the change of variables of Definition~\hb{def:question}, which uses only $K_j(e^{2z})=e^{\ip\one z}\Lambda_j(z)$ and $\Psi_j(e^{2z})=e^{\ip\one z}\Lambda_j(-z)$, gives $P_F(\lambda q)=2^{n-1}\int_{\cT}e^{\lambda X(z)}dz$, absolutely convergent as $|e^{\lambda X}|=e^{\re\lambda\,X}$; here $X$ is the exponent of $q$, and the exponent of $\lambda q$ is $\lambda X$, the exponent being linear in the question.

(i) By Proposition~\hb{prop:evenfloor}(i) and $dz=dv$ on each sector, with $h_\pi=\prod_jk_{j,\pi}^{-s_j}\psi_{j,\pi}^{-t_j}$,
\[
P_F(\lambda q)=2^{n-1}\sum_\pi\int_{[0,\infty)^{n-1}}e^{-2\lambda\sum_k\kappa_{U_k}v_k}\,h_\pi(e^{-2v})^{\lambda}\,dv ,
\]
and with $v=w/\lambda$,
\[
\lambda^{n-1}P_F(\lambda q)=2^{n-1}\sum_\pi\int_{[0,\infty)^{n-1}}e^{-2\sum_k\kappa_{U_k}w_k}\,h_\pi(e^{-2w/\lambda})^{\lambda}\,dw .
\]
On the cube $|\log h_\pi|\le C$ (Proposition~\hb{prop:evenfloor}(i)), so for $0<\lambda\le1$ the integrand lies between $e^{\mp C}e^{-2\sum_k\kappa_{U_k}w_k}$, whose integral is $\prod_k(2\kappa_{U_k})^{-1}$; it tends to $e^{-2\sum_k\kappa_{U_k}w_k}$, and by domination $\lambda^{n-1}P_F(\lambda q)\to2^{n-1}\sum_\pi\prod_k(2\kappa_{U_k})^{-1}=\sum_\pi\prod_k\kappa_{U_k}^{-1}$. For the next order, $|h^\lambda-1-\lambda\log h|\le\frac12\lambda^2C^2e^{\lambda C}$, and for $w$ in the open orthant $e^{-2w/\lambda}\to0$, so $\log h_\pi(e^{-2w/\lambda})\to\log h_\pi(0)$, boundedly; hence
\[
\lambda^{-1}\Big(\lambda^{n-1}P_F(\lambda q)-\sum_\pi\prod_k\kappa_{U_k}^{-1}\Big)\to\sum_\pi\prod_k\kappa_{U_k}^{-1}\,\log h_\pi(0),
\]
and $h_\pi(0)=\prod_j\theta_{j,T^+_{j,\pi}}^{-s_j}\theta_{j,T^-_{j,\pi}}^{-t_j}$ is the sector integrand at $q=0$, so that $\log h_\pi(0)=-\sum_j\big(s_j\log\theta_{j,T^+_{j,\pi}}+t_j\log\theta_{j,T^-_{j,\pi}}\big)$, the coefficient stated. Finally $\kappa$ is positively homogeneous of degree one and positive on $\cT\setminus0$ (Proposition~\hb{prop:evenfloor}(ii)), so $\operatorname{vol}\{\kappa\le r\}=r^{n-1}\operatorname{vol}\{\kappa\le1\}$ and
\[
2^{n-1}\int_{\cT}e^{-2\kappa(z)}dz=2^{n-1}\operatorname{vol}\{\kappa\le1\}\,(n-1)\int_0^\infty e^{-2r}r^{n-2}dr=(n-1)!\operatorname{vol}\{\kappa\le1\},
\]
while on the sectors the same integral is $2^{n-1}\sum_\pi\prod_k(2\kappa_{U_k})^{-1}$. (For $s_j,t_j\ge0$, $\{\kappa\le1\}$ is the polar of $\mathcal N_F-\mvec$ in $\one^\perp$.)

(ii) Let $s_j,t_j\ge0$ and $f=-X$ on $\cT$, of dimension $n-1$ with the measure $dz$. The hypotheses of Lemma~\ref{lem:highstrength} hold. (a) Every derivative of $\log\Lambda_j$ of order $k\ge1$ is a joint cumulant of order $k$ of $\chi_T\in\{\pm1\}^{\cE}$ under $\mathbb P_{j,e^{2z}}$ (Theorem~\hb{thm:innerMGF}), bounded by a constant depending on $k$ alone, and $f=\sum_j\big(s_j\log\Lambda_j(z)+t_j\log\Lambda_j(-z)\big)-2\ip{\bmu_\perp}z$. (b) $v^{\mathsf T}\nabla^2f(z)v=\sum_j\big(s_j\operatorname{Var}_{\mathbb P_{j,e^{2z}}}\ip{\chi_T}v+t_j\operatorname{Var}_{\mathbb P_{j,e^{-2z}}}\ip{\chi_T}v\big)$; a question with $s_j,t_j\ge0$ inside its walls is on a connected family, as $\kappa_S+\kappa_{\cE\setminus S}=2\sum_j\sigma_jc_j(S)>0$ (Lemma~\ref{lem:walls}) puts some active table of positive width on every $S$, and the argument of Proposition~\hb{prop:evenfloor}(iii) makes the form positive definite on $\cT$ at every $z$; so $f$ is strictly convex. (c) $f\ge2\kappa-C\ge2\kappa_{\min}(\max_ez_e-\min_ez_e)-C$ by Proposition~\hb{prop:evenfloor}(i)--(ii), and $\max_ez_e-\min_ez_e$ is a norm on $\cT$. A strictly convex $f$ with (c) attains its minimum at one point $z^\star$, its only critical point. By Theorem~\hb{thm:innerMGF}, $\nabla\log\Lambda_j(z)=2R_j(\alpha)-\one$ and $\nabla[\log\Lambda_j(-z)]=\one-2R_j(\alpha^{-1})=2R^\ast_j(\alpha)-\one$ at $\alpha=e^{2z}$, so the derivative of $f$ along $v\in\one^\perp$ is $2\ip{\mathbf u(\alpha)-\mvec}v$, $\mathbf u=\sum_j(s_jR_j+t_jR^\ast_j)$ the usage, $\mvec-\sigma\one-\bmu_\perp$ being a multiple of $\one$. It vanishes for every $v$ exactly when $\mathbf u(\alpha)-\mvec\in\mathbb R\one$, that is $\mathbf u(\alpha)=\mvec$, the two having the total $\sum_j(s_jr_j+t_jr^\ast_j)$ on the balance wall; and $\mvec$ is the target of $q$ there, so $z^\star$ is its native gauge (Definition~\hb{def:native}). The usage and the masses of $\lambda q$ are $\lambda$ times those of $q$, so $z^\star$ is the native gauge of every $\lambda q$, $\lambda>0$. Lemma~\ref{lem:highstrength} with $d=n-1$ gives the expansion, the determinant taken on a chart $z_g=0$, that is of the Hessian of the invariant form $\tilde X$ of the proof of Proposition~\hb{prop:evenfloor} with the row and the column of $g$ deleted. With $y=2z$, $\Phi(y)=\sigma\ip\one y+\Phi^\circ(z)$ (Theorem~\hb{thm:saddle}), so $\partial_{z_e}\partial_{z_f}(-\tilde X)=\partial_{z_e}\partial_{z_f}\Phi^\circ=4\partial_{y_e}\partial_{y_f}\Phi$, rung two of the still of the realization $\big((s_j,t_j)_j;e^{2z^\star}\big)$ of the active tables (Definition~\hb{def:still}). At the centre $f=\sum_j\sigma_jE_j$ is even and strictly convex, so $z^\star=0$, $f(0)=2\sum_j\sigma_j\log K_j(\one)$, and $\nabla^2f(0)=\sum_j\sigma_j\cdot2\operatorname{Cov}_{\mathbb P_{j,\one}}(\chi_T)=8\sum_j\sigma_j\big(\kappa_{j,\one}(X_e,X_f)\big)_{e,f}$, as $\chi_T=2X-\one$, read on the chart $z_g=0$ with the row and the column of $g$ deleted.
\end{proof}

\begin{sremark}
At one table and the centre of strength one, $\kappa=\frac12\wdt$ and $(n-1)!\operatorname{vol}\{\kappa\le1\}=2^{n-1}(n-1)!\operatorname{vol}\{\wdt\le1\}$, the reference's low end; its high end is (ii) at offsets zero. Divided by the same limit for the list $\binom{\cE}1$ at a full basis lean, whose walls are the masses, (i) is the limit of Proposition~\hb{prop:movingtarget}: there $\Gst{n}{1}{\lambda\btau}{\lambda\tau_{\cE}}$ has $\lambda^{n-1}\Gst{n}{1}{\lambda\btau}{\lambda\tau_{\cE}}\to\tau_{\cE}/\prod_e\tau_e=\sum_\pi\prod_k\tau_{U_k}^{-1}$. In (ii) the real ray is Proposition~\hb{prop:laplace} written in the signed coordinates; what the lemma adds is the whole half-plane, uniformly.
\end{sremark}

\subsection{How far a question can lean, and the tails}

\begin{hproposition}[\hb{prop:reach} (How far a question can lean; the values of the statistics)]
As stated.
\end{hproposition}

\begin{proof}
The decomposition of the walls is Lemma~\ref{lem:walls}. Every proper wall is positive exactly when, for every $S$, $\kappa^{\mathbf E}_S+\kappa^{\mathbf O}_S>0$ and $\kappa^{\mathbf E}_S-\kappa^{\mathbf O}_S>0$, that is $|\kappa^{\mathbf O}_S|<\kappa^{\mathbf E}_S$. The fold of Theorem~\hb{thm:fold} sends the pair coordinates $(\bmu_\perp;\bsig,\bdel)$ to $(-\bmu_\perp;\bsig,-\bdel)$, negating the offsets and the leans and keeping the strengths (Definition~\hb{def:realization}); so it negates $\kappa^{\mathbf O}_S$ and swaps $\kappa_S$ with $\kappa_{\cE\setminus S}$. When every $\kappa^{\mathbf E}_S(\bsig)$ is positive the conditions read $\kappa^{\mathbf E}_S/|\kappa^{\mathbf O}_S|>1$ wherever $\kappa^{\mathbf O}_S\ne0$, that is $\ell^\ast_{\bsig}(\bmu_\perp,\bdel)>1$: finitely many strict linear inequalities in the odd part, each with its negative, an open convex polyhedron symmetric under the fold. It is bounded exactly when the linear map $(c,\mathbf b)\mapsto(\kappa^{\mathbf O}_S(c,\mathbf b))_S$ is injective: then it lies in the preimage of a bounded box, and otherwise a line through $0$ in the kernel lies in it. The homogeneities and the evenness in $(c,\mathbf b)$ are read off the definition. For one connected table at offsets zero, $\kappa^{\mathbf E}_S=\sigma c(S)$ and $\kappa^{\mathbf O}_S=-\delta\varepsilon^\perp(S)$, so the conditions are $\sigma>0$ and $|\delta|<\sigma\min c(S)/|\varepsilon^\perp(S)|=\ell^\ast_1(0,1)\,\sigma$.

(i) The exponent is linear in the question, so with $q'=(c;\bsig+\mathbf a,\mathbf b)$ and the exponent $X_{q'}$ of $q'$, the expectation is $2^{n-1}P_F(\sigma\one;\bsig,\bsig)^{-1}\int_{\cT}e^{X_{q'}(z)}dz$, an integral of a positive function. By Proposition~\hb{prop:evenfloor}(i), whose identity holds at every real exponents, $e^{X_{q'}}$ lies between constant multiples of $e^{-2\kappa'(z)}$, $\kappa'$ the wall at the tilt of $q'$, and $\int_{H_\pi}e^{-2\kappa'(z)}dz=\int_{[0,\infty)^{n-1}}e^{-2\sum_k\kappa_{U_k}(q')v_k}dv$ is finite exactly when every $\kappa_{U_k}(q')$ is positive. Every nonempty proper $S$ is a tail of some order, so the expectation is finite exactly when every proper wall of $q'$ is positive, and it is then $P_F(q')/P_F(\sigma\one;\bsig,\bsig)$ by Definition~\hb{def:question}. The walls of $q'$ are $\kappa^{\mathbf E}_S(\bsig+\mathbf a)\pm\kappa^{\mathbf O}_S(c,\mathbf b)$; at $\mathbf a=0$ their positivity is $\ell^\ast_{\bsig}(c,\mathbf b)>1$.

(ii) With $|\delta_j|\le\sigma_j$ the exponents are nonnegative, and by the proof of Theorem~\hb{thm:severity}(ii) $-X$ is strictly convex on $\cT$ with its least value at the native gauge $z^\star$, at which alone it is attained. $-X(z)$ is the value of the linear form $-2\ip{\bmu_\perp}x+\sum_j(\sigma_jy_j+\delta_jw_j)$ at the point $(z;E_j(z),O_j(z))_j$, and $z$ is a coordinate of that point, so the least value of the form on the set of values of the statistics is taken at $\big(z^\star;E_j(z^\star),O_j(z^\star)\big)_j$ and only there. At the centre $z^\star=0$ (Proposition~\hb{prop:evenfloor}(iii)), $E_j(0)=2\log K_j(\one)$ and $O_j(0)=0$.
\end{proof}

\begin{hproposition}[\hb{prop:tails} (The tails of the statistics are read at the walls)]
As stated.
\end{hproposition}

\begin{proof}
By Lemma~\ref{lem:walls} the walls are linear in the pair coordinates, and for the direction $v$ of pair coordinates $(c;\mathbf a,\mathbf b)$, $\kappa_S(v)=\kappa^{\mathbf E}_S(\mathbf a)+\kappa^{\mathbf O}_S(c,\mathbf b)$, the offset direction $c$ entering through $c_S=\sum_{e\in S}c_e$. With $X_v(z)=2\ip cz-\sum_j(a_jE_j(z)+b_jO_j(z))$, $Y=X_v(Z)$ and $X_q+\lambda X_v=X_{q+\lambda v}$, so for real $\lambda\ge0$, $\oE[e^{\lambda Y}]=2^{n-1}P_F(q)^{-1}\int_{\cT}e^{X_{q+\lambda v}}dz$, finite, by the argument of Proposition~\hb{prop:reach}(i), exactly when every $\kappa_S(q)+\lambda\kappa_S(v)$ is positive, that is for $\lambda<\lambda^\ast$, and then the ratio of the periods.

\emph{Upper bound.} For $0<\lambda<\lambda^\ast$, $\mathbb P(Y>x)\le e^{-\lambda x}\oE[e^{\lambda Y}]$, so $\limsup_x x^{-1}\log\mathbb P(Y>x)\le-\lambda$ for every $\lambda<\lambda^\ast$.

\emph{Lower bound.} Let $\lambda^\ast<\infty$, let $S^\ast$ attain it, and let $\pi$ list $\cE\setminus S^\ast$ first, so that $S^\ast=U_{k^\ast}$ is a tail of $\pi$. Write $\kappa_q(z)$ and $\kappa_v(z)=\sum_k\kappa_{U_k}(v)\big(z_{\pi(k-1)}-z_{\pi(k)}\big)$ for the walls at the tilt of $q$ and of $v$, the latter read as balanced data with the exponents $(\mathbf a+\mathbf b,\mathbf a-\mathbf b)$. By Proposition~\hb{prop:evenfloor}(i), whose identity holds at any balanced real data, $|X_q+2\kappa_q|\le C_q$ and $|X_v+2\kappa_v|\le C_v$ on $\cT$. For $\varepsilon>0$ put $a_x=(1+\varepsilon)x/(-2\kappa_{S^\ast}(v))$ and let $A_x\subset H_\pi$ be the set of the tilts whose difference along $\pi$ at the tail $S^\ast$ lies in $[a_x,a_x+1]$ and whose other differences lie in $[0,1]$, of $dz$-volume $1$. On $A_x$, as $\kappa_{S^\ast}(v)<0$,
\[
Y\ge-2\kappa_v(z)-C_v\ge(1+\varepsilon)x-2\sum_k|\kappa_{U_k}(v)|-C_v>x\quad\text{for large }x,
\]
\[
X_q\ge-2\kappa_q(z)-C_q\ge-2\kappa_{S^\ast}(q)(a_x+1)-2\sum_k|\kappa_{U_k}(q)|-C_q=-(1+\varepsilon)\lambda^\ast x-C',
\]
since $2\kappa_{S^\ast}(q)a_x=(1+\varepsilon)\lambda^\ast x$. So $\mathbb P(Y>x)\ge2^{n-1}P_F(q)^{-1}e^{-(1+\varepsilon)\lambda^\ast x-C'}$ for large $x$, $\liminf_xx^{-1}\log\mathbb P(Y>x)\ge-(1+\varepsilon)\lambda^\ast$ for every $\varepsilon>0$, and with the upper bound the limit is $-\lambda^\ast$.

If $\lambda^\ast=\infty$, every $\kappa_S(v)\ge0$, so $\kappa_v\ge0$ on every sector and $Y\le C_v$.

At the centre of strengths $\bsig$ and the direction of pair coordinates $(c;\mathbf 0,\mathbf b)$, $\kappa_S(q)=\kappa^{\mathbf E}_S(\bsig)$ and $\kappa_S(v)=\kappa^{\mathbf O}_S(c,\mathbf b)$. The first is unchanged and the second negated by $S\mapsto\cE\setminus S$, so the sets on which $\kappa_S(v)<0$ are the complements of those on which it is positive, and $\lambda^\ast=\min_S\kappa^{\mathbf E}_S/|\kappa^{\mathbf O}_S|=\ell^\ast_{\bsig}(c,\mathbf b)$, for $Y$ and for $-Y$ alike.
\end{proof}

\begin{sremark}
The one-table rate of both tails of $O(Z)$ at the centre of strength $\sigma_0$ is $\ell^\ast_{\sigma_0}(0,1)=\sigma_0\,\ell^\ast_1(0,1)$. The same proposition, applied to the family with the list $\binom{\cE}1$ added, gives the tails of the Grey law (Proposition~\hb{prop:greylaw}(iii), proved in Section~\ref{sec:ch39more}).
\end{sremark}

\section{Chapter~\hb{ch:law}: the Breta law and the port on the balance wall}\label{sec:S1417-ch17}

\begin{hproposition}[\hb{prop:stick} (Scale and shapes along a bunch)]
As stated: off the balance wall and at complex exponents, the scale integral then carrying the appended list $\binom{\cE}1$.
\end{hproposition}

\begin{proof}
\emph{The display.} The masses are positive and every proper wall has positive real part, so $P_F(\mvec;\mathbf s,\mathbf t)=\int_{\Delta_{\cE}}u^{\mvec-\one}W(u)\,d\omega_{\cE}(u)$ converges (\hb{thm:projective}). By \hb{thm:aggregation} with the blocks $B$ and $R$, $u=(w_Bv_B,w_Rv_R)$ with $(w_B,w_R)\in\Delta_{\{B,R\}}$, $v_B\in\Delta_B$ and $v_R\in\Delta_R$, and
\[
u^{\mvec-\one}d\omega_{\cE}=w_B^{m_B-1}w_R^{m_R-1}dw_B\cdot v_B^{\mvec|_B-\one}d\omega_B\cdot v_R^{\mvec|_R-\one}d\omega_R .
\]
With $x=w_B/w_R\in(0,\infty)$, $w_B=x/(1+x)$, $w_R=1/(1+x)$ and $dw_B=(1+x)^{-2}dx$, so $w_B^{m_B-1}w_R^{m_R-1}dw_B=x^{m_B-1}(1+x)^{-m_{\cE}}dx$. The word $W$ is homogeneous of degree $\kappa_{\cE}-m_{\cE}$ (\hb{def:family}), so $W(w_Bv_B,w_Rv_R)=w_R^{\kappa_{\cE}-m_{\cE}}W(xv_B,v_R)=(1+x)^{m_{\cE}-\kappa_{\cE}}W(xv_B,v_R)$, and the powers of $1+x$ collect to $(1+x)^{-\kappa_{\cE}}$. By Tonelli on the modulus, which is the same integrand at the real parts of the exponents, and Fubini,
\[
P_F(\mvec;\mathbf s,\mathbf t)=\int_{\Delta_B}\int_{\Delta_R}v_B^{\mvec|_B-\one}v_R^{\mvec|_R-\one}\int_0^\infty x^{m_B-1}(1+x)^{-\kappa_{\cE}}W(xv_B,v_R)\,dx\,d\omega_R(v_R)\,d\omega_B(v_B).
\]
The densities of $\Dir_{\mvec|_B}$ and $\Dir_{\mvec|_R}$ are $v_B^{\mvec|_B-\one}/\Gst{|B|}1{\mvec|_B}{m_B}$ and $v_R^{\mvec|_R-\one}/\Gst{|R|}1{\mvec|_R}{m_R}$ (\hb{def:dirlaw}), which gives the display. At $(xU'_B,U'_R)$ the letter $|\alpha|$ of the list $\binom{\cE}1$ is $x|U'_B|+|U'_R|=1+x$, so the inner integral is the scale integral (\hb{thm:scale}) of the family with that table appended at the exponents $(\kappa_{\cE},0)$ (\hb{prop:sectiontable}). On the balance wall the factor $(1+x)^{-\kappa_{\cE}}$ is $1$, the inner integral is $\operatorname{Sc}_B(U'_B,U'_R;\mvec;\mathbf s,\mathbf t)$, and the double integral above is \hb{cor:fiber} at $h=|\alpha_R|$.

\emph{The conditional density.} For real data with every proper wall positive the same computation, which needs no sign of the masses, shows that under $\mathrm{Breta}^\circ(F;\mvec;\mathbf s,\mathbf t)$ the variables $\rho=|U_B|/|U_R|$, $V_B$ and $V_R$ have the joint density
\[
\rho^{m_B-1}(1+\rho)^{-\kappa_{\cE}}W(\rho V_B,V_R)\,V_B^{\mvec|_B-\one}V_R^{\mvec|_R-\one}\big/P_F(\mvec;\mathbf s,\mathbf t)
\]
against $d\rho\,d\omega_B\,d\omega_R$. Given $(V_B,V_R)$, the density of $\rho$ is proportional to $\rho^{m_B-1}(1+\rho)^{-\kappa_{\cE}}W(\rho V_B,V_R)$. On the balance wall this is the reference, proof of \hb{prop:stick}.
\end{proof}

\begin{hcorollary}[\hb{cor:conjugacyoff} (Conjugacy for the continued period)]
As stated.
\end{hcorollary}

\begin{proof}
Let $\Omega$ be the set of the data at which every nonempty proper wall has positive real part, open and convex, and not empty: at $\mathbf s=\mathbf t=\mathbf 0$ the walls are $\kappa_S=m_S$. Adding $\one_T$ to the masses and $\mathbf e_j$ to $\mathbf s$ adds $r_j$ to $m_{\cE}$ and to $\sum_i(s_ir_i+t_ir^\ast_i)$, so it keeps the imbalance and maps the balance wall into itself, and it adds $|T\cap S|-d_j(S)\ge0$ to the wall $\kappa_S$ (\hb{def:family}), so it maps $\Omega$ into $\Omega$; adding $\one_{\cE\setminus T'}$ to the masses and $\mathbf e_j$ to $\mathbf t$ does the same, adding $|S\setminus T'|-d^\ast_j(S)\ge0$ to $\kappa_S$. At data in $\Omega$ the three periods converge absolutely as integrals over $\Delta_{\cE}$ (\hb{thm:projective}). On $\Delta_{\cE}$, $K_j(u)=\sum_T\theta_{j,T}u^{\one_T}$ and $\Psi_j(u)=\sum_{T'}\theta_{j,T'}u^{\one_{\cE\setminus T'}}$ (\hb{def:table}), so
\[
u^{\mvec-\one}\prod_iK_i(u)^{-s_i}\Psi_i(u)^{-t_i}=\sum_T\theta_{j,T}\,u^{\mvec+\one_T-\one}K_j(u)^{-s_j-1}\prod_{i\neq j}K_i(u)^{-s_i}\prod_i\Psi_i(u)^{-t_i},
\]
the powers of the positive letters being principal, and integrating term by term gives the first identity on $\Omega$; the second is the same with $\Psi_j$. By \hb{lem:periodmero} the three sides are meromorphic in the masses and the exponents, the shifted ones being $P_F$ composed with translations, and they agree on the nonempty open set $\Omega$, so they agree as meromorphic functions. Their restrictions to the balance wall, meromorphic by \hb{lem:periodmero} and mapped into themselves by the translations, then agree as well.
\end{proof}

\begin{slemma}\label{lem:ch17-laplace}
Let $G$ be a convex function of class $C^2$ on $\mathbb R^d$ with $G(0)=0$, $\nabla G(0)=0$ and $\nabla^2G(0)$ positive definite, and let $b$ be a measurable function on $\mathbb R^d$, continuous at $0$, with $|b(y)|\le e^{c_0(1+|y|)}$ for a constant $c_0\ge0$. Then $\int_{\mathbb R^d}|b|e^{-\lambda G}dy<\infty$ for every large $\lambda$, and for every bounded continuous function $H$ on $\mathbb R^d$
\[
\lim_{\lambda\to+\infty}\lambda^{d/2}\int_{\mathbb R^d}b(y)H(\sqrt\lambda\,y)e^{-\lambda G(y)}dy=b(0)\int_{\mathbb R^d}H(x)e^{-\frac12\ip x{\nabla^2G(0)x}}dx .
\]
At $H=1$ the right side is $b(0)(2\pi)^{d/2}\det\big(\nabla^2G(0)\big)^{-1/2}$.
\end{slemma}

\begin{proof}
Let $\epsilon>0$ be the least eigenvalue of $\nabla^2G(0)$. By Taylor's formula $G(y)=\frac12\ip y{\nabla^2G(0)y}+o(|y|^2)$; choose $\delta>0$ with $G(y)\ge\frac14\ip y{\nabla^2G(0)y}\ge\frac\epsilon4|y|^2$ and $|b(y)|\le|b(0)|+1$ for $|y|\le\delta$.

\emph{Linear growth.} For $|y|\ge\delta$ put $t=\delta/|y|\in(0,1]$. Convexity and $G(0)=0$ give $G(ty)\le tG(y)$, so $G(y)\ge t^{-1}G(ty)\ge\frac{|y|}\delta\cdot\frac\epsilon4\delta^2=\gamma|y|$, with $\gamma=\frac{\epsilon\delta}4$.

\emph{The outer part.} For $\lambda\gamma\ge2c_0$ and $|y|\ge\delta$, $|b(y)|e^{-\lambda G(y)}\le e^{c_0-\frac12\lambda\gamma|y|}\le e^{c_0-\frac14\lambda\gamma\delta}e^{-\frac14\lambda\gamma|y|}$. So $\int_{|y|\ge\delta}|b|e^{-\lambda G}dy\le e^{c_0-\frac14\lambda\gamma\delta}\big(\tfrac4{\lambda\gamma}\big)^d\int_{\mathbb R^d}e^{-|x|}dx$, and $\lambda^{d/2}$ times it tends to $0$. With the bound on $|y|\le\delta$, the integral of $|b|e^{-\lambda G}$ is finite.

\emph{The inner part.} With $y=x/\sqrt\lambda$,
\[
\lambda^{d/2}\int_{|y|<\delta}b(y)H(\sqrt\lambda\,y)e^{-\lambda G(y)}dy=\int_{\mathbb R^d}\one\big[|x|<\delta\sqrt\lambda\big]\,b(x/\sqrt\lambda)H(x)e^{-\lambda G(x/\sqrt\lambda)}dx .
\]
As $\lambda\to\infty$ the integrand tends to $b(0)H(x)e^{-\frac12\ip x{\nabla^2G(0)x}}$ at every $x$, by the continuity of $b$ at $0$ and $\lambda G(x/\sqrt\lambda)\to\frac12\ip x{\nabla^2G(0)x}$, and it is bounded by $(|b(0)|+1)\sup|H|\,e^{-\epsilon|x|^2/4}$, which is integrable. Dominated convergence gives the limit. At $H=1$, $\int e^{-\frac12\ip x{Ax}}dx=(2\pi)^{d/2}(\det A)^{-1/2}$ for positive definite $A$.
\end{proof}

\begin{hproposition}[\hb{prop:laplace} (The period along a ray of data)]
As stated: a base point of complex data on the balance wall added to the scaled data.
\end{hproposition}

\begin{proof}
Write $\mathbf d_\lambda=(\mvec_0+\lambda\btau;\mathbf s_0+\lambda\mathbf s,\mathbf t_0+\lambda\mathbf t)$, and in this proof $G(y)=\Phi(y)-\ip{\btau}y$, $\Phi$ the potential at $(\mathbf s,\mathbf t)$.

\emph{The data.} The walls and the imbalance are affine in the data with real coefficients. At $(\btau;\mathbf s,\mathbf t)$, real data with $\btau>0$, $s_j,t_j\ge0$ not all zero and a native gauge on a connected family, every proper wall is positive (\hb{thm:saddle}(iii)), and both $(\btau;\mathbf s,\mathbf t)$ and $(\mvec_0;\mathbf s_0,\mathbf t_0)$ lie on the balance wall. So $\kappa_{\cE}(\mathbf d_\lambda)=0$, and $\re\kappa_S(\mathbf d_\lambda)=\re\kappa_S(\mvec_0;\mathbf s_0,\mathbf t_0)+\lambda\kappa_S(\btau;\mathbf s,\mathbf t)\to+\infty$ for every nonempty proper $S$: for $\lambda$ large every proper wall of $\mathbf d_\lambda$ has positive real part.

\emph{The chart integral.} For such $\lambda$, $P_F(\mathbf d_\lambda)$ is the integral over the chart $\alpha_g=1$ (\hb{lem:chart}). In $y=\log\alpha$, $y_g=0$, the measure is $\alpha^{\mvec-\one}d\alpha_{\cE-g}=e^{\ip{\mvec}y}dy_{\cE-g}$, and the family word at the exponents $(\mathbf s_0+\lambda\mathbf s,\mathbf t_0+\lambda\mathbf t)$ is $W_0(e^y)e^{-\lambda\Phi(y)}$ (\hb{def:potential}). So
\[
P_F(\mathbf d_\lambda)=\int_{\{y_g=0\}}b(y)\,e^{-\lambda G(y)}\,dy_{\cE-g},\qquad b(y)=e^{\ip{\mvec_0}y}W_0(e^y).
\]

\emph{The minimum.} $\Phi$ is convex (\hb{thm:saddle}(ii)), and $\nabla G(y)=\nabla\Phi(y)-\btau$ is the usage at $e^y$ less the target, which vanishes at $y^\star=\log\alpha^\star$ (\hb{thm:saddle}(i)); normalize $\alpha^\star_g=1$, so that $y^\star$ lies on the chart. For $v\in\mathbb R^{\cE}$ and $\alpha=e^y$,
\[
\ip v{\nabla^2\Phi(y)v}=\sum_j\Big(s_j\operatorname{Var}_{T\sim\mathbb P_{j,\alpha}}\ip v{\one_T}+t_j\operatorname{Var}_{T\sim\mathbb P^\ast_{j,\alpha}}\ip v{\one_T}\Big),
\]
since $\partial_c^2\log K_j(e^{y+cv})$ at $c=0$ is the variance of $\ip v{\one_T}$ under $\mathbb P_{j,\alpha}(T)=\theta_{j,T}\alpha^T/K_j(\alpha)$ (\hb{thm:innerMGF}), and likewise for $\Psi_j$ under $\mathbb P^\ast_{j,\alpha}$, with $\ip v{\one_{\cE\setminus T}}=\ip v\one-\ip v{\one_T}$. Both laws charge every basis. So if $\ip v{\nabla^2\Phi(y)v}=0$, then $\ip v{\one_T}$ is constant on the bases of every active table, $v$ is constant on every molecule of the active family (\hb{lem:circuit}), and the active family of a connected family is one molecule (Lemma~\ref{lem:families-separators}); so $v\in\mathbb R\one$. Conversely $\nabla^2\Phi\,\one=0$, $\Phi(y+c\one)-\Phi(y)$ being linear in $c$. So the restriction of $\nabla^2\Phi(y^\star)$ to the chart $\{v_g=0\}$, the matrix $\big(\nabla^2\Phi(y^\star)\big)_{(g)}$, is positive definite; it is the Hessian at $y^\star$ of $G$ on the chart.

\emph{The amplitude.} For every basis $T$ of table $j$, $|\ip{\one_T}y|\le r_j|y|$, so $\log\min_T\theta_{j,T}-r_j|y|\le\log K_j(e^y)\le\log K_j(\one)+r_j|y|$, and likewise for $\Psi_j$ with $n-r_j$. Hence $|b(y)|=e^{\ip{\re\mvec_0}y}\prod_jK_j(e^y)^{-\re s_{0,j}}\Psi_j(e^y)^{-\re t_{0,j}}\le e^{c_0(1+|y|)}$ for a constant $c_0$, and $b$ is continuous and nowhere $0$.

\emph{Laplace's method.} Lemma~\ref{lem:ch17-laplace} on the chart, with $d=n-1$, the function $G(y^\star+x)-G(y^\star)$ and the amplitude $b(y^\star+x)$, bounded by $e^{c_0(1+|y^\star|+|x|)}$, gives
\[
P_F(\mathbf d_\lambda)=e^{-\lambda G(y^\star)}\Big(\frac{2\pi}\lambda\Big)^{\frac{n-1}2}\det\big(\nabla^2\Phi(y^\star)\big)_{(g)}^{-1/2}\big(b(y^\star)+o(1)\big),
\]
with $b(y^\star)=(\alpha^\star)^{\mvec_0}W_0(\alpha^\star)\neq0$ and $G(y^\star)=\Phi(\log\alpha^\star)-\ip{\btau}{\log\alpha^\star}$: the display. The function $\alpha\mapsto\alpha^{\mvec_0}W_0(\alpha)$ is homogeneous of degree $\kappa_{\cE}(\mvec_0;\mathbf s_0,\mathbf t_0)=0$, so $b(y^\star)$ does not depend on the scale of $\alpha^\star$; likewise $G(y+c\one)=G(y)$, the data $(\btau;\mathbf s,\mathbf t)$ being balanced, and $\nabla^2\Phi$ does not change along $\one$. By \hb{def:still}, rung two of the still of $((s_j,t_j)_j;\alpha^\star)$ is $4\nabla^2\Phi(\log\alpha^\star)$.

\emph{The normal limit.} For real $(\mvec_0;\mathbf s_0,\mathbf t_0)$ and $\lambda$ large, $\mathrm{Breta}^\circ(F;\mathbf d_\lambda)$ is defined (\hb{def:shapelaw}). For a bounded continuous function $H$ of $y$ modulo $\mathbb R\one$, the function $u\mapsto H(\log u)u^{\mvec-\one}W(u)$ is homogeneous of degree $-n$, so its integral over $\Delta_{\cE}$ is its integral over the chart, as in the proof of Lemma~\ref{lem:ch16-compound}(ii). So the law of $Y=\log U$ modulo $\mathbb R\one$, read on the chart, has the density $b(y)e^{-\lambda G(y)}/P_F(\mathbf d_\lambda)$, with $b>0$ continuous. Applying Lemma~\ref{lem:ch17-laplace} as above to $H(\sqrt\lambda(y-y^\star))$ and to $1$, and dividing,
\[
\oE\big[H\big(\sqrt\lambda(Y-y^\star)\big)\big]\to\int H(x)e^{-\frac12\ip x{\nabla^2\Phi(y^\star)_{(g)}x}}dx\Big/\int e^{-\frac12\ip x{\nabla^2\Phi(y^\star)_{(g)}x}}dx,
\]
the expectation of $H$ under the centred normal law on the chart with covariance $\big(\nabla^2\Phi(y^\star)_{(g)}\big)^{-1}$: the centred normal law on $\mathbb R^{\cE}/\mathbb R\one$ with covariance $(\nabla^2\Phi)^{-1}$, the quadratic form of $\nabla^2\Phi(y^\star)$ on $\mathbb R^{\cE}/\mathbb R\one$ being read on the chart. At $(\mvec_0;\mathbf s_0,\mathbf t_0)=0$, $b=1$ and this is the reference, proof of \hb{prop:laplace}.
\end{proof}

\begin{sremark}
At $(\mvec_0;\mathbf s_0,\mathbf t_0)=0$ let $q$ be the question of the data $(\btau;\mathbf s,\mathbf t)$ (\hb{def:question}), with the exponent $X(z)=2\ip{\bmu_\perp}z-\sum_j\big(\sigma_jE_j(z)+\delta_jO_j(z)\big)$, read at the representative of $z$ in $\one^\perp$. By \hb{thm:saddle}, $\Phi(y)=\sigma\ip\one y+\sum_j\big(\sigma_jE_j(z)+\delta_jO_j(z)\big)$ in $y=2z$, and $\ip{\btau}y=2\ip{\bmu}z+\sigma\ip\one y$, $\btau$ being the masses of the question; so $\ip{\btau}y-\Phi(y)=X(z)$ for $z\in\one^\perp$, and both sides are unchanged along $\one$. Hence $e^{-\lambda(\Phi(y^\star)-\ip{\btau}{y^\star})}=e^{\lambda X(z^\star)}$, $z^\star$ the tilt of $\frac12y^\star$. The Hessian in $z$ is four times the Hessian in $y$, so on a chart $\det\big(-\nabla^2X(z^\star)\big)=4^{n-1}\det\big(\nabla^2\Phi(y^\star)\big)_{(g)}$, and
\[
2^{n-1}e^{\lambda X(z^\star)}\Big(\frac{2\pi}\lambda\Big)^{\frac{n-1}2}\det\big(-\nabla^2X(z^\star)\big)^{-\frac12}=e^{-\lambda(\Phi(y^\star)-\ip{\btau}{y^\star})}\Big(\frac{2\pi}\lambda\Big)^{\frac{n-1}2}\det\big(\nabla^2\Phi(y^\star)\big)_{(g)}^{-1/2}:
\]
the leading term of \hb{thm:severity}(ii) at real $\lambda$ is the display of \hb{prop:laplace}.
\end{sremark}

\begin{hdefinition}[\hb{def:skellaw} (Skeleton law)]
As stated, at real nonzero letter exponents.
\end{hdefinition}

\begin{proof}[Proof of the last clause]
For a real coefficient $\xi_\ell>0$ and $y>0$, $(y+\xi_\ell)^{-w_\ell}>0$. A coefficient that is not real comes with its conjugate and the same real letter exponent $w_\ell$; $y+\xi_\ell$ and $y+\bar\xi_\ell$ are conjugate and off $(-\infty,0]$, so their principal powers are conjugate and $(y+\xi_\ell)^{-w_\ell}(y+\bar\xi_\ell)^{-w_\ell}=|y+\xi_\ell|^{-2w_\ell}>0$. So $y^{\kappa-1}\prod_\ell(y+\xi_\ell)^{-w_\ell}$ is positive and continuous on $(0,\infty)$. As $y\to0$ it is asymptotic to $y^{\kappa-1}\prod_\ell\xi_\ell^{-w_\ell}$, the product being positive by the same pairing, and as $y\to\infty$ to $y^{\kappa-1-|\mathbf w|}$; so it is integrable on $(0,\infty)$ exactly when $0<\kappa<|\mathbf w|$. Its integral is the value of \hb{prop:scalebreta}, which asks no sign of the letter exponents; so the normalizer is that value, positive, and the density is that of a law.
\end{proof}

\begin{htheorem}[\hb{thm:skellaw} (The Skeleton law)]
As stated: (i) off the balance wall and at exponents of any sign; (ii) at real nonzero letter exponents; (iii) about every point $r\ge\max_\ell\xi_\ell$ of the axis.
\end{htheorem}

\begin{proof}
(i) By \hb{prop:stick}, given the internal shapes $(V_B,V_R)$ the scale $\rho$ of $B$ has a density proportional to $\rho^{m_B-1}(1+\rho)^{-\kappa_{\cE}}\prod_jK_j(\rho V_B,V_R)^{-s_j}\Psi_j(\rho V_B,V_R)^{-t_j}$. At the gauge $\alpha=(V_B,V_R)$, \hb{def:skelpoly} gives
\[
K_j(\rho\alpha_B,\alpha_R)=\rho^{d_j(B)}K_j(\alpha)\prod_i\frac{\rho+\xi_{j,i}}{1+\xi_{j,i}},\qquad\Psi_j(\rho\alpha_B,\alpha_R)=\rho^{d^\ast_j(B)}\Psi_j(\alpha)\prod_i\frac{\rho+\xi^\ast_{j,i}}{1+\xi^\ast_{j,i}},
\]
over the Skeleton values $\xi_{j,i}$ and $\xi^\ast_{j,i}$ of the two letters. The Skeleton polynomials are real, so their values not real come in conjugate pairs of equal multiplicity; the exponents are real, so the principal powers of conjugate factors are conjugate, and $K_j(\rho\alpha_B,\alpha_R)^{-s_j}=\rho^{-s_jd_j(B)}K_j(\alpha)^{-s_j}\prod_i(\rho+\xi_{j,i})^{-s_j}(1+\xi_{j,i})^{s_j}$, likewise for $\Psi_j$. The power of $\rho$ collects to $\rho^{m_B-1-\sum_j(s_jd_j(B)+t_jd^\ast_j(B))}=\rho^{\kappa_B-1}$ (\hb{def:family}), and the linear factors collect to $\prod_\ell(\rho+\xi_\ell)^{-w_\ell}$ over the distinct points of the joint spectrum of the letters with nonzero exponent, with the letter exponents $w_\ell$ of \hb{thm:scale}. The factor $(1+\rho)^{-\kappa_{\cE}}$ is the coefficient $1$ with the letter exponent $\kappa_{\cE}$, added to the letter exponent of $1$ when $1$ is a point of the joint spectrum; discarding the coefficients of letter exponent $0$ changes nothing. So the conditional density is proportional to $\rho^{\kappa_B-1}\prod_\ell(\rho+\xi'_\ell)^{-w'_\ell}$, with coefficients in $\mathbb C\setminus(-\infty,0]$ closed under conjugation and real nonzero letter exponents, equal on conjugate coefficients. Their total is $|\mathbf w|+\kappa_{\cE}=\kappa_B+\kappa^{\mathrm{top}}_B+\kappa_{\cE}$ (\hb{thm:scale}), which is $\kappa_B+\kappa_{\cE\setminus B}$ (\hb{thm:wallscremona}); $B$ and $\cE\setminus B$ are nonempty proper sets and every proper wall is positive, so $0<\kappa_B<|\mathbf w'|$, and the conditional law is $\mathrm{Skel}_{\mathbf w',\kappa_B}(\boldsymbol\xi')$ (\hb{def:skellaw}). On the balance wall this is the reference, proof of \hb{thm:skellaw}(i).

(ii) For $-\kappa<\re k<|\mathbf w|-\kappa$, $\oE[Y^k]$ is $\int_0^\infty y^{\kappa+k-1}\prod_\ell(y+\xi_\ell)^{-w_\ell}dy$ over $\int_0^\infty y^{\kappa-1}\prod_\ell(y+\xi_\ell)^{-w_\ell}dy$, both convergent, the values of \hb{prop:scalebreta} at the indices $\kappa+k$ and $\kappa$; for positive letter exponents the factors $\Gamma(w_\ell)$ of the two stacked Gammas cancel, which is the display. For the inversion put $y=1/z$: $dy=z^{-2}dz$, and $y+\xi=\xi z^{-1}(z+\xi^{-1})$ with $\operatorname{Log}(y+\xi)=\operatorname{Log}\xi-\log z+\operatorname{Log}(z+\xi^{-1})$, the principal arguments of $\xi$ and of $z+\xi^{-1}$ having opposite signs and sum in $(-\pi,\pi)$. So the density of $1/Y$ is proportional to $z^{1-\kappa}z^{-2}z^{|\mathbf w|}\prod_\ell(z+\xi_\ell^{-1})^{-w_\ell}=z^{|\mathbf w|-\kappa-1}\prod_\ell(z+\xi_\ell^{-1})^{-w_\ell}$, the constant $\prod_\ell\xi_\ell^{-w_\ell}$ being positive; with $0<|\mathbf w|-\kappa<|\mathbf w|$ this is $\mathrm{Skel}_{\mathbf w,|\mathbf w|-\kappa}(\boldsymbol\xi^{-1})$. This is the reference, proof of \hb{thm:skellaw}(ii), with \hb{cor:inversion}, which uses no sign of the letter exponents.

(iii) The coefficients are real, hence positive, and $r\ge\max_\ell\xi_\ell$, so $|\xi_\ell-r|<r$ and $0\le g_\ell<1$. By the expansion in the proof of \hb{thm:gaps},
\[
y^{\kappa-1}\prod_\ell(y+\xi_\ell)^{-w_\ell}=\sum_{\mathbf n\in\mathbb N^J}\prod_\ell\frac{(w_\ell)_{n_\ell}}{n_\ell!}(r-\xi_\ell)^{n_\ell}\,y^{\kappa-1}(y+r)^{-|\mathbf w|-|\mathbf n|},
\]
with nonnegative terms, the letter exponents being positive. For $a>\kappa>0$, the law $\mathrm{Skel}_{a,\kappa}(1)$ has the density $y^{\kappa-1}(1+y)^{-a}/\Gst21{\kappa,\ a-\kappa}a$, the Breta of one coefficient $1$ being $1$ (the paragraph after \hb{prop:scalebreta}); so $rY_a$, $Y_a\sim\mathrm{Skel}_{a,\kappa}(1)$, has the density $r^{a-\kappa}y^{\kappa-1}(y+r)^{-a}/\Gst21{\kappa,\ a-\kappa}a$. With $a=|\mathbf w|+|\mathbf n|$ and $(r-\xi_\ell)^{n_\ell}=r^{n_\ell}g_\ell^{n_\ell}$, the density of $Y$ is the sum over $\mathbf n$ of $r^{\kappa-|\mathbf w|}\prod_\ell\frac{(w_\ell)_{n_\ell}}{n_\ell!}g_\ell^{n_\ell}\Gst21{\kappa,\ |\mathbf w|+|\mathbf n|-\kappa}{|\mathbf w|+|\mathbf n|}$ times the density of $rY_{\mathbf n}$, divided by the normalizer; the coefficients are nonnegative and, by Tonelli, sum to $1$. So $Y$ is the stated mixture. At $r=\max_\ell\xi_\ell$ the largest coefficient has $g_\ell=0$, and only $n_\ell=0$ contributes.

(iv) On the balance wall with $s_j,t_j\ge0$, given the shapes the density of the scale $\rho$ is proportional to the integrand of the scale integral at $\alpha=(V_B,V_R)$, by (i). By the proof of \hb{thm:sectors}(i) it is, on $(0,1)$, proportional to $\rho^{\kappa_B-1}\prod_i\sQ_{B,i}(\rho;\alpha)^{-w_i}$, and the density of $1/\rho$ on $(0,1)$ is proportional to $y^{\kappa^{\mathrm{top}}_B-1}\prod_i\sQ_{R,i}(y;\alpha)^{-w_i}$. The expansions of \hb{thm:sectors}(ii) have nonnegative terms, each the term of $\mathbf n$ and $k$ times $y^{\kappa+k-1}(1-y)^{|\mathbf n|}/\Gst21{\kappa+k,\ |\mathbf n|+1}{\kappa+k+|\mathbf n|+1}$, the density of the first coordinate of $\Dir_{(\kappa+k,|\mathbf n|+1)}$ (\hb{thm:dirichlet}); by Tonelli the conditional laws are the stated mixtures, and the two sectors carry probabilities proportional to their Trailmoments. This is the reference, proof of \hb{thm:skellaw}(v).
\end{proof}

\begin{hdefinition}[\hb{def:wallcontour} (Wall contour; joint period)]
As stated, at complex data on the balance wall with $\re s_j>0$ or $s_j=0$, and $\re t_j>0$ or $t_j=0$.
\end{hdefinition}

\begin{sremark}
A wall contour is fixed by real parts alone, and the fibre integral at $B$ has complex pins, whose admissible real points (\hb{def:fibre}) are the wall contours. The two descriptions agree: $\operatorname{load}_e=\sigma+\langle\boldsymbol\delta\rangle_e$ on the share sums and $m_e=\sigma+\mu_e$, so $\re\load_e=\re m_e$ exactly when $\re\langle\boldsymbol\delta\rangle_e=\re\mu_e$. For real data with $s_j,t_j\ge0$ these are the conditions of the reference.
\end{sremark}

\begin{htheorem}[\hb{thm:wallsplit} (The port on the balance wall)]
As stated, at complex data on the balance wall with $\re s_j>0$ or $s_j=0$, and $\re t_j>0$ or $t_j=0$.
\end{htheorem}

\begin{proof}
A letter with exponent $0$ is $1$ and is not opened; the shares, the stacked Gammas and the products below run over the letters with nonzero exponent, whose exponents have positive real part.

\emph{The real parts.} The walls and the imbalance are affine in the data with real coefficients. So the data $(\re\mvec;\re\mathbf s,\re\mathbf t)$ lie on the balance wall, and their proper walls are the real parts of those of $(\mvec;\mathbf s,\mathbf t)$, positive since $P_F(\mvec;\mathbf s,\mathbf t)$ converges (\hb{thm:projective}). Their exponents are nonnegative and not all zero: otherwise $\kappa_S=\re m_S$ for every $S$, and $0=\kappa_{\cE}=\kappa_B+\kappa_R>0$. Their masses are positive: $\re m_e=\kappa_{\{e\}}+\sum_j\big(\re s_j\,d_j(\{e\})+\re t_j\,d^\ast_j(\{e\})\big)>0$, $\{e\}$ being a nonempty proper set ($n\ge2$, as $B$ is proper). A nonempty proper $S$ with $c_j(S)=0$ in every active table would have $\kappa_S+\kappa_{\cE\setminus S}=\kappa_{\cE}=0$ at the real parts (\hb{prop:compoundwall}), against the positivity of both walls; so $F$ is connected, and its active family, whose tables are those of the real parts, is one molecule (Lemma~\ref{lem:families-separators}). By \hb{thm:saddle}(iii) the real parts have a native gauge.

\emph{(i), convergence and the room integral.} By \hb{thm:tworeadings}(i), $K_j=e^{\ip\one z}\Lambda_j(z)$ and $\Psi_j=e^{\ip\one z}\Lambda_j(-z)$ at $\alpha=e^{2z}$, so with $\sigma=\sum_j\sigma_j$ the integrand of $P_F$ is $f(\alpha)=\alpha^{\bmu-\one}\prod_j\Lambda_j(z)^{-s_j}\Lambda_j(-z)^{-t_j}$, homogeneous of degree $\kappa_{\cE}-n=-n$ (\hb{thm:tworeadings}(iii), \hb{thm:pairform}(iii)), and $|f|$ is the same function at the real parts of the data. At the real parts, \hb{thm:projective} on the section $|\alpha_R|=1$, whose measure is $d\alpha_B\,d\omega_R$ (proof of \hb{cor:fiber} above), gives $\int_{\Delta_R}\int_{\mathbb R^B_{>0}}|f(\alpha_B,v)|\,d\alpha_B\,d\omega_R(v)<\infty$; and $f(\alpha_B,v)$ is $v^{\bmu|_R-\one}$ times the integrand of $\Xi_B(v;\bmu|_B;\mathbf s,\mathbf t)$. By Tonelli that integral converges absolutely at almost every $v\in\Delta_R$. For two room points $\alpha_R,\alpha'_R$ there is $C\ge1$ with $\Lambda_{j,(p)}(\pm z_R)\le C\Lambda_{j,(p)}(\pm z'_R)$ for all $j,p$, these being finite sums of positive terms; by \hb{thm:taylorcube}, $\Lambda_j(z)=\sum_pe^{\ip{z_B}{\chi_p}}\Lambda_{j,(p)}(z_R)$, so $\Lambda_j(z_B,z_R)\ge C^{-1}\Lambda_j(z_B,z'_R)$, likewise at $-z$, and the modulus of the integrand of $\Xi_B$ at $\alpha_R$ is at most $C^{\sum_j(\re s_j+\re t_j)}$ times that at $\alpha'_R$. So $\Xi_B$ converges at every $\alpha_R$, and Fubini gives $P_F=\int_{\Delta_R}v^{\bmu|_R-\one}\Xi_B(v)\,d\omega_R(v)$.

\emph{Independence of the pins.} Let $R_1,\dots,R_q$ be the molecules of the children (for $q=1$ there are no pins $\kappa_{R_i}$ below), and write $\kappa_{R_i}=m_{R_i}-\ell_i$ with $\ell_i(\mathbf a,\mathbf b)=\sum_{j,p}\big(a_{j,p}r_{j,(p)}(R_i)+b_{j,p}(|R_i|-r_{j,(p)}(R_i))\big)$. The linear forms $|\mathbf a_j|$ and $|\mathbf b_j|$ of the letters, $\langle\boldsymbol\delta\rangle_e$ for $e\in B$, and $\ell_i$ for $i<q$, are linearly independent. Indeed, let $\sum_{e\in B}\gamma_e\langle\boldsymbol\delta\rangle_e+\sum_{i<q}w_i\ell_i=\sum_j(u_j|\mathbf a_j|+u'_j|\mathbf b_j|)$, and put $V\in\mathbb R^{\cE}$ equal to $\gamma$ on $B$, to $w_i$ on $R_i$ for $i<q$ and to $0$ on $R_q$. For a basis $T$ of table $j$ with $T\cap B=p$, $T\setminus B$ is a basis of the child $(j,p)$, which meets each $R_i$ in $r_{j,(p)}(R_i)$ marks, $R_i$ being a union of molecules of that child (\hb{prop:directsum}); and $\frac12\ip\gamma{\chi_p}=\ip\gamma{\one_p}-\frac12\ip\gamma{\one_B}$. Comparing the coefficients of $a_{j,p}$ gives $\ip V{\one_T}=u_j+\frac12\ip\gamma{\one_B}$, and comparing those of $b_{j,p}$ gives $\ip V{\one_T}=\ip V{\one_R}+\frac12\ip\gamma{\one_B}-u'_j$, for every basis $T$ of table $j$ whenever that letter is opened. So $\ip V{\one_T}$ is constant on the bases of every active table; by \hb{lem:circuit} $V$ is constant on the one molecule $\cE$ of the active family, and as it vanishes on $R_q\neq\varnothing$, $V=0$: $\gamma=0$ and every $w_i=0$.

\emph{(i), the fibre formula.} Fix $\alpha_R$ and a wall contour. Put $X_{j,p}=e^{\ip{z_B}{\chi_p}}\Lambda_{j,(p)}(z_R)$ and $Y_{j,p}=e^{-\ip{z_B}{\chi_p}}\Lambda_{j,(p)}(-z_R)$, positive, so that $\Lambda_j(z)=\sum_pX_{j,p}$ and $\Lambda_j(-z)=\sum_pY_{j,p}$. For a letter with more than one pattern, \hb{thm:E1} on the contour $\re a_{j,p}=c_{j,p}$, admissible since every $c_{j,p}>0$ and $\sum_pc_{j,p}=\re s_j$, gives $\Lambda_j(z)^{-s_j}=\int_{|\mathbf a_j|=s_j}\Gst{|\pat_j(B)|}1{\mathbf a_j}{s_j}\prod_pX_{j,p}^{-a_{j,p}}$; for one pattern this is the evaluation $a_{j,p}=s_j$. Likewise for $\Lambda_j(-z)^{-t_j}$ with the shares $\mathbf b_j$. Since $\prod_{j,p}e^{-\ip{z_B}{\chi_p}(a_{j,p}-b_{j,p})}=e^{-2\ip{z_B}{\langle\boldsymbol\delta\rangle}}=\alpha_B^{-\langle\boldsymbol\delta\rangle}$ (\hb{def:shares}), in $y=\log\alpha_B$, where $\alpha_B^{\bmu|_B-\one}d\alpha_B=e^{\ip{\bmu|_B}y}dy$,
\begin{gather*}
\Xi_B(\alpha_R)=\int_{\mathbb R^B}\Big(\int_{\mathcal X_0}\mathcal G(\mathbf a,\mathbf b)\,e^{\ip{\bmu|_B-\langle\boldsymbol\delta\rangle}{y}}\Big)dy,\\ \mathcal G=\prod_j\Gst{|\pat_j(B)|}1{\mathbf a_j}{s_j}\Gst{|\pat_j(B)|}1{\mathbf b_j}{t_j}\prod_{j,p}\Lambda_{j,(p)}(z_R)^{-a_{j,p}}\Lambda_{j,(p)}(-z_R)^{-b_{j,p}},
\end{gather*}
the inner integral the product of the share integrals, over the set $\mathcal X_0$ of shares on the contour with the share sums pinned. On a wall contour $\re\langle\boldsymbol\delta\rangle_e=\re\mu_e$ for $e\in B$, so $e^{\ip{\bmu|_B-\langle\boldsymbol\delta\rangle}y}=e^{i\ip{\im\bmu|_B-\im\langle\boldsymbol\delta\rangle}y}$. By Stirling's formula $|\Gamma(c+i\eta)|\le C_c(1+|\eta|)^{c-1/2}e^{-\pi|\eta|/2}$, so $|\mathcal G|$, whose factors $\Lambda_{j,(p)}^{-a_{j,p}}$ have moduli constant on the contour, decays exponentially along every imaginary direction of $\mathcal X_0$, uniformly on parallel affine subspaces through a compact set. The map $(\mathbf a,\mathbf b)\mapsto\im\langle\boldsymbol\delta\rangle\in\mathbb R^B$ is affine on $\mathcal X_0$, with linear part onto $\mathbb R^B$ by the independence of the pins; so the image of $\mathcal G$ times the measure of the share integrals has a density $g$ on $\mathbb R^B$, continuous and integrable, and the inner integral is $e^{i\ip{\im\bmu|_B}y}\hat g(y)$, $\hat g(y)=\int g(\nu)e^{-i\ip\nu y}d\nu$. The modulus $|\hat g(y)|$ is the modulus of the integrand of $\Xi_B$ in $y$, integrable by the first step. By Fourier inversion, $\Xi_B(\alpha_R)=\int\hat g(y)e^{i\ip{\im\bmu|_B}y}dy=(2\pi)^{|B|}g(\im\bmu|_B)$, which is the fibre integral of $\mathcal G$ over the fibre $\mathcal X_1$ cut by the pins $|\mathbf a_j|=s_j$, $|\mathbf b_j|=t_j$ and $\langle\boldsymbol\delta\rangle=\bmu|_B$ (\hb{def:fibre}: the $|B|$ further pins raise $(2\pi)^{q-N}$ by $(2\pi)^{|B|}$ and add the factor $\delta^{|B|}(\im\langle\boldsymbol\delta\rangle-\im\bmu|_B)$), with the wall contour as its admissible real point. This is the display of (i). For real data this is the reference, proof of \hb{thm:wallsplit}(i), where $\im\bmu|_B=0$.

\emph{(ii).} On the pins of (i), $\operatorname{load}_e=m_e$ for $e\in B$, and by \hb{lem:minor} the child $(j,p)$ has rank $r_j-|p|$; so
\[
\sum_{i=1}^q\kappa_{R_i}=m_R-\sum_{j,p}\big(a_{j,p}(r_j-|p|)+b_{j,p}(|R|-r_j+|p|)\big)=m_R-\sum_j\big(s_jr_j+t_jr^\ast_j\big)+\sum_{e\in B}\operatorname{load}_e=\kappa_{\cE}=0,
\]
using $\sum_{e\in B}\operatorname{load}_e=\sum_{j,p}\big(a_{j,p}|p|+b_{j,p}(|B|-|p|)\big)$; so $\kappa_{R_q}=0$ holds on $\mathcal X$. Take the section $\{|\alpha_{R_q}|=1\}$, with the measure $d\alpha_B\prod_{i<q}d\alpha_{R_i}\,d\omega_{R_q}$, and as in (i) $P_F=\int\alpha_R^{\bmu|_R-\one}\,\Xi_B(\alpha_R)\prod_{i<q}d\alpha_{R_i}\,d\omega_{R_q}$, absolutely convergent. Write $\alpha_{R_i}=e^{y_i}u_i$, $u_i\in\Delta_{R_i}$, for $i<q$, and $u_q=\alpha_{R_q}$; then $\alpha_{R_i}^{\bmu|_{R_i}-\one}d\alpha_{R_i}=e^{y_i\mu_{R_i}}u_i^{\bmu|_{R_i}-\one}dy_i\,d\omega_{R_i}$. Each $R_i$ is a union of molecules of every child, so every basis $X$ of the child $(j,p)$ has $\ip{\chi_X}{\one_{R_i}}=2r_{j,(p)}(R_i)-|R_i|$, and $\Lambda_{j,(p)}(z_R+\frac12y_i\one_{R_i})=e^{\frac12y_i(2r_{j,(p)}(R_i)-|R_i|)}\Lambda_{j,(p)}(z_R)$. Inserting the fibre formula of (i) and using $\sum_{j,p}\sigma_{j,p}=\sigma$, the exponent of $e^{y_i}$ collects to $\mu_{R_i}-\sum_{j,p}\delta_{j,p}\big(2r_{j,(p)}(R_i)-|R_i|\big)=\kappa_{R_i}(\mathbf a,\mathbf b)$, and the rest is $\mathcal G$ at $u$ times $u^{\bmu|_R-\one}$, that is, by \hb{thm:tworeadings} and $\sum_{j,p}\sigma_{j,p}=\sigma$, the stacked Gammas times the integrand of the joint period of the children at the masses $\mvec|_R$. For fixed $\mathbf y=(y_i)_{i<q}$, on a wall contour with $\re\kappa_{R_i}=0$ for $i<q$, the modulus of this integrand over $u\in\prod_i\Delta_{R_i}$ and the shares is the modulus of the stacked Gammas times the joint integrand at the real parts $(\re\mvec|_R;\mathbf c,\mathbf c')$, integrable by the hypothesis on $P^\otimes$ and by Stirling; so by Fubini the integral over $u$ and the shares is
\[
\hat g_1(\mathbf y)=\int_{\mathcal X_1}\prod_j\Gst{|\pat_j(B)|}1{\mathbf a_j}{s_j}\Gst{|\pat_j(B)|}1{\mathbf b_j}{t_j}\,P^\otimes_{F_B(\mathbf a,\mathbf b)}(\mvec|_R;\mathbf a,\mathbf b)\,e^{i\sum_{i<q}y_i\im\kappa_{R_i}(\mathbf a,\mathbf b)},
\]
$\mathcal X_1$ the fibre of (i), and $\int|\hat g_1(\mathbf y)|\,d\mathbf y$ is at most $P_F$ at the real parts, finite. The joint period is continuous in the shares and bounded by its value at the real parts, by dominated convergence. The map $(\mathbf a,\mathbf b)\mapsto(\im\kappa_{R_i})_{i<q}$ is affine on $\mathcal X_1$ with linear part onto $\mathbb R^{q-1}$, by the independence of the pins. So, as in (i), Fourier inversion in $\mathbf y$ gives $P_F=\int\hat g_1(\mathbf y)\,d\mathbf y$ as the fibre integral over $\mathcal X$, the pins $\kappa_{R_i}=0$, $i<q$, holding in real part on the contour and in imaginary part by inversion. For $q=1$ there is no $\mathbf y$, and $P_F=\hat g_1$ by Fubini.

\emph{(iii).} The real parts have a native gauge (first step). For real data the statement is the reference, proof of \hb{thm:wallsplit}(iii): the native split is a wall contour, lies on the real fibre, and $P^\otimes$ converges there. A wall contour, and the real parts of the pins of $\mathcal X$, $\re|\mathbf a_j|=\re s_j$, $\re|\mathbf b_j|=\re t_j$, $\re\langle\boldsymbol\delta\rangle=\re\bmu|_B$ and $\re\kappa_{R_i}=\kappa_{R_i}(\re\mvec|_R;\mathbf c,\mathbf c')$, involve the data only through their real parts. So the native split of the real parts is a wall contour for the data, an admissible real point of $\mathcal X$, and $P^\otimes_{F_B(\mathbf c,\mathbf c')}(\re\mvec|_R;\mathbf c,\mathbf c')$ converges there, as (ii) asks.
\end{proof}

\begin{hproposition}[\hb{prop:ct-pins} (The pins reach every integer point)]
As stated.
\end{hproposition}

\begin{proof}
The pins are a linear map $\mathbb Z^{\text{shares}}\to\mathbb Z^{\text{rows}}$, the rows being the share sums of the active letters and the loads at the marks $e\neq g$; a share $a_{j,T}$ has $1$ in the row of its sum and in the loads of the marks of $T-g$, and a share $b_{j,T}$ has $1$ in the row of its sum and in the loads of the marks of $\cE\setminus T$ other than $g$. Write $\mathbf e_x$ for the unit vector of the load of $x\neq g$ and $\mathbf e_g=0$. If $x\neq y$ lie in a common circuit of table $j$, \hb{lem:linkexchange} gives a basis $T\ni x$, $T\not\ni y$ with $T'=(T\setminus\{x\})\cup\{y\}$ a basis; if $K_j$ is active the difference of the columns of $a_{j,T}$ and $a_{j,T'}$ is $\mathbf e_x-\mathbf e_y$, and if $\Psi_j$ is active that of $b_{j,T}$ and $b_{j,T'}$ is $\mathbf e_y-\mathbf e_x$. Along the sequences of the hypothesis every $\mathbf e_x-\mathbf e_y$ lies in the image, and with $y=g$ every $\mathbf e_x$, $x\neq g$. Subtracting from the column of a share the loads it carries leaves the unit vector of its share sum. So the image is all of $\mathbb Z^{\text{rows}}$.

At integer data the target point is an integer point, so it has an integer preimage $x_0$, and the integer points of the fibre are $x_0+\Lambda$ with $\Lambda$ the kernel of the map on $\mathbb Z^{\text{shares}}$. The map is onto a free module of the rank of the number of rows, so the rows are independent and $\Lambda$, a subgroup of $\mathbb Z^{\text{shares}}$, is free of rank the number of shares less the number of rows, the dimension of the fibre. With a basis $v_1,\dots,v_D$ of $\Lambda$ the point $x_0+\sum_iz_iv_i$ is an integer point exactly when $z\in\mathbb Z^D$.
\end{proof}

\section{Chapter~\hb{ch:blocks}: bunches of width one: pooling and cleaning}\label{sec:ch18}

Throughout this section, for a bunch $B$ of width one of table $j$ whose weights factor across $B$ and $R$, $\theta_{j,T}=\theta'_{j,T\cap B}\,\theta''_{j,T\cap R}$, the minors $\cB_j|B$ and $\cB_j/R$ on $B$ carry $\theta'_j$ and the minors $\cB_j|R$ and $\cB_j/B$ on $R$ carry $\theta''_j$ (\hb{def:readings}), and every integral of a positive function is read in $[0,\infty]$ (Tonelli).

\begin{hproposition}[\hb{prop:classreadings} (Hinge values of parallel and series bunches)]
As stated, for classes of any table whose weights factor across them.
\end{hproposition}

\begin{proof}
That a parallel or series bunch has width at most one, and width zero exactly when it is a union of molecules, is the reference, proof of prop:classreadings. Let $B$ be a parallel bunch of width one: $r(B)=1$ and $c(B)=r(B)-d(B)=1$, so $d(B)=0$; $\cB/R$ has rank $0$, with the single basis $\varnothing$ and the weight $\theta'_\varnothing$, and $\cB|B$ has rank one, with the singletons of $B$ as bases, none of its marks being a loop, and the weights $\theta'_{\{e\}}=\theta'_\varnothing w_e$. Then $K_{\cB/R}=\theta'_\varnothing$, $K_{\cB|B}=\sum_e\theta'_{\{e\}}\alpha_e=\theta'_\varnothing\sum_ew_e\alpha_e$, $\Psi_{\cB/R}=\theta'_\varnothing\alpha^B=\theta'_\varnothing e_{|B|}(\alpha_B)$ and $\Psi_{\cB|B}=\sum_e\theta'_{\{e\}}\alpha^{B-e}=\theta'_\varnothing e_{|B|}(\alpha_B)\sum_ew_e\alpha_e^{-1}$, and the hinge values are the quotients of \hb{def:readings}. Let $B$ be a series bunch of width one: no mark of $B$ is a coloop, so $d(B)\ge|B|-1$, and width one gives $r(B)=|B|$ and $d(B)=|B|-1$; $\cB|B$ has the single basis $B$, with the weight $\theta'_B$, and $\cB/R$ the sets $B-e$, every one of them occurring since a basis avoids $e$, with the weights $\theta'_{B-e}=\theta'_Bw_e$. Then $K_{\cB|B}=\theta'_Be_{|B|}(\alpha_B)$, $K_{\cB/R}=\theta'_B\sum_ew_e\alpha^{B-e}=\theta'_Be_{|B|}(\alpha_B)\sum_ew_e\alpha_e^{-1}$, $\Psi_{\cB|B}=\theta'_B$ and $\Psi_{\cB/R}=\theta'_B\sum_ew_e\alpha_e$, which give $\gamma^\Psi_B$ and $\gamma^K_B$. In a list with $\theta'=\theta''=1$ every $w_e$ is $1$, and for a pair with the shape $(x,1-x)$: parallel, $\gamma^K_B=1$ and $\gamma^\Psi_B=(x^{-1}+(1-x)^{-1})^{-1}=x(1-x)$; series, $\gamma^\Psi_B=1$ and $\gamma^K_B=x(1-x)$.
\end{proof}

\begin{sremark}
A change of the factorization (\hb{def:readings}) multiplies the factors $w_e$ of a parallel bunch by $c$, the weights of its level $r(B)=1$ over those of its level $d(B)=0$, and those of a series bunch by $c^{-1}$, its patterns $B-e$ lying on the level $d(B)$; either way $\gamma^K_B$ is multiplied by $c$ and $\gamma^\Psi_B$ by $c^{-1}$. In exact arithmetic (receipt \texttt{check\_classes\_realization.py}): a parallel bunch of three marks whose cell of the bunch is $\binom{\{1,2,3,g\}}1$ with the weights $2,\frac37,5$ on the singletons of the bunch and $\frac43$ on $\{g\}$, and a series bunch of three marks whose cell is $\binom{\{1,2,3,g\}}3$ with the weights $\frac32$ on $\{1,2,3\}$ and $\frac15,7,\frac23$ on the sets through $g$, each glued to a weighted room cell, have the hinge values of the statement as rational functions, and $\varrho_B=(\sum_ew_e)^{-2}$ at the centre of the shape simplex of the parallel bunch.
\end{sremark}

\begin{htheorem}[\hb{thm:widthone} (A bunch of width one is one mark to its room)]
As stated; the converse holds with any number of products.
\end{htheorem}

\begin{proof}
All but the converse is the reference, proof of thm:widthone. On the open orthant let $K(\alpha)=\sum_{l=0}^kA_l(\alpha_B)C_l(\alpha_R)$, $C_l$ homogeneous of degree $D_l$. Grouping the bases by their level $i=|T\cap B|$, which takes every value $d(B)\le i\le r(B)$ (\hb{thm:window}), $K=\sum_iK^{(i)}$ with $K^{(i)}(\alpha_B,\lambda\alpha_R)=\lambda^{r-i}K^{(i)}(\alpha)$ and $K^{(i)}>0$ on the open orthant. At a fixed $\alpha$, with $\lambda\alpha_R$ for $\alpha_R$,
\[
\sum_{i=d(B)}^{r(B)}\lambda^{r-i}K^{(i)}(\alpha)=\sum_D\lambda^D\sum_{l:\ D_l=D}A_l(\alpha_B)\,C_l(\alpha_R)\qquad(\lambda>0),
\]
the sum on the right over the distinct values $D$ among the $D_l$. The functions $\lambda\mapsto\lambda^D$ with distinct real $D$ are linearly independent on $(0,\infty)$, so each of the $c(B)+1$ exponents $r-i$ on the left, whose coefficient is positive, is one of the $D_l$: $c(B)+1\le k+1$. With $k=1$ this is the converse of the reference.
\end{proof}

\begin{hlemma}[\hb{lem:mot-onehinge} (One hinge value of any degree is width at most one)]
As stated.
\end{hlemma}

\begin{proof}
Write $d=d(B)$, $c=c(B)$ and $K^{(k)}=\sum_{p\in\pat^{(k)}(B)}\alpha_B^{\,p}K_{(p)}(\alpha_R)$ for the level letters (\hb{def:levels}), nonzero exactly for $d\le k\le d+c$ (\hb{thm:window}); each child $K_{(p)}$ has positive coefficients, the weights. Let $a$ be the degree of $A$ and $P(\alpha_R,y)=\sum_{i=0}^DP_i(\alpha_R)y^i$. For $\tau>0$,
\[
\sum_k\tau^kK^{(k)}(\alpha)=K(\tau\alpha_B,\alpha_R)=\sum_{i=0}^D\tau^{a+i}A(\alpha_B)\gamma(\alpha_B)^iP_i(\alpha_R),
\]
and at a fixed $\alpha$ the functions $\tau\mapsto\tau^x$ with distinct real $x$ are linearly independent on $(0,\infty)$. So $a$ is an integer, $K^{(d)}$ being nonzero, and $K^{(k)}=A\gamma^{k-a}P_{k-a}$ for every $k$, the right side read as $0$ when $k-a\notin\{0,\dots,D\}$.

Let $d\le k\le d+c$. Then $P_{k-a}\neq0$; fixing $\alpha_R$ where $P_{k-a}(\alpha_R)\neq0$ shows that $A\gamma^{k-a}$ is a polynomial $\widehat f_k$ in $\alpha_B$, so $K^{(k)}(\alpha)=\widehat f_k(\alpha_B)P_{k-a}(\alpha_R)$ as polynomials. The coefficient of $\alpha_B^{\,p}\alpha_R^{\,X}$ in $K^{(k)}$ is the weight of $p\sqcup X$, distinct patterns giving distinct monomials $\alpha_B^{\,p}$, so the matrix of these coefficients, rows indexed by the patterns of the level and columns by the monomials of the room, has rank one: the children at the patterns of the level are positive multiples of one polynomial $C_k$, $K_{(p)}=\lambda_pC_k$ with $\lambda_p>0$, and $K^{(k)}=f_k(\alpha_B)C_k(\alpha_R)$ with $f_k=\sum_{p\in\pat^{(k)}(B)}\lambda_p\alpha_B^{\,p}$. Comparing with $A\gamma^{k-a}P_{k-a}$ at a fixed $\alpha_R$ where $C_k$ and $P_{k-a}$ do not vanish, $f_k=\mu_kA\gamma^{k-a}$ with a constant $\mu_k>0$.

Suppose $c\ge2$. Then $f_{d+1}^2=\mu\,f_df_{d+2}$ on the open orthant with $\mu=\mu_{d+1}^2/(\mu_d\mu_{d+2})>0$, an identity of polynomials. For a pattern $p$ of $d+1$ marks the monomial $\alpha_B^{\,2\one_p}$ has the coefficient $\lambda_p^2>0$ in $f_{d+1}^2$, since a product $\alpha_B^{\,p'}\alpha_B^{\,p''}$ of two monomials of sets equals it only for $p'=p''=p$; in $f_df_{d+2}$ it would need two sets $q,q'$ with $q=q'=p$, against $|q|=d\neq d+2=|q'|$. So $c\le1$.

Conversely, with the weights factoring across $B$ and $R$: at $c=0$ the one level consists of the unions of a pattern and a room side, all combinations occurring (\hb{thm:levels}), so $K=K_{\cB|B}(\alpha_B)K_{\cB|R}(\alpha_R)$ with the factored weights; at $c=1$ this is \hb{thm:widthone}. For $\Psi$, apply the statement to the dual table (\hb{prop:identity}), in which the window of $B$ is read on the cobases and $c(B)$ is unchanged (\hb{def:skeleton}, \hb{thm:window}).
\end{proof}

\begin{hproposition}[\hb{prop:hingefactor} (The freedom of the factorization)]
As stated.
\end{hproposition}

\begin{proof}
By \hb{thm:widthone} the bases of the level $d(B)$ are the sets $X\sqcup Y$ with $X\in\cB/R$ and $Y\in\cB|R$, and those of the level $r(B)$ the sets $X\sqcup Y$ with $X\in\cB|B$ and $Y\in\cB/B$, every such set being a basis. Let $\theta_{X\sqcup Y}=\theta'_X\theta''_Y=\phi'_X\phi''_Y$ be two factorizations and fix a basis $X_d\sqcup Y_d$ of the level $d(B)$. For every pattern $X$ of this level $X\sqcup Y_d$ is a basis, so $\phi'_X=\lambda_d\theta'_X$ with $\lambda_d=\theta''_{Y_d}/\phi''_{Y_d}$; for every room side $Y$ of this level $X_d\sqcup Y$ is a basis, so $\phi''_Y=\theta'_{X_d}\theta''_Y/\phi'_{X_d}=\lambda_d^{-1}\theta''_Y$. Likewise at the level $r(B)$ with a constant $\lambda_r$. Conversely any $\lambda_d,\lambda_r>0$ give a factorization, the patterns of the two levels having different sizes, and so have their room sides. The new factors multiply $K_{\cB/R}$ and $\Psi_{\cB/R}$ by $\lambda_d$, $K_{\cB|B}$ and $\Psi_{\cB|B}$ by $\lambda_r$, $K_{\cB|R}$ and $\Psi_{\cB|R}$ by $\lambda_d^{-1}$, and $K_{\cB/B}$ and $\Psi_{\cB/B}$ by $\lambda_r^{-1}$. With $c=\lambda_r/\lambda_d$,
\[
\gamma^K_B=\frac{K_{\cB|B}}{K_{\cB/R}}\mapsto c\,\gamma^K_B,\quad \gamma^\Psi_B=\frac{\Psi_{\cB/R}}{\Psi_{\cB|B}}\mapsto c^{-1}\gamma^\Psi_B,\quad \gamma^K_R=\frac{K_{\cB|R}}{K_{\cB/B}}\mapsto c\,\gamma^K_R,\quad \gamma^\Psi_R=\frac{\Psi_{\cB/B}}{\Psi_{\cB|R}}\mapsto c^{-1}\gamma^\Psi_R,
\]
so $\varrho_B=\gamma^\Psi_B/\gamma^K_B$ and $\varrho_R=\mathbf E_{\cB/B}/\mathbf E_{\cB|R}=\gamma^\Psi_R/\gamma^K_R$ are multiplied by $c^{-2}$, and $\gamma^K_R/\gamma^K_B$, $\gamma^\Psi_R/\gamma^\Psi_B$ and $\varrho_B/\varrho_R$ are unchanged.
\end{proof}

\begin{sremark}
The quotients $\gamma^K_R/\gamma^K_B$ and $\gamma^\Psi_R/\gamma^\Psi_B$ are the Skeleton values of \hb{cor:widthonestage}, which belong to the table and its gauge alone. What depends on the factorization is the split of a Skeleton value into a hinge value of the bunch and one of the room, and with it the cells, the cell families and the pooled families; the identities of the chapter hold at every factorization, each side changing with it.
\end{sremark}

Two lemmas are used below.

\begin{slemma}\label{lem:ch18-widthone}
Let a bunch $B$ of a table $(\cB,\theta)$ of rank $r$ on $\cE$, $n=|\cE|$, have width one, the weights factoring across $B$ and $R$. With $\alpha_B=\rho u$,
\begin{gather*}
\rho^{-d(B)}K(\rho u,\alpha_R)=K_{\cB/R}(u)\,K_{\cB|R}(\alpha_R)+\rho\,K_{\cB|B}(u)\,K_{\cB/B}(\alpha_R),\\
\rho^{-d^\ast(B)}\Psi(\rho u,\alpha_R)=\Psi_{\cB|B}(u)\,\Psi_{\cB/B}(\alpha_R)+\rho\,\Psi_{\cB/R}(u)\,\Psi_{\cB|R}(\alpha_R);
\end{gather*}
the cells have the ranks $r(B)=d(B)+1$ and $r-d(B)$ and the letters
\begin{gather*}
K_{\cB^+_B}=K_{\cB/R}\,(\alpha_g+\gamma^K_B),\quad \Psi_{\cB^+_B}=\Psi_{\cB|B}\,(\alpha_g+\gamma^\Psi_B),\\ K_{\cB^+_R}=K_{\cB|R}+\alpha_gK_{\cB/B},\quad \Psi_{\cB^+_R}=\alpha_g\Psi_{\cB|R}+\Psi_{\cB/B};
\end{gather*}
and $d(B)+d(R)=r-1$, $d^\ast(B)+d^\ast(R)=n-r-1$. For a family with masses $\mvec$ in which $B$ has width one in every table, $\kappa_B+\kappa_R=\kappa_{\cE}+2\sigma$, $\sigma$ the total strength.
\end{slemma}

\begin{proof}
By \hb{thm:widthone} the bases of the level $d(B)$ are the $X\sqcup Y$ with $X\in\cB/R$, $Y\in\cB|R$, and those of the level $r(B)=d(B)+1$ the $X\sqcup Y$ with $X\in\cB|B$, $Y\in\cB/B$, with $\theta_{X\sqcup Y}=\theta'_X\theta''_Y$ and the complement $(B\setminus X)\sqcup(R\setminus Y)$. Summing over the two levels, $K=K_{\cB/R}K_{\cB|R}+K_{\cB|B}K_{\cB/B}$ and $\Psi=\Psi_{\cB/R}\Psi_{\cB|R}+\Psi_{\cB|B}\Psi_{\cB/B}$, and on $B$ the polynomials $K_{\cB/R}$, $K_{\cB|B}$, $\Psi_{\cB|B}$, $\Psi_{\cB/R}$ are homogeneous of the degrees $d(B)$, $d(B)+1$, $|B|-r(B)=d^\ast(B)$ and $|B|-d(B)=d^\ast(B)+1$. The deletions and the contractions of the hinge are $\cB^+_B\setminus g=\cB|B$, $\cB^+_B/g=\cB/R$, $\cB^+_R\setminus g=\cB|R$ and $\cB^+_R/g=\cB/B$ (\hb{thm:widthone}), which gives the ranks and, by \hb{def:delcon} and \hb{def:readings}, the letters. By \hb{def:skeleton}, $d(S)=r-r(\cE\setminus S)$, so $d(B)+d(R)=2r-r(B)-r(R)=r-c(B)=r-1$, and $d^\ast(B)+d^\ast(R)=n-r(B)-r(R)=n-r-1$. Table by table, $\kappa_B+\kappa_R=m_{\cE}-\sum_j\big(s_j(r_j-1)+t_j(n-r_j-1)\big)=\kappa_{\cE}+\sum_j(s_j+t_j)$.
\end{proof}

\begin{slemma}\label{lem:ch18-section}
Let a family $F$ on $\cE=B\sqcup R$, $B$ and $R$ nonempty, have real data and be on its balance wall, and let $f(\alpha)=\alpha^{\mvec-\one}\prod_jK_j(\alpha)^{-s_j}\Psi_j(\alpha)^{-t_j}$. Then
\[
P_F(\mvec;\mathbf s,\mathbf t)=\int_{\Delta_R}\int_{\mathbb R^B_{>0}}f(x,v)\,dx\,d\omega_R(v)=\int_{\Delta_R}\int_{\Delta_B}\int_0^\infty\rho^{|B|-1}f(\rho u,v)\,d\rho\,d\omega_B(u)\,d\omega_R(v).
\]
\end{slemma}

\begin{proof}
On the balance wall $f$ is homogeneous of degree $\kappa_{\cE}-n=-n$, $n=|\cE|$. Let $\varphi$ be the indicator of $[1,e]$, so that $\int_0^\infty\varphi(\lambda)\,d\lambda/\lambda=1$, and compute $I=\int_{\mathbb R^{\cE}_{>0}}f(\alpha)\varphi(|\alpha_R|)\,d\alpha$ in two ways. With $\alpha=\lambda w$, $w\in\Delta_{\cE}$, $d\alpha=\lambda^{n-1}d\lambda\,d\omega_{\cE}(w)$ (\hb{lem:scaleshape}) and $f(\lambda w)=\lambda^{-n}f(w)$, so $I=\int_{\Delta_{\cE}}f(w)\int_0^\infty\varphi(\lambda|w_R|)\frac{d\lambda}\lambda\,d\omega_{\cE}(w)=\int_{\Delta_{\cE}}f\,d\omega_{\cE}=P_F$ (\hb{thm:projective}). With $\alpha_R=\lambda v$, $v\in\Delta_R$, and $\alpha_B=\lambda x$, $x\in\mathbb R^B_{>0}$, $d\alpha=\lambda^{|R|-1}d\lambda\,d\omega_R(v)\,\lambda^{|B|}dx$ and $f(\lambda x,\lambda v)=\lambda^{-n}f(x,v)$, so $I=\int_0^\infty\varphi(\lambda)\frac{d\lambda}\lambda\int_{\Delta_R}\int_{\mathbb R^B_{>0}}f(x,v)\,dx\,d\omega_R(v)$, the first form. The second is \hb{lem:scaleshape} for $x=\rho u$. For $|B|=1$, $\Delta_B$ is a point with the Dirac measure (\hb{def:conv}(iii)).
\end{proof}

\begin{htheorem}[\hb{thm:wallsglue} (Walls glue)]
As stated: (i) and (ii) for every family; in (iii) the identities for every family and the gluing when every active letter reads $B$ alike.
\end{htheorem}

\begin{proof}
The gluing of the counts of one table, the counts of its cells,
\begin{gather*}
d_{\cB^+_B}(S_B)=d_{\cB/R}(S_B),\quad d_{\cB^+_B}(S_B+g)=d_{\cB|B}(S_B),\\ d^\ast_{\cB^+_B}(S_B)=d^\ast_{\cB|B}(S_B),\quad d^\ast_{\cB^+_B}(S_B+g)=d^\ast_{\cB/R}(S_B),\\
d_{\cB^+_R}(S_R)=d_{\cB/B}(S_R),\quad d_{\cB^+_R}(S_R+g)=d_{\cB|R}(S_R),\\ d^\ast_{\cB^+_R}(S_R)=d^\ast_{\cB|R}(S_R),\quad d^\ast_{\cB^+_R}(S_R+g)=d^\ast_{\cB/B}(S_R),
\end{gather*}
and the inequalities $d_{\cB/R}\le d_{\cB|B}\le d_{\cB/R}+1$ and $d^\ast_{\cB|B}\le d^\ast_{\cB/R}\le d^\ast_{\cB|B}+1$ on the subsets of $B$, $d_{\cB/B}\le d_{\cB|R}\le d_{\cB/B}+1$ and $d^\ast_{\cB|R}\le d^\ast_{\cB/B}\le d^\ast_{\cB|R}+1$ on those of $R$, are the reference, proof of thm:wallsglue. For table $j$ and $S=S_B\sqcup S_R$ call
\begin{gather*}
a^K_j=d_{\cB_j|B}(S_B)-d_{\cB_j/R}(S_B),\quad b^K_j=d_{\cB_j|R}(S_R)-d_{\cB_j/B}(S_R),\\ a^\Psi_j=d^\ast_{\cB_j/R}(S_B)-d^\ast_{\cB_j|B}(S_B),\quad b^\Psi_j=d^\ast_{\cB_j/B}(S_R)-d^\ast_{\cB_j|R}(S_R),
\end{gather*}
all in $\{0,1\}$, in this proof the \emph{jumps} of the letters $K_j$ and $\Psi_j$ on the side of the bunch and on the side of the room. The gluing of the counts reads
\[
d_j(S)=d_{\cB_j/R}(S_B)+d_{\cB_j/B}(S_R)+\min(a^K_j,b^K_j),\qquad d^\ast_j(S)=d^\ast_{\cB_j|B}(S_B)+d^\ast_{\cB_j|R}(S_R)+\min(a^\Psi_j,b^\Psi_j).
\]
At $S_B=\varnothing$ the jumps on the side of the bunch are $0$, and at $S_R=\varnothing$ those on the side of the room. At $S_B=B$ those on the side of the bunch are $1$: the ranks give $d_{\cB_j|B}(B)=d_j(B)+1$, $d_{\cB_j/R}(B)=d_j(B)$, $d^\ast_{\cB_j/R}(B)=|B|-d_j(B)=d^\ast_j(B)+1$ and $d^\ast_{\cB_j|B}(B)=d^\ast_j(B)$. At $S_R=R$ those on the side of the room are $1$: $d_{\cB_j|R}(R)=r_j-d_j(B)=d_{\cB_j/B}(R)+1$ and $d^\ast_{\cB_j/B}(R)=|R|-r_j+d_j(B)+1=d^\ast_{\cB_j|R}(R)+1$.

(i) The tables of $F_{B\to g}(u)$ of index $j$ have the bases of $\cB^+_{j,R}$, at $(s_j,0)$ and at $(0,t_j)$, so its wall at $S\subseteq R+g$ is $m'_S-\sum_j\big(s_jd_{\cB^+_{j,R}}(S)+t_jd^\ast_{\cB^+_{j,R}}(S)\big)$, $\mvec'=(\kappa_B,\mvec|_R)$. At $S=S_R$ the counts of the cell are $d_{\cB_j/B}(S_R)$ and $d^\ast_{\cB_j|R}(S_R)$, which are $d_j(S_R)$ and $d^\ast_j(S_R)$ by the gluing at $S_B=\varnothing$: the wall is $\kappa_{S_R}$. At $S=S_R+g$ they are $d_{\cB_j|R}(S_R)$ and $d^\ast_{\cB_j/B}(S_R)$, and the gluing at $S_B=B$ gives $d_j(B\sqcup S_R)=d_j(B)+d_{\cB_j|R}(S_R)$ and $d^\ast_j(B\sqcup S_R)=d^\ast_j(B)+d^\ast_{\cB_j/B}(S_R)$; with $\kappa_B=m_B-\sum_j(s_jd_j(B)+t_jd^\ast_j(B))$ the wall is $\kappa_{B\sqcup S_R}$.

(ii) With $\alpha_{S_B}$ replaced by $\varepsilon\alpha_{S_B}$, $K_{\cB_j/R}$ is $\varepsilon^{d_{\cB_j/R}(S_B)}$ times a polynomial positive at $\varepsilon=0$ on the open orthant, its monomials having at least $d_{\cB_j/R}(S_B)$ marks in $S_B$ and some exactly that many (\hb{thm:window}); likewise $\Psi_{\cB_j|B}$ with $d^\ast_{\cB_j|B}(S_B)$. By the gluing at $S_R=\varnothing$ these are $d_j(S_B)$ and $d^\ast_j(S_B)$, so the exponent of the density at the face is $\kappa_{S_B}$.

(iii) The walls of $G_B$ and $G_R$ are read with the counts of the cells and the masses $(\mvec|_B,2\sigma-\kappa_B)$ and $(\kappa_B,\mvec|_R)$. With
\[
X=m_{S_B}+m_{S_R}-\sum_j\Big(s_j\big(d_{\cB_j/R}(S_B)+d_{\cB_j/B}(S_R)\big)+t_j\big(d^\ast_{\cB_j|B}(S_B)+d^\ast_{\cB_j|R}(S_R)\big)\Big),
\]
the counts of the cells give $\kappa^{G_B}_{S_B}+\kappa^{G_R}_{S_R}=X$ and, the two hinge masses adding up to $2\sigma=\sum_j(s_j+t_j)$, $\kappa^{G_B}_{S_B+g}+\kappa^{G_R}_{S_R+g}=X+\sum_j\big(s_j(1-a^K_j-b^K_j)+t_j(1-a^\Psi_j-b^\Psi_j)\big)$, while $\kappa_{S_B\sqcup S_R}=X-\sum_j\big(s_j\min(a^K_j,b^K_j)+t_j\min(a^\Psi_j,b^\Psi_j)\big)$. Since $1-a-b=(1-\max(a,b))-\min(a,b)$ for $a,b\in\{0,1\}$,
\begin{equation}\label{eq:S1820-glue}
\begin{aligned}
\kappa^{G_B}_{S_B}+\kappa^{G_R}_{S_R}-\kappa_{S_B\sqcup S_R}&=\sum_j\big(s_j\min(a^K_j,b^K_j)+t_j\min(a^\Psi_j,b^\Psi_j)\big),\\
\kappa^{G_B}_{S_B+g}+\kappa^{G_R}_{S_R+g}-\kappa_{S_B\sqcup S_R}&=\sum_j\big(s_j(1-\max(a^K_j,b^K_j))+t_j(1-\max(a^\Psi_j,b^\Psi_j))\big),
\end{aligned}
\end{equation}
at all masses and exponents. At $S_R=\varnothing$ the first line gives $\kappa^{G_B}_{S_B}=\kappa_{S_B}$, the wall at the empty set being $0$; at $S_B=\varnothing$ it gives $\kappa^{G_R}_{S_R}=\kappa_{S_R}$. The imbalance of $G_B$ is
\[
\kappa^{G_B}_{B+g}=m_B+2\sigma-\kappa_B-\sum_j\big(s_j(d_j(B)+1)+t_j(d^\ast_j(B)+1)\big)=2\sigma-\sum_j(s_j+t_j)=0,
\]
the cell $\cB^+_{j,B}$ having rank $d_j(B)+1$ on $|B|+1=d_j(B)+d^\ast_j(B)+2$ marks: $G_B$ is on its balance wall at all data. At $S_B=B$ the second line then gives $\kappa^{G_R}_{S_R+g}=\kappa_{B\sqcup S_R}$; at $S_R=R$ in particular $\kappa^{G_R}_{R+g}=\kappa_{\cE}$, so $G_R$ is on its balance wall exactly when $F$ is, and then the second line at $S_R=R$ gives $\kappa^{G_B}_{S_B+g}=\kappa_{S_B\sqcup R}$. These are the four identities.

Let now $s_j,t_j\ge0$ and let every active letter read $B$ alike. The jump of a letter on the side of the bunch at $S_B$ is the order to which its hinge value vanishes at the face $u_{S_B}=0$: by the expansion in (ii) applied to $\cB_j|B$ and $\cB_j/R$, with $\varepsilon\alpha_{S_B}$ for $\alpha_{S_B}$, $\gamma^K_{j,B}=K_{\cB_j|B}/K_{\cB_j/R}$ is $\varepsilon^{a^K_j}$ times a function positive at $\varepsilon=0$, and $\gamma^\Psi_{j,B}=\Psi_{\cB_j/R}/\Psi_{\cB_j|B}$ is $\varepsilon^{a^\Psi_j}$ times one. Constant multiples of one function vanish to one order at each face, so at $S_B$ the active letters have one jump $a\in\{0,1\}$ on the side of the bunch. If $a=0$ every term of the first line of \eqref{eq:S1820-glue} vanishes, a letter of exponent zero contributing nothing; if $a=1$ every term of the second does. Both lines are nonnegative, the exponents being nonnegative, so $\kappa_{S_B\sqcup S_R}$ is the smaller of the two sums. One active letter reads $B$ alike by itself.

Positivity. If both cell families have every proper wall positive, let $S=S_B\sqcup S_R$ be nonempty and proper. In the first sum a nonempty part is a proper subset of the marks of its cell family, $S_B\subseteq B\subsetneq B+g$ and $S_R\subseteq R\subsetneq R+g$, and an empty part has the wall $0$: the first sum is positive. In the second both $S_B+g$ and $S_R+g$ are nonempty, and at most one of them is all the marks of its cell family, $S$ being proper, its wall then being the imbalance $0$: the second sum is positive. So $\kappa_S>0$. Conversely, by the identities, $\kappa^{G_B}_{S_B}=\kappa_{S_B}$ for $\varnothing\neq S_B\subseteq B$ and $\kappa^{G_B}_{S_B+g}=\kappa_{S_B\sqcup R}$ for $S_B\subsetneq B$ are proper walls of $F$, and likewise for $G_R$; this direction does not use the hypothesis.
\end{proof}

\begin{sremark}
Without the reading alike, (iii) keeps the identities and loses the gluing. On $\cE=\{1,2,3,4\}$ with $B=\{1,2\}$, the family of the lists $\binom{\cE}3$ and $\binom{\cE}1$, each at $(0,1)$, has $B$ of width one in both tables, with the cells $\binom{\{1,2,g\}}2$, $\binom{\{3,4,g\}}2$ and $\binom{\{1,2,g\}}1$, $\binom{\{3,4,g\}}1$, and the hinge values $\gamma^\Psi_{1,B}=\alpha_1+\alpha_2$ and $\gamma^\Psi_{2,B}=\alpha_1\alpha_2/(\alpha_1+\alpha_2)$, which are not constant multiples of one function. At the masses $(\frac25,\frac85,\frac25,\frac85)$, on the balance wall, $\kappa_B=1$ and both cell families have every proper wall positive, the least being $\frac25$, while $\kappa_{\{1,3\}}=\frac25+\frac25-1=-\frac15$, against $\frac45$ from the gluing; the walls at $\{1,4\}$, $\{2,3\}$ and $\{2,4\}$ do not glue either. In exact rational arithmetic (receipt \texttt{check\_walls.py}): this family; the four identities at every subset at two partial leans of the same tables, where four walls do not glue; and the gluing at every subset for two tables sharing one weighted cell of the bunch at cobasis full leans, three pairs of exponents and three sets of masses, for a table at $(s,0)$ with the cell $\binom{\{1,2,g\}}1$ together with a table at $(0,t)$ with the cell $\binom{\{1,2,g\}}2$, whose hinge values $\gamma^K_{1,B}$ and $\gamma^\Psi_{2,B}$ are both $\alpha_1+\alpha_2$, and for $\binom{\cE}2$ with $B$ three of its four marks at two partial leans, where $\varrho_B=1$.
\end{sremark}

\begin{htheorem}[\hb{thm:cornerfactor} (Pooling factors at one hinge value)]
As stated: at real exponents of any sign, every active letter having the one hinge value $\gamma$.
\end{htheorem}

\begin{proof}
The walls $\kappa_{\{e\}}$ at the singletons are positive and the exponents real, so the masses are real, and the integrals below are of positive functions.

\emph{Bounds.} $\kappa_B$ and $\kappa_R$ are proper walls of $F$, hence positive, and $\kappa_B+\kappa_R=2\sigma$ on the balance wall (Lemma~\ref{lem:ch18-widthone}): $0<\kappa_B<2\sigma$.

\emph{The cell families.} By the identities of \hb{thm:wallsglue}(iii), which hold for every family on the balance wall, $G_B$ and $G_R$ are on their balance walls and each of their proper walls is a proper wall of $F$, hence positive; their periods converge (\hb{thm:projective}).

\emph{The family.} Put $\phi(u)=u^{\mvec|_B-\one}\prod_jK_{\cB_j/R}(u)^{-s_j}\Psi_{\cB_j|B}(u)^{-t_j}$ on $\Delta_B$. By \hb{thm:widthone}, with $K_{\cB_j/R}(\rho u)=\rho^{d_j(B)}K_{\cB_j/R}(u)$, $\Psi_{\cB_j|B}(\rho u)=\rho^{d^\ast_j(B)}\Psi_{\cB_j|B}(u)$ and the hinge values homogeneous of degree one, Lemma~\ref{lem:ch18-section} gives
\begin{multline*}
P_F(\mvec;\mathbf s,\mathbf t)=\int_{\Delta_B}\phi(u)\int_{\Delta_R}\int_0^\infty\rho^{\kappa_B-1}v^{\mvec|_R-\one}\\\times\prod_jK_{\cB^+_{j,R}}\big(\rho\gamma^K_{j,B}(u),v\big)^{-s_j}\Psi_{\cB^+_{j,R}}\big(\rho\gamma^\Psi_{j,B}(u),v\big)^{-t_j}d\rho\,d\omega_R(v)\,d\omega_B(u),
\end{multline*}
the powers of $\rho$ collecting to $\rho^{m_B-1-\sum_j(s_jd_j(B)+t_jd^\ast_j(B))}=\rho^{\kappa_B-1}$. Every active letter has the hinge value $\gamma(u)$, and a letter of exponent zero is the factor $1$. The substitution $x=\rho\gamma(u)$ at fixed $u$, $\rho^{\kappa_B-1}d\rho=\gamma(u)^{-\kappa_B}x^{\kappa_B-1}dx$, and Lemma~\ref{lem:ch18-section} for $G_R$ on $\{g\}\sqcup R$ at the masses $(\kappa_B,\mvec|_R)$ give
\[
P_F(\mvec;\mathbf s,\mathbf t)=P_{G_R}\big((\kappa_B,\mvec|_R);\mathbf s,\mathbf t\big)\int_{\Delta_B}\phi(u)\,\gamma(u)^{-\kappa_B}\,d\omega_B(u).
\]
\emph{The cell family of the bunch.} By Lemma~\ref{lem:ch18-widthone}, with $y$ the variable of $g$, $K_{\cB^+_{j,B}}(u,y)=K_{\cB_j/R}(u)\big(y+\gamma^K_{j,B}(u)\big)$ and $\Psi_{\cB^+_{j,B}}(u,y)=\Psi_{\cB_j|B}(u)\big(y+\gamma^\Psi_{j,B}(u)\big)$; every active letter has the factor $y+\gamma(u)$, and the exponents of the active letters add up to $\sum_j(s_j+t_j)=2\sigma$. Lemma~\ref{lem:ch18-section} for $G_B$ on $\{g\}\sqcup B$ at the masses $(\mvec|_B,2\sigma-\kappa_B)$, and the integral of a linear letter (the paragraph after \hb{prop:scalebreta}, with $c_0=\gamma(u)$, $c_1=1$ and $0<2\sigma-\kappa_B<2\sigma$), give
\begin{multline*}
P_{G_B}\big((\mvec|_B,2\sigma-\kappa_B);\mathbf s,\mathbf t\big)=\int_{\Delta_B}\phi(u)\int_0^\infty y^{2\sigma-\kappa_B-1}\big(y+\gamma(u)\big)^{-2\sigma}dy\,d\omega_B(u)\\=\Gst21{\kappa_B,\ 2\sigma-\kappa_B}{2\sigma}\int_{\Delta_B}\phi(u)\,\gamma(u)^{-\kappa_B}\,d\omega_B(u).
\end{multline*}
Multiplying the display for $P_F$ by $\Gst21{\kappa_B,\ 2\sigma-\kappa_B}{2\sigma}$ gives the statement, every factor being finite and positive. The parenthesis of the statement is \hb{def:readings}: the constants $\lambda_d,\lambda_r$ of table $j$ multiply $\gamma^K_{j,B}$ by $c$ and $\gamma^\Psi_{j,B}$ by $c^{-1}$, so a table with one active letter is normalized by its own constant; at full leans of all tables only the hinge values of different tables are compared; one table with both letters active and $\varrho_B$ constant is normalized by $c=\varrho_B^{1/2}$, and $\gamma^K_B=\gamma^\Psi_B$ forces $\varrho_B=1$. One table at $(0,t)$ has $s=0$ and $2\sigma=t$, and at $(s,0)$, $2\sigma=s$.
\end{proof}

\begin{sremark}
One table at $(0,t)$ has the one active letter $\Psi$, whose hinge value $\gamma^\Psi_B$ serves as $\gamma$ at any factorization; the cell families are the cells $(\cB^+_B,\theta')$ and $(\cB^+_R,\theta'')$ at $(0,t)$, and the display is that of the reference, $\Gst21{\kappa_B,\ t-\kappa_B}{t}P_F(\mvec;0,t)=P_{\cB^+_B}\big((\mvec|_B,t-\kappa_B);0,t\big)P_{\cB^+_R}\big((\kappa_B,\mvec|_R);0,t\big)$; at $(s,0)$ the same holds with $\gamma=\gamma^K_B$ and $s$ for $t$. The proof above is the proof of the reference with $\gamma^\Psi_B$ replaced by the common hinge value of the active letters. A family at full leans, every table at $(s_j,0)$ or $(0,t_j)$, has one active letter in each table, normalized by the factorization of its table; the hypothesis then asks only that the hinge values of the active letters of different tables be constant multiples of one another, a condition across the tables, and \hb{cor:fullleancore} is the same condition at every cell of a core. One table with $s,t>0$ has both letters active; the hypothesis is $\varrho_B$ constant, reached at the factorization with $\varrho_B=1$, and the display then holds at every pair $(s,t)$ with $2\sigma=s+t$, the separation of \hb{prop:readingsagree} completed by the Beta corner of the hinge.
\end{sremark}

\begin{hproposition}[\hb{prop:cleanrealization} (Pooling a realization)]
As stated, for every family and every realization.
\end{hproposition}

\begin{proof}
The masses of $\Upsilon$ are $\tau_e=\sum_j\big(s_j\alpha_e\partial_e\log K_j(\alpha)+t_j\alpha_e\partial_e\log\Psi_j(\alpha)\big)$ (\hb{def:realization}). With $u=\alpha_B/|\alpha_B|$ fixed, the tables of index $j$ of $F_{B\to g}(u)$, read as one table with two systems of weights, have the letters $K^{u}_j(\alpha_g,\alpha_R)=K_{\cB^+_{j,R}}(\gamma^K_{j,B}(u)\alpha_g,\alpha_R)$ and $\Psi^{u}_j(\alpha_g,\alpha_R)=\Psi_{\cB^+_{j,R}}(\gamma^\Psi_{j,B}(u)\alpha_g,\alpha_R)$ (\hb{thm:pooling}). By \hb{thm:widthone} and the homogeneity of the hinge values, $\gamma^K_{j,B}(\alpha_B)=\rho\gamma^K_{j,B}(u)$ with $\rho=|\alpha_B|$, and
\[
K_j(\alpha)=K_{\cB_j/R}(\alpha_B)\,K^{u}_j(\rho,\alpha_R),\qquad \Psi_j(\alpha)=\Psi_{\cB_j|B}(\alpha_B)\,\Psi^{u}_j(\rho,\alpha_R).
\]
For $e\in R$ the first factors do not depend on $\alpha_e$, so $\tau_e$ is the usage of $e$ in $F_{B\to g}(u)$ at the gauge $(\rho,\alpha_R)$ with the exponents $(s_j,t_j)$. Summing over $e\in B$, Euler's identity gives $\sum_{e\in B}\alpha_e\partial_e\log K_{\cB_j/R}(\alpha_B)=d_j(B)$ and $\sum_{e\in B}\alpha_e\partial_e\log\Psi_{\cB_j|B}(\alpha_B)=d^\ast_j(B)$, and, $\gamma^K_{j,B}$ being homogeneous of degree one,
\begin{multline*}
\sum_{e\in B}\alpha_e\partial_e\log K_{\cB^+_{j,R}}\big(\gamma^K_{j,B}(\alpha_B),\alpha_R\big)=\big[y\partial_y\log K_{\cB^+_{j,R}}(y,\alpha_R)\big]_{y=\gamma^K_{j,B}(\alpha_B)}\\=\big[\alpha_g\partial_{\alpha_g}\log K^{u}_j(\alpha_g,\alpha_R)\big]_{\alpha_g=\rho},
\end{multline*}
and likewise for $\Psi$. So $\tau_B=\sum_j\big(s_jd_j(B)+t_jd^\ast_j(B)\big)+\tau^u_g$, with $\tau^u_g$ the usage of $g$ in $F_{B\to g}(u)$ at $(\rho,\alpha_R)$, and $\tau^u_g$ is the wall $\kappa_B$ at the masses $\btau$. The exponents are those of a realization and $(\rho,\alpha_R)$ is a gauge, so $((s_j,t_j)_j;(\rho,\alpha_R))$ is a realization of $F_{B\to g}(u)$ with these masses, its usage. For one table at $(0,t)$, $\Psi^u(\rho,\alpha_R)=\Psi_{\cB^+_R}(\gamma^\Psi_B(\alpha_B),\alpha_R)$, so the usage of $(\cB^+_R,\theta'')$ at $(0,t)$ and at the gauge $(\gamma^\Psi_B(\alpha_B),\alpha_R)$ is $\btau|_R$ on $R$ and, by the chain rule above, $\tau_B-t\,d^\ast(B)$ at $g$.
\end{proof}

\begin{sremark}
The reference is one list at $(0,t)$, its last sentence. In exact rational arithmetic (receipt \texttt{check\_classes\_realization.py}): two tables on five marks, each the two-sum of a weighted cell of the bunch on $\{1,2,g\}$ and a weighted cell of the room on $\{3,4,5,g\}$, at the partial leans $(\frac23,\frac54)$ and $(3,\frac12)$ and a rational gauge; the usage of $F_{B\to g}(u)$ at $(|\alpha_B|,\alpha_R)$ is $(\kappa_B(\btau),\btau|_R)$.
\end{sremark}

\begin{hproposition}[\hb{prop:readingsagree} (Hinge values that agree)]
As stated, for families.
\end{hproposition}

\begin{proof}
As in the proof of \hb{thm:cornerfactor}, with $\phi(u)=u^{\mvec|_B-\one}\prod_jK_{\cB_j/R}(u)^{-s_j}\Psi_{\cB_j|B}(u)^{-t_j}$,
\begin{multline*}
P_F(\mvec;\mathbf s,\mathbf t)=\int_{\Delta_B}\phi(u)\int_{\Delta_R}\int_0^\infty\rho^{\kappa_B-1}v^{\mvec|_R-\one}\\\times\prod_jK_{\cB^+_{j,R}}\big(\rho\gamma^K_{j,B}(u),v\big)^{-s_j}\Psi_{\cB^+_{j,R}}\big(\rho\gamma^\Psi_{j,B}(u),v\big)^{-t_j}d\rho\,d\omega_R(v)\,d\omega_B(u).
\end{multline*}
For $s_j\neq0$, $K_{\cB^+_{j,R}}(\rho c_j\gamma(u),v)=K_{\cB|R}(v)+\rho c_j\gamma(u)K_{\cB/B}(v)$ is the letter $K$ of $\cB^+_{j,R}$ with its weights multiplied by $c_j$ on its bases through $g$, read at $(\rho\gamma(u),v)$; for $t_j\neq0$, $\Psi_{\cB^+_{j,R}}(\rho c^\ast_j\gamma(u),v)$ is the letter $\Psi$ of $\cB^+_{j,R}$ with its weights multiplied by $c^\ast_j$ on its bases avoiding $g$, read at $(\rho\gamma(u),v)$ (Lemma~\ref{lem:ch18-widthone}). These are the letters of $G^{\mathbf c}_R$, and the letters of exponent zero are the factor $1$. The substitution $x=\rho\gamma(u)$ and Lemma~\ref{lem:ch18-section} for $G^{\mathbf c}_R$, which has the bases and so the walls of $G_R$ and is on its balance wall by \hb{thm:wallsglue}(iii), give the display. Its period converges, each proper wall of $G_R$ being a proper wall of $F$, so the shape integral is finite.

For one table with $\varrho_B$ constant, \hb{prop:hingefactor} reaches $\varrho_B=1$ by $c=\varrho_B^{1/2}$; then $\gamma^K_B=\gamma^\Psi_B$, $c=c^\ast=1$ with $\gamma=\gamma^K_B$, and $G^{\mathbf c}_R$ is the cell $(\cB^+_R,\theta'')$ at $(s,t)$, which is the reference.
\end{proof}

\begin{htheorem}[\hb{thm:twosides} (Gluing across a bunch of width one)]
As stated: at nonnegative exponents, a letter of exponent zero being opened by nothing.
\end{htheorem}

\begin{proof}
\emph{The letters on the scale.} For table $j$ and $(u,v)\in\Delta_B\times\Delta_R$ put
\begin{gather*}
A^K_j=K_{\cB_j/R}(u)K_{\cB_j|R}(v),\quad C^K_j=K_{\cB_j|B}(u)K_{\cB_j/B}(v),\\ A^\Psi_j=\Psi_{\cB_j|B}(u)\Psi_{\cB_j/B}(v),\quad C^\Psi_j=\Psi_{\cB_j/R}(u)\Psi_{\cB_j|R}(v),
\end{gather*}
all positive, so that $\rho^{-d_j(B)}K_j(\rho u,v)=A^K_j+\rho C^K_j$ and $\rho^{-d^\ast_j(B)}\Psi_j(\rho u,v)=A^\Psi_j+\rho C^\Psi_j$ (Lemma~\ref{lem:ch18-widthone}). The minors of the level $d$ sit in $A^K_j$ and in $C^\Psi_j$, those of the level $r$ in $C^K_j$ and in $A^\Psi_j$.

\emph{The scale.} By Lemma~\ref{lem:ch18-section}, the powers of $\rho$ collecting to $\rho^{\kappa_B-1}$,
\begin{gather*}
P_F(\mvec;\mathbf s,\mathbf t)=\int_{\Delta_R}\int_{\Delta_B}u^{\mvec|_B-\one}v^{\mvec|_R-\one}\,I(u,v)\,d\omega_B(u)\,d\omega_R(v),\\ I(u,v)=\int_0^\infty\rho^{\kappa_B-1}\prod_{\ell=1}^p(A_\ell+\rho C_\ell)^{-w_\ell}d\rho ,
\end{gather*}
where $\ell$ runs over the $p$ active letters, with their exponents $w_\ell>0$ and their coefficients $(A_\ell,C_\ell)$, a letter of exponent zero being the factor $1$. Here $\sum_\ell w_\ell=\sum_j(s_j+t_j)=2\sigma$, and $\kappa_B$ and $\kappa_R=2\sigma-\kappa_B$ (Lemma~\ref{lem:ch18-widthone}, $\kappa_{\cE}=0$) are proper walls: $0<\kappa_B<2\sigma$.

\emph{Opening the scale.} If $p=1$, the integral of a linear letter (the paragraph after \hb{prop:scalebreta}) gives $I=\Gst21{w_1-\kappa_B,\ \kappa_B}{w_1}A_1^{\kappa_B-w_1}C_1^{-\kappa_B}$. If $p\ge2$, choose $c_\ell\in(0,w_\ell)$ for $\ell<p$ with $0<\kappa_B-\sum_{\ell<p}c_\ell<w_p$, possible as $0<\kappa_B<\sum_\ell w_\ell$, and open each letter $\ell<p$ by \hb{thm:E1} with $k=2$ on the line $\re x_\ell=c_\ell$:
\[
(A_\ell+\rho C_\ell)^{-w_\ell}=\frac1{2\pi i}\int_{c_\ell+i\mathbb R}\Gst21{w_\ell-x_\ell,\ x_\ell}{w_\ell}A_\ell^{x_\ell-w_\ell}(\rho C_\ell)^{-x_\ell}\,dx_\ell .
\]
The modulus of the integrand is its value at the real parts times the modulus of the stacked Gamma, which decays exponentially along the line (Stirling), and $\int_0^\infty\rho^{\kappa_B-\sum_{\ell<p}c_\ell-1}(A_p+\rho C_p)^{-w_p}d\rho$ is finite; so Fubini puts the integral over $\rho$ inside, where the integral of a linear letter at the complex index $x_p=\kappa_B-\sum_{\ell<p}x_\ell$, $0<\re x_p<w_p$, gives $\Gst21{w_p-x_p,\ x_p}{w_p}A_p^{x_p-w_p}C_p^{-x_p}$. Hence
\begin{gather*}
I(u,v)=\frac1{(2\pi i)^{p-1}}\int_{c_1+i\mathbb R}\!\!\cdots\int_{c_{p-1}+i\mathbb R}\ \prod_{\ell=1}^p\Gst21{w_\ell-x_\ell,\ x_\ell}{w_\ell}A_\ell^{x_\ell-w_\ell}C_\ell^{-x_\ell}\,dx_1\cdots dx_{p-1},\\ x_p=\kappa_B-\sum_{\ell<p}x_\ell ,
\end{gather*}
the share integral $\int_{x_1+\cdots+x_p=\kappa_B}$ of \hb{def:fibre}: the factor $(2\pi i)^{-1}$ of each of the $p-1$ openings is its normalization $(2\pi i)^{1-p}$, and its dimension is $p-1$. At $p=1$ the share integral has dimension zero, one share and the pin $P=(1)$, and \hb{def:fibre} gives the value of the integrand at the one admissible point $x_1=\kappa_B$: the formula for $p=1$ above.

\emph{The level shares.} For the letter $K_j$ write $x_\ell=a^r_j$ and $w_\ell-x_\ell=a^d_j$, for $\Psi_j$ write $x_\ell=b^d_j$ and $w_\ell-x_\ell=b^r_j$, and give the letters of exponent zero the shares $0$. The pin $\sum_\ell x_\ell=\kappa_B$ is $\sum_j(a^r_j+b^d_j)=\kappa_B$; the factor of $K_j$ is $\Gst21{a^d_j,\ a^r_j}{s_j}(A^K_j)^{-a^d_j}(C^K_j)^{-a^r_j}$ and that of $\Psi_j$ is $\Gst21{b^d_j,\ b^r_j}{t_j}(C^\Psi_j)^{-b^d_j}(A^\Psi_j)^{-b^r_j}$. The product of the powers over all letters, times $u^{\mvec|_B-\one}v^{\mvec|_R-\one}$, is $f(u)\,g(v)$, with
\[
f(u)=u^{\mvec|_B-\one}\prod_jK_{\cB_j/R}^{-a^d_j}\Psi_{\cB_j/R}^{-b^d_j}K_{\cB_j|B}^{-a^r_j}\Psi_{\cB_j|B}^{-b^r_j},\qquad g(v)=v^{\mvec|_R-\one}\prod_jK_{\cB_j|R}^{-a^d_j}\Psi_{\cB_j|R}^{-b^d_j}K_{\cB_j/B}^{-a^r_j}\Psi_{\cB_j/B}^{-b^r_j},
\]
the integrands of the periods of $F^{\mathrm{bunch}}(\mathbf a,\mathbf b)$ on $\Delta_B$ and of $F^{\mathrm{room}}(\mathbf a,\mathbf b)$ on $\Delta_R$ (\hb{thm:projective}).

\emph{Balance on the pin.} The tables $\cB_j/R$ and $\cB_j|B$ on $B$ have the ranks $d_j(B)$ and $d_j(B)+1$, and $\cB_j|R$ and $\cB_j/B$ on $R$ the ranks $r_j-d_j(B)$ and $r_j-d_j(B)-1$. With $a^d_j+a^r_j=s_j$, $b^d_j+b^r_j=t_j$ and $|B|-d_j(B)-1=d^\ast_j(B)$, the imbalance of $F^{\mathrm{bunch}}(\mathbf a,\mathbf b)$ at $\mvec|_B$ is
\[
m_B-\sum_j\big(a^d_jd_j(B)+b^d_j(|B|-d_j(B))+a^r_j(d_j(B)+1)+b^r_jd^\ast_j(B)\big)=\kappa_B-\sum_j(a^r_j+b^d_j);
\]
with $|R|+1-r_j+d_j(B)=n-r_j-d^\ast_j(B)$, that of $F^{\mathrm{room}}(\mathbf a,\mathbf b)$ at $\mvec|_R$ is
\[
m_R-\sum_j\big(s_j(r_j-d_j(B))+t_j(|R|+1-r_j+d_j(B))\big)+\sum_j(a^r_j+b^d_j)=\kappa_{\cE}-\kappa_B+\sum_j(a^r_j+b^d_j).
\]
Both vanish on the pin, $F$ being on its balance wall.

\emph{Fubini.} On a contour of the statement the modulus of $f(u)g(v)$ is $f(u)g(v)$ at the real parts of the shares, which do not move along the contour, and its integral over $\Delta_B\times\Delta_R$ is the product of the two periods at the real parts, finite by hypothesis; the stacked Gammas decay exponentially along the contour. So Fubini exchanges the share integral with the integral over $\Delta_B\times\Delta_R$, and the inner integrals are $P_{F^{\mathrm{bunch}}(\mathbf a,\mathbf b)}(\mvec|_B;\mathbf a,\mathbf b)$ and $P_{F^{\mathrm{room}}(\mathbf a,\mathbf b)}(\mvec|_R;\mathbf a,\mathbf b)$: the general display.

\emph{Pair coordinates.} When every $s_j$ and $t_j$ is positive, $p=2N$ and the shares are $(a^r_j,b^d_j)_j$. With $a^d_j=\sigma^d_j+\delta^d_j$, $b^d_j=\sigma^d_j-\delta^d_j$ and $a^r_j=s_j-a^d_j$, the affine map $(\sigma^d_j,\delta^d_j)_j\mapsto(a^r_j,b^d_j)_j$ has in each table the linear part $(y_\sigma,y_\delta)\mapsto(-y_\sigma-y_\delta,\ y_\sigma-y_\delta)$, of determinant $2$. So the imaginary parts $\mathbf y$ of the shares and $\mathbf y'$ of the pair coordinates have $d\mathbf y=2^Nd\mathbf y'$ and $\sum_\ell y_\ell=-2\sum_jy'_{\delta,j}$, and with $\im\kappa_B=0$ and $\delta^1(-2z)=\frac12\delta^1(z)$, \hb{def:fibre} gives
\[
\int_{\sum_\ell x_\ell=\kappa_B}h=(2\pi)^{1-2N}\int_{\mathbb R^{2N}}h\,\delta^1\Big(\sum_\ell y_\ell\Big)d\mathbf y=2^{N-1}(2\pi)^{1-2N}\int_{\mathbb R^{2N}}h\,\delta^1\Big(\sum_jy'_{\delta,j}\Big)d\mathbf y',
\]
$2^{N-1}$ times the fibre integral in the pair coordinates pinned by $\sum_j\delta^d_j=\frac12(\sum_js_j-\kappa_B)$, the image of the pin $\sum_j(a^r_j+b^d_j)=\sum_j(s_j-2\delta^d_j)=\kappa_B$, the admissible real points corresponding under the map. With $\Gst21{a^d_j,\ a^r_j}{s_j}\Gst21{b^d_j,\ b^r_j}{t_j}=\Gst42{a^d_j,\ b^d_j,\ a^r_j,\ b^r_j}{s_j,\ t_j}$ this is the first display, whose $N$ coordinates $\sigma^d_j$ are not pinned. For one table the pin fixes $\delta^d$, $2^{N-1}=1$, and the fibre integral is $(2\pi)^{-1}\int dy=(2\pi i)^{-1}\int d\sigma^d$ along the line in $\sigma^d$. The other forms of the pin, the value $2\delta^d=s-\kappa_B=\delta(1+\varepsilon(B))-\mu_B$ (\hb{thm:pairform}(iii) with $c(B)=1$) and the swap of $B$ and $R$ are the reference, proof of thm:twosides.
\end{proof}

\begin{sremark}
The first display needs both letters of every table active: a letter of exponent zero would put $\Gamma(0)$ into the stacked Gamma of its table. The general display carries one Beta corner $\Gst21{\cdot,\ \cdot}{w}$ for each active letter, and at a full lean of one table it has no free share, which is \hb{cor:pinnedproduct}.
\end{sremark}

\begin{hcorollary}[\hb{cor:pinnedproduct} (One Skeleton value pins the share)]
As stated: real exponents of any sign, one Skeleton value for all active letters.
\end{hcorollary}

\begin{proof}
$0<\kappa_B<2\sigma$ as in the proof of \hb{thm:cornerfactor}. With the coefficients of the proof of \hb{thm:twosides}, Lemma~\ref{lem:ch18-section} gives, all functions being positive,
\begin{multline*}
P_F(\mvec;\mathbf s,\mathbf t)=\int_{\Delta_R}\int_{\Delta_B}u^{\mvec|_B-\one}v^{\mvec|_R-\one}\\\times\int_0^\infty\rho^{\kappa_B-1}\prod_j(A^K_j+\rho C^K_j)^{-s_j}(A^\Psi_j+\rho C^\Psi_j)^{-t_j}d\rho\,d\omega_B(u)\,d\omega_R(v).
\end{multline*}
The Skeleton values are $A^K_j/C^K_j=\gamma^K_{j,R}(v)/\gamma^K_{j,B}(u)$ and $A^\Psi_j/C^\Psi_j=\gamma^\Psi_{j,R}(v)/\gamma^\Psi_{j,B}(u)$ (\hb{cor:widthonestage}), and by the hypothesis, the constants cancelling, every active letter has $A_\ell=\xi C_\ell$ with $\xi=\gamma_R(v)/\gamma(u)$. The exponents of the active letters add up to $2\sigma$, so the integral of a linear letter gives $\int_0^\infty\rho^{\kappa_B-1}(\xi+\rho)^{-2\sigma}d\rho=\Gst21{\kappa_B,\ 2\sigma-\kappa_B}{2\sigma}\xi^{\kappa_B-2\sigma}$, and
\begin{multline*}
P_F(\mvec;\mathbf s,\mathbf t)=\Gst21{\kappa_B,\ 2\sigma-\kappa_B}{2\sigma}\int_{\Delta_R}\int_{\Delta_B}u^{\mvec|_B-\one}v^{\mvec|_R-\one}\\\times\prod_{s_j\neq0}(C^K_j)^{-s_j}\prod_{t_j\neq0}(C^\Psi_j)^{-t_j}\,\xi^{\kappa_B-2\sigma}\,d\omega_B(u)\,d\omega_R(v).
\end{multline*}
For level shares as in the statement, $(A^K_j)^{-a^d_j}(C^K_j)^{-a^r_j}=(C^K_j)^{-s_j}\xi^{-a^d_j}$ for an active $K_j$, $(A^\Psi_j)^{-b^r_j}(C^\Psi_j)^{-b^d_j}=(C^\Psi_j)^{-t_j}\xi^{-b^r_j}$ for an active $\Psi_j$, and an inactive letter contributes $1$; since $\sum_j(a^d_j+b^r_j)=2\sigma-\sum_j(a^r_j+b^d_j)=2\sigma-\kappa_B$, the integrand above is $f(u)g(v)$ with $f$ and $g$ the integrands of the periods of $F^{\mathrm{bunch}}(\mathbf a,\mathbf b)$ and $F^{\mathrm{room}}(\mathbf a,\mathbf b)$ (proof of \hb{thm:twosides}), whatever the shares. At complex shares $|f(u)g(v)|$ is $f(u)g(v)$ at the real parts, which satisfy the same constraints; its integral is finite by the display at the real parts, and both factors are positive, so $|f|$ and $|g|$ are integrable and Fubini gives the statement, the product of the two periods not depending on the shares. With one active letter the hypothesis holds with $\gamma$ and $\gamma_R$ its hinge values of $B$ and of $R$. One table at $(0,t)$: $\mathbf a=0$, $b^d+b^r=t$ and $b^d=\kappa_B$, so the level shares are $(0,\kappa_B,0,t-\kappa_B)$ and the two families have the forms of the statement; at $(s,0)$, $a^r=\kappa_B$ and $a^d=s-\kappa_B$.
\end{proof}

\begin{sremark}
The reference, rem:pinnedproduct, is one table at a full lean, the case of one active letter. The hypothesis does not depend on the factorizations: a change of the factorization of table $j$ multiplies the hinge values of a letter of table $j$ at $B$ and at $R$ by one constant (\hb{def:readings}), so the Skeleton values are unchanged.
\end{sremark}

\begin{hlemma}[\hb{lem:slopesrange} (Slopes and the range of the ratios)]
As stated: for every table of a family, both letters of every table active.
\end{hlemma}

\begin{proof}
With $a^\ell_j=\sigma^\ell_j+\delta^\ell_j$ and $b^\ell_j=\sigma^\ell_j-\delta^\ell_j$, $K^{-a}\Psi^{-b}=\mathbf E^{-\sigma}(K/\Psi)^{-\delta}$ for every minor, $\mathbf E=K\Psi$. So the factors of table $j$ in the word of the bunch family (proof of \hb{thm:twosides}) are
\[
\mathbf E_{\cB_j/R}^{-\sigma^d_j}\mathbf E_{\cB_j|B}^{-\sigma^r_j}\Big(\frac{K_{\cB_j/R}}{\Psi_{\cB_j/R}}\Big)^{-\delta^d_j}\Big(\frac{K_{\cB_j|B}}{\Psi_{\cB_j|B}}\Big)^{-\delta^r_j},\qquad \mathbf E_{\cB_j/R}^{-\sigma^d_j}\mathbf E_{\cB_j|B}^{-\sigma_j+\sigma^d_j}=\varrho_{j,B}^{-\sigma^d_j}\mathbf E_{\cB_j|B}^{-\sigma_j},
\]
and in the word of the room family $\mathbf E_{\cB_j|R}^{-\sigma^d_j}\mathbf E_{\cB_j/B}^{-\sigma_j+\sigma^d_j}=\varrho_{j,R}^{\sigma^d_j}\mathbf E_{\cB_j/B}^{-\sigma_j}$ (\hb{def:readings}), the other factors not involving $\sigma^d_j$ at fixed $\delta^d_j$: $\sigma^d_j$ enters through $\varrho_{j,B}(u)^{-\sigma^d_j}\varrho_{j,R}(v)^{\sigma^d_j}$.

Slopes. The wall at $S\subseteq B$ of the bunch family is $m_S-\sum_j\big(a^d_jd_{\cB_j/R}(S)+b^d_jd^\ast_{\cB_j/R}(S)+a^r_jd_{\cB_j|B}(S)+b^r_jd^\ast_{\cB_j|B}(S)\big)$, affine in the shares; raising $\sigma^d_j$ by one at fixed $\delta^d_j$ raises $a^d_j,b^d_j$ and lowers $a^r_j,b^r_j$ by one, so the slope is $c_{\cB_j/R}(S)-c_{\cB_j|B}(S)$, $c=|S|-d-d^\ast$. With $r^+$ the rank function of $\cB^+_{j,B}$, $\cB_j|B=\cB^+_{j,B}\setminus g$ has the rank function $r^+$ on $B$ and $\cB_j/R=\cB^+_{j,B}/g$ the rank function $Z\mapsto r^+(Z+g)-1$, and $c_Z(S)=r_Z(S)+r_Z(B\setminus S)-r_Z(B)$ for a table $Z$ on $B$ (\hb{def:skeleton}, \hb{thm:window}); the difference is the first displayed slope. On $R$ the table of the level $d$ is $\cB_j|R=\cB^+_{j,R}\setminus g$ and that of the level $r$ is $\cB_j/B=\cB^+_{j,R}/g$, which gives the second.

Range. On the sector of an order of $B$, $\alpha_{\pi(1)}=\rho_0$ and $\alpha_{\pi(k)}=\rho_0q_2\cdots q_k$ with $q_k\in(0,1]$ and tails $U_k$, the letters of a table $Z$ on $B$ are $K_Z=\rho_0^{r_Z}\prod_kq_k^{d_Z(U_k)}k_{Z,\pi}$ and $\Psi_Z=\rho_0^{|B|-r_Z}\prod_kq_k^{d^\ast_Z(U_k)}\psi_{Z,\pi}$, with sector letters positive on the closed cube (\hb{thm:exist}(ii)), hence between positive constants. So
\[
\varrho_{j,B}=\frac{\mathbf E_{\cB_j/R}}{\mathbf E_{\cB_j|B}}=\prod_kq_k^{\,c_{\cB_j|B}(U_k)-c_{\cB_j/R}(U_k)}\times(\text{a function between positive constants}),
\]
the exponent of $q_k$ being minus the slope of $U_k$. Every nonempty proper subset of $B$ is a tail of some order, the $q_k$ range over $(0,1]$ independently, and finitely many sectors cover $\Delta_B$: $\varrho_{j,B}$ is bounded exactly when no slope on $B$ is positive (a positive slope at a tail $U_k$ sends it to $\infty$ as $q_k\to0$), and bounded below by a positive constant exactly when none is negative. On $R$, $\varrho_{j,R}=\mathbf E_{\cB_j/B}/\mathbf E_{\cB_j|R}$ carries $q_k$ to the power of the slope of $U_k$: $\inf\varrho_{j,R}>0$ exactly when no slope on $R$ is positive and $\sup\varrho_{j,R}<\infty$ exactly when none is negative. A wall falls as $\sigma^d_j$ decreases, toward the level $d$, exactly when its slope is positive, which forces $\sup\varrho_{j,B}=\infty$ or $\inf\varrho_{j,R}=0$, so $\sup\varrho_{j,B}/\inf\varrho_{j,R}=\infty$; toward the level $r$, a negative slope forces $\inf\varrho_{j,B}/\sup\varrho_{j,R}=0$. A change of the factorization of table $j$ multiplies $\varrho_{j,B}$ and $\varrho_{j,R}$ by one constant $c^{-2}$ (\hb{prop:hingefactor}), which keeps the quotient, the finiteness of the suprema, the positivity of the infima and the two ratios; the slopes are counts.
\end{proof}

\begin{hproposition}[\hb{prop:gaugeclose} (The same shares closed at the gauge)]
As stated: at complex exponents and masses, with the bound in absolute value.
\end{hproposition}

\begin{proof}
At width one the letters on the scale are linear in $\rho$ (Lemma~\ref{lem:ch18-widthone}), and divided by their values at $\rho=1$ they are the Skeleton polynomials of the statement (\hb{def:scaleletters}): with the coefficients of the proof of \hb{thm:twosides}, $\pi_{j,d}=A^K_j/(A^K_j+C^K_j)$, $\pi_{j,r}=1-\pi_{j,d}$, $\pi^\ast_{j,r}=A^\Psi_j/(A^\Psi_j+C^\Psi_j)$ and $\pi^\ast_{j,d}=1-\pi^\ast_{j,r}$, the chances of the levels under the basis law and the cobasis law at $(u,v)$, all in $(0,1)$. Each polynomial is $1-(1-x)\epsilon$ with a constant exceedance polynomial $\epsilon$: $\pi_{j,r}$ for $K_j$ and $\pi^\ast_{j,d}$ for $\Psi_j$ on the bunch, and, read from the top of the scale, $\pi_{j,d}$ and $\pi^\ast_{j,r}$ on the room. In \hb{thm:sectors}(ii) only $k=0$ remains, $[x^k]\prod_i\epsilon_i^{n_i}=0$ for $k>0$, and the term of $\mathbf n$ is $\prod_i\frac{(w_i)_{n_i}}{n_i!}\epsilon_i^{n_i}\Gst21{\kappa,\ |\mathbf n|+1}{\kappa+|\mathbf n|+1}$; the two sectors of \hb{thm:sectors}(i) give the display, of nonnegative terms for real data with $s_j,t_j\ge0$. For complex exponents and masses with $\re\kappa_B>0$ and $\re\kappa^{\mathrm{top}}_B>0$ the series converges absolutely to the same value (\hb{thm:sectors}, last sentence; \hb{thm:scale} for the convergence of the scale integral). The bound: $|(w)_n|\le(|w|)_n$; $\prod_j\pi_{j,r}^{n_j}\pi^{\ast\,n'_j}_{j,d}\le\lambda_B^M$; $\big|\Gst21{\kappa_B,\ M+1}{\kappa_B+M+1}\big|=\big|\int_0^1x^{\kappa_B-1}(1-x)^Mdx\big|\le\int_0^1x^{\re\kappa_B-1}dx=1/\re\kappa_B$; and $\sum_{|\mathbf n|+|\mathbf n'|=M}\prod_j\frac{(|s_j|)_{n_j}(|t_j|)_{n'_j}}{n_j!\,n'_j!}=\frac{(\sum_j(|s_j|+|t_j|))_M}{M!}$, the coefficient of $z^M$ in $\prod_j(1-z)^{-|s_j|-|t_j|}$. Likewise on the room with $\lambda_R$ and $\re\kappa^{\mathrm{top}}_B$. The bounds are summable over $M$, $\sum_M\frac{(W)_M}{M!}\lambda^M=(1-\lambda)^{-W}$ for $\lambda<1$.
\end{proof}

\begin{hproposition}[\hb{prop:loopsimplex} (A loop or a coloop on the simplex)]
As stated: a mark that is a loop or a coloop in every table of a family.
\end{hproposition}

\begin{proof}
In a table where $e$ is a loop no basis contains $e$, so $K_j=K_{j,\setminus e}$ and $\Psi_j=\alpha_e\Psi_{j,\setminus e}$; where $e$ is a coloop every basis contains $e$, so $K_j=\alpha_eK_{j,/e}$ and $\Psi_j=\Psi_{j,/e}$ (\hb{def:delcon}). The letters of $F'$ (\hb{cor:content}) are these minors, of the ranks $r'_j=r_j$ at a loop and $r_j-1$ at a coloop on $n-1$ marks, so the family word is $\alpha_e^{-\sum_{\eta_j=+1}s_j-\sum_{\eta_j=-1}t_j}\,W'(\alpha_{\cE-e})$, with $W'$ the word of $F'$, homogeneous of degree $\kappa'-m_{\cE-e}$, $\kappa'$ the imbalance of $F'$ at $\mvec|_{\cE-e}$. On $\Delta_{\cE}$ put $u=(u_e,(1-u_e)v)$, $u_e\in(0,1)$, $v\in\Delta_{\cE-e}$, so that $u^{\mvec-\one}d\omega_{\cE}(u)=u_e^{m_e-1}(1-u_e)^{m_{\cE-e}-1}du_e\,v^{\mvec|_{\cE-e}-\one}d\omega_{\cE-e}(v)$ (\hb{thm:aggregation}). Then
\[
P_F(\mvec;\mathbf s,\mathbf t)=\int_0^1u_e^{\nu_e-1}(1-u_e)^{\kappa'-1}du_e\ P_{F'}(\mvec|_{\cE-e};\mathbf s,\mathbf t),\qquad \nu_e=m_e-\sum_{\eta_j=+1}s_j-\sum_{\eta_j=-1}t_j ,
\]
for real data in $[0,\infty]$, and the integral over $u_e$ is $\Gst21{\nu_e,\ \kappa'}{\nu_e+\kappa'}$ where $\nu_e,\kappa'>0$. This $\nu_e$ is the wall $\kappa_{\{e\}}$, since $d_j(\{e\})=1$ and $d^\ast_j(\{e\})=0$ at a coloop and the reverse at a loop, and by \hb{thm:pairform}(iii), with $c_j(\{e\})=0$ and $\varepsilon_j(\{e\})=\eta_j$, it is $\mu_e-\sum_j\eta_j\delta_j$. The imbalance of $F'$ is $m_{\cE}-m_e-\sum_{\eta_j=+1}\big(s_j(r_j-1)+t_j(n-r_j)\big)-\sum_{\eta_j=-1}\big(s_jr_j+t_j(n-1-r_j)\big)=\kappa_{\cE}-\nu_e$, so $\nu_e+\kappa'=\kappa_{\cE}$. Several such marks are removed one at a time, each keeping its kind in the minors, and the product of the Beta corners is $\Gst{k+1}1{(\nu_{e_i})_i,\ \kappa'}{\kappa_{\cE}}$, split by \hb{thm:aggregation} as stated. On the balance wall $\kappa'=-\nu_e$, and $\nu_e$ and $\kappa'$ are not both positive.

The set $S=\{e\}$ has $c_j(S)=0$ in every table, as \hb{thm:offdomain} asks; there $\kappa_{\{e\}}=\nu_e$ and, by \hb{thm:pairform}(iii) and \hb{thm:window} ($c_j(\cE-e)=c_j(\{e\})=0$, $\varepsilon_j(\cE-e)=2r_j-n-\eta_j$), $\kappa_{\cE\setminus e}=\kappa_{\cE}-\nu_e=\kappa'$. The integral of \hb{thm:offdomain} over $\Delta_{\{e\}}\times\Delta_{\cE-e}$, $\Delta_{\{e\}}$ the point $u_e=1$ with the Dirac measure, is the integral over $\Delta_{\cE-e}$ of $v^{\mvec|_{\cE-e}-\one}$ times the family word at $(1,v)$, which is $W'(v)$: it is $P_{F'}(\mvec|_{\cE-e};\mathbf s,\mathbf t)$, and the compound formula at $S=\{e\}$ is the display.
\end{proof}

\begin{htheorem}[\hb{thm:poolingmove} (Pooling is a move)]
As stated: for tables with positive rational weights, and in (ii) whenever the active letters have one hinge value at rational factorizations.
\end{htheorem}

\begin{proof}
(i) is the reference, proof of thm:cleaningmove(i): the map $(\rho,u,\alpha_R)\mapsto(\rho u,\alpha_R)$ and its inverse are defined over $\mathbb Q$ whatever the weights, and the faces correspond as stated. After the change of variables the form is $\rho^{\kappa_B-1}\phi(u)$ times the form of $F_{B\to g}(u)$ at the gauge $(\rho,\alpha_R)$ (proof of \hb{thm:pooling}), with $\phi$ as in the proof of \hb{thm:cornerfactor}; its coefficients lie in the field of the weights and of the factors, which is $\mathbb Q$. The walls are matched by \hb{thm:wallsglue}(i), (ii) and the gluing of the counts, for families.

(ii) The common hinge value $\gamma$ is the hinge value of an active letter, a quotient of polynomials with rational coefficients, positive on the open simplex. The change $(\rho,u,\alpha_R)\mapsto(\rho\gamma(u),u,\alpha_R)$ is birational over $\mathbb Q$, with the inverse $(x,u,\alpha_R)\mapsto(x/\gamma(u),u,\alpha_R)$, and bijective on the open domain; by the proof of \hb{thm:cornerfactor} it carries the form of (i) to the product of the form $\phi(u)\gamma(u)^{-\kappa_B}d\omega_B(u)$ of the shape and the form of $G_R$ at the masses $(\kappa_B,\mvec|_R)$. In the datum of $G_B$, read on the section $|\alpha_B|=1$ (Lemma~\ref{lem:ch18-section}), the change $y=\gamma(u)z$, birational over $\mathbb Q$ and bijective on the open domain, carries the form to $\phi(u)\gamma(u)^{-\kappa_B}d\omega_B(u)$ times $z^{2\sigma-\kappa_B-1}(1+z)^{-2\sigma}dz$, the form of the list on two marks whose bases are the two singletons, at the exponents $(0,2\sigma)$ and the masses $(\kappa_B,2\sigma-\kappa_B)$, read in its chart at the first mark, where its letter is $1+z$. So the data of $F$ and of that list, and the data of $G_B$ and $G_R$, are related by these moves to the same product of three data, with rational masses $\kappa_B$ and $2\sigma-\kappa_B$. If $P_F$ converges, every proper wall of $F$ is positive (\hb{thm:projective}) and \hb{thm:cornerfactor} gives the convergence of the periods of $G_B$ and $G_R$, the Beta corner, and the shape integral, which is $P_F$ over $P_{G_R}$; and $P_F$ is the first integral of the chain.
\end{proof}

\begin{sremark}
Positive rational weights that factor across $B$ and $R$ factor with rational factors: fix a basis $X_\ell\sqcup Y_\ell$ at each level $\ell\in\{d(B),r(B)\}$ and put $\theta'_X=\theta_{X\sqcup Y_\ell}$ and $\theta''_Y=\theta_{X_\ell\sqcup Y}/\theta_{X_\ell\sqcup Y_\ell}$ for the patterns $X$ and the room sides $Y$ of the level $\ell$. For any factorization $\theta=\vartheta'\vartheta''$, $\theta_{X\sqcup Y_\ell}\theta_{X_\ell\sqcup Y}/\theta_{X_\ell\sqcup Y_\ell}=\vartheta'_X\vartheta''_{Y_\ell}\vartheta'_{X_\ell}\vartheta''_Y/(\vartheta'_{X_\ell}\vartheta''_{Y_\ell})=\theta_{X\sqcup Y}$. The hypothesis of (ii) asks more, rational factors at which the active letters share one hinge value; for one table with both letters active it asks a rational factorization with $\varrho_B=1$.
\end{sremark}

\begin{hproposition}[\hb{prop:cleanmaps} (Cleaning commutes with maps)]
As stated, for maps between any tables.
\end{hproposition}

\begin{proof}
A bijection $g$ that carries $(\cB,\theta)$ to $(\cB_1,\theta_1)$ maps bases to bases with their weights. Parallel and series bunches, the counts and the widths are defined by the bases, so they are carried to those of the image; a factorization $\theta'$, $\theta''$ across a class $B$ is carried to $\theta'_{1,g(p)}=\theta'_p$, $\theta''_{1,g(Y)}=\theta''_Y$, a factorization across $g(B)$, and the minors of $g(B)$ are the images of those of $B$ with their weights, so $\gamma^K_{g(B)}(\alpha\circ g^{-1})=\gamma^K_B(\alpha)$ and likewise for $\Psi$. An exchange maps $T$ to $\cE_1\setminus g(T)$, with the weight $\theta_T$. A parallel bunch $B$ has $|T\cap B|\le1$ for every basis and no loop, so $g(B)$ has $|g(B)\cap T_1|\ge|B|-1$ for every basis $T_1$ of the image and no coloop: a series bunch, and conversely. The patterns of $g(B)$ are the sets $g(B\setminus p)$, $p$ a pattern of $B$, and putting $\theta'_{1,g(B\setminus p)}=\theta'_p$ and $\theta''_{1,g(R\setminus Y)}=\theta''_Y$ factors the weights of the image. The level of $g(B\setminus p)$ is $|B|-|p|$, so the restriction of the image to $g(B)$ has the patterns $g(B\setminus p)$, $p\in\cB/R$, and its contraction of the room the patterns $g(B\setminus p)$, $p\in\cB|B$; hence $K_{\cB_1|g(B)}(\alpha\circ g^{-1})=\sum_{p\in\cB/R}\theta'_p\alpha^{B\setminus p}=\Psi_{\cB/R}(\alpha)$ and $K_{\cB_1/g(R)}(\alpha\circ g^{-1})=\Psi_{\cB|B}(\alpha)$, so the hinge value by $K$ of $g(B)$ is $\gamma^\Psi_B$; likewise its hinge value by $\Psi$ is $\gamma^K_B$. The statement on the cleanings of lists is the reference, proof of prop:cleanmaps.
\end{proof}

\section{Chapter~\hb{ch:words}: words of a tree}\label{sec:ch19}

\begin{hdefinition}[\hb{def:sideshape} (Side shape; half-sum and difference; pattern coefficients)]
As stated, for every table of a family.
\end{hdefinition}

\begin{proof}
The one claim is the identity $L_{j,w}(z)=\sum_pe^{\ip{z_B}{\chi_p}}c_{j,p}(w;z_R)$. Theorem~\hb{thm:taylorcube} for table $j$ gives $\Lambda_j(z)=\sum_{p\subseteq B}e^{\ip{z_B}{\chi_p}}\Lambda_{j,(p)}(z_R)$, with $\Lambda_{j,(p)}=0$ off the patterns of table $j$. Replacing $z$ by $-z$ gives
\[
\Lambda_j(-z)=\sum_pe^{-\ip{z_B}{\chi_p}}\Lambda_{j,(p)}(-z_R).
\]
Since $\chi_{B\setminus p}=-\chi_p$ on $B$, the substitution $p\mapsto B\setminus p$ turns this into
\[
\Lambda_j(-z)=\sum_pe^{\ip{z_B}{\chi_p}}\Lambda_{j,(B\setminus p)}(-z_R)=\sum_pe^{\ip{z_B}{\chi_p}}\Lambda^\ast_{j,(p)}(z_R).
\]
So
\[
w\Lambda_j(z)+(1-w)\Lambda_j(-z)=\sum_pe^{\ip{z_B}{\chi_p}}\big(w\Lambda_{j,(p)}(z_R)+(1-w)\Lambda^\ast_{j,(p)}(z_R)\big).
\]
\end{proof}

\begin{hproposition}[\hb{prop:sideexchange} (The side shape, and what an exchange does to it)]
As stated, for every table of a family and every bijection that carries or exchanges the family.
\end{hproposition}

\begin{proof}
For $T\subseteq\cE$ and $z\in\mathbb R^{\cE}$, with $(g^\ast z)_e=z_{g(e)}$,
\[
\ip z{\chi_{g(T)}}=\sum_{e\in\cE}z_{g(e)}\,\chi_{g(T)}(g(e))=\sum_{e\in\cE}(g^\ast z)_e\,\chi_T(e)=\ip{g^\ast z}{\chi_T}.
\]
If $g$ carries $F$, then $g(\cB_j)=\cB_{h(j)}$ with $\theta_{h(j),g(T)}=\theta_{j,T}$ (Definition~\hb{def:carry}), so
\[
\Lambda_{h(j)}(z)=\sum_{T\in\cB_j}\theta_{j,T}\,e^{\ip z{\chi_{g(T)}}}=\Lambda_j(g^\ast z),
\]
and $\Lambda_{h(j)}(-z)=\Lambda_j(-g^\ast z)$, as $g^\ast$ is linear; hence $L_{h(j),w}(z)=L_{j,w}(g^\ast z)$. If $g$ exchanges $F$, then $T\mapsto\cE\setminus g(T)$ maps $\cB_j$ onto $\cB_{h(j)}$ with $\theta_{h(j),\cE\setminus g(T)}=\theta_{j,T}$, and $\chi_{\cE\setminus g(T)}=-\chi_{g(T)}$, so $\Lambda_{h(j)}(z)=\Lambda_j(-g^\ast z)$ and $\Lambda_{h(j)}(-z)=\Lambda_j(g^\ast z)$; hence
\[
L_{h(j),w}(z)=w\,\Lambda_j(-g^\ast z)+(1-w)\,\Lambda_j(g^\ast z)=L_{j,1-w}(g^\ast z).
\]
The display holds for any two positive numbers $a,b$ in place of $\Lambda_j(z),\Lambda_j(-z)$ and any $s_j,t_j>0$: it is Theorem~\hb{thm:tower} at the unit configuration on a two-element set with the coarsest partition, whose Zamma is $(wa+(1-w)b)^{-(s_j+t_j)}=\sZ_{s_j+t_j;(1)}(wa+(1-w)b)$ at the shape $(w,1-w)$, that is, the two-factor form $a^{-s_j}b^{-t_j}=\Gst12{s_j+t_j}{s_j,\ t_j}\int_0^1w^{s_j-1}(1-w)^{t_j-1}(wa+(1-w)b)^{-s_j-t_j}dw$ (the reference, proof of prop:sideexchange). Taking $a=\Lambda_j(z)$ and $b=\Lambda_j(-z)$, $wa+(1-w)b=L_{j,w}(z)$.

For one table an exchange has $h$ the identity, so that $L_w(z)=L_{1-w}(g^\ast z)$ for every $w$. Hence
\[
\sZ_{s+t;(1)}\big(L_w(g^\ast z)\big)=\sZ_{s+t;(1)}\big(L_{1-w}(z)\big):
\]
the exchange together with $w\mapsto1-w$ keeps the Zamma. The substitution $w\mapsto1-w$ takes $w^{s-1}(1-w)^{t-1}$ to $w^{t-1}(1-w)^{s-1}$, which is the same function on $(0,1)$ exactly when $s=t$.
\end{proof}

\begin{htheorem}[\hb{thm:junctionpattern} (The junction in pattern form)]
As stated, for every family.
\end{htheorem}

\begin{proof}
The proof is that of the reference (proof of thm:junctionpattern), with a product over the tables in each step.

\emph{Signed coordinates.} By Theorem~\hb{thm:projective} the period is the integral of $\alpha^{\mvec-\one}\prod_jK_j^{-s_j}\Psi_j^{-t_j}\,d\alpha$ over the section $\alpha_g=1$. Put $\alpha=e^{2z}$, so that $d\alpha_e=2\alpha_e\,dz_e$ for each of the $n-1$ marks $e\neq g$. Theorem~\hb{thm:tworeadings}(i) for each table gives $K_j(\alpha)=e^{\ip\one z}\Lambda_j(z)$ and $\Psi_j(\alpha)=e^{\ip\one z}\Lambda_j(-z)$, so, with $\sum_j(s_j+t_j)=2\sigma$ and $\bmu=\mvec-\sigma\one$ (Definition~\hb{def:family}),
\begin{align*}
\alpha^{\mvec}\prod_jK_j^{-s_j}\Psi_j^{-t_j}&=e^{2\ip{\mvec}{z}-2\sigma\ip\one z}\prod_j\Lambda_j(z)^{-s_j}\Lambda_j(-z)^{-t_j}\\
&=e^{2\ip{\bmu}{z}}\prod_j\Lambda_j(z)^{-s_j}\Lambda_j(-z)^{-t_j},
\end{align*}
and, with $e^{2\ip{\bmu}{z}}=e^{2\ip{\bmu|_B}{z_B}}e^{2\ip{\bmu|_R}{z_R}}$,
\[
P_F(\mvec;\mathbf s,\mathbf t)=2^{n-1}\int_{\mathbb R^{\cE-g}}e^{2\ip{\bmu}{z}}\prod_j\Lambda_j(z)^{-s_j}\Lambda_j(-z)^{-t_j}\,dz_{\cE-g}\Big|_{z_g=0}.
\]

\emph{Side shapes.} For $j$ with $s_j,t_j>0$, Proposition~\hb{prop:sideexchange} writes the factor of table $j$ as
\[
\Lambda_j(z)^{-s_j}\Lambda_j(-z)^{-t_j}=\Gst12{s_j+t_j}{s_j,\ t_j}\int_0^1w_j^{s_j-1}(1-w_j)^{t_j-1}L_{j,w_j}(z)^{-(s_j+t_j)}\,dw_j .
\]
For $t_j=0<s_j$ the factor is $\Lambda_j(z)^{-s_j}=L_{j,1}(z)^{-s_j}$, for $s_j=0<t_j$ it is $\Lambda_j(-z)^{-t_j}=L_{j,0}(z)^{-t_j}$, and an inactive table contributes $1$. So every active table contributes $L_{j,w_j}(z)^{-(s_j+t_j)}$, integrated over $w_j$ when $s_j,t_j>0$ and with the value of $w_j$ fixed in the statement otherwise.

\emph{Scales.} For each active $j$, $L_{j,w_j}(z)>0$ and $L^{-a}=\Gamma(a)^{-1}\int_0^\infty y^{a-1}e^{-yL}dy$ with $a=s_j+t_j>0$. The constants combine to $\Gst12{s_j+t_j}{s_j,\ t_j}\Gamma(s_j+t_j)^{-1}=\Gamma(s_j)^{-1}\Gamma(t_j)^{-1}$ when $s_j,t_j>0$, and to $\Gamma(s_j)^{-1}$ or $\Gamma(t_j)^{-1}$ for a table at a full lean: the prefactor of the statement.

\emph{Patterns.} By Definition~\hb{def:sideshape},
\[
\sum_{j\text{ active}}y_jL_{j,w_j}(z)=\sum_{p\subseteq B}e^{\ip{z_B}{\chi_p}}C_p,\qquad C_p=\sum_{j\text{ active}}y_j\,c_{j,p}(w_j;z_R)\ge0,
\]
each $c_{j,p}(w;z_R)$ being a combination with the nonnegative coefficients $w$ and $1-w$ of the nonnegative functions $\Lambda_{j,(p)}(z_R)$ and $\Lambda^\ast_{j,(p)}(z_R)$.

\emph{The Gramma.} With $x=e^{z_B}$, Definition~\hb{def:gramma} in its logarithmic form gives
\[
\int_{\mathbb R^B}e^{2\ip{\bmu|_B}{z_B}}\exp\Big(-\sum_pC_p\,e^{\ip{\chi_p}{z_B}}\Big)dz_B=\sG_{2\bmu|_B;(\chi_p)_p}\big((C_p)_p\big),
\]
the configuration having one column $\chi_p$ for each $p$ with $C_p\neq0$.

\emph{Order of integration.} Every integrand is nonnegative, so Tonelli's theorem allows the integrals to be taken in the order of the statement: the side shapes outermost, then $z_{R-g}$ at $z_g=0$, then the scales $y_j$, then $z_B$. The two sides agree in $[0,\infty]$ and are finite, every proper wall being positive (Theorem~\hb{thm:projective}).
\end{proof}

\begin{sremark}
With one table the formula is that of the reference, and at a full lean of that table the integral over the side shape is absent, which is the reference's last sentence.
\end{sremark}

\begin{hproposition}[\hb{prop:mapcoeff} (What a map keeping the bunch gives the coefficients)]
As stated, for every family.
\end{hproposition}

\begin{proof}
Since $\tilde g$ keeps $B$, it maps $R$ onto $R$, so $(\tilde g^\ast z)_R=\tilde g_R^\ast z_R$, and for $p\subseteq B$, substituting $e'=\tilde g(e)$,
\[
\ip{(\tilde g^\ast z)_B}{\chi_p}=\sum_{e\in B}z_{\tilde g(e)}\chi_p(e)=\sum_{e'\in B}z_{e'}\chi_{\tilde g(p)}(e')=\ip{z_B}{\chi_{\tilde g(p)}}.
\]
If $\tilde g$ carries $F$, Proposition~\hb{prop:sideexchange} and Definition~\hb{def:sideshape} give
\begin{align*}
\sum_{q\subseteq B}e^{\ip{z_B}{\chi_q}}c_{h(j),q}(w;z_R)&=L_{h(j),w}(z)=L_{j,w}(\tilde g^\ast z)\\
&=\sum_{p\subseteq B}e^{\ip{z_B}{\chi_{\tilde g(p)}}}c_{j,p}(w;\tilde g_R^\ast z_R).
\end{align*}
The vectors $\chi_q$, $q\subseteq B$, are distinct, and at fixed $z_R$ the exponentials $e^{\ip{z_B}{\chi_q}}$ with distinct exponents are linearly independent functions of $z_B$: along a generic direction the largest exponent dominates (the reference, proof of prop:mapcoeff). Comparing the coefficients at $q=\tilde g(p)$ gives $c_{h(j),\tilde g(p)}(w;z_R)=c_{j,p}(w;\tilde g_R^\ast z_R)$. If $\tilde g$ exchanges $F$, the same comparison with $L_{h(j),w}(z)=L_{j,1-w}(\tilde g^\ast z)$ gives $c_{h(j),\tilde g(p)}(w;z_R)=c_{j,p}(1-w;\tilde g_R^\ast z_R)$ for every $w$, and $w\mapsto1-w$ gives the statement.
\end{proof}

\begin{hproposition}[\hb{prop:levelformQ} (Level form of the half-sum)]
As stated, for every table of the family of Theorem~\hb{thm:junctionpattern}.
\end{hproposition}

\begin{proof}
Every claim concerns one table $j$, and the reference's proof (proof of prop:levelformQ) applies to it. The two displayed coefficients are $c_{j,p}(w;z_R)$ at $w=\frac12$ and its derivative in $w$ (Definition~\hb{def:sideshape}).

(i) $c_{j,B\setminus p}(\frac12;-z_R)=\frac12\big(\Lambda_{j,(B\setminus p)}(-z_R)+\Lambda_{j,(p)}(z_R)\big)=c_{j,p}(\frac12;z_R)$, and $\partial_wc_{j,B\setminus p}(w;-z_R)=\Lambda_{j,(B\setminus p)}(-z_R)-\Lambda_{j,(p)}(z_R)=-\partial_wc_{j,p}(w;z_R)$.

(ii) Let $\tau$ be a transposition of two marks of $B$ fixing the rest. As in the proof of Proposition~\hb{prop:mapcoeff}, $L_{j,1/2}(\tau^\ast z)=\sum_pe^{\ip{z_B}{\chi_{\tau(p)}}}c_{j,p}(\frac12;z_R)$. If $c_{j,p}(\frac12;\cdot)$ depends on $|p|$ alone, this is $L_{j,1/2}(z)$, since $|\tau(p)|=|p|$. Conversely, if $L_{j,1/2}(\tau^\ast z)=L_{j,1/2}(z)$ for every such $\tau$, the independence of the exponentials gives $c_{j,\tau(p)}(\frac12;\cdot)=c_{j,p}(\frac12;\cdot)$, and these transpositions act transitively on the subsets of $B$ of each size. So level form is the invariance of $L_{j,1/2}$ under these transpositions. Further,
\begin{gather*}
L_{j,1/2}(z)=\tfrac12\sum_{T\in\cB_j}\theta_{j,T}\big(e^{\ip z{\chi_T}}+e^{\ip z{\chi_{\cE\setminus T}}}\big),\\
L_{j,1/2}(\tau^\ast z)=\tfrac12\sum_{T\in\cB_j}\theta_{j,T}\big(e^{\ip z{\chi_{\tau(T)}}}+e^{\ip z{\chi_{\cE\setminus\tau(T)}}}\big),
\end{gather*}
by the computation of the proof of Proposition~\hb{prop:sideexchange}. The exponentials $e^{\ip z{\chi_S}}$, $S\subseteq\cE$, are linearly independent, so the invariance under $\tau$ says that $\tau$ maps the multiset $\cB_j\uplus\{\cE\setminus T:\ T\in\cB_j\}$, with its weights, onto itself. When $2r_j\neq n$ a basis and a complement of a basis have different sizes, $\tau$ keeps sizes, and the condition says that $\tau$ maps $\cB_j$ onto itself with its weights, that is, carries table $j$ (Definition~\hb{def:carry}), as in Proposition~\hb{prop:levelform}.

(iii) If $\tau$ carries table $j$, it keeps $\Lambda_j(z)$ and $\Lambda_j(-z)$ under $z\mapsto\tau^\ast z$; if it exchanges table $j$, it swaps them (Proposition~\hb{prop:sideexchange}, the bijection of the tables being the identity). Either way $L_{j,1/2}(\tau^\ast z)=L_{j,1/2}(z)$, so the half-sum has level form by (ii), and $\partial_wL_{j,w}(\tau^\ast z)=\pm\partial_wL_{j,w}(z)$, so the square of the difference is unchanged by $\tau$. The transpositions of $B$ generate the permutations of $B$.
\end{proof}

\begin{htheorem}[\hb{thm:buildup} (Building back up)]
As stated; the last sentence of (iv) is stated for every molecule.
\end{htheorem}

\begin{proof}
Parts (i) to (iii), and (iv) up to its last sentence, are proved in the reference (proof of thm:buildup). The last sentence of (iv) makes three claims.

\emph{Dimension.} In (iv) every table has $s_j,t_j>0$ and is active. The unknowns are the $c_{j,T}$, $T\in\cB_j$, $\sum_j|\cB_j|$ of them, subject to the $N$ share sums $\sum_Tc_{j,T}=t_j-s_j$ and the $n-1$ pins $-\frac12\sum_{j,T}c_{j,T}\chi_T(e)=\mu_e$, $e\in\cE-g$. Their linear parts are independent. Indeed, let $\sum_j\lambda_j\sum_Tc_{j,T}+\sum_{e\neq g}d_e\sum_{j,T}c_{j,T}\chi_T(e)$ vanish identically, and put $d_g=0$. The coefficient of $c_{j,T}$ is $\lambda_j+\ip d{\chi_T}=\lambda_j+2\ip d{\one_T}-\ip d\one$, so $\ip d{\one_T}$ is constant over the bases of every table. By Lemma~\hb{lem:circuit}, $d$ is constant on every molecule; the family is a molecule, so $d$ is constant, and $d_g=0$ makes it $0$; then every $\lambda_j=0$. So the fibre of differences has dimension $\sum_j|\cB_j|-N-(n-1)=\sum_j(|\cB_j|-1)-(n-1)$.

\emph{The pins at the centre.} At the centre $t_j=s_j$ and $\bmu=0$, so the share sums read $\sum_Tc_{j,T}=0$. With $\chi_T=2\one_T-\one$,
\[
\sum_Tc_{j,T}\chi_T(e)=2\big(M_j^{\mathsf T}\mathbf c_j\big)_e-\sum_Tc_{j,T},
\]
so on the share sums the pin of $e\neq g$ reads $(\sum_jM_j^{\mathsf T}\mathbf c_j)_e=0$. The equation at $g$ follows, since $\sum_{e\in\cE}(M_j^{\mathsf T}\mathbf c_j)_e=\sum_Tc_{j,T}|T|=r_j\sum_Tc_{j,T}=0$. So the pins and the share sums together are equivalent to $\sum_Tc_{j,T}=0$ for every $j$ and $\sum_jM_j^{\mathsf T}\mathbf c_j=0$.

\emph{The contour.} At the centre the gauge $\one$ is native (Proposition~\hb{prop:centrefamily}(i)). The leaves are opened by the atom formula with one share per basis and letter (Corollary~\hb{cor:atomwall}), and the native split there (Definition~\hb{def:split}) gives the basis $T$ of table $j$ the shares $a_{j,T}=s_j\theta_{j,T}/K_j(\one)$ and $b_{j,T}=t_j\theta_{j,T}/\Psi_j(\one)$, the probabilities of $T$ under the basis law and the cobasis law at $\one$ times the exponents. As $\Psi_j(\one)=\sum_T\theta_{j,T}=K_j(\one)$ and $s_j=t_j$, $c_{j,T}=b_{j,T}-a_{j,T}=0$ there, so the wall contour of the native split (Theorem~\hb{thm:wallsplit}(iii)) is $\re\mathbf c=0$.
\end{proof}

\begin{sremark}
With one table $\sum_Tc_T=\frac1r\sum_e(M^{\mathsf T}\mathbf c)_e$ follows from $M^{\mathsf T}\mathbf c=0$, and the dimension is $|\cB_1|-n$: the reference's case of one table at the centre.
\end{sremark}

\begin{hproposition}[\hb{prop:jointdraw} (The draw series)]
As stated; the reading of its last part is stated for every family.
\end{hproposition}

\begin{proof}
The first part, through the Trailmoment and the Dirichlet moments, is the reference's proposition (proof of prop:jointdraw). For the last part take one block, $q=2N$, $P_{2j-1}=K_j$ with the index $s_j$ and $P_{2j}=\Psi_j$ with the index $t_j$; an index $0$ contributes the factor $1$.

\emph{Constants.} $K_j$ is homogeneous of degree $r_j$, with the squarefree monomials $u^T$, $T\in\cB_j$, and the coefficients $\theta_{j,T}$. For $\mathbf a=\one_T$, $b_{\mathbf a}=r_j!/\prod_ea_e!=r_j!$, and every other $\mathbf a$ has $[u^{\mathbf a}]K_j=0$; so $C_{2j-1}=\theta_{j,\max}/r_j!$. Likewise $\Psi_j$ has degree $n-r_j$ and the monomials $u^{\cE\setminus T}$ with the coefficients $\theta_{j,T}$, so $C_{2j}=\theta_{j,\max}/(n-r_j)!$.

\emph{The period.} $\oE_{U\sim\Dir_{\mvec}}[f(U)]=\Gst{n}{1}{\mvec}{m_{\cE}}^{-1}\int_{\Delta_{\cE}}u^{\mvec-\one}f\,d\omega_{\cE}$ for $f\ge0$ (Theorem~\hb{thm:dirichlet}), so the left side times $\Gst{n}{1}{\mvec}{m_{\cE}}$ is the stated integral, which on the balance wall is $P_F(\mvec;\mathbf s,\mathbf t)$ (Theorem~\hb{thm:projective}).

\emph{The urn.} For $u$ in the simplex, $(\sum_eu_e)^{g}=\sum_{|\mathbf a|=g}\frac{g!}{\mathbf a!}u^{\mathbf a}=1$ is the total chance of $g$ independent draws from $u$, grouped by their counts $\mathbf a$. For $P_{2j-1}=K_j$,
\[
\frac{P_{2j-1}}{C_{2j-1}}=\sum_{T\in\cB_j}\frac{\theta_{j,T}}{\theta_{j,\max}}\,r_j!\,u^T
\]
is the chance that a round of $r_j$ draws is a basis $T$ of table $j$ and does not fail, a round that is the basis $T$ failing with probability $1-\theta_{j,T}/\theta_{j,\max}$; so $Q_{2j-1}=1-P_{2j-1}/C_{2j-1}$ is the chance that the round fails, and likewise $Q_{2j}$ for a round of $n-r_j$ draws and the cobases. Given $u$, independent rounds all fail with the chance $\prod_iQ_i(u)^{k_i}$. The urn that draws a mark with probability proportional to its current mass and then adds one to that mass gives a sequence of draws with the counts $\mathbf b$ the chance $\prod_e(m_e)_{b_e}/(m_{\cE})_{|\mathbf b|}=\oE_{U\sim\Dir_{\mvec}}[U^{\mathbf b}]$ (Theorem~\hb{thm:dirmoments}), the chance of the same sequence when $U\sim\Dir_{\mvec}$ is drawn first and the marks are then drawn independently from $U$. Summing over the sequences whose rounds all fail, $\oE_{U\sim\Dir_{\mvec}}[\prod_iQ_i(U)^{k_i}]$ is the urn's chance of the event of the statement; the order of the rounds does not change it, the product being symmetric in its factors. For lists every ratio $\theta_{j,T}/\theta_{j,\max}$ is $1$, the $Q_i$ have integer coefficients, and the expectation is a finite sum of Dirichlet moments.

\emph{Theorem~\hb{thm:peak}(i).} With the references $M_i=C_i$, on the simplex
\[
M_i^{-w_i-n_i}(M_i-L_i)^{n_i}=C_i^{-w_i}(1-L_i/C_i)^{n_i}=C_i^{-w_i}Q_i^{n_i},
\]
so the series of Theorem~\hb{thm:peak}(i), divided by $\Gst{n}{1}{\mvec}{m_{\cE}}$, is the present series term by term.
\end{proof}

\begin{hcorollary}[\hb{cor:cremonatrail} (At $\delta=0$ and at a full lean the period reads the law of one letter for each table)]
As stated, for every family.
\end{hcorollary}

\begin{proof}
\emph{Letters.} For a table riding the line, $K_j^{-\sigma_j}\Psi_j^{-\sigma_j}=\mathbf E_j^{-\sigma_j}$; at $(s_j,0)$ the factor is $K_j^{-s_j}$ and at $(0,t_j)$ it is $\Psi_j^{-t_j}$. So on the open simplex
\[
\prod_jK_j^{-s_j}\Psi_j^{-t_j}=\prod_jL_j^{-\nu_j}=\sZ_{\boldsymbol\nu;I_N}\big((L_j)_j\big),
\]
the Zamma at the unit configuration being $\prod_j\xi_j^{-\nu_j}$ (Definition~\hb{def:zamma}).

\emph{The period.} By Theorem~\hb{thm:projective} and the Dirichlet density (Theorem~\hb{thm:dirichlet}),
\begin{align*}
P_F(\mvec;\mathbf s,\mathbf t)&=\int_{\Delta_{\cE}}u^{\mvec-\one}\prod_jK_j^{-s_j}\Psi_j^{-t_j}\,d\omega_{\cE}\\
&=\Gst{n}{1}{\mvec}{m_{\cE}}\,\oE_{U\sim\Dir_{\mvec}}\Big[\sZ_{\boldsymbol\nu;I_N}\big((L_j(U))_j\big)\Big].
\end{align*}
The expectation is the integral of the kernel against the image $\phi_{\mvec}$ of $\Dir_{\mvec}$ under $(L_j)_j$, the Trailmoment $(\sT_{\boldsymbol\nu;I_N}[\mathrm{id}]\,\phi_{\mvec})((0,\infty)^N)$ with the identity as coefficient function (Definition~\hb{def:jointtrail}), as in the display of Theorem~\hb{thm:normalizer}(i). With one table it is the Trailmoment $\Tx{\nu}{\phi_{\mvec}}{0}$ of a law on the half-line, at $\delta=0$ the image of $\Dir_{\mvec}$ under the even letter (Theorem~\hb{thm:transportfree}).

\emph{The law.} Each $L_j$ is a polynomial, positive on the open simplex and continuous on the closed one, and $\Dir_{\mvec}$ charges only the open simplex; so $\phi_{\mvec}$ is carried by $\prod_j(0,\max_{\Delta_{\cE}}L_j]$. A law carried by a compact set is determined by its moments, polynomials being dense in the continuous functions on that set. Every moment $\oE[\prod_jL_j(U)^{k_j}]$ is the expectation of a polynomial in $U$, a finite sum of Dirichlet moments (Theorem~\hb{thm:dirmoments}).
\end{proof}

\begin{htheorem}[\hb{thm:drawchain} (The draw series along a chain)]
As stated, for every family of tables.
\end{htheorem}

\begin{proof}
The proof is that of the reference (proof of thm:drawchain), with the acceptances and a product over the tables.

\emph{The failures split.} Fix table $j$ and a round of $r_j$ draws. A round that extends on $S_{i+1}$ extends on $S_i$, with the trace $P_i=P_{i+1}\cap S_i$, and every round extends on $S_0=\varnothing$ with the trace $\varnothing$. So the events $A_i$, that the round extends on $S_i$ and is accepted at $S_1,\dots,S_i$, decrease with $i$, $A_0$ is certain, and $\mathfrak f_{j,i}$ is the chance of $A_i\setminus A_{i+1}$; hence $\sum_{i=0}^d\mathfrak f_{j,i}=1-\Pr(A_{d+1})$.

\emph{The acceptances telescope.} A round extends on $S_{d+1}=\cE$ exactly when it is a basis $T$ of table $j$ with distinct draws, which has the chance $r_j!\,u^T$; its traces are then $P_i=T\cap S_i$, and its acceptances multiply to
\[
\prod_{i=0}^d\frac{\theta_j(S_{i+1};T\cap S_{i+1})}{\theta_j(S_i;T\cap S_i)}=\frac{\theta_j(\cE;T)}{\theta_j(\varnothing;\varnothing)}=\frac{\theta_{j,T}}{\theta_{j,\max}} .
\]
So $\Pr(A_{d+1})=\sum_Tr_j!\,u^T\theta_{j,T}/\theta_{j,\max}=r_j!\,K_j(u)/\theta_{j,\max}$. The cobasis rounds are the basis rounds of the dual table, whose bases $\cE\setminus T$ carry the weights $\theta_{j,T}$, whose largest weight is $\theta_{j,\max}$ and whose basis polynomial is $\Psi_j$ (Proposition~\hb{prop:identity}); so $\sum_i\mathfrak f^\ast_{j,i}=1-(n-r_j)!\,\Psi_j/\theta_{j,\max}$.

\emph{Polynomials.} Each $\mathfrak f_{j,i}(u)$ is a finite sum, over the sequences of $r_j$ draws, of the monomial of the sequence times a product of acceptance and rejection probabilities in $[0,1]$: a polynomial with nonnegative coefficients, with integer coefficients for lists, whose rounds are always accepted.

\emph{The series.} Proposition~\hb{prop:jointdraw}, with one block and the constants $C_{2j-1}=\theta_{j,\max}/r_j!$, $C_{2j}=\theta_{j,\max}/(n-r_j)!$ of its last part, gives
\begin{multline*}
\oE_{U\sim\Dir_{\mvec}}\Big[\prod_jK_j^{-s_j}\Psi_j^{-t_j}\Big]=\prod_jC_{2j-1}^{-s_j}C_{2j}^{-t_j}\sum_{\mathbf k\in\mathbb N^{2N}}\prod_j\frac{(s_j)_{k_{2j-1}}(t_j)_{k_{2j}}}{k_{2j-1}!\,k_{2j}!}\\
\times\oE_{U\sim\Dir_{\mvec}}\Big[\prod_j\Big(\sum_i\mathfrak f_{j,i}\Big)^{k_{2j-1}}\Big(\sum_i\mathfrak f^\ast_{j,i}\Big)^{k_{2j}}\Big],
\end{multline*}
both sides finite or infinite together. The multinomial theorem expands
\[
\frac{(s_j)_k}{k!}\Big(\sum_i\mathfrak f_{j,i}\Big)^k=\sum_{|\mathbf k_j|=k}\frac{(s_j)_{|\mathbf k_j|}}{\mathbf k_j!}\prod_i\mathfrak f_{j,i}^{k_{j,i}},
\]
with nonnegative terms, and likewise for the cobasis rounds; Tonelli's theorem exchanges the sums and the expectation, and multiplying by $\Gst{n}{1}{\mvec}{m_{\cE}}$ (Theorem~\hb{thm:dirichlet}) gives the display. Each expectation is that of a polynomial with nonnegative coefficients, a finite nonnegative combination of Dirichlet moments (Theorem~\hb{thm:dirmoments}).

\emph{Nesting.} $\sum_{i'\ge i}k_{j,i'}$ is a tail sum of nonnegative integers, so it decreases with $i$.

\emph{Bijections.} Let $g$ keep each $B_i$, hence each $S_i$, and the masses. A sequence of draws $D$ from $u$ has the chance of the sequence $g(D)$ from $g_\ast u$. If $g$ maps the bases of table $j$ onto those of table $h(j)$ with their weights, then $g(T\cap S_i)=g(T)\cap S_i$, so $D$ extends on $S_i$ for table $j$ exactly when $g(D)$ extends on $S_i$ for table $h(j)$, the traces corresponding under $g$, and $\theta_j(S_i;P)=\theta_{h(j)}(S_i;g(P))$, so the acceptances agree: $\mathfrak f_{h(j),i}(g_\ast u)=\mathfrak f_{j,i}(u)$, and likewise for the cobasis rounds, $g$ carrying the dual tables as well. If $T\mapsto\cE\setminus g(T)$ maps the bases of table $j$ onto those of table $h(j)$ with their weights, then $g(T)$ is a cobasis of table $h(j)$ for every basis $T$ of table $j$, the round lengths agree, $r_j=n-r_{h(j)}$, the largest weights over the traces agree in the same way, and $\mathfrak f^\ast_{h(j),i}(g_\ast u)=\mathfrak f_{j,i}(u)$, $\mathfrak f_{h(j),i}(g_\ast u)=\mathfrak f^\ast_{j,i}(u)$. Since the masses are kept, $g_\ast U\sim\Dir_{\mvec}$ when $U\sim\Dir_{\mvec}$, so in the second case
\[
\oE\Big[\prod_{j,i}\mathfrak f_{j,i}(U)^{k_{j,i}}\mathfrak f^\ast_{j,i}(U)^{l_{j,i}}\Big]=\oE\Big[\prod_{j,i}\mathfrak f^\ast_{h(j),i}(U)^{k_{j,i}}\mathfrak f_{h(j),i}(U)^{l_{j,i}}\Big],
\]
the expectation at the indices $(\mathbf k'_{j'},\mathbf l'_{j'})=(\mathbf l_{h^{-1}(j')},\mathbf k_{h^{-1}(j')})$. With one table $h$ is the identity, an exchange maps the bases onto complements of bases, so $r=n-r$, and the array is symmetric under $\mathbf k\leftrightarrow\mathbf l$.
\end{proof}

\begin{htheorem}[\hb{thm:drawtree} (The draw series along a tree)]
As stated, for every molecule whose children's weights factor across the childcules.
\end{htheorem}

\begin{proof}
The proof is that of the reference (proof of thm:drawtree), the weights entering through the acceptances; the coordinates of the tree and their independence (Theorem~\hb{thm:aggregation} at every node) are unchanged.

(i) Let $\cB'$ have rank $r'$ at the node $v$, and write $R_i=R_{v,i}$, $x=x^{(v)}$, $u=u^{(v)}$ and $v^{(i)}$ for the shape of $R_i$. By Proposition~\hb{prop:directsum} the bases of $\cB'$ are the sets $T=p\sqcup\bigsqcup_iX_i$, $p$ a pattern at $B_v$ and $X_i$ a basis of $\cB'_{(p)}|R_i$ of size $r'_{p,i}$, and by the hypothesis $\theta_T=\prod_i\theta_{(p),i}(X_i)$. The shape of $M_v$ is $x_{B_v}u$ on $B_v$ and $x_{R_i}v^{(i)}$ on $R_i$, so by homogeneity
\[
K_{\cB'}=\sum_px_{B_v}^{|p|}u^p\prod_ix_{R_i}^{r'_{p,i}}K_{\cB'_{(p)}|R_i}\big(v^{(i)}\big),
\]
the restrictions carrying the weights $\theta_{(p),i}$. Multiply by $r'!/\theta_{\max}(\cB')$, write $r'!=\frac{r'!}{\prod_ir'_{p,i}!}\prod_ir'_{p,i}!$, and multiply and divide the term of $p$ by
\[
\prod_i\theta_{\max}\big(\cB'_{(p)}|R_i\big);
\]
this gives $\mathfrak s_{\cB'}=\sum_p\mathfrak p_{\cB',p}\prod_i\mathfrak s_{\cB'_{(p)}|R_i}$. The hypothesis fixes the weights of the restrictions up to constant factors whose product is $1$, which change neither the $\mathfrak s$ of the restrictions nor the first factor of $\mathfrak p_{\cB',p}$. That factor is at most $1$, since
\[
\theta_{\max}(\cB')=\max_p\ \max_{(X_i)_i}\prod_i\theta_{(p),i}(X_i)=\max_p\prod_i\theta_{\max}\big(\cB'_{(p)}|R_i\big).
\]
The chance that $r'$ independent draws from the shape of $M_v$ are the basis $T$ with distinct draws is $r'!$ times its monomial, and the acceptance with probability $\theta_T/\theta_{\max}(\cB')$ gives the chance $r'!K_{\cB'}/\theta_{\max}(\cB')$ of an accepted basis. The round passes $v$ showing $p$, its draws in $B_v$ being the marks of $p$ and its numbers of draws in the $R_i$ the $r'_{p,i}$, with the chance $\frac{r'!}{\prod_ir'_{p,i}!}x_{B_v}^{|p|}u^p\prod_ix_{R_i}^{r'_{p,i}}$, and given this its draws in $R_i$ are $r'_{p,i}$ independent draws from $v^{(i)}$. The acceptance splits as
\[
\frac{\theta_T}{\theta_{\max}(\cB')}=\frac{\prod_i\theta_{\max}\big(\cB'_{(p)}|R_i\big)}{\theta_{\max}(\cB')}\ \prod_i\frac{\theta_{(p),i}(X_i)}{\theta_{\max}\big(\cB'_{(p)}|R_i\big)},
\]
an acceptance at $v$ followed by independent acceptances in the childcules. So $\mathfrak p_{\cB',p}$ is the chance of passing $v$ showing $p$ and being accepted there, and $\mathfrak s_{\cB'_{(p)}|R_i}$ the chance of an accepted basis of the restriction in $R_i$.

(ii) By (i) the round fails at $v$ with the chance $1-\sum_p\mathfrak p_{\cB',p}$, and otherwise passes $v$ showing one pattern $p$ and is accepted there; it then first fails at $w$ in the subtree of $R_i$ exactly when it succeeds in $R_1,\dots,R_{i-1}$ and first fails at $w$ inside $R_i$, which is the recursion. Write $\mathfrak p_{\cB',p}=a_p\,\mathfrak p^\circ_{\cB',p}$ with $a_p\le1$ the first factor. Then
\[
1-\sum_p\mathfrak p_{\cB',p}=\Big(1-\sum_p\mathfrak p^\circ_{\cB',p}\Big)+\sum_p(1-a_p)\,\mathfrak p^\circ_{\cB',p}.
\]
The bracket is $1=(x_{B_v}\sum_{e\in B_v}u_e+\sum_ix_{R_i})^{r'}$, expanded by the multinomial theorem, less some of its terms, distinct patterns giving distinct terms (the reference, proof of thm:drawtree, part (ii)); it has nonnegative coefficients in the coordinates of $v$, and so has the sum. With $1-\prod_iS_i=\sum_i\prod_{i'<i}S_{i'}(1-S_i)$ and induction on the subtree, every $\mathfrak f_{\cB'\to w}$ is a finite sum of products, over the blocks of the tree, of polynomials with nonnegative coefficients in one block each.

(iii) On the simplex $r_j!K_j/\theta_{j,\max}=\mathfrak s_{\cB_j}=1-\sum_w\mathfrak f_{j,w}$, with the sum in $[0,1)$ on the interior. The letter $\Psi_j$ is the basis polynomial of the dual table (Proposition~\hb{prop:identity}), whose bases $\cE\setminus T$ carry the weights $\theta_{j,T}$, so that its largest weight is $\theta_{j,\max}$; its children are the duals of the children with the same weights (Lemma~\hb{lem:descent}(ii) with the identity), and the dual of a direct sum is the direct sum of the duals of its parts with the product weights, so the factoring passes to the dual and (i), (ii) apply to its rounds. Then $K_j^{-s_j}=(r_j!/\theta_{j,\max})^{s_j}(1-\sum_w\mathfrak f_{j,w})^{-s_j}$, expanded as in Proposition~\hb{prop:jointdraw} and by the multinomial theorem as in Theorem~\hb{thm:drawchain}, with nonnegative terms; Tonelli's theorem, the independence of the blocks under $\Dir_{\mvec}$ (Theorem~\hb{thm:aggregation}) and Theorem~\hb{thm:dirmoments} give the series and the form of its expectations. On the balance wall the left side is the period (Theorem~\hb{thm:projective}).

(iv) A round that first fails in the subtree of $w$ passes every node above $w$, which gives the nesting. On a chain every node has one childcule. The table at the node $M_k$ that a round of table $j$ reaches is the child of table $j$ at its trace $P_k$ on $S_k$, whose largest weight is $\theta_j(S_k;P_k)$ in the notation of Theorem~\hb{thm:drawchain}; so the first factor of $\mathfrak p$ at that node is $\theta_j(S_{k+1};P_{k+1})/\theta_j(S_k;P_k)$, the acceptance at $S_{k+1}$, and passing the node is extending on $S_{k+1}$. The chances $\mathfrak f$ of the two theorems agree, and the series is that of Theorem~\hb{thm:drawchain}.
\end{proof}

\begin{hproposition}[\hb{prop:centrerates} (Rates at the centre)]
As stated, for every family whose active tables are connected.
\end{hproposition}

\begin{proof}
(i) is Proposition~\hb{prop:centrefamily}(ii), with $c_j(\cE\setminus S)=c_j(S)$ (Theorem~\hb{thm:window}).

(ii) The rate is $\varpi=\min_i\varpi_i$ over the letters, $\varpi_i$ the least of $\kappa_S/d_i(S)$ over the nonempty proper $S$ with $d_i(S)>0$ (Theorem~\hb{thm:peak}(ii)); an inactive table has $\sigma_j=0$ and no letters. A connected table has $c_j(S)\ge1$ on every nonempty proper $S$ (Definition~\hb{def:family}), and so no loops and no coloops, a loop or a coloop $e$ having $c_j(\{e\})=r_j(\{e\})-d_j(\{e\})=0$ (Theorem~\hb{thm:window}) with $\{e\}$ proper as $n\ge3$. So $\kappa_S=\sum_j\sigma_jc_j(S)\ge\sum_j\sigma_j=\sigma$.

For the letter $K_j$, $d_j(S)=|S|-r^\ast_j(S)=r_j-r_j(\cE\setminus S)\le r_j-1$ (Definition~\hb{def:skeleton}), as $\cE\setminus S$ is nonempty and contains no loop. For $\Psi_j$, $d^\ast_j(S)=|S|-r_j(S)$ is the nullity of $S$, which increases with the set, and $S\subseteq\cE-e$ for some $e$, of nullity $n-1-r_j(\cE-e)=n-1-r_j$, as $e$ is not a coloop. So $\kappa_S/d_j(S)\ge\sigma/(r_j-1)$ wherever $d_j(S)>0$, and $\kappa_S/d^\ast_j(S)\ge\sigma/(n-r_j-1)$ wherever $d^\ast_j(S)>0$.

At $S=\cE-e$, every table without loops or coloops has $c_j(\cE-e)=c_j(\{e\})=r_j(\{e\})-d_j(\{e\})=1$ (Theorem~\hb{thm:window}), so $\kappa_{\cE-e}=\sigma$; and $d_j(\cE-e)=r_j-r_j(\{e\})=r_j-1$, $d^\ast_j(\cE-e)=n-1-r_j$. Hence $\varpi_{K_j}=\sigma/(r_j-1)$ when $r_j\ge2$ and $\varpi_{\Psi_j}=\sigma/(n-r_j-1)$ when $n-r_j\ge2$, attained at the faces $\cE-e$; a letter with no $S$ of positive count does not enter the minimum. As $(r_j-1)+(n-r_j-1)=n-2\ge1$, at least one letter of every active table enters, and $\varpi=\sigma/\max_j\max(r_j-1,n-r_j-1)$.

(iii) By Definition~\hb{def:delcon}, $K_j=K_{j,\setminus e}+u_eK_{j,/e}$ and $\Psi_j=u_e\Psi_{j,\setminus e}+\Psi_{j,/e}$, the deletion and the contraction being tables since $e$ is neither a loop nor a coloop of the active table $j$. They are homogeneous in $u_{\cE-e}$ of the degrees $r_j$, $r_j-1$, $n-1-r_j$ and $n-r_j$; substituting $u_e=1-x$ and $u_{\cE-e}=xv$ gives the display, and as $x\to0$ the terms of $x^{r_j-1}K_{j,/e}(v)$ and $x^{n-r_j-1}\Psi_{j,\setminus e}(v)$ dominate.

(iv) Let $\beta$ carry the family, with $v'_{\beta(e')}=v_{e'}$ on $\cE-e$. It maps the bases of table $j$ onto those of table $h(j)$ with their weights, the bases through $e$ onto those through $f$, so
\[
K_{h(j),/f}(v')=\sum_{T\ni e}\theta_{j,T}\,v'^{\,\beta(T)-f}=\sum_{T\ni e}\theta_{j,T}\,v^{T-e}=K_{j,/e}(v),
\]
and in the same way $\Psi_{h(j),\setminus f}(v')=\Psi_{j,\setminus e}(v)$, the bases avoiding $e$ going to those avoiding $f$. Let $\beta$ exchange the family: $T\mapsto\cE\setminus\beta(T)$ maps the bases of table $j$ through $e$ onto the bases of table $h(j)$ avoiding $f$, with their weights, so
\[
\Psi_{h(j),\setminus f}(v')=\sum_{T\ni e}\theta_{j,T}\,v'^{\,(\cE-f)\setminus(\cE\setminus\beta(T))}=\sum_{T\ni e}\theta_{j,T}\,v'^{\,\beta(T)-f}=K_{j,/e}(v),
\]
and in the same way $K_{h(j),/f}(v')=\Psi_{j,\setminus e}(v)$. Either way the vertex product $K_{j,/e}\Psi_{j,\setminus e}$ goes to $K_{h(j),/f}\Psi_{h(j),\setminus f}$. With one table, a bijection taking $e$ to $f$ that carries or exchanges it exists exactly when $e$ and $f$ lie in one class at the centre, by Corollary~\hb{cor:certifyone} at the table's gauge, the gauge of the centre.
\end{proof}

\begin{sremark}
Connectedness of every active table is needed for (ii). On the marks $\{0,1,2,3\}$, for the family of the list whose bases are all the two-element sets and the disconnected list with the bases $\{a,b\}$, $a\in\{0,2\}$, $b\in\{1,3\}$, at the strengths $(1,3)$, the face $S=\{0,2\}$ has $\kappa_S=1\cdot2+3\cdot0=2$ and $d_2(S)=1$, so $\varpi\le2$, while $\sigma/\max_j\max(r_j-1,n-r_j-1)=4$.
\end{sremark}

\begin{hcorollary}[\hb{cor:mergedvalues} (Merged Skeleton values)]
As stated, at every pair of exponents and every mass on the balance wall.
\end{hcorollary}

\begin{proof}
\emph{The letters.} With $\theta_{\cE\setminus T}=\theta_T$, $\Psi(\alpha)=\sum_T\theta_T\alpha^{\cE\setminus T}=\sum_T\theta_{\cE\setminus T}\alpha^{\cE\setminus T}=K(\alpha)$, the last sum running over the same bases. For a mark $e$ the bases $T\not\ni e$ correspond to the bases $T'=\cE\setminus T\ni e$, with $(\cE-e)\setminus T=T'-e$ and $\theta_T=\theta_{T'}$, so
\[
\Psi_{\setminus e}=\sum_{T\not\ni e}\theta_T\,\alpha^{(\cE-e)\setminus T}=\sum_{T'\ni e}\theta_{T'}\,\alpha^{T'-e}=K_{/e}
\]
(Definition~\hb{def:delcon}), and likewise $\Psi_{/e}=K_{\setminus e}$. The Skeleton values of the scale of $e$ are $K_{\setminus e}/(\alpha_eK_{/e})$ and $\Psi_{/e}/(\alpha_e\Psi_{\setminus e})$ (Corollary~\hb{cor:stage}), which coincide.

\emph{The table.} Bases and complements of bases are both bases, so $r=n-r$. A coloop lies in every basis, hence in no complement of one, and a loop lies in no basis, hence in every complement; both contradict the complementation, so no mark is a loop or a coloop, and the deletion $(\cB\setminus e,\theta)$ is a table of rank $r$ on $\cE-e$.

\emph{The walls at $\{e\}$ and at $\cE-e$.} $d(\{e\})=d^\ast(\{e\})=0$, so $\kappa_{\{e\}}=m_e$. Also $d(\cE-e)=r-r(\{e\})=r-1$ and $d^\ast(\cE-e)=n-1-r(\cE-e)=n-1-r=r-1$. On the balance wall $m_{\cE}=sr+t(n-r)=(s+t)r$, so $\kappa_{\cE-e}=m_{\cE}-m_e-(s+t)(r-1)=s+t-m_e$. The proper walls being positive, $0<m_e<s+t$.

\emph{Integrating the mark.} Let $g\neq e$ and $R=\cE-e$. By Theorem~\hb{thm:projective} the period is the integral of $\alpha^{\mvec-\one}K^{-s}\Psi^{-t}=\alpha^{\mvec-\one}K^{-(s+t)}$ over the section $\alpha_g=1$. By Tonelli's theorem the integral over $\alpha_e=x$ may be taken first, at fixed $\alpha_R$: with $K=K_{\setminus e}+xK_{/e}$, the Beta corner of a linear letter (Proposition~\hb{prop:scalebreta} with one coefficient) gives
\[
\int_0^\infty x^{m_e-1}\big(K_{\setminus e}+xK_{/e}\big)^{-(s+t)}dx=\Gst21{m_e,\ s+t-m_e}{s+t}K_{\setminus e}^{-(s+t-m_e)}K_{/e}^{-m_e},
\]
convergent as $0<m_e<s+t$, and $K_{/e}=\Psi_{\setminus e}$.

\emph{The deletion.} What remains is $\Gst21{m_e,\ s+t-m_e}{s+t}$ times the integral, over the section $\alpha_g=1$ of the orthant of $R$, of
\[
\alpha_R^{\mvec|_R-\one}K_{\setminus e}^{-(s+t-m_e)}\Psi_{\setminus e}^{-m_e}\,d\alpha_{R-g},
\]
the family word of the deletion at the exponents $(s+t-m_e,m_e)$. Its balance wall holds: $m_{\cE-e}=(s+t)r-m_e$, which is $(s+t-m_e)r+m_e(n-1-r)$ as $n-1-r=r-1$. By Theorem~\hb{thm:projective} the section integral is $P_{\cB\setminus e}(\mvec|_{\cE-e};s+t-m_e,m_e)$; it is finite because $P_F$ is, so the walls of the deletion are positive there. Since $K=\Psi$, $P_F(\mvec;s,t)$ depends on the exponents only through $s+t$. At $s=t=\sigma$ and $\mvec=\sigma\one$, $m_e=\sigma$, which is the last identity.
\end{proof}

\begin{htheorem}[\hb{thm:atomforward} (The scale integral of one mark)]
As stated; parts (i) and (ii) are stated for every family.
\end{htheorem}

\begin{proof}
(i) By Theorem~\hb{thm:scale} with $B=\{e\}$ and $m_B=m_e$, at the gauge $(\alpha_e,\alpha_R)$,
\[
\operatorname{Sc}_{\{e\}}(\alpha;\mvec;\mathbf s,\mathbf t)=\int_0^\infty x^{m_e-1}\prod_jK_j(x\alpha_e,\alpha_R)^{-s_j}\Psi_j(x\alpha_e,\alpha_R)^{-t_j}\,dx ,
\]
convergent since the walls $\kappa_{\{e\}}$ and $\kappa^{\mathrm{top}}_{\{e\}}$ are positive. At $\alpha_e=1$, Definition~\hb{def:delcon} gives $K_j=K_{j,\setminus e}(\alpha_R)+xK_{j,/e}(\alpha_R)$ and $\Psi_j=x\Psi_{j,\setminus e}(\alpha_R)+\Psi_{j,/e}(\alpha_R)$ for every active table, the deletion and the contraction being tables on $R$ as $e$ is neither a loop nor a coloop there. At $\alpha_R=e^{2z_R}$, Theorem~\hb{thm:tworeadings}(i) for these tables gives $K_{j,\star}(\alpha_R)=e^{\ip\one{z_R}}\Lambda_{j,\star}(z_R)$ and $\Psi_{j,\star}(\alpha_R)=e^{\ip\one{z_R}}\Lambda_{j,\star}(-z_R)$ for $\star\in\{\setminus e,/e\}$. So
\begin{gather*}
K_j=e^{\ip\one{z_R}}\big(\Lambda_{j,\setminus e}(z_R)+x\Lambda_{j,/e}(z_R)\big),\\
\Psi_j=e^{\ip\one{z_R}}\big(\Lambda_{j,/e}(-z_R)+x\Lambda_{j,\setminus e}(-z_R)\big),
\end{gather*}
and the powers carry the factor $e^{-\sum_j(s_j+t_j)\ip\one{z_R}}=e^{-2\sigma\ip\one{z_R}}$, an inactive table contributing $1$. This is the display, and the rest of the integrand depends on the room only through the functions $\Lambda_{j,\setminus e}$ and $\Lambda_{j,/e}$. Corollary~\hb{cor:fiber} at the bunch $\{e\}$, whose shape simplex is the point $u_e=1$, gives $P_F(\mvec;\mathbf s,\mathbf t)=\int\alpha_R^{\mvec|_R-\one}\operatorname{Sc}_{\{e\}}((1,\alpha_R);\mvec;\mathbf s,\mathbf t)\,d\alpha_{R-g}$ on the section $\alpha_g=1$, $g\in R$. With $\alpha_R=e^{2z_R}$, $\alpha_R^{\mvec|_R-\one}d\alpha_{R-g}=2^{n-2}e^{2\ip{\mvec|_R}{z_R}}dz_{R-g}$ there, $R-g$ having $n-2$ marks.

(ii) A realization that is an atom is a realization of a molecule (Definition~\hb{def:view}), so $F$ is connected, and Theorem~\hb{thm:gaugestep} gives $\alpha^\star\propto\one$. By Theorem~\hb{thm:normalizer}(ii), in $y=\log u$ modulo $\mathbb R\one$ the law $\mathrm{Breta}^\circ(F;\btau;\mathbf s,\mathbf t)$ has a density proportional to $e^{\ip\btau y-\Phi(y)}$. The target is the usage at $\alpha^\star$ (Definition~\hb{def:native}),
\[
\tau_e=\sum_j\big(s_jR_{j,e}(\alpha^\star)+t_jR^\ast_{j,e}(\alpha^\star)\big)=\partial_{y_e}\Phi(y^\star),\qquad y^\star=\log\alpha^\star,
\]
with the potential $\Phi$ of Definition~\hb{def:potential}. With $y=y^\star+v$ the density is proportional to
\[
\exp\big(\ip{\nabla\Phi(y^\star)}{v}-\Phi(y^\star+v)\big).
\]
By Proposition~\hb{prop:onewriting} there are orders $\mathbf g_e$ of the rooms in which the potential written from $e$ at the gauge $\alpha^\star$ (Definition~\hb{def:view}),
\[
(v_0,v_1,\dots,v_{n-1})\longmapsto\Phi_{\mathbf X}\Big(y^\star+v_0\one_e+\sum_kv_k\one_{g_k}\Big)\qquad(\mathbf X\in\{\mathbf E,\mathbf O\}),
\]
is one pair of functions for every $e$. Hence $v\mapsto\Phi(y^\star+v)$ read in the coordinates $(v_e,v_{g_1},\dots,v_{g_{n-1}})$ is one function for every $e$, and so is its gradient at $0$; the density read in these coordinates is one function for every $e$, and $dv$ is invariant under the renaming. The density is invariant along $\one$, so it is a function of the differences $v_{g_k}-v_e$, and the law of these differences does not depend on $e$. As $y^\star$ is a multiple of $\one$, $y_{g_k}-y_e=v_{g_k}-v_e$.

(iii) The reference (proof of thm:atomforward) proves this part, a statement for one table.
\end{proof}

\begin{htheorem}[\hb{thm:uniformjunction} (The junction under level form)]
As stated, for every family at the centre.
\end{htheorem}

\begin{proof}
Parts (ii) and (iii) concern the Gramma $\sG_{0;(\chi_p)_p}((\upsilon_{|p|})_p)$ alone and are proved in the reference (proof of thm:uniformjunction, parts (ii) and (iii)).

\emph{The preamble.} By Proposition~\hb{prop:levelformQ}(iii) for every table, every half-sum has level form and the square of every difference is unchanged by the permutations of $B$; $q_{j,b-k}(-z_R)=q_{j,k}(z_R)$ is Proposition~\hb{prop:levelformQ}(i). On $z_B$ each $(\partial_wL_{j,w})^2$ has exponents in $\{-2,0,2\}^B$, being a sum of products of two terms $e^{\ip{z_B}{\chi_p}}$. A function of $z_B$ of this kind that is unchanged by the permutations of $B$ is a combination of the products $e^{-2\ip\one{z_B}}e_k(e^{2z_B})e_{k'}(e^{2z_B})$, $k\le k'$: such a product is the symmetric sum of the monomials with $k$ exponents $2$ and $b-k'$ exponents $-2$, plus terms with fewer of both, and induction removes the terms (the reference, proof of thm:uniformjunction, part (i)). The coefficients are Laurent polynomials in $e^{z_R}$, the $\Lambda_{j,(p)}$ being such. The product over $j$ of $(\partial_wL_{j,w})^{2l_j}$ therefore expands in products of $2|\mathbf l|$ of the $e_k$, with Laurent polynomial coefficients.

(i) At the centre $\bmu=0$, and the first step of the proof of Theorem~\hb{thm:junctionpattern} gives
\[
P_F(\sigma\one;\boldsymbol\sigma,\boldsymbol\sigma)=2^{n-1}\int_{\mathbb R^{\cE-g}}\prod_j\big(\Lambda_j(z)\Lambda_j(-z)\big)^{-\sigma_j}\,dz_{\cE-g}\Big|_{z_g=0}.
\]
For each table $\Lambda_j(z)\Lambda_j(-z)=L_{j,1/2}(z)^2-\frac14(\partial_wL_{j,w}(z))^2$, and Proposition~\hb{prop:pairladder} with $a=\Lambda_j(z)$, $b=\Lambda_j(-z)$, $s=t=\sigma_j$, the midpoint $L_{j,1/2}(z)$ and $x=\partial_wL_{j,w}(z)/(2L_{j,1/2}(z))\in(-1,1)$ gives
\[
\big(\Lambda_j(z)\Lambda_j(-z)\big)^{-\sigma_j}=\sum_{l_j\ge0}\frac{(\sigma_j)_{l_j}}{l_j!\,4^{l_j}}\big(\partial_wL_{j,w}(z)\big)^{2l_j}L_{j,1/2}(z)^{-2\sigma_j-2l_j},
\]
a series of nonnegative terms. Open each power of the half-sum,
\[
L_{j,1/2}^{-2\sigma_j-2l_j}=\Gamma(2\sigma_j+2l_j)^{-1}\int_0^\infty y_j^{2\sigma_j+2l_j-1}e^{-y_jL_{j,1/2}}\,dy_j ;
\]
by Tonelli's theorem the sums and integrals may be taken in the order of the statement. With $\upsilon_k=\sum_jy_jq_{j,k}(z_R)$,
\[
\sum_jy_jL_{j,1/2}(z)=\sum_k\upsilon_k\,e^{-\ip\one{z_B}}e_k\big(e^{2z_B}\big),
\]
and inserting $e^{-\ip\one{z_B}}e_k(e^{2z_B})$ under the exponential of minus this sum is $-\partial_{\upsilon_k}$. With the expansion of the preamble and Definition~\hb{def:gramma} (index $0$, the integral over $z_B$),
\begin{multline*}
\int_{\mathbb R^B}\prod_j\big(\partial_wL_{j,w}(z)\big)^{2l_j}e^{-\sum_jy_jL_{j,1/2}(z)}\,dz_B\\
=\sum_{|\mathbf h|=2|\mathbf l|}c^{(\mathbf l)}_{\mathbf h}(z_R)\Big((-\partial_{\boldsymbol\upsilon})^{\mathbf h}\,\sG_{0;(\chi_p)_p}\big((\upsilon_{|p|})_p\big)\Big)\Big|_{\boldsymbol\upsilon=\sum_jy_j\mathbf q_j(z_R)} .
\end{multline*}
Each term converges. The family being a molecule, some active table $j$ has $c_j(B)\ge1$ (Definition~\hb{def:family}), so its patterns at $B$ fill at least two consecutive levels $k,k+1$ (Theorem~\hb{thm:taylorcube}), and $q_{j,k},q_{j,k+1}>0$, each being a sum of nonnegative terms one of which is positive; by $q_{j,b-k}(-z_R)=q_{j,k}(z_R)$ also $q_{j,b-k},q_{j,b-k-1}>0$. The vectors $\chi_p$ of these four levels are closed under negation, $\chi_{B\setminus p}=-\chi_p$, and span $\mathbb R^B$, since $\chi_{p\cup\{e\}}-\chi_p=2\mathbf e_e$ for $|p|=k$, $e\notin p$; so they positively span $\mathbb R^B$, the exponential of $-\sum_jy_jL_{j,1/2}(z)$ decays in $z_B$ faster than every exponential, and the inserted exponentials are integrable against it.

(iv) For each table the swap $\tau$ of $e$ and $f$ keeps $B$, fixes $R$ and exchanges table $j$, the bijection of the tables being the identity. Proposition~\hb{prop:mapcoeff} gives $c_{j,\tau(p)}(1-w;z_R)=c_{j,p}(w;z_R)$. At $w=\frac12$, $c_{j,\{f\}}(\frac12;\cdot)=c_{j,\{e\}}(\frac12;\cdot)$; in the derivative, $\partial_wc_{j,\varnothing}=\partial_wc_{j,B}=0$ and $\partial_wc_{j,\{f\}}=-\partial_wc_{j,\{e\}}$. With $c_{j,B}(\frac12;z_R)=c_{j,\varnothing}(\frac12;-z_R)$ (Proposition~\hb{prop:levelformQ}(i)) this gives the two displayed forms. Evenness: $c_{j,\{e\}}(\frac12;-z_R)=c_{j,\{f\}}(\frac12;z_R)=c_{j,\{e\}}(\frac12;z_R)$ and $\partial_wc_{j,\{e\}}(w;-z_R)=-\partial_wc_{j,\{f\}}(w;z_R)=\partial_wc_{j,\{e\}}(w;z_R)$, by Proposition~\hb{prop:levelformQ}(i). In (i) change $(z_e,z_f)$ to $(z_+,z_-)$, with $dz_e\,dz_f=\frac12dz_+dz_-$. Then $\prod_j(\partial_wL_{j,w})^{2l_j}=(e^{z_-}-e^{-z_-})^{2|\mathbf l|}\prod_j(\partial_wc_{j,\{e\}})^{2l_j}$ and
\[
\sum_jy_jL_{j,1/2}(z)=e^{-z_+}C_\varnothing(z_R)+e^{z_+}C_\varnothing(-z_R)+\big(e^{z_-}+e^{-z_-}\big)C_e(z_R).
\]
The integral over $z_+$ is $\sG_{0;(1,-1)}\big(C_\varnothing(-z_R),C_\varnothing(z_R)\big)$ (Definition~\hb{def:gramma} on one scale $e^{z_+}$), and the integral over $z_-$ is Lemma~\hb{lem:hn} with $\xi=C_e(z_R)$ and the power $|\mathbf l|$, $4^{|\mathbf l|}(\frac12)_{|\mathbf l|}C_e^{-|\mathbf l|}\sG_{|\mathbf l|;(1,-1)}(C_e,C_e)$. The constants collect to $2^{n-1}\cdot\frac12\cdot4^{|\mathbf l|}\prod_j4^{-l_j}=2^{n-2}$, which is the display; every term is nonnegative.
\end{proof}

\begin{sremark}
With one table, $(\frac12)_l(\sigma)_l/(l!\,\Gamma(2\sigma+2l))=\frac1{l!}\Gst23{\sigma+l,\ \frac12+l}{2\sigma+2l,\ \sigma,\ \frac12}$ and $y^{2\sigma+2l-1}(yc_{\{e\}})^{-l}=y^{2\sigma+l-1}c_{\{e\}}^{-l}$, which is the reference's form of (iv).
\end{sremark}

\begin{htheorem}[\hb{thm:classcloses} (A series connection closes at its reach)]
As stated, for series connections and for series connections of the dual table.
\end{htheorem}

\begin{proof}
(i) By Proposition~\hb{prop:seriesconnection}(i), $f_B=C^{-s}C'^{-t}\prod_if^\circ_i$, and by Proposition~\hb{prop:seriesconnection}(ii), which is Lemma~\hb{lem:latticeproduct} applied $q-1$ times, $f_B$ has a lattice expansion of reach $(\sum_i\sqrt{x_i})^2$ whose terms are the products of the terms of the parts collected by exponent. A function has at most one lattice expansion (Definition~\hb{def:latticeexp}), so these are the coefficients $\beta_{\eta,j}$ of Theorem~\hb{thm:nodeclosing}, and $\operatorname{reach}(B)\le(\sum_i\sqrt{x_i})^2<\inf\varrho_R$, the reach being an infimum. Theorem~\hb{thm:nodeclosing}(ii) applies: the even share closes toward the level $d$ on the lattice of continued periods of the bunch family, with the bound for every $x>\operatorname{reach}(B)$, in particular for every $x>(\sum_i\sqrt{x_i})^2$. Under the hypotheses of Proposition~\hb{prop:seriesconnection}(iii), $\operatorname{reach}(B)=(\sum_i\sqrt{\sup\varrho_i})^2$, and Theorem~\hb{thm:nodeintervals}(ii) gives the root rate $\operatorname{reach}(B)/\inf\varrho_R$. For a series bunch the parts are the marks (Proposition~\hb{prop:seriesconnection}), with
\begin{align*}
f^\circ_e(y)&=\int_0^\infty x^{m_e-s-1}e^{-w_e/x-yw_ex}\,dx=\sG_{m_e-s;(1,-1)}(yw_e,w_e)\\
&=(yw_e)^{s-m_e}\sG_{m_e-s;(1,-1)}\big(1,w_e^2y\big)
\end{align*}
by the torus homogeneity of the Gramma; by Lemma~\hb{lem:onemark}(iii) it has a lattice expansion of reach $w_e^2$, a power of $y$ keeping the reach (Lemma~\hb{lem:latticeproduct}). Proposition~\hb{prop:seriesproduct}(iii) gives $\operatorname{reach}(B)=\inf\varrho_B=(\sum_ew_e)^2$ exactly, and Proposition~\hb{prop:seriesproduct}(i) gives the coefficients.

(ii) The reference (proof of thm:classcloses, part (ii)) gives the argument for a parallel bunch, and it applies to a connection. The dual table at the exponents $(t,s)$ has the same period (Proposition~\hb{prop:dualfamily}). Since $K^\ast=\Psi$ and $\Psi^\ast=K$ (Proposition~\hb{prop:identity}), the dual node has the bunch $B$, of width one in the dual table, the room $R$, the ratios $1/\varrho_B$ and $1/\varrho_R$, and its levels swapped: its level $d$ carries the even letter of the level $r$ of the given node. No slope on the room being negative, $\sup\varrho_R<\infty$ (Lemma~\hb{lem:slopesrange}), so $\inf(1/\varrho_R)>0$ and no slope on the room of the dual node is positive (the same lemma for the dual node). Part (i) for the dual node, whose bunch is a series connection with the reaches $x^\ast_i$, closes the even share toward its level $d$, the level $r$ of the given node, when $(\sum_i\sqrt{x^\ast_i})^2<1/\sup\varrho_R$, at a root rate at most $(\sum_i\sqrt{x^\ast_i})^2\sup\varrho_R$. For a parallel bunch, the dual pattern weights are $\theta'^{\ast}_B=\theta'_\varnothing$ and $\theta'^{\ast}_{B-e}=\theta'_{\{e\}}$, so the $w_e$ of the dual series bunch are $\theta'_{\{e\}}/\theta'_\varnothing$, $x^\ast_e=w_e^2$, and the root rate is exact by (i). The ratio of the bunch of the dual node is $1/\varrho_B=(\sum_ew_e/u_e)(\sum_ew_eu_e)$, whose least value $(\sum_ew_e)^2$ is taken at the point of equal coordinates of $\Delta_B$ (Proposition~\hb{prop:seriesproduct}(ii)); so $\sup\varrho_B=(\sum_ew_e)^{-2}$, the condition reads $\sup\varrho_R<\sup\varrho_B$, and the root rate is $(\sum_ew_e)^2\sup\varrho_R=\sup\varrho_R/\sup\varrho_B$.
\end{proof}

\begin{htheorem}[\hb{thm:onetree} (One tree suffices)]
As stated; the chain clause of (iii) is stated whenever some active list is connected.
\end{htheorem}

\begin{proof}
All of the theorem but the chain clause of (iii) is the reference's (proof of thm:onetree). For the chain clause, let some active list of the molecule $F$ be connected. The weights of lists factor across every split, so Corollary~\hb{cor:leafchain} applies: every peel leaves one childcule, the Foreman tree is a chain (Corollary~\hb{cor:treechain}), and at every realization, the centre and every lean, $|B_0|\le|B_1|\le\cdots\le|B_{d-1}|\le|L|$, $L$ the leaf. The shape of the bunch $B_k$ then has $|B_k|-1\le|L|-1$ integrations, and its node adds one, the total of the bunch against its room, by (ii). A molecule of one list is such an $F$, its list being connected.
\end{proof}

\section{Chapter~\hb{ch:glue}: gluing and what crosses it}\label{sec:ch20}

\begin{htheorem}[\hb{thm:block} (Replacement of marks by blocks)]
As stated: every mark is replaced, by a block of any size and by either rule.
\end{htheorem}

\begin{proof}
\emph{One mark, first rule.} Let $h$ be replaced by a block $X$ of $q$ marks by the first rule, the other marks left in place: the new sets are the bases $T\not\ni h$ and the sets $(T-h)\sqcup\{x\}$, $T\ni h$, $x\in X$, each with the weight $\theta_T$. If $h$ is neither a loop nor a coloop, $(\cB,\theta)$ is a cell on $(\cE-h)\sqcup\{h\}$ with the hinge $h$, and the parallel block $\binom{X\sqcup\{g\}}1$ is a cell on $X\sqcup\{g\}$ whose deletion of $g$ has the singletons of $X$ as bases and whose contraction of $g$ has the single basis $\varnothing$. By \hb{lem:gluemark} for these two cells, the sets $\{x\}\sqcup(T-h)$, $T\ni h$, and $\varnothing\sqcup T$, $T\not\ni h$, are the bases of a table of rank $1+r-1=r$: they are the new sets. If $h$ is a loop, no basis contains $h$, and the new sets are the bases of $\cB$ read on $\cE-h$, the marks of $X$ being loops: a table of rank $r$. If $h$ is a coloop, every basis contains $h$, and the new sets are the sets $(T-h)\sqcup\{x\}$: the direct sum of the contraction of $h$ and the table $\binom X1$, a table of rank $r$, the exchange property of a direct sum holding in each summand.

\emph{One mark, complementary rule.} The dual table has the bases $\cE\setminus T$ (\hb{prop:identity}). The first rule at $h$ on the dual table gives the sets $\cE\setminus T$ for $T\ni h$ and $(\cE\setminus T)-h+x$ for $T\not\ni h$, $x\in X$, the bases of a table on $(\cE-h)\sqcup X$ by the first case. Their complements in $(\cE-h)\sqcup X$ are $(T-h)\sqcup X$ for $T\ni h$ and $T\sqcup(X-x)$ for $T\not\ni h$, the sets of the complementary rule, and the complements of the bases of a table are the bases of its dual. Its rank is $(|\cE|-1+q)-(|\cE|-r)=r+q-1$.

\emph{All marks.} The rule at a mark $h$ decides the part of a new set in $X_h$ from whether $h\in T$ alone, so replacing the marks one at a time, in any order, gives the sets of the statement; by the two cases each step gives a table, and the step at $h\in\cE_2$ raises the rank by $q_h-1$. A new set $T'$ determines $T=\{h\in\cE_1:\ X_h\cap T'\neq\varnothing\}\sqcup\{h\in\cE_2:\ X_h\subseteq T'\}$ and the choice made in every block, so it carries the one weight $\theta_T$. A block of one mark gives back its mark under either rule.

\emph{Substitution.} When no mark is a loop or a coloop, $(\cB,\theta)$ is a core (\hb{def:substitution}), and the parallel block $\binom{X_h\sqcup\{g_h\}}1$ and the series block $\binom{X_h\sqcup\{g_h\}}{q_h}$ are cells, $g_h$ being neither a loop nor a coloop of either. In the parallel block the deletion of $g_h$ has the singletons of $X_h$ as bases and the contraction the single basis $\varnothing$; in the series block the deletion has the single basis $X_h$ and the contraction the sets $X_h-x$. So the sets of the substitution contain, for a basis $T$ of the core, one mark of $X_h$ for $h\in T\cap\cE_1$, none for $h\in\cE_1\setminus T$, all of $X_h$ for $h\in T\cap\cE_2$ and all but one for $h\in\cE_2\setminus T$, each with the weight $\theta_T$, every weight of a block being $1$: the table of the statement.

\emph{Letters.} By the correspondence between the new sets and the pairs of a basis $T$ and a choice in every block, $K'(\alpha)$ is the sum over $T$ of $\theta_T$ times the product over the blocks of the sum, over the choices, of the monomials of the chosen marks: $e_1(\alpha_{X_h})$ for $h\in T\cap\cE_1$, $1$ for $h\in\cE_1\setminus T$, $e_{q_h}(\alpha_{X_h})$ for $h\in T\cap\cE_2$ and $e_{q_h-1}(\alpha_{X_h})$ for $h\in\cE_2\setminus T$. Taking out $\prod_{h\in\cE_2}e_{q_h-1}$,
\[
K'=\prod_{h\in\cE_2}e_{q_h-1}\sum_T\theta_T\prod_{h\in T\cap\cE_1}e_1\prod_{h\in T\cap\cE_2}\frac{e_{q_h}}{e_{q_h-1}},
\]
the first display. The complement of a new set meets $X_h$ in all of $X_h$ but the chosen mark for $h\in T\cap\cE_1$, in all of $X_h$ for $h\in\cE_1\setminus T$, in nothing for $h\in T\cap\cE_2$ and in the missing mark for $h\in\cE_2\setminus T$, so
\[
\Psi'=\sum_T\theta_T\prod_{h\in T\cap\cE_1}e_{q_h-1}\prod_{h\in\cE_1\setminus T}e_{q_h}\prod_{h\in\cE_2\setminus T}e_1=\prod_{h\in\cE_1}e_{q_h-1}\sum_T\theta_T\prod_{h\in\cE_1\setminus T}\frac{e_{q_h}}{e_{q_h-1}}\prod_{h\in\cE_2\setminus T}e_1 ,
\]
which is $\prod_{h\in\cE_1}e_{q_h-1}\,\Psi(\beta^\ast)$ with $\beta^\ast_h=e_{q_h}/e_{q_h-1}$ on $\cE_1$ and $\beta^\ast_h=e_1$ on $\cE_2$, since $\Psi(\beta^\ast)=\sum_T\theta_T(\beta^\ast)^{\cE\setminus T}$. For one block at $h$, every other block of one mark, $K$ and $\Psi$ are affine in $\alpha_h$, $K=K_{\setminus h}+\alpha_hK_{/h}$ and $\Psi=\alpha_h\Psi_{\setminus h}+\Psi_{/h}$ (\hb{def:delcon}), and the two displays become the four formulas of the statement. In scale and shape, $e_k$ is homogeneous of degree $k$, $e_1(u)=|u|=1$ on $\Delta_X$, and the measure is \hb{lem:scaleshape}.
\end{proof}

\begin{htheorem}[\hb{thm:blockcounts} (Counts, walls and atoms under replacement)]
As stated: (i) and (ii) for blocks of any sizes and either rule at each mark; (iii) at every positive exponent pairs.
\end{htheorem}

\begin{proof}
(i) The marks and the rank are \hb{thm:block}. At $\alpha=\one$, $e_1=q_h$, $e_{q_h}=1$ and $e_{q_h-1}=q_h$ on a block of $q_h$ marks, so the first display of \hb{thm:block} gives the total weight $K'(\one)=\prod_{h\in\cE_2}q_h\cdot K\big((q_h)_{h\in\cE_1},(q_h^{-1})_{h\in\cE_2}\big)$, and expanding $K$ over the bases gives the sum displayed; it counts $q_h$ choices at $h\in T\cap\cE_1$ and $q_h$ at $h\in\cE_2\setminus T$. With $\cE_2=\varnothing$ it is $K((q_h)_h)$; with $\cE_1=\varnothing$ it is $\prod_hq_h\,K((q_h^{-1})_h)=\Psi((q_h)_h)$ (\hb{prop:identity}).

(ii) By \hb{thm:window}, $d'(Y)=\min_{T'}|T'\cap Y|$ and $d^{\ast\prime}(Y)=\min_{T'}|Y\setminus T'|$ over the new bases $T'$, which are the pairs of a basis $T$ and a choice in every block (proof of \hb{thm:block}). Fix $T$ and minimize over the choices, block by block. In $|T'\cap Y|$ a block contributes, at its least: for $h\in T\cap\cE_1$, $1$ if $X_h\subseteq Y$ and $0$ otherwise, the chosen mark avoiding $Y$ when it can; for $h\in\cE_1\setminus T$, $0$; for $h\in T\cap\cE_2$, $|X_h\cap Y|$; for $h\in\cE_2\setminus T$, $\max(|X_h\cap Y|-1,0)$, the missing mark taken in $Y$ when it can. Writing $|X_h\cap Y|=\max(|X_h\cap Y|-1,0)+[h\in\lceil Y\rceil]$ for $h\in\cE_2$, the least value at $T$ is
\[
\sum_{h\in\cE_2}\max\big(|X_h\cap Y|-1,0\big)+\big|T\cap\big((\lfloor Y\rfloor\cap\cE_1)\sqcup(\lceil Y\rceil\cap\cE_2)\big)\big|,
\]
and the least over $T$ is $d'(Y)$ as stated (\hb{thm:window}). In $|Y\setminus T'|$ a block contributes, at its least: for $h\in T\cap\cE_1$, $\max(|X_h\cap Y|-1,0)$; for $h\in\cE_1\setminus T$, $|X_h\cap Y|$; for $h\in T\cap\cE_2$, $0$; for $h\in\cE_2\setminus T$, $1$ if $X_h\subseteq Y$ and $0$ otherwise. In the same way $d^{\ast\prime}(Y)$ is as stated, and $c'(Y)=|Y|-d'(Y)-d^{\ast\prime}(Y)$ is \hb{def:skeleton}. Under the first rule alone, $r'(Y)=\max_{T'}|T'\cap Y|=\max_T|T\cap\lceil Y\rceil|=r(\lceil Y\rceil)$, $d'(Y)=d(\lfloor Y\rfloor)=r-r(\cE\setminus\lfloor Y\rfloor)$ (\hb{def:skeleton}) and $c'(Y)=r'(Y)-d'(Y)$ (\hb{thm:window}). Under the complementary rule alone, $|Y|=\sum_h\max(|X_h\cap Y|-1,0)+|\lceil Y\rceil|$, so
\[
c'(Y)=|\lceil Y\rceil|-d(\lceil Y\rceil)-d^\ast(\lfloor Y\rfloor)=|\lceil Y\rceil|-|\lfloor Y\rfloor|+r(\cE\setminus\lceil Y\rceil)+r(\lfloor Y\rfloor)-r .
\]
If $Y=\bigcup_{h\in S}X_h$, then $\lceil Y\rceil=\lfloor Y\rfloor=S$, $d'(Y)=\sum_{h\in S\cap\cE_2}(q_h-1)+d(S)$, $d^{\ast\prime}(Y)=\sum_{h\in S\cap\cE_1}(q_h-1)+d^\ast(S)$ and $|Y|=\sum_{h\in S}q_h$, so $c'(Y)=|S|-d(S)-d^\ast(S)=c(S)$. If $\varnothing\neq Y\subsetneq X_h$, then $\lceil Y\rceil=\{h\}$ and $\lfloor Y\rfloor=\varnothing$: for $h\in\cE_1$, $d'(Y)=0$ and $d^{\ast\prime}(Y)=|Y|-1+d^\ast(\{h\})$, and for $h\in\cE_2$, $d'(Y)=|Y|-1+d(\{h\})$ and $d^{\ast\prime}(Y)=0$; a seed table connected on at least two marks has $c(\{h\})=1-d(\{h\})-d^\ast(\{h\})>0$, so $d(\{h\})=d^\ast(\{h\})=0$ and $c'(Y)=1$.

(iii) Write $X_e=\{e\}\times\{1,\dots,q\}$. For $h\in H$, carrying table $j$ to table $\tau(j)$ (\hb{def:carry}), put $\hat h(e,i)=(h(e),i)$; the rule and $q$ are the same at every mark, so $\hat h$ maps the replaced table $j$ onto the replaced table $\tau(j)$ with the same weights, and the two tables have the same pairs: $\hat h$ carries the replaced family to itself. A permutation of $X_e$ fixing the other new marks maps every replaced table onto itself with its weights, the sets of the rule being symmetric in the marks of $X_e$. The group generated by these bijections is transitive on the new marks, $H$ being transitive on $\cE$.

The replaced family is connected. The seed family is one molecule: a molecule $M\neq\cE$ would be a union of molecules of every table (\hb{prop:directsum}), with $|T\cap M|$ constant on the bases of every table and so $c_j(M)=0$ for every $j$ (\hb{thm:window}), all tables being active at positive pairs. Let $e\neq f$ be linked in table $j$: by \hb{lem:linkexchange} some basis $T$ of table $j$ has $e\in T$, $f\notin T$ and $T-e+f$ a basis. Under the first rule, a new basis $T'$ from $T$ choosing $x$ in $X_e$ gives $T'-x+x'$, the new basis from $T$ choosing $x'$, and $T'-x+y$, $y\in X_f$, the new basis from $T-e+f$ choosing $y$ with the same choices elsewhere; under the complementary rule, the new basis $T'$ from $T-e+f$ missing $x'$ in $X_e$ gives $T'-x+x'$, the one missing $x$, and the new basis $T''$ from $T$ missing $y$ in $X_f$ gives $T''-x+y$, the new basis from $T-e+f$ missing $x$. By \hb{lem:linkexchange} any two marks of $X_e$ are linked, and every mark of $X_e$ is linked to every mark of $X_f$. Every mark is linked to another mark, the seed family being one molecule on at least two marks, so the new marks form one molecule. Then no nonempty proper set $Y$ of new marks has $c'_j(Y)=0$ in every table: by \hb{thm:window} $|T'\cap Y|$ would be constant on the bases of every table, and by \hb{lem:circuit} no circuit would join $Y$ to its complement.

Finally $((s_j,t_j)_j;\one)$ is a realization of the replaced family, and every bijection above fixes the table's gauge, $\hat h^\ast\log\one=\log\one$. By \hb{prop:certify} the reading of the image of a mark equals the reading of the mark, so the reading is constant on the new marks, which form one molecule: the realization is an atom (\hb{def:view}).
\end{proof}

\begin{sremark}
In exact rational arithmetic (receipt \texttt{check\_blocks\_splice.py}): on the seed of rank two on $\{a,b,c,d\}$ whose bases are the pairs other than $\{c,d\}$, with weights $1,2,\frac32,5,\frac13$, blocks of $2,3,1,2$ marks, the first rule at $a,c$ and the complementary rule at $b,d$, the new sets satisfy the exchange property, the rank is $r+\sum_{h\in\cE_2}(q_h-1)$, the total weight is that of (i), $d'(Y)$ and $d^{\ast\prime}(Y)$ are those of (ii) at every one of the $2^8$ sets $Y$, and unions of full blocks keep the width; with the first rule at every mark and the same sizes, $r'(Y)=r(\lceil Y\rceil)$, $d'(Y)=r-r(\cE\setminus\lfloor Y\rfloor)$ and the total weight is $K((q_h)_h)$.
\end{sremark}

\begin{hproposition}[\hb{prop:mot-dualletters} (Dual pairs)]
As stated.
\end{hproposition}

\begin{proof}
Write $\theta$ for the one system of weights of $\mathcal A$. (i) The bases of $\mathcal A$ are the $X\in\mathcal A\setminus g$ and the $Y+g$ with $Y\in\mathcal A/g$ (\hb{def:delcon}). If $\mathcal A$ is a dual pair, the complement $(B\setminus X)+g$ of $X$ is a basis with the weight $\theta_X$, and the complement $B\setminus Y$ of $Y+g$ is a basis of $\mathcal A\setminus g$ with the weight $\theta_{Y+g}$: the bases are closed under complement with equal weights. Conversely, if they are, $X\in\mathcal A\setminus g$ gives the basis $(B\setminus X)+g$ with the weight $\theta_X$, so $B\setminus X\in\mathcal A/g$, and $Y\in\mathcal A/g$ gives $B\setminus Y\in\mathcal A\setminus g$: $\mathcal A/g$ is the dual of $\mathcal A\setminus g$ with its weights. Since $g$ is neither a loop nor a coloop, $\mathcal A\setminus g$ has the rank $r_{\mathcal A}$ and $\mathcal A/g$ the rank $r_{\mathcal A}-1$, while the dual of $\mathcal A\setminus g$ has the rank $|B|-r_{\mathcal A}$; so $|B|=2r_{\mathcal A}-1$.

(ii) $K_{\mathcal A/g}=\sum_Y\theta_{Y+g}\alpha^Y=\sum_X\theta_X\alpha^{B\setminus X}=\psi$ and $\Psi_{\mathcal A/g}=\sum_Y\theta_{Y+g}\alpha^{B\setminus Y}=\sum_X\theta_X\alpha^X=k$. So $\gamma^K_{\mathcal A}=K_{\mathcal A\setminus g}/K_{\mathcal A/g}=k/\psi$ and $\gamma^\Psi_{\mathcal A}=\Psi_{\mathcal A/g}/\Psi_{\mathcal A\setminus g}=k/\psi$ (\hb{def:substitution}).

(iii) The cell $\binom{\{h,g\}}1$ has the deletion with the basis $\{h\}$ and the contraction with the basis $\varnothing$, the dual of $\{h\}$ on $\{h\}$, so it is a dual pair with $k=\alpha_h$ and $\psi=1$. By \hb{thm:substletters} and (ii), $K=\prod_iK_{\mathcal A_i/g_i}\,K_{\mathcal H}\big((k_i/\psi_i)_i\big)=\pi K_{\mathcal H}\big((k_i/\psi_i)_i\big)$ and $\Psi=\prod_i\Psi_{\mathcal A_i\setminus g_i}\,\Psi_{\mathcal H}\big((k_i/\psi_i)_i\big)=\pi\Psi_{\mathcal H}\big((k_i/\psi_i)_i\big)$. $K_{\mathcal H}$ and $\Psi_{\mathcal H}$ are homogeneous of the degrees $r_{\mathcal H}$ and $q-r_{\mathcal H}$, so multiplying the arguments by $\pi$ gives the two identities; both sides are polynomials, $r_{\mathcal H}\ge1$ and $q-r_{\mathcal H}\ge1$ since no mark of the core is a loop or a coloop.
\end{proof}

\begin{hcorollary}[\hb{cor:mot-widthbound} (The width of the marks left in place)]
As stated.
\end{hcorollary}

\begin{proof}
In \hb{thm:substwalls} take $S\subseteq L$: the part of $S$ in the cell $\binom{\{l,g\}}1$ of a mark $l\in S$ is $\{l\}$, with $d(\{l\})=d^\ast(\{l\})=0$ and $d(\{l,g\})=d^\ast(\{l,g\})=1$, the bases being $\{l\}$ and $\{g\}$ (\hb{thm:window}); every other cell meets $S$ in $\varnothing$, and its hinge, neither a loop nor a coloop, has $d(\{g\})=d^\ast(\{g\})=0$. So $D(S)=D^\ast(S)=S$, read on the marks of the core, and $d(S)=d_{\mathcal H}(S)$, $d^\ast(S)=d^\ast_{\mathcal H}(S)$, whence $c(S)=|S|-d(S)-d^\ast(S)=c_{\mathcal H}(S)$ (\hb{def:skeleton}). By \hb{def:skeleton}, $d_{\mathcal H}(S)=r_{\mathcal H}-r_{\mathcal H}(L\sqcup H'\setminus S)$, so
\[
c_{\mathcal H}(S)=r_{\mathcal H}(S)+r_{\mathcal H}(L\sqcup H'\setminus S)-r_{\mathcal H}\le r_{\mathcal H}(L\sqcup H'\setminus S)\le|L\setminus S|+|H'|,
\]
the rank of a set being at most its size and $r_{\mathcal H}(S)\le r_{\mathcal H}$. In $\binom{B\sqcup\{h_1,\dots,h_c\}}c$ the bases are the sets of $c$ marks, which meet $B$ in every number from $0$ to $c$, so $B$ has the width $c$ (\hb{thm:window}); no $h_i$ is a loop or a coloop, and by the first part the width stays $c$ after any cells are substituted at the $h_i$.
\end{proof}

\begin{hlemma}[\hb{lem:twosumpieces} (Two-sums of cells; hinge walls)]
As stated, for cell families glued table by table.
\end{hlemma}

\begin{proof}
That $\mathcal A\oplus_g\mathcal A'$ is a table of rank $r_{\mathcal A}+r_{\mathcal A'}-1$, plain when the cells are, in which $B$ is a bunch of width one with the cells $\mathcal A$ and $\mathcal A'$, and that a table with a bunch of width one and factoring weights is the two-sum of its two cells, is the reference, proof of lem:twosumpieces(i), the weights multiplied letter by letter. The two-sum is $\binom{\{h_1,h_2\}}1\langle\mathcal A,\mathcal A'\rangle$: the bases $\{h_1\}$ and $\{h_2\}$ of the core give the two kinds of sets of \hb{lem:gluemark} (\hb{def:substitution}). So $F$ is substituted table by table, with the core $\binom{\{h_1,h_2\}}1$ in every table and the cells $\mathcal A_j$ at $h_1$ and $\mathcal A'_j$ at $h_2$, and \hb{thm:substwalls} at $S=B$ gives the wall $m_B-\sum_j\big(s_j(r_{\mathcal A_j}-1)+t_j(|B|-r_{\mathcal A_j})\big)=\kappa_{F_B}$, and at $S=R$ the wall $\kappa_{F_R}$. Table $j$ of $F$ has the rank $r_j=r_{\mathcal A_j}+r_{\mathcal A'_j}-1$ on $n=|B|+|R|$ marks, so
\[
\kappa_{F_B}+\kappa_{F_R}=m_{\cE}-\sum_j\big(s_j(r_j-1)+t_j(n-r_j-1)\big)=\kappa_{\cE}+\sum_j(s_j+t_j)=\kappa_{\cE}+2\sigma .
\]
With the mass $2\sigma-\kappa_{F_B}$ on $g$, the imbalance of $F_B$ on $B\sqcup\{g\}$ is
\[
m_B+2\sigma-\kappa_{F_B}-\sum_j\big(s_jr_{\mathcal A_j}+t_j(|B|+1-r_{\mathcal A_j})\big)=2\sigma-\sum_j(s_j+t_j)=0,
\]
on the balance wall of $F$ or off it.
\end{proof}

\begin{htheorem}[\hb{thm:hingepairs} (Splicing hinges in pairs)]
As stated; (iii) holds for any number of cells and any graph of pairs.
\end{htheorem}

\begin{proof}
(i) and (ii) are the reference, proof of thm:hingepairs(i),(ii). (iii) In the direct sum of the cells the letters are $K=\prod_iK_{\mathcal A_i}$ and $\Psi=\prod_i\Psi_{\mathcal A_i}$ (\hb{prop:directsum}), and by \hb{def:channels}, with $G=\bigsqcup_iG_i$,
\[
K=\sum_{(p_i)_i}\alpha_G^{\,\sqcup_ip_i}\prod_iK_{\mathcal A_i,(p_i)},\qquad \Psi=\sum_{(p_i)_i}\alpha_G^{\,G\setminus\sqcup_ip_i}\prod_i\Psi_{\mathcal A_i,(p_i)},
\]
over $p_i\subseteq G_i$, absent channels counting as zero. Splicing a pair $\{a,b\}$ acts on a monomial $\alpha_G^P$ by $(\partial_{\alpha_a}+\partial_{\alpha_b})\alpha_G^P|_{\alpha_a=\alpha_b=0}$, which is $\alpha_G^{P-a}$ when $a\in P\not\ni b$, $\alpha_G^{P-b}$ when $b\in P\not\ni a$, and $0$ otherwise; on $\alpha_G^{G\setminus P}$ it is likewise nonzero exactly when $P$ contains one of $a$, $b$. The splicings of disjoint pairs commute and act on different variables, so splicing all pairs keeps the terms in which $P=\sqcup_ip_i$ is $p\sqcup J$ with $p\subseteq G'$ and $J$ containing exactly one hinge of every pair, and leaves $\alpha_{G'}^p$ in $K$ and $\alpha_{G'}^{G'\setminus p}$ in $\Psi$. By (ii) the result is the letters of the spliced cell $\mathcal A$, and by the Taylor cube (\hb{thm:taylorcube}) its channels are the coefficients of these monomials:
\[
K_{\mathcal A,(p)}=\sum_J\prod_iK_{\mathcal A_i,((p\sqcup J)\cap G_i)},\qquad \Psi_{\mathcal A,(p)}=\sum_J\prod_i\Psi_{\mathcal A_i,((p\sqcup J)\cap G_i)} .
\]
The marks other than hinges are $\bigsqcup_iB_i$, and the rank is $\sum_ir_{\mathcal A_i}-k$ by (ii). The degrees add up: $\sum_i\big(r_{\mathcal A_i}-|(p\sqcup J)\cap G_i|\big)=\sum_ir_{\mathcal A_i}-k-|p|$, since $|J|=k$, and $\sum_i\big(n_{\mathcal A_i}-r_{\mathcal A_i}+|(p\sqcup J)\cap G_i|\big)=n_{\mathcal A}-r_{\mathcal A}+|p|$, since $n_{\mathcal A}-r_{\mathcal A}=\sum_i(n_{\mathcal A_i}-r_{\mathcal A_i})+k$. Each channel of $\mathcal A$ is a sum of products of one channel of each cell, so the map is multilinear. For two cells with the pairs named by $I$, a set $J$ is given by the set $I'\subseteq I$ of the pairs whose hinge on the side of $\mathcal A_1$ it contains, $(p\sqcup J)\cap G_1=p_1\sqcup I'$ and $(p\sqcup J)\cap G_2=p_2\sqcup(I\setminus I')$: the bilinear map displayed. With $k=1$ and no other hinges, $K_{\mathcal A,(\varnothing)}=K_{\mathcal A_1\setminus g_1}K_{\mathcal A_2/g_2}+K_{\mathcal A_1/g_1}K_{\mathcal A_2\setminus g_2}$ and $\Psi_{\mathcal A,(\varnothing)}=\Psi_{\mathcal A_1\setminus g_1}\Psi_{\mathcal A_2/g_2}+\Psi_{\mathcal A_1/g_1}\Psi_{\mathcal A_2\setminus g_2}$, the letters of \hb{cor:blocksclosures}(iii).
\end{proof}

\begin{sremark}
In exact arithmetic (receipt \texttt{check\_blocks\_splice.py}): three cells in a closed path, a weighted $\binom{\{x_1,x_2\}\sqcup G_1}3$ with three hinges, one of them unpaired, a weighted $\binom{\{y_1,y_2\}\sqcup G_2}2$ and a weighted $\binom{\{z_1\}\sqcup G_3}1$ with two hinges each, spliced along three pairs, one between each two cells: both channels of the spliced cell at both values of the unpaired hinge are the sums over $J$ of the products of the channels of the three cells.
\end{sremark}

\begin{htheorem}[\hb{thm:interfacerank} (Multivectors bound what crosses)]
As stated; (iii) holds at any tensor rank.
\end{htheorem}

\begin{proof}
(i) and (ii) are the reference, proof of thm:cutrank(i),(ii), with the interface rank for the cut rank; (ii) is proved there at any tensor rank, the interface rank being subadditive. (iii) By (i) and (ii), $\operatorname{rank}\Theta^{(k)}_B\le\operatorname{ir}_k(\boldsymbol\omega)^2\le i_k^2$ with $i_k=\sum_l\binom{c_l}{k-d_l}$, and summing over the levels bounds the interface dimension $\sum_k\operatorname{rank}\Theta^{(k)}_B$ (\hb{def:levelweights}); for a real multivector (i) gives $\operatorname{rank}\Theta^{(k)}_B\le\frac12\operatorname{ir}_k(\boldsymbol\omega)(\operatorname{ir}_k(\boldsymbol\omega)+1)$, and $x\mapsto\frac12x(x+1)$ increases on $[0,\infty)$. The bounds involve only the widths and the floors of the slices. A determinantal table has a multivector of tensor rank one, $\boldsymbol\omega_T=\det A_T$ (\hb{def:multivector}); a set $T$ is a basis exactly when $\det A_T\neq0$, so the column space of $A_S$ has the dimension $r(S)$, the column spaces of $A_B$ and $A_R$ span $\mathbb F^r$ and meet in dimension $r(B)+r(R)-r=c(B)$, and $d_1=r-r(R)=d(B)$ (\hb{def:skeleton}). So $\operatorname{ir}_k(\boldsymbol\omega)\le\binom{c(B)}{k-d(B)}$, and $\sum_k\binom{c(B)}{k-d(B)}^2=\sum_i\binom{c(B)}i^2=\binom{2c(B)}{c(B)}$. At width one $\binom1{k-d(B)}=1$ at $k=d(B)$ and $k=d(B)+1$: one bunch letter at each level.
\end{proof}

\section{Chapter~\hb{ch:construct}: periods of constructions}\label{sec:ch21}

\begin{hdefinition}[\hb{def:twist} (Twists)]
As stated: the torus relation for every twist, every vector $\boldsymbol\lambda$ on the marks and all constants $c_j,c^\ast_j$.
\end{hdefinition}

\begin{proof}
Put $\mathbf x'=(c_j\boldsymbol\lambda\circ\mathbf x^K_j,\ c^\ast_j\boldsymbol\lambda\circ\mathbf x^\Psi_j)_j$. Every basis of table $j$ has $r_j$ marks and its complement $r^\ast_j$, so $K_j$ is homogeneous of degree $r_j$ and $\Psi_j$ of degree $r^\ast_j$, and the letters of $F[\mathbf x']$ at $\alpha$ are
\[
K_j(c_j\boldsymbol\lambda\circ\mathbf x^K_j\circ\alpha)=c_j^{r_j}K_j\big(\mathbf x^K_j\circ(\boldsymbol\lambda\circ\alpha)\big),\qquad \Psi_j(c^\ast_j\boldsymbol\lambda\circ\mathbf x^\Psi_j\circ\alpha)=c_j^{\ast\,r^\ast_j}\Psi_j\big(\mathbf x^\Psi_j\circ(\boldsymbol\lambda\circ\alpha)\big).
\]
Write $\omega_{\mathbf x}=\beta^{\mvec-\one}\prod_jK_j(\mathbf x^K_j\circ\beta)^{-s_j}\Psi_j(\mathbf x^\Psi_j\circ\beta)^{-t_j}\,d\beta$ for the form of $F[\mathbf x]$. The map $\Phi:\alpha\mapsto\boldsymbol\lambda\circ\alpha$ is a linear bijection of the open orthant onto itself that commutes with positive scalings, and with $\beta_e=\lambda_e\alpha_e$ one has $\beta^{\mvec-\one}d\beta=\boldsymbol\lambda^{\mvec}\alpha^{\mvec-\one}d\alpha$. So the form of $F[\mathbf x']$ is
\[
\omega_{\mathbf x'}=\boldsymbol\lambda^{-\mvec}\prod_jc_j^{-r_js_j}c_j^{\ast\,-r^\ast_jt_j}\ \Phi^\ast\omega_{\mathbf x}.
\]
A twist does not change degrees, so on the balance wall both forms are homogeneous of degree zero, and each period is the integral of its form over the section $\{h=1\}$ of any positive function $h$ homogeneous of degree one (\hb{thm:projective}). Since $\int_{\{h=1\}}\Phi^\ast\omega_{\mathbf x}=\int_{\{h\circ\Phi^{-1}=1\}}\omega_{\mathbf x}$ and $h\circ\Phi^{-1}$ is again positive and homogeneous of degree one, the two periods differ by the constant factor of the display. The moduli of the two integrands are related by the same change of variables and a positive constant, so the two sides converge together; for real data the integrands are nonnegative and the identity holds in $[0,\infty]$.

At $\boldsymbol\lambda=\mu\one$ and $c_j=c^\ast_j=\mu^{-1}$ every coordinate of $\mathbf x'$ equals that of $\mathbf x$, and the factor is $\mu^{-m_{\cE}}\mu^{\sum_j(r_js_j+r^\ast_jt_j)}=1$, since $m_{\cE}=\sum_j(r_js_j+r^\ast_jt_j)$ on the balance wall. With $\boldsymbol\lambda$ equal to $1$ off $g$ and to $\lambda$ at $g$, and $c_j=c^\ast_j=1$, the factor is $\lambda^{-m_g}$; with $\lambda=1/x^K_1$ this reads every twist at $g$ with $x^K_1=1$, up to the factor $(x^K_1)^{-m_g}$, and for one table it leaves the one number $x^\Psi_1/x^K_1$, carried by the cobasis letter. With $\boldsymbol\lambda=\one$, one $c_j$ or one $c^\ast_j$ different from $1$ and the other constants $1$, the factor is $c_j^{-r_js_j}$ or $c_j^{\ast\,-r^\ast_jt_j}$; for one table twisted in its cobasis letter by $\mu\mathbf x$ it is $\mu^{-r^\ast t}$.
\end{proof}

\begin{sremark}
The reference, \hb{def:twist}, rescales a twist at one mark, normalized by the scale of the mark, and the relation for one table twisted in its cobasis letter is used with \hb{def:hold}. The exact check of the scalar case is check~11 of S21-23\_work/checks.py.
\end{sremark}

\begin{htheorem}[\hb{thm:coreshares} (The cells ride on the core's fibre)]
As stated: at complex data, every exponent $0$ or of positive real part, on contours on which the load families converge absolutely.
\end{htheorem}

\begin{proof}
The steps are those of the reference, proof of \hb{thm:coreshares}, with the moduli carried along. By \hb{thm:substletters}, table $j$ has the letters
\[
K_j=\prod_iK_{\mathcal A_{ij}/g_i}(\alpha_{B_i})\,K_{\mathcal H_j}\big((\gamma^K_{\mathcal A_{ij}}(\alpha_{B_i}))_i\big),\qquad \Psi_j=\prod_i\Psi_{\mathcal A_{ij}\setminus g_i}(\alpha_{B_i})\,\Psi_{\mathcal H_j}\big((\gamma^\Psi_{\mathcal A_{ij}}(\alpha_{B_i}))_i\big).
\]
For $\re s_j>0$, \hb{thm:E1} opens $K_{\mathcal H_j}(X)^{-s_j}=\big(\sum_T\theta^{\mathcal H_j}_TX^T\big)^{-s_j}$ with one share $a_{j,T}$ per basis on a contour of positive real parts, and at $X_i=\gamma^K_{\mathcal A_{ij}}$ the monomials give $\prod_T(X^T)^{-a_{j,T}}=\prod_i(\gamma^K_{\mathcal A_{ij}})^{-a_j(h_i)}$; in the same way $\Psi_{\mathcal H_j}(Y)^{-t_j}=\big(\sum_T\theta^{\mathcal H_j}_TY^{\{h\}\setminus T}\big)^{-t_j}$ is opened with the shares $b_{j,T}$, and $\prod_T(Y^{\{h\}\setminus T})^{-b_{j,T}}=\prod_i(\gamma^\Psi_{\mathcal A_{ij}})^{-b_j(h_i)}$. A letter with exponent $0$ is the constant $1$ and is opened by nothing. Since $\gamma^K_{\mathcal A}=K_{\mathcal A\setminus g}/K_{\mathcal A/g}$ and $\gamma^\Psi_{\mathcal A}=\Psi_{\mathcal A/g}/\Psi_{\mathcal A\setminus g}$ (\hb{def:substitution}),
\[
K_{\mathcal A_{ij}/g_i}^{-s_j}(\gamma^K_{\mathcal A_{ij}})^{-a_j(h_i)}=K_{\mathcal A_{ij}/g_i}^{-(s_j-a_j(h_i))}K_{\mathcal A_{ij}\setminus g_i}^{-a_j(h_i)},\qquad \Psi_{\mathcal A_{ij}\setminus g_i}^{-t_j}(\gamma^\Psi_{\mathcal A_{ij}})^{-b_j(h_i)}=\Psi_{\mathcal A_{ij}\setminus g_i}^{-(t_j-b_j(h_i))}\Psi_{\mathcal A_{ij}/g_i}^{-b_j(h_i)},
\]
and $e^{-|\alpha|}=\prod_ie^{-|\alpha_{B_i}|}$, so for fixed shares the integrand of $J_F$ is the product over the cells of the integrands of the interests of the load families $F_i(\mathbf a,\mathbf b)$, times the share kernels and the powers of the weights of the cores.

It remains to exchange the share integral and the integral over the orthant. For $X>0$, $|X^{-a}|=X^{-\re a}$, and $|\alpha^{\mvec-\one}|=\alpha^{\re\mvec-\one}$; so the modulus of the joint integrand is
\[
\prod_j\big|\Gst{|\mathcal H_j|}1{\mathbf a_j}{s_j}\big|\,\big|\Gst{|\mathcal H_j|}1{\mathbf b_j}{t_j}\big|\prod_{j,T}\big(\theta^{\mathcal H_j}_T\big)^{-\re(a_{j,T}+b_{j,T})}
\]
times the product of the integrands of the interests of the load families at the real parts of their data and exponents. Integrated over the orthant, the second factor gives the product of these interests at the real parts, which are finite by the hypothesis and constant along the contour, since they depend on the shares only through their real parts. On the contour, $a_{j,T}=c_{j,T}+i\tau_{j,T}$ with $c_{j,T}>0$ fixed and $\sum_T\tau_{j,T}=\im s_j$, and the classical bound $|\Gamma(c+i\tau)|\le C_c\,e^{-\pi|\tau|/2}(1+|\tau|)^{c-1/2}$ makes $\prod_T|\Gamma(a_{j,T})|$ integrable over the imaginary parts of all shares but one; the same holds for the shares $b_{j,T}$. So the joint integrand is absolutely integrable, the integrals may be taken in either order, and the display follows. The identity $J_{F_i(\mathbf a,\mathbf b)}=\Gamma(\kappa_{B_i}-\sum_j(a_j(h_i)+b_j(h_i)))\,P_{F_i(\mathbf a,\mathbf b)}$ is \hb{thm:projective} at the imbalance of the load family: its tables $\mathcal A_{ij}\setminus g_i$ and $\mathcal A_{ij}/g_i$ have the degrees $r_{\mathcal A_{ij}}$, $|B_i|-r_{\mathcal A_{ij}}$, $r_{\mathcal A_{ij}}-1$ and $|B_i|-r_{\mathcal A_{ij}}+1$, and the masses $\mvec|_{B_i}$ exceed their degree by $\kappa_{B_i}-\sum_j(a_j(h_i)+b_j(h_i))$, with $\kappa_{B_i}$ of \hb{thm:substwalls}.
\end{proof}

\begin{htheorem}[\hb{thm:pinned} (A cell read alike is pinned)]
As stated: the hypothesis on every active letter, whatever the sign of its exponent.
\end{htheorem}

\begin{proof}
By \hb{thm:corereads}, which holds for every substituted family at all data, both sides converging together,
\[
P_F(\mvec;\mathbf s,\mathbf t)=\int\prod_i\beta_i(du_i)\ P_{F_{\mathcal H}[\boldsymbol\gamma(u)]}\big((\kappa_{B_i})_i;\mathbf s,\mathbf t\big),
\]
$\beta_i$ the cell kernels and $F_{\mathcal H}[\boldsymbol\gamma(u)]$ the core family read at $u$. A letter with exponent $0$ is the constant $1$ in every period, so the twists of the inactive letters do not enter the period of the core family read at $u$; give them the values $c_{ij}=1$, $c^\ast_{ij}=1$. Under the hypothesis the twists of the active letters at $h_i$ are $c_{ij}\gamma_i(u_i)$ in $K_j$ and $c^\ast_{ij}\gamma_i(u_i)$ in $\Psi_j$, so $F_{\mathcal H}[\boldsymbol\gamma(u)]$ has the period of $F_{\mathcal H}[\mathbf c]$ twisted further by the vector $\boldsymbol\lambda=(\gamma_i(u_i))_i$ on the marks of the core, in both letters of every table. The core family read at $u$ has the imbalance of $F$ (\hb{thm:corereads}), so it is on its balance wall, and the torus relation of \hb{def:twist} with this $\boldsymbol\lambda$ and constants $1$ gives
\[
P_{F_{\mathcal H}[\boldsymbol\gamma(u)]}\big((\kappa_{B_i})_i;\mathbf s,\mathbf t\big)=\prod_i\gamma_i(u_i)^{-\kappa_{B_i}}\,P_{F_{\mathcal H}[\mathbf c]}\big((\kappa_{B_i})_i;\mathbf s,\mathbf t\big).
\]
Inserting this and integrating cell by cell gives the display, by Tonelli for real data and by Fubini where the integrals converge absolutely. The step is the reference proof of \hb{thm:pinned}, which nowhere uses the sign of an exponent; the hypothesis there is written for the letters with positive exponent.
\end{proof}

\begin{hcorollary}[\hb{cor:fullleancore} (Read by one hinge value, the period factors through the core)]
As stated: for every substituted family with real data at every cell of which the active letters read the cell by one hinge value.
\end{hcorollary}

\begin{proof}
\emph{(1) The first display.} The hypothesis is that of \hb{thm:pinned} with all constants $c_{ij}=c^\ast_{ij}=1$, and $F_{\mathcal H}[\mathbf c]=F_{\mathcal H}$; \hb{thm:pinned} gives the first display, in $[0,\infty]$ since the data are real. At a mark $h$ left in place, the cell $\binom{\{h,g\}}1$ has both hinge values equal to $\alpha_h$ (\hb{def:substitution}), so the hypothesis holds there with $\gamma=\alpha_h$.

\emph{(2) The cells.} The cell family $F_i$ has the tables $\mathcal A_{ij}$, cells on $|B_i|+1$ marks of rank $r_{\mathcal A_{ij}}$, and by \hb{thm:substwalls}
\[
\kappa_{B_i}=m_{B_i}-\sum_j\big(s_j(r_{\mathcal A_{ij}}-1)+t_j(|B_i|-r_{\mathcal A_{ij}})\big),
\]
so $m_{B_i}+2\sigma-\kappa_{B_i}=\sum_j\big(s_jr_{\mathcal A_{ij}}+t_j(|B_i|+1-r_{\mathcal A_{ij}})\big)$, the degree of the family word of $F_i$: at the masses $(\mvec|_{B_i},2\sigma-\kappa_{B_i})$ it is on its balance wall. Read its period on the section $\{|\alpha_{B_i}|=1\}$ (\hb{thm:projective}), with $\alpha_{B_i}=u\in\Delta_{B_i}$ and $\alpha_{g_i}=v>0$. By \hb{def:substitution}, $K_{\mathcal A_{ij}}=K_{\mathcal A_{ij}/g_i}(u)\big(\gamma^K_{\mathcal A_{ij}}(u)+v\big)$ and $\Psi_{\mathcal A_{ij}}=\Psi_{\mathcal A_{ij}\setminus g_i}(u)\big(v+\gamma^\Psi_{\mathcal A_{ij}}(u)\big)$. Under the hypothesis every active letter carries the factor $(v+\gamma_i(u))$ to the power of its exponent, and an inactive letter has exponent $0$; so the word of $F_i$ is the density of the cell kernel without $u^{\mvec|_{B_i}-\one}$ times $(v+\gamma_i(u))^{-2\sigma}$, $2\sigma=\sum_j(s_j+t_j)$. For $0<\kappa<2\sigma$ and $\gamma>0$, the substitution $v=\gamma w$ and the Beta integral give
\[
\int_0^\infty v^{2\sigma-\kappa-1}(v+\gamma)^{-2\sigma}\,dv=\gamma^{-\kappa}\int_0^\infty w^{2\sigma-\kappa-1}(1+w)^{-2\sigma}\,dw=\Gst21{\kappa,\ 2\sigma-\kappa}{2\sigma}\gamma^{-\kappa},
\]
the Beta corner of a linear letter (\hb{prop:scalebreta} with one coefficient). Integrating $v$ first (Tonelli),
\[
P_{F_i}\big((\mvec|_{B_i},2\sigma-\kappa_{B_i});\mathbf s,\mathbf t\big)=\Gst21{\kappa_{B_i},\ 2\sigma-\kappa_{B_i}}{2\sigma}\int_{\Delta_{B_i}}\gamma_i^{-\kappa_{B_i}}\,d\beta_i ,
\]
which identifies the pinned shape integral.

\emph{(3) The walls.} The marks of a core are neither loops nor coloops (\hb{def:substitution}), so a core has at least two marks, $q\ge2$, and $B_i$ is a nonempty proper set of marks; it has width one in every table (\hb{prop:substitution}(i)). Width one in table $j$ means $d_j(B_i)+d_j(\cE\setminus B_i)=r_j-1$ and $d^\ast_j(B_i)+d^\ast_j(\cE\setminus B_i)=r^\ast_j-1$, so
\[
\kappa_{B_i}+\kappa_{\cE\setminus B_i}=m_{\cE}-\sum_j\big(s_j(r_j-1)+t_j(r^\ast_j-1)\big)=\kappa_{\cE}+2\sigma ,
\]
which is $2\sigma$ on the balance wall (\hb{lem:twosumpieces}; check~2 of S21-23\_work/checks.py). If every proper wall is positive, $\kappa_{B_i}>0$ and $\kappa_{\cE\setminus B_i}>0$, so $0<\kappa_{B_i}<2\sigma$.

\emph{(4) The second display.} Multiply the first display by $\prod_i\Gst21{\kappa_{B_i},\ 2\sigma-\kappa_{B_i}}{2\sigma}\in(0,\infty)$ and use step~(2). The left side is finite (\hb{thm:projective}), every factor on the right lies in $(0,\infty]$, so every factor is finite. The core family at $(\kappa_{B_i})_i$ has the imbalance of $F$ (\hb{thm:corereads}), the cell families are balanced by step~(2), and a family on its balance wall with a finite period has every proper wall positive (\hb{thm:projective}).

\emph{(5) One cell.} With one cell at $h_1$ and the other marks of the core left in place, the left-in-place cell family at $h$, the list $\binom{\{h,g\}}1$ in every table with the masses $(m_h,2\sigma-m_h)$, has the word $(\alpha_h+\alpha_g)^{-2\sigma}$ and the period $\Gst21{m_h,\ 2\sigma-m_h}{2\sigma}$, while $\kappa_{\{h\}}=m_h$; these factors cancel the Beta corners of the marks left in place. The core family is then the family of the cells of the room at the masses $(\kappa_B,\mvec|_R)$ and $F_1$ that of the cells of the bunch (\hb{prop:substitution}(iii)), and the second display is the identity of \hb{thm:cornerfactor}. In the first display the marks left in place contribute $1$, their simplex being a point and their kernel the constant $1$, and it is the identity of \hb{prop:readingsagree} with all constants $1$.
\end{proof}

\begin{sremark}
The reference proof (\hb{cor:fullleancore}, one table at $(0,t)$) uses of the full lean only that the letter which sees the cells sees each through one function. A family at mixed full leans is covered: on $B=\{1,2\}$ let table~1 at $(s_1,0)$ carry the series block $\binom{\{1,2,g\}}2$ and table~2 at $(0,t_2)$ the parallel block $\binom{\{1,2,g\}}1$. Both active letters read $B$ through $e_2/e_1$ (\hb{cor:blocksclosures}(i)), $\kappa_B=m_B-2\sigma$, the cell family has the word $e_2(\alpha_1,\alpha_2,\alpha_g)^{-2\sigma}$ and the period $\Gamma(2\sigma-m_1)\Gamma(2\sigma-m_2)\Gamma(\kappa_B)/\Gamma(2\sigma)$ (\hb{prop:spleans}), and the pinned shape integral is $\int_0^1u^{m_1-1}(1-u)^{m_2-1}(u(1-u))^{-\kappa_B}du=\Gamma(2\sigma-m_1)\Gamma(2\sigma-m_2)/\Gamma(4\sigma-m_B)$, whose quotient is $\Gst21{\kappa_B,\ 2\sigma-\kappa_B}{2\sigma}$, as step~(2) says. At $(2\sigma,m_1,m_2)=(3,2,2)$ the period is $\frac12$, and the identity holds exactly there and at $(\frac72,\frac32,\frac52)$ (checks~7a to~7d of S21-23\_work/checks.py).
\end{sremark}

\begin{hcorollary}[\hb{cor:mot-loadpinned} (The pinned shape integral is the period of a load family)]
As stated.
\end{hcorollary}

\begin{proof}
A load $a_{ij}\neq0$ occurs only where $s_j\neq0$, where $\gamma^K_{\mathcal A_{ij}}=K_{\mathcal A_{ij}\setminus g_i}/K_{\mathcal A_{ij}/g_i}=\gamma_i$, and a load $b_{ij}\neq0$ only where $t_j\neq0$, where $\gamma^\Psi_{\mathcal A_{ij}}=\Psi_{\mathcal A_{ij}/g_i}/\Psi_{\mathcal A_{ij}\setminus g_i}=\gamma_i$ (\hb{def:substitution}). Since $\sum_j(a_{ij}+b_{ij})=\kappa_{B_i}$,
\[
\gamma_i^{-\kappa_{B_i}}\prod_jK_{\mathcal A_{ij}/g_i}^{-s_j}\Psi_{\mathcal A_{ij}\setminus g_i}^{-t_j}=\prod_jK_{\mathcal A_{ij}\setminus g_i}^{-a_{ij}}\Psi_{\mathcal A_{ij}\setminus g_i}^{-(t_j-b_{ij})}K_{\mathcal A_{ij}/g_i}^{-(s_j-a_{ij})}\Psi_{\mathcal A_{ij}/g_i}^{-b_{ij}},
\]
the word of the load family, whose imbalance $\kappa_{B_i}-\sum_j(a_{ij}+b_{ij})$ is zero (\hb{thm:coreshares}). Its period is the integral of $u^{\mvec|_{B_i}-\one}$ times its word over the shape simplex (\hb{thm:projective}), the pinned shape integral, and the product formula is \hb{cor:fullleancore}. If every proper wall of $F$ is positive, then $0<\kappa_{B_i}<2\sigma=\sum_j(s_j+t_j)$ (\hb{cor:fullleancore}); with $s_j,t_j\ge0$, raising the loads one after another from $0$, each up to its bound $s_j$ or $t_j$ and only on the active letters, reaches the total $\kappa_{B_i}$, and at integer data $\kappa_{B_i}$ is an integer (\hb{thm:substwalls}) and the loads can be integers. The exponents $a_{ij}$, $t_j-b_{ij}$, $s_j-a_{ij}$ and $b_{ij}$ are then nonnegative. When $P_F$ is finite every factor of the product formula is finite, the factors being positive, and the load family, on its balance wall, has every proper wall positive (\hb{thm:projective}).
\end{proof}

\begin{htheorem}[\hb{thm:chainoperator} (A chain is the power of one operator)]
As stated: at real exponents, the pinned case for the active letters.
\end{htheorem}

\begin{proof}
The reference proof of \hb{thm:chainoperator} applies verbatim at real exponents. It uses pooling (\hb{thm:pooling}, for every family at all data, both sides converging together), the nonnegativity of the integrands (at real exponents the letters are positive, and so is every power of them), the scalar twist of a mark (\hb{def:twist}) and the Beta integral $\int_0^1v^{m_{g_{\mathrm{in}}}-1}(1-v)^{m_{g_{\mathrm{out}}}-1}(1-v+\lambda v)^{-2\sigma}dv=\lambda^{-m_{g_{\mathrm{in}}}}\Gst21{m_{g_{\mathrm{in}}},\ m_{g_{\mathrm{out}}}}{2\sigma}$ at $m_{g_{\mathrm{in}}}+m_{g_{\mathrm{out}}}=2\sigma$; none of these uses the sign of an exponent. In the pinned case the reference assumes that all the active letters read the cell through one function; a letter of exponent $0$ is absent from every period, which is the hypothesis as stated.
\end{proof}

\begin{htheorem}[\hb{thm:lettertransfer} (Letters by transfer)]
As stated: the reduction of the states to linearly independent ones, and the bound on the interface dimension for a multivector of any tensor rank.
\end{htheorem}

\begin{proof}
The products of transfer matrices are the reference's. \emph{The stations.} The stations of level $k$ of $S_l$ give a level factorization with as many terms (the children factorization, \hb{def:levelweights}), so their number is at least $\operatorname{rank}\Theta^{(k)}_{S_l}$, and the number of states is at least the interface dimension.

\emph{The reduction.} The child of a pattern $p$ of $S_l$ at level $k$ is $\sum_X\theta_{p\sqcup X}\alpha^X$, the polynomial of the row $p$ of $\Theta^{(k)}_{S_l}$, so the children of level $k$ span a space of dimension $\operatorname{rank}\Theta^{(k)}_{S_l}$. Splitting $X$ along the step $B_{l-1}$, the child at $p$ of $S_{l-1}$ is
\[
\sum_X\theta_{p\sqcup X}\alpha^X=\sum_{p'\subseteq B_{l-1}}\alpha^{p'}\sum_{X'}\theta_{(p\sqcup p')\sqcup X'}\alpha^{X'},
\]
the sum over $p'$ of $\alpha^{p'}$ times the child at $p\sqcup p'$ of $S_l$. Choose at every $S_l$ and every level a basis of the span of the children and write every child in it; the display expresses each basis element at $S_{l-1}$, a combination of children of $S_{l-1}$, through the basis elements at $S_l$, with coefficients polynomial in $\alpha_{B_{l-1}}$. These are transfer matrices on $\sum_k\operatorname{rank}\Theta^{(k)}_{S_l}$ states at $S_l$, and the products give $K$; the same coefficients, read on the complements, give $\Psi$ (\hb{def:levelweights}).

\emph{The bound.} By \hb{thm:interfacerank}(i), $\operatorname{rank}\Theta^{(k)}_{S_l}\le\operatorname{ir}_k(\boldsymbol\omega)^2$, and $\le\frac12\operatorname{ir}_k(\boldsymbol\omega)(\operatorname{ir}_k(\boldsymbol\omega)+1)$ when $\boldsymbol\omega$ is real; by (ii), $\operatorname{ir}_k(\boldsymbol\omega)\le\sum_i\binom{c_i}{k-d_i}$. The functions $q\mapsto q^2$ and $q\mapsto\frac12q(q+1)$ increase on $[0,\infty)$, and summing over $k$ gives the bound. For a determinantal table, of tensor rank one with $c_1=c(S_l)$ and $d_1=d(S_l)$ (\hb{thm:interfacerank}(ii)), the bound is $\sum_m\binom{c}{m}^2=\binom{2c}{c}$, $c=c(S_l)$, counting the $c$-subsets of a set of $2c$ elements by their intersection with a fixed half.
\end{proof}

\begin{htheorem}[\hb{thm:momenttransfer} (Integer joint moments by fixed matrices)]
As stated: all the letters of a family at once, along a chain of bunches with any steps and any independent measures on them.
\end{htheorem}

\begin{proof}
This is the reference proof of \hb{thm:momenttransfer} with the tensor product taken over all the tables and with one matrix per step. A product of scalars $\mathbf u_i^{\mathsf T}A_{i,0}\cdots A_{i,L-1}\mathbf v_i$, $i=1,\dots,n$, equals $(\bigotimes_i\mathbf u_i)^{\mathsf T}\big(\bigotimes_iA_{i,0}\big)\cdots\big(\bigotimes_iA_{i,L-1}\big)(\bigotimes_i\mathbf v_i)$ by the mixed-product rule $(\bigotimes_iA_i)(\bigotimes_iB_i)=\bigotimes_i(A_iB_i)$. Apply it to the $a_j$ copies of $K_j$ and the $b_j$ copies of $\Psi_j$ for every $j$: the factor of the step $l$ is $\bigotimes_j\big(\mathbf M^K_{j,l}(\alpha_{B_l})^{\otimes a_j}\otimes\mathbf M^\Psi_{j,l}(\alpha_{B_l})^{\otimes b_j}\big)$, a matrix function of $\alpha_{B_l}$ alone. Every entry of the product is a finite sum of products of one entry of each step factor; the measures of the steps are independent, so the integral of such a product is the product of the integrals, by Tonelli for nonnegative entries and by Fubini under the stated finiteness, and integrating entrywise replaces the factor of the step $l$ by $\mathbf A_l$.

For identical steps every $\mathbf A_l$ is one matrix $\mathbf A$ of the stated size, and the scalars $\mathbf u^{\mathsf T}\mathbf A^L\mathbf v$ satisfy the linear recurrence whose coefficients are those of the characteristic polynomial of $\mathbf A$, which annihilates $\mathbf A$. With independent laws $\operatorname{Gamma}_{m_x}$ on the marks, $\prod_jK_j^{a_j}\Psi_j^{b_j}$ is homogeneous of degree $d=\sum_j(a_jr_j+b_jr^\ast_j)$, and the draw splits into an independent total of law $\operatorname{Gamma}_{m_{\cE}}$ and a shape of law $\Dir_{\mvec}$ (\hb{prop:gammadraws}(ii)), the total having the moment $(m_{\cE})_d$ of order $d$. With $\Dir_{\mvec|_{B_l}}$ on a step, the entries of $\mathbf A_l$ are the Dirichlet moments of the monomials (\hb{thm:dirmoments}).
\end{proof}

\begin{htheorem}[\hb{thm:anywidth} (Gluing across a bunch of any width)]
As stated: the share display for one table at a full lean at every width, and the count $2N-1$ of the free shares at width one.
\end{htheorem}

\begin{proof}
The main display is the reference's. For one table at the full lean $(0,t)$ in level form, one station per level, the active letter $\Psi$ has the $c(B)+1$ terms of the levels $k=d(B),\dots,r(B)$, so the stacked Gamma of the display is $\prod_k\Gamma(b_k)/\Gamma(t)$ and the pins are $\sum_kb_k=t$ and $\sum_k(r(B)-k)b_k=\kappa_B$, the main pin without the shares of $K$. The fibre integral (\hb{def:fibre}) has $c(B)+1$ shares and two pins. Integrating the delta of the pins against the imaginary parts of $b_{d(B)}$ and $b_{r(B)}$ divides by $|\det|$ of the matrix of the two pins in these shares,
\[
\begin{pmatrix}1&1\\c(B)&0\end{pmatrix},\qquad|\det|=c(B)
\]
(check~8a of S21-23\_work/checks.py), and leaves $(2\pi)^{2-(c(B)+1)}$ times the integral over the imaginary parts of the $c(B)-1$ free shares, which is $\frac1{(2\pi i)^{c(B)-1}}\int\prod db_k$. The pins solved for $b_{d(B)}$ and $b_{r(B)}$ give the stated values; at $c(B)=2$ they are $\frac{\kappa_B-z}2$ and $\frac{2t-\kappa_B-z}2$ with $b_{d(B)+1}=z$ (check~8b). At $c(B)=1$ the fibre has dimension $0$, the matrix of the pins has $|\det|=1$, and the integral is the value at $(b_{d(B)},b_{r(B)})=(\kappa_B,t-\kappa_B)$, the Beta corner times the two inner periods, which is \hb{cor:pinnedproduct}: the bunch and room letters of the two levels are the minors $\Psi_{\cB/R},\Psi_{\cB|B}$ and $\Psi_{\cB|R},\Psi_{\cB/B}$ up to constant factors of the factorization that cancel in the product. At width one with factoring weights each level has one station (\hb{thm:levels}), so each of the $2N$ active letters has two terms, and the free shares are $4N-2N-1=2N-1$, the dimension of the fibre integral of \hb{thm:twosides}, of which the main display is then the formula.
\end{proof}

\begin{sremark}
The paragraph after \hb{thm:anywidth} gives the Mellin transform of the scale integral $I(\boldsymbol\zeta)=\int_0^\infty x^{\kappa_B-1}\big(1+\sum_{i=1}^{c-1}\zeta_ix^i+x^c\big)^{-t}dx$, $c=c(B)$. By Tonelli and the substitutions $\zeta_i=\eta_ix^{-i}$, then $x^c=w$,
\[
\int\prod_i\zeta_i^{z_i-1}I(\boldsymbol\zeta)\,d\boldsymbol\zeta=\frac{\prod_i\Gamma(z_i)\,\Gamma(t-\sum_iz_i)}{\Gamma(t)}\int_0^\infty x^{\kappa_B-\sum_iiz_i-1}(1+x^c)^{\sum_iz_i-t}dx=\frac{\prod_{k=d(B)}^{r(B)}\Gamma(b_k)}{c\,\Gamma(t)}
\]
at $b_{r(B)-i}=z_i$, $b_{d(B)}=(\kappa_B-\sum_iiz_i)/c$ and $b_{r(B)}=t-b_{d(B)}-\sum_iz_i$, wherever these have positive real parts: the share kernel of the theorem with the factor $1/c$ of the pin.
\end{sremark}

\section{Chapter~\hb{ch:hinged}: the law of a cell}\label{sec:ch22}

\begin{hdefinition}[\hb{def:hinged} (The law of a cell family)]
As stated: the Mellin transform of the law of a cell family, and the transforms and the hinge measures of a cell family, those of one cell being their readings.
\end{hdefinition}

\begin{proof}
The definition asserts three things. \emph{The Mellin transform.} For complex $\mathbf z$ and real positive coordinates, $|\boldsymbol\varrho_F^{-\mathbf z}|=\boldsymbol\varrho_F^{-\re\mathbf z}$, so the integral converges absolutely where it is finite at the real part, with modulus at most that value; for one cell $\boldsymbol\varrho_F^{-\mathbf z}=(\gamma^\Psi_{\mathcal A}/\gamma^K_{\mathcal A})^{-z_{1,\Psi}}=\varrho_{\mathcal A}^{-z}$, and the second form is the image of the measure. The base form of one cell is $\phi_F$ with $(\gamma^K_{\mathcal A})^{-\kappa}=K_{\mathcal A/g}^{\kappa}K_{\mathcal A\setminus g}^{-\kappa}$.

\emph{Homogeneity.} Substitute $\alpha=c\beta$ in the orthant forms: $\alpha^{\mvec-\one}d\alpha=c^{m_B}\beta^{\mvec-\one}d\beta$, the letters of the deletions have the degrees $r_{\mathcal A_j}$ and $|B|-r_{\mathcal A_j}$, those of the contractions $r_{\mathcal A_j}-1$ and $|B|-r_{\mathcal A_j}+1$, and the hinge values have degree one. So
\[
f_F(\mathbf x,\mathbf y)=c^{m_B-\sum_j(s_jr_{\mathcal A_j}+t_j(|B|-r_{\mathcal A_j}))}f_F(\mathbf x/c,c\mathbf y)=c^{\kappa_F-|\mathbf s|}f_F(\mathbf x/c,c\mathbf y),
\]
and in the same way $\check f_F(\mathbf x,\mathbf y)=c^{\kappa_F-|\mathbf t|}\check f_F(\mathbf x/c,c\mathbf y)$. With $c=x_1$ each transform is fixed by its values at $x_1=1$.

\emph{One cell.} At $N=1$, $f_F(1,y)$ and $\check f_F(1,y)$ are the integrals called the orthant forms of the transforms of one cell, and the images of $e^{-\lambda}\mathfrak r_F$ and $e^{-\lambda}\check{\mathfrak r}_F$ under $(\lambda,\mu)\mapsto\mu$ are the images of $\alpha^{\mvec-\one}K_{\mathcal A\setminus g}^{-s}\Psi_{\mathcal A\setminus g}^{-t}e^{-1/\gamma^K_{\mathcal A}}d\alpha$ under $\gamma^\Psi_{\mathcal A}$ and of $\alpha^{\mvec-\one}K_{\mathcal A/g}^{-s}\Psi_{\mathcal A/g}^{-t}e^{-1/\gamma^\Psi_{\mathcal A}}d\alpha$ under $\gamma^K_{\mathcal A}$.
\end{proof}

\begin{htheorem}[\hb{thm:seriesclosure} (The series closure at positive exponents)]
As stated: the closure through the transforms of the cell families, each transform one Gramma of index $|\mathbf s|-\kappa_{B_i}$ at every data, and its Breta shape reading wherever the family of the hinge minors is balanced.
\end{htheorem}

\begin{proof}
\emph{The closure.} The responses of the reference, proof and statement of \hb{thm:seriesclosure}, are the integrals over $\mathbb R^{B_i}_{>0}$ of $\alpha^{\mvec|_{B_i}-\one}\prod_jK_{\mathcal A_{ij}\setminus g_i}^{-s_j}\Psi_{\mathcal A_{ij}\setminus g_i}^{-t_j}$ against $\exp\big(-\sum_jx_jK_{\mathcal A_{ij}/g_i}/K_{\mathcal A_{ij}\setminus g_i}-\sum_jy_j\Psi_{\mathcal A_{ij}/g_i}/\Psi_{\mathcal A_{ij}\setminus g_i}\big)$ with $y_1=1$; since $K_{\mathcal A/g}/K_{\mathcal A\setminus g}=1/\gamma^K_{\mathcal A}$ and $\Psi_{\mathcal A/g}/\Psi_{\mathcal A\setminus g}=\gamma^\Psi_{\mathcal A}$, they are the transforms $f_{F_i}(\mathbf x,(1,y_2,\dots,y_N))$ of \hb{def:hinged}, and the first display is the reference's. The statement on repeated cells is read off the product.

\emph{One table.} By the homogeneity of \hb{def:hinged} with $c=x$, $f_{F_i}(x,1)=x^{\kappa_{B_i}-s}f_{\mathcal A_i}(x)$. The core $\binom{\{h_1,\dots,h_q\}}{q-1}$ has rank $q-1$ and corank $1$, and the core family read at $u$ is balanced (\hb{thm:corereads}), so $\sum_i\kappa_{B_i}=(q-1)s+t$. The power of $x$ is then $x^{s-1+(q-1)s+t-qs}=x^{t-1}$, and the closure is $\frac1{\Gamma(s)\Gamma(t)}\int_0^\infty x^{t-1}\prod_if_{\mathcal A_i}(x)\,dx$.

\emph{The Gramma.} Fix $i$ and write $\alpha=\rho u$, $u\in\Delta_{B_i}$, $\rho>0$, so that $\alpha^{\mvec|_{B_i}-\one}d\alpha=\rho^{m_{B_i}-1}u^{\mvec|_{B_i}-\one}d\rho\,d\omega$ (\hb{lem:scaleshape}). The marks $g_i$ being neither loops nor coloops, $K_{\mathcal A_{ij}\setminus g_i}$ is homogeneous of degree $r_{\mathcal A_{ij}}$ and $\Psi_{\mathcal A_{ij}\setminus g_i}$ of degree $|B_i|-r_{\mathcal A_{ij}}$, and the hinge values of degree one. So the power of $\rho$ is $\rho^{\nu-1}$ with
\[
\nu=m_{B_i}-\sum_j\big(s_jr_{\mathcal A_{ij}}+t_j(|B_i|-r_{\mathcal A_{ij}})\big)=\kappa_{B_i}-|\mathbf s|,
\]
and the exponential is $\exp(-a(u)/\rho-b(u)\rho)$ with $a=\sum_jx_j/\gamma^K_{\mathcal A_{ij}}$ and $b=\sum_jy_j\gamma^\Psi_{\mathcal A_{ij}}$. By \hb{def:gramma} and the reflection (\hb{prop:grammareflect}),
\[
\int_0^\infty\rho^{\nu-1}e^{-b\rho-a/\rho}\,d\rho=\sG_{\nu;(1,-1)}(b,a)=\sG_{-\nu;(1,-1)}(a,b),
\]
and Tonelli gives the display of the transform.

\emph{Index $0$.} The family of the hinge minors has the degree $\sum_j\big(s_jr_{\mathcal A_{ij}}+t_j(|B_i|-r_{\mathcal A_{ij}})\big)$ at the masses $\mvec|_{B_i}$, so it is on its balance wall exactly when $\kappa_{B_i}=|\mathbf s|$; at the centre of one table, every mass $\sigma$ and $s=t=\sigma$, $\kappa_{B_i}=\sigma|B_i|-\sigma(r-1)-\sigma(|B_i|-r)=s$. Torus homogeneity (\hb{thm:grammalaws}(iii)) with $\lambda=a$ gives $\sG_{0;(1,-1)}(a,b)=\sG_{0;(1,-1)}(1,ab)$. With every proper wall positive, the form of the family of the hinge minors divided by its period is its Breta shape law (\hb{def:shapelaw}), which gives the expectation; for one table, $ab=xy\,\gamma^\Psi_{\mathcal A_i}/\gamma^K_{\mathcal A_i}=xy\varrho_{\mathcal A_i}$. The period of $k$ copies is the first display with all $F_i$ equal.
\end{proof}

\begin{htheorem}[\hb{thm:node} (The law of a substitution)]
As stated: (i) and (ii) for families at real data; the present (i) was the gluing of a family, and (ii) for one table at $s,t>0$.
\end{htheorem}

\begin{proof}
\emph{(ii)} This is the reference proof of \hb{thm:node}(ii) with the index $j$, the steps being those of \hb{lem:absorb} applied to one cell family after another. Write $B=\cE_{\mathcal H}\sqcup\bigsqcup_iB_i$ and $\mathcal A_j=\mathcal H_j\langle\mathcal A_{1j},\dots,\mathcal A_{kj}\rangle$, a cell with the hinge $g_0$ (\hb{lem:nodesum}). For a Borel $H\ge0$, $\int H\,d((\boldsymbol\varrho_F)_\ast\phi_F)$ is the integral over the section $\{|\alpha|=1\}$ of the orthant form
\[
\omega=\alpha^{\mvec-\one}H(\boldsymbol\varrho_F)\,K_{\mathcal A_1/g_0}^{\kappa_F-s_1}K_{\mathcal A_1\setminus g_0}^{-\kappa_F}\prod_{j\ge2}K_{\mathcal A_j/g_0}^{-s_j}\prod_j\Psi_{\mathcal A_j\setminus g_0}^{-t_j}\,d\alpha ,
\]
homogeneous of degree $m_B-\sum_j\big(s_j(r_{\mathcal A_j}-1)+t_j(|B|-r_{\mathcal A_j})\big)-\kappa_F=0$, the base form being $u^{\mvec-\one}\prod_jK_{\mathcal A_j/g_0}^{-s_j}\Psi_{\mathcal A_j\setminus g_0}^{-t_j}(K_{\mathcal A_1\setminus g_0}/K_{\mathcal A_1/g_0})^{-\kappa_F}d\omega$. By \hb{lem:nodesum}, with the hinge values of the cells at their own marks,
\[
K_{\mathcal A_j\setminus g_0}=\prod_iK_{\mathcal A_{ij}/g_i}\,K_{\mathcal H_j\setminus g_0}\big(\alpha_{\cE_{\mathcal H}},(\gamma^K_{\mathcal A_{ij}})_i\big),\qquad K_{\mathcal A_j/g_0}=\prod_iK_{\mathcal A_{ij}/g_i}\,K_{\mathcal H_j/g_0}\big(\alpha_{\cE_{\mathcal H}},(\gamma^K_{\mathcal A_{ij}})_i\big),
\]
and the same for $\Psi$ with $\prod_i\Psi_{\mathcal A_{ij}\setminus g_i}$ and the hinge values $\gamma^\Psi_{\mathcal A_{ij}}$. The products over the cells cancel in the quotients, so the hinge values of $\mathcal A_j$ at $g_0$ are those of $\mathcal H_j$ read at $(\alpha_{\cE_{\mathcal H}},(\gamma^K_{\mathcal A_{ij}})_i)$ and $(\alpha_{\cE_{\mathcal H}},(\gamma^\Psi_{\mathcal A_{ij}})_i)$.

Apply \hb{lem:absorb} to $F_1$ with $Z=\cE_{\mathcal H}\sqcup\bigsqcup_{i\ge2}B_i$. Assign to the letter $K_{\mathcal A_{1j}}$ the functions $\prod_{i\ge2}K_{\mathcal A_{ij}/g_i}\cdot K_{\mathcal H_j/g_0}$ and $\prod_{i\ge2}K_{\mathcal A_{ij}/g_i}\cdot K_{\mathcal H_j\setminus g_0}$, read at the hinge value of the cell at $g_1$, with the exponents $s_j$ and $0$ for $j\ge2$ and $s_1-\kappa_F$ and $\kappa_F$ for $j=1$; and to the letter $\Psi_{\mathcal A_{1j}}$ the functions $\prod_{i\ge2}\Psi_{\mathcal A_{ij}\setminus g_i}\cdot\Psi_{\mathcal H_j\setminus g_0}$ and $\prod_{i\ge2}\Psi_{\mathcal A_{ij}\setminus g_i}\cdot\Psi_{\mathcal H_j/g_0}$ with the exponents $t_j$ and $0$. The exponents assigned to each letter add to its exponent, the letters of the hinge minors of $F_1$ collect to $K_{\mathcal A_{1j}/g_1}^{-s_j}\Psi_{\mathcal A_{1j}\setminus g_1}^{-t_j}$, the product of the assigned powers is the rest of $\omega$, and $H(\boldsymbol\varrho_F)$ is a Borel function of the values of the assigned functions, which give the hinge values of every $\mathcal A_j$ as quotients. The lemma collapses $B_1$ to $g_1$ with the mass $\kappa_{F_1}$ and replaces the hinge values of the cells at $g_1$ by $x_{1,j,K}\alpha_{g_1}$ and $x_{1,j,\Psi}\alpha_{g_1}$, the twist at $g_1$ by the point $\mathbf x_1$ of the law of $F_1$. Repeat for $F_2,\dots,F_k$, moving the integrals over the laws outside (Tonelli). What remains is the integral over a section of the form of the cell family $F_{\mathcal H}[\mathbf x]$, its hinge minors at $g_0$ read with the twists, its masses $\mvec|_{\cE_{\mathcal H}}$ and $\kappa_{F_i}$, the power $\kappa_F$ in place of its hinge wall, and $H$ of its ratio vector inserted. Its hinge wall is
\[
m_{\cE_{\mathcal H}}+\sum_i\kappa_{F_i}-\sum_j\big(s_j(r_{\mathcal H_j}-1)+t_j(|\cE_{\mathcal H}|+k-r_{\mathcal H_j})\big),
\]
and it equals $\kappa_F$: insert $\kappa_{F_i}=m_{B_i}-\sum_j(s_j(r_{\mathcal A_{ij}}-1)+t_j(|B_i|-r_{\mathcal A_{ij}}))$, $r_{\mathcal A_j}=r_{\mathcal H_j}+\sum_i(r_{\mathcal A_{ij}}-1)$ and $|B|=|\cE_{\mathcal H}|+\sum_i|B_i|$ (\hb{lem:nodesum}); the coefficients of $s_j$ and of $t_j$ agree term by term (check~1 of S21-23\_work/checks.py). So the remaining integral is $\int H\,d((\boldsymbol\varrho_{F_{\mathcal H}[\mathbf x]})_\ast\phi_{F_{\mathcal H}[\mathbf x]})$, and the mixture follows; the hinge wall does not depend on $\mathbf x$. With $H=\boldsymbol\varrho^{-\mathbf z}$, a nonnegative Borel function for real $\mathbf z$, the Mellin transforms mix in the same way.

\emph{(i)} The same steps with the family word of $F$ in place of the base form: on the balance wall the form of $F$ is homogeneous of degree zero, its letters are $\prod_iK_{\mathcal A_{ij}/g_i}\cdot K_{\mathcal H_j}$ and $\prod_i\Psi_{\mathcal A_{ij}\setminus g_i}\cdot\Psi_{\mathcal H_j}$ read at the hinge values (\hb{lem:nodesum}), the factors $K_{\mathcal H_j}$ and $\Psi_{\mathcal H_j}$ are assigned to the letters of each cell family in turn with the exponents $s_j$ and $t_j$, and \hb{lem:absorb} collapses every cell to its marker with the mass $\kappa_{F_i}$ and the twist by the point of its law. This is the reference's remark on families (rem:familiesnodes): \hb{lem:absorb} holds at real exponents, its proof collecting degrees and using Tonelli. For the core $\binom{\{g_1,g_2\}}1$, twisted at $g_1$ by $\mathbf x_1$ and at $g_2$ by $\mathbf x_2$, table $j$ has $K_j=x_{1,j,K}\alpha_{g_1}+x_{2,j,K}\alpha_{g_2}$ and $\Psi_j=x_{1,j,\Psi}\alpha_{g_1}+x_{2,j,\Psi}\alpha_{g_2}$, and on the section $\alpha_{g_2}=1$ its period is the displayed integral, a Breta of one bunch (\hb{prop:scalebreta}); the substitution $v\mapsto1/v$ multiplies the integrand by $v^{2\sigma}$ and turns $v^{\kappa_{F_1}-1}dv$ into $v^{-\kappa_{F_1}-1}dv$, which gives the swap since $\kappa_{F_1}+\kappa_{F_2}=2\sigma$.
\end{proof}

\begin{hproposition}[\hb{prop:faces} (Faces of one law)]
As stated, for every cell family with plain tables and $s_j,t_j>0$.
\end{hproposition}

\begin{proof}
Write $\kappa=\kappa_F$ and $\gamma^K_j$, $\gamma^\Psi_j$ for the hinge values of $\mathcal A_j$. The reference proof of \hb{prop:faces} is the case $N=1$; the steps below are that proof with the index $j$, the changes being the bookkeeping of the powers of $\gamma^K_1$ and $\gamma^\Psi_1$, which pass the integrals from the simplex to the law.

\emph{(i)} In the orthant form of $f_F$ write $\alpha=\rho u$ (\hb{lem:scaleshape}). As in the proof of \hb{thm:seriesclosure} in this part, the power of $\rho$ is $\rho^{\kappa-|\mathbf s|-1}$ and
\[
f_F(\mathbf x,\mathbf y)=\int_{\Delta_B}\sG_{|\mathbf s|-\kappa;(1,-1)}\Big(\sum_j\frac{x_j}{\gamma^K_j(u)},\ \sum_jy_j\gamma^\Psi_j(u)\Big)\,u^{\mvec-\one}\prod_jK_{\mathcal A_j\setminus g}^{-s_j}\Psi_{\mathcal A_j\setminus g}^{-t_j}\,d\omega(u).
\]
Since $K_{\mathcal A_j\setminus g}=\gamma^K_jK_{\mathcal A_j/g}$, the measure is $(\gamma^K_1)^{\kappa}\prod_j(\gamma^K_j)^{-s_j}\phi_F=(\gamma^K_1)^{\kappa-|\mathbf s|}\prod_j(\gamma^K_j/\gamma^K_1)^{-s_j}\phi_F$. Torus homogeneity (\hb{thm:grammalaws}(iii)) with $\lambda=1/\gamma^K_1$, $\sG_{\nu;(1,-1)}(\lambda\xi_1,\xi_2/\lambda)=\lambda^{-\nu}\sG_{\nu;(1,-1)}(\xi_1,\xi_2)$, gives, at $\mathbf z=\boldsymbol\varrho_F(u)$,
\[
\sG_{|\mathbf s|-\kappa;(1,-1)}\Big(\sum_j\frac{x_j}{\gamma^K_j},\ \sum_jy_j\gamma^\Psi_j\Big)=(\gamma^K_1)^{|\mathbf s|-\kappa}\,\sG_{|\mathbf s|-\kappa;(1,-1)}\Big(\sum_j\frac{x_j}{z_{j,K}},\ \sum_jy_jz_{j,\Psi}\Big).
\]
The powers of $\gamma^K_1$ cancel, and the image under $\boldsymbol\varrho_F$ gives the first law form. For $\check f_F$ the contracted letters have the degrees $r_{\mathcal A_j}-1$ and $|B|-r_{\mathcal A_j}+1$, so the power of $\rho$ is $\rho^{\kappa-|\mathbf t|-1}$, the index is $|\mathbf t|-\kappa$, and the coefficients are $\sum_jx_j/\gamma^\Psi_j$ and $\sum_jy_j\gamma^K_j$; the measure is $u^{\mvec-\one}\prod_jK_{\mathcal A_j/g}^{-s_j}\Psi_{\mathcal A_j/g}^{-t_j}d\omega=(\gamma^K_1)^{\kappa}\prod_j(\gamma^\Psi_j)^{-t_j}\phi_F$, since $\Psi_{\mathcal A_j/g}=\gamma^\Psi_j\Psi_{\mathcal A_j\setminus g}$. Torus homogeneity with $\lambda=1/\gamma^\Psi_1$ extracts $(\gamma^\Psi_1)^{|\mathbf t|-\kappa}$, and $(\gamma^K_1)^{\kappa}\prod_j(\gamma^\Psi_j)^{-t_j}(\gamma^\Psi_1)^{|\mathbf t|-\kappa}=z_{1,\Psi}^{-\kappa}\prod_j(z_{j,\Psi}/z_{1,\Psi})^{-t_j}$, with $z_{1,\Psi}/z_{j,\Psi}=\gamma^\Psi_1/\gamma^\Psi_j$ and $z_{j,K}/z_{1,\Psi}=\gamma^K_j/\gamma^\Psi_1$: the second law form. At $N=1$, $z_{1,K}=1$, $z_{1,\Psi}=\varrho_{\mathcal A}$ and $\mathbf x=1$ give $\sG_{s-\kappa;(1,-1)}(1,\varrho_{\mathcal A}y)$ and $\sG_{t-\kappa;(1,-1)}(1,y/\varrho_{\mathcal A})\varrho_{\mathcal A}^{-\kappa}$.

\emph{(ii)} The Laplace forms are the definitions of the images. For the law form of $\mathfrak r_F$ write $\alpha=\rho u$ and $p=\rho\gamma^K_1(u)$: then $1/\gamma^K_j(\alpha)=1/(pz_{j,K})$ and $\gamma^\Psi_j(\alpha)=pz_{j,\Psi}$ at $\mathbf z=\boldsymbol\varrho_F(u)$, and $\rho^{\kappa-|\mathbf s|-1}d\rho=(\gamma^K_1)^{|\mathbf s|-\kappa}p^{\kappa-|\mathbf s|-1}dp$, which with the measure of (i) leaves $\prod_jz_{j,K}^{-s_j}\phi_F(du)\otimes p^{\kappa-|\mathbf s|-1}dp$. A Laplace transform of a positive measure is completely monotone on the interior of the set where it is finite, by differentiation under the integral, the derivatives being dominated near each such point by the integrand at a point with smaller coordinates, and it is log-convex along segments by Hölder's inequality. For one cell, with $\lambda=1/p$ and $p=1/q$, the measure $e^{-\lambda}\mathfrak r_F$ becomes $e^{-q}q^{s-\kappa-1}dq$ against the law, and the coordinate $\mu=pz$ becomes $z/q$; for $\check{\mathfrak r}_F$, with $p=\rho\gamma^\Psi(u)$, the same computation gives the image of $z^{-\kappa}(\varrho_{\mathcal A})_\ast\phi_{\mathcal A}(dz)\otimes p^{\kappa-t-1}dp$ under $(z,p)\mapsto(1/p,p/z)$, and $e^{-\lambda}$ and $p=1/q$ give the image of $z^{-\kappa}(\varrho_{\mathcal A})_\ast\phi_{\mathcal A}(dz)\otimes q^{t-\kappa-1}e^{-q}dq$ under $(z,q)\mapsto1/(zq)$.

\emph{(iii)} Insert the Laplace form of (ii) and integrate first in $x_j$, $j\ge2$, and $y_j$ (Tonelli): $\int_0^\infty x^{a-1}e^{-x\lambda}dx=\Gamma(a)\lambda^{-a}$. With $\lambda_j=1/(pz_{j,K})$, $\mu_j=pz_{j,\Psi}$ and $x_1=1$,
\[
\prod_{j\ge2}\Gamma(a_j)\prod_j\Gamma(b_j)\iint e^{-1/p}\prod_{j\ge2}(pz_{j,K})^{a_j}\prod_j(pz_{j,\Psi})^{-b_j}\prod_jz_{j,K}^{-s_j}\,p^{\kappa-|\mathbf s|-1}dp\,(\boldsymbol\varrho_F)_\ast\phi_F(d\mathbf z).
\]
The power of $p$ is $p^{-a_1-1}$ with $a_1=\sum_jb_j-\sum_{j\ge2}a_j-\kappa+|\mathbf s|$, the given $a_1$ by the hyperplane, and $\int_0^\infty e^{-1/p}p^{-a_1-1}dp=\Gamma(a_1)$ for $a_1>0$; the power of $\mathbf z$ is $\prod_jz_{j,K}^{-(s_j-a_j)}z_{j,\Psi}^{-b_j}$, $z_{1,K}=1$ (check~5). For $\check f_F$ the law form of $\check{\mathfrak r}_F$, obtained as in (ii) with $p=\rho\gamma^\Psi_1$, is the image of $z_{1,\Psi}^{-\kappa}\prod_j(z_{j,\Psi}/z_{1,\Psi})^{-t_j}(\boldsymbol\varrho_F)_\ast\phi_F(d\mathbf z)\otimes p^{\kappa-|\mathbf t|-1}dp$ under $(\mathbf z,p)\mapsto\big((z_{1,\Psi}/(pz_{j,\Psi}))_j,(pz_{j,K}/z_{1,\Psi})_j\big)$; the same integration leaves $p^{-a_1-1}$ with $a_1=\sum_jb_j-\sum_{j\ge2}a_j-\kappa+|\mathbf t|$, and the total power of $z_{1,\Psi}$ is $-\sum_{j\ge2}a_j+\sum_jb_j-\kappa+\sum_{j\ge2}t_j=a_1-t_1$, so the power of $\mathbf z$ is $\prod_jz_{j,K}^{-b_j}z_{j,\Psi}^{-(t_j-a_j)}$. Both points satisfy $\sum_j(z_{j,K}+z_{j,\Psi})=\kappa$. At $N=1$, $b=z$ and $a=z+s-\kappa$, the first gives $\Gamma(z)\Gamma(z+s-\kappa)\int\varrho_{\mathcal A}^{-z}d\phi_{\mathcal A}$, and the second, with $a=z+t-\kappa$, $\Gamma(z)\Gamma(z+t-\kappa)\int\varrho_{\mathcal A}^{-(\kappa-z)}d\phi_{\mathcal A}$.

\emph{(iv)} Table by table, the dual table $\mathcal A^\ast$ has the bases $(B\sqcup\{g\})\setminus T$ with the weights $\theta_T$, and splitting by whether $g\in T$,
\[
K_{\mathcal A^\ast/g}=\Psi_{\mathcal A\setminus g},\quad K_{\mathcal A^\ast\setminus g}=\Psi_{\mathcal A/g},\quad \Psi_{\mathcal A^\ast\setminus g}=K_{\mathcal A/g},\quad \Psi_{\mathcal A^\ast/g}=K_{\mathcal A\setminus g},
\]
its rank is $|B|+1-r_{\mathcal A}$, and $g$ is neither a loop nor a coloop of it. So $\gamma^K_{\mathcal A^\ast}=\gamma^\Psi_{\mathcal A}$ and $\gamma^\Psi_{\mathcal A^\ast}=\gamma^K_{\mathcal A}$, and the hinge wall at $(t_j,s_j)$ is $m_B-\sum_j\big(t_j(|B|-r_{\mathcal A_j})+s_j(r_{\mathcal A_j}-1)\big)=\kappa$. The base form of $F^\ast$ is $u^{\mvec-\one}\prod_j\Psi_{\mathcal A_j\setminus g}^{-t_j}K_{\mathcal A_j/g}^{-s_j}(\gamma^\Psi_1)^{-\kappa}=(\gamma^\Psi_1/\gamma^K_1)^{-\kappa}\phi_F$, its ratio vector is $(\gamma^\Psi_j/\gamma^\Psi_1,\gamma^K_j/\gamma^\Psi_1)_j=(z_{j,\Psi}/z_{1,\Psi},z_{j,K}/z_{1,\Psi})_j$ at $\mathbf z=\boldsymbol\varrho_F$, which gives the law, and the orthant forms of $f_{F^\ast}$ and $\check f_{F^\ast}$ are those of $\check f_F$ and $f_F$. At $N=1$ the ratio is $1/\varrho_{\mathcal A}$, and $\int\varrho_{\mathcal A^\ast}^{-z}d\phi_{\mathcal A^\ast}=\int\varrho_{\mathcal A}^{z}\varrho_{\mathcal A}^{-\kappa}d\phi_{\mathcal A}$.

\emph{(v)} By (iii) for $f_F$ at $(\mathbf a,\mathbf b)$ and for $\check f_F$ at $(\mathbf t-\mathbf b,\mathbf s-\mathbf a)$, whose hyperplane condition $\sum_j(s_j-a_j)-\sum_j(t_j-b_j)=\kappa-|\mathbf t|$ is that of $(\mathbf a,\mathbf b)$, the Mellin transform of the law appears at the same point $(\mathbf s-\mathbf a,\mathbf b)$ on both sides, with the Gamma products $\prod_j\Gamma(a_j)\Gamma(b_j)$ and $\prod_j\Gamma(t_j-b_j)\Gamma(s_j-a_j)$; all these Gammas are finite and positive, which gives the relation in $(0,\infty]$. At $N=1$, with $z=s-a$ and $b=\kappa-z$, the ratio of the Gammas is that of the displayed one-cell relation (check~6). The expansion clause concerns the node of width one of one table and is the reference's.
\end{proof}

\begin{slemma}\label{lem:hinged-inversion}
Let $F$ be a cell family with plain tables, $s_j,t_j>0$, and $\mathbf a,\mathbf b\in(0,\infty)^N$ real with $\sum_jb_j-\sum_ja_j=\kappa_F-|\mathbf s|$ and $\int\boldsymbol\varrho_F^{-(\mathbf s-\mathbf a,\mathbf b)}d\phi_F<\infty$. Write $\mathbf u=(x_2,\dots,x_N,y_1,\dots,y_N)\in(0,\infty)^{2N-1}$, $g(\mathbf u)=f_F((1,x_2,\dots,x_N),\mathbf y)$, $\mathbf c=(a_2,\dots,a_N,b_1,\dots,b_N)$, and for $\boldsymbol\tau=(\tau_2,\dots,\tau_N,\tau'_1,\dots,\tau'_N)\in\mathbb R^{2N-1}$ put $\tau_1=\sum_j\tau'_j-\sum_{j\ge2}\tau_j$ and
\[
G(\boldsymbol\tau)=\prod_j\Gamma(a_j+i\tau_j)\Gamma(b_j+i\tau'_j)\int\boldsymbol\varrho_F^{-(\mathbf s-\mathbf a-i\boldsymbol\tau_K,\ \mathbf b+i\boldsymbol\tau_\Psi)}d\phi_F,
\]
$\boldsymbol\tau_K=(\tau_j)_j$ and $\boldsymbol\tau_\Psi=(\tau'_j)_j$. Then $g$ is finite and continuous on $(0,\infty)^{2N-1}$, $G$ is the Mellin transform of $g$ on the vertical subspace through $\mathbf c$ and is integrable, and
\[
g(\mathbf u)=(2\pi)^{1-2N}\int_{\mathbb R^{2N-1}}G(\boldsymbol\tau)\,\mathbf u^{-\mathbf c-i\boldsymbol\tau}\,d\boldsymbol\tau\qquad(\mathbf u\in(0,\infty)^{2N-1}).
\]
\end{slemma}

\begin{proof}
The computation of \hb{prop:faces}(iii) at the complex exponents $\mathbf c+i\boldsymbol\tau$ converges absolutely, its integrand having at every point the modulus of the integrand at $\mathbf c$, so it computes the Mellin transform of $g$ at $\mathbf c+i\boldsymbol\tau$ as $G(\boldsymbol\tau)$; the point $(\mathbf s-\mathbf a-i\boldsymbol\tau_K,\mathbf b+i\boldsymbol\tau_\Psi)$ stays on the hyperplane, which fixes $\tau_1$. By \hb{def:hinged}, $|G(\boldsymbol\tau)|\le\prod_j|\Gamma(a_j+i\tau_j)\Gamma(b_j+i\tau'_j)|\int\boldsymbol\varrho_F^{-(\mathbf s-\mathbf a,\mathbf b)}d\phi_F$, and the bound $|\Gamma(c+i\tau)|\le C_c\,e^{-\pi|\tau|/2}(1+|\tau|)^{c-1/2}$ makes $G$ integrable on $\mathbb R^{2N-1}$. The Mellin transform of $g$ at the real point $\mathbf c$ is finite, so $g<\infty$ almost everywhere. By \hb{prop:faces}(ii), $g$ is the Laplace transform of a positive measure in the variables $\mathbf u$, hence nonincreasing in each of them: if $g$ were infinite at a point it would be infinite on the set of points below it, which has positive measure. So $g$ is finite everywhere, and continuous by dominated convergence, the integrand near a point being dominated by the integrand at a point with smaller coordinates. In the coordinates $\mathbf v=\log\mathbf u$ the function $g(e^{\mathbf v})e^{\ip{\mathbf c}{\mathbf v}}$ is integrable with the integrable Fourier transform $G$, so the Fourier inversion theorem gives the display almost everywhere, and everywhere since both sides are continuous.
\end{proof}

\begin{htheorem}[\hb{thm:glue} (Two cell families glued)]
As stated, for two cell families with plain tables and $s_j,t_j>0$.
\end{htheorem}

\begin{proof}
\emph{(i)} \hb{thm:node}(i) with the core $\binom{\{g_1,g_2\}}1$ in every table, $F_1$ at $g_1$ and $F_2$ at $g_2$, the period of the twisted core family being $\mathcal G$.

\emph{(ii), first form.} The gluing of two cells along the hinge is the series closure $\binom{\{h_1,h_2\}}1\langle\mathcal A_1,\mathcal A_2\rangle$ (\hb{cor:blocksclosures}(iii)), so \hb{thm:seriesclosure} at $q=2$ gives
\[
P=C\int\prod_jx_j^{s_j-1}\prod_{j\ge2}y_j^{t_j-1}\,f_{F_1}f_{F_2}\big(\mathbf x,(1,\mathbf y')\big)\,d\mathbf x\,d\mathbf y',\qquad C=\prod_j\frac1{\Gamma(s_j)\Gamma(t_j)} .
\]
Substitute $\mathbf x=x_1(1,\boldsymbol\xi)$, $d\mathbf x=x_1^{N-1}dx_1\,d\boldsymbol\xi$, and use the homogeneity of \hb{def:hinged} with $c=x_1$, $f_{F_i}(\mathbf x,(1,\mathbf y'))=x_1^{\kappa_{F_i}-|\mathbf s|}f_{F_i}((1,\boldsymbol\xi),x_1(1,\mathbf y'))$; then substitute $\boldsymbol\eta=x_1(1,\mathbf y')$, $dx_1\,d\mathbf y'=x_1^{1-N}d\boldsymbol\eta$. The powers of $x_1=\eta_1$ collect to
\[
(|\mathbf s|-N)+(N-1)+(\kappa_{F_1}+\kappa_{F_2}-2|\mathbf s|)+(1-N)-(|\mathbf t|-t_1)+(N-1)=t_1-1,
\]
since $\kappa_{F_1}+\kappa_{F_2}=|\mathbf s|+|\mathbf t|$ (check~10 of S21-23\_work/checks.py): the first form.

\emph{(ii), second form.} The complements of the bases of a two-sum of cells are the bases of the two-sum of the dual cells: for $X$ a basis of $\mathcal A\setminus g$ and $Y\sqcup\{g\}$ one of $\mathcal A'$, $B\setminus X$ is a basis of $\mathcal A^\ast/g$ and $R\setminus Y$ one of $\mathcal A'^\ast\setminus g$, and symmetrically (\hb{lem:gluemark}), with the product weights. By \hb{prop:dualfamily} the period of $F_1\oplus_gF_2$ at $(\mathbf s,\mathbf t)$ is that of $F_1^\ast\oplus_gF_2^\ast$ at $(\mathbf t,\mathbf s)$, and by \hb{prop:faces}(iv) the dual cell families have the hinge walls $\kappa_{F_1},\kappa_{F_2}$ and the transforms $\check f_{F_1},\check f_{F_2}$. The first form for them is the second form.

\emph{(iii)} Insert the Laplace forms of \hb{prop:faces}(ii) into (ii) and integrate first in $x_j$, $j\ge2$, and $y_j$ (Tonelli): $\int_0^\infty x^{s-1}e^{-x(\lambda+\lambda')}dx=\Gamma(s)(\lambda+\lambda')^{-s}$, and $x_1=1$ leaves $e^{-\lambda_1-\lambda'_1}$; the factors $\Gamma(s_j)$, $j\ge2$, and $\Gamma(t_j)$ cancel against $C$. The period converges exactly when every proper wall is positive (\hb{thm:projective}).

\emph{(iv)} Let the contour have the real point $(\mathbf a^d,\mathbf b^d,\mathbf a^r,\mathbf b^r)$, all coordinates positive, with $\mathbf a^d+\mathbf a^r=\mathbf s$, $\mathbf b^d+\mathbf b^r=\mathbf t$ and $\sum_j(a^r_j+b^d_j)=\kappa_{F_1}$, and put $\mathbf a=\mathbf a^d$, $\mathbf b=\mathbf b^d$; then $\sum_jb_j-\sum_ja_j=\kappa_{F_1}-|\mathbf s|$, and the Mellin transform of the law of $F_1$ at $(\mathbf s-\mathbf a,\mathbf b)=(\mathbf a^r,\mathbf b^d)$ is finite. By Lemma~\ref{lem:hinged-inversion}, $f_{F_1}$ is the inverse Mellin integral of $G_1$ through $\mathbf c$. Insert it into the first form of (ii). The double integral converges absolutely: integrating first in $\mathbf u$ gives $\int|G_1|\,d\boldsymbol\tau$ times $\int\prod_{j\ge2}x_j^{s_j-a_j-1}\prod_jy_j^{t_j-b_j-1}f_{F_2}\,d\mathbf x'\,d\mathbf y$, which is \hb{prop:faces}(iii) for $F_2$ at $(\mathbf s-\mathbf a,\mathbf t-\mathbf b)$: its first coordinate is $|\mathbf t|-|\mathbf b|-(|\mathbf s|-s_1-\sum_{j\ge2}a_j)-\kappa_{F_2}+|\mathbf s|=s_1-a_1$, by $\kappa_{F_2}=2\sigma-\kappa_{F_1}$, so it equals $\prod_j\Gamma(s_j-a_j)\Gamma(t_j-b_j)\int\boldsymbol\varrho_{F_2}^{-(\mathbf a^d,\mathbf b^r)}d\phi_{F_2}<\infty$. By Fubini,
\[
P=C\,(2\pi)^{1-2N}\int_{\mathbb R^{2N-1}}\prod_j\Gamma(a^d_j)\Gamma(b^d_j)\Gamma(a^r_j)\Gamma(b^r_j)\int\boldsymbol\varrho_{F_1}^{-(\mathbf a^r,\mathbf b^d)}d\phi_{F_1}\int\boldsymbol\varrho_{F_2}^{-(\mathbf a^d,\mathbf b^r)}d\phi_{F_2}\,d\boldsymbol\tau,
\]
the shares moving with the imaginary parts $\boldsymbol\tau$ of $(a^d_2,\dots,a^d_N,b^d_1,\dots,b^d_N)$, the others fixed by the pins, and $C\prod_j\Gamma(\cdot)$ is $\prod_j\Gst42{a^d_j,\ b^d_j,\ a^r_j,\ b^r_j}{s_j,\ t_j}$. This is the fibre integral of \hb{def:fibre} over the level shares with their $2N+1$ pins: solved for $a^r_j$, $b^r_j$ and $a^d_1$ the matrix of the pins has $|\det|=1$, so the fibre integral is $(2\pi)^{1-2N}$ times the integral over $\boldsymbol\tau$. In the pair coordinates of \hb{thm:twosides}, $\sigma^d_j=\frac12(a^d_j+b^d_j)$ and $\delta^d_j=\frac12(a^d_j-b^d_j)$ with the one pin $\sum_j\delta^d_j=\frac12(|\mathbf s|-\kappa_{F_1})$, the linear map from $\boldsymbol\tau$ to the imaginary parts of $(\sigma^d_1,\dots,\sigma^d_N,\delta^d_2,\dots,\delta^d_N)$ has $|\det|=2^{1-N}$: the map $(a,b)\mapsto(\frac{a+b}2,\frac{a-b}2)$ of $\mathbb R^{2N}$ has $|\det|=2^{-N}$, it sends the normal coordinate $\sum_j(a_j-b_j)$ of the pin to twice the normal coordinate $\sum_j\delta_j$, and the quotient is $2^{1-N}$. So the integral is $2^{N-1}$ times the fibre integral in the pair coordinates, the formula of \hb{thm:twosides}. That the two Mellin transforms are the periods of the bunch family and of the room family at the level shares is the identity $u^{\mvec-\one}\prod_jK_{\mathcal A_j/g}^{-a^d_j}\Psi_{\mathcal A_j/g}^{-b^d_j}K_{\mathcal A_j\setminus g}^{-a^r_j}\Psi_{\mathcal A_j\setminus g}^{-b^r_j}=\boldsymbol\varrho_{F_1}^{-(\mathbf a^r,\mathbf b^d)}(\gamma^K_{\mathcal A_1})^{\kappa_{F_1}-\sum_j(a^r_j+b^d_j)}$ times the density of $\phi_{F_1}$, the last power being $1$ on the pin, with $\cB_j|B=\mathcal A_j\setminus g$ and $\cB_j/R=\mathcal A_j/g$ (\hb{lem:twosumpieces}), and its analogue for the room, whose pin $\sum_j(a^d_j+b^r_j)=2\sigma-\kappa_{F_1}=\kappa_{F_2}$ follows from the letter totals.

\emph{One table.} At $N=1$ the laws are measures on $(0,\infty)$, the kernel is that of \hb{lem:twomarks}, (ii) and (iii) have no variables $x_j$, the image of $e^{-\lambda}\mathfrak r_{F_i}$ under $\mu$ being $\mathfrak r^\Psi$, and in (iv) the fibre is the line of $b^d=z$, with $a^d=s-\kappa_{\mathcal A}+z$, $a^r=\kappa_{\mathcal A}-z$, $b^r=t-z$; the conditions on the contour are $\max(0,\kappa_{\mathcal A}-s)<z_0<\min(t,\kappa_{\mathcal A})$ and the finiteness of $\int\varrho_{\mathcal A}^{-z_0}d\phi_{\mathcal A}$ and $\int\varrho_{\mathcal A'}^{z_0-t}d\phi_{\mathcal A'}$, and $d\sigma^d=dz$. The first form of (ii) is \hb{thm:onescale}(i) with the room transform $y^tf_{\mathcal A'}(y)$.
\end{proof}

\begin{hdefinition}[\hb{def:onemark} (One weighted mark)]
As stated: the one-mark cell family of a ratio vector, with two systems of weights in every table.
\end{hdefinition}

\begin{proof}
Table $j$ has the bases $\{e\}$ and $\{g\}$, with $\theta^K=(x_{0,j,K},1)$ and $\theta^\Psi=(1,x_{0,j,\Psi})$, so $K_j=x_{0,j,K}\alpha_e+\alpha_g$ and $\Psi_j=\theta^\Psi_{\{e\}}\alpha_g+\theta^\Psi_{\{g\}}\alpha_e=\alpha_g+x_{0,j,\Psi}\alpha_e$, two systems of weights being allowed (\hb{def:substitution}, \hb{def:family}). The hinge minors are $K_{\setminus g}=x_{0,j,K}\alpha_e$, $K_{/g}=1$, $\Psi_{\setminus g}=1$ and $\Psi_{/g}=x_{0,j,\Psi}\alpha_e$, so the hinge values are $x_{0,j,K}\alpha_e$ and $x_{0,j,\Psi}\alpha_e$ and the ratio vector is the constant $\mathbf x_0$, $x_{0,1,K}=1$. The rank is $1$ and $|B|=1$, so the hinge wall is $m_e$. The simplex $\Delta_{\{e\}}$ is the point $u_e=1$, carrying the unit mass, and there the density of the base form is $1\cdot1\cdot(x_{0,1,K})^{-m_e}=1$: the law is $\delta_{\mathbf x_0}$.
\end{proof}

\begin{htheorem}[\hb{thm:retained} (The law is what a cell family retains)]
As stated, for every pair of cell families with plain tables and $s_j,t_j>0$.
\end{htheorem}

\begin{proof}
\emph{(a)$\Rightarrow$(b).} By \hb{thm:node}(i) with the core $\binom{\{g_1,g_2\}}1$, the period of $F\oplus_gF''$ depends on $F$ only through $\kappa$ and its law; this needs no plainness of $F''$. (b)$\Rightarrow$(c) is a special case.

\emph{(c)$\Rightarrow$(d).} By \hb{thm:node}(i) and the law $\delta_{\mathbf x_0}$ of the one-mark cell family (\hb{def:onemark}), the period against it is $\Pi_F(\mathbf x_0)=\int\mathcal G(\mathbf z,\mathbf x_0)\,d((\boldsymbol\varrho_F)_\ast\phi_F)(\mathbf z)$, $\mathcal G(\mathbf z,\mathbf x_0)=\int_0^\infty v^{\kappa-1}\prod_j(z_{j,K}v+x_{0,j,K})^{-s_j}(z_{j,\Psi}v+x_{0,j,\Psi})^{-t_j}dv$. Since $(zv+x)/(zv+x')$ lies between $\min(1,x/x')$ and $\max(1,x/x')$, $\mathcal G(\cdot,\mathbf x_0)$ and $\mathcal G(\cdot,\mathbf x'_0)$ are comparable within constants that depend on $\mathbf x_0,\mathbf x'_0$ only, so $\Pi_F$ is finite everywhere on $(0,\infty)^{2N-1}$ or nowhere.

Take the Mellin transform of $\Pi_F$ in the coordinates of $\mathbf x_0$ other than $x_{0,1,K}=1$, with the exponents $s_j-z_{j,K}$ at $x_{0,j,K}$, $j\ge2$, and $t_j-z_{j,\Psi}$ at $x_{0,j,\Psi}$. By Tonelli and $\int_0^\infty q^{w-1}(c+q)^{-s}dq=c^{w-s}\Gst21{w,\ s-w}{s}$, $0<w<s$, each coordinate gives a Beta corner and the power $(z_{j,K}v)^{-z_{j,K}}$ or $(z_{j,\Psi}v)^{-z_{j,\Psi}}$, and the remaining integral is $\int_0^\infty v^{z_{1,K}-1}(1+v)^{-s_1}dv=\Gst21{z_{1,K},\ s_1-z_{1,K}}{s_1}$, the power of $v$ collecting to $v^{z_{1,K}-1}$ by the hyperplane. So the Mellin transform is
\[
\prod_{j\ge2}\Gst21{s_j-z_{j,K},\ z_{j,K}}{s_j}\,\Gst21{z_{1,K},\ s_1-z_{1,K}}{s_1}\prod_j\Gst21{t_j-z_{j,\Psi},\ z_{j,\Psi}}{t_j}\int\boldsymbol\varrho_F^{-\mathbf z}d\phi_F<\infty .
\]
In particular $\Pi_F$ is finite almost everywhere, hence everywhere. For complex exponents with these real parts the same computation converges absolutely, and the Beta corners at complex arguments of positive real part do not vanish. Where finite, $\Pi_F$ extends holomorphically to the product of right half-planes: $|(zv+\zeta)^{-s}|\le(zv+\re\zeta)^{-s}$ for $\re\zeta>0$ and $s>0$, so dominated convergence and Morera's theorem give holomorphy in each coordinate, the extension being locally bounded, and hence jointly. So $\Pi_F$ and $\Pi_{F'}$ are real-analytic on $(0,\infty)^{2N-1}$, equal on a nonempty open set, and equal everywhere; their Mellin transforms on the vertical subspace through $\mathbf z$ agree, which is (d).

\emph{(d)$\Rightarrow$(a).} The finite measures $\boldsymbol\varrho^{-\mathbf z}(\boldsymbol\varrho_F)_\ast\phi_F$ and $\boldsymbol\varrho^{-\mathbf z}(\boldsymbol\varrho_{F'})_\ast\phi_{F'}$ on $(0,\infty)^{2N-1}$, read in the logarithms of the coordinates, have the Fourier transforms $\boldsymbol\tau\mapsto\int\boldsymbol\varrho^{-(\mathbf z+i\boldsymbol\tau)}d\phi$, equal by (d); a finite measure on $\mathbb R^{2N-1}$ is determined by its Fourier transform, so the two measures agree, and multiplying by the positive continuous function $\boldsymbol\varrho^{\mathbf z}$ the laws agree. This replaces the Gaussian smoothing of the reference proof of \hb{thm:retained}, which is the same argument on the line.

\emph{One table.} The plain one-mark cell of ratio $x_0$ has the law $x_0^{(m_e-s)/2}\delta_{x_0}$, a positive multiple of the point law, and its ratios fill $(0,\infty)$, so (c) with the plain cells gives the equality of $\Pi_F$ on an open interval; (d) is the case $N=1$, $\mathbf w=(w_{1,K},w)$ with $w_{1,K}$ not entering.
\end{proof}

\begin{hcorollary}[\hb{cor:piecesmarks} (Cells that are marks)]
As stated: (i) for every cell family with plain tables and $s_j,t_j>0$.
\end{hcorollary}

\begin{proof}
\emph{(i)} This is the reference proof of \hb{cor:piecesmarks}(i) with the ratio vector in place of the ratio. If $\boldsymbol\varrho_F\equiv\mathbf x_0$, the image of $\phi_F$ is $\phi_F(\Delta_B)\delta_{\mathbf x_0}$. Conversely, $\phi_F$ has a positive continuous density on the open simplex and $\boldsymbol\varrho_F$ is continuous, so if the law is carried by $\mathbf x_0$ the open set $\{\boldsymbol\varrho_F\neq\mathbf x_0\}$ has $\phi_F$-measure $0$ and is empty. The mass $\phi_F(\Delta_B)$ is positive. By \hb{thm:node}(i), the period of a gluing $F\oplus_gF''$ is $\phi_F(\Delta_B)\int\mathcal G(\mathbf x_0,\mathbf z')\,d((\boldsymbol\varrho_{F''})_\ast\phi_{F''})(\mathbf z')$, a positive multiple of $\phi_F(\Delta_B)$; so a finite period forces $\phi_F(\Delta_B)<\infty$, and the integral is the period of the gluing of the one-mark cell family of ratio vector $\mathbf x_0$ and mass $\kappa_F$ with $F''$, whose law is $\delta_{\mathbf x_0}$ and whose hinge wall is $\kappa_F$. For one plain cell, the one-mark cell of ratio $x_0$ and mass $\kappa_{\mathcal A}$ has the law $x_0^{(\kappa_{\mathcal A}-s)/2}\delta_{x_0}$, which gives the factor. (ii) and (iii) are the reference's.
\end{proof}

\begin{htheorem}[\hb{thm:seriesparallel} (Series and parallel nodes)]
As stated: (i) and (ii) for cell families with plain tables and $s_j,t_j>0$.
\end{htheorem}

\begin{proof}
This is the reference proof of \hb{thm:seriesparallel}(i),(ii), table by table. \emph{(i)} The series core $\binom Gq$, $G=\{g_0,\dots,g_q\}$, has rank $q$; its deletion at $g_0$ has the one basis $\{g_1,\dots,g_q\}$, so $K_{\binom Gq\setminus g_0}(\mathbf v)=\prod_iv_i$ and $\Psi_{\binom Gq\setminus g_0}=1$, and its contraction has the bases of size $q-1$ of $\{g_1,\dots,g_q\}$, so $K_{\binom Gq/g_0}=e_{q-1}(\mathbf v)$ and $\Psi_{\binom Gq/g_0}(\mathbf v)=\sum_iv_i$. By \hb{lem:nodesum},
\[
K_{\mathcal A_j\setminus g_0}=\prod_iK_{\mathcal A_{ij}/g_i}\gamma^K_{\mathcal A_{ij}}=\prod_iK_{\mathcal A_{ij}\setminus g_i},\qquad \Psi_{\mathcal A_j\setminus g_0}=\prod_i\Psi_{\mathcal A_{ij}\setminus g_i},
\]
and the hinge values of $\mathcal A_j$ are $\prod_i\gamma^K_{\mathcal A_{ij}}/e_{q-1}((\gamma^K_{\mathcal A_{ij}})_i)$, whose reciprocal is $\sum_i1/\gamma^K_{\mathcal A_{ij}}$, and $\sum_i\gamma^\Psi_{\mathcal A_{ij}}$. In the orthant form of $f_F$ the measure $\alpha^{\mvec-\one}\prod_jK_{\mathcal A_j\setminus g_0}^{-s_j}\Psi_{\mathcal A_j\setminus g_0}^{-t_j}d\alpha$ is the product over $i$ of the measures of the cell families, and the exponential is the product of their exponentials, so Tonelli gives $f_F=\prod_if_{F_i}$, and the image of the product measure under the sum of the maps of the factors is the convolution: $\mathfrak r_F=\ast_i\mathfrak r_{F_i}$. The hinge-wall formula of \hb{thm:node}(ii), with $\cE_{\mathcal H}=\varnothing$, $k=q$ and $r_{\mathcal H_j}=q$, gives $\kappa_F=\sum_i\kappa_{F_i}-(q-1)|\mathbf s|$ (check~3a of S21-23\_work/checks.py). For one table, $f_{\mathcal A}(y)=f_F(1,y)=\prod_if_{F_i}(1,y)$, and since $e^{-\lambda}=\prod_ie^{-\lambda_i}$ for $\lambda=\sum_i\lambda_i$, the measure $e^{-\lambda}\mathfrak r_F$ is the convolution of the $e^{-\lambda_i}\mathfrak r_{F_i}$, whose images under $\mu$ convolve: $\mathfrak r^\Psi_{\mathcal A}=\ast_i\mathfrak r^\Psi_{\mathcal A_i}$. The twisted series core is the reference's.

\emph{(ii)} The parallel core $\binom G1$ has rank $1$: $K_{\binom G1\setminus g_0}=\sum_iv_i$, $K_{\binom G1/g_0}=1$, $\Psi_{\binom G1\setminus g_0}=e_{q-1}(\mathbf v)$ and $\Psi_{\binom G1/g_0}=\prod_iv_i$. So $K_{\mathcal A_j/g_0}=\prod_iK_{\mathcal A_{ij}/g_i}$, $\Psi_{\mathcal A_j/g_0}=\prod_i\Psi_{\mathcal A_{ij}\setminus g_i}\gamma^\Psi_{\mathcal A_{ij}}=\prod_i\Psi_{\mathcal A_{ij}/g_i}$, the hinge values combine as stated, the orthant form of $\check f_F$ factors, and $\kappa_F=\sum_i\kappa_{F_i}-(q-1)|\mathbf t|$ (check~3b). Equally, \hb{prop:faces}(iv) exchanges (i) and (ii) through the dual cell families, the dual of a parallel core being a series core. (iii) and (iv) are the reference's.
\end{proof}

\begin{hdefinition}[\hb{def:hingetree} (Hinge tree)]
As stated: the common hinge tree of the tables of a family, with its node families and its cell families.
\end{hdefinition}

\begin{sremark}
The added objects are definitions. Each $\mathcal A_{v,j}$ is a cell with the hinge $g_{v,\mathfrak e_v}$ on the marks below $v$ (\hb{def:hingetree}), the same for every $j$, so the $\mathcal A_{v,j}$ form a cell family in the sense of \hb{def:hinged}; and the substitution along the tree does not depend on the order of the two-sums, table by table (\hb{lem:nodesum}).
\end{sremark}

\begin{htheorem}[\hb{thm:treeperiod} (The period along a hinge tree)]
As stated: for every family whose tables have a common hinge tree, at real data.
\end{htheorem}

\begin{proof}
Induction from the leaves. At a leaf $v$ the cell family $F_v$ is the node family with its one marker as hinge. At an inner vertex $v\neq o$, table $j$ of $F_v$ is the substitution of the node $\mathcal H_{v,j}$, a hinged core with the hinge $g_{v,\mathfrak e_v}$, with the subtree cells $\mathcal A_{w,j}$ of the vertices $w$ directly below it at the markers of their edges (\hb{def:hingetree}), so \hb{thm:node}(ii) gives (i), the masses on $\cE_v$ being those of $F$. At the root, the node has no hinge, and \hb{thm:node}(i) gives (ii). (iii) is the recursion of (i): the law at $v$ is computed from the nodes and masses at $v$ and the laws of the vertices directly below, and by induction from the data below $v$. This is the reference proof of \hb{thm:treeperiod} (thm:contraction), with \hb{thm:node} for families, which the reference's remark on families (rem:familiesnodes) gives verbatim.
\end{proof}

\begin{hproposition}[\hb{prop:work} (What the work depends on)]
As stated, for a family whose tables have a common hinge tree.
\end{hproposition}

\begin{proof}
\emph{(i)} A bijection of marks carrying the subtree cells onto each other in every table, with their weights, masses and hinge, carries the letters onto the letters, hence the base forms and the ratio vectors onto each other; the ranks, sizes and masses are kept, and so is the hinge wall.

\emph{(ii)} At a series node of unit weights, $F_v$ is the series connection of the cell families directly below (\hb{def:spnodes}), so \hb{thm:seriesparallel}(i) gives $f_{F_v}=\prod_wf_{F_w}$ and $\kappa_{F_v}=\sum_w\kappa_{F_w}-(q-1)|\mathbf s|$, $q=\sum_lk_l$; members of one orbit have equal transforms and hinge walls by (i), and collecting them gives the powers. A weight $\theta_{g_i}$ on the basis of a series node that omits the marker $g_i$ is carried, by the weight formula of \hb{def:substitution}, by the cell at $g_i$ as the factor $\theta_{g_i}/\theta_{g_0}$ on its bases through its hinge, the weight $\theta_{g_0}$ of the basis omitting the hinge remaining as a constant factor of the table; so a series node with weights is a series connection of cells so weighted, and the orbits are taken of these. The parallel node is the same with \hb{thm:seriesparallel}(ii).

\emph{(iii)} By \hb{def:hinged} the hinge wall of $F_v$ is $m_{B_v}-\sum_j\big(s_j(r_{\mathcal A_{v,j}}-1)+t_j(|B_v|-r_{\mathcal A_{v,j}})\big)$, $B_v$ the marks below $v$, and by \hb{lem:nodesum} the ranks $r_{\mathcal A_{v,j}}$ are sums of the ranks of the nodes below $v$ less one per edge. A change at $v_0$ keeping the total mass at $v_0$, the ranks and the numbers of marks of the nodes, keeps every $m_{B_v}$, rank and size, since $B_v$ contains all of $\cE_{v_0}$ or none of it; so no hinge wall changes, and by \hb{thm:treeperiod}(iii) the law of $F_v$ changes only if $v_0$ lies below $v$, that is, for $v$ on the path from $v_0$ to the root. A change of masses at $v_0$ and $v_1$ keeping the total changes $m_{B_v}$ exactly when $B_v$ contains one of $\cE_{v_0},\cE_{v_1}$ and not the other, by the change of the total mass below $\mathfrak e_v$.
\end{proof}

\begin{hdefinition}[\hb{def:treeproc} (The hinge tree procedure)]
As stated: for a family whose tables have a common hinge tree, the representations (a) and (b) in $2N-1$ variables, (c) for one table.
\end{hdefinition}

\begin{sremark}
Each operation rests on a statement for families. The laws and hinge walls of the orbits are those of \hb{thm:treeperiod} and \hb{prop:work}(i). (1) The one-mark leaf: by \hb{prop:faces}(i) at the point law $\delta_{\mathbf x_0}$ with $\kappa=m_e$, $f(\mathbf x,\mathbf y)=\prod_jx_{0,j,K}^{-s_j}\sG_{|\mathbf s|-m_e;(1,-1)}(\sum_jx_j/x_{0,j,K},\sum_jy_jx_{0,j,\Psi})$, and the second law form with torus homogeneity at $\lambda=x_{0,1,\Psi}$ gives $\check f(\mathbf x,\mathbf y)=\prod_jx_{0,j,\Psi}^{-t_j}\sG_{|\mathbf t|-m_e;(1,-1)}(\sum_jx_j/x_{0,j,\Psi},\sum_jy_jx_{0,j,K})$, the powers of $x_{0,1,\Psi}$ cancelling; the orthant forms give the same directly. For one table, $\mathbf x_0=(1,x_0)$ and the plain cell's law $x_0^{(m_e-s)/2}\delta_{x_0}$ give the displayed one-table transforms, with $\sG_{t-m_e;(1,-1)}(x_0^{-1},y)=x_0^{t-m_e}\sG_{t-m_e;(1,-1)}(1,y/x_0)$. (2) \hb{thm:seriesparallel}(i),(ii), with \hb{prop:work}(ii) for the weights; a pointwise product of transforms is a Mellin convolution of their Mellin transforms of \hb{prop:faces}(iii). (3) \hb{prop:faces}(v) on the Mellin side, and Lemma~\ref{lem:hinged-inversion} for the inverse Mellin transform, applied to the dual cell family for $\check f$ (\hb{prop:faces}(iv)). (4) \hb{thm:node}(ii) for families. (5) \hb{thm:node}(i), and for a series or parallel root \hb{thm:glue}(ii) with the series or parallel connection of all but one cell family on one side.
\end{sremark}

\begin{hcorollary}[\hb{cor:sptables} (Series-parallel families)]
As stated, for a family whose tables have a common hinge tree of series and parallel nodes over one-mark cell families.
\end{hcorollary}

\begin{proof}
This is the reference proof of \hb{cor:sptables} with the family statements. The leaves have their transforms in closed form, Grammas of the marks (the remark after \hb{def:treeproc} in this part). At a series node the transform of the subtree is the product of those below, and at a parallel node the dual transform (\hb{thm:seriesparallel}), so no integral is taken there; repeated subtrees enter as powers (\hb{prop:work}(ii)). At an edge joining a series node to a parallel node the transform of the subtree below is changed into its dual transform by \hb{prop:faces}(v): the Mellin transform of $f$ on a vertical subspace through a real point $(\mathbf a,\mathbf b)$, $0<a_j<s_j$, $0<b_j<t_j$, of its domain, an integral on $(0,\infty)^{2N-1}$; multiplication by the ratio of Gammas; and the inverse Mellin transform of the result, which is the Mellin transform of $\check f$ (\hb{prop:faces}(iii) at $(\mathbf t-\mathbf b,\mathbf s-\mathbf a)$), an integral on $(0,\infty)^{2N-1}$ that converges absolutely and returns $\check f$ by Lemma~\ref{lem:hinged-inversion} applied to the dual cell family (\hb{prop:faces}(iv)). The hypothesis is the finiteness of the Mellin transform of the law at $\mathbf z=(\mathbf s-\mathbf a,\mathbf b)$, which these points require. At the root one integral on $(0,\infty)^{2N-1}$ remains, \hb{thm:glue}(ii). The count of the integrations depends only on the edges where the kind of node changes and on the root. For one table the variables are one.
\end{proof}

\section{Chapter~\hb{ch:repetition}: repetition and the change of query}\label{sec:ch23}

\begin{htheorem}[\hb{thm:repeatwall} (How often a cell family repeats)]
As stated: for every cell family with plain tables, at real exponents.
\end{htheorem}

\begin{proof}
\emph{(i)} The series connection of $k$ copies of $\mathcal A_j$ has $k|B|$ marks besides the hinge, the masses repeated, and the rank $k+k(r_{\mathcal A_j}-1)=kr_{\mathcal A_j}$ (\hb{lem:nodesum}, the series core of rank $k$). By \hb{def:hinged}, since $s_j(kr_{\mathcal A_j}-1)=ks_j(r_{\mathcal A_j}-1)+(k-1)s_j$ and $t_j(k|B|-kr_{\mathcal A_j})=kt_j(|B|-r_{\mathcal A_j})$,
\[
\kappa_{F^{\mathrm s}_k}=km_B-\sum_j\big(s_j(kr_{\mathcal A_j}-1)+t_j(k|B|-kr_{\mathcal A_j})\big)=k\kappa_F-(k-1)|\mathbf s|=|\mathbf s|+k(\kappa_F-|\mathbf s|).
\]
The parallel connection has the rank $1+k(r_{\mathcal A_j}-1)$, and the same count gives $k\kappa_F-(k-1)|\mathbf t|$ (checks~3a and~3b of S21-23\_work/checks.py, which are this count at the exponents as symbols). No sign of an exponent enters.

\emph{(ii)} The walls of $F^{\mathrm s}_k\oplus_gF'$ at the marks of the copies and at those of $F'$ are the hinge walls $\kappa_{F^{\mathrm s}_k}$ and $\kappa_{F'}$, and on the balance wall they add to $2\sigma$ (\hb{lem:twosumpieces}). A convergent period has every proper wall positive (\hb{thm:projective}, at any complex exponents), so $0<\kappa_{F^{\mathrm s}_k}<2\sigma$. If $\kappa_F\neq|\mathbf s|$ the middle term is unbounded in $k$, which bounds $k$: $|\mathbf s|-k(|\mathbf s|-\kappa_F)>0$ gives $k<|\mathbf s|/(|\mathbf s|-\kappa_F)$ when $\kappa_F<|\mathbf s|$, and $|\mathbf s|+k(\kappa_F-|\mathbf s|)<|\mathbf s|+|\mathbf t|$ gives $k<|\mathbf t|/(\kappa_F-|\mathbf s|)$ when $\kappa_F>|\mathbf s|$. The parallel case is the same with the roles of $|\mathbf s|$ and $|\mathbf t|$ exchanged.
\end{proof}

\begin{htheorem}[\hb{thm:seriesrepeat} (Series repetition)]
As stated: (i) for cell families with plain tables and $s_j,t_j>0$; (ii) to (iv) are the reference's, for one cell.
\end{htheorem}

\begin{proof}
\emph{(i)} By \hb{thm:seriesparallel}(i) with $k$ copies of $F$, $f_{F^{\mathrm s}_k}=f_F^k$; by \hb{thm:repeatwall}(i), $\kappa_{F^{\mathrm s}_k}=|\mathbf s|$, which lies in $(0,2\sigma)$ since $|\mathbf s|,|\mathbf t|>0$. So the first form of \hb{thm:glue}(ii) applies to $F^{\mathrm s}_k$ and $F'$ and gives the display. With $\nu$ the measure of the statement, $P(k)=\int f_F^k\,d\nu=\int w^k\,d((f_F)_\ast\nu)(w)$. Where finite, $P(k)^2=\big(\int f_F^{(k-1)/2}f_F^{(k+1)/2}d\nu\big)^2\le P(k-1)P(k+1)$ by the Cauchy--Schwarz inequality, and for $z\ge0$ Tonelli gives $\sum_kP(k)z^k/k!=\int e^{zf_F}d\nu$.

\emph{The value $P(0)$.} Insert the orthant form of $f_{F'}$ and integrate first in $x_j$, $j\ge2$, and $y_j$ (Tonelli): $\int_0^\infty x^{s_j-1}e^{-x/\gamma^K}dx=\Gamma(s_j)(\gamma^K)^{s_j}$ and $\int_0^\infty y^{t_j-1}e^{-y\gamma^\Psi}dy=\Gamma(t_j)(\gamma^\Psi)^{-t_j}$. Since $K_{\mathcal A'_j\setminus g}(\gamma^K_{\mathcal A'_j})^{-1}=K_{\mathcal A'_j/g}$ and $\Psi_{\mathcal A'_j\setminus g}\gamma^\Psi_{\mathcal A'_j}=\Psi_{\mathcal A'_j/g}$, what remains is
\[
\prod_{j\ge2}\Gamma(s_j)\prod_j\Gamma(t_j)\int\alpha^{\mvec|_R-\one}e^{-1/\gamma^K_{\mathcal A'_1}}K_{\mathcal A'_1\setminus g}^{-s_1}\prod_{j\ge2}K_{\mathcal A'_j/g}^{-s_j}\prod_j\Psi_{\mathcal A'_j/g}^{-t_j}\,d\alpha .
\]
In $\alpha=\rho u$ the power of $\rho$ is $\rho^{\kappa_{F'}-|\mathbf t|-s_1-1}=\rho^{-s_1-1}$, by $\kappa_{F'}=|\mathbf t|$ (check~9 of S21-23\_work/checks.py), and $\int_0^\infty\rho^{-s_1-1}e^{-1/(\rho\gamma)}d\rho=\Gamma(s_1)\gamma^{s_1}$ turns $K_{\mathcal A'_1\setminus g}^{-s_1}$ into $K_{\mathcal A'_1/g}^{-s_1}$. So $P(0)$ is the period at $\mvec|_R$ of the family of the contractions, whose degree is $\sum_j(s_j(r_{\mathcal A'_j}-1)+t_j(|R|-r_{\mathcal A'_j}+1))$, so that it is balanced exactly when $\kappa_{F'}=|\mathbf t|$. For one cell, $\int\varrho_{\mathcal A'}^{-t}d\phi_{\mathcal A'}$ with $\kappa_{\mathcal A'}=t$ has the density $u^{\mvec-\one}K_{\mathcal A'/g}^{t-s}K_{\mathcal A'\setminus g}^{-t}\Psi_{\mathcal A'\setminus g}^{-t}(\gamma^\Psi/\gamma^K)^{-t}=u^{\mvec-\one}K_{\mathcal A'/g}^{-s}\Psi_{\mathcal A'/g}^{-t}$, the same period.

(ii) to (iv) are the reference's, for one cell.
\end{proof}

\begin{sremark}
For a family the asymptotics of $P(k)$ is a Laplace problem in $2N-1$ variables: a series-balanced cell family has $f_F=\int\sG_{0;(1,-1)}(1,ab)\,d\phi_{F\setminus g}$ by \hb{thm:seriesclosure}, of the order of $P_{F\setminus g}$ times the logarithm of the scale of $\mathbf y$ plus a term that depends on the direction of $\mathbf y$ and on $\mathbf x'$, and the uniformity of the error in these directions does not follow from the one-variable argument. So (ii) to (iv) are stated for one cell.
\end{sremark}

\begin{hcorollary}[\hb{cor:parallelrepeat} (Parallel repetition)]
As stated: the first display for cell families with plain tables and $s_j,t_j>0$; the rest is the reference's, for one cell.
\end{hcorollary}

\begin{proof}
By \hb{thm:seriesparallel}(ii), $\check f_{F^{\mathrm p}_k}=\check f_F^k$, and by \hb{thm:repeatwall}(i), $\kappa_{F^{\mathrm p}_k}=|\mathbf t|\in(0,2\sigma)$; the second form of \hb{thm:glue}(ii) gives the display. For $P(0)$, integrate the orthant form of $\check f_{F'}$ in $x_j$, $j\ge2$, and $y_j$: $\int x^{t_j-1}e^{-x/\gamma^\Psi}dx=\Gamma(t_j)(\gamma^\Psi)^{t_j}$ and $\int y^{s_j-1}e^{-y\gamma^K}dy=\Gamma(s_j)(\gamma^K)^{-s_j}$ turn $\Psi_{\mathcal A'_j/g}^{-t_j}$ into $\Psi_{\mathcal A'_j\setminus g}^{-t_j}$, $j\ge2$, and every $K_{\mathcal A'_j/g}^{-s_j}$ into $K_{\mathcal A'_j\setminus g}^{-s_j}$; in $\alpha=\rho u$ the power of $\rho$ is $\rho^{\kappa_{F'}-|\mathbf s|-t_1-1}=\rho^{-t_1-1}$, by $\kappa_{F'}=|\mathbf s|$, and $\int\rho^{-t_1-1}e^{-1/(\rho\gamma^\Psi_1)}d\rho=\Gamma(t_1)(\gamma^\Psi_1)^{t_1}$ turns $\Psi_{\mathcal A'_1/g}^{-t_1}$ into $\Psi_{\mathcal A'_1\setminus g}^{-t_1}$. What remains is the period of the family of the deletions at $\mvec|_R$, balanced exactly when $\kappa_{F'}=|\mathbf s|$. Equally, the dual cell families at $(\mathbf t,\mathbf s)$ turn the statement into (i) of \hb{thm:seriesrepeat} (\hb{prop:faces}(iv), \hb{prop:dualfamily}), the dual of the contractions being the deletions of the duals. The rest is the reference's, for one cell.
\end{proof}

\begin{hproposition}[\hb{prop:chain} (A chain is a power of one operator)]
As stated, for chain cell families at real exponents, (iv) with plain tables and $s_j,t_j>0$.
\end{hproposition}

\begin{proof}
This is the reference proof of \hb{prop:chain} with the laws on $(0,\infty)^{2N-1}$. \emph{(i)} \hb{thm:node}(ii) with the one cell family $F_l$ at $g_{\mathrm{in}}$: its hinge-wall formula, $\mathcal U_j$ having the marks $\cE_{\mathcal U}\sqcup\{g_{\mathrm{in}}\}$ besides its hinge, is $\kappa_{F_{l+1}}=\kappa_{F_l}+\kappa_U$. By \hb{lem:twosumpieces}, $F_l\oplus_gF'$ is on its balance wall exactly when $\kappa_{F_l}+\kappa_{F'}=2\sigma$; with $\kappa_{F_l}=\kappa_{F_0}+l\kappa_U$ this holds for every $l$ exactly when it holds at $l=0$ and $\kappa_U=0$.

\emph{(ii)} \hb{thm:node}(ii) with the one cell family $F_l$ at $g_{\mathrm{in}}$ gives $\lambda_{l+1}=\int\lambda^U_{\mathbf x}\,d\lambda_l(\mathbf x)=\mathcal Q\lambda_l$, the twisted family carrying the mass $\kappa_{F_l}=\kappa_{F_0}$ at $g_{\mathrm{in}}$; by induction $\lambda_l=\mathcal Q^l\lambda_0$. \hb{thm:node}(i) with the core $\binom{\{g_1,g_2\}}1$ gives $P_{F_l\oplus_gF'}=\int h_{F'}\,d\lambda_l$, and $\int H\,d(\mathcal Q\mu)=\int(\mathcal Q^\ast H)\,d\mu$ for $H\ge0$ by Tonelli, so $\int h_{F'}\,d\lambda_l=\int(\mathcal Q^{\ast l}h_{F'})\,d\lambda_0$. The series in $z\ge0$ is summed under the integral by Tonelli.

\emph{(iii)} $\mathcal Q^\ast$ preserves inequalities between nonnegative functions, so $\mathcal Q^\ast H\le MH$ gives $\mathcal Q^{\ast l}H\le M^lH$, and $h_{F'}\le CH$ gives $\mathcal Q^{\ast l}h_{F'}\le CM^lH$; integrate against $\lambda_0$. The lower bound is the same.

\emph{(iv)} For the series core $\binom{\{g_{\mathrm{out}},g_{\mathrm{in}},g'\}}2$ with $F''$ at $g'$, $F_{l+1}$ is the series connection of $F_l$ and $F''$, so \hb{thm:seriesparallel}(i) gives $f_{F_{l+1}}=f_{F''}f_{F_l}$, and the hinge-wall formula with the mass $0$ at $g_{\mathrm{in}}$, the marks $\{g_{\mathrm{in}}\}\sqcup B''$ and the ranks $2+r_{\mathcal A''_j}-1$, gives $\kappa_U=m_{B''}-\sum_j\big(s_jr_{\mathcal A''_j}+t_j(|B''|-r_{\mathcal A''_j})\big)=\kappa_{F''}-|\mathbf s|$. For the parallel core the ranks are $r_{\mathcal A''_j}$, the formula gives $\kappa_{F''}-|\mathbf t|$, and \hb{thm:seriesparallel}(ii) gives the product of the dual transforms.
\end{proof}

\begin{htheorem}[\hb{thm:hingeweight} (The weight of a hinge)]
As stated: (i) for cell families with plain tables, one weight per table; (ii) to (iv) are the reference's, for one cell.
\end{htheorem}

\begin{proof}
\emph{(i)} The bases through $g$ carry the monomials of $K_{\mathcal A_j/g}$ and of $\Psi_{\mathcal A_j/g}$, which are multiplied by $w_j$; the deletions keep their weights (the reference proof of \hb{thm:hingeweight}(i), table by table). So the hinge wall is kept, $\gamma^K_{\mathcal A_j}=K_{\mathcal A_j\setminus g}/K_{\mathcal A_j/g}$ becomes $\gamma^K_{\mathcal A_j}/w_j$ and $\gamma^\Psi_{\mathcal A_j}=\Psi_{\mathcal A_j/g}/\Psi_{\mathcal A_j\setminus g}$ becomes $w_j\gamma^\Psi_{\mathcal A_j}$; the coordinates of the ratio vector become $(w_1/w_j)\gamma^K_{\mathcal A_j}/\gamma^K_{\mathcal A_1}$ and $w_1w_j\gamma^\Psi_{\mathcal A_j}/\gamma^K_{\mathcal A_1}$, and the base form $u^{\mvec-\one}\prod_jK_{\mathcal A_j/g}^{-s_j}\Psi_{\mathcal A_j\setminus g}^{-t_j}(\gamma^K_{\mathcal A_1})^{-\kappa_F}d\omega$ is multiplied by $\prod_jw_j^{-s_j}\cdot w_1^{\kappa_F}$. In the orthant forms, $x_j/\gamma^K_{\mathcal A_j}$ becomes $w_jx_j/\gamma^K_{\mathcal A_j}$ and $y_j\gamma^\Psi_{\mathcal A_j}$ becomes $w_jy_j\gamma^\Psi_{\mathcal A_j}$ with the measure of $f_F$ unchanged, and for $\check f_F$ the measure is multiplied by $\prod_jw_j^{-s_j-t_j}$ while $x_j/\gamma^\Psi_{\mathcal A_j}$ and $y_j\gamma^K_{\mathcal A_j}$ are divided by $w_j$. For one cell, $w_1=w$: the ratio becomes $w^2\varrho_{\mathcal A}$, the base form $w^{\kappa_{\mathcal A}-s}\phi_{\mathcal A}$, and the homogeneity of \hb{def:hinged} with $c=w$ gives $f_{w\mathcal A}(y)=f_F(w,wy)=w^{\kappa_{\mathcal A}-s}f_{\mathcal A}(w^2y)$, and with $c=1/w$, $\check f_{w\mathcal A}(y)=w^{-s-t}\check f_F(1/w,y/w)=w^{-s-\kappa_{\mathcal A}}\check f_{\mathcal A}(y/w^2)$. (ii) to (iv) are the reference's.
\end{proof}

\begin{sremark}
For a family the weights act on the law through an $N$-dimensional subtorus of $(0,\infty)^{2N-1}$, and the Mellin transform in $\mathbf w$ is a restricted Mellin transform of the law against the transform of the room, without the product formula of (ii); (iii) and (iv) concern the one-variable expansions of the node of width one. So (ii) to (iv) are stated for one cell.
\end{sremark}

\begin{hproposition}[\hb{prop:local} (Carrying a calculation through a change)]
As stated, for a family whose tables have a common hinge tree.
\end{hproposition}

\begin{proof}
(i) and (iii) are \hb{prop:work}(iii) with the recursion of \hb{thm:treeperiod}(i): the law at a vertex is recomputed from the new law of the vertex directly below it on the path and the kept laws of the others, and the root by \hb{thm:treeperiod}(ii). (ii) is \hb{thm:hingeweight}(i) for the cell family $F_v$: the weight of the hinge multiplies the base form by $w_1^{\kappa_{F_v}}\prod_jw_j^{-s_j}$ and moves the ratio vector by the stated map, so the law is multiplied by that factor and carried by that map; the cell families below $v$ are unchanged, and above $v$ the change of the law of $F_v$ is a change at a vertex directly below, carried by (i).
\end{proof}

\section{Chapter~\hb{ch:holding}: holds and the lean convolution}\label{sec:ch24}

In this section $F=((\cB_j,\theta_j);s_j,t_j)_{j=1}^N$ is a family as in Definition~\hb{def:hold}: its tables are each one molecule on $\cE$, $n=|\cE|\ge2$, and its exponents are positive, unless a statement says otherwise. Its $2N$ letters are written $X$, running over $K_1,\dots,K_N,\Psi_1,\dots,\Psi_N$. The letter $X$ of table $j$ has the exponent $w_X$, which is $s_j$ or $t_j$, the degree $\deg X$, which is $r_j$ or $r^\ast_j$, the twist $\mathbf x^X$, which is $\mathbf x^K_j$ or $\mathbf x^\Psi_j$, and the Newton polytope $\mathcal N_X$, which is $P(\cB_j)$ or $P(\cB_j^\ast)$; for masses $\mathbf a$ with $|\mathbf a|=w_X\deg X$ the period of its one-letter word $X^{-w_X}$ is written $P_X(\mathbf a)$, which is $P_{\cB_j}(\mathbf a;s_j,0)$ or $P_{\cB_j}(\mathbf a;0,t_j)$. $\mathcal X'$ is the set of the $2N-1$ letters other than $K_1$, and in a split $\mvec=\sum_X\mathbf a^X$ the load of $X$ is $\mathbf a^X$, so that $\mathbf a^{K_j}=\mathbf a_j$ and $\mathbf a^{\Psi_j}=\mathbf b_j$.

Two facts are used throughout. First, the base polytope of a table that is one molecule, and that of its dual table, have dimension $n-1$: if $w\in\mathbb R^{\cE}$ is orthogonal to every difference $\one_T-\one_{T'}$ of bases, then $\ip w{\one_T}$ is constant over the bases and $w$ is constant by Lemma~\hb{lem:circuit}; so these differences span $\one^\perp$, and so do the differences $\one_{\cE\setminus T'}-\one_{\cE\setminus T}$. Second, a point $\sum_kp_kv_k$ with every $p_k>0$ and $\sum_kp_k=1$ lies in the relative interior of the convex hull of the finitely many points $v_k$: for $x=\sum_kq_kv_k$ in the hull and $0<\varepsilon\le\min_kp_k$, the point $\sum_kp_kv_k+\varepsilon\big(\sum_kp_kv_k-x\big)=\sum_k\big((1+\varepsilon)p_k-\varepsilon q_k\big)v_k$ has nonnegative coefficients of sum $1$; so every segment from a point of the hull to the given point extends beyond it within the hull, which characterizes the relative interior of a convex set.

\begin{hdefinition}[\hb{def:posword} (Positive letters; Newton words and sums)]
As stated: the period of a product of powers of positive letters at complex exponents, wherever it converges absolutely, a Newton word being the case of positive exponents.
\end{hdefinition}

\begin{sremark}
The independence of the section is the proof of Lemma~\hb{lem:chart}, which uses of the integrand $f=\alpha^{\mvec-\one}W$ only that it is measurable, homogeneous, $f(\lambda\alpha)=\lambda^{-n}f(\alpha)$ for $\lambda>0$ when $m_{\cE}=D$, and absolutely integrable over one section; here $|f|=\alpha^{\re\mvec-\one}\prod_iL_i^{-\re w_i}$. The family word at any exponents is a product of powers of its letters with $D=\sum_j(s_jr_j+t_jr^\ast_j)$, and on the balance wall its period is that of Theorem~\hb{thm:projective}.
\end{sremark}

\begin{hproposition}[\hb{prop:oneletter} (One positive letter)]
As stated.
\end{hproposition}

\begin{proof}
Put $D=\deg L$ and $t=(t_k)_k$. Every $t_k$ has $|t_k|=D$, and $\mathcal N_L$ has dimension $n-1$ in the hyperplane $|x|=D$, which does not contain $0$; so the $t_k$ span $\mathbb R^{\cE}$.

(i) By Lemma~\hb{lem:scaleshape}, $x=\rho u$ with $\rho>0$ and $u\in\Delta_{\cE}$ gives $x^{\mathbf a-\one}dx=\rho^{|\mathbf a|-1}u^{\mathbf a-\one}d\rho\,d\omega_{\cE}(u)$, and $L(\rho u)=\rho^DL(u)$. For $\re|\mathbf a|>0$ the substitution $\sigma=\rho^DL(u)$ gives
\[
\int_0^\infty\rho^{|\mathbf a|-1}e^{-\rho^DL(u)}\,d\rho=\frac{\Gamma(w)}D\,L(u)^{-w},\qquad w=|\mathbf a|/D .
\]
So $\sG_{\mathbf a;t}((c_k)_k)=\frac{\Gamma(w)}D\int_{\Delta_{\cE}}u^{\mathbf a-\one}L(u)^{-w}d\omega_{\cE}(u)=\frac{\Gamma(w)}D\,P_{L^{-w}}(\mathbf a)$, the last integral being the period over the section $\Delta_{\cE}$ (Definition~\hb{def:posword}). The same computation at $\re\mathbf a$, by Tonelli's theorem, shows that the two integrals converge absolutely together.

(ii) $|u^{\mathbf a-\one}L(u)^{-w}|=u^{\re\mathbf a-\one}L(u)^{-\re w}$, so $P_{L^{-w}}(\mathbf a)$ converges absolutely exactly when the period of the Newton word $L^{-\re w}$, $\re w>0$, converges at $\re\mathbf a$. Its Newton sum is $(\re w)\mathcal N_L$, of dimension $n-1$, and Theorem~\hb{thm:poswordconv} gives the condition $\re\mathbf a\in\mathrm{ri}\big((\re w)\mathcal N_L\big)=(\re w)\,\mathrm{ri}\,\mathcal N_L$.

(iii) For real $\mathbf a$ and $|\boldsymbol\eta|=0$ the exponent $w$ is the same real number at $\mathbf a$ and at $\mathbf a+i\boldsymbol\eta$, so by (i) $P_{L^{-w}}(\mathbf a+i\boldsymbol\eta)=D\,\sG_{\mathbf a+i\boldsymbol\eta;t}((c_k)_k)/\Gamma(w)$. If $P_{L^{-w}}(\mathbf a)$ converges, so does $\sG_{\mathbf a;t}$, and $\mathbf a\in\interior\cone(t)$ (Theorem~\hb{thm:grammadomain}). Theorem~\hb{thm:decay} gives $|\sG_{\mathbf a+i\boldsymbol\eta;t}((c_k)_k)|\le\sG_{\mathbf a;t}(\lambda(c_k)_k)\,e^{-(1-\varepsilon)h_t(\boldsymbol\eta)}$ with $\lambda=\sin\frac{\pi\varepsilon}2$, and $h_t=h_L$ is a norm since the $t_k$ span. Every $t_k$ having total $D$, the substitution $x=\lambda^{-1/D}y$ gives $\sG_{\mathbf a;t}(\lambda(c_k)_k)=\lambda^{-w}\sG_{\mathbf a;t}((c_k)_k)$; multiplying by $D/\Gamma(w)$ gives the second bound. The bound by $P_{L^{-w}}(\mathbf a)$ itself is that of the modulus of the integrand by its value at the real parts.
\end{proof}

\begin{sremark}
At the letter $K_{\mathcal H}$ of a table that is one molecule, with the support $(\one_T)_T$, the coefficients $\boldsymbol\theta$ and $\mathcal N_{K_{\mathcal H}}=P(\mathcal H)$, items (i), (ii) and (iii) are Proposition~\hb{prop:fullleans}(i), (ii) and (iv) at the full basis lean, and at $\Psi_{\mathcal H}$ at the full cobasis lean. At a simplicial letter, (i) with Corollary~\hb{cor:square} at the square configuration $t$ is Proposition~\hb{prop:simplicial}: $\sG_{\mathbf a;t}((c_k)_k)=|\det t|^{-1}\prod_k\Gamma(b_k)c_k^{-b_k}$ with $\mathbf b=t^{-1}\mathbf a$, and $|\mathbf b|=w$ since $\one^{\mathsf T}t=D\one^{\mathsf T}$.
\end{sremark}

\begin{hlemma}[\hb{lem:concat} (Concatenation of Grammas)]
As stated, for any number of parts.
\end{hlemma}

\begin{proof}
For $i\ge2$ the configuration $t^{(i)}$ spans $\mathbb R^{\cE}$, so by Theorem~\hb{thm:decay} the Gramma $\sG_{\mathbf w;t^{(i)}}(\boldsymbol\xi^{(i)})$ on $\mathbf c_i+i\mathbb R^{\cE}$ is at most a constant times $e^{-(1-\varepsilon)h_{t^{(i)}}(\im\mathbf w)}$, $h_{t^{(i)}}$ a norm, and is integrable there. In $\mathbf y=\log x$ it is the Fourier transform of the continuous integrable function $e^{\ip{\mathbf c_i}{\mathbf y}}\exp\big(-\sum_k\xi^{(i)}_ke^{\ip{t^{(i)}_k}{\mathbf y}}\big)$ (Definition~\hb{def:gramma}), so Fourier inversion gives, for every $x\in\mathbb R^{\cE}_{>0}$,
\[
\exp\Big(-\sum_k\xi^{(i)}_kx^{t^{(i)}_k}\Big)=\frac1{(2\pi i)^n}\int_{\mathbf c_i+i\mathbb R^{\cE}}\sG_{\mathbf w_i;t^{(i)}}(\boldsymbol\xi^{(i)})\,x^{-\mathbf w_i}\,d\mathbf w_i .
\]
Multiply these for $i=2,\dots,p$ by $x^{\nu-\one}\exp\big(-\sum_k\xi^{(1)}_kx^{t^{(1)}_k}\big)$ and integrate over $x$. The modulus of the integrand is $x^{\re\nu-\sum_{i\ge2}\mathbf c_i-\one}\exp\big(-\sum_k\xi^{(1)}_kx^{t^{(1)}_k}\big)\prod_{i\ge2}|\sG_{\mathbf w_i;t^{(i)}}(\boldsymbol\xi^{(i)})|$, whose integral is $\sG_{\re\nu-\sum_{i\ge2}\mathbf c_i;t^{(1)}}(\boldsymbol\xi^{(1)})\prod_{i\ge2}\int|\sG_{\mathbf w_i;t^{(i)}}(\boldsymbol\xi^{(i)})|\,|d\mathbf w_i|$, finite by Theorem~\hb{thm:grammadomain} since $\re\nu-\sum_{i\ge2}\mathbf c_i\in\interior\cone(t^{(1)})$. So Fubini's theorem exchanges the integrals; the integral over $x$ is $\sG_{\nu-\sum_{i\ge2}\mathbf w_i;t^{(1)}}(\boldsymbol\xi^{(1)})$, and the left side is the Gramma of the concatenation (Definition~\hb{def:gramma}).
\end{proof}

\begin{hdefinition}[\hb{def:hold} (Holding)]
As stated: the hold of a family, the Mellin transform in the twists of its letters other than $K_1$.
\end{hdefinition}

\begin{sremark}
That the bases of every table span $\mathbb R^{\cE}$ is the first fact above. By the torus relation of Definition~\hb{def:twist}, a factor $\mu>0$ on every coordinate of one twist $\mathbf x^X$ multiplies the period by $\mu^{-w_X\deg X}$, and a twist of every letter by one vector $\boldsymbol\lambda$ multiplies it by $\boldsymbol\lambda^{-\mvec}$; so the Mellin transform in all the twists of all the letters at all the marks diverges, and the hold keeps the twist of $K_1$ at $\one$ and, in Lemma~\hb{lem:holdall}, the twists at one mark.
\end{sremark}

\begin{hdefinition}[\hb{def:fullleanfns} (Loads; admissible points)]
As stated: the loads and the admissible points of a family, an admissible point of one table being named by its $\Psi$-load.
\end{hdefinition}

\begin{hproposition}[\hb{prop:fullleans} (The periods at the two full leans)]
As stated; (iii) is stated for every family, and the closing sentence is new.
\end{hproposition}

\begin{proof}
Items (i), (ii) and (iv) concern one table $\mathcal H$ of $F$, one molecule, and are the reference, proof of prop:fullleans. The closing sentence is Proposition~\hb{prop:oneletter}, proved above; at the letter $K_{\mathcal H}$, with the support $(\one_T)_T$, the coefficients $\boldsymbol\theta$ and the Newton polytope $P(\mathcal H)$ of dimension $n-1$, it gives back (i), (ii) and (iv) at the full basis lean, and at $\Psi_{\mathcal H}$ at the full cobasis lean.

(iii) Let $\alpha$ be a gauge. By Theorem~\hb{thm:innerMGF}, $R_j(\alpha)=\oE_{T\sim\mathbb P_{j,\alpha}}[\one_T]=\sum_{T\in\cB_j}\mathbb P_{j,\alpha}(T)\,\one_T$ with $\mathbb P_{j,\alpha}(T)=\theta_{j,T}\alpha^T/K_j(\alpha)>0$ for every basis, and $R^\ast_{j,e}(\alpha)=\mathbb P^\ast_{j,\alpha}(e\notin T)$, so $R^\ast_j(\alpha)=\sum_T\mathbb P^\ast_{j,\alpha}(T)\,\one_{\cE\setminus T}$ with $\mathbb P^\ast_{j,\alpha}(T)=\theta_{j,T}\alpha^{\cE\setminus T}/\Psi_j(\alpha)>0$. By the second fact above $R_j(\alpha)\in\mathrm{ri}\,P(\cB_j)$ and $R^\ast_j(\alpha)\in\mathrm{ri}\,P(\cB_j^\ast)$, the two polytopes being the convex hulls of the $\one_T$ and of the $\one_{\cE\setminus T}$.

Let $\mvec$ be real on the balance wall with every proper wall positive. A mark of a table that is one molecule on at least two marks is neither a loop nor a coloop, so $d_j(\{e\})=d^\ast_j(\{e\})=0$ and $m_e=\kappa_{\{e\}}>0$ for every $e$ (Definition~\hb{def:family}). A nonempty proper $S$ with $c_j(S)=0$ would have $|T\cap S|$ constant over the bases of table $j$, and so would be a union of molecules of it by Lemma~\hb{lem:separator}, which is impossible, table $j$ being one molecule; so $c_j(S)>0$ in every table, every table being active, and $F$ is connected. By Theorem~\hb{thm:saddle}(ii) and (iii) a native gauge $\alpha^\star$ exists, the target being $\mvec$ on the balance wall (Definition~\hb{def:native}) and the walls at the target positive. Its usage is the target, $\sum_j\big(s_jR_j(\alpha^\star)+t_jR^\ast_j(\alpha^\star)\big)=\mvec$: the native loads are a split of $\mvec$ among the letters, with $|\mathbf a^\star_j|=s_jr_j$ and $|\mathbf b^\star_j|=t_jr^\ast_j$, every $\one_T$ having $r_j$ marks and every $\one_{\cE\setminus T}$ having $r^\ast_j$. By the preceding paragraph $\mathbf a^\star_j\in s_j\,\mathrm{ri}\,P(\cB_j)$ and $\mathbf b^\star_j\in t_j\,\mathrm{ri}\,P(\cB_j^\ast)$, so the native loads are admissible (Definition~\hb{def:fullleanfns}).
\end{proof}

\begin{hlemma}[\hb{lem:holdall} (Holding every mark separates the letters)]
As stated, for every family.
\end{hlemma}

\begin{proof}
Let the data be real. By Theorem~\hb{thm:projective} the period of $F[\mathbf x]$ is the integral over the section $\{\alpha_{e_0}=1\}$,
\[
P_{F[\mathbf x]}(\mvec;\mathbf s,\mathbf t)=\int_{(0,\infty)^{\cE\setminus e_0}}\alpha^{\mvec-\one}K_1(\alpha)^{-s_1}\prod_{X\in\mathcal X'}X(\mathbf x^X\circ\alpha)^{-w_X}\,d\alpha',\qquad\alpha=(\alpha',1),
\]
the twists at $e_0$ being $1$ and the twist of $K_1$ being $\one$. Insert this into the hold and integrate first over the twists, by Tonelli's theorem, every function being positive. For fixed $\alpha'$ substitute $\mathbf v^X=\mathbf x^X\circ\alpha'$ on $\cE\setminus e_0$ for every $X\in\mathcal X'$: then $(\mathbf x^X)^{\boldsymbol\beta^X-\one}d\mathbf x^X=\alpha'^{-\boldsymbol\beta^X}(\mathbf v^X)^{\boldsymbol\beta^X-\one}d\mathbf v^X$ and $X(\mathbf x^X\circ\alpha)=X(\mathbf v^X,1)$, so the integral over $\mathbf x^X$ is $\alpha'^{-\boldsymbol\beta^X}$ times
\[
\int_{(0,\infty)^{\cE\setminus e_0}}(\mathbf v^X)^{\boldsymbol\beta^X-\one}X(\mathbf v^X,1)^{-w_X}\,d\mathbf v^X=P_X(\mathbf a^X),
\]
the integral over the same section of the form of the one-letter word $X^{-w_X}$ at the extension $\mathbf a^X$ of $\boldsymbol\beta^X$ of total $w_X\deg X$ (Theorem~\hb{thm:projective}); here $\boldsymbol\beta^{K_j}=\boldsymbol\beta^K_j$ and $\boldsymbol\beta^{\Psi_j}=\boldsymbol\beta^\Psi_j$, and $\mathbf a^X$ is the $\mathbf a_j$ or $\mathbf b_j$ of the statement. What is left is $\int\alpha'^{\mvec'-\one-\sum_X\boldsymbol\beta^X}K_1(\alpha',1)^{-s_1}d\alpha'$, $\mvec'=\mvec|_{\cE\setminus e_0}$, which is $P_{K_1}(\mathbf a_1)$ for $\mathbf a_1=\mvec-\sum_{X\in\mathcal X'}\mathbf a^X$: its coordinates off $e_0$ are $\mvec'-\sum_X\boldsymbol\beta^X$, and its total is $m_{\cE}-\sum_{X\in\mathcal X'}w_X\deg X=r_1s_1$ on the balance wall. This is the identity in $[0,\infty]$. For complex $\mvec$ and $\boldsymbol\beta$ the moduli of all the integrands are their values at the real parts, where the integrals converge by assumption, and Fubini's theorem replaces Tonelli's.
\end{proof}

\begin{slemma}\label{lem:ch24-pins}
Let shares $\mathbf v\in\mathbb C^M$ be pinned by $P\mathbf v=\boldsymbol\pi$, with $q$ pins, as in Definition~\hb{def:fibre}.
\begin{enumerate}[label=(\roman*)]
\item The fibre integral does not change when the shares are replaced by $A^{-1}\mathbf v$, the integrand by $f\circ A$ and the pins by $PA$, for a real $M\times M$ matrix $A$ with $|\det A|=1$, nor when the pins are replaced by $CP\mathbf v=C\boldsymbol\pi$, for a real $q\times q$ matrix $C$ with $|\det C|=1$.
\item Let $\mathbf v=(\mathbf u,\mathbf w)$ with $\mathbf w\in\mathbb C^q$, so that $P=(P_{\mathbf u},P_{\mathbf w})$, and let $|\det P_{\mathbf w}|=1$. Then
\[
\int_{P\mathbf v=\boldsymbol\pi}f=\frac1{(2\pi i)^{M-q}}\int_{\mathbf c_{\mathbf u}+i\mathbb R^{M-q}}f\big(\mathbf u,P_{\mathbf w}^{-1}(\boldsymbol\pi-P_{\mathbf u}\mathbf u)\big)\,d\mathbf u,
\]
$\mathbf c_{\mathbf u}$ the $\mathbf u$-part of the admissible real point, the two sides converging absolutely together.
\item Let $\mathbf v=(\mathbf v_0,\mathbf v_1,\dots,\mathbf v_k)$ and let the pins be the union of pins $P_l\mathbf v_l=\boldsymbol\pi_l$, each on one block. If the fibre integral of $f$ converges absolutely, it is the integral over the fibre $P_0\mathbf v_0=\boldsymbol\pi_0$ of the fibre integral over $P_1\mathbf v_1=\boldsymbol\pi_1,\dots,P_k\mathbf v_k=\boldsymbol\pi_k$ at fixed $\mathbf v_0$, the inner one converging absolutely for almost every point of the outer imaginary translate.
\end{enumerate}
\end{slemma}

\begin{proof}
For $\mathbf y_0\in\mathbb R^q$ let $\mu_{P,\mathbf y_0}=(\det PP^{\mathsf T})^{-1/2}\mathrm{vol}_{M-q}$ on $\{\mathbf y\in\mathbb R^M:\ P\mathbf y=\mathbf y_0\}$, so that the fibre integral is $(2\pi)^{q-M}\int f(\mathbf c+i\mathbf y)\,d\mu_{P,\im\boldsymbol\pi}(\mathbf y)$. By the coarea formula for the linear map $P$, whose Jacobian is $(\det PP^{\mathsf T})^{1/2}$, these measures disintegrate Lebesgue measure, $\int_{\mathbb R^M}\varphi\,d\mathbf y=\int_{\mathbb R^q}\big(\int\varphi\,d\mu_{P,\mathbf y_0}\big)d\mathbf y_0$ for integrable $\varphi$, and $\mathbf y_0\mapsto\int\varphi\,d\mu_{P,\mathbf y_0}$ is continuous for continuous $\varphi$ of compact support; these two properties determine the family $(\mu_{P,\mathbf y_0})_{\mathbf y_0}$.

(i) The substitution $\mathbf y=A\mathbf y'$ preserves Lebesgue measure and carries the disintegration along $PA$ to that along $P$; the substitution $\mathbf y_0=C^{-1}\mathbf y'_0$ preserves Lebesgue measure on $\mathbb R^q$ and carries the disintegration along $CP$ to that along $P$, with $\mu_{CP,C\mathbf y_0}=\mu_{P,\mathbf y_0}$. The admissible real point and the factor $(2\pi)^{q-M}$ do not change.

(ii) The map $\mathbf y\mapsto(\mathbf y_{\mathbf u},P\mathbf y)$ has the block triangular matrix $\begin{pmatrix}I&0\\P_{\mathbf u}&P_{\mathbf w}\end{pmatrix}$, of determinant of modulus $1$; so $d\mathbf y=d\mathbf y_{\mathbf u}\,d\mathbf y_0$, and $\mu_{P,\mathbf y_0}$ is the image of Lebesgue measure on $\mathbb R^{M-q}$ under $\mathbf y_{\mathbf u}\mapsto\big(\mathbf y_{\mathbf u},P_{\mathbf w}^{-1}(\mathbf y_0-P_{\mathbf u}\mathbf y_{\mathbf u})\big)$. With $\mathbf u=\mathbf c_{\mathbf u}+i\mathbf y_{\mathbf u}$ and $d\mathbf u=i^{M-q}d\mathbf y_{\mathbf u}$ this is the display; both sides are one integral against one measure.

(iii) The pin matrix is block diagonal, so the disintegration along it is the product of the disintegrations along the blocks, $\mu_{P,(\mathbf y_{0,l})_l}=\bigotimes_l\mu_{P_l,\mathbf y_{0,l}}$, and the factor $(2\pi)^{q-M}$ is the product of the factors of the blocks. Fubini's theorem gives the iterated integral.
\end{proof}

\begin{sremark}
In Theorem~\hb{thm:leanconv}(ii) the shares are the loads of the $2N$ letters, $M=2Nn$, with the $q=n+2N-1$ pins $\sum_X\mathbf a^X=\mvec$ and $|\mathbf a^X|=w_X\deg X$ for $X\in\mathcal X'$, and the eliminated coordinates are $\mathbf a_1$ and the loads $a^X_{e_0}$, $X\in\mathcal X'$. With the rows ordered as the mass pins at $e\neq e_0$, the mass pin at $e_0$ and the totals, and the columns as $a_{1,e}$ for $e\neq e_0$, $a_{1,e_0}$ and the $a^X_{e_0}$, the block $P_{\mathbf w}$ is
\[
\begin{pmatrix}I_{n-1}&0&0\\0&1&\one^{\mathsf T}\\0&0&I_{2N-1}\end{pmatrix},
\]
of determinant $1$, and Lemma~\ref{lem:ch24-pins}(ii) gives the coordinate form stated there, with the factor $(2\pi i)^{-(2N-1)(n-1)}$. With the further pins on $G$ of Theorem~\hb{thm:pinnedhold} the loads $a^X_e$, $X\in\mathcal X'$, $e\in G$, are eliminated as well, each by its own pin with the coefficient $1$, and the block stays unit upper triangular. The function \texttt{eliminated\_block\_det} of \texttt{checks/bookkeeping.py} computes this determinant exactly for $n\le4$ and up to six letters.
\end{sremark}

\begin{slemma}\label{lem:ch24-euler}
Let $A,B>0$, $w>0$ and $c\in\mathbb C$. If $0<\re c<w$,
\[
\int_0^\infty x^{c-1}(A+xB)^{-w}\,dx=\Gst21{c,\ w-c}{w}A^{c-w}B^{-c},
\]
absolutely convergent; for real $c\notin(0,w)$ the integral is $+\infty$.
\end{slemma}

\begin{proof}
The substitution $x=(A/B)y$ gives $A^{c-w}B^{-c}\int_0^\infty y^{c-1}(1+y)^{-w}dy$. The last integral is the integral over the section $\{u_2=1\}$ of the form of the word $(u_1+u_2)^{-w}$ on two marks at the masses $(c,w-c)$: the family of the list $\binom{\{1,2\}}1$ with unit weights at the exponents $(w,0)$, on its balance wall, whose walls are $c$ and $w-c$. For $0<\re c<w$ its period converges absolutely (Theorem~\hb{thm:projective}), and by Lemma~\hb{lem:chart} the integral equals that over $\Delta_{\{1,2\}}$, where $u_1+u_2=1$; by Theorem~\hb{thm:dirichlet} it is $\int_{\Delta}u^{(c,w-c)-\one}d\omega=\Gst21{c,\ w-c}{w}$. For real $c\le0$ the integrand is at least $2^{-w}y^{c-1}$ on $(0,1)$, and for real $c\ge w$ at least $2^{-w}y^{c-w-1}$ on $(1,\infty)$; both integrals diverge.
\end{proof}

\begin{htheorem}[\hb{thm:leanconv} (The lean convolution)]
As stated, for every family; the closing sentence takes the active letters of a family with nonnegative exponents.
\end{htheorem}

\begin{proof}
(i) Let $\mathbf c$ be admissible. By Proposition~\hb{prop:fullleans}(ii), $\mathbf c^K_j\in s_j\,\mathrm{ri}\,P(\cB_j)$ is the condition $(\mathbf c^K_j)_S>s_jd_j(S)$ for every nonempty proper $S$, and likewise $(\mathbf c^\Psi_j)_S>t_jd^\ast_j(S)$; summing over the letters, $\re\kappa_S=\re m_S-\sum_j\big(s_jd_j(S)+t_jd^\ast_j(S)\big)>0$ (Definition~\hb{def:family}). Conversely, if every proper wall of $\re\mvec$ is positive, the native loads of $\re\mvec$ are admissible for $(\mvec;\mathbf s,\mathbf t)$ by Proposition~\hb{prop:fullleans}(iii), admissibility involving $\re\mvec$ only. The admissible points are cut out by linear equations and the convex conditions $\mathbf c^X\in w_X\,\mathrm{ri}\,\mathcal N_X$, so they form a convex set.

(ii) Apply Theorem~\hb{thm:factors} to the family word $W=\prod_XX^{-w_X}$ split into the $2N$ factors $W_X=X^{-w_X}$, of degrees $w_X\deg X$. The Newton polytope of $W_X$ is $w_X\mathcal N_X$, of dimension $n-1$ by the first fact above, so its relative interior is $w_X\,\mathrm{ri}\,\mathcal N_X$, and the points $(\mathbf c_i)_i$ of Theorem~\hb{thm:factors} are exactly the admissible points. With the twists $\mathbf x^X$ of the factors, the left side of Theorem~\hb{thm:factors} is the integral over the section $\{\alpha_{e_0}=1\}$ of the word of $F[\mathbf x]$, whose letters are $X(\mathbf x^X\circ\alpha)$ (Definition~\hb{def:twist}), that is $P_{F[\mathbf x]}(\mvec;\mathbf s,\mathbf t)$ (Theorem~\hb{thm:projective}); the periods of the factors are the $P_X(\mathbf a^X)$. This is the display, absolutely convergent, of dimension $(2N-1)(n-1)$; $K_1$ is taken as the last factor of Theorem~\hb{thm:factors}, whose total pin is omitted, and that pin holds by the others and the balance wall. The coordinate form is the remark after Lemma~\ref{lem:ch24-pins}.

(iii) For a real point $\mathbf a$ at which $P_X(\mathbf a)$ converges, $|P_X(\mathbf a+i\mathbf y)|\le P_X(\mathbf a)$, the modulus of the integrand of a period being its value at the real parts, and $|(\mathbf x^X)^{-\mathbf a^X}|=(\mathbf x^X)^{-\re\mathbf a^X}$. On the imaginary translate through $\mathbf c$ the imaginary part $\mathbf y^X$ of every load has $|\mathbf y^X|=0$, by the pins, and Proposition~\hb{prop:fullleans}(iv) bounds $|P_X(\mathbf c^X+i\mathbf y^X)|$ by $\big(\sin\frac{\pi\varepsilon}2\big)^{-w_X}e^{-(1-\varepsilon)h_X(\mathbf y^X)}P_X(\mathbf c^X)$, $h_X$ the norm of that proposition for the letter $X$.

The closing sentence is Theorem~\hb{thm:factors} with one factor for each active letter, the inactive letters being the constant $1$; any factor may be taken last, and the walls and the admissible points are read with the active letters alone.
\end{proof}

\begin{hcorollary}[\hb{cor:retaintable} (What a family retains)]
As stated, for every family.
\end{hcorollary}

\begin{proof}
By Theorem~\hb{thm:leanconv}(ii) and its closing sentence every period of $F$ and of its twists, at exponents $s_j,t_j\ge0$ not all zero and at masses on the balance wall with every proper wall positive, is a fibre integral of products of the periods $P_X$ of its active letters at their loads, times the twists. By Proposition~\hb{prop:fullleans}(i), $P_{\cB_j}(\mathbf a;s_j,0)=r_j\,\sG_{\mathbf a;(\one_T)_T}(\boldsymbol\theta_j)/\Gamma(s_j)$ for $s_j>0$, and by the same item and Cremona $P_{\cB_j}(\mathbf b;0,t_j)=P_{\cB_j}(t_j\one-\mathbf b;t_j,0)=r_j\,\sG_{t_j\one-\mathbf b;(\one_T)_T}(\boldsymbol\theta_j)/\Gamma(t_j)$ for $t_j>0$. The loads run over $|\mathbf a|=r_js_j$ and $|\mathbf b|=r^\ast_jt_j$, that is $|t_j\one-\mathbf b|=nt_j-r^\ast_jt_j=r_jt_j$. The converse is Proposition~\hb{prop:fullleans}(i) for table $j$.
\end{proof}

\begin{htheorem}[\hb{thm:pinnedhold} (Pinned holds)]
As stated, for every family.
\end{htheorem}

\begin{proof}
Put $d=(2N-1)|G|$. In the coordinates of Theorem~\hb{thm:leanconv}(ii), the loads of the letters of $\mathcal X'$ at the marks other than $e_0$, split the coordinates into their values $\mathbf z_G$ on $G$ and the rest $\mathbf z'$, and let $\mathcal I(\mathbf z_G,\mathbf z')$ be the integrand at $\mathbf x=\one$, the product of the periods $P_X$ at the loads that $(\mathbf z_G,\mathbf z')$ and the pins determine. Let $\mathbf c$ be the admissible point of the statement, $\boldsymbol\beta=\mathbf c_G$, and for $\mathbf y\in\mathbb R^d$ put
\[
\Phi(\mathbf y)=\frac1{(2\pi i)^{(2N-1)(n-1-|G|)}}\int\mathcal I(\mathbf c_G+i\mathbf y,\mathbf z')\,d\mathbf z',
\]
over the imaginary translate through $\mathbf c$ in $\mathbf z'$. By Lemma~\ref{lem:ch24-pins}(ii) and the remark after it, $\Phi(\boldsymbol\eta)$ is the fibre integral of the statement at $\boldsymbol\beta+i\boldsymbol\eta$.
\begin{enumerate}[label=(\alph*)]
\item By Theorem~\hb{thm:leanconv}(iii), $|\mathcal I|$ on the imaginary translate is at most a constant times $\prod_{X\in\mathcal X'}e^{-(1-\varepsilon)h_X(\mathbf y^X)}$, $\mathbf y^X$ the imaginary part of the load of $X$, a vector of total $0$ determined by its coordinates off $e_0$; since $h_X$ is a norm on the vectors of total $0$, this bound is integrable in all the imaginary parts. So $\Phi$ is defined everywhere, continuous by dominated convergence, and integrable, with $\int|\Phi|\le(2\pi)^{(2N-1)(n-1-|G|)}\int|\mathcal I|$.
\item For a twist at the marks of $G$ alone, $\mathbf x^X=e^{\mathbf u^X}$ on $G$ with $\mathbf u^X\in\mathbb R^G$ for $X\in\mathcal X'$ and $\mathbf x^{K_1}=\one$, the factor $\prod_X(\mathbf x^X)^{-\mathbf a^X}$ of Theorem~\hb{thm:leanconv}(ii) is $e^{-\ip{\mathbf c_G+i\mathbf y}{\mathbf u}}$ on the translate, $\mathbf u=(\mathbf u^X)_X$. By Fubini's theorem in $(\mathbf y,\mathbf z')$, $d\mathbf z_G=i^d\,d\mathbf y$,
\[
e^{\ip{\mathbf c_G}{\mathbf u}}P_{F[e^{\mathbf u}]}(\mvec;\mathbf s,\mathbf t)=\frac1{(2\pi)^d}\int_{\mathbb R^d}\Phi(\mathbf y)\,e^{-i\ip{\mathbf y}{\mathbf u}}\,d\mathbf y=:\widehat\Phi(\mathbf u).
\]
\item Let first $\mvec$ be real. Then $\widehat\Phi\ge0$, and for $\varepsilon>0$, by Fubini's theorem and the Gaussian integral,
\[
\int_{\mathbb R^d}\widehat\Phi(\mathbf u)\,e^{-\varepsilon|\mathbf u|^2/2}\,d\mathbf u=\int_{\mathbb R^d}\Phi(\mathbf y)\,\frac{e^{-|\mathbf y|^2/(2\varepsilon)}}{(2\pi\varepsilon)^{d/2}}\,d\mathbf y\longrightarrow\Phi(\mathbf 0)\qquad(\varepsilon\downarrow0),
\]
$\Phi$ being integrable and continuous at $\mathbf 0$. By monotone convergence $\int\widehat\Phi\,d\mathbf u=\Phi(\mathbf 0)<\infty$, and with $d\mathbf x/\mathbf x=d\mathbf u$ this integral is $\operatorname{Hold}_G(\mathbf c_G)$: the statement at $\boldsymbol\eta=\mathbf 0$.
\item For complex $\mvec$, $|P_{F[\mathbf x]}(\mvec;\mathbf s,\mathbf t)|\le P_{F[\mathbf x]}(\re\mvec;\mathbf s,\mathbf t)$, and $\mathbf c$ is admissible for $\re\mvec$ as well; by (c) at $\re\mvec$ the hold converges absolutely at every $\mathbf c_G+i\boldsymbol\eta$. Then $\widehat\Phi$ of (b) is integrable, $\Phi$ is integrable and continuous, and Fourier inversion gives, at every $\boldsymbol\eta$,
\[
\operatorname{Hold}_G(\mathbf c_G+i\boldsymbol\eta)=\int_{\mathbb R^d}e^{\ip{\mathbf c_G+i\boldsymbol\eta}{\mathbf u}}P_{F[e^{\mathbf u}]}(\mvec;\mathbf s,\mathbf t)\,d\mathbf u=\int_{\mathbb R^d}\widehat\Phi(\mathbf u)\,e^{i\ip{\boldsymbol\eta}{\mathbf u}}\,d\mathbf u=\Phi(\boldsymbol\eta).
\]
\end{enumerate}
At $G=\cE\setminus e_0$ the fibre has dimension $0$, and its value is the product of Lemma~\hb{lem:holdall}.
\end{proof}

\begin{hproposition}[\hb{prop:holdone} (Holding one mark)]
As stated, for every family.
\end{hproposition}

\begin{proof}
Let the data be real. Take the section $\{h=1\}$ of Theorem~\hb{thm:projective} with $h(\alpha)=\sum_{e\neq g}\alpha_e$: on it $\alpha=(\alpha',\alpha_g)$ with $\alpha'\in\Delta_{\cE\setminus g}$ and $\alpha_g\in(0,\infty)$, and the measure is $d\omega_{\cE\setminus g}(\alpha')\,d\alpha_g$. By Definitions~\hb{def:twist} and~\hb{def:substitution} the twist at $g$ acts on the letters as
\[
K_j(\mathbf x@g\circ\alpha)=K_{\cB_j\setminus g}(\alpha')+x^K_j\alpha_gK_{\cB_j/g}(\alpha'),\qquad\Psi_j(\mathbf x@g\circ\alpha)=x^\Psi_j\alpha_g\Psi_{\cB_j\setminus g}(\alpha')+\Psi_{\cB_j/g}(\alpha'),
\]
the minors carrying the weights of table $j$, all four nonzero since $g$ is neither a loop nor a coloop. Insert the period into the hold and integrate first over the twists of the letters of $\mathcal X'$, by Tonelli's theorem. By Lemma~\ref{lem:ch24-euler}, for $j\ge2$, respectively for every $j$,
\begin{gather*}
\int_0^\infty x^{a_j-1}\big(K_{\cB_j\setminus g}+x\alpha_gK_{\cB_j/g}\big)^{-s_j}dx=\Gst21{a_j,\ s_j-a_j}{s_j}K_{\cB_j\setminus g}^{a_j-s_j}\big(\alpha_gK_{\cB_j/g}\big)^{-a_j},\\
\int_0^\infty x^{b_j-1}\big(\Psi_{\cB_j/g}+x\alpha_g\Psi_{\cB_j\setminus g}\big)^{-t_j}dx=\Gst21{b_j,\ t_j-b_j}{t_j}\Psi_{\cB_j/g}^{b_j-t_j}\big(\alpha_g\Psi_{\cB_j\setminus g}\big)^{-b_j},
\end{gather*}
finite exactly when $0<a_j<s_j$, respectively $0<b_j<t_j$, and $+\infty$ otherwise. The power of $\alpha_g$ left is $\alpha_g^{m_g-1-\sum_{j\ge2}a_j-\sum_jb_j}=\alpha_g^{a_1-1}$, and by the same lemma
\[
\int_0^\infty\alpha_g^{a_1-1}\big(K_{\cB_1\setminus g}+\alpha_gK_{\cB_1/g}\big)^{-s_1}d\alpha_g=\Gst21{a_1,\ s_1-a_1}{s_1}K_{\cB_1\setminus g}^{a_1-s_1}K_{\cB_1/g}^{-a_1},
\]
finite exactly when $0<a_1<s_1$. What remains is
\[
\int_{\Delta_{\cE\setminus g}}\alpha'^{\mvec|_{\cE\setminus g}-\one}\prod_jK_{\cB_j\setminus g}^{-(s_j-a_j)}\,\Psi_{\cB_j\setminus g}^{-b_j}\,K_{\cB_j/g}^{-a_j}\,\Psi_{\cB_j/g}^{-(t_j-b_j)}\,d\omega_{\cE\setminus g},
\]
the period of the children family $F_{\{g\}}\big((s_j-a_j,a_j),(b_j,t_j-b_j)\big)_j$ at $\mvec|_{\cE\setminus g}$, which is on its balance wall (the paragraph before the proposition); it is finite exactly when every proper wall of it is positive (Theorem~\hb{thm:projective}). The integrands being positive, the hold is the product of these factors in $[0,\infty]$ and is finite exactly when all of them are. For complex loads and masses whose real parts are such data the moduli are dominated by the real computation, Lemma~\ref{lem:ch24-euler} holds in the strip, and Fubini's theorem gives the same identity. The sentence on admissible points is Theorem~\hb{thm:pinnedhold} with $G=\{g\}$.

For the limits, let the letter $Y$ have the exponent $u>m_g>0$ and let the twists at $g$ of the other letters decrease to $0$. Each other letter decreases pointwise to one of its minors, $K_{\cB_j\setminus g}$ or $\Psi_{\cB_j/g}$, and its power with a negative exponent increases; by monotone convergence the period increases to the integral of the limit. In the limit only $Y$ contains $\alpha_g$, as $A+\alpha_gB$ with $A$ and $B$ its two minors, and by Lemma~\ref{lem:ch24-euler} the integral over $\alpha_g$ of $\alpha_g^{m_g-1}(A+\alpha_gB)^{-u}$ is $\Gst21{m_g,\ u-m_g}{u}A^{m_g-u}B^{-m_g}$. What is left is the period of the children family at the shares of the statement: for $Y=K_j$, $A=K_{\cB_j\setminus g}$ carries the exponent $s_j-m_g$ and $B=K_{\cB_j/g}$ the exponent $m_g$; for $Y=\Psi_j$, $A=\Psi_{\cB_j/g}$ carries $t_j-m_g$ and $B=\Psi_{\cB_j\setminus g}$ carries $m_g$; every other $K_{j'}$ is $K_{\cB_{j'}\setminus g}$ at $s_{j'}$ and every other $\Psi_{j'}$ is $\Psi_{\cB_{j'}/g}$ at $t_{j'}$.
\end{proof}

\begin{sremark}
For one table a twist at $g$ in the cobasis letter is $x=x^\Psi_1$ with $x^K_1=1$ (Definition~\hb{def:twist}), and $x\downarrow0$ is the limit with the letter $K$ chosen. With the letter $\Psi$ chosen the twist $y$ of $K$ decreases to $0$ with that of $\Psi$ at $1$; by the torus relation of Definition~\hb{def:twist} at the mark $g$, $P_{\mathcal H[(y,1)@g]}=y^{-m_g}P_{\mathcal H[(1,1/y)@g]}=x^{m_g}P_{\mathcal H[x@g]}$ with $x=1/y$. So the two limits are those of $P_{\mathcal H[x@g]}$ as $x\downarrow0$ and of $x^{m_g}P_{\mathcal H[x@g]}$ as $x\uparrow\infty$, as the paragraph after the proposition reads them.
\end{sremark}

\begin{hlemma}[\hb{lem:hingesplit} (The Mellin transform of a law is the period of a children family)]
As stated, for every cell family.
\end{hlemma}

\begin{proof}
By Definitions~\hb{def:hinged} and~\hb{def:substitution}, with $\sum_j(z_{j,K}+z_{j,\Psi})=\kappa_F$ the power of $\gamma^K_{\mathcal A_1}$ in $\boldsymbol\varrho_F^{-\mathbf z}\,d\phi_F$ is $-\kappa_F+\sum_{j\ge2}z_{j,K}+\sum_jz_{j,\Psi}=-z_{1,K}$. With $\gamma^K_{\mathcal A_j}=K_{\mathcal A_j\setminus g}/K_{\mathcal A_j/g}$ and $\gamma^\Psi_{\mathcal A_j}=\Psi_{\mathcal A_j/g}/\Psi_{\mathcal A_j\setminus g}$,
\begin{align*}
\boldsymbol\varrho_F^{-\mathbf z}\,d\phi_F&=u^{\mvec-\one}\prod_jK_{\mathcal A_j/g}^{-s_j}\Psi_{\mathcal A_j\setminus g}^{-t_j}\big(\gamma^K_{\mathcal A_j}\big)^{-z_{j,K}}\big(\gamma^\Psi_{\mathcal A_j}\big)^{-z_{j,\Psi}}\,d\omega_B\\
&=u^{\mvec-\one}\prod_jK_{\mathcal A_j\setminus g}^{-z_{j,K}}\,\Psi_{\mathcal A_j\setminus g}^{-(t_j-z_{j,\Psi})}\,K_{\mathcal A_j/g}^{-(s_j-z_{j,K})}\,\Psi_{\mathcal A_j/g}^{-z_{j,\Psi}}\,d\omega_B .
\end{align*}
This is the form of the children family of $\{g\}$ at the shares of the statement, its tables $\mathcal A_j\setminus g$, of rank $r_{\mathcal A_j}$ and corank $|B|-r_{\mathcal A_j}$, at $(z_{j,K},t_j-z_{j,\Psi})$ and $\mathcal A_j/g$, of rank $r_{\mathcal A_j}-1$ and corank $|B|-r_{\mathcal A_j}+1$, at $(s_j-z_{j,K},z_{j,\Psi})$. Its degree is
\begin{multline*}
\sum_j\Big(r_{\mathcal A_j}z_{j,K}+(|B|-r_{\mathcal A_j})(t_j-z_{j,\Psi})+(r_{\mathcal A_j}-1)(s_j-z_{j,K})+(|B|-r_{\mathcal A_j}+1)z_{j,\Psi}\Big)\\=\kappa_F+\sum_j\big(s_j(r_{\mathcal A_j}-1)+t_j(|B|-r_{\mathcal A_j})\big)=m_B,
\end{multline*}
so the children family is on its balance wall, and the integral over $\Delta_B$ is its period (Theorem~\hb{thm:projective}), in $[0,\infty]$ for real $\mathbf z$. This is the first identity.

With the mass $\sum_j(s_j+t_j)-\kappa_F$ on $g$ the family of the tables $\mathcal A_j$ is on its balance wall (Lemma~\hb{lem:twosumpieces}). Apply Proposition~\hb{prop:holdone} to it at the mark $g$, a hinge and so neither a loop nor a coloop, with $a_j=s_j-z_{j,K}$ for $j\ge2$ and $b_j=t_j-z_{j,\Psi}$. Then
\[
a_1=\sum_j(s_j+t_j)-\kappa_F-\sum_{j\ge2}(s_j-z_{j,K})-\sum_j(t_j-z_{j,\Psi})=s_1-z_{1,K},
\]
the children family of Proposition~\hb{prop:holdone} at the shares $(s_j-a_j,a_j)=(z_{j,K},s_j-z_{j,K})$ and $(b_j,t_j-b_j)=(t_j-z_{j,\Psi},z_{j,\Psi})$ is that of the first identity, its Beta corners are $\Gst21{s_j-z_{j,K},\ z_{j,K}}{s_j}$ and $\Gst21{t_j-z_{j,\Psi},\ z_{j,\Psi}}{t_j}$, and its hold is the left side of the second identity. At the real part of $\mathbf z$ the conditions of Proposition~\hb{prop:holdone} hold: $0<a_j<s_j$ and $0<b_j<t_j$ are $0<\re z_{j,K}<s_j$ and $0<\re z_{j,\Psi}<t_j$, and the walls of the children family are positive since its period, $\int\boldsymbol\varrho_F^{-\re\mathbf z}d\phi_F$ by the first identity, is finite. So Proposition~\hb{prop:holdone} gives the second identity for complex $\mathbf z$ with this real part.
\end{proof}

\begin{sremark}
The balance of the children family, the balance of the family of the tables $\mathcal A_j$ with the mass $\sum_j(s_j+t_j)-\kappa_F$ on $g$, and the identity $a_1=s_1-z_{1,K}$ are computed exactly by \texttt{check\_hingesplit\_balance} in \texttt{checks/bookkeeping.py}, and the balance of the children family of Proposition~\hb{prop:holdone} by \texttt{check\_holdone\_balance}.
\end{sremark}

\section{Chapter~\hb{ch:leantree}: periods at the full leans and networks along a hinge tree}\label{sec:ch25}

The notation of Section~\ref{sec:ch24} is kept. When table $j$ of a family carries a hinge tree $\mathfrak T_j$ as in Definition~\hb{def:branches}, its marker loads and node loads are read in table $j$ (Definition~\hb{def:markerloads}), at its loads $(\mathbf a_j,\mathbf b_j)$ and its exponents $(s_j,t_j)$.

\begin{hlemma}[\hb{lem:twosummol} (Substitutions of molecules)]
As stated, for any number of cells.
\end{hlemma}

\begin{proof}
The case of two cells is the reference, proof of lem:twosummol. By Proposition~\hb{prop:substitution}(i) the substitution is the gluing of the cells to $\mathcal H$ along its marks one at a time: $\cB_0=\mathcal H$, and $\cB_i$ is the gluing of $\mathcal A_i$ and $\cB_{i-1}$ along $h_i$, the marks $h_{i+1},\dots,h_q$ left in place. By Definition~\hb{def:substitution} and Proposition~\hb{prop:substitution}(iii) this gluing is the two-sum of $\mathcal A_i$ and $\cB_{i-1}$, the latter read as a cell with the hinge $h_i$. This $h_i$ is neither a loop nor a coloop of $\cB_{i-1}$: $\mathcal H$ has bases through $h_i$ and bases avoiding it, and each extends to a basis of $\cB_{i-1}$ by bases of the hinge minors of the cells glued before, which exist since their hinges are neither loops nor coloops. By the case of two cells $\cB_i$ is one molecule exactly when $\mathcal A_i$ and $\cB_{i-1}$ are, and by induction $\cB_q$ is one molecule exactly when $\mathcal H$ and all the $\mathcal A_i$ are.
\end{proof}

\begin{hcorollary}[\hb{cor:nodeconv} (Convergence is decided at the nodes)]
As stated, for every family.
\end{hcorollary}

\begin{proof}
Every table is one molecule by Lemma~\hb{lem:twosummol} applied along its hinge tree. By Theorem~\hb{thm:leantree} in table $j$, at real loads, $P_{\cB_j}(\mathbf a_j;s_j,0)$ is finite exactly when every marker load $k_{j;v,\mathfrak e}$ lies in $(0,s_j)$ and the period of every node at its full basis lean is finite at its node loads. By Proposition~\hb{prop:fullleans}(ii) at the node, one molecule, the latter is $(\mathbf a_{j,v})_S>s_jd_{j,v}(S)$ for every nonempty proper $S$ of its marks and markers; at $S=\{g\}$ and at the complement of a marker $g$ these inequalities are $k>0$ and $k<s_j$ for the marker load $k$ at $g$, since $g$ is neither a loop nor a coloop of the node and $|\mathbf a_{j,v}|=r_{j,v}s_j$ (Lemma~\hb{lem:nodeloads}(ii)). Likewise at the full cobasis lean. So the node inequalities hold at a split exactly when every load is admissible for its letter (Proposition~\hb{prop:fullleans}(ii) in table $j$), that is when the split is an admissible point, and by Theorem~\hb{thm:leanconv}(i) such a point exists exactly when every proper wall is positive; the native loads are one. The unknowns are the loads of the $2N$ letters at $n$ marks less $n+2N-1$ independent pins, $(2N-1)(n-1)$ of them, and the node $v$ of table $j$ contributes $2(2^{|\cE_{j,v}|+|G_{j,v}|}-2)$ inequalities.
\end{proof}

\begin{hproposition}[\hb{prop:spleans} (Series and parallel nodes at their full leans)]
As stated; the last sentence, which reads the four closed forms off Proposition~\hb{prop:simplicial}, is new.
\end{hproposition}

\begin{proof}
The closed forms are the reference, proof of prop:spleans. The letter $K$ of $\binom{\mathcal X}1$ is $\sum_x\theta_x\alpha_x$, and so is the letter $\Psi$ of $\binom{\mathcal X}q$, whose basis $\mathcal X\setminus x$ has the cobasis $\{x\}$: both are simplicial with the configuration $I=(\mathbf e_x)_x$, of degree $1$, $|\det I|=1$ and the shares $\mathbf b=\mathbf a$, and Proposition~\hb{prop:simplicial} gives $\Gst{q+1}1{\mathbf a}{|\mathbf a|}\prod_x\theta_x^{-a_x}$, the forms of $P_{\binom{\mathcal X}1}(\mathbf a;s,0)$ and $P_{\binom{\mathcal X}q}(\mathbf b;0,t)$. The letter $K$ of $\binom{\mathcal X}q$ and the letter $\Psi$ of $\binom{\mathcal X}1$ are $\sum_x\theta_x\alpha^{\mathcal X\setminus x}$, simplicial with the configuration $J-I=(\one-\mathbf e_x)_x$ of degree $q$, $J$ the matrix of ones. $J-I$ has the eigenvalues $q$, once, and $-1$, $q$ times, so $|\det(J-I)|=q$, and $(J-I)^{-1}=\frac1qJ-I$, so the shares of $\mathbf a$ with $|\mathbf a|=qs$ are $\frac{|\mathbf a|}q\one-\mathbf a=(s-a_x)_x$, of sum $s$. Proposition~\hb{prop:simplicial} gives $\frac qq\Gst{q+1}1{(s-a_x)_x}{s}\prod_x\theta_x^{a_x-s}$, the form of $P_{\binom{\mathcal X}q}(\mathbf a;s,0)$, and in the same way that of $P_{\binom{\mathcal X}1}(\mathbf b;0,t)$ with $|\mathbf b|=qt$. The criteria of convergence agree, the positivity of the real parts of the shares being that of the arguments of the Gammas.
\end{proof}

\begin{sremark}
The function \texttt{check\_spleans\_is\_simplicial} of \texttt{checks/bookkeeping.py} verifies exactly, at $q=4$, the shares $s-a_x$ and the prefactor $q/|\det(J-I)|=1$.
\end{sremark}

\begin{hcorollary}[\hb{cor:spfulllean} (Series-parallel tables at their full leans)]
As stated, with any weights on the bases of the series and parallel nodes.
\end{hcorollary}

\begin{proof}
By Theorem~\hb{thm:leantree} the period of $\cB$ at its full basis lean is the product over the nodes of their periods at their node loads over the Beta corners $\Gst21{k_{v,\mathfrak e},\ k_{w,\mathfrak e}}{s}$ of the edges, the weights of $\cB$ being the products of those of the nodes (Definition~\hb{def:branches}). Proposition~\hb{prop:spleans} gives each node: $\Gst{n_v}1{(s-(\mathbf a_v)_x)_x}{s}\prod_x\theta_{v,x}^{(\mathbf a_v)_x-s}$ at a series node and $\Gst{n_v}1{((\mathbf a_v)_x)_x}{s}\prod_x\theta_{v,x}^{-(\mathbf a_v)_x}$ at a parallel node, $n_v=|\cE_v|+|G_v|$. The Gammas, the factors $\Gamma(s)$ and the corners combine as in the reference, proof of cor:spfulllean: a tree has one node more than edges, which leaves $1/\Gamma(s)$; a mark $e$ contributes $\Gamma(s-a_e)$ at a series node and $\Gamma(a_e)$ at a parallel node; at an edge $\mathfrak e=\{v,w\}$, where $k_{v,\mathfrak e}+k_{w,\mathfrak e}=s$ (Lemma~\hb{lem:nodeloads}(i)), the two markers contribute $\Gamma(s-k_{v,\mathfrak e})\Gamma(s-k_{w,\mathfrak e})=\Gamma(k_{w,\mathfrak e})\Gamma(k_{v,\mathfrak e})$ between two series nodes and $\Gamma(k_{v,\mathfrak e})\Gamma(k_{w,\mathfrak e})$ between two parallel nodes, which the corner cancels, and $\Gamma(s-k_{v,\mathfrak e})\Gamma(k_{w,\mathfrak e})$ when $v$ is a series and $w$ a parallel node, which the corner leaves as $\Gamma(s-k_{\mathfrak e})/\Gamma(k_{\mathfrak e})$, $k_{\mathfrak e}=k_{v,\mathfrak e}$. The monomials of the nodes stand as they are. The full cobasis lean is the same with the second pair of forms of Proposition~\hb{prop:spleans}: $\Gamma(b_e)$ at series and $\Gamma(t-b_e)$ at parallel nodes, and at an edge with the series end $v$ the quotient $\Gamma(l_{v,\mathfrak e})\Gamma(t-l_{w,\mathfrak e})/\big(\Gamma(l_{v,\mathfrak e})\Gamma(l_{w,\mathfrak e})\big)=\Gamma(l_{\mathfrak e})/\Gamma(t-l_{\mathfrak e})$. The criterion of finiteness does not see the weights, which enter through monomials, positive and finite at real loads; it is the reference's.
\end{proof}

\begin{hcorollary}[\hb{cor:primelean} (Periods along a hinge tree read by one hinge value)]
As stated: for a family whose tables have a common hinge tree, at real data, wherever the active letters read every subtree cell by one hinge value.
\end{hcorollary}

\begin{proof}
Induction on the number of edges of $\mathfrak T$. With one vertex there is nothing to prove.

\emph{Removing a leaf.} A tree with an edge has two leaves, so some leaf $v$ differs from $o$; let $p$ be its parent and $\mathfrak e=\mathfrak e_v$. The subtree cell $\mathcal A_{v,j}$ is the node $\mathcal H_{v,j}$ with the hinge $g_{v,\mathfrak e}$ (\hb{def:hingetree}). The two-sums along a tree may be taken in any order (\hb{lem:nodesum}, \hb{prop:substitution}(i),(ii)), so table $j$ of $F$ is the substitution $\mathcal H'_j\langle\mathcal A_{v,j}\rangle$ of the cell $\mathcal A_{v,j}$ at the mark $g_{p,\mathfrak e}$ of the table $\mathcal H'_j$ obtained from the other nodes along $\mathfrak T-v$, every other mark of $\mathcal H'_j$ left in place. The cell on $\cE_v$ is read by the active letters through $\gamma_v$, and a mark left in place through $\alpha_h$, so \hb{cor:fullleancore} applies; the factors of the marks left in place cancel (its step~(5)), and with $\kappa_{\mathfrak e}=\kappa_{\cE_v}$,
\[
P_F(\mvec;\mathbf s,\mathbf t)\,\Gst21{\kappa_{\mathfrak e},\ 2\sigma-\kappa_{\mathfrak e}}{2\sigma}=P_{F'}(\mvec';\mathbf s,\mathbf t)\,P_{F_{\mathcal H_v}}(\mvec_v;\mathbf s,\mathbf t),
\]
$F'$ the family of the tables $\mathcal H'_j$ at the masses $\mvec'$ equal to $\mvec$ on its marks and to $\kappa_{\cE_v}$ at $g_{p,\mathfrak e}$, and $\mvec_v$ equal to $\mvec|_{\cE_v}$ and to $2\sigma-\kappa_{\cE_v}=\kappa_{\cE\setminus\cE_v}$ at $g_{v,\mathfrak e}$; every factor is on its balance wall with every proper wall positive.

\emph{The walls of $F'$.} By \hb{thm:substwalls}, a set $S'$ of marks of $F'$ containing $g_{p,\mathfrak e}$ has the wall of $(S'\setminus\{g_{p,\mathfrak e}\})\sqcup\cE_v$ in $F$, the cell entering with its whole set $\cE_v$, of wall $\kappa_{\cE_v}$, and the marker in $D_j$ and $D^\ast_j$; a set avoiding $g_{p,\mathfrak e}$ has the same wall in both. So the walls across the edges of $\mathfrak T-v$ are those of $F$.

\emph{The hypothesis passes to $F'$.} $F'$ has the common hinge tree $\mathfrak T-v$, with the same nodes and $g_{p,\mathfrak e}$ a mark of the node $p$. Its subtree cells are those of $F$ except at the vertices $w$ on the path from $p$ to $o$, where $\mathcal A_{w,j}=\mathcal A'_{w,j}\langle\mathcal A_{v,j}\rangle$, the cell at the mark $g_{p,\mathfrak e}$ (\hb{prop:substitution}(ii)). By \hb{thm:substletters} the hinge minors of $\mathcal A_{w,j}$ are those of $\mathcal A'_{w,j}$ with $\alpha_{g_{p,\mathfrak e}}$ replaced by $\gamma^K_{\mathcal A_{v,j}}$ in the letters $K$ and by $\gamma^\Psi_{\mathcal A_{v,j}}$ in the letters $\Psi$, times factors of $\mathcal A_{v,j}$ that cancel in the quotients, so
\[
\gamma^K_{\mathcal A_{w,j}}(\alpha)=\gamma^K_{\mathcal A'_{w,j}}\big(\alpha',\gamma^K_{\mathcal A_{v,j}}(\alpha_{\cE_v})\big),\qquad \gamma^\Psi_{\mathcal A_{w,j}}(\alpha)=\gamma^\Psi_{\mathcal A'_{w,j}}\big(\alpha',\gamma^\Psi_{\mathcal A_{v,j}}(\alpha_{\cE_v})\big),
\]
$\alpha'$ the variables of the other marks. For the active letters both inner hinge values are $\gamma_v(\alpha_{\cE_v})$, and $\gamma_v$, positive and homogeneous of degree one, takes every value in $(0,\infty)$. So for every value $c>0$ at $g_{p,\mathfrak e}$ the hinge values of $\mathcal A'_{w,j}$ by the active letters at $(\alpha',c)$ all equal $\gamma_w(\alpha)$ for any $\alpha$ with $\gamma_v(\alpha_{\cE_v})=c$: the active letters read $\mathcal A'_{w,j}$ by one hinge value, and the induction hypothesis applies to $F'$. Its product formula and the display give that of $F$, with the masses on the markers the walls across the edges.

\emph{One table at a full lean.} One letter is active, and the hypothesis is empty. At $(s,0)$ on the balance wall $|\mvec|=rs$, so the loads $\mathbf a=\mvec$ are admissible in \hb{thm:leantree}; the marker load $k_{v,\mathfrak e}=sr^{\mathfrak e}_v-m_{\cE^{\mathfrak e}_v}$ is $s-\kappa_{\cE^{\mathfrak e}_v}$, the wall of the branch at $v$ being $m_{\cE^{\mathfrak e}_v}-s(r^{\mathfrak e}_v-1)$, and the two walls across an edge add to $s$; so the node loads are the masses $\mvec_v$ and the two identities are the same (\hb{lem:nodeloads}). At $(0,t)$ the same holds with $\mathbf b$ and $t$. The series and parallel nodes are \hb{prop:spleans}.

\emph{Tables glued along a tree.} The glued table has the hinge tree whose nodes are the tables $\mathcal A_v$ and whose markers are the glued marks; the glued marks of connected tables are neither loops nor coloops, and a two-sum has the product weights (\hb{lem:twosumpieces}). The masses of the glued marks are positive, every singleton being a proper wall of its table, so every $\kappa_{\mathfrak e}$ lies in $(0,t)$. Each table $\mathcal A_v$ is balanced, $m_{\cE_{\mathcal A_v}}=tr^\ast_{\mathcal A_v}$, and the glued table, of corank $\sum_vr^\ast_{\mathcal A_v}-(k-1)$ and total mass $\sum_vtr^\ast_{\mathcal A_v}-(k-1)t$, is on its balance wall. Removing a leaf as above, with the first display of \hb{cor:fullleancore} in $[0,\infty]$ and its step~(2), the cell at a leaf $v$ glued along its mark $g$ has the wall $t-m_g$ and its cell family is $\mathcal A_v$ at its own masses, so the period of the glued table is $\prod_vP_{\mathcal A_v}/\prod_{\mathfrak e}\Gst21{\kappa_{\mathfrak e},\ t-\kappa_{\mathfrak e}}{t}$ in $[0,\infty]$; it is finite, and the glued table, balanced, has every proper wall positive (\hb{thm:projective}).

\emph{A chain.} By \hb{lem:twosumpieces} the chain $F^{\rhd k}$ of \hb{thm:chainoperator} has the hinge tree whose nodes are the copies, a path, the glued marks its markers; by the propagation of the hypothesis above, its subtree cells are read by the active letters through one hinge value when the copy is, and the walls across its edges are $m_{g_{\mathrm{in}}}$ and $m_{g_{\mathrm{out}}}$, so with every proper wall positive the display gives the product $P_F^k/\Gst21{m_{g_{\mathrm{in}}},\ m_{g_{\mathrm{out}}}}{2\sigma}^{k-1}$.
\end{proof}

\begin{htheorem}[\hb{thm:blockagg} (Blocks read by one hinge value aggregate exactly)]
As stated: in (ii) blocks read through $e_\ell/e_{\ell-1}$ as well as blocks read through their totals, and the active letters at nonzero exponents of either sign.
\end{htheorem}

\begin{proof}
\emph{(i) and (iii).} These are the reference proof of \hb{thm:blockagg}(i),(iii), the letters with positive exponent replaced by the active letters: at a parallel block $K_{\mathcal A/g}=1$ and $\gamma^K_{\mathcal A}=e_1$, at a series block $\Psi_{\mathcal A\setminus g}=1$ and $\gamma^\Psi_{\mathcal A}=e_1$ (\hb{cor:blocksclosures}(i)), so by \hb{thm:substletters} every active letter is the letter of its core read at the block totals, and the totals of $U\sim\Dir_{\mvec}$ have the law $\Dir_{(m_{B_i})_i}$ (\hb{def:dirlaw}). The sign of an exponent does not enter.

\emph{(ii).} When a block has $\ell_i\ge2$ marks a table with both letters active would need a block that is both parallel and series; so at such a block each table has at most one active letter, and the inactive letters, of exponent $0$, enter neither the walls nor the kernels. By \hb{cor:blocksclosures}(i), for $i\in I$ the active hinge values are $\gamma^K=e_1$ at the parallel blocks and $\gamma^\Psi=e_1$ at the series blocks, and for $i\notin I$ they are $\gamma^K=e_{\ell_i}/e_{\ell_i-1}$ at the series blocks and $\gamma^\Psi=e_{\ell_i}/e_{\ell_i-1}$ at the parallel blocks. So \hb{thm:pinned} applies with all constants $1$ and $\gamma_i=e_1$ for $i\in I$, $\gamma_i=e_{\ell_i}/e_{\ell_i-1}$ for $i\notin I$:
\[
P_F(\mvec;\mathbf s,\mathbf t)=P_{F_{\mathcal H}}\big((\kappa_{B_i})_i;\mathbf s,\mathbf t\big)\prod_iS_i,\qquad S_i=\int_{\Delta_{B_i}}\gamma_i^{-\kappa_{B_i}}\,d\beta_i ,
\]
in $[0,\infty]$. By \hb{thm:substwalls}, $\kappa_{B_i}=m_{B_i}-\sum_j\big(s_j(r_{\mathcal A_{ij}}-1)+t_j(\ell_i-r_{\mathcal A_{ij}})\big)$ with the ranks $1$ of a parallel block and $\ell_i$ of a series block.

For $i\in I$ the active letters see ranks that make every term vanish, so $\kappa_{B_i}=m_{B_i}$; the density of the cell kernel is $u^{\mvec|_{B_i}-\one}$, the active minors $K_{\mathcal A/g}$ and $\Psi_{\mathcal A\setminus g}$ being $1$; and $e_1=1$ on the simplex, so $S_i=\int_{\Delta_{B_i}}u^{\mvec|_{B_i}-\one}d\omega$, which is $\Gst{\ell_i}1{\mvec|_{B_i}}{m_{B_i}}$ when every $m_x$ is positive and $+\infty$ otherwise, by the shape integral (\hb{thm:dirmoments}).

For $i\notin I$ each active letter contributes its exponent times $\ell_i-1$, so $\kappa_{B_i}=m_{B_i}-(\ell_i-1)2\sigma$. The active minors are $K_{\mathcal A/g}=e_{\ell_i-1}$ at a series block and $\Psi_{\mathcal A\setminus g}=e_{\ell_i-1}$ at a parallel block, so the density of the cell kernel is $u^{\mvec|_{B_i}-\one}e_{\ell_i-1}(u)^{-2\sigma}$, and with $e_{\ell_i}=\prod_{x\in B_i}u_x$,
\[
S_i=\int_{\Delta_{B_i}}\Big(\frac{e_{\ell_i}}{e_{\ell_i-1}}\Big)^{-\kappa_{B_i}}u^{\mvec|_{B_i}-\one}e_{\ell_i-1}^{-2\sigma}\,d\omega=\int_{\Delta_{B_i}}u^{\mvec|_{B_i}-\kappa_{B_i}\one-\one}\,e_{\ell_i-1}(u)^{-(2\sigma-\kappa_{B_i})}\,d\omega .
\]
This is the period of the series node $\binom{B_i}{\ell_i-1}$ with unit weights, whose letter $K$ is $e_{\ell_i-1}$, at the full basis lean $(2\sigma-\kappa_{B_i},0)$ and the masses $m_x-\kappa_{B_i}$; they add to $m_{B_i}-\ell_i\kappa_{B_i}=(\ell_i-1)(2\sigma-\kappa_{B_i})$, its balance wall. By \hb{prop:spleans} it is $\Gst{\ell_i}1{(2\sigma-\kappa_{B_i}-m_x+\kappa_{B_i})_x}{2\sigma-\kappa_{B_i}}=\Gst{\ell_i}1{(2\sigma-m_x)_x}{2\sigma-\kappa_{B_i}}$, convergent exactly when every $2\sigma-m_x$ is positive, and $+\infty$ otherwise, the integrand being positive. At $\ell_i=1$ both kinds give $\kappa_{B_i}=m_x$ and $S_i=1$.

So the display holds in $[0,\infty]$ with the stated reading of the stacked Gammas, and when these are finite the two sides are finite together.

\emph{The last sentence.} At $\ell\ge2$ the functions $e_1$ and $e_\ell/e_{\ell-1}$ are not proportional on the simplex, their quotient $e_1e_{\ell-1}/e_\ell=(\sum_xu_x)(\sum_xu_x^{-1})$ being unbounded; so \hb{thm:pinned} does not apply to a block read both ways, and \hb{thm:corereads} keeps it as one mark with its kernel and both hinge values.
\end{proof}

\begin{sremark}
Checks~7a to~7d of S21-23\_work/checks.py are the case $\ell=2$ of the blocks with $i\notin I$, with one table at $(s_1,0)$ and one at $(0,t_2)$: $S=\Gamma(2\sigma-m_1)\Gamma(2\sigma-m_2)/\Gamma(4\sigma-m_B)$, and the period of the cell family is $\Gst21{\kappa,\ 2\sigma-\kappa}{2\sigma}S$.
\end{sremark}

\begin{htheorem}[\hb{thm:network} (The network)]
As stated, for every family.
\end{htheorem}

\begin{proof}
Every table is one molecule by Lemma~\hb{lem:twosummol} along its hinge tree, and every proper wall is positive, so Theorem~\hb{thm:leanconv}(ii) writes $P_{F[\mathbf x]}$ as the fibre integral of $\prod_jP_{\cB_j}(\mathbf a_j;s_j,0)P_{\cB_j}(\mathbf b_j;0,t_j)$ times the twists, on the imaginary translate through an admissible point $\mathbf c$. At $\mathbf c$ every $P_{\cB_j}(\mathbf c^K_j;s_j,0)$ and $P_{\cB_j}(\mathbf c^\Psi_j;0,t_j)$ is finite (Proposition~\hb{prop:fullleans}(ii)), so by Theorem~\hb{thm:leantree} in table $j$ every marker load lies in $(0,s_j)$, respectively $(0,t_j)$, every node period is finite at its node loads, and both factorizations of Theorem~\hb{thm:leantree} hold on the tube over these real loads, which contains the translate. Inserting them gives the integrand of the statement. The native loads are admissible by Proposition~\hb{prop:fullleans}(iii). The node periods are functions of the node loads alone, the exponent of the one active letter of a node being its total load divided by its degree (Definition~\hb{def:fullleanfns}).
\end{proof}

\begin{hcorollary}[\hb{cor:spnetwork} (Series-parallel periods)]
As stated, for every family and any weights on the bases of the nodes.
\end{hcorollary}

\begin{proof}
Insert Corollary~\hb{cor:spfulllean} for both full leans of every table into Theorem~\hb{thm:network} and multiply by $\prod_j\Gamma(s_j)\Gamma(t_j)$. At a series node the monomials of the two leans multiply to $\theta_{j,v,x}^{(\mathbf a_{j,v})_x-s_j}\theta_{j,v,x}^{-(\mathbf b_{j,v})_x}$, at a parallel node to $\theta_{j,v,x}^{-(\mathbf a_{j,v})_x}\theta_{j,v,x}^{(\mathbf b_{j,v})_x-t_j}$. At a real point at which every Gamma above the line has an argument of positive real part all the periods of the tables at their full leans are finite (Corollary~\hb{cor:spfulllean}), so the point is admissible (Proposition~\hb{prop:fullleans}(ii)) and Theorem~\hb{thm:network} applies on it. The dimension is that of Theorem~\hb{thm:network}, and the integrand does not involve the bases of the tables.
\end{proof}

\begin{hproposition}[\hb{prop:poolpair} (Pooling a pair)]
As stated, with any weights on the bases of the node.
\end{hproposition}

\begin{proof}
For one table the pins are $\mathbf a+\mathbf b=\mvec$ and $|\mathbf b|=r^\ast t$, the twist $\mathbf x$ is on the cobasis letter, and the integrand of Corollary~\hb{cor:spnetwork} is a function of $\mathbf b$, with $\Gamma(s-a_e)=\Gamma(c_e+b_e)$ at a mark $e$ of a series node. The marks $e$ and $f$ of one node $v$ lie on one side of every edge, so every marker load, an affine function of sums of loads over branches (Definition~\hb{def:markerloads}), depends on $b_e$ and $b_f$ through $b_{\{e,f\}}$ alone; so do the Gammas of the other marks, the monomials of the other nodes and those of the markers of $v$. What is left at a series node is $\Gamma(c_e+b_e)\Gamma(b_e)\Gamma(c_f+b_f)\Gamma(b_f)$ times
\[
\theta_{v,e}^{a_e-b_e-s}\theta_{v,f}^{a_f-b_f-s}x_e^{-b_e}x_f^{-b_f}=\theta_{v,e}^{m_e-s}\theta_{v,f}^{m_f-s}\big(x_e\theta_{v,e}^2\big)^{-b_e}\big(x_f\theta_{v,f}^2\big)^{-b_f},
\]
a function of $b_{\{e,f\}}$ when $x_e\theta_{v,e}^2=x_f\theta_{v,f}^2$; at a parallel node it is $\Gamma(m_e-b_e)\Gamma(t-b_e)\Gamma(m_f-b_f)\Gamma(t-b_f)$ times $\theta_{v,e}^{-m_e-t}\theta_{v,f}^{-m_f-t}\big(x_e\theta_{v,e}^{-2}\big)^{-b_e}\big(x_f\theta_{v,f}^{-2}\big)^{-b_f}$, from the monomials $\theta_{v,x}^{b_x-a_x-t}$, a function of $b_{\{e,f\}}$ when $x_e\theta_{v,e}^{-2}=x_f\theta_{v,f}^{-2}$. The rest is the reference, proof of prop:poolpair, with $x_e\theta_{v,e}^{\pm2}$ in place of $x_e$: at fixed $b_{\{e,f\}}$ the integral over $b_e$ is the integral of four Gammas displayed in the statement, and disjoint pairs occupy disjoint coordinates.
\end{proof}

\begin{htheorem}[\hb{thm:holdnode} (Nodes in hold form)]
As stated, for every substituted family.
\end{htheorem}

\begin{proof}
The cores $\mathcal H_j$ and the cells $\mathcal A_{ij}$ are molecules by Lemma~\hb{lem:twosummol}, and the markers are neither loops nor coloops of their cores and cells. Table $j$ carries the hinge tree whose nodes are $\mathcal H_j$ and the $\mathcal A_{ij}$, each cell read as a core on $B_i\sqcup\{g_i\}$ and joined to $\mathcal H_j$ by an edge $\mathfrak e_i$ (Definition~\hb{def:hingetree}), the table being obtained by these two-sums (Proposition~\hb{prop:substitution}(i)). Each $B_i$ is nonempty, a hinge being neither a loop nor a coloop; fix $e_i\in B_i$.

(i) \emph{Step 1: the network.} Let $\mathbf c$ be an admissible point of $F$. By Theorem~\hb{thm:network} with these trees, $P_{F[\mathbf x]}(\mvec;\mathbf s,\mathbf t)$ is the fibre integral, over the loads of $F$ with the pins $\sum_j(\mathbf a_j+\mathbf b_j)=\mvec$, $|\mathbf a_j|=r_js_j$ ($j\ge2$) and $|\mathbf b_j|=r^\ast_jt_j$, on the translate through $\mathbf c$, of
\[
\prod_j\frac{P_{\mathcal H_j}(\mathbf a_{j,\mathcal H};s_j,0)P_{\mathcal H_j}(\mathbf b_{j,\mathcal H};0,t_j)\prod_iP_{\mathcal A_{ij}}(\mathbf a_{j,i};s_j,0)P_{\mathcal A_{ij}}(\mathbf b_{j,i};0,t_j)}{\prod_i\Gst21{a_{j,g_i},\ s_j-a_{j,g_i}}{s_j}\Gst21{b_{j,g_i},\ t_j-b_{j,g_i}}{t_j}}\,(\mathbf x^K_j)^{-\mathbf a_j|_{\cE_{\mathcal H}}}(\mathbf x^\Psi_j)^{-\mathbf b_j|_{\cE_{\mathcal H}}}.
\]
Here, by Definition~\hb{def:markerloads} with the branch $B_i$ on the side of the cell, of rank $r_{\mathcal A_{ij}}$ and corank $|B_i|+1-r_{\mathcal A_{ij}}$, and by Lemma~\hb{lem:nodeloads}(i), the marker loads of the core at $g_i$ are
\[
a_{j,g_i}=a_{j,B_i}-s_j\big(r_{\mathcal A_{ij}}-1\big),\qquad b_{j,g_i}=b_{j,B_i}-t_j\big(|B_i|-r_{\mathcal A_{ij}}\big),
\]
those of the cell at $g_i$ are $s_j-a_{j,g_i}$ and $t_j-b_{j,g_i}$, the node loads of the core are $\mathbf a_{j,\mathcal H}=(\mathbf a_j|_{\cE_{\mathcal H}},(a_{j,g_i})_i)$ and $\mathbf b_{j,\mathcal H}=(\mathbf b_j|_{\cE_{\mathcal H}},(b_{j,g_i})_i)$, and those of the cell $\mathcal A_{ij}$ are $\mathbf a_{j,i}=(\mathbf a_j|_{B_i},s_j-a_{j,g_i})$ and $\mathbf b_{j,i}=(\mathbf b_j|_{B_i},t_j-b_{j,g_i})$. The twists are $1$ off $\cE_{\mathcal H}$.

\emph{Step 2: the marker loads as coordinates.} In every table $j$ and for every $i$ replace the coordinate $a_{j,e_i}$ by $a_{j,g_i}$ and $b_{j,e_i}$ by $b_{j,g_i}$. Each new coordinate is the old one plus other old coordinates plus a constant, so the change of coordinates is unit triangular, of determinant $1$, and by Lemma~\ref{lem:ch24-pins}(i) the fibre integral does not change. The new coordinates are the \emph{core coordinates}, the loads $(\mathbf a_{j,\mathcal H},\mathbf b_{j,\mathcal H})_j$ of the cores on $\cE_{\mathcal H}\sqcup\{g_1,\dots,g_k\}$, and for each $i$ the \emph{coordinates of $B_i$}, the loads $a_{j,e}$ and $b_{j,e}$ at the marks $e\in B_i\setminus e_i$, $j=1,\dots,N$.

\emph{Step 3: the core masses.} Replace the mass pin at $e_i$ by the sum $C_i$ of the mass pins over $B_i$, a row operation with the coefficient $1$ on the replaced pin, of determinant $1$ (Lemma~\ref{lem:ch24-pins}(i)). In the new coordinates
\[
C_i:\qquad\sum_j\big(a_{j,g_i}+b_{j,g_i}\big)=m_{B_i}-\sum_j\Big(s_j\big(r_{\mathcal A_{ij}}-1\big)+t_j\big(|B_i|-r_{\mathcal A_{ij}}\big)\Big)=\kappa_{F_i}
\]
by Definition~\hb{def:hinged}: the mass of the cores at $g_i$ is the hinge wall of the cell family at $g_i$.

\emph{Step 4: the totals, and the split of the pins.} By Proposition~\hb{prop:substitution}(i), $r_j=r_{\mathcal H_j}+\sum_i(r_{\mathcal A_{ij}}-1)$, and so $r^\ast_j=n-r_j=r^\ast_{\mathcal H_j}+\sum_i(|B_i|-r_{\mathcal A_{ij}})$ with $r^\ast_{\mathcal H_j}=|\cE_{\mathcal H}|+k-r_{\mathcal H_j}$ and $n=|\cE_{\mathcal H}|+\sum_i|B_i|$. In the new coordinates therefore
\[
|\mathbf a_j|=\big|\mathbf a_j|_{\cE_{\mathcal H}}\big|+\sum_i\Big(a_{j,g_i}+s_j\big(r_{\mathcal A_{ij}}-1\big)\Big),\qquad\text{so}\qquad|\mathbf a_j|=r_js_j\iff|\mathbf a_{j,\mathcal H}|=r_{\mathcal H_j}s_j,
\]
and likewise $|\mathbf b_j|=r^\ast_jt_j\iff|\mathbf b_{j,\mathcal H}|=r^\ast_{\mathcal H_j}t_j$. The pins are now the \emph{core pins}, namely the mass pins at $\cE_{\mathcal H}$, the $C_i$ and the totals of the cores, on the core coordinates alone, which are the pins of $\mathcal F$; and for each $i$ the \emph{pins of $B_i$}, the mass pins at $e\in B_i\setminus e_i$, on the coordinates of $B_i$ alone. The pin matrix is block diagonal, and since the fibre integral converges absolutely (Theorem~\hb{thm:leanconv}(ii)), Lemma~\ref{lem:ch24-pins}(iii) writes it as the integral over $\mathcal F$, on the translate through the core node loads of $\mathbf c$, of the fibre integrals over the coordinates of the $B_i$ at fixed core coordinates. The corners and the twists depend on the core coordinates alone and leave the inner integrals, and the inner integral for $B_i$ is
\[
\Phi_i=\int\prod_jP_{\mathcal A_{ij}}(\mathbf a_{j,i};s_j,0)\,P_{\mathcal A_{ij}}(\mathbf b_{j,i};0,t_j)
\]
over the coordinates of $B_i$ with the pins of $B_i$.

\emph{Step 5: the inner integrals are holds.} Let $F_i^\natural$ be the family of the tables $\mathcal A_{ij}$ on $B_i\sqcup\{g_i\}$ with the masses $\mvec|_{B_i}$ and $\sum_j(s_j+t_j)-\kappa_{F_i}$ on $g_i$, on its balance wall (Lemma~\hb{lem:twosumpieces}). The node loads $(\mathbf a_{j,i},\mathbf b_{j,i})_j$ of the cells are loads of its letters: for every $j$, $|\mathbf a_{j,i}|=a_{j,B_i}+s_j-a_{j,g_i}=r_{\mathcal A_{ij}}s_j$ and $|\mathbf b_{j,i}|=t_j(|B_i|+1-r_{\mathcal A_{ij}})$, their masses at $e\in B_i$ are the $m_e$, and at $g_i$ they add to $\sum_j(s_j+t_j-a_{j,g_i}-b_{j,g_i})=\sum_j(s_j+t_j)-\kappa_{F_i}$ by $C_i$. At $\mathbf c$ every period of a cell at its node loads is finite (Theorem~\hb{thm:leantree}), so the node loads of the cells at $\mathbf c$ are an admissible point of $F_i^\natural$ (Proposition~\hb{prop:fullleans}(ii)), and every proper wall of $F_i^\natural$ is positive (Theorem~\hb{thm:leanconv}(i)). In Theorem~\hb{thm:pinnedhold} for $F_i^\natural$ take $e_i$ as the fixed mark and $G=\{g_i\}$: the pinned fibre integral at the pins $s_j-a_{j,g_i}$ ($j\ge2$) and $t_j-b_{j,g_i}$ at $g_i$ is, in the coordinates of the loads of the letters other than $K_1$ at the marks of $B_i\setminus e_i$, the integral with the factor $(2\pi i)^{-(2N-1)(|B_i|-1)}$, the other loads solved from its pins. The integral $\Phi_i$, with $a_{1,e}$ for $e\in B_i\setminus e_i$ solved from the mass pin at $e$, has the same coordinates, the same factor by Lemma~\ref{lem:ch24-pins}(ii), and the same integrand, since the loads $(\mathbf a_{j,i},\mathbf b_{j,i})_j$ satisfy all the pins of $F_i^\natural$. The real parts of the core coordinates being those of $\mathbf c$, Theorem~\hb{thm:pinnedhold} gives
\[
\Phi_i=\operatorname{Hold}_{\{g_i\}}\big((s_j-a_{j,g_i})_{j\ge2},(t_j-b_{j,g_i})_j\big)\qquad\text{for }F_i^\natural .
\]

\emph{Step 6: the corners cancel.} Apply Lemma~\hb{lem:hingesplit} to the cell family $F_i$ at $\mathbf z_i=(a_{j,g_i},b_{j,g_i})_j$, of sum $\kappa_{F_i}$ by $C_i$. Its conditions hold: $0<\re a_{j,g_i}<s_j$ and $0<\re b_{j,g_i}<t_j$, the marker loads at $\mathbf c$ lying in these intervals (Theorem~\hb{thm:leantree}); and the hold of Step 5 at the real parts is finite, so by Proposition~\hb{prop:holdone} and the first identity of Lemma~\hb{lem:hingesplit} $\int\boldsymbol\varrho_{F_i}^{-\re\mathbf z_i}d\phi_{F_i}$ is finite. So
\[
\Phi_i=\prod_j\Gst21{s_j-a_{j,g_i},\ a_{j,g_i}}{s_j}\Gst21{t_j-b_{j,g_i},\ b_{j,g_i}}{t_j}\int\boldsymbol\varrho_{F_i}^{-\mathbf z_i}d\phi_{F_i},
\]
and these are the corners of the edge $\mathfrak e_i$ in Step 1, $\Gst21{a,\ b}{a+b}$ being symmetric in $a$ and $b$. They cancel, and what is left is the integral of (i) over $\mathcal F$. Its dimension is the number $2N(|\cE_{\mathcal H}|+k)$ of core coordinates less the $|\cE_{\mathcal H}|+k+2N-1$ core pins, $(2N-1)(|\cE_{\mathcal H}|+k-1)$.

(ii) Let $F'^\natural$ be the family of the tables of $F'$ with the mass $\sum_j(s_j+t_j)-\kappa_{F'}$ on $g_0$, on its balance wall (Lemma~\hb{lem:twosumpieces}), and $\mathbf c$ an admissible point of it as in the statement. Its loads at $g_0$ lie in $(0,s_j)$ and $(0,t_j)$, $g_0$ being neither a loop nor a coloop, so $\re\mathbf z$ meets the conditions of Lemma~\hb{lem:hingesplit} for $F'$; the finiteness of $\int\boldsymbol\varrho_{F'}^{-\re\mathbf z}d\phi_{F'}$ follows as in Step 6 from that of the hold of $F'^\natural$ at $g_0$ at the loads of $\mathbf c$ (Theorem~\hb{thm:pinnedhold}). By the second identity of Lemma~\hb{lem:hingesplit} the left side of (ii) is the hold of $F'^\natural$ at $g_0$ at the loads $(s_j-z_{j,K})_{j\ge2}$ and $(t_j-z_{j,\Psi})_j$, and by Theorem~\hb{thm:pinnedhold} with $G=\{g_0\}$ it is the fibre integral of Theorem~\hb{thm:leanconv}(ii) for $F'^\natural$ at $\mathbf x=\one$ with the pins at $g_0$. Steps 1 to 6 apply to it with $g_0$ a mark of the cores, its pins being core pins, and give the integral over $\mathcal F_{\mathbf z}$, of dimension $(2N-1)(|\cE_{\mathcal H}|+k)-(2N-1)=(2N-1)(|\cE_{\mathcal H}|+k-1)$.
\end{proof}

\begin{sremark}
Let the core of every table be $\binom{\{g_1,g_2\}}1$ with unit weights, the weights of the tables being those of the cells (Lemma~\hb{lem:twosumpieces}), so that $\cE_{\mathcal H}=\varnothing$ and $B=B_1$ is a bunch of width one in every table. By Proposition~\hb{prop:spleans} at $q=1$,
\[
P_{\mathcal H_j}(\mathbf a_{j,\mathcal H};s_j,0)\,P_{\mathcal H_j}(\mathbf b_{j,\mathcal H};0,t_j)=\Gst21{a_{j,g_1},\ a_{j,g_2}}{s_j}\Gst21{b_{j,g_2},\ b_{j,g_1}}{t_j}=\Gst42{a^d_j,\ b^d_j,\ a^r_j,\ b^r_j}{s_j,\ t_j}
\]
of Theorem~\hb{thm:twosides}, with the level shares $a^r_j=a_{j,g_1}$, $a^d_j=a_{j,g_2}$, $b^d_j=b_{j,g_1}$ and $b^r_j=b_{j,g_2}$. With $\mathbf z_1=(a^r_j,b^d_j)_j$ the children family of $F_1$ in Lemma~\hb{lem:hingesplit} has the tables $\mathcal A_{1j}\setminus g_1=\cB_j|B$ at $(a^r_j,t_j-b^d_j)=(a^r_j,b^r_j)$ and $\mathcal A_{1j}/g_1=\cB_j/R$ at $(s_j-a^r_j,b^d_j)=(a^d_j,b^d_j)$: it is the bunch family of Theorem~\hb{thm:twosides}; with $\mathbf z_2=(a^d_j,b^r_j)_j$ the children family of $F_2$ is the room family. The core pins are $|\mathbf a_{j,\mathcal H}|=s_j$ ($j\ge2$), $|\mathbf b_{j,\mathcal H}|=t_j$ and the masses $\kappa_{F_1}$ and $\kappa_{F_2}$; solving the totals for $a^d_j$ and $b^r_j$ by Lemma~\ref{lem:ch24-pins}(ii), with the block of the totals and of the mass $\kappa_{F_2}$ unit triangular after the row operations of Lemma~\ref{lem:ch24-pins}(i), leaves the share integral over $(a^r_j,b^d_j)_j$ pinned by $\sum_j(a^r_j+b^d_j)=\kappa_{F_1}=\kappa_B$. This is the integral of Theorem~\hb{thm:twosides} over the scale shares, with the same integrand. In the pair coordinates $\sigma^d_j=\frac12(a^d_j+b^d_j)$ and $\delta^d_j=\frac12(a^d_j-b^d_j)$, the map $(a^r_j,b^d_j)\mapsto(\sigma^d_j,\delta^d_j)$, with $a^d_j=s_j-a^r_j$, has the determinant $\frac12$ in each table, so $d(\mathbf a^r,\mathbf b^d)=2^N\,d(\boldsymbol\sigma^d,\boldsymbol\delta^d)$; and the pin $\sum_j(a^r_j+b^d_j)=\kappa_B$ reads $-2\sum_j\delta^d_j=\kappa_B-\sum_js_j$, its row being $-2$ times that of $\sum_j\delta^d_j=\frac12(\sum_js_j-\kappa_B)$, which by Definition~\hb{def:fibre} multiplies the fibre measure by $\frac12$. The two factors give the $2^{N-1}$ of the first display of Theorem~\hb{thm:twosides}, a consistency check of the bookkeeping above. In \texttt{checks/bookkeeping.py}, \texttt{check\_twosides\_pairs} computes $2^N\cdot\frac12-2^{N-1}=0$ exactly; \texttt{check\_core\_family\_balance} computes the balance of the core family at the masses $(\mvec|_{\cE_{\mathcal H}},\kappa_{F_1},\dots,\kappa_{F_k})$ and the corank identity of Step 4; and \texttt{check\_holdnode\_pins} computes, for small families, the unit Jacobian of Step 2, the split of the pins of Steps 3 and 4 with the core masses $\kappa_{F_i}$ and the totals of the cores, and the totals and the hinge mass of the node loads of the cells in Step 5.
\end{sremark}

\begin{hcorollary}[\hb{cor:holdrules} (The glue and the series and parallel rules in hold form)]
As stated; in (i) the reading for a family is new.
\end{hcorollary}

\begin{proof}
For one table the point $\mathbf z_i$ of Theorem~\hb{thm:holdnode} is $(\kappa_{\mathcal A_i}-b_{g_i},b_{g_i})$, by the core mass at $g_i$, and $\boldsymbol\varrho_{F_i}^{-\mathbf z_i}=\varrho_{\mathcal A_i}^{-b_{g_i}}$, the ratio vector of one cell being its ratio (Definition~\hb{def:hinged}); so (i) for one table, (ii) and (iii) are the reference, proof of cor:holdrules. The reading of (i) for a family is the remark after the proof of Theorem~\hb{thm:holdnode}.
\end{proof}

\begin{hdefinition}[\hb{def:networkproc} (The network procedure)]
As stated, for a family whose tables carry hinge trees; pooling and the passage to the transforms are for one table, and the recursive Mellin evaluation asks the tables to share one rooted tree.
\end{hdefinition}

\begin{sremark}
The operations rest on the statements they cite, proved above for families: (1) is Corollary~\hb{cor:nodeconv}, infeasibility meaning that some proper wall is not positive, so that the period diverges (Theorem~\hb{thm:projective}); (2) is Theorem~\hb{thm:network} and Corollary~\hb{cor:spnetwork}; (4) uses that the admissible points form an open subset of the real points of the pins, the conditions $\mathbf c^X\in w_X\,\mathrm{ri}\,\mathcal N_X$ being open there, and the bounds of Theorem~\hb{thm:leanconv}(iii), uniform on compact sets of admissible points; (5) is Theorem~\hb{thm:holdnode}, the substituted family at a vertex $v$ having the cores $\mathcal H_{j,v}$ and, as its cell families, the subtree cell families of the vertices directly below $v$. The reusable data are those of Corollary~\hb{cor:retaintable}, one Gramma for each orbit of nodes.
\end{sremark}

\begin{hproposition}[\hb{prop:networkcost} (What the network costs)]
As stated, for a family.
\end{hproposition}

\begin{proof}
(i) The fibre of Theorem~\hb{thm:network} has dimension $(2N-1)(n-1)$, and for one table each pooled pair removes one coordinate (Proposition~\hb{prop:poolpair}). Opening the period of a prime node $\mathcal H_{j,v}$ at one full lean on its share side adds $|\mathcal H_{j,v}|-n_{j,v}$ coordinates (Corollary~\hb{cor:nodecost}), so the two full leans add $2(|\mathcal H_{j,v}|-n_{j,v})$; series and parallel nodes have closed forms (Proposition~\hb{prop:spleans}). The atom formula (Corollary~\hb{cor:atomwall}) has the shares $\mathbf a_j,\mathbf b_j\in\mathbb C^{\cB_j}$ against the $2N$ total pins and the $n-1$ pins $\langle\boldsymbol\delta\rangle_e=\mu_e$, $e\neq g$: dimension $\sum_j2|\cB_j|-2N-n+1$. After its even shares close (Theorem~\hb{thm:buildup}(iv)) the fibre of differences has dimension $\sum_j(|\cB_j|-1)-(n-1)$, and each table keeps one scale integral: $\sum_j|\cB_j|-n+1$ in all.

(ii) is Corollary~\hb{cor:retaintable}: one Gramma of $n_v$ variables for each orbit of nodes, the orbits taken over all the tables.

(iii) By Theorem~\hb{thm:holdnode}(ii) a point of the Mellin transform of the law of a subtree cell family with the top node $v$ costs a fibre integral of dimension $(2N-1)(|\cE_v|+k_v-1)$, $k_v=|G_v|-1$ the number of vertices directly below $v$, that is $(2N-1)(|\cE_v|+|G_v|-2)$; at the root Theorem~\hb{thm:holdnode}(i) costs $(2N-1)(|\cE_o|+|G_o|-1)$. Their integrands need the periods of the tables of the node at their full leans at the loads of the fibre, closed at series and parallel nodes and on the share side at a prime node (Corollary~\hb{cor:nodecost}).
\end{proof}

\begin{hproposition}[\hb{prop:weightstwists} (Product weights are twists)]
As stated, for every family.
\end{hproposition}

\begin{proof}
$K_{F^{\boldsymbol\zeta},j}(\alpha)=\sum_T\theta_{j,T}\boldsymbol\zeta_j^T\alpha^T=K_j(\boldsymbol\zeta_j\circ\alpha)$, and $\Psi_{F^{\boldsymbol\zeta},j}(\alpha)=\sum_T\theta_{j,T}\boldsymbol\zeta_j^T\alpha^{\cE\setminus T}=\boldsymbol\zeta_j^{\cE}\sum_T\theta_{j,T}(\alpha/\boldsymbol\zeta_j)^{\cE\setminus T}=\boldsymbol\zeta_j^{\cE}\Psi_j(\alpha/\boldsymbol\zeta_j)$. These are the letters of $F[\mathbf x]$ with $\mathbf x^K_j=\boldsymbol\zeta_j$ and $\mathbf x^\Psi_j=\boldsymbol\zeta_j^{-1}$ (Definition~\hb{def:twist}), the cobasis letters multiplied by $\boldsymbol\zeta_j^{\cE}$, and $(\boldsymbol\zeta_j^{\cE})^{-t_j}=\boldsymbol\zeta_j^{-t_j\one}$ gives the first identity. By the torus relation of Definition~\hb{def:twist} for the family of the one letter $K_j$, the period of $K_j(\boldsymbol\zeta_j\circ\alpha)^{-s_j}$ at $\mathbf a$ is $\boldsymbol\zeta_j^{-\mathbf a}P_{\cB_j}(\mathbf a;s_j,0)$, and for the one letter $\Psi_j$ the period of $\big(\boldsymbol\zeta_j^{\cE}\Psi_j(\alpha/\boldsymbol\zeta_j)\big)^{-t_j}$ at $\mathbf b$ is $\boldsymbol\zeta_j^{-t_j\one}\boldsymbol\zeta_j^{\mathbf b}P_{\cB_j}(\mathbf b;0,t_j)$. In the lean convolution of $F^{\boldsymbol\zeta}$ the integrand is therefore that of $F$ times $\prod_j\boldsymbol\zeta_j^{-t_j\one}\boldsymbol\zeta_j^{\mathbf b_j-\mathbf a_j}$. With one $\boldsymbol\zeta$ for every table, the torus relation with the vector $\boldsymbol\zeta^{-1}$ on every letter gives $P_{F[\mathbf x]}=\boldsymbol\zeta^{-\mvec}P_{F[\mathbf x']}$ for $\mathbf x^K_j=\boldsymbol\zeta$, $\mathbf x^\Psi_j=\boldsymbol\zeta^{-1}$ and $\mathbf x'^K_j=\one$, $\mathbf x'^\Psi_j=\boldsymbol\zeta^{-2}$, and $\prod_j\boldsymbol\zeta^{-t_j\one}=\boldsymbol\zeta^{-|\mathbf t|\one}$. Weights at one mark $g$ are the $\boldsymbol\zeta_j$ equal to $1$ off $g$.
\end{proof}

\begin{hproposition}[\hb{prop:carry} (Carrying a calculation through a change, in hold form)]
As stated, for every family.
\end{hproposition}

\begin{proof}
In Theorem~\hb{thm:network} the integrand is the product of the node periods at the node loads over the Beta corners at the marker loads, times the twists. (i) The exponents of table $j$ enter the pins $|\mathbf a_j|=r_js_j$ and $|\mathbf b_j|=r^\ast_jt_j$ and the marker loads $k_{j;v,\mathfrak e}=s_jr^{\mathfrak e}_{j,v}-a_{j,\cE^{\mathfrak e}_v}$ and $l_{j;v,\mathfrak e}=t_jr^{\ast\mathfrak e}_{j,v}-b_{j,\cE^{\mathfrak e}_v}$, affine in $(s_j,t_j)$, and the corners through them; the node periods are functions of the node loads alone (Definition~\hb{def:fullleanfns}), and the native loads of the new data are admissible (Proposition~\hb{prop:fullleans}(iii)). (ii) In the coordinates of Theorem~\hb{thm:leanconv}(ii) the masses enter $\mathbf a_1=\mvec-\sum_{j\ge2}\mathbf a_j-\sum_j\mathbf b_j$ alone, the loads at $e_0$ of the other letters being fixed by their totals. (iii) is Proposition~\hb{prop:weightstwists}, for one node read for its table. (iv) The node periods are the only data of a node in the integrand; a replaced node of the same rank keeps the ranks of the branches and so the marker loads, and Theorem~\hb{thm:network} applies to the new trees, whose nodes are molecules with markers neither loops nor coloops. (v) By Theorem~\hb{thm:holdnode}(ii) the transform at a vertex is computed from the node there and the transforms of the vertices directly below it, so a change at one node changes the transforms on its path to the root only.
\end{proof}

\begin{hproposition}[\hb{prop:leanmoment} (The exponents are moment indices)]
As stated, for every family, any of its letters being closed.
\end{hproposition}

\begin{proof}
The walls are affine in $(\mathbf s,\mathbf t)$ (Definition~\hb{def:family}), so $I$, the intersection of the hyperplane with the open orthant and the open half-spaces $\kappa_S>0$, is relatively open and convex. On $I$ the period converges (Theorem~\hb{thm:projective}), and $\mvec>0$, every singleton wall $m_e-\sum_j(s_jd_j(\{e\})+t_jd^\ast_j(\{e\}))$ being positive and the counts nonnegative. Fix $i$ and apply Lemma~\hb{lem:chart} to the form of $F$ with $h=L_i^{1/g_i}$, positive and homogeneous of degree one, and $\varphi(\rho)=\rho^{w_ig_i}e^{-\rho^{g_i}}$, for which $\int_0^\infty\varphi(\rho)\,d\rho/\rho=\Gamma(w_i)/g_i$: the integrand $x^{\mvec-\one}\prod_{i'}L_{i'}^{-w_{i'}}$ times $\varphi(h)=L_i^{w_i}e^{-L_i}$ is $x^{\mvec-\one}e^{-L_i}\prod_{i'\neq i}L_{i'}^{-w_{i'}}$, which gives the first identity. Every letter being positive on the orthant and every $w_{i'}>0$, $L_{i'}(x)^{-w_{i'}}=\Gamma(w_{i'})^{-1}\int_0^\infty\lambda^{w_{i'}-1}e^{-\lambda L_{i'}(x)}d\lambda$, and Tonelli's theorem gives the second identity, the integral over $x$ being $\sG_i(\mvec;\boldsymbol\lambda)$, the Gramma of the concatenated configuration (Definition~\hb{def:gramma}). The derivative $\partial^k_{\lambda_{i'}}\sG_i$ is $(-1)^k$ times the integral with the further factor $L_{i'}^k$, which gives the complete monotonicity. By Proposition~\hb{prop:gammadraws}(iii), $P_F(\mvec;\mathbf s,\mathbf t)/\Gst n1{\mvec}{m_{\cE}}$ is the moment generating function, in the exponents, of $\big(-\log K_j(U),-\log\Psi_j(U)\big)_j$ for $U\sim\Dir_{\mvec}$, and by H\"older's inequality a moment generating function is log-convex on the convex set where it is finite; $\Gst n1{\mvec}{m_{\cE}}$ does not depend on the exponents.
\end{proof}

\section{Chapter~\hb{ch:interfaces}: finite interfaces and certified enclosures}\label{sec:ch26}

\begin{hdefinition}[\hb{def:channellaw} (The channel law; moment vectors)]
As stated; cells with the same hinges on the same marks are kept together by reading their channels jointly.
\end{hdefinition}

\begin{hproposition}[\hb{prop:channellaw} (The channel law composes)]
As stated, for cells kept together and for every cell family.
\end{hproposition}

\begin{proof}
By Theorem~\hb{thm:aggregation} with the blocks $B_1$ and $B_2$, $U\sim\Dir_{(\mvec_1,\mvec_2)}$ on $\Delta_{B_1\sqcup B_2}$ is $U|_{B_l}=V_lU_l$ with $(V_1,V_2)\sim\Dir_{(m_{B_1},m_{B_2})}$ and $U_l\sim\Dir_{\mvec_l}$, independent. In every table $j$ the channels of the spliced cell are the bilinear map of Theorem~\hb{thm:hingepairs}(iii) applied to the channels of the two cells of table $j$, and a channel $c$ has $c(V_lU_l)=V_l^{\deg c}c(U_l)$. So $\operatorname{ch}(\mathcal A_{12})(U)$ is the image of the replaced channel vectors under the bilinear map applied table by table, and its law is the channel law of $\mathcal A_{12}$. For the tables of a family kept together without hinges the channels are the letters $K_j$ and $\Psi_j$, and Proposition~\hb{prop:gammadraws}(i) is the stated reading.

For the cell family $F$, by Definitions~\hb{def:hinged} and~\hb{def:substitution} at the point $\operatorname{ch}(\mathcal A)(u)$, $\gamma^K_{\mathcal A_j}=x_{j,0}/x_{j,1}$ and $\gamma^\Psi_{\mathcal A_j}=x^\ast_{j,1}/x^\ast_{j,0}$, so $\gamma^K_{\mathcal A_j}/\gamma^K_{\mathcal A_1}=x_{j,0}x_{1,1}/(x_{j,1}x_{1,0})$ and $\gamma^\Psi_{\mathcal A_j}/\gamma^K_{\mathcal A_1}=x^\ast_{j,1}x_{1,1}/(x^\ast_{j,0}x_{1,0})$: $\boldsymbol\varrho_F=R_F\circ\operatorname{ch}(\mathcal A)$. The base form is $\phi_F(du)=\varpi(\operatorname{ch}(\mathcal A)(u))\,u^{\mvec-\one}d\omega_B(u)$ with $\varpi=\prod_jx_{j,1}^{-s_j}x_{j,0}^{\ast\,-t_j}(x_{1,0}/x_{1,1})^{-\kappa_F}$, the factor $\gamma^K_{\mathcal A_1}(u)^{-\kappa_F}$ being $(x_{1,0}/x_{1,1})^{-\kappa_F}$, and $u^{\mvec-\one}d\omega_B=\Gst{|B|}1{\mvec}{m_B}\Dir_{\mvec}(du)$ (Definition~\hb{def:dirlaw}). For a positive measure $\mu$ and positive functions, $(R\circ c)_\ast\big((\varpi\circ c)\,\mu\big)=R_\ast\big(\varpi\,c_\ast\mu\big)$, both sides giving $\int H(R(c(u)))\,\varpi(c(u))\,d\mu(u)$ against every Borel $H\ge0$. With $c=\operatorname{ch}(\mathcal A)$ and $\mu=\Gst{|B|}1{\mvec}{m_B}\Dir_{\mvec}$ this is the display.
\end{proof}

\begin{htheorem}[\hb{thm:gammalift} (Gamma draws close the moments)]
As stated, for cells kept together.
\end{htheorem}

\begin{proof}
The reference, proof of thm:gammalift, uses of the channel vector only that it is a vector of polynomials on the marks, homogeneous of known degrees, that the Gamma draws on disjoint sets of marks are independent, and that a splicing acts on the channel vectors by a bilinear map whose every product has the degree of the channel it gives. For cells kept together the channel vector is the joint vector $(\operatorname{ch}(\mathcal A_j))_j$, and a splicing acts table by table, so on the joint vectors by the direct sum over the tables of the maps of Theorem~\hb{thm:hingepairs}(iii), again bilinear with that property. So (i), (iii) and (iv) hold word for word. In (ii) the joint channels of $\mathcal A_{12}$ at $\alpha$ are $\operatorname{sp}\big(\operatorname{ch}(\mathcal A_1)(\alpha|_{B_1}),\operatorname{ch}(\mathcal A_2)(\alpha|_{B_2})\big)$, each a sum of products of one joint channel of each cell, so a monomial of degree $|\mathbf k|$ in them expands into monomials of degree $|\mathbf k|$ in the joint channels of each cell, and the independence of $\alpha|_{B_1}$ and $\alpha|_{B_2}$ gives the formula. In (v) the joint channels are bounded on $\Delta_B$, so by (iii) their moments determine their law, and Proposition~\hb{prop:channellaw} composes it.
\end{proof}

\begin{htheorem}[\hb{thm:respwalls} (Taylor coefficients in the weights are periods with walls at least as high)]
As stated, at real exponents of any sign.
\end{htheorem}

\begin{proof}
For $u\in\Delta_{\cE}$ put $y_i(u)=\sum_j\zeta_jD_{ij}(u)/L_i(u)$. Write $L_i=\sum_{\boldsymbol\gamma}c_{i,\boldsymbol\gamma}u^{\boldsymbol\gamma}$ with $c_{i,\boldsymbol\gamma}>0$; $D_{ij}$ has its support in that of $L_i$, so $D_{ij}=\sum_{\boldsymbol\gamma}c_{i,\boldsymbol\gamma}v_{i,\boldsymbol\gamma,j}u^{\boldsymbol\gamma}$ with real $v_{i,\boldsymbol\gamma,j}$, and with $V_{ij}=\max_{\boldsymbol\gamma}|v_{i,\boldsymbol\gamma,j}|$, $|D_{ij}|\le V_{ij}L_i$ on the orthant and $|y_i|\le r_i(\boldsymbol\zeta):=\sum_j|\zeta_j|V_{ij}$; let every $r_i(\boldsymbol\zeta)<1$. For real $w$ and $|y|<1$, $(1+y)^{-w}=\sum_{l\ge0}(-1)^l\frac{(w)_l}{l!}y^l$, and $|(w)_l|\le(|w|)_l$ since $|w+k|\le|w|+k$, so the series of the moduli is at most $\sum_l\frac{(|w|)_l}{l!}|y|^l=(1-|y|)^{-|w|}$. With the multinomial expansion of $y_i^{l_i}$,
\[
L_i(\boldsymbol\zeta)^{-w_i}=L_i^{-w_i}(1+y_i)^{-w_i}=\sum_{(l_{ij})_j}\frac{(-1)^{l_i}(w_i)_{l_i}}{\prod_jl_{ij}!}\,L_i^{-w_i-l_i}\prod_j\big(\zeta_jD_{ij}\big)^{l_{ij}},\qquad l_i=\sum_jl_{ij},
\]
and the product over $i$ of the series of the moduli is at most $\prod_iL_i^{-w_i}(1-r_i(\boldsymbol\zeta))^{-|w_i|}$. Against $u^{\mvec-\one}d\omega_{\cE}$ this bound is integrable, $P(\mathbf 0)$ converging at the real data; so termwise integration is justified, $P$ is a convergent power series on $\{r_i(\boldsymbol\zeta)<1\ \forall i\}$, and its coefficients are the stated ones. After the numerators are expanded every term is a coefficient times $u^{\mvec+\boldsymbol\gamma-\one}\prod_iL_i^{-w_i-l_i}$, $\boldsymbol\gamma$ a sum, over $i$, of $l_i$ exponents of the support of $L_i$. So $|\boldsymbol\gamma|=\sum_il_ig_i$ and the term is balanced, and $\gamma_S\ge\sum_il_id_i(S)$ with $d_i(S)$ the least value of $\ip{\one_S}{\cdot}$ on the support of $L_i$, which is the count $d_j(S)$ or $d^\ast_j(S)$ of the letter in the walls (Definition~\hb{def:family}). So $\kappa'_S=\kappa_S+\gamma_S-\sum_il_id_i(S)\ge\kappa_S>0$, and the term converges by Theorem~\hb{thm:projective}, whose criterion holds for exponents of any sign.
\end{proof}

\begin{htheorem}[\hb{thm:weightbank} (A uniform certificate for weights)]
As stated, at real exponents of any sign.
\end{htheorem}

\begin{proof}
Let $|\zeta_j|\le R_j$, so that $|y_i|\le r_i<1$ in the notation of the proof of Theorem~\hb{thm:respwalls}. For real $w_i$, $|(1+y_i)^{-w_i}|=|1+y_i|^{-w_i}$, the real part of $1+y_i$ being positive, and $1-r_i\le|1+y_i|\le1+r_i$, so $|L_i(\boldsymbol\zeta)^{-w_i}|\le L_i^{-w_i}\max\big((1-r_i)^{-w_i},(1+r_i)^{-w_i}\big)$; integrating, $|P(\boldsymbol\zeta)|\le M$. The radii can be enlarged a little with every $r_i<1$ kept, so $P$ is holomorphic near the closed polydisc by Theorem~\hb{thm:respwalls}, and the Cauchy estimates give $|[\boldsymbol\zeta^{\mathbf k}]P|\le M\mathbf R^{-\mathbf k}$; the bounds on the Taylor remainder follow from these as in the reference, proof of thm:weightbank, there being $\binom{k+p-1}{p-1}$ multi-indices of degree $k$. For the explicit bound on $P(\mathbf 0)$: by Theorem~\hb{thm:exist}(ii) and Theorem~\hb{thm:projective}, off the balance wall and then on it by continuity, $P(\mathbf 0)=\sum_\pi\int_{[0,1]^{n-1}}\prod_kq_k^{\kappa_{U_{\pi,k}}-1}\prod_iL_{i,\pi}(q)^{-w_i}\,dq$. Each sector letter has nonnegative coefficients and is increasing in every $q_k$, so $L_{i,\pi}(\mathbf 0)\le L_{i,\pi}(q)\le L_{i,\pi}(\one)$ and $L_{i,\pi}(q)^{-w_i}\le\max\big(L_{i,\pi}(\mathbf 0)^{-w_i},L_{i,\pi}(\one)^{-w_i}\big)$ for every real $w_i$, and $\int_{[0,1]^{n-1}}\prod_kq_k^{\kappa_{U_{\pi,k}}-1}dq=\prod_k\kappa_{U_{\pi,k}}^{-1}$. The enclosure of Proposition~\hb{prop:sectorcompiler} holds for real data of any sign, its error bound being written with $\sum_j(|s_j|+|t_j|)$, so its upper end bounds $\prod_jK_j(\one)^{s_j+t_j}P(\mathbf 0)$ from above.
\end{proof}

\section{Chapter~\hb{ch:width}: bunches of any width}\label{sec:ch27}

\begin{hdefinition}[\hb{def:seam} (Seam)]
As stated: the pair constant does not depend on the choice.
\end{hdefinition}

\begin{proof}
Let $p\in C$ and $X\in D$ with $p\sqcup X\in\cB$. The child at $p$ carries the weight $\theta_{p\sqcup X}$ on its basis $X$ (\hb{lem:minor}) and is $w_p$ times the child of the station $C$, so $\theta_{p\sqcup X}=w_p\,\theta_{(C),X}$; the room child at $X$ carries $\theta_{p\sqcup X}$ on its basis $p$ and is $w'_X$ times the room child of the station $D$, so $\theta_{p\sqcup X}=w'_X\,\theta_{(D),p}$. Hence $\theta_{p\sqcup X}/(w_pw'_X)$ is $\theta_{(C),X}/w'_X$, which does not depend on $p$, and $\theta_{(D),p}/w_p$, which does not depend on $X$. The children in a station are positive multiples of one another and so have the same bases: $p'\sqcup X\in\cB$ for every $p'\in C$ and $p\sqcup X'\in\cB$ for every $X'\in D$. Changing $p$ and then $X$ one at a time, the quotient takes one value at all pairs of $C\times D$.
\end{proof}

\begin{hproposition}[\hb{prop:seamidentity} (The seam identity)]
As stated; in (v) the seam matrix of an extreme excess is a single entry exactly when the weights of that level factor.
\end{hproposition}

\begin{proof}
Parts (i)--(iv) and (vi), and in (v) the room children and the seam letters at the two extremes, are the reference, proof of prop:seamidentity. Let $i\in\{0,c(B)\}$. By \hb{thm:levels} every pattern of excess $i$ and every room side of excess $i$ make a basis, so every station of patterns of excess $i$ is compatible with every station of room sides of excess $i$.

Let the weights of the level factor along these products, $\theta_{p\sqcup X}=\theta'_p\,\theta''_X$ for the patterns $p$ and the room sides $X$ of excess $i$, with positive $\theta'$ and $\theta''$. The child at a pattern $p$ of excess $i$ is the table on $R$ whose bases are all the room sides of excess $i$, with the weights $\theta'_p\,\theta''_X$: it is $\theta'_p$ times one table. So the patterns of excess $i$ form one station and $H_i$ has one row. The columns of a matrix with one row are proportional, and by (iv) no two stations of $\mathcal D_i$ have proportional columns; so $H_i$ has one column, and its entry is positive, the pair being compatible.

Conversely let $H_i$ be a single entry $h_{CD}$. The patterns of excess $i$ then form the one station $C$ and its room sides the one station $D$, and (ii) gives $\theta_{p\sqcup X}=h_{CD}\,w_p\,w'_X$ for every such $p$ and $X$: the weights of the level factor along these products, with the factors $h_{CD}\,w_p$ and $w'_X$.
\end{proof}

\begin{hproposition}[\hb{prop:ladder} (The width ladder)]
As stated; in (i) and (ii) the seam is one rectangle at each excess exactly when the weights factor across $B$ and $R$, and the closing sentence holds at every width.
\end{hproposition}

\begin{proof}
(i) That $c(B)=0$ exactly when $B$ is a union of molecules, with one excess, every pattern meeting every room side, and the letters $K=\mathsf x_0^{\mathsf T}H_0\,\mathsf y_0$ and $\Psi=\mathsf x_0^{\ast\mathsf T}H_0\,\mathsf y^\ast_0$, is the reference, proof of prop:ladder(i). The one excess is extreme, so by \hb{prop:seamidentity}(v) the seam is one rectangle exactly when the weights of its level factor along the products $p\sqcup X$. Every basis lies in this level, so this is a factoring $\theta_T=\theta'_{T\cap B}\,\theta''_{T\cap R}$ across $B$ and $R$, and conversely. With one station $C$ and one station $D$, \hb{prop:seamidentity}(iii) reads $K=h_{CD}\,\mathsf x_C\,\mathsf y_D$ and $\Psi=h_{CD}\,\mathsf x^\ast_C\,\mathsf y^\ast_D$.

(ii) That $c(B)=1$ exactly when there are two excesses, both extreme, is the reference, proof of prop:ladder(ii). By \hb{prop:seamidentity}(v) the excess $i$ is one rectangle exactly when the weights of its level factor, $\theta_{p\sqcup X}=\theta'_{i,p}\,\theta''_{i,X}$ for the patterns $p$ and the room sides $X$ of excess $i$. If both levels factor, put $\theta'_p=\theta'_{i(p),p}$ and $\theta''_X=\theta''_{i(X),X}$, each pattern and each room side having one excess; the pattern and the room side of a basis have one excess, so $\theta_T=\theta'_{T\cap B}\,\theta''_{T\cap R}$ for every basis $T$. Conversely a factoring across $B$ and $R$ restricts to a factoring of each level. The seam letters, the room station letters and the room-station factorization are the reference's (proof of prop:ladder(ii), and exercise ex:widthone).

(iii) and the closing sentence. A matrix with one row has proportional columns, and a matrix with one column proportional rows; so by \hb{prop:seamidentity}(iv) an excess with one station on one side has one station on the other, at every width. The criterion for level form is \hb{prop:levelform}, which does not restrict the width. In level form the children at the patterns of one excess are equal, so these patterns form one station; the excess then has one station on each side and is one rectangle. That an excess $0<i<c(B)$ can carry several rectangles is the reference's (prop:ladder(iii) and exercise ex:prismseam).
\end{proof}

\begin{hproposition}[\hb{prop:seamsym} (Symmetries of a seam)]
As stated; in (iii) an exchanging bijection is read at every lean, and it carries the law of $B$ at the lean $\boldsymbol\delta$ to its law at $-\boldsymbol\delta$.
\end{hproposition}

\begin{proof}
Parts (i) and (ii), the first sentence of (iii), the carrying bijections and the sentence on the potential are the reference, proof of prop:seamsym.

Let $g$ keep $B$ and the masses and exchange the family, and let $j\mapsto j'$ be the bijection of its tables of \hb{def:carry}. The pairs are kept, $(\sigma_{j'},\delta_{j'})=(\sigma_j,\delta_j)$, so at the lean $-\boldsymbol\delta$ table $j'$ carries the exponents $(t_{j'},s_{j'})=(t_j,s_j)$. Read every function of $\alpha_B$ at $g^\ast\alpha$. This permutes the coordinates of $\Delta_B$ and keeps $u^{\mvec_B-\one}d\omega_B$, since $g$ keeps $B$ and the masses. By \hb{lem:windowturn}, $c_{j'}(B)=c_j(B)$, $d_{j'}(B)=|B|-r_j(B)=d^\ast_j(B)$ and $d^\ast_{j'}(B)=d_j(B)$.
\begin{enumerate}[label=(\alph*)]
\item By (ii) the seam letter $V_{jD}$ of $K_j$ at the excess $i(D)$ is, at $g^\ast\alpha$, the seam letter $V^\ast_{j'D'}$ of $\Psi_{j'}$ at the excess $c_j-i(D)$, and $V^\ast_{jD}$ is $V_{j'D'}$. Summing over the stations of one excess, $V^{(i)}_j=V^{\ast(c_j-i)}_{j'}$ and $V^{\ast(i)}_j=V^{(c_j-i)}_{j'}$.
\item $K_j$ is active at $\boldsymbol\delta$ exactly when $s_j>0$, that is exactly when $\Psi_{j'}$ is active at $-\boldsymbol\delta$; likewise $\Psi_j$ and $K_{j'}$. The sum $\kappa_B+\kappa^{\mathrm{top}}_B=\sum_j(s_j+t_j)c_j$ is the same at the two leans.
\item By (a) and (b) the scale of the piece at $\boldsymbol\delta$ (\hb{def:seamlaw}) is
\[
\prod_j\Big(\frac{V^{\ast(0)}_{j'}}{V^{\ast(c_j)}_{j'}}\Big)^{s_j/(\kappa_B+\kappa^{\mathrm{top}}_B)}\Big(\frac{V^{(c_j)}_{j'}}{V^{(0)}_{j'}}\Big)^{t_j/(\kappa_B+\kappa^{\mathrm{top}}_B)},
\]
which is the scale at $-\boldsymbol\delta$, where $\Psi_{j'}$ carries the exponent $s_j$ and $K_{j'}$ the exponent $t_j$.
\item The wall at $-\boldsymbol\delta$ is $m_B-\sum_j\big(t_jd_{j'}(B)+s_jd^\ast_{j'}(B)\big)=m_B-\sum_j\big(s_jd_j(B)+t_jd^\ast_j(B)\big)=\kappa_B$.
\item By (a), (c) and (d) the scale form at $\boldsymbol\delta$,
\[
u^{\mvec_B-\one}\prod_j(V^{(0)}_j)^{-s_j}(V^{\ast(c_j)}_j)^{-t_j}(\gamma^\rho_B)^{-\kappa_B}d\omega_B,
\]
is the scale form at $-\boldsymbol\delta$.
\item For a station $D$ of excess $i(D)$ and its image $D'$ of excess $c_j-i(D)$,
\[
\varrho^\rho_{jD}=\frac{V_{jD}}{V^{(0)}_j\,(\gamma^\rho_B)^{i(D)}}=\frac{V^\ast_{j'D'}}{V^{\ast(c_j)}_{j'}\,(\gamma^\rho_B)^{c_j-i(D')}},
\]
the coordinate $\varrho^{\rho\ast}_{j'D'}$ at $-\boldsymbol\delta$; likewise $\varrho^{\rho\ast}_{jD}$ is the coordinate $\varrho^\rho_{j'D'}$ at $-\boldsymbol\delta$.
\end{enumerate}
So the seam vector at $\boldsymbol\delta$, followed by the map of the coordinates, is the seam vector at $-\boldsymbol\delta$ read at $g^\ast\alpha$, and by (e) the scale forms correspond. The image of the law at $\boldsymbol\delta$ is therefore the law at $-\boldsymbol\delta$. By (a) and (b) the relations defining the seam torus at $\boldsymbol\delta$ go over into those at $-\boldsymbol\delta$. At the centre $\boldsymbol\delta=-\boldsymbol\delta$, and the law is kept.
\end{proof}

\begin{hlemma}[\hb{lem:absorbwidth} (Absorbing a piece of any width)]
As stated, with functions $f_l$ homogeneous of degree one in their second argument.
\end{hlemma}

\begin{proof}
Take the section $h=\gamma^\rho_B(\alpha_B)+|\alpha_Z|$, a positive function of degree one on $(0,\infty)^{B\sqcup Z}$, and write $\alpha_B=\lambda u$ with $\lambda>0$ and $u\in\Delta_B$ (\hb{lem:scaleshape}).
\begin{itemize}
\item \emph{The seam letters along the scale.} Let $f_l$ be assigned to $K_j$. By \hb{def:seamlaw}, $V_{jD}=V^{(0)}_j\,(\gamma^\rho_B)^{i(D)}\varrho^\rho_{jD}$ on $\Delta_B$, and $V_{jD}$ is homogeneous of degree $d_j+i(D)$. Hence
\begin{gather*}
V_{jD}(\lambda u)=\lambda^{d_j}V^{(0)}_j(u)\,\big(\lambda\gamma^\rho_B(u)\big)^{i(D)}\varrho^\rho_{jD}(u),\\
\text{that is}\quad\mathbf V_l(\lambda u)=V_l(\lambda u)\,\mathbf x_l(\alpha_g),
\end{gather*}
at $\alpha_g=\lambda\gamma^\rho_B(u)$ and $\mathbf x=\boldsymbol\varrho^\rho_B(u)$, with $V_l(\lambda u)=\lambda^{d_j}V^{(0)}_j(u)>0$. For $\Psi_j$ the same holds with $V^\ast_{jD}=V^{\ast(c_j)}_j(\gamma^\rho_B)^{c_j-i(D)}\varrho^{\rho\ast}_{jD}$, of degree $d^\ast_j+c_j-i(D)$, and $V_l(\lambda u)=\lambda^{d^\ast_j}V^{\ast(c_j)}_j(u)$.
\item \emph{Homogeneity.} Since $f_l$ is homogeneous of degree one in its second argument,
\[
f_l\big(\alpha_Z,\mathbf V_l(\lambda u)\big)=V_l(\lambda u)\,f_l\big(\alpha_Z,\mathbf x_l(\alpha_g)\big).
\]
So the arguments of $H$ are the $f_l(\alpha_Z,\mathbf x_l(\alpha_g))$. The exponents assigned to each active letter add to its exponent, so
\begin{multline*}
\prod_lf_l\big(\alpha_Z,\mathbf V_l(\lambda u)\big)^{-a_l}\\
=\lambda^{-\sum_j(s_jd_j+t_jd^\ast_j)}\prod_jV^{(0)}_j(u)^{-s_j}V^{\ast(c_j)}_j(u)^{-t_j}\prod_lf_l\big(\alpha_Z,\mathbf x_l(\alpha_g)\big)^{-a_l}.
\end{multline*}
\item \emph{The scale and the stand-in.} With $\alpha_B^{\mvec_B-\one}d\alpha_B=\lambda^{m_B-1}u^{\mvec_B-\one}\,d\lambda\,d\omega_B(u)$, the power of $\lambda$ is $\lambda^{\kappa_B-1}$. At fixed $u$ substitute $\alpha_g=\lambda\gamma^\rho_B(u)$, so that $\lambda^{\kappa_B-1}d\lambda=\gamma^\rho_B(u)^{-\kappa_B}\alpha_g^{\kappa_B-1}d\alpha_g$. In the coordinates $(u,\alpha_g,\alpha_Z)$ the form $\omega$ is then $\phi^\rho_B(du)$ times $\omega_{\boldsymbol\varrho^\rho_B(u)}$ (\hb{def:seamlaw}), and $h=\alpha_g+|\alpha_Z|$.
\item \emph{The integral.} At fixed $u$ the map $(\lambda,\alpha_Z)\mapsto(\alpha_g,\alpha_Z)$ intertwines the scalings of the two orthants, so $\omega_{\boldsymbol\varrho^\rho_B(u)}$ is homogeneous of degree zero together with $\omega$. By \hb{lem:chart}, the integral of a form homogeneous of degree zero over a section does not depend on the section. Tonelli's theorem, applied to the nonnegative integrand, then gives
\begin{align*}
\int_{\{h=1\}}\omega&=\int_{\Delta_B}\phi^\rho_B(du)\int_{\{\alpha_g+|\alpha_Z|=1\}}\omega_{\boldsymbol\varrho^\rho_B(u)}\\
&=\int\Big(\int_{\{h'=1\}}\omega_{\mathbf x}\Big)\,d\big((\boldsymbol\varrho^\rho_B)_\ast\phi^\rho_B\big)(\mathbf x).
\end{align*}
\end{itemize}
This is the reference's proof (proof of lem:absorbwidth). There, linearity of the $f_l$ enters only through the identity $f_l(\alpha_Z,\mathbf V_l(\lambda u))=V_l(\lambda u)\,f_l(\alpha_Z,\mathbf x_l(\alpha_g))$, which homogeneity of degree one gives.
\end{proof}

\begin{sremark}
At width one with factoring weights, normalized as in \hb{prop:widthonelaw}, the seam letters of $K_j$ are $(K_{\cB_j/R},K_{\cB_j|B})$ and those of $\Psi_j$ are $(\Psi_{\cB_j/R},\Psi_{\cB_j|B})$ (Lemma~\ref{lem:ch27-normal}). The positive continuous functions homogeneous of degree one in the second argument are, for $K_j$, $f_l(\alpha_Z,(y_0,y_1))=y_0\,f(\alpha_Z,y_1/y_0)$, and for $\Psi_j$, $f_l(\alpha_Z,(y_0,y_1))=y_1\,f(\alpha_Z,y_0/y_1)$, with $f$ positive and continuous on $(0,\infty)^Z\times(0,\infty)$. At the seam letters they are $K_{\cB_j/R}\,f(\alpha_Z,\gamma^K_{j,B})$ and $\Psi_{\cB_j|B}\,f(\alpha_Z,\gamma^\Psi_{j,B})$. These are the factors $L_i\,f_i(\alpha_Z,\gamma_i)$ of \hb{lem:absorb} for the cell family $F^+_B$, with arbitrary $f_i$. Linear functions $f_l$ give only the $f_i$ that are affine in the hinge value; homogeneity of degree one gives all of them.
\end{sremark}

\begin{slemma}\label{lem:ch27-normal}
Let a bunch $B$ have width one in table $j$ of a family, the weights factoring across $B$ and $R$, $\theta_{j,T}=\theta'_{j,T\cap B}\,\theta''_{j,T\cap R}$, with the stations $D_0$ and $D_1$ of the excesses $0$ and $1$ (\hb{prop:ladder}(ii)). In this statement write $\hat\theta_i$ for the largest value of $\theta''_j$ on the room sides of excess $i$, and read $\cB_j/R$ and $\cB_j|B$ with the weights $\theta'_j$ and $\cB_j|R$ and $\cB_j/B$ with the weights $\theta''_j$.
\begin{enumerate}[label=(\roman*)]
\item $K_{j,(D_0)}=\hat\theta_0K_{\cB_j/R}$, $\Psi_{j,(D_0)}=\hat\theta_0\Psi_{\cB_j/R}$, $K_{j,(D_1)}=\hat\theta_1K_{\cB_j|B}$ and $\Psi_{j,(D_1)}=\hat\theta_1\Psi_{\cB_j|B}$; and $\mathsf y_{jD_0}=K_{\cB_j|R}/\hat\theta_0$, $\mathsf y^\ast_{jD_0}=\Psi_{\cB_j|R}/\hat\theta_0$, $\mathsf y_{jD_1}=K_{\cB_j/B}/\hat\theta_1$ and $\mathsf y^\ast_{jD_1}=\Psi_{\cB_j/B}/\hat\theta_1$.
\item For $i=0,1$, divide $\theta''_j$ by $\hat\theta_i$ on the room sides of excess $i$ and multiply $\theta'_j$ by $\hat\theta_i$ on the patterns of excess $i$. This gives a factorization of the same weights with $\hat\theta_0=\hat\theta_1=1$, in which $\theta''_j$ is the system of the multipliers of the room sides (\hb{def:seam}). The seam letters and the room station letters do not depend on the factorization. In the normalized factorization
\[
\gamma^K_{j,B}=\frac{K_{j,(D_1)}}{K_{j,(D_0)}},\qquad\gamma^\Psi_{j,B}=\frac{\Psi_{j,(D_0)}}{\Psi_{j,(D_1)}},\qquad\gamma^K_{j,R}=\frac{\mathsf y_{jD_0}}{\mathsf y_{jD_1}},\qquad\gamma^\Psi_{j,R}=\frac{\mathsf y^\ast_{jD_1}}{\mathsf y^\ast_{jD_0}}.
\]
In a factorization with the constants $\hat\theta_0$ and $\hat\theta_1$, the four ratios on the right are $(\hat\theta_1/\hat\theta_0)\gamma^K_{j,B}$, $(\hat\theta_0/\hat\theta_1)\gamma^\Psi_{j,B}$, $(\hat\theta_1/\hat\theta_0)\gamma^K_{j,R}$ and $(\hat\theta_0/\hat\theta_1)\gamma^\Psi_{j,R}$.
\end{enumerate}
\end{slemma}

\begin{proof}
(i) By \hb{thm:widthone} the bases of table $j$ are the $X\sqcup Y$ with $X\in\cB_j|B$ and $Y\in\cB_j/B$, at the level $r_j(B)$, and with $X\in\cB_j/R$ and $Y\in\cB_j|R$, at the level $d_j(B)$.
\begin{itemize}
\item So the patterns and the room sides of excess $0$ are the bases of $\cB_j/R$ and of $\cB_j|R$, those of excess $1$ the bases of $\cB_j|B$ and of $\cB_j/B$, and every pattern meets every room side of its excess.
\item The room child at a room side $X$ of excess $0$ (\hb{def:seam}) is the table on $B$ with the bases of $\cB_j/R$ and the weights $\theta'_{j,p}\,\theta''_{j,X}$, that is $\theta''_{j,X}$ times $\cB_j/R$ with the weights $\theta'_j$.
\item These room children are positive multiples of one another, so the room sides of excess $0$ form the one station $D_0$. Its multipliers are $w'_X=\theta''_{j,X}/\hat\theta_0$, the largest equal to $1$, and the room child of the station is $\hat\theta_0$ times $\cB_j/R$ with $\theta'_j$.
\item Hence $K_{j,(D_0)}=\hat\theta_0K_{\cB_j/R}$ and $\Psi_{j,(D_0)}=\hat\theta_0\Psi_{\cB_j/R}$, and the station letters are $\mathsf y_{jD_0}=\sum_Xw'_X\alpha_R^X=K_{\cB_j|R}/\hat\theta_0$ and $\mathsf y^\ast_{jD_0}=\sum_Xw'_X\alpha_R^{R\setminus X}=\Psi_{\cB_j|R}/\hat\theta_0$.
\item The excess $1$ is the same with $\cB_j|B$, $\cB_j/B$ and $\hat\theta_1$.
\end{itemize}

(ii) A basis has its pattern and its room side at one excess $i$, so the new weights multiply to $(\hat\theta_i\,\theta'_{j,p})(\theta''_{j,X}/\hat\theta_i)=\theta_{j,p\sqcup X}$. The new $\theta''_j$ has the largest value $1$ on each excess, so it is the system of the multipliers $w'_X$. The stations, the multipliers, the seam letters and the station letters are defined from the weights $\theta_j$ alone (\hb{def:seam}), so they do not change. In the normalized factorization, (i) holds with $\hat\theta_0=\hat\theta_1=1$, and the four ratios are the hinge values of \hb{def:readings}, read with the minors carrying the normalized weights. In a factorization with the constants $\hat\theta_0$ and $\hat\theta_1$, (i) gives the four ratios as stated.
\end{proof}

\begin{hproposition}[\hb{prop:widthonelaw} (Width one is the hinge)]
As stated: for a family, with $\theta''_j$ normalized on each excess, and with part (iv) for the family.
\end{hproposition}

\begin{proof}
\emph{The letters.} With $\theta''_j$ of largest value $1$ on the room sides of each excess, Lemma~\ref{lem:ch27-normal} gives the seam letters of the piece of the bunch and the room letters of its room, read as its partner (the paragraph after \hb{def:seampiece}):
\[
V_{jD_0}=K_{\cB_j/R},\quad V_{jD_1}=K_{\cB_j|B},\quad V^\ast_{jD_0}=\Psi_{\cB_j/R},\quad V^\ast_{jD_1}=\Psi_{\cB_j|B},
\]
\[
L_{jD_0}=K_{\cB_j|R},\quad L_{jD_1}=K_{\cB_j/B},\quad L^\ast_{jD_0}=\Psi_{\cB_j|R},\quad L^\ast_{jD_1}=\Psi_{\cB_j/B},
\]
the minors on $B$ carrying the weights $\theta'_j$ and those on $R$ the weights $\theta''_j$. The hinge values of \hb{def:readings} are $\gamma^K_{j,B}=V_{jD_1}/V_{jD_0}$, $\gamma^\Psi_{j,B}=V^\ast_{jD_0}/V^\ast_{jD_1}$, $\gamma^K_{j,R}=L_{jD_0}/L_{jD_1}$ and $\gamma^\Psi_{j,R}=L^\ast_{jD_1}/L^\ast_{jD_0}$. Every letter is active and $c_j=1$, so $\kappa_B+\kappa^{\mathrm{top}}_B=\sum_j(s_j+t_j)=2\sigma$, and on the balance wall of the glue $\kappa_B+\kappa_R=2\sigma$.

\emph{The cells.}
\begin{itemize}
\item The cell of the bunch $\cB^+_{j,B}$ has the bases of $\cB_j|B$ avoiding $g$ and the $X+g$, $X\in\cB_j/R$, through $g$, with the weights $\theta'_j$ (\hb{thm:widthone}). So $\cB^+_{j,B}\setminus g=\cB_j|B$ and $\cB^+_{j,B}/g=\cB_j/R$, its hinge values are $\gamma^K_{j,B}$ and $\gamma^\Psi_{j,B}$ (\hb{def:substitution}), and its rank is $d_j(B)+1$.
\item Since $d^\ast_j(B)=|B|-d_j(B)-1$, the hinge wall of $F^+_B$ (\hb{def:hinged}) is $m_B-\sum_j\big(s_jd_j(B)+t_j(|B|-d_j(B)-1)\big)=\kappa_B$.
\item Likewise $\cB^+_{j,R}\setminus g=\cB_j|R$ and $\cB^+_{j,R}/g=\cB_j/B$, with the weights $\theta''_j$ and the hinge values $\gamma^K_{j,R}$ and $\gamma^\Psi_{j,R}$.
\item The hinge wall of $F^+_R$ is $m_R-\sum_j\big(s_j(r(\cB_j|R)-1)+t_j(|R|-r(\cB_j|R))\big)$. This is the wall $\kappa_R$ of the partner, whose letters at $D_0$ have the degrees $r(\cB_j|R)=e_j+1$ and $|R|-r(\cB_j|R)=e^\ast_j$ (\hb{def:seampiece}).
\end{itemize}

(i) By \hb{def:seamlaw},
\[
\gamma^\rho_B=\prod_j\Big(\frac{V_{jD_1}}{V_{jD_0}}\Big)^{s_j/(2\sigma)}\Big(\frac{V^\ast_{jD_0}}{V^\ast_{jD_1}}\Big)^{t_j/(2\sigma)}=\prod_j(\gamma^K_{j,B})^{s_j/(2\sigma)}(\gamma^\Psi_{j,B})^{t_j/(2\sigma)}.
\]
As $\sum_j(s_j+t_j)/(2\sigma)=1$, this is $\gamma^K_{1,B}\prod_j(\gamma^K_{j,B}/\gamma^K_{1,B})^{s_j/(2\sigma)}(\gamma^\Psi_{j,B}/\gamma^K_{1,B})^{t_j/(2\sigma)}=\gamma^K_{1,B}\,\psi(\boldsymbol\varrho_{F^+_B})$, the ratio vector of $F^+_B$ being $(\gamma^K_{j,B}/\gamma^K_{1,B},\gamma^\Psi_{j,B}/\gamma^K_{1,B})_j$ (\hb{def:hinged}). The scale form is
\[
\phi^\rho_B(du)=u^{\mvec_B-\one}\prod_jK_{\cB_j/R}^{-s_j}\Psi_{\cB_j|B}^{-t_j}\big(\gamma^K_{1,B}\psi\big)^{-\kappa_B}d\omega_B(u)=\psi^{-\kappa_B}\phi_{F^+_B}(du),
\]
by the base form of \hb{def:hinged}, whose letters are $K_{\cB^+_{j,B}/g}=K_{\cB_j/R}$ and $\Psi_{\cB^+_{j,B}\setminus g}=\Psi_{\cB_j|B}$, with the hinge wall $\kappa_B$.

(ii) By \hb{def:seamlaw}, $\varrho^\rho_{jD_0}=V_{jD_0}/V^{(0)}_j=1$ and $\varrho^{\rho\ast}_{jD_1}=V^\ast_{jD_1}/V^{\ast(1)}_j=1$. Further $\varrho^\rho_{jD_1}=V_{jD_1}/(V_{jD_0}\gamma^\rho_B)=\gamma^K_{j,B}/(\gamma^K_{1,B}\psi)=x_{j,K}/\psi(\mathbf x)$ and $\varrho^{\rho\ast}_{jD_0}=V^\ast_{jD_0}/(V^\ast_{jD_1}\gamma^\rho_B)=x_{j,\Psi}/\psi(\mathbf x)$, at $\mathbf x=\boldsymbol\varrho_{F^+_B}$. Write $\Phi(\mathbf x)=\big(x_{j,K}/\psi(\mathbf x),\,x_{j,\Psi}/\psi(\mathbf x)\big)_j$.
\begin{itemize}
\item The seam torus is the set of the points with $x_{jD_0}=x^\ast_{jD_1}=1$ and $\prod_jx_{jD_1}^{s_j}x^{\ast\,t_j}_{jD_0}=1$. Since $\prod_j(x_{j,K}/\psi)^{s_j}(x_{j,\Psi}/\psi)^{t_j}=\psi^{2\sigma}\psi^{-2\sigma}=1$, $\Phi$ maps $(0,\infty)^{2N-1}$ into it.
\item Let a point of the torus have $x_{jD_1}=a_j$ and $x^\ast_{jD_0}=b_j$. The vector $\mathbf x$ with $x_{j,K}=a_j/a_1$ and $x_{j,\Psi}=b_j/a_1$ has $\psi(\mathbf x)=\big(\prod_ja_j^{s_j}b_j^{t_j}\big)^{1/(2\sigma)}/a_1=1/a_1$, so $\Phi(\mathbf x)$ is the point.
\item Conversely $\Phi(\mathbf x)$ determines $\psi(\mathbf x)$ as the inverse of its coordinate $x_{1,K}/\psi(\mathbf x)$, and then $\mathbf x$. So $\Phi$ is a bijection.
\end{itemize}
As $\boldsymbol\varrho^\rho_B=\Phi\circ\boldsymbol\varrho_{F^+_B}$ and $\phi^\rho_B=\psi(\boldsymbol\varrho_{F^+_B})^{-\kappa_B}\phi_{F^+_B}$, the law $(\boldsymbol\varrho^\rho_B)_\ast\phi^\rho_B$ is the image under $\Phi$ of $\psi^{-\kappa_B}$ times $(\boldsymbol\varrho_{F^+_B})_\ast\phi_{F^+_B}$.

(iii) The partner with its stand-in at $\Phi(\mathbf x)$ (\hb{thm:gluewidth}(i)) has the letters
\begin{gather*}
\sum_Dx_{jD}\,\alpha_g^{\,i(D)}L_{jD}=K_{\cB_j|R}+\frac{x_{j,K}}{\psi(\mathbf x)}\,\alpha_g\,K_{\cB_j/B},\\
\sum_Dx^\ast_{jD}\,\alpha_g^{\,1-i(D)}L^\ast_{jD}=\frac{x_{j,\Psi}}{\psi(\mathbf x)}\,\alpha_g\,\Psi_{\cB_j|R}+\Psi_{\cB_j/B}.
\end{gather*}
After $\alpha_g\mapsto\psi(\mathbf x)\,\alpha_g$ these are $K_{\cB_j|R}+x_{j,K}\,\alpha_gK_{\cB_j/B}$ and $x_{j,\Psi}\,\alpha_g\Psi_{\cB_j|R}+\Psi_{\cB_j/B}$. These are the letters of $\cB^+_{j,R}$ twisted at $g$ by $x_{j,K}$ in $K_j$ and by $x_{j,\Psi}$ in $\Psi_j$: by \hb{def:twist}, the twist multiplies the weights in $K_j$ of the bases through $g$ and the weights in $\Psi_j$ of the bases avoiding $g$. The mass $\kappa_B$ on $g$ turns $\alpha_g^{\kappa_B-1}d\alpha_g$ into $\psi(\mathbf x)^{\kappa_B}\alpha_g^{\kappa_B-1}d\alpha_g$, which cancels the density of (ii). The functions of \hb{lem:absorbwidth}, homogeneous of degree one in the vector $(y_0,y_1)$ of the seam letters of a letter, are $y_0\,f(\alpha_Z,y_1/y_0)$ for $K_j$ and $y_1\,f(\alpha_Z,y_0/y_1)$ for $\Psi_j$, with $f$ positive and continuous. At the seam letters they are $K_{\cB_j/R}\,f(\alpha_Z,\gamma^K_{j,B})$ and $\Psi_{\cB_j|B}\,f(\alpha_Z,\gamma^\Psi_{j,B})$, the factors $L_i\,f_i(\alpha_Z,\gamma_i)$ of \hb{lem:absorb} for $F^+_B$. At $\mathbf x_l(\alpha_g)$, after the change of variable, they are $f(\alpha_Z,x_{j,K}\alpha_g)$ and $f(\alpha_Z,x_{j,\Psi}\alpha_g)$, the factors of $\omega_{\mathbf x}$ in \hb{lem:absorb}. So the two lemmas are one, read through $\Phi$.

(iv) \emph{The room.} Read with the reversed excesses, the station $D_1$ has the partner excess $0$ and $D_0$ the partner excess $1$. By \hb{def:seamlaw} for the partner:
\begin{itemize}
\item its scale is
\[
\prod_j\Big(\frac{L_{jD_0}}{L_{jD_1}}\Big)^{s_j/(2\sigma)}\Big(\frac{L^\ast_{jD_1}}{L^\ast_{jD_0}}\Big)^{t_j/(2\sigma)}=\prod_j(\gamma^K_{j,R})^{s_j/(2\sigma)}(\gamma^\Psi_{j,R})^{t_j/(2\sigma)}=\gamma^K_{1,R}\,\psi(\boldsymbol\varrho_{F^+_R}),
\]
as in (i);
\item its scale form is $v^{\mvec_R-\one}\prod_jK_{\cB_j/B}^{-s_j}\Psi_{\cB_j|R}^{-t_j}(\gamma^K_{1,R}\psi)^{-\kappa_R}d\omega_R(v)=\psi^{-\kappa_R}\phi_{F^+_R}$, since the base form of $F^+_R$ has the letters $K_{\cB^+_{j,R}/g}=K_{\cB_j/B}$ and $\Psi_{\cB^+_{j,R}\setminus g}=\Psi_{\cB_j|R}$ and the hinge wall $\kappa_R$;
\item with $\gamma_R$ its scale, its seam vector has $\varrho^\rho_{jD_1}=\varrho^{\rho\ast}_{jD_0}=1$, $\varrho^\rho_{jD_0}=L_{jD_0}/(L_{jD_1}\gamma_R)=x'_{j,K}/\psi(\mathbf x')$ and $\varrho^{\rho\ast}_{jD_1}=L^\ast_{jD_1}/(L^\ast_{jD_0}\gamma_R)=x'_{j,\Psi}/\psi(\mathbf x')$ at $\mathbf x'=\boldsymbol\varrho_{F^+_R}$;
\item its law follows as in (ii).
\end{itemize}
The coordinate ratios are $x_{j,\Psi}/x_{j,K}=\gamma^\Psi_{j,B}/\gamma^K_{j,B}=\varrho_{j,B}$ and $x'_{j,\Psi}/x'_{j,K}=\varrho_{j,R}$ (\hb{def:readings}).

\emph{The kernel.} At $c_j=1$ the configuration $t_{\mathcal D}$ has the column $\mathbf e_{K_j}$ at $(j,D_0)$ and $\mathbf e_g+\mathbf e_{K_j}$ at $(j,D_1)$ for $K_j$, and $\mathbf e_g+\mathbf e_{\Psi_j}$ at $D_0$ and $\mathbf e_{\Psi_j}$ at $D_1$ for $\Psi_j$. The display of \hb{thm:gluewidth}(ii) gives
\begin{multline*}
\Gamma(\kappa_B)\,\sZ_{(\kappa_B,\mathbf s,\mathbf t);t_{\mathcal D}}(\mathbf p\circ\mathbf p')\\
=\int_0^\infty v^{\kappa_B-1}\prod_j\big(p_{jD_0}p'_{jD_0}+p_{jD_1}p'_{jD_1}v\big)^{-s_j}\big(p^\ast_{jD_0}p'^\ast_{jD_0}v+p^\ast_{jD_1}p'^\ast_{jD_1}\big)^{-t_j}dv,
\end{multline*}
and the coordinates of (ii) and of the partner give the first line of (iv). Substitute $v=(\psi(\mathbf x)/\psi(\mathbf x'))\,w$. Each factor $x_{j,K}v/\psi(\mathbf x)+x'_{j,K}/\psi(\mathbf x')$ becomes $(x_{j,K}w+x'_{j,K})/\psi(\mathbf x')$, and likewise for $\Psi_j$. Since $\sum_j(s_j+t_j)-\kappa_B=\kappa_R$, the integral is
\begin{multline*}
\Big(\frac{\psi(\mathbf x)}{\psi(\mathbf x')}\Big)^{\kappa_B}\psi(\mathbf x')^{\sum_j(s_j+t_j)}\int_0^\infty w^{\kappa_B-1}\prod_j(x_{j,K}w+x'_{j,K})^{-s_j}(x_{j,\Psi}w+x'_{j,\Psi})^{-t_j}dw\\
=\psi(\mathbf x)^{\kappa_B}\psi(\mathbf x')^{\kappa_R}\,P_{F_{\mathcal H}[(\mathbf x,\mathbf x')]}\big((\kappa_B,\kappa_R);\mathbf s,\mathbf t\big).
\end{multline*}
The last step is the period of the twisted core of \hb{thm:node}(iii) with the core $\binom{\{g,g'\}}1$ in every table, read on the section $\alpha_{g'}=1$ (\hb{lem:chart}); its letters are $x_{j,K}\alpha_g+x'_{j,K}\alpha_{g'}$ and $x_{j,\Psi}\alpha_g+x'_{j,\Psi}\alpha_{g'}$ (\hb{def:twist}).

\emph{The reductions.}
\begin{itemize}
\item \hb{thm:gluewidth}(i) is part (iii).
\item \hb{thm:gluewidth}(ii): insert the laws of (ii) and of the partner. The densities $\psi(\mathbf x)^{-\kappa_B}$ and $\psi(\mathbf x')^{-\kappa_R}$ cancel the factor of the kernel, and the period becomes
\[
\int\!\!\int P_{F_{\mathcal H}[(\mathbf x,\mathbf x')]}\big((\kappa_B,\kappa_R);\mathbf s,\mathbf t\big)\,d\big((\boldsymbol\varrho_{F^+_B})_\ast\phi_{F^+_B}\big)(\mathbf x)\,d\big((\boldsymbol\varrho_{F^+_R})_\ast\phi_{F^+_R}\big)(\mathbf x').
\]
This is \hb{thm:node}(iii) for the family whose table $j$ is $\binom{\{g,g'\}}1\langle\cB^+_{j,B},\cB^+_{j,R}\rangle$, which is the table $(\cB_j,\theta_j)$ by \hb{prop:substitution}(iii).
\item \hb{thm:gluewidth}(iii): on the pins the seam word of the piece is $\prod_jK_{\cB_j/R}^{-a_{jD_0}}K_{\cB_j|B}^{-a_{jD_1}}\Psi_{\cB_j/R}^{-b_{jD_0}}\Psi_{\cB_j|B}^{-b_{jD_1}}$. This is the word of the bunch family $F^{\mathrm{bunch}}(\mathbf a,\mathbf b)$ of \hb{thm:twosides} at the level shares $(a^d_j,a^r_j,b^d_j,b^r_j)=(a_{jD_0},a_{jD_1},b_{jD_0},b_{jD_1})$, and the room word is that of $F^{\mathrm{room}}(\mathbf a,\mathbf b)$.
  \begin{itemize}
  \item The seam functions on the pins are the periods of these words (\hb{lem:seamword}), and
  \[
  \Gst21{a_{jD_0},\ a_{jD_1}}{s_j}\Gst21{b_{jD_0},\ b_{jD_1}}{t_j}=\Gst42{a^d_j,\ b^d_j,\ a^r_j,\ b^r_j}{s_j,\ t_j}.
  \]
  \item The scale pin $\sum_j(a_{jD_1}+b_{jD_0})=\kappa_B$ is the pin of \hb{thm:twosides}.
  \item In the notation of \hb{def:fibre}, eliminating the letter totals leaves the shares $(a^d_j,b^d_j)_j$ with the one pin $\sum_j(b^d_j-a^d_j)=\kappa_B-\sum_js_j$, with the same power of $2\pi$.
  \item In the pair coordinates $\sigma^d_j=\frac12(a^d_j+b^d_j)$ and $\delta^d_j=\frac12(a^d_j-b^d_j)$, Lebesgue measure gains the factor $2^N$, and the pin becomes $-2\sum_j\delta^d_j$, whose delta is $\frac12$ that of $\sum_j\delta^d_j$.
  \end{itemize}
  So the fibre integral of \hb{thm:gluewidth}(iii) is $2^{N-1}$ times the fibre integral in the pair coordinates, which is \hb{thm:twosides}.
\end{itemize}
\end{proof}

\begin{sremark}
The constants match exactly under the normalization. In a factorization with the constants $\hat\theta_{0j},\hat\theta_{1j}$ of Lemma~\ref{lem:ch27-normal}, the scale of the piece, its seam vector and its law are unchanged, since they are defined from the seam letters, while the hinge values move. Then (i) holds up to the constant factor $\prod_j(\hat\theta_{1j}/\hat\theta_{0j})^{\delta_j/\sigma}$, and the coordinate $\varrho^\rho_{jD_1}$ of (ii) is $(\hat\theta_{1j}/\hat\theta_{0j})\prod_l(\hat\theta_{1l}/\hat\theta_{0l})^{-\delta_l/\sigma}$ times $x_{j,K}/\psi(\mathbf x)$. Normalizing $\theta''_j$ on each excess removes these constants and costs no generality, by Lemma~\ref{lem:ch27-normal}(ii). For one table, $\mathbf x=(1,\varrho_B)$ and $\mathbf x'=(1,\varrho_R)$. The family $F_{\mathcal H}[(\mathbf x,\mathbf x')]$ is then the list $\binom{\{g,g'\}}1$ twisted in its cobasis letter by $\varrho_B$ at $g$ and by $\varrho_R$ at $g'$, and $\psi(\mathbf x)^{\kappa_B}\psi(\mathbf x')^{\kappa_R}=\varrho_B^{\kappa_Bt/(s+t)}\varrho_R^{\kappa_Rt/(s+t)}$. So (iv) is the identity with \hb{lem:twomarks}, and the reductions are \hb{thm:glue}(i) and (iv) together with \hb{thm:twosides}.
\end{sremark}

\begin{hproposition}[\hb{prop:relgrammaforeman} (The Gramma of a set of pairs in the Foreman's order)]
As stated, for a set of pairs with any number of parts.
\end{hproposition}

\begin{proof}
The configuration $\partial\mathfrak R$ of \hb{def:relgramma} has at the pair $(C,D)$ the column $\mathbf e_C+\mathbf e_D$ on the kept stations, the station not kept in each part read as $0$. A pair of $\mathfrak R_l$ has both its stations in $\mathcal K_l\sqcup\mathcal D_l$, so its column is supported on the kept stations of $\mathfrak R_l$, and different parts have different kept stations. By \hb{thm:grammaproduct}, applied to one part after another, the Gramma of $\mathfrak R$ is the product of the Grammas of its parts,
\[
\sG_{(\mathbf a,\mathbf a');\partial\mathfrak R}(\mathbf h)=\prod_{l=1}^{\nu}\sG_{(\mathbf a_{\mathcal K_l},\mathbf a'_{\mathcal D_l});\partial\mathfrak R_l}(\mathbf h),
\]
and it converges exactly when every factor does, its domain being the product of their domains. For one part the statement is the reference, proof of prop:relgrammaforeman. There the scales of $\mathcal K_l$ are integrated first, $\int_0^\infty x^{a_C-1}e^{-xL_C(v)}dx=\Gamma(a_C)L_C(v)^{-a_C}$. The form left on the stations of $\mathcal D_l$ is homogeneous of degree $|\mathbf a'_{\mathcal D_l}|-|\mathcal D_l|-|\mathbf a_{\mathcal K_l}|=-|\mathcal D_l|$, so its integral over a chart is its integral over $\Delta_{\mathcal D_l}$ (\hb{lem:chart}), and the convergence criterion is \hb{thm:poswordconv}. Multiplying over the parts gives the statement.
\end{proof}

\begin{hproposition}[\hb{prop:relgrammaclose} (Closing the Gramma of a set of pairs)]
As stated: any number of parts, and (ii) at real and complex loads.
\end{hproposition}

\begin{proof}
(i) On each part the reference's proof (proof of prop:relgrammaclose(i)) applies on the simplex $\Delta_{\mathcal D_l}$. On its interior $0\le\bar L_C<1$ almost everywhere, and $L_C^{-a_C}=M_C^{-a_C}\sum_n\frac{(a_C)_n}{n!}\bar L_C^{\,n}$ with nonnegative terms. On the simplex $\bar L_C$ is a linear function with nonnegative coefficients, and Tonelli's theorem with \hb{thm:dirmoments} gives the series. By \hb{prop:relgrammaforeman} the Gramma is $\prod_C\Gamma(a_C)$ times the product of these integrals over the parts.

(ii) The exchange relation of \hb{thm:grammalaws}(ii) holds at every index of the domain, real or complex.
\begin{itemize}
\item At a kept station of $\mathcal K$ or of $\mathcal D$ it is the stated relation, the column of $(C,D)$ being $\mathbf e_C+\mathbf e_D$ with the coefficient $h_{CD}$.
\item At the station not kept of a part it is the sum of the relations at the stations of $\mathcal K$ of that part less those at its other stations of $\mathcal D$, the totals of a part being equal (the reference, proof of prop:relgrammaclose(ii)).
\item A shifted load $(\mathbf a+\mathbf e_C,\mathbf a'+\mathbf e_{D'})$ with $(C,D')\in\mathfrak R$ keeps the totals of every part equal, since $C$ and $D'$ lie in one part, and it lies in the domain with $(\mathbf a,\mathbf a')$ (\hb{thm:grammalaws}).
\end{itemize}
Neither (i) nor (ii) uses one part, and (ii) does not use real loads.

(iii) concerns one part and is the reference, proof of prop:relgrammaclose(iii).
\end{proof}

\begin{htheorem}[\hb{thm:retainwidth} (The law is what a piece retains, at every width)]
As stated; the plain tinkers suffice under the stated condition on the tables with both letters active.
\end{htheorem}

\begin{proof}
The equivalence of (a)--(d) is the reference, proof of thm:retainwidth. Its step (c)$\Rightarrow$(d) reads as follows in the notation of \hb{def:fibre}.

\emph{The step (c)$\Rightarrow$(d).} Let $V$ be the real space of the directions of the pins of the seam (\hb{lem:seamword}): the vectors $\boldsymbol\vartheta=(\vartheta_{jD},\vartheta^\ast_{jD})_{j,D}$, one coordinate for each coordinate of the seam vector, with
\begin{gather*}
\sum_D\vartheta_{jD}=0\ \text{for every active }K_j,\qquad\sum_D\vartheta^\ast_{jD}=0\ \text{for every active }\Psi_j,\\
\sum_{j,D}\big(i(D)\vartheta_{jD}+(c_j-i(D))\vartheta^\ast_{jD}\big)=0.
\end{gather*}
For a tinker $\mathbf v$ of $\mathcal Y_0$ with the logarithms $\boldsymbol\ell=\log\mathbf v$, the seam word of $\mathcal Y_{0,\mathbf v}$ is $\mathbf v^{-(\mathbf a,\mathbf b)}$ times that of $\mathcal Y_0$. So, by \hb{lem:seamword} and by \hb{thm:gluewidth}(iii) on the contour through $\mathbf c$,
\[
e^{\ip{\mathbf c}{\boldsymbol\ell}}\,P_{\mathcal X\sqcup\mathcal Y_{0,\mathbf v}}=(2\pi)^{-n_{\mathcal D}}\int_Ve^{-i\ip{\boldsymbol\ell}{\boldsymbol\vartheta}}\,f(\boldsymbol\vartheta)\,d\boldsymbol\vartheta,
\]
where $f(\boldsymbol\vartheta)$ is the product of the stacked Gammas of \hb{thm:gluewidth}(iii) and of the seam functions of $\mathcal X$ and of $\mathcal Y_0$ at $\mathbf c+i\boldsymbol\vartheta$. The stacked Gammas have no zero, and $f$ is continuous and integrable on $V$ (the reference, proof of thm:retainwidth). The proof uses (c) only to know the Fourier transform of $f$ at every point of the dual of $V$; then $f$ is determined. The rest is unchanged: the stacked Gammas have no zero, the seam function of $\mathcal Y_0$ is positive at $\mathbf c$, the identity theorem applies, and the projection is injective on the seam torus.

\emph{The plain tinkers.} Their logarithms form the space $\Lambda$ of the $\boldsymbol\ell$ with $\ell_{jD}=\ell^\ast_{jD}$ for every table $j$ with both letters active (\hb{def:tinkers}). A table with one active letter has one coordinate for each station, and on it a plain tinker is a general one. The restrictions to $V$ of the functionals $\boldsymbol\vartheta\mapsto\ip{\boldsymbol\ell}{\boldsymbol\vartheta}$, $\boldsymbol\ell\in\Lambda$, exhaust the dual of $V$ exactly when no nonzero $\boldsymbol\vartheta\in V$ is orthogonal to $\Lambda$.

\emph{No orthogonal vector.} Let $\boldsymbol\vartheta\in V$ be orthogonal to $\Lambda$.
\begin{itemize}
\item Then $\vartheta_{jD}=0$ in a table with only $K_j$ active, $\vartheta^\ast_{jD}=0$ in a table with only $\Psi_j$ active, and $\vartheta^\ast_{jD}=-\vartheta_{jD}$ in a table with both letters active.
\item The letter totals give $\sum_D\vartheta_{jD}=0$ in every table with both letters active. The scale pin becomes $\sum_{j,D}\big(2i(D)-c_j\big)\vartheta_{jD}=0$, the sum running over the tables with both letters active.
\item A table with $|\mathcal D_j|=1$ has $\vartheta_{jD}=0$ by its total.
\item The exceptional table has $c_j\ge1$, and $\mathcal D_{j,0}$ and $\mathcal D_{j,c_j}$ are nonempty (\hb{def:seampiece}), so its two stations $D'$ and $D''$ have the excesses $0$ and $c_j$. Its total gives $\vartheta_{jD'}=-\vartheta_{jD''}=\tau$, and the scale pin reads $-c_j\tau-c_j\tau=0$, so $\tau=0$.
\end{itemize}
Hence $\boldsymbol\vartheta=0$, and the plain tinkers suffice in (c).
\end{proof}

\begin{sremark}
In general, let $J_2$ be the tables with both letters active. The orthogonal of $\Lambda$ in $V$ has the dimension $\sum_{j\in J_2}(|\mathcal D_j|-1)$, less one when some table of $J_2$ has positive width. Indeed, the vectors with $\sum_D\vartheta_{jD}=0$ on the tables of $J_2$ and zero elsewhere form a space of the first dimension, and on it the scale pin $\sum_{j,D}(2i(D)-c_j)\vartheta_{jD}$ vanishes identically exactly when every table of $J_2$ has width zero. With every letter active and some table of positive width, $n_{\mathcal D}=2|\mathcal D|-2N-1$ and the orthogonal has the dimension $|\mathcal D|-N-1$, so the plain tinkers of one partner move $|\mathcal D|-N$ of the directions of the seam torus, as the paragraph after the theorem says.
\end{sremark}

\begin{hlemma}[\hb{lem:mot-probe} (Partners on one mark see a point mass)]
As stated.
\end{hlemma}

\begin{proof}
\emph{The one-mark piece.} On one mark the seam letters have $d_j=d^\ast_j=0$, the wall is the mass $\kappa$, and by the relations of the seam torus the excess sums are $V^{(0)}_j=\sum_{D\in\mathcal D_{j,0}}x_{jD}=1$, $V^{\ast(c_j)}_j=1$, $V^{(c_j)}_j=\alpha_g^{\,c_j}\sum_{D\in\mathcal D_{j,c_j}}x_{jD}$ and $V^{\ast(0)}_j=\alpha_g^{\,c_j}\sum_{D\in\mathcal D_{j,0}}x^\ast_{jD}$. So the scale and the seam vector of \hb{def:seamlaw} are
\begin{gather*}
\alpha_g\Big(\prod_j\Big(\sum_{D\in\mathcal D_{j,c_j}}x_{jD}\Big)^{s_j}\Big(\sum_{D\in\mathcal D_{j,0}}x^\ast_{jD}\Big)^{t_j}\Big)^{1/w}=\alpha_g,\\
\Big(\frac{x_{jD}\,\alpha_g^{\,i(D)}}{\alpha_g^{\,i(D)}},\ \frac{x^\ast_{jD}\,\alpha_g^{\,c_j-i(D)}}{\alpha_g^{\,c_j-i(D)}}\Big)_{j,D}=\mathbf x,
\end{gather*}
and the scale form on the one point of $\Delta_{\{g\}}$ has the mass $1$: the law is $\delta_{\mathbf x}$, and the station transform is the kernel at $\mathbf x$ (\hb{def:seamtransform}). In the chart $\alpha_b=1$ the period of the glue of $\mathcal X$ with the partner on $b$ is $\int_{(0,\infty)^B}\alpha^{\mvec_B-\one}\prod_j(\sum_Dp_{jD}V_{jD})^{-s_j}(\sum_Dp^\ast_{jD}V^\ast_{jD})^{-t_j}d\alpha$, the second expression of \hb{def:seamtransform}; the partner has the degrees $e_j=e^\ast_j=0$, the wall $w-\kappa$, and its room letters are those of the partner with the coefficients $1$ tinkered by $\mathbf p$ (\hb{def:tinkers}).

\emph{Coordinates.} Index the coordinates of the seam vector by the pairs of an active letter $l$ and a station of its table, $n_{\mathcal X}$ of them; write $\mathbf p=e^{\mathbf P}$ with $\mathbf P\in\mathbb R^{n_{\mathcal X}}$, and pair shares $(\mathbf a,\mathbf b)\in\mathbb R^{n_{\mathcal X}}$, one for each coordinate, with $\mathbf P$ by $\ip{(\mathbf a,\mathbf b)}{\mathbf P}$. Let $\mathbf u_l$ be the indicator of the coordinates of the letter $l$, $U'$ the span of the $\mathbf u_l$, $\mathbf u_0$ the vector with the entries $i(D)$ at the coordinates of the letters $K_j$ and $c_j-i(D)$ at those of the letters $\Psi_j$, and $U=U'+\mathbb R\mathbf u_0$. Since $w>0$ some active letter has $c_j>0$, and on its coordinates $\mathbf u_0$ takes the values $0$ and $c_j$ at the excesses $0$ and $c_j$, which have stations; so $\mathbf u_0\notin U'$. The pins of the seam (\hb{lem:seamword}) are $\ip{(\mathbf a,\mathbf b)}{\mathbf u_l}$ equal to the exponent of $l$ and $\ip{(\mathbf a,\mathbf b)}{\mathbf u_0}=\kappa$; let $\mathcal A$ be the set of the shares with positive coordinates on the pins. For one letter of width $c_j>0$ and exponent $s$, the sums $\sum_Di(D)a_{jD}$ over positive shares with $\sum_Da_{jD}=s$ fill $(0,sc_j)$, and they are $0$ at width zero; so on the pins the scale pin takes every value in $(0,w)$, and $\mathcal A$ is not empty, open in the affine space of the pins, whose direction space is $V=U^\perp$.

\emph{The station transform opens the seam function.} Fix $(\mathbf a,\mathbf b)\in\mathcal A$ and put $f(\mathbf P)=e^{\ip{(\mathbf a,\mathbf b)}{\mathbf P}}\,\Gamma(\kappa)\sZ_{(\kappa,\mathbf s,\mathbf t);t_{\mathcal D}}(e^{\mathbf P})$. By the kernel of \hb{thm:gluewidth}(ii) with $v=e^\tau$,
\[
f(\mathbf P)=\int_{\mathbb R}g(\mathbf P+\tau\mathbf u_0)\,d\tau,\qquad g(\mathbf Q)=e^{\ip{(\mathbf a,\mathbf b)}{\mathbf Q}}\prod_{j:\,s_j>0}\Big(\sum_De^{Q_{jD}}\Big)^{-s_j}\prod_{j:\,t_j>0}\Big(\sum_De^{Q^\ast_{jD}}\Big)^{-t_j},
\]
using $\ip{(\mathbf a,\mathbf b)}{\mathbf u_0}=\kappa$. Translation by $\mathbf u_l$ multiplies the letter $l$ by $e^\lambda$ and the exponential by $e^{\lambda s}$, $s$ its exponent: $g$ is invariant under $U'$, and $f$ under $U$, a translation by $\mathbf u_0$ being absorbed by $\tau$. Let $W$ be a linear complement of $U$ and $W'$ the subspace on which the coordinate of one chosen station $D_l$ of every active letter vanishes, a complement of $U'$. Since $\mathbf u_0\notin U'$, $(\mathbf P,\tau)\mapsto\mathbf P+\tau\mathbf u_0$ is a linear isomorphism of $W\times\mathbb R$ onto $\mathbb R^{n_{\mathcal X}}/U'\cong W'$, and by Tonelli and the invariance of $g$
\[
\int_Wf\,d\mathbf P=c_W\int_{W'}g=c_W\prod_{j:\,s_j>0}\Gst{|\mathcal D_j|}1{\mathbf a_j}{s_j}\prod_{j:\,t_j>0}\Gst{|\mathcal D_j|}1{\mathbf b_j}{t_j}=:\beta(\mathbf a,\mathbf b)\in(0,\infty),
\]
with $c_W>0$ fixed by $W$, $W'$ and their Lebesgue measures. On $W'$ the integral factors over the letters; for a letter $K_j$, with $y_D=e^{Q_{jD}}$ and $a_{jD_l}=s_j-\sum_{D\neq D_l}a_{jD}>0$, it is $\int_{(0,\infty)^{\mathcal D_j\setminus D_l}}\prod_{D\neq D_l}y_D^{a_{jD}-1}\big(1+\sum_{D\neq D_l}y_D\big)^{-s_j}d\mathbf y$, which is the stacked Gamma, as $(1+|\mathbf y|)^{-s_j}=\Gamma(s_j)^{-1}\int_0^\infty x^{s_j-1}e^{-x(1+|\mathbf y|)}dx$ and Tonelli integrate it to $\prod_D\Gamma(a_{jD})/\Gamma(s_j)$; likewise for a letter $\Psi_j$. For $\mathbf x$ on the seam torus and $\mathbf X=\log\mathbf x$, $e^{\ip{(\mathbf a,\mathbf b)}{\mathbf P}}\Gamma(\kappa)\sZ(\mathbf x\circ e^{\mathbf P})=e^{-\ip{(\mathbf a,\mathbf b)}{\mathbf X}}f(\mathbf P+\mathbf X)$, and $\int_Wf(\mathbf P+\mathbf X)\,d\mathbf P=\int_Wf\,d\mathbf P$ by the invariance of $f$ under $U$, Lebesgue measure on $W$ being invariant under translation. By the first expression of \hb{def:seamtransform} and Tonelli,
\[
\int_We^{\ip{(\mathbf a,\mathbf b)}{\mathbf P}}\,\mathcal I_{\mathcal X}(e^{\mathbf P})\,d\mathbf P=\beta(\mathbf a,\mathbf b)\int\mathbf x^{-(\mathbf a,\mathbf b)}\,d\mu(\mathbf x),
\]
$\mu$ the law of $\mathcal X$. The hypothesis makes the left side $c\,\beta(\mathbf a,\mathbf b)\,\mathbf x_0^{-(\mathbf a,\mathbf b)}$, so $\int\mathbf x^{-(\mathbf a,\mathbf b)}d\mu=c\,\mathbf x_0^{-(\mathbf a,\mathbf b)}$ for every $(\mathbf a,\mathbf b)\in\mathcal A$: the seam function is finite on all of $\mathcal A$.

\emph{A point mass.} Fix $(\mathbf a_0,\mathbf b_0)\in\mathcal A$ and let $\nu=\mathbf x^{-(\mathbf a_0,\mathbf b_0)}\mu$, a finite measure of the total mass $C=c\,\mathbf x_0^{-(\mathbf a_0,\mathbf b_0)}$. For $\boldsymbol\delta\in V$ near $0$, $(\mathbf a_0,\mathbf b_0)\pm\boldsymbol\delta\in\mathcal A$, and so $\int(\mathbf x/\mathbf x_0)^{\mp\boldsymbol\delta}d\nu=C$ and
\[
\int\Big((\mathbf x/\mathbf x_0)^{\boldsymbol\delta/2}-(\mathbf x/\mathbf x_0)^{-\boldsymbol\delta/2}\Big)^2d\nu(\mathbf x)=C-2C+C=0 .
\]
So $\ip{\log\mathbf x-\log\mathbf x_0}{\boldsymbol\delta}=0$ for $\nu$-almost every $\mathbf x$, for every $\boldsymbol\delta$ of a basis of $V$, and $\log\mathbf x-\log\mathbf x_0\in V^\perp=U$ for $\mu$-almost every $\mathbf x$, $\mu$ and $\nu$ having the same null sets. Two points $\mathbf x,\mathbf x'$ of the seam torus with $\log\mathbf x-\log\mathbf x'\in U$ are equal: $\log x_{jD}-\log x'_{jD}=\lambda_j+\lambda i(D)$ and $\log x^\ast_{jD}-\log x'^\ast_{jD}=\lambda^\ast_j+\lambda(c_j-i(D))$; the relations at the excess $0$ of the letters $K_j$ and at the excess $c_j$ of the letters $\Psi_j$ give $\lambda_j=\lambda^\ast_j=0$, and then the product relation gives $e^{\lambda w}=1$, so $\lambda=0$. Hence $\mu$ is carried by $\mathbf x_0$, and $\mu=c\,\delta_{\mathbf x_0}$.
\end{proof}

\begin{htheorem}[\hb{thm:storeone} (A bunch is stored on one mark exactly at width one)]
As stated, the clause at width one read with the normalization of \hb{prop:widthonelaw}.
\end{htheorem}

\begin{proof}
Everything except the clause at width one is the reference, proof of thm:storeone. At width one with factoring weights, Lemma~\ref{lem:ch27-normal}(i) writes the letters of the partner with its stand-in at a seam point $\mathbf x$, where $x_{jD_0}=x^\ast_{jD_1}=1$, as
\[
\frac1{\hat\theta_0}\Big(K_{\cB_j|R}+\frac{\hat\theta_0}{\hat\theta_1}\,x_{jD_1}\,\alpha_g\,K_{\cB_j/B}\Big),\qquad\frac1{\hat\theta_1}\Big(\frac{\hat\theta_1}{\hat\theta_0}\,x^\ast_{jD_0}\,\alpha_g\,\Psi_{\cB_j|R}+\Psi_{\cB_j/B}\Big),
\]
the constants $\hat\theta_0,\hat\theta_1$ being those of table $j$. Under the normalization of \hb{prop:widthonelaw} the constants are $1$. The letters are then those of $\cB^+_{j,R}$ twisted at $g$ by $x_{jD_1}$ in $K_j$ and by $x^\ast_{jD_0}$ in $\Psi_j$ (\hb{thm:widthone}, \hb{def:twist}). At the seam point of a value $\mathbf y$ of the ratio vector of $F^+_B$ these twists are $y_{j,K}/\psi(\mathbf y)$ and $y_{j,\Psi}/\psi(\mathbf y)$ (\hb{prop:widthonelaw}(ii)), and the change of variable $\alpha_g\mapsto\psi(\mathbf y)\alpha_g$ of \hb{prop:widthonelaw}(iii) turns them into the twist by $\mathbf y$, as stated. Without the normalization the two letters carry the constant factors $1/\hat\theta_0$ and $1/\hat\theta_1$ and the twists $(\hat\theta_0/\hat\theta_1)x_{jD_1}$ and $(\hat\theta_1/\hat\theta_0)x^\ast_{jD_0}$. The last sentence of the clause is \hb{prop:widthonelaw}(i)--(iii).
\end{proof}

\begin{htheorem}[\hb{thm:mot-pointmass} (A point mass is a pinned set)]
As stated.
\end{htheorem}

\begin{proof}
\emph{The walls.} On the balance wall $\kappa_S+\kappa_{S'}=2\sum_j\sigma_jc_j(S)=w$ (\hb{prop:compoundwall}), and both walls are positive; so $w>0$ and $0<\kappa_S<w$. By the paragraph after \hb{def:seampiece}, $S$ is a piece with the seam letters $V_{jD}=K_{j,(D)}$ and $V^\ast_{jD}=\Psi_{j,(D)}$ of the stations $D$ of its room, $S'$ is a partner, and their glue is $F$ (\hb{prop:seamidentity}(iii)); the wall of the piece is $\kappa_S$ and its top wall $w-\kappa_S=\kappa_{S'}$.

\emph{(i) and the seam vector.} The scale form has a positive continuous density on the open simplex $\Delta_S$, the excess sums and the scale being positive there, and the seam vector is continuous. So the law is carried by one point exactly when the seam vector is constant, $\mathbf x_0$ say, and then it is $c\,\delta_{\mathbf x_0}$ with $c=\int_{\Delta_S}\phi^\rho_S>0$. By \hb{thm:gluewidth}(i) with the partner $S'$, $P_F=c\,P_{S'^{(\mathbf x_0)}}$ in $[0,\infty]$; $P_F$ is finite and $P_{S'^{(\mathbf x_0)}}>0$, so $c$ and $P_{S'^{(\mathbf x_0)}}$ are finite.

\emph{(ii)$\Rightarrow$(i).} An active table of width zero at $S$ with its weights factoring across $S$ and $S'$ has one station of room sides (\hb{prop:ladder}(i)), whose coordinates of the seam vector are $V_{jD}/V^{(0)}_j=1$ and $V^\ast_{jD}/V^{\ast(c_j)}_j=1$. An active table of width one with factoring weights has the stations $D_0$ and $D_1$ and the seam letters $K_{\cB_j/S'},K_{\cB_j|S}$ and $\Psi_{\cB_j/S'},\Psi_{\cB_j|S}$ (\hb{prop:ladder}(ii), the paragraph after \hb{def:seamlaw}), so its coordinates are $1$, $\gamma^K_{j,S}/\gamma^\rho_S$, $\gamma^\Psi_{j,S}/\gamma^\rho_S$ and $1$, the hinge values read at the factorization that normalizes the multipliers. Under (ii) the scale is $\gamma^\rho_S=\prod_j(c_j\gamma)^{s_j/w}(c^\ast_j\gamma)^{t_j/w}=C\gamma$, the product over the tables of width one with $C=\prod_jc_j^{s_j/w}c_j^{\ast\,t_j/w}$, and every coordinate is a constant: $c_j/C$, $c^\ast_j/C$ or $1$.

\emph{(i)$\Rightarrow$(ii).} Let the seam vector be the constant $\mathbf x_0$. For every active letter $K_j$, $V_{jD}=x_{0,jD}V^{(0)}_j(\gamma^\rho_S)^{i(D)}$, and by the room-station factorization (\hb{prop:seamidentity}(iii))
\[
K_j=\sum_DV_{jD}\,\mathsf y_{jD}=V^{(0)}_j(\alpha_S)\sum_Dx_{0,jD}\,\gamma^\rho_S(\alpha_S)^{i(D)}\,\mathsf y_{jD}(\alpha_{S'}),
\]
so $K_j$ reads $S$ through the one function $\gamma^\rho_S$, positive and homogeneous of degree one, with $V^{(0)}_j$ positive and homogeneous: by \hb{lem:mot-onehinge}, $c_j(S)\le1$. Likewise for every active letter $\Psi_j$, with $V^{\ast(c_j)}_j$ and $(\gamma^\rho_S)^{c_j-i(D)}$. Two stations $D\neq D'$ of one excess would have proportional seam letters, $V_{jD}/V_{jD'}=x_{0,jD}/x_{0,jD'}$ (or the same with $V^\ast$), so proportional room children, against \hb{def:seam}(i); so every excess of every active table has one station of room sides, hence one station of patterns (\hb{prop:ladder}), and is one rectangle, which at the extreme excesses, and so at widths zero and one, means that the weights factor across $S$ and $S'$ (\hb{prop:ladder}(i),(ii)). At width one the coordinates above are constant, so the active hinge values are positive multiples of $\gamma^\rho_S$: they read $S$ alike, with $\gamma=\gamma^\rho_S$.

\emph{(i)$\Rightarrow$(iii)$\Rightarrow$(iv)$\Rightarrow$(i).} The first is \hb{thm:gluewidth}(i) with the law $c\,\delta_{\mathbf x_0}$. For the second take the partner on one mark $b$ with the room letters $p_{jD}\alpha_b^{\,c_j-i(D)}$ and $p^\ast_{jD}\alpha_b^{\,i(D)}$ and the mass $w-\kappa_S$: the period of its glue with $S$ is the station transform at $\mathbf p$ (\hb{lem:mot-probe}), and the partner with its stand-in at $\mathbf x_0$ is the glue of the one-mark piece at $\mathbf x_0$ with it, whose period is the kernel $\Gamma(\kappa_S)\sZ_{(\kappa_S,\mathbf s,\mathbf t);t_{\mathcal D}}(\mathbf x_0\circ\mathbf p)$ (\hb{thm:gluewidth}(ii)). The third is \hb{lem:mot-probe}.

\emph{Lists.} With the weights $1$ the seam letters are the letters themselves and the multipliers are $1$ (\hb{def:seam}). Let two active hinge values of tables of width one be proportional, $\gamma^K_{j,S}=\mu\,\gamma^\Psi_{j',S}$ say, that is $K_{\cB_j|S}\Psi_{\cB_{j'}|S}=\mu\,K_{\cB_j/S'}\Psi_{\cB_{j'}/S'}$ (\hb{def:readings}). Along a maximal chain of subsets of $S$ the terms of least degree of each of these four letters, at every set of the chain at once, form one monomial with the coefficient $1$ (\hb{prop:mot-chainbottom}); for a product of two polynomials with nonnegative coefficients the least degrees add and the terms of least degree are the products of those of the factors, so the two sides have such terms $1\cdot\alpha^{\mathbf a}$ and $\mu\,\alpha^{\mathbf a'}$, and $\mu=1$. The other pairs of letters are alike. So reading alike is equality, every $c_j$ and $c^\ast_j$ is $1$, $C=1$, and $\mathbf x_0=\one$.

\emph{One list.} With both letters active, $w=(s+t)c(S)>0$, so pinned means $c(S)=1$ and $\gamma^K_S$ and $\gamma^\Psi_S$ equal, by the above, that is $\varrho_S=1$, which is $\mathbf E_{\cB|S}=\mathbf E_{\cB/S'}$ (\hb{def:readings}). The cell of the bunch $\cB^+_S$ is a list with the hinge minors $\cB|S$ and $\cB/S'$ (\hb{thm:widthone}) and the ratio $\varrho_S$ (\hb{def:substitution}).
\end{proof}

\begin{hcorollary}[\hb{cor:mot-whatsplits} (Which sets split off)]
As stated.
\end{hcorollary}

\begin{proof}
(i) A pinned set has width at most one in every active table; apply \hb{thm:mot-pointmass}. (ii) A prime table is connected, so a set with at least two marks on each side has positive width (\hb{lem:separator}), and it has no bunch of width one with two marks on each side (\hb{thm:primetree}); so such a set has width at least two, and (i) applies. (iii) The set $S=\cE\setminus\{e\}$ has the width $c(\{e\})=r(\{e\})-d(\{e\})=1$ (\hb{thm:window}), and one letter is active, so $S$ is pinned with $\gamma=\gamma^K_S$, and it splits off. At $(s,0)$, $d(\{e\})=0$ gives $\kappa_{\{e\}}=m_e$ and $\kappa_S=s-m_e$, as $\kappa_S+\kappa_{\{e\}}=s$ (\hb{prop:compoundwall}). \hb{cor:pinnedproduct} with the bunch $S$ and the room $\{e\}$, at the level shares $(s-\kappa_S,0,\kappa_S,0)$, gives the Beta corner $\Gst21{s-m_e,\ m_e}{s}$ times the period of the bunch family, the tables $\cB/\{e\}=\cB/e$ at $(m_e,0)$ and $\cB|S=\cB\setminus e$ at $(s-m_e,0)$ (\hb{def:minors}), times the period of the room family on one mark, which is $1$. At $(0,t)$, $d^\ast(\{e\})=0$ gives $\kappa_{\{e\}}=m_e$, $\kappa_S=t-m_e$, and the level shares $(0,\kappa_S,0,t-\kappa_S)$ give the stated family. (iv) $S=\cE\setminus\{e\}$ has width one and $|S|\ge2$, and its hinge minors are $\cB|S=\cB\setminus e$ and $\cB/\{e\}=\cB/e$; apply the last part of \hb{thm:mot-pointmass}.
\end{proof}

\section{Chapter~\hb{ch:composebunches}: composing bunches}\label{sec:ch28}

\begin{htheorem}[\hb{thm:lawreadings} (Two transforms of a law, and their rules)]
As stated; the first rule of (iii) holds at any number of sub-pieces.
\end{htheorem}

\begin{proof}
Parts (i), (ii) and (iv), and the second and third rules of (iii), are the reference, proof of thm:lawreadings. For the first rule:
\begin{itemize}
\item Let $S_1,\dots,S_k$ be as in \hb{thm:nestedlaw}, and $\mathbf p$ a point with positive coordinates, one for each coordinate of the seam vector of $S$.
\item By \hb{def:seamtransform}, $\mathcal I_S(\mathbf p)=\int H\,d\lambda_S$, with the law $\lambda_S=(\boldsymbol\varrho^\rho_S)_\ast\phi^\rho_S$ and the Borel function $H(\mathbf z)=\Gamma(\kappa_S)\,\sZ_{(\kappa_S,\mathbf s,\mathbf t);t_{\mathcal D}}(\mathbf z\circ\mathbf p)\ge0$, $\mathcal D$ being the stations of room sides of $S$.
\item \hb{thm:nestedlaw}(ii) gives $\int H\,d\lambda_S=\int\big(\int H\,d\lambda_{\mathbf x}\big)\prod_ld\lambda_{S_l}(\mathbf x_l)$, with $\lambda_{\mathbf x}$ the law of the node $S^{(\mathbf x_1,\dots,\mathbf x_k)}$.
\item The node is a piece with the stations $\mathcal D$, the exponents of the family and the wall $\kappa_S$ (\hb{thm:nestedlaw}(ii)), so by the same definition $\int H\,d\lambda_{\mathbf x}=\mathcal I_{S^{(\mathbf x_1,\dots,\mathbf x_k)}}(\mathbf p)$.
\end{itemize}
This is the reference's argument for one sub-piece, which does not use $k=1$. The second and third rules use the transfer matrices from one sub-piece: \hb{thm:nestedlaw}(i) gives linearity in the seam letters of each $S_l$ separately, with coefficients that depend on the marks of the other sub-pieces.
\end{proof}

\begin{htheorem}[\hb{thm:cyclicchain} (Cyclic chains)]
As stated: (i) holds for a family, and in (ii) each letter is opened over the starting states at which its term is not zero.
\end{htheorem}

\begin{proof}
(i) Fix a table $j$.
\begin{itemize}
\item By \hb{thm:hingepairs}(ii), the bases of table $j$ are the unions of the unpaired parts of the choices of one basis in every cell $\mathcal A_{j,i}$ that occupy exactly one hinge of every pair. Its weight in each letter at a basis is the sum, over the choices that give it, of the products of the weights chosen.
\item A pair of $I_{i-1}$ occupied on the side of the cell $i-1$ is free on the side of the cell $i$, and conversely. So, with $J_i\subseteq I_i$ the pairs occupied on the side of the cell $i$, the basis chosen in $\mathcal A_{j,i}$ occupies the hinges $(I_{i-1}\setminus J_{i-1})\sqcup J_i$.
\item Its part $X_i\subseteq B_i$ is a basis of the child of the bunch of its hinges at that pattern. With their weights, these choices sum to the channel $K_{\mathcal A_{j,i},((I_{i-1}\setminus J_{i-1})\sqcup J_i)}(\alpha_{B_i})$ (\hb{def:channels}).
\item A basis $\bigsqcup_iX_i$ leaves the cobasis $\bigsqcup_i(B_i\setminus X_i)$, so $\Psi_j$ is summed the same way with the channels of $\Psi$.
\item Summing over the sequences $(J_0,\dots,J_{k-1})$, $J_k=J_0$, gives $K_j=\operatorname{tr}\prod_i\mathbf M_{K_j,i}$ and $\Psi_j=\operatorname{tr}\prod_i\mathbf M_{\Psi_j,i}$. The entries are polynomials in $\alpha_{B_i}$ with nonnegative coefficients and no $\varepsilon$, which is the top layer $0$.
\item The pairs are named alike in every table, so the sets of states $\mathcal J_i=\{J:\ J\subseteq I_i\}$ serve every letter, as \hb{def:cyclicchain} asks.
\end{itemize}
This is the reference's proof (proof of thm:cyclicchain(i)) applied table by table.

(ii) and (iii) are the reference, proof of thm:cyclicchain. In (ii), $L$ is the sum over the starting states $J$ of the terms $\mathbf e_{(J,0)}^{\mathsf T}\prod_i\widetilde{\mathbf M}_{L,i}\,\mathbf e_{(J,h)}$, each a sum over the paths from $(J,0)$ to $(J,h)$ of products of entries with nonnegative coefficients. \hb{thm:E1} opens a sum of positive terms, one share for each. A starting state whose term vanishes, the state of no path, therefore carries no share and no flow.
\end{proof}

\section{Chapter~\hb{ch:fusion}: Foreman fusion}\label{sec:chfusion}

The proofs use the rank function of a table, $r(Y)=\max_T|T\cap Y|$ (Definition~\hb{def:molecule}), through its standard properties: $r(\varnothing)=0$, $r(Y)\le r(Y+e)\le r(Y)+1$ and $r(Y\cup Y')+r(Y\cap Y')\le r(Y)+r(Y')$; conversely a function on the subsets of a finite set with these three properties is the rank function of the table whose bases are the sets $Y$ with $|Y|=r(Y)$ equal to the value at the whole set (the rank axioms, which hold for the bases of Definition~\hb{def:table} by the arena property, the exchange property of the bases). A set is \emph{closed} when it is maximal for its rank, that is when adding any mark raises its rank; the \emph{closure} of $Y$ is the set of the marks whose addition does not raise the rank of $Y$, it is closed, and it lies in every closed set that contains $Y$. A mark $e\notin Y$ lies in the closure of $Y$ exactly when some circuit $C$ has $e\in C\subseteq Y+e$: a basis $I$ of $Y$ with $I+e$ dependent contains, with $e$, a circuit through $e$, and a circuit $C$ through $e$ inside $Y+e$ puts $e$ in the closure of $C-e$. For a minor contracting an independent set $C$ and deleting $D$, the rank of a set $Y$ of the remaining marks is $r(Y\cup C)-r(C)$ (Definition~\hb{def:minors}). The dual rank function is $r^\ast(Y)=|Y|+r(\cE\setminus Y)-r$ (Definition~\hb{def:skeleton}), and $d(Y)=|Y|-r^\ast(Y)$.

\begin{hdefinition}[\hb{def:quotientlift} (Quotients; lifts; the level cell of a pair)]
As stated.
\end{hdefinition}

\begin{proof}
That the lifts and the level cell are tables is Proposition~\hb{prop:paircell}(ii) and (iii). For a bunch $B$ with the room $R$, $\cB/R$ is a quotient of $\cB|B$: for $X\subseteq Y\subseteq B$, submodularity gives $r(Y)-r(X)\ge r(Y\cup R)-r(X\cup R)$, which is the inequality of Proposition~\hb{prop:paircell}(i) for the rank functions $r(\cdot)$ of $\cB|B$ and $r(\cdot\cup R)-r(R)$ of $\cB/R$. The independent sets of $\mathcal L_j(\cB|B,\cB/R)$ are the $X$ with $|X|\le r(X)$ and $|X|\le r_{\cB/R}(X)+j$, the independent sets of Proposition~\hb{prop:higgs}(ii) for $i=j$, so the two lifts are one table, and the level cells of Definition~\hb{def:quotientlift} and Lemma~\hb{lem:higgsmajor} have the same bases. Level form at $G$ is Proposition~\hb{prop:levelform}: a transposition of two marks of $G$ fixes $X$ and $|Q|$.
\end{proof}

\begin{hproposition}[\hb{prop:paircell} (The level cell of a pair)]
As stated.
\end{hproposition}

\begin{proof}
(i) Let the inequality hold, $A$ be closed in $\mathcal N$ and $e\notin A$. Then $r_{\mathcal N}(A+e)=r_{\mathcal N}(A)+1$, the inequality at $X=A$, $Y=A+e$ gives $r_{\mathcal M}(A+e)=r_{\mathcal M}(A)+1$, and $A$ is closed in $\mathcal M$. Conversely, let every closed set of $\mathcal N$ be closed in $\mathcal M$. Along a chain from $X$ to $Y$ adding one mark at a time it suffices to take $Y=X+e$. If $r_{\mathcal N}(X+e)=r_{\mathcal N}(X)$ the inequality holds, $r_{\mathcal M}$ being monotone. Otherwise $e$ is not in the closure $A$ of $X$ in $\mathcal N$; $A$ is closed in $\mathcal M$ and contains $X$, so the closure of $X$ in $\mathcal M$ lies in $A$ and avoids $e$, and $r_{\mathcal M}(X+e)=r_{\mathcal M}(X)+1$. At $X=\varnothing$, $Y=\cE$ the inequality gives $c\ge0$; at $X=\varnothing$ it gives $r_{\mathcal M}(Y)\ge r_{\mathcal N}(Y)$, so a set independent in $\mathcal N$ is independent in $\mathcal M$; at $Y=\cE$ it gives $r_{\mathcal N}(\cE)-r_{\mathcal N}(X)\le r_{\mathcal M}(\cE)-r_{\mathcal M}(X)$, so a set spanning $\mathcal M$ spans $\mathcal N$; and two such inequalities add.

(ii) Put $f_1(Y\sqcup P)=r_{\mathcal M}(Y)+|P|$ and $f_2(Y\sqcup P)=r_{\mathcal N}(Y)+c$. Each is monotone, rises by at most one when a mark is added, and is submodular, being a rank function plus a function that is modular; $f_1(\varnothing)=0\le c=f_2(\varnothing)$; and $f_1-f_2=r_{\mathcal M}(Y)-r_{\mathcal N}(Y)+|P|-c$ is nondecreasing under inclusion by (i). The minimum of two submodular functions with a nondecreasing difference is submodular: for sets $Z,Z'$ the cases where one function attains the minimum at both are its own submodularity, and if $\min(Z)=f_1(Z)$ and $\min(Z')=f_2(Z')$, then
\[
f_1(Z)+f_2(Z')-f_1(Z\cap Z')-f_2(Z\cup Z')\ge\big(f_2(Z)-f_2(Z\cap Z')\big)-\big(f_2(Z\cup Z')-f_2(Z')\big)\ge0,
\]
the first step because $(f_1-f_2)(Z)\ge(f_1-f_2)(Z\cap Z')$ and the second by the submodularity of $f_2$, so that $\min(Z)+\min(Z')\ge f_1(Z\cap Z')+f_2(Z\cup Z')\ge\min(Z\cap Z')+\min(Z\cup Z')$; the case with $f_1$ and $f_2$ exchanged is the same with the roles of $Z$ and $Z'$ exchanged. So $r^\natural=\min(f_1,f_2)$ satisfies the rank axioms. Its value at $\cE\sqcup G$ is $\min(r_{\mathcal M}+c,r_{\mathcal N}+c)=r_{\mathcal M}$, and $Y\sqcup P$ with $|Y|+|P|=r_{\mathcal M}$ is a basis exactly when $r_{\mathcal M}(Y)\ge|Y|$ and $r_{\mathcal N}(Y)\ge|Y|+|P|-c=r_{\mathcal N}$: $Y$ independent in $\mathcal M$ and spanning $\mathcal N$, with $|Y|=r_{\mathcal N}+c-|P|$. For $Q\subseteq G$, $r^\natural(Q)=|Q|$, so the minor contracting $Q$ and deleting $G\setminus Q$ has the rank function
\[
Y\mapsto r^\natural(Y\sqcup Q)-|Q|=\min\big(r_{\mathcal M}(Y),\ r_{\mathcal N}(Y)+c-|Q|\big),
\]
that of $\mathcal L_{c-|Q|}$, provided some basis contains $Q$ and avoids $G\setminus Q$. Such bases exist for every $Q$: a basis $Y_0$ of $\mathcal N$ is independent in $\mathcal M$ by (i), it extends to a basis of $\mathcal M$ of $r_{\mathcal N}+c$ marks, and every set between the two is independent in $\mathcal M$ and spans $\mathcal N$. So every $Q$ is a pattern, $d(G)=0$, $r(G)=c$ and $c(G)=c$, and the child at $Q$ is $\mathcal L_{c-|Q|}$ (Lemma~\hb{lem:minor}). At $Q=\varnothing$ the rank function is $r_{\mathcal M}$, since $r_{\mathcal M}(Y)\le r_{\mathcal N}(Y)+c$ by (i) at $X=Y$ and $\cE$; at $Q=G$ it is $r_{\mathcal N}$, since $r_{\mathcal N}\le r_{\mathcal M}$. The rank function depends on $P$ through $|P|$ alone, so every permutation of $G$ maps the bases onto themselves.

(iii) By (ii) $\mathcal L_j$ is the child at any $Q$ of $c-j$ marks, a table, with the stated rank and bases, and $\mathcal L_0=\mathcal N$, $\mathcal L_c=\mathcal M$. For $j\le j'$ take $Q'\subseteq Q$ with $|Q|=c-j$ and $|Q'|=c-j'$ and a chain of patterns from $Q'$ to $Q$ adding one mark at a time; each child is a quotient of the one before (Proposition~\hb{prop:adjacent}), and by (i) $\mathcal L_j$ is a quotient of $\mathcal L_{j'}$. With the rank difference $b-a$ of $(\mathcal L_b,\mathcal L_a)$,
\[
\min\big(r_{\mathcal L_b},\ r_{\mathcal L_a}+i\big)=\min\big(r_{\mathcal M},\ r_{\mathcal N}+b,\ r_{\mathcal M}+i,\ r_{\mathcal N}+a+i\big)=\min\big(r_{\mathcal M},\ r_{\mathcal N}+a+i\big),
\]
since $i\ge0$ and $a+i\le b$.

(iv) $C$ is a pattern, so the minor is defined, and its bases are the sets $Y\sqcup(Q\setminus C)$ with $C\subseteq Q$, $Q\cap D=\varnothing$ and $Y$ a basis of $\mathcal L_{c-|Q|}$. With $Q''=Q\setminus C\subseteq G''=G\setminus(C\cup D)$ and $c''=c-|C|-|D|=|G''|$, $c-|Q|=|D|+(c''-|Q''|)$, and by (iii) $\mathcal L_{|D|+i}=\mathcal L_i(\mathcal L_{c-|C|},\mathcal L_{|D|})$ for $0\le i\le c''$, the rank difference of that pair being $c''$.
\end{proof}

\begin{hlemma}[\hb{lem:liftcomplement} (Complements of the lifts toward the dual)]
As stated.
\end{hlemma}

\begin{proof}
Write $r(Y)$ and $r^\ast(Y)$ for the rank functions of $\cB$ and $\cB^\ast$, so that $r(\cE\setminus X)=|\cE\setminus X|+r^\ast(X)-r^\ast$ and $r^\ast(\cE\setminus X)=|\cE\setminus X|+r(X)-r$. By Proposition~\hb{prop:paircell}(iii) $\mathcal L_j$ is a table of rank $r+j$, with the rank function $r_{\mathcal L_j}(Y)=\min\big(r^\ast(Y),\ r(Y)+j\big)$ (\hb{def:quotientlift}), and the rank function of its dual table is
\begin{align*}
|X|+r_{\mathcal L_j}(\cE\setminus X)-r-j&=|X|+\min\big(n-|X|+r(X)-r,\ n-|X|+r^\ast(X)-r^\ast+j\big)-r-j\\
&=\min\big(r(X)+c-j,\ r^\ast(X)\big)=r_{\mathcal L_{c-j}}(X),
\end{align*}
using $n-2r=c$ and $n-r^\ast-r=0$. Two tables with one rank function have the same bases, the sets $Y$ with $|Y|=r_{\mathcal L}(Y)$ equal to the rank; so the dual table of $\mathcal L_j$, whose bases are the complements of those of $\mathcal L_j$, has the bases of $\mathcal L_{c-j}$.
\end{proof}

\begin{hdefinition}[\hb{def:fusion} (Fusion width; Foreman fusion; majors)]
As stated.
\end{hdefinition}

\begin{proof}
$c\ge0$ for a quotient is Proposition~\hb{prop:paircell}(i). A quotient $\mathcal N$ of $\mathcal M$ of the same rank has $r_{\mathcal N}=r_{\mathcal M}$: (i) of that proposition gives $r_{\mathcal M}(Y)\ge r_{\mathcal N}(Y)$ at $X=\varnothing$, and $r_{\mathcal N}(X)\ge r_{\mathcal M}(X)$ at $Y=\cE$; so the two have the same bases, the sets $Y$ with $|Y|=r_{\mathcal N}(Y)=r_{\mathcal N}(\cE)$. At $c=0$ a fusable table therefore has $\cB=\cB^\ast$ and $\theta^\ast=\theta$, its bases closed under complement with $\theta_{\cE\setminus T}=\theta_T$; a list has $\theta^\ast=\theta$, so at $c=0$ it is fusable exactly when $\cB=\cB^\ast$. The lifts are then the one table $\mathcal L_0=\mathcal L_c$, with $\theta^{(0)}=\theta=\theta^\ast=\theta^{(c)}$. The condition on the weights pairs the bases of $\mathcal L_j$ with those of $\mathcal L_{c-j}$ by complement, which is a bijection by Lemma~\hb{lem:liftcomplement}; at $j=0$ it is the definition of $\theta^\ast$. $L_0=K_{\cB^\ast}=\Psi$ (Proposition~\hb{prop:identity}) and $L_c=K_{\cB}=K$. Every permutation of $G$ maps the bases of $\cB^\sharp$ onto themselves (Proposition~\hb{prop:paircell}(ii)) and keeps their weights, which depend on $X$ and $|Q|$ alone, so the partial fusion depends on $G'$ only through $|G'|$, up to the names of its marks.
\end{proof}

\begin{htheorem}[\hb{thm:fusion} (Foreman fusion)]
As stated.
\end{htheorem}

\begin{proof}
(i) Proposition~\hb{prop:paircell}(ii) with $(\mathcal M,\mathcal N)=(\cB^\ast,\cB)$, $r_{\mathcal M}=r^\ast$, $r_{\mathcal N}=r$, gives the table, its rank, its rank function and the two minors; $n+c=2n-2r=2r^\ast$. The bases $X\sqcup G$ carry $\theta^{(0)}_X=\theta_X$ and the bases $X\sqcup\varnothing$ carry $\theta^{(c)}_X=\theta^\ast_X$, so the contraction of $G$ is $(\cB,\theta)$ and the deletion is $(\cB^\ast,\theta^\ast)$ (Definition~\hb{def:minors}).

(ii) By Lemma~\hb{lem:liftcomplement}, $X$ is a basis of $\mathcal L_j$ exactly when $\cE\setminus X$ is a basis of $\mathcal L_{c-j}$, and with the weights $\theta^{(c-j)}_{\cE\setminus X}=\theta^{(j)}_X$ the dual table of $(\mathcal L_j,\theta^{(j)})$ is $(\mathcal L_{c-j},\theta^{(c-j)})$. A basis $X\sqcup Q$ of $\cB^\sharp$, $X$ a basis of $\mathcal L_{c-|Q|}$, has the complement $(\cE\setminus X)\sqcup(G\setminus Q)$ with $\cE\setminus X$ a basis of $\mathcal L_{|Q|}=\mathcal L_{c-|G\setminus Q|}$ and the weight $\theta^{(|Q|)}_{\cE\setminus X}=\theta^{(c-|Q|)}_X$: complementation maps the bases onto themselves with their weights. Summing $\theta^\sharp\alpha^{X\sqcup Q}$ over $Q$ by its size gives $K^\sharp=\sum_le_l(\alpha_G)K_{\mathcal L_{c-l}}$; $\Psi^\sharp=K^\sharp$ by Corollary~\hb{cor:mergedvalues}; and $\Psi_{\mathcal L_l}=K_{\mathcal L_l^\ast}=K_{\mathcal L_{c-l}}$ by Proposition~\hb{prop:identity}.

(iii) Let $L$ be a major and $X\sqcup Q$ one of its bases; at $c=0$, $L=(\cB,\theta)=\cB^\sharp$. In $L$ the children of the bunch $G$ at $\varnothing$ and at $G$ are $\cB^\ast$ and $\cB$, so $\varnothing$ and $G$ are patterns, $d(G)=0$ and $G$ is independent; the contraction of the room onto $G$ has the rank $d(G)=0$, so by Proposition~\hb{prop:higgs}(i) every subset of $G$ is a pattern of $G$ in $L$, and along chains of patterns from $\varnothing$ to $Q$ and from $Q$ to $G$ the child at $Q$, of which $X$ is a basis, is a quotient of $\cB^\ast$ and has $\cB$ as a quotient (Propositions~\hb{prop:adjacent} and~\hb{prop:paircell}(i)). So $X$ is independent in $\cB^\ast$ and spans $\cB$ (Proposition~\hb{prop:paircell}(i)); it has $r^\ast-|Q|=r+c-|Q|$ marks, so it is a basis of $\mathcal L_{c-|Q|}$ (Proposition~\hb{prop:paircell}(iii)) and $X\sqcup Q$ is a basis of $\cB^\sharp$.

(iv) The weights of $\cB^\sharp$ depend on $X$ and $|Q|$, so $G$ is in level form in it (Definition~\hb{def:quotientlift}). Let $L$ be a major in which $G$ is in level form. By Proposition~\hb{prop:gluenec}(d) for $L$ and its bunch $G$, the room sides of $r+j$ marks are the bases of the $j$-th lift of the table of the room sides of $r$ marks, those of $L/G=\cB$, toward the table of the room sides of $r+c$ marks, those of $L\setminus G=\cB^\ast$: the bases of $\mathcal L_j$. A room side $X$ of $r+j$ marks makes a basis with some pattern of $c-j$ marks, and in level form the children at all patterns of $c-j$ marks are one table, so $X\sqcup Q$ is a basis of $L$ for every $Q$ of $c-j$ marks: every basis of $\cB^\sharp$ is a basis of $L$, and by (iii) the two have the same bases. The weights of $L$ at $X\sqcup G$ and $X\sqcup\varnothing$ are those of its minors, $\theta$ and $\theta^\ast$, and level form makes the others depend on $X$ and $|Q|$ alone (Proposition~\hb{prop:levelform}); $L$ is identically self-dual exactly when the weight of $X\sqcup Q$ equals that of $(\cE\setminus X)\sqcup(G\setminus Q)$, that is when the middle weights are paired as in Definition~\hb{def:fusion}(ii), and $L$ is then the Foreman fusion whose middle weight $\theta^{(c-|Q|)}_X$ is the weight of $X\sqcup Q$. When $L$ and $\cB$ are lists all these weights are $1$, and $L=\cB^\sharp$. The transpositions generate the permutations of $G$, which gives the equivalent form (Proposition~\hb{prop:levelform}). A bijection $\beta$ carrying $(\cB,\theta)$ carries $(\cB^\ast,\theta^\ast)$, keeps $r$ and $r^\ast$, and so keeps the rank function of every $\mathcal L_j$; with $\theta^{(j)}\circ\beta=\theta^{(j)}$, $\beta\sqcup\pi$ maps each basis $X\sqcup Q$ to the basis $\beta(X)\sqcup\pi(Q)$ of the same weight.

(v) At $c=1$, $\mathcal L_0=\cB$ and $\mathcal L_1=\cB^\ast$: the bases are $X+g$, $X\in\cB$, with $\theta_X$, and $X\in\cB^\ast$ with $\theta^\ast_X$, that is $\cE\setminus T$ with $\theta_T$; $K^\sharp=e_0L_0+e_1L_1=\Psi+\alpha_gK$. A major $L$ has $L/g=(\cB,\theta)$ and $L\setminus g=(\cB^\ast,\theta^\ast)$, which fixes its bases through $g$ and avoiding $g$ with their weights: $L=\cB^\sharp$.
\end{proof}

\begin{htheorem}[\hb{thm:fusiontell} (The tell)]
As stated.
\end{htheorem}

\begin{proof}
Three facts about a table $\cB$ of rank $r$. ($\alpha$) A set $A$ is closed in $\cB^\ast$ exactly when $\cE\setminus A$ is a union of circuits of $\cB$: for $e\notin A$, $r^\ast(A+e)-r^\ast(A)=1+r(\cE\setminus A-e)-r(\cE\setminus A)$, which is $1$ exactly when $e$ lies in the closure of $\cE\setminus A-e$, that is in a circuit inside $\cE\setminus A$. ($\beta$) The cocircuits are the complements of the closed sets of rank $r-1$: $D$ meets every basis exactly when $\cE\setminus D$ contains none, that is $r(\cE\setminus D)\le r-1$, and $D$ is minimal with this property exactly when $\cE\setminus D$ is maximal with it, a closed set of rank $r-1$. ($\gamma$) The complement of a closed set $A$ is a union of cocircuits: for $e\notin A$, a basis $I$ of $A$ has $I+e$ independent, $I+e$ extends to a basis $T$, and the closure $H$ of $T-e$ has rank $r-1$, contains $A$, which lies in the closure of $I$, and avoids $e$, $T$ being independent; so $\cE\setminus H$ is a cocircuit through $e$ inside $\cE\setminus A$, by ($\beta$).

(a)$\Rightarrow$(b): the complement of a cocircuit is closed in $\cB$ by ($\beta$), so in $\cB^\ast$, and the cocircuit is a union of circuits by ($\alpha$). (b)$\Rightarrow$(a): the complement of a closed set $A$ of $\cB$ is a union of cocircuits by ($\gamma$), each a union of circuits by (b), so $A$ is closed in $\cB^\ast$ by ($\alpha$). (b)$\Rightarrow$(c): if $e$ is not in the closure of $A$, ($\gamma$) applied to that closure gives a cocircuit $D$ with $e\in D$ and $D\cap A=\varnothing$, so $D-e\subseteq A'$; by (b) a circuit $C$ has $e\in C\subseteq D$, and $e$ lies in the closure of $C-e\subseteq A'$. (c)$\Rightarrow$(b): for a cocircuit $D$ and $e\in D$ split $\cE\setminus\{e\}$ into $A=\cE\setminus D$ and $A'=D-e$; $A$ is closed by ($\beta$) and avoids $e$, so by (c) $e$ lies in the closure of $D-e$, in a circuit inside $D$. By Proposition~\hb{prop:paircell}(i) (a) implies $c\ge0$.

By Definition~\hb{def:fusion} the table is fusable exactly when (a) holds and, at $c=0$, $\theta^\ast=\theta$, which a list always has. At $c=0$, $G$ is empty, a major is a table equal to $(\cB,\theta)$ and to $(\cB^\ast,\theta^\ast)$, and one exists exactly when $\cB=\cB^\ast$ and $\theta^\ast=\theta$, which is fusability, a quotient of the same rank having the same bases (proof of Definition~\hb{def:fusion}). Let $c\ge1$. A fusable table has a major, its Foreman fusion (Theorem~\hb{thm:fusion}(i), whose proof uses Proposition~\hb{prop:paircell} alone). Conversely, in a major the children of the bunch $G$ at $\varnothing$ and at $G$ are $\cB^\ast$ and $\cB$, so $\varnothing$ and $G$ are patterns, $d(G)=0$ and $G$ is independent; the contraction of the room onto $G$ has rank $d(G)=0$, and by Proposition~\hb{prop:higgs}(i) every subset of $G$ is a pattern. Along a chain from $\varnothing$ to $G$ adding one mark at a time each child is a quotient of the one before (Proposition~\hb{prop:adjacent}), and $\cB$ is a quotient of $\cB^\ast$ by Proposition~\hb{prop:paircell}(i): the table is fusable. A table on $\cE\sqcup G$ whose children at $\varnothing$ and at $G$ have the bases of $\cB^\ast$ and of $\cB$ is a major once its bases $X\sqcup G$ carry $\theta_X$ and its bases $X\sqcup\varnothing$ carry $\theta^\ast_X$, the two sets of bases being disjoint as $G\neq\varnothing$; this is the equivalent form of (d).

A coloop $e$ makes $\{e\}$ a cocircuit, and a circuit inside $\{e\}$ would make $e$ a loop; so (b) fails. The circuits of $\cB^\ast$ are the cocircuits of $\cB$ and its cocircuits the circuits of $\cB$, which gives (b) for the dual.
\end{proof}

\begin{hcorollary}[\hb{cor:fusionisd} (The fused tables are the identically self-dual tables)]
As stated.
\end{hcorollary}

\begin{proof}
(i) By Corollary~\hb{cor:mergedvalues} $\mathcal H$ has rank $|\cE'|/2$ and no loop or coloop; $C$ is independent in $\mathcal H$ and so in $\mathcal H^\ast=\mathcal H$. $\mathcal H/C$ has rank $|\cE'|/2-|C|$ on $|\cE'|-|C|$ marks, hence the fusion width $|C|$. Its dual has the bases $(\cE'\setminus C)\setminus(T\setminus C)=\cE'\setminus T$, $T\supseteq C$, with the weights $\theta_T=\theta_{\cE'\setminus T}$: the bases of $\mathcal H$ avoiding $C$ with their weights, so $(\mathcal H/C)^\ast=\mathcal H\setminus C$ and $\mathcal H$ is a major of $\mathcal H/C$. By Theorem~\hb{thm:fusiontell} $\mathcal H/C$ is fusable; $C$ is in level form in $\mathcal H$ (Proposition~\hb{prop:levelform}), and $\mathcal H$ is identically self-dual, so by Theorem~\hb{thm:fusion}(iv) it is the Foreman fusion of $\mathcal H/C$ with the middle weights $\theta^{(j)}_X=\theta_{X\sqcup Q}$, $|Q|=c-j$. A one-mark set is independent, $\mathcal H$ having no loop, and it has no transposition. With Theorem~\hb{thm:fusion}(ii) the fused tables are exactly the identically self-dual ones.

(ii) $c(S)=|S|-d(S)-d^\ast(S)$ does not change under duality, which exchanges $d$ and $d^\ast$, so $\cB^\ast$ is connected and its circuits link all of $\cE$ (Proposition~\hb{prop:ladder}(i)). The circuits of $\cB^\ast=\cB^\sharp\setminus G$ are circuits of $\cB^\sharp$, so $\cE$ lies in one molecule of $\cB^\sharp$. For $g\in G$ take a basis $X\sqcup Q$ with $g\notin Q$; the circuit through $g$ inside $X\sqcup Q+g$ does not lie in the independent set $G$, so it meets $\cE$, and $g$ lies in the molecule of $\cE$. So $\cB^\sharp$ is one molecule, connected by Proposition~\hb{prop:ladder}(i).
\end{proof}

\begin{hproposition}[\hb{prop:fusionminors} (Partial fusions and inner lifts)]
As stated.
\end{hproposition}

\begin{proof}
(i) is Proposition~\hb{prop:paircell}(iv) for $(\cB^\ast,\cB)$, the minor carrying the weights of $\cB^\sharp$ (Definition~\hb{def:minors}). (ii) With $C=G\setminus G'$ and $D=\varnothing$, (i) gives the level cell of $(\mathcal L_i,\cB)$, of rank $r+i$ on $n+i$ marks, so of the fusion width $c-i$. Being identically self-dual, $\cB^\sharp$ has $(\cB^\sharp/(G\setminus G'))^\ast=\cB^\sharp\setminus(G\setminus G')$ with the weights, as in the proof of Corollary~\hb{cor:fusionisd}(i); $G\setminus G'$ is independent and in level form in $\cB^\sharp$ (Proposition~\hb{prop:paircell}(ii), Definition~\hb{def:quotientlift}); so Corollary~\hb{cor:fusionisd}(i) makes $\cB^\sharp_i$ fusable with the Foreman fusion $\cB^\sharp$, its middle weights read off $\theta^\sharp$. (iii) $\mathcal L_j$ has the rank $r+j$ and the corank $n-r-j=r+c-j$, so the fusion width $c-2j$. For $2j\le c$, $\mathcal L_j$ is a quotient of $\mathcal L_{c-j}=\mathcal L_j^\ast$ (Proposition~\hb{prop:paircell}(iii) and Theorem~\hb{thm:fusion}(ii)), and at $2j=c$, where $\mathcal L_j^\ast=\mathcal L_j$, its weights are those of its dual, $\theta^{(j)}_{\cE\setminus X}=\theta^{(j)}_X$; so it is fusable, and its lifts toward its dual are the $\mathcal L_{j+i}$, $0\le i\le c-2j$ (Proposition~\hb{prop:paircell}(iii)). Its level cell is $\cB^\sharp/C\setminus D$ with $|C|=|D|=j$ by (i), whose weights are paired by complement as those of $\cB^\sharp$ are, so it is the Foreman fusion of $(\mathcal L_j,\theta^{(j)})$ with the middle weights read off $\theta^\sharp$. At $j=c/2$ the lift is identically self-dual; at $j=(c-1)/2$ the fusion width is one.
\end{proof}

\begin{hproposition}[\hb{prop:fusionseam} (The seam of the fusion bunch)]
As stated.
\end{hproposition}

\begin{proof}
(i) $\pat(G)=2^G$, $d(G)=0$, $c(G)=c$ and level form are Proposition~\hb{prop:paircell}(ii) and Definition~\hb{def:quotientlift}. A pattern $Q$ of $l$ marks has the child $(\mathcal L_{c-l},\theta^{(c-l)})$, the same for all of them, so they form one station with the multipliers $1$ (Definition~\hb{def:seam}); its station letters are $\sum_{|Q|=l}\alpha_G^Q=e_l(\alpha_G)$ and $\sum_{|Q|=l}\alpha_G^{G\setminus Q}=e_{c-l}(\alpha_G)$. A room side $X$ of excess $l$ is a basis of $\mathcal L_{c-l}$, its size fixing $l$, and its room child has the bases of $\binom Gl$, each with the weight $\theta^{(c-l)}_X$: all the room children of the excess $l$ are positive multiples of $\binom Gl$, one station $D_l$, with the multipliers $w'_X=\theta^{(c-l)}_X/\mu_l$ and the room child of the station $\mu_l\binom Gl$, of the letters $\mu_le_l$ and $\mu_le_{c-l}$. The room-station letters are $\sum_Xw'_X\alpha^X=L_l/\mu_l$ and $\sum_Xw'_X\alpha^{\cE\setminus X}=\Psi_{\mathcal L_{c-l}}/\mu_l=L_{c-l}/\mu_l$ (Theorem~\hb{thm:fusion}(ii)), and the pair constant is $\theta^{(c-l)}_X/(1\cdot w'_X)=\mu_l$. One station on each side at each excess is one rectangle, and the room-station factorization of Proposition~\hb{prop:seamidentity}(iii) is the stated one. The weight of $X\sqcup Q$ depends on $X$ alone, $|X|$ fixing $|Q|$, so the weights factor across $G$ and $\cE$ and Theorem~\hb{thm:envelope} applies: with one rectangle at every excess the table is its level envelope.

(ii) By Definition~\hb{def:interaction}, with $K_{(Q)}=K_{\mathcal L_{c-|Q|}}$ and $\Psi_{(Q)}=\Psi_{\mathcal L_{c-|Q|}}=K_{\mathcal L_{|Q|}}$,
\[
\iota_G=\sum_{l=0}^c\binom cl(-1)^{c-l}\log K_{\mathcal L_{c-l}},\qquad\iota^\ast_G=\sum_{l=0}^c\binom cl(-1)^l\log K_{\mathcal L_l},
\]
both equal to $\sum_j(-1)^j\binom cj\log K_{\mathcal L_j}$; at $c=2$ this is $\log K_{\mathcal L_0}-2\log K_{\mathcal L_1}+\log K_{\mathcal L_2}=-\log z_1$.

(iii) At the table's gauge the pattern probability of $Q$ is $K_{(Q)}(\one)/K^\sharp(\one)=L_{|Q|}(\one)/K^\sharp(\one)$ (Definition~\hb{def:split}); the $\binom cl$ patterns of $l$ marks give $\Pi_l$. $L_l(\one)$ is the total weight of $(\mathcal L_{c-l},\theta^{(c-l)})$, which equals that of its dual $(\mathcal L_l,\theta^{(l)})$, $L_{c-l}(\one)$; so $\Pi_l=\Pi_{c-l}$. The peel (Definition~\hb{def:peel}) gives the child at $Q$ the exponents $\pi_Q(\one)(s,t)$, and merging the equal children of each excess gives $(\mathcal L_{c-l};\Pi_ls,\Pi_lt)$, that is $(\mathcal L_j;\Pi_js,\Pi_jt)$. Its word is $\prod_jK_{\mathcal L_j}^{-\Pi_js}\Psi_{\mathcal L_j}^{-\Pi_jt}=\prod_jK_{\mathcal L_j}^{-\Pi_js}\prod_jK_{\mathcal L_{c-j}}^{-\Pi_jt}=\prod_jK_{\mathcal L_j}^{-\Pi_j(s+t)}$, reindexing the second product by $j\mapsto c-j$.
\end{proof}

\begin{hdefinition}[\hb{def:universalpiece} (The universal piece; the lift partner)]
As stated.
\end{hdefinition}

\begin{proof}
The seam letters $e_l$ and $e_{c-l}$ of the station $D_l$ are nonzero, with nonnegative coefficients, homogeneous of the degrees $l=d+l$ and $c-l=d^\ast+c-l$ with $d=d^\ast=0$, and positive letters where the degree is positive; the room letters $L_l=K_{\mathcal L_{c-l}}$ and $L_{c-l}$ have the degrees $r+c-l=e+c-l$ and $r+l=e^\ast+l$ with $e=e^\ast=r$, as Definition~\hb{def:seampiece} asks of a partner read with the reversed excesses.
\end{proof}

\begin{htheorem}[\hb{thm:fusionglue} (The glue of a fused table)]
As stated.
\end{htheorem}

\begin{proof}
The family $((\cB^\sharp,\theta^\sharp);2\sigma,0)$ has the word $K^{\sharp\,-2\sigma}=K^{\sharp\,-s}\Psi^{\sharp\,-t}$ (Theorem~\hb{thm:fusion}(ii)), and it is the glue of the universal piece and the lift partner with the one active letter $K$ (Definition~\hb{def:seampiece}, Proposition~\hb{prop:fusionseam}(i)).

(i) In Definition~\hb{def:seamlaw} with the one active letter at the exponent $2\sigma$: $\kappa_G=m_G$, $\kappa_G+\kappa^{\mathrm{top}}_G=2\sigma c$, $V^{(0)}=e_0=1$ and $V^{(c)}=e_c$, so $\gamma^\rho_G=e_c^{1/c}$, the scale form is $\phi^\rho_G(du)=u^{\mvec|_G-\one}e_c(u)^{-m_G/c}\,d\omega_G(u)$, and the seam vector at $D_l$ is $e_l/(\gamma^\rho_G)^l$. The form $\nu_G=\alpha^{\mvec|_G-\one}e_c(\alpha)^{-m_G/c}\,d\alpha$ on $(0,\infty)^G$ and the seam vector are homogeneous of degree zero. For a Borel $\psi\ge0$ on the seam torus and $\one_{[1,e]}$ the indicator of $[1,e]$, with $\int_0^\infty\one_{[1,e]}(\rho)\,d\rho/\rho=1$, the integral of $\psi(\boldsymbol\varrho^\rho_G(\alpha))\,\one_{[1,e]}(h(\alpha))\,\nu_G$ over the orthant is, for every positive $h$ of degree one, the integral of $\psi(\boldsymbol\varrho^\rho_G)$ against the law: in the coordinates $\alpha=\rho u$ of scale and shape (Lemma~\hb{lem:scaleshape}) the form is $\one_{[1,e]}(\rho)\,d\rho/\rho$ times $\phi^\rho_G$, the integrand not depending on $\rho$, and a change of the section $h$ is a change of $\rho$ along each ray, which leaves $\int\one_{[1,e]}(\rho)\,d\rho/\rho$ unchanged (the computation of Lemma~\hb{lem:chart}, in $[0,\infty]$ by Tonelli). Take $h=e_c^{1/c}$ and write $\alpha=\rho\,e^w$ with $\rho>0$, $w\in H_G$. In the coordinates $y=\log\alpha=\lambda\one+w$, $\lambda=\log\rho$, the map $(\lambda,w_{g_1},\dots,w_{g_{c-1}})\mapsto y$, $w_{g_c}=-\sum_{i<c}w_{g_i}$, is linear with determinant $\pm c$: adding all rows to the last turns it into $(c,0,\dots,0)$. So $d\alpha=\alpha^{\one}c\,d\lambda\,d\lambda_G(w)$, and since $e_c(\rho e^w)=\rho^c$,
\[
\nu_G=e^{m_G\lambda+\ip{\mvec|_G}w}\,e^{-m_G\lambda}\,c\,d\lambda\,d\lambda_G(w)=c\,e^{\ip{\mvec|_G}w}\,\frac{d\rho}\rho\,d\lambda_G(w),
\]
while $\boldsymbol\varrho^\rho_G(\rho e^w)=\mathbf e(e^w)$. So the integral of $\psi$ against the law is $c\int_{H_G}\psi(\mathbf e(e^w))\,e^{\ip{\mvec|_G}w}\,d\lambda_G(w)$ for every $\psi$, which is the statement.

(ii) By Lemma~\hb{lem:chart}, in $[0,\infty]$ by Tonelli, $P_{\cB^\sharp}(\mvec;s,t)$ is the integral over the orthant of $\alpha^{\mvec-\one}K^{\sharp\,-2\sigma}\one_{[1,e]}(h(\alpha))\,d\alpha$ with $h=e_c(\alpha_G)^{1/c}$. In the coordinates $(\alpha_{\cE},\rho,w)$, $\alpha_G=\rho e^w$, as in (i), $\alpha_G^{\mvec|_G-\one}d\alpha_G=c\,e^{\ip{\mvec|_G}w}\rho^{m_G-1}d\rho\,d\lambda_G(w)$ and $K^\sharp=\sum_l\rho^le_l(e^w)L_l(\alpha_{\cE})$ (Theorem~\hb{thm:fusion}(ii)). By Tonelli the integral is $c\int_{H_G}e^{\ip{\mvec|_G}w}I(w)\,d\lambda_G(w)$ with
\[
I(w)=\int_{(0,\infty)^{\cE}\times(0,\infty)}\alpha_{\cE}^{\mvec|_{\cE}-\one}\rho^{m_G-1}\Big(\sum_le_l(e^w)\rho^lL_l(\alpha_{\cE})\Big)^{-2\sigma}\one_{[1,e]}(\rho)\,d\alpha_{\cE}\,d\rho .
\]
The letter $\sum_lx_l\alpha_g^lL_l$ is homogeneous of degree $r+c=r^\ast$ on $\cE\sqcup\{g\}$, so the form of $W[\mathbf x]$ at the masses $(\mvec|_{\cE},m_G)$ has degree $m_{\cE}+m_G-2\sigma r^\ast=0$ on the balance wall, and $I(w)$ is its integral against $\one_{[1,e]}$ of the section function $\alpha_g$: by Lemma~\hb{lem:chart} it is $P_{W[\mathbf e(e^w)]}(\mvec|_{\cE},m_G)$.

(iii) $\sum_lx_lL_ly^l$ is a polynomial of degree $c$ in $y$ with positive coefficients, $x_cL_c\neq0$, so its roots $-\xi_i$ avoid $[0,\infty)$; read as in Definition~\hb{def:skelpoly} on the bunch $\{g\}$, whose floor is $0$, its Skeleton polynomial at $\alpha_g=1$ is $\prod_i(x+\xi_i)/(1+\xi_i)$. Non-real roots come in conjugate pairs, and for $\rho>0$ the product of the principal powers $(\rho+\xi_i)^{-2\sigma}$ is the positive power of the product. Proposition~\hb{prop:scalebreta} with $J=\{1,\dots,c\}$, the letter exponents $w_i=2\sigma$, $|\mathbf w|=2\sigma c$ and $\kappa=m_G$ gives the formula for $0<m_G<2\sigma c$, coincident values merging.

(iv) This is Theorem~\hb{thm:gluewidth}(iii) for the universal piece and the lift partner with the one active letter: one share $a_l=a_{D_l}$ per station, the letter total $\sum_la_l=2\sigma$ and the scale pin $\sum_li(D_l)a_l=\kappa_G=m_G$, the share kernel $\Gst{c+1}1{\mathbf a}{2\sigma}$, and the two seam functions, which on the pins are the periods of the seam word $\prod_le_l^{-a_l}$ at $\mvec|_G$ and of the room word $\prod_lL_l^{-a_l}$ at $\mvec|_{\cE}$ (Lemma~\hb{lem:seamword}), $e_0=1$ dropping out. Two pins on $c+1$ shares leave a fibre of dimension $c-1$.

Let the masses be real and every proper wall of $\cB^\sharp$ positive; we show that both periods converge at the native split at the stations. The walls of $((\cB^\sharp,\theta^\sharp);2\sigma,0)$ are $\kappa_U=m_U-2\sigma d^\sharp(U)$, with $d^\sharp(U)=r^\ast-r^\sharp(\cE^\sharp\setminus U)$ on $\cE^\sharp=\cE\sqcup G$; at one mark, no mark being a coloop (Corollary~\hb{cor:mergedvalues}), $d^\sharp$ is $0$, so the masses are positive. A set $U$ of the width zero in $\cB^\sharp$ has $r^\sharp(U)+r^\sharp(\cE^\sharp\setminus U)=r^\ast$, so $d^\sharp(U)=r^\sharp(U)$ and $d^\sharp(\cE^\sharp\setminus U)=r^\sharp(\cE^\sharp\setminus U)$, and $\kappa_U+\kappa_{\cE^\sharp\setminus U}=m_{\cE^\sharp}-2\sigma r^\ast=0$ on the balance wall; so no nonempty proper $U$ has the width zero, $\cB^\sharp$ is connected (Proposition~\hb{prop:ladder}(i)), and it has a native gauge (Theorem~\hb{thm:saddle}(iii)). By Lemma~\hb{lem:nativestations}, $c(G)=c\ge1$, the native split at the stations is a point of the pins with positive shares $a_l$ at which $\mvec|_G$ and $\mvec|_{\cE}$ lie in the relative interiors of the Newton sums of the seam word and of the room word, and the two words converge there once these Newton sums have the full dimensions $c-1$ and $n-1$; the seam letters $\mu_le_l$ and the room-station letters $L_l/\mu_l$ of the bunch $G$ (Proposition~\hb{prop:fusionseam}(i)) differ from $e_l$ and $L_l$ by constant factors, which change neither the Newton sums nor the convergence. The Newton sum $\sum_{l\ge1}a_l\mathcal N_{e_l}$ of the seam word contains a translate of $a_1\operatorname{conv}\{\one_{\{g\}}:\ g\in G\}$, of the dimension $c-1$. The Newton sum $\sum_la_l\mathcal N_{L_l}$ of the room word, $\mathcal N_{L_l}$ the base polytope of $\mathcal L_{c-l}$, has as its directions the sums of the directions of its terms, so a vector $w\in\mathbb R^{\cE}$ orthogonal to them has $\ip w{\one_X}$ constant over the bases $X$ of every lift, and $w$ is constant on every molecule of the family of the lifts (Lemma~\hb{lem:circuit}). That family is connected. Otherwise a proper nonempty $S\subset\cE$ is a union of its molecules, hence of molecules of every lift (Proposition~\hb{prop:directsum}), and has the width zero in $\mathcal L_0=\cB$ and in $\mathcal L_1$, whose rank function is $r(Y)+\min(-\varepsilon(Y),1)$, $\varepsilon(Y)=r(Y)-r^\ast(Y)\le0$ the slant of $Y$ in $\cB$ (Definition~\hb{def:skeleton}, Proposition~\hb{prop:paircell}(i)). The two widths give $r(S)+r(\cE\setminus S)=r$ and $\min(-\varepsilon(S),1)+\min(-\varepsilon(\cE\setminus S),1)=1$, so $\varepsilon(S)=0$ or $\varepsilon(\cE\setminus S)=0$, say $\varepsilon(S)=0$, exchanging $S$ and $\cE\setminus S$ if needed. Then by Theorem~\hb{thm:fusion}(i) $r^\sharp(S)=\min(r^\ast(S),r(S)+c)=r(S)$ and $r^\sharp((\cE\setminus S)\sqcup G)=\min(r^\ast(\cE\setminus S)+c,\ r(\cE\setminus S)+c)=r(\cE\setminus S)+c$, whose sum is $r+c=r^\ast$: $S$ has the width zero in $\cB^\sharp$, which is excluded. So $w$ is constant, the Newton sum of the room word has the dimension $n-1$, and both words converge at the native split (Lemma~\hb{lem:nativestations}, Theorem~\hb{thm:poswordconv}).

(v) $e_c(\alpha_G)=\alpha_G^{\one}$, so $\alpha_G^{\mvec|_G-\one}\prod_{l=1}^ce_l^{-a_l}=\alpha_G^{(\mvec|_G-a_c\one)-\one}\prod_{l=1}^{c-1}e_l^{-a_l}$: the two forms are one, and so are their periods; $e_l$ is the letter of $\binom Gl$. At $c=1$ the pins give $a_1=m_g$, and the form $\alpha_g^{m_g-1}\alpha_g^{-a_1}d\alpha_g$ on one mark has the period $1$. At $c=2$, $u_1+u_2=1$ on $\Delta_G$, so the period of $e_1^{-a_1}$ at $(m_1-a_2,m_2-a_2)$ is $\int_{\Delta_G}u^{(\mvec|_G-a_2\one)-\one}d\omega_G=\Gst21{m_1-a_2,\ m_2-a_2}{a_1}$ (Theorem~\hb{thm:dirichlet}), $m_1+m_2-2a_2=a_1$ on the pins.
\end{proof}

\begin{hcorollary}[\hb{cor:fusionone} (Width one)]
As stated.
\end{hcorollary}

\begin{proof}
$\cB^\sharp$ is identically self-dual (Theorem~\hb{thm:fusion}(ii)), and Corollary~\hb{cor:mergedvalues} at the mark $g$, which allows the exponents $(2\sigma,0)$, gives $P_{\cB^\sharp}=\Gst21{m_g,\ 2\sigma-m_g}{2\sigma}P_{\cB^\sharp\setminus g}(\mvec|_{\cE};2\sigma-m_g,m_g)$. $\cB^\sharp\setminus g=(\cB^\ast,\theta^\ast)$, whose letters are $\Psi$ and $K$ (Proposition~\hb{prop:identity}), so its period at $(2\sigma-m_g,m_g)$ is $P_{\cB}(\mvec|_{\cE};m_g,2\sigma-m_g)$. At $c=1$, $H_G=\{0\}$ and the seam vector of the universal piece is $(1,1)$; that of the lift partner, $\big(L_0/(L_1\gamma_{\cE}),\,L_1/L_1\big)$ with $\gamma_{\cE}=L_0/L_1$ (Definition~\hb{def:seamlaw}, read with the reversed excesses), is $(1,1)$; and the two pins fix $a_0=2\sigma-m_g$ and $a_1=m_g$.
\end{proof}

\begin{hcorollary}[\hb{cor:fusiontwo} (Width two)]
As stated.
\end{hcorollary}

\begin{proof}
(i) Theorem~\hb{thm:fusionglue}(iv) and (v) at $c=2$: the pins give $a_1=m_1+m_2-2a_2$ and $a_0=2\sigma-m_1-m_2+a_2$, solved for $(a_0,a_1)$ by a matrix of determinant one, and the kernel times the universal factor is $\frac{\Gamma(a_0)\Gamma(a_1)\Gamma(a_2)}{\Gamma(2\sigma)}\cdot\frac{\Gamma(m_1-a_2)\Gamma(m_2-a_2)}{\Gamma(a_1)}$. (ii) Theorem~\hb{thm:fusionglue}(i) and (ii) at $c=2$: $H_G=\{(w,-w)\}$ with $d\lambda_G=dw$, $v=e^w$, $\mathbf e(e^w)=(1,v+v^{-1},1)$ and $2e^{(m_1-m_2)w}dw=2v^{m_1-m_2-1}dv$. (iii) $K^\sharp=A+\alpha_{g_2}B$ with $A=L_0+\alpha_{g_1}L_1$ and $B=L_1+\alpha_{g_1}L_2$, and $\int_0^\infty\rho^{\nu-1}(A+B\rho)^{-2\sigma}d\rho=\Gst21{\nu,\ 2\sigma-\nu}{2\sigma}A^{\nu-2\sigma}B^{-\nu}$ for $0<\nu<2\sigma$ (Proposition~\hb{prop:scalebreta}) at $\nu=m_2$. With $\alpha_{g_1}=(L_0/L_1)x$, $A=L_0(1+x)$ and $B=L_1(1+z_1^{-1}x)$, which gives the display by Tonelli; the $x$-integral is $z_1^{m_2}\int_0^\infty x^{m_1-1}(x+1)^{-(2\sigma-m_2)}(x+z_1)^{-m_2}dx$, the scale integral of Proposition~\hb{prop:scalebreta} with the coefficients $1,z_1$, the letter exponents $2\sigma-m_2,m_2$ and $\kappa=m_1$, finite exactly for $0<m_1<2\sigma$. At $m_1=m_2=1$, $2\sigma=2$, $\int_0^\infty\frac{dx}{(1+x)(1+x/z_1)}=\frac{z_1\log z_1}{z_1-1}$ for $z_1\neq1$ and $1$ for $z_1=1$, and $L_1^{-2}z_1\log z_1/(z_1-1)=\log z_1/(L_1^2-L_0L_2)$. At integer data and $z_1\neq1$ the integrand is a rational function of $x$ with poles at $-1$ and $-z_1$ only and of order at most $x^{-2}$ at infinity, $m_1<2\sigma$; its partial fractions give an antiderivative $R(x)+A'\log(1+x)+B'\log(1+x/z_1)$ with $R$ rational and $A'+B'=0$, all coefficients rational in $z_1$, whose increment over $(0,\infty)$ is $R(\infty)-R(0)+A'\log z_1$ (receipt \texttt{t10b\_weight\_exact.py} for $m_1=m_2=k$, $2\sigma=2k$, $k\le4$).
\end{proof}

\begin{hlemma}[\hb{lem:liftproducts} (Products of lifts)]
As stated.
\end{hlemma}

\begin{proof}
Let $A$ be a basis of $\mathcal L_i$. Since $\mathcal L_i$ is a quotient of $\mathcal L_j$ (Proposition~\hb{prop:paircell}(iii)), $A$ is independent in $\mathcal L_j$ (Proposition~\hb{prop:paircell}(i)) and lies in a basis $C$ of $\mathcal L_j$. The monomial $\alpha^{\one_A+\one_C}=\alpha^{2\one_A+\one_{C\setminus A}}$ occurs in $K_{\mathcal L_i}K_{\mathcal L_j}$ with a positive coefficient, every coefficient being nonnegative, and exactly $|A|=r_{\mathcal N}+i$ of its exponents equal $2$. A monomial $\alpha^{\one_{A'}+\one_{C'}}$ of $K_{\mathcal L_{i'}}K_{\mathcal L_{j'}}$ has the exponent $2$ exactly on $A'\cap C'$, at most $|A'|=r_{\mathcal N}+i'<r_{\mathcal N}+i$ marks. Both products are nonzero with nonnegative coefficients, so a positive constant multiple of one has the support of the one, and they are not proportional. The ratio is the case $(i',i,j,j')=(j-1,j,j,j+1)$, $0<j<c$; a constant ratio would make $K_{\mathcal L_j}^2-\lambda K_{\mathcal L_{j-1}}K_{\mathcal L_{j+1}}$ vanish on the orthant, hence identically.
\end{proof}

\begin{htheorem}[\hb{thm:fusionprice} (The price of width)]
As stated.
\end{htheorem}

\begin{proof}
(i) By Theorem~\hb{thm:fusionglue}(i) the universal law is the image of a measure with a positive continuous density on $H_G$, of dimension $c-1$, under the continuous map $w\mapsto\mathbf e(e^w)$. At $c=1$ it is the unit point mass at $(1,1)$. At $c\ge2$, $e_1(e^w)=\sum_ge^{w_g}$ is $c$ at $w=0$ and larger at every other point of $H_G$ by the inequality of the arithmetic and geometric means, $\prod_ge^{w_g}=1$; a law carried by one point would make it constant almost everywhere, hence everywhere by continuity. The law of the lift partner is the image of its scale form, which has a positive continuous density on the open simplex $\Delta_{\cE}$, under its continuous seam vector; it is carried by one point exactly when the seam vector is constant. Read with the reversed excesses, $V^{(0)}=L_c$, $V^{(c)}=L_0$, the scale is $\gamma_{\cE}=(L_0/L_c)^{1/c}$ and the seam vector at $D_l$ is $L_l/(L_c\gamma_{\cE}^{\,c-l})=L_lL_0^{-(c-l)/c}L_c^{-l/c}$ (Definition~\hb{def:seamlaw}). Put $f_l=\log L_l-\log L_0-\frac lc(\log L_c-\log L_0)$, so that $f_0=f_c=0$ and the second differences of $(f_l)_l$ are those of $(\log L_l)_l$, which are $-\log z_{c-l}$, since $L_l=K_{\mathcal L_{c-l}}$. The seam vector is constant exactly when every $f_l$ is constant; then every lift ratio is constant, and conversely constant second differences with $f_0=f_c=0$ determine $f_l$ as constants. At $c=1$ there is no lift ratio and the seam vector is $(1,1)$; at $c\ge2$ Lemma~\hb{lem:liftproducts} excludes constant lift ratios.

(ii) For $a\in\mathbb C^c$ and a set $S$ of $l$ indices, $\pi(S)$ runs over the sets of $l$ indices as $\pi$ runs over $S_c$, each reached $l!\,(c-l)!$ times, so
\[
\sum_{\pi\in S_c}\prod_{i=1}^c\big(1+a_i\eta_{\pi(i)}\big)=\sum_{S}a^S\sum_\pi\prod_{i\in S}\eta_{\pi(i)}=\sum_{l=0}^cl!\,(c-l)!\,e_l(a)\,e_l(\boldsymbol\eta),
\]
and the left side is $\operatorname{per}(\one\one^{\mathsf T}+a\boldsymbol\eta^{\mathsf T})$ by the definition of the permanent. With $L_l=L_0e_l(\boldsymbol\eta)/\binom cl$, read off $\sum_l\binom clL_lx^l=L_0\prod_i(1+\eta_ix)$, division by $c!$ gives $\sum_le_l(\alpha_G)L_l=K^\sharp$. The Skeleton polynomial of $G$ at $(\one_G,\alpha_{\cE})$ is $K^\sharp(x\one_G,\alpha_{\cE})/K^\sharp(\one_G,\alpha_{\cE})=\sum_l\binom clL_lx^l/K^\sharp$, $d(G)=0$, and with $L_0>0$ its roots $-\xi_i$ give $\sum_l\binom clL_lx^l=L_0\prod_i(1+x/\xi_i)$: $\eta_i=1/\xi_i$. At a gauge all $\eta_i$ equal $\eta$ exactly when $\sum_l\binom clL_lx^l=L_0(1+\eta x)^c$, that is $L_l=L_0\eta^l$ for every $l$, which holds exactly when $L_l^2=L_{l-1}L_{l+1}$ for $0<l<c$ (then $L_l=L_0(L_1/L_0)^l$ by induction), that is every lift ratio is $1$. Coincidence on an open set would make $K_{\mathcal L_j}^2-K_{\mathcal L_{j-1}}K_{\mathcal L_{j+1}}$ vanish there and so identically, against Lemma~\hb{lem:liftproducts}; and $K^\sharp=L_0\prod_g(1+\alpha_g\eta)=L_0\sum_le_l(\alpha_G)\eta^l$ would give $L_l=L_0\eta^l$ everywhere.

(iii) By Proposition~\hb{prop:fusionminors}(ii) the partial fusion $\mathcal A$ at $G'$ is the level cell of $(\mathcal L_i,\cB)$ on $\cE\sqcup G'$, in which $g$ lies in some bases and avoids others, so it is a cell with the hinge $g$; by Proposition~\hb{prop:paircell}(iv) its hinge minors are $\mathcal A/g$, the level cell of $(\mathcal L_{i-1},\cB)$, and $\mathcal A\setminus g$, the level cell of $(\mathcal L_i,\mathcal L_1)$, both on $\cE\sqcup(G'-g)$. At $i=c$, $\mathcal A=\cB^\sharp$ is identically self-dual, $\mathcal A/g$ is the dual of $\mathcal A\setminus g$ with the same weights (as in the proof of Corollary~\hb{cor:fusionisd}(i)), so $K_{\mathcal A/g}=\Psi_{\mathcal A\setminus g}$ and $\Psi_{\mathcal A/g}=K_{\mathcal A\setminus g}$ (Proposition~\hb{prop:identity}) and the ratio is $1$; and Corollary~\hb{cor:mergedvalues} at $g$, with $P_{\mathcal A\setminus g}$ at $(2\sigma-m_g,m_g)$ equal to $P_{\mathcal A/g}$ at $(m_g,2\sigma-m_g)$, gives the display. Let $i<c$. The letters of a level cell of $(\mathcal M',\mathcal N')$ on $\cE\sqcup G''$ are $K=\sum_Q\alpha^QK_{\mathcal L'_{c''-|Q|}}$ and $\Psi=\sum_Q\alpha^{G''\setminus Q}\Psi_{\mathcal L'_{c''-|Q|}}$, which at $\alpha_{G''}=0$ are $K_{\mathcal M'}$ and $\Psi_{\mathcal N'}$. So at $\alpha_{G'-g}=0$ the even letters of the hinge minors are
\[
\mathbf E_{\mathcal A/g}=K_{\mathcal L_{i-1}}\Psi_{\cB}=K_{\mathcal L_{i-1}}K_{\mathcal L_c},\qquad \mathbf E_{\mathcal A\setminus g}=K_{\mathcal L_i}\Psi_{\mathcal L_1}=K_{\mathcal L_i}K_{\mathcal L_{c-1}}.
\]
A constant ratio $\varrho_{\mathcal A}=\mathbf E_{\mathcal A/g}/\mathbf E_{\mathcal A\setminus g}$ on the orthant would make $\mathbf E_{\mathcal A/g}-\lambda\mathbf E_{\mathcal A\setminus g}$ vanish identically, also at $\alpha_{G'-g}=0$, against Lemma~\hb{lem:liftproducts} with $(i',i,j,j')=(i-1,i,c-1,c)$. The law of the cell is the image of its base form, with a positive continuous density on the open simplex, under the continuous ratio, and it is carried by one point only if the ratio is constant (Corollary~\hb{cor:piecesmarks}(i)).
\end{proof}

\begin{hproposition}[\hb{prop:freeextension} (The free extensions by one mark)]
As stated.
\end{hproposition}

\begin{proof}
Write $K_{\setminus},\Psi_{\setminus}$ and $K_{/},\Psi_{/}$ for the letters of $\mathcal A\setminus g$ and $\mathcal A/g$, of the degrees $r_{\mathcal A}$, $|\cE|-r_{\mathcal A}$, $r_{\mathcal A}-1$ and $|\cE|-r_{\mathcal A}+1$; $\mathcal A$ being plain, each hinge minor carries one system of weights in its two letters. If $\mathcal A/g$ is the dual of $\mathcal A\setminus g$ with its weights multiplied by $\mu$, then $K_/=\mu\Psi_\setminus$ and $\Psi_/=\mu K_\setminus$ (Proposition~\hb{prop:identity}), and $\varrho_{\mathcal A}=\mathbf E_{\mathcal A/g}/\mathbf E_{\mathcal A\setminus g}=\mu^2$. Conversely let $\varrho_{\mathcal A}$ be a constant $\lambda$. Connected tables on at least two marks are single molecules (Proposition~\hb{prop:ladder}(i)), so their letters are irreducible (Theorem~\hb{thm:irreducible}), and $K_/\Psi_/=\lambda K_\setminus\Psi_\setminus$ matches the irreducible factors on the two sides up to constant factors. $K_/$ and $K_\setminus$ have different degrees, so $K_/=\mu\Psi_\setminus$ with a constant $\mu$, positive since both have positive coefficients; comparing monomials, the bases of $\mathcal A/g$ are the complements in $\cE$ of the bases $T$ of $\mathcal A\setminus g$, with $\theta_{(\cE\setminus T)+g}=\mu\theta_T$. The degrees give $r_{\mathcal A}-1=|\cE|-r_{\mathcal A}$, so $\mathcal A/g$ has rank $r_{\mathcal A}-1$ on $2r_{\mathcal A}-1$ marks, the fusion width one; it is fusable by Theorem~\hb{thm:fusiontell}(d), $\mathcal A$ having $\mathcal A/g$ and $\mathcal A\setminus g$, with the bases of $(\mathcal A/g)^\ast$, as its children at the full and at the empty pattern of $\{g\}$. By Theorem~\hb{thm:fusion}(v) its fusion has the letter $K^\sharp=\Psi_/+\alpha_gK_/=\mu K_\setminus+\alpha_gK_/$, and $K^\sharp(\alpha_{\cE},\mu\alpha_g)=\mu K_\setminus+\mu\alpha_gK_/=\mu K_{\mathcal A}(\alpha_{\cE},\alpha_g)$: $\mathcal A$ has the bases of that fusion, the weights $\theta_{T'+g}$ on the bases $T'+g$ and $\theta_{T'+g}/\mu$ on the bases $\cE\setminus T'$, $T'$ a basis of $\mathcal A/g$. For a list every coefficient is $1$, so $\mu=1$: the complement in $\cE\sqcup\{g\}$ of a basis $T$ avoiding $g$ is the basis $(\cE\setminus T)+g$ and conversely, $\mathcal A$ is identically self-dual, $\mathcal A\setminus g=(\mathcal A/g)^\ast$, and $\mathcal A$ is the only major of $\mathcal A/g$, its fusion; an identically self-dual $\mathcal A$ has the ratio $1$, its hinge minors being dual with the same weights. At exponents $s,t>0$ a ratio is constant exactly when the law is carried by one point, the base form having a positive continuous density on the open simplex (Corollary~\hb{cor:piecesmarks}(i)).
\end{proof}

\begin{hcorollary}[\hb{cor:mot-pinneddual} (A pinned set of one list is a dual pair)]
As stated.
\end{hcorollary}

\begin{proof}
By Theorem~\hb{thm:mot-pointmass}, $S$ is pinned exactly when it has width one and $\varrho_S=1$. The cell of the bunch $\cB^+_S$ is a list with the hinge minors $\cB|S$ and $\cB/S'$ (\hb{thm:widthone}) and the ratio $\varrho_S$ (\hb{def:substitution}). Its hinge minors being connected, its ratio is constant exactly when it is identically self-dual (Proposition~\hb{prop:freeextension}), that is a dual pair (Proposition~\hb{prop:mot-dualletters}(i)), and then the ratio is $1$ (Proposition~\hb{prop:mot-dualletters}(ii)). So $\varrho_S=1$ exactly when $\cB^+_S$ is a dual pair, that is $\cB/S'=(\cB|S)^\ast$.
\end{proof}

\begin{hproposition}[\hb{prop:cousins} (Cousins)]
As stated.
\end{hproposition}

\begin{proof}
(i) is Corollary~\hb{cor:fusionisd}(i) with $C=\{x\}$, and its proof gives $\cB^\sharp\setminus x=(\cB^\sharp/x)^\ast$; for $x\in G$, $\cB^\sharp/x$ is the partial fusion at $G-x$, and at $c=1$, $\cB^\sharp/g=\cB$. (ii) Corollary~\hb{cor:mergedvalues} at $x$, which allows every pair of exponents with $s+t=2\sigma$, gives $P_{\cB^\sharp}=\Gst21{m_x,\ 2\sigma-m_x}{2\sigma}P_{\cB^\sharp\setminus x}(\mvec_{-x};2\sigma-m_x,m_x)$, and the dual $\cB^\sharp\setminus x$ of $\cB^\sharp/x$ has the letters exchanged (Proposition~\hb{prop:identity}); the same holds at $y$. (iii) $K^\sharp(\beta^\ast\alpha)=K^\sharp(\alpha)$, and the change of coordinates $u\mapsto\beta^\ast u$ permutes the coordinates of $\Delta_{\cE\sqcup G}$ and keeps $d\omega$, so the period at $\beta^\ast\mvec$ is the period at $\mvec$. At $c=1$ Corollary~\hb{cor:fusionone} writes $P_{\cB^\sharp}$ on its balance wall, where every proper wall is positive, as the stated function of $(\mvec|_{\cE},s')$, $s'=m_g$; an automorphism permutes the proper walls, $\kappa_U(\beta^\ast\mvec)=\kappa_{\beta(U)}(\mvec)$, so it keeps that part of the hyperplane, and the function inherits the invariance; at $\beta^\ast\mvec$ the mass of $g$ is $m_{\beta(g)}$ and the mass of $\beta^{-1}(g)$ is $m_g=s'$.
\end{proof}

\begin{htheorem}[\hb{thm:fusiontwosum} (Fusion across a two-sum)]
As stated.
\end{htheorem}

\begin{proof}
(i) $\mathcal S$ has the rank $r_1+r_2-1$ on $n_1+n_2-2$ marks (Lemma~\hb{lem:twosumpieces}), so the fusion width $c_1+c_2$. By Lemma~\hb{lem:gluemark} its bases are the $X\sqcup Y$ with $X$ a basis of $\mathcal A_1\setminus g$ and $Y$ of $\mathcal A_2/g$, or $X$ of $\mathcal A_1/g$ and $Y$ of $\mathcal A_2\setminus g$; since independent sets of a table extend to bases inside any set that contains them and has full rank, a set $I_1\sqcup I_2$ is independent exactly when $I_1$ and $I_2+g$ are independent in $\mathcal A_1$ and $\mathcal A_2$, or $I_1+g$ and $I_2$ are. A dependent set inside $\cE_1$ is dependent in $\mathcal A_1$, $\{g\}$ being independent in $\mathcal A_2$, so the circuits inside $\cE_1$ are the circuits of $\mathcal A_1$ avoiding $g$, and likewise for $\cE_2$. A circuit $C_1'\sqcup C_2'$ meeting both sides contains no such circuit, so $C_1'$ and $C_2'$ are independent; it is dependent exactly when $C_1'+g$ and $C_2'+g$ are dependent, that is when $g$ lies in the closures of $C_1'$ and of $C_2'$; and removing $e\in C_1'$ leaves an independent set exactly when $C_1'-e+g$ is independent, $C_2'+g$ being dependent. So the minimal ones are those with $C_1'+g$ and $C_2'+g$ circuits, and every such set is a circuit. The complement of $X\sqcup Y$ is $(\cE_1\setminus X)\sqcup(\cE_2\setminus Y)$, and $(\mathcal A_1\setminus g)^\ast=\mathcal A_1^\ast/g$, $(\mathcal A_2/g)^\ast=\mathcal A_2^\ast\setminus g$, so $\mathcal S^\ast=\mathcal A_1^\ast\oplus_g\mathcal A_2^\ast$, $g$ being neither a loop nor a coloop of $\mathcal A_i^\ast$; its circuits, the cocircuits of $\mathcal S$, are read off the same description.

(ii) A cocircuit $D$ of $\mathcal A_1$ avoiding $g$ contains no circuit of $\mathcal A_2$ and no set $(C_1\setminus g)\cup(C_2\setminus g)$, $C_2\setminus g$ being a nonempty part of $\cE_2$; so it is a union of circuits of $\mathcal S$ exactly when it is one of $\mathcal A_1$, and these cocircuits answer $a_1$; likewise $a_2$. For $D=(D_1\setminus g)\cup(D_2\setminus g)$ and $e\in D_1\setminus g$, the circuits of $\mathcal S$ inside $D$ through $e$ are the circuits of $\mathcal A_1$ inside $D_1$ avoiding $g$ through $e$ and the sets $(C_1\setminus g)\cup(C_2\setminus g)$ with $e\in C_1\subseteq D_1$ and $C_2\subseteq D_2$ circuits through $g$. So $e$ is covered in $D$ exactly when $D_1$ covers it by a circuit avoiding $g$, or by one through $g$ while $D_2$ covers $g$. $D_1$ and $D_2$ run independently over the cocircuits through $g$. If $h_2$ holds, the marks of the side $\cE_1$ are covered for every $D_2$ exactly when $u_1$; if $h_2$ fails, some $D_2$ does not cover $g$, and they are covered exactly when $s_1$, which implies $u_1$. So the side $\cE_1$ asks $u_1$ and ($s_1$ or $h_2$), the side $\cE_2$ asks $u_2$ and ($s_2$ or $h_1$), and Theorem~\hb{thm:fusiontell}(b) gives the rule. For one cell, its cocircuits avoiding $g$ answer $a$, and those through $g$ are unions of circuits exactly when their marks other than $g$ are covered, $u$, and $g$ is, $h$. Two fusable cells have $a_i,u_i,h_i$, and the rule holds.

(iii) By the bases of (i), $r_{\mathcal S}(\cE_2)=r_2$ and $d_{\mathcal S}(\cE_2)=r_2-1$, so $r^\ast_{\mathcal S}(\cE_2)=|\cE_2|-d_{\mathcal S}(\cE_2)=n_2-r_2$. A fusable $\mathcal S$ has $r_{\mathcal S}\le r^\ast_{\mathcal S}$ pointwise, $\mathcal S$ being a quotient of $\mathcal S^\ast$ (Proposition~\hb{prop:paircell}(i)): $r_2\le n_2-r_2$, $c_2\ge0$, and likewise $c_1\ge0$. In $\mathcal S^\sharp$, $r^\sharp(\cE_2)=r^\ast_{\mathcal S}(\cE_2)=n_2-r_2$, the rank of a set avoiding $G$ being its rank in $\mathcal S^\sharp\setminus G=\mathcal S^\ast$. An identically self-dual table has $r^\ast(A)=r(A)$, so $d(A)=|A|-r(A)$ and $c(A)=2r(A)-|A|$, which for $A=\cE_2$ is $2(n_2-r_2)-(n_2-1)=c_2+1$; likewise for $\cE_1$.

(iv) In $\mathcal A_i^\sharp$ the mark $g$ lies in a basis $X\sqcup G_i$, $X$ a basis of $\mathcal A_i$ through $g$, and avoids a basis $X$ of $\mathcal A_i^\ast$, so $\mathcal A_i^\sharp$ is a cell with the hinge $g$ and $W=\mathcal A_1^\sharp\oplus_g\mathcal A_2^\sharp$ is a table on $\cE_1\sqcup\cE_2\sqcup G_1\sqcup G_2$. By (i) $W^\ast=\mathcal A_1^{\sharp\ast}\oplus_g\mathcal A_2^{\sharp\ast}=W$. Contracting $G=G_1\sqcup G_2$ commutes with the minors at $g$, marks of different sets (Definition~\hb{def:minors}), so the bases of $W$ containing $G$, less $G$, are those of $(\mathcal A_1^\sharp/G_1)\oplus_g(\mathcal A_2^\sharp/G_2)=\mathcal S$, and likewise $W\setminus G=\mathcal A_1^\ast\oplus_g\mathcal A_2^\ast=\mathcal S^\ast$: $W$ is an identically self-dual major. If $c_2=0$, $\mathcal A_2$ is a quotient of $\mathcal A_2^\ast$ of the same rank, so the two rank functions agree (Proposition~\hb{prop:paircell}(i) at $X\subseteq Y$ and at $Y\subseteq\cE$) and $\mathcal A_2=\mathcal A_2^\ast=\mathcal A_2^\sharp$; every transposition of two marks of $G_1$ carries $\mathcal A_1^\sharp$, hence $W$, and $W=\mathcal S^\sharp$ by Theorem~\hb{thm:fusion}(iv); likewise for $c_1=0$. Let $c_1,c_2\ge1$, $x_1\in G_1$, $x_2\in G_2$, and let $r^\ast_1=n_1-r_1$, the rank of $\mathcal A_1^\sharp$. A basis $Z_1\sqcup Z_2$ of $W$ has $|Z_1|\in\{r^\ast_1,r^\ast_1-1\}$. At the pattern $\{x_1\}$ of $G$, $|Z_1\cap\cE_1|=|Z_1|-1\le r^\ast_1-1$, with equality for $Z_1=X_1\sqcup\{x_1\}$, $X_1$ a basis of $\mathcal L_{c_1-1}(\mathcal A_1^\ast,\mathcal A_1)$ avoiding $g$, and $Z_2=X_2-g$, $X_2$ a basis of $\mathcal A_2^\ast$ through $g$; such bases exist, $g$ being independent in $\mathcal A_1$ and so in its lift $\mathcal L_1=\mathcal L_{c_1-1}^\ast$, and not a loop of $\mathcal A_2^\ast$. At the pattern $\{x_2\}$, $Z_1$ a basis of $\mathcal A_1^\ast$ avoiding $g$ and $Z_2=(X_2\sqcup\{x_2\})-g$, $X_2$ a basis of $\mathcal L_{c_2-1}(\mathcal A_2^\ast,\mathcal A_2)$ through $g$, give $|Z_1\cap\cE_1|=r^\ast_1$. So the two children have the ranks $r^\ast_1-1$ and $r^\ast_1$ on $\cE_1$, they differ, and $G$ is not in level form in $W$, while it is in $\mathcal S^\sharp$.

(v) If both cells are fusable, $\mathcal S$ is fusable by (ii), and for $c_2=0$ the fusable list $\mathcal A_2$ of the width zero equals its dual, as in (iv), so it is identically self-dual, and $\mathcal S^\sharp=\mathcal A_1^\sharp\oplus_g\mathcal A_2$ by (iv). Conversely let $\mathcal S$ be fusable; by (iii) $\{c_1,c_2\}=\{0,1\}$; let $c_2=0$. By (iii) $\cE_2$ is a bunch of width one of $\mathcal S^\sharp$, a list, and by Theorem~\hb{thm:widthone} and Proposition~\hb{prop:substitution}(iii) $\mathcal S^\sharp$ is the two-sum of the cell $\mathcal N_2$ of the bunch $\cE_2$ and the cell $\mathcal N_1$ of its room $\cE_1\sqcup G$, along a hinge $h$. The cell of $\cE_2$ in the dual table has the bases $\mathcal S^{\sharp\ast}|\cE_2=(\mathcal S^\sharp/(\cE_1\sqcup G))^\ast$ and $\mathcal S^{\sharp\ast}/(\cE_1\sqcup G)+h=(\mathcal S^\sharp|\cE_2)^\ast+h$, the complements in $\cE_2+h$ of the bases of $\mathcal N_2$; $\mathcal S^\sharp$ being identically self-dual, $\mathcal N_2$ is. In $\mathcal S=\mathcal S^\sharp/G$ the cell of $\cE_2$ is $\mathcal A_2$ (Lemma~\hb{lem:gluemark}); its bases through the hinge, $\mathcal S/\cE_1=\mathcal S^\sharp/(\cE_1\sqcup G)$, are those of $\mathcal N_2$, and its bases avoiding the hinge are those of $(\mathcal S^\sharp/G)|\cE_2$. Here $r^\sharp(\cE_2)=n_2-r_2=r_2$ by (iii) and $c_2=0$, while $r^\sharp(\cE_2\sqcup G)=r_{\mathcal S}(\cE_2)+|G|=r_2+|G|$; by submodularity every independent $I\subseteq\cE_2$ of $\mathcal S^\sharp$ has $r^\sharp(I\sqcup G)\ge|I|+r^\sharp(\cE_2\sqcup G)-r^\sharp(\cE_2)=|I|+|G|$, so $(\mathcal S^\sharp/G)|\cE_2=\mathcal S^\sharp|\cE_2$ and $\mathcal A_2=\mathcal N_2$, identically self-dual. The cell of the room $\cE_1$ in $\mathcal S$ has the bases of $\mathcal S|\cE_1=(\mathcal S^\sharp|(\cE_1\sqcup G))/G$ and of $\mathcal S/\cE_2+h=(\mathcal S^\sharp/\cE_2)/G+h$, so $\mathcal A_1=\mathcal N_1/G$; $\mathcal N_1$ is identically self-dual by the same argument as $\mathcal N_2$, so $\mathcal N_1\setminus G=(\mathcal N_1/G)^\ast$ and $\mathcal N_1$ is a major of $\mathcal A_1$: $\mathcal A_1$ is fusable (Theorem~\hb{thm:fusiontell}), of the width $|G|=1$, and $\mathcal N_1=\mathcal A_1^\sharp$ (Theorem~\hb{thm:fusion}(v)). So $\mathcal S^\sharp=\mathcal A_1^\sharp\oplus_g\mathcal A_2$, as (iv) also gives: $\mathcal S^\sharp$ is the two-sum along the hinge $g$ of $\mathcal S$, with $\cE_2$ a bunch of width one.
\end{proof}

\begin{hcorollary}[\hb{cor:isdcell} (A two-sum with an identically self-dual cell)]
As stated.
\end{hcorollary}

\begin{proof}
Every proper wall of $\mathcal S$ positive makes $P_{\mathcal S}$ finite (Theorem~\hb{thm:projective}); $\kappa$ and $2\sigma-\kappa$ are the walls of $\mathcal S$ at $\cE_1$ and $\cE_2$ (Lemma~\hb{lem:twosumpieces}), so $0<\kappa<2\sigma$. $\mathcal A_1$ is identically self-dual, so $\mathcal A_1/g$ is the dual of $\mathcal A_1\setminus g$ with the same weights, its ratio is $1$, and $\phi_{\mathcal A_1}(\Delta_{\cE_1})=P_{\mathcal A_1\setminus g}(\mvec|_{\cE_1};\kappa,2\sigma-\kappa)$ (Corollary~\hb{cor:piecesmarks}(ii)), finite since $\mathcal A_1$ sits in a gluing of finite period. By Corollary~\hb{cor:piecesmarks}(i) $\mathcal A_1$ acts in $\mathcal S$ as this number times the one-mark cell of ratio $1$ and mass $\kappa$, whose two-sum with $\mathcal A_2$ is $\mathcal A_2$ with its hinge carrying $\kappa$: $P_{\mathcal S}=\phi_{\mathcal A_1}(\Delta_{\cE_1})\,P_{\mathcal A_2}(\mvec|_{\cE_2},\kappa;s,t)$. The Beta corner of Corollary~\hb{cor:piecesmarks}(iii) at $g$ with the mass $2\sigma-\kappa$, which puts $\mathcal A_1$ on its balance wall (Lemma~\hb{lem:twosumpieces}), is the integral over the scale of $g$, valid in $[0,\infty]$ for $0<\kappa<2\sigma$ by Tonelli: $P_{\mathcal A_1}(\mvec|_{\cE_1},2\sigma-\kappa;s,t)=\Gst21{2\sigma-\kappa,\ \kappa}{2\sigma}P_{\mathcal A_1\setminus g}(\mvec|_{\cE_1};\kappa,2\sigma-\kappa)$. When the cells are lists and $\mathcal A_2$ is fusable of the fusion width one, $\mathcal A_1$, identically self-dual, is fusable of the width zero, and Theorem~\hb{thm:fusiontwosum}(v), with the roles of the two cells exchanged, gives $\mathcal S^\sharp=\mathcal A_1\oplus_g\mathcal A_2^\sharp$.
\end{proof}

\begin{hproposition}[\hb{prop:repair} (Uncovered marks and repair)]
As stated.
\end{hproposition}

\begin{proof}
(i) A cocircuit is a union of circuits exactly when it covers each of its marks, so $\cB$ has no uncovered mark exactly when (b) of Theorem~\hb{thm:fusiontell} holds, equivalently (a). By that theorem a list, or a table of the width $c\ge1$, is fusable exactly when (a) holds; at $c<0$ neither (a) holds nor is the table fusable; and at $c=0$ the table is fusable exactly when it is identically self-dual (Definition~\hb{def:fusion}). (ii) By Theorem~\hb{thm:blockcounts}(i) and (ii), under the first rule alone $\cB_S$ has $n+|S|$ marks, the rank $r$ and the rank function $Y\mapsto r(\lceil Y\rceil)$, $\lceil Y\rceil$ the marks of $\cB$ whose blocks meet $Y$. So $Y$ is independent exactly when it meets each block in at most one mark and $\lceil Y\rceil$ is independent: the circuits of $\cB_S$ are the pairs $\{h,h'\}$, $h\in S$, and the sets meeting each block in at most one mark whose shadow $\lceil\cdot\rceil$ is a circuit of $\cB$. A closed set contains with a mark its whole block, adding a mark of a block already met not raising the rank; so the closed sets are the unions of the blocks over the closed sets of $\cB$, and since the cocircuits are the complements of the closed sets of rank $r-1$ (fact ($\beta$) in the proof of Theorem~\hb{thm:fusiontell}) the cocircuits of $\cB_S$ are $D^+=D\cup\{h':\ h\in D\cap S\}$, $D$ a cocircuit of $\cB$. In $D^+$ the marks $h,h'$ of a block are covered by $\{h,h'\}$, and a mark $e\in D\setminus S$ is covered exactly when a circuit of $\cB_S$ through $e$ lies in $D^+$, that is when a circuit of $\cB$ through $e$, its shadow, lies in $D$. So $\cB_S$ has no uncovered mark exactly when every mark of $\cE\setminus S$ is covered in every cocircuit of $\cB$ through it: when $S$ contains the uncovered marks; and by (i) for $\cB_S$, a list or a table of the width $c+|S|\neq0$, it is then fusable exactly when it has no uncovered mark. Its fusion width is $n+|S|-2r$. The bases of $\cB_S$ avoiding $S'$ are the bases $T$ of $\cB$, each with its weight $\theta_T$ (Theorem~\hb{thm:block}), so $(\cB_S,\theta)\setminus S'=(\cB,\theta)$. (iii) $(\cB_S)^\sharp/G=(\cB_S,\theta)$ and $(\cB_S,\theta)\setminus S'=(\cB,\theta)$, the two minors being taken at disjoint sets. Let $c\ge0$, $(\cB,\theta)$ not fusable, and $U$ its set of uncovered marks, none a loop, since a loop lies in no cocircuit. If $U\neq\varnothing$, $S=U$ gives the width $c+|U|\ge1$ and no uncovered mark, and $\cB_U$ is fusable by (ii). If $U=\varnothing$, (i) leaves $c=0$ and $(\cB,\theta)$ not identically self-dual, while $\cB=\cB^\ast$ by (a) and the proof of Definition~\hb{def:fusion}; a loop of $\cB$ would be a coloop of $\cB^\ast=\cB$, so there is none, and $\cE\neq\varnothing$, since on no marks the one basis $\varnothing$ is its own complement with its own weight. For any mark $h$, $S=\{h\}$ gives the width one and no uncovered mark, and $\cB_S$ is fusable by (ii). For a list $\cB$, the set of the marks that are not loops contains $U$, and the list $\cB_S$ it doubles has no uncovered mark and is fusable by (i). (iv) The copy $h'$ is neither a loop nor a coloop, so its wall at $t=0$ is $\kappa_{\{h'\}}=m_{h'}$. Theorem~\hb{thm:wallres} at the bunch $\{h'\}$, whose bottom level holds the bases avoiding $h'$, with the empty pattern alone, so that the weights factor on it, gives the residue as the product of the periods of the faces: on $\{h'\}$ the contraction of the room, of rank $d(\{h'\})=0$ and letter $1$, whose period on one mark is $1$; on the room, $(\cB_S|(\cE_S\setminus h');s,0)=(\cB_{S-h};s,0)$.
\end{proof}


\section{Chapter~\hb{ch:closeglue}: closing a glue along a direction}\label{sec:ch29}

\begin{htheorem}[\hb{thm:closecone} (Closing the glue along a direction)]
As stated, the largest values of the monomials of the letter vector of the partner being read on its law when its room word has a Newton sum of full dimension.
\end{htheorem}

\begin{proof}
Apart from the sentence on the partner, the statement is the reference's, with its proof (the reference, proof of thm:closecone). Let the room word of the partner have a Newton sum of full dimension. The partner is a piece on $R$ with the stations $\mathcal D$, read with the reversed excesses $c_j-i(D)$ (\hb{def:seampiece}); its scale $\gamma^\rho_R$ and its seam vector $\boldsymbol\varrho^\rho_R=(\varrho^\rho_{R,jD},\varrho^{\rho\ast}_{R,jD})_{j,D}$ are those of \hb{def:seamlaw} with its room letters and these excesses. In this proof write $L^{(0)}_{j,R}$ for the sum of the $L_{jD}$ over the stations of excess $c_j$, which have the reversed excess $0$, and $L^{\ast(c)}_{j,R}$ for the sum of the $L^\ast_{jD}$ over the stations of excess $0$, which have the reversed excess $c_j$. By \hb{def:seamlaw},
\[
L_{jD}=L^{(0)}_{j,R}\,(\gamma^\rho_R)^{c_j-i(D)}\,\varrho^\rho_{R,jD},\qquad L^\ast_{jD}=L^{\ast(c)}_{j,R}\,(\gamma^\rho_R)^{i(D)}\,\varrho^{\rho\ast}_{R,jD}.
\]
A direction $\mathbf v\in V$ keeps the letter totals, $\sum_Dv_{jD}=0$ and $\sum_Dv^\ast_{jD}=0$ for every active letter. It also keeps the pin of \hb{thm:transversepins} for every vector of $M$, the pin of $\one$ holding by the letter totals on the balance wall, and $\one_R\in\vec{\mathcal N}_R^\perp\subseteq M$, the room letters being homogeneous. The exponents in $\mathbf n(\mathbf a,\mathbf b)$ of the monomials of $V_{jD}$ and $V^\ast_{jD}$ vanish on $R$, and those of the monomials of $L_{jD}$ and $L^\ast_{jD}$ have the totals $\deg L_{jD}=e_j+c_j-i(D)$ and $\deg L^\ast_{jD}=e^\ast_j+i(D)$ on $R$ (\hb{def:seampiece}); so, by the letter totals,
\[
0=\ip{\mathbf n(\mathbf v)}{\one_R}=\sum_{j,D}\big(v_{jD}\deg L_{jD}+v^\ast_{jD}\deg L^\ast_{jD}\big)=\sum_{j,D}\big((c_j-i(D))v_{jD}+i(D)v^\ast_{jD}\big).
\]
So in $\mathbf V_R^{\mathbf v}=\prod_{j,D}L_{jD}^{v_{jD}}L^\ast_{jD}{}^{v^\ast_{jD}}$ the sums $L^{(0)}_{j,R}$ and $L^{\ast(c)}_{j,R}$ and the scale $\gamma^\rho_R$ enter with the exponent $0$, and $\mathbf V_R^{\mathbf v}=(\boldsymbol\varrho^\rho_R)^{\mathbf v}$ on the open orthant of $R$, where these sums are positive, both sides being homogeneous of degree zero; hence $\sup\mathbf V_R^{\mathbf v}=\sup_{v\in\Delta_R}\boldsymbol\varrho^\rho_R(v)^{\mathbf v}$, the supremum over the open simplex. The scale form $\phi^\rho_R$ has a positive continuous density on the open simplex and the seam vector is continuous there, so the law $(\boldsymbol\varrho^\rho_R)_\ast\phi^\rho_R$ is carried by the closure of the image of the open simplex, and every point of that image lies in its support: the preimage of a neighbourhood of an image point is a nonempty open set, of positive mass. The supremum of the continuous function $\mathbf x^{\mathbf v}$ over the image is therefore its largest value on the support of the law, in $(0,\infty]$.
\end{proof}

\begin{slemma}\label{lem:ch29-fan}
Let $L_1,\dots,L_p$ be positive letters on a finite set $B$ (\hb{def:posword}), $e_0\in B$, $e_1,\dots,e_p$ real exponents and $\mvec$ real masses on $B$ with $m_B=\sum_\ell e_\ell\deg L_\ell$; the pairings below are taken on $B\setminus e_0$.
\begin{enumerate}[label=(\roman*)]
\item There are finitely many rational simplicial cones of dimension $|B|-1$ with disjoint interiors covering $\mathbb R^{B\setminus e_0}$, each contained, for every $\ell$, in a set $\{\mathbf v:\ \ip{\mathbf t}{\mathbf v}\le\ip{\mathbf t'}{\mathbf v}\ \text{for every exponent}\ \mathbf t'\ \text{of the support of}\ L_\ell\}$ for one exponent $\mathbf t$ of that support.
\item On such a cone $\pi$, with the primitive integral generators $\mathbf v_{\pi,1},\dots,\mathbf v_{\pi,|B|-1}$, write $x=\log\alpha_{B\setminus e_0}=-\sum_ku_k\mathbf v_{\pi,k}$ on the section $\alpha_{e_0}=1$, $u_k\ge0$, and $q_k=e^{-u_k}$. The exponent $\mathbf t_{\ell,\pi}$ of (i) is unique, and $L_\ell=\prod_kq_k^{\ip{\mathbf t_{\ell,\pi}}{\mathbf v_{\pi,k}}}\,k_{\ell,\pi}(q)$ on the section, with $k_{\ell,\pi}$ a polynomial in $q$ with nonnegative coefficients whose constant term is the coefficient of $\alpha^{\mathbf t_{\ell,\pi}}$ in $L_\ell$.
\item For a generator $\mathbf v$ of $\pi$ the number $\kappa_{\mathbf v}=\ip{\mvec-\sum_\ell e_\ell\mathbf t_{\ell,\pi}}{\mathbf v}$ is the same for every cone $\pi$ of which $\mathbf v$ is a generator, and, in $[0,\infty]$,
\begin{multline*}
\int_{\alpha_{e_0}=1}\alpha^{\mvec-\one}\prod_\ell L_\ell^{-e_\ell}\,d\alpha_{B\setminus e_0}\\
=\sum_\pi|\det(\mathbf v_{\pi,1},\dots,\mathbf v_{\pi,|B|-1})|\int_{[0,1]^{|B|-1}}\prod_kq_k^{\kappa_{\mathbf v_{\pi,k}}-1}\prod_\ell k_{\ell,\pi}(q)^{-e_\ell}\,dq ;
\end{multline*}
both sides are finite exactly when every $\kappa_{\mathbf v}$ is positive, and they are then the period of the word $\prod_\ell L_\ell^{-e_\ell}$ at $\mvec$, the integral over a section of \hb{def:posword}.
\end{enumerate}
\end{slemma}

\begin{proof}
(i) Identify $\mathbb R^{B\setminus e_0}$ with the hyperplane $v_{e_0}=0$ of $\mathbb R^B$. For each $\ell$ the sets $C_{\ell,\mathbf t}=\{\mathbf v:\ \ip{\mathbf t}{\mathbf v}\le\ip{\mathbf t'}{\mathbf v}\ \text{for every}\ \mathbf t'\}$, $\mathbf t$ in the support of $L_\ell$, are rational polyhedral cones, defined by the inequalities $\ip{\mathbf t-\mathbf t'}{\mathbf v}\le0$: the cones of the inner normal fan of the Newton polytope of $L_\ell$ cut by that hyperplane, which form a fan covering $\mathbb R^{B\setminus e_0}$, every $\mathbf v$ having an exponent of least pairing. The cones $\bigcap_\ell C_{\ell,\mathbf t_\ell}\cap O$, $O$ running over the closed coordinate orthants of $\mathbb R^{B\setminus e_0}$, are pointed rational polyhedral cones which with their faces form a fan covering $\mathbb R^{B\setminus e_0}$, the common refinement of these fans and of the fan of the orthants. Order its rays and triangulate each cone by pulling: the cone is the union of the cones spanned by its first ray and by the triangulated faces that do not contain that ray. A face $F$ of a pointed cone $C$ has $\operatorname{span}F\cap C=F$, so the first ray of $C$ lies off $\operatorname{span}F$ for such a face, and the cones produced are simplicial; the pulling triangulation of a cone restricts on each face to the pulling triangulation of that face, which depends on the face alone, so the triangulations of two cones agree on their common face. This gives finitely many rational simplicial cones of dimension $|B|-1$ with disjoint interiors covering $\mathbb R^{B\setminus e_0}$, each in some $\bigcap_\ell C_{\ell,\mathbf t_\ell}$.

(ii) If $\mathbf t$ and $\mathbf t''$ both pair least with every $\mathbf v_{\pi,k}$, then $\ip{\mathbf t-\mathbf t''}{\mathbf v_{\pi,k}}=0$ for every $k$; the generators span $\mathbb R^{B\setminus e_0}$, so $\mathbf t$ and $\mathbf t''$ agree on $B\setminus e_0$, and as $|\mathbf t|=|\mathbf t''|=\deg L_\ell$ they are equal. On the section $\alpha^{\mathbf t'}=e^{\ip{\mathbf t'}{x}}=\prod_kq_k^{\ip{\mathbf t'}{\mathbf v_{\pi,k}}}$, so with $L_\ell=\sum_{\mathbf t'}c_{\mathbf t'}\alpha^{\mathbf t'}$ and $\mathbf t=\mathbf t_{\ell,\pi}$,
\[
L_\ell=\prod_kq_k^{\ip{\mathbf t}{\mathbf v_{\pi,k}}}\sum_{\mathbf t'}c_{\mathbf t'}\prod_kq_k^{\ip{\mathbf t'-\mathbf t}{\mathbf v_{\pi,k}}}.
\]
Every exponent $\ip{\mathbf t'-\mathbf t}{\mathbf v_{\pi,k}}$ is an integer, the exponents and the generators being integral, and it is nonnegative by the choice of $\mathbf t$; by the uniqueness just proved, all of them vanish only at $\mathbf t'=\mathbf t$. So the sum is a polynomial $k_{\ell,\pi}$ with nonnegative coefficients and the constant term $c_{\mathbf t}>0$.

(iii) A cone $\pi$ of which $\mathbf v$ is a generator has $\ip{\mathbf t_{\ell,\pi}}{\mathbf v}=\min_{\mathbf t'}\ip{\mathbf t'}{\mathbf v}$, the minimum over the support of $L_\ell$, which does not depend on $\pi$; so neither does $\kappa_{\mathbf v}$. On $\pi$ the map $(u_k)_k\mapsto x=-\sum_ku_k\mathbf v_{\pi,k}$ is a linear bijection of $[0,\infty)^{|B|-1}$ onto $-\pi$ with the Jacobian $|\det(\mathbf v_{\pi,1},\dots,\mathbf v_{\pi,|B|-1})|$; on the section $d\alpha_e=\alpha_e\,dx_e$ for $e\neq e_0$, $\alpha^{\mvec}=e^{\ip{\mvec}{x}}$, and $du_k=-dq_k/q_k$. By (ii), on the image of $-\pi$, the form $\alpha^{\mvec-\one}\prod_\ell L_\ell^{-e_\ell}\,d\alpha_{B\setminus e_0}=e^{\ip{\mvec}{x}}\prod_\ell L_\ell^{-e_\ell}\,dx$ is
\[
|\det(\mathbf v_{\pi,1},\dots,\mathbf v_{\pi,|B|-1})|\prod_kq_k^{\ip{\mvec-\sum_\ell e_\ell\mathbf t_{\ell,\pi}}{\mathbf v_{\pi,k}}-1}\prod_\ell k_{\ell,\pi}(q)^{-e_\ell}\,dq .
\]
The cones $-\pi$, like the cones $\pi$, cover $\mathbb R^{B\setminus e_0}$ with disjoint interiors, their boundaries are null and the integrands are positive, so the two sides agree in $[0,\infty]$ (Tonelli's theorem). On $[0,1]^{|B|-1}$ each $k_{\ell,\pi}$ lies between its constant term and the sum of its coefficients, so $\prod_\ell k_{\ell,\pi}^{-e_\ell}$ is bounded above and below by positive constants, and the integral over the cube is finite exactly when every exponent $\kappa_{\mathbf v_{\pi,k}}$ is positive; every generator is a generator of some cone. The form $\alpha^{\mvec-\one}\prod_\ell L_\ell^{-e_\ell}d\alpha$ is projective, the masses lying on the balance hyperplane, and $\alpha_{e_0}$ is positive and homogeneous of degree one, so the integral over the section is the period (\hb{def:posword}, \hb{lem:chart}).
\end{proof}

\begin{htheorem}[\hb{thm:pastwalls} (Closing past the walls)]
As stated, for every piece whose seam function converges at a real point of the pins of the seam with positive shares, with the sector form of (ii) for every fan of the seam word and with part (iv).
\end{htheorem}

\begin{proof}
\emph{The hypothesis.} At a real point of the pins with positive shares the seam word $W_B(\mathbf a,\mathbf b)$ is a Newton word (\hb{def:posword}), and the masses $\mvec_B$ lie on its balance hyperplane: by the letter totals and the scale pin,
\[
\sum_{j,D}\Big(a_{jD}\big(d_j+i(D)\big)+b_{jD}\big(d^\ast_j+c_j-i(D)\big)\Big)=\sum_j\big(s_jd_j+t_jd^\ast_j\big)+\kappa_B=m_B
\]
(\hb{def:seampiece}, \hb{lem:seamword}). Its period, the seam function, converges there, so its Newton sum has the dimension $|B|-1$ (\hb{thm:poswordconv}).

\emph{(i).} The reference proves (i) for a piece of tables (the reference, proof of thm:pastwalls(i)). Of the piece it uses only that its seam word has a Newton sum of full dimension, so that its reduced period in \hb{thm:transversepins} is its seam function at the sums over the monomials, the seam letters repeating over the monomials of the partner; the reference obtains that fact from the connectedness of the family of tables, and here it is the hypothesis. The rest of that proof concerns the partner, whose room letters are read in their monomials, so that $\vec{\mathcal N}_R=0$ and its reduced period is $\prod\xi^{-\mathbf a}\prod\xi^{\ast-\mathbf b}$ (\hb{thm:transversepins}); with $M=\mathbb R\one_B\oplus\mathbb R^R$ and the lattice basis $\one$, $\mathbf e_r$ ($r\in R$), the pins of \hb{thm:transversepins} are the displayed ones, the scale pin following by the balance.

\emph{(ii), the sector form.} On the pins of the seam the exponents $\bar a_{jD}$, $\bar b_{jD}$ of the seam word are affine in the shares with rational linear parts, and the masses $\mvec_B$ lie on its balance hyperplane at every share, by the display above with the sums in place of the shares. Lemma~\ref{lem:ch29-fan}, applied to the seam letters of the active letters with the exponents $\bar a_{jD}$, $\bar b_{jD}$, gives the fans of the seam word, the sector coordinates, the sector letters with their exponents, the walls and the identity at every real point of the pins, and shows that the seam function converges at a real point exactly when every wall is positive there. The coefficients $\ip{\mathbf t_{jD,\pi}}{\mathbf v}$ and $\ip{\mathbf t^\ast_{jD,\pi}}{\mathbf v}$ of the shares in a wall are integers, so the wall is affine with a rational linear part, and it does not depend on the sector (Lemma~\ref{lem:ch29-fan}(iii)).

\emph{(ii), continuation and bound.} On a sector the integrand is $\prod_kq_k^{\kappa_{\mathbf v_{\pi,k}}(\mathbf z)-1}h(q,\mathbf z)$ with $h=\prod_{j,D}k_{jD,\pi}^{-\bar a_{jD}}\psi_{jD,\pi}^{-\bar b_{jD}}$. The sector letters are at least their positive constant terms on the closed cube, so their real parts are positive on a neighbourhood of it, where their principal logarithms are holomorphic; so $h$ is holomorphic in $q$ near $[0,1]^{|B|-1}$ and entire in the shares. \hb{lem:cube}, with $x_k=\kappa_{\mathbf v_{\pi,k}}(\mathbf z)$ and the shares as its parameter, continues each sector integral meromorphically to all complex shares, with its poles on the hyperplanes on which a wall of a generator of the sector lies in $-\mathbb N$; a wall constant on the pins is positive there, the seam function converging, and gives no hyperplane. The formula of \hb{lem:cube}, repeated in the $|B|-1$ coordinates, writes each term with at most $|B|-1$ factors $(\kappa_{\mathbf v_{\pi,k}}+l)^{-1}$, which bounds the orders of the poles by $|B|-1$. For the bound, the proof of \hb{lem:cubegrowth} is run on each sector integral with the several parameters entering through the powers of the sector letters and through the walls, $|\re\mathbf z|_1$ and $|\im\mathbf z|_1$ in place of $|\re z|$ and $|\im z|$ (the reference, proof of thm:pastwalls(ii)): each of the at most $|B|-1$ factors $(\kappa+l)^{-1}$ of a term is at most a constant times $\delta^{-1}$, and nothing else depends on $\delta$.

\emph{(ii), a piece of tables.} Let the seam letters at the station $D$ of the active letters be the basis and the cobasis polynomials of a table $\cB_{jD}$ on $B$ with positive weights, and consider the family of these tables at the exponents $\bar a_{jD},\bar b_{jD}$ at a point with positive shares. Its potential $\sum_{j,D}\big(\bar a_{jD}\log K_{\cB_{jD}}(e^y)+\bar b_{jD}\log\Psi_{\cB_{jD}}(e^y)\big)$ (\hb{def:potential}) is affine along a direction $\mathbf w$ exactly when, in every letter, all the exponents of the support have one pairing with $\mathbf w$, a logarithm of a sum of exponentials with positive coefficients being strictly convex along every direction that separates two of its exponents: exactly when $\mathbf w$ is orthogonal to the directions of the Newton sum. If the family is connected, that space is $\mathbb R\one$ (\hb{thm:saddle}(iii), a space that does not involve the masses), and the Newton sum has the dimension $|B|-1$. If it is not, some nonempty proper $S\subset B$ has the width $c_{jD}(S)=0$ in every table (\hb{def:family}), so that $|T\cap S|=d_{jD}(S)$ for every basis $T$ of $\cB_{jD}$ (\hb{thm:window}). The base polytope of $\cB_{jD}$ then lies in $\{x_S=d_{jD}(S)\}$ and that of its dual in $\{x_S=|S|-d_{jD}(S)\}$, so the Newton sum lies in a hyperplane $\{x_S=\text{const}\}$, and as $\one_S\notin\mathbb R\one$ its dimension is less than $|B|-1$.

In the coordinates of \hb{thm:exist}(ii) on the sector $H_\pi$ of an order $\pi$ of $B$, $\alpha_{\pi(1)}=\rho$ and $\alpha_{\pi(k)}=\rho q_2\cdots q_k$, so that, with $u_k=-\log q_k$,
\[
\log\alpha_{\pi(m)}-\log\alpha_{\pi(1)}=\sum_{k=2}^m\log q_k=-\sum_{k=2}^{|B|}u_k\,(\one_{U_k})_{\pi(m)},
\]
because $\pi(m)\in U_k$ exactly when $k\le m$. So, modulo $\one$, $H_\pi$ is the cone with the generators $\one_{U_k}$, $2\le k\le|B|$, read on the section as $\one_{U_k}-(\one_{U_k})_{e_0}\one$. These $|B|-1$ vectors form a basis of $\mathbb Z^B/\mathbb Z\one$, since the differences $\one_{U_k}-\one_{U_{k+1}}=\mathbf e_{\pi(k)}$ and $\one_{U_1}=\one$ span $\mathbb Z^B$; so the cone is simplicial of dimension $|B|-1$ with the determinant $\pm1$, and the sectors of the orders cover $\mathbb R^B/\mathbb R\one$ with disjoint interiors. For a basis $T$, $\alpha^T/\alpha^{T^+_\pi}=\prod_kq_k^{|T\cap U_k|-|T^+_\pi\cap U_k|}$, and \hb{lem:greedy}, $\alpha^T\le\alpha^{T^+_\pi}$ on $H_\pi$ for every basis, gives, with one $q_k$ small and the others $1$, $\ip{\one_{T^+_\pi}}{\one_{U_k}}\le\ip{\one_T}{\one_{U_k}}$ for every $k$. As $\one_T$ and $\one_{T^+_\pi}$ have one total, their pairings on $B\setminus e_0$ with the generator $\one_{U_k}-(\one_{U_k})_{e_0}\one$ differ by the difference of their pairings with $\one_{U_k}$. So the sector lies in the set of (ii) for $K_{\cB_{jD}}$ with $\mathbf t=\one_{T^+_\pi}$, and likewise for $\Psi_{\cB_{jD}}$ with $\mathbf t=\one_{B\setminus T^-_\pi}$, the greedy bases being those of $\cB_{jD}$; the sectors of the orders form a fan of the seam word. Its sector letters are $\sum_T\theta_{jD,T}\prod_kq_k^{|T\cap U_k|-d_{jD}(U_k)}$ and their cobasis counterparts, the sector letters of \hb{thm:exist}(ii). The vector $\mvec_B-\sum_{j,D}(\bar a_{jD}\one_{T^+}+\bar b_{jD}\one_{B\setminus T^-})$ has the total $0$ on the pins, so its pairing with the generator of a tail $U$ is its pairing with $\one_U$; by \hb{lem:greedy}, $|T^+\cap U|=d_{jD}(U)$ and $|U\setminus T^-|=d^\ast_{jD}(U)$, and the wall of that generator is $m_U-\sum_{j,D}(\bar a_{jD}d_{jD}(U)+\bar b_{jD}d^\ast_{jD}(U))$, the proper wall $\kappa_U$. Every nonempty proper $U$ is a tail of some order.

\emph{(iii).} The reference proves (iii) for a piece of tables (the reference, proof of thm:pastwalls(iii)), using three facts about the piece. The first is the continuation with the bound of (ii). The second is that the seam function converges at a real point exactly when its walls are positive there, so that it converges along a ray along which no wall falls, the walls staying above their positive values at $\bar{\mathbf c}$, and diverges past a wall that falls. The third is that at real points it is $\int\mathbf x^{-\mathbf w}\,d\big((\boldsymbol\varrho^\rho_B)_\ast\phi^\rho_B\big)$, with the scale form positive and continuous on the open simplex and the seam vector continuous there, which with \hb{thm:closecone} gives $g(\boldsymbol\sigma;0)=\sup\mathbf V_B^{-\bar{\boldsymbol\sigma}}$. The first two are (ii), and the third holds for every piece (\hb{def:seamlaw}, \hb{lem:seamword}). The rest of that proof (the Gamma bounds, Stirling's formula with $\sum_l\sigma_l=0$, the crossings and the null set) does not use the tables.

\emph{(iv).} The partner is a piece on $R$ read with the reversed excesses. The pins of its seam, the letter totals and its scale pin $\sum_{j,D}\big((c_j-i(D))a_{jD}+i(D)b_{jD}\big)=\kappa_R$, are on the balance wall $\kappa_R=\kappa^{\mathrm{top}}_B$ those of the piece (\hb{thm:gluewidth}(iii)), and its seam function converges at $\mathbf c$; so (ii) applies to the piece and to the partner. On the share side of \hb{thm:gluewidth}(iii) the integrand at $\mathbf c+\tau\boldsymbol\sigma+i\boldsymbol\vartheta$ is $\prod_{j:\,s_j>0}\Gst{|\mathcal D_j|}1{\mathbf a_j}{s_j}\prod_{j:\,t_j>0}\Gst{|\mathcal D_j|}1{\mathbf b_j}{t_j}$ times the two seam functions. For $M_B>g_B(\boldsymbol\sigma;\theta_B)$ and $M_R>g_R(\boldsymbol\sigma;\theta_R)$ the definitions of the growths give $C$, $n_1$ and $n_2$ with
\[
\big|P_{W_B}\,P_{W_R}\big|\le C^2\delta^{-2n_2}(1+\tau+|\boldsymbol\vartheta|)^{2n_1}(M_BM_R)^\tau e^{(\theta_B+\theta_R)|\boldsymbol\vartheta|_1}
\]
for $\tau\ge0$, $\delta$ the smaller of $1$ and the distance from the polar hyperplanes of both seam functions. From here the proof of (iii) applies with $\theta=\theta_B+\theta_R<\pi/8$ and $M_BM_R$ in place of $M$, the coefficients sitting inside the seam functions:
\begin{itemize}
\item the contour moves along times at distance at least $\tau^{-3}$ from the crossings, where $\delta^{-2n_2}$ is a power of $\tau$, and the power of $1+|\boldsymbol\vartheta|$ is absorbed into $e^{\eta|\boldsymbol\vartheta|_1}$ with $\eta$ small;
\item a share whose real part $x$ grows has $|\Gamma(x+iy)|e^{\theta|y|}\le\Gamma(x)e^{2\theta^2x/\pi}$, and a share that falls or stays absorbs $e^{\theta|y|}$ against $e^{-\pi|y|/2}$;
\item integrating over the directions of the shares, Stirling's formula and $\sum_l\sigma_l=0$, by the letter totals, give the factor $e^{\tau(h_{\theta_B,\theta_R}(\boldsymbol\sigma)+\epsilon)}$ as $M_B\downarrow g_B$ and $M_R\downarrow g_R$.
\end{itemize}
The ray crosses the hyperplanes of the falling shares and the hyperplanes on which a falling wall of either word lies in $-\mathbb N$, and each crossing adds the fibre integral of the residue of the whole integrand (\hb{lem:pathclose}), a Laurent coefficient where several factors are polar. The pins and the walls of both words have rational linear parts, so crossings on parallel hyperplanes coincide or keep a fixed positive distance. Crossings on hyperplanes that are not parallel, closer than $T^{-3}$ near $\tau=T$, occur only for $\mathbf c$ in a set of measure $O(T^{-3})$, summable over $T$ (Borel and Cantelli). The local decay (H1) of \hb{lem:pathclose} holds because both seam functions grow at most like $e^{\theta'|\im\mathbf z|_1}$ for every $\theta'>0$, by (ii), against the decay of the Gammas at bounded real parts. (H2) holds at the times above when $h_{\theta_B,\theta_R}(\boldsymbol\sigma)<0$, and \hb{lem:pathclose} gives the closing with the stated bounds.

For the last sentence of (iv), let no wall of the word of a side fall along $\boldsymbol\sigma$. Both words have Newton sums of full dimension, their seam functions converging at $\mathbf c$ (\hb{thm:poswordconv}), so the largest value of $\mathbf x^{-\boldsymbol\sigma}$ on the support of the law of the side is $\sup\mathbf V^{-\boldsymbol\sigma}$ of its letter vector (\hb{thm:closecone}, for the partner by its sentence on the partner). At $\theta=0$ the proof of (iii), run on that side, gives the growth $\sup\mathbf V^{-\boldsymbol\sigma}$, finite, and the growth does not increase with the angle, a bound with a smaller angle being a bound with a larger one. On the real ray, where $\boldsymbol\vartheta=0$, the walls stay above their positive values at $\mathbf c$, so $\delta$ is bounded below there, and the seam function at $\mathbf c+\tau\boldsymbol\sigma$ is at least a positive constant times $e^{\tau h'}$ for every $h'<\log\sup\mathbf V^{-\boldsymbol\sigma}$ (the proof of (iii)); so every admissible $M$, at every angle, is at least $\sup\mathbf V^{-\boldsymbol\sigma}$.
\end{proof}

\begin{hdefinition}[\hb{def:chamber} (Chambers of a closing)]
As stated, for every fibre integral with finitely many affine polar functions.
\end{hdefinition}

\begin{proof}
The chamber is convex. If $\boldsymbol\sigma',\boldsymbol\sigma''\in C(\boldsymbol\sigma)$ and $0\le s\le1$, then $\partial_{s\boldsymbol\sigma'+(1-s)\boldsymbol\sigma''}\ell=s\,\partial_{\boldsymbol\sigma'}\ell+(1-s)\,\partial_{\boldsymbol\sigma''}\ell$ equals $\partial_{\boldsymbol\sigma}\ell$ for every polar function that $\boldsymbol\sigma$ falls and is nonnegative for every other. The ray of $s\boldsymbol\sigma'+(1-s)\boldsymbol\sigma''$ consists of the points $s(\mathbf c+\tau\boldsymbol\sigma')+(1-s)(\mathbf c+\tau\boldsymbol\sigma'')$ of the convex $\Omega$.

The instances:
\begin{itemize}
\item In \hb{thm:closecone} the integrand is the Gammas times the reduced periods of the two words. These are holomorphic where they converge absolutely, on a set convex by H\"older's inequality (the reference, def:chamber). On the pins the integrand of a reduced period at $\mathbf c'+\tau\boldsymbol\sigma$ is its integrand at $\mathbf c'$ times $\mathbf V^{-\tau\boldsymbol\sigma}$ read at the point of the orthant (the reference, proof of thm:closecone), at most $(\sup\mathbf V^{-\boldsymbol\sigma})^\tau$; so when both suprema are finite a ball about $\mathbf c$ inside the set of convergence, translated by $\tau\boldsymbol\sigma$, stays inside it, and the ray lies in $\Omega$.
\item In \hb{thm:pastwalls}(iii) the seam function continues to all complex shares, with its poles where a wall lies in $-\mathbb N$ (\hb{thm:pastwalls}(ii)), the walls read through the sums of the shares being affine in the shares.
\item For a fibre integral of Gammas and coefficients the integrand $\prod_l\Gamma(a_l)\xi_l^{-a_l}$, up to a constant, is meromorphic on all of $A+iV$ with its poles where a share lies in $-\mathbb N$. So the polar functions are the shares, $\Omega=A$, and $\partial_{\boldsymbol\sigma'}a_l=\sigma'_l$, which gives the chamber displayed. When the pins fix the total of the shares of every letter, every $\boldsymbol\sigma'\in V$ has $\sum_{l\in K_j}\sigma'_l=0$ and $\sum_{l\in\Psi_j}\sigma'_l=0$, so on the chamber $0\le\sigma'_l\le\sum|\sigma_{l'}|$, the sum over the falling shares $l'$ of the letter of $l$: the chamber is a compact convex polytope.
\item A source factor that is a cube integral continued by $\Reg$ is meromorphic on all of $A+iV$, with its poles where one of its walls lies in $-\mathbb N$ (\hb{lem:cube}).
\end{itemize}
\end{proof}

\begin{htheorem}[\hb{thm:chamber} (A closing depends only on its chamber)]
As stated, for every fibre integral of \hb{def:chamber} with the local decay (H1) of \hb{lem:pathclose}.
\end{htheorem}

\begin{proof}
A polar function $\ell$ that $\boldsymbol\sigma$ falls has $\ell(\mathbf c+\tau\boldsymbol\sigma)=\ell(\mathbf c)+\tau\partial_{\boldsymbol\sigma}\ell=\ell(\mathbf c+\tau\boldsymbol\sigma')$. So its hyperplanes $\ell=-n$ are met by the two rays at the same times, and $\ell$ is constant on the segment joining their points at a time $\tau$. A polar function that $\boldsymbol\sigma$ does not fall is not fallen by $\boldsymbol\sigma'$ either; it is at least $\ell(\mathbf c)>0$ at both points, hence on the segment, and does not lie in $-\mathbb N$ there.

So the two rays meet the same polar hyperplanes at the same times. At a crossing time $\tau^\ast$ the polar hyperplanes through $\mathbf c+\tau^\ast\boldsymbol\sigma$ and through $\mathbf c+\tau^\ast\boldsymbol\sigma'$ are the same, with $\partial_{\boldsymbol\sigma}\ell_H=\partial_{\boldsymbol\sigma'}\ell_H$, so (H0) holds along one ray exactly when it holds along the other.

At a time $\tau$ that is no crossing, the segment avoids every polar hyperplane at a positive distance and lies in the convex $\Omega$. By (H1) on a compact neighbourhood of the segment in $\Omega$, step (1) of the proof of \hb{lem:pathclose} (Cauchy's theorem on rectangles whose horizontal sides vanish by (H1), then Fubini's theorem) equates the fibre integrals through its two ends. A crossing term is the difference of the fibre integrals just before and just after the crossing (step (2) of that proof), so the crossing terms of the two rays at a common crossing are equal. Equal terms and equal remainders carry convergence and bounds from one closing to the other.

The rest is the reference's (the reference, proof of thm:chamber). The integrands of \hb{thm:closecone} and of \hb{thm:pastwalls}(iii) have the local decay (H1), by Stirling's formula and the bounds of their proofs, so the first part applies to them. For every $\epsilon>0$ some $\boldsymbol\sigma'\in C(\boldsymbol\sigma)$, with some $\theta$ in the case of \hb{thm:pastwalls}(iii), has a rate below $h'+\epsilon/2$. The null set of \hb{thm:closecone} is the same for every direction of the chamber, its crossings being those of the falling polar functions, and the bounds of \hb{thm:closecone} or of \hb{thm:pastwalls}(iii) along $\boldsymbol\sigma'$ bound the terms and the remainders along $\boldsymbol\sigma$. On the chamber the falling shares are fixed, $x\log x$ is convex on $[0,\infty)$, and $\log\sup\mathbf V_B^{-\boldsymbol\sigma'}$ and $\log\sup\mathbf V_R^{-\boldsymbol\sigma'}$ are suprema of linear functions of $\boldsymbol\sigma'$, so the rate $h$ is convex there.
\end{proof}

\begin{htheorem}[\hb{thm:averaging} (Averaging the laws)]
As stated, for every share side whose integrand is the Gammas of the shares times a Laplace transform $\int e^{-\ip{\mathbf a}{y}}\,d\lambda$ of a source measure, the kinds (A) and (B) among them.
\end{htheorem}

\begin{proof}
\emph{(i).} At a source point $\zeta$ the kernel is a fibre integral of Gammas and the positive coefficients $e^{y_l(\zeta)}$: the share side of a Gramma (\hb{prop:shareform}) divided by the constant $\prod_j\Gamma(s_j)\prod_j\Gamma(t_j)$. At the real point $\mathbf c$ the Gammas are integrable over the fibre, $|e^{-\ip{\mathbf c+i\boldsymbol\vartheta}{y}}|=e^{-\ip{\mathbf c}y}$, and $\int e^{-\ip{\mathbf c}y}\,d\lambda<\infty$. So the fibre and the source integrals converge absolutely together, and by Fubini's theorem the fibre integral of the share side is the integral of the kernel against $\lambda$. In (A) and (B) that fibre integral is $P_{B\sqcup R}(\mvec)$ (\hb{thm:gluewidth}(iii), \hb{thm:pastwalls}(i)). In (A) the kernel is $\Gamma(\kappa_B)\sZ_{(\kappa_B,\mathbf s,\mathbf t);t_{\mathcal D}}(\mathbf x\circ\mathbf x')$ (the reference, proof of thm:pool(i)): each active letter of the stand-in pair of \hb{thm:gluewidth}(ii) is opened by \hb{thm:E1} with one share per station, and the integral over the scale is the pin of the scale.

The bound on the crossing terms is the reference's (the reference, proof of thm:pool(i)). It uses three things:
\begin{itemize}
\item \hb{thm:chamber} for the kernel, whose polar functions are the shares, with $\Omega=A$ and (H1) holding by Stirling's formula, so that the closings along $\boldsymbol\sigma$ and along every $\boldsymbol\sigma'\in C$ are one series;
\item Stirling's formula for the Gammas, uniformly on the compact $C$, the terms $\tau x\log\tau$ cancelling over each letter by its total;
\item at a crossing, the expansion of $e^{-\ip{\mathbf a}y}=e^{-\ip{\mathbf c+\tau\boldsymbol\sigma'}y}e^{-i\ip{\boldsymbol\vartheta'}y}e^{-w\ip{\mathbf e}y}$ in the transverse variable $w$ to the order $p-1$, with $|\ip{\mathbf e}y|\le C|y|_V$, against the Laurent coefficients of the Gammas, which Cauchy's estimate bounds off the null set.
\end{itemize}
None of these steps uses the kind of the share side.

\emph{(ii).} The source points with $y(\zeta)\in S_\lambda$ carry all of $\lambda$, and there $h_C(y)\le\hat h_C$. By (i) the averaged terms are absolutely convergent integrals, at most $C_\epsilon e^{\tau_k(\hat h_C+\epsilon)}\int w\,d(y_\ast\lambda)$, which is finite by the certificate of \hb{thm:averaging}. At each source point the kernel is the sum of its crossing terms before a time $T$ separating the crossings and of the integral left at $T$ (\hb{lem:pathclose}, whose identity at a finite time needs no limit), and that integral is at most $C_\epsilon e^{-\ip{\mathbf c}y}e^{T(\hat h_C+\epsilon)}$. Integrating against $\lambda$, the fibre integral of the share side is the sum of the averaged terms before $T$ up to $C_\epsilon e^{T(\hat h_C+\epsilon)}\int e^{-\ip{\mathbf c}y}\,d\lambda$; the crossings up to $T$ number $O(T)$, and with $\hat h_C+\epsilon<0$ both the sum and the bound follow. At a crossing point, along some $\boldsymbol\sigma'\in C$, inside the set where $\int e^{-\ip{\mathbf a}y}\,d\lambda$ converges absolutely, with its derivatives in the shares, Fubini's theorem (the Gammas decaying by Stirling's formula) exchanges the fibre and the source integrals in the crossing term of the share side. The inner integral is then the crossing term of the kernel at that point, which is $\mathcal R_k(\zeta)$ by \hb{thm:chamber}.

\emph{(iii).} The function $h(\boldsymbol\sigma';y)$ is continuous, convex in $\boldsymbol\sigma'$ on $C$ (the falling coordinates being fixed and $x\log x$ convex on $[0,\infty)$), and affine in $y$; $C$ is compact and convex. Sion's minimax theorem on $C\times\operatorname{conv}S_\lambda$ gives the identity, the supremum of an affine function over $\operatorname{conv}S_\lambda$ being its supremum over $S_\lambda$.

For the reading through the log-support,
\[
\Big|\int e^{-\ip{\mathbf c+\tau\boldsymbol\sigma'+i\boldsymbol\vartheta}{y}}\,d\lambda\Big|\le e^{\tau\sup_{y\in S_\lambda}(-\ip{\boldsymbol\sigma'}y)}\int e^{-\ip{\mathbf c}y}\,d\lambda .
\]
With the Gamma bounds of the proof of \hb{thm:closecone}, this bounds the closing of the share side along $\boldsymbol\sigma'$ at the rate $\sum_l\sigma'_l\log|\sigma'_l|+\sup_{S_\lambda}(-\ip{\boldsymbol\sigma'}y)=\sup_{y\in S_\lambda}h(\boldsymbol\sigma';y)$. In (A) and (B) this is the rate $h(\boldsymbol\sigma')$ of \hb{thm:closecone} (the reference, proof of thm:pool(iii)):
\begin{itemize}
\item the seam words of the sides that are laws have Newton sums of full dimension, their seam functions converging at $\mathbf c$ (\hb{thm:poswordconv});
\item so the largest values of the monomials of their letter vectors are read on the supports of their laws (\hb{thm:closecone}; for the partner in (A), its sentence on the partner);
\item in (B), $\sup\mathbf V_R^{-\boldsymbol\sigma'}=\boldsymbol\xi^{-\boldsymbol\sigma'}$;
\item the largest value of $-\ip{\boldsymbol\sigma'}{y}$ on the log-support of a product of laws is the sum of its largest values on the log-supports of the factors.
\end{itemize}
In a bundle $G$, Jensen's inequality for the logarithm, in the proportions $\sigma'_l/g$, gives $\sum_G\sigma'_l\log(\sigma'_le^{-y_l})\ge-g\log\big(g^{-1}\sum_Ge^{y_l}\big)$, with equality exactly when $e^{y_l}/\sigma'_l$ is constant on $G$.
\end{proof}

\begin{hdefinition}[\hb{def:dirkeyhole} (Keyhole continuations along a direction; ceiling and phase)]
As stated, for every piece of \hb{thm:pastwalls} and the sectors of any fan of its seam word.
\end{hdefinition}

\begin{proof}
The displayed sum is \hb{thm:pastwalls}(ii) at the sums of the shares over the monomials of the partner. A set $Y$ of a box is compact and is the product of a box, of intervals, and of keyholes, a closed disc about $0$ joined to a segment of the positive axis (\hb{def:keyhole}); each factor is star-shaped, about the centre of its disc for a keyhole, so $Y$ is star-shaped, and so is $Y$ plus a small open polydisc about $0$. The sector letters do not vanish on $Y$, by the definition of a keyhole continuation, hence on such a neighbourhood of $Y$, which is simply connected, so they have logarithms there. The real points of $Y$ form a product of intervals, connected, containing points with nonnegative coordinates, where the letters, with their nonnegative coefficients and positive constant terms, are positive; real and without zeros on that set, they are positive on it, and the logarithms can be taken real at the real points. For a piece of tables and the sectors of the orders, the walls of the generators are the proper walls of the tails and the sector letters are those of \hb{thm:exist}(ii) (\hb{thm:pastwalls}(ii)).
\end{proof}

\begin{htheorem}[\hb{thm:keyholebound} (The travelling bound past the walls)]
As stated, for every piece of \hb{thm:pastwalls}.
\end{htheorem}

\begin{proof}
This is the reference's proof (the reference, proof of thm:keyholebound), with the sectors of the fan in place of the orders of $B$ and the generators in place of the tails.

(i) Fix a sector $\pi$, a box $\Xi$ of it, the set $J$ of its falling variables with keyholes and their radii $r_k$, and choose $r_k^-\in(0,r_k)$ with $\sum_{k\in J}(-\dot\kappa_k)\log(r_k/r_k^-)<\epsilon/2$. The term of $\pi$ is $|\det(\mathbf v_{\pi,1},\dots,\mathbf v_{\pi,|B|-1})|\int\prod_{k\in J}q_k^{\kappa_k(\bar{\mathbf z})-1}\Phi\,dq$, with $\Phi=\prod_{k\notin J}q_k^{\kappa_k(\bar{\mathbf z})-1}\mathbf k_\pi(q)^{-\bar{\mathbf z}}$ holomorphic in the variables of $J$ on a neighbourhood of their keyholes, because the sector letters do not vanish on the sets $Y$. The following steps apply as in the reference:
\begin{itemize}
\item \hb{lem:keyhole}, applied in the variables of $J$ one at a time with $r_k^-$ and $r_k$;
\item the maximum modulus principle in the variables of $J\setminus A$ that lie in their discs;
\item the exact integration of the falling monomials, which carries no phase;
\item the identity $|k^{-z}|=|k|^{-\re z}e^{\im z\arg k}$ at the points of the outline;
\item the bounds of the letters above and below on the compact outline;
\item the integrability of $q_k^{\kappa_k(\bar{\mathbf c})-1}$ over the base variables, the walls being positive at $\bar{\mathbf c}$;
\item $|\kappa_k(\bar{\mathbf z})+l_k|\ge c\,\delta$ for at most $|B|-1$ factors.
\end{itemize}
Each of these uses only the sector form of \hb{thm:pastwalls}(ii): walls affine in the shares and positive at $\bar{\mathbf c}$, and sector letters polynomial, positive on the closed cube and nonvanishing on the sets $Y$. The determinant is a constant, and the boxes, the sets $A$ and the sectors are finitely many.

(ii) The integrand of the share side of \hb{thm:pastwalls}(i) at $\mathbf c+\tau\boldsymbol\sigma+i\boldsymbol\vartheta$ is the Gammas and the coefficients times the seam function, which (i) bounds. The steps are those of the reference:
\begin{itemize}
\item for a growing share, the exact bound $|\Gamma(x+iy)|e^{\varphi y}\le\Gamma(x)(\cos\varphi)^{-x}$ for $|\varphi|<\frac\pi2$, from $\Gamma(z)=e^{i\varphi z}\int_0^\infty u^{z-1}e^{-ue^{i\varphi}}du$, applied with $\varphi=\upsilon_l(q)+\eta\operatorname{sign}y$;
\item falling and staying shares absorb $e^{\upsilon_l(q)y}$ against $e^{-\pi|y|/2}$, as $|\upsilon_l|\le\frac\pi2-\epsilon_0$;
\item Stirling's formula with $\sum_l\sigma_l=0$, the supremum over $q$ being taken after the integration over the directions of the shares.
\end{itemize}
Together these give the travelling bound at the rate $h_{\boldsymbol\upsilon}(\boldsymbol\sigma)$. Then \hb{lem:pathclose} applies: (H0) off the null set; (H1) by \hb{thm:pastwalls}(ii) with any $\theta<\pi/2$ and Stirling's formula; (H2) at times at distance at least $\tau^{-3}$ from the crossings.
\end{proof}

\begin{hcorollary}[\hb{cor:phasefree} (When the phase costs nothing)]
As stated, for every piece of \hb{thm:pastwalls}.
\end{hcorollary}

\begin{proof}
In (i) the price $-\sum_{\sigma_l>0}\sigma_l\log\cos\upsilon_l$ vanishes. In (ii), $\log|\varrho_{\pi,\bar{\boldsymbol\sigma}}|-\sum_{\sigma_l>0}\sigma_l\log\cos\upsilon_l\le\log H_{\mathcal O}$ on the outline, with equality approached where $|\varrho_{\pi,\bar{\boldsymbol\sigma}}|$ approaches $H_{\mathcal O}$, the price being nonnegative.

With no falling variables, the outline is the union of the real cubes of the sectors, where the sector letters are positive and the phase vanishes. Differentiating the wall of \hb{thm:pastwalls}(ii) gives $\dot\kappa_k=-\ip{\sum_{j,D}(\bar\sigma_{jD}\mathbf t_{jD,\pi}+\bar\sigma^\ast_{jD}\mathbf t^\ast_{jD,\pi})}{\mathbf v_{\pi,k}}$. So on the section $\varrho_{\pi,\bar{\boldsymbol\sigma}}=\prod_{j,D}\big(\prod_kq_k^{\ip{\mathbf t_{jD,\pi}}{\mathbf v_{\pi,k}}}k_{jD,\pi}\big)^{-\bar\sigma_{jD}}\big(\prod_kq_k^{\ip{\mathbf t^\ast_{jD,\pi}}{\mathbf v_{\pi,k}}}\psi_{jD,\pi}\big)^{-\bar\sigma^\ast_{jD}}=\mathbf V_B^{-\bar{\boldsymbol\sigma}}$, read in the sector coordinates. The cubes of the sectors cover the section, so the ceiling is $\sup\mathbf V_B^{-\bar{\boldsymbol\sigma}}$, the monomial being homogeneous of degree zero, and the rate is that of \hb{thm:closecone} for the glue read in the monomials of the partner.
\end{proof}

\begin{hproposition}[\hb{prop:pathcompose} (How travelling bounds compose)]
As stated, the face coefficient of (iii) being read on the sectors of any fan of the seam word.
\end{hproposition}

\begin{proof}
(i) and (ii) are the reference's (the reference, proof of prop:pathcompose). For (iii), the reference's proof applies with one sentence changed: a wall $\kappa_{\mathbf v}$ is the exponent of $q_k$ in each sector of which $\mathbf v=\mathbf v_{\pi,k}$ is a generator (\hb{thm:pastwalls}(ii)).
\begin{itemize}
\item The pole is simple, as opposed to the multiple poles of the statement, when no other factor and no second wall of a sector is polar on the hyperplane. Then, by \hb{lem:cube}, the residue of the sector integral in the direction $\mathbf e$ at $\kappa_{\mathbf v}=-n$ is $|\det(\mathbf v_{\pi,1},\dots,\mathbf v_{\pi,|B|-1})|$ times the integral over the other coordinates of the $n$-th Taylor coefficient in $q_k$ of the rest of the integrand, divided by $\partial_{\mathbf e}\kappa_{\mathbf v}$. Summing over the sectors gives the face coefficient.
\item The other powers $q_{k'}^{\kappa_{k'}-1}$ do not depend on $q_k$, so the Taylor coefficient is that of the product of the powers of the sector letters: a finite sum of products of derivatives of the letters and powers of the letters at $q_k=0$, with coefficients polynomial in the exponents.
\item Where two walls of one sector lie in $-\mathbb N$ on the hyperplane, the expansion of \hb{lem:cube} in both variables gives a pole of the second order. A Laurent coefficient of a product is a finite sum of products of Laurent coefficients, and the derivatives in the shares of $q^{\kappa(\mathbf z)}$ and of $k^{-z}$ bring $\log q$ and $-\log k$.
\item The wall crossed falls along the direction, so $q_k$ is a falling variable and $q_k=0$ lies in its keyhole, where the letters do not vanish; Cauchy's estimate on $|q_k|=r_k$ bounds the $n$-th Taylor coefficient by $r_k^{-n}$ times the largest modulus on the circle.
\end{itemize}
\end{proof}

\begin{hdefinition}[\hb{def:retainedcohort} (A rational retained cohort)]
As stated: the fibre integrals of Gammas, coefficients and source factors whose pins carry a grading.
\end{hdefinition}

\begin{proof}
The share side of \hb{thm:pastwalls}(i) and the kernel of \hb{thm:averaging} are graded by their letter totals: the sum of the rows of the pins that fix the totals of the letters takes the value $1$ on the column of every share, each share belonging to one letter. By \hb{thm:pastwalls}(ii), the sector integrals of a fan of a seam word are cube integrals $\int_{[0,1]^k}q^{\boldsymbol\kappa(\mathbf a)-\one}\prod_iP_i(q)^{-w_i(\mathbf a)}dq$. Here the $P_i$ are the sector letters, with nonnegative coefficients and positive constant terms, and the $\kappa_i$ and $w_i$ are affine with rational linear parts; for tables these are the sector integrals of \hb{cor:pastwalls}.
\end{proof}

\begin{sremark}
The grading is part of the definition because \hb{thm:retainedchain} extracts the scalar of a joint word through the change of logarithmic scales $z=\rho\mathbf v+E\log X$, which needs a vector $\mathbf v$ taking the value $1$ on every column. Without one the form of that theorem can fail, as follows.

Take the Gramma $\sG_{\nu;(1,2)}(\xi_1,\xi_2)=\int_0^\infty x^{\nu-1}e^{-\xi_1x-\xi_2x^2}dx$, $\nu,\xi_1,\xi_2>0$. Its configuration has no grading, since $\ip{1}{v}=\ip{2}{v}=1$ is impossible. Its share side is pinned by $b_1+2b_2=\nu$, and its closing along $(-2,1)$ crosses $b_1=-N$, $N\in\mathbb N$, with the crossing terms
\[
\frac{(-\xi_1)^N}{N!}\cdot\frac12\Gamma\Big(\frac{\nu+N}2\Big)\xi_2^{-(\nu+N)/2},
\]
the terms of the expansion of the integral in powers of $\xi_1$, using $\int_0^\infty x^{\nu+N-1}e^{-\xi_2x^2}dx=\frac12\Gamma(\frac{\nu+N}2)\xi_2^{-(\nu+N)/2}$.

The scalar $\Gamma(\frac{\nu+N}2)/\Gamma(N+1)$ has an $N$-th root tending to $0$: its logarithm divided by $N$ is $-\frac12\log N+O(1)$. A joint word with the scalar of \hb{thm:retainedchain}, here with $r_G=1$, would in this positive case have the root rate $\operatorname{ess\,sup}c>0$, or vanishing terms (\hb{cor:retainedheight}).
\end{sremark}

\begin{htheorem}[\hb{thm:retainedchain} (A finite procedure for a retained cohort)]
As stated, for every retained cohort, the scalar being written through the grading of its pins.
\end{htheorem}

\begin{proof}
This is the reference's proof (the reference, proof of thm:retainedchain); only the radial step involves the totals of the shares.
\begin{itemize}
\item At a Gamma crossing, the crossed share is fixed at $-N(n)$, and its Gamma is replaced by the residue $(-1)^{N(n)}/\Gamma(N(n)+1)$.
\item The crossed share and its added pin are eliminated together. This leaves the columns $t_l$, $l\neq f$, of the original pins, with the values $\boldsymbol\nu+N(n)t_f$, $t_f$ the column of the crossed share.
\item By \hb{lem:gammastriplift} each share of $\mathcal S_-$ is translated, $b_l=a_l+\delta_l(n)$, with the factor $(-1)^{\delta_l(n)}\xi_l^{\delta_l(n)}/\Gamma(\delta_l(n))$, the power of $\xi_l$ staying in the joint input.
\end{itemize}
The completed retained fibre is then the Gramma $\sG_{\boldsymbol\mu_n;t}$ (\hb{prop:shareform}), with the remaining columns and the index $\boldsymbol\mu_n=\boldsymbol\nu+N(n)t_f+\sum_{l\in\mathcal S_-}\delta_l(n)t_l$. Every remaining column is a column of the original pins, so $\ip{t_l}{\mathbf v}=1$ for the grading $\mathbf v$, and
\[
S_n=\ip{\mathbf v}{\boldsymbol\mu_n}=w+N(n)+\sum_{l\in\mathcal S_-}\delta_l(n).
\]
The change of logarithmic scales $z=\rho\mathbf v+E\log X$, with $E$ a rational matrix whose columns complement $\mathbf v$, gives $\sG_{\boldsymbol\mu_n;t}=J_{\mathbf v}\Gamma(S_n)\int X^{E^{\mathsf T}\boldsymbol\mu_n-\one}\mathcal L(X)^{-S_n}dX$: the integral over $\rho$ is $\Gamma(S_n)\mathcal L^{-S_n}$. Collecting the residue, the factors of the translations and $\Gamma(S_n)$ gives the stated $G_n$. Its slopes in $n$ agree in the numerator and the denominator, $d_0+\sum_{\mathcal S_-}d_l=\bar d$, so the terms $n\log n$ of Stirling's formula cancel, and the rest of Stirling's formula gives $r_G$ and the polynomial bound. On the share sides of \hb{thm:pastwalls} and \hb{thm:averaging}, $\mathbf v$ is the sum of the letter-total rows and $w=|\mathbf s|+|\mathbf t|$. The other steps of the procedure (the source strips, the circles, the joint read, the resolution of the endpoints and the fan) do not involve the grading.
\end{proof}

\begin{htheorem}[\hb{thm:retainedchamberrate} (The exact positive rate of a retained chamber)]
As stated, with a real index.
\end{htheorem}

\begin{proof}
The reference's proof (the reference, proof of thm:retainedchamberrate) uses rationality only of the configuration and of the direction, through the cohort and its integer step.
\begin{itemize}
\item \hb{def:retainedcohort} allows real constant terms and indices. The crossing times of a falling share are $(a^0_l+N)/d_l$, $N\in\mathbb N$, and passing to the integer step makes the changes of all falling shares integral, whatever the constant terms.
\item The positive construction uses \hb{lem:gammastriplift} with $\delta_l(n)=\delta_l(0)+d_ln$ and real starting values.
\item The rate is computed by \hb{lem:positivecapacity}, at the slope $\sum_{f}d_ft_f$ of the index, which does not involve its constant term.
\item The root limit is that of \hb{cor:retainedheight}, whose positive case allows real constant terms.
\end{itemize}
The chamber polytope of the statement is the chamber of \hb{def:chamber} of the share side, whose polar functions are the shares (\hb{def:chamber}), for the direction per step.
\end{proof}

\begin{htheorem}[\hb{thm:averagingsharp} (Averaging is sharp)]
As stated, for every share side of Theorem~\hb{thm:averaging} and every direction with rational coordinates whose crossings are simple and nonresonant.
\end{htheorem}

\begin{proof}
At each source point the kernel of \hb{thm:averaging}(i) is, up to the constant factor $\prod_j\Gamma(s_j)^{-1}\prod_j\Gamma(t_j)^{-1}$, the share side of the Gramma whose configuration is the matrix of the pins, rational by hypothesis, with the grading of the letter totals; whose index is the vector of the values of the pins, real; and whose coefficients are $e^{y_l(\zeta)}$, positive and continuous in $\zeta$ (\hb{prop:shareform}). Let $\boldsymbol\sigma$ be rational and the crossings of its ray simple and nonresonant. A falling share $l$ crosses its poles at the times $(c_l+N)/|\sigma_l|$, $N\in\mathbb N$; for a rational $\Delta>0$ with every $\Delta\sigma_l$ an integer these times fall into finitely many progressions of step $\Delta$, so the crossing terms form finitely many cohorts with the integer step of length $\Delta$ in $\tau$ (\hb{def:retainedcohort}), the direction per step being $\Delta\boldsymbol\sigma$, from the positive point $\mathbf c$.
\begin{itemize}
\item \hb{thm:retainedchamberrate}, with a real index, gives every retained term a sign at each index that does not depend on the coefficients, hence not on $\zeta$, so the averaged term has that sign.
\item \hb{cor:retainedaveragedrate} applies to each cohort: one averaged term is absolutely integrable by \hb{thm:averaging}(i) and its certificate, $\hat h_C<\infty$, and the coefficient map is continuous, so the essential supremum of the chamber rate is its supremum on $S_\lambda$.
\item The chamber of the step direction is $\Delta C$, and since $\sum_l\sigma'_l=0$ on $V$, $h(\Delta\boldsymbol\sigma';y)=\Delta h(\boldsymbol\sigma';y)$. So the root rate per step is the limit $e^{\Delta\hat h_C}$.
\end{itemize}
Hence for every $\epsilon>0$ and every cohort, $|\int\mathcal R_n\,d\lambda|\ge e^{n\Delta(\hat h_C-\epsilon/2)}$ for all large $n$. The crossing time being $\tau_n=\tau_0+n\Delta$, this exceeds $e^{\tau_n(\hat h_C-\epsilon)}$ for all large $n$, the quotient being $e^{n\Delta\epsilon/2-\tau_0(\hat h_C-\epsilon)}$; the cohorts being finitely many, the bound holds for all large $k$. With \hb{thm:averaging}(ii), the terms tend to zero when $\hat h_C<0$ and do not when $\hat h_C>0$.

Now let every falling share be alone in its bundle and the merged fibre have dimension one.
\begin{itemize}
\item The merged pins are rational, so the merged directions form a line spanned by a rational vector $\mathbf r$. The merged direction of $\boldsymbol\sigma$, nonzero since $\boldsymbol\sigma$ falls a share, which is alone in its bundle, is $\mu\mathbf r$ with $\mu\neq0$, and $\mu>0$ after the choice of the sign of $\mathbf r$. The falling coordinates of the positive multiple $\mu^{-1}\boldsymbol\sigma$ of $\boldsymbol\sigma$ are coordinates of $\mathbf r$, rational.
\item The chamber of $\mu^{-1}\boldsymbol\sigma$ is a nonempty polytope defined by rational equations and inequalities, so its vertices have rational coordinates. By \hb{thm:chamber} the closing along $\mu^{-1}\boldsymbol\sigma$, which is the closing along $\boldsymbol\sigma$ after a change of scale of $\tau$, is the closing along that rational direction.
\item At a crossing time the growing shares are positive at the real point of the ray, so the crossing is simple and nonresonant exactly when no other falling share lies in $-\mathbb N$ there, that is, when no two lattices of the falling shares meet.
\end{itemize}
The form of the crossing terms is the reference's (the reference, proof of thm:pool(iv)). The fibre is a bundle over the merged line whose fibres are the splits of the bundles at fixed sums, and \hb{thm:E1} completes each split, $\int_{\sum_{l\in G}a_l=\bar a}\prod_{l\in G}\Gamma(a_l)e^{-a_ly_l}=\Gamma(\bar a)e^{-\bar ay_G}$. A crossing term on the merged line is a fibre integral of dimension zero: the residue of the Gamma crossed, times the other Gammas at the crossing point and $\prod_Ge^{-\bar a_{k,G}y_G}$, divided by the Gammas of the exponents and by the determinant of the pins.
\end{proof}

\section{Chapter~\hb{ch:lawreach}: laws at their ends: regions of equality and reach}\label{sec:ch30}

\begin{htheorem}[\hb{thm:dirreach} (The reach along a direction is read on the law)]
Parts (i) to (iv) as stated, with part (v).
\end{htheorem}

\begin{proof}
Parts (i) to (iv) are the reference's (the reference, proof of thm:dirreach); the piece of tables named in (iii) has a seam word with a Newton sum of full dimension when its tables form a connected family, by \hb{thm:pastwalls}(ii). For (v) fix $\theta\ge0$ and let $g_\theta$ be the infimum of (v).

First, $r\le g_\theta$. On the half-strip $\Sigma$ of \hb{def:dirreach} the imaginary parts are bounded, say by $Y$, so $e^{\theta|\im z|}\le e^{\theta Y}$ there. A bound $|\ell(z)|\le C(1+|z|)^{n'}M^{\re z}e^{\theta|\im z|}$ on the half-plane, off the $\delta$-neighbourhoods of the candidate poles, therefore gives on $\Sigma$ the bound of \hb{def:dirreach} with $Ce^{\theta Y}$, and $r\le M$.

Second, $g_\theta\le r$. By (i), on $(r,\infty)$ the measure $\lambda$ has a density that is a wall series of radius $r$. By (ii) with $x_0=r$, for every $x'>r$ there is $C$ with $|\ell(z)|\le C\delta^{-J-1}x'^{\,\re z}$ at every $z$ with $\re z\ge0$ at distance at least $\delta\le1$ from the poles, whatever $\im z$. The poles with $\re z\ge0$ lie among the candidate poles, and those with $\re z<0$ keep a distance at least some $\epsilon_0>0$ from the half-plane, $\ell$ converging absolutely on a strip about the imaginary axis (\hb{def:dirreach}). So at every $z$ with $\re z\ge0$ at distance at least $\delta$ from the candidate poles, with $\delta_1=\min(\delta,1,\epsilon_0)$, $|\ell(z)|\le C\delta_1^{-J-1}x'^{\,\re z}$: the bound of (v) with the constant $C\delta_1^{-J-1}$ and $n'=0$. So $M=x'$ is admissible in (v) at $\theta=0$, hence at every $\theta\ge0$, the factor $e^{\theta|\im z|}$ being at least $1$, and $g_\theta\le x'$ for every $x'>r$.
\end{proof}

\begin{htheorem}[\hb{thm:dirheights} (Ceilings, simple tops and fibres along a direction)]
As stated, for every piece whose seam function converges at a real point of the pins with positive shares, with the sectors of any fan of its seam word.
\end{htheorem}

\begin{proof}
By \hb{thm:pastwalls}(ii) the piece has walls and growth bounds (\hb{def:wallsgrowth}), its walls being those of the generators of the fan, so \hb{def:dirreach} applies with these walls. The proof is the reference's (the reference, proof of thm:dirheights), with the sectors of the fan in place of the orders of $B$, the generators in place of the tails, and the constant $|\det(\mathbf v_{\pi,1},\dots,\mathbf v_{\pi,|B|-1})|$ in the sector forms. At the point $\mathbf w_0+i\bar{\boldsymbol\vartheta}+z\bar{\boldsymbol\sigma}$ the integrand of the sector integral of $\pi$ is
\[
|\det(\mathbf v_{\pi,1},\dots,\mathbf v_{\pi,|B|-1})|\prod_kq_k^{\kappa_k(\mathbf w_0)+i\partial_{\bar{\boldsymbol\vartheta}}\kappa_k-1}\,\mathbf k_\pi^{-\mathbf w_0-i\bar{\boldsymbol\vartheta}}\,\varrho_{\pi,\bar{\boldsymbol\sigma}}^{\,z}.
\]
(i) The proof of \hb{thm:keyholebound}(i) concerns the seam function alone. It holds at the points $\mathbf w_0+\re z\,\bar{\boldsymbol\sigma}+i(\bar{\boldsymbol\vartheta}+\im z\,\bar{\boldsymbol\sigma})$, with $\delta$ the distance of the complex point from the hyperplanes on which a wall lies in $-\mathbb N$, which is at least a constant times the distance of $z$ from the candidate poles (the reference, proof of cor:dirtube). On $\Sigma$ the phase is bounded, the imaginary parts being bounded there. So $|\ell(z)|\le C\delta^{-(|B|-1)}e^{\re z(\log H_{\mathcal O}+\epsilon)}$, and $r\le H_{\mathcal O}e^\epsilon$ for every $\epsilon>0$.

(ii) and (iii) follow from the proof of \hb{thm:simpletop}(i) and (iii), with the point $\mathbf w_0+m\bar{\boldsymbol\sigma}$ in place of $b^d=-k$, $\varrho_{\pi,\bar{\boldsymbol\sigma}}$ in place of $\varrho_\pi$ and $\phi_{\pi,\mathbf w_0}$ in place of $\phi_\pi$.
\begin{itemize}
\item The falling variables carry $q_k^{\kappa_k(\mathbf w_0)-m|\dot\kappa_k|-1}$ with $|\dot\kappa_k|$ an integer, so the factor $(e^{2\pi i\kappa_k(\mathbf w_0)}-1)^{-1}$ of their circles is the same at every $m$. The letters whose exponents are nonnegative integers at every large $m$ carry polynomials and need not be kept from vanishing.
\item Every wall is the wall of a generator of some sector, so its falling rate is one of the integers $|\dot\kappa_k|$. The points $m$ lie at a distance bounded below from the candidate poles, the falling walls being off the integers at $\mathbf w_0$, so $|\ell(m)|\le C(1+m)^Dr'^{\,m}$ for every $r'>r$ by \hb{def:dirreach}, and $\operatorname{growth}_{\mathbb Z}\le r$.
\item Under (iii), $\operatorname{growth}_{\mathbb Z}=H_{\mathcal O}$ as in \hb{thm:simpletop}(iii), and a keyhole continuation of ceiling $H_{\mathcal O}$ gives $r\le H_{\mathcal O}$ by (i).
\end{itemize}
(iv) follows from the proofs of \hb{thm:fibrebound} and \hb{thm:mixedbound}, applied to the positive directional law, with \hb{thm:dirreach}(i) and (ii) in place of \hb{thm:lawreach}.
\begin{itemize}
\item A keyhole box is handled as in (i).
\item In a fibre box, $\varrho_{\pi,\bar{\boldsymbol\sigma}}$ is meromorphic in each falling variable on its disc, the exponents $\bar\sigma_L$ and $\dot\kappa_j$ being integers. At a point above the ceiling and the stationary bound, the map is a submersion on its stratum: by \hb{lem:degreeone} when one coordinate of falling rate one lies inside its disc, and by the definition of the stationary bound otherwise.
\item The Mellin transform of a fibre law is a cube integral as in \hb{lem:cubegrowth}, with its poles in finitely many cosets of $(\operatorname{lcm}_j|\dot\kappa_j|)^{-1}\mathbb Z$, so \hb{lem:wallcontinue} applies.
\end{itemize}
None of these steps uses more of a sector integral than the form displayed above, which \hb{thm:pastwalls}(ii) supplies for every fan of the seam word of every such piece.
\end{proof}

\begin{htheorem}[\hb{thm:dirclosing} (The closing past the walls is read along its direction)]
As stated, with the last sentence of (i), for the share side of \hb{thm:pastwalls}(iv).
\end{htheorem}

\begin{proof}
Apart from the last sentence of (i), the statement is the reference's (the reference, proof of thm:dirclosing). For that sentence, the piece and the partner have walls and growth bounds (\hb{thm:pastwalls}(ii), applied to each), and $\boldsymbol\sigma$ has rational coordinates, so \hb{cor:dirtube} applies to each of the two seam functions at $\mathbf c$, over the whole space $V$ of the directions of the shares. There the growth $g_X(\boldsymbol\sigma;\theta)$ of \hb{thm:pastwalls}(iv), $X\in\{B,R\}$, is the function $g_V(\boldsymbol\sigma;\theta)$ of \hb{cor:dirtube} for that seam function: the two definitions coincide, the shares carrying no sums and the walls of \hb{thm:pastwalls}(ii) serving as the walls of (W1), so that $\delta$ is measured from the same hyperplanes, those of the polar functions of the seam function (\hb{def:chamber}). It does not increase with $\theta$ and tends to $r^\sharp_X$ as $\theta\downarrow0$, so $g_X(\boldsymbol\sigma;\theta)\le r^\sharp_X$ for every $\theta>0$. Then
\[
h_{\theta/2,\theta/2}(\boldsymbol\sigma)\le\sum_l\sigma_l\log|\sigma_l|+\log r^\sharp_B+\log r^\sharp_R+\frac{2\theta^2}\pi\sum_{l:\,\sigma_l>0}\sigma_l,
\]
which is negative for $\theta$ small when $\sum_l\sigma_l\log|\sigma_l|+\log r^\sharp_B+\log r^\sharp_R<0$. \hb{thm:pastwalls}(iv) gives the closing, its terms and remainders being $O(e^{\tau(h_{\theta/2,\theta/2}(\boldsymbol\sigma)+\epsilon)})$, and letting $\theta\to0$ gives the stated bound.
\end{proof}

\begin{hcorollary}[\hb{cor:anyfibre} (A law with log expansions bounds a closing of any fibre dimension)]
As stated, with its last sentence, for the share side of \hb{thm:pastwalls}(iv).
\end{hcorollary}

\begin{proof}
Apart from the last sentence, the statement is the reference's (the reference, proof of cor:anyfibre). For that sentence, apply \hb{prop:logexpansions} to the law of the piece and to that of the partner, a piece with the reversed excesses. On the contour through $\mathbf c+\tau\boldsymbol\sigma'$ each seam function is at most $C\delta^{-n}X^X_{\mathrm{in}}(\boldsymbol\sigma')^\tau$, with no factor in the imaginary parts and no power of $\tau$, so $g_X(\boldsymbol\sigma';0)\le X^X_{\mathrm{in}}(\boldsymbol\sigma')$ for every direction $\boldsymbol\sigma'$. Hence
\[
h_{0,0}(\boldsymbol\sigma')\le\sum_l\sigma'_l\log|\sigma'_l|+\log X^B_{\mathrm{in}}(\boldsymbol\sigma')+\log X^R_{\mathrm{in}}(\boldsymbol\sigma'),\qquad\boldsymbol\sigma'\in C(\boldsymbol\sigma).
\]
The polar functions of the share side of \hb{thm:pastwalls}(iv) are the shares and the walls of both words, with $\Omega=A$, and (H1) holds by \hb{thm:pastwalls}(ii) and Stirling's formula. So \hb{thm:chamber} makes the closings along the directions of the chamber one series; their crossings and the near coincidences of these are those of the falling polar functions, the same for every direction of the chamber, so $\mathbf c$ lies off the null set of \hb{thm:pastwalls}(iv) for each of them; and \hb{thm:pastwalls}(iv) at $\theta_B=\theta_R=0$, along a direction of the chamber whose bound is within $\epsilon$ of the infimum, gives the closing along $\boldsymbol\sigma$ with the stated bound. On the chamber the falling coordinates are fixed and the others are nonnegative, where $x\log x$ is convex, and the logarithms of the inner bounds are convex (\hb{def:logexpansions}).
\end{proof}


\begin{sremark}[The directional ratio at width one, in the paragraph after \hb{def:dirreach}]
At width one with one table, factoring weights normalized as in \hb{prop:widthonelaw}, and $s,t>0$, the directional ratio $x^\ast_{D_0}/x_{D_1}$ along $\bar{\boldsymbol\sigma}=(-1,1,-1,1)$ is the ratio $\varrho_B$ at the seam points, and every coordinate of the seam torus there is a power of $\varrho_B$.
\end{sremark}

\begin{proof}
On the seam torus $x_{D_0}=x^\ast_{D_1}=1$, so $\mathbf x^{-\bar{\boldsymbol\sigma}}=x_{D_0}\,x_{D_1}^{-1}\,x^\ast_{D_0}\,x^{\ast\,-1}_{D_1}=x^\ast_{D_0}/x_{D_1}$. At the seam point of $u$, \hb{def:seamlaw} gives $\varrho^{\rho\ast}_{D_0}/\varrho^\rho_{D_1}=V^\ast_{D_0}V_{D_0}/(V^\ast_{D_1}V_{D_1})$, the scale cancelling. By Lemma~\ref{lem:ch27-normal}(i) this is
\[
\Big(\frac{\hat\theta_0}{\hat\theta_1}\Big)^2\,\frac{\Psi_{\cB/R}\,K_{\cB/R}}{\Psi_{\cB|B}\,K_{\cB|B}}=\Big(\frac{\hat\theta_0}{\hat\theta_1}\Big)^2\varrho_B,
\]
with $\varrho_B=\gamma^\Psi_B/\gamma^K_B$ of \hb{def:readings}. Under the normalization $\hat\theta_0=\hat\theta_1=1$, so it is $\varrho_B$. By \hb{prop:widthonelaw}(ii) with $N=1$, $\mathbf x=(1,\varrho_B)$ and $\psi=\varrho_B^{\,t/(s+t)}$, the coordinates of the seam point are $1$, $\varrho_B^{-t/(s+t)}$, $\varrho_B^{\,s/(s+t)}$ and $1$, all powers of $\varrho_B$.
\end{proof}

% Proofs of Chapter ch:motive. No new macros.
\section{Chapter~\hb{ch:motive}: the motive of a family and the coaction along a tree}\label{sec:chmotive}

Five outside works enter these proofs, each named where it is used together with the property of the setting it needs: the theory of toric varieties (W.~Fulton, \emph{Introduction to toric varieties}, Princeton University Press, 1993, Chapters~1 to~3), for the braid fan; F.~Brown, \emph{Notes on motivic periods}, Commun.\ Number Theory Phys.\ 11 (2017), 557--655 (\emph{Brown's notes}), for motivic periods, the coaction and the weight-graded face maps; F.~Brown, \emph{Feynman amplitudes, coaction principle, and cosmic Galois group}, Commun.\ Number Theory Phys.\ 11 (2017), 453--556, \S5, for the affine model; A.~Huber and S.~M\"uller-Stach, \emph{Periods and Nori motives}, Springer, 2017, for the functoriality of relative cohomology in Nori's category, whose objects and morphisms realize in Brown's category; and H.~Hironaka, \emph{Resolution of singularities of an algebraic variety over a field of characteristic zero}, Ann.\ of Math.\ 79 (1964), 109--326. In this section a \emph{pair} $(V,D)$ is a smooth variety $V$ over $\mathbb Q$ of dimension $m$ with a divisor $D\subset V$ with simple normal crossings and smooth components. Its relative cohomology $H^m(V,D)$ is an object of Brown's category, a morphism of pairs giving a morphism of Nori's category and so of Brown's (Huber and M\"uller-Stach). A regular form $\eta$ of degree $m$ on $V$ is closed and restricts to zero on the components of $D$, which have smaller dimension; $[\eta]$ is its class in $H^m_{\mathrm{dR}}(V,D)$. A compact semialgebraic $\sigma\subset V(\mathbb R)$ that is the closure of an oriented open semialgebraic set $\sigma^\circ\subseteq(V\setminus D)(\mathbb R)$ of dimension $m$, with $\sigma\setminus\sigma^\circ\subseteq D$, is a \emph{relative cycle}: a semialgebraic triangulation makes the sum of its oriented top simplices a chain whose boundary lies in $D(\mathbb R)$, and its class in $H_m(V(\mathbb C),D(\mathbb C);\mathbb Q)$ is written $\sigma$ again. The comparison of the two realizations is integration, so the period of $[H^m(V,D),\sigma,[\eta]]^{\mathrm{mot}}$ is $\int_\sigma\eta$. The positive cells $\widetilde\sigma_Z$ are oriented as in Definition~\hb{def:mot-mperiod}.

\begin{hdefinition}[\hb{def:mot-model} (The model of a family)]
As stated.
\end{hdefinition}

\begin{proof}
$\Omega_\cE$ is the contraction of $d\alpha_{e_1}\wedge\cdots\wedge d\alpha_{e_n}$ with the vector field $\sum_e\alpha_e\,\partial/\partial\alpha_e$, which does not depend on the order, while the volume form changes by the sign of the permutation. The coefficient $f=\alpha^{\mvec-\one}\prod_jK_j^{-s_j}\Psi_j^{-t_j}$ of $\omega_F=f\,\Omega_\cE$ is homogeneous of the degree $m_\cE-\sum_j(s_jr_j+t_jr^\ast_j)-n=\kappa_\cE-n=-n$ (Definition~\hb{def:family}), the masses being on the balance wall. With the affine coordinates $x_i=\alpha_{e_i}/\alpha_{e_n}$, $i<n$, of the chart $\alpha_{e_n}\neq0$ of $\mathbb P(\mathbb Q^\cE)$, a direct computation gives $dx_1\wedge\cdots\wedge dx_{n-1}=(-1)^{n-1}\alpha_{e_n}^{-n}\Omega_\cE$, so that $f\,\Omega_\cE=(-1)^{n-1}f(x_1,\dots,x_{n-1},1)\,dx_1\wedge\cdots\wedge dx_{n-1}$ is the pull-back of a rational form on that chart. These forms agree on the overlaps of the charts, all pulling back to $f\,\Omega_\cE$ on $\mathbb Q^\cE\setminus\{0\}$; so $\omega_F$ is a rational form of degree $n-1$ on $\mathbb P(\mathbb Q^\cE)$, and on $X_\cE$, which contains the torus of $\mathbb P(\mathbb Q^\cE)$ as a dense open set.
\end{proof}

\begin{hdefinition}[\hb{def:mot-flag} (Flags and steps)]
As stated.
\end{hdefinition}

\begin{proof}
The steps are $\ell+1$ nonempty sets that partition $\cE$, so $\sum_i(|Z_i|-1)=n-1-\ell$. For $\ell\ge1$ the step lists at $Z_i$ depend on $S_{i-1}$ and $S_i$ alone (Proposition~\hb{prop:mot-chainbottom}), and the $\alpha_J$ with $J\subseteq Z_i$ on $Z_i$ alone; so does $\operatorname{Let}(Z_i;\mathfrak f)$. When $|Z_i|=1$, $X_{Z_i}$ is the toric variety of the fan $\{0\}$ in $\mathbb R^{Z_i}/\mathbb R\one=0$, a point, and the step has no letter: the basis polynomials of the lists on one mark are $1$ and $\alpha_e$, monomials, and no $J\subseteq Z_i$ has two marks.
\end{proof}

\begin{hlemma}[\hb{lem:mot-charts} (Charts, strata and the letters in a chart)]
As stated.
\end{hlemma}

\begin{proof}
(i) The fan lies in $N_{\mathbb R}$, $N=\mathbb Z^\cE/\mathbb Z\one$, with the dual lattice $M=\{\mathbf a\in\mathbb Z^\cE:\ a_\cE=0\}$, whose elements are the characters $\alpha^{\mathbf a}$ of the torus. For an order $e_1,\dots,e_n$ the vectors $\one_{S_1},\dots,\one_{S_{n-1}},\one$ form a basis of $\mathbb Z^\cE$, their matrix in the basis $(\mathbf e_{e_k})_k$ being triangular with ones on the diagonal; so the images of $\one_{S_1},\dots,\one_{S_{n-1}}$ form a basis of $N$, with the dual basis $\mathbf e_{e_i}-\mathbf e_{e_{i+1}}$, $i<n$, of $M$, since $\ip{\one_{S_k}}{\mathbf e_{e_i}-\mathbf e_{e_{i+1}}}$ is $1$ for $k=i$ and $0$ otherwise. Every $\mathbf x\in\mathbb R^\cE$ lies in such a cone modulo $\one$: ordering $\cE$ with $x_{e_1}\ge\cdots\ge x_{e_n}$, $\mathbf x=\sum_{i<n}(x_{e_i}-x_{e_{i+1}})\one_{S_i}+x_{e_n}\one$. The cones of two chains meet in the cone of their common sets, a face of both. So the braid fan is a complete fan whose cones are spanned by parts of bases of $N$; it is the normal fan of the polytope $\sum_{\{a,b\}}\operatorname{conv}\{\mathbf e_a,\mathbf e_b\}$, the sum over the pairs of marks, since on the interior of the cone of an order a linear form $\mathbf x$ has $x_{e_1}>\cdots>x_{e_n}$ and attains its least value on that polytope at the one vertex $\sum_{\{a,b\}}\mathbf e_{\max(a,b)}$, the maximum taken in the order, which is the same on the whole cone and differs for different orders. By the dictionary between fans and toric varieties (Fulton, Chapters~1 to~3), which applies to such a fan, $X_\cE$ is smooth and complete, and projective, the fan being a normal fan; the chart of the cone of the chain is the affine space with the coordinates $\alpha^{\mathbf e_{e_i}-\mathbf e_{e_{i+1}}}=\varepsilon_i$; the charts of the maximal cones cover $X_\cE$; and the closure $D_S$ of the orbit of the ray $\rho_S$ meets the chart exactly when $\rho_S$ is a ray of its cone, in the coordinate hyperplane of the dual coordinate, here $\{\varepsilon_i=0\}$ for $S=S_i$, several of these closures meeting exactly when their rays span a cone, here when their sets form a chain (the orbit--cone correspondence). In every chart $D$ is a union of coordinate hyperplanes, so it has simple normal crossings and smooth components. In every chart the torus is $\{\varepsilon_i\neq0\ \text{for all }i\}$ and $\widetilde\sigma^\circ_\cE=(0,\infty)^{n-1}$, since $\alpha_{e_j}/\alpha_{e_n}=\prod_{i\ge j}\varepsilon_i$; its closure in the open set $U(\mathbb R)$ is $[0,\infty)^{n-1}=\widetilde\sigma_\cE\cap U$. The points of $\widetilde\sigma^\circ_\cE$ with $\alpha_{e_1}<\cdots<\alpha_{e_n}$ are $(0,1)^{n-1}$ in the chart, with the closure $[0,1]^{n-1}$, compact and so closed in $X_\cE(\mathbb R)$. Every point of $\widetilde\sigma^\circ_\cE$ is a limit of points with distinct coordinates, so $\widetilde\sigma^\circ_\cE$ lies in the union of the sectors, a finite union of compact sets, and so does its closure $\widetilde\sigma_\cE$, which contains every sector. The open cube $(0,1)^{n-1}$ of one order is open in $X_\cE(\mathbb R)$ and misses the open sector of every other order, hence its closure; so two sectors meet in the boundary of a cube, of measure zero.

(ii) For a basis $T$ of a table of rank $\delta$, $\alpha^T/\alpha_{e_n}^\delta=\alpha^{\one_T-\delta\mathbf e_{e_n}}=\prod_{i<n}\varepsilon_i^{\,|T\cap S_i|}$, by the dual basis of (i) and $e_n\notin S_i$. So $K_\cB/\alpha_{e_n}^\delta$ is the sum of these monomials over the bases. Every exponent $|T\cap S_i|$ is at least $d(S_i)$ (Theorem~\hb{thm:window}), and along the maximal chain $(S_i)_{i<n}$ exactly one basis has $|T\cap S_i|=d(S_i)$ for every $i$, its term having the coefficient $1$ (Proposition~\hb{prop:mot-chainbottom}); dividing by its monomial leaves $\widetilde f$, a polynomial with nonnegative integer coefficients, the numbers of bases with given traces, and with the constant term $1$. For $\Psi_\cB=K_{\cB^\ast}$ the same holds with the dual table, whose floors are the $d^\ast(S)$ (Proposition~\hb{prop:identity}, Theorem~\hb{thm:window}), and for $\alpha_J$ with the list $\binom J1$ on $\cE$, whose floor at $S$ is $1$ or $0$ as $J\subseteq S$ or not. In the sector coordinates of Theorem~\hb{thm:exist}(ii), with $\pi(k)=e_{n+1-k}$, the coordinates are $q_k=\varepsilon_{n+1-k}$ and the tails are $U_k=S_{n+1-k}$, and $\widetilde f$ is the sector letter there. On $\varepsilon\ge0$, $\widetilde f\ge1$.

On $\{\varepsilon_i=0:\ i\in I\}$ only the bases $T$ with $|T\cap S_i|=d(S_i)$ for every $i\in I$ leave a term. By Proposition~\hb{prop:mot-chainbottom} for the chain $(S_i)_{i\in I}$ they are the unions $\bigsqcup_ZX_Z$ of one basis $X_Z$ of the step table at the bottom at each step $Z$ of that chain. For $i\notin I$ let $Z=S_b\setminus S_a$ be the step with $S_a\subsetneq S_i\subsetneq S_b$, where $S_a$ may be $\varnothing$ and $S_b$ may be $\cE$; then $|T\cap S_i|=d(S_a)+|X_Z\cap(S_i\setminus S_a)|$. The basis at the bottom of the maximal chain is one of these unions and has the least traces, so $d(S_i)-d(S_a)$ is the floor $d_Z(S_i\setminus S_a)$ of $S_i\setminus S_a$ in the step table at $Z$, and
\[
\widetilde f\big|_{\{\varepsilon_i=0:\,i\in I\}}=\prod_Z\ \sum_{X_Z}\ \prod_{a<k<b}\varepsilon_k^{\,|X_Z\cap(S_k\setminus S_a)|-d_Z(S_k\setminus S_a)} .
\]
The factor of $Z$ is the polynomial that the rule above attaches to $K$ of the step table in the chart of $X_Z$ of the order $e_{a+1},\dots,e_b$, whose initial segments are the $S_k\setminus S_a$ and whose coordinates are the $\varepsilon_k$, $a<k<b$. For $\Psi$ the same holds with the tops and the dual tables, and for $\alpha_J$ with its step letters (Proposition~\hb{prop:mot-chainbottom}). Finally, on the torus part of $U$, $f=0$ exactly where $\widetilde f=0$. No irreducible component of $\{\widetilde f=0\}$ lies in a coordinate hyperplane $\{\varepsilon_i=0\}$, since $\varepsilon_i$ does not divide $\widetilde f$, whose constant term is $1$; so each component is the closure of its part in the torus, and $\{\widetilde f=0\}$ is the closure in $U$ of the zero set of $f$ in the torus, that is $\widetilde Y_f\cap U$.

(iii) In the affine coordinates $x_j=\alpha_{e_j}/\alpha_{e_n}=\prod_{i\ge j}\varepsilon_i$, $\Omega_\cE/\alpha^\cE=(-1)^{n-1}\bigwedge_{j<n}dx_j/x_j$ (proof of Definition~\hb{def:mot-model}), and $\bigwedge_jd\log x_j=\bigwedge_id\log\varepsilon_i$, the matrix of the $\log x_j$ in the $\log\varepsilon_i$ being triangular with ones on the diagonal; so $\Omega_\cE/\alpha^\cE=(-1)^{n-1}\bigwedge_{i<n}d\varepsilon_i/\varepsilon_i$. By (ii), $\alpha^{\mvec}/\alpha_{e_n}^{m_\cE}=\prod_i\varepsilon_i^{m_{S_i}}$, $K_j/\alpha_{e_n}^{r_j}=\prod_i\varepsilon_i^{d_j(S_i)}\widetilde K_j$ and $\Psi_j/\alpha_{e_n}^{r^\ast_j}=\prod_i\varepsilon_i^{d^\ast_j(S_i)}\widetilde\Psi_j$. On the balance wall the powers of $\alpha_{e_n}$ cancel, and the exponent of $\varepsilon_i$ is $m_{S_i}-\sum_j\big(s_jd_j(S_i)+t_jd^\ast_j(S_i)\big)-1=\kappa_{S_i}-1$.
\end{proof}

\begin{slemma}[Integrals over the positive cell]\label{lem:motive-sections}
Let $\phi$ be a measurable function on $\mathbb R^\cE_{>0}$, homogeneous of degree zero, $n=|\cE|\ge2$. Then, in $[0,\infty]$ when $\phi\ge0$, and whenever one side converges absolutely,
\[
\int_{\widetilde\sigma^\circ_\cE}\phi\,\frac{\Omega_\cE}{\alpha^\cE}=\int_{\Delta_\cE}\phi(u)\,u^{-\one}\,d\omega_\cE(u)=\int_{\Delta_B\times\Delta_R}\int_0^\infty\phi(\rho u,u')\,\frac{d\rho}\rho\ u^{-\one}u'^{-\one}\,d\omega_B(u)\,d\omega_R(u')
\]
for every partition $\cE=B\sqcup R$ into nonempty sets; in the chart of an order the left side is $\int_{(0,\infty)^{n-1}}\phi\prod_id\varepsilon_i/\varepsilon_i$. If $\phi\,\Omega_\cE/\alpha^\cE$ is the restriction of a form that is regular on an open set containing $\widetilde\sigma_\cE$, the integrals converge absolutely and are its integral over $\widetilde\sigma_\cE$.
\end{slemma}

\begin{proof}
In the logarithmic coordinates $y=\log\alpha$, $\alpha^{-\one}d\alpha=dy$. For an order of $\cE$ the linear map $y\mapsto\big((y_{e_i}-y_{e_{i+1}})_{i<n},y_{e_n}\big)$ has the determinant $\pm1$, and its first $n-1$ coordinates are the $\log\varepsilon_i$, on which $\phi$ depends alone. For a positive $h$ homogeneous of degree one and a measurable $\varphi\ge0$ on $(0,\infty)$, at fixed $\varepsilon$ the substitution $\rho=e^{y_{e_n}}h(\alpha/\alpha_{e_n})$ has $d\rho/\rho=dy_{e_n}$; by Tonelli
\[
\int_{\mathbb R^\cE_{>0}}\phi(\alpha)\,\varphi\big(h(\alpha)\big)\,\alpha^{-\one}d\alpha=\int_0^\infty\varphi(\rho)\frac{d\rho}\rho\ \int_{(0,\infty)^{n-1}}\phi\prod_i\frac{d\varepsilon_i}{\varepsilon_i},
\]
the same holds with $|\phi|$, and with $\phi$ whenever the integrals of $|\phi|$ are finite. With $h=|\alpha|$, Lemma~\hb{lem:scaleshape} on $\cE$ writes the left side as $\int_0^\infty\varphi(\rho)\frac{d\rho}\rho\int_{\Delta_\cE}\phi(u)u^{-\one}d\omega_\cE(u)$. With $h=|\alpha_R|$, Lemma~\hb{lem:scaleshape} on $B$ and on $R$, $\alpha_B=\rho'u$ and $\alpha_R=\lambda u'$, gives $\alpha^{-\one}d\alpha=u^{-\one}u'^{-\one}\frac{d\rho'}{\rho'}\frac{d\lambda}\lambda d\omega_B(u)d\omega_R(u')$ and $\phi(\rho'u,\lambda u')=\phi(\rho u,u')$ with $\rho'=\lambda\rho$, so the left side is the third expression times $\int\varphi(\lambda)d\lambda/\lambda$. A $\varphi$ with $\int_0^\infty\varphi\,d\rho/\rho=1$ gives the identities. In the chart $\Omega_\cE/\alpha^\cE=(-1)^{n-1}\bigwedge_id\varepsilon_i/\varepsilon_i$ (proof of Lemma~\hb{lem:mot-charts}(iii)), which with the orientation of $\widetilde\sigma^\circ_\cE=(0,\infty)^{n-1}$ is the measure $\prod_id\varepsilon_i/\varepsilon_i$. If the form is regular near the compact $\widetilde\sigma_\cE$, it is $g\,d\varepsilon$ on each sector cube with $g$ continuous on the cube; so its integral over each sector converges absolutely, the sectors cover $\widetilde\sigma_\cE$ and meet in sets of measure zero, and $\widetilde\sigma_\cE\setminus\widetilde\sigma^\circ_\cE$ lies in the coordinate hyperplanes of the charts (Lemma~\hb{lem:mot-charts}(i)).
\end{proof}

\begin{slemma}[The continued period off the balance wall]\label{lem:motive-continued}
Let $F$ be a family of lists on $\cE$, $n\ge2$. For an order $\pi$ with the chain $(S^\pi_i)_{i<n}$ let $\widetilde K_j,\widetilde\Psi_j$ and $\widetilde L_\pi$ be the sector letters of $K_j,\Psi_j$ and of the linear letter $\alpha_\cE=|\alpha|$ in its chart (Lemma~\hb{lem:mot-charts}(ii)), and $h_\pi(\varepsilon;\kappa,\mathbf s,\mathbf t)=\widetilde L_\pi^{-\kappa}\prod_j\widetilde K_j^{-s_j}\widetilde\Psi_j^{-t_j}$. As meromorphic functions of $(\mvec,\mathbf s,\mathbf t)\in\mathbb C^\cE\times\mathbb C^N\times\mathbb C^N$,
\[
\frac{J_F(\mvec;\mathbf s,\mathbf t)}{\Gamma(\kappa_\cE)}=\sum_\pi\operatorname{cube}\Big(\big(\kappa_{S^\pi_i}\big)_{i<n},\ (\kappa_\cE,\mathbf s,\mathbf t);\ h_\pi\Big),
\]
with $\operatorname{cube}$ the cube integral of Lemma~\hb{lem:cube}; the right side is holomorphic off the hyperplanes $\kappa_S=-j$, $\varnothing\neq S\subsetneq\cE$, $j\in\mathbb N$. So the period $P_F=J_F/\Gamma(\kappa_\cE)$ (Theorem~\hb{thm:projective}) has its poles on these hyperplanes only, and on the balance wall it is the right side at $\kappa_\cE=0$, where $h_\pi=\prod_j\widetilde K_j^{-s_j}\widetilde\Psi_j^{-t_j}$.
\end{slemma}

\begin{proof}
On the region where $\re\kappa_S>0$ for every nonempty $S\subseteq\cE$, which contains the points with $\mathbf s=\mathbf t=0$ and $\re\mvec>0$, Theorem~\hb{thm:exist}(ii) is the display, read in the chart coordinates as in the proof of Lemma~\hb{lem:mot-charts}(ii): its sector letters $k_{j,\pi},\psi_{j,\pi},L_\pi$ are $\widetilde K_j,\widetilde\Psi_j,\widetilde L_\pi$. The sector letters are at least $1$ on the cube, so $h_\pi$ is smooth in $\varepsilon$ near $[0,1]^{n-1}$, with every $\varepsilon$-derivative continuous in all the variables and holomorphic in $(\kappa,\mathbf s,\mathbf t)\in\mathbb C^{1+2N}$. The map $(\mvec,\mathbf s,\mathbf t)\mapsto\big((\kappa_{S^\pi_i})_{i<n},\kappa_\cE,\mathbf s,\mathbf t\big)$ is a linear isomorphism of $\mathbb C^\cE\times\mathbb C^{2N}$, since $\mvec\mapsto(m_{S^\pi_1},\dots,m_{S^\pi_{n-1}},m_\cE)$ is triangular. Lemma~\hb{lem:cube} continues each term to a meromorphic function on the whole space with poles on the hyperplanes $\kappa_{S^\pi_i}=-j$ only. The sum agrees with the meromorphic $J_F/\Gamma(\kappa_\cE)$ (Theorems~\hb{thm:exist}(iii) and~\hb{thm:projective}) on a nonempty open set, hence everywhere. At $\kappa_\cE=0$, $\widetilde L_\pi^{-\kappa_\cE}=1$.
\end{proof}

\begin{hproposition}[\hb{prop:mot-model} (The model)]
As stated.
\end{hproposition}

\begin{proof}
(i) $\widetilde\sigma_\cE$ is the union of the finitely many sectors, cubes (Lemma~\hb{lem:mot-charts}(i)), so it is compact. In a chart $\widetilde\sigma_\cE\cap U=[0,\infty)^{n-1}$ and $\widetilde\sigma^\circ_\cE=(0,\infty)^{n-1}$, whose difference lies in the coordinate hyperplanes, that is in $D$; and the charts cover $X_\cE$. For a letter $f$ of a list or a linear letter, $\widetilde Y_f\cap U=\{\widetilde f=0\}$ and $\widetilde f\ge1$ on $[0,\infty)^{n-1}$ (Lemma~\hb{lem:mot-charts}(ii)).

(ii) Let $\varnothing\neq S\subsetneq\cE$ and take an order whose chain contains $S=S_i$. $D_S$ meets $V_F$: the face $\{\varepsilon_i=0\}\cap[0,1]^{n-1}$ of the sector cube lies in $D_S\cap\widetilde\sigma_\cE$, and $\widetilde\sigma_\cE\subseteq V_F(\mathbb R)$ by (i). On $V_F\cap U$ every $\widetilde K_j$ with $s_j\neq0$ and every $\widetilde\Psi_j$ with $t_j\neq0$ is a unit, $V_F$ omitting their zero sets, so by Lemma~\hb{lem:mot-charts}(iii) $\omega_F$ is a unit times $\prod_k\varepsilon_k^{\kappa_{S_k}-1}d\varepsilon_1\wedge\cdots\wedge d\varepsilon_{n-1}$ there. Its order along $D_S$, read at a point of $D_S\cap V_F\cap U$, is $\kappa_S-1$, and on the torus, where every $\varepsilon_k$ is a unit, it has neither zeros nor poles. A prime divisor of $V_F$ either meets the torus or is one of the $D_S\cap V_F$, each irreducible as a nonempty open subset of the irreducible $D_S$. So the divisor of $\omega_F$ on the smooth $V_F$ is $\sum_S(\kappa_S-1)(D_S\cap V_F)$, and $\omega_F$ is regular exactly when this divisor is effective, that is when $F$ is convergent. On the cube of a sector, $\omega_F=\pm\prod_k\varepsilon_k^{\kappa_{S^\pi_k}-1}h_\pi\,d\varepsilon$ with $h_\pi=\prod_j\widetilde K_j^{-s_j}\widetilde\Psi_j^{-t_j}$ positive and continuous on the compact cube, so bounded between two positive constants; as $\int_0^1\varepsilon^{\kappa-1}d\varepsilon$ converges for $\kappa\ge1$, the integral over $\widetilde\sigma_\cE$ converges when $F$ is convergent. By Lemma~\hb{lem:motive-sections} it is the integral over $\Delta_\cE$ of $u^{\mvec-\one}\prod_jK_j^{-s_j}\Psi_j^{-t_j}$, the period (Theorem~\hb{thm:projective}), positive with the orientation in which $\Omega_\cE/\alpha^\cE$ is positive on $\widetilde\sigma^\circ_\cE$.

(iii) $V_F$ is open in the smooth $X_\cE$. The complement in $X_\cE$ of the $\widetilde Y_{\alpha_J}$, $|J|\ge2$, is Brown's affine model (Brown, \emph{Feynman amplitudes}, \S5), which applies because $X_\cE$ is the iterated blow-up of $\mathbb P(\mathbb Q^\cE)$ along its coordinate linear subspaces, the toric variety of the subdivision of its fan by all chains, and the removed hypersurfaces are the strict transforms of the hyperplanes $\alpha_J=0$ of all the sets $J$ of at least two marks. In toric terms the divisor $A=\sum_{|J|\ge2}\widetilde Y_{\alpha_J}$ is ample. In the chart of an order, $\alpha_J/\alpha_y=\widetilde\alpha_J$ for the last mark $y$ of $J$ in the order (Lemma~\hb{lem:mot-charts}(ii)), so $\widetilde Y_{\alpha_J}$ is the divisor of the section $\sum_{e\in J}\alpha^{\mathbf e_e}$ of the line bundle of the simplex $\Delta_J=\operatorname{conv}\{\mathbf e_e:\ e\in J\}$, whose vertex of least value on the cone of the order is $\mathbf e_y$; $A$ is then the divisor of a section of the line bundle of the polytope $\sum_J\Delta_J$, whose vertex of least value on the interior of the cone of the order $e_1,\dots,e_n$ is $\sum_J\mathbf e_{y_J}=\sum_k(2^{k-1}-1)\,\mathbf e_{e_k}$, the mark $e_k$ being the last mark of $2^{k-1}-1$ sets $J$; it is the only point of least value there and differs for different orders. So the normal fan of $\sum_J\Delta_J$ is the braid fan, the divisor $A$ is ample (Fulton, Chapter~3), a multiple of it is a hyperplane section in a projective embedding, and its complement is affine. A finite set $\mathcal G$ of polynomials that contains the $\alpha_J$ removes in addition the $\widetilde Y_g$, $g\in\mathcal G$, divisors of the smooth $X_\cE$, hence Cartier divisors; the complement of a Cartier divisor in an affine variety is affine, the inclusion being locally the complement of the zero set of one function. $D_F$ is the restriction of $D$ to the open $V_F$, with simple normal crossings and smooth components (Lemma~\hb{lem:mot-charts}(i)). A stratum is an intersection of components meeting transversally, so smooth, and closed in the affine $V_F$, so affine.

(iv) Let $\tau$ be the cone of the flag. By the orbit--cone correspondence (Fulton, Chapter~3), $D_{S_1}\cap\cdots\cap D_{S_\ell}$ is the toric variety of the star of $\tau$, the fan in $N/(N\cap\mathbb R\tau)$ of the images of the cones that contain $\tau$. Restriction to the steps, $\mathbf x\mapsto(\mathbf x|_{Z_i}\bmod\mathbb Z\one_{Z_i})_i$, maps $\mathbb Z^\cE$ onto $\bigoplus_i\mathbb Z^{Z_i}/\mathbb Z\one_{Z_i}$ with the kernel spanned by $\one_{Z_1},\dots,\one_{Z_{\ell+1}}$, that is by $\one_{S_1},\dots,\one_{S_\ell}$ and $\one$; a set $S$ with $S_{i-1}\subsetneq S\subsetneq S_i$ goes to $\one_{S\setminus S_{i-1}}$ in the summand of $Z_i$ and to $0$ in the others. So the cones of the chains that refine the flag go to the products of cones of chains of the steps, the star of $\tau$ is the product of the braid fans of the steps, and $D_{S_1}\cap\cdots\cap D_{S_\ell}\cong\prod_iX_{Z_i}$. The charts that meet the stratum are those of the orders whose chains contain the flag; in such a chart the stratum is $\{\varepsilon_i=0:\ i\in I\}$, $I$ the positions of the sets of the flag, and the remaining coordinates are the chart coordinates of the steps. By Lemma~\hb{lem:mot-charts}(ii) $\widetilde Y_f$ meets it in the union over the steps $Z$ of the zero sets of the polynomials attached to the step letters of $f$ at $Z$, and by the same lemma on $Z$ these are the $\widetilde Y_g\cap U_Z$ for the step letters $g$, empty for a monomial. The step letters at $Z$ of the augmented letters of $F$ that are not monomials form $\operatorname{Let}(Z;\mathfrak f)$: the step letter of $\alpha_J$ is $\alpha_{J\cap Z}$ at the first step $Z$ with $J\subseteq S_b$ and $1$ elsewhere, and every $\alpha_{J'}$ with $J'\subseteq Z$, $|J'|\ge2$, arises from $J=J'$. So the stratum of $\mathfrak f$ is $\prod_iV_{Z_i;\mathfrak f}$. Its boundary, its meeting with the $D_S$ of the sets $S$ outside the flag, is $\{\varepsilon_k=0\}$ for the other positions $k$, the union of the pull-backs of the toric boundaries of the factors; and $\widetilde\sigma_\cE$ meets it in $[0,\infty)^{n-1}\cap\{\varepsilon_i=0:\ i\in I\}=\prod_Z[0,\infty)^{|Z|-1}$, chart by chart, that is in $\prod_i\widetilde\sigma_{Z_i}$.

(v) Where every proper wall has positive real part the period converges (Theorem~\hb{thm:projective}), and by Lemma~\hb{lem:motive-sections}, the sectors and Lemma~\hb{lem:mot-charts}(iii) it is $\sum_\pi\int_{[0,1]^{n-1}}\prod_i\varepsilon_i^{\kappa_{S^\pi_i}-1}h_\pi\,d\varepsilon=\sum_\pi\operatorname{cube}(\boldsymbol\kappa_\pi(\mvec);h_\pi)$, each term converging absolutely. By Lemma~\hb{lem:motive-continued} the right side is, on the whole balance wall, the restriction of $J_F/\Gamma(\kappa_\cE)$, meromorphic there jointly in the masses and the exponents, with poles on the hyperplanes $\kappa_S=-j$ only.
\end{proof}

\begin{hdefinition}[\hb{def:mot-mperiod} (Motivic periods; Galois conjugates; open spaces)]
As stated.
\end{hdefinition}

\begin{proof}
The motivic periods, the de Rham periods, their relations and the coaction are those of Brown's notes (\S2 and \S3.1), where the coaction is shown to be well defined. For a convergent family, $(V_F,D_F)$ is a pair with $V_F$ affine (Proposition~\hb{prop:mot-model}(iii)), so $H_F$ is an object of Brown's category; $\widetilde\sigma_\cE$ is a relative cycle of it, compact, the closure of the open cell, with $\widetilde\sigma_\cE\setminus\widetilde\sigma^\circ_\cE\subseteq D_F$ and $\widetilde\sigma_\cE\subseteq V_F(\mathbb R)$ (Proposition~\hb{prop:mot-model}(i)); and $\omega_F$ is a regular form of degree $n-1$ on $V_F$ (Proposition~\hb{prop:mot-model}(ii)), with its class $[\omega_F]$. The period of $P^{\mathrm{mot}}_F$ is $\int_{\widetilde\sigma_\cE}\omega_F=P_F(\mvec;\mathbf s,\mathbf t)$ (Proposition~\hb{prop:mot-model}(ii)). For a step with $|Z|\ge2$ the letters of the step are letters of lists on $Z$ and the $\alpha_J$, $J\subseteq Z$, $|J|\ge2$; so by Proposition~\hb{prop:mot-model}(i) and~(iii) on $Z$, $(V_{Z;\mathfrak f},\partial V_{Z;\mathfrak f})$ is a pair with $V_{Z;\mathfrak f}$ affine and $\widetilde\sigma_Z$ is a relative cycle of it.
\end{proof}

\begin{hlemma}[\hb{lem:mot-conj} (Galois conjugates)]
As stated.
\end{hlemma}

\begin{proof}
$\mathcal P^{\mathrm{dR}}$ is the ring of functions on the de Rham Galois group, $[M,f,\omega]^{\mathrm{dR}}$ being the function $g\mapsto f(g\,\omega)$ (Brown's notes, \S2); evaluation at the identity is a $\mathbb Q$-linear map $\varphi_1:\mathcal P^{\mathrm{dR}}\to\mathbb Q$ with $\varphi_1([M,f,\omega]^{\mathrm{dR}})=f(\omega)$.

(i) For $\xi=[M,\gamma,\omega]^{\mathrm{mot}}$, $(\mathrm{id}\otimes\varphi_1)\Delta\xi=\sum_i[M,\gamma,e_i]^{\mathrm{mot}}\,e_i^\vee(\omega)=[M,\gamma,\omega]^{\mathrm{mot}}$ by bilinearity; by linearity the same holds for every $\xi$.

(ii) $(\mathrm{id}\otimes\varphi)\Delta[M,\gamma,\omega]^{\mathrm{mot}}=\big[M,\gamma,\sum_i\varphi([M,e_i^\vee,\omega]^{\mathrm{dR}})\,e_i\big]^{\mathrm{mot}}\in L(M,\gamma)$, and the left factors of the coaction of an element of $L(M,\gamma)$ lie in $L(M,\gamma)$.

(iii) $\mathrm C(\xi)$ is the image of the linear map $\varphi\mapsto(\mathrm{id}\otimes\varphi)\Delta\xi$, a space; so $\mathrm C(\xi+\xi')\subseteq\mathrm C(\xi)+\mathrm C(\xi')$ and $\mathrm C(q\,\xi)=q\,\mathrm C(\xi)=\mathrm C(\xi)$ for $q\neq0$. Write $\Delta\xi=\sum_i\xi_i\otimes\xi'_i$ with the $\xi'_i$ linearly independent; a $\varphi$ with $\varphi(\xi'_k)$ equal to $1$ for $k=i$ and to $0$ otherwise shows $\xi_i\in\mathrm C(\xi)$. Likewise $\Delta\xi'=\sum_k\eta_k\otimes\eta'_k$ with $\eta_k\in\mathrm C(\xi')$. The coaction is a homomorphism of rings (Brown's notes, \S2), so $(\mathrm{id}\otimes\varphi)\Delta(\xi\xi')=\sum_{i,k}\xi_i\eta_k\,\varphi(\xi'_i\eta'_k)\in\mathrm C(\xi)\cdot\mathrm C(\xi')$.
\end{proof}

\begin{hlemma}[\hb{lem:mot-birational} (Birational invariance)]
As stated.
\end{hlemma}

\begin{proof}
\emph{The resolution.} Let $\Gamma\subseteq X_1\times X_2$ be the closure of the graph of $\Phi$, a projective variety. The graph is closed in $U_1\times X_2$, so $\Gamma\cap(U_1\times X_2)$ is the graph; it is also the graph of $\Phi^{-1}$, closed in $X_1\times U_2$, so $\Gamma\cap(X_1\times U_2)$ is the graph too; and the graph, isomorphic to $U_1$, is a smooth dense open subset of $\Gamma$. By resolution of singularities in characteristic zero (Hironaka), in the form that gives for a projective variety and a smooth dense open subset of it a smooth projective variety with a proper birational morphism onto it that is an isomorphism over that subset and under which the complement pulls back to a divisor with simple normal crossings and smooth components, there is such a $\varpi:\widetilde X\to\Gamma$ for the graph; the property used is that the graph of $\Phi$ is smooth and the ground field $\mathbb Q$ has characteristic zero. Let $q_i:\widetilde X\to X_i$ be the two projections; by the above $q_1^{-1}(U_1)=q_2^{-1}(U_2)=\varpi^{-1}(\text{graph})$, on which $q_1$ and $q_2=\Phi\circ q_1$ are isomorphisms onto $U_1$ and $U_2$. Put $\widetilde D_i=q_i^{-1}(D_i)$ and $\widetilde Y=q_1^{-1}(Y_1)\cup q_2^{-1}(Y_2)$, both closed. A point of $\widetilde D_1$ is not over the graph, $D_1$ missing $U_1$; so $\widetilde D_1$, the support of the pull-back of a Cartier divisor under the dominant $q_1$, is a divisor contained in the divisor $\varpi^{-1}(\Gamma\setminus\text{graph})$ and so a union of its components; likewise $\widetilde D_2$ and $\widetilde D_1\cup\widetilde D_2$. All three have simple normal crossings and smooth components.

\emph{The cycle.} Let $\widetilde\sigma$ be the closure in $\widetilde X(\mathbb R)$ of $\widetilde\sigma^\circ=q_1^{-1}(\sigma_1^\circ)$, which $q_1$ maps isomorphically onto $\sigma^\circ_1$ and $q_2$ onto $\sigma^\circ_2$; orient it by $q_1$. It is compact, and $q_i(\widetilde\sigma)=\sigma_i$, which misses $Y_i$; so $\widetilde\sigma$ misses $\widetilde Y$. If $x\in\widetilde\sigma\setminus\widetilde\sigma^\circ$ had $q_1(x)\in\sigma^\circ_1\subseteq U_1$, $x$ would lie over the graph, where $q_1$ is injective, and so in $\widetilde\sigma^\circ$; hence $q_1(x)\in\sigma_1\setminus\sigma^\circ_1\subseteq D_1$ and $x\in\widetilde D_1$. The same argument with $q_2$ and $\Phi(\sigma^\circ_1)=\sigma^\circ_2$ gives $x\in\widetilde D_2$. So $\widetilde\sigma$ is a relative cycle of $(\widetilde X\setminus\widetilde Y,\widetilde D_i\setminus\widetilde Y)$ for $i=1,2$ and of $(\widetilde X\setminus\widetilde Y,(\widetilde D_1\cup\widetilde D_2)\setminus\widetilde Y)$, and $q_i$ carries its class to that of $\sigma_i$, mapping $\widetilde\sigma^\circ$ onto $\sigma^\circ_i$ with the orientation, $\Phi$ preserving it, and $\widetilde\sigma\setminus\widetilde\sigma^\circ$ into $D_i$.

\emph{The form.} $q_1^\ast\omega_1$ and $q_2^\ast\omega_2$ are rational forms on $\widetilde X$ that agree on the dense open $q_1^{-1}(U_1)$, where $q_2=\Phi\circ q_1$ and $\Phi^\ast\omega_2=\omega_1$; so they are one form $\widetilde\omega$, regular off $q_1^{-1}(Y_1)$ and off $q_2^{-1}(Y_2)$, hence on $\widetilde X\setminus\widetilde Y$.

\emph{The chain of pairs.} The maps
\[
(X_1^\circ,D_1^\circ)\xleftarrow{\ q_1\ }(\widetilde X\setminus\widetilde Y,\widetilde D_1\setminus\widetilde Y)\xrightarrow{\ \mathrm{id}\ }\big(\widetilde X\setminus\widetilde Y,(\widetilde D_1\cup\widetilde D_2)\setminus\widetilde Y\big)\xleftarrow{\ \mathrm{id}\ }(\widetilde X\setminus\widetilde Y,\widetilde D_2\setminus\widetilde Y)\xrightarrow{\ q_2\ }(X_2^\circ,D_2^\circ)
\]
are morphisms of pairs, so their maps on relative cohomology are morphisms of Brown's category (Huber and M\"uller-Stach). Each pulls back the class of the form of degree $m$ to the class of the pulled-back form, and its transpose carries the class of $\widetilde\sigma$, or of $\sigma_1$, $\sigma_2$ at the ends, to the class of the cycle at its target. The relation $[M,\phi^{\mathsf t}\gamma',\omega]^{\mathrm{mot}}=[M',\gamma',\phi(\omega)]^{\mathrm{mot}}$ along the four morphisms gives
\[
\big[H^m(X_1^\circ,D_1^\circ),\sigma_1,\omega_1\big]^{\mathrm{mot}}=\big[H^m(\widetilde X\setminus\widetilde Y,\widetilde D_1\setminus\widetilde Y),\widetilde\sigma,\widetilde\omega\big]^{\mathrm{mot}}=\cdots=\big[H^m(X_2^\circ,D_2^\circ),\sigma_2,\omega_2\big]^{\mathrm{mot}} .
\]

\emph{The particular cases.} With $X_1=X_2$, $D_1=D_2$, $Y_2\supseteq Y_1$ closed and missing $\sigma_1$, $U_1=U_2=X_1\setminus(D_1\cup Y_2)$ and $\Phi$ the identity, the hypotheses hold, $\sigma^\circ_1$ lying in $U_1(\mathbb R)$. For a convergent family, $\omega_F$ is regular on $X_\cE\setminus\bigcup_{f\in\operatorname{Let}(F)}\widetilde Y_f$, by the proof of Proposition~\hb{prop:mot-model}(ii), which uses only that the letters of $F$ are units there; and $V_F$ removes in addition the $\widetilde Y_{\alpha_J}$, which miss $\widetilde\sigma_\cE$ (Proposition~\hb{prop:mot-model}(i)). Two convergent families $F,F'$ on $\cE$ with $\omega_F=\omega_{F'}$ give $X_1=X_2=X_\cE$, $D_1=D_2=D$, $Y_1$ and $Y_2$ the unions of the $\widetilde Y_f$ over the augmented letters of $F$ and of $F'$, both missing $\widetilde\sigma_\cE$, $U_1=U_2$ the torus less $Y_1\cup Y_2$ and $\Phi$ the identity.
\end{proof}

\begin{slemma}[Products of pairs]\label{lem:motive-product}
For $i=1,\dots,p$ let $(V_i,D_i)$ be a pair of dimension $m_i$ with a relative cycle $\sigma_i$ and a regular form $\eta_i$ of degree $m_i$ on $V_i$. Let $V=\prod_iV_i$, $D=\bigcup_i\mathrm{pr}_i^{-1}(D_i)$, $\sigma=\prod_i\sigma_i$ with the product orientation, and $\eta=\mathrm{pr}_1^\ast\eta_1\wedge\cdots\wedge\mathrm{pr}_p^\ast\eta_p$. Then $(V,D)$ is a pair, $\sigma$ is a relative cycle of it, and, with $m=\sum_im_i$,
\[
\big[H^m(V,D),\sigma,[\eta]\big]^{\mathrm{mot}}=\prod_{i=1}^p\big[H^{m_i}(V_i,D_i),\sigma_i,[\eta_i]\big]^{\mathrm{mot}} .
\]
\end{slemma}

\begin{proof}
In products of local coordinates $D$ is a union of coordinate hyperplanes of the factors, so $(V,D)$ is a pair; $\sigma$ is the closure of the oriented $\prod_i\sigma^\circ_i$, of dimension $m$, and $\sigma\setminus\prod_i\sigma^\circ_i\subseteq D$. The cross product $\mu:\bigotimes_iH^{m_i}(V_i,D_i)\to H^m(V,D)$ is a morphism of Nori's category (Huber and M\"uller-Stach), hence of Brown's; on de Rham realizations it carries $\bigotimes_i[\eta_i]$ to $[\eta]$, and its transpose carries $\sigma$ to $\bigotimes_i\sigma_i$, the class of a product of relative cycles being the cross product of their classes; on forms this is Fubini, $\int_\sigma\mathrm{pr}_1^\ast\theta_1\wedge\cdots\wedge\mathrm{pr}_p^\ast\theta_p=\prod_i\int_{\sigma_i}\theta_i$. The relation of Definition~\hb{def:mot-mperiod} for $\mu$, with the product of motivic periods defined by the tensor product of objects (Brown's notes, \S2), gives the identity.
\end{proof}

\begin{hcorollary}[\hb{cor:mot-cremona} (Cremona between motivic periods)]
As stated.
\end{hcorollary}

\begin{proof}
\emph{The walls.} With $d(S)=r-r(\cE\setminus S)$, $d^\ast(S)=|S|-r(S)$, $d(\cE\setminus S)=r-r(S)$ and $d^\ast(\cE\setminus S)=|\cE\setminus S|-r(\cE\setminus S)$ in every table (Definition~\hb{def:skeleton}), and $m_\cE=\sum_j(s_jr_j+t_j(n-r_j))$ on the balance wall,
\begin{gather*}
\kappa_S(F^{\mathrm c})=2\sigma|S|-m_S-\sum_jt_j\big(r_j-r_j(\cE\setminus S)\big)-\sum_js_j\big(|S|-r_j(S)\big),\\
\kappa_{\cE\setminus S}(F)=m_\cE-m_S-\sum_js_j\big(r_j-r_j(S)\big)-\sum_jt_j\big(n-|S|-r_j(\cE\setminus S)\big),
\end{gather*}
and both equal $-m_S+\sum_js_jr_j(S)+\sum_jt_j\big(|S|-r_j+r_j(\cE\setminus S)\big)$, using $2\sigma=\sum_j(s_j+t_j)$. At $S=\cE$ the same computation gives $\kappa_\cE(F^{\mathrm c})=-\kappa_\cE(F)=0$. So $F^{\mathrm c}$ is on its balance wall, at integer data, and convergent with $F$.

\emph{The forms.} Let $\Phi$ be the inversion $\alpha\mapsto\alpha^{-1}$ of the torus, an automorphism that maps $\widetilde\sigma^\circ_\cE$ onto itself. By Proposition~\hb{prop:identity}, $K_j(\alpha^{-1})=\alpha^{-\cE}\Psi_j(\alpha)$ and $\Psi_j(\alpha^{-1})=\alpha^{-\cE}K_j(\alpha)$, and $\Phi^\ast(\Omega_\cE/\alpha^\cE)=(-1)^{n-1}\Omega_\cE/\alpha^\cE$, every $d\log\alpha_e$ changing sign. So
\[
\Phi^\ast\omega_{F^{\mathrm c}}=\alpha^{\mvec-2\sigma\one}\,\alpha^{2\sigma\one}\prod_jK_j^{-s_j}\Psi_j^{-t_j}\,(-1)^{n-1}\frac{\Omega_\cE}{\alpha^\cE}=(-1)^{n-1}\omega_F ,
\]
the factor $\alpha^{2\sigma\one}$ collecting the powers $\alpha^{\cE}$ of the letters, $\sum_j(s_j+t_j)=2\sigma$. Let $U_1$ be the torus less the zero sets of the augmented letters of $F$ and the preimages under $\Phi$ of those of $F^{\mathrm c}$, and $U_2=\Phi(U_1)$; these polynomials are positive at the positive points, so $\widetilde\sigma^\circ_\cE\subseteq U_1(\mathbb R)$. The periods $P_F$ and $P_{F^{\mathrm c}}=\int\omega_{F^{\mathrm c}}$ over the positively oriented cell are both positive, so $\Phi$ maps the oriented $\widetilde\sigma^\circ_\cE$ onto itself with the sign $(-1)^{n-1}$ of the forms; reversing the orientation of the target cell and the sign of $\omega_{F^{\mathrm c}}$ together when that sign is $-1$, Lemma~\hb{lem:mot-birational} applies to $(X_\cE,D,Y_F)$ and $(X_\cE,D,Y_{F^{\mathrm c}})$, $Y$ the unions of the $\widetilde Y_f$ over the augmented letters, and gives $P^{\mathrm{mot}}_{F^{\mathrm c}}=P^{\mathrm{mot}}_F$ by bilinearity.
\end{proof}

\begin{hdefinition}[\hb{def:mot-attach} (Attachment)]
As stated.
\end{hdefinition}

\begin{proof}
Every vertex of a tree of bunches is a split vertex, a node that is not a leaf, or a leaf (Definition~\hb{def:bunchtree}), and the first two have children: the molecules of a split vertex, and the room of a node. Along the descent $Z\subseteq M_w$ at every vertex $w$ reached: this holds at the root, a split vertex passes $Z$ to the child that contains it, and a node passes $Z$ to its room only when $Z\subseteq M_w\setminus B_w$, the marks of the room. At each vertex exactly one rule applies, $Z$ being nonempty, and the vertex reached next has fewer marks: a child of a split vertex is one of at least two molecules, and $B_w\neq\varnothing$. So the descent ends, and it ends exactly where $Z$ is attached.
\end{proof}

\begin{hlemma}[\hb{lem:mot-attach} (The letters of an attached step)]
As stated.
\end{hlemma}

\begin{proof}
For a vertex $w$ on the descent of $Z$ say that \emph{$Z$ is read at $w$} when every letter of the step other than the linear letters is $K$, respectively $\Psi$, of the child at a pattern of $M_w\setminus Z$ of a table at $w$ that carries $K$, respectively $\Psi$. At the root it holds: the step lists of $\cB_j$ at $Z$ are the children $\cB_{j,(T_0\setminus Z)}$ and $\cB_{j,(T_1\setminus Z)}$ at patterns of $\cE\setminus Z$ (Proposition~\hb{prop:mot-chainbottom}), or $\cB_j$ itself when $Z=\cE$, and the tables at the root are the active tables of $F$, $\cB_j$ carrying $K$ exactly when $s_j\neq0$ and $\Psi$ exactly when $t_j\neq0$. Let it hold at $w$, and let the descent pass to a child $w'$.

At a split vertex, $Z\subseteq M'$ for a molecule $M'$, and every table $N$ at $w$ is the direct sum of its restrictions to the molecules (Proposition~\hb{prop:directsum}). So the bases of $N$ with a given trace $p$ on $M_w\setminus Z$ have their other parts fixed, and the child of $N$ at $p$ is the child of $N|M'$ at the pattern $p\cap M'$ of $M'\setminus Z$; $N|M'$ is a table at $w'$ with the exponents of $N$.

At a node that is not a leaf, $Z\subseteq M_w\setminus B_w$. A pattern $p$ of $M_w\setminus Z$ in a table $N$ at $w$ is $p_B\sqcup p_R$ with $p_B=p\cap B_w$ a pattern of $B_w$ and $p_R\subseteq(M_w\setminus B_w)\setminus Z$, and the bases $T$ of $N$ with $T\setminus Z=p$ are those with $T\cap B_w=p_B$ whose parts $T\setminus B_w$, bases of the child $N_{(p_B)}$, have the trace $p_R$ on $(M_w\setminus B_w)\setminus Z$. So the child of $N$ at $p$ is the child of $N_{(p_B)}$ at $p_R$ (Proposition~\hb{prop:twosteps}). $N_{(p_B)}$ is a table made at $w$, carrying the letters $N$ carries with exponents that are positive multiples of those of $N$ (Definition~\hb{def:bunchtree}); it is a table of the room, of the one childcule when the room is a node and of the split vertex otherwise.

So $Z$ is read at the vertex where it is attached, which is the statement.
\end{proof}

\begin{hlemma}[\hb{lem:mot-openall} (One form carries every period)]
As stated.
\end{hlemma}

\begin{proof}
Number the components $D_1,\dots,D_K$ of $D$ and put $D_I=\bigcap_{k\in I}D_k$, $D_\varnothing=V$; each $D_I$ is smooth of dimension $m-|I|$ or empty, and closed in the affine $V$, so affine. The relative de Rham cohomology $H^\bullet_{\mathrm{dR}}(V,D)$ is the hypercohomology of the double complex $\Omega^\bullet_V\to\bigoplus_{|I|=1}\Omega^\bullet_{D_I}\to\bigoplus_{|I|=2}\Omega^\bullet_{D_I}\to\cdots$ with the alternating restrictions (Brown's notes, \S10; Huber and M\"uller-Stach). Its terms are coherent sheaves on the affine $V$, with no higher cohomology, so it is the cohomology of the total complex of global sections, $\operatorname{Tot}^q=\bigoplus_I\Omega^{q-|I|}(D_I)$. Since $\Omega^{q}(D_I)=0$ for $q>m-|I|$, $\operatorname{Tot}^{m+1}=0$ and every $(\eta_I)_I\in\operatorname{Tot}^m$ is a cocycle, $\eta_I$ a form of top degree on $D_I$. Let $\ell$ be the largest $|I|$ with $\eta_I\neq0$, and suppose $\ell\ge1$.

\emph{Restriction is onto.} For $|I'|=\ell-1$, restriction maps $\Omega^{m-\ell}(D_{I'})$ onto $\bigoplus_{k\notin I'}\Omega^{m-\ell}(D_{I'}\cap D_k)$, the top forms on the boundary components of $D_{I'}$. In local coordinates of $D_{I'}$ in which the components $D_{I'}\cap D_k$ through a point are $\{z_1=0\},\dots,\{z_r=0\}$ and $dw$ is the product of the differentials of the other coordinates, the form $\sum_t\widetilde g_t\,dz_1\wedge\cdots\widehat{dz_t}\cdots\wedge dz_r\wedge dw$ restricts to $\pm g_t\,dw$ on $\{z_t=0\}$, $\widetilde g_t$ any local extension of $g_t$, and to $0$ on the others; so the map of coherent sheaves is onto, and so is its map on global sections, $D_{I'}$ being affine.

\emph{The deepest level removed.} For $|I|=\ell$ let $I'=I\setminus\{\max I\}$. Choose $\gamma_{I'}\in\Omega^{m-\ell}(D_{I'})$ for every $|I'|=\ell-1$ restricting on $D_{I'\cup\{k\}}$, $k>\max I'$, to the form that cancels $\eta_{I'\cup\{k\}}$ in the total differential, and to $0$ on the $D_{I'\cup\{k\}}$ with $k<\max I'$. In $(\eta)-d\gamma$ the entry at $|I|=\ell$ is $\eta_I$ less the restrictions of the $\gamma_{I\setminus\{k\}}$, $k\in I$, of which only $k=\max I$ contributes, and it is $0$; the entries above $\ell$ stay $0$, and those at $\ell-1$ change by $d\gamma_{I'}$. By induction every class is the class of $(\eta_\varnothing,0,\dots,0)$, a regular $m$-form on $V$.

For a convergent family, $(V_F,D_F)$ is such a pair (Proposition~\hb{prop:mot-model}(iii)), and $V_{\cE;\varnothing\subsetneq\cE}=V_F$, the letters of the one step $\cE$ being the augmented letters of $F$ that are not monomials, whose $\widetilde Y_f$ are empty; so $L(H_F,\widetilde\sigma_\cE)=\operatorname{Open}(\cE;\varnothing\subsetneq\cE)$. For a step with $|Z|\ge2$ the same holds on $Z$ (proof of Definition~\hb{def:mot-mperiod}); for $|Z|=1$ both sides are $\mathbb Q$.
\end{proof}

\begin{hproposition}[\hb{prop:mot-weights} (Weight-graded faces)]
As stated.
\end{hproposition}

\begin{proof}
For a chain $\mathfrak c$ of nonempty proper sets let $D_{\mathfrak c}$ be its stratum in $V_F$, a smooth affine variety of dimension $n-1-|\mathfrak c|$, nonempty since it meets $\widetilde\sigma_\cE$ (Proposition~\hb{prop:mot-model}(iii),(iv)), with its boundary $\partial D_{\mathfrak c}$, its meeting with the other components of $D_F$. The summands of the double complex of the proof of Lemma~\hb{lem:mot-openall} over the chains refining $\mathfrak c$ form the double complex of $(D_{\mathfrak c},\partial D_{\mathfrak c})$, shifted by $|\mathfrak c|$, and their inclusion is the \emph{face map} $\phi_{\mathfrak c}:H^{n-1-|\mathfrak c|}(D_{\mathfrak c},\partial D_{\mathfrak c})\to H^{n-1}(V_F,D_F)$, a morphism of Brown's category whose transpose takes a relative cycle that is a manifold with corners, with its faces on the components of $D_F$, to its face on $D_{\mathfrak c}$, with a sign (Brown's notes, \S10). Here $\widetilde\sigma_\cE$ is such a cycle, $[0,\infty)^{n-1}$ in every chart, and $\phi_{\mathfrak c}^{\mathsf t}\widetilde\sigma_\cE=\pm\widetilde\sigma_\cE\cap D_{\mathfrak c}$. By Brown's weight-graded face maps (Brown's notes, \S10), for $k<n-1$ the sum of the face maps of the chains with $n-1-|\mathfrak c|=k$ maps $\bigoplus_{\mathfrak c}W_kH^k(D_{\mathfrak c},\partial D_{\mathfrak c})$ onto $W_kH^{n-1}(V_F,D_F)$; the property it needs is that every stratum is affine, so that its cohomology vanishes above its dimension. Write $e=\sum_{\mathfrak c}\phi_{\mathfrak c}(e_{\mathfrak c})$. By Lemma~\hb{lem:mot-openall} on the pair $(D_{\mathfrak c},\partial D_{\mathfrak c})$, $e_{\mathfrak c}=[\eta_{\mathfrak c}]$ for a regular $k$-form on $D_{\mathfrak c}$, and by the relation for $\phi_{\mathfrak c}$
\[
[H_F,\widetilde\sigma_\cE,e]^{\mathrm{mot}}=\sum_{\mathfrak c}\pm\big[H^k(D_{\mathfrak c},\partial D_{\mathfrak c}),\widetilde\sigma_\cE\cap D_{\mathfrak c},[\eta_{\mathfrak c}]\big]^{\mathrm{mot}} .
\]
The chain $\mathfrak c$ is a flag $\mathfrak f$ of dimension $k$, and by Proposition~\hb{prop:mot-model}(iv) its stratum with its boundary is the product of the pairs $(V_{Z_i;\mathfrak f},\partial V_{Z_i;\mathfrak f})$ over its steps, with $\widetilde\sigma_\cE\cap D_{\mathfrak c}=\prod_i\widetilde\sigma_{Z_i}$. The factors are affine, so the regular functions on the product are the tensor product of those on the factors, the top forms likewise, and $\eta_{\mathfrak c}$ is a finite sum of products of top forms on the factors. By Lemma~\hb{lem:motive-product} each term is a product over the steps of elements of $\operatorname{Open}(Z_i;\mathfrak f)$, a step of one mark giving a rational number; the signs are absorbed, $\operatorname{Open}(Z;\mathfrak f)$ being closed under negation.
\end{proof}

\begin{htheorem}[\hb{thm:mot-smallsteps} (Small steps along the tree)]
As stated.
\end{htheorem}

\begin{proof}
Proposition~\hb{prop:mot-weights} at $e=e_i$ is the inclusion. A flag of dimension $k$ has $\sum_i(|Z_i|-1)=k$, so each of its steps has at most $k+1$ marks. Each step is attached to exactly one vertex (Definition~\hb{def:mot-attach}), and its letters are those of Lemma~\hb{lem:mot-attach}.
\end{proof}

\begin{hproposition}[\hb{prop:mot-convergentdata} (Laurent coefficients are read on convergent data)]
As stated.
\end{hproposition}

\begin{proof}
\emph{Meromorphy.} The line $\mvec+\varepsilon\mathbf v$ lies in the balance wall, $v_\cE=0$, and $\kappa_S(\mvec+\varepsilon\mathbf v)=\kappa_S+\varepsilon v_S$ with $v_S\neq0$, so the line meets every hyperplane $\kappa_S=-j$ in at most one point, and the continued period is meromorphic on it (Proposition~\hb{prop:mot-model}(v)).

\emph{A relation at each wall.} For lists every weight is $1$. Let $\varnothing\neq S\subsetneq\cE$. Summing the mark relations of Theorem~\hb{thm:confJ}(iii) over $e\in S$, and subtracting $s_jd_j(S)$ times the first letter relation of table $j$ and $t_jd^\ast_j(S)$ times the second, for every $j$, gives, as meromorphic functions of the data,
\begin{multline*}
\kappa_S\,P_F(\mvec;\mathbf s,\mathbf t)=\kappa_\cE\sum_{e\in S}P_F(\mvec+\mathbf e_e;\mathbf s,\mathbf t)+\sum_js_j\sum_{T\in\cB_j}\big(|T\cap S|-d_j(S)\big)P_F(\mvec+\one_T;\mathbf s+\mathbf e_j,\mathbf t)\\
+\sum_jt_j\sum_{T\in\cB_j}\big(|S\setminus T|-d^\ast_j(S)\big)P_F(\mvec+\one_{\cE\setminus T};\mathbf s,\mathbf t+\mathbf e_j).
\end{multline*}
Each term has its poles on the hyperplanes $\kappa_{S'}=-j$, $S'$ proper, of its own data (Lemma~\hb{lem:motive-continued}), none of which contains the balance wall; so on the balance wall the first term vanishes, and the data of the other terms stay on it, $|T|=r_j$ and $|\cE\setminus T|=r^\ast_j$. On the line $\mvec+\varepsilon\mathbf v$ every remaining term is meromorphic, as above, and
\begin{multline*}
P_F(\mvec+\varepsilon\mathbf v;\mathbf s,\mathbf t)=\frac1{\kappa_S+\varepsilon v_S}\Big(\sum_js_j\sum_T\big(|T\cap S|-d_j(S)\big)P_F(\mvec+\one_T+\varepsilon\mathbf v;\mathbf s+\mathbf e_j,\mathbf t)\\
+\sum_jt_j\sum_T\big(|S\setminus T|-d^\ast_j(S)\big)P_F(\mvec+\one_{\cE\setminus T}+\varepsilon\mathbf v;\mathbf s,\mathbf t+\mathbf e_j)\Big),
\end{multline*}
with $1/(\kappa_S+\varepsilon v_S)\in\mathbb Q(\varepsilon)$ and integer coefficients.

\emph{The descent.} Put $\delta(F)=\sum_{\varnothing\neq S\subsetneq\cE}\max(0,1-\kappa_S)$, a nonnegative integer, zero exactly when $F$ is convergent. If $\delta(F)>0$, choose $S$ with $\kappa_S\le0$ and apply the relation at $S$, keeping the terms with nonzero coefficients. Their families are on the same tables, at integer data on their balance walls, with the exponents larger by one unit vector; their walls are $\kappa_{S'}+|T\cap S'|-d_j(S')\ge\kappa_{S'}$, respectively $\kappa_{S'}+|S'\setminus T|-d^\ast_j(S')\ge\kappa_{S'}$, for every $S'$ (Theorem~\hb{thm:window}), and at $S$ itself, the coefficient being nonzero, at least $\kappa_S+1$. As $1-\kappa_S\ge1$, the term of $S$ in $\delta$ drops by one, and $\delta$ drops by at least one. Repeating on every family with $\delta>0$ ends after finitely many steps, with convergent families $F_k$ and coefficients $R_k\in\mathbb Q(\varepsilon)$, products of the factors above. The exponents only grow, so with $\mathbf s,\mathbf t\ge0$ every $F_k$ has nonnegative exponents, and then positive masses, $m^{(k)}_e=\kappa_{\{e\}}(F_k)+\sum_j\big(s^{(k)}_jd_j(\{e\})+t^{(k)}_jd^\ast_j(\{e\})\big)\ge1$, the singletons being proper as $n\ge2$.

\emph{The coefficients.} At the masses $\mvec^{(k)}+\varepsilon\mathbf v$ the walls of $F_k$ have positive real parts for $|\varepsilon|$ small, so its period is the integral over $\Delta_\cE$, converging absolutely and locally uniformly, holomorphic in $\varepsilon$, with the derivatives under the integral sign, $u^{\varepsilon\mathbf v}=\exp(\varepsilon\ip{\mathbf v}{\log u})$. The displayed integrals converge absolutely: on $\Delta_\cE$, $|\ip{\mathbf v}{\log u}|^l\le C\prod_eu_e^{-\eta}$ for every $\eta>0$, and the walls at $\mvec^{(k)}-\eta\one$ are positive for small $\eta$, so that the integral of $u^{\mvec^{(k)}-\eta\one-\one}\prod_jK_j^{-s^{(k)}_j}\Psi_j^{-t^{(k)}_j}$ converges (Theorem~\hb{thm:projective}, which reads the integral over $\Delta_\cE$ off the balance wall as well). The Laurent coefficients of $\sum_kR_kP_{F_k}$ at $0$ are rational combinations of the Taylor coefficients of the $P_{F_k}$, those of the $R_k$ being rational.
\end{proof}

\begin{hcorollary}[\hb{cor:mot-openfamilies} (Open periods from families of lists)]
As stated.
\end{hcorollary}

\begin{proof}
\emph{The forms.} The torus part of $V_{Z;\mathfrak f}$ is the torus less the zero sets of the letters of the step, a principal open subset of the torus, so a regular function on $V_{Z;\mathfrak f}$ is there a homogeneous Laurent polynomial divided by a product of powers of the letters, of total degree zero; and $\Omega_Z/\alpha^Z$ is a nowhere vanishing form of top degree on the torus. Splitting the Laurent polynomial into its monomials writes $\eta$ as a finite sum of terms $c_t\,\alpha^{\mathbf a_t}\prod_gg^{-k_{t,g}}\,\Omega_Z/\alpha^Z$, $k_{t,g}\in\mathbb N$, each the form $\omega_{F_t}$ of the family of the statement, a letter $K$ of a step list being $K$ of the table $(N;k_{t,g},0)$, a letter $\Psi$ being $\Psi$ of the table $(N;0,k_{t,g})$, and $\alpha_J$ being $K$ of $(\binom J1;k_{t,g},0)$ (Definition~\hb{def:mot-model}); $\alpha^{\mathbf a_t}\prod_gg^{-k_{t,g}}$ has degree zero, so $F_t$ is on its balance wall.

\emph{The period.} For $\mathbf x\in\mathbb C^Z$ with $x_Z=0$ near $0$ put $G(\mathbf x)=\int_{\widetilde\sigma_Z}\alpha^{\mathbf x}\eta$. In the chart of an order $\pi$, $\alpha^{\mathbf x}=\prod_i\varepsilon_i^{x_{S_i}}$ (proof of Lemma~\hb{lem:mot-charts}(ii)), and $\eta$, regular near $\widetilde\sigma_Z\subseteq V_{Z;\mathfrak f}(\mathbb R)$ (proof of Definition~\hb{def:mot-mperiod}), integrates over the sector as $g_\pi\,d\varepsilon$ on $[0,1]^{|Z|-1}$ with $g_\pi$ smooth near the cube. So $G(\mathbf x)=\sum_\pi\operatorname{cube}(\mathbf x_\pi+\one;g_\pi)$, $\mathbf x_\pi=(x_{S^\pi_i})_i$, holomorphic near $\mathbf x=0$, and $G(0)=\int_{\widetilde\sigma_Z}\eta$ (Lemma~\hb{lem:motive-sections}). On the open cube Lemma~\hb{lem:mot-charts}(iii) gives, with the same sign for every $t$,
\[
g_\pi=\sum_tc_t\,\varepsilon^{\boldsymbol\kappa_\pi(F_t)-\one}h_{t,\pi},\qquad h_{t,\pi}=\prod_g\widetilde g^{\,-k_{t,g}},
\]
$\boldsymbol\kappa_\pi(F_t)$ the walls of $F_t$ at $\mathbf a_t$ along the chain of $\pi$. Choose $A\in\mathbb N$ with $\boldsymbol\kappa_\pi(F_t)+A\one\ge\one$ for every $t$ and $\pi$. The cube integrals have $\operatorname{cube}(\mathbf y+\mathbf b;h)=\operatorname{cube}(\mathbf y;\varepsilon^{\mathbf b}h)$ for $\mathbf b\in\mathbb N^{|Z|-1}$, where both converge and so, both sides being meromorphic (Lemma~\hb{lem:cube}), everywhere; they are linear in $h$. So, as meromorphic functions of $\mathbf x$,
\begin{align*}
\sum_tc_t\operatorname{cube}\big(\mathbf x_\pi+\boldsymbol\kappa_\pi(F_t);h_{t,\pi}\big)&=\sum_tc_t\operatorname{cube}\big(\mathbf x_\pi+(1-A)\one;\varepsilon^{\boldsymbol\kappa_\pi(F_t)+(A-1)\one}h_{t,\pi}\big)\\
&=\operatorname{cube}\big(\mathbf x_\pi+(1-A)\one;\varepsilon^{A\one}g_\pi\big)=\operatorname{cube}(\mathbf x_\pi+\one;g_\pi),
\end{align*}
the functions $\varepsilon^{\boldsymbol\kappa_\pi(F_t)+(A-1)\one}h_{t,\pi}$ being smooth near the cube. Summing over $\pi$ and using Proposition~\hb{prop:mot-model}(v) with $\kappa_S(\mathbf a_t+\mathbf x)=\kappa_S(\mathbf a_t)+x_S$, $G(\mathbf x)=\sum_tc_tP_{F_t}(\mathbf a_t+\mathbf x)$ as meromorphic functions near $\mathbf x=0$ on $\{x_Z=0\}$. On the line $\mathbf x=\varepsilon\mathbf v$, every $v_S$ being nonzero, each term is meromorphic in $\varepsilon$, as in the proof of Proposition~\hb{prop:mot-convergentdata}, the sum is $G(\varepsilon\mathbf v)$, holomorphic at $0$ with the value $\int_{\widetilde\sigma_Z}\eta$, and the constant term is linear.

\emph{The rest.} Every $F_t$ has nonnegative exponents, so Proposition~\hb{prop:mot-convergentdata} writes its constant term through convergent families of lists on $Z$ with nonnegative integer exponents and positive integer masses, whose periods are values of the serif family (Proposition~\hb{prop:gammadraws}(i)), and their derivatives along $\mathbf v$. The period of $[H^{|Z|-1}(V_{Z;\mathfrak f},\partial V_{Z;\mathfrak f}),\widetilde\sigma_Z,\eta]^{\mathrm{mot}}$ is $\int_{\widetilde\sigma_Z}\eta$, and taking periods is a homomorphism of rings; Theorem~\hb{thm:mot-smallsteps} gives the last sentence.
\end{proof}

\begin{htheorem}[\hb{thm:mot-node} (The node form)]
As stated.
\end{htheorem}

\begin{proof}
(i) A letter $K$ of the step is $K$ of a step list, which by Lemma~\hb{lem:mot-attach} is the child $N'$ of a table $N$ at $v$ at a pattern $p$ of $M_v\setminus Z$, a list being determined by its basis polynomial; a letter $\Psi$ is $\Psi$ of such a child $N'$, that is $K$ of $N'^\ast$ (Proposition~\hb{prop:identity}); a letter $\alpha_J$ is $K$ of $\binom J1$ read on $Z$, its marks outside $J$ loops. Lemma~\hb{lem:mot-childwidth}, for the table $N$ on $M_v$ with the bunch $M_v\setminus Z$, its room $Z$, the pattern $p$, $A=Z_B$ and $W=B_v\setminus Z$, so that $A\cup W=B_v$, gives $c_{N'}(Z_B)\le c_N(B_v)\le c$; the dual has the same widths (Theorem~\hb{thm:window}). $\binom J1$ has rank at most one, so the width of $Z_B$ in it is at most one. Finally $c\ge1$: if every table at $v$ had $c(B_v)=0$, $B_v$ would be a union of molecules of every table (Lemma~\hb{lem:separator}), no circuit of a table at $v$ would meet both $B_v$ and $M_v\setminus B_v$, and $B_v$ would be a union of molecules of the tables at $v$ (Definition~\hb{def:molecule}), against $\varnothing\neq B_v\subsetneq M_v$ and the tables at the node $v$ forming one molecule.

(ii) For $g=K_{N_g}$ and the gauge $(u,u')$, $g(\rho u,u')=\sum_k\rho^kK^{(k)}_{Z_B}(u,u')$ over the levels $d_g\le k\le d_g+c_g$ (Definition~\hb{def:levels}, Theorem~\hb{thm:window}). The bases of the top level are the unions of a basis of $N_g|Z_B$ and one of $N_g/Z_B$, each of weight $1$ (Theorem~\hb{thm:levels}), so the top level letter is $K_{N_g|Z_B}(u)\,K_{N_g/Z_B}(u')$. The Skeleton polynomial $\sQ^K_{Z_B}(x;(u,u'))=x^{-d_g}g(xu,u')/g(u,u')=\prod_i(x+\xi_{g,i})/(1+\xi_{g,i})$ (Definition~\hb{def:skelpoly}) shows that $\rho^{-d_g}g(\rho u,u')$ is a polynomial of degree $c_g$ in $\rho$ with the roots $-\xi_{g,i}$, and comparing leading coefficients gives the display. The Skeleton values of $K_{N''^\ast}$ along $Z_B$ are those of $\Psi_{N''}$ (Definition~\hb{def:skelpoly}), so the union over the letters is the joint spectrum of $Z_B$ in the sense of Definition~\hb{def:pairsk} and Theorem~\hb{thm:scale}.

(iii) By Corollary~\hb{cor:mot-openfamilies}, on the torus $\eta=\phi\,\Omega_Z/\alpha^Z$ with $\phi=\sum_tc_t\alpha^{\mathbf a_t}\prod_gg^{-k_{t,g}}$, homogeneous of degree zero. By Lemma~\hb{lem:motive-sections} for the partition $Z=Z_B\sqcup Z_R$, $\int_{\widetilde\sigma_Z}\eta$ is the absolutely convergent integral over $\Delta_{Z_B}\times\Delta_{Z_R}\times(0,\infty)$ of $\phi(\rho u,u')\rho^{-1}u^{-\one}u'^{-\one}$. By (ii) the term $t$ of this integrand is
\[
A_t(u)\,B_t(u')\,\rho^{\nu_t-1}\prod_{g,i}\big(\rho+\xi_{g,i}\big)^{-k_{t,g}},
\]
with $A_t=c_t\,u^{\mathbf a_t|_{Z_B}-\one}\prod_gK_{N_g|Z_B}^{-k_{t,g}}$, $B_t=u'^{\,\mathbf a_t|_{Z_R}-\one}\prod_gK_{N_g/Z_B}^{-k_{t,g}}$ and $\nu_t=\sum_{e\in Z_B}a_{t,e}-\sum_gk_{t,g}d_g$. By Tonelli, for almost every $(u,u')$ the integral over $\rho$ of the sum converges absolutely, and the integral over $\Delta_{Z_B}\times\Delta_{Z_R}$ of it is the period, with an integrable modulus. At such a point the integral over $\rho$ is its own finite part (Proposition~\hb{prop:Regconv}): each term is $\rho^{\nu_t-1}$ times a rational function of $\rho$ without poles on $(0,\infty)$, the $\xi_{g,i}$ lying off $(-\infty,0]$, whose primitive, a sum of rational functions and of logarithms $\log\rho$ and $\log(\rho+\xi_{g,i})$, has log-power expansions at both ends. Finite parts are linear, so the integral over $\rho$ is $\sum_tA_tB_t\,\mathrm{Sc}_t$. A term in which no letter with $k_{t,g}>0$ has a Skeleton value integrates $\rho^{\nu_t-1}$, whose primitive $\rho^{\nu_t}/\nu_t$ or $\log\rho$ has no constant term at either end, so $\mathrm{Sc}_t=0$. Otherwise, by Theorem~\hb{thm:RegCT}, $\mathrm{Sc}_t$ is the constant term at $\kappa=\nu_t$ of the integral $\int_0^\infty\rho^{\kappa-1}\prod_{g,i}(\rho+\xi_{g,i})^{-k_{t,g}}d\rho$, continued from the strip $0<\re\kappa<|\mathbf w_t|$, where it converges, the product having degree $|\mathbf w_t|=\sum_gk_{t,g}c_g$ in $\rho$; on the strip it is the stacked Gamma times the Breta of the statement (Proposition~\hb{prop:scalebreta}, coincident Skeleton values merging), and both sides are meromorphic in $\kappa$, so their constant terms at $\nu_t$ agree.

(iv) At $c=1$ every $c_g$ is at most one by (i). A letter with $c_g=1$ has the one Skeleton value $\gamma^g_R(u')/\gamma^g_B(u)$, positive (Corollary~\hb{cor:widthonestage}, the weights $1$ factoring). Fix $\kappa$ in the strip. Corollary~\hb{cor:sharepercoefficient}, one coefficient $\xi_g$ for each letter of the term with a Skeleton value, writes the stacked Gamma times the Breta as an integral over $|J_t|-1$ shares $a_\ell$ on the lines $\re a_\ell=c_\ell$ of a product of Beta corners in the shares and of the powers $\xi_g^{-b_g}$, $b_\ell=w_{t,\ell}-a_\ell$ for the shared coefficients and $b_0=w_{t,0}-\kappa+|\mathbf a|$ for the last, affine in the shares and in $\kappa$, with $\re b_g$ constant on the contour. The Beta corners are absolutely integrable over the contour: those in one share are bounded by a constant times $e^{-\pi|\im a_\ell|}(1+|\im a_\ell|)^{w_{t,\ell}-1}$ by the classical bound $|\Gamma(c+i\tau)|\le C_c\,e^{-\pi|\tau|/2}(1+|\tau|)^{c-1/2}$, and the last one by its value at the real parts of its entries, since $|\Gamma(z)|\le\Gamma(\re z)$ for $\re z>0$. As $|\xi_g^{-b_g}|=\gamma^g_R(u')^{-\re b_g}\gamma^g_B(u)^{\re b_g}$, the integral of the modulus of $A_tB_t$ times the share integrand over the contour, $\Delta_{Z_B}$ and $\Delta_{Z_R}$ is the integral of the modulus of the share kernel times the two integrals of the statement, finite; by Fubini the integrals over $\Delta_{Z_B}$ and $\Delta_{Z_R}$ move inside the share integral, where at fixed shares the integrand is a function of $u$ times a function of $u'$, and they separate. If $0<\nu_t<|\mathbf w_t|$, the integral over $\rho$ of the term converges at $\kappa=\nu_t$ and is $\mathrm{Sc}_t$ (Proposition~\hb{prop:Regconv}), and the argument runs at $\kappa=\nu_t$.

(v) By Lemma~\hb{lem:mot-openall}, every element of $L(H_F,\widetilde\sigma_\cE)$ is $[H_F,\widetilde\sigma_\cE,[\eta]]^{\mathrm{mot}}$ for a regular form $\eta$ on $V_F=V_{\cE;\varnothing\subsetneq\cE}$. The step $\cE$ meets the bunch of the root and its complement, so it is attached to the root as a node step, and (iii) applies to $\eta$. The left factors are elements of $L(H_F,\widetilde\sigma_\cE)$ by Definition~\hb{def:mot-mperiod}.

(vi) At a split vertex the tables are direct sums over the molecules (Proposition~\hb{prop:directsum}), so the children at patterns of $M_v\setminus Z$ are direct sums over the parts of $Z$ in the molecules (proof of Lemma~\hb{lem:mot-attach}), and $Z_B=Z\cap M_1$ is a union of molecules of each of them, of width zero (Lemma~\hb{lem:separator}); the same holds for the duals. The list $\binom J1$ has width one at $Z_B$ exactly when it has bases on both sides, that is when $J$ meets $Z_B$ and $Z_R$, and width zero otherwise. The proofs of (ii) to (iv) use only the widths of the letters at $Z_B$, here at most one.
\end{proof}

\begin{slemma}[Families on two marks]\label{lem:motive-twomarks}
Every convergent family $F$ on a set $\cE$ of two marks (Definition~\hb{def:mot-model}) has $H_F\cong\mathbb Q(0)$, and $P^{\mathrm{mot}}_F=P_F$, the Beta corner $\Gst21{\kappa_{\{1\}},\ \kappa_{\{2\}}}{\kappa_{\{1\}}+\kappa_{\{2\}}}$ of its walls, a positive rational number.
\end{slemma}

\begin{proof}
The lists on $\cE=\{1,2\}$ have the basis polynomials $1$, $\alpha_1$, $\alpha_2$, $\alpha_1\alpha_2$ and $\alpha_1+\alpha_2$, and their cobasis polynomials are among these. $X_\cE$ is the projective line, its boundary the two points $\alpha_1=0$ and $\alpha_2=0$, the monomials have no zeros on the torus, and the linear letter $\alpha_1+\alpha_2$ is always removed; so $V_F$ is the projective line less one point, an affine line, and $D_F$ two points of it. In the long exact sequence of the pair, a sequence of morphisms of Brown's category (Huber and M\"uller-Stach), $H^0(V_F)\to H^0(D_F)\to H^1(V_F,D_F)\to H^1(V_F)$, the first map is the diagonal $\mathbb Q(0)\to\mathbb Q(0)^2$ and $H^1$ of the affine line is $0$, every form $f(x)\,dx$ being exact; so $H_F\cong\mathbb Q(0)$. By the relation of Definition~\hb{def:mot-mperiod} for this isomorphism every motivic period of $H_F$ is a rational multiple of the unit $[\mathbb Q(0),1,1]^{\mathrm{mot}}=1$, the multiple being its period. On $\Delta_\cE$ the letter $\alpha_1+\alpha_2$ is $1$ and the monomials shift the masses to the walls, so $P_F=\int_0^1u^{\kappa_{\{1\}}-1}(1-u)^{\kappa_{\{2\}}-1}du$, with $\kappa_{\{1\}},\kappa_{\{2\}}\ge1$ integers.
\end{proof}

\begin{htheorem}[\hb{thm:mot-pinned} (Pinned cells split off motivically)]
As stated.
\end{htheorem}

\begin{proof}
Write $\mathbf h=\{h_1,\dots,h_q\}$ and $r_{ij}=r_{\mathcal A_{ij}}$, so that $K_{\mathcal A_{ij}/g_i}$ and $\Psi_{\mathcal A_{ij}\setminus g_i}$ have the degrees $r_{ij}-1$ and $|R_i|-r_{ij}$; the core $\mathcal H_j$ has at least two marks, as no mark of a table on one mark is neither a loop nor a coloop, so $\cE\setminus R_i\neq\varnothing$.

(v) is Corollary~\hb{cor:mot-widthbound}.

\emph{The walls of the cells.} $R_i$ has width one in every active table (Proposition~\hb{prop:substitution}(i)), so on the balance wall $\kappa_{R_i}+\kappa_{\cE\setminus R_i}=\sum_j(s_j+t_j)=2\sigma$ (Proposition~\hb{prop:compoundwall}); both are at least $1$, $F$ being convergent, so $1\le\kappa_i\le2\sigma-1$.

\emph{The letters.} By Theorem~\hb{thm:substletters} and the hypothesis every active letter reads the core at one point:
\begin{gather*}
K_j=\prod_iK_{\mathcal A_{ij}/g_i}(\alpha_{R_i})\,K_{\mathcal H_j}(\mathbf x),\qquad\Psi_j=\prod_i\Psi_{\mathcal A_{ij}\setminus g_i}(\alpha_{R_i})\,\Psi_{\mathcal H_j}(\mathbf x),\\
\mathbf x=\big(\alpha_L,\gamma_1(\alpha_{R_1}),\dots,\gamma_q(\alpha_{R_q})\big),
\end{gather*}
a mark of $L$ standing at its own variable, both hinge values of $\binom{\{l,g\}}1$ being $\alpha_l$.

\emph{The map.} Each $\gamma_i$ is a quotient of polynomials, homogeneous of degree one, so
\[
\Phi:\ \mathbb P(\mathbb Q^\cE)\dashrightarrow\mathbb P(\mathbb Q^{L\sqcup\mathbf h})\times\prod_i\mathbb P(\mathbb Q^{R_i}),\qquad[\alpha]\mapsto\big([\mathbf x],[\alpha_{R_1}],\dots,[\alpha_{R_q}]\big),
\]
is well defined, between varieties of the same dimension $|\cE|-1$; off the zero sets of the numerators and denominators of the $\gamma_i$ it is an isomorphism between open subsets of the tori, with the inverse $\alpha_L=x_L$, $\alpha_{R_i}=\big(x_{h_i}/\gamma_i(v_i)\big)v_i$ at $([\mathbf x],[v_1],\dots,[v_q])$, well defined by the same homogeneity. Let $U_1$ be the torus less these zero sets, those of the augmented letters of $F$, and the preimages under $\Phi$ of those of the augmented letters of the core family and of the pinned families, and $U_2=\Phi(U_1)$. All these polynomials are positive at the positive points, so $\Phi$ maps $\widetilde\sigma^\circ_\cE$ onto $\widetilde\sigma^\circ_{L\sqcup\mathbf h}\times\prod_i\widetilde\sigma^\circ_{R_i}$.

\emph{The forms.} On the positive points write $\alpha_{R_i}=\lambda_iv_i$ with $\lambda_i>0$ and $v_i$ on a section of the rays of $\mathbb R^{R_i}_{>0}$, so that $x_{h_i}=\lambda_i\gamma_i(v_i)$. In logarithmic coordinates $\Phi$ keeps $\log\alpha_L$ and the $v_i$ and adds $\log\gamma_i(v_i)$ to $\log\lambda_i$, a map with Jacobian determinant $1$; with $\nu_X=\Omega_X/\alpha^X$, the form $\Omega_X/\alpha^X$ being the invariant volume of the logarithmic coordinates modulo $\one$ (proof of Lemma~\hb{lem:motive-sections}), $\Phi^\ast\big(\nu_{L\sqcup\mathbf h}\wedge\bigwedge_i\nu_{R_i}\big)=\pm\nu_\cE$, an identity of rational forms since it holds at the Zariski dense positive points. The core family $F_{\mathcal H}$ is on its balance wall: by Proposition~\hb{prop:substitution}(i), $r_j=r_{\mathcal H_j}+\sum_i(r_{ij}-1)$, and
\[
m_L+\sum_i\kappa_i=m_\cE-\sum_{i,j}\big(s_j(r_{ij}-1)+t_j(|R_i|-r_{ij})\big)=\sum_j\Big(s_jr_{\mathcal H_j}+t_j\big(|L|+q-r_{\mathcal H_j}\big)\Big),
\]
with $\kappa_i=m_{R_i}-\sum_j\big(s_j(r_{ij}-1)+t_j(|R_i|-r_{ij})\big)$ (Theorem~\hb{thm:substwalls}). The word of $G_i$ is that of Corollary~\hb{cor:mot-loadpinned}, $\gamma_i^{-\kappa_i}\prod_jK_{\mathcal A_{ij}/g_i}^{-s_j}\Psi_{\mathcal A_{ij}\setminus g_i}^{-t_j}$, homogeneous of the degree $-m_{R_i}$, so $G_i$ is on its balance wall. Pulling back $\omega_{F_{\mathcal H}}$ and $\omega_{G_i}$ by $\Phi$ substitutes $\mathbf x$ and $v_i=\alpha_{R_i}$ in their coefficients, and the powers $\gamma_i(\alpha_{R_i})^{\pm\kappa_i}$ cancel:
\[
\Phi^\ast\big(\omega_{F_{\mathcal H}}\boxtimes\omega_{G_1}\boxtimes\cdots\boxtimes\omega_{G_q}\big)=\pm\,\alpha^{\mvec}\prod_jK_j^{-s_j}\Psi_j^{-t_j}\,\nu_\cE=\pm\,\omega_F .
\]

(i) By the change of variables the period of $F$, of the positive integrand, is the product of the periods of $F_{\mathcal H}$ and of the $G_i$ in $(0,\infty]$; this is Corollary~\hb{cor:fullleancore} with Corollary~\hb{cor:mot-loadpinned}. $P_F$ is finite, so every factor is finite, and at integer data on its balance wall every proper wall of each factor is then at least $1$ (Theorem~\hb{thm:projective}): $F_{\mathcal H}$ and the $G_i$ are convergent, their tables being lists.

(ii) Apply Lemma~\hb{lem:mot-birational} to $(X_\cE,D,Y_F)$ and to the product of the toric varieties of $F_{\mathcal H}$ and of the $G_i$ with the union of the pull-backs of their boundaries and of their removed sets, a smooth projective variety with a divisor with simple normal crossings and smooth components, with the cells $\widetilde\sigma_\cE$ and $\widetilde\sigma_{L\sqcup\mathbf h}\times\prod_i\widetilde\sigma_{R_i}$ and the forms $\omega_F$ and $\omega_{F_{\mathcal H}}\boxtimes\omega_{G_1}\boxtimes\cdots\boxtimes\omega_{G_q}$, regular on $V_{F_{\mathcal H}}\times\prod_iV_{G_i}$ by (i); the two periods are positive, so the signs of the orientations and of the forms match as in the proof of Corollary~\hb{cor:mot-cremona}. With Lemma~\hb{lem:motive-product},
\[
P^{\mathrm{mot}}_F=\Big[H^{n-1}\Big(V_{F_{\mathcal H}}\times\prod_iV_{G_i},\partial\Big),\widetilde\sigma_{L\sqcup\mathbf h}\times\prod_i\widetilde\sigma_{R_i},\omega_{F_{\mathcal H}}\boxtimes\cdots\boxtimes\omega_{G_q}\Big]^{\mathrm{mot}}=P^{\mathrm{mot}}_{F_{\mathcal H}}\prod_iP^{\mathrm{mot}}_{G_i},
\]
and Lemma~\hb{lem:mot-conj}(iii) gives the conjugates.

(iii) Each cell is itself a substitution: $\mathcal A_{ij}=\binom{\{x,h\}}1\langle\binom{\{g_i,y\}}1,\mathcal A'_{ij}\rangle$, the mark $g_i$ left in place at $x$ and $\mathcal A'_{ij}$, a copy of $\mathcal A_{ij}$ with its hinge renamed, at $h$: the basis $\{h\}$ of the core contributes the bases of $\mathcal A_{ij}\setminus g_i$ and the basis $\{x\}$ the sets $Y+g_i$, $Y$ a basis of $\mathcal A_{ij}/g_i$ (Definition~\hb{def:substitution}). The hinge values of $\mathcal A'_{ij}$ are those of $\mathcal A_{ij}$, so the cell family $F_i$ is read through the one hinge value $\gamma_i$, its wall at $R_i$ is $\kappa_i$ (Theorem~\hb{thm:substwalls}, with the masses $\mvec|_{R_i}$), and it is on its balance wall with the mass $2\sigma-\kappa_i$ on $g_i$ (Lemma~\hb{lem:twosumpieces}). Its core family $F'_i$ has the tables $\binom{\{g_i,h\}}1$, whose two letters are both $\alpha_{g_i}+\alpha_h$, with the masses $2\sigma-\kappa_i$ on $g_i$ and $\kappa_i$ on $h$, and its pinned family at the same loads is $G_i$. By the change of variables as in (i), in $(0,\infty]$, $P_{F_i}=P_{F'_i}P_{G_i}$ with $P_{F'_i}=\int_0^1u^{2\sigma-\kappa_i-1}(1-u)^{\kappa_i-1}du=\Gst21{\kappa_i,\ 2\sigma-\kappa_i}{2\sigma}$, finite since $1\le\kappa_i\le2\sigma-1$. So $P_{F_i}$ is finite, $F_i$ is convergent (Theorem~\hb{thm:projective}), and (ii) for $F_i$ gives $P^{\mathrm{mot}}_{F_i}=P^{\mathrm{mot}}_{F'_i}P^{\mathrm{mot}}_{G_i}$, with $P^{\mathrm{mot}}_{F'_i}$ the Beta corner by Lemma~\hb{lem:motive-twomarks}.

(iv) The form of $G_i$ was computed above without reference to the loads; its period is the integral over $\Delta_{R_i}$ of its coefficient, the pinned shape integral of Theorem~\hb{thm:pinned} with all its constants $1$; and two choices of loads give two convergent families on $R_i$ with one form, hence one motivic period (Lemma~\hb{lem:mot-birational}).
\end{proof}

\begin{hcorollary}[\hb{cor:mot-dualfull} (Dual pairs and full leans)]
As stated.
\end{hcorollary}

\begin{proof}
(a) By Proposition~\hb{prop:mot-dualletters}(ii) each cell has $\gamma^K_{\mathcal A_i}=\gamma^\Psi_{\mathcal A_i}=k_i/\psi_i$, so the one active table reads it through one hinge value whatever $s$ and $t$, and Theorem~\hb{thm:mot-pinned} applies. With $K_{\mathcal A_i/g_i}=\Psi_{\mathcal A_i\setminus g_i}=\psi_i$, the form of the pinned family (Theorem~\hb{thm:mot-pinned}(iv)) is $(k_i/\psi_i)^{-\kappa_i}\psi_i^{-s}\psi_i^{-t}\alpha^{\mvec|_{R_i}}\nu_{R_i}=k_i^{-\kappa_i}\psi_i^{-(s+t-\kappa_i)}\alpha^{\mvec|_{R_i}}\nu_{R_i}$, the form of $(\mathcal A_i\setminus g_i;\kappa_i,s+t-\kappa_i)$ at the masses $\mvec|_{R_i}$, and $\mathcal A_i\setminus g_i=\cB|R_i$, the cell of the bunch $R_i$ (Proposition~\hb{prop:substitution}(i), Theorem~\hb{thm:widthone}). By Proposition~\hb{prop:mot-dualletters}(i), $|R_i|-r_{\mathcal A_i}=r_{\mathcal A_i}-1$ with $r_{\mathcal A_i}=r(R_i)$, so Theorem~\hb{thm:substwalls} gives $\kappa_i=m_{R_i}-(s+t)(r(R_i)-1)$. The family $(\cB|R_i;\kappa_i,s+t-\kappa_i)$ at the masses $\mvec|_{R_i}$ has the form of the convergent pinned family, which is homogeneous of degree zero, so it is on its balance wall; its walls are the orders of the form along the boundary divisors plus one (Proposition~\hb{prop:mot-model}(ii)), those of the pinned family, so it is convergent, and by Lemma~\hb{lem:mot-birational} the two have one motivic period. Theorem~\hb{thm:mot-pinned}(ii) is the display.

(b) Let $t=0$; the case $s=0$ is the same with $\Psi$. If $s=0$ too, $\kappa_S+\kappa_{\cE\setminus S}=0$ for every $S$ on the balance wall (Proposition~\hb{prop:compoundwall}), against convergence; so $s\neq0$, and only $K$ is active: every cell of every substitution of $\cB$ is read through one hinge value, its $\gamma^K$. By Proposition~\hb{prop:compoundwall}, $\kappa_S+\kappa_{\cE\setminus S}=s\,c(S)\ge2$ for every nonempty proper $S$, so $\cB$ is connected. We induct on $|\cE|$.

If $\cB$ has a set $R$ of width one with $|R|,|\cE\setminus R|\ge2$, then $\cB=\mathcal H\langle\mathcal A\rangle$, the marks of $L=\cE\setminus R$ left in place, $\mathcal H$ the cell of the room and $\mathcal A$ the cell of the bunch $R$ (Proposition~\hb{prop:substitution}(iii)), connected lists on fewer marks (Proposition~\hb{prop:substitution}(iv)). Theorem~\hb{thm:mot-pinned}(ii),(iii) with $q=1$ give
\[
P^{\mathrm{mot}}_F=\Gst21{\kappa_R,\ s-\kappa_R}{s}^{-1}P^{\mathrm{mot}}_{(\mathcal H;s,0)}\,P^{\mathrm{mot}}_{(\mathcal A;s,0)},
\]
the two families convergent at the masses $(\mvec|_L,\kappa_R)$ and $(\mvec|_R,s-\kappa_R)$, where $s-\kappa_R=\kappa_L$ (Proposition~\hb{prop:compoundwall}); the marker of each carries the wall of $F$ at the marks on the other side of the split. A wall of $(\mathcal H;s,0)$ at a set containing the marker is the wall of $F$ at that set with the marker replaced by $R$, and at a set of marks of $L$ it is the wall of $F$ there (Theorem~\hb{thm:substwalls}, the hinge $g$ of $\mathcal A$ having $d_{\mathcal A}(R+g)=d_{\mathcal A}(R)+1$ and $d^\ast_{\mathcal A}(R+g)=d^\ast_{\mathcal A}(R)+1$); likewise for $(\mathcal A;s,0)$. So the induction, applied to the two families, writes $P^{\mathrm{mot}}_F$ as a positive rational number times the product over the prime nodes of $P^{\mathrm{mot}}_{F_v}$, with the masses of the statement.

Otherwise $\cB$ is prime, its own one node, or it has two or three marks. A connected list on two marks is $\binom\cE1$, and one on three marks is $\binom\cE1$ or $\binom\cE2$, having no loop and no coloop. On two marks $P^{\mathrm{mot}}_F$ is a positive rational number by Lemma~\hb{lem:motive-twomarks}. On three marks, for a mark $g$, neither a loop nor a coloop, $\cB$ is the cell $\binom{\{g,h\}}1\langle g,\mathcal A'\rangle$ of the proof of Theorem~\hb{thm:mot-pinned}(iii) with the two other marks as $R$, and Theorem~\hb{thm:mot-pinned}(ii) writes $P^{\mathrm{mot}}_F$ as the product of the motivic periods of a convergent family on the two marks $g,h$ and of a convergent pinned family on the two marks of $R$, each a positive rational number by Lemma~\hb{lem:motive-twomarks}. Lemma~\hb{lem:mot-conj}(iii) gives the conjugates, a nonzero rational factor changing nothing.
\end{proof}

\begin{htheorem}[\hb{thm:mot-splitoff} (A pinned set splits off motivically)]
As stated.
\end{htheorem}

\begin{proof}
$F$ is a family of lists, so by Theorem~\hb{thm:mot-pointmass} the active hinge values at $S$ are equal, $\gamma=\gamma^K_{j,S}=\gamma^\Psi_{j',S}$ for all active $K_j$ and $\Psi_{j'}$. Every active table $\cB_j$ has width one at $S$ with factoring weights, so it is the gluing of its cell of the bunch $\cB^+_{j,S}$ and its cell of the room $\cB^+_{j,S'}$ along their hinges (Proposition~\hb{prop:substitution}(iii)), that is the substitution $\cB^+_{j,S'}\langle\cB^+_{j,S}\rangle$ with the marks of $S'$ left in place (Definition~\hb{def:substitution}); the hinge values of $\cB^+_{j,S}$ are $\gamma^K_{j,S}$ and $\gamma^\Psi_{j,S}$. So the hypothesis of Theorem~\hb{thm:mot-pinned} holds with $q=1$, $L=S'$, $R_1=S$ and $\gamma_1=\gamma$, $|S|\ge2$. Its core family is $((\cB^+_{j,S'};s_j,t_j))_j$ at the masses $(\mvec|_{S'},\kappa_S)$. By Proposition~\hb{prop:ladder}(ii) the room station letters of $S'$ are $\mathsf y_{jD_0}=K_{\cB_j|S'}$, $\mathsf y_{jD_1}=K_{\cB_j/S}$, $\mathsf y^\ast_{jD_0}=\Psi_{\cB_j|S'}$ and $\mathsf y^\ast_{jD_1}=\Psi_{\cB_j/S}$, so the partner with its stand-in at $\one$ has the letters $K_{\cB_j|S'}+\alpha_gK_{\cB_j/S}$ and $\alpha_g\Psi_{\cB_j|S'}+\Psi_{\cB_j/S}$ and the mass $\kappa_S$ on $g$ (Theorem~\hb{thm:gluewidth}(i)), the letters of the cell of the room (Theorem~\hb{thm:widthone}): it is the core family. The period of the pinned family is $\int_{\Delta_S}\gamma^{-\kappa_S}u^{\mvec|_S-\one}\prod_jK_{\cB_j/S'}^{-s_j}\Psi_{\cB_j|S}^{-t_j}\,d\omega_S$ (Theorem~\hb{thm:mot-pinned}(iv)), and by the display after Definition~\hb{def:seamlaw}, with the weights factoring, this is $\int_{\Delta_S}\phi^\rho_S$: there $\gamma^\rho_S=\prod_j(\gamma^K_{j,S})^{s_j/w}(\gamma^\Psi_{j,S})^{t_j/w}=\gamma$, $w=\sum_j(s_j+t_j)$ being the sum of the exponents of the active tables, all of width one at $S$. Theorem~\hb{thm:mot-pinned}(ii) is the display, and its period is $P_F=c\,P_{S'^{(\one)}}$.
\end{proof}

\section{Chapter~\hb{ch:positive}: positive words}\label{sec:ch32}

\begin{htheorem}[\hb{thm:jets} (Positive coefficient jets)]
As stated, at positive weights that need not be rational.
\end{htheorem}

\begin{proof}
Write $P=P_F(n\mvec;n\mathbf s,n\mathbf t)$, finite and positive because its walls $n\kappa_S$ are positive (\hb{thm:projective}), and $\oE$ for the expectation of $U\sim\mathrm{Breta}^\circ(F;n\mvec;n\mathbf s,n\mathbf t)$, whose density on the simplex is $P^{-1}u^{n\mvec-\one}\prod_jK_j(u)^{-ns_j}\Psi_j(u)^{-nt_j}$, positive on the open simplex. For $v\in\mathbb R^k$ put $X_j=v^{\mathsf T}\partial_{\mathbf b}K_j/K_j$ and $X^\ast_j=v^{\mathsf T}\partial_{\mathbf b}\Psi_j/\Psi_j$. With $u_{j,T}=\mathbf b_{j,T}/\theta_{j,T}$,
\[
\frac{\partial_{\mathbf b}K_j}{K_j}=\oE_{T\sim\mathbb P_{j,\alpha}}[u_{j,T}],\qquad\frac{\partial_{\mathbf b}\Psi_j}{\Psi_j}=\oE_{T\sim\mathbb P^\ast_{j,\alpha}}[u_{j,T}],
\]
so $|X_j|$ and $|X^\ast_j|$ are at most $\beta_j=\max_T|v^{\mathsf T}u_{j,T}|$ on the whole orthant.

(i) For $|\varepsilon|\max_j\beta_j<\frac12$ the weights $\theta_{j,T}\pm\varepsilon v^{\mathsf T}\mathbf b_{j,T}=\theta_{j,T}(1\pm\varepsilon v^{\mathsf T}u_{j,T})$ are positive, and the letters of the two copies of table $j$ are $K_j(1\pm\varepsilon X_j)$ and $\Psi_j(1\pm\varepsilon X^\ast_j)$. The copies have the levels of table $j$, each counted twice at half the exponents, so the doubled family at the masses $n\mvec$ has the imbalance $0$ and the walls $n\kappa_S$, and its word is
\[
\prod_jK_j^{-ns_j}\Psi_j^{-nt_j}\cdot B_\varepsilon,\qquad B_\varepsilon=\prod_j\big(1-\varepsilon^2X_j^2\big)^{-\frac{ns_j}2}\big(1-\varepsilon^2X_j^{\ast2}\big)^{-\frac{nt_j}2}.
\]
Its period is therefore $P\,\oE[B_\varepsilon(U)]$, the display of (i). For $|\varepsilon|\le\varepsilon_0<(2\max_j\beta_j)^{-1}$, $B_\varepsilon$ is a power series in $\varepsilon^2$ whose coefficients are polynomials in the bounded functions $X_j^2,X_j^{\ast2}$, converging uniformly on the orthant; so $\oE[B_\varepsilon(U)]$ is analytic in $\varepsilon$ and expands term by term, with the coefficient $\sum_j\big(\frac{ns_j}2\oE[X_j^2]+\frac{nt_j}2\oE[X_j^{\ast2}]\big)$ of $\varepsilon^2$. Since $X_j^2=v^{\mathsf T}\big(\partial_{\mathbf b}K_j\,\partial_{\mathbf b}K_j^{\mathsf T}/K_j^2\big)v$, $P$ times this coefficient is $v^{\mathsf T}\mathbf H_nv$. It is a sum of nonnegative terms. If it vanishes, then for every $j$ with $s_j>0$, $X_j(U)=0$ almost surely; the law of $U$ has a positive density on the open simplex, so the polynomial $\sum_T(v^{\mathsf T}\mathbf b_{j,T})\alpha^T$ vanishes on an open set, and every $v^{\mathsf T}\mathbf b_{j,T}$ is $0$, the monomials of distinct bases being distinct. For $t_j>0$ the same holds with $\alpha^{\cE\setminus T}$. So $v$ is orthogonal to every direction of an active table, and $v=0$ when these span $\mathbb R^k$.

(ii) Write $Q=R^{\mathsf T}R$ with $R$ real. Then
\[
\operatorname{tr}(Q\mathbf H_n)=\frac n2\,P\,\oE\Big[\sum_j\Big(s_j\frac{|R\,\partial_{\mathbf b}K_j|^2}{K_j^2}+t_j\frac{|R\,\partial_{\mathbf b}\Psi_j|^2}{\Psi_j^2}\Big)(U)\Big],
\]
a sum of nonnegative terms, which vanishes exactly when, as in (i), $R\mathbf b_{j,T}=0$, that is $\mathbf b_{j,T}^{\mathsf T}Q\mathbf b_{j,T}=0$, for every basis of every active table. If $Q_{11}>0$, then $Q-Q_{11}^{-1}Q_{\cdot1}Q_{\cdot1}^{\mathsf T}$ is rational, positive semidefinite, with first row and column zero; if $Q_{11}=0$, the first row and column of $Q$ vanish, $Q$ being positive semidefinite. Induction on $k$ gives $Q=\sum_lc_lv_lv_l^{\mathsf T}$ with rational $c_l>0$ and rational $v_l$, and then $\operatorname{tr}(Q\mathbf H_n)=\sum_lc_l\,v_l^{\mathsf T}\mathbf H_nv_l$.

(iii) $\partial_{\mathbf b}K_j^{\mathsf T}Q\,\partial_{\mathbf b}K_j/K_j^2=\sum_{T,T'}\mathbf b_{j,T}^{\mathsf T}Q\,\mathbf b_{j,T'}\,\alpha^{\one_T+\one_{T'}}K_j^{-2}$, and $P\,\oE\big[(\alpha^{\one_T+\one_{T'}}K_j^{-2})(U)\big]$ is the period whose integrand is that of $P$ times $u^{\one_T+\one_{T'}}K_j(u)^{-2}$, the period at $\big(n\mvec+\one_T+\one_{T'};n\mathbf s+2\mathbf e_j,n\mathbf t\big)$. Its imbalance is that of $(n\mvec;n\mathbf s,n\mathbf t)$ plus $2r_j-2r_j=0$, and its wall of $S$ is $n\kappa_S+|T\cap S|+|T'\cap S|-2d_j(S)\ge n\kappa_S$, by \hb{thm:window}. The cobasis part is the same with $u^{\one_{\cE\setminus T}+\one_{\cE\setminus T'}}\Psi_j^{-2}$ and the change $|S\setminus T|+|S\setminus T'|-2d^\ast_j(S)\ge0$. This is the case of order two of \hb{thm:respwalls} for the doubled family, the two copies of each table merged.

(iv) $x\mapsto x^{\mathsf T}Qx$ is convex, so $\big(\oE_{T}[u_{j,T}]\big)^{\mathsf T}Q\,\oE_T[u_{j,T}]\le\oE_T\big[u_{j,T}^{\mathsf T}Qu_{j,T}\big]\le\max_Tu_{j,T}^{\mathsf T}Qu_{j,T}$ under either law, and $\operatorname{tr}(Q\mathbf H_n)\le\frac n2P\sum_j(s_j+t_j)\max_T\mathbf b_{j,T}^{\mathsf T}Q\,\mathbf b_{j,T}/\theta_{j,T}^2=n\,C(Q)\,P$; positivity is (ii).

(v) With the basis weights $\theta^K_{j,T}$ and the cobasis weights $\theta^\Psi_{j,T}$ of every table independent (\hb{def:family}), $\partial^2K_j^{-ns_j}/\partial\theta^K_{j,T}\partial\theta^K_{j,T'}=ns_j(ns_j+1)K_j^{-ns_j-2}\alpha^{T}\alpha^{T'}$. The differentiated integrals are the periods of (iii), convergent locally uniformly in the weights, so differentiation under the integral gives
\[
\sum_{T,T'}\mathbf b_{j,T}\mathbf b_{j,T'}^{\mathsf T}\,\frac{\partial^2P}{\partial\theta^K_{j,T}\partial\theta^K_{j,T'}}=ns_j(ns_j+1)\,P\,\oE\Big[\frac{\partial_{\mathbf b}K_j\,\partial_{\mathbf b}K_j^{\mathsf T}}{K_j^2}(U)\Big],
\]
which is $2(ns_j+1)$ times the part $\frac n2s_jP\,\oE[\cdots]$ of $\mathbf H_n$ carried by $K_j$; the cobasis part is the same with $t_j$. On the share side of \hb{cor:atomwall} the weights enter only through $\prod_{j,T}(\theta^K_{j,T})^{-a_{j,T}}(\theta^\Psi_{j,T})^{-b_{j,T}}$: a derivative in $\theta^K_{j,T}$ multiplies the integrand by $-a_{j,T}/\theta^K_{j,T}$, one in $\theta^\Psi_{j,T}$ by $-b_{j,T}/\theta^\Psi_{j,T}$, and neither the contour, nor the pins, nor the share kernels involve the weights.

(vi) At the centre the two copies of table $j$ carry the equal exponents $(\frac{ns_j}2,\frac{ns_j}2)$, and the masses $n\sigma\one$ are the total strength of the doubled family times $\one$. \hb{prop:centrefamily}(i) holds at every family at its centre, whatever its weights: at $\one$ the usage is $\sum_j\sigma_j(R_j+R^\ast_j)(\one)=\sigma\one$, since $R_{j,e}(\one)+R^\ast_{j,e}(\one)$ is the chance that $e$ lies in the basis plus the chance that it does not.
\end{proof}

\begin{sremark}
No step uses the rationality of the weights; rational weights are what the arithmetic of the jets asks for, through the rows of the periods of (iii).
\end{sremark}

\begin{htheorem}[\hb{thm:twolists} (Two bases generate positive contrast words of any depth)]
As stated, at every point of the balance wall with real data and positive proper walls, and at positive weights.
\end{htheorem}

\begin{proof}
Write $P_n=P_F(n\mvec;n\mathbf s,n\mathbf t)$, finite and positive (\hb{thm:projective}), $\oE$ for the expectation of $U\sim\mathrm{Breta}^\circ(F;n\mvec;n\mathbf s,n\mathbf t)$, and $c=\mathbb P_{j,\alpha}(T)-\mathbb P_{j,\alpha}(T')$, $c^\ast=\mathbb P^\ast_{j,\alpha}(T)-\mathbb P^\ast_{j,\alpha}(T')$. Each chance lies in $(0,1)$ and two chances of distinct bases sum to at most one, so $|c|,|c^\ast|<1$ on the open orthant. $c$ vanishes only where $\theta_{j,T}\alpha^{\one_T}=\theta_{j,T'}\alpha^{\one_{T'}}$, a hypersurface, since $\one_T\neq\one_{T'}$; so $c^{2h}>0$ almost everywhere, and the contrast word, $\frac12P_n$ times the expectation of a function that is positive almost everywhere under a law with a positive density, is positive.

(i) $c^{2h}=\sum_{l=0}^{2h}(-1)^l\binom{2h}l\theta_{j,T}^{2h-l}\theta_{j,T'}^{\,l}\,\alpha^{v_l}K_j^{-2h}$ with $v_l=(2h-l)\one_T+l\one_{T'}$, each term homogeneous of degree $0$, and $P_n\,\oE\big[(\alpha^{v_l}K_j^{-2h})(U)\big]$ is the period whose integrand is that of $P_n$ times $u^{v_l}K_j(u)^{-2h}$, the period at $\big(n\mvec+v_l;n\mathbf s+2h\mathbf e_j,n\mathbf t\big)$. Its imbalance changes by $|v_l|-2hr_j=0$, and its wall of $S$ by $(2h-l)|T\cap S|+l|T'\cap S|-2hd_j(S)\ge0$ (\hb{thm:window}). The second half is the same with $\mathbb P^\ast_{j,\alpha}(T)=\theta_{j,T}\alpha^{\cE\setminus T}/\Psi_j$, the masses raised by $(2h-l)\one_{\cE\setminus T}+l\one_{\cE\setminus T'}$, the cobasis exponent by $2h$, and the wall of $S$ by $(2h-l)|S\setminus T|+l|S\setminus T'|-2hd^\ast_j(S)\ge0$.

From $\Psi_j(\alpha)=\alpha^{\one}K_j(\alpha^{-1})$, $K_j(\alpha^{-1})=\alpha^{-\one}\Psi_j(\alpha)$ and $\Psi_j(\alpha^{-1})=\alpha^{-\one}K_j(\alpha)$, so $\mathbb P^\ast_{j,\alpha^{-1}}(T)=\mathbb P_{j,\alpha}(T)$. On the chart $y_g=0$ of \hb{thm:projective} the integrand of a period at data $(\mvec';\mathbf s',\mathbf t')$ is $e^{\ip{\mvec'}y}\prod_jK_j(e^y)^{-s'_j}\Psi_j(e^y)^{-t'_j}$, and $y\mapsto-y$, which keeps $dy$, carries it to $e^{\ip{2\sigma'\one-\mvec'}y}\prod_jK_j(e^y)^{-t'_j}\Psi_j(e^y)^{-s'_j}$, $\sigma'$ the total strength of $(\mathbf s',\mathbf t')$. Applied to the second half at $(n\mvec;n\mathbf s,n\mathbf t)$ it gives the first half at $\big(n(2\sigma\one-\mvec);n\mathbf t,n\mathbf s\big)$, the depth $n$ of the data $(2\sigma\one-\mvec;\mathbf t,\mathbf s)$ of \hb{thm:cremona}, whose wall of $S$ is the wall of $\cE\setminus S$ at $(\mvec;\mathbf s,\mathbf t)$, positive. At the centre these data are $(\mvec;\mathbf s,\mathbf t)$ themselves, and the two halves are equal.

(ii) By \hb{def:question}, with $X(z)=2\ip{\bmu_\perp}z-\Phi^\circ(z)$ read on $\one^\perp$ and the number of marks written $|\cE|$, $P_n=2^{|\cE|-1}\int_{\mathcal T}e^{nX}dz$ and the contrast word is $2^{|\cE|-2}\int_{\mathcal T}e^{nX}\big(c^{2h}+c^{\ast2h}\big)dz$, with $c,c^\ast$ at $\alpha=e^{2z}$. For $h=\eta n$ put $\psi=-X-2\eta\log\max(|c|,|c^\ast|)$, with values in $(-\infty,\infty]$, continuous where finite and lower semicontinuous; then $e^{-n\psi}\le e^{nX}(c^{2h}+c^{\ast2h})\le2e^{-n\psi}$. By \hb{prop:evenfloor}(i) and (ii), $-X(z)\ge2\kappa_{\min}(\max_ez_e-\min_ez_e)-C$, and the logarithmic term is nonnegative, so $\psi$ grows at least linearly on $\one^\perp$. Laplace's principle follows: outside a large ball the integral of $e^{-n\psi}$ is at most $e^{-n(\inf\psi+1)}$ for large $n$, inside it is at most $e^{-n\inf\psi}$ times the volume, and near a point $z_0$ with $\psi(z_0)<\inf\psi+\epsilon$ the integrand is at least $e^{-n(\inf\psi+2\epsilon)}$ on a ball. So $\frac1n\log$ of the contrast word tends to $-\inf\psi$, and the same with $\eta=0$ and without the contrasts gives $\frac1n\log P_n\to\max X=-\min(\Phi^\circ-2\ip{\bmu_\perp}\cdot)$, the maximum attained, $X$ being continuous and tending to $-\infty$. The set $K=\{z\in\one^\perp:\ 2\kappa_{\min}(\max_ez_e-\min_ez_e)-C\le\inf\psi+1\}$ is compact and contains every $z$ with $\psi(z)\le\inf\psi+1$; on it $\max(|c|,|c^\ast|)\le1-\delta_K<1$, so $\inf\psi\ge\min(-X)+2\eta|\log(1-\delta_K)|>\min(-X)$. For $A>0$ let $B_A$ be a ball of $\one^\perp$ outside which $2\kappa_{\min}(\max_ez_e-\min_ez_e)-C\ge A$; on $B_A$, $\psi\ge\min(-X)+2\eta|\log(1-\delta_{B_A})|$, which exceeds $A$ for large $\eta$, so $\inf\psi\ge A$ for large $\eta$.

(iii) Let $s_j>0$ and $q=2h$. In the share form of $P_F$ used by \hb{cor:atomwall}, the letter $K_j^{-ns_j-q}$ of the term $l$ of (i) is opened by one share $a'_{U}$ per basis, $\sum_Ua'_U=ns_j+q$, with the share kernel $\prod_U\Gamma(a'_U)/\Gamma(ns_j+q)$ and the factors $(\theta_{j,U}\alpha^{\one_U})^{-a'_U}$ (\hb{thm:E1}). The insertion $\theta_{j,T}^{q-l}\theta_{j,T'}^{\,l}\alpha^{v_l}$ turns the factors of $T$ and $T'$ into $(\theta_{j,T}\alpha^{\one_T})^{-a_T}$ and $(\theta_{j,T'}\alpha^{\one_{T'}})^{-a_{T'}}$ with $a_T=a'_T-q+l$, $a_{T'}=a'_{T'}-l$, the other shares unchanged as $a_U=a'_U$. In the shares $a$ the total is $ns_j$ and the factors in $\alpha$ and in the weights are those of $P_F(n\mvec;n\mathbf s,n\mathbf t)$, so the integral over $\alpha$, which produces the pins, gives the pins of $(n\mvec;n\mathbf s,n\mathbf t)$, and the kernel is
\[
\frac{\Gamma(a_T+q-l)\,\Gamma(a_{T'}+l)\prod_{U\neq T,T'}\Gamma(a_U)}{\Gamma(ns_j+q)}=\frac{\prod_U\Gamma(a_U)}{\Gamma(ns_j)}\cdot\frac{(a_T)_{q-l}\,(a_{T'})_l}{(ns_j)_q}.
\]
The factor $\Gamma(a_T+q-l)\Gamma(a_{T'}+l)$ is holomorphic for $\re a_T>-(q-l)$ and $\re a_{T'}>-l$, so the contour returns to a wall contour of $(n\mvec;n\mathbf s,n\mathbf t)$ without crossing a pole, the share kernels decaying along the imaginary directions. Summing over $l$ with the weights $\frac12(-1)^l\binom ql$ gives the share side of $P_F(n\mvec;n\mathbf s,n\mathbf t)$ with the insertion $H_q(a_{j,T},a_{j,T'})/\big(2(ns_j)_q\big)$. For $t_j>0$ the second half is the same with $\Psi_j$, the cobasis shares $b$ and $nt_j$. The insertions are polynomials in the shares and change neither the pins nor the share kernels, hence no pole hyperplane. At the centre the two halves are equal by (i), and their insertions add to the one of the reference.
\end{proof}

\begin{htheorem}[\hb{thm:portsequence} (A positive sequence of Beta ports on an interval of a selected mark)]
As stated, at every mass $m_e\in(0,1]$.
\end{htheorem}

\begin{proof}
In $x$, $d\alpha_e=(\lambda_1-\lambda_0)\,dx$, and the constant goes into the room density. The factor of the mass is $\alpha_e^{m_e-1}=L_0^{-(1-m_e)}$ with $L_0=\lambda_0+(\lambda_1-\lambda_0)x$, a linear letter with $c_0=\lambda_0$, $d_0=\lambda_1-\lambda_0>0$ and the exponent $w_0=1-m_e\ge0$; it satisfies $c_0\ge\lambda_0d_0/(\lambda_1-\lambda_0)$ with equality, so the read is $\langle\cdot\rangle_{1,1;\mathbf w}$ over the letters of the chart and $L_0$. For a letter with $d_i>0$ put $\mu=\lambda_0/(\lambda_1-\lambda_0)$; then $x(1-x)d_i/L_i\le x(1-x)/(\mu+x)$, whose derivative vanishes on $[0,1]$ only where $x^2+2\mu x-\mu=0$, at $x=\sqrt{\mu(1+\mu)}-\mu$, with the value $\big(\sqrt{1+\mu}-\sqrt\mu\big)^2=(\sqrt{\lambda_1}-\sqrt{\lambda_0})^2/(\lambda_1-\lambda_0)=\varepsilon_{01}$, the ratio of the statement; for $d_i=0$ the quotient is $0$. By \hb{thm:portreads}(i) at $a=b=1$,
\[
\big\langle(-1)^n\Pi^{1,1}_n\big\rangle_{1,1;\mathbf w}=\sum_{|v|=n}\prod_i\frac{(w_i)_{v_i}}{v_i!}\Big\langle\prod_i\Big(\frac{x(1-x)d_i}{L_i}\Big)^{v_i}\Big\rangle_{1,1;\mathbf w},
\]
every term between $0$ and $\prod_i\frac{(w_i)_{v_i}}{v_i!}\varepsilon_{01}^{\,n}\langle1\rangle_{1,1;\mathbf w}$, and $\sum_{|v|=n}\prod_i(w_i)_{v_i}/v_i!=(|\mathbf w|)_n/n!$ by the Vandermonde identity for rising factorials. A letter $L_i$ that sees the mark makes the term $v=n\mathbf e_i$ positive (\hb{thm:portreads}(ii) with $J=\{i\}$), and for $n=0$ the read is $\langle1\rangle_{1,1;\mathbf w}>0$. When $m_e<1$, $w_0d_0=(1-m_e)(\lambda_1-\lambda_0)>0$: the letter of the mass sees the mark.
\end{proof}

\begin{htheorem}[\hb{thm:portrecurrence} (The finite-band recurrence in the order of the Beta ports)]
As stated, at every pair $(a_n,b_n)$ with $a_n,b_n>0$.
\end{htheorem}

\begin{proof}
Write $a=a_n$, $b=b_n$, $w_i=w_i(n)$. The letters are at least $\min_ic_i>0$ on $[0,1]$, so every read of a polynomial converges and the flux $x^a(1-x)^b\,A\prod_iL_i^{-w_i}$ of a polynomial $A$ vanishes at both ends. \hb{thm:markrelation}(a) with the polynomial $Lg$ reads, by $x(1-x)\partial_x(Lg)=x(1-x)(L\partial_xg+g\partial_xL)$ and $\partial_x\log L_i=d_i/L_i$,
\[
x(1-x)\partial_x(Lg)+\big(a-(a+b)x\big)Lg-x(1-x)Lg\sum_i\frac{w_id_i}{L_i}=L\,x(1-x)\,\partial_xg+A_ng .
\]

The coefficient of $x^q$ in $\Pi^{a,b}_q$ is $\binom{a+b+2q-2}q$, a product of the $q$ positive numbers $a+b+q-1,\dots,a+b+2q-2$ over $q!$, so $\Pi^{a,b}_0,\dots,\Pi^{a,b}_N$ are a basis of the polynomials of degree at most $N$, and the conditions $a+b+2q-i\neq0$ of \hb{thm:betaport}(v) hold. By its finite rules, multiplication by $x$ maps $\Pi^{a,b}_q$ into the span of $\Pi^{a,b}_{q-1},\Pi^{a,b}_q,\Pi^{a,b}_{q+1}$, and $x(1-x)\partial_x$ maps it into the span of $\Pi^{a,b}_{q-1},\Pi^{a,b}_q$ and $x\Pi^{a,b}_q$, with coefficients rational in $(a,b,q)$; for $q=0$, $x\Pi^{a,b}_0=x$ and $\partial_x\Pi^{a,b}_0=0$. So multiplication by a polynomial of degree $d$ moves $\Pi^{a,b}_q$ into the span of $\Pi^{a,b}_{q-d},\dots,\Pi^{a,b}_{q+d}$. $L$ has degree $h$ and $A_n$ degree at most $h+1$, so at $g=(-1)^r\Pi^{a,b}_r$,
\[
L\,x(1-x)\,\partial_xg+A_ng=\sum_{s=-h-1}^{h+1}e_s\,\Pi^{a,b}_{r+s},
\]
$e_s=0$ when $r+s<0$, with $e_s$ rational in $(a,b,r)$ and polynomial in the coefficients of $L$ and of $A_n$, hence in the room data and the $w_i$; the mark relation is the displayed relation with $\mathfrak b_{n,r,s}=(-1)^{r+s}e_s$.

$\Pi^{a,b}_{r+h+1}$ is the only member of the expansion of degree $r+h+1$, so $e_{h+1}$ is the coefficient of $x^{r+h+1}$ divided by $\binom{a+b+2r+2h}{r+h+1}$. The coefficient of $x^r$ in $g$ is $(-1)^r\binom{a+b+2r-2}r$. The coefficient of $x^{h+1}$ in $A_n$ is $-(a+b)p_h-hp_h+|\mathbf w(n)|\,p_h$, the last from $-x(1-x)\sum_iw_id_i\prod_{i'\neq i}L_{i'}$, whose leading coefficient is $\sum_iw_id_i\prod_{i'\neq i}d_{i'}=|\mathbf w(n)|\,p_h$; and the coefficient of $x^{r+h+1}$ in $L\,x(1-x)\,\partial_xx^r$ is $-rp_h$. So the coefficient of $x^{r+h+1}$ is $(-1)^r\binom{a+b+2r-2}r\,p_h\big(|\mathbf w(n)|-a-b-h-r\big)$, and $\mathfrak b_{n,r,h+1}=(-1)^{r+h+1}e_{h+1}$ is the stated value.
\end{proof}

At $a_n=b_n=n+1$ the furthest coefficient is
\[
\big[|\mathbf w(n)|-2n-2-h-r\big]\,p_h\,(-1)^{h+1}\binom{2n+2r}r\Big/\binom{2n+2r+2h+2}{r+h+1},
\]
that of the reference.

\section{Chapter~\hb{ch:share}: the critical-line law}\label{sec:chshare}

\begin{hproposition}[\hb{prop:ra-binomial} (Binomial moments of the critical-line law)]
As stated.
\end{hproposition}

\begin{proof}
(i) The density $p_{\mathrm{crit}}(y)=2\pi/(e^{\pi y}+e^{-\pi y})^2$ is at most $2\pi e^{-2\pi|y|}$, so $\oE[e^{tX}]=e^{t/2}\oE[e^{itY}]$ is holomorphic in $t$ on the strip $|\im t|<2\pi$, as is $te^t/(e^t-1)$, whose poles $2\pi ik$, $k\neq0$, lie outside it; they agree for real $t$ (\hb{thm:sharelaw}(iii)), hence on the strip. For $|z|<1$, $t=\log(1+z)$ has $|\im t|<\frac\pi2$, and $te^t/(e^t-1)=(1+z)\log(1+z)/z$.

(ii) $|\binom xi|\le\binom{|x|+i-1}{i}$ termwise, and $\sum_i\binom{|x|+i-1}i|z|^i=(1-|z|)^{-|x|}$. For $|z|<1-e^{-2\pi}$, $(1-|z|)^{-|X|}\le e^{c(\frac12+|Y|)}$ with $c=-\log(1-|z|)<2\pi$, which has a finite expectation, so $\oE[(1+z)^X]=\sum_i\oE[\binom Xi]z^i$ there. The expansion $\log(1+z)/z=\sum_{i\ge0}(-1)^iz^i/(i+1)$ times $1+z$ has the coefficient $1$ at $i=0$ and $(-1)^i\big(\frac1{i+1}-\frac1i\big)=\frac{(-1)^{i+1}}{i(i+1)}$ at $i\ge1$.

(iii) Differentiating $(1+z)^u=\sum_i\binom uiz^i$ in $u$ gives $\log(1+z)(1+z)^u=\sum_i\frac d{du}\binom uiz^i$, so $\frac d{du}\binom ui=\sum_{m=1}^i\frac{(-1)^{m-1}}m\binom u{i-m}$, an identity of polynomials. By (ii), $\oE[\frac d{du}\binom ui|_X]$ is the coefficient of $z^i$ in $\log(1+z)\cdot(1+z)\log(1+z)/z$. From $\frac d{dz}\log^2(1+z)=2\log(1+z)/(1+z)=2\sum_{n\ge1}(-1)^{n-1}H_nz^n$, $\log^2(1+z)=2\sum_{n\ge2}(-1)^nH_{n-1}z^n/n$, and the coefficient of $z^i$ in $(1+z)\log^2(1+z)/z$ is
\[
2(-1)^{i+1}\frac{H_i}{i+1}+2(-1)^i\frac{H_{i-1}}i=2(-1)^i\Big(\frac{H_{i-1}}i-\frac{H_{i-1}+\frac1i}{i+1}\Big)=2(-1)^i\frac{H_{i-1}-1}{i(i+1)} .
\]
By \hb{prop:critical} at $k=1$, for $R=(1-2u)Q$, $R(\frac12)=0$ and $\operatorname{crit}_1[R]=\oE[Q'(X)]$, which is the sum.
\end{proof}

\begin{hproposition}[\hb{prop:ra-halfpoles} (Poles at the half-integers)]
As stated.
\end{hproposition}

\begin{proof}
Put $S(u)=\sum_m(u+m)^{-(2k+1)}$. From $\Gst40{u,\ u,\ 1-u,\ 1-u}{}=\sum_m(u+m)^{-2}$ (\hb{lem:resonantpair}), $\partial_u^{2k-1}\Gst40{u,\ u,\ 1-u,\ 1-u}{}=-(2k)!\,S(u)$, so $\operatorname{crit}_k[R]=\frac2{2\pi i}\int_{(1/2)}R(u)S(u)\,du/(1-2u)$. $S$ has the period $1$, poles of order $2k+1$ at the integers, it decays exponentially as $|\im u|\to\infty$ uniformly in $\re u$, and $S(1-u)=-S(u)$, so $S(\frac12)=0$.

For $p\notin\mathbb Z$ with $\re p>\frac12$ let $I(p)=\frac1{2\pi i}\int_{(1/2)}S(u)\,du/(u-p)$. Move the line to $\re u=N+\frac12$: the horizontal segments vanish by the decay of $S$, and on the new line $|1/(u-p)|\le1/(N+\frac12-\re p)\to0$ while $\int|S|$ there is that on $\re u=\frac12$. So $I(p)$ is minus the sum of the residues to the right of the line: at $u=p$ the residue $S(p)$, and at an integer $i\ge1$ the residue of $(u-i)^{-2k-1}/(u-p)$, which is $(i-p)^{-2k-1}$. At $p=\frac12+l$, $l\ge1$, $S(p)=\sum_{x\in\frac12+\mathbb Z}x^{-2k-1}=0$, and
\[
I(p)=-\sum_{i\ge1}\big(i-\tfrac12-l\big)^{-2k-1}=-\sum_{x\in\frac12+\mathbb Z,\ x>l}x^{-2k-1}=-\Tx{2k+1}{\Ilat{\frac12+\mathbb Z}}{l},
\]
the points $x$ with $|x|<l$ cancelling in pairs. At $l=0$ the integrand $S(u)/(u-\frac12)$ is regular on the line; for real $p=\frac12+\varepsilon$, $\varepsilon\downarrow0$, $|S(u)/(u-p)|\le|S(u)/(u-\frac12)|$ on the line, an integrable function, so the integral is $\lim I(p)=-S(\tfrac12)-\sum_{i\ge1}(i-\frac12)^{-2k-1}=-\Tx{2k+1}{\Ilat{\frac12+\mathbb Z}}{0}$.

For the word, $\frac1{(u-p)(1-2u)}=\frac1{1-2p}\big(\frac1{u-p}-\frac1{u-\frac12}\big)$, so at $p=\frac12+l$
\begin{multline*}
\operatorname{crit}_k\Big[\frac1{u-p}\Big]=\frac2{1-2p}\Big(\Tx{2k+1}{\Ilat{\frac12+\mathbb Z}}{0}-\Tx{2k+1}{\Ilat{\frac12+\mathbb Z}}{l}\Big)\\
=-\frac1l\,\Th{2k+1}{\Ilat{\frac12+\mathbb Z}}{0}{l}
\end{multline*}
by the cut relation (\hb{prop:cutlaw}). Finally $u\mapsto1-u$ maps the line onto itself, reversing it, and with $S(1-u)=-S(u)$ it gives $\operatorname{crit}_k[R(1-u)]=\operatorname{crit}_k[R(u)]$; for $R(u)=1/(u-\frac12+l)$, $R(1-u)=-1/(u-\frac12-l)$.
\end{proof}

\section{Chapter~\hb{ch:arith}: the gap and the shape words}\label{sec:ch34}

\begin{hproposition}[\hb{prop:kdepth} (The root rate of the shape words)]
As stated, with the maximizer at every $k\ge2$.
\end{hproposition}

\begin{proof}
(i), the limit in (ii) and the last sentence are the reference's (the reference, proof of prop:kdepth). For the maximum, fix $x_1,\dots,x_k$ and put $y=x_1\cdots x_k$. The factor of $f_k$ in $x_0$ is $x_0(1-x_0)/\big(1-(1-y)x_0\big)$; its derivative vanishes where $1-2x_0+(1-y)x_0^2=0$, at $x_0=1/(1+\sqrt y)$, where $1-x_0=\sqrt y/(1+\sqrt y)$ and $1-(1-y)x_0=\sqrt y$, with the value $(1+\sqrt y)^{-2}$. So $M_k(0)=\sup\prod_{i\ge1}x_i(1-x_i)\big/\big(1+\sqrt{x_1\cdots x_k}\big)^2$. At a fixed product the numerator is largest when the $x_i$ are equal, $t\mapsto\log(1-e^t)$ being concave, so $M_k(0)=\max_{0<a<1}\phi(a)$ with $\phi(a)=a^k(1-a)^k/(1+a^{k/2})^2$, and
\[
\frac{\phi'(a)}{\phi(a)}=k\Big(\frac1a-\frac1{1-a}-\frac{a^{k/2-1}}{1+a^{k/2}}\Big)=\frac{k\,\big(1-2a-a^{k/2+1}\big)}{a(1-a)(1+a^{k/2})}.
\]
The numerator decreases strictly from $1$ at $a=0$ to $-2^{-k/2-1}$ at $a=\frac12$, so $\phi$ increases up to the unique root $a\in(0,\frac12)$ of $a^{k/2+1}+2a=1$ and decreases after it. At $\eta=0$ the equation of \hb{thm:portgauge} is $(1-2a)^2=a^{k+2}$, whose root in $(0,\frac12)$ solves $1-2a=a^{k/2+1}$: the same $a$, with $x_2=\dots=x_k=a$ and $x_0=a/(1-a)=1/(1+a^{k/2})$, since $a^{k/2}=(1-2a)/a$. The interval follows from $\phi(\frac12)\le M_k(0)\le4^{-k}$. At $k=2$ the root is $\sqrt2-1$ and at $k=4$ the root of $a^3+2a-1$, the maximizers of the reference.
\end{proof}

\begin{hproposition}[\hb{prop:portindexstates} (The insertion index supplies auxiliary states)]
As stated, for every $k\ge3$.
\end{hproposition}

\begin{proof}
By \hb{thm:fourshapeladder}(ii), which holds for every $k\ge2$, the share function of $P_r$ is
\[
g_{n,r}(a)=\binom{n+r}r\,\frac{(n!)^{k-2}\,(1-a)_n^2\,(n+1-a)_r}{(a)_{n+1}^k\,(n+1+a)_r}.
\]
Its poles in $a\le0$ are at $a=-h$, $0\le h\le n$, of order exactly $k$, the numerator being nonzero there, and at $a=-h$, $n<h\le n+r$, simple. For $k\ge3$ the coefficients of order $k-1\ge2$ come from the first poles only, so $C_{n,k}(P_r)=\sum_{h=0}^nR_h(-h)$ and $C_{n,k-1}(P_r)=\sum_{h=0}^nR_h'(-h)$, with $R_h(a)=(a+h)^kg_{n,r}(a)$, holomorphic near $-h$. First,
\[
R_h(-h)=\binom{n+r}r\frac{(n!)^{k-2}(1+h)_n^2}{\prod_{j\neq h}(j-h)^k}\cdot\frac{(n+1+h)_r}{(n+1-h)_r},\qquad\frac{(n+1+h)_r}{(n+1-h)_r}=\frac{(n+1-h+r)_{2h}}{(n+1-h)_{2h}},
\]
the product over $0\le j\le n$; $\binom{n+r}r=(r+1)_n/n!$ and $(n+1-h+r)_{2h}$ are polynomials in $r$, $n+1-h\ge1$, so $R_h(-h)\in\mathbb Q[r]$ and $C_{n,k}(P_r)\in\mathbb Q[r]$. Second, $R_h'(-h)=R_h(-h)\,\rho_h$ with
\[
\rho_h=-2\sum_{i=1}^n\frac1{i+h}-\sum_{i=n+1}^{n+r}\Big(\frac1{i+h}+\frac1{i-h}\Big)-k\sum_{j\neq h}\frac1{j-h},
\]
the logarithmic derivatives at $-h$ of $(1-a)_n^2$, $(n+1-a)_r$, $(n+1+a)_r^{-1}$ and $\prod_{j\neq h}(a+j)^{-k}$. The finite Trailmoment $\Th{1}{\IZ}{0}{N}$ is the sum of $m^{-1}$ over $0<m\le N$, so $\sum_{i=n+1}^{n+r}\frac1{i+h}=\Th{1}{\IZ}{0}{n+r}-\Th{1}{\IZ}{0}{n+h}+\sum_{i=1}^h\frac1{n+r+i}$ and $\sum_{i=n+1}^{n+r}\frac1{i-h}=\Th{1}{\IZ}{0}{n+r}-\Th{1}{\IZ}{0}{n-h}-\sum_{i=0}^{h-1}\frac1{n+r-i}$, and $\rho_h=-2\,\Th{1}{\IZ}{0}{n+r}+\tilde\rho_h(r)$ with $\tilde\rho_h\in\mathbb Q(r)$. Summing, $C_{n,k-1}(P_r)=-2\,\Th{1}{\IZ}{0}{n+r}\,C_{n,k}(P_r)+q_n(r)$ with $q_n=\sum_hR_h(-h)\tilde\rho_h\in\mathbb Q(r)$.

Let $\sum_{i=0}^db_i(r)\,C_{n,k-1}(P_{r+i})=0$ for all large $r$, $b_i\in\mathbb Q(r)$. Since $\Th{1}{\IZ}{0}{n+r+i}=\Th{1}{\IZ}{0}{n+r}+\sum_{l=1}^i\frac1{n+r+l}$, this reads $A(r)\,\Th{1}{\IZ}{0}{n+r}+B(r)=0$ for all large $r$ with $A=-2\sum_ib_i(r)C_{n,k}(P_{r+i})$ and $B$ in $\mathbb Q(r)$. If $A\neq0$, then $\Th{1}{\IZ}{0}{n+r}=-B(r)/A(r)$ for all large $r$; a rational function is asymptotic to $cr^e$ with $e\in\mathbb Z$, while $\Th{1}{\IZ}{0}{n+r}-\log r$ has a limit, which is impossible. So $A=0$: the recurrence annihilates $C_{n,k}(P_r)$.
\end{proof}

\begin{htheorem}[\hb{thm:portbudget} (An arithmetic budget for growing Beta ports)]
As stated, with the coefficient of the top coordinate at every $k\ge2$.
\end{htheorem}

\begin{proof}
The two limits and the clearing are the reference's (the reference, proof of thm:portbudget). The coefficient of $\Tx{k+1}{\IZ}{0}$ in $I_n(k;P)$ is $k\,C_{n,k}(P)$ (\hb{thm:insertclosing}(iii)), and $P=(-1)^{q}\Pi^{n+1,n+1}_q(x_1)$ is free of $x_0$, so the theorem applies. As in the proof of \hb{prop:portindexstates}, at $r=q$, $C_{n,k}(P)=\sum_{h=0}^nR_h(-h)$ with $\prod_{j\neq h}(j-h)=(-1)^hh!\,(n-h)!$ and $(1+h)_n=(n+h)!/h!$:
\begin{align*}
R_h(-h)&=\binom{n+q}q\frac{(n!)^{k-2}}{(h!(n-h)!)^k}\frac{((n+h)!)^2}{(h!)^2}(-1)^{hk}\frac{(n+1+h)_q}{(n+1-h)_q}\\&=\binom{n+q}q(-1)^{hk}\binom nh^k\binom{n+h}n^2\frac{(n+1+h)_q}{(n+1-h)_q}.
\end{align*}
For even $k$ every term is positive.
\end{proof}

\begin{hcorollary}[\hb{cor:growingport} (A growing Beta port with a certified gap)]
As stated, for every $k\ge2$.
\end{hcorollary}

\begin{proof}
The port factor is at most one, since $x_0x_1\bar x_1x_2\cdots x_k\le x_0x_1\cdots x_k\le L$, so $M_k(\eta)\le M_k(0)$ for $\eta\ge0$ and
\[
\ell_k(\eta)+(k+1)(1+\eta)\,c_{\mathrm{lcm}}\le u(\eta)=\log M_k(0)+(1+\eta)\log(1+\eta)-\eta\log\eta+(k+1)(1+\eta)\,c_{\mathrm{lcm}} .
\]
$u$ is continuous on $[0,1]$ with $u(0)=\log M_k(0)+(k+1)c_{\mathrm{lcm}}<0$ (\hb{prop:kdepth}), so $u<0$ on $[0,\eta_0]$ for a rational $\eta_0>0$. At a rational $\eta$ the equation of \hb{thm:portgauge}, cleared of denominators, is an integer polynomial equation in $a$ with the unique root in $(0,\frac1{2+\eta})$; the difference of its two sides is $1$ at $a=0$ and $-(1+\eta)^k(2+\eta)^{-k-2}<0$ at $a=\frac1{2+\eta}$, so bisection at rational points certifies the root in rational intervals as small as asked. The two displays of \hb{thm:portbudget} give the bound, and its last part the sign of the coefficient of $\Tx{k+1}{\IZ}{0}$ for even $k$.
\end{proof}

\section{Chapter~\hb{ch:rays}: rays, native gauges and tiers}\label{sec:ch35}

The proofs of this chapter rest on two readings on the sectors of the orthant: of an insertion along a ray (Lemma~\ref{lem:rays-sector}) and of the direction of a ray (Lemma~\ref{lem:ch35-direction}).

\begin{slemma}\label{lem:rays-sector}
Let a ray of a family on $n\ge2$ marks be given, $g$ an insertion as in Proposition~\hb{prop:rays-strength}, and $\pi$ an order of $\cE$ with the tails $U_l=\{\pi(l),\dots,\pi(n)\}$ and the sector coordinates $q=(q_2,\dots,q_n)\in(0,1]^{n-1}$ of Theorem~\hb{thm:exist}(ii) on its sector $H_\pi$. On $H_\pi$
\[
g=\prod_{l=2}^nq_l^{\kappa_{U_l}(g)}\,g_\pi(q),\qquad f=\prod_{l=2}^nq_l^{\kappa_{1,U_l}}\,\varphi_\pi(q),\qquad \frac{\omega_0}{dy}=\prod_{l=2}^nq_l^{\kappa_{0,U_l}}\,\omega_\pi(q),
\]
with $g_\pi=\prod_jk_{j,\pi}^{-a_j}\psi_{j,\pi}^{-b_j}$, $\varphi_\pi=\prod_jk_{j,\pi}^{-s_{1j}}\psi_{j,\pi}^{-t_{1j}}$ and $\omega_\pi=\prod_jk_{j,\pi}^{-s_{0j}}\psi_{j,\pi}^{-t_{0j}}$, the powers principal. These are smooth on a neighbourhood of $[0,1]^{n-1}$, and $g_\pi$, $\varphi_\pi$ and $|\omega_\pi|$ lie between positive constants on the cube. For a function $F$ on the orthant, homogeneous of degree $0$ and integrable on the chart, $\int_{y_{e_0}=0}F\,dy=\sum_\pi\int_{(0,1]^{n-1}}F\prod_l\frac{dq_l}{q_l}$.
\end{slemma}

\begin{proof}
In the sector coordinates $\alpha_{\pi(1)}=\rho$, $\alpha_{\pi(i)}=\rho\,q_2\cdots q_i$, one has $\alpha^u=\rho^{u_{\cE}}\prod_lq_l^{u_{U_l}}$, and by Theorem~\hb{thm:exist}(ii) $K_j=\rho^{r_j}\prod_lq_l^{d_j(U_l)}k_{j,\pi}$ and $\Psi_j=\rho^{n-r_j}\prod_lq_l^{d^\ast_j(U_l)}\psi_{j,\pi}$. The three functions are homogeneous of degree $0$, the data of $g$, of the direction and of the origin lying on the balance wall, so $\rho$ cancels, and the exponents of the $q_l$ are the walls at $U_l$ of the data $(u;\mathbf a,\mathbf b)$, $(\mvec_1;\mathbf s_1,\mathbf t_1)$ and $(\mvec_0;\mathbf s_0,\mathbf t_0)$. The sector letters are polynomials with nonnegative coefficients and positive constant terms, so on the cube they are at least their constant terms, and they are positive on a neighbourhood of it; their principal powers are smooth there, the modulus of a principal power $x^{-w}$, $x>0$, is $x^{-\re w}$, and compactness gives the bounds. The sectors cover the orthant and meet in null sets. On the chart the map $y\mapsto(y_{\pi(l-1)}-y_{\pi(l)})_{l=2}^n$ is linear with an integral inverse, $y$ being recovered from the differences and from $y_{e_0}=0$, so it keeps the Lebesgue measure, and $\log q_l=y_{\pi(l)}-y_{\pi(l-1)}$ on the sector gives $dy=\prod_ldq_l/q_l$.
\end{proof}

\begin{hproposition}[\hb{prop:rays-strength} (The relations at every strength)]
As stated.
\end{hproposition}

\begin{proof}
(i) Put $Y=g\,f^k\omega_0/dy$, homogeneous of degree $0$. By Lemma~\ref{lem:rays-sector}, on the sector of $\pi$, $|Y|=\prod_lq_l^{c_l}\,|g_\pi\varphi_\pi^k\omega_\pi|$ with $c_l=\re\kappa_{0,U_l}+\kappa_{U_l}(g)+\re(k)\kappa_{1,U_l}>0$ and the last factor bounded, every nonempty proper $U$ being a tail of some order. So $Y$ is integrable on the chart, and it tends to $0$ at infinity: a sequence leaving every compact set of the chart has, along its terms in any one sector, $\min_lq_l\to0$, and $\prod_{l'}q_{l'}^{c_{l'}}\le q_l^{c_l}$ for every $l$. By \hb{lem:insertions}(ii), $\partial_{y_e}Y=(c^{(0)}_e+k\,c^{(1)}_e+\partial_{y_e})g\cdot f^k\omega_0/dy$, and by \hb{lem:insertions}(i) $(c^{(0)}_e+k\,c^{(1)}_e+\partial_{y_e})g$ is $(m_{0,e}+k\,m_{1,e}+u_e)\,g$ less the insertions $(s_{0j}+k\,s_{1j}+a_j)\,\theta_{j,T}\,\alpha^{\one_T}K_j^{-1}g$, $T\ni e$, and $(t_{0j}+k\,t_{1j}+b_j)\,\theta_{j,T}\,\alpha^{\one_{\cE\setminus T}}\Psi_j^{-1}g$, $T\not\ni e$, whose walls are $\kappa_U(g)+|T\cap U|-d_j(U)\ge\kappa_U(g)$ and $\kappa_U(g)+|U\setminus T|-d^\ast_j(U)\ge\kappa_U(g)$ (\hb{thm:window}). So what precedes applies to them: $\partial_{y_e}Y$ is integrable, and its value converges absolutely. On the chart $\alpha_{e_0}=1$, for $e\neq e_0$, Fubini's theorem and $\int_{-R}^R\partial_{y_e}Y\,dy_e=Y|_{y_e=R}-Y|_{y_e=-R}\to0$ give $\int\partial_{y_e}Y\,dy=0$; $Y$ being homogeneous of degree $0$, $\sum_e\partial_{y_e}Y=0$ (Euler's identity), and $\partial_{y_{e_0}}Y=-\sum_{e\neq e_0}\partial_{y_e}Y$ integrates to $0$ as well. For the certificate, at such a $k$ every value involved converges absolutely and is linear in the insertion, $\langle f^jg'\rangle_k=\langle g'\rangle_{k+j}$, and the right side has the value $0$.

(ii) By Lemma~\ref{lem:rays-sector}, for $\re k$ large,
\[
\langle g\rangle_k=\sum_\pi\int_{[0,1]^{n-1}}\prod_lq_l^{x_l(k)-1}h_\pi(q,k)\,dq,\qquad x_l(k)=\kappa_{0,U_l}+\kappa_{U_l}(g)+k\,\kappa_{1,U_l},\quad h_\pi=g_\pi\varphi_\pi^k\omega_\pi .
\]
$h_\pi$ is smooth in $q$ on a neighbourhood of the cube and entire in $k$, with $q$-derivatives continuous in $(q,k)$ and holomorphic in $k$. By \hb{lem:cube} each term extends meromorphically in $(x,k)$, and its recursion, applied in every coordinate, writes it on $\{\re x_l>-M\ \forall l\}$ as a finite sum of products of factors $(x_l+i)^{-1}$, at most one for each $l$, with functions holomorphic there. Along $x=x(k)$ each $x_l(k)+i$ is affine in $k$ with the slope $\kappa_{1,U_l}>0$, so the composite is meromorphic on the half-plane where every $\re x_l(k)>-M$, with poles only where some $x_l(k)$ is a nonpositive integer, of order at most $n-1$; $M$ being arbitrary, this continues $\langle g\rangle_k$ to $\mathbb C$. On the half-plane where the integral converges absolutely it is holomorphic, and it agrees there with the continuation. An element of $\mathcal L$ is a finite sum $\sum_{e,i}r_{e,i}(k)(c^{(0)}_e+k\,c^{(1)}_e+\partial_{y_e})g_{e,i}$ with insertions $g_{e,i}$ of the above kind; its value is the $\mathbb C(k)$-linear extension of $g\mapsto\langle g\rangle_k$, which does not depend on the way the element is written, and by (i), the walls of the direction being positive, it vanishes for $\re k$ large off the poles of the $r_{e,i}$, hence identically. $\langle fg\rangle_k$ and $\langle g\rangle_{k+1}$ are one integral for $\re k$ large. A recurrence that follows from the relations is the value of the element $\sum_jp_j(k)f^jg$ of $\mathcal L$.
\end{proof}

\begin{hproposition}[\hb{prop:rays-certtop} (The top and the bottom of a certificate)]
As stated.
\end{hproposition}

\begin{proof}
For a polynomial $G=\sum_ik^iG^{[i]}$, with $G^{[-1]}=0$, the coefficient of $k^m$ in $(c^{(0)}_e+k\,c^{(1)}_e+\partial_{y_e})G$ is $(c^{(0)}_e+\partial_{y_e})G^{[m]}+c^{(1)}_eG^{[m-1]}$. At $m=d$, where every $G^{[d]}_e$ vanishes, this gives the first identity, the left side having the coefficient $\sum_jc_jf^jg=\chi(f)\,g$; at $m=0$ it gives the second. At a native gauge every $c^{(1)}_e$ vanishes (Definition~\hb{def:raygauge}), so $\chi\big(f(\alpha_c)\big)g(\alpha_c)=0$. For the origin $0$, $c^{(0)}=0$.
\end{proof}

\begin{slemma}\label{lem:ch35-direction}
Let a ray of a family on $n\ge2$ marks have $\kappa_{1,S}>0$ for every nonempty proper $S$, with its stride $f$ on the chart $y_{e_0}=0$.
\begin{enumerate}[label=(\roman*)]
\item $\mvec_1>0$, the exponents of the direction are not all zero, and the tables with $s_{1j}+t_{1j}>0$ form a connected family.
\item The direction has a native gauge $\alpha^\star\in\mathbb R^{\cE}_{>0}$, unique up to scale; with $\alpha^\star_{e_0}=1$, $0<f\le f(\alpha^\star)$ on the chart with equality only at $\log\alpha^\star$, and $f\le C'e^{-c'\|y\|}$ there for some $c',C'>0$.
\item The Hessian of $-\log f$ at $\log\alpha^\star$ on the chart is
\[
H=\Big(\sum_j\big(s_{1j}\,\kappa_{j,\alpha^\star}(X_e,X_{e'})+t_{1j}\,\kappa_{j,(\alpha^\star)^{-1}}(X_e,X_{e'})\big)\Big)_{e,e'\neq e_0},
\]
and it is positive definite.
\item $\langle1\rangle_k=P_F(\mvec_0+k\mvec_1;\mathbf s_0+k\mathbf s_1,\mathbf t_0+k\mathbf t_1)$, and it converges absolutely exactly when $\re\kappa_{0,S}+k\kappa_{1,S}>0$ for every nonempty proper $S$.
\end{enumerate}
\end{slemma}

\begin{proof}
(i) Every $\{e\}$ is proper, and $\kappa_{1,\{e\}}=m_{1,e}-\sum_j\big(s_{1j}d_j(\{e\})+t_{1j}d^\ast_j(\{e\})\big)$ with $s_{1j},t_{1j}\in\mathbb N$ and $d_j,d^\ast_j\ge0$, so $m_{1,e}\ge\kappa_{1,\{e\}}>0$. If every exponent of the direction were $0$, the balance wall would give $m_{1,\cE}=0$. On the balance wall \hb{prop:compoundwall} gives $\kappa_{1,S}+\kappa_{1,\cE\setminus S}=\sum_j(s_{1j}+t_{1j})c_j(S)$, so every nonempty proper $S$ has $c_j(S)>0$ for a table with $s_{1j}+t_{1j}>0$.

(ii) The direction is real data on its balance wall, where the target is the masses, with nonnegative exponents not all zero and positive masses. By (i) and \hb{thm:saddle}(iii) it has a native gauge $\alpha^\star\in\mathbb R^{\cE}_{>0}$, unique up to scale, the unique point up to scale at which $\alpha^{\mvec_1}\prod_jK_j^{-s_{1j}}\Psi_j^{-t_{1j}}=f$ is largest. At $y=2z$ the stride is $e^{X_1(z)}$, $X_1$ the exponent of the question of the direction (\hb{def:question}), and \hb{prop:evenfloor}(i) and (ii) give $X_1(z)\le-2\kappa_{1,\min}(\max_ez_e-\min_ez_e)+C$; on the chart $|y_e|\le\max y-\min y$, since $y_{e_0}=0$.

(iii) The Hessian of $\log K_j(e^y)$ in $y$ is the covariance of the indicators of the basis under the basis law at $\alpha=e^y$ (\hb{thm:innerMGF}), and $\Psi_j(e^y)=e^{\ip\one y}K_j(e^{-y})$, so the Hessian of $\log\Psi_j(e^y)$ is the covariance under the law at $\alpha^{-1}$; at $\alpha^\star$ this gives $H$. For $v$ with $v_{e_0}=0$,
\[
v^{\mathsf T}Hv=\sum_j\Big(s_{1j}\operatorname{Var}_{\mathbb P_{j,\alpha^\star}}\ip{\one_T}v+t_{1j}\operatorname{Var}_{\mathbb P_{j,(\alpha^\star)^{-1}}}\ip{\one_T}v\Big)\ge0,
\]
both laws charging every basis. Let $v^{\mathsf T}Hv=0$ with $v\neq0$. Then $v$ is not constant, and $\ip{\one_T}v$ is the same at every basis of every table with $s_{1j}+t_{1j}>0$. Let $S=\{e:v_e=\max v\}$, nonempty and proper, and take by (i) such a table with $c_j(S)>0$ and a basis $T$ with $|T\cap S|<r_j(S)$. $T\cap S$ is independent in the restriction to $S$ and not a basis of it, so some $f\in S\setminus T$ keeps it independent; the circuit of $T+f$ through $f$ is not contained in $S$, and with $e$ a mark of it outside $S$, $T-e+f$ is a basis with $\ip{\one_{T-e+f}}v-\ip{\one_T}v=v_f-v_e>0$, a contradiction.

(iv) On the chart $\alpha_{e_0}=1$, $\alpha^{\mvec-\one}\prod_{e\neq e_0}d\alpha_e=\alpha^{\mvec}dy$, so the value $\langle1\rangle_k=\int f^k\omega_0$ of \hb{def:ray} is the period at the point $k$ of the ray in the form of \hb{thm:projective}, which converges exactly when the real parts of its proper walls $\kappa_{0,S}+k\kappa_{1,S}$ are positive; the modulus of its integrand is the integrand at the real parts of the data.
\end{proof}

\begin{hproposition}[\hb{prop:nativeorder} (The native value bounds the order)]
As stated: every ray with positive direction walls, its origin complex.
\end{hproposition}

\begin{proof}
By Lemma~\ref{lem:ch35-direction}(iv), $\langle1\rangle_k$ converges absolutely for $k\ge k_1$. The amplitude $g_0=\omega_0/dy=\alpha^{\mvec_0}\prod_jK_j^{-s_{0j}}\Psi_j^{-t_{0j}}$, the powers of the positive letters principal, is continuous and nowhere zero on the chart. Let $y^\star=\log\alpha^\star$ with $\alpha^\star_{e_0}=1$. By Lemma~\ref{lem:ch35-direction}(ii) and (iii), $f/f(\alpha^\star)$ is real-analytic, has its unique maximum $1$ at $y^\star$, nondegenerate with the Hessian $H$, and tends to $0$ at infinity; so off a ball $B$ about $y^\star$, $f\le e^{-\delta}f(\alpha^\star)$ for some $\delta>0$, and for $k\ge k_1$
\[
\Big|\int_{y_{e_0}=0,\ y\notin B}f^kg_0\,dy\Big|\le e^{-\delta(k-k_1)}f(\alpha^\star)^{k-k_1}\int f^{k_1}|g_0|\,dy .
\]
On $B$, Laplace's method with the continuous amplitude $g_0$ gives
\[
\int_Bf^kg_0\,dy=f(\alpha^\star)^k\Big(\frac{2\pi}k\Big)^{\frac{n-1}2}(\det H)^{-\frac12}\big(g_0(y^\star)+o(1)\big).
\]
Hence $\langle1\rangle_k=f(\alpha^\star)^k\big(\frac{2\pi}k\big)^{\frac{n-1}2}(\det H)^{-\frac12}g_0(y^\star)(1+o(1))$ with $g_0(y^\star)\neq0$: $\langle1\rangle_k\neq0$ for all large $k$ and $\langle1\rangle_{k+1}/\langle1\rangle_k\to f(\alpha^\star)$. For a real origin $g_0>0$, and $\langle1\rangle_k>0$ wherever it converges.

Let $\sum_{j=0}^hp_j(k)\langle1\rangle_{k+j}=0$ for all large $k$, the $p_j\in\mathbb Q[k]$ not all zero, $d$ their largest degree and $c_j$ the coefficient of $k^d$ in $p_j$. Dividing by $k^d\langle1\rangle_k$ and letting $k\to\infty$ gives $\sum_jc_jf(\alpha^\star)^j=0$, and the symbol $\sum_jc_j\lambda^j$ is a nonzero rational polynomial of degree at most $h$. So $f(\alpha^\star)$ is algebraic, its minimal polynomial divides the symbol, every conjugate is a root of it, and $h\ge[\mathbb Q(f(\alpha^\star)):\mathbb Q]$. For $h=1$ this degree is one, and a transcendental native value admits no $h$.
\end{proof}

\begin{hproposition}[\hb{prop:rays-conjugates} (The conjugates of the native value of the direction)]
As stated.
\end{hproposition}

\begin{proof}
On the chart $\alpha_{e_0}=1$ the letter complement $X_F$ is the affine variety $\{\alpha\in(\mathbb C^\times)^{\cE-e_0}:\ \prod_jK_j(\alpha)\Psi_j(\alpha)\neq0\}$, and the defects $c^{(1)}_e=m_{1,e}-\sum_j\big(s_{1j}R_{j,e}+t_{1j}R^\ast_{j,e}\big)$ are regular functions on it with coefficients in $\Bbbk$, the direction being integral. Their sum over $e$ is $\sum_e\partial_{y_e}\log f=0$, $f$ being homogeneous of degree $0$, so the native gauges are the common zeros in $X_F$ of the $c^{(1)}_e$, $e\neq e_0$: an algebraic set defined over $\Bbbk$. Its Jacobian matrix in the coordinates $y_e$, $e\neq e_0$, is the Hessian of $\log f$, nonsingular at $\alpha^\star$ (Lemma~\ref{lem:ch35-direction}(iii)); by the holomorphic inverse function theorem $\alpha^\star$ is an isolated point of it, and an isolated point of an algebraic set defined over $\Bbbk$ has coordinates algebraic over $\Bbbk$, the irreducible components of the set being defined over the algebraic closure of $\Bbbk$. An embedding $\tau$ over $\Bbbk$ keeps every polynomial identity with coefficients in $\Bbbk$: $\tau(\alpha^\star)$ satisfies the equations of nativity, no letter vanishes at it, the Hessian determinant, a rational function over $\Bbbk$, takes there the value $\tau$ of a nonzero number, and $f\big(\tau(\alpha^\star)\big)=\tau\big(f(\alpha^\star)\big)$. Every conjugate of $f(\alpha^\star)$ over $\Bbbk$ is $\sigma\big(f(\alpha^\star)\big)$ for an embedding $\sigma$ of $\Bbbk\big(f(\alpha^\star)\big)$, which extends to an embedding of the finite extension $\Bbbk(\alpha^\star)$ in the algebraically closed $\mathbb C$. The $[\Bbbk(\alpha^\star):\Bbbk]$ embeddings of $\Bbbk(\alpha^\star)$ differ on a coordinate of $\alpha^\star$, so they give as many distinct nondegenerate native gauges, and Theorem~\hb{thm:gaugesolution}(iii) bounds the dimension. Corollary~\hb{cor:symbol} at $g=1$ gives the last statement.
\end{proof}

\begin{hproposition}[\hb{prop:ra-framegrowth} (The growth of the frame)]
As stated.
\end{hproposition}

\begin{proof}
By \hb{prop:frame}, $C_k$ is $C_{k_0}((-1)^hc)^{k-k_0}$ times $\prod_{i\le d_0}\Gamma(k-a_i)/\prod_{i\le d_h}\Gamma(k-b_i)$ times a nonzero constant. Stirling's formula, which applies to $\Gamma(k-a)$ for fixed complex $a$ as $k\to+\infty$ along the reals, the argument staying in a sector about the positive axis, gives $\log|\Gamma(k-a)|=(k-\re a-\frac12)\log k-k+O(1)$; summing, the Gammas contribute $(d_0-d_h)(k\log k-k)+O(\log k)$. When $d_0=d_h=d$ is the largest degree, $c_0$ and $c_h$ are the leading coefficients of $p_0$ and $p_h$, both nonzero, the symbol has degree $h$, and the product of its roots is $(-1)^hc_0/c_h=(-1)^hc$; the term in $k\log k$ is absent and $\frac1k\log|C_k|\to\log|c|$.
\end{proof}

\begin{hdefinition}[\hb{def:centreray} (The centre ray)]
As stated, for a family connected at strengths $\sigma_j\in\mathbb N$.
\end{hdefinition}

\begin{proof}
The direction $(\sigma\one;\boldsymbol\sigma,\boldsymbol\sigma)$ is integral, and it is on the balance wall: $m_{1,\cE}=n\sigma=\sum_j\sigma_j(r_j+r^\ast_j)$. The origin $0$ is on it too. From $\Psi_j(\alpha)=\alpha^{\one}K_j(\alpha^{-1})$, $K_j\Psi_j=\alpha^{\one}K_j(\alpha)K_j(\alpha^{-1})$, and with $\sum_j\sigma_j=\sigma$ the stride is $\alpha^{\sigma\one}\prod_j(K_j\Psi_j)^{-\sigma_j}=\prod_j\big(K_j(e^y)K_j(e^{-y})\big)^{-\sigma_j}$ at $\alpha=e^y$. The origin form is $dy$, so $\langle1\rangle_k=\int_{y_{e_0}=0}f^kdy=P_F(k\sigma\one;k\boldsymbol\sigma,k\boldsymbol\sigma)$ by Lemma~\ref{lem:ch35-direction}(iv). The direction walls are $\sigma|S|-\sum_j\sigma_j\big(d_j(S)+d^\ast_j(S)\big)=\sum_j\sigma_jc_j(S)$, positive on every nonempty proper $S$ exactly when the family is connected at the strengths (\hb{prop:centrefamily}(ii)).
\end{proof}

\begin{htheorem}[\hb{thm:pencil} (A ray is one function)]
As stated, for every ray with real origin and positive direction walls.
\end{htheorem}

\begin{proof}
(i) The bounds on $f$ are Lemma~\ref{lem:ch35-direction}(ii). For $k\ge k_0$ the walls $\kappa_{0,S}+k\kappa_{1,S}\ge\kappa_{0,S}+k_0\kappa_{1,S}$ are real and positive, so $\langle1\rangle_k$ converges (Lemma~\ref{lem:ch35-direction}(iv)), and its integrand $f^k\omega_0/dy$ is positive.

(ii) For $|q|<f(\alpha^\star)^{-1}$, $|qf|\le|q|f(\alpha^\star)<1$ on the chart, so $\sum_{k\ge k_0}q^{k-k_0}f^k=f^{k_0}/(1-qf)$, the partial sums dominated by $f^{k_0}/(1-|q|f(\alpha^\star))$ times the positive density $\omega_0/dy$, which is integrable at $k=k_0$. Summing under the integral gives the first form of $G$, and the second is the first with numerator and denominator multiplied by $\prod_jK_j^{s_{1j}}\Psi_j^{t_{1j}}$. On smaller discs the same domination allows differentiation under the integral, $\partial_q^l(1-qf)^{-1}=l!\,f^l(1-qf)^{-l-1}$. For real $q<f(\alpha^\star)^{-1}$, $1-qf\ge\min\big(1,1-qf(\alpha^\star)\big)>0$, so the integrand is positive and at most a multiple of $f^{k_0}\omega_0/dy$.

(iii) With the positive measure $\mu=f^k\omega_0$ on the chart, $\langle1\rangle_{k+i+i'}=\int f^{i+i'}d\mu$, and the identity of Andr\'eief gives
\[
\det\Big[\int f^{i+i'}d\mu\Big]_{i,i'<h}=\frac1{h!}\int\det\big[f(y_a)^i\big]_{i<h,\,a\le h}^2\prod_ad\mu(y_a)=\frac1{h!}\int\prod_{a<b}\big(f(y_a)-f(y_b)\big)^2\prod_ad\mu(y_a),
\]
positive because $f$ is real-analytic and not constant, so that the set where two values $f(y_a)$ agree is null. If $G=N_1/N_2$ were rational, $N_2(0)\neq0$, $d=\deg N_2$, then $\sum_{i=0}^d[q^i]N_2\,\langle1\rangle_{k_0+m-i}=0$ for $m>\deg N_1$, a relation with $[q^0]N_2\neq0$ among the rows of the Hankel matrix $\big[\langle1\rangle_{k+l+i'}\big]_{l,i'\le d}$ for all large $k$, whose determinant would vanish, against (iii) at $h=d+1$.

On the centre ray $\kappa_{0,S}=0<\kappa_{1,S}$, so $k_0=1$, and $\alpha^\star=\one$ by \hb{prop:centrefamily}(i).
\end{proof}

\begin{htheorem}[\hb{thm:branch} (The branch at the native gauge of the direction)]
As stated, for every ray with real origin and positive direction walls.
\end{htheorem}

\begin{proof}
Put $y^\star=\log\alpha^\star$, $\varphi=-\log\big(f/f(\alpha^\star)\big)\ge0$ on the chart, zero only at $y^\star$, with the Hessian $H$ there, positive definite (Lemma~\ref{lem:ch35-direction}(iii)), and $g=f^{k_0}\omega_0/dy$, positive, real-analytic and integrable. At $q=A(1-\varepsilon)$, $1-qf=1-(1-\varepsilon)e^{-\varphi}$, so for $0<\varepsilon\le1$
\[
G(q)=\int\frac{g\,dy}{1-(1-\varepsilon)e^{-\varphi}} .
\]
Fix a small $t_1>0$. On $\{\varphi\ge t_1\}$, $|1-(1-\varepsilon)e^{-\varphi}|\ge1-|1-\varepsilon|e^{-t_1}>0$ for $|\varepsilon|$ small, so that part of the integral is holomorphic at $\varepsilon=0$. Near $y^\star$ the real-analytic Morse lemma gives coordinates $v\in\mathbb R^{2a}$, $2a=n-1$, with $\varphi=\frac12|v|^2$ and $|\det\partial y/\partial v|=(\det H)^{-1/2}$ at $v=0$; for $t_1$ small, $\{\varphi<t_1\}$ lies in their domain. With $\tilde g=g\,|\det\partial y/\partial v|$ and
\[
\beta(t)=2^{a-1}\int_{S^{2a-1}}\tilde g\big(\sqrt{2t}\,\omega\big)\,d\omega ,
\]
analytic near $[0,t_1]$, the spherical mean being an even analytic function of the radius, $\int_{\varphi<t}g\,dy$ has the derivative $t^{a-1}\beta(t)$. In $w=e^t-1$, $1-(1-\varepsilon)e^{-t}=e^{-t}(w+\varepsilon)$ and $e^tdt=dw$, so the part over $\{\varphi<t_1\}$ is
\[
\int_0^{w_1}\frac{w^{a-1}\gamma(w)}{w+\varepsilon}\,dw,\qquad\gamma(w)=\Big(\frac{\log(1+w)}w\Big)^{a-1}\beta\big(\log(1+w)\big),\quad\gamma(0)=\beta(0),
\]
with $\gamma=\sum_m\gamma_mw^m$ on a disc of radius $R'>w_1$, $t_1$ being small, so that $|\gamma_m|\le C{R'}^{-m}$. For $c>0$ not an integer and $|\varepsilon|<w_1$, $\varepsilon\notin(-\infty,0]$,
\[
\int_0^{w_1}\frac{w^{c-1}\,dw}{w+\varepsilon}=\frac\pi{\sin\pi c}\,\varepsilon^{c-1}-\sum_{l\ge0}\frac{(-\varepsilon)^lw_1^{c-1-l}}{l+1-c},
\]
from $\int_0^\infty w^{c-1}dw/(w+\varepsilon)=\frac\pi{\sin\pi c}\varepsilon^{c-1}$ at $0<\re c<1$ and continuation in $c$; for an integer $c=N\ge1$, $w^{N-1}=\big((w+\varepsilon)-\varepsilon\big)^{N-1}$ gives the integral as $(-1)^N\varepsilon^{N-1}\log\varepsilon$ plus $(-\varepsilon)^{N-1}\log(w_1+\varepsilon)$ plus a polynomial in $\varepsilon$, the two last of modulus $O\big((w_1+2\rho)^{N-1}\big)$ on $|\varepsilon|\le\rho$. Summed with the weights $\gamma_m$ at $c=a+m$, the singular terms give $\varepsilon^{a-1}V(\varepsilon)$ for $n$ even and $\varepsilon^{a-1}\log\varepsilon\,V(\varepsilon)$ for $n$ odd, with
\[
V(\varepsilon)=\frac\pi{\sin\pi a}\,\gamma(-\varepsilon)\quad(n\text{ even}),\qquad V(\varepsilon)=(-1)^a\gamma(-\varepsilon)\quad(n\text{ odd}),
\]
and the remainders converge to a function holomorphic near $0$, since $|l+1-a-m|\ge\frac12$ for $n$ even and $w_1+2\rho<R'$ for small $\rho$. This is the form of $G$, with $U$ the sum of the remainders and the part over $\{\varphi\ge t_1\}$. Further $\beta(0)=2^{a-1}|S^{2a-1}|\,g(\alpha^\star)(\det H)^{-1/2}=(2\pi)^ag(\alpha^\star)/\big(\Gamma(a)\sqrt{\det H}\big)$, $\frac\pi{\sin\pi a}=\Gamma(a)\Gamma(1-a)$, and $g(\alpha^\star)=f(\alpha^\star)^{k_0}(\omega_0/dy)(\alpha^\star)=(f^{k_0-1}\omega_0/dy)(\alpha^\star)/A$, which gives $V(0)$.

$G^{(l)}(A)=l!\int f^{k_0+l}\omega_0/(1-Af)^{l+1}$ is a positive integral, and near $y^\star$, $1-Af=1-e^{-\varphi}$ is comparable with $|v|^2$, so it is finite exactly when $2(l+1)<2a$. For $l<a-1$, $G^{(l)}(q)\to G^{(l)}(A)$ as $q\uparrow A$ by monotone convergence, while the $l$-th derivative of $\varepsilon^{a-1}V$, resp.\ of $\varepsilon^{a-1}\log\varepsilon\,V$, tends to $0$; so the Taylor coefficients of $U$ of order $l<a-1$ are those of $G$ at $A$, with $q-A=-A\varepsilon$, and the rest of $U$ and $\varepsilon^{a-1}(V-V(0))$ are $O(\varepsilon^{\lceil a-1\rceil})$. On the centre ray $k_0=1$, $\alpha^\star=\one=(\alpha^\star)^{-1}$, $A=f(\one)^{-1}=\prod_jK_j(\one)^{2\sigma_j}$, $f^0\omega_0/dy=1$ and $H=2\sum_j\sigma_j\big(\kappa_{j,\one}(X_e,X_{e'})\big)_{e,e'\neq e_0}$.
\end{proof}

\begin{hproposition}[\hb{prop:rays-centrenative} (The native gauges of a centre ray)]
As stated.
\end{hproposition}

\begin{proof}
The direction is $(\sigma\one;\boldsymbol\sigma,\boldsymbol\sigma)$, so $c^{(1)}_e=\sigma-\sum_j\sigma_j\big(R_{j,e}+R^\ast_{j,e}\big)$ (Definition~\hb{def:ray}). By \hb{prop:identity}, $\Psi_j(\alpha)=\alpha^{\cE}K_j(\alpha^{-1})$, so $R^\ast_{j,e}(\alpha)=\alpha_e\partial_{\alpha_e}\log\Psi_j(\alpha)=1-R_{j,e}(\alpha^{-1})$ as rational functions, the identity of \hb{thm:innerMGF} at real gauges; with $\sum_j\sigma_j=\sigma$ this gives the defects. The stride is $\prod_j\big(K_j(\alpha)K_j(\alpha^{-1})\big)^{-\sigma_j}$, a rational function homogeneous of degree $0$ (Definition~\hb{def:centreray}).
\end{proof}

\begin{htheorem}[\hb{thm:signgauges} (Native gauges at the signs)]
As stated, for the centre ray of a family.
\end{htheorem}

\begin{proof}
(i) On the chart $\alpha_{e_0}=1$, $f(\alpha)=\prod_j\big(K_j(\alpha)K_j(\alpha^{-1})\big)^{-\sigma_j}$ (\hb{def:centreray}), so $f(\alpha^{-1})=f(\alpha)$; $\alpha\mapsto\alpha^{-1}$ keeps the letter complement, since $K_j(\alpha^{-1})\Psi_j(\alpha^{-1})=\alpha^{-2\one}\Psi_j(\alpha)K_j(\alpha)$. A point of $(\mathbb C^\times)^{\cE}/\mathbb C^\times$ fixed by the inversion has $\alpha_e^2$ independent of $e$, hence $\alpha_e=\pm1$ on the chart: the $2^{n-1}$ signs. By \hb{lem:insertions}(ii) the native gauges of the ray are the critical points of $\log f$ on $X_F$, and the inversion, which keeps $f$, permutes them and keeps their native values.

(ii) At a sign $\eta$, $\Psi_j(\eta)=\eta^{\one}K_j(\eta)$, so $K_j(\eta)\neq0$ for every $j$ puts $\eta$ in $X_F$. In the chart $\alpha=\eta e^y$, $y_{e_0}=0$, $\alpha^{-1}=\eta e^{-y}$, so $y\mapsto\log f(\eta e^y)$ is even and its gradient vanishes at $y=0$: $\eta$ is a native gauge, of native value $\prod_j\big(K_j(\eta)K_j(\eta)\big)^{-\sigma_j}$. With $K_j(\eta e^y)=\sum_T\theta_{j,T}\eta^Te^{\ip{\one_T}y}=K_j(\eta)\sum_T\mathbb P_{j,\eta}(T)e^{\ip{\one_T}y}$, the Hessian of $\log K_j(\eta e^y)$ at $0$ is the covariance of the indicators under the signed measure $\mathbb P_{j,\eta}$ of total mass one, and $\log K_j(\eta e^{-y})$ has the same Hessian; so the Hessian of $-\log f$ is $2\sum_j\sigma_j\kappa_{j,\eta}$.

(iii) $W^{-1}=\prod_j(K_j\Psi_j)^{\sigma_j}$ is a polynomial, and $\partial_{y_e}\log f=\big(\sigma W^{-1}-\alpha_e\partial_eW^{-1}\big)/W^{-1}$. For a sign $\eta$, $h(y)=W^{-1}(\eta e^y)\,e^{-\sigma\ip\one y}=\prod_j\big(\eta^{\one}K_j(\eta e^y)K_j(\eta e^{-y})\big)^{\sigma_j}$ is even in $y$, so at $y=0$
\[
\partial_{y_e}h=\big(\alpha_e\partial_eW^{-1}-\sigma W^{-1}\big)(\eta)=0
\]
at every sign, whether or not the $K_j(\eta)$ vanish, and $W^{-1}(\eta)=\eta^{\sigma\one}\prod_jK_j(\eta)^{2\sigma_j}$. The ring $\mathbb Q(\boldsymbol\theta)[\alpha_e:e\neq e_0]/(\alpha_e^2-1)$ is the product of copies of $\mathbb Q(\boldsymbol\theta)$ indexed by the signs, a class being its values at the signs, with the idempotents $2^{1-n}\prod_{e\neq e_0}(1+\eta_e\alpha_e)$, which take the value $1$ at $\eta$ and $0$ at the other signs. So the critical polynomials vanish in it, and $W^{-1}\equiv\alpha^{\sigma\one}\prod_jK_j^{2\sigma_j}$. Localizing at $\prod_jK_j$ keeps exactly the factors of the signs with $\prod_jK_j(\eta)\neq0$, whose idempotents form a basis, and multiplication by $W^{-1}/\alpha^{\sigma\one}$, $\alpha^{\sigma\one}$ being a unit, acts on the idempotent of $\eta$ by its value $\prod_jK_j(\eta)^{2\sigma_j}$.
\end{proof}

\begin{htheorem}[\hb{thm:separatesigns} (Weights that separate the signs)]
As stated, for a family of $N$ tables connected at its strengths.
\end{htheorem}

\begin{proof}
Treat the weights of all tables as variables. For a sign $\eta$, $K_j(\eta)=\sum_T\theta_{j,T}\eta^T$ is a nonzero linear form. Write $M_j=K_j(\eta)^2\big(\kappa_{j,\eta}(X_e,X_{e'})\big)_{e,e'\neq e_0}$, whose entries $K_j(\eta)\sum_{T\ni e,e'}\theta_{j,T}\eta^T-\sum_{T\ni e}\theta_{j,T}\eta^T\sum_{T\ni e'}\theta_{j,T}\eta^T$ are quadratic, and
\[
D_\eta=\det\Big[\sum_j\sigma_jM_j\prod_{i\neq j}K_i(\eta)^2\Big]=\prod_iK_i(\eta)^{2(n-1)}\det\Big[\sum_j\sigma_j\kappa_{j,\eta}\Big],
\]
of degree $2N(n-1)$. At $\eta=\one$ and unit weights $D_{\one}\neq0$: the Hessian $2\sum_j\sigma_j\kappa_{j,\one}$ is positive definite, by the argument of Lemma~\ref{lem:ch35-direction}(iii) at the strengths. The substitution $\theta_{j,T}\mapsto\theta_{j,T}\eta^T$ carries $K_j(\one)$ to $K_j(\eta)$ and $\mathbb P_{j,\one}$ to $\mathbb P_{j,\eta}$, so $D_\eta$ is $D_{\one}$ after an invertible linear substitution, and $D_\eta\neq0$. For signs $\eta\neq\eta'$, $\prod_jK_j(\eta)^{2\sigma_j}-\prod_jK_j(\eta')^{2\sigma_j}$, of degree $2\sigma$, is not zero: otherwise unique factorization, the linear forms $K_j(\eta)$ being irreducible and in disjoint variables for distinct $j$, makes $K_j(\eta)$ a multiple of $K_j(\eta')$ for every $j$ with $\sigma_j>0$, that is $\eta^T\eta'^T$ constant on the bases of such a table; with $S=\{e:\eta_e\neq\eta'_e\}$, nonempty and proper since $\eta_{e_0}=\eta'_{e_0}$, the parity of $|T\cap S|$ is then constant, so the window $\{d_j(S),\dots,r_j(S)\}$ of \hb{thm:window} is one value and $c_j(S)=0$ for every $j$ with $\sigma_j>0$, against connectedness. The product
\[
\prod_\eta\Big(D_\eta\prod_jK_j(\eta)\Big)\prod_{\{\eta,\eta'\}}\Big(\prod_jK_j(\eta)^{2\sigma_j}-\prod_jK_j(\eta')^{2\sigma_j}\Big)
\]
is a nonzero polynomial of degree $2^{n-1}\big(N+2N(n-1)\big)+2\sigma\binom{2^{n-1}}2=2^{n-1}\big((2n-1)N+(2^{n-1}-1)\sigma\big)$, so by the Schwartz--Zippel bound it does not vanish at some point of $\{1,\dots,D+1\}$ in every variable, $D$ that degree. At such weights every $K_j(\eta)\neq0$, every sign is a native gauge by \hb{thm:signgauges}(ii), nondegenerate since $D_\eta\neq0$, and the native values are distinct. \hb{thm:gaugesolution}(iii) bounds the dimension below by $2^{n-1}$, and \hb{cor:symbol} at $g=1$ the order of every recurrence that follows from the relations.
\end{proof}

\begin{hlemma}[\hb{lem:wallgauges} (Where native gauges can reach a wall)]
As stated, for a family.
\end{hlemma}

\begin{proof}
(i) On the balance wall the target of the direction is $\mvec_1$, so nativity at $\alpha$ reads $\mvec_1=\sum_j\big(s_{1j}R_j(\alpha)+t_{1j}R^\ast_j(\alpha)\big)$ (\hb{def:native}). $R_{j,e}=\alpha_e\partial_eK_j/K_j=\sum_T\one_T(e)\,\mathbb P_{j,\alpha}(T)$ and $R^\ast_{j,e}=\sum_T\one_{\cE\setminus T}(e)\,\mathbb P^\ast_{j,\alpha}(T)$, and summing over $e\in B$ gives (i).

(ii) In the coordinates $(\rho,u,\alpha_R)$, $\alpha^T=\rho^{|T\cap B|}u^{T\cap B}\alpha_R^{T\cap R}$, so $\sum_T|T\cap B|\,\theta_{j,T}\alpha^T=\rho\,\partial_\rho K_j$, and with $\widetilde K_j=\sum_l\rho^lK^{[l]}_j$, $\widetilde\Psi_j=\sum_l\rho^l\Psi^{[l]}_j$,
\[
\kappa_{1,B}=\sum_j\Big(s_{1j}\,\frac{\rho\,\partial_\rho\widetilde K_j}{\widetilde K_j}+t_{1j}\,\frac{\rho\,\partial_\rho\widetilde\Psi_j}{\widetilde\Psi_j}\Big)
\]
at every native gauge of the direction, by (i) and $\kappa_{1,B}=m_{1,B}-\sum_j\big(s_{1j}d_j(B)+t_{1j}d^\ast_j(B)\big)$. Along the sequence the exponents stay bounded, and $\rho_i\partial_\rho\widetilde K_j=O(\rho_i)$. If no bottom letter vanished at $(u^\ast,\alpha^\ast_R)$, every quotient would be $O(\rho_i)$ and the right side would tend to $0$, while the left side tends to $\kappa_{1,B}(\mathbf v)>0$. So some bottom letter vanishes there. Let it be the only one, say $K^{[0]}_{j_0}$. Then the terms of the other letters tend to $0$, so $s_{1j_0}(\mathbf v_i)\,\rho_i\partial_\rho\widetilde K_{j_0}/\widetilde K_{j_0}\to\kappa_{1,B}(\mathbf v)$, and since $\rho_i\partial_\rho\widetilde K_{j_0}\to0$, $|s_{1j_0}(\mathbf v_i)/\widetilde K_{j_0}|\to\infty$. Nativity in a logarithmic coordinate $D$ of $u$ or of $\alpha_R$ reads
\[
D\log\big(u^{\mvec_{1}|_B}\alpha_R^{\mvec_1|_R}\big)=\sum_j\Big(s_{1j}\frac{D\widetilde K_j}{\widetilde K_j}+t_{1j}\frac{D\widetilde\Psi_j}{\widetilde\Psi_j}\Big);
\]
the left side and the terms of the other letters stay bounded, so $s_{1j_0}D\widetilde K_{j_0}/\widetilde K_{j_0}$ does, and $D\widetilde K_{j_0}(\rho_i,u_i,\alpha_{R,i})\to DK^{[0]}_{j_0}(u^\ast,\alpha^\ast_R)$ must vanish. The bottom letter is homogeneous in $u$ and in $\alpha_R$, so with $K^{[0]}_{j_0}(u^\ast,\alpha^\ast_R)=0$ Euler's identity gives the derivatives along the scalings of $u$ and of $\alpha_R$ as well: the whole gradient vanishes. The same holds for a cobasis letter. If two or more bottom letters vanish, every term of the gradient of the product contains a vanishing factor; if one does, its gradient vanishes. In both cases the product vanishes at $(u^\ast,\alpha^\ast_R)$ with its gradient.
\end{proof}

\begin{hproposition}[\hb{prop:oneslot} (Small lower coefficients on one parameter)]
As stated, with $p$ letters, for real $C_i$ bounded below by a positive constant on $[0,\infty)$.
\end{hproposition}

\begin{proof}
For $\rho\ge0$ and $\lambda>0$, $\lambda\xi_i+\rho C_i(\rho)$ is $\lambda\xi_i$ plus a nonnegative number, so it stays in $\mathbb C\setminus(-\infty,0]$, the powers are principal, and the integrand is continuous on $(0,\infty)$. Put $F_{\mathbf a}=G\prod_iC_i^{-w_i-a_i}$ and $c_{\mathbf a}(\lambda)=\prod_i\binom{-w_i}{a_i}(\lambda\xi_i)^{a_i}$, so that, wherever $|\lambda\xi_i|<|\rho C_i(\rho)|$ for every $i$,
\[
\rho^{\kappa-1}G\prod_i\big(\lambda\xi_i+\rho C_i\big)^{-w_i}=\rho^{\kappa^\flat-1}F_{\mathbf 0}\prod_i\Big(1+\frac{\lambda\xi_i}{\rho C_i}\Big)^{-w_i}=\sum_{\mathbf a}c_{\mathbf a}(\lambda)\,\rho^{\kappa^\flat-|\mathbf a|-1}F_{\mathbf a},
\]
absolutely and uniformly where $|\lambda\xi_i|\le\frac12|\rho C_i(\rho)|$. Choose $b>0$ with $G$ and the $C_i$ holomorphic and the $C_i$ zero-free on a neighbourhood of $|\rho|\le b$, and let $|C_i|\ge c>0$ on $[0,\infty)$ and on $|\rho|\le b$. For $|\lambda|\le\lambda_0=bc/(2\max_i|\xi_i|)$ each letter $\lambda\xi_i+\rho C_i(\rho)$ has exactly one root in $|\rho|<b$ (Rouch\'e, against $\rho C_i(\rho)$) and none on $|\rho|=b$ or on $[b,\infty)$.

\emph{The analytic part.} Split the integral at $b$. On $[b,\infty)$ the expansion above holds uniformly and its terms are dominated by $\sum_{\mathbf a}\prod_i|\binom{-w_i}{a_i}|2^{-|\mathbf a|}$ times $|\rho^{\kappa^\flat-1}F_{\mathbf 0}|$, integrable at $\infty$, so $\int_b^\infty=\sum_{\mathbf a}c_{\mathbf a}(\lambda)\int_b^\infty\rho^{\kappa^\flat-|\mathbf a|-1}F_{\mathbf a}\,d\rho$, holomorphic for $|\lambda|<\lambda_0$. Let $\oint$ start at $\rho=b$ on the principal branches and run once around $|\rho|=b$; the integrand continues along it, and the expansion holds on the circle, each $\rho^{\kappa^\flat-|\mathbf a|-1}$ continued along it and $F_{\mathbf a}$ single-valued on the disc. For $\re c'>0$ and $F$ holomorphic near the closed disc,
\[
\oint_{|\rho|=b}\rho^{c'-1}F\,d\rho=\big(e^{2\pi ic'}-1\big)\int_0^b\rho^{c'-1}F\,d\rho ,
\]
by Cauchy's theorem on the disc cut along $[0,b]$; both sides are holomorphic in $c'$ off $-\mathbb N$, the right side read with the finite part $\Reg\int_0^b$ (\hb{def:Reg}, \hb{thm:RegCT}), so the identity holds at every $c'\notin\mathbb Z$. At $c'=\kappa^\flat-|\mathbf a|$, $e^{2\pi ic'}=e^{2\pi i\kappa^\flat}\neq1$. Hence
\[
A_b(\lambda)=\frac1{e^{2\pi i\kappa^\flat}-1}\oint_{|\rho|=b}\rho^{\kappa-1}G\prod_i\big(\lambda\xi_i+\rho C_i\big)^{-w_i}d\rho=\sum_{\mathbf a}c_{\mathbf a}(\lambda)\,\Reg\!\int_0^b\rho^{\kappa^\flat-|\mathbf a|-1}F_{\mathbf a}\,d\rho
\]
is holomorphic for $|\lambda|<\lambda_0$, the circle integral of an integrand holomorphic in $\lambda$ along a contour that no root reaches, and the series converges there. Since $\Reg\int_0^\infty=\Reg\int_0^b+\int_b^\infty$, $A(\lambda)=A_b(\lambda)+\int_b^\infty\cdots$ is the series of the statement, holomorphic at $0$.

\emph{The branch part.} For $\lambda>0$ put $\rho=\lambda u$ in $\int_0^b$ and in the circle integral; $(\lambda X)^{-w}=\lambda^{-w}X^{-w}$ on the principal branches at $\lambda>0$, so both carry the factor $\lambda^{\kappa^\flat}$:
\[
\int_0^b\cdots\,d\rho-A_b(\lambda)=\lambda^{\kappa^\flat}D_\lambda(b/\lambda),\qquad D_\lambda(R)=\int_0^R\psi_\lambda\,du-\frac1{e^{2\pi i\kappa^\flat}-1}\oint_{|u|=R}\psi_\lambda\,du,
\]
with $\psi_\lambda(u)=u^{\kappa-1}G(\lambda u)\prod_i\big(\xi_i+uC_i(\lambda u)\big)^{-w_i}$. The roots of the letters in $u$ are $\rho_i(\lambda)/\lambda$, where $\rho_i(\lambda)$ is the root in $|\rho|<b$; from $|\rho_iC_i(\rho_i)|=|\lambda\xi_i|$ they have modulus at most $|\xi_i|/c$. Take $R_0>2\max_i|\xi_i|/c$, so that $|uC_i(\lambda u)|>2|\xi_i|$ on $|u|=R$ for $R\ge R_0$. For $R_0\le R\le b/|\lambda|$, one turn around $|u|=R$ multiplies $\psi_\lambda$ by $e^{2\pi i\kappa^\flat}$, the factor $u^{\kappa-1}$ giving $e^{2\pi i\kappa}$ and each letter, with one root inside, $e^{-2\pi iw_i}$; so $\frac d{dR}\oint_{|u|=R}\psi_\lambda\,du=\big[\psi_\lambda(Re^{i\vartheta})e^{i\vartheta}\big]_{\vartheta=0}^{2\pi}=(e^{2\pi i\kappa^\flat}-1)\psi_\lambda(R)$, and $\frac d{dR}D_\lambda(R)=0$. Thus $B(\lambda)=D_\lambda(R_0)$, with $R_0$ fixed, and $\int_0^\infty=A(\lambda)+\lambda^{\kappa^\flat}B(\lambda)$ for small $\lambda>0$. For complex $\lambda$ near $0$, $\psi_\lambda$ is holomorphic in $\lambda$ on $[0,R_0]$ and on $|u|=R_0$: there $|\lambda u|\le|\lambda|R_0<b$, the letters do not vanish, $\xi_i+uC_i(\lambda u)$ stays off $(-\infty,0]$ on $[0,R_0]$ for $\lambda$ small by compactness, and $u^{\kappa-1}$ is integrable at $0$ since $\re\kappa>0$. So $B$ is holomorphic at $0$, and
\[
B(0)=G(0)\Big[\int_0^{R_0}u^{\kappa-1}\prod_i\big(\xi_i+uC_i(0)\big)^{-w_i}du-\frac1{e^{2\pi i\kappa^\flat}-1}\oint_{|u|=R_0}u^{\kappa-1}\prod_i\big(\xi_i+uC_i(0)\big)^{-w_i}du\Big],
\]
independent of $R_0\ge2\max_i|\xi_i|/c$ by the same derivative. When $\re\kappa^\flat<0$ the circle integral is $O(R_0^{\re\kappa^\flat})$ and the bracket tends to the convergent integral $\int_0^\infty$; the bracket is holomorphic in $(\kappa,\mathbf w)$ off $\kappa^\flat\in\mathbb Z$, so it is the continuation of that integral in $\kappa$ and $\mathbf w$. Where every $\re w_i>0$, $C_i(0)>0$ gives $\prod_i(\xi_i+uC_i(0))^{-w_i}=\prod_iC_i(0)^{-w_i}\prod_i(u+\xi_i/C_i(0))^{-w_i}$, and \hb{prop:scalebreta} with $J=\{1,\dots,p\}$ and the coefficients $\xi_i/C_i(0)$ gives the stacked Gamma times the Breta.

(ii) When $G$ and the $C_i$ are holomorphic on a neighbourhood of $[0,\infty)$, the cluster port of \hb{def:clusterport} with $r_i=1$ and the lower coefficients $\lambda\xi_i$, at the $b$ chosen above, is
\[
\frac1{2\pi i}\Big[\big(e^{2\pi i\kappa^\flat}-1\big)\int_b^\infty\cdots\,d\rho+\oint_{|\rho|=b}\cdots\,d\rho\Big]=\frac{e^{2\pi i\kappa^\flat}-1}{2\pi i}\Big(\int_b^\infty\cdots\,d\rho+A_b(\lambda)\Big)=\frac{e^{2\pi i\kappa^\flat}-1}{2\pi i}\,A(\lambda),
\]
and it does not depend on $b$ (\hb{thm:clusterport}(i)).
\end{proof}

\begin{sremark}
With $p=1$, $\xi_1=1$ and $C=C_1$, $\int_0^\infty u^{\kappa-1}(1+uC(0))^{-w}du=C(0)^{-\kappa}\Gst21{\kappa,\ -\kappa^\flat}{w}$, the branch part of the reference.
\end{sremark}

\begin{htheorem}[\hb{thm:rayspace} (The values along a ray)]
As stated, over a subfield $\Bbbk$ of $\mathbb C$ containing the weights and the origin.
\end{htheorem}

\begin{proof}
(i) is the reference's (the reference, proof of thm:rayspace). The letters have their coefficients in $\Bbbk$, so $X_F$, the occupations and the stride are defined over $\Bbbk$, and with the origin in $\Bbbk$ and the direction integral the operators $c^{(0)}_e+k\,c^{(1)}_e+\partial_{y_e}$ map $\Bbbk[X_F]\otimes\Bbbk(k)$ to itself. Let $\mathcal L_\Bbbk$ be the $\Bbbk(k)$-span of their images and $M_\Bbbk=\big(\Bbbk[X_F]\otimes\Bbbk(k)\big)/\mathcal L_\Bbbk$. Since $\mathbb C[X_F]=\Bbbk[X_F]\otimes_\Bbbk\mathbb C$, $M_\Bbbk\otimes_{\Bbbk(k)}\mathbb C(k)$ is the quotient of \hb{def:raygauge}, a cokernel commuting with the extension of scalars; so $\dim_{\Bbbk(k)}M_\Bbbk$ is the dimension $d$ of the ray, finite by (i).

(ii) The classes of $g,fg,\dots,f^dg$ in $M_\Bbbk$ are $d+1$ vectors of a space of dimension $d$, so there are $p_0,\dots,p_d\in\Bbbk(k)$, not all zero, with $\sum_jp_j(k)\,f^jg\in\mathcal L_\Bbbk$; clearing denominators puts them in $\Bbbk[k]$. With $h$ the largest $j$ with $p_j\neq0$, $h\le d$, and \hb{cor:symbol} gives $\sum_{j\le h}p_j(k)\langle g\rangle_{k+j}=0$ for all large $k$.

(iii) Let $k_1$ exceed the integer zeros of $p_h$ and the bound of (ii). For $k\ge k_1$, $\langle g\rangle_{k+h}=-\sum_{j<h}\frac{p_j(k)}{p_h(k)}\langle g\rangle_{k+j}$ with $p_j(k)/p_h(k)\in\Bbbk$, the values of polynomials over $\Bbbk$ at integers lying in $\Bbbk$; induction on $k$ writes every $\langle g\rangle_k$, $k\ge k_1$, as a $\Bbbk$-combination of $\langle g\rangle_{k_1},\dots,\langle g\rangle_{k_1+h-1}$.
\end{proof}

\begin{hlemma}[\hb{lem:polecount} (Index by pole count)]
As stated, over a subfield $\Bbbk$ of $\mathbb C$.
\end{hlemma}

\begin{proof}
The poles of $r$ are rational, so its partial fractions have coefficients in $\Bbbk$: $r(k)=\sum_{h\ge1}\sum_{s=1}^pc_{h,s}(k+h)^{-s}$, finitely many $h$, $c_{h,s}\in\Bbbk$, with no polynomial part since $r\to0$, and $\sum_hc_{h,1}=0$ since $r(k)=O(k^{-2})$. \hb{prop:terminalrow} at $t_0=0$, whose identities are linear in the coefficients and hold at coefficients in any field, gives
\[
\sum_{k\ge0}r(k)=\sum_{s\ge2}C_s\,\Tx{s}{\IZ}{0}-\sum_{h,s}c_{h,s}\,\Th{s}{\IZ}{0}{h-1},\qquad C_s=\sum_hc_{h,s}\in\Bbbk,
\]
the index-one terms regularized and their coefficient $C_1=0$, and every $\Th{s}{\IZ}{0}{h-1}$, a finite sum of $m^{-s}$, rational. If $r(-M-k)=-r(k)$, uniqueness of the partial fractions gives $(-1)^sc_{M-h,s}=-c_{h,s}$, so $C_s=(-1)^{s+1}C_s$ and $C_s=0$ for even $s$. For poles in one coset $-b+\mathbb Z$ the same computation runs in the lattice input $\Ilat{b+\mathbb Z}$, the sums $\sum_{k\ge0}(k+b+h)^{-s}$ being its Trailmoments at its cuts; with several cosets the index-one sums of the cosets are regularized one by one, and only their total coefficient vanishes, which leaves the differences of the index-one values of the cosets. The coefficients are in $\Bbbk$ throughout.
\end{proof}

% ---- Proofs of the new chapters ----
\section{Chapter~\hb{ch:twoends}: the two ends of a ray}\label{sec:chtwoends}

The proofs of this chapter use three facts about the symmetries of Proposition~\hb{prop:raysym}: in the coordinates $y=\log\alpha$ they are linear, $y\mapsto\pm\beta^\ast y$; they keep the stride; and an expansion at a native gauge depends only on the Taylor series of the insertion there, which insertions realize to every finite order, $\mathbb C[X_F]$ containing the Laurent polynomials of degree $0$.

\begin{hdefinition}[\hb{def:raymodule} (The module of a ray; masters; the one-step matrix)]
As stated.
\end{hdefinition}

\begin{proof}
$\mathcal L$ is spanned over $\mathbb C(k)$ by $\rho_e(g)=(c^{(0)}_e+k\,c^{(1)}_e+\partial_{y_e})g$. On $\mathbb C[X_F]\otimes\mathbb C(k)$ put $\operatorname{Sh}(g\otimes r)=fg\otimes\varsigma r$, $(\varsigma r)(k)=r(k+1)$; then $\operatorname{Sh}(r'x)=\varsigma(r')\operatorname{Sh}(x)$. By \hb{lem:insertions}(ii), $\partial_{y_e}(fg)=c^{(1)}_efg+f\,\partial_{y_e}g$, so
\[
f\,\big(c^{(0)}_e+(k+1)\,c^{(1)}_e+\partial_{y_e}\big)g=\big(c^{(0)}_e+k\,c^{(1)}_e+\partial_{y_e}\big)(fg),
\]
and $\operatorname{Sh}(r(k)\rho_e(g))=r(k+1)\,\rho_e(fg)\in\mathcal L$. The stride is a unit of $\mathbb C[X_F]$, its exponents being integers and the letters units there, so $\operatorname{Sh}$ has the inverse $x\mapsto f^{-1}\varsigma^{-1}x$, which maps $\mathcal L$ into itself by the same identity read backwards: $\operatorname{Sh}$ descends to a bijection of $\mathfrak M$ with $\operatorname{Sh}(rm)=\varsigma(r)\operatorname{Sh}(m)$. A semilinear bijection maps bases of subspaces to bases of their images, so the classes $\operatorname{Sh}^{-1}[b_i]=[f^{-1}b_i]$ of masters of a submodule form a basis of it, and $\mathbf B_{\mathbf b}$ is invertible.

Let the cyclic submodule of $[1]$ be finite-dimensional, and let $h'$ be the least index with $\operatorname{Sh}^{h'}[1]$ in the span $W$ of $[1],\dots,\operatorname{Sh}^{h'-1}[1]$. Then $\operatorname{Sh}(W)\subseteq W$, and $\operatorname{Sh}(W)$ has the dimension of $W$, so $\operatorname{Sh}(W)=W$ and $\operatorname{Sh}^{-1}(W)=W$: $W$ is the cyclic submodule of $[1]$, and $h'=h$. Writing $\operatorname{Sh}^h[1]$ in the basis and clearing denominators gives a relation with $p_h\neq0$. If $p_0=0$, applying $\operatorname{Sh}^{-1}$ to $\sum_{i=1}^hp_i(k)\operatorname{Sh}^i[1]=0$ gives $\sum_{i=1}^hp_i(k-1)\operatorname{Sh}^{i-1}[1]=0$, a relation among $[1],\dots,\operatorname{Sh}^{h-1}[1]$, which is impossible. Two relations of length $h$ are proportional over $\mathbb C(k)$, since otherwise a combination of them eliminating $\operatorname{Sh}^h[1]$ would be a nonzero relation among the basis; with coefficients in $\mathbb C[k]$ without a common factor the proportionality factor is a constant.
\end{proof}

\begin{hdefinition}[\hb{def:raysolution} (Solutions)]
As stated.
\end{hdefinition}

\begin{proof}
On germs $\varsigma$ is bijective and $\varsigma(rx)=r(k+1)\varsigma(x)$. On $\lambda^kk^\rho\,\mathbb C((1/k))$, a module of rank one over $\mathbb C((1/k))$ with the generator $\lambda^kk^\rho$, $\varsigma$ is the composite of the substitution $k\mapsto k+1$, the automorphism $\varsigma_0$ of the field $\mathbb C((1/k))$ with $\frac1k\mapsto\frac1k\big(1+\frac1k\big)^{-1}$, with the multiplication by the unit $\lambda(1+\frac1k)^\rho=\lambda\sum_i\binom\rho ik^{-i}$; so it is bijective and $\varsigma(rx)=r(k+1)\varsigma(x)$ for $r\in\mathbb C(k)\subset\mathbb C((1/k))$.
\end{proof}

\begin{hlemma}[\hb{lem:types} (Solutions of distinct types)]
As stated.
\end{hlemma}

\begin{proof}
With $\varsigma_0$ the substitution $k\mapsto k+1$ on $\mathbb C((1/k))$, $\hat A_a(\operatorname{Sh}m)=\lambda_a(1+\frac1k)^{\rho_a}\varsigma_0\big(\hat A_a(m)\big)$. Suppose $\sum_{a\in S}c_a\hat A_a=0$ on $\mathfrak M$ with $c_a\in\mathbb C((1/k))$, every $c_a\neq0$, and $S$ nonempty and minimal. Applying the relation at $\operatorname{Sh}m$, and applying $\varsigma_0$ to it at $m$, gives
\[
\sum_{a\in S}c_a\lambda_a\big(1+\tfrac1k\big)^{\rho_a}\,\varsigma_0\big(\hat A_a(m)\big)=0,\qquad \sum_{a\in S}\varsigma_0(c_a)\,\varsigma_0\big(\hat A_a(m)\big)=0
\]
for every $m$. The coefficient vectors of relations among the maps $\varsigma_0\circ\hat A_a$, $a\in S$, form a space over $\mathbb C((1/k))$ of dimension one: two independent vectors in it would combine to a nonzero one with a zero coordinate, and its image under $\varsigma_0^{-1}$ would be a relation among fewer $\hat A_a$. So $c_a\lambda_a(1+\frac1k)^{\rho_a}=\mu\,\varsigma_0(c_a)$ for one $\mu$ and every $a\in S$. With $c_a=c^0_ak^{v_a}\big(1+c'_ak^{-1}+O(k^{-2})\big)$, $c^0_a\neq0$, $v_a\in\mathbb Z$, the two sides are $c^0_ak^{v_a}\lambda_a\big(1+(\rho_a+c'_a)k^{-1}+\cdots\big)$ and $\mu\,c^0_ak^{v_a}\big(1+(v_a+c'_a)k^{-1}+\cdots\big)$, so $\mu=\lambda_a\big(1+(\rho_a-v_a)k^{-1}+O(k^{-2})\big)$ for every $a\in S$: one $\lambda$, and $\rho_a-v_a$ independent of $a$, so all $a\in S$ have one type. Fix $\rho$ of that type and write every $A_a$, $a\in S$, in $\lambda^kk^\rho\,\mathbb C((1/k))$, which changes the $c_a$ by integral powers of $k$; then $c_a\lambda(1+\frac1k)^\rho=\mu\,\varsigma_0(c_a)$ makes $c_a/\varsigma_0(c_a)$ independent of $a$, so every quotient $c_a/c_b$ is fixed by $\varsigma_0$. A series $r=\sum_{i\ge i_0}r_ik^{-i}$ with $\varsigma_0r=r$ is constant: if $r_i\neq0$ for some $i\neq0$, the least such $i$ gives in $\varsigma_0r-r$ the term $-i\,r_ik^{-i-1}$ and no other of that order. So the relation is $c(k)\sum_{a\in S}\gamma_a\hat A_a=0$ with constants $\gamma_a\neq0$, and $\sum_a\gamma_aA_a=0$ is a relation over $\mathbb C$ among solutions of one type: excluded. Extended $\mathbb C((1/k))$-linearly to $\mathfrak N\otimes\mathbb C((1/k))$, the normalizations are $r$ independent linear forms on a space of dimension $\dim_{\mathbb C(k)}\mathfrak N$; the argument runs on $\mathfrak N$ in place of $\mathfrak M$, $\operatorname{Sh}$ mapping $\mathfrak N$ onto itself.
\end{proof}

\begin{htheorem}[\hb{thm:determine} (What solutions determine)]
As stated.
\end{htheorem}

\begin{proof}
(i) $V[f^{-1}b_i]=V(\operatorname{Sh}^{-1}[b_i])=\varsigma^{-1}V[b_i]$, and $V[f^{-1}b_i]=\sum_l\mathbf B_{il}V[b_l]$ by $\mathbb C(k)$-linearity. For the values, $\langle f^{-1}b_i\rangle_k=\langle b_i\rangle_{k-1}$ (\hb{thm:raylaws}).

(ii) $\operatorname{Ann}(\mathcal S)$ is a $\mathbb C(k)$-subspace, and $V(\operatorname{Sh}^{\pm1}m)=\varsigma^{\pm1}V(m)$ vanishes with $V(m)$.

(iii) By (i), $\mathbf B'v_V-\varsigma^{-1}v_V=(\mathbf B'-\mathbf B)v_V$, whose $i$-th entry is $V(m_i)$, $m_i=\sum_l(\mathbf B'-\mathbf B)_{il}[b_l]\in\mathfrak N$. If a nonzero $m=\sum_lw_l[b_l]$ lies in $\operatorname{Ann}(\mathcal S)\cap\mathfrak N$, adding $w$ to the first row of $\mathbf B$ gives another such matrix.

(iv) For a germ the relation of (i) holds at every integer beyond some bound, and so do the equations of a pair $(cP_0,cQ_0)$; for a formal series they hold as identities. Conversely, let $(P,Q)$ satisfy all the equations. By (i), $Q\,\varsigma^{-1}v_V-Pv_V=(Q\mathbf B-P)v_V$, so every row of $Q\mathbf B-P$, read in $\mathfrak N$, lies in $\operatorname{Ann}(\mathcal S)\cap\mathfrak N=0$, and $P=Q\mathbf B=QP_0/Q_0$. If $Q=0$ then $P=0$. Otherwise $Q_0$ divides every entry of $QP_0$; with $d=Q_0/\gcd(Q_0,Q)$, $d$ divides every entry of $P_0$ and $Q_0$, whose greatest common divisor is $1$, $Q_0$ being the least common denominator, so $d=1$, $Q=cQ_0$ and $P=cP_0$, and the degree bounds give $\deg c\le D-D_0$. The equations are linear and homogeneous in the coefficients; for each $D$ the solution spaces of the first $N$ of them, in any enumeration, decrease with $N$ in a space of finite dimension, so they are constant from some $N$ on.
\end{proof}

\begin{hcorollary}[\hb{cor:onesolution} (One solution of an irreducible module)]
As stated.
\end{hcorollary}

\begin{proof}
$\operatorname{Ann}(\{V\})\cap\mathfrak N$ is a submodule of $\mathfrak N$ (Theorem~\hb{thm:determine}(ii)) other than $\mathfrak N$, hence $0$. For the cyclic submodule of $[1]$, $\sum_ip_i(k)\varsigma^iV[1]=V\big(\sum_ip_i(k)\operatorname{Sh}^i[1]\big)=0$. A relation $\sum_{i<h}q_i(k)\varsigma^iV[1]=0$ is $V(m)=0$ with $m=\sum_{i<h}q_i(k)\operatorname{Sh}^i[1]$, nonzero as the $\operatorname{Sh}^i[1]$, $i<h$, are a basis; then $\operatorname{Ann}(\{V\})$ meets the cyclic submodule in a nonzero submodule, all of it, and $V[1]=0$.
\end{proof}

\begin{hproposition}[\hb{prop:rays-lawrows} (The law of {$[1]$} read from a one-step matrix)]
As stated.
\end{hproposition}

\begin{proof}
(i) $\operatorname{Sh}^{-1}$ is a bijection of $\mathfrak M$ with $\operatorname{Sh}^{-1}(r(k)\,m)=r(k-1)\operatorname{Sh}^{-1}(m)$ (Definition~\hb{def:raymodule}). For $i=0$, $[1]=[b_1]$. If $\operatorname{Sh}^{-(i-1)}[1]=\sum_lw_{i-1,l}(k)[b_l]$, then
\[
\operatorname{Sh}^{-i}[1]=\sum_lw_{i-1,l}(k-1)\,[f^{-1}b_l]=\sum_{l,l'}w_{i-1,l}(k-1)\,\mathbf B_{ll'}(k)\,[b_{l'}],
\]
and induction gives the product form. A solution has $V\circ\operatorname{Sh}^{-i}=\varsigma^{-i}\circ V$ and is $\mathbb C(k)$-linear. For the values, $f^{-i}-\sum_lw_{i,l}(k)\,b_l$ lies in $\mathcal L$, so its value vanishes identically, and $\langle f^{-i}\rangle_k=\langle1\rangle_{k-i}$ (Proposition~\hb{prop:rays-strength}(ii)).

(ii) The $w_i$ lie in $\mathbb C(k)^{\mathfrak m}$, so $h\le\mathfrak m$. By (i), $\sum_iq_i(k)\operatorname{Sh}^{-i}[1]=0$, and $\operatorname{Sh}^h$, with $\operatorname{Sh}^h(r(k)\,m)=r(k+h)\operatorname{Sh}^h(m)$, turns it into $\sum_iq_i(k+h)\operatorname{Sh}^{h-i}[1]=0$, which is the stated relation, $\operatorname{Sh}^j[1]=[f^j]$; the shift of the variable keeps the coefficients without a common factor. $p_0(k)=q_h(k+h)\neq0$, since $w_0,\dots,w_{h-1}$ are independent. If $\mathbf B$ is invertible and $q_0=0$, then $\sum_{i\ge1}q_i(k)\,w_{i-1}(k-1)\,\mathbf B(k)=0$ gives $\sum_{i\ge1}q_i(k)\,w_{i-1}(k-1)=0$, and at $k+1$ a dependence among $w_0,\dots,w_{h-1}$ with coefficients not all zero: so $q_0\neq0$ and $p_h\neq0$. Let $h'$ be the dimension of the cyclic submodule of $[1]$; its basis $[1],\operatorname{Sh}[1],\dots,\operatorname{Sh}^{h'-1}[1]$ (Definition~\hb{def:raymodule}) is carried by the semilinear bijection $\operatorname{Sh}^{-(h'-1)}$ to the basis $\operatorname{Sh}^{-i}[1]$, $0\le i<h'$. A dependence among $w_0,\dots,w_h$ gives by (i) a dependence among $[1],\dots,\operatorname{Sh}^{-h}[1]$ with the same coefficients, so $h\ge h'$. When the $[b_l]$ are a basis, a dependence among the $\operatorname{Sh}^{-i}[1]$ is one among their coordinate rows $w_i$, so $h=h'$, and the relation, of length $h'$ with polynomial coefficients without a common factor, is the law of $[1]$ up to a constant factor (Definition~\hb{def:raymodule}).

(iii) $M=[w_0;\dots;w_h]_J$ is an $(h+1)\times h$ matrix of rank $h$, so its left kernel is a line. It contains $(q_i)_i$, and also the vector of the signed maximal minors $\mu_i=(-1)^i\det M_{\hat i}$, $M_{\hat i}$ being $M$ without its row $i$: for every column $l$, $\sum_i(-1)^iM_{il}\det M_{\hat i}$ is the expansion along its first column of the square matrix $[M_{\cdot l}\,|\,M]$, which has two equal columns; and $\mu_h=(-1)^h\det\mathbf W_J\neq0$. For the divisibility, put $D_i(k)=\prod_{j=0}^{i-1}Q(k-j)$, so that $D_iw_i$ is a row of polynomials, $w_i$ being $w_0$ times $\mathbf B(k-i+1)\cdots\mathbf B(k)$. The polynomial matrix $\tilde M=\operatorname{diag}(D_0,\dots,D_h)\,M$ has the left kernel spanned by $(q_i/D_i)_i$, and also by its signed maximal minors $\tilde\mu_i$, polynomials with $\tilde\mu_h=(-1)^hN$, $N=\det\mathbf W_J\prod_{i<h}D_i$. Let $\tilde q$ be a polynomial vector without a common factor spanning that kernel; then $\tilde\mu=\Lambda\tilde q$ with $\Lambda$ a polynomial, since a denominator of $\Lambda$ would divide every $\tilde q_i$. With $D=D_h$, the polynomial vector $(q_iD/D_i)_i$ spans the kernel, so $\tilde q=(q_iD/D_i)_i/g$ with $g$ the greatest common divisor of its entries, and $\tilde q_h=q_h/g$ divides $N$. Since $g$ divides every $q_iD$ and the $q_i$ have no common factor, $g$ divides $D$. So $q_h=g\,\tilde q_h$ divides $DN=\det\mathbf W_J\prod_{i=0}^hD_i=\det\mathbf W_J\prod_{j=0}^{h-1}Q(k-j)^{h-j}$. At $k_0$ with $Q(k_0-j)\neq0$ for $j<h$, every $D_i(k_0)\neq0$ and the rows $w_i$, $i<h$, are regular, so the order of $q_h$ at $k_0$ is at most that of $\det\mathbf W_J$, for every $J$ with $\det\mathbf W_J\neq0$: if $q_h(k_0)=0$, every maximal minor of $\mathbf W(k_0)$ vanishes.

(iv) Take $J$ all the columns. Since $w_i(k)=w_{i-1}(k-1)\mathbf B(k)$, $[w_1;\dots;w_h](k)=\mathbf W(k-1)\,\mathbf B(k)$, so $\mu_0=\det\mathbf W(k-1)\det\mathbf B(k)$ and $\mu_h=(-1)^h\det\mathbf W(k)$, and $q_0\mu_h=q_h\mu_0$ by (iii).

(v) $\operatorname{Sh}^{-1}\big(\sum_lc_l(k)[b_l]\big)=\sum_{l,l'}c_l(k-1)\,\mathbf B_{ll'}(k)\,[b_{l'}]$, so the rows obey the same recursion, and the proofs of (i) to (iv) run with $m$ in place of $[1]$.
\end{proof}

\begin{hproposition}[\hb{prop:raysym} (The symmetries of a ray)]
As stated.
\end{hproposition}

\begin{proof}
$\phi_\beta\circ\phi_{\beta'}=\phi_{\beta'\circ\beta}$, the Cremona commutes with every $\phi_\beta$ and its square is the identity, so the maps $\phi_\beta$ and $\bar\phi_\beta$ form a group, and the two conditions defining a symmetry pass to composites and inverses.

(i) If $\beta$ carries the tables, $K_j(\alpha\circ\beta)=\sum_{T\in\cB_j}\theta_{j,T}\alpha^{\beta(T)}=\sum_{T'\in\cB_{h(j)}}\theta_{h(j),T'}\alpha^{T'}=K_{h(j)}(\alpha)$, and the same with the complements gives $\Psi$. If $\beta$ exchanges them, the sets $\beta(T)$, $T\in\cB_j$, are the complements of the bases $\cE\setminus\beta(T)$ of table $h(j)$, with the same weights, so $K_j(\alpha\circ\beta)=\Psi_{h(j)}(\alpha)$ and $\Psi_j(\alpha\circ\beta)=K_{h(j)}(\alpha)$. By \hb{prop:identity}, $K_j(\alpha^{-1})=\alpha^{-\cE}\Psi_j(\alpha)$ and $\Psi_j(\alpha^{-1})=\alpha^{-\cE}K_j(\alpha)$; with $(\alpha\circ\beta)^{\cE}=\alpha^{\cE}$, $K_j\big((\alpha\circ\beta)^{-1}\big)=\alpha^{-\cE}\Psi_j(\alpha\circ\beta)=\alpha^{-\cE}K_{h(j)}(\alpha)$, and the same for $\Psi_j$.

(ii) With $(\alpha\circ\beta)^u=\alpha^{u\circ\beta^{-1}}$ and $(\alpha^{-1})^u=\alpha^{-u}$, (i) gives in the four cases
\begin{gather*}
f\circ\phi_\beta=\alpha^{\mvec_1\circ\beta^{-1}}\prod_jK_{h(j)}^{-s_{1j}}\Psi_{h(j)}^{-t_{1j}},\qquad f\circ\phi_\beta=\alpha^{\mvec_1\circ\beta^{-1}}\prod_j\Psi_{h(j)}^{-s_{1j}}K_{h(j)}^{-t_{1j}},\\
f\circ\operatorname{cr}=\alpha^{-\mvec_1+2\sigma_1\one}\prod_j\Psi_j^{-s_{1j}}K_j^{-t_{1j}},\qquad f\circ\bar\phi_\beta=\alpha^{-\mvec_1\circ\beta^{-1}+2\sigma_1\one}\prod_jK_{h(j)}^{-s_{1j}}\Psi_{h(j)}^{-t_{1j}},
\end{gather*}
each equal to $f$ under the conditions with $i=1$, after reindexing the tables by $h$; the same computation with the data of the origin gives $(\omega_0/dy)\circ\phi=\omega_0/dy$ under the conditions with $i=0$, on $X_F$ for an integral origin and on the positive orthant with the principal powers otherwise. In the coordinates $y$ the maps are linear, $y\mapsto\pm\beta^\ast y$, and they keep the form $dy$ on $\mathbb C^{\cE}/\mathbb C\one$ up to the sign of the permutation and the factor $(-1)^{n-1}$ of the inversion; so $\phi^\ast\omega_0=\pm\omega_0$ for an integral origin. The letters go to letters times monomials, so $X_F$ is mapped onto itself, and positive gauges go to positive gauges. The value $\langle g'\circ\phi\rangle_k$ is the integral of $(g'\circ\phi)f^k(\omega_0/dy)$ against the Lebesgue measure of the quotient, which does not depend on the chart (\hb{def:question}); the linear change of variables $y\mapsto\pm\beta^\ast y$ keeps that measure, and with $f\circ\phi=f$ and $(\omega_0/dy)\circ\phi=\omega_0/dy$ it gives $\langle g'\rangle_k$.

(iii) Substitute (i) in $\alpha^u\prod_jK_j^{-a_j}\Psi_j^{-b_j}$ and reindex the tables by $h$.
\end{proof}

\begin{htheorem}[\hb{thm:invariant} (The invariant module)]
As stated.
\end{htheorem}

\begin{proof}
(i) In the coordinates $y$, $\phi$ is $y\mapsto M_\phi y$ with $M_\phi$ linear. For a constant vector field $\xi$ and a function $u$, $(\xi u)\circ\phi=(M_\phi^{-1}\xi)(u\circ\phi)$. The relation space is spanned by $\rho_\xi(g')=\xi\big(\log(f^k\omega_0/dy)\big)g'+\xi g'=\sum_e\xi_e\big(c^{(0)}_e+k\,c^{(1)}_e+\partial_{y_e}\big)g'$, and since $\log(f^k\omega_0/dy)\circ\phi$ differs from $\log(f^k\omega_0/dy)$ by a constant (Proposition~\hb{prop:raysym}), $\rho_\xi(g')\circ\phi=\rho_{M_\phi^{-1}\xi}(g'\circ\phi)\in\mathcal L$. So $\phi^\ast$ is well defined on $\mathfrak M$, and with $(fg')\circ\phi=f\,(g'\circ\phi)$ it commutes with $\operatorname{Sh}$. Then $e_G$ is an idempotent commuting with $\operatorname{Sh}$, its image and kernel are submodules with $\mathfrak M=e_G\mathfrak M\oplus(1-e_G)\mathfrak M$, $e_G[g']=[e_Gg']$ with $e_Gg'$ invariant, and the class of an invariant insertion is fixed by $e_G$.

(ii) $V=V\circ e_G$ vanishes on the kernel of $e_G$. The values are invariant by Proposition~\hb{prop:raysym}(ii). $1$ is invariant, so $[1]\in\mathfrak M^G$, a submodule, which then contains the cyclic submodule of $[1]$.

(iii) $\phi$ fixes $\alpha_c$ and keeps $f$, so it maps a steepest-descent disc $D_c$ through $\alpha_c$ to a disc through $\alpha_c$ on which $|f|<|f(\alpha_c)|$ off $\alpha_c$, with $|f|\le|f(\alpha_c)|e^{-\varepsilon}$ on its boundary. In Morse coordinates $v$ at $\alpha_c$, $\log f=\log f(\alpha_c)-\frac12\sum_iv_i^2$, the homology of a small ball relative to the part where $|f|\le|f(\alpha_c)|e^{-\varepsilon}$ is $\mathbb Z$ in the dimension $n-1$, generated by the real ball; so $\phi(D_c)$ is $\pm D_c$ up to a boundary and pieces where $|f|\le|f(\alpha_c)|e^{-\varepsilon}$. Fix a branch of $\omega_0$ near $\alpha_c$; $(\omega_0/dy)\circ\phi$ and $\omega_0/dy$ are branches of one function there and $\phi^\ast dy=\pm dy$, so $\phi^\ast\omega_0=c\,\omega_0$ near $\alpha_c$ for a constant $c\neq0$. Then $\int_{D_c}(g\circ\phi)f^k\omega_0=c^{-1}\int_{\phi(D_c)}gf^k\omega_0$, and since $gf^k\omega_0$ is a closed form of top degree, $A_c\circ\phi^\ast=c'A_c$ for a constant $c'$. The leading coefficients of $A_c(g\circ\phi)$ and of $A_c(g)$ are $(2\pi)^{(n-1)/2}(\omega_0/dv)(\alpha_c)$ times $(g\circ\phi)(\alpha_c)=g(\alpha_c)$ (\hb{thm:gaugesolution}(i)), and an insertion with $g(\alpha_c)\neq0$ gives $c'=1$.

For $\phi\in G$ and $\psi\in G_c$, $A_c\circ(\phi\psi)^\ast=A_c\circ\psi^\ast\circ\phi^\ast=A_c\circ\phi^\ast$, so $A_O$ is well defined; for $\chi\in G$, $A_O\circ\chi^\ast=\sum A_c\circ(\chi\phi)^\ast$ runs over the same cosets, so $A_O$ is invariant, and $A_O\circ e_G=A_O$. By the argument above, $A_c\circ\phi^\ast$ is a constant multiple of the expansion over the disc $\phi(D_c)$, a steepest-descent disc through the native gauge $\phi(\alpha_c)$. The nondegenerate native gauges are isolated points of the critical locus, a variety, so finitely many; take an insertion $g$ with $g(\alpha_c)=1$ that vanishes to order $N$ at every other one. Its expansion there vanishes to the relative order $k^{-\lceil N/2\rceil}$, since $a_i(g)$ involves derivatives of $g$ of order at most $2i$, so $A_O(e_Gg)=A_O(g)$ has the nonzero leading coefficient of $A_c(g)$, and $A_O$ does not vanish on $\mathfrak M^G$. The same insertions, applied to a relation $\sum_O\gamma_OA_O=0$ on $\mathfrak M^G$ among orbits of one native value, give $\gamma_O=0$ orbit by orbit. Lemma~\hb{lem:types} gives the bound.
\end{proof}

\begin{htheorem}[\hb{thm:reversed} (A reversed exchange at a lean)]
As stated.
\end{htheorem}

\begin{proof}
(i) $T\mapsto\cE\setminus\beta(T)$ is a bijection of $\cB_j$ onto $\cB_{h(j)}$ keeping the weights, so $K_j(\one)=K_{h(j)}(\one)$ and, under the weight laws, $R_{h(j),\beta(e)}(\one)=\mathbb P_{j,\one}\big(\beta(e)\notin\beta(T)\big)=\mathbb P_{j,\one}(e\notin T)=1-R_{j,e}(\one)$. By \hb{prop:massesgauge}, $\tau_e=\sum_j\big(\sigma_j+\delta_j(2R_{j,e}(\one)-1)\big)$; reindexing by $h$, with $(\sigma_{h(j)},\delta_{h(j)})=(\sigma_j,\delta_j)$ (\hb{def:carry}), $\tau_{\beta(e)}=\sum_j\big(\sigma_j+\delta_j(1-2R_{j,e}(\one))\big)=2\sigma-\tau_e$. The exponents $(s_j,t_j)$ are carried by $h$, so Proposition~\hb{prop:raysym}(ii) applies to $\bar\phi_\beta$ with the masses $\btau$, and to the ray with origin $0$ and direction $(c\btau;c\mathbf s,c\mathbf t)$. For $\phi_\beta$ and $\operatorname{cr}$, Proposition~\hb{prop:raysym}(i) turns the integrand into $\alpha^{\btau\circ\beta^{-1}}\prod_j\Psi_{h(j)}^{-s_j}K_{h(j)}^{-t_j}$ and $\alpha^{-\btau+2\sigma\one}\prod_j\Psi_j^{-s_j}K_j^{-t_j}$. Since $\tau_{\beta^{-1}(e)}=2\sigma-\tau_e=\sum_j\big(\sigma_j-\delta_j(2R_{j,e}(\one)-1)\big)$, the masses of $\Upsilon_{\boldsymbol\sigma,-\boldsymbol\delta;\one}$, whose exponents are $(t_j,s_j)$, both are its integrand.

(ii) This is the exchange clause of \hb{prop:certify} at the table's gauge, read on the views. In the coordinates $y$, Proposition~\hb{prop:raysym}(i) gives $(K_j/\Psi_j)\circ\bar\phi_\beta=K_{h(j)}/\Psi_{h(j)}$ and $(K_j\Psi_j)\circ\bar\phi_\beta=\alpha^{-2\cE}K_{h(j)}\Psi_{h(j)}$, so, with $(\sigma_{h(j)},\delta_{h(j)})=(\sigma_j,\delta_j)$,
\[
\Phi_{\mathbf O}(-\beta^\ast y)=\Phi_{\mathbf O}(y),\qquad \Phi_{\mathbf E}(-\beta^\ast y)=\Phi_{\mathbf E}(y)-2\sigma\ip\one y .
\]
At the table's gauge $\Psi_j(e^y)=e^{\ip\one y}K_j(e^{-y})$ (\hb{prop:identity}), so $\Phi_{\mathbf O}(-y)=-\Phi_{\mathbf O}(y)$ and $\Phi_{\mathbf E}(-y)=\Phi_{\mathbf E}(y)-2\sigma\ip\one y$. For $w=v_0\one_e+\sum_iv_i\one_{g_i}$ and $w'=v_0\one_{\beta(e)}+\sum_iv_i\one_{\beta(g_i)}$, $\beta^\ast w'=w$ and $\ip\one{w'}=\ip\one w$, so
\[
\Phi_{\mathbf O}(w')=\Phi_{\mathbf O}(-w)=-\Phi_{\mathbf O}(w),\qquad \Phi_{\mathbf E}(w')=\Phi_{\mathbf E}(-w)+2\sigma\ip\one w=\Phi_{\mathbf E}(w),
\]
the identities of the views at the table's gauge. $\Phi_{\mathbf O}$ is linear in $\boldsymbol\delta$, so $-\Phi^{\boldsymbol\delta}_{\mathbf O}=\Phi^{-\boldsymbol\delta}_{\mathbf O}$, and the views, hence the readings, of $\beta(e)$ at $\boldsymbol\delta$ are those of $e$ at $-\boldsymbol\delta$, the partners matched by $\beta$. At the centre $\Phi_{\mathbf O}=0$ and $e$ and $\beta(e)$ read alike. At a lean the partition at $-\boldsymbol\delta$ is that at $\boldsymbol\delta$ (\hb{prop:leanray}): if $e$ and $e'$ lie in one class at $\boldsymbol\delta$, they lie in one class at $-\boldsymbol\delta$, so $\beta(e)$ and $\beta(e')$ read alike at $\boldsymbol\delta$. So $\beta$ maps classes into classes, and so does $\beta^{-1}$, which exchanges the family as well: $\beta$ permutes the classes.

(iii) By \hb{cor:exchangecumulants}, $\beta$ maps every class of the centre onto itself. Let $C$ split at the lean into $C_K$ and $C_\Psi$, with $Q$ the polynomial written as $K$ from every mark of $C_K$ and as $\Psi$ from every mark of $C_\Psi$ (\hb{thm:acrosscentre}). With one table, $h$ is the identity and Proposition~\hb{prop:raysym}(i) gives $K(\alpha\circ\beta)=\Psi(\alpha)$ and $\Psi(\alpha\circ\beta)=K(\alpha)$. For $e\in C_K$ and an order $\mathbf g$ of its room with $K_{e,\mathbf g}=Q$, put the position variables $x_0$ at $\beta(e)$ and $x_i$ at $\beta(g_i)$; then $\alpha\circ\beta$ carries $x_0$ at $e$ and $x_i$ at $g_i$, so $\Psi_{\beta(e),\beta(\mathbf g)}=K_{e,\mathbf g}=Q$ and $\beta(e)\in C_\Psi$. For $e\in C_\Psi$ with $\Psi_{e,\mathbf g}=Q$, $K(\alpha)=\Psi(\alpha\circ\beta)$ gives $K_{\beta(e),\beta(\mathbf g)}=\Psi_{e,\mathbf g}=Q$ and $\beta(e)\in C_K$. $C_K$ and $C_\Psi$ have equally many marks, so $\beta(C_K)=C_\Psi$ and $\beta(C_\Psi)=C_K$.
\end{proof}

\begin{hdefinition}[\hb{def:raybox} (Levels, wall degree, boxes)]
As stated.
\end{hdefinition}

\begin{proof}
The products of the monomials $\alpha^v$, $v_{\cE}=0$, and of the chances $\theta_{j,T}\alpha^T/K_j$ and $\theta_{j,T}\alpha^{\cE\setminus T}/\Psi_j$ that generate $\mathbb C[X_F]$ (Definition~\hb{def:ray}) are constant multiples of evaluations of symbols, which therefore span it. The wall degree is the largest of finitely many numbers $-\kappa_U(g)/\kappa_{1,U}$, and $\kappa_U(g)+x\,\kappa_{1,U}\ge0$ for every $U$ exactly when $x$ is at least each of them. By \hb{lem:insertions}(i) and Definition~\hb{def:ray}, with $R_{j,e}=\sum_{T\ni e}\theta_{j,T}\alpha^T/K_j$ and $R^\ast_{j,e}=\sum_{T\not\ni e}\theta_{j,T}\alpha^{\cE\setminus T}/\Psi_j$, the operator $c^{(0)}_e+k\,c^{(1)}_e+\partial_{y_e}$ multiplies $\alpha^u\prod_jK_j^{-a_j}\Psi_j^{-b_j}$ by
\[
m_{0,e}+k\,m_{1,e}+u_e-\sum_j\big((s_{0j}+k\,s_{1j}+a_j)R_{j,e}+(t_{0j}+k\,t_{1j}+b_j)R^\ast_{j,e}\big),
\]
and $R_{j,e}$ times the evaluation of $(u;\mathbf a,\mathbf b)$ is the evaluation of $\sum_{T\ni e}\theta_{j,T}(u+\one_T;\mathbf a+\mathbf e_j,\mathbf b)$, and likewise for $R^\ast_{j,e}$; this is the evaluation of the mark relation. The letter identities evaluate to the evaluation of $g$ times $1-K_j/K_j$, or $1-\Psi_j/\Psi_j$. The maps of the four kinds carry each letter to a letter times $1$ or $\alpha^{-\cE}$ (Proposition~\hb{prop:raysym}(i)), so do their composites, and the action of Proposition~\hb{prop:raysym}(iii) is the substitution into the evaluation.
\end{proof}

\begin{hproposition}[\hb{prop:box-walldegree} (The wall degree)]
As stated.
\end{hproposition}

\begin{proof}
(i) By Lemma~\ref{lem:rays-sector}, on the sector of an order $\pi$, $|g|\,f^x=\prod_lq_l^{\kappa_{U_l}(g)+x\kappa_{1,U_l}}\,|g_\pi|\varphi_\pi^x$ with the last factor between positive constants, so $|g|f^x$ is bounded on the sector exactly when every exponent is nonnegative; every nonempty proper $U$ is a tail of some order, which gives the first statement. Likewise $|g\,f^k\omega_0/dy|$ is, on the sector of $\pi$, $\prod_lq_l^{c_l}$ times a factor between positive constants, $c_l=\re\kappa_{0,U_l}+\kappa_{U_l}(g)+\re(k)\kappa_{1,U_l}$, and $\int_{(0,1]^{n-1}}\prod_lq_l^{c_l-1}\,dq$ is finite exactly when every $c_l>0$; so the value converges absolutely exactly when $\re\kappa_{0,U}+\kappa_U(g)+\re(k)\kappa_{1,U}>0$ for every $U$; for the origin $0$ this reads $\re k>-\kappa_U(g)/\kappa_{1,U}$ for every $U$, and in general $\kappa_U(g)\ge-\deg_\kappa g\,\kappa_{1,U}$ gives the sufficient condition.

(ii) $\kappa_U(\alpha^u)=u_U$. For $u\neq0$ with $u_{\cE}=0$ some $u_e<0$, and $U=\{e\}$ is proper, $n\ge2$, so the wall degree is positive; for $u=0$ it is $0$. A monomial of $f^{-1}=\alpha^{-\mvec_1}\prod_jK_j^{s_{1j}}\Psi_j^{t_{1j}}$ is $\alpha^w$ with $w=-\mvec_1+\sum_j\big(\sum_{i\le s_{1j}}\one_{T_{j,i}}+\sum_{i\le t_{1j}}\one_{\cE\setminus T'_{j,i}}\big)$ for bases $T_{j,i},T'_{j,i}$ of table $j$, and by \hb{thm:window} $w_U\ge-m_{1,U}+\sum_j\big(s_{1j}d_j(U)+t_{1j}d^\ast_j(U)\big)=-\kappa_{1,U}$. The walls of Laurent symbols add, so $\kappa_U(\alpha^{u+w})+(\deg_\kappa\alpha^u+1)\kappa_{1,U}\ge0$.

(iii) The terms are $g$ and the symbols $(u+\one_T;\mathbf a+\mathbf e_j,\mathbf b)$ and $(u+\one_{\cE\setminus T};\mathbf a,\mathbf b+\mathbf e_j)$, whose walls are $\kappa_U(g)+|T\cap U|-d_j(U)$ and $\kappa_U(g)+|U\setminus T|-d^\ast_j(U)$, at least $\kappa_U(g)$ (\hb{thm:window}).

(iv) For a mark $e$, $\{e\}$ is proper, and $\kappa_{\{e\}}(g)\ge-x\,\kappa_{1,\{e\}}$ bounds $u_e$ from below at a fixed level; with $u_{\cE}$ fixed by the level, $u$ is bounded from above as well.

(v) A symmetry maps the positive orthant onto itself and keeps $f$ (Proposition~\hb{prop:raysym}(ii)), and the evaluation of $\phi^\ast g$ is that of $g$ composed with $\phi$, so $|\phi^\ast g|f^x$ and $|g|f^x$ have the same bound, and (i) gives the same wall degree. The action of Proposition~\hb{prop:raysym}(iii) maps the level $(\mathbf 0,\mathbf 0)$ to itself.
\end{proof}

\begin{hlemma}[\hb{lem:box-laurent} (The Laurent part)]
As stated.
\end{hlemma}

\begin{proof}
$\mathfrak M$ is finite-dimensional (Theorem~\hb{thm:rayspace}(i)). The direction has $s_{1j},t_{1j}\in\mathbb N$, so $f^{-1}$ is a Laurent polynomial and $\operatorname{Sh}^{-1}[\alpha^v]=[f^{-1}\alpha^v]$ lies in $\mathfrak M_{\mathrm L}$: $\operatorname{Sh}^{-1}$ maps $\mathfrak M_{\mathrm L}$ into itself, injectively and semilinearly, hence onto itself, and $\mathfrak M_{\mathrm L}$ is a submodule. It contains $[1]$, and $[g]=\operatorname{Sh}^l[f^{-l}g]$. When $s_{1j},t_{1j}\ge1$ for every $j$, $f^{-l}\alpha^u\prod_jK_j^{-a_j}\Psi_j^{-b_j}$ is a Laurent polynomial as soon as $l$ is at least every $a_j$ and every $b_j$. The symmetries map Laurent insertions to Laurent insertions (Proposition~\hb{prop:box-walldegree}(v)), so $e_G\mathfrak M_{\mathrm L}\subseteq\mathfrak M_{\mathrm L}\cap\mathfrak M^G$, and an element $x$ of the intersection is $e_Gx$; the intersection of two submodules is one (Theorem~\hb{thm:invariant}(i)), and $e_G$ of a Laurent insertion is invariant.
\end{proof}

\begin{hproposition}[\hb{prop:box-relations} (The relations of a box)]
As stated.
\end{hproposition}

\begin{proof}
(i) The terms of a letter identity of a symbol $(u;\mathbf a,\mathbf b)$ of the box for $K_j$, $a_j<\bar a_j$, are at levels at most the bounds and of wall degree at most that of the symbol (Proposition~\hb{prop:box-walldegree}(iii)); likewise for $\Psi_j$. A term of the mark relation of a stem at the level $(\bar{\mathbf a}^-+\mathbf e_j,\bar{\mathbf b}^-)$ has the coefficient $s_{0j}+k\,s_{1j}+\bar a^-_j$; if $\bar a_j\ge1$ its level is at most $(\bar{\mathbf a},\bar{\mathbf b})$, and if $\bar a_j=0$ then $K_j$ is not active and $\bar a^-_j=0$, so the coefficient vanishes identically and the term is omitted; likewise for $\Psi_j$, and the wall degrees do not grow. The evaluations are those of Definition~\hb{def:raybox}(iii), and Proposition~\hb{prop:rays-strength}(ii) gives the identity in $k$.

(ii) A letter identity of a symbol $x=(u;\mathbf a,\mathbf b)$ of the box for $K_j$, $a_j<\bar a_j$, makes $x$ congruent to a combination of symbols of the box at the level $(\mathbf a+\mathbf e_j,\mathbf b)$, by (i). Repeating, every formal combination $y$ is congruent modulo the letter identities to a combination $y'=\sum_uc_u(u;\bar{\mathbf a},\bar{\mathbf b})$ of symbols at the top level, with distinct $u$. The letter identities evaluate to $0$, so if $y$ evaluates to $0$, so does $y'$, and $\prod_jK_j^{\bar a_j}\Psi_j^{\bar b_j}$ times its evaluation is the Laurent polynomial $\sum_uc_u\alpha^u$, which is $0$ only when every $c_u$ is: $y'=0$, and $y$ is a combination of letter identities.

(iii) Let $x=(u;\mathbf a,\mathbf b)$ be such a symbol and $\prod_jK_j^{\bar a^-_j-a_j}\Psi_j^{\bar b^-_j-b_j}=\sum_v\lambda_v\alpha^v$, $\lambda_v\in\Bbbk$. Then $x$ and $\sum_v\lambda_v\,s_v$, $s_v=(u+v;\bar{\mathbf a}^-,\bar{\mathbf b}^-)$, have one evaluation. Each $v$ is a sum of $\bar a^-_j-a_j$ indicators of bases and $\bar b^-_j-b_j$ of complements of bases of every table $j$, so $\kappa_U(s_v)\ge\kappa_U(x)$ (\hb{thm:window}) and the $s_v$ are stems. The formal combination $\mathrm{MR}_e(x)-\sum_v\lambda_v\,\mathrm{MR}_e(s_v)$ evaluates to $(c^{(0)}_e+k\,c^{(1)}_e+\partial_{y_e})$ of $0$, and its terms lie in the box: those of $\mathrm{MR}_e(s_v)$ by (i), and those of $\mathrm{MR}_e(x)$ because their wall degree is at most $\deg_\kappa x\le\mathfrak t$ and a term at the level $(\mathbf a+\mathbf e_j,\mathbf b)$, with $a_j\le\bar a^-_j$, either has a level within the bounds, when $\bar a_j\ge1$, or the coefficient $s_{0j}+k\,s_{1j}$, identically $0$, when $\bar a_j=0$; likewise for $\Psi_j$. By (ii) it is a combination of letter identities, and $\mathrm{MR}_e(x)\in\mathcal R_{\mathfrak B}$.

(iv) $\phi$ permutes the letters keeping the bounds, so it maps the levels within the bounds onto themselves and the level of the stems onto itself, and it keeps the wall degree (Proposition~\hb{prop:box-walldegree}(v)); so it permutes the symbols of the box and the stems. The image of a letter identity of the box evaluates to $0$ and has its terms in the box, so it lies in the span of the letter identities by (ii). The image of the mark relation of a stem $s$ at $e$ evaluates to $\big((c^{(0)}_e+k\,c^{(1)}_e+\partial_{y_e})s\big)\circ\phi$, which is $\sum_{e'}\mu_{e'}(c^{(0)}_{e'}+k\,c^{(1)}_{e'}+\partial_{y_{e'}})(s\circ\phi)$ for constants $\mu_{e'}$, $\phi$ being linear in the coordinates $y$ and keeping $\log(f^k\omega_0/dy)$ up to a constant (proof of Theorem~\hb{thm:invariant}(i)); so the image less $\sum_{e'}\mu_{e'}\mathrm{MR}_{e'}(\phi^\ast s)$, with $\phi^\ast s$ a stem, evaluates to $0$ and lies in the span of the letter identities by (ii). So $\phi$ maps $\mathcal R_{\mathfrak B}$ into itself, and onto itself, $G$ being finite. Then $e_G\mathcal R_{\mathfrak B}\subseteq\mathcal R_{\mathfrak B}\cap\mathcal F^G_{\mathfrak B}$, and an element $x$ of the intersection is $e_Gx$.
\end{proof}

\begin{hproposition}[\hb{prop:box-closing} (The closing relations have more solutions than the ray)]
As stated.
\end{hproposition}

\begin{proof}
With $\mathbf s_0=\mathbf t_0=\mathbf 0$, $\omega_0/dy=\alpha^{\mvec_0}$ and $c^{(0)}_e=m_{0,e}$, and $c^{(1)}_e=\partial_{y_e}\log f=-\partial_{y_e}f^{-1}/f^{-1}$ (\hb{lem:insertions}(ii)), so
\[
\big(c^{(0)}_e+k\,c^{(1)}_e+\partial_{y_e}\big)\big(\alpha^wf^{-1}\big)=\alpha^w\big((m_{0,e}+w_e)\,f^{-1}+(1-k)\,\partial_{y_e}f^{-1}\big),
\]
which is the stated Laurent polynomial, $\partial_{y_e}\alpha^v=v_e\alpha^v$; it lies in $\mathcal L$, $\alpha^wf^{-1}$ being an insertion. (i) The evaluation at $z$ is multiplicative on Laurent polynomials and gives $z^w\big((m_{0,e}+w_e)f^{-1}(z)+(1-k)\,(\partial_{y_e}f^{-1})(z)\big)=0$. (ii) It gives $q_0+\sum_{i\ge1}q_i\,f^{-1}(z)^i=q_0$. Applying $\operatorname{Sh}^{-h}$, with $\operatorname{Sh}^{-h}(r(k)\,m)=r(k-h)\operatorname{Sh}^{-h}(m)$, to the law of $[1]$ gives $\sum_jp_j(k-h)[f^{j-h}]=0$, the element of the statement, whose coefficient at $f^0$ is $p_h(k-h)\neq0$. A linear form that vanishes on $\mathcal L$ and agrees with the evaluation at $z$ on the Laurent insertions would vanish on that element and take there the value $p_h(k-h)$.
\end{proof}

\begin{hdefinition}[\hb{def:box-reduction} (The reduction of a box)]
As stated.
\end{hdefinition}

\begin{proof}
By Proposition~\hb{prop:box-relations}(iv) each $\phi\in G$ permutes the symbols of the box, so $e_Gg$ is the mean of the symbols of the orbit of $g$; orbit averages of distinct orbits have disjoint supports, and they span $\mathcal F^G_{\mathfrak B}$. The wall degree is constant on an orbit and $\phi^\ast$ keeps the Laurent symbols, as do the inverses (Proposition~\hb{prop:box-walldegree}(v)). Gaussian elimination in the order of the ranking gives a basis $(r_p)$ of $\mathcal R^G_{\mathfrak B}$, indexed by the pivots, in which $r_p$ is led by $p$ with the coefficient $1$ and involves no other pivot. For $x\in\mathcal F^G_{\mathfrak B}$, $x-\sum_px_p\,r_p$, $x_p$ the coefficient of the average of $p$ in $x$, involves no pivot and is congruent to $x$; two such combinations differ by an element of $\mathcal R^G_{\mathfrak B}$ that involves no pivot, which is $0$, since a nonzero element is led by a pivot. Every $r_p$ involves only orbits ranked at most $p$, which gives the last statement of (ii). No orbit ranks below that of $1$: the other Laurent orbits have positive wall degree (Proposition~\hb{prop:box-walldegree}(ii)). If it were a pivot, $r_{1}=1$ would lie in $\mathcal R^G_{\mathfrak B}\subseteq\mathcal R_{\mathfrak B}$ (Proposition~\hb{prop:box-relations}(iv)), and $1$ would evaluate into $\mathcal L$ (Proposition~\hb{prop:box-relations}(i)), $[1]=0$; but the values vanish on $\mathcal L$ (Proposition~\hb{prop:rays-strength}(ii)), and $\langle1\rangle_k\neq0$ for all large $k$ (Proposition~\hb{prop:nativeorder}). The monomials of $f^{-1}b_i$ have wall degree at most $\mathfrak s+1\le\mathfrak t$ (Proposition~\hb{prop:box-walldegree}(ii)), and $f^{-1}$ has its coefficients in $\Bbbk$, so $f^{-1}b_i\in\mathcal F_{\mathfrak B}$; on Laurent combinations the formal and the evaluated products agree, distinct Laurent monomials being independent functions, and $f\circ\phi=f$, so $f^{-1}b_i$ is invariant.
\end{proof}

\begin{htheorem}[\hb{thm:raybox} (The box gives a spanning set, the one-step matrix and a certificate)]
As stated.
\end{htheorem}

\begin{proof}
(i) $f^{-1}b_i-\sum_l\mathbf B_{il}b_l$ lies in $\mathcal R^G_{\mathfrak B}\subseteq\mathcal R_{\mathfrak B}$, so its evaluation lies in $\mathcal L$ (Proposition~\hb{prop:box-relations}(i), (iv)), and $[f^{-1}b_i]=\sum_l\mathbf B_{il}[b_l]$. With $\mathfrak N$ the span of the $[b_l]$, $\operatorname{Sh}^{-1}\mathfrak N\subseteq\mathfrak N$; $\operatorname{Sh}^{-1}$ is injective and semilinear, so it keeps dimensions, $\operatorname{Sh}^{-1}\mathfrak N=\mathfrak N$, and $\mathfrak N$ is a submodule; it contains $[b_1]=[1]$. The $b_l$ are invariant insertions, so $\mathfrak N\subseteq\mathfrak M^G$ (Theorem~\hb{thm:invariant}(i)). An invariant combination of Laurent symbols of wall degree at most $\mathfrak s$ is a combination of averages of Laurent orbits of wall degree at most $\mathfrak s$, and by the ranking every orbit ranked at most the highest of these is a Laurent orbit of wall degree at most $\mathfrak s$; so its reduction is a combination of the $b_l$ (Definition~\hb{def:box-reduction}(ii)), congruent to it modulo $\mathcal R^G_{\mathfrak B}$, and it has a class in $\mathfrak N$. When the $[b_l]$ are independent they are a basis of $\mathfrak N$, and the matrix of $\operatorname{Sh}^{-1}$ on a basis is unique.

(ii) As an element of $\mathcal R_{\mathfrak B}$, $f^{-1}b_i-\sum_l\mathbf B_{il}b_l$ is a combination of letter identities and of mark relations $\mathrm{MR}_e(s)$ of stems $s$, with coefficients $\nu_{i,e,s}\in\Bbbk(k)$. Evaluated, the letter identities give $0$ and $\mathrm{MR}_e(s)$ gives $(c^{(0)}_e+k\,c^{(1)}_e+\partial_{y_e})s$; put $G_{i,e}=\sum_s\nu_{i,e,s}\,s$.

(iii) $V[f^{-1}b_i]=V(\operatorname{Sh}^{-1}[b_i])=\varsigma^{-1}V[b_i]$, and by (i) and the $\mathbb C(k)$-linearity of $V$ it is $\sum_l\mathbf B_{il}V[b_l]$. The certificate is an element of $\mathcal L$, so its value vanishes identically, and $\langle f^{-1}b_i\rangle_k=\langle b_i\rangle_{k-1}$ (Proposition~\hb{prop:rays-strength}(ii)). At a real $k$ as stated, the stems, the $b_l$ and the monomials of $f^{-1}b_i$ have wall degree at most $\mathfrak t$, so their values converge absolutely (Proposition~\hb{prop:box-walldegree}(i)), and Proposition~\hb{prop:rays-strength}(i) gives the value $0$ for every term $(c^{(0)}_e+k\,c^{(1)}_e+\partial_{y_e})s$.

(iv) Write $\rho^{(k)}_e=c^{(0)}_e+k\,c^{(1)}_e+\partial_{y_e}$; by the identity in the proof of Definition~\hb{def:raymodule}, $f^{-1}\rho^{(k-1)}_e(x)=\rho^{(k)}_e(f^{-1}x)$ and $f^h\rho^{(k)}_e(x)=\rho^{(k-h)}_e(f^hx)$. With the $G_{l,e}$ of (ii), put $G^{(0)}_e=0$ and $G^{(i)}_e(k)=\sum_lw_{i-1,l}(k-1)\,G_{l,e}(k)+f^{-1}G^{(i-1)}_e(k-1)$ for $i\ge1$, and $R_i(k)=\sum_e\rho^{(k)}_e\big(G^{(i)}_e(k)\big)\in\mathcal L$. Then $f^{-i}=\sum_lw_{i,l}(k)\,b_l+R_i(k)$: for $i=0$ this is $1=b_1$, and multiplying the identity for $i-1$ at the strength $k-1$ by $f^{-1}$ and inserting $f^{-1}b_l=\sum_{l'}\mathbf B_{ll'}(k)b_{l'}+\sum_e\rho^{(k)}_e(G_{l,e}(k))$ gives it for $i$. So $\sum_iq_i(k)f^{-i}=\sum_iq_i(k)R_i(k)$, and multiplying by $f^h$ and replacing $k$ by $k+h$ gives
\[
\sum_{j=0}^hp_j(k)\,f^j=\sum_e\rho^{(k)}_e\Big(\sum_{i=0}^hq_i(k+h)\,f^hG^{(i)}_e(k+h)\Big),
\]
a certificate. Proposition~\hb{prop:rays-lawrows}(ii) gives the rest.
\end{proof}

\begin{htheorem}[\hb{thm:rayboxstop} (The search stops)]
As stated.
\end{htheorem}

\begin{proof}
Let $\mathcal L_\Bbbk$ be the $\Bbbk(k)$-span of the elements $(c^{(0)}_e+k\,c^{(1)}_e+\partial_{y_e})x$, $x\in\Bbbk[X_F]$, and $M_\Bbbk=(\Bbbk[X_F]\otimes\Bbbk(k))/\mathcal L_\Bbbk$. As in the proof of Theorem~\hb{thm:rayspace}, $M_\Bbbk\otimes_{\Bbbk(k)}\mathbb C(k)=\mathfrak M$, so $M_\Bbbk$ has the finite dimension of the ray and embeds in $\mathfrak M$, and a $\Bbbk(k)$-subspace $S$ of $M_\Bbbk$ has $(S\otimes\mathbb C(k))\cap M_\Bbbk=S$, a basis of $S$ extending to one of $M_\Bbbk$. The symmetries act on $M_\Bbbk$ as on $\mathfrak M$, the weights being in $\Bbbk$. Let $\mathfrak N_{\mathfrak s}\subseteq M_\Bbbk$ be the span of the classes of the averages of the Laurent orbits of wall degree at most $\mathfrak s$. The $\mathfrak N_{\mathfrak s}$ increase with $\mathfrak s$ in a space of finite dimension, so they are constant from some $\mathfrak s_0$ on, and their union is the span of the classes of the invariant Laurent insertions, whose $\mathbb C(k)$-span in $\mathfrak M$ is $e_G\mathfrak M_{\mathrm L}$ (Lemma~\hb{lem:box-laurent}).

The averages $\bar g$ of the Laurent orbits of wall degree at most $\mathfrak s_0+1$ are finitely many (Proposition~\hb{prop:box-walldegree}(iv)), and each has its class in $\mathfrak N_{\mathfrak s_0}$:
\[
\bar g-\sum_{\bar h}\lambda_{\bar g,\bar h}\,\bar h=\sum_{e,i}\nu_{\bar g,e,i}\,(c^{(0)}_e+k\,c^{(1)}_e+\partial_{y_e})G_{\bar g,i},
\]
with $\bar h$ running over the averages of the Laurent orbits of wall degree at most $\mathfrak s_0$, $\lambda,\nu\in\Bbbk(k)$, and finitely many symbols $G_{\bar g,i}$, the evaluations of the symbols spanning $\Bbbk[X_F]$ over $\Bbbk$. Take an integer $\mathfrak t\ge\mathfrak s_0+1$ at least the wall degree of every $G_{\bar g,i}$, and the level bounds $\bar\ell(L)=1+\max\ell(L)$ for every letter $L$, the maximum over the levels $\ell$ of the symbols $\phi^\ast G_{\bar g,i}$, $\phi\in G$, at $L$. These bounds are at least $1$, and kept by $G$, the set of the $\phi^\ast G_{\bar g,i}$ being kept by $G$; every $G_{\bar g,i}$ is at a level at most $(\bar{\mathbf a}^-,\bar{\mathbf b}^-)$, here $\bar a^-_j=\bar a_j-1$ and $\bar b^-_j=\bar b_j-1$. By Proposition~\hb{prop:box-relations}(iii) every $\mathrm{MR}_e(G_{\bar g,i})$ lies in $\mathcal R_{\mathfrak B}$, and the formal combination $\bar g-\sum\lambda_{\bar g,\bar h}\bar h-\sum\nu_{\bar g,e,i}\mathrm{MR}_e(G_{\bar g,i})$ has its terms in the box (Proposition~\hb{prop:box-walldegree}(iii)) and evaluates to $0$, so it lies in $\mathcal R_{\mathfrak B}$ (Proposition~\hb{prop:box-relations}(ii)). Hence $\bar g-\sum\lambda_{\bar g,\bar h}\bar h$, which is invariant, lies in $\mathcal R_{\mathfrak B}\cap\mathcal F^G_{\mathfrak B}=\mathcal R^G_{\mathfrak B}$ (Proposition~\hb{prop:box-relations}(iv)).

Take $\mathfrak s=\mathfrak s_0<\mathfrak t$. A target $f^{-1}b_i$ is an invariant combination of Laurent symbols of wall degree at most $\mathfrak s_0+1$ (Proposition~\hb{prop:box-walldegree}(ii)), so a combination of the $\bar g$, and it is congruent modulo $\mathcal R^G_{\mathfrak B}$ to a combination of averages of Laurent orbits of wall degree at most $\mathfrak s_0$. The reduction of that combination, which is the reduction of the target, involves only free Laurent orbits of wall degree at most $\mathfrak s_0$ (Definition~\hb{def:box-reduction}(ii)): the spanning set. So the box is accepted at $\mathfrak s_0$, and its spanning set spans $\mathfrak N_{\mathfrak s_0}\otimes\mathbb C(k)=e_G\mathfrak M_{\mathrm L}$, by Theorem~\hb{thm:raybox}(i); this is $\mathfrak M^G$ when every letter exponent of the direction is positive (Lemma~\hb{lem:box-laurent}). A search through all the quadruples reaches this one, and acceptance is decided by finitely many operations in $\Bbbk(k)$.
\end{proof}

\begin{hdefinition}[\hb{def:source} (Sources; regular points; the normal density)]
As stated.
\end{hdefinition}

\begin{proof}
$Z$ is irreducible, so its smooth locus is connected; there $d\log f$ vanishes along $Z$, so $f$ is constant on it, and on $Z$ by continuity. Differentiating $d\log f=0$ along $Z$ shows that the tangent space of $Z$ at a smooth point lies in the kernel of the Hessian, whose rank is then at most $n-1-d$; the conditions defining a regular point are open, so the regular points form an open set, dense when it is nonempty, $Z$ being irreducible. For another choice of coordinates $(u',w')$ with $Z=\{w'=0\}$, at $w=0$ one has $u'=U(u)$ and $w'=J(u)w+O(|w|^2)$, the Jacobian of $(u,w)\mapsto(u',w')$ being block triangular there, so $du'\wedge dw'=\det(\partial U/\partial u)\det J\,du\wedge dw$ at $w=0$, $A=J^{\mathsf T}A'J$, and $du'=\det(\partial U/\partial u)\,du$ on $Z$. So $\det A'^{-1/2}=\pm\det J\,\det A^{-1/2}$ and $\frac{\omega_0}{du'\wedge dw'}=\big(\det(\partial U/\partial u)\det J\big)^{-1}\frac{\omega_0}{du\wedge dw}$, and the product of these with $du'$ is $\pm\det A^{-1/2}\frac{\omega_0}{du\wedge dw}\,du$: $\nu$ is defined up to sign. For $d=0$ a point is regular exactly when the Hessian has rank $n-1$, and an isolated point of the critical locus lies on no other component.
\end{proof}

\begin{htheorem}[\hb{thm:closedcycle} (A closed cycle on a source carries a solution)]
As stated.
\end{htheorem}

\begin{proof}
(i) In the coordinates $y$ the torus is, near each point, an open set of the vector space $\mathbb C^{\cE}/\mathbb C\one$, in which $z+v$ has a meaning for small $v$. At $z\in Z^\circ$ call a real subspace $V$ of dimension $m$ of the tangent space a \emph{descent subspace} when the real part of the Hessian of $-\log f$ is positive definite on it; then $V$ meets the tangent space of $Z$ only at $0$, the Hessian vanishing there. The projection to the normal space $\mathbb C^m$, in coordinates in which the normal form is $\sum_iw_i^2$, maps $V$ onto a real subspace $V'$ with $V'\cap i\mathbb R^m=0$, the form being negative on $i\mathbb R^m$; so $V'=\{x+iBx:x\in\mathbb R^m\}$ for a real matrix $B$, and the form has the real part $|x|^2-|Bx|^2$ on it, positive definite exactly when $\|B\|<1$. So the descent subspaces at $z$ are the graphs over such $V'$ of real linear maps into the tangent space of $Z$: a contractible set, whose oriented members form two contractible components, one for each sheet of $\nu$. Over the compact $\gamma\subset\tilde Z^\circ$ there is therefore a continuous oriented choice $z\mapsto V_z$, and by Taylor's formula and compactness there are $r,c>0$ with $|f(z+v)|\le|f_Z|e^{-c|v|^2}$ for $v\in V_z$, $|v|\le r$. Put $D_z=z+\{v\in V_z:|v|\le r\}$ and $N(\gamma)=\bigcup_{z\in\gamma}D_z$; as $\partial\gamma=0$, $\partial N(\gamma)$ is the union of the boundaries of the discs, on which $|f|\le|f_Z|e^{-cr^2}$. For two choices over $\gamma$, or for $\gamma-\gamma'=\partial\Gamma$, the same construction over $\gamma\times[0,1]$ or over $\Gamma$, with discs of the same kind, gives a chain of dimension $n$ whose boundary is the difference of the two chains together with a union of disc boundaries; the discs need not be affine, any smooth family of embedded discs tangent at $z$ to a descent subspace serving as well for $r$ small.

(ii) The integral over $N(\gamma)$ is the integral over $\gamma\times\{|v|\le r\}$ of the pull-back by $(z,v)\mapsto z+v$, $v\in V_z$. At each $z$, Laplace's method in $v$, uniform in $z$ on the compact $\gamma$, expands the integral over $D_z$ in powers of $1/k$ after the factor $f_Z^{\,k}k^{-m/2}$, with coefficients that are continuous $d$-forms along $\gamma$; integrating them over $\gamma$ gives the $c_i(\gamma;g)$. The leading coefficient is $(2\pi)^{m/2}\det\big[H_z(\mathsf f_a,\mathsf f_b)\big]^{-1/2}g\,\omega_0(\,\cdot\,,\mathsf f_1,\dots,\mathsf f_m)$, $H_z$ the Hessian of $-\log f$ and $\mathsf f_1,\dots,\mathsf f_m$ a basis of $V_z$, the square root fixed by the orientation; it does not change under a complex linear change of $\mathsf f$, and $V_z$ spans over $\mathbb C$ a complement of the tangent space of $Z$, which projects onto the normal space as the $w$-space does in coordinates $(u,w)$, so it is $(2\pi)^{m/2}g\nu$ on the sheet. Where holomorphic normal slices $\{u=\mathrm{const}\}$ cover $Z^\circ$, take every $D_z$ in the slice through $z$, a real descent disc there, which (i) allows. The integral over the slice at $u$ has the phase $-\log f=-\log f_Z+\frac12w^{\mathsf T}A(u)w+O(|w|^3)$ and the amplitude $g\,\omega_0/(du\wedge dw)$, and Laplace's method with the holomorphic parameter $u$ gives $f_Z^{\,k}k^{-m/2}\sum_ik^{-i}a_i(u;g)$, where $a_i(u;g)$ is $du$ times $(2\pi)^{m/2}\det A(u)^{-1/2}$ times a polynomial in finitely many $w$-derivatives at $w=0$ of the amplitude and of the phase and in the entries of $A(u)^{-1}$. These coefficients are holomorphic in $u$, so each $a_i(\cdot\,;g)$ is a holomorphic $d$-form, $a_0(\cdot\,;g)=(2\pi)^{m/2}g\,\nu$, and $c_i(\gamma;g)=\int_\gamma a_i(\cdot\,;g)$.

(iii) By \hb{lem:insertions}(ii), $(c^{(0)}_e+kc^{(1)}_e+\partial_{y_e})g\cdot f^k\omega_0/dy=\partial_{y_e}\big(gf^k\omega_0/dy\big)$, so $(c^{(0)}_e+kc^{(1)}_e+\partial_{y_e})g\cdot f^k\omega_0=d\big(\iota_{\partial_{y_e}}(gf^k\omega_0)\big)$, the form being of top degree and $\partial_{y_e}$ a constant field. By Stokes its integral over $N(\gamma)$ is an integral over $\partial N(\gamma)$, of modulus at most a constant times $|f_Z|^ke^{-kcr^2}$; so the expansion of every relation vanishes to all orders, and $A_\gamma$, extended $\mathbb C(k)$-linearly, vanishes on $\mathcal L$. The integral of $fg\,f^k\omega_0$ is that of $g\,f^{k+1}\omega_0$, so $A_\gamma(fg)(k)=A_\gamma(g)(k+1)$ after re-expansion, and $A_\gamma$ is a solution of type $(f_Z,-\frac m2+\mathbb Z)$. The forms $gf^k\omega_0$ are closed, so by Stokes on the chains of (i) homologous cycles give integrals differing by terms of modulus at most a constant times $|f_Z|^ke^{-kcr^2}$, hence the same expansion. The deck involution reverses the orientation of every disc. For an origin that is not integral the same construction runs with a branch of $\omega_0$ carried along the chains, that is, for cycles with coefficients in the local system of the branches of $\nu$.

(iv) A relation $\sum_ac_a\hat A_{\gamma_a}=0$ with $c_a\in\mathbb C((1/k))$ not all zero, scaled so that the least order of the $c_a$ is $0$ with constant terms $c^0_a$ not all zero, gives at the leading order, applied to $g_b$, $(2\pi)^{m/2}\sum_ac^0_a\int_{\gamma_a}g_b\nu=0$ for every $b$, so $c^0=0$ by the determinant: a contradiction. The relations over $\mathbb C$ are among them.
\end{proof}

\begin{hcorollary}[\hb{cor:curvesources} (A curve of native gauges: residues and periods)]
As stated.
\end{hcorollary}

\begin{proof}
The first homology of $\tilde Z^\circ$, a compact curve less finitely many points, is generated by small loops around the punctures and by the cycles of $\bar Z$, and $A_\gamma$ depends only on the class of $\gamma$ (Theorem~\hb{thm:closedcycle}(iii)). For the family $\alpha=\alpha_0b^{\mathbf v}e^w$ the coordinates $(\log b,w)$ are linear in $y$, so $dy$ is a constant multiple of $d\log b\wedge dw$. With an integral origin $f$, $g$ and $\omega_0/dy$ are rational functions of $\alpha$, and so are their derivatives in $y$, $\partial_{y_e}=\alpha_e\partial_{\alpha_e}$; their $w$-derivatives at $w=0$ are constant combinations of these at $\alpha=\alpha_0b^{\mathbf v}$, so rational in $b$, and so are the entries of $A(b)$ and $A(b)^{-1}$; the slices $\{b=\mathrm{const}\}$ cover $Z^\circ$, and the Laplace coefficients of Theorem~\hb{thm:closedcycle}(ii) are polynomials in these times $\det A(b)^{-1/2}$, so $a_i=r_i(b;g)\det A(b)^{-1/2}\,db/b$ with $r_i$ rational. On the curve $Y^2=\det A(b)$ these are meromorphic forms, and the integral over a small loop around a puncture is $2\pi i$ times the residue there.
\end{proof}

\begin{htheorem}[\hb{thm:amu} (A native gauge of corank at most one)]
As stated.
\end{htheorem}

\begin{proof}
The Hessian of $\log f$ at $\alpha_c$ has rank at least $n-2$. By the splitting lemma, the holomorphic Morse lemma in the directions where the Hessian is nondegenerate with the remaining coordinate as a parameter, there are coordinates $(X,w)$ with $\log f=\log f(\alpha_c)-\frac12\sum_iw_i^2+R(X)$, $R(0)=R'(0)=0$; when the Hessian is nondegenerate, $X$ is one of the Morse coordinates and $R=-\frac12X^2$. $R$ does not vanish identically, as otherwise $\{w=0\}$ would be a curve of native gauges through $\alpha_c$. With $\mu+1$ its order at $0$, $R=cX^{\mu+1}(1+O(X))$ with $c\neq0$, and $X(1+O(X))^{1/(\mu+1)}$ is a coordinate in which $R=cX^{\mu+1}$; $\mu=1$ exactly when the Hessian is nondegenerate, and otherwise $R''(0)=0$ and $\partial_X$ spans the kernel.

(i) Let $\varrho_0,\dots,\varrho_\mu$ be the directions with $c\,\varrho^{\mu+1}<0$, $\varrho_b=\varrho_0\zeta^b$, $\zeta=e^{2\pi i/(\mu+1)}$; around them lie the sectors on which $\re(cX^{\mu+1})<0$. For the path coming in along $\varrho_0\,\mathbb R_{>0}$ and going out along $\varrho_b\,\mathbb R_{>0}$,
\[
\int X^je^{cX^{\mu+1}}dX=\big(\varrho_b^{\,j+1}-\varrho_0^{\,j+1}\big)\frac{\Gamma\big(\frac{j+1}{\mu+1}\big)}{(\mu+1)|c|^{(j+1)/(\mu+1)}},
\]
and these $\mu$ paths generate the homology of the plane relative to the part where $\re(cX^{\mu+1})<-1$. The $\mu\times\mu$ matrix of the numbers $\zeta^{b(j+1)}-1$, $1\le b\le\mu$, $0\le j<\mu$, is invertible: its columns, extended by the entry $0$ at $b=0$, are the differences of the columns $l=j+1$ and $l=0$ of the invertible matrix $(\zeta^{bl})_{0\le b,l\le\mu}$. So there are complex combinations $\Gamma_i$ of the paths with $\int_{\Gamma_i}X^je^{cX^{\mu+1}}dX=\delta_{ij}$, $0\le j<\mu$. Integration by parts gives $\int_\Gamma X^{j+\mu+1}e^{cX^{\mu+1}}dX=-\frac{j+1}{(\mu+1)c}\int_\Gamma X^je^{cX^{\mu+1}}dX$, and $\int_\Gamma X^\mu e^{cX^{\mu+1}}dX=0$, the integrand being a derivative; so $\int_{\Gamma_i}X^je^{cX^{\mu+1}}dX=0$ unless $j\equiv i\bmod\mu+1$. On the chain $C_i=D\times\Gamma_i$, $D$ the real ball $|w|\le\varepsilon$ and $\Gamma_i$ cut off at $|X|=\varepsilon$, the boundary lies where $|f|\le|f(\alpha_c)|e^{-\varepsilon'}$, and the relations and the shift hold up to exponentially smaller terms as in Theorem~\hb{thm:closedcycle}(iii). Expand the amplitude $\mathsf a=g\,\omega_0/(dX\wedge dw)=\sum\mathsf a_{j,\beta}X^jw^\beta$; after $X\mapsto k^{-1/(\mu+1)}X$ and $w\mapsto k^{-1/2}w$ the term $(j,\beta)$ contributes $k^{-(n-2)/2-|\beta|/2-(j+1)/(\mu+1)}$ times a Gaussian moment, zero for odd $|\beta|$, times $\int_{\Gamma_i}X^je^{cX^{\mu+1}}dX$, and the usual bounds of Laplace's method make the sum an asymptotic expansion. Only $j\equiv i\bmod\mu+1$ survive, the exponents are $-\frac{n-2}2-\frac{i+1}{\mu+1}$ less integers, and the leading one, $j=i$ and $\beta=0$, has the coefficient $(2\pi)^{(n-2)/2}\mathsf a_{i,0}$. Since $0<\frac{i+1}{\mu+1}<1$, distinct $i$ give distinct types; each $A_{c,i}$ is nonzero, an insertion with $\mathsf a_{i,0}\neq0$ existing, and Lemma~\hb{lem:types} gives the independence. For $\mu=1$, $R=-\frac12X^2$, the normalization $\int_{\Gamma_0}e^{-X^2/2}dX=1$ makes $\Gamma_0$ the real line divided by $(2\pi)^{1/2}$, and $C_0$ is a steepest-descent disc divided by $(2\pi)^{1/2}$: $A_{c,0}=(2\pi)^{-1/2}A_c$.

(ii) In the coordinates $y$, $\phi$ is $y\mapsto M_\phi y$, linear, of finite order, fixing $\log\alpha_c$; it keeps $\log f$, so its Hessian and the kernel: $M_\phi\xi=\zeta_\phi\xi$ for $\xi$ spanning the kernel, with $\zeta_\phi$ a root of unity, and $M_\phi$ maps an invariant complement $N$ of the kernel onto itself. In linear coordinates $(X',w')$ along $\xi$ and $N$ the curve $\{\partial_{w'}\log f=0\}$ through $\alpha_c$ is smooth and parametrized by $X'$, the Hessian being nondegenerate on $N$; $\phi$ maps it onto itself by $X'\mapsto\zeta_\phi X'$, and $\log f$ on it is $\log f(\alpha_c)$ plus a function of order $\mu+1$ in $X'$ with nonzero leading coefficient, the reduced function of the splitting lemma. Its invariance gives $\zeta_\phi^{\mu+1}=1$.

As in Theorem~\hb{thm:invariant}(iii), $\phi^\ast\omega_0=c_\phi\,\omega_0$ near $\alpha_c$ for a constant $c_\phi\neq0$, so $A_{c,i}(g\circ\phi)$, the expansion over $C_i$ of $(g\circ\phi)f^k\omega_0=c_\phi^{-1}\phi^\ast(gf^k\omega_0)$, is $c_\phi^{-1}$ times the expansion over $\phi(C_i)$, a chain near $\alpha_c$ with its boundary where $|f|$ is smaller. The relative homology near $\alpha_c$ has rank $\mu$ over $\mathbb C$ with the basis $C_0,\dots,C_{\mu-1}$, so $A_{c,i}\circ\phi^\ast=\sum_{i'}\kappa_{i'}A_{c,i'}$. By (i) applied to the insertion $g\circ\phi$, $A_{c,i}\circ\phi^\ast$ has the type of $A_{c,i}$, and the $A_{c,i'}$ have distinct types, so $A_{c,i}\circ\phi^\ast=\kappa_iA_{c,i}$. The expansions depend only on the Taylor series of the insertion at $\alpha_c$, and insertions realize every finite jet there, so the identity holds on formal power series at $\alpha_c$. Let $\ell$ be the linear function of $y-\log\alpha_c$ that vanishes on $N$ with $\ell(\xi)=1$; then $\ell\circ\phi=\zeta_\phi\ell$. In the coordinates of (i), $\ell(X,0)=\ell_1X+O(X^2)$ with $\ell_1\neq0$, since $\partial_X$ is along $\xi$ at $\alpha_c$, so the coefficient of $X^i$ in $\ell^i\,\omega_0/(dX\wedge dw)$ at $w=0$ is $\ell_1^{\,i}\,(\omega_0/(dX\wedge dw))(\alpha_c)\neq0$, and $A_{c,i}(\ell^i)\neq0$. Then $\kappa_iA_{c,i}(\ell^i)=A_{c,i}(\ell^i\circ\phi)=\zeta_\phi^{\,i}A_{c,i}(\ell^i)$ gives $\kappa_i=\zeta_\phi^{\,i}$.
\end{proof}

\begin{hcorollary}[\hb{cor:signparity} (Signs at a centre)]
As stated.
\end{hcorollary}

\begin{proof}
The centre ray has origin $0$ and direction $(\sigma\one;\boldsymbol\sigma,\boldsymbol\sigma)$, so the Cremona is a symmetry by Proposition~\hb{prop:raysym}(ii). It fixes $\eta=\eta^{-1}$, and in the chart $\alpha=\eta e^y$ it is $y\mapsto-y$. (i) is Theorem~\hb{thm:invariant}(iii). In (ii), $\zeta=-1$, so $\mu+1$ is even, and corank one makes $\mu\ge2$, so $\mu\ge3$; Theorem~\hb{thm:amu}(ii) gives $A_{\eta,i}\circ\operatorname{cr}^\ast=(-1)^iA_{\eta,i}$. For a Cremona-invariant insertion $g$, $A_{\eta,i}(g)=A_{\eta,i}(g\circ\operatorname{cr})=(-1)^iA_{\eta,i}(g)$ vanishes for odd $i$; for even $i$ and an insertion $g'$ with $A_{\eta,i}(g')\neq0$, $\frac12(g'+g'\circ\operatorname{cr})$ is invariant with the same value. The exponents are those of Theorem~\hb{thm:amu}(i) at $\mu=3$.
\end{proof}

\begin{hproposition}[\hb{prop:sourcesym} (Symmetries on the cycles of a source)]
As stated.
\end{hproposition}

\begin{proof}
$\phi$ is an automorphism of $X_F$ keeping $f$, so it maps the critical locus onto itself, sources onto sources, regular points onto regular points, and descent discs at $z$ onto descent discs at $\phi(z)$, with their orientations; the sheet of $\nu$ at $\phi(z)$ that matches the image orientation defines $\tilde\phi$, and $\phi(N(\gamma))=N(\tilde\phi_\ast\gamma)$ as oriented chains. With $f\circ\phi=f$ and $\varepsilon_\phi^2=1$, $(g\circ\phi)f^k\omega_0=\varepsilon_\phi\,\phi^\ast(gf^k\omega_0)$, so $\int_{N(\gamma)}(g\circ\phi)f^k\omega_0=\varepsilon_\phi\int_{N(\tilde\phi_\ast\gamma)}gf^k\omega_0$. On $\mathfrak M^G$, $A_\gamma=A_\gamma\circ e_G=|G|^{-1}\sum_\phi\varepsilon_\phi A_{\tilde\phi_\ast\gamma}$, which is linear in the cycle. Grouping $\phi\in G$ by the cosets $\psi G_Z$, with $\varepsilon_{\psi\chi}=\varepsilon_\psi\varepsilon_\chi$ and $\widetilde{\psi\chi}_\ast=\tilde\psi_\ast\tilde\chi_\ast$, each coset contributes $\varepsilon_\psi A_{\tilde\psi_\ast z}$ with $z=\sum_{\chi\in G_Z}\varepsilon_\chi\tilde\chi_\ast\gamma$, zero when $[z]=0$.
\end{proof}

\begin{htheorem}[\hb{thm:invariantcount} (What the invariant module sees at the high end)]
As stated.
\end{htheorem}

\begin{proof}
For $\phi\in G$ and $\psi\in G_c$, $A_{c,i}\circ(\phi\psi)^\ast=A_{c,i}\circ\psi^\ast\circ\phi^\ast=\zeta_\psi^{\,i}A_{c,i}\circ\phi^\ast$ (Theorem~\hb{thm:amu}(ii), and Theorem~\hb{thm:invariant}(iii) for $\mu=1$), so
\[
A_{c,i}\circ e_G=|G|^{-1}\sum_{\phi G_c}\Big(\sum_{\psi\in G_c}\zeta_\psi^{\,i}\Big)A_{c,i}\circ\phi^\ast .
\]
When $\psi\mapsto\zeta_\psi^{\,i}$ is a nontrivial character of $G_c$ its sum over $G_c$ vanishes and $A_{c,i}$ vanishes on $\mathfrak M^G=e_G\mathfrak M$. Otherwise $A_{c,i}\circ e_G$ is $|G_c|/|G|$ times $A_{O,i}=\sum_{\phi G_c}A_{c,i}\circ\phi^\ast$, which is invariant, and each $A_{c,i}\circ\phi^\ast$ is a constant multiple of an expansion over the chain $\phi(C_i)$ near $\phi(\alpha_c)$. The isolated native gauges are finitely many. Take an insertion with a Taylor jet at $\alpha_c$ that makes the leading coefficient of $A_{c,i}$ nonzero (Theorem~\hb{thm:amu}(i)) and that vanishes to a high order at every other isolated native gauge; on it $A_{O,i}$ has that leading coefficient, so $A_{O,i}$ does not vanish on $\mathfrak M^G$, and the same insertions show that the $A_{O,i}$ of one type, over the orbits and the classes, are linearly independent over $\mathbb C$ on $\mathfrak M^G$, the solutions of one type at one point being a single $A_{c,i}$. Lemma~\hb{lem:types} gives the bound.
\end{proof}

\begin{hdefinition}[\hb{def:closings} (Closings)]
As stated.
\end{hdefinition}

\begin{proof}
$\mathbb C[X_F]$ is spanned by the insertions $g=\alpha^u\prod_jK_j^{-a_j}\Psi_j^{-b_j}$ with $a_j,b_j\in\mathbb N$, products of the generators of \hb{def:ray}, and for them $g\,f^{-j}\omega_0/dy=\alpha^{u-j\mvec_1+\mvec_0}\prod_lK_l^{\,js_{1l}-s_{0l}-a_l}\Psi_l^{\,jt_{1l}-t_{0l}-b_l}$, a Laurent polynomial when every exponent of a letter is nonnegative, which holds for large $j$ as $s_{1l},t_{1l}\ge1$; a finite sum takes the largest of the bounds. The function is homogeneous of degree $0$, so its monomials $\alpha^u$ have $u_{\cE}=0$, and on the chart $\alpha_{e_0}=1$ its constant term, the coefficient of the monomials with $u_e=0$ for $e\neq e_0$, is the coefficient of $\alpha^0$, whatever $e_0$. On the chart $dy=\prod_{e\neq e_0}d\alpha_e/\alpha_e$, which gives the integral over the torus.
\end{proof}

\begin{hproposition}[\hb{prop:rays-pencil} (One form on the orthant and on the torus)]
As stated.
\end{hproposition}

\begin{proof}
On the orthant $f>0$ and $1/(f^{-1}-q)=f/(1-qf)$, so the first integral is the generating function of Theorem~\hb{thm:pencil}(ii) with $k_0=1$, which the origin $0$ allows. The Laurent polynomial $f^{-1}=\alpha^{-\mvec_1}\prod_jK_j^{s_{1j}}\Psi_j^{t_{1j}}$ has nonnegative coefficients, so on the torus through $\alpha^\star$, $|f^{-1}(\alpha)|\le f^{-1}(\alpha^\star)=A$. For $|q|>A$ the series $1/(f^{-1}-q)=-\sum_{j\ge0}f^{-j}q^{-j-1}$ converges uniformly on that torus, and term by term $(2\pi i)^{-(n-1)}\oint f^{-j}\,dy$ is the constant term, $\oint\alpha^v\,dy=(2\pi i)^{n-1}$ for $v=0$ and $0$ for the other monomials of degree $0$. For every letter exponent positive this is $\mathsf C_j(1)$ (Definition~\hb{def:closings}, the origin form being $dy$).
\end{proof}

\begin{hproposition}[\hb{prop:rays-signclosings} (The constant terms of a centre ray and its signs)]
As stated.
\end{hproposition}

\begin{proof}
By Definition~\hb{def:centreray}, $f^{-1}=\prod_l\big(K_l(\alpha)K_l(\alpha^{-1})\big)^{\sigma_l}$, a Laurent polynomial with nonnegative coefficients, homogeneous of degree $0$; at a sign $\eta^{-1}=\eta$, so $f^{-1}(\eta)=\prod_lK_l(\eta)^{2\sigma_l}$. On the chart $\alpha_{e_0}=1$ a monomial $\alpha^v$, $v_{\cE}=0$, takes at $\eta$ the value $\prod_{e\neq e_0}\eta_e^{v_e}$, whose mean over the signs is $1$ when every $v_e$, $e\neq e_0$, is even, and then $v_{e_0}=-\sum_{e\neq e_0}v_e$ is even too, and $0$ otherwise. The constant term is the coefficient at $v=0$, one of the nonnegative coefficients summed.
\end{proof}

\begin{htheorem}[\hb{thm:closings} (The closings are a solution)]
As stated.
\end{htheorem}

\begin{proof}
By \hb{lem:insertions}(ii), $c^{(0)}_e-j\,c^{(1)}_e=\partial_{y_e}\log\big(f^{-j}\omega_0/dy\big)$, so $\big(c^{(0)}_e-jc^{(1)}_e+\partial_{y_e}\big)g\cdot f^{-j}\omega_0/dy=\partial_{y_e}\big(gf^{-j}\omega_0/dy\big)$. For $j\ge j(g)$ the right side is $\partial_{y_e}$ of a Laurent polynomial, which multiplies each monomial $\alpha^u$ by $u_e$ and so has no constant term. $\mathsf C_j(fg)=\CT\big(gf^{-(j-1)}\omega_0/dy\big)=\mathsf C_{j-1}(g)$ for $j-1\ge j(g)$. An element of $\mathcal L$ is a finite sum $\sum r_{e,l}(k)\big(c^{(0)}_e+kc^{(1)}_e+\partial_{y_e}\big)h_{e,l}$, and $V_{\mathsf C}$ of it at an integer $k=-j$ off the poles of the $r_{e,l}$ with $j>j(h_{e,l})$ for all $e,l$, where the closings of $c^{(1)}_eh_{e,l}$ and of $(c^{(0)}_e+\partial_{y_e})h_{e,l}$ are defined, $c^{(0)}_e$, $c^{(1)}_e$ and $\partial_{y_e}$ costing at most one power of each letter and $f^{-1}$ supplying at least one, is $\sum r_{e,l}(-j)\,\mathsf C_j\big((c^{(0)}_e-jc^{(1)}_e+\partial_{y_e})h_{e,l}\big)=0$; so $V_{\mathsf C}$ is well defined on $\mathfrak M$ as a germ at $-\infty$. It is $\mathbb C(k)$-linear, and $V_{\mathsf C}(\operatorname{Sh}(g\otimes r))(k)=r(k+1)\,\mathsf C_{-k}(fg)=r(k+1)\,\mathsf C_{-k-1}(g)=\big(\varsigma V_{\mathsf C}(g\otimes r)\big)(k)$. Theorem~\hb{thm:determine}(i) with $v(k)=\mathsf C_{-k}(\mathbf b)$ reads $\mathsf C_{1-k}(\mathbf b)=\mathbf B_{\mathbf b}(k)\,\mathsf C_{-k}(\mathbf b)$, which at $k=-j$ is the stated identity, and the law of $[1]$ applied to $V_{\mathsf C}$ gives $\sum_ip_i(k)\,\mathsf C_{-k-i}(1)=0$, at $k=-j$ the last statement.
\end{proof}

\begin{hproposition}[\hb{prop:principal} (The closings are the top principal parts)]
As stated.
\end{hproposition}

\begin{proof}
By \hb{lem:insertions}(iii), $\langle g\rangle_k$ is a finite combination of periods $P_F$ at data $(k\mvec_1+v;k\mathbf s_1+\mathbf a,k\mathbf t_1+\mathbf b)$ on the balance wall. On the balance wall, where every proper wall has positive real part, \hb{thm:exist}(ii) and \hb{thm:projective}, $J_F=\Gamma(\kappa_{\cE})P_F$, give
\[
P_F(\mvec;\mathbf s,\mathbf t)=\sum_\pi\int_{[0,1]^{n-1}}\prod_{l=2}^nq_l^{\kappa_{U_l}-1}\prod_jk_{j,\pi}(q)^{-s_j}\psi_{j,\pi}(q)^{-t_j}\,dq,
\]
the limit $\kappa_{\cE}\downarrow0$ of the same identity with the factor $L_\pi^{-\kappa_{\cE}}$, $L_\pi\in[1,n]$. The sector letters are positive on the closed cube, so the integrand is smooth in $q$ and holomorphic in the exponents, and \hb{lem:cube} continues each term meromorphically in the walls with the poles on $\kappa_{U_l}=-i$, $i\in\mathbb N$, as a finite sum of products of at most $n-1$ factors $(\kappa_{U_l}+i_l)^{-1}$ with functions holomorphic near the point. Along the ray every wall is $k\kappa_{1,U_l}$ plus a constant, with $\kappa_{1,U_l}>0$, so $\langle g\rangle_k$ is meromorphic in $k$ with poles of order at most $n-1$.

For an integer $j\ge j(g)$ put $\epsilon=k+j$; then $\langle g\rangle_k=\langle P\rangle_\epsilon$ with the Laurent polynomial $P=gf^{-j}$, an identity for large $\re k$ that continues. For a monomial $\alpha^u$ with $u\neq0$ choose $e$ with $u_e\neq0$. The relation of \hb{thm:raylaws} at the strength $\epsilon$, with $c^{(0)}=0$ for the origin $0$, holds for complex $\epsilon$ of large real part by the same integration by parts, the integrals converging absolutely there, and gives $u_e\langle\alpha^u\rangle_\epsilon=-\epsilon\,\langle c^{(1)}_e\alpha^u\rangle_\epsilon$, then as meromorphic functions; the right side has a pole of order at most $n-2$ at $\epsilon=0$. So $\epsilon^{n-1}\langle P\rangle_\epsilon$ and $\CT(P)\,\epsilon^{n-1}\langle1\rangle_\epsilon$ have the same value at $\epsilon=0$. Finally $\langle1\rangle_\epsilon=P_F(\epsilon\mvec_1;\epsilon\mathbf s_1,\epsilon\mathbf t_1)$ is the period of the question of the direction at the strength $\epsilon$, whose walls are positive, and \hb{thm:severity}(i) gives $\epsilon^{n-1}\langle1\rangle_\epsilon\to V_1$ as $\epsilon\downarrow0$; the function being holomorphic at $0$, that is its value.
\end{proof}

\begin{hproposition}[\hb{prop:rays-lowpoles} (The low end, for every direction)]
As stated.
\end{hproposition}

\begin{proof}
(i) For an integral origin $\kappa_{0,U}$ and $\kappa_U(g)$ are integers, so $\kappa_{1,U}k$ is one at such a pole; Proposition~\hb{prop:rays-strength}(ii) gives the rest.

(ii) The argument is that of the proof of Proposition~\hb{prop:principal}, which uses of the direction only that $P=f^{-j}g$ is a Laurent polynomial. By Proposition~\hb{prop:rays-strength}(ii), $\langle g\rangle_k=\langle f^jP\rangle_k=\langle P\rangle_{k+j}$ as meromorphic functions; put $\epsilon=k+j$. For a monomial $\alpha^u$, $u\neq0$, take $e$ with $u_e\neq0$; with the origin $0$, $(k\,c^{(1)}_e+\partial_{y_e})\alpha^u$ lies in $\mathcal L$ and its value vanishes identically, so $u_e\langle\alpha^u\rangle_\epsilon=-\epsilon\langle c^{(1)}_e\alpha^u\rangle_\epsilon$, whose pole at $\epsilon=0$ has order at most $n-2$. So $\epsilon^{n-1}\langle P\rangle_\epsilon$ and $\CT(P)\,\epsilon^{n-1}\langle1\rangle_\epsilon$ have the same value at $\epsilon=0$. $\langle1\rangle_\epsilon=P_F(\epsilon\mvec_1;\epsilon\mathbf s_1,\epsilon\mathbf t_1)$ (Lemma~\ref{lem:ch35-direction}(iv)) is the period of the question of the direction (Definition~\hb{def:question}) at the strength $\epsilon$, whose walls are positive, and Theorem~\hb{thm:severity}(i) gives $\epsilon^{n-1}\langle1\rangle_\epsilon\to V_1$ as $\epsilon\downarrow0$; $\epsilon^{n-1}\langle1\rangle_\epsilon$ is holomorphic at $0$, so that is its value.
\end{proof}

\begin{hcorollary}[\hb{cor:reflected} (The law read at the low end)]
As stated.
\end{hcorollary}

\begin{proof}
The law says that $\sum_ip_i(k)f^i$ lies in $\mathcal L$, a finite sum $\sum r_{e,l}(k)\big(c^{(0)}_e+kc^{(1)}_e+\partial_{y_e}\big)h_{e,l}$. Integrating against $f^kdy$, by integration by parts for complex $k$ of large real part, where the integrals converge absolutely, $\sum_ip_i(k)\langle1\rangle_{k+i}=0$ there off the poles of the $r_{e,l}$, and so identically as meromorphic functions. At $k=-j$, $j\ge0$, the term $i\le j$ has $\langle1\rangle_{k+i}$ with a pole of order at most $n-1$ and the coefficient $V_1\mathsf C_{j-i}(1)$ of $(k+j)^{-(n-1)}$ (Proposition~\hb{prop:principal}, $j(1)=0$), and $p_i(k)=p_i(-j)+O(k+j)$; the terms $i>j$ are holomorphic there, $k+i$ being near the positive integer $i-j$, where the walls are positive. So the coefficient of $(k+j)^{-(n-1)}$ is $V_1\sum_{i\le j}p_i(-j)\mathsf C_{j-i}(1)=0$, with $V_1>0$. On the centre ray $f=\prod_j\big(K_j(\alpha)K_j(\alpha^{-1})\big)^{-\sigma_j}$ (\hb{def:centreray}).
\end{proof}

\begin{hdefinition}[\hb{def:rays-lowend} (The low-end coefficients)]
As stated.
\end{hdefinition}

\begin{proof}
By Proposition~\hb{prop:rays-strength}(ii) the value of $1$ is meromorphic with a pole of order at most $n-1$ at $0$, and by Proposition~\hb{prop:rays-lowpoles}(ii) at $j=0$ the coefficient of $k^{1-n}$ is $V_1$, a sum of products of positive walls.
\end{proof}

\begin{hcorollary}[\hb{cor:rays-lowend} (The low-end relation)]
As stated.
\end{hcorollary}

\begin{proof}
(i) By Proposition~\hb{prop:rays-strength}(ii) the value of $\sum_ip_i(k)f^i$, which is $\sum_ip_i(k)\langle1\rangle_{k+i}$, vanishes identically, and every $\langle1\rangle_{k+i}$ converges for $\re k>0$, the walls $(k+i)\kappa_{1,U}$ having positive real parts. Near $k=-j$, for $i\le j$, $\lim(k+j)^{n-1}\langle1\rangle_{k+i}=V_1\CT(f^{i-j})$ by Proposition~\hb{prop:rays-lowpoles}(ii), while for $i>j$ the value $\langle1\rangle_{k+i}$ is holomorphic at $k=-j$, $k+i$ being near a positive integer. Multiplying by $(k+j)^{n-1}$ and letting $k\to-j$ gives $V_1\sum_{i\le\min(h,j)}p_i(-j)\CT(f^{i-j})=0$, and $V_1>0$.

(ii) Near $k=0$ the terms $i\ge1$ are holomorphic, so $p_0(k)\langle1\rangle_k=p_0(k)\,k^{1-n}\big(\mathsf L_0+O(k)\big)$ is holomorphic at $0$; as $\mathsf L_0\neq0$, $p_0$ vanishes to order at least $n-1$ there, and the value at $0$ of $p_0(k)\langle1\rangle_k$ is $[k^{n-1}]p_0\cdot\mathsf L_0$.

(iii) With the origin $0$, $c^{(0)}=0$; at $k=0$ the certificate gives $\sum_ip_i(0)f^i=\sum_e\partial_{y_e}G_e(0)$, and $p_0(0)=0$ by (ii), $n\ge2$. The $G_e(0)$ are homogeneous of degree $0$, so $\partial_{y_{e_0}}G_{e_0}(0)=-\sum_{e\neq e_0}\partial_{y_e}G_{e_0}(0)$ on the chart, and $\sum_e\partial_{y_e}G_e(0)$ is the divergence of the stated field in the coordinates $y_e$, $e\neq e_0$. By the divergence theorem its integral over the cube is the flux out of the cube, and as $R\to\infty$ the integral of $\sum_{i\ge1}p_i(0)f^i$ over the cube tends to $\sum_{i\ge1}p_i(0)\langle1\rangle_i$, each $f^i$, $i\ge1$, being integrable; by (ii) this is $-[k^{n-1}]p_0\cdot\mathsf L_0$.
\end{proof}

\begin{hproposition}[\hb{prop:rays-faces} (The low-end coefficients are integrals over the faces of the sectors)]
As stated.
\end{hproposition}

\begin{proof}
By Lemma~\ref{lem:rays-sector} with $g=1$ and the origin $0$, $\langle1\rangle_k=\sum_\pi\int_{[0,1]^{n-1}}\prod_lq_l^{k\kappa_l-1}\varphi_\pi^k\,dq$ for $\re k>0$. For $h$ smooth on a neighbourhood of the cube, $\sum_Ah_{[A]}=h$: the restriction of $h$ to $\{q_l=0,\ l\notin B\}$ enters with the coefficient $\sum_{A\supseteq B}(-1)^{|A\setminus B|}$, which is $1$ for $B=\{2,\dots,n\}$ and $0$ otherwise. $h_{[A]}$ depends on $q_A$ only, and where $q_l=0$, $l\in A$, the terms of $B$ and $B\cup\{l\}$, $l\notin B$, cancel; so, one coordinate at a time (Hadamard's lemma), $h_{[A]}=\prod_{l\in A}q_l\cdot\tilde h_A$ with $\tilde h_A(q_A)=\int_{[0,1]^A}\big(\prod_{l\in A}\partial_{q_l}\big)h_{[A]}(t\,q_A)\,dt$ smooth, $t\,q_A=(t_lq_l)_{l\in A}$. With $h=\varphi_\pi^k$, holomorphic in $k$ together with its $q$-derivatives,
\[
\int_{[0,1]^{n-1}}\prod_lq_l^{k\kappa_l-1}\big(\varphi^k_\pi\big)_{[A]}\,dq=\prod_{l\notin A}\frac1{k\kappa_l}\,J_{\pi,A}(k),\qquad J_{\pi,A}(k)=\int_{[0,1]^A}\prod_{l\in A}q_l^{k\kappa_l}\,\tilde h_A(q_A,k)\,dq_A ,
\]
and $J_{\pi,A}$ is holomorphic for $\re k>-1/\max_l\kappa_l$. So $k^{n-1}\langle1\rangle_k=\sum_\pi\sum_Ak^{|A|}\prod_{l\notin A}\kappa_l^{-1}J_{\pi,A}(k)$ near $0$, and $\mathsf L_i=\sum_\pi\sum_A\prod_{l\notin A}\kappa_l^{-1}[k^{i-|A|}]J_{\pi,A}$. The Taylor coefficients of $J_{\pi,A}$ at $0$ are obtained by differentiating under the integral, the $k$-derivatives of the integrand being dominated on the cube by constants times $\prod_{l\in A}q_l^{-\varepsilon}|\log q_l|^{N}$ for $|k|$ small; by Leibniz's rule, with $\partial^m_k\varphi^k_\pi|_{k=0}=\log^m\varphi_\pi$, $\partial^\gamma_kq^{k\kappa}|_{k=0}=(\kappa\log q)^\gamma$ and $h\mapsto h_{[A]}$, $h\mapsto\tilde h_A$ linear and independent of $k$, the coefficient of $k^{i'}$ is the sum over $m+|\boldsymbol\gamma|=i'$ of the integrals of the statement, $\big(\log^m\varphi_\pi\big)_{[A]}/\prod_{l\in A}q_l$ being bounded. For $i=0$ only $A=\varnothing$ and $m=0$ enter, giving $V_1$; for $i=1$, $A=\varnothing$ with $m=1$ gives $\sum_\pi\prod_l\kappa_l^{-1}\log\varphi_\pi(0)$, the terms with $|A|=1$ and $m=0$ vanishing, $1_{[A]}=0$; the constant terms of the sector letters are $\theta_{j,T^\pm_{j,\pi}}$ (Theorem~\hb{thm:exist}(ii)). For $A=\{2,\dots,n\}$ the exponent $m$ must be at least $1$, so the whole cube enters only from $i=n$ on.
\end{proof}

\begin{hproposition}[\hb{prop:formaltype} (The high end reads the formal type)]
As stated.
\end{hproposition}

\begin{proof}
By Theorem~\hb{thm:determine}(i), $\varsigma^{-1}v_V=\mathbf Bv_V$. Divide by $\lambda^kk^\rho$: $\lambda^{-1}(1-\frac1k)^\rho\big(a_0+a_1(k-1)^{-1}+\cdots\big)=\big(\mathbf B_0+\mathbf B_1k^{-1}+\cdots\big)\big(a_0+a_1k^{-1}+\cdots\big)$. The order $0$ gives $\lambda^{-1}a_0=\mathbf B_0a_0$, and the order $1$ gives $\lambda^{-1}(a_1-\rho a_0)=\mathbf B_0a_1+\mathbf B_1a_0$; multiplied by $w$ on the left, with $w\mathbf B_0=\lambda^{-1}w$, it gives $-\lambda^{-1}\rho\,wa_0=w\mathbf B_1a_0$. For a simple eigenvalue the left and right eigenvectors pair to a nonzero number. A source whose solution does not vanish on $\mathfrak N$ has a solution vector of the form of the statement after an integral shift of $\rho$, the exponents of Theorems~\hb{thm:closedcycle} and~\hb{thm:amu} being defined modulo $\mathbb Z$.
\end{proof}

\begin{hproposition}[\hb{prop:rays-spectrum} (The spectrum at infinite strength; changes of masters)]
As stated.
\end{hproposition}

\begin{proof}
By Theorem~\hb{thm:gaugesolution}(i) and (ii), $\hat A_c$ is $\mathbb C(k)$-linear, vanishes on $\mathcal L$, and $\hat A_c(g)=C_c\,g(\alpha_c)+O(1/k)$ with $C_c=(2\pi)^{(n-1)/2}(\omega_0/dv)(\alpha_c)\neq0$, the same for every insertion $g$.

(i) Let $\sum_l\lambda_l(k)[b_l]=0$ with $\lambda\neq0$, and multiply $\lambda$ by a power of $k$ so that every $\lambda_l$ is regular at $\infty$ and $\lambda(\infty)\neq0$. Applying $\hat A_c$ and reading the coefficient of $k^0$ gives $C_c\sum_l\lambda_l(\infty)\,b_l(\alpha_c)=0$ for every $\alpha_c\in C$; as the $\beta_c$ span, $\lambda(\infty)=0$, a contradiction. So the $[b_l]$ are a basis of $\mathfrak N$, and the matrix of $\operatorname{Sh}^{-1}$ on a basis is unique.

(ii) As in Theorem~\hb{thm:determine}(i), $V[f^{-1}b_i]=\varsigma^{-1}V[b_i]$ and $V[f^{-1}b_i]=\sum_l\mathbf B_{il}V[b_l]$ for every solution $V$, whatever the independence of the $b_l$: $\varsigma^{-1}v_V=\mathbf Bv_V$ with $v_V=(V[b_l])_l$. For $V=A_c$, of type $\big(f(\alpha_c),-\frac{n-1}2+\mathbb Z\big)$, the solution vector is $\lambda^kk^\rho\big(a_0+O(1/k)\big)$ with $\lambda=f(\alpha_c)$, $\rho=-\frac{n-1}2$ and $a_0=C_c\beta_c$, and $\varsigma^{-1}$ multiplies $\lambda^kk^\rho$ by $\lambda^{-1}\big(1+O(1/k)\big)$ (Definition~\hb{def:raysolution}); the coefficient of $\lambda^kk^\rho$ in $\varsigma^{-1}v_V=\mathbf Bv_V$ gives $\lambda^{-1}a_0=\mathbf B(\infty)\,a_0$. Eigenvectors of distinct eigenvalues are independent, and a matrix with an eigenbasis taken among the $\beta_c$, $\alpha_c\in C$, is diagonalizable with the eigenvalues of those vectors, every $\beta_c$ being an eigenvector.

(iii) $\operatorname{Sh}^{-1}\tilde b_i=\sum_lT_{il}(k-1)[f^{-1}b_l]=\sum_{l,l'}T_{il}(k-1)\mathbf B_{ll'}(k)[b_{l'}]$, which is $\sum_{i'}\tilde{\mathbf B}_{ii'}(k)\,\tilde b_{i'}$, and the solution vector of $A_c$ on the $\tilde b_i$ is $T(k)\,v_{A_c}$, which satisfies $\varsigma^{-1}\tilde v=\tilde{\mathbf B}\tilde v$ by the same computation. $T(k)$ is invertible and $\hat A_c(b_1)$ has the nonzero constant term $C_c$, so $T(k)\big(\hat A_c(b_l)\big)_l=k^\nu\big(\tilde\beta_c+O(1/k)\big)$ with $\nu\in\mathbb Z$ and $\tilde\beta_c\neq0$, and the computation of (ii) with $\rho=-\frac{n-1}2+\nu$ gives $\tilde{\mathbf B}(\infty)\tilde\beta_c=f(\alpha_c)^{-1}\tilde\beta_c$. Finally $T(k)\,\mathbf B(k)\,T(k)^{-1}=\big(T(k-1)T(k)^{-1}\big)^{-1}\tilde{\mathbf B}(k)$, which tends to $\tilde{\mathbf B}(\infty)$ when $T(k-1)T(k)^{-1}\to I$; it has the characteristic polynomial of $\mathbf B(k)$, whose coefficients therefore tend to those of $\tilde{\mathbf B}(\infty)$, and the roots of a monic polynomial depend continuously on its coefficients.

(iv) $T(k-1)\,\mathbf B(k)\,T(k)^{-1}=S(k-1)^{-1}S(k-1)\,\mathbf B_\circ\,S(k)^{-1}S(k)=\mathbf B_\circ$, and $T(k-1)T(k)^{-1}=S(k-1)^{-1}S(k)=\begin{bmatrix}1&0\\-(k-1)&1\end{bmatrix}\begin{bmatrix}1&0\\k&1\end{bmatrix}=\begin{bmatrix}1&0\\1&1\end{bmatrix}$. With $S(k)^{-1}=\begin{bmatrix}1&0\\-k&1\end{bmatrix}$,
\[
\mathbf B(k)=\begin{bmatrix}x_{11}-x_{12}k&x_{12}\\(k-1)(x_{11}-x_{12}k)+x_{21}-x_{22}k&(k-1)x_{12}+x_{22}\end{bmatrix},
\]
whose entry $(2,1)$ is $-x_{12}k^2+O(k)$, whose trace is $x_{11}+x_{22}-x_{12}$ and whose determinant is $\det\mathbf B_\circ$; the trace of $\mathbf B_\circ$ is $x_{11}+x_{22}$.
\end{proof}

\begin{htheorem}[\hb{thm:separate} (Separating the sources)]
As stated.
\end{htheorem}

\begin{proof}
For new masters $\mathbf b'=T\mathbf b$, $T$ invertible over $\mathbb C((1/k))$, $\operatorname{Sh}^{-1}b'_i=\sum_lT_{il}(k-1)\operatorname{Sh}^{-1}b_l$, so the matrix of $\operatorname{Sh}^{-1}$ becomes $T(k-1)\mathbf B(k)T(k)^{-1}$. After a constant change of basis, $\mathbf B_0=\operatorname{diag}(\mathbf B^{(1)}_0,\dots,\mathbf B^{(s)}_0)$ with $\operatorname{spec}\mathbf B^{(a)}_0=\Lambda_a$. Seek $T=I+\sum_{i\ge1}T_ik^{-i}$ and a block diagonal $C=\sum_{i\ge0}C_ik^{-i}$, $C_0=\mathbf B_0$, with $T(k-1)\mathbf B(k)=C(k)T(k)$. Write $T(k-1)=\sum_i\tilde T_ik^{-i}$, where $\tilde T_i-T_i$ depends on $T_1,\dots,T_{i-1}$ only. The coefficient of $k^{-i}$ is
\[
T_i\mathbf B_0-\mathbf B_0T_i=C_i-\mathbf B_i+E_i,\qquad E_i=\sum_{0<a<i}\big(C_aT_{i-a}-\tilde T_a\mathbf B_{i-a}\big)-(\tilde T_i-T_i)\mathbf B_0,
\]
with $E_i$ known from the lower orders. Take $C_i$ the block diagonal part of $\mathbf B_i-E_i$, the diagonal blocks of $T_i$ zero, and the block $(a,b)$, $a\neq b$, of $T_i$ the solution $X$ of $X\mathbf B^{(b)}_0-\mathbf B^{(a)}_0X=(E_i-\mathbf B_i)^{(ab)}$, unique since the spectra are disjoint. The blocks of the new masters span the $\hat{\mathfrak N}_a$.

If $\bigoplus_a\hat{\mathfrak N}'_a$ is another such decomposition, the inclusion of $\hat{\mathfrak N}'_a$ followed by the projection onto $\hat{\mathfrak N}_b$, $b\neq a$, commutes with the shift; in bases it is a matrix $X$ over $\mathbb C((1/k))$ with $C'^{(a)}(k)X(k)=X(k-1)C^{(b)}(k)$. If $X\neq0$, $X=k^v(X_0+O(k^{-1}))$ with $X_0\neq0$, and the leading order $C'^{(a)}_0X_0=X_0C^{(b)}_0$ forces $X_0=0$, the spectra being disjoint. So $\hat{\mathfrak N}'_a\subseteq\hat{\mathfrak N}_a$ for every $a$, and the sums being direct and equal to the whole space, $\hat{\mathfrak N}'_a=\hat{\mathfrak N}_a$. A solution of type $(\lambda,\rho+\mathbb Z)$, extended $\mathbb C((1/k))$-linearly and restricted to $\hat{\mathfrak N}_b$, has in the basis of $\hat{\mathfrak N}_b$ the solution vector $\lambda^kk^\rho w$; if it is not zero, $w=k^{-v}(w_0+\cdots)$ with $w_0\neq0$, and the first statement of Proposition~\hb{prop:formaltype}, whose proof uses only the expansion of the matrix at $\infty$, gives $C^{(b)}(\infty)w_0=\lambda^{-1}w_0$, so $\lambda^{-1}\in\Lambda_b$. A symmetry commutes with the shift, so it carries the decomposition to one with the same property, the matrix of $\operatorname{Sh}^{-1}$ on the image basis being the same; by uniqueness it maps each $\hat{\mathfrak N}_a$ onto itself.
\end{proof}

\begin{hcorollary}[\hb{cor:glue} (Gluing at the two ends)]
As stated.
\end{hcorollary}

\begin{proof}
With $\mathcal S_{\mathrm f}$ the formal solutions in $\mathcal S$, $\operatorname{Ann}(\mathcal S)\cap\mathfrak N=\operatorname{Ann}(\mathcal S_{\mathrm f})\cap\mathfrak N_{\mathsf C}$, a submodule of $\mathfrak N_{\mathsf C}$ (Theorem~\hb{thm:determine}(ii)), and Theorem~\hb{thm:determine}(iii) applies. If $\mathfrak N$ is irreducible and $V_{\mathsf C}$ does not vanish on it, $\mathfrak N_{\mathsf C}$ is a submodule other than $\mathfrak N$, hence $0$. If $\mathfrak m$ formal solutions are linearly independent over $\mathbb C$ on $\mathfrak N$, type by type, Lemma~\hb{lem:types} makes their normalizations $\mathfrak m$ independent linear forms on $\mathfrak N\otimes\mathbb C((1/k))$, of dimension $\mathfrak m$, so no nonzero element of $\mathfrak N$ is annihilated by all of them.
\end{proof}

\begin{htheorem}[\hb{thm:onekernel} (The closings, the values and the native gauge)]
As stated.
\end{htheorem}

\begin{proof}
The origin form is $u\,dy$ with $u=\omega_0/dy$ a unit of $\mathbb C[X_F]$, and $\partial_{y_e}\log u=c^{(0)}_e$ (\hb{lem:insertions}(ii)), so $(k\,c^{(1)}_e+\partial_{y_e})(ug)=u\,(c^{(0)}_e+k\,c^{(1)}_e+\partial_{y_e})g$. Multiplication by $u$ therefore carries $\mathcal L$ onto the relation space of the ray from the origin with the same direction, commutes with $\operatorname{Sh}$, and carries the values, the closings ($\mathsf C_j(g)=\CT(ug\,f^{-j})$) and the expansions at $\alpha^\star$ of the one ray to those of the other; so we may take the origin to be $0$, $\omega_0=dy$.

By Theorem~\hb{thm:pencil}(i), $0<f\le f(\alpha^\star)$ on the positive orthant with equality only at $\alpha\propto\alpha^\star$, and $f\le C'e^{-c'\|y\|}$ on the chart. The family is one molecule: a union $S$ of molecules has $c_j(S)=0$ for every $j$ (\hb{lem:separator}), and then $\kappa_{1,S}+\kappa_{1,\cE\setminus S}=\kappa_{1,\cE}=0$ (\hb{prop:compoundwall}), against the positive walls. The Hessian of $\log K_j(e^y)$ in $y$ is the covariance of the basis indicator $\one_T$ under the basis law at $e^y$ (\hb{thm:innerMGF}), every basis having positive weight at a real point, and that of $\log\Psi_j(e^y)$ the covariance of $\one_{\cE\setminus T}$; so a real direction $u'$ in the kernel of the Hessian of $\log f$ at a real point has $\ip{u'}{\one_T}$ constant over the bases of every table, and $u'$ is constant on the one molecule (\hb{lem:circuit}), $u'\in\mathbb R\one$. The Hessian $H$ of $\log f$ at $\alpha^\star$ is negative definite on $\mathbb R^{\cE}/\mathbb R\one$, and $\alpha^\star$ is a nondegenerate native gauge. Write $f^\star=f(\alpha^\star)$, $\hat A_\star=(f^\star)^{-k}k^{(n-1)/2}A_\star$, and, for $m=\sum_lr_l(k)[g_l]$, $\langle m\rangle_k=\sum_lr_l(k)\langle g_l\rangle_k$, $V_{\mathsf C}(m)(-j)=\sum_lr_l(-j)\,\mathsf C_j(g_l)$ for all large $j$, and $Q(k)$ for a common denominator of the $r_l$. By Theorem~\hb{thm:gaugesolution}(i), the orthant near $\alpha^\star$ being a steepest-descent disc and the rest of it contributing $O((f^\star-\delta)^k)$ for some $\delta>0$, $\langle g\rangle_k$ has the asymptotic expansion $A_\star(g)$, and $\langle m\rangle_k$ the expansion $A_\star(m)$.

(ii)$\Rightarrow$(iii). If $\langle m\rangle_k=0$ at all large integers, $A_\star(m)$ is the asymptotic expansion of $0$ along the integers, so $A_\star(m)=0$.

(iii)$\Rightarrow$(ii), and the continuation. For $0<t<f^\star$ the level sets $\Sigma_t=\{f=t\}$ of the real quotient are compact real-analytic hypersurfaces without critical points. By the coarea formula, $\langle g\rangle_k=\int_0^{f^\star}t^k\rho_g(t)\,dt$ for $\re k$ large, with $\rho_g(t)=\int_{\Sigma_t}g\,dy/df$ real-analytic on $(0,f^\star)$, and by the Morse lemma at $\log\alpha^\star$, $\rho_g(t)=(f^\star-t)^{(n-3)/2}h_g(t)$ with $h_g$ real-analytic at $f^\star$. With $t=f^\star e^{-s}$, $\langle g\rangle_k=(f^\star)^{k+1}\mathcal L[p_g](k+1)$, $\mathcal L$ the Laplace transform and $p_g(s)=\rho_g(f^\star e^{-s})$, real-analytic for $s>0$ and $s^{(n-3)/2}$ times a function analytic at $0$. Write $Q(k)r_l(k)=R_l(k+1)$ with polynomials $R_l$. On distributions supported in $[0,\infty)$, $\mathcal L[Dv](\kappa)=\kappa\,\mathcal L[v](\kappa)$ with $D=d/ds$, so $Q(k)\langle m\rangle_k=(f^\star)^{k+1}\mathcal L[E](k+1)$ for $\re k$ large, with $E=\sum_lR_l(D)\big(p_{g_l}\one_{s>0}\big)$. On $(0,\infty)$, $E$ is the real-analytic function $\sum_lR_l(d/ds)\,p_{g_l}$; near $0$ it is a convergent sum $\sum_\lambda c_\lambda s^\lambda$ times $\one_{s>0}$, a finite part where needed, with $\lambda\in\frac{n-3}2+\mathbb Z$ and never a negative integer (for odd $n$ the function part is analytic at $0$), plus finitely many derivatives $\delta^{(b)}$ of the Dirac measure at $0$. By Watson's lemma
\[
\mathcal L[E](\kappa)\ \sim\ \sum_\lambda c_\lambda\,\Gamma(\lambda+1)\,\kappa^{-\lambda-1}+\sum_bd_b\,\kappa^b ,
\]
in distinct powers, with $\Gamma(\lambda+1)\neq0$, and this is the expansion of $(f^\star)^{-k-1}Q(k)\langle m\rangle_k$, which is $(f^\star)^{-1}Q(k)k^{-(n-1)/2}\hat A_\star(m)$. So if $A_\star(m)=0$, every $c_\lambda$ and every $d_b$ vanishes, $E$ vanishes near $0$, hence on $(0,\infty)$ by analyticity, hence $E=0$: $Q(k)\langle m\rangle_k=0$ for $\re k$ large, and the continuation of $\langle m\rangle_k$ vanishes identically.

(iii)$\Rightarrow$(i). By what precedes the continuation of $\langle m\rangle_k$ vanishes, and by Proposition~\hb{prop:principal}, $V_1V_{\mathsf C}(m)(-j)=\lim_{k\to-j}(k+j)^{n-1}\langle m\rangle_k=0$ for all large $j$.

(i)$\Rightarrow$(iii), and the expansion of the closings. The constant term of a Laurent polynomial is its mean over every torus $|\alpha_e|=\rho_e$. On the one through $\alpha^\star$, $\alpha=\alpha^\star e^{i\theta}$ with $\theta\in\mathbb R^{\cE}/(2\pi\mathbb Z^{\cE}+\mathbb R\one)$ and $d\theta$ the measure of the chart $\theta_{e_0}=0$,
\[
\mathsf C_j(g)=(2\pi)^{-(n-1)}\int g(\alpha^\star e^{i\theta})\,f(\alpha^\star e^{i\theta})^{-j}\,d\theta,\qquad j\ge j(g).
\]
The letters have nonnegative coefficients, so $|K_j(\alpha^\star e^{i\theta})|\le K_j(\alpha^\star)$, $|\Psi_j(\alpha^\star e^{i\theta})|\le\Psi_j(\alpha^\star)$ and $|f^{-1}(\alpha^\star e^{i\theta})|\le1/f^\star$, with equality only when $\ip\theta{\one_T}$ is the same modulo $2\pi$ for all bases $T$ of every table. For linked marks $e,e'$ some basis $T$ has $e\in T$, $e'\notin T$ and $(T\setminus\{e\})\cup\{e'\}$ a basis (\hb{lem:linkexchange}), so equality gives $\theta_e\equiv\theta_{e'}$; the family being one molecule, $\theta$ is constant modulo $2\pi$, and $\alpha=\alpha^\star$. So the torus outside a neighbourhood of $\alpha^\star$ contributes $O\big((1-\delta)^j(f^\star)^{-j}\big)$ for some $\delta>0$, $g\,f^{-j(g)}$ being bounded on it. Near $\theta=0$, $\log f(\alpha^\star e^{i\theta})-\log f^\star=-\frac12\theta^{\mathsf T}H\theta+O(|\theta|^3)$, whose real part has a nondegenerate minimum at $0$, and Laplace's method there is the formal expansion at $\alpha^\star$ under the substitution $y=i\theta$, $k=-j$. The coefficient of $k^{-i}$ in $\hat A_\star(g)$ is a sum of Gaussian moments of covariance $(-H)^{-1}$ and of degree $2i+2L$, each times one Taylor coefficient of $g$ and $L$ of $\log f$; the substitution multiplies such a term by $i^{2i+2L}(-1)^L=(-1)^i$ and keeps the covariance, while $(2\pi i)^{-(n-1)}\prod_{e\neq e_0}d\alpha_e/\alpha_e=(2\pi)^{-(n-1)}d\theta$. This gives the expansion of $j^{(n-1)/2}(f^\star)^j\,\mathsf C_j(g)$ in the statement, and the same holds for $m$, with $\hat A_\star(m)=\sum_lr_l(k)\hat A_\star(g_l)$ read at $k=-j$ as a series in $1/j$. If $V_{\mathsf C}(m)=0$, that series is the expansion of $0$, so $\hat A_\star(m)=0$.
\end{proof}

\begin{hcorollary}[\hb{cor:seeall} (When the closings miss nothing)]
As stated.
\end{hcorollary}

\begin{proof}
By Theorem~\hb{thm:onekernel}, $\operatorname{Ann}(V_{\mathsf C})\cap\mathfrak N$, $\operatorname{Ann}(\langle\cdot\rangle)\cap\mathfrak N$ and $\operatorname{Ann}(A_\star)\cap\mathfrak N$ are one submodule (Theorem~\hb{thm:determine}(ii)), and it is $0$ exactly when $\alpha^\star$ sees every nonzero submodule, being itself a submodule on which $A_\star$ vanishes. A symmetry maps the positive orthant onto itself and keeps $f$ (Proposition~\hb{prop:raysym}(ii)), so it fixes $\alpha^\star$, the one point there of largest value, and $A_\star\circ\phi^\ast=A_\star$ (Theorem~\hb{thm:invariant}(iii)); so $A_\star$, and with it the values and the closings, vanish on $(1-e_G)\mathfrak M$.

For the rest, $\mathbf B$ regular at $\infty$ makes the $\mathbb C[[1/k]]$-lattice $\Lambda$ spanned by the masters stable under $\operatorname{Sh}^{-1}$. For a submodule $\mathfrak N'$, $\Lambda'=\Lambda\cap(\mathfrak N'\otimes\mathbb C((1/k)))$ is a lattice of $\mathfrak N'\otimes\mathbb C((1/k))$, stable as well and saturated in $\Lambda$, so in a basis of $\Lambda'$ the matrix of $\operatorname{Sh}^{-1}$ is regular at $\infty$, with value there the restriction of $\mathbf B(\infty)$ to $\Lambda'/k^{-1}\Lambda'\subseteq\Lambda/k^{-1}\Lambda$. The proof of Theorem~\hb{thm:separate} runs for any basis over $\mathbb C((1/k))$ in which the matrix is regular at $\infty$, so it applies to $\mathfrak N'$, and each of its pieces maps by zero to the pieces of $\mathfrak N$ of other spectra, as in the uniqueness there, so $\mathfrak N'\otimes\mathbb C((1/k))$ is the direct sum of its intersections with the $\hat{\mathfrak N}_a$. $A_\star$ vanishes on the pieces other than $\hat{\mathfrak N}_\star$ (Theorem~\hb{thm:separate}), and not on $\hat{\mathfrak N}_\star$, since it does not vanish on $\mathfrak N$; $\hat{\mathfrak N}_\star$ being a line, $A_\star$ vanishes on $\mathfrak N'$ exactly when $\hat{\mathfrak N}_\star\not\subseteq\mathfrak N'\otimes\mathbb C((1/k))$. If every nonzero submodule contains $\hat{\mathfrak N}_\star$, any two of them meet in a nonzero submodule, completion commuting with intersection, so the intersection of all of them is nonzero and is the only minimal one. Of two nonzero submodules with intersection $0$ at most one contains $\hat{\mathfrak N}_\star$, and the closings annihilate the other.
\end{proof}

\begin{hcorollary}[\hb{cor:valuemodule} (What the closings determine)]
As stated.
\end{hcorollary}

\begin{proof}
The kernel of the values is a submodule, so $\operatorname{Sh}$ induces on the value module a bijection with $\operatorname{Sh}(r\,m)=\varsigma(r)\operatorname{Sh}(m)$, and by Theorem~\hb{thm:onekernel}, $V_{\mathsf C}$ induces on it a solution that annihilates nothing; the proofs of Theorem~\hb{thm:determine}(i), (iii) and~(iv) use no more than this. For the recurrences, $\sum_ip_i(k)\langle1\rangle_{k+i}$ is the value of $m=\sum_ip_i(k)\operatorname{Sh}^i[1]$, and $V_{\mathsf C}(m)(-j)=\sum_ip_i(-j)\,\mathsf C_{j-i}(1)$; by Theorem~\hb{thm:onekernel} the one vanishes for all large arguments exactly when the other does, so the least orders agree.
\end{proof}

\begin{hproposition}[\hb{prop:blinddim} (The closings measure what they miss)]
As stated.
\end{hproposition}

\begin{proof}
The closings are a solution (Theorem~\hb{thm:closings}), and as in the proof of Theorem~\hb{thm:determine}(iii), $(P,Q)\in\mathcal S_D$ exactly when every row of $Q\mathbf B_{\mathbf b}-P$, read in $\mathfrak N$, lies in $\mathfrak A=\operatorname{Ann}(V_{\mathsf C})\cap\mathfrak N$. In the coordinates of the masters $\mathfrak A$ is a $d$-dimensional subspace of $\mathbb C(k)^{1\times\mathfrak m}$; let $\mathfrak A_D$ be its rows of polynomials of degree at most $D$. Projecting to $d$ coordinates on which $\mathfrak A$ projects isomorphically gives $\dim\mathfrak A_D\le d(D+1)$, and polynomial rows $a_1,\dots,a_d$ spanning $\mathfrak A$, of degrees $\nu_i$, give the independent rows $\sum_ic_ia_i$, $\deg c_i\le D-\nu_i$, so $\dim\mathfrak A_D=dD+O(1)$. On $\mathcal S_D$ the map $(P,Q)\mapsto Q$ has fibres that are empty or translates of $\mathfrak A_D^{\,\mathfrak m}$, each row of $P$ being fixed modulo $\mathfrak A$ by $Q$; its image lies in the polynomials of degree at most $D$ and contains $cQ_0$ for $\deg c\le D-D_0$. So $\mathfrak m\dim\mathfrak A_D+D-D_0+1\le\dim\mathcal S_D\le\mathfrak m\dim\mathfrak A_D+D+1$ for $D\ge D_0$, which is $(\mathfrak md+1)D+O(1)$; for $d=0$, Theorem~\hb{thm:determine}(iv) gives $\dim\mathcal S_D=\max(0,D-D_0+1)$, whose steps are $0$ and $1$.
\end{proof}

\begin{hproposition}[\hb{prop:onemark} (A table one mark larger)]
As stated.
\end{hproposition}

\begin{proof}
By Corollary~\hb{cor:mergedvalues}, $\cB'$ has no loop or coloop and its letters agree, $\Psi'=K'$. The bases of $\cB'$ through $g$, less $g$, are the bases of $\cB$, and complementation in $\cE'$ carries the bases through $g$ onto the bases avoiding $g$, with the same weights; so $K'=\sum_{T\not\ni g}\theta'_T\alpha^T+\alpha_g\sum_{T\ni g}\theta'_T\alpha^{T\setminus g}=\Psi+\alpha_gK$, and $K(\one)=\Psi(\one)$. The deletion $\cB'\setminus g$ is the dual of the contraction, and a connected table has a connected deletion or a connected contraction at each mark (Tutte), duality keeping connectedness; so $\cB$ is connected, and both centre rays have positive walls (Definition~\hb{def:centreray}). Write $x_{a,j}(h)=\alpha_g^aK'^{-j}h$; the monomials and the chances of $\cB'$ are of this form, so these span $\mathbb C[X_{F'}]$.

(i) For $\alpha$ in the positive orthant, $\int_0^\infty\rho^{\nu-1}(c_0+c_1\rho)^{-s}\,d\rho=\Gst21{\nu,\ s-\nu}{s}c_0^{\nu-s}c_1^{-\nu}$ gives
\[
\int_0^\infty\alpha_g^{k+a}(\Psi+\alpha_gK)^{-2k-j}\,\frac{d\alpha_g}{\alpha_g}=\Gst21{k+a,\ k+j-a}{2k+j}\,\Psi^{a-j-k}K^{-a-k},\qquad k>\max(-a,a-j),
\]
and with $f'^k=\alpha^{k\one}\alpha_g^kK'^{-2k}$, $f^k=\alpha^{k\one}(K\Psi)^{-k}$ and $\Gst21{k+a,\ k+j-a}{2k+j}=r_{a,j}(k)\Gst21{k,\ k}{2k}$ this is the identity for $g'=x_{a,j}(h)$. The integral is linear in $g'$, so on a combination of the $x_{a,j}(h)$ the defining formula depends only on the function they sum to, both sides being rational in $k$ and equal for all large $k$; so $\pi_\ast$ is well defined and $\mathbb C(k)$-linear on $\mathbb C[X_{F'}]\otimes\mathbb C(k)$, and the identity holds for every insertion. On the chart $y_{e_0}=0$ the form of $F'$ is $dy\,d\alpha_g/\alpha_g$, and its positive orthant is that of $F$ times $(0,\infty)$; Fubini gives the values. The constant term in $\alpha_g$ of $x_{a,j}(h)f'^{-J}=\alpha_g^{a-J}\alpha^{-J\one}h\,(\Psi+\alpha_gK)^{2J-j}$ is $\binom{2J-j}{J-a}\Psi^{J-j+a}K^{J-a}\alpha^{-J\one}h$, and for $J\ge|a|+j$, $\Gamma(x-J)/\Gamma(-J)=(-1)^xJ!/(J-x)!$ gives $r_{a,j}(-J)=\binom{2J-j}{J-a}\big/\binom{2J}J$; so that constant term is $\binom{2J}J\,\pi_\ast\big(x_{a,j}(h)\big)\big|_{k=-J}f^{-J}$, and the constant term in the remaining scales gives the closings.

On $X_{F'}$, $c'^{(1)}_g=\partial_{y_g}\log f'=1-2\alpha_gK/K'$, which is $-1$ where $\Psi=0$ and $1$ where $K=0$; where $K\Psi\neq0$ it vanishes exactly at $\alpha_g=\Psi/K$, where $K'=2\Psi$, $f'=f/4$ and $\partial_{y_g}^2\log f'=-2\,\frac{\alpha_gK}{K'}\,\frac{\Psi}{K'}=-\frac12$. So $\log f(\alpha)=\log4+\log f'(\alpha,\Psi/K)$ on $X_F$, the defects of $F'$ at $(\alpha,\Psi/K)$ in the marks of $\cE$ are those of $F$ at $\alpha$, the native gauges of $F'$ are the points $\alpha_{c'}$, and the Hessian of $\log f$ at $\alpha_c$ is the Schur complement of the entry $-\frac12$ in the Hessian of $\log f'$ at $\alpha_{c'}$, whose determinant is $-\frac12$ times that of the complement. Over a steepest-descent disc $D_c$, with Morse coordinates $v$, put $\alpha_g=e^s\,\Psi/K$; then $d\alpha_g/\alpha_g=ds$ on the fibre and
\[
\log f'=\log f(\alpha_c)-\log4-\tfrac12\textstyle\sum_iv_i^2-2\log\cosh\frac s2 ,
\]
so with the real-analytic coordinate $w$ of the sign of $s$ and $\frac12w^2=2\log\cosh\frac s2$, the image of $D_c\times\{|s|\le\varepsilon\}$ is a steepest-descent disc at $\alpha_{c'}$. The integral over all $s$ is the Beta corner, the part with $|s|>\varepsilon$ is exponentially smaller, and the expansion at $\alpha_{c'}$ is that of the Beta corner times $A_c(\pi_\ast g')$.

(ii) Let $\rho'_e(g')=(k\,c'^{(1)}_e+\partial_{y_e})g'$, so that $f'^k\rho'_e(g')=\partial_{y_e}(f'^kg')$, and likewise $\rho_e$ for $F$. Differentiating the identity of (i) under the integral, $\pi_\ast\rho'_e=\rho_e\pi_\ast$ for $e\in\cE$; and $\pi_\ast\rho'_g=0$, the integral of $\partial_{y_g}(f'^kg')$ over $(0,\infty)$ vanishing for large $k$. Both are identities between elements of $\mathbb C[X_F]\otimes\mathbb C(k)$ that hold for all large $k$, so they hold. Next,
\[
\rho'_g\big(x_{a,j}(h)\big)=(k+a)\,x_{a,j}(h)-(2k+j)\,x_{a+1,j+1}(Kh),\qquad x_{a,j}(h)=x_{a,j+1}(\Psi h)+x_{a+1,j+1}(Kh),
\]
so modulo the span $\mathcal L'_g$ of the $\rho'_g(g')$
\[
x_{a,j}(h)\equiv\frac{2k+j}{k+a}\,x_{a+1,j+1}(Kh)\equiv\frac{2k+j}{k+j-a}\,x_{a,j+1}(\Psi h).
\]
A finite sum $z=\sum_ic_ix_{a_i,j_i}(h_i)$ is therefore congruent modulo $\mathcal L'_g$ to one term $x_{A,A+B}(H)$, for any $A\ge\max_ia_i$ and $B\ge\max_i(j_i-a_i)$, with $H$ a Laurent polynomial, and $\pi_\ast z=r_{A,A+B}K^{-A}\Psi^{-B}H$; if $\pi_\ast z=0$ then $H=0$ and $z\in\mathcal L'_g$. And $\pi_\ast$ is onto: $\pi_\ast x_{a,a+b}(\alpha^u)=r_{a,a+b}\,\alpha^uK^{-a}\Psi^{-b}$ for $a,b\ge0$, with $r_{a,a+b}\neq0$, and these span $\mathbb C[X_F]$. So $\pi_\ast$ is a bijection of $(\mathbb C[X_{F'}]\otimes\mathbb C(k))/\mathcal L'_g$ onto $\mathbb C[X_F]\otimes\mathbb C(k)$ that carries $\mathcal L'=\mathcal L'_g+\sum_{e\in\cE}\rho'_e(\cdot)$ onto $\sum_{e\in\cE}\rho_e(\cdot)=\mathcal L$, and it induces a bijection $\mathfrak M'\to\mathfrak M$. Finally $\Gst21{k+1,\ k+1}{2k+2}=\frac k{2(2k+1)}\Gst21{k,\ k}{2k}$, so $\pi_\ast(f'g')=\frac k{2(2k+1)}\,f\,\varsigma(\pi_\ast g')$.

(iii) By (ii) $T_\beta$ is a $\mathbb C(k)$-linear bijection of $\mathfrak M$. $\phi_\beta$ is a symmetry of $F'$ (Proposition~\hb{prop:raysym}(ii)), so $\phi_\beta^\ast$ commutes with $\operatorname{Sh}'$ (Theorem~\hb{thm:invariant}(i)); with $c=\frac k{2(2k+1)}$, $\operatorname{Sh}=c^{-1}\pi_\ast\operatorname{Sh}'\pi_\ast^{-1}$, and then $T_\beta\operatorname{Sh}=c^{-1}\pi_\ast\phi_\beta^\ast\operatorname{Sh}'\pi_\ast^{-1}=c^{-1}\pi_\ast\operatorname{Sh}'\phi_\beta^\ast\pi_\ast^{-1}=\operatorname{Sh}T_\beta$. If $\beta(g)=g$, comparing the coefficients of $\alpha_g^0$ and $\alpha_g^1$ in $K'\circ\phi_\beta=K'$ gives $K\circ\phi_\beta=K$ and $\Psi\circ\phi_\beta=\Psi$, so $\phi_\beta$ restricted to $\cE$ is a symmetry of $F$ and $\pi_\ast\phi_\beta^\ast=\phi_\beta^\ast\pi_\ast$ on the $x_{a,j}(h)$. For the transposition of $e$ and $g$, $\phi_\beta^\ast\alpha^u=x_{u_e,0}(\alpha^{u-u_e\one_e})$, and $\pi_\ast$ of it is the stated class, with $r_{u_e,0}(k)=\Gamma(k+u_e)\Gamma(k-u_e)/\Gamma(k)^2$. For the Cremona, $K'\circ\operatorname{cr}=\alpha^{-\one}\alpha_g^{-1}K'$, $K\circ\operatorname{cr}=\alpha^{-\one}\Psi$ and $\Psi\circ\operatorname{cr}=\alpha^{-\one}K$ (Proposition~\hb{prop:raysym}(i)), so $\operatorname{cr}^\ast x_{a,j}(h)=x_{j-a,j}\big(\alpha^{j\one}\,h\circ\operatorname{cr}\big)$, and both $\pi_\ast\operatorname{cr}^\ast x_{a,j}(h)$ and $\operatorname{cr}^\ast\pi_\ast x_{a,j}(h)$ are $r_{a,j}\,\Psi^{-a}K^{a-j}\alpha^{j\one}\,h\circ\operatorname{cr}$, as $r_{j-a,j}=r_{a,j}$. The group generated is the image of the finite group generated by the $\phi_\beta$ and the Cremona of $F'$.

(iv) The values and the closings of $F'$ are kept by every $\phi_\beta$ and by its Cremona (Proposition~\hb{prop:raysym}(ii); the torus and the constant term are kept by every permutation of the scales and by their inversion). By (i) the values and the closings of $F$ are those of $F'$ composed with $\pi_\ast^{-1}$, divided by $\Gst21{k,\ k}{2k}$ and by $\binom{2j}j$, which do not depend on the insertion; so they are invariant under $G'$. The native gauge of the direction of $F'$ on the positive orthant is unique up to scale (Theorem~\hb{thm:pencil}(i)), so it is fixed by every $\phi_\beta$ and by the Cremona, and it is $\one$; on the orthant $f'\le f/4$, with equality exactly at $\alpha_g=\Psi/K$, so it lies over $\alpha^\star$, and $\alpha^\star=\one$. Then $A'_\star\circ\phi_\beta^\ast=A'_\star$ (Theorem~\hb{thm:invariant}(iii)), and $A_\star\circ T_\beta=A_\star$ by (i). A solution $V$ invariant under $G'$ has $V=V\circ e_{G'}$ and vanishes on $(1-e_{G'})\mathfrak M$. If $G\subseteq G'$, then $e_{G'}e_G=e_{G'}$, so $e_{G'}$ maps $\mathfrak M^G$ onto $\mathfrak M^{G'}\subseteq\mathfrak M^G$, and $\mathfrak M^G=e_{G'}\mathfrak M^G\oplus(1-e_{G'})\mathfrak M^G$, both submodules since $e_{G'}$ commutes with $\operatorname{Sh}$. Last, $\pi_\ast$ intertwines the action on $\mathfrak M'$ with $G'$, so it carries the invariant module of $F'$ onto $\mathfrak M^{G'}$.
\end{proof}

\begin{hcorollary}[\hb{cor:cousinsonemark} (The cousins of a fused table)]
As stated.
\end{hcorollary}

\begin{proof}
$\cB^\sharp$ is connected by Corollary~\hb{cor:fusionisd}(ii), and identically self-dual. For a mark $x$ of $\cE\sqcup G$ the cousin at $x$ is the contraction $\cB^\sharp/x$, and Proposition~\hb{prop:onemark} applies with $(\cB',\theta')=(\cB^\sharp,\theta^\sharp)$, its $n+c-1\ge2$ marks other than $x$, and the mark $x$ in the place of $g$.
\end{proof}

\begin{hproposition}[\hb{prop:onemarkhidden} (What a table one mark larger hides from the closings)]
As stated.
\end{hproposition}

\begin{proof}
(i) For a nondegenerate native gauge $\alpha_{c'}$ of $F'$ and $\gamma\in\Gamma'$, $A'_{c'}\circ\gamma^\ast$ is a nonzero constant multiple of the expansion over the steepest-descent disc $\gamma(D_{c'})$ at $\gamma(\alpha_{c'})$, and $A'_{c'}\circ\psi^\ast=A'_{c'}$ when $\psi\in\Gamma'$ fixes $\alpha_{c'}$ (proof of Theorem~\hb{thm:invariant}(iii)). Two $\gamma$ that carry $\alpha_{c'_0}$ to one point differ by an element of its stabilizer, and $(\gamma\psi)^\ast=\psi^\ast\gamma^\ast$, so the transported $A'_{c'}$ is well defined, and $A'_{c'}\circ\chi^\ast=A'_{\chi(c')}$ for every $\chi\in\Gamma'$. By Proposition~\hb{prop:onemark}(i), $A_c$ is a nonzero constant multiple of the expansion of $F$ at $\alpha_c$: a solution of the type $\big(f(\alpha_c),-\frac{n-1}2+\mathbb Z\big)$, $n=|\cE|$, whose leading coefficient on an insertion $h$ is $C_c\,h(\alpha_c)$ with $C_c\neq0$ (Theorem~\hb{thm:gaugesolution}(i)). The native gauges below one orbit $O'_a$ have one native value, $4f'$ at the points of $O'_a$, which $\Gamma'$ keeps, and so one type.

Every element of $G'$ is $T_\gamma=\pi_\ast\gamma^\ast\pi_\ast^{-1}$ for some $\gamma\in\Gamma'$ (Proposition~\hb{prop:onemark}(iii)), and $A_c\circ T_\gamma=\Gst21{k,\ k}{2k}^{-1}A'_{\gamma(c')}\circ\pi_\ast^{-1}=A_{\gamma\cdot c}$, with $\gamma\cdot c$ the native gauge below $\gamma(\alpha_{c'})$. For $\phi\in G$, $\phi^\ast=T_\gamma$ with $\gamma$ the map of the same kind on the torus of $F'$, and the lift commutes with $\phi$: $\phi_\beta$, $\beta(g)=g$, keeps $K$ and $\Psi$, and $\Psi(\alpha^{-1})/K(\alpha^{-1})=K(\alpha)/\Psi(\alpha)$ (Proposition~\hb{prop:identity}). So $\gamma\cdot c=\phi(\alpha_c)$ and $A_c\circ\phi^\ast=A_{\phi(\alpha_c)}$. Let $O_{a,1},\dots,O_{a,s_a}$ be the orbits of $G$ below $O'_a$ and $V_{a,i}=\sum_{c\in O_{a,i}}A_c$; each $\phi\in G$ permutes $O_{a,i}$, so $V_{a,i}$ is invariant under $G$. For $x\in\mathfrak M^{G'}$ and native gauges $\alpha_c$, $\alpha_{\tilde c}$ below $O'_a$, with $\gamma(\alpha_{c'})=\alpha_{\tilde c'}$, $A_{\tilde c}(x)=A_c(T_\gamma x)=A_c(x)$; so $V_{a,i}(x)=|O_{a,i}|\,A_c(x)$ for any $\alpha_c$ below $O'_a$, and
\[
W_{a,i}=|O_{a,i}|^{-1}V_{a,i}-|O_{a,1}|^{-1}V_{a,1},\qquad 2\le i\le s_a,
\]
vanish on $\mathfrak M^{G'}$.

The $V_{a,i}$ of one type are linearly independent over $\mathbb C$ on $\mathfrak M^G$. Let $\sum\gamma_{a,i}V_{a,i}$ vanish on $\mathfrak M^G$, the sum over those of one type, and fix $(a,i)$ and $\alpha_c\in O_{a,i}$. As in the proof of Theorem~\hb{thm:invariant}(iii), take an insertion $h$ with $h(\alpha_c)=1$ that vanishes to the order $N\ge2$ at the other native gauges below $O'_1,\dots,O'_q$, which are finitely many; the expansion at such a point is smaller than the leading power of the type by the factor $k^{-\lceil N/2\rceil}$. The class of $e_Gh$ lies in $\mathfrak M^G$, and every $V_{a',i'}$, being invariant, takes the same value at $e_Gh$ and at $h$; so the leading coefficient of the relation at $e_Gh$ is $\gamma_{a,i}C_c$, and $\gamma_{a,i}=0$. Each $W_{a,i}$ involves its own $V_{a,i}$, so the $W_{a,i}$ of one type are linearly independent over $\mathbb C$ on $\mathfrak M^G$. They vanish on $\mathfrak M^{G'}$, and $\mathfrak M^G=\mathfrak M^{G'}\oplus(1-e_{G'})\mathfrak M^G$ (Proposition~\hb{prop:onemark}(iv)), so a combination of them that vanishes on $(1-e_{G'})\mathfrak M^G$ vanishes on $\mathfrak M^G$: they are linearly independent over $\mathbb C$ on $(1-e_{G'})\mathfrak M^G$, type by type. Lemma~\hb{lem:types} gives the bound; $(1-e_{G'})\mathfrak M^G$ is a submodule, $e_{G'}$ commuting with $\operatorname{Sh}$, and the closings, the values and $A_\star$ vanish on it by Proposition~\hb{prop:onemark}(iv). The solution of the last statement is $-W_{a,2}$ for an indexing with $O_{a,1}=O_1$ and $O_{a,2}=O_2$.

(ii) A nonempty proper $S\subseteq[2r+2]$ meets the bases of $\binom{[2r+2]}{r+1}$ in every number of marks from $\max(0,|S|-r-1)$ to $\min(|S|,r+1)$, so $c(S)>0$ and the table is connected (Definition~\hb{def:family}); the complement of a set of $r+1$ of its marks has $r+1$ marks, so it is identically self-dual; and its contraction at $g$ (Definition~\hb{def:delcon}) has the bases $T\setminus\{g\}$, $T\ni g$, the sets of $r$ marks of $[2r+1]$, with $K=e_r$ and $\Psi=e_{r+1}$, the elementary symmetric polynomials of $\alpha_1,\dots,\alpha_{2r+1}$. So Proposition~\hb{prop:onemark} applies, $|\cE|=2r+1\ge5$. Every bijection of $[2r+2]$ is an automorphism, so $G$ consists of the $\phi_\beta$, $\beta$ a bijection of $[2r+1]$, and their composites with the Cremona. These are the symmetries of the four kinds of the centre ray (Proposition~\hb{prop:raysym}(ii), the masses being $\one$ and the exponents $s=t=1$): a bijection that exchanged the list would carry its bases, of $r$ marks, onto complements of bases, of $r+1$ marks, so there are no exchanges and no reversed exchanges.

At a sign $\eta$ with $i$ coordinates $-1$ among $l+i$, $e_m(\eta)$ is the coefficient of $z^m$ in $(1-z)^i(1+z)^l$. With $\binom{2r}{r-1}=\binom{2r}{r+1}=rC_r$, $\binom{2r-1}{r-1}=\binom{2r-1}r=\frac{r+1}2C_r$ and $\binom{2r-1}{r-2}=\frac{r-1}2C_r$,
\[
[z^r](1-z)(1+z)^{2r}=C_r,\qquad [z^r](1-z)^2(1+z)^{2r-1}=-C_r,
\]
so $K(\eta)=e_r(\eta)=\pm C_r\neq0$ at the signs of $F$ with one and with two scales $-1$, which are native gauges of the native value $K(\eta)^{-2}=C_r^{-2}$ (Theorem~\hb{thm:signgauges}(ii)). The bijections of $[2r+1]$ act transitively on the marks and on the pairs of marks, the Cremona fixes every sign, and a sign with one coordinate $-1$ is not a multiple of one with two, $2r+1\ge5$: the signs form two orbits of $G$. At a sign $\eta^{-1}=\eta$, so $\Psi(\eta)=\eta^{\cE}K(\eta)$ (Proposition~\hb{prop:identity}), and the native gauge above $\eta$ is $(\eta,\eta^{\cE})$, with $\eta^{\cE}=-1$ for one scale $-1$ and $\eta^{\cE}=1$ for two: the native gauges above the two kinds are the $\binom{2r+2}2$ signs of $F'$ with two scales $-1$, each once, and the bijections of $[2r+2]$ act transitively on them.

At such a sign $\eta'$, with $D=\{a_1,a_2\}$ and $C$ as in the statement, $a\in D$ and $c,c'\in C$, the Hessian of $-\log f'$ is $2\kappa$, $\kappa$ the covariances of the indicators $X_e$ of the marks under the signed measure $\mathbb P(T)=\eta'^T/K'(\eta')$ on the sets $T$ of $r+1$ marks (Theorem~\hb{thm:signgauges}(ii), $K'=\Psi'$); as a matrix on $\mathbb C^{\cE'}$ it has zero row sums, $\sum_eX_e=r+1$. From $[z^{r+1}](1-z)^2(1+z)^{2r}=2rC_r-2(r+1)C_r$, $[z^{r-1}](1-z)(1+z)^{2r-1}=C_r$ and, with $\binom{2r-2}{r-1}=\frac{r(r+1)}{2(2r-1)}C_r$,
\[
[z^{r-1}](1-z)^2(1+z)^{2r-2}=\binom{2r-2}{r-1}\Big(1-\frac{2(r-1)}r+\frac{(r-1)(r-2)}{r(r+1)}\Big)=-\frac{r-2}{2r-1}\,C_r ,
\]
\begin{gather*}
K'(\eta')=-2C_r,\qquad \mathbb P(e\in T)=\tfrac12,\qquad \mathbb P(a_1,a_2\in T)=\frac{rC_r}{-2C_r}=-\frac r2,\\
\mathbb P(a,c\in T)=\frac12,\qquad \mathbb P(c,c'\in T)=\frac{r-2}{2(2r-1)},
\end{gather*}
the first for every mark, as $\eta'_ee_r(\eta'_{\hat e})=-C_r$ for $e\in D$ and for $e\in C$. So $\kappa_{ee}=\frac14$, $\kappa_{a_1a_2}=-\frac{2r+1}4$, $\kappa_{ac}=\frac14$ and $\kappa_{cc'}=-\frac3{4(2r-1)}$, and
\begin{gather*}
\kappa(\one_{a_1}-\one_{a_2})=\tfrac{r+1}2(\one_{a_1}-\one_{a_2}),\qquad \kappa(\one_c-\one_{c'})=\tfrac{r+1}{2(2r-1)}(\one_c-\one_{c'}),\\
\kappa(r\one_D-\one_C)=-\tfrac{r+1}2(r\one_D-\one_C),
\end{gather*}
the last from its coordinates $r(\frac14-\frac{2r+1}4)-\frac{2r}4=-\frac{r(r+1)}2$ at $a$ and $\frac{2r}4-\frac14+\frac34=\frac{r+1}2$ at $c$, where $r\one_D-\one_C$ has the coordinates $r$ and $-1$. With $c_1\in C$, the vectors $\one_{a_1}-\one_{a_2}$, $r\one_D-\one_C$ and $\one_{c_1}-\one_c$, $c\in C\setminus\{c_1\}$, are $2r+1$ independent vectors orthogonal to $\one$, which gives the eigenvalues of $-H'=2\kappa$ modulo $\one$; they do not vanish, so $\eta'$ is nondegenerate, and so is the sign below it (Proposition~\hb{prop:onemark}(i)). A symmetric matrix $A$ of size $N'$ with $A\one=0$ and rank $N'-1$ has an adjugate of rank one killed by $A$, so equal to $\lambda\,\one\one^{\mathsf T}/N'$ with $\lambda$ the sum of its principal minors of size $N'-1$, the product of the eigenvalues of $A$ modulo $\one$; so every such minor is $\lambda/N'$, and on a chart
\[
\det(-H')=\frac{(r+1)\cdot(-(r+1))\cdot\big(\frac{r+1}{2r-1}\big)^{2r-1}}{2r+2}=-\frac{(r+1)^{2r}}{2(2r-1)^{2r-1}} .
\]
The Hessian of $\log f$ at $\eta$ is the Schur complement of the entry $-\frac12$ in that of $\log f'$ at $\eta'$ (proof of Proposition~\hb{prop:onemark}(i)), so on charts with $e_0\in\cE$, $\det H=-2\det H'$ and $\det(-H)=2\det(-H')$, which is the stated value. Part (i) with the one orbit of the signs of $F'$ with two scales $-1$, below which lie two orbits of $G$, gives $(1-e_{G'})\mathfrak M^G\neq0$, and the closings vanish on it.
\end{proof}

\begin{hproposition}[\hb{prop:wheelholonomy} (The holonomy of a wheel)]
As stated.
\end{hproposition}

\begin{proof}
Put $S(a)=\mathbf U(a)$ and $R(b)=\begin{bmatrix}b&0\\1&b\end{bmatrix}=b\,\mathbf L(b^{-1})$, and $M(\alpha)=\prod_iS(\alpha_{s_i})R(\alpha_{r_i})=\prod_i\alpha_{r_i}\cdot\mathbf h(\alpha)$. Index rows and columns by the states $0,1$ and read the factors from left to right as steps around the rim, $S(\alpha_{s_i})$ at the vertex $i$ and $R(\alpha_{r_i})$ at the rim $r_i$: $S$ allows $0\to0$, $0\to1$ with the weight $\alpha_{s_i}$, and $1\to1$; $R$ allows $0\to0$ and $1\to1$ with the weight $\alpha_{r_i}$, and $1\to0$ with the weight $1$. Then $\operatorname{tr}M$ is the sum over the closed walks of states of the products of the weights. A spanning tree of the wheel omits at least one rim; the rims it omits cut the rim cycle into arcs of consecutive vertices, joined by the rims it contains, and it contains exactly one spoke at a vertex of each arc, at least one for connection and at most one, two closing a cycle through the hub; every such choice is a spanning tree. To a tree corresponds the closed walk in state $0$ at the start of each arc and in state $1$ from the vertex of its spoke on, leaving $1$ for $0$ exactly at the omitted rims; its weight is $\alpha^T$. The closed walks that never go from $1$ to $0$ stay in one state all around, taking every rim and no spoke: two walks of weight $\prod_i\alpha_{r_i}$. So $\operatorname{tr}M=K+2\prod_i\alpha_{r_i}$, which is the formula for $K$. $\det S=1$ and $\det R(b)=b^2$, so $\det\mathbf h=1$. By \hb{prop:identity}, $\Psi(\alpha)=\alpha^{\cE}K(\alpha^{-1})=\prod_i\alpha_{s_i}\alpha_{r_i}\cdot\prod_i\alpha_{r_i}^{-1}\big(\operatorname{tr}\mathbf h(\alpha^{-1})-2\big)$.

The transpose of a product is the product of the transposes in the reverse order, $\mathbf U(a)^{\mathsf T}=\mathbf L(a)$ and $\mathbf L(b)^{\mathsf T}=\mathbf U(b)$, so
\[
\operatorname{tr}\mathbf h(\alpha^{-1})=\operatorname{tr}\Big(\prod_{i=1}^p\mathbf U(\alpha_{s_i}^{-1})\mathbf L(\alpha_{r_i})\Big)^{\mathsf T}=\operatorname{tr}\big(\mathbf U(\alpha_{r_p})\mathbf L(\alpha_{s_p}^{-1})\cdots\mathbf U(\alpha_{r_1})\mathbf L(\alpha_{s_1}^{-1})\big)=\operatorname{tr}\mathbf h(\alpha\circ x),
\]
the factor $i$ of the last product being $\mathbf U(\alpha_{r_{p+1-i}})\mathbf L(\alpha_{s_{p+1-i}}^{-1})=\mathbf U\big((\alpha\circ x)_{s_i}\big)\mathbf L\big((\alpha\circ x)_{r_i}^{-1}\big)$. Then $K(\alpha\circ x)=\prod_i\alpha_{x(r_i)}\big(\operatorname{tr}\mathbf h(\alpha\circ x)-2\big)=\prod_i\alpha_{s_i}\big(\operatorname{tr}\mathbf h(\alpha^{-1})-2\big)=\Psi(\alpha)$. Comparing the monomials, all with coefficient $1$, the sets $x(T)$ of the bases $T$ are the complements of the bases, so $T\mapsto\cE\setminus x(T)$ maps the bases onto the bases: $x$ exchanges the list. At $\alpha=\one$, $\mathbf U(1)\mathbf L(1)=\begin{bmatrix}2&1\\1&1\end{bmatrix}$, of trace $3$ and determinant $1$, with the eigenvalues $\lambda_\pm$, the squares of the golden ratio and of its inverse, and $\mathbf h(\one)$ is its $p$-th power, of trace $\lambda_+^p+\lambda_-^p$, the Lucas number of index $2p$, the Lucas numbers being the sums of the powers of the golden ratio and of minus its inverse.
\end{proof}

\begin{hproposition}[\hb{prop:wheelcentre} (The centre ray of a wheel)]
As stated.
\end{hproposition}

\begin{proof}
By \hb{def:centreray}, $f=\alpha^{\sigma\cE}(K\Psi)^{-\sigma}$, and by Proposition~\hb{prop:wheelholonomy}, $K\Psi=\alpha^{\cE}\,\mathfrak d\,\mathfrak d^\ast$; so $f=(\mathfrak d\,\mathfrak d^\ast)^{-\sigma}$, and $X_F=\{K\Psi\neq0\}=\{\mathfrak d\,\mathfrak d^\ast\neq0\}$ on the torus. (i) The native gauges are the critical points of $\log f=-\sigma(\log\mathfrak d+\log\mathfrak d^\ast)$ on $X_F$ (\hb{lem:insertions}(ii), \hb{def:raygauge}), where $\mathfrak d^\ast d\mathfrak d+\mathfrak d\,d\mathfrak d^\ast=0$, and then $d\mathfrak d^\ast=-(\mathfrak d^\ast/\mathfrak d)\,d\mathfrak d$. (ii) $\eta^{-1}=\eta$, and $\mathfrak d(\one)=K(\one)$. (iii) Since $\det\mathbf h\equiv1$, $\operatorname{tr}(\mathbf h^{-1}d\mathbf h)=d\log\det\mathbf h=0$; where $\mathbf h=-I$ this is $-\operatorname{tr}d\mathbf h=-d\mathfrak d=0$, and the same holds for $\mathfrak d^\ast$ where $\mathbf h(\alpha^{-1})=-I$. On the torsion stratum $\mathfrak d=\mathfrak d^\ast=-4$, so it lies in $X_F$, $d\log f=0$ there and $f=16^{-\sigma}$. Each of $\mathbf h(\alpha)=-I$ and $\mathbf h(\alpha^{-1})=-I$ is three equations, $\mathbf h_{12}=\mathbf h_{21}=0$ and $\mathbf h_{11}=-1$, the fourth following from $\det\mathbf h=1$. They do not change under scaling, $\mathbf h(c\alpha)$ being $\mathbf h(\alpha)$ conjugated by $\operatorname{diag}(c^{1/2},c^{-1/2})$; on the chart $\alpha_{e_0}=1$, a torus of dimension $2p-1$, every component of the zero set of six functions has dimension at least $2p-7$.
\end{proof}

\begin{hproposition}[\hb{prop:wheelclasses} (The classes of a wheel)]
As stated.
\end{hproposition}

\begin{proof}
The number of spanning trees that contain $s_1$ is the derivative of $K$ in $\alpha_{s_1}$ at $\one$. With $S$, $R$ and $M$ as in the proof of Proposition~\hb{prop:wheelholonomy}, $K=\operatorname{tr}M-2\prod_i\alpha_{r_i}$, and the derivative of $S$ is $\begin{bmatrix}0&1\\0&0\end{bmatrix}$, so that number is the entry $(2,1)$ of $R(1)M_1^{\,p-1}$, $M_1=S(1)R(1)=\begin{bmatrix}2&1\\1&1\end{bmatrix}$. With the Fibonacci numbers $\mathrm{Fib}_m$, $M_1=\begin{bmatrix}1&1\\1&0\end{bmatrix}^2$ and $\begin{bmatrix}1&1\\1&0\end{bmatrix}^m=\begin{bmatrix}\mathrm{Fib}_{m+1}&\mathrm{Fib}_m\\\mathrm{Fib}_m&\mathrm{Fib}_{m-1}\end{bmatrix}$; the second row of $R(1)$ is $(1,1)$, so the entry is $\mathrm{Fib}_{2p-1}+\mathrm{Fib}_{2p-2}=\mathrm{Fib}_{2p}=(\lambda_+^p-\lambda_-^p)/\sqrt5$. Dividing by $K(\one)$ gives $R_{\mathrm s}$. With the Lucas numbers $\mathrm{Luc}_m=\mathrm{Fib}_{m-1}+\mathrm{Fib}_{m+1}$, $\mathrm{Luc}_m-2\mathrm{Fib}_m=\mathrm{Fib}_{m-3}$, so $2R_{\mathrm s}=1$ reads $\mathrm{Fib}_{2p-3}=2$, that is $2p-3=3$ and $p=3$. The rotation $i\mapsto i+1$ carries the list, so all spokes have one occupation and all rims one, and Theorem~\hb{thm:reversed}(i) with $x$ gives $R_{\mathrm r}=1-R_{\mathrm s}$; \hb{prop:massesgauge} gives the masses. By \hb{cor:certifyone} the classes at the centre are the orbits of the bijections that carry or exchange the list, and the rotations with $x$ make one class. At a lean they are the orbits of the bijections that carry the list. For $p\ge4$ the circuits of three marks are the triangles $\{s_i,s_{i+1},r_i\}$, so a spoke lies in two of them and a rim in one, and a bijection that carries the list, carrying circuits to circuits, maps spokes to spokes: the classes at a lean are the spokes and the rims, the rotations acting transitively on each. By Theorem~\hb{thm:reversed}(i), $\bar\phi_x$ keeps the integrand at every lean, and $x$ maps the spokes onto the rims. For $p=3$ the permutations of the four vertices of $K_4$ carry the list and act transitively on its six edges.
\end{proof}


\section{Chapter~\hb{ch:strideslaw}: the stride's law and the gauge end}\label{sec:chstrideslaw}

In this section the chart is $\alpha_{e_0}=1$ with the coordinates $y=\log\alpha\in\mathbb R^{\cE\setminus e_0}$, $f^\star=f(\alpha^\star)$, and a differential operator $A$ with polynomial coefficients on $(0,\infty)$ has the formal adjoint $A^\dagger$, with $\ip{AT}\phi=\ip T{A^\dagger\phi}$ for distributions $T$ and test functions $\phi$; for $\theta=t\frac d{dt}$, $\theta^\dagger=-1-\theta$, so that $p(-1-i-\theta)^\dagger=p(\theta-i)$ and
\[
\mathfrak D_{\mathbf p}^\dagger=\sum_ip_i(\theta-i)\,t^i,\qquad\mathfrak D_{\mathbf p}^\dagger\,t^k=\sum_ip_i(k)\,t^{k+i}.
\]
Five lemmas carry the proofs: a moment lemma, which turns recurrences of moments into equations for measures; the volumes of the sublevel sets of a Morse function; the uniqueness of expansions in powers and logarithms; the formal solutions of a recurrence at a simple root of its symbol; and the point masses that a cut germ produces.

\begin{slemma}\label{lem:sl-moment}
Let $\nu_1,\dots,\nu_r$ be complex measures on $(0,R]$ with $\int t^{K_0}\,d|\nu_l|<\infty$ for a real $K_0$, let $A_1,\dots,A_r$ be differential operators $A_l=\sum_it^{a_{l,i}}q_{l,i}(\theta)$ with $a_{l,i}\in\mathbb N$ and polynomials $q_{l,i}$, and let $E=\sum_lA_l\nu_l$, a distribution on $(0,\infty)$. For $\re k\ge K_0$ the numbers $M_E(k)=\sum_l\int A_l^\dagger(t^k)\,d\nu_l=\sum_{l,i}q_{l,i}(-1-k-a_{l,i})\int t^{k+a_{l,i}}\,d\nu_l$ are defined, holomorphic for $\re k>K_0$.
\begin{enumerate}[label=(\roman*)]
\item If $M_E(k)=0$ for all large integers $k$, then $E=0$.
\item If $E=0$, then $M_E(k)=0$ for every $k$ with $\re k\ge K_0$.
\end{enumerate}
\end{slemma}

\begin{proof}
$A_l^\dagger=\sum_iq_{l,i}(-1-\theta)\,t^{a_{l,i}}$, and $\theta\,t^m=m\,t^m$ gives the formula for $M_E$; the integrals converge absolutely for $\re k\ge K_0$, $t$ being at most $R$, and are holomorphic in $k$ for $\re k>K_0$ by dominated convergence on compact sets.

(i) For $\psi\in C^\infty(\mathbb R)$ and an integer $K\ge K_0$ put $\Lambda_K(\psi)=\sum_l\int A_l^\dagger(t^K\psi)\,d\nu_l$. Since $\theta(t^m\psi)=t^m(m+\theta)\psi$ and $\theta^c=\sum_{b\le c}S(c,b)\,t^b\partial^b$ with integers $S(c,b)$, $A_l^\dagger(t^K\psi)$ is a finite combination of the functions $t^{K+a+b}\psi^{(b)}$, $b$ at most the largest order $d$ of the $A_l$, with coefficients that depend on $K$; so $|\Lambda_K(\psi)|\le C_K\max_{b\le d}\sup_{[0,R]}|\psi^{(b)}|$. For $\psi=t^m$, $\Lambda_K(t^m)=M_E(K+m)$, which vanishes when $K$ is large. Polynomials are dense in $C^d([0,R])$ (the Weierstrass approximation theorem applied to $\psi^{(d)}$, followed by $d$ integrations from $0$ with the initial values $\psi^{(b)}(0)$), so $\Lambda_K(\psi)=0$ for every $\psi\in C^\infty(\mathbb R)$ and every large $K$. A test function $\phi$ on $(0,\infty)$ is $t^K\psi$ with $\psi=t^{-K}\phi$, smooth with compact support in $(0,\infty)$, and $\ip E\phi=\sum_l\int A_l^\dagger\phi\,d\nu_l=\Lambda_K(\psi)=0$.

(ii) Let $\chi\in C^\infty(\mathbb R)$ vanish on $(-\infty,1]$ and equal $1$ on $[2,\infty)$, put $\chi_\epsilon(t)=\chi(t/\epsilon)$, and let $\tilde\chi$ be a test function equal to $1$ on $[\epsilon,R]$. Then $0=\ip E{t^k\chi_\epsilon\tilde\chi}=\sum_l\int A_l^\dagger(t^k\chi_\epsilon)\,d\nu_l$, the measures being carried by $(0,R]$. The difference $A_l^\dagger(t^k\chi_\epsilon)-\chi_\epsilon A_l^\dagger(t^k)$ involves only the functions $\theta^c\chi_\epsilon(t)=(\theta^c\chi)(t/\epsilon)$ with $c\ge1$, which are bounded uniformly in $\epsilon$ and vanish off $[\epsilon,2\epsilon]$; so it is at most $C\,t^{\re k}$ on $[\epsilon,2\epsilon]$ and $0$ elsewhere, and its integral tends to $0$. By dominated convergence $\int\chi_\epsilon A_l^\dagger(t^k)\,d\nu_l\to\int A_l^\dagger(t^k)\,d\nu_l$, so $M_E(k)=0$.
\end{proof}

\begin{slemma}\label{lem:sl-morse}
Let $F$ be real-analytic near $0\in\mathbb R^{d}$, $d\ge1$, with $F(0)=0$, $\nabla F(0)=0$ and $\nabla^2F(0)=H$ positive definite, and let $w$ be real-analytic near $0$.
\begin{enumerate}[label=(\roman*)]
\item There are $\epsilon>0$ and a function $A$ holomorphic at $0$ with $\int_{\{F\le\sigma\}\cap B_\epsilon}w\,dy=\sigma^{\frac d2}A(\sigma)$ for small $\sigma>0$, $B_\epsilon$ the ball of radius $\epsilon$, and $A(0)=w(0)\,(2\pi)^{\frac d2}\big/\big(\Gamma(\frac d2+1)\sqrt{\det H}\big)$.
\item Let $F$ be even and $w=1$. Then $F$ extends holomorphically to a neighbourhood of $0$ in $\mathbb C^{d}$, $F_{\mathbb T}(\vartheta)=-F(i\vartheta)$ is real-analytic and even for real $\vartheta$ near $0$, with the Hessian $H$ at $0$, and with the same $A$, $\int_{\{F_{\mathbb T}\le\sigma\}\cap B_\epsilon}d\vartheta=\sigma^{\frac d2}A(-\sigma)$ for small $\sigma>0$.
\end{enumerate}
\end{slemma}

\begin{proof}
(i) Write $F(y)=\frac12y^{\mathsf T}Q(y)y$ with $Q(y)=2\int_0^1(1-\tau)\nabla^2F(\tau y)\,d\tau$, symmetric, real-analytic, $Q(0)=H$; near $0$, $Q(y)$ has an analytic positive square root, and $v(y)=Q(y)^{1/2}y$ is an analytic change of coordinates with $F=\frac12|v|^2$ and $\partial v/\partial y(0)=H^{1/2}$ (the Morse lemma). With $J(v)=w(y(v))\,|\det\partial y/\partial v|=\sum_\beta J_\beta v^\beta$, convergent on a ball, $J(0)=w(0)(\det H)^{-1/2}$, and for $\epsilon$ small the set $\{F\le\sigma\}\cap B_\epsilon$ is the preimage of the ball $|v|^2\le2\sigma$ for small $\sigma$. The monomials of odd degree integrate to $0$ over the ball, and $\int_{|v|^2\le2\sigma}v^\beta dv=c_\beta(2\sigma)^{\frac{d+|\beta|}2}$, so the integral is $\sigma^{\frac d2}\sum_m\big(\sum_{|\beta|=2m}J_\beta c_\beta2^{\frac d2+m}\big)\sigma^m$, with the series convergent for small $\sigma$; at $m=0$, $c_0$ is the volume $\pi^{\frac d2}/\Gamma(\frac d2+1)$ of the unit ball.

(ii) $F$ is even, so $\nabla^2F$ and $Q$ are even and $v$ is odd, with real Taylor coefficients in the monomials of odd degree; its inverse $y(v)$ is odd too, and $J$ is even. A function analytic near $0$ extends holomorphically to a polydisc, and $F(i\vartheta)$ is real for real $\vartheta$ since only even monomials occur, $i^{2m}=(-1)^m$. Put $v_{\mathbb T}(\vartheta)=-i\,v(i\vartheta)$, real for real $\vartheta$ since only odd monomials occur, $-i\cdot i^{2m+1}=(-1)^m$. Then $\frac12|v_{\mathbb T}|^2=\frac12(-i)^2\,v(i\vartheta)\cdot v(i\vartheta)=-F(i\vartheta)=F_{\mathbb T}(\vartheta)$, and $v_{\mathbb T}$ is an analytic change of coordinates with the inverse $\vartheta=-i\,y(i\,v_{\mathbb T})$, whose Jacobian matrix is $(\partial y/\partial v)(i\,v_{\mathbb T})$; so $d\vartheta=J(i\,v_{\mathbb T})\,dv_{\mathbb T}$ with $J(iv)=\sum_\beta J_\beta i^{|\beta|}v^\beta=\sum_\beta(-1)^{|\beta|/2}J_\beta v^\beta$, $|\beta|$ even. The computation of (i) gives $\sigma^{\frac d2}\sum_m(-1)^m\big(\sum_{|\beta|=2m}J_\beta c_\beta2^{\frac d2+m}\big)\sigma^m=\sigma^{\frac d2}A(-\sigma)$. The Hessian of $F_{\mathbb T}$ at $0$ is $-i^2H=H$.
\end{proof}

\begin{slemma}\label{lem:sl-asymp}
\begin{enumerate}[label=(\roman*)]
\item Let $(z_c,\tau_c)$, $c=1,\dots,m$, be distinct pairs with $|z_c|=1$ and $\tau_c\in\mathbb R$, and $a_c\in\mathbb C$. If $\sum_ca_cz_c^k\,k^{i\tau_c}\to0$ as $k\to\infty$ through the integers, every $a_c$ is $0$.
\item Let $F(k)=\sum a_{z,\mu,j}\,z^kk^{-\mu}(\log k)^j$ be a finite sum over distinct triples with $|z|=1$, $\mu\in\mathbb C$ and $j\in\mathbb N$. If $F(k)=O(k^{-N})$ along the integers for some $N>\max\re\mu$, every $a_{z,\mu,j}$ is $0$.
\end{enumerate}
\end{slemma}

\begin{proof}
(i) Let $x_k=\sum_ca_cz_c^kk^{i\tau_c}$ and take the logarithmic means $W_M=\frac1{\log M}\sum_{k=1}^M|x_k|^2/k$, which tend to $0$ when $x_k\to0$. In $|x_k|^2$ the term of $(c,c')$ is $a_c\bar a_{c'}w^kk^{i\Delta}$ with $w=z_c\bar z_{c'}$ and $\Delta=\tau_c-\tau_{c'}$. For $c=c'$ it contributes $|a_c|^2\frac1{\log M}\sum_{k\le M}\frac1k\to|a_c|^2$. If $w\neq1$, the partial sums of $w^k$ are bounded by $2/|1-w|$ and $k\mapsto k^{i\Delta-1}$ has bounded total variation, $|(k+1)^{i\Delta-1}-k^{i\Delta-1}|\le(1+|\Delta|)k^{-2}$; summation by parts bounds $\sum_{k\le M}w^kk^{i\Delta-1}$ uniformly in $M$. If $w=1$, then $\Delta\neq0$ and $\sum_{k\le M}k^{i\Delta-1}=\int_1^Mx^{i\Delta-1}dx+O(1)=(M^{i\Delta}-1)/(i\Delta)+O(1)$ is bounded. Divided by $\log M$ these tend to $0$, so $W_M\to\sum_c|a_c|^2$, and every $a_c$ vanishes.

(ii) Let $\sigma_0$ be the least $\re\mu$ among the triples with $a\neq0$ and $j_0$ the largest $j$ among those with $\re\mu=\sigma_0$. Then $k^{\sigma_0}(\log k)^{-j_0}F(k)$ is $\sum a_{z,\mu,j_0}z^kk^{-i\im\mu}$, over $\re\mu=\sigma_0$, plus terms that tend to $0$, and it tends to $0$ itself since $N>\sigma_0$; the pairs $(z,-\im\mu)$ are distinct, so (i) makes these $a$ vanish, against the choice of $\sigma_0$ and $j_0$.
\end{proof}

\begin{slemma}\label{lem:sl-formal}
Let $\mathbf p=(p_0,\dots,p_h)$ be as in Definition~\hb{def:sl-densityop}, $D\ge1$, and $\lambda\neq0$ a simple root of $\chi$. Consider formal sums $\mathcal S=\sum a_{\mu,j}k^\mu(\log k)^j$ over $j\le J$ and $\mu$ in finitely many classes modulo $\mathbb Z$ with $\re\mu$ bounded above, with the operator $L=\sum_ip_i(k)\lambda^i\varsigma^i$, where $\varsigma^i\big(k^\mu(\log k)^j\big)=k^\mu(1+\frac ik)^\mu\big(\log k+\log(1+\frac ik)\big)^j$ is expanded in powers of $k^{-1}$. If $L\mathcal S=0$, then $\mathcal S=a\,k^{\rho}\big(1+\sum_{m\ge1}e_mk^{-m}\big)$ with $\rho=-\chi_1(\lambda)/(\lambda\chi'(\lambda))$, $a\in\mathbb C$, and $e_m$ fixed by $\mathbf p$ and $\lambda$; and this series is a formal solution.
\end{slemma}

\begin{proof}
Each coefficient of $L\mathcal S$ is a finite sum, since $L$ raises the exponent by at most $D$ and $\varsigma$ does not raise it. With $p_i=c_ik^D+c'_ik^{D-1}+O(k^{D-2})$ and $\varsigma^i(k^\mu(\log k)^j)=k^\mu(\log k)^j+i\mu\,k^{\mu-1}(\log k)^j+ij\,k^{\mu-1}(\log k)^{j-1}+\cdots$, the dots of lower exponent,
\begin{multline*}
L\big(k^\mu(\log k)^j\big)=\chi(\lambda)k^{\mu+D}(\log k)^j+\big(\mu\lambda\chi'(\lambda)+\chi_1(\lambda)\big)k^{\mu+D-1}(\log k)^j\\
+j\lambda\chi'(\lambda)k^{\mu+D-1}(\log k)^{j-1}+\cdots,
\end{multline*}
using $\sum_ic_ii\lambda^i=\lambda\chi'(\lambda)$, and $\chi(\lambda)=0$. Let $\sigma$ be the largest $\re\mu$ in $\mathcal S$ and, for $\mu$ with $\re\mu=\sigma$, $j_\mu$ the largest $j$ with $a_{\mu,j}\neq0$. No term of $\mathcal S$ has an exponent $\mu+r$, $r\ge1$, so the coefficient of $k^{\mu+D-1}(\log k)^{j_\mu}$ in $L\mathcal S$ is $a_{\mu,j_\mu}(\mu\lambda\chi'(\lambda)+\chi_1(\lambda))$, and $\mu=\rho$. The coefficient of $k^{\rho+D-1}(\log k)^{j_\rho-1}$ is then $j_\rho\lambda\chi'(\lambda)a_{\rho,j_\rho}$, so $j_\rho=0$: the terms of $\mathcal S$ of largest real exponent are $a\,k^\rho$. For the series $k^\rho\sum_me_mk^{-m}$, $e_0=1$, the coefficient of $k^{\rho-M+D-1}$ in $L$ of it is $-M\lambda\chi'(\lambda)e_M$ plus a combination of $e_0,\dots,e_{M-1}$, which fixes $e_M$; so it is a formal solution, and $\mathcal S-a\,k^\rho\sum_me_mk^{-m}$ is a formal solution whose terms all have $\re\mu<\re\rho$, hence $0$ by what precedes.
\end{proof}

\begin{slemma}\label{lem:sl-pointmass}
Let $\mathcal D=\sum_{j=0}^Db_j(s)\,\partial_s^j$ have polynomial coefficients with $b_{D,0}=0$, and let $u=\sum_ru_rs^r$ be holomorphic near $s=0$. Then, with $\one_+=\one_{s>0}$ and $\mathsf p_r$ as in Theorem~\hb{thm:gaugeend}(iv),
\[
\mathcal D(u\,\one_+)=(\mathcal Du)\,\one_++\sum_{r=0}^{D-2}\mathsf p_r(u)\,\delta^{(D-2-r)},
\]
and the coefficient of $u_r$ in $\mathsf p_r(u)$ is $r!\,\big(b_{D-1,0}-(D-1-r)\,b_{D,1}\big)$.
\end{slemma}

\begin{proof}
For $i<j$, $\partial^j(s^i\one_+)=i!\,\delta^{(j-1-i)}$, and the Taylor remainder of order $j$ of $u$ times $\one_+$ has the classical $j$-th derivative; so $\partial^j(u\one_+)=(\partial^ju)\one_++\sum_{i<j}i!\,u_i\,\delta^{(j-1-i)}$. Further $s^l\delta^{(m)}=(-1)^l[m]_l\,\delta^{(m-l)}$ for $l\le m$ and $0$ for $l>m$, with $[m]_l=\binom ml\,l!$. So the term $b_{j,l}s^l\partial^j$ sends $u_i\,s^i$, $i<j$, to $(-1)^l\binom{j-1-i}l\,i!\,l!\,b_{j,l}\,u_i\,\delta^{(j-1-i-l)}$. The level $r=D-2-(j-1-i-l)$ has $l=j+r+1-D-i$, and $l\le j-1-i$ exactly when $r\le D-2$; $l\ge0$ and $j\le D$ give $i\le r+1$, with $i=r+1$ only for $j=D$, $l=0$, excluded by $b_{D,0}=0$. At $i=r$ the terms are $j=D-1$, $l=0$, giving $r!\,b_{D-1,0}$, and $j=D$, $l=1$, giving $-(D-1-r)\,r!\,b_{D,1}$.
\end{proof}

\begin{hdefinition}[\hb{def:strideslaw} (The stride's law)]
As stated.
\end{hdefinition}

\begin{proof}
The density of the origin form, $\alpha^{\mvec_0}\prod_jK_j^{-s_{0j}}\Psi_j^{-t_{0j}}$ with real data, is positive on the positive orthant, the letters being positive there; so $\omega_0$ is a positive measure on the chart and its image under the measurable map $y\mapsto f(e^y)$ is a measure on $(0,\infty)$.
\end{proof}

\begin{hproposition}[\hb{prop:sl-moments} (A ray is one Trailmoment on a line and one Gramma)]
As stated.
\end{hproposition}

\begin{proof}
(i) By Theorem~\hb{thm:pencil}(i), $0<f\le f^\star$ on the chart and $f\le C'e^{-c'\|y\|}$, so $\{f\ge t\}$ is compact for $t>0$ and carries a finite mass of the continuous density $\omega_0/dy$; $\mu_f$ is carried by $(0,f^\star]$. For the image measure $\int t^k\mu_f(dt)=\int f^k\omega_0=\langle1\rangle_k$, the integrals converging absolutely together. With the kernel $\sZ_{k;(1)}(x)=x^{-k}$ (Definition~\hb{def:lattice}) the Trailmoment of the input $\mu_f$ (Definition~\hb{def:inputs}) with the coefficient function $t^{-1}$ over $(0,\infty)$ is $\int t^k\mu_f(dt)$, and Theorem~\hb{thm:transportfree} with $\tau=\iota$ writes it as $\Tx{k}{\iota_\ast\mu_f}{0}$.

(ii) If $0$ were not an interior point of the cone of $\mathbf v$, a nonzero $y$ would have $\ip{\mathbf v_i}y\le0$ for every $i$ (the separation of a point from a closed convex cone), and then $f^{-1}(e^{\lambda y})\le\sum_iw_i$ for $\lambda>0$, against $f\le C'e^{-c'\|y\|}$; so the cone is $\mathbb R^{\cE\setminus e_0}$, and Theorem~\hb{thm:grammadomain} makes the Gramma entire. For $X>0$ and $\re k>0$, $\Gamma(k)X^{-k}=\int_0^\infty u^{k-1}e^{-uX}du$. At $X=f^{-1}$, integrated against $\omega_0=\alpha^{\mvec_0}dy$, Fubini's theorem applies, the double integral of the modulus being $\Gamma(\re k)\langle1\rangle_{\re k}$ at real data, and the inner integral is $\int_{\mathbb R^{\cE\setminus e_0}_{>0}}x^{\mvec_0|_{\cE\setminus e_0}-\one}\exp\big(-\sum_iuw_ix^{\mathbf v_i}\big)dx=\sG_{\mvec_0|_{\cE\setminus e_0};\mathbf v}(u\mathbf w)$ (Definition~\hb{def:gramma}), $\alpha^{\mvec_0}$ being $x^{\mvec_0|_{\cE\setminus e_0}}$ on the chart. This is Theorem~\hb{thm:laplaceside} for the Trailmoment of (i), whose Laplace transform beyond the cut $0$ is $\int e^{-ux}\,\iota_\ast\mu_f(dx)=\int e^{-uf^{-1}}\omega_0$.

(iii) The shape law has the density $u^{\mvec-\one}\prod_jK_j(u)^{-s_j}\Psi_j(u)^{-t_j}$ against $d\omega_{\cE}$, divided by the period (Definition~\hb{def:shapelaw}), and the change of variables of Theorem~\hb{thm:projective} that carries the period to the chart, $\langle1\rangle_k=\int f^k\omega_0$ (Lemma~\hb{lem:ch35-direction}(iv)), carries the expectation of every bounded measurable function $\varphi$ of degree $0$ to $\int\varphi\,f^k\omega_0/\langle1\rangle_k$. At $\varphi=\psi\circ f$ this is $\int\psi(t)\,t^k\mu_f(dt)/\langle1\rangle_k$. So $\log\langle1\rangle_k=\log\int t^k\mu_f(dt)$ is the log-partition of the family $t^k\mu_f(dt)/\langle1\rangle_k$; it is real-analytic in $k$, and differentiation under the integral gives the first derivative $\oE[\log f(U)]$ and the second $\oE[\log^2f(U)]-\oE[\log f(U)]^2$, which is positive because $\log f(U)$ is not almost surely constant: $U$ has a positive continuous density on the chart, and $f$ is continuous there and not constant, $f(\alpha^\star)=f^\star>0$ while $f\le C'e^{-c'\|y\|}$ (Theorem~\hb{thm:pencil}(i)).
\end{proof}

\begin{hproposition}[\hb{prop:sl-ends} (The two ends of the stride's law)]
As stated.
\end{hproposition}

\begin{proof}
(i) At $y=2z$ the stride is $e^{X_1(z)}$, $X_1$ the exponent of the question of the direction (proof of Lemma~\hb{lem:ch35-direction}(ii)), and $-X_1$ is strictly convex (Theorem~\hb{thm:severity}(ii)); so $\log f$ is strictly concave on the chart, its only critical point is its maximum $\log\alpha^\star$, and $\nabla f\neq0$ elsewhere. For $0<t<f^\star$ the level set $\Sigma_t=\{f=t\}$ is then a compact real-analytic hypersurface, compact because $f\to0$ at infinity of the chart, and not empty, the chart being connected. Near $t_1\in(0,f^\star)$ the flow $\Phi$ of the real-analytic field $\nabla f/|\nabla f|^2$, along which $f$ grows at the unit rate, maps $\Sigma_{t_1}\times(t_1-\epsilon,t_1+\epsilon)$ real-analytically onto a neighbourhood of $\Sigma_{t_1}$, with $f(\Phi(x,\tau))=\tau$, and $dy=J(x,\tau)\,d\sigma(x)\,d\tau$ with $J$ real-analytic and positive. So $\mu_f$ has the density $\varrho_f(\tau)=\int_{\Sigma_{t_1}}(\omega_0/dy)(\Phi(x,\tau))\,J(x,\tau)\,d\sigma(x)$ there, real-analytic and positive; this is the coarea formula.

(ii) Apply Lemma~\ref{lem:sl-morse}(i) to $F(y)=\log f^\star-\log f(e^y)$ at $\log\alpha^\star$, with the Hessian $H$, positive definite (Lemma~\hb{lem:ch35-direction}(iii)), and $w=\omega_0/dy$. Since $f<f^\star-\delta'$ off a neighbourhood of $\log\alpha^\star$, $\mu_f([t,f^\star])=\int_{\{F\le L\}\cap B_\epsilon}w\,dy=L^{\frac{n-1}2}A(L)$ for $t$ near $f^\star$, with $L=\log(f^\star/t)$. Differentiating, with $dL/dt=-1/t$,
\[
\varrho_f(t)=-\frac d{dt}\mu_f([t,f^\star])=\frac1t\,L^{\frac{n-3}2}B(L),\qquad B(L)=\tfrac{n-1}2A(L)+LA'(L).
\]
Put $\lambda(t)=L/(f^\star-t)$, holomorphic and positive near $f^\star$ with $\lambda(f^\star)=1/f^\star$. Then $\varrho_f(t)=(f^\star-t)^{\frac{n-3}2}g(f^\star-t)$ with $g(f^\star-t)=t^{-1}\lambda(t)^{\frac{n-3}2}B(L)$, holomorphic at $t=f^\star$, and $g(0)=(f^\star)^{-\frac{n-1}2}\frac{n-1}2A(0)$; with $A(0)$ from Lemma~\ref{lem:sl-morse}(i) at $d=n-1$ and $\frac{n-1}2/\Gamma(\frac{n+1}2)=1/\Gamma(\frac{n-1}2)$ this is the stated value.

(iii) By Proposition~\hb{prop:evenfloor}(i) for the question of the direction, $2\kappa(z)+\sum_j(s_{1j}+t_{1j})\log\theta_{j,\min}\le-X_1(z)\le2\kappa(z)+\sum_j(s_{1j}+t_{1j})\log K_j(\one)$, and $2\kappa(y/2)=\kappa(y)$, $\kappa$ being positively homogeneous of degree one; with $\log f(e^y)=X_1(y/2)$ this is the squeeze. On the sector of an order $\pi$ (Definition~\hb{def:tiltwall}) the differences $u_l=y_{\pi(l-1)}-y_{\pi(l)}\ge0$, $2\le l\le n$, are coordinates in which $dy$ is the Lebesgue measure and $\kappa(y)=\sum_l\kappa_{1,U_l}u_l$, so $\int_{\text{sector}}e^{-\kappa}dy=\prod_l\kappa_{1,U_l}^{-1}$; the sectors cover the chart and meet in null sets, so $\int e^{-\kappa}dy=V_1$. By homogeneity $\operatorname{vol}\{\kappa\le\sigma\}=\sigma^{n-1}\operatorname{vol}\{\kappa\le1\}$ and $V_1=\int_0^\infty e^{-\sigma}\,d\operatorname{vol}\{\kappa\le\sigma\}=(n-1)!\operatorname{vol}\{\kappa\le1\}$, as in Theorem~\hb{thm:severity}(i). For $t\le C_-$ the squeeze gives $\{\kappa\le\log(C_-/t)\}\subseteq\{f\ge t\}\subseteq\{\kappa\le\log(C_+/t)\}$, and $\mu_f([t,f^\star])=\operatorname{vol}\{f\ge t\}$ for the origin $0$. Last, $\int e^{-k\kappa}dy=k^{1-n}V_1$ for $k>0$, by homogeneity, and $C_-^ke^{-k\kappa}\le f^k\le C_+^ke^{-k\kappa}$.
\end{proof}

\begin{hdefinition}[\hb{def:sl-densityop} (The density operator of a recurrence)]
As stated.
\end{hdefinition}

\begin{proof}
For a polynomial $p$ of degree at most $D$, $p(-1-i-\theta)=c(-1)^D\theta^D+\cdots$ with $c=[k^D]p$, and $\theta^D=t^D\partial^D+\cdots$; so the coefficient of $\partial^D$ in $t^ip_i(-1-i-\theta)$ is $(-1)^Dc_it^{D+i}$, and $a_D=(-1)^Dt^D\chi(t)$. Extension by zero involves only $u$ near $\lambda$: two intervals give the same condition on their common part, $\mathfrak D_{\mathbf p}$ being a differential operator.
\end{proof}

\begin{hproposition}[\hb{prop:sl-density} (A recurrence of the values is an equation for the density)]
As stated.
\end{hproposition}

\begin{proof}
Let $k_\ast$ be the bound beyond which $\langle1\rangle_k$ converges absolutely, $\re k>k_\ast$ (Lemma~\hb{lem:ch35-direction}(iv)), and $K_0>k_\ast$. Apply Lemma~\ref{lem:sl-moment} to $\nu_1=\mu_f$, carried by $(0,f^\star]$ with $\int t^{K_0}d\mu_f=\langle1\rangle_{K_0}<\infty$, and $A_1=\mathfrak D_{\mathbf p}$, for which $A_1^\dagger t^k=\sum_ip_i(k)t^{k+i}$ and $M_E(k)=\sum_ip_i(k)\langle1\rangle_{k+i}$, $E=\mathfrak D_{\mathbf p}\mu_f$. (a) gives (c) by Lemma~\ref{lem:sl-moment}(i), and (c) gives the vanishing for $\re k\ge K_0$ by (ii), for every $K_0>k_\ast$. By Proposition~\hb{prop:rays-strength}(ii) the values continue meromorphically to $\mathbb C$, so the identity holds as one of meromorphic functions, and at every $k$ where its terms converge; (b) gives (a).

On $(0,f^\star)$, $\mu_f=\varrho_f\,dt$ with $\varrho_f$ real-analytic (Proposition~\hb{prop:sl-ends}(i)), so $E$ restricted there is $(\mathfrak D_{\mathbf p}\varrho_f)\,dt$, and (c) gives $\mathfrak D_{\mathbf p}\varrho_f=0$. On $(f^\star-\delta,f^\star+\delta)$, $\mu_f=\varrho_f\one_{(f^\star-\delta,f^\star)}dt$, the point $f^\star$ having no mass, and $\varrho_f$ is integrable there by Proposition~\hb{prop:sl-ends}(ii), $\frac{n-3}2>-1$; so the germ extends by zero. For the origin $0$, the proof of Corollary~\hb{cor:rays-lowend}(i) and~(ii) uses of $p_0,\dots,p_h$ only that the identity holds as one of meromorphic functions, with Proposition~\hb{prop:rays-lowpoles}(ii) at $k=-j$ and at $k=0$, $\mathsf L_0=V_1\neq0$ (Definition~\hb{def:rays-lowend}).
\end{proof}

\begin{htheorem}[\hb{thm:gaugeend} (The gauge end)]
As stated.
\end{htheorem}

\begin{proof}
(i) If $D=0$, every $p_i$ is a constant $c_i$, $\mathfrak D_{\mathbf p}$ is the multiplication by $\chi(t)$, not identically zero, and $\chi\varrho_f=0$ on $(0,f^\star)$ (Proposition~\ref{prop:sl-density}) contradicts $\varrho_f>0$. So $D\ge1$. With $b_j(s)=(-1)^ja_j(f^\star-s)$, $b_D(s)=(f^\star-s)^D\chi(f^\star-s)$, so $b_{D,0}=0$ and $b_{D,1}=-(f^\star)^D\chi'(f^\star)$. Write $p_i=c_ik^D+c'_ik^{D-1}+\cdots$; $p_i(-1-i-\theta)=(-1)^Dc_i\big(\theta^D+D(1+i)\theta^{D-1}\big)+(-1)^{D-1}c'_i\theta^{D-1}+\cdots$, and $\theta^m=t^m\partial^m+\binom m2t^{m-1}\partial^{m-1}+\cdots$; so
\begin{gather*}
a_{D-1}=(-1)^Dt^{D-1}\Big(\big(\tbinom D2+D\big)\chi(t)+D\,t\,\chi'(t)-\chi_1(t)\Big),\\
b_{D-1,0}=-(f^\star)^{D-1}\big(Df^\star\chi'(f^\star)-\chi_1(f^\star)\big),
\end{gather*}
which gives the two forms of $\gamma$. Since $s^j\partial_s^j=[s\partial_s]_j$, $s^{D-1}\mathfrak D_{\mathbf p}=\sum_{j,l}b_{j,l}s^{l+D-1-j}[s\partial_s]_j$; the exponent $m=l+D-1-j$ is nonnegative, as $l\ge1$ for $j=D$, and the terms with $m=0$ are $(j,l)=(D-1,0)$ and $(D,1)$, which give $R_0$. So $s=0$ is a regular singular point, $R_0$ having the degree $D$ of the operator, and its exponents are the roots of $R_0$.

(ii) Let $u$ extend by zero on $(f^\star-\delta,f^\star)$. By Fuchs's theorem on regular singular points, which applies by (i), $u=\sum_{\lambda,j}s^\lambda(\log s)^jh_{\lambda,j}(s)$ near $s=0^+$, a finite sum with $\lambda$ among the exponents, $j\le D$ and $h_{\lambda,j}$ holomorphic at $0$; so $u$ has a convergent expansion $\sum c_{\mu,j}s^\mu(\log s)^j$, $\mu\in\{0,\dots,D-2,\gamma\}+\mathbb N$, and integrability forces $\re\mu>-1$ for its terms: the terms of least real part $\re\mu$ and largest $j$ are $s^{\re\mu}(\log s)^j$ times a trigonometric polynomial in $\log s$, not identically $0$, whose modulus has a positive mean, and this is not integrable at $0$ when $\re\mu\le-1$. Fix $\psi$ smooth, $0$ on $(-\infty,f^\star-\frac23\delta]$ and $1$ on $[f^\star-\frac13\delta,\infty)$, and put $T=\psi\,u\,\one_{(f^\star-\delta,f^\star)}$. On $(f^\star-\delta,f^\star+\delta)$, $\mathfrak D_{\mathbf p}T=\psi\,\mathfrak D_{\mathbf p}(u\one)+\sum_{a\ge1}(\ldots)\psi^{(a)}$, where the first term vanishes and the others are smooth functions carried by $[f^\star-\frac23\delta,f^\star-\frac13\delta]$, where $u$ is real-analytic; so $\mathfrak D_{\mathbf p}T=\mathcal E$ is a smooth function with that support, and, $T$ having compact support in $(0,\infty)$,
\[
\sum_ip_i(k)\,M(k+i)=\int t^k\,\mathcal E(t)\,dt=O\big((f^\star-\tfrac13\delta)^k\big),\qquad M(k)=\int t^kT(t)\,dt,
\]
for real $k\to\infty$. In $t=f^\star e^{-\sigma}$, $M(k)=f^\star(f^\star)^k\int_0^\infty e^{-(k+1)\sigma}\,T(f^\star e^{-\sigma})\,d\sigma$, and $s=f^\star(1-e^{-\sigma})=f^\star\sigma(1+O(\sigma))$; so near $\sigma=0$, $T(f^\star e^{-\sigma})$ is a convergent sum of $\sigma^\mu(\log\sigma)^j$ times holomorphic functions, with the same leading terms as $u$ up to the factors $(f^\star)^\mu$, and the part away from $0$ contributes $O\big((f^\star-\frac13\delta)^k\big)$. Watson's lemma with logarithms, $\int_0^{c}e^{-\kappa\sigma}\sigma^\mu(\log\sigma)^jd\sigma=\partial_\mu^j\big[\Gamma(\mu+1)\kappa^{-\mu-1}\big]+O(e^{-c\kappa/2})$ for $\re\mu>-1$, gives $A(k)=(f^\star)^{-k}M(k)$ an asymptotic expansion $\mathcal S_u=\sum a_{\mu,j}k^{-\mu-1}(\log k)^j$ along the integers, in finitely many classes of exponents and with $\re\mu$ bounded below. Its leading terms are $f^\star(f^\star)^\mu c_{\mu,j_0}\Gamma(\mu+1)(-1)^{j_0}k^{-\mu-1}(\log k)^{j_0}$, over the $\mu$ of least real part and the largest $j_0$ among them, with $\Gamma(\mu+1)\neq0$; by Lemma~\ref{lem:sl-asymp}(ii) they do not cancel, so $\mathcal S_u\neq0$ when $u\neq0$. Applying the recurrence to truncations of the expansion, $\sum_ip_i(k)(f^\star)^iA(k+i)=O(r^k)$, $r<1$, and Lemma~\ref{lem:sl-asymp}(ii), order by order, give $L\mathcal S_u=0$ with $L=\sum_ip_i(k)(f^\star)^i\varsigma^i$, so by Lemma~\ref{lem:sl-formal}, $\mathcal S_u$ is a multiple of one fixed series. The map $u\mapsto\mathcal S_u$ is linear and injective, so the germs that extend by zero form a space of dimension at most one.

(iii) By Proposition~\ref{prop:sl-density} the germ of $\varrho_f$ extends by zero, and it is not zero; so it spans the space of (ii). Put $a=\frac{n-3}2$; by Proposition~\ref{prop:sl-ends}(ii), $\varrho_f=s^ag(s)$ with $g(0)\neq0$, and it solves $\mathfrak D_{\mathbf p}$ on $s>0$. If $a\notin\mathbb Z$, the leading coefficient of $s^{D-1}\mathfrak D_{\mathbf p}(s^ag)=\sum_ms^mR_m(s\partial_s)(s^ag)$ is $R_0(a)g(0)$, so $a$ is a root of $R_0$ other than $0,\dots,D-2$, and $a=\gamma$. If $a\in\mathbb N$, $\varrho_f$ is holomorphic at $s=0$, with $u_r=0$ for $r<a$ and $u_a=g(0)$. When $a\le D-2$, the germ extends by zero, and by Lemma~\ref{lem:sl-pointmass} $\mathsf p_a(\varrho_f)=a!\,b_{D,1}(a-\gamma)\,g(0)=0$, the lower coefficients vanishing; so $\gamma=a$. When $a\ge D-1$, the coefficient of $s^a$ in $s^{D-1}\mathfrak D_{\mathbf p}\varrho_f=0$ is $R_0(a)g(0)$, so $a$ is a root of $R_0$ and again $\gamma=a$.

(iv) If $\gamma\notin\{0,\dots,D-2\}$, then $\gamma=a$ is not an integer or is an integer at least $D-1$, and in both cases $R_0(\gamma+N)\neq0$ for $N\ge1$, the roots of $R_0$ being $0,\dots,D-2$ and $\gamma$. Inserting $s^\gamma\sum_Ng_Ns^N$ into $\sum_ms^mR_m(s\partial_s)$ and using $s\partial_s(s^{\gamma+N})=(\gamma+N)s^{\gamma+N}$, the coefficient of $s^{\gamma+N}$ is $\sum_{m\le N}R_m(\gamma+N-m)g_{N-m}$; its vanishing is the stated recurrence, so $g=g(0)\Phi$. If $\gamma\in\{0,\dots,D-2\}$, then by (iii) $\varrho_f$ is holomorphic at $s=0$ with $u_r=0$ for $r<\gamma$ and $u_\gamma=g(0)$; it extends by zero, so $\mathsf p_r(\varrho_f)=0$ for $r\le D-2$ (Lemma~\ref{lem:sl-pointmass}, the distributions $\delta^{(q)}$ being linearly independent), and for $\gamma<r\le D-2$ the coefficient $r!\,b_{D,1}(r-\gamma)$ of $u_r$ is not $0$, so these equations fix $u_r$ from the lower coefficients. For $N\ge D-1$, $R_0(N)\neq0$, and the coefficient of $s^N$ in $s^{D-1}\mathfrak D_{\mathbf p}\varrho_f$ gives the last recurrence.

(v) Each point of $(0,f^\star)$ has a disc in $\mathbb C$ to which $\varrho_f$ extends holomorphically (Proposition~\ref{prop:sl-ends}(i)); two such discs meet in a convex set that meets the real line, so the extensions agree there, and they define a holomorphic function on the union $U$, on which $\mathfrak D_{\mathbf p}\varrho_f=0$ by the identity theorem. The germ at $f^\star$ is that of (iv), and its continuation along a path in $U$ that follows $(0,f^\star)$ is $\varrho_f$; near a root of $\chi$ on the segment, $\varrho_f$ is holomorphic, so the continuations around it on either side agree. For $\re k>k_\ast$, $t^k(\log t)^j$ is dominated on $(0,f^\star]$ by $C\,t^{K}$ for every $K$ with $k_\ast<K<\re k$, locally uniformly in $k$, so the derivatives pass under $\int t^k\mu_f(dt)$.
\end{proof}

\begin{hproposition}[\hb{prop:sl-exponents} (The law knows the dimension of its chart)]
As stated.
\end{hproposition}

\begin{proof}
Let $\alpha_c$ be the native gauge with $f(\alpha_c)=\lambda$ and $A_c$ its solution (Theorem~\hb{thm:gaugesolution}, Definition~\hb{def:raysolution}). The class of $\sum_ip_i(k)f^i$ in $\mathfrak M$ is $\sum_ip_i(k)\operatorname{Sh}^i[1]=0$, so $\sum_ip_i(k)\varsigma^iA_c[1]=0$, and $\lambda^{-k}A_c[1]=k^{-\frac{n-1}2}\big(a_0(1)+\cdots\big)$ with $a_0(1)=(2\pi)^{\frac{n-1}2}(\omega_0/dv)(\alpha_c)\neq0$ is a formal solution of $L=\sum_ip_i(k)\lambda^i\varsigma^i$. By Lemma~\ref{lem:sl-formal} every formal solution is a multiple of $k^{\rho_\lambda}\big(1+\sum_me_mk^{-m}\big)$, and comparing leading exponents, $\rho_\lambda=-\frac{n-1}2$. The computation of Theorem~\ref{thm:gaugeend}(i) uses of $f^\star$ only that it is a simple nonzero root of $\chi$, so at $\lambda$ the point $s=0$ is regular singular with the exponents $0,\dots,D-2$ and $-1+\chi_1(\lambda)/(\lambda\chi'(\lambda))=-1-\rho_\lambda=\frac{n-3}2$.
\end{proof}

\begin{hproposition}[\hb{prop:sl-jet} (The germ at the gauge end is local data at the native gauge)]
As stated.
\end{hproposition}

\begin{proof}
(i) The potential of the realization is $\Phi(y)=\sum_j\big(s_{1j}\log K_j(e^y)+t_{1j}\log\Psi_j(e^y)\big)$ (Definition~\hb{def:potential}) and $\log f(e^y)=\ip{\mvec_1}y-\Phi(y)$, so for $m\ge2$ the Taylor tensor of $\log f$ is $-\partial^m\Phi=-2^{-m}\kappa_\Upsilon$ (Definition~\hb{def:still}), and Proposition~\hb{prop:deviation} at $y_0=\log\alpha^\star$ gives the cumulants.

(ii) Put $\eta(y)=\log f-\log f^\star+\frac12H_{ab}y^ay^b$, of order three. Outside a ball $|y|<\epsilon$ the stride is at most $f^\star e^{-c}$, so that part of $\langle1\rangle_k$ is $O\big((f^\star e^{-c})^{k-1}\langle1\rangle_1\big)$. Inside, with $y=v/\sqrt k$, $\langle1\rangle_k=(f^\star)^kk^{-\frac{n-1}2}\int_{|v|<\epsilon\sqrt k}e^{-\frac12H_{ab}v^av^b}e^{k\eta(v/\sqrt k)}dv$, and $k\eta(v/\sqrt k)=k^{-1/2}h_3(v)+k^{-1}h_4(v)+k^{-3/2}h_5(v)+O(k^{-2}|v|^6)$ with $h_3=\frac16h_{abc}v^av^bv^c$, $h_4=\frac1{24}h_{abcd}v^av^bv^cv^d$ and $h_5$ the quintic term; for $\epsilon$ small $|k\eta(v/\sqrt k)|\le\frac14H_{ab}v^av^b$ on the ball. Expanding the exponential to the order $k^{-3/2}$, with the remainder bounded by $Ck^{-2}(1+|v|)^{12}e^{\frac14Hvv}$, the terms of order $k^{-1/2}$ and $k^{-3/2}$ are odd in $v$ and integrate to $0$ over the ball, and extending the Gaussian integrals to $\mathbb R^{n-1}$ costs an exponentially small amount. So $\langle1\rangle_k=(f^\star)^k(2\pi/k)^{\frac{n-1}2}(\det H)^{-1/2}\big(1+k^{-1}\oE[h_4+\frac12h_3^2]+O(k^{-2})\big)$ for the centred Gaussian with covariance $G$. Its moments are $\oE[v^av^bv^cv^d]=G^{ab}G^{cd}+G^{ac}G^{bd}+G^{ad}G^{bc}$ and, for six factors, the sum over the $15$ pairings; by the symmetry of $h_{abc}$, $9$ of the pairings of $h_{abc}h_{def}$ join one pair inside each factor and $6$ join every index across. So $\oE[h_4]=\frac18h_{abcd}G^{ab}G^{cd}$ and $\frac12\oE[h_3^2]=\frac1{72}\big(9h_{abc}h_{def}G^{ab}G^{cd}G^{ef}+6h_{abc}h_{def}G^{ad}G^{be}G^{cf}\big)$.

(iii) On $(f^\star-\delta,f^\star)$, $\varrho_f=\sum_mg_m(f^\star-t)^{\frac{n-3}2+m}$, and $\int_0^{f^\star-\delta}t^k\mu_f(dt)=O\big((f^\star-\delta)^{k-1}\langle1\rangle_1\big)$. For each $M$, $\int_{f^\star-\delta}^{f^\star}t^k\varrho_f\,dt$ is the sum over $m<M$ of $g_m\int_{f^\star-\delta}^{f^\star}t^k(f^\star-t)^{\frac{n-3}2+m}dt$ plus $O\big(\int_0^{f^\star}t^k(f^\star-t)^{\frac{n-3}2+M}dt\big)$, and $\int_0^{f^\star}t^k(f^\star-t)^{b-1}dt=(f^\star)^{k+b}\mathrm B(k+1,b)$, the part over $(0,f^\star-\delta)$ being exponentially smaller. With $\mathrm B(k+1,b)=\Gamma(b)\Gamma(k+1)/\Gamma(k+1+b)=\Gamma(b)k^{-b}(1+O(k^{-1}))$ the term of $g_m$, $b=\frac{n-1}2+m$, begins at $(f^\star)^{k+b}\Gamma(b)g_mk^{-b}$, and the remainder is $O\big((f^\star)^kk^{-\frac{n-1}2-M}\big)$. So the expansion of $(f^\star)^{-k}k^{\frac{n-1}2}\langle1\rangle_k$ is triangular in the $g_m$, with the nonzero diagonal $(f^\star)^{\frac{n-1}2+m}\Gamma(\frac{n-1}2+m)$.
\end{proof}

\begin{hproposition}[\hb{prop:sl-slopes} (Slopes from the gauge end)]
As stated.
\end{hproposition}

\begin{proof}
(i) Differentiate $j$ times the identity of meromorphic functions $\sum_ip_i(k)\langle1\rangle_{k+i}=0$ (Proposition~\ref{prop:sl-density}) by the Leibniz rule.

(ii) By Theorem~\ref{thm:gaugeend}(v), $\partial_k^j\langle1\rangle_k|_{k=K}=\int_0^{f^\star}t^K(\log t)^j\varrho_f(t)\,dt$, and on $(0,f^\star)$ the series gives $\varrho_f=\sum_mg_m(f^\star-t)^{\frac{n-3}2+m}$, convergent there. Let $K'<K''<K$. On $(0,f^\star]$, $t^K|\log t|^j\le C\,t^{K''}$, so $\sum_m|g_m|\int_0^{f^\star}t^K|\log t|^j(f^\star-t)^{\frac{n-3}2+m}dt\le C\sum_m|g_m|(f^\star)^{K''+\frac{n-1}2+m}\mathrm B\big(K''+1,\frac{n-1}2+m\big)$, which is finite because $\mathrm B(K''+1,b)\le\mathrm B(K'+1,b)$, the Beta function being decreasing in its first argument. So the integral of the sum is the sum of the integrals (Fubini's theorem), and $\int_0^{f^\star}t^a(\log t)^j(f^\star-t)^{b-1}dt=\partial_a^j\big[(f^\star)^{a+b}\mathrm B(a+1,b)\big]$, differentiation under the integral being allowed for $a$ near $K$.
\end{proof}

\begin{hdefinition}[\hb{def:sl-gaugeread} (The procedure that reads a ray from its gauge end)]
As stated.
\end{hdefinition}

\begin{proof}
(1) $\log f$ is strictly concave on the positive part of the chart (proof of Proposition~\ref{prop:sl-ends}(i)), so its critical point is unique, the native gauge of the direction there (Theorem~\hb{thm:pencil}(i)). (2) and (3) are Definition~\ref{def:sl-densityop} and Theorem~\ref{thm:gaugeend}(i), (iii) and (iv); the coefficients $b_{j,l}$ lie in $\Bbbk(f^\star)$, and so do the ratios $u_r/u_\gamma$ and $\Phi_N$. (4) $\epsilon_k=\int_{(0,f^\star-\delta)}t^k\mu_f(dt)$ is nonnegative and, for real $k'\le k$, at most $(f^\star-\delta)^{k-k'}\int t^{k'}\mu_f(dt)$; the bound on $\langle1\rangle_{k'}$ is Proposition~\ref{prop:sl-ends}(iii). On $[f^\star-\delta,f^\star]$, $\varrho_f$ is its germ, by Theorem~\ref{thm:gaugeend}(iv). (5) The recurrence holds at every positive integer (Proposition~\ref{prop:sl-density}), and $\langle1\rangle_k=-\sum_{i\ge1}p_i(k)\langle1\rangle_{k+i}/p_0(k)$ where $p_0(k)\neq0$; the values are moments of $\varrho_f$ by Theorem~\ref{thm:gaugeend}(v). (6) is Proposition~\ref{prop:sl-slopes}, and (7) is Corollary~\hb{cor:rays-lowend}(ii), which holds for $\mathbf p$ by Proposition~\ref{prop:sl-density}.
\end{proof}

\begin{hdefinition}[\hb{def:torusvalue} (The torus value)]
As stated.
\end{hdefinition}

\begin{proof}
On $\mathbb T$, $\alpha^{-1}=\bar\alpha$, and the letters have real coefficients, so $K_j(\alpha^{-1})=\overline{K_j(\alpha)}$ and $f=\prod_j|K_j(\alpha)|^{-2\sigma_j}$; the coefficients being nonnegative, $|K_j(\alpha)|\le K_j(\one)$, and $f\ge\prod_jK_j(\one)^{-2\sigma_j}=f(\one)=f^\star$, the native gauge of the direction being $\one$ (Theorem~\hb{thm:pencil}). The zero set of an active letter has measure zero, the letter being a Laurent polynomial that is not zero, so a real-analytic function on the connected torus that is not identically zero, so $f$ is finite almost everywhere and $\nu^{\mathbb T}_f$ is a probability measure on $[f^\star,\infty)$. For $\re k\le0$, $|f^k|=f^{\re k}\le(f^\star)^{\re k}$, read as $0$ where $f=\infty$ and $\re k<0$, so the integral converges. When no $K_j$ vanishes on $\mathbb T$, $\mathbb T$ is a compact cycle in $X_F$, $\Psi_j(\alpha)=\alpha^{\cE}K_j(\alpha^{-1})$ (Proposition~\hb{prop:identity}) does not vanish either, and on it $dy=i^{n-1}d\vartheta$, the coordinates being $y_e=i\vartheta_e$.
\end{proof}

\begin{hproposition}[\hb{prop:sl-torus} (The torus is the second chain)]
As stated.
\end{hproposition}

\begin{proof}
(i) Equality $|K_j(e^{i\vartheta})|=K_j(\one)$ holds exactly when the numbers $\ip{\one_T}\vartheta$, $T$ the bases of table $j$, agree modulo $2\pi$, all the coefficients being positive. The active tables form a connected family (Definition~\hb{def:centreray}), so one molecule, and for linked marks $e,e'$ some basis $T$ has $e\in T$, $e'\notin T$ and $(T\setminus\{e\})\cup\{e'\}$ a basis (Lemma~\hb{lem:linkexchange}); equality for every active table therefore gives $\vartheta_e\equiv\vartheta_{e'}$ for all marks, and $\alpha=\one$ on the chart. So $f>f^\star$ off $\one$, as in the proof of Theorem~\hb{thm:onekernel}. The function $|K_j(e^{i\vartheta})|^2$ is a trigonometric polynomial, real-analytic and not identically $0$ since its value at $\vartheta=0$ is $K_j(\one)^2$; the zero set of such a function on the connected torus has measure zero. For $\re k\le0$ the integrand is bounded by $(f^\star)^{\re k}$; dominated convergence gives the continuity, and Morera's theorem the holomorphy on $\re k<0$. At $k=-j$, $f^{-j}=\prod_j|K_j|^{2j\sigma_j}$ is the Laurent polynomial $\prod_j(K_j(\alpha)K_j(\alpha^{-1}))^{j\sigma_j}$ restricted to $\mathbb T$, whose mean is its constant term; it is $\mathsf C_j(1)$ when every letter exponent of the direction is at least $1$ (Definition~\hb{def:closings}).

(ii) The hypotheses of Theorem~\hb{thm:onekernel} hold for the centre ray with every $\sigma_j\ge1$, so by Corollary~\hb{cor:valuemodule} $\sum_ip_i(-j)\mathsf C_{j-i}(1)=0$ for all large $j$, that is, by (i), $\sum_ip_i(k)\langle1\rangle^{\mathbb T}_{k+i}=0$ at all large negative integers $k$. With $x=1/t$, $\nu^{\mathbb T}_f$ becomes a probability measure $\tilde\nu$ on $(0,1/f^\star]$ with no mass at $0$, $\theta=-x\frac d{dx}$, and $\mathfrak D_{\mathbf p}^\dagger=\sum_ip_i(\theta-i)\,t^i$ becomes $\tilde A=\sum_ip_i(-x\frac d{dx}-i)\,x^{-i}$. For a test function $\phi(t)=x^{K}\psi(x)$, $\tilde A(x^K\psi)=\sum_ix^{K-i}p_i(-x\frac d{dx}-K)\psi$ is bounded on $(0,1/f^\star]$ by $C_K\max_{b\le D}\sup|\psi^{(b)}|$ for $K\ge h$, and at $\psi=x^m$ its integral against $\tilde\nu$ is $\sum_ip_i(-K-m)\langle1\rangle^{\mathbb T}_{-K-m+i}$, which vanishes for $K$ large. As in the proof of Lemma~\ref{lem:sl-moment}(i), on $[0,1/f^\star]$, the Weierstrass approximation theorem gives $\ip{\mathfrak D_{\mathbf p}\nu^{\mathbb T}_f}\phi=\int\tilde A(x^K\psi)\,d\tilde\nu=0$ for every test function. Conversely, for $\re k\le-h$ the functions $t^{k+i}$ are bounded on $[f^\star,\infty)$, and the cutoff argument of Lemma~\ref{lem:sl-moment}(ii), at $x\to0$ in place of $t\to0$, gives $\sum_ip_i(k)\langle1\rangle^{\mathbb T}_{k+i}=\int\mathfrak D_{\mathbf p}^\dagger(t^k)\,d\nu^{\mathbb T}_f=0$, the cutoff terms being bounded and carried by $x\in[\epsilon,2\epsilon]$, whose mass under $\tilde\nu$ tends to $0$.

(iii) The stride is holomorphic and nonzero near $y=0$ in $\mathbb C^{\cE\setminus e_0}$, and $F(y)=\log f^\star-\log f(e^y)$ is even, $f(\alpha^{-1})=f(\alpha)$ on the centre ray (Theorem~\hb{thm:signgauges}(i)), with the Hessian $H$ at $0$ positive definite (Lemma~\hb{lem:ch35-direction}(iii)). On $\mathbb T$, $\log f(e^{i\vartheta})-\log f^\star=-F(i\vartheta)=F_{\mathbb T}(\vartheta)$, real, and by (i) and compactness $f\ge f^\star+\delta'$ on $\mathbb T$ off a neighbourhood of $\one$; so for $t$ near $f^\star$, $\nu^{\mathbb T}_f([f^\star,t])=(2\pi)^{1-n}\operatorname{vol}\{F_{\mathbb T}\le\log(t/f^\star)\}\cap B_\epsilon$. By Lemma~\ref{lem:sl-morse}, with $L=\log(f^\star/t)$, $\mu_f([t,f^\star])=L^{\frac{n-1}2}A(L)$ for $t<f^\star$ and $\nu^{\mathbb T}_f([f^\star,t])=(2\pi)^{1-n}(-L)^{\frac{n-1}2}A(L)$ for $t>f^\star$, with one $A$, $w=1$. Differentiating, with $B=\frac{n-1}2A+LA'$, $\varrho_f(t)=t^{-1}L^{\frac{n-3}2}B(L)$ and
\[
\varrho^{\mathbb T}_f(t)=\frac d{dt}\Big[(2\pi)^{1-n}(-L)^{\frac{n-1}2}A(L)\Big]=(2\pi)^{1-n}\,t^{-1}(-L)^{\frac{n-3}2}B(L).
\]
With $\lambda(t)=L/(f^\star-t)$, positive near $f^\star$, $L=(f^\star-t)\lambda$ and $-L=(t-f^\star)\lambda$, so both densities carry the factor $t^{-1}\lambda(t)^{\frac{n-3}2}B(L)$, which is $g(f^\star-t)$ by the proof of Proposition~\ref{prop:sl-ends}(ii).

(iv) By Proposition~\ref{prop:sl-density} and Theorem~\ref{thm:gaugeend}(iii) the germ of $\varrho_f$ spans the line. By (ii), $\mathfrak D_{\mathbf p}\nu^{\mathbb T}_f=0$; near $f^\star$, $\nu^{\mathbb T}_f=\varrho^{\mathbb T}_f\one_{(f^\star,f^\star+\delta)}dt$, integrable by (iii), so its germ extends by zero on the torus's side. (iii) gives the common germ.
\end{proof}

\section{Chapter~\hb{ch:readray}: reading a ray}\label{sec:chreadray}

In this section $f^\star=f(\alpha^\star)$, and the symbol $\chi$ and the second symbol $\chi_1$ of polynomials $p_0,\dots,p_h$ are those of Definition~\hb{def:sl-densityop}. One lemma carries the proofs of the selection: the solutions of a recurrence whose symbol has distinct roots, sorted by their growth.

\begin{slemma}\label{lem:rd-growth}
Let $p_0,\dots,p_h\in\mathbb C[k]$, $h\ge1$, have $\deg p_0=\deg p_h=D\ge\deg p_i$ for every $i$, let $\chi$ have distinct roots $\lambda_1,\dots,\lambda_h$, and let $k_1$ be an integer with $p_0(k)p_h(k)\neq0$ for every integer $k\ge k_1$. Let $\mathcal V$ be the space of the sequences $(x(k))_{k\ge k_1}$ with $\sum_ip_i(k)x(k+i)=0$ for every integer $k\ge k_1$.
\begin{enumerate}[label=(\roman*)]
\item $\dim\mathcal V=h$, and $\mathcal V$ has a basis $x_1,\dots,x_h$ with $x_c(k)=\lambda_c^kk^{\rho_c}\big(1+O(1/k)\big)$ as $k\to\infty$, $\rho_c=-\chi_1(\lambda_c)/(\lambda_c\chi'(\lambda_c))$.
\item For $R\ge0$, the $x\in\mathcal V$ with $x(k)=O(r^k)$ for every $r>R$ are the combinations of the $x_c$ with $|\lambda_c|\le R$, and those with $x(k)=O(r^k)$ for some $r<R$ are the combinations of the $x_c$ with $|\lambda_c|<R$.
\end{enumerate}
\end{slemma}

\begin{proof}
(i) The map $x\mapsto(x(k_1),\dots,x(k_1+h-1))$ is a bijection of $\mathcal V$ onto $\mathbb C^h$, the recurrence at $k$ giving $x(k+h)$ from the $h$ values before it, $p_h(k)\neq0$; so $\dim\mathcal V=h$. The leading coefficients $c_i=[k^D]p_i$ have $c_0c_h\neq0$, so $\chi(0)\neq0$ and the roots are nonzero. If $D=0$ the recurrence has constant coefficients, $\rho_c=0$, and $x_c(k)=\lambda_c^k$. Let $D\ge1$. Divided by $k^D$ the recurrence has the coefficients $p_i(k)/k^D$, holomorphic at $k=\infty$ with the values $c_i$ there; written as a system for $(x(k),\dots,x(k+h-1))$, its matrix at $k=\infty$ is the companion matrix of $\chi/c_h$, whose eigenvalues are the $\lambda_c$. This is the normal case of Birkhoff's theorem on linear difference equations, which applies because the coefficients are holomorphic at infinity and the eigenvalues of the matrix there are distinct and nonzero: for every $c$ there is a solution, defined at all large integers, with an asymptotic expansion $\lambda_c^kk^{\rho}\big(1+\sum_{m\ge1}e_mk^{-m}\big)$ as $k\to\infty$ that is a formal solution of the recurrence. Divided by $\lambda_c^k$ it is a formal solution of $\sum_ip_i(k)\lambda_c^i\varsigma^i$, so $\rho=\rho_c$ by Lemma~\hb{lem:sl-formal}, $\lambda_c$ being a simple root. Solving the recurrence for $x(k)$ downward, $p_0(k)\neq0$, extends each of these solutions to $k\ge k_1$; call them $x_c$.

(ii) and the independence. Let $x=\sum_ca_cx_c$ with not all $a_c$ zero, $R_x=\max\{|\lambda_c|:a_c\neq0\}$, $\sigma=\max\{\re\rho_c:a_c\neq0,\ |\lambda_c|=R_x\}$, and $C_0$ the set of the $c$ with $a_c\neq0$, $|\lambda_c|=R_x$ and $\re\rho_c=\sigma$. Then
\[
\frac{x(k)}{R_x^kk^\sigma}=\sum_{c\in C_0}a_c\Big(\frac{\lambda_c}{R_x}\Big)^kk^{i\im\rho_c}+o(1),
\]
the other terms being $O\big((|\lambda_c|/R_x)^kk^{\re\rho_c-\sigma}\big)$ with $|\lambda_c|<R_x$, or $O(k^{\re\rho_c-\sigma})$ with $\re\rho_c<\sigma$, and the remainders $O(1/k)$. The pairs $(\lambda_c/R_x,\im\rho_c)$, $c\in C_0$, are distinct, the $\lambda_c$ being distinct, so by Lemma~\hb{lem:sl-asymp}(i) the right side does not tend to $0$, and $|x(k)|\ge\eta R_x^kk^\sigma$ for some $\eta>0$ and infinitely many $k$. So $x\neq0$, and the $x_c$ are independent, a basis by (i). And $x(k)=O(r^k)$ fails for every $r<R_x$. If $x(k)=O(r^k)$ for every $r>R$, then $R_x\le R$; if for some $r<R$, then $R_x<R$. Conversely $x_c(k)=O(r^k)$ for every $r>|\lambda_c|$.
\end{proof}

\begin{hproposition}[\hb{prop:rd-second} (The low-end relations at every pole)]
As stated.
\end{hproposition}

\begin{proof}
By Proposition~\hb{prop:sl-density} the identity $\sum_{i=0}^hp_i(k)\langle1\rangle_{k+i}=0$ holds as one of meromorphic functions of $k$. For $i>j$ the value $\langle1\rangle_{k+i}$ is holomorphic at $k=-j$, $k+i$ being near the positive integer $i-j$, where it converges absolutely (Lemma~\hb{lem:ch35-direction}(iv), the origin being $0$). So $S(k)=\sum_{i\le\min(h,j)}p_i(k)\langle1\rangle_{k+i}=-\sum_{i=j+1}^hp_i(k)\langle1\rangle_{k+i}$ is holomorphic at $-j$ with the value of (i). For $i\le j$, Proposition~\hb{prop:rays-lowpoles}(ii) at the integer $j-i$ gives $(k+j)^{n-1}\langle1\rangle_{k+i}\to V_1\CT(f^{i-j})$ as $k\to-j$, and $p_i(k)=(k+j)^{m_0}\big([(k+j)^{m_0}]p_i+O(k+j)\big)$. So
\[
(k+j)^{n-1-m_0}S(k)=\sum_{i\le\min(h,j)}\frac{p_i(k)}{(k+j)^{m_0}}\,(k+j)^{n-1}\langle1\rangle_{k+i}\ \longrightarrow\ V_1\sum_{i\le\min(h,j)}[(k+j)^{m_0}]p_i\cdot\CT\big(f^{i-j}\big)
\]
as $k\to-j$. For $m_0<n-1$ the left side tends to $0$, $S$ being holomorphic at $-j$, and $V_1>0$ gives (ii). For $m_0=n-1$ the left side tends to $S(-j)$, which (i) gives, and this is (iii). The reciprocal stride $f^{-1}=\alpha^{-\mvec_1}\prod_lK_l^{s_{1l}}\Psi_l^{t_{1l}}$ is a Laurent polynomial whose coefficients are polynomials with integer coefficients in the weights, so every $\CT(f^{i-j})$ lies in the field of the weights.
\end{proof}

\begin{hproposition}[\hb{prop:rd-slit} (The values live on the slit plane)]
As stated.
\end{hproposition}

\begin{proof}
(i) By Proposition~\hb{prop:sl-moments}(i), $\mu_f$ is carried by $(0,f^\star]$ and $\langle1\rangle_k=\int t^k\mu_f(dt)$. For $|q|<1/f^\star$, $\sum_{k\ge k_0}|q|^{k-k_0}\int t^k\mu_f(dt)=\int t^{k_0}\mu_f(dt)/(1-|q|t)\le\langle1\rangle_{k_0}/(1-|q|f^\star)$, so Fubini's theorem exchanges the sum and the integral and gives the formula. For $q$ in a compact set $Q\subseteq\mathbb C\setminus[1/f^\star,\infty)$ the function $(q,t)\mapsto|1-qt|$ is continuous and positive on $Q\times[0,f^\star]$, since $qt=1$ with $t\in(0,f^\star]$ means $q=1/t\in[1/f^\star,\infty)$; so it is at least some $\delta_Q>0$ there, the integrand is holomorphic in $q$ and at most $t^{k_0}/\delta_Q$, which is $\mu_f$-integrable, and the integral is holomorphic on $\mathbb C\setminus[1/f^\star,\infty)$ (dominated convergence and Morera's theorem).

(ii) Put $t_0=1/x\in(0,f^\star)$. By Proposition~\hb{prop:sl-ends}(i), $\varrho_f$ is real-analytic on $(0,f^\star)$, so there is $\delta>0$ with $[t_0-2\delta,t_0+2\delta]\subseteq(0,f^\star)$ such that $\phi(t)=t^{k_0}\varrho_f(t)$ extends holomorphically to the disc $|t-t_0|\le2\delta$. Write $G=G_1+G_2$, with $G_1(q)=\int_{|t-t_0|\ge\delta}t^{k_0}\mu_f(dt)/(1-qt)$ and $G_2(q)=\int_{t_0-\delta}^{t_0+\delta}\phi(t)\,dt/(1-qt)$. For $q$ close enough to $x$, $1/q$ lies within $\delta/2$ of $t_0$, so $1-qt$ does not vanish on the support of the first integral, and $G_1$ is holomorphic near $x$ as in (i). Let $\Gamma_+$ and $\Gamma_-$ be the half circles $t_0+\delta e^{i\vartheta}$ from $t_0-\delta$ to $t_0+\delta$ through the upper and through the lower half plane. For $q$ near $x$ with $\im q>0$ the pole $t=1/q$ of the integrand lies in the disc and below the real axis, $\im(1/q)=-\im q/|q|^2$, so Cauchy's theorem moves the segment of $G_2$ onto $\Gamma_+$: $G_2(q)=\int_{\Gamma_+}\phi(t)\,dt/(1-qt)$, and the right side is holomorphic in $q$ on the whole neighbourhood of $x$ where $1/q$ stays within $\delta/2$ of $t_0$. So $G$ continues from the upper half plane across a neighbourhood of $x$, with $G(x+i0)=G_1(x)+\int_{\Gamma_+}\phi(t)\,dt/(1-xt)$; likewise from below with $\Gamma_-$. The difference is the integral of $\phi(t)/(1-xt)$ along $\Gamma_+$ followed by $\Gamma_-$ reversed, a clockwise circle around $t_0$, inside which the integrand has one simple pole, at $t_0=1/x$, with the residue $-\phi(1/x)/x$, as $1-xt=-x(t-1/x)$. So $G(x+i0)-G(x-i0)=-2\pi i\cdot\big(-\phi(1/x)/x\big)=2\pi i\,x^{-1-k_0}\varrho_f(1/x)$.

(iii) On $\mathbb T$, $f=\prod_j|K_j(\alpha)|^{-2\sigma_j}\ge f^\star$ (proof of Definition~\hb{def:torusvalue}), finite at the points of $X_F$; so the native value of a native gauge on $\mathbb T$ lies in $[f^\star,\infty)$.

(iv) By Proposition~\hb{prop:rays-conjugates} every conjugate of $f^\star$ over the field $\Bbbk$ of the weights is the native value of a nondegenerate native gauge. Let $\sum_ip_i(k)\langle1\rangle_{k+i}=0$ for all large integers $k$, with $p_i\in\Bbbk[k]$ not all zero, of largest degree $D$. By Proposition~\hb{prop:nativeorder}, $\langle1\rangle_k>0$ for large $k$ and $\langle1\rangle_{k+i}/\langle1\rangle_k\to(f^\star)^i$; dividing the recurrence by $k^D\langle1\rangle_k$ and letting $k\to\infty$ gives $\sum_i[k^D]p_i\cdot(f^\star)^i=\chi(f^\star)=0$. The symbol is a nonzero polynomial over $\Bbbk$, so the minimal polynomial of $f^\star$ over $\Bbbk$ divides it, and every conjugate of $f^\star$ is a root. A real conjugate in $(0,f^\star)$ lies on the slit by its definition.
\end{proof}

\begin{htheorem}[\hb{thm:positiveselection} (Positive selection)]
As stated.
\end{htheorem}

\begin{proof}
As in the proof of Proposition~\ref{prop:rd-slit}(iv), $\chi(f^\star)=0$; the hypotheses on the degrees make $\chi$ of degree $h$ with $\chi(0)=[k^D]p_0\neq0$. The walls $\kappa_{0,S}+k\kappa_{1,S}$ increase with $k$, so the values converge at every integer $k\ge k_1$, and by Proposition~\hb{prop:sl-density} the recurrence holds there: the values are a solution. Lemma~\ref{lem:rd-growth}(i) gives $\dim\mathcal V=h$ and a basis $x_1,\dots,x_h$.

(i) For each $i$, with $m=k+i$,
\begin{align*}
\sum_{k\ge k_1}p_i(k)x(k+i)\,q^{k-k_1+h}&=q^{h-i}\sum_{m\ge k_1+i}p_i(m-i)x(m)\,q^{m-k_1}\\
&=q^{h-i}p_i\Big(q\frac d{dq}+k_1-i\Big)\Big(X-\sum_{m=k_1}^{k_1+i-1}x(m)\,q^{m-k_1}\Big),
\end{align*}
since $q\frac d{dq}$ multiplies $q^{m-k_1}$ by $m-k_1$. Summed over $i$ the left sides vanish, so $L_qX=R_x$ with
\[
R_x=\sum_{i=0}^h\ \sum_{m=k_1}^{k_1+i-1}p_i(m-i)\,x(m)\,q^{h-i+m-k_1},
\]
a polynomial of degree at most $h-1$, linear in $x(k_1),\dots,x(k_1+h-1)$. With $c_i=[k^D]p_i$, $p_i(q\frac d{dq}+c)=c_i(q\frac d{dq})^D+\cdots$ and $(q\frac d{dq})^D=q^D\frac{d^D}{dq^D}+\cdots$, the dots of lower order in $\frac d{dq}$; so the coefficient of $\frac{d^D}{dq^D}$ in $L_q$ is $\sum_iq^{h-i}c_iq^D=q^{D+h}\chi(1/q)$, which vanishes in $\mathbb C^\times$ exactly at the $1/\lambda_c$. For the values, $X$ is the generating function of Proposition~\ref{prop:rd-slit}(i) with $k_0=k_1$.

(ii) For $k\ge k_1$, $\langle1\rangle_k=\int t^k\mu_f(dt)\le(f^\star)^{k-k_1}\langle1\rangle_{k_1}$, $\mu_f$ being carried by $(0,f^\star]$ (Proposition~\hb{prop:sl-moments}(i)); so the values lie in the subspace of the statement, which Lemma~\ref{lem:rd-growth}(ii) at $R=f^\star$ identifies. By Proposition~\hb{prop:sl-ends}(ii), with $a=\frac{n-3}2$ and a small $\delta>0$,
\[
\langle1\rangle_k=\int_0^\delta(f^\star-s)^k\,s^ag(s)\,ds+O\big((f^\star-\delta)^{k-k_1}\langle1\rangle_{k_1}\big).
\]
In $\sigma=-\log(1-s/f^\star)$ the integral is $(f^\star)^k\int_0^{\sigma_\delta}e^{-k\sigma}h(\sigma)\,d\sigma$ with $h(\sigma)=s^ag(s)\,f^\star e^{-\sigma}=g(0)(f^\star)^{a+1}\sigma^a\big(1+O(\sigma)\big)$, and Watson's lemma, which applies because $h(\sigma)\sigma^{-a}$ is continuous at $0$ and $a>-1$, gives $\langle1\rangle_k=C(f^\star)^kk^{-\frac{n-1}2}(1+o(1))$ with $C=g(0)\Gamma(\frac{n-1}2)(f^\star)^{\frac{n-1}2}>0$. Last, $f^\star$ is a simple root of $\chi$ and the recurrence meets the conditions of Proposition~\hb{prop:sl-density}, so Theorem~\hb{thm:gaugeend}(iii) gives $\frac{n-3}2=-1+\chi_1(f^\star)/(f^\star\chi'(f^\star))$.

(iii) By Lemma~\ref{lem:rd-growth}(ii) at $R=\lambda'$, the solutions with $x(k)=O(r^k)$ for some $r<\lambda'$ are the multiples of $x_1$, $\lambda_1=f^\star$ being the only root of modulus less than $\lambda'$; the values are such a solution, $f^\star<\lambda'$, and they are not $0$, so they span that line. A generating function converges on $|q|<1/f^\star$ exactly when $\limsup|x(k)|^{1/k}\le f^\star$, that is when $x(k)=O(r^k)$ for every $r>f^\star$, and Lemma~\ref{lem:rd-growth}(ii) at $R=f^\star$ gives the same line. For the origin $0$, Proposition~\hb{prop:sl-density} gives the low-end relation $a\mathsf L_0+\sum_{i\ge1}p_i(0)\langle1\rangle_i=0$ with $\mathsf L_0=V_1$ (Definition~\hb{def:rays-lowend}), that is $\ell(\langle1\rangle)=V_1$; a nonzero solution of least growth is $x=c\langle1\rangle$ with $c\neq0$, and $\ell(x)=cV_1$.

(iv) $\ell$ is a linear form on $\mathcal V$ with $\ell(\langle1\rangle)=V_1\neq0$, so $\{\ell=0\}$ is a hyperplane that does not contain $\langle1\rangle$, and $\langle1\rangle-x^{(0)}\in\{\ell=0\}$ gives $M$. Take $x_1=\langle1\rangle$, which (iii) allows. The $z_c=x_c-\frac{\ell(x_c)}{V_1}x_1$, $2\le c\le h$, form a basis of $\{\ell=0\}$, and $x^{(i)}=\sum_cS_{ic}z_c$ with $S$ invertible. The equations of the statement read $\mathbf X_KM^{(K)}=-\big(x^{(0)}(K+l)\big)_l$ with $\mathbf X_K=\big[x^{(i)}(K+l)\big]_{l,i}$, while $\mathbf X_KM=\big(\langle1\rangle_{K+l}-x^{(0)}(K+l)\big)_l$; so, when $\mathbf X_K$ is invertible, $M^{(K)}-M=-\mathbf X_K^{-1}v_K$ with $v_K=(\langle1\rangle_{K+l})_l$. Now $\mathbf X_K=\mathbf Z_KS^{\mathsf T}$ with $\mathbf Z_K=[z_c(K+l)]_{l,c}$, and $z_c(K+l)=\lambda_c^KK^{\rho_c}\big(\lambda_c^l+o(1)\big)$, since $x_c(K+l)=\lambda_c^{K+l}(K+l)^{\rho_c}\big(1+O(1/K)\big)$ and $x_1(K+l)=O\big((f^\star)^K\big)$ with $f^\star<|\lambda_c|$. So $\mathbf Z_K=\mathbf N_K\Lambda_K$ with $\Lambda_K=\operatorname{diag}(\lambda_c^KK^{\rho_c})_{c\ge2}$ and $\mathbf N_K\to[\lambda_c^l]_{0\le l\le h-2,\ 2\le c\le h}$, the Vandermonde matrix of the distinct $\lambda_2,\dots,\lambda_h$, which is invertible. For large $K$, $\mathbf X_K$ is invertible, $M^{(K)}$ is unique, and with $|v_K|=O\big((f^\star)^KK^{-\frac{n-1}2}\big)$ by (ii),
\begin{align*}
|M^{(K)}-M|&\le\big\|(S^{\mathsf T})^{-1}\big\|\,\big\|\Lambda_K^{-1}\big\|\,\big\|\mathbf N_K^{-1}\big\|\,|v_K|\\
&=O\Big(\max_{c\ge2}\,|\lambda_c|^{-K}K^{-\re\rho_c}\,(f^\star)^KK^{-\frac{n-1}2}\Big)=O\big((f^\star/\lambda')^KK^{C'}\big).
\end{align*}
Over $\Bbbk$: the recurrence solved forward from initial values in $\Bbbk$, $p_h(k)\neq0$, gives solutions with values in $\Bbbk$, and $\mathcal V$ has a basis of them; $\ell$ has its coefficients in $\Bbbk$ and $V_1\in\Bbbk$, so $x^{(0)}$ and a basis of $\{\ell=0\}$ can be taken with values in $\Bbbk$, and the equations for $M^{(K)}$ then have their coefficients in $\Bbbk$.
\end{proof}

\begin{hproposition}[\hb{prop:rd-moment} (Positivity is the moment problem of the stride's law)]
As stated.
\end{hproposition}

\begin{proof}
(i) Every $\langle1\rangle_{K+m}$, $m\in\mathbb N$, converges with $\langle1\rangle_K$, $\mu_f$ being carried by $(0,f^\star]$, and $L_{K,p}(\langle1\rangle)=\sum_mp_m\int t^{K+m}\mu_f(dt)=\int t^Kp\,d\mu_f$ (Proposition~\hb{prop:sl-moments}(i)). The integrand is nonnegative, and positive off the finitely many zeros of $p$ on $(0,f^\star)$, where $\mu_f$ has a positive density (Proposition~\hb{prop:sl-ends}(i)); so the integral is positive.

(ii) By the Markov--Lukács theorem on polynomials nonnegative on an interval, which applies to $[0,f^\star]$, a real polynomial of degree $2m$ nonnegative there is $A^2+t(f^\star-t)B^2$ with real $A$ and $B$ of degrees at most $m$ and $m-1$, and one of degree $2m+1$ is $tA^2+(f^\star-t)B^2$ with real $A$ and $B$ of degrees at most $m$; every term has degree at most that of the polynomial. For $q=\sum_iq_it^i$, $L_{K,q^2}(x)=\sum_{i,i'}q_iq_{i'}\,x(K+i+i')$, and likewise $L_{K,tq^2}$, $L_{K,(f^\star-t)q^2}$ and $L_{K,t(f^\star-t)q^2}$ are the quadratic forms of the other three Hankel matrices. The four kinds have degree at most $N$ exactly when $\deg q$ is at most $\lfloor\frac N2\rfloor$, $\lfloor\frac{N-1}2\rfloor$, $\lfloor\frac{N-1}2\rfloor$ and $\lfloor\frac{N-2}2\rfloor$, which gives the sizes. If $L_{K,p}(x)\ge0$ for every $p\ge0$ on $[0,f^\star]$ of degree at most $N$, the four quadratic forms are nonnegative, these $p$ being among them; conversely, if the forms are nonnegative, $L_{K,p}(x)\ge0$ for every $p$ of the four kinds, and by linearity for every sum of them.

(iii) If $x(K+m)=\int t^m\nu(dt)$ for a positive $\nu$ on $[0,f^\star]$, then $L_{K,p}(x)=\int p\,d\nu\ge0$. Conversely, the linear functional $\Lambda(p)=L_{K,p}(x)$ on real polynomials is nonnegative on those nonnegative on $[0,f^\star]$, so $|\Lambda(p)|\le\Lambda(1)\max_{[0,f^\star]}|p|$, the polynomials $\max|p|\pm p$ being nonnegative there. Polynomials are dense in $C([0,f^\star])$ (the Weierstrass approximation theorem), so $\Lambda$ extends continuously to $C([0,f^\star])$, and the extension is positive: a continuous $\varphi\ge0$ is the uniform limit of polynomials $p_l$, and $p_l+\max|\varphi-p_l|\ge0$. The Riesz representation theorem gives a positive measure $\nu$ on $[0,f^\star]$ with $\Lambda(\varphi)=\int\varphi\,d\nu$, and at $\varphi=t^m$ this is $x(K+m)=\int t^m\nu(dt)$. Two finite measures on $[0,f^\star]$ with the same moments agree on polynomials, hence on $C([0,f^\star])$, hence are equal. For the values, $t^K\mu_f$ is a finite positive measure on $(0,f^\star]$ with the moments $\langle1\rangle_{K+m}$.
\end{proof}

\begin{htheorem}[\hb{thm:rd-rate} (The rate of positivity off the segment)]
As stated.
\end{htheorem}

\begin{proof}
Put $A(p)=L_{K,p}(\langle1\rangle)=\int t^Kp\,d\mu_f$ and $B(p)=L_{K,p}(u)=\int_{I_\lambda}t^Kp(t)\nu(t)\,dt$, both convergent, $|t|^{K+m}\le C_m|t|^K$ on the bounded $I_\lambda$. For $p\ge0$ on $[0,f^\star]$, $A(p)\ge0$ (Proposition~\ref{prop:rd-moment}(i)), and $A(p)+\varepsilon B(p)\ge0$ holds for every $\varepsilon$ when $B(p)=0$ and otherwise on a closed half line that contains $0$. So $E_N$, the intersection over the $p$ of degree at most $N$, is a closed interval that contains $0$, and if $p_1,p_2$ of degree at most $N$ have $B(p_1)$ and $B(p_2)$ of opposite signs, its length is at most $A(p_1)/|B(p_1)|+A(p_2)/|B(p_2)|$.

Let $w(t)=2t/f^\star-1$, which maps $[0,f^\star]$ onto $[-1,1]$, and $T_m$ the polynomial of degree $m$ with $T_m(\cos\vartheta)=\cos m\vartheta$; then $|T_m(w)|\le1$ for $|w|\le1$ and $|T_m(w)|=\cosh\big(m\operatorname{arccosh}|w|\big)$ for $|w|\ge1$. On the part of $I_\lambda$ off $[0,f^\star]$ put $\Gamma(t)=\operatorname{arccosh}|w(t)|$: it is $C^1$, it grows as $t$ moves from the segment toward $\lambda$, $|w(t)|$ doing so, and $\Gamma(\lambda)=\log\big(|w|+\sqrt{w^2-1}\big)=g(\lambda)$ at $w=w(\lambda)$, the branch of the statement; $g(\lambda)>0$ as $|w(\lambda)|>1$. Let $\eta(t)=f^\star-t$ if $\lambda>f^\star$ and $\eta(t)=t$ if $\lambda<0$, so that $\eta\ge0$ on $[0,f^\star]$ and $\eta(\lambda)<0$, and take $p_1=T_m(w)^2$ and $p_2=\eta\,T_m(w)^2$, nonnegative on $[0,f^\star]$, of degrees $2m$ and $2m+1$. Then $A(p_1)\le\langle1\rangle_K$ and $A(p_2)\le f^\star\langle1\rangle_K$.

Fix $\delta>0$ so small that $I'=\{t\in I_\lambda:\ |t-\lambda|\le\delta\}$ lies off $[0,f^\star]$ and on $I'$: $t^K$ and $\eta$ have the signs of their values at $\lambda$, $|t^K|\ge\frac12|\lambda|^K$, $|\eta(t)|\ge\frac12|\eta(\lambda)|$, $\nu(t)=c|\lambda-t|^a(1+\epsilon(t))$ with $|\epsilon|\le\frac12$, and $\Gamma(t)\ge g(\lambda)-\Gamma_1|\lambda-t|$ for a constant $\Gamma_1$. On $I'$ the integrand of $B(p_i)$ has the constant sign of $c\lambda^K\eta(\lambda)^{i-1}$, and by $\cosh^2\ge\frac14e^{2m\Gamma}$ its modulus is at least $c_1e^{2mg(\lambda)}e^{-2m\Gamma_1|\lambda-t|}|\lambda-t|^a$ for a constant $c_1>0$; so its integral over $I'$ has modulus at least
\[
c_1e^{2mg(\lambda)}\int_0^\delta e^{-2m\Gamma_1s}s^a\,ds=c_1e^{2mg(\lambda)}(2m\Gamma_1)^{-a-1}\big(\Gamma(a+1)+o(1)\big)\ge c_2\,m^{-a-1}e^{2mg(\lambda)}
\]
for large $m$. Off $I'$, $T_m(w)^2\le\cosh^2\big(m\Gamma(\lambda_\delta)\big)\le e^{2m(g(\lambda)-\eta')}$, with $\lambda_\delta$ the end of $I'$ other than $\lambda$ and $\eta'=g(\lambda)-\Gamma(\lambda_\delta)>0$, by the monotonicity of $\Gamma$ and $|T_m|\le1$ on $[0,f^\star]$; $\int_{I_\lambda}|t|^K|\nu|\,dt$ being finite and $\eta$ bounded on $I_\lambda$, that part of $B(p_i)$ is at most $C_3e^{2m(g(\lambda)-\eta')}$ in modulus. So for large $m$, $B(p_1)$ and $B(p_2)$ have the opposite signs of $c\lambda^K$ and $-c\lambda^K$, and moduli at least $\frac{c_2}2m^{-a-1}e^{2mg(\lambda)}$. With $m=\lfloor N/2\rfloor$ for $p_1$ and $m=\lfloor(N-1)/2\rfloor$ for $p_2$, so that $2m\ge N-2$, the length of $E_N$ is at most $C\,N^{a+1}e^{-Ng(\lambda)}$ for large $N$.

For $\lambda>f^\star$, $\rho=f^\star/\lambda\in(0,1)$, $w=\frac2\rho-1>1$, $w^2-1=4(1-\rho)/\rho^2$ and $w+\sqrt{w^2-1}=\big(2-\rho+2\sqrt{1-\rho}\big)/\rho=\big(1+\sqrt{1-\rho}\big)^2/\rho$. So $e^{-g(\lambda)}=\rho/(1+\sqrt{1-\rho})^2=(1-\sqrt{1-\rho})/(1+\sqrt{1-\rho})$, using $\rho=(1-\sqrt{1-\rho})(1+\sqrt{1-\rho})$, and $\sqrt{1-\rho}=1-\frac\rho2+O(\rho^2)$ gives $\frac\rho4+O(\rho^2)$.
\end{proof}

\begin{hproposition}[\hb{prop:rd-blind} (What positivity cannot see)]
As stated.
\end{hproposition}

\begin{proof}
(i) $\mu_f+\varepsilon\nu\ge\mu_f-|\varepsilon|\,|\nu|\ge(1-|\varepsilon|C)\mu_f\ge0$. For real $k\ge K$, $t^k\le(f^\star)^{k-K}t^K$ on $(0,f^\star]$, so $\int t^k\,d|\nu|\le C(f^\star)^{k-K}\langle1\rangle_K$, and $\langle1\rangle_k+\varepsilon u(k)=\int t^k\,d(\mu_f+\varepsilon\nu)$ (Proposition~\hb{prop:sl-moments}(i)). For $K'\ge K$ and $p\ge0$ on $[0,f^\star]$, $L_{K',p}(\langle1\rangle+\varepsilon u)=\int t^{K'}p\,d(\mu_f+\varepsilon\nu)\ge0$.

(ii) $t^K(\mu_f+\varepsilon\nu)$ is a finite measure on $[0,f^\star]$, of total variation at most $(1+|\varepsilon|C)\langle1\rangle_K$, with the stated moments; two finite signed measures on $[0,f^\star]$ with the same moments agree on polynomials, hence on $C([0,f^\star])$ by the Weierstrass approximation theorem, hence are equal. On $J$ its density is $t^K(\varrho_f+\varepsilon\nu')$, $\nu'$ the density of $\nu$. If it agreed almost everywhere on $(0,f^\star)$ with a real-analytic function $\varphi$, then on $J$, $\nu'=(t^{-K}\varphi-\varrho_f)/\varepsilon$ almost everywhere, a real-analytic function on $J$, $\varrho_f$ being real-analytic on $(0,f^\star)$ (Proposition~\hb{prop:sl-ends}(i)); this is excluded. $t^K\mu_f$ has the real-analytic density $t^K\varrho_f$.

(iii) Let $|\nu'|\le B$ on $(0,\delta]$. For $K_0>-1$, $\int t^{K_0}d|\nu|\le B\int_0^\delta t^{K_0}dt+\max_{[\delta,f^\star]}t^{K_0}\cdot C\mu_f([\delta,f^\star])<\infty$ (Proposition~\hb{prop:sl-moments}(i)), so $u(k)$ converges for $\re k\ge K_0$. Lemma~\hb{lem:sl-moment}(ii), with the measure $\nu$ on $(0,f^\star]$, the operator $\mathfrak D_{\mathbf p}$, for which $\mathfrak D_{\mathbf p}^\dagger t^k=\sum_ip_i(k)t^{k+i}$, and $E=\mathfrak D_{\mathbf p}\nu=0$, gives $\sum_ip_i(k)u(k+i)=0$ for $\re k\ge K_0$, for every $K_0>-1$. At $k=0$, $p_0(0)=0$, $p_0$ vanishing to order at least $n-1\ge1$ at $0$ (Proposition~\hb{prop:sl-density}), and what remains is $\sum_{i\ge1}p_i(0)u(i)=0$.
\end{proof}

\begin{hproposition}[\hb{prop:rd-onemaster} (A recurrence of order two is the volume of the direction times a continued fraction)]
As stated.
\end{hproposition}

\begin{proof}
(i) $\langle1\rangle_{k+1}=\int t^{k+1}\mu_f(dt)<f^\star\int t^k\mu_f(dt)$, $\mu_f$ having a positive density on $(0,f^\star)$ and no mass at $f^\star$, whose preimage is the point $\log\alpha^\star$ (Proposition~\hb{prop:sl-ends}(i)); and $\langle1\rangle_k>0$. By Proposition~\hb{prop:sl-density} the recurrence holds at every integer $k\ge1$, the origin being $0$; divided by $\langle1\rangle_{k+1}$ it reads $p_0(k)/\rho_k+p_1(k)+p_2(k)\rho_{k+1}=0$.

(ii) The low-end relation (Proposition~\hb{prop:sl-density}, Definition~\hb{def:rays-lowend}) is $aV_1+p_1(0)\langle1\rangle_1+p_2(0)\langle1\rangle_2=0$, and $\langle1\rangle_2=\rho_1\langle1\rangle_1$. Let $\varphi(\rho)=-a/(p_1(0)+p_2(0)\rho)$. If $\varphi(0)$ and $\varphi(f^\star)$ are finite and positive, the affine denominator has the sign of $-a$ at both ends of $[0,f^\star]$, so it does not vanish there, $a\neq0$, and $\varphi$ is continuous and monotone on $[0,f^\star]$, strictly unless $p_2(0)=0$, when the two values agree. As $\rho_1\in(0,f^\star)$, $\langle1\rangle_1/V_1=\varphi(\rho_1)$ lies between $\varphi(0)$ and $\varphi(f^\star)$, strictly when they differ.

(iii) The symbol $\chi(\lambda)=c_0+c_1\lambda+c_2\lambda^2$ has real coefficients and the real root $f^\star$, so $\lambda_2$ is real, with $|\lambda_2|=\lambda'>f^\star$. By Lemma~\ref{lem:rd-growth}(i) and the proof of Theorem~\ref{thm:positiveselection}(iii), $\mathcal V$ has the basis $\langle1\rangle$, $x_2$, with $x_2(k)=\lambda_2^kk^{\rho_2}(1+O(1/k))$. For $r\in[0,f^\star]$ and $K\ge1$ put
\[
\psi=\langle1\rangle_{K+2}-r\langle1\rangle_{K+1},\qquad\phi=x_2(K+2)-r\,x_2(K+1),\qquad y=\phi\,\langle1\rangle-\psi\,x_2 ,
\]
a solution with $y(K+2)=r\,y(K+1)$. Let $\tau_k=x_2(k+1)/x_2(k)$, defined for $k\ge J$ and tending to $\lambda_2$, and $\rho_k\to f^\star$ (Proposition~\hb{prop:nativeorder}). Then $\phi=x_2(K+1)(\tau_{K+1}-r)$ with $|\tau_{K+1}-r|\ge\frac12(\lambda'-f^\star)$ for large $K$, so $|\phi|\ge c\,(\lambda')^KK^{\rho_2}$, and $|\psi|\le2f^\star\langle1\rangle_{K+1}=O\big((f^\star)^KK^{-\frac{n-1}2}\big)$ by Theorem~\ref{thm:positiveselection}(ii), uniformly in $r$. So $|\psi/\phi|\le C_1(f^\star/\lambda')^KK^{C'}$. Write $y(k)=\phi\,\langle1\rangle_k\big(1-\epsilon_k\big)$ with $\epsilon_k=\frac\psi\phi\frac{x_2(k)}{\langle1\rangle_k}$. For $k<J$, $|\epsilon_k|\le C_J|\psi/\phi|\to0$. For $J\le k\le K+1$,
\[
\epsilon_k=\frac{\rho_{K+1}-r}{\tau_{K+1}-r}\prod_{l=k}^K\frac{\rho_l}{\tau_l},
\]
and choosing $J$ larger if needed, $|\rho_l/\tau_l|\le\bar q=\frac12(1+f^\star/\lambda')<1$ for $l\ge J$, while $|(\rho_{K+1}-r)/(\tau_{K+1}-r)|\to(f^\star-r)/|\lambda_2-r|\le(f^\star-r)/(\lambda'-r)\le f^\star/\lambda'<\bar q$ uniformly in $r\in[0,f^\star]$. So there is $K_0$ such that for $K\ge K_0$ and every $r$, $|\epsilon_k|<1$ for $1\le k\le K+1$, hence $y(k)\neq0$ there. For such $K$ put $\sigma_k=y(k+1)/y(k)$, $1\le k\le K+1$; $\sigma_{K+1}=r$, and the recurrence at $k$, divided by $y(k+1)$, gives $p_1(k)+p_2(k)\sigma_{k+1}=-p_0(k)/\sigma_k\neq0$, so $\sigma_k=\mathsf M_k(\sigma_{k+1})$ for $1\le k\le K$. So $\mathsf M_1\circ\cdots\circ\mathsf M_K(r)=\sigma_1$ is defined, and
\[
\sigma_1=\frac{\langle1\rangle_2-(\psi/\phi)x_2(2)}{\langle1\rangle_1-(\psi/\phi)x_2(1)}=\rho_1+O\big((f^\star/\lambda')^KK^{C'}\big),
\]
uniformly in $r$. By (ii) and $a\neq0$, $p_1(0)+p_2(0)\rho_1\neq0$ and $\langle1\rangle_1=-aV_1/(p_1(0)+p_2(0)\rho_1)$. Let $\Sigma_K$ be the set of the $\mathsf M_1\circ\cdots\circ\mathsf M_K(r)$ over the $r\in(0,f^\star)$ whose intermediate values lie in $(0,f^\star)$. An element of $\Sigma_{K+1}$, built from $r$, is $\mathsf M_1\circ\cdots\circ\mathsf M_K(r')$ with $r'=\mathsf M_{K+1}(r)\in(0,f^\star)$ and the same intermediate values, so it lies in $\Sigma_K$; and $\rho_1\in\Sigma_K$, built from $r=\rho_{K+1}$, its intermediate values being the $\rho_k\in(0,f^\star)$ by (i).
\end{proof}

\begin{hdefinition}[\hb{def:readprocedure} (The procedure that reads a ray from its list)]
As stated.
\end{hdefinition}

\begin{proof}
(1) $V_1$ is the coefficient of Proposition~\hb{prop:principal}; the maps of Proposition~\hb{prop:raysym}(ii) that are symmetries generate a group of symmetries, and a finite one is chosen. (2) The search stops by Theorem~\hb{thm:rayboxstop}, every box of it being finite (Proposition~\hb{prop:box-walldegree}(iv)); Proposition~\ref{prop:rd-certain}(iii) and~(v) bound the strengths and primes at which an elimination misjudges acceptance. (3) The reductions from the relations the targets need are those of the box by Proposition~\ref{prop:rd-certain}(vi); the evaluations give the values of the entries of $\mathbf B$ at the strengths where the elimination has the generic pivots (Proposition~\ref{prop:rd-certain}(i)), the reconstruction is Proposition~\ref{prop:rd-certain}(ii), the check is Proposition~\ref{prop:rd-certain}(iv) and~(v), and the certificate and the one-step relation are Theorem~\hb{thm:raybox}(ii) and~(iii). (4) By Theorem~\hb{thm:raybox}(iv) Proposition~\hb{prop:rays-lawrows} applies to the spanning set and its matrix: the first dependence $\sum_iq_iw_i=0$ gives $\sum_jp_j(k)f^j\in\mathcal L$ (Proposition~\hb{prop:rays-lawrows}(ii)), and the ratios $q_i/q_h$ are ratios of signed maximal minors (Proposition~\hb{prop:rays-lawrows}(iii)), rational functions whose values at a point $k_0$ where $\mathbf B(k_0-h+1),\dots,\mathbf B(k_0)$ are finite and $w_0(k_0),\dots,w_{h-1}(k_0)$ independent are the coefficients of the dependence of the rows at $k_0$, so Proposition~\ref{prop:rd-certain}(ii) reconstructs them. The relation is the law of $[1]$ when the spanning set is independent (Theorem~\hb{thm:raybox}(iv)), as under Proposition~\hb{prop:rays-spectrum}(i). The zeros of $p_0(k)=q_h(k+h)$ are those described, by the divisibility of Proposition~\hb{prop:rays-lawrows}(iii). (5) is Definition~\hb{def:sl-gaugeread}, which applies by Proposition~\hb{prop:sl-density}, or Theorem~\ref{thm:positiveselection}(iv) with Proposition~\ref{prop:rd-moment}. (6) The closings are a solution with values in germs at $-\infty$ (Theorem~\hb{thm:closings}), so Theorem~\hb{thm:raybox}(iii) gives $\varsigma^{-1}v=\mathbf Bv$ for their vector $v(k)=\mathsf C_{-k}(\mathbf b)$, that is $\mathsf C_{j+1}(\mathbf b)=\mathbf B(-j)\mathsf C_j(\mathbf b)$ for all large $j$; the reflected recurrence and the low-end relation hold by Proposition~\hb{prop:sl-density} with Corollary~\hb{cor:rays-lowend}(i) and~(ii); the eigenvalues are Proposition~\hb{prop:rays-spectrum}(ii); the exponent is Theorem~\hb{thm:gaugeend}(iii).
\end{proof}

\begin{hproposition}[\hb{prop:rd-lattice} (The relations of a box are difference equations)]
As stated.
\end{hproposition}

\begin{proof}
(i) By Definition~\hb{def:raybox}(i) the symbols at the level $(\mathbf a,\mathbf b)$ are the $(u;\mathbf a,\mathbf b)$ with $u\in\mathbb Z^{\cE}$ and $u_{\cE}=c$, $c=\sum_j(a_jr_j+b_jr^\ast_j)$: a translate of the lattice $\{u_{\cE}=0\}$, of rank $n-1$. As $\kappa_{1,U}>0$, the wall degree of $g$ is at most $\mathfrak t$ exactly when $\kappa_U(g)+\mathfrak t\kappa_{1,U}\ge0$ for every nonempty proper $U$ (Definition~\hb{def:raybox}(ii)). With $L(U)=\sum_j\big(a_jd_j(U)+b_jd^\ast_j(U)\big)$ and $\kappa_U(g)=u_U-L(U)$, the sets $U=\{e\}$ and $U=\cE\setminus\{e\}$, nonempty and proper as $n\ge2$, give
\[
L(\{e\})-\mathfrak t\kappa_{1,\{e\}}\le u_e\le c-L(\cE\setminus\{e\})+\mathfrak t\kappa_{1,\cE\setminus\{e\}},
\]
an interval of length at most $C(1+\mathfrak t)$ with $C=\max_e\max\big(|c-L(\cE\setminus\{e\})-L(\{e\})|,\ \kappa_{1,\{e\}}+\kappa_{1,\cE\setminus\{e\}}\big)$, which holds at most $C(1+\mathfrak t)+1$ integers; and $u_{e_0}=c-\sum_{e\neq e_0}u_e$.

(ii) By Definition~\hb{def:raybox}(iii) the letter identity of $g=(u;\mathbf a,\mathbf b)$ for $K_j$ has the terms $g$ and $\theta_{j,T}(u+\one_T;\mathbf a+\mathbf e_j,\mathbf b)$, $T\in\cB_j$, and that for $\Psi_j$ the terms $g$ and $\theta_{j,T}(u+\one_{\cE\setminus T};\mathbf a,\mathbf b+\mathbf e_j)$; the mark relation at $e$ has $g$ and, for every $j$, the terms of the $T\in\cB_j$ with $T\ni e$ at the level $(\mathbf a+\mathbf e_j,\mathbf b)$ and of the $T\not\ni e$ at the level $(\mathbf a,\mathbf b+\mathbf e_j)$, at most $|\cB_j|$ for each $j$, with coefficients affine in $k$. The shifts are the $(\one_T;\mathbf e_j,\mathbf 0)$ and the $(\one_{\cE\setminus T};\mathbf 0,\mathbf e_j)$, $T\in\cB_j$.

(iii) The letter identities have the coefficients $1$ and $-\theta_{j,T}$, and the mark relations $m_{0,e}+k\,m_{1,e}+u_e$, $-(s_{0j}+k\,s_{1j}+a_j)\theta_{j,T}$ and $-(t_{0j}+k\,t_{1j}+b_j)\theta_{j,T}$, affine in $k$, with the stated parts in $k$. For a relation $\sum_gc_gg$, $e_G\sum_gc_gg=\sum_gc_g\,e_Gg$, and $e_Gg$ is the orbit average of the orbit of $g$ (Definition~\hb{def:box-reduction}(i)); so in the basis of the orbit averages the entry of an orbit is the sum of the $c_g$ over its symbols, again affine in $k$, the part in $k$ vanishing for the letter identities.
\end{proof}

\begin{hproposition}[\hb{prop:rd-certain} (What a reading by evaluation certifies)]
As stated.
\end{hproposition}

\begin{proof}
(i) Gaussian elimination of the rows of $A$ over $\Bbbk(k)$, the columns taken in the order of the ranking, highest first, gives a basis of the row space $\mathcal R^G_{\mathfrak B}$ in echelon form, with distinct leading orbits, and the leading orbit of every nonzero element of the row space is one of them; so these are the pivots $P$ of Definition~\hb{def:box-reduction}(ii), and $r=|P|$ is the rank. An element of the row space with zero coordinates on $P$ would have a leading orbit outside $P$, so it is $0$: the projection of $\mathcal R^G_{\mathfrak B}$ onto the coordinates $P$ is injective, hence bijective. The chosen rows are a basis of $\mathcal R^G_{\mathfrak B}$, so $A_P(k)$ is invertible, and $E=A_P^{-1}[A_P\,|\,A_F]$, written on the columns $P$ and $F$, is the basis of $\mathcal R^G_{\mathfrak B}$ with the identity on $P$, the unique one by the bijectivity; it is therefore the reduced echelon basis, whose row of a pivot $p$ vanishes on the orbits ranked above $p$. For $x$ with coefficients in $\Bbbk$, $x-x_PE$ is congruent to $x$ and vanishes on $P$, so it is the reduction, $x_F-x_PA_P^{-1}A_F$ on $F$. The entries of $A_P$ are affine in $k$, so $\det A_P$, not $0$, has degree at most $r$. At $k_0$: the minors of order $r+1$ of $A$ vanish identically, so $A(k_0)$ has rank at most $r$. If the minor of order $r$ of the rows $J'$ on the columns $P$ does not vanish at $k_0$, then $E=A_{J',P}^{-1}A_{J'}$ as well, by the uniqueness, so $E$ is regular at $k_0$, and $E(k_0)$ has the identity on $P$ and the zeros of an echelon form, identities in $k$; its $r$ independent rows lie in the row space of $A(k_0)$ and span it, so the elimination at $k_0$ has the pivots $P$ and gives the reduction $x_F-x_PE_F(k_0)$, the value at $k_0$. Conversely, if the elimination at $k_0$ has the pivots $P$, the projection of the row space of $A(k_0)$ onto the coordinates $P$ is injective, as above, so the columns $P$ of $A(k_0)$ have rank $r$, and a minor of order $r$ on them does not vanish.

(ii) The targets $f^{-1}b_i$ have their coefficients in $\Bbbk$. By (i) their reductions are $\big(x_F\det A_P-x_P\operatorname{adj}(A_P)A_F\big)/\det A_P$, the adjugate having entries of degree at most $r-1$ and $A_F$ entries of degree at most $1$; the entries of $\mathbf B$ are coefficients of these reductions (Definition~\hb{def:box-reduction}(iii)). If $N/D$ and $N'/D'$ have numerators of degree at most $d_1$ and denominators of degree at most $d_2$ and agree at $d_1+d_2+1$ points where $D$ and $D'$ do not vanish, then $ND'-N'D$, of degree at most $d_1+d_2$, vanishes at these points and is $0$. With a common denominator, two numerators of degree at most $d_1$ that agree at $d_1+1$ points are equal.

(iii) Write the coefficient of a free orbit $c$ in the reduction of a target as $N_c/\det A_P$, $\deg N_c\le r$, by (ii). If the box is accepted at $\mathfrak s$ over $\Bbbk(k)$, every $N_c$ with $c$ outside the free Laurent orbits of wall degree at most $\mathfrak s$ vanishes, and at every $k_0$ with $\det A_P(k_0)\neq0$ the elimination has the pivots $P$, so the same free orbits, and gives the values of the reductions: it accepts. If not, some such $N_c$ is a nonzero polynomial, and at every $k_0$ with $\det A_P(k_0)N_c(k_0)\neq0$ the elimination has the pivots $P$ and its reduction has a nonzero coefficient on $c$: it rejects. The exceptions lie among the at most $2r$ zeros of $\det A_P$ and $N_c$.

(iv) Let $\mathbf B_{il}\neq\mathbf B'_{il}$, write $\mathbf B_{il}=N_{il}/\det A_P$ by (ii) and $\mathbf B'_{il}=N'_{il}/Q'$, and put $\Delta=N_{il}Q'-N'_{il}\det A_P$, a nonzero polynomial of degree at most $r+d$. At a strength $k_0$ with $Q'(k_0)\neq0$ at which the elimination has the pivots $P$, it gives $\mathbf B(k_0)$, $\mathbf B$ being regular at $k_0$ by (i). If $\det A_P(k_0)\neq0$, agreement there gives $\Delta(k_0)=0$; if $\det A_P(k_0)=0$, then $N_{il}=\mathbf B_{il}\det A_P$ vanishes at $k_0$, $\mathbf B_{il}$ being regular there, and $\Delta(k_0)=0$ again. So the strengths of agreement are among the at most $r+d$ zeros of $\Delta$, and a uniformly drawn strength of $S$ is one of them with probability at most $(r+d)/|S|$.

(v) The entries of $A$ are $q$-integral, and reduction modulo $q$ is a ring map, so the minors of order $r+1$ of $A$ vanish modulo $q$ and $\det A_P$ modulo $q$ is the determinant of $A_P$ modulo $q$, which is not $0$; the entries of $E=\operatorname{adj}(A_P)A_{J}/\det A_P$, $A_J$ the chosen rows, reduce to those of $\bar E=\operatorname{adj}(\bar A_P)\bar A_J/\det\bar A_P$, which has the identity on $P$ and the zeros of an echelon form. As in (i), $\bar A$ has the rank $r$ and the pivots $P$ over $\mathbb F_q(k)$, and its reductions $x_F-x_P\bar E_F$ are those over $\mathbb Q(k)$ taken modulo $q$; the proofs of (i) to (iv) use only linear algebra over a field and hold over $\mathbb F_q$. For the last claim, take an entry with $\mathbf B_{il}\neq\mathbf B'_{il}$ and write $\mathbf B_{il}=\tilde N/\tilde D$ and $\mathbf B'_{il}=N'/Q'$ with $\tilde N,\tilde D,N',Q'\in\mathbb Z[k]$ and $\tilde D=c_A\det A_P$ for a positive integer $c_A$; then $\Delta=\tilde NQ'-N'\tilde D\in\mathbb Z[k]$ is not $0$. The finite set consists of the primes excluded above and those that divide $c_A$, every coefficient of $Q'$, or every coefficient of $\Delta$. For a prime outside it, $\mathbf B'$ modulo $q$ is $N'/Q'$ taken modulo $q$, $\Delta$ modulo $q$ is a nonzero polynomial of degree at most $r+d$, and the argument of (iv) over $\mathbb F_q$ bounds the probability by $(r+d)/q$.

(vi) An element of $\mathcal R'$ lies in $\mathcal R^G_{\mathfrak B}$ with the same leading orbit, so the pivots of $\mathcal R'$ are pivots of $\mathcal R^G_{\mathfrak B}$, and the free orbits of $\mathcal R^G_{\mathfrak B}$ are free for $\mathcal R'$. So $\bar x$ is a combination of free orbits of $\mathcal R'$ with $x-\bar x\in\mathcal R'$. The free orbits of $\mathcal R'$ are independent modulo $\mathcal R'$, an element of $\mathcal R'$ that combines them having a free leading orbit and so being $0$; so $\bar x$ is the reduction modulo $\mathcal R'$.
\end{proof}

\begin{hproposition}[\hb{prop:rd-rational} (Data with rational masses)]
As stated.
\end{hproposition}

\begin{proof}
(i) Both data lie on the balance wall with the same exponents, so the origin $(\mvec-\mvec';\mathbf 0,\mathbf 0)$ has $(\mvec-\mvec')_{\cE}=0$ and lies on it too. The walls are linear in the data, so the walls along the ray at $k$ are $\kappa_U(\mvec;\mathbf s,\mathbf t)+(k-1)\kappa_U(\mvec';\mathbf s,\mathbf t)$, positive for real $k\ge1$; by Lemma~\hb{lem:ch35-direction}(iv) the values converge there and $\langle1\rangle_1=P_F(\mvec;\mathbf s,\mathbf t)$, the data at $k=1$ being $(\mvec;\mathbf s,\mathbf t)$.

(ii) The direction $N(\mvec;\mathbf s,\mathbf t)$ is integral, on the balance wall, with the walls $N\kappa_U>0$; its stride is $\alpha^{N\mvec}\prod_jK_j^{-Ns_j}\Psi_j^{-Nt_j}=f_{\mvec}^N$ on the positive orthant, and by Lemma~\hb{lem:ch35-direction}(iv) at $k=\frac1N$, where the walls are the $\kappa_U>0$, $\langle1\rangle_{1/N}=P_F(\mvec;\mathbf s,\mathbf t)$. The wall degree is that of Definition~\hb{def:raybox}(ii) with $\kappa_{1,U}=N\kappa_U$.
\end{proof}

\begin{hproposition}[\hb{prop:rd-lift} (One factorization serves every coefficient)]
As stated.
\end{hproposition}

\begin{proof}
$A_P$ is affine in $k$, so $A_P(k_0+\epsilon)=M+\epsilon E$. For $|\epsilon|\,\|EM^{-1}\|<1$, $M+\epsilon E=(I+\epsilon EM^{-1})M$ and $(M+\epsilon E)^{-1}=M^{-1}\sum_{j\ge0}(-\epsilon EM^{-1})^j$, so $x_P(M+\epsilon E)^{-1}=\sum_jx_PM^{-1}(-EM^{-1})^j\epsilon^j$, which is the recursion. With $A_F(k_0+\epsilon)=A_F(k_0)+\epsilon E_F$, $E_F=A_F(k_0+1)-A_F(k_0)$, the coefficient of $\epsilon^j$, $j\ge1$, in the reduction $x_F-x_PA_P^{-1}A_F$ is $-\lambda_jA_F(k_0)-\lambda_{j-1}E_F$. Each $\lambda_j$ is a product of a row by $E$ and one solve with the factorization of $M$. If $N/D$ and $N'/D'$, with the degree bounds and $D(k_0)D'(k_0)\neq0$, have the same Taylor coefficients up to the order $d_1+d_2$ at $k_0$, then $ND'-N'D=DD'(N/D-N'/D')$ vanishes to the order $d_1+d_2+1$ at $k_0$ and has degree at most $d_1+d_2$, so it is $0$.
\end{proof}

\begin{hproposition}[\hb{prop:rd-masterdensities} (The densities of a spanning set)]
As stated.
\end{hproposition}

\begin{proof}
(i) A Laurent insertion $b$ with real coefficients is real and continuous on the positive part of the chart, and $b\,\omega_0$ has the density $b\,\omega_0/dy$; the sets $\{f\ge t\}$, $t>0$, are compact (proof of Proposition~\hb{prop:sl-moments}(i)), so $\mu_{f,b}$ is finite on $[t,f^\star]$, with $|\mu_{f,b}|$ at most the image of $|b|\,\omega_0$. Where $\int|b|f^{\re k}\omega_0<\infty$, $\langle b\rangle_k=\int bf^k\omega_0=\int t^k\mu_{f,b}(dt)$ by the change of variables of an image measure. The proof of Proposition~\hb{prop:sl-ends}(i), with the weight $b\,\omega_0/dy$ in place of $\omega_0/dy$, gives a real-analytic density on $(0,f^\star)$; its positivity is not used there. The proof of Proposition~\hb{prop:sl-ends}(ii) uses Lemma~\hb{lem:sl-morse}(i), which allows every real-analytic weight $w$, here $w=b\,\omega_0/dy$, with $w(\log\alpha^\star)=b(\alpha^\star)\,(\omega_0/dy)(\alpha^\star)$; so the germ has the stated form with $g_b(0)=b(\alpha^\star)g(0)$.

(ii) Every monomial of $b_l$ has an absolutely convergent value for large real $k$ (Proposition~\hb{prop:rays-strength}(i)), so $\int t^{K_0}d|\mu_{f,b_l}|<\infty$ and $\int t^{K_0}d|t^{-1}\mu_{f,b_i}|<\infty$ for some $K_0$. Apply Lemma~\hb{lem:sl-moment} to the measures $t^{-1}\mu_{f,b_i}$ and $\mu_{f,b_l}$ on $(0,f^\star]$ and the operators $Q(-1-\theta)$ and $-\tilde P_{il}(-1-\theta)$, whose adjoints are $Q(\theta)$ and $-\tilde P_{il}(\theta)$, $\theta^\dagger=-1-\theta$, with $Q(\theta)t^k=Q(k)t^k$. Its numbers are $M_E(k)=Q(k)\langle b_i\rangle_{k-1}-\sum_l\tilde P_{il}(k)\langle b_l\rangle_k$, which vanish for large $k$; by Lemma~\hb{lem:sl-moment}(i) the distribution $E$ vanishes, and that is the system.

(iii) The identity $v(k-1)=\mathbf B(k)v(k)$ of meromorphic functions, iterated, gives $v(k-K+1)=\mathbf B(k-K+2)\cdots\mathbf B(k)\,v(k)$. Near $k=K$ the right side is holomorphic, the factors $\mathbf B(m)$, $2\le m\le K$, being finite and $v$ converging in a neighbourhood of $K$, the walls at $K$ being positive; so the left side, $v$ near $1$, is holomorphic at $1$ with the stated value.
\end{proof}

\begin{hproposition}[\hb{prop:rd-sizes} (The native gauges of all exponents)]
As stated.
\end{hproposition}

\begin{proof}
(i) $K_j$ and $\Psi_j$ are homogeneous of the degrees $r_j$ and $r^\ast_j$, so by Euler's identity $\sum_eR_{j,e}=r_j$ and $\sum_eR^\ast_{j,e}=r^\ast_j$, and $\mathbf u(\alpha;\mathbf s,\mathbf t)_{\cE}=\sum_j(s_jr_j+t_jr^\ast_j)$: the balance wall. A native gauge of the ray is a point of $X_F$ with $c^{(1)}_e=m_{1,e}-\sum_j\big(s_{1j}R_{j,e}+t_{1j}R^\ast_{j,e}\big)=0$ for every $e$ (Definitions~\hb{def:ray} and~\hb{def:raygauge}), that is $\mvec_1=\mathbf u(\alpha;\mathbf s_1,\mathbf t_1)$; for data on the balance wall with $m_{\cE}\neq0$ the target of Definition~\hb{def:native} is $\mvec$. The set is the graph of the regular map $(\alpha,\mathbf s,\mathbf t)\mapsto\mathbf u(\alpha;\mathbf s,\mathbf t)$, isomorphic by the projection to $X_F\times\mathbb C^{2N}$; $X_F$ is the complement in the torus $(\mathbb C^\times)^{\cE}/\mathbb C^\times$, of dimension $n-1$, of the hypersurface $\prod_jK_j\Psi_j=0$, which does not contain $\one$, so it is a nonempty open subset of an irreducible smooth variety, irreducible and smooth itself.

(ii) Let $\alpha_1,\dots,\alpha_\nu$ be the nondegenerate native gauges of the ray, $w_1$ its direction, and $W$ the data on the balance wall, a vector space of dimension $n-1+2N$ with the coordinates $(m_e)_{e\neq e_0}$, $\mathbf s$, $\mathbf t$. In the coordinates $y_e=\log\alpha_e$, $e\neq e_0$, of the chart near each $\alpha_c$, the map $(y,\mathbf s,\mathbf t)\mapsto\big(\mathbf u(e^y;\mathbf s,\mathbf t);\mathbf s,\mathbf t\big)$ into $W$ is holomorphic, and its derivative is block triangular, with the identity on $(\mathbf s,\mathbf t)$ and the matrix $\big(\partial_{y_{e'}}\mathbf u_e\big)_{e,e'\neq e_0}=-\big(\partial_{y_{e'}}c^{(1)}_e\big)$, minus the Hessian of $\log f$ on the chart, which is nonsingular at $\alpha_c$ (Definition~\hb{def:raygauge}). By the holomorphic inverse function theorem there are disjoint neighbourhoods $U_c$ of the $\alpha_c$ in $X_F$ and a neighbourhood $\mathcal W$ of $w_1$ in $W$ such that every $w\in\mathcal W$ has exactly one point native for it in each $U_c$. The exceptional set of Theorem~\hb{thm:rayspace}(i), whose count concerns the critical points of the product of powers of the letters with the exponents of the direction, is a countable union of proper subvarieties; a proper subvariety is closed with empty interior in $W$, and its real points have empty interior in the real data, a polynomial that vanishes on an open set of real points being $0$. By the Baire category theorem $\mathcal W$ contains data, real if need be, outside the exceptional set; at such data the native points are $(-1)^{n-1}\operatorname{eu}(X_F)$ in number, and they include the $\nu$ points in the $U_c$.
\end{proof}

\begin{hdefinition}[\hb{def:rd-continuation} (The procedure that finds native gauges by continuation)]
As stated.
\end{hdefinition}

\begin{proof}
(1) is Proposition~\ref{prop:rd-sizes}(i). Steps (2) and (3) only propose points; nothing in the output rests on them. (4) In the coordinates $\alpha_e$, $e\neq e_0$, of the chart the map $F=(c^{(1)}_e)_{e\neq e_0}$ is rational, with no pole on a box $Y$ on which no letter vanishes; read as a map of a real space of dimension $2(n-1)$, it is continuously differentiable there. Krawczyk's test asks for a point $\tilde y\in Y$, a nonsingular real matrix $\mathcal Y$ and an interval enclosure $F'(Y)$ of the derivative on $Y$ with $\tilde y-\mathcal YF(\tilde y)+(I-\mathcal YF'(Y))(Y-\tilde y)\subseteq Y$. By the mean value theorem in integral form, $y-\mathcal YF(y)=\tilde y-\mathcal YF(\tilde y)+\int_0^1\big(I-\mathcal YF'(\tilde y+\tau(y-\tilde y))\big)d\tau\,(y-\tilde y)$, and the integral of matrices in the convex set $F'(Y)$ lies in it; so $y\mapsto y-\mathcal YF(y)$ maps the compact convex set $Y$ into itself and has a fixed point (Brouwer's fixed point theorem), a zero of $F$ since $\mathcal Y$ is nonsingular. A zero of $F$ in $Y$ is a native gauge of the ray (Definition~\hb{def:raygauge}, the defects summing to $0$ as $f$ has degree $0$, so that the defect at $e_0$ vanishes with the others). The Jacobian of $F$ in the $\alpha_e$ is the Hessian of $\log f$ in the $y_e$ times $\operatorname{diag}(\alpha_e^{-1})$, so an enclosure of the Hessian determinant that excludes $0$ makes every native gauge in $Y$ nondegenerate. Disjoint boxes hold distinct points, and distinct values of $f$ at native gauges are distinct native values.
\end{proof}

\section{Chapter~\hb{ch:rayinputs}: inputs along a ray}\label{sec:chrayinputs}

In this section $\vartheta=y\,d/dy$ on $(0,\infty)$, and for a function $F$ on $(0,\infty)$, $\hat F(s)=\int_0^\infty y^{s-1}F(y)\,dy$ wherever it converges. Two facts about this Mellin read are used throughout.

\begin{slemma}\label{lem:ri-mellin}
Let $F$ be smooth on $(0,\infty)$ with every $\vartheta^aF$ of size $O\big((1+|\log y|)^N\big)$ as $y\to0$ and $O\big(y^Ce^{-b\sqrt y}\big)$ as $y\to\infty$, for some $N,C$ and $b>0$. Then $\hat F$ is holomorphic on $\re s>0$, and:
\begin{enumerate}[label=(\roman*)]
\item for every polynomial $r$ and $k\in\mathbb N$, $\int_0^\infty y^{s-1}\,y^kr(\vartheta)F\,dy=r(-s-k)\,\hat F(s+k)$ on $\re s>0$;
\item if $\hat F$ vanishes on a vertical line in $\re s>0$, then $F=0$.
\end{enumerate}
\end{slemma}

\begin{proof}
The bounds give absolute convergence, locally uniform on $\re s>0$, of every integral involved. (i) For $G=\vartheta^aF$ and $\re a'>0$, $\int_0^\infty y^{a'-1}\vartheta G\,dy=\int_0^\infty y^{a'}G'\,dy=\big[y^{a'}G\big]_0^\infty-a'\int_0^\infty y^{a'-1}G\,dy$, and the bracket vanishes at both ends by the bounds. So $\vartheta$ acts as $-a'$ under $\int_0^\infty y^{a'-1}(\cdot)\,dy$, and with $a'=s+k$ this gives the claim for every monomial $r(\vartheta)=\vartheta^a$, hence for every polynomial. (ii) With $y=e^v$ and $s=c+i\tau$, $\hat F(c+i\tau)=\int_{\mathbb R}e^{i\tau v}\,e^{cv}F(e^v)\,dv$ is the Fourier transform of the integrable continuous function $e^{cv}F(e^v)$; it vanishes for every $\tau$, so that function vanishes (the uniqueness of the Fourier transform on $L^1(\mathbb R)$), and $F=0$.
\end{proof}

\begin{hdefinition}[\hb{def:ri-inputs} (Recurrences; the values input and the slope input)]
As stated.
\end{hdefinition}

\begin{proof}
The objects are well defined: $L$, $L^\ast$ and $L'$ act on functions on a coset of $\mathbb Z$, or on meromorphic functions, the coefficients being polynomials; the symbol is that of Corollary~\hb{cor:symbol} for the coefficients $c_i$; the two inputs are inputs in the sense of Definition~\hb{def:inputs}, a lattice input with the densities $p$ and $\tilde p$ at the points of the coset beyond the cut, where both are finite because $p$ is holomorphic there.
\end{proof}

\begin{hproposition}[\hb{prop:ri-rayinput} (The values of a ray as an input)]
As stated.
\end{hproposition}

\begin{proof}
(i) By Theorem~\hb{thm:pencil}(i), $0<f\le C'e^{-c'\|y\|}$ on the chart, so $\int_{y_{e_0}=0}|f^k|\,dy=\int f^{\re k}\,dy$ converges for $\re k>0$, locally uniformly, and $\langle1\rangle_k$ is holomorphic there. By Proposition~\hb{prop:principal} it continues to a meromorphic function $p$ on $\mathbb C$, with $\lim_{k\to-j}(k+j)^{n-1}p(k)=V_1\,\mathsf C_j(1)$ for every $j\ge j(1)=0$. The class $m=\sum_ic_i(k)\operatorname{Sh}^i[1]$ of the module (Definition~\hb{def:raymodule}) has the values $\langle m\rangle_k=\sum_ic_i(k)\langle f^i\rangle_k=\sum_ic_i(k)\langle1\rangle_{k+i}$, which vanish for all large integers $k$; by Theorem~\hb{thm:onekernel}, whose hypotheses are those of this statement, the continuation of $\langle m\rangle_k$ vanishes identically, and that continuation is $(Lp)(k)$. Multiply $(Lp)(k)=0$ by $(k+j)^{n-1}$ and let $k\to-j$: for $i\le j$, $(k+j)^{n-1}p(k+i)\to V_1\,\mathsf C_{j-i}(1)$, and for $i>j$, $p(k+i)$ is holomorphic at $k=-j$, since $-j+i\ge1$, so its product with $(k+j)^{n-1}$ tends to $0$, $n\ge2$. So $V_1\sum_{i\le\min(h,j)}c_i(-j)\,\mathsf C_{j-i}(1)=0$, and $V_1>0$.

(ii) Differentiating $\sum_ic_i(x)p(x+i)=0$ gives $\sum_ic_i'(x)p(x+i)+\sum_ic_i(x)p'(x+i)=0$, that is $L\tilde p=L'p$.

(iii) By Proposition~\hb{prop:nativeorder}, $\langle1\rangle_k>0$ for large $k$ and $\langle1\rangle_{k+1}/\langle1\rangle_k\to f(\alpha^\star)$, so the root rate of $p$ along the integers is $f(\alpha^\star)$. For real $k\ge1$, $\tilde p(k)=\int f^k\log(1/f)\,dy$ by differentiation under the integral, which the bound of (i) allows, and $|\tilde p(k)|\le f(\alpha^\star)^{k-1}\int f\,|\log f|\,dy$ because $0<f\le f(\alpha^\star)$ on the chart (Theorem~\hb{thm:pencil}(i)). The last integral converges: $\log f$ is a combination of the coordinates $y_e$ and of logarithms of the letters, polynomials with nonnegative coefficients in $e^{y}$, so $|\log f|\le C(1+\|y\|)$, while $f\le C'e^{-c'\|y\|}$. So the root rate of $\tilde p$ is at most $f(\alpha^\star)$.
\end{proof}

\begin{hproposition}[\hb{prop:ri-sbp} (Summation by parts at the cut)]
As stated.
\end{hproposition}

\begin{proof}
For a function $\varphi$ on the integers $x>-h$ with $\sum_x|R(x)|\,|\varphi(x)|$ and $\sum_x|R(x)\,c_i(x)\,\varphi(x+i)|$ finite over $x\ge1$, which holds for $\varphi\in\{p,\tilde p\}$ since $R$ and the $c_i$ are polynomials and $p$, $\tilde p$ have root rate below one,
\[
\sum_{x\ge1}(L^\ast R)(x)\,\varphi(x)=\sum_{i=0}^h\sum_{x\ge1}c_i(x-i)R(x-i)\,\varphi(x)=\sum_{i=0}^h\sum_{x'\ge1-i}c_i(x')R(x')\,\varphi(x'+i),
\]
with $x'=x-i$. The terms with $x'\ge1$ add up to $\sum_{x'\ge1}R(x')\,(L\varphi)(x')$, and those with $1-i\le x'\le0$, written $x'=-j$ with $0\le j<i$, to $\sum_{j=0}^{h-1}R(-j)\,B_j(\varphi)$, $B_j(\varphi)=\sum_{i=j+1}^hc_i(-j)\,\varphi(i-j)$. So
\[
\Tx0{(L^\ast R)\,\varphi\,\IZ}0=\sum_{x\ge1}R(x)\,(L\varphi)(x)+\sum_{j=0}^{h-1}R(-j)\,B_j(\varphi),
\]
the Trailmoment at index zero on the counting input being the sum over the integers beyond the cut (Definition~\hb{def:ratword}). For $\varphi=p$, $Lp=0$ and $B_j(p)=\sum_{i>j}c_i(x)p(x+i)\big|_{x=-j}=-\beta_j(-j)$, which is the first identity. For $\varphi=\tilde p$, $L\tilde p=L'p$ (Proposition~\hb{prop:ri-rayinput}(ii), whose proof uses only $Lp=0$), and the same rearrangement with $L'$ in place of $L$ gives $\sum_{x\ge1}R(x)(L'p)(x)=\Tx0{(L'^\ast R)\,p\,\IZ}0-\sum_jR(-j)B_j'(p)$ with $B_j'(p)=\sum_{i>j}c_i'(-j)\,p(i-j)$. Hence the second identity holds with the boundary $B_j(\tilde p)-B'_j(p)=\sum_{i>j}\big(c_i(x)\tilde p(x+i)-c_i'(x)p(x+i)\big)\big|_{x=-j}=-\frac d{dx}\sum_{i>j}c_i(x)p(x+i)\big|_{x=-j}=\beta_j'(-j)$.

If every $c_i$ with $i>j\ge1$ vanishes to order two at $-j$, then near $x=-j$ each term of $\beta_j=-\sum_{i>j}c_i(x)p(x+i)$ is a polynomial with a double zero at $-j$ times a function holomorphic there, so $\beta_j(-j)=\beta_j'(-j)=0$ for $j\ge1$.
\end{proof}

\begin{hproposition}[\hb{prop:ri-block} (The fixed block)]
As stated.
\end{hproposition}

\begin{proof}
(i) For $R$ of degree $d$ with leading coefficient $\rho\neq0$, $c_i(x-i)R(x-i)$ has degree at most $D+d$ and coefficient $[x^D]c_i\,\rho$ at $x^{D+d}$, so $L^\ast R$ has degree at most $D+d$ and coefficient $\chi(1)\rho\neq0$ there. So $L^\ast$ is injective on $\Bbbk[x]$, and for every $d$ it maps the polynomials of degree at most $d$, a space of dimension $d+1$, into those of degree at most $D+d$, meeting the polynomials of degree below $D$ only in $0$. The dimensions add up, $(d+1)+D=D+d+1$, so the polynomials of degree at most $D+d$ are the direct sum of the two, for every $d$; this is the decomposition of $\Bbbk[x]$.

(ii) Write $W=L^\ast B+r$ with $\deg r<D$ by (i), and $L'^\ast B=L^\ast B_2+r_2$ in the same way; $B$, $r$, $B_2$, $r_2$ have coefficients in $\Bbbk$. By Proposition~\hb{prop:ri-sbp},
\begin{align*}
\Tx0{W\,p\,\IZ}0&=\Tx0{r\,p\,\IZ}0-\sum_jB(-j)\,\beta_j(-j),\\
\Tx0{W\,\tilde p\,\IZ}0&=\Tx0{r\,\tilde p\,\IZ}0+\Tx0{r_2\,p\,\IZ}0\\
&\qquad-\sum_jB_2(-j)\,\beta_j(-j)+\sum_jB(-j)\,\beta_j'(-j).
\end{align*}
A polynomial of degree below $D$ is a $\Bbbk$-combination of the $(x-1)^a$, $a<D$, and the Trailmoment at index zero is linear in the density, which gives the claim.
\end{proof}

\begin{hcorollary}[\hb{cor:ri-openmark} (A mark opened at an integer strength)]
As stated.
\end{hcorollary}

\begin{proof}
The bunch $B=\{e\}$ has width one, $c(\{e\})=r(\{e\})+r(R)-r(\cE)=1$ since $e$ is neither a loop nor a coloop, and the unit weights factor across $B$ and $R$. Its minors on $B$ are $\cB|B$, with the one basis $\{e\}$, and $\cB/R$, with the one basis $\varnothing$; so the bunch family of Theorem~\hb{thm:twosides} has the period $1$, and the hinge ratio of $B$ is $\varrho_B=\mathbf E_{\cB/R}/\mathbf E_{\cB|B}=\alpha_e/\alpha_e=1$ (Definition~\hb{def:readings}). The room family is the deletion $\cB\setminus e=\cB|R$ at the level $d$ and the contraction $\cB/e=\cB/B$ at the level $r$, and $\varrho_R=\mathbf E_{\cB/e}/\mathbf E_{\cB\setminus e}=X_//X_\setminus$, since $\mathbf E=K\Psi=\alpha^RK(\alpha)K(\alpha^{-1})$ for each of them. At the centre $(s\one;s,s)$ the wall of $B$ is $\kappa_B=s\,c(B)=s$ (Proposition~\hb{prop:centrefamily}(ii)), so the odd share is fixed at $2\delta^d=s-\kappa_B=0$ and the level shares are $(a^d,b^d,a^r,b^r)=(d,d,s-d,s-d)$, $d$ the even share. The one-table display of Theorem~\hb{thm:twosides} reads
\[
\Gamma(s)^2P_{\cB}(s\one;s,s)=\frac1{2\pi i}\int\Gamma(d)^2\,\Gamma(s-d)^2\,\Phi_s(d)\,dd,\qquad \Phi_s(d)=\int_{y_{e_0}=0}X_\setminus^{-d}X_/^{-(s-d)}\,dy,
\]
the room period at the masses $s\one_R$, whose integrand is $\alpha^{s\one_R}K_\setminus^{-d}\Psi_\setminus^{-d}K_/^{-(s-d)}\Psi_/^{-(s-d)}=X_\setminus^{-d}X_/^{-(s-d)}$.

When $X_\flat=X_\setminus$, $\sup\varrho_B=1<\inf\varrho_R$, and Corollary~\hb{cor:closetwosides}(i) gives $P_{\cB}(s\one;s,s)=\sum\operatorname{Res}$ over the poles of $\Gamma(d)^2$ at $d=-x$, $x\in\mathbb N$, where the two lattices of the level $d$ coincide, every room period there converging. Then $\Phi_s(-x)=q_s(x)$. When $X_\flat=X_/$, $\inf\varrho_B>\sup\varrho_R$, and the same corollary closes on the level $r$ with the opposite sign; the substitution $d=s-\tau$ exchanges the two levels and the two minors and turns that closing into the closing at $\tau=-x$, with $\Phi_s(s-\tau)=\int X_/^{-\tau}X_\setminus^{-(s-\tau)}\,dy$ and $q_s(x)=\Phi_s(s+x)$. In both cases, writing $\Phi^\flat(\tau)=\int X_\flat^{-\tau}X_\sharp^{-(s-\tau)}\,dy$, so that $q_s(x)=\Phi^\flat(-x)$, the room periods $\Phi^\flat(-x)$, $x\in\mathbb N$, converge. The walls of the room family with the minor of $X_\flat$ at the share $-x$ and that of $X_\sharp$ at $s+x$, at the masses $s\one_R$, are $\kappa_S=-x\,c_\flat(S)+(s+x)\,c_\sharp(S)$ with the widths $c_\flat$, $c_\sharp$ of the two minors, by the computation of Proposition~\hb{prop:centrefamily}(ii), which is linear in the strengths; they are affine in $x$ and positive at every point of $\mathbb N$ (Theorem~\hb{thm:projective}), hence on $[0,\infty)$ and on a neighbourhood of it, so $q_s$ converges and is holomorphic there. The closing reads
\[
\Gamma(s)^2P_{\cB}(s\one;s,s)=\sum_{x\in\mathbb N}\operatorname*{Res}_{\tau=-x}\Big(\Gamma(\tau)^2\,\Gamma(s-\tau)^2\,\Phi^\flat(\tau)\Big).
\]
At the integer $s$, $\Gamma(s-\tau)=(s-1-\tau)\cdots(1-\tau)\,\Gamma(1-\tau)$, so $\Gamma(\tau)^2\Gamma(s-\tau)^2=V_s(-\tau)\,\Gst40{\tau,\ \tau,\ 1-\tau,\ 1-\tau}{}$, and by Lemma~\hb{lem:resonantpair} the last factor has at $\tau=-x$ the principal part $(\tau+x)^{-2}$ and no simple part. The residue is therefore the derivative of $V_s(-\tau)\Phi^\flat(\tau)$ at $\tau=-x$,
\[
-V_s'(x)\,q_s(x)+V_s(x)\,\Phi^{\flat\prime}(-x)=V_s(x)\,\tilde q_s(x)-V_s'(x)\,q_s(x),
\]
since $\Phi^{\flat\prime}(-x)=-q_s'(x)$. The ratio $X_\flat/X_\sharp$ is at most $\rho$, so $0<q_s(x)\le\rho^{x}q_s(0)$ and, by differentiation under the integral, $|\tilde q_s(x)|\le\rho^{x-1}\int\log(X_\sharp/X_\flat)\,(X_\flat/X_\sharp)\,X_\sharp^{-s}\,dy$ for $x\ge1$, the last integral being $\tilde q_s(1)$, finite by the holomorphy at $1$; with $V_s$ a polynomial the terms have root rate at most $\rho$.
\end{proof}

\begin{hproposition}[\hb{prop:ri-boundary} (Boundary relations)]
As stated.
\end{hproposition}

\begin{proof}
The identity $Lp=0$ reads $\beta_{x_0,j}(x)=-\sum_{i>j}c_i(x)\,p(x+i)$ near $x_0$, and the right side is holomorphic at $x_0$ by the hypothesis; so is $\beta_{x_0,j}$, and the first relation is this identity at $x_0$. Differentiating it at $x_0$ gives $\beta'_{x_0,j}(x_0)=-\sum_{i>j}\big(c_i'(x_0)p(x_0+i)+c_i(x_0)p'(x_0+i)\big)$, which is the second relation since $\tilde p=-p'$. $\beta_{x_0,j}=\sum_{i\le j}c_i\,p(\cdot+i)$ is a finite sum of products of the Taylor expansions of the $c_i$ at $x_0$ with the Laurent expansions of $p$ at $x_0+i$; its value and its derivative at $x_0$ are the coefficients of $(x-x_0)^0$ and $(x-x_0)^1$ of that sum, finitely many products of these coefficients.

When $c_h$ has no zero at the integers $x\ge1$, $L\tilde p=L'p$ at $x$ expresses $\tilde p(x+h)$ through $\tilde p(x),\dots,\tilde p(x+h-1)$ and values of $p$, so by induction every $\tilde p(x)$, $x\ge1$, is an affine function of $\tilde p(1),\dots,\tilde p(h)$; for an integer $x_0$ with $x_0+j\ge0$ the slopes in the second relation are at the integers $x_0+i\ge1$, $i>j$, and the relation is affine in the initial slopes.

At the cut of the values of a ray, Proposition~\hb{prop:ri-rayinput}(i) at $j=0$ gives $p(x)=V_1x^{-(n-1)}(1+O(x))$, $\mathsf C_0(1)=1$, so $\mathsf L_0=V_1$, and $p$ has a pole of order exactly $n-1$ at $0$. $\beta_{0,0}=c_0\,p$ is holomorphic at $0$, so $c_0$ vanishes to order at least $n-1$ there, and $c_0(x)p(x)=\gamma_{n-1}\mathsf L_0+(\gamma_{n-1}\mathsf L_1+\gamma_n\mathsf L_0)\,x+O(x^2)$. The two relations at $x_0=0$, $j=0$ are therefore the displayed ones; the first is the low-end relation of Corollary~\hb{cor:lowendrelation}. When every table has unit weights, the values are $\langle1\rangle_k=P_F(kq)$ for the question $q$ of the direction (Definition~\hb{def:question}), the direction being a point of the balance wall with real data, and Theorem~\hb{thm:severity}(i) gives $k^{n-1}P_F(kq)=V_1+o(k)$ as $k\downarrow0$, every $\log\theta_{j,T}$ being $0$; since $k^{n-1}p(k)=\mathsf L_0+\mathsf L_1k+O(k^2)$ is holomorphic at $0$, $\mathsf L_1=0$.
\end{proof}

\begin{hproposition}[\hb{prop:ri-minimal} (Minimality)]
As stated.
\end{hproposition}

\begin{proof}
(i) Let $D=\deg c_h$. Divided by $c_h(x)$, which has no zero at the integers $x\ge1$, the recurrence reads $u(x+h)+\sum_{i<h}a_i(x)\,u(x+i)=0$ with $a_i=c_i/c_h\to[x^D]c_i/[x^D]c_h$ as $x\to\infty$, and $a_0(x)\neq0$ at these integers. The limit characteristic polynomial is $\chi(\lambda)/[x^D]c_h$, of degree $h$, with $\chi(0)=[x^D]c_0\neq0$ and roots $\lambda_c$ of distinct moduli. Perron's theorem on such recurrences (coefficients with limits, $a_0$ without zeros on the range, roots of the limit polynomial of distinct moduli) gives a basis $u_1,\dots,u_h$ of the solutions on the integers $x\ge1$ with $u_c(x+1)/u_c(x)\to\lambda_c$, so $|u_c(x)|^{1/x}\to|\lambda_c|$. Let $u=\sum_cb_cu_c\neq0$, and let $d$ be the index with the largest $|\lambda_d|$ among the $c$ with $b_c\neq0$. For the other such $c$, $u_c(x)/u_d(x)\to0$, the root rates being smaller; so $u(x)=b_du_d(x)\,(1+o(1))$, $u(x+1)/u(x)\to\lambda_d$, and $u$ has the root rate $|\lambda_d|$. Hence $u$ has root rate at most $\lambda$ exactly when $b_c=0$ for every $c$ with $|\lambda_c|>\lambda$, and $U_\lambda$ is spanned by the $u_c$ with $|\lambda_c|\le\lambda$, which are independent.

(ii) $\tilde p$ satisfies $L\tilde p=L'p$ (Proposition~\hb{prop:ri-rayinput}(ii), whose proof uses only $Lp=0$) and has root rate at most $\lambda$; a $\varphi$ with the same two properties differs from it by $u=\varphi-\tilde p$ with $Lu=0$ and root rate at most $\lambda$, the root rate of a difference being at most the larger of the two; conversely every $\tilde p+u$, $u\in U_\lambda$, has both properties. The same holds for $L\varphi=0$ with $p$. A second relation of Proposition~\hb{prop:ri-boundary} at an integer $x_0$ with $x_0+j\le0$ holds for $\tilde p+u$ exactly when its left side vanishes on $u$, since its right side is read off the Laurent data of $p$ at $x_0,\dots,x_0+j\le0$ and involves no slope; relations at other points are consequences of $L\tilde p=L'p$ and single out nothing; so if these left sides vanish together on no nonzero $u\in U_\lambda$, $\tilde p$ is the only solution of root rate at most $\lambda$ that satisfies the relations. The first relations single out $p$ in the same way.
\end{proof}

\begin{hproposition}[\hb{prop:ri-pairing} (The two ends paired)]
As stated.
\end{hproposition}

\begin{proof}
(i) For $\sigma<0$ every term vanishes. For $\sigma\ge0$, $c_i(\sigma-i)=\epsilon\,c_i(-\sigma)$ by the hypothesis read at $x=-\sigma$, so the sum is $\epsilon\sum_{i=0}^{\min(h,\sigma)}c_i(-\sigma)\,\mathsf C_{\sigma-i}$, which vanishes by Proposition~\hb{prop:ri-rayinput}(i) at $j=\sigma$.

(ii) Write $p(x)=\langle1\rangle_x$, and $(L^\ast\mathsf C)(\sigma)=\sum_ic_i(\sigma-i)\mathsf C_{\sigma-i}$. Shifting $j$ by one,
\[
B(\sigma+1)=\sum_{i=1}^h\sum_{j=0}^{i-1}c_i(\sigma-j)\,\mathsf C_{\sigma-j}\,p(\sigma-j+i),
\]
and the difference with $B(\sigma)$ keeps the term $j=0$ of the first sum and the term $j=i$ of the second:
\[
B(\sigma+1)-B(\sigma)=\sum_{i=1}^h\big(c_i(\sigma)\,\mathsf C_\sigma\,p(\sigma+i)-c_i(\sigma-i)\,\mathsf C_{\sigma-i}\,p(\sigma)\big)=\mathsf C_\sigma\,(Lp)(\sigma)-p(\sigma)\,(L^\ast\mathsf C)(\sigma),
\]
the terms $i=0$ cancelling. For integers $\sigma\ge1$ both vanish, by Proposition~\hb{prop:ri-rayinput}(i) and by (i). At $\sigma=1$ only $j=1$ has $\mathsf C_{1-j}\neq0$, with $\mathsf C_0=1$, so $B(1)=\sum_{i\ge1}c_i(0)\,p(i)$, which is $-\gamma_{n-1}V_1$ by Proposition~\hb{prop:ri-boundary}.

(iii) By Proposition~\hb{prop:ri-boundary}, $c_0$ vanishes to order at least $n-1$ at $0$; with $\deg c_0\le n-1$ and $c_0\neq0$, $c_0=\gamma_{n-1}x^{n-1}$, $\gamma_{n-1}\neq0$, and the coefficients of $x^{n-1}$ give $\chi(\lambda)=\sum_i[x^{n-1}]c_i\lambda^i$ with $\chi(0)=\gamma_{n-1}$. On the centre ray of a family, $\alpha^\star=\one$ and $\langle1\rangle_k=P_F(kq)$ for the centre question $q$; Theorem~\hb{thm:severity}(ii) at the centre, with $-\nabla^2X(0)=4H$ on the chart in the coordinates $z=y/2$ ($H$ as in Theorem~\hb{thm:branch}) and $e^{X(0)}=f(\one)$, gives
\[
\langle1\rangle_k=f(\one)^k\Big(\frac{2\pi}k\Big)^{\frac{n-1}2}(\det H)^{-1/2}\big(1+O(k^{-1})\big),
\]
the factor $2^{n-1}$ of that theorem cancelling against $\det(4H)^{1/2}=2^{n-1}(\det H)^{1/2}$. Theorem~\hb{thm:onekernel} at $g=1$ gives $j^{(n-1)/2}f(\one)^j\mathsf C_j(1)\to(2\pi)^{-(n-1)}a_0(1)$, where $a_0(1)=(2\pi)^{(n-1)/2}(dy/dv)(\one)=(2\pi)^{(n-1)/2}(\det H)^{-1/2}$ by Theorem~\hb{thm:gaugesolution}(i), the coordinates $v$ having $\frac12\sum v_i^2=\frac12y^{\mathsf T}Hy+O(|y|^3)$. So for fixed $i$, as $a\to\infty$, $\mathsf C_a\,\langle1\rangle_{a+i}=\lambda_\star^i(\det H)^{-1}a^{-(n-1)}(1+o(1))$, with $\lambda_\star=f(\one)$, while $c_i(a)=[x^{n-1}]c_i\,a^{n-1}+O(a^{n-2})$. Every term of $B(\sigma)$ with $a=\sigma-j$ therefore tends to $[x^{n-1}]c_i\,\lambda_\star^i(\det H)^{-1}$, and
\[
\lim_{\sigma\to\infty}B(\sigma)=(\det H)^{-1}\sum_{i=1}^hi\,[x^{n-1}]c_i\,\lambda_\star^i=(\det H)^{-1}\lambda_\star\chi'(\lambda_\star).
\]
By (ii) the limit is $-\chi(0)V_1$, so $V_1\det H=-\lambda_\star\chi'(\lambda_\star)/\chi(0)$, which is positive. The symbol vanishes at $\lambda_\star$: $\langle1\rangle_{k+1}/\langle1\rangle_k\to\lambda_\star$ (Proposition~\hb{prop:nativeorder}, a limit that needs no hypothesis on the coefficients of $L$), so dividing $\sum_ic_i(k)\langle1\rangle_{k+i}=0$ by $k^{n-1}\langle1\rangle_k$ and letting $k\to\infty$ gives $\chi(\lambda_\star)=0$; writing $\chi(\lambda)=\chi(0)\prod_c(1-\lambda/\lambda_c)$ over its roots with multiplicity, a multiple root at $\lambda_\star$ would give $\chi'(\lambda_\star)=0$, so $\lambda_\star$ is simple, and $\lambda_\star\chi'(\lambda_\star)=-\chi(0)\prod_{c\neq\star}(1-\lambda_\star/\lambda_c)$.
\end{proof}

The rest of this section reads the list $\binom{\cE}1$. There $\mathsf g(y)=\sG_{0;(1,-1)}(1,y)=\int_0^\infty x^{-1}e^{-x-y/x}\,dx$ and $\mathsf h(y)=\sum_{m\ge0}y^m/m!^2$. The configuration $(1,-1)$ on one scale has $\cone(t)=\mathbb R$, so $\sG_{\nu;(1,-1)}(1,y)$ is entire in $\nu$ (Theorem~\hb{thm:grammadomain}); its ladder and its exchange relation (Theorem~\hb{thm:grammalaws}(i),(ii)) read $\partial_y\sG_{\nu;(1,-1)}(1,y)=-\sG_{\nu-1;(1,-1)}(1,y)$ and $\nu\,\sG_{\nu;(1,-1)}(1,y)=\sG_{\nu+1;(1,-1)}(1,y)-y\,\sG_{\nu-1;(1,-1)}(1,y)$. So $\vartheta\mathsf g=-y\,\sG_{-1;(1,-1)}(1,y)$, the exchange relation at $\nu=-1$ gives $y\,\sG_{-2;(1,-1)}(1,y)=\mathsf g+\sG_{-1;(1,-1)}(1,y)$, and
\[
\vartheta^2\mathsf g=-y\,\sG_{-1;(1,-1)}(1,y)+y^2\sG_{-2;(1,-1)}(1,y)=y\,\mathsf g ;
\]
and $\vartheta^2\mathsf h=\sum_{m\ge1}y^m/(m-1)!^2=y\,\mathsf h$. By Lemma~\hb{lem:onemark}(i) at $l=0$, $\mathsf g=-\mathsf h\log y+\tilde{\mathsf g}$ with the entire function $\tilde{\mathsf g}(y)=\sum_m2(\log\Gamma)'(1+m)\,y^m/m!^2$. The solutions of $\vartheta^2u=yu$ on $(0,\infty)$ form a space $\mathcal V$ of dimension two with the basis $\mathsf g,\mathsf h$, the two being independent since $\mathsf g/\mathsf h$ is not constant; the $n+1$ products $\mathsf g^{n-j}\mathsf h^j$ are independent as well, $\mathsf h^n$ times the powers of the nonconstant function $\mathsf g/\mathsf h$.

\begin{hproposition}[\hb{prop:ri-signs} (The native gauges of $\binom{\cE}1$)]
As stated.
\end{hproposition}

\begin{proof}
(i) The occupation of $e$ is $R_e=\alpha_e\partial_{\alpha_e}\log K=u_e$, and since $\Psi(\alpha)=\alpha^{\cE}K(\alpha^{-1})$ the dual occupation is $R^\ast_e=\alpha_e\partial_{\alpha_e}\log\Psi=1-v_e$, $v_e=\alpha_e^{-1}/K(\alpha^{-1})$. The defect of the direction (Definition~\hb{def:ray}) is $c^{(1)}_e=m-s\,u_e-t\,(1-v_e)=\frac{s-t}n-s\,u_e+t\,v_e$, by the balance $m-t=\frac{s-t}n$. On $X_F$ the $\alpha_e$, $K(\alpha)$ and $K(\alpha^{-1})$ are nonzero, so $u_e,v_e\neq0$, $u_ev_e=\lambda(\alpha)\neq0$ for every $e$, and $\sum_eu_e=\sum_ev_e=1$.

Let $\alpha$ be a native gauge, $c^{(1)}_e=0$ for every $e$. With $v_e=\lambda/u_e$, $s\,u_e^2-\frac{s-t}nu_e-t\lambda=0$ for every $e$: the occupations are roots of one quadratic, whose roots $u_+,u_-$ have $u_++u_-=\frac{s-t}{ns}$ and $u_+u_-=-t\lambda/s\neq0$. Let $u_+$ be taken on $n-j$ marks and $u_-$ on $j$ marks, $u_-$ on none when the occupations are all equal; $\sum_eu_e=1$ gives $(n-j)u_++ju_-=1$, and with $u_-=\frac{s-t}{ns}-u_+$, $(n-2j)u_+=1-j\frac{s-t}{ns}$. At $2j=n$ this would read $1=\frac{s-t}{2s}$, that is $s=-t$, which is excluded. Conversely let the occupations of $\alpha\in X_F$ take the value $u_+$ on $n-j$ marks and $u_-$ on $j$ marks with the three relations. Adding $(n-j)u_++ju_-=1$, which follows from them, to $(n-j)u_-+ju_+$ gives $n(u_++u_-)=\frac{s-t}s$, so $(n-j)u_-+ju_+=-\frac ts$ and
\[
\lambda^{-1}=\sum_e\frac{1}{u_e}=\frac{n-j}{u_+}+\frac j{u_-}=\frac{(n-j)u_-+ju_+}{u_+u_-}=-\frac{t}{s\,u_+u_-},
\]
using $\sum_ev_e=1$, $v_e=\lambda/u_e$. Then $s\,u_e-t\,v_e=s\,u_e+s\,u_+u_-/u_e=s(u_++u_-)=\frac{s-t}n$ at both values, so every defect vanishes. Since $\alpha_e=u_eK(\alpha)$, such a point is $(u_e)_e$ up to scale; conversely, for given $u_\pm$ with the three relations and a set of $j$ marks, the point $u=(u_e)_e$ has $K(u)=1$ and $\Psi(u)=\prod_eu_e\sum_eu_e^{-1}=-\prod_eu_e\,t/(s\,u_+u_-)\neq0$, so it lies in $X_F$. The relations for the complement of the set, its $n-j$ marks at the value $u_+$ in the role of $u_-$ and the other $j$ at $u_-$ in the role of $u_+$, hold as well: the second is symmetric, and $(n-j)u_++ju_-=1$ with $u_+=\frac{s-t}{ns}-u_-$ gives $(2j-n)u_-=1-(n-j)\frac{s-t}{ns}$; so a set and its complement give the same point. The stride $f=\alpha^{m\one}K^{-s}\Psi^{-t}=\alpha^{(m-t)\one}K(\alpha)^{-s}K(\alpha^{-1})^{-t}$ is homogeneous of degree $0$, and at $u$, where $K(u)=1$ and $K(u^{-1})=\lambda^{-1}$, it is $\big(\prod_eu_e\big)^{m-t}\lambda^t$.

(ii) At $s=t=\sigma$ the balance gives $m=\sigma$, the relations give $u_-=-u_+$ and $u_+=1/(n-2j)$, so $u$ is $1/(n-2j)$ times the sign $\eta$ with $\eta_e=-1$ exactly on the $j$ marks; on the chart $\alpha_{e_0}=1$ the native gauge is the sign $\eta_{e_0}\eta$, with $K\neq0$. Conversely every sign with $K(\eta)\neq0$ is a native gauge (Theorem~\hb{thm:signgauges}(ii)). The native value is $\lambda^\sigma=u_+^{2\sigma}=(n-2j)^{-2\sigma}=K(\eta)^{-2\sigma}$, and $|n-2j|$ for $0\le j\le n$, $2j\neq n$, takes the $\lceil n/2\rceil$ values $n,n-2,\dots$ down to $1$ or $2$. By Theorem~\hb{thm:signgauges}(ii) the Hessian is $2\sigma\big(\kappa_\eta(X_e,X_{e'})\big)_{e,e'\neq e_0}$, the covariances of the indicators under $\mathbb P_\eta(\{e\})=\pi_e=\eta_e/K(\eta)$. The indicator $X_e$ takes the value $1$ at the basis $\{e\}$ and $0$ at the others, so $\kappa_\eta(X_e,X_{e'})=\delta_{ee'}\pi_e-\pi_e\pi_{e'}$. The determinant of a diagonal matrix less a matrix of rank one, $\det\big(\operatorname{diag}(\pi_e)-\pi\pi^{\mathsf T}\big)_{e,e'\neq e_0}=\prod_{e\neq e_0}\pi_e\big(1-\sum_{e\neq e_0}\pi_e\big)$, is $\prod_e\pi_e$ since $\sum_e\pi_e=1$, and the Hessian has the determinant $(2\sigma)^{n-1}\prod_e\eta_e/K(\eta)^n\neq0$.
\end{proof}

\begin{hproposition}[\hb{prop:ri-parallel} (A parallel class enters through its list)]
As stated.
\end{hproposition}

\begin{proof}
(i) Theorem~\hb{thm:block} with every mark replaced by the first rule gives $K_{\cB}=K_{\mathcal H}\big((e_1(\alpha_{P_i}))_i\big)=K_{\mathcal H}(A)$, the weights of a list being one. Write $K_{\mathcal H}(A)^k=\sum_w\big([A^w]K_{\mathcal H}^k\big)A^w$ and $A^w=\prod_i\sum_{v_i\in\mathbb N^{P_i},\ |v_i|=w_i}\binom{w_i}{v_i}\alpha^{v_i}$. The classes are disjoint, so a monomial $\alpha^v$ determines $w_i=v_{P_i}$, and it occurs only in the term of that $w$, with the coefficient $\big([A^w]K_{\mathcal H}^k\big)\prod_i\binom{w_i}{v|_{P_i}}$.

(ii) On the centre ray at the strength one, $f^{-1}=K_{\cB}(\alpha)K_{\cB}(\alpha^{-1})$ (Corollary~\hb{cor:reflected}), and $\mathsf C^{\cB}_k(1)$ is the constant term of $K_{\cB}(\alpha)^kK_{\cB}(\alpha^{-1})^k$ (Definition~\hb{def:closings}, the origin form being $dy$), which is $\sum_v\big([\alpha^v]K_{\cB}^k\big)^2$, the coefficients being real. By (i) this is $\sum_w\big([A^w]K_{\mathcal H}^k\big)^2\prod_i\sum_{|v_i|=w_i}\binom{w_i}{v_i}^2$, and $\sum_{|v_i|=w_i}\binom{w_i}{v_i}^2=\sum_{v_i}\big([\alpha^{v_i}]K_{P_i}^{w_i}\big)^2$ is the constant term of $\big(K_{P_i}(\alpha)K_{P_i}(\alpha^{-1})\big)^{w_i}$, $K_{P_i}=\sum_{e\in P_i}\alpha_e$, the closing of $\binom{P_i}1$; for $|P_i|=1$ it is $1$.

(iii) At a sign, (i) gives $K_{\cB}(\eta)=K_{\mathcal H}\big((\textstyle\sum_{e\in P_i}\eta_e)_i\big)$, and the sum of $s_i$ signs takes the values $s_i,s_i-2,\dots,-s_i$; the native values are those of Theorem~\hb{thm:signgauges}(ii) for one table at the strength $\sigma$.
\end{proof}

\begin{hproposition}[\hb{prop:ri-addmark} (A mark added is a closing of the critical-line kernel)]
As stated.
\end{hproposition}

\begin{proof}
(i) By the display of Section~\hb{sec:ri-centre} at the centre, from Corollary~\hb{cor:cycle} and Proposition~\hb{prop:dualfamily}, $\Gamma(s)^2p_{n+1}(s)$ is, for real $s>0$, the share integral $\int_{a_1+\cdots+a_{n+1}=s}\prod_e\Gamma(a_e)^2$ of Corollary~\hb{cor:cycle}, which converges absolutely, $|\Gamma(c+iy)|$ decaying like $e^{-\pi|y|/2}$. In the coordinates $(a_1,\dots,a_{n-1},d)$ with $d=a_1+\cdots+a_n$, of Jacobian one, Fubini's theorem splits it into the integral over $\re d=c$ of $\Gamma(s-d)^2$ times the share integral of $n$ marks at the total $d$. That inner integral is holomorphic in $d$ on $\re d>0$, by absolute and locally uniform convergence, and equals $\Gamma(d)^2p_n(d)$ at real $d$ by the same display at $n$ marks; so it equals $\Gamma(d)^2p_n(d)$ on $\re d>0$, $p_n$ being holomorphic there by the bound of Theorem~\hb{thm:pencil}(i), as in the proof of Proposition~\hb{prop:ri-rayinput}(i). This gives the formula at real $s$; both sides are holomorphic in $s$ on $\re s>c$, so it holds on $\re s>0$. At $s=1$ and $c=\frac12$, with $d=\frac12+iy$, the kernel $\Gamma(d)^2\Gamma(1-d)^2$ is $2\pi\,p_{\mathrm{crit}}(y)$ (Definition~\hb{def:pairlaw}) and $dd/(2\pi i)=dy/(2\pi)$.

(ii) At the integer $s$, $\Gamma(s-d)=(s-1-d)\cdots(1-d)\,\Gamma(1-d)$, so the integrand of (i) is $\Gst40{d,\ d,\ 1-d,\ 1-d}{}\,G(d)$ with $G=W_s\,p_n$, which has double zeros at $1,\dots,s-1$. On $\re d\ge1$, $|p_n(d)|\le p_n(\re d)\le4^{1-\re d}p_n(1)$, since $0<(K(\alpha)K(\alpha^{-1}))^{-1}\le n^{-2}\le\frac14$ on the chart (Theorem~\hb{thm:pencil}(i)); and $\big|\Gst40{d,\ d,\ 1-d,\ 1-d}{}\big|=\pi^2/|\sin\pi d|^2$ is at most $\pi^2/\cosh^2(\pi\im d)$ on the lines $\re d=N+\frac12$ and at most $Ce^{-2\pi|\im d|}$ for $|\im d|\ge1$. So the integrals over the lines $\re d=N+\frac12$ and over the horizontal segments $|\im d|=Y$ tend to $0$, and moving the line to the right (Theorem~\hb{thm:closing}) gives $-\sum_{m\ge1}\operatorname{Res}_{d=m}$. By Lemma~\hb{lem:resonantpair} the kernel has at $d=m$ the principal part $(d-m)^{-2}$ and no simple part, so the residue is $G'(m)$, zero for $m<s$, and
\[
\Gamma(s)^2p_{n+1}(s)=-\sum_{m\ge1}\big(W_s\,p_n\big)'(m)=\sum_{m\ge1}\big(W_s(m)\,\tilde p_n(m)-W_s'(m)\,p_n(m)\big),
\]
the series converging geometrically; it is the index-zero Trailmoment cut at $0$ (Definition~\hb{def:ratword}). For the second identity, $\int_0^\infty y^{s-1}\mathsf g^n\,dy=\Gamma(s)^2p_n(s)$ at every real $s>0$, and the terms of $\mathsf h$ are positive, so by monotone convergence
\[
\int_0^\infty y^{s-1}\mathsf g^n\mathsf h\,dy=\sum_{k\ge0}\frac{\Gamma(s+k)^2}{k!^2}\,p_n(s+k)=\sum_{k\ge0}W_s(s+k)\,p_n(s+k)=\sum_{m\ge1}W_s(m)\,p_n(m),
\]
since $\Gamma(s+k)/k!=(k+1)\cdots(k+s-1)$ and $W_s$ vanishes at $1,\dots,s-1$.

(iii) The share integral of (i) is symmetric in the $n+1$ shares, the fibre integral not depending on the coordinates in which it is written (Definition~\hb{def:fibre}), and $\sum_ea_e=s$ on the fibre; so for each $e$, $\int a_e\prod_{e'}\Gamma(a_{e'})^2=\frac s{n+1}\Gamma(s)^2p_{n+1}(s)$. For $e$ the last mark, $a_e=s-d$ in the coordinates of (i), and at $s=1$ the closing of (ii), with $G(d)=(1-d)p_n(d)$, gives
\[
\frac1{n+1}\,p_{n+1}(1)=-\sum_{m\ge1}\big((1-d)p_n(d)\big)'\big|_{d=m}=\sum_{m\ge1}\big(p_n(m)-(m-1)\,\tilde p_n(m)\big).
\]
With $p_{n+1}(1)=\Tx0{\tilde p_n\,\IZ}0$ from (ii) this is the stated relation.
\end{proof}

The following bounds give the Mellin reads of the products of $\mathsf g$ and $\mathsf h$ the hypotheses of Lemma~\hb{lem:ri-mellin} and fix their growth along the integers.

\begin{slemma}\label{lem:ri-growth}
\begin{enumerate}[label=(\roman*)]
\item $\mathsf g>0$ and $\mathsf h\ge1$ on $(0,\infty)$, and for $y\ge1$
\[
\frac1{2e}\,y^{-1/2}e^{-2\sqrt y}\le\mathsf g(y)\le\sqrt\pi\,y^{-1/4}e^{-2\sqrt y},\qquad e^{-4}y^{-1}e^{2\sqrt y}\le\mathsf h(y)\le e^{2\sqrt y}.
\]
\item For integers $0\le j<n/2$, every $\vartheta^a\big(\mathsf g^{n-j}\mathsf h^j\big)$ is $O\big((1+|\log y|)^n\big)$ as $y\to0$ and $O\big(y^Ce^{-2(n-2j)\sqrt y}\big)$ as $y\to\infty$. For integers $n/2\le j\le n$, $\int_1^\infty y^{s-1}\mathsf g^{n-j}\mathsf h^j\,dy=\infty$ for real $s\ge n$.
\item For these $j$, $\lim_{k\to\infty}\Big(\Gamma(k)^{-2}\int_0^\infty y^{k-1}\mathsf g^{n-j}\mathsf h^j\,dy\Big)^{1/k}=(n-2j)^{-2}$ along the integers.
\end{enumerate}
\end{slemma}

\begin{proof}
(i) $\mathsf g$ is the integral of a positive function and $\mathsf h\ge1$ term by term. The upper bound on $\mathsf g$ is Lemma~\hb{lem:onemark}(ii) at $\nu=0$. For the lower bound restrict the integral to $x=\sqrt y+\tau$, $0\le\tau\le1$, where $x+y/x=2\sqrt y+\tau^2/(\sqrt y+\tau)\le2\sqrt y+1$ and $x^{-1}\ge(2\sqrt y)^{-1}$. For $\mathsf h$, $\sum_my^m/m!^2\le\big(\sum_my^{m/2}/m!\big)^2=e^{2\sqrt y}$; and with $m=\lfloor\sqrt y\rfloor\ge1$ and $m!\le e\,m^{m+1}e^{-m}$ (from $\log m!\le\int_1^m\log x\,dx+\log m$),
\[
\mathsf h(y)\ge\frac{y^m}{m!^2}\ge\frac{e^{2m}}{e^2m^2}\Big(\frac{\sqrt y}m\Big)^{2m}\ge\frac{e^{2\sqrt y-2}}{e^2\,y}.
\]
(ii) By the ladder, $\vartheta\mathsf g=-y\,\sG_{-1;(1,-1)}(1,y)$, at most $\sqrt\pi e^{1/4}y^{1/4}e^{-2\sqrt y}$ for $y\ge1$ by Lemma~\hb{lem:onemark}(ii) at $\nu=-1$; and $0\le\vartheta\mathsf h=\sum_mmy^m/m!^2\le y\,\mathsf h$. Since $\vartheta^2=y$ on both, every $\vartheta^a\mathsf g$ and $\vartheta^a\mathsf h$ is a combination of the function and its $\vartheta$ with polynomial coefficients, of size $O(y^Ce^{\mp2\sqrt y})$ at infinity; near $0$, $\mathsf g=-\mathsf h\log y+\tilde{\mathsf g}$ with $\mathsf h$ and $\tilde{\mathsf g}$ entire, so every $\vartheta^a\mathsf g$ is $O(1+|\log y|)$ and every $\vartheta^a\mathsf h$ bounded. By the Leibniz rule $\vartheta^a(\mathsf g^{n-j}\mathsf h^j)$ is a sum of products of $n-j$ derivatives of $\mathsf g$ and $j$ of $\mathsf h$, which gives the bounds, $e^{-2(n-j)\sqrt y}e^{2j\sqrt y}=e^{-2(n-2j)\sqrt y}$. For $n/2\le j\le n$, (i) gives on $[1,\infty)$
\[
\mathsf g^{n-j}\mathsf h^j\ge(2e)^{-(n-j)}e^{-4j}\,y^{-(n+j)/2}\,e^{2(2j-n)\sqrt y}\ge(2e)^{-n}e^{-4n}\,y^{-n},
\]
and $\int_1^\infty y^{s-1-n}\,dy=\infty$ for $s\ge n$.

(iii) Put $b=n-2j\ge1$ and $I(k)=\int_0^\infty y^{k-1}\mathsf g^{n-j}\mathsf h^j\,dy$. The part over $(0,1]$ is at most $\int_0^1C(1+|\log y|)^n\,dy$, bounded in $k\ge1$. By (i), on $[1,\infty)$ the integrand lies between $c_1y^{k-1-n}e^{-2b\sqrt y}$ and $c_2y^{k-1}e^{-2b\sqrt y}$ with $c_1,c_2>0$, and for a fixed real $a$, with $y=\tau^2$,
\[
\int_1^\infty y^{k-1+a}e^{-2b\sqrt y}\,dy=2\int_1^\infty\tau^{2k+2a-1}e^{-2b\tau}\,d\tau=2\,\Gamma(2k+2a)\,(2b)^{-2k-2a}+O(1).
\]
$\Gamma(2k+2a)/\Gamma(2k)$ has root rate one and $\Gamma(2k)(2b)^{-2k}$ grows faster than every exponential, so $\big(I(k)/\Gamma(2k)\big)^{1/k}\to(2b)^{-2}$; and $\Gamma(2k)/\Gamma(k)^2=\frac k2\binom{2k}k$ has the root rate $4$, so $\big(I(k)/\Gamma(k)^2\big)^{1/k}\to4(2b)^{-2}=b^{-2}$.
\end{proof}

\begin{htheorem}[\hb{thm:ri-law} (The law of the values of $\binom{\cE}1$)]
As stated.
\end{htheorem}

\begin{proof}
(i) \emph{Annihilation.} Let $u\in\mathcal V$ and $v=\vartheta u$, so that $\vartheta v=yu$. For $0\le i\le n$, $\mathbf R_i(u^n)=\frac{n!}{(n-i)!}u^{n-i}v^i$: this holds at $i=0,1$, and $\vartheta(u^{n-i}v^i)=(n-i)u^{n-i-1}v^{i+1}+i\,y\,u^{n-i+1}v^{i-1}$ gives
\[
\vartheta\mathbf R_i(u^n)=\frac{n!}{(n-i-1)!}u^{n-i-1}v^{i+1}+i(n+1-i)\,y\,\frac{n!}{(n-i+1)!}u^{n-i+1}v^{i-1},
\]
whose last term is $i(n+1-i)\,y\,\mathbf R_{i-1}(u^n)$; so the recursion gives the claim at $i+1\le n$, and at $i=n$ the first term carries the factor $n-i=0$ and $\mathbf R_{n+1}(u^n)=0$. A product of $n$ solutions is a combination of $n$-th powers of solutions, $u_1\cdots u_n=\frac1{2^nn!}\sum_{\epsilon\in\{\pm1\}^n}\epsilon_1\cdots\epsilon_n\,(\epsilon_1u_1+\cdots+\epsilon_nu_n)^n$, so $\mathbf Q_n$ annihilates every such product.

\emph{Normal form and leading terms.} Since $\vartheta y=y(\vartheta+1)$, every operator built from $y$ and $\vartheta$ is $\sum_ky^kr_k(\vartheta)$ with polynomials $r_k$, uniquely: applied to $y^z$ it gives $\sum_ky^{k+z}r_k(z)$. Write $\mathbf R_i=\sum_ky^k\ell_{i,k}(\vartheta)$. By induction $\ell_{i,0}=\vartheta^i$ and $\deg\ell_{i,k}\le i-2k$: $\vartheta\,y^k\ell_{i,k}(\vartheta)=y^k(\vartheta+k)\ell_{i,k}(\vartheta)$ has degree at most $i+1-2k$, and so has $y\cdot y^{k-1}\ell_{i-1,k-1}(\vartheta)$. The coefficients $\alpha_{i,k}$ of $\vartheta^{i-2k}$ in $\ell_{i,k}$ therefore satisfy $\alpha_{i+1,k}=\alpha_{i,k}-i(n+1-i)\alpha_{i-1,k-1}$, the shift by $k$ not changing a leading coefficient, and $\sigma_i(\xi,\eta)=\sum_k\alpha_{i,k}\eta^k\xi^{i-2k}$ satisfies $\sigma_{i+1}=\xi\sigma_i-i(n+1-i)\,\eta\,\sigma_{i-1}$, $\sigma_0=1$, $\sigma_1=\xi$. Let $J$ be the matrix, indexed by $0,\dots,n$, of the derivation $\Delta$ of the forms of degree $n$ in $U,V$ with $\Delta U=V$ and $\Delta V=\eta U$, in the basis $w_i=U^{n-i}V^i$: $\Delta w_i=(n-i)w_{i+1}+i\eta\,w_{i-1}$. The leading principal minors of $\xi-J$, expanded along their last row, satisfy the same recursion, $\det_{i+1}=\xi\det_i-(n+1-i)\,i\eta\,\det_{i-1}$, so $\sigma_{n+1}=\det(\xi-J)$. Over $\mathbb C(\eta^{1/2})$ the forms $E_\pm=V\pm\eta^{1/2}U$ have $\Delta E_\pm=\pm\eta^{1/2}E_\pm$, so the basis $E_+^{n-j}E_-^j$ diagonalizes $\Delta$ with the eigenvalues $(n-2j)\eta^{1/2}$, and $\det(\xi-J)=\prod_{j=0}^n\big(\xi-(n-2j)\eta^{1/2}\big)$; the factors $j$ and $n-j$ pair into $\xi^2-(n-2j)^2\eta$, and at even $n$ the factor $j=n/2$ is $\xi$. With $q_{n,k}=\ell_{n+1,k}$ and $a_{n,k}=\alpha_{n+1,k}$ this is the display; $a_{n,h}$ is the coefficient of the top power $\eta^h$, $\prod_{0\le j<n/2}\big(-(n-2j)^2\big)\neq0$.

\emph{The adjoint.} Expanding the recursion, $\mathbf R_m$ is the sum over the tilings of $\{1,\dots,m\}$ by monomers $\{a\}$ and dimers $\{a,a+1\}$ of the products of $\vartheta$ for each monomer and $-a(n+1-a)\,y$ for each dimer, taken from left to right in decreasing order of position: the tilings of $\{1,\dots,m+1\}$ end with the monomer $\{m+1\}$ or the dimer $\{m,m+1\}$, which are the two terms of the recursion. The formal adjoint for $dy/y$ reverses products, takes $\vartheta$ to $-\vartheta$ and keeps $y$; a tiling of $\{1,\dots,m\}$ with $d$ dimers has $m-2d$ monomers, so $\mathbf R_m^\ast$ is $(-1)^m$ times the same sum with the products taken in increasing order of position. The reflection $a\mapsto m+1-a$ of the positions maps tilings to tilings, reverses the order, and takes the dimer $\{a,a+1\}$ to $\{a',a'+1\}$ with $a'=m-a$, of weight $-(m-a')(n+1-m+a')\,y$. At $m=n+1$ this weight is $-a'(n+1-a')\,y$, so $\mathbf Q_n^\ast=\mathbf R_{n+1}^\ast=(-1)^{n+1}\mathbf R_{n+1}$. Since $q(-\vartheta)\,y^k=y^kq(-\vartheta-k)$, $\mathbf Q_n^\ast=\sum_kq_{n,k}(-\vartheta)y^k=\sum_ky^kq_{n,k}(-\vartheta-k)$, and the uniqueness of the normal form gives $q_{n,k}(-\vartheta-k)=(-1)^{n+1}q_{n,k}(\vartheta)$.

(ii) For $k\ge1$, $c_{n,k}$ has degree at most $(n+1-2k)+2(k-1)=n-1$ with the coefficient $(-1)^{n+1}a_{n,k}$ of $s^{n-1}$, as $c_{n,0}=(-1)^{n+1}s^{n-1}$ has with $a_{n,0}=1$. Setting $\xi=1$ in (i), $\sum_ka_{n,k}\lambda^k=\prod_{0\le j<n/2}\big(1-(n-2j)^2\lambda\big)$, whose coefficients $(-1)^ke_k\big((n-2j)^2\big)$ are nonzero for $0\le k\le h$; so every $c_{n,k}$ has degree $n-1$ and the symbol is $\chi_n$. By (i), $c_{n,k}(-s-k)=q_{n,k}(s)\big((-s-k+1)\cdots(-s-1)\big)^2=(-1)^{n+1}q_{n,k}(-s-k)\big((s+1)\cdots(s+k-1)\big)^2=(-1)^{n+1}c_{n,k}(s)$ for $k\ge1$, and $c_{n,0}(-s)=s^{n-1}=(-1)^{n+1}c_{n,0}(s)$. The double zeros are those of the squared factor. For odd $n$, $q_{n,h}$ is the nonzero constant $a_{n,h}$; for even $n$ it has degree one, $q_{n,h}(\vartheta)=a_{n,h}(\vartheta-\vartheta_0)$, and $q_{n,h}(-\vartheta-h)=-q_{n,h}(\vartheta)$ forces $\vartheta_0=-h/2$, so $q_{n,h}(-s-h)=-a_{n,h}(s+h/2)$; in both cases $c_{n,h}(s)\neq0$ at the integers $s\ge0$. For $n\ge2$ let $M(s)=\Gamma(s)^2P_{\binom{\cE}1}(s\one;s,s)=\int_0^\infty y^{s-1}\mathsf g^n\,dy$ (Section~\hb{sec:ri-centre}). By Lemma~\hb{lem:ri-growth}(ii), Lemma~\hb{lem:ri-mellin}(i) applies to $\mathsf g^n$, and $\mathbf Q_n(\mathsf g^n)=0$ gives $\sum_kq_{n,k}(-s-k)\,M(s+k)=0$ on $\re s>0$. Dividing by $\Gamma(s+1)^2$, with $\Gamma(s)^2/\Gamma(s+1)^2=s^{-2}$ and $\Gamma(s+k)^2/\Gamma(s+1)^2=\big((s+1)\cdots(s+k-1)\big)^2$, and $q_{n,0}(-s)=(-s)^{n+1}$, gives the identity on $\re s>0$, and on $\mathbb C$ by the meromorphic continuation (Proposition~\hb{prop:principal}).

(iii) Write $P(k)=P_{\binom{\cE}1}(k\one;k,k)$ and suppose $\sum_{i=0}^mr_i(k)P(k+i)=0$ for all large integers $k$, with $r_i\in\mathbb C[k]$, $r_0r_m\neq0$ and $m<h$. The centre ray at the strength one has origin $0$, letter exponents one and walls $\kappa_{1,S}=1$, so by Proposition~\hb{prop:ri-rayinput}(i) the relation holds identically. On $\re s>0$ multiply it by $\Gamma(s+m)^2$: $\sum_i\tilde r_i(s)M(s+i)=0$ with the polynomials $\tilde r_i(s)=r_i(s)\big((s+i)\cdots(s+m-1)\big)^2$, $\tilde r_m=r_m$. Put $T=\sum_iy^i\,\tilde r_i(-\vartheta-i)$. By Lemma~\hb{lem:ri-mellin}(i), with Lemma~\hb{lem:ri-growth}(ii), the Mellin read of $T(\mathsf g^n)$ is $\sum_i\tilde r_i(s)M(s+i)=0$, and by Lemma~\hb{lem:ri-mellin}(ii) $T(\mathsf g^n)=0$.

On the ring of polynomials in an indeterminate $L$ with entire coefficients let $\vartheta$ act by $\vartheta\big(\sum_lE_lL^l\big)=\sum_l(\vartheta E_l)L^l+\sum_llE_lL^{l-1}$, and put $\Phi(y,L)=T\big[(\tilde{\mathsf g}-\mathsf h\,L)^n\big]=\sum_{l\le n}E_l(y)L^l$, $E_l$ entire. Substituting $L=\log y-c$, for a constant $c$, commutes with this $\vartheta$, so $T\big[(\mathsf g+c\,\mathsf h)^n\big]=\Phi(y,\log y-c)$ on $(0,\infty)$, and $\Phi(y,\log y)=T(\mathsf g^n)=0$. The function $y\mapsto\Phi(y,\log y)$ continues along every path in $\mathbb C^\times$, and once around $0$ it becomes $\Phi(y,\log y+2\pi i)$; the continuation of $0$ is $0$, so $\Phi(y,\log y+2\pi i\ell)=0$ for every $\ell\in\mathbb Z$ and $y>0$. For fixed $y$ the polynomial $L\mapsto\Phi(y,L)$ then has infinitely many zeros, so $\Phi=0$, and $T\big[(\mathsf g+c\mathsf h)^n\big]=0$ for every $c\in\mathbb C$; the coefficients of the powers of $c$ give $T\big(\mathsf g^{n-j}\mathsf h^j\big)=0$ for $0\le j\le n$.

For $0\le j<n/2$ let $I_j(k)=\int_0^\infty y^{k-1}\mathsf g^{n-j}\mathsf h^j\,dy$. Lemma~\hb{lem:ri-mellin}(i) for $\mathsf g^{n-j}\mathsf h^j$ (Lemma~\hb{lem:ri-growth}(ii)) gives $\sum_i\tilde r_i(k)I_j(k+i)=0$ for the integers $k\ge1$. Take $k_1$ beyond the zeros of $r_m$; on the integers $k\ge k_1$ a solution of this recurrence of order $m$ is fixed by its first $m$ values, so the $h>m$ sequences $(I_j(k))_{k\ge k_1}$ are linearly dependent. They are not: if $\sum_jb_jI_j(k)=0$ for $k\ge k_1$ with some $b_j\neq0$, let $j_0$ be the largest such $j$; by Lemma~\hb{lem:ri-growth}(iii), $I_j(k)/\Gamma(k)^2$ has the exact root rate $(n-2j)^{-2}$, largest at $j_0$, so $|b_{j_0}|I_{j_0}(k)>\sum_{j<j_0}|b_j|I_j(k)$ for large $k$, a contradiction. So the tier is at least $h$, and the recurrence of (ii) has the order $h$ and $c_{n,0}c_{n,h}\neq0$.

By Proposition~\hb{prop:ri-rayinput}(i) the closings, here the constant terms of $(\sum_e\alpha_e)^j(\sum_e\alpha_e^{-1})^j$ (Corollary~\hb{cor:reflected}), that is $\sum_{|v|=j}\binom jv^2$, satisfy the reflected recurrence for every $j\in\mathbb N$, and by Corollary~\hb{cor:valuemodule} they satisfy no recurrence that the values do not satisfy, read at $k$ and at $-j$; so none of order less than $h$. In the value module, a relation $\sum_{i<h}r_i(k)\operatorname{Sh}^i[1]\in\operatorname{Ann}(\langle\cdot\rangle)$ with rational $r_i$ not all zero would be a recurrence of order less than $h$ for the values, so the classes of $\operatorname{Sh}^i[1]$, $0\le i<h$, are independent there; the element $\sum_{i=0}^hc_{n,i}(k)\operatorname{Sh}^i[1]$, $k$ the strength, has vanishing values and lies in $\operatorname{Ann}(\langle\cdot\rangle)$, a submodule (Theorem~\hb{thm:determine}(ii)), and with $c_{n,0}c_{n,h}\neq0$ it expresses $\operatorname{Sh}^h[1]$ and $\operatorname{Sh}^{-1}[1]$ through these classes; so their span is the cyclic submodule of $[1]$ in the value module, of dimension $h$.
\end{proof}

\begin{hproposition}[\hb{prop:ri-markadded} (A mark added carries the law of the values)]
As stated.
\end{hproposition}

\begin{proof}
(i) For a fixed solution $u$ both sides are linear differential operators in $F$ of order $n+2$, with coefficients analytic on $(0,\infty)$. By Theorem~\hb{thm:ri-law}(i), $\mathbf Q_{n+1}=\vartheta^{n+2}+\sum_{k\ge1}y^kq_{n+1,k}(\vartheta)$ with $\deg q_{n+1,k}\le n$, so $\mathbf Q_{n+1}(uF)=u\,\vartheta^{n+2}F+(n+2)(\vartheta u)\,\vartheta^{n+1}F+(\text{order}\le n)$; and $\mathbf Q_n=\vartheta^{n+1}+\sum_{k\ge1}y^kq_{n,k}(\vartheta)$ with $\deg q_{n,k}\le n-1$, so the right side is $u\,\vartheta^{n+2}F+(n+2)(\vartheta u)\,\vartheta^{n+1}F+(\text{order}\le n)$. Their difference is an operator of order at most $n$ that annihilates every $F$ in the space spanned by the products of $n$ solutions, of dimension $n+1$: the right side vanishes there by Theorem~\hb{thm:ri-law}(i), and the left side because $uF$ is a product of $n+1$ solutions. An operator of order at most $n$ whose highest nonzero coefficient does not vanish on an interval has there a space of solutions of dimension at most $n$, and the $n+1$ analytic functions stay independent on every interval; so the difference is $0$.

(ii) Apply (i) with $u=\mathsf h$ to $F=y^z$, $z\in\mathbb C$, using $\vartheta^ay^w=w^ay^w$ and $\vartheta\mathsf h=\sum_mm\,y^m/m!^2$. Both sides are $y^z$ times entire functions of $y$, and comparing the coefficients of $y^{z+N}$,
\[
\sum_{k+m=N}\frac{q_{n+1,k}(z+m)}{m!^2}=\sum_{j+m=N}\frac{\big(z+j+(n+2)m\big)\,q_{n,j}(z)}{m!^2}\qquad(N\in\mathbb N,\ z\in\mathbb C).
\]
By Theorem~\hb{thm:ri-law}(i), $q_{n+1,k}(z+m)=(-1)^nq_{n+1,k}(-z-N)$ when $k+m=N$ and $q_{n,j}(z)=(-1)^{n+1}q_{n,j}(w+m)$ with $w=-z-N$ when $j+m=N$; and $z+j+(n+2)m=(n+1)m-w$. So
\[
\sum_{k+m=N}\frac{q_{n+1,k}(w)}{m!^2}=\sum_{j+m=N}\frac{\big(w-(n+1)m\big)\,q_{n,j}(w+m)}{m!^2},
\]
and applying the same symmetry once more, $q_{n+1,k}(w)=(-1)^nq_{n+1,k}(-w-k)$ and $q_{n,j}(w+m)=(-1)^{n+1}q_{n,j}(-w-N)$,
\[
\sum_{k+m=N}\frac{q_{n+1,k}(-s-k)}{m!^2}=\sum_{j+m=N}\frac{\big((n+1)m-s\big)\,q_{n,j}(-s-N)}{m!^2}\qquad(s=w).
\]
Multiply by $\Gamma(s+N)^2/\Gamma(s+1)^2$ and put $x=s+N$. By the definitions, $c_{n+1,k}(s)=q_{n+1,k}(-s-k)\,\Gamma(s+k)^2/\Gamma(s+1)^2$ for every $k\ge0$, so the left side becomes $\sum_kc_{n+1,k}(s)\,\Gamma(x)^2/\big(\Gamma(s+k)^2\Gamma(x-s-k+1)^2\big)=\sum_kc_{n+1,k}(s)\binom{x-1}{s+k-1}^2$. With $c_{n,j}(x-j)=q_{n,j}(-x)\,\Gamma(x)^2/\Gamma(x-j+1)^2$, $(n+1)m-s=(n+1)(x-j)-(n+2)s$ and $A_s(x-j)=\big((n+1)(x-j)-(n+2)s\big)\Gamma(x-j+1)^2/\big(\Gamma(s+1)^2\Gamma(x-j-s+1)^2\big)$, the right side becomes $\sum_jc_{n,j}(x-j)A_s(x-j)$. After division by $\binom xs^2$ both sides are rational functions of $(x,s)$, each binomial quotient being a quotient of rising factorials, and they agree on the lines $x-s=N$, $N\in\mathbb N$; a polynomial in $(x,s)$ vanishing on infinitely many of these lines is divisible by infinitely many distinct factors $x-s-N$, so it is zero, and the identity holds. At integers $s,x\ge1$, $W_{s+k}(x)/\Gamma(s+k)^2=\prod_{l=1}^{s+k-1}(x-l)^2/(s+k-1)!^2=\binom{x-1}{s+k-1}^2$, both sides vanishing for $x<s+k$; the right side is $(L_n^\ast A_s)(x)$ for the recurrence $L_n$ with the coefficients $c_{n,j}$ (Definition~\hb{def:ri-inputs}), and $\binom xs=x(x-1)\cdots(x-s+1)/s!$ vanishes at $x=0,1,\dots,s-1$.
\end{proof}

\begin{hproposition}[\hb{prop:ri-lowered} (The lowered settings)]
As stated.
\end{hproposition}

\begin{proof}
(i) By Theorem~\hb{thm:ri-law}(i), $\mathbf Q_n(\mathsf g^{n-j}\mathsf h^j)=0$; Lemma~\hb{lem:ri-mellin}(i), which applies by Lemma~\hb{lem:ri-growth}(ii), gives $\sum_kq_{n,k}(-s-k)\,\Gamma(s+k)^2\xi_{n,j}(s+k)=0$ on $\re s>0$, and the division of the proof of Theorem~\hb{thm:ri-law}(ii) gives the recurrence. For $j=0$ the display of Section~\hb{sec:ri-centre} gives $\xi_{n,0}(s)=P_{\binom{\cE}1}(s\one;s,s)$.

(ii) On $(0,1]$, $0<\mathsf g^{n-j}\mathsf h^j\le C(1+|\log y|)^{n-j}$ by Lemma~\hb{lem:ri-growth}(ii) and $\mathsf g=-\mathsf h\log y+\tilde{\mathsf g}$, so for real $s\in(0,1]$, $\int_0^\infty y^{s-1}\mathsf g^{n-j}\mathsf h^j\,dy\le C\int_0^1y^{s-1}(1+|\log y|)^{n-j}dy+C'=O(s^{-(n-j+1)})$, and with $\Gamma(s)^{-2}=O(s^2)$, $s^{n-1}\xi_{n,j}(s)=O(s^j)\to0$ for $j\ge1$. In the recurrence on $(0,\infty)$ the term $c_{n,0}(s)\xi_{n,j}(s)=(-1)^{n+1}s^{n-1}\xi_{n,j}(s)$ then tends to $0$ as $s\downarrow0$, while the others tend to $c_{n,k}(0)\xi_{n,j}(k)$, $\xi_{n,j}$ being continuous on $\re s>0$; so $\sum_{k\ge1}c_{n,k}(0)\xi_{n,j}(k)=0$. For $j=0$, $s^{n-1}\xi_{n,0}(s)\to V_1=n!$ by Proposition~\hb{prop:principal}, the walls of the direction being $\kappa_{1,S}=c(S)=1$ for every nonempty proper $S$, so that $V_1=\sum_\pi1=n!$; and the same limit gives $\sum_{k\ge1}c_{n,k}(0)\xi_{n,0}(k)=-(-1)^{n+1}n!$.

(iii) By Lemma~\hb{lem:ri-growth}(iii), $\xi_{n,j}$ has the root rate $\lambda=(n-2j)^{-2}$ along the integers. The recurrence has $c_{n,0}(k)c_{n,h}(k)\neq0$ at the integers $k\ge1$, every $c_{n,k}$ of degree $n-1$, and the roots $(n-2j')^{-2}$, $0\le j'<n/2$, of its symbol are positive and distinct (Theorem~\hb{thm:ri-law}(ii)); so the proof of Proposition~\hb{prop:ri-minimal}(i), which uses only these properties, gives $\xi_{n,j}(k+1)/\xi_{n,j}(k)\to\lambda_c$ for a root $\lambda_c$ of modulus $\lambda$, that is $\lambda_c=(n-2j)^{-2}$.

(iv) $c_{n,h}(k)\neq0$ at the integers $k\ge1$ (Theorem~\hb{thm:ri-law}(ii)), so a solution on these integers is fixed by its first $h$ values, and the solutions form a space of dimension $h$. The $h$ sequences are solutions by (i) and are independent by (iii), as in the proof of Theorem~\hb{thm:ri-law}(iii): in a vanishing combination the term with the largest $j$ would dominate. So they are a basis. The linear form $x\mapsto\sum_{k=1}^hc_{n,k}(0)x(k)$ is $(-1)^nn!$ on $\xi_{n,0}$ and $0$ on the others by (ii), which gives the last assertion.
\end{proof}

\begin{hcorollary}[\hb{cor:ri-span} (The value span of $\binom{\cE}1$)]
As stated.
\end{hcorollary}

\begin{proof}
(i) The first relation is Proposition~\hb{prop:ri-lowered}(ii) at $j=0$. The recursion of Theorem~\hb{thm:ri-law} has integer coefficients, so every $q_{n,k}$ and every $c_{n,k}$ has integer coefficients. $c_{n,h}(0)\neq0$, so the first relation gives $p_n(h)\in\mathbb Q+\sum_{i<h}\mathbb Q\,p_n(i)$, and for $k>h$ the recurrence of Theorem~\hb{thm:ri-law}(ii) at $s=k-h\ge1$, with $c_{n,h}(k-h)\neq0$, writes $p_n(k)$ through $p_n(k-h),\dots,p_n(k-1)$; by induction every $p_n(k)$ lies in $\mathbb Q+\sum_{i=1}^{h-1}\mathbb Q\,p_n(i)$. The values of the centre ray at the strength $\sigma$ are $p_n(k\sigma)$ (Definition~\hb{def:centreray}).

(ii) For even $n$ every $n-2j$, $0\le j<n/2$, is at least $2$, so $\chi_n(1)=(-1)^{n+1}\prod_j\big(1-(n-2j)^2\big)\neq0$. The values of the centre ray have the root rate $n^{-2}<1$ and the slopes at most $n^{-2}$ (Proposition~\hb{prop:ri-rayinput}(iii), $f(\alpha^\star)=f(\one)=n^{-2}$). Every $c_{n,k}$ with $k>j\ge1$ has a double zero at $-j$ (Theorem~\hb{thm:ri-law}(ii)), so by Proposition~\hb{prop:ri-sbp} the only boundary numbers are $\beta_0(0)$ and $\beta_0'(0)$ of $\beta_0=c_{n,0}\,p_n=(-1)^{n+1}s^{n-1}p_n$: by Proposition~\hb{prop:ri-boundary}, $\beta_0(0)=(-1)^{n+1}V_1=(-1)^{n+1}n!$, and $\beta_0'(0)=(-1)^{n+1}\mathsf L_1=0$, the list having unit weights. Proposition~\hb{prop:ri-block}(ii) with $\Bbbk=\mathbb Q$ and $D=n-1$ gives the claim. The two last numbers are such Trailmoments by Proposition~\hb{prop:ri-addmark}(ii), with the rational polynomials $W_s$ and $W_s'$.
\end{proof}

\begin{hcorollary}[\hb{cor:ri-pairing} (The exact pairing of $\binom{\cE}1$)]
As stated.
\end{hcorollary}

\begin{proof}
The centre ray of $\binom{\cE}1$ at the strength one satisfies the hypotheses of Proposition~\hb{prop:ri-rayinput}, and the recurrence of Theorem~\hb{thm:ri-law}(ii) holds for its values with $c_{n,k}(-s-k)=(-1)^{n+1}c_{n,k}(s)$; its closings are the $\mathsf C_l$ of the statement (proof of Theorem~\hb{thm:ri-law}(iii)). Proposition~\hb{prop:ri-pairing}(ii), with $\gamma_{n-1}=(-1)^{n+1}$ and $V_1=n!$, gives the first display. Every $c_{n,k}$ has degree $n-1$, so Proposition~\hb{prop:ri-pairing}(iii) applies: the roots of $\chi_n$ are the simple roots $(n-2j)^{-2}$, $0\le j<n/2$, with $\lambda_\star=f(\one)=n^{-2}$ at $j=0$, and $1-(n-2j)^2/n^2=4j(n-j)/n^2$. The Hessian $H$ of Theorem~\hb{thm:branch} at the centre is $2\big(\kappa_{\one}(X_e,X_{e'})\big)_{e,e'\neq e_0}$, which is the Hessian of Proposition~\hb{prop:ri-signs}(ii) at the sign $\one$ and $\sigma=1$, of determinant $2^{n-1}/n^n$. As a check, the two sides agree directly: for $n=2l+1$, $h-1=l$ and $\prod_{j=1}^lj(n-j)=l!\,(2l)!/l!=(n-1)!$, and for $n=2l+2$, $h-1=l$ and $\prod_{j=1}^lj(n-j)=l!\,(2l+1)!/(l+1)!=2(n-1)!/n$; in both cases $\prod_{j=1}^{h-1}4j(n-j)/n^2=2^{n-1}(n-1)!/n^{n-1}=n!\,2^{n-1}/n^n$.
\end{proof}

\begin{hproposition}[\hb{prop:ri-mixture} (The power of one mark as a mixture)]
As stated.
\end{hproposition}

\begin{proof}
Write $p_n(s)=P_{\binom{\cE}1}(s\one;s,s)$ for $|\cE|=n$. By Theorem~\hb{thm:pencil}(i), $0<f_n\le f_n(\one)=n^{-2}$ on the chart, so $\mu_{f_n}$ is carried by $(0,n^{-2}]$, and $\int t^s\mu_{f_n}(dt)=\int_{y_{e_0}=0}f_n^s\,dy=p_n(s)$ for $\re s>0$, the image of a measure under a map (Definition~\hb{def:conv}(vi)); the same holds at $n+1$ marks. The measure $\mu_{f_n}$ has infinite mass near $0$, and every moment of positive order is finite. By Fubini's theorem and $y=ty'$, $\int_0^\infty y^{s-1}\mathsf g(y/t)\,dy=t^s\int_0^\infty y'^{s-1}\mathsf g(y')\,dy'=t^s\Gamma(s)^2$, the last because $\int_0^\infty y^{s-1}e^{-y/x}dy=\Gamma(s)x^s$ and $\int_0^\infty x^{s-1}e^{-x}dx=\Gamma(s)$.

(i) Put $G(y)=\int\mathsf g(y/t)\,\mu_{f_n}(dt)$. For real $s>0$, by Tonelli's theorem $\int_0^\infty y^{s-1}G(y)\,dy=\Gamma(s)^2\int t^s\mu_{f_n}(dt)=\Gamma(s)^2p_n(s)$, finite, which is $\int_0^\infty y^{s-1}\mathsf g^n\,dy$ by the display of Section~\hb{sec:ri-centre}. So $G$ is finite almost everywhere; since $\mathsf g$ is decreasing ($\partial_y\mathsf g=-\sG_{-1;(1,-1)}(1,y)<0$), $\mathsf g(y/t)\le\mathsf g(y_0/t)$ for $y\ge y_0$, and dominated convergence makes $G$ finite and continuous on $(0,\infty)$. For $s=c+i\tau$ the same Fubini argument applies, so the integrable functions $e^{cv}G(e^v)$ and $e^{cv}\mathsf g(e^v)^n$ have the same Fourier transform; they are equal, by the uniqueness of the Fourier transform on $L^1(\mathbb R)$ as in Lemma~\hb{lem:ri-mellin}(ii), and both being continuous, $G=\mathsf g^n$.

(ii) The table on $\{1,2\}$ with the bases $\{1\}$ and $\{2\}$ is the table of Corollary~\hb{cor:cycle} on two marks, the basis $\{2\}=\{1,2\}-1$ carrying $\theta_1=1$ and $\{1\}$ carrying $\theta_2=t^{-1/2}$; its letters are $K=\alpha_2+t^{-1/2}\alpha_1$ and $\Psi=\alpha_1+t^{-1/2}\alpha_2$, and its centre stride on the chart $\alpha_1=1$, $\alpha_2=x$, is $\alpha_1\alpha_2/(K\Psi)=x/\big((x+t^{-1/2})(1+t^{-1/2}x)\big)=\phi_t(x)$, and $dy=dx/x$ on that chart. So $\int u^s\mu_{\phi_t}(du)=P(s\one;s,s)$ for that table, and Corollary~\hb{cor:cycle} at $\mvec=s\one$, $\nu_e=0$, gives $\Gamma(s)^2\int u^s\mu_{\phi_t}(du)=\int_0^\infty y^{s-1}\mathsf g(y)\,\mathsf g(\theta_2^2y)\,dy=\int_0^\infty y^{s-1}\mathsf g(y)\,\mathsf g(y/t)\,dy$. By (i) and Tonelli's theorem, for real $s>0$,
\begin{align*}
\Gamma(s)^2\int u^s\mu_{f_{n+1}}(du)&=\int_0^\infty y^{s-1}\mathsf g(y)\,\mathsf g(y)^n\,dy=\int\mu_{f_n}(dt)\int_0^\infty y^{s-1}\mathsf g(y)\,\mathsf g(y/t)\,dy\\
&=\Gamma(s)^2\int\mu_{f_n}(dt)\int u^s\mu_{\phi_t}(du),
\end{align*}
and by Fubini's theorem the same holds on $\re s>0$. The measures $\mu_{f_{n+1}}$ and $\int\mu_{\phi_t}\,\mu_{f_n}(dt)$ therefore have the same moments $\int u^s(\cdot)$ on the line $\re s=c>0$; their images under $u\mapsto v=\log u$, weighted by $e^{cv}$, are finite measures on $\mathbb R$ with the same Fourier transform, hence equal, and so are the two measures.
\end{proof}

\begin{hproposition}[\hb{prop:ri-directions} (Balanced directions of $\binom{\cE}1$)]
As stated.
\end{hproposition}

\begin{proof}
The walls of $\binom{\cE}1$ at $(\mvec;s,t)$ are $\kappa_S=m_S-t(|S|-1)$, every proper $S$ meeting a basis $\{e\}$ in at most one mark and the cobases $\cE-e$ in at least $|S|-1$, so the hypothesis says that every proper wall is positive, at $(k\mvec;ks,kt)$ as well. The display of Section~\hb{sec:ri-centre} at these data, from Corollary~\hb{cor:cycle} for the dual table at the exponents $(kt,ks)$ and Proposition~\hb{prop:dualfamily}, gives $\Gamma(sk)\Gamma(tk)P(k)=\int_0^\infty y^{tk-1}\prod_e\sG_{c_ek;(1,-1)}(1,y)\,dy$. By the torus homogeneity (Theorem~\hb{thm:grammalaws}(iii)) at $\lambda=y^{1/2}$, $\sG_{\nu;(1,-1)}(\lambda,\lambda^{-1}y)=\lambda^{-\nu}\sG_{\nu;(1,-1)}(1,y)$, that is $\sG_{\nu;(1,-1)}(1,y)=y^{\nu/2}\,\sG_{\nu;(1,-1)}(\sqrt y,\sqrt y)$; the product of the powers is $y^{k\sum_ec_e/2}=y^{(s-t)k/2}$, and $tk-1+\frac{(s-t)k}2=\frac{(s+t)k}2-1$. The reflection of Proposition~\hb{prop:grammareflect} gives $\sG_{\nu;(1,-1)}(a,b)=\sG_{-\nu;(-1,1)}(a,b)=\sG_{-\nu;(1,-1)}(b,a)$, so at equal coefficients the Gramma of one mark is even in its index, and the integral depends only on $s+t$ and the multiset $\{|c_e|\}$. Two directions with the same $s+t$ and multiset therefore have $\Gamma(sk)\Gamma(tk)P(k)=\Gamma(s'k)\Gamma(t'k)P'(k)$, which is the stated quotient of Beta corners, a quotient of factorials at integer $k$. The image of the integrand $\alpha^{k\mvec}K^{-ks}\Psi^{-kt}$ under $\alpha\mapsto\alpha^{-1}$, with $K(\alpha^{-1})=\alpha^{-\cE}\Psi(\alpha)$ and $\Psi(\alpha^{-1})=\alpha^{-\cE}K(\alpha)$, is $\alpha^{k((s+t)\one-\mvec)}K^{-kt}\Psi^{-ks}$, the integrand at $((s+t)\one-\mvec;t,s)$, whose offsets are $(s+t)-m_e-s=-c_e$.
\end{proof}

\section{Chapter~\hb{ch:rayarith}: arithmetic along a ray}\label{sec:chrayarith}

Throughout, Stirling's formula is $\log\Gamma(z)=(z-\frac12)\log z-z+\frac12\log2\pi+O(|z|^{-1})$ as $|z|\to\infty$ with $|\arg z|\le\pi-\delta$; it applies below to $\Gamma(k-a)$ and to $N!=\Gamma(N+1)$ with $a$ fixed and $k,N\to+\infty$ along the reals, so that the arguments stay in a sector around the positive axis.

\begin{slemma}\label{lem:rayarith-cyclic}
Let $\Bbbk$ be a subfield of $\mathbb C$. Let $X(n+1)=M(n)\,X(n)$ for the integers $n\ge n_2$, with $X(n)\in\mathbb C^m$ and $M(n)$ an $m\times m$ matrix whose entries are rational functions of $n$ with coefficients in $\Bbbk$ and whose determinant is not the zero function, and let $\psi(n)=\ell\,X(n)$ for a fixed row vector $\ell\in\Bbbk^{1\times m}$. Then $\psi$ satisfies, for all large $n$, a recurrence $\sum_{i=0}^{h}a_i(n)\,\psi(n+i)=0$ with $a_i\in\Bbbk[n]$, $h\le m$ and $a_0a_h\neq0$.
\end{slemma}

\begin{proof}
If $\ell=0$, $\psi=0$ satisfies $\psi(n)=0$. Otherwise put $\mathbf r_0=\ell$ and $\mathbf r_i(n)=\ell\,M(n+i-1)\cdots M(n)$, row vectors over the field $\Bbbk(n)$; then $\psi(n+i)=\mathbf r_i(n)X(n)$ wherever the entries are defined, and $\mathbf r_i(n)=\mathbf r_{i-1}(n+1)M(n)$. The $m+1$ vectors $\mathbf r_0,\dots,\mathbf r_m$ are dependent over $\Bbbk(n)$; let $h\ge1$ be the least index with $\mathbf r_0,\dots,\mathbf r_h$ dependent, and $\sum_{i=0}^ha_i\mathbf r_i=0$ with $a_i\in\Bbbk[n]$ after clearing denominators. By minimality $a_h\neq0$. If $a_0=0$, then $\sum_{i=1}^ha_i(n)\mathbf r_{i-1}(n+1)M(n)=0$, and $M(n)$ is invertible over $\Bbbk(n)$, so $\sum_{i=1}^ha_i(n)\mathbf r_{i-1}(n+1)=0$; replacing $n$ by $n-1$, an automorphism of $\Bbbk(n)$, gives a dependence among $\mathbf r_0,\dots,\mathbf r_{h-1}$ with the coefficient $a_h(n-1)\neq0$, against minimality. So $a_0\neq0$, and $\sum_ia_i(n)\psi(n+i)=\big(\sum_ia_i(n)\mathbf r_i(n)\big)X(n)=0$ for all $n$ beyond the poles of the entries.
\end{proof}

\begin{htheorem}[\hb{thm:ra-oneshare} (Closing one free share along a ray)]
As stated.
\end{htheorem}

\begin{proof}
(i) For an integer $r\ge0$ and $x$ off the poles, $\Gamma(x+r)=(x)_r\Gamma(x)$ (\hb{def:conv}(v)) and $\Gamma(x-r)=\Gamma(x)/(x-r)_r=(x)_{-r}\Gamma(x)$ (\hb{prop:stacked}(iii)). Applied to every Gamma of $W_k$ with $x=\epsilon_\gamma a+b_\gamma$ and $r=ku_\gamma$ of either sign this is $W_k=W_0\,g_k$; the factors of $g_k$ have integer coefficients.

(ii) The upper Gammas have poles where $\epsilon_\alpha a+b_\alpha+ku_\alpha\in-\mathbb N$, at integers since $\epsilon_\alpha=\pm1$; the lower ones contribute the entire factors $1/\Gamma$. For $n>n_1$ the arguments at $a=-n+\varepsilon$ are $-N_\gamma(n)+\varepsilon$ on $A_+$, $N_\gamma(n)+1-\varepsilon$ on $A_-$, $-N_\gamma(n)+\epsilon_\gamma\varepsilon$ on $B_+$ and $N_\gamma(n)+1-|\epsilon_\gamma|\varepsilon$ on $B_-$, with $N_\gamma(n)\in\mathbb N$, and $N_\gamma(n)$ is $|\epsilon_\gamma|n$ plus an integer. By \hb{prop:resonance}, $\Gamma(-N+\delta)=(-1)^N\delta^{-1}\Gamma(1+\delta)\Gamma(1-\delta)/\Gamma(1+N-\delta)$ for $N\in\mathbb N$ and $|\delta|<1$. Put $s_\gamma=1$ on $A_-\cup B_+$ and $s_\gamma=-1$ on $A_+\cup B_-$. Then
\[
W_k(-n+\varepsilon)=(-1)^{\sum_{A_+\cup B_+}N_\gamma}\prod_{B_+}\epsilon_\gamma\cdot\varepsilon^{-w}\prod_{\gamma}\Gamma\big(1+N_\gamma-|\epsilon_\gamma|\varepsilon\big)^{s_\gamma}\,\frac{\prod_{A_+}\Gamma(1+\varepsilon)\Gamma(1-\varepsilon)}{\prod_{B_+}\Gamma(1+\epsilon_\gamma\varepsilon)\Gamma(1-\epsilon_\gamma\varepsilon)}.
\]
At $\varepsilon=0$ the factor after $\varepsilon^{-w}$ is $\prod_\gamma(N_\gamma!)^{s_\gamma}$, which gives $C_n$ as stated, nonzero and rational. By \hb{prop:headstails} at the integer $b=N\ge0$, where the lattice input $\Ilat{N+\mathbb Z}$ is $\IZ$, $\log\Gamma(1+N-\delta)=\log N!+\sum_{i\ge1}\frac{\delta^i}i\Tx{i}{\IZ}{N}$, the index-one value regularized, and at $N=0$ with $\pm\delta$, $\log\Gamma(1+\delta)+\log\Gamma(1-\delta)=\sum_{j\ge1}\frac{\delta^{2j}}j\Tx{2j}{\IZ}{0}$. So $W_k(-n+\varepsilon)=C_n\varepsilon^{-w}\exp\big(\sum_i\Lambda_{n,i}\varepsilon^i\big)$ with
\[
\Lambda_{n,i}=\frac1i\sum_\gamma s_\gamma|\epsilon_\gamma|^i\,\Tx{i}{\IZ}{N_\gamma(n)}+[i\text{ even}]\,\frac2i\Big(|A_+|-\sum_{B_+}\epsilon_\gamma^{\,i}\Big)\Tx{i}{\IZ}{0}.
\]
The cut relation (\hb{prop:cutlaw}, with the regularized values at index one) gives $\Tx{i}{\IZ}{N}=\Tx{i}{\IZ}{0}-\Th{i}{\IZ}{0}{N}$. The balance of the slopes gives
\[
\sum_\gamma s_\gamma|\epsilon_\gamma|=\#A_--\#A_++\sum_{B_+}\epsilon_\gamma-\sum_{B_-}|\epsilon_\gamma|=-\sum_A\epsilon_\alpha+\sum_B\epsilon_\beta=0,
\]
so at $i=1$ the regularized value drops out and $\Lambda_{n,1}=-\sum_\gamma s_\gamma|\epsilon_\gamma|\,\Th{1}{\IZ}{0}{N_\gamma(n)}$, while for $i\ge2$, $\Lambda_{n,i}$ is $\Tx{i}{\IZ}{0}$ times $\frac1i\sum_\gamma s_\gamma|\epsilon_\gamma|^i+[i\text{ even}]\frac2i(|A_+|-\sum_{B_+}\epsilon_\gamma^i)$, less $\frac1i\sum_\gamma s_\gamma|\epsilon_\gamma|^i\Th{i}{\IZ}{0}{N_\gamma(n)}$. These coefficients depend on the slopes alone.

For $n,n+1>n_1$, $N_\gamma(n+1)=N_\gamma(n)+|\epsilon_\gamma|$, so $N_\gamma(n+1)!/N_\gamma(n)!=(N_\gamma(n)+1)_{|\epsilon_\gamma|}$ and
\[
\rho(n)=\frac{C_{n+1}}{C_n}=(-1)^{|A_+|+\sum_{B_+}\epsilon_\gamma}\prod_\gamma\big(N_\gamma(n)+1\big)_{|\epsilon_\gamma|}^{\,s_\gamma},
\]
a rational function with integer coefficients; with $\rho=\rho_1/\rho_2$ in polynomials, $\rho_2(n)C_{n+1}-\rho_1(n)C_n=0$ is a recurrence of order one with both coefficients nonzero, so the tier is at most one (\hb{prop:frame}). By Stirling's formula $\log N!=N\log N-N+O(\log(N+2))$, and with $N_\gamma(n)=|\epsilon_\gamma|n+O(1)$,
\[
\log|C_n|=\sum_\gamma s_\gamma\big(|\epsilon_\gamma|n\log n+|\epsilon_\gamma|n\log|\epsilon_\gamma|-|\epsilon_\gamma|n\big)+O(\log n)=n\log r_W+O(\log n),
\]
the terms in $n\log n$ and in $n$ cancelling by $\sum_\gamma s_\gamma|\epsilon_\gamma|=0$, and $|\epsilon_\gamma|\log|\epsilon_\gamma|=0$ on $A$. So $|C_n|^{1/n}\to r_W$.

The residue at $-n$ is the coefficient of $\varepsilon^{-1}$, $C_n[\varepsilon^{w-1}]\exp(\sum_i\Lambda_{n,i}\varepsilon^i)$, which is $0$ when $w\le0$, the exponential having no negative powers. For $w\ge1$ it is $C_n\sum\prod_i\Lambda_{n,i}^{m_i}/m_i!$ over $\sum_iim_i=w-1$, a polynomial with rational coefficients in $\Lambda_{n,1},\dots,\Lambda_{n,w-1}$; inserting the combinations above gives the stated polynomial, each monomial of total index $w-1$.

(iii) Moving the line to the left, the arguments of the Gammas of $A_+$ fall and cross their poles, those of $A_-$ rise and cross none, and the lower Gammas contribute entire factors; under the hypothesis of the closing, \hb{thm:closing}, with the Laurent coefficients at coincident poles, gives the integral as the sum of the residues at the integers $-n<c$. At each of the finitely many integers $-n<c$ with $n\le n_1$ the residue is a Laurent coefficient of finitely many Gammas at integer arguments: by \hb{prop:resonance} and \hb{prop:headstails} each is a nonzero rational number times the exponential of a series whose coefficient of $\varepsilon^i$ is a rational combination of $\Tx{i}{\IZ}{0}$ and of heads at integer cuts, which are rational; the coefficient of $\Tx{1}{\IZ}{0}$ is $-\epsilon_\alpha\varepsilon$ for an upper Gamma and $\epsilon_\beta\varepsilon$ for a lower one, whatever the sign of the argument, and these cancel by the balance. So these residues are rational combinations of products of $\Tx{i}{\IZ}{0}$, $i\ge2$. For $n>n_1$, expand the polynomial of (ii) into monomials: each is a rational number times a product of Trailmoments $\Tx{i}{\IZ}{0}$ times $\psi(n)$ for one of the finitely many products $\psi$ of the statement, the heads of the monomial having total index at most $w-1$. Since $|\Th{i}{\IZ}{0}{N}|\le1+\log(N+1)$ for $i\ge1$ and $N_\gamma(n)$ grows linearly, $|\psi(n)|\le|C_n|\,(c'\log n)^{w}$ for large $n$, of root rate $r_W<1$, so $\sum_{n>n_1}|\psi(n)|<\infty$ (\hb{def:conv}(viii)); the sum of the residues, a finite sum of absolutely convergent series, is the stated combination.

Each $\psi$ is live: $C_n$ is a product of factorials of the moving data, a stacked Gamma, and the heads are Trailmoments at the moving cut. For the recurrence put $H_l(n)=\Th{i_l}{\IZ}{0}{N_{\gamma_l}(n)}$; then $H_l(n+1)=H_l(n)+\eta_l(n)$ with $\eta_l(n)=\sum_{N_{\gamma_l}(n)<j\le N_{\gamma_l}(n)+|\epsilon_{\gamma_l}|}j^{-i_l}$, a rational function of $n$ with rational coefficients. The vector $X(n)=\big(C_n\prod_{l\in T}H_l(n)\big)_{T\subseteq\{1,\dots,q\}}$ satisfies
\[
X(n+1)_T=\rho(n)\sum_{U\subseteq T}\Big(\prod_{l\in T\setminus U}\eta_l(n)\Big)X(n)_U,
\]
a matrix $\rho(n)$ times a unitriangular one in the order of inclusion, of determinant $\rho(n)^{2^q}\neq0$. Lemma~\ref{lem:rayarith-cyclic} with $\ell$ the coordinate of $T=\{1,\dots,q\}$ gives the recurrence of order at most $2^q$ with nonzero extreme coefficients.

(iv) At the steps $k$ and $k'$, $N_\gamma(n)$ differs by $-(k-k')u_\gamma$ on $A_+\cup B_+$ and by $(k-k')u_\gamma$ on $A_-\cup B_-$, a constant $t_\gamma$. So the sign of $C_n$ changes by a constant, and $N_\gamma(n)!$ at $k$ is $N_\gamma(n)!$ at $k'$ times $(N_\gamma(n)+1)_{t_\gamma}$ or divided by $(N_\gamma(n)+t_\gamma+1)_{-t_\gamma}$, the data at $k'$; the quotient of the two $C_n$ is a rational function $R_0(n)$ with rational coefficients. Likewise $\Th{i}{\IZ}{0}{N+t}$ is $\Th{i}{\IZ}{0}{N}$ plus or minus the finite sum of $j^{-i}$ between $N$ and $N+t$, a rational function $\eta'(n)$. Hence, with $H'_l$ the heads at $k'$,
\[
\psi^{(k)}(n)=R_0(n)\,C^{(k')}_n\prod_l\big(H'_l(n)+\eta'_l(n)\big)=\sum_{T\subseteq\{1,\dots,q\}}R_0(n)\prod_{l\notin T}\eta'_l(n)\;C^{(k')}_n\prod_{l\in T}H'_l(n),
\]
and each $C^{(k')}_n\prod_{l\in T}H'_l(n)$ is an input of (iii) at the step $k'$, its heads a sub-multiset of those of $\psi$. Inserting this into (iii) at the step $k$, and adding the finitely many terms between the two bounds $n_1$, which are rational combinations of products of $\Tx{i}{\IZ}{0}$, $i\ge2$, by the argument for $n\le n_1$ in (iii), gives the last statement; the sums $\sum_nR(n)\psi^{(k')}(n)$ converge absolutely, $\psi^{(k')}$ having root rate $r_W<1$.

Closing to the right of $W_k(a)$ is closing to the left of $a\mapsto W_k(-a)$, the word with every slope negated, which is again balanced.
\end{proof}

\begin{hproposition}[\hb{prop:ra-paired} (Paired words)]
As stated.
\end{hproposition}

\begin{proof}
Write $N_\gamma(n)=|\epsilon_\gamma|n+e_\gamma$ with $e_\gamma\in\mathbb Z$. By the proof of \hb{thm:ra-oneshare}(ii),
\[
\rho(n)=\frac{C_{n+1}}{C_n}=(-1)^{|A_+|+\sum_{B_+}\epsilon_\gamma}\,r_W\prod_\gamma\prod_{t=1}^{|\epsilon_\gamma|}\big(n+\theta_{\gamma,t}\big)^{s_\gamma},\qquad\theta_{\gamma,t}=\frac{e_\gamma+t}{|\epsilon_\gamma|},
\]
using $(N+1)_v=v^v\prod_{t=1}^v(n+(e+t)/v)$ for $N=vn+e$ and $\prod_\gamma|\epsilon_\gamma|^{|\epsilon_\gamma|s_\gamma}=r_W$. For a fixed $\gamma$ with $|\epsilon_\gamma|=v$ the $\theta_{\gamma,t}$ are $v$ consecutive multiples of $\frac1v$, so modulo $\mathbb Z$ they meet each class $j/v$ once. For $v\ge1$ put $m_v=\sum_{|\epsilon_\gamma|=v}s_\gamma$, so that $m_v=\#\{\beta:\epsilon_\beta=v\}-\#\{\beta:\epsilon_\beta=-v\}$ for $v\ge2$, and $\sum_vv\,m_v=\sum_\gamma s_\gamma|\epsilon_\gamma|=0$ by the balance. The signed number of the $-\theta_{\gamma,t}$ in a class of $\mathbb Q/\mathbb Z$ of exact order $v'$ is $\sum_{v'\mid v}m_v$.

(b)$\Rightarrow$(a). If $C_n=(-1)^{\varsigma n}R(n)$, then $\rho(n)=(-1)^\varsigma R(n+1)/R(n)$ as rational functions. Writing $R=c\prod_x(n-x)^{\mu_x}$, $R(n+1)/R(n)=\prod_x\big((n-x+1)/(n-x)\big)^{\mu_x}$, whose zeros less poles, counted with multiplicity, add to zero in every class of $\mathbb C/\mathbb Z$. Comparing with the zeros and poles $-\theta_{\gamma,t}$ of $\rho$, of multiplicities $s_\gamma$, $\sum_{v'\mid v}m_v=0$ for every $v'\ge2$, and descending from the largest slope, $m_v=0$ for every $v\ge2$: that is (a). It does not involve $k$.

(a)$\Rightarrow$(b). Under (a), $m_v=0$ for $v\ge2$, hence $m_1=0$ by $\sum_vvm_v=0$, and $r_W=\prod_{v\ge2}v^{vm_v}=1$. Every class of $\mathbb Q/\mathbb Z$ then holds as many $\theta_{\gamma,t}$ with $s_\gamma=1$ as with $s_\gamma=-1$; pair them within classes, $\theta=\theta'+j$ with $j\in\mathbb Z$, and put $R_{\theta,\theta'}(n)=\Gamma(n+\theta)/\Gamma(n+\theta')$, which is $(n+\theta')_j$ or its reciprocal by \hb{prop:stacked}(iii), a rational function with $R_{\theta,\theta'}(n+1)/R_{\theta,\theta'}(n)=(n+\theta)/(n+\theta')$. With $\tilde R$ the product of the $R_{\theta,\theta'}$, $\rho(n)=(-1)^{\varsigma}\tilde R(n+1)/\tilde R(n)$, $\varsigma\equiv|A_+|+\sum_{B_+}\epsilon_\gamma$. For $n>n_1$ the numbers $n+\theta_{\gamma,t}$ are positive, so $\tilde R$ has no zero or pole there, and $C_n=(-1)^{\varsigma n}R(n)$ with $R(n)=(-1)^{\varsigma(n_1+1)}C_{n_1+1}\tilde R(n)/\tilde R(n_1+1)$. As $|A_+|+\sum_{B_+}\epsilon_\gamma=p+2\#\{\epsilon_\beta=1\}+\sum_{B^{\ge2}}\epsilon_\beta$, $\varsigma\equiv p+\sum_{B^{\ge2}}\epsilon_\beta$.

The factorization. By (i) of \hb{thm:ra-oneshare}, $W_k=W_0g_k$, so it suffices to factor $W_0$. For integers $b,b'$, $\Gamma(\pm a+b)/\Gamma(\pm a+b')$ and $\Gamma(\pm va+b)/\Gamma(\pm va)$ are rational functions of $a$ with rational coefficients (\hb{prop:stacked}(iii)). The balance and (a) give $|A_+|-|A_-|=\#\{\epsilon_\beta=1\}-\#\{\epsilon_\beta=-1\}$, so $p=|A_-|-\#\{\epsilon_\beta=-1\}$ as well. Matching the Gammas of $A_+$ with those of slope $1$ below, and the Gammas of $A_-$ with those of slope $-1$, leaves a rational function times $\Gamma(a)^p\Gamma(1-a)^p=\pi^p/\sin(\pi a)^p$ (\hb{prop:stacked}(v)), read as $\sin(\pi a)^{|p|}/\pi^{|p|}$ when $p<0$; each pair $\beta,\beta'$ with $\epsilon_\beta=v=-\epsilon_{\beta'}$, $v\ge2$, leaves a rational function times $1/(\Gamma(va)\Gamma(1-va))=\sin(\pi va)/\pi$. There are $|B^{\ge2}|$ pairs, and $w=|A_+|-|B_+|=p-|B^{\ge2}|$, which gives the display, the sign and $R_k$ collecting the rational factors.

The residue. At $a=-n+\varepsilon$, $\sin(\pi a)=(-1)^n\sin(\pi\varepsilon)$ and $\sin(\pi va)=(-1)^{vn}\sin(\pi v\varepsilon)$, so the trigonometric factor is $(-1)^{n(\sum_{B^{\ge2}}\epsilon_\beta-p)}\pi^w\prod_\beta\sin(\pi\epsilon_\beta\varepsilon)/\sin(\pi\varepsilon)^p$. Since $\sin(\pi x)/(\pi x)$ is an even power series in $\pi x$ with rational coefficients and constant term $1$, this is $(-1)^{\varsigma n}\varepsilon^{-w}\sum_{l\ge0}\tau_l\pi^{2l}\varepsilon^{2l}$ with rational $\tau_l$ depending on the slopes and $\tau_0=\prod_{B^{\ge2}}\epsilon_\beta$. For $n>n_1$ the pole of $W_k$ at $-n$ has the order $w$ exactly, $C_n\neq0$, and the trigonometric factor has a pole of order $w$ there, so $R_k$ is holomorphic at $-n$; the coefficient of $\varepsilon^{-1}$ is the displayed sum.

The closing. The finitely many residues at $-n<c$ with $n\le n_1$ are, by the same factorization, combinations of the $\tau_l\pi^{2l}$ with Laurent coefficients of $R_k$, so they lie in $\mathbb Q[\pi^2]$. For $n>n_1$ the residue is $\pm(-1)^{\varsigma n}r(n)$ with $r(n)=\sum_l\tau_l\pi^{2l}R_k^{(w-1-2l)}(-n)/(w-1-2l)!$, a rational function of $n$ whose partial fractions have coefficients in $\mathbb Q[\pi^2]$, the poles of $R_k$ being rational. The closing converges; when $\varsigma=0$ its terms are $r(n)$, so $r(n)=O(n^{-2})$, and when $\varsigma=1$, $r(n)\to0$, so $r(n_1+1+2m)-r(n_1+2+2m)=O(m^{-2})$, and summing consecutive pairs does not change the sum. \hb{lem:polecount} over $\Bbbk=\mathbb Q(\pi^2)$, whose identities are linear in the coefficients of the partial fractions (its proof applies \hb{prop:terminalrow} to them), writes the sum over $1$, the Trailmoments of the lattice inputs on the cosets of the poles at their cuts, and the differences of their index-one values, with coefficients in $\mathbb Q[\pi^2]$. Finally $\Tx{2}{\IZ}{0}=\pi^2/6$ by \hb{thm:hparity}(iv) at $j=1$.
\end{proof}

\begin{htheorem}[\hb{thm:ra-frameclearing} (The frame bounds the clearing from below)]
As stated.
\end{htheorem}

\begin{proof}
(i) As $p_h(k)\neq0$ for $k\ge k_0$, a solution is fixed by its values at $k_0,\dots,k_0+h-1$, so independent solutions have independent initial vectors and $C_{k_0}\neq0$; by \hb{prop:frame}, $C_{k+1}=(-1)^h\frac{p_0(k)}{p_h(k)}C_k$ with $p_0(k)\neq0$, so $C_k\neq0$ for every $k\ge k_0$. Then $D_k^{\,h}C_k=\det\big[D_kx^{(i)}_{k+l}\big]$ is a nonzero integer and $D_k^{\,h}|C_k|\ge1$. By \hb{prop:ra-framegrowth}, $\log|C_k|=k\log|c|+(d_0-d_h)(k\log k-k)+O(\log k)$, and $h\log D_k\ge-\log|C_k|$ gives the inequality. When $d_0=d_h$ the term in $k\log k$ is absent, and the same proposition identifies $|c|$.

(ii) The coefficients of $v_l$ and of $1$ in $\sum_jp_j(k)x_{k+j}=0$ are $\sum_jp_j(k)A_{k+j,l}$ and $-\sum_jp_j(k)B_{k+j}$, rational since $p_j\in\mathbb Q[k]$ and $k$ is an integer; as $1,v_1,\dots,v_m$ are independent over $\mathbb Q$, all of them vanish. A $D_k$ that clears the rows of $x_k,\dots,x_{k+h-1}$ makes the values at $k,\dots,k+h-1$ of all $m+1$ sequences integers, in particular those of $h$ independent ones, and (i) applies to them.
\end{proof}

\begin{hdefinition}[\hb{def:ra-merit} (The merit of rows along a ray)]
As stated.
\end{hdefinition}

\begin{proof}
$f(\alpha^\star)$ is positive and finite (\hb{thm:pencil}(i)), and $D_k\ge1$ gives $\limsup_k\frac1k\log D_k\ge0$, so the merit lies in $[-\infty,-\log f(\alpha^\star)]$.
\end{proof}

\begin{htheorem}[\hb{thm:ra-merit} (The merit decides the gap along a ray)]
As stated.
\end{htheorem}

\begin{proof}
(i) By \hb{thm:pencil}(i), $0<f\le f(\alpha^\star)$ on the orthant and $\langle1\rangle_k=\int_{y_{e_0}=0}f^k\omega_0$ converges and is positive for $k\ge k_1$; the origin is real, so $\omega_0/dy=\alpha^{\mvec_0}\prod_jK_j^{-s_{0j}}\Psi_j^{-t_{0j}}$ is positive on the orthant. Hence $\langle1\rangle_k=\int f^{k-k_1}f^{k_1}\omega_0\le f(\alpha^\star)^{k-k_1}\langle1\rangle_{k_1}$, and $D_k\ge1$.

(ii) If the merit is positive, choose $\eta>0$ below it; then $D_k\le f(\alpha^\star)^{-k}e^{-k\eta}$ for all large $k$, and by (i) $0<D_k\langle1\rangle_k\le e^{-k\eta}f(\alpha^\star)^{-k_1}\langle1\rangle_{k_1}\to0$. With $x_k=\langle1\rangle_k$, \hb{thm:gap} gives an irrational $v_l$, and for $m=1$, $A_{k,1}\neq0$ for all large $k$; then $v_1$ is the irrational one.

(iii) $\frac1k\log D_k\le\frac{\log C}k+\log\gamma+\frac wk\log d_{\lfloor ak+b\rfloor}$. For $a>0$, $N=\lfloor ak+b\rfloor\to\infty$ with $N/k\to a$, and $\limsup_N\frac1N\log d_N\le c_{\mathrm{lcm}}$ (\hb{lem:lcm}), so $\limsup_k\frac1k\log D_k\le\log\gamma+awc_{\mathrm{lcm}}$; for $a=0$ the last term tends to $0$. The merit is then at least $-\log f(\alpha^\star)-\log\gamma-awc_{\mathrm{lcm}}$, positive when $\gamma e^{awc_{\mathrm{lcm}}}f(\alpha^\star)<1$.

(iv) By \hb{thm:ra-frameclearing}(ii) and (i), $\liminf_k\frac1k\log D_k\ge\frac1h\log\frac1{|c|}$, so the $\limsup$ is at least as large, and the merit is at most $-\log f(\alpha^\star)+\frac1h\log|c|$.
\end{proof}

\begin{hproposition}[\hb{prop:ra-lowend} (The low end of the law)]
As stated.
\end{hproposition}

\begin{proof}
(i) The law says that $\sum_ip_i(k)f^i$ lies in the relation space $\mathcal L$ (\hb{def:raygauge}), a finite sum $\sum r_{e,l}(k)(c^{(0)}_e+kc^{(1)}_e+\partial_{y_e})g_{e,l}$ with $r_{e,l}\in\mathbb C(k)$. Integration by parts against $f^kdy$, for complex $k$ of large real part where the integrals converge absolutely (\hb{thm:raylaws}), gives $\sum_ip_i(k)\langle1\rangle_{k+i}=0$ there off the poles of the $r_{e,l}$, hence identically for the meromorphic continuation of \hb{prop:principal}. At an integer $k\ge1$ every $\langle1\rangle_{k+i}$ is a convergent integral, the walls $(k+i)\kappa_{1,S}$ being positive, and the law holds there. Near $k=0$ the terms $i\ge1$ are holomorphic, and by \hb{prop:principal} at $j=0$, with $j(1)=0$ and $\mathsf C_0(1)=\CT(1)=1$, $\langle1\rangle_k=V_1k^{-(n-1)}+O(k^{-(n-2)})$ with $V_1>0$. So $p_0(k)\langle1\rangle_k$ is holomorphic at $0$, which forces $p_0$ to vanish there to the order $n-1$ at least, and its value at $0$ is $\varpi V_1$. The law at $k=0$ is the low-end relation. The walls of the integral direction are integers, $\kappa_{1,S}=m_{1,S}-\sum_j(s_{1j}d_j(S)+t_{1j}d^\ast_j(S))$ (\hb{def:family}), so $V_1$ is rational.

(ii) $c$ and $r$ exist and are rational because $p_2(k)\neq0$ for $k\ge1$ and $p_0,p_1,p_2,\varpi V_1$ are rational. The low-end coefficient of $c$ is $-p_1(0)-p_2(0)\big(-p_1(0)/p_2(0)\big)=0$ and that of $r$ is $-p_2(0)\big(-\varpi V_1/p_2(0)\big)=\varpi V_1$. The sequence $x_k-x_1c_k-r_k$ solves the law on $k\ge1$; it vanishes at $k=1$, and at $k=2$ it is $\big(p_2(0)x_2+p_1(0)x_1+\varpi V_1\big)/p_2(0)=0$ by the low-end relation; so it vanishes identically.
\end{proof}

\begin{hdefinition}[\hb{def:ra-reflexive} (Reflexive laws)]
As stated.
\end{hdefinition}

\begin{proof}
Let $c_i$ be the coefficient of $k^d$ in $p_i$, so the symbol is $\sigma(\varrho)=\sum_ic_i\varrho^i$ with $c_0c_h\neq0$ (\hb{cor:symbol}). The coefficient of $k^d$ in $p_{h-i}(-k-h+1)$ is $(-1)^dc_{h-i}$, so $(-1)^dc_{h-i}=\gamma\lambda^ic_i$, and
\[
\varrho^h\sigma(\lambda/\varrho)=\sum_ic_i\lambda^i\varrho^{h-i}=(-1)^d\gamma^{-1}\sum_ic_{h-i}\varrho^{h-i}=(-1)^d\gamma^{-1}\sigma(\varrho).
\]
As $\sigma(0)=c_0\neq0$, $\varrho\mapsto\lambda/\varrho$ maps the roots of $\sigma$ to roots, with multiplicities. At $h=2$ the identities at $i=0$ and $i=2$ give $c_0/c_2=(-1)^d\gamma^{-1}$ and $c_0/c_2=(-1)^d\gamma\lambda^2$, so $\gamma^2\lambda^2=1$, $\gamma=\pm\lambda^{-1}$, and the product of the roots, $c_0/c_2$, is $\pm\lambda$.
\end{proof}

\begin{htheorem}[\hb{thm:ra-closingcoordinate} (The closings read the coordinate of a reflexive law)]
As stated.
\end{htheorem}

\begin{proof}
(i) By \hb{cor:reflected}, $\sum_{i=0}^{\min(h,j)}p_i(-j)\mathsf C_{j-i}(1)=0$ for every $j\ge0$. For $k\ge1$ put $j=k+h-1\ge h$. Reflexivity, $p_i(k)=\gamma^{-1}\lambda^{-i}p_{h-i}(-k-h+1)$, gives
\[
\sum_{i=0}^hp_i(k)\,\lambda^{k+i-1}\mathsf C_{k+i-1}(1)=\gamma^{-1}\lambda^{k-1}\sum_{i=0}^hp_{h-i}(-j)\,\mathsf C_{j-(h-i)}(1)=0 .
\]
For the low-end coefficient, reflexivity at $k=0$ gives $p_i(0)=\gamma^{-1}\lambda^{-i}p_{h-i}(-h+1)$, so
\[
\sum_{i=1}^hp_i(0)\,\lambda^{i-1}\mathsf C_{i-1}(1)=\gamma^{-1}\lambda^{-1}\sum_{i'=0}^{h-1}p_{i'}(-(h-1))\,\mathsf C_{h-1-i'}(1)=0,
\]
by \hb{cor:reflected} at $j=h-1$.

(ii) At $h=2$, $y_1=\mathsf C_0(1)=1=c_1$, and the vanishing of the low-end coefficient, $p_1(0)+p_2(0)y_2=0$, gives $y_2=-p_1(0)/p_2(0)=c_2$; both solve the law and $p_2(k)\neq0$ for $k\ge1$, so $y=c$, and \hb{prop:ra-lowend}(ii) gives the display. With $\langle1\rangle_1=A_1v-B_1$ and $r_k$ rational, $\langle1\rangle_k=A_1\lambda^{k-1}\mathsf C_{k-1}(1)\,v-\big(B_1\lambda^{k-1}\mathsf C_{k-1}(1)-r_k\big)$.

(iii) With integer weights, $f^{-1}=\alpha^{-\mvec_1}\prod_jK_j^{s_{1j}}\Psi_j^{t_{1j}}$ is a Laurent polynomial with integer coefficients, and so is each power; at origin $0$, $\omega_0/dy=1$, and $\mathsf C_j(1)=\CT(f^{-j})\in\mathbb Z$. The coefficient is $a^{k-1}\mathsf C_{k-1}(1)/b^{k-1}$; as $\gcd(a,b)=1$, its reduced denominator is $b^{k-1}/\gcd\big(b^{k-1},a^{k-1}\mathsf C_{k-1}(1)\big)=b^{k-1}/\gcd\big(b^{k-1},\mathsf C_{k-1}(1)\big)$.
\end{proof}

\begin{hcorollary}[\hb{cor:ra-meritbound} (The merit of a reflexive ray of order two)]
As stated.
\end{hcorollary}

\begin{proof}
The ray has origin $0$, real, and $\kappa_{1,S}>0$, so Definition~\hb{def:ra-merit} applies with $m=1$, $v_1=\langle1\rangle_1$, $A_{k,1}=\lambda^{k-1}\mathsf C_{k-1}(1)$ and $B_k=-r_k$. A clearing $D_k$ makes $D_kA_{k,1}$ an integer, so by \hb{thm:ra-closingcoordinate}(iii) the integer $b^{k-1}/g_k$, $g_k=\gcd(b^{k-1},\mathsf C_{k-1}(1))$, divides $D_k$, and $\frac1k\log D_k\ge\frac{k-1}k\log b-\frac1k\log g_k$. Taking the upper limit, $\limsup_k\frac1k\log D_k\ge\log b-\liminf_k\frac1k\log g_k$, which is the bound.
\end{proof}

\begin{hproposition}[\hb{prop:ra-hingerow} (The hinge closing as a row over the room)]
As stated.
\end{hproposition}

\begin{proof}
(i) By \hb{cor:closetwosides}(i) the closed value is the absolutely convergent sum of the residues, and grouping them by $n$ gives $\sum_n\ip{\Phi_n}{\mathcal X_n}$. Then $\ip{\Phi_n}{\mathcal X_n}=\sum_i\chi_i\ip{\Phi_n}{\mathbf c_i(n)}$, and a finite sum of absolutely convergent series may be summed term by term.

(ii) By \hb{cor:closetwosides}(ii) the two lattices of the level $d$ meet exactly when $2\delta^d\in\mathbb Z$, so otherwise every pole is simple and the residue at $\sigma^d=\pm\delta^d-n$ is a coefficient times the room period there. The sequence $\sum_ix^\pm_ic^\pm_i(n)$ solves the recurrence and agrees with $x^\pm$ at $0,\dots,h_\pm-1$; as $p^\pm_{h_\pm}(n)\neq0$ for $n\ge0$, the two agree for every $n$. The $c^\pm_i$ are computed from rational initial values by the recurrence, so they are rational when the coefficients are. So $\mathcal X_n=\sum_ix^+_i\big(c^+_i(n),0\big)+\sum_ix^-_i\big(0,c^-_i(n)\big)$, and (i) applies.

(iii) For $2\delta^d\in\mathbb Z$ the lattices $\delta^d-\mathbb N$ and $-\delta^d-\mathbb N$ are both contained in $|\delta^d|-\mathbb N$, and by \hb{cor:closetwosides}(ii) the residue at a common point carries the room period and its derivative in $\sigma^d$, at a simple point the room period alone. The room period converges at the points of the lattice (\hb{cor:closetwosides}(i)), and near them it is an absolutely convergent integral over the simplex of a function holomorphic in the level shares (\hb{thm:projective}), so it is holomorphic in the even share there. The state $U_n(\sigma)=(x_n(\sigma),\dots,x_{n+h-1}(\sigma))$ satisfies $U_{n+1}(\sigma)=\mathbf A(n,\sigma)U_n(\sigma)$ with the companion matrix of the recurrence, whose entries $-p_i(n,\sigma)/p_h(n,\sigma)$ are regular at $\sigma=\sigma_0$ for $n\ge0$. By \hb{prop:jetrecurrence} with $\varepsilon=\sigma-\sigma_0$ and $J=1$, $U_{n+1,0}=\mathbf A_0(n)U_{n,0}$ and $U_{n+1,1}=\mathbf A_0(n)U_{n,1}+\mathbf A_1(n)U_{n,0}$, so $(U_{n,0},U_{n,1})$ is a product of block-triangular matrices, which depend on the recurrence alone, applied to $(U_{0,0},U_{0,1})$, the $2h$ initial data. This writes $\mathcal X_n$ as in (i), which applies when the resulting sums converge absolutely.

With the bunch and the room exchanged, \hb{cor:closetwosides}(i) closes on the lattices of the level $r$ with the opposite sign, and the same argument applies.
\end{proof}

\begin{hdefinition}[\hb{def:ra-criticalprime} (Critical primes)]
As stated.
\end{hdefinition}

\begin{proof}
A native gauge lies in $X_F$, so every letter is nonzero there, and $f(\alpha_c)\in\Bbbk^\times$ since the weights and the coordinates lie in $\Bbbk$; its valuation at $\mathfrak p$ is defined.
\end{proof}

\begin{htheorem}[\hb{thm:ra-degenerate} (The letters degenerate at a critical prime)]
As stated.
\end{htheorem}

\begin{proof}
Let $\mathcal O$ be the local ring of $\Bbbk$ at $\mathfrak p$, a discrete valuation ring, with a uniformizer $\varpi_{\mathfrak p}$, and $v=v_{\mathfrak p}$. The letters have coefficients in $\mathcal O$ and the coordinates of $\alpha_c$ are units, so every $L(\alpha_c)$ and every $\partial_eL(\alpha_c)$ lies in $\mathcal O$, and $L(\alpha_c)\neq0$. Since $\alpha_c^{\mvec_1}$ is a unit, $v(f(\alpha_c))=-\sum_Lw_Lv(L(\alpha_c))<0$, so the largest valuation $v_{\max}$ is positive, and the letters of $S$ vanish at $\bar\alpha_c$. The native gauge has $c^{(1)}_e=\partial_{y_e}\log f=0$ for $e\neq e_0$ (\hb{def:raygauge}, \hb{lem:insertions}(ii)), that is $m_{1,e}=\sum_Lw_L\,\alpha_e\partial_eL/L$ at $\alpha_c$; by Euler's identity for the homogeneous letters, $\sum_e\alpha_e\partial_eL=(\deg L)L$, and the balance $m_{1,\cE}=\sum_Lw_L\deg L$ of a direction ($f$ has degree $0$), the same holds at $e=e_0$. Write $L(\alpha_c)=\varpi_{\mathfrak p}^{v_L}u'_L$ with $u'_L\in\mathcal O^\times$ and multiply the equation of the mark $e$ by $\varpi_{\mathfrak p}^{v_{\max}}$:
\[
\sum_Lw_L\,\alpha_{c,e}\,\partial_eL(\alpha_c)\,\varpi_{\mathfrak p}^{v_{\max}-v_L}u_L'^{-1}=m_{1,e}\,\varpi_{\mathfrak p}^{v_{\max}} .
\]
Every term lies in $\mathcal O$; modulo $\mathfrak p$ the right side and the terms with $v_L<v_{\max}$ vanish, and reduction commutes with $\partial_e$. Dividing by the unit $\bar\alpha_{c,e}$ gives (i) with $u_L=u_L'^{-1}$. Under the hypothesis of (ii) every coefficient $\overline{w_Lu_L}$ is nonzero, so the gradients of the reduced letters of $S$ at $\bar\alpha_c$ satisfy a linear relation with all coefficients nonzero: for one letter its gradient vanishes at a point of its zero set, which is singular there, and for several the gradients are linearly dependent, whereas smooth hypersurfaces meeting transversally have independent gradients.
\end{proof}

\begin{hproposition}[\hb{prop:ra-symbolprimes} (Critical primes and the symbol)]
As stated.
\end{hproposition}

\begin{proof}
(i) $c_h^{h-1}\sigma(\lambda)=\sum_jc_jc_h^{h-1-j}(c_h\lambda)^j$ is a monic polynomial with integer coefficients in $c_h\lambda$, which vanishes at a root.

(ii) If $p\nmid c_h$, $c_h$ is a unit at every prime $\mathfrak P$ of $K$ above $p$ and $\lambda=c_h\lambda/c_h$ has $v_{\mathfrak P}(\lambda)\ge0$ by (i). Conversely let $p\mid c_h$ and let $\mathfrak P$ be a prime above $p$. If every root had $v_{\mathfrak P}\ge0$, the elementary symmetric functions of the roots, $\pm c_{h-j}/c_h$, would be $\mathfrak P$-integral, so $v_{\mathfrak P}(c_{h-j})\ge v_{\mathfrak P}(c_h)>0$ for every $j$, and $p$ would divide every coefficient, against their coprimality. So some root has negative valuation at $\mathfrak P$.

(iii) Under the hypothesis, a root $\lambda=f(\alpha_c)$ with $v_{\mathfrak P}(\lambda)<0$ at a prime $\mathfrak P$ above $p$ makes $\mathfrak P$ a critical prime of $\alpha_c$ (\hb{def:ra-criticalprime} with $\Bbbk=K$), and conversely; (ii) gives the equivalence. When $p_0$ and $p_h$ have the largest degree $d$, $c_0$ and $c_h$ are their leading coefficients (\hb{cor:symbol}), and $c_0/c_h$ is the quotient $c$ of \hb{prop:frame}.
\end{proof}

\section{Chapter~\hb{ch:relations}: relations and separation}\label{sec:ch36}

\begin{hproposition}[\hb{prop:relationmatrix} (A finite matrix of mark relations)]
As stated, for several letters $D_1,\dots,D_h$.
\end{hproposition}

\begin{proof}
In the coordinate $x_i$, with the other coordinates as room and the room density $\prod_{l\neq i}x_l^{a_l-1}(1-x_l)^{b_l-1}$, \hb{thm:markrelation}(a) applies to the letters $D_k$, positive on the closed box, at the exponents $p_k+1$, and to the polynomial $AD$. With $\partial_iD=\sum_k\partial_iD_k\prod_{k'\neq k}D_{k'}$ its expression is
\begin{multline*}
x_i(1-x_i)(D\,\partial_iA+A\,\partial_iD)+\big(a_i-(a_i+b_i)x_i\big)AD-x_i(1-x_i)A\sum_k(p_k+1)\,\partial_iD_k\prod_{k'\neq k}D_{k'}\\
=D\big[x_i(1-x_i)\,\partial_iA+\big(a_i-(a_i+b_i)x_i\big)A\big]-x_i(1-x_i)\,A\sum_kp_k\,\partial_iD_k\prod_{k'\neq k}D_{k'}=\operatorname{mr}_i(A),
\end{multline*}
and the read is the stated integral. At $A=x^\beta$, $x_i(1-x_i)\partial_ix^\beta+(a_i-(a_i+b_i)x_i)x^\beta=\big((\beta_i+a_i)-(\beta_i+a_i+b_i)x_i\big)x^\beta$, and $x_i(1-x_i)x^\beta\sum_kp_k\partial_iD_k\prod_{k'\neq k}D_{k'}=(1-x_i)x^\beta\sum_\gamma D^{(i)}_\gamma x^\gamma$; expanding $D=\sum_\gamma D_\gamma x^\gamma$ gives the formula. Its coefficients are rational in $(a_i,b_i)$ over the rational coefficients of the $D_k$ and the $p_k$, so for a degree bound the vectors $\operatorname{mr}_i(x^\beta)$ form a finite matrix, rational at rational $a_i,b_i$; by linearity a numerator $\sum_i\operatorname{mr}_i(A_i)$ has the read $0$. With one letter, $D^{(i)}_\gamma=p\gamma_iD_\gamma$, the formula of the reference.
\end{proof}

\begin{hproposition}[\hb{prop:criticalstage} (The critical polynomial block)]
As stated, at every order $k\ge1$.
\end{hproposition}

\begin{proof}
For $R(u)=P(2u)$, $R(\frac12)=P(1)$ and $Q_P$ is the polynomial $(R(u)-R(\frac12))/(1-2u)$, so \hb{prop:critical} at the order $k$ gives the display. $Q_P^{(2k-1)}$ has rational coefficients and the moments of $X=\frac12+iY$ are rational (\hb{thm:sharelaw}(iii)), so $h(P)\in\mathbb Q$. The value of $c+\operatorname{crit}_k[P(2u)]$ is $\big(c+h(P)\big)+(2^{2k+1}-1)P(1)\,\Tx{2k+1}{\IZ}{0}$; the coordinate map is onto $\mathbb Q^2$, the constant $c$ giving the first coordinate and a constant $P$, for which $Q_P=0$, the second; and $a+b\,\Tx{2k+1}{\IZ}{0}=0$ with $(a,b)\in\mathbb Q^2$ not zero forces $b\neq0$ and $\Tx{2k+1}{\IZ}{0}=-a/b$. At the cut $N$, $\Tx{2k+1}{\IZ}{N}=\Tx{2k+1}{\IZ}{0}-\Th{2k+1}{\IZ}{0}{N}$, a rational head.
\end{proof}

\section{Chapter~\hb{ch:content}: the content and its coaction}\label{sec:chcontent}

Every identity proved below between admitted words is a proved relation, so it lies in the relation ideal $\mathfrak I$ (\hb{def:words}) and holds between the classes in $\mathfrak W/\mathfrak I$; in particular the words $\zeta=e^{2\pi ir/M}$ multiply as the roots of unity do, by the functional equation of the exponential, and a constant factor $\zeta$ in front of a trail is a factor of its density, by linearity.

\begin{hdefinition}[\hb{def:ct-content} (The content)]
As stated.
\end{hdefinition}

\begin{proof}
For $\nu_d\ge2$ the trail converges absolutely: its terms are bounded by $\prod_in_i^{-\nu_i}$, and $\sum_{0<n_1<\cdots<n_d}\prod_in_i^{-\nu_i}\le\sum_{n_d}(1+\log n_d)^{d-1}n_d^{-2}<\infty$; so the limit of the heads is the Trailmoment. For $\nu_d=1$ and $e^{2\pi ia_d}\neq1$ the limit exists by \hb{thm:ct-trailchain}. The data $(\boldsymbol\nu,\mathbf a)$ are finite and rational, so the trails and the words $\zeta$ are admitted words, and the spans are defined.
\end{proof}

\begin{hproposition}[\hb{prop:ct-algebra} (The content is an algebra)]
As stated.
\end{hproposition}

\begin{proof}
(ii) For an integer $N\ge1$ let $H_N(T)$ be the head of $T$ over $\{0<y_1<\cdots<y_d\le N\}$. The product $H_N(T)H_N(T')$ is the finite sum, over pairs of increasing tuples $0<n_1<\cdots<n_d\le N$ and $0<n'_1<\cdots<n'_{d'}\le N$, of $\prod_ie^{2\pi ia_in_i}n_i^{-\nu_i}\prod_je^{2\pi ia'_jn'_j}{n'_j}^{-\nu'_j}$. The distinct values among the $n_i$ and $n'_j$, in increasing order $m_1<\cdots<m_e$, and the positions of the points of each tuple among them give a surjection $\phi$, increasing on each part; conversely $\phi$ and $m_1<\cdots<m_e$ give back the pair. The term is $\prod_re^{2\pi i(\sum_{\phi(x)=r}a_x)m_r}m_r^{-\sum_{\phi(x)=r}\nu_x}$, the term of $T_\phi$ at $(m_r)_r$. So $H_N(T)H_N(T')=\sum_\phi H_N(T_\phi)$. Every $T_\phi$ is a trail of level $M$: its frequencies lie in $\frac1M\mathbb Z$, and since $\phi$ is onto and increasing on each part, its last point holds the last point of $T$ or of $T'$; if it holds only one of them it has that trail's last index and frequency, and if both, its index is at least $2$. As $N\to\infty$ every head converges (\hb{def:ct-content}), and $TT'=\sum_\phi T_\phi$. The depth $e$ lies between $\max(d,d')$ and $d+d'$ and the indices add to $|\boldsymbol\nu|+|\boldsymbol\nu'|$.

(i) and (iii). $1$ is the word $\zeta=1$. At level $1$ the frequencies are integers, so the trails of depth one of level $1$ are the $\Tx{k}{\IZ}{0}$ with $k\ge2$, the index one being excluded by $e^{2\pi ia_1}=1$. For $\zeta,\zeta'\in\mu_M$, $(\zeta T)(\zeta'T')=\zeta\zeta'\sum_\phi T_\phi$ by (ii), with $\zeta\zeta'\in\mu_M$, so $\mathcal C^{(M)}$ is closed under products, with weights adding, and it contains every product of the $\Tx{k}{\IZ}{0}$. When $M\mid M'$, $\frac1M\mathbb Z\subseteq\frac1{M'}\mathbb Z$ and $\mu_M\subseteq\mu_{M'}$, so $\mathcal C^{(M)}\subseteq\mathcal C^{(M')}$; as $\mathcal C^{(M)}$ and $\mathcal C^{(M')}$ both lie in $\mathcal C^{(MM')}$, the union $\mathcal C$ is an algebra.
\end{proof}

\begin{hproposition}[\hb{prop:ct-level} (Cosets, cuts and index one)]
As stated.
\end{hproposition}

\begin{proof}
(i) \emph{Cosets.} Put $n_i=My_i$, so that $y_i\in b_i+\mathbb Z$ exactly when the integer $n_i$ is congruent to $Mb_i$ modulo $M$, and the indicator of that congruence is $\frac1M\sum_{s=0}^{M-1}e^{2\pi is(n_i-Mb_i)/M}$ (the lattice relation of \hb{thm:latticefree}, read on one coset). Then $y_i^{-\nu_i}=M^{\nu_i}n_i^{-\nu_i}$ and $e^{2\pi ia_iy_i}=e^{2\pi i(a_i/M)n_i}$, and the heads of the trail up to $N'$ are, term by term,
\[
\sum_{\mathbf s\in\{0,\dots,M-1\}^d}M^{|\boldsymbol\nu|-d}\prod_ie^{-2\pi is_ib_i}\ \sum_{MN<n_1<\cdots<n_d\le MN'}\ \prod_ie^{2\pi i\frac{a_i+s_i}Mn_i}\,n_i^{-\nu_i}.
\]
Here $e^{-2\pi is_ib_i}=e^{-2\pi is_i(Mb_i)/M}\in\mu_M$, and the frequencies $(a_i+s_i)/M$ lie in $\frac1{MM'}\mathbb Z$. If $\nu_d=1$, then $a_d\notin\mathbb Z$, so $(a_d+s_d)/M\notin\mathbb Z$: every inner sum is the head of a trail on $\mathbb Z^d$ at the cut $c=MN\in\mathbb N$ whose last index and frequency satisfy the condition of \hb{def:ct-content}, and the trail is the finite combination of their limits.

\emph{Cuts.} For data $D=(\nu_i,a''_i)_{i\le d}$ on $\mathbb Z^d$ let $T_c(D)$ be the trail at the cut $c\in\mathbb N$, read as the limit of its heads, $T_c(\varnothing)=1$, and $D_{>k}$ the last $d-k$ data. Splitting the points $0<n_1<\cdots<n_d$ by the number $k$ of them at most $c$,
\[
T_0(D)=\sum_{k=0}^dh_k(D)\,T_c(D_{>k}),\qquad h_k(D)=\sum_{0<n_1<\cdots<n_k\le c}\ \prod_{i\le k}e^{2\pi ia''_in_i}n_i^{-\nu_i},
\]
first for the heads and then in the limit. So $T_c(D)=T_0(D)-\sum_{k\ge1}h_k(D)T_c(D_{>k})$, and by induction on the depth $T_c(D)$ is a polynomial in the finite sums $h_k(D_{>l})$ times $1$ or a trail $T_0(D_{>l})$, which keeps the last data of $D$ and so is a trail of level $MM'$. Each $h_k$ is a finite sum of rational multiples of roots of unity in $\mu_{MM'}$, a combination of the words $\zeta$ by linearity; with \hb{prop:ct-algebra} the class lies in $\mathcal C^{(MM')}$, and the value is as stated. A head of the trail is a finite sum of terms $e^{2\pi i\sum_ia_iy_i}\prod_iy_i^{-\nu_i}$ with $a_iy_i\in\frac1{MM'}\mathbb Z$, a rational multiple of an element of $\mu_{MM'}$.

(ii) By \hb{def:lattice}, $\Th{1}{\Ilat{b+\mathbb Z}}{N}{X}-\log(X-N)\to\Tx{1}{\Ilat{b+\mathbb Z}}{N}$ as $X\to\infty$, and $\log(X-N)-\log(X-N')\to0$, so the difference $D$ of the two index-one values is the limit of $\Th{1}{\Ilat{b+\mathbb Z}}{N}{X}-\Th{1}{\Ilat{b'+\mathbb Z}}{N'}{X}$. With $n=My$ and the indicator above,
\[
\Th{1}{\Ilat{b+\mathbb Z}}{N}{X}=\sum_{s=0}^{M-1}e^{-2\pi isb}\sum_{MN<n\le MX}\frac{e^{2\pi isn/M}}n .
\]
In the difference the terms $s=0$ leave the finite rational sum of $1/n$ between $MN$ and $MN'$, with its sign. For $s\neq0$, $\sum_{MN<n\le MX}e^{2\pi isn/M}/n$ is the head up to $MX$ of the trail $T_s$ of level $M$ of depth one, index $1$ and frequency $s/M$, with $e^{2\pi is/M}\neq1$, less the finite sum over $0<n\le MN$. So $D=\sum_{s=1}^{M-1}\big(e^{-2\pi isb}-e^{-2\pi isb'}\big)T_s$ plus finite sums of rational multiples of roots of unity in $\mu_M$, and $e^{-2\pi isb}\in\mu_M$ as $Mb\in\mathbb Z$.
\end{proof}

\begin{hdefinition}[\hb{def:ct-chain} (Chains)]
As stated.
\end{hdefinition}

\begin{proof}
No letter lies on the open segment, so the integrand is continuous on the open simplex, and its only possible singularities on the closure are at $t_i\in\{0,1\}$; the convergence is \hb{prop:ct-between}(i).
\end{proof}

\begin{hproposition}[\hb{prop:ct-chains} (The content as chains)]
As stated.
\end{hproposition}

\begin{proof}
(i) is \hb{thm:ct-trailchain} with $e^{2\pi ia_i}\in\mu_M$: the letters $c_i^{-1}$ lie in $\mu_M$, the first letter is $c_1^{-1}\neq0$, the last is $0$ when $\nu_d\ge2$ and $c_d^{-1}=e^{-2\pi ia_d}\neq1$ when $\nu_d=1$, and no point of $\{0\}\cup\mu_M$ lies on the open segment from $0$ to $1$.

(ii) The ratios $c_i/c_{i+1}$ lie in $\mu_M$, so frequencies $a'_i\in\frac1M\mathbb Z$ with $e^{2\pi ia'_i}=c_i/c_{i+1}$ exist, and $e^{2\pi i(a'_i+\cdots+a'_d)}=c_i/c_{d+1}=c_i$. As $a_1\neq0$, $j_1=1$, and the word of (i) for the data $(\nu_i,a'_i)$, the blocks $(c_i^{-1},0^{\nu_i-1})$, is $a_1,\dots,a_w$. If $\nu_d=1$ then $j_d=w$ and $e^{2\pi ia'_d}=c_d=a_w^{-1}\neq1$, so the trail is of level $M$, and (i) gives the identity.

(iii) The values of $\mathcal C^{(M)}$ are the $\mathbb Q$-span of the $\zeta$ and the $\zeta T$; the $\mathbb Q$-span of $\mu_M$ is $\mathbb Q(\mu_M)$, and by (i) and (ii) the values of the trails of level $M$ are, up to sign, exactly the chains described, the weight $|\boldsymbol\nu|$ of a trail being the number $w$ of letters of its chain.
\end{proof}

\begin{hproposition}[\hb{prop:ct-between} (Chains with an end at zero)]
As stated.
\end{hproposition}

\begin{proof}
(i) Write $\phi_i(t)=\gamma'(t)/(\gamma(t)-a_i)$. As no letter lies on the open segment, $|\phi_i(t)|\le C/t$ if $a_i=a_0$, $|\phi_i(t)|\le C/(1-t)$ if $a_i=a_{w+1}$, and $|\phi_i|\le C$ otherwise. Split the simplex by the number $k$ of the $t_i$ below $\frac12$: the part is a product of $\{0<t_1<\cdots<t_k<\frac12\}$ and $\{\frac12\le t_{k+1}<\cdots<t_w<1\}$, on which the factors singular at the other end are bounded. On the first, put $G_j(t)=\int_{0<t_1<\cdots<t_j<t}\prod_{i\le j}|\phi_i(t_i)|\,dt$ for $t\le\frac12$, where $|\phi_i(s)|\le C(1+1/s)$; since $a_1\neq a_0$, $G_1(t)\le Ct$, and if $G_{j-1}(s)\le C_{j-1}s$ then $G_j(t)\le\int_0^tC_{j-1}s\,C(1+1/s)\,ds\le C_jt$; so the first factor is finite. The second is the same after $t\mapsto1-t$, with $a_w\neq a_{w+1}$. So the integral of the modulus is finite.

(ii) For $\tilde\gamma=\lambda\gamma+\eta$, $\tilde\gamma'(t)/(\tilde\gamma(t)-\lambda a_i-\eta)=\gamma'(t)/(\gamma(t)-a_i)$, and the open segment is mapped onto the open segment of the new ends.

(iii) The segment from $a_{w+1}$ to $a_0$ is $\tilde\gamma(t)=\gamma(1-t)$, with $\tilde\gamma'=-\gamma'$; the substitution $t_i=1-t'_{w+1-i}$ maps the simplex onto itself, and each of the $w$ factors changes sign.

(iv) For $s\in\mu_M$, (ii) with $\lambda=s^{-1}$ gives $\mathrm I(0;a_1,\dots,a_w;s)=\mathrm I(0;a_1/s,\dots,a_w/s;1)$, with letters in $\{0\}\cup\mu_M$, first letter $a_1/s\neq0$ and last letter $a_w/s\neq1$. For a chain from $s$ to $0$, (iii) gives $(-1)^w$ times a chain from $0$ to $s$ whose first letter is $a_w\neq0$ and whose last letter is $a_1\neq s$, and the first case applies. \hb{prop:ct-chains}(ii) and (iii) place the value in the values of $\mathcal C^{(M)}$.
\end{proof}

\begin{hdefinition}[\hb{def:ct-symbols} (Symbols of chains)]
As stated.
\end{hdefinition}

\begin{proof}
The variables are the formal chains of weight at least one with distinct ends; the expressions with no letter or with equal ends are not variables, so assigning them $1$ and $0$ involves no relation among the variables. The symbol of a trail depends on the frequencies only through the $c_i$, which do not change when a frequency moves by an integer, as the trail does not; and the realization converges by \hb{prop:ct-between}(i), no point of $\{0\}\cup\mu_M$ lying on the open segment between $0$ and a root of unity.
\end{proof}

\begin{hdefinition}[\hb{def:ct-cutting} (The cutting rule)]
As stated.
\end{hdefinition}

\begin{proof}
$\mathfrak S_M$ is a polynomial algebra on the formal chains, so a homomorphism is fixed by arbitrary values on them; the right side of the display is an element of $\mathfrak S_M\otimes\mathfrak S_M$, its degenerate entries read by the conventions of \hb{def:ct-symbols}. The rotation $\rho_\lambda$ and the reversal map formal chains to formal chains, up to sign, keeping the ends distinct, and the conventions to the conventions; $\rho_{\lambda^{-1}}$ and the reversal itself are their inverses.
\end{proof}

\begin{htheorem}[\hb{thm:ct-coaction} (The cutting rule is a coaction)]
As stated.
\end{htheorem}

\begin{proof}
All maps in (i) and (iii) are homomorphisms of algebras, so it suffices to compare them on a formal chain $x=\mathrm I^{\mathfrak s}(a_0;a_1,\dots,a_w;a_{w+1})$. For a set $U\subseteq\{1,\dots,w\}$ the \emph{gaps} of $U$ are the pairs $t<t'$ of consecutive elements of $U\cup\{0,w+1\}$; for $T\subseteq S$ put
\[
\tau(T,S)=\mathrm I^{\mathfrak s}(a_0;a_T;a_{w+1})\otimes\prod_{(t,t')\text{ gap of }T}\mathrm I^{\mathfrak s}\big(a_t;a_{S\cap(t,t')};a_{t'}\big)\otimes\prod_{(s,s')\text{ gap of }S}\mathrm I^{\mathfrak s}\big(a_s;a_{(s,s')};a_{s'}\big),
\]
with $a_{(s,s')}$ the letters strictly between the positions $s$ and $s'$, the degenerate entries read by the conventions.

First, for any pair of ends and letters, the display of \hb{def:ct-cutting} gives $\Delta$ of the corresponding expression even when it is not a formal chain. With no letter the expression is $1$, and the display has the single term $1\otimes1=\Delta1$. With equal ends and at least one letter the expression is $0$; in the display every term with a kept letter has a left factor with equal ends, which is $0$, and the term with no kept letter has the right factor equal to the expression itself, which is $0$; so the display gives $0=\Delta0$.

(i) $(\Delta\otimes\mathrm{id})\Delta x=\sum_S\Delta\,\mathrm I^{\mathfrak s}(a_0;a_S;a_{w+1})\otimes\prod_{\text{gaps of }S}(\cdots)$, and the display for the chain of the letters $a_S$ keeps the letters of $T\subseteq S$ and reads the letters of $S$ between consecutive kept ones: the sum is $\sum_{T\subseteq S}\tau(T,S)$. Next, $(\mathrm{id}\otimes\Delta)\Delta x=\sum_T\mathrm I^{\mathfrak s}(a_0;a_T;a_{w+1})\otimes\prod_{(t,t')}\Delta\,\mathrm I^{\mathfrak s}(a_t;a_{(t,t')};a_{t'})$, and by the remark above each factor is the display, a sum over the sets $S_{(t,t')}$ of letters kept inside the gap $(t,t')$. A family of such sets is the same as one set $S\supseteq T$, the middle factors are the chains of the letters of $S$ inside the gaps of $T$, and the right factors are the gaps of the $S_{(t,t')}$ inside the gaps of $T$, which are the gaps of $S$: the sum is again $\sum_{T\subseteq S}\tau(T,S)$. For the counit, $\varepsilon$ kills every formal chain, so in $(\varepsilon\otimes\mathrm{id})\Delta x$ only $S=\varnothing$ survives, with the right factor $x$; in $(\mathrm{id}\otimes\varepsilon)\Delta x$ a right factor contributes $1$ only when it has no letter, so only $S=\{1,\dots,w\}$ survives, with the left factor $x$.

(ii) The term of $S$ has the left weight $|S|$ and the right weight $w-|S|$, the runs between kept letters together holding the letters not kept. The terms $S=\{1,\dots,w\}$ and $S=\varnothing$ are $x\otimes1$ and $1\otimes x$, and every other term has $0<|S|<w$.

(iii) Multiplying every entry by $\lambda$ commutes with keeping a subsequence and with cutting out runs, so $\Delta\rho_\lambda x=(\rho_\lambda\otimes\rho_\lambda)\Delta x$. The reversal maps the term of $S$ for $x$ to the term of the mirror set $S^\ast=\{w+1-s:s\in S\}$ for the reversed chain: its left factor is $(-1)^{|S|}$ times the reversal of the kept chain read backwards, each run is reversed with the sign $(-1)^{\text{its length}}$, and the runs appear in the opposite order, which does not matter in the commutative algebra. The signs multiply to $(-1)^w$, the sign of $\varrho x$, so $\Delta\varrho x=(\varrho\otimes\varrho)\Delta x$.
\end{proof}

\begin{hdefinition}[\hb{def:ct-letter} (The letter of an input)]
As stated.
\end{hdefinition}

\begin{proof}
The Laplace transform beyond the cut $n_0-1$ is $\mathcal F_{n_0-1}(u)=\sum_{n\ge n_0}\psi(n)e^{-un}=L_\psi(e^{-u})$ (\hb{def:inputs}), and $L_\psi(1)=\sum_n\psi(n)$ where the series converges at $1$, which is the Trailmoment at index zero (\hb{def:ratword}). For $\psi(n)=e^{2\pi ian}$, $n\ge1$, $L_\psi(z)=\sum_{n\ge1}(e^{2\pi ia}z)^n=e^{2\pi ia}z/(1-e^{2\pi ia}z)$, with its pole at $e^{-2\pi ia}$, which is the letter $c_1^{-1}$ of the depth-one trail at that frequency in \hb{prop:ct-chains}(i).
\end{proof}

\begin{hproposition}[\hb{prop:ct-letter} (The letter of a live input)]
As stated.
\end{hproposition}

\begin{proof}
(i) $p_i(\theta-i)z^n=p_i(n-i)z^n$, so
\[
\sum_iz^{h-i}p_i(\theta-i)L_\psi=\sum_i\sum_{n\ge n_0}p_i(n-i)\psi(n)z^{n+h-i}=\sum_mz^{m+h}\sum_{i:\,m+i\ge n_0}p_i(m)\psi(m+i).
\]
For $m\ge n_0$ the inner sum is the recurrence, zero; the finitely many $m$ with $n_0-h\le m<n_0$ leave a polynomial $Q$ of degree less than $n_0+h$. Since $\theta^j$ is $z^j\frac{d^j}{dz^j}$ plus terms of lower order, the coefficient of $\frac{d^d}{dz^d}$ is $\sum_iz^{h-i}c_iz^d$, nonzero as some $c_i\neq0$; for $z\neq0$ it vanishes exactly when $\sigma(1/z)=0$.

(ii) For large $n$, $p_h(n)\neq0$, and the state $U(n)=(\psi(n),\dots,\psi(n+h-1))$ satisfies $U(n+1)=\mathbf A(n)U(n)$ with the companion matrix of the recurrence, whose entries $-p_i(n)/p_h(n)$ tend to $-c_i/c_h$, as $p_h$ has the largest degree: $\mathbf A(n)\to\mathbf A_\infty$, the companion matrix of $\sigma/c_h$, of spectral radius $\varrho$. For $\eta>0$ there is a norm on $\mathbb C^h$ whose operator norm of $\mathbf A_\infty$ is at most $\varrho+\eta/2$ (a vector norm adapted to the Jordan form of $\mathbf A_\infty$, which exists for every square matrix); then $\|\mathbf A(n)\|\le\varrho+\eta$ for $n\ge n_3$, and $\|U(n)\|\le\|U(n_3)\|(\varrho+\eta)^{n-n_3}$, which with the finitely many earlier terms gives the bound. So $L_\psi$ converges on $|z|<1/(\varrho+\eta)$ for every $\eta$. By (i) it satisfies, near any point of that disc other than $0$, a linear differential equation with polynomial coefficients and polynomial right side, whose leading coefficient vanishes only at $0$ and at the reciprocals of the nonzero roots of $\sigma$; the solutions of such an equation continue analytically along every path on which the leading coefficient does not vanish (the existence theorem for linear differential equations with holomorphic coefficients, which applies because the coefficients are polynomials and the leading one is nonzero along the path). The cut plane $\Omega$ is simply connected and star-shaped about $0$, and contains the disc; a path in $\Omega$ meets the singular points only inside the disc, at $0$, where $L_\psi$ is holomorphic, so $L_\psi$ continues along every path in $\Omega$, and by the monodromy theorem it defines a holomorphic function on $\Omega$.

(iii) Let the two recurrences have orders $h$ and $h'$, with companion matrices $\mathbf A(n)$ and $\mathbf A'(n)$, invertible as rational matrices since their determinants are $\pm p_0(n)/p_h(n)$ and $\pm p'_0(n)/p'_{h'}(n)$. The vector $X(n)=U(n)\otimes U'(n)$ satisfies $X(n+1)=(\mathbf A(n)\otimes\mathbf A'(n))X(n)$ for large $n$, with an invertible Kronecker product, and the product sequence is the coordinate of $X(n)$ that pairs the first coordinates. Lemma~\ref{lem:rayarith-cyclic} over the field $\Bbbk$ generated by the coefficients gives the recurrence of order at most $hh'$ with nonzero extreme coefficients.

(iv) By (ii), $|\psi(n)\psi'(n)|\le C(\varrho+\eta)^n(\varrho'+\eta)^n$, so the letter of $\psi\psi'$, the termwise product of the coefficients of $L_\psi$ and $L_{\psi'}$, converges on $|z|<1/(\varrho\varrho')$. By (ii), $L_\psi$ and $L_{\psi'}$ are holomorphic on the planes cut along the rays from the reciprocals of the nonzero roots of $\sigma$ and of $\sigma'$. Hadamard's multiplication theorem, which applies to two power series holomorphic on planes cut along the outward rays from finitely many nonzero points, makes the series of the termwise products holomorphic on the plane cut along the outward rays from the products of those points, the points $1/(\lambda\lambda')$.
\end{proof}

\begin{hproposition}[\hb{prop:ct-tiertwo} (The row of a hinge closing read by letters)]
As stated.
\end{hproposition}

\begin{proof}
(i) The $c^\pm_i$ solve the room's recurrence, whose last coefficient has the largest degree, so by \hb{prop:ct-letter}(ii) $|c^\pm_i(n)|\le C(\varrho^\pm+\eta)^n$; likewise $|\Phi^\pm_n|\le C(\varrho^\pm_\Phi+\eta)^n$ for large $n$. With $\eta$ small, $(\varrho^\pm+\eta)(\varrho^\pm_\Phi+\eta)<1$, and the sums converge absolutely; the radius of convergence of $L_{\psi^\pm_i}$ exceeds $1$, so $\Xi^\pm_i=L_{\psi^\pm_i}(1)$. \hb{prop:ra-hingerow}(ii) gives the row.

(ii) By \hb{prop:ct-letter}(iii) the product $\psi^\pm_i$ satisfies a recurrence with polynomial coefficients for large $n$, so by \hb{prop:ct-letter}(i) its letter, up to a polynomial from the finitely many first terms, satisfies a linear differential equation with polynomial coefficients; and \hb{prop:ct-letter}(iv), applied to the tails of the two factors, gives the continuation, the finitely many first terms adding a polynomial.
\end{proof}

\section{Chapter~\hb{ch:finitealgebra}: finite algebra for scale reads}\label{sec:ch37}

The letters of this chapter have complex coefficients and no roots on $[0,\infty)$. Two facts carry the proofs. First, along $[0,\infty)$ a letter $P_j$ has a continuous argument with limits at $0$ and at $\infty$, so for $x\to\infty$, $P_j(x)^{-w_j}=C_jx^{-w_j\deg P_j}(1+O(x^{-1}))$ with $C_j\neq0$, and $|P_j^{-w_j}|$ is $|P_j|^{-\re w_j}$ up to factors bounded above and below on $[0,\infty)$: a read converges absolutely exactly where it does for positive letters of the same degrees. Second, the reduction of \hb{thm:scalereduce} is algebraic in $\mathbb F'$ and does not use the values of the letters.

\begin{hdefinition}[\hb{def:insertions} (Insertions and reads on one scale)]
As stated, for letters without roots on $[0,\infty)$.
\end{hdefinition}

\begin{proof}
$P_j(0)\neq0$ and $P_j$ does not vanish on $[0,\infty)$, so a continuous logarithm of $P_j$ along $[0,\infty)$ starting at the principal logarithm of $P_j(0)$ exists and is unique, and it defines $P_j^{-w_j}$; on $(0,\infty)$ it is locally a holomorphic branch, so $\big(P_j^{-w_j}\big)'=-w_jP_j'P_j^{-w_j-1}$, and the product rule gives $\big(x^\kappa\prod_jP_j^{-w_j}g\big)'=x^{\kappa-1}\prod_jP_j^{-w_j}\,\eth g$. The squarefree decomposition and the gcd splitting are computed in $\mathbb F[x]$ by gcds and exact divisions, so the $q_k$ lie in $\mathbb F[x]$, and they divide $P_1\cdots P_N$, so none vanishes at $0$ or on $[0,\infty)$. For $\xi_\ell\in\mathbb C\setminus(-\infty,0]$, $x+\xi_\ell$ stays in $\mathbb C\setminus(-\infty,0]$ for $x\ge0$, so the continuous branch is the principal one, that of \hb{prop:scalebreta} and \hb{def:shifts}.
\end{proof}

\begin{htheorem}[\hb{thm:scalereduce} (Reduction by the scale exchange)]
As stated, for letters without roots on $[0,\infty)$.
\end{htheorem}

\begin{proof}
Write $\mu_k=\sum_jm_{jk}w_j$, so that $\eth=x\frac d{dx}+\kappa-\sum_k\mu_k\frac{xq_k'}{q_k}$ and $\mathfrak A=\mathbb F'[x,x^{-1},Q^{-1}]$. Every $R\in\mathfrak A$ is a Laurent polynomial in $x$ plus principal parts at the $q_k$.

\emph{Three reductions.} (a) If $R$ has a pole of order $h+1\ge2$ at $q_k$ with leading numerator $s$, $\deg s<\deg q_k$, then
\[
\eth\big(r\,q_k^{-h}\big)=-(h+\mu_k)\,\frac{xq_k'\,r}{q_k^{h+1}}+O\big(q_k^{-h}\big),
\]
plus simple poles at the other $q_{k'}$ and no polynomial part and no pole at $0$, $r\,q_k^{-h}$ vanishing at $\infty$ and being regular at $0$; $xq_k'$ is prime to $q_k$, $q_k$ being squarefree and $q_k(0)\neq0$, so the extended Euclidean algorithm gives $r$ with $-(h+\mu_k)xq_k'r\equiv s$ mod $q_k$, one division by $h+\mu_k$. (b) If $R$ has a pole $x^{-r}$, $r\ge1$, at $0$, $\eth(x^{-r})=(\kappa-r)x^{-r}-\sum_k\mu_kx^{1-r}q_k'/q_k$, whose other terms have poles of order below $r$ at $0$, simple poles at the $q_k$ and no polynomial part. (c) If the polynomial part of $R$ has degree $a\ge0$, $\eth(x^a)=\big(\kappa+a-\sum_jw_j\deg P_j\big)x^a$ plus a polynomial of degree below $a$ plus simple poles at the $q_k$, since $xq_k'/q_k=\deg q_k+O(x^{-1})$ and $\sum_k\mu_k\deg q_k=\sum_jw_j\deg P_j$. Subtracting $\eth$ of a multiple at each step lowers the order of the pole, resp.\ the degree, and leaves $\sum_{i<D}c_ix^i/Q$ once the poles at the $q_k$ are simple and there is no pole at $0$ and no polynomial part; each step divides by one of the three numbers of the hypothesis.

\emph{Uniqueness.} Let $\sum_ic_ix^i/Q=\eth g$ with $g\in\mathfrak A$, $g\neq0$. A pole of order $h\ge1$ of $g$ at $q_k$ gives $\eth g$ a pole of order $h+1$ by (a), but the left side has simple poles; so $g$ is a Laurent polynomial. A pole $x^{-r}$ of $g$ at $0$ gives $\eth g$ the pole $(\kappa-r)x^{-r}$, but $Q(0)\neq0$; so $g$ is a polynomial, of degree $a\ge0$, and by (c) $\eth g$ is of order $x^a$ at $\infty$, while the left side is $O(x^{-1})$. So $g=0$ and every $c_i=0$: $\dim_{\mathbb F'}\mathfrak A/\eth\mathfrak A=D$.

\emph{Reads.} Let the reads of $R$ and of the $x^i/Q$ converge absolutely, so that $\Phi=x^\kappa\prod_jP_j^{-w_j}g$ has the absolutely integrable derivative $x^{\kappa-1}\prod_jP_j^{-w_j}\big(R-\sum_ic_ix^i/Q\big)$ and limits at $0$ and at $\infty$. With $\nu_0$ the order of $g$ at $0$ and $a$ its degree at $\infty$, $\Phi=c_0x^{\kappa+\nu_0}(1+o(1))$ at $0$ and $\Phi=c_\infty x^{\kappa+a-\sum_jw_j\deg P_j}(1+o(1))$ at $\infty$, with $c_0,c_\infty\neq0$ by the first fact above; a limit exists only if the exponent has negative real part at $\infty$, positive at $0$, or is $0$. At $0$: if $\nu_0\le-1$ the exponent $\kappa-|\nu_0|$ is not $0$ by hypothesis, and if $\nu_0\ge0$ its real part is positive since the read of $1/Q$ converges, $\re\kappa>0$. At $\infty$: if $a\ge0$ the exponent is not $0$ by hypothesis, and if $a\le-1$ its real part is negative since the read of $x^{D-1}/Q$ converges, $\re(\kappa-\sum_jw_j\deg P_j)<1$. So both limits are $0$, $\langle\eth g\rangle=0$, and $\langle R\rangle=\sum_ic_i\langle x^i/Q\rangle$; at symbolic exponents the identity holds on the open set of absolute convergence and continues. For the linear letters, $Q=\prod_\ell(x+\xi_\ell)$ is squarefree, the $x^i/Q$, $i<|J|$, and the $(x+\xi_\ell)^{-1}$ span the same space by partial fractions, and the three numbers are $\kappa-r$, $\kappa+a-|\mathbf w|$ and $w_\ell+h$, nonzero in $\mathbb K$.
\end{proof}

\begin{hdefinition}[\hb{def:masters} (Deformation; masters)]
As stated, along every admissible direction.
\end{hdefinition}

\begin{proof}
Along the deformation the three numbers of \hb{thm:scalereduce} are $\kappa-r+\lambda_0\varepsilon$, $\big(\kappa+a-\sum_jw_j\deg P_j\big)+\varepsilon\big(\lambda_0-\sum_j\lambda_j\deg P_j\big)$ and $(\mu_k+h)+\varepsilon\sum_jm_{jk}\lambda_j$, polynomials in $\varepsilon$ with nonzero slopes, hence nonzero in $\mathbb F'(\lambda_0,\boldsymbol\lambda)(\varepsilon)$. At large real $\varepsilon$ the exponent of a read at $0$ has real part growing like $\lambda_0\varepsilon$ and at $\infty$ falling like $(\sum_j\lambda_j\deg P_j-\lambda_0)\varepsilon$, the moduli of the powers being those of positive letters up to bounded factors; so finitely many reads converge absolutely together, the coefficients of the reductions have no pole there, and \hb{thm:scalereduce} gives the identities there, hence as meromorphic functions of $\varepsilon$, the reads continuing meromorphically by the expansions of the integrands at $0$ and at $\infty$ (\hb{lem:cube}). $Q$ divides $P_1\cdots P_N$. The conditions are strict inequalities and finitely many inequations, invariant under positive scaling, so the admissible directions form an open cone; $(\frac12;\one)$ is in it, since $\sum_j\deg P_j\ge1>\frac12$ and $\sum_jm_{jk}\ge1$ for every $k$, $q_k$ dividing some $P_j$.
\end{proof}

\begin{hlemma}[\hb{lem:simplepoles} (Simple poles at a corner)]
As stated, along every admissible direction and for letters without roots on $[0,\infty)$.
\end{hlemma}

\begin{proof}
The three reductions of the proof of \hb{thm:scalereduce} act on separate parts of an insertion: (a) on the poles of order at least two at one $q_k$, (b) on the pole at $0$, (c) on the polynomial part, and each step feeds only the lower orders of its own part and the simple poles, which form the remainder $\sum_ic_ix^i/Q$. In (c) the step at the degree $a$ divides by $\big(\kappa+a-\sum_jw_j\deg P_j\big)+\varepsilon\big(\lambda_0-\sum_j\lambda_j\deg P_j\big)$, which vanishes at $\varepsilon=0$ for at most one $a$; in (b) the step at $r$ divides by $\kappa-r+\lambda_0\varepsilon$, vanishing at $0$ for at most one $r$; in (a), for each $k$, the step at $h$ divides by $\mu_k+h+\varepsilon\sum_jm_{jk}\lambda_j$, vanishing at $0$ for at most one $h$. Along each chain every other division is by a number holomorphic and nonzero near $\varepsilon=0$. So for $R$ holomorphic in $\varepsilon$, whose principal parts are then holomorphic in $\varepsilon$, every coefficient is a sum of terms with at most one division vanishing at $\varepsilon=0$: at most a simple pole.

For a master write $H_i=\int_0^1+\int_1^\infty$. At $0$ the integrand is $x^{\kappa+\lambda_0\varepsilon-1}\phi_\varepsilon(x)$ with $\phi_\varepsilon$ analytic at $0$ and holomorphic in $\varepsilon$, and \hb{lem:cube} gives poles at $\kappa+\lambda_0\varepsilon+l=0$, $l\in\mathbb N$, all simple and at distinct points, at most one of them at $\varepsilon=0$. At $\infty$, by the first fact above, the integrand is $x^{e(\varepsilon)-1}$ times a power series in $x^{-1}$ holomorphic in $\varepsilon$, $e(\varepsilon)=\kappa-\sum_jw_j\deg P_j+i-D+\varepsilon\big(\lambda_0-\sum_j\lambda_j\deg P_j\big)$, and $x\mapsto1/x$ gives simple poles at $e(\varepsilon)=l$, $l\in\mathbb N$, at most one at $\varepsilon=0$. So $H_i$ has at most a simple pole there.
\end{proof}

\begin{htheorem}[\hb{thm:corner} (Finitely many numbers at every corner)]
As stated, along every admissible direction and for letters without roots on $[0,\infty)$.
\end{htheorem}

\begin{proof}
Absolute convergence of $\langle R\rangle_{\kappa,\mathbf w;\mathbf P}$ is a strict inequality at $0$ and one at $\infty$, so the deformed read converges absolutely for $\varepsilon$ near $0$ and is holomorphic there, with the value $\langle R\rangle_{\kappa,\mathbf w;\mathbf P}$ at $0$. By \hb{def:masters} it equals $\sum_ic_i(\varepsilon)H_i(\varepsilon)$ as a meromorphic function, and by \hb{lem:simplepoles} each product has at most a double pole, with the constant term $c_{i,-1}h_{i,1}+c_{i,0}h_{i,0}+c_{i,1}h_{i,-1}$. The constant term of a function holomorphic at $0$ is its value. The masters do not involve $R$. The derivative $\partial_\varepsilon^b$ of $x^{\lambda_0\varepsilon}\prod_jP_j^{-\varepsilon\lambda_j}$ is $\big(\lambda_0\log x-\sum_j\lambda_j\log P_j\big)^b$ times it, the logarithms continuous along $[0,\infty)$, so the read of $R\big(\lambda_0\log x-\sum_j\lambda_j\log P_j\big)^b$ is $b!$ times the coefficient of $\varepsilon^b$ of $\sum_ic_iH_i$, which uses $h_{i,-1},\dots,h_{i,b+1}$: $(b+3)D$ numbers. For the last sentence write $\ell_0=\log x$ and $\ell_j=\log P_j$. The $b$-th powers of the linear forms $\lambda_0\ell_0-\sum_j\lambda_j\ell_j$, for $(\lambda_0;\boldsymbol\lambda)$ in a dense subset of a nonempty open set, span the forms of degree $b$ in $(\ell_0,\dots,\ell_N)$: a linear functional vanishing on them is a polynomial in $(\lambda_0;\boldsymbol\lambda)$ vanishing on a dense subset of an open set, hence zero, and its coefficients are, up to multinomial factors, its values on the monomials. The rational admissible directions are such a set, and a linear system with rational coefficients that has a solution has a rational one; so $\ell_0^{k_0}\prod_j\ell_j^{k_j}$ is a rational combination of finitely many such powers, and its read, absolutely convergent with that of $R$, is the same combination of reads along these directions.
\end{proof}

\begin{htheorem}[\hb{thm:rayrec} (Integer rays have certified recurrences)]
As stated, for letters without roots on $[0,\infty)$.
\end{htheorem}

\begin{proof}
Over $\mathbb F'(n,\varepsilon)$ the three numbers of \hb{thm:scalereduce} at the data $\big(\kappa+an+\frac\varepsilon2,\mathbf w+n\mathbf b+\varepsilon\one\big)$ have the nonzero slopes $\frac12$, $\frac12-\sum_j\deg P_j$ and $\sum_jm_{jk}$ in $\varepsilon$. Reduce the $D+1$ insertions $y_k=x^{ak}\prod_jP_j^{-b_jk}$: $y_k=\eth_ng_k+\sum_ic_{ki}x^i/Q$, where the denominators of $g_k$ and of $c_{ki}$ are products of the three numbers. The $(D+1)\times D$ matrix $(c_{ki})$ has a nonzero vector in its left kernel, found by elimination; clearing denominators gives polynomials $p_k(n,\varepsilon)$, and dividing by the largest power of $\varepsilon$ that divides all of them makes the $p_k(n,0)$ not all zero. Then $g=\sum_kp_kg_k$, divided by that power of $\varepsilon$, has $\eth_ng=\sum_kp_ky_k$, and its denominators are products of the three numbers and a power of $\varepsilon$, none vanishing identically in $\varepsilon$ at an integer $n$. Fix an integer $n$ at which the $D+1$ displayed reads converge, absolutely since the reads of a ray are read where they converge absolutely; they converge absolutely for small $\varepsilon$ too. For all but finitely many small $\varepsilon\neq0$, $g$ is defined and the exponents of $x^{\kappa+an+\varepsilon/2}\prod_jP_j^{-w_j-nb_j-\varepsilon}g$ at $0$ and at $\infty$, which move with $\varepsilon$, are not $0$; the argument of \hb{thm:scalereduce} then gives $\langle\eth_ng\rangle=0$, so $\sum_kp_k(n,\varepsilon)\langle1\rangle_{\kappa+a(n+k)+\varepsilon/2,\,\mathbf w+(n+k)\mathbf b+\varepsilon\one}=0$. Letting $\varepsilon\to0$, the reads being continuous there, gives the recurrence at $n$.
\end{proof}

\begin{htheorem}[\hb{thm:scalejets} (Coefficient jets on one scale)]
As stated, for letters without roots on $[0,\infty)$, with the powers of the changed letters continued in $u$.
\end{htheorem}

\begin{proof}
$P_{ja}/P_j$ is continuous on $[0,\infty]$, $P_j$ having no roots on $[0,\infty)$ and $\deg P_{ja}\le\deg P_j$, so $M_{ja}$ is finite. On the polydisc $z_j=\sum_au_aP_{ja}/P_j$ has $|z_j|\le\sum_a|u_a|M_{ja}<1$, so $P_j(x;u)=P_j(x)(1+z_j)$ has no roots on $[0,\infty)$, and the power continued in $u$ from $u=0$, the polydisc being convex, is $P_j(x)^{-w_j}(1+z_j)^{-w_j}$ with the principal branch of the second factor. Then
\[
(1+z_j)^{-w_j}=\sum_{m\ge0}\frac{(-1)^m(w_j)_m}{m!}z_j^m,\qquad z_j^m=\sum_{|n_j|=m}\frac{m!}{\prod_an_{ja}!}\prod_a\Big(u_a\frac{P_{ja}}{P_j}\Big)^{n_{ja}},
\]
and the product over $j$ is the displayed series inside the read. It is dominated by $\prod_j\sum_m\frac{|(w_j)_m|}{m!}\rho_j^m$, $\rho_j=\sum_a|u_a|M_{ja}<1$, finite by the ratio test, times the absolutely integrable integrand of $\langle1\rangle_{\kappa,\mathbf w;\mathbf P}$; so the read of the sum is the sum of the reads, absolutely, uniformly on smaller polydiscs, and holomorphic in $u$. Each coefficient is the read of the rational insertion $\prod_{j,a}(P_{ja}/P_j)^{n_{ja}}$, bounded on $[0,\infty)$, absolutely convergent with $\langle1\rangle_{\kappa,\mathbf w;\mathbf P}$, and \hb{thm:corner} gives it. The continuation in $u$ fixes the branch: the continuous branch of $P_j(x;u)^{-w_j}$ along $[0,\infty)$ from its principal value at $0$ can differ from it by a factor $e^{2\pi i\ell w_j}$, $\ell\in\mathbb Z$.
\end{proof}

\section{Chapter~\hb{ch:inference}: the family as inference}\label{sec:ch39}

\subsection{Samples of bases}

Every statement of this section rests on one factorization: the chance of a sample of a family is the product over its draws of the chances of single bases, each of which factors over the patterns of a bunch, and the product over all the draws is $\theta_{\mathrm{draws}}\alpha^{\mvec}/\prod_jK_j(\alpha)^{s_j}\Psi_j(\alpha)^{t_j}$.

\begin{hproposition}[\hb{prop:windowsdata} (What the data see of the bases)]
As stated.
\end{hproposition}

\begin{proof}
(i) is a property of one table. For $e\in C$, $C\setminus e$ is independent by the minimality of $C$ and lies in a basis $T$, which does not contain $C$; so $T\cap C=C\setminus e$. Hence $r(C)=k-1$, every $C\setminus e$ is a pattern of $C$, and $C$, not independent, is not. For $k\ge3$ two marks $e,f$ of $C$ lie in $C\setminus g$ for a third $g\in C$, hence in a common basis.

(ii) The level of a basis draw of table $j$ lies in the window $[d_j(S),r_j(S)]$ of $S$ (Theorem~\hb{thm:window}), and the colevel of a cobasis draw is $|S|-|T'\cap S|\in[|S|-r_j(S),|S|-d_j(S)]=[d^\ast_j(S),|S|-d_j(S)]$. Some basis reaches $r_j(S)$ and every basis has positive chance, so $p_j>0$; the $s_j$ basis draws of table $j$ are independent, and none reaches $r_j(S)$ with the chance $(1-p_j)^{s_j}$. The same holds at $d_j(S)$, and for the cobasis draws with $\mathbb P^\ast_{j,\alpha}$ and $t_j$.
\end{proof}

\begin{hproposition}[\hb{prop:samplelaws} (The exact laws of a sample of bases)]
As stated.
\end{hproposition}

\begin{proof}
(i) With $p=T\cap B$ and the weights $\theta_{j,(p),T\setminus p}=\theta_{j,T}$ of the child (Definition~\hb{def:patterns}),
\[
\mathbb P_{j,\alpha}(T)=\frac{\theta_{j,T}\alpha^T}{K_j(\alpha)}=\frac{\alpha_B^pK_{j,(p)}(\alpha_R)}{K_j(\alpha)}\cdot\frac{\theta_{j,(p),T\setminus p}\,\alpha_R^{T\setminus p}}{K_{j,(p)}(\alpha_R)},
\]
and, since $\alpha^{\cE\setminus T}=\alpha_B^{B\setminus p}\alpha_R^{R\setminus(T\setminus p)}$,
\[
\mathbb P^\ast_{j,\alpha}(T)=\frac{\alpha_B^{B\setminus p}\Psi_{j,(p)}(\alpha_R)}{\Psi_j(\alpha)}\cdot\frac{\theta_{j,(p),T\setminus p}\,\alpha_R^{R\setminus(T\setminus p)}}{\Psi_{j,(p)}(\alpha_R)} ,
\]
the first factors the pattern probabilities (Theorem~\hb{thm:tensor}) and the second the basis and cobasis laws of the child at $\alpha_R$. The draws are independent, so the patterns of the basis draws of table $j$ are independent with the chances $\pi_{j,p}(\alpha)$ and their numbers $a_{j,p}$ form a multinomial vector of size $s_j$, likewise the $b_{j,p}$ of size $t_j$; and given all the patterns the room parts are independent with the laws of the children. A basis draw with pattern $p$ uses $e\in B$ exactly when $e\in p$, and a cobasis draw with pattern $p$ uses $e$, by missing it, exactly when $e\notin p$; so $m_e=\sum_j\big(\sum_{p\ni e}a_{j,p}+\sum_{p\not\ni e}b_{j,p}\big)=\operatorname{load}_e$ (Definition~\hb{def:shares}). In the room a basis draw uses the marks of $T\setminus p$, and a cobasis draw the marks of $R\setminus(T'\setminus p)$, which it misses: what the draws of the children use. The mean of a multinomial vector is its size times its chances, which at a native gauge is the native split of Definition~\hb{def:split}.

(ii) The chance of a sample is
\[
\prod_{j,i}\frac{\theta_{j,T_{j,i}}\alpha^{T_{j,i}}}{K_j(\alpha)}\prod_{j,i}\frac{\theta_{j,T'_{j,i}}\alpha^{\cE\setminus T'_{j,i}}}{\Psi_j(\alpha)}=\frac{\theta_{\mathrm{draws}}\,\alpha^{\mvec}}{\prod_jK_j(\alpha)^{s_j}\Psi_j(\alpha)^{t_j}} .
\] Expanding $\prod_jK_j^{s_j}\Psi_j^{t_j}$ draw by draw, every sample contributes the monomial $\theta_{\mathrm{draws}}\alpha^{\mvec}$ of its counts, so the samples with the counts $\mvec$ have the total chance $\Theta_F(\mvec)\alpha^{\mvec}/\prod_jK_j^{s_j}\Psi_j^{t_j}$, $\Theta_F(\mvec)$ the coefficient of $\alpha^{\mvec}$; the ratio is $\theta_{\mathrm{draws}}/\Theta_F(\mvec)$.

(iii) $\kappa_S$ is the sum over the draws of $|T_{j,i}\cap S|-d_j(S)$ and $|S\setminus T'_{j,i}|-d^\ast_j(S)$ (Proposition~\hb{prop:samplequestion}(iii)), independent, with the generating functions $\sQ^K_{j,S}(x;\alpha)=\oE_{T\sim\mathbb P_{j,\alpha}}[x^{N_S-d_j(S)}]$ and $\sQ^\Psi_{j,S}(x;\alpha)$ (Definition~\hb{def:skelpoly}); so the generating function of $\kappa_S$ is the product. Each factor has degree $c_j(S)$ and every coefficient from degree $0$ to $c_j(S)$ positive, every level of the window being reached by some basis (Theorem~\hb{thm:window}); so the values of $\kappa_S$ fill the stated range and the product at $x=0$ is positive. The mean of a level excess is $\sum_{e\in S}R_{j,e}(\alpha)-d_j(S)$ and that of a colevel excess $\sum_{e\in S}\mathbb P^\ast_{j,\alpha}(e\notin T)-d^\ast_j(S)=\sum_{e\in S}R^\ast_{j,e}(\alpha)-d^\ast_j(S)$; the cumulants of order $k\ge2$ of an excess are those of its level, $\sum_i\varkappa_k$ of the coins (Proposition~\hb{prop:sklaws}(iv)), and summed with the sizes they are $\partial_S^k\Phi(\log\alpha)$ (Proposition~\hb{prop:stillcoins}). With a positive spectrum each excess is a sum of independent coins (Proposition~\hb{prop:sklaws}(iv)), and $\kappa_S$ is the sum of all of them.
\end{proof}

\begin{hproposition}[\hb{prop:fibretests} (Tests that do not see the gauge)]
As stated.
\end{hproposition}

\begin{proof}
(i) is Proposition~\hb{prop:samplelaws}(ii); for lists $\theta_{\mathrm{draws}}=1$. (ii) $\mathbb P(A)=\sum_{\mvec}\mathbb P(\mvec)\,\mathbb P(A\mid\mvec)\le\sum_{\mvec}\mathbb P(\mvec)\,\varepsilon=\varepsilon$, the conditional chances being free of the gauge by (i). (iii) The samples of $F'$ with the counts $\mvec$ are those of $F$ whose every draw is a basis of its table in $F'$, with the same $\theta_{\mathrm{draws}}$; by (i) their conditional chance under $F$ is $\Theta_{F'}(\mvec)/\Theta_F(\mvec)$. The rejection region, the union over the $\mvec$ with $p_{F'|F}(\mvec)\le\varepsilon$ of these samples, carries at most $\varepsilon\Theta_F(\mvec)$ of the weight at every $\mvec$, and (ii) applies. A draw that is not a basis of its table in $F'$ has chance zero under $F'$ at every gauge. (iv) The true family is rejected with chance at most $\varepsilon$.
\end{proof}

\begin{hcorollary}[\hb{cor:evidence} (The evidence of a family)]
As stated.
\end{hcorollary}

\begin{proof}
The centre of strengths $\bsig_0$ is a question with every proper wall positive on a connected family. If every draw is a basis of its table of $F$, Proposition~\hb{prop:posterior}(ii) with the prior question $(\sigma_0\one;\bsig_0,\bsig_0)$ gives the chance of the sample with the tilt integrated out, which on the balance wall is the gauge integrated out up to its scale, on which the chance of the sample does not depend (Proposition~\hb{prop:samplequestion}(i)). Otherwise the sample has chance zero under $F$ at every gauge. The posterior over the collection is the prior times this chance, normalized.
\end{proof}

\begin{sremark}
At one table, (i) and (ii) of Proposition~\hb{prop:samplelaws} are the reference's statement, with $s=q$ and $t=0$, and Proposition~\hb{prop:fibretests} and Corollary~\hb{cor:evidence} at one list are the reference's. Part (iii) of Proposition~\hb{prop:samplelaws} is the family reading of Proposition~\hb{prop:windowsdata}(ii): the floors of the windows of every table are seen together through the one wall $\kappa_S$ of the sample, whose chance of vanishing is $\prod_j\sQ^K_{j,S}(0;\alpha)^{s_j}\sQ^\Psi_{j,S}(0;\alpha)^{t_j}$.
\end{sremark}

\subsection{Statistics from a centre}

\begin{hproposition}[\hb{prop:greycentre} (The even and odd statistics from a centre)]
As stated.
\end{hproposition}

\begin{proof}
(i) The law of $Z$ is that of the question $(0;\bsig,\mathbf 0)$, whose exponent is $-\sum_j\sigma_jE_j$ (Definition~\hb{def:question}). The centre is inside its walls, so the family is connected, and by Proposition~\hb{prop:evenfloor}(iii) $\sum_j\sigma_jE_j$ is strictly convex and even, least only at $0$, with $\sum_j\sigma_jE_j(0)=2\sum_j\sigma_j\log K_j(\one)$. Its Hessian at $0$ is $\sum_j\sigma_j\cdot2\operatorname{Cov}_{\mathbb P_{j,\one}}(\chi_T)=8\sum_j\sigma_j\kappa_{j,\one}(X_e,X_f)$, which is rung two of the still of the realization $\big((\sigma_j,\sigma_j)_j;\one\big)$ by Proposition~\hb{prop:stillgauge}. For one table, $E_j$ is convex and even with its least value $2\log K_j(\one)$ at $0$. If $v$ is constant on the molecules of table $j$, $\ip{\chi_T}v=c$ for every basis (Proposition~\hb{prop:directsum}), so $\Lambda_j(z+v)=e^c\Lambda_j(z)$, $\Lambda_j(-z-v)=e^{-c}\Lambda_j(-z)$, and $E_j(z+v)=E_j(z)$; in particular $E_j(v)=E_j(0)$. Conversely if $E_j(z)=E_j(0)$, the convex function $t\mapsto E_j(tz)$ has its least value at $t=0$ and at $t=1$ and so is constant on $[0,1]$; its second derivative at $0$, $2\operatorname{Var}_{\mathbb P_{j,\one}}\ip{\chi_T}z$, vanishes, and $z$ is constant on the molecules of table $j$ (Lemma~\hb{lem:circuit}). The Laplace transform is Proposition~\hb{prop:reach}(i) at $c=0$, $\mathbf a=\mathbf k$, $\mathbf b=0$: the question $(0;\bsig+\mathbf k,\mathbf 0)$ is the centre of strengths $\bsig+\mathbf k$, with the data $\big((\sigma+|\mathbf k|)\one;\bsig+\mathbf k,\bsig+\mathbf k\big)$, and its walls are $\kappa^{\mathbf E}_S(\bsig+\mathbf k)=\kappa^{\mathbf E}_S(\bsig)+\kappa^{\mathbf E}_S(\mathbf k)$. Along $\mathbf k=\lambda\bsig$ they are $(1+\lambda)\kappa^{\mathbf E}_S(\bsig)$, all zero at $\lambda=-1$; along $\lambda\mathbf k$ the first to vanish is at the least ratio stated.

(ii) The density $e^{-\sum_j\sigma_jE_j(z)}$ is even, and $z\mapsto-z$ keeps every $E_j$ and negates every $O_j$; so the two triples have one law, and a joint moment of odd total degree in $Z$ and the $O_j(Z)$, finite as the density decays exponentially (Proposition~\hb{prop:evenfloor}) while the statistics grow linearly, equals its own negative. The Laplace transform and its domain are Proposition~\hb{prop:reach}(i) at $\mathbf a=0$. When $|b_j|\le\sigma_j$ and $\ell^\ast_{\bsig}(c,\mathbf b)>1$, the question $(c;\bsig,\mathbf b)$ has nonnegative exponents and every proper wall positive, so it has a native gauge (Theorem~\hb{thm:saddle}(iii)), its target is its masses, and its period over $P_{\mathrm c}$ is its pseudoexpectation (Definition~\hb{def:threeE}(iii)). The tails are Proposition~\hb{prop:tails} at the centre; when $\ell^\ast_{\bsig}(c,\mathbf b)=\infty$ both $Y$ and $-Y$ are bounded above, and $Y$ is bounded.

(iii) By Theorem~\hb{thm:tworeadings}(i), $K_j(\alpha)=e^{\ip\one z}\Lambda_j(z)$ and $\Psi_j(\alpha)=e^{\ip\one z}\Lambda_j(-z)$, and $\log\Lambda_j(\pm z)=\frac12\big(E_j(z)\pm O_j(z)\big)$. So $K_j^{-s_j}=e^{-s_j\ip\one z}e^{-\frac{s_j}2(E_j+O_j)}$ is the factor of table $j$ at $(\sigma_j,\delta_j)=(\frac{s_j}2,\frac{s_j}2)$, its factor $e^{-s_j\ip\one z}=\alpha^{-\frac{s_j}2\one}$ the shift of the masses, and likewise $\Psi_j^{-t_j}$. Splitting each table into its two letters, at $(s_j,0)$ and $(0,t_j)$, keeps the family word, and the pair coordinates add: $(\frac{s_j}2+\frac{t_j}2,\frac{s_j}2-\frac{t_j}2)=(\sigma_j,\delta_j)$.
\end{proof}

\begin{hproposition}[\hb{prop:twostatistics} (What the two statistics are)]
As stated.
\end{hproposition}

\begin{proof}
For scores $a_T$ with weights $\theta_T>0$, $\max_T(a_T+\log\theta_T)\le\log\sum_T\theta_Te^{a_T}\le\max_Ta_T+\log\sum_T\theta_T$: a soft maximum; and $-\log\sum_T\theta_Te^{-a_T}$ is a soft minimum. (i) $E_j=\log\Lambda_j(z)+\log\Lambda_j(-z)$ and $O_j=\log\Lambda_j(z)-\log\Lambda_j(-z)$ are the soft maximum less, and plus, the soft minimum, and $\Phi^\circ=\sum_j\big(s_j\log\Lambda_j(z)+t_j\log\Lambda_j(-z)\big)$ (Theorem~\hb{thm:saddle}). (ii) By Theorem~\hb{thm:innerMGF} at $\alpha=\one$, $\Lambda_j(z)/\Lambda_j(0)=\oE_{T\sim\mathbb P_{j,\one}}[e^{Y_j}]$, so near $0$, $\log\Lambda_j(z)=\log K_j(\one)+\sum_{k\ge1}\frac1{k!}\kappa_{j,\one}(Y_j^{\times k})$ and $\log\Lambda_j(-z)$ is the same with $(-1)^k$; their sum and difference are the two expansions. $\kappa_{j,\one}(Y_j)=\sum_ez_e\oE[\chi_{T,e}]=\sum_e(2R_{j,e}(\one)-1)z_e$, and for $k\ge2$ multilinearity and $\chi_T=2X-\one$ give $\kappa_{j,\one}(Y_j^{\times k})=\sum_wz^w\kappa_{j,\one}(\chi_{T,w_1},\dots,\chi_{T,w_k})=2^k\sum_wz^w\kappa_{j,\one}(w)$, a constant shift not entering a cumulant of order $k\ge2$. The terms of order $k\ge2$ of $\Phi^\circ$ are then $\frac1{k!}\sum_w2^k\sum_j\big(s_j+(-1)^kt_j\big)\kappa_{j,\one}(w)z^w$, which at the exponents of a realization is $\frac1{k!}\sum_w\kappa_\Upsilon(w)z^w$ (Proposition~\hb{prop:stillgauge}). (iii) $|\log\Lambda_j(z)-h_j(z)|\le\max(|\log\theta_{j,\min}|,|\log K_j(\one)|)$, so $E_j(\lambda z)-\lambda\wdt_j(z)$ and $O_j(\lambda z)-\lambda(h_j(z)-h_j(-z))$ are bounded; divide by $\lambda$. On the sector of $\pi$ the largest score is at $T^+_{j,\pi}$ and the least at $T^-_{j,\pi}$ (Lemma~\hb{lem:greedy}), and $\sum_j\big(s_jh_j(z)+t_jh_j(-z)\big)=\ip{\sum_j(s_j\chi_{T^+_{j,\pi}}+t_j\chi_{\cE\setminus T^-_{j,\pi}})}z=\ip{2v_\pi-2\sigma\one}z$; for $s_j,t_j\ge0$, in the interior of the sector $\one_{T^+_{j,\pi}}$ and $\one_{\cE\setminus T^-_{j,\pi}}$ are the only points of $P_j$ and $P^\ast_j$ at which $\ip{\cdot}z$ is largest, so $v_\pi$ is the vertex of $\mathcal N_F$ at which it is largest. By Proposition~\hb{prop:evenfloor}(i), $|X+2\kappa|$ is bounded, so $X(\lambda z)/\lambda\to-2\kappa(z)$, and on $\one^\perp$, where $\ip{\one}z=0$ and $\ip{\bmu_\perp}z=\ip{\mvec}z$, $2\kappa(z)=\ip{2v_\pi-2\sigma\one}z-2\ip{\bmu_\perp}z=2\ip{v_\pi-\mvec}z$. At the tilt $-\one_{U_k}+\frac{|U_k|}n\one\in H_\pi\cap\one^\perp$ this reads $\kappa_{U_k}=m_{U_k}-(v_\pi)_{U_k}$, the term in $\one$ dropping because $(v_\pi)_{\cE}=\sum_j(s_jr_j+t_jr^\ast_j)=m_{\cE}$ on the balance wall.
\end{proof}

\begin{sremark}
At one table and the centre of strength $\sigma_0$, Proposition~\hb{prop:greycentre} is the reference's, with $\ell^\ast_{\sigma_0}(0,b)=\sigma_0\,\ell^\ast_1(0,1)/|b|$; Proposition~\hb{prop:twostatistics}(i)--(ii) for one table is the reference's statement, and the family adds the still and, in (iii), the greedy vertex of the Newton sum, whose tail sums against the masses are the walls along the order: the reading of the walls of Theorem~\hb{thm:saddle}(iii) at the vertex that the order selects.
\end{sremark}

\subsection{The Grey law, the dial and the sorted draw}\label{sec:ch39more}

\begin{hproposition}[\hb{prop:greylaw}(iii) (The tails of the Grey law)]
As stated, for every real direction of the exponents.
\end{hproposition}

\begin{proof}
Let $F_0$ be the family of the list $\binom{\cE}1$, whose bases are the singletons and whose letter is $K_0=|\alpha|$, and let $F'$ be the family of the tables of $F$ at the exponents $(0,0)$ and the list $\binom{\cE}1$ at $(m_{\cE},0)$. These data are on the balance wall of $F'$, the list $\binom{\cE}1$ having rank one; that table has least level $0$ on every proper set, so the walls are $\kappa_S=m_S>0$; and the law of the question is $\Dir_{\mvec}$ (Proposition~\hb{prop:sectiontable}). Let $v$ be the balanced direction of the data that keeps the masses, moves the exponents of table $j$ of $F$ by $(a_j,b_j)$ and those of the list $\binom{\cE}1$ by $(-d,0)$, $d=\sum_j(a_jr_j+b_jr^\ast_j)$. Along $q+\lambda v$ the walls are $m_S-\lambda g(S)$, so the $\lambda^\ast$ of Proposition~\hb{prop:tails} is $k^\ast$. The statistic $Y$ of that proposition is $X_v(Z)$, the logarithm of the ratio of the integrands of the periods at $q+v$ and at $q$ (Definition~\hb{def:question}, and the proof of Proposition~\hb{prop:evenfloor}(i)), which at the tilt $Z$ of $U$ is $\log\big(\prod_jK_j(U)^{-a_j}\Psi_j(U)^{-b_j}K_0(U)^{d}\big)=\sum_j(a_jG_j+b_jG^\ast_j)$, as $K_0(U)=|U|=1$ on the simplex. Proposition~\hb{prop:tails} gives the tail when $k^\ast<\infty$, and the bound above when $k^\ast=\infty$.
\end{proof}

\begin{sremark}
The reference reads (iii), for $a_j,b_j\ge0$, off the abscissa of convergence of the Laplace transform of (i), which gives the limit superior; the argument above gives the limit, at every real direction. In the direction $(\mathbf s,\mathbf t)$ of a word it is the tail of Proposition~\hb{prop:dial}(iv), which the reference proves from the collapse of a set attaining $k^\ast$.
\end{sremark}

\begin{hproposition}[\hb{prop:dial} (The exponent is a dial between the mean and the maximum)]
As stated, for real exponents of any signs, not all zero.
\end{hproposition}

\begin{proof}
(i) is a property of power means and does not involve the exponents. (ii) By Proposition~\hb{prop:twolayers}(ii), $\oE[e^{kA_k}]=\frac1q\sum_i\oE[W(U_i)^k]=P_F(\mvec;k\mathbf s,k\mathbf t)\big/\Gst{n}{1}{\mvec}{m_{\cE}}$ wherever every proper wall at $(\mvec;k\mathbf s,k\mathbf t)$, which is $m_S-kg(S)$, is positive, that is for $0<k<k^\ast$; and $e^{kA_k}=\frac1q\sum_iW(U_i)^k$ has a finite variance exactly when $\oE[W(U)^{2k}]<\infty$, that is $2k<k^\ast$ (Proposition~\hb{prop:twolayers}(i) at the exponents $(2k\mathbf s,2k\mathbf t)$). (iii) For $k\ge k^\ast$ some proper wall at $(\mvec;k\mathbf s,k\mathbf t)$ is not positive, and $\oE[W(U)^k]=\infty$ by Proposition~\hb{prop:twolayers}(i); the law of the maximum of independent draws is the power of the law of one. (iv) $Y=\sum_j(s_jG_j+t_jG^\ast_j)$ with $(G_j,G^\ast_j)_j\sim\mathrm{Grey}(F;\mvec)$, and Proposition~\hb{prop:greylaw}(iii) in the direction $(\mathbf s,\mathbf t)$ gives the tail and the bound. For $s_j,t_j\ge0$: $(|U_S|,|U_{\cE\setminus S}|)\sim\Dir_{(m_S,\,m_{\cE\setminus S})}$ by the aggregation of the coordinates of a Dirichlet draw; on the simplex every coordinate in $S$ is at most $\rho$ and every coordinate at most $1$, so $K_j(U)\le K_j(\one)\rho^{d_j(S)}$ and $\Psi_j(U)\le\Psi_j(\one)\rho^{d^\ast_j(S)}=K_j(\one)\rho^{d^\ast_j(S)}$, and raising these to the powers $-s_j\le0$, $-t_j\le0$ gives the bound on $W(U)$.
\end{proof}

\begin{hproposition}[\hb{prop:sortedrow} (A sorted draw lives on the sector of its order)]
As stated, for real exponents of any signs.
\end{hproposition}

\begin{proof}
The family is connected: on the balance wall $\kappa_S+\kappa_{\cE\setminus S}=2\sum_j\sigma_jc_j(S)$ (Lemma~\ref{lem:walls}), so positive walls put on every nonempty proper $S$ a table with $\sigma_j>0$, hence active, and $c_j(S)>0$. The tilt $Z$ of $U$ has the density $2^{n-1}e^{X(z)}/P_F(\mvec;\mathbf s,\mathbf t)$ with respect to $dz$ (Definition~\hb{def:question}), and on the sector $H_\pi$, where $dz$ is the Lebesgue measure in the differences $v_k$, Proposition~\hb{prop:evenfloor}(i) gives $e^{X(z)}=e^{-2\sum_k\kappa_kv_k}\prod_jk_{j,\pi}(q)^{-s_j}\psi_{j,\pi}(q)^{-t_j}$ with $q_k=e^{-2v_k}$, at any signs of the exponents. In the gaps $u_k=2v_k=\log(U_{\pi(k-1)}/U_{\pi(k)})$, $du=2^{n-1}dv$, so on $H_\pi$ the gaps have the density $P_F(\mvec;\mathbf s,\mathbf t)^{-1}e^{-\sum_k\kappa_ku_k}\prod_jk_{j,\pi}(e^{-u})^{-s_j}\psi_{j,\pi}(e^{-u})^{-t_j}$ on $[0,\infty)^{n-1}$, whose total is $I_\pi/P_F(\mvec;\mathbf s,\mathbf t)$ (Corollary~\hb{cor:pastwalls}); the boundaries of the sectors are null, and summing over $\pi$ gives $P_F(\mvec;\mathbf s,\mathbf t)=\sum_\pi I_\pi$. This is (i) and (ii). (iii) The bounds on the sector letters are those of Proposition~\hb{prop:evenfloor}(i); raised to the powers $-s_j$ and $-t_j$, of any signs, they bound $\prod_jk_{j,\pi}^{-s_j}\psi_{j,\pi}^{-t_j}$ on the cube between two positive constants, and for $s_j,t_j\ge0$ between $\prod_jK_j(\one)^{-s_j-t_j}$ and $\prod_j\theta_{j,T^+_{j,\pi}}^{-s_j}\theta_{j,T^-_{j,\pi}}^{-t_j}$; and $\int_{[0,\infty)^{n-1}}e^{-\sum_k\kappa_ku_k}du=\prod_k\kappa_k^{-1}$. (iv) is Lemma~\hb{lem:greedy}, which does not involve the exponents.
\end{proof}

\begin{sremark}
Proposition~\hb{prop:movingtarget} is the ratio of two limits of Theorem~\hb{thm:severity}(i): by Proposition~\hb{prop:twolayers}(i), $\oE_{V\sim\Dir_{\lambda\btau}}\big[\prod_jK_j(V)^{-\lambda s_j}\Psi_j(V)^{-\lambda t_j}\big]=P_F(\lambda\btau;\lambda\mathbf s,\lambda\mathbf t)\big/\Gst{n}{1}{\lambda\btau}{\lambda\tau_{\cE}}$, the numerator is the period of $\lambda$ times the question of data $(\btau;\mathbf s,\mathbf t)$, with the walls $\tau_S-g(S)$, and the denominator is the period of the list $\binom{\cE}1$ at a full basis lean at the masses $\lambda\btau$, with the walls $\lambda\tau_S$ (Proposition~\hb{prop:sectiontable}); multiplied by $\lambda^{n-1}$ they tend to $\sum_\pi\prod_k(\tau_{U_k}-g(U_k))^{-1}$ and $\sum_\pi\prod_k\tau_{U_k}^{-1}$.
\end{sremark}

\subsection{The interaction of a bunch}

The one-table proofs integrate the joint cumulant of the indicators of a bunch, one logarithm of a Taylor cube; a family integrates a sum of such logarithms with its exponents, one for each letter, so one lemma on one logarithm serves both.

\begin{slemma}[An alternating sum over the subsets of a bunch]\label{lem:cubesum}
Let $B$ be a finite set with $k=|B|\ge2$, let $(c_p)_{p\subseteq B}$ be positive, and $f(y)=\log\sum_{p\subseteq B}c_p\,e^{\ip{\one_p}y}$ on $\mathbb R^B$, and write $\partial^B=\prod_{e\in B}\partial_{y_e}$. Then $|\partial^Bf(y)|\le C\,e^{-\max_{e\in B}|y_e|}$, so $\partial^Bf$ is integrable on $\mathbb R^B$, and
\[
\int_{\mathbb R^B}\partial^Bf(y)\,dy=\sum_{p\subseteq B}(-1)^{|B\setminus p|}\log c_p .
\]
\end{slemma}

\begin{proof}
\emph{The bound.} $f(y+w)-f(y)$ is the logarithm of the moment generating function at $w$ of the indicators $(X_e)_{e\in B}$ of a random set $P\subseteq B$ with the law $\pi_y(p)\propto c_pe^{\ip{\one_p}y}$, so $\partial^Bf(y)$ is their joint cumulant. Fix $e\in B$. A joint cumulant of $k\ge2$ variables does not change when one of them is shifted by a constant, and it is a sum, with coefficients depending on $k$ alone, of products of joint moments, one for each block of a partition of the $k$ places, exactly one block holding the place of $e$. With $X_e-\oE X_e$ in that place and $|X_f|\le1$, $|\partial^Bf(y)|\le C_k\,\oE|X_e-\oE X_e|=2C_kp_e(1-p_e)$, $p_e=\pi_y(e\in P)$. Pairing each $p\not\ni e$ with $p+e$, the odds $p_e/(1-p_e)=\sum_{p\not\ni e}c_{p+e}e^{y_e}w_p\big/\sum_{p\not\ni e}c_pw_p$, $w_p=e^{\ip{\one_p}y}>0$, lie between $\rho_-e^{y_e}$ and $\rho_+e^{y_e}$, $\rho_\pm$ the largest and the least of the ratios $c_{p+e}/c_p$; so $p_e(1-p_e)\le\min(\mathrm{odds},\mathrm{odds}^{-1})\le\max(\rho_+,\rho_-^{-1})\,e^{-|y_e|}$. The best $e$ gives the bound.

\emph{The value}, by induction on $k$. Fix $e\in B$, put $B'=B-e$, $y'=y_{B'}$, and let $f^{/e}$ and $f^{\setminus e}$ be the functions of the same form on $B'$ with the coefficients $c_{p+e}$ and $c_p$, $p\subseteq B'$: the contraction and the deletion of $e$ in the cube. With $\ell(x)=\log(1+e^x)$,
\[
f(y)=f^{/e}(y')+y_e+\ell\big(f^{\setminus e}(y')-f^{/e}(y')-y_e\big)=f^{\setminus e}(y')+\ell\big(y_e+f^{/e}(y')-f^{\setminus e}(y')\big).
\]
Every derivative of $\ell$ is $O(e^x)$ as $x\to-\infty$, and every derivative of order at least one of $f^{/e}$ and $f^{\setminus e}$ is bounded; since $B'\neq\varnothing$, $\partial^{B'}y_e=0$, so $\partial^{B'}f(y)\to\partial^{B'}f^{/e}(y')$ as $y_e\to+\infty$ and $\partial^{B'}f(y)\to\partial^{B'}f^{\setminus e}(y')$ as $y_e\to-\infty$. Hence $\int_{\mathbb R}\partial^Bf\,dy_e=\partial^{B'}f^{/e}(y')-\partial^{B'}f^{\setminus e}(y')$ for every $y'$, and by the bound and Fubini
\[
\int_{\mathbb R^B}\partial^Bf\,dy=\int_{\mathbb R^{B'}}\big(\partial^{B'}f^{/e}-\partial^{B'}f^{\setminus e}\big)\,dy'.
\]
If $k\ge3$, each term is integrable by the bound, and the induction gives $\sum_{p\subseteq B'}(-1)^{|B'\setminus p|}(\log c_{p+e}-\log c_p)=\sum_{p\subseteq B}(-1)^{|B\setminus p|}\log c_p$. If $k=2$, $B'=\{e'\}$, the integrand is the derivative of $f^{/e}-f^{\setminus e}=\log\big((c_e+c_{ee'}e^{y'})/(c_\varnothing+c_{e'}e^{y'})\big)$, whose limits at $\pm\infty$ are $\log(c_{ee'}/c_{e'})$ and $\log(c_e/c_\varnothing)$; their difference is the alternating sum.
\end{proof}

\begin{htheorem}[\hb{thm:interactiontotal} (The interaction totals the still over the bunch's scales)]
As stated.
\end{htheorem}

\begin{proof}
Fix $\alpha_R$. By the Taylor cube (Theorem~\hb{thm:taylorcube})
\[
\log K_j(e^{y_B},\alpha_R)=\log\sum_{p\subseteq B}K_{j,(p)}(\alpha_R)\,e^{\ip{\one_p}{y_B}},
\]
with every $K_{j,(p)}(\alpha_R)>0$ as $\pat_j(B)=2^B$; Lemma~\ref{lem:cubesum} gives $\int\partial^B\log K_j\,dy_B=\iota_{j,B}(\alpha_R)$, absolutely. Likewise $\log\Psi_j(e^{y_B},\alpha_R)=\log\sum_{P\subseteq B}\Psi_{j,(B\setminus P)}(\alpha_R)\,e^{\ip{\one_P}{y_B}}$, and the lemma gives $\sum_{P}(-1)^{|B\setminus P|}\log\Psi_{j,(B\setminus P)}=\sum_p(-1)^{|p|}\log\Psi_{j,(p)}=\iota^\ast_{j,B}(\alpha_R)$. The parts of the potential are these logarithms with the coefficients $\sigma_j(1,1)$ and $\delta_j(1,-1)$, and a table with both exponents zero does not enter; integrating term by term gives $\iota^{\mathbf E}_B$ and $\iota^{\mathbf O}_B$. At a realization $\kappa_\Upsilon(e_1,\dots,e_k)=2^k\partial_{y_{e_1}}\cdots\partial_{y_{e_k}}\Phi$ (Definition~\hb{def:still}), and Proposition~\hb{prop:deviation} writes it through the cumulants at $\alpha$ and at $\alpha^{-1}$. For one table at $(1,0)$, $\partial^B\log K(e^y)=\kappa_\alpha(e_1,\dots,e_k)$, the joint cumulant of the indicators under $\mathbb P_\alpha$ (Definition~\hb{def:weightlaw}).
\end{proof}

\begin{hproposition}[\hb{prop:interactionroom} (The interaction belongs to the room)]
As stated.
\end{hproposition}

\begin{proof}
(i) The child at $p$ has rank $r_j-|p|$ on $R$ (Lemma~\hb{lem:minor}), so $K_{j,(p)}$ is homogeneous of degree $r_j-|p|$ and $\Psi_{j,(p)}$ of degree $|R|-r_j+|p|$. Under $\alpha_R\mapsto c\alpha_R$, $\iota_{j,B}$ moves by $\log c\sum_p(-1)^{|B\setminus p|}(r_j-|p|)$ and $\iota^\ast_{j,B}$ by $\log c\sum_p(-1)^{|p|}(|R|-r_j+|p|)$; both vanish, as $\sum_{p\subseteq B}(-1)^{|p|}=0$ for $|B|\ge1$ and $\sum_{p\subseteq B}(-1)^{|p|}|p|=0$ for $|B|\ge2$.

(ii) A basis draw of table $j$ meets $B$ in $P$ with the chance $\pi_{j,P}(\alpha)$, and a cobasis draw misses exactly $P$ on $B$ when it meets $B$ in $B\setminus P$, with the chance $\pi^\ast_{j,B\setminus P}(\alpha)=\alpha_B^{P}\Psi_{j,(B\setminus P)}(\alpha_R)/\Psi_j(\alpha)$ (Theorem~\hb{thm:tensor}); independence gives $w(P)$. So
\[
\log w(P)=\sum_j(s_j+t_j)\ip{y_B}{\one_P}+\sum_j\big(s_j\log K_{j,(P)}(\alpha_R)+t_j\log\Psi_{j,(B\setminus P)}(\alpha_R)\big)-\Phi(\log\alpha).
\]
The coefficients of the expansion are $\lambda_A=\sum_{P\subseteq A}(-1)^{|A\setminus P|}\log w(P)$, and for $|A|\ge2$ the terms constant in $P$ and linear in $\one_P$ drop out:
\[
\lambda_A=\sum_j\Big(s_j\sum_{P\subseteq A}(-1)^{|A\setminus P|}\log K_{j,(P)}+t_j\sum_{P\subseteq A}(-1)^{|A\setminus P|}\log\Psi_{j,(B\setminus P)}\Big),
\]
free of $\alpha_B$. At $A=B$ the second inner sum is $\iota^\ast_{j,B}$, and $\lambda_B=\iota^{\mathbf E}_B+\iota^{\mathbf O}_B$. For $A\subsetneq B$: the bases of table $j$ that avoid $B\setminus A$ form its deletion of $B\setminus A$, a table as some basis avoids $B$, in which every subset of $A$ is a pattern, with the children $K_{j,(P)}$, $P\subseteq A$; the bases that contain $B\setminus A$, less $B\setminus A$, form its contraction, a table as $B$ is independent, in which every subset of $A$ is a pattern and the pattern $P'\subseteq A$ has the child of the pattern $P'\cup(B\setminus A)$ of $B$. With $P'=A\setminus P$, $B\setminus P=P'\cup(B\setminus A)$, so the second inner sum is $\sum_{P'\subseteq A}(-1)^{|P'|}\log\Psi_{j,(P'\cup(B\setminus A))}$, the dual interaction of $A$ in the contraction. The deletion is the restriction to $\cE\setminus(B\setminus A)$, some basis avoiding $B\setminus A$, and these are the tables of the face of the wall at $B\setminus A$ on its room (Theorem~\hb{thm:wallres}), at the exponents $(s_j,0)$ and $(0,t_j)$: every basis of the restriction and of the contraction comes from one basis of the table, whose weight it keeps, so the weights factor on the bottom levels as that theorem asks. So $\lambda_A$ is $\iota^{\mathbf E}_A+\iota^{\mathbf O}_A$ computed in that face. At one table and $(1,0)$, $w(P)=\pi_P(\alpha)$.
\end{proof}

\begin{hproposition}[\hb{prop:interactionparity} (Both sides, and parity)]
As stated.
\end{hproposition}

\begin{proof}
The child at $p$ is a table on $R$, so $\Psi_{(p)}(\alpha_R)=\alpha_R^{R}K_{(p)}(\alpha_R^{-1})$ (Proposition~\hb{prop:identity}); in $z_R$, $\log\Psi_{(p)}=2\ip\one{z_R}+\log K_{(p)}(e^{-2z_R})$, and as $\sum_p(-1)^{|p|}=0$, $\iota^\ast_B(z_R)=\sum_p(-1)^{|p|}\log K_{(p)}(e^{-2z_R})=(-1)^{|B|}\iota_B(-z_R)$. Substituting in Definition~\hb{def:interaction} gives $\iota^{\mathbf E}_B$ and $\iota^{\mathbf O}_B$; the parities in $z_R$ and the values at $z_R=0$ are read off. The total for $-\Phi$ is $-(\iota^{\mathbf E}_B+\iota^{\mathbf O}_B)=-\sum_j(s_j\iota_{j,B}+t_j\iota^\ast_{j,B})$.
\end{proof}

\begin{sremark}
At one table and $(1,0)$ these are the reference's statements: the interaction, its homogeneity and log-linear reading, its total as the integral of the joint cumulant, and the reading of the dual table. In a family every subset of the bunch must be a pattern in every active table, and the family reads the alternating sum over the patterns of its whole potential: the term of the pattern $P$ is the family of the children $K_{j,(P)}$ at $s_j$ and $\Psi_{j,(B\setminus P)}$ at $t_j$, which is where the potential goes as the scales of $P$ grow and those of $B\setminus P$ fall. Proposition~\hb{prop:design} is one table and is unchanged.
\end{sremark}

\section{Receipts}
The identities below were checked on small families before or after they were written; the checks are receipts, not part of any proof.

\subsection{Chapters~\hb{ch:positive} to~\hb{ch:finitealgebra}}
In exact rational arithmetic at $(a,b)=(\frac73,\frac54)$, with two linear letters of rational coefficients and the exponents $\frac97$ and $-\frac23$, and at $r=3$: the expansion of $L\,x(1-x)\,\partial_xg+A_ng$ at $g=(-1)^r\Pi^{a,b}_r$ in the Beta ports had no component below the order $r-h-1$, and its coefficient of the order $r+h+1$ matched Theorem~\hb{thm:portrecurrence}. At $(n,k,r)\in\{(2,3,1),(2,3,2),(1,4,3),(2,5,1)\}$ the column sums $C_{n,k}(P_r)$, read off the Laurent expansions of the share functions, matched the sums of the proof of Theorem~\hb{thm:portbudget}, the signs $(-1)^{jk}$ included, and $C_{n,k-1}(P_r)+2\,\Th{1}{\IZ}{0}{n+r}\,C_{n,k}(P_r)$ was rational. For two letters at $\kappa=0.7$, $\mathbf w=(0.4,0.5)$, $\boldsymbol\xi=(1,1+i)$, $G=e^{-\rho}$, $C_1=1+\rho$ and $C_2=2+\rho^2$, the quantity $\big(\int_0^\infty\cdots\,d\rho-A(0)\big)\lambda^{-\kappa^\flat}$ of Proposition~\hb{prop:oneslot} was $4.10646-0.33755i$, $4.10314-0.33828i$ and $4.10276-0.33837i$ at $\lambda=10^{-3},10^{-4},10^{-5}$, against $B(0)=4.10272-0.33838i$.

\subsection{Chapters~\hb{ch:questions} and~\hb{ch:inference}}
On four marks, with the tables $\binom{\cE}2$, the rank-two table with one parallel pair, the direct sum of two rank-one tables and a rank-three table with a coloop, all with random weights: the identity of Proposition~\hb{prop:evenfloor}(i) held to $10^{-13}$ at random tilts and random questions with exponents of both signs, and the two forms of the wall at a tilt agreed to $10^{-13}$; the period computed on the sectors agreed with a direct cubature of the tilt integral to $10^{-8}$; $\lambda^{3}P_F(\lambda q)$ approached $\sum_\pi\prod_k\kappa_{U_k}^{-1}=1.87097$, and $(\sum_\pi\prod_k\kappa_{U_k}^{-1}-\lambda^3P_F(\lambda q))/\lambda$ approached $\sum_\pi\prod_k\kappa_{U_k}^{-1}\sum_j(s_j\log\theta_{j,T^+_{j,\pi}}+t_j\log\theta_{j,T^-_{j,\pi}})=2.75547$ ($2.7548$ at $\lambda=3\cdot10^{-4}$). On three marks with three tables: the relative error of Theorem~\hb{thm:severity}(ii) times $|\lambda|$ was $0.172$--$0.184$ along $\lambda\in\{2,\dots,40\}$ and along $\lambda=2+i\eta$, $\eta\in\{5,10\}$, the same constant in both directions; the usage at the computed native gauge equalled the masses; and the local slopes of $\log\mathbb P(\pm Y>x)$ for an odd statistic at the centre were between $-1.4985$ and $-1.5014$ against $-\ell^\ast_{\bsig}(c,\mathbf b)=-1.5$. In exact rational arithmetic, on four marks with sample sizes $\mathbf s=(1,0,1,0)$, $\mathbf t=(0,1,1,1)$: the conditional law of a sample given its counts was $\theta_{\mathrm{draws}}/\Theta_F(\mvec)$, and the law of the wall $\kappa_S$ was $\prod_j\sQ^K_{j,S}{}^{s_j}\sQ^\Psi_{j,S}{}^{t_j}$ for every $S$. The total of Theorem~\hb{thm:interactiontotal} matched $\iota^{\mathbf E}_B+\iota^{\mathbf O}_B$ to $10^{-12}$ for $|B|=2$ in three tables, and to $10^{-10}$ for $|B|=3$, a Hamiltonian path of $K_4$ in the graphic table of $K_4$ together with $\binom{\cE}3$ on six marks; the relation $\iota^\ast_B(z_R)=(-1)^{|B|}\iota_B(-z_R)$ held to $10^{-12}$ in each table. All of these were rerun on 20 October 2026 and reproduced (the cubature for $|B|=2$ on a grid of $1601$ points a side, matching to $6\cdot10^{-13}$).

In exact rational arithmetic, on six marks with two weighted graphic tables of $K_4$ and a weighted $\binom{\cE}3$, at the exponents $(2,1)$, $(1,3)$, $(0,2)$, a rational gauge and the bunch $B$ of a Hamiltonian path: for every $A\subseteq B$ with $|A|\ge2$ the coefficient $\lambda_A$ of Proposition~\hb{prop:interactionroom}(ii) satisfied $e^{\lambda_A}=e^{\iota^{\mathbf E}_A+\iota^{\mathbf O}_A}$ exactly, the right side computed from the deletions and contractions of $B\setminus A$ themselves.

\part{Exercises}\label{part:exercises}

Each exercise is a computation on a particular list, table or ray, set as a guided example. Its solution follows the procedures and theorems of the chapters step by step: every step names the statement it uses, and every quantity that the next step needs is displayed, in exact arithmetic. A step that is a machine computation too large to display is given by its exact input, its method, the output that the next step uses, the check that was run on it, and the receipt that reproduces it. Decimal values appear only as renderings of exact ones.

Lists are named by an upright capital with the number of marks as a subscript, and each is defined by its marks and bases where it first appears.

\chapter{Exercises on tables, families and classification}\label{ch:exclass}

The exercises of this chapter work the statements of Parts~II and~III, and the first statements of the reading of rays, on particular lists. The list on two marks is opened on its share side and read along its centre ray, where its values, its closings and the matrix of its branch fit together exactly. The pins of the atom formula of six uniform lists have the elementary divisors $1$, and two of them still pin some shares at fractions. Four Foreman trees are drawn: one that is an atom at every lean, one that peels at a bunch of width one into a cell whose contraction is the dual of its deletion, and one whose Foreman bunch at a lean has width two and whose sets do not split off at the centre. Along a set of width two the Skeleton values are shown not to be rational functions of the point, and a family with exponents of both signs is peeled with and without its record. Every number displayed is exact; the receipts that recompute them are the scripts named in each solution.

\section{The list on two marks}\label{sec:exc-two}

\begin{exercise}[The list $\binom{[2]}1$ on its share side]\label{ex:cls-share}
Let $\cE=\{1,2\}$ and $F=\big(\binom{\cE}1;s,t\big)$, the list whose bases are the two singletons, with $\re s,\re t>0$ and masses $\mvec=(m_1,m_2)$ on the balance wall with $\re m_1,\re m_2>0$.
\begin{enumerate}[label=(\roman*)]
\item Write the atom formula of $P_F(\mvec;s,t)$ (Corollary~\ref{cor:atomwall}): its shares, its pins, its fibre and its measure, and close it.
\item At $(\one;1,1)$ find the integrand on the line $\re a=\frac12$ and the period.
\end{enumerate}
\end{exercise}

\begin{solution}
Receipt: \texttt{ex01\_share\_side.py}.

(i) \emph{The letters.} The bases are $T_1=\{1\}$ and $T_2=\{2\}$, so (Definition~\ref{def:table})
\[
K(\alpha)=\alpha^{T_1}+\alpha^{T_2}=\alpha_1+\alpha_2,\qquad \Psi(\alpha)=\alpha^{\cE\setminus T_1}+\alpha^{\cE\setminus T_2}=\alpha_2+\alpha_1 ,
\]
the rank and the corank are $1$, and the balance wall is $\kappa_\cE=m_1+m_2-(s+t)=0$ (Definition~\ref{def:family}), with the strength $\sigma=\frac{s+t}2$ and the offset $\mu_1=m_1-\sigma$.

\emph{The shares and the pins.} Corollary~\ref{cor:atomwall}, with $g=2$, opens each letter with one share per basis: $a_1,a_2$ for $K$, $a_i$ the share of $\alpha^{T_i}$, and $b_1,b_2$ for $\Psi$, $b_i$ the share of $\alpha^{\cE\setminus T_i}$. Its pins are the two letter totals and the odd load at the mark $1$ (Definition~\ref{def:shares}), with $\chi_{T_1}(1)=1$ and $\chi_{T_2}(1)=-1$:
\[
a_1+a_2=s,\qquad b_1+b_2=t,\qquad \tfrac12(a_1-b_1)-\tfrac12(a_2-b_2)=m_1-\tfrac{s+t}2 .
\]
With $a=a_1$ they give
\[
a_2=s-a,\qquad b_1=a+t-m_1,\qquad b_2=m_1-a ,
\]
and the load at the mark $1$, $a_1+b_2$, is $m_1$, as $\load_1=\sigma+\langle\bdel\rangle_1$ requires (Definition~\ref{def:shares}). The weights are $1$, so the factor $\prod_T\theta_T^{-a_T-b_T}$ is $1$, and the fibre has the dimension $4-3=1$.

\emph{The measure.} The matrix of the pins, the shares in the order $(a_1,a_2,b_1,b_2)$, and its Gram matrix are
\[
P=\begin{pmatrix}1&1&0&0\\0&0&1&1\\ \frac12&-\frac12&-\frac12&\frac12\end{pmatrix},\qquad PP^{\mathsf T}=\operatorname{diag}(2,2,1),\qquad \sqrt{\det PP^{\mathsf T}}=2 ,
\]
and the fibre runs along $(1,-1,1,-1)$, of length $2$: the imaginary part $y$ of $a$ moves the point by the volume $2\,dy$. So the fibre integral of Definition~\ref{def:fibre} is $(2\pi)^{3-4}\cdot\frac12\cdot2\int dy=\frac1{2\pi i}\int da$:
\[
P_F(\mvec;s,t)=\frac1{2\pi i}\int_{(c)}\Gst21{a,\ s-a}{s}\,\Gst21{a+t-m_1,\ m_1-a}{t}\,da ,
\]
on a line $\re a=c$ on which all four shares have positive real parts, $\max(0,\re(m_1-t))<c<\min(\re s,\re m_1)$; such lines exist, since $\re s,\re t,\re m_1>0$ and $\re(m_1-t)<\re s$ is $\re m_2>0$.

\emph{The closing.} The integrand is $\Gamma(a)\,\Gamma(t-m_1+a)\,\Gamma(s-a)\,\Gamma(m_1-a)$ over $\Gamma(s)\Gamma(t)$: four Gammas on a line that separates the poles of the first two from those of the last two. Corollary~\ref{cor:fourgammas}, with its four entries $0$, $t-m_1$, $s$, $m_1$ and the variable $a$, adds them in pairs, $0+s$, $(t-m_1)+m_1$, $0+m_1$ and $(t-m_1)+s=m_2$:
\[
\frac1{2\pi i}\int_{(c)}\Gamma(a)\,\Gamma(t-m_1+a)\,\Gamma(s-a)\,\Gamma(m_1-a)\,da=\Gst41{s,\ t,\ m_1,\ m_2}{s+t}.
\]
Dividing by $\Gamma(s)\Gamma(t)$ (Proposition~\ref{prop:stacked}(i),(ii)),
\[
P_F(\mvec;s,t)=\Gst21{m_1,\ m_2}{m_1+m_2}.
\]
This agrees with the period read on the simplex (Theorem~\ref{thm:projective}): there $K=\Psi=u_1+u_2=1$, and the period is the shape integral $\int_{\Delta_\cE}u^{\mvec-\one}d\omega_\cE=\Gst21{m_1,\ m_2}{m_1+m_2}$ (Theorem~\ref{thm:dirichlet}).

(ii) At $(\one;1,1)$, $m_1=s=t=1$, and the shares are $a_1=b_1=a$ and $a_2=b_2=1-a$: the basis share and the cobasis share of each basis are equal. By Proposition~\ref{prop:stacked}(i),(v) the integrand is
\[
\Gst21{a,\ 1-a}{1}\,\Gst21{a,\ 1-a}{1}=\Gst40{a,\ a,\ 1-a,\ 1-a}{}=\Big(\frac{\pi}{\sin\pi a}\Big)^2 ,
\]
regular on $0<\re a<1$. On $a=\frac12+iy$, where $\sin\pi a=\cosh\pi y$,
\[
\frac1{2\pi i}\Gst40{a,\ a,\ 1-a,\ 1-a}{}\,da=\frac{\pi}{2\cosh^2\pi y}\,dy=\frac{2\pi}{(e^{\pi y}+e^{-\pi y})^2}\,dy=p_{\mathrm{crit}}(y)\,dy ,
\]
the density of the critical-line law (Definition~\ref{def:pairlaw}). Its primitive is $\frac12\tanh\pi y$, which runs from $-\frac12$ to $\frac12$, so its integral is $1$; and (i) gives the same, $P_F(\one;1,1)=\Gst21{1,\ 1}{2}=1$ (Theorem~\ref{thm:sharelaw}).

\emph{Answer.} The atom formula is the integral over one share
\[
P_F(\mvec;s,t)=\frac1{2\pi i}\int_{(c)}\Gst21{a,\ s-a}{s}\Gst21{a+t-m_1,\ m_1-a}{t}\,da=\Gst21{m_1,\ m_2}{m_1+m_2}.
\]
At $(\one;1,1)$ its integrand on the line $\re a=\frac12$ is $\Gst40{a,\ a,\ 1-a,\ 1-a}{}=\pi^2/\cosh^2\pi y$, the density of the critical-line law times $2\pi$, and the period is $1$.
\end{solution}

\begin{exercise}[The centre ray of $\binom{[2]}1$]\label{ex:cls-centre}
For the centre ray of the list $\binom{\cE}1$, $\cE=\{1,2\}$, at the strength $1$ (Definition~\ref{def:centreray}):
\begin{enumerate}[label=(\roman*)]
\item find its stride, its native gauge and native value, its values $\langle1\rangle_k$ and a recurrence of order one that they satisfy, and the integrand of the atom formula at step $k$ against the one at step $1$;
\item find its closings $\mathsf C_j(1)$ (Definition~\ref{def:closings}) and the volume $V_1$ of its direction, and check Proposition~\ref{prop:principal} and the reflected recurrence at every $j$;
\item find the matrix $H$ of Theorem~\ref{thm:branch} and show that $k\,\langle1\rangle_k\,\mathsf C_k(1)=2=(\det H)^{-1}$ for every $k\ge1$; compare with the leading terms of Theorems~\ref{thm:gaugesolution} and~\ref{thm:onekernel}.
\end{enumerate}
\end{exercise}

\begin{solution}
Receipt: \texttt{ex01\_share\_side.py}.

(i) \emph{Stride and native gauge.} Definition~\ref{def:centreray} with one table and $\sigma_1=1$ gives the origin $0$ and the direction $(\one;1,1)$; on the chart $\alpha_2=1$, with $x=\alpha_1$,
\[
f=\frac{\alpha^{\one}}{K\Psi}=\frac{\alpha_1\alpha_2}{(\alpha_1+\alpha_2)^2}=\frac{x}{(1+x)^2}.
\]
The gauge $\one$ is native (Proposition~\ref{prop:centrefamily}(i)), and $f\le f(\one)=\frac14$ on the orthant with equality only at $x=1$ (Theorem~\ref{thm:pencil}(i)): the native value of the direction is $\frac14$.

\emph{Values.} By Definition~\ref{def:centreray}, $\langle1\rangle_k=P_F(k\one;k,k)$, which Exercise~\ref{ex:cls-share}(i) at $m_1=m_2=s=t=k$ gives:
\[
\langle1\rangle_k=\Gst21{k,\ k}{2k}=\frac{((k-1)!)^2}{(2k-1)!}:\qquad 1,\ \tfrac16,\ \tfrac1{30},\ \tfrac1{140},\ \tfrac1{630},\ \tfrac1{2772}\quad(k=1,\dots,6).
\]
By the shift (Proposition~\ref{prop:stacked}(iii)) $\langle1\rangle_{k+1}/\langle1\rangle_k=k^2/\big(2k(2k+1)\big)=k/\big(2(2k+1)\big)$, the ratios $\frac16,\frac15,\frac3{14},\frac29,\frac5{22}$, so
\[
2(2k+1)\,\langle1\rangle_{k+1}-k\,\langle1\rangle_k=0 ,
\]
a recurrence of order one. Its symbol, read off the coefficients of $k$, is $4\lambda-1$, which vanishes at the native value $\frac14$, as Proposition~\ref{prop:nativeorder} requires; and the ratios tend to $\frac14$.

\emph{The integrand at step $k$.} At $m_1=s=t=k$ the shares of Exercise~\ref{ex:cls-share}(i) are $a_1=b_1=a$ and $a_2=b_2=k-a$, and the integrand is $\Gamma(a)^2\Gamma(k-a)^2/\Gamma(k)^2$. The shift applied $k-1$ times to $\Gamma(k-a)$ and to $\Gamma(k)$ gives
\[
\Gst42{a,\ a,\ k-a,\ k-a}{k,\ k}=\Gst40{a,\ a,\ 1-a,\ 1-a}{}\,\bigg(\frac{(1-a)(2-a)\cdots(k-1-a)}{(k-1)!}\bigg)^2 :
\]
at $k=2,3,4$ the factor is $(a-1)^2$, $\frac14(a-1)^2(a-2)^2$ and $\frac1{36}(a-1)^2(a-2)^2(a-3)^2$. The step multiplies the kernel of step $1$ by the square of a polynomial in the share, which adds zeros at $a=1,\dots,k-1$ and no poles; the values are rational numbers.

(ii) \emph{Closings.} The origin is integral and both letters have the exponent $1$ in the direction, so Definition~\ref{def:closings} applies, with $\omega_0/dy=1$ and $\mathsf C_j(1)=\CT(f^{-j})$ (Corollary~\ref{cor:reflected}). On the chart $f^{-j}=x^{-j}(1+x)^{2j}$, whose constant term is the coefficient of $x^j$ in $(1+x)^{2j}$:
\[
\mathsf C_j(1)=\binom{2j}j:\qquad 1,\ 2,\ 6,\ 20,\ 70,\ 252,\ 924,\ 3432\quad(j=0,\dots,7).
\]
\emph{The volume.} The two orders of $\cE$ have the tails $U_2=\{2\}$ and $U_2=\{1\}$, and the direction walls are $\kappa_{1,\{e\}}=\sigma_1c(\{e\})=1$ (Proposition~\ref{prop:centrefamily}(ii); $c(\{e\})=r(\{e\})+r(\cE\setminus\{e\})-r=1$ by Definition~\ref{def:skeleton} and Theorem~\ref{thm:window}). So $V_1=1+1=2$ (Proposition~\ref{prop:principal}).

\emph{The principal parts.} $\langle1\rangle_k=\Gamma(k)^2/\Gamma(2k)$ continues meromorphically in $k$. At $k=-j$, $j\in\mathbb N$, the residues of Proposition~\ref{prop:stacked}(iv) give $\Gamma(k)=\frac{(-1)^j}{j!}(k+j)^{-1}+O(1)$ and, the pole of $\Gamma$ at $2k=-2j$ having the residue $\frac1{(2j)!}$, $\Gamma(2k)=\frac1{2(2j)!}(k+j)^{-1}+O(1)$. So the pole has order $1=n-1$ and
\[
\lim_{k\to-j}(k+j)\,\langle1\rangle_k=\frac{(j!)^{-2}}{(2(2j)!)^{-1}}=2\binom{2j}j=V_1\,\mathsf C_j(1):\qquad 2,\ 4,\ 12,\ 40,\ 140,\ 504\quad(j=0,\dots,5),
\]
Proposition~\ref{prop:principal} at every $j$.

\emph{The reflected recurrence.} Corollary~\ref{cor:valuemodule} reads the recurrence of (i), $p_0(k)=-k$ and $p_1(k)=2(2k+1)$, at $k=-j$ on the closings: $j\,\mathsf C_j(1)+2(1-2j)\,\mathsf C_{j-1}(1)=0$, that is $\mathsf C_j(1)=\frac{2(2j-1)}j\,\mathsf C_{j-1}(1)$ for $j\ge1$, which $\binom{2j}j$ satisfies; at $j=0$ the sum is $p_0(0)\,\mathsf C_0(1)=0$, as in Corollary~\ref{cor:reflected}.

(iii) \emph{The matrix $H$.} In Theorem~\ref{thm:branch}, $\alpha^\star=\one$, $e_0=2$, and $s_{11}=t_{11}=1$, so $H$ is the $1\times1$ matrix $\kappa_{\one}(X_1,X_1)+\kappa_{\one}(X_1,X_1)$, the cobasis law at $(\alpha^\star)^{-1}=\one$ being the weight law (Theorem~\ref{thm:innerMGF}). Under the weight law, uniform on $\{1\},\{2\}$ (Definition~\ref{def:weightlaw}), $X_1$ is $1$ with the probability $\frac12$ and has the variance $\frac14$; so $H=\frac12$. Directly: $-\log f=2\log(1+e^y)-y$ on the chart $x=e^y$, with the second derivative $2e^y/(1+e^y)^2=\frac12$ at $y=0$.

\emph{The product.} For every $k\ge1$,
\[
k\,\langle1\rangle_k\,\mathsf C_k(1)=k\,\frac{((k-1)!)^2}{(2k-1)!}\cdot\frac{(2k)!}{(k!)^2}=k\cdot\frac{2k}{k^2}=2=(\det H)^{-1}.
\]
\emph{The leading terms.} The native gauge $\one$ is nondegenerate (Theorem~\ref{thm:onekernel}), and near it $\log f=\log\frac14-\frac12v^2$ with $v=H^{1/2}y+O(y^2)$, so $(\omega_0/dv)(\one)=(dy/dv)(0)=H^{-1/2}=\sqrt2$. Theorem~\ref{thm:gaugesolution}(i) gives $\langle1\rangle_k\sim4^{-k}k^{-1/2}a_0(1)$ with $a_0(1)=(2\pi)^{1/2}\sqrt2=2\sqrt\pi$, and Theorem~\ref{thm:onekernel}, in which $\hat A_\star(1)$ is that expansion divided by $f(\one)^kk^{-1/2}=4^{-k}k^{-1/2}$, with the leading term $a_0(1)$, gives $\mathsf C_j(1)\sim4^jj^{-1/2}(2\pi)^{-1}2\sqrt\pi=4^j(\pi j)^{-1/2}$. Their product at $j=k$ is $k^{-1}\,a_0(1)^2/(2\pi)=2k^{-1}$ to the leading order, $(\det H)^{-1}k^{-(n-1)}$; here the product is $2k^{-1}$ exactly, at every $k$.

\emph{Answer.} The stride is $\alpha_1\alpha_2/(\alpha_1+\alpha_2)^2$, with the native gauge $\one$ and the native value $\frac14$; the values are $\langle1\rangle_k=\Gst21{k,\ k}{2k}$, with $2(2k+1)\langle1\rangle_{k+1}=k\langle1\rangle_k$, and the integrand at step $k$ is the one at step $1$ times $\big((1-a)_{k-1}/(k-1)!\big)^2$. The closings are $\mathsf C_j(1)=\binom{2j}j$, $V_1=2$, and $\lim_{k\to-j}(k+j)\langle1\rangle_k=2\binom{2j}j$ at every $j$. $H=\frac12$, and $k\,\langle1\rangle_k\,\mathsf C_k(1)=2=(\det H)^{-1}$ at every $k\ge1$.
\end{solution}

\section{The pins of the atom formula}\label{sec:exc-pins}

\begin{exercise}[The pins of the atom formula of six lists]\label{ex:cls-pins}
For a list $(\cE,\cB)$ on $n$ marks at data on the balance wall, the atom formula (Corollary~\ref{cor:atomwall}) is a fibre integral over the shares $a_T,b_T$, $T\in\cB$, pinned by $|\mathbf a|=s$, $|\mathbf b|=t$ and $\langle\bdel\rangle_e=\mu_e$ for $e\neq g$. For the six lists $\binom{[2]}1$, $\binom{[3]}1$, $\binom{[3]}2$, $\binom{[4]}2$, $\binom{[4]}3$ and $\binom{[5]}2$, with $g=n$:
\begin{enumerate}[label=(\roman*)]
\item write the pins with integer coefficients;
\item find the elementary divisors of their matrix;
\item decide for which of the six lists, at integer data, every point of the fibre at which the shares outside $n+1$ shares with independent columns are integers has all its shares integers, and for the others give such a point with a share that is not an integer.
\end{enumerate}
\end{exercise}

\begin{solution}
Receipt: \texttt{ex02\_pins.py}.

(i) By Definition~\ref{def:shares} the load at a mark is $\load_e=\sum_{T\ni e}a_T+\sum_{T\not\ni e}b_T=\sigma+\langle\bdel\rangle_e$, with $\sigma=\frac12(|\mathbf a|+|\mathbf b|)$. So, with the two letter totals, the pins $\langle\bdel\rangle_e=\mu_e$ are the pins $\load_e=m_e$, $e\neq g$, which have the coefficients $0$ and $1$. Passing from the one set of rows to the other adds to the row of $e$ half the sum of the two rows of the totals: a triangular change with determinant $1$, which leaves the delta function of Definition~\ref{def:fibre}, and so the fibre integral, unchanged. The matrix of the pins, $M$, has $n+1$ rows, the two totals and the loads at $e\neq g$, and $2|\cB|$ columns:
\[
\text{column of }a_T:\ \big(1,\ 0,\ (\one_T(e))_{e\neq g}\big),\qquad \text{column of }b_T:\ \big(0,\ 1,\ (\one_{\cE\setminus T}(e))_{e\neq g}\big).
\]

(ii) An integer matrix of rank $\rho$ has integer matrices $U$, $V$ of determinant $\pm1$ with $UMV$ diagonal, with the entries $d_1\mid d_2\mid\cdots\mid d_\rho$ and zeros, and $d_1\cdots d_k$ is the greatest common divisor of the $k\times k$ minors of $M$ (the Smith normal form of an integer matrix, which applies here because the pins have integer coefficients). So the elementary divisors are all $1$ as soon as one maximal minor is $\pm1$, which also gives the rank $n+1$. With the rows in the order (total of $K$, total of $\Psi$, loads at $1,\dots,n-1$), these minors are $\pm1$:
\[
\begin{array}{ccc}
\binom{[2]}1 & \binom{[3]}1 & \binom{[3]}2\\[2pt]
(a_1,a_2,b_1) & (a_1,a_2,a_3,b_1) & (a_{12},a_{13},a_{23},b_{12})\\[4pt]
\begin{pmatrix}1&1&0\\0&0&1\\1&0&0\end{pmatrix} & \begin{pmatrix}1&1&1&0\\0&0&0&1\\1&0&0&0\\0&1&0&1\end{pmatrix} & \begin{pmatrix}1&1&1&0\\0&0&0&1\\1&1&0&0\\1&0&1&0\end{pmatrix}\\[4pt]
\det=1 & \det=1 & \det=-1\\[14pt]
\binom{[4]}3 & \binom{[4]}2 & \binom{[5]}2\\[2pt]
(a_{123},a_{124},a_{134},a_{234},b_{123}) & (a_{12},a_{13},a_{14},a_{23},b_{12}) & (a_{12},a_{13},a_{14},a_{15},a_{23},b_{12})\\[4pt]
\begin{pmatrix}1&1&1&1&0\\0&0&0&0&1\\1&1&1&0&0\\1&1&0&1&0\\1&0&1&1&0\end{pmatrix} & \begin{pmatrix}1&1&1&1&0\\0&0&0&0&1\\1&1&1&0&0\\1&0&0&1&0\\0&1&0&1&1\end{pmatrix} & \begin{pmatrix}1&1&1&1&1&0\\0&0&0&0&0&1\\1&1&1&1&0&0\\1&0&0&0&1&0\\0&1&0&0&1&1\\0&0&1&0&0&1\end{pmatrix}\\[4pt]
\det=-1 & \det=1 & \det=-1
\end{array}
\]
So every elementary divisor is $1$ for the six lists; the receipt computes the Smith normal forms, $\operatorname{diag}(1,\dots,1)$ in each case. Then $UMV=(I\mid0)$: for every integer right side $(s,t,(m_e)_{e\neq g})$ the pins have integer solutions, and the integer points of the fibre form a translate of the integer vectors of the kernel of $M$, a lattice of rank $2|\cB|-n-1$, the dimension of the fibre.

(iii) Fix a set $J$ of $n+1$ shares whose columns have the minor $\det M_J\neq0$, and give the other shares integer values. The pins then fix the shares of $J$ by Cramer's rule, with the denominator $\det M_J$. The absolute values of all maximal minors, with their numbers, are
\[
\begin{array}{l|cccc}
 & 0 & 1 & 2 & 3\\\hline
\binom{[2]}1 & 0 & 4 & 0 & 0\\
\binom{[3]}1,\ \binom{[3]}2 & 3 & 12 & 0 & 0\\
\binom{[4]}3 & 24 & 32 & 0 & 0\\
\binom{[4]}2 & 288 & 480 & 24 & 0\\
\binom{[5]}2 & 17940 & 19170 & 1620 & 30
\end{array}
\]
so for $\binom{[2]}1$, $\binom{[3]}1$, $\binom{[3]}2$ and $\binom{[4]}3$ every such point at integer data has integer shares. For $\binom{[4]}2$ the minor on $(a_{12},a_{34},b_{13},b_{14},b_{24})$ is
\[
\det\begin{pmatrix}1&1&0&0&0\\0&0&1&1&1\\1&0&0&0&1\\1&0&1&1&0\\0&1&0&1&1\end{pmatrix}=2 ,
\]
and at the centre $(\one;1,1)$ with the other seven shares $0$ the pins $a_{12}+a_{34}=1$, $b_{13}+b_{14}+b_{24}=1$, $a_{12}+b_{24}=1$ (the load at $1$), $a_{12}+b_{13}+b_{14}=1$ (at $2$) and $a_{34}+b_{14}+b_{24}=1$ (at $3$) give
\[
a_{12}=a_{34}=b_{13}=b_{24}=\tfrac12,\qquad b_{14}=0 .
\]
For $\binom{[5]}2$ the minor on $(a_{12},a_{13},a_{45},b_{14},b_{15},b_{23})$ is
\[
\det\begin{pmatrix}1&1&1&0&0&0\\0&0&0&1&1&1\\1&1&0&0&0&1\\1&0&0&1&1&0\\0&1&0&1&1&0\\0&0&1&0&1&1\end{pmatrix}=-3 ,
\]
and at the centre $(\one;1,1)$ with the other fourteen shares $0$ the pins give $a_{12}=a_{13}=a_{45}=b_{14}=b_{15}=b_{23}=\frac13$: every total and every load is $3\cdot\frac13=1$.

\emph{Answer.} The pins are the two letter totals and the loads $\load_e=m_e$, $e\neq g$, a $0/1$ matrix with $n+1$ rows and $2|\cB|$ columns. For all six lists its elementary divisors are $1$: at integer data the integer points of the fibre form a lattice of full rank in it. For $\binom{[2]}1$, $\binom{[3]}1$, $\binom{[3]}2$ and $\binom{[4]}3$ every maximal minor is $0$ or $\pm1$, and every point with integer shares outside $n+1$ independent ones has integer shares. For $\binom{[4]}2$ and $\binom{[5]}2$ some maximal minors are $\pm2$, and $\pm3$ for $\binom{[5]}2$, and at the integer data $(\one;1,1)$ the points displayed have the seven, respectively fourteen, shares outside such a minor at the integer $0$ and shares $\frac12$, respectively $\frac13$, on it. So the pins of these two atom formulas are not unimodular in the sense of the paragraph after Theorem~\ref{thm:onetree}: a closing (Theorem~\ref{thm:closing}) that pins the shares outside such a minor at poles of their Gammas can reach shares that are not integers, even at integer data, although the elementary divisors are $1$.
\end{solution}

\section{Foreman trees}\label{sec:exc-trees}

\begin{exercise}[The Foreman trees of $\binom{[4]}1$]\label{ex:cls-u14tree}
Let $F=\big(\binom{[4]}1;s,t\big)$, the list on $\cE=\{1,2,3,4\}$ whose bases are its singletons, read at its realizations $\Upsilon_{\sigma,\delta;\one}$ at the table's gauge, $s=\sigma+\delta$, $t=\sigma-\delta$ (Definition~\ref{def:realization}). Find its classes at the centre and at every lean, and its three Foreman trees (Theorem~\ref{thm:threetrees}).
\end{exercise}

\begin{solution}
Receipt: \texttt{ex03\_tree\_u14.py}.

\emph{The letters and the molecule.} $K=\alpha_1+\alpha_2+\alpha_3+\alpha_4$ and $\Psi=\alpha_1\alpha_2\alpha_3+\alpha_1\alpha_2\alpha_4+\alpha_1\alpha_3\alpha_4+\alpha_2\alpha_3\alpha_4$ (Definition~\ref{def:table}). Every nonempty $S$ has $r(S)=1$, so every proper nonempty $S$ has $c(S)=r(S)+r(\cE\setminus S)-1=1>0$ (Definition~\ref{def:skeleton}, Theorem~\ref{thm:window}): the family is connected (Definition~\ref{def:family}) and one molecule.

\emph{The masses.} Each mark lies in one of the four bases, so the occupation at the table's gauge is $R_e(\one)=\frac14$ (Theorem~\ref{thm:innerMGF}), and Proposition~\ref{prop:massesgauge} gives the realization $\Upsilon_{\sigma,\delta;\one}$ the masses
\[
\tau_e=\sigma+\delta\big(2\cdot\tfrac14-1\big)=\sigma-\tfrac\delta2\qquad\text{at every mark},
\]
with the total $4\sigma-2\delta=s\cdot1+t\cdot3$, on the balance wall.

\emph{The writings.} By Proposition~\ref{prop:viewform}, for a mark $e$ and an order $\mathbf g=(g_1,g_2,g_3)$ of its room, the pattern of the basis $\{h\}$ is the position of $h$, and at the table's gauge each pattern of one position has the weight $1$:
\[
K_{e,\mathbf g}(x)=x_0+x_1+x_2+x_3,\qquad \Psi_{e,\mathbf g}(x)=x_1x_2x_3+x_0x_2x_3+x_0x_1x_3+x_0x_1x_2 ,
\]
the same two polynomials for every mark and every order (the receipt finds one pair among the $4\cdot6$ writings). So one pair of functions, $\Phi^{\mathbf E}_{e,\mathbf g}=\sigma\log(K_{e,\mathbf g}\Psi_{e,\mathbf g})$ and $\Phi^{\mathbf O}_{e,\mathbf g}=\delta\log(K_{e,\mathbf g}/\Psi_{e,\mathbf g})$ with $x_i=e^{v_i}$, is written from every mark. With no partner, the reading of $e$ (Definition~\ref{def:view} at $k=0$) is the pair of $K=e^{v}+3$ and $\Psi=3e^{v}+1$ at every mark.

\emph{The classes.} By Proposition~\ref{prop:onewriting} the reading is constant on $\cE$ at every lean $\delta$, the centre included: $\operatorname{Cl}(\Upsilon_{\sigma,\delta;\one})=\{\cE\}$. Corollary~\ref{cor:certifyone} says the same: all $24$ permutations of $\cE$ carry the table, mapping singletons to singletons, and they take every mark to every other; none exchanges it, a basis having one mark and a cobasis three.

\emph{The trees.} Every realization is atomic (Definition~\ref{def:view}), so by Definition~\ref{def:tree} the molecule is an atom and its Foreman tree is that one leaf, with no bunch and no peel. Theorem~\ref{thm:threetrees} gives the list at its one pair the three trees $\operatorname{Tree}(F;\sigma,0)$ and $\operatorname{Tree}(F;\sigma,\pm\delta)$, and each is this atom. The molecule answers yes and the label is $(1,1)$ (Definition~\ref{def:label}).

\emph{Answer.} At the centre and at every lean the one class is $\cE$; the three Foreman trees of $\binom{[4]}1$ are one atom, with no bunch.
\end{solution}

\begin{exercise}[A tree that peels at a bunch of width one]\label{ex:cls-five}
Let $\mathrm P_5$ be the list on $\{1,\dots,5\}$ whose bases are the nine pairs other than $\{1,2\}$, read at its realizations $\Upsilon_{\sigma,\delta;\one}$. Find its classes and its Foreman bunch at the centre and at a lean, the width of the bunch, its patterns with their probabilities and children, its record and its peel, and its three Foreman trees. Write $\mathrm P_5$ as a substitution through the two cells of Theorem~\ref{thm:widthone}, and show that the contraction of the cell of the room is the dual of its deletion.
\end{exercise}

\begin{solution}
Receipt: \texttt{ex04\_five\_marks.py}.

\emph{The letters.} With $e_k$ the elementary symmetric polynomials of $\alpha$, the rank is $2$, the corank $3$, and
\[
K=e_2(\alpha)-\alpha_1\alpha_2,\qquad \Psi=e_3(\alpha)-\alpha_3\alpha_4\alpha_5 .
\]

\emph{The classes.} A bijection of the marks carries $\mathrm P_5$ exactly when it maps the one pair that is not a basis, $\{1,2\}$, to itself: these are the $12$ products of a permutation of $\{1,2\}$ and one of $\{3,4,5\}$, with the orbits $\{1,2\}$ and $\{3,4,5\}$. No bijection exchanges $\mathrm P_5$, a basis having two marks and a cobasis three. By Corollary~\ref{cor:certifyone} the classes at the centre and at every lean are
\[
\operatorname{Cl}=\big\{\{1,2\},\ \{3,4,5\}\big\}.
\]
The occupations confirm the split at the centre: $R_1(\one)=R_2(\one)=\frac39=\frac13$ and $R_3(\one)=R_4(\one)=R_5(\one)=\frac49$, and two marks of one class at the centre have equal or complementary occupations (Proposition~\ref{prop:massclass}(ii)), while $\frac49\notin\{\frac13,\frac23\}$.

\emph{The Foreman bunch.} By Definition~\ref{def:bunch} the Foreman bunch at the centre is the class with the fewest marks, $B=\{1,2\}$; at a lean the key $(|C|,|C^\circ|,C^\circ,C)$ is least at $C=\{1,2\}$, of two marks, so $B(\mathrm P_5;\sigma,\pm\delta)=\{1,2\}$ as well. Its window is $\{0,1\}$ (a basis meets $B$ in at most one mark, and $\{3,4\}$ avoids it), so $d(B)=0$, $r(B)=1$ and
\[
c(B)=r(B)+r(R)-r=1+2-2=1
\]
(Definition~\ref{def:skeleton}, Theorem~\ref{thm:window}), $R=\{3,4,5\}$; no basis contains both of its marks, so $B$ is a parallel bunch (Definition~\ref{def:width}).

\emph{Patterns, probabilities, children.} The patterns of $B$ are $\varnothing$, $\{1\}$, $\{2\}$ (Definition~\ref{def:patterns}), each the trace of three of the nine bases, so each has the probability $\pi_p(\one)=\frac13$ at the table's gauge (Definition~\ref{def:split}). The children on $R$ are
\[
\cB_{(\varnothing)}=\{34,35,45\}=\binom R2,\qquad \cB_{(\{1\})}=\cB_{(\{2\})}=\{3,4,5\}=\binom R1 ,
\]
with $K_{\binom R2}=e_2(\alpha_R)$, $\Psi_{\binom R2}=e_1(\alpha_R)$, $K_{\binom R1}=e_1(\alpha_R)$ and $\Psi_{\binom R1}=e_2(\alpha_R)$.

\emph{Record and peel.} The record (Definition~\ref{def:record}) is the three indexed children $\big(p;\cB_{(p)};\frac s3,\frac t3\big)$, $p=\varnothing,\{1\},\{2\}$. The peel (Definition~\ref{def:peel}) merges the two equal children:
\[
\Big(\binom R2;\tfrac s3,\tfrac t3\Big),\quad\Big(\binom R1;\tfrac{2s}3,\tfrac{2t}3\Big),\qquad\text{word}\quad e_1(\alpha_R)^{-\frac{2s+t}3}\,e_2(\alpha_R)^{-\frac{s+2t}3} .
\]
At the centre, $s=t=\sigma$, the word is $\big(e_1(\alpha_R)e_2(\alpha_R)\big)^{-\sigma}$, the word of the list $\binom R1$ at its own centre, and the gauge $\one_R$ is native for the masses $\sigma\one_R$ (Proposition~\ref{prop:herecentre}); the balance is $\frac\sigma3(2+1)+\frac{2\sigma}3(1+2)=3\sigma$.

\emph{The trees.} The peel is one molecule, the circuit $R$ of $\binom R2$ linking its marks (Definition~\ref{def:molecule}), and every permutation of $R$ carries both of its tables, so it reads every mark alike (Proposition~\ref{prop:certify}) and is an atom at every lean (Definition~\ref{def:view}). So the three Foreman trees (Theorem~\ref{thm:threetrees}) have one shape: the root bunch $\{1,2\}$ of width one and the room $\{3,4,5\}$, an atom, at the exponents of the peel.

\emph{The cells.} The weights are $1$ and factor across $B$ and $R$. The minors (Definition~\ref{def:minors}) are
\[
\mathrm P_5|B=\binom B1,\qquad \mathrm P_5/R=\{\varnothing\},\qquad \mathrm P_5|R=\binom R2,\qquad \mathrm P_5/B=\binom R1 ,
\]
$\mathrm P_5/R$ taking the bases inside $R$ and $\mathrm P_5/B$ those meeting $B$. The cell of the bunch (Theorem~\ref{thm:widthone}) has the bases of $\mathrm P_5|B$ and $\varnothing+g$: it is $\binom{\{1,2,g\}}1$. The cell of the room has the bases of $\binom R2$ and $\{x\}+g$, $x\in R$: it is $\mathcal A=\binom{\{3,4,5,g\}}2$. The hinge values and their ratio (Definition~\ref{def:readings}) are
\[
\gamma^K_B=\frac{K_{\mathrm P_5|B}}{K_{\mathrm P_5/R}}=\alpha_1+\alpha_2,\qquad \gamma^\Psi_B=\frac{\Psi_{\mathrm P_5/R}}{\Psi_{\mathrm P_5|B}}=\frac{\alpha_1\alpha_2}{\alpha_1+\alpha_2},\qquad \varrho_B=\frac{\alpha_1\alpha_2}{(\alpha_1+\alpha_2)^2}=u(1-u)
\]
on the shape $(u,1-u)$, as for every parallel pair (Proposition~\ref{prop:classreadings}).

\emph{The substitution.} By Proposition~\ref{prop:substitution}(iii) $\mathrm P_5$ is the gluing of the two cells along $g$: in Definition~\ref{def:substitution},
\[
\mathrm P_5=\binom{\{1,2,h\}}1\big\langle1,2,\mathcal A\big\rangle,\qquad \mathcal A=\binom{\{3,4,5,g\}}2\ \text{at }h,
\]
the marks $1,2$ left in place. Checked on the bases: the core bases $\{1\}$ and $\{2\}$ avoid $h$ and take $X\in\mathcal A/g=\binom R1$, giving the six pairs $\{1,x\}$, $\{2,x\}$; the core basis $\{h\}$ takes $X\in\mathcal A\setminus g=\binom R2$, giving the three pairs inside $R$. These are the nine bases.

\emph{The dual.} $\mathcal A\setminus g=\binom R2$ and $\mathcal A/g=\binom R1$, whose bases are the complements in $R$ of the pairs: $\mathcal A/g$ is the dual table of $\mathcal A\setminus g$, with the weights $1$. So $\varrho_{\mathcal A}\equiv1$ (Corollary~\ref{cor:piecesmarks}(ii)): $\gamma^K_{\mathcal A}=e_2(\alpha_R)/e_1(\alpha_R)=\gamma^\Psi_{\mathcal A}$.

\emph{Answer.} The classes are $\{1,2\}$ and $\{3,4,5\}$ at the centre and at every lean, and the Foreman bunch is $\{1,2\}$, a parallel bunch of width one. Its patterns $\varnothing,\{1\},\{2\}$ have the probability $\frac13$ each and the children $\binom R2,\binom R1,\binom R1$; the peel is $(\binom R2;\frac s3,\frac t3)$ with $(\binom R1;\frac{2s}3,\frac{2t}3)$, at the centre the word $(e_1e_2)^{-\sigma}$ on $R$, an atom. The three Foreman trees are the bunch $\{1,2\}$ over that atom. And $\mathrm P_5=\binom{\{1,2,h\}}1\langle1,2,\binom{\{3,4,5,g\}}2\rangle$, the cell of the room having its contraction $\binom R1$ dual to its deletion $\binom R2$, with the ratio $1$.
\end{solution}

\begin{exercise}[A Foreman bunch of width two]\label{ex:cls-sixtree}
Let $\mathrm T_6$ be the list of rank three on $\{1,\dots,6\}$ whose bases are the nineteen sets of three marks other than $\{1,2,3\}$. Show that it is prime (Theorem~\ref{thm:primetree}); find its classes at the centre and at a lean, its Foreman bunches at a basis lean and at a cobasis lean, and their widths; and draw its three Foreman trees.
\end{exercise}

\begin{solution}
Receipt: \texttt{ex05\_six\_marks.py}.

\emph{Ranks and widths.} Every set of at most three marks other than $\{1,2,3\}$ is independent, $r(\{1,2,3\})=2$, and every set of at least four marks contains a basis and has rank $3$. With $c(S)=r(S)+r(\cE\setminus S)-3$ (Definition~\ref{def:skeleton}, Theorem~\ref{thm:window}):
\[
\begin{array}{l|c}
S & c(S)\\\hline
\text{one mark} & 1+3-3=1\\
\text{two marks} & 2+3-3=2\\
\text{three marks},\ S\neq\{1,2,3\},\{4,5,6\} & 3+3-3=3\\
\{1,2,3\}\ \text{or}\ \{4,5,6\} & 2+3-3=2\\
\text{four marks} & 3+2-3=2
\end{array}
\]
Every proper nonempty set has positive width, so the list is connected (Definition~\ref{def:family}); the $50$ sets with at least two marks on each side have the widths $2$ ($32$ of them) and $3$ ($18$), so none has width one: $\mathrm T_6$ is prime (Theorem~\ref{thm:primetree}).

\emph{Classes.} A bijection carries $\mathrm T_6$ exactly when it maps the one triple that is not a basis, $\{1,2,3\}$, to itself: the $36$ products of a permutation of $\{1,2,3\}$ and one of $\{4,5,6\}$, with the orbits $C=\{1,2,3\}$ and $D=\{4,5,6\}$. The cobases are the triples other than $\{4,5,6\}$, so a bijection exchanges $\mathrm T_6$ exactly when it maps $\{1,2,3\}$ onto $\{4,5,6\}$, as $i\mapsto i+3$ modulo $6$ does: $36$ more bijections, and the $72$ together take every mark to every other. By Corollary~\ref{cor:certifyone}, at the centre the one class is $\cE$, an atom, and at every lean the classes are $C$ and $D$, with the occupations $R_e(\one)=\frac9{19}$ on $C$ and $\frac{10}{19}$ on $D$ (Proposition~\ref{prop:massclass}(ii)).

\emph{The Foreman bunch at a lean.} Both classes have three marks and lie in the one class $\cE$ of the centre, so the key of Definition~\ref{def:bunch} is decided by the order $\prec$ at the lean. $C$ and $D$ make up the molecule, so their levels are mirrored (Proposition~\ref{prop:mirror}, with $k_j=r_j=3$). Under the weight law, uniform on the $19$ bases (Definition~\ref{def:weightlaw}), the level $N_C=|T\cap C|$ is $0$ at $\{4,5,6\}$, $1$ at the $9$ bases with one mark of $C$, and $2$ at the $9$ with two:
\[
\mathbb P(N_C=0,1,2)=\tfrac1{19},\ \tfrac9{19},\ \tfrac9{19};\qquad \oE N_C=\tfrac{27}{19}.
\]
The deviations $-\frac{27}{19},-\frac8{19},\frac{11}{19}$ give the cumulants
\begin{gather*}
\kappa_{\one}(N_C,N_C)=\frac{9\cdot64+9\cdot121+729}{19^3}=\frac{126}{361},\\
\kappa_{\one}(N_C^{\times3})=\frac{-19683-9\cdot512+9\cdot1331}{19^4}=-\frac{648}{6859}.
\end{gather*}
The least odd order $a\ge3$ with $\delta\,\kappa_{\one}(N_C^{\times a})\neq0$ is $a=3$, so by Proposition~\ref{prop:mirror}(ii) $C\prec D$ exactly when $\delta\,\kappa_{\one}(N_C^{\times3})<0$, that is at a basis lean, $\delta>0$, and $D\prec C$ at a cobasis lean. (In Corollary~\ref{cor:firststages}(i) the coefficient of $v_0^3$ in the odd part of stage $0$ is $\frac2{3!}\delta\,\kappa_{\one}(N_C^{\times3})=-\frac{216}{6859}\delta$ for $C$ and $+\frac{216}{6859}\delta$ for $D$, the even parts being equal.) So
\[
B(\mathrm T_6;\sigma,\delta)=\{1,2,3\}\quad(\delta>0),\qquad B(\mathrm T_6;\sigma,\delta)=\{4,5,6\}\quad(\delta<0),
\]
each of width $2$, and the exchange $i\mapsto i+3$ maps the one onto the other (Theorem~\ref{thm:order}(iii)).

\emph{The trees.} At the centre the molecule is an atom (Definition~\ref{def:tree}): one leaf. At a basis lean the root bunch is $C$, with the patterns (Definition~\ref{def:patterns}) and probabilities (Definition~\ref{def:split})
\[
\begin{array}{l|c|l}
p & \pi_p(\one) & \text{child on }R=\{4,5,6\}\\\hline
\varnothing & \frac1{19} & \text{the one basis }R\\
\{b\},\ b\in C & \frac3{19}\text{ each} & \binom R2\\
\{b,b'\}\subset C & \frac3{19}\text{ each} & \binom R1
\end{array}
\]
$C$ itself not being a pattern. The peel (Definition~\ref{def:peel}) is the table with the one basis $R$ at $(\frac s{19},\frac t{19})$, $\binom R2$ at $(\frac{9s}{19},\frac{9t}{19})$ and $\binom R1$ at $(\frac{9s}{19},\frac{9t}{19})$; the circuit $R$ of $\binom R2$ makes $R$ one molecule, and every permutation of $R$ carries the three tables, so the room is an atom (Proposition~\ref{prop:certify}). At a cobasis lean the root bunch is $D$, with the patterns $\{x\}$ (child $\binom C2$), $\{x,y\}$ (child $\binom C1$), each of probability $\frac3{19}$, and $D$ (child: the one basis $\varnothing$ on $C$) of probability $\frac1{19}$, $\varnothing$ not being a pattern; the room $C$ is again an atom. The exchange $i\mapsto i+3$ maps the tree at $\delta$ onto the tree at $-\delta$ (Theorem~\ref{thm:mirrortree}).

\emph{Answer.} $\mathrm T_6$ is connected and prime. At the centre it is an atom. At a lean its classes are $\{1,2,3\}$ and $\{4,5,6\}$, and its Foreman bunch is $\{1,2,3\}$ at a basis lean and $\{4,5,6\}$ at a cobasis lean, both of width two. Its three Foreman trees are one atom and two trees whose root bunch, of width two, stands over a room of three marks that is an atom.
\end{solution}

\begin{exercise}[What a bunch of width two shows, and what splits off]\label{ex:cls-sixsplit}
For the list $\mathrm T_6$ of Exercise~\ref{ex:cls-sixtree} at its centre:
\begin{enumerate}[label=(\roman*)]
\item find the levels of $K$ along the scale of $C=\{1,2,3\}$, the Skeleton polynomial of $C$ at the table's gauge, its discriminant and the Skeleton values there, and show that the Skeleton values along $C$ are not rational functions of the point;
\item show that no set with at least two marks on each side and no complement of a single mark splits off (Theorem~\ref{thm:mot-pointmass}).
\end{enumerate}
\end{exercise}

\begin{solution}
Receipt: \texttt{ex05\_six\_marks.py}.

(i) \emph{The levels.} $d(C)=0$, the basis $\{4,5,6\}$ avoiding $C$, and $r(C)=2$. By the Taylor cube (Theorem~\ref{thm:taylorcube}) and Definition~\ref{def:levels}, the level $0$ is the pattern $\varnothing$ with the child of the one basis $R$, the level $1$ the patterns $\{b\}$ with the child $\binom R2$, and the level $2$ the patterns $\{b,b'\}$ with the child $\binom R1$:
\[
K(x\alpha_C,\alpha_R)=\alpha_4\alpha_5\alpha_6+x\,e_1(\alpha_C)\,e_2(\alpha_R)+x^2\,e_2(\alpha_C)\,e_1(\alpha_R),
\]
each level a product of a bunch letter and a room letter.

\emph{At the table's gauge.} $K(\one)=19$, and the Skeleton polynomial (Definition~\ref{def:skelpoly}) is
\[
\sQ^K_C(x;\one)=\frac{K(x\one_C,\one_R)}{K(\one)}=\frac{1+9x+9x^2}{19},
\]
of degree $c(C)=2$. The discriminant of $9x^2+9x+1$ is $81-36=45$, its roots are $x=\frac{-9\pm3\sqrt5}{18}$, and the Skeleton values are their negatives,
\[
\xi_{1,2}=\frac{3\mp\sqrt5}6 ,
\]
positive and irrational.

\emph{Not rational functions.} The Skeleton values at $\alpha$ are the negatives of the roots of $K(x\alpha_C,\alpha_R)$ in $x$, so
\[
(\xi_1-\xi_2)^2\,e_2(\alpha_C)^2\,e_1(\alpha_R)^2=\Delta(\alpha):=e_1(\alpha_C)^2e_2(\alpha_R)^2-4\,e_2(\alpha_C)\,e_1(\alpha_R)\,\alpha_4\alpha_5\alpha_6 .
\]
If $\xi_1,\xi_2$ were rational functions of $\alpha$, $\Delta$ would be the square of a rational function, hence, the polynomial ring being a unique factorization domain, the square of a polynomial, and so would its restriction to every line. On $\alpha=(1,1,1,1,1,\lambda)$, $e_1(\alpha_C)=e_2(\alpha_C)=3$, $e_1(\alpha_R)=2+\lambda$, $e_2(\alpha_R)=1+2\lambda$ and $\alpha_4\alpha_5\alpha_6=\lambda$:
\[
\Delta=9(1+2\lambda)^2-12\lambda(2+\lambda)=3\,(8\lambda^2+4\lambda+3),
\]
and $8\lambda^2+4\lambda+3$ has the discriminant $16-96=-80$: two simple roots, so $\Delta$ is not a square on this line. The value $45$ at the gauge shows only that the Skeleton values are irrational numbers there; the line shows that they are not rational functions of the point.

(ii) \emph{Sets with two marks on each side.} $\mathrm T_6$ is prime (Exercise~\ref{ex:cls-sixtree}), and at its centre, masses $\sigma\one$, every proper wall is positive (Proposition~\ref{prop:centrefamily}(ii)); so by Corollary~\ref{cor:mot-whatsplits}(ii) no set with at least two marks on each side splits off. Each such set has width two or three and enters every period of the list through a law carried by no point (Corollary~\ref{cor:mot-whatsplits}(i)).

\emph{Complements of single marks.} $\cE\setminus\{e\}$ has the width $c(\{e\})=1$ (Theorem~\ref{thm:window}); with both letters active, Corollary~\ref{cor:mot-whatsplits}(iv) says it splits off exactly when the deletion and the contraction at $e$ have one even letter, $\mathbf E_{\mathrm T_6\setminus e}=\mathbf E_{\mathrm T_6/e}$ (Definition~\ref{def:delcon}). For $e\in C$ the deletion has as its bases all ten triples of $\cE\setminus\{e\}$ and the contraction the nine pairs $\{x,y\}$ of $\cE\setminus\{e\}$ with $\{e,x,y\}\neq\{1,2,3\}$; for $e\in D$ the deletion has nine bases and the contraction ten. At the table's gauge $\mathbf E=K\Psi$ is the square of the number of bases, so
\[
\mathbf E_{\mathrm T_6\setminus e}(\one)=100,\ \ \mathbf E_{\mathrm T_6/e}(\one)=81\quad(e\in C);\qquad \mathbf E_{\mathrm T_6\setminus e}(\one)=81,\ \ \mathbf E_{\mathrm T_6/e}(\one)=100\quad(e\in D):
\]
the two even letters differ, and no complement of a mark splits off at the centre.

\emph{Answer.} Along $C$ the levels are $\alpha_4\alpha_5\alpha_6$, $e_1(\alpha_C)e_2(\alpha_R)$ and $e_2(\alpha_C)e_1(\alpha_R)$; at the table's gauge $K(x\one_C,\one_R)=1+9x+9x^2$, with the discriminant $45$ and the Skeleton values $\frac{3\mp\sqrt5}6$; on the line $(1,1,1,1,1,\lambda)$ the discriminant is $3(8\lambda^2+4\lambda+3)$, not a square, so the Skeleton values are not rational functions of the point. At the centre no set with at least two marks on each side and no complement of a mark splits off: $\mathrm T_6$ is its own and only factor. (At a full lean each complement of a mark does split off, Corollary~\ref{cor:mot-whatsplits}(iii).)
\end{solution}

\section{Skeleton values at width two}\label{sec:exc-skel}

\begin{exercise}[Skeleton values of $\binom{[4]}2$ along a pair]\label{ex:cls-u24skel}
For the list $\binom{[4]}2$ and its bunch $B=\{1,2\}$, find the width of $B$, the levels of $K$ along the scale of $B$, the Skeleton polynomial and the Skeleton values at the table's gauge, and show that the Skeleton values are not rational functions of the point. Why does the value of the discriminant at the table's gauge not settle this?
\end{exercise}

\begin{solution}
Receipt: \texttt{ex06\_skeleton\_u24.py}.

\emph{The width.} The basis $\{3,4\}$ avoids $B$ and the basis $\{1,2\}$ is $B$, so the window of $B$ is $\{0,1,2\}$: $d(B)=0$, $r(B)=2$ and $c(B)=2$ (Theorem~\ref{thm:window}).

\emph{The levels.} The patterns are $\varnothing$ (child: the one basis $\{3,4\}$), $\{1\}$ and $\{2\}$ (child $\binom{\{3,4\}}1$, with $K=\alpha_3+\alpha_4$) and $\{1,2\}$ (child: the one basis $\varnothing$). By Theorem~\ref{thm:taylorcube} and Definition~\ref{def:levels},
\[
K(\rho\alpha_1,\rho\alpha_2,\alpha_3,\alpha_4)=\alpha_3\alpha_4+\rho\,(\alpha_1+\alpha_2)(\alpha_3+\alpha_4)+\rho^2\alpha_1\alpha_2 .
\]
The Skeleton polynomial is $\sQ^K_B(x;\alpha)=K(x\alpha_B,\alpha_R)/K(\alpha)$ (Definition~\ref{def:skelpoly}), of degree $2$, whose roots $-\xi_1,-\xi_2$ have
\[
\xi_1+\xi_2=\frac{(\alpha_1+\alpha_2)(\alpha_3+\alpha_4)}{\alpha_1\alpha_2},\qquad \xi_1\xi_2=\frac{\alpha_3\alpha_4}{\alpha_1\alpha_2}.
\]
The list is identically self-dual and $d^\ast(B)=0$, so $\Psi$ has the same levels along $B$ and the same Skeleton values.

\emph{At the table's gauge.} $K(x\one_B,\one_R)=1+4x+x^2$ and $K(\one)=6$; the discriminant is $16-4=12$, and the Skeleton values are $\xi=2\pm\sqrt3$, irrational.

\emph{What the value at the gauge decides.} $(\xi_1-\xi_2)^2(\alpha_1\alpha_2)^2=\Delta(\alpha):=(\alpha_1+\alpha_2)^2(\alpha_3+\alpha_4)^2-4\alpha_1\alpha_2\alpha_3\alpha_4$, and $\Delta(\one)=12$ is not the square of a rational number. That shows that the Skeleton values are irrational numbers at the gauge, and no more: if $\Delta$ were $3q(\alpha)^2$ for a polynomial $q$ with rational coefficients, $\xi_{1,2}=\frac{(\alpha_1+\alpha_2)(\alpha_3+\alpha_4)\mp\sqrt3\,q}{2\alpha_1\alpha_2}$ would be rational functions of the point, with the value $12$ at the gauge.

\emph{Not rational functions.} If $\xi_1,\xi_2$ are rational functions of $\alpha$, $\Delta$ is the square of a rational function, so of a polynomial (unique factorization in the polynomial ring), and so is its restriction to every line: $\Delta$ would be a constant times the square of a polynomial with rational coefficients. On $\alpha=(1,1,1,\lambda)$,
\[
\Delta=4(1+\lambda)^2-4\lambda=4(\lambda^2+\lambda+1),
\]
whose roots $\frac{-1\pm\sqrt{-3}}2$ are simple. So $\Delta$ is not a constant times a square, and the Skeleton values along $B$ are not rational functions of the point.

\emph{Answer.} $B$ has width two, with the levels $\alpha_3\alpha_4$, $(\alpha_1+\alpha_2)(\alpha_3+\alpha_4)$ and $\alpha_1\alpha_2$; at the table's gauge the Skeleton polynomial is $(1+4x+x^2)/6$ and the Skeleton values are $2\pm\sqrt3$. The discriminant is $4(\lambda^2+\lambda+1)$ on the line $(1,1,1,\lambda)$, not a constant times a square, so the Skeleton values are not rational functions of the point. The value $12$ at the gauge only excludes a rational square, not a constant times a square. At width one, by contrast, the one Skeleton value is the quotient of two hinge values (Corollary~\ref{cor:widthonestage}).
\end{solution}

\section{The record of a peel}\label{sec:exc-record}

\begin{exercise}[Peeling with and without the record]\label{ex:cls-record}
Let
\[
F=\Big(\Big(\binom{[5]}2;1,0\Big),\Big(\binom{[5]}3;-1,0\Big),\Big(\binom{[5]}4;4,0\Big)\Big)
\]
on $\cE=\{1,\dots,5\}$, with the masses $3\one$. Show that it is on its balance wall with every proper wall positive and that the table's gauge is native; peel it at $\{1\}$ with the native split of the table's gauge, write its record and its peel, and compare the letters they carry. What changes when every exponent is nonnegative?
\end{exercise}

\begin{solution}
Receipt: \texttt{ex07\_record.py}.

\emph{Walls.} The exponents are $s_j=1,-1,4$ and $t_j=0$; the ranks are $2,3,4$. The balance is $\sum_js_jr_j=2-3+16=15=m_\cE$. In $\binom{\cE}r$ a set $S$ meets every basis in at least $d(S)=\max\big(0,|S|-(5-r)\big)$ marks (Theorem~\ref{thm:window}), so the walls $\kappa_S=m_S-\sum_js_jd_j(S)$ (Definition~\ref{def:family}) depend on $|S|$ alone:
\[
\begin{array}{c|ccccc}
|S| & 1&2&3&4&5\\\hline
(d_1,d_2,d_3) & (0,0,0)&(0,0,1)&(0,1,2)&(1,2,3)&(2,3,4)\\
\kappa_S & 3&2&2&1&0
\end{array}
\]
every proper wall is positive and $\kappa_\cE=0$.

\emph{The native gauge.} On the balance wall the target is $\mvec$ (Definition~\ref{def:native}). At the table's gauge the occupation in $\binom{\cE}r$ is $\frac r5$ at every mark, so the usage is $1\cdot\frac25-1\cdot\frac35+4\cdot\frac45=3=m_e$: the gauge $\one$ is native. An exponent is negative, so $F$ has no realization (Definition~\ref{def:realization}) and no Foreman tree (Definition~\ref{def:tree}); its native split (Definition~\ref{def:split}), its record and its peel are defined at the native gauge $\one$.

\emph{The record.} The patterns of $\{1\}$ in $\binom{\cE}r$ are $\varnothing$, with the child $\binom Rr$ on $R=\{2,3,4,5\}$ and the probability $\binom4r/\binom5r=\frac{5-r}5$, and $\{1\}$, with the child $\binom R{r-1}$ and the probability $\frac r5$ (Definitions~\ref{def:patterns} and~\ref{def:split}). The record (Definition~\ref{def:record}) is the six indexed children, each with the exponents $\pi_{j,p}(\one)(s_j,0)$:
\[
\begin{array}{c|c|c|c}
\text{table }j & p & \text{child} & \text{exponents}\\\hline
1 & \varnothing & \binom R2 & (\frac35,0)\\
1 & \{1\} & \binom R1 & (\frac25,0)\\
2 & \varnothing & \binom R3 & (-\frac25,0)\\
2 & \{1\} & \binom R2 & (-\frac35,0)\\
3 & \varnothing & \binom R4 & (\frac45,0)\\
3 & \{1\} & \binom R3 & (\frac{16}5,0)
\end{array}
\]
On $R$ the children lie on their balance wall at the masses $3\one_R$, $\frac35\cdot2+\frac25\cdot1-\frac25\cdot3-\frac35\cdot2+\frac45\cdot4+\frac{16}5\cdot3=12$, and their usage at $\one_R$ is $3$ at every mark (Theorem~\ref{thm:hered}). The record determines $F$ (Proposition~\ref{prop:record}): table $j$ is the gluing of its two children, and their exponents add to $1$, $-1$ and $4$.

\emph{The peel.} Definition~\ref{def:peel} merges the children with the same bases, adding their exponents:
\[
\Big(\binom R1;\tfrac25,0\Big),\qquad\Big(\binom R2;\tfrac35-\tfrac35,0\Big)=\Big(\binom R2;0,0\Big),\qquad\Big(\binom R3;\tfrac{14}5,0\Big),\qquad\Big(\binom R4;\tfrac45,0\Big).
\]
The table $\binom R2$ has both exponents zero: it is not active, and the peel carries no letter $e_2(\alpha_R)$ (Definition~\ref{def:family}). The record carries it twice, in the child $(1,\varnothing)$ at the exponent $\frac35$ and in the child $(2,\{1\})$ at $-\frac35$. The two words agree,
\[
e_1(\alpha_R)^{-2/5}\,e_3(\alpha_R)^{-14/5}\,e_4(\alpha_R)^{-4/5},
\]
since $e_2^{-3/5}e_2^{3/5}=1$; what the peel loses is which pattern of which table each share came from, which the record keeps (the paragraph after Proposition~\ref{prop:record}), and with it the letter $e_2(\alpha_R)$ that the tables $1$ and $2$ hand to the room.

\emph{Nonnegative exponents.} When every $s_j,t_j\ge0$, every child carries $\pi_{j,p}(\one)\,(s_j,t_j)$ with $\pi_{j,p}(\one)>0$ (Proposition~\ref{prop:herecentre}); a merged exponent is a sum of nonnegative terms, zero only when each term is, so merging loses no letter, and a full lean is handed down full, table by table. This is the scope of the realizations, where the Foreman tree is read.

\emph{Answer.} $F$ is on its balance wall with the walls $3,2,2,1$ at the sets of $1,2,3,4$ marks, and the table's gauge is native. Its record at $\{1\}$ is the six children of the table displayed; its peel is $\binom R1$, $\binom R3$, $\binom R4$ at $\frac25,\frac{14}5,\frac45$, the child $\binom R2$ cancelling between the tables $1$ and $2$: merging loses the letter $e_2$ on the room, which the record keeps. With nonnegative exponents no letter is lost.
\end{solution}

\chapter{Exercises on evaluation and construction}\label{ch:exconstruct}

The exercises of this chapter work the statements of Parts~IV and~V on particular lists. A set of width two whose stations of room sides are single room sides cannot be opened with one share per station, the room's seam function being infinite at every share, and it is opened with the transverse pins instead. The width of the marks left in place by a substitution is read off the core, bounded by the number of the other cells, and every width occurs. The uniform list of rank two on five marks is fused by the procedure of Foreman fusion: its fusion, its one major, its cousins and the symmetry the fusion adds. Every number displayed is exact; the receipts that recompute them are the scripts named in each solution.

\section{One share per station}\label{sec:exk-stations}

\begin{exercise}[One share per station at a set of width two]\label{ex:con-transverse}
Let $\mathrm Q_5$ be the list on $\{1,\dots,5\}$ whose bases are the eight sets of three marks other than $\{1,2,5\}$ and $\{3,4,5\}$, at its centre $(\one;1,1)$, and let $S=\{1,2,3\}$, with the room $R=\{4,5\}$.
\begin{enumerate}[label=(\roman*)]
\item Find the width of $S$, its patterns and room sides at every excess, its stations, its seam letters and the room station letters (Definition~\ref{def:seam}), and check the room-station factorization.
\item Read the period with one share per station (Theorem~\ref{thm:gluewidth}(iii)): write the pins of the seam and the room word, and show that the seam function of the room is $+\infty$ at every share.
\item Read it with the transverse pins of Theorem~\ref{thm:transversepins}: find the lattice, the pins, the reduced periods and the share side, and check that the native split at the stations lies on the pins, with the masses of $S$ inside the Newton sum of the seam word there.
\end{enumerate}
\end{exercise}

\begin{solution}
Receipt: \texttt{ex08\_transverse.py}.

(i) \emph{The list.} Its circuits (Definition~\ref{def:molecule}) are the two triples $\{1,2,5\}$ and $\{3,4,5\}$ that are not bases and the set $\{1,2,3,4\}$, all of whose triples are bases; the rank is $3$ and the corank $2$. In $S$ every basis has one to three marks: $r(S)=3$, $S$ being a basis, $r(R)=2$, and by Definition~\ref{def:skeleton} and Theorem~\ref{thm:window}
\begin{gather*}
d(S)=r-r(R)=1,\qquad d^\ast(S)=|S|-r(S)=0,\qquad c(S)=r(S)-d(S)=2,\\
d(R)=d^\ast(R)=0,\qquad c(R)=2 .
\end{gather*}

\emph{Patterns and room sides.} An excess $i$ is the level $d(S)+i=1+i$ (Definition~\ref{def:seam}); its patterns have $1+i$ marks and its room sides $2-i$. With the children (Definition~\ref{def:patterns}) and the room children:
\[
\begin{array}{c|l|l}
i & \text{patterns : children on }R & \text{room sides : room children on }S\\\hline
0 & \{1\},\{2\}\ :\ \{45\} & \{4,5\}\ :\ \{1\},\{2\}\\
1 & \{1,2\}\ :\ \{4\};\quad \{1,3\},\{2,3\}\ :\ \{4\},\{5\} & \{4\}\ :\ \{12\},\{13\},\{23\};\quad \{5\}\ :\ \{13\},\{23\}\\
2 & \{1,2,3\}\ :\ \{\varnothing\} & \varnothing\ :\ \{123\}
\end{array}
\]
The two room sides of the excess $1$ have room children with three and with two bases, not multiples of one another; so every room side is its own station: $D_0=\{\{4,5\}\}$, $D_1=\{\{4\}\}$, $D'_1=\{\{5\}\}$ and $D_2=\{\varnothing\}$. The stations of patterns are $\{\{1\},\{2\}\}$, $\{\{1,2\}\}$, $\{\{1,3\},\{2,3\}\}$ and $\{\{1,2,3\}\}$. For a list every multiplier and pair constant is $1$.

\emph{Seam letters and room station letters.} The seam letters are the letters of the room children, on $S$, and the room station letters are $\mathsf y_D=\alpha_R^{X}$ and $\mathsf y^\ast_D=\alpha_R^{R\setminus X}$ for the one room side $X$ of $D$ (Definition~\ref{def:seam}(iii)):
\[
\begin{array}{c|c|c|c|c}
D & V_D=K_{(D)} & V^\ast_D=\Psi_{(D)} & \mathsf y_D & \mathsf y^\ast_D\\\hline
D_0 & \alpha_1+\alpha_2 & \alpha_3(\alpha_1+\alpha_2) & \alpha_4\alpha_5 & 1\\
D_1 & e_2(\alpha_S) & e_1(\alpha_S) & \alpha_4 & \alpha_5\\
D'_1 & \alpha_3(\alpha_1+\alpha_2) & \alpha_1+\alpha_2 & \alpha_5 & \alpha_4\\
D_2 & \alpha_1\alpha_2\alpha_3 & 1 & 1 & \alpha_4\alpha_5
\end{array}
\]
\emph{The room-station factorization} (Proposition~\ref{prop:seamidentity}(iii)):
\begin{gather*}
K=(\alpha_1+\alpha_2)\,\alpha_4\alpha_5+e_2(\alpha_S)\,\alpha_4+\alpha_3(\alpha_1+\alpha_2)\,\alpha_5+\alpha_1\alpha_2\alpha_3,\\
\Psi=\alpha_3(\alpha_1+\alpha_2)+e_1(\alpha_S)\,\alpha_5+(\alpha_1+\alpha_2)\,\alpha_4+\alpha_4\alpha_5 ,
\end{gather*}
the eight bases and the eight cobases, each once.

\emph{The walls.} At $(\one;1,1)$, $\kappa_S=m_S-d(S)-d^\ast(S)=2$, $\kappa_R=m_R-d(R)-d^\ast(R)=2$, and the top wall of $S$ is $(s+t)\,c(S)-\kappa_S=2=\kappa_R$ (Definition~\ref{def:seampiece}): the glue is on its balance wall.

(ii) \emph{The pins of the seam.} One share per station and letter: $a_D$ for $K$ and $b_D$ for $\Psi$, $D\in\{D_0,D_1,D'_1,D_2\}$. By Lemma~\ref{lem:seamword} the pins are the two letter totals and the scale pin, whose coefficients are the excesses $i(D)$ for $K$ and $c(S)-i(D)$ for $\Psi$:
\begin{gather*}
a_{D_0}+a_{D_1}+a_{D'_1}+a_{D_2}=1,\qquad b_{D_0}+b_{D_1}+b_{D'_1}+b_{D_2}=1,\\
a_{D_1}+a_{D'_1}+2a_{D_2}+2b_{D_0}+b_{D_1}+b_{D'_1}=\kappa_S=2 .
\end{gather*}
\emph{The room word.} The room is the partner with the room letters $L_D=\mathsf y_D$ and $L^\ast_D=\mathsf y^\ast_D$ (the paragraph after Definition~\ref{def:seampiece}), so its word is the monomial
\begin{gather*}
W_R(\mathbf a,\mathbf b)=\prod_D\mathsf y_D^{-a_D}\,\mathsf y^\ast_D{}^{-b_D}=\alpha_4^{-p_4}\alpha_5^{-p_5},\\
p_4=a_{D_0}+a_{D_1}+b_{D'_1}+b_{D_2},\qquad p_5=a_{D_0}+a_{D'_1}+b_{D_1}+b_{D_2},
\end{gather*}
and $p_4+p_5=2\big(|\mathbf a|+|\mathbf b|\big)-\big(a_{D_1}+a_{D'_1}+2a_{D_2}+2b_{D_0}+b_{D_1}+b_{D'_1}\big)=4-2=2=m_4+m_5$ on the pins.

\emph{Its seam function.} The seam function of the room is the period of its word at $\mvec_R=(1,1)$ (Lemma~\ref{lem:seamword}, read on the partner, whose pins on the balance wall are those of the seam; Theorem~\ref{thm:gluewidth}(iii)). On the section $\alpha_5=1$ (Definition~\ref{def:posword}) it is
\[
\int_0^\infty\alpha_4^{\,1-p_4}\,\frac{d\alpha_4}{\alpha_4}=+\infty
\]
at every real value of $p_4$: the integrand is a power, which is not integrable at $0$ when its exponent is at most $-1$ and not at $\infty$ when it is at least $-1$. Every room letter is a monomial because every room side is its own station, so the Newton sum of the room word is one point, of dimension $0<|R|-1$, and Theorem~\ref{thm:poswordconv} excludes convergence at every share. No point of the pins has both seam functions finite: the display of Theorem~\ref{thm:gluewidth}(iii), a fibre integral of the dimension $2\cdot4-2-1=5$, does not hold here. Its parts (i) and (ii) hold in every case.

(iii) \emph{The Newton sums.} At shares with positive real parts the seam word of $S$,
\begin{multline*}
W_S(\mathbf a,\mathbf b)=\prod_DV_D^{-a_D}V^\ast_D{}^{-b_D}\\
=(\alpha_1+\alpha_2)^{-a_{D_0}-b_{D'_1}}\big(\alpha_3(\alpha_1+\alpha_2)\big)^{-b_{D_0}-a_{D'_1}}e_2(\alpha_S)^{-a_{D_1}}e_1(\alpha_S)^{-b_{D_1}}(\alpha_1\alpha_2\alpha_3)^{-a_{D_2}},
\end{multline*}
has a Newton sum that contains $b_{D_1}$ times the triangle $\operatorname{conv}\{\one_{\{1\}},\one_{\{2\}},\one_{\{3\}}\}$ of $e_1(\alpha_S)$ (Definition~\ref{def:posword}): its space of directions is $\vec{\mathcal N}_S=\{x\in\mathbb R^S:\ x_1+x_2+x_3=0\}$. The room word is a monomial, $\vec{\mathcal N}_R=0$.

\emph{The lattice and the pins.} In Theorem~\ref{thm:transversepins},
\[
M=\vec{\mathcal N}_S^\perp\oplus\vec{\mathcal N}_R^\perp=\mathbb R\one_S\oplus\mathbb R^{\{4,5\}},
\]
of dimension $3$, containing $\one$; $M\cap\mathbb Z^5=\{(a,a,a,b,c)\}$ has the basis $\one,\mathbf e_4,\mathbf e_5$, since $\one_S=\one-\mathbf e_4-\mathbf e_5$. So $q=2$, $\mathbf v_1=\mathbf e_4$ and $\mathbf v_2=\mathbf e_5$. The fixed monomials of the seam letters live on $S$, and those of the room letters are the room letters themselves, so $\ip{\mathbf n(\mathbf a,\mathbf b)}{\mathbf e_4}=p_4$ and $\ip{\mathbf n(\mathbf a,\mathbf b)}{\mathbf e_5}=p_5$. The pins are the letter totals and
\[
a_{D_0}+a_{D_1}+b_{D'_1}+b_{D_2}=m_4=1,\qquad a_{D_0}+a_{D'_1}+b_{D_1}+b_{D_2}=m_5=1 :
\]
four pins on eight shares, a fibre of the dimension $4$. With the totals they imply the scale pin, by the identity for $p_4+p_5$ in (ii).

\emph{The reduced periods.} $\vec{\mathcal N}_S^\perp=\mathbb R\one_S$, so the reduced period of the seam word is its integral over $\mathbb R^S/\mathbb R\one_S$, that is over a section: the period $P_{W_S(\mathbf a,\mathbf b)}(\one_S)$ (Definition~\ref{def:posword}). $\vec{\mathcal N}_R^\perp=\mathbb R^R$, so the quotient is a point, and the reduced period of the room word is the value of $\alpha_4^{m_4-p_4}\alpha_5^{m_5-p_5}$, which is $1$ on the pins, the room letters having the coefficient $1$.

\emph{The share side.} The share kernels are $\Gst41{\mathbf a}{1}=\prod_D\Gamma(a_D)$ and $\Gst41{\mathbf b}{1}=\prod_D\Gamma(b_D)$, $\Gamma(1)=1$ (Definition~\ref{def:stacked}). So, on every contour of the pins with positive real parts at which $P_{W_S}$ converges,
\[
P_{\mathrm Q_5}(\one;1,1)=\int_{\text{pins}}\ \prod_D\Gamma(a_D)\,\Gamma(b_D)\ P_{W_S(\mathbf a,\mathbf b)}(1,1,1),
\]
the pins being $|\mathbf a|=|\mathbf b|=1$ and $p_4=p_5=1$.

\emph{The native split.} At the table's gauge the basis law is uniform on the eight bases and the cobasis law is the same law (Theorem~\ref{thm:innerMGF}). $T\cap R$ is $\{4,5\}$ at $145,245$, $\{4\}$ at $124,134,234$, $\{5\}$ at $135,235$ and $\varnothing$ at $123$, so Lemma~\ref{lem:nativestations} gives
\[
a^\star_D=b^\star_D=\tfrac14,\ \tfrac38,\ \tfrac14,\ \tfrac18\qquad(D=D_0,D_1,D'_1,D_2),
\]
with the totals $1$ and $1$, $p_4=\frac14+\frac38+\frac14+\frac18=1$ and $p_5=\frac14+\frac14+\frac38+\frac18=1$: the native split lies on the transverse pins, with positive coordinates. There $\one_S$ lies in the relative interior of the Newton sum of the seam word (Lemma~\ref{lem:nativestations}); the receipt checks it exactly, the support function of the Newton sum exceeding $\ip u{\one_S}$ by at least $\frac98$ at every edge normal $u$ of its summands. So $P_{W_S(\mathbf a^\star,\mathbf b^\star)}(\one_S)$ converges (Theorem~\ref{thm:poswordconv}), and the contour through the native split carries the display.

\emph{Answer.} $S$ has width two; its four stations of room sides are the single room sides $\{4,5\}$, $\{4\}$, $\{5\}$, $\varnothing$, the room letters are monomials, and the seam function of the room is $+\infty$ at every share, so the period does not open with one share per station. With the transverse pins $p_4=p_5=1$ in place of the scale pin it opens as the fibre integral of dimension four displayed, the seam word of $S$ against the share kernels, the room contributing $1$.
\end{solution}

\section{The width of the marks left in place}\label{sec:exk-width}

\begin{exercise}[The width of a bunch of a substitution]\label{ex:con-widthbound}
Let $\mathrm M_8=\binom{\{b_1,b_2,h_1,h_2\}}2\big\langle b_1,b_2,\mathcal A_1,\mathcal A_2\big\rangle$, the marks $b_1,b_2$ left in place and the cells $\mathcal A_i=\binom{\{x_i,y_i,z_i,g_i\}}2$ at $h_i$, with $R_i=\{x_i,y_i,z_i\}$ (Definition~\ref{def:substitution}).
\begin{enumerate}[label=(\roman*)]
\item Find the rank and the bases of $\mathrm M_8$, and the widths of $\{b_1\}$, $\{b_2\}$ and $B=\{b_1,b_2\}$, from the bases and from the core; compare with the bound of Corollary~\ref{cor:mot-widthbound}.
\item Write $K$ along the scale of $B$.
\item For $c=1,2,3$ let $B=\{b_1,\dots,b_c\}$ be left in place in the core $\binom{B\sqcup\{h_1,\dots,h_c\}}c$, with the cells $\binom{\{x_i,y_i,z_i,g_i\}}2$ at the $h_i$. Find the width of $B$.
\end{enumerate}
\end{exercise}

\begin{solution}
Receipt: \texttt{ex09\_width\_subst.py}.

(i) \emph{Rank and bases.} The core has the rank $2$, each cell $\mathcal A_i$ the rank $2$, and a mark left in place the cell $\binom{\{b,g\}}1$ of rank $1$; so $\mathrm M_8$ is a table of rank $2+(2-1)+(2-1)=4$ on eight marks (Proposition~\ref{prop:substitution}(i)). Its bases (Definition~\ref{def:substitution}) are read off the six bases $T$ of the core: $b\in T\cap B$ gives the mark $b$, and at $h_i$ the cell gives a basis of $\mathcal A_i\setminus g_i=\binom{R_i}2$ when $h_i\in T$ and one of $\mathcal A_i/g_i=\binom{R_i}1$ when $h_i\notin T$, three choices each: $6\cdot9=54$ bases.

\emph{Widths.} From the bases, the numbers of marks of $\{b_1\}$, $\{b_2\}$ and $B$ in a basis fill the windows $\{0,1\}$, $\{0,1\}$ and $\{0,1,2\}$ (Theorem~\ref{thm:window}): the widths are $1$, $1$ and $2$. In the core $\binom{\{b_1,b_2,h_1,h_2\}}2$ a pair meets one mark in $0$ or $1$ marks and a pair of marks in $0$, $1$ or $2$: the same widths. That is Corollary~\ref{cor:mot-widthbound}: a set of marks left in place has the width it has in the core. In Theorem~\ref{thm:substwalls} the cell $\binom{\{b,g\}}1$ of a mark $b$ of $S\subseteq\{b_1,b_2\}$ has $d(\{b\})=d^\ast(\{b\})=0$ and $d(\{b,g\})=d^\ast(\{b,g\})=1$, the cells $\mathcal A_i$ meet $S$ in $\varnothing$ with $d(\{g_i\})=d^\ast(\{g_i\})=0$, so $D(S)=D^\ast(S)=S$ and $d(S)=d_{\mathcal H}(S)=0$, $d^\ast(S)=d^\ast_{\mathcal H}(S)=0$.

\emph{The bound.} With $L=\{b_1,b_2\}$ and the two cells, $c(S)\le|L\setminus S|+2$: $1\le3$ for $\{b_1\}$ and $\{b_2\}$, and $2\le2$ for $B$, where it is attained. For the list $\mathrm P_5=\binom{\{1,2,h\}}1\langle1,2,\binom{\{3,4,5,g\}}2\rangle$ of Exercise~\ref{ex:cls-five} the bound at its Foreman bunch $\{1,2\}$ is $0+1=1$, attained as well.

(ii) \emph{The letter along $B$.} The cells have $K_{\mathcal A_i/g_i}=e_1(\alpha_{R_i})$ and the hinge values $\gamma_i:=\gamma^K_{\mathcal A_i}=e_2(\alpha_{R_i})/e_1(\alpha_{R_i})$ (Definition~\ref{def:substitution}); their contractions $\binom{R_i}1$ are the duals of their deletions $\binom{R_i}2$, so also $\gamma^\Psi_{\mathcal A_i}=\gamma_i$ (Corollary~\ref{cor:piecesmarks}(ii)). The core letter is $e_2$ on four marks, so by Theorem~\ref{thm:substletters}
\[
K(\tau\alpha_B,\alpha_R)=e_1(\alpha_{R_1})\,e_1(\alpha_{R_2})\Big(\gamma_1\gamma_2+\tau\,(\alpha_{b_1}+\alpha_{b_2})(\gamma_1+\gamma_2)+\tau^2\alpha_{b_1}\alpha_{b_2}\Big):
\]
three levels, the width two, and the room sees $B$ only through the two hinge values of the cells and the marks of $B$.

(iii) In the core $\binom{B\sqcup\{h_1,\dots,h_c\}}c$, $|B|=c$, the sets of $c$ marks meet $B$ in every number of marks from $0$ to $c$, so $B$ has the width $c$ (Theorem~\ref{thm:window}); no $h_i$ is a loop or a coloop, and by Corollary~\ref{cor:mot-widthbound} the width stays $c$ after the cells are substituted. From the bases:
\[
\begin{array}{c|c|c|c|c|c}
c & \text{marks} & \text{rank} & \text{bases} & \text{window of }B & c(B)\\\hline
1 & 4 & 2 & 6 & \{0,1\} & 1\\
2 & 8 & 4 & 54 & \{0,1,2\} & 2\\
3 & 12 & 6 & 540 & \{0,1,2,3\} & 3
\end{array}
\]
At $c=1$ the substitution $\binom{\{b_1,h_1\}}1\langle b_1,\mathcal A_1\rangle$ is the list $\binom{\{b_1,x_1,y_1,z_1\}}2$ itself, and at $c=2$ it is $\mathrm M_8$. In each case the bound $|L\setminus B|+c=c$ is attained.

\emph{Answer.} $\mathrm M_8$ is a list of rank four on eight marks with $54$ bases; $\{b_1\}$ and $\{b_2\}$ have width one and $B=\{b_1,b_2\}$ width two, as in the core, against the bounds $3$, $3$ and $2$. Along $B$, $K=e_1(\alpha_{R_1})e_1(\alpha_{R_2})(\gamma_1\gamma_2+\tau(\alpha_{b_1}+\alpha_{b_2})(\gamma_1+\gamma_2)+\tau^2\alpha_{b_1}\alpha_{b_2})$. In $\binom{B\sqcup\{h_1,\dots,h_c\}}c$ with cells at the $h_i$ the bunch $B$ has the width $c$ for $c=1,2,3$: a bunch that is all the marks left in place has at most as much width as there are cells, and every width occurs.
\end{solution}

\section{Fusing a uniform list}\label{sec:exk-fusion}

\begin{exercise}[Fusing $\binom{[5]}2$]\label{ex:con-fusion}
Run the procedure of Section~\ref{sec:fusionproc} on the list $\binom{[5]}2$, read at its centre: decide each step, find its fusion width, its tell, its Foreman fusion with its letter, its lifts and its majors, and its cousins; read the period of the fusion through two of its cousins; and find the index of the bijections carrying $\binom{[5]}2$, times the permutations of the fusion bunch, among those carrying its fusion.
\end{exercise}

\begin{solution}
Receipt: \texttt{ex10\_fusion\_u25.py}.

\emph{(1) Peel.} Every permutation of $[5]$ maps pairs to pairs and carries the list, so at the centre, and at every lean, all marks lie in one class (Corollary~\ref{cor:certifyone}): the realization is an atom (Definition~\ref{def:view}) and nothing is peeled (Definition~\ref{def:tree}).

\emph{(2) Go down.} The complement of a basis is a triple, not a basis: the list is not identically self-dual (Corollary~\ref{cor:piecesmarks}(iii)).

\emph{(3) Go up.} The rank is $r=2$ and the corank $r^\ast=3$, so the fusion width is $c=r^\ast-r=1$ (Definition~\ref{def:fusion}). The circuits (Definition~\ref{def:molecule}) are the ten triples, every pair being independent; the cocircuits, the least sets meeting every pair, are the five sets of four marks. Each cocircuit is the union of the four triples inside it, so the tell (Theorem~\ref{thm:fusiontell}(b)) holds. In its form (c): for a mark $e$ and a split of the other four marks into two halves, one half has at least two marks, the rank $2$, which $e$ does not raise. The form (a), $r^\ast(Y)-r^\ast(X)\ge r(Y)-r(X)$ for all $X\subseteq Y$, the receipt checks on the rank functions. At $c=1$ the weights play no part: $\binom{[5]}2$ is fusable.

\emph{The fusion.} With the fusion bunch $G=\{g\}$, Theorem~\ref{thm:fusion}(v) gives the bases $T+g$, $T$ a pair of $[5]$, and $[5]\setminus T$, the triples of $[5]$: ten and ten, all twenty triples of $[5]\sqcup\{g\}$. So
\[
\cB^\sharp=\binom{[5]\sqcup\{g\}}3,\qquad K^\sharp=\Psi+\alpha_gK=e_3(\alpha_{[5]})+\alpha_g\,e_2(\alpha_{[5]})=e_3(\alpha_{[5]\sqcup\{g\}})=\Psi^\sharp ,
\]
identically self-dual, the complement of a triple of six marks being a triple. The lifts (Definition~\ref{def:fusion}(ii)) are $\mathcal L_0=\binom{[5]}2$ and $\mathcal L_1=\binom{[5]}3$, its dual, with the lift letters $L_0=\Psi=e_3$ and $L_1=K=e_2$; at $c=1$ there are no middle weights and no lift ratios, and $K^\sharp=L_0+\alpha_gL_1$ is Theorem~\ref{thm:fusion}(ii).

\emph{The majors.} A table on $[5]\sqcup\{g\}$ whose contraction of $g$ is $\binom{[5]}2$ has as its bases through $g$ the sets $T+g$, $T\in\binom{[5]}2$, and one whose deletion of $g$ is $\binom{[5]}3$ has as its bases avoiding $g$ the triples of $[5]$ (Definition~\ref{def:delcon}). A major (Definition~\ref{def:fusion}(i)) has both, so its bases are the twenty triples, and $\cB^\sharp$ is the only major, as Theorem~\ref{thm:fusion}(v) says.

\emph{One Beta corner.} On the balance wall of the fusion, $m_{[5]}+m_g=2\sigma r^\ast=6\sigma$, with every proper wall positive, Corollary~\ref{cor:fusionone} gives
\[
P_{\binom{[6]}3}\big(\mvec|_{[5]},m_g;s,t\big)=\Gst21{m_g,\ 2\sigma-m_g}{2\sigma}\,P_{\binom{[5]}2}\big(\mvec|_{[5]};m_g,2\sigma-m_g\big),
\]
$[6]$ written for $[5]\sqcup\{g\}$. It can be read off the letter: the fused table is the family $\big(\binom{[6]}3;2\sigma,0\big)$ (the paragraph after Definition~\ref{def:universalpiece}), and on the section $|\alpha_{[5]}|=1$ the scale of $g$ integrates out, with $\alpha_g=(e_3/e_2)v$,
\begin{align*}
\int_0^\infty\alpha_g^{\,m_g-1}\big(e_3+\alpha_g\,e_2\big)^{-2\sigma}d\alpha_g&=e_3^{\,m_g-2\sigma}e_2^{-m_g}\int_0^\infty v^{m_g-1}(1+v)^{-2\sigma}dv\\
&=\Gst21{m_g,\ 2\sigma-m_g}{2\sigma}\,e_3^{\,m_g-2\sigma}e_2^{-m_g},
\end{align*}
$e_k=e_k(\alpha_{[5]})$, the last integral being the shape integral of two marks (Theorem~\ref{thm:dirichlet}) after $v=u/(1-u)$; what is left is the form of $\binom{[5]}2$ with $K=e_2$ at the exponent $m_g$ and $\Psi=e_3$ at $2\sigma-m_g$. The masses of $\binom{[5]}2$ are then on its own balance wall, $m_{[5]}=2m_g+3(2\sigma-m_g)$.

\emph{The cousins.} The cousins of $\binom{[6]}3$ (Proposition~\ref{prop:cousins}) are its six one-mark contractions $\binom{[6]\setminus\{x\}}2$, each a copy of $\binom{[5]}2$, the cousin at $g$ being $\binom{[5]}2$ itself; the deletion $\binom{[6]\setminus\{x\}}3$ is the dual of each, so each has the fusion width one and fuses back into $\binom{[6]}3$ with the fusion bunch $\{x\}$ (Proposition~\ref{prop:cousins}(i)). For any two marks the period is read through either cousin (Proposition~\ref{prop:cousins}(ii)); at $x=g$ and $y=5$, writing $s'=m_g$ and renaming $g$ as $5$ in the second cousin,
\begin{multline*}
\Gst21{s',\ 2\sigma-s'}{2\sigma}\,P_{\binom{[5]}2}\big(m_1,\dots,m_5;s',2\sigma-s'\big)\\
=\Gst21{m_5,\ 2\sigma-m_5}{2\sigma}\,P_{\binom{[5]}2}\big(m_1,m_2,m_3,m_4,s';m_5,2\sigma-m_5\big)
\end{multline*}
on $m_{[5]}+s'=6\sigma$: the basis exponent $s'$ of $\binom{[5]}2$ and the mass of one mark trade places (Proposition~\ref{prop:cousins}(iii)). At the centre, $\sigma=1$ and every mass $1$, both sides are $\Gst21{1,\ 1}2P_{\binom{[5]}2}(\one;1,1)=P_{\binom{[5]}2}(\one;1,1)$, and $P_{\binom{[6]}3}(\one;1,1)=P_{\binom{[5]}2}(\one;1,1)$.

\emph{The symmetry the fusion adds.} All $720$ permutations of $[5]\sqcup\{g\}$ carry $\binom{[6]}3$; those carrying $\binom{[5]}2$, times the permutations of $G=\{g\}$, are the $120\cdot1$ permutations fixing $g$. The index is $720/120=6=\binom{n+c}c=\binom61$ (Remark~\ref{rem:fusionsymmetry}).

\emph{(4) Two-sums and (5) repair.} Every set of two or three marks has the width $2+2-2=2$ (Definition~\ref{def:skeleton}), so $\binom{[5]}2$ is prime (Theorem~\ref{thm:primetree}) and no two-sum, and being fusable it needs no repair.

\emph{Answer.} The procedure stops at step (3): $\binom{[5]}2$ has the fusion width one, passes the tell, and fuses into $\binom{[5]\sqcup\{g\}}3$, with the letter $e_3=\Psi+\alpha_gK$, its only major. Its cousins are the six copies $\binom{[6]\setminus\{x\}}2$ of $\binom{[5]}2$, and the period of the fusion is one Beta corner times the period of any cousin, the basis exponent of $\binom{[5]}2$ and the masses of its marks trading places under the permutations of six marks; the fusion adds the index $6$ to the symmetry.
\end{solution}

\chapter{Exercises on the motive and the arithmetic of values}\label{ch:exmotive}

The exercises of this chapter work the statements of Chapter~\ref{ch:motive} on small lists. The list $\binom{[3]}2$ at its centre has a pair of width one that does not split off; read at that pair as its root bunch, its period is a node form with two Skeleton values, and it is three times the Trailmoment at the index two of the lattice input that counts $1+3\mathbb Z$ and subtracts $2+3\mathbb Z$. The list $\binom{[4]}2$ is, at the complement of each mark, a core on two marks with a dual pair substituted, and it has the motivic period of $\binom{[3]}2$; its pairs have width two and do not split off. The list $\binom{[4]}1$ has rational periods and rational motivic periods at its full leans; at its centre a pair of width one has a law with a density, and the period is the critical-line word $\operatorname{crit}_1[1]$, seven times the Trailmoment at the index three of the counting input. Substitutions of dual pairs on five and on eight marks have the motivic periods $(P^{\mathrm{mot}}_{\binom{[3]}2})^2$ and $(P^{\mathrm{mot}}_{\binom{[3]}2})^3$. At the exponents $(0,2)$ with unit masses every hinge wall and every Beta corner is $1$, and the motivic period of a list is the product of the motivic periods of its prime nodes. Every number displayed is exact, and the receipts that recompute them are the scripts named in each solution.

Throughout, a list $\cB$ \emph{at its centre} is the family $(\cB;1,1)$ at the masses $\one$ (Proposition~\ref{prop:centrefamily}), and when it is convergent $P^{\mathrm{mot}}_{\cB}$ denotes its motivic period (Definition~\ref{def:mot-mperiod}); a family at a lean is written out with its exponents and its masses.

\section{Uniform lists at their centres}\label{sec:exmot-centres}

\begin{exercise}[The list $\binom{[3]}2$ at its centre]\label{ex:mot-series}
Let $\cE=\{1,2,3\}$ and $F=\big(\binom{\cE}2;1,1\big)$ at the masses $\one$.
\begin{enumerate}[label=(\roman*)]
\item Show that $F$ is convergent (Definition~\ref{def:mot-model}), that the pair $S=\{1,2\}$ has width one with the ratio $\varrho_S=1/(u(1-u))$ on its shape $(u,1-u)$, and that $S$ does not split off (Theorem~\ref{thm:mot-pointmass}).
\item Take a tree of bunches whose root carries the bunch $S$. Find the Skeleton values of the letters along $S$, the node form of $\omega_F$ at the root (Theorem~\ref{thm:mot-node}) and its share side.
\item Evaluate the period $P_{\binom{\cE}2}(\one;1,1)$ as a Trailmoment of a lattice input.
\end{enumerate}
\end{exercise}

\begin{solution}
Receipt: \texttt{ex1\_series\_u23.py}.

(i) \emph{Letters and walls.} The bases are the three pairs, so $K=e_2(\alpha)=\alpha_1\alpha_2+\alpha_1\alpha_3+\alpha_2\alpha_3$ and $\Psi=e_1(\alpha)$ (Definition~\ref{def:table}), with the rank $2$ and the corank $1$. The imbalance is $\kappa_\cE=3-(1\cdot2+1\cdot1)=0$ (Definition~\ref{def:family}). By Theorem~\ref{thm:window} and Definition~\ref{def:skeleton}, $d(S)=\min_T|T\cap S|$ and $d^\ast(S)=|S|-r(S)$: a mark has $d=d^\ast=0$, and a pair has $d=1$ and $d^\ast=0$. So the walls $\kappa_S=m_S-d(S)-d^\ast(S)$ are
\[
\kappa_{\{e\}}=1-0-0=1,\qquad \kappa_{\{e,f\}}=2-1-0=1,
\]
the widths of these sets, as Proposition~\ref{prop:centrefamily}(ii) says at a centre. Every proper wall is at least $1$, so $F$ is convergent, its motivic period $P^{\mathrm{mot}}_{\binom{[3]}2}$ is defined (Definition~\ref{def:mot-mperiod}), and its period is
\[
P_{\binom{\cE}2}(\one;1,1)=\int_{\Delta_\cE}\frac{d\omega_\cE(u)}{e_1(u)\,e_2(u)}
\]
(Proposition~\ref{prop:mot-model}(ii), Theorem~\ref{thm:projective}).

\emph{The pair.} The bases meet $S$ in one or two marks, so $d(S)=1$, $r(S)=2$ and $c(S)=1$ (Theorem~\ref{thm:window}); every basis contains all but at most one mark of $S$, so $S$ is a series bunch (Definition~\ref{def:width}). With the room $R=\{3\}$, the restriction $\cB|S$ has the one basis $S$ and the contraction $\cB/R$ has the bases $\{1\}$ and $\{2\}$ (Definition~\ref{def:minors}). The hinge values and their ratio (Definition~\ref{def:readings}) are
\[
\gamma^K_S=\frac{K_{\cB|S}}{K_{\cB/R}}=\frac{\alpha_1\alpha_2}{\alpha_1+\alpha_2},\qquad
\gamma^\Psi_S=\frac{\Psi_{\cB/R}}{\Psi_{\cB|S}}=\frac{\alpha_1+\alpha_2}{1},\qquad
\varrho_S=\frac{\gamma^\Psi_S}{\gamma^K_S}=\frac{(\alpha_1+\alpha_2)^2}{\alpha_1\alpha_2}=\frac1{u(1-u)}
\]
on the shape $(u,1-u)$, as Proposition~\ref{prop:classreadings} gives for a series pair.

\emph{The pair does not split off.} For one list with $s,t>0$ and $|S|\ge2$, Theorem~\ref{thm:mot-pointmass} says that $S$ is pinned exactly when it has width one and $\mathbf E_{\cB|S}=\mathbf E_{\cB/R}$. Here
\[
\mathbf E_{\cB|S}=\alpha_1\alpha_2\cdot1,\qquad \mathbf E_{\cB/R}=(\alpha_1+\alpha_2)(\alpha_1+\alpha_2),
\]
which differ: their ratio $\varrho_S$ is at least $4$, with equality only at $u=\frac12$. So $S$ is not pinned; by Theorem~\ref{thm:mot-pointmass}, (ii), (iii) and (i), it does not split off and its law is carried by no point. This is Corollary~\ref{cor:mot-whatsplits}(iv) at the mark $3$, whose deletion and contraction are $\cB\setminus3=\cB|S$ and $\cB/3=\cB/R$ (Definition~\ref{def:delcon}).

(ii) \emph{The step at the root.} The list is one molecule, its circuit $\{1,2,3\}$ linking its marks (Definition~\ref{def:molecule}), so the root of a tree of bunches is a node, and it may carry the bunch $S$, of width one, its room being the one mark $3$ (Definition~\ref{def:bunchtree}). The flag $\varnothing\subsetneq\cE$ has the one step $\cE$, which meets $S$ and $R$: it is a node step of the root (Definition~\ref{def:mot-attach}(ii)), with $Z_B=S$ and $Z_R=R$. The form $\omega_F$ is regular on $V_F$ (Proposition~\ref{prop:mot-model}(ii)), and its period is read by its node form at the root (Theorem~\ref{thm:mot-node}(v)). Read $\alpha_S=\rho\,(u,1-u)$ with $\rho>0$ and $\alpha_3=1$, the one point of $\Delta_R$, and put $Q=u(1-u)$.

\emph{Letters and Skeleton values.} By Theorem~\ref{thm:mot-node}(i) the letter $K$ is the basis polynomial of $\binom\cE2$, in which $S$ has the floor $1$, and $\Psi$ is that of the dual list $\binom\cE1$ (Proposition~\ref{prop:identity}), in which $S$ has the floor $0$; $S$ has the width $1$ in both. Along the scale (Theorem~\ref{thm:mot-node}(ii)),
\[
K(\rho u,1)=\rho^2u(1-u)+\rho\,\big(u+(1-u)\big)=\rho\,Q\,\big(\rho+Q^{-1}\big),\qquad \Psi(\rho u,1)=\rho+1,
\]
so the joint spectrum of $S$ is $\{Q^{-1},1\}$: the Skeleton value of $K$ is $\xi_K=Q^{-1}=\varrho_S$ and that of $\Psi$ is $\xi_\Psi=1$. At width one each is the hinge value of the room over that of the bunch (Theorem~\ref{thm:mot-node}(iv), Corollary~\ref{cor:widthonestage}): with $\gamma^K_R=K_{\cB|R}/K_{\cB/S}=\alpha_3/1$ and $\gamma^\Psi_R=\Psi_{\cB/S}/\Psi_{\cB|R}=\alpha_3/1$ (Definition~\ref{def:readings}), on $\Delta_S\times\Delta_R$
\[
\xi_K=\frac{\gamma^K_R}{\gamma^K_S}=\frac1Q,\qquad \xi_\Psi=\frac{\gamma^\Psi_R}{\gamma^\Psi_S}=\frac11 .
\]

\emph{The node form.} On the section $\alpha_3=1$ the masses are $\one$ and $d\alpha_1\,d\alpha_2=\rho\,d\rho\,du$ (Lemma~\ref{lem:scaleshape}), so (Theorem~\ref{thm:projective})
\[
P_{\binom\cE2}(\one;1,1)=\int_0^1\!\!\int_0^\infty\frac{\rho\,d\rho}{K(\rho u,1)\,\Psi(\rho u,1)}\,du=\int_0^1Q^{-1}\int_0^\infty\big(\rho+Q^{-1}\big)^{-1}(\rho+1)^{-1}\,d\rho\,du .
\]
This is the node form of Theorem~\ref{thm:mot-node}(iii) with one term: $A_t=Q^{-1}$, $B_t=1$, $\nu_t=1$, the letter exponents $k_{t,K}=k_{t,\Psi}=1$, and $\mathbf w_t=(1,1)$ on the joint spectrum, $|\mathbf w_t|=2$. Since $0<\nu_t<|\mathbf w_t|$, $\mathrm{Sc}_t$ is the convergent scale integral, the stacked Gamma times the Breta of one bunch (Proposition~\ref{prop:scalebreta}), and at positive coefficients the Breta is the integral over the shape simplex of $J$ (Definition~\ref{def:breta}, Proposition~\ref{prop:mot-bretadomain}):
\[
\mathrm{Sc}_t=\Gst22{1,\ 1}{1,\ 1}\,\sB^{\{J\}}_{(1,1),\,1;\,I_J}\big(Q^{-1},1\big)=\int_0^1\frac{dv}{Q^{-1}v+(1-v)}=\frac{\log Q^{-1}}{Q^{-1}-1}=\frac{Q\log Q}{Q-1}.
\]
So
\[
P_{\binom\cE2}(\one;1,1)=\int_0^1A_t\,\mathrm{Sc}_t\,du=\int_0^1\frac{\log Q}{Q-1}\,du,\qquad Q=u(1-u).
\]

\emph{The share side.} Corollary~\ref{cor:sharepercoefficient}, with $\xi_0=1$, $\xi_1=Q^{-1}$, $w_0=w_1=\kappa=1$ and one share $a$ on a line $\re a=c$, $0<c<1$, writes the term as in Theorem~\ref{thm:mot-node}(iv):
\begin{align*}
A_t\,\mathrm{Sc}_t&=Q^{-1}\,\frac1{2\pi i}\int_{(c)}\Gst21{a,\ 1-a}{1}\Gst21{1-a,\ a}{1}\big(Q^{-1}\big)^{a-1}\,da\\
&=\frac1{2\pi i}\int_{(c)}\Gst40{a,\ a,\ 1-a,\ 1-a}{}\,Q^{-a}\,da .
\end{align*}
Here $\xi_K^{-b}=(\gamma^K_S)^{b}$ with $b=1-a$, and the condition of Theorem~\ref{thm:mot-node}(iv) on the bunch, $\int_0^1Q^{-1}Q^{1-c}\,du<\infty$, holds at $c=\frac12$, where it is $\int_0^1(u(1-u))^{-1/2}du=\Gst21{\frac12,\ \frac12}{1}$ (Theorem~\ref{thm:dirichlet}); the room is one point. Closing to the left (Lemma~\ref{lem:resonantpair}(ii), with $X=Q\le\frac14$), the stacked Gamma has at $a=-j$, $j\in\mathbb N$, the principal part $(a+j)^{-2}$ and no simple part, so the residue there is the derivative of $Q^{-a}$ at $a=-j$, which is $-Q^j\log Q$, and
\[
\frac1{2\pi i}\int_{(1/2)}\Gst40{a,\ a,\ 1-a,\ 1-a}{}\,Q^{-a}\,da=-\sum_{j\ge0}Q^j\log Q=\frac{\log Q}{Q-1},
\]
the value of the node form above.

(iii) \emph{The reflection of the shape.} $\log Q=\log u+\log(1-u)$, and $1-Q=1-u+u^2$ is unchanged by $u\mapsto1-u$; so the two logarithms give equal integrals and
\[
P_{\binom\cE2}(\one;1,1)=-2\int_0^1\frac{\log u}{1-u+u^2}\,du .
\]
\emph{The logarithmic coordinate.} With $u=e^{-y}$, $du=-e^{-y}dy$, and multiplying numerator and denominator by $1+e^{-y}$, since $(1-z+z^2)(1+z)=1+z^3$,
\[
P_{\binom\cE2}(\one;1,1)=2\int_0^\infty y\,\frac{e^{-y}}{1-e^{-y}+e^{-2y}}\,dy=2\int_0^\infty y\,\frac{e^{-y}+e^{-2y}}{1+e^{-3y}}\,dy .
\]
\emph{The lattice input.} Let $\phi=\Ilat{1+6\mathbb Z}+\Ilat{2+6\mathbb Z}-\Ilat{4+6\mathbb Z}-\Ilat{5+6\mathbb Z}$, an input with the density $\pm1$ on four lattice inputs (Definitions~\ref{def:lattice} and~\ref{def:inputs}). It is tame beyond the cut $0$, and its Laplace transform there (Definition~\ref{def:inputs}) sums its four cosets as geometric series: with $z=e^{-y}$,
\[
\mathcal F_0(y)=\sum_{k\ge0}\big(z^{6k+1}+z^{6k+2}-z^{6k+4}-z^{6k+5}\big)=\frac{z+z^2-z^4-z^5}{1-z^6}=\frac{(z+z^2)(1-z^3)}{(1-z^3)(1+z^3)}=\frac{z+z^2}{1+z^3}.
\]
Since $\int_{(0,\infty)}x^{-2}\,|\phi|(dx)<\infty$, the Laplace side (Theorem~\ref{thm:laplaceside}) at one unit exponent and the index $2$ reads $\Gamma(2)\,\Tx{2}{\phi}{0}=\int_0^\infty y\,\mathcal F_0(y)\,dy$, and
\[
P_{\binom\cE2}(\one;1,1)=2\,\Tx{2}{\big(\Ilat{1+6\mathbb Z}+\Ilat{2+6\mathbb Z}-\Ilat{4+6\mathbb Z}-\Ilat{5+6\mathbb Z}\big)}{0}.
\]
\emph{Cosets and dilations.} $1+3\mathbb Z=(1+6\mathbb Z)\sqcup(4+6\mathbb Z)$ and $2+3\mathbb Z=(2+6\mathbb Z)\sqcup(5+6\mathbb Z)$, so
\[
\phi=\Ilat{1+3\mathbb Z}-\Ilat{2+3\mathbb Z}+2\big(\Ilat{2+6\mathbb Z}-\Ilat{4+6\mathbb Z}\big),
\]
and the Trailmoment, an integral against its input (Definition~\ref{def:jointtrail}), splits accordingly, every part converging. The dilation $\tau(x)=2x$ carries $\Ilat{1+3\mathbb Z}$ to $\Ilat{2+6\mathbb Z}$ and $\Ilat{2+3\mathbb Z}$ to $\Ilat{4+6\mathbb Z}$, so transport at one unit exponent (Theorem~\ref{thm:transportfree} and the paragraph after it) gives
\[
\Tx{2}{\big(\Ilat{2+6\mathbb Z}-\Ilat{4+6\mathbb Z}\big)}{0}=2^{-2}\,\Tx{2}{\big(\Ilat{1+3\mathbb Z}-\Ilat{2+3\mathbb Z}\big)}{0},
\]
and
\[
P_{\binom\cE2}(\one;1,1)=2\,\big(1+2\cdot2^{-2}\big)\,\Tx{2}{\big(\Ilat{1+3\mathbb Z}-\Ilat{2+3\mathbb Z}\big)}{0}=3\,\Tx{2}{\big(\Ilat{1+3\mathbb Z}-\Ilat{2+3\mathbb Z}\big)}{0},
\]
which is $2.343907238689458890\ldots$ to eighteen places. The receipt checks the identities of rational functions exactly, the decomposition of the input at every $m\le600$, the decomposition and the dilation together on the exact heads $\sum_{m\le6N}\phi(m)m^{-2}$, and the integral of the node form against the Trailmoment to twenty-five places.

\emph{Answer.} $F$ is convergent, with every proper wall $1$. The pair $S=\{1,2\}$ is a series bunch of width one with $\gamma^K_S=\alpha_1\alpha_2/(\alpha_1+\alpha_2)$, $\gamma^\Psi_S=\alpha_1+\alpha_2$ and $\varrho_S=1/(u(1-u))$; it is not pinned, so it does not split off. At the root bunch $S$ the Skeleton values are $Q^{-1}$ and $1$, $Q=u(1-u)$; the node form is $\int_0^1Q^{-1}\Gst22{1,\ 1}{1,\ 1}\sB^{\{J\}}_{(1,1),1;I_J}(Q^{-1},1)\,du=\int_0^1\frac{\log Q}{Q-1}\,du$, and its share side reads the stacked Gamma $\Gst40{a,\ a,\ 1-a,\ 1-a}{}$ against $Q^{-a}$. The period is
\[
P_{\binom{[3]}2}(\one;1,1)=3\,\Tx{2}{\big(\Ilat{1+3\mathbb Z}-\Ilat{2+3\mathbb Z}\big)}{0}.
\]
\end{solution}

\begin{exercise}[The list $\binom{[4]}2$ at its centre]\label{ex:mot-dualpair}
Let $\cE=\{1,2,3,4\}$ and $F=\big(\binom{\cE}2;1,1\big)$ at the masses $\one$.
\begin{enumerate}[label=(\roman*)]
\item Find the widths and the walls of the sets of marks. Show that the sets of width one with at least two marks are the complements $R=\cE\setminus\{e\}$ of single marks, that the cell of each is a dual pair (Definition~\ref{def:substitution}), and that $\binom{\cE}2=\binom{\{e,h\}}1\langle e,\mathcal A_R\rangle$ with the mark $e$ left in place.
\item Show that $P^{\mathrm{mot}}_{\binom{[4]}2}=P^{\mathrm{mot}}_{\binom{[3]}2}$, and find the period.
\item Show that no pair of marks splits off.
\end{enumerate}
\end{exercise}

\begin{solution}
Receipt: \texttt{ex2\_dualpair\_u24.py}.

(i) \emph{Letters, widths, walls.} The complement of a pair is a pair, so $K=\Psi=e_2(\alpha)$ (Definition~\ref{def:table}). A set $S$ has $r(S)=\min(|S|,2)$, and by Definition~\ref{def:skeleton} its width is $c(S)=r(S)+r(\cE\setminus S)-2$:
\[
c(\{e\})=1+2-2=1,\qquad c(\{e,f\})=2+2-2=2,\qquad c(\cE\setminus\{e\})=2+1-2=1 .
\]
At the centre the walls are the widths (Proposition~\ref{prop:centrefamily}(ii)): $1$, $2$ and $1$. Every proper wall is at least $1$, and $F$ is convergent (Definition~\ref{def:mot-model}).

\emph{The cells.} The sets of width one with at least two marks are the triples $R=\cE\setminus\{e\}$. The restriction $\cB|R$ has the three pairs of $R$ and the contraction $\cB/\{e\}$ the three marks of $R$, the bases through $e$ less $e$ (Definition~\ref{def:minors}). The cell of the bunch $R$ (Theorem~\ref{thm:widthone}) has the bases of $\cB|R=\binom R2$ and the sets $X+g$, $X\in\cB/\{e\}=\binom R1$:
\[
\mathcal A_R=\cB^+_R=\binom{R+g}2,\qquad \mathcal A_R\setminus g=\binom R2,\qquad \mathcal A_R/g=\binom R1 .
\]
The complements in $R$ of the pairs of $R$ are its marks, so $\mathcal A_R/g$ is the dual table of $\mathcal A_R\setminus g$, the weights being $1$: the cell is a dual pair (Definition~\ref{def:substitution}). The cell of the room $\{e\}$ has the basis $\{e\}$ of $\cB|\{e\}$ and the set $\varnothing+g$, $\varnothing$ being the one basis of $\cB/R$: it is $\binom{\{e,g\}}1$. By Proposition~\ref{prop:substitution}(iii) the list is the gluing of the two cells, that is $\binom{\{e,h\}}1\langle e,\mathcal A_R\rangle$ with $\mathcal A_R$ at $h$ (Definition~\ref{def:substitution}). On the bases: the core basis $\{e\}$ avoids $h$ and takes $X\in\mathcal A_R/g$, giving the three pairs through $e$; the core basis $\{h\}$ takes $X\in\mathcal A_R\setminus g$, giving the three pairs of $R$.

\emph{The letters of the substitution.} For $e=4$ and $R=\{1,2,3\}$ the cell has $k=K_{\mathcal A_R\setminus g}=e_2(\alpha_R)$ and $\psi=\Psi_{\mathcal A_R\setminus g}=e_1(\alpha_R)$, and Proposition~\ref{prop:mot-dualletters}(ii) gives the one hinge value $\gamma:=\gamma^K_{\mathcal A_R}=\gamma^\Psi_{\mathcal A_R}=e_2(\alpha_R)/e_1(\alpha_R)$, with $\varrho_{\mathcal A_R}=1$. The mark $4$ left in place has $k=\alpha_4$ and $\psi=1$, so $\pi=e_1(\alpha_R)$, and Proposition~\ref{prop:mot-dualletters}(iii), with $r_{\mathcal H}=1$ and $q=2$, reads
\[
K_{\mathcal H}\big(\pi\alpha_4,\ e_2(\alpha_R)\big)=e_1(\alpha_R)\,\alpha_4+e_2(\alpha_R)=e_2(\alpha)=\pi^{0}K,\qquad \Psi_{\mathcal H}\big(\pi\alpha_4,\ e_2(\alpha_R)\big)=e_2(\alpha)=\pi^{0}\Psi,
\]
$K_{\mathcal H}$ and $\Psi_{\mathcal H}$ both being $x_4+x_h$.

(ii) \emph{The dual-pair factorization.} Corollary~\ref{cor:mot-dualfull}(a) applies with $q=1$ at $s=t=1$. The wall of the cell is
\[
\kappa_R=m_R-(s+t)\big(r(R)-1\big)=3-2\,(2-1)=1,
\]
the wall of $F$ at $R$ found above. The pinned family has the form of $(\cB|R;\kappa_R,s+t-\kappa_R)=\big(\binom R2;1,1\big)$ at the masses $\one_R$, and the core is $\big(\binom{\{e,h\}}1;1,1\big)$ at the masses $m_e=1$ and $\kappa_R=1$ at $h$:
\[
P^{\mathrm{mot}}_{\binom{[4]}2}=P^{\mathrm{mot}}_{(\binom{\{e,h\}}1;1,1)}\cdot P^{\mathrm{mot}}_{(\binom R2;1,1)} .
\]
The coefficients of the forms show the factorization: with $\gamma$ as above,
\[
\frac{\alpha^{\one}}{K\,\Psi}=\frac{\alpha_4\,\gamma}{(\alpha_4+\gamma)^2}\cdot\frac{\alpha^{\one_R}}{e_1(\alpha_R)\,e_2(\alpha_R)},
\]
the coefficient of the core at $(\alpha_4,\gamma)$ times that of $\big(\binom R2;1,1\big)$, an identity of rational functions.

\emph{The core.} The core is a family on two marks with both letters $\alpha_e+\alpha_h$ and the walls $\kappa_{\{e\}}=\kappa_{\{h\}}=1$; its form is $\Omega/(\alpha_e+\alpha_h)^2$, and by Lemma~\ref{lem:motive-twomarks} its motivic period is its Beta corner,
\[
P^{\mathrm{mot}}_{(\binom{\{e,h\}}1;1,1)}=\Gst21{1,\ 1}{2}=\frac{\Gamma(1)\,\Gamma(1)}{\Gamma(2)}=1 .
\]
The pinned family $\big(\binom R2;1,1\big)$ at $\one_R$ is $\binom{[3]}2$ at its centre with its marks named by $R$; renaming the coordinates is an isomorphism of the models that maps the positive cells and the forms onto each other, so the two have one motivic period (Lemma~\ref{lem:mot-birational}). Hence
\[
P^{\mathrm{mot}}_{\binom{[4]}2}=P^{\mathrm{mot}}_{\binom{[3]}2},\qquad \mathrm C\big(P^{\mathrm{mot}}_{\binom{[4]}2}\big)=\mathrm C\big(P^{\mathrm{mot}}_{\binom{[3]}2}\big),
\]
and by Exercise~\ref{ex:mot-series} the period is $P_{\binom{[4]}2}(\one;1,1)=3\,\Tx{2}{\big(\Ilat{1+3\mathbb Z}-\Ilat{2+3\mathbb Z}\big)}{0}$.

(iii) \emph{The pairs.} A pair has width two in the one active table, so by Corollary~\ref{cor:mot-whatsplits}(i) it does not split off, and its law is carried by no point. The list is connected and its only sets with two marks on each side are the pairs, so it is prime (Theorem~\ref{thm:primetree}), and Corollary~\ref{cor:mot-whatsplits}(ii) says the same. The complements of the marks do split off, as in (ii): $\mathbf E_{\cB\setminus e}=e_2(\alpha_R)\,e_1(\alpha_R)=\mathbf E_{\cB/e}$ (Corollary~\ref{cor:mot-whatsplits}(iv)).

\emph{Answer.} The marks and the triples have width one and the pairs width two, and at the centre these are the walls. Every triple $R=\cE\setminus\{e\}$ has the cell $\binom{R+g}2$, a dual pair, and $\binom{\cE}2=\binom{\{e,h\}}1\langle e,\binom{R+g}2\rangle$. With $\kappa_R=1$, the core on two marks has the motivic period $\Gst21{1,\ 1}{2}=1$, so
\[
P^{\mathrm{mot}}_{\binom{[4]}2}=P^{\mathrm{mot}}_{\binom{[3]}2},\qquad P_{\binom{[4]}2}(\one;1,1)=3\,\Tx{2}{\big(\Ilat{1+3\mathbb Z}-\Ilat{2+3\mathbb Z}\big)}{0}.
\]
No pair splits off: pairs have width two.
\end{solution}

\begin{exercise}[The list $\binom{[4]}1$ at its full leans and at its centre]\label{ex:mot-critical}
Let $\cE=\{1,2,3,4\}$ and $\cB=\binom{\cE}1$.
\begin{enumerate}[label=(\roman*)]
\item Show that at its full leans $(s,0)$ and $(0,t)$, at positive integer masses on the balance wall with every proper wall at least $1$, the periods are rational, and so are the motivic periods.
\item At the centre, show that $S=\{3,4\}$ has width one, find its hinge values and its law (Definition~\ref{def:seamlaw}), and show that the law is not a point mass.
\item Read the period at the root bunch $S$, write it on the share side of the node form, show that it is the critical-line word $\operatorname{crit}_1[1]$ (Definition~\ref{def:critical}), and evaluate it.
\end{enumerate}
\end{exercise}

\begin{solution}
Receipt: \texttt{ex3\_critical\_u14.py}.

(i) \emph{Counts.} The bases are the four marks, so $K=e_1(\alpha)$ and $\Psi=e_3(\alpha)$ (Definition~\ref{def:table}), of the rank $1$ and the corank $3$. A proper set $S$ has $d(S)=0$, a basis avoiding it, $r(S)=1$, $d^\ast(S)=|S|-1$ and $c(S)=1$ (Theorem~\ref{thm:window}, Definition~\ref{def:skeleton}).

\emph{The full lean $(s,0)$.} The balance wall is $m_\cE=s$, and the walls are $\kappa_S=m_S-s\,d(S)=m_S$ (Definition~\ref{def:family}), at least $1$ at positive integer masses. On the shape simplex $K(u)=|u|=1$, so (Theorems~\ref{thm:projective} and~\ref{thm:dirichlet})
\[
P_{\cB}(\mvec;s,0)=\int_{\Delta_\cE}u^{\mvec-\one}\,d\omega_\cE(u)=\Gst41{m_1,\ m_2,\ m_3,\ m_4}{s},
\]
a positive rational number; at $(4,0)$ and the masses $\one$ it is $\Gst41{1,\ 1,\ 1,\ 1}{4}=\frac16$.

\emph{The full lean $(0,t)$.} The balance wall is $m_\cE=3t$. By Corollary~\ref{cor:mot-cremona}, with $2\sigma=t$, the family $(\cB;0,t)$ at $\mvec$ has the motivic period of $(\cB;t,0)$ at $t\one-\mvec$, whose total is $4t-3t=t$ and whose walls are those of $(\cB;0,t)$ at the complements, so
\[
P_{\cB}(\mvec;0,t)=\Gst41{t-m_1,\ t-m_2,\ t-m_3,\ t-m_4}{t},
\]
rational again; the walls of $(\cB;0,t)$ at the complements of the marks are $t-m_e\ge1$.

\emph{The motivic periods.} The set $\{3,4\}$ has width one with two marks on each side, and by Proposition~\ref{prop:substitution}(iii) the list is the gluing of its two cells: the cell of the bunch has the bases $\{3\},\{4\}$ of $\cB|\{3,4\}$ and $g$, the one basis of $\cB/\{1,2\}$ being $\varnothing$, and the cell of the room has the bases $\{1\},\{2\},\{g\}$. Both are the parallel nodes $\binom{\{3,4,g\}}1$ and $\binom{\{1,2,g\}}1$ on three marks, so the splitting of Corollary~\ref{cor:mot-dualfull}(b) has no prime node, and at either full lean $P^{\mathrm{mot}}_F$ is a positive rational number $c$. Its period is $c$, so $c=P_F$: the motivic periods are the rational numbers displayed.

(ii) \emph{The centre.} At $(1,1)$ and the masses $\one$ every proper wall is the width $1$ (Proposition~\ref{prop:centrefamily}(ii)), and $F=(\cB;1,1)$ is convergent. The pair $S=\{3,4\}$, with the room $S'=\{1,2\}$, has width one, and no basis contains both of its marks: it is a parallel bunch (Definition~\ref{def:width}). The restriction $\cB|S$ has the bases $\{3\},\{4\}$, and the contraction $\cB/S'$, the bases meeting $S'$ in one mark less that mark, has the one basis $\varnothing$ (Definition~\ref{def:minors}), whose complement in $S$ is $S$. So (Definition~\ref{def:readings}, Proposition~\ref{prop:classreadings})
\[
\gamma^K_S=\frac{\alpha_3+\alpha_4}{1},\qquad \gamma^\Psi_S=\frac{\alpha_3\alpha_4}{\alpha_3+\alpha_4},\qquad \varrho_S=\frac{\alpha_3\alpha_4}{(\alpha_3+\alpha_4)^2}=u(1-u)
\]
on the shape $(u,1-u)$ of $S$, and likewise $\varrho_{S'}=u'(1-u')$ on the shape $(u',1-u')$ of $S'$.

\emph{The law.} At width one the stations of the room are $D_0$ and $D_1$, of the excesses $0$ and $1$, and the seam letters are the letters of $\cB/S'$ and of $\cB|S$ (Proposition~\ref{prop:ladder}(ii)):
\begin{gather*}
V_{D_0}=K_{\cB/S'}=1,\qquad V_{D_1}=K_{\cB|S}=\alpha_3+\alpha_4,\\
V^\ast_{D_0}=\Psi_{\cB/S'}=\alpha_3\alpha_4,\qquad V^\ast_{D_1}=\Psi_{\cB|S}=\alpha_3+\alpha_4 .
\end{gather*}
With $d(S)=0$ and $d^\ast(S)=1$ the wall of the piece is $\kappa_S=2-(1\cdot0+1\cdot1)=1$ and its top wall is $(1+1)\cdot1-1=1$ (Definition~\ref{def:seampiece}). Definition~\ref{def:seamlaw} gives the scale, the scale form on the shape $(u,1-u)$, and the seam vector $(\varrho^\rho_{D_0},\varrho^\rho_{D_1},\varrho^{\rho\ast}_{D_0},\varrho^{\rho\ast}_{D_1})$:
\begin{gather*}
\gamma^\rho_S=\Big(\frac{\alpha_3+\alpha_4}{1}\Big)^{1/2}\Big(\frac{\alpha_3\alpha_4}{\alpha_3+\alpha_4}\Big)^{1/2}=(\alpha_3\alpha_4)^{1/2},\\
\phi^\rho_S(du)=1^{-1}\,\big(u+(1-u)\big)^{-1}\big(u(1-u)\big)^{-1/2}\,du=\varrho_S^{-1/2}\,du,\\
\boldsymbol\varrho^\rho_S=\Big(1,\ \frac{\alpha_3+\alpha_4}{(\alpha_3\alpha_4)^{1/2}},\ \frac{\alpha_3\alpha_4}{(\alpha_3+\alpha_4)(\alpha_3\alpha_4)^{1/2}},\ 1\Big)=\big(1,\ \varrho_S^{-1/2},\ \varrho_S^{1/2},\ 1\big).
\end{gather*}
The seam torus is cut out by $x_{D_0}=1$, $x^\ast_{D_1}=1$ and $x_{D_1}x^\ast_{D_0}=1$: it is the curve of the points $(1,x,x^{-1},1)$, $x>0$. The law is the image of $\varrho_S^{-1/2}du$ under $x=\varrho_S^{-1/2}=(u(1-u))^{-1/2}$. Since $u(1-u)\le\frac14$ it is carried by $x\ge2$; each $x>2$ has the two preimages $u$ and $1-u$, with $(1-2u)\,du=-2x^{-3}dx$ and $|1-2u|=(1-4x^{-2})^{1/2}$, so the law is
\[
\frac{4\,dx}{x\,(x^2-4)^{1/2}}\quad\text{on }x>2,\qquad\text{of total mass }\int_0^1\big(u(1-u)\big)^{-1/2}du=\Gst21{\frac12,\ \frac12}{1}
\]
(Theorem~\ref{thm:dirichlet}). It has a density, so it is not a point mass. Theorem~\ref{thm:mot-pointmass} says the same through the hinge values: $\mathbf E_{\cB|S}=(\alpha_3+\alpha_4)^2$ and $\mathbf E_{\cB/S'}=\alpha_3\alpha_4$ differ, $\varrho_S$ not being constant, so $S$ is not pinned and does not split off.

(iii) \emph{The node form at the root bunch $S$.} As in Exercise~\ref{ex:mot-series}(ii), a tree of bunches whose root carries $S$ has the step $\cE$ as a node step of the root (Definition~\ref{def:mot-attach}(ii)), and the period is read by the node form there (Theorem~\ref{thm:mot-node}(v)). Read $\alpha_S=\rho\,(u,1-u)$ and $\alpha_{S'}=(u',1-u')$ on the section $\alpha_1+\alpha_2=1$; then $d\alpha_3\,d\alpha_4=\rho\,d\rho\,du$ (Lemma~\ref{lem:scaleshape}) and
\[
K(\rho u,u')=\rho+1,\qquad \Psi(\rho u,u')=\rho\,u'(1-u')+\rho^2u(1-u)=\rho\,\varrho_S\Big(\rho+\frac{\varrho_{S'}}{\varrho_S}\Big):
\]
the Skeleton values (Theorem~\ref{thm:mot-node}(ii)) are $\xi_K=1$ and $\xi_\Psi=\varrho_{S'}/\varrho_S=\gamma^\Psi_{S'}/\gamma^\Psi_S$, the hinge value of the room over that of the bunch on $\Delta_S\times\Delta_{S'}$ (Theorem~\ref{thm:mot-node}(iv)). By Theorem~\ref{thm:projective},
\[
P_{\cB}(\one;1,1)=\int_{\Delta_S\times\Delta_{S'}}\int_0^\infty\frac{d\rho}{(\rho+1)\,(\varrho_S\,\rho+\varrho_{S'})}\,d\omega_S(u)\,d\omega_{S'}(u'),
\]
the node form with one term: $A_t=\varrho_S^{-1}$, $B_t=1$, $\nu_t=1$ and $\mathbf w_t=(1,1)$, so that $\mathrm{Sc}_t=\Gst22{1,\ 1}{1,\ 1}\sB^{\{J\}}_{(1,1),\,1;\,I_J}\big(1,\varrho_{S'}/\varrho_S\big)$: the bunch and the room meet through the one Skeleton value $\varrho_{S'}/\varrho_S$.

\emph{The share side.} Corollary~\ref{cor:sharepercoefficient}, with $\xi_0=1$, $\xi_1=\varrho_{S'}/\varrho_S$ and one share $a$, writes the term as
\begin{align*}
A_t\,\mathrm{Sc}_t&=\varrho_S^{-1}\,\frac1{2\pi i}\int_{(c)}\Gst40{a,\ a,\ 1-a,\ 1-a}{}\Big(\frac{\varrho_{S'}}{\varrho_S}\Big)^{a-1}da\\
&=\frac1{2\pi i}\int_{(c)}\Gst40{a,\ a,\ 1-a,\ 1-a}{}\,\varrho_S^{-a}\,\varrho_{S'}^{a-1}\,da .
\end{align*}
On the line $\re a=\frac12$ the two conditions of Theorem~\ref{thm:mot-node}(iv) are $\int_{\Delta_S}\varrho_S^{-1/2}d\omega_S<\infty$ and $\int_{\Delta_{S'}}\varrho_{S'}^{-1/2}d\omega_{S'}<\infty$, both $\Gst21{\frac12,\ \frac12}{1}$; so the bunch and the room separate under the share integral, and by Theorem~\ref{thm:dirichlet}, for $0<\re a<1$,
\[
\int_0^1\big(u(1-u)\big)^{-a}du=\Gst21{1-a,\ 1-a}{2-2a},\qquad \int_0^1\big(u'(1-u')\big)^{a-1}du'=\Gst21{a,\ a}{2a} .
\]
So
\[
P_{\cB}(\one;1,1)=\frac1{2\pi i}\int_{(1/2)}\Gst40{a,\ a,\ 1-a,\ 1-a}{}\,\Gst21{1-a,\ 1-a}{2-2a}\,\Gst21{a,\ a}{2a}\,da :
\]
the moments of the orders $-a$ and $a-1$ of the two ratios of hinge values, paired over one share. This is the one-share display of Theorem~\ref{thm:twosides} at the level shares $a^d=b^d=a$ and $a^r=b^r=1-a$, the bunch family and the room family having these two Beta corners as periods: the even share $\sigma^d=a$ is free and the odd share is $0$.

\emph{The critical-line word.} By the reflection at $a$ and at $2a$ and the shift $\Gamma(2-2a)=(1-2a)\,\Gamma(1-2a)$ (Proposition~\ref{prop:stacked}(iii),(v)),
\[
\Gst21{a,\ a}{2a}\Gst21{1-a,\ 1-a}{2-2a}=\frac{\Gst40{a,\ a,\ 1-a,\ 1-a}{}}{(1-2a)\,\Gst20{2a,\ 1-2a}{}}=\frac{\pi^2}{\sin^2\pi a}\cdot\frac{\sin2\pi a}{\pi\,(1-2a)}=\frac{2\pi\cot\pi a}{1-2a},
\]
and $\partial_a\Gst40{a,\ a,\ 1-a,\ 1-a}{}=\partial_a\big(\pi^2/\sin^2\pi a\big)=-2\pi\cot\pi a\cdot\Gst40{a,\ a,\ 1-a,\ 1-a}{}$. So
\[
P_{\cB}(\one;1,1)=-\frac1{2\pi i}\int_{(1/2)}\frac{\partial_a\Gst40{a,\ a,\ 1-a,\ 1-a}{}}{1-2a}\,da=\operatorname{crit}_1[1],
\]
Definition~\ref{def:critical} at $k=1$ and $R=1$, where $-\frac2{(2k)!}=-1$.

\emph{Closing.} The integrand is holomorphic on $0<\re a<1$, regular at $a=\frac12$ (Definition~\ref{def:critical}), and on every vertical line of the strip it is $O(e^{-2\pi|\im a|})$, since $|\sin\pi a|^2=\sin^2\pi\re a+\sinh^2\pi\im a$; so the line moves to $\re a=c$ with $\frac12<c<1$, and an integration by parts, with no boundary terms, gives
\[
P_{\cB}(\one;1,1)=\frac1{2\pi i}\int_{(c)}\Gst40{a,\ a,\ 1-a,\ 1-a}{}\,g(a)\,da,\qquad g(a)=\partial_a\frac1{1-2a}=\frac{2}{(1-2a)^2}.
\]
The function $g$ is rational, holomorphic on $\re a\ge c$, and $O(|a|^{-2})$, so the right closing of Lemma~\ref{lem:resonantpair}(i) applies: with $g'(a)=8(1-2a)^{-3}$,
\[
P_{\cB}(\one;1,1)=-\sum_{m\ge1}g'(m)=\sum_{m\ge1}\frac{8}{(2m-1)^3}=\sum_{m\ge1}\Big(m-\frac12\Big)^{-3}.
\]
(Closed directly, the integrand $-\partial_a\Gst40{a,\ a,\ 1-a,\ 1-a}{}/(1-2a)$ has at $a=m\ge1$ a pole of order three with the residue $8(1-2m)^{-3}$; the receipt checks $m\le6$.)

\emph{The Trailmoment.} At one unit exponent and the coefficient function $x$ the last sum is $\Tx{3}{\Ilat{\frac12+\mathbb Z}}{0}$ (Definition~\ref{def:lattice}). The lattice relation at the step $2$ (Theorem~\ref{thm:latticefree}(i), $q=2$, $b=0$, the cut $0$) assembles $\frac12\mathbb Z$ from the cosets $\mathbb Z$ and $\frac12+\mathbb Z$, the coefficient function $x/2$ multiplying the kernel by $2^3$:
\[
\Tx{3}{\IZ}{0}+\Tx{3}{\Ilat{\frac12+\mathbb Z}}{0}=2^3\,\Tx{3}{\IZ}{0}.
\]
So $P_{\cB}(\one;1,1)=\Tx{3}{\Ilat{\frac12+\mathbb Z}}{0}=7\,\Tx{3}{\IZ}{0}$, which is $\operatorname{crit}_1[1]$ of Theorem~\ref{thm:positiveslot}; it is $8.414398322117159997\ldots$ to eighteen places. The receipt checks the relation on the exact heads $\sum_{m\le N}(m-\frac12)^{-3}+\sum_{m\le N}m^{-3}=8\sum_{k\le2N}k^{-3}$ and the share integral against $7\,\Tx{3}{\IZ}{0}$ to twenty-five places. At a full lean one letter is active with one Skeleton value, no share is free, and the period is rational, by (i); at the centre the one free share runs over the critical line and carries a Trailmoment of the odd index three.

\emph{Answer.} At $(s,0)$ the period is $\Gst41{\mvec}{s}$ and at $(0,t)$ it is $\Gst41{t\one-\mvec}{t}$, rational, and the motivic periods are these rational numbers, the list having no prime node. At the centre $S=\{3,4\}$ is a parallel bunch of width one with $\gamma^K_S=\alpha_3+\alpha_4$, $\gamma^\Psi_S=\alpha_3\alpha_4/(\alpha_3+\alpha_4)$ and $\varrho_S=u(1-u)$; its law is $4\,dx/(x(x^2-4)^{1/2})$ on the points $(1,x,x^{-1},1)$, $x>2$, of its seam torus, not a point mass. Read at the root bunch $S$,
\begin{align*}
P_{\binom{[4]}1}(\one;1,1)&=\frac1{2\pi i}\int_{(1/2)}\Gst40{a,\ a,\ 1-a,\ 1-a}{}\Gst21{1-a,\ 1-a}{2-2a}\Gst21{a,\ a}{2a}\,da\\
&=\operatorname{crit}_1[1]=7\,\Tx{3}{\IZ}{0}.
\end{align*}
\end{solution}

\section{Substitutions of dual pairs}\label{sec:exmot-subst}

\begin{exercise}[Substitutions on five and on eight marks]\label{ex:mot-substitution}
Let $\mathrm P_5=\binom{\{1,2,h\}}1\langle1,2,\mathcal A\rangle$, the marks $1,2$ left in place and the cell $\mathcal A=\binom{\{3,4,5,g\}}2$ at $h$, and let $\mathrm M_8=\binom{\{b_1,b_2,h_1,h_2\}}2\langle b_1,b_2,\mathcal A_1,\mathcal A_2\rangle$, the marks $b_1,b_2$ left in place and the cells $\mathcal A_i=\binom{\{x_i,y_i,z_i,g_i\}}2$ at $h_i$, with $R=\{3,4,5\}$, $R_i=\{x_i,y_i,z_i\}$ and $B=\{b_1,b_2\}$ (Definition~\ref{def:substitution}).
\begin{enumerate}[label=(\roman*)]
\item Find the bases of $\mathrm P_5$, show that its cell is a dual pair and that $\mathrm P_5$ is convergent at its centre, and find its motivic period there.
\item Find the rank of $\mathrm M_8$, the widths of $R_1$, $R_2$ and $B$, and its letters; show that it is convergent at its centre and find its motivic period there. What becomes of the bunch $B$?
\end{enumerate}
\end{exercise}

\begin{solution}
Receipt: \texttt{ex4\_substitution.py}, with the helper \texttt{subst.py}.

(i) \emph{The bases.} By Definition~\ref{def:substitution}, the core bases $\{1\}$ and $\{2\}$ avoid $h$ and take a basis of $\mathcal A/g=\binom R1$, and the core basis $\{h\}$ takes a basis of $\mathcal A\setminus g=\binom R2$: the bases of $\mathrm P_5$ are the six pairs $\{1,x\},\{2,x\}$, $x\in R$, and the three pairs of $R$, all the pairs of $\{1,\dots,5\}$ but $\{1,2\}$, and its rank is $1+(2-1)=2$ (Proposition~\ref{prop:substitution}(i)). The complements in $R$ of the pairs of $R$ are its marks, so $\mathcal A/g$ is the dual of $\mathcal A\setminus g$: the cell is a dual pair, and by Proposition~\ref{prop:mot-dualletters}(ii) both letters read it through the one hinge value $\gamma=e_2(\alpha_R)/e_1(\alpha_R)$. By Theorem~\ref{thm:substletters}, with $K_{\mathcal A/g}=\Psi_{\mathcal A\setminus g}=e_1(\alpha_R)$,
\[
K=e_1(\alpha_R)\,\big(\alpha_1+\alpha_2+\gamma\big),\qquad \Psi=e_1(\alpha_R)\,\big(\alpha_1\alpha_2+(\alpha_1+\alpha_2)\,\gamma\big).
\]

\emph{Walls.} $R$ is a bunch of width one with the cell $\mathcal A$ (Proposition~\ref{prop:substitution}(i)); its wall at the centre is $\kappa_R=m_R-(s+t)(r(R)-1)=3-2\cdot1=1$, $r(R)=2$ (Corollary~\ref{cor:mot-dualfull}(a)), and $2\sigma-\kappa_R=1$. Both letters read $R$ through $\gamma$, so by Theorem~\ref{thm:wallsglue}(iii) $\mathrm P_5$ has every proper wall positive exactly when its two cell families do: $(\mathcal A;1,1)$ at the masses $(\one_R,1)$, which is $\binom{[4]}2$ at its centre, and $\big(\binom{\{1,2,h\}}1;1,1\big)$ at the masses $(1,1,\kappa_R)=\one$, which is $\binom{[3]}1$ at its centre; both are connected, and their walls are their widths (Proposition~\ref{prop:centrefamily}(ii)). So every proper wall of $\mathrm P_5$ is positive, an integer and so at least $1$ (the receipt finds the values $1$ and $2$), and $\mathrm P_5$ is convergent at its centre.

\emph{The motivic period.} Corollary~\ref{cor:mot-dualfull}(a) with $q=1$:
\[
P^{\mathrm{mot}}_{\mathrm P_5}=P^{\mathrm{mot}}_{(\binom{\{1,2,h\}}1;1,1)}\cdot P^{\mathrm{mot}}_{(\binom R2;1,1)},
\]
the core at the masses $\one$ and the pinned family $(\cB|R;\kappa_R,2-\kappa_R)=\big(\binom R2;1,1\big)$ at $\one_R$. The core has the letters $e_1$ and $e_2$ of $\alpha_1,\alpha_2,\alpha_h$, so its form is $\Omega/(e_1e_2)$, the form of $\binom{[3]}2$ at its centre, whose letters are $e_2$ and $e_1$; the pinned family is $\binom{[3]}2$ at its centre on the marks of $R$. Two convergent families with one form have one motivic period (Lemma~\ref{lem:mot-birational}, with the renaming of the coordinates), so
\[
P^{\mathrm{mot}}_{\mathrm P_5}=\big(P^{\mathrm{mot}}_{\binom{[3]}2}\big)^2,
\]
the two factors being one motivic period. In coefficients,
\[
\frac{\alpha^{\one}}{K\,\Psi}=\frac{\alpha_1\alpha_2\,\gamma}{e_1(\alpha_1,\alpha_2,\gamma)\,e_2(\alpha_1,\alpha_2,\gamma)}\cdot\frac{\alpha^{\one_R}}{e_1(\alpha_R)\,e_2(\alpha_R)},
\]
the core at $(\alpha_1,\alpha_2,\gamma)$ times $\binom{[3]}2$ on $R$.

(ii) \emph{Rank and widths.} By Proposition~\ref{prop:substitution}(i) $\mathrm M_8$ is a list of rank $2+(2-1)+(2-1)=4$ on eight marks, and each $R_i$ is a bunch of width one with the cell $\mathcal A_i$; from the six core bases, with three choices at each cell, it has $6\cdot3\cdot3=54$ bases. The pairs of the core meet $B$ in $0$, $1$ or $2$ marks, so $B$ has the width two in the core, and by Corollary~\ref{cor:mot-widthbound} in $\mathrm M_8$.

\emph{Letters.} Both cells are dual pairs with $k_i=e_2(\alpha_{R_i})$, $\psi_i=e_1(\alpha_{R_i})$ and the hinge values $\gamma_i=e_2(\alpha_{R_i})/e_1(\alpha_{R_i})$, and the marks left in place have $k=\alpha_b$, $\psi=1$ (Proposition~\ref{prop:mot-dualletters}(ii),(iii)). With $\pi=e_1(\alpha_{R_1})\,e_1(\alpha_{R_2})$, $r_{\mathcal H}=2$ and $q=4$, Proposition~\ref{prop:mot-dualletters}(iii) gives, $K_{\mathcal H}=\Psi_{\mathcal H}=e_2$ being homogeneous of degree two,
\[
\pi^2\,e_2(\alpha_{b_1},\alpha_{b_2},\gamma_1,\gamma_2)=\pi\,K,\qquad \pi^2\,e_2(\alpha_{b_1},\alpha_{b_2},\gamma_1,\gamma_2)=\pi^{4-2-1}\,\Psi,
\]
so $K=\Psi=e_1(\alpha_{R_1})\,e_1(\alpha_{R_2})\,e_2(\alpha_{b_1},\alpha_{b_2},\gamma_1,\gamma_2)$: $\mathrm M_8$ is identically self-dual, its cobases being its bases.

\emph{Walls.} The walls of the cells are $\kappa_{R_i}=3-2\,(2-1)=1$ (Corollary~\ref{cor:mot-dualfull}(a)). Theorem~\ref{thm:wallsglue}(iii), at $R_1$ and then at $R_2$ in the cell of the room, both letters reading each $R_i$ through $\gamma_i$, reduces the positivity of the walls to that of the three families $(\mathcal A_1;1,1)$, $(\mathcal A_2;1,1)$ and $(\binom{\{b_1,b_2,h_1,h_2\}}2;1,1)$, all at the masses $\one$ (the hinges carrying $2\sigma-\kappa_{R_i}=1$ and $\kappa_{R_i}=1$): three copies of $\binom{[4]}2$ at its centre, convergent by Exercise~\ref{ex:mot-dualpair}. So $\mathrm M_8$ is convergent at its centre; its $254$ proper walls are positive integers (the receipt lists them, the least being $1$).

\emph{The motivic period.} Corollary~\ref{cor:mot-dualfull}(a) with $q=2$, the pinned families $\big(\binom{R_i}2;1,1\big)$ at $\one_{R_i}$ and the core $\binom{\{b_1,b_2,h_1,h_2\}}2$ at its centre:
\[
P^{\mathrm{mot}}_{\mathrm M_8}=P^{\mathrm{mot}}_{\binom{[4]}2}\cdot\big(P^{\mathrm{mot}}_{\binom{[3]}2}\big)^2=\big(P^{\mathrm{mot}}_{\binom{[3]}2}\big)^3,
\]
the second equality by Exercise~\ref{ex:mot-dualpair}(ii), every factor at its centre. In coefficients,
\[
\frac{\alpha^{\one}}{K\,\Psi}=\frac{\alpha_{b_1}\alpha_{b_2}\,\gamma_1\gamma_2}{e_2(\alpha_{b_1},\alpha_{b_2},\gamma_1,\gamma_2)^2}\cdot\prod_{i=1}^2\frac{\alpha^{\one_{R_i}}}{e_1(\alpha_{R_i})\,e_2(\alpha_{R_i})} .
\]

\emph{The bunch $B$.} $B$ is left in place and keeps in $\mathrm M_8$ the width two it has in the core (Theorem~\ref{thm:mot-pinned}(v)); by Corollary~\ref{cor:mot-whatsplits}(i) it does not split off and enters the period through a law carried by no point. The cells split off around it, and $B$ goes whole into the core $\binom{[4]}2$, which splits further only at the complement of a single mark, a set of width one (Exercise~\ref{ex:mot-dualpair}). Along $B$, $K(\tau\alpha_B,\alpha_{R_1},\alpha_{R_2})=\pi\,e_2(\tau\alpha_{b_1},\tau\alpha_{b_2},\gamma_1,\gamma_2)$, so the Skeleton values of $B$ in $\mathrm M_8$ are those of the pair $\{1,2\}$ of $\binom{[4]}2$ at the point $(\alpha_{b_1},\alpha_{b_2},\gamma_1,\gamma_2)$.

\emph{Answer.} $\mathrm P_5$ has as its bases the pairs of $\{1,\dots,5\}$ other than $\{1,2\}$, its cell $\binom{\{3,4,5,g\}}2$ is a dual pair, and $P^{\mathrm{mot}}_{\mathrm P_5}=(P^{\mathrm{mot}}_{\binom{[3]}2})^2$. $\mathrm M_8$ has the rank $4$; $R_1,R_2$ have width one and $B$ width two; $K=\Psi=e_1(\alpha_{R_1})e_1(\alpha_{R_2})\,e_2(\alpha_{b_1},\alpha_{b_2},\gamma_1,\gamma_2)$, and $P^{\mathrm{mot}}_{\mathrm M_8}=(P^{\mathrm{mot}}_{\binom{[3]}2})^3$. So
\begin{gather*}
P_{\mathrm P_5}(\one;1,1)=\Big(3\,\Tx{2}{\big(\Ilat{1+3\mathbb Z}-\Ilat{2+3\mathbb Z}\big)}{0}\Big)^2,\\
P_{\mathrm M_8}(\one;1,1)=\Big(3\,\Tx{2}{\big(\Ilat{1+3\mathbb Z}-\Ilat{2+3\mathbb Z}\big)}{0}\Big)^3,
\end{gather*}
with $\mathrm C(P^{\mathrm{mot}}_{\mathrm P_5})\subseteq\mathrm C(P^{\mathrm{mot}}_{\binom{[3]}2})\cdot\mathrm C(P^{\mathrm{mot}}_{\binom{[3]}2})$, and likewise with three factors for $\mathrm M_8$ (Theorem~\ref{thm:mot-pinned}(ii)). The bunch $B$ of width two does not split off and goes whole into the core.
\end{solution}

\section{Full leans}\label{sec:exmot-lean}

\begin{exercise}[Full leans at the exponents $(0,2)$]\label{ex:mot-fulllean}
Let $F=(\cB;0,2)$ be a convergent family of one list on $\cE$, $|\cE|\ge3$, at the masses $\one$ (Definition~\ref{def:mot-model}).
\begin{enumerate}[label=(\roman*)]
\item Show that $|\cE|=2r$ and that $\kappa_S=2r(S)-|S|$. At a set $R$ of width one with at least two marks on each side, show that the two hinge walls $\kappa_R$ and $\kappa_{\cE\setminus R}$ are $1$, that both cell families have the masses $\one$, and that the Beta corner of the split is $1$.
\item Show that every node of the splitting of Corollary~\ref{cor:mot-dualfull}(b) is prime and that $P^{\mathrm{mot}}_F=\prod_vP^{\mathrm{mot}}_{F_v}$, every $F_v=(\mathcal H_v;0,2)$ at the masses $\one$.
\item Let $\mathrm T_6$ be the list of rank three on $\{1,\dots,6\}$ whose bases are the nineteen triples other than $\{1,2,3\}$, let $\mathcal T$ be the cell $\mathrm T_6$ on $\{1,\dots,5\}\sqcup\{g\}$ with the mark $6$ as its hinge $g$, let $\mathcal T'$ be a copy of it on $\{1',\dots,5'\}\sqcup\{g\}$, and let $\mathrm D_{10}=\mathcal T\oplus_g\mathcal T'$ (Lemma~\ref{lem:twosumpieces}). Find the motivic period of $(\mathrm D_{10};0,2)$ at the masses $\one$, and decide whether $\{1',\dots,5'\}$ splits off at the centre and at $(0,2)$.
\end{enumerate}
\end{exercise}

\begin{solution}
Receipt: \texttt{ex5\_fulllean.py}, with the helper \texttt{subst.py}.

(i) \emph{Balance and walls.} At $(0,2)$ only $\Psi$ is active, $\sigma=\frac{0+2}2=1$, and the imbalance is $\kappa_\cE=n-2\,(n-r)=2r-n$ (Definition~\ref{def:family}); on the balance wall $n=2r$. The walls are $\kappa_S=m_S-2\,d^\ast(S)$ with $d^\ast(S)=|S|-r(S)$ (Definition~\ref{def:skeleton}), so
\[
\kappa_S=|S|-2\big(|S|-r(S)\big)=2r(S)-|S| .
\]

\emph{The hinge walls.} For a set $R$ of width one, Proposition~\ref{prop:compoundwall} gives
\[
\kappa_R+\kappa_{\cE\setminus R}=\kappa_\cE+2\sigma\,c(R)=0+2\cdot1\cdot1=2 ,
\]
and both walls are integers at least $1$, $F$ being convergent; so $\kappa_R=\kappa_{\cE\setminus R}=1$. By Proposition~\ref{prop:substitution}(iii) the list is $\mathcal H\langle\mathcal A\rangle$, the cell of the bunch $\mathcal A$ on $R\sqcup\{g\}$ substituted at the marker $h$ of the cell of the room $\mathcal H$, the marks of $L=\cE\setminus R$ left in place. One letter is active, so the cell is read through the one hinge value $\gamma^\Psi_{\mathcal A}$, and Theorem~\ref{thm:mot-pinned} applies with $q=1$, $2\sigma=2$, $\kappa_1=\kappa_R=1$ and the loads $a=0$, $b=1$. The core family $(\mathcal H;0,2)$ carries the masses $\one_L$ and $\kappa_R=1$ at $h$, and the cell family $(\mathcal A;0,2)$ the masses $\one_R$ and $2\sigma-\kappa_R=1$ at $g$: both have the masses $\one$, and both are convergent (Theorem~\ref{thm:mot-pinned}(i),(iii)). The pinned family $G$ has the tables $\mathcal A\setminus g$ at $(0,2-1)=(0,1)$ and $\mathcal A/g$ at $(0,1)$, at the masses $\one_R$. The Beta corner of the split is
\[
\Gst21{\kappa_1,\ 2\sigma-\kappa_1}{2\sigma}=\Gst21{1,\ 1}{2}=\frac{\Gamma(1)\,\Gamma(1)}{\Gamma(2)}=1 .
\]
So Theorem~\ref{thm:mot-pinned}(ii),(iii) give $P^{\mathrm{mot}}_F=P^{\mathrm{mot}}_{(\mathcal H;0,2)}P^{\mathrm{mot}}_G$ and $P^{\mathrm{mot}}_{(\mathcal A;0,2)}=1\cdot P^{\mathrm{mot}}_G$, that is
\[
P^{\mathrm{mot}}_F=P^{\mathrm{mot}}_{(\mathcal H;0,2)}\,P^{\mathrm{mot}}_{(\mathcal A;0,2)},
\]
both at the masses $\one$: the split costs no factor.

(ii) \emph{The nodes.} Each of the two families of (i) is again a convergent family of one list at $(0,2)$ with the masses $\one$, so (i) applies to it: it has an even number $2r_v$ of marks, and every further split costs no factor. A node made by a split has at least three marks, two of one side and its marker, and an even number of them, hence at least four; so no node on two or three marks occurs. A node at which the splitting stops has no set of width one with two marks on each side, and it is connected, since $\kappa_S+\kappa_{\cE_v\setminus S}=2c(S)$ with both walls at least $1$ gives $c(S)\ge1$ (Proposition~\ref{prop:compoundwall}, Definition~\ref{def:family}); so it is prime (Theorem~\ref{thm:primetree}), as Corollary~\ref{cor:mot-dualfull}(b) says of the nodes on at least four marks. By induction on the number of marks,
\[
P^{\mathrm{mot}}_F=\prod_vP^{\mathrm{mot}}_{F_v},\qquad F_v=(\mathcal H_v;0,2)\ \text{at the masses }\one,\qquad \mathrm C\big(P^{\mathrm{mot}}_F\big)\subseteq\prod_v\mathrm C\big(P^{\mathrm{mot}}_{F_v}\big)
\]
(Lemma~\ref{lem:mot-conj}(iii)): the rational constant of Corollary~\ref{cor:mot-dualfull}(b) is $1$, and every marker carries the mass $1$. By Corollary~\ref{cor:mot-cremona} the same holds at $(2,0)$, whose masses $2\sigma\one-\one$ are again $\one$.

(iii) \emph{The list $\mathrm T_6$.} Every set of at most three marks other than $\{1,2,3\}$ lies in a basis, $r(\{1,2,3\})=2$, and every set of at least four marks contains a basis; $n=6=2r$, so $\mathrm T_6$ at $(0,2)$ with the masses $\one$ is on its balance wall. By (i) its walls are $2r(S)-|S|$:
\[
\begin{array}{l|cccccc}
S & |S|=1 & |S|=2 & |S|=3,\ S\neq\{1,2,3\} & S=\{1,2,3\} & |S|=4 & |S|=5\\\hline
r(S) & 1 & 2 & 3 & 2 & 3 & 3\\
\kappa_S & 1 & 2 & 3 & 1 & 2 & 1
\end{array}
\]
all at least $1$: $(\mathrm T_6;0,2)$ is convergent. Its widths $c(S)=r(S)+r(\cE\setminus S)-3$ (Definition~\ref{def:skeleton}) are $1$ at one mark, $2$ at two marks, $3$ at a triple other than $\{1,2,3\}$ and $\{4,5,6\}$ and $2$ at these two: the sets with two marks on each side have widths two and three, and $\mathrm T_6$ is connected and prime (Theorem~\ref{thm:primetree}). The mark $6$ is neither a loop nor a coloop, so $\mathcal T$ is a cell, with the hinge minors
\[
\mathcal T\setminus g:\ \text{the nine triples of }\{1,\dots,5\}\text{ other than }\{1,2,3\},\qquad \mathcal T/g:\ \text{the ten pairs of }\{1,\dots,5\},
\]
of the ranks $3$ and $2$ (Definition~\ref{def:delcon}).

\emph{The list $\mathrm D_{10}$.} By Lemma~\ref{lem:twosumpieces} it is a list on ten marks of rank $3+3-1=5$, with the $9\cdot10+10\cdot9=180$ bases $X\sqcup Y$, $X\in\mathcal T\setminus g$ and $Y\in\mathcal T'/g$ or $X\in\mathcal T/g$ and $Y\in\mathcal T'\setminus g$, and $S=\{1',\dots,5'\}$ is a bunch of width one with the cell of the bunch $\mathcal T'$ and the cell of the room $\mathcal T$. At $(0,2)$ with the masses $\one$ it is on its balance wall, $10=2\cdot5$, and its hinge walls are (Lemma~\ref{lem:twosumpieces})
\[
\kappa_S=m_S-2\,\big(|S|-r_{\mathcal T'}\big)=5-2\,(5-3)=1,\qquad \kappa_{\cE\setminus S}=5-2\,(5-3)=1,
\]
adding to $2\sigma=2$. One letter is active, so by Theorem~\ref{thm:wallsglue}(iii) every proper wall of $\mathrm D_{10}$ is positive exactly when both cell families are, $(\mathcal T';0,2)$ and $(\mathcal T;0,2)$ at the masses $\one$, the hinge carrying $2\sigma-\kappa=1$: two copies of $(\mathrm T_6;0,2)$ at $\one$, whose walls are displayed above. So $(\mathrm D_{10};0,2)$ is convergent; the receipt finds the least of its $1022$ proper walls to be $1$.

\emph{The motivic period.} The split at $S$ gives the nodes $\mathcal T'$ and $\mathcal T$, copies of $\mathrm T_6$ with the marker in place of the mark $6$, both prime; by (ii), and Lemma~\ref{lem:mot-birational} for the renaming of the marks,
\[
P^{\mathrm{mot}}_{(\mathrm D_{10};0,2)}=\big(P^{\mathrm{mot}}_{(\mathrm T_6;0,2)}\big)^2,\qquad \mathrm C\big(P^{\mathrm{mot}}_{(\mathrm D_{10};0,2)}\big)\subseteq\mathrm C\big(P^{\mathrm{mot}}_{(\mathrm T_6;0,2)}\big)\cdot\mathrm C\big(P^{\mathrm{mot}}_{(\mathrm T_6;0,2)}\big),
\]
all at the masses $\one$, with the Beta corner $1$ and no other factor.

\emph{The set $S$ at the centre.} At $(1,1)$ with the masses $\one$ the walls are the widths (Proposition~\ref{prop:centrefamily}(ii)), all positive, $\mathrm D_{10}$ being connected by (ii); so Theorem~\ref{thm:mot-pointmass} applies. With both letters active and $|S|=5$, $S$ is pinned exactly when $\mathbf E_{\mathrm D_{10}|S}=\mathbf E_{\mathrm D_{10}/(\cE\setminus S)}$, and these hinge minors are $\mathcal T'\setminus g$ and $\mathcal T'/g$ (Theorem~\ref{thm:widthone}), with $9$ and $10$ bases. At the gauge $\one$ the even letter of a list is the square of its number of bases, so the two even letters take the values $81$ and $100$ there and differ: at the centre $S$ is not pinned, it does not split off, and its law is carried by no point. At $(0,2)$ one letter is active, so $S$ is pinned (Theorem~\ref{thm:mot-pointmass}(ii)) and splits off, with the Beta corner $1$, as above.

\emph{Answer.} At $(0,2)$ with the masses $\one$, $|\cE|=2r$ and $\kappa_S=2r(S)-|S|$; at a set of width one the two hinge walls add to $2$ and are $1$, the cell families keep the masses $\one$, and the Beta corner is $\Gst21{1,\ 1}{2}=1$. Every node of the splitting is prime, and $P^{\mathrm{mot}}_F=\prod_vP^{\mathrm{mot}}_{F_v}$ with every $F_v=(\mathcal H_v;0,2)$ at the masses $\one$. For $\mathrm D_{10}=\mathcal T\oplus_g\mathcal T'$, a list of rank five on ten marks with $180$ bases,
\[
P^{\mathrm{mot}}_{(\mathrm D_{10};0,2)}=\big(P^{\mathrm{mot}}_{(\mathrm T_6;0,2)}\big)^2\quad\text{at the masses }\one ;
\]
the set $\{1',\dots,5'\}$ splits off at $(0,2)$ and does not split off at the centre, where its hinge minors have the even letters $81$ and $100$ at the gauge $\one$.
\end{solution}

% Part XI, first half of XI.4: exercises on rays, the small centres (the list binom([3],1)).
% No new macros. Receipts: the scripts ex18_centre_bracket.py to ex23_directions.py of the exercise receipts.
\chapter{Exercises on rays: the small centres}\label{ch:exrays}

The exercises of this chapter read the list $\binom{[3]}1$ along its rays. Its centre ray is read four ways. A certificate proves its law at every strength, the low-end relation fixes the scale, and positivity with the native value of the direction brackets the value $\ell=P_{\binom{[3]}1}(\one;1,1)$ to $10^{-7}$. The box of its walls gives a spanning set of two insertions, the one-step matrix, a certificate supported on the stems and the law, with no closing computed, and shows why the closing relations alone do not give the law. The gauge end gives the values and their slopes to a hundred digits with no datum from the low end, and returns the low-end coefficients. The torus through the native gauge carries a second solution of the law; the torus value, the values and the torus read backwards are tied by an explicit connection, which reads the mean of $\log|K|$ on the torus and writes $\ell$ as three times a Trailmoment of index two of the lattice inputs of step three. A lean of three marks and the balanced directions of small size have their values in $\mathbb Q+\mathbb Q\ell$, read as values of insertions along the centre ray and reduced in its boxes. Every number displayed is exact unless it is labelled as a rendering; the receipts that recompute them are the scripts named in the solutions.

\section{The centre ray, bracketed}\label{sec:smallc-bracket}

\begin{exercise}[The centre ray of $\binom{[3]}1$, bracketed]\label{ex:smallc-bracket}
Let $F$ be the list $\binom{[3]}1$: the marks $\cE=\{1,2,3\}$, the bases $\{1\},\{2\},\{3\}$ and unit weights, so that its letters are
\[
K=\alpha_1+\alpha_2+\alpha_3,\qquad \Psi=\alpha_1\alpha_2+\alpha_1\alpha_3+\alpha_2\alpha_3 .
\]
Let $P(k)=P_F(k\one;k,k)$ be the values of its centre ray at the strength $1$, and $\ell=P(1)$.
\begin{enumerate}[label=(\roman*)]
\item From the closings of the ray find a candidate recurrence of the values, and prove it by a certificate.
\item Show that it holds at every $k$ with $\re k>0$, and read its low-end relation.
\item Find a rational $q$ with $0<\ell-q<10^{-6}$, and prove both inequalities.
\end{enumerate}
Exact computation may be done by machine; floating point, quadrature and integer-relation searches serve only as checks.
\end{exercise}

\begin{solution}
Receipt: \texttt{ex18\_centre\_bracket.py}. Throughout, $e_1,e_2,e_3$ are the elementary symmetric polynomials of $\alpha_1,\alpha_2,\alpha_3$, so that $K=e_1$ and $\Psi=e_2$, and $D_e=\alpha_e\,\partial/\partial\alpha_e$, which is $\partial_{y_e}$ in the coordinates $y_e=\log\alpha_e$.

\emph{Step 1: the ray.} The list has rank $1$. For a nonempty proper $S$, $r(S)=1$ and $r^\ast(S)=|S|+r(\cE\setminus S)-1=|S|$, so $d(S)=0$, $d^\ast(S)=|S|-1$ and the width is $c(S)=1$ (Definition~\ref{def:skeleton}). The family is connected at the strength $1$, and its centre ray (Definition~\ref{def:centreray}) has the origin $0$, the direction $(\one;1,1)$ and the walls $\kappa_{1,S}=c(S)=1$ for every nonempty proper $S$. On the chart $\alpha_3=1$ of the mark $e_0=3$, with $\alpha_1=e^{y_1}$, $\alpha_2=e^{y_2}$ and $dy=d\alpha_1\,d\alpha_2/(\alpha_1\alpha_2)$,
\[
f=\frac{\alpha^{\one}}{K\Psi}=\frac{e_3}{e_1e_2},\qquad P(k)=\int_{\mathbb R^2}f^k\,dy,\qquad
\ell=\int_0^\infty\!\!\int_0^\infty\frac{d\alpha_1\,d\alpha_2}{(1+\alpha_1+\alpha_2)(\alpha_1+\alpha_2+\alpha_1\alpha_2)} .
\]
The measure is $dy$: with $d\alpha_1\,d\alpha_2$ in its place the integrand is of the order of $1/(\alpha_1+\alpha_2)$ at infinity and the integral diverges. The volume of the direction (Proposition~\ref{prop:principal}) is $V_1=\sum_\pi\prod_{l=2}^3\kappa_{1,U_l}^{-1}=3!\cdot1=6$.

\emph{Step 2: the stride.} (a) Both sides of
\[
e_1e_2-9e_3=\alpha_1(\alpha_2-\alpha_3)^2+\alpha_2(\alpha_3-\alpha_1)^2+\alpha_3(\alpha_1-\alpha_2)^2
\]
equal $\sum_{i\neq j}\alpha_i^2\alpha_j-6e_3$, the left one because $e_1e_2=\sum_{i\neq j}\alpha_i^2\alpha_j+3e_3$. So $0<f\le\frac19$ on the orthant, with equality exactly on the diagonal, which on the chart is the point $y=0$: the native gauge of the direction is $\alpha^\star=\one$ and $f(\alpha^\star)=\frac19$ (Theorem~\ref{thm:pencil}(i)). Every $P(k)$, $k>0$, converges and is positive, and since $f<\frac19$ off one point,
\[
P(k+1)=\int f\cdot f^k\,dy<\tfrac19\,P(k)\qquad(k>0).
\]
(b) Order the coordinates, $\alpha_{(1)}\ge\alpha_{(2)}\ge\alpha_{(3)}$. Then $\alpha_{(1)}\le e_1\le3\alpha_{(1)}$, $\alpha_{(1)}\alpha_{(2)}\le e_2\le3\alpha_{(1)}\alpha_{(2)}$ and $e_3=\alpha_{(1)}\alpha_{(2)}\alpha_{(3)}$, so $\frac19\,\alpha_{(3)}/\alpha_{(1)}\le f\le\alpha_{(3)}/\alpha_{(1)}$. On the chart $\alpha_{(3)}/\alpha_{(1)}=e^{-\varphi(y)}$ with $\varphi(y)=\max(y_1,y_2,0)-\min(y_1,y_2,0)$. The lines $y_1=y_2$, $y_1=0$, $y_2=0$ cut the plane into six cones, one for each order of the entries of $(y_1,y_2,0)$, and on each $\varphi$ is the largest entry less the smallest; for real $k>0$
\[
\int_{y_1\ge y_2\ge0}e^{-ky_1}dy=\int_{y_1\ge0\ge y_2}e^{-k(y_1-y_2)}dy=\int_{0\ge y_1\ge y_2}e^{ky_2}dy=\frac1{k^2},
\]
and the three cones with $y_1$ and $y_2$ exchanged give the same, so $\int e^{-k\varphi}dy=6/k^2$ and
\[
\frac{6\cdot9^{-k}}{k^2}\le P(k)\le\frac6{k^2}\qquad(k>0).
\]
The six cones are the six orders $\pi$ of Proposition~\ref{prop:principal}, each contributing $\prod_l\kappa_{1,U_l}^{-1}=1$. On $[k_1,\infty)$ the integrands $f^k$ are dominated by $f^{k_1}$, as $0<f<1$, so $P$ is continuous on $(0,\infty)$.

\emph{Step 3: the candidate recurrence.} Here $f^{-1}=e_1e_2/e_3=(\alpha_1+\alpha_2+\alpha_3)(\alpha_1^{-1}+\alpha_2^{-1}+\alpha_3^{-1})$, a Laurent polynomial, and the closings of $1$ (Definition~\ref{def:closings}) are its constant terms,
\begin{gather*}
x_j=\mathsf C_j(1)=\CT\big(f^{-j}\big)=\sum_{a_1+a_2+a_3=j}\Big(\frac{j!}{a_1!\,a_2!\,a_3!}\Big)^2,\\
(x_j)_{j\ge0}=1,\ 3,\ 15,\ 93,\ 639,\ 4653,\ 35169,\ 272835,\ 2157759,\ \dots,
\end{gather*}
since the constant term pairs the monomial $\alpha^a$ of $(\sum_e\alpha_e)^j$ with the monomial $\alpha^{-a}$ of $(\sum_e\alpha_e^{-1})^j$, both with the coefficient $j!/(a_1!a_2!a_3!)$. Exact linear algebra with the ansatz $\sum_{i=0}^2p'_i(j)\,x_{j+i}=0$, $\deg p'_i\le2$, on $x_0,\dots,x_{15}$ ($14$ equations in $9$ unknowns) has a space of solutions of dimension one,
\[
(j+2)^2x_{j+2}-(10j^2+30j+23)\,x_{j+1}+9(j+1)^2x_j=0 .
\]
By Corollary~\ref{cor:reflected} a law $\sum_ip_i(k)\operatorname{Sh}^i[1]=0$ gives $\sum_ip_i(-j)\,x_{j-i}=0$; read back, $p_i(k)=p'_{2-i}(-k-2)$ is the candidate
\begin{equation}\tag{L}
k^2P(k)-(10k^2+10k+3)\,P(k+1)+9(k+1)^2\,P(k+2)=0 .
\end{equation}
At this point (L) is a candidate: sixteen closings do not prove their own recurrence, and Corollary~\ref{cor:reflected} runs from a law to its closings.

\emph{Step 4: the certificate.} The defects of the direction (Definition~\ref{def:ray}, Lemma~\ref{lem:insertions}(ii)) are
\[
c_e=D_e\log f=1-R_e-R^\ast_e,\qquad R_e=\frac{\alpha_e}{e_1},\qquad R^\ast_e=\frac{\alpha_e(e_1-\alpha_e)}{e_2},
\]
the occupations of $e$ under the basis law and under the cobasis law; $\sum_eR_e=1$ and $\sum_eR^\ast_e=2$, so $\sum_ec_e=0$. The origin is $0$, so its defects vanish, and a certificate of (L) (Proposition~\ref{prop:rays-strength}(i): an explicit element of the relation space $\mathcal L$ of Definition~\ref{def:raygauge} whose image is the recurrence) is a choice of insertions $g_1,g_2,g_3$ with
\begin{equation}\tag{C}
k^2-(10k^2+10k+3)\,f+9(k+1)^2f^2=\sum_{e=1}^3\big(k\,c_e+D_e\big)g_e .
\end{equation}
The ansatz: $g_e=N_e/(e_1^2e_2)$, $N_e$ homogeneous of degree $4$ with coefficients in $\mathbb Q+\mathbb Qk$, $N_1$ symmetric in $\alpha_2,\alpha_3$ and $g_e$ the image of $g_1$ under the exchange of the marks $1$ and $e$. Multiplied by $e_1^3e_2^2$, (C) is a polynomial identity of degree $7$ in $\alpha$ and $2$ in $k$, linear in the $18$ coefficients of $N_1$: $99$ equations of rank $10$. One solution is
\begin{equation}\tag{G}
g_1=-\frac{\alpha_1^2\big(k\,\alpha_1(\alpha_2+\alpha_3)+(2+3k)(\alpha_2^2+\alpha_3^2)-(1-k)\,\alpha_2\alpha_3\big)}{e_1^2e_2},
\end{equation}
and within the ansatz the others differ from it by one insertion $h$ added at every mark, $h$ with a symmetric numerator, an $8$-dimensional space; such an $h$ contributes $k\,h\sum_ec_e+\sum_eD_eh=0$, by Euler's identity for $h$ of degree $0$. Each $g_e$ is a polynomial over $K^2\Psi$, homogeneous of degree $0$, so a regular function on $X_F$: an insertion (Definition~\ref{def:ray}).

\emph{Verification.} Write $g_e=g^{(0)}_e+k\,g^{(1)}_e$, with $e_1^2e_2\,g^{(0)}_1=-\alpha_1^2(2\alpha_2^2-\alpha_2\alpha_3+2\alpha_3^2)$ and $e_1^2e_2\,g^{(1)}_1=-\alpha_1^2(\alpha_1\alpha_2+\alpha_1\alpha_3+3\alpha_2^2+\alpha_2\alpha_3+3\alpha_3^2)$. The powers of $k$ split (C) into
\[
\sum_ec_eg^{(1)}_e=1-10f+9f^2,\qquad\sum_e\big(D_eg^{(1)}_e+c_eg^{(0)}_e\big)=-10f+18f^2,\qquad\sum_eD_eg^{(0)}_e=-3f+9f^2 .
\]
Times $e_1^3e_2^2$ every sum over the marks is a symmetric polynomial of degree $7$, and in the basis of the products of $e_1,e_2,e_3$ exact expansion gives
\begin{gather*}
\textstyle\sum_ee_1^3e_2^2\,c_eg^{(1)}_e=e_1^3e_2^2-10\,e_1^2e_2e_3+9\,e_1e_3^2,\qquad
\sum_ee_1^3e_2^2\,D_eg^{(0)}_e=-3\,e_1^2e_2e_3+9\,e_1e_3^2,\\
\textstyle\sum_ee_1^3e_2^2\,D_eg^{(1)}_e=-2\,e_1e_2^3-3\,e_1^2e_2e_3+9\,e_1e_3^2,\qquad
\sum_ee_1^3e_2^2\,c_eg^{(0)}_e=2\,e_1e_2^3-7\,e_1^2e_2e_3+9\,e_1e_3^2 ,
\end{gather*}
which, with $e_1^3e_2^2f=e_1^2e_2e_3$ and $e_1^3e_2^2f^2=e_1e_3^2$, are the three right sides times $e_1^3e_2^2$. So (C) holds identically in $\alpha$ and $k$: $\sum_ip_i(k)f^i\in\mathcal L$, with $(p_0,p_1,p_2)=\big(k^2,\,-(10k^2+10k+3),\,9(k+1)^2\big)$. To test another certificate on one point: at $\alpha=(1,2,4)$ and $k=\frac13$, $f=\frac4{49}$, the left side of (C) is $-\frac{2809}{7203}$, and (G) gives $(g_1,g_2,g_3)=\big(-\frac{85}{1029},-\frac{310}{1029},-\frac{424}{1029}\big)$; agreement at a point is a test, the identity above is the proof.

\emph{Step 5: the recurrence at every strength.} The terms of $g_e$ are the insertions $\alpha^vK^{-2}\Psi^{-1}$ for the monomials $\alpha^v$ of $N_e$, with the data $(v;2,1)$ and the walls $\kappa_U=v_U-2d(U)-d^\ast(U)=v_U-|U|+1$. The monomials of $N_1$ are $\alpha_1^3\alpha_2$, $\alpha_1^3\alpha_3$, $\alpha_1^2\alpha_2^2$, $\alpha_1^2\alpha_3^2$ and $\alpha_1^2\alpha_2\alpha_3$; each has two marks in its support, so $v_U\ge1$ for every pair $U$, and every wall is at least $0$ (the least is $0$). With the walls $k\kappa_{1,U}=k$ of the ray, the condition of Proposition~\ref{prop:rays-strength}(i) holds at every $k$ with $\re k>0$, where the values converge, and the certificate gives (L) there. The integration by parts behind it handles the mark $e_0=3$ on its own chart: by Euler's identity $D_3=-D_1-D_2$ on functions of degree $0$, so $\sum_eD_e(g_ef^k)=\partial_{y_1}\big((g_1-g_3)f^k\big)+\partial_{y_2}\big((g_2-g_3)f^k\big)$; $y_3$ is not a coordinate of the chart. Corollary~\ref{cor:rays-lowend}(i) gives the same as an identity of meromorphic functions of $k$, for the element (C) of $\mathcal L$.

\emph{Step 6: the native values; what positivity selects.} The symbol of (L) (Corollary~\ref{cor:symbol}) is $\chi(\lambda)=1-10\lambda+9\lambda^2=(1-\lambda)(1-9\lambda)$. By Theorem~\ref{thm:signgauges} every sign $\eta$, $\eta_3=1$, with $K(\eta)\neq0$ is a native gauge of native value $K(\eta)^{-2}$, with the Hessian $2\big(\kappa_\eta(X_e,X_{e'})\big)_{e,e'\neq3}$ of the indicators under $\mathbb P_\eta(\{e\})=\eta_e/K(\eta)$: at $\eta=(1,1,1)$, $K(\eta)=3$, the native value is $\frac19$ and the Hessian is $2\big(\begin{smallmatrix}2/9&-1/9\\-1/9&2/9\end{smallmatrix}\big)$, of determinant $\frac4{27}$; at $(-1,1,1)$, $(1,-1,1)$, $(-1,-1,1)$, $K(\eta)=\pm1$, the native value is $1$ and the determinant is $-4$ (at $(-1,1,1)$ the matrix is $2\big(\begin{smallmatrix}-2&1\\1&0\end{smallmatrix}\big)$). All four are nondegenerate, so Corollary~\ref{cor:symbol} at $g=1$ makes $\frac19$ and $1$ roots of $\chi$; $\chi$ has degree $2$, so they are all its roots. (The corollary runs from native values to roots; a root of a symbol need not be a native value.) The root other than $f(\alpha^\star)=\frac19$ has the modulus $1>\frac19$, so Theorem~\ref{thm:positiveselection}(iii) applies: up to a constant factor the values are the only solution of (L) whose generating function has a radius of convergence of at least $9$, and the low-end relation fixes the factor. The rational solutions of (L) grow like the native value $1$ of the signs and the values decay like $9^{-k}$, so a bracket from positivity shrinks by about $\frac19$ per step of the strength.

\emph{Step 7: the low-end relation.} By Step 2(b), $6\cdot9^{-k}\le k^2P(k)\le6$ for real $k>0$, so $k^2P(k)\to6$ as $k\to0^+$: the low-end coefficient $\mathsf L_0$ is $V_1=6$ (Definition~\ref{def:rays-lowend}). Corollary~\ref{cor:rays-lowend}(ii) applies to the element (C): $p_0=k^2$ vanishes to the order $n-1=2$ at $0$, $[k^2]p_0=1$, $p_1(0)=-3$, $p_2(0)=9$, and
\begin{equation}\tag{S}
6-3\ell+9P(2)=0,\qquad P(2)=\frac{\ell-2}3 .
\end{equation}
Directly: take (L) at real $k\in(0,1)$ and let $k\to0^+$; $k^2P(k)\to6$, $P(k+1)\to\ell$ and $P(k+2)\to P(2)$ by the continuity of Step 2, while $10k^2+10k+3\to3$ and $9(k+1)^2\to9$.

A solution that looks right goes wrong here. Reading (L) at $k=0$ as $0\cdot P(0)-3P(1)+9P(2)=0$ treats $0\cdot P(0)$ as $0$; but $P(0)=\int dy=\infty$, and the product $k^2P(k)$ tends to $V_1=6$. The misreading gives $P(2)=\ell/3$, which Step 2(a) refutes at once: $P(2)<\frac19P(1)$ would say $\frac\ell3<\frac\ell9$. Every $P(K)$ would then be $\ell$ times a positive rational, positivity would fix no digit of $\ell$, and the bracket of Step 9 would be empty. Passing the limit under the integral, $k^2f^k\to0$ pointwise, fails for the same reason: there is no integrable majorant, $\int dy$ being infinite.

\emph{Step 8: the rational solutions and their signs.} For integers $k\ge1$, $9(k+1)^2\neq0$, so (L) gives $P(k+2)$ from $P(k)$ and $P(k+1)$. Let $c$ and $r$ be the solutions with $c(1)=1$, $c(2)=\frac13$ and $r(1)=0$, $r(2)=-\frac23$. From $P(1)=\ell$, (S) and induction,
\[
P(K)=c(K)\,\ell+r(K)\qquad(K\ge1).
\]
\emph{Signs at every $K$.} Put $x'_j=9^jc(j+1)$. Then (L) at $k=j+1$, times $9^{j+1}$, reads
\[
(j+2)^2x'_{j+2}=(10j^2+30j+23)\,x'_{j+1}-9(j+1)^2x'_j,\qquad x'_0=1,\ x'_1=3 .
\]
If $x'_{j+1}\ge3x'_j>0$, then $(j+2)^2x'_{j+2}\ge\big(10j^2+30j+23-3(j+1)^2\big)x'_{j+1}=(7j^2+24j+20)\,x'_{j+1}\ge3(j+2)^2x'_{j+1}$, because $7j^2+24j+20-3(j+2)^2=4(j+1)(j+2)\ge0$. So $x'_{j+1}\ge3x'_j>0$ for every $j$, $x'_j\ge3^j$, and
\[
c(K)\ge3^{1-K}>0,\qquad 9c(K)\ge3c(K-1)>c(K-1)\qquad(K\ge2).
\]
The closings satisfy the same recurrence with the same first two terms (Corollary~\ref{cor:rays-lowend}(i) for the element (C)), so $x'_j=x_j$ and $c(K)=9^{1-K}x_{K-1}$ at every $K$; the signs do not use this. In exact arithmetic:
\begin{center}
\small
\renewcommand{\arraystretch}{1.12}
\begin{tabular}{@{}r r r r@{}}
\toprule
$K$ & $c(K)$ & $r(K)$ & $a_K=-r(K)/c(K)$\\ \midrule
1 & $1$ & $0$ & $0$\\
2 & $1/3$ & $-2/3$ & $2$\\
3 & $5/27$ & $-23/54$ & $23/10$\\
4 & $31/243$ & $-145/486$ & $145/62$\\
5 & $71/729$ & $-1331/5832$ & $1331/568$\\
6 & $517/6561$ & $-242353/1312200$ & $242353/103400$\\
7 & $11723/177147$ & $-5495507/35429400$ & $5495507/2344600$\\
8 & $10105/177147$ & $-46422941/347208120$ & $46422941/19805800$\\
9 & $26639/531441$ & $-2447622221/20832487200$ & $2447622221/1044248800$\\
10 & $5773279/129140163$ & $-530455590781/5062294389600$ & $530455590781/226312536800$\\
\bottomrule
\end{tabular}
\end{center}
The frame of $c$ and $r$ (Proposition~\ref{prop:frame}, $h=2$) satisfies $C_{k+1}=\frac{p_0(k)}{p_2(k)}C_k=\frac{k^2}{9(k+1)^2}C_k$ with $C_1=c(1)r(2)-c(2)r(1)=-\frac23$, so
\[
c(k)\,r(k+1)-c(k+1)\,r(k)=-\frac23\cdot\frac{9^{1-k}}{k^2}\qquad(k\ge1).
\]

\emph{Step 9: the bracket and the answer.} (a) \emph{Below, by positivity.} $P(K)=c(K)\ell+r(K)>0$ and $c(K)>0$ give $\ell>a_K=-r(K)/c(K)$. (b) \emph{Above, by the native value.} $P(K)<\frac19P(K-1)$ (Step 2(a)) is $\big(9c(K)-c(K-1)\big)\ell<r(K-1)-9r(K)$, and $9c(K)-c(K-1)>0$ (Step 8), so
\[
\ell<b_K=\frac{r(K-1)-9\,r(K)}{9\,c(K)-c(K-1)} .
\]
At $K=2$ this is the first bracket: $P(2)=\frac{\ell-2}3>0$ gives $\ell>2$, and $P(2)<\frac19P(1)$ gives $\ell<3$. (c) \emph{The width.} With the frame at $k=K-1$ and $c(K)=9^{1-K}x_{K-1}$,
\begin{align*}
b_K-a_K&=\frac{r(K-1)\,c(K)-r(K)\,c(K-1)}{c(K)\big(9c(K)-c(K-1)\big)}=\frac{2\cdot9^{2-K}}{3(K-1)^2\,c(K)\big(9c(K)-c(K-1)\big)}\\
&=\frac{2\cdot9^{K-1}}{3(K-1)^2\,x_{K-1}(x_{K-1}-x_{K-2})},
\end{align*}
which the bounds of Step 8 make at most $1/(K-1)^2$: the brackets close on one number. In exact arithmetic, with the width also as a rendering:
\begin{center}
\small
\renewcommand{\arraystretch}{1.12}
\begin{tabular}{@{}r r r l@{}}
\toprule
$K$ & $b_K$ & $b_K-a_K$ & rendering\\ \midrule
2 & $3$ & $1$ & \\
3 & $19/8$ & $3/40$ & $7.5\cdot10^{-2}$\\
4 & $61/26$ & $3/403$ & $7.44\cdot10^{-3}$\\
5 & $3413/1456$ & $81/103376$ & $7.84\cdot10^{-4}$\\
6 & $104539/44600$ & $243/2882275$ & $8.43\cdot10^{-5}$\\
7 & $149014/63575$ & $2187/238492712$ & $9.17\cdot10^{-6}$\\
8 & $909876251/388187800$ & $19683/19613188595$ & $1.0036\cdot10^{-6}$\\
9 & $6414407843/2736630400$ & $6561/59511099776$ & $1.1025\cdot10^{-7}$\\
10 & $232184895407/99058909600$ & $177147/14589151085627$ & $1.214\cdot10^{-8}$\\
\bottomrule
\end{tabular}
\end{center}
At $K=8$ the certified width exceeds $10^{-6}$; at $K=9$ it is $\frac{6561}{59511099776}<\frac{111}{10^9}$.

\emph{Answer.} $q=a_9=\dfrac{2447622221}{1044248800}$, with
\[
\frac{2447622221}{1044248800}<\ell<\frac{6414407843}{2736630400},\qquad b_9-a_9=\frac{6561}{59511099776}<1.11\cdot10^{-7},
\]
so $0<\ell-q<10^{-6}$.

\emph{What does not certify.} The value $a_8$ is also a correct answer, but the bracket at $K=8$ is too wide to show it; the upper bound at $K=9$ does, $\ell-a_8<b_9-a_8=\frac{19683}{22574408320}<10^{-6}$. A rate of convergence is no bound, and neither is the last step: $a_9-a_8$ is smaller than $\ell-a_8$, since $a_9<\ell$. The bound $P(K)\le9^{-K}\ell$ does not follow from $f\le\frac19$ (it would need $\int dy<\infty$); iterating Step 2(a) gives $P(K)<9^{1-K}P(1)=9^{1-K}\ell$, that is $\big(c(K)-9^{1-K}\big)\ell<-r(K)$ with $c(K)-9^{1-K}=9^{1-K}(x_{K-1}-1)>0$ for $K\ge2$: a valid and cruder upper bound, which needs $K=10$ to certify six digits. A sign of $c(K)$ or of $9c(K)-c(K-1)$ that is not proved can reverse an inequality: here Step 8 proves both at every $K$.

\emph{Check} (floating point, not part of the solution). On the chart the integral over $\alpha_2$ is elementary, $\int_0^\infty\frac{d\alpha_2}{(\alpha_2+A)(A\alpha_2+B)}=\frac{\log(A^2/B)}{A^2-B}$ with $A=1+\alpha_1$, $B=\alpha_1$, so
\[
\ell=\int_0^\infty\frac{\log\big((1+x)^2/x\big)}{1+x+x^2}\,dx ,
\]
which renders as $2.343907238689458890601562288872277069095\ldots$; then $\ell-a_9$ and $b_9-\ell$ render as $9.6062\cdot10^{-8}$ and $1.4186\cdot10^{-8}$. Exercise~\ref{ex:smallc-gauge} computes $\ell$ to a hundred digits without quadrature.
\end{solution}

\section{The centre ray by its box}\label{sec:smallc-box}

\begin{exercise}[The centre of $\binom{[3]}1$ by its box]\label{ex:smallc-box}
Let $F$, $P$ and $\ell$ be as in Exercise~\ref{ex:smallc-bracket}. Read the centre ray of $F$ by steps (1) to (4) and (6) of the procedure of Definition~\ref{def:readprocedure}, with the box $\mathfrak B_2(1,1)$ (Definition~\ref{def:raybox}) modulo the group $G$ generated by the permutations of the marks and the Cremona: give its symbols, orbits and relations, the pivots and the free orbits, the spanning set at $\mathfrak s=1$, the matrix $\mathbf B(k)$, a certificate supported on the stems, the law of $[1]$, $\mathbf B(\infty)$ and the closings of the spanning set. Show that the closing relations alone do not give the law.
\end{exercise}

\begin{solution}
Receipt: \texttt{ex19\_centre\_box.py}. The elimination is small enough to run once over $\mathbb Q(k)$, so the evaluation at random strengths of step (2) of the procedure is not needed.

\emph{Step 1: walls and symmetries.} The walls of the direction are $\kappa_{1,S}=1$ and $V_1=6$ (Exercise~\ref{ex:smallc-bracket}, Step 1). Every permutation $\beta$ of the marks carries the bases and keeps the masses $\one$ and the exponents, and the Cremona is a symmetry because the masses are $\sigma\one$ and $s=t$ (Proposition~\ref{prop:raysym}(ii)). So $G=S_3\times\{1,\operatorname{cr}\}$, of order $12$, acting on the symbols by $(u;a,b)\mapsto(u\circ\beta^{-1};a,b)$ and $(u;a,b)\mapsto(-u+(a+b)\one;b,a)$ (Proposition~\ref{prop:raysym}(iii)).

\emph{Step 2: the box.} Both letters are active, so the box with one unit of each active letter at $\mathfrak t=2$ is $\mathfrak B=\mathfrak B_2(1,1)$: the symbols $(u;a,b)$ with $a,b\in\{0,1\}$, $u_{\cE}=a+2b$, and wall degree $\deg_\kappa=\max_U(-\kappa_U)\le2$, where $\kappa_U=u_U-a\,d(U)-b\,d^\ast(U)=u_U-b(|U|-1)$. It has $110$ symbols, $19$ at the level $(0,0)$, $27$ at each of $(1,0)$ and $(0,1)$ and $37$ at $(1,1)$. $G$ keeps it, the Cremona exchanging the levels $(1,0)$ and $(0,1)$, and the symbols fall into $16$ orbits:
\begin{center}
\small
\renewcommand{\arraystretch}{1.15}
\begin{tabular}{@{}l l r r l@{\qquad}l l r r l@{}}
\toprule
orbit of & level & $\deg_\kappa$ & size & & orbit of & level & $\deg_\kappa$ & size & \\ \midrule
$\alpha_1^4/(\alpha_3K\Psi)$ & $(1,1)$ & $2$ & $12$ & pivot & $\alpha_1^3/(\alpha_2\alpha_3K)$ & $(1,0),(0,1)$ & $2$ & $6$ & pivot\\
$\alpha_1^4\alpha_2/(\alpha_3^2K\Psi)$ & $(1,1)$ & $2$ & $6$ & pivot & $\alpha_1^2\alpha_2/(\alpha_3^2K)$ & $(1,0),(0,1)$ & $2$ & $12$ & pivot\\
$\alpha_1^3/(K\Psi)$ & $(1,1)$ & $1$ & $6$ & pivot & $\alpha_1^3/(\alpha_3^2K)$ & $(1,0),(0,1)$ & $2$ & $12$ & pivot\\
$\alpha_1^3\alpha_2/(\alpha_3K\Psi)$ & $(1,1)$ & $1$ & $6$ & pivot & $\alpha_1\alpha_2/(\alpha_3K)$ & $(1,0),(0,1)$ & $1$ & $6$ & pivot\\
$\alpha_1^2\alpha_2/(K\Psi)$ & $(1,1)$ & $0$ & $6$ & pivot & $\alpha_1^2/(\alpha_3K)$ & $(1,0),(0,1)$ & $1$ & $12$ & pivot\\
$\alpha_1\alpha_2\alpha_3/(K\Psi)$ & $(1,1)$ & $-1$ & $1$ & free & $\alpha_1/K$ & $(1,0),(0,1)$ & $0$ & $6$ & pivot\\
$\alpha_1^2/(\alpha_2\alpha_3)$ & $(0,0)$ & $2$ & $6$ & pivot & $\alpha_1/\alpha_3$ & $(0,0)$ & $1$ & $6$ & free\\
$\alpha_1^2/\alpha_3^2$ & $(0,0)$ & $2$ & $6$ & pivot & $1$ & $(0,0)$ & $0$ & $1$ & free\\
\bottomrule
\end{tabular}
\end{center}
The stems are the $19$ Laurent symbols. The relations (Definition~\ref{def:raybox}(iv)) are the letter identities for $K$ of the $46$ symbols with $a=0$, those for $\Psi$ of the $46$ symbols with $b=0$, and the mark relations of the $19$ stems at the $3$ marks: $149$ relations, which give $22$ distinct combinations of orbit averages (Definition~\ref{def:box-reduction}(i)), of rank $13$ over $\mathbb Q(k)$. With the ranking of Definition~\ref{def:box-reduction}(ii) (positive levels above the Laurent orbits, then the total level, then the wall degree), the elimination makes the $13$ orbits marked in the table pivots; the free orbits are those of $f=\alpha_1\alpha_2\alpha_3/(K\Psi)$, of $\alpha_1/\alpha_3$ and of $1$.

\emph{Step 3: the spanning set and the matrix.} The free Laurent orbits of wall degree at most $\mathfrak s=1$ give the spanning set $\mathbf b=(1,b_1)$, with $b_1=\frac16\sum_{i\neq j}\alpha_i/\alpha_j$, the average of the orbit of $\alpha_1/\alpha_3$ (Definition~\ref{def:box-reduction}(iii)). The first target is reduced already: $f^{-1}=(\sum_e\alpha_e)(\sum_e\alpha_e^{-1})=3+6b_1$ identically. The second, $f^{-1}b_1$, has $19$ monomials, and its reduction is a combination of $1$ and $b_1$. So the box is accepted at $\mathfrak s=1$, with
\[
\mathbf B(k)=\begin{pmatrix}3&6\\[2pt] \dfrac{2(k-1)}{k-2}&\dfrac{(k-1)(7k-11)}{(k-2)^2}\end{pmatrix},
\]
and by Theorem~\ref{thm:raybox}(i) and~(iii) $[f^{-1}b_i]=\sum_l\mathbf B_{il}(k)[b_l]$ in the module, and $\langle b_i\rangle_{k-1}=\sum_l\mathbf B_{il}(k)\langle b_l\rangle_k$ as meromorphic functions of $k$.

\emph{Step 4: the certificate.} Theorem~\ref{thm:raybox}(ii) provides combinations $G_e$ of stems with $f^{-1}b_1-\mathbf B_{21}-\mathbf B_{22}\,b_1=\sum_e\big(k\,c_e+D_e\big)G_e$, $c_e$ and $D_e$ as in Exercise~\ref{ex:smallc-bracket}. With $G_1$ a combination of stems symmetric under the exchange of the marks $2$ and $3$ and $G_e$ its image under the exchange of $1$ and $e$, the identity is linear in $11$ unknowns, $37$ equations, with a family of solutions of dimension $6$; one of them is
\begin{multline*}
G_1=\frac{\alpha_2^2+\alpha_3^2}{6(k-2)\,\alpha_1^2}+\frac{(5k^2-6k-4)\,\alpha_2\alpha_3}{6(k+1)(k-2)^2\,\alpha_1^2}+\frac{3k-4}{6(k+1)(k-2)^2}\Big(\frac{\alpha_3^2}{\alpha_1\alpha_2}+\frac{\alpha_2^2}{\alpha_1\alpha_3}\Big)\\
+\frac{5k-9}{6(k-2)^2}\,\frac{\alpha_2+\alpha_3}{\alpha_1}+\frac{3k-4}{6(k+1)(k-2)}\Big(\frac{\alpha_3^2}{\alpha_2^2}+\frac{\alpha_2^2}{\alpha_3^2}\Big),
\end{multline*}
whose terms are stems, of wall degree $2$, $2$, $2$, $1$ and $2$. The identity holds as one of rational functions in $\alpha$ and $k$ (exact expansion); at $\alpha=(1,2,4)$, $k=\frac13$, both sides are $\frac{11903}{800}$. Its coefficients have poles at $k=2$ and $k=-1$, and by Theorem~\ref{thm:raybox}(iii) it integrates term by term at every real $k>2$.

\emph{Step 5: the law.} By Proposition~\ref{prop:rays-lawrows} the rows are $w_0=(1,0)$, $w_1(k)=w_0\mathbf B(k)=(3,6)$ and
\[
w_2(k)=w_1(k-1)\,\mathbf B(k)=\Big(\frac{3(7k-10)}{k-2},\ \frac{6(10k^2-30k+23)}{(k-2)^2}\Big),
\]
and their first dependence is $q_0w_0+q_1w_1+q_2w_2=0$ with $(q_0,q_1,q_2)=\big(9(k-1)^2,\,-(10k^2-30k+23),\,(k-2)^2\big)$. Then $p_j(k)=q_{2-j}(k+2)$ gives $\big(k^2,\,-(10k^2+10k+3),\,9(k+1)^2\big)$: the recurrence (L) of Exercise~\ref{ex:smallc-bracket}, now with no closing computed, and with a certificate assembled from those of the rows (Theorem~\ref{thm:raybox}(iv)). As a check of Proposition~\ref{prop:rays-lawrows}(iv), $\det\mathbf B(k)=9(k-1)^2/(k-2)^2$ and $\det[w_0;w_1]=6$, so $q_0\cdot6=q_2\cdot6\cdot\det\mathbf B(k)$.

The spanning set is independent: at the native gauges $(1,1,1)$ and $(-1,1,1)$ (Exercise~\ref{ex:smallc-bracket}, Step 6) the vectors of values $\beta_c=(1,b_1(\alpha_c))$ are $(1,1)$ and $(1,-\frac13)$, which span $\mathbb C^2$, so by Proposition~\ref{prop:rays-spectrum}(i) $1,b_1$ are masters of the submodule $\mathfrak N$ they span and $\mathbf B=\mathbf B_{\mathbf b}$. By Proposition~\ref{prop:rays-lawrows}(ii) the order $2$ is then the dimension of the cyclic submodule of $[1]$, which is $\mathfrak N$, and (L) is the law of $[1]$. Finally
\[
\mathbf B(\infty)=\begin{pmatrix}3&6\\2&7\end{pmatrix},\qquad \mathbf B(\infty)\begin{pmatrix}1\\1\end{pmatrix}=9\begin{pmatrix}1\\1\end{pmatrix},\qquad \mathbf B(\infty)\begin{pmatrix}1\\-\frac13\end{pmatrix}=\begin{pmatrix}1\\-\frac13\end{pmatrix},
\]
the reciprocals of the native values $\frac19$ and $1$, with the vectors $\beta_c$ as eigenvectors (Proposition~\ref{prop:rays-spectrum}(ii)).

\emph{Step 6: checks.} The closings of the spanning set are $\mathsf C_j(1)=x_j$ and, by the first row, $\mathsf C_j(b_1)=\frac16(x_{j+1}-3x_j)=0,1,8,60,456,3535,27888,\dots$; they satisfy $\mathsf C_{j+1}(\mathbf b)=\mathbf B(-j)\,\mathsf C_j(\mathbf b)$ (Theorem~\ref{thm:closings}) for $0\le j<12$. The low-end relation of (L) returns $V_1=6$ (Exercise~\ref{ex:smallc-bracket}, Step 7).

\emph{Step 7: the closing relations alone.} The closing relations $(k\,c_e+D_e)(\alpha^wf^{-1})$ involve Laurent insertions only (Proposition~\ref{prop:box-closing}). Let $\omega=e^{2\pi i/3}$ and $z=(1,\omega,\omega^2)$. Then $K(z)=1+\omega+\omega^2=0$ and $\Psi(z)=\omega+\omega^2+\omega^3=0$, so $f^{-1}=K\Psi/e_3$ and every $D_ef^{-1}=(\Psi D_eK+KD_e\Psi)/e_3-f^{-1}$ vanish at $z$: $z$ is a singular point of the zero set of $f^{-1}$ on the torus. By Proposition~\ref{prop:box-closing}(i) the evaluation at $z$ vanishes on every closing relation, and by (ii) it takes the value $p_2(k-2)=9(k-1)^2\neq0$ on the element $\sum_{i=0}^2p_{2-i}(k-2)f^{-i}$ of $\mathcal L$ that the law gives: no combination of closing relations gives the law. Concretely, at the strength $9$ the elimination of the closing relations of the Laurent monomials of wall degree at most $R-1$, on the invariant Laurent combinations of wall degree at most $R$, leaves three free orbits, $1$, $b_1$ and $b_2=\frac16\sum_{i\neq j}\alpha_i^2/\alpha_j^2$, for every $R$ from $3$ to $7$; the law at $k=7$, read at the strength $9$ as $49f^{-2}-563f^{-1}+576$, reduces for each such $R$ to
\[
504+248\,b_1-392\,b_2\neq0,
\]
and at $z$, where $b_1(z)=b_2(z)=-\frac12$, this is $576=9(k+1)^2$. At $k=10$ and $k=\frac{23}3$ the reductions are $880+\frac{946}3b_1-\frac{2200}3b_2$ and $\frac{8671}{15}+\frac{1313}5b_1-\frac{6877}{15}b_2$, with the values $1089$ and $676$ at $z$, again $9(k+1)^2$. The point $z$ lies off $X_F$; the box supplies the relations that the closing relations lack through its units of the letters, by the letter identities and the mark relations of the occupations.

\emph{Answer.} The box $\mathfrak B_2(1,1)$ modulo $G$ has $110$ symbols in $16$ orbits and $149$ relations of rank $13$ modulo $G$; it is accepted at $\mathfrak s=1$ with the masters $1$ and $b_1=\frac16\sum_{i\neq j}\alpha_i/\alpha_j$ and the matrix $\mathbf B_{\mathbf b}(k)$ above, with the certificate $G_1$ on the stems; the law of $[1]$ is (L), and $\mathbf B_{\mathbf b}(\infty)$ has the eigenvalues $9$ and $1$. The closing relations alone leave the evaluation at $(1,\omega,\omega^2)$, on which the law takes the value $9(k+1)^2$.
\end{solution}

\section{The centre ray from its gauge end}\label{sec:smallc-gauge}

\begin{exercise}[The centre of $\binom{[3]}1$ from its gauge end]\label{ex:smallc-gauge}
Let $F$, $P$, $\ell$ and the law (L) be as in Exercise~\ref{ex:smallc-bracket}, and let $\mu_f$ be the stride's law of the centre ray (Definition~\ref{def:strideslaw}), with the density $\varrho_f$.
\begin{enumerate}[label=(\roman*)]
\item Write the density operator of (L), its exponents at $f(\alpha^\star)=\frac19$, and the germs there that extend by zero.
\item Find $\varrho_f(\frac19)$ and the Taylor coefficients of the germ of $\varrho_f$ at $\frac19$.
\item Read $P(1)$, $P'(1)$ and $P''(1)$ from the gauge end (Theorem~\ref{thm:gaugeend}) to sixty digits, find the low-end coefficients $\mathsf L_0,\mathsf L_1,\mathsf L_2$ from the faces of the sectors, and check the output against the low-end relation and its first two derivatives.
\item Find the arc integrals $S_j=-\int_0^{2\pi/3}\log^{j-1}\big(2\sin\frac\vartheta2\big)\,d\vartheta$ for $j=1,2$.
\end{enumerate}
\end{exercise}

\begin{solution}
Receipt: \texttt{ex20\_gauge\_end.py}. Theorem~\ref{thm:gaugeend} applies: the ray has the origin $0$ and positive walls, and (L) holds at every $k$ with $\re k>0$ (Exercise~\ref{ex:smallc-bracket}, Step 5), in particular at all large integers; $f(\alpha^\star)=\frac19$ is a simple root of the symbol $\chi(\lambda)=(1-\lambda)(1-9\lambda)$. Write $f^\star=\frac19$.

\emph{Step 1: the density operator.} With $\theta=t\,d/dt$ and $(p_0,p_1,p_2)=\big(k^2,-(10k^2+10k+3),9(k+1)^2\big)$,
\begin{align*}
\sum_it^ip_i(-1-i-\theta)&=(\theta+1)^2-t\,(10\theta^2+30\theta+23)+9t^2(\theta+2)^2\\
&=t^2(t-1)(9t-1)\,\partial_t^2+t\,(45t^2-40t+3)\,\partial_t+(36t^2-23t+1),
\end{align*}
of order $D=2$, with the leading coefficient $(-1)^Dt^D\chi(t)$; its singular points are $0$, $\frac19$, $1$ and $\infty$. The second symbol is $\chi_1(\lambda)=-10\lambda+18\lambda^2$.

\emph{Step 2: the exponents at $f^\star$.} In $s=f^\star-t$ the operator is $b_2\partial_s^2+b_1\partial_s+b_0$ with
\[
b_2=\tfrac1{81}s(9s-1)^2(9s+8),\qquad b_1=\tfrac1{81}(9s-1)(405s^2+270s-8),\qquad b_0=\tfrac19(324s^2+135s-10),
\]
so $b_{2,0}=0$, $b_{2,1}=\frac8{81}=-(f^\star)^2\chi'(f^\star)$ and $b_{1,0}=\frac8{81}$. Theorem~\ref{thm:gaugeend}(i) gives $\gamma=D-1-b_{1,0}/b_{2,1}=0$, which is $-1+\chi_1(f^\star)/(f^\star\chi'(f^\star))=-1+\frac{-8/9}{-8/9}$, and $\frac{n-3}2=0$ as Theorem~\ref{thm:gaugeend}(iii) requires. The polynomials $R_m(x)=\sum_jb_{j,m-1+j}[x]_j$ are
\[
R_0=\tfrac8{81}x^2,\qquad R_1=-\tfrac19(15x^2+23x+10),\qquad R_2=(2x+3)(3x+5),\qquad R_3=9(x+2)^2,
\]
and $R_m=0$ for $m\ge4$. The exponents at $s=0$ are $0$ and $0$: one solution $\phi=\sum_me_ms^m$, $e_0=1$, holomorphic at $s=0$, with $R_0(N)e_N=-\sum_{m=1}^3R_m(N-m)e_{N-m}$, and one of the form $\phi\log s+\psi$, $\psi$ holomorphic with $\psi(0)=0$.

\emph{Step 3: the line.} For a germ $v$ on $s>0$ with $(b_2\partial_s^2+b_1\partial_s+b_0)v=0$, integration by parts on $(\epsilon,\infty)$ shows that the distribution $\mathfrak D(v\one_{s>0})$ pairs with a test function $w$ as
\[
\lim_{\epsilon\to0^+}\Big[b_2\,(v'w-vw')+(b_1-b_2')\,vw\Big]_{s=\epsilon} .
\]
For $v=\phi$ this is $(b_{1,0}-b_{2,1})\,\phi(0)w(0)=0$: $\phi$ extends by zero (Definition~\ref{def:sl-densityop}). For $v=\phi\log s+\psi$ the term $b_2v'w$ tends to $b_{2,1}w(0)$, while the terms in $\log\epsilon$ have the coefficient $(b_{1,0}-b_{2,1})w(0)=0$, so the pairing is $\frac8{81}w(0)$: the germ leaves a point mass $\frac8{81}\delta_{f^\star}$ and does not extend by zero, and neither does any germ with a logarithmic term. So the germs that extend by zero form the line of $\phi$ (Theorem~\ref{thm:gaugeend}(ii)), the germ of $\varrho_f$ spans it (Theorem~\ref{thm:gaugeend}(iii)), and $\varrho_f=\varrho_f(f^\star)\,\phi(f^\star-t)$ near $f^\star$, the case $\gamma\in\{0,\dots,D-2\}$ of Theorem~\ref{thm:gaugeend}(iv). In exact arithmetic
\[
(e_m)_{m\ge0}=1,\ \tfrac{45}4,\ \tfrac{3645}{32},\ \tfrac{70713}{64},\ \tfrac{21500397}{2048},\ \tfrac{807258879}{8192},\ \tfrac{15034997331}{16384},\ \tfrac{139237010559}{16384},\ \dots,
\]
all positive.

\emph{Step 4: the value at the native value.} By Theorem~\ref{thm:branch}, on the centre ray the Hessian of $-\log f$ at $y=0$ is twice the covariance matrix of the indicators of the marks $1,2$ under the uniform law on the three marks,
\[
H=2\begin{pmatrix}\frac29&-\frac19\\-\frac19&\frac29\end{pmatrix},\qquad\det H=\tfrac4{27} .
\]
By Definition~\ref{def:raygauge} there are coordinates $v$ near $y=0$ with $\log f=\log f^\star-\frac12|v|^2$, here real, with $dy=J(v)\,dv$ and $J(0)=(\det H)^{-1/2}$. So $\mu_f([t,f^\star])=\int_{|v|\le\rho}J\,dv$ with $\rho^2=2\log(f^\star/t)$, whose derivative in $\rho$ is $2\pi\rho J(0)+O(\rho^2)$, and $\varrho_f(t)=-\frac d{dt}\mu_f([t,f^\star])=\frac1{\rho t}\big(2\pi\rho J(0)+O(\rho^2)\big)$, which tends to
\[
\varrho_f(f^\star)=\frac{2\pi}{f^\star\sqrt{\det H}}=2\pi\cdot9\cdot\frac{3\sqrt3}2=27\sqrt3\,\Gamma(\tfrac12)^2 .
\]

\emph{Step 5: the values and the slopes.} The other singular points of the operator lie at $s=\frac19$ ($t=0$) and $s=-\frac89$, so the series $\phi$ converges on $|s|<\frac19$, which contains $t\in(0,f^\star]$, and $\varrho_f(t)=27\sqrt3\,\Gamma(\frac12)^2\sum_me_m(\frac19-t)^m$ there: this series is the continuation of the germ along $(0,f^\star)$ in Theorem~\ref{thm:gaugeend}(v). Its terms are positive, so by monotone convergence Theorem~\ref{thm:gaugeend}(v) gives, for every real $K>0$ and $j\in\mathbb N$,
\[
P^{(j)}(K)=27\sqrt3\,\Gamma(\tfrac12)^2\sum_{m\ge0}e_m\,\frac{\partial^j}{\partial a^j}\Big[9^{-a-m-1}\,\Gst21{a+1,\ m+1}{a+m+2}\Big]_{a=K},
\]
the Beta corner $\int_0^{1/9}t^a(\frac19-t)^m\,dt=9^{-a-m-1}\Gst21{a+1,\ m+1}{a+m+2}$ and its derivatives in $a$ through the digamma function and its derivatives. At $K=150$ the terms fall below $10^{-150}$ times the first after $520$ terms. Differentiated $j$ times, (L) is $\sum_{l\le j}\binom jl\sum_ip_i^{(j-l)}(k)\,P^{(l)}(k+i)=0$, an identity of meromorphic functions, solved for $P^{(j)}(k)$ because $p_0(k)=k^2\neq0$ at $k\ge1$. The run from $K=150,151$ down to $k=1$, at $140$ digits, agrees with the run from $K=200,201$ to $138$ digits. The output, rendered to sixty digits:
\begin{gather*}
P(1)=\ell=2.343907238689458890601562288872277069095403009635856260666\ldots,\\
P'(1)=-8.075152263531644010269651241467376501347357821642322192221\ldots,\\
P''(1)=31.24865700624850564562564881961470385267564378216839540571\ldots,
\end{gather*}
and $P(2)=0.114635746229819630200520762957425689698467669878618753555\ldots$, $P'(2)=-0.316647537315232940933633840990335691962273491820393503777\ldots$ No datum of the low end enters.

\emph{Step 6: the low-end coefficients.} $\mathsf L_0=V_1=6$ (Definition~\ref{def:rays-lowend}), and $\mathsf L_1=0$ by Proposition~\ref{prop:rays-faces}, the weights being $1$. For $\mathsf L_2$ use the sector coordinates of an order $\pi$, $\alpha_{\pi(1)}=\rho$, $\alpha_{\pi(2)}=\rho q_2$, $\alpha_{\pi(3)}=\rho q_2q_3$: then $K=\rho(1+q_2+q_2q_3)$, $\Psi=\rho^2q_2(1+q_3+q_2q_3)$ and $e_3=\rho^3q_2^2q_3$, so
\[
f=q_2q_3\,\varphi_\pi,\qquad\varphi_\pi=\big((1+q_2+q_2q_3)(1+q_3+q_2q_3)\big)^{-1},\qquad\kappa_2=\kappa_3=1,
\]
the same for all six orders. In the formula of Proposition~\ref{prop:rays-faces} at $i=2$: $A=\varnothing$ gives $\frac12\log^2\varphi_\pi(0)=0$; $A=\{2\}$ gives, with $m=1$, $\int_0^1(\log\varphi_\pi)_{[\{2\}]}\,dq_2/q_2=-\int_0^1\log(1+q)\,dq/q$ and, with $m=0$, nothing, $1_{[A]}$ being $0$; $A=\{3\}$ gives the same as $\{2\}$; $A=\{2,3\}$ gives nothing. Expanding $\log(1+q)$,
\[
\int_0^1\log(1+q)\,\frac{dq}q=\sum_{m\ge1}\frac{(-1)^{m-1}}{m^2}=\frac12\Tx{2}{\IZ}{0},
\]
so each order contributes $-\Tx{2}{\IZ}{0}$ and $\mathsf L_2=-6\,\Tx{2}{\IZ}{0}=-\pi^2=-\Gamma(\frac12)^4$, by the classical value $\sum_{m\ge1}m^{-2}=\pi^2/6$.

\emph{Step 7: the output against the low end.} By Definition~\ref{def:rays-lowend}, $k^2P(k)=\mathsf L_0+\mathsf L_1k+\mathsf L_2k^2+O(k^3)$ near $0$, while $P(k+1)$ and $P(k+2)$ are holomorphic there. Expanding (L) at $k=0$ to the orders $0$, $1$ and $2$:
\begin{gather*}
\mathsf L_0-3P(1)+9P(2)=0,\qquad \mathsf L_1-10P(1)-3P'(1)+18P(2)+9P'(2)=0,\\
\mathsf L_2-10P(1)-10P'(1)-\tfrac32P''(1)+9P(2)+18P'(2)+\tfrac92P''(2)=0 .
\end{gather*}
The first is the low-end relation (Corollary~\ref{cor:rays-lowend}(ii)). With $\mathsf L_0=6$, $\mathsf L_1=0$ and $\mathsf L_2=-\Gamma(\frac12)^4$, the three left sides, read on the output of Step 5, are below $2\cdot10^{-138}$ in modulus: the gauge end returns the low-end coefficients it did not use. The expansion at the orders $3$ and $4$ reads $\mathsf L_3$ and $\mathsf L_4$ off the output; they render as $4.808227612638377141598952646045799963\ldots$ and $15.484819811652584020251232019547781376\ldots$, and agree with $4\,\Tx{3}{\IZ}{0}$ and $\frac7{60}\Gamma(\frac12)^8+\frac34\ell^2$ to $138$ digits (a check).

\emph{Step 8: the arc integrals.} $S_1=-\frac{2\pi}3=-\frac23\Gamma(\frac12)^2$. For $S_2$, expand $-\log\big(2\sin\frac\vartheta2\big)=\sum_{m\ge1}\frac{\cos m\vartheta}m$ on $(0,2\pi)$, a series that converges in $L^2(0,2\pi)$ and so may be integrated term by term over $(0,x)$:
\[
-\int_0^x\log\big(2\sin\tfrac\vartheta2\big)\,d\vartheta=\sum_{m\ge1}\frac{\sin mx}{m^2}\qquad(0\le x\le2\pi).
\]
At $x=\frac{2\pi}3$, $\sin\frac{2\pi m}3=\frac{\sqrt3}2\chi(m)$ with $\chi(m)=1,-1,0$ for $m\equiv1,2,0$ modulo $3$, so
\[
S_2=\frac{\sqrt3}2\Big(\Tx{2}{\Ilat{1+3\mathbb Z}}{0}-\Tx{2}{\Ilat{2+3\mathbb Z}}{0}\Big),
\]
which renders as $0.676627737606435750014135036183013523961\ldots$; $2\sqrt3\,S_2$ agrees with the value of $\ell$ of Step 5 to $140$ digits, and Exercise~\ref{ex:smallc-torus} proves $\ell=2\sqrt3\,S_2$.

\emph{Answer.} The germ at $\frac19$ is $\varrho_f(t)=27\sqrt3\,\Gamma(\frac12)^2\sum_me_m(\frac19-t)^m$, holomorphic, the line of the germs that extend by zero; it holds on all of $(0,\frac19]$. The values and slopes at $1$ are the renderings above, the low-end coefficients are $\mathsf L_0=6$, $\mathsf L_1=0$, $\mathsf L_2=-\Gamma(\frac12)^4$, and $S_1=-\frac23\Gamma(\frac12)^2$, $S_2=\frac{\sqrt3}2\big(\Tx{2}{\Ilat{1+3\mathbb Z}}{0}-\Tx{2}{\Ilat{2+3\mathbb Z}}{0}\big)$.
\end{solution}

\section{The torus at three marks}\label{sec:smallc-torus}

\begin{exercise}[The torus at three marks]\label{ex:smallc-torus}
Let $F$, $P$, $\ell$, (L), the solution $c$ and the closings $x_j$ be as in Exercise~\ref{ex:smallc-bracket}, and $S_2$ as in Exercise~\ref{ex:smallc-gauge}. On the torus $\alpha_3=1$, $\alpha_1=e^{i\vartheta_1}$, $\alpha_2=e^{i\vartheta_2}$, let
\[
T(k)=\frac1{(2\pi)^2}\int_{[0,2\pi]^2}|K|^{-2k}\,d\vartheta ,
\]
the torus value of the centre ray (Definition~\ref{def:torusvalue}), wherever the integral converges.
\begin{enumerate}[label=(\roman*)]
\item Show that $T$ converges for $\re k<1$, satisfies (L), continues meromorphically, has $T(-j)=x_j$ and has a simple pole at $k=1$ with the residue $R=-1/(\sqrt3\,\Gamma(\frac12)^2)$.
\item Show that $U(k)=9^{1-k}T(1-k)$ satisfies (L), that $U(K)=c(K)$, and that $P(k)U(k+1)-P(k+1)U(k)=\frac23\cdot9^{1-k}/k^2$ for every $k$.
\item Show that $T=\dfrac{P}{2\Gamma(\frac12)^4}-\dfrac{\cot(\pi k)}{\sqrt3}\,U$.
\item Find the torus mean of $\log|K|$, and show that $\ell=2\sqrt3\,S_2=3\big(\Tx{2}{\Ilat{1+3\mathbb Z}}{0}-\Tx{2}{\Ilat{2+3\mathbb Z}}{0}\big)$.
\end{enumerate}
\end{exercise}

\begin{solution}
Receipt: \texttt{ex21\_torus.py}.

\emph{Step 1: the stride on the torus.} On the torus $\alpha^{-1}=\bar\alpha$, and $\Psi(\alpha)=e_3\,K(\alpha^{-1})$, so $\Psi=e_3\bar K$ and $f=e_3/(K\Psi)=|K|^{-2}$, real and at least $\frac19$. Put $F=|K|^2=3+2\cos\vartheta_1+2\cos\vartheta_2+2\cos(\vartheta_1-\vartheta_2)$. Its gradient $\big(-2\sin\vartheta_1-2\sin(\vartheta_1-\vartheta_2),\,-2\sin\vartheta_2+2\sin(\vartheta_1-\vartheta_2)\big)$ vanishes exactly when $\sin\vartheta_1=-\sin\vartheta_2=-\sin(\vartheta_1-\vartheta_2)$: either $\vartheta_2=-\vartheta_1$ and $\sin\vartheta_1(1+2\cos\vartheta_1)=0$, or $\vartheta_2=\vartheta_1+\pi$ and $\sin\vartheta_1=0$. So $F$ has six critical points: $(0,0)$, the gauge $\one$, with $F=9$, its maximum; $(\pi,0)$, $(0,\pi)$, $(\pi,\pi)$, the signs, with $F=1$; and the two zeros $z_\pm=(\frac{2\pi}3,\frac{4\pi}3),(\frac{4\pi}3,\frac{2\pi}3)$ of $K$, where $F=Q(\delta)+O(|\delta|^3)$ in $\delta=\vartheta-\vartheta(z_\pm)$ with $Q(\delta)=\delta_1^2-\delta_1\delta_2+\delta_2^2$, of determinant $\frac34$. In $u=\log f=-\log F$ the Hessian at $(0,0)$ is $H=\big(\begin{smallmatrix}4/9&-2/9\\-2/9&4/9\end{smallmatrix}\big)$, of determinant $\frac4{27}$, the Hessian of Exercise~\ref{ex:smallc-gauge}, Step 4, read on the torus; at the signs it has the determinant $-4$ (Theorem~\ref{thm:signgauges}(ii) read on the torus): they are nondegenerate saddles.

\emph{Step 2: convergence, the closings and the law.} Near $z_\pm$, $|K|^{-2k}$ is of the size $|\delta|^{-2\re k}$, integrable in two variables exactly when $\re k<1$; so $T$ is holomorphic on $\re k<1$, with $|T(k)|\le T(\re k)$. For $j\in\mathbb N$, $T(-j)$ is the mean of the Laurent polynomial $f^{-j}=|K|^{2j}$ over the torus, its constant term: $T(-j)=\CT(f^{-j})=x_j$ (Definition~\ref{def:closings}). Multiply the certificate (C) of Exercise~\ref{ex:smallc-bracket} by $f^k=\exp(k\log f)$, $\log f$ real on the torus off $z_\pm$: $\sum_e(k\,c_e+D_e)g_e\cdot f^k=\sum_eD_e(g_ef^k)=D_1H_1+D_2H_2$ with $H_i=(g_i-g_3)f^k$ (Euler's identity, as in Step 5 there), and on the torus $D_e=-i\,\partial_{\vartheta_e}$. There $|\Psi|=|K|$, so $|g_e|\le C|K|^{-3}$, $|D_eg_e|\le C|K|^{-4}$ and $|c_e|\le C|K|^{-1}$; for $\re k\le-3$ the $H_i$ are $C^1$ on the torus and $\int\partial_{\vartheta_i}H_i\,d\vartheta=0$. So (L) holds for $T$ at $\re k\le-3$, and by the identity theorem wherever $T(k)$, $T(k+1)$, $T(k+2)$ are holomorphic, $\re k<-1$. Solved for $T(k+2)$, (L) continues $T$ meromorphically to $\mathbb C$, with poles only at the positive integers; on $\re k<2$ the only possible pole is at $1$, the argument of $T(k+2)$ at the zero $k=-1$ of $9(k+1)^2$.

\emph{Step 3: the residue at $1$.} Cut out the regions $Q(\delta)<\varepsilon^2$ around $z_\pm$. Off them the integrand is bounded for $\re k\le\frac32$, and on them $F^{-k}-Q^{-k}=O(|\delta|^{1-2\re k})$ is integrable for $\re k<\frac32$; both give functions holomorphic near $k=1$. And
\[
\int_{Q(\delta)<\varepsilon^2}Q(\delta)^{-k}\,d\delta=\frac{2\pi}{\sqrt{3/4}}\int_0^\varepsilon r^{1-2k}\,dr=\frac{4\pi}{\sqrt3}\cdot\frac{\varepsilon^{2-2k}}{2-2k}.
\]
Two zeros and the factor $(2\pi)^{-2}$ give $T(k)=\frac1{\sqrt3\,\pi}\cdot\frac1{1-k}+O(1)$ near $1$: a simple pole with the residue $R=-1/(\sqrt3\,\pi)=-1/(\sqrt3\,\Gamma(\frac12)^2)$. Write $T(k)=R/(k-1)+T_0+O(k-1)$.

\emph{Step 4: the torus read backwards.} Read at $-k-1$, (L) for $T$ is $(k+1)^2T(-k-1)-(10k^2+10k+3)\,T(-k)+9k^2T(1-k)=0$. With $U(k)=9^{1-k}T(1-k)$, $U(k+1)=9^{-k}T(-k)$ and $U(k+2)=9^{-k-1}T(-k-1)$,
\begin{multline*}
k^2U(k)-(10k^2+10k+3)\,U(k+1)+9(k+1)^2U(k+2)\\
=9^{-k}\big(9k^2T(1-k)-(10k^2+10k+3)\,T(-k)+(k+1)^2T(-k-1)\big)=0 .
\end{multline*}
So $U$ solves (L), and $U(K)=9^{1-K}x_{K-1}=c(K)$ (Exercise~\ref{ex:smallc-bracket}, Step 8). Its poles are at $0,-1,-2,\dots$, and $U(k)=-9R/k+O(1)$ at $0$. It is the solution of the type of the signs: $c(K)$ decays like the native value $1$, not like $9^{-K}$.

\emph{Step 5: the frame of the chart and the torus.} For meromorphic solutions $X,Y$ of (L), $W(X,Y)(k)=X(k)Y(k+1)-X(k+1)Y(k)$ satisfies $p_2(k)W(k+1)=p_0(k)W(k)$ (substitute $p_2X(k+2)=-p_0X(k)-p_1X(k+1)$ and the same for $Y$; Proposition~\ref{prop:frame} with $h=2$), so $k^29^{k-1}W(X,Y)(k)$ is periodic of period $1$. A periodic function $h$ of period $1$, holomorphic on the strip $|\re k|\le\frac12$ with $|h(k)|\le C(1+|k|)^N$ there, is constant: $h(k)=\mathcal H(e^{2\pi ik})$ with $\mathcal H$ holomorphic on $\mathbb C^\times$ and $|\mathcal H(z)|=O(|\log|z||^N)$ at $0$ and at $\infty$, so both singularities are removable and $\mathcal H$ is constant by Liouville's theorem. On the strip, $P$ is holomorphic except for its pole at $0$, $P(k)=6/k^2+O(1)$ ($\mathsf L_0=6$, $\mathsf L_1=0$: Exercise~\ref{ex:smallc-gauge}, Step 6), and polynomially bounded, being continued from $\re k>0$, where $|P(k)|\le P(\re k)$, by (L) solved for $P(k)$; $T(k)$ and $T(-k)$ are bounded; $T(k+1)$ and $T(1-k)$ have a simple pole at $0$ only and are polynomially bounded, by (L).

(a) In $k^29^{k-1}W(P,U)$ the pole of $P(k)$ at $0$ meets the holomorphic $U(k+1)$ and the simple pole of $U(k)$ meets the holomorphic $P(k+1)$, so the product with $k^2$ is holomorphic at $0$: the function is constant, and at $k=1$ it is $P(1)c(2)-P(2)c(1)=\frac\ell3-\frac{\ell-2}3=\frac23$. So
\[
P(k)U(k+1)-P(k+1)U(k)=\frac23\cdot\frac{9^{1-k}}{k^2}\qquad\text{for every }k,
\]
the frame of the chart and the torus read backwards; at the integers it is the frame of $r$ and $c$ of Exercise~\ref{ex:smallc-bracket}, Step 8.

(b) Let $m_1=(2\pi)^{-2}\int\log|K|\,d\vartheta$, so that $T(-k)=1+2m_1k+O(k^2)$ and $U(k+1)=1+(2m_1-\log9)\,k+O(k^2)$. Near $0$, $P(k)U(k+1)=6/k^2+6(2m_1-\log9)/k+O(1)$ and $P(k+1)U(k)=-9R\ell/k+O(1)$, while $\frac23\cdot9^{1-k}/k^2=6/k^2-6\log9/k+O(1)$. The terms in $1/k$ give $12m_1+9R\ell=0$:
\[
m_1=-\tfrac34R\,\ell=\frac{\sqrt3\,\ell}{4\,\Gamma(\frac12)^2} .
\]

\emph{Step 6: the connection up to one constant.} $W(P,U)\neq0$, so $P$ and $U$ are independent solutions and Cramer's rule on $T(k),T(k+1)$ gives $T=\alpha P+\beta U$ with $\alpha=W(T,U)/W(P,U)$ and $\beta=W(P,T)/W(P,U)$, periodic of period $1$ by (a). First,
\[
\alpha(k)=\tfrac{k^2}6\,T(k)T(-k)-\tfrac{3k^2}2\,T(k+1)T(1-k),
\]
holomorphic on the strip, the double pole $T(k+1)T(1-k)=-R^2/k^2+O(1)$ being cancelled by $k^2$, and polynomially bounded: $\alpha$ is constant, equal to its limit at $0$, $\alpha=\frac32R^2=\frac1{2\pi^2}=\frac1{2\Gamma(\frac12)^4}$. Next,
\[
\beta(k)=\tfrac{3k^2}2\,9^{k-1}\big(P(k)T(k+1)-P(k+1)T(k)\big)=\frac Rk+T_0+R\log9+O(k)\qquad\text{near }0,
\]
from $P(k)=6/k^2+O(1)$ and $T(k+1)=R/k+T_0+O(k)$. The function $R\,\pi\cot(\pi k)$ has simple poles at the integers with the residue $R$, so $\beta-R\,\pi\cot(\pi k)$ is periodic, entire and polynomially bounded on the strip, $\cot(\pi k)$ being bounded away from the real axis: it is a constant $b_0=T_0+R\log9$. So
\[
T=\frac P{2\Gamma(\frac12)^4}+\Big(b_0-\frac{\cot(\pi k)}{\sqrt3}\Big)U .
\]

\emph{Step 7: $b_0=0$, by the Fourier transform at the frequency $0$.} Put $t=e^u$. The image $\nu$ of $d\vartheta/(2\pi)^2$ under $u=-\log F$ is carried by $[u_0,\infty)$, $u_0=-\log9$, and has a density $n(u)$ with the following properties, by the Morse lemma at the nondegenerate critical points of Step 1 and the regularity of the other level sets:
(a) near the maximum, $u-u_0=\frac12|v|^2$ in coordinates $v$ with $d\vartheta=J\,dv$, $J(0)=(\det H)^{-1/2}$, so $n$ is continuous on $[u_0,u_0+\delta)$ with $n(u_0)=(2\pi)^{-2}\,2\pi J(0)=\frac{3\sqrt3}{4\pi}>0$; in $t$, the density of the torus law at $\frac19$ is $9\,n(u_0)=\frac{27\sqrt3}{4\pi}=(2\pi)^{-2}\varrho_f(\frac19)$, with $\varrho_f(\frac19)$ of Exercise~\ref{ex:smallc-gauge};
(b) near each saddle, $u=v_1v_2$ in coordinates with $d\vartheta=\tilde J\,dv$, $\tilde J(0)=\frac12$, the inverse square root of the modulus $4$ of the determinant of the Hessian,, and the density of $v_1v_2$ on $|v_1|,|v_2|<\delta$ at $s\neq0$ is $\int_{|s|/\delta<|v_1|<\delta}\tilde J(v_1,s/v_1)\,dv_1/|v_1|$; writing $\tilde J(v)=\tilde J(0)+v_1A(v)+v_2B(v)$, this is $2\tilde J(0)\log(\delta^2/|s|)$ plus two integrals that converge, as $s\to0$ from either side, to the same limits; so near $u=0$, with the three saddles, $n(u)=-a\log|u|+b(u)$ with $a=3\cdot2\tilde J(0)/(2\pi)^2=\frac3{4\pi^2}$ and $b$ continuous at $0$;
(c) near $z_\pm$, the nondegenerate minima of $F$, $n(u)=O(e^{-u})$, so $e^{cu}n(u)$ is integrable on $[u_0,\infty)$ for $0<c<1$.

Fix $c\in(0,1)$ and put $k=c+iy$, $y>0$. Then $T(k)=\int e^{iyu}\psi(u)\,du$ with $\psi=e^{cu}n$; $P(k)=\int e^{iyu}\varphi(u)\,du$ with $\varphi(u)=e^{(c+1)u}\varrho_f(e^u)$, carried by $u\le u_0<0$; and, with $w=u_0-u$ in $T(1-k)$, $U(k)=9^{1-c}\int e^{iyw}\tilde\psi(w)\,dw$ with $\tilde\psi(w)=e^{(1-c)(u_0-w)}n(u_0-w)$, carried by $w\le0$. And $\beta(k)=b_0+\frac i{\sqrt3}+O(e^{-2\pi y})$, as $\cot(\pi k)\to-i$.

Let $\chi$ be smooth with support in $[1,2]$ and $\int\chi=1$, $g(s)=s\chi(s)$, $\hat g(w)=\int g(s)e^{isw}ds$, rapidly decreasing, and for a function $\Phi$ of $y$ put $M_Y[\Phi]=Y^{-1}\int\Phi(y)\,y\,\chi(y/Y)\,dy$. For $\Phi(y)=\int e^{iyu}\psi(u)\,du$ with $\psi$ integrable, Fubini gives $M_Y[\Phi]=\int\psi(w/Y)\,\hat g(w)\,dw$. Now $\int\hat g=2\pi g(0)=0$, $\int_0^\infty\hat g=i\int g(s)\,s^{-1}ds=i$, $\int_{-\infty}^0\hat g=-i$, and $\int\log|w|\,\hat g(w)\,dw=-\pi$ (with $\hat g=-i\hat\chi'$, integrate by parts and use $\mathrm{PV}\!\int e^{isw}dw/w=i\pi$ for $s>0$). So, as $Y\to\infty$: if $\psi$ is bounded near $0$ with one-sided limits $\psi(0\pm)$, $M_Y[\Phi]\to i\big(\psi(0+)-\psi(0-)\big)$; if $\psi=-a\log|u|+b(u)$ near $0$ with $b$ continuous at $0$, $M_Y[\Phi]\to\pi a$.

Apply $M_Y$ to the connection along $k=c+iy$. The torus value gives $\pi a$, by (b). The values give $0$, $\varphi$ vanishing near $0$. The last term gives $\big(b_0+\frac i{\sqrt3}\big)\,9^{1-c}\,i\big(0-\tilde\psi(0-)\big)$, the factor $O(e^{-2\pi y})$ of $\beta$ contributing nothing in the limit, with $9^{1-c}\tilde\psi(0-)=9^{1-c}e^{(1-c)u_0}n(u_0)=n(u_0)$. With $Q=n(u_0)>0$:
\[
\pi a=\Big(b_0+\frac i{\sqrt3}\Big)(-i)\,Q=\frac Q{\sqrt3}-i\,b_0Q .
\]
$a$, $b_0$ and $Q$ are real, so $b_0=0$ (and $\pi a=Q/\sqrt3$: both are $\frac3{4\pi}$). Hence (iii), and $T_0=-R\log9=\dfrac{2\log3}{\sqrt3\,\Gamma(\frac12)^2}$.

\emph{Step 8: what the connection carries.} (a) At $k=1+j$, $j\ge0$, $P$ and $U$ are holomorphic and $-\cot(\pi k)/\sqrt3$ has the residue $-1/(\sqrt3\,\pi)$, so $\operatorname{Res}_{k=1+j}T=-\frac{9^{-j}x_j}{\sqrt3\,\Gamma(\frac12)^2}$: the poles of the torus value carry its values at the negative integers. (b) Near $0$, $P=2\Gamma(\frac12)^4\big(T+\frac{\cot(\pi k)}{\sqrt3}U\big)$ and $\frac{\cot(\pi k)}{\sqrt3}U=\frac1{\sqrt3\pi k}\cdot\frac{-9R}k+O(\frac1k)=\frac3{\pi^2k^2}+O(\frac1k)$: the low-end coefficient $\mathsf L_0=6$ is the residue of the torus at the zeros of $K$. (c) For real $a\ge0$, $a\neq1$, the mean of $\log|a+e^{i\vartheta}|$ over $\vartheta$ is $\log\max(a,1)$ (for $a>1$ write $\log|a+e^{i\vartheta}|=\log a+\re\log(1+e^{i\vartheta}/a)$, whose second term is a power series in $e^{i\vartheta}$ without constant term; for $a<1$ the same with $\log|1+ae^{-i\vartheta}|$). Integrating over $\vartheta_1$ first, with $a=|1+e^{i\vartheta_2}|=2|\cos\frac{\vartheta_2}2|$, which exceeds $1$ exactly when $|\vartheta_2|<\frac{2\pi}3$,
\[
m_1=\frac1\pi\int_0^{2\pi/3}\log\big(2\cos\tfrac\vartheta2\big)\,d\vartheta=-\frac1\pi\int_0^{\pi/3}\log\big(2\sin\tfrac\varphi2\big)\,d\varphi=\frac1\pi\sum_{m\ge1}\frac{\sin(m\pi/3)}{m^2},
\]
by $\vartheta=\pi-\varphi$, $\int_0^\pi\log(2\sin\frac\varphi2)\,d\varphi=0$ and the Fourier series of Exercise~\ref{ex:smallc-gauge}, Step 8. Here $\sin\frac{m\pi}3=\frac{\sqrt3}2\varepsilon(m)$ with $\varepsilon(m)=1,1,0,-1,-1,0$ for $m\equiv1,\dots,6$ modulo $6$, and $\varepsilon(m)=\chi(m)$ for odd $m$, $-\chi(m)$ for even $m$, $\chi$ as in Exercise~\ref{ex:smallc-gauge}; $\chi(2m)=-\chi(m)$, so $\sum_m\varepsilon(m)m^{-2}=\sum_m\chi(m)m^{-2}+\frac12\sum_m\chi(m)m^{-2}$, and $m_1=\frac{3\sqrt3}{4\pi}\sum_m\chi(m)m^{-2}=\frac3{2\pi}S_2$. With Step 5(b):
\[
\ell=\frac{4\pi m_1}{\sqrt3}=2\sqrt3\,S_2=3\Big(\Tx{2}{\Ilat{1+3\mathbb Z}}{0}-\Tx{2}{\Ilat{2+3\mathbb Z}}{0}\Big).
\]

\emph{Checks} (floating point, not part of the solution). By the mean over one circle, $T(k)=\frac2\pi\int_0^2I_k(a)\,\frac{da}{\sqrt{4-a^2}}$ with $I_k(a)$ the mean of $|a+e^{i\vartheta}|^{-2k}$, a hypergeometric series in $a^2$ or $a^{-2}$; this gives $T(-j)=x_j$, (L) on $T$ at $k=-1.3$ and $k=-2.7+0.4i$ to $20$ digits, and, with $P$ from the gauge end, the connection of (iii) at $k=0.2,0.3,\dots,0.8$ to between $17$ and $23$ digits, $k^29^{k-1}W(P,U)=\frac23$ at the same points to $20$ digits, $\beta(\frac12)=b_0$ below $2\cdot10^{-23}$, and $T_0$ rendering as $0.40379779973879724450\ldots$; $m_1$ renders as $0.323065947219450514093636510724\ldots$

\emph{Answer.} $T$ continues to a meromorphic solution of (L) with $T(-j)=x_j$ and the residue $-1/(\sqrt3\,\Gamma(\frac12)^2)$ at $1$; $U(k)=9^{1-k}T(1-k)$ is the solution $c$ of Exercise~\ref{ex:smallc-bracket}, with $P(k)U(k+1)-P(k+1)U(k)=\frac23\cdot9^{1-k}/k^2$; the connection is $T=\frac P{2\Gamma(\frac12)^4}-\frac{\cot(\pi k)}{\sqrt3}U$; the torus mean of $\log|K|$ is $\frac{\sqrt3\,\ell}{4\Gamma(\frac12)^2}$, and $\ell=2\sqrt3\,S_2=3\big(\Tx{2}{\Ilat{1+3\mathbb Z}}{0}-\Tx{2}{\Ilat{2+3\mathbb Z}}{0}\big)$.
\end{solution}

\section{A lean of three marks}\label{sec:smallc-lean}

\begin{exercise}[The lean $(2\one;4,1)$ of $\binom{[3]}1$]\label{ex:smallc-lean}
Let $F$ and $\ell$ be as in Exercise~\ref{ex:smallc-bracket}, and consider the ray of $F$ with the origin $0$ and the direction $(2\one;4,1)$, with the values $P_\lambda(k)=P_F(2k\one;4k,k)$.
\begin{enumerate}[label=(\roman*)]
\item Find its walls, $V_1$, its stride and native value, and its closings.
\item Reduce its box $\mathfrak B_2(1,1)$ modulo the permutations of the marks: the spanning set at $\mathfrak s=1$, the matrix, the law of $[1]$, its symbol and its native gauges, and $\mathbf B(\infty)$.
\item Read the low-end relation, and show that $54\,P_\lambda(1)=11-4\ell$.
\item Write $P_\lambda(1)$ on the share side, and say why the gap theorem cannot be applied to the rational coordinates of this ray.
\end{enumerate}
\end{exercise}

\begin{solution}
Receipt: \texttt{ex22\_lean.py}.

\emph{Step 1: the ray.} By Definition~\ref{def:family} and the counts of Exercise~\ref{ex:smallc-bracket}, Step 1, the walls of the direction are $\kappa_{1,S}=m_S-4d(S)-d^\ast(S)=m_S-(|S|-1)$: $2$ on the marks and $3$ on the pairs. Every order $\pi$ has the tails $U_2$, a pair, and $U_3$, a mark, so $V_1=6\cdot\frac1{3\cdot2}=1$ (Proposition~\ref{prop:principal}). The stride is $f=\alpha^{2\one}K^{-4}\Psi^{-1}=e_3^2/(e_1^4e_2)$. The defects $c_e=2-4R_e-R^\ast_e$ vanish at $\one$ ($2-\frac43-\frac23=0$), so $\alpha^\star=\one$ and $f(\alpha^\star)=\frac1{3^4\cdot3}=\frac1{243}$. The constant term of $f^{-j}=e_1^{4j}\big(\sum_e\alpha^{\cE\setminus e}\big)^j\alpha^{-2j\one}$ pairs the monomial $\alpha^{(j-b_1,j-b_2,j-b_3)}$ of the second factor with the monomial $\alpha^{(j+b_1,j+b_2,j+b_3)}$ of $e_1^{4j}$, so
\begin{gather*}
x_j=\CT(f^{-j})=(4j)!\,j!\sum_{b_1+b_2+b_3=j}\ \prod_{e=1}^3\frac1{b_e!\,(b_e+j)!},\\
(x_j)_{j\ge0}=1,\ 36,\ 4620,\ 762300,\ 140216076,\ 27409661136,\ \dots
\end{gather*}

\emph{Step 2: the box.} The permutations of the marks are symmetries; the Cremona is not, the exponents $4$ and $1$ being different (Proposition~\ref{prop:raysym}(ii)). Modulo $S_3$ the box $\mathfrak B_2(1,1)$ (Definition~\ref{def:raybox}), the box with one unit of each letter at $\mathfrak t=2$, has $398$ symbols in $79$ orbits ($17$, $19$, $20$ and $23$ at the levels $(0,0)$, $(1,0)$, $(0,1)$, $(1,1)$) and $82$ stems. Its relations are $184$ letter identities for $K$, $178$ for $\Psi$ and $246$ mark relations of stems; modulo $S_3$ they give $117$ distinct combinations of orbit averages, of rank $76$. The free orbits are those of $\alpha_1\alpha_2\alpha_3/(K\Psi)$ (wall degree $-\frac13$), of $\alpha_1/\alpha_3$ (wall degree $\frac12$) and of $1$, so the spanning set at $\mathfrak s=1$ is again $(1,b_1)$, $b_1=\frac16\sum_{i\neq j}\alpha_i/\alpha_j$, and the box is accepted there (Definition~\ref{def:box-reduction}) with
\begin{gather*}
\mathbf B_{11}=\frac{199k^3-420k^2+299k-72}{(k-1)^2(3k-2)},\qquad\mathbf B_{12}=\frac{(2k-1)(265k^2-424k+171)}{(k-1)^2(3k-2)},\\
\mathbf B_{21}=\frac{3710k^4-12349k^3+15022k^2-7997k+1584}{3(k-1)(2k-3)(3k-4)(3k-2)},\\
\mathbf B_{22}=\frac{(2k-1)(4706k^3-13323k^2+12238k-3681)}{3(k-1)(2k-3)(3k-4)(3k-2)}.
\end{gather*}
Certificates on the stems (Theorem~\ref{thm:raybox}(ii)): for each row, $G_1$ a combination of the stems symmetric under the exchange of the marks $2$ and $3$, $G_e$ its images; the identity is linear in $44$ unknowns with $118$ equations and a family of solutions of dimension $24$; a solution with $20$ nonzero coefficients, rational functions of $k$, satisfies it exactly for each row. It is not displayed here; the receipt prints it.

\emph{Step 3: the law.} The rows of Proposition~\ref{prop:rays-lawrows} give the first dependence with
\begin{gather*}
q_0=-972(k-1)^2(4k-3)(4k-1)(265k^2-954k+860),\\
q_1=-1721705k^6+13084958k^5-40594556k^4+65762882k^3\\
-58662047k^2+27325016k-5197140,\\
q_2=3(k-2)^2(3k-5)(3k-4)(265k^2-424k+171),
\end{gather*}
and $p_j(k)=q_{2-j}(k+2)$, divided by their common factor, is the recurrence
\begin{multline*}
-3k^2(3k+1)(3k+2)(265k^2+636k+383)\,P_\lambda(k)\\
+(1721705k^6+7575502k^5+13047276k^4+11068046k^3+4766659k^2+965100k+75600)\,P_\lambda(k+1)\\
+972(k+1)^2(4k+5)(4k+7)(265k^2+106k+12)\,P_\lambda(k+2)=0,
\end{multline*}
of degree $6$, with the symbol $265(243\lambda-1)(64\lambda+27)$; its reflection is satisfied by the closings of Step 1 (Corollary~\ref{cor:rays-lowend}(i)). The native gauges: at $\one$ the native value $\frac1{243}$, and the Hessian of $\log f$ on the chart $\alpha_3=1$ has the determinant $\frac{25}{27}$; at $\alpha=(3,3,-2)$ and its permutations all three defects vanish (exact evaluation), the native value is $\frac{(-18)^2}{4^4\cdot(-3)}=-\frac{27}{64}$, and the Hessian has the determinant $-\frac{75}2$. Both are nondegenerate, their native values are the roots of the symbol (Corollary~\ref{cor:symbol}), and the vectors $(1,b_1(\alpha_c))$ are $(1,1)$ and $(1,-\frac7{18})$, which span: by Proposition~\ref{prop:rays-spectrum}(i) the spanning set is a set of masters and the recurrence is the law of $[1]$. Finally
\[
\mathbf B(\infty)=\begin{pmatrix}\frac{199}3&\frac{530}3\\[2pt]\frac{1855}{27}&\frac{4706}{27}\end{pmatrix},
\]
with the eigenvalues $243$ and $-\frac{64}{27}$ and the eigenvectors $(1,1)$ and $(1,-\frac7{18})$ (Proposition~\ref{prop:rays-spectrum}(ii)).

\emph{Step 4: the low-end relation.} $p_0$ vanishes to the order $2$ at $0$ with $[k^2]p_0=-3\cdot1\cdot2\cdot383=-2298$, $p_1(0)=75600$ and $p_2(0)=972\cdot5\cdot7\cdot12=408240$; with $\mathsf L_0=V_1=1$, Corollary~\ref{cor:rays-lowend}(ii) gives $-2298+75600\,P_\lambda(1)+408240\,P_\lambda(2)=0$, that is
\[
12600\,P_\lambda(1)+68040\,P_\lambda(2)=383 .
\]

\emph{Step 5: the value, read on the centre ray.} At $k=2$ the centre ray of Exercise~\ref{ex:smallc-bracket} has the stride $f_{\mathrm c}=\alpha^{\one}/(K\Psi)$, and
\[
\big\langle\Psi K^{-2}\big\rangle_2=\int\Psi K^{-2}\,\frac{\alpha^{2\one}}{K^2\Psi^2}\,dy=\int\frac{\alpha^{2\one}}{K^4\Psi}\,dy=P_\lambda(1):
\]
the value of the lean is the value of the insertion $\Psi K^{-2}=\sum_e\alpha^{\cE\setminus e}K^{-2}$ along the centre ray. In the box $\mathfrak B_2(2,2)$ of the centre ray modulo $S_3$ and the Cremona ($339$ symbols, $43$ orbits, $499$ relations of rank $40$, the free orbits those of $f_{\mathrm c}^2$, of $\alpha_1/\alpha_3$ and of $1$) the orbit average of $\Psi K^{-2}$ reduces to
\[
\Psi K^{-2}\equiv-\frac{2(k-1)}{3(k+1)}+\frac{(k-1)^2}{k(k+1)}\,b_1 .
\]
The relations of a box evaluate into $\mathcal L$ and their values vanish identically in $k$ (Proposition~\ref{prop:box-relations}(i)); the symmetries keep the values (Proposition~\ref{prop:raysym}(ii)); and at $k=2$ every value involved converges. So, with $\langle b_1\rangle_k=\frac16\big(\langle f_{\mathrm c}^{-1}\rangle_k-3\langle1\rangle_k\big)=\frac16\big(P(k-1)-3P(k)\big)$ from $f_{\mathrm c}^{-1}=3+6b_1$, $\langle b_1\rangle_2=\frac16\big(\ell-(\ell-2)\big)=\frac13$, and
\[
P_\lambda(1)=-\frac29\,P(2)+\frac16\cdot\frac13=-\frac29\cdot\frac{\ell-2}3+\frac1{18}=\frac{11-4\ell}{54} .
\]
Then Step 4 gives $P_\lambda(2)=\frac{2800\ell-6551}{204120}$, and the law gives every $P_\lambda(K)$ in $\mathbb Q+\mathbb Q\ell$.

\emph{Step 6: the share side.} By Proposition~\ref{prop:dualfamily}, $P_F(\mvec;s,t)$ is the period of the dual table, whose bases are the sets $\cE\setminus e$, at the exponents $(t,s)$; that table is one circuit, and Corollary~\ref{cor:cycle} with the exponents $(1,4)$, $\nu_e=m_e-1=1$ and unit weights (its condition $m_S>|S|-1$ holds: $2>0$, $4>1$) gives
\[
P_\lambda(1)=\frac1{\Gamma(4)\Gamma(1)}\int_0^\infty\sG_{1;(1,-1)}(1,y)^3\,dy ,
\]
that is $\frac16\int_0^\infty\sG_{1;(1,-1)}(1,y)^3\,dy$.
A check by quadrature: both this integral and $\frac{11-4\ell}{54}$ render as $0.030080945282262304399884274898349847\ldots$

\emph{Step 7: why the gap does not apply.} Write $P_\lambda(K)=c_\lambda(K)\,\ell+r_\lambda(K)$ with rationals from Steps 3 to 5, and let $D_K$ be the least common denominator of $c_\lambda(K)$ and $r_\lambda(K)$, the least integer that clears the row. Theorem~\ref{thm:gap} needs $D_KP_\lambda(K)\to0$. But the values decay no faster than $243^{-K}$: as in Exercise~\ref{ex:smallc-bracket}, Step 2(b), $f\ge\frac1{243}\,\alpha_{(2)}\alpha_{(3)}^2/\alpha_{(1)}^3$, and on the cone of an order the exponent $3y_{(1)}-y_{(2)}-2y_{(3)}=3(y_{(1)}-y_{(2)})+2(y_{(2)}-y_{(3)})$ has the walls $3$ and $2$ as coefficients, so $P_\lambda(K)\ge243^{-K}\cdot6\cdot\frac1{6K^2}=243^{-K}/K^2$. The exact denominators satisfy $D_K>243^KK^2$ for every $K$ with $3\le K\le240$, so $D_KP_\lambda(K)>1$ throughout that range; $\frac1K\log D_K$ renders as about $10.3$ at $K=240$, against $\log243=5.49$. The clearing integers of this ray grow faster than its values decay, and the hypothesis of Theorem~\ref{thm:gap} fails on every computed row.

\emph{Answer.} Walls $2$ and $3$, $V_1=1$, $f(\alpha^\star)=\frac1{243}$; the masters $1,b_1$, the law above with the symbol $265(243\lambda-1)(64\lambda+27)$ and the native gauges $\one$ and $(3,3,-2)$; the low-end relation $12600P_\lambda(1)+68040P_\lambda(2)=383$; and $54\,P_\lambda(1)=11-4\ell$, with the share side $\frac16\int_0^\infty\sG_{1;(1,-1)}(1,y)^3dy$.
\end{solution}

\section{The balanced directions}\label{sec:smallc-directions}

\begin{exercise}[The balanced directions of $\binom{[3]}1$]\label{ex:smallc-directions}
Let $F$ and $\ell$ be as in Exercise~\ref{ex:smallc-bracket}. For a direction $(\mvec;s,t)$ on the balance wall put $c_e=m_e-t$.
\begin{enumerate}[label=(\roman*)]
\item List the balanced directions with $1\le s,t\le5$, positive walls and $m_1\ge m_2\ge m_3$, and sort them by $S=s+t$ and the multiset $\{|c_e|\}$.
\item Show that $\Gamma(sk)\Gamma(tk)\,P_F(k\mvec;ks,kt)$ depends on the direction only through $S$ and $\{|c_e|\}$, so that within a class the values are in the ratio of two Beta corners.
\item For every direction with $|c_e|$ constant find $P_F(\mvec;s,t)$ in $\mathbb Q+\mathbb Q\ell$, and its native value $f(\alpha^\star)$.
\end{enumerate}
\end{exercise}

\begin{solution}
Receipt: \texttt{ex23\_directions.py}.

\emph{Step 1: the directions.} The balance wall is $m_{\cE}=s+2t$, and the walls are $\kappa_S=m_S-t(|S|-1)$ (Exercise~\ref{ex:smallc-lean}, Step 1): $m_e>0$ and $m_e+m_{e'}>t$. There are $123$ such directions. The multiset $\{|c_e|\}$, with $\sum_ec_e=s-t$, has one value on $17$ of them, two on $70$ and three on $36$; the classes of equal $S$ and $\{|c_e|\}$ number $54$: $9$, $31$ and $14$. The $17$ directions with $|c_e|$ constant fall into nine classes:
\begin{center}
\small
\renewcommand{\arraystretch}{1.15}
\begin{tabular}{@{}l l l@{}}
\toprule
$|c_e|$, $S$ & directions $(\mvec;s,t)$ & $P_F(\mvec;s,t)$\\ \midrule
$0$, $2c'$ & $(c'\one;c',c')$, $c'=1,\dots,5$ & $\ell,\ \frac{\ell-2}3,\ \frac{10\ell-23}{54},\ \frac{62\ell-145}{486},\ \frac{568\ell-1331}{5832}$\\
$1$, $5$ & $(4,2,2;2,3)$, $(3,3,1;3,2)$ & $\frac{11-4\ell}{18}$\\
 & $(3,3,3;1,4)$, $(2,2,2;4,1)$ & $\frac{11-4\ell}{54}$\\
$1$, $7$ & $(5,3,3;3,4)$, $(4,4,2;4,3)$ & $\frac{19-8\ell}{54}$\\
 & $(4,4,4;2,5)$, $(3,3,3;5,2)$ & $\frac{19-8\ell}{108}$\\
$1$, $9$ & $(6,4,4;4,5)$, $(5,5,3;5,4)$ & $\frac{751-320\ell}{2916}$\\
$2$, $8$ & $(7,3,3;3,5)$, $(5,5,1;5,3)$ & $\frac{80\ell-181}{648}$\\
\bottomrule
\end{tabular}
\end{center}
the last column being found in Step 3.

\emph{Step 2: the classes.} The walls are positive, so Proposition~\ref{prop:dualfamily} and Corollary~\ref{cor:cycle} apply as in Exercise~\ref{ex:smallc-lean}, Step 6, at the data $(k\mvec;ks,kt)$, with $\nu_e=km_e-kt=kc_e$:
\[
\Gamma(sk)\Gamma(tk)\,P_F(k\mvec;ks,kt)=\int_0^\infty y^{tk-1}\prod_e\sG_{kc_e;(1,-1)}(1,y)\,dy .
\]
Proposition~\ref{prop:grammareflect} with $A=1$ and $\mathbf h=\frac12\log y$ gives $\sG_{\nu;(1,-1)}(1,y)=y^{\nu/2}\,\sG_{\nu;(1,-1)}(\sqrt y,\sqrt y)$, and with $A=-1$ it gives $\sG_{\nu;(1,-1)}(\xi_1,\xi_2)=\sG_{-\nu;(1,-1)}(\xi_2,\xi_1)$, so at equal coefficients the Gramma is even in $\nu$. With $\sum_ec_e=s-t$,
\[
\Gamma(sk)\Gamma(tk)\,P_F(k\mvec;ks,kt)=\int_0^\infty y^{\frac{Sk}2-1}\prod_e\sG_{k|c_e|;(1,-1)}(\sqrt y,\sqrt y)\,dy ,
\]
which depends on $S$ and $\{|c_e|\}$ only. So two directions $d,d'$ of one class have
\[
\frac{P_{d'}(k)}{P_d(k)}=\frac{\Gamma(sk)\Gamma(tk)}{\Gamma(s'k)\Gamma(t'k)}=\Gst21{sk,\ tk}{Sk}\Big/\Gst21{s'k,\ t'k}{Sk},
\]
two Beta corners at the same total, rational at the integers; for instance $\frac{\Gamma(4k)\Gamma(k)}{\Gamma(3k)\Gamma(2k)}=3,7,\frac{33}2,39$ at $k=1,2,3,4$ for $(4,2,2;2,3)$ against $(3,3,3;1,4)$. A class has one value span.

\emph{Step 3: the values, read on the centre ray.} Along the centre ray of Exercise~\ref{ex:smallc-bracket}, with $f_{\mathrm c}=\alpha^{\one}/(K\Psi)$, the insertion $\alpha^uK^{-a}\Psi^{-b}$ has the value $\langle\alpha^uK^{-a}\Psi^{-b}\rangle_k=P_F(k\one+u;k+a,k+b)$ (Definition~\ref{def:ray}). So
\[
P_F(k\one+2\mathbf e_1;k,k+1)=\langle\alpha_1^2\Psi^{-1}\rangle_k,\qquad P_F(3\one+4\mathbf e_1;3,5)=\langle\alpha_1^4\Psi^{-2}\rangle_3,
\]
and $(4,2,2;2,3)$, $(5,3,3;3,4)$, $(6,4,4;4,5)$ are $\langle\alpha_1^2\Psi^{-1}\rangle_k$ at $k=2,3,4$. Reduced in the boxes of the centre ray modulo $S_3$ and the Cremona, as in Exercise~\ref{ex:smallc-box} (the orbit averages have the same values, Proposition~\ref{prop:raysym}(ii)),
\begin{gather*}
\alpha_1^2\Psi^{-1}\equiv-\frac23+\frac{k-1}k\,b_1\quad\text{in }\mathfrak B_2(1,1),\\
\alpha_1^4\Psi^{-2}\equiv\frac{7k+5}{9(k+1)}-\frac{(k-1)(2k^2-6k-5)}{3k(k-2)(k+1)}\,b_1\quad\text{in }\mathfrak B_2(2,2),
\end{gather*}
and the relations of a box hold at every strength where the values converge (Proposition~\ref{prop:box-relations}(i)). With $P(K)=c(K)\ell+r(K)$ (Exercise~\ref{ex:smallc-bracket}, Step 8) and $\langle b_1\rangle_k=\frac16\big(P(k-1)-3P(k)\big)$, so that $\langle b_1\rangle_2=\frac13$, $\langle b_1\rangle_3=\frac{11-4\ell}{108}$ and $\langle b_1\rangle_4=\frac{19-8\ell}{243}$,
\begin{gather*}
P_F(4,2,2;2,3)=-\tfrac23\cdot\tfrac{\ell-2}3+\tfrac12\cdot\tfrac13=\tfrac{11-4\ell}{18},\qquad
P_F(5,3,3;3,4)=-\tfrac23\cdot\tfrac{10\ell-23}{54}+\tfrac23\cdot\tfrac{11-4\ell}{108}=\tfrac{19-8\ell}{54},\\
P_F(6,4,4;4,5)=-\tfrac23\cdot\tfrac{62\ell-145}{486}+\tfrac34\cdot\tfrac{19-8\ell}{243}=\tfrac{751-320\ell}{2916},\\
P_F(7,3,3;3,5)=\tfrac{13}{18}\cdot\tfrac{10\ell-23}{54}+\tfrac5{18}\cdot\tfrac{11-4\ell}{108}=\tfrac{80\ell-181}{648}.
\end{gather*}
The centre's class is the centre ray read at $c'k$: $P_F(c'\one;c',c')=P(c')=c(c')\ell+r(c')$. The other members follow by the ratios of Step 2 at $k=1$: $\frac{\Gamma(2)\Gamma(3)}{\Gamma(3)\Gamma(2)}=1$ for $(3,3,1;3,2)$, $\frac{\Gamma(2)\Gamma(3)}{\Gamma(1)\Gamma(4)}=\frac13$ for $(3,3,3;1,4)$ and $(2,2,2;4,1)$, $1$ for $(4,4,2;4,3)$, $\frac{\Gamma(3)\Gamma(4)}{\Gamma(2)\Gamma(5)}=\frac12$ for $(4,4,4;2,5)$ and $(3,3,3;5,2)$, and $1$ for $(5,5,3;5,4)$ and $(5,5,1;5,3)$. The value of $(2,2,2;4,1)$ agrees with Exercise~\ref{ex:smallc-lean}, where it was read on another insertion.

\emph{Step 4: the native values.} The native gauge on the orthant is unique (Theorem~\ref{thm:pencil}(i)), so a symmetry of the masses fixes it: for $m_2=m_3$ it is $(x,1,1)$, for $m_1=m_2$ it is $(1,1,x)$, with $x>0$ the root of the defect equation $m_e=s\,\alpha_e/K+t\,\alpha_e(K-\alpha_e)/\Psi$ at one mark. Exactly: $f(\alpha^\star)=9^{-c'}$ on the centre's class; $\frac1{243}$ at $\one$ for $(3,3,3;1,4)$ and $(2,2,2;4,1)$, and $\frac{64}{6561}$ at $(4,1,1)$ and $(1,1,\frac14)$ for $(4,2,2;2,3)$ and $(3,3,1;3,2)$; $\frac1{2187}$ at $\one$ for $(4,4,4;2,5)$ and $(3,3,3;5,2)$, and $\frac{3125}{3779136}$ at $(\frac52,1,1)$ and $(1,1,\frac25)$; $\frac1{12500}$ at $(2,1,1)$ and $(1,1,\frac12)$; $\frac{823543}{553584375}$ at $(7,1,1)$ and $(1,1,\frac17)$. Within a class the native values are in the ratio of the rates of the Beta corners, $\frac{64}{27}=\frac{4^4}{3^32^2}$ and $\frac{3125}{1728}=\frac{5^52^2}{4^43^3}$.

\emph{Check} (not part of the solution). The least recurrences of the closings $\CT(f^{-j})$, searched by linear algebra modulo the prime $2^{61}-1$ on $60$ closings to the order $4$ and the degree $24$, have the order $2$ on $16$ of the $17$ directions with $|c_e|$ constant; on the seventeenth, $(5\one;5,5)$, the closings are those of the centre at $5j$, whose recurrence of order $2$ has the degree $18$, found with $70$ closings. They have the order $3$ on $5$ directions and the order $4$ on $2$, all with two values of $|c_e|$, and none within these bounds on the other $100$.

\emph{Answer.} $123$ directions in $54$ classes; within a class $\Gamma(sk)\Gamma(tk)P$ is one function of $k$, and the values are in the ratio of two Beta corners at the same total; the $17$ directions with $|c_e|$ constant have the values of the table, all in $\mathbb Q+\mathbb Q\ell$, and the native values of Step 4.
\end{solution}

% Exercises of chapter XI.4, second half: the centres of four and five marks, W_3 at (1;0,2), and
% positivity against the gauge end. Appended to ch:exrays. No new macros.
% Receipts: /home/claude/v2/new/exercises4b/scripts/ (ex24a_halflattice.py ... ex27_positivity.py).

\section{The centre of \texorpdfstring{$\binom{[4]}1$}{([4] 1)}}\label{sec:exsb-four}

The list $\binom{[4]}1$ has the marks $1,2,3,4$ and the four singletons as its bases, every weight $1$, so that $K(\alpha)=\alpha_1+\alpha_2+\alpha_3+\alpha_4$ and $\Psi(\alpha)=\alpha^{[4]}K(\alpha^{-1})$. Its centre ray at the strength one (Definition~\ref{def:centreray}) has the stride $f=\big(K(\alpha)K(\alpha^{-1})\big)^{-1}$, so that $f^{-1}=\sum_{e,e'}\alpha_e/\alpha_{e'}$; the walls of the direction are the widths, $\kappa_{1,S}=r(S)+r([4]\setminus S)-1=1$ for every nonempty proper $S$, and the volume of the direction is $V_1=\sum_\pi\prod_{l=2}^4\kappa_{1,U_l}^{-1}=4!=24$ (Proposition~\ref{prop:principal}). Its values are $P(k)=\langle1\rangle_k=P_{\binom{[4]}1}(k\one;k,k)$, and throughout this section $v=P(1)$. The five exercises close $v$ on the half-lattice, derive the law of the values and its rows, write the row at every step in closed form, read the coefficient of the coordinate as a closing in two ways, and clear the rows.

\begin{exercise}[The closing on the half-lattice]\label{ex:cfour-halflattice}
Write $v=P(1)$ on its share side, merge the marks in pairs, and close the remaining share. Show that
\[
v=\Tx{3}{\Ilat{\frac12+\mathbb Z}}{0}=7\,\Tx3\IZ0 .
\]
\end{exercise}

\begin{solution}
Receipt: \texttt{ex24a\_halflattice.py}.

\emph{The share side.} The dual list $\binom{[4]}3$ has the bases $[4]-e$, and replacing the table by its dual exchanges the exponents (Proposition~\ref{prop:dualfamily}), which at the centre are equal: the two lists have the same values. Corollary~\ref{cor:cycle} applies to $\binom{[4]}3$ at $s=t=1$ with the masses $\one$: they lie on the balance wall, $m_{[4]}=4=3s+t$, they meet $m_S=|S|>s(|S|-1)$, and $\nu_e=m_e-s=0$. So
\[
v=\frac1{\Gamma(1)^2}\int_{a_1+a_2+a_3+a_4=1}\ \prod_{e=1}^4\Gamma(a_e)^2 ,
\]
a fibre integral of dimension $3$.

\emph{Merging the pairs.} Put $c=a_1+a_2$. Corollary~\ref{cor:fourgammas} with the entries $a=b=0$, $c=d=c$ and the variable $a_1$ merges the pair $\{1,2\}$,
\[
\frac1{2\pi i}\int\Gamma(a_1)^2\,\Gamma(c-a_1)^2\,da_1=\Gst41{c,\ c,\ c,\ c}{2c},
\]
and the same with the total $1-c$ merges $\{3,4\}$. One share is left:
\[
v=\frac1{2\pi i}\int_{(\frac12)}\Gst41{c,\ c,\ c,\ c}{2c}\,\Gst41{1-c,\ 1-c,\ 1-c,\ 1-c}{2-2c}\,dc .
\]

\emph{The integrand.} By the reflection and the shift (Proposition~\ref{prop:stacked}(iii),(v)), $\Gamma(c)^4\Gamma(1-c)^4=\pi^4/\sin^4\pi c$ and $\Gamma(2c)\Gamma(2-2c)=(1-2c)\,\Gamma(2c)\Gamma(1-2c)=(1-2c)\,\pi/\sin2\pi c$, so the integrand is
\[
w(c)=\frac{\pi^4}{\sin^4\pi c}\cdot\frac{\sin2\pi c}{(1-2c)\,\pi}=\frac{2\pi^3\cos\pi c}{(1-2c)\sin^3\pi c}.
\]
At $c=\frac12$ the zero of $\cos\pi c$ cancels that of $1-2c$, so the poles of $w$ are the integers. The factor $2\pi^3\cos\pi c/\sin^3\pi c$ is $-\partial_c\Gst40{c,\ c,\ 1-c,\ 1-c}{}=2\sum_{m\in\mathbb Z}(c+m)^{-3}$ (Lemma~\ref{lem:resonantpair}), of period one, and at $c=m+\varepsilon$
\[
\frac{2\pi^3\cos\pi\varepsilon}{\sin^3\pi\varepsilon}=\frac2{\varepsilon^3}-\frac{2\pi^4}{15}\,\varepsilon+O(\varepsilon^3):
\]
no terms in $\varepsilon^{-2}$, $\varepsilon^{-1}$ or $\varepsilon^{0}$.

\emph{The closing.} At $c=N+\frac12+iy$, $|\cos\pi c|=|\sinh\pi y|$ and $|\sin\pi c|=\cosh\pi y$, so $|w(c)|\le4\pi^3e^{-2\pi|y|}/(2N)$: the integrals over the lines $\re c=N+\frac12$ tend to zero, and so do those over the horizontal segments $\im c=\pm Y$ as $Y\to\infty$. The line closes to the right (Theorem~\ref{thm:closing}) on the triple poles at the integers $m\ge1$, with the residues
\[
\operatorname*{Res}_{c=m}w=2\,[\varepsilon^2]\,\frac1{1-2m-2\varepsilon}=\frac{2\cdot4}{(1-2m)^3}=-\frac{8}{(2m-1)^3},
\]
and closing to the right takes minus their sum:
\[
v=\sum_{m\ge1}\frac{8}{(2m-1)^3}=\sum_{x\in\frac12+\mathbb Z,\ x>0}x^{-3}=\Tx{3}{\Ilat{\frac12+\mathbb Z}}{0}
\]
(Definition~\ref{def:lattice}).

\emph{The coset relation.} Theorem~\ref{thm:latticefree}(i) at $q=2$ and $b=0$, with the kernel $x^{-3}$ and the cut $0$, assembles $\frac12\mathbb Z$ from its two cosets, the rescaled kernel being $(x/2)^{-3}=8x^{-3}$:
\[
\Tx3\IZ0+\Tx{3}{\Ilat{\frac12+\mathbb Z}}{0}=8\,\Tx3\IZ0 .
\]

So $v=\Tx{3}{\Ilat{\frac12+\mathbb Z}}{0}=7\,\Tx3\IZ0$, which is $8.4143983221171599978\ldots$ to $19$ places; the receipt evaluates the contour integral by quadrature as a check. The coset $\frac12+\mathbb Z$ is the doubled share $2c$ of a merged pair, the lower entry of $\Gst41{c,\ c,\ c,\ c}{2c}$: the word has one lower Gamma of slope $2$ and one of slope $-2$, which pair off (Proposition~\ref{prop:ra-paired}), and its rational factor $(1-2c)^{-1}$ has its pole on that coset.
\end{solution}

\begin{exercise}[The law, the low-end relation and the rows]\label{ex:cfour-law}
\begin{enumerate}[label=(\roman*)]
\item With $\vartheta=y\frac{d}{dy}$, compute the operators $\mathbf R_0,\dots,\mathbf R_5$ of Theorem~\ref{thm:ri-law} for $n=4$, and from them the law of the values, $\sum_{k=0}^2c_{4,k}(s)\,P(s+k)=0$.
\item Give the closings $\mathsf C_j(1)$ for $j\le8$, and check the law read at $k=-j$ on them.
\item Give the symbol, its roots and the native gauges they belong to.
\item Write the low-end relation.
\item Write the rows $P(k)=a_kv-b_k$ for $k\le6$ (Definition~\ref{def:row}), and the value span of the ray.
\end{enumerate}
\end{exercise}

\begin{solution}
Receipt: \texttt{ex24b\_law.py}.

(i) \emph{One mark.} At the strength $s$ the masses are $s\one$ and the exponents $(s,s)$, and Corollary~\ref{cor:cycle}, as in Exercise~\ref{ex:cfour-halflattice}, gives the share side in its first form,
\[
\Gamma(s)^2P(s)=\int_0^\infty y^{s-1}\,\mathsf g(y)^4\,dy,\qquad \mathsf g(y)=\sG_{0;(1,-1)}(1,y),\quad \re s>0 .
\]
The Gramma $\mathsf g$ satisfies $\vartheta^2\mathsf g=y\,\mathsf g$. Write $\mathsf g_\nu=\sG_{\nu;(1,-1)}(1,y)$. The ladder (Theorem~\ref{thm:grammalaws}(i)) in the coefficient $y$ of the exponent $-1$ gives $\vartheta\mathsf g_\nu=-y\,\mathsf g_{\nu-1}$, so $\vartheta^2\mathsf g_0=-y\,\mathsf g_{-1}+y^2\mathsf g_{-2}$; the exchange relation (Theorem~\ref{thm:grammalaws}(ii)) at $\nu=-1$ reads $-\mathsf g_{-1}=\mathsf g_0-y\,\mathsf g_{-2}$, and substituting $y\,\mathsf g_{-2}=\mathsf g_0+\mathsf g_{-1}$ gives $\vartheta^2\mathsf g_0=y\,\mathsf g_0$.

\emph{The operators.} Theorem~\ref{thm:ri-law}(i) at $n=4$ has $\mathbf R_0=1$, $\mathbf R_1=\vartheta$ and $\mathbf R_{i+1}=\vartheta\mathbf R_i-i(5-i)\,y\,\mathbf R_{i-1}$, with the coefficients $i(5-i)=4,6,6,4$; an operator is written $\sum_ky^kq(\vartheta)$, and $\vartheta\,y^kq(\vartheta)=y^k(\vartheta+k)\,q(\vartheta)$. So
\begin{align*}
\mathbf R_2&=\vartheta^2-4y,\\
\mathbf R_3&=\vartheta^3-2y(5\vartheta+2),\\
\mathbf R_4&=\vartheta^4-2y(8\vartheta^2+7\vartheta+2)+24y^2,\\
\mathbf R_5&=\vartheta^5-2y(2\vartheta+1)(5\vartheta^2+5\vartheta+2)+64y^2(\vartheta+1),
\end{align*}
and $\mathbf Q_4=\mathbf R_5$ annihilates every product of four solutions of $\vartheta^2u=yu$, $\mathsf g^4$ among them; the receipt checks it on the basis $u^{4-i}(\vartheta u)^i$. The leading coefficients $1,-20,64$ of $q_{4,0},q_{4,1},q_{4,2}$ are those of $\xi(\xi^2-16\eta)(\xi^2-4\eta)=\xi^5-20\eta\xi^3+64\eta^2\xi$, as Theorem~\ref{thm:ri-law}(i) says.

\emph{The law.} Integrated against $y^{s-1}$, a term $y^kq(\vartheta)$ becomes $q(-s-k)$ times the transform at $s+k$: integration by parts turns $\vartheta$ into $-s$, the boundary terms vanishing because $\mathsf g$ is $O(|\log y|)$ at $0$ (Lemma~\ref{lem:onemark}(i)) and at most $\sqrt\pi\,y^{-1/4}e^{-2\sqrt y}$ for $y\ge1$ (Lemma~\ref{lem:onemark}(ii)). With $\Gamma(s+k)^2=\Gamma(s)^2(s)_k^2$ and a division by $\Gamma(s)^2s^2$ this is Theorem~\ref{thm:ri-law}(ii):
\begin{gather*}
c_{4,0}(s)=-s^3,\qquad c_{4,1}(s)=q_{4,1}(-s-1)=2(2s+1)(5s^2+5s+2),\\
c_{4,2}(s)=q_{4,2}(-s-2)\,(s+1)^2=-64(s+1)^3,
\end{gather*}
\[
-s^3P(s)+2(2s+1)(5s^2+5s+2)\,P(s+1)-64(s+1)^3P(s+2)=0
\]
identically in $s$. At $s=1$ the coefficients are $-1$, $72$, $-512$.

(ii) \emph{The closings.} By Corollary~\ref{cor:reflected} $\mathsf C_j(1)$ is the constant term of $f^{-j}=K(\alpha)^jK(\alpha^{-1})^j$; the multinomial expansions of the two factors pair the monomial $\alpha^v$ with $\alpha^{-v}$, so
\[
\mathsf C_j(1)=\sum_{|v|=j}\binom jv^2=1,\ 4,\ 28,\ 256,\ 2716,\ 31504,\ 387136,\ 4951552,\ 65218204\qquad(j=0,\dots,8).
\]
The law read at $k=-j$ (Corollary~\ref{cor:reflected}, Theorem~\ref{thm:ri-law}(iii)) is $\sum_{i\le\min(2,j)}c_{4,i}(-j)\,\mathsf C_{j-i}(1)=0$: at $j=1$, $1\cdot4+(-4)\cdot1=0$; at $j=2$, $8\cdot28+(-72)\cdot4+64\cdot1=0$; the receipt checks it for $j\le40$, and also that $\mathbf Q_4$ annihilates $\mathsf h^4=\sum_j\mathsf C_j(1)\,y^j/j!^2$, $\mathsf h(y)=\sum_my^m/m!^2$, through $y^{40}$.

(iii) \emph{The symbol.} The coefficients of $s^3$ in $c_{4,0},c_{4,1},c_{4,2}$ are $-1,20,-64$, so the symbol (Corollary~\ref{cor:symbol}) is
\[
\chi_4(\lambda)=-1+20\lambda-64\lambda^2=-(1-4\lambda)(1-16\lambda),
\]
with the roots $\frac1{16}$ and $\frac14$. By Proposition~\ref{prop:ri-signs}(ii) the native gauges of the ray are the signs $\eta$ with $K(\eta)\neq0$, of native value $K(\eta)^{-2}$: $\eta=\one$, with $K=4$ and the native value $\frac1{16}=f(\alpha^\star)$, $\alpha^\star=\one$ being the native gauge of the direction (Theorem~\ref{thm:pencil}); and the signs with one or three marks at $-1$, with $K=\pm2$ and the native value $\frac14$. The signs with two marks at $-1$ have $K=0$ and lie off $X_F$.

(iv) \emph{The low-end relation.} By Corollary~\ref{cor:rays-lowend}(ii), with $[s^3]c_{4,0}=-1$, $\mathsf L_0=V_1=24$, $c_{4,1}(0)=4$ and $c_{4,2}(0)=-64$,
\[
-24+4P(1)-64P(2)=0,\qquad P(2)=\frac{P(1)-6}{16}.
\]

(v) \emph{The rows.} $P(1)=v$ and $P(2)=\frac1{16}v-\frac38$, and the law at $s=k\ge1$, whose last coefficient $c_{4,2}(k)$ does not vanish, gives $P(k+2)=-\big(c_{4,0}(k)P(k)+c_{4,1}(k)P(k+1)\big)/c_{4,2}(k)$; at $k=1$, $P(3)=\big(72P(2)-P(1)\big)/512=\frac7{1024}v-\frac{27}{512}$. The rows $(a_k;b_k)$:
\[
\begin{array}{c|cccccc}
k&1&2&3&4&5&6\\\hline
a_k&1&\frac1{16}&\frac7{1024}&\frac1{1024}&\frac{679}{4194304}&\frac{1969}{67108864}\\[3pt]
b_k&0&\frac38&\frac{27}{512}&\frac{37}{4608}&\frac{25555}{18874368}&\frac{9304913}{37748736000}
\end{array}
\]
So every value lies in $\mathbb Q+\mathbb Q\,v$, $v=\Tx{3}{\Ilat{\frac12+\mathbb Z}}{0}$ (Exercise~\ref{ex:cfour-halflattice}), and this is the value span of the ray (Corollary~\ref{cor:ri-span}(i), with $h=2$): the order $h=2$ leaves one value free after the low-end relation. As renderings, $P(2)=0.15089989513232249986\ldots$, $P(3)=0.0047859260300977734224\ldots$ and $P(4)=0.00018767197255365017146\ldots$; the receipt checks $P(2)$ against the quadrature of $\int_0^\infty y\,\mathsf g(y)^4dy$.
\end{solution}

\begin{exercise}[The row at every step]\label{ex:cfour-row}
For $n\ge0$ put $p_l=l+\frac12$ and $j=n-l$ for $0\le l\le n$,
\[
c_l=\binom{2l}l\binom{2j}j,\qquad S_l=\frac{c_l}{\binom nl}=\frac{(2l)!\,(2j)!}{l!\,j!\,n!},\qquad w_l=S_l^2c_l,\qquad F(\sigma)=\binom{\sigma-1}{n}^{\!2},
\]
$F[\sigma,p]=\big(F(\sigma)-F(p)\big)/(\sigma-p)$, and for a polynomial $g$
\[
\mathcal J[g]=\frac1{2\pi i}\int_{(\frac12)}g(\sigma)\sum_{m\in\mathbb Z}\frac{d\sigma}{(\sigma+m)^3}.
\]
Show that
\[
256^n\,P(n+1)=\sum_{l=0}^nw_l\,\Tx{3}{\Ilat{\frac12+\mathbb Z}}{l}-16^n\sum_{l=0}^nc_l\,\mathcal J\big[F[\,\cdot\,,p_l]\big],
\]
evaluate $\mathcal J$ on the binomial polynomials $\binom\sigma i$, and read the rows of $P(2)$ and $P(3)$ off the formula.
\end{exercise}

\begin{solution}
Receipt: \texttt{ex24c\_row.py}.

(a) \emph{The word.} At the strength $k=n+1$ the share side of Exercise~\ref{ex:cfour-halflattice}, with the shares of total $k$ and the pairs merged by Corollary~\ref{cor:fourgammas} at the totals $\sigma$ and $k-\sigma$, is
\[
P(k)=\frac1{\Gamma(k)^2}\,\frac1{2\pi i}\int\frac{\Gamma(\sigma)^4\,\Gamma(k-\sigma)^4}{\Gamma(2\sigma)\,\Gamma(2k-2\sigma)}\,d\sigma\qquad(0<\re\sigma<k).
\]
With $\Gamma(k-\sigma)=(1-\sigma)_n\Gamma(1-\sigma)$ and $\Gamma(2k-2\sigma)=(2-2\sigma)_{2n}(1-2\sigma)\Gamma(1-2\sigma)$ (Proposition~\ref{prop:stacked}(iii)), the reflections of Exercise~\ref{ex:cfour-halflattice} give the factor $2\pi^3\cos\pi\sigma/\sin^3\pi\sigma=2\sum_m(\sigma+m)^{-3}$ against $(1-\sigma)_n^4/\big((1-2\sigma)(2-2\sigma)_{2n}\big)$. Splitting $(2-2\sigma)_{2n}$ into its factors $2(1-\sigma+t)$ and $2(\frac32-\sigma+t)$, $0\le t<n$, gives $(1-2\sigma)(2-2\sigma)_{2n}=2\cdot4^n(\frac12-\sigma)_{n+1}(1-\sigma)_n$, so
\[
P(n+1)=\frac1{4^n\,n!^2}\,\frac1{2\pi i}\int\rho(\sigma)\sum_m\frac{d\sigma}{(\sigma+m)^3},\qquad \rho(\sigma)=\frac{(1-\sigma)_n^{\,3}}{(\frac12-\sigma)_{n+1}} .
\]
The triple poles of the kernel at $1,\dots,n$ meet the triple zeros of $\rho$, and the simple poles of $\rho$ at the $p_l$ meet the zeros of the kernel, so the integrand is regular on $0<\re\sigma<k$ and the line may be $\re\sigma=\frac12$.

(b) \emph{Partial fractions.} With $\Omega(\sigma)=\prod_{i=1}^n(\sigma-i)=n!\binom{\sigma-1}n$, $(1-\sigma)_n=(-1)^n\Omega$ and $(\frac12-\sigma)_{n+1}=(-1)^{n+1}\prod_{i=0}^n(\sigma-p_i)$, so $\rho/\Omega^2=-\Omega/\prod_i(\sigma-p_i)$, whose residue at $p_l$ is $-\Omega(p_l)/\prod_{i\neq l}(l-i)$, with $\prod_{i\neq l}(l-i)=(-1)^jl!\,j!$ and
\[
\Omega(p_l)=\prod_{t=0}^{l-1}\big(t+\tfrac12\big)\prod_{t=1}^{j}\big(\tfrac12-t\big)=\frac{(-1)^j(2l-1)!!\,(2j-1)!!}{2^n}.
\]
As $(2l-1)!!/l!=\binom{2l}l/2^l$, the residue is $-c_l/4^n$:
\[
\frac{\rho}{\Omega^2}=-\frac1{4^n}\sum_{l=0}^n\frac{c_l}{\sigma-p_l}.
\]
Multiplying by $\Omega^2=n!^2F$ and writing $\Omega(\sigma)^2=\Omega(p_l)^2+(\sigma-p_l)\,n!^2F[\sigma,p_l]$ splits $\rho$ into a polynomial and its poles,
\[
\rho=\mathcal P+\sum_{l=0}^n\frac{r_l}{\sigma-p_l},\qquad \mathcal P=-\frac{n!^2}{4^n}\sum_lc_l\,F[\sigma,p_l],\qquad r_l=-\frac{c_l\,\Omega(p_l)^2}{4^n}=-\frac{n!^2\,w_l}{64^n},
\]
the last because $(2l-1)!!\,(2j-1)!!=S_l\,n!/2^n$. The receipt checks the partial fractions and the $r_l$ for $n\le3$.

(c) \emph{The pole terms.} The kernel is $\pi^3\cos\pi\sigma/\sin^3\pi\sigma$, with the principal part $(\sigma-i)^{-3}$ and no terms of orders $-2,-1,0$ at every integer $i$ (Exercise~\ref{ex:cfour-halflattice}), so the residue of $(\sigma-p_l)^{-1}$ times the kernel at $i$ is $\frac12\partial^2_\sigma(\sigma-p_l)^{-1}|_{\sigma=i}=(i-p_l)^{-3}$; at $p_l$ the kernel vanishes. On the lines $\re\sigma=N+\frac12$ the integrand is $O(N^{-1}e^{-2\pi|\im\sigma|})$, so the line closes to the right (Theorem~\ref{thm:closing}):
\[
\frac1{2\pi i}\int_{(\frac12)}\frac1{\sigma-p_l}\sum_m\frac{d\sigma}{(\sigma+m)^3}=-\sum_{i\ge1}\frac1{(i-p_l)^3}=-\Tx{3}{\Ilat{\frac12+\mathbb Z}}{l},
\]
the terms $i\le l$, at the negative points $i-p_l$, cancelling the head $\Th{3}{\Ilat{\frac12+\mathbb Z}}{0}{l}$ of the terms at the positive points (Proposition~\ref{prop:cutlaw}). The polynomial part gives $\mathcal J[\mathcal P]$.

(d) \emph{Assembling.} By (a) to (c),
\[
P(n+1)=\frac1{4^nn!^2}\Big(-\frac{n!^2}{4^n}\sum_lc_l\,\mathcal J\big[F[\cdot,p_l]\big]+\frac{n!^2}{64^n}\sum_lw_l\,\Tx{3}{\Ilat{\frac12+\mathbb Z}}{l}\Big),
\]
which is the stated formula after multiplication by $256^n$. With $\Tx{3}{\Ilat{\frac12+\mathbb Z}}{l}=v-\Th{3}{\Ilat{\frac12+\mathbb Z}}{0}{l}$ the row of $P(n+1)$ is
\[
a_{n+1}=\frac1{256^n}\sum_{l=0}^nw_l,\qquad b_{n+1}=\frac1{256^n}\Big(\sum_{l=0}^nw_l\,\Th{3}{\Ilat{\frac12+\mathbb Z}}{0}{l}+16^n\sum_{l=0}^nc_l\,\mathcal J\big[F[\cdot,p_l]\big]\Big).
\]

(e) \emph{The functional $\mathcal J$.} By Lemma~\ref{lem:resonantpair}, $\sum_m(\sigma+m)^{-3}=-\frac12\partial_\sigma\Gst40{\sigma,\ \sigma,\ 1-\sigma,\ 1-\sigma}{}$, and on the line the stacked Gamma is $2\pi\,p_{\mathrm{crit}}(y)$, $\sigma=\frac12+iy$, decaying like $e^{-2\pi|y|}$ (Definition~\ref{def:pairlaw}). Integrating by parts,
\[
\mathcal J[g]=\frac12\,\oE\big[g'(X)\big],\qquad X=\tfrac12+iY,\quad Y\sim\mathbb P_{\mathrm{crit}} .
\]
By Theorem~\ref{thm:sharelaw}(iii), $\oE[e^{tX}]=te^t/(e^t-1)$ for $|t|<2\pi$, so $\oE[X^m]=1,\frac12,\frac16,0,-\frac1{30},0,\frac1{42},0,-\frac1{30}$ for $m=0,\dots,8$, and at $t=\log(1+z)$, $\oE[(1+z)^X]=(1+z)\log(1+z)/z$. Since $\sum_i\binom\sigma iz^i=(1+z)^\sigma$, whose derivative in $\sigma$ is $\log(1+z)\,(1+z)^\sigma$,
\[
\sum_{i\ge0}\mathcal J\Big[\binom\sigma i\Big]z^i=\frac12\log(1+z)\,\oE\big[(1+z)^X\big]=\frac{(1+z)\log^2(1+z)}{2z}.
\]
With $\log^2(1+z)=2\sum_{m\ge2}(-1)^m\Th{1}{\IZ}{0}{m-1}\,z^m/m$, the head $\Th{1}{\IZ}{0}{m}=\sum_{r=1}^m\frac1r$ (Definition~\ref{def:lattice}), the coefficients are
\[
\lambda_0=0,\qquad \lambda_i=\mathcal J\Big[\binom\sigma i\Big]=(-1)^i\,\frac{\Th{1}{\IZ}{0}{i-1}-1}{i(i+1)}\quad(i\ge1),
\]
that is $0,\frac12,0,-\frac1{24},\frac1{24},-\frac{13}{360},\frac{11}{360},-\frac{29}{1120},\frac{223}{10080}$ for $i=0,\dots,8$; the receipt checks them against the moments for $i<40$. A polynomial written in the binomial basis, $g=\sum_i(\Delta^ig)(0)\binom\sigma i$, has $\mathcal J[g]=\sum_i(\Delta^ig)(0)\,\lambda_i$.

(f) \emph{The first steps.} At $n=1$: $c_0=c_1=2$, $S_0=S_1=2$, $w_0=w_1=8$; $F(\sigma)=(\sigma-1)^2$, $F[\sigma,p]=\sigma+p-2$, so $F[\sigma,\frac12]=\sigma-\frac32$ and $F[\sigma,\frac32]=\sigma-\frac12$, each with $\mathcal J=\frac12$; with $\Th{3}{\Ilat{\frac12+\mathbb Z}}{0}{1}=8$,
\[
256\,P(2)=16v-8\cdot8-16\Big(2\cdot\frac12+2\cdot\frac12\Big)=16v-96,\qquad P(2)=\frac1{16}v-\frac38 .
\]
At $n=2$: $(c_0,c_1,c_2)=(6,4,6)$, $(S_0,S_1,S_2)=(6,2,6)$, $(w_0,w_1,w_2)=(216,16,216)$, and
\begin{align*}
F[\sigma,\tfrac12]&=\tfrac14\sigma^3-\tfrac{11}8\sigma^2+\tfrac{41}{16}\sigma-\tfrac{55}{32},&\mathcal J&=\tfrac{21}{32},\\
F[\sigma,\tfrac32]&=\tfrac14\sigma^3-\tfrac98\sigma^2+\tfrac{25}{16}\sigma-\tfrac{21}{32},&\mathcal J&=\tfrac9{32},\\
F[\sigma,\tfrac52]&=\tfrac14\sigma^3-\tfrac78\sigma^2+\tfrac{17}{16}\sigma-\tfrac{11}{32},&\mathcal J&=\tfrac5{32},
\end{align*}
for instance $\mathcal J\big[F[\cdot,\frac12]\big]=\frac12\oE\big[\frac34X^2-\frac{11}4X+\frac{41}{16}\big]=\frac12\big(\frac18-\frac{11}8+\frac{41}{16}\big)=\frac{21}{32}$. With the heads $\Th{3}{\Ilat{\frac12+\mathbb Z}}{0}{1}=8$ and $\Th{3}{\Ilat{\frac12+\mathbb Z}}{0}{2}=8+\frac8{27}$,
\[
256^2P(3)=448\,v-\Big(16\cdot8+216\cdot\frac{224}{27}\Big)-256\Big(6\cdot\frac{21}{32}+4\cdot\frac9{32}+6\cdot\frac5{32}\Big)=448\,v-1920-1536 ,
\]
so $P(3)=\frac7{1024}v-\frac{27}{512}$. Both rows are those of Exercise~\ref{ex:cfour-law}.

So the row of $P(n+1)$ at every step is $a_{n+1}=256^{-n}\sum_lw_l$ and $b_{n+1}$ as in (d), with $\mathcal J$ read through the $\lambda_i$. For $n=0,1,2,3$ the sums are $\sum_lw_l=1,16,448,16384$, $256^nb_{n+1}=0,96,3456,\frac{1212416}9$ and $\sum_lc_l\mathcal J[F[\cdot,p_l]]=0,2,6,\frac{853}{54}$; the receipt checks that the formula reproduces the rows of the law for $n\le40$.
\end{solution}

\begin{exercise}[The coefficient of the coordinate is a closing]\label{ex:cfour-closing}
Show that $a_k=\mathsf C_{k-1}(1)/64^{k-1}$ for every $k\ge1$, and give the denominator of $a_k$:
\begin{enumerate}[label=(\roman*)]
\item from the law of Exercise~\ref{ex:cfour-law}, by Theorem~\ref{thm:ra-closingcoordinate};
\item directly from the row of Exercise~\ref{ex:cfour-row}, by a telescoping identity in the summation index.
\end{enumerate}
\end{exercise}

\begin{solution}
Receipt: \texttt{ex24d\_closing.py}.

(i) \emph{The law is reflexive.} With $h=2$, Definition~\ref{def:ra-reflexive} asks for $c_{4,2-i}(-k-1)=\gamma\lambda^ic_{4,i}(k)$, $i=0,1,2$. Here
\begin{gather*}
c_{4,2}(-k-1)=64k^3=-64\,c_{4,0}(k),\qquad c_{4,0}(-k-1)=(k+1)^3,\\
c_{4,1}(-k-1)=-2(2k+1)(5k^2+5k+2)=-c_{4,1}(k),
\end{gather*}
so $\gamma=-64$, $\gamma\lambda=-1$ gives the ratio $\lambda=\frac1{64}$, and $\gamma\lambda^2c_{4,2}(k)=-\frac{64}{4096}\cdot(-64)(k+1)^3=(k+1)^3$ checks the third. The ray has the origin $0$, every letter exponent positive and positive walls, as Proposition~\ref{prop:ra-lowend} asks. By Theorem~\ref{thm:ra-closingcoordinate}(i), $y_k=64^{1-k}\mathsf C_{k-1}(1)$ solves the law on $k\ge1$ with the low-end coefficient $-\big(c_{4,1}(0)y_1+c_{4,2}(0)y_2\big)=-\big(4\cdot1-64\cdot\frac4{64}\big)=0$, and since $c_{4,2}(k)\neq0$ for $k\ge0$, Theorem~\ref{thm:ra-closingcoordinate}(ii) gives
\[
P(k)=\frac{\mathsf C_{k-1}(1)}{64^{k-1}}\,P(1)+r_k,\qquad a_k=\frac{\mathsf C_{k-1}(1)}{64^{k-1}} .
\]
By Theorem~\ref{thm:ra-closingcoordinate}(iii), with $\lambda=\frac1{64}$, the denominator of $a_k$ is $64^{k-1}/\gcd\big(64^{k-1},\mathsf C_{k-1}(1)\big)$:
\[
\begin{array}{c|cccccccc}
k&1&2&3&4&5&6&7&8\\\hline
\mathsf C_{k-1}(1)&1&4&28&256&2716&31504&387136&4951552\\
v_2\big(\mathsf C_{k-1}(1)\big)&0&2&2&8&2&4&6&9\\
\log_2\operatorname{den}(a_k)&0&4&10&10&22&26&30&33
\end{array}
\]
with $v_2$ the exponent of $2$.

(ii) \emph{The closings across the pairs.} Split $f^{-1}=(X+Y)(X'+Y')$ with $X=\alpha_1+\alpha_2$, $X'=\alpha_1^{-1}+\alpha_2^{-1}$, and $Y,Y'$ the same on $\{3,4\}$. In $(X+Y)^n(X'+Y')^n=\sum_{l,l'}\binom nl\binom n{l'}X^lY^{n-l}X'^{l'}Y'^{n-l'}$ a term has a constant term in $\alpha_1,\alpha_2$ only if $l=l'$, $X^l$ and $X'^{l'}$ being homogeneous of degrees $l$ and $-l'$, and $\CT\big((XX')^l\big)=\sum_i\binom li^2=\binom{2l}l$. So
\[
\mathsf C_n(1)=\sum_{l=0}^n\binom nl^2\binom{2l}l\binom{2j}j=:U(n),\qquad j=n-l .
\]
Put $T(n)=4^{-n}\sum_lw_l$, so that $a_{n+1}=4^nT(n)/256^n$ by Exercise~\ref{ex:cfour-row}. It remains to show $T=U$.

\emph{The telescoping identities.} Both $T$ and $U$ satisfy
\[
(n+2)^3X_{n+2}-2(2n+3)(5n^2+15n+12)\,X_{n+1}+64(n+1)^3X_n=0\qquad(n\ge0),
\]
which is the law of the closings of Exercise~\ref{ex:cfour-law}(ii) at $j=n+2$: $c_{4,0}(-n-2)=(n+2)^3$, $c_{4,1}(-n-2)=-2(2n+3)(5n^2+15n+12)$, $c_{4,2}(-n-2)=64(n+1)^3$. Write $t(n,l)$ for the summand of either sum, $t_T(n,l)=4^{-n}S_l^2c_l$ and $t_U(n,l)=\binom nl^2c_l$, zero outside $0\le l\le n$. Their ratios are rational:
\begin{align*}
\frac{t_T(n,l+1)}{t_T(n,l)}&=\frac{(2l+1)^3\,j}{(l+1)(2j-1)^3},&\frac{t_T(n+1,l)}{t_T(n,l)}&=\frac{2(2j+1)^3}{(j+1)(n+1)^2},\\
\frac{t_U(n,l+1)}{t_U(n,l)}&=\frac{j^3(2l+1)}{(l+1)^3(2j-1)},&\frac{t_U(n+1,l)}{t_U(n,l)}&=\frac{2(n+1)^2(2j+1)}{(j+1)^3}.
\end{align*}
With $G=R\,t$ and the rational functions
\[
R_T=\frac{4l\,(2j+1)^3\,\kappa_T}{(n+1)^2(j+1)(j+2)},\qquad R_U=-\frac{4(n+1)^2\,l^3\,(2j+1)\,\kappa_U}{(j+1)^3(j+2)^3},
\]
\begin{align*}
\kappa_T&=12j^3+10lj^2+2l^2j+59j^2+33lj+4l^2+94j+26l+49,\\
\kappa_U&=12j^3+10lj^2+2l^2j+62j^2+35lj+3l^2+104j+30l+56,
\end{align*}
each summand satisfies, for $0\le l\le n-1$,
\[
64(n+1)^3t(n,l)-2(2n+3)(5n^2+15n+12)\,t(n+1,l)+(n+2)^3t(n+2,l)=G(n,l+1)-G(n,l).
\]
Divided by $t(n,l)$ and written with the ratios this is an identity of rational functions of $(n,l)$; cleared of denominators it is an identity of polynomials, which the receipt verifies coefficient by coefficient for both $T$ and $U$. Summing it over $0\le l\le n-1$, with $G(n,0)=0$ (the factor $l$), leaves $G(n,n)$, and the full sum over $l$ of the left side is $G(n,n)$ plus the left side at $l=n,n+1,n+2$, where only the summands with $l\le n+1$ at $n+1$ and $l\le n+2$ at $n+2$ survive. Divided by $t(n,n)$, with $t(n+1,n)/t(n,n)$, $t(n+1,n+1)/t(n,n)$, $t(n+2,n)/t(n,n)$, $t(n+2,n+1)/t(n,n)$ and $t(n+2,n+2)/t(n,n)$ read off the ratios, this boundary sum is a rational function of $n$ that vanishes identically (receipt). So both sums satisfy the recurrence; $T(0)=U(0)=1$, $T(1)=\frac14(8+8)=4=U(1)$, and the leading coefficient $(n+2)^3$ does not vanish, so $T=U$.

Hence $\sum_lw_l=4^n\mathsf C_n(1)$ and $a_{n+1}=4^n\mathsf C_n(1)/256^n=\mathsf C_n(1)/64^n$, as in (i). The coefficient of the coordinate $v$ in $P(k)$ is the closing $\mathsf C_{k-1}(1)$ over $64^{k-1}$, with no lattice sum: $64^{k-1}a_k\in\mathbb Z$, and the denominator of $a_k$ is $64^{k-1}/\gcd\big(64^{k-1},\mathsf C_{k-1}(1)\big)$.
\end{solution}

\begin{exercise}[The clearing]\label{ex:cfour-clearing}
Show that for every $k\ge1$ the part of the denominator of $b_k$ prime to $2$ divides $d_{k-1}^{\,3}$, and that $256^{k-1}d_{k-1}^{\,3}\,b_k\in\mathbb Z$; so $256^{k-1}d_{k-1}^{\,3}$ clears the row of $P(k)$. Compare it with the bound of the frame.
\end{exercise}

\begin{solution}
Receipt: \texttt{ex24e\_clearing.py}.

Write $n=k-1$. By Exercise~\ref{ex:cfour-row}(d), $256^nb_{n+1}$ is the sum of the heads and the moments,
\[
\sum_lw_l\,\Th{3}{\Ilat{\frac12+\mathbb Z}}{0}{l}\qquad\text{and}\qquad16^n\sum_lc_l\,\mathcal J\big[F[\cdot,p_l]\big];
\]
and by Exercise~\ref{ex:cfour-closing} $64^na_{n+1}=\mathsf C_n(1)\in\mathbb Z$. Two facts from outside the handbook are used: Legendre's formula, $v_q(m!)=\sum_{t\ge1}\lfloor m/q^t\rfloor$ for a prime $q$, and Kummer's theorem, that $v_q\binom{a+b}a$ is the number of carries when $a$ and $b$ are added in base $q$; $v_q$ is the exponent of $q$.

\emph{The integers.} $S_l\in\mathbb Z$: by Legendre's formula, $v_q(S_l)=\sum_{t\ge1}\big(\lfloor2l/q^t\rfloor+\lfloor2j/q^t\rfloor-\lfloor l/q^t\rfloor-\lfloor j/q^t\rfloor-\lfloor n/q^t\rfloor\big)$, and every term is nonnegative by $\lfloor2x\rfloor+\lfloor2y\rfloor\ge\lfloor x\rfloor+\lfloor y\rfloor+\lfloor x+y\rfloor$ (with $x=a+\alpha$, $y=b+\beta$, $0\le\alpha,\beta<1$, it reads $\lfloor2\alpha\rfloor+\lfloor2\beta\rfloor\ge\lfloor\alpha+\beta\rfloor$, true since $\alpha+\beta\ge1$ forces $\max(\alpha,\beta)\ge\frac12$). So $w_l=S_l^2c_l\in\mathbb Z$.

Fix an odd prime $q$ and $e\ge1$ with $q^{e-1}\le n<q^e$, so that $v_q(d_n^{\,3})=3(e-1)$; it suffices that both sums have $v_q\ge-3(e-1)$.

\emph{The heads.} $\Th{3}{\Ilat{\frac12+\mathbb Z}}{0}{l}=8\sum_{i=1}^l(2i-1)^{-3}$, and $2l-1<2n<2q^e$, so the only odd number up to $2l-1$ with $v_q\ge e$ is $q^e$, when $q^e\le2l-1$. In that case, at the level $q^e$ of Legendre's formula, $\lfloor l/q^e\rfloor=\lfloor j/q^e\rfloor=\lfloor n/q^e\rfloor=0$ and $\lfloor2l/q^e\rfloor=1$, so $v_q(S_l)\ge1$ and $v_q\binom{2l}l\ge1$, the other levels contributing $\ge0$; then $v_q(w_l)\ge3$, which pays for $(q^e)^{-3}$. Otherwise the head already has $v_q\ge-3(e-1)$.

\emph{The moments.} Vandermonde's identity twice, $\binom\sigma m=\sum_t\binom{\sigma-p}t\binom p{m-t}$ and $\binom{\sigma-p-1}{t-1}=\sum_i\binom\sigma i\binom{-p-1}{t-1-i}$, with $\binom{\sigma-p}t=\frac{\sigma-p}t\binom{\sigma-p-1}{t-1}$, gives the binomial coefficients of the divided difference,
\[
\big(\Delta^iF[\cdot,p]\big)(0)=\sum_mf_m\sum_{t=i+1}^m\frac1t\binom p{m-t}\binom{-p-1}{t-1-i},\qquad F=\sum_mf_m\binom\sigma m .
\]
Here $f_m=\sum_{i,i'\le n}(-1)^{i+i'}\binom mi\binom i{m-i'}$, from $\binom{\sigma-1}n=\sum_{i\le n}(-1)^{n-i}\binom\sigma i$ and $\binom\sigma i\binom\sigma{i'}=\sum_m\binom mi\binom i{m-i'}\binom\sigma m$; so $f_m\in\mathbb Z$, and $q\mid f_m$ for $m\ge q^e$: a nonzero term has $i<q^e$ and $m-i\le i'<q^e$, and adding $i$ and $m-i$ in base $q$ carries into the place $q^e$, so $q\mid\binom mi$ (Kummer's theorem). At $p=p_l$, a $q$-adic integer, the binomials of $p$ and of $-p-1$ are $q$-integral. Since $t\le2n<2q^e$, the only $t$ with $v_q(t)\ge e$ is $t=q^e$, where $m\ge q^e$ and $q\mid f_m$. So $v_q\big((\Delta^iF[\cdot,p_l])(0)\big)\ge1-e$ for every $i$, and $\ge2-e$ for $i\ge q^e$, where every $t$ and $m$ exceed $q^e$.

For $1\le i\le2n-1$: $v_q(\lambda_i)\ge2-2e$ for $i<q^e$ and $v_q(\lambda_i)\ge1-2e$ for $i\ge q^e$. The heads of the counting input at index one have $v_q\big(\Th{1}{\IZ}{0}{m}\big)\ge1-e$ for $m<q^e$ and $\ge-e$ for $q^e\le m<2q^e$, while
\[
v_q\big(\Th{1}{\IZ}{0}{q^e-1}\big)\ge2-e:
\]
grouping the $r<q^e$ by $v_q(r)=s$, $\Th{1}{\IZ}{0}{q^e-1}=\sum_{s<e}q^{-s}\sum_{r'<q^{e-s},\ q\nmid r'}\frac1{r'}$, and $\sum_{r'<q^a,\ q\nmid r'}\frac1{r'}$ has $v_q\ge a$, pairing $r'$ with $q^a-r'$, $\frac1{r'}+\frac1{q^a-r'}=\frac{q^a}{r'(q^a-r')}$; each term has $v_q\ge e-2s\ge2-e$. The integers $i$ and $i+1$ are coprime and below $2q^e$, so $v_q\big(i(i+1)\big)\le e-1$ unless $q^e\in\{i,i+1\}$, where it is $e$. For $i<q^e-1$ this gives $v_q(\lambda_i)\ge(1-e)-(e-1)$. At $i=q^e-1$, $\Th{1}{\IZ}{0}{q^e-2}-1=\Th{1}{\IZ}{0}{q^e-1}-\frac{q^e}{q^e-1}$ has $v_q\ge\min(2-e,e)$, so $v_q(\lambda_i)\ge\min(2-e,e)-e\ge2-2e$. At $i=q^e$, $\Th{1}{\IZ}{0}{q^e-1}-1$ has $v_q\ge\min(2-e,0)$, so $v_q(\lambda_i)\ge\min(2-e,0)-e\ge1-2e$; and for $q^e<i\le2n-1$, $v_q(\lambda_i)\ge-e-(e-1)$. In each case the product of the coefficient and $\lambda_i$ has $v_q\ge3-3e$, and $16^n$ and $c_l$ are $q$-integral. The receipt checks every one of these estimates for $n\le30$ and every odd prime $q\le2n$; for instance $v_3\big(\Th{1}{\IZ}{0}{26}\big)=-1$, $v_5\big(\Th{1}{\IZ}{0}{24}\big)=1$.

\emph{The prime $2$.} For $p\in\frac12+\mathbb Z$, $v_2\binom pr=-r-v_2(r!)\ge1-2r$ for $r\ge1$; the product of the two binomials in the coefficient, of total order $m-1-i\le2n-1$, has $v_2\ge3-4n$; $v_2(1/t)\ge-\lfloor\log_22n\rfloor$; and $v_2(\lambda_i)\ge-2\lfloor\log_22n\rfloor$, from $v_2\big(\Th{1}{\IZ}{0}{i-1}\big)\ge-\lfloor\log_2(i-1)\rfloor$ and $v_2(i(i+1))\le\lfloor\log_2(i+1)\rfloor$. So each term of $16^n\sum_lc_l\mathcal J[F[\cdot,p_l]]$ has $v_2\ge-3\lfloor\log_2n\rfloor=-v_2(d_n^{\,3})$, and the heads, $w_l$ times $8$ times odd denominators, are $2$-integral. Hence $256^nd_n^{\,3}b_{n+1}\in\mathbb Z$.

\emph{The numbers.} The receipt checks both statements for $k\le80$. The first denominators, with their parts at $2$:
\[
\begin{array}{c|ccccccc}
k&2&3&4&5&6&7&8\\\hline
\operatorname{den}(b_k)&2^3&2^9&2^9\cdot9&2^{21}\cdot9&2^{25}\cdot1125&2^{29}\cdot125&2^{32}\cdot42875\\
d_{k-1}^{\,3}&1&2^3&2^3\cdot27&2^6\cdot27&2^6\cdot3375&2^6\cdot3375&2^6\cdot1157625
\end{array}
\]

\emph{The frame.} The sequences $(a_k)$ and $(b_k)$ both solve the law of Exercise~\ref{ex:cfour-law}: $(a_k)$ by Exercise~\ref{ex:cfour-closing}(i), and $b_k=a_kv-P(k)$ as a difference of two solutions. They are independent, $(a_1,a_2)=(1,\frac1{16})$ and $(b_1,b_2)=(0,\frac38)$. Their frame (Proposition~\ref{prop:frame}) $C_k=a_kb_{k+1}-a_{k+1}b_k$ has $C_1=\frac38$ and $C_{k+1}=\frac{c_{4,0}(k)}{c_{4,2}(k)}C_k=\frac{k^3}{64(k+1)^3}C_k$, so
\[
C_k=\frac{3}{8\cdot64^{k-1}k^3}.
\]
By Theorem~\ref{thm:ra-frameclearing}(i), every integer $D_k$ that clears the rows of $P(k)$ and $P(k+1)$ has $D_k^2\ge|C_k|^{-1}=\frac83\,64^{k-1}k^3$, so $\liminf_k\frac1k\log D_k\ge\log8$. The clearing found grows like $256^k$ times $d_{k-1}^{\,3}$; its merit (Definition~\ref{def:ra-merit}) is at most $-\log f(\alpha^\star)-\log256=\log16-\log256=-\log16$, so Theorem~\ref{thm:ra-merit}(ii) does not apply to these rows. For every clearing the frame bounds the merit by $\log16-\log8=\log2$, and Corollary~\ref{cor:ra-meritbound}, with the ratio $\frac1{64}$, by $-\log4+\liminf_k\frac1k\log\gcd\big(64^{k-1},\mathsf C_{k-1}(1)\big)$.

So the row of $P(k)$ is cleared by $256^{k-1}d_{k-1}^{\,3}$: the coefficient $a_k$ by $64^{k-1}$, the rational part $b_k$ by $256^{k-1}$ at the prime $2$ and by $d_{k-1}^{\,3}$ at the odd primes.
\end{solution}

\section{The centre of \texorpdfstring{$\binom{[5]}1$}{([5] 1)}}\label{sec:exsb-five}

The list $\binom{[5]}1$ has the marks $1,\dots,5$ and the five singletons as its bases, every weight $1$. Its centre ray at the strength one has the stride $f=\big(K(\alpha)K(\alpha^{-1})\big)^{-1}$, $K(\alpha)=\sum_{e=1}^5\alpha_e$, the walls $\kappa_{1,S}=1$ for every nonempty proper $S$, the volume of the direction $V_1=5!=120$, and the values $P(k)=\langle1\rangle_k=P_{\binom{[5]}1}(k\one;k,k)$. In this section $P_4$ denotes the values of the centre ray of $\binom{[4]}1$, continued in the strength, and $v=P_4(1)=7\,\Tx3\IZ0$ (Section~\ref{sec:exsb-four}); $\mathsf g(y)=\sG_{0;(1,-1)}(1,y)$ and $\mathsf h(y)=\sum_{m\ge0}y^m/m!^2$, both solutions of $\vartheta^2u=yu$, $\vartheta=y\frac d{dy}$ (Exercise~\ref{ex:cfour-law}; $\vartheta^2y^m=m^2y^m$ gives $\vartheta^2\mathsf h=y\,\mathsf h$).

\begin{exercise}[The law, the low-end relation and the value span]\label{ex:cfive-law}
\begin{enumerate}[label=(\roman*)]
\item Compute the operators $\mathbf R_2,\dots,\mathbf R_6$ of Theorem~\ref{thm:ri-law} for $n=5$ and the law of the values $\sum_{k=0}^3c_{5,k}(s)\,P(s+k)=0$.
\item Give the closings $\mathsf C_j(1)$ for $j\le6$, the symbol, its roots and their native gauges.
\item Write the low-end relation and the value span, and write $\Gamma(s)^2P(s)$ for $s=3,4$ in $1$, $P(1)$ and $P(2)$.
\end{enumerate}
\end{exercise}

\begin{solution}
Receipt: \texttt{ex25a\_law.py}.

(i) As in Exercise~\ref{ex:cfour-law}, Corollary~\ref{cor:cycle} and Proposition~\ref{prop:dualfamily} give $\Gamma(s)^2P(s)=\int_0^\infty y^{s-1}\mathsf g(y)^5dy$. Theorem~\ref{thm:ri-law}(i) at $n=5$ has the coefficients $i(6-i)=5,8,9,8,5$ in $\mathbf R_{i+1}=\vartheta\mathbf R_i-i(6-i)\,y\,\mathbf R_{i-1}$:
\begin{align*}
\mathbf R_2&=\vartheta^2-5y,\\
\mathbf R_3&=\vartheta^3-y(13\vartheta+5),\\
\mathbf R_4&=\vartheta^4-y(22\vartheta^2+18\vartheta+5)+45y^2,\\
\mathbf R_5&=\vartheta^5-y(30\vartheta^3+40\vartheta^2+23\vartheta+5)+y^2(149\vartheta+130),\\
\mathbf R_6&=\vartheta^6-y(35\vartheta^4+70\vartheta^3+63\vartheta^2+28\vartheta+5)+y^2(259\vartheta^2+518\vartheta+285)-225y^3 .
\end{align*}
The leading coefficients $1,-35,259,-225$ are those of $(\xi^2-25\eta)(\xi^2-9\eta)(\xi^2-\eta)$. Theorem~\ref{thm:ri-law}(ii) gives
\begin{gather*}
c_{5,0}(s)=s^4,\qquad c_{5,1}(s)=-(35s^4+70s^3+63s^2+28s+5),\\
c_{5,2}(s)=(s+1)^2(259s^2+518s+285),\qquad c_{5,3}(s)=-225(s+1)^2(s+2)^2,
\end{gather*}
and the law $\sum_{k=0}^3c_{5,k}(s)P(s+k)=0$ identically in $s$; every $c_{5,k}$ has degree $4$ and $c_{5,k}(-s-k)=c_{5,k}(s)$. At $s=1$ and $s=2$ the coefficients are $(1,-201,4248,-8100)$ and $(16,-1433,21213,-32400)$.

(ii) The closings are $\mathsf C_j(1)=\sum_{|v|=j}\binom jv^2=1,5,45,545,7885,127905,2241225$ for $j\le6$, and they satisfy $\sum_{i\le\min(3,j)}c_{5,i}(-j)\,\mathsf C_{j-i}(1)=0$ (Corollary~\ref{cor:reflected}; the receipt checks $j\le40$). The coefficients of $s^4$ give the symbol
\[
\chi_5(\lambda)=1-35\lambda+259\lambda^2-225\lambda^3=(1-\lambda)(1-9\lambda)(1-25\lambda),
\]
with the roots $\frac1{25},\frac19,1$. By Proposition~\ref{prop:ri-signs}(ii) they are the native values $K(\eta)^{-2}$ of the signs: $\eta=\one$, with $K=5$, the native gauge of the direction; the signs with one or four marks at $-1$, $K=\pm3$; and those with two or three marks at $-1$, $K=\pm1$.

(iii) \emph{The low-end relation.} By Corollary~\ref{cor:rays-lowend}(ii), with $[s^4]c_{5,0}=1$, $\mathsf L_0=V_1=120$ and $\big(c_{5,1}(0),c_{5,2}(0),c_{5,3}(0)\big)=(-5,285,-900)$,
\[
120-5P(1)+285P(2)-900P(3)=0,\qquad P(3)=\frac2{15}-\frac1{180}P(1)+\frac{19}{60}P(2).
\]
\emph{The value span.} The law has order $3$ and its last coefficient does not vanish at the positive integers, so every value is a rational combination of $1$, $P(1)$ and $P(2)$: the value span of the ray is $\mathbb Q+\mathbb Q\,P(1)+\mathbb Q\,P(2)$ (Corollary~\ref{cor:ri-span}(i), $h=3$). The law at $s=1$ gives $P(4)=\big(P(1)-201P(2)+4248P(3)\big)/8100$, so
\begin{align*}
\Gamma(3)^2P(3)&=\frac8{15}-\frac1{45}P(1)+\frac{19}{15}P(2),\\
\Gamma(4)^2P(4)&=\frac{944}{375}-\frac{113}{1125}P(1)+\frac{1907}{375}P(2).
\end{align*}
As renderings, $P(1)=39.94558801199624569894\ldots$, $P(2)=0.2946961074182751313903\ldots$ and $P(3)=0.004733833949141315501685\ldots$ (quadrature of the share side, a check; Exercise~\ref{ex:cfive-gauge} reads them exactly from the gauge end).
\end{solution}

\begin{exercise}[The exact pairing and the Hessian]\label{ex:cfive-pairing}
Show that for every integer $\sigma\ge1$ the concomitant
\[
B(\sigma)=\sum_{i=1}^3\sum_{j=1}^ic_{5,i}(\sigma-j)\,\mathsf C_{\sigma-j}(1)\,P(\sigma-j+i),
\]
with the closings at negative indices read as $0$, equals $-120$. Deduce $\det H$, $H$ the Hessian of $-\log f$ at $\one$ on the chart $\alpha_5=1$, and compute $H$ directly.
\end{exercise}

\begin{solution}
Receipt: \texttt{ex25b\_pairing.py}.

\emph{The adjoint.} The law is self-adjoint under the reflection of the strength: $c_{5,i}(-\sigma-i)=c_{5,i}(\sigma)$ for $i=0,\dots,3$ (Exercise~\ref{ex:cfive-law}), so $c_{5,i}(\sigma-i)=c_{5,i}(-\sigma)$, and $\sum_ic_{5,i}(\sigma-i)\,\mathsf C_{\sigma-i}(1)=\sum_ic_{5,i}(-\sigma)\,\mathsf C_{\sigma-i}(1)$ is the law of the closings at $j=\sigma$, which vanishes for $\sigma\ge0$; for $\sigma<0$ every closing in it is $0$. The closings solve the adjoint recurrence (Proposition~\ref{prop:ri-pairing}(i)).

\emph{The concomitant is constant.} Shifting $j$ by one in $B(\sigma+1)$,
\[
B(\sigma+1)-B(\sigma)=\mathsf C_\sigma(1)\sum_{i=1}^3c_{5,i}(\sigma)P(\sigma+i)-P(\sigma)\sum_{i=1}^3c_{5,i}(\sigma-i)\,\mathsf C_{\sigma-i}(1),
\]
and by the law and the adjoint recurrence the two sums are $-c_{5,0}(\sigma)P(\sigma)$ and $-c_{5,0}(\sigma)\mathsf C_\sigma(1)$, so the difference vanishes (the receipt checks the identity symbolically). At $\sigma=1$ only $j=1$ survives, and by the low-end relation of Exercise~\ref{ex:cfive-law}
\[
B(1)=-5P(1)+285P(2)-900P(3)=-120 .
\]
At $\sigma=2$, with $\mathsf C_1(1)=5$, the five terms are $-201\cdot5\,P(2)$, $4248\cdot5\,P(3)$, $285\,P(2)$, $-8100\cdot5\,P(4)$ and $-900\,P(3)$, so $B(2)=-720P(2)+20340P(3)-40500P(4)$; with the rows of Exercise~\ref{ex:cfive-law},
\[
20340P(3)=2712-113P(1)+6441P(2),\qquad 40500P(4)=2832-113P(1)+5721P(2),
\]
and $B(2)=-120$, the coefficients of $P(1)$ and $P(2)$ cancelling. The receipt checks $B(\sigma)=-120$ exactly in the value span for $\sigma\le50$. This is Corollary~\ref{cor:ri-pairing} at $n=5$, with $(-1)^nn!=-120$.

\emph{The limit.} Every $c_{5,i}$ has degree $4=n-1$, $\chi_5(0)=1$, and $\lambda_\star=\frac1{25}$ is a simple root. By Proposition~\ref{prop:ri-pairing}(iii),
\begin{gather*}
V_1\det H=-\frac{\lambda_\star\chi_5'(\lambda_\star)}{\chi_5(0)}=\prod_{c\neq\star}\Big(1-\frac{\lambda_\star}{\lambda_c}\Big):\\
\chi_5'\big(\tfrac1{25}\big)=-25\cdot\frac{24}{25}\cdot\frac{16}{25}=-\frac{384}{25},\qquad \Big(1-\frac9{25}\Big)\Big(1-\frac1{25}\Big)=\frac{384}{625},
\end{gather*}
so $120\det H=\frac{384}{625}$ and $\det H=\frac{16}{3125}=\frac{2^4}{5^5}$.

\emph{Directly.} By Theorem~\ref{thm:branch} at the centre, $H=2\big(\kappa_{\one}(X_e,X_{e'})\big)_{e,e'\neq5}$, the covariances of the indicators under the weight law at $\one$, which picks one mark of five with the chance $\frac15$: $\operatorname{Var}X_e=\frac15-\frac1{25}$ and $\operatorname{Cov}(X_e,X_{e'})=-\frac1{25}$. So
\[
H=2\Big(\frac15I-\frac1{25}J\Big)=\frac1{25}\begin{pmatrix}8&-2&-2&-2\\-2&8&-2&-2\\-2&-2&8&-2\\-2&-2&-2&8\end{pmatrix},
\]
$J$ the matrix of ones, with the eigenvalues $\frac25$ on the vectors of sum zero and $\frac2{25}$ on $\one$: $\det H=\big(\frac25\big)^3\frac2{25}=\frac{16}{3125}$, as Proposition~\ref{prop:ri-signs}(ii) gives at $\eta=\one$, $(2\sigma)^{n-1}/K(\one)^n=2^4/5^5$.

So $B(\sigma)=-120$ at every strength, the low-end relation carried to every $\sigma$, and its limit at high strength gives $\det H=2^4/5^5$, the value of the Hessian computed directly.
\end{solution}

\begin{exercise}[The lowered settings and the frame]\label{ex:cfive-lowered}
For $j=0,1,2$ the lowered setting with $j$ marks lowered is $\xi_{5,j}(s)=\Gamma(s)^{-2}\int_0^\infty y^{s-1}\mathsf g(y)^{5-j}\mathsf h(y)^j\,dy$ (Proposition~\ref{prop:ri-lowered}); put $x^{(j)}(s)=\Gamma(s)^2\xi_{5,j}(s)$, so that $x^{(0)}(s)=\Gamma(s)^2P(s)$.
\begin{enumerate}[label=(\roman*)]
\item Show that the $x^{(j)}$ satisfy the law of Exercise~\ref{ex:cfive-law} in this normalization, and find their low-end coefficients and their rates.
\item Write the frame $\Phi(s)=\det\big[x^{(i)}(s+l)\big]_{i,l=0}^2$ in terms of $\Phi(1)$, and find $\Phi(1)$ from the behaviour of the $x^{(j)}$ at large $s$.
\item Deduce the value of $x^{(1)}(1)\,x^{(2)}(2)-x^{(1)}(2)\,x^{(2)}(1)$.
\end{enumerate}
\end{exercise}

\begin{solution}
Receipt: \texttt{ex25c\_lowered.py}.

(i) \emph{The law.} $\mathsf g^{5-j}\mathsf h^j$ is a product of five solutions of $\vartheta^2u=yu$, so $\mathbf Q_5=\mathbf R_6$ annihilates it (Theorem~\ref{thm:ri-law}(i)), and integration against $y^{s-1}$ turns it into the law for $x(s)=\Gamma(s)^2\xi(s)$:
\[
s^6x(s)-(35s^4+70s^3+63s^2+28s+5)\,x(s+1)+(259s^2+518s+285)\,x(s+2)-225\,x(s+3)=0,
\]
the coefficients being $q_{5,k}(-s-k)$; this is the law of Exercise~\ref{ex:cfive-law} multiplied by $\Gamma(s)^2s^2$ (Proposition~\ref{prop:ri-lowered}(i)). The boundary terms vanish: near $0$, $\mathsf g^{5-j}\mathsf h^j=O(|\log y|^{5-j})$ (Lemma~\ref{lem:onemark}(i)), and for $y\ge1$, $\mathsf g\le\sqrt\pi\,y^{-1/4}e^{-2\sqrt y}$ (Lemma~\ref{lem:onemark}(ii)) and $\mathsf h\le e^{2\sqrt y}$, so the integrand decays like $e^{-2(5-2j)\sqrt y}$, $j<\frac52$.

\emph{The low-end coefficients.} Near $y=0$, $\mathsf g(y)=-\log y+2(\log\Gamma)'(1)+O(y|\log y|)$ (Lemma~\ref{lem:onemark}(i)) and $\mathsf h(y)=1+O(y)$, and $\int_0^1y^{s-1}(-\log y)^{5-j}dy=\Gamma(6-j)\,s^{j-6}$; so $x^{(j)}(s)=\Gamma(6-j)\,s^{j-6}\big(1+O(s)\big)$ and $s^4\xi_{5,j}(s)\to120,0,0$ for $j=0,1,2$, $\Gamma(s)^{-2}$ being $s^2(1+O(s))$. Read at $s\to0$, the law gives
\[
-5x^{(j)}(1)+285x^{(j)}(2)-225x^{(j)}(3)=-120,\ 0,\ 0\qquad(j=0,1,2):
\]
the lowered settings with $j\ge1$ have no low-end coefficient (Proposition~\ref{prop:ri-lowered}(ii)).

\emph{The rates.} $x^{(j)}(s+1)/\big(s^2x^{(j)}(s)\big)\to(5-2j)^{-2}=\frac1{25},\frac19,1$ (Proposition~\ref{prop:ri-lowered}(iii)): the setting with $j$ marks lowered carries the native value of the signs with $j$ marks at $-1$. This also follows from the amplitudes below.

(ii) \emph{The frame.} The end coefficients of the law in $x$ are $s^6$ and $-225$, with no zero at the integers $s\ge1$, so Proposition~\ref{prop:frame} gives $\Phi(s+1)=(-1)^3\frac{s^6}{-225}\Phi(s)=\frac{s^6}{225}\Phi(s)$ and
\[
\Phi(s)=\Phi(1)\,\Gamma(s)^6\,225^{1-s}.
\]
\emph{One mark at large $y$.} By the torus homogeneity (Theorem~\ref{thm:grammalaws}(iii)) $\mathsf g(y)=\sG_{0;(1,-1)}(\lambda,\lambda)$ with $\lambda=\sqrt y$, and Theorem~\ref{thm:configgauge}(iii) at $t=(1,-1)$, $\boldsymbol\xi=(1,1)$, $\nu=\nu'=0$ applies: $\Phi_{\boldsymbol\xi}(y')=e^{y'}+e^{-y'}$, $y^\star=0$, $\Phi_{\boldsymbol\xi}(0)=2$, $\sum_kw^\star_kt_k^2=2$, so
\[
\mathsf g(y)=e^{-2\sqrt y}\Big(\frac{2\pi}{\sqrt y}\Big)^{\frac12}2^{-\frac12}\big(1+O(y^{-\frac12})\big)=\sqrt\pi\,y^{-\frac14}e^{-2\sqrt y}\big(1+O(y^{-\frac12})\big).
\]
$\mathsf h(y)$ is the constant term of $e^{\sqrt y(z+z^{-1})}$, that is $\frac1{2\pi}\int_{-\pi}^{\pi}e^{2\sqrt y\cos\theta}d\theta$, and Laplace's method (an outside result; it applies because $\cos\theta$ has one maximum on the circle, at $\theta=0$, where its second derivative is $-1$) gives $\mathsf h(y)=e^{2\sqrt y}(4\pi\sqrt y)^{-1/2}\big(1+O(y^{-1/2})\big)$. So, with $\rho_j=2(5-2j)=10,6,2$,
\[
\mathsf g^{5-j}\mathsf h^j=\pi^{\frac{5-2j}2}2^{-j}\,y^{-\frac54}e^{-\rho_j\sqrt y}\big(1+O(y^{-\frac12})\big).
\]
\emph{The amplitudes.} With $y=u^2$, $x^{(j)}(s)=2\int_0^\infty u^{2s-1}\mathsf g(u^2)^{5-j}\mathsf h(u^2)^j\,du$; the part $u\le u_0$ is $O(u_0^{2s})$, negligible against $\Gamma(2s)\rho_j^{-2s}$, and the rest gives
\[
x^{(j)}(s)=A_j\,\Gamma\big(2s-\tfrac52\big)\,\rho_j^{-2s}\big(1+O(s^{-1})\big),\qquad A_j=2\pi^{\frac{5-2j}2}2^{-j}\rho_j^{\frac52},
\]
$A_0=200\sqrt{10}\,\pi^{5/2}$, $A_1=36\sqrt6\,\pi^{3/2}$, $A_2=2\sqrt2\,\pi^{1/2}$, $\prod_jA_j=\pi^{9/2}\,120^{5/2}$. Factoring $A_i\rho_i^{-2s}$ out of the rows and $\Gamma(2s+2l-\frac52)$ out of the columns,
\[
\Phi(s)=\prod_iA_i\rho_i^{-2s}\prod_{l=0}^2\Gamma\big(2s+2l-\tfrac52\big)\big(V+O(s^{-1})\big),\qquad V=\det\big[u_i^{\,l}\big]=\prod_{i<i'}(u_{i'}-u_i)=\frac{16}{16875},
\]
$u_i=\rho_i^{-2}=\frac1{100},\frac1{36},\frac14$. Theorem~\ref{thm:configgauge}(iii) for $\Gamma(\lambda+a)=\lambda^{\lambda+a}\sG_{\lambda+a;(1)}(\lambda)$, with $\Phi_{\xi}(y')=e^{y'}-y'$, $y^\star=0$ and $w^\star=1$, gives $\Gamma(\lambda+a)=\lambda^{\lambda+a}e^{-\lambda}(2\pi/\lambda)^{1/2}\big(1+O(\lambda^{-1})\big)$, hence $\prod_l\Gamma(2s+2l-\frac52)=\Gamma(2s)^3(2s)^{-3/2}\big(1+O(s^{-1})\big)$ and $\Gamma(2s)=2^{2s-1}\pi^{-1/2}s^{1/2}\Gamma(s)^2\big(1+O(s^{-1})\big)$, so
\[
\Phi(s)=\prod_jA_j\cdot V\cdot2^{-\frac92}\pi^{-\frac32}\,\Gamma(s)^6\Big(\frac{64}{\prod_j\rho_j^2}\Big)^s\big(1+O(s^{-1})\big),\qquad \frac{64}{\prod_j\rho_j^2}=\frac{64}{14400}=\frac1{225}.
\]
Comparing with $\Phi(s)=225\,\Phi(1)\,\Gamma(s)^6\,225^{-s}$,
\[
\Phi(1)=\frac{\pi^{9/2}\,120^{5/2}}{225}\cdot\frac{16}{16875}\cdot\frac{1}{2^{9/2}\pi^{3/2}}=\frac{128\sqrt{15}}{16875}\,\pi^3=\frac{128\sqrt{15}}{16875}\,\Gamma\big(\tfrac12\big)^6,
\]
which is $0.9108805647982030328903972\ldots$ to $25$ places; the receipt computes the nine settings by quadrature and reproduces it, and checks the law, the low-end coefficients and the amplitudes.

(iii) \emph{The block.} By (i), $x^{(j)}(3)=\big(-5x^{(j)}(1)+285x^{(j)}(2)+120\,\delta_{j0}\big)/225$, so subtracting $-\frac5{225}$ times the first column and $\frac{285}{225}$ times the second from the third leaves the column $(\frac{120}{225},0,0)$:
\[
\Phi(1)=\frac{120}{225}\Big(x^{(1)}(1)\,x^{(2)}(2)-x^{(1)}(2)\,x^{(2)}(1)\Big),\qquad x^{(1)}(1)\,x^{(2)}(2)-x^{(1)}(2)\,x^{(2)}(1)=\frac{16\sqrt{15}}{1125}\,\Gamma\big(\tfrac12\big)^6,
\]
which is $1.707901058996630686669495\ldots$ to $24$ places. The values $P(1),P(2)$ do not enter $\Phi$: their coefficients are minors that the vanishing low-end coefficients of the lowered settings remove.
\end{solution}

\begin{exercise}[The reading from the gauge end]\label{ex:cfive-gauge}
Read $P(1)$, $P(2)$ and $P(3)$ from the gauge end (Theorem~\ref{thm:gaugeend}, Definition~\ref{def:sl-gaugeread}): the native gauge and the constant $g(0)$; the density operator of the law and its exponents at $f^\star=\frac1{25}$; the germs that extend by zero; the Taylor coefficients of the germ; the values at a large strength and the law run back; and the low-end relation read on the output.
\end{exercise}

\begin{solution}
Receipt: \texttt{ex25d\_gauge.py}.

(1) \emph{The native gauge.} $\alpha^\star=\one$ (Theorem~\ref{thm:pencil}), $f^\star=f(\one)=\frac1{25}$, $\det H=\frac{16}{3125}$ (Exercise~\ref{ex:cfive-pairing}) and $(\omega_0/dy)(\one)=1$, the origin being $0$. By Proposition~\ref{prop:sl-ends}(ii), with $n=5$,
\[
g(0)=\frac1{\Gamma(2)\sqrt{\det H}}\Big(\frac{2\pi}{f^\star}\Big)^2=\frac{25\sqrt5}4\,(50\pi)^2=15625\sqrt5\,\Gamma\big(\tfrac12\big)^4 .
\]

(2) \emph{The operator.} The density operator of the law $\mathbf p=(c_{5,0},\dots,c_{5,3})$ is $\mathfrak D_{\mathbf p}=\sum_it^ic_{5,i}(-1-i-\theta)$, $\theta=t\frac d{dt}$ (Definition~\ref{def:sl-densityop}). With $\theta^m=\sum_jS(m,j)\,t^j\partial_t^j$, $S(m,j)$ the number of partitions of a set of $m$ elements into $j$ nonempty parts, $\mathfrak D_{\mathbf p}=\sum_{j=0}^4a_j(t)\partial_t^j$ with
\begin{align*}
a_0&=1-201t+4248t^2-8100t^3,\\
a_1&=-t\,(24300t^3-16965t^2+1232t-15),\\
a_2&=-t^2(16650t^3-14271t^2+1358t-25),\\
a_3&=-2t^3(1800t^3-1813t^2+210t-5),\\
a_4&=-t^4(t-1)(9t-1)(25t-1)=t^4\chi_5(t).
\end{align*}
In $s=\frac1{25}-t$, $\mathfrak D_{\mathbf p}=\sum_jb_j(s)\partial_s^j$ with $b_{4,0}=0$ and $b_{4,1}=-(f^\star)^4\chi_5'(f^\star)=\frac{384}{25^5}$. The second symbol is $\chi_1(\lambda)=-70\lambda+1036\lambda^2-1350\lambda^3$, and Theorem~\ref{thm:gaugeend}(i) gives
\[
\gamma=-1+\frac{\chi_1(f^\star)}{f^\star\chi_5'(f^\star)}=-1+\frac{-768/625}{-384/625}=1=\frac{n-3}2,\qquad R_0(x)=\frac{384}{25^5}\,x(x-1)^2(x-2):
\]
the exponents at $f^\star$ are $0,1,2$ and $\gamma=1$, a double root. The further polynomials $R_m(x)=\sum_jb_{j,m-3+j}[x]_j$ are
\begin{align*}
R_1&=-\tfrac{8}{25^4}\,x(x-1)(163x^2-67x+72),\qquad R_2=\tfrac1{25^3}\,x(1385x^3+1840x^2+1631x+576),\\
R_3&=-\tfrac1{625}(180x^4+1290x^3+2129x^2+1619x+476),\\
R_4&=-\tfrac1{25}(490x^4+2530x^3+5454x^2+5789x+2499),\\
R_5&=(x+2)^2(196x^2+802x+819),\qquad R_6=225(x+2)^2(x+3)^2,
\end{align*}
and $R_m=0$ for $m\ge7$.

(3) \emph{The line.} As $\gamma=1\in\{0,1,2\}$, the germ is $\varrho_f=\sum_ru_rs^r$ with $u_0=0$, $u_1=g(0)$ (Theorem~\ref{thm:gaugeend}(iv)). The germs that extend by zero are fixed by the conditions $\mathsf p_r(u)=0$, $r\le2$, the point masses of $\mathfrak D_{\mathbf p}$ applied to the germ cut at $f^\star$:
\[
\mathsf p_0(u)=-\frac{384}{25^5}u_0,\qquad \mathsf p_1(u)=-\frac{576}{25^4}u_0,\qquad \mathsf p_2(u)=\frac{576}{25^3}u_0-\frac{1344}{25^4}u_1+\frac{768}{25^5}u_2,
\]
the coefficient of $u_r$ in $\mathsf p_r$ being $r!\,b_{4,1}(r-\gamma)$. They leave one line, $u_0=0$ and $u_2=\frac{1344\cdot25}{768}u_1=\frac{175}4u_1$ (Theorem~\ref{thm:gaugeend}(ii)): the germ vanishes to the first order, as the exponent $\gamma=1$ says, and the receipt also derives the three conditions from the distributional derivatives of the cut germ.

(4) \emph{The germ.} For $N\ge3$, $R_0(N)u_N=-\sum_{m=1}^{\min(N,6)}R_m(N-m)\,u_{N-m}$. At $N=3$, $R_0(3)=\frac{4608}{25^5}$, $R_1(2)=-\frac{9440}{25^4}$, $R_2(1)=\frac{5432}{25^3}$, so
\[
u_3=-\frac{25^5}{4608}\Big(-\frac{9440}{25^4}\cdot\frac{175}4+\frac{5432}{25^3}\Big)u_1=\frac{48125}{32}u_1,\qquad u_4=\frac{35703125}{768}u_1 ,
\]
and the receipt continues to $u_{400}$, all rational multiples of $u_1$, and checks that the series solves $\mathfrak D_{\mathbf p}\varrho=0$ to the order kept.

(5) \emph{Large strengths.} The singular points of $\mathfrak D_{\mathbf p}$ in $t$ are $0,\frac1{25},\frac19,1$, so the series converges for $|s|<\frac1{25}$, on $0<t\le f^\star$, and by Proposition~\ref{prop:sl-slopes}(ii) at $j=0$
\[
P(K)=\int_0^{f^\star}t^K\varrho_f(t)\,dt=g(0)\sum_{r\ge1}\frac{u_r}{u_1}\,(f^\star)^{K+r+1}\,\mathrm B(K+1,r+1),
\]
the terms falling like $\prod_{i\le r}\frac i{K+1+i}$; at $K=150$ the term $r=400$ is $3\cdot10^{-136}$ of the sum. These are the values at $K=150,151,152$.

(6) \emph{The run back.} The law holds at every positive integer (Proposition~\ref{prop:sl-density}), and $c_{5,0}(k)=k^4\neq0$, so $P(k)=-\big(c_{5,1}(k)P(k+1)+c_{5,2}(k)P(k+2)+c_{5,3}(k)P(k+3)\big)/k^4$ from $k=149$ down to $k=1$. In $150$-digit arithmetic, with the renderings to $60$ places,
\begin{align*}
P(1)&=39.9455880119962456989440284984239958918497829247698595723666\ldots,\\
P(2)&=0.294696107418275131390303208655999730074422376946254727936899\ldots,\\
P(3)&=0.00473383394914131550168474663871104845773495867314811066687036\ldots,
\end{align*}
and the reading from $K=250$ agrees to $10^{-133}$.

(7) \emph{Checks.} The low-end relation read on the output (Definition~\ref{def:sl-gaugeread} step (7)) gives $5P(1)-285P(2)+900P(3)=120.000\ldots$ to all the places computed: the output returns $V_1=120$, which the reading did not use. The values agree with the quadrature of Exercise~\ref{ex:cfive-law} to its $22$ places.

So $P_{\binom{[5]}1}(\one;1,1)=39.94558801199624569894402849842\ldots$, read from the law, the native gauge $\one$ and $\det H$ alone, with the germ of the density $\varrho_f$ exact over $\mathbb Q$ up to the factor $15625\sqrt5\,\Gamma(\frac12)^4$.
\end{solution}

\begin{exercise}[The value as a Trailmoment of the slope input of four marks]\label{ex:cfive-slope}
Let $\tilde p_4=-P_4'$ be the slope input of the centre ray of $\binom{[4]}1$ (Definition~\ref{def:ri-inputs}). Show that $P(1)=\Tx{0}{\tilde p_4\,\IZ}{0}=\sum_{x\ge1}\tilde p_4(x)$, find the recurrence of the slope input and its relation at the cut, and write $\tilde p_4(m)$, $m\le5$, in $1$, $v$ and $\tilde p_4(1)$.
\end{exercise}

\begin{solution}
Receipt: \texttt{ex25e\_slope.py}.

\emph{The mark added.} By Proposition~\ref{prop:ri-addmark}(i), merging four marks, $\Gamma(s)^2P(s)=\frac1{2\pi i}\int_{(c)}\Gamma(d)^2\Gamma(s-d)^2P_4(d)\,dd$; at $s=1$ the kernel is $\Gst40{d,\ d,\ 1-d,\ 1-d}{}$, whose right closing reads the derivative of the rest at the positive integers (Lemma~\ref{lem:resonantpair}(i)). Proposition~\ref{prop:ri-addmark}(ii) at $s=1$, with $W_1=1$, gives
\[
P(1)=\Tx{0}{\tilde p_4\,\IZ}{0}=\sum_{x\ge1}\tilde p_4(x):
\]
the centre of five marks at the strength one is the slope input of four marks summed over the positive integers.

\emph{The recurrence.} By Proposition~\ref{prop:ri-rayinput}(ii), $L\tilde p_4=L'P_4$ for the law $L$ of Exercise~\ref{ex:cfour-law}:
\begin{gather*}
\sum_{i=0}^2c_{4,i}(\sigma)\,\tilde p_4(\sigma+i)=\sum_{i=0}^2c_{4,i}'(\sigma)\,P_4(\sigma+i),\\
c_{4,0}'=-3\sigma^2,\qquad c_{4,1}'=6(10\sigma^2+10\sigma+3),\qquad c_{4,2}'=-192(\sigma+1)^2 .
\end{gather*}

\emph{The cut.} Proposition~\ref{prop:ri-boundary} at $x_0=0$, $j=0$, with $c_{4,0}=-\sigma^3$, so $\gamma_3=-1$ and $\gamma_4=0$, and with $\ell_1=0$, every weight being $1$:
\[
\sum_{i=1}^2\big(c_{4,i}(0)\,\tilde p_4(i)-c_{4,i}'(0)\,P_4(i)\big)=\gamma_3\ell_1+\gamma_4V_1=0,\qquad 4\tilde p_4(1)-64\tilde p_4(2)-18P_4(1)+192P_4(2)=0 .
\]
With $P_4(1)=v$ and $P_4(2)=\frac1{16}v-\frac38$ (Exercise~\ref{ex:cfour-law}),
\[
\tilde p_4(2)=\frac1{16}\tilde p_4(1)-\frac3{32}v-\frac98 .
\]

\emph{The rows.} The recurrence, whose coefficient $c_{4,2}(m)=-64(m+1)^3$ does not vanish, gives every slope from the first two:
\begin{align*}
\tilde p_4(3)&=\tfrac7{1024}\tilde p_4(1)-\tfrac{57}{4096}v-\tfrac{279}{2048},\\
\tilde p_4(4)&=\tfrac1{1024}\tilde p_4(1)-\tfrac{73}{32768}v-\tfrac{2819}{147456},\\
\tilde p_4(5)&=\tfrac{679}{4194304}\tilde p_4(1)-\tfrac{13081}{33554432}v-\tfrac{463499}{150994944},
\end{align*}
the coefficient of $\tilde p_4(1)$ being $a_m$ of Exercise~\ref{ex:cfour-law}, since the recurrence has the law as its homogeneous part. So every $\tilde p_4(m)$ lies in $\mathbb Q+\mathbb Q\,v+\mathbb Q\,\tilde p_4(1)$.

\emph{The value.} The slopes have the root rate $\frac1{16}$ along the positive integers (Proposition~\ref{prop:ri-rayinput}(iii)). As renderings, $\tilde p_4(1)=-P_4'(1)=39.3815821231221656249398093451\ldots$ by quadrature of $-\int_0^\infty\log y\,\mathsf g(y)^4dy+2(\log\Gamma)'(1)\,v$, then $\tilde p_4(2)=0.547499039997\ldots$, $\tilde p_4(3)=0.0158846667692\ldots$, $\tilde p_4(4)=0.00059555377139\ldots$, $\tilde p_4(5)=0.0000253984887\ldots$ from the rows; the sum over $m\le70$ is $39.94558801199624569894402849842\ldots$, the value of Exercise~\ref{ex:cfive-gauge} to all $36$ places computed, and the steps $m\ge2$ add $0.564005888874080074\ldots$ to the first slope. The receipt also checks Proposition~\ref{prop:ri-addmark}(iii) at $n=4$, $5\,\Tx{0}{P_4\IZ}{0}-5\,\Tx{0}{(x-1)\,\tilde p_4\IZ}{0}=\Tx{0}{\tilde p_4\IZ}{0}$.

So $P_{\binom{[5]}1}(\one;1,1)=\Tx{0}{\tilde p_4\,\IZ}{0}$, a Trailmoment at the index zero of the slope input of four marks, whose density lies in $\mathbb Q+\mathbb Q\,\Tx3\IZ0+\mathbb Q\,\tilde p_4(1)$ at every point.
\end{solution}

\section{\texorpdfstring{$W_3$ at $(\one;0,2)$}{W3 at (1;0,2)}}\label{sec:exsb-wheel}

The wheel $W_3$ (Proposition~\ref{prop:wheelholonomy}) has the hub, the rim vertices $1,2,3$, the spokes $s_1,s_2,s_3$ and the rims $r_1,r_2,r_3$, $r_i$ joining $i$ to $i+1$; its bases are the $16$ spanning trees, the triples of marks other than the four triangles $\{r_1,r_2,r_3\}$, $\{s_1,s_2,r_1\}$, $\{s_2,s_3,r_2\}$, $\{s_3,s_1,r_3\}$, its rank is $3$, its cobases are the complements of the bases and every weight is $1$. The ray of this section has the origin $0$ and the direction $(\one;0,2)$: masses $1$, exponent $0$ for $K$ and $2$ for $\Psi$, on the balance wall, $m_{\cE}=6=0\cdot3+2\cdot3$. Its stride is $f=\alpha^{\one}/\Psi^2$, only $\Psi$ being active; its walls are $\kappa_{1,S}=|S|-2\big(|S|-r(S)\big)=2r(S)-|S|$, with the values $1,2,3$ on the nonempty proper $S$; and its volume of the direction is $V_1=\sum_\pi\prod_{l=2}^6\kappa_{1,U_l}^{-1}=84$. At $\one$ every mark lies in $8$ of the $16$ bases, so every dual occupation is $\frac12$, the defects $1-2\cdot\frac12$ vanish, and $\alpha^\star=\one$ with $f^\star=\Psi(\one)^{-2}=\frac1{256}$. The values are $P(k)=\langle1\rangle_k=P_{W_3}(k\one;0,2k)$. Since $K$ has the exponent $0$, the constant terms of the powers of $f^{-1}=\Psi^2/\alpha^{\one}$ are written explicitly: $\CT(f^{-j})=[\alpha^{j\one}]\,\Psi^{2j}$.

\begin{exercise}[The box, the law and the spectrum]\label{ex:wthree-box}
\begin{enumerate}[label=(\roman*)]
\item Read the box $\mathfrak B_2$ with one unit of $\Psi$ (Definition~\ref{def:raybox}) modulo the automorphisms of $W_3$, find its free orbits, and show that it is accepted at $\mathfrak s=1$ with the spanning set $b_1=1$, $b_2=b=\frac14\sum_T\alpha^T/\alpha^{\cE\setminus T}$, the sum over the four triangles $T$; give its matrix $\mathbf B(k)$.
\item Derive the law of $[1]$ from the rows, its symbol, and show that $b_1,b_2$ are masters.
\item Show that $\mathbf B$ is not regular at $k=\infty$, and repair it by a change of masters.
\end{enumerate}
\end{exercise}

\begin{solution}
Receipt: \texttt{ex26a\_box.py} (with \texttt{w3common.py}).

(i) \emph{Symmetries.} The $24$ permutations of the four vertices of the graph induce permutations of the marks that carry the spanning trees onto themselves and keep the masses and the exponents: maps $\phi_\beta$ of the first kind of Proposition~\ref{prop:raysym}(ii), forming a group $G$. The level bounds are $\bar a=0$, $K$ being inactive, and $\bar b=1$; $G$ keeps $\Psi$, so it keeps the box.

\emph{The box.} $\mathfrak B_2(0,1)$ holds the symbols of wall degree at most $2$ at the levels $(0,0)$ and $(0,1)$ (Definition~\ref{def:raybox}(iv)): $1167$ Laurent symbols $\alpha^u$, $u_{\cE}=0$, which are the stems, and $2848$ symbols at the level $(0,1)$, $u_{\cE}=3$. Its relations are the letter identities of $\Psi$, $\alpha^u-\sum_{T}(u+\one_{\cE\setminus T};0,1)$, of the $1167$ Laurent symbols, and the mark relations of the stems at the six marks (Definition~\ref{def:raybox}(iii)), here, with the origin $0$, $m_{1,e}=1$, $s_1=0$ and $t_1=2$,
\[
\mathrm{MR}_e(\alpha^u)=(k+u_e)\,\alpha^u-2k\sum_{T\not\ni e}(u+\one_{\cE\setminus T};0,1),
\]
$7002$ of them. Modulo $G$ (Definition~\ref{def:box-reduction}(i)) there are $66$ Laurent orbits, of wall degrees $0$ (one orbit), $1$ (eight) and $2$ (fifty-seven), and $145$ orbits at the level $(0,1)$: $211$ orbit averages, on which the averaged relations, $389$ distinct ones, have the rank $200$ over $\mathbb Q(k)$.

\emph{The reduction.} The ranking puts the orbits at the level $(0,1)$ highest, then the Laurent orbits by decreasing wall degree, with the orbit of $\alpha^{r_1r_2r_3}/\alpha^{s_1s_2s_3}$ last among those of wall degree $1$ (Definition~\ref{def:box-reduction}(ii)). The free orbits are $11$: nine at the level $(0,1)$ and the Laurent orbits of $1$ and of $\alpha^{r_1r_2r_3}/\alpha^{s_1s_2s_3}$. The latter orbit has four members, $\alpha^T/\alpha^{\cE\setminus T}$ for the four triangles $T$, each with its complement the star at the opposite vertex; so the spanning set at $\mathfrak s=1$ is $b_1=1$, $b_2=b$ (Definition~\ref{def:box-reduction}(iii)). The targets $f^{-1}b_1=\Psi^2/\alpha^{\one}$ and $f^{-1}b_2$ involve $9$ and $22$ orbits.

\emph{Acceptance.} The reductions of the two targets were computed by elimination at integer strengths and reconstructed over $\mathbb Q(k)$ (Proposition~\ref{prop:rd-certain}(i),(ii)); with the common denominator $Q_c=(k-1)^2(k-2)^5(k^2-k+1)$ every $Q_c\mathbf B_{il}$ is a polynomial of degree at most $10$:
\begin{align*}
\mathbf B_{11}&=-\frac{2(2k-1)(k^2-4k+2)(7k^2-8k+3)}{(k-1)^2(k^2-k+1)},\qquad \mathbf B_{12}=\frac{14k^2(2k-1)}{k^2-k+1},\\
\mathbf B_{21}&=-\frac{2k(k-3)(k-1)(k+1)(2k-1)(7k^2-14k+8)}{(k-2)^4(k^2-k+1)},\\
\mathbf B_{22}&=\frac{2k^2(k-1)(2k-1)(k^2-2)(7k^2-20k+15)}{(k-2)^5(k^2-k+1)}.
\end{align*}
The proof that these are the reductions over $\mathbb Q(k)$: for each target $x_i$ the row $x_i-\sum_l\mathbf B_{il}b_l$, times $Q_c$, stacked under the relation matrix $A(k)$, whose entries are affine in $k$ (Proposition~\ref{prop:rd-lattice}(iii)), has the rank $200$ at $230$ integer strengths; its minors of order $201$ are polynomials of degree at most $219$, so they vanish identically, the row lies in $\mathcal R^G_{\mathfrak B}$ over $\mathbb Q(k)$, and $\sum_l\mathbf B_{il}b_l$, a combination of free orbits, is the reduction (Definition~\ref{def:box-reduction}(ii)). So the box is accepted at $\mathfrak s=1$ with the matrix $\mathbf B$, the classes of $b_1,b_2$ span a submodule $\mathfrak N\ni[1]$, there is a certificate supported on the stems, and $\langle b_i\rangle_{k-1}=\sum_l\mathbf B_{il}(k)\langle b_l\rangle_k$ at every strength (Theorem~\ref{thm:raybox}(i)--(iii)).

(ii) \emph{The law.} With $w_0=(1,0)$ and $w_i(k)=w_{i-1}(k-1)\mathbf B(k)$, the first dependence is $q_0w_0+q_1w_1+q_2w_2=0$ with
\begin{gather*}
q_0=16(k-1)(2k-3)(2k-1)(4k-5)(4k-3),\qquad q_2=(k-2)^5,\\
q_1=-2(2k-3)(65k^4-390k^3+885k^2-900k+346),
\end{gather*}
and $p_j(k)=q_{2-j}(k+2)$ (Proposition~\ref{prop:rays-lawrows}(ii), Theorem~\ref{thm:raybox}(iv)) gives
\begin{multline*}
k^5P(k)-2(2k+1)(65k^4+130k^3+105k^2+40k+6)\,P(k+1)\\
+16(k+1)(2k+1)(2k+3)(4k+3)(4k+5)\,P(k+2)=0
\end{multline*}
at every strength (Corollary~\ref{cor:rays-lowend}(i)). Its symbol is $1-260\lambda+1024\lambda^2=(1-4\lambda)(1-256\lambda)$, and at both roots the exponent $-\chi_1(\lambda)/(\lambda\chi'(\lambda))$, with the second symbol $\chi_1(\lambda)=10\lambda(512\lambda-65)$, is $-\frac52=-\frac{n-1}2$ (Proposition~\ref{prop:sl-exponents}). The constant terms $\CT(f^{-j})=1,12,756,78960,10451700,1587790512,263964176784$ ($j\le6$) satisfy $\sum_ip_i(-j)\CT(f^{i-j})=0$ (Corollary~\ref{cor:rays-lowend}(i)), a check.

\emph{Masters.} Two native gauges: $\one$, with $\Psi=16$, the native value $\frac1{256}$ and the determinant $\frac3{256}$ of the Hessian of $-\log f$ on the chart; and $\alpha_c=(1,\omega,\omega^2,\omega^2,1,\omega)$ on $(s_1,s_2,s_3,r_1,r_2,r_3)$, $\omega=e^{2\pi i/3}$, at which every triangle and every star has the product $1$, $\Psi(\alpha_c)=-2$, every dual occupation is $\frac12$ (so $c^{(1)}=0$), the native value is $\prod_e\alpha_{c,e}/\Psi^2=\frac14$, and the Hessian has the determinant $-\frac32$. Both are nondegenerate with distinct native values, so by Corollary~\ref{cor:symbol} at $g=1$ every recurrence of $\langle1\rangle$ that follows from the relations has order at least $2$: the cyclic submodule of $[1]$, inside $\mathfrak N$ of dimension at most $2$, has dimension $2$. So $b_1,b_2$ are masters, $\mathbf B=\mathbf B_{\mathbf b}$, and the relation above is the law of $[1]$.

(iii) \emph{The spectrum.} $b(\alpha_c)=1=b(\one)$, so the vectors $\beta_c=(1,1)$ of the two native gauges are proportional with distinct native values, and $\mathbf B$ is not regular at $k=\infty$ (Proposition~\ref{prop:rays-spectrum}(ii)); indeed $\mathbf B(k)/k\to\begin{bmatrix}-28&28\\-28&28\end{bmatrix}$. The change $T(k)=\begin{bmatrix}1&0\\-k&k\end{bmatrix}$, to the elements $[1]$ and $k([b]-[1])$, gives $\tilde{\mathbf B}(k)=T(k-1)\mathbf B(k)T(k)^{-1}$ with
\[
\tilde{\mathbf B}(\infty)=\begin{bmatrix}88&28\\504&172\end{bmatrix},\qquad \det\big(\lambda-\tilde{\mathbf B}(\infty)\big)=(\lambda-256)(\lambda-4),
\]
the eigenvectors $(\frac16,1)$ and $(-\frac13,1)$: the reciprocals of the native values $\frac1{256}$ and $\frac14$ (Proposition~\ref{prop:rays-spectrum}(iii)). As $T(k-1)T(k)^{-1}=\operatorname{diag}\big(1,\frac{k-1}k\big)\to I$, the eigenvalues of $\mathbf B(k)$ tend to $256$ and $4$: $\operatorname{tr}\mathbf B(k)\to260$ and
\[
\det\mathbf B(k)=\frac{16k^2(2k-1)^2(4k-5)(4k-3)(k^2-3k+3)}{(k-2)^5(k-1)(k^2-k+1)}\to1024 .
\]
So the box gives the masters $1,b$, the matrix $\mathbf B_{\mathbf b}$ above and the law of $[1]$ of symbol $(1-4\lambda)(1-256\lambda)$; $\mathbf B_{\mathbf b}$ grows like $k$, and the masters $1$ and $k(b-1)$ repair it.
\end{solution}

\begin{exercise}[The low-end relation, the brackets and the value]\label{ex:wthree-bracket}
From the law of Exercise~\ref{ex:wthree-box}: write the low-end relation; show $7<P(1)<\frac{64}7$; give the brackets of the two conditions $P(K)>0$ and $P(K)<f^\star P(K-1)$ at every $K$, their widths and their rate; and give $P(1)$.
\end{exercise}

\begin{solution}
Receipt: \texttt{ex26b\_bracket.py}.

\emph{The low-end relation.} By Corollary~\ref{cor:rays-lowend}(ii), with $[k^5]p_0=1$, $\mathsf L_0=V_1=84$ (Definition~\ref{def:rays-lowend}, Proposition~\ref{prop:rays-lowpoles}(ii), which allows the exponent $0$ of $K$), $p_1(0)=-12$ and $p_2(0)=720$:
\[
84-12P(1)+720P(2)=0,\qquad P(2)=\frac{P(1)-7}{60},
\]
and every value lies in $\mathbb Q+\mathbb Q\,P(1)$, $p_2$ having no zero at $k\ge0$.

\emph{The first bracket.} Proposition~\ref{prop:rd-onemaster}(ii), with $a=1$: $-aV_1/p_1(0)=7$ and $-aV_1/\big(p_1(0)+f^\star p_2(0)\big)=-84/(-12+\frac{720}{256})=-84/(-\frac{147}{16})=\frac{64}7$, both positive, so $7<P(1)<\frac{64}7$.

\emph{The brackets.} Let $c$ and $r$ solve the law with $(c(1),c(2))=(1,\frac1{60})$, of low-end coefficient $0$, and $(r(1),r(2))=(0,-\frac7{60})$, of low-end coefficient $84$; then $P(K)=c(K)P(1)+r(K)$ (Proposition~\ref{prop:ra-lowend}(ii)):
\[
\begin{array}{c|cccccc}
K&1&2&3&4&5&6\\\hline
c(K)&1&\frac1{60}&\frac1{900}&\frac{47}{400400}&\frac{79}{5105100}&\frac{2273}{977728752}\\[3pt]
r(K)&0&-\frac7{60}&-\frac{173}{21600}&-\frac{73219}{86486400}&-\frac{140653}{1260230400}&-\frac{3541028833}{211189410432000}
\end{array}
\]
Positivity at $K$ (Proposition~\ref{prop:rd-moment}(i) with $p=1$) and at $K-1$ with $p=f^\star-t$, where $c(K)>0$ and $256c(K)-c(K-1)>0$ (true for $2\le K\le90$), give
\[
a_K=-\frac{r(K)}{c(K)}<P(1)<b_K=\frac{r(K-1)-256\,r(K)}{256\,c(K)-c(K-1)} ;
\]
$K=2$ is the first bracket. The widths are $2.143$, $1.297\cdot10^{-2}$, $1.513\cdot10^{-4}$, $2.042\cdot10^{-6}$ at $K=2,\dots,5$, $3.50\cdot10^{-19}$ at $K=12$ and $1.161\cdot10^{-33}$ at $K=20$; the brackets are nested, and the ratio of consecutive widths, $0.01524$ at $K=10$ and $0.01562$ at $K=80$, tends to $\frac1{64}=f^\star/\lambda'$, $\lambda'=\frac14$ (Theorem~\ref{thm:positiveselection}(iv)). At $K=12$
\[
\frac{125311627268888255929639}{17374611099178960704000}<P(1)<\frac{290907750860286043092257947}{40334717113590600816192000}.
\]

\emph{The value.} The roots $\frac1{256}$ and $\frac14$ are distinct, the second outside the disc $|\lambda|\le f^\star$, and $p_0,p_2$ have the largest degree $5$ and no zero at $k\ge1$: by Theorem~\ref{thm:positiveselection}(iii) the values are the minimal solution, and by (iv), with $x^{(0)}=r$ and $x^{(1)}=c$, the approximants are $M^{(K)}=a_K$. At $K=90$, of width $4\cdot10^{-160}$, as a rendering,
\[
P(1)=7.21234141895756571239842896906869994458991775404299329075362\ldots,
\]
and $P(2)=(P(1)-7)/60$. The continued fraction of Proposition~\ref{prop:rd-onemaster}(iii), $\rho_1=\mathsf M_1\circ\cdots\circ\mathsf M_K(0)$ and $P(1)=-aV_1/(p_1(0)+p_2(0)\rho_1)$, gives the same number, to $20$ places at depth $10$ and to $39$ at depth $40$.
\end{solution}

\section{Positivity against the gauge end}\label{sec:exsb-positivity}

The centre ray of $\binom{[3]}1$ at the strength one has the stride $\big(K(\alpha)K(\alpha^{-1})\big)^{-1}$, $K=\alpha_1+\alpha_2+\alpha_3$, $V_1=3!=6$, and the values $P(k)=P_{\binom{[3]}1}(k\one;k,k)$; $\ell=P(1)$. By Theorem~\ref{thm:ri-law} at $n=3$ its law is $s^2P(s)-(10s^2+10s+3)P(s+1)+9(s+1)^2P(s+2)=0$, of symbol $(1-\lambda)(1-9\lambda)$ with the roots $f^\star=\frac19$ and $1$ (Proposition~\ref{prop:ri-signs}), and the low-end relation (Corollary~\ref{cor:rays-lowend}(ii)) is $6-3\ell+9P(2)=0$. So $P(K)=c(K)\,\ell+r(K)$ with $(c(1),c(2))=(1,\frac13)$ and $(r(1),r(2))=(0,-\frac23)$; the law is reflexive with the ratio $\frac19$, and $c(K)=9^{1-K}\mathsf C_{K-1}(1)$ (Theorem~\ref{thm:ra-closingcoordinate}(ii)), $\mathsf C_j(1)=1,3,15,93,639,4653,\dots$.

\begin{exercise}[All positivity conditions on the centre of $\binom{[3]}1$]\label{ex:cthree-positivity}
\begin{enumerate}[label=(\roman*)]
\item Read $\ell$ from the gauge end.
\item For $N\le22$ compare the bracket of the two conditions $P(K)>0$, $P(K)<\frac19P(K-1)$ at $K=N+1$ with the bracket of all positivity conditions of degree at most $N$ at the strength $1$ (Proposition~\ref{prop:rd-moment}), both read on $P(1),\dots,P(N+1)$.
\item Find the rate of the second (Theorem~\ref{thm:rd-rate}).
\end{enumerate}
\end{exercise}

\begin{solution}
Receipt: \texttt{ex27\_positivity.py}.

(i) \emph{The gauge end.} $n=3$, $f^\star=\frac19$, $\det H=(2\sigma)^2/K(\one)^3=\frac4{27}$ (Proposition~\ref{prop:ri-signs}(ii)), so $g(0)=\frac{2\pi}{f^\star\sqrt{\det H}}=27\sqrt3\,\Gamma(\frac12)^2$ (Proposition~\ref{prop:sl-ends}(ii)). The density operator (Definition~\ref{def:sl-densityop}) is $a_0+a_1\partial_t+a_2\partial_t^2$ with $a_0=1-23t+36t^2$, $a_1=t(3-40t+45t^2)$, $a_2=t^2(t-1)(9t-1)$. At $\frac19$, $b_{2,1}=\frac8{81}$, $\gamma=0=\frac{n-3}2$, $R_0=\frac8{81}x^2$, $R_1=-\frac19(15x^2+23x+10)$, $R_2=(2x+3)(3x+5)$, $R_3=9(x+2)^2$: the exponents are $0$ and $\gamma=0$, and the line is the holomorphic germ, $u_0=g(0)$ and $R_0(N)u_N=-\sum_{m=1}^{3}R_m(N-m)u_{N-m}$ (Theorem~\ref{thm:gaugeend}(iv)), so $u_r/u_0=1,\frac{45}4,\frac{3645}{32},\frac{70713}{64},\frac{21500397}{2048}$ for $r\le4$. The Beta sums at $K=200,201$ with $500$ terms and the law run back give, as a rendering,
\[
\ell=2.34390723868945889060156228887227706909540300963585626066630\ldots;
\]
the reading from $K=300$ agrees to $10^{-129}$, and $3P(1)-9P(2)=6.000\ldots$ returns $V_1$.

(ii) \emph{The two conditions.} At $K=N+1$, $a_K=-r(K)/c(K)<\ell<b_K=\big(r(K-1)-9r(K)\big)/\big(9c(K)-c(K-1)\big)$: at $K=2$, $2<\ell<3$; at $K=3$, with $c(3)=\frac5{27}$ and $r(3)=-\frac{23}{54}$, $\frac{23}{10}<\ell<\frac{19}8$.

\emph{All conditions.} By Proposition~\ref{prop:rd-moment}(ii) the conditions of degree at most $N$ at the strength $1$ are the positive semidefiniteness of the Hankel matrices $[P(1+i+i')]$, $[P(2+i+i')]$, $[\frac19P(1+i+i')-P(2+i+i')]$ and $[\frac19P(2+i+i')-P(3+i+i')]$ of the sizes allowed by $N$. Each is affine in $\ell$; the set where one is positive semidefinite is an interval, its least eigenvalue being concave in $\ell$, whose ends are the roots of its determinant nearest to a point where it is positive definite, such as the gauge-end $\ell$. The determinants are computed exactly as polynomials in $\ell$ and their real roots isolated in ball arithmetic. At $N=1$ the conditions are $\frac13\ell-\frac23\ge0$ and $\frac23-\frac29\ell\ge0$: $2\le\ell\le3$. At $N=2$ the binding ones are
\[
\det\begin{bmatrix}P(1)&P(2)\\P(2)&P(3)\end{bmatrix}=\frac{4\ell^2+\ell-24}{54}\ge0,\qquad \frac19P(2)-P(3)=\frac{19}{54}-\frac4{27}\ell\ge0,
\]
so $\frac{\sqrt{385}-1}8\le\ell\le\frac{19}8$, of width $\frac{20-\sqrt{385}}8=0.0473\ldots$. The widths:
\[
\begin{array}{c|cc}
N&\text{two conditions}&\text{all conditions}\\\hline
1&1&1\\
2&7.50\cdot10^{-2}&4.73\cdot10^{-2}\\
4&7.84\cdot10^{-4}&6.90\cdot10^{-5}\\
8&1.10\cdot10^{-7}&9.52\cdot10^{-11}\\
12&1.64\cdot10^{-11}&1.04\cdot10^{-16}\\
16&2.46\cdot10^{-15}&1.03\cdot10^{-22}\\
22&4.58\cdot10^{-21}&9.10\cdot10^{-32}
\end{array}
\]
and the ratio of the widths of all conditions at $N$ and $N-1$ is $0.0382$, $0.0335$, $0.0320$, $0.0314$, $0.0308$ at $N=4,8,12,16,22$. Every bracket contains the gauge-end $\ell$.

(iii) \emph{The rate.} The free direction is $u=c$. By Proposition~\ref{prop:sl-torus}(i), $\mathsf C_j(1)=\int t^{-j}\nu^{\mathbb T}_f(dt)$, the torus law (Definition~\ref{def:torusvalue}) being carried by $[\frac19,\infty)$; with $\tau=1/(9t)$, $c(K)=\int_0^1\tau^K\nu(\tau)\,d\tau$, $\nu(\tau)$ being $\tau^{-1}$ times the density of the image of $\nu^{\mathbb T}_f$, and by Proposition~\ref{prop:sl-torus}(iii) $\nu(\tau)\to\frac19\cdot\frac{g(0)}{(2\pi)^2}=\frac{3\sqrt3}{4\pi}\neq0$ as $\tau\uparrow1$: the exponent $a=0$ of Theorem~\ref{thm:rd-rate} at $\lambda=1$. With $w=2\lambda/f^\star-1=17$ the bracket of degree $N$ has length at most $C\,N\,e^{-Ng(1)}$,
\[
e^{-g(1)}=\frac1{17+\sqrt{288}}=17-12\sqrt2=0.02943\ldots,
\]
against the factor $\frac19$ per step of the two conditions (Theorem~\ref{thm:positiveselection}(iv)): $-\log_{10}(17-12\sqrt2)=1.531$ digits of $\ell$ per value against $\log_{10}9=0.954$. The ratios of the table approach $17-12\sqrt2$ from above, as the factor $N$ allows; with the nine values $P(1),\dots,P(9)$ all conditions give $\ell$ to $10^{-10}$, the two conditions to $10^{-7}$.
\end{solution}

\section{What the closings see}\label{sec:exrays-closings}

% Exercise for the exercises chapter on rays (Part XI): what the closings see at small centres.
% Source: items (i) to (iii) of rem:blindcentres of the v1 chapter ch:twoends; item (iv) is
% Proposition prop:onemarkhidden.  Lists named here: H_4, W_3, N_6 and the uniform lists binom([n],r).
% Receipts: /home/claude/v2/new/twoends_ed/scripts/ (boxexact.py, boxes_all.py, three_exact.py, run_blind2.sh with
% blind2.py of receipts_src/twoends/closings_see_scripts, signs_span.py, n_exchange.py, u25_core.py, u36_lift.py,
% binf_exact.py, oddcentre_hessian.py, ell_check.py) and the Singular files u25.sing, u36.sing.  No new macros.

\begin{exercise}[What the closings see at small centres]\label{ex:closingssee}
Take the following lists, with unit weights, on the marks $[n]=\{1,\dots,n\}$: the lists $\binom{[n]}1$, $n=3,4,5$, whose bases are the single marks; the list $\mathrm H_4$ on $[4]$, whose bases are the pairs other than $\{1,2\}$; the wheel $W_3$ of Proposition~\ref{prop:wheelholonomy}, whose marks are the six edges $e_1=01$, $e_2=02$, $e_3=03$, $e_4=12$, $e_5=13$, $e_6=23$ of the complete graph on the vertices $0,1,2,3$ and whose bases are its sixteen spanning trees; the list $\mathrm N_6$ on $[6]$, whose bases are the nineteen triples other than $\{1,2,3\}$; and the uniform lists $\binom{[5]}2$ and $\binom{[6]}3$. Read each on its centre ray at the strength $1$ (Definition~\ref{def:centreray}), with origin $0$, direction $(\one;1,1)$ and stride $f=\alpha^{\one}(K\Psi)^{-1}$, and read $\mathrm H_4$ also on the ray with origin $0$ and direction $(1,1,3,3;3,1)$. Let $G$ be generated by the automorphisms that keep the masses and, at a centre, the Cremona, with the exchanges where they are named. For a box accepted at $\mathfrak s$ (Definition~\ref{def:box-reduction}), with the spanning set $\mathbf b=(b_1,\dots,b_{\mathfrak m})$, $b_1=1$, and the matrix $\mathbf B$, write $\mathbf B=P_0/Q_0$ with $Q_0$ the monic least common denominator of its entries, $D_0=\max(\deg P_0,\deg Q_0)$, and let $\mathcal S_D$ be the space of the pairs $(P,Q)$ of a matrix and a scalar of polynomials of degree at most $D$ with $Q(-j)\,\mathsf C_{j+1}(\mathbf b)=P(-j)\,\mathsf C_j(\mathbf b)$ for all large $j$, the closing equations of Proposition~\ref{prop:blinddim}.
\begin{enumerate}[label=(\roman*)]
\item At the centres of $\binom{[3]}1$, $\binom{[4]}1$, $\binom{[5]}1$, $\mathrm H_4$, $W_3$, with and without its exchanges, and $\mathrm N_6$ with its exchanges, and on the ray $(1,1,3,3;3,1)$ of $\mathrm H_4$, show that $\dim\mathcal S_D=0,1,2,3$ at $D=D_0-1,\dots,D_0+2$, the line at $D_0$ being spanned by $(P_0,Q_0)$, and find a degree up to which no element of the span of $\mathbf b$ with polynomial coordinates is annihilated by the closings. Work the centre of $\binom{[3]}1$ in full.
\item At the centre of $\mathrm N_6$, with $G$ generated by the automorphisms and the Cremona but not the exchanges, show that the closings miss a nonzero submodule of the span of the spanning set.
\item At the centre of $\binom{[5]}2$, with every symmetry of the four kinds in $G$: find the native gauges; show that the spanning set of the box $\mathfrak B_3(1,1)$ at $\mathfrak s=2$ is a set of masters of its span $\mathfrak N$; show that $\mathfrak N$ is the direct sum of the cyclic submodule of $[1]$ and a line $L$ on which the closings, the values and $A_\star$ vanish, reading $\binom{[5]}2$ as the contraction of $\binom{[6]}3$ (Proposition~\ref{prop:onemark}); find which solutions at the native gauges see $L$; and read the centre of $\binom{[6]}3$ by its box.
\end{enumerate}
\end{exercise}

\begin{solution}
Every ray here has origin $0$, positive letter exponents and positive walls, so Definition~\ref{def:closings}, Theorem~\ref{thm:closings}, Proposition~\ref{prop:principal} and Theorem~\ref{thm:onekernel} apply. Three facts are used throughout.

(a) If an element $m=\sum_lp_l(k)[b_l]$ with polynomial coefficients has values that vanish identically, then $\sum_lp_l(-j)\,\mathsf C_j(b_l)=0$ for every $j\ge0$, not only for large $j$: the $b_l$ are Laurent insertions, so their closings are defined from $j=0$ on, and Proposition~\ref{prop:principal} gives $V_1\sum_lp_l(-j)\mathsf C_j(b_l)=\lim_{k\to-j}(k+j)^{n-1}\langle m\rangle_k=0$. This holds for the elements of $\mathcal L$ (Proposition~\ref{prop:rays-strength}(ii)) and for the elements on which the closings vanish (Theorem~\ref{thm:onekernel}). The box gives $Q_0(k)\big(f^{-1}b_i-\sum_l\mathbf B_{il}(k)\,b_l\big)\in\mathcal L$ (Theorem~\ref{thm:raybox}(i)), and $\mathsf C_j(f^{-1}b_i)=\mathsf C_{j+1}(b_i)$, so
\[
\mathsf C_{j+1}(\mathbf b)=\mathbf B(-j)\,\mathsf C_j(\mathbf b)\qquad\text{for every }j\ge0\text{ with }Q_0(-j)\neq0 .
\]

(b) By (a), $(P,Q)\in\mathcal S_D$ exactly when every row of $Q\mathbf B-P$, read as an element of the span of $\mathbf b$, is annihilated by the closings; $Q_0$ times such a row is a polynomial row $p$ with $p(-j)\cdot\mathsf C_j(\mathbf b)=0$ for every $j\ge0$, so every pair of $\mathcal S_D$ satisfies its equation at every $j\ge0$ with $Q_0(-j)\neq0$. The equations at finitely many such $j$ cut out a space that contains $\mathcal S_D$, and so does that system reduced modulo a prime, whose rank is at most its rank over $\mathbb Q$. The pairs $(cP_0,cQ_0)$, $\deg c\le D-D_0$, lie in $\mathcal S_D$. So when a computed dimension is $\max(0,D-D_0+1)$, it is $\dim\mathcal S_D$.

(c) Likewise the polynomial rows of degree at most $e$ that the closings annihilate at finitely many $j$ contain the coordinate rows, of degree at most $e$, of the elements of the span of $\mathbf b$ on which the closings vanish.

\emph{(i) The centre of $\binom{[3]}1$.} The letters are $K=\alpha_1+\alpha_2+\alpha_3$ and $\Psi=\alpha_1\alpha_2+\alpha_1\alpha_3+\alpha_2\alpha_3$, and
\[
f^{-1}=\frac{K\Psi}{\alpha_1\alpha_2\alpha_3}=3+\sum_{i\neq j}\frac{\alpha_i}{\alpha_j}.
\]
For a nonempty proper $U$, $d(U)=0$ and $d^\ast(U)=|U|-1$, so every wall of the direction is $\kappa_{1,U}=1$ and the wall degree of $\alpha^u$ is $\max_U(-u_U)$ (Definition~\ref{def:raybox}(ii)). $G$ is generated by the six bijections of $[3]$ and the Cremona; it holds every symmetry of the four kinds, the list having no exchange (its bases have one mark, its cobases two). The box $\mathfrak B_2(1,1)$ has $110$ symbols; modulo $G$ its relations, taken at one symbol of each orbit, are $26$ combinations of $16$ orbit averages, of rank $13$ over $\mathbb Q(k)$ (receipts \texttt{boxexact.py}, \texttt{boxes\_all.py}, exact elimination over $\mathbb Z[k]$). The free orbits are those of $1$, of $\alpha_3/\alpha_1$, of wall degree $1$, and of $(\one;1,1)$, at a positive level; so at $\mathfrak s=1$ the spanning set is $b_1=1$, $b_2=e_G(\alpha_3/\alpha_1)=\frac16\sum_{i\neq j}\alpha_i/\alpha_j$, both targets reduce to it, and
\[
\mathbf B(k)=\begin{bmatrix}3&6\\[2pt]\dfrac{2(k-1)}{k-2}&\dfrac{(k-1)(7k-11)}{(k-2)^2}\end{bmatrix},
\]
with $Q_0=(k-2)^2$ and
\[
P_0=\begin{bmatrix}3(k-2)^2&6(k-2)^2\\2(k-1)(k-2)&(k-1)(7k-11)\end{bmatrix},
\]
so $D_0=2$; the first row is the Laurent identity $f^{-1}b_1=3b_1+6b_2$.

The native gauge of the direction, $\one$, has the native value $\frac19$, and the three signs with one scale $-1$ have $K(\eta)=1$, so they are native gauges of the native value $1$ (Theorem~\ref{thm:signgauges}(ii)); on a chart $\det(-H)=\frac4{27}$ at $\one$ and $-4$ at the signs (receipt \texttt{signs\_span.py}), so all four are nondegenerate. Their solutions, summed over the two orbits, have distinct types and do not vanish at $[1]$, so the span $\mathfrak N$ of $b_1,b_2$ has dimension $2$ (Lemma~\ref{lem:types}): $\mathbf b$ is a set of masters and $\mathbf B=\mathbf B_{\mathbf b}$. The rows $w_0=(1,0)$ and $w_1=(3,6)$ are independent, so $\mathfrak N$ is the cyclic submodule of $[1]$ (Proposition~\ref{prop:rays-lawrows}(ii)).

The closings are the constant terms $\mathsf C_j(b)=\CT(b\,f^{-j})$; $\mathsf C_j(1)$ is the coefficient of $\alpha_1^j\alpha_2^j\alpha_3^j$ in $(K\Psi)^j$. By direct expansion,
\[
\begin{array}{c|ccccccccc}
j&0&1&2&3&4&5&6&7&8\\ \hline
\mathsf C_j(b_1)&1&3&15&93&639&4653&35169&272835&2157759\\
\mathsf C_j(b_2)&0&1&8&60&456&3535&27888&223209&1807760
\end{array}
\]
for instance $\mathsf C_1(b_2)=\CT\big(b_2(3+6b_2)\big)=6\cdot\frac6{36}=1$, the six products $\frac{\alpha_i}{\alpha_j}\frac{\alpha_j}{\alpha_i}$ in $b_2^2$ being the only constant ones. $Q_0$ has no zero at a nonpositive integer, so (a) gives $\mathsf C_{j+1}=\mathbf B(-j)\,\mathsf C_j$ for every $j\ge0$: at $j=1$ and $j=2$,
\[
\mathbf B(-1)=\begin{bmatrix}3&6\\ \frac43&4\end{bmatrix},\quad \mathbf B(-1)\begin{bmatrix}3\\1\end{bmatrix}=\begin{bmatrix}15\\8\end{bmatrix};\qquad \mathbf B(-2)=\begin{bmatrix}3&6\\ \frac32&\frac{75}{16}\end{bmatrix},\quad \mathbf B(-2)\begin{bmatrix}15\\8\end{bmatrix}=\begin{bmatrix}93\\60\end{bmatrix},
\]
and the direct expansion agrees for every $j\le79$ (receipt \texttt{three\_exact.py}).

The closing equations at the degree $D$ have $5(D+1)$ unknowns, the coefficients of the four entries of $P$ and of $Q$, and two equations for each $j$. With $1\le j\le79$, $158$ equations, exact elimination over $\mathbb Q$ gives the solution dimensions $0,1,2,3$ at $D=1,2,3,4$; by (b) these are $\dim\mathcal S_D$, and the line at $D=2$ is spanned by $(P_0,Q_0)$. So the closings alone single out $\mathbf B_{\mathbf b}$ among the matrices of degree at most $2$, with no source supplied. The rows $(p_1,p_2)$ of degree at most $24$ with $p_1(-j)\mathsf C_j(b_1)+p_2(-j)\mathsf C_j(b_2)=0$ for $1\le j\le80$, $50$ unknowns and $80$ equations, are only $0$; by (c), no nonzero element of $\mathfrak N$ with polynomial coordinates of degree at most $24$ is annihilated by the closings.

\emph{(i) The other rays.} The same computation runs on the other boxes (receipt \texttt{blind2.py}, driven by \texttt{run\_blind2.sh}): the closings at $j\le7$ are read on the torus as constant terms and continued by (a), and the equations are reduced modulo the prime $2^{61}-1$; by (b) and (c) the dimensions found are upper bounds, attained here. The matrices $\mathbf B$ are those of the boxes built from Definition~\ref{def:raybox} and eliminated exactly over $\mathbb Q(k)$ (receipt \texttt{boxes\_all.py}); at $\mathrm N_6$ they are the rayreader's, checked against an elimination of the whole box at a random prime and strength. The spanning sets, with exponents in the order of the marks, are: $1,e_G(\alpha_4/\alpha_1)$ for $\binom{[4]}1$; $1,e_G(\alpha_5/\alpha_1),e_G(\alpha_5^2/\alpha_1^2)$ for $\binom{[5]}1$; $1,e_G\alpha^u$ with $u=(-1,0,0,1),(-1,1,-1,1),(-1,1,0,0)$ for $\mathrm H_4$ at its centre, and $u=(0,1,-1,0),(0,2,-2,0),(0,2,-1,-1)$ on its ray $(1,1,3,3;3,1)$; $1,e_G\alpha^u$ with $u=(-1,-1,0,0,1,1),(-1,-1,0,1,0,1),(-1,-1,1,1,-1,1),(-1,-1,1,1,0,0)$ for $W_3$, with or without the exchanges; and for $\mathrm N_6$ with its exchanges $1,e_G\alpha^u$ with the twelve $u=(-1,-1,0,0,1,1)$, $(-1,-1,1,-1,1,1)$, $(-1,-1,1,0,0,1)$, $(-1,0,0,0,0,1)$, $(-1,0,1,-1,0,1)$, $(-1,0,1,0,0,0)$, $(-1,-1,-1,1,1,1)$, $(-2,-2,0,0,2,2)$, $(-2,-2,0,1,1,2)$, $(-2,-2,1,-1,2,2)$, $(-2,-2,1,0,1,2)$, $(-2,-2,1,1,1,1)$. In the table, $\dim\mathcal S_D$ is listed at $D=D_0-1,\dots,D_0+2$, $h$ is the first index at which the rows $w_0,\dots,w_h$ of $[1]$ are dependent (Proposition~\ref{prop:rays-lawrows}(ii)), and the last but one column gives the degree up to which no nonzero row is annihilated, at the $j$ from $10$ to the bound shown.
\begin{center}
\begin{tabular}{lccccccc}
\toprule
ray&$|G|$&box, $\mathfrak s$&$\mathfrak m$&$D_0$&$\dim\mathcal S_D$&rows&$h$\\
\midrule
$\binom{[4]}1$&$48$&$(2;1,1)$, $1$&$2$&$3$&$0,1,2,3$&$24$, $j\le120$&$2$\\
$\binom{[5]}1$&$240$&$(3;1,1)$, $2$&$3$&$6$&$0,1,2,3$&$24$, $j\le120$&$3$\\
$\mathrm H_4$&$8$&$(2;1,1)$, $1$&$4$&$3$&$0,1,2,3$&$24$, $j\le150$&$4$\\
$\mathrm H_4$, $(1,1,3,3;3,1)$&$4$&$(2;1,1)$, $1$&$4$&$10$&$0,1,2,3$&$24$, $j\le150$&$4$\\
$W_3$&$48$, $96$&$(2;1,1)$, $1$&$5$&$9$&$0,1,2,3$&$40$, $j\le260$&$4$\\
$\mathrm N_6$, exchanges&$144$&$(3;1,1)$, $2$&$13$&$16$&$0,1,2,3$&$32$, $j\le480$&$11$\\
\bottomrule
\end{tabular}
\end{center}
The denominators are $Q_0=(k-2)^3$, $(k-3)^4$, $(k-2)^3$, $k(k-1)^3(k-\frac43)^3(k-\frac53)^3$, $(k-2)^5(k-\frac32)(k^3-3k^2+\frac32k+2)$, and for $\mathrm N_6$ a polynomial of degree $16$ with the factors $(k-4)(k-3)^5(k-2)^3(k-\frac52)$ and an irreducible sextic; only that of the ray $(1,1,3,3;3,1)$ vanishes at a nonpositive integer, at $k=0$, below the $j$ used. At $\binom{[4]}1$ the native gauges $\one$ and the four signs with one scale $-1$, of the native values $\frac1{16}$ and $\frac14$, and at $\binom{[5]}1$ the native gauges $\one$ and the signs with one and with two scales $-1$, of the native values $\frac1{25}$, $\frac19$ and $1$, are nondegenerate (receipt \texttt{signs\_span.py}), so by Lemma~\ref{lem:types} their spanning sets are sets of masters. At $\mathrm H_4$, $W_3$ and $\mathrm N_6$ the signs of the largest native value are degenerate native gauges, and the computations are read for the spanning set as it stands; then $h$ bounds the dimension of the cyclic submodule of $[1]$ from above. At each of these rays the closings alone single out the box's matrix among the matrices of degree at most $D_0$; at $W_3$ and at $\mathrm N_6$ the rows of $[1]$ become dependent before the size of the spanning set is reached, and the closings single out the whole matrix on the spanning set, not only the law that the rows of $[1]$ give.

\emph{(ii) $\mathrm N_6$ without the exchanges.} Now $|G|=72$, and the box $\mathfrak B_3(1,1)$ is accepted at $\mathfrak s=2$ with $15$ elements: the thirteen above, here averaged over the smaller group, and $e_G\alpha^u$ for $u=(-1,0,0,-1,1,1)$ and $(-2,-2,2,-2,2,2)$. The bijection $x$ with $x(i)=i+3$ and $x(i+3)=i$, $i\le3$, carries the bases, the triples other than $\{1,2,3\}$, onto the complements of the bases, the triples other than $\{4,5,6\}$: it exchanges the list, and at the centre it is a symmetry of the ray (Proposition~\ref{prop:raysym}(ii)). With $G_x$ generated by $G$ and $x$, $G$ is normal of index $2$ in $G_x$, and for $b=e_G\alpha^{(-1,-1,1,0,0,1)}$, an element of the spanning set, $x^\ast b=e_G\alpha^{(-1,0,0,-1,1,1)}$, another (receipt \texttt{n\_exchange.py}); so $(1-e_{G_x})b=\frac12(b-x^\ast b)$, on which the closings, the values and $A_\star$ vanish (Corollary~\ref{cor:seeall}, applied to $G_x$). It is not $0$: at the sign $\eta$ with the scales $-1$ on $\{4,5\}$, $K(\eta)=-5$, a nondegenerate native gauge of the native value $\frac1{25}$ with $\det(-H)=-\frac{9216}{15625}$ on a chart, $b(\eta)=-\frac13$ and $x^\ast b(\eta)=1$, so the leading coefficient of $A_\eta(b-x^\ast b)$ is $-\frac43$ times that of $A_\eta(1)$ (Theorem~\ref{thm:gaugesolution}(i)). So the closings miss the nonzero submodule generated by $b-x^\ast b$, the part of the span that the exchanges negate. The computation of (i) agrees: the rows of degree at most $e$ annihilated at $10\le j\le400$ have the dimension $e+1$ for $e\le8$, so by (c) they are exactly the multiples $p(k)(b-x^\ast b)$, $\deg p\le e$; and the solution spaces of the closing equations at $D=17,18,19,20$ have dimensions at most $400,431,462,493$, growing by $31=15\cdot2+1$, the slope of Proposition~\ref{prop:blinddim} for a missed submodule of dimension two.

\emph{(iii) The centre of $\binom{[5]}2$.} Here $K=e_2$, $\Psi=e_3$, $f^{-1}=e_2e_3/(\alpha_1\cdots\alpha_5)$, and $G$ is generated by the bijections of $[5]$ and the Cremona, $|G|=240$: the list has no exchange, its bases having two marks and its cobases three, so this is every symmetry of the four kinds.

\emph{The native gauges.} On the chart $\alpha_5=1$ the defects times $e_2e_3$ generate an ideal whose saturation by the scales and the letters is zero-dimensional and radical of degree $66$ (receipt \texttt{u25.sing}, an exact Gr\"obner basis computation): there are $66$ native gauges, all simple zeros of the defects and so nondegenerate. Five points, checked exactly to be native gauges (receipt \texttt{u25\_core.py}), and their orbits under $G$ account for all $66$; with $\zeta=e^{2\pi i/3}$ and the spanning set $\mathbf b$ found below,
\[
\begin{array}{lccl}
\text{point}&\text{native value}&\text{orbit}&(b_1,\dots,b_5)\text{ at the point}\\ \hline
\one&\frac1{100}&1&(1,1,1,1,1)\\
(1,1,1,1,-1)&\frac14&5&(1,-\frac35,\frac15,1,\frac15)\\
(1,1,1,-1,-1)&\frac14&10&(1,\frac15,-\frac15,1,-\frac15)\\
(\sqrt3-2,\sqrt3-2,\sqrt3-2,1,1)&-\frac18&20&(1,1,-\frac45,\frac{113}5,\frac{23}5)\\
(1,1,\zeta,\zeta,\zeta^2)&1&30&(1,\frac1{10},-\frac15,\frac1{10},-\frac15)
\end{array}
\]
On the chart $\det(-H)$ is $\frac{81}{3125}$, $-3$, $-3$, $-\frac{81}{16}$ and $81$ at these points. The orbits are the table's gauge, the signs with one and with two scales $-1$, the points with three scales $\sqrt3-2$ and two scales $1$ or two and three, which the Cremona joins, and the points with two scales $1$, two $\zeta$ and one $\zeta^2$.

\emph{The box.} The box $\mathfrak B_3(1,1)$ has $10322$ symbols; modulo $G$ its relations at the orbit representatives are $196$ combinations of $114$ orbit averages, of rank $103$, eliminated exactly over $\mathbb Q(k)$. The free Laurent orbits of wall degree at most $2$ give the spanning set
\begin{gather*}
b_1=1,\qquad b_2=e_G\alpha^{(-1,-1,0,1,1)},\qquad b_3=e_G\frac{\alpha_5}{\alpha_1}=\frac1{20}\sum_{i\neq j}\frac{\alpha_i}{\alpha_j},\\
b_4=e_G\alpha^{(-2,-2,0,2,2)},\qquad b_5=e_G\alpha^{(-2,-2,1,1,2)},
\end{gather*}
of wall degrees $0,1,1,2,2$, and the box is accepted at $\mathfrak s=2$. With $q(k)=3k^3+4k^2-61k+130$, its matrix has the rows $\mathbf B_{1\bullet}=(10,30,60,0,0)$, the Laurent identity $f^{-1}=10b_1+30b_2+60b_3$, and
\begin{align*}
\tfrac{q}{k-1}\mathbf B_{2\bullet}&=\Big(27k^2+111k-130,\ \tfrac{2(51k^3+137k^2-756k+580)}{k-2},\ \tfrac{16(9k^3+33k^2-155k+108)}{k-2},\\
&\qquad -3(3k^2-9k-10),\ 12(3k^2-k+2)\Big),\\
\tfrac{q}{k-1}\mathbf B_{3\bullet}&=\Big(6(5k^2+8k-65),\ \tfrac{6(15k^3+6k^2-313k+440)}{k-2},\ \tfrac{4(33k^3+74k^2-971k+1340)}{k-2},\\
&\qquad -12(k-5)(k-2),\ 12(5k^2-14k+5)\Big),
\end{align*}
\begin{align*}
\tfrac{q\,(k-3)^4}{k-1}\mathbf B_{41}&=k(27k^5+245k^4-1789k^3+3709k^2-5388k+5820),\\
\tfrac{q\,(k-3)^4}{k-1}\mathbf B_{42}&=\tfrac{2(45k^7+225k^6-2310k^5+3606k^4+1139k^3-157k^2-7152k+2460)}{k-2},\\
\tfrac{q\,(k-3)^4}{k-1}\mathbf B_{43}&=\tfrac{4(216k^7-552k^6-925k^5-1703k^4+18473k^3-28953k^2+33288k-33444)}{3(k-2)},\\
\tfrac{q\,(k-3)^4}{k-1}\mathbf B_{44}&=3k^6-283k^5+3037k^4-13753k^3+33126k^2-43566k+25020,\\
\tfrac{q\,(k-3)^4}{k-1}\mathbf B_{45}&=-4(k-2)(27k^5-373k^4+1321k^3-1507k^2-450k+1878),\\
\tfrac{q\,(k-3)^3}{k-1}\mathbf B_{51}&=2k(15k^4+52k^3-232k^2+153k-420),\\
\tfrac{q\,(k-3)^3}{k-1}\mathbf B_{52}&=\tfrac{2(45k^6+111k^5-963k^4+461k^3+936k^2+478k-540)}{k-2},\\
\tfrac{q\,(k-3)^3}{k-1}\mathbf B_{53}&=\tfrac{4(117k^6+225k^5-2281k^4+1565k^3+1564k^2-1562k+5028)}{3(k-2)},\\
\tfrac{q\,(k-3)^3}{k-1}\mathbf B_{54}&=-2(6k^5-17k^4-185k^3+1052k^2-2368k+2280),\\
\tfrac{q\,(k-3)^3}{k-1}\mathbf B_{55}&=4(9k^5+32k^4-250k^3+412k^2-65k-522),
\end{align*}
so $Q_0=(k-3)^4(k-2)\,q(k)/3$ and $D_0=8$; it equals the rayreader's matrix.

\emph{Masters.} The five orbit sums of the solutions $A_c$ (Theorem~\ref{thm:invariant}(iii)) have the types $\frac1{100},\frac14,\frac14,-\frac18,1$. Those alone in their type do not vanish at $[1]$. On $b_1,b_2$ the two of the type $\frac14$ have the leading coefficients $5C(1,-\frac35)$ and $10C'(1,\frac15)$, $C,C'\neq0$, and $\det\begin{bmatrix}1&-\frac35\\1&\frac15\end{bmatrix}=\frac45\neq0$, so they are linearly independent on $\mathfrak N$. By Lemma~\ref{lem:types}, $\dim\mathfrak N\ge5$: the spanning set is a set of masters of $\mathfrak N$, of dimension exactly $5$, and $\mathbf B=\mathbf B_{\mathbf b}$. The five rows of values in the table have the determinant $\frac{3888}{625}$, so $\mathbf B(\infty)$ is diagonalizable with the eigenvalues $100,4,4,-8,1$ (Proposition~\ref{prop:rays-spectrum}(i), (ii)); indeed
\[
\mathbf B(\infty)=\begin{bmatrix}10&30&60&0&0\\9&34&48&-3&12\\10&30&44&-4&20\\9&30&96&1&-36\\10&30&52&-4&12\end{bmatrix},\qquad \det(\lambda-\mathbf B(\infty))=(\lambda-100)(\lambda-4)^2(\lambda+8)(\lambda-1)
\]
(receipt \texttt{binf\_exact.py}).

\emph{The law and the line.} The rows $w_i(k)=w_{i-1}(k-1)\mathbf B(k)$, $w_0=(1,0,0,0,0)$, are independent for $i\le3$ and dependent at $i=4$, so the cyclic submodule $\mathfrak N_1$ of $[1]$ has dimension $4$ and the law of $[1]$ (Proposition~\ref{prop:rays-lawrows}(ii)) is $\sum_{j=0}^4p_j(k)\operatorname{Sh}^j[1]=0$ with
\begin{align*}
p_0&=k^4(33k^2+88k+60),\\
p_1&=-(3201k^6+14938k^5+27699k^4+26214k^3+13598k^2+3736k+440),\\
p_2&=-4(2772k^6+18480k^5+48699k^4+63150k^3+40515k^2+11296k+1140),\\
p_3&=16(k+2)^2(7491k^4+34958k^3+51724k^2+24507k+4430),\\
p_4&=-3200(k+2)^2(k+3)^2(33k^2+22k+5).
\end{align*}
Put
\begin{multline*}
m=k(k-5)(k-2)(k+2)\,[b_1]-(3k^4+11k^3-71k^2+227k-210)\,[b_2]\\
-\tfrac23(3k^4-11k^3+119k^2-125k-90)\,[b_3]-(k-5)(k-2)^3\,[b_4]+(k-2)^2(5k^2-14k+5)\,[b_5]
\end{multline*}
and $\tilde m=\frac{k-1}{q(k)}\,m$. With $\tilde w$ the coordinate row of $\tilde m$, the identity $\tilde w(k-1)\mathbf B(k)=4\,\tilde w(k)$ holds exactly, that is $\operatorname{Sh}^{-1}\tilde m=4\tilde m$, so $L=\mathbb C(k)\tilde m$ is a submodule; and at $k=\frac{37}3$ the rows $w_0,\dots,w_3,\tilde w$ have rank $5$. So
\[
\mathfrak N=\mathfrak N_1\oplus L .
\]

\emph{The table one mark larger.} $\binom{[5]}2$ is the contraction at the mark $6$ of $\binom{[6]}3$, which is connected and identically self-dual, so Proposition~\ref{prop:onemark} applies, with $G'$ generated by the $T_\beta$, $\beta$ a bijection of $[6]$, and the Cremona, and $G\subseteq G'$. The native gauges lift by $\alpha_6=\Psi/K$: the five points of the table to $\alpha_6=1,-1,1,1,\zeta^2$ (receipt \texttt{u36\_lift.py}). The lifts are native gauges of the centre of $\binom{[6]}3$, which has $66$ of them, all simple (receipt \texttt{u36.sing}), in orbits of $S_6$ and the Cremona of $1,15,20,30$ points with the native values $\frac1{400},\frac1{16},-\frac1{32},\frac14$, a quarter of those below them: the two kinds of signs lift into one orbit of $15$, the signs of $\binom{[6]}3$ with two scales $-1$, and this is Proposition~\ref{prop:onemarkhidden}(ii) at $r=2$, with $C_2=2$.

\emph{The element $\ell$.} Put $S=\sum_{i\neq j\le5}\alpha_i/\alpha_j=20\,b_3$ and $T=\frac K\Psi\sum_i\alpha_i+\frac\Psi K\sum_i\alpha_i^{-1}$. Since $S$ is invariant under $G$, the image of the stabilizer of the mark $6$, and the transpositions $(e\,6)$, $e\le5$, with the identity, represent its cosets, $e_{G'}[S]=\frac16\big([S]+\sum_{e\le5}T_{(e\,6)}[S]\big)$. By Proposition~\ref{prop:onemark}(iii), with $\frac{\Gamma(k+1)\Gamma(k-1)}{\Gamma(k)^2}=\frac k{k-1}$,
\[
T_{(e\,6)}\Big[\frac{\alpha_i}{\alpha_j}\Big]=\Big[\frac{\alpha_i}{\alpha_j}\Big]\ (e\notin\{i,j\}),\qquad \frac k{k-1}\Big[\frac\Psi{K\alpha_j}\Big]\ (e=i),\qquad \frac k{k-1}\Big[\frac{K\alpha_i}\Psi\Big]\ (e=j).
\]
Each pair $(i,j)$ is left alone by three of the five $e$, and each $\alpha_j^{-1}$ and each $\alpha_i$ arises four times, so $\sum_eT_{(e\,6)}[S]=3[S]+\frac{4k}{k-1}[T]$, $e_{G'}[S]=\frac23[S]+\frac{2k}{3(k-1)}[T]$, and
\[
\ell=3(1-e_{G'})[S]=[S]-\frac{2k}{k-1}\,[T].
\]
The symbols of $T$, at the levels $(0,1)$ and $(1,0)$ and of wall degrees $0$ and $1$, lie in the box, which reduces $T$ to the spanning set (receipt \texttt{u25\_core.py}):
\[
[T]=\frac{10(k-1)}k\,[b_3]-\frac{10}{k\,q(k)}\,m .
\]
So
\[
\ell=20[b_3]-\frac{2k}{k-1}\Big(\frac{10(k-1)}k[b_3]-\frac{10}{k\,q(k)}m\Big)=\frac{20}{(k-1)\,q(k)}\,m=\frac{20}{(k-1)^2}\,\tilde m .
\]
Since $e_{G'}(1-e_{G'})=0$, $e_{G'}\ell=0$, hence $e_{G'}\tilde m=0$ and $L\subseteq(1-e_{G'})\mathfrak M^G$: the closings, the values and $A_\star$ vanish on $L$ (Proposition~\ref{prop:onemark}(iv)). The values can be seen to vanish directly as well: by the formula for $\pi_\ast$ in Proposition~\ref{prop:onemark}, $\pi_\ast(\alpha_i/\alpha_6)=\frac k{k-1}\frac K\Psi\alpha_i$ and $\pi_\ast(\alpha_6/\alpha_i)=\frac k{k-1}\frac\Psi K\alpha_i^{-1}$, so $\ell=\pi_\ast\ell'$ with $\ell'=\sum_{i\neq j\le5}\frac{\alpha_i}{\alpha_j}-2\sum_{i\le5}\big(\frac{\alpha_i}{\alpha_6}+\frac{\alpha_6}{\alpha_i}\big)$, three times the part of $\sum_{i\neq j\le5}\alpha_i/\alpha_j$ that the average over the bijections of $[6]$ removes; the values of $\binom{[6]}3$ are kept by those bijections (Proposition~\ref{prop:raysym}(ii)), so $\langle\ell'\rangle'_k=0$ and $\langle\ell\rangle_k=\Gst21{k,\ k}{2k}^{-1}\langle\ell'\rangle'_k=0$ (Proposition~\ref{prop:onemark}(i)). As a numerical check, a trapezoidal quadrature over the chart gives $|\langle\ell\rangle_k|/\langle1\rangle_k$ equal to $1.0\cdot10^{-13}$ at $k=\frac{13}2$ and $2.8\cdot10^{-13}$ at $k=8$, the accuracy of the quadrature (receipt \texttt{ell\_check.py}). So $L$ is a relation among the values of the five masters that the relations of the ray do not produce, $m$ being nonzero in $\mathfrak N$, nor its symmetries of the four kinds, $m$ lying in $\mathfrak M^G$.

\emph{What sees $L$.} A solution $V$ of the type $(\lambda,\rho+\mathbb Z)$ has $\varsigma^{-1}V(\tilde m)=V(\operatorname{Sh}^{-1}\tilde m)=4V(\tilde m)$, which forces $V(\tilde m)=0$ unless $\lambda^{-1}=4$: the solutions at the orbits of the native values $\frac1{100}$, $-\frac18$ and $1$ vanish on $L$. At a sign, $A_c(b_l)/A_c(1)=b_l(\alpha_c)+O(1/k)$ (Theorem~\ref{thm:gaugesolution}(i)), and the coefficients of $k^4$ in $m$ are $(1,-3,-2,-1,5)$, so $A_c(m)/A_c(1)$ begins with
\[
\big(1+\tfrac95-\tfrac25-1+1\big)k^4=\tfrac{12}5k^4\quad\text{(one scale }-1\text{)},\qquad \big(1-\tfrac35+\tfrac25-1-1\big)k^4=-\tfrac65k^4\quad\text{(two scales }-1\text{)}:
\]
the solutions at both kinds of signs see $L$. With the solutions at the fifteen signs transported from one of them by $G'$ as in Proposition~\ref{prop:onemarkhidden}(i), $A_c(1)$ is the same for all fifteen, $T_\gamma[1]=[1]$; their sum, invariant under $G'$, has on $m$ the leading coefficient $5\cdot\frac{12}5-10\cdot\frac65=0$ times that of $A_c(1)$, as it must, while the mean over the five less the mean over the ten has $\frac{12}5+\frac65=\frac{18}5$ and sees $L$.

\emph{The closings.} The rows of degree at most $e$ annihilated by the closings at $10\le j\le150$ have the dimension $\max(0,e-3)$ for $e\le24$, which is the dimension of the multiples $p(k)\,w_m$, $\deg p\le e-4$, of the primitive coordinate row $w_m$ of $3m$, of degree $4$; by (c) the closings annihilate $L$ and no element of $\mathfrak N$ outside $L$ with polynomial coordinates of degree at most $24$. The solution spaces of the closing equations at $D=7,8,9,10$ have dimensions at most $24,30,36,42$, growing by $6=5\cdot1+1$, the slope of Proposition~\ref{prop:blinddim} for a missed line; a step of more than one already shows that something is missed.

\emph{The centre of $\binom{[6]}3$ by its box.} Modulo $S_6$ and the Cremona, of order $1440$, the box $\mathfrak B_3(1,1)$, with $78296$ symbols, has $406$ relations at orbit representatives on $226$ orbit averages, of rank $186$; it is accepted at $\mathfrak s=2$ with the spanning set $1$, $e\,\alpha^{(-1,-1,-1,1,1,1)}$, $e\,\alpha^{(-1,-1,0,0,1,1)}$, $e\,\alpha^{(-2,-2,-2,2,2,2)}$, $e$ the average over that group, and $Q_0=(k-3)^5(k-2)^2(k+1)(k^3-8k^2+\frac{71}3k-\frac{76}3)$, $D_0=11$. The four orbits of native gauges give solutions of four distinct types, none vanishing at $[1]$, so the spanning set is a set of masters of a submodule $\mathfrak N'$ of dimension $4$ (Lemma~\ref{lem:types}), and $\mathbf B'(\infty)$ has the simple spectrum $400,16,4,-32$. Its rows of $[1]$ are independent up to $h=4$, so $\mathfrak N'$ is the cyclic submodule of $[1]$. Since $\pi_\ast[1]=[1]$ and $\pi_\ast\operatorname{Sh}'=\frac k{2(2k+1)}\operatorname{Sh}\pi_\ast$, $\pi_\ast\mathfrak N'=\mathfrak N_1$, and the law of $[1]$ at the centre of $\binom{[6]}3$ is the law of $\binom{[5]}2$ for $\Gst21{k,\ k}{2k}\langle1\rangle_k$: its coefficients are proportional to $p_i(k)\prod_{l=i}^{3}\frac{k+l}{2(2k+2l+1)}$, which is the law that the box gives, exactly (receipt \texttt{u36\_lift.py}). So
\[
\mathfrak N\cap\mathfrak M^{G'}=\mathfrak N_1=\pi_\ast\mathfrak N',\qquad \mathfrak N\cap(1-e_{G'})\mathfrak M^G=L ,
\]
$\mathfrak N_1$ lying in $\mathfrak M^{G'}$, which contains $[1]$, and $L$ in $(1-e_{G'})\mathfrak M^G$. At the centre of $\binom{[6]}3$ the closing equations have the dimensions $0,1,2,3$ at $D=10,\dots,13$, and no nonzero row of degree at most $40$ is annihilated at $10\le j\le200$: the closings there single out $\mathbf B'$.

So: at the small centres of (i) the closings alone single out the box's matrix at the degree $D_0$, with nothing of low degree missed; at the centre of $\mathrm N_6$ without its exchanges they miss the part that the exchanges negate; and at the centre of $\binom{[5]}2$, with every symmetry of the four kinds, the span of the box is the cyclic submodule of $[1]$, which the box of $\binom{[6]}3$ reads, plus the line $L$ that the symmetric group of $\binom{[6]}3$ hides from the closings, the values and $A_\star$, and that the signs see.
\end{solution}

\chapter{Exercises on rays: other lists}\label{ch:exlists}

These exercises read rays of lists other than the lists $\binom{\cE}1$: the list $\mathrm H_4$ on four marks at its centre and on a ray strictly inside its walls whose native value has a conjugate on the slit, and the five-mark lists along their hinge trees. Each exercise defines its list by its marks and its bases; every weight is $1$. A recurrence of the closings is proved for every index by an explicit telescoping identity, and Theorems~\ref{thm:closings} and~\ref{thm:onekernel}, with Corollary~\ref{cor:valuemodule}, carry it to the values. A \emph{certificate} of a recurrence is an explicit element of the relation space $\mathcal L$ (Definition~\ref{def:raygauge}) that equals the recurrence applied to an insertion, as in Proposition~\ref{prop:rays-strength}(i); a telescoping identity among binomial sums is not one. Numbers are exact; a decimal is a rendering of an exact value, truncated. The receipts are the scripts \texttt{ex28\_h4\_centre.py}, \texttt{ex29\_h4\_lean.py} and \texttt{ex30\_trees.py}, with the shared tools in \texttt{common.py}; each prints the numbers its solution displays.

\section{The list $\mathrm H_4$}\label{sec:lista-h}

In this section $\mathrm H_4$ is the list on the marks $a,b,2,3$ whose bases are the five pairs other than $\{a,b\}$: $\{a,2\},\{a,3\},\{b,2\},\{b,3\},\{2,3\}$. With $A=\alpha_a+\alpha_b$ and $B=\alpha_2+\alpha_3$ its letters are
\[
K=AB+\alpha_2\alpha_3,\qquad \Psi=AB+\alpha_a\alpha_b ,
\]
the second from the cobases $\{b,3\},\{b,2\},\{a,3\},\{a,2\},\{a,b\}$, and $\Psi(\alpha)=\alpha_a\alpha_b\alpha_2\alpha_3\,K(\alpha^{-1})$. The pair $\{a,b\}$ is a parallel class.

\begin{exercise}[$\mathrm H_4$ at its centre]\label{ex:lista-hcentre}
Let $P(k)=\langle1\rangle_k=P_{\mathrm H_4}(k\one;k,k)$ be the values of the centre ray of $\mathrm H_4$ at the strength one (Definition~\ref{def:centreray}).
\begin{enumerate}[label=(\roman*)]
\item Write the stride, the walls of the direction and $V_1$.
\item Find every native gauge of the ray with its native value, and decide which are nondegenerate.
\item Find the closings and a law of the values of order four, prove the law, and give its symbol.
\item Write $\Gamma(\sigma)^2P(\sigma)$ on the share side of the hinge tree of $\mathrm H_4$, close it at the integer strengths, and decide which of its two closings converges.
\item Read the low-end relation, the principal parts of $P$ at $k=0$ and at $k=-1$, and the relation that the law gives at $k=-1$.
\item Give $P(1)$ and $P(2)$, and show that every value $P(k)$, $k\ge1$, lies in $\mathbb Q+\mathbb Q\,P(1)+\mathbb Q\,P(2)$.
\end{enumerate}
\end{exercise}

\begin{solution}
\emph{(i) Stride, walls, $V_1$.} By Definition~\ref{def:centreray} the stride and its reciprocal are
\[
f=\frac{\alpha_a\alpha_b\alpha_2\alpha_3}{K\Psi}=\frac1{K(\alpha)\,K(\alpha^{-1})},\qquad f^{-1}=\frac{(AB+\alpha_2\alpha_3)(AB+\alpha_a\alpha_b)}{\alpha_a\alpha_b\alpha_2\alpha_3},
\]
and the walls of the direction are the widths, $\kappa_{1,S}=c(S)=r(S)+r(\cE\setminus S)-2$ (Definition~\ref{def:skeleton}, Proposition~\ref{prop:centrefamily}(ii)). A single mark has rank one and its complement rank two; $\{a,b\}$ has rank one and $\{2,3\}$ rank two; the other pairs have rank two with complements of rank two; a triple has rank two and its complement rank one. So $\kappa_{1,S}=2$ for the four pairs $\{a,2\},\{a,3\},\{b,2\},\{b,3\}$, and $\kappa_{1,S}=1$ for every other nonempty proper $S$: all walls are positive. In $V_1=\sum_\pi\prod_{l=2}^4\kappa_{1,U_l}^{-1}$ (Proposition~\ref{prop:principal}) the tails $U_2$ and $U_4$ have three marks and one mark, of wall $1$, so an order contributes $\kappa_{1,U_3}^{-1}$, $U_3$ being its last two marks. Each of the six pairs ends four of the $24$ orders, so
\[
V_1=4\,(1+1)+4\cdot4\cdot\tfrac12=16 .
\]

\emph{(ii) Native gauges.} By Proposition~\ref{prop:rays-centrenative} a point $\alpha$ of $X_F$ is a native gauge exactly when the occupations $R_e=\alpha_e\partial_{\alpha_e}\log K$ agree at $\alpha$ and at $\alpha^{-1}$, and its native value is $\big(K(\alpha)K(\alpha^{-1})\big)^{-1}$. Here
\[
R_a(\alpha)=\frac{\alpha_aB}K,\quad R_a(\alpha^{-1})=\frac{\alpha_bB}\Psi,\qquad R_2(\alpha)=\frac{\alpha_2(A+\alpha_3)}K,\quad R_2(\alpha^{-1})=\frac{\alpha_3A+\alpha_a\alpha_b}\Psi,
\]
and $R_b$, $R_3$ follow by exchanging $a\leftrightarrow b$ and $2\leftrightarrow3$. So the native gauges are the solutions in $X_F$ of
\begin{gather*}
B(\alpha_a\Psi-\alpha_bK)=0,\qquad B(\alpha_b\Psi-\alpha_aK)=0,\\
\alpha_2(A+\alpha_3)\Psi=(\alpha_3A+\alpha_a\alpha_b)K,\qquad \alpha_3(A+\alpha_2)\Psi=(\alpha_2A+\alpha_a\alpha_b)K .
\end{gather*}
The occupations add up to the rank $2$ at $\alpha$ and at $\alpha^{-1}$, so any three of the four equations imply the fourth.

If $B=0$, then $K=\alpha_2\alpha_3$ and $\Psi=\alpha_a\alpha_b$, and with $\alpha_3=-\alpha_2$ the third equation reads $\alpha_2A(\alpha_a\alpha_b-\alpha_2^2)=0$: either $A=0$ or $\alpha_a\alpha_b=\alpha_2^2$. If $B\neq0$, the first two equations give $(\alpha_a-\alpha_b)(\Psi+K)=0$ and $A(\Psi-K)=0$. For $A=0$ they leave $K=\alpha_2\alpha_3$, $\Psi=\alpha_a\alpha_b=-\alpha_a^2$ and $\Psi+K=0$, that is $\alpha_a^2=\alpha_2\alpha_3$, and the third equation, $\alpha_2\alpha_3\Psi=\alpha_a\alpha_bK$, holds there. For $A\neq0$ they give $\Psi=K$, that is $\alpha_a\alpha_b=\alpha_2\alpha_3$, and then $\alpha_a=\alpha_b$, since $\Psi+K=2K\neq0$; the third equation becomes $A(\alpha_2-\alpha_3)=0$, so $\alpha_2=\alpha_3$ and $\alpha_2^2=\alpha_a^2$. Up to scale the native gauges are therefore:
\begin{itemize}
\item the table's gauge $\one$, where $K=\Psi=5$: native value $\frac1{25}$;
\item the sign $(1,1,-1,-1)$, where $K=\Psi=-3$: native value $\frac19$;
\item the three curves $\{A=0,\ \alpha_a^2=\alpha_2\alpha_3\}$, $\{B=0,\ \alpha_2^2=\alpha_a\alpha_b\}$ and $\{A=B=0\}$, on which $K=\alpha_2\alpha_3$ and $\Psi=\alpha_a\alpha_b$: native value $1$.
\end{itemize}
No letter vanishes at these points of the torus, so they lie in $X_F$. On the chart $y_3=0$, in the coordinates $y_a,y_b,y_2$, the Hessian of $-\log f$ at $\one$ is twice the covariance matrix of the indicators under the uniform weight law on the five bases (Theorem~\ref{thm:branch}), and at the sign it is twice the covariance matrix under the signed measure $\mathbb P_\eta$ with the masses $\frac13,\frac13,\frac13,\frac13,-\frac13$ on $\{a,2\},\{a,3\},\{b,2\},\{b,3\},\{2,3\}$ (Theorem~\ref{thm:signgauges}(ii)):
\[
H_{\one}=\frac2{25}\begin{bmatrix}6&-4&-1\\-4&6&-1\\-1&-1&6\end{bmatrix},\quad \det H_{\one}=\frac{32}{625};\qquad H_\eta=\frac29\begin{bmatrix}2&-4&1\\-4&2&1\\1&1&2\end{bmatrix},\quad\det H_\eta=-\frac{32}{81}.
\]
Both are nonsingular: the table's gauge and the sign are nondegenerate native gauges (Definition~\ref{def:raygauge}). The three curves are sources of dimension one (Definition~\ref{def:source}), along which $f$ is constant and the Hessian is singular.

\emph{(iii) Closings and the law.} The origin is $0$ and both letters have the exponent one, so the closings are $\mathsf C_j(1)=\CT(f^{-j})$ (Definition~\ref{def:closings}). In
\[
(K\Psi)^j=\sum_{u,v}\binom ju\binom jv(AB)^{u+v}(\alpha_2\alpha_3)^{j-u}(\alpha_a\alpha_b)^{j-v}
\]
the monomial $(\alpha_a\alpha_b\alpha_2\alpha_3)^j$ takes $\alpha_a^v\alpha_b^v$ from $A^{u+v}$ and $\alpha_2^u\alpha_3^u$ from $B^{u+v}$, so $u=v$ and
\[
\mathsf C_j(1)=\sum_{u=0}^jF(j,u),\qquad F(j,u)=\binom ju^2\binom{2u}u^2;\qquad \mathsf C_j(1)=1,\ 5,\ 53,\ 761,\ 12661,\ 229705,\ \dots
\]
Exact linear algebra on the first $150$ closings gives one line of recurrences $\sum_{i=0}^4p_i(j)\,\mathsf C_{j+i}(1)=0$ with $\deg p_i\le7$, spanned by
\begin{align*}
p_0(j)&=225(j+1)^2(j+2)^2\,w(j),\qquad w(j)=144j^3+1296j^2+3853j+3790,\\
p_1(j)&=-2(j+2)^2(34848j^5+435600j^4+2135762j^3+5129622j^2+6027508j+2768385),\\
p_2(j)&=2(21168j^7+381024j^6+2917047j^5+12308044j^4\\
&\qquad+30893608j^3+46094457j^2+37814839j+13140768),\\
p_3(j)&=-2(j+3)(2592j^6+42768j^5+289602j^4+1028216j^3+2015044j^2+2063287j+861336),\\
p_4(j)&=(j+3)(j+4)^3(144j^3+864j^2+1693j+1089);
\end{align*}
no recurrence of order at most three and degree at most $30$ holds on the first $200$ closings. The recurrence holds for every $j\in\mathbb N$, by a telescoping identity. For $0\le u\le j+4$ put
\[
T(j,u)=\frac{((j+4)!)^2\,((2u)!)^2}{(u!)^6\,((j+4-u)!)^2}>0,\qquad G(j,u)=\frac{T(j,u)\,A(j,u)}{\prod_{l=1}^4(j+l)^2},\qquad G(j,j+5)=0,
\]
with $A(j,u)=u^4(j+1)^2(j+2)^2\sum_{q=0}^4Q_q(j)\,u^q$ and
\begin{align*}
Q_0&=(j+4)(6912j^6+125568j^5+938496j^4+3695236j^3+8085594j^2+9321748j+4422321),\\
Q_1&=-4(5760j^6+106272j^5+809680j^4+3260831j^3+7321143j^2+8687769j+4256385),\\
Q_2&=2(13248j^5+198864j^4+1186148j^3+3514356j^2+5173077j+3027087),\\
Q_3&=-4(22j+63)\,w(j),\qquad Q_4=15\,w(j).
\end{align*}
For $0\le u\le j+4$, $F(j+i,u)=T(j,u)\prod_{l=i+1}^4\big((j+l-u)/(j+l)\big)^2$, which vanishes for $u>j+i$ as $F$ does, and $T(j,u+1)=T(j,u)\,4(2u+1)^2(j+4-u)^2/(u+1)^4$, which vanishes at $u=j+4$. So, divided by $T(j,u)\big/\big((u+1)^4\prod_l(j+l)^2\big)$, the identity
\[
\sum_{i=0}^4p_i(j)\,F(j+i,u)=G(j,u+1)-G(j,u)\qquad(0\le u\le j+4)
\]
is the identity of polynomials
\[
(u+1)^4\sum_{i=0}^4p_i(j)\prod_{l\le i}(j+l)^2\prod_{l>i}(j+l-u)^2=4(2u+1)^2(j+4-u)^2A(j,u+1)-(u+1)^4A(j,u),
\]
which holds in $\mathbb Q[j,u]$, both sides expanded. Summed over $0\le u\le j+4$, the right side telescopes to $G(j,j+5)-G(j,0)=0$, as $A(j,0)=0$, and the left side is $\sum_ip_i(j)\,\mathsf C_{j+i}(1)$. So the closings satisfy the recurrence for every $j$.

The direction has the integral origin $0$, both letter exponents $1$ and positive walls. Written as $\sum_ic_i(-j)\,\mathsf C_{j-i}(1)=0$ with $c_i(k)=-p_{4-i}(-k-4)$, the recurrence says that the closings vanish on $\sum_ic_i(k)\operatorname{Sh}^i[1]$ (Theorem~\ref{thm:closings}); by Theorem~\ref{thm:onekernel} the values vanish on it as well, and the continuation of its value vanishes identically; Corollary~\ref{cor:valuemodule} states the same for the recurrences of $\langle1\rangle$. With $\langle f^ig\rangle_k=\langle g\rangle_{k+i}$ (Proposition~\ref{prop:rays-strength}(ii)) this is the law
\[
\sum_{i=0}^4c_i(k)\,P(k+i)=0,
\]
an identity of meromorphic functions of $k$, with
\begin{align*}
c_0(k)&=k^3(k+1)(144k^3+864k^2+1693k+1075),\\
c_1(k)&=-2(k+1)(2592k^6+19440k^5+56322k^4+80296k^3+60004k^2+23017k+3580),\\
c_2(k)&=2(21168k^7+211680k^6+884919k^5+2003456k^4\\
&\qquad+2646984k^3+2031207k^2+830703k+137260),\\
c_3(k)&=-2(k+2)^2(34848k^5+261360k^4+741842k^3+984642k^2+598948k+127215),\\
c_4(k)&=225(k+2)^2(k+3)^2(144k^3+432k^2+397k+102).
\end{align*}
Its symbol (Corollary~\ref{cor:symbol}) is read off the coefficients of $k^7$:
\[
\chi(\lambda)=144-5184\lambda+42336\lambda^2-69696\lambda^3+32400\lambda^4=144(25\lambda-1)(9\lambda-1)(\lambda-1)^2 .
\]
Its roots are the native values $\frac1{25}$ and $\frac19$ of the two nondegenerate native gauges and the native value $1$ of the curves, doubled; the root $\frac1{25}=f(\one)$ is forced by Proposition~\ref{prop:nativeorder}.

\emph{(iv) The share side and its closings.} A set of two marks has width one exactly for $\{a,b\}$ and $\{2,3\}$ (part (i)), so $\mathrm H_4$ is the two-sum, at a hinge $g$, of the parallel cell $\mathcal P$ of $\{a,b\}$, with the core $\binom{\{g,a,b\}}1$, and the series cell $\mathcal S$ of $\{2,3\}$, with the core $\binom{\{g,2,3\}}2$: the bases $\{2,3\}$ and $\{x,y\}$, $x\in\{a,b\}$, $y\in\{2,3\}$, are those of the two-sum (Definition~\ref{def:hingetree}, Theorem~\ref{thm:primetree}). At the centre of real strength $\sigma>0$, with the masses $\sigma$ and $s=t=\sigma$, every mark is a one-mark cell of ratio $1$, whose transform and dual transform are $\mathsf g(y)=\sG_{0;(1,-1)}(1,y)$ (Definition~\ref{def:treeproc}(1)), and Theorem~\ref{thm:seriesparallel} gives
\[
f_{\mathcal S}=\mathsf g^2,\quad \kappa_{\mathcal S}=2\sigma-\sigma=\sigma;\qquad \check f_{\mathcal P}=\mathsf g^2,\quad \kappa_{\mathcal P}=2\sigma-\sigma=\sigma .
\]
The root of two one-mark leaves at the strength $z$ (Definition~\ref{def:treeproc}(5)), whose transforms are again $\mathsf g$, and the Beta integral give
\begin{gather*}
\int_0^\infty y^{z-1}\mathsf g(y)^2\,dy=\Gamma(z)^2P_2(z),\\
P_2(z)=P_{\binom{[2]}1}(z\one;z,z)=\int_0^\infty x^{z-1}(1+x)^{-2z}\,dx=\frac{\Gamma(z)^2}{\Gamma(2z)}\qquad(z>0).
\end{gather*}
By Proposition~\ref{prop:faces}(iii), with $s=t=\kappa=\sigma$, the two Mellin transforms of the laws are $\int\varrho_{\mathcal S}^{-z}d\phi_{\mathcal S}=P_2(z)$ and $\int\varrho_{\mathcal P}^{z-\sigma}d\phi_{\mathcal P}=P_2(z)$ for $z>0$. Theorem~\ref{thm:glue}(iv) for one table, with $\mathcal S$ in the place of $\mathcal A$ and $\mathcal P$ in the place of $\mathcal A'$, on a line $\re z=c$ with $0<c<\sigma$, gives
\[
\Gamma(\sigma)^2P(\sigma)=\frac1{2\pi i}\int_{(c)}\Gamma(z)^2\,\Gamma(\sigma-z)^2\,P_2(z)^2\,dz=\frac1{2\pi i}\int_{(c)}\mathcal F(z,\sigma)\,dz,\qquad \mathcal F=\frac{\Gamma(z)^6\,\Gamma(\sigma-z)^2}{\Gamma(2z)^2},
\]
and both sides are holomorphic in $\sigma$ on $\re\sigma>c$.

Closing (Theorem~\ref{thm:closing}) moves the line across the poles of the Gammas on one side and needs the contribution of the moved line to tend to zero. At an integer strength $k\ge1$, by the shift and the reflection of Proposition~\ref{prop:stacked}(iii),(v),
\begin{gather*}
\Gamma(z)\,\Gamma(k-z)=\frac{\pi\,(z-1)(z-2)\cdots(z-k+1)}{\sin\pi(z-k+1)},\\
\mathcal F(z,k)=\frac{\pi^2\,V_k(z)\,P_2(z)^2}{\sin^2\pi(z-k+1)},\qquad V_k(z)=\prod_{j=1}^{k-1}(z-j)^2 .
\end{gather*}
On the line $\re z=k+N+\frac12$, $|\sin\pi(z-k+1)|=\cosh(\pi\im z)$, $|P_2(z)|\le\int_0^1\big(w(1-w)\big)^{\re z-1}dw\le4^{1-\re z}$ and $V_k(z)$ is a polynomial, so the contribution of that line is $O(N^{2k-2}16^{-N})$ and tends to zero. To the right the poles are the integers $x\ge k$, double, with the principal part $(z-x)^{-2}$ of $\pi^2/\sin^2$ and no simple part (Lemma~\ref{lem:resonantpair}), and closing to the right gives
\begin{gather*}
\Gamma(k)^2P(k)=-\sum_{x\ge k}\big(V_kP_2^2\big)'(x),\\
P(1)=\sum_{m\ge0}\frac{4\,(m!)^4}{((2m+1)!)^2}\big(h_{2m+1}-h_m\big)=4+\frac5{54}+\frac{47}{13500}+\frac{319}{2058000}+\cdots,
\end{gather*}
with $h_n=\sum_{i=1}^ni^{-1}$, since $P_2'/P_2=2(\log\Gamma)'(x)-2(\log\Gamma)'(2x)$ and $(\log\Gamma)'(n+1)-(\log\Gamma)'(1)=h_n$ (Proposition~\ref{prop:stacked}(vi)). The terms are rational and decrease like $16^{-x}$: the share side of $\mathrm H_4$ closes on the lattice of the integers to its right.

To the left the poles are at $z=-m$, $m\in\mathbb N$, of order four: $\Gamma(z)^6$ has order six there and $\Gamma(2z)^{-2}$ a double zero. With the duplication formula $\Gamma(2z)=\pi^{-1/2}2^{2z-1}\Gamma(z)\Gamma(z+\frac12)$ and the reflection,
\[
\mathcal F(z,\sigma)=4\pi^3\,16^{-z}\,\frac{\cos^2\pi z}{\sin^4\pi z}\,\mathcal B(z)^2,\qquad \mathcal B(z)=\frac{\Gamma(\frac12-z)\,\Gamma(\sigma-z)}{\Gamma(1-z)^2},
\]
so the residue at $-m$ is $16^m\,\mathcal B(-m)^2\big(c+O(m^{-1})\big)$ with $c=\frac2{3\pi}\log16\,\big(2\pi^2-(\log16)^2\big)=7.0911\ldots$: it is $4\pi^{-1}16^m/3!$ times the third derivative at $0$ of $16^{-\varepsilon}\mathcal B(\varepsilon-m)^2(\pi\varepsilon)^4\cos^2(\pi\varepsilon)\sin^{-4}(\pi\varepsilon)$, in which the derivatives of $\log\mathcal B$ at $-m$ are $O(m^{-1})$; and $\mathcal B(-m)^2$ grows like $m^{2\sigma-3}$. At $\sigma=1$ the ratios of consecutive residues are $14.5771$, $15.2453$, $15.4869$, $15.6114$ at $m=10,20,30,40$, against $16\big((m+\frac12)/(m+1)\big)^2=14.5785$, $15.2472$, $15.4880$, $15.6121$, and the residue over $16^m\mathcal B(-m)^2$ is $7.1273$, $7.1173$, $7.1104$, $7.1063$. The residues grow like $16^m$, and the closing to the left diverges.

\emph{(v) The low end.} By (iii) the law holds as an identity of meromorphic functions. Near $k=0$ the terms $i\ge1$ are holomorphic, $c_0$ vanishes to order three, $[k^3]c_0=1075$, and $\lim_{k\to0}k^3P(k)=V_1\,\mathsf C_0(1)=16$ (Proposition~\ref{prop:principal}); the value of the law at $0$ is the low-end relation, by the argument of Corollary~\ref{cor:rays-lowend}(ii), which uses of the law only that it holds as an identity of meromorphic functions:
\[
17200-7160\,P(1)+274520\,P(2)-1017720\,P(3)+826200\,P(4)=0 .
\]
At $k=-1$ the coefficients $c_0$ and $c_1$ vanish, simply, while $c_2(-1)=-342$, $c_3(-1)=4842$, $c_4(-1)=-6300$, so the value of the law at $-1$ is
\[
\CT_{k=-1}\big(c_0(k)P(k)+c_1(k)P(k+1)\big)-342\,P(1)+4842\,P(2)-6300\,P(3)=0,
\]
and the first term needs the principal parts of $P$ at $-1$ and at $0$. The share side gives them. On the line $\re z=\frac12$ the bound $|\Gamma(c'+i\tau)|\le C_{c'}e^{-\pi|\tau|/2}(1+|\tau|)^{c'-1/2}$, with $1/\Gamma(2z)=\Gamma(1-2z)\sin(2\pi z)/\pi$, makes $|\mathcal F(z,\sigma)|$ at most a constant times $(1+|\im z|)^{2\re\sigma}e^{-2\pi|\im z|}$ for $|\im z|\ge1$, uniformly for $\sigma$ in compact sets; so $J(\sigma)=\frac1{2\pi i}\int_{(1/2)}\mathcal F\,dz$ is holomorphic wherever no pole of $\Gamma(\sigma-z)^2$ lies on the line, in particular on the strips $|\re\sigma|<\frac12$ and $-\frac32<\re\sigma<-\frac12$. For $\re\sigma>\frac12$, $\Gamma(\sigma)^2P(\sigma)=J(\sigma)$. As $\sigma$ moves left the double poles $z=\sigma$ and then $z=\sigma+1$ of $\Gamma(\sigma-z)^2$ cross the line, and the identity of meromorphic functions that continues it is
\begin{gather*}
\Gamma(\sigma)^2P(\sigma)=J(\sigma)-R_0(\sigma)\qquad(|\re\sigma|<\tfrac12),\\
\Gamma(\sigma)^2P(\sigma)=J(\sigma)-R_0(\sigma)-R_1(\sigma)\qquad(-\tfrac32<\re\sigma<-\tfrac12),
\end{gather*}
with the residues, by Proposition~\ref{prop:resonance} at $p=2$,
\[
R_m(\sigma)=\operatorname*{Res}_{z=\sigma+m}\mathcal F=\frac{\Gamma(\sigma+m)^6}{(m!)^2\,\Gamma(2\sigma+2m)^2}\Big(6(\log\Gamma)'(\sigma+m)-4(\log\Gamma)'(2\sigma+2m)-2(\log\Gamma)'(m+1)\Big).
\]
$\Gamma(\sigma)^{-2}$ vanishes to order two at $0$ and at $-1$, so $J$ enters $P$ there only from the order $(\sigma-\sigma_0)^2$ on, and the principal parts of $P$ are those of $-R_0\,\Gamma^{-2}$ and of $-(R_0+R_1)\,\Gamma^{-2}$. Expanding every Gamma factor at its pole or regular point by Proposition~\ref{prop:resonance}, in $(\log\Gamma)'(1)$ and the Trailmoments $\Tx{n}{\IZ}{0}$ of the counting input, gives $R_0(x)=-16x^{-5}-32(\log\Gamma)'(1)\,x^{-4}+O(x^{-3})$ near $0$ and $(R_0+R_1)(-1+x)=-80x^{-5}+O(x^{-4})$ near $-1$, and, the terms in $(\log\Gamma)'(1)$ cancelling,
\begin{gather*}
P(x)=\frac{16}{x^3}-\frac{24\,\Tx{2}{\IZ}{0}}{x}+O(1),\\
P(-1+x)=\frac{80}{x^3}-\frac{152}{x^2}-\frac{120\,\Tx{2}{\IZ}{0}}{x}+O(1),
\end{gather*}
with $\Tx{2}{\IZ}{0}=\sum_{n\ge1}n^{-2}=\frac16\Gamma(\frac12)^4$ (Lemma~\ref{lem:resonantpair} at $a=\frac12$). The top coefficients are $V_1\,\mathsf C_0(1)=16$ and $V_1\,\mathsf C_1(1)=80$ (Proposition~\ref{prop:principal}), and the coefficient $0$ of $x^{-2}$ at $0$ is the low-end coefficient $\mathsf L_1$, which vanishes for a list (Proposition~\ref{prop:rays-faces}, every weight being $1$). With
\[
c_0(-1+x)=-102x-91x^2+453x^3+O(x^4),\qquad c_1(-1+x)=510x-514x^2-3056x^3+O(x^4),
\]
and $T=\Tx{2}{\IZ}{0}$, the coefficients of $x^{-2}$ and $x^{-1}$ in $c_0(k)P(k)+c_1(k)P(k+1)$ at $k=-1+x$ are $-102\cdot80+510\cdot16=0$ and $-102\cdot(-152)-91\cdot80-514\cdot16=0$, as the law asks, and its constant term is
\[
-102\cdot(-120T)-91\cdot(-152)+453\cdot80+510\cdot(-24T)-514\cdot0-3056\cdot16=1176 :
\]
the multiples of $\Gamma(\frac12)^4$ cancel. Hence $-342\,P(1)+4842\,P(2)-6300\,P(3)=-1176$, that is
\[
57\,P(1)-807\,P(2)+1050\,P(3)=196,
\]
the second low-end relation.

\emph{(vi) Values and their span.} The share side at $\sigma=1,\dots,5$, integrated on the line $\re z=\frac12$ in ball arithmetic, and the closing of (iv), summed over its first $80$ terms at $k=1,2,3$, agree to $80$ digits and give
\begin{gather*}
P(1)=4.0962369983378828500555532283512\ldots,\\
P(2)=0.047716498691513110297412080890663\ldots;
\end{gather*}
the law at $k=1$, the low-end relation and the second low-end relation hold at these values to $75$ digits. For $k\ge1$, $c_4(k)>0$, so the law gives $P(k+4)$ as a rational combination of $P(k),\dots,P(k+3)$; the second low-end relation gives $P(3)$, and then the low-end relation $P(4)$, in $\mathbb Q+\mathbb Q\,P(1)+\mathbb Q\,P(2)$. So every value lies in $\mathbb Q+\mathbb Q\,P(1)+\mathbb Q\,P(2)$. An integer-relation search at $290$ bits finds no relation among $1$, $P(1)$, $P(2)$ with coefficients below $10^{25}$.

\emph{Answer.} $V_1=16$. The native gauges are $\one$ (native value $\frac1{25}$, $\det H=\frac{32}{625}$), the sign $(1,1,-1,-1)$ ($\frac19$, $\det H=-\frac{32}{81}$) and three curves of native value $1$. The closings are $\sum_u\binom ju^2\binom{2u}u^2$; the law is the display of (iii), of symbol $144(25\lambda-1)(9\lambda-1)(\lambda-1)^2$, proved by the telescoping identity and Theorem~\ref{thm:onekernel}. On the share side $\Gamma(\sigma)^2P(\sigma)=\frac1{2\pi i}\int\Gamma(z)^6\Gamma(\sigma-z)^2\Gamma(2z)^{-2}dz$; it closes to the right, $\Gamma(k)^2P(k)=-\sum_{x\ge k}(V_kP_2^2)'(x)$, and not to the left. The low-end relation is $17200-7160P(1)+274520P(2)-1017720P(3)+826200P(4)=0$, the second low-end relation $57P(1)-807P(2)+1050P(3)=196$, and the values lie in $\mathbb Q+\mathbb Q\,P(1)+\mathbb Q\,P(2)$.
\end{solution}

\begin{exercise}[$\mathrm H_4$ on a ray with a conjugate on the slit]\label{ex:lista-hslit}
On $\mathrm H_4$ take the masses $\mvec_1=(1,1,3,3)$ on $(a,b,2,3)$ and the letter exponents $s_1=3$, $t_1=1$, and let $P(k)=\langle1\rangle_k=P_{\mathrm H_4}\big(k\mvec_1;3k,k\big)$, $k\ge1$, be the values of the ray with the origin $0$ and the direction $(\mvec_1;3,1)$.
\begin{enumerate}[label=(\roman*)]
\item Show that the direction lies on the balance wall and strictly inside its walls, and compute $V_1$.
\item Find the native gauge $\alpha^\star$ of the direction on the positive orthant, $f^\star=f(\alpha^\star)$ and its conjugate over $\mathbb Q$, and show that the conjugate is the native value of a nondegenerate native gauge and lies on the slit $(0,f^\star)$.
\item Find the closings and a law of the values of order four, prove the law, and give its symbol.
\item Write the low-end relation, and show that Theorem~\ref{thm:positiveselection}(iii) does not apply.
\item Read $P(1)$ and $P(2)$ from the gauge end (Theorem~\ref{thm:gaugeend}), and check them against the low-end relation.
\item Write $\Gamma(3k)\Gamma(k)P(k)$ on the share side, read the second low-end relation from it, and give the value span.
\end{enumerate}
\end{exercise}

\begin{solution}
\emph{(i) Walls and $V_1$.} The direction has $m_{1,\cE}=8=3r+1\cdot r^\ast=3\cdot2+2$, on the balance wall. Its walls are $\kappa_{1,S}=m_{1,S}-3d(S)-d^\ast(S)$ (Definition~\ref{def:family}), with $d(S)=|S|-r^\ast(S)$ and $d^\ast(S)=|S|-r(S)$ (Definition~\ref{def:skeleton}). A single mark has $d=d^\ast=0$; $\{a,b\}$ has $d=0$, $d^\ast=1$; $\{2,3\}$ has $d=1$, $d^\ast=0$; the four pairs $\{x,y\}$, $x\in\{a,b\}$, $y\in\{2,3\}$, have $d=d^\ast=0$; every triple has $d=d^\ast=1$. So
\begin{gather*}
\kappa_{1,S}=1\quad\text{for }S\in\big\{\{a\},\{b\},\{a,b\},\{a,b,2\},\{a,b,3\}\big\},\\
\kappa_{1,S}=3\quad\text{for }S\in\big\{\{2\},\{3\},\{2,3\},\{a,2,3\},\{b,2,3\}\big\},
\end{gather*}
and $\kappa_{1,S}=4$ for the four pairs $\{x,y\}$: all positive. An order of the marks is the chain of its tails $U_2\supset U_3\supset U_4$, and $V_1=\sum\big(\kappa_{1,U_2}\kappa_{1,U_3}\kappa_{1,U_4}\big)^{-1}$ (Proposition~\ref{prop:principal}). Below $U_2=\{a,b,y\}$ the pairs and their marks give $1\cdot(1+1)+2\cdot\frac14(1+\frac13)=\frac83$; below $U_2=\{x,2,3\}$ they give $2\cdot\frac14(1+\frac13)+\frac13(\frac13+\frac13)=\frac89$. So
\[
V_1=2\cdot\frac83+2\cdot\frac13\cdot\frac89=\frac{160}{27}.
\]

\emph{(ii) The native gauge and its conjugate.} The exchanges $a\leftrightarrow b$ and $2\leftrightarrow3$ keep the direction, and $\alpha^\star$ is unique up to scale (Theorem~\ref{thm:pencil}(i)), so $\alpha^\star=(x,x,1,1)$ for some $x>0$. On that slice $K=4x+1$, $\Psi=x(x+4)$ and
\[
f=\frac{\alpha_a\alpha_b\alpha_2^3\alpha_3^3}{K^3\Psi}=\frac x{(x+4)(4x+1)^3},\qquad \frac{d}{dx}\log f=\frac{-4(3x^2+8x-1)}{x(x+4)(4x+1)},
\]
so $x^\star=\frac{\sqrt{19}-4}3$, where the gradient of $\log f$ vanishes in all four coordinates, and
\[
f^\star=\frac{7643+1064\sqrt{19}}{1366875}=0.0089846317\ldots,\qquad \det H=\frac{376096\sqrt{19}-1614848}{151875}=0.16142523\ldots,
\]
$H$ the Hessian of $-\log f$ on the chart $y_3=0$; $f^\star$ is a root of $1366875\lambda^2-15286\lambda+27$. By Proposition~\ref{prop:rays-conjugates} the conjugate point $\alpha'=(x',x',1,1)$, $x'=\frac{-\sqrt{19}-4}3=-2.786\ldots$, where $K=\frac{-13-4\sqrt{19}}3$ and $\Psi=x'(x'+4)$ do not vanish, is a nondegenerate native gauge, with
\begin{gather*}
f(\alpha')=\frac{7643-1064\sqrt{19}}{1366875}=0.0021985415\ldots,\\
\det H(\alpha')=\frac{-376096\sqrt{19}-1614848}{151875}=-21.4269\ldots ,
\end{gather*}
and $0<f(\alpha')<f^\star$: the conjugate lies on the slit, at a gauge where the class $\{a,b\}$ is flipped.

\emph{(iii) Closings and the law.} Every letter exponent of the direction is positive and the origin is $0$, so $\mathsf C_j(1)=\CT(f^{-j})=[\alpha^{j\mvec_1}]K^{3j}\Psi^j$ (Definition~\ref{def:closings}). In
\[
K^{3j}\Psi^j=\sum_{u,v}\binom{3j}u\binom jv(AB)^{4j-u-v}(\alpha_2\alpha_3)^u(\alpha_a\alpha_b)^v
\]
the monomial $\alpha_a^j\alpha_b^j\alpha_2^{3j}\alpha_3^{3j}$ takes $(\alpha_a\alpha_b)^{j-v}$ from $A^{4j-u-v}$ and $(\alpha_2\alpha_3)^{3j-u}$ from $B^{4j-u-v}$, so $u=2j+v$, with the coefficient $\binom{2j-2v}{j-v}^2$; with $w=j-v$,
\begin{gather*}
\mathsf C_j(1)=\sum_{w=0}^jF(j,w),\qquad F(j,w)=\binom{3j}w\binom jw\binom{2w}w^2;\\
\mathsf C_j(1)=1,\ 13,\ 589,\ 37597,\ 2791949,\ \dots
\end{gather*}
On the first $300$ closings there is no recurrence of order one, two, three with degree at most $138$, $91$, $68$, and none of order four with degree at most $24$; at order four and degree $25$ exact linear algebra on the first $170$ closings gives one line, spanned by $p_0,\dots,p_4\in\mathbb Z[j]$ with coefficients of up to $36$ digits, and the recurrence holds on all $300$. It holds for every $j$. For $0\le w\le j+4$ put
\begin{gather*}
T(j,w)=\frac{((2w)!)^2}{(w!)^6\,(3j+12-w)!\,(j+4-w)!},\\
G(j,w)=(3j)!\,j!\,T(j,w)\,A(j,w),\qquad G(j,j+5)=0,
\end{gather*}
with $A\in\mathbb Q[j,w]$, a polynomial of $416$ terms, of degree $16$ in $w$ and divisible by $w^4$, printed by the receipt. As $F(j+i,w)=(3j)!\,j!\,T(j,w)\prod_{l=1}^{3i}(3j+l)\prod_{l=1}^{i}(j+l)\prod_{l=3i+1}^{12}(3j+l-w)\prod_{l=i+1}^4(j+l-w)$ and $T(j,w+1)=T(j,w)\,4(2w+1)^2(3j+12-w)(j+4-w)/(w+1)^4$, the identity $\sum_ip_i(j)F(j+i,w)=G(j,w+1)-G(j,w)$ for $0\le w\le j+4$ is the identity of polynomials
\[
(w+1)^4N(j,w)=4(2w+1)^2(3j+12-w)(j+4-w)A(j,w+1)-(w+1)^4A(j,w),
\]
$N=\sum_ip_i(j)\prod_{l=1}^{3i}(3j+l)\prod_{l=1}^{i}(j+l)\prod_{l=3i+1}^{12}(3j+l-w)\prod_{l=i+1}^4(j+l-w)$, which holds in $\mathbb Q[j,w]$; summed over $w$ it telescopes to $0$, as $A(j,0)=0$. As in Exercise~\ref{ex:lista-hcentre}(iii), Theorems~\ref{thm:closings} and~\ref{thm:onekernel} turn it into the law
\[
\sum_{i=0}^4c_i(k)\,P(k+i)=0,\qquad c_i(k)=-p_{4-i}(-k-4),
\]
an identity of meromorphic functions of $k$. The coefficients factor as
\begin{gather*}
c_0=3k^3(k+1)(3k+1)^3(3k+2)^3a_0,\qquad c_1=-2(k+1)a_1,\qquad c_2=2(3k+4)(3k+5)a_2,\\
c_3=-2(3k+4)(3k+5)(3k+7)(3k+8)(k+2)^2a_3,\\
c_4=151875(3k+4)(3k+5)(3k+7)(3k+8)(3k+10)(3k+11)(k+2)^2(k+3)^2a_4,
\end{gather*}
with $a_i\in\mathbb Z[k]$ of degrees $15,24,23,19,15$, printed by the receipt, and the symbol is
\[
\chi(\lambda)=168132676253771376\,(\lambda-1)^2\,(1366875\lambda^2-15286\lambda+27):
\]
$f^\star$, its conjugate, and $1$, doubled. The root $1$ is the native value of the line $\{A=B=0\}$, a source of dimension one (Definition~\ref{def:source}): there $K=\alpha_2\alpha_3$ and $\Psi=\alpha_a\alpha_b$, the occupations of $K$ and of $\Psi$ are $(0,0,1,1)$ and $(1,1,0,0)$ on $(a,b,2,3)$, so the usage $3R+R^\ast$ is $\mvec_1$, and $f=1$.

\emph{(iv) The low end and positive selection.} $c_0$ vanishes to order three at $0$ with $[k^3]c_0=656679698393897858976$, and with $V_1=\frac{160}{27}$ the argument of Corollary~\ref{cor:rays-lowend}(ii) gives the low-end relation, divided by $15360$,
\begin{multline*}
253348649071719853+204923553451974000\,P(1)-272918959559705480400\,P(2)\\
+1177478827066019041200\,P(3)-1011712684533185250000\,P(4)=0 .
\end{multline*}
Theorem~\ref{thm:positiveselection} asks the roots of the symbol to be distinct, which the double root $1$ already violates, and (iii) asks every root other than $f^\star$ to have modulus greater than $f^\star$. The conjugate $f(\alpha')$ lies in $(0,f^\star)$, so a solution of the law that grows like $f(\alpha')^k$ grows more slowly than the values, and growth does not separate it from them: the hypothesis fails. Set up anyway, the equations of Theorem~\ref{thm:positiveselection}(iv), with $P(1),P(2),P(3)$ free, $P(4)$ from the low-end relation and the solution required to vanish at $K,K+1,K+2$, return $P(1)=-2.6\cdot10^{23}$ at $K=40$ and $-7.4\cdot10^{47}$ at $K=80$.

\emph{(v) The gauge end.} $f^\star$ is a simple root of $\chi$, and the density operator $\mathfrak D_{\mathbf p}$ of the law (Theorem~\ref{thm:gaugeend}) has order $D=25$; at $f^\star$ its coefficient $b_{D,0}$ vanishes and the extra exponent is $\gamma=D-1-b_{D-1,0}/b_{D,1}=\frac12=\frac{n-3}2$, as Theorem~\ref{thm:gaugeend}(iii) requires. As $\gamma\notin\{0,\dots,23\}$, the germ of the density is $\varrho_f(t)=s^{1/2}g(0)\Phi(s)$, $s=f^\star-t$, with $\Phi$ fixed by the recursion of Theorem~\ref{thm:gaugeend}(iv), $1800$ terms computed over $\mathbb Q(\sqrt{19})$ in ball arithmetic, and the scale read off the Hessian at $\alpha^\star$ (the end of the stride's law at $f^\star$, Definition~\ref{def:strideslaw}, with $(\omega_0/dy)(\alpha^\star)=1$):
\[
g(0)=\frac1{\Gamma(\frac32)\sqrt{\det H}}\Big(\frac{2\pi}{f^\star}\Big)^{3/2}=51938.4198592508155915\ldots
\]
By Theorem~\ref{thm:gaugeend}(v) the density continues holomorphically along $(0,f^\star)$, through the root $f(\alpha')$ of the symbol, so the series of $\Phi$ converges on $[0,f^\star)$: the coefficients $g_m$ of $g=g(0)\Phi$ have $|g_m|(f^\star)^m\le5.2\cdot10^5$ for $m\le1800$. The Beta integral gives
\[
P(k)=\int_0^{f^\star}t^k\varrho_f(t)\,dt=\sum_{m\ge0}g_m\,(f^\star)^{k+\frac32+m}\,\mathrm B\big(k+1,m+\tfrac32\big),
\]
read at $k=K,\dots,K+3$, and the law, with $c_0(k)\neq0$ at every integer $1\le k\le200$, runs back to $k=1$. From $K=80$ and from $K=120$ the readings agree to $92$ digits:
\begin{gather*}
P(1)=0.48830319489526888388541464396010900874354031922042\ldots,\\
P(2)=0.00132119179498773222717641714721510065075969610885\ldots
\end{gather*}
The low-end relation, which the reading does not use, returns $V_1=\frac{160}{27}$ from these values to $191$ digits.

\emph{(vi) The share side and the second low-end relation.} The widths of (i) of Exercise~\ref{ex:lista-hcentre} do not depend on the direction, so the hinge tree is the same: the series cell $\mathcal S$ of $\{2,3\}$, whose two marks have the masses $3k$, and the parallel cell $\mathcal P$ of $\{a,b\}$, with the masses $k$, at $s=3k$, $t=k$. Every mark is again a one-mark cell whose transform, or dual transform, is $\mathsf g$ (Definition~\ref{def:treeproc}(1): $s=m_2=m_3$, $t=m_a=m_b$), the hinge walls are $\kappa_{\mathcal S}=6k-3k=3k$ and $\kappa_{\mathcal P}=2k-k=k$, the walls of $\{2,3\}$ and $\{a,b\}$ (Theorem~\ref{thm:seriesparallel}), and Proposition~\ref{prop:faces}(iii) gives $\int\varrho_{\mathcal S}^{-z}d\phi_{\mathcal S}=P_2(z)=\int\varrho_{\mathcal P}^{z-k}d\phi_{\mathcal P}$. Theorem~\ref{thm:glue}(iv), on a line $\re z=c$, $0<c<k$, gives
\[
\Gamma(3k)\Gamma(k)\,P(k)=\frac1{2\pi i}\int_{(c)}\frac{\Gamma(z)^6\,\Gamma(k-z)\,\Gamma(3k-z)}{\Gamma(2z)^2}\,dz,
\]
which at $k=1$, on $c=\frac12$ in ball arithmetic, agrees with (v) to $60$ digits. At $k=-1$ the coefficients $c_0$ and $c_1$ vanish, simply, and the second low-end relation needs the principal parts of $P$ at $0$ and $-1$. As in Exercise~\ref{ex:lista-hcentre}(v), the integral on $\re z=\frac12$ is holomorphic near $k=0$ and $k=-1$, $\Gamma(3k)\Gamma(k)$ has a double pole at both, and the principal parts of $P$ come from the poles that cross the line, $z=k$ and $z=3k$ near $0$, and $z=k,k+1$ and $z=3k,\dots,3k+3$ near $-1$, each simple, with the residues $(-1)^{m+1}\Gamma(2k-m)\,\Gamma(k+m)^6/\big(m!\,\Gamma(2k+2m)^2\big)$ at $z=k+m$ and $(-1)^{m+1}\Gamma(-2k-m)\,\Gamma(3k+m)^6/\big(m!\,\Gamma(6k+2m)^2\big)$ at $z=3k+m$ (Proposition~\ref{prop:stacked}(iv)). Expanded by Proposition~\ref{prop:resonance}, with $T=\Tx{2}{\IZ}{0}$,
\[
P(x)=\frac{160}{27x^3}-\frac{208\,T}{9x}+O(1),\qquad P(-1+x)=\frac{2080}{27x^3}-\frac{21118}{81x^2}-\frac{2704\,T+\frac{887}9}{9x}+O(1),
\]
with the top coefficients $V_1\mathsf C_0(1)$ and $V_1\mathsf C_1(1)=\frac{160}{27}\cdot13$ (Proposition~\ref{prop:principal}). With the expansions of $c_0$ and $c_1$ at $-1$ to the order three, the multiples of $T$ cancel in the constant term at $k=-1$, which is $-2077616120857361856$, and $(c_2,c_3,c_4)(-1)=625190400\cdot(7005765879,-74409017489,88989468750)$; so
\[
7005765879\,P(1)-74409017489\,P(2)+88989468750\,P(3)=\frac{5981712159277}{1800},
\]
which the values of (v) satisfy to $109$ digits. For $k\ge1$, $c_4(k)>0$ (its real roots are below $\frac19$), so, as in Exercise~\ref{ex:lista-hcentre}(vi), every value lies in $\mathbb Q+\mathbb Q\,P(1)+\mathbb Q\,P(2)$.

\emph{Answer.} $V_1=\frac{160}{27}$; $f^\star=\frac{7643+1064\sqrt{19}}{1366875}$ at $\alpha^\star=(x^\star,x^\star,1,1)$, $x^\star=\frac{\sqrt{19}-4}3$, $\det H=\frac{376096\sqrt{19}-1614848}{151875}$; its conjugate $\frac{7643-1064\sqrt{19}}{1366875}$ is the native value of the nondegenerate native gauge $(x',x',1,1)$, $x'=\frac{-\sqrt{19}-4}3$, and lies on the slit. The closings are $\sum_w\binom{3j}w\binom jw\binom{2w}w^2$, the law has order four, degree $25$ and the symbol $(\lambda-1)^2(1366875\lambda^2-15286\lambda+27)$ up to a constant, positive selection does not apply, and the gauge end gives $P(1)=0.4883031948952688838854\ldots$ The second low-end relation is $7005765879P(1)-74409017489P(2)+88989468750P(3)=\frac{5981712159277}{1800}$, and the value span is $\mathbb Q+\mathbb Q\,P(1)+\mathbb Q\,P(2)$.
\end{solution}

\section{The five-mark lists}\label{sec:lista-five}

In this section the marks are $1,\dots,5$, every weight is $1$, and $\ell=P_{\binom{[3]}1}(\one;1,1)$ is the value at the strength one of the centre ray of the list $\binom{[3]}1$, $P_3(\sigma)=P_{\binom{[3]}1}(\sigma\one;\sigma,\sigma)$ $P_2(\sigma)=P_{\binom{[2]}1}(\sigma\one;\sigma,\sigma)=\Gamma(\sigma)^2/\Gamma(2\sigma)$, and $\mathsf g(y)=\sG_{0;(1,-1)}(1,y)$ as in Exercise~\ref{ex:lista-hcentre}(iv). The lists are $\binom{[5]}1$ and $\binom{[5]}2$, and three lists of rank two named by their parallel classes:
\begin{itemize}
\item $\mathrm P_5$, whose bases are the nine pairs other than $\{1,2\}$: the classes $\{1,2\},\{3\},\{4\},\{5\}$;
\item $\mathrm R_5$, whose bases are the eight pairs $\{x,y\}$ with $x,y$ in different classes of $\{1,2\},\{3,4\},\{5\}$;
\item $\mathrm T_5$, whose bases are the seven pairs $\{x,y\}$ with $x,y$ in different classes of $\{1,2,3\},\{4\},\{5\}$.
\end{itemize}

\begin{exercise}[The five-mark lists along their hinge trees]\label{ex:lista-trees}
\begin{enumerate}[label=(\roman*)]
\item Find, for each of the five lists, the sets of width one with at least two marks on each side, and its canonical hinge tree.
\item Show that $P_{\mathrm P_5}(\sigma\one;\sigma,\sigma)=P_3(\sigma)^2$. Prove the law of $P_3$, and derive from it the law of the centre ray of $\mathrm P_5$, its symbol, its low-end relation and its value span.
\item Show that $q(\sigma)=P_{\binom{[4]}2}(\sigma\one;\sigma,\sigma)=P_2(\sigma)P_3(\sigma)$, and find the law of $q$.
\item Write the centres of $\mathrm T_5$ and $\mathrm R_5$ on their share sides, close $\mathrm T_5$ at the integer strengths, and write $P_{\mathrm T_5}(\one;1,1)$ in the coordinates $1$, $\ell$ and $\ell'=\partial_\sigma P_3(\sigma)|_{\sigma=1}$.
\end{enumerate}
\end{exercise}

\begin{solution}
\emph{(i) Widths and trees.} The width of $S$ is $c(S)=r(S)+r(\cE\setminus S)-r$ (Definition~\ref{def:skeleton}). In $\binom{[5]}1$ every nonempty set has rank one, so every proper $S$ has width $1$, and the list is one parallel node. In $\binom{[5]}2$ every set of two or three marks has rank two and so has its complement: width $2$, and the list is prime (Theorem~\ref{thm:primetree}); so is $\binom{[4]}2$, whose pairs have width $2$. In a list of rank two whose classes are the parallel classes, a set of two or three marks has width one exactly when it, or its complement, lies in one class, and width two otherwise. So the sets of width one with at least two marks on each side are $\{1,2\}$ for $\mathrm P_5$; $\{1,2\}$ and $\{3,4\}$ for $\mathrm R_5$; $\{1,2,3\}$ and its pairs for $\mathrm T_5$; each with its complement. The canonical hinge trees (Definition~\ref{def:hingetree}, Theorem~\ref{thm:primetree}) are:
\begin{itemize}
\item $\binom{[5]}1$: one parallel node; $\binom{[5]}2$: one prime node;
\item $\mathrm P_5$: the parallel node $\binom{\{g,1,2\}}1$ and the prime node $\binom{\{g,3,4,5\}}2$;
\item $\mathrm R_5$: the series node $\binom{\{g_1,g_2,5\}}2$ with the parallel nodes $\binom{\{g_1,1,2\}}1$ and $\binom{\{g_2,3,4\}}1$;
\item $\mathrm T_5$: the series node $\binom{\{g,4,5\}}2$ with the parallel node $\binom{\{g,1,2,3\}}1$ of four marks.
\end{itemize}
For $\mathrm P_5$, for instance, the bases of the two-sum are $\{x,y\}\subseteq\{3,4,5\}$, from $\{g\}$ and a pair of $\{3,4,5\}$, and $\{1,y\}$, $\{2,y\}$, $y\in\{3,4,5\}$, from $\{1\}$ or $\{2\}$ and $\{g,y\}$: the nine pairs other than $\{1,2\}$.

\emph{(ii) The doubled mark of $\binom{[4]}2$.} The cell $\mathcal A=\binom{\{g,3,4,5\}}2$ of the bunch $\{3,4,5\}$ has the hinge minors $\mathcal A\setminus g=\binom{\{3,4,5\}}2$ and $\mathcal A/g=\binom{\{3,4,5\}}1$, dual with the identity and the same weights, so by Corollary~\ref{cor:piecesmarks}(ii) its ratio is identically $1$ and $\phi_{\mathcal A}(\Delta)=P_{\binom{[3]}2}(\sigma\one;\kappa_{\mathcal A},2\sigma-\kappa_{\mathcal A})$, with the hinge wall $\kappa_{\mathcal A}=\sigma\,c(\{3,4,5\})=\sigma(2+1-2)=\sigma$. By Corollary~\ref{cor:piecesmarks}(i) the cell acts in the gluing as $\phi_{\mathcal A}(\Delta)$ times one mark $e$ of ratio $1$ and mass $\sigma$, and the parallel cell of $\{1,2\}$ glued to that mark is the list $\binom{\{1,2,e\}}1$. By Proposition~\ref{prop:dualfamily} the dual table $\binom{[3]}1$ at the swapped exponents $(\sigma,\sigma)$ has the same period, so
\[
P_{\mathrm P_5}(\sigma\one;\sigma,\sigma)=P_{\binom{[3]}2}(\sigma\one;\sigma,\sigma)\,P_{\binom{\{1,2,e\}}1}(\sigma\one;\sigma,\sigma)=P_3(\sigma)^2 .
\]

The law of $P_3$. The centre ray of $\binom{[3]}1$ has $f^{-1}=(\alpha_1+\alpha_2+\alpha_3)(\alpha_1^{-1}+\alpha_2^{-1}+\alpha_3^{-1})$, and its closings are (Definition~\ref{def:closings})
\[
\mathsf C_j(1)=\sum_{a+b+c=j}\Big(\frac{j!}{a!\,b!\,c!}\Big)^2=\sum_{u=0}^j\binom ju^2\binom{2u}u=1,\ 3,\ 15,\ 93,\ 639,\ 4653,\ \dots,
\]
grouping $a+b=u$. No recurrence of order one and degree at most $40$ holds on the first $120$ closings, and exact linear algebra gives
\[
(j+2)^2\,\mathsf C_{j+2}(1)-(10j^2+30j+23)\,\mathsf C_{j+1}(1)+9(j+1)^2\,\mathsf C_j(1)=0 .
\]
It holds for every $j$: with $F(j,u)=\binom ju^2\binom{2u}u$, $T(j,u)=((j+2)!)^2(2u)!/\big((u!)^4((j+2-u)!)^2\big)$ for $0\le u\le j+2$, $A(j,u)=u^3(j+1)^2(3u-4j-8)$ and $G(j,u)=T(j,u)A(j,u)/((j+1)^2(j+2)^2)$, $G(j,j+3)=0$, one has $F(j+i,u)=T(j,u)\prod_{l=i+1}^2\big((j+l-u)/(j+l)\big)^2$ and $T(j,u+1)=T(j,u)\,2(2u+1)(j+2-u)^2/(u+1)^3$, so the identity $\sum_ip_i(j)F(j+i,u)=G(j,u+1)-G(j,u)$, $p_0=9(j+1)^2$, $p_1=-(10j^2+30j+23)$, $p_2=(j+2)^2$, is the identity of polynomials
\[
(u+1)^3\sum_{i=0}^2p_i(j)\prod_{l\le i}(j+l)^2\prod_{l>i}(j+l-u)^2=2(2u+1)(j+2-u)^2A(j,u+1)-(u+1)^3A(j,u),
\]
which holds in $\mathbb Q[j,u]$; summed over $u$ it telescopes to $0$. As in Exercise~\ref{ex:lista-hcentre}(iii), Theorems~\ref{thm:closings} and~\ref{thm:onekernel} give the law of the values, an identity of meromorphic functions,
\[
k^2P_3(k)-(10k^2+10k+3)\,P_3(k+1)+9(k+1)^2\,P_3(k+2)=0,
\]
of symbol $9\lambda^2-10\lambda+1=(9\lambda-1)(\lambda-1)$. The walls of the direction are all $1$, so $V_1=3!=6$ (Proposition~\ref{prop:principal}), and the low-end relation (Corollary~\ref{cor:rays-lowend}(ii)) is $6-3P_3(1)+9P_3(2)=0$. So $P_3(1)=\ell$, $P_3(2)=\frac\ell3-\frac23$, and the law gives
\[
P_3(3)=\frac{5\ell}{27}-\frac{23}{54},\qquad P_3(4)=\frac{31\ell}{243}-\frac{145}{486},\qquad P_3(k)\in\mathbb Q+\mathbb Q\,\ell\quad(k\ge1).
\]

The law of the squares. Write the law as $P_3(k+2)=a(k)P_3(k)+b(k)P_3(k+1)$; then $P_3(k+3)=a^{(1)}P_3(k)+b^{(1)}P_3(k+1)$ with $a^{(1)}=b(k+1)a(k)$ and $b^{(1)}=a(k+1)+b(k+1)b(k)$. The squares $x(k)=P_3(k)^2,\dots,x(k+3)$ are quadratic forms in $P_3(k),P_3(k+1)$, and the combination $d_0x(k)+d_1x(k+1)+d_2x(k+2)+x(k+3)$ with $d_2=-a^{(1)}b^{(1)}/(ab)$, $d_0=-d_2a^2-(a^{(1)})^2$, $d_1=-d_2b^2-(b^{(1)})^2$ vanishes identically. Cleared of denominators, the law of the centre ray of $\mathrm P_5$ is $\sum_{i=0}^3d_i(k)\,x(k+i)=0$ with
\begin{gather*}
d_0=k^4(10k^2+30k+23),\qquad d_3=-729(k+2)^4(10k^2+10k+3),\\
d_1=-(10k^2+10k+3)\,w(k),\qquad d_2=9(10k^2+30k+23)\,w(k),\\
w(k)=91k^4+364k^3+506k^2+284k+60,
\end{gather*}
which the squares of the values above satisfy identically in $\ell$ at $k=1,\dots,5$. Its symbol is $-10(\lambda-1)(9\lambda-1)(81\lambda-1)$, the products of the roots $\frac19$ and $1$ in pairs; $\frac1{81}=K(\one)^{-2}$, $K(\one)=9$ the number of bases. The walls of $\mathrm P_5$ are $1$ on single marks, on $\{1,2\}$ and on their complements, and $2$ on the other sets of two or three marks, and $V_1=36$, the square of the $6$ of $\binom{[3]}1$, as $P_3(k)^2\sim36k^{-4}$ at $0$ requires. $d_0$ vanishes to order four at $0$, $[k^4]d_0=23$, and the low-end relation
\[
23\cdot36+d_1(0)\,x(1)+d_2(0)\,x(2)+d_3(0)\,x(3)=828-180\,x(1)+12420\,x(2)-34992\,x(3)=0
\]
holds identically in $\ell$ at $x(i)=P_3(i)^2$. Every value $x(k)=P_3(k)^2$ lies in $\mathbb Q+\mathbb Q\,\ell+\mathbb Q\,\ell^2$. As a check, the procedure that reads a ray from its list by a box (Definition~\ref{def:readprocedure}, Theorem~\ref{thm:raybox}), run on $\mathrm P_5$ by the supplied reader, returns the same law and $P_{\mathrm P_5}(\one;1,1)=5.4939011435808440124442576078206911958\ldots$, which is $\ell^2$ to $87$ digits.

\emph{(iii) $\binom{[4]}2$ at its centre.} $\binom{[4]}2$ is identically self-dual: the complement of a pair is a pair, and the weights are $1$. Corollary~\ref{cor:mergedvalues} at a mark, and Proposition~\ref{prop:dualfamily} for the deletion $\binom{[3]}2$, give $q(\sigma)=\Gst21{\sigma,\ \sigma}{2\sigma}P_3(\sigma)=P_2(\sigma)P_3(\sigma)$. With $P_2(\sigma+1)/P_2(\sigma)=\sigma/(2(2\sigma+1))$ (Proposition~\ref{prop:stacked}(iii)), the law of $P_3$ multiplied through by $\sigma P_2(\sigma)$ is the law of $q$:
\[
\sigma^3q(\sigma)-2(2\sigma+1)(10\sigma^2+10\sigma+3)\,q(\sigma+1)+36(\sigma+1)(2\sigma+1)(2\sigma+3)\,q(\sigma+2)=0,
\]
with the symbol $(4\lambda-1)(36\lambda-1)$ and $V_1=12=2\cdot6$, the walls of $\binom{[4]}2$ being $1,2,1$ on sets of one, two and three marks.

\emph{(iv) The triangle rows.} $\mathrm T_5$ is $\mathrm H_4$ of Section~\ref{sec:lista-h} with the class enlarged to $\{1,2,3\}$: the two-sum of the series cell $\mathcal S$ of $\{4,5\}$, with $f_{\mathcal S}=\mathsf g^2$ and $\kappa_{\mathcal S}=\sigma$, and the parallel cell $\mathcal P$ of $\{1,2,3\}$, with $\check f_{\mathcal P}=\mathsf g^3$ and $\kappa_{\mathcal P}=3\sigma-2\sigma=\sigma$ (Theorem~\ref{thm:seriesparallel}). The root of three one-mark leaves (Definition~\ref{def:treeproc}(5)) gives $\int_0^\infty y^{z-1}\mathsf g^3dy=\Gamma(z)^2P_3(z)$, so Proposition~\ref{prop:faces}(iii) gives $\int\varrho_{\mathcal P}^{z-\sigma}d\phi_{\mathcal P}=P_3(z)$, and $\int\varrho_{\mathcal S}^{-z}d\phi_{\mathcal S}=P_2(z)$ as in Exercise~\ref{ex:lista-hcentre}(iv). Theorem~\ref{thm:glue}(iv) gives, for $0<c<\sigma$,
\[
\Gamma(\sigma)^2P_{\mathrm T_5}(\sigma\one;\sigma,\sigma)=\frac1{2\pi i}\int_{(c)}\Gamma(z)^2\,\Gamma(\sigma-z)^2\,q(z)\,dz .
\]
$|q(z)|\le q(\re z)$, the modulus of the integrand of a period being the integrand at the real parts, and $q(x)\le4^{1-x}P_3(1)$ for $x\ge1$; so the closing to the right of Exercise~\ref{ex:lista-hcentre}(iv) applies, the moved line contributing $O(N^{2s-2}4^{-N})$, and at an integer strength $s$
\[
\Gamma(s)^2P_{\mathrm T_5}(s\one;s,s)=-\sum_{x\ge s}\big(V_s\,q\big)'(x),\qquad V_s(x)=\prod_{j=1}^{s-1}(x-j)^2 .
\]
At $s=1$ this needs the slopes $q'(x)=P_2(x)\big(P_3'(x)+2(h_{x-1}-h_{2x-1})P_3(x)\big)$. The law of $P_3$ holds as an identity of meromorphic functions, so its derivative in $k$ holds as well and gives $P_3'(k+2)$ from $P_3'(k)$, $P_3'(k+1)$ and the values, starting from $P_3'(1)=\ell'$ and $P_3'(2)$. Near $k=0$, $k^2P_3(k)=6+\mathsf L_1k+O(k^2)$ with $\mathsf L_1=0$ (Proposition~\ref{prop:rays-faces}, every weight being $1$), so the derivative of the law at $k=0$ reads $-10\ell-3\ell'+18P_3(2)+9P_3'(2)=0$:
\[
P_3'(2)=\frac{4\ell}9+\frac{\ell'}3+\frac43,\qquad P_3'(3)=\frac{26\ell}{81}+\frac{5\ell'}{27}+\frac{13}{18},\qquad P_3'(k)\in\mathbb Q+\mathbb Q\,\ell+\mathbb Q\,\ell'.
\]
So $P_{\mathrm T_5}(\one;1,1)=-\sum_{x\ge1}q'(x)=a+b\,\ell+c\,\ell'$, where $a$, $b$, $c$ are series of rational numbers whose terms decrease like $4^{-x}$; summed to $x=65$,
\begin{gather*}
a=-0.46159963274916253008918\ldots,\qquad b=2.01706417220768446716615\ldots,\\
c=-1.06282955174549612502149\ldots
\end{gather*}
With $\Gamma(\sigma)^2P_3(\sigma)=\int_0^\infty y^{\sigma-1}\mathsf g(y)^3\,dy$, so that $\ell'=\int_0^\infty\log y\,\mathsf g^3dy-2(\log\Gamma)'(1)\,\ell$, quadrature gives
\begin{gather*}
\ell=2.3439072386894588906015622888722770690954\ldots,\\
\ell'=-8.0751522635316440102696512414673765013\ldots,
\end{gather*}
and
\[
P_{\mathrm T_5}(\one;1,1)=12.848722141915555957312339931253534636619\ldots,
\]
which the reading of $\mathrm T_5$ by the box procedure (Definition~\ref{def:readprocedure}) reproduces to $53$ digits. For $\mathrm R_5$ the series root has three cells, the mark $5$, with the transform $\mathsf g$, and the parallel cells of $\{1,2\}$ and $\{3,4\}$, each with the transform $f_{\mathcal P_2}$ whose Mellin transform is $\Gamma(w)^2P_2(\sigma-w)$, $0<w<\sigma$ (Proposition~\ref{prop:faces}(iii), with $\check f_{\mathcal P_2}=\mathsf g^2$); by Definition~\ref{def:treeproc}(5)
\[
\Gamma(\sigma)^2P_{\mathrm R_5}(\sigma\one;\sigma,\sigma)=\int_0^\infty y^{\sigma-1}\,\mathsf g(y)\,f_{\mathcal P_2}(y)^2\,dy,
\]
and the box procedure reads $P_{\mathrm R_5}(\one;1,1)=7.8658373683640215924219746974461682109\ldots$

\emph{Answer.} The trees are those listed in (i). $P_{\mathrm P_5}=P_3^2$, with the law $\sum_{i=0}^3d_i(k)x(k+i)=0$ of symbol $-10(\lambda-1)(9\lambda-1)(81\lambda-1)$, $V_1=36$, the low-end relation $828-180x(1)+12420x(2)-34992x(3)=0$, and the value span $\mathbb Q+\mathbb Q\,\ell+\mathbb Q\,\ell^2$. $P_{\binom{[4]}2}=P_2P_3$ obeys $\sigma^3q(\sigma)-2(2\sigma+1)(10\sigma^2+10\sigma+3)q(\sigma+1)+36(\sigma+1)(2\sigma+1)(2\sigma+3)q(\sigma+2)=0$. $\Gamma(\sigma)^2P_{\mathrm T_5}(\sigma)=\frac1{2\pi i}\int\Gamma(z)^2\Gamma(\sigma-z)^2q(z)\,dz$ closes to the right, $\Gamma(s)^2P_{\mathrm T_5}(s)=-\sum_{x\ge s}(V_sq)'(x)$, and $P_{\mathrm T_5}(\one;1,1)=a+b\,\ell+c\,\ell'=12.8487221419\ldots$ with the series $a,b,c$ of (iv).
\end{solution}

\section{Larger lists and the cost of a reading}\label{sec:exlists-larger}

% Exercises of Part XI, chapter ch:exlists, second half (items 33 to 36 of EXERCISES_PLAN.md):
% N_6 at (1;0,2) and at its centre, V_8 at (1;0,2), sizing a reading (Q_6 and N_7), a table of rays.
% Lists named here: N_6, V_8, Q_6 (alias: the whirl), N_7, H_4, the lists binom([n],1), binom([n],2) and the wheel W_3.
% No new macros.

\begin{exercise}[The list $\mathrm N_6$: letters and constant terms]\label{ex:listsb-nletters}
Let $\mathrm N_6$ be the list of rank three on $\cE=\{1,\dots,6\}$ whose bases are the nineteen triples other than $a=\{1,2,3\}$, and put $b=\{4,5,6\}$. On $a$ write $K_a=\alpha_1+\alpha_2+\alpha_3$, $\Psi_a=\alpha_1\alpha_2+\alpha_1\alpha_3+\alpha_2\alpha_3$ and $\alpha^a=\alpha_1\alpha_2\alpha_3$, the letters and the monomial of the list $\binom a1$, and $K_b$, $\Psi_b$, $\alpha^b$ the same on $b$.
\begin{enumerate}[label=(\roman*)]
\item Show that the letters of $\mathrm N_6$ are $K=Q+\alpha^b$ and $\Psi=Q+\alpha^a$, with $Q=\Psi_aK_b+K_a\Psi_b$.
\item For $i,l\in\mathbb N$ with $2i+l=3D$, $D\in\mathbb N$, let $\gamma_{i,l}$ be the coefficient of $\alpha_1^D\alpha_2^D\alpha_3^D$ in $\Psi_a^i\,K_a^l$. Show that
\[
\gamma_{i,l}=\sum_{w}\binom iw\binom l{w+\frac{l-i}3}\binom{2D-2w}{D-w}.
\]
\item The centre ray of $\mathrm N_6$ at the strength $1$ (Definition~\ref{def:centreray}) has the direction $(\one;1,1)$, and the ray with origin $0$ and direction $(\one;0,2)$, the full cobasis lean, has the stride $f=\alpha^{\cE}\Psi^{-2}$. Write the constant terms $\CT(f^{-j})$ of the two reciprocal strides as sums of products of two numbers $\gamma_{i,l}$, and compute them for $j\le4$.
\end{enumerate}
\end{exercise}

\begin{solution}
(i) A triple $T$ other than $a$ meets $a$ in at most two marks. If it meets $a$ in two marks it is a pair of $a$ and a mark of $b$, and these monomials $\alpha^T$ add up to $\Psi_aK_b$; if it meets $a$ in one mark it is a mark of $a$ and a pair of $b$, which add up to $K_a\Psi_b$; if it meets $a$ in no mark it is $b$. So $K=\sum_{T\in\cB}\alpha^T=\Psi_aK_b+K_a\Psi_b+\alpha^b=Q+\alpha^b$. The cobases are the complements of the bases; the complement of a triple is a triple, and the complement of $a$ is $b$, so the cobases are the triples other than $b$, and the same count with $a$ and $b$ exchanged gives $\Psi=\sum_{T\in\cB}\alpha^{\cE\setminus T}=\Psi_bK_a+K_b\Psi_a+\alpha^a=Q+\alpha^a$.

(ii) Write $\Psi_a=\alpha_1\alpha_2+\alpha_3(\alpha_1+\alpha_2)$ and $K_a=\alpha_3+(\alpha_1+\alpha_2)$. Take $w$ of the $i$ factors of $\Psi_a^i$ equal to $\alpha_1\alpha_2$ and $v$ of the $l$ factors of $K_a^l$ equal to $\alpha_3$:
\[
\Psi_a^iK_a^l=\sum_{w,v}\binom iw\binom lv(\alpha_1\alpha_2)^w\,\alpha_3^{\,i-w+v}\,(\alpha_1+\alpha_2)^{\,i-w+l-v}.
\]
The power $D$ of $\alpha_3$ asks $v=D-i+w=w+\frac{l-i}3$, since $D-i=\frac{2i+l}3-i$. Then $i-w+l-v=2i+l-D-2w=2D-2w$, and the coefficient of $\alpha_1^{D-w}\alpha_2^{D-w}$ in $(\alpha_1+\alpha_2)^{2D-2w}$ is $\binom{2D-2w}{D-w}$, which gives the formula. The first values are
\[
\gamma_{0,0}=1,\quad \gamma_{1,1}=3,\quad \gamma_{3,0}=\gamma_{0,3}=6,\quad \gamma_{2,2}=15,\quad \gamma_{4,1}=\gamma_{1,4}=36,\quad \gamma_{3,3}=93 ;
\]
for instance $\gamma_{2,2}=\binom20\binom20\binom42+\binom21\binom21\binom21+\binom22\binom22\binom00=6+8+1$. Each $\gamma_{i,l}$ is the constant term $\CT\big(\alpha^{-D\one}\Psi_a^{\,i}K_a^{\,l}\big)$ on $\binom a1$: the first constant term of the three-mark ray with direction $(D\one;l,i)$.

(iii) \emph{The centre.} The stride is $f=\alpha^{\cE}(K\Psi)^{-1}$ (Definition~\ref{def:centreray}), so $\CT(f^{-j})$ is the coefficient of $\alpha^{j\one}$ in $(K\Psi)^j$; every letter exponent of the direction is $1$, and this is the closing $\mathsf C_j(1)$ of Definition~\ref{def:closings}. Expand $K^j=(\Psi_aK_b+K_a\Psi_b+\alpha^b)^j$ with $u_1,u_2,u_3$ factors of the three kinds and $\Psi^j=(K_a\Psi_b+\Psi_aK_b+\alpha^a)^j$ with $v_1,v_2,v_3$. The product splits into a part on $a$ and a part on $b$,
\[
\Psi_a^{\,u_1+v_2}K_a^{\,u_2+v_1}(\alpha^a)^{v_3}\cdot\Psi_b^{\,u_2+v_1}K_b^{\,u_1+v_2}(\alpha^b)^{u_3},
\]
and the coefficient of $\alpha^{j\one}$ is the product of the coefficient of $\alpha_1^{j-v_3}\alpha_2^{j-v_3}\alpha_3^{j-v_3}$ in the first factor without $(\alpha^a)^{v_3}$ and the like coefficient on $b$. Put $p=j-u_3=u_1+u_2$, $q=j-v_3=v_1+v_2$, $i=u_1+v_2$ and $l=u_2+v_1$. The degrees ask $2i+l=3q$ and $i+2l=3p$, that is $i=2q-p$ and $l=2p-q$, and the two coefficients are $\gamma_{2q-p,\,2p-q}$ and $\gamma_{2p-q,\,2q-p}$. The multinomial weights summed over $u_1$, with $u_2=p-u_1$, $v_2=i-u_1$, $v_1=q-i+u_1$, give
\[
\sum_{u_1}\frac{j!}{u_1!\,u_2!\,u_3!}\cdot\frac{j!}{v_1!\,v_2!\,v_3!}=\binom jp\binom jq\sum_{u_1}\binom p{u_1}\binom q{i-u_1}=\binom jp\binom jq\binom{p+q}{2q-p}
\]
by Vandermonde's identity. So
\[
\CT(f^{-j})=\sum_{p,q=0}^j\binom jp\binom jq\binom{p+q}{2q-p}\,\gamma_{2q-p,\,2p-q}\,\gamma_{2p-q,\,2q-p}.
\]
At $j=1$ the terms $(p,q)=(0,0)$ and $(1,1)$ survive, $1+2\cdot3\cdot3=19$, the number of bases. At $j=2$ the terms $(0,0)$, $(1,1)$, $(1,2)$, $(2,1)$, $(2,2)$ are
\[
1+2\cdot2\cdot2\cdot3\cdot3+2\cdot1\cdot1\cdot6\cdot6+1\cdot2\cdot1\cdot6\cdot6+1\cdot1\cdot6\cdot15\cdot15=1+72+72+72+1350=1567 ,
\]
and at $j=3$ the eight terms $(0,0),(1,1),(1,2),(2,1),(2,2),(2,3),(3,2),(3,3)$ are $1+162+324+324+12150+19440+19440+172980=224821$.

\emph{The full cobasis lean.} Here $f^{-1}=\Psi^2/\alpha^{\cE}$, and $\CT(f^{-j})$ is the coefficient of $\alpha^{j\one}$ in $\Psi^{2j}$; the direction has the letter exponent $s=0$, so Definition~\ref{def:closings} does not apply and the constant term is taken as written. With $v_1,v_2,v_3$ factors $K_a\Psi_b$, $\Psi_aK_b$, $\alpha^a$ out of $2j$ and $p=v_1+v_2$, the part on $b$ is $\Psi_b^{\,v_1}K_b^{\,v_2}$, of degree $3j$, so $2v_1+v_2=3j$: $v_1=3j-p$ and $v_2=2p-3j$; the part on $a$ is $K_a^{\,v_1}\Psi_a^{\,v_2}(\alpha^a)^{v_3}$ with $v_3=2j-p$. So
\[
\CT(f^{-j})=\sum_{p}\binom{2j}p\binom p{2p-3j}\,\gamma_{2p-3j,\,3j-p}\,\gamma_{3j-p,\,2p-3j}.
\]
At $j=1$ only $p=2$ survives: $\binom22\binom21\gamma_{1,1}^2=18$, the bases whose complement is a basis, all but $b$. At $j=2$, $p=3$ and $p=4$ give $\binom43\binom33\gamma_{0,3}\gamma_{3,0}+\binom44\binom42\gamma_{2,2}^2=144+1350=1494$, and at $j=3$, $p=5$ and $p=6$ give $\binom65\binom51\gamma_{1,4}\gamma_{4,1}+\binom66\binom63\gamma_{3,3}^2=38880+172980=211860$.

The closing values for $j\le4$, by the two formulas and, as a check, by expanding $(K\Psi)^j$ and $\Psi^{2j}$ with every exponent above $j$ dropped (receipt \texttt{nsix\_exact.py}), are
\[
\begin{array}{r|rrrrr}
j&0&1&2&3&4\\\hline
\text{centre }(\one;1,1)&1&19&1567&224821&41485807\\
\text{lean }(\one;0,2)&1&18&1494&211860&38486070
\end{array}
\]
So the constant terms of $\mathrm N_6$ are sums of products of two first constant terms of the three-mark list $\binom{[3]}1$ along its directions $(D\one;l,i)$, one factor for each of the triples $a$ and $b$; the sums cost $O(j^2)$ and $O(j)$ products of numbers $\gamma_{i,l}$, each a sum of $O(j)$ terms.
\end{solution}

\begin{exercise}[$\mathrm N_6$ at its full cobasis lean]\label{ex:listsb-nlean}
Read the ray of $\mathrm N_6$ (Exercise~\ref{ex:listsb-nletters}) with origin $0$ and direction $(\one;0,2)$, whose stride is $f=\alpha^{\cE}\Psi^{-2}$, by the procedure of Definition~\ref{def:readprocedure}:
\begin{enumerate}[label=(\roman*)]
\item its walls and the volume $V_1$ of its direction;
\item the native gauge $\alpha^\star$ of the direction, $f^\star=f(\alpha^\star)$ and the Hessian $H$ there;
\item the box with one unit of each active letter, its spanning set and its matrix $\mathbf B$, the matrix at infinite strength and its spectrum, and that the spanning set is a set of masters;
\item the law of $[1]$, its two end coefficients, its symbol, and the strengths at which its rows line up;
\item the low-end relation and the constant terms;
\item the value $\langle1\rangle_1=P_{\mathrm N_6}(\one;0,2)$.
\end{enumerate}
\end{exercise}

\begin{solution}
Throughout, $S$ is a nonempty proper subset of $\cE$, $r(S)$ the largest trace of a basis on $S$, and the \emph{type} of $S$ is the pair $(i,j)=(|S\cap a|,|S\cap b|)$.

(i) \emph{Walls.} By Definition~\ref{def:family} the walls of the direction are $\kappa_{1,S}=m_{1,S}-s\,d(S)-t\,d^\ast(S)$, with $d(S)$ and $d^\ast(S)$ the least traces of a basis and of a cobasis on $S$. Here $m_1=\one$, $s=0$, $t=2$, and the least trace of a cobasis is $d^\ast(S)=|S|-r(S)$, so
\[
\kappa_{1,S}=2\,r(S)-|S| .
\]
The rank of $\mathrm N_6$ is $r(S)=\min(|S|,3)$, except $r(a)=2$, so $\kappa_{1,S}$ depends on the type alone:
\[
\begin{array}{c|ccccccc}
(i,j)&(1,0)&(0,1)&(2,0)&(1,1)&(0,2)&(3,0)&(2,1)\\\hline
\kappa_{1,S}&1&1&2&2&2&1&3\\
F(i,j)&1&1&1&1&1&3&1
\end{array}
\]
\[
\begin{array}{c|ccccccc}
(i,j)&(1,2)&(0,3)&(3,1)&(2,2)&(1,3)&(3,2)&(2,3)\\\hline
\kappa_{1,S}&3&3&2&2&2&1&1\\
F(i,j)&1&1&3&2&2&12&10
\end{array}
\]
Every wall is positive, the least is $1$: the direction is strictly inside its walls, and counted with the sizes of the types the $62$ walls are $1$ thirteen times, $2$ thirty times and $3$ nineteen times.

\emph{The volume.} By Proposition~\ref{prop:principal}, $V_1=\sum_\pi\prod_{l=2}^6\kappa_{1,U_l}^{-1}$ over the orders $\pi$ of $\cE$ with their tails $U_l=\{\pi(l),\dots,\pi(6)\}$. An order is a chain of tails $U_6\subset\dots\subset U_2$, so $V_1=\sum_eF(\cE\setminus\{e\})$ with $F(\{e\})=\kappa_{1,\{e\}}^{-1}$ and $F(U)=\kappa_{1,U}^{-1}\sum_{e\in U}F(U\setminus\{e\})$. $F$ depends on the type alone, $F(i,j)=\kappa_{1,(i,j)}^{-1}\big(i\,F(i-1,j)+j\,F(i,j-1)\big)$, which gives the last row of the table: for instance $F(3,0)=1^{-1}\cdot3\,F(2,0)=3$, $F(3,1)=2^{-1}\big(3F(2,1)+F(3,0)\big)=3$, $F(3,2)=1^{-1}\big(3F(2,2)+2F(3,1)\big)=12$ and $F(2,3)=1^{-1}\big(2F(1,3)+3F(2,2)\big)=10$. So
\[
V_1=3\,F(2,3)+3\,F(3,2)=30+36=66 .
\]

(ii) \emph{The native gauge.} By Theorem~\ref{thm:pencil}(i), $f\le f(\alpha^\star)$ on the positive orthant with equality only at $\alpha\propto\alpha^\star$. The $36$ automorphisms of $\mathrm N_6$, the permutations of $a$ and of $b$, keep the masses $\one$ and the exponents, so they are symmetries of the ray (Proposition~\ref{prop:raysym}(ii)); they keep $f$ and the orthant and therefore fix $\alpha^\star$ up to scale. On the chart $\alpha_6=1$ this puts $\alpha^\star$ on the slice $\alpha=(x,x,x,1,1,1)$, where $Q=3x^2\cdot3+3x\cdot3$, $\Psi=x^3+9x^2+9x$ and
\[
f=\frac{x^3}{x^2(x^2+9x+9)^2}=\frac x{(x^2+9x+9)^2},\qquad \frac d{dx}\log f=-\frac{3(x^2+3x-3)}{x(x^2+9x+9)} .
\]
So $\alpha^\star=(x^\star,x^\star,x^\star,1,1,1)$ with $x^\star=\frac{\sqrt{21}-3}2$, and with $x^2=3-3x$, $x^2+9x+9=6(x+2)$ and $f=x/(36(x+2)^2)$ there:
\[
f^\star=\frac{7\sqrt{21}-27}{1800}=0.0028211277026060444701\ldots
\]
At a complex point fixed by the automorphisms the differential of $\log f$ in $y=\log\alpha$ is fixed by them too, so it is constant on $a$ and on $b$, $(u,u,u,v,v,v)$, and $u+v=0$ because $f$ is homogeneous of degree $0$; it vanishes as soon as its derivative along the slice, $3u$, does. So the other root $x'=-\frac{\sqrt{21}+3}2$ gives a native gauge $\alpha'=(x',x',x',1,1,1)$ too (Definition~\ref{def:raygauge}), at which $K=9x'^2+9x'+1$ and $\Psi=6x'(x'+2)$ do not vanish, of the conjugate native value $-\frac{7\sqrt{21}+27}{1800}$. The Hessian of $-\log f$ at $\log\alpha^\star$ on the chart is $H=2\big(\kappa_{\alpha^\star}(X_e,X_{e'})\big)_{e,e'\neq6}$, twice the covariances of the indicators under the cobasis law at $\alpha^\star$ (Theorem~\ref{thm:branch}); computed in $\mathbb Q(\sqrt{21})$ (receipt \texttt{hessians.py}),
\[
\det H=\frac{13049\sqrt{21}-59339}{37500}=0.0122408064927668332\ldots
\]

(iii) \emph{The box} (Procedure~\ref{def:readprocedure}, the masters). The letter $K$ has the exponent $0$ along the ray and $\Psi$ the exponent $2k$, so only $\Psi$ is active, and the box with one unit of each active letter is $\mathfrak B_2(0,1)$ of Definition~\ref{def:raybox}(iv): the symbols of wall degree at most $2$ at the levels $(0,0)$ and $(0,1)$; its stems are the Laurent symbols of wall degree at most $2$. The $36$ automorphisms keep it (Definition~\ref{def:raybox}(vi)), and modulo them (Definition~\ref{def:box-reduction}(i)) there are $104$ orbits of stems and $329$ orbits of symbols; the relations are the $104$ letter identities $\alpha^u-\sum_{T}\alpha^{u+\one_{\cE\setminus T}}\Psi^{-1}$ of the stems and their $104\cdot6=624$ mark relations, $728$ in all. One elimination, the orbits ranked as in Definition~\ref{def:box-reduction}(ii), done modulo a prime of $61$ bits at a random strength, which gives the pivots over $\mathbb Q(k)$ off finitely many strengths, leaves free among the Laurent orbits of wall degree at most $1$ those of the seven orbit averages
\begin{gather*}
b_1=1,\quad b_2=\Big[\frac{\alpha_3\alpha_5\alpha_6}{\alpha_1\alpha_2\alpha_4}\Big],\quad b_3=\Big[\frac{\alpha_3\alpha_6}{\alpha_1\alpha_2}\Big],\quad b_4=\Big[\frac{\alpha_5\alpha_6}{\alpha_1\alpha_4}\Big],\\
b_5=\Big[\frac{\alpha_6}{\alpha_1}\Big],\quad b_6=\Big[\frac{\alpha_3\alpha_6}{\alpha_1\alpha_4}\Big],\quad b_7=\Big[\frac{\alpha_2\alpha_3\alpha_6}{\alpha_1\alpha_4\alpha_5}\Big],
\end{gather*}
$[\,\cdot\,]$ the average over the $36$ automorphisms, of wall degree $0$ and $1$ (Proposition~\ref{prop:box-walldegree}(ii): the wall degree of $\alpha^u$ is $\max_U(-u_U)/\kappa_{1,U}$; for $b_2,\dots,b_7$ it is $1$, reached at $U=\{1\}$), and every target $f^{-1}b_i$, of wall degree at most $2$, reduces to them. So the box is accepted at $\mathfrak s=1$ (Definition~\ref{def:box-reduction}(iii)), and by Theorem~\ref{thm:raybox}(i) $[f^{-1}b_i]=\sum_l\mathbf B_{il}(k)[b_l]$ in the module, the classes of $b_1,\dots,b_7$ spanning a submodule $\mathfrak N$ that contains $[1]$.

\emph{The matrix} (Procedure~\ref{def:readprocedure}, the matrix over $\mathbb Q(k)$). Input: the $728$ relations, with coefficients $c_0+c_1k$, $c_0,c_1\in\mathbb Z$. Method: the targets need $292$ of them, of rank $292$; eliminations modulo primes of $62$ bits at random strengths, $24$ strengths on the first prime, every entry read as a quotient of polynomials, and Chinese remaindering with rational reconstruction over three primes. Output: the common denominator of the entries is $Q_c=(k-2)^5(k+1)$, and the row $i$ of $\mathbf B$ is a factor $\phi_i$ times a row of polynomials:
{\small
\begin{longtable}{@{}c >{$}l<{$} >{\raggedright\arraybackslash}p{0.66\textwidth}@{}}
\toprule
$i$&\phi_i&$\mathbf B_{il}/\phi_i$, $l=1,\dots,7$\\\midrule
$1$&\frac{6(2k-1)}{(k-2)(k+1)}&$6$;\ $4(k-2)$;\ $6(k-1)$;\ $6(k-1)$;\ $8(k+1)$;\ $-24$;\ $-(k+1)$\\[3pt]
$2$&\frac{2(k-1)(2k-1)}{(k-2)^5(k+1)}&$k^4+20k^3-36k^2+40k-16$;\ $2k(k-2)(8k^2-19k+10)$;\ $2k(12k^3-37k^2+42k-20)$;\ $2k(12k^3-37k^2+42k-20)$;\ $4k(k+1)(8k^2-12k+9)$;\ $-4(k^4+20k^3-48k^2+36k-8)$;\ $-8(k-1)^3(k+1)$\\[3pt]
$3$&\frac{2(k-1)(2k-1)}{(k-2)^4(k+1)}&$-2k^3+33k^2-42k+16$;\ $(k-2)(15k^2-33k+14)$;\ $25k^3-75k^2+70k-16$;\ $19k^3-57k^2+70k-40$;\ $2k(k+1)(14k-11)$;\ $4(2k^3-27k^2+44k-20)$;\ $-(k-1)(k+1)(5k-6)$\\[3pt]
$4$&\frac{2(k-1)(2k-1)}{(k-2)^4(k+1)}&the row $3$ with its third and fourth entries exchanged\\[3pt]
$5$&\frac{2(k-1)(2k-1)}{(k-2)^3(k+1)}&$24(k-1)$;\ $4(k-2)(3k-5)$;\ $2(9k^2-23k+16)$;\ $2(9k^2-23k+16)$;\ $(k+1)(33k-32)$;\ $4(3k^2-23k+22)$;\ $-(k+1)(3k-4)$\\[3pt]
$6$&\frac{2(k-1)(2k-1)}{(k-2)^4(k+1)}&$2k^3+11k^2-26k+16$;\ $(k-2)(k-1)(7k-10)$;\ $13k^3-41k^2+40k-8$;\ $13k^3-41k^2+40k-8$;\ $(k+1)(33k^2-73k+48)$;\ $4(k^3-13k^2+23k-14)$;\ $-(k+1)(k^2-2)$\\[3pt]
$7$&\frac{2(k-1)(2k-1)}{(k-2)^5(k+1)}&$-2k^4+23k^3-44k^2+53k-16$;\ $(k-2)(k-1)^2(16k-7)$;\ $27k^4-94k^3+99k^2-34k+22$;\ $27k^4-94k^3+99k^2-34k+22$;\ $(k+1)(8k^3+17k^2-35k+32)$;\ $-13k^4-39k^3+231k^2-273k+22$;\ $-(k+1)(17k^3-80k^2+131k-73)$\\
\bottomrule
\end{longtable}}
Check: $\mathbf B$ agrees with an elimination of the whole box at a prime and a strength drawn at random, and again at three fresh ones. Receipt: \texttt{rayreader.py N "1 1 1 1 1 1 ; 0 2"}, read by \texttt{analyze\_run.py} and \texttt{matrix\_rows\_tex.py}; the re-check by \texttt{tests/check\_run.py}.

\emph{Infinite strength.} Every entry has a numerator of degree at most that of its denominator, so $\mathbf B$ is regular at $k=\infty$, with
\[
\mathbf B(\infty)=\begin{pmatrix}0&48&72&72&96&0&-12\\4&64&96&96&128&-16&-32\\-8&60&100&76&112&32&-20\\-8&60&76&100&112&32&-20\\0&48&72&72&132&48&-12\\8&28&52&52&132&16&-4\\-8&64&108&108&32&-52&-68\end{pmatrix},
\]
\[
\det\big(\lambda-\mathbf B(\infty)\big)=(\lambda-24)^2\,(\lambda^2-324\lambda-10800)\,(\lambda^3+28\lambda^2-320\lambda+1728).
\]
By Proposition~\ref{prop:rays-spectrum}(ii), at every nondegenerate native gauge $\alpha_c$ the vector $\beta_c=(b_l(\alpha_c))_l$ is an eigenvector of $\mathbf B(\infty)$ for $f(\alpha_c)^{-1}$. At $\alpha^\star$ every monomial of the orbit of $\alpha^u$ takes the value $(x^\star)^{u_1+u_2+u_3}$, so $\beta_\star=(1,\frac1{x^\star},\frac1{x^\star},\frac1{x^\star},\frac1{x^\star},1,x^\star)$ with $\frac1{x^\star}=\frac{\sqrt{21}+3}6$; the first row gives
\[
\frac{48+72+72+96}{x^\star}-12x^\star=48(\sqrt{21}+3)-6(\sqrt{21}-3)=162+42\sqrt{21}=\frac{1800}{7\sqrt{21}-27}=\frac1{f^\star},
\]
and the other rows the same, exactly. The eigenvalue $24$ belongs to two rational native gauges: at $\alpha=(-1,\frac32,-1,1,1,1)$ and at $\alpha=(-1,-1,-1,1,-\frac23,1)$ every dual occupation is $R^\ast_e=\frac12$, so every defect $1-2R^\ast_e$ vanishes, $f=\frac1{24}$, and the Hessian of $\log f$ on the chart has the determinant $\frac{625}{864}$ at both; their vectors
\[
\beta=\big(1,\tfrac{17}{18},\tfrac{17}{18},-\tfrac49,-\tfrac49,-\tfrac7{18},\tfrac{11}9\big),\qquad \beta'=\big(1,\tfrac{17}{18},-\tfrac49,\tfrac{17}{18},-\tfrac49,-\tfrac7{18},\tfrac{11}9\big)
\]
satisfy $\mathbf B(\infty)\beta=24\beta$ and $\mathbf B(\infty)\beta'=24\beta'$ exactly, and are independent. Receipt: \texttt{nsix\_gauges24.py}.

\emph{The native gauges and the masters.} Input: the critical equations $1-2R^\ast_e(\alpha)=0$ on the chart, saturated by the coordinates and the letters. Method: a Groebner basis over $\mathbb Q$, and, separately, the critical points of $\alpha^{\mathbf a}\Psi^{-b}$ for random complex exponents counted by monodromy and carried to $(\one;2)$. Output: the ideal is zero-dimensional with $35$ points counted with multiplicity, and the $35$ native gauges found by the paths are distinct and nondegenerate, in $7$ orbits of the automorphisms of sizes $1,1,3,3,9,9,9$: $\alpha^\star$, $\alpha'$, the two rational orbits of native value $\frac1{24}$, and three orbits whose native values are the roots of $1728\lambda^3-320\lambda^2+28\lambda+1$, the real one $-0.0265230928\ldots$ and $0.1058541390\ldots\pm0.1030231201\ldots\,i$. Check: Newton's method at $40$ digits at each point, residual below $10^{-38}$; the seven vectors $\beta_c$, one per orbit, have the singular values from $8.93$ down to $0.358$, and each is an eigenvector of $\mathbf B(\infty)$ to $10^{-38}$. Receipts: \texttt{gb\_direction.py}, \texttt{critpoints.py}, \texttt{polish.py}, \texttt{betas.py}. So the vectors $\beta_c$ span $\mathbb C^7$, and by Proposition~\ref{prop:rays-spectrum}(i) the spanning set is a set of masters of $\mathfrak N$: $\dim_{\mathbb C(k)}\mathfrak N=7$, $\mathbf B=\mathbf B_{\mathbf b}$, and $\mathbf B(\infty)$ is diagonalizable with the eigenvalues $f(\alpha_c)^{-1}$, $24$ twice for the two orbits of native value $\frac1{24}$.

(iv) \emph{The law} (Procedure~\ref{def:readprocedure}, the law). Input: $\mathbf B$. Method: the rows $w_0=(1,0,\dots,0)$, $w_i(k)=w_{i-1}(k-1)\,\mathbf B(k)$ of Proposition~\ref{prop:rays-lawrows}, evaluated modulo primes at consecutive strengths; their first dependence; the coefficients reconstructed over $\mathbb Q$ from four primes and made integral and primitive. Output: $w_0,\dots,w_5$ are independent and $w_0,\dots,w_6$ dependent, so $h=6$, and the relation $\sum_{j=0}^6c_j(k)\langle1\rangle_{k+j}=0$ of Proposition~\ref{prop:rays-lawrows}(ii) has $c_j\in\mathbb Z[k]$ of degree $24$, with coefficients of up to $31$ digits. Check: the relation holds at a prime and a point drawn at random, and at three fresh ones. Since $b_1,\dots,b_7$ are masters, Proposition~\ref{prop:rays-lawrows}(ii) makes it the law of $[1]$ (Definition~\ref{def:raymodule}): the cyclic submodule of $[1]$ has dimension $6$, inside $\mathfrak N$ of dimension $7$. Factored over $\mathbb Q$,
\begin{gather*}
c_0=k^5(k+1)^3(k+2)(k+3)\,r(k),\\
c_6=6998400\,(2k+1)(2k+3)(2k+5)(2k+7)(2k+9)(2k+11)(k+4)(k+5)^3\,r(k-1),
\end{gather*}
with $r$ irreducible of degree $14$ and integer coefficients of up to $18$ digits, and the symbol (Corollary~\ref{cor:symbol}), the coefficients of $k^{24}$,
\[
\chi(\lambda)=7276255875\,(24\lambda-1)(10800\lambda^2+324\lambda-1)(1728\lambda^3-320\lambda^2+28\lambda+1).
\]
Its six roots are $\frac1{24}$, $f^\star$ and its conjugate (the roots of $10800\lambda^2+324\lambda-1$, $\frac{-27\pm7\sqrt{21}}{1800}$), and the three roots of the cubic: the six distinct native values of (iii), as Corollary~\ref{cor:symbol} asks, and the reciprocals of the eigenvalues of $\mathbf B(\infty)$, $24$ counted once.

\emph{Line-up strengths.} In the notation of Proposition~\ref{prop:rays-lawrows}, $p_j=c_j$ and $q_6(k)=c_0(k-6)$, and Proposition~\ref{prop:rays-lawrows}(iii) gives
\[
c_0(k-6)\ \Big|\ \det\mathbf W(k)_J\prod_{j=0}^{5}Q_c(k-j)^{6-j},\qquad Q_c=(k-2)^5(k+1),
\]
for every set $J$ of six columns with $\det\mathbf W_J\neq0$. The irreducible factor $r(k-6)$, of degree $14$, divides no linear factor $Q_c(k-j)$, so it divides $\det\mathbf W(k)_J$: at its fourteen roots, none a root of a $Q_c(k-j)$, the rows $w_0,\dots,w_5$ are dependent. So $14$ of the $24$ degrees of $c_0$ are strengths at which $\langle1\rangle_k,\dots,\langle1\rangle_{k-5}$, read as linear forms in the values of the masters, line up; the other $10$ are the integers $0,-1,-2,-3$.

(v) \emph{The low end.} By Corollary~\ref{cor:rays-lowend}(ii), $c_0$ vanishes to order $n-1=5$ at $k=0$, as it does, and with $\mathsf L_0=V_1=66$ (Definition~\ref{def:rays-lowend}) the low-end relation $[k^5]c_0\cdot66+\sum_{i=1}^6c_i(0)\langle1\rangle_i=0$, divided by $91238400$, reads
\begin{multline*}
64402904777-12632598108\,\langle1\rangle_1-2849737197672\,\langle1\rangle_2+380837271432240\,\langle1\rangle_3\\
+14853466134400320\,\langle1\rangle_4-369781615117353600\,\langle1\rangle_5+3400346988339840000\,\langle1\rangle_6=0 .
\end{multline*}
By Corollary~\ref{cor:rays-lowend}(i) the constant terms obey $\sum_{i=0}^{\min(6,j)}c_i(-j)\CT(f^{i-j})=0$ for every $j\in\mathbb N$; with the constant terms of Exercise~\ref{ex:listsb-nletters} this holds exactly for $0\le j\le60$ (receipt \texttt{nsix\_ctcheck.py}), a check of the law that uses neither the box nor the values.

(vi) \emph{The value} (Procedure~\ref{def:readprocedure}, the value). By Theorem~\ref{thm:gaugeend} the values are the moments $\langle1\rangle_k=\int_0^{f^\star}t^k\varrho_f(t)\,dt$ of the density $\varrho_f$ of the stride's law (Definition~\ref{def:strideslaw}); $f^\star$ is a simple root of $\chi$, the density operator of the law has order $24$, its exponent at $f^\star$ beyond $0,\dots,22$ is $\frac{n-3}2=\frac32$, and the germ of $\varrho_f$ is $(f^\star-t)^{3/2}g(f^\star-t)$ with $g$ fixed by $g(0)=\big(2\Gamma(\frac12)^2/f^\star\big)^{5/2}\big/\big(\Gamma(\frac52)\sqrt{\det H}\big)$, the constant of the branch of Theorem~\ref{thm:branch}. Input: the law, $f^\star$ and $\det H$. Method: the germ to $168$ terms, the Beta sums $\int_0^{f^\star}t^K(f^\star-t)^{\frac32+m}dt=(f^\star)^{K+\frac52+m}\Gst21{K+1,\ \frac52+m}{K+\frac72+m}$ at the six strengths from $K=57$ on, and the law run back to $k=1$, in ball arithmetic at $1188$ bits; again from $K=82$. Output:
\[
P_{\mathrm N_6}(\one;0,2)=\langle1\rangle_1=4.75680256375822944340803728605708\ldots,
\]
the two starting strengths agreeing to $50$ digits. Check: the exponent read from the law at $f^\star$ is $\frac32$ to $16$ digits; the reading at $700$ digits of Exercise~\ref{ex:listsb-nspan} agrees with this value, and on its values $\langle1\rangle_1,\dots,\langle1\rangle_6$ the low-end relation of (v), which the reading does not use, holds to $731$ digits relative to its largest term. Receipt: \texttt{rayreader.py N "1 1 1 1 1 1 ; 0 2" --digits 30}.

So $\mathrm N_6$ at $(\one;0,2)$ has the volume $V_1=66$, the native value $f^\star=\frac{7\sqrt{21}-27}{1800}$, seven masters with $\mathbf B(\infty)$ of spectrum $\{24,24\}$, the reciprocals of the roots of $10800\lambda^2+324\lambda-1$ and of $1728\lambda^3-320\lambda^2+28\lambda+1$, a law of order $6$ and degree $24$, and the value $P_{\mathrm N_6}(\one;0,2)=4.7568025637582294434\ldots$, with $P/V_1=0.0720727\ldots$
\end{solution}

\begin{exercise}[$\mathrm N_6$ at its centre]\label{ex:listsb-ncentre}
Read the centre ray of $\mathrm N_6$ at the strength $1$, with the direction $(\one;1,1)$ and the stride $f=\alpha^{\cE}(K\Psi)^{-1}$, by the procedure of Definition~\ref{def:readprocedure}: its walls and $V_1$; the native gauge of the direction and the signs; the box, the spanning set and its matrix at infinite strength; the law of $[1]$, its ends, its symbol and its line-up strengths; the low-end relation, the closings and the value.
\end{exercise}

\begin{solution}
\emph{Walls and volume.} With $m_1=\one$ and $s=t=1$, Definition~\ref{def:family} and $d(S)=3-r(\cE\setminus S)$, $d^\ast(S)=|S|-r(S)$ give $\kappa_{1,S}=r(S)+r(\cE\setminus S)-3$, which depends on the type $(i,j)=(|S\cap a|,|S\cap b|)$ alone. With the recursion of Exercise~\ref{ex:listsb-nlean}(i):
\[
\begin{array}{c|ccccccc}
(i,j)&(1,0)&(0,1)&(2,0)&(1,1)&(0,2)&(3,0)&(2,1)\\\hline
\kappa_{1,S}&1&1&2&2&2&2&3\\
F(i,j)&1&1&1&1&1&\frac32&1
\end{array}
\]
\[
\begin{array}{c|ccccccc}
(i,j)&(1,2)&(0,3)&(3,1)&(2,2)&(1,3)&(3,2)&(2,3)\\\hline
\kappa_{1,S}&3&2&2&2&2&1&1\\
F(i,j)&1&\frac32&\frac94&2&\frac94&\frac{21}2&\frac{21}2
\end{array}
\]
Every wall is positive: $1$ twelve times, $2$ thirty-two times and $3$ eighteen times, the triples $a$ and $b$ having the wall $2$. By Proposition~\ref{prop:principal},
\[
V_1=3\,F(2,3)+3\,F(3,2)=\tfrac{63}2+\tfrac{63}2=63 .
\]

\emph{The native gauge and the signs.} On the centre ray $\alpha^\star=\one$ and $f^\star=K(\one)^{-2}=19^{-2}=\frac1{361}$ (Theorem~\ref{thm:pencil}, Definition~\ref{def:centreray}), and $H=2\big(\kappa_{\one}(X_e,X_{e'})\big)_{e,e'\neq6}$, the covariances under the weight law, which is uniform on the nineteen bases (Theorem~\ref{thm:branch}): each mark of $a$ lies in $9$ bases and each mark of $b$ in $10$, and
\[
H=\frac1{361}\begin{pmatrix}180&-48&-48&-28&-28\\-48&180&-48&-28&-28\\-48&-48&180&-28&-28\\-28&-28&-28&180&-48\\-28&-28&-28&-48&180\end{pmatrix},\qquad\det H=\frac{580608}{361^3}.
\]
By Theorem~\ref{thm:signgauges}(ii) every sign $\eta$ with $K(\eta)\neq0$ is a native gauge of native value $K(\eta)^{-2}$. A sign is fixed by its partition of the marks into the scales $1$ and $-1$, and the $32$ signs of the chart $\eta_6=1$ fall into eight orbits of the automorphisms: $K(\eta)=19$ when all scales agree; $|K(\eta)|=5$ when two marks of $a$, or two marks of $b$, stand against the other four, two orbits of three signs; $|K(\eta)|=3$ when one mark of $a$ and one of $b$ stand against the other four, nine signs; and $|K(\eta)|=1$ on the other four orbits, one mark against five ($3$ and $3$ signs), $a$ against $b$ ($1$), and two marks of $a$ with one of $b$ against the rest ($9$). The Hessian of $\log f$ at a sign, in the chart $\alpha_e=\eta_ee^{y_e}$, has the determinants $-\frac{580608}{361^3}$, $\frac{9216}{15625}$, $\frac{1024}{729}$ at the first three kinds and $0$ at the last: the native values $\frac1{361}$, $\frac1{25}$ (two orbits) and $\frac19$ are nondegenerate, and the native value $1$ is degenerate at every sign (receipt \texttt{nsix\_signs.py}). On the slice $(x,x,x,1,1,1)$, fixed by the automorphisms, $f=x^2/\big((x^2+9x+9)(9x^2+9x+1)\big)$ and $\frac d{dx}\log f$ has the numerator $-18(x-1)(x+1)(x^2+5x+1)$; as at the lean, a critical point of the slice is a native gauge, and $x^2+5x+1=0$ gives the two native gauges $\big(\frac{-5\pm\sqrt{21}}2\big)$ of native value $-\frac1{80}$, exchanged by the Cremona.

\emph{The box} (Procedure~\ref{def:readprocedure}, the masters). Both letters are active, and the symmetries are the $36$ automorphisms and the Cremona (Proposition~\ref{prop:raysym}), $72$ maps, which keep every box with equal level bounds. The box $\mathfrak B_2(1,1)$ is not accepted: at $\mathfrak s=1$ eight Laurent orbits are free and seven of their targets do not reduce to them, and at $\mathfrak s=0$ the target of $1$ does not reduce. The box $\mathfrak B_3(1,1)$ has $252$ orbits of stems and $1909$ orbits of symbols, and $3774$ relations: the letter identities of the symbols at the levels $(0,0)$, $(1,0)$ and $(0,1)$ and the $252\cdot6$ mark relations of the stems. It is accepted at $\mathfrak s=2$, with the spanning set of the fifteen averages, over the $72$ maps, of $1$ and of $\alpha^u$ for
\begin{gather*}
u=(-1,-1,0,0,1,1),\ (-1,-1,1,-1,1,1),\ (-1,-1,1,0,0,1),\ (-1,0,0,-1,1,1),\\
(-1,0,0,0,0,1),\ (-1,0,1,-1,0,1),\ (-1,0,1,0,0,0),\ (-1,-1,-1,1,1,1),\\
(-2,-2,0,0,2,2),\ (-2,-2,0,1,1,2),\ (-2,-2,1,-1,2,2),\\
(-2,-2,1,0,1,2),\ (-2,-2,1,1,1,1),\ (-2,-2,2,-2,2,2),
\end{gather*}
of the wall degrees $1$ (the first seven), $\frac32$ (the eighth, at $U=a$) and $2$ (Theorem~\ref{thm:raybox}(i)).

\emph{The matrix.} Input: the $3774$ relations. Method: as in Exercise~\ref{ex:listsb-nlean}; the targets need $1023$ relations, of rank $1023$; $55$ strengths on the first prime; three primes. Output: $\mathbf B\in\operatorname{Mat}_{15}(\mathbb Q(k))$ with $8$ vanishing entries, numerators and denominators of degree at most $17$ in lowest terms, and the common denominator
\[
Q_c=(2k-5)(k-2)^2(k-3)^5(k-4)\,S_9(k),
\]
$S_9$ irreducible of degree $9$, $S_9=6555027k^9-143215495k^8+\dots-6615463200$. Check: $\mathbf B$ agrees with eliminations of the whole box at four primes and strengths drawn at random. Receipt: \texttt{rayreader.py N "1 1 1 1 1 1 ; 1 1"}, read by \texttt{analyze\_run.py}. $\mathbf B$ is regular at $k=\infty$, and
\[
\det\big(\lambda-\mathbf B(\infty)\big)=(\lambda-361)\,(\lambda-25)^2\,(\lambda-9)\,(\lambda-1)^8\,(\lambda+80)\,\big(\lambda^2+\tfrac{34}3\lambda-\tfrac{361}3\big).
\]
By Proposition~\ref{prop:rays-spectrum}(ii) the vectors $\beta_\eta$ of the nondegenerate signs and $\beta$ of the slice gauge are eigenvectors: computed exactly, $\mathbf B(\infty)\beta_\eta=K(\eta)^2\beta_\eta$ at the signs of native values $\frac1{361}$, $\frac1{25}$ (both orbits) and $\frac19$, and $\mathbf B(\infty)\beta=-80\beta$ at $\big(\frac{-5-\sqrt{21}}2\big)$, five independent vectors (receipt \texttt{nsix\_centre\_exact.py}). The roots of $\lambda^2+\frac{34}3\lambda-\frac{361}3$ are $\frac{-17\pm14\sqrt7}3$, the reciprocals of $\frac{17\pm14\sqrt7}{361}$, which are native values of gauges off the signs and the slice: the monodromy count of Exercise~\ref{ex:listsb-nlean}, run at this direction, reaches eighteen nondegenerate native gauges at each (receipt \texttt{critpoints.py}).

\emph{The law.} The rows of Proposition~\ref{prop:rays-lawrows} first depend at $h=11$, and the relation $\sum_{j=0}^{11}c_j(k)\langle1\rangle_{k+j}=0$ has $c_j\in\mathbb Z[k]$ of degree $52$ with coefficients of up to $88$ digits, from ten primes; it holds at random points of four primes. By Proposition~\ref{prop:rays-lawrows}(ii) this relation, of order $11$ with coprime coefficients, annihilates $[1]$, so the values satisfy it; its order is at least the dimension of the cyclic submodule of $[1]$, with equality when the spanning set is a set of masters (Definition~\ref{def:raymodule}). Factored over $\mathbb Q$,
\begin{gather*}
c_0=3\,(2k+1)\,k^5(k+1)^2(k+2)^2(k+3)\,S_{41}(k),\\
c_{11}=-4691556000\,(k+4)(k+5)(k+6)(k+7)^2(k+8)^2(k+9)^2(k+10)^2\,S_{41}(k-1),
\end{gather*}
with $S_{41}$ irreducible of degree $41$, and the symbol is
\[
\chi(\lambda)\propto(361\lambda-1)(25\lambda-1)(9\lambda-1)(\lambda-1)^5(80\lambda+1)(361\lambda^2-34\lambda-3),
\]
whose roots are the native values found: $\frac1{361},\frac1{25},\frac19$, the degenerate $1$ five times, $-\frac1{80}$ and $\frac{17\pm14\sqrt7}{361}$. The eigenvalues of $\mathbf B(\infty)$ are the reciprocals of the roots of $\chi$, with their multiplicities, and one more $25$ and three more $1$.

\emph{Line-up strengths.} $q_{11}(k)=c_0(k-11)$ divides $\det\mathbf W(k)_J\prod_{j=0}^{10}Q_c(k-j)^{11-j}$ (Proposition~\ref{prop:rays-lawrows}(iii)). The irreducible factors of $Q_c$ have degrees $1$ and $9$, so $S_{41}(k-11)$ divides $\det\mathbf W(k)_J$: $41$ of the $52$ degrees of $c_0$ are strengths at which the rows $w_0,\dots,w_{10}$ line up, and the other $11$ are the rational roots $0,-\frac12,-1,-2,-3$.

\emph{The low end and the closings.} $c_0$ vanishes to order $5$ at $0$, and the low-end relation of Corollary~\ref{cor:rays-lowend}(ii), $[k^5]c_0\cdot63+\sum_{i=1}^{11}c_i(0)\langle1\rangle_i=0$, has integer coefficients of up to $69$ digits after division by their greatest common divisor $11377315676160$ (receipt \texttt{lowend\_forms.py}). The closings $\mathsf C_j(1)=\CT(f^{-j})$ of Exercise~\ref{ex:listsb-nletters} satisfy $\sum_{i=0}^{\min(11,j)}c_i(-j)\,\mathsf C_{j-i}(1)=0$ exactly for $0\le j\le60$ (Corollary~\ref{cor:reflected}; receipt \texttt{nsix\_ctcheck.py}).

\emph{The value.} By Theorem~\ref{thm:gaugeend}, with $f^\star=\frac1{361}$ a simple root of $\chi$, the density operator of order $52$, the exponent $\frac32$ at $f^\star$ (read from the law to $16$ digits), and $g(0)=\big(2\Gamma(\frac12)^2\cdot361\big)^{5/2}\big/\big(\Gamma(\frac52)\sqrt{\det H}\big)$: Beta sums from $K=57$ and from $K=82$ with $168$ terms of the germ at $1188$ bits, the law run back to $k=1$. Output:
\[
P_{\mathrm N_6}(\one;1,1)=\langle1\rangle_1=4.57891247888240480638020373022855\ldots,
\]
the two starts agreeing to $50$ digits. Receipt: \texttt{rayreader.py N "1 1 1 1 1 1 ; 1 1" --digits 30}.

So the centre of $\mathrm N_6$ has $V_1=63$, $f^\star=\frac1{361}$, a box accepted at $\mathfrak B_3(1,1)$ with fifteen spanning elements, the law of $[1]$ of order $11$ and degree $52$, and $P_{\mathrm N_6}(\one;1,1)=4.5789124788824048\ldots$, with $P/V_1=0.0726811\ldots$
\end{solution}

\begin{exercise}[The value span of $\mathrm N_6$]\label{ex:listsb-nspan}
The \emph{value span} of a ray is the $\mathbb Q$-span of its values $\langle1\rangle_k$, $k\ge1$. For the two rays of $\mathrm N_6$ of Exercises~\ref{ex:listsb-nlean} and~\ref{ex:listsb-ncentre}:
\begin{enumerate}[label=(\roman*)]
\item bound the value span by the law of $[1]$ and the low-end relation;
\item find the integer relations of small height among $1$ and the first values, read to $700$ digits, by lattice reduction, and among them and the numbers $\ell=P_{\binom{[3]}1}(\one;1,1)$ and $\Tx3{\IZ}0$.
\end{enumerate}
\end{exercise}

\begin{solution}
(i) \emph{The lean.} The leading coefficient of the law, $c_6=6998400\,(2k+1)(2k+3)(2k+5)(2k+7)(2k+9)(2k+11)(k+4)(k+5)^3\,r(k-1)$, has its rational roots at negative numbers, and $r$, irreducible of degree $14$, has none; so $c_6(k)\neq0$ for every integer $k\ge0$. The law holds at every $k\ge1$ (Corollary~\ref{cor:rays-lowend}(i)), so $\langle1\rangle_{k+6}=-c_6(k)^{-1}\sum_{i<6}c_i(k)\langle1\rangle_{k+i}$ for $k\ge1$, and by induction every value lies in the $\mathbb Q$-span of $\langle1\rangle_1,\dots,\langle1\rangle_6$. The low-end relation of Exercise~\ref{ex:listsb-nlean}(v) has the constant term $64402904777\neq0$ and the coefficient $3400346988339840000\neq0$ at $\langle1\rangle_6$, so it puts $1$ in the value span and $\langle1\rangle_6$ in $\mathbb Q+\mathbb Q\langle1\rangle_1+\dots+\mathbb Q\langle1\rangle_5$. So the value span of the lean is
\[
\mathbb Q+\mathbb Q\langle1\rangle_1+\mathbb Q\langle1\rangle_2+\dots+\mathbb Q\langle1\rangle_5,
\]
of dimension at most $6$. \emph{The centre.} The same with $c_{11}=-4691556000\,(k+4)(k+5)(k+6)(k+7)^2(k+8)^2(k+9)^2(k+10)^2\,S_{41}(k-1)$, nonzero at every integer $k\ge0$, and the low-end relation of Exercise~\ref{ex:listsb-ncentre}, whose constant term and coefficient of $\langle1\rangle_{11}$ are nonzero: the value span of the centre is $\mathbb Q+\mathbb Q\langle1\rangle_1+\dots+\mathbb Q\langle1\rangle_{10}$, of dimension at most $11$.

(ii) \emph{The values.} Input: the laws and the native gauges of the two rays. Method: Theorem~\ref{thm:gaugeend} at $700$ digits, the Beta sums from the strengths $861$ and $886$ with $2580$ terms of the germ, at $20224$ bits, and the laws run back; this gives $\langle1\rangle_1,\dots,\langle1\rangle_7$ at the lean and $\langle1\rangle_1,\dots,\langle1\rangle_{12}$ at the centre. Check: the two starting strengths agree to $700$ digits; the low-end relations, which the reading does not use, hold on the values to $731$ digits relative to their largest terms; and the values agree with those of Exercises~\ref{ex:listsb-nlean} and~\ref{ex:listsb-ncentre}. Receipt: \texttt{valuespan\_N.py}.

\emph{The relations.} Method: for numbers $v_0=1,v_1,\dots,v_N$, the Lenstra--Lenstra--Lov\'asz reduction of the lattice spanned by the rows $\big(\mathbf e_i,\ \lfloor10^{660}v_i\rfloor\big)$; a reduced row $(x,\ast)$ with a small residual $\sum_ix_iv_i$ is a candidate integer relation, and the first rows that are not relations show the size of the generic vectors. Output at the lean, for $1,\langle1\rangle_1,\dots,\langle1\rangle_6$: three vectors of heights $3.0\cdot10^9$, $5.2\cdot10^9$ and $1.4\cdot10^{10}$ with residuals below $10^{-724}$,
\begin{align*}
&15476-3451\langle1\rangle_1+556832\langle1\rangle_2-17388081\langle1\rangle_3\\
&\qquad+111849150\langle1\rangle_4-425617416\langle1\rangle_5-2986275600\langle1\rangle_6,\\
&84581-18857\langle1\rangle_1+3035312\langle1\rangle_2-97811235\langle1\rangle_3\\
&\qquad+798867930\langle1\rangle_4-5212068120\langle1\rangle_5+1459458000\langle1\rangle_6,\\
&5489069-1222045\langle1\rangle_1+190733734\langle1\rangle_2-4628283534\langle1\rangle_3\\
&\qquad+13506396540\langle1\rangle_4-1721929104\langle1\rangle_5-44906400\langle1\rangle_6,
\end{align*}
and then vectors of height $3\cdot10^{157}$ and residual $10^{-503}$. Check: the low-end relation, which is exact, lies in the rational span of the three (exact linear algebra; receipt \texttt{valuespan\_exact.py}). Solved for the last values, the three read
\begin{gather*}
\langle1\rangle_4=\tfrac{556429709}{1602852300}\langle1\rangle_3-\tfrac{137833049}{9617113800}\langle1\rangle_2+\tfrac{26497699}{288513414000}\langle1\rangle_1-\tfrac{170029763}{412162020000},
\end{gather*}
and like expressions of $\langle1\rangle_5$ and $\langle1\rangle_6$ with denominators of up to $20$ digits. Among $1,\langle1\rangle_1,\langle1\rangle_2,\langle1\rangle_3$ and $\ell$ the reduction returns one vector of height $1.1\cdot10^7$ and residual $10^{-727}$, the next of height $10^{163}$:
\[
3906\,\langle1\rangle_1-506772\,\langle1\rangle_2+10837800\,\langle1\rangle_3=15263+1047\,\ell .
\]
$\ell$ is the value of the minor $\binom{\{1,2,3,5\}}2=(\mathrm N_6/4)\setminus6$ of $\mathrm N_6$ at its own full lean $(\one;0,2)$ and at its centre: the procedure reads $\binom{[4]}2$ at both, with laws of order $2$ and degrees $3$ and $3$, and returns $\ell$ to the $32$ digits printed (Exercise~\ref{ex:listsb-table}).

At the centre, for $1,\langle1\rangle_1,\dots,\langle1\rangle_{12}$, four vectors of heights $5.4\cdot10^{23}$ to $1.1\cdot10^{24}$ with residuals below $10^{-714}$, the next of height $2.5\cdot10^{62}$ and residual $10^{-598}$; the law read at $k=1$ and the low-end relation, both exact, lie in their rational span, and they express $\langle1\rangle_9,\dots,\langle1\rangle_{12}$ through $1,\langle1\rangle_1,\dots,\langle1\rangle_8$. Among $1,\langle1\rangle_1,\dots,\langle1\rangle_8$, $\ell$ and $\Tx3{\IZ}0$ the reduction returns two vectors, of heights $1.4\cdot10^{25}$ and $1.6\cdot10^{25}$ and residuals below $10^{-710}$, the next of height $2.9\cdot10^{67}$; their two coefficients at $\ell$ and $\Tx3{\IZ}0$ form the nonsingular matrix $\begin{pmatrix}146909488464924911235&-115139161649314013904\\-10928078403398636616&24473573045860222080\end{pmatrix}$, so they write $\ell$ and $\Tx3{\IZ}0$ through the values. Here $7\Tx3{\IZ}0$ is the value of the minor $\binom{\{1,2,3,6\}}1=\mathrm N_6/\{4,5\}$ at its centre, and $\ell$ that of $\binom{\{1,2,3,5\}}2$.

So the value span of $\mathrm N_6$ is $\mathbb Q+\sum_{i\le5}\mathbb Q\langle1\rangle_i$ at the lean and $\mathbb Q+\sum_{i\le10}\mathbb Q\langle1\rangle_i$ at the centre, of dimensions at most $6$ and $11$; the integer relations that lattice reduction finds at $700$ digits, holding there to $724$ and $714$ digits, cut them to $\mathbb Q+\mathbb Q\langle1\rangle_1+\mathbb Q\langle1\rangle_2+\mathbb Q\langle1\rangle_3$, containing $\ell$, and to $\mathbb Q+\sum_{i\le8}\mathbb Q\langle1\rangle_i$, containing $\ell$ and $\Tx3{\IZ}0$: at the lean the value of the minor $\binom{\{1,2,3,5\}}2$ at its own lean, at the centre the values of both minors on four marks at their centres, and besides them two numbers at the lean and six at the centre.
\end{solution}

\begin{exercise}[$\mathrm V_8$ at its full cobasis lean, read from its list]\label{ex:listsb-vlean}
Let $\mathrm V_8$ be the list of rank four on $\cE=\{1,\dots,8\}$ with the pairs $P_1=\{1,2\}$, $P_2=\{3,4\}$, $P_3=\{5,6\}$, $P_4=\{7,8\}$, whose bases are the sixty-five sets of four marks other than the five unions $P_i\cup P_j$ with $\{i,j\}\neq\{3,4\}$. Read its ray with origin $0$ and direction $(\one;0,2)$, of stride $f=\alpha^{\cE}\Psi^{-2}$, by the procedure of Definition~\ref{def:readprocedure}:
\begin{enumerate}[label=(\roman*)]
\item the walls, $V_1$, the native gauge of the direction, $f^\star$ and $\det H$;
\item the box, its spanning set, the matrix $\mathbf B$ and its common denominator;
\item the relation of the rows and its line-up strengths;
\item the native gauges of the ray, and what they say about the dimension, the masters, the order of the law of $[1]$ and $\mathbf B$ at infinite strength;
\item the value $\langle1\rangle_1=P_{\mathrm V_8}(\one;0,2)$.
\end{enumerate}
\end{exercise}

\begin{solution}
(i) \emph{Walls and volume.} As for $\mathrm N_6$ at its lean (Exercise~\ref{ex:listsb-nlean}(i)), $s=0$, $t=2$ and $d^\ast(S)=|S|-r(S)$ give $\kappa_{1,S}=2r(S)-|S|$. Input: the $254$ nonempty proper subsets of $\cE$ and their ranks. Method: $\kappa_{1,S}$ for each, and the recursion over chains $F(U)=\kappa_{1,U}^{-1}\sum_{e\in U}F(U\setminus\{e\})$, $V_1=\sum_eF(\cE\setminus\{e\})$ of Exercise~\ref{ex:listsb-nlean}(i) (Proposition~\ref{prop:principal}), in exact rationals. Output: the walls are $1$ sixteen times, $2$ sixty-one times, $3$ one hundred and twelve times and $4$ sixty-five times, all positive, and
\[
V_1=300 .
\]
Check: the same program gives $V_1=66$ and $63$ for the two rays of $\mathrm N_6$, as computed by hand in Exercises~\ref{ex:listsb-nlean} and~\ref{ex:listsb-ncentre}. Receipt: \texttt{exact\_data.py}.

\emph{The automorphisms.} The non-bases are the unions of two pairs that involve $P_1$ or $P_2$, so an automorphism keeps $\{P_1,P_2\}$ and $\{P_3,P_4\}$; the $64$ permutations that exchange the marks within pairs, $P_1$ with $P_2$ and $P_3$ with $P_4$ are automorphisms, and they keep the masses $\one$, so they are symmetries of the ray (Proposition~\ref{prop:raysym}(ii)).

\emph{The native gauge.} The points fixed by the $64$ automorphisms are, on the chart $\alpha_8=1$, the slice $\alpha=(x,x,x,x,1,1,1,1)$. A cobasis is the complement of a basis, so the cobases are the sets of four other than the complements $P_3\cup P_4$, $P_2\cup P_4$, $P_2\cup P_3$, $P_1\cup P_4$, $P_1\cup P_3$ of the non-bases. On the slice a set $C$ of four has $\alpha^C=x^{|C\cap(P_1\cup P_2)|}$, the sets of four give $\sum_j\binom4j\binom4{4-j}x^j=1+16x+36x^2+16x^3+x^4$, and the five sets left out give $1+4x^2$, so
\begin{gather*}
\Psi=x\,(x^3+16x^2+32x+16),\qquad f=\frac{x^4}{\Psi^2}=\frac{x^2}{(x^3+16x^2+32x+16)^2},\\
\frac d{dx}\log f=\frac{-4\,(x^3+8x^2-8)}{x\,(x^3+16x^2+32x+16)} .
\end{gather*}
At a point fixed by the automorphisms the differential of $\log f$ in $y=\log\alpha$ is $(u,u,u,u,v,v,v,v)$ with $u+v=0$, $f$ being homogeneous of degree $0$, and it vanishes when its derivative along the slice does, as in Exercise~\ref{ex:listsb-nlean}(ii). So the three roots of $x^3+8x^2-8$ give native gauges (Definition~\ref{def:raygauge}); at a root $x^3+16x^2+32x+16=8(x+1)(x+3)$, and the native value is
\[
f=\Big(\frac x{8(x+1)(x+3)}\Big)^2 .
\]
The positive root $x^\star=0.94566781799051109644\ldots$ gives $\alpha^\star$ (Theorem~\ref{thm:pencil}(i)) and
\[
f^\star=2.3709304548895617852\ldots\cdot10^{-4};
\]
the roots $-1.0748031540504515212\ldots$ and $-7.8708646639400595752\ldots$ give the native values $0.87033768220817625690\ldots$ and $8.6423132047649619559\ldots\cdot10^{-4}$. The Hessian of $-\log f$ at $\log\alpha^\star$ on the chart is twice the covariances of the indicators under the cobasis law at $\alpha^\star$ (Theorem~\ref{thm:branch}); reduced in $\mathbb Q(x^\star)$ (receipt \texttt{hessians.py}),
\[
\det H=\frac{395719931704\,x^{\star2}+2812581517252\,x^\star-3012741822955}{379727508532},
\]
which is $0.0024055146373001934835\ldots$

(ii) \emph{The box} (Procedure~\ref{def:readprocedure}, the masters). Only $\Psi$ is active. Input: the boxes $\mathfrak B_{\mathfrak t}(0,1)$ modulo the $64$ automorphisms (Definition~\ref{def:raybox}(iv), Definition~\ref{def:box-reduction}(i)). Method: as in Exercise~\ref{ex:listsb-nlean}(iii). Output: $\mathfrak B_2(0,1)$ has $1155$ orbits of stems, $4745$ orbits of symbols and $10395$ relations; at $\mathfrak s=1$ it leaves $36$ free Laurent orbits of which $33$ targets do not reduce to them, and at $\mathfrak s=0$ the target of $1$ does not reduce: not accepted. $\mathfrak B_3(0,1)$ has $9723$ orbits of stems, $32810$ orbits of symbols and $87507=9723\cdot9$ relations, the letter identities of the stems and their eight mark relations; one elimination leaves $60$ free Laurent orbits of wall degree at most $2$, every target reduces to them, and the box is accepted at $\mathfrak s=2$ (Theorem~\ref{thm:raybox}(i)), with a spanning set of $60$ orbit averages. Receipt: \texttt{rayreader.py V "1 1 1 1 1 1 1 1 ; 0 2" --digits 30} (log \texttt{runlogs/V\_lean.stdout}).

\emph{The matrix} (the matrix over $\mathbb Q(k)$). Input: the $87507$ relations. Method: the targets need $19581$ of them, of rank $19581$; eliminations modulo primes of $62$ bits at random strengths; every entry read as a quotient of polynomials; Chinese remaindering with rational reconstruction. Output: on the first prime $178$ strengths determine every entry, the largest sum of the degrees of a numerator and a denominator being $163$; all the denominators divide one polynomial of degree $85$, over which the numerators have degree at most $88$, so each further prime needs $94$ strengths; the entries become reconstructible over $\mathbb Q$ at seven primes, $432$ bits. The common denominator is, up to a constant factor,
\[
Q_c=k\,(k-1)^2(2k-3)(k-2)^8(2k-5)(k-3)^7\,S_{65}(k),
\]
$S_{65}$ irreducible of degree $65$; $72$ entries vanish, and in lowest terms the numerators have degree at most $83$ and the denominators at most $80$, so $d_{\mathbf B}=83$. Check: $\mathbf B$ agrees with an elimination of the whole box at a prime and a strength drawn at random.

(iii) \emph{The relation of the rows.} Input: $\mathbf B$. Method: the rows $w_i$ of Proposition~\ref{prop:rays-lawrows} modulo primes, their first dependence, reconstructed over $\mathbb Q$. Output (receipt as in (ii)): $w_0,\dots,w_{23}$ are independent and $w_{24}$ depends on them, so $h=24$, and the relation $\sum_{j=0}^{24}p_j(k)\langle1\rangle_{k+j}=0$ has $p_j\in\mathbb Z[k]$ of degree $377$, with coefficients of up to $753$ digits, from $63$ primes, and over $\mathbb Q$
\[
p_0=k^7(k+1)^5(k+2)^5\prod_{j=3}^7(k+j)^3\prod_{j=8}^{13}(k+j)\,S_{339}(k),
\]
$S_{339}$ irreducible of degree $339$, while $p_{24}$ is a product of $38$ linear factors, with roots at integers, halves and quarters from $-\frac12$ to $-\frac{47}2$, and $S_{339}(k-1)$. Check: the relation holds at a prime and a point drawn at random. \emph{Line-up strengths.} With $q_{24}(k)=p_0(k-24)$, Proposition~\ref{prop:rays-lawrows}(iii) gives
\[
p_0(k-24)\ \Big|\ \det\mathbf W(k)_J\prod_{j=0}^{23}Q_c(k-j)^{24-j}
\]
for every set $J$ of $24$ columns with $\det\mathbf W_J\neq0$. The irreducible factors of $Q_c$ have degree at most $65$, so the irreducible $S_{339}(k-24)$ divides none of the $Q_c(k-j)$ and divides $\det\mathbf W(k)_J$: $339$ of the $377$ degrees of $p_0$ are strengths at which the rows $w_0,\dots,w_{23}$ line up, and the other $38$ are the integers $0,\dots,-13$.

(iv) \emph{The native gauges.} A native gauge of the ray is a point of $X_F$ with $1=2R^\ast_e(\alpha)$ for every mark (Definition~\ref{def:raygauge}). Input: these equations, and the same with exponents $(\mathbf a;b)$ in place of $(\one;2)$, $|\mathbf a|=4b$, which are linear in the exponents. Method: as in Exercise~\ref{ex:listsb-sizing}(iii): Newton's method at random complex exponents, the solutions carried around random triangles of exponents until several turns in a row add nothing, run from two seeds; then every solution carried to $(\one;2)$ along four routes, the straight segment and three broken ones through random exponents. Output: from both seeds, $839$ solutions at random exponents, the count reaching $839$ at the sixteenth and seventeenth turns and then staying there for four and six turns; at the direction, along every route, the same $331$ nondegenerate native gauges, in $31$ orbits of the $64$ automorphisms, of sizes $1,1,1,2,2,2,2$, eight of $4$, fourteen of $16$ and two of $32$, with $18$ distinct native values:
\[
\begin{array}{rl}
\text{one orbit:}&f^\star,\ 0.000818456\ (16),\ 0.000864231\ (1),\ 0.124846403\ (16),\ 0.870337682\ (1),\\
&-0.592832429\pm0.713227829\,i\ (16\text{ each}),\\[2pt]
\text{two orbits:}&0.004487298\ (4),\ -0.013533898\ (8),\ -0.026238829\ (8),\ 0.066968510\ (32),\\
&0.007758455\pm0.013607932\,i\ (32\text{ each}),\ 0.109146811\ (32),\ \tfrac18\ (64),\\
&0.870512702\ (4),\\[2pt]
\text{three orbits:}&0.021107646\ (24),\ 0.416392354\ (24),
\end{array}
\]
the number of native gauges in brackets (renderings). The other $508$ paths stop before the direction, at $\tau$ between $0.598$ and $1-1.8\cdot10^{-6}$ along the segment, where the step collapses. Check: one native gauge of each orbit polished by Newton's method at $40$ digits, with the Hessian determinants of modulus at least $\det H=0.0024055\ldots$, reached at $\alpha^\star$; the three native gauges of the slice in (i) are the three orbits of size $1$. Receipts: \texttt{critpoints.py} (\texttt{crit\_V8\_lean.json}, \texttt{crit\_V8\_lean\_seed2.json}), \texttt{census\_paths.py} (\texttt{census\_V8\_routes.json}), \texttt{escape\_check.py} (\texttt{escape\_V8.out}), \texttt{betas.py} (\texttt{betas\_V8.out}).

So, by Theorem~\ref{thm:gaugesolution}(iii), the ray has dimension at least $331$; by Theorem~\ref{thm:invariant}(iii) the invariant module under the $64$ automorphisms has dimension at least $31$; and by Corollary~\ref{cor:symbol} the law of $[1]$ has order at least $18$, so that with (iii) and Proposition~\ref{prop:rays-lawrows}(ii) it lies between $18$ and $24$. The vectors $\beta_c=(b_l(\alpha_c))_l$ of the spanning set at one native gauge of each orbit have rank $31$, their singular values running from $1262$ down to $0.036$ (receipt \texttt{betas.py}): they do not span $\mathbb C^{60}$, and Proposition~\ref{prop:rays-spectrum}(i) does not decide whether the $60$ are masters. At large strength the entries of $\mathbf B$ grow like $k^3$ (output, receipt as in (ii)): $\mathbf B$ is not regular at $k=\infty$, and Proposition~\ref{prop:rays-spectrum}(ii) does not apply to it.

(v) \emph{The value} (Theorem~\ref{thm:gaugeend}). Input: the relation of (iii), $f^\star$ and $\det H$; $f^\star$ is a simple root of the symbol, and the exponent of the density at $f^\star$ is $\frac{n-3}2=\frac52$. Method: as in Exercise~\ref{ex:listsb-nlean}(vi). Output (receipt as in (ii)):
\[
P_{\mathrm V_8}(\one;0,2)=\langle1\rangle_1=7.12823937562968675472229135353893457156078898940720766490704\ldots,
\]
two starting strengths agreeing to $2.4\cdot10^{-82}$, and $P/V_1=0.0237607979187656\ldots$ Check: the low-end relation of Corollary~\ref{cor:rays-lowend}(ii), with $V_1=300$, holds on the values to $73$ digits.

So $\mathrm V_8$ at $(\one;0,2)$ has $V_1=300$, $f^\star=\big(x^\star/(8(x^\star+1)(x^\star+3))\big)^2$ with $x^{\star3}+8x^{\star2}=8$, a box accepted at $\mathfrak B_3(0,1)$ with $60$ spanning elements, a relation of order $24$ and degree $377$ of which $339$ degrees are line-up strengths, at least $331$ native gauges in $31$ orbits with $18$ native values, and the value $7.1282393756296867547\ldots$
\end{solution}

\begin{exercise}[Sizing a reading: $\mathrm Q_6$ and $\mathrm N_7$]\label{ex:listsb-sizing}
Let $\mathrm Q_6$ be the list of rank three on $\{1,\dots,6\}$ whose bases are the seventeen triples other than $\{1,3,6\}$, $\{2,3,4\}$ and $\{2,5,6\}$, and $\mathrm N_7$ the list of rank three on $\{1,\dots,7\}$ whose bases are the thirty-four triples other than $\{1,2,5\}$. For the rays with origin $0$ of $\mathrm Q_6$ along $(1,3,2,3,1,2;2,2)$ and of $\mathrm N_7$ along $(1,2,2,1,2,1,1;2,1)$, find what sizes their reading before any relation is written: the walls, $V_1$, the native gauge of the direction, the native gauges of the ray and their native values, with the bounds these give on the dimension of the ray and on the order of a law of $[1]$; then the size of the box with one unit of each letter at $\mathfrak t=2$, and whether it is accepted.
\end{exercise}

\begin{solution}
(i) \emph{Walls and volume} (Definition~\ref{def:family}, Proposition~\ref{prop:principal}). For $\mathrm Q_6$, $s=t=2$ and the masses add up to $12=2\cdot3+2\cdot3$, on the balance wall; the walls $\kappa_{1,S}=m_{1,S}-2d(S)-2d^\ast(S)$ of the $62$ sets $S$ are, by height and count,
\[
1(4),\ 2(7),\ 3(13),\ 4(13),\ 5(12),\ 6(9),\ 7(3),\ 8(1),
\]
all positive, and the recursion over chains of Exercise~\ref{ex:listsb-nlean}(i), run over the $62$ sets (no automorphism of $\mathrm Q_6$ fixes the masses), gives $V_1=\frac{16493}{2430}$. For $\mathrm N_7$, $s=2$, $t=1$ and $10=2\cdot3+1\cdot4$; the $126$ walls are $1(7),\ 2(17),\ 3(29),\ 4(39),\ 5(29),\ 6(5)$, and $V_1=\frac{3519}{128}$. The automorphisms of $\mathrm N_7$ fixing the masses are the $12$ permutations that exchange the marks $2$ and $5$ of the triple $\{1,2,5\}$ or not and permute the marks $4,6,7$ of mass $1$. Receipt: \texttt{exact\_data.py}.

(ii) \emph{The native gauge of the direction} (Theorem~\ref{thm:pencil}(i)), by Newton's method on the chart: for $\mathrm Q_6$, $\alpha^\star=(0.45796,\,3.83890,\,1.55835,\,3.85580,\,0.54507,\,1)$, $f^\star=8.3244647577707389\ldots\cdot10^{-5}$ and $\det H=0.13385450029314584\ldots$; for $\mathrm N_7$, $\alpha^\star=(1.15660,\,3.63195,\,3.14019,\,1,\,3.63195,\,1,\,1)$, constant on the orbits of the automorphisms as it must be, $f^\star=7.2002200376070769\ldots\cdot10^{-5}$ and $\det H=0.027073861743142843\ldots$ (renderings; receipt \texttt{hessians.py}).

(iii) \emph{The native gauges.} A native gauge is a point of $X_F$ at which $m_{1,e}=s\,R_e+t\,R^\ast_e$ for every mark (Definition~\ref{def:raygauge}). Input: these equations, and the same with the exponents $(\mathbf a;b_K,b_\Psi)$ in place of $(\mvec_1;s,t)$, $|\mathbf a|=3b_K+(n-3)b_\Psi$, which are linear in the exponents. Method: for random exponents, a Groebner basis of the equations, multiplied by $K\Psi$, on the chart and saturated by the coordinates and the letters, modulo the prime $32003$, and, separately, Newton's method at random complex exponents with the solutions carried around random triangles of exponents until several turns add nothing; then each solution carried along the segment to $(\mvec_1;s,t)$, a path that either arrives at a point of $X_F$ or runs to its boundary; and a Groebner basis at $(\mvec_1;s,t)$ itself. Output: for $\mathrm Q_6$, $114$ solutions at random exponents by both methods; $90$ of the paths arrive, at $90$ distinct native gauges, all nondegenerate, with $90$ distinct native values, and $24$ run off; at the direction the Groebner basis has exactly $90$ points counted with multiplicity, so these are all the native gauges. For $\mathrm N_7$, $1174$ solutions at random exponents; all $1174$ arrive, at distinct nondegenerate native gauges, in $268$ orbits of the $12$ automorphisms ($32$ of size $1$, $28$ of size $2$, $86$ of size $3$, $106$ of size $6$, $16$ of size $12$), with $83$ distinct native values. Check: the native gauges of $\mathrm Q_6$ polished at $40$ digits, residuals below $10^{-37}$, the Hessian determinants at least $0.13$, the native values at least $1.6\cdot10^{-4}$ apart, and the same $90$ values from two further runs with other random exponents. Receipts: \texttt{gb\_count.py}, \texttt{gb\_direction.py}, \texttt{critpoints.py}, \texttt{polish.py}.

So, by Theorem~\ref{thm:gaugesolution}(iii), the dimension of the ray is at least $90$ for $\mathrm Q_6$ and at least $1174$ for $\mathrm N_7$; by Theorem~\ref{thm:invariant}(iii) the invariant module of $\mathrm N_7$ under its $12$ automorphisms has dimension at least $268$; and by Corollary~\ref{cor:symbol} every relation $\sum_jp_j(k)f^j\in\mathcal L$, the law of $[1]$ among them, has order at least $90$ for $\mathrm Q_6$ and $83$ for $\mathrm N_7$, the number of distinct native values. These counts are read off the list and the direction alone.

(iv) \emph{The box of $\mathrm Q_6$.} Both letters are active, and the box with one unit of each is $\mathfrak B_2(1,1)$ (Definition~\ref{def:raybox}(iv)). Input: the box, with no symmetry. Output: $361682$ relations on $225874$ symbols, the mark relations of $31785$ stems and the letter identities, each of $1+|\cB|=18$ terms (receipt \texttt{box\_test.py Q6.json "1 3 2 3 1 2 ; 2 2" 2 1 1}). Method for the rest: the procedure, as in Exercise~\ref{ex:listsb-nlean}(iii) and (iv); receipt \texttt{rayreader.py Q6.json "1 3 2 3 1 2 ; 2 2" --digits 30}. Output: one elimination modulo a prime at a random strength leaves $114$ free Laurent symbols of wall degree at most $1$, and every target reduces to them, so the box is accepted at $\mathfrak s=1$, and the spanning set of $114$ spans a submodule that contains $[1]$ (Theorem~\ref{thm:raybox}(i)). The targets need $78921$ of the relations. Modulo the first prime $\mathbf B$ has a common denominator of degree $126$ and numerators of degree up to $129$, so it is not regular at $k=\infty$; it needs $271$ strengths on that prime and $135$ on each further one; and modulo that prime the rows first depend at $h=90$, where $12288$ consecutive strengths do not determine the relation, so that its degree exceeds $6100$. By (iii) and Proposition~\ref{prop:rays-lawrows}(ii), $90\le h\le114$ whatever the spanning set, and the rows meet the lower bound.

\emph{The box of $\mathrm N_7$.} The box $\mathfrak B_2(1,1)$ modulo the $12$ automorphisms has $71320$ orbits of symbols, $7613$ orbits of stems and $100719$ relations of up to $1+|\cB|=35$ terms. One elimination modulo $2^{61}-1$ at a random strength has rank $63883$; of the $359$ Laurent orbits of wall degree at most $1$, $247$ are free, and their targets all lie in the box, but $246$ of them do not reduce to the $247$; at $\mathfrak s=0$ the target of $1$ does not reduce. So the box with one unit of each letter at $\mathfrak t=2$ is not accepted (Definition~\ref{def:box-reduction}(iii)), and Theorem~\ref{thm:rayboxstop} sends the search to larger boxes. By Definition~\ref{def:box-reduction}(ii) the class of every invariant Laurent symbol of wall degree at most $1$ is a combination of the $247$ free ones, fewer than the $268$ dimensions of the invariant module found in (iii). Receipt: \texttt{box\_test.py N7.json "1 2 2 1 2 1 1 ; 2 1" 2 1 1}. The next boxes of the search are larger: $\mathfrak B_2(2,1)$ has $236600$ relations on $158724$ symbols, and $\mathfrak B_3(1,1)$ has $707989$ relations of up to $35$ terms (receipt for the search, \texttt{rayreader.py N7.json "1 2 2 1 2 1 1 ; 2 1" --digits 30}).

So before any relation is written, the walls give $V_1=\frac{16493}{2430}$ and $\frac{3519}{128}$, and the critical equations give at least $90$ and $1174$ dimensions and laws of order at least $90$ and $83$; the box of $\mathrm Q_6$ at $\mathfrak t=2$ is accepted with $114$ spanning elements, and the box of $\mathrm N_7$ with one unit of each letter at $\mathfrak t=2$ is not accepted.
\end{solution}

\begin{exercise}[A table of rays read by the procedure]\label{ex:listsb-table}
Let $\binom{[n]}1$ be the list of rank one on $\{1,\dots,n\}$ whose bases are the single marks, $\binom{[n]}2$ the list of rank two whose bases are all the pairs, $\mathrm H_4$ the list of rank two on $\{1,2,3,4\}$ whose bases are the five pairs other than $\{1,2\}$, and $W_3$ the wheel with three spokes (Proposition~\ref{prop:wheelholonomy}), of rank three on its six spokes and rims, whose bases are the sixteen spanning trees of its graph. Read by the procedure of Definition~\ref{def:readprocedure} the rays with origin $0$ of
\begin{gather*}
\tbinom{[3]}1,\ \tbinom{[4]}1,\ \tbinom{[5]}1,\ \tbinom{[4]}2,\ \tbinom{[5]}2,\ \mathrm H_4,\ \mathrm N_6\ \text{along }(\one;1,1);\qquad \tbinom{[4]}2,\ W_3,\ \mathrm N_6,\ \mathrm V_8\ \text{along }(\one;0,2);\\
\tbinom{[3]}1\ \text{along }(2\one;4,1);\qquad \mathrm H_4\ \text{along }(1,1,3,3;3,1),
\end{gather*}
$\mathrm N_6$ and $\mathrm V_8$ as in Exercises~\ref{ex:listsb-nletters} and~\ref{ex:listsb-vlean}, and the ray of $\binom{[3]}1$ with origin $(u_0;0,0)$, $u_0=(-\frac12,0,\frac12)$, and direction $(\one;1,1)$, whose value at $k=1$ is $P_{\binom{[3]}1}(\frac12,1,\frac32;1,1)$. For each ray give:
\begin{enumerate}[label=(\roman*)]
\item the accepted box with its stems, symbols and relations modulo the symmetries of the ray, the bound $\mathfrak s$ at which it is accepted, the size $\mathfrak m$ of its spanning set, the largest degree $d_{\mathbf B}$ of a numerator or a denominator of an entry of $\mathbf B$ in lowest terms, and the order $h$ and the degree $D$ of the relation of the rows;
\item whether the spanning set is a set of masters and that relation the law of $[1]$, and the spectrum of $\mathbf B(\infty)$ where $\mathbf B$ is regular at $k=\infty$;
\item $V_1$ and the value $\langle1\rangle_1$, with the checks that the reading allows.
\end{enumerate}
\end{exercise}

\begin{solution}
(i) \emph{The boxes, the matrices and the rows.} Input: the list and the direction. Method: the procedure as in Exercise~\ref{ex:listsb-nlean}(iii) and (iv). The boxes $\mathfrak B_{\mathfrak t}$ with one unit of each active letter are tried for $\mathfrak t=2,3,\dots$ (Definition~\ref{def:raybox}(iv), Theorem~\ref{thm:rayboxstop}), each reduced modulo the symmetries of the ray: the automorphisms that fix the masses and the origin, and the Cremona at a centre with origin $0$ (Proposition~\ref{prop:raysym}, Definition~\ref{def:box-reduction}(i)). One elimination modulo a prime of $61$ bits at a random strength decides acceptance (Definition~\ref{def:box-reduction}(iii)); $\mathbf B$ over $\mathbb Q(k)$ comes from eliminations of the relations that the targets need, modulo primes of $62$ bits at random strengths, with rational reconstruction; and the rows $w_i$ of Proposition~\ref{prop:rays-lawrows} give the first dependence, reconstructed over $\mathbb Q$ from several primes. Output:
\begin{center}
\small
\begin{tabular}{@{}l c r r r c r r c@{}}
\toprule
ray&box&stems&symbols&relations&$\mathfrak s$&$\mathfrak m$&$d_{\mathbf B}$&$(h,D)$\\\midrule
$\binom{[3]}1$, $(\one;1,1)$&$\mathfrak B_2(1,1)$&$4$&$30$&$32$&$1$&$2$&$2$&$(2,2)$\\
$\binom{[4]}1$, $(\one;1,1)$&$\mathfrak B_2(1,1)$&$5$&$42$&$46$&$1$&$2$&$3$&$(2,3)$\\
$\binom{[5]}1$, $(\one;1,1)$&$\mathfrak B_3(1,1)$&$10$&$100$&$112$&$2$&$3$&$5$&$(3,4)$\\
$\binom{[3]}1$, $(2\one;4,1)$&$\mathfrak B_2(1,1)$&$17$&$79$&$124$&$1$&$2$&$4$&$(2,6)$\\
$\binom{[3]}1$, $(\one;1,1)$, $u_0$&$\mathfrak B_2(1,1)$&$19$&$110$&$149$&$1$&$4$&$3$&$(3,4)$\\
$\binom{[4]}2$, $(\one;0,2)$&$\mathfrak B_2(0,1)$&$8$&$20$&$40$&$1$&$2$&$4$&$(2,3)$\\
$\binom{[4]}2$, $(\one;1,1)$&$\mathfrak B_2(1,1)$&$7$&$33$&$66$&$1$&$2$&$4$&$(2,3)$\\
$\binom{[5]}2$, $(\one;1,1)$&$\mathfrak B_3(1,1)$&$21$&$234$&$245$&$2$&$5$&$8$&$(4,6)$\\
$\mathrm H_4$, $(\one;1,1)$&$\mathfrak B_2(1,1)$&$17$&$102$&$180$&$1$&$4$&$3$&$(4,7)$\\
$\mathrm H_4$, $(1,1,3,3;3,1)$&$\mathfrak B_2(1,1)$&$115$&$631$&$1000$&$1$&$4$&$9$&$(4,25)$\\
$W_3$, $(\one;0,2)$&$\mathfrak B_2(0,1)$&$66$&$211$&$462$&$1$&$2$&$8$&$(2,5)$\\
$\mathrm N_6$, $(\one;0,2)$&$\mathfrak B_2(0,1)$&$104$&$329$&$728$&$1$&$7$&$6$&$(6,24)$\\
$\mathrm N_6$, $(\one;1,1)$&$\mathfrak B_3(1,1)$&$252$&$1909$&$3774$&$2$&$15$&$17$&$(11,52)$\\
$\mathrm V_8$, $(\one;0,2)$&$\mathfrak B_3(0,1)$&$9723$&$32810$&$87507$&$2$&$60$&$83$&$(24,377)$\\
\bottomrule
\end{tabular}
\end{center}
At the centres of $\binom{[5]}1$, $\binom{[5]}2$ and $\mathrm N_6$ the box $\mathfrak B_2(1,1)$ is not accepted: at $\mathfrak s=1$ it leaves $2$, $3$ and $8$ free Laurent orbits of which $1$, $2$ and $7$ targets do not reduce to them, and at $\mathfrak s=0$ the target of $1$ does not reduce; at the lean of $\mathrm V_8$, $\mathfrak B_2(0,1)$ leaves $36$ free orbits at $\mathfrak s=1$, of which $33$ targets do not reduce, and fails at $\mathfrak s=0$ the same way. Every other box in the table is the first one tried. Check: every $\mathbf B$ agrees with an elimination of the whole box at a prime and a strength drawn at random, and every relation of the rows holds at a random prime and point; on the thirteen rays other than $\mathrm V_8$ both checks are repeated at three fresh primes and points. Receipts: \texttt{runs\_table.sh}, \texttt{table\_summary.py}, \texttt{tests/check\_run.py} (\texttt{table\_checkrun.out}).

(ii) \emph{Masters and the law of $[1]$.} By Proposition~\ref{prop:rays-lawrows}(ii), with Theorem~\ref{thm:raybox}(iii), $h$ is at least the dimension of the cyclic submodule of $[1]$, with equality when the spanning set is a set of masters, and then the relation of the rows is the law of $[1]$; by Corollary~\ref{cor:symbol} that dimension is at least the number of distinct native values of nondegenerate native gauges. So the relation is the law of $[1]$ as soon as either the vectors $\beta_c=(b_l(\alpha_c))_l$ at nondegenerate native gauges span $\mathbb C^{\mathfrak m}$ (Proposition~\ref{prop:rays-spectrum}(i)), or $h$ distinct native values are found; in the second case, when moreover $h=\mathfrak m$, the cyclic submodule of $[1]$ fills the span of the $\mathfrak m$ classes, and the spanning set is a set of masters. Input: the critical equations of each direction. Method: the census of Exercise~\ref{ex:listsb-sizing}(iii), monodromy at random exponents and paths to the direction; each native gauge polished at $40$ digits with its Hessian; and the vectors $\beta_c$ at one native gauge of each orbit, each $b_l$ averaged over the symmetries of the ray. Output and conclusion:
\begin{itemize}
\item $\binom{[3]}1$, $\binom{[4]}1$ at their centres: the table's gauge, of native value $\frac19$ and $\frac1{16}$, and the signs with one scale $-1$, of native value $1$ and $\frac14$ (Theorem~\ref{thm:signgauges}(ii)), $4$ and $5$ native gauges in all, every one nondegenerate; the two vectors $\beta_c$ are independent: masters, and the law of $[1]$.
\item $\binom{[3]}1$ along $(2\one;4,1)$: $4$ native gauges, of native values $\frac1{243}$ and $-\frac{27}{64}$, two independent vectors: masters. The ray with origin $u_0$: the $4$ native gauges of the centre direction, which no symmetry of the ray permutes, give four independent vectors: masters.
\item $\binom{[5]}2$ at its centre: $66$ native gauges in $6$ orbits of the $120$ automorphisms, of native values $\frac1{100},\frac14,-\frac18,1$, all nondegenerate, with vectors of rank $5$: masters, and $h=4<\mathfrak m$, the cyclic submodule of $[1]$ of dimension $4$ inside $\mathfrak N$ of dimension $5$ (Exercise~\ref{ex:closingssee}).
\item $\binom{[5]}1$ at its centre: $16$ native gauges, the table's gauge ($\frac1{25}$), the five signs with one scale $-1$ ($\frac19$) and the ten with two ($1$); three distinct native values and $h=\mathfrak m=3$: masters and the law of $[1]$, although the three vectors $\beta_c$ lie in a plane. $\binom{[4]}2$ at both directions: $4$ native gauges, of native values $\frac1{36}$ (the table's gauge) and $\frac14$ (an orbit of three), with $h=\mathfrak m=2$, and $W_3$ at $(\one;0,2)$: $3$ native gauges, of native values $\frac1{256}$ and $\frac14$, with $h=\mathfrak m=2$: masters and the law of $[1]$, the two vectors $\beta_c$ being proportional.
\item $\mathrm N_6$: Exercises~\ref{ex:listsb-nlean} and~\ref{ex:listsb-ncentre}. $\mathrm V_8$: Exercise~\ref{ex:listsb-vlean}.
\item $\mathrm H_4$: of the $6$ critical points at random exponents, two arrive at nondegenerate native gauges at each direction. At the centre they are the table's gauge and the sign $(-1,-1,1,1)$, of native values $\frac1{25}$ and $\frac19$ and Hessian determinants $-\frac{32}{625}$ and $\frac{32}{81}$ on the chart; along $(1,1,3,3;3,1)$ they are the points $(x,x,1,1)$ with $3x^2+8x-1=0$, of native values $\frac{7643\pm1064\sqrt{19}}{1366875}$, the larger one, at $x=\frac{\sqrt{19}-4}3>0$, being $f^\star$. The other four paths meet in pairs at the points $(\mp i,\pm i,-1,1)$ at the centre and $(\mp\frac i{\sqrt3},\pm\frac i{\sqrt3},-1,1)$ along $(1,1,3,3;3,1)$, native gauges of native value $1$ whose Hessian is singular (exact; receipt \texttt{h4\_gauges.py}). So Corollary~\ref{cor:symbol} bounds the dimension of the cyclic submodule of $[1]$ below by $2$ only, and the two vectors $\beta_c$ span a plane: the relation of the rows is a relation of order $4$ in $\mathcal L$, of order at least that of the law of $[1]$ (Theorem~\ref{thm:raybox}(iii)), which lies between $2$ and $4$.
\end{itemize}

\emph{Infinite strength.} $\mathbf B$ is regular at $k=\infty$ on nine of the rays, with the spectra in the table of (iii), and at every nondegenerate native gauge of the census $\beta_c$ is an eigenvector of $\mathbf B(\infty)$ for $f(\alpha_c)^{-1}$, to $10^{-38}$ at $40$ digits (Proposition~\ref{prop:rays-spectrum}(ii)). On $\binom{[5]}1$, $\binom{[4]}2$ and $W_3$ some entries of $\mathbf B$ have numerators of larger degree than their denominators, and Proposition~\ref{prop:rays-spectrum}(ii) says why they must: nondegenerate native gauges of distinct native values have dependent vectors $\beta_c$ there, three in a plane for $\binom{[5]}1$ and two on a line for the others, which a regular $\mathbf B$ would make independent eigenvectors. For $\mathrm V_8$ see Exercise~\ref{ex:listsb-vlean}.

(iii) \emph{The values.} $V_1$ is computed from the walls by Proposition~\ref{prop:principal} with the recursion over chains of Exercise~\ref{ex:listsb-nlean}(i) (receipt \texttt{exact\_data.py}); for the ray with origin $u_0$ it is that of the direction $(\one;1,1)$. The values come from Theorem~\ref{thm:gaugeend}. Input: the relation of the rows, $f^\star$ and $\det H$. Method: as in Exercise~\ref{ex:listsb-nlean}(vi), asked to $60$ digits. Output:
\begin{center}
\small
\begin{tabular}{@{}l c l l@{}}
\toprule
ray&$V_1$&eigenvalues of $\mathbf B(\infty)$&$\langle1\rangle_1$\\\midrule
$\binom{[3]}1$, $(\one;1,1)$&$6$&$9,\ 1$&$2.3439072386894588906\ldots$\\
$\binom{[4]}1$, $(\one;1,1)$&$24$&$16,\ 4$&$8.4143983221171599977\ldots$\\
$\binom{[5]}1$, $(\one;1,1)$&$120$&not regular&$39.945588011996245698\ldots$\\
$\binom{[3]}1$, $(2\one;4,1)$&$1$&$243,\ -\frac{64}{27}$&$0.030080945282262304399\ldots$\\
$\binom{[3]}1$, $(\one;1,1)$, $u_0$&$6$&$9,\ 1,\ 1,\ 1$&$4.7773947923309302997\ldots$\\
$\binom{[4]}2$, $(\one;0,2)$&$12$&not regular&$2.3439072386894588906\ldots$\\
$\binom{[4]}2$, $(\one;1,1)$&$12$&not regular&$2.3439072386894588906\ldots$\\
$\binom{[5]}2$, $(\one;1,1)$&$30$&$100,\ 4,\ 4,\ 1,\ -8$&$4.0297141707110496073\ldots$\\
$\mathrm H_4$, $(\one;1,1)$&$16$&$25,\ 9,\ 1,\ 1$&$4.0962369983378828500\ldots$\\
$\mathrm H_4$, $(1,1,3,3;3,1)$&$\frac{160}{27}$&$\frac{7643\pm1064\sqrt{19}}{27},\ 1,\ 1$&$0.48830319489526888388\ldots$\\
$W_3$, $(\one;0,2)$&$84$&not regular&$7.2123414189575657123\ldots$\\
$\mathrm N_6$, $(\one;0,2)$&$66$&Exercise~\ref{ex:listsb-nlean}&$4.7568025637582294434\ldots$\\
$\mathrm N_6$, $(\one;1,1)$&$63$&Exercise~\ref{ex:listsb-ncentre}&$4.5789124788824048063\ldots$\\
$\mathrm V_8$, $(\one;0,2)$&$300$&not regular&$7.1282393756296867547\ldots$\\
\bottomrule
\end{tabular}
\end{center}
Check: on every ray the two starting strengths agree to at least $82$ digits. On the thirteen rays with origin $0$ the low-end relation of Corollary~\ref{cor:rays-lowend}(ii), $[k^{n-1}]p_0\cdot V_1+\sum_{i=1}^hp_i(0)\langle1\rangle_i=0$, which the reading does not use and which holds for the relation of the rows as well as for the law of $[1]$, holds on the values $\langle1\rangle_1,\dots,\langle1\rangle_h$ to at least $89$ digits relative to its largest term on the twelve read here at $60$ digits, and to $73$ digits on $\mathrm V_8$ (Exercise~\ref{ex:listsb-vlean}(v)). Five values agree with numbers computed otherwise: those of $\binom{[3]}1$ at its centre and of $\binom{[4]}2$ at both directions with $\frac13\big(\Tx{2}{\Ilat{\frac13+\mathbb Z}}{0}-\Tx{2}{\Ilat{\frac23+\mathbb Z}}{0}\big)$, summed as the series of Definition~\ref{def:lattice}, to $90$ digits; those of $\binom{[4]}1$ at its centre and of $W_3$ at $(\one;0,2)$ with $7\,\Tx3{\IZ}0$ and $6\,\Tx3{\IZ}0$, to $89$ digits; and that of the ray with origin $u_0$ with
\[
\int_0^\infty2\log\frac{(1+u^2)^2}{u^2}\,\frac{du}{1+u^2+u^4},
\]
computed by quadrature at $90$ digits, to $84$ digits. Receipts: \texttt{table\_check.py} (\texttt{table\_check.out}), \texttt{table\_masters.py} (\texttt{table\_masters.out}).

So the procedure reads every ray of the table from its list and its direction, with boxes of up to $87507$ relations; the relation of the rows is the law of $[1]$ on all of them but the two rays of $\mathrm H_4$ and $\mathrm V_8$, where the census bounds the law's order below by $2$ and $18$; every value passes the low-end relation, and the five that can be compared with numbers computed otherwise agree with them to $84$ digits or more. In particular $\binom{[4]}2$ has the value $\ell$ of $\binom{[3]}1$ at its centre both at $(\one;0,2)$ and at $(\one;1,1)$, to $90$ digits.
\end{solution}

\chapter{Exercises on the content}\label{ch:excontent}

The exercises of this chapter work the statements of Chapter~\ref{ch:content} on the two smallest centre values met in the exercises on rays. The value $\ell$ of the list of three marks is placed in the content of level $3$ as a combination of two trails of depth one and of two chains; the value of the half-lattice met at four marks is placed in the content of level $2$, and the coset relation returns it to level $1$. The cutting rule is then computed on the symbols of two trails, subset by subset, and the three parts of Theorem~\ref{thm:ct-coaction} are checked on them. Throughout, $\omega=e^{2\pi i/3}$, so that $\omega^3=1$, $\omega^2+\omega+1=0$, $\bar\omega=\omega^2=\omega^{-1}$ and $\mu_3=\{1,\omega,\omega^2\}$. Every number displayed is exact; decimal values are labelled renderings, and the receipts that recompute them are the scripts named in each solution.

\section{Two centre values in the content}\label{sec:exct-values}

\begin{exercise}[The values of three and four marks as trails and chains]\label{ex:ct-values}
Let $\ell=P_{\binom{[3]}1}(\one;1,1)=3\big(\Tx{2}{\Ilat{1+3\mathbb Z}}{0}-\Tx{2}{\Ilat{2+3\mathbb Z}}{0}\big)$ (Exercise~\ref{ex:smallc-torus}) and $v=\Tx{3}{\Ilat{\frac12+\mathbb Z}}{0}$ (Exercise~\ref{ex:cfour-halflattice}).
\begin{enumerate}[label=(\roman*)]
\item Write $\ell$ as a combination, with coefficients in $\mathbb Q(\mu_3)$, of the trails of level $3$ of depth one at the index $2$ and the frequencies $\frac13$ and $\frac23$ (Definition~\ref{def:ct-content}), and show that $\ell$ is a value of $\mathcal C^{(3)}_2$.
\item Write these two trails as chains of level $3$ from $0$ to $1$ (Proposition~\ref{prop:ct-chains}(i)), check the first by integrating it, and write $\ell$ as a combination of chains.
\item Show that $v$ is a value of $\mathcal C^{(2)}_3$ by writing it with the trails of level $2$ at the index $3$ and the frequencies $0$ and $\frac12$. Then use the coset relation to show that $v$ is a value of $\mathcal C^{(1)}_3$, find the trail at the frequency $\frac12$ in terms of the one at the frequency $0$, and write $v$ as a chain.
\end{enumerate}
\end{exercise}

\begin{solution}
Receipt: \texttt{ct\_values.py}.

(i) \emph{The cosets as frequencies.} For an integer $m$, $\omega^m-\omega^{-m}$ is $0$, $\omega-\omega^2$ and $\omega^2-\omega$ for $m$ in $3\mathbb Z$, $1+3\mathbb Z$ and $2+3\mathbb Z$, as $\omega^{-1}=\omega^2$. So the difference of the two indicators is
\[
\one_{1+3\mathbb Z}(m)-\one_{2+3\mathbb Z}(m)=\frac{\omega^m-\omega^{-m}}{\omega-\omega^2}\qquad(m\in\mathbb Z).
\]
At the index $2$ every sum below converges absolutely, so the value of a lattice input at the cut $0$ is the sum of its density times $m^{-2}$ over the points $m>0$ (Definition~\ref{def:lattice}), and the sums may be split and recombined. With $\omega^{-m}=e^{2\pi i\cdot 2m/3}$,
\begin{multline*}
\Tx{2}{\Ilat{1+3\mathbb Z}}{0}-\Tx{2}{\Ilat{2+3\mathbb Z}}{0}=\sum_{m\ge1}\frac{\omega^m-\omega^{-m}}{\omega-\omega^2}\,m^{-2}\\
=\frac{1}{\omega-\omega^2}\Big(\Tx{2}{e^{2\pi ix/3}\IZ}{0}-\Tx{2}{e^{4\pi ix/3}\IZ}{0}\Big),
\end{multline*}
the lattice inputs $e^{2\pi iax}\IZ$ at the frequencies $a=\frac13$ and $a=\frac23$ on the counting lattice (Definition~\ref{def:lattice}).

\emph{The constant.} $(\omega-\omega^2)^2=\omega^2-2\omega^3+\omega^4=\omega^2+\omega-2=-3$. So $\frac{1}{\omega-\omega^2}=-\frac{\omega-\omega^2}{3}$ and $\frac{3}{\omega-\omega^2}=\omega^2-\omega$, and
\[
\ell=(\omega^2-\omega)\Big(\Tx{2}{e^{2\pi ix/3}\IZ}{0}-\Tx{2}{e^{4\pi ix/3}\IZ}{0}\Big).
\]

\emph{The trails.} For $d=1$, $\nu_1=2$ and $a_1\in\{\frac13,\frac23\}\subseteq\frac13\mathbb Z$, the trail of level $3$ of Definition~\ref{def:ct-content} is $\big(\sT_{2;I_1}[y]\,e^{2\pi ia_1y}\,\Ilat{\mathbb Z}\big)(\{0<y\})=\sum_{m\ge1}e^{2\pi ia_1m}m^{-2}=\Tx{2}{e^{2\pi ia_1x}\IZ}{0}$: it has the depth $1$, the positive integer index $2$ and a frequency in $\frac13\mathbb Z$, and the condition at index one does not arise. Write $T'$ and $T''$ for the trails at $\frac13$ and $\frac23$ in this solution. Then
\[
\ell=\omega^2T'-\omega T'-\omega^2T''+\omega T'' ,
\]
the value of the combination $[\omega^2T']-[\omega T']-[\omega^2T'']+[\omega T'']$, with the rational coefficients $\pm1$, of the classes of the words $\zeta T$ with $\zeta\in\mu_3$ and $T$ a trail of level $3$ of weight $2$. So $\ell$ is a value of $\mathcal C^{(3)}_2$. Since $\bar\omega=\omega^2$, $T''=\overline{T'}$, and $\ell=(\omega^2-\omega)(T'-\overline{T'})=(-i\sqrt3)(2i\im T')=2\sqrt3\im T'$, with $\im T'=\sum_{m\ge1}\sin(\frac{2\pi m}3)m^{-2}=\frac{\sqrt3}2\big(\Tx{2}{\Ilat{1+3\mathbb Z}}{0}-\Tx{2}{\Ilat{2+3\mathbb Z}}{0}\big)$, the $S_2$ of Exercise~\ref{ex:smallc-torus}.

Rendering: both sides of the display for $\ell$ are $2.3439072386894588\ldots$.

(ii) \emph{The chains.} Proposition~\ref{prop:ct-chains}(i) with $d=1$ and $\nu_1=2$ gives $c_1=e^{2\pi ia_1}$, the letters $c_1^{-1}$ and one $0$, and the sign $(-1)^1$. At $a_1=\frac13$, $c_1=\omega$ and $c_1^{-1}=\omega^2$; at $a_1=\frac23$, $c_1=\omega^2$ and $c_1^{-1}=\omega$:
\[
T'=-\mathrm I(0;\omega^2,0;1),\qquad T''=-\mathrm I(0;\omega,0;1).
\]

\emph{The check.} For the ends $0$ and $1$ the path is $\gamma(t)=t$, and for a letter $b$ with $|b|=1$, $b\neq1$ (Definition~\ref{def:ct-chain}),
\[
\mathrm I(0;b,0;1)=\int_{0<t_1<t_2<1}\frac{dt_1}{t_1-b}\,\frac{dt_2}{t_2}.
\]
For $0\le t<1$, $1-t/b$ lies in the open disc of radius $1$ about $1$, where the principal logarithm is holomorphic; $\log(1-t/b)$ vanishes at $t=0$ and has the derivative $\frac{-1/b}{1-t/b}=\frac{1}{t-b}$, so the inner integral is $\log(1-t_2/b)=-\sum_{n\ge1}b^{-n}t_2^n/n$. The series of absolute values, $\sum_nt_2^{n-1}/n$, integrates over $(0,1)$ to $\sum_nn^{-2}<\infty$, so the series integrates term by term, $\int_0^1t_2^{n-1}\,dt_2=\frac1n$:
\[
\mathrm I(0;b,0;1)=-\sum_{n\ge1}b^{-n}n^{-2}.
\]
At $b=\omega^2$, $b^{-n}=\omega^{-2n}=\omega^n$, and $\mathrm I(0;\omega^2,0;1)=-\sum_n\omega^nn^{-2}=-T'$, as Proposition~\ref{prop:ct-chains}(i) states.

Rendering: the double integral is $0.5483113556160754\ldots-0.6766277376064357\ldots i$, and so is $-T'$.

\emph{The value.} Substituting,
\[
\ell=(\omega^2-\omega)\big(-\mathrm I(0;\omega^2,0;1)+\mathrm I(0;\omega,0;1)\big)=(\omega-\omega^2)\big(\mathrm I(0;\omega^2,0;1)-\mathrm I(0;\omega,0;1)\big),
\]
a combination of chains of level $3$ from $0$ to $1$ with first letter not $0$ and last letter not $1$, as Proposition~\ref{prop:ct-chains}(iii) describes the values of $\mathcal C^{(3)}_2$.

(iii) \emph{Level $2$.} The positive points of $\frac12+\mathbb Z$ are the $\frac{2m-1}2$, $m\ge1$, so $v=\sum_{m\ge1}8(2m-1)^{-3}$ (Definition~\ref{def:lattice}), absolutely convergent. The indicator of the odd integers is $\frac{1-(-1)^n}2=\frac{1-e^{\pi in}}2$, so
\[
v=8\sum_{n\ge1}\frac{1-e^{\pi in}}2\,n^{-3}=4\,\Tx{3}{\IZ}{0}-4\,\Tx{3}{e^{\pi ix}\IZ}{0}.
\]
$\Tx{3}{\IZ}{0}$ and $\Tx{3}{e^{\pi ix}\IZ}{0}$ are the trails of depth one at the index $3$ and the frequencies $0$ and $\frac12$, both in $\frac12\mathbb Z$: trails of level $2$. So $v$ is a value of $\mathcal C^{(2)}_3$, as Proposition~\ref{prop:ct-level}(i) gives directly for the coset $b=\frac12\in\frac12\mathbb Z$ at the frequency $0$, with $M=2$ and $M'=1$.

\emph{Level $1$.} The coset relation of Exercise~\ref{ex:cfour-halflattice}, Theorem~\ref{thm:latticefree}(i) at $q=2$ and $b=0$, is $\Tx{3}{\IZ}{0}+v=8\,\Tx{3}{\IZ}{0}$. So
\[
v=7\,\Tx{3}{\IZ}{0},
\]
a rational multiple of a trail of level $1$, and $v$ is a value of $\mathcal C^{(1)}_3$. Comparing the two expressions of $v$, $4\,\Tx{3}{e^{\pi ix}\IZ}{0}=4\,\Tx{3}{\IZ}{0}-7\,\Tx{3}{\IZ}{0}$:
\[
\Tx{3}{e^{\pi ix}\IZ}{0}=-\tfrac34\,\Tx{3}{\IZ}{0}.
\]

\emph{The chain.} Proposition~\ref{prop:ct-chains}(i) with $d=1$, $\nu_1=3$ and $a_1=0$ gives $c_1=1$, the letters $1,0,0$, and $\Tx{3}{\IZ}{0}=-\mathrm I(0;1,0,0;1)$; so $v=-7\,\mathrm I(0;1,0,0;1)$.

Rendering: $v=8.4143983221171599\ldots$, the value of both expressions above; the ratio of the trail at the frequency $\frac12$ to the trail at the frequency $0$ renders as $-0.75$, and the triple integral $\mathrm I(0;1,0,0;1)$ as $-1.2020569031595942\ldots$.

\emph{Answer.}
\begin{align*}
\ell&=(\omega^2-\omega)\Big(\Tx{2}{e^{2\pi ix/3}\IZ}{0}-\Tx{2}{e^{4\pi ix/3}\IZ}{0}\Big)\\
&=(\omega-\omega^2)\big(\mathrm I(0;\omega^2,0;1)-\mathrm I(0;\omega,0;1)\big),
\end{align*}
a value of $\mathcal C^{(3)}_2$. $v=4\,\Tx{3}{\IZ}{0}-4\,\Tx{3}{e^{\pi ix}\IZ}{0}$ is a value of $\mathcal C^{(2)}_3$, and by the coset relation $v=7\,\Tx{3}{\IZ}{0}=-7\,\mathrm I(0;1,0,0;1)$, a value of $\mathcal C^{(1)}_3$, with $\Tx{3}{e^{\pi ix}\IZ}{0}=-\frac34\,\Tx{3}{\IZ}{0}$.
\end{solution}

\section{The cutting rule on two trails}\label{sec:exct-cutting}

\begin{exercise}[The cutting rule on $\Tx{3}{\IZ}{0}$ and on the trail of index $2$ at the frequency $\frac13$]\label{ex:ct-cutting}
Let $\sigma_3$ be the symbol of the trail $\Tx{3}{\IZ}{0}$ of level $1$, and $\sigma'$ that of the trail $\Tx{2}{e^{2\pi ix/3}\IZ}{0}$ of level $3$ (Definition~\ref{def:ct-symbols}).
\begin{enumerate}[label=(\roman*)]
\item Compute $\Delta\sigma_3$ and $\Delta\sigma'$ (Definition~\ref{def:ct-cutting}), subset by subset.
\item Check the three parts of Theorem~\ref{thm:ct-coaction} on $\sigma_3$ and $\sigma'$: coassociativity and the counit, the weights, and the compatibility with the reversal and the rotations.
\item Read the left factors of positive weight as symbols of trails (Proposition~\ref{prop:ct-chains}(ii)), and compute the chains of the two factors of weight one in $\Delta\sigma'$.
\end{enumerate}
\end{exercise}

\begin{solution}
Receipts: \texttt{ct\_cutting.py} for (i) and (ii), \texttt{ct\_values.py} for the renderings in (iii).

\emph{The symbols.} Both trails have depth $d=1$. For $\Tx{3}{\IZ}{0}$, $\nu_1=3$ and $c_1=1$; for $\Tx{2}{e^{2\pi ix/3}\IZ}{0}$, $\nu_1=2$ and $c_1=\omega$, $c_1^{-1}=\omega^2$. So, with $x=\mathrm I^{\mathfrak s}(0;1,0,0;1)$ and $y=\mathrm I^{\mathfrak s}(0;\omega^2,0;1)$,
\[
\sigma_3=-x,\qquad\sigma'=-y .
\]
$\Delta$ is $\mathbb Q(\mu_M)$-linear, so $\Delta\sigma_3=-\Delta x$ and $\Delta\sigma'=-\Delta y$. In $\mathfrak S_M$ a run with no letters is $1$, and a run of at least one letter between equal ends is $0$ (Definition~\ref{def:ct-symbols}); these two rules decide every right factor below.

(i) \emph{The cutting rule on $x$.} The entries are $a_0=0$, $a_1=1$, $a_2=0$, $a_3=0$, $a_4=1$, and $w=3$. For each $S\subseteq\{1,2,3\}$, the left factor keeps the letters of $S$ between the ends $0$ and $1$, and the right factors are the runs between consecutive kept positions $s_0=0<s_1<\cdots<s_{|S|+1}=4$:
\begin{center}
\begin{tabular}{lll}
\toprule
$S$ & left factor & runs, and the term\\
\midrule
$\varnothing$ & $1$ & $\mathrm I^{\mathfrak s}(0;1,0,0;1)=x$: the term $1\otimes x$\\
$\{1\}$ & $\mathrm I^{\mathfrak s}(0;1;1)$ & $\mathrm I^{\mathfrak s}(0;\varnothing;1)=1$, $\mathrm I^{\mathfrak s}(1;0,0;1)=0$: no term\\
$\{2\}$ & $\mathrm I^{\mathfrak s}(0;0;1)$ & $\mathrm I^{\mathfrak s}(0;1;0)=0$: no term\\
$\{3\}$ & $\mathrm I^{\mathfrak s}(0;0;1)$ & $\mathrm I^{\mathfrak s}(0;1,0;0)=0$: no term\\
$\{1,2\}$ & $\mathrm I^{\mathfrak s}(0;1,0;1)$ & $1$, $1$, $\mathrm I^{\mathfrak s}(0;0;1)$\\
$\{1,3\}$ & $\mathrm I^{\mathfrak s}(0;1,0;1)$ & $1$, $\mathrm I^{\mathfrak s}(1;0;0)$, $1$\\
$\{2,3\}$ & $\mathrm I^{\mathfrak s}(0;0,0;1)$ & $\mathrm I^{\mathfrak s}(0;1;0)=0$: no term\\
$\{1,2,3\}$ & $x$ & $1,1,1,1$: the term $x\otimes1$\\
\bottomrule
\end{tabular}
\end{center}
With $z=\mathrm I^{\mathfrak s}(0;1,0;1)$ and $R=\mathrm I^{\mathfrak s}(0;0;1)+\mathrm I^{\mathfrak s}(1;0;0)$,
\[
\Delta x=x\otimes1+1\otimes x+z\otimes R .
\]

\emph{The cutting rule on $y$.} The entries are $a_0=0$, $a_1=\omega^2$, $a_2=0$, $a_3=1$, and $w=2$:
\begin{center}
\begin{tabular}{lll}
\toprule
$S$ & left factor & runs, and the term\\
\midrule
$\varnothing$ & $1$ & $y$: the term $1\otimes y$\\
$\{1\}$ & $\mathrm I^{\mathfrak s}(0;\omega^2;1)$ & $\mathrm I^{\mathfrak s}(0;\varnothing;\omega^2)=1$, $\mathrm I^{\mathfrak s}(\omega^2;0;1)$\\
$\{2\}$ & $\mathrm I^{\mathfrak s}(0;0;1)$ & $\mathrm I^{\mathfrak s}(0;\omega^2;0)=0$: no term\\
$\{1,2\}$ & $y$ & $1,1,1$: the term $y\otimes1$\\
\bottomrule
\end{tabular}
\end{center}
With $A=\mathrm I^{\mathfrak s}(0;\omega^2;1)$ and $B=\mathrm I^{\mathfrak s}(\omega^2;0;1)$,
\[
\Delta y=y\otimes1+1\otimes y+A\otimes B .
\]

(ii) \emph{The pieces are primitive.} A formal chain $u$ of weight one has two subsets: $S=\varnothing$ gives $1\otimes u$ and $S=\{1\}$ gives $u\otimes1$, its two runs being empty. So $\Delta u=u\otimes1+1\otimes u$ for $u=\mathrm I^{\mathfrak s}(0;0;1)$, $\mathrm I^{\mathfrak s}(1;0;0)$, $A$ and $B$, and $\Delta R=R\otimes1+1\otimes R$ by linearity. For $z$, with the entries $0,1,0,1$: $S=\{1\}$ leaves the run $\mathrm I^{\mathfrak s}(1;0;1)=0$ and $S=\{2\}$ the run $\mathrm I^{\mathfrak s}(0;1;0)=0$, so $\Delta z=z\otimes1+1\otimes z$ as well.

\emph{Coassociativity.} $\Delta$ is an algebra homomorphism with $\Delta1=1\otimes1$, so on $y$
\[
(\Delta\otimes\mathrm{id})\Delta y=1\otimes1\otimes y+(y\otimes1+1\otimes y+A\otimes B)\otimes1+(A\otimes1+1\otimes A)\otimes B ,
\]
\[
(\mathrm{id}\otimes\Delta)\Delta y=y\otimes1\otimes1+1\otimes(y\otimes1+1\otimes y+A\otimes B)+A\otimes(B\otimes1+1\otimes B),
\]
and both are the sum of the six tensors $1\otimes1\otimes y$, $1\otimes y\otimes1$, $y\otimes1\otimes1$, $A\otimes B\otimes1$, $A\otimes1\otimes B$ and $1\otimes A\otimes B$. On $x$ the same expansion, with $z$ and $R$ primitive, gives on both sides the six tensors $1\otimes1\otimes x$, $1\otimes x\otimes1$, $x\otimes1\otimes1$, $z\otimes R\otimes1$, $z\otimes1\otimes R$ and $1\otimes z\otimes R$.

\emph{The counit.} $\varepsilon(1)=1$ and $\varepsilon$ vanishes on the formal chains, so $(\varepsilon\otimes\mathrm{id})\Delta x=x+\varepsilon(z)R=x$ and $(\mathrm{id}\otimes\varepsilon)\Delta x=x+z\,\varepsilon(R)=x$, and in the same way on $y$. This is Theorem~\ref{thm:ct-coaction}(i) on the two symbols.

\emph{The weights.} The terms of $\Delta x$ have the weights $(3,0)$, $(0,3)$ and $(2,1)$, those of $\Delta y$ the weights $(2,0)$, $(0,2)$ and $(1,1)$; the sums are $3$ and $2$, and $\Delta x-x\otimes1-1\otimes x=z\otimes R$ and $\Delta y-y\otimes1-1\otimes y=A\otimes B$ have both factors of positive weight: Theorem~\ref{thm:ct-coaction}(ii).

\emph{The reversal.} $\varrho y=(-1)^2\,\mathrm I^{\mathfrak s}(1;0,\omega^2;0)$. Its entries $1,0,\omega^2,0$ give: $S=\varnothing$, $1\otimes\varrho y$; $S=\{1\}$, the left factor $\mathrm I^{\mathfrak s}(1;0;0)$ and the run $\mathrm I^{\mathfrak s}(0;\omega^2;0)=0$, no term; $S=\{2\}$, the left factor $\mathrm I^{\mathfrak s}(1;\omega^2;0)$ and the runs $\mathrm I^{\mathfrak s}(1;0;\omega^2)$ and $1$; $S=\{1,2\}$, $\varrho y\otimes1$. On the other side $\varrho A=-\mathrm I^{\mathfrak s}(1;\omega^2;0)$ and $\varrho B=-\mathrm I^{\mathfrak s}(1;0;\omega^2)$, so
\[
\Delta\varrho y=\varrho y\otimes1+1\otimes\varrho y+\mathrm I^{\mathfrak s}(1;\omega^2;0)\otimes\mathrm I^{\mathfrak s}(1;0;\omega^2)=(\varrho\otimes\varrho)\Delta y .
\]
$\varrho x=(-1)^3\,\mathrm I^{\mathfrak s}(1;0,0,1;0)$. Of the eight subsets of the entries $1,0,0,1,0$, those with a term are $\varnothing$, $\{1,2,3\}$, $\{1,3\}$, with the left factor $\mathrm I^{\mathfrak s}(1;0,1;0)$ and the run $\mathrm I^{\mathfrak s}(0;0;1)$, and $\{2,3\}$, with the same left factor and the run $\mathrm I^{\mathfrak s}(1;0;0)$; the other four leave a run between equal ends. So $\Delta\varrho x=\varrho x\otimes1+1\otimes\varrho x-\mathrm I^{\mathfrak s}(1;0,1;0)\otimes R$. On the other side $\varrho z=(-1)^2\,\mathrm I^{\mathfrak s}(1;0,1;0)$ and $\varrho R=-\mathrm I^{\mathfrak s}(1;0;0)-\mathrm I^{\mathfrak s}(0;0;1)=-R$, so $(\varrho\otimes\varrho)(z\otimes R)=-\mathrm I^{\mathfrak s}(1;0,1;0)\otimes R$, and $\Delta\varrho x=(\varrho\otimes\varrho)\Delta x$.

\emph{The rotations.} $\rho_\lambda$, $\lambda\in\mu_3$, multiplies every entry by $\lambda$ and fixes the entry $0$. It changes neither which ends of a run are equal nor which runs are empty, so the same subsets give the terms of $\Delta\rho_\lambda y$, with every entry multiplied by $\lambda$: $\Delta\rho_\lambda y=\rho_\lambda y\otimes1+1\otimes\rho_\lambda y+\mathrm I^{\mathfrak s}(0;\lambda\omega^2;\lambda)\otimes\mathrm I^{\mathfrak s}(\lambda\omega^2;0;\lambda)=(\rho_\lambda\otimes\rho_\lambda)\Delta y$. At level $1$ the only rotation of $x$ is the identity. This is Theorem~\ref{thm:ct-coaction}(iii) on the two symbols.

(iii) \emph{The left factor of $\Delta x$.} The chain $\mathrm I(0;1,0;1)$ has the first letter $1\neq0$ and the last letter $0\neq1$. In Proposition~\ref{prop:ct-chains}(ii) its one nonzero letter is at $j_1=1$, with $j_2=w+1=3$, so $d=1$, $c_1=1$, $\nu_1=2$, $c_2=1$, and the frequency $a'_1=0$: $\mathrm I(0;1,0;1)=-\Tx{2}{\IZ}{0}$, and $z=-\sigma_2$ with $\sigma_2$ the symbol of $\Tx{2}{\IZ}{0}$. The right factor $R=\mathrm I^{\mathfrak s}(0;0;1)-\varrho\,\mathrm I^{\mathfrak s}(0;0;1)$ is a formal chain of weight one minus its reversal; the first letter of $\mathrm I^{\mathfrak s}(0;0;1)$ is its first end, so it has no realization (Definition~\ref{def:ct-symbols}), and $\int_0^1dt/t$ diverges. Rendering: $\mathrm I(0;1,0;1)=-1.6449340668482264\ldots$.

\emph{The factors of $\Delta y$.} $A$ has the letter $\omega^2\neq0,1$. In Proposition~\ref{prop:ct-chains}(ii), $d=1$, $c_1=\omega$, $\nu_1=1$ and $c_2=1$, with the frequency $a'_1=\frac13$, $e^{2\pi ia'_1}=\omega$: $\mathrm I(0;\omega^2;1)$ is minus the trail of level $3$ of depth one at the index $1$ and the frequency $\frac13$, $\lim_N\sum_{m\le N}\omega^mm^{-1}$, which is a trail because $e^{2\pi i/3}\neq1$ (Definition~\ref{def:ct-content}); $A$ is minus its symbol. On the path $\gamma(t)=t$, $t-\omega^2=t+\frac12+\frac{\sqrt3}2i$ has the imaginary part $\frac{\sqrt3}2>0$, so the principal logarithm is continuous along it and
\[
\mathrm I(0;\omega^2;1)=\int_0^1\frac{dt}{t-\omega^2}=\log\Big(\frac32+\frac{\sqrt3}2i\Big)-\log\Big(\frac12+\frac{\sqrt3}2i\Big)=\Big(\frac12\log3+\frac{\pi i}6\Big)-\frac{\pi i}3=\frac12\log3-\frac{\pi i}6 .
\]
$B$ has the ends $\omega^2$ and $1$, so it is a chain between two singular points with no end at $0$, read as it stands. Its letter $0$ is not on the open segment from $\omega^2$ to $1$ and differs from both ends, so the chain is defined and converges (Definition~\ref{def:ct-chain}, Proposition~\ref{prop:ct-between}(i)). On $\gamma(t)=\omega^2+t(1-\omega^2)=-\frac12+\frac32t-\frac{\sqrt3}2(1-t)i$ the imaginary part is negative for $t<1$, so the argument of $\gamma(t)$ runs continuously from $-\frac{2\pi}3$ to $0$ and the modulus from $1$ to $1$:
\[
\mathrm I(\omega^2;0;1)=\int_0^1\frac{\gamma'(t)\,dt}{\gamma(t)}=\log1-\log\omega^2=0-\Big(-\frac{2\pi i}3\Big)=\frac{2\pi i}3 .
\]
Rendering: the integrals render as $0.5493061443340548\ldots-0.5235987755982988\ldots i$ and $2.0943951023931954\ldots i$.

\emph{Answer.} With $\sigma_2$ the symbol of $\Tx{2}{\IZ}{0}$ and $\sigma'_1=-\mathrm I^{\mathfrak s}(0;\omega^2;1)$ the symbol of the trail of index $1$ at the frequency $\frac13$,
\[
\Delta\sigma_3=\sigma_3\otimes1+1\otimes\sigma_3+\sigma_2\otimes\big(\mathrm I^{\mathfrak s}(0;0;1)+\mathrm I^{\mathfrak s}(1;0;0)\big),\qquad\Delta\sigma'=\sigma'\otimes1+1\otimes\sigma'+\sigma'_1\otimes\mathrm I^{\mathfrak s}(\omega^2;0;1).
\]
Both are coassociative and counital, graded by the weight, and compatible with the reversal and the rotations, as Theorem~\ref{thm:ct-coaction} states. The left factors $\sigma_2$ and $\sigma'_1$ are symbols of trails; the right factor $\mathrm I^{\mathfrak s}(\omega^2;0;1)$ is a chain between the singular points $\omega^2$ and $1$ with the value $\frac{2\pi i}3$, and $\mathrm I(0;\omega^2;1)=\frac12\log3-\frac{\pi i}6$.
\end{solution}

\chapter*{Aliases}
\addcontentsline{toc}{chapter}{Aliases}
\label{ch:aliases}

The text uses the program's own names. This list gives, for a reader who arrives from the literature, the classical names under which some of its objects and results appear there. An alias is a name and nothing more: the statement of the handbook, with its hypotheses, is the one that holds, and where the classical result is narrower or wider the handbook's statement decides.

\section*{Lists, tables and their letters}

\begin{longtable}{p{0.45\textwidth}p{0.49\textwidth}}
\toprule
In this handbook & Classical name\\
\midrule
\endhead
a table, its marks and bases (Definition~\ref{def:table}) & a matroid with weights on its bases, its ground set and bases\\
a list (the weights all $1$) & a matroid\\
the list $\binom{\cE}r$, $|\cE|=n$ & the uniform matroid $U_{r,n}$\\
the dual table & the dual matroid\\
a molecule; a connected table & a connected component; a connected matroid\\
the width $c(S)$ of a set (Definition~\ref{def:skeleton}) & the connectivity $r(S)+r(\cE\setminus S)-r(\cE)$ of the set\\
the two-sum $\mathcal A\oplus_g\mathcal A'$ of two cells (Lemma~\ref{lem:gluemark}) & the $2$-sum of matroids\\
for the list of the spanning trees of a graph: the basis polynomial $K$ and the cobasis polynomial $\Psi$ & the spanning-tree polynomial, and the graph polynomial, also called the first Symanzik or the Kirchhoff polynomial\\
the Cremona (Theorem~\ref{thm:cremona}) & the Cremona transformation $\alpha\mapsto\alpha^{-1}$\\
the period at the full cobasis lean $(\one;0,2)$ of the list of the spanning trees of a primitive logarithmically divergent graph & its Feynman period $\int\Omega/\Psi^2$; the volume of the direction $V_1$ of that ray is its Hepp bound\\
the list $\binom{[n]}1$ & the cycle matroid of the banana graph with $n$ edges\\
the wheel $W_p$ (Proposition~\ref{prop:wheelholonomy}); $W_3$ & the cycle matroid of the wheel with $p$ spokes; $W_3$ is the cycle matroid of $K_4$\\
$\mathrm H_4$ & the cycle matroid of a triangle with one edge doubled\\
$\mathrm N_6$ & $U_{3,6}$ with one basis removed: six points of a plane, exactly three of them collinear\\
$\mathrm V_8$ & the V\'amos matroid\\
the whirl list on six marks (rank three, three dependent sets of three marks meeting pairwise) & the rank-three whirl\\
$\mathrm F_7$, $\mathrm A_9$, $\mathrm Z_9$ (Chapter~\ref{ch:fusion}) & the Fano plane, the affine plane $AG(2,3)$, the Pappus configuration\\
\bottomrule
\end{longtable}

\section*{Values}

\begin{longtable}{p{0.45\textwidth}p{0.49\textwidth}}
\toprule
In this handbook & Classical name\\
\midrule
\endhead
$\Tx m{\IZ}0$ & $\zeta(m)$\\
$\Gamma(\frac12)^2$ & $\pi$\\
$\Tx2{(\Ilat{1+3\mathbb Z}-\Ilat{2+3\mathbb Z})}0$; so $P_{\binom{[3]}1}(\one;1,1)=3\,\Tx2{(\Ilat{1+3\mathbb Z}-\Ilat{2+3\mathbb Z})}0$ & $L(\chi_{-3},2)$; the value is $3L(\chi_{-3},2)$\\
$P_{\binom{[4]}1}(\one;1,1)=7\,\Tx3{\IZ}0$; the period of $W_3$ at $(\one;0,2)$, $6\,\Tx3{\IZ}0$ & $7\zeta(3)$; $6\zeta(3)$\\
the signed sums of Trailmoments of lattice inputs at the cosets prime to $|D|$, weighted by the Kronecker signs & $L(\chi_D,2)$\\
the arc integrals of the three-mark centre (Part~\ref{part:exercises}) & the log-sine integrals $\mathrm{Ls}_j(\frac{2\pi}3)$; the second is the Clausen value $\mathrm{Cl}_2(\frac{2\pi}3)$\\
the stacked Gammas at fifteenths, at twentieths & Chowla--Selberg periods of $\mathbb Q(\sqrt{-15})$, of $\mathbb Q(\sqrt{-5})$\\
the Gramma of one mark, $\sG_{\nu;(1,-1)}(1,y)=2y^{\nu/2}K_\nu(2\sqrt y)$ & the modified Bessel function $K_\nu$\\
the centre values of $\binom{[n]}1$, $P_{\binom{[n]}1}(s\one;s,s)=2^{n+1-2s}\Gamma(s)^{-2}\int_0^\infty x^{2s-1}K_0(x)^n\,dx$ & Bessel moments; at $n=3$ the equal-mass two-loop vacuum (sunrise) integral with propagator powers $s$ in dimension $2s$\\
\bottomrule
\end{longtable}

\section*{Rays}

\begin{longtable}{p{0.45\textwidth}p{0.49\textwidth}}
\toprule
In this handbook & Classical name\\
\midrule
\endhead
the relations along a ray, reduced in a box (Definition~\ref{def:raybox}) & integration-by-parts reduction on the parametric form, in the manner of Laporta\\
the strength & the parameter that plays the dimension\\
masters (Definition~\ref{def:raymodule}) & master integrals\\
the one-step matrix & the dimension-shift recurrence written as a first-order difference system\\
the law of $[1]$ & the scalar recurrence obtained from a cyclic vector\\
a certificate (Theorem~\ref{thm:raybox}(ii)) & a creative-telescoping certificate\\
the wall degree & the tropical degree of an insertion\\
the search through boxes (Theorem~\ref{thm:rayboxstop}) & Laporta's algorithm, its seeds grown until the targets reduce\\
the dimension of a ray (Definition~\ref{def:raygauge}) & the number of master integrals\\
the spectrum at infinite strength (Proposition~\ref{prop:rays-spectrum}) & the characteristic roots of a difference system (Birkhoff); a change of masters is a gauge transformation\\
the evaluation modulo primes in the reading procedure (Definition~\ref{def:readprocedure}) & finite-field reconstruction\\
the closings (Definition~\ref{def:closings}) & the classical periods of the Laurent polynomial $f^{-1}$; the law read at the low end is the Picard--Fuchs operator of its pencil, read on coefficients\\
the closings of the centre ray of $\binom{[3]}1$ & the sums of squares of trinomial coefficients\\
positive selection (Theorem~\ref{thm:positiveselection}) & Ap\'ery limits; the minimal solution of Pincherle and Perron\\
positivity as the moment problem of the stride's law (Proposition~\ref{prop:rd-moment}); its rate (Theorem~\ref{thm:rd-rate}); what it cannot see (Proposition~\ref{prop:rd-blind}) & the Hausdorff moment problem, its four families of Markov and Luk\'acs; the Green's function of the complement of a segment; the indeterminacy of a moment problem along a Stieltjes transform\\
the stride's law (Definition~\ref{def:strideslaw}) & the pushforward of Lebesgue measure under the stride, a density of states; the values are its moments\\
the gauge end (Theorem~\ref{thm:gaugeend}) & Laplace's method with the Morse lemma at the end point, and the Frobenius method at a regular singular point; the germ is the Borel sum of the expansion at large strength\\
the torus value (Definition~\ref{def:torusvalue}) & the zeta Mahler measure of the letter; at three marks the moments of the planar walk of three unit steps\\
the merit (Definition~\ref{def:ra-merit}, Theorem~\ref{thm:ra-merit}) & Ap\'ery's criterion\\
the law of the values of $\binom{\cE}1$ (Theorem~\ref{thm:ri-law}) & the recurrences of the Bessel moments $\int_0^\infty x^{2s-1}K_0(x)^n\,dx$\\
\bottomrule
\end{longtable}

\section*{The motive and the content}

\begin{longtable}{p{0.45\textwidth}p{0.49\textwidth}}
\toprule
In this handbook & Classical name\\
\midrule
\endhead
the motive and the motivic period of a family (Definitions~\ref{def:mot-model} and~\ref{def:mot-mperiod}) & motivic periods in Brown's affine model\\
small steps along the tree (Theorem~\ref{thm:mot-smallsteps}) & Brown's small graphs principle in its affine form, here for families of lists\\
dual pairs and full leans (Corollary~\ref{cor:mot-dualfull}) & at $(0,2)$, the product identity of Feynman periods at a two-vertex cut (Schnetz), as an identity of motivic periods\\
chains (Definition~\ref{def:ct-chain}); the content (Definition~\ref{def:ct-content}) & iterated integrals on $\{0\}$ and the roots of unity of order $M$; the multiple polylogarithms at roots of unity\\
the cutting rule (Definition~\ref{def:ct-cutting}, Theorem~\ref{thm:ct-coaction}) & Goncharov's coaction on iterated integrals\\
critical primes (Definition~\ref{def:ra-criticalprime}) & primes of bad reduction\\
\bottomrule
\end{longtable}

\section*{The family as inference}

\begin{longtable}{p{0.45\textwidth}p{0.49\textwidth}}
\toprule
In this handbook & Classical name\\
\midrule
\endhead
the interaction of a split of the prior marks of a regression design (Proposition~\ref{prop:design}) & minus the logarithm of the generalized variance inflation factor\\
\bottomrule
\end{longtable}

\end{document}
